\IfFileExists{stacks-project-book.cls}{%
\documentclass{stacks-project-book}
}{%
\documentclass{amsbook}
}

\usepackage{verbatim}
\newenvironment{reference}{\comment}{\endcomment}
\newenvironment{slogan}{\comment}{\endcomment}
\newenvironment{history}{\comment}{\endcomment}

\usepackage[all]{xy}

\xyoption{2cell}
\UseAllTwocells

\usepackage{multicol}

\usepackage{lmodern}
\usepackage[T1]{fontenc}
\usepackage[french]{babel}
\AtBeginDocument{\shorthandoff{?}}


\usepackage{hyperref}


\theoremstyle{plain}
\newtheorem{theorem}[subsection]{Théorème}
\newtheorem{proposition}[subsection]{Proposition}
\newtheorem{lemma}[subsection]{Lemme}

\theoremstyle{definition}
\newtheorem{definition}[subsection]{Définition}
\newtheorem{example}[subsection]{Exemple}
\newtheorem{exercise}[subsection]{Exercice}
\newtheorem{situation}[subsection]{Situation}

\theoremstyle{remark}
\newtheorem{remark}[subsection]{Remarque}
\newtheorem{remarks}[subsection]{Remarques}

\renewcommand{\proofname}{Démonstration}

\numberwithin{equation}{subsection}

\def\lim{\mathop{\mathrm{lim}}\nolimits}
\def\colim{\mathop{\mathrm{colim}}\nolimits}
\def\Spec{\mathop{\mathrm{Spec}}}
\def\Hom{\mathop{\mathrm{Hom}}\nolimits}
\def\Ext{\mathop{\mathrm{Ext}}\nolimits}
\def\SheafHom{\mathop{\mathcal{H}\!\mathit{om}}\nolimits}
\def\SheafExt{\mathop{\mathcal{E}\!\mathit{xt}}\nolimits}
\def\Sch{\mathit{Sch}}
\def\Mor{\mathop{\mathrm{Mor}}\nolimits}
\def\Ob{\mathop{\mathrm{Ob}}\nolimits}
\def\Sh{\mathop{\mathit{Sh}}\nolimits}
\def\NL{\mathop{N\!L}\nolimits}
\def\CH{\mathop{\mathrm{CH}}\nolimits}
\def\proetale{{pro\text{-}\acute{e}tale}}
\def\etale{{\acute{e}tale}}
\def\QCoh{\mathit{QCoh}}
\def\Ker{\mathop{\mathrm{Ker}}}
\def\Im{\mathop{\mathrm{Im}}}
\def\Coker{\mathop{\mathrm{Coker}}}
\def\Coim{\mathop{\mathrm{Coim}}}

\DeclareMathSymbol{\boxtimes}{\mathbin}{AMSa}{"02}

\def\QCohstack{\mathcal{QC}\!\mathit{oh}}
\def\Cohstack{\mathcal{C}\!\mathit{oh}}
\def\Spacesstack{\mathcal{S}\!\mathit{paces}}
\def\Quotfunctor{\mathrm{Quot}}
\def\Hilbfunctor{\mathrm{Hilb}}
\def\Curvesstack{\mathcal{C}\!\mathit{urves}}
\def\Polarizedstack{\mathcal{P}\!\mathit{olarized}}
\def\Complexesstack{\mathcal{C}\!\mathit{omplexes}}
\def\Pic{\mathop{\mathrm{Pic}}\nolimits}
\def\Picardstack{\mathcal{P}\!\mathit{ic}}
\def\Picardfunctor{\mathrm{Pic}}
\def\Deformationcategory{\mathcal{D}\!\mathit{ef}}
\begin{document}
\title{Le projet Stacks --- édition française, partie 1}
\maketitle
\tableofcontents

% OK, start here.
%

\chapter{Introduction}



\phantomsection
\label{introduction-section-phantom}




\section{Vue d'ensemble}
\label{introduction-section-overview}

\noindent
Outre le livre de Laumon et Moret-Bailly (voir \cite{LM-B}) et les travaux
(en cours) de Fulton et al., nous pensons qu'un manuel libre consacr\'e aux
champs alg\'ebriques et \`a la g\'eom\'etrie alg\'ebrique n\'ecessaire \`a leur
d\'efinition a sa place. Le projet Champs tente d'y parvenir en construisant les
fondements \`a partir de l'alg\`ebre commutative, puis en passant par la
th\'eorie des sch\'emas et des espaces alg\'ebriques pour aboutir \`a des
fondements complets de la th\'eorie des champs alg\'ebriques.

\medskip\noindent
Nous pr\'evoyons que ce contenu sera lu en ligne, car l'une de ses
caract\'eristiques essentielles est la pr\'esence d'hyperliens donnant un
acc\`es rapide aux r\'ef\'erences internes r\'eparties sur de nombreuses pages.
Si vous utilisez dans votre navigateur une visionneuse PDF ou DVI int\'egr\'ee,
les liens entre fichiers devraient fonctionner.

\medskip\noindent
Ce projet est le fruit d'un travail collaboratif et nous vous encourageons
\`a y contribuer. Veuillez signaler par courriel toute coquille ou erreur
que vous relevez au cours de votre lecture, ainsi que toute suggestion,
tout contenu suppl\'ementaire ou tout exemple, \`a l'adresse
\href{mailto:stacks.project@gmail.com}{stacks.project@gmail.com}.
Vous pouvez t\'el\'echarger une archive tar contenant tous les fichiers sources,
l'extraire, ex\'ecuter la commande make et utiliser localement une visionneuse DVI ou PDF.
N'h\'esitez pas \`a modifier les fichiers LaTeX et \`a envoyer vos am\'eliorations
par courriel.


\section{Attribution}
\label{introduction-section-attribution}

\noindent
Vu l'ampleur de cet ouvrage, rendre justice, de façon à la fois exacte et
succincte, à tous les mathématiciens dont
les travaux ont conduit au d\'eveloppement de la th\'eorie que nous tentons
d'exposer ici est une tâche redoutable. Nous esp\'erons susciter \`a terme suffisamment d'int\'er\^et au
sein de la communaut\'e pour trouver des contributeurs dispos\'es \`a r\'ediger,
pour chacun des chapitres, des sections contenant des remarques historiques.

\medskip\noindent
Les personnes ayant contribu\'e \`a cet ouvrage sont mentionn\'ees sur la page
de titre de sa version sous forme de livre ainsi que
\href{https://stacks.math.columbia.edu/tex/CONTRIBUTORS}{sur cette page en ligne}.
Nous souhaitons ici signaler une s\'election de contributions majeures :
\begin{enumerate}
\item Jarod Alper a rédigé un chapitre consacré à la bibliographie relative
aux champs algébriques, voir
Guide bibliographique, section \ref{guide-section-short-introductions}.
\item Bhargav Bhatt a r\'edig\'e la premi\`ere version d'un chapitre
sur les morphismes \'etales de sch\'emas, voir
Morphismes \'etales, section \ref{etale-section-introduction}.
\item Bhargav Bhatt a r\'edig\'e la premi\`ere version de
Compl\'ements d'alg\`ebre, section \ref{more-algebra-section-formal-glueing}.
\item Kiran Kedlaya a fourni la premi\`ere r\'edaction de
Descente, section \ref{descent-section-descent-universally-injective}.
\item Les premi\`eres versions de
\begin{enumerate}
\item Alg\`ebre, section \ref{algebra-section-oka-families},
\item Objets injectifs, section \ref{injectives-section-baer}, et
\item le chapitre sur les corps, voir
Corps, section \ref{fields-section-introduction},
\end{enumerate}
proviennent du
\href{https://math.uchicago.edu/~amathew/cr.html}{The CRing Project},
avec l'aimable concours d'Akhil Mathew et al.
\item Alex Perry a r\'edig\'e le contenu consacr\'e aux modules projectifs
et aux modules de Mittag-Leffler, y compris la preuve de
Alg\`ebre, th\'eor\`eme \ref{algebra-theorem-ffdescent-projectivity}.
\item Alex Perry a r\'edig\'e le chapitre sur la th\'eorie des d\'eformations
selon l'approche de Schlessinger et Rim, voir
Th\'eorie formelle des d\'eformations, section \ref{formal-defos-section-introduction}.
\item Thibaut Pugin, Zachary Maddock et Min Lee ont pris des notes pour un
cours qui a servi de base \`a un chapitre sur la cohomologie \'etale
et \`a un chapitre sur la formule des traces. Voir
Cohomologie \'etale, section \ref{etale-cohomology-section-introduction}
et
La formule des traces, section \ref{trace-section-introduction}.
\item David Rydh a apport\'e de nombreux commentaires utiles, signal\'e plusieurs
erreurs, apport\'e une aide essentielle au contenu consacr\'e aux gerbes r\'esiduelles
et est \`a l'origine du contenu de
Compl\'ements sur les groupo\"ides dans les espaces, sections
\ref{spaces-more-groupoids-section-finite} et
\ref{spaces-more-groupoids-section-etale-localize}.
\item Burt Totaro a apport\'e les contributions suivantes : Exemples, sections
\ref{examples-section-non-descending-property-projective},
\ref{examples-section-non-effective-descent-projective},
ainsi que
Propri\'et\'es des champs algébriques, section \ref{stacks-properties-section-dimension}.
\item Le chapitre sur la cohomologie pro-\'etale, voir
Cohomologie pro-\'etale, section \ref{proetale-section-introduction},
est tir\'e d'un article de Bhargav Bhatt et Peter Scholze.
\item Bhargav Bhatt a apport\'e les contributions suivantes : Exemples, sections
\ref{examples-section-non-algebraic-hom-stack} et
\ref{examples-section-flat-not-colimit-flat-finitely-presented}.
\item Ofer Gabber a relev\'e des erreurs et apport\'e des corrections ; il a
\'egalement apport\'e les contributions suivantes :
Vari\'et\'es, lemme \ref{varieties-lemma-image-connected-component},
Espaces algébriques formels, lemme \ref{formal-spaces-lemma-completion-countably-indexed},
le contenu de
Compl\'ements sur les schémas en groupoïdes, section \ref{more-groupoids-section-ind-quasi-affine},
le r\'esultat principal de
Propri\'et\'es des espaces algébriques, section \ref{spaces-properties-section-fpqc},
ainsi que la preuve de
Compl\'ements sur la platitude, proposition
\ref{flat-proposition-finite-type-injective-into-flat-mod-m}.
\item J\'anos Koll\'ar a apport\'e les contributions suivantes :
Alg\`ebre, lemme \ref{algebra-lemma-hart-serre-loc-thm}, et
Cohomologie locale, proposition \ref{local-cohomology-proposition-kollar}.
\item Kiran Kedlaya a r\'edig\'e la premi\`ere version de
Compl\'ements d'alg\`ebre, section \ref{more-algebra-section-beauville-laszlo}.
\item Matthew Emerton, Toby Gee et Brandon Levin ont apport\'e
certains r\'esultats sur les \'epaississements, en particulier
Compl\'ements sur les morphismes de champs, lemmes
\ref{stacks-more-morphisms-lemma-reduced-diagonal},
\ref{stacks-more-morphisms-lemma-thickening-diagonals}, et
\ref{stacks-more-morphisms-lemma-thickening-properties}.
\item Lena Min Ji a r\'edig\'e la premi\`ere version de
Compl\'ements d'alg\`ebre, section \ref{more-algebra-section-principal-radical-ideals}.
\item Matthew Emerton et Toby Gee ont r\'edig\'e les premi\`eres versions de
G\'eom\'etrie des champs,
sections \ref{stacks-geometry-section-multiplicities} et
\ref{stacks-geometry-section-dimension-of-algebraic-stacks}.
\end{enumerate}







% OK, start here
%

\chapter{Conventions}



\phantomsection
\label{conventions-section-phantom}


\section{Remarques}
\label{conventions-section-comments}

\noindent
La philosophie qui sous-tend les conventions adoptées dans ces textes consiste
à choisir celles qui fonctionnent.

\section{Théorie des ensembles}
\label{conventions-section-sets}

\noindent
Nous utilisons la théorie des ensembles de Zermelo--Fraenkel avec l'axiome du
choix. Voir \cite{Kunen}. Nous n'utilisons pas d'univers (contrairement à
SGA~4). Nous n'insistons pas sur les questions ensemblistes, mais nous veillons
à ce que tout soit correct (bien entendu) et ne les éludons donc pas davantage.


\section{Catégories}
\label{conventions-section-categories}

\noindent
Une catégorie $\mathcal{C}$ est constituée d'un ensemble d'objets et, pour
chaque couple d'objets, d'un ensemble de morphismes entre eux. Autrement dit,
c'est ce que d'autres textes appellent une catégorie « petite ». Nous
utiliserons également des catégories « grandes » (des catégories dont les
objets forment une classe propre), mais seulement celles qui sont énumérées
dans Catégories, remarque \ref{categories-remark-big-categories}.

\section{Algèbre}
\label{conventions-section-algebra}

\noindent
Dans ces notes, un anneau est un anneau commutatif muni d'une unité $1$. La catégorie des
anneaux possède donc un objet initial $\mathbf{Z}$ et un objet final $\{0\}$
(c'est l'unique anneau dans lequel $1 = 0$). Les modules sont supposés
unitaires. Voir \cite{Eisenbud}.

\section{Notations}
\label{conventions-section-notation}

\noindent
Les entiers naturels sont les éléments de
$\mathbf{N} = \{1, 2, 3, \ldots\}$.
Les entiers relatifs sont les éléments de
$\mathbf{Z} = \{\ldots, -2, -1, 0, 1, 2, \ldots\}$.
Le corps des nombres rationnels est noté $\mathbf{Q}$.
Le corps des nombres réels est noté $\mathbf{R}$.
Le corps des nombres complexes est noté $\mathbf{C}$.





% OK, start here.
%

\chapter{Théorie des ensembles}



\phantomsection
\label{sets-section-phantom}


\section{Introduction}
\label{sets-section-introduction}

\noindent
Nous avons besoin de quelques notions de théorie des ensembles de temps à autre.
Nous utilisons la théorie des ensembles de Zermelo--Fraenkel avec l'axiome du
choix (ZFC), telle qu'elle est décrite dans \cite{Kunen} et \cite{Jech}.

\section{Tout est un ensemble}
\label{sets-section-sets-everything}

\noindent
La plupart des mathématiciens considèrent que la théorie des ensembles fournit
les fondements de base des mathématiques. Mais comment cela fonctionne-t-il
réellement ? Par exemple, comment traduit-on la phrase
« $X$ est un schéma » en théorie des ensembles ? Il suffit de dérouler les
définitions : un schéma est un espace localement annelé tel que tout point
possède un voisinage ouvert qui est un schéma affine. Un espace localement
annelé est un espace annelé en anneaux locaux. Un espace annelé est une paire
$(X, \mathcal{O}_X)$ constituée
d'un espace topologique $X$ et d'un faisceau d'anneaux $\mathcal{O}_X$ sur
celui-ci. Un espace topologique est une paire $(X, \tau)$ constituée d'un
ensemble $X$ et d'un ensemble de sous-ensembles
$\tau \subset \mathcal{P}(X)$ satisfaisant les axiomes d'une topologie. Et
ainsi de suite.

\medskip\noindent
Ainsi, étant donné un ensemble $S$, comment reconnaîtrions-nous s'il s'agit
d'un schéma ? La première chose que nous cherchons à déterminer est si
l'ensemble $S$ est une paire ordonnée. Cela signifie par définition (voir
\cite{Jech}, page 7) que $S$ est de la forme
$(a, b) := \{\{a\}, \{a, b\}\}$ pour certains ensembles $a, b$. Si c'est le
cas, nous examinerions alors si $a$ est une paire ordonnée $(c, d)$. Si oui,
nous vérifierions si $d \subset \mathcal{P}(c)$, puis, dans l'affirmative, si
$d$ constitue la collection d'ensembles d'une topologie sur l'ensemble $c$.
Et ainsi de suite.

\medskip\noindent
Ainsi, même si la rédaction d'une formule complète $\phi_{scheme}(x)$ à une
variable libre $x$ en théorie des ensembles, exprimant la notion « $x$ est un
schéma », demanderait un travail considérable, elle est possible. Il devrait en
être de même pour tout objet mathématique.

\section{Classes}
\label{sets-section-classes}

\noindent
Nous utilisons informellement la notion de {\it classe}.\par\noindent
Étant donnée une formule $\phi(x, p_1, \ldots, p_n)$, nous appelons
$$
C = \{x : \phi(x, p_1, \ldots, p_n)\}
$$
une {\it classe}. Une classe est plus facile à manipuler que la formule qui la
définit, mais ce n'est pas, à proprement parler, un objet mathématique. Par
exemple, si $R$ est un anneau, nous pouvons considérer la classe de tous les
$R$-modules (puisque, après tout, nous pouvons traduire la phrase « $M$ est un
$R$-module » en une formule de théorie des ensembles, laquelle définit alors
une classe). Une {\it classe propre} est une classe qui n'est pas un ensemble.

\medskip\noindent
Nous pouvons ainsi considérer la catégorie des $R$-modules, qui est une
« grande » catégorie---autrement dit, elle possède une classe propre d'objets.
De même, nous pouvons considérer la « grande » catégorie des schémas, la
« grande » catégorie des anneaux, etc.



\section{Ordinaux}
\label{sets-section-ordinals}

\noindent
Un ensemble $T$ est {\it transitif} si $x\in T$ implique $x\subset T$.
Un ensemble $\alpha$ est un {\it ordinal} s'il est transitif et bien ordonné
par $\in$. Dans ce cas, nous définissons
$\alpha + 1 = \alpha \cup \{\alpha\}$, qui est un autre ordinal appelé le
{\it successeur} de $\alpha$. Un ordinal $\alpha$ est appelé un
{\it ordinal successeur} s'il existe un ordinal $\beta$ tel que
$\alpha = \beta + 1$. Le plus petit ordinal est $\emptyset$, également noté
$0$. Si $\alpha$ n'est ni $0$ ni un ordinal successeur, alors $\alpha$ est
appelé un {\it ordinal limite} et nous avons
$$
\alpha
=
\bigcup\nolimits_{\gamma \in \alpha} \gamma.
$$
Le premier ordinal limite est $\omega$ ; c'est aussi le premier ordinal
infini. Le premier ordinal non dénombrable $\omega_1$ est l'ensemble de tous
les ordinaux dénombrables. La collection de tous les ordinaux est une classe
propre. Elle est bien ordonnée par $\in$ au sens suivant : tout ensemble non
vide (et même toute classe non vide) d'ordinaux possède un plus petit élément.
Étant donné un ensemble $A$ d'ordinaux, nous définissons la {\it borne supérieure} de
$A$ par $\sup_{\alpha \in A} \alpha =
\bigcup_{\alpha \in A} \alpha$. C'est le plus petit ordinal supérieur ou égal
à tout $\alpha \in A$. Étant donné un ensemble bien ordonné $(S, <)$, il existe
un unique ordinal $\alpha$ tel que $(S, <) \cong (\alpha, \in)$ ; on l'appelle
le {\it type d'ordre} de l'ensemble bien ordonné.

\section{La hiérarchie des ensembles}
\label{sets-section-sets-hierarchy}

\noindent
Nous définissons par récurrence transfinie $V_0 = \emptyset$,
$V_{\alpha + 1} = P(V_\alpha)$ (ensemble des parties), et, pour un ordinal
limite $\alpha$,
$$
V_\alpha = \bigcup\nolimits_{\beta < \alpha} V_\beta.
$$
Remarquons que chaque $V_\alpha$ est un ensemble transitif.

\begin{lemma}
\label{sets-lemma-axiom-regularity}
Tout ensemble est un élément de $V_\alpha$ pour un certain ordinal $\alpha$.
\end{lemma}

\begin{proof}
Voir \cite[lemme 6.3]{Jech}.
\end{proof}

\noindent
Dans \cite[chapitre III]{Kunen}, il est expliqué que ce lemme équivaut à
l'axiome de fondation. Le {\it rang} d'un ensemble $S$ est le plus petit ordinal
$\alpha$ tel que $S \in V_{\alpha + 1}$. Par {\it univers partiel}, nous
entendrons un ensemble suffisamment grand de la forme $V_\alpha$, lequel sera
clair d'après le contexte.

\section{Cardinalité}
\label{sets-section-cardinals}

\noindent
La {\it cardinalité} d'un ensemble $A$ est le plus petit ordinal $\alpha$ tel
qu'il existe une bijection entre $A$ et $\alpha$. Nous utilisons parfois la
notation $\alpha = |A|$ pour l'indiquer. Nous disons qu'un ordinal $\alpha$ est
un {\it cardinal} si et seulement s'il apparaît comme cardinalité d'un certain
ensemble $A$---autrement dit, si $\alpha = |A|$. Nous utilisons les lettres
grecques $\kappa$, $\lambda$ pour les cardinaux. Le premier cardinal infini est
$\omega$ et, dans ce contexte, il est noté $\aleph_0$. Un ensemble est
{\it dénombrable} si sa cardinalité est $\leq \aleph_0$. Si $\alpha$ est un
ordinal, nous notons $\alpha^+$ le plus petit cardinal $> \alpha$. On peut
l'utiliser pour définir $\aleph_1 = \aleph_0^+$,
$\aleph_2 = \aleph_1^+$, etc., et l'on peut en fait définir $\aleph_\alpha$
pour tout ordinal $\alpha$ par récurrence transfinie. Nous relevons l'égalité
$\aleph_1 = \omega_1$.

\medskip\noindent
L'{\it addition} des cardinaux $\kappa, \lambda$ est notée
$\kappa \oplus \lambda$ ; c'est la cardinalité de
$\kappa \amalg \lambda$. La {\it multiplication} des cardinaux
$\kappa, \lambda$ est notée $\kappa \otimes \lambda$ ; c'est la cardinalité
de $\kappa \times \lambda$. Si $\kappa$ et $\lambda$ sont des cardinaux
infinis, alors
$\kappa \oplus \lambda = \kappa \otimes \lambda = \max(\kappa, \lambda)$.
L'{\it exponentiation} des cardinaux $\kappa, \lambda$ est notée
$\kappa^\lambda$ ; c'est la cardinalité de l'ensemble des applications
(ensemblistes) de $\lambda$ vers $\kappa$. Étant donné un ensemble $K$ de
cardinaux, la {\it borne supérieure} de $K$ est
$\sup_{\kappa \in K} \kappa = \bigcup_{\kappa \in K} \kappa$, qui est
également un cardinal.

\section{Cofinalité}
\label{sets-section-cofinality}

\noindent
Un {\it sous-ensemble cofinal} $S$ d'un ensemble bien ordonné $T$ est un
sous-ensemble $S \subset T$ tel que
$\forall t \in T \exists s\in S (t \leq s)$. Remarquons qu'un sous-ensemble
d'un ensemble bien ordonné est lui-même bien ordonné (pour l'ordre induit).
Étant donné un ordinal $\alpha$, la {\it cofinalité} $\text{cf}(\alpha)$ de
$\alpha$ est le plus petit ordinal $\beta$ qui apparaît comme type d'ordre
d'un sous-ensemble cofinal de $\alpha$. La cofinalité d'un ordinal est toujours
un cardinal. Nous pouvons donc aussi définir la cofinalité de $\alpha$ comme la
plus petite cardinalité d'un sous-ensemble cofinal de $\alpha$.

\begin{lemma}
\label{sets-lemma-map-from-set-lifts}
Supposons que $T = \colim_{\alpha < \beta} T_\alpha$ soit une colimite
d'ensembles indexée par les ordinaux strictement inférieurs à un ordinal donné
$\beta$. Supposons que $\varphi : S \to T$ soit une application d'ensembles.
Alors $\varphi$ se relève en une application à valeurs dans $T_\alpha$ pour un
certain $\alpha < \beta$, pourvu que $\beta$ ne soit pas la limite d'une
famille d'ordinaux indexée par $S$, autrement dit si $\beta$ est un ordinal tel que
$\text{cf}(\beta) > |S|$.
\end{lemma}

\begin{proof}
Pour chaque élément $s \in S$, choisissons un $\alpha_s < \beta$ et un élément
$t_s \in T_{\alpha_s}$ qui s'envoie sur $\varphi(s)$ dans $T$. Par hypothèse,
$\alpha = \sup_{s \in S} \alpha_s$ est strictement inférieur à $\beta$.
L'application $\varphi_\alpha : S \to T_\alpha$ qui associe à $s$ l'image de
$t_s$ dans $T_\alpha$ est donc une solution.
\end{proof}

\noindent
Le résultat suivant est essentiellement l'argument de Grothendieck établissant
l'existence d'ordinaux de cofinalité arbitrairement grande, argument qu'il a
utilisé pour démontrer l'existence d'assez d'injectifs dans certaines
catégories abéliennes ; voir \cite{Tohoku}.

\begin{proposition}
\label{sets-proposition-exist-ordinals-large-cofinality}
Soit $\kappa$ un cardinal. Il existe alors un ordinal dont la cofinalité est
strictement supérieure à $\kappa$.
\end{proposition}

\begin{proof}
Si $\kappa$ est fini, alors $\omega = \text{cf}(\omega)$ convient. Supposons
donc $\kappa$ infini. Considérons le plus petit ordinal $\alpha$ dont la
cardinalité est strictement supérieure à $\kappa$. Nous affirmons que
$\text{cf}(\alpha) > \kappa$. Remarquons que $\alpha$ est un ordinal limite,
car si $\alpha = \beta + 1$, alors $|\alpha| = |\beta|$ (puisque $\alpha$ et
$\beta$ sont infinis), ce qui contredit la minimalité de $\alpha$. (Bien sûr,
$\alpha$ est aussi un cardinal, mais nous n'en avons pas besoin.) Pour obtenir
une contradiction, supposons que $S \subset \alpha$ soit un sous-ensemble
cofinal tel que $|S| \leq \kappa$. Pour $\beta \in S$, c'est-à-dire
$\beta < \alpha$, nous avons $|\beta| \leq \kappa$ par minimalité de $\alpha$.
Comme $\alpha$ est un ordinal limite et $S$ est cofinal dans $\alpha$, nous
obtenons $\alpha = \bigcup_{\beta \in S} \beta$. Par conséquent,
$|\alpha| \leq |S| \otimes \kappa \leq \kappa \otimes \kappa \leq \kappa$,
ce qui contredit notre choix de $\alpha$.
\end{proof}



\section{Principe de réflexion}
\label{sets-section-reflection-principle}

\noindent
Une partie de ce contenu se trouve dans le chapitre de \cite{Kunen} intitulé
« Easy consistency proofs ».

\medskip\noindent
Soit $\phi(x_1, \ldots, x_n)$ une formule de théorie des ensembles. Adoptons la
convention selon laquelle cette notation implique que toutes les variables
libres de $\phi$ figurent parmi $x_1, \ldots, x_n$. Soit $M$ un ensemble. La
formule $\phi^M(x_1, \ldots, x_n)$ est celle que l'on obtient à partir de
$\phi(x_1, \ldots, x_n)$ en remplaçant respectivement tous les $\forall x$ et
$\exists x$ par $\forall x\in M$ et $\exists x\in M$. Ainsi, la formule
$\phi(x_1, x_2) = \exists x (x\in x_1 \wedge x\in x_2)$ devient
$\phi^M(x_1, x_2) = \exists x \in M (x\in x_1 \wedge x\in x_2)$.
La formule $\phi^M$ est appelée la {\it relativisation de $\phi$ à $M$}.

\begin{theorem}
\label{sets-theorem-reflection-principle}
Considérons une famille {\bf finie} de formules de théorie des ensembles,
notées $\phi_1(x_1, \ldots, x_n), \ldots, \phi_m(x_1, \ldots, x_n)$.
Soit $M_0$ un ensemble. Il existe un ensemble $M$ tel
que $M_0 \subset M$ et que l'on ait, $\forall x_1, \ldots, x_n \in M$,
$$
\forall i = 1, \ldots, m, \ \phi_i^{M}(x_1, \ldots, x_n)
\Leftrightarrow
\forall i = 1, \ldots, m, \ \phi_i(x_1, \ldots, x_n).
$$
En fait, nous pouvons prendre $M = V_\alpha$ pour un certain ordinal limite
$\alpha$.
\end{theorem}

\begin{proof}
Voir \cite[théorème 12.14]{Jech} ou \cite[théorème 7.4]{Kunen}.
\end{proof}

\noindent
Nous interprétons ce théorème comme suit : étant donnés
$x_1, \ldots, x_n \in M$, les formules sont vraies lorsque les variables liées
parcourent tous les ensembles si et seulement si elles sont vraies lorsque les
variables liées parcourent les éléments de $V_\alpha$. Ce théorème est un
métathéorème, car il porte directement sur les formules de la théorie des
ensembles. Il affirme en réalité qu'étant donnée la liste finie de formules
$\phi_1, \ldots, \phi_m$, dont toutes les variables libres figurent parmi
$x_1, \ldots, x_n$, la phrase
$$
\begin{matrix}
\forall M_0\ \exists M, \ M_0 \subset M\ \forall x_1, \ldots, x_n \in M \\
\phi_1(x_1, \ldots, x_n) \wedge \ldots \wedge \phi_m(x_1, \ldots, x_n)
\leftrightarrow
\phi_1^M(x_1, \ldots, x_n) \wedge \ldots \wedge \phi_m^M(x_1, \ldots, x_n)
\end{matrix}
$$
est démontrable dans ZFC. Autrement dit, chaque fois que nous écrivons
effectivement une liste finie de formules $\phi_i$, nous obtenons un théorème.

\medskip\noindent
Il est assez difficile d'utiliser ce théorème dans les « mathématiques
ordinaires », car le sens des formules $\phi_i^M(x_1, \ldots, x_n)$ n'est pas
très clair ! Nous utiliserons plutôt l'idée de la démonstration du principe de
réflexion pour établir directement les résultats d'existence dont nous avons
besoin.

\section{Construction de catégories de schémas}
\label{sets-section-categories-schemes}

\noindent
Nous allons expliquer comment appliquer ceci pour produire, à partir d'un
ensemble initial de schémas, une « petite » catégorie de schémas stable par une
liste d'opérations naturelles. Avant cela, introduisons la taille d'un schéma.
Étant donné un schéma $S$, nous définissons
$$
\text{size}(S) = \max(\aleph_0, \kappa_1, \kappa_2),
$$
où nous définissons les nombres cardinaux $\kappa_1$ et $\kappa_2$ comme suit :
\begin{enumerate}
\item Nous prenons pour $\kappa_1$ la cardinalité de l'ensemble des ouverts
affines de $S$.
\item Pour $\kappa_2$, nous prenons la borne supérieure des cardinalités des
$\Gamma(U, \mathcal{O}_S)$, lorsque $U \subset S$ est un ouvert affine.
\end{enumerate}

\begin{lemma}
\label{sets-lemma-bounded-size}
Pour tout cardinal $\kappa$, il existe un ensemble $A$ tel que chaque élément
de $A$ soit un schéma et que, pour tout schéma $S$ vérifiant
$\text{size}(S) \leq \kappa$, il existe un élément $X \in A$ tel que
$X \cong S$ (isomorphisme de schémas).
\end{lemma}

\begin{proof}
La démonstration est omise. Indication : réfléchir à la manière dont tout schéma est isomorphe à un
schéma obtenu par recollement de schémas affines.
\end{proof}

\noindent
Nous notons $Bound$ la fonction qui associe à chaque cardinal $\kappa$
\begin{equation}
\label{sets-equation-bound}
Bound(\kappa) = \max\{\kappa^{\aleph_0}, \kappa^+\}.
\end{equation}
Nous pourrions faire croître cette fonction beaucoup plus rapidement ; par
exemple, nous pourrions poser $Bound(\kappa) = \kappa^\kappa$, et le résultat
ci-dessous resterait vrai. Pour tout ordinal $\alpha$, nous notons
$\Sch_\alpha$ la sous-catégorie pleine de la catégorie des schémas dont les
objets sont des éléments de $V_\alpha$. Voici le résultat que nous allons
démontrer.

\begin{lemma}
\label{sets-lemma-construct-category}
Avec les notations $\text{size}$, $Bound$ et $\Sch_\alpha$ ci-dessus, soit
$S_0$ un ensemble de schémas. Il existe un ordinal limite $\alpha$ possédant
les propriétés suivantes :
\begin{enumerate}
\item
\label{sets-item-inclusion}
Nous avons $S_0 \subset V_\alpha$ ; autrement dit,
$S_0 \subset \Ob(\Sch_\alpha)$.
\item
\label{sets-item-bounded}
Pour tout $S \in \Ob(\Sch_\alpha)$ et tout schéma $T$ tel que
$\text{size}(T) \leq Bound(\text{size}(S))$, il existe un schéma
$S' \in \Ob(\Sch_\alpha)$ tel que $T \cong S'$.
\item
\label{sets-item-limit}
Pour toute catégorie d'indices dénombrable\footnote{L'ensemble des objets
et les ensembles de morphismes sont tous dénombrables. On peut en fait
démontrer le lemme en remplaçant $\aleph_0$ par un cardinal quelconque dans
(3) et (4).} $\mathcal{I}$ et tout foncteur
$F : \mathcal{I} \to \Sch_\alpha$, la limite $\lim_\mathcal{I} F$ existe dans
$\Sch_\alpha$ si et seulement si elle existe dans $\Sch$ et, de plus, dans ce
cas, le morphisme naturel entre elles est un isomorphisme.
\item
\label{sets-item-colimit}
Pour toute catégorie d'indices dénombrable $\mathcal{I}$ et tout foncteur
$F : \mathcal{I} \to \Sch_\alpha$, la colimite $\colim_\mathcal{I} F$ existe
dans $\Sch_\alpha$ si et seulement si elle existe dans $\Sch$ et, de plus,
dans ce cas, le morphisme naturel entre elles est un isomorphisme.
\end{enumerate}
\end{lemma}

\begin{proof}
Nous définissons, par récurrence transfinie, une fonction $f$ qui associe un
ordinal à chaque ordinal comme suit. Posons $f(0) = 0$. Étant donné
$f(\alpha)$, nous définissons $f(\alpha + 1)$ comme le plus petit ordinal
$\beta$ pour lequel les conditions suivantes sont satisfaites :
\begin{enumerate}
\item Nous avons $\alpha + 1 \leq \beta$ et $f(\alpha) \leq \beta$.
\item Soient $S \in \Ob(\Sch_{f(\alpha)})$ et $T$ un schéma tel que
$\text{size}(T) \leq Bound(\text{size}(S))$. Il existe alors un schéma
$S' \in \Ob(\Sch_\beta)$ tel que $T \cong S'$.
\item Pour toute catégorie d'indices dénombrable $\mathcal{I}$ et tout
foncteur $F : \mathcal{I} \to \Sch_{f(\alpha)}$, si la limite
$\lim_\mathcal{I} F$ ou la colimite $\colim_\mathcal{I} F$ existe dans
$\Sch$, alors elle est isomorphe à un schéma de $\Sch_\beta$.
\end{enumerate}
Pour voir que $\beta$ existe, raisonnons comme suit. Puisque
$\Ob(\Sch_{f(\alpha)})$ est un ensemble, nous voyons que
$\kappa =
\sup_{S \in \Ob(\Sch_{f(\alpha)})} Bound(\text{size}(S))$
existe et est un cardinal. Soit $A$ un ensemble de schémas obtenu à partir de
$\kappa$ comme dans le lemme \ref{sets-lemma-bounded-size}. Il existe un ensemble
$CountCat$ de catégories dénombrables tel que toute catégorie dénombrable soit
isomorphe à un élément de $CountCat$. Nous pouvons donc supposer, au point (3)
ci-dessus, que $\mathcal{I}$ est un élément de $CountCat$. Cela signifie que
les paires $(\mathcal{I}, F)$ du point (3) parcourent un ensemble. Il existe
donc un ensemble $B$ dont les éléments sont des schémas tel que, pour toute
paire $(\mathcal{I}, F)$ comme au point (3), si la limite ou la colimite existe,
elle soit isomorphe à un élément de $B$. Par conséquent, si nous choisissons un
$\beta$ quelconque tel que $A \cup B \subset V_\beta$ et
$\beta > \max\{\alpha + 1, f(\alpha)\}$, alors (1)--(3) sont satisfaites.
Comme toute collection non vide d'ordinaux possède un plus petit élément, nous
voyons que $f(\alpha + 1)$ est bien définie. Enfin, si $\alpha$ est un ordinal
limite, nous posons
$f(\alpha) = \sup_{\alpha' < \alpha} f(\alpha')$.

\medskip\noindent
Choisissons $\beta_0$ tel que $S_0 \subset V_{\beta_0}$. Par construction,
$f(\beta) \geq \beta$, et nous voyons que
$S_0 \subset V_{f(\beta_0)}$ également. De plus, comme $f$ est croissante au
sens large, nous voyons que $S_0 \subset V_{f(\beta)}$ est vrai pour tout
$\beta \geq \beta_0$. Choisissons ensuite un ordinal quelconque
$\beta_1 > \beta_0$ de cofinalité
$\text{cf}(\beta_1) > \omega = \aleph_0$. C'est possible, car la cofinalité des
ordinaux devient arbitrairement grande ; voir la proposition
\ref{sets-proposition-exist-ordinals-large-cofinality}. Nous affirmons que
$\alpha = f(\beta_1)$ est une solution au problème posé dans le lemme.

\medskip\noindent
La première propriété du lemme découle de notre choix de
$\beta_1 > \beta_0$ ci-dessus.

\medskip\noindent
Puisque $\beta_1$ est un ordinal limite (sa cofinalité étant infinie), nous
obtenons $f(\beta_1) = \sup_{\beta < \beta_1} f(\beta)$. Ainsi,
$\{f(\beta) \mid \beta < \beta_1\} \subset f(\beta_1)$ est un sous-ensemble
cofinal. Nous voyons donc que
$$
V_\alpha = V_{f(\beta_1)} = \bigcup\nolimits_{\beta < \beta_1} V_{f(\beta)}.
$$
Soit maintenant $S \in \Ob(\Sch_\alpha)$. Nous définissons $\beta(S)$ comme le
plus petit ordinal $\beta$ tel que
$S \in \Ob(\Sch_{f(\beta)})$. D'après ce qui précède, nous voyons que
$\beta(S) < \beta_1$ toujours. Puisque
$\Ob(\Sch_{f(\beta + 1)}) \subset
\Ob(\Sch_\alpha)$, la construction de $f$ ci-dessus montre que la deuxième
propriété du lemme est satisfaite.

\medskip\noindent
Supposons que $\{S_1, S_2, \ldots\} \subset \Ob(\Sch_\alpha)$ soit une
collection dénombrable. Considérons la fonction
$\omega \to \beta_1$, $n \mapsto \beta(S_n)$. Comme la cofinalité de $\beta_1$
est $> \omega$, l'image de cette fonction ne peut pas être un sous-ensemble
cofinal. Il existe donc un $\beta < \beta_1$ tel que
$\{S_1, S_2, \ldots\} \subset \Ob(\Sch_{f(\beta)})$. Il s'ensuit que tout
foncteur $F : \mathcal{I} \to \Sch_\alpha$ se factorise par l'une des
sous-catégories $\Sch_{f(\beta)}$. Ainsi, s'il existe un schéma $X$ qui est la
colimite ou la limite du diagramme $F$, la construction de $f$ montre que $X$
est isomorphe à un objet de $\Sch_{f(\beta + 1)}$, qui est une sous-catégorie
de $\Sch_\alpha$. Cela démontre les deux dernières assertions du lemme.
\end{proof}

\begin{remark}
\label{sets-remark-how-to-use-reflection}
Le lemme ci-dessus peut également être démontré à l'aide du principe de
réflexion. Il faut toutefois être prudent. Supposons en effet que la phrase
$\phi_{scheme}(X)$ exprime la propriété « $X$ est un schéma » ; que signifie
alors la formule $\phi_{scheme}^{V_\alpha}(X)$ ? Il est vrai que le principe de
réflexion affirme que nous pouvons trouver $\alpha$ tel que, pour tout
$X \in V_\alpha$, nous ayons
$\phi_{scheme}(X) \leftrightarrow \phi_{scheme}^{V_\alpha}(X)$,
mais cela ne sert absolument à rien. C'est seulement en combinant deux énoncés
de ce type que quelque chose d'intéressant se produit. Supposons par exemple
que $\phi_{red}(X, Y)$ exprime la propriété « $X$, $Y$ sont des schémas et $Y$
est la réduction de $X$ » (voir Schémas, définition
\ref{schemes-definition-reduced-induced-scheme}). Supposons que nous appliquions
le principe de réflexion à la paire de formules
$\phi_1(X, Y) = \phi_{red}(X, Y)$,
$\phi_2(X) = \exists Y, \phi_1(X, Y)$. Il est alors facile de voir que tout
$\alpha$ produit par le principe de réflexion possède la propriété suivante :
pour tout $X \in \Ob(\Sch_\alpha)$, la réduction de $X$ est également un objet
de $\Sch_\alpha$ (vérification laissée en exercice).
\end{remark}

\begin{lemma}
\label{sets-lemma-bound-affine}
Soit $S$ un schéma affine. Soit $R = \Gamma(S, \mathcal{O}_S)$. Alors la taille
de $S$ est égale à $\max\{ \aleph_0, |R|\}$.
\end{lemma}

\begin{proof}
Il y a au plus $\max\{|R|, \aleph_0\}$ ouverts affines de $\Spec(R)$. Cela est
clair, car tout ouvert affine $U \subset \Spec(R)$ est une union finie
d'ouverts principaux $D(f_1) \cup \ldots \cup D(f_n)$ ; le nombre d'ouverts
affines est donc au plus
$\sup_n |R|^n = \max\{|R|, \aleph_0\}$, voir
\cite[chap. I, 10.13]{Kunen}. D'autre part, nous voyons que
$\Gamma(U, \mathcal{O}) \subset R_{f_1} \times \ldots \times R_{f_n}$ et donc
$|\Gamma(U, \mathcal{O})| \leq
\max\{\aleph_0, |R_{f_1}|, \ldots, |R_{f_n}|\}$. Il suffit donc de démontrer
que $|R_f| \leq \max\{\aleph_0, |R|\}$, ce que nous omettons.
\end{proof}

\begin{lemma}
\label{sets-lemma-bound-size}
Soit $S$ un schéma. Soit $S = \bigcup_{i \in I} S_i$ un recouvrement ouvert.
Alors
$\text{size}(S) \leq \max\{|I|, \sup_i\{\text{size}(S_i)\}\}$.
\end{lemma}

\begin{proof}
Soit $U \subset S$ un ouvert affine quelconque. Comme $U$ est quasi-compact,
il existe un nombre fini d'éléments $i_1, \ldots, i_n \in I$ et des ouverts
affines $U_i \subset U \cap S_i$ tels que
$U = U_1 \cup U_2 \cup \ldots \cup U_n$. Ainsi
$$
|\Gamma(U, \mathcal{O}_U)|
\leq
|\Gamma(U_1, \mathcal{O})|
\otimes
\ldots
\otimes
|\Gamma(U_n, \mathcal{O})|
\leq \sup\nolimits_i\{\text{size}(S_i)\}
$$
De plus, cela montre que la cardinalité de l'ensemble des ouverts affines de
$S$ est inférieure ou égale à celle de l'ensemble
$$
\coprod_{n \in \omega}
\coprod_{i_1, \ldots, i_n \in I}
\{\text{ouverts affines de }S_{i_1}\}
\times
\ldots
\times
\{\text{ouverts affines de }S_{i_n}\}.
$$
La cardinalité de chacun des ensembles figurant dans l'union disjointe est
majorée par $\sup_i\{\text{size}(S_i)\}$. Celle de l'ensemble d'indices est
au plus égale à $\max\{|I|, \aleph_0\}$, voir
\cite[chap. I, 10.13]{Kunen}. Par conséquent, d'après
\cite[lemme 5.8]{Jech}, la cardinalité du coproduit est au plus
$$
\max\{\aleph_0, |I|\}
\otimes \sup_i\{\text{size}(S_i)\}
$$
Le lemme en résulte.
\end{proof}

\begin{lemma}
\label{sets-lemma-bound-size-fibre-product}
Soient $f : X \to S$, $g : Y \to S$ des morphismes de schémas. Alors nous
avons
$\text{size}(X \times_S Y) \leq \max\{\text{size}(X), \text{size}(Y)\}$.
\end{lemma}

\begin{proof}
Soit $S = \bigcup_{k \in K} S_k$ un recouvrement ouvert affine. Soient
$X = \bigcup_{i \in I} U_i$, $Y = \bigcup_{j \in J} V_j$ des recouvrements
ouverts affines, où $I$, $J$ ont des cardinalités
$\leq \text{size}(X), \text{size}(Y)$. Pour chaque $i \in I$, il existe un
ensemble fini $K_i$ d'éléments $k \in K$ tel que
$f(U_i) \subset \bigcup_{k \in K_i} S_k$. Pour chaque $j \in J$, il existe un
ensemble fini $K_j$ d'éléments $k \in K$ tel que
$g(V_j) \subset \bigcup_{k \in K_j} S_k$. Par conséquent, $f(X), g(Y)$ sont
contenus dans $S' = \bigcup_{k \in K'} S_k$, où
$K' = \bigcup_{i \in I} K_i \cup \bigcup_{j \in J} K_j$. Remarquons que la
cardinalité de $K'$ est au plus
$\max\{\aleph_0, |I|, |J|\}$. En appliquant le lemme
\ref{sets-lemma-bound-size}, nous voyons qu'il suffit de démontrer que
$\text{size}(f^{-1}(S_k) \times_{S_k} g^{-1}(S_k))
\leq \max\{\text{size}(X), \text{size}(Y)\}$ pour $k \in K'$. Autrement dit,
nous pouvons supposer que $S$ est affine.

\medskip\noindent
Supposons $S$ affine. Soient
$X = \bigcup_{i \in I} U_i$, $Y = \bigcup_{j \in J} V_j$ des recouvrements
ouverts affines, où $I$, $J$ ont des cardinalités
$\leq \text{size}(X), \text{size}(Y)$. De nouveau, d'après le lemme
\ref{sets-lemma-bound-size}, il suffit de démontrer le lemme pour les produits
$U_i \times_S V_j$. D'après le lemme \ref{sets-lemma-bound-affine}, nous voyons
qu'il suffit de montrer que
$$
|A \otimes_C B| \leq \max\{\aleph_0, |A|, |B|\}.
$$
Nous omettons la démonstration de cette inégalité.
\end{proof}

\begin{lemma}
\label{sets-lemma-bound-finite-type}
Soit $S$ un schéma. Soit $f : X \to S$ un morphisme localement de type fini,
où $X$ est quasi-compact. Alors
$\text{size}(X) \leq \text{size}(S)$.
\end{lemma}

\begin{proof}
Nous pouvons trouver un recouvrement ouvert affine fini
$X = \bigcup_{i = 1, \ldots n} U_i$ tel que chaque $U_i$ s'envoie dans un
ouvert affine $S_i$ de $S$. Ainsi, d'après le lemme
\ref{sets-lemma-bound-size}, nous nous ramenons au cas où $S$ et $X$ sont tous deux
affines. Dans ce cas, le lemme \ref{sets-lemma-bound-affine} montre qu'il suffit de
prouver
$$
|A[x_1, \ldots, x_n]| \leq \max\{\aleph_0, |A|\}.
$$
Nous omettons la démonstration de cette inégalité.
\end{proof}

\noindent
Dans Algèbre, lemme \ref{algebra-lemma-epimorphism-cardinality}, nous
montrerons que si $A \to B$ est un épimorphisme d'anneaux, alors
$|B| \leq \max(|A|, \aleph_0)$. L'analogue pour les schémas est le lemme
suivant.

\begin{lemma}
\label{sets-lemma-bound-monomorphism}
Soit $f : X \to Y$ un monomorphisme de schémas. Si au moins l'une des
propriétés suivantes est satisfaite, alors
$\text{size}(X) \leq \text{size}(Y)$ :
\begin{enumerate}
\item $f$ est quasi-compact,
\item $f$ est localement de présentation finie,
\item ajouter ici d'autres cas selon les besoins.
\end{enumerate}
Mais cette borne n'est pas valable pour les monomorphismes localement de type
fini.
\end{lemma}

\begin{proof}
Soit $Y = \bigcup_{j \in J} V_j$ un recouvrement ouvert affine de $Y$ tel que
$|J| \leq \text{size}(Y)$. D'après le lemme \ref{sets-lemma-bound-size}, il suffit
de borner la taille de l'image réciproque de $V_j$ dans $X$. Nous nous ramenons
donc au cas où $Y$ est affine, disons $Y = \Spec(B)$. Pour tout ouvert affine
$\Spec(A) \subset X$, nous avons
$|A| \leq \max(|B|, \aleph_0) = \text{size}(Y)$ ; voir la remarque ci-dessus
et le lemme \ref{sets-lemma-bound-affine}. Il suffit donc de montrer que $X$ possède
au plus $\text{size}(Y)$ ouverts affines. Cela est clair si $X$ est
quasi-compact, d'où le cas (1). Dans le cas (2), le nombre de classes
d'isomorphisme de $B$-algèbres $A$ susceptibles d'apparaître est borné par
$\text{size}(B)$, car chaque $A$ est de type fini sur $B$, donc isomorphe à une
algèbre $B[x_1, \ldots, x_n]/(f_1, \ldots, f_m)$ pour certains $n, m$ et
$f_j \in B[x_1, \ldots, x_n]$. Cependant, comme $X \to Y$ est un
monomorphisme, il existe au plus un morphisme $\Spec(A) \to X$ au-dessus de
$Y = \Spec(B)$. Le nombre d'ouverts affines de $X$ est donc borné par le nombre
de ces classes d'isomorphisme.

\medskip\noindent
Pour démontrer la dernière assertion du lemme, considérons l'anneau
$B = \prod_{n \in \mathbf{N}} \mathbf{F}_2$ et posons $Y = \Spec(B)$. Pour
chaque ultrafiltre $\mathcal{U}$ sur $\mathbf{N}$, nous obtenons un idéal
maximal $\mathfrak m_\mathcal{U}$ de corps résiduel $\mathbf{F}_2$ ;
l'application $B \to \mathbf{F}_2$ envoie l'élément $(x_n)$ sur
$\lim_\mathcal{U} x_n$. Les détails sont omis. Le morphisme de schémas
$X = \coprod_\mathcal{U} \Spec(\mathbf{F}_2) \to Y$ est un monomorphisme,
puisque tous les points sont distincts. Cependant, la cardinalité de l'ensemble
des sous-schémas ouverts affines de $X$ est égale à celle de l'ensemble des
ultrafiltres sur $\mathbf{N}$, qui vaut $2^{2^{\aleph_0}}$. Nous concluons grâce
à $|B| = 2^{\aleph_0} < 2^{2^{\aleph_0}}$.
\end{proof}
\begin{lemma}
\label{sets-lemma-what-is-in-it}
Soit $\alpha$ un ordinal comme dans le lemme \ref{sets-lemma-construct-category}
ci-dessus. La catégorie $\Sch_\alpha$ vérifie les propriétés suivantes :
\begin{enumerate}
\item Si $X, Y, S \in \Ob(\Sch_\alpha)$, alors, pour tous morphismes
$f : X \to S$, $g : Y \to S$, le produit fibré $X \times_S Y$ dans
$\Sch_\alpha$ existe et est un produit fibré dans la catégorie des schémas.
\item Étant donnée toute famille au plus dénombrable $S_1, S_2, \ldots$
d'éléments de $\Ob(\Sch_\alpha)$, le coproduit $\coprod_i S_i$ existe dans
$\Ob(\Sch_\alpha)$ et est un coproduit dans la catégorie des schémas.
\item Pour tout $S \in \Ob(\Sch_\alpha)$ et toute immersion ouverte
$U \to S$, il existe un $V \in \Ob(\Sch_\alpha)$ tel que $V \cong U$.
\item Pour tout $S \in \Ob(\Sch_\alpha)$ et toute immersion fermée
$T \to S$, il existe un $S' \in \Ob(\Sch_\alpha)$ tel que $S' \cong T$.
\item Pour tout $S \in \Ob(\Sch_\alpha)$ et tout morphisme de type fini
$T \to S$, il existe un $S' \in \Ob(\Sch_\alpha)$ tel que $S' \cong T$.
\item Supposons que $S$ soit un schéma possédant un recouvrement ouvert
$S = \bigcup_{i \in I} S_i$ tel qu'il existe un $T \in \Ob(\Sch_\alpha)$
avec (a) $\text{size}(S_i) \leq \text{size}(T)^{\aleph_0}$ pour tout
$i \in I$, et (b) $|I| \leq \text{size}(T)^{\aleph_0}$.
Alors $S$ est isomorphe à un objet de $\Sch_\alpha$.
\item Pour tout $S \in \Ob(\Sch_\alpha)$ et tout morphisme $f : T \to S$
localement de type fini tel que $T$ puisse être recouvert par au plus
$\text{size}(S)^{\aleph_0}$ ouverts affines, il existe un
$S' \in \Ob(\Sch_\alpha)$ tel que $S' \cong T$.
Par exemple, cela vaut si $T$ peut être recouvert par au plus
$|\mathbf{R}| = 2^{\aleph_0} = \aleph_0^{\aleph_0}$ ouverts affines.
\item Pour tout $S \in \Ob(\Sch_\alpha)$ et tout monomorphisme $T \to S$
qui est soit localement de présentation finie, soit quasi-compact, il existe
un $S' \in \Ob(\Sch_\alpha)$ tel que $S' \cong T$.
\item Supposons que $T \in \Ob(\Sch_\alpha)$ soit affine. Écrivons
$R = \Gamma(T, \mathcal{O}_T)$. Alors chacun des schémas suivants est
isomorphe à un schéma de $\Sch_\alpha$ :
\begin{enumerate}
\item Pour tout idéal $I \subset R$, si $R^* = \lim_n R/I^n$ est le
complété, le schéma $\Spec(R^*)$.
\item Pour toute $R$-algèbre de type fini $R'$, le schéma $\Spec(R')$.
\item Pour toute localisation $S^{-1}R$, le schéma $\Spec(S^{-1}R)$.
\item Pour tout idéal premier $\mathfrak p \subset R$, le schéma
$\Spec(\overline{\kappa(\mathfrak p)})$.
\item Pour tout sous-anneau $R' \subset R$, le schéma $\Spec(R')$.
\item Tout schéma de type fini sur un anneau de cardinalité au plus
$|R|^{\aleph_0}$.
\item Et ainsi de suite.
\end{enumerate}
\end{enumerate}
\end{lemma}

\begin{proof}
Les assertions (1) et (2) découlent directement des définitions.
L'assertion (3) découle du fait que la taille d'un sous-schéma ouvert $U$ de
$S$ est clairement inférieure ou égale à celle de $S$.
L'assertion (4) découle de (5).
L'assertion (5) découle de (7).
L'assertion (6) découle du fait que la taille de $S$ est
$\leq \max\{|I|, \sup_i \text{size}(S_i)\} \leq \text{size}(T)^{\aleph_0}$
d'après le lemme \ref{sets-lemma-bound-size}. L'assertion (7) découle de (6).
En effet, pour tout ouvert affine $V \subset T$, nous avons
$\text{size}(V) \leq \text{size}(S)$ d'après le
lemme \ref{sets-lemma-bound-finite-type}.
Nous voyons donc que (6) s'applique dans la situation de (7).
La partie (8) découle du lemme \ref{sets-lemma-bound-monomorphism}.

\medskip\noindent
L'assertion (9) se traduit, via le lemme \ref{sets-lemma-bound-affine}, par une
borne supérieure pour la cardinalité des anneaux
$R^*$, $S^{-1}R$, $\overline{\kappa(\mathfrak p)}$, $R'$, etc.
Le cas le plus intéressant est peut-être l'anneau $R^*$. En tant qu'ensemble,
il est l'image d'une application surjective $R^{\mathbf{N}} \to R^*$.
Puisque $|R^{\mathbf{N}}| = |R|^{\aleph_0}$, la borne voulue résulte de ce que
nous avons choisi $Bound(\kappa)$ supérieur ou égal à $\kappa^{\aleph_0}$.
Ouf ! (La cardinalité de la clôture algébrique d'un corps
est la même que celle du corps, ou bien elle vaut $\aleph_0$.)
\end{proof}

\begin{remark}
\label{sets-remark-what-is-not-in-it}
Soit $R$ un anneau. Supposons que nous considérions l'anneau
$\prod_{\mathfrak p \in \Spec(R)} \kappa(\mathfrak p)$.
La cardinalité de cet anneau est bornée par $|R|^{2^{|R|}}$, mais n'est pas
bornée par $|R|^{\aleph_0}$ en général.
Par exemple, si $R = \mathbf{C}[x]$, elle n'est pas bornée par
$|R|^{\aleph_0}$, et si $R = \prod_{n \in \mathbf{N}} \mathbf{F}_2$,
elle n'est pas bornée par $|R|^{|R|}$.
Ainsi, le « Et ainsi de suite » du lemme \ref{sets-lemma-what-is-in-it} ci-dessus
doit être pris avec prudence. Bien entendu, s'il devient un jour nécessaire
de considérer ces anneaux dans des arguments relatifs à la cohomologie
fppf/\'etale, nous pourrons remplacer la fonction $Bound$ ci-dessus par la
fonction $\kappa \mapsto \kappa^{2^\kappa}$.
\end{remark}

\noindent
Dans le lemme suivant, nous utilisons la notion de recouvrement fpqc introduite
dans Topologies, section \ref{topologies-section-fpqc}.

\begin{lemma}
\label{sets-lemma-bound-by-covering}
Soit $f : X \to Y$ un morphisme de schémas. Supposons qu'il existe un
recouvrement fpqc $\{g_j : Y_j \to Y\}_{j \in J}$ tel que $g_j$ se factorise
par $f$. Alors $\text{size}(Y) \leq \text{size}(X)$.
\end{lemma}

\begin{proof}
Soit $V \subset Y$ un ouvert affine. Par définition, il existe un entier
$n \geq 0$, une application $a : \{1, \ldots, n\} \to J$ et des ouverts
affines $V_i \subset Y_{a(i)}$ tels que
$V = g_{a(1)}(V_1) \cup \ldots \cup g_{a(n)}(V_n)$.
Notons $h_j : Y_j \to X$ un morphisme tel que $f \circ h_j = g_j$.
Alors $h_{a(1)}(V_1) \cup \ldots \cup h_{a(n)}(V_n)$ est une partie
quasi-compacte de $f^{-1}(V)$. Nous pouvons donc trouver un ouvert
quasi-compact $W \subset f^{-1}(V)$ qui contient $h_{a(i)}(V_i)$ pour
$i = 1, \ldots, n$. En particulier, $V = f(W)$.

\medskip\noindent
D'une part, cela montre que la cardinalité de l'ensemble des ouverts affines
de $Y$ est au plus celle de l'ensemble $S$ des ouverts quasi-compacts de $X$.
Puisque tout ouvert quasi-compact de $X$ est une réunion finie d'ouverts
affines, nous voyons que la cardinalité de cet ensemble est au plus
$\sup |S|^n = \max(\aleph_0, |S|)$.
D'autre part, nous avons
$\mathcal{O}_Y(V) \subset \prod_{i = 1, \ldots, n} \mathcal{O}_{Y_{a(i)}}(V_i)$
car $\{V_i \to V\}$ est un recouvrement fpqc. Par conséquent,
$\mathcal{O}_Y(V) \subset \mathcal{O}_X(W)$ puisque $V_i \to V$ se factorise
par $W$. Comme $W$ possède lui aussi un recouvrement fini par des ouverts
affines de $X$, nous concluons que $|\mathcal{O}_Y(V)|$ est borné par la
taille de $X$. Le lemme découle maintenant de la définition de la taille d'un
schéma.
\end{proof}

\noindent
Dans le lemme suivant, nous utilisons la notion de recouvrement fppf introduite
dans Topologies, section \ref{topologies-section-fppf}.

\begin{lemma}
\label{sets-lemma-bound-fppf-covering}
Soit $\{f_i : X_i \to X\}_{i \in I}$ un recouvrement fppf d'un schéma.
Il existe un recouvrement fppf $\{W_j \to X\}_{j \in J}$ qui raffine
$\{X_i \to X\}_{i \in I}$ et tel que
$\text{size}(\coprod W_j) \leq \text{size}(X)$.
\end{lemma}

\begin{proof}
Choisissons un recouvrement ouvert affine $X = \bigcup_{a \in A} U_a$ avec
$|A| \leq \text{size}(X)$. Pour chaque $a$, nous pouvons choisir une partie
finie $I_a \subset I$ et, pour $i \in I_a$, un ouvert quasi-compact
$W_{a, i} \subset X_i$ tels que
$U_a = \bigcup_{i \in I_a} f_i(W_{a, i})$.
Alors $\text{size}(W_{a, i}) \leq \text{size}(X)$ d'après le
lemme \ref{sets-lemma-bound-finite-type}.
Nous concluons que
$$
\text{size}(\coprod_a \coprod_{i \in I_a} W_{i, a}) \leq \text{size}(X)
$$
d'après le lemme \ref{sets-lemma-bound-size}.
\end{proof}





\section{Ensembles munis d'une action de groupe}
\label{sets-section-sets-with-group-action}

\noindent
Soit $G$ un groupe. Nous notons $G\textit{-Ensembles}$ la « grande » catégorie
des $G$-ensembles. Pour tout ordinal $\alpha$, nous notons
$G\textit{-Ensembles}_\alpha$ la sous-catégorie pleine de $G\textit{-Ensembles}$ dont
les objets appartiennent à $V_\alpha$. Comme notion de taille d'un
$G$-ensemble, nous prenons
$\text{size}(S) = \max\{\aleph_0, |G|, |S|\}$ (où $|G|$ et $|S|$ sont les
cardinalités des ensembles sous-jacents). Comme ci-dessus, nous utilisons la
fonction $Bound(\kappa) = \kappa^{\aleph_0}$.

\begin{lemma}
\label{sets-lemma-sets-with-group-action}
Avec les notations $G$, $G\textit{-Ensembles}_\alpha$, $\text{size}$ et $Bound$
ci-dessus, soit $S_0$ un ensemble de $G$-ensembles.
Il existe un ordinal limite $\alpha$ ayant les propriétés suivantes :
\begin{enumerate}
\item Nous avons $S_0 \cup \{{}_GG\} \subset \Ob(G\textit{-Ensembles}_\alpha)$.
\item Soit $S \in \Ob(G\textit{-Ensembles}_\alpha)$.\par\noindent
Pour tout $G$-ensemble $T$ tel
que $\text{size}(T) \leq Bound(\text{size}(S))$, il existe un
$S' \in \Ob(G\textit{-Ensembles}_\alpha)$ qui est isomorphe à $T$.
\item Soit $\mathcal{I}$ une catégorie d'indices dénombrable. Pour tout foncteur
$F : \mathcal{I} \to G\textit{-Ensembles}_\alpha$, la limite
$\lim_\mathcal{I} F$ et la colimite $\colim_\mathcal{I} F$ existent dans
$G\textit{-Ensembles}_\alpha$ et coïncident avec ceux calculés dans $G\textit{-Ensembles}$.
\end{enumerate}
\end{lemma}

\begin{proof}
Omis. La démonstration est analogue à celle du
lemme \ref{sets-lemma-construct-category} ci-dessus, mais plus facile.
\end{proof}

\begin{lemma}
\label{sets-lemma-what-is-in-it-G-sets}
Soit $\alpha$ un ordinal comme dans le
lemme \ref{sets-lemma-sets-with-group-action} ci-dessus. La catégorie
$G\textit{-Ensembles}_\alpha$ vérifie les propriétés suivantes :
\begin{enumerate}
\item Le $G$-ensemble ${}_GG$ est un objet de $G\textit{-Ensembles}_\alpha$.
\item Les (co)produits, les produits fibrés et les sommes amalgamées existent
dans $G\textit{-Ensembles}_\alpha$ et coïncident avec ceux calculés dans
$G\textit{-Ensembles}$.
\item Étant donné un objet $U$ de $G\textit{-Ensembles}_\alpha$, toute partie
$G$-stable $O \subset U$ est isomorphe à un objet de
$G\textit{-Ensembles}_\alpha$.
\end{enumerate}
\end{lemma}

\begin{proof}
Omis.
\end{proof}

\section{Recouvrements d'un site}
\label{sets-section-coverings-site}

\noindent
Supposons que $\mathcal{C}$ soit une catégorie (comme dans Catégories,
définition \ref{categories-definition-category}) et que
$\text{Cov}(\mathcal{C})$ soit une classe propre de recouvrements satisfaisant
les propriétés (1), (2) et (3) de Sites,
définition \ref{sites-definition-site}. Nous les rappelons ici :
\begin{enumerate}
\item Si $V \to U$ est un isomorphisme, alors $\{V \to U\} \in
\text{Cov}(\mathcal{C})$.
\item Si $\{U_i \to U\}_{i\in I} \in \text{Cov}(\mathcal{C})$ et si, pour
chaque $i$, nous avons
$\{V_{ij} \to U_i\}_{j\in J_i} \in \text{Cov}(\mathcal{C})$, alors
$\{V_{ij} \to U\}_{i \in I, j\in J_i} \in \text{Cov}(\mathcal{C})$.
\item Si $\{U_i \to U\}_{i\in I}\in \text{Cov}(\mathcal{C})$ et si
$V \to U$ est un morphisme de $\mathcal{C}$, alors $U_i \times_U V$ existe
pour tout $i$ et
$\{U_i \times_U V \to V \}_{i\in I} \in \text{Cov}(\mathcal{C})$.
\end{enumerate}
Pour un ordinal $\alpha$, nous posons
$\text{Cov}(\mathcal{C})_\alpha = \text{Cov}(\mathcal{C}) \cap V_\alpha$.
Étant donnés un ordinal $\alpha$ et un cardinal $\kappa$, nous définissons
$\text{Cov}(\mathcal{C})_{\kappa, \alpha}$ comme l'ensemble des éléments
$\mathcal{U} =
\{\varphi_i : U_i \to U\}_{i\in I} \in \text{Cov}(\mathcal{C})_\alpha$
tels que $|I| \leq \kappa$.

\medskip\noindent
Rappelons la notion suivante ; voir Sites,
définition \ref{sites-definition-combinatorial-tautological}.
Deux familles de morphismes, $\{\varphi_i : U_i \to U\}_{i\in I}$ et
$\{\psi_j : W_j \to U\}_{j\in J}$, ayant la même cible dans $\mathcal{C}$,
sont dites {\it combinatoirement équivalentes} s'il existe des applications
$\alpha : I \to J$ et $\beta : J\to I$ telles que
$\varphi_i = \psi_{\alpha(i)}$ et $\psi_j = \varphi_{\beta(j)}$.
Cela définit une relation d'équivalence sur les familles de morphismes ayant
une cible fixée.

\begin{lemma}
\label{sets-lemma-coverings-site}
Avec les notations ci-dessus, soit
$\text{Cov}_0 \subset \text{Cov}(\mathcal{C})$ un ensemble contenu dans
$\text{Cov}(\mathcal{C})$.
Il existe un cardinal $\kappa$ et un ordinal limite $\alpha$ ayant les
propriétés suivantes :
\begin{enumerate}
\item Nous avons
$\text{Cov}_0 \subset \text{Cov}(\mathcal{C})_{\kappa, \alpha}$.
\item L'ensemble de recouvrements
$\text{Cov}(\mathcal{C})_{\kappa, \alpha}$ satisfait les propriétés (1), (2)
et (3) de Sites, définition \ref{sites-definition-site} (voir ci-dessus).
Autrement dit, $(\mathcal{C}, \text{Cov}(\mathcal{C})_{\kappa, \alpha})$
est un site.
\item Tout recouvrement de $\text{Cov}(\mathcal{C})$ est combinatoirement
équivalent à un recouvrement de
$\text{Cov}(\mathcal{C})_{\kappa, \alpha}$.
\end{enumerate}
\end{lemma}

\begin{proof}
Pour le démontrer, considérons d'abord l'ensemble $\mathcal{S}$ de tous les
ensembles de morphismes de $\mathcal{C}$ ayant une cible fixée.
Autrement dit, un élément de $\mathcal{S}$ est une partie $T$ de
$\text{Arrows}(\mathcal{C})$ telle que tous les éléments de $T$ aient la même
cible. Étant donnée une famille
$\mathcal{U} = \{\varphi_i : U_i \to U\}_{i\in I}$ de morphismes ayant une
cible fixée, nous définissons
$Supp(\mathcal{U}) = \{ \varphi \in \text{Arrows}(\mathcal{C})
\mid \exists i\in I, \varphi = \varphi_i\}$.
Remarquons que deux familles
$\mathcal{U} =  \{\varphi_i : U_i \to U\}_{i\in I}$ et
$\mathcal{V} = \{V_j \to V\}_{j \in J}$ sont combinatoirement équivalentes
si et seulement si $Supp(\mathcal{U}) = Supp(\mathcal{V})$.
Ensuite, nous définissons $\mathcal{S}_\tau \subset \mathcal{S}$ comme la
partie
$\mathcal{S}_\tau = \{ T \in \mathcal{S} \mid
\exists\ \mathcal{U} \in \text{Cov}(\mathcal{C}) \ T = Supp(\mathcal{U})\}$.
Pour tout élément $T \in \mathcal{S}_\tau$, posons $\beta(T)$ égal au plus
petit ordinal $\beta$ tel qu'il existe un
$\mathcal{U} \in \text{Cov}(\mathcal{C})_\beta$ tel que
$T = \text{Supp}(\mathcal{U})$. Enfin, posons
$\beta_0 = \sup_{T \in S_\tau} \beta(T)$.
Il en résulte que tout $\mathcal{U} \in \text{Cov}(\mathcal{C})$ est
combinatoirement équivalent à un élément de
$\text{Cov}(\mathcal{C})_{\beta_0}$.

\medskip\noindent
Soit $\kappa$ le maximum de $\aleph_0$, de la cardinalité
$|\text{Arrows}(\mathcal{C})|$,
$$
\sup\nolimits_{\{U_i \to U\}_{i\in I} \in \text{Cov}(\mathcal{C})_{\beta_0}}
|I|,
\quad\text{et}\quad
\sup\nolimits_{\{U_i \to U\}_{i\in I} \in \text{Cov}_0} |I|.
$$
Puisque $\kappa$ est un cardinal infini, nous avons
$\kappa \otimes \kappa = \kappa$. Remarquons qu'évidemment
$\text{Cov}(\mathcal{C})_{\beta_0} =
\text{Cov}(\mathcal{C})_{\kappa, \beta_0}$.

\medskip\noindent
Nous définissons, par récurrence transfinie, une fonction $f$ qui associe un
ordinal à tout ordinal de la manière suivante. Posons $f(0) = 0$.
Étant donné $f(\alpha)$, nous définissons $f(\alpha + 1)$ comme le plus petit
ordinal $\beta$ tel que les conditions suivantes soient satisfaites :
\begin{enumerate}
\item Nous avons $\alpha + 1 \leq \beta$ et $f(\alpha) \leq \beta$.
\item Si $\{U_i \to U\}_{i\in I}
\in \text{Cov}(\mathcal{C})_{\kappa, f(\alpha)}$ et si, pour tout $i$, nous
avons
$\{W_{ij} \to U_i\}_{j\in J_i}
\in \text{Cov}(\mathcal{C})_{\kappa, f(\alpha)}$,
alors
$\{W_{ij} \to U\}_{i \in I, j\in J_i}
\in \text{Cov}(\mathcal{C})_{\kappa, \beta}$.
\item Si $\{U_i \to U\}_{i\in I}
\in \text{Cov}(\mathcal{C})_{\kappa, \alpha}$ et si $W \to U$ est un
morphisme de $\mathcal{C}$, alors
$\{U_i \times_U W \to W \}_{i\in I}
\in \text{Cov}(\mathcal{C})_{\kappa, \beta}$.
\end{enumerate}
Pour voir que $\beta$ existe, remarquons que la collection de tous les
recouvrements $\{W_{ij} \to U\}$ et $\{U_i \times_U W \to W \}$ qui
interviennent dans (2) et (3) forme clairement un ensemble. Il existe donc un
ordinal $\beta$ tel que $V_\beta$ contienne tous ces recouvrements. En outre,
l'ensemble d'indices du recouvrement $\{W_{ij} \to U\}$ a pour cardinalité
$\sum_{i \in I} |J_i| \leq \kappa \otimes \kappa = \kappa$, et ces
recouvrements sont donc contenus dans
$\text{Cov}(\mathcal{C})_{\kappa, \beta}$.
Puisque toute collection non vide d'ordinaux possède un plus petit élément,
nous voyons que $f(\alpha + 1)$ est bien défini. Enfin, si $\alpha$ est un
ordinal limite, nous posons
$f(\alpha) = \sup_{\alpha' < \alpha} f(\alpha')$.

\medskip\noindent
Choisissons un ordinal $\beta_1$ tel que
$\text{Arrows}(\mathcal{C}) \subset V_{\beta_1}$,
$\text{Cov}_0 \subset V_{\beta_0}$ et $\beta_1 \geq \beta_0$.
Par construction, $f(\beta_1) \geq \beta_1$, et nous voyons que les mêmes
propriétés sont vérifiées pour $V_{f(\beta_1)}$. De plus, comme $f$ est
croissante au sens large, cela reste vrai pour tout $\beta \geq \beta_1$.
Choisissons ensuite un ordinal $\beta_2 > \beta_1$ de cofinalité
$\text{cf}(\beta_2) > \kappa$. C'est possible puisque la cofinalité des
ordinaux devient arbitrairement grande ; voir la
proposition \ref{sets-proposition-exist-ordinals-large-cofinality}.
Nous affirmons que la paire formée par $\kappa$ et $\alpha = f(\beta_2)$ est
une solution au problème posé dans le lemme.

\medskip\noindent
La première et la troisième propriétés du lemme découlent de nos choix de
$\kappa$, $\beta_2 > \beta_1 > \beta_0$ ci-dessus.

\medskip\noindent
Puisque $\beta_2$ est un ordinal limite (car sa cofinalité est infinie), nous
obtenons $f(\beta_2) = \sup_{\beta < \beta_2} f(\beta)$.
Ainsi, $\{f(\beta) \mid \beta < \beta_2\} \subset f(\beta_2)$ est une partie
cofinale. Nous voyons donc que
$$
V_\alpha = V_{f(\beta_2)} = \bigcup\nolimits_{\beta < \beta_2} V_{f(\beta)}.
$$
Maintenant, soit $\mathcal{U} \in \text{Cov}_{\kappa, \alpha}$.
Nous définissons $\beta(\mathcal{U})$ comme le plus petit ordinal $\beta$ tel
que $\mathcal{U} \in \text{Cov}_{\kappa, f(\beta)}$. D'après ce qui précède,
nous voyons que $\beta(\mathcal{U}) < \beta_2$ dans tous les cas.

\medskip\noindent
Nous devons montrer que les propriétés (1), (2) et (3) qui définissent un site
sont vérifiées pour la paire
$(\mathcal{C}, \text{Cov}_{\kappa, \alpha})$.
La première l'est parce que, par notre choix de $\beta_2$, toutes les flèches
de $\mathcal{C}$ appartiennent à $V_{f(\beta_2)}$.
Pour la troisième, étant donné un recouvrement
$\mathcal{U} = \{U_i \to U\}_{i \in I}
\in \text{Cov}(\mathcal{C})_{\kappa, \alpha}$, nous utilisons le fait que
$\beta(\mathcal{U}) < \beta_2$ ; tout changement de base de $\mathcal{U}$
appartient alors, par construction de $f$, à
$\text{Cov}(\mathcal{C})_{\kappa, f(\beta + 1)}$, et donc à
$\text{Cov}(\mathcal{C})_{\kappa, \alpha}$.

\medskip\noindent
Enfin, pour la deuxième condition, supposons que
$\{U_i \to U\}_{i\in I}
\in \text{Cov}(\mathcal{C})_{\kappa, f(\alpha)}$ et que, pour chaque $i$,
nous ayons
$\mathcal{W}_i = \{W_{ij} \to U_i\}_{j\in J_i}
\in \text{Cov}(\mathcal{C})_{\kappa, f(\alpha)}$.
Considérons la fonction
$I \to \beta_2$, $i \mapsto \beta(\mathcal{W}_i)$. Puisque la cofinalité
de $\beta_2$ est $> \kappa \geq |I|$, l'image de cette fonction ne peut être
une partie cofinale. Il existe donc un $\beta < \beta_1$ tel que
$\mathcal{W}_i \in \text{Cov}_{\kappa, f(\beta)}$ pour tout $i \in I$.
Il en résulte que le recouvrement
$\{W_{ij} \to U\}_{i\in I, j \in J_i}$ est un élément de
$\text{Cov}(\mathcal{C})_{\kappa, f(\beta + 1)}
\subset \text{Cov}(\mathcal{C})_{\kappa, \alpha}$, comme souhaité.
\end{proof}

\begin{remark}
\label{sets-remark-better}
Il est probable que, pour un certain ordinal limite $\alpha$, l'ensemble de
recouvrements $\text{Cov}(\mathcal{C})_\alpha$ satisfasse les conditions du
lemme. C'est après tout ce qu'une application du principe de réflexion semble
donner (aux réserves près décrites à la fin de la
section \ref{sets-section-reflection-principle} et dans la
remarque \ref{sets-remark-how-to-use-reflection}).
\end{remark}








\section{Catégories abéliennes et injectifs}
\label{sets-section-abelian-categories-injectives}

\noindent
Le lemme suivant s'applique à la catégorie des modules sur un faisceau
d'anneaux sur un site.

\begin{lemma}
\label{sets-lemma-abelian-injectives}
Supposons donnée une grande catégorie $\mathcal{A}$ (voir Catégories,
remarque \ref{categories-remark-big-categories}).
Supposons que $\mathcal{A}$ soit abélienne et possède suffisamment d'injectifs.
Voir Homologie, définitions \ref{homology-definition-abelian-category}
et \ref{homology-definition-enough-injectives}.
Alors, pour tout ensemble donné d'objets $\{A_s\}_{s\in S}$ de
$\mathcal{A}$, il existe une sous-catégorie abélienne
$\mathcal{A}' \subset \mathcal{A}$ ayant les propriétés suivantes :
\begin{enumerate}
\item le foncteur d'inclusion $\mathcal{A}' \to \mathcal{A}$ est exact,
\item $\Ob(\mathcal{A}')$ est un ensemble,
\item $\Ob(\mathcal{A}')$ contient $A_s$ pour tout $s \in S$,
\item $\mathcal{A}'$ possède suffisamment d'injectifs, et
\item un objet de $\mathcal{A}'$ est injectif si et seulement s'il est un
objet injectif de $\mathcal{A}$.
\end{enumerate}
\end{lemma}

\begin{proof}
Omis.
\end{proof}














% OK, start here.
%

\chapter{Catégories}



\phantomsection
\label{categories-section-phantom}


\section{Introduction}
\label{categories-section-introduction}

\noindent
Les catégories ont été introduites dans \cite{GenEqui}.
La catégorie des catégories (qui est une classe propre)
est une $2$-catégorie. De même, la catégorie des champs
forme une $2$-catégorie. Si vous connaissez déjà
les catégories, mais pas les $2$-catégories, vous devriez lire
la section \ref{categories-section-formal-cat-cat}
comme introduction aux définitions formelles données plus loin.

\section{Définitions}
\label{categories-section-definition-categories}

\noindent
Nous rappelons les définitions, en partie pour fixer les notations.

\begin{definition}
\label{categories-definition-category}
Une {\it catégorie} $\mathcal{C}$ est constituée des données suivantes :
\begin{enumerate}
\item Un ensemble d'objets $\Ob(\mathcal{C})$.
\item Pour chaque paire $x, y \in \Ob(\mathcal{C})$, un ensemble de morphismes
$\Mor_\mathcal{C}(x, y)$.
\item Pour chaque triplet $x, y, z\in \Ob(\mathcal{C})$, une application de
composition $ \Mor_\mathcal{C}(y, z) \times \Mor_\mathcal{C}(x, y)
\to \Mor_\mathcal{C}(x, z) $, notée $(\phi, \psi) \mapsto
\phi \circ \psi$.
\end{enumerate}
Ces données doivent satisfaire aux règles suivantes :
\begin{enumerate}
\item Pour tout élément $x\in \Ob(\mathcal{C})$, il existe un
morphisme $\text{id}_x\in \Mor_\mathcal{C}(x, x)$ tel que
$\text{id}_x \circ \phi = \phi$ et $\psi \circ \text{id}_x = \psi $ dès que
ces compositions ont un sens.
\item La composition est associative, c'est-à-dire
$(\phi \circ \psi) \circ \chi =
\phi \circ ( \psi \circ \chi)$ dès que ces compositions ont un sens.
\end{enumerate}
\end{definition}

\noindent
Il est d'usage d'exiger que tous les ensembles de morphismes
$\Mor_\mathcal{C}(x, y)$ soient disjoints.
Ainsi, un morphisme $\phi : x \to y$ possède une {\it source} $x$
unique et un {\it but} $y$ unique. Cela n'est pas strictement nécessaire,
bien qu'il faille alors prendre garde à la formulation de la condition (2)
ci-dessus. Cette convention est commode et nous la supposerons souvent
satisfaite. Dans ce cas, nous disons que $\phi$ et $\psi$ sont
{\it composables} si la source de $\phi$ est égale au
but de $\psi$ ; alors $\phi \circ \psi$ est défini.
Une définition équivalente consisterait à définir une catégorie
comme un quintuplet $(\text{Ob}, \text{Flèches}, s, t, \circ)$
formé d'un ensemble d'objets, d'un ensemble de morphismes (flèches),
d'une source, d'un but et d'une composition, soumis à une longue liste
d'axiomes. Nous adopterons occasionnellement ce point de vue.

\begin{remark}
\label{categories-remark-big-categories}
Grandes catégories. Dans certains textes, une catégorie peut avoir une classe
propre d'objets. Nous l'autoriserons également dans ces notes, mais seulement
dans les cas de la liste suivante (qui sera complétée au fur et à mesure).
En particulier, lorsque nous disons : « Soit $\mathcal{C}$ une catégorie »,
il est entendu que $\Ob(\mathcal{C})$ est un ensemble.
\begin{enumerate}
\item La catégorie $\textit{Ensembles}$ des ensembles.
\item La catégorie $\textit{Ab}$ des groupes abéliens.
\item La catégorie $\textit{Groups}$ des groupes.
\item Étant donné un groupe $G$, la catégorie $G\textit{-Ensembles}$ des
ensembles munis d'une action à gauche de $G$.
\item Étant donné un anneau $R$, la catégorie $\text{Mod}_R$ des $R$-modules.
\item Étant donné un corps $k$, la catégorie des espaces vectoriels sur $k$.
\item La catégorie des anneaux.
\item La catégorie des anneaux à puissances divisées, voir
Algèbre à puissances divisées, section \ref{dpa-section-divided-power-rings}.
\item La catégorie des schémas.
\item La catégorie $\textit{Top}$ des espaces topologiques.
\item Étant donné un espace topologique $X$, la catégorie
$\textit{PSh}(X)$ des préfaisceaux d'ensembles sur $X$.
\item Étant donné un espace topologique $X$, la catégorie
$\Sh(X)$ des faisceaux d'ensembles sur $X$.
\item Étant donné un espace topologique $X$, la catégorie
$\textit{PAb}(X)$ des préfaisceaux de groupes abéliens sur $X$.
\item Étant donné un espace topologique $X$, la catégorie
$\textit{Ab}(X)$ des faisceaux de groupes abéliens sur $X$.
\item Étant donné une petite catégorie $\mathcal{C}$, la catégorie des foncteurs
de $\mathcal{C}$ dans $\textit{Ensembles}$.
\item Étant donné une catégorie $\mathcal{C}$, la catégorie des préfaisceaux
d'ensembles sur $\mathcal{C}$.
\item Étant donné un site $\mathcal{C}$, la catégorie des faisceaux
d'ensembles sur $\mathcal{C}$.
\end{enumerate}
L'une des raisons de cette énumération est d'éviter de travailler
avec quelque chose comme la « collection » des « grandes » catégories,
ce qui reviendrait à travailler avec la collection de toutes les classes
et constituerait, à notre avis, un objet définitivement métamathématique.
\end{remark}

\begin{remark}
\label{categories-remark-unique-identity}
Il découle directement de la définition que deux morphismes identités
d'un objet $x$ de $\mathcal{A}$ sont égaux. Nous pouvons donc parler,
et parlerons, {\it du} morphisme identité $\text{id}_x$ de $x$.
\end{remark}

\begin{definition}
\label{categories-definition-isomorphism}
Un morphisme $\phi : x \to y$ est un {\it isomorphisme} de la catégorie
$\mathcal{C}$ s'il existe un morphisme $\psi : y \to x$
tel que $\phi \circ \psi = \text{id}_y$ et
$\psi \circ \phi = \text{id}_x$.
\end{definition}

\noindent
Un isomorphisme $\phi$ est aussi parfois appelé morphisme
{\it inversible}, et le morphisme $\psi$ de la définition est appelé
l'{\it inverse} et noté $\phi^{-1}$. Il est unique s'il existe. Remarquons que,
pour un objet $x$ d'une catégorie $\mathcal{A}$, l'ensemble des éléments
inversibles $\text{Aut}_\mathcal{A}(x)$
de $\Mor_\mathcal{A}(x, x)$ forme un groupe pour la composition.
Ce groupe est appelé le groupe des {\it automorphismes} de $x$ dans
$\mathcal{A}$.

\begin{definition}
\label{categories-definition-groupoid}
Un {\it groupoïde} est une catégorie dans laquelle tout morphisme est
un isomorphisme.
\end{definition}

\begin{example}
\label{categories-example-group-groupoid}
Un groupe $G$ définit un groupoïde ayant un unique objet $x$
et pour morphismes $\Mor(x, x) = G$, la loi de composition
étant donnée par la loi de groupe de $G$. Tout groupoïde ayant un unique
objet est de cette forme.
\end{example}

\begin{example}
\label{categories-example-set-groupoid}
Un ensemble $C$ définit un groupoïde $\mathcal{C}$ de la manière suivante :
on pose $\Ob(\mathcal{C}) := C$ et, pour les morphismes,
on prend $\Mor(x, y)$ vide si $x\neq y$ et égal à
$\{\text{id}_x\}$ si $x = y$.
\end{example}

\begin{definition}
\label{categories-definition-functor}
Un {\it foncteur} $F : \mathcal{A} \to \mathcal{B}$
entre deux catégories $\mathcal{A}, \mathcal{B}$ est défini par les
données suivantes :
\begin{enumerate}
\item Une application $F : \Ob(\mathcal{A}) \to \Ob(\mathcal{B})$.
\item Pour tous $x, y \in \Ob(\mathcal{A})$, une application
$F : \Mor_\mathcal{A}(x, y) \to \Mor_\mathcal{B}(F(x), F(y))$,
notée $\phi \mapsto F(\phi)$.
\end{enumerate}
Ces données doivent être compatibles avec la composition et les morphismes
identités de la manière suivante : $F(\phi \circ \psi) =
F(\phi) \circ F(\psi)$ pour toute paire composable $(\phi, \psi)$ de
morphismes de $\mathcal{A}$, et $F(\text{id}_x) = \text{id}_{F(x)}$.
\end{definition}

\noindent
Toute catégorie $\mathcal{A}$ possède un foncteur
{\it identité} $\text{id}_\mathcal{A}$.
De plus, étant donnés un foncteur $G : \mathcal{B} \to \mathcal{C}$
et un foncteur $F : \mathcal{A} \to \mathcal{B}$, il existe
un foncteur {\it composé} $G \circ F : \mathcal{A} \to \mathcal{C}$
défini de manière évidente.

\begin{definition}
\label{categories-definition-faithful}
Soit $F : \mathcal{A} \to \mathcal{B}$ un foncteur.
\begin{enumerate}
\item Nous disons que $F$ est {\it fidèle} si,
pour tous objets $x, y \in \Ob(\mathcal{A})$, l'application
$$
F : \Mor_\mathcal{A}(x, y) \to \Mor_\mathcal{B}(F(x), F(y))
$$
est injective.
\item Si toutes ces applications sont bijectives, $F$ est dit
{\it pleinement fidèle}.
\item
Le foncteur $F$ est dit {\it essentiellement surjectif} si, pour tout
objet $y \in \Ob(\mathcal{B})$, il existe un objet
$x \in \Ob(\mathcal{A})$ tel que $F(x)$ soit isomorphe à $y$ dans
$\mathcal{B}$.
\end{enumerate}
\end{definition}

\begin{definition}
\label{categories-definition-subcategory}
Une {\it sous-catégorie} d'une catégorie $\mathcal{B}$ est une catégorie
$\mathcal{A}$ dont les objets et les flèches forment des sous-ensembles des
objets et des flèches de $\mathcal{B}$, telle que la source, le but et la
composition dans $\mathcal{A}$ coïncident avec ceux de $\mathcal{B}$ et que
le morphisme identité d'un objet de $\mathcal{A}$ coïncide avec celui de
$\mathcal{B}$. Nous disons que $\mathcal{A}$ est une
{\it sous-catégorie pleine} de $\mathcal{B}$ si
$\Mor_\mathcal{A}(x, y)
= \Mor_\mathcal{B}(x, y)$ pour tous
$x, y \in \Ob(\mathcal{A})$.
Nous disons que $\mathcal{A}$ est une sous-catégorie
{\it strictement pleine} de $\mathcal{B}$ si elle est pleine et si,
pour tout $x \in \Ob(\mathcal{A})$, tout objet de $\mathcal{B}$
isomorphe à $x$ appartient également à $\mathcal{A}$.
\end{definition}

\noindent
Si $\mathcal{A} \subset \mathcal{B}$ est une sous-catégorie, l'application
identité est un foncteur de $\mathcal{A}$ dans $\mathcal{B}$.
De plus, une sous-catégorie $\mathcal{A} \subset \mathcal{B}$
est pleine si et seulement si le foncteur d'inclusion est pleinement fidèle.
Remarquons que, pour une catégorie $\mathcal{B}$, l'ensemble des
sous-catégories pleines de $\mathcal{B}$ s'identifie à l'ensemble des sous-ensembles de
$\Ob(\mathcal{B})$.

\begin{remark}
\label{categories-remark-functor-into-sets}
Supposons que $\mathcal{A}$ soit une catégorie.
Un foncteur $F$ de $\mathcal{A}$ dans $\textit{Ensembles}$
est un objet mathématique (c'est-à-dire un ensemble, et non une classe
ni une formule de la théorie des ensembles ; voir
Ensembles, section \ref{sets-section-sets-everything}),
bien que la catégorie des ensembles soit « grande ».
En effet, l'image de $F$ sur les objets est
un ensemble $F(\Ob(\mathcal{A}))$ et nous pouvons alors
considérer $F$ comme un foncteur entre
$\mathcal{A}$ et la sous-catégorie pleine
de la catégorie des ensembles dont les objets
sont les éléments de $F(\Ob(\mathcal{A}))$.
\end{remark}

\begin{example}
\label{categories-example-group-homomorphism-functor}
Un homomorphisme de groupes $p : G\to H$ définit un foncteur entre
les groupoïdes associés de l'exemple \ref{categories-example-group-groupoid}. Il est
fidèle (resp.\ pleinement fidèle) si et seulement si $p$ est injectif
(resp.\ est un isomorphisme).
\end{example}

\begin{example}
\label{categories-example-category-over-X}
Étant donnés une catégorie $\mathcal{C}$ et un objet
$X\in \Ob(\mathcal{C})$, nous définissons la
{\it catégorie des objets au-dessus de $X$},
notée $\mathcal{C}/X$, comme suit.
Les objets de $\mathcal{C}/X$ sont les morphismes $Y\to X$,
où $Y\in \Ob(\mathcal{C})$. Les morphismes entre les objets
$Y\to X$ et $Y'\to X$ sont les morphismes $Y\to Y'$ de $\mathcal{C}$
qui rendent commutatif le diagramme évident. Il existe un foncteur
$p_X : \mathcal{C}/X\to \mathcal{C}$ qui oublie simplement le morphisme.
De plus, tout morphisme $f : X'\to X$ de $\mathcal{C}$ induit un foncteur
$F : \mathcal{C}/X' \to \mathcal{C}/X$ obtenu par composition avec $f$,
et $p_X\circ F = p_{X'}$.
\end{example}

\begin{example}
\label{categories-example-category-under-X}
Étant donnés une catégorie $\mathcal{C}$ et un objet
$X\in \Ob(\mathcal{C})$, nous définissons la
{\it catégorie des objets au-dessous de $X$},
notée $X/\mathcal{C}$, comme suit.
Les objets de $X/\mathcal{C}$ sont les morphismes $X\to Y$,
où $Y\in \Ob(\mathcal{C})$. Les morphismes entre les objets
$X\to Y$ et $X\to Y'$ sont les morphismes $Y\to Y'$ de $\mathcal{C}$
qui rendent commutatif le diagramme évident. Il existe un foncteur
$p_X : X/\mathcal{C}\to \mathcal{C}$ qui oublie simplement le morphisme.
De plus, tout morphisme $f : X'\to X$ de $\mathcal{C}$ induit un foncteur
$F : X/\mathcal{C} \to X'/\mathcal{C}$
obtenu par composition avec $f$,
et $p_{X'}\circ F = p_X$.
\end{example}




\begin{definition}
\label{categories-definition-transformation-functors}
Soient $F, G : \mathcal{A} \to \mathcal{B}$ des foncteurs.
Une {\it transformation naturelle}, ou un {\it morphisme de foncteurs},
$t : F \to G$, est une famille $\{t_x\}_{x\in \Ob(\mathcal{A})}$
telle que
\begin{enumerate}
\item $t_x : F(x) \to G(x)$ soit un morphisme de la catégorie $\mathcal{B}$, et
\item pour tout morphisme $\phi : x \to y$ de $\mathcal{A}$, le diagramme
suivant soit commutatif :
$$
\xymatrix{
F(x) \ar[r]^{t_x} \ar[d]_{F(\phi)} & G(x) \ar[d]^{G(\phi)} \\
F(y) \ar[r]^{t_y} & G(y) }
$$
\end{enumerate}
\end{definition}

\noindent
Nous utilisons parfois le diagramme
$$
\xymatrix{
\mathcal{A}
\rtwocell^F_G{t}
&
\mathcal{B}
}
$$
pour indiquer que $t$ est un morphisme de $F$ vers $G$.

\medskip\noindent
Tout foncteur $F$ est muni de la transformation
{\it identité} $\text{id}_F : F \to F$. De plus, étant donnés un morphisme de
foncteurs $t : F \to G$ et un morphisme de foncteurs $s : E \to F$,
leur {\it composé} $t \circ s$ est défini par la règle
$$
(t \circ s)_x = t_x \circ s_x : E(x) \to G(x)
$$
pour $x \in \Ob(\mathcal{A})$.
On vérifie facilement qu'il s'agit bien d'un morphisme de foncteurs
de $E$ vers $G$.
Ainsi, à partir de deux catégories
$\mathcal{A}$ et $\mathcal{B}$, nous obtenons une nouvelle catégorie,
à savoir la catégorie des foncteurs de $\mathcal{A}$ dans $\mathcal{B}$.

\begin{remark}
\label{categories-remark-functors-sets-sets}
C'est l'un des cas où le même énoncé ne vaut pas si
$\mathcal{A}$ est une « grande » catégorie. Considérons par exemple les
foncteurs $\textit{Ensembles} \to \textit{Ensembles}$. Avec notre définition actuelle,
un tel foncteur est une classe et non un ensemble. Autrement dit, il est
donné par une formule de la théorie des ensembles (dont certaines variables
sont égales à des ensembles prescrits) ! Il n'est pas judicieux d'essayer
d'intégrer toutes les formules possibles de la théorie des ensembles à la
définition d'un objet mathématique. Le même problème se présente, par exemple,
lorsqu'on considère les faisceaux sur la catégorie des schémas.
Nous reviendrons plus tard sur ce point.
\end{remark}

\begin{definition}
\label{categories-definition-equivalence-categories}
Une {\it équivalence de catégories}
$F : \mathcal{A} \to \mathcal{B}$ est un foncteur tel qu'il existe
un foncteur $G : \mathcal{B} \to \mathcal{A}$ pour lequel
les composés $F \circ G$ et $G \circ F$ sont respectivement isomorphes aux
foncteurs identité $\text{id}_\mathcal{B}$ et
$\text{id}_\mathcal{A}$.
Dans ce cas, nous disons que $G$ est un {\it quasi-inverse} de $F$.
\end{definition}

\begin{lemma}
\label{categories-lemma-construct-quasi-inverse}
Soit $F : \mathcal{A} \to \mathcal{B}$ un foncteur pleinement fidèle.
Supposons donnés, pour tout $X \in \Ob(\mathcal{B})$, un
objet $j(X)$ de $\mathcal{A}$ et un isomorphisme $i_X : X \to F(j(X))$.
Il existe alors un unique foncteur $j : \mathcal{B} \to \mathcal{A}$
tel que l'action de $j$ sur les objets soit celle qui vient d'être prescrite et
pour lequel les isomorphismes
$i_X$ définissent un isomorphisme de foncteurs
$\text{id}_\mathcal{B} \to F \circ j$. De plus, $j$ et $F$
sont des équivalences de catégories quasi-inverses.
\end{lemma}

\begin{proof}
La construction de $j : \mathcal{B} \to \mathcal{A}$ se fait en deux étapes.
Premièrement, nous définissons l'application
$j : \Ob(\mathcal{B}) \to \Ob(\mathcal{A})$ qui associe
$j(X)$ à $X \in \mathcal{B}$. Deuxièmement, si
$X,Y \in \Ob(\mathcal{B})$ et $\phi : X \to Y$, nous considérons
$\phi' := i_Y \circ\phi\circ i_X^{-1}$. Il existe un unique
$\varphi$ tel que $F(\varphi) = \phi'$, puisque $F$ est pleinement fidèle.
Nous posons $j(\phi) = \varphi$. Nous omettons la vérification que $j$
est un foncteur. Par construction, le diagramme
$$
\xymatrix{
X \ar[r]_-{i_X} \ar[d]_{\phi}
&
F(j(X)) \ar[d]^{F(j(\phi))}
\\
Y \ar[r]^-{i_Y}
&
F(j(Y))
}
$$
est commutatif. Puisque chaque $i_X$ est un isomorphisme, la famille
$\{i_X\}_X$ est donc un isomorphisme de foncteurs
$\text{id}_\mathcal{B} \to F\circ j$. Pour conclure, il reste à montrer que
$j \circ F$ est isomorphe à $\text{id}_\mathcal{A}$.
Cependant, puisque $F$ est pleinement fidèle, il suffit pour cela de le
démontrer après composition à gauche par $F$, c'est-à-dire qu'il suffit de
montrer que $F \circ j \circ F$ est isomorphe à
$F \circ \text{id}_\mathcal{A}$
(nous omettons ce petit détail). Cela est clair puisque
$F \circ j \cong \text{id}_\mathcal{B}$.
\end{proof}

\begin{lemma}
\label{categories-lemma-equivalence-categories}
Un foncteur est une équivalence de catégories si et seulement s'il est
pleinement fidèle et essentiellement surjectif.
\end{lemma}

\begin{proof}
Soit $F : \mathcal{A} \to \mathcal{B}$ essentiellement surjectif et pleinement
fidèle. Puisque, par convention, toutes les catégories sont petites et que
$F$ est essentiellement surjectif, l'axiome du choix permet de choisir, pour
tout $X \in \Ob(\mathcal{B})$, un objet $j(X)$ de $\mathcal{A}$ et un
isomorphisme $i_X : X \to F(j(X))$. Nous appliquons alors le
lemme \ref{categories-lemma-construct-quasi-inverse},
en utilisant le fait que $F$ est pleinement fidèle.
\end{proof}

\begin{definition}
\label{categories-definition-product-category}
Soient $\mathcal{A}$ et $\mathcal{B}$ des catégories.
Nous définissons la {\it catégorie produit}
$\mathcal{A} \times \mathcal{B}$ comme la catégorie dont les objets sont
$\Ob(\mathcal{A} \times \mathcal{B}) =
\Ob(\mathcal{A}) \times \Ob(\mathcal{B})$
et telle que
$$
\Mor_{\mathcal{A} \times \mathcal{B}}((x, y), (x', y'))
:=
\Mor_\mathcal{A}(x, x')\times
\Mor_\mathcal{B}(y, y').
$$
La composition est définie composante par composante.
\end{definition}


\section{Catégories opposées et lemme de Yoneda}
\label{categories-section-opposite}

\begin{definition}
\label{categories-definition-opposite}
Étant donnée une catégorie $\mathcal{C}$, la {\it catégorie opposée}
$\mathcal{C}^{opp}$ est la catégorie qui a les mêmes objets
que $\mathcal{C}$, mais dont tous les morphismes sont renversés.
\end{definition}

\noindent
Autrement dit,
$\Mor_{\mathcal{C}^{opp}}(x, y) = \Mor_\mathcal{C}(y, x)$.
La composition dans $\mathcal{C}^{opp}$ est celle de $\mathcal{C}$,
mais prise à l'envers : si $\phi : y \to z$ et $\psi : x \to y$
sont des morphismes de $\mathcal{C}^{opp}$, c'est-à-dire des flèches
$z \to y$ et $y \to x$ dans $\mathcal{C}$,
alors $\phi \circ^{opp} \psi$ est le morphisme $x \to z$
de $\mathcal{C}^{opp}$ qui correspond au composé
$z \to y \to x$ dans $\mathcal{C}$.

\begin{definition}
\label{categories-definition-contravariant}
Soient $\mathcal{C}$ et $\mathcal{S}$ des catégories.
Un foncteur {\it contravariant} $F$
de $\mathcal{C}$ dans $\mathcal{S}$
est un foncteur $\mathcal{C}^{opp}\to \mathcal{S}$.
\end{definition}

\noindent
Concrètement, un foncteur contravariant $F$ est donné
par une application $F : \Ob(\mathcal{C}) \to
\Ob(\mathcal{S})$ et, pour tout morphisme
$\psi : x \to y$ dans $\mathcal{C}$, par un morphisme
$F(\psi) : F(y) \to F(x)$. Ces données doivent vérifier la propriété
suivante : étant donné un autre morphisme
$\phi : y \to z$, on a $F(\phi \circ \psi)
= F(\psi) \circ F(\phi)$ comme morphismes $F(z) \to F(x)$.
(Remarquer l'inversion de l'ordre.)

\begin{definition}
\label{categories-definition-presheaf}
Soit $\mathcal{C}$ une catégorie.
\begin{enumerate}
\item Un {\it préfaisceau d'ensembles sur $\mathcal{C}$},
ou simplement un {\it préfaisceau}, est un foncteur contravariant
$F$ de $\mathcal{C}$ dans $\textit{Ensembles}$.
\item La catégorie des préfaisceaux est notée $\textit{PSh}(\mathcal{C})$.
\end{enumerate}
\end{definition}

\noindent
Bien entendu, la catégorie des préfaisceaux est une classe propre.

\begin{example}
\label{categories-example-hom-functor}
Foncteur des points.
Pour tout $U\in \Ob(\mathcal{C})$, il existe un foncteur contravariant
$$
\begin{matrix}
h_U & : & \mathcal{C}
&
\longrightarrow
&
\textit{Ensembles} \\
& &
X
&
\longmapsto
&
\Mor_\mathcal{C}(X, U)
\end{matrix}
$$
qui associe à un objet $X$ l'ensemble
$\Mor_\mathcal{C}(X, U)$. Autrement dit, $h_U$ est un préfaisceau.
À un morphisme $f : X\to Y$, l'application correspondante
$h_U(f) : \Mor_\mathcal{C}(Y, U)\to \Mor_\mathcal{C}(X, U)$
envoie $\phi$ sur $\phi\circ f$. Nous noterons toujours ce préfaisceau
$h_U : \mathcal{C}^{opp} \to \textit{Ensembles}$.
On l'appelle le {\it préfaisceau représentable} associé à $U$.
Si $\mathcal{C}$ est la catégorie des schémas, ce foncteur est
parfois appelé le \emph{foncteur des points} de $U$.
\end{example}

\noindent
Remarquons qu'à un morphisme $\phi : U \to V$ de $\mathcal{C}$ correspond
une transformation naturelle de foncteurs $h(\phi) : h_U \to h_V$,
définie par composition avec le morphisme $U \to V$. Ainsi,
la composition des morphismes dans $\mathcal{C}$ devient la composition
des transformations de foncteurs. Autrement dit, nous obtenons un foncteur
$$
h :
\mathcal{C}
\longrightarrow
\textit{PSh}(\mathcal{C})
$$
Remarquons que le but est une « grande » catégorie ; voir la
remarque \ref{categories-remark-big-categories}. En revanche,
$h$ est bien un objet mathématique (c'est-à-dire\ un ensemble) ; comparer
la remarque \ref{categories-remark-functor-into-sets}.

\begin{lemma}[Lemme de Yoneda]
\label{categories-lemma-yoneda}
\begin{reference}
Est apparu sous une certaine forme dans \cite{Yoneda-homology}. Utilisé par
Grothendieck sous une forme généralisée dans \cite{Gr-II}.
\end{reference}
Soient $U, V \in \Ob(\mathcal{C})$.
Pour tout morphisme de foncteurs $s : h_U \to h_V$,
il existe un unique morphisme $\phi : U \to V$
tel que $h(\phi) = s$. Autrement dit, le
foncteur $h$ est pleinement fidèle. Plus généralement,
pour tout foncteur contravariant $F$ et tout objet
$U$ de $\mathcal{C}$, on a une bijection naturelle
$$
\Mor_{\textit{PSh}(\mathcal{C})}(h_U, F) \longrightarrow F(U),
\quad
s \longmapsto s_U(\text{id}_U).
$$
\end{lemma}

\begin{proof}
Pour le premier énoncé, il suffit de prendre
$\phi = s_U(\text{id}_U) \in \Mor_\mathcal{C}(U, V)$.
Pour le second, étant donné $\xi \in F(U)$, on définit
$s$ par $s_V : h_U(V) \to F(V)$ en envoyant l'élément $f : V \to U$
de $h_U(V) = \Mor_\mathcal{C}(V, U)$ sur $F(f)(\xi)$.
\end{proof}

\begin{definition}
\label{categories-definition-representable-functor}
Un foncteur contravariant $F : \mathcal{C}\to \textit{Ensembles}$ est dit
{\it représentable} s'il est isomorphe au foncteur des points
$h_U$ pour un certain objet $U$ de $\mathcal{C}$.
\end{definition}

\noindent
Soit $\mathcal{C}$ une catégorie et soit
$F : \mathcal{C}^{opp} \to \textit{Ensembles}$ un foncteur représentable.
Choisissons un objet $U$ de $\mathcal{C}$ et un isomorphisme $s : h_U \to F$.
Le lemme de Yoneda garantit que le couple $(U, s)$
est unique à isomorphisme unique près. L'objet
$U$ est appelé un objet {\it représentant} $F$.
Par le lemme de Yoneda, la transformation $s$ correspond à un unique
élément $\xi \in F(U)$. Cet élément est appelé l'{\it objet universel}.
Il possède la propriété que, pour $V \in \Ob(\mathcal{C})$, l'application
$$
\Mor_\mathcal{C}(V, U) \longrightarrow F(V),\quad
(f : V \to U) \longmapsto F(f)(\xi)
$$
est une bijection. Ainsi, $\xi$ est universel en ce sens que tout élément
de $F(V)$ est l'image de $\xi$ par $F(f)$ pour un unique morphisme
$f : V \to U$ dans $\mathcal{C}$.






\section{Produits de deux objets}
\label{categories-section-products-pairs}

\begin{definition}
\label{categories-definition-products}
Soient $x, y\in \Ob(\mathcal{C})$.
Un {\it produit} de $x$ et $y$ est
un objet $x \times y \in \Ob(\mathcal{C})$
muni de morphismes
$p\in \Mor_{\mathcal C}(x \times y, x)$ et
$q\in\Mor_{\mathcal C}(x \times y, y)$
satisfaisant la propriété universelle suivante : pour tout
$w\in \Ob(\mathcal{C})$ et tous morphismes
$\alpha \in \Mor_{\mathcal C}(w, x)$ et
$\beta \in \Mor_\mathcal{C}(w, y)$,
il existe un unique
$\gamma\in \Mor_{\mathcal C}(w, x \times y)$ qui rend le diagramme
$$
\xymatrix{
w \ar[rrrd]^\beta \ar@{-->}[rrd]_\gamma \ar[rrdd]_\alpha & & \\
& & x \times y \ar[d]_p \ar[r]_q & y \\
& & x &
}
$$
commutatif.
\end{definition}

\noindent
S'il existe, un produit est unique à isomorphisme unique près.
Cela découle du lemme de Yoneda, car la définition exige que
$x \times y$ soit un objet de $\mathcal{C}$ tel que
$$
h_{x \times y}(w) = h_x(w) \times h_y(w)
$$
fonctoriellement en $w$. Autrement dit, le produit $x \times y$
est un objet représentant le foncteur
$w \mapsto h_x(w) \times h_y(w)$.

\begin{definition}
\label{categories-definition-has-products-of-pairs}
Nous disons que la catégorie $\mathcal{C}$ {\it admet les produits de deux
objets} si un produit $x \times y$ existe pour tous
$x, y \in \Ob(\mathcal{C})$.
\end{definition}

\noindent
Nous employons cette terminologie pour distinguer cette notion de celles
d'« admettre les produits » ou d'« admettre les produits finis », qui ont
habituellement un autre sens (elles impliquent notamment toujours l'existence
d'un objet final).





\section{Coproduits de deux objets}
\label{categories-section-coproducts-pairs}

\begin{definition}
\label{categories-definition-coproducts}
Soient $x, y \in \Ob(\mathcal{C})$.
Un {\it coproduit}, ou une {\it somme}, de $x$ et $y$ est
un objet $x \amalg y \in \Ob(\mathcal{C})$
muni de morphismes
$i \in \Mor_{\mathcal C}(x, x \amalg y)$ et
$j \in \Mor_{\mathcal C}(y, x \amalg y)$
satisfaisant la propriété universelle suivante : pour tout
$w \in \Ob(\mathcal{C})$ et tous morphismes
$\alpha \in \Mor_{\mathcal C}(x, w)$ et
$\beta \in \Mor_\mathcal{C}(y, w)$,
il existe un unique
$\gamma \in \Mor_{\mathcal C}(x \amalg y, w)$ qui rend le diagramme
$$
\xymatrix{
& y \ar[d]^j \ar[rrdd]^\beta \\
x \ar[r]^i \ar[rrrd]_\alpha & x \amalg y \ar@{-->}[rrd]^\gamma \\
& & & w
}
$$
commutatif.
\end{definition}

\noindent
S'il existe, un coproduit est unique à isomorphisme unique près.
Cela découle du lemme de Yoneda (appliqué à la catégorie opposée),
car la définition exige que $x \amalg y$ soit un objet
de $\mathcal{C}$ tel que
$$
\Mor_\mathcal{C}(x \amalg y, w) =
\Mor_\mathcal{C}(x, w) \times \Mor_\mathcal{C}(y, w)
$$
fonctoriellement en $w$.

\begin{definition}
\label{categories-definition-has-coproducts-of-pairs}
Nous disons que la catégorie $\mathcal{C}$ {\it admet les coproduits de deux
objets} si un coproduit $x \amalg y$ existe pour tous
$x, y \in \Ob(\mathcal{C})$.
\end{definition}

\noindent
Nous employons cette terminologie pour distinguer cette notion de celles
d'« admettre les coproduits » ou d'« admettre les coproduits finis », qui ont
habituellement un autre sens (elles impliquent notamment toujours l'existence
d'un objet initial dans $\mathcal{C}$).





\section{Produits fibrés}
\label{categories-section-fibre-products}

\begin{definition}
\label{categories-definition-fibre-products}
Soient $x, y, z\in \Ob(\mathcal{C})$,
$f\in \Mor_\mathcal{C}(x, y)$
et $g\in \Mor_{\mathcal C}(z, y)$.
Un {\it produit fibré} de $f$ et $g$ est
un objet $x \times_y z\in \Ob(\mathcal{C})$
muni de morphismes
$p \in \Mor_{\mathcal C}(x \times_y z, x)$ et
$q \in \Mor_{\mathcal C}(x \times_y z, z)$ qui rendent le diagramme
$$
\xymatrix{
x \times_y z \ar[r]_q \ar[d]_p & z \ar[d]^g \\
x \ar[r]^f & y
}
$$
commutatif et qui satisfont la propriété universelle suivante : pour tout
$w\in \Ob(\mathcal{C})$ et tous morphismes
$\alpha \in \Mor_{\mathcal C}(w, x)$ et
$\beta \in \Mor_\mathcal{C}(w, z)$ tels que
$f \circ \alpha = g \circ \beta$,
il existe un unique
$\gamma \in \Mor_{\mathcal C}(w, x \times_y z)$ qui rend le diagramme
$$
\xymatrix{
w \ar[rrrd]^\beta \ar@{-->}[rrd]_\gamma \ar[rrdd]_\alpha & & \\
& & x \times_y z \ar[d]^p \ar[r]_q & z \ar[d]^g \\
& & x \ar[r]^f & y
}
$$
commutatif.
\end{definition}

\noindent
S'il existe, un produit fibré est unique à isomorphisme unique près.
Cela découle du lemme de Yoneda, car la définition exige que
$x \times_y z$ soit un objet de $\mathcal{C}$ tel que
$$
h_{x \times_y z}(w) = h_x(w) \times_{h_y(w)} h_z(w)
$$
fonctoriellement en $w$. Autrement dit, le produit fibré $x \times_y z$
est un objet représentant le foncteur
$w \mapsto h_x(w) \times_{h_y(w)} h_z(w)$.

\begin{definition}
\label{categories-definition-cartesian}
Nous disons qu'un diagramme commutatif
$$
\xymatrix{
w \ar[r] \ar[d] &
z \ar[d] \\
x \ar[r] &
y
}
$$
dans une catégorie est {\it cartésien} si $w$ et les morphismes
$w \to x$ et $w \to z$ forment un produit fibré des morphismes
$x \to y$ et $z \to y$.
\end{definition}

\begin{definition}
\label{categories-definition-has-fibre-products}
Nous disons que la catégorie $\mathcal{C}$ {\it admet les produits fibrés}
si le produit fibré existe pour tous
$f\in \Mor_{\mathcal C}(x, y)$ et
$g\in \Mor_{\mathcal C}(z, y)$.
\end{definition}

\begin{definition}
\label{categories-definition-representable-morphism}
Un morphisme $f : x \to y$ d'une catégorie $\mathcal{C}$ est dit
{\it représentable} si, pour tout morphisme $z \to y$
dans $\mathcal{C}$, le produit fibré $x \times_y z$ existe.
\end{definition}

\begin{lemma}
\label{categories-lemma-composition-representable}
Soit $\mathcal{C}$ une catégorie.
Soient $f : x \to y$ et $g : y \to z$ des morphismes représentables.
Alors $g \circ f : x \to z$ est représentable.
\end{lemma}

\begin{proof}
Soient $t \in \Ob(\mathcal C)$ et
$ \varphi \in \Mor_{\mathcal C}(t,z) $.
Puisque $g$ et $f$ sont représentables, nous obtenons des diagrammes
commutatifs
$$
\xymatrix{
y \times_z t \ar[r]_q \ar[d]_p &
t \ar[d]^{\varphi} \\
y \ar[r]^{g} & z
}
\quad\quad
\xymatrix{
x \times_y (y\! \times_z\! t) \ar[r]_{q'} \ar[d]_{p'} &
y \times_z t \ar[d]^p \\
x \ar[r]^f & y
}
$$
possédant la propriété universelle de la
définition \ref{categories-definition-fibre-products}.
Nous affirmons que $x \times_z t = x \times_y (y \times_z t)$,
muni des morphismes $q \circ q' : x \times_z t \to t$ et
$p' : x \times_z t \to x$, est un produit fibré.
La commutativité des diagrammes ci-dessus donne d'abord
$\varphi \circ q \circ q' = g \circ f \circ p'$.
Pour vérifier la propriété universelle, soient
$w \in \Ob(\mathcal C)$ et des morphismes $\alpha : w \to x$ et
$\beta : w \to t$ tels que
$\varphi \circ \beta = g \circ f \circ \alpha$.
Par définition du produit fibré, il existe des morphismes uniques
$\delta$ et $\gamma$ tels que
$$
\xymatrix{
w \ar[rrrd]^\beta \ar@{-->}[rrd]_\delta \ar[rrdd]_{f\circ\alpha} & & \\
& & y \times_z t \ar[d]_p \ar[r]_q & t \ar[d]^{\varphi} \\
& & y \ar[r]^{g} & z
}
$$
et
$$
\xymatrix{
w \ar[rrrd]^\delta \ar@{-->}[rrd]_\gamma \ar[rrdd]_{\alpha} & & \\
& & x \times_y (y\!\times_z\! t) \ar[d]_{p'} \ar[r]_{q'} &
y \times_z t \ar[d]^{p} \\
& & x \ar[r]^{f} & y
}
$$
soient commutatifs. Alors $\gamma$ rend commutatif le diagramme
$$
\xymatrix{
w \ar[rrrd]^\beta \ar@{-->}[rrd]_\gamma \ar[rrdd]_{\alpha} & & \\
& & x \times_z t \ar[d]_{p'} \ar[r]_{q\circ q'} & t \ar[d]^{\varphi} \\
& & x \ar[r]^{g\circ f} & z
}
$$
Pour montrer son unicité, soit $\gamma'$ tel que
$ q\circ q'\circ\gamma' =
\beta $ et
$ p'\circ \gamma' = \alpha $. Par l'unicité de $\gamma$, il suffit de
prouver que $ q'\circ\gamma' = \delta $ et
$ p'\circ\gamma' = \alpha $. La seconde égalité est supposée.
Pour la première, nous utilisons aussi l'unicité de $\delta$.
En effet, $\delta$ est l'unique morphisme vérifiant
$ q\circ\delta = \beta $ et
$ p\circ\delta = f\circ\alpha $. Nous avons déjà supposé
$ q\circ (q'\circ\gamma') = \beta $. De plus, par définition du produit
fibré, $ f\circ p' = p\circ q' $. Par conséquent :
$$
p\circ (q'\circ\gamma') =
(p\circ q')\circ\gamma' =
(f\circ p')\circ\gamma' =
f\circ (p'\circ\gamma') = f\circ\alpha.
$$
Ainsi $ q'\circ\gamma' = \delta $, ce qui achève la preuve.
\end{proof}

\begin{lemma}
\label{categories-lemma-base-change-representable}
Soit $\mathcal{C}$ une catégorie.
Soit $f : x \to y$ un morphisme représentable.
Soit $y' \to y$ un morphisme de $\mathcal{C}$.
Alors le morphisme $x' := x \times_y y' \to y'$ est lui aussi représentable.
\end{lemma}

\begin{proof}
Soit $z \to y'$ un morphisme. Le produit fibré
$x' \times_{y'} z$ doit représenter le foncteur
\begin{eqnarray*}
w & \mapsto & h_{x'}(w)\times_{h_{y'}(w)} h_z(w) \\
& = & (h_x(w) \times_{h_y(w)} h_{y'}(w)) \times_{h_{y'}(w)} h_z(w) \\
& = & h_x(w) \times_{h_y(w)} h_z(w)
\end{eqnarray*}
qui est représentable par hypothèse.
\end{proof}

\section{Exemples de produits fibrés}
\label{categories-section-example-fibre-products}

\noindent
Dans cette section, nous énumérons des exemples de produits fibrés
et les décrivons.

\medskip\noindent
Comme premier exemple tout à fait trivial, observons
que la catégorie des ensembles admet les produits fibrés et que tout
morphisme y est donc représentable. En effet, si $f : X \to Y$
et $g : Z \to Y$ sont des applications d'ensembles, nous définissons
$X \times_Y Z$ comme le sous-ensemble de $X \times Z$ formé
des couples $(x, z)$ tels que $f(x) = g(z)$. Les morphismes
$p : X \times_Y Z \to X$ et $q : X \times_Y Z \to Z$ sont
les projections $(x, z) \mapsto x$ et $(x, z) \mapsto z$.
Enfin, si $\alpha : W \to X$ et $\beta : W \to Z$
sont des morphismes tels que $f \circ \alpha = g \circ \beta$,
alors l'application $W \to X \times Z$, $w\mapsto (\alpha(w), \beta(w))$,
est manifestement à valeurs dans $X \times_Y Z$, comme voulu.

\medskip\noindent
Dans de nombreuses catégories dont les objets sont des ensembles munis
de certains types de structures algébriques, le produit fibré des ensembles
sous-jacents fournit aussi le produit fibré dans la catégorie. Supposons,
par exemple, que $X$, $Y$ et $Z$ ci-dessus soient des groupes et que
$f$, $g$ soient des homomorphismes de groupes. Le produit fibré ensembliste
$X \times_Y Z$ hérite alors d'une structure de groupe, simplement en
définissant le produit de deux couples par la formule
$(x, z) \cdot (x', z') = (xx', zz')$. Voici les catégories
pour lesquelles un raisonnement analogue s'applique.
\begin{enumerate}
\item La catégorie $\textit{Groups}$ des groupes.
\item La catégorie $G\textit{-Ensembles}$ des ensembles
munis d'une action à gauche de $G$, pour un groupe fixé $G$.
\item La catégorie des anneaux.
\item La catégorie des $R$-modules, pour un anneau $R$ donné.
\end{enumerate}




\section{Produits fibrés et représentabilité}
\label{categories-section-representable-map-presheaves}

\noindent
Dans cette section, nous déterminons les produits fibrés dans la
catégorie des foncteurs contravariants d'une catégorie
dans la catégorie des ensembles. Cette discussion sera remplacée plus loin
par celle des sites, des préfaisceaux et des faisceaux. La notion de
« morphisme représentable » entre de tels foncteurs présente un certain
intérêt.

\begin{lemma}
\label{categories-lemma-fibre-product-presheaves}
Soit $\mathcal{C}$ une catégorie.
Soient $F, G, H : \mathcal{C}^{opp} \to \textit{Ensembles}$
des foncteurs. Soient $a : F \to G$ et $b : H \to G$ des
transformations de foncteurs. Alors le produit fibré
$F \times_{a, G, b} H$ dans la catégorie
$\textit{PSh}(\mathcal{C})$
existe et est donné par la formule
$$
(F \times_{a, G, b} H)(X) =
F(X) \times_{a_X, G(X), b_X} H(X)
$$
pour tout objet $X$ de $\mathcal{C}$.
\end{lemma}

\begin{proof}
Omis.
\end{proof}

\noindent
Comme cas particulier, supposons donnés un morphisme
$a : F \to G$, un objet $U \in \Ob(\mathcal{C})$
et un élément $\xi \in G(U)$. D'après le
lemme de Yoneda \ref{categories-lemma-yoneda}, cela définit une transformation
$\xi : h_U \to G$. Dans ce cas, le produit fibré est décrit par la règle
$$
(h_U \times_{\xi, G, a} F)(X) =
\{ (f, \xi') \mid f : X \to U, \ \xi' \in F(X), \ G(f)(\xi) = a_X(\xi')\}
$$
Si $F$ et $G$ sont eux aussi représentables, ce foncteur représente le
produit fibré, lorsque celui-ci existe ; voir la
section \ref{categories-section-fibre-products}.
L'analogie avec la définition \ref{categories-definition-representable-morphism}
nous conduit à définir une notion de transformations représentables.

\begin{definition}
\label{categories-definition-representable-map-presheaves}
Soit $\mathcal{C}$ une catégorie.
Soient $F, G : \mathcal{C}^{opp} \to \textit{Ensembles}$
des foncteurs. Nous disons qu'un morphisme $a : F \to G$ est
{\it représentable}, ou que {\it $F$ est relativement représentable
au-dessus de $G$}, si, pour tout $U \in \Ob(\mathcal{C})$
et tout $\xi \in G(U)$, le foncteur
$h_U \times_{\xi, G, a} F$ est représentable.
\end{definition}

\begin{lemma}
\label{categories-lemma-representable-over-representable}
Soit $\mathcal{C}$ une catégorie.
Soit $a : F \to G$ un morphisme de foncteurs contravariants
de $\mathcal{C}$ dans $\textit{Ensembles}$. Si $a$ est représentable
et si $G$ est un foncteur représentable, alors $F$ est représentable.
\end{lemma}

\begin{proof}
Omis.
\end{proof}

\begin{lemma}
\label{categories-lemma-representable-diagonal}
Soit $\mathcal{C}$ une catégorie.
Soit $F : \mathcal{C}^{opp} \to \textit{Ensembles}$ un foncteur.
Supposons que $\mathcal{C}$ admette les produits de deux objets et les
produits fibrés. Les assertions suivantes sont équivalentes :
\begin{enumerate}
\item la diagonale $\Delta : F \to F \times F$ est représentable ;
\item pour tout $U$ dans $\mathcal{C}$
et tout $\xi \in F(U)$, le morphisme $\xi : h_U \to F$ est représentable ;
\item pour toute paire $U, V$ dans $\mathcal{C}$
et tous $\xi \in F(U)$, $\xi' \in F(V)$, le produit fibré
$h_U \times_{\xi, F, \xi'} h_V$ est représentable.
\end{enumerate}
\end{lemma}

\begin{proof}
Nous continuerons d'utiliser le lemme de Yoneda pour identifier $F(U)$
aux transformations de foncteurs $h_U \to F$.

\medskip\noindent
Équivalence de (2) et (3). Soient $U, \xi, V, \xi'$ comme dans (3).
Les assertions (2) et (3) disent toutes deux exactement que
$h_U \times_{\xi, F, \xi'} h_V$ est représentable ; la seule différence
est que l'énoncé (3) est symétrique en $U$ et $V$, tandis que (2) ne l'est pas.

\medskip\noindent
Supposons la condition (1). Soient $U, \xi, V, \xi'$
comme dans (3). Remarquons que
$h_U \times h_V = h_{U \times V}$ est représentable.
Notons $\eta : h_{U \times V} \to F \times F$ le morphisme
correspondant au produit $\xi \times \xi' : h_U \times h_V \to F \times F$.
Le produit fibré $F \times_{\Delta, F \times F, \eta} h_{U \times V}$
est alors représentable par hypothèse. Cela signifie qu'il existe
$W \in \Ob(\mathcal{C})$, des morphismes
$W \to U$, $W \to V$ et $h_W \to F$ tels que
$$
\xymatrix{
h_W \ar[d] \ar[r] & h_U \times h_V \ar[d]^{\xi \times \xi'} \\
F \ar[r] & F \times F
}
$$
soit cartésien. La description explicite des produits fibrés donnée dans le
lemme \ref{categories-lemma-fibre-product-presheaves} montre alors que
$h_W = h_U \times_{\xi, F, \xi'} h_V$, comme voulu.

\medskip\noindent
Supposons les conditions équivalentes (2) et (3). Soit $U$ un objet
de $\mathcal{C}$ et soit $(\xi, \xi') \in (F \times F)(U)$.
Par (3), le produit fibré $h_U \times_{\xi, F, \xi'} h_U$ est
représentable. Choisissons un objet $W$ et un isomorphisme
$h_W \to h_U \times_{\xi, F, \xi'} h_U$. Les deux projections
$\text{pr}_i : h_U \times_{\xi, F, \xi'} h_U \to h_U$
correspondent, par Yoneda, à des morphismes $p_i : W \to U$.
Considérons
$W' = W \times_{(p_1, p_2), U \times U} U$. Un argument formel montre
que $W'$ représente $F \times_{\Delta, F \times F} h_U$, car
$$
h_{W'} = h_W \times_{h_U \times h_U} h_U
= (h_U \times_{\xi, F, \xi'} h_U) \times_{h_U \times h_U} h_U
= F \times_{F \times F} h_U.
$$
Ainsi, $\Delta$ est représentable, ce qui achève la preuve.
\end{proof}








\section{Sommes amalgamées}
\label{categories-section-pushouts}

\noindent
La notion duale de celle de produit fibré est celle de somme amalgamée.

\begin{definition}
\label{categories-definition-pushouts}
Soient $x, y, z\in \Ob(\mathcal{C})$,
$f\in \Mor_\mathcal{C}(y, x)$
et $g\in \Mor_{\mathcal C}(y, z)$.
Une {\it somme amalgamée} de $f$ et $g$ est
un objet $x\amalg_y z\in \Ob(\mathcal{C})$
muni de morphismes
$p\in \Mor_{\mathcal C}(x, x\amalg_y z)$ et
$q\in\Mor_{\mathcal C}(z, x\amalg_y z)$ qui rendent le diagramme
$$
\xymatrix{
y \ar[r]_g \ar[d]_f & z \ar[d]^q \\
x \ar[r]^p & x\amalg_y z
}
$$
commutatif et qui satisfont la propriété universelle suivante :
pour tout $w\in \Ob(\mathcal{C})$ et tous morphismes
$\alpha \in \Mor_{\mathcal C}(x, w)$ et
$\beta \in \Mor_\mathcal{C}(z, w)$ tels que
$\alpha \circ f = \beta \circ g$, il existe un unique
$\gamma\in \Mor_{\mathcal C}(x\amalg_y z, w)$ qui rend le diagramme
$$
\xymatrix{
y \ar[r]_g \ar[d]_f & z \ar[d]^q \ar[rrdd]^\beta & & \\
x \ar[r]^p \ar[rrrd]^\alpha & x \amalg_y z \ar@{-->}[rrd]^\gamma & & \\
& & & w
}
$$
commutatif.
\end{definition}

\noindent
On peut démontrer directement, et sans difficulté, l'unicité du triplet
$(x\amalg_y z, p, q)$ à isomorphisme unique près (lorsqu'il existe).
On peut aussi considérer la somme amalgamée comme le
produit fibré dans la catégorie opposée ; son unicité découle alors
immédiatement de la discussion de la section \ref{categories-section-fibre-products}.

\begin{definition}
\label{categories-definition-cocartesian}
Nous disons qu'un diagramme commutatif
$$
\xymatrix{
y \ar[r] \ar[d] & z \ar[d] \\
x \ar[r] & w
}
$$
dans une catégorie est {\it cocartésien} si $w$ et les morphismes
$x \to w$ et $z \to w$ forment une somme amalgamée des morphismes
$y \to x$ et $y \to z$.
\end{definition}


\section{Égalisateurs}
\label{categories-section-equalizers}

\begin{definition}
\label{categories-definition-equalizers}
Supposons que $X$ et $Y$ soient des objets d'une catégorie $\mathcal{C}$
et que $a, b : X \to Y$ soient des morphismes. Un morphisme
$e : Z \to X$ est appelé un {\it égalisateur} de la paire $(a, b)$ si
$a \circ e = b \circ e$ et si $(Z, e)$ satisfait la propriété universelle
suivante : pour tout morphisme $t : W \to X$
dans $\mathcal{C}$ tel que $a \circ t = b \circ t$, il existe
un unique morphisme $s : W \to Z$ tel que $t = e \circ s$.
\end{definition}

\noindent
Comme pour les produits fibrés ci-dessus, lorsqu'ils existent, les
égalisateurs sont uniques à isomorphisme unique près. Cette définition
se généralise immédiatement au cas où l'on a plus de $2$ morphismes.

\section{Coégalisateurs}
\label{categories-section-coequalizers}

\begin{definition}
\label{categories-definition-coequalizers}
Supposons que $X$ et $Y$ soient des objets d'une catégorie $\mathcal{C}$
et que $a, b : X \to Y$ soient des morphismes. Un morphisme
$c : Y \to Z$ est appelé un {\it coégalisateur} de la paire $(a, b)$ si
$c \circ a = c \circ b$ et si $(Z, c)$ satisfait la propriété universelle
suivante : pour tout morphisme $t : Y \to W$
dans $\mathcal{C}$ tel que $t \circ a = t \circ b$, il existe
un unique morphisme $s : Z \to W$ tel que $t = s \circ c$.
\end{definition}

\noindent
Comme pour les sommes amalgamées ci-dessus, lorsqu'ils existent, les
coégalisateurs sont uniques à isomorphisme unique près ; cela résulte
de l'unicité des égalisateurs en considérant la catégorie opposée.
Cette définition se généralise immédiatement au cas où l'on a plus de
$2$ morphismes.

\section{Objets initiaux et finaux}
\label{categories-section-initial-final}

\begin{definition}
\label{categories-definition-initial-final}
Soit $\mathcal{C}$ une catégorie.
\begin{enumerate}
\item Un objet $x$ de la catégorie $\mathcal{C}$ est dit
{\it initial} si, pour tout objet $y$ de $\mathcal{C}$,
il existe exactement un morphisme $x \to y$.
\item Un objet $x$ de la catégorie $\mathcal{C}$ est dit
{\it final} si, pour tout objet $y$ de $\mathcal{C}$,
il existe exactement un morphisme $y \to x$.
\end{enumerate}
\end{definition}

\noindent
Dans la catégorie des ensembles, l'ensemble vide $\emptyset$ est un
objet initial, et c'est même le seul objet initial.
Tout {\it singleton}, c'est-à-dire tout ensemble à un élément,
est par ailleurs un objet final (il n'est donc pas unique).




\section{Monomorphismes et épimorphismes}
\label{categories-section-mono-epi}

\begin{definition}
\label{categories-definition-mono-epi}
Soit $\mathcal{C}$ une catégorie et soit $f : X \to Y$
un morphisme de $\mathcal{C}$.
\begin{enumerate}
\item Nous disons que $f$ est un {\it monomorphisme} si, pour tout objet
$W$ et toute paire de morphismes $a, b : W \to X$ tels que
$f \circ a = f \circ b$, on a $a = b$.
\item Nous disons que $f$ est un {\it épimorphisme} si, pour tout objet
$W$ et toute paire de morphismes $a, b : Y \to W$ tels que
$a \circ f = b \circ f$, on a $a = b$.
\end{enumerate}
\end{definition}

\begin{example}
\label{categories-example-mono-epi-sets}
Dans la catégorie des ensembles, les monomorphismes sont les applications
injectives et les épimorphismes sont les applications surjectives.
\end{example}

\begin{lemma}
\label{categories-lemma-characterize-mono-epi}
Soit $\mathcal{C}$ une catégorie et soit $f : X \to Y$
un morphisme de $\mathcal{C}$. Alors
\begin{enumerate}
\item $f$ est un monomorphisme si et seulement si la diagonale
$X \to X \times_Y X$ est un isomorphisme, et
\item $f$ est un épimorphisme si et seulement si la codiagonale
$Y \amalg_X Y \to Y$ est un isomorphisme.
\end{enumerate}
\end{lemma}

\begin{proof}
Supposons que $f$ soit un monomorphisme. Soit $W$ un objet de $\mathcal C$
et soient $\alpha, \beta \in \Mor_\mathcal C(W,X)$ tels que
$f\circ \alpha = f\circ
\beta$. On a donc $\alpha = \beta$,
puisque $f$ est un monomorphisme. De plus, on a le diagramme commutatif
$$
\xymatrix{
X \ar[r]^{\text{id}_X} \ar[d]_{\text{id}_X} & X \ar[d]^{f} \\
X \ar[r]^{f} & Y
}
$$
qui vérifie la propriété universelle avec $\gamma := \alpha = \beta$.
Ainsi, $X$ est bien le produit fibré $X\times_Y X$.

\medskip\noindent
Supposons que la diagonale $X \to X \times_Y X$ soit un isomorphisme.
Le diagramme
$$
\xymatrix{
X \ar[r]^{\text{id}_X} \ar[d]_{\text{id}_X} & X \ar[d]^{f} \\
X \ar[r]^{f} & Y }
$$
est commutatif et, si $W \in \Ob(\mathcal C)$ et
$\alpha, \beta : W \to X$ sont tels que
$f \circ \alpha = f \circ \beta$, il existe un unique
$\gamma$ vérifiant
$$
\gamma = \text{id}_X\circ\gamma = \alpha = \beta
$$
ce qui montre que $\alpha = \beta$.

\medskip\noindent
La preuve du second point est exactement la même, en utilisant la codiagonale
$Y\amalg_X Y \to Y$.
\end{proof}





\section{Limites et colimites}
\label{categories-section-limits}

\noindent
Soit $\mathcal{C}$ une catégorie. Un {\it diagramme} dans $\mathcal{C}$ est
simplement un foncteur $M : \mathcal{I} \to \mathcal{C}$. Nous disons que
$\mathcal{I}$ est la {\it catégorie d'indices}, ou que $M$ est un
diagramme indexé par $\mathcal{I}$. Nous noterons $M_i$ l'image de l'objet
$i$ de $\mathcal{I}$. Ainsi, pour un morphisme $\phi : i \to i'$
dans $\mathcal{I}$, on a $M(\phi) : M_i \to M_{i'}$.

\begin{definition}
\label{categories-definition-limit}
Une {\it limite} du diagramme indexé par $\mathcal{I}$ que définit $M$ dans la
catégorie
$\mathcal{C}$ est la donnée d'un objet $\lim_\mathcal{I} M$ de $\mathcal{C}$
et de morphismes $p_i : \lim_\mathcal{I} M \to M_i$ tels que
\begin{enumerate}
\item pour tout morphisme $\phi : i \to i'$
dans $\mathcal{I}$, on ait $p_{i'} = M(\phi) \circ p_i$, et
\item pour tout objet $W$ de $\mathcal{C}$ et toute famille de
morphismes $q_i : W \to M_i$ (indexée par $i \in \Ob(\mathcal{I})$)
telle que, pour tout $\phi : i \to i'$
dans $\mathcal{I}$, on ait $q_{i'} = M(\phi) \circ q_i$, il existe
un unique morphisme $q : W \to \lim_\mathcal{I} M$ tel que
$q_i = p_i \circ q$ pour tout objet $i$ de $\mathcal{I}$.
\end{enumerate}
\end{definition}

\noindent
Lorsqu'elles existent, les limites
$(\lim_\mathcal{I} M, (p_i)_{i\in \Ob(\mathcal{I})})$
sont uniques à isomorphisme unique près, par la condition d'unicité
de la définition. Les produits de deux objets, les produits fibrés et les
égalisateurs sont des exemples de limites. La limite du diagramme vide
est un objet final de $\mathcal{C}$.
Dans la catégorie des ensembles, toutes les limites existent.
La notion duale est celle de colimite.

\begin{definition}
\label{categories-definition-colimit}
Une {\it colimite} du diagramme indexé par $\mathcal{I}$ que définit $M$ dans la
catégorie
$\mathcal{C}$ est la donnée d'un objet $\colim_\mathcal{I} M$ de $\mathcal{C}$
et de morphismes $s_i : M_i \to \colim_\mathcal{I} M$ tels que
\begin{enumerate}
\item pour tout morphisme $\phi : i \to i'$
dans $\mathcal{I}$, on ait $s_i = s_{i'} \circ M(\phi)$, et
\item pour tout objet $W$ de $\mathcal{C}$ et toute famille de
morphismes $t_i : M_i \to W$ (indexée par $i \in \Ob(\mathcal{I})$)
telle que, pour tout $\phi : i \to i'$
dans $\mathcal{I}$, on ait $t_i = t_{i'} \circ M(\phi)$, il existe
un unique morphisme $t : \colim_\mathcal{I} M \to W$ tel que
$t_i = t \circ s_i$ pour tout objet $i$ de $\mathcal{I}$.
\end{enumerate}
\end{definition}

\noindent
Lorsqu'elles existent, les colimites
$(\colim_\mathcal{I} M, (s_i)_{i\in \Ob(\mathcal{I})})$
sont uniques à isomorphisme unique près, par la condition d'unicité
de la définition. Les coproduits de deux objets, les sommes amalgamées et
les coégalisateurs sont des exemples de colimites. La colimite du diagramme
vide est un objet initial de $\mathcal{C}$. Dans la catégorie des ensembles,
toutes les colimites existent.

\begin{remark}
\label{categories-remark-diagram-small}
La catégorie d'indices d'une (co)limite ne pourra jamais avoir une classe
propre d'objets. Dans ce projet, cela signifie qu'elle ne peut être l'une
des catégories énumérées dans la remarque \ref{categories-remark-big-categories}.
\end{remark}

\begin{remark}
\label{categories-remark-limit-colim}
Nous écrivons souvent $\lim_i M_i$, $\colim_i M_i$,
$\lim_{i\in \mathcal{I}} M_i$ ou $\colim_{i\in \mathcal{I}} M_i$
à la place des versions indexées par $\mathcal{I}$.
Avec cette notation et la description des limites et colimites d'ensembles
donnée dans la section \ref{categories-section-limit-sets} ci-dessous,
nous pouvons énoncer ce qui suit.
Soit $M : \mathcal{I} \to \mathcal{C}$ un diagramme.
\begin{enumerate}
\item S'il existe, l'objet $\lim_i M_i$ vérifie la propriété suivante :
$$
\Mor_\mathcal{C}(W, \lim_i M_i)
=
\lim_i \Mor_\mathcal{C}(W, M_i)
$$
où la limite du membre de droite est prise dans la catégorie des ensembles.
\item S'il existe, l'objet $\colim_i M_i$ vérifie la propriété suivante :
$$
\Mor_\mathcal{C}(\colim_i M_i, W)
=
\lim_{i \in \mathcal{I}^{opp}} \Mor_\mathcal{C}(M_i, W)
$$
où la limite du membre de droite, prise sur la catégorie opposée, est calculée
dans la catégorie des ensembles.
\end{enumerate}
Par le lemme de Yoneda (et son dual), cette formule détermine complètement
la limite, respectivement la colimite.
\end{remark}

\begin{remark}
\label{categories-remark-cones-and-cocones}
Soit $M : \mathcal{I} \to \mathcal{C}$ un diagramme. Dans ce contexte,
un {\it cône} sur $M$ est la donnée d'un objet $W$ et d'une famille de
morphismes $q_i : W \to M_i$, $i \in \Ob(\mathcal{I})$, tels que, pour tout
morphisme $\phi : i \to i'$ de $\mathcal{I}$, le diagramme
$$
\xymatrix{
& W \ar[dl]_{q_i} \ar[dr]^{q_{i'}} \\
M_i \ar[rr]^{M(\phi)} & & M_{i'}
}
$$
soit commutatif. Les cônes forment une catégorie dont la notion de morphisme
est évidente. Manifestement, la limite de $M$, si elle existe, est un objet
final de la catégorie des cônes. Dualement, un {\it cocône} sur $M$ est la
donnée d'un objet $W$ et d'une famille de morphismes $t_i : M_i \to W$ tels
que, pour tout morphisme $\phi : i \to i'$ dans $\mathcal{I}$, le diagramme
$$
\xymatrix{
M_i \ar[rr]^{M(\phi)} \ar[dr]_{t_i} & & M_{i'} \ar[dl]^{t_{i'}} \\
& W
}
$$
soit commutatif. Les cocônes forment une catégorie dont la notion de
morphisme est évidente. Comme ci-dessus, la colimite de $M$ existe si et
seulement si la catégorie des cocônes possède un objet initial.
\end{remark}

\noindent
Comme application des notions de limite et de colimite,
nous définissons les produits et les coproduits.

\begin{definition}
\label{categories-definition-product}
Supposons que $I$ soit un ensemble et que, pour tout $i \in I$, soit donné
un objet $M_i$ de la catégorie $\mathcal{C}$. Un {\it produit}
$\prod_{i\in I} M_i$ est, par définition, $\lim_\mathcal{I} M$
(s'il existe), où $\mathcal{I}$ est la catégorie dont les objets
sont les éléments de $I$ et dont les seuls morphismes sont les identités.
\end{definition}

\noindent
Un cas particulier important est $I = \emptyset$ ; le
produit est alors un objet final de la catégorie.
Les morphismes $p_i : \prod M_i \to M_i$ sont appelés les
{\it morphismes de projection}.

\begin{definition}
\label{categories-definition-coproduct}
Supposons que $I$ soit un ensemble et que, pour tout $i \in I$, soit donné
un objet $M_i$ de la catégorie $\mathcal{C}$. Un {\it coproduit}
$\coprod_{i\in I} M_i$ est, par définition, $\colim_\mathcal{I} M$
(s'il existe), où $\mathcal{I}$ est la catégorie dont les objets
sont les éléments de $I$ et dont les seuls morphismes sont les identités.
\end{definition}

\noindent
Un cas particulier important est $I = \emptyset$ ; le
coproduit est alors un objet initial de la catégorie.
Remarquons que le coproduit est muni de morphismes
$M_i \to \coprod M_i$. Ceux-ci sont parfois appelés les
{\it coprojections}.

\begin{lemma}
\label{categories-lemma-functorial-colimit}
Supposons que $M : \mathcal{I} \to \mathcal{C}$
et $N : \mathcal{J} \to \mathcal{C}$ soient des diagrammes
dont les colimites existent. Supposons que
$H : \mathcal{I} \to \mathcal{J}$ soit
un foncteur et que $t : M \to N \circ H$
soit une transformation de foncteurs.
Il existe alors un unique morphisme
$$
\theta :
\colim_\mathcal{I} M
\longrightarrow
\colim_\mathcal{J} N
$$
tel que tous les diagrammes
$$
\xymatrix{
M_i \ar[d]_{t_i} \ar[r]
&
\colim_\mathcal{I} M \ar[d]^{\theta}
\\
N_{H(i)} \ar[r]
&
\colim_\mathcal{J} N
}
$$
soient commutatifs.
\end{lemma}

\begin{proof}
Omis.
\end{proof}

\begin{lemma}
\label{categories-lemma-functorial-limit}
Supposons que $M : \mathcal{I} \to \mathcal{C}$
et $N : \mathcal{J} \to \mathcal{C}$ soient des diagrammes
dont les limites existent. Supposons que
$H : \mathcal{I} \to \mathcal{J}$ soit un foncteur et que
$t : N \circ H \to M$ soit une transformation de foncteurs.
Il existe alors un unique morphisme
$$
\theta :
\lim_\mathcal{J} N
\longrightarrow
\lim_\mathcal{I} M
$$
tel que tous les diagrammes
$$
\xymatrix{
\lim_\mathcal{J} N \ar[d]^{\theta} \ar[r]
&
N_{H(i)} \ar[d]_{t_i}
\\
\lim_\mathcal{I} M \ar[r]
&
M_i
}
$$
soient commutatifs.
\end{lemma}

\begin{proof}
Omis.
\end{proof}


\begin{lemma}
\label{categories-lemma-colimits-commute}
Soient $\mathcal{I}$ et $\mathcal{J}$ des catégories d'indices. Soit
$M : \mathcal{I} \times \mathcal{J} \to \mathcal{C}$ un foncteur.
Supposons que $M_{i, \infty} = \colim_j M_{i,j}$
existe pour tout $i$. Alors le foncteur ainsi obtenu
$M_{-, \infty} : \mathcal{I} \to \mathcal{C}$
admet une colimite si et seulement si $M$ en admet une, et les deux
colimites coïncident. En particulier,
$$
\colim_i \colim_j M_{i, j}
=
\colim_{i, j} M_{i, j}
=
\colim_j \colim_i M_{i, j}
$$
pourvu que toutes les colimites indiquées existent. Il en va de même
pour les limites.
\end{lemma}

\begin{proof}
Omis.
\end{proof}

\begin{lemma}
\label{categories-lemma-limits-products-equalizers}
Soit $M : \mathcal{I} \to \mathcal{C}$ un diagramme.
Écrivons $I = \Ob(\mathcal{I})$ et $A = \text{Flèches}(\mathcal{I})$.
Notons $s, t : A \to I$ les applications source et but.
Supposons que $\prod_{i \in I} M_i$ et $\prod_{a \in A} M_{t(a)}$
existent. Supposons que l'égalisateur de
$$
\xymatrix{
\prod_{i \in I} M_i
\ar@<1ex>[r]^\phi \ar@<-1ex>[r]_\psi
&
\prod_{a \in A} M_{t(a)}
}
$$
existe, les morphismes étant déterminés par leurs composantes
comme suit : $p_a \circ \psi = M(a) \circ p_{s(a)}$
et $p_a \circ \phi = p_{t(a)}$. Cet égalisateur est alors la
limite du diagramme.
\end{lemma}

\begin{proof}
Omis.
\end{proof}


\begin{lemma}
\label{categories-lemma-colimits-coproducts-coequalizers}
\begin{slogan}
Si tous les coproduits et tous les coégalisateurs existent, toutes les
colimites existent.
\end{slogan}
Soit $M : \mathcal{I} \to \mathcal{C}$ un diagramme.
Écrivons $I = \Ob(\mathcal{I})$ et $A = \text{Flèches}(\mathcal{I})$.
Notons $s, t : A \to I$ les applications source et but.
Supposons que $\coprod_{i \in I} M_i$ et
$\coprod_{a \in A} M_{s(a)}$ existent. Supposons que le coégalisateur de
$$
\xymatrix{
\coprod_{a \in A} M_{s(a)}
\ar@<1ex>[r]^\phi \ar@<-1ex>[r]_\psi
&
\coprod_{i \in I} M_i
}
$$
existe, les morphismes étant déterminés par leurs composantes
comme suit : la composante $M_{s(a)}$ s'envoie par $\psi$
dans la composante $M_{t(a)}$ au moyen du morphisme $M(a)$.
La composante $M_{s(a)}$ s'envoie par $\phi$ dans la composante
$M_{s(a)}$ au moyen du morphisme identité. Ce coégalisateur est alors la
colimite du diagramme.
\end{lemma}

\begin{proof}
Omis.
\end{proof}








\section{Limites et colimites dans la catégorie des ensembles}
\label{categories-section-limit-sets}

\noindent
Non seulement les limites et les colimites existent dans $\textit{Ensembles}$,
mais elles sont également faciles à décrire. Soit en effet
$M : \mathcal{I}
\to \textit{Ensembles}$, $i \mapsto M_i$, un diagramme
d'ensembles. Notons $I = \Ob(\mathcal{I})$.
La limite s'écrit
$$
\lim_\mathcal{I} M
=
\{
(m_i)_{i\in I} \in \prod\nolimits_{i\in I} M_i
\mid
\forall \phi : i \to i' \text{ dans }\mathcal{I},
M(\phi)(m_i) = m_{i'}
\}.
$$
Nous considérons donc un élément de la limite comme un système compatible
d'éléments de tous les ensembles $M_i$.

\medskip\noindent
D'autre part, la colimite est
$$
\colim_\mathcal{I} M
=
(\coprod\nolimits_{i\in I} M_i)/\sim
$$
où la relation d'équivalence $\sim$ est engendrée en imposant
$m_i \sim m_{i'}$ si $m_i \in M_i$,
$m_{i'} \in M_{i'}$ et $M(\phi)(m_i) = m_{i'}$ pour un certain
$\phi : i \to i'$. Autrement dit, $m_i \in M_i$
et $m_{i'} \in M_{i'}$ sont équivalents s'il existe une
chaîne de morphismes dans $\mathcal{I}$
$$
\xymatrix{
&
i_1 \ar[ld] \ar[rd] & &
i_3 \ar[ld] & &
i_{2n-1} \ar[rd] & \\
i = i_0 & &
i_2 & &
\ldots & &
i_{2n} = i'
}
$$
et des éléments $m_{i_j} \in M_{i_j}$ dont les images coïncident par
les applications $M_{i_{2k-1}} \to M_{i_{2k-2}}$ et
$M_{i_{2k-1}}
\to M_{i_{2k}}$ induites par les flèches de
$\mathcal{I}$ ci-dessus.

\medskip\noindent
Ce type d'objet n'est pas très commode à manipuler.
Mais si le diagramme est filtrant, sa description est beaucoup plus simple.
Nous l'expliquerons dans la section \ref{categories-section-directed-colimits}.



\section{Limites connexes}
\label{categories-section-connected-limits}

\noindent
Une (co)limite est dite connexe si sa catégorie d'indices est connexe.

\begin{definition}
\label{categories-definition-category-connected}
Nous disons qu'une catégorie $\mathcal{I}$ est {\it connexe}
si la relation d'équivalence engendrée par
$x \sim y \Leftrightarrow \Mor_\mathcal{I}(x, y) \not = \emptyset$
possède exactement une classe d'équivalence.
\end{definition}

\noindent
Nous suivons ici la convention de
Topologie, définition \ref{topology-definition-connected-components},
selon laquelle les espaces connexes sont non vides.
Le résultat suivant caractérise, en un sens quelque peu vague,
les limites connexes.

\begin{lemma}
\label{categories-lemma-connected-limit-over-X}
Soit $\mathcal{C}$ une catégorie.
Soit $X$ un objet de $\mathcal{C}$.
Soit $M : \mathcal{I} \to \mathcal{C}/X$ un diagramme
dans la catégorie des objets au-dessus de $X$.
Si la catégorie d'indices $\mathcal{I}$ est connexe
et si la limite de $M$ existe dans $\mathcal{C}/X$,
alors la limite du composé
$\mathcal{I} \to \mathcal{C}/X \to \mathcal{C}$
existe et lui est égale.
\end{lemma}

\begin{proof}
Soit $L \to X$ un objet représentant la limite dans $\mathcal{C}/X$.
Considérons le foncteur
$$
W \longmapsto \lim_i \Mor_\mathcal{C}(W, M_i).
$$
Soit $(\varphi_i)$ un élément de l'ensemble de droite.
Comme chaque $M_i$ est muni d'un morphisme $s_i : M_i \to X$, nous
obtenons des morphismes $f_i = s_i \circ \varphi_i : W \to X$.
La connexité de $\mathcal{I}$ montre que tous les $f_i$ sont égaux.
Puisque $\mathcal{I}$ est non vide, il existe au moins un $f_i$.
Leur valeur commune $W \to X$ confère donc à $W$ une structure
d'objet de $\mathcal{C}/X$, et $(\varphi_i)$ définit un
élément de $\lim_i \Mor_{\mathcal{C}/X}(W, M_i)$.
Nous obtenons ainsi un unique morphisme $\phi : W \to L$ tel que
$\varphi_i$ soit le composé de $\phi$ et de $L \to M_i$, comme voulu.
\end{proof}

\begin{lemma}
\label{categories-lemma-connected-colimit-under-X}
Soit $\mathcal{C}$ une catégorie.
Soit $X$ un objet de $\mathcal{C}$.
Soit $M : \mathcal{I} \to X/\mathcal{C}$ un diagramme
dans la catégorie des objets au-dessous de $X$.
Si la catégorie d'indices $\mathcal{I}$ est connexe
et si la colimite de $M$ existe dans $X/\mathcal{C}$,
alors la colimite du composé
$\mathcal{I} \to X/\mathcal{C} \to \mathcal{C}$
existe et lui est égale.
\end{lemma}

\begin{proof}
Omis. Indication : ce lemme est dual du
lemme \ref{categories-lemma-connected-limit-over-X}.
\end{proof}




\section{Catégories cofinales et initiales}
\label{categories-section-cofinal}

\noindent
Dans la littérature, le mot « final » est parfois employé à la place de
« cofinal » dans la définition suivante.

\begin{definition}
\label{categories-definition-cofinal}
Soit $H : \mathcal{I} \to \mathcal{J}$ un foncteur entre catégories.
Nous disons que {\it $\mathcal{I}$ est cofinale dans $\mathcal{J}$},
ou que $H$ est {\it cofinal}, si
\begin{enumerate}
\item pour tout $y \in \Ob(\mathcal{J})$, il existe
$x \in \Ob(\mathcal{I})$ et un morphisme $y \to H(x)$, et
\item étant donnés $y \in \Ob(\mathcal{J})$,
$x, x' \in \Ob(\mathcal{I})$ et des morphismes
$y \to H(x)$ et $y \to H(x')$, il existe une suite de morphismes
$$
x = x_0 \leftarrow x_1 \rightarrow x_2 \leftarrow x_3 \rightarrow \ldots
\rightarrow x_{2n} = x'
$$
dans $\mathcal{I}$ et des morphismes $y \to H(x_i)$ dans $\mathcal{J}$,
les morphismes $y \to H(x_0)$ et $y \to H(x_{2n})$ étant ceux donnés,
tels que les diagrammes
$$
\xymatrix{
& y \ar[ld] \ar[d] \ar[rd] \\
H(x_{2k}) & H(x_{2k + 1}) \ar[l] \ar[r] & H(x_{2k + 2})
}
$$
soient commutatifs pour $k = 0, \ldots, n - 1$.
\end{enumerate}
Autrement dit, fixons un objet $y$ de $\mathcal{J}$ et considérons
l'ensemble $S$ des couples $(x, b)$, où $x$ est un objet de $\mathcal{I}$
et $b : y \to H(x)$ un morphisme. Considérons sur $S$ la relation
d'équivalence engendrée par $(x, b) \sim (x', b')$ s'il existe un
morphisme $a : x \to x'$ tel que $b' = H(a) \circ b$. Alors $S$ doit
être formé d'une seule classe d'équivalence.
\end{definition}

\begin{lemma}
\label{categories-lemma-cofinal}
Soit $H : \mathcal{I} \to \mathcal{J}$ un foncteur entre catégories.
Supposons $\mathcal{I}$ cofinale dans $\mathcal{J}$. Alors, pour tout
diagramme $M : \mathcal{J} \to \mathcal{C}$, on a un isomorphisme canonique
$$
\colim_\mathcal{I} M \circ H
=
\colim_\mathcal{J} M
$$
dès que l'un des deux membres existe.
\end{lemma}

\begin{proof}
Omis.
\end{proof}

\begin{definition}
\label{categories-definition-initial}
Soit $H : \mathcal{I} \to \mathcal{J}$ un foncteur entre catégories.
Nous disons que {\it $\mathcal{I}$ est initiale dans $\mathcal{J}$},
ou que $H$ est {\it initial}, si
\begin{enumerate}
\item pour tout $y \in \Ob(\mathcal{J})$, il existe
$x \in \Ob(\mathcal{I})$ et un morphisme $H(x) \to y$ ;
\item pour tout $y \in \Ob(\mathcal{J})$, tous
$x , x' \in \Ob(\mathcal{I})$ et tous morphismes
$H(x) \to y$, $H(x') \to y$ dans $\mathcal{J}$,
il existe une suite de morphismes
$$
x = x_0 \leftarrow x_1 \rightarrow x_2 \leftarrow x_3 \rightarrow \ldots
\rightarrow x_{2n} = x'
$$
dans $\mathcal{I}$ et des morphismes $H(x_i) \to y$ dans $\mathcal{J}$
tels que les diagrammes
$$
\xymatrix{
H(x_{2k}) \ar[rd] &
H(x_{2k + 1}) \ar[l] \ar[r] \ar[d] &
H(x_{2k + 2}) \ar[ld] \\
& y
}
$$
soient commutatifs pour $k = 0, \ldots, n - 1$.
\end{enumerate}
\end{definition}

\noindent
C'est simplement la notion duale de celle de foncteur « cofinal ».

\begin{lemma}
\label{categories-lemma-initial}
Soit $H : \mathcal{I} \to \mathcal{J}$ un foncteur de catégories.
Supposons $\mathcal{I}$ initiale dans $\mathcal{J}$.
Alors, pour tout diagramme $M : \mathcal{J} \to \mathcal{C}$,
on a un isomorphisme canonique
$$
\lim_\mathcal{I} M \circ H = \lim_\mathcal{J} M
$$
dès que l'un des deux membres existe.
\end{lemma}

\begin{proof}
Omis.
\end{proof}

\begin{lemma}
\label{categories-lemma-colimit-constant-connected-fibers}
Soit $F : \mathcal{I} \to \mathcal{I}'$ un foncteur.
Supposons que
\begin{enumerate}
\item les catégories fibres (voir la
définition \ref{categories-definition-fibre-category})
de $\mathcal{I}$ au-dessus de $\mathcal{I}'$ soient toutes connexes, et que
\item pour tout morphisme $\alpha' : x' \to y'$ dans $\mathcal{I}'$, il
existe un morphisme $\alpha : x \to y$ dans $\mathcal{I}$ tel que
$F(\alpha) = \alpha'$.
\end{enumerate}
Alors, pour tout diagramme $M : \mathcal{I}' \to \mathcal{C}$,
la colimite $\colim_\mathcal{I} M \circ F$ existe si et seulement si
$\colim_{\mathcal{I}'} M$ existe ; lorsqu'elles existent, ces colimites
coïncident.
\end{lemma}

\begin{proof}
On peut le démontrer en établissant que $\mathcal{I}$ est cofinale dans
$\mathcal{I}'$, puis en appliquant le lemme \ref{categories-lemma-cofinal}.
Mais on peut aussi donner la preuve directe suivante.
Il suffit de montrer que, pour tout objet $T$ de $\mathcal{C}$, on a
$$
\lim_{\mathcal{I}^{opp}} \Mor_\mathcal{C}(M_{F(i)}, T)
=
\lim_{(\mathcal{I}')^{opp}} \Mor_\mathcal{C}(M_{i'}, T)
$$
Si $(g_{i'})_{i' \in \Ob(\mathcal{I}')}$ est un élément du
membre de droite, alors, en posant $f_i = g_{F(i)}$, on obtient un
élément $(f_i)_{i \in \Ob(\mathcal{I})}$ du membre de gauche.
Réciproquement, soit $(f_i)_{i \in \Ob(\mathcal{I})}$ un élément du
membre de gauche. Sur chaque catégorie fibre (connexe)
$\mathcal{I}_{i'}$, le foncteur $M \circ F$
est constant de valeur $M_{i'}$. Les morphismes
$f_i$, pour $i \in \Ob(\mathcal{I})$ tel que $F(i) = i'$,
sont donc tous égaux et déterminent un morphisme bien défini
$g_{i'} : M_{i'} \to T$. Par l'hypothèse (2), la famille
$(g_{i'})_{i' \in \Ob(\mathcal{I}')}$ définit un élément
du membre de droite.
\end{proof}

\begin{lemma}
\label{categories-lemma-product-with-connected}
Soient $\mathcal{I}$ et $\mathcal{J}$ des catégories et notons
$p : \mathcal{I} \times \mathcal{J} \to \mathcal{J}$ la projection.
Si $\mathcal{I}$ est connexe, alors, pour un diagramme
$M : \mathcal{J} \to \mathcal{C}$, la colimite
$\colim_\mathcal{J} M$ existe si et seulement si
$\colim_{\mathcal{I} \times \mathcal{J}} M \circ p$ existe et,
dans ce cas, ces colimites sont égales.
\end{lemma}

\begin{proof}
C'est un cas particulier du
lemme \ref{categories-lemma-colimit-constant-connected-fibers}.
\end{proof}





\section{Limites et colimites finies}
\label{categories-section-finite-limits}

\noindent
Une (co)limite est dite {\it finie} si sa catégorie d'indices est finie,
c'est-à-dire si celle-ci possède un nombre fini d'objets et un nombre fini
de morphismes. Une (co)limite est dite {\it non vide} si sa catégorie
d'indices est non vide. Une (co)limite est dite {\it connexe} si sa catégorie
d'indices est connexe ; voir la
définition \ref{categories-definition-category-connected}.
Il existe en fait « suffisamment » de catégories d'indices finies.

\begin{lemma}
\label{categories-lemma-finite-diagram-category}
Soit $\mathcal{I}$ une catégorie telle que
\begin{enumerate}
\item $\Ob(\mathcal{I})$ soit fini, et
\item il existe un nombre fini de morphismes
$f_1, \ldots, f_m \in \text{Flèches}(\mathcal{I})$ tels
que tout morphisme de $\mathcal{I}$ soit un composé
$f_{j_1} \circ f_{j_2} \circ \ldots \circ f_{j_k}$.
\end{enumerate}
Il existe alors un foncteur $F : \mathcal{J} \to \mathcal{I}$ tel que
\begin{enumerate}
\item[(a)] $\mathcal{J}$ soit une catégorie finie, et
\item[(b)] pour tout diagramme $M : \mathcal{I} \to \mathcal{C}$, la
(co)limite de $M$ sur $\mathcal{I}$ existe si et seulement si
la (co)limite de $M \circ F$ sur $\mathcal{J}$ existe ; dans ce cas,
les (co)limites sont canoniquement isomorphes.
\end{enumerate}
De plus, $\mathcal{J}$ est connexe (resp.\ non vide) si et seulement si
$\mathcal{I}$ l'est.
\end{lemma}

\begin{proof}
Écrivons $\Ob(\mathcal{I}) = \{x_1, \ldots, x_n\}$.
Notons $s, t : \{1, \ldots, m\} \to \{1, \ldots, n\}$ les applications
telles que $f_j : x_{s(j)} \to x_{t(j)}$.
Posons $\Ob(\mathcal{J}) = \{y_1, \ldots, y_n, z_1, \ldots, z_n\}$.
Outre les identités, introduisons des morphismes
$g_j : y_{s(j)} \to z_{t(j)}$, $j = 1, \ldots, m$, et des morphismes
$h_i : y_i \to z_i$, $i = 1, \ldots, n$. Comme tous les morphismes de
$\mathcal{J}$ autres que les identités vont d'un $y$ vers un $z$, il n'y a
aucune composition à définir ni aucune associativité à vérifier.
Posons $F(y_i) = F(z_i) = x_i$. Posons
$F(g_j) = f_j$ et $F(h_i) = \text{id}_{x_i}$.
Il est clair que $F$ est un foncteur.
Il est clair que $\mathcal{J}$ est finie.
Il est clair que $\mathcal{J}$ est connexe, resp.\ non vide,
si et seulement si $\mathcal{I}$ l'est.

\medskip\noindent
Soit $M : \mathcal{I} \to \mathcal{C}$ un diagramme.
Considérons un objet $W$ de $\mathcal{C}$ et des morphismes
$q_i : W \to M(x_i)$ comme dans la
définition \ref{categories-definition-limit}.
En prenant $q_i : W \to M(F(y_i)) = M(F(z_i)) = M(x_i)$, nous obtenons
une famille de morphismes comme dans la
définition \ref{categories-definition-limit}
pour le diagramme $M \circ F$.
Réciproquement, supposons donnés des morphismes
$q_{y_i} : W \to M(F(y_i))$ et $q_{z_i} : W \to M(F(z_i))$
comme dans la définition \ref{categories-definition-limit}
pour le diagramme $M \circ F$. Puisque
$$
M(F(h_i)) = \text{id} : M(F(y_i)) = M(x_i) \longrightarrow M(x_i) = M(F(z_i))
$$
on en déduit que $q_{y_i} = q_{z_i}$ pour tout $i$. Posons $q_i$ égal à
cette valeur commune. La compatibilité de
$q_{s(j)} = q_{y_{s(j)}}$ et de $q_{t(j)} = q_{z_{t(j)}}$ avec le morphisme
$M(f_j)$ garantit que la famille $q_i$ est compatible avec tous les
morphismes de $\mathcal{I}$, puisque, par hypothèse, chacun d'eux est un
composé de morphismes $f_j$. Nous avons ainsi obtenu une bijection canonique
$$
\lim_{B \in \Ob(\mathcal{J})} \Mor_\mathcal{C}(W, M(F(B)))
=
\lim_{A \in \Ob(\mathcal{I})} \Mor_\mathcal{C}(W, M(A))
$$
qui entraîne l'assertion du lemme concernant les limites. L'assertion
concernant les colimites se démontre de la même manière (preuve omise).
\end{proof}

\begin{lemma}
\label{categories-lemma-fibre-products-equalizers-exist}
Soit $\mathcal{C}$ une catégorie.
Les assertions suivantes sont équivalentes :
\begin{enumerate}
\item Les limites finies connexes existent dans $\mathcal{C}$.
\item Les égalisateurs et les produits fibrés existent dans $\mathcal{C}$.
\end{enumerate}
\end{lemma}

\begin{proof}
Puisque les égalisateurs et les produits fibrés sont des limites finies
connexes, (1) implique (2). Réciproquement, soit $\mathcal{I}$
une catégorie d'indices finie et connexe. Soit
$F : \mathcal{J} \to \mathcal{I}$
le foncteur de catégories d'indices construit dans la preuve du
lemme \ref{categories-lemma-finite-diagram-category}.
Nous pouvons alors remplacer $\mathcal{I}$ par $\mathcal{J}$.
Nous pouvons donc supposer que
$\Ob(\mathcal{I}) = \{x_1, \ldots, x_n\} \amalg \{y_1, \ldots, y_m\}$,
avec $n, m \geq 1$, et que tous les morphismes de $\mathcal{I}$ autres que les
identités sont des morphismes $f : x_i \to y_j$ pour certains
$i$ et $j$.

\medskip\noindent
Supposons $n > 1$. Puisque $\mathcal{I}$ est connexe, il existe
des indices $i_1, i_2$ et $j_0$, ainsi que des morphismes
$a : x_{i_1} \to y_{j_0}$ et $b : x_{i_2} \to y_{j_0}$.
Considérons la catégorie
$$
\mathcal{I}' =
\{x\} \amalg \{x_1, \ldots, \hat x_{i_1}, \ldots, \hat x_{i_2}, \ldots, x_n\}
\amalg \{y_1, \ldots, y_m\}
$$
avec
$$
\Mor_{\mathcal{I}'}(x, y_j) = \Mor_\mathcal{I}(x_{i_1}, y_j)
\amalg \Mor_\mathcal{I}(x_{i_2}, y_j)
$$
et tous les autres ensembles de morphismes égaux à ceux de
$\mathcal{I}$. Pour tout foncteur
$M : \mathcal{I} \to \mathcal{C}$, nous pouvons construire un foncteur
$M' : \mathcal{I}' \to \mathcal{C}$ en posant
$$
M'(x) = M(x_{i_1}) \times_{M(a), M(y_{j_0}), M(b)} M(x_{i_2})
$$
et, pour un morphisme $f' : x \to y_j$ correspondant, disons, à
$f : x_{i_1} \to y_j$, en posant $M'(f') = M(f) \circ \text{pr}_1$.
Le foncteur $M$ admet alors une limite si et seulement si le foncteur
$M'$ en admet une (preuve omise). Par récurrence, on se ramène donc
au cas $n = 1$.

\medskip\noindent
Si $n = 1$, la limite de tout $M : \mathcal{I} \to \mathcal{C}$ est
l'égalisateur successif des paires de morphismes $x_1 \to y_j$ ;
elle existe donc par hypothèse.
\end{proof}

\begin{lemma}
\label{categories-lemma-almost-finite-limits-exist}
Soit $\mathcal{C}$ une catégorie.
Les assertions suivantes sont équivalentes :
\begin{enumerate}
\item Les limites finies non vides existent dans $\mathcal{C}$.
\item Les produits de deux objets et les égalisateurs existent dans
$\mathcal{C}$.
\item Les produits de deux objets et les produits fibrés existent dans
$\mathcal{C}$.
\end{enumerate}
\end{lemma}

\begin{proof}
Puisque les produits de deux objets, les produits fibrés et les égalisateurs
sont des limites dont les catégories d'indices sont non vides, (1) implique
(2) et (3). Supposons (2). Les produits finis non vides et les égalisateurs
existent alors. Le lemme \ref{categories-lemma-limits-products-equalizers}
montre donc que les limites finies non vides existent, c'est-à-dire que
(1) est vraie. Supposons (3).
Si $a, b : A \to B$ sont des morphismes de $\mathcal{C}$, l'égalisateur
de $a, b$ est
$$
(A \times_{a, B, b} A)\times_{(\text{pr}_1, \text{pr}_2), A \times A, \Delta} A.
$$
Ainsi, (3) implique (2), ce qui démontre le lemme.
\end{proof}

\begin{lemma}
\label{categories-lemma-finite-limits-exist}
Soit $\mathcal{C}$ une catégorie.
Les assertions suivantes sont équivalentes :
\begin{enumerate}
\item Les limites finies existent dans $\mathcal{C}$.
\item Les produits finis et les égalisateurs existent.
\item La catégorie possède un objet final et les produits fibrés existent.
\end{enumerate}
\end{lemma}

\begin{proof}
Puisque les produits finis, les produits fibrés, les égalisateurs et les
objets finaux sont des limites sur des catégories d'indices finies, (1)
implique (2) et (3). Le lemme \ref{categories-lemma-limits-products-equalizers}
ci-dessus montre que (2) implique (1). Supposons (3).
Le produit $A \times B$ est le produit fibré au-dessus de l'objet final.
Si $a, b : A \to B$ sont des morphismes de $\mathcal{C}$, l'égalisateur
de $a, b$ est
$$
(A \times_{a, B, b} A)\times_{(\text{pr}_1, \text{pr}_2), A \times A, \Delta} A.
$$
Ainsi, (3) implique (2), ce qui démontre le lemme.
\end{proof}

\begin{lemma}
\label{categories-lemma-push-outs-coequalizers-exist}
Soit $\mathcal{C}$ une catégorie.
Les assertions suivantes sont équivalentes :
\begin{enumerate}
\item Les colimites finies connexes existent dans $\mathcal{C}$.
\item Les coégalisateurs et les sommes amalgamées existent dans $\mathcal{C}$.
\end{enumerate}
\end{lemma}

\begin{proof}
Omis. Indication : c'est l'énoncé dual du
lemme \ref{categories-lemma-fibre-products-equalizers-exist}.
\end{proof}

\begin{lemma}
\label{categories-lemma-almost-finite-colimits-exist}
Soit $\mathcal{C}$ une catégorie.
Les assertions suivantes sont équivalentes :
\begin{enumerate}
\item Les colimites finies non vides existent dans $\mathcal{C}$.
\item Les coproduits de deux objets et les coégalisateurs existent dans
$\mathcal{C}$.
\item Les coproduits de deux objets et les sommes amalgamées existent dans
$\mathcal{C}$.
\end{enumerate}
\end{lemma}

\begin{proof}
Omis. Indication : c'est l'énoncé dual du
lemme \ref{categories-lemma-almost-finite-limits-exist}.
\end{proof}

\begin{lemma}
\label{categories-lemma-colimits-exist}
Soit $\mathcal{C}$ une catégorie.
Les assertions suivantes sont équivalentes :
\begin{enumerate}
\item Les colimites finies existent dans $\mathcal{C}$.
\item Les coproduits finis et les coégalisateurs existent dans $\mathcal{C}$.
\item La catégorie possède un objet initial et les sommes amalgamées existent.
\end{enumerate}
\end{lemma}

\begin{proof}
Omis. Indication : c'est l'énoncé dual du
lemme \ref{categories-lemma-finite-limits-exist}.
\end{proof}





\section{Limites inductives filtrantes}
\label{categories-section-directed-colimits}

\noindent
Les colimites sont plus faciles à calculer ou à décrire lorsqu'elles
portent sur un diagramme filtrant. En voici la définition.

\begin{definition}
\label{categories-definition-directed}
Nous disons qu'un diagramme $M : \mathcal{I} \to \mathcal{C}$ est
{\it dirigé}, ou {\it filtrant}, si les conditions suivantes sont satisfaites :
\begin{enumerate}
\item la catégorie $\mathcal{I}$ possède au moins un objet ;
\item pour toute paire d'objets $x, y$ de $\mathcal{I}$,
il existe un objet $z$ et des morphismes $x \to z$,
$y \to z$ ; et
\item pour toute paire d'objets $x, y$ de $\mathcal{I}$
et toute paire de morphismes $a, b : x \to y$ de $\mathcal{I}$,
il existe un morphisme $c : y \to z$ de $\mathcal{I}$
tel que $M(c \circ a) = M(c \circ b)$ comme morphismes dans $\mathcal{C}$.
\end{enumerate}
Nous disons qu'une catégorie d'indices $\mathcal{I}$ est
{\it dirigée}, ou {\it filtrante}, si
$\text{id} : \mathcal{I} \to \mathcal{I}$ est filtrant
(autrement dit, on supprime le $M$ dans la partie (3) ci-dessus).
\end{definition}

\noindent
Tout diagramme dont la catégorie d'indices est filtrante est filtrant ;
c'est ainsi que les limites inductives filtrantes se présentent habituellement.
En fait, si $M : \mathcal{I} \to \mathcal{C}$ est un diagramme filtrant,
on peut factoriser $M$ sous la forme
$\mathcal{I} \to \mathcal{I}' \to \mathcal{C}$,
où $\mathcal{I}'$ est une catégorie d'indices
filtrante\footnote{On prend pour $\mathcal{I}'$ les mêmes objets que pour
$\mathcal{I}$, mais $\Mor_{\mathcal{I}'}(x, y)$ est le quotient de
$\Mor_\mathcal{I}(x, y)$ par la relation d'équivalence qui identifie
$a, b : x \to y$ si $M(a) = M(b)$.},
de telle sorte que $\colim_\mathcal{I} M$ existe si et seulement si
$\colim_{\mathcal{I}'} M'$ existe ; dans ce cas, les colimites sont
canoniquement isomorphes.

\medskip\noindent
Supposons que $M : \mathcal{I} \to \textit{Ensembles}$ soit un diagramme filtrant.
Dans ce cas, la relation d'équivalence de la formule
$$
\colim_\mathcal{I} M
=
(\coprod\nolimits_{i\in I} M_i)/\sim
$$
se décrit simplement comme suit :
$$
m_i \sim m_{i'}
\Leftrightarrow
\exists i'', \phi : i \to i'', \phi': i' \to i'',
M(\phi)(m_i) = M(\phi')(m_{i'}).
$$
Autrement dit, deux éléments sont égaux dans la colimite si et seulement
s'ils « finissent par devenir égaux ».

\begin{lemma}
\label{categories-lemma-directed-commutes}
Soient $\mathcal{I}$ et $\mathcal{J}$ des catégories d'indices.
Supposons que $\mathcal{I}$ soit filtrante et que $\mathcal{J}$ soit finie.
Soit $M : \mathcal{I} \times \mathcal{J} \to \textit{Ensembles}$,
$(i, j) \mapsto M_{i, j}$, un diagramme de diagrammes d'ensembles.
Alors
$$
\colim_i \lim_j M_{i, j}
=
\lim_j \colim_i M_{i, j}.
$$
En particulier, les colimites sur $\mathcal{I}$ commutent aux produits finis,
aux produits fibrés et aux égalisateurs d'ensembles.
\end{lemma}

\begin{proof}
Omis. C'est en fait un exercice agréable de montrer qu'une catégorie est
filtrante si et seulement si les colimites sur cette catégorie commutent
aux limites finies (à valeurs dans la catégorie des ensembles).
\end{proof}

\noindent
Donnons un contre-exemple au lemme lorsque $\mathcal{J}$ est infinie.
Prenons pour $\mathcal{I}$ la catégorie dont les objets sont
$\mathbf{N} = \{1, 2, 3, \ldots\}$ et dans laquelle il existe un unique
morphisme $i \to i'$ dès que $i \leq i'$.
Prenons pour $\mathcal{J}$ la catégorie discrète
$\mathbf{N} = \{1, 2, 3, \ldots\}$ (ses seuls morphismes sont les identités).
Posons $M_{i, j} = \{1, 2, \ldots, i\}$, avec les applications d'inclusion
évidentes $M_{i, j} \to M_{i', j}$ lorsque $i \leq i'$. Dans ce cas,
$\colim_i M_{i, j} = \mathbf{N}$, et donc
$$
\lim_j \colim_i M_{i, j}
=
\prod\nolimits_j \mathbf{N}
=
\mathbf{N}^\mathbf{N}
$$
D'autre part, $\lim_j M_{i, j} = \prod\nolimits_j M_{i, j}$, et donc
$$
\colim_i \lim_j M_{i, j}
=
\bigcup\nolimits_i \{1, 2, \ldots, i\}^{\mathbf{N}}
$$
qui est plus petit que l'autre limite.

\begin{lemma}
\label{categories-lemma-cofinal-in-filtered}
\begin{reference}
\cite[Proposition 3.2.4]{KS}
\end{reference}
Soit $\mathcal{I}$ une catégorie. Soit $\mathcal{J}$ une sous-catégorie
pleine. Supposons que $\mathcal{I}$ soit filtrante. Supposons en outre que,
pour tout objet $i$ de $\mathcal{I}$, il existe un morphisme $i \to j$
vers un objet $j$ de $\mathcal{J}$. Alors $\mathcal{J}$
est filtrante et cofinale dans $\mathcal{I}$.
\end{lemma}

\begin{proof}
Omis. Exercice agréable sur les notions en jeu.
\end{proof}

\noindent
Nous avons parfois besoin d'un contrôle plus fin des
conditions que l'on peut imposer aux catégories d'indices. Nous ajoutons
donc quelques lemmes sur les diverses conditions possibles.

\begin{lemma}
\label{categories-lemma-preserve-products}
Soit $\mathcal{I}$ une catégorie d'indices, c'est-à-dire une catégorie.
Supposons que, pour toute paire d'objets $x, y$ de $\mathcal{I}$,
il existe un objet $z$ et des morphismes $x \to z$ et $y \to z$.
Alors
\begin{enumerate}
\item si $M$ et $N$ sont des diagrammes d'ensembles sur $\mathcal{I}$,
alors $\colim (M_i \times N_i) \to \colim M_i \times \colim N_i$
est surjective ;
\item en général, les colimites de diagrammes d'ensembles sur $\mathcal{I}$
ne commutent pas aux produits finis non vides.
\end{enumerate}
\end{lemma}

\begin{proof}
Preuve de (1). Soit $(\overline{m}, \overline{n})$
un élément de $\colim M_i \times \colim N_i$.
Il existe alors $m \in M_x$ et $n \in N_y$, pour certains
$x, y \in \Ob(\mathcal{I})$, tels que $m$ s'envoie sur
$\overline{m}$ et $n$ sur $\overline{n}$. Voir la
section \ref{categories-section-limit-sets}.
Choisissons $a : x \to z$ et $b : y \to z$
dans $\mathcal{I}$. Alors $(M(a)(m), N(b)(n))$ est un élément de
$(M \times N)_z$ dont l'image dans $\colim (M_i \times N_i)$
s'envoie sur $(\overline{m}, \overline{n})$, comme voulu.

\medskip\noindent
Preuve de (2). Soit $G$ un groupe non trivial et
soit $\mathcal{I}$ la catégorie à un objet dont le monoïde des
endomorphismes est $G$. Alors $\mathcal{I}$ satisfait trivialement la
condition du lemme. Faisons agir $G$ sur lui-même par translations et
considérons le $G$-ensemble $G$ comme un diagramme indexé par $\mathcal{I}$
à valeurs dans les ensembles. Alors
$$
\colim_\mathcal{I} G \times \colim_\mathcal{I} G \cong G/G \times G/G
$$
n'est pas isomorphe à
$$
\colim_\mathcal{I} (G \times G) \cong (G \times G)/G
$$
Cet exemple montre qu'on ne peut pas simplement supprimer la condition
supplémentaire du lemme \ref{categories-lemma-directed-commutes},
même si l'on ne s'intéresse qu'aux produits finis.
\end{proof}

\begin{lemma}
\label{categories-lemma-colimits-abelian-as-sets}
Soit $\mathcal{I}$ une catégorie d'indices, c'est-à-dire une catégorie.
Supposons que, pour toute paire d'objets $x, y$ de $\mathcal{I}$,
il existe un objet $z$ et des morphismes $x \to z$ et $y \to z$.
Soit $M : \mathcal{I} \to \textit{Ab}$ un diagramme de groupes abéliens
sur $\mathcal{I}$. Alors l'application naturelle de la colimite de $M$
dans la catégorie des ensembles vers la colimite de $M$ dans la catégorie
des groupes abéliens est surjective.
\end{lemma}

\begin{proof}
Rappelons que la colimite dans la catégorie des ensembles est le quotient
de la réunion disjointe $\coprod M_i$ par une relation d'équivalence ; voir la
section \ref{categories-section-limit-sets}.
De même, la colimite dans la catégorie des groupes abéliens est un quotient
de la somme directe $\bigoplus M_i$.
L'hypothèse du lemme signifie que, étant donnés
$i, j \in \Ob(\mathcal{I})$, $m \in M_i$ et $n \in M_j$, on peut trouver
un objet $k$ et des morphismes $a : i \to k$ et $b : j \to k$.
Ainsi, $m + n$ est représenté dans la colimite par l'élément
$M(a)(m) + M(b)(n)$ de $M_k$. L'application naturelle de la réunion
disjointe $\coprod M_i$ vers la colimite est donc surjective.
\end{proof}

\begin{lemma}
\label{categories-lemma-split-into-connected}
Soit $\mathcal{I}$ une catégorie d'indices, c'est-à-dire une catégorie.
Supposons que, pour tout diagramme en traits pleins
$$
\xymatrix{
x \ar[d] \ar[r] & y \ar@{..>}[d] \\
z \ar@{..>}[r] & w
}
$$
dans $\mathcal{I}$, il existe un objet $w$ et des flèches pointillées
qui rendent le diagramme commutatif. Alors $\mathcal{I}$ est soit vide,
soit une réunion disjointe non vide de catégories connexes possédant
la même propriété.
\end{lemma}

\begin{proof}
Si $\mathcal{I}$ est la catégorie vide, le lemme est vrai.
Sinon, définissons une relation sur les objets de $\mathcal{I}$ en posant
$x \sim y$ s'il existe $z$ et des morphismes
$x \to z$ et $y \to z$. C'est une relation d'équivalence
par l'hypothèse du lemme. Ainsi, $\Ob(\mathcal{I})$
est une réunion disjointe de classes d'équivalence. Soient
$\mathcal{I}_j$ les sous-catégories pleines correspondant à ces classes.
Alors $\mathcal{I} = \coprod \mathcal{I}_j$, où les $\mathcal{I}_j$
sont non vides, comme voulu.
\end{proof}

\begin{lemma}
\label{categories-lemma-preserve-injective-maps}
Soit $\mathcal{I}$ une catégorie d'indices, c'est-à-dire une catégorie.
Supposons que, pour tout diagramme en traits pleins
$$
\xymatrix{
x \ar[d] \ar[r] & y \ar@{..>}[d] \\
z \ar@{..>}[r] & w
}
$$
dans $\mathcal{I}$, il existe un objet $w$ et des flèches pointillées
qui rendent le diagramme commutatif. Alors
\begin{enumerate}
\item tout morphisme injectif $M \to N$ de diagrammes d'ensembles sur
$\mathcal{I}$ induit une application injective
$\colim M_i \to \colim N_i$ d'ensembles ;
\item en général, il n'en va pas de même pour les diagrammes de groupes
abéliens et leurs colimites.
\end{enumerate}
\end{lemma}

\begin{proof}
Si $\mathcal{I}$ est la catégorie vide, le lemme est vrai.
Nous pouvons donc supposer $\mathcal{I}$ non vide. Dans ce cas,
on peut écrire $\mathcal{I} = \coprod \mathcal{I}_j$, où chaque
$\mathcal{I}_j$ est non vide et satisfait la même propriété ; voir le
lemme \ref{categories-lemma-split-into-connected}. Puisque
$\colim_\mathcal{I} M = \coprod_j \colim_{\mathcal{I}_j} M|_{\mathcal{I}_j}$,
la preuve de (1) se ramène au cas connexe.

\medskip\noindent
Supposons $\mathcal{I}$ connexe et $M \to N$ injectif, c'est-à-dire
que toutes les applications $M_i \to N_i$ sont injectives.
Identifions $M_i$ à l'image de $M_i \to N_i$ ; autrement dit,
considérons $M_i$ comme un sous-ensemble de $N_i$.
Nous utiliserons sans plus de mention la description des colimites donnée
dans la section \ref{categories-section-limit-sets}.
Soient $s, s' \in \colim M_i$ ayant la même image dans $\colim N_i$.
Disons que $s$ provient d'un élément $m$ de $M_i$ et $s'$ d'un
élément $m'$ de $M_{i'}$. Il existe alors une suite
$i = i_0, i_1, \ldots, i_{2n} = i'$ d'objets de $\mathcal{I}$
et des morphismes
$$
\xymatrix{
&
i_1 \ar[ld] \ar[rd] & &
i_3 \ar[ld] & &
i_{2n-1} \ar[rd] & \\
i = i_0 & &
i_2 & &
\ldots & &
i_{2n} = i'
}
$$
ainsi que des éléments $n_{i_j} \in N_{i_j}$ dont les images coïncident
par les applications $N_{i_{2k-1}} \to N_{i_{2k-2}}$ et
$N_{i_{2k-1}}
\to N_{i_{2k}}$ induites par les flèches de
$\mathcal{I}$ ci-dessus, avec $n_{i_0} = m$ et $n_{i_{2n}} = m'$.
Nous allons montrer par récurrence sur $n$ que cela entraîne $s = s'$.
Le cas initial $n = 0$ est trivial. Supposons $n \geq 1$.
L'hypothèse sur $\mathcal{I}$ fournit un diagramme commutatif
$$
\xymatrix{
& i_1 \ar[ld] \ar[rd] \\
i_0 \ar[rd] & & i_2 \ar[ld] \\
& w
}
$$
Nous en déduisons que $m$ et $n_{i_2}$ ont la même image dans $N_w$,
puisqu'ils sont tous deux l'image de l'élément $n_{i_1}$.
En particulier, cet élément est un élément $m'' \in M_w$ qui
définit le même élément que $s$ dans $\colim M_i$.
Nous obtenons alors la chaîne
$$
\xymatrix{
&
i_3 \ar[ld] \ar[rd] & &
i_5 \ar[ld] & &
i_{2n-1} \ar[rd] & \\
w & &
i_4 & &
\ldots & &
i_{2n} = i'
}
$$
et les éléments $n_{i_j}$ pour $j \geq 3$ ; cette chaîne est plus courte
que la chaîne initiale. Cela démontre l'étape de récurrence et achève
la preuve de (1).

\medskip\noindent
Soit $G$ un groupe et soit $\mathcal{I}$ la catégorie à un objet dont
le monoïde des endomorphismes est $G$. Alors $\mathcal{I}$ satisfait
la condition du lemme, car, étant donnés $g_1, g_2 \in G$, on peut trouver
$h_1, h_2 \in G$ tels que $h_1 g_1 = h_2 g_2$.
Un diagramme $M$ sur $\mathcal{I}$ à valeurs dans $\textit{Ab}$ est la même
chose qu'un groupe abélien $M$ muni d'une action de $G$, et
$\colim_\mathcal{I} M$ est le groupe des coinvariants $M_G$ de $M$.
Prenons pour $G$ le groupe d'ordre $2$ agissant trivialement sur
$M = \mathbf{Z}/2\mathbf{Z}$, plongé dans le premier facteur de
$N = \mathbf{Z}/2\mathbf{Z} \times \mathbf{Z}/2\mathbf{Z}$,
où l'élément non trivial de $G$ agit par
$(x, y) \mapsto (x + y, y)$. Alors $M_G \to N_G$ est nul.
\end{proof}

\begin{lemma}
\label{categories-lemma-split-into-directed}
Soit $\mathcal{I}$ une catégorie d'indices, c'est-à-dire une catégorie.
Supposons que
\begin{enumerate}
\item pour toute paire de morphismes $a : w \to x$ et $b : w \to y$
dans $\mathcal{I}$, il existe un objet $z$ et des morphismes
$c : x \to z$ et $d : y \to z$ tels que $c \circ a = d \circ b$, et
\item pour toute paire de morphismes $a, b : x \to y$, il existe
un morphisme $c : y \to z$ tel que $c \circ a = c \circ b$.
\end{enumerate}
Alors $\mathcal{I}$ est une réunion (éventuellement vide)
de catégories d'indices filtrantes disjointes $\mathcal{I}_j$.
\end{lemma}

\begin{proof}
Si $\mathcal{I}$ est la catégorie vide, le lemme est vrai.
Sinon, définissons une relation sur les objets de $\mathcal{I}$ en posant
$x \sim y$ s'il existe $z$ et des morphismes
$x \to z$ et $y \to z$. C'est une relation d'équivalence
par la première hypothèse du lemme. Ainsi, $\Ob(\mathcal{I})$
est une réunion disjointe de classes d'équivalence. Soient
$\mathcal{I}_j$ les sous-catégories pleines correspondant à ces classes.
Le reste découle immédiatement des définitions.
\end{proof}

\begin{lemma}
\label{categories-lemma-almost-directed-commutes-equalizers}
Soit $\mathcal{I}$ une catégorie d'indices satisfaisant les hypothèses du
lemme \ref{categories-lemma-split-into-directed} ci-dessus. Alors les colimites sur
$\mathcal{I}$ commutent aux produits fibrés et aux égalisateurs d'ensembles
(et, plus généralement, aux limites finies connexes).
\end{lemma}

\begin{proof}
D'après le lemme \ref{categories-lemma-split-into-directed},
nous pouvons écrire $\mathcal{I} = \coprod \mathcal{I}_j$, chaque
$\mathcal{I}_j$ étant filtrante. D'après le
lemme \ref{categories-lemma-directed-commutes}, les colimites sur $\mathcal{I}_j$
commutent aux égalisateurs et aux produits fibrés. Il suffit donc de montrer
que les égalisateurs et les produits fibrés commutent aux coproduits dans la
catégorie des ensembles (y compris aux coproduits vides).
Autrement dit, étant donnés un ensemble $J$, des ensembles
$A_j, B_j, C_j$ et des applications
$A_j \to B_j$, $C_j \to B_j$ pour $j \in J$, il faut montrer que
$$
(\coprod\nolimits_{j \in J} A_j)
\times_{(\coprod\nolimits_{j \in J} B_j)}
(\coprod\nolimits_{j \in J} C_j)
=
\coprod\nolimits_{j \in J} A_j \times_{B_j} C_j
$$
et, étant données $a_j, a'_j : A_j \to B_j$, que
$$
\text{Égalisateur}(
\coprod\nolimits_{j \in J} a_j,
\coprod\nolimits_{j \in J} a'_j)
=
\coprod\nolimits_{j \in J}
\text{Égalisateur}(a_j, a'_j)
$$
Cela reste vrai si $J = \emptyset$. Les détails sont omis.
\end{proof}



\section{Limites projectives cofiltrantes}
\label{categories-section-codirected-limits}

\noindent
Les limites projectives sont plus faciles à calculer ou à décrire lorsqu'elles
portent sur un diagramme cofiltrant. En voici la définition.

\begin{definition}
\label{categories-definition-codirected}
On dit qu'un diagramme $M : \mathcal{I} \to \mathcal{C}$ est
{\it codirigé} ou {\it cofiltrant} si les conditions suivantes sont satisfaites :
\begin{enumerate}
\item la catégorie $\mathcal{I}$ possède au moins un objet ;
\item pour toute paire d'objets $x, y$ de $\mathcal{I}$, il existe un objet
$z$ et des morphismes $z \to x$, $z \to y$ ;
\item pour toute paire d'objets $x, y$ de $\mathcal{I}$ et toute paire de
morphismes $a, b : x \to y$ de $\mathcal{I}$, il existe un morphisme
$c : w \to x$ de $\mathcal{I}$ tel que
$M(a \circ c) = M(b \circ c)$ comme morphismes dans $\mathcal{C}$.
\end{enumerate}
On dit qu'une catégorie d'indices $\mathcal{I}$ est {\it codirigée}, ou
{\it cofiltrante}, si $\text{id} : \mathcal{I} \to \mathcal{I}$ est
cofiltrant (autrement dit, on efface le $M$ dans la partie (3) ci-dessus).
\end{definition}

\noindent
On observe que tout diagramme dont la catégorie d'indices est cofiltrante
est cofiltrant ; c'est sous cette forme que la situation se présente
habituellement.

\medskip\noindent
Pour illustrer en quoi les limites projectives cofiltrantes d'ensembles sont
« plus faciles » que les limites projectives générales, signalons qu'un
diagramme cofiltrant d'ensembles finis non vides a une limite non vide
(lemme \ref{categories-lemma-nonempty-limit}). Ce résultat n'est pas vrai pour une
limite générale d'ensembles finis non vides.











\section{Limites projectives et inductives sur des ensembles préordonnés}
\sectionmark{Limites sur les ensembles préordonnés}
\label{categories-section-posets-limits}

\noindent
Les systèmes sur des ensembles préordonnés constituent un cas particulier
de diagrammes.

\begin{definition}
\label{categories-definition-directed-set}
Soit $I$ un ensemble et soit $\leq$ une relation binaire sur $I$.
\begin{enumerate}
\item On dit que $\leq$ est un {\it préordre} si cette relation est
transitive (si $i \leq j$ et $j \leq k$, alors $i \leq k$) et réflexive
($i \leq i$ pour tout $i \in I$).
\item Un {\it ensemble préordonné} est un ensemble muni d'un préordre.
\item Un {\it ensemble ordonné filtrant} est un ensemble préordonné
$(I, \leq)$ tel que $I$ soit non vide et tel que $\forall i, j \in I$,
il existe $k \in I$ vérifiant $i \leq k, j \leq k$.
\item On dit que $\leq$ est un {\it ordre partiel} si c'est un préordre
antisymétrique (si $i \leq j$ et $j \leq i$, alors $i = j$).
\item Un {\it ensemble partiellement ordonné} est un ensemble muni d'un
ordre partiel.
\item Un {\it ensemble partiellement ordonné filtrant} est un ensemble
ordonné filtrant dont l'ordre est partiel.
\end{enumerate}
\end{definition}

\noindent
Il est d'usage d'omettre $\leq$ de la notation lorsqu'on parle d'ensembles
préordonnés, c'est-à-dire de parler de l'ensemble préordonné $I$ plutôt que
de l'ensemble préordonné $(I, \leq)$. Étant donné un ensemble préordonné
$I$, le symbole $\geq$ est défini par la règle
$i \geq j \Leftrightarrow j \leq i$ pour tous $i, j \in I$.
L'abréviation anglaise « poset » est parfois employée pour désigner un
ensemble partiellement ordonné.

\medskip\noindent
À tout ensemble préordonné $I$, on peut associer une catégorie : ses objets
sont les éléments de $I$, il existe exactement un morphisme $i \to i'$ si
$i \leq i'$, et il n'en existe aucun sinon. Réciproquement, étant donnée une
catégorie $\mathcal{C}$ ayant au plus une flèche entre deux objets quelconques,
l'ensemble $\Ob(\mathcal{C})$ est muni du préordre défini par la règle
$x \leq y \Leftrightarrow \Mor_\mathcal{C}(x, y) \not = \emptyset$.

\begin{definition}
\label{categories-definition-system-over-poset}
Soit $(I, \leq)$ un ensemble préordonné. Soit $\mathcal{C}$ une catégorie.
\begin{enumerate}
\item Un {\it système sur $I$ à valeurs dans $\mathcal{C}$}, parfois appelé
{\it système inductif sur $I$ à valeurs dans $\mathcal{C}$}, est la donnée d'objets
$M_i$ de $\mathcal{C}$ et, pour chaque $i \leq i'$, d'un morphisme
$f_{ii'} : M_i \to M_{i'}$ tel que $f_{ii} = \text{id}$ et
$f_{ii''} = f_{i'i''} \circ f_{i i'}$ dès que $i \leq i' \leq i''$.
\item Un {\it système inverse sur $I$ à valeurs dans $\mathcal{C}$}, parfois appelé
{\it système projectif sur $I$ à valeurs dans $\mathcal{C}$}, est la donnée d'objets
$M_i$ de $\mathcal{C}$ et, pour chaque $i' \leq i$, d'un morphisme
$f_{ii'} : M_i \to M_{i'}$ tel que $f_{ii} = \text{id}$ et
$f_{ii''} = f_{i'i''} \circ f_{i i'}$ dès que $i'' \leq i' \leq i$.
(Remarquer le renversement des inégalités.)
\end{enumerate}
On écrira $(M_i, f_{ii'})$ pour désigner un système de l'un ou l'autre type sur $I$.
Les morphismes $f_{ii'}$ sont parfois appelés les
{\it morphismes de transition}.
\end{definition}

\noindent
Autrement dit, un système sur $I$ n'est autre qu'un diagramme
$M : \mathcal{I} \to \mathcal{C}$, où $\mathcal{I}$ est la catégorie
associée ci-dessus à $I$ : ses objets sont les éléments de $I$ et il existe
une unique flèche $i \to i'$ dans $\mathcal{I}$ si et seulement si
$i \leq i'$. Un système inverse est un diagramme
$M : \mathcal{I}^{opp} \to \mathcal{C}$. De ce point de vue, on pourrait
prendre les limites et colimites de tout système, inverse ou non, sur $I$. Toutefois, il est
d'usage de ne prendre {\it que les limites inductives des systèmes sur $I$}
et {\it que les limites projectives des systèmes inverses sur $I$}.
Plus précisément, étant donné un système $(M_i, f_{ii'})$ sur $I$, la limite
inductive du système $(M_i, f_{ii'})$ est définie par
$$
\colim_{i \in I} M_i = \colim_\mathcal{I} M,
$$
c'est-à-dire comme la limite inductive du diagramme correspondant.
Étant donné un système inverse $(M_i, f_{ii'})$ sur $I$, la limite projective
du système inverse $(M_i, f_{ii'})$ est définie par
$$
\lim_{i \in I} M_i = \lim_{\mathcal{I}^{opp}} M,
$$
c'est-à-dire comme la limite projective du diagramme correspondant.

\begin{remark}
\label{categories-remark-preorder-versus-partial-order}
Soit $I$ un ensemble préordonné. À partir de $I$, on peut construire un
ensemble partiellement ordonné canonique $\overline{I}$ et une application
croissante $\pi : I \to \overline{I}$. En effet, on peut définir une relation
d'équivalence $\sim$ sur $I$ par la règle
$$
i \sim j \Leftrightarrow (i \leq j\text{ et }j \leq i).
$$
On pose $\overline{I} = I/\sim$ et l'on note
$\pi : I \to \overline{I}$ l'application quotient. Enfin, $\overline{I}$
est muni d'un unique ordre partiel tel que
$\pi(i) \leq \pi(j) \Leftrightarrow i \leq j$.
Remarquons que, si $I$ est un ensemble ordonné filtrant, alors
$\overline{I}$ est un ensemble partiellement ordonné filtrant.
Étant donné un système, inverse ou non, $N$ sur $\overline{I}$, on obtient
un système du même type $M$ sur $I$ en posant $M_i = N_{\pi(i)}$.
Cette construction définit un foncteur entre les catégories de systèmes,
inverses ou non, sur $I$ et sur $\overline{I}$. C'est en fait une équivalence.
La raison en est que, si $i \sim j$, alors, pour tout système $M$ sur $I$,
les morphismes $M_i \to M_j$ et $M_j \to M_i$ sont des isomorphismes
réciproques. Plus précisément, si l'on choisit une section
$s : \overline{I} \to I$ de $\pi$, un quasi-inverse du foncteur précédent
envoie $M$ sur $N$ défini par
$N_{\overline{i}} = M_{s(\overline{i})}$.
Enfin, cette correspondance est compatible avec les limites inductives des
systèmes : si $M$ et $N$ sont reliés comme ci-dessus et si l'une des limites
$\colim_{\overline{I}} N$ ou $\colim_I M$ existe, alors l'autre existe aussi et
$\colim_{\overline{I}} N = \colim_I M$.
Des résultats analogues valent pour les systèmes inverses et leurs limites
projectives.
\end{remark}

\noindent
Il résulte de la remarque \ref{categories-remark-preorder-versus-partial-order} que,
lorsqu'on calcule la limite inductive d'un système ou la limite projective
d'un système inverse, on peut toujours supposer que le préordre est un ordre
partiel.

\begin{definition}
\label{categories-definition-directed-system}
Soit $I$ un ensemble préordonné. On dit qu'un système (resp.\ un système
inverse) $(M_i, f_{ii'})$ est un {\it système filtrant} (resp.\ un
{\it système inverse filtrant}) si $I$ est un ensemble ordonné filtrant
(définition \ref{categories-definition-directed-set}) : $I$ est non vide et, pour tous
$i_1, i_2 \in I$, il existe $i\in I$ tel que
$i_1 \leq i$ et $i_2 \leq i$.
\end{definition}

\noindent
Dans ce cas, la limite inductive est parfois appelée (malheureusement)
« limite directe ». Nous n'emploierons pas cette dernière terminologie.
En fait, les diagrammes sur une catégorie filtrante ne sont pas plus
généraux que les systèmes filtrants, au sens suivant.

\begin{lemma}
\label{categories-lemma-directed-category-system}
Soit $\mathcal{I}$ une catégorie d'indices filtrante. Il existe un ensemble
ordonné filtrant $I$ et un système $(x_i, \varphi_{ii'})$ sur $I$ à valeurs dans
$\mathcal{I}$ possédant les propriétés suivantes :
\begin{enumerate}
\item Pour toute catégorie $\mathcal{C}$ et tout diagramme
$M : \mathcal{I} \to \mathcal{C}$ à valeurs dans $\mathcal{C}$, notons
$(M(x_i), M(\varphi_{ii'}))$ le système correspondant sur $I$. Si
$\colim_{i \in I} M(x_i)$ existe, alors $\colim_\mathcal{I} M$ existe aussi
et la transformation
$$
\theta :
\colim_{i \in I} M(x_i)
\longrightarrow
\colim_\mathcal{I} M
$$
du lemme \ref{categories-lemma-functorial-colimit} est un isomorphisme.
\item Pour toute catégorie $\mathcal{C}$ et tout diagramme
$M : \mathcal{I}^{opp} \to \mathcal{C}$ à valeurs dans $\mathcal{C}$, notons
$(M(x_i), M(\varphi_{ii'}))$ le système inverse correspondant sur $I$.
Si $\lim_{i \in I} M(x_i)$ existe, alors $\lim_{\mathcal{I}^{opp}} M$
existe aussi et la transformation
$$
\theta :
\lim_{\mathcal{I}^{opp}} M
\longrightarrow
\lim_{i \in I} M(x_i)
$$
du lemme \ref{categories-lemma-functorial-limit} est un isomorphisme.
\end{enumerate}
\end{lemma}

\begin{proof}
Comme expliqué dans le texte qui suit la
définition \ref{categories-definition-system-over-poset}, on peut considérer les
ensembles préordonnés comme des catégories et les systèmes comme des
foncteurs. Tout au long de la preuve, nous passerons librement d'un point de
vue à l'autre. Nous démontrons le premier énoncé en construisant une catégorie
$\mathcal{I}_0$ correspondant à un ensemble ordonné
filtrant\footnote{En fait, notre construction produira un ensemble
partiellement ordonné filtrant.}, et un foncteur cofinal
$M_0 : \mathcal{I}_0 \to \mathcal{I}$. Alors, d'après le
lemme \ref{categories-lemma-cofinal}, la limite inductive d'un diagramme
$M : \mathcal{I} \to \mathcal{C}$ coïncide avec celle du diagramme
$M \circ M_0 : \mathcal{I}_0 \to \mathcal{C}$, d'où l'énoncé. Le second
énoncé est dual du premier et peut se démontrer en interprétant une limite
projective dans $\mathcal{C}$ comme une limite inductive dans
$\mathcal{C}^{opp}$. Nous omettons les détails.

\medskip\noindent
Une catégorie $\mathcal{F}$ est dite {\em finiment engendrée} s'il existe un
ensemble fini $F$ de flèches de $\mathcal{F}$ tel que toute flèche de
$\mathcal{F}$ s'obtienne en composant des flèches de $F$. Cela implique en
particulier que $\mathcal{F}$ possède un nombre fini d'objets. Commençons par
nous ramener au cas où $\mathcal{I}$ possède la propriété suivante : toute
sous-catégorie finiment engendrée de $\mathcal{I}$ est contenue dans une
sous-catégorie finiment engendrée ayant un unique objet final.

\medskip\noindent
Notons $\omega$ l'ensemble ordonné filtrant des ordinaux finis, que nous
considérons comme une catégorie filtrante. On vérifie aisément que la
catégorie produit $\mathcal{I}\times \omega$ est elle aussi filtrante et que
la projection
$\Pi : \mathcal{I} \times \omega \to \mathcal{I}$
est cofinale.

\medskip\noindent
Soit maintenant $\mathcal{F}$ une sous-catégorie finiment engendrée
quelconque de $\mathcal{I}\times \omega$. À l'aide des axiomes d'une catégorie
filtrante et d'une récurrence simple sur un ensemble fini de générateurs de
$\mathcal{F}$, on peut construire un cocône $(\{f_i\}, i_\infty)$ dans
$\mathcal{I} \times \omega$ pour le diagramme
$\mathcal{F} \to \mathcal{I} \times \omega$. Autrement dit, on dispose d'un
morphisme $f_i : i \to i_\infty$ pour chaque objet $i$ de $\mathcal{F}$,
tel que, pour toute flèche $f : i \to i'$ de $\mathcal{F}$, on ait
$f_i = f_{i'} \circ f$. On peut en outre choisir $i_\infty$ de sorte qu'il
n'existe aucune flèche de $i_\infty$ vers un objet de $\mathcal{F}$. Cela est
possible puisqu'on peut toujours composer à gauche les flèches $f_i$ avec
une flèche qui est l'identité sur la composante $\mathcal{I}$ et strictement
croissante sur la composante $\omega$.
Notons maintenant $\mathcal{F}^+$ la catégorie formée de tous les objets et
de toutes les flèches de $\mathcal{F}$, auxquels on adjoint l'objet
$i_\infty$, la flèche identité $\text{id}_{i_\infty}$ et les flèches $f_i$.
Comme aucune flèche de $i_\infty$ dans $\mathcal{F}^+$ n'aboutit à un objet de
$\mathcal{F}$, l'ensemble des flèches de $\mathcal{F}^+$ est stable par
composition ; $\mathcal{F}^+$ est donc bien une catégorie. Par construction,
c'est une sous-catégorie finiment engendrée de $\mathcal{I}\times\omega$ dont
$i_\infty$ est l'unique objet final. Puisque, d'après le
lemme \ref{categories-lemma-cofinal}, la limite inductive de tout diagramme
$M : \mathcal{I} \to \mathcal{C}$ coïncide avec celle de $M\circ\Pi$, on
obtient ainsi la réduction voulue.

\medskip\noindent
L'ensemble de toutes les sous-catégories finiment engendrées de
$\mathcal{I}$ ayant un unique objet final est naturellement ordonné par
inclusion. On prend pour $\mathcal{I}_0$ la catégorie correspondant à cet
ensemble. On dispose aussi d'un foncteur
$M_0 : \mathcal{I}_0 \to \mathcal{I}$ qui envoie une flèche
$\mathcal{F} \subset \mathcal{F'}$ de $\mathcal{I}_0$ sur l'unique morphisme
de l'objet final de $\mathcal{F}$ vers l'objet final de $\mathcal{F}'$.
Étant données deux sous-catégories finiment engendrées quelconques de
$\mathcal{I}$, la catégorie qu'elles engendrent est elle aussi finiment
engendrée. D'après l'hypothèse faite sur $\mathcal{I}$, elle est donc contenue
dans une sous-catégorie finiment engendrée de $\mathcal{I}$ ayant un unique
objet final. Cela montre que $\mathcal{I}_0$ est filtrante.

\medskip\noindent
Vérifions enfin que $M_0$ est cofinal. Tout objet de $\mathcal{I}$ étant
l'objet final de la sous-catégorie formée de ce seul objet et de sa flèche
identité, le foncteur $M_0$ est surjectif sur les objets. La condition (1) de
la définition \ref{categories-definition-cofinal} est donc satisfaite. Étant donnés un
objet $i$ de $\mathcal{I}$, des objets $\mathcal{F}_1, \mathcal{F}_2$ de
$\mathcal{I}_0$ et des morphismes
$\varphi_1 : i \to M_0(\mathcal{F}_1)$ et
$\varphi_2 : i \to M_0(\mathcal{F}_2)$ dans $\mathcal{I}$, on peut prendre
pour $\mathcal{F}_{12}$ une catégorie finiment engendrée ayant un unique
objet final et contenant $\mathcal{F}_1$, $\mathcal{F}_2$ ainsi que les
morphismes $\varphi_1, \varphi_2$. Le diagramme obtenu est commutatif :
$$
\xymatrix{
& M_0(\mathcal{F}_{12}) & \\
M_0(\mathcal{F}_{1}) \ar[ru] & & M_0(\mathcal{F}_{2}) \ar[lu] \\
& i \ar[lu] \ar[ru]
}
$$
car il se trouve dans la catégorie $\mathcal{F}_{12}$ et
$M_0(\mathcal{F}_{12})$ est final dans cette catégorie. La condition (2) est
donc elle aussi satisfaite, ce qui achève la preuve.
\end{proof}

\begin{remark}
\label{categories-remark-trick-needed}
Remarquons qu'un ensemble ordonné filtrant fini $(I, \leq)$ possède toujours
un plus grand élément $i_\infty$. Toute limite inductive d'un système
$(M_i, f_{ii'})$ sur un tel ensemble est donc triviale, au sens où elle est
égale à $M_{i_\infty}$. En revanche, une limite inductive indexée par une
catégorie filtrante finie n'est pas nécessairement triviale. Par exemple,
soit $\mathcal{I}$ la catégorie ayant un unique objet $i$ et un unique
morphisme non trivial $e$ satisfaisant $e = e \circ e$. La limite inductive
d'un diagramme $M : \mathcal{I} \to Ensembles$ est l'image de l'idempotent $M(e)$.
Cela montre qu'un artifice tel que le passage à
$\mathcal{I}\times \omega$ dans la preuve du
lemme \ref{categories-lemma-directed-category-system} est essentiel.
\end{remark}

\begin{lemma}
\label{categories-lemma-nonempty-limit}
Si $S : \mathcal{I} \to \textit{Ensembles}$ est un diagramme cofiltrant
d'ensembles et si tous les $S_i$ sont finis et non vides, alors
$\lim_i S_i$ est non vide. Autrement dit, la limite projective d'un système
inverse filtrant d'ensembles finis non vides est non vide.
\end{lemma}

\begin{proof}
Les deux énoncés sont équivalents d'après le
lemme \ref{categories-lemma-directed-category-system}. Soit $I$ un ensemble ordonné
filtrant et soit $(S_i)_{i \in I}$ un système inverse d'ensembles finis non
vides sur $I$. Appelons {\it sous-système} $T$ une famille
$T = (T_i)_{i \in I}$ de parties non vides $T_i \subset S_i$ telle que
$T_{i'}$ soit envoyé dans $T_i$ par le morphisme de transition
$S_{i'} \to S_i$ pour tout $i' \geq i$. Notons $\mathcal{T}$ l'ensemble des
sous-systèmes et ordonnons $\mathcal{T}$ par inclusion. Supposons que les
$T_\alpha$, $\alpha \in A$, forment une famille totalement ordonnée d'éléments
de $\mathcal{T}$. Écrivons
$T_\alpha = (T_{\alpha, i})_{i \in I}$. On peut alors trouver un minorant
$T = (T_i)_{i \in I}$ en posant
$T_i = \bigcap_{\alpha \in A} T_{\alpha, i}$ ; c'est manifestement une partie
finie non vide de $S_i$, puisque tous les $T_{\alpha, i}$ sont non vides et
que les $T_\alpha$ forment une famille totalement ordonnée. Le lemme de Zorn
montre donc que $\mathcal{T}$ possède des éléments minimaux.

\medskip\noindent
Analysons la forme d'un élément minimal $T \in \mathcal{T}$. Remarquons
d'abord que toutes les applications $T_{i'} \to T_i$ sont surjectives.
En effet, puisque $I$ est un ensemble ordonné filtrant et que $T_i$ est fini,
l'intersection $T'_i = \bigcap_{i' \geq i} \Im(T_{i'} \to T_i)$ est non vide.
Ainsi, $T' = (T'_i)$ est un sous-système contenu dans $T$ et, par minimalité,
$T' = T$. Enfin, affirmons que $T_i$ est un singleton pour tout $i$. En effet,
si $x \in T_i$, on peut définir
$T'_{i'} = (T_{i'} \to T_i)^{-1}(\{x\})$ pour $i' \geq i$ et
$T'_j = T_j$ si $j \not \geq i$. Il s'agit encore d'un sous-système, car on
a vu ci-dessus que les morphismes de transition du sous-système $T$ sont
surjectifs. Par minimalité, $T = T'$, ce qui implique bien que $T_i$ est un
singleton. Cela vaut pour chaque $i \in I$ ; on a donc
$T_i = \{x_i\}$ pour un certain $x_i \in S_i$, avec $x_{i'} \mapsto x_i$
par l'application $S_{i'} \to S_i$ pour tout $i' \geq i$. Autrement dit,
$(x_i) \in \lim S_i$, ce qui démontre le lemme.
\end{proof}






\section{Systèmes essentiellement constants}
\label{categories-section-essentially-constant}

\noindent
Soit $M : \mathcal{I} \to \mathcal{C}$ un diagramme à valeurs dans une catégorie
$\mathcal{C}$. Supposons la catégorie d'indices $\mathcal{I}$ filtrante.
Il existe alors trois notions, de plus en plus fortes, qui distinguent un
objet $X$ de $\mathcal{C}$. La première est simplement
$$
X = \colim_{i \in \mathcal{I}} M_i.
$$
L'objet $X$ est alors muni des coprojections $M_i \to X$. Une condition plus
forte consisterait à exiger que $X$ soit la limite inductive et qu'il existe
$i \in \mathcal{I}$ et un morphisme $X \to M_i$ tels que la composée
$X \to M_i \to X$ soit $\text{id}_X$. Voici une condition encore plus forte.

\begin{definition}
\label{categories-definition-essentially-constant-diagram}
Soit $M : \mathcal{I} \to \mathcal{C}$ un diagramme à valeurs dans une catégorie
$\mathcal{C}$.
\begin{enumerate}
\item Supposons la catégorie d'indices $\mathcal{I}$ filtrante et soit
$(X, \{M_i \to X\}_i)$ un cocône de $M$, voir la
remarque \ref{categories-remark-cones-and-cocones}. On dit que $M$ est
{\it essentiellement constant} de {\it valeur} $X$ s'il existe un
$i \in \mathcal{I}$ et un morphisme $X \to M_i$ tels que
\begin{enumerate}
\item $X \to M_i \to X$ soit $\text{id}_X$ ;
\item pour tout $j$, il existe $k$ et des morphismes $i \to k$ et
$j \to k$ tels que le morphisme $M_j \to M_k$ soit égal à la composée
$M_j \to X \to M_i \to M_k$.
\end{enumerate}
\item Supposons la catégorie d'indices $\mathcal{I}$ cofiltrante et soit
$(X, \{X \to M_i\}_i)$ un cône de $M$, voir la
remarque \ref{categories-remark-cones-and-cocones}. On dit que $M$ est
{\it essentiellement constant} de {\it valeur} $X$ s'il existe un
$i \in \mathcal{I}$ et un morphisme $M_i \to X$ tels que
\begin{enumerate}
\item $X \to M_i \to X$ soit $\text{id}_X$ ;
\item pour tout $j$, il existe $k$ et des morphismes $k \to i$ et
$k \to j$ tels que le morphisme $M_k \to M_j$ soit égal à la composée
$M_k \to M_i \to X \to M_j$.
\end{enumerate}
\end{enumerate}
Lorsqu'on utilise cette définition, il faut garder à l'esprit le
lemme \ref{categories-lemma-essentially-constant-is-limit-colimit}.
\end{definition}

\noindent
Le contexte indiquera laquelle des deux versions est considérée. En cas
d'ambiguïté, on les distinguera en disant, dans le premier cas, que $M$ est
essentiellement constant en tant qu'{\it ind-objet} et, dans le second, qu'il
est essentiellement constant en tant que {\it pro-objet}. Cette terminologie
est expliquée plus avant dans les remarques \ref{categories-remark-ind-category} et
\ref{categories-remark-pro-category}. En fait, nous emploierons souvent l'expression
« système essentiellement constant », laquelle, à strictement parler, n'est
définie que pour les systèmes sur des ensembles ordonnés filtrants.

\begin{definition}
\label{categories-definition-essentially-constant-system}
Soit $\mathcal{C}$ une catégorie. Un système filtrant
$(M_i, f_{ii'})$ est un {\it système essentiellement constant} si $M$,
considéré comme un foncteur $I \to \mathcal{C}$, définit un diagramme
essentiellement constant. Un système inverse filtrant $(M_i, f_{ii'})$ est
un {\it système inverse essentiellement constant} si $M$, considéré comme un
foncteur $I^{opp} \to \mathcal{C}$, définit un diagramme inverse
essentiellement constant.
\end{definition}

\noindent
Si $(M_i, f_{ii'})$ est un système essentiellement constant et si les
morphismes $f_{ii'}$ sont des monomorphismes, alors, pour tous
$i \leq i'$ suffisamment grands, les morphismes $f_{ii'}$ sont des
isomorphismes. Considérons en
revanche le système
$$
\mathbf{Z}^2 \to \mathbf{Z}^2 \to \mathbf{Z}^2 \to \ldots
$$
dont les applications sont données par $(a, b) \mapsto (a + b, 0)$. Ce
système est essentiellement constant de valeur $\mathbf{Z}$, mais tout
morphisme de transition possède un noyau.

\medskip\noindent
Voici un exemple de système qui n'est pas essentiellement constant. Posons
$M = \bigoplus_{n \geq 0} \mathbf{Z}$ et soit $S : M \to M$ l'opérateur de
décalage $(a_0, a_1, \ldots) \mapsto (a_1, a_2, \ldots)$. Alors le système
$M \to M \to M \to \ldots$ dont les morphismes de transition sont $S$ a
pour limite inductive $0$, et la composée $0 \to M \to 0$ est l'identité,
mais ce système n'est pas essentiellement constant.

\medskip\noindent
Le lemme suivant est un contrôle de cohérence élémentaire.

\begin{lemma}
\label{categories-lemma-essentially-constant-is-limit-colimit}
Soit $M : \mathcal{I} \to \mathcal{C}$ un diagramme.
Si $\mathcal{I}$ est filtrante et si $M$ est essentiellement constant en tant
qu'ind-objet, alors $X = \colim M_i$ existe et $M$ est essentiellement
constant de valeur $X$.
Si $\mathcal{I}$ est cofiltrante et si $M$ est essentiellement constant en
tant que pro-objet, alors $X = \lim M_i$ existe et $M$ est essentiellement
constant de valeur $X$.
\end{lemma}

\begin{proof}
Omis. C'est un bon exercice sur les définitions.
\end{proof}

\begin{remark}
\label{categories-remark-ind-category}
Soit $\mathcal{C}$ une catégorie. Il existe une grande catégorie
$\text{Ind-}\mathcal{C}$ des {\it ind-objets} de $\mathcal{C}$. En effet, si
$F : \mathcal{I} \to \mathcal{C}$ et $G : \mathcal{J} \to \mathcal{C}$
sont des diagrammes filtrants à valeurs dans $\mathcal{C}$, on peut définir
$$
\Mor_{\text{Ind-}\mathcal{C}}(F, G) =
\lim_i \colim_j \Mor_\mathcal{C}(F(i), G(j)).
$$
Il existe un foncteur canonique $\mathcal{C} \to \text{Ind-}\mathcal{C}$ qui
envoie $X$ sur le {\it système constant} de valeur $X$. C'est un plongement
pleinement fidèle. Dans ce langage, on voit qu'un diagramme $F$ est
essentiellement constant si et seulement si $F$ est isomorphe à un système
constant. Si nous avons un jour besoin de ce résultat, nous l'énoncerons sous
forme de lemme et le démontrerons ici.
\end{remark}

\begin{remark}
\label{categories-remark-pro-category}
Soit $\mathcal{C}$ une catégorie. Il existe une grande catégorie
$\text{Pro-}\mathcal{C}$ des {\it pro-objets} de $\mathcal{C}$. En effet, si
$F : \mathcal{I} \to \mathcal{C}$ et $G : \mathcal{J} \to \mathcal{C}$
sont des diagrammes cofiltrants à valeurs dans $\mathcal{C}$, on peut définir
$$
\Mor_{\text{Pro-}\mathcal{C}}(F, G) =
\lim_j \colim_i \Mor_\mathcal{C}(F(i), G(j)).
$$
Il existe un foncteur canonique $\mathcal{C} \to \text{Pro-}\mathcal{C}$ qui
envoie $X$ sur le {\it système constant} de valeur $X$. C'est un plongement
pleinement fidèle. Dans ce langage, on voit qu'un diagramme $F$ est
essentiellement constant si et seulement si $F$ est isomorphe à un système
constant. Si nous avons un jour besoin de ce résultat, nous l'énoncerons sous
forme de lemme et le démontrerons ici.
\end{remark}

\begin{example}
\label{categories-example-pro-morphism-inverse-systems}
Soit $\mathcal{C}$ une catégorie. Soient $(X_n)$ et $(Y_n)$ des systèmes
inverses dans $\mathcal{C}$ sur $\mathbf{N}$ muni de l'ordre usuel.
Schématiquement :
$$
\ldots \to X_3 \to X_2 \to X_1
\quad\text{et}\quad
\ldots \to Y_3 \to Y_2 \to Y_1
$$
Soit $a : (X_n) \to (Y_n)$ un morphisme de pro-objets de $\mathcal{C}$.
Quelles données représente $a$ ? Pour chaque $n \in \mathbf{N}$, il devrait
exister un $m(n)$ et un morphisme $a_n : X_{m(n)} \to Y_n$. Ces morphismes
devraient être compatibles au sens suivant : pour tout $n' \geq n$, il existe
un $m(n', n) \geq m(n'), m(n)$ tel que le diagramme
$$
\xymatrix{
X_{m(n', n)} \ar[rr] \ar[d] & & X_{m(n)} \ar[d]^{a_n} \\
X_{m(n')} \ar[r]^{a_{n'}} & Y_{n'} \ar[r] & Y_n
}
$$
soit commutatif. Après avoir remplacé $m(n)$ par
$\max_{l \leq k \leq n}\{m(n, k), m(k, l)\}$, on voit que l'on obtient
$\ldots \geq m(3) \geq m(2) \geq m(1)$ et un diagramme commutatif
$$
\xymatrix{
\ldots \ar[r] &
X_{m(3)} \ar[d]^{a_3} \ar[r] &
X_{m(2)} \ar[d]^{a_2} \ar[r] &
X_{m(1)} \ar[d]^{a_1} \\
\ldots \ar[r] &
Y_3 \ar[r] &
Y_2 \ar[r] &
Y_1
}
$$
Étant donnée une application croissante
$m' : \mathbf{N} \to \mathbf{N}$ telle que $m' \geq m$, et en posant
$a'_i : X_{m'(i)} \to X_{m(i)} \to Y_i$, la paire $(m', a')$ définit le même
morphisme de pro-systèmes. Réciproquement, deux paires $(m_1, a_1)$ et
$(m_2, a_2)$ comme ci-dessus définissent le même morphisme de pro-objets si et
seulement s'il existe $m' \geq m_1, m_2$ tel que $a'_1 = a'_2$.
\end{example}

\begin{remark}
\label{categories-remark-pro-category-copresheaves}
Soit $\mathcal{C}$ une catégorie. Soient
$F : \mathcal{I} \to \mathcal{C}$ et $G : \mathcal{J} \to \mathcal{C}$
des diagrammes cofiltrants à valeurs dans $\mathcal{C}$. Considérons les foncteurs
$A, B : \mathcal{C} \to \textit{Ensembles}$ définis par
$$
A(X) = \colim_i \Mor_\mathcal{C}(F(i), X)
\quad\text{et}\quad
B(X) = \colim_j \Mor_\mathcal{C}(G(j), X)
$$
Nous affirmons qu'un morphisme de pro-systèmes de $F$ vers $G$ équivaut à une
transformation de foncteurs $t : B \to A$. En effet, étant donné $t$, on peut
appliquer $t$ à la classe de $\text{id}_{G(j)}$ dans $B(G(j))$ et obtenir un
système compatible d'éléments
$\xi_j \in A(G(j)) = \colim_i \Mor_\mathcal{C}(F(i), G(j))$ ; c'est exactement
notre définition d'un morphisme dans $\text{Pro-}\mathcal{C}$ donnée dans la
remarque \ref{categories-remark-pro-category}. Nous omettons la construction d'une
transformation $B \to A$ à partir d'un morphisme de pro-objets de $F$ vers $G$.
\end{remark}

\begin{lemma}
\label{categories-lemma-image-essentially-constant}
Soit $\mathcal{C}$ une catégorie. Soit $M : \mathcal{I} \to \mathcal{C}$
un diagramme dont la catégorie d'indices $\mathcal{I}$ est filtrante
(resp.\ cofiltrante). Soit $F : \mathcal{C} \to \mathcal{D}$ un foncteur.
Si $M$ est essentiellement constant en tant qu'ind-objet (resp.\ pro-objet),
alors il en va de même de $F \circ M : \mathcal{I} \to \mathcal{D}$.
\end{lemma}

\begin{proof}
Si $X$ est une valeur de $M$, il découle immédiatement de la définition que
$F(X)$ est une valeur de $F \circ M$.
\end{proof}

\begin{lemma}
\label{categories-lemma-characterize-essentially-constant-ind}
Soit $\mathcal{C}$ une catégorie. Soit $M : \mathcal{I} \to \mathcal{C}$
un diagramme dont la catégorie d'indices $\mathcal{I}$ est filtrante.
Les assertions suivantes sont équivalentes :
\begin{enumerate}
\item $M$ est un ind-objet essentiellement constant ;
\item il existe un cocône $(X, \{M_i \to X\}_i)$ tel que, pour tout $W$ dans
$\mathcal{C}$, l'application
$\colim_i \Mor_\mathcal{C}(W, M_i) \to \Mor_\mathcal{C}(W, X)$
soit bijective ;
\item $X = \colim_i M_i$ existe et, pour tout $W$ dans $\mathcal{C}$,
l'application
$\colim_i \Mor_\mathcal{C}(W, M_i) \to \Mor_\mathcal{C}(W, X)$
est bijective ;
\item il existe un objet $X$ de $\mathcal{C}$, un $i$ dans $\mathcal{I}$ et un
morphisme $X \to M_i$ tels que, pour tout $W$ dans $\mathcal{C}$, l'application
$\Mor_\mathcal{C}(W, X) \to
\colim_{j \in \mathcal{I}} \Mor_\mathcal{C}(W, M_j)$
soit bijective.
\end{enumerate}
Dans les cas (2), (3) et (4), la valeur du système essentiellement constant
est $X$.
\end{lemma}

\begin{proof}
Il est clair que (3) implique (2). Supposons (2). Alors
$\text{id}_X \in \Mor_\mathcal{C}(X, X)$ provient d'un morphisme
$X \to M_i$ pour un certain $i$, c'est-à-dire que $X \to M_i \to X$ est
l'identité. Les deux applications
$$
\Mor_\mathcal{C}(W, X)
\longrightarrow
\colim_i \Mor_\mathcal{C}(W, M_i)
\longrightarrow
\Mor_\mathcal{C}(W, X)
$$
sont alors bijectives pour tout $W$ ; la première est induite par le morphisme
$X \to M_i$ trouvé ci-dessus et leur composée est l'identité. Cela signifie
que la composée
$$
\colim_i \Mor_\mathcal{C}(W, M_i)
\longrightarrow
\Mor_\mathcal{C}(W, X)
\longrightarrow
\colim_i \Mor_\mathcal{C}(W, M_i)
$$
est elle aussi l'identité. En prenant $W = M_j$ et en partant de
$\text{id}_{M_j}$ dans la limite inductive, on voit que
$M_j \to X \to M_i \to M_k$ est égal à $M_j \to M_k$ pour un certain $k$
suffisamment loin dans le système. Cela démontre (1).

\medskip\noindent
Supposons (4). Soit $k$ un objet de $\mathcal{I}$. En prenant $W = M_k$,
on déduit qu'il existe un unique morphisme $M_k \to X$ tel qu'il existe un
$j$ et des morphismes $k \to j$ et $i \to j$ dans $\mathcal{I}$ pour lesquels
$M_k \to X \to M_i \to M_j$ est égal à $M_k \to M_j$. L'unicité garantit
que l'on obtient ainsi un cocône $(X, \{M_k \to X\})$. On voit de cette
manière que (4) implique (2) ; certains détails sont omis.

\medskip\noindent
Nous omettons la preuve que (1) implique les autres assertions.
\end{proof}

\begin{lemma}
\label{categories-lemma-characterize-essentially-constant-pro}
Soit $\mathcal{C}$ une catégorie. Soit $M : \mathcal{I} \to \mathcal{C}$
un diagramme dont la catégorie d'indices $\mathcal{I}$ est cofiltrante.
Les assertions suivantes sont équivalentes :
\begin{enumerate}
\item $M$ est un pro-objet essentiellement constant ;
\item il existe un cône $(X, \{X \to M_i\})$ tel que, pour tout $W$ dans
$\mathcal{C}$, l'application
$\colim_{i \in \mathcal{I}^{opp}} \Mor_\mathcal{C}(M_i, W) \to
\Mor_\mathcal{C}(X, W)$ soit bijective ;
\item $X = \lim_i M_i$ existe et, pour tout $W$ dans $\mathcal{C}$,
l'application
$\colim_{i \in \mathcal{I}^{opp}} \Mor_\mathcal{C}(M_i, W) \to
\Mor_\mathcal{C}(X, W)$ est bijective ;
\item il existe un objet $X$ de $\mathcal{C}$, un $i$ dans $\mathcal{I}$ et un
morphisme $M_i \to X$ tels que, pour tout $W$ dans $\mathcal{C}$, l'application
$\Mor_\mathcal{C}(X, W) \to
\colim_{j \in \mathcal{I}^{opp}} \Mor_\mathcal{C}(M_j, W)$ soit bijective.
\end{enumerate}
Dans les cas (2), (3) et (4), la valeur du système essentiellement constant
est $X$.
\end{lemma}

\begin{proof}
Ce lemme est dual du lemme \ref{categories-lemma-characterize-essentially-constant-ind}.
\end{proof}

\begin{lemma}
\label{categories-lemma-cofinal-essentially-constant}
Soit $\mathcal{C}$ une catégorie. Soit $H : \mathcal{I} \to \mathcal{J}$ un
foncteur entre catégories d'indices filtrantes. Si $H$ est cofinal, alors tout
diagramme $M : \mathcal{J} \to \mathcal{C}$ est essentiellement constant si
et seulement si $M \circ H$ est essentiellement constant.
\end{lemma}

\begin{proof}
Cela résulte formellement des
lemmes \ref{categories-lemma-characterize-essentially-constant-ind} et
\ref{categories-lemma-cofinal}.
\end{proof}

\begin{lemma}
\label{categories-lemma-essentially-constant-over-product}
Soient $\mathcal{I}$ et $\mathcal{J}$ des catégories filtrantes et notons
$p : \mathcal{I} \times \mathcal{J} \to \mathcal{J}$ la projection. Alors
$\mathcal{I} \times \mathcal{J}$ est filtrante, et un diagramme
$M : \mathcal{J} \to \mathcal{C}$ est essentiellement constant si et
seulement si $M \circ p : \mathcal{I} \times \mathcal{J} \to \mathcal{C}$
est essentiellement constant.
\end{lemma}

\begin{proof}
Nous omettons la vérification que $\mathcal{I} \times \mathcal{J}$ est
filtrante. L'équivalence résulte du
lemme \ref{categories-lemma-cofinal-essentially-constant}, car $p$ est cofinal
(vérification omise).
\end{proof}

\begin{lemma}
\label{categories-lemma-initial-essentially-constant}
Soit $\mathcal{C}$ une catégorie. Soit $H : \mathcal{I} \to \mathcal{J}$ un
foncteur entre catégories d'indices cofiltrantes. Si $H$ est initial, alors
tout diagramme $M : \mathcal{J} \to \mathcal{C}$ est essentiellement constant
si et seulement si $M \circ H$ est essentiellement constant.
\end{lemma}

\begin{proof}
Cela résulte formellement des
lemmes \ref{categories-lemma-characterize-essentially-constant-pro},
\ref{categories-lemma-initial}, \ref{categories-lemma-cofinal}, et du fait que, si $\mathcal{I}$ est
initiale dans $\mathcal{J}$, alors $\mathcal{I}^{opp}$ est cofinale dans
$\mathcal{J}^{opp}$.
\end{proof}
\section{Foncteurs exacts}
\label{categories-section-exact-functor}

\noindent
Dans cette section, nous définissons les foncteurs exacts.

\begin{definition}
\label{categories-definition-exact}
Soit $F : \mathcal{A} \to \mathcal{B}$ un foncteur.
\begin{enumerate}
\item Supposons que toutes les limites finies existent dans $\mathcal{A}$.
Nous disons que $F$ est {\it exact à gauche} s'il commute avec toutes
les limites finies.
\item Supposons que toutes les colimites finies existent dans $\mathcal{A}$.
Nous disons que $F$ est {\it exact à droite} s'il commute avec toutes les
colimites finies.
\item Nous disons que $F$ est {\it exact} s'il est exact à la fois à gauche et
à droite.
\end{enumerate}
\end{definition}

\begin{lemma}
\label{categories-lemma-characterize-left-exact}
Soit $F : \mathcal{A} \to \mathcal{B}$ un foncteur. Supposons que toutes les
limites finies existent dans $\mathcal{A}$, voir le lemme
\ref{categories-lemma-finite-limits-exist}. Les assertions suivantes sont équivalentes :
\begin{enumerate}
\item $F$ est exact à gauche,
\item $F$ commute avec des produits finis et des égalisateurs, et
\item $F$ transforme un objet final de $\mathcal{A}$ en un objet final de
$\mathcal{B}$ et commute avec des produits fibrés.
\end{enumerate}
\end{lemma}

\begin{proof}
Le lemme \ref{categories-lemma-limits-products-equalizers} montre que (2) implique (1).
Supposons que (3) soit vrai. Le produit fibré sur l'objet final est le
produit. Si $a, b : A \to B$ sont des morphismes de $\mathcal{A}$, alors
l'égalisateur de $a, b$ est
$$
(A \times_{a, B, b} A)\times_{(\text{pr}_1, \text{pr}_2), A \times A, \Delta} A.
$$
Ainsi (3) implique (2). Enfin (1) implique (3) car la limite vide est un objet
final et les produits fibrés sont des limites.
\end{proof}

\begin{lemma}
\label{categories-lemma-characterize-right-exact}
Soit $F : \mathcal{A} \to \mathcal{B}$ un foncteur. Supposons que toutes les
colimites finies existent dans $\mathcal{A}$, voir le lemme
\ref{categories-lemma-colimits-exist}. Les assertions suivantes sont équivalentes :
\begin{enumerate}
\item $F$ est exact à droite,
\item $F$ commute avec des coproduits finis et des coégalisateurs, et
\item $F$ transforme un objet initial de $\mathcal{A}$ en un objet initial de
$\mathcal{B}$ et commute avec les sommes amalgamées.
\end{enumerate}
\end{lemma}

\begin{proof}
Dual du lemme \ref{categories-lemma-characterize-left-exact}.
\end{proof}




\section{Foncteurs adjoints}
\label{categories-section-adjoint}

\begin{definition}
\label{categories-definition-adjoint}
Soient $\mathcal{C}$, $\mathcal{D}$ des catégories. Soient $u : \mathcal{C} \to \mathcal{D}$ et $v : \mathcal{D} \to \mathcal{C}$ des foncteurs. On dit que
$u$ est un {\it adjoint à gauche} de $v$, ou que $v$ est un {\it adjoint à droite}
de $u$ s'il existe des bijections
$$
\Mor_\mathcal{D}(u(X), Y)
\longrightarrow
\Mor_\mathcal{C}(X, v(Y))
$$
fonctorielles en $X \in \Ob(\mathcal{C})$ et $Y \in \Ob(\mathcal{D})$.
\end{definition}

\noindent
En d'autres termes, cela signifie que l'on s'est {\it donné} un isomorphisme entre
les foncteurs $\mathcal{C}^{opp} \times \mathcal{D} \to \textit{Ensembles}$, allant de
$\Mor_\mathcal{D}(u(-), -)$ à $\Mor_\mathcal{C}(-, v(-))$. Pour tout objet $X$
de $\mathcal{C}$ on obtient un morphisme $X \to v(u(X))$ correspondant à
$\text{id}_{u(X)}$. De même, pour tout objet $Y$ de $\mathcal{D}$ on obtient
un morphisme $u(v(Y)) \to Y$ correspondant à $\text{id}_{v(Y)}$. Ces
morphismes sont appelés {\it morphismes d'adjonction}. Les morphismes
d'adjonction sont fonctoriels en $X$ et $Y$ ; on obtient donc des morphismes
de foncteurs
$$
\eta : \text{id}_\mathcal{C} \to v \circ u\quad (\text{unité})
\quad\text{et}\quad
\epsilon : u \circ v \to \text{id}_\mathcal{D}\quad (\text{counité}).
$$
De plus, si $\alpha : u(X) \to Y$ et $\beta : X \to v(Y)$ sont des morphismes,
alors les conditions suivantes sont équivalentes
\begin{enumerate}
\item $\alpha$ et $\beta$ se correspondent via la bijection de la définition,
\item $\beta$ est la composition $X \to v(u(X)) \xrightarrow{v(\alpha)} v(Y)$,
et
\item $\alpha$ est la composition $u(X) \xrightarrow{u(\beta)} u(v(Y)) \to Y$.
\end{enumerate}
On peut ainsi reformuler la notion de foncteurs adjoints en termes de
morphismes d'adjonction.

\begin{lemma}
\label{categories-lemma-adjoint-exists}
Soit $u : \mathcal{C} \to \mathcal{D}$ un foncteur entre catégories. Si pour
chaque $y \in \Ob(\mathcal{D})$ le foncteur $x \mapsto \Mor_\mathcal{D}(u(x), y)$ est représentable, alors $u$ a un adjoint à droite.
\end{lemma}

\begin{proof}
Pour chaque $y$, choisissons un objet $v(y)$ et un isomorphisme
$\Mor_\mathcal{C}(-, v(y)) \to \Mor_\mathcal{D}(u(-), y)$ de foncteurs. Par le
lemme de Yoneda (lemme \ref{categories-lemma-yoneda}) pour tout morphisme $g : y \to y'$
la transformation de foncteurs
$$
\Mor_\mathcal{C}(-, v(y)) \to \Mor_\mathcal{D}(u(-), y) \to
\Mor_\mathcal{D}(u(-), y') \to \Mor_\mathcal{C}(-, v(y'))
$$
correspond à un morphisme unique $v(g) : v(y) \to v(y')$. Nous omettons de
vérifier que $v$ est un foncteur et qu'il est adjoint à droite de $u$.
\end{proof}

\begin{lemma}
\label{categories-lemma-left-adjoint-composed-fully-faithful}
\begin{reference}
Bhargav Bhatt, communication privée.
\end{reference}
Soit $u$ un adjoint à gauche de $v$ comme dans la définition
\ref{categories-definition-adjoint}.
\begin{enumerate}
\item Si $v \circ u$ est pleinement fidèle, alors $u$ est pleinement fidèle.
\item Si $u \circ v$ est pleinement fidèle, alors $v$ est pleinement fidèle.
\end{enumerate}
\end{lemma}

\begin{proof}
Preuve de (2). Supposons que $u \circ v$ soit pleinement fidèle. Soient
$X$ et $Y$ des objets de $\mathcal{D}$. Alors l'application composite
naturelle
$$
\Mor(X,Y) \to \Mor(v(X),v(Y)) \to \Mor(u(v(X)), u(v(Y)))
$$
est une bijection, donc $v$ est au moins fidèle. Pour montrer qu'il est
pleinement fidèle, il faut montrer que la deuxième application ci-dessus est injective.
Mais l'adjonction entre $u$ et $v$ dit que
$$
\Mor(v(X), v(Y)) \to \Mor(u(v(X)), u(v(Y))) \to \Mor(u(v(X)), Y)
$$
est une bijection, où la première application est naturelle et la deuxième
est induite par la counité $u(v(Y)) \to Y$ de l'adjonction. Cela signifie
donc que $\Mor(v(X), v(Y)) \to \Mor(u(v(X)), u(v(Y)))$ est également injectif,
comme souhaité. La preuve de (1) est duale.
\end{proof}

\begin{lemma}
\label{categories-lemma-adjoint-fully-faithful}
Soit $u$ un adjoint à gauche de $v$ comme dans la définition
\ref{categories-definition-adjoint}. Alors
\begin{enumerate}
\item $u$ est pleinement fidèle $\Leftrightarrow$ $\text{id} \cong v \circ u$
$\Leftrightarrow$ $\eta : \text{id} \to v \circ u$ est un isomorphisme,
\item $v$ est pleinement fidèle $\Leftrightarrow$ $u \circ v \cong \text{id}$
$\Leftrightarrow$ $\epsilon : u \circ v \to \text{id}$ est un isomorphisme.
\end{enumerate}
\end{lemma}

\begin{proof}
Preuve de (1). Supposons que $u$ soit pleinement fidèle. Nous allons montrer
que $\eta_X : X \to v(u(X))$ est un isomorphisme. Considérons un morphisme
quelconque $X' \to v(u(X))$. Par adjonction cela correspond à un morphisme $u(X') \to u(X)$. Par la pleine fidélité de $u$, cela correspond à un morphisme unique $X' \to X$. On voit ainsi que la post-composition par $\eta_X$ définit une
bijection $\Mor(X', X) \to \Mor(X', v(u(X)))$. Donc $\eta_X$ est un
isomorphisme. S'il existe un isomorphisme $\text{id} \cong v \circ u$ de
foncteurs, alors $v \circ u$ est pleinement fidèle. Par le lemme
\ref{categories-lemma-left-adjoint-composed-fully-faithful} nous voyons que $u$ est
pleinement fidèle. D'après ce qui précède, cela implique que $\eta$ est un
isomorphisme. Ainsi, les $3$ conditions sont équivalentes (et ces
conditions sont également équivalentes au fait que $v \circ u$ soit
pleinement fidèle).

\medskip\noindent
La partie (2) est duale de la partie (1).
\end{proof}

\begin{lemma}
\label{categories-lemma-adjoint-exact}
Soit $u$ un adjoint à gauche de $v$ comme dans la définition
\ref{categories-definition-adjoint}.
\begin{enumerate}
\item Supposons que $M : \mathcal{I} \to \mathcal{C}$ soit un diagramme et
supposons que $\colim_\mathcal{I} M$ existe dans $\mathcal{C}$. Alors
$u(\colim_\mathcal{I} M) =
\colim_\mathcal{I} u \circ M$. En d'autres termes,
$u$ commute avec les colimites existantes.
\item Supposons que $M : \mathcal{I} \to \mathcal{D}$ soit un diagramme et
supposons que $\lim_\mathcal{I} M$ existe dans $\mathcal{D}$. Alors
$v(\lim_\mathcal{I} M) =
\lim_\mathcal{I} v \circ M$. En d'autres termes, $v$
commute avec les limites existantes.
\end{enumerate}
\end{lemma}

\begin{proof}
La donnée d'un morphisme d'une colimite vers un objet équivaut à celle d'un
système compatible de morphismes des constituants de la colimite vers l'objet, voir
remarque \ref{categories-remark-limit-colim}. Alors
$$
\begin{matrix}
\Mor_\mathcal{D}(u(\colim_{i \in \mathcal{I}} M_i), Y) &
= & \Mor_\mathcal{C}(\colim_{i \in \mathcal{I}} M_i, v(Y)) \\
& = &
\lim_{i \in \mathcal{I}^{opp}}
\Mor_\mathcal{C}(M_i, v(Y)) \\
& = &
\lim_{i \in \mathcal{I}^{opp}}
\Mor_\mathcal{D}(u(M_i), Y)
\end{matrix}
$$
prouve que $u(\colim_{i \in \mathcal{I}} M_i)$ est la colimite que nous
recherchons. Un argument analogue s'applique au second énoncé.
\end{proof}

\begin{lemma}
\label{categories-lemma-exact-adjoint}
Soit $u$ un adjoint à gauche de $v$ comme dans la définition
\ref{categories-definition-adjoint}.
\begin{enumerate}
\item Si $\mathcal{C}$ a des colimites finies, alors $u$ est exact à droite.
\item Si $\mathcal{D}$ a des limites finies, alors $v$ est exact à gauche.
\end{enumerate}
\end{lemma}

\begin{proof}
Évident d'après les définitions et le lemme \ref{categories-lemma-adjoint-exact}.
\end{proof}

\begin{lemma}
\label{categories-lemma-unit-counit-relations}
Soit $u : \mathcal{C} \to \mathcal{D}$ un adjoint à gauche du foncteur $v : \mathcal{D} \to \mathcal{C}$. Soit $\eta_X : X \to v(u(X))$ l'unité et
$\epsilon_Y : u(v(Y)) \to Y$ la counité. Alors
$$
u(X) \xrightarrow{u(\eta_X)} u(v(u(X)))
\xrightarrow{\epsilon_{u(X)}} u(X)
\quad\text{et}\quad
v(Y) \xrightarrow{\eta_{v(Y)}} v(u(v(Y))) \xrightarrow{v(\epsilon_Y)}
v(Y)
$$
sont les identités.
\end{lemma}

\begin{proof}
Omis.
\end{proof}



\begin{lemma}
\label{categories-lemma-transformation-between-functors-and-adjoints}
Soient $u_1, u_2 : \mathcal{C} \to \mathcal{D}$ des foncteurs avec adjoints
droits $v_1, v_2 : \mathcal{D} \to \mathcal{C}$. Soit $\beta : u_2 \to u_1$
une transformation de foncteurs. Soit $\beta^\vee : v_1 \to v_2$ la
transformation correspondante des foncteurs adjoints. Alors
$$
\xymatrix{
u_2 \circ v_1 \ar[r]_\beta \ar[d]_{\beta^\vee} &
u_1 \circ v_1 \ar[d] \\
u_2 \circ v_2 \ar[r] & \text{id}
}
$$
est commutatif, où les flèches sans étiquette sont les transformations de
counité.
\end{lemma}

\begin{proof}
Cela est vrai car $\beta^\vee_D : v_1D \to v_2D$ est le morphisme unique tel
que l'application induite $\Mor(C, v_1D) \to \Mor(C, v_2D)$ soit
l'application $\Mor(u_1C, D) \to \Mor(u_2C, D)$ induite par $\beta_C : u_2C \to u_1C$. Autrement dit, cela signifie que l'application
$$
\Mor(u_1 v_1 D, D') \to \Mor(u_2 v_1 D, D')
$$
induite par $\beta_{v_1 D}$ est la même que l'application
$$
\Mor(v_1 D, v_1 D') \to \Mor(v_1 D, v_2 D')
$$
induite par $\beta^\vee_{D'}$. En prenant $D' = D$, nous constatons que la counité
$u_1 v_1 D \to D$ précomposée par $\beta_{v_1D}$ correspond à $\beta^\vee_D$
sous adjonction. Cela signifie exactement que le diagramme est commutatif
lorsqu'il est évalué sur $D$.
\end{proof}

\begin{lemma}
\label{categories-lemma-compose-counits}
Soient $\mathcal{A}$, $\mathcal{B}$ et $\mathcal{C}$ des catégories. Soient $v : \mathcal{A} \to \mathcal{B}$ et $v' : \mathcal{B} \to \mathcal{C}$ des
foncteurs ayant pour adjoints à gauche $u$ et $u'$ respectivement. Alors
\begin{enumerate}
\item Le foncteur $v'' = v' \circ v$ a un adjoint à gauche égal à $u'' = u \circ u'$.
\item Étant donné $X$ dans $\mathcal{A}$, nous avons
\begin{equation}
\label{categories-equation-compose-counits}
\epsilon_X^v \circ u(\epsilon^{v'}_{v(X)}) = \epsilon^{v''}_X :
u''(v''(X)) \to X
\end{equation}
où $\epsilon$ est la counité des adjonctions.
\end{enumerate}
\end{lemma}

\begin{proof}
Déroulons la formule de (2) car cela prouvera aussi immédiatement (1).
Premièrement, les counités des adjonctions correspondant aux paires $(u, v)$ et $(u', v')$
sont les morphismes $\epsilon_X^v : u(v(X)) \to X$ et $\epsilon_Y^{v'} : u'(v'(Y)) \to Y$, voir la discussion qui suit la définition
\ref{categories-definition-adjoint}. Avec $u''$ et $v''$ comme en (1) on déroule tout
$$
u''(v''(X)) = u(u'(v'(v(X)))) \xrightarrow{u(\epsilon_{v(X)}^{v'})}
u(v(X)) \xrightarrow{\epsilon_X^v} X
$$
pour obtenir le morphisme du membre de gauche de
(\ref{categories-equation-compose-counits}). Notons-le $\epsilon_X^{v''}$ pour
l'instant. Pour voir qu'il s'agit de la counité d'une paire de foncteurs adjoints $(u'', v'')$
il faut montrer qu'étant donné $Z$ dans $\mathcal{C}$ la règle qui envoie un
morphisme $\beta : Z \to v''(X)$ à $\alpha = \epsilon_X^{v''} \circ u''(\beta) : u''(Z) \to X$ est une bijection entre ensembles de morphismes. C'est vrai
car il s'agit de la composition de la règle envoyant $\beta$ à
$\epsilon_{v(X)}^{v'} \circ u'(\beta)$ qui est une bijection par hypothèse sur
$(u', v')$ puis envoyant celle-ci à $\epsilon_X^v \circ u(\epsilon_{v(X)}^{v'} \circ u'(\beta))$ qui est une bijection par hypothèse sur $(u, v)$.
\end{proof}






\section{Un critère de représentabilité}
\label{categories-section-representable}

\noindent
Le lemme suivant est souvent utile pour prouver l'existence d'objets
universels dans de grandes catégories ; voir la discussion dans
la remarque \ref{categories-remark-how-to-use-it}.

\begin{lemma}
\label{categories-lemma-a-version-of-brown}
Soit $\mathcal{C}$ une grande\footnote{Voir remarque
\ref{categories-remark-big-categories}.} catégorie qui a des limites. Soit $F : \mathcal{C} \to \textit{Ensembles}$ un foncteur. Supposons que
\begin{enumerate}
\item $F$ commute avec les limites,
\item il existe une famille $\{x_i\}_{i \in I}$ d'objets de $\mathcal{C}$ et
pour chaque $i \in I$ un élément $f_i \in F(x_i)$ tel que pour $y \in \Ob(\mathcal{C})$ et $g \in F(y)$ il existe un $i$ et un morphisme $\varphi : x_i \to y$ avec $F(\varphi)(f_i) = g$.
\end{enumerate}
Alors $F$ est représentable, c'est-à-dire qu'il existe un objet $x$ de
$\mathcal{C}$ tel que
$$
F(y) = \Mor_\mathcal{C}(x, y)
$$
fonctoriellement en $y$.
\end{lemma}

\begin{proof}
Soit $\mathcal{I}$ la catégorie dont les objets sont les couples $(x_i, f_i)$
et dont les morphismes $(x_i, f_i) \to (x_{i'}, f_{i'})$ sont des morphismes
$\varphi : x_i \to x_{i'}$ dans $\mathcal{C}$ tels que $F(\varphi)(f_i) = f_{i'}$. Posons
$$
x = \lim_{(x_i, f_i) \in \mathcal{I}} x_i
$$
(ce ne sera pas le $x$ que nous recherchons, voir ci-dessous). La limite
existe par hypothèse. Comme $F$ commute avec les limites, nous avons
$$
F(x) = \lim_{(x_i, f_i) \in \mathcal{I}} F(x_i).
$$
Il existe donc un élément universel $f \in F(x)$ dont l'image est
$f_i \in F(x_i)$ lorsque l'on applique $F$ au morphisme de projection
$x \to x_i$. En
utilisant $f$ on obtient une transformation de foncteurs
$$
\xi : \Mor_\mathcal{C}(x, - ) \longrightarrow F(-)
$$
voir la section \ref{categories-section-opposite}. Soit $y$ un objet quelconque de
$\mathcal{C}$ et soit $g \in F(y)$. Choisissons $x_i \to y$ de telle sorte que
$f_i$ corresponde à $g$, ce qui est possible par hypothèse. Alors $F$ appliqué
aux morphismes
$$
x \longrightarrow x_i \longrightarrow y
$$
(le premier étant le morphisme de projection de la limite définissant $x$)
envoie $f$ sur $g$. La transformation $\xi$ est donc surjective.

\medskip\noindent
Afin de trouver l'objet représentant $F$, prenons pour $e : x' \to x$
l'égalisateur de tous les endomorphismes $\varphi : x \to x$ avec
$F(\varphi)(f) = f$. Puisque $F$ commute avec les limites, il commute avec les
égalisateurs, et nous voyons qu'il existe un élément $f' \in F(x')$ dont l'image
est $f$ dans $F(x)$. Puisque $\xi$ est surjective et que $f'$
correspond à $f$, nous voyons que $\xi' : \Mor_\mathcal{C}(x', -) \to F(-)$
est également surjective. Enfin, supposons que $a, b : x' \to y$ soient deux
morphismes tels que $F(a)(f') = F(b)(f')$. Nous devons montrer $a = b$.
Considérons l'égalisateur $e' : x'' \to x'$. De nouveau, nous trouvons un
$f'' \in F(x'')$ dont l'image est $f'$. Choisissons un morphisme $\psi : x \to x''$ tel que $F(\psi)(f) = f''$. On voit alors que $e \circ e' \circ \psi : x \to x$ est un morphisme avec $F(e \circ e' \circ \psi)(f) = f$. D'où $e \circ e' \circ \psi \circ e = e$. Puisque $e$ est un monomorphisme, cela
implique que $e'$ est un épimorphisme, donc $a = b$ comme souhaité.
\end{proof}

\begin{remark}
\label{categories-remark-how-to-use-it}
Le lemme ci-dessus est souvent utilisé pour construire quelque chose de libre
sur quelque chose. Par exemple le groupe abélien libre sur un ensemble, le
groupe libre sur un ensemble, etc. L'idée, disons dans le cas du groupe libre
sur un ensemble $E$, est de considérer le foncteur
$$
F : \textit{Groups} \to \textit{Ensembles},\quad
G \longmapsto \text{Map}(E, G)
$$
Ce foncteur commute avec les limites. Comme famille d'objets, nous pouvons
prendre une famille de structures $E \to G_i$ constituées de groupes $G_i$
de cardinalité au plus $\max(\aleph_0, |E|)$ et d'applications ensemblistes
$E \to G_i$, de telle sorte que chaque classe d'isomorphisme d'une telle
structure soit représentée
au moins une fois. En effet, si $E \to G$ est une application de $E$ vers un
groupe $G$, alors le sous-groupe $G'$ généré par l'image a une cardinalité au
plus $\max(\aleph_0, |E|)$. Le lemme nous dit que le foncteur est
représentable, il existe donc un groupe $F_E$ tel que
$\Mor_{\textit{Groups}}(F_E, G) = \text{Map}(E, G)$. En particulier, le
morphisme identité de $F_E$ correspond à une application $E \to F_E$ et on
peut montrer que $F_E$ est généré par l'image sans imposer de relations.

\medskip\noindent
Une autre application typique est que nous pouvons utiliser le lemme pour
construire des colimites dès que l'on sait que des limites existent. Nous
l'illustrons à l'aide de la catégorie des espaces topologiques, qui admet des
limites d'après Topologie, Lemme \ref{topology-lemma-limits}. Supposons
que $\mathcal{I} \to \textit{Top}$, $i \mapsto X_i$ soit un foncteur. On peut
alors considérer
$$
F : \textit{Top} \longrightarrow \textit{Ensembles},\quad
Y \longmapsto \lim_\mathcal{I} \Mor_{\textit{Top}}(X_i, Y)
$$
Ce foncteur commute avec les limites. De plus, étant donné tout espace
topologique $Y$ et un élément $(\varphi_i : X_i \to Y)$ de $F(Y)$, il existe
un sous-espace $Y' \subset Y$ de cardinalité au plus $|\coprod X_i|$ tel que
les morphismes $\varphi_i$ se factorisent par $Y'$. En effet, nous pouvons prendre
la topologie induite sur l'union des images des $\varphi_i$. Ainsi il est clair
que les hypothèses du lemme sont satisfaites et on retrouve un espace
topologique $X$ représentant le foncteur $F$, ce qui signifie précisément que
$X$ est la colimite du diagramme $i \mapsto X_i$.
\end{remark}

\begin{theorem}[Théorème du foncteur adjoint]
\label{categories-theorem-adjoint-functor}
Soit $G : \mathcal{C} \to \mathcal{D}$ un foncteur entre grandes catégories.
Supposons que $\mathcal{C}$ admette des limites, que $G$ commute avec elles et que,
pour chaque objet $y$ de $\mathcal{D}$, il existe un ensemble de paires $(x_i, f_i)_{i \in I}$ avec $x_i \in \Ob(\mathcal{C})$, $f_i \in \Mor_\mathcal{D}(y, G(x_i))$ tel que pour toute paire $(x, f)$ avec $x \in \Ob(\mathcal{C})$, $f \in \Mor_\mathcal{D}(y, G(x))$ il existe un $i$ et un morphisme $h : x_i \to x$ tel que $f = G(h) \circ f_i$. Alors $G$ a un adjoint à gauche $F$.
\end{theorem}

\begin{proof}
Les hypothèses impliquent que pour chaque objet $y$ de $\mathcal{D}$ le
foncteur $x \mapsto \Mor_\mathcal{D}(y, G(x))$ satisfait les hypothèses du
lemme \ref{categories-lemma-a-version-of-brown}. Il est donc représentable par un objet,
appelons-le $F(y)$. Une application du lemme de Yoneda (Lemme
\ref{categories-lemma-yoneda}) permet de faire de la règle $y \mapsto F(y)$ un foncteur qui
par construction est un adjoint à gauche de $G$. Nous omettons les détails.
\end{proof}





\section{Objets catégoriquement compacts}
\label{categories-section-compact}

\noindent
Un peu sur les « petits » objets d'une catégorie.

\begin{definition}
\label{categories-definition-compact-object}
Soit $\mathcal{C}$ une grande catégorie \footnote{Voir la remarque
\ref{categories-remark-big-categories}.}. Un objet $X$ de $\mathcal{C}$ est dit
{\it catégoriquement compact} si on a
$$
\Mor_\mathcal{C}(X, \colim_i M_i) =
\colim_i \Mor_\mathcal{C}(X, M_i)
$$
pour chaque diagramme filtrant $M : \mathcal{I} \to \mathcal{C}$ tel que
$\colim_i M_i$ existe.
\end{definition}

\noindent
Souvent, cette définition n'est donnée que sous l'hypothèse que $\mathcal{C}$
admet toutes les limites inductives filtrantes.

\begin{lemma}
\label{categories-lemma-extend-functor-by-colim}
Soient $\mathcal{C}$ et $\mathcal{D}$ deux grandes catégories admettant des limites
inductives filtrantes. Soit $\mathcal{C}' \subset \mathcal{C}$ une petite sous-catégorie
pleine composée d'objets catégoriquement compacts de $\mathcal{C}$, telle
que chaque objet de $\mathcal{C}$ soit une limite inductive filtrante d'objets
de $\mathcal{C}'$. Alors chaque foncteur $F' : \mathcal{C}' \to \mathcal{D}$ a une extension unique $F : \mathcal{C} \to \mathcal{D}$ qui commute
avec les limites inductives filtrantes.
\end{lemma}

\begin{proof}
Pour chaque objet $X$ de $\mathcal{C}$, nous pouvons écrire $X$ comme limite
inductive filtrante $X = \colim X_i$ avec $X_i \in \Ob(\mathcal{C}')$. Posons
$$
F(X) = \colim F'(X_i)
$$
dans $\mathcal{D}$. Nous montrerons ci-dessous que cette construction ne
dépend pas du choix d'une telle présentation de $X$.

\medskip\noindent
Soit $\alpha : X \to Y$ un morphisme de $\mathcal{C}$, et supposons que $X = \colim_{i \in I} X_i$ et $Y = \colim_{j \in J} Y_j$ soient écrits comme
limites inductives filtrantes d'objets de $\mathcal{C}'$. Pour chaque $i \in I$, puisque
$X_i$ est un objet catégoriquement compact de $\mathcal{C}$, nous pouvons
trouver un $j \in J$ et un diagramme commutatif
$$
\xymatrix{
X_i \ar[r] \ar[d] & X \ar[d]^\alpha \\
Y_j \ar[r] & Y
}
$$
On obtient alors un morphisme $F'(X_i) \to F'(Y_j) \to F(Y)$ où le deuxième
morphisme est la coprojection dans $F(Y) = \colim F'(Y_j)$. La flèche $\beta_i : F'(X_i) \to F(Y)$ ne dépend pas du choix de $j$. Pour $i \leq i'$, la
composition
$$
F'(X_i) \to F'(X_{i'}) \xrightarrow{\beta_{i'}} F(Y)
$$
est égale à $\beta_i$. On obtient ainsi une flèche bien définie
$$
F(\alpha) : F(X) = \colim F(X_i) \to F(Y)
$$
par la propriété universelle de la colimite. Si $\alpha' : Y \to Z$ est un
deuxième morphisme de $\mathcal{C}$ et que $Z = \colim Z_k$ s'écrit également
comme limite inductive filtrante d'objets de $\mathcal{C}'$, alors c'est un exercice
agréable de montrer que les morphismes induits $F(\alpha) : F(X) \to F(Y)$ et
$F(\alpha') : F(Y) \to F(Z)$ ont pour composé le morphisme $F(\alpha' \circ \alpha)$.
Détails omis.

\medskip\noindent
En particulier, si on nous donne deux présentations $X = \colim X_i$ et $X = \colim X'_{i'}$ comme limites inductives filtrantes de systèmes dans $\mathcal{C}'$,
alors nous obtenons des flèches mutuellement inverses $\colim F'(X_i) \to \colim F'(X'_{i'})$ et $\colim F'(X'_{i'}) \to \colim F'(X_i)$. Autrement dit,
la valeur $F(X)$ est bien définie indépendamment du choix de la présentation
de $X$ comme limite inductive filtrante d'objets de $\mathcal{C}'$. Avec la
fonctorialité de $F$ discutée dans le paragraphe précédent, nous constatons
que $F$ est un foncteur. En outre, il est clair que $F(X) = F'(X)$ si $X \in \Ob(\mathcal{C}')$.

\medskip\noindent
L'énoncé d'unicité du lemme est clair, à condition de montrer que
$F$ commute avec les limites inductives filtrantes (car cet énoncé n'a
aucun sens autrement). Pour cela, supposons que $X = \colim_{\lambda \in \Lambda} X_\lambda$ soit une limite inductive filtrante dans $\mathcal{C}$. Puisque
$F$ est un foncteur, nous obtenons un morphisme naturel
$$
\colim_\lambda F(X_\lambda) \longrightarrow F(X)
$$
D'autre part, écrivons $X = \colim X_i$ comme limite inductive filtrante d'objets de
$\mathcal{C}'$. Comme ci-dessus, pour chaque $i \in I$ on peut choisir un
$\lambda \in \Lambda$ et un diagramme commutatif
$$
\xymatrix{
X_i \ar[rr] \ar[rd] & & X_\lambda \ar[ld] \\
& X
}
$$
Comme ci-dessus cela détermine un morphisme bien défini $F'(X_i) \to \colim_\lambda F(X_\lambda)$ compatible avec les morphismes de transition et
donc un morphisme
$$
F(X) = \colim_i F'(X_i) \longrightarrow \colim_\lambda F(X_\lambda)
$$
Ce morphisme est inverse du morphisme ci-dessus (détails omis) et prouve que
$F(X) = \colim_\lambda F(X_\lambda)$ comme souhaité.
\end{proof}







\section{Localisation dans les catégories}
\label{categories-section-localization}

\noindent
L'idée de base de cette section est la suivante : étant donnés une catégorie $\mathcal{C}$ et
un ensemble de flèches $S$, construire un foncteur $F : \mathcal{C} \to S^{-1}\mathcal{C}$ tel que tous les éléments de $S$ deviennent inversibles
dans $S^{-1}\mathcal{C}$ et tel que $F$ soit universel parmi tous les foncteurs
possédant cette propriété. Les références pour cette section sont
\cite[chapitre I, section 2]{GZ} et \cite[chapitre II, section 2]{Verdier}.

\begin{definition}
\label{categories-definition-multiplicative-system}
Soit $\mathcal{C}$ une catégorie. Un ensemble de flèches $S$ de $\mathcal{C}$
est appelé un {\it système multiplicatif à gauche} s'il possède les propriétés
suivantes :
\begin{enumerate}
\item[LMS1] L'identité de chaque objet de $\mathcal{C}$ est dans $S$ et la
composition de deux éléments composables de $S$ est dans $S$.
\item[LMS2] Tout diagramme en traits pleins
$$
\xymatrix{
X \ar[d]_t \ar[r]_g & Y \ar@{..>}[d]^s \\
Z \ar@{..>}[r]^f & W
}
$$
avec $t \in S$ peut être complété en un carré commutatif par des flèches pointillées avec $s \in S$.
\item[LMS3] Pour chaque paire de morphismes $f, g : X \to Y$ et $t \in S$ de
but $X$ tel que $f \circ t = g \circ t$, il existe un $s \in S$ de
source $Y$ tel que $s \circ f = s \circ g$.
\end{enumerate}
Un ensemble de flèches $S$ de $\mathcal{C}$ est appelé un {\it système
multiplicatif à droite} s'il possède les propriétés suivantes :
\begin{enumerate}
\item[RMS1] L'identité de chaque objet de $\mathcal{C}$ est dans $S$ et la
composition de deux éléments composables de $S$ est dans $S$.
\item[RMS2] Tout diagramme en traits pleins
$$
\xymatrix{
X \ar@{..>}[d]_t \ar@{..>}[r]_g & Y \ar[d]^s \\
Z \ar[r]^f & W
}
$$
avec $s \in S$ peut être complété en un carré commutatif par des flèches pointillées avec $t \in S$.
\item[RMS3] Pour chaque paire de morphismes $f, g : X \to Y$ et $s \in S$ de
source $Y$ tel que $s \circ f = s \circ g$, il existe un $t \in S$ de
but $X$ tel que $f \circ t = g \circ t$.
\end{enumerate}
Un ensemble de flèches $S$ de $\mathcal{C}$ est appelé un {\it système
multiplicatif} s'il s'agit à la fois d'un système multiplicatif à gauche et d'un
système multiplicatif à droite. En d'autres termes, cela signifie que MS1, MS2,
MS3 sont vérifiées, où MS1 $=$ LMS1 $+$ RMS1, MS2 $=$ LMS2 $+$ RMS2 et MS3 $=$
LMS3 $+$ RMS3. (Cela dit, bien sûr LMS1 $=$ RMS1 $=$ MS1.)
\end{definition}

\noindent
Ces conditions sont utiles pour construire les catégories $S^{-1}\mathcal{C}$
comme suit.

\medskip\noindent
{\bf Calcul des fractions à gauche.} Soit $\mathcal{C}$ une catégorie et $S$
un système multiplicatif à gauche. Nous définissons une nouvelle catégorie
$S^{-1}\mathcal{C}$ comme suit (nous vérifions que cela fonctionne dans la
preuve du lemme \ref{categories-lemma-left-localization}) :
\begin{enumerate}
\item Nous définissons $\Ob(S^{-1}\mathcal{C}) = \Ob(\mathcal{C})$.
\item Les morphismes $X \to Y$ de $S^{-1}\mathcal{C}$ sont représentés par des paires
$(f : X \to Y', s : Y \to Y')$ avec $s \in S$ à équivalence près.
(L'équivalence est définie ci-dessous. Considérons la classe d'équivalence d'une
paire $(f, s)$ comme $s^{-1}f : X \to Y$.)
\item Deux paires $(f_1 : X \to Y_1, s_1 : Y \to Y_1)$ et $(f_2 : X \to Y_2, s_2 : Y \to Y_2)$ sont dites équivalentes s'il existe une troisième paire
$(f_3 : X \to Y_3, s_3 : Y \to Y_3)$ ainsi que des morphismes $u : Y_1 \to Y_3$ et $v : Y_2 \to Y_3$ de $\mathcal{C}$ qui s'insèrent dans le diagramme commutatif
$$
\xymatrix{
 & Y_1 \ar[d]^u & \\
X \ar[ru]^{f_1} \ar[r]^{f_3} \ar[rd]_{f_2} &
Y_3 &
Y \ar[lu]_{s_1} \ar[l]_{s_3} \ar[ld]^{s_2} \\
& Y_2 \ar[u]_v &
}
$$
\item La composition des classes d'équivalence des couples $(f : X \to Y', s : Y \to Y')$ et $(g : Y \to Z', t : Z \to Z')$ est définie comme la classe
d'équivalence d'un couple $(h \circ f : X \to Z'', u \circ t : Z \to Z'')$ où
$h$ et $u \in S$ sont choisis pour s'insérer dans un diagramme commutatif
$$
\xymatrix{
Y \ar[d]_s \ar[r]_g & Z' \ar[d]^u \\
Y' \ar[r]^h & Z''
}
$$
qui existe par hypothèse.
\item Le morphisme identité $X \to X$ dans $S^{-1} \mathcal{C}$ est la
classe d'équivalence de la paire $(\text{id} : X \to X,
\text{id} : X \to X)$.
\end{enumerate}

\begin{lemma}
\label{categories-lemma-left-localization}
Soit $\mathcal{C}$ une catégorie et $S$ un système multiplicatif à gauche.
\begin{enumerate}
\item La relation sur les paires définie ci-dessus est une relation d'équivalence.
\item La règle de composition donnée ci-dessus est bien définie sur les
classes d'équivalence.
\item La composition est associative (et les morphismes identité satisfont
les axiomes d'identité), et donc $S^{-1}\mathcal{C}$ est une catégorie.
\end{enumerate}
\end{lemma}

\begin{proof}
Preuve de (1). Disons que deux paires $p_1 = (f_1 : X \to Y_1, s_1 : Y \to Y_1)$ et $p_2 = (f_2 : X \to Y_2, s_2 : Y \to Y_2)$ sont élémentairement
équivalentes s'il existe un morphisme $a : Y_1 \to Y_2$ de $\mathcal{C}$ tel
que $a \circ f_1 = f_2$ et $a \circ s_1 = s_2$. Le diagramme :
$$
\xymatrix{
X \ar@{=}[d] \ar[r]_{f_1} & Y_1 \ar[d]^a & Y \ar[l]^{s_1} \ar@{=}[d] \\
X \ar[r]^{f_2} & Y_2 & Y \ar[l]_{s_2}
}
$$
Notons cette propriété en disant $p_1Ep_2$. Notez que $pEp$ et $aEb, bEc \Rightarrow aEc$. (Malgré son nom, $E$ n'est pas une relation d'équivalence.)
La partie (1) affirme que la relation $p \sim p' \Leftrightarrow \exists q: pEq \wedge p'Eq$ (où $q$ est une paire satisfaisant les mêmes
conditions que $p$ et $p'$) est une relation d'équivalence. Un simple argument
formel, utilisant les propriétés de $E$ ci-dessus, montre qu'il suffit de
prouver $p_3Ep_1, p_3Ep_2 \Rightarrow p_1 \sim p_2$. Supposons donc qu'on nous
donne un diagramme commutatif
$$
\xymatrix{
 & Y_1 & \\
X \ar[ru]^{f_1} \ar[r]^{f_3} \ar[rd]_{f_2} &
Y_3 \ar[u]_{a_{31}} \ar[d]^{a_{32}} &
Y \ar[lu]_{s_1} \ar[l]_{s_3} \ar[ld]^{s_2} \\
& Y_2 &
}
$$
avec $s_i \in S$. Nous appliquons d'abord LMS2 pour obtenir un diagramme
commutatif
$$
\xymatrix{
Y \ar[d]_{s_1} \ar[r]_{s_2} & Y_2 \ar@{..>}[d]^{a_{24}} \\
Y_1 \ar@{..>}[r]^{a_{14}} & Y_4
}
$$
avec $a_{24} \in S$. Ensuite, nous avons
$$
a_{14} \circ a_{31} \circ s_3 =
a_{14} \circ s_1 =
a_{24} \circ s_2 =
a_{24} \circ a_{32} \circ s_3.
$$
Ainsi, par LMS3, il existe un morphisme $s_{44} : Y_4 \to Y'_4$ tel que
$s_{44} \in S$ et $s_{44} \circ a_{14} \circ a_{31}
= s_{44} \circ a_{24} \circ a_{32}$. Ainsi, après avoir remplacé $Y_4$, $a_{14}$ et $a_{24}$ par
$Y'_4$, $s_{44} \circ a_{14}$ et $s_{44} \circ a_{24}$, nous pouvons supposer
que $a_{14} \circ a_{31} = a_{24} \circ a_{32}$ (et nous avons toujours
$a_{24} \in S$ et $a_{14} \circ s_1 = a_{24} \circ s_2$). Posons
$$
f_4 =
a_{14} \circ f_1 =
a_{14} \circ a_{31} \circ f_3 =
a_{24} \circ a_{32} \circ f_3 =
a_{24} \circ f_2
$$
et $s_4 = a_{14} \circ s_1 = a_{24} \circ s_2$. Ensuite, le diagramme
$$
\xymatrix{
X \ar@{=}[d] \ar[r]_{f_1} & Y_1 \ar[d]^{a_{14}} & Y \ar[l]^{s_1} \ar@{=}[d] \\
X \ar[r]^{f_4} & Y_4 & Y \ar[l]_{s_4}
}
$$
est commutatif, et nous avons $s_4 \in S$ (par LMS1). Ainsi, $p_1 E p_4$, où
$p_4 = (f_4, s_4)$. De même, $p_2 E p_4$. Ces deux faits donnent $p_1 \sim p_2$.

\medskip\noindent
Preuve de (2). Soient $p = (f : X \to Y', s : Y \to Y')$ et $q = (g : Y \to Z', t : Z \to Z')$ des paires comme dans la définition de la composition
ci-dessus. Pour composer on choisit un diagramme
$$
\xymatrix{
Y \ar[d]_s \ar[r]_g & Z' \ar[d]^{u_2} \\
Y' \ar[r]^{h_2} & Z_2
}
$$
avec $u_2 \in S$. Nous montrons d'abord que la classe d'équivalence du couple
$r_2 = (h_2 \circ f : X \to Z_2, u_2 \circ t : Z \to Z_2)$ est indépendante du
choix de $(Z_2, h_2, u_2)$. Supposons en effet que $(Z_3, h_3, u_3)$ soit un
autre choix avec la composition correspondante $r_3 = (h_3 \circ f : X \to Z_3, u_3 \circ t : Z \to Z_3)$. Alors, par LMS2 on peut choisir un diagramme
$$
\xymatrix{
Z' \ar[d]_{u_2} \ar[r]_{u_3} & Z_3 \ar[d]^{u_{34}} \\
Z_2 \ar[r]^{h_{24}} & Z_4
}
$$
avec $u_{34} \in S$. Nous avons $h_2 \circ s = u_2 \circ g$ et de même $h_3 \circ s = u_3 \circ g$. De plus,
$$
u_{34} \circ h_3 \circ s
= u_{34} \circ u_3 \circ g
= h_{24} \circ u_2 \circ g
= h_{24} \circ h_2 \circ s.
$$
Ainsi, LMS3 montre qu'il existe un $Z'_4$ et un $s_{44} : Z_4 \to Z'_4$ tel
que $s_{44} \circ u_{34} \circ h_3 = s_{44} \circ h_{24} \circ h_2$. En
remplaçant $Z_4$, $h_{24}$ et $u_{34}$ par $Z'_4$, $s_{44} \circ h_{24}$ et
$s_{44} \circ u_{34}$, nous pouvons supposer que $u_{34} \circ h_3 = h_{24} \circ h_2$. Pendant ce temps, les relations $u_{34} \circ u_3 = h_{24} \circ u_2$ et $u_{34} \in S$ restent satisfaites. Nous pouvons maintenant définir
$h_4 = u_{34} \circ h_3 = h_{24} \circ h_2$ et $u_4 = u_{34} \circ u_3 = h_{24} \circ u_2$. On a alors un diagramme commutatif
$$
\xymatrix{
X \ar@{=}[d] \ar[r]_{h_2\circ f} &
Z_2 \ar[d]^{h_{24}} &
Z \ar[l]^{u_2 \circ t} \ar@{=}[d] \\
X \ar@{=}[d] \ar[r]^{h_4\circ f} &
Z_4 &
Z \ar@{=}[d] \ar[l]_{u_4 \circ t} \\
X \ar[r]^{h_3 \circ f} &
Z_3 \ar[u]^{u_{34}} &
Z \ar[l]_{u_3 \circ t}
}
$$
Nous obtenons donc une paire $r_4 =
(h_4 \circ f : X \to Z_4, u_4 \circ t : Z \to Z_4)$ et le diagramme ci-dessus montre que nous avons $r_2Er_4$ et
$r_3Er_4$, d'où $r_2 \sim r_3$, comme souhaité. Il est donc désormais légitime
de définir $p \circ q$ comme la classe d'équivalence de toutes les paires
possibles $r$ obtenues comme ci-dessus.

\medskip\noindent
Pour terminer la preuve de (2), nous devons montrer que, pour des paires
$p_1, p_2, q$ telles que $p_1Ep_2$, on a $p_1 \circ q = p_2 \circ q$ et $q \circ p_1 = q \circ p_2$ chaque fois que les compositions ont un sens. Pour ce
faire, écrivons $p_1 = (f_1 : X \to Y_1, s_1 : Y \to Y_1)$ et $p_2 = (f_2 : X \to Y_2, s_2 : Y \to Y_2)$, et soit $a : Y_1 \to Y_2$ un morphisme de
$\mathcal{C}$ tel que $f_2 = a \circ f_1$ et $s_2 = a \circ s_1$. Supposons
d'abord que $q = (g : Y \to Z', t : Z \to Z')$. Dans ce cas, choisissez un
diagramme commutatif comme celui de gauche
$$
\vcenter{
\xymatrix{
Y \ar[d]_{s_2} \ar[r]^g & Z' \ar[d]^u \\
Y_2 \ar[r]^h & Z''
}
}
\quad
\Rightarrow
\quad
\vcenter{
\xymatrix{
Y \ar[d]_{s_1} \ar[r]^g & Z' \ar[d]^u \\
Y_1 \ar[r]^{h \circ a} & Z''
}
}
$$
(avec $u \in S$), ce qui implique que le diagramme de droite est également
commutatif. À l'aide de ces diagrammes, nous voyons que les deux compositions
$q \circ p_1$ et $q \circ p_2$ représentent la classe d'équivalence de $(h \circ a \circ f_1 : X \to Z'', u \circ t : Z \to Z'')$. Ainsi $q \circ p_1 = q \circ p_2$. La preuve de l'autre cas, dans lequel nous devons montrer $p_1 \circ q = p_2 \circ q$, est omise. (C'est similaire au cas que nous avons
traité.)

\medskip\noindent
Preuve de (3). Nous devons prouver l'associativité de la composition.
Considérons un diagramme en traits pleins
$$
\xymatrix{
& & & Z \ar[d] \\
& & Y \ar[d] \ar[r] & Z' \ar@{..>}[d] \\
& X \ar[d] \ar[r] & Y' \ar@{..>}[d] \ar@{..>}[r] & Z'' \ar@{..>}[d] \\
W \ar[r] & X' \ar@{..>}[r] & Y'' \ar@{..>}[r] & Z'''
}
$$
(dont les flèches verticales appartiennent à $S$) qui détermine trois
paires composables. En utilisant LMS2 on peut choisir les flèches pointillées
qui rendent les carrés commutatifs et telles que les flèches verticales soient dans
$S$. Il est alors clair que la composition des trois paires est la classe
d'équivalence de la paire $(W \to Z''', Z \to Z''')$ obtenue en composant les
flèches horizontales sur la ligne du bas et les flèches verticales sur la
colonne de droite.

\medskip\noindent
Nous laissons au lecteur le soin de vérifier les axiomes d'identité.
\end{proof}

\begin{remark}
\label{categories-remark-motivation-localization}
La motivation pour la construction de $S^{-1} \mathcal{C}$ est de « forcer »
les morphismes de $S$ à être inversibles en créant artificiellement des
inverses (au prix de certains morphismes existants qui peuvent s'identifier les
uns aux autres). Ceci est analogue à la localisation d'un anneau commutatif
par rapport à une partie multiplicative, et plus généralement à la
localisation d'un anneau non commutatif par rapport à un ensemble de
dénominateurs à droite (voir \cite[section 10A]{Lam}). C'est plus qu'une simple
similitude : la construction de $S^{-1} \mathcal{C}$ (ou, plus précisément, sa
version pour les catégories additives $\mathcal{C}$) généralise en fait ce
dernier type de localisation. En effet, un anneau non commutatif peut être
considéré comme une catégorie préadditive avec un seul objet (les morphismes
étant les éléments de l'anneau) ; une partie multiplicative de cet anneau
devient alors un ensemble $S$ de morphismes satisfaisant LMS1 (alias RMS1).
Ensuite, les conditions RMS2 et RMS3 pour cette catégorie et cet ensemble
$S$ se traduisent par les deux conditions (« permutable à droite » et
« réversible à droite ») d'un ensemble de dénominateurs à droite (et de même
pour les ensembles de dénominateurs à gauche et LMS), et $S^{-1} \mathcal{C}$
(munie d'une structure additive convenablement définie) est la catégorie à un objet
correspondant à la localisation de l'anneau.
\end{remark}

\begin{definition}
\label{categories-definition-left-localization-as-fraction}
Soit $\mathcal{C}$ une catégorie et soit $S$ un système multiplicatif à gauche
de morphismes de $\mathcal{C}$. Étant donné tout morphisme $f : X \to Y'$ dans
$\mathcal{C}$ et tout morphisme $s : Y \to Y'$ dans $S$, on note {\it $s^{-1} f$} la classe d'équivalence de la paire $(f : X \to Y', s : Y \to Y')$. Il
s'agit d'un morphisme de $X$ à $Y$ dans $S^{-1} \mathcal{C}$.
\end{definition}

\noindent
Cette notation est suggestive, et les choses qu'elle suggère sont vraies :
étant donnés un morphisme $f : X \to Y'$ dans $\mathcal{C}$ et deux
morphismes quelconques $s : Y \to Y'$ et $t : Y' \to Y''$ dans $S$, nous avons
$\left(t \circ s\right)^{-1} \left(t \circ f\right) = s^{-1} f$. De plus, pour
tous les $f : X \to Y'$ et $g : Y' \to Z'$ dans $\mathcal{C}$ et tous les $s : Z \to Z'$ dans $S$, nous avons $s^{-1} \left(g \circ f\right) = \left(s^{-1} g\right) \circ
\left(\text{id}_{Y'}^{-1} f\right)$. Enfin, pour tout $f : X \to Y'$ dans $\mathcal{C}$, tous les $s : Y \to Y'$ dans $S$ et $t : Z \to Y$
dans $S$, nous avons $\left(s \circ t\right)^{-1} f
= \left(t^{-1} \text{id}_Y\right)
\circ \left(s^{-1} f\right)$. Tout cela ressort clairement
de la définition. Nous pouvons « écrire toute famille finie de
morphismes de même but sous forme de fractions à dénominateur commun ».

\begin{lemma}
\label{categories-lemma-morphisms-left-localization}
Soit $\mathcal{C}$ une catégorie et soit $S$ un système multiplicatif à gauche
de morphismes de $\mathcal{C}$. Étant donnée toute famille finie $g_i : X_i \to Y$ de morphismes de $S^{-1}\mathcal{C}$ (indexée par $i$), on peut trouver
un élément $s : Y \to Y'$ de $S$ et une famille de morphismes $f_i : X_i \to Y'$ de $\mathcal{C}$ telle que chaque $g_i$ soit la classe d'équivalence de la
paire $(f_i : X_i \to Y', s : Y \to Y')$.
\end{lemma}

\begin{proof}
Pour chaque $i$, choisissons un représentant $(X_i \to Y_i, s_i : Y \to Y_i)$
de $g_i$. Il suffit de trouver un morphisme $s : Y \to Y'$ dans
$S$ tel que pour chaque $i$ il existe un morphisme $a_i : Y_i \to Y'$ avec
$a_i \circ s_i = s$. Si nous avons deux indices $i = 1, 2$, alors nous pouvons
le faire en complétant le carré
$$
\xymatrix{
Y \ar[d]_{s_1} \ar[r]_{s_2} & Y_2 \ar[d]^{t_2} \\
Y_1 \ar[r]^{a_1} & Y'
}
$$
avec $t_2 \in S$ ce qui est possible par la définition
\ref{categories-definition-multiplicative-system}. Ensuite, $s = t_2 \circ s_2 \in S$
convient. Si nous avons $n > 2$ morphismes, alors nous utilisons
l'argument ci-dessus pour nous ramener au cas de $n - 1$ morphismes, et nous concluons
par induction.
\end{proof}

\noindent
Il existe une caractérisation facile de l'égalité des morphismes s'ils ont le
même dénominateur.

\begin{lemma}
\label{categories-lemma-equality-morphisms-left-localization}
Soit $\mathcal{C}$ une catégorie et soit $S$ un système multiplicatif à gauche
de morphismes de $\mathcal{C}$. Soient $A, B : X \to Y$ des morphismes de
$S^{-1}\mathcal{C}$ qui sont les classes d'équivalence de $(f : X \to Y', s : Y \to Y')$ et $(g : X \to Y', s : Y \to Y')$. Les assertions suivantes sont
équivalentes
\begin{enumerate}
\item $A = B$
\item il existe un morphisme $t : Y' \to Y''$ dans $S$ avec $t \circ f = t \circ g$, et
\item il existe un morphisme $a : Y' \to Y''$ tel que $a \circ f = a \circ g$
et $a \circ s \in S$.
\end{enumerate}
\end{lemma}

\begin{proof}
Nous utiliserons le fait que $S^{-1}\mathcal{C}$ est une catégorie (Lemme
\ref{categories-lemma-left-localization}) et nous utiliserons la notation de la définition
\ref{categories-definition-left-localization-as-fraction} ainsi que la discussion qui suit
cette définition pour identifier certains morphismes dans $S^{-1}\mathcal{C}$.
Ainsi nous écrivons $A = s^{-1}f$ et $B = s^{-1}g$.

\medskip\noindent
Si $A = B$ alors $(\text{id}_{Y'}^{-1}s) \circ A = (\text{id}_{Y'}^{-1}s) \circ B$. Nous avons $(\text{id}_{Y'}^{-1}s) \circ A = \text{id}_{Y'}^{-1}f$
et $(\text{id}_{Y'}^{-1}s) \circ B = \text{id}_{Y'}^{-1}g$. L'égalité de
$\text{id}_{Y'}^{-1}f$ et $\text{id}_{Y'}^{-1}g$ signifie par définition qu'il
existe un diagramme commutatif
$$
\xymatrix{
 & Y' \ar[d]^u & \\
X \ar[ru]^f \ar[r]^h \ar[rd]_g &
Z &
Y' \ar[lu]_{\text{id}_{Y'}} \ar[l]_t \ar[ld]^{\text{id}_{Y'}} \\
& Y' \ar[u]_v &
}
$$
avec $t \in S$. En particulier $u = v = t \in S$ et $t \circ f = t\circ g$.
Ainsi (1) implique (2).

\medskip\noindent
L'implication (2) $\Rightarrow$ (3) est immédiate. Supposons que $a$ soit
comme dans (3). Posons $s' = a \circ s \in S$. Alors $\text{id}_{Y''}^{-1}s'$
est un isomorphisme dans la catégorie $S^{-1}\mathcal{C}$ (avec l'inverse
$(s')^{-1}\text{id}_{Y''}$). Ainsi pour vérifier $A = B$ il suffit de vérifier
que $\text{id}_{Y''}^{-1}s' \circ A = \text{id}_{Y''}^{-1}s' \circ B$. Nous
calculons en utilisant les règles discutées dans le texte suivant la
définition \ref{categories-definition-left-localization-as-fraction} que
$\text{id}_{Y''}^{-1}s' \circ A =
\text{id}_{Y''}^{-1}(a \circ s) \circ s^{-1}f =
\text{id}_{Y''}^{-1}(a \circ f) =
\text{id}_{Y''}^{-1}(a \circ g) =
\text{id}_{Y''}^{-1}(a \circ s) \circ s^{-1}g =
\text{id}_{Y''}^{-1}s' \circ B$ et nous voyons que (1) est vrai.
\end{proof}

\begin{remark}
\label{categories-remark-left-localization-morphisms-colimit}
Soit $\mathcal{C}$ une catégorie. Soit $S$ un système multiplicatif à gauche.
Étant donné un objet $Y$ de $\mathcal{C}$ on note $Y/S$ la catégorie dont les
objets sont $s : Y \to Y'$ avec $s \in S$ et dont les morphismes sont des
diagrammes commutatifs
$$
\xymatrix{
& Y \ar[ld]_s \ar[rd]^t & \\
Y' \ar[rr]^a & & Y''
}
$$
où $a : Y' \to Y''$ est arbitraire. Nous affirmons que la catégorie $Y/S$ est
filtrante (voir définition \ref{categories-definition-directed}). En effet, LMS1 implique
que $\text{id}_Y : Y \to Y$ se trouve dans $Y/S$ ; donc $Y/S$ n'est pas vide.
LMS2 implique qu'étant donnés deux morphismes $s_1 : Y \to Y_1$ et $s_2 : Y \to Y_2$ nous
pouvons trouver un diagramme
$$
\xymatrix{
Y \ar[d]_{s_1} \ar[r]_{s_2} & Y_2 \ar[d]^t \\
Y_1 \ar[r]^a & Y_3
}
$$
avec $t \in S$. Par conséquent, $s_1 : Y \to Y_1$ et $s_2 : Y \to Y_2$ ont
tous deux des morphismes vers $t \circ s_2 : Y \to Y_3$ dans $Y/S$. Enfin, étant
donnés deux morphismes $a, b$ de $s_1 : Y \to Y_1$ à $s_2 : Y \to Y_2$ dans
$Y/S$, nous voyons que $a \circ s_1 = b \circ s_1$ ; donc par LMS3 il existe
un $t : Y_2 \to Y_3$ dans $S$ tel que $t \circ a = t \circ b$. Les
résultats conjoints des lemmes \ref{categories-lemma-morphisms-left-localization} et
\ref{categories-lemma-equality-morphisms-left-localization} nous disent que
\begin{equation}
\label{categories-equation-left-localization-morphisms-colimit}
\Mor_{S^{-1}\mathcal{C}}(X, Y) =
\colim_{(s : Y \to Y') \in Y/S} \Mor_\mathcal{C}(X, Y')
\end{equation}
Cette formule exprimant les ensembles de morphismes dans $S^{-1}\mathcal{C}$
comme limite inductive filtrante d'ensembles de morphismes dans $\mathcal{C}$
est parfois utile.
\end{remark}

\begin{lemma}
\label{categories-lemma-properties-left-localization}
Soit $\mathcal{C}$ une catégorie et soit $S$ un système multiplicatif à gauche
de morphismes de $\mathcal{C}$.
\begin{enumerate}
\item Les règles $X \mapsto X$ et $(f : X \to Y) \mapsto (f : X \to Y, \text{id}_Y : Y \to Y)$ définissent un foncteur $Q : \mathcal{C} \to S^{-1}\mathcal{C}$.
\item Pour tout $s \in S$ le morphisme $Q(s)$ est un isomorphisme dans
$S^{-1}\mathcal{C}$.
\item Si $G : \mathcal{C} \to \mathcal{D}$ est un foncteur tel que $G(s)$ est
inversible pour chaque $s \in S$, alors il existe un foncteur unique $H : S^{-1}\mathcal{C} \to \mathcal{D}$ tel que $H \circ Q = G$.
\end{enumerate}
\end{lemma}

\begin{proof}
Les assertions (1) et (2) sont claires. (Dans (2), l'inverse de $Q(s)$ est la
classe d'équivalence de la paire $(\text{id}_Y, s)$.) Pour voir (3),
définissons simplement $H(X) = G(X)$ et $H((f : X \to Y', s : Y \to Y')) = G(s)^{-1} \circ G(f)$. Détails omis.
\end{proof}

\begin{lemma}
\label{categories-lemma-left-localization-limits}
Soit $\mathcal{C}$ une catégorie et soit $S$ un système multiplicatif à gauche
de morphismes de $\mathcal{C}$. Le foncteur de localisation $Q : \mathcal{C} \to S^{-1}\mathcal{C}$ commute avec les colimites finies.
\end{lemma}

\begin{proof}
Soit $\mathcal{I}$ une catégorie finie et soit $\mathcal{I} \to \mathcal{C}$,
$i \mapsto X_i$ un foncteur dont la colimite existe. Alors, en utilisant
(\ref{categories-equation-left-localization-morphisms-colimit}), le fait que $Y/S$ soit
filtrante et le lemme \ref{categories-lemma-directed-commutes}, nous avons
\begin{align*}
\Mor_{S^{-1}\mathcal{C}}(Q(\colim X_i), Q(Y))
& =
\colim_{(s : Y \to Y') \in Y/S} \Mor_\mathcal{C}(\colim X_i, Y') \\
& =
\colim_{(s : Y \to Y') \in Y/S} \lim_i \Mor_\mathcal{C}(X_i, Y') \\
& =
\lim_i \colim_{(s : Y \to Y') \in Y/S} \Mor_\mathcal{C}(X_i, Y') \\
& =
\lim_i \Mor_{S^{-1}\mathcal{C}}(Q(X_i), Q(Y))
\end{align*}
et cet isomorphisme commute avec les projections des deux côtés vers
l'ensemble $\Mor_{S^{-1}\mathcal{C}}(Q(X_j), Q(Y))$ pour chaque $j \in \Ob(\mathcal{I})$. Ainsi, $Q(\colim X_i)$ satisfait la propriété universelle
de la colimite du foncteur $i \mapsto Q(X_i)$ ; c'est donc cette colimite,
comme souhaité.
\end{proof}

\begin{lemma}
\label{categories-lemma-left-localization-diagram}
Soit $\mathcal{C}$ une catégorie. Soit $S$ un système multiplicatif à gauche.
Si $f : X \to Y$, $f' : X' \to Y'$ sont deux morphismes de $\mathcal{C}$ et si
$$
\xymatrix{
Q(X) \ar[d]_{Q(f)} \ar[r]_a & Q(X') \ar[d]^{Q(f')} \\
Q(Y) \ar[r]^b & Q(Y')
}
$$
est un diagramme commutatif dans $S^{-1}\mathcal{C}$, alors il existe un
morphisme $f'' : X'' \to Y''$ dans $\mathcal{C}$ et un diagramme commutatif
$$
\xymatrix{
X \ar[d]_f \ar[r]_g & X'' \ar[d]^{f''} & X' \ar[d]^{f'} \ar[l]^s \\
Y \ar[r]^h & Y'' & Y' \ar[l]_t
}
$$
dans $\mathcal{C}$ avec $s, t \in S$ et $a = s^{-1}g$, $b = t^{-1}h$.
\end{lemma}

\begin{proof}
Nous choisissons les morphismes et les objets de la manière suivante :
Écrivons d'abord $a = s^{-1}g$ pour un certain $s : X' \to X''$ dans $S$ et un morphisme $g : X \to X''$. Par LMS2 on peut trouver $t : Y' \to Y''$ dans $S$ et $f'' : X'' \to Y''$ de sorte que
$$
\xymatrix{
X \ar[d]_f \ar[r]_g & X'' \ar[d]^{f''} & X' \ar[d]^{f'} \ar[l]^s \\
Y & Y'' & Y' \ar[l]_t
}
$$
est commutatif. À ce stade, dans ce diagramme, nous
allons modifier à plusieurs reprises notre choix de
$$
X'' \xrightarrow{f''} Y'' \xleftarrow{t} Y'
$$
en postcomposant à la fois $t$ et $f''$ par un morphisme $d : Y'' \to Y'''$
avec la propriété $d \circ t \in S$. D'après la remarque
\ref{categories-remark-left-localization-morphisms-colimit}, nous pouvons, après un tel
remplacement, supposer qu'il existe un morphisme $h : Y \to Y''$ tel que $b = t^{-1}h$\footnote{Voici une manière plus terre-à-terre de voir
cela : écrivons $b = q^{-1}i$ pour un certain $q : Y' \to Z$ dans $S$ et un certain
$i : Y \to Z$. Par LMS2, nous pouvons trouver $r : Y'' \to Y'''$ dans $S$ et
$j : Z \to Y'''$ tel que $j \circ q = r \circ t$. Posons alors $d = r$ et $h = j \circ i$.}. À ce stade, nous avons toutes les données du lemme, sauf
que nous ne savons pas que le carré de gauche du diagramme est commutatif.
Mais la définition de composition dans $S^{-1} \mathcal{C}$ montre que $b \circ Q\left(f\right)$ est la classe d'équivalence de la paire $(h \circ f : X \to Y'', t : Y' \to Y'')$ (puisque $b$ est la classe d'équivalence de la paire
$(h : Y \to Y'', t : Y' \to Y'')$, tandis que $Q\left(f\right)$ est la classe
d'équivalence de la paire $(f : X \to Y, \text{id} : Y \to Y)$), tandis que
$Q\left(f'\right) \circ a$ est la classe d'équivalence de la paire $(f'' \circ g : X \to Y'', t : Y' \to Y'')$ (puisque $a$ est la classe d'équivalence de la
paire $(g : X \to X'', s : X' \to X'')$, tandis que $Q\left(f'\right)$ est la
classe d'équivalence de la paire $(f' : X' \to Y', \text{id} : Y' \to Y')$). Puisque l'on sait que $b \circ Q\left(f\right) = Q\left(f'\right) \circ a$, on conclut que les classes d'équivalence des couples $(h \circ f : X \to Y'', t : Y' \to Y'')$ et $(f'' \circ g : X \to Y'', t : Y' \to Y'')$ sont
égales. Ainsi en utilisant le lemme
\ref{categories-lemma-equality-morphisms-left-localization} nous pouvons trouver un
morphisme $d : Y'' \to Y'''$ tel que $d \circ t \in S$ et $d \circ h \circ f = d \circ f'' \circ g$. Nous effectuons donc un remplacement supplémentaire du
type décrit ci-dessus et nous concluons.
\end{proof}

\noindent
{\bf Calcul des fractions à droite.} Soit $\mathcal{C}$ une catégorie et $S$ un
système multiplicatif à droite. Nous définissons une nouvelle catégorie
$S^{-1}\mathcal{C}$ comme suit (nous vérifions que cela fonctionne dans la
preuve du lemme \ref{categories-lemma-right-localization}) :
\begin{enumerate}
\item Nous définissons $\Ob(S^{-1}\mathcal{C}) = \Ob(\mathcal{C})$.
\item Les morphismes $X \to Y$ de $S^{-1}\mathcal{C}$ sont représentés par des couples
$(f : X' \to Y, s : X' \to X)$ avec $s \in S$ à équivalence près.
(L'équivalence est définie ci-dessous. Considérons la classe d'équivalence d'une
paire $(f, s)$ comme $fs^{-1} : X \to Y$.)
\item Deux paires $(f_1 : X_1 \to Y, s_1 : X_1 \to X)$ et $(f_2 : X_2 \to Y, s_2 : X_2 \to X)$ sont dites équivalentes s'il existe une troisième paire
$(f_3 : X_3 \to Y, s_3 : X_3 \to X)$ ainsi que des morphismes $u : X_3 \to X_1$ et $v : X_3 \to X_2$ de $\mathcal{C}$ qui s'insèrent dans le diagramme commutatif
$$
\xymatrix{
 & X_1 \ar[ld]_{s_1} \ar[rd]^{f_1} & \\
X &
X_3 \ar[l]_{s_3} \ar[u]_u \ar[d]^v \ar[r]^{f_3} &
Y \\
& X_2 \ar[lu]^{s_2} \ar[ru]_{f_2} &
}
$$
\item La composition des classes d'équivalence des couples $(f : X' \to Y, s : X' \to X)$ et $(g : Y' \to Z, t : Y' \to Y)$ est définie comme la classe
d'équivalence d'un couple $(g \circ h : X'' \to Z, s \circ u : X'' \to X)$ où
$h$ et $u \in S$ sont choisis pour s'insérer dans un diagramme commutatif
$$
\xymatrix{
X'' \ar[d]_u \ar[r]^h & Y' \ar[d]^t \\
X' \ar[r]^f & Y
}
$$
qui existe par hypothèse.
\item Le morphisme identité $X \to X$ dans $S^{-1} \mathcal{C}$ est la
classe d'équivalence de la paire $(\text{id} : X \to X,
\text{id} : X \to X)$.
\end{enumerate}

\begin{lemma}
\label{categories-lemma-right-localization}
Soit $\mathcal{C}$ une catégorie et $S$ un système multiplicatif à droite.
\begin{enumerate}
\item La relation sur les paires définie ci-dessus est une relation d'équivalence.
\item La règle de composition donnée ci-dessus est bien définie sur les
classes d'équivalence.
\item La composition est associative (et les morphismes identité satisfont
les axiomes d'identité), et donc $S^{-1}\mathcal{C}$ est une catégorie.
\end{enumerate}
\end{lemma}

\begin{proof}
Ce lemme est dual du lemme \ref{categories-lemma-left-localization}. Il découle
formellement de ce lemme en remplaçant $\mathcal{C}$ par sa catégorie opposée
dans laquelle $S$ est un système multiplicatif à gauche.
\end{proof}

\begin{definition}
\label{categories-definition-right-localization-as-fraction}
Soit $\mathcal{C}$ une catégorie et soit $S$ un système multiplicatif à droite
de morphismes de $\mathcal{C}$. Étant donné tout morphisme $f : X' \to Y$ dans
$\mathcal{C}$ et tout morphisme $s : X' \to X$ dans $S$, on note {\it $f s^{-1}$} la classe d'équivalence de la paire $(f : X' \to Y, s : X' \to X)$.
Il s'agit d'un morphisme de $X$ à $Y$ dans $S^{-1} \mathcal{C}$.
\end{definition}

\noindent
On a des identités analogues (en fait, duales) à celles de la définition
\ref{categories-definition-left-localization-as-fraction}. Nous pouvons « écrire
toute famille finie de morphismes de même source sous forme de
fractions à dénominateur commun ».

\begin{lemma}
\label{categories-lemma-morphisms-right-localization}
Soit $\mathcal{C}$ une catégorie et soit $S$ un système multiplicatif à droite
de morphismes de $\mathcal{C}$. Étant donnée toute famille finie $g_i : X \to Y_i$ de morphismes de $S^{-1}\mathcal{C}$ (indexée par $i$), on peut
trouver un élément $s : X' \to X$ de $S$ et une famille de morphismes $f_i : X' \to Y_i$ de $\mathcal{C}$ tels que chaque $g_i$ soit la classe d'équivalence du
couple $(f_i : X' \to Y_i, s : X' \to X)$.
\end{lemma}

\begin{proof}
Ce lemme est le dual du lemme \ref{categories-lemma-morphisms-left-localization} et
découle formellement de ce lemme en remplaçant toutes les catégories considérées
par leurs opposés.
\end{proof}

\noindent
Il existe une caractérisation facile de l'égalité des morphismes s'ils ont le
même dénominateur.

\begin{lemma}
\label{categories-lemma-equality-morphisms-right-localization}
Soit $\mathcal{C}$ une catégorie et soit $S$ un système multiplicatif à droite
de morphismes de $\mathcal{C}$. Soient $A, B : X \to Y$ des morphismes de
$S^{-1}\mathcal{C}$ qui sont les classes d'équivalence de $(f : X' \to Y, s : X' \to X)$ et $(g : X' \to Y, s : X' \to X)$. Les assertions suivantes sont
équivalentes
\begin{enumerate}
\item $A = B$,
\item il existe un morphisme $t : X'' \to X'$ dans $S$ avec $f \circ t = g \circ t$, et
\item il existe un morphisme $a : X'' \to X'$ avec $f \circ a = g \circ a$ et
$s \circ a \in S$.
\end{enumerate}
\end{lemma}

\begin{proof}
Ce lemme est dual du lemme \ref{categories-lemma-equality-morphisms-left-localization}.
\end{proof}

\begin{remark}
\label{categories-remark-right-localization-morphisms-colimit}
Soit $\mathcal{C}$ une catégorie. Soit $S$ un système multiplicatif à droite.
Étant donné un objet $X$ de $\mathcal{C}$ on note $S/X$ la catégorie dont les
objets sont $s : X' \to X$ avec $s \in S$ et dont les morphismes sont des
diagrammes commutatifs
$$
\xymatrix{
X' \ar[rd]_s \ar[rr]_a & & X'' \ar[ld]^t \\
& X
}
$$
où $a : X' \to X''$ est arbitraire. La catégorie $S/X$ est cofiltrante (voir
définition \ref{categories-definition-codirected}). (Ceci est le dual de l'énoncé
correspondant dans la remarque
\ref{categories-remark-left-localization-morphisms-colimit}.) Les résultats
conjoints des lemmes \ref{categories-lemma-morphisms-right-localization} et
\ref{categories-lemma-equality-morphisms-right-localization} nous disent que
\begin{equation}
\label{categories-equation-right-localization-morphisms-colimit}
\Mor_{S^{-1}\mathcal{C}}(X, Y) =
\colim_{(s : X' \to X) \in (S/X)^{opp}} \Mor_\mathcal{C}(X', Y)
\end{equation}
Cette formule exprimant les ensembles de morphismes dans $S^{-1}\mathcal{C}$
comme limite inductive filtrante d'ensembles de morphismes dans $\mathcal{C}$
est parfois utile.
\end{remark}

\begin{lemma}
\label{categories-lemma-properties-right-localization}
Soit $\mathcal{C}$ une catégorie et soit $S$ un système multiplicatif à droite
de morphismes de $\mathcal{C}$.
\begin{enumerate}
\item Les règles $X \mapsto X$ et $(f : X \to Y) \mapsto (f : X \to Y, \text{id}_X : X \to X)$ définissent un foncteur $Q : \mathcal{C} \to S^{-1}\mathcal{C}$.
\item Pour tout $s \in S$ le morphisme $Q(s)$ est un isomorphisme dans
$S^{-1}\mathcal{C}$.
\item Si $G : \mathcal{C} \to \mathcal{D}$ est un foncteur tel que $G(s)$ est
inversible pour chaque $s \in S$, alors il existe un foncteur unique $H : S^{-1}\mathcal{C} \to \mathcal{D}$ tel que $H \circ Q = G$.
\end{enumerate}
\end{lemma}

\begin{proof}
Ce lemme est le dual du lemme \ref{categories-lemma-properties-left-localization} et
découle formellement de ce lemme en remplaçant toutes les catégories considérées
par leurs opposés.
\end{proof}

\begin{lemma}
\label{categories-lemma-right-localization-limits}
Soit $\mathcal{C}$ une catégorie et soit $S$ un système multiplicatif à droite
de morphismes de $\mathcal{C}$. Le foncteur de localisation $Q : \mathcal{C} \to S^{-1}\mathcal{C}$ commute avec les limites finies.
\end{lemma}

\begin{proof}
Ce lemme est dual du lemme \ref{categories-lemma-left-localization-limits}.
\end{proof}

\begin{lemma}
\label{categories-lemma-right-localization-diagram}
Soit $\mathcal{C}$ une catégorie. Soit $S$ un système multiplicatif à droite.
Si $f : X \to Y$, $f' : X' \to Y'$ sont deux morphismes de $\mathcal{C}$ et si
$$
\xymatrix{
Q(X) \ar[d]_{Q(f)} \ar[r]_a & Q(X') \ar[d]^{Q(f')} \\
Q(Y) \ar[r]^b & Q(Y')
}
$$
est un diagramme commutatif dans $S^{-1}\mathcal{C}$, alors il existe un
morphisme $f'' : X'' \to Y''$ dans $\mathcal{C}$ et un diagramme commutatif
$$
\xymatrix{
X \ar[d]_f & X'' \ar[l]^s \ar[d]^{f''} \ar[r]_g & X' \ar[d]^{f'} \\
Y & Y'' \ar[l]_t \ar[r]^h & Y'
}
$$
dans $\mathcal{C}$ avec $s, t \in S$ et $a = gs^{-1}$, $b = ht^{-1}$.
\end{lemma}

\begin{proof}
Ce lemme est dual du lemme \ref{categories-lemma-left-localization-diagram}.
\end{proof}

\noindent
{\bf Systèmes multiplicatifs et calcul bilatéral des fractions.} Si $S$ est un
système multiplicatif, alors les calculs de fractions à gauche et à droite donnent des
catégories canoniquement isomorphes.

\begin{lemma}
\label{categories-lemma-multiplicative-system}
Soit $\mathcal{C}$ une catégorie et $S$ un système multiplicatif. Les catégories
de fractions à gauche et de fractions à droite
$S^{-1}\mathcal{C}$ sont canoniquement isomorphes.
\end{lemma}

\begin{proof}
Notons $\mathcal{C}_{left}$, $\mathcal{C}_{right}$ les deux catégories de
fractions. Par les propriétés universelles des lemmes
\ref{categories-lemma-properties-left-localization} et
\ref{categories-lemma-properties-right-localization} nous obtenons les foncteurs
$\mathcal{C}_{left} \to \mathcal{C}_{right}$ et $\mathcal{C}_{right} \to \mathcal{C}_{left}$. D'après l'énoncé d'unicité dans les propriétés
universelles, ces foncteurs sont inverses les uns des autres.
\end{proof}

\begin{definition}
\label{categories-definition-saturated-multiplicative-system}
Soit $\mathcal{C}$ une catégorie et $S$ un système multiplicatif. On dit que
$S$ est {\it saturé} si, en plus de MS1, MS2, MS3, on a aussi
\begin{enumerate}
\item[MS4] Étant donnés trois morphismes composables $f, g, h$, si $fg, gh \in S$, alors $g \in S$.
\end{enumerate}
\end{definition}

\noindent
Notez qu'un système multiplicatif saturé contient tous les isomorphismes. De
plus, si $f, g, h$ sont des morphismes composables dans une catégorie et que
$fg, gh$ sont des isomorphismes, alors $g$ est un isomorphisme (car alors $g$
a à la fois un inverse gauche et un inverse droit, donc est inversible).

\begin{lemma}
\label{categories-lemma-what-gets-inverted}
Soit $\mathcal{C}$ une catégorie et $S$ un système multiplicatif. Notons $Q : \mathcal{C} \to S^{-1}\mathcal{C}$ le foncteur de localisation. L'ensemble
$$
\hat S = \{f \in \text{Flèches}(\mathcal{C}) \mid
Q(f) \text{ est un isomorphisme}\}
$$
est égal à
$$
S' = \{f \in \text{Flèches}(\mathcal{C}) \mid
\text{il existe }g, h\text{ tels que }gf, fh \in S\}
$$
et est le plus petit système multiplicatif saturé contenant $S$. En
particulier, si $S$ est saturé, alors $\hat S = S$.
\end{lemma}

\begin{proof}
Il est clair que $S \subset S' \subset \hat S$ car les éléments de $S'$
correspondent aux morphismes de $S^{-1}\mathcal{C}$ qui ont à la fois des
inverses à gauche et à droite. Notez que $S'$ satisfait à MS4 et que $\hat S$
satisfait à MS1. Montrons ensuite que $S' = \hat S$.

\medskip\noindent
Soit $f \in \hat S$. Soit $s^{-1}g = ht^{-1}$ le morphisme inverse dans
$S^{-1}\mathcal{C}$. (Nous pouvons utiliser à la fois des fractions gauches et
des fractions droites pour décrire les morphismes dans $S^{-1}\mathcal{C}$,
voir le lemme \ref{categories-lemma-multiplicative-system}.) La relation $\text{id}_X = s^{-1}gf$ dans $S^{-1}\mathcal{C}$ signifie qu'il existe un diagramme
commutatif.
$$
\xymatrix{
 & X' \ar[d]^u & \\
X \ar[ru]^{gf} \ar[r]^{f'} \ar[rd]_{\text{id}_X} &
X'' &
X \ar[lu]_s \ar[l]_{s'} \ar[ld]^{\text{id}_X} \\
& X \ar[u]_v &
}
$$
pour certains morphismes $f', u, v$ et $s' \in S$. D'où $ugf = s' \in S$. De
même, en utilisant l'égalité $\text{id}_Y = fht^{-1}$ on prouve que $fhw \in S$ pour
un certain $w$. Nous concluons que $f \in S'$. Ainsi $S' = \hat S$. Une fois
établi que $S' = \hat S$ est un système multiplicatif, il est
clair que $S' = \hat S$ est le plus petit système
saturé contenant $S$.

\medskip\noindent
Nos remarques ci-dessus établissent MS1 et MS4, donc pour terminer la
preuve du lemme nous devons montrer que LMS2, RMS2, LMS3, RMS3 sont valables
pour $\hat S$. Vérifions que LMS2 est satisfaite pour $\hat S$. Supposons que
nous ayons un diagramme en traits pleins
$$
\xymatrix{
X \ar[d]_t \ar[r]_g & Y \ar@{..>}[d]^s \\
Z \ar@{..>}[r]^f & W
}
$$
avec $t \in \hat S$. Choisissons un morphisme $a : Z \to Z'$ tel que $at \in S$. Nous pouvons alors appliquer LMS2 à $S$ pour trouver un diagramme
commutatif
$$
\xymatrix{
X \ar[d]_t \ar[r]_g & Y \ar[dd]^s \\
Z \ar[d]_a \\
Z' \ar[r]^{f'} & W
}
$$
et, en posant $f = f' \circ a$, nous concluons. La preuve de RMS2 est duale. Enfin,
supposons qu'une paire de morphismes $f, g : X \to Y$ et $t \in \hat S$ soient
donnés, où ce dernier est de but $X$ et vérifie $ft = gt$. Choisissons alors un
morphisme $b$ tel que $tb \in S$. Alors $ftb = gtb$, ce qui implique par LMS3 pour
$S$ qu'il existe un $s \in S$ de source $Y$ tel que $sf = sg$, comme
souhaité. La preuve de RMS3 est duale.
\end{proof}








\section{Propriétés formelles}
\label{categories-section-formal-cat-cat}

\noindent
Dans cette section, nous discutons de certaines propriétés formelles de la
$2$-catégorie des catégories. Cela nous amènera plus tard à la définition d'une
$2$-catégorie (stricte).

\medskip\noindent
Notons $\Ob(\textit{Cat})$ la classe de toutes les catégories. Pour chaque
paire de catégories $\mathcal{A}, \mathcal{B} \in \Ob(\textit{Cat})$, nous
avons la « petite » catégorie des foncteurs $\text{Fun}(\mathcal{A}, \mathcal{B})$. La composition de transformations de foncteurs comme
$$
\xymatrix{
\mathcal{A}
\rruppertwocell^{F''}{t'}
\ar[rr]_(.3){F'}
\rrlowertwocell_F{t}
& &
\mathcal{B}
}
\text{ a pour composé }
\xymatrix{
\mathcal{A}
\rrtwocell^{F''}_F{\ \ t \circ t'}
& &
\mathcal{B}
}
$$
est appelée composition {\it verticale}. Nous utiliserons pour cela le symbole
habituel $\circ$. Nous définirons ensuite la composition {\it horizontale}.
Pour ce faire, nous expliquons un peu plus la structure en question.

\medskip\noindent
En effet, pour chaque triplet de catégories $\mathcal{A}$, $\mathcal{B}$ et
$\mathcal{C}$, il existe une loi de composition
$$
\circ : \Ob(\text{Fun}(\mathcal{B}, \mathcal{C}))
\times
\Ob(\text{Fun}(\mathcal{A}, \mathcal{B}))
\longrightarrow
\Ob(\text{Fun}(\mathcal{A}, \mathcal{C}))
$$
provenant de la composition des foncteurs. Cette loi de composition est
associative et les foncteurs identité agissent comme des unités. Autrement
dit -- en oubliant les transformations des foncteurs -- on voit que
$\textit{Cat}$ forme une catégorie. Comment cette structure interagit-elle
avec les morphismes entre foncteurs ?

\medskip\noindent
Soit $t : F \to F'$ une transformation de foncteurs, où $F, F' : \mathcal{A} \to \mathcal{B}$, et soit $G : \mathcal{B} \to \mathcal{C}$ un foncteur.
Nous pouvons définir une transformation de foncteurs $G\circ F \to G \circ F'$, que nous noterons ${}_Gt$. Elle est donnée par la formule
$({}_Gt)_x = G(t_x) : G(F(x)) \to G(F'(x))$ pour tous les $x \in \mathcal{A}$.
De cette façon, la composition avec $G$ devient un foncteur
$$
\text{Fun}(\mathcal{A}, \mathcal{B})
\longrightarrow
\text{Fun}(\mathcal{A}, \mathcal{C}).
$$
Pour le voir, il suffit de vérifier que ${}_G(\text{id}_F) = \text{id}_{G \circ F}$ et que ${}_G(t_1 \circ t_2) = {}_Gt_1 \circ {}_Gt_2$.
Bien sûr, nous avons aussi ${}_{\text{id}_\mathcal{B}}t = t$.

\medskip\noindent
De même, soit $s : G \to G'$ une transformation de foncteurs, où $G, G' : \mathcal{B} \to \mathcal{C}$, et soit $F : \mathcal{A} \to \mathcal{B}$ un foncteur.
Nous pouvons définir $s_F$ comme étant la transformation de foncteurs $G\circ F \to G' \circ F$ donnée par $(s_F)_x = s_{F(x)} : G(F(x)) \to G'(F(x))$ pour
tous les $x \in \mathcal{A}$. De cette façon, la composition avec $F$ devient
un foncteur
$$
\text{Fun}(\mathcal{B}, \mathcal{C})
\longrightarrow
\text{Fun}(\mathcal{A}, \mathcal{C}).
$$
Pour le voir, il suffit de vérifier que $(\text{id}_G)_F = \text{id}_{G\circ F}$ et que $(s_1 \circ s_2)_F = s_{1, F} \circ s_{2, F}$.
Bien sûr, nous avons aussi $s_{\text{id}_\mathcal{B}} = s$.

\medskip\noindent
Ces constructions satisfont en outre aux propriétés
$$
{}_{G_1}({}_{G_2}t) = {}_{G_1\circ G_2}t,
\ (s_{F_1})_{F_2} = s_{F_1 \circ F_2},
\text{ et }{}_H(s_F) = ({}_Hs)_F
$$
chaque fois que cela a du sens. Enfin, étant donnés les foncteurs $F, F' : \mathcal{A} \to \mathcal{B}$ et $G, G' : \mathcal{B} \to \mathcal{C}$ ainsi que les
transformations $t : F \to F'$ et $s : G \to G'$, le diagramme suivant est
commutatif
$$
\xymatrix{
G \circ F \ar[r]^{{}_Gt} \ar[d]_{s_F}
&
G \circ F' \ar[d]^{s_{F'}} \\
G' \circ F \ar[r]_{{}_{G'}t}
&
G' \circ F'
}
$$
en d'autres termes ${}_{G'}t \circ s_F = s_{F'}\circ {}_Gt$. Pour le prouver,
considérons simplement ce qui se passe sur tout objet $x \in \Ob(\mathcal{A})$ :
$$
\xymatrix{
G(F(x)) \ar[r]^{G(t_x)} \ar[d]_{s_{F(x)}}
&
G(F'(x)) \ar[d]^{s_{F'(x)}} \\
G'(F(x)) \ar[r]_{G'(t_x)}
&
G'(F'(x))
}
$$
qui est commutatif car $s$ est une transformation de foncteurs. Cette relation
de compatibilité permet de définir une composition horizontale.

\begin{definition}
\label{categories-definition-horizontal-composition}
Étant donné un diagramme comme dans le premier diagramme ci-dessous :
$$
\xymatrix{
\mathcal{A}
\rtwocell^F_{F'}{t}
&
\mathcal{B}
\rtwocell^G_{G'}{s}
&
\mathcal{C}
}
\text{ donne }
\xymatrix{
\mathcal{A}
\rrtwocell^{G \circ F} _{G' \circ F'}{\ \ s \star t}
& &
\mathcal{C}
}
$$
nous définissons la composition {\it horizontale} $s \star t$ comme étant la
transformation de foncteurs ${}_{G'}t \circ s_F = s_{F'}\circ {}_Gt$.
\end{definition}

\noindent
Nous voyons maintenant que l'on peut retrouver nos transformations
précédemment construites ${}_Gt$ et $s_F$ sous les noms $ {}_Gt = \text{id}_G \star t $ et $ s_F = s \star \text{id}_F $. De plus, toutes les règles que
nous avons trouvées ci-dessus sont des conséquences des propriétés énoncées
dans le lemme qui suit.

\begin{lemma}
\label{categories-lemma-properties-2-cat-cats}
Les compositions horizontales et verticales ont les propriétés suivantes
\begin{enumerate}
\item $\circ$ et $\star$ sont associatifs,
\item les transformations identité $\text{id}_F$ sont des unités pour
$\circ$,
\item les transformations identité des foncteurs identité
$\text{id}_{\text{id}_\mathcal{A}}$ sont des unités pour $\star$ et $\circ$,
et
\item étant donné un diagramme
$$
\xymatrix{
\mathcal{A}
\rruppertwocell^F{t}
\ar[rr]_(.3){F'}
\rrlowertwocell_{F''}{t'}
& &
\mathcal{B}
\rruppertwocell^G{s}
\ar[rr]_(.3){G'}
\rrlowertwocell_{G''}{s'}
& &
\mathcal{C}
}
$$
nous avons $ (s' \circ s) \star (t' \circ t) = (s' \star t') \circ (s \star t)$.
\end{enumerate}
\end{lemma}

\begin{proof}
La dernière assertion devient, dans la notation précédente, l'équation
suivante
$$
s'_{F''}
\circ
{}_{G'}t'
\circ
s_{F'}
\circ
{}_Gt
=
(s' \circ s)_{F''}
\circ
{}_G(t' \circ t).
$$
D'après notre résultat ci-dessus appliqué à la composition du milieu, nous
pouvons réécrire le côté gauche comme $
s'_{F''}
\circ
s_{F''}
\circ
{}_Gt'
\circ
{}_Gt
$ qui est facilement reconnu égal au côté droit.
\end{proof}

\noindent
Une autre façon de formuler la condition (4) du lemme est que la composition
des foncteurs et la composition horizontale des transformations de foncteurs
définissent un foncteur
$$
(\circ, \star) :
\text{Fun}(\mathcal{B}, \mathcal{C})
\times
\text{Fun}(\mathcal{A}, \mathcal{B})
\longrightarrow
\text{Fun}(\mathcal{A}, \mathcal{C})
$$
dont la source est la catégorie produit, voir Définition
\ref{categories-definition-product-category}.

\section{2-catégories}
\label{categories-section-2-categories}

\noindent
Nous donnerons une définition des $2$-catégories (strictes) telles qu'elles
apparaissent dans le contexte des champs. Avant de poursuivre, consulter
la section \ref{categories-section-formal-cat-cat} et l'exemple
\ref{categories-example-2-1-category-of-categories}. Il s'agit essentiellement de partir de cet
exemple et d'expliciter toutes les règles satisfaites par les objets, les
$1$-morphismes et les $2$-morphismes dans cet exemple.

\begin{definition}
\label{categories-definition-2-category}
Une (stricte) {\it $2$-catégorie} $\mathcal{C}$ se compose des données suivantes
\begin{enumerate}
\item Un ensemble d'objets $\Ob(\mathcal{C})$.
\item Pour chaque paire $x, y \in \Ob(\mathcal{C})$ une catégorie
$\Mor_\mathcal{C}(x, y)$. Les objets de $\Mor_\mathcal{C}(x, y)$ seront
appelés {\it $1$-morphismes} et notés $F : x \to y$. Les morphismes entre ces
$1$-morphismes seront appelés {\it $2$-morphismes} et notés $t : F' \to F$. La
composition des $2$-morphismes dans $\Mor_\mathcal{C}(x, y)$ sera appelée
composition {\it verticale} et sera notée $t \circ t'$ pour $t : F' \to F$ et
$t' : F'' \to F'$.
\item Pour chaque triple $x, y, z\in \Ob(\mathcal{C})$ un foncteur
$$
(\circ, \star) :
\Mor_\mathcal{C}(y, z) \times \Mor_\mathcal{C}(x, y)
\longrightarrow
\Mor_\mathcal{C}(x, z).
$$
L'image du couple de $1$-morphismes $(F, G)$ dans le membre de gauche sera appelée
la {\it composition} de $F$ et $G$, et notée $F\circ G$. L'image de la paire
de $2$-morphismes $(t, s)$ sera appelée composition {\it horizontale} et notée
$t \star s$.
\end{enumerate}
Ces données doivent satisfaire aux règles suivantes :
\begin{enumerate}
\item L'ensemble des objets ainsi que l'ensemble des $1$-morphismes dotés
d'une composition de $1$-morphismes forment une catégorie.
\item La composition horizontale des $2$-morphismes est associative.
\item Le $2$-morphisme identité $\text{id}_{\text{id}_x}$ de l'identité
$1$-morphisme $\text{id}_x$ est une unité pour la composition horizontale.
\end{enumerate}
\end{definition}

\noindent
Ce type de structure n'est évidemment pas très commode à manier. D'un autre
côté, il existe de nombreux exemples où la manière de procéder est
assez claire. Le seul exemple que nous ayons jusqu'à présent est
celui de la $2$-catégorie dont les objets forment une famille donnée de
catégories, les $1$-morphismes sont des foncteurs entre ces catégories, et les
$2$-morphismes sont des transformations naturelles de foncteurs, voir la
section \ref{categories-section-formal-cat-cat}. En ce qui concerne ce texte, toutes les
$2$-catégories seront des sous-$2$-catégories de cet exemple. Voici ce que
signifie être une sous-$2$-catégorie.

\begin{definition}
\label{categories-definition-sub-2-category}
Soit $\mathcal{C}$ une $2$-catégorie. Une {\it sous-$2$-catégorie}
$\mathcal{C}'$ de $\mathcal{C}$ est donnée par un sous-ensemble
$\Ob(\mathcal{C}')$ de $\Ob(\mathcal{C})$ et des sous-catégories
$\Mor_{\mathcal{C}'}(x, y)$ des catégories $\Mor_\mathcal{C}(x, y)$ pour tous
les $x, y \in \Ob(\mathcal{C}')$ de sorte que ces données, avec les opérations
$\circ$ (composition des $1$-morphismes), $\circ$ (composition verticale des
$2$-morphismes) et $\star$ (composition horizontale) forment une
$2$-catégorie.
\end{definition}

\begin{remark}
\label{categories-remark-big-2-categories}
Grandes $2$-catégories. Dans de nombreux textes, une $2$-catégorie est
autorisée à avoir une classe d'objets (mais, espérons-le, une « classe de
classes » n'est pas autorisée). Nous autoriserons également ces « grandes »
$2$-catégories, mais uniquement dans la liste de cas suivante (à mettre à jour
au fur et à mesure) :
\begin{enumerate}
\item La $2$-catégorie des catégories $\textit{Cat}$.
\item La $(2, 1)$-catégorie des catégories $\textit{Cat}$.
\item La $2$-catégorie des groupoïdes $\textit{Groupoïdes}$ ; il s'agit d'une
$(2, 1)$-catégorie.
\item La $2$-catégorie des catégories fibrées au-dessus d'une catégorie fixe.
\item La $(2, 1)$-catégorie des catégories fibrées au-dessus d'une catégorie fixe.
\item La $2$-catégorie des catégories fibrées en groupoïdes au-dessus d'une catégorie
fixe ; il s'agit d'une $(2, 1)$-catégorie.
\item La $2$-catégorie des champs sur un site fixe.
\item La $(2, 1)$-catégorie des champs sur un site fixe.
\item La $2$-catégorie des champs en groupoïdes sur un site fixe ; il s'agit
d'une $(2, 1)$-catégorie.
\item La $2$-catégorie des champs en setoïdes sur un site fixe ; il s'agit d'une
$(2, 1)$-catégorie.
\item La $2$-catégorie des champs algébriques sur un schéma fixe ; il s'agit
d'une $(2, 1)$-catégorie.
\end{enumerate}
Voir la définition \ref{categories-definition-2-1-category}. Notez que dans chaque cas,
la classe d'objets de la $2$-catégorie $\mathcal{C}$ est une classe propre,
mais pour tous les objets $x, y \in \Ob(\mathcal{C})$, la catégorie $\Mor_\mathcal{C}(x, y)$ est « petite » (selon nos conventions).
\end{remark}

\noindent
La notion d'équivalence des catégories que nous avons définie dans la section
\ref{categories-section-definition-categories} s'étend au cadre plus général des
$2$-catégories comme suit.

\begin{definition}
\label{categories-definition-equivalence}
Deux objets $x, y$ d'une $2$-catégorie sont {\it équivalents} s'il existe des
$1$-morphismes $F : x \to y$ et $G : y \to x$ tels que $F \circ G$ est
$2$-isomorphe à $\text{id}_y$ et $G \circ F$ est $2$-isomorphe à
$\text{id}_x$.
\end{definition}

\noindent
Parfois, nous devons dire ce que signifie avoir un foncteur d'une catégorie
dans une $2$-catégorie.

\begin{definition}
\label{categories-definition-functor-into-2-category}
Soit $\mathcal{A}$ une catégorie et $\mathcal{C}$ une $2$-catégorie.
\begin{enumerate}
\item Un {\it foncteur} d'une catégorie ordinaire dans une $2$-catégorie
ignorera les $2$-morphismes, sauf indication contraire. Autrement dit, ce sera
un foncteur « habituel » dans la catégorie obtenue à partir de cette 2-catégorie en oubliant
tous les 2-morphismes.
\item Un {\it foncteur faible}, ou {\it pseudo-foncteur} $\varphi$ de
$\mathcal{A}$ dans la 2-catégorie $\mathcal{C}$ est défini par les données
suivantes :
\begin{enumerate}
\item une application $\varphi : \Ob(\mathcal{A}) \to \Ob(\mathcal{C})$,
\item pour chaque paire $x, y\in \Ob(\mathcal{A})$ et chaque morphisme $f : x \to y$, un $1$-morphisme $\varphi(f) : \varphi(x) \to \varphi(y)$,
\item pour chaque $x\in \Ob(\mathcal{A})$ un $2$-morphisme $\alpha_x : \text{id}_{\varphi(x)} \to \varphi(\text{id}_x)$, et
\item pour chaque paire de morphismes composables $f : x \to y$, $g : y \to z$
de $\mathcal{A}$, un $2$-morphisme $\alpha_{g, f} : \varphi(g \circ f) \to \varphi(g) \circ \varphi(f)$.
\end{enumerate}
Ces données sont soumises aux conditions suivantes :
\begin{enumerate}
\item les $2$-morphismes $\alpha_x$ et $\alpha_{g, f}$ sont tous des
isomorphismes,
\item pour tout morphisme $f : x \to y$ dans $\mathcal{A}$ nous avons
$\alpha_{\text{id}_y, f} = \alpha_y \star \text{id}_{\varphi(f)}$ :
$$
\xymatrix{
\varphi(x)
\rrtwocell^{\varphi(f)}_{\varphi(f)}{\ \ \ \ \text{id}_{\varphi(f)}}
& &
\varphi(y)
\rrtwocell^{\text{id}_{\varphi(y)}}_{\varphi(\text{id}_y)}{\alpha_y}
& &
\varphi(y)
}
=
\xymatrix{
\varphi(x)
\rrtwocell^{\varphi(f)}_{\varphi(\text{id}_y) \circ \varphi(f)}{\ \ \ \ \alpha_{\text{id}_y, f}}
& &
\varphi(y)
}
$$
\item pour tout morphisme $f : x \to y$ dans $\mathcal{A}$ nous avons
$\alpha_{f, \text{id}_x} = \text{id}_{\varphi(f)} \star \alpha_x$,
\item pour tout triple de morphismes composables $f : w \to x$, $g : x \to y$
et $h : y \to z$ de $\mathcal{A}$ nous avons
$$
(\text{id}_{\varphi(h)} \star \alpha_{g, f})
\circ
\alpha_{h, g \circ f}
=
(\alpha_{h, g} \star \text{id}_{\varphi(f)})
\circ
\alpha_{h \circ g, f}
$$
en d'autres termes, le diagramme suivant, dont les objets sont les $1$-morphismes
et les flèches les $2$-morphismes, est commutatif
$$
\xymatrix{
\varphi(h \circ g \circ f)
\ar[d]_{\alpha_{h, g \circ f}}
\ar[rr]_{\alpha_{h \circ g, f}}
& &
\varphi(h \circ g) \circ \varphi(f)
\ar[d]^{\alpha_{h, g} \star \text{id}_{\varphi(f)}} \\
\varphi(h) \circ \varphi(g \circ f)
\ar[rr]^{\text{id}_{\varphi(h)} \star \alpha_{g, f}}
& &
\varphi(h) \circ \varphi(g) \circ \varphi(f)
}
$$
\end{enumerate}
\end{enumerate}
\end{definition}

\noindent
Encore une fois, ce n'est pas une notion très pratique, mais elle revient
parfois. Il existe un théorème qui dit que tout pseudo-foncteur est isomorphe
à un foncteur. Enfin, il y a les notions de {\it foncteur entre
$2$-catégories} et de {\it pseudo-foncteur entre $2$-catégories}. Cette dernière
notion nous amène dans le domaine des $3$-catégories. Nous aimerions
éviter de devoir définir cela à tout prix ou presque !

\section{(2, 1)-catégories}
\label{categories-section-2-1-categories}

\noindent
Certaines $2$-catégories ont la propriété que tous les $2$-morphismes sont des
isomorphismes. Celles-ci joueront un rôle important dans la suite et seront plus
faciles à utiliser.

\begin{definition}
\label{categories-definition-2-1-category}
Une (stricte) {\it $(2, 1)$-catégorie} est une $2$-catégorie dans laquelle
tous les $2$-morphismes sont des isomorphismes.
\end{definition}

\begin{example}
\label{categories-example-2-1-category-of-categories}
La $2$-catégorie $\textit{Cat}$, voir la remarque
\ref{categories-remark-big-2-categories}, peut être transformée en $(2, 1)$-catégorie en
ne conservant comme $2$-morphismes que les isomorphismes de foncteurs.
\end{example}

\noindent
En fait, plus généralement toute $2$-catégorie $\mathcal{C}$ donne une
$(2, 1)$-catégorie en considérant la sous-$2$-catégorie $\mathcal{C}'$ avec
les mêmes objets et $1$-morphismes mais dont les $2$-morphismes sont les
$2$-morphismes inversibles de $\mathcal{C}$. Dans cette situation, nous dirons
« {\it soit $\mathcal{C}'$ la $(2, 1)$-catégorie associée à
$\mathcal{C}$} » ou une formulation analogue. Par exemple, la $(2, 1)$-catégorie de
groupoïdes désigne la $2$-catégorie dont les objets sont des groupoïdes, dont
les $1$-morphismes sont des foncteurs et dont les $2$-morphismes sont des
isomorphismes de foncteurs. Cet exemple est toutefois mal choisi, car une
transformation entre foncteurs entre groupoïdes est automatiquement un
isomorphisme !

\begin{remark}
\label{categories-remark-other-2-categories}
Il existe donc des variantes de la construction de l'exemple
\ref{categories-example-2-1-category-of-categories} ci-dessus dans lesquelles on considère la
$2$-catégorie des groupoïdes, ou des catégories fibrées en groupoïdes au-dessus d'une
catégorie fixe, ou des champs. Et ainsi de suite.
\end{remark}






\section{2-produits fibrés}
\label{categories-section-2-fibre-products}

\noindent
Dans cette section, nous présentons les $2$-produits fibrés. Supposons que
$\mathcal{C}$ soit une 2-catégorie. On dit qu'un diagramme
$$
\xymatrix{
w \ar[r] \ar[d] & y \ar[d] \\
x \ar[r] & z }
$$
est 2-commutatif si les deux 1-morphismes $w \to y \to z$ et $w \to x \to z$ sont
2-isomorphes. Dans une 2-catégorie, il est plus naturel de demander la
2-commutativité des diagrammes que leur commutativité stricte. (En
effet, certains diront peut-être que nous ne devrions pas du tout travailler
avec des 2-catégories strictes, et dans une 2-catégorie « faible » la notion
de diagramme commutatif de 1-morphismes n'a même pas de sens.) En conséquence,
la notion de produit fibré doit être ajustée.

\medskip\noindent
Soit $\mathcal{C}$ une $2$-catégorie. Soient $x, y, z\in \Ob(\mathcal{C})$ et
$f\in \Mor_\mathcal{C}(x, z)$ et $g\in \Mor_{\mathcal C}(y, z)$. Afin de
définir le 2-produit fibré de $f$ et $g$, nous allons examiner les
diagrammes 2-commutatifs
$$
\xymatrix{
& w \ar[r]_a \ar[d]_b & x \ar[d]^{f} \\
& y \ar[r]^{g} & z. }
$$
Dans le cas des catégories ordinaires, le produit fibré est un objet final
dans la catégorie de tels diagrammes. En conséquence, un 2-produit fibré
est un objet final dans une 2-catégorie (voir la définition ci-dessous). La {\it
$2$-catégorie des diagrammes $2$-commutatifs sur $f$ et $g$} est la
$2$-catégorie définie comme suit :
\begin{enumerate}
\item Les objets sont des quadruples $(w, a, b, \phi)$ comme ci-dessus où
$\phi$ est un 2-morphisme inversible $\phi : f \circ a \to g \circ b$,
\item Les 1-morphismes de $(w', a', b', \phi')$ à $(w, a, b, \phi)$ sont
les triplets $(k : w' \to w, \alpha : a' \to a \circ k, \beta : b' \to b \circ k)$ tels que
$$
\xymatrix{
f \circ a'
\ar[rr]_{\text{id}_f \star \alpha}
\ar[d]_{\phi'}
& &
f \circ a \circ k
\ar[d]^{\phi \star \text{id}_k}
\\
g \circ b'
\ar[rr]^{\text{id}_g \star \beta}
& &
g \circ b \circ k
}
$$
est commutatif,
\item étant donné un deuxième $1$-morphisme $(k', \alpha', \beta') : (w'', a'', b'', \phi'') \to
(w', a', b', \phi')$, la composition des
$1$-morphismes est donnée par la règle
$$
(k, \alpha, \beta) \circ (k', \alpha', \beta') =
(k \circ k',
(\alpha \star \text{id}_{k'}) \circ \alpha',
(\beta \star \text{id}_{k'}) \circ \beta'),
$$
\item un 2-morphisme entre les $1$-morphismes $(k_i, \alpha_i, \beta_i)$, $i = 1, 2$ avec la même source et le même but est donné par un 2-morphisme
$\delta : k_1 \to k_2$ tel que
$$
\xymatrix{
a'
\ar[rd]_{\alpha_2}
\ar[r]_{\alpha_1} &
a \circ k_1
\ar[d]^{\text{id}_a \star \delta} &
&
b \circ k_1
\ar[d]_{\text{id}_b \star \delta} &
b'
\ar[l]^{\beta_1}
\ar[ld]^{\beta_2}
\\
&
a \circ k_2
&
&
b \circ k_2
&
}
$$
est commutatif,
\item la composition verticale des $2$-morphismes est donnée par la
composition verticale des morphismes $\delta$ dans $\mathcal{C}$, et
\item la composition horizontale du diagramme
$$
\xymatrix{
(w'', a'', b'', \phi'')
\rrtwocell^{(k'_1, \alpha'_1, \beta'_1)}_{(k'_2, \alpha'_2, \beta'_2)}{\delta'}
& &
(w', a', b', \phi')
\rrtwocell^{(k_1, \alpha_1, \beta_1)}_{(k_2, \alpha_2, \beta_2)}{\delta}
& &
(w, a, b, \phi)
}
$$
est donnée par le diagramme
$$
\xymatrix@C=12pc{
(w'', a'', b'', \phi'')
\rtwocell^{(k_1 \circ k'_1, (\alpha_1 \star \text{id}_{k'_1}) \circ \alpha'_1, (\beta_1 \star \text{id}_{k'_1}) \circ \beta'_1)}_{(k_2 \circ k'_2, (\alpha_2 \star \text{id}_{k'_2}) \circ \alpha'_2, (\beta_2 \star \text{id}_{k'_2}) \circ \beta'_2)}{\ \ \ \delta \star \delta'}
&
(w, a, b, \phi)
}
$$
\end{enumerate}
Notez que si $\mathcal{C}$ est en fait une $(2, 1)$-catégorie, les morphismes
$\alpha$ et $\beta$ dans (2) ci-dessus sont automatiquement également des
isomorphismes\footnote{En fait, il semble que dans le cas d'une $2$-catégorie,
l'on pourrait définir une autre 2-catégorie de diagrammes 2-commutatifs où le
sens des flèches $\alpha$, $\beta$ est inversé, ou encore où le sens d'une
seule d'entre elles est inversé. C'est pourquoi nous nous limiterons plus tard
aux $(2, 1)$-catégories.}. De plus, la $2$-catégorie des diagrammes
$2$-commutatifs est également une $(2, 1)$-catégorie si $\mathcal{C}$ est une
$(2, 1)$-catégorie.

\begin{definition}
\label{categories-definition-final-object-2-category}
Un {\it objet final} d'une $(2, 1)$-catégorie $\mathcal{C}$ est un objet $x$
tel que
\begin{enumerate}
\item pour chaque $y \in \Ob(\mathcal{C})$ il existe un morphisme $y \to x$,
et
\item deux morphismes quelconques $y \to x$ sont isomorphes par un 2-morphisme
unique.
\end{enumerate}
\end{definition}

\noindent
Il est probable que dans le cas plus général des $2$-catégories, il existe
différentes versions d'objets finaux. Nous ne voulons pas entrer dans ce sujet
et c'est pourquoi nous définissons uniquement les $2$-produits fibrés dans
le cas des $(2, 1)$-catégories.

\begin{definition}
\label{categories-definition-2-fibre-products}
Soit $\mathcal{C}$ une $(2, 1)$-catégorie. Soient $x, y, z\in \Ob(\mathcal{C})$ et $f\in \Mor_\mathcal{C}(x, z)$ et $g\in \Mor_{\mathcal C}(y, z)$. Un {\it 2-produit fibré de $f$ et $g$} est un objet final dans
la $2$-catégorie des diagrammes $2$-commutatifs décrite ci-dessus. Si un 2-produit
fibré existe, nous le noterons $x \times_z y\in \Ob(\mathcal{C})$, et
noterons les morphismes correspondants $p\in \Mor_{\mathcal C}(x \times_z y, x)$ et
$q\in \Mor_{\mathcal C}(x \times_z y, y)$ rendant le diagramme
$$
\xymatrix{
& x \times_z y \ar[r]^{p} \ar[d]_q & x \ar[d]^{f} \\
& y \ar[r]^{g} & z }
$$
2-commutatif ; nous noterons le 2-morphisme inversible qui en témoigne
$\psi : f \circ p \to g \circ q$.
\end{definition}

\noindent
On a donc la propriété universelle suivante : pour tout $w\in \Ob(\mathcal{C})$, tous les morphismes $a \in \Mor_{\mathcal C}(w, x)$ et $b \in \Mor_\mathcal{C}(w, y)$ munis d'un 2-isomorphisme $\phi : f \circ a \to g\circ b$
donné, il existe un $\gamma \in \Mor_{\mathcal C}(w, x \times_z y)$ tel que le
diagramme
$$
\xymatrix{
w\ar[rrrd]^a \ar@{-->}[rrd]_\gamma \ar[rrdd]_b & & \\
& & x \times_z y \ar[r]_p \ar[d]_q & x \ar[d]^{f} \\
& & y \ar[r]^{g} & z }
$$
soit 2-commutatif, de telle sorte que, pour des choix appropriés de $a \to p \circ \gamma$ et $b \to q \circ \gamma$, le diagramme
$$
\xymatrix{
f \circ a \ar[r] \ar[d]_\phi &
f \circ p \circ \gamma
\ar[d]^{\psi \star \text{id}_\gamma}
\\
g\circ b
\ar[r]
&
g \circ q \circ \gamma
}
$$
est commutatif. De plus, $\gamma$ est unique à
l'isomorphisme près. Bien entendu, les propriétés précises sont plus fines que
cela. Tous les cas de 2-produits fibrés dont nous aurons besoin plus tard
proviennent de l'exemple suivant de 2-produits fibrés dans la 2-catégorie de
catégories.

\begin{example}
\label{categories-example-2-fibre-product-categories}
Soient $\mathcal{A}$, $\mathcal{B}$ et $\mathcal{C}$ des catégories. Soient $F : \mathcal{A} \to \mathcal{C}$ et $G : \mathcal{B} \to \mathcal{C}$ des
foncteurs. Nous définissons une catégorie $\mathcal{A} \times_\mathcal{C} \mathcal{B}$ comme suit :
\begin{enumerate}
\item un objet de $\mathcal{A} \times_\mathcal{C} \mathcal{B}$ est un triple
$(A, B, f)$, où $A\in \Ob(\mathcal{A})$, $B\in \Ob(\mathcal{B})$ et $f : F(A) \to G(B)$ est un isomorphisme dans $\mathcal{C}$,
\item un morphisme $(A, B, f) \to (A', B', f')$ est donné par une paire $(a, b)$, où $a : A \to A'$ est un morphisme dans $\mathcal{A}$, et $b : B \to B'$
est un morphisme dans $\mathcal{B}$ tel que le diagramme
$$
\xymatrix{
F(A) \ar[r]^f \ar[d]^{F(a)} & G(B) \ar[d]^{G(b)} \\
F(A') \ar[r]^{f'} & G(B')
}
$$
est commutatif.
\end{enumerate}
De plus, nous définissons les foncteurs $p : \mathcal{A} \times_\mathcal{C}\mathcal{B} \to \mathcal{A}$ et $q : \mathcal{A} \times_\mathcal{C}\mathcal{B} \to \mathcal{B}$ en posant
$$
p(A, B, f) = A, \quad q(A, B, f) = B,
$$
en d'autres termes, ce sont les foncteurs d'oubli. Nous définissons une
transformation de foncteurs $\psi : F \circ p \to G \circ q$. Sur l'objet $\xi = (A, B, f)$ elle est donnée par $\psi_\xi = f : F(p(\xi)) = F(A) \to G(B) = G(q(\xi))$.
\end{example}

\begin{lemma}
\label{categories-lemma-2-fibre-product-categories}
Dans la $(2, 1)$-catégorie des catégories, des $2$-produits fibrés existent
et sont donnés par la construction de l'exemple
\ref{categories-example-2-fibre-product-categories}.
\end{lemma}

\begin{proof}
Vérifions la propriété universelle : soit $\mathcal{W}$ une catégorie, soit $a : \mathcal{W} \to \mathcal{A}$ et $b : \mathcal{W} \to \mathcal{B}$ des
foncteurs, et soit $t : F \circ a \to G \circ b$ un isomorphisme de foncteurs.

\medskip\noindent
Considérons le foncteur $\gamma : \mathcal{W} \to \mathcal{A} \times_\mathcal{C}\mathcal{B}$ donné par $W \mapsto (a(W), b(W), t_W)$.
(La vérification qu'il s'agit d'un foncteur est omise.) De plus, considérons $\alpha : a \to p \circ \gamma$ et $\beta : b \to q \circ \gamma$ obtenus à partir des
identités $p \circ \gamma = a$ et $q \circ \gamma = b$. Il est alors clair que
$(\gamma, \alpha, \beta)$ est un morphisme de $(\mathcal{W}, a, b, t)$ vers $(\mathcal{A} \times_\mathcal{C} \mathcal{B}, p, q, \psi)$.

\medskip\noindent
Soit $(k, \alpha', \beta') :
(\mathcal{W}, a, b, t) \to (\mathcal{A} \times_\mathcal{C} \mathcal{B}, p, q, \psi)$ un deuxième morphisme de ce type. Pour un objet $W$
de $\mathcal{W}$ écrivons $k(W) = (a_k(W), b_k(W), t_{k, W})$. D'où $p(k(W)) = a_k(W)$, et ainsi de suite. La transformation $\alpha'$ correspond à une
famille fonctorielle de morphismes $\alpha' : a(W) \to a_k(W)$. Puisque nous
travaillons dans la $(2, 1)$-catégorie des catégories, en fait chacun de ces
morphismes $a(W) \to a_k(W)$ est un isomorphisme. Nous pouvons les utiliser
(ainsi que leurs analogues $b(W) \to b_k(W)$) pour obtenir des isomorphismes
$$
\delta_W :
\gamma(W) = (a(W), b(W), t_W)
\longrightarrow
(a_k(W), b_k(W), t_{k, W}) = k(W).
$$
Il est immédiat de montrer que $\delta$ définit un $2$-isomorphisme entre
$\gamma$ et $k$ dans la $2$-catégorie des diagrammes $2$-commutatifs, comme
souhaité.
\end{proof}

\begin{remark}
\label{categories-remark-other-description-2-fibre-product}
Soient $\mathcal{A}$, $\mathcal{B}$ et $\mathcal{C}$ des catégories. Soient $F : \mathcal{A} \to \mathcal{C}$ et $G : \mathcal{B} \to \mathcal{C}$ des
foncteurs. Une autre construction, légèrement plus symétrique, d'un
$2$-produit fibré $\mathcal{A} \times_\mathcal{C} \mathcal{B}$ est la suivante. Un
objet est un quintuple $(A, B, C, a, b)$ où $A, B, C$ sont des objets de
$\mathcal{A}, \mathcal{B}, \mathcal{C}$ et où $a : F(A) \to C$ et $b : G(B) \to C$ sont des isomorphismes. Un morphisme $(A, B, C, a, b) \to (A', B', C', a', b')$ est donné par un triplet de morphismes $A \to A', B \to B', C \to C'$
compatible avec les morphismes $a, b, a', b'$. Nous pouvons prouver
directement que cette construction donne un $2$-produit fibré. Cependant, il est plus
facile d'observer que le foncteur $(A, B, C, a, b) \mapsto (A, B, b^{-1} \circ a)$ donne une équivalence entre la catégorie des quintuples et la catégorie
construite dans l'exemple \ref{categories-example-2-fibre-product-categories}.
\end{remark}

\begin{lemma}
\label{categories-lemma-functoriality-2-fibre-product}
Soit le diagramme
$$
\xymatrix{
& \mathcal{Y} \ar[d]_I \ar[rd]^K & \\
\mathcal{X} \ar[r]^H \ar[rd]^L &
\mathcal{Z} \ar[rd]^M & \mathcal{B} \ar[d]^G \\
& \mathcal{A} \ar[r]^F & \mathcal{C}
}
$$
$2$-commutatif de catégories. Un choix d'isomorphismes
$\alpha : G \circ K \to M \circ I$ et $\beta : M \circ H \to F \circ L$
détermine un morphisme
$$
\mathcal{X} \times_\mathcal{Z} \mathcal{Y}
\longrightarrow
\mathcal{A} \times_\mathcal{C} \mathcal{B}
$$
entre les $2$-produits fibrés associés à cette situation.
\end{lemma}

\begin{proof}
On utilise simplement le foncteur
$$
(X, Y, \phi) \longmapsto (L(X), K(Y),
\alpha^{-1}_Y \circ M(\phi) \circ \beta^{-1}_X)
$$
sur les objets et
$$
(a, b) \longmapsto (L(a), K(b))
$$
sur les morphismes.
\end{proof}

\begin{lemma}
\label{categories-lemma-equivalence-2-fibre-product}
Plaçons-nous sous les hypothèses du lemme \ref{categories-lemma-functoriality-2-fibre-product}.
\begin{enumerate}
\item Si $K$ et $L$ sont fidèles alors le morphisme $\mathcal{X} \times_\mathcal{Z} \mathcal{Y} \to
\mathcal{A} \times_\mathcal{C} \mathcal{B}$
est fidèle.
\item Si $K$ et $L$ sont pleinement fidèles et que $M$ est fidèle alors le
morphisme $\mathcal{X} \times_\mathcal{Z} \mathcal{Y} \to
\mathcal{A} \times_\mathcal{C} \mathcal{B}$ est pleinement fidèle.
\item Si $K$ et $L$ sont des équivalences et que $M$ est pleinement fidèle
alors le morphisme $\mathcal{X} \times_\mathcal{Z} \mathcal{Y} \to
\mathcal{A} \times_\mathcal{C} \mathcal{B}$ est une équivalence.
\end{enumerate}
\end{lemma}

\begin{proof}
Soient $(X, Y, \phi)$ et $(X', Y', \phi')$ des objets de $\mathcal{X} \times_\mathcal{Z} \mathcal{Y}$. Posons $Z = H(X)$ et identifions-le à
$I(Y)$ au moyen de $\phi$. Identifions également $M(Z)$ à $F(L(X))$ au moyen de $\beta_X$
et $M(Z)$ à $G(K(Y))$ au moyen de $\alpha_Y$. De même pour $Z' = H(X')$
et $M(Z')$. L'application sur les morphismes est l'application
$$
\xymatrix{
\Mor_\mathcal{X}(X, X')
\times_{\Mor_\mathcal{Z}(Z, Z')}
\Mor_\mathcal{Y}(Y, Y')
\ar[d] \\
\Mor_\mathcal{A}(L(X), L(X'))
\times_{\Mor_\mathcal{C}(M(Z), M(Z'))}
\Mor_\mathcal{B}(K(Y), K(Y'))
}
$$
Les parties (1) et (2) suivent donc. De plus, si $K$ et $L$ sont des
équivalences et que $M$ est pleinement fidèle, alors tout objet $(A, B, \phi)$ appartient à l'image essentielle pour les raisons suivantes : choisissons
$X$, $Y$ tels que $L(X) \cong A$ et $K(Y) \cong B$. Alors la pleine fidélité de
$M$ garantit que l'on peut trouver un isomorphisme $H(X) \cong I(Y)$. Certains
détails sont omis.
\end{proof}

\begin{lemma}
\label{categories-lemma-associativity-2-fibre-product}
Soit le diagramme
$$
\xymatrix{
\mathcal{A} \ar[rd] & & \mathcal{C} \ar[ld] \ar[rd] & & \mathcal{E} \ar[ld] \\
& \mathcal{B} & & \mathcal{D}
}
$$
de catégories et de foncteurs. Alors il existe un
isomorphisme canonique
$$
(\mathcal{A} \times_\mathcal{B} \mathcal{C}) \times_\mathcal{D} \mathcal{E}
\cong
\mathcal{A} \times_\mathcal{B} (\mathcal{C} \times_\mathcal{D} \mathcal{E})
$$
de catégories.
\end{lemma}

\begin{proof}
On utilise simplement le foncteur
$$
((A, C, \phi), E, \psi)
\longmapsto
(A, (C, E, \psi), \phi)
$$
ce qui définit le foncteur voulu, si l'on voit ce que cela signifie.
\end{proof}

\noindent
Désormais, nous omettrons les parenthèses lorsqu'il s'agit de produits
fibrés de plus de 2 catégories.

\begin{lemma}
\label{categories-lemma-triple-2-fibre-product-pr02}
Soit le diagramme
$$
\xymatrix{
\mathcal{A} \ar[rd] & & \mathcal{C} \ar[ld] \ar[rd] & & \mathcal{E} \ar[ld] \\
& \mathcal{B} \ar[rd]_F & & \mathcal{D} \ar[ld]^G \\
& & \mathcal{F} &
}
$$
commutatif de catégories et de foncteurs. Alors il existe un
foncteur canonique
$$
\text{pr}_{02} :
\mathcal{A} \times_\mathcal{B} \mathcal{C} \times_\mathcal{D} \mathcal{E}
\longrightarrow
\mathcal{A} \times_\mathcal{F} \mathcal{E}
$$
de catégories.
\end{lemma}

\begin{proof}
Si nous écrivons $\mathcal{A} \times_\mathcal{B} \mathcal{C}
\times_\mathcal{D} \mathcal{E}$ comme $(\mathcal{A} \times_\mathcal{B} \mathcal{C})
\times_\mathcal{D} \mathcal{E}$ alors nous pouvons simplement
utiliser le foncteur
$$
((A, C, \phi), E, \psi)
\longmapsto
(A, E, G(\psi) \circ F(\phi))
$$
ce qui définit le foncteur voulu, si l'on voit ce que cela signifie.
\end{proof}

\begin{lemma}
\label{categories-lemma-2-fibre-product-erase-factor}
Soit le diagramme
$$
\mathcal{A} \to
\mathcal{B} \leftarrow \mathcal{C} \leftarrow \mathcal{D}
$$
de catégories et de foncteurs. Alors il existe un
isomorphisme canonique
$$
\mathcal{A} \times_\mathcal{B} \mathcal{C} \times_\mathcal{C} \mathcal{D}
\cong
\mathcal{A} \times_\mathcal{B} \mathcal{D}
$$
de catégories.
\end{lemma}

\begin{proof}
Omis.
\end{proof}

\noindent
Nous affirmons que cela signifie que l'on peut travailler avec ces
$2$-produits fibrés comme avec des produits fibrés ordinaires. Voici quelques
autres lemmes qui reviendront plus tard.

\begin{lemma}
\label{categories-lemma-diagonal-1}
Soit le diagramme
$$
\xymatrix{
\mathcal{C}_3 \ar[r] \ar[d] & \mathcal{S} \ar[d]^\Delta \\
\mathcal{C}_1 \times \mathcal{C}_2 \ar[r]^{G_1 \times G_2} &
\mathcal{S} \times \mathcal{S}
}
$$
un $2$-produit fibré de catégories. Il existe alors un isomorphisme
canonique $\mathcal{C}_3 \cong
\mathcal{C}_1 \times_{G_1, \mathcal{S}, G_2} \mathcal{C}_2$.
\end{lemma}

\begin{proof}
Nous pouvons supposer que $\mathcal{C}_3$ est la catégorie $(\mathcal{C}_1 \times \mathcal{C}_2)\times_{\mathcal{S} \times \mathcal{S}}
\mathcal{S}$
construite dans l'exemple \ref{categories-example-2-fibre-product-categories}. Un objet
est donc un triple $((X_1, X_2), S, \phi)$ où $\phi = (\phi_1, \phi_2) : (G_1(X_1), G_2(X_2)) \to (S, S)$ est un isomorphisme. On peut ainsi y associer
le triple $(X_1, X_2, \phi_2^{-1} \circ \phi_1)$. À l'inverse, si $(X_1, X_2, \psi)$ est un objet de $\mathcal{C}_1 \times_{G_1, \mathcal{S}, G_2} \mathcal{C}_2$, alors on peut lui associer le triple $((X_1, X_2), G_2(X_2), (\psi, \text{id}_{G_2(X_2)}))$. Nous affirmons que ces constructions
définissent des foncteurs mutuellement inverses. Nous omettons de décrire comment
traiter les morphismes et de montrer qu'ils sont mutuellement inverses.
\end{proof}

\begin{lemma}
\label{categories-lemma-diagonal-2}
Soit le diagramme
$$
\xymatrix{
\mathcal{C}' \ar[r] \ar[d] & \mathcal{S} \ar[d]^\Delta \\
\mathcal{C} \ar[r]^{(G_1, G_2)} &
\mathcal{S} \times \mathcal{S}
}
$$
un $2$-produit fibré de catégories. Alors il existe un isomorphisme
canonique
$$
\mathcal{C}' \cong
(\mathcal{C} \times_{G_1, \mathcal{S}, G_2} \mathcal{C})
\times_{(p, q), \mathcal{C} \times \mathcal{C}, \Delta}
\mathcal{C}.
$$
\end{lemma}

\begin{proof}
Un objet du côté droit est donné par $((C_1, C_2, \phi), C_3, \psi)$ où $\phi : G_1(C_1) \to G_2(C_2)$ est un isomorphisme et $\psi = (\psi_1, \psi_2) : (C_1, C_2) \to (C_3, C_3)$ est un isomorphisme. On peut donc y associer le
triple $(C_3, G_1(C_1), (G_1(\psi_1^{-1}), \phi^{-1} \circ G_2(\psi_2^{-1})))$
qui est un objet de $\mathcal{C}'$. Détails omis.
\end{proof}

\begin{lemma}
\label{categories-lemma-fibre-product-after-map}
Soient $\mathcal{A} \to \mathcal{C}$, $\mathcal{B} \to \mathcal{C}$ et
$\mathcal{C} \to \mathcal{D}$ des foncteurs entre catégories. Alors le
diagramme
$$
\xymatrix{
\mathcal{A} \times_\mathcal{C} \mathcal{B} \ar[d] \ar[r] &
\mathcal{A} \times_\mathcal{D} \mathcal{B} \ar[d] \\
\mathcal{C} \ar[r]^-{\Delta_{\mathcal{C}/\mathcal{D}}} &
\mathcal{C} \times_\mathcal{D} \mathcal{C}
}
$$
est $2$-cartésien.
\end{lemma}

\begin{proof}
Omis.
\end{proof}

\begin{lemma}
\label{categories-lemma-base-change-diagonal}
Soit le diagramme
$$
\xymatrix{
\mathcal{U} \ar[d] \ar[r] & \mathcal{V} \ar[d] \\
\mathcal{X} \ar[r] & \mathcal{Y}
}
$$
un $2$-produit fibré de catégories. Alors le diagramme
$$
\xymatrix{
\mathcal{U} \ar[d] \ar[r] &
\mathcal{U} \times_\mathcal{V} \mathcal{U} \ar[d] \\
\mathcal{X} \ar[r] &
\mathcal{X} \times_\mathcal{Y} \mathcal{X}
}
$$
est $2$-cartésien.
\end{lemma}

\begin{proof}
Il s'agit d'un énoncé purement $2$-catégorique, valable dans
toute $(2, 1)$-catégorie avec des $2$-produits fibrés. Explicitement, cela
découle de la chaîne d'équivalences suivante :
\begin{align*}
\mathcal{X} \times_{(\mathcal{X} \times_\mathcal{Y} \mathcal{X})}
(\mathcal{U} \times_\mathcal{V} \mathcal{U})
& =
\mathcal{X} \times_{(\mathcal{X} \times_\mathcal{Y} \mathcal{X})}
((\mathcal{X} \times_\mathcal{Y} \mathcal{V})
\times_\mathcal{V} (\mathcal{X} \times_\mathcal{Y} \mathcal{V})) \\
& =
\mathcal{X} \times_{(\mathcal{X} \times_\mathcal{Y} \mathcal{X})}
(\mathcal{X} \times_\mathcal{Y} \mathcal{X}
\times_\mathcal{Y} \mathcal{V}) \\
& =
\mathcal{X} \times_\mathcal{Y} \mathcal{V} = \mathcal{U}
\end{align*}
voir les lemmes \ref{categories-lemma-associativity-2-fibre-product} et
\ref{categories-lemma-2-fibre-product-erase-factor}.
\end{proof}








\section{Catégories au-dessus de catégories}
\label{categories-section-categories-over-categories}

\noindent
Dans cette section, nous avons un foncteur $p : \mathcal{S} \to \mathcal{C}$.
Nous considérons $\mathcal{S}$ comme placée au-dessus et $\mathcal{C}$ comme
placée au-dessous. Pour nous
assurer que tout le monde sait de quoi nous parlons, nous définissons la
$2$-catégorie des catégories au-dessus de $\mathcal{C}$.

\begin{definition}
\label{categories-definition-categories-over-C}
Soit $\mathcal{C}$ une catégorie. La {\it $2$-catégorie des catégories au-dessus de
$\mathcal{C}$} est la $2$-catégorie définie comme suit :
\begin{enumerate}
\item Ses objets seront des foncteurs $p : \mathcal{S} \to \mathcal{C}$.
\item Ses $1$-morphismes $(\mathcal{S}, p) \to (\mathcal{S}', p')$ seront des
foncteurs $G : \mathcal{S} \to \mathcal{S}'$ tels que $p' \circ G = p$.
\item Ses $2$-morphismes $t : G \to H$ pour $G, H : (\mathcal{S}, p) \to (\mathcal{S}', p')$ seront des morphismes de foncteurs tels que $p'(t_x) = \text{id}_{p(x)}$ pour tous $x \in \Ob(\mathcal{S})$.
\end{enumerate}
Dans cette situation nous désignerons par
$$
\Mor_{\textit{Cat}/\mathcal{C}}(\mathcal{S}, \mathcal{S}')
$$
la catégorie des $1$-morphismes entre $(\mathcal{S}, p)$ et $(\mathcal{S}', p')$.
\end{definition}

\noindent
Dans cette $2$-catégorie, nous définissons la composition horizontale et
verticale exactement comme cela est fait pour $\textit{Cat}$ dans la section
\ref{categories-section-formal-cat-cat}. Les axiomes d'une $2$-catégorie sont satisfaits
pour la même raison que ceux de $\textit{Cat}$. Pour voir cela, on peut
également utiliser le fait que les axiomes sont valables dans $\textit{Cat}$
et vérifier des assertions telles que « la composition verticale des
$2$-morphismes au-dessus de $\mathcal{C}$ donne un autre $2$-morphisme
au-dessus de $\mathcal{C}$ ». C'est
clair.

\medskip\noindent
Par analogie avec la fibre d'une application d'espaces, nous avons la
notion de catégorie fibre, ainsi que quelques notions de relèvement associées à
cette situation.

\begin{definition}
\label{categories-definition-fibre-category}
Soit $\mathcal{C}$ une catégorie. Soit $p : \mathcal{S} \to \mathcal{C}$ une
catégorie au-dessus de $\mathcal{C}$.
\begin{enumerate}
\item La {\it catégorie fibre} au-dessus d'un objet $U\in \Ob(\mathcal{C})$ est la
catégorie $\mathcal{S}_U$ avec des objets
$$
\Ob(\mathcal{S}_U) = \{x\in \Ob(\mathcal{S}) :
p(x) = U\}
$$
et morphismes
$$
\Mor_{\mathcal{S}_U}(x, y) = \{ \phi \in \Mor_\mathcal{S}(x, y) :
p(\phi) = \text{id}_U\}.
$$
\item Un {\it relèvement} d'un objet $U \in \Ob(\mathcal{C})$ est un objet
$x\in \Ob(\mathcal{S})$ tel que $p(x) = U$, c'est-à-dire que $x\in \Ob(\mathcal{S}_U)$. Nous dirons aussi parfois que {\it $x$ est au-dessus de
$U$}.
\item De même, un {\it relèvement} d'un morphisme $f : V \to U$ dans $\mathcal{C}$
est un morphisme $\phi : y \to x$ dans $\mathcal{S}$ tel que $p(\phi) = f$. On
dit parfois que {\it $\phi$ est au-dessus de $f$}.
\end{enumerate}
\end{definition}

\noindent
Nous pourrions faire ici quelques observations. Par exemple, si $F : (\mathcal{S}, p) \to (\mathcal{S}', p')$ est un $1$-morphisme de catégories
au-dessus de $\mathcal{C}$, alors $F$ induit des foncteurs entre catégories
fibres $F : \mathcal{S}_U \to \mathcal{S}'_U$. Il en va de même pour les
$2$-morphismes.

\medskip\noindent
Voici le lemme incontournable décrivant le $2$-produit fibré dans la
$(2, 1)$-catégorie des catégories au-dessus de $\mathcal{C}$.

\begin{lemma}
\label{categories-lemma-2-product-categories-over-C}
Soit $\mathcal{C}$ une catégorie. La $(2, 1)$-catégorie des catégories
au-dessus de $\mathcal{C}$ admet des 2-produits fibrés. Supposons que $F : \mathcal{X} \to \mathcal{S}$ et $G : \mathcal{Y} \to \mathcal{S}$ soient des
morphismes de catégories au-dessus de $\mathcal{C}$. Un 2-produit fibré explicite
$\mathcal{X} \times_\mathcal{S}\mathcal{Y}$ est donné par la description
suivante :
\begin{enumerate}
\item un objet de $\mathcal{X} \times_\mathcal{S} \mathcal{Y}$ est un
quadruple $(U, x, y, f)$, où $U \in \Ob(\mathcal{C})$, $x\in \Ob(\mathcal{X}_U)$, $y\in \Ob(\mathcal{Y}_U)$ et $f : F(x) \to G(y)$ est un
isomorphisme dans $\mathcal{S}_U$,
\item un morphisme $(U, x, y, f) \to (U', x', y', f')$ est donné par une paire
$(a, b)$, où $a : x \to x'$ est un morphisme dans $\mathcal{X}$, et $b : y \to y'$ est un morphisme dans $\mathcal{Y}$ tel que
\begin{enumerate}
\item $a$ et $b$ induisent le même morphisme $U \to U'$, et
\item le diagramme
$$
\xymatrix{
F(x) \ar[r]^f \ar[d]^{F(a)} & G(y) \ar[d]^{G(b)} \\
F(x') \ar[r]^{f'} & G(y')
}
$$
est commutatif.
\end{enumerate}
\end{enumerate}
Les foncteurs $p : \mathcal{X} \times_\mathcal{S}\mathcal{Y} \to \mathcal{X}$
et $q : \mathcal{X} \times_\mathcal{S}\mathcal{Y} \to \mathcal{Y}$ sont les
foncteurs d'oubli dans ce cas. La transformation $\psi : F \circ p \to
G \circ q$ est donnée sur l'objet $\xi = (U, x, y, f)$ par $\psi_\xi = f : F(p(\xi)) = F(x) \to G(y) = G(q(\xi))$.
\end{lemma}

\begin{proof}
Vérifions la propriété universelle : soit $p_\mathcal{W} : \mathcal{W}\to \mathcal{C}$ une catégorie au-dessus de $\mathcal{C}$, soient $X : \mathcal{W} \to \mathcal{X}$ et $Y : \mathcal{W} \to \mathcal{Y}$ des foncteurs au-dessus de
$\mathcal{C}$, et soit $t : F \circ X \to G \circ Y$ un isomorphisme de
foncteurs au-dessus de $\mathcal{C}$. Le foncteur recherché $\gamma : \mathcal{W} \to \mathcal{X} \times_\mathcal{S} \mathcal{Y}$ est donné par $W \mapsto (p_\mathcal{W}(W), X(W), Y(W), t_W)$. Détails omis ; comparer avec le lemme
\ref{categories-lemma-2-fibre-product-categories}.
\end{proof}

\begin{example}
\label{categories-example-different-2-fibre-products}
Les constructions de $2$-produits fibrés dans la 2-catégorie des catégories
au-dessus d'une catégorie, données dans le lemme \ref{categories-lemma-2-product-categories-over-C},
et dans la 2-catégorie des catégories, données dans le lemme
\ref{categories-lemma-2-fibre-product-categories} (comme dans l'exemple
\ref{categories-example-2-fibre-product-categories}) produisent des résultats non
équivalents en général. Plus précisément, soit $\mathcal{S}$ le groupoïde
avec un objet et deux flèches, et soit $\mathcal{X}$ la catégorie discrète
avec un objet. En prenant le $2$-produit fibré $\mathcal{X} \times_\mathcal{S} \mathcal{X}$ en tant que catégories, on obtient la catégorie
discrète avec deux objets. Cependant, si nous considérons toutes ces catégories
comme des catégories au-dessus de $\mathcal{S}$, le $2$-produit fibré $\mathcal{X} \times_\mathcal{S} \mathcal{X}$ en tant que catégories au-dessus de $\mathcal{S}$ est
la catégorie discrète avec un objet. La différence est que (dans la notation
du lemme \ref{categories-lemma-2-product-categories-over-C}), nous étions autorisés à
choisir n'importe quel isomorphisme de comparaison $f$ dans la première
situation, mais nous ne pouvions choisir la flèche d'identité que dans la
seconde situation.
\end{example}

\begin{lemma}
\label{categories-lemma-fibre-2-fibre-product-categories-over-C}
Soit $\mathcal{C}$ une catégorie. Soient $f : \mathcal{X} \to \mathcal{S}$ et $g : \mathcal{Y} \to \mathcal{S}$ des morphismes de catégories au-dessus de $\mathcal{C}$.
Pour tout objet $U$ de $\mathcal{C}$ nous avons l'identité suivante entre les
catégories fibres
$$
\left(\mathcal{X} \times_\mathcal{S}\mathcal{Y}\right)_U
=
\mathcal{X}_U \times_{\mathcal{S}_U} \mathcal{Y}_U
$$
\end{lemma}

\begin{proof}
Omis.
\end{proof}

\section{Catégories fibrées}
\label{categories-section-fibred-categories}

\noindent
Il convient de donner une très brève présentation des catégories fibrées.

\medskip\noindent
Soit $p : \mathcal{S} \to \mathcal{C}$ une catégorie au-dessus de
$\mathcal{C}$. Étant donné un objet $x \in \mathcal{S}$ tel que $p(x) = U$,
et un morphisme $f : V \to U$, nous pouvons essayer de former une sorte de
« produit fibré $V \times_U x$ » (ou un {\it changement de base} de $x$
par $V \to U$). En effet, un morphisme d'un objet $z \in \mathcal{S}$
vers « $V \times_U x$ » devrait être donné par une paire
$(\varphi, g)$, où
$\varphi : z \to x$, $g : p(z) \to V$ et
$p(\varphi) = f \circ g$. Voici la situation :
$$
\xymatrix{
z \ar@{~>}[d]^p \ar@{-}[r] &
? \ar[r] \ar@{~>}[d]^p &
x \ar@{~>}[d]^p \\
p(z) \ar[r] & V \ar[r]^f & U
}
$$
Si un tel morphisme $V \times_U x \to x$ existe, on l'appelle un morphisme
fortement cartésien.

\begin{definition}
\label{categories-definition-cartesian-over-C}
Soit $\mathcal{C}$ une catégorie.
Soit $p : \mathcal{S} \to \mathcal{C}$ une catégorie au-dessus de
$\mathcal{C}$. Un {\it morphisme fortement cartésien}, ou plus précisément
un {\it morphisme fortement $\mathcal{C}$-cartésien}, est un
morphisme $\varphi : y \to x$ de $\mathcal{S}$ tel que, pour tout
$z \in \Ob(\mathcal{S})$, l'application
$$
\Mor_\mathcal{S}(z, y)
\longrightarrow
\Mor_\mathcal{S}(z, x)
\times_{\Mor_\mathcal{C}(p(z), p(x))}
\Mor_\mathcal{C}(p(z), p(y)),
$$
donnée par $\psi \longmapsto (\varphi \circ \psi, p(\psi))$,
soit bijective.
\end{definition}

\noindent
Notons que, par le lemme de Yoneda \ref{categories-lemma-yoneda}, étant donnés
$x \in \Ob(\mathcal{S})$ au-dessus de $U \in \Ob(\mathcal{C})$
et un morphisme $f : V \to U$ de $\mathcal{C}$, s'il existe un
morphisme fortement cartésien $\varphi : y \to x$ tel que
$p(\varphi) = f$, alors $(y, \varphi)$ est unique à isomorphisme unique près.
Cela résulte clairement de la définition ci-dessus, puisque le foncteur
$$
z
\longmapsto
\Mor_\mathcal{S}(z, x)
\times_{\Mor_\mathcal{C}(p(z), U)}
\Mor_\mathcal{C}(p(z), V)
$$
ne dépend que des données $(x, U, f : V \to U)$. Ainsi, nous emploierons
parfois $V \times_U x \to x$ ou $f^*x \to x$ pour désigner un morphisme
fortement cartésien qui relève $f$.

\begin{lemma}
\label{categories-lemma-composition-cartesian}
Soit $\mathcal{C}$ une catégorie.
Soit $p : \mathcal{S} \to \mathcal{C}$ une catégorie au-dessus de
$\mathcal{C}$.
\begin{enumerate}
\item Le composé de deux morphismes fortement cartésiens est fortement
cartésien.
\item Tout isomorphisme de $\mathcal{S}$ est fortement cartésien.
\item Tout morphisme fortement cartésien $\varphi$ tel que $p(\varphi)$
soit un isomorphisme est un isomorphisme.
\end{enumerate}
\end{lemma}

\begin{proof}
Démonstration de (1). Soient $\varphi : y \to x$ et
$\psi : z \to y$ des morphismes fortement cartésiens. Soit $t$ un objet
quelconque de $\mathcal{S}$. Alors
\begin{align*}
& \Mor_\mathcal{S}(t, z) \\
& =
\Mor_\mathcal{S}(t, y)
\times_{\Mor_\mathcal{C}(p(t), p(y))}
\Mor_\mathcal{C}(p(t), p(z)) \\
& =
\Mor_\mathcal{S}(t, x)
\times_{\Mor_\mathcal{C}(p(t), p(x))}
\Mor_\mathcal{C}(p(t), p(y))
\times_{\Mor_\mathcal{C}(p(t), p(y))}
\Mor_\mathcal{C}(p(t), p(z)) \\
& =
\Mor_\mathcal{S}(t, x)
\times_{\Mor_\mathcal{C}(p(t), p(x))}
\Mor_\mathcal{C}(p(t), p(z))
\end{align*}
et donc $z \to x$ est fortement cartésien.

\medskip\noindent
Démonstration de (2). Soit $y \to x$ un isomorphisme. Alors
$p(y) \to p(x)$ est également un isomorphisme. Par conséquent,
$\Mor_\mathcal{C}(p(z), p(y)) \to
\Mor_\mathcal{C}(p(z), p(x))$
est une bijection. Ainsi,
$\Mor_\mathcal{S}(z, x)
\times_{\Mor_\mathcal{C}(p(z), p(x))}
\Mor_\mathcal{C}(p(z), p(y))$ est en bijection avec
$\Mor_\mathcal{S}(z, x)$.
L'application affichée dans la
définition \ref{categories-definition-cartesian-over-C}
est donc bijective puisque $y \to x$ est un isomorphisme, et nous en
concluons que $y \to x$ est fortement cartésien.

\medskip\noindent
Démonstration de (3). Supposons que $\varphi : y \to x$ soit fortement
cartésien et que $p(\varphi) : p(y) \to p(x)$ soit un isomorphisme.
En appliquant la définition avec $z = x$, on voit que
$(\text{id}_x, p(\varphi)^{-1})$ provient d'un unique morphisme
$\chi : x \to y$. Nous omettons la vérification que $\chi$ est l'inverse
de $\varphi$.
\end{proof}

\begin{lemma}
\label{categories-lemma-cartesian-over-cartesian}
Soient $F : \mathcal{A} \to \mathcal{B}$ et
$G : \mathcal{B} \to \mathcal{C}$ des foncteurs composables entre
catégories. Soit $x \to y$ un morphisme de $\mathcal{A}$. Si
$x \to y$ est fortement $\mathcal{B}$-cartésien et si
$F(x) \to F(y)$ est fortement $\mathcal{C}$-cartésien, alors
$x \to y$ est fortement $\mathcal{C}$-cartésien.
\end{lemma}

\begin{proof}
Cela résulte directement de la définition.
\end{proof}

\begin{lemma}
\label{categories-lemma-strongly-cartesian-fibre-product}
Soit $\mathcal{C}$ une catégorie.
Soit $p : \mathcal{S} \to \mathcal{C}$ une catégorie au-dessus de
$\mathcal{C}$. Soient $x \to y$ et $z \to y$ des morphismes de
$\mathcal{S}$. Supposons que
\begin{enumerate}
\item $x \to y$ soit fortement cartésien,
\item $p(x) \times_{p(y)} p(z)$ existe, et
\item il existe un morphisme fortement cartésien $a : w \to z$ dans
$\mathcal{S}$ tel que $p(w) = p(x) \times_{p(y)} p(z)$ et
$p(a) = \text{pr}_2 : p(x) \times_{p(y)} p(z) \to p(z)$.
\end{enumerate}
Alors le produit fibré $x \times_y z$ existe et est isomorphe à $w$.
\end{lemma}

\begin{proof}
Comme $x \to y$ est fortement cartésien, il existe un unique morphisme
$b : w \to x$ tel que $p(b) = \text{pr}_1$. Pour voir que $w$ est le
produit fibré, calculons
\begin{align*}
& \Mor_\mathcal{S}(t, w) \\
& = \Mor_\mathcal{S}(t, z)
\times_{\Mor_\mathcal{C}(p(t), p(z))}
\Mor_\mathcal{C}(p(t), p(w)) \\
& = \Mor_\mathcal{S}(t, z)
\times_{\Mor_\mathcal{C}(p(t), p(z))}
(\Mor_\mathcal{C}(p(t), p(x))
\times_{\Mor_\mathcal{C}(p(t), p(y))}
\Mor_\mathcal{C}(p(t), p(z))) \\
& = \Mor_\mathcal{S}(t, z)
\times_{\Mor_\mathcal{C}(p(t), p(y))}
\Mor_\mathcal{C}(p(t), p(x)) \\
& = \Mor_\mathcal{S}(t, z)
\times_{\Mor_\mathcal{S}(t, y)}
\Mor_\mathcal{S}(t, y)
\times_{\Mor_\mathcal{C}(p(t), p(y))}
\Mor_\mathcal{C}(p(t), p(x)) \\
& = \Mor_\mathcal{S}(t, z)
\times_{\Mor_\mathcal{S}(t, y)}
\Mor_\mathcal{S}(t, x)
\end{align*}
comme voulu. La première égalité vaut parce que $a : w \to z$ est
fortement cartésien, et la dernière parce que $x \to y$ est fortement
cartésien.
\end{proof}

\begin{definition}
\label{categories-definition-fibred-category}
Soit $\mathcal{C}$ une catégorie.
Soit $p : \mathcal{S} \to \mathcal{C}$ une catégorie au-dessus de
$\mathcal{C}$. On dit que $\mathcal{S}$ est une
{\it catégorie fibrée au-dessus de $\mathcal{C}$} si, pour tout
$x \in \Ob(\mathcal{S})$ au-dessus de
$U \in \Ob(\mathcal{C})$ et tout morphisme $f : V \to U$ de
$\mathcal{C}$, il existe un morphisme fortement cartésien
$f^*x \to x$ au-dessus de $f$.
\end{definition}

\noindent
Supposons que $p : \mathcal{S} \to \mathcal{C}$ soit une catégorie fibrée.
Pour tous $f : V \to U$ et $x\in \Ob(\mathcal{S}_U)$ comme dans la
définition, nous pouvons choisir un morphisme fortement cartésien
$f^\ast x \to x$ au-dessus de $f$. Par l'axiome du choix, nous pouvons
choisir simultanément $f^*x \to x$ pour tous les
$f: V \to U = p(x)$. Nous affirmons que, pour tout morphisme
$\phi : x \to x'$ dans $\mathcal{S}_U$ et tout $f : V \to U$, il existe
un unique morphisme $f^\ast \phi : f^\ast x \to f^\ast x'$ dans
$\mathcal{S}_V$ tel que
$$
\xymatrix{
f^\ast x \ar[r]_{f^\ast \phi} \ar[d] & f^\ast x' \ar[d] \\
x \ar[r]^{\phi} & x' }
$$
soit commutatif. En effet, la flèche existe et est unique parce que
$f^*x' \to x'$ est fortement cartésien. L'unicité de cette flèche garantit
que $f^\ast$ (désormais également défini sur les morphismes) est un
foncteur $ f^\ast : \mathcal{S}_U \to \mathcal{S}_V$.

\begin{definition}
\label{categories-definition-pullback-functor-fibred-category}
Supposons que $p : \mathcal{S} \to \mathcal{C}$ soit une catégorie fibrée.
\begin{enumerate}
\item Un {\it choix de changements de base}\footnote{Cette terminologie
n'est probablement pas standard. Dans certains textes, on parle de
« clivage », mais ce terme évoque une image inadéquate. Peut-être
« opération de clivage » conviendrait-il mieux. Une notion connexe est
celle de « scindage », mais, dans beaucoup de textes, un « scindage »
désigne un choix de changements de base tel que
$g^*f^* = (f \circ g)^*$ pour toute paire de morphismes composables.
Comparer aussi à la définition \ref{categories-definition-split-fibred-category}.}
pour $p : \mathcal{S} \to \mathcal{C}$
consiste à choisir un morphisme fortement cartésien
$f^\ast x \to x$ au-dessus de $f$ pour tout morphisme
$f: V \to U$ de $\mathcal{C}$ et tout $x \in \Ob(\mathcal{S}_U)$.
\item Étant donné un choix de changements de base, pour tout morphisme
$f : V \to U$ de $\mathcal{C}$, le foncteur
$f^* : \mathcal{S}_U \to \mathcal{S}_V$ décrit ci-dessus est appelé
{\it foncteur de changement de base} (associé aux choix
$f^*x \to x$ faits ci-dessus).
\end{enumerate}
\end{definition}

\noindent
Nous pouvons évidemment toujours supposer que notre choix de changements
de base vérifie $\text{id}_U^*x = x$, bien qu'en pratique cette propriété
soit inutile si l'on n'impose pas d'autres hypothèses aux changements de base.

\begin{lemma}
\label{categories-lemma-fibred}
Supposons que $p : \mathcal{S} \to \mathcal{C}$ soit une catégorie fibrée.
Supposons que l'on ait choisi des changements de base pour
$p : \mathcal{S} \to \mathcal{C}$.
\begin{enumerate}
\item Pour toute paire de morphismes composables $f : V \to U$,
$g : W \to V$, il existe un unique isomorphisme
$$
\alpha_{g, f} :
(f \circ g)^\ast
\longrightarrow
g^\ast \circ f^\ast
$$
de foncteurs $\mathcal{S}_U \to \mathcal{S}_W$ tel que, pour tout
$y\in \Ob(\mathcal{S}_U)$, le diagramme suivant soit commutatif :
$$
\xymatrix{
g^\ast f^\ast y \ar[r]
&
f^\ast y \ar[d] \\
(f \circ g)^\ast y \ar[r]
\ar[u]^{(\alpha_{g, f})_y}
&
y
}
$$
\item Si $f = \text{id}_U$, il existe un isomorphisme canonique
$\alpha_U : \text{id} \to (\text{id}_U)^*$ de foncteurs
$\mathcal{S}_U \to \mathcal{S}_U$.
\item Le quadruplet
$(U \mapsto \mathcal{S}_U, f \mapsto f^*, \alpha_{g, f}, \alpha_U)$
définit un pseudo-foncteur de $\mathcal{C}^{opp}$ vers la
$(2, 1)$-catégorie des catégories, voir la
définition \ref{categories-definition-functor-into-2-category}.
\end{enumerate}
\end{lemma}

\begin{proof}
En fait, il est clair que le diagramme commutatif de la partie (1)
détermine de manière unique le morphisme
$(\alpha_{g, f})_y$ dans la catégorie fibre
$\mathcal{S}_W$. C'est un isomorphisme, puisque tant le morphisme
$(f \circ g)^*y \to y$ que le composé
$g^*f^*y \to f^*y \to y$ sont des morphismes fortement cartésiens
relevant $f \circ g$ (voir la discussion qui suit la
définition \ref{categories-definition-cartesian-over-C} et le
lemme \ref{categories-lemma-composition-cartesian}). De même, comme
$\text{id}_x : x \to x$ est manifestement fortement cartésien
au-dessus de $\text{id}_U$ (avec $U = p(x)$), nous voyons qu'il existe
un isomorphisme $(\alpha_U)_x : x \to (\text{id}_U)^*x$.
(Nous aurions évidemment pu supposer au préalable que $f^*x = x$
chaque fois que $f$ est un morphisme identité, mais il vaut mieux,
par souci de généralité, ne pas faire cette hypothèse.)
Nous omettons la vérification que $\alpha_{g, f}$ et
$\alpha_U$ ainsi obtenus sont des transformations de foncteurs.
Nous omettons aussi la vérification de (3).
\end{proof}

\begin{lemma}
\label{categories-lemma-fibred-equivalent}
Soit $\mathcal{C}$ une catégorie.
Soient $\mathcal{S}_1$, $\mathcal{S}_2$ des catégories au-dessus de
$\mathcal{C}$. Supposons que $\mathcal{S}_1$ et $\mathcal{S}_2$ soient
équivalentes comme catégories au-dessus de $\mathcal{C}$.
Alors $\mathcal{S}_1$ est fibrée au-dessus de $\mathcal{C}$ si et seulement
si $\mathcal{S}_2$ est fibrée au-dessus de $\mathcal{C}$.
\end{lemma}

\begin{proof}
Notons $p_i : \mathcal{S}_i \to \mathcal{C}$ les foncteurs donnés.
Soient $F : \mathcal{S}_1 \to \mathcal{S}_2$,
$G : \mathcal{S}_2 \to \mathcal{S}_1$ des foncteurs au-dessus de
$\mathcal{C}$, et soient
$i : F \circ G \to \text{id}_{\mathcal{S}_2}$,
$j : G \circ F \to \text{id}_{\mathcal{S}_1}$ des isomorphismes de
foncteurs au-dessus de $\mathcal{C}$.
Nous affirmons que, dans ce cas, $F$ envoie les morphismes fortement
cartésiens sur des morphismes fortement cartésiens. En effet, supposons que
$\varphi : y \to x$ soit fortement cartésien dans $\mathcal{S}_1$.
Posons $f : V \to U$ ; il s'agit de $p_1(\varphi)$. Supposons que
$z' \in \Ob(\mathcal{S}_2)$, avec $W = p_2(z')$, et que l'on se soit donné
$g : W \to V$ et $\psi' : z' \to F(x)$ tels que
$p_2(\psi') = f \circ g$. Alors
$$
\psi = j \circ G(\psi') : G(z') \to G(F(x)) \to x
$$
est un morphisme dans $\mathcal{S}_1$ tel que
$p_1(\psi) = f \circ g$.
Par hypothèse, il existe donc un unique morphisme $\xi : G(z') \to y$
au-dessus de $g$ tel que $\psi = \varphi \circ \xi$. Celui-ci fournit à son
tour un morphisme
$$
\xi' = F(\xi) \circ i^{-1} : z' \to F(G(z')) \to F(y)
$$
au-dessus de $g$ tel que $\psi' = F(\varphi) \circ \xi'$. Nous omettons
la vérification que $\xi'$ est unique.
\end{proof}

\noindent
Le lemme \ref{categories-lemma-fibred-equivalent} permet de conclure que les équivalences
envoient les morphismes fortement cartésiens sur des morphismes fortement
cartésiens. Cela peut toutefois être faux pour un foncteur quelconque entre
catégories fibrées au-dessus de $\mathcal{C}$. Nous définissons donc la
$2$-catégorie des catégories fibrées comme suit.

\begin{definition}
\label{categories-definition-fibred-categories-over-C}
Soit $\mathcal{C}$ une catégorie.
La {\it $2$-catégorie des catégories fibrées au-dessus de $\mathcal{C}$}
est la sous-$2$-catégorie de la $2$-catégorie des catégories
au-dessus de $\mathcal{C}$ (voir la
définition \ref{categories-definition-categories-over-C}) définie comme suit :
\begin{enumerate}
\item Ses objets sont les catégories fibrées
$p : \mathcal{S} \to \mathcal{C}$.
\item Ses $1$-morphismes $(\mathcal{S}, p) \to (\mathcal{S}', p')$
sont les foncteurs $G : \mathcal{S} \to \mathcal{S}'$ tels que
$p' \circ G = p$ et tels que $G$ envoie les morphismes fortement cartésiens
sur des morphismes fortement cartésiens.
\item Ses $2$-morphismes $t : G \to H$, pour
$G, H : (\mathcal{S}, p) \to (\mathcal{S}', p')$, sont les morphismes de
foncteurs tels que $p'(t_x) = \text{id}_{p(x)}$
pour tout $x \in \Ob(\mathcal{S})$.
\end{enumerate}
Dans cette situation, nous noterons
$$
\Mor_{\textit{Fib}/\mathcal{C}}(\mathcal{S}, \mathcal{S}')
$$
la catégorie des $1$-morphismes entre
$(\mathcal{S}, p)$ et $(\mathcal{S}', p')$.
\end{definition}

\noindent
Notons la condition imposée aux $1$-morphismes.
Notons aussi qu'il s'agit d'une véritable $2$-catégorie et non d'une
$(2, 1)$-catégorie. Ainsi, lorsque nous prenons des $2$-produits fibrés,
nous passons d'abord à la $(2, 1)$-catégorie associée.

\begin{lemma}
\label{categories-lemma-2-product-fibred-categories-over-C}
Soit $\mathcal{C}$ une catégorie.
La $(2, 1)$-catégorie des catégories fibrées au-dessus de $\mathcal{C}$
possède des 2-produits fibrés, qui sont décrits comme dans le
lemme \ref{categories-lemma-2-product-categories-over-C}.
\end{lemma}

\begin{proof}
Il s'agit essentiellement de montrer que, si
$F : \mathcal{X} \to \mathcal{S}$ et
$G : \mathcal{Y} \to \mathcal{S}$ sont des morphismes de catégories
fibrées au-dessus de $\mathcal{C}$, alors la catégorie
$\mathcal{X} \times_\mathcal{S} \mathcal{Y}$
décrite dans le lemme \ref{categories-lemma-2-product-categories-over-C} est fibrée.
Montrons que $\mathcal{X} \times_\mathcal{S} \mathcal{Y}$
possède suffisamment de morphismes fortement cartésiens.
Soit donc $(U, x, y, \phi)$ un objet de
$\mathcal{X} \times_\mathcal{S} \mathcal{Y}$.
Soit $f : V \to U$ un morphisme de $\mathcal{C}$.
Choisissons des morphismes fortement cartésiens $a : f^*x \to x$ dans
$\mathcal{X}$ au-dessus de $f$ et $b : f^*y \to y$ dans $\mathcal{Y}$
au-dessus de $f$.
Par hypothèse, $F(a)$ et $G(b)$ sont fortement cartésiens.
Comme $\phi : F(x) \to G(y)$ est un isomorphisme, l'unicité des morphismes
fortement cartésiens fournit un unique isomorphisme
$f^*\phi : F(f^*x) \to G(f^*y)$ tel que
$G(b) \circ f^*\phi = \phi \circ F(a)$. Autrement dit,
$(a, b) : (V, f^*x, f^*y, f^*\phi) \to (U, x, y, \phi)$
est un morphisme dans $\mathcal{X} \times_\mathcal{S} \mathcal{Y}$.
Nous omettons la vérification que c'est un morphisme fortement cartésien
(et que ce sont en fait les seuls morphismes fortement cartésiens).
\end{proof}

\begin{lemma}
\label{categories-lemma-cute}
Soit $\mathcal{C}$ une catégorie. Soit $U \in \Ob(\mathcal{C})$.
Si $p : \mathcal{S} \to \mathcal{C}$ est une catégorie fibrée
et si $p$ se factorise par $p' : \mathcal{S} \to \mathcal{C}/U$,
alors $p' : \mathcal{S} \to \mathcal{C}/U$ est une catégorie fibrée.
\end{lemma}

\begin{proof}
Supposons que $\varphi : x' \to x$ soit fortement cartésien relativement
à $p$. Nous affirmons que $\varphi$ est également fortement cartésien
relativement à $p'$. Posons $g = p'(\varphi)$, de sorte que
$g : V'/U \to V/U$ pour certains morphismes $f : V \to U$ et
$f' : V' \to U$.
Soit $z \in \Ob(\mathcal{S})$. Posons $p'(z) = (W \to U)$.
Pour montrer que $\varphi$ est fortement cartésien pour $p'$, nous devons
montrer que l'application
$$
\Mor_\mathcal{S}(z, x')
\longrightarrow
\Mor_\mathcal{S}(z, x)
\times_{\Mor_{\mathcal{C}/U}(W/U, V/U)}
\Mor_{\mathcal{C}/U}(W/U, V'/U),
$$
donnée par $\psi' \longmapsto (\varphi \circ \psi', p'(\psi'))$,
est bijective. Soit $(\psi, h)$ un élément du membre de droite. Alors, en
particulier, $g \circ h = p'(\psi)$ et, puisque $\varphi$ est fortement
cartésien, nous obtenons un unique morphisme $\psi' : z \to x'$ tel que
$\psi = \varphi \circ \psi'$ et tel que $p(\psi')$ égale l'application
sous-jacente à $h$. Or $p'(\psi') : W/U \to V'/U$ est un morphisme
dont l'application sous-jacente $W \to V'$ est l'application sous-jacente
de $h$ ; il est donc égal à $h$ comme morphisme dans $\mathcal{C}/U$.
Ainsi, $\psi'$ est l'unique morphisme $z \to x'$ qui s'envoie sur la paire
donnée $(\psi, h)$. Cela démontre l'affirmation.

\medskip\noindent
Enfin, soient $g : V'/U \to V/U$ et $x$ tel que $p'(x) = V/U$.
Puisque $p : \mathcal{S} \to \mathcal{C}$ est une catégorie fibrée,
il existe un morphisme fortement cartésien $\varphi : x' \to x$
tel que $p(\varphi)$ soit l'application sous-jacente de $g$.
Le même argument que ci-dessus montre que
$p'(\varphi) = g : V'/U \to V/U$. Comme nous l'avons vu plus haut,
le morphisme $\varphi$ est fortement cartésien. Les conditions de la
définition \ref{categories-definition-fibred-category} sont donc satisfaites, ce qui
achève la preuve.
\end{proof}

\begin{lemma}
\label{categories-lemma-fibred-over-fibred}
Soient $\mathcal{A} \to \mathcal{B} \to \mathcal{C}$ des foncteurs entre
catégories. Si $\mathcal{A}$ est fibrée au-dessus de $\mathcal{B}$ et si
$\mathcal{B}$ est fibrée au-dessus de $\mathcal{C}$, alors $\mathcal{A}$
est fibrée au-dessus de $\mathcal{C}$.
\end{lemma}

\begin{proof}
Cela résulte des définitions et du
lemme \ref{categories-lemma-cartesian-over-cartesian}.
\end{proof}

\begin{lemma}
\label{categories-lemma-fibred-category-representable-goes-up}
Soit $p : \mathcal{S} \to \mathcal{C}$ une catégorie fibrée.
Soient $x \to y$ et $z \to y$ des morphismes de $\mathcal{S}$,
où $x \to y$ est fortement cartésien. Si
$p(x) \times_{p(y)} p(z)$ existe, alors $x \times_y z$ existe,
$p(x \times_y z) = p(x) \times_{p(y)} p(z)$, et
$x \times_y z \to z$ est fortement cartésien.
\end{lemma}

\begin{proof}
Choisissons un morphisme fortement cartésien
$\text{pr}_2^*z \to z$ au-dessus de
$\text{pr}_2 : p(x) \times_{p(y)} p(z) \to p(z)$. Alors
$\text{pr}_2^*z = x \times_y z$ par le
lemme \ref{categories-lemma-strongly-cartesian-fibre-product}.
\end{proof}

\begin{lemma}
\label{categories-lemma-ameliorate-morphism-fibred-categories}
Soit $\mathcal{C}$ une catégorie. Soit
$F : \mathcal{X} \to \mathcal{Y}$ un $1$-morphisme de catégories fibrées
au-dessus de $\mathcal{C}$.
Il existe des $1$-morphismes de catégories fibrées au-dessus de
$\mathcal{C}$
$$
\xymatrix{
\mathcal{X} \ar@<1ex>[r]^u &
\mathcal{X}' \ar[r]^v \ar@<1ex>[l]^w & \mathcal{Y}
}
$$
tels que $F = v \circ u$ et tels que
\begin{enumerate}
\item $u : \mathcal{X} \to \mathcal{X}'$ soit pleinement fidèle,
\item $w$ soit adjoint à gauche de $u$, et
\item le foncteur $v : \mathcal{X}' \to \mathcal{Y}$ définisse une
catégorie fibrée.
\end{enumerate}
\end{lemma}

\begin{proof}
Notons $p : \mathcal{X} \to \mathcal{C}$ et
$q : \mathcal{Y} \to \mathcal{C}$ les foncteurs structuraux.
Nous construisons explicitement $\mathcal{X}'$ comme suit.
Un objet de $\mathcal{X}'$ est un quadruplet $(U, x, y, f)$ où
$x \in \Ob(\mathcal{X}_U)$, $y \in \Ob(\mathcal{Y}_U)$
et $f : y \to F(x)$ est un morphisme dans $\mathcal{Y}_U$.
Un morphisme $(a, b) : (U, x, y, f) \to (U', x', y', f')$ est donné
par $a : x \to x'$ et $b : y \to y'$ tels que
$p(a) = q(b) : U \to U'$ et tels que
$f' \circ b = F(a) \circ f$.

\medskip\noindent
Faisons un choix de changements de base pour $p$ et pour $q$, et employons
la même notation pour les désigner.
Soit $(U, x, y, f)$ un objet et soit $h : V \to U$ un morphisme.
Considérons le morphisme
$c : (V, h^*x, h^*y, h^*f) \to (U, x, y, f)$
provenant des morphismes fortement cartésiens donnés
$h^*x \to x$ et $h^*y \to y$.
Nous affirmons que $c$ est fortement cartésien dans $\mathcal{X}'$
au-dessus de $\mathcal{C}$.
Supposons en effet que l'on se soit donné un objet
$(W, x', y', f')$ de $\mathcal{X}'$, un morphisme
$(a, b) : (W, x', y', f') \to (U, x, y, f)$ au-dessus de
$W \to U$, et une factorisation $W \to V \to U$ de $W \to U$ par $h$.
Comme $h^*x \to x$ et $h^*y \to y$ sont fortement cartésiens, nous obtenons
des morphismes $a' : x' \to h^*x$ et $b' : y' \to h^*y$ au-dessus du
morphisme donné $W \to V$. Considérons le diagramme
$$
\xymatrix{
y' \ar[d]_{f'} \ar[r] & h^*y \ar[r] \ar[d]_{h^*f} & y \ar[d]_f \\
F(x') \ar[r] & F(h^*x) \ar[r] & F(x)
}
$$
Le rectangle extérieur et le carré de droite sont commutatifs.
Puisque $F$ est un $1$-morphisme de catégories fibrées, le morphisme
$F(h^*x) \to F(x)$ est fortement cartésien.
Le carré de gauche est donc commutatif par la propriété universelle des
morphismes fortement cartésiens. Cela montre que $\mathcal{X}'$
est fibrée au-dessus de $\mathcal{C}$.

\medskip\noindent
Le foncteur $u : \mathcal{X} \to \mathcal{X}'$ est donné par
$x \mapsto (p(x), x, F(x), \text{id})$. Il est pleinement fidèle.
Le foncteur $\mathcal{X}' \to \mathcal{Y}$ est donné par
$(U, x, y, f) \mapsto y$. Le foncteur
$w : \mathcal{X}' \to \mathcal{X}$ est donné par
$(U, x, y, f) \mapsto x$. Chacun de ces foncteurs est un
$1$-morphisme de catégories fibrées au-dessus de $\mathcal{C}$ d'après
notre description des morphismes fortement cartésiens de $\mathcal{X}'$
au-dessus de $\mathcal{C}$. Dire que $w$ est adjoint à gauche de $u$ signifie que
$$
\Mor_\mathcal{X}(x, x') =
\Mor_{\mathcal{X}'}((U, x, y, f), (p(x'), x', F(x'), \text{id})),
$$
ce qui résulte immédiatement des définitions.

\medskip\noindent
Enfin, nous devons montrer que $\mathcal{X}' \to \mathcal{Y}$ est une
catégorie fibrée. Soit $c : y' \to y$ un morphisme dans $\mathcal{Y}$
et soit $(U, x, y, f)$ un objet de $\mathcal{X}'$ au-dessus de $y$.
Posons $V = q(y')$ et $h = q(c) : V \to U$. Soient
$a : h^*x \to x$ et $b : h^*y \to y$ les morphismes fortement cartésiens
au-dessus de $h$.
Comme $F$ est un $1$-morphisme de catégories fibrées, nous pouvons identifier
$h^*F(x) = F(h^*x)$, le morphisme fortement cartésien correspondant étant
$F(a) : F(h^*x) \to F(x)$. Par la propriété universelle de
$b : h^*y \to y$, il existe un morphisme $c' : y' \to h^*y$ dans
$\mathcal{Y}_V$ tel que $c = b \circ c'$. Nous affirmons que
$$
(a, c) : (V, h^*x, y', h^*f \circ c') \longrightarrow (U, x, y, f)
$$
est fortement cartésien dans $\mathcal{X}'$ au-dessus de $\mathcal{Y}$.
Pour le voir, soit $(W, x_1, y_1, f_1)$ un objet de $\mathcal{X}'$, soit
$(a_1, b_1) : (W, x_1, y_1, f_1) \to (U, x, y, f)$ un morphisme,
et supposons que $b_1 = c \circ b_1'$ pour un morphisme
$b_1' : y_1 \to y'$. Alors
$$
(a_1', b_1') : (W, x_1, y_1, f_1) \longrightarrow (V, h^*x, y', h^*f \circ c')
$$
(où $a_1' : x_1 \to h^*x$ est l'unique morphisme au-dessus du morphisme donné
$q(b_1') : W \to V$ tel que $a_1 = a \circ a_1'$)
est le morphisme recherché.
\end{proof}
\section{Inertie}
\label{categories-section-inertia}

\noindent
Étant données des catégories fibrées $p : \mathcal{S} \to \mathcal{C}$ et
$p' : \mathcal{S}' \to \mathcal{C}$ au-dessus d'une catégorie $\mathcal{C}$
ainsi qu'un $1$-morphisme $F : \mathcal{S} \to \mathcal{S}'$,
nous disposons du morphisme diagonal
$$
\Delta = \Delta_{\mathcal{S}/\mathcal{S}'} :
\mathcal{S} \longrightarrow \mathcal{S} \times_{\mathcal{S}'} \mathcal{S}
$$
dans la $(2, 1)$-catégorie des catégories fibrées au-dessus de $\mathcal{C}$.

\begin{lemma}
\label{categories-lemma-inertia-fibred-category}
Soit $\mathcal{C}$ une catégorie. Soient
$p : \mathcal{S} \to \mathcal{C}$ et
$p' : \mathcal{S}' \to \mathcal{C}$ des catégories fibrées.
Soit $F : \mathcal{S} \to \mathcal{S}'$ un $1$-morphisme de
catégories fibrées au-dessus de $\mathcal{C}$. Considérons la catégorie
$\mathcal{I}_{\mathcal{S}/\mathcal{S}'}$ au-dessus de $\mathcal{C}$ dont
\begin{enumerate}
\item les objets sont les paires $(x, \alpha)$ où $x \in \Ob(\mathcal{S})$
et $\alpha : x \to x$ est un automorphisme tel que $F(\alpha) = \text{id}$,
\item les morphismes $(x, \alpha) \to (y, \beta)$ sont donnés par les morphismes
$\phi : x \to y$ tels que le diagramme
$$
\xymatrix{
x\ar[r]_\phi\ar[d]_\alpha &
y\ar[d]^{\beta} \\
x\ar[r]^\phi &
y \\
}
$$
soit commutatif, et
\item le foncteur $\mathcal{I}_{\mathcal{S}/\mathcal{S}'} \to \mathcal{C}$
est donné par $(x, \alpha) \mapsto p(x)$.
\end{enumerate}
Alors
\begin{enumerate}
\item il existe une équivalence
$$
\mathcal{I}_{\mathcal{S}/\mathcal{S}'} \longrightarrow
\mathcal{S}
\times_{\Delta, (\mathcal{S} \times_{\mathcal{S}'} \mathcal{S}), \Delta}
\mathcal{S}
$$
dans la $(2, 1)$-catégorie des catégories au-dessus de $\mathcal{C}$, et
\item $\mathcal{I}_{\mathcal{S}/\mathcal{S}'}$ est une catégorie fibrée au-dessus de
$\mathcal{C}$.
\end{enumerate}
\end{lemma}
\begin{proof}
Remarquons que (2) découle de (1) par les
lemmes \ref{categories-lemma-2-product-fibred-categories-over-C} et
\ref{categories-lemma-fibred-equivalent}. Il suffit donc de prouver (1).
Nous utiliserons sans autre mention la construction du $2$-produit fibré
donnée dans le
lemme \ref{categories-lemma-2-product-fibred-categories-over-C}.
En particulier, un objet de
$\mathcal{S}
\times_{\Delta, (\mathcal{S} \times_{\mathcal{S}'} \mathcal{S}), \Delta}
\mathcal{S}$
est un triplet $(x, y, (\iota, \kappa))$ où $x$ et $y$ sont des objets de
$\mathcal{S}$, et
$(\iota, \kappa) : (x, x, \text{id}_{F(x)}) \to (y, y, \text{id}_{F(y)})$
est un isomorphisme dans $\mathcal{S} \times_{\mathcal{S}'} \mathcal{S}$.
Cela signifie simplement que $\iota, \kappa : x \to y$ sont des isomorphismes et que
$F(\iota) = F(\kappa)$. Considérons le foncteur
$$
\mathcal{I}_{\mathcal{S}/\mathcal{S}'}
\longrightarrow
\mathcal{S}
\times_{\Delta, (\mathcal{S} \times_{\mathcal{S}'} \mathcal{S}), \Delta}
\mathcal{S}
$$
qui, à un objet $(x, \alpha)$ du membre de gauche, associe l'objet
$(x, x, (\alpha, \text{id}_x))$ du membre de droite
et, à un morphisme $\phi$ du membre de gauche,
associe le morphisme $(\phi, \phi)$ du membre de droite.
Nous affirmons qu'un quasi-inverse de ce foncteur est donné par le
foncteur
$$
\mathcal{S}
\times_{\Delta, (\mathcal{S} \times_{\mathcal{S}'} \mathcal{S}), \Delta}
\mathcal{S}
\longrightarrow
\mathcal{I}_{\mathcal{S}/\mathcal{S}'}
$$
qui, à un objet $(x, y, (\iota, \kappa))$ du membre de gauche,
associe l'objet $(x, \kappa^{-1} \circ \iota)$ du membre de droite
et, à un morphisme
$(\phi, \phi') : (x, y, (\iota, \kappa)) \to (z, w, (\lambda, \mu))$
du membre de gauche, associe le morphisme $\phi$.
En effet, l'endofoncteur de $\mathcal{I}_{\mathcal{S}/\mathcal{S}'}$ obtenu
en composant les deux foncteurs ci-dessus est strictement l'identité, et
l'endofoncteur induit sur
$\mathcal{S}
\times_{\Delta, (\mathcal{S} \times_{\mathcal{S}'} \mathcal{S}), \Delta}
\mathcal{S}$
est isomorphe à
l'identité par l'isomorphisme naturel
$$
(\text{id}_x, \kappa) :
(x, x, (\kappa^{-1} \circ \iota, \text{id}_x))
\longrightarrow
(x, y, (\iota, \kappa)).
$$
Certains détails sont omis.
\end{proof}

\begin{definition}
\label{categories-definition-inertia-fibred-category}
Soit $\mathcal{C}$ une catégorie.
\begin{enumerate}
\item Soit $F : \mathcal{S} \to \mathcal{S}'$ un $1$-morphisme de
catégories fibrées au-dessus de $\mathcal{C}$. L'{\it inertie relative
de $\mathcal{S}$ au-dessus de $\mathcal{S}'$} est la catégorie fibrée
$\mathcal{I}_{\mathcal{S}/\mathcal{S}'} \to \mathcal{C}$ du
lemme \ref{categories-lemma-inertia-fibred-category}.
\item Par {\it catégorie fibrée d'inertie $\mathcal{I}_\mathcal{S}$
de $\mathcal{S}$}, nous entendons
$\mathcal{I}_\mathcal{S} = \mathcal{I}_{\mathcal{S}/\mathcal{C}}$.
\end{enumerate}
\end{definition}

\noindent
Remarquons qu'il existe des $1$-morphismes canoniques
\begin{equation}
\label{categories-equation-inertia-structure-map}
\mathcal{I}_{\mathcal{S}/\mathcal{S}'} \longrightarrow \mathcal{S}
\quad\text{et}\quad
\mathcal{I}_\mathcal{S} \longrightarrow \mathcal{S}
\end{equation}
de catégories fibrées au-dessus de $\mathcal{C}$. Dans la description du
lemme \ref{categories-lemma-inertia-fibred-category},
ils envoient simplement l'objet $(x, \alpha)$ sur l'objet $x$ et le morphisme
$\phi : (x, \alpha) \to (y, \beta)$ sur le morphisme $\phi : x \to y$.
Il existe aussi une {\it section neutre}
\begin{equation}
\label{categories-equation-neutral-section}
e : \mathcal{S} \to \mathcal{I}_{\mathcal{S}/\mathcal{S}'}
\quad\text{et}\quad
e : \mathcal{S} \to \mathcal{I}_\mathcal{S}
\end{equation}
définie par les règles $x \mapsto (x, \text{id}_x)$ et
$(\phi : x \to y) \mapsto \phi$. C'est un inverse à droite de
(\ref{categories-equation-inertia-structure-map}). Étant donné un carré $2$-commutatif
$$
\xymatrix{
\mathcal{S}_1 \ar[d]_{F_1} \ar[r]_G & \mathcal{S}_2 \ar[d]^{F_2} \\
\mathcal{S}'_1 \ar[r]^{G'} & \mathcal{S}'_2
}
$$
il existe des {\it morphismes de fonctorialité}
\begin{equation}
\label{categories-equation-functorial}
\mathcal{I}_{\mathcal{S}_1/\mathcal{S}'_1}
\longrightarrow
\mathcal{I}_{\mathcal{S}_2/\mathcal{S}'_2}
\quad\text{et}\quad
\mathcal{I}_{\mathcal{S}_1}
\longrightarrow
\mathcal{I}_{\mathcal{S}_2}
\end{equation}
définis par les règles $(x, \alpha) \mapsto (G(x), G(\alpha))$
et $\phi \mapsto G(\phi)$. En particulier, il existe toujours un
morphisme de comparaison
\begin{equation}
\label{categories-equation-comparison}
\mathcal{I}_{\mathcal{S}/\mathcal{S}'}
\longrightarrow
\mathcal{I}_\mathcal{S}
\end{equation}
et tous les morphismes ci-dessus sont compatibles avec celui-ci.

\begin{lemma}
\label{categories-lemma-relative-inertia-as-fibre-product}
Soit $F : \mathcal{S} \to \mathcal{S}'$ un $1$-morphisme de catégories
fibrées au-dessus d'une catégorie $\mathcal{C}$. Alors le diagramme
$$
\xymatrix{
\mathcal{I}_{\mathcal{S}/\mathcal{S}'}
\ar[d]_{F \circ (\ref{categories-equation-inertia-structure-map})}
\ar[rr]_{(\ref{categories-equation-comparison})} & &
\mathcal{I}_\mathcal{S} \ar[d]^{(\ref{categories-equation-functorial})} \\
\mathcal{S}' \ar[rr]^e & &
\mathcal{I}_{\mathcal{S}'}
}
$$
est un $2$-produit fibré.
\end{lemma}

\begin{proof}
Omis.
\end{proof}





\section{Catégories fibrées en groupoïdes}
\label{categories-section-fibred-groupoids}

\noindent
Dans cette section, nous expliquons comment concevoir les catégories fibrées en groupoïdes
et nous voyons qu'elles reviennent essentiellement aux foncteurs à valeurs dans
la $(2, 1)$-catégorie des groupoïdes.

\begin{definition}
\label{categories-definition-fibred-groupoids}
Soit $p : \mathcal{S} \to \mathcal{C}$ un foncteur.
On dit que $\mathcal{S}$ est une {\it catégorie fibrée en groupoïdes} au-dessus de $\mathcal{C}$ si
les deux conditions suivantes sont satisfaites :
\begin{enumerate}
\item Pour tout morphisme $f : V \to U$ dans $\mathcal{C}$ et tout
relèvement $x$ de $U$, il existe un relèvement $\phi : y \to x$ de $f$
de but $x$.
\item Pour toute paire de morphismes $\phi : y \to x$ et $ \psi : z \to x$
et tout morphisme $f : p(z) \to p(y)$ tel que $p(\phi) \circ f = p(\psi)$,
il existe un unique relèvement $\chi : z \to y$ de $f$ tel que
$\phi \circ \chi = \psi$.
\end{enumerate}
\end{definition}

\noindent
Formulée autrement, la condition (2) dit que
l'application induite par le foncteur $p$ est une bijection entre les ensembles
de flèches pointillées dans le diagramme commutatif suivant :
$$
\xymatrix{
y \ar[r] & x & p(y) \ar[r] & p(x) \\
z \ar@{-->}[u] \ar[ru] & & p(z) \ar@{-->}[u]\ar[ru] & \\
}
$$
Voici une autre manière d'interpréter la deuxième condition.
Supposons que $g : W \to V$ et $f : V \to U$ soient des morphismes de $\mathcal{C}$.
Soit $x \in \Ob(\mathcal{S}_U)$. Par la première condition, nous pouvons relever
$f$ en $ \phi : y \to x$, puis relever $g$ en $\psi : z \to y$.
Au lieu de procéder en deux étapes, nous pouvons relever directement $f \circ g$ en
$\gamma : z' \to x$. On obtient ainsi les flèches pleines du diagramme
\begin{equation}
\label{categories-equation-fibred-groupoids}
\vcenter{
\xymatrix{
z' \ar@{-->}[d]\ar[rrd]^\gamma & & \\
z \ar@{-->}[u] \ar[r]^\psi \ar@{~>}[d]^p &
y \ar[r]^\phi \ar@{~>}[d]^p &
x \ar@{~>}[d]^p
\\
W \ar[r]^g & V \ar[r]^f & U \\
}
}
\end{equation}
où les flèches ondulées représentent non pas des morphismes, mais le foncteur $p$.
En appliquant la deuxième condition aux flèches $\phi \circ \psi$, $\gamma$
et $\text{id}_W$, nous concluons qu'il existe un unique morphisme
$\chi : z \to z'$ dans $\mathcal{S}_W$ tel que
$\gamma \circ \chi = \phi \circ \psi$. De même, il existe un unique morphisme
$z' \to z$. L'unicité implique que les morphismes $z' \to z$ et
$z\to z'$ sont inverses l'un de l'autre, autrement dit des isomorphismes.

\medskip\noindent
Cette discussion montre clairement qu'une
catégorie fibrée en groupoïdes est très étroitement liée
à une catégorie fibrée. Voici le résultat.

\begin{lemma}
\label{categories-lemma-fibred-groupoids}
Soit $p : \mathcal{S} \to \mathcal{C}$ un foncteur.
Les assertions suivantes sont équivalentes :
\begin{enumerate}
\item $p : \mathcal{S} \to \mathcal{C}$ est une catégorie
fibrée en groupoïdes, et
\item toutes les catégories fibres sont des groupoïdes et
$\mathcal{S}$ est une catégorie fibrée au-dessus de $\mathcal{C}$.

\end{enumerate}
De plus, dans ce cas, tout morphisme de $\mathcal{S}$ est
fortement cartésien. En outre, étant donnés des $f^\ast x \to x$
au-dessus de $f$ pour tous les $f: V \to U = p(x)$, les données
$(U \mapsto \mathcal{S}_U, f \mapsto f^*, \alpha_{g, f}, \alpha_U)$
construites dans le lemme \ref{categories-lemma-fibred}
définissent un pseudo-foncteur de $\mathcal{C}^{opp}$ vers
la $(2, 1)$-catégorie des groupoïdes.
\end{lemma}

\begin{proof}
Supposons que le foncteur $p : \mathcal{S} \to \mathcal{C}$ définisse une catégorie fibrée en groupoïdes.
Pour montrer que toutes les catégories fibres $\mathcal{S}_U$, pour
$U \in \Ob(\mathcal{C})$,
sont des groupoïdes, nous devons exhiber, pour tout $f : y \to x$ dans $\mathcal{S}_U$, un
morphisme inverse. Le diagramme de gauche (dans $\mathcal{S}_U$) est envoyé par
$p$ sur le diagramme de droite :
$$
\xymatrix{
y \ar[r]^f & x & U \ar[r]^{\text{id}_U} & U \\
x \ar@{-->}[u] \ar[ru]_{\text{id}_x} & &
U \ar@{-->}[u]\ar[ru]_{\text{id}_U} & \\
}
$$
Puisque seul $\text{id}_U$ rend commutatif le diagramme de droite, il existe un
unique $g : x \to y$ rendant commutatif le diagramme de gauche, donc
$fg = \text{id}_x$. Par un argument analogue, il existe un unique $h : y \to x$ tel
que $gh = \text{id}_y$. Alors $fgh = f : y \to x$. Comme $fg = \text{id}_x$,
on a $h = f$. La condition (2) de la définition \ref{categories-definition-fibred-groupoids} dit
exactement que tout morphisme de $\mathcal{S}$ est fortement cartésien. Ainsi,
la condition (1) de la définition \ref{categories-definition-fibred-groupoids} implique que
$\mathcal{S}$ est une catégorie fibrée au-dessus de $\mathcal{C}$.
\medskip\noindent
Réciproquement, supposons que toutes les catégories fibres soient des groupoïdes et que
$\mathcal{S}$ soit une catégorie fibrée au-dessus de $\mathcal{C}$.
Il faut vérifier les conditions (1) et (2) de la
définition \ref{categories-definition-fibred-groupoids}.
La première condition est immédiate. Soient $\phi : y \to x$,
$\psi : z \to x$ et $f : p(z) \to p(y)$ tels que
$p(\phi) \circ f = p(\psi)$, comme dans la condition (2) de la
définition \ref{categories-definition-fibred-groupoids}.
Écrivons $U = p(x)$, $V = p(y)$, $W = p(z)$, $p(\phi) = g : V \to U$,
$p(\psi) = h : W \to U$. Choisissons un morphisme fortement cartésien $g^*x \to x$
au-dessus de $g$. Nous obtenons alors un morphisme $i : y \to g^*x$ dans
$\mathcal{S}_V$, qui est donc un isomorphisme. Nous
obtenons aussi un morphisme $j : z \to g^*x$ correspondant à
la paire $(\psi, f)$ puisque $g^*x \to x$ est fortement cartésien.
On vérifie alors que $\chi = i^{-1} \circ j$ est une solution.

\medskip\noindent
Nous avons vu dans la preuve de (1) $\Rightarrow$ (2) que
tout morphisme de $\mathcal{S}$ est fortement cartésien.
La dernière assertion découle directement du lemme \ref{categories-lemma-fibred}.
\end{proof}

\begin{lemma}
\label{categories-lemma-fibred-gives-fibred-groupoids}
Soit $\mathcal{C}$ une catégorie.
Soit $p : \mathcal{S} \to \mathcal{C}$ une catégorie fibrée.
Soit $\mathcal{S}'$ la sous-catégorie de $\mathcal{S}$ définie
comme suit :
\begin{enumerate}
\item $\Ob(\mathcal{S}') = \Ob(\mathcal{S})$, et
\item pour $x, y \in \Ob(\mathcal{S}')$, l'ensemble des morphismes entre $x$
et $y$ dans $\mathcal{S}'$ est l'ensemble des morphismes fortement cartésiens entre
$x$ et $y$ dans $\mathcal{S}$.
\end{enumerate}
Soit $p' : \mathcal{S}' \to \mathcal{C}$ la restriction de $p$
à $\mathcal{S}'$. Alors $p' : \mathcal{S}' \to \mathcal{C}$ est une catégorie
fibrée en groupoïdes.
\end{lemma}

\begin{proof}
Remarquons que la construction a un sens puisque, d'après le
lemme \ref{categories-lemma-composition-cartesian},
le morphisme identité de tout objet de $\mathcal{S}$ est fortement cartésien,
et la composée de morphismes fortement cartésiens est fortement cartésienne.
La première propriété de relèvement de la
définition \ref{categories-definition-fibred-groupoids}
découle de la condition selon laquelle, dans une catégorie fibrée,
pour tout morphisme $f : V \to U$ et tout $x$ au-dessus de $U$, il existe
un morphisme fortement cartésien $\varphi : y \to x$ au-dessus de $f$.
Vérifions la deuxième propriété de relèvement de la
définition \ref{categories-definition-fibred-groupoids}
pour la catégorie $p' : \mathcal{S}' \to \mathcal{C}$ au-dessus de $\mathcal{C}$.
Pour ce faire, raisonnons comme dans la discussion qui suit la
définition \ref{categories-definition-fibred-groupoids}.
Ainsi, dans le diagramme \ref{categories-equation-fibred-groupoids}, les
morphismes $\phi$, $\psi$ et $\gamma$ sont des morphismes fortement cartésiens
de $\mathcal{S}$.
Par conséquent, $\gamma$ et $\phi \circ \psi$ sont des morphismes fortement cartésiens
de $\mathcal{S}$ au-dessus de la même flèche de $\mathcal{C}$ et
ayant le même but dans $\mathcal{S}$. D'après la discussion qui suit la
définition \ref{categories-definition-cartesian-over-C},
ces deux flèches sont donc isomorphes, comme voulu (nous utilisons aussi ici
le fait que tout isomorphisme de $\mathcal{S}$ est fortement cartésien, d'après le
lemme \ref{categories-lemma-composition-cartesian}).
\end{proof}

\begin{example}
\label{categories-example-group-homomorphism-fibreedingroupoids}
Un homomorphisme de groupes $p : G \to H$ donne naissance à un foncteur
$p : \mathcal{S}\to\mathcal{C}$ comme dans l'exemple
\ref{categories-example-group-homomorphism-functor}. Ce foncteur
$p : \mathcal{S}\to\mathcal{C}$ définit une catégorie fibrée en groupoïdes si et seulement si
$p$ est surjectif. La catégorie fibre $\mathcal{S}_U$ au-dessus de l'unique
objet $U\in \Ob(\mathcal{C})$ est la catégorie associée au
noyau de $p$, comme dans l'exemple \ref{categories-example-group-groupoid}.
\end{example}

\noindent
Étant donné $p : \mathcal{S} \to \mathcal{C}$, on peut se demander : si la catégorie
fibre $\mathcal{S}_U$ est un groupoïde pour tout $U \in \Ob(\mathcal{C})$,
$\mathcal{S}$ est-elle nécessairement fibrée en groupoïdes au-dessus de $\mathcal{C}$ ?
La réponse est négative, comme on peut le voir ainsi. Partons d'une catégorie fibrée en
groupoïdes $p : \mathcal{S} \to \mathcal{C}$. Modifier les morphismes de $\mathcal{S}$
qui ne s'envoient pas sur le morphisme identité d'un objet ne modifie pas les
catégories $\mathcal{S}_U$. On peut donc faire échouer les conditions d'existence et
d'unicité des relèvements. Un exemple est fourni par le foncteur de l'exemple
\ref{categories-example-group-homomorphism-fibreedingroupoids} lorsque $G \to H$ n'est pas
surjectif. Voici un autre exemple.

\begin{example}
\label{categories-example-not-fibred-in-groupoids-but-fibre-cats-are}
Posons $\Ob(\mathcal{C}) = \{A, B, T\}$ et
$\Mor_\mathcal{C}(A, B) = \{f\}$, $\Mor_\mathcal{C}(B, T) = \{g\}$,
$\Mor_\mathcal{C}(A, T) = \{h\} = \{gf\}, $ auxquels s'ajoute le morphisme identité
de chaque objet. Le diagramme ci-dessous représente cette catégorie. Posons maintenant
$\Ob(\mathcal{S}) = \{A', B', T'\}$ et
$\Mor_\mathcal{S}(A', B') = \emptyset$,
$\Mor_\mathcal{S}(B', T') = \{g'\}$,
$\Mor_\mathcal{S}(A', T') = \{h'\}, $ auxquels s'ajoutent les morphismes identité. Le
foncteur $p : \mathcal{S} \to \mathcal{C}$ est évident. Alors, pour tout
$U \in \Ob(\mathcal{C})$, $\mathcal{S}_U$ est la catégorie ayant un seul
objet et le morphisme identité de cet objet, donc un groupoïde, mais le
morphisme $f: A \to B$ ne peut être relevé. De même, si l'on pose
$\Mor_\mathcal{S}(A', B') = \{f'_1, f'_2\}$ et
$ \Mor_\mathcal{S}(A', T') = \{h'\} = \{g'f'_1 \} = \{g'f'_2\}$, alors
les catégories fibres sont les mêmes et $f: A \to B$ dans le diagramme ci-dessous
possède deux relèvements.
$$
\xymatrix{
B' \ar[r]^{g'} & T' & & B \ar[r]^g & T & \\
A' \ar@{-->}[u]^{??} \ar[ru]_{h'} & & \ar@{}[u]^{\text{ci-dessus}} &
A \ar[u]^f \ar[ru]_{gf = h} & \\
}
$$
\end{example}

\noindent
Plus loin, nous voudrons formuler des assertions telles que « toute catégorie fibrée en
groupoïdes au-dessus de $\mathcal{C}$ est équivalente à une catégorie scindée », ou
« toute catégorie fibrée en groupoïdes dont les catégories fibres sont assimilables à des ensembles
est équivalente à une catégorie fibrée en ensembles ». La notion d'équivalence
dépend de la $2$-catégorie dans laquelle nous travaillons.

\begin{definition}
\label{categories-definition-categories-fibred-in-groupoids-over-C}
Soit $\mathcal{C}$ une catégorie.
La {\it $2$-catégorie des catégories fibrées en groupoïdes au-dessus de $\mathcal{C}$}
est la sous-$2$-catégorie de la $2$-catégorie des catégories fibrées
au-dessus de $\mathcal{C}$ (voir la définition \ref{categories-definition-fibred-categories-over-C})
définie comme suit :
\begin{enumerate}
\item Ses objets sont les catégories
$p : \mathcal{S} \to \mathcal{C}$ fibrées en groupoïdes.
\item Ses $1$-morphismes $(\mathcal{S}, p) \to (\mathcal{S}', p')$
sont les foncteurs $G : \mathcal{S} \to \mathcal{S}'$ tels que
$p' \circ G = p$ (puisque tout morphisme est fortement cartésien,
$G$ les préserve automatiquement).
\item Ses $2$-morphismes $t : G \to H$ pour
$G, H : (\mathcal{S}, p) \to (\mathcal{S}', p')$
sont les morphismes de foncteurs
tels que $p'(t_x) = \text{id}_{p(x)}$
pour tout $x \in \Ob(\mathcal{S})$.
\end{enumerate}
\end{definition}

\noindent
Remarquons que tout $2$-morphisme est automatiquement un isomorphisme !
Il s'agit donc en fait d'une $(2, 1)$-catégorie, et non simplement d'une
$2$-catégorie. Voici l'inévitable lemme sur les $2$-produits fibrés.

\begin{lemma}
\label{categories-lemma-2-product-fibred-categories}
Soit $\mathcal{C}$ une catégorie.
La $2$-catégorie des catégories fibrées en groupoïdes
au-dessus de $\mathcal{C}$ admet des $2$-produits fibrés, et ceux-ci sont décrits comme dans le
lemme \ref{categories-lemma-2-product-categories-over-C}.
\end{lemma}

\begin{proof}
D'après le lemme \ref{categories-lemma-2-product-fibred-categories-over-C},
le produit fibré décrit dans le
lemme \ref{categories-lemma-2-product-categories-over-C} est une catégorie fibrée.
Il suffit donc de prouver que les catégories fibres sont des
groupoïdes ; voir le lemme \ref{categories-lemma-fibred-groupoids}.
D'après le lemme \ref{categories-lemma-fibre-2-fibre-product-categories-over-C},
il suffit de montrer que le 2-produit fibré de groupoïdes
est un groupoïde, ce qui est clair (par exemple d'après la construction du
lemme \ref{categories-lemma-2-fibre-product-categories}).
\end{proof}

\begin{remark}
\label{categories-remark-2-product-fibred-categories-same}
Soit $\mathcal{C}$ une catégorie. Soient $f : \mathcal{X} \to \mathcal{S}$
et $g : \mathcal{Y} \to \mathcal{S}$ des $1$-morphismes de
catégories fibrées en groupoïdes au-dessus de $\mathcal{C}$.
Soit $p : \mathcal{S} \to \mathcal{C}$ le foncteur donné.
Nous affirmons que le $2$-produit fibré du lemme \ref{categories-lemma-2-product-fibred-categories}
est canoniquement équivalent (comme catégorie) à celui de
l'exemple \ref{categories-example-2-fibre-product-categories}.
Les objets du premier sont les quadruplets $(U, x, y, \alpha)$
où $p(\alpha) = \text{id}_U$
(voir le lemme \ref{categories-lemma-2-product-categories-over-C}),
et les objets du second sont les triplets $(x, y, \alpha)$ (voir
l'exemple \ref{categories-example-2-fibre-product-categories}).
L'équivalence entre les deux catégories est donnée par
les règles $(U, x, y, \alpha) \mapsto (x, y, \alpha)$
et $(x, y, \alpha) \mapsto (p(f(x)), x, y', \alpha')$,
où $\alpha' = g(\gamma)^{-1} \circ \alpha$ et $\gamma : y' \to y$
est un relèvement de la flèche $p(\alpha) : p(f(x)) \to p(g(y))$.
Les détails sont omis.
\end{remark}

\begin{lemma}
\label{categories-lemma-equivalence-fibred-categories}
Soient $p : \mathcal{S}\to \mathcal{C}$ et
$p' : \mathcal{S'}\to \mathcal{C}$ des catégories fibrées en groupoïdes, et
supposons que $G : \mathcal{S}\to \mathcal {S}'$ soit un foncteur au-dessus de
$\mathcal{C}$.
\begin{enumerate}
\item Alors $G$ est fidèle (resp.\ pleinement fidèle, resp.\ une équivalence)
si et seulement si, pour chaque $U\in\Ob(\mathcal{C})$, le foncteur induit
$G_U : \mathcal{S}_U\to \mathcal{S}'_U$ est fidèle
(resp.\ pleinement fidèle, resp.\ une équivalence).
\item Si $G$ est une équivalence, alors $G$ est une équivalence dans la
$2$-catégorie des catégories fibrées en groupoïdes au-dessus de $\mathcal{C}$.
\end{enumerate}
\end{lemma}

\begin{proof}
Soient $x, y$ des objets de $\mathcal{S}$ au-dessus du même objet $U$.
Considérons le diagramme commutatif
$$
\xymatrix{
\Mor_\mathcal{S}(x, y) \ar[rd]_p \ar[rr]_G & &
\Mor_{\mathcal{S}'}(G(x), G(y)) \ar[ld]^{p'} \\
& \Mor_\mathcal{C}(U, U) &
}
$$
Ce diagramme montre clairement que, si $G$ est fidèle (resp.\ pleinement fidèle),
alors chaque $G_U$ l'est aussi.

\medskip\noindent
Supposons que $G$ soit une équivalence. Pour tout objet
$x'$ de $\mathcal{S}'$, il existe un objet $x$ de $\mathcal{S}$
tel que $G(x)$ soit isomorphe à $x'$. Supposons que $x'$ soit
au-dessus de $U'$ et que $x$ soit au-dessus de $U$. Il existe alors un isomorphisme
$f : U' \to U$ dans $\mathcal{C}$, à savoir l'image par $p'$ de
l'isomorphisme $x' \to G(x)$. Par les axiomes d'une catégorie fibrée
en groupoïdes, il existe une flèche $f^*x \to x$ de $\mathcal{S}$
au-dessus de $f$. Il existe donc un isomorphisme
$\alpha : x' \to G(f^*x)$ tel que $p'(\alpha) = \text{id}_{U'}$
(cette fois par les axiomes de $\mathcal{S}'$). En définitive, nous concluons
que, pour tout objet $x'$ de $\mathcal{S}'$, nous pouvons choisir
une paire $(o_{x'}, \alpha_{x'})$ formée d'un objet
$o_{x'}$ de $\mathcal{S}$ et d'un isomorphisme $\alpha_{x'} : x' \to G(o_{x'})$
tel que $p'(\alpha_{x'}) = \text{id}_{p'(x')}$. À partir de là,
nous procédons comme d'habitude (voir la preuve du
lemme \ref{categories-lemma-equivalence-categories}) pour construire un foncteur
inverse $F : \mathcal{S}' \to \mathcal{S}$, en prenant
$x' \mapsto o_{x'}$ et $\varphi' : x' \to y'$ pour l'unique
flèche $\varphi_{\varphi'} : o_{x'} \to o_{y'}$ telle que
$\alpha_{y'}^{-1} \circ G(\varphi_{\varphi'}) \circ \alpha_{x'} = \varphi'$.
Avec ces choix, $F$ est un foncteur au-dessus de $\mathcal{C}$.
Nous omettons de vérifier que $G \circ F$ et $F \circ G$ sont
$2$-isomorphes aux foncteurs identité respectifs
(dans la $2$-catégorie des catégories fibrées en groupoïdes au-dessus de $\mathcal{C}$).

\medskip\noindent
Supposons que $G_U$ soit fidèle (resp.\ pleinement fidèle)
pour tout $U\in\Ob(\mathcal C)$. Pour
montrer que $G$ est fidèle (resp.\ pleinement fidèle),
il faut montrer, pour tous objets
$x, y\in\Ob(\mathcal{S})$, que $G$ induit une
injection (resp.\ bijection) entre
$\Mor_\mathcal{S}(x, y)$ et
$\Mor_{\mathcal{S}'}(G(x), G(y))$.
Posons $U = p(x)$ et $V = p(y)$.
Il suffit de prouver que $G$
induit une injection (resp.\ bijection) entre les morphismes
$x \to y$ au-dessus de $f$ et les morphismes $G(x) \to G(y)$ au-dessus de $f$,
pour tout morphisme $f : U \to V$.
Fixons maintenant $f : U \to V$. Choisissons un morphisme fortement cartésien
$f^*y \to y$ au-dessus de ce morphisme. Alors $G(f^*y) \to G(y)$ est lui aussi
fortement cartésien.
L'ensemble des morphismes de $x$ vers $y$ au-dessus de $f$
est en bijection avec l'ensemble des morphismes entre
$x$ et $f^*y$ au-dessus de $\text{id}_U$. (Par le deuxième axiome
d'une catégorie fibrée en groupoïdes.) De même,
l'ensemble des morphismes de $G(x)$ vers $G(y)$ au-dessus de $f$
est en bijection avec l'ensemble des morphismes entre
$G(x)$ et $G(f^*y)$ au-dessus de $\text{id}_U$.
Le fait que $G_U$ soit fidèle (resp.\ pleinement fidèle)
donne donc le résultat voulu.

\medskip\noindent
Enfin, supposons que tous les $G_U$ soient des équivalences pour tout $U$ ; ils sont donc
pleinement fidèles et essentiellement surjectifs. Nous avons vu que cela implique que $G$ est
pleinement fidèle ; pour montrer que c'est une équivalence, il reste donc à montrer qu'il
est essentiellement surjectif. C'est clair : si $z'\in
\Ob(\mathcal{S}')$, alors $z'\in \Ob(\mathcal{S}'_U)$, où
$U = p'(z')$. Puisque $G_U$ est essentiellement surjectif, nous savons que
$z'$ est isomorphe, dans $\mathcal{S}'_U$, à un objet de la forme
$G_U(z)$ pour un certain $z\in \Ob(\mathcal{S}_U)$. Or les morphismes
de $\mathcal{S}'_U$ sont des morphismes de $\mathcal{S}'$ ; ainsi $z'$ est
isomorphe à $G(z)$ dans $\mathcal{S}'$.
\end{proof}

\begin{lemma}
\label{categories-lemma-fully-faithful-diagonal-equivalence}
Soit $\mathcal{C}$ une catégorie. Soient $p : \mathcal{S}\to \mathcal{C}$ et
$p' : \mathcal{S'}\to \mathcal{C}$ des catégories fibrées en groupoïdes.
Soit $G : \mathcal{S}\to \mathcal {S}'$ un foncteur au-dessus de $\mathcal{C}$.
Alors $G$ est pleinement fidèle si et seulement si la diagonale
$$
\Delta_G :
\mathcal{S}
\longrightarrow
\mathcal{S} \times_{G, \mathcal{S}', G} \mathcal{S}
$$
est une équivalence.
\end{lemma}

\begin{proof}
D'après le
lemme \ref{categories-lemma-equivalence-fibred-categories},
il suffit de considérer les catégories fibres au-dessus d'un objet $U$ de $\mathcal{C}$.
Un objet du membre de droite est un triplet $(x, x', \alpha)$ où
$\alpha : G(x) \to G(x')$ est un morphisme de $\mathcal{S}'_U$.
Le foncteur $\Delta_G$ envoie l'objet $x$ de $\mathcal{S}_U$
sur le triplet $(x, x, \text{id}_{G(x)})$. Remarquons que $(x, x', \alpha)$
appartient à l'image essentielle de $\Delta_G$ si et seulement si $\alpha = G(\beta)$
pour un certain morphisme $\beta : x \to x'$ dans $\mathcal{S}_U$ (détails omis).
Par conséquent, pour que $\Delta_G$ soit une équivalence, tout $\alpha$ doit
être l'image d'un morphisme $\beta : x \to x'$, et deux morphismes
distincts $\beta, \beta' : x \to x'$ doivent aussi donner des morphismes distincts
$G(\beta), G(\beta')$. Cela prouve le lemme.
\end{proof}

\begin{lemma}
\label{categories-lemma-morphisms-equivalent-fibred-groupoids}
Soit $\mathcal{C}$ une catégorie.
Soient $\mathcal{S}_i$, $i = 1, 2, 3, 4$, des catégories fibrées en
groupoïdes au-dessus de $\mathcal{C}$.
Supposons que $\varphi : \mathcal{S}_1 \to \mathcal{S}_2$ et
$\psi : \mathcal{S}_3 \to \mathcal{S}_4$ soient des équivalences
au-dessus de $\mathcal{C}$. Alors
$$
\Mor_{\textit{Cat}/\mathcal{C}}(\mathcal{S}_2, \mathcal{S}_3)
\longrightarrow
\Mor_{\textit{Cat}/\mathcal{C}}(\mathcal{S}_1, \mathcal{S}_4),
\quad \alpha \longmapsto \psi \circ \alpha \circ \varphi
$$
est une équivalence de catégories.
\end{lemma}

\begin{proof}
C'est un fait général, valable dans toute $2$-catégorie.
\end{proof}

\begin{lemma}
\label{categories-lemma-inertia-fibred-groupoids}
Soit $\mathcal{C}$ une catégorie.
Si $p : \mathcal{S} \to \mathcal{C}$ est une catégorie fibrée en groupoïdes, alors
la catégorie fibrée d'inertie $\mathcal{I}_\mathcal{S} \to \mathcal{C}$ l'est aussi.
\end{lemma}

\begin{proof}
Cela ressort de la construction du
lemme \ref{categories-lemma-inertia-fibred-category},
ou bien du fait (établi dans le même lemme) que
$\mathcal{I}_\mathcal{S} \to \mathcal{S}
\times_{\Delta, \mathcal{S} \times_\mathcal{C} \mathcal{S}, \Delta}\mathcal{S}$
est une équivalence, combiné au
lemme \ref{categories-lemma-2-product-fibred-categories}.
\end{proof}

\begin{lemma}
\label{categories-lemma-cute-groupoids}
Soit $\mathcal{C}$ une catégorie. Soit $U \in \Ob(\mathcal{C})$.
Si $p : \mathcal{S} \to \mathcal{C}$ est une catégorie fibrée en groupoïdes
et si $p$ se factorise par $p' : \mathcal{S} \to \mathcal{C}/U$,
alors $p' : \mathcal{S} \to \mathcal{C}/U$ est une catégorie fibrée en groupoïdes.
\end{lemma}

\begin{proof}
Nous avons déjà vu dans le lemme \ref{categories-lemma-cute} que $p'$ est une catégorie
fibrée. Il suffit donc de prouver que les catégories fibres sont des groupoïdes ;
voir le lemme \ref{categories-lemma-fibred-groupoids}.
Pour $V \in \Ob(\mathcal{C})$, on a
$$
\mathcal{S}_V = \coprod\nolimits_{f : V \to U} \mathcal{S}_{(f : V \to U)}
$$
où le membre de gauche est la catégorie fibre de $p$ et le membre de droite
est la réunion disjointe des catégories fibres de $p'$.
D'où le résultat.
\end{proof}

\begin{lemma}
\label{categories-lemma-fibred-in-groupoids-over-fibred-in-groupoids}
Soient $\mathcal{A} \to \mathcal{B} \to \mathcal{C}$ des foncteurs entre
catégories. Si $\mathcal{A}$ est fibrée en groupoïdes au-dessus de $\mathcal{B}$
et $\mathcal{B}$ est fibrée en groupoïdes au-dessus de $\mathcal{C}$, alors
$\mathcal{A}$ est fibrée en groupoïdes au-dessus de $\mathcal{C}$.
\end{lemma}

\begin{proof}
On peut le prouver directement à partir de la définition. Toutefois, nous raisonnerons
à l'aide du critère du lemme \ref{categories-lemma-fibred-groupoids}.
D'après le lemme \ref{categories-lemma-fibred-over-fibred}, $\mathcal{A}$
est fibrée au-dessus de $\mathcal{C}$. Pour achever la preuve, montrons que la catégorie
fibre $\mathcal{A}_U$ est un groupoïde pour tout objet $U$ de $\mathcal{C}$.
En effet, si $x \to y$ est un morphisme de $\mathcal{A}_U$, alors son
image dans $\mathcal{B}$ est un isomorphisme puisque $\mathcal{B}_U$ est
un groupoïde. Mais alors $x \to y$ est un isomorphisme, par exemple d'après le
lemme \ref{categories-lemma-composition-cartesian} et le fait que tout
morphisme de $\mathcal{A}$ est fortement $\mathcal{B}$-cartésien
(voir le lemme \ref{categories-lemma-fibred-groupoids}).
\end{proof}

\begin{lemma}
\label{categories-lemma-fibred-groupoids-fibre-product-goes-up}
Soit $p : \mathcal{S} \to \mathcal{C}$ une catégorie fibrée en groupoïdes.
Soient $x \to y$ et $z \to y$ des morphismes de $\mathcal{S}$.
Si $p(x) \times_{p(y)} p(z)$ existe, alors
$x \times_y z$ existe et $p(x \times_y z) = p(x) \times_{p(y)} p(z)$.
\end{lemma}

\begin{proof}
Cela découle du
lemme \ref{categories-lemma-fibred-category-representable-goes-up}.
\end{proof}

\begin{lemma}
\label{categories-lemma-ameliorate-morphism-categories-fibred-groupoids}
Soit $\mathcal{C}$ une catégorie. Soit $F : \mathcal{X} \to \mathcal{Y}$
un $1$-morphisme de catégories fibrées en groupoïdes au-dessus de $\mathcal{C}$.
Il existe une factorisation $\mathcal{X} \to \mathcal{X}' \to \mathcal{Y}$
par des $1$-morphismes de catégories fibrées en groupoïdes au-dessus de $\mathcal{C}$,
telle que $\mathcal{X} \to \mathcal{X}'$ soit une équivalence au-dessus de $\mathcal{C}$
et que $\mathcal{X}'$ soit une catégorie fibrée en groupoïdes au-dessus de
$\mathcal{Y}$.
\end{lemma}

\begin{proof}
Notons $p : \mathcal{X} \to \mathcal{C}$ et $q : \mathcal{Y} \to \mathcal{C}$
les foncteurs structuraux. Construisons explicitement $\mathcal{X}'$ comme suit.
Un objet de $\mathcal{X}'$ est un quadruplet $(U, x, y, f)$ où
$x \in \Ob(\mathcal{X}_U)$, $y \in \Ob(\mathcal{Y}_U)$
et $f : F(x) \to y$ est un isomorphisme dans $\mathcal{Y}_U$.
Un morphisme $(a, b) : (U, x, y, f) \to (U', x', y', f')$ est donné
par $a : x \to x'$ et $b : y \to y'$, avec $p(a) = q(b)$ et
tels que $f' \circ F(a) = b \circ f$. Autrement dit,
$\mathcal{X}' = \mathcal{X} \times_{F, \mathcal{Y}, \text{id}} \mathcal{Y}$
avec la construction du $2$-produit fibré du
lemme \ref{categories-lemma-2-product-categories-over-C}.
D'après le
lemme \ref{categories-lemma-2-product-fibred-categories},
$\mathcal{X}'$ est une catégorie fibrée en groupoïdes au-dessus de
$\mathcal{C}$ et $\mathcal{X}' \to \mathcal{Y}$ est un morphisme de
catégories au-dessus de $\mathcal{C}$. Comme foncteur $\mathcal{X} \to \mathcal{X}'$, prenons
$x \mapsto (p(x), x, F(x), \text{id}_{F(x)})$ sur les objets et
$(a : x \to x') \mapsto (a, F(a))$ sur les morphismes. Il est clair que
la composée $\mathcal{X} \to \mathcal{X}' \to \mathcal{Y}$
est égale à $F$. Nous omettons de vérifier que
$\mathcal{X} \to \mathcal{X}'$ est une équivalence de catégories fibrées au-dessus de
$\mathcal{C}$.

\medskip\noindent
Enfin, il reste à montrer que $\mathcal{X}' \to \mathcal{Y}$ est une catégorie
fibrée en groupoïdes. Soit $b : y' \to y$ un morphisme de $\mathcal{Y}$
et soit $(U, x, y, f)$ un objet de $\mathcal{X}'$ au-dessus de $y$.
Puisque $\mathcal{X}$ est fibrée en groupoïdes au-dessus de $\mathcal{C}$, nous
pouvons trouver un morphisme $a : x' \to x$ au-dessus de $U' = q(y') \to q(y) = U$.
Puisque $\mathcal{Y}$ est fibrée en groupoïdes au-dessus de $\mathcal{C}$ et que
$F(x') \to F(x)$ et $y' \to y$ sont tous deux au-dessus du même morphisme $U' \to U$,
nous pouvons trouver $f' : F(x') \to y'$ au-dessus de $\text{id}_{U'}$ tel que
$f \circ F(a) = b \circ f'$. Nous obtenons ainsi
$(a, b) : (U', x', y', f') \to (U, x, y, f)$.
Ceci vérifie la première condition (1) de la
définition \ref{categories-definition-fibred-groupoids}.
Pour vérifier (2), soient
$(a, b) : (U', x', y', f') \to (U, x, y, f)$ et
$(a', b') : (U'', x'', y'', f'') \to (U, x, y, f)$ des morphismes de
$\mathcal{X}'$, et soit $b'' : y' \to y''$ un morphisme de $\mathcal{Y}$
tel que $b' \circ b'' = b$. Il faut montrer qu'il existe
un unique morphisme $a'' : x' \to x''$ tel que
$f'' \circ F(a'') = b'' \circ f'$ et tel que
$(a', b') \circ (a'', b'') = (a, b)$. Puisque $\mathcal{X}$ est fibrée
en groupoïdes, nous savons qu'il existe un unique morphisme
$a'' : x' \to x''$ tel que $a' \circ a'' = a$ et $p(a'') = q(b'')$.
Puisque $\mathcal{Y}$ est fibrée en groupoïdes, nous voyons que
$F(a'')$ est l'unique morphisme $F(x') \to F(x'')$ tel que
$F(a') \circ F(a'') = F(a)$ et $q(F(a'')) = q(b'')$. La relation
$f'' \circ F(a'') = b'' \circ f'$ découle alors de cette unicité et des relations données
$f \circ F(a) = b \circ f'$ et $f \circ F(a') = b' \circ f''$.
\end{proof}

\begin{lemma}
\label{categories-lemma-amelioration-unique}
Soit $\mathcal{C}$ une catégorie. Soit $F : \mathcal{X} \to \mathcal{Y}$
un $1$-morphisme de catégories fibrées en groupoïdes au-dessus de $\mathcal{C}$.
Supposons donné un diagramme $2$-commutatif
$$
\xymatrix{
\mathcal{X}' \ar[rd]_f &
\mathcal{X} \ar[l]^a \ar[d]^F \ar[r]_b &
\mathcal{X}'' \ar[ld]^g \\
& \mathcal{Y}
}
$$
où $a$ et $b$ sont des équivalences de catégories au-dessus de $\mathcal{C}$
et où $f$ et $g$ définissent des catégories fibrées en groupoïdes. Il existe alors
une équivalence $h : \mathcal{X}'' \to \mathcal{X}'$ de catégories au-dessus de
$\mathcal{Y}$ telle que $h \circ b$ soit $2$-isomorphe à $a$ comme $1$-morphismes
de catégories au-dessus de $\mathcal{C}$. Si le diagramme ci-dessus est effectivement
commutatif, on peut faire en sorte que $h \circ b$ soit $2$-isomorphe à $a$ comme
$1$-morphismes de catégories au-dessus de $\mathcal{Y}$.
\end{lemma}

\begin{proof}
Nous montrerons que $\mathcal{X}'$ et $\mathcal{X}''$, au-dessus de $\mathcal{Y}$,
sont toutes deux équivalentes à la catégorie fibrée en groupoïdes
$\mathcal{X} \times_{F, \mathcal{Y}, \text{id}} \mathcal{Y}$
au-dessus de $\mathcal{Y}$ ; voir la preuve du
lemme \ref{categories-lemma-ameliorate-morphism-categories-fibred-groupoids}.
Choisissons un quasi-inverse $b^{-1} : \mathcal{X}'' \to \mathcal{X}$ dans la
$2$-catégorie des catégories au-dessus de $\mathcal{C}$.
Puisque le triangle de droite du diagramme est $2$-commutatif, le diagramme
$$
\xymatrix{
\mathcal{X} \ar[d]_F & \mathcal{X}'' \ar[l]^{b^{-1}} \ar[d]^g \\
\mathcal{Y} & \mathcal{Y} \ar[l]
}
$$
est $2$-commutatif. Nous obtenons donc un $1$-morphisme
$c : \mathcal{X}'' \to
\mathcal{X} \times_{F, \mathcal{Y}, \text{id}} \mathcal{Y}$
par la propriété universelle du $2$-produit fibré. De plus, $c$
est un morphisme de catégories au-dessus de $\mathcal{Y}$ (!) et une équivalence
(puisque $b$ est supposé être une équivalence ; voir le
lemme \ref{categories-lemma-equivalence-2-fibre-product}).
Ainsi, $c$ est une équivalence dans la $2$-catégorie des catégories fibrées
en groupoïdes au-dessus de $\mathcal{Y}$ d'après le
lemme \ref{categories-lemma-equivalence-fibred-categories}.

\medskip\noindent
Il reste à construire un $2$-isomorphisme entre $c \circ b$ et
le foncteur $d : \mathcal{X} \to
\mathcal{X} \times_{F, \mathcal{Y}, \text{id}} \mathcal{Y}$,
$x \mapsto (p(x), x, F(x), \text{id}_{F(x)})$,
construit dans la preuve du
lemme \ref{categories-lemma-ameliorate-morphism-categories-fibred-groupoids}.
Soient $\alpha : F \to g \circ b$ et $\beta : b^{-1} \circ b \to \text{id}$
des $2$-isomorphismes entre $1$-morphismes de catégories au-dessus de $\mathcal{C}$.
Remarquons que $c \circ b$ est donné sur les objets par la règle
$$
x \mapsto (p(x), b^{-1}(b(x)), g(b(x)), \alpha_x \circ F(\beta_x))
$$
On voit alors que
$$
(\beta_x, \alpha_x) :
(p(x), x, F(x), \text{id}_{F(x)})
\longrightarrow
(p(x), b^{-1}(b(x)), g(b(x)), \alpha_x \circ F(\beta_x))
$$
est un isomorphisme fonctoriel qui donne notre $2$-morphisme
$d \to c \circ b$. Enfin, si le diagramme est commutatif, alors
$\alpha_x$ est l'identité pour tout $x$, et ce
$2$-morphisme est donc un $2$-morphisme dans la $2$-catégorie des catégories
au-dessus de $\mathcal{Y}$.
\end{proof}


































\section{Préfaisceaux de catégories}
\label{categories-section-presheaves-categories}

\noindent
Dans cette section, nous comparons la notion de catégorie fibrée
à la notion étroitement apparentée de « préfaisceau de catégories ».
La construction de base est expliquée dans l'exemple suivant.

\begin{example}
\label{categories-example-functor-categories}
Soit $\mathcal{C}$ une catégorie.
Supposons que $F : \mathcal{C}^{opp} \to \textit{Cat}$ soit un foncteur
vers la $2$-catégorie des catégories ; voir la
définition \ref{categories-definition-functor-into-2-category}.
Pour $f : V \to U$ dans $\mathcal{C}$, nous écrirons de manière suggestive
$F(f) = f^\ast$ pour le foncteur de $F(U)$ vers $F(V)$.
À partir de là, nous pouvons construire une catégorie fibrée $\mathcal{S}_F$ au-dessus de
$\mathcal{C}$ comme suit. Posons
$$
\Ob(\mathcal{S}_F) =
\{(U, x) \mid U\in \Ob(\mathcal{C}), x\in \Ob(F(U))\}.
$$
Pour $(U, x), (V, y) \in \Ob(\mathcal{S}_F)$, posons
\begin{align*}
\Mor_{\mathcal{S}_F}((V, y), (U, x)) & =
\{ (f, \phi) \mid f \in \Mor_\mathcal{C}(V, U),
\phi \in \Mor_{F(V)}(y, f^\ast x)\} \\
& =
\coprod\nolimits_{f \in \Mor_\mathcal{C}(V, U)}
\Mor_{F(V)}(y, f^\ast x)
\end{align*}
Pour définir la composition, nous utilisons l'égalité $g^\ast \circ f^\ast =
(f \circ g)^\ast$ pour une paire de morphismes composables de $\mathcal{C}$
(par définition d'un foncteur vers une $2$-catégorie).
Plus précisément, nous définissons la composée de $\psi : z \to g^\ast y$ et
$ \phi : y \to f^\ast x$ comme $ g^\ast(\phi) \circ \psi$. Le foncteur
$p_F : \mathcal{S}_F \to \mathcal{C}$ est donné par la règle
$(U, x) \mapsto U$.
Vérifions qu'il s'agit bien d'une catégorie fibrée.
Étant donnés $f: V \to U$ dans $\mathcal{C}$ et $(U, x)$, objet au-dessus de $U$, nous
affirmons que $(f, \text{id}_{f^\ast x}): (V, {f^\ast x}) \to (U, x)$ est un
relèvement fortement cartésien de $f$.
Il faut montrer qu'un $h$ dans le diagramme de gauche
détermine $(h, \nu)$ dans celui de droite :
$$
\xymatrix{
V \ar[r]^f &
U &
(V, f^*x) \ar[r]^{(f, \text{id}_{f^*x})} &
(U, x) \\
W \ar@{-->}[u]^h \ar[ru]_g & &
(W, z) \ar@{-->}[u]^{(h, \nu)} \ar[ru]_{(g, \psi)} &
}
$$
Il suffit de prendre $\nu = \psi$, ce qui convient puisque $f \circ h = g$
et donc $g^*x = h^*f^*x$. De plus, c'est l'unique relèvement
qui rende commutatif le diagramme (de droite).
\end{example}

\begin{definition}
\label{categories-definition-split-fibred-category}
Soit $\mathcal{C}$ une catégorie.
Supposons que $F : \mathcal{C}^{opp} \to \textit{Cat}$ soit un foncteur
vers la $2$-catégorie des catégories.
Nous noterons $p_F : \mathcal{S}_F \to \mathcal{C}$ la
catégorie fibrée construite dans
l'exemple \ref{categories-example-functor-categories}.
Une {\it catégorie fibrée scindée} est une catégorie fibrée isomorphe (!),
au-dessus de $\mathcal{C}$, à l'une de ces catégories {\it $\mathcal{S}_F$}.
\end{definition}

\begin{lemma}
\label{categories-lemma-when-split}
Soit $\mathcal{C}$ une catégorie.
Soit $\mathcal{S}$ une catégorie fibrée au-dessus de $\mathcal{C}$.
Alors $\mathcal{S}$ est scindée si et seulement si, pour un certain clivage
(voir la définition \ref{categories-definition-pullback-functor-fibred-category}),
les foncteurs de changement de base
$(f \circ g)^*$ et $g^* \circ f^*$ sont égaux.
\end{lemma}

\begin{proof}
Cela découle immédiatement des définitions.
\end{proof}

\begin{lemma}
\label{categories-lemma-fibred-strict}
Soit $ p : \mathcal{S} \to \mathcal{C}$ une catégorie fibrée.
Il existe un foncteur contravariant $F : \mathcal{C} \to \textit{Cat}$
tel que $\mathcal{S}$ soit équivalente à $\mathcal{S}_F$
dans la $2$-catégorie des catégories fibrées au-dessus de $\mathcal{C}$. Autrement
dit, toute catégorie fibrée est équivalente à une catégorie scindée.
\end{lemma}

\begin{proof}
Choisissons un clivage (voir la
définition \ref{categories-definition-pullback-functor-fibred-category}).
Le lemme \ref{categories-lemma-fibred} fournit des foncteurs de changement de base $f^*$ pour
tout morphisme $f$ de $\mathcal{C}$.

\medskip\noindent
Construisons une nouvelle catégorie $\mathcal{S}'$ comme suit.
Les objets de $\mathcal{S}'$ sont les paires $(x, f)$
formées d'un morphisme $f : V \to U$ de $\mathcal{C}$
et d'un objet $x$ de $\mathcal{S}$ au-dessus de $U$, c'est-à-dire
$x\in \Ob(\mathcal{S}_U)$. Le foncteur
$p' : \mathcal{S}' \to \mathcal{C}$ envoie la paire $(x, f)$ sur la source
du morphisme $f$ ; autrement dit, $p'(x, f : V\to U) = V$. Un morphisme
$\varphi : (x_1, f_1: V_1 \to U_1) \to (x_2, f_2 : V_2 \to U_2)$ est donné par une
paire $(\varphi, g)$ formée d'un morphisme $g : V_1 \to V_2$ et d'un morphisme
$\varphi : f_1^\ast x_1 \to f_2^\ast x_2$ tel que $p(\varphi) = g$. Il est
immédiat de définir la loi de composition : $(\varphi, g) \circ (\psi, h) =
(\varphi \circ \psi, g\circ h)$ pour toute paire de morphismes composables.
Il existe un foncteur naturel $\mathcal{S} \to \mathcal{S}'$ qui envoie simplement
$x$ au-dessus de $U$ sur la paire $(x, \text{id}_U)$.

\medskip\noindent
Il nous faut à présent vérifier que $p'$ fait de $\mathcal{S}'$ une
catégorie fibrée au-dessus de $\mathcal{C}$, et que
$\mathcal{S} \to \mathcal{S}'$ est une équivalence de catégories au-dessus de
$\mathcal{C}$ qui envoie les morphismes fortement cartésiens sur des morphismes fortement
cartésiens. Nous omettons ces vérifications.

\medskip\noindent
Enfin, nous pouvons définir des foncteurs de changement de base sur $\mathcal{S}'$
en posant $g^\ast(x, f) = (x, f \circ g)$ sur les objets si
$g : V' \to V$ et $f : V \to U$. Sur les morphismes
$(\varphi, \text{id}_V) : (x_1, f_1) \to (x_2, f_2)$
de $\mathcal{S}'_V$, nous posons $g^\ast(\varphi, \text{id}_V) =
(g^\ast\varphi, \text{id}_{V'})$, où nous utilisons les identifications uniques
$g^\ast f_i^\ast x_i = (f_i \circ g)^\ast x_i$ du lemme
\ref{categories-lemma-fibred} pour considérer $g^\ast\varphi$ comme un morphisme de
$(f_1 \circ g)^\ast x_1$ vers $(f_2 \circ g)^\ast x_2$. Manifestement, ces foncteurs
de changement de base $g^\ast$ vérifient
$g_1^\ast \circ g_2^\ast = (g_2\circ g_1)^\ast$ ; autrement dit, $\mathcal{S}'$
est scindée, comme voulu.
\end{proof}



















\section{Préfaisceaux de groupoïdes}
\label{categories-section-presheaves-groupoids}

\noindent
Dans cette section, nous comparons la notion de catégorie fibrée en groupoïdes
à la notion étroitement apparentée de « préfaisceau de groupoïdes ». La construction
de base est expliquée dans l'exemple suivant.

\begin{example}
\label{categories-example-functor-groupoids}
Cet exemple est l'analogue de
l'exemple \ref{categories-example-functor-categories},
pour les « préfaisceaux de groupoïdes » plutôt que pour les « préfaisceaux de catégories ».
Le résultat sera une catégorie fibrée en groupoïdes plutôt qu'une catégorie fibrée.
Supposons que $F : \mathcal{C}^{opp} \to \textit{Groupoïdes}$ soit un foncteur
vers la $2$-catégorie de groupoïdes ; voir la
définition \ref{categories-definition-functor-into-2-category}.
Pour $f : V \to U$ dans $\mathcal{C}$, nous écrirons de manière suggestive
$F(f) = f^\ast$ pour le foncteur de $F(U)$ vers $F(V)$.
Construisons une catégorie $\mathcal{S}_F$ fibrée en groupoïdes au-dessus de $\mathcal{C}$
comme suit. Posons
$$
\Ob(\mathcal{S}_F) =
\{(U, x) \mid U\in \Ob(\mathcal{C}), x\in \Ob(F(U))\}.
$$
Pour $(U, x), (V, y) \in \Ob(\mathcal{S}_F)$, posons
\begin{align*}
\Mor_{\mathcal{S}_F}((V, y), (U, x))
& =
\{ (f, \phi) \mid f \in \Mor_\mathcal{C}(V, U),
\phi \in \Mor_{F(V)}(y, f^\ast x)\} \\
& =
\coprod\nolimits_{f \in \Mor_\mathcal{C}(V, U)}
\Mor_{F(V)}(y, f^\ast x)
\end{align*}
Pour définir la composition, nous utilisons l'égalité $g^\ast \circ f^\ast =
(f \circ g)^\ast$ pour une paire de morphismes composables de $\mathcal{C}$
(par définition d'un foncteur vers une $2$-catégorie).
Plus précisément, nous définissons la composée de $\psi : z \to g^\ast y$ et
$ \phi : y \to f^\ast x$ comme $ g^\ast(\phi) \circ \psi$. Le foncteur
$p_F : \mathcal{S}_F \to \mathcal{C}$ est donné par la règle $(U, x) \mapsto U$.
La condition selon laquelle $F(U)$ est un groupoïde pour tout $U$ garantit que
$\mathcal{S}_F$ est fibrée en groupoïdes au-dessus de $\mathcal{C}$, puisque nous avons
déjà vu dans
l'exemple \ref{categories-example-functor-categories}
que $\mathcal{S}_F$ est une catégorie fibrée ; voir le
lemme \ref{categories-lemma-fibred-groupoids}.
Mais nous pouvons aussi démontrer directement les conditions (1), (2) de la
définition \ref{categories-definition-fibred-groupoids}
comme suit : (1) Les relèvements de
morphismes existent : étant donnés $f: V \to U$ dans $\mathcal{C}$ et $(U, x)$,
objet de $\mathcal{S}_F$ au-dessus de $U$,
$(f, \text{id}_{f^\ast x}): (V, {f^\ast x}) \to (U, x)$ est un relèvement de $f$.
(2) Supposons données les flèches pleines des diagrammes suivants :
$$
\xymatrix{
V \ar[r]^f & U & (V, y) \ar[r]^{(f, \phi)} & (U, x) \\
W \ar@{-->}[u]^h \ar[ru]_g & &
(W, z) \ar@{-->}[u]^{(h, \nu)} \ar[ru]_{(g, \psi)} & \\
}
$$
Alors, pour les flèches pointillées, on a $\nu = (h^\ast \phi)^{-1} \circ \psi$ ;
ainsi, $h$ étant donné, il existe un $\nu$, unique par unicité des inverses.
\end{example}

\begin{definition}
\label{categories-definition-split-category-fibred-in-groupoids}
Soit $\mathcal{C}$ une catégorie.
Supposons que $F : \mathcal{C}^{opp} \to \textit{Groupoïdes}$ soit un foncteur
vers la $2$-catégorie des groupoïdes.
Nous noterons $p_F : \mathcal{S}_F \to \mathcal{C}$ la
catégorie fibrée en groupoïdes construite dans
l'exemple \ref{categories-example-functor-groupoids}.
Une {\it catégorie fibrée en groupoïdes scindée} est une
catégorie fibrée en groupoïdes isomorphe (!),
au-dessus de $\mathcal{C}$, à l'une de ces catégories {\it $\mathcal{S}_F$}.
\end{definition}

\begin{lemma}
\label{categories-lemma-fibred-groupoids-strict}
Soit $ p : \mathcal{S} \to \mathcal{C}$ une catégorie fibrée en groupoïdes.
Il existe un foncteur contravariant $F : \mathcal{C} \to \textit{Groupoïdes}$
tel que $\mathcal{S}$ soit équivalente à $\mathcal{S}_F$ au-dessus de $\mathcal{C}$.
Autrement dit, toute catégorie fibrée en groupoïdes est équivalente à une catégorie scindée.
\end{lemma}

\begin{proof}
Choisissons un clivage (voir la
définition \ref{categories-definition-pullback-functor-fibred-category}).
Les lemmes \ref{categories-lemma-fibred} et \ref{categories-lemma-fibred-groupoids}
fournissent des foncteurs de changement de base $f^*$ pour
tout morphisme $f$ de $\mathcal{C}$.

\medskip\noindent
Construisons une nouvelle catégorie $\mathcal{S}'$ comme suit.
Les objets de $\mathcal{S}'$ sont les paires $(x, f)$
formées d'un morphisme $f : V \to U$ de $\mathcal{C}$
et d'un objet $x$ de $\mathcal{S}$ au-dessus de $U$, c'est-à-dire
$x\in \Ob(\mathcal{S}_U)$. Le foncteur
$p' : \mathcal{S}' \to \mathcal{C}$ envoie la paire $(x, f)$ sur la source
du morphisme $f$ ; autrement dit, $p'(x, f : V\to U) = V$. Un morphisme
$\varphi : (x_1, f_1: V_1 \to U_1) \to (x_2, f_2 : V_2 \to U_2)$ est donné par une
paire $(\varphi, g)$ formée d'un morphisme $g : V_1 \to V_2$ et d'un morphisme
$\varphi : f_1^\ast x_1 \to f_2^\ast x_2$ tel que $p(\varphi) = g$. Il est
immédiat de définir la loi de composition : $(\varphi, g) \circ (\psi, h) =
(\varphi \circ \psi, g\circ h)$ pour toute paire de morphismes composables.
Il existe un foncteur naturel $\mathcal{S} \to \mathcal{S}'$ qui envoie simplement
$x$ au-dessus de $U$ sur la paire $(x, \text{id}_U)$.

\medskip\noindent
Il nous faut à présent vérifier que $p'$ fait de $\mathcal{S}'$ une catégorie
fibrée en groupoïdes au-dessus de $\mathcal{C}$, et que
$\mathcal{S} \to \mathcal{S}'$ est une équivalence de catégories au-dessus de
$\mathcal{C}$. Nous omettons ces vérifications.

\medskip\noindent
Enfin, nous pouvons définir des foncteurs de changement de base sur $\mathcal{S}'$
en posant $g^\ast(x, f) = (x, f \circ g)$ sur les objets si
$g : V' \to V$ et $f : V \to U$. Sur les morphismes
$(\varphi, \text{id}_V) : (x_1, f_1) \to (x_2, f_2)$
de $\mathcal{S}'_V$, nous posons $g^\ast(\varphi, \text{id}_V) =
(g^\ast\varphi, \text{id}_{V'})$, où nous utilisons les identifications uniques
$g^\ast f_i^\ast x_i = (f_i \circ g)^\ast x_i$ du lemme
\ref{categories-lemma-fibred-groupoids} pour considérer $g^\ast\varphi$ comme un morphisme de
$(f_1 \circ g)^\ast x_1$ vers $(f_2 \circ g)^\ast x_2$. Manifestement, ces foncteurs
de changement de base $g^\ast$ vérifient
$g_1^\ast \circ g_2^\ast = (g_2\circ g_1)^\ast$ ; autrement dit, $\mathcal{S}'$
est scindée, comme voulu.
\end{proof}

\noindent
Nous verrons une autre preuve de ce lemme dans la
section \ref{categories-section-representable-1-morphisms}.





\section{Catégories fibrées en ensembles}
\label{categories-section-fibred-in-sets}

\begin{definition}
\label{categories-definition-discrete}
Une catégorie est dite {\it discrète} si ses seuls morphismes sont les
identités.
\end{definition}

\noindent
Une catégorie discrète ne possède qu'une seule donnée intéressante :
son ensemble d'objets. Ainsi, nous confondons parfois les catégories discrètes
avec les ensembles.

\begin{definition}
\label{categories-definition-category-fibred-sets}
Soit $\mathcal{C}$ une catégorie.
Une {\it catégorie fibrée en ensembles}, ou une {\it catégorie fibrée
en catégories discrètes}, est une catégorie fibrée en groupoïdes dont
toutes les catégories fibres sont discrètes.
\end{definition}

\noindent
Nous voulons préciser la relation entre les catégories fibrées en ensembles
et les préfaisceaux (voir la définition \ref{categories-definition-presheaf}).
Pour ce faire, il est naturel de commencer par la définition suivante.

\begin{definition}
\label{categories-definition-categories-fibred-in-sets-over-C}
Soit $\mathcal{C}$ une catégorie.
La {\it $2$-catégorie des catégories fibrées en ensembles au-dessus de
$\mathcal{C}$} est la sous-$2$-catégorie de la $2$-catégorie des catégories
fibrées en groupoïdes au-dessus de $\mathcal{C}$ (voir la
définition \ref{categories-definition-categories-fibred-in-groupoids-over-C})
définie comme suit :
\begin{enumerate}
\item Ses objets sont les catégories
$p : \mathcal{S} \to \mathcal{C}$ fibrées en ensembles.
\item Ses $1$-morphismes $(\mathcal{S}, p) \to (\mathcal{S}', p')$
sont les foncteurs $G : \mathcal{S} \to \mathcal{S}'$ tels que
$p' \circ G = p$ (puisque tout morphisme est fortement cartésien,
$G$ les préserve automatiquement).
\item Ses $2$-morphismes $t : G \to H$ pour
$G, H : (\mathcal{S}, p) \to (\mathcal{S}', p')$
sont les morphismes de foncteurs
tels que $p'(t_x) = \text{id}_{p(x)}$
pour tout $x \in \Ob(\mathcal{S})$.
\end{enumerate}
\end{definition}

\noindent
Notons que tout $2$-morphisme est automatiquement un isomorphisme.
Cette $2$-catégorie est donc en fait une $(2, 1)$-catégorie.
Voici le lemme obligé sur l'existence des $2$-produits fibrés.

\begin{lemma}
\label{categories-lemma-2-product-categories-fibred-sets}
Soit $\mathcal{C}$ une catégorie.
La 2-catégorie des catégories fibrées en ensembles au-dessus de $\mathcal{C}$
admet des 2-produits fibrés. Plus précisément, le 2-produit fibré décrit dans le
lemme \ref{categories-lemma-2-product-categories-over-C}
est une catégorie fibrée en ensembles lorsque les données de départ le sont.
\end{lemma}

\begin{proof}
Omis.
\end{proof}

\begin{example}
\label{categories-example-presheaf}
Cet exemple est l'analogue des
exemples \ref{categories-example-functor-categories} et
\ref{categories-example-functor-groupoids}
pour les préfaisceaux au lieu des « préfaisceaux de catégories ».
Le résultat sera une catégorie fibrée en ensembles plutôt qu'une catégorie fibrée.
Supposons que $F : \mathcal{C}^{opp} \to \textit{Ensembles}$ soit un préfaisceau.
Pour $f : V \to U$ dans $\mathcal{C}$, nous écrirons de manière suggestive
$F(f) = f^\ast : F(U) \to F(V)$.
Nous construisons une catégorie $\mathcal{S}_F$ fibrée en ensembles au-dessus de
$\mathcal{C}$ comme suit. Posons
$$
\Ob(\mathcal{S}_F) =
\{(U, x) \mid U \in \Ob(\mathcal{C}), x \in F(U)\}.
$$
Pour $(U, x), (V, y) \in \Ob(\mathcal{S}_F)$, nous posons
\begin{align*}
\Mor_{\mathcal{S}_F}((V, y), (U, x))
& =
\{f \in \Mor_\mathcal{C}(V, U) \mid f^*x = y\}
\end{align*}
La composition est héritée de celle de $\mathcal{C}$ ; cela fonctionne car
$g^\ast \circ f^\ast = (f \circ g)^\ast$
pour une paire de morphismes composables de $\mathcal{C}$.
Le foncteur $p_F : \mathcal{S}_F \to \mathcal{C}$
est donné par la règle $(U, x) \mapsto U$.
Puisque toute catégorie fibre $\mathcal{S}_{F, U}$ est discrète, d'ensemble
sous-jacent $F(U)$, et que nous avons déjà vu dans
l'exemple \ref{categories-example-functor-groupoids}
que $\mathcal{S}_F$ est une catégorie fibrée en groupoïdes,
nous concluons que $\mathcal{S}_F$ est fibrée en ensembles.
\end{example}

\begin{lemma}
\label{categories-lemma-2-category-fibred-sets}
\begin{slogan}
Les catégories fibrées en ensembles sont exactement les préfaisceaux.
\end{slogan}
Soit $\mathcal{C}$ une catégorie.
Les seuls $2$-morphismes entre catégories fibrées en ensembles sont les identités.
Autrement dit, la $2$-catégorie des catégories fibrées en ensembles est une catégorie.
De plus, il existe une équivalence de catégories
$$
\left\{
\begin{matrix}
\text{la catégorie des préfaisceaux}\\
\text{d'ensembles sur }\mathcal{C}
\end{matrix}
\right\}
\leftrightarrow
\left\{
\begin{matrix}
\text{la catégorie des catégories}\\
\text{fibrées en ensembles au-dessus de }\mathcal{C}
\end{matrix}
\right\}
$$
Le foncteur de gauche à droite est la construction
$F \to \mathcal{S}_F$ étudiée dans
l'exemple \ref{categories-example-presheaf}.
Le foncteur de droite à gauche associe à $p : \mathcal{S} \to \mathcal{C}$
le préfaisceau d'objets $U \mapsto \Ob(\mathcal{S}_U)$.
\end{lemma}

\begin{proof}
La première assertion est claire, puisque les seuls morphismes des catégories
fibres sont les identités.

\medskip\noindent
Supposons que $p :
\mathcal{S} \to \mathcal{C}$ soit fibrée en ensembles. Soit $f : V \to U$
un morphisme de $\mathcal{C}$ et soit $x \in \Ob(\mathcal{S}_U)$.
Il existe alors exactement un choix pour l'objet $f^\ast x$. Nous voyons donc que
$(f \circ g)^\ast x = g^\ast(f^\ast x)$ pour $f, g$ comme dans le lemme
\ref{categories-lemma-fibred-groupoids}. Il s'ensuit que l'on peut considérer les
applications $U \mapsto \Ob(\mathcal{S}_U)$ et $f \mapsto f^\ast$
comme un préfaisceau sur $\mathcal{C}$.
\end{proof}

\noindent
Voici un exemple important de catégorie fibrée en ensembles.

\begin{example}
\label{categories-example-fibred-category-from-functor-of-points}
Soit $\mathcal{C}$ une catégorie. Soit $X \in \Ob(\mathcal{C})$.
Considérons le préfaisceau représentable $h_X = \Mor_\mathcal{C}(-, X)$
(voir l'exemple \ref{categories-example-hom-functor}).
D'autre part, considérons la catégorie $p : \mathcal{C}/X \to \mathcal{C}$
de l'exemple \ref{categories-example-category-over-X}.
La catégorie fibre $(\mathcal{C}/X)_U$ a pour objets les morphismes
$h : U \to X$, et pour seuls morphismes les identités. Nous voyons donc que,
sous la correspondance du
lemme \ref{categories-lemma-2-category-fibred-sets},
nous avons
$$
h_X \longleftrightarrow \mathcal{C}/X.
$$
Autrement dit, la catégorie $\mathcal{C}/X$ est canoniquement équivalente
à la catégorie $\mathcal{S}_{h_X}$ associée
à $h_X$ dans
l'exemple \ref{categories-example-presheaf}.
\end{example}

\noindent
Pour cette raison, il est tentant de définir un objet « représentable » de la
2-catégorie des catégories fibrées en groupoïdes comme une catégorie fibrée en
ensembles dont le préfaisceau associé est représentable. Ce ne serait toutefois
pas une bonne définition, car nous préférons disposer d'une notion qui soit
invariante par équivalences. Pour préciser ce point, nous étudions exactement
quelles catégories fibrées en groupoïdes sont équivalentes à des catégories
fibrées en ensembles.








\section{Catégories fibrées en setoïdes}
\label{categories-section-fibred-in-setoids}

\begin{definition}
\label{categories-definition-setoid}
Appelons {\it setoïde}\footnote{Un ensemble sous stéroïdes !?} toute catégorie
qui est un groupoïde dans lequel tout objet
possède exactement un automorphisme : l'identité.
\end{definition}

\noindent
Si $C$ est un ensemble muni d'une relation d'équivalence $\sim$, on peut en faire
un setoïde $\mathcal{C}$ comme suit : $\Ob(\mathcal{C}) = C$ et
$\Mor_\mathcal{C}(x, y) = \emptyset$ sauf si $x \sim y$, auquel
cas on pose $\Mor_\mathcal{C}(x, y) = \{1\}$. La transitivité de
$\sim$ signifie que l'on peut composer les morphismes. Réciproquement, tout
setoïde définit sur ses objets la relation d'équivalence « être isomorphes »,
à partir de laquelle on retrouve la catégorie (à isomorphisme
unique près, et non simplement à équivalence près) par le procédé qui vient d'être décrit.

\medskip\noindent
Les catégories discrètes sont des setoïdes. Pour tout setoïde $\mathcal{C}$,
il existe un procédé canonique qui produit une catégorie
discrète équivalente : on remplace $\Ob(\mathcal{C})$ par l'ensemble des classes
d'isomorphisme (et l'on ajoute les identités). En termes d'ensembles
munis d'une relation d'équivalence, cela correspond au passage au quotient
par la relation d'équivalence.

\begin{definition}
\label{categories-definition-category-fibred-setoids}
Soit $\mathcal{C}$ une catégorie. Une {\it catégorie fibrée en setoïdes} est
une catégorie fibrée en groupoïdes dont toutes les catégories fibres sont des
setoïdes.
\end{definition}

\noindent
Nous préciserons ci-dessous la relation entre les catégories fibrées en
setoïdes et les catégories fibrées en ensembles.

\begin{definition}
\label{categories-definition-categories-fibred-in-setoids-over-C}
Soit $\mathcal{C}$ une catégorie.
La {\it $2$-catégorie des catégories fibrées en setoïdes
au-dessus de $\mathcal{C}$} est la sous-$2$-catégorie de la $2$-catégorie des
catégories fibrées en groupoïdes au-dessus de $\mathcal{C}$ (voir la
définition \ref{categories-definition-categories-fibred-in-groupoids-over-C})
définie comme suit :
\begin{enumerate}
\item Ses objets sont les catégories
$p : \mathcal{S} \to \mathcal{C}$ fibrées en setoïdes.
\item Ses $1$-morphismes $(\mathcal{S}, p) \to (\mathcal{S}', p')$
sont les foncteurs $G : \mathcal{S} \to \mathcal{S}'$ tels que
$p' \circ G = p$ (puisque tout morphisme est fortement cartésien,
$G$ les préserve automatiquement).
\item Ses $2$-morphismes $t : G \to H$ pour
$G, H : (\mathcal{S}, p) \to (\mathcal{S}', p')$
sont les morphismes de foncteurs
tels que $p'(t_x) = \text{id}_{p(x)}$
pour tout $x \in \Ob(\mathcal{S})$.
\end{enumerate}
\end{definition}

\noindent
Notons que tout $2$-morphisme est automatiquement un isomorphisme.
Cette $2$-catégorie est donc en fait une $(2, 1)$-catégorie.

\noindent
Voici le lemme obligé sur l'existence des $2$-produits fibrés.

\begin{lemma}
\label{categories-lemma-2-product-categories-fibred-setoids}
Soit $\mathcal{C}$ une catégorie.
La 2-catégorie des catégories fibrées en setoïdes au-dessus de
$\mathcal{C}$ admet des 2-produits fibrés. Plus précisément, le 2-produit fibré
décrit dans le lemme \ref{categories-lemma-2-product-categories-over-C} est une catégorie
fibrée en setoïdes lorsque les données de départ le sont.
\end{lemma}

\begin{proof}
Omis.
\end{proof}

\begin{lemma}
\label{categories-lemma-setoid-fibres}
Soit $\mathcal{C}$ une catégorie. Soit $\mathcal{S}$ une catégorie
au-dessus de $\mathcal{C}$.
\begin{enumerate}
\item Si $\mathcal{S} \to \mathcal{S}'$ est une équivalence
au-dessus de $\mathcal{C}$, avec $\mathcal{S}'$ fibrée en ensembles au-dessus
de $\mathcal{C}$, alors
\begin{enumerate}
\item $\mathcal{S}$ est fibrée en setoïdes au-dessus de
$\mathcal{C}$, et
\item pour tout $U \in \Ob(\mathcal{C})$, l'application
$\Ob(\mathcal{S}_U) \to \Ob(\mathcal{S}'_U)$
identifie le but à l'ensemble des classes d'isomorphisme de la source.
\end{enumerate}
\item Si $p : \mathcal{S} \to \mathcal{C}$ est une catégorie fibrée en
setoïdes, alors il existe une catégorie fibrée en ensembles
$p' : \mathcal{S}' \to \mathcal{C}$ et une équivalence
$\text{can} : \mathcal{S} \to \mathcal{S}'$ au-dessus de $\mathcal{C}$.
\end{enumerate}
\end{lemma}

\begin{proof}
Démontrons (2).
Un objet de la catégorie $\mathcal{S}'$ est une paire $(U, \xi)$, où
$U \in \Ob(\mathcal{C})$ et $\xi$ est une classe d'isomorphisme d'objets
de $\mathcal{S}_U$. Un morphisme $(U, \xi) \to (V , \psi)$ est donné par un
morphisme $x \to y$, où $x \in \xi$ et $y \in \psi$. Ici, nous identifions
deux morphismes $x \to y$ et $x' \to y'$ s'ils induisent le même morphisme
$U \to V$, et si, pour certains choix d'isomorphismes $x \to x'$ dans
$\mathcal{S}_U$ et $y \to y'$ dans $\mathcal{S}_V$, les composées
$x \to x' \to y'$ et $x \to y \to y'$ coïncident. Par construction, les
applications sur les objets et les morphismes induites par $\mathcal{S} \to
\mathcal{S}'$ sont surjectives. Nous définissons la composition des morphismes dans
$\mathcal{S}'$ comme l'unique loi qui fasse de $\mathcal{S} \to \mathcal{S}'$
un foncteur. Certains détails sont omis.
\end{proof}

\noindent
Ainsi, les catégories fibrées en setoïdes sont exactement les
catégories fibrées en groupoïdes qui sont équivalentes à des catégories fibrées
en ensembles. De plus, une équivalence de catégories fibrées en ensembles est
un isomorphisme d'après le lemme \ref{categories-lemma-2-category-fibred-sets}.

\begin{lemma}
\label{categories-lemma-2-category-fibred-setoids}
Soit $\mathcal{C}$ une catégorie. La construction du
lemme \ref{categories-lemma-setoid-fibres},
partie (2), donne un foncteur
$$
F :
\left\{
\begin{matrix}
\text{la 2-catégorie des catégories}\\
\text{fibrées en setoïdes au-dessus de }\mathcal{C}
\end{matrix}
\right\}
\longrightarrow
\left\{
\begin{matrix}
\text{la catégorie des catégories}\\
\text{fibrées en ensembles au-dessus de }\mathcal{C}
\end{matrix}
\right\}
$$
(voir la
définition \ref{categories-definition-functor-into-2-category}).
Ce foncteur est une équivalence au sens suivant :
\begin{enumerate}
\item pour deux 1-morphismes quelconques
$f, g : \mathcal{S}_1 \to \mathcal{S}_2$
tels que $F(f) = F(g)$, il existe un unique 2-isomorphisme $f \to g$,
\item pour tout morphisme $h : F(\mathcal{S}_1) \to F(\mathcal{S}_2)$,
il existe un 1-morphisme $f : \mathcal{S}_1 \to \mathcal{S}_2$
tel que $F(f) = h$, et
\item toute catégorie fibrée en ensembles $\mathcal{S}$ est égale à
$F(\mathcal{S})$.
\end{enumerate}
En particulier, si l'on définit $F_i \in \textit{PSh}(\mathcal{C})$ par la
règle $F_i(U) = \Ob(\mathcal{S}_{i, U})/\cong$, on a
$$
\Mor_{\textit{Cat}/\mathcal{C}}(\mathcal{S}_1, \mathcal{S}_2)
\Big/
2\text{-isomorphisme}
=
\Mor_{\textit{PSh}(\mathcal{C})}(F_1, F_2)
$$
Plus précisément, pour tout morphisme $\phi : F_1 \to F_2$, il existe un
$1$-morphisme $f : \mathcal{S}_1 \to \mathcal{S}_2$ qui induit
$\phi$ sur les classes d'isomorphisme d'objets et
qui est unique à un unique $2$-isomorphisme près.
\end{lemma}

\begin{proof}
D'après le
lemme \ref{categories-lemma-2-category-fibred-sets},
le but de $F$ est une catégorie ; l'assertion a donc un sens.
La construction du
lemme \ref{categories-lemma-setoid-fibres}, partie (2),
associe à $\mathcal{S}$ la catégorie fibrée en ensembles dont la valeur au-dessus
de $U$ est l'ensemble des classes d'isomorphisme dans $\mathcal{S}_U$. Il est
donc clair qu'elle définit un foncteur comme indiqué.
Soient $f, g : \mathcal{S}_1 \to \mathcal{S}_2$
tels que $F(f) = F(g)$ comme en (1). Pour tout objet $U$ de $\mathcal{C}$
et tout objet $x$ de $\mathcal{S}_{1, U}$, l'hypothèse donne
$f(x) \cong g(x)$. Comme $\mathcal{S}_2$ est fibrée en setoïdes,
il existe un unique isomorphisme $t_x : f(x) \to g(x)$ dans
$\mathcal{S}_{2, U}$.
Il est clair que la règle $x \mapsto t_x$ donne le $2$-isomorphisme souhaité
$f \to g$. Nous omettons les démonstrations de (2) et (3).
Pour voir la dernière assertion, appliquons le
lemme \ref{categories-lemma-2-category-fibred-sets}
pour identifier le membre de droite à
$\Mor_{\textit{Cat}/\mathcal{C}}(F(\mathcal{S}_1), F(\mathcal{S}_2))$
et appliquons (1) et (2) ci-dessus.
\end{proof}

\noindent
Voici une autre caractérisation des catégories fibrées en setoïdes parmi
toutes les catégories fibrées en groupoïdes.

\begin{lemma}
\label{categories-lemma-characterize-fibred-setoids-inertia}
Soit $\mathcal{C}$ une catégorie.
Soit $p : \mathcal{S} \to \mathcal{C}$ une catégorie fibrée en groupoïdes.
Les conditions suivantes sont équivalentes :
\begin{enumerate}
\item $p : \mathcal{S} \to \mathcal{C}$ est une catégorie fibrée en setoïdes,
et
\item le $1$-morphisme canonique $\mathcal{I}_\mathcal{S} \to \mathcal{S}$,
voir (\ref{categories-equation-inertia-structure-map}), est une équivalence (de catégories
au-dessus de $\mathcal{C}$).
\end{enumerate}
\end{lemma}

\begin{proof}
Supposons (2). La catégorie $\mathcal{I}_\mathcal{S}$ a pour objets
les $(x, \alpha)$ où $x \in \mathcal{S}$, avec, disons, $p(x) = U$, et où
$\alpha : x \to x$ est un morphisme de $\mathcal{S}_U$. Ainsi, si
$\mathcal{I}_\mathcal{S} \to \mathcal{S}$ est une équivalence au-dessus de
$\mathcal{C}$, alors les deux objets de toute paire
$(x, \alpha)$, $(x, \alpha')$ sont isomorphes dans la catégorie fibre de
$\mathcal{I}_\mathcal{S}$ au-dessus de $U$.
En examinant la définition des morphismes de $\mathcal{I}_\mathcal{S}$,
nous concluons que $\alpha$, $\alpha'$ sont conjugués dans le groupe
des automorphismes de $x$. En prenant $\alpha' = \text{id}_x$, nous concluons
donc que tout automorphisme de $x$ est égal à l'identité.
Puisque $\mathcal{S} \to \mathcal{C}$ est fibrée en groupoïdes, cela implique
que $\mathcal{S} \to \mathcal{C}$ est fibrée en setoïdes.
Nous omettons la démonstration de (1) $\Rightarrow$ (2).
\end{proof}

\begin{lemma}
\label{categories-lemma-category-fibred-setoids-presheaves-products}
Soit $\mathcal{C}$ une catégorie.
La construction du
lemme \ref{categories-lemma-2-category-fibred-setoids},
qui associe un préfaisceau à une catégorie fibrée en setoïdes,
est compatible avec les produits, au sens où le préfaisceau associé à un
$2$-produit fibré $\mathcal{X} \times_\mathcal{Y} \mathcal{Z}$
est le produit fibré des préfaisceaux associés à
$\mathcal{X}, \mathcal{Y}, \mathcal{Z}$.
\end{lemma}

\begin{proof}
Soit $U \in \Ob(\mathcal{C})$. Le lemme affirme simplement que
$$
\Ob((\mathcal{X} \times_\mathcal{Y} \mathcal{Z})_U)/\!\cong
\quad \text{est égal à} \quad
\Ob(\mathcal{X}_U)/\!\cong
\ \times_{\Ob(\mathcal{Y}_U)/\!\cong}
\ \Ob(\mathcal{Z}_U)/\!\cong
$$
et nous en omettons la démonstration. (Notons cependant que ce ne serait pas
vrai en général si la catégorie $\mathcal{Y}_U$ n'était pas un setoïde.)
\end{proof}










\section{Catégories fibrées en groupoïdes représentables}
\label{categories-section-representable-fibred-groupoids}

\noindent
Voici notre définition d'une catégorie fibrée en groupoïdes représentable.
Comme promis, elle est invariante par équivalences.

\begin{definition}
\label{categories-definition-representable-fibred-category}
Soit $\mathcal{C}$ une catégorie.
Une catégorie fibrée en groupoïdes $p : \mathcal{S} \to \mathcal{C}$ est dite
{\it représentable} s'il existe un objet
$X$ de $\mathcal{C}$ et une équivalence $j : \mathcal{S} \to \mathcal{C}/X$
(dans la $2$-catégorie des catégories fibrées en groupoïdes au-dessus de
$\mathcal{C}$).
\end{definition}

\noindent
L'abus de notation habituel consiste à dire que
{\it $X$ représente $\mathcal{S}$} sans mentionner l'équivalence $j$.
Explicitons ce que cela implique.

\begin{lemma}
\label{categories-lemma-characterize-representable-fibred-category}
Soit $\mathcal{C}$ une catégorie.
Soit $p : \mathcal{S} \to \mathcal{C}$ une catégorie fibrée en groupoïdes.
\begin{enumerate}
\item $\mathcal{S}$ est représentable si et seulement si
les conditions suivantes sont satisfaites :
\begin{enumerate}
\item $\mathcal{S}$ est fibrée en setoïdes, et
\item le préfaisceau $U \mapsto \Ob(\mathcal{S}_U)/\cong$ est
représentable.
\end{enumerate}
\item Si $\mathcal{S}$ est représentable, la paire $(X, j)$, où $j$ est
l'équivalence $j : \mathcal{S} \to \mathcal{C}/X$, est déterminée de manière
unique à isomorphisme près.
\end{enumerate}
\end{lemma}

\begin{proof}
La première assertion découle immédiatement du
lemme \ref{categories-lemma-setoid-fibres}.
Pour la seconde, supposons que $(X', j')$ soit
une seconde paire de ce type. Choisissons un $1$-morphisme
$t' : \mathcal{C}/X' \to \mathcal{S}$ tel que
$j' \circ t' \cong \text{id}_{\mathcal{C}/X'}$ et
$t' \circ j' \cong \text{id}_\mathcal{S}$. Alors
$j \circ t' : \mathcal{C}/X' \to \mathcal{C}/X$ est une équivalence.
C'est donc un isomorphisme, voir le lemme \ref{categories-lemma-2-category-fibred-sets}.
Par le lemme de Yoneda \ref{categories-lemma-yoneda} (via, par exemple,
l'exemple \ref{categories-example-fibred-category-from-functor-of-points}),
il est donc donné par un isomorphisme $X' \to X$.
\end{proof}

\begin{lemma}
\label{categories-lemma-morphisms-representable-fibred-categories}
Soit $\mathcal{C}$ une catégorie.
Soient $\mathcal{X}$, $\mathcal{Y}$ des catégories fibrées en groupoïdes
au-dessus de $\mathcal{C}$. Supposons que $\mathcal{X}$, $\mathcal{Y}$
soient représentées par des objets $X$, $Y$ de $\mathcal{C}$.
Alors
$$
\Mor_{\textit{Cat}/\mathcal{C}}(\mathcal{X}, \mathcal{Y})
\Big/
2\text{-isomorphisme}
=
\Mor_\mathcal{C}(X, Y)
$$
Plus précisément, pour tout $\phi : X \to Y$, il existe un
$1$-morphisme $f : \mathcal{X} \to \mathcal{Y}$ qui induit
$\phi$ sur les classes d'isomorphisme d'objets et
qui est unique à un unique $2$-isomorphisme près.
\end{lemma}

\begin{proof}
D'après
l'exemple \ref{categories-example-fibred-category-from-functor-of-points},
nous avons $\mathcal{C}/X = \mathcal{S}_{h_X}$ et
$\mathcal{C}/Y = \mathcal{S}_{h_Y}$. D'après le
lemme \ref{categories-lemma-2-category-fibred-setoids},
nous avons
$$
\Mor_{\textit{Cat}/\mathcal{C}}(\mathcal{X}, \mathcal{Y})
\Big/
2\text{-isomorphisme}
=
\Mor_{\textit{PSh}(\mathcal{C})}(h_X, h_Y)
$$
Par le lemme de Yoneda
\ref{categories-lemma-yoneda},
nous avons $\Mor_{\textit{PSh}(\mathcal{C})}(h_X, h_Y)
= \Mor_\mathcal{C}(X, Y)$.
\end{proof}








\section{Le lemme de 2-Yoneda}
\label{categories-section-2-yoneda}

\noindent
Soit $\mathcal{C}$ une catégorie. La $2$-catégorie des catégories fibrées
au-dessus de $\mathcal{C}$ a été définie dans la
définition \ref{categories-definition-fibred-categories-over-C}.
Si $\mathcal{S}$, $\mathcal{S}'$ sont des catégories fibrées
au-dessus de $\mathcal{C}$, alors
$$
\Mor_{\textit{Fib}/\mathcal{C}}(\mathcal{S}, \mathcal{S}')
$$
désigne la catégorie des $1$-morphismes dans cette $2$-catégorie.
Voici l'analogue $2$-catégorique du lemme de Yoneda
dans le cadre des catégories fibrées.

\begin{lemma}[Lemme de 2-Yoneda pour les catégories fibrées]
\label{categories-lemma-yoneda-2category-fibred}
Soit $\mathcal{C}$ une catégorie.
Soit $\mathcal{S} \to \mathcal{C}$ une catégorie fibrée au-dessus de
$\mathcal{C}$.
Soit $U \in \Ob(\mathcal{C})$.
Le foncteur
$$
\Mor_{\textit{Fib}/\mathcal{C}}(\mathcal{C}/U, \mathcal{S})
\longrightarrow
\mathcal{S}_U
$$
donné par $G \mapsto G(\text{id}_U)$ est une équivalence.
\end{lemma}

\begin{proof}
Choisissons un clivage de $\mathcal{S}$
(voir la définition \ref{categories-definition-pullback-functor-fibred-category}).
Nous définissons un foncteur
$$
\mathcal{S}_U
\longrightarrow
\Mor_{\textit{Fib}/\mathcal{C}}(\mathcal{C}/U, \mathcal{S})
$$
comme suit. Étant donné
$x \in \Ob(\mathcal{S}_U)$,
le foncteur associé est défini
\begin{enumerate}
\item sur les objets par $(f : V \to U) \mapsto f^*x$, et
\item sur les morphismes en envoyant la flèche
$(g : V'/U \to V/U)$ sur la composée
$$
(f \circ g)^*x \xrightarrow{(\alpha_{g, f})_x} g^*f^*x \rightarrow f^*x
$$
où $\alpha_{g, f}$ est comme dans le lemme \ref{categories-lemma-fibred}.
\end{enumerate}
Nous omettons la vérification que ce foncteur est un inverse de celui du lemme.
\end{proof}

\noindent
Soit $\mathcal{C}$ une catégorie. La $2$-catégorie des catégories
fibrées en groupoïdes au-dessus de $\mathcal{C}$ est une sous-$2$-catégorie
« pleine » de la $2$-catégorie des catégories au-dessus de
$\mathcal{C}$ (voir la
définition \ref{categories-definition-categories-fibred-in-groupoids-over-C}).
Ainsi, si $\mathcal{S}$, $\mathcal{S}'$ sont fibrées en groupoïdes
au-dessus de $\mathcal{C}$, alors
$$
\Mor_{\textit{Cat}/\mathcal{C}}(\mathcal{S}, \mathcal{S}')
$$
désigne la catégorie des $1$-morphismes dans cette $2$-catégorie
(voir la définition \ref{categories-definition-categories-over-C}).
Ce sont tous des groupoïdes, voir les remarques qui suivent la
définition \ref{categories-definition-categories-fibred-in-groupoids-over-C}.
Voici l'analogue $2$-catégorique du lemme de Yoneda.

\begin{lemma}[Lemme de 2-Yoneda]
\label{categories-lemma-yoneda-2category}
Soit $\mathcal{S}\to \mathcal{C}$ fibrée en groupoïdes.
Soit $U \in \Ob(\mathcal{C})$.
Le foncteur
$$
\Mor_{\textit{Cat}/\mathcal{C}}(\mathcal{C}/U, \mathcal{S})
\longrightarrow
\mathcal{S}_U
$$
donné par $G \mapsto G(\text{id}_U)$ est une équivalence.
\end{lemma}

\begin{proof}
Choisissons un clivage de $\mathcal{S}$
(voir la définition \ref{categories-definition-pullback-functor-fibred-category}).
Nous définissons un foncteur
$$
\mathcal{S}_U
\longrightarrow
\Mor_{\textit{Cat}/\mathcal{C}}(\mathcal{C}/U, \mathcal{S})
$$
comme suit. Étant donné
$x \in \Ob(\mathcal{S}_U)$,
le foncteur associé est défini
\begin{enumerate}
\item sur les objets par $(f : V \to U) \mapsto f^*x$, et
\item sur les morphismes en envoyant la flèche
$(g : V'/U \to V/U)$ sur la composée
$$
(f \circ g)^*x \xrightarrow{(\alpha_{g, f})_x} g^*f^*x \rightarrow f^*x
$$
où $\alpha_{g, f}$ est comme dans le lemme \ref{categories-lemma-fibred-groupoids}.
\end{enumerate}
Nous omettons la vérification que ce foncteur est un inverse de celui du lemme.
\end{proof}

\begin{remark}
\label{categories-remark-alternative-fibred-groupoids-strict}
On peut utiliser le lemme de $2$-Yoneda pour donner une autre démonstration du
lemme \ref{categories-lemma-fibred-groupoids-strict}.
Soit $p : \mathcal{S} \to \mathcal{C}$ une catégorie fibrée en groupoïdes.
Nous définissons un foncteur contravariant $F$ de $\mathcal{C}$ vers la
catégorie des groupoïdes comme suit : pour $U\in \Ob(\mathcal{C})$,
posons
$$
F(U) = \Mor_{\textit{Cat}/\mathcal{C}}(\mathcal{C}/U, \mathcal{S}).
$$
Si $f : U \to V$, le foncteur induit $\mathcal{C}/U \to \mathcal{C}/V$
induit le morphisme $F(f) : F(V) \to F(U)$. Il est clair que $F$ est un foncteur.
Soit $\mathcal{S}'$ la catégorie fibrée en groupoïdes associée par
l'exemple \ref{categories-example-functor-groupoids}.
Il existe un foncteur évident $G : \mathcal{S}' \to \mathcal{S}$
au-dessus de $\mathcal{C}$, qui envoie la paire $(U, x)$, où
$U \in \Ob(\mathcal{C})$ et $x \in F(U)$, sur
$x(\text{id}_U) \in \mathcal{S}$. Le
lemme \ref{categories-lemma-yoneda-2category}
implique maintenant que, pour tout $U$,
$$
G_U : \mathcal{S}'_U = F(U)=
\Mor_{\textit{Cat}/\mathcal{C}}(\mathcal{C}/U, \mathcal{S})
\to
\mathcal{S}_U
$$
est une équivalence ; ainsi, $G$ est une équivalence entre $\mathcal{S}$ et
$\mathcal{S}'$ d'après le lemme \ref{categories-lemma-equivalence-fibred-categories}.
\end{remark}





\section{1-flèches représentables}
\label{categories-section-representable-1-morphisms}

\noindent
Soit $\mathcal{C}$ une catégorie.
Dans cette section, nous expliquons ce que signifie, pour une $1$-flèche
entre catégories fibrées en groupoïdes au-dessus de $\mathcal{C}$,
d'être représentable.

\medskip\noindent
Soit $\mathcal{C}$ une catégorie.
Soient $\mathcal{X}$, $\mathcal{Y}$ des catégories fibrées en groupoïdes
au-dessus de $\mathcal{C}$.
Soit $U \in \Ob(\mathcal{C})$.
Soient $F : \mathcal{X} \to \mathcal{Y}$ et
$G : \mathcal{C}/U \to \mathcal{Y}$ des $1$-morphismes de catégories
fibrées en groupoïdes au-dessus de $\mathcal{C}$.
Nous voulons décrire
le $2$-produit fibré
$$
\xymatrix{
(\mathcal{C}/U) \times_\mathcal{Y} \mathcal{X} \ar[r] \ar[d] &
\mathcal{X} \ar[d]^F \\
\mathcal{C}/U \ar[r]^G &
\mathcal{Y}
}
$$
Soit $y = G(\text{id}_U) \in \mathcal{Y}_U$.
Choisissons un clivage de $\mathcal{Y}$
(voir la définition \ref{categories-definition-pullback-functor-fibred-category}).
Alors $G$ est isomorphe au foncteur $(f : V \to U) \mapsto f^*y$,
voir le lemme \ref{categories-lemma-yoneda-2category} et sa démonstration.
On peut considérer un objet de
$(\mathcal{C}/U) \times_\mathcal{Y} \mathcal{X}$
comme un quadruplet $(V, f : V \to U, x, \phi)$, voir le
lemme \ref{categories-lemma-2-product-categories-over-C}.
Avec la description de $G$ ci-dessus, on peut considérer $\phi$ comme
un isomorphisme $\phi : f^*y \to F(x)$ dans $\mathcal{Y}_V$.

\begin{lemma}
\label{categories-lemma-identify-fibre-product}
Dans la situation ci-dessus, la catégorie fibre de
$(\mathcal{C}/U) \times_\mathcal{Y} \mathcal{X}$ au-dessus
d'un objet $f : V \to U$ de $\mathcal{C}/U$
est la catégorie décrite comme suit :
\begin{enumerate}
\item ses objets sont les paires $(x, \phi)$,
où $x \in \Ob(\mathcal{X}_V)$ et
$\phi : f^*y \to F(x)$ est un morphisme de $\mathcal{Y}_V$ ;
\item l'ensemble des morphismes entre $(x, \phi)$ et $(x', \phi')$
est l'ensemble des morphismes $\psi : x \to x'$ dans $\mathcal{X}_V$
tels que $F(\psi) = \phi' \circ \phi^{-1}$.
\end{enumerate}
\end{lemma}

\begin{proof}
Voir la discussion ci-dessus.
\end{proof}

\begin{lemma}
\label{categories-lemma-prepare-representable-map-stack-in-groupoids}
Soit $\mathcal{C}$ une catégorie.
Soient $\mathcal{X}$, $\mathcal{Y}$ des catégories fibrées en groupoïdes
au-dessus de $\mathcal{C}$.
Soit $F : \mathcal{X} \to \mathcal{Y}$ un $1$-morphisme.
Soit $G : \mathcal{C}/U \to \mathcal{Y}$ un $1$-morphisme.
Alors
$$
(\mathcal{C}/U) \times_\mathcal{Y} \mathcal{X}
\longrightarrow
\mathcal{C}/U
$$
est une catégorie fibrée en groupoïdes.
\end{lemma}

\begin{proof}
Nous avons déjà vu dans le lemme \ref{categories-lemma-2-product-fibred-categories}
que la composée
$$
(\mathcal{C}/U) \times_\mathcal{Y} \mathcal{X}
\longrightarrow
\mathcal{C}/U
\longrightarrow
\mathcal{C}
$$
est une catégorie fibrée en groupoïdes. Le lemme découle alors du
lemme \ref{categories-lemma-cute-groupoids}.
\end{proof}

\begin{definition}
\label{categories-definition-representable-map-categories-fibred-in-groupoids}
Soit $\mathcal{C}$ une catégorie.
Soient $\mathcal{X}$, $\mathcal{Y}$ des catégories fibrées en groupoïdes
au-dessus de $\mathcal{C}$.
Soit $F : \mathcal{X} \to \mathcal{Y}$ un $1$-morphisme.
Nous disons que $F$ est {\it représentable}, ou que
{\it $\mathcal{X}$ est relativement représentable au-dessus de $\mathcal{Y}$},
si, pour tout $U \in \Ob(\mathcal{C})$
et tout $G : \mathcal{C}/U \to \mathcal{Y}$,
la catégorie fibrée en groupoïdes
$$
(\mathcal{C}/U) \times_\mathcal{Y} \mathcal{X}
\longrightarrow
\mathcal{C}/U
$$
est représentable.
\end{definition}

\begin{lemma}
\label{categories-lemma-spell-out-representable-map-stack-in-groupoids}
Soit $\mathcal{C}$ une catégorie.
Soient $\mathcal{X}$, $\mathcal{Y}$ des catégories fibrées en groupoïdes
au-dessus de $\mathcal{C}$.
Soit $F : \mathcal{X} \to \mathcal{Y}$ un $1$-morphisme.
Si $F$ est représentable, alors chacun des foncteurs
$$
F_U : \mathcal{X}_U \longrightarrow \mathcal{Y}_U
$$
entre catégories fibres est fidèle.
\end{lemma}

\begin{proof}
Cela résulte de la description des catégories fibres du
lemme \ref{categories-lemma-identify-fibre-product} et de la caractérisation
des catégories fibrées représentables du
lemme \ref{categories-lemma-characterize-representable-fibred-category}.
\end{proof}

\begin{lemma}
\label{categories-lemma-criterion-representable-map-stack-in-groupoids}
Soit $\mathcal{C}$ une catégorie.
Soient $\mathcal{X}$, $\mathcal{Y}$ des catégories fibrées en groupoïdes
au-dessus de $\mathcal{C}$.
Soit $F : \mathcal{X} \to \mathcal{Y}$ un $1$-morphisme.
Choisissons un clivage de $\mathcal{Y}$.
Supposons que
\begin{enumerate}
\item tout foncteur $F_U : \mathcal{X}_U \longrightarrow \mathcal{Y}_U$
entre catégories fibres soit fidèle, et
\item pour tout $U$ et tout $y \in \mathcal{Y}_U$, le préfaisceau
$$
(f : V \to U)
\longmapsto
\{(x, \phi) \mid x \in \mathcal{X}_V, \phi : f^*y \to F(x)\}/\cong
$$
soit un préfaisceau représentable sur $\mathcal{C}/U$.
\end{enumerate}
Alors $F$ est représentable.
\end{lemma}

\begin{proof}
Cela résulte de la description des catégories fibres du
lemme \ref{categories-lemma-identify-fibre-product} et de la caractérisation
des catégories fibrées représentables du
lemme \ref{categories-lemma-characterize-representable-fibred-category}.
\end{proof}

\noindent
Avant d'énoncer le lemme suivant, signalons que la $2$-catégorie
des catégories fibrées en groupoïdes est une $(2, 1)$-catégorie, et que nous
savons donc ce que signifie pour elle d'avoir un objet final (voir la
définition \ref{categories-definition-final-object-2-category}). Elle possède effectivement
un objet final, à savoir $\text{id} : \mathcal{C} \to \mathcal{C}$.
Nous définissons donc les {\it $2$-produits} de catégories fibrées en groupoïdes
au-dessus de $\mathcal{C}$ comme les $2$-produits fibrés
$$
\mathcal{X} \times \mathcal{Y} :=
\mathcal{X} \times_\mathcal{C} \mathcal{Y}.
$$
Avec cette définition, le lemme suivant a un sens.

\begin{lemma}
\label{categories-lemma-representable-diagonal-groupoids}
Soit $\mathcal{C}$ une catégorie.
Soit $\mathcal{S} \to \mathcal{C}$ une catégorie fibrée en groupoïdes.
Supposons que $\mathcal{C}$ admette les produits de paires d'objets et les
produits fibrés.
Les conditions suivantes sont équivalentes :
\begin{enumerate}
\item La diagonale $\mathcal{S} \to \mathcal{S} \times \mathcal{S}$
est représentable.
\item Pour tout $U$ dans $\mathcal{C}$, tout
$G : \mathcal{C}/U \to \mathcal{S}$ est représentable.
\end{enumerate}
\end{lemma}

\begin{proof}
Supposons la diagonale représentable, et fixons $U, G$.
Considérons un objet quelconque $V \in \Ob(\mathcal{C})$ et un
$G' : \mathcal{C}/V \to \mathcal{S}$ quelconque.
Notons que $\mathcal{C}/U \times \mathcal{C}/V = \mathcal{C}/(U \times V)$
est représentable. Par hypothèse, le produit fibré
$$
\xymatrix{
(\mathcal{C}/(U \times V))
\times_{(\mathcal{S} \times \mathcal{S})}
\mathcal{S}
\ar[r] \ar[d] &
\mathcal{S} \ar[d] \\
\mathcal{C}/(U \times V) \ar[r]^{(G, G')} &
\mathcal{S} \times \mathcal{S}
}
$$
est représentable.
Cela signifie qu'il existe $W \to U \times V$ dans $\mathcal{C}$,
tel que
$$
\xymatrix{
\mathcal{C}/W \ar[d] \ar[r] & \mathcal{S} \ar[d] \\
\mathcal{C}/U \times \mathcal{C}/V \ar[r] & \mathcal{S} \times \mathcal{S}
}
$$
soit cartésien. Cela implique que
$\mathcal{C}/W \cong \mathcal{C}/U \times_\mathcal{S} \mathcal{C}/V$
(voir le lemme \ref{categories-lemma-diagonal-1} et la
remarque \ref{categories-remark-2-product-fibred-categories-same}),
comme voulu.

\medskip\noindent
Supposons (2). Considérons un objet quelconque $V \in \Ob(\mathcal{C})$
et un $(G, G') : \mathcal{C}/V \to \mathcal{S} \times \mathcal{S}$
quelconque. Nous devons montrer que
$\mathcal{C}/V \times_{\mathcal{S} \times \mathcal{S}} \mathcal{S}$
est représentable. Nous savons que
$\mathcal{C}/V \times_{G, \mathcal{S}, G'} \mathcal{C}/V$
est représentable, disons par $a : W \to V$ dans $\mathcal{C}/V$.
L'équivalence
$$
\mathcal{C}/W \to \mathcal{C}/V \times_{G, \mathcal{S}, G'} \mathcal{C}/V
$$
suivie de la seconde projection vers $\mathcal{C}/V$ donne un
second morphisme $a' : W \to V$. Considérons
$W' = W \times_{(a, a'), V \times V} V$.
Il existe une équivalence
$$
\mathcal{C}/W' \cong
\mathcal{C}/V \times_{\mathcal{S} \times \mathcal{S}} \mathcal{S}
$$
à savoir
\begin{eqnarray*}
\mathcal{C}/W' & \cong &
\mathcal{C}/W \times_{(\mathcal{C}/V \times \mathcal{C}/V)} \mathcal{C}/V \\
& \cong &
\left(\mathcal{C}/V \times_{(G, \mathcal{S}, G')} \mathcal{C}/V\right)
\times_{(\mathcal{C}/V \times \mathcal{C}/V)} \mathcal{C}/V \\
& \cong &
\mathcal{C}/V \times_{(\mathcal{S} \times \mathcal{S})} \mathcal{S}
\end{eqnarray*}
(pour le dernier isomorphisme, voir le lemme \ref{categories-lemma-diagonal-2} et la
remarque \ref{categories-remark-2-product-fibred-categories-same}),
ce qui démontre le lemme.
\end{proof}

\noindent
{\bf Notes bibliographiques :}
Certaines parties de ce qui précède sont tirées des notes de Vistoli \cite{Vis2}.





\section{Catégories monoïdales}
\label{categories-section-monoidal}

\noindent
Soit $\mathcal{C}$ une catégorie. Supposons donné un foncteur
$$
\otimes : \mathcal{C} \times \mathcal{C} \longrightarrow \mathcal{C}
$$
On cherche souvent à savoir si $\otimes$ satisfait une loi d'associativité
et s'il existe une unité pour $\otimes$.

\medskip\noindent
Une {\it contrainte d'associativité} pour $(\mathcal{C}, \otimes)$ est
un isomorphisme fonctoriel
$$
\phi_{X, Y, Z} : X \otimes (Y \otimes Z) \to (X \otimes Y) \otimes Z
$$
tel que, pour tous objets $X, Y, Z, W$, le diagramme
\begin{equation}
\label{categories-equation-pentagon}
\vcenter{
\xymatrix{
X \otimes (Y \otimes ( Z \otimes W)) \ar[r] \ar[d] &
X \otimes ((Y \otimes Z) \otimes W) \ar[d] \\
(X \otimes Y) \otimes (Z \otimes W) \ar[d] &
(X \otimes (Y \otimes Z)) \otimes W \ar[ld] \\
((X \otimes Y) \otimes Z) \otimes W
}
}
\end{equation}
soit commutatif, toute flèche étant obtenue en appliquant convenablement
$\phi$ et en utilisant la fonctorialité de $\otimes$.

\medskip\noindent
Dans \cite[Théorème 3.1]{associativity}, il est expliqué en quel sens un triplet
$(\mathcal{C}, \otimes, \phi)$ comme ci-dessus est cohérent. Une reformulation
consiste à dire qu'un tel triplet donne lieu à un système de foncteurs bien définis
$$
\mathcal{C} \times \ldots \times \mathcal{C} \longrightarrow \mathcal{C},
\quad
(X_1, \ldots, X_n) \longmapsto X_1 \otimes \ldots \otimes X_n
$$
pour tout $n \geq 1$, ainsi qu'à des isomorphismes fonctoriels, pour
$n, m \geq 1$,
$$
(X_1 \otimes \ldots \otimes X_n) \otimes (Y_1 \otimes \ldots \otimes Y_m)
\to
X_1 \otimes \ldots \otimes X_n \otimes Y_1 \otimes \ldots \otimes Y_m
$$
tels que tous les diagrammes que l'on peut former à partir de ces isomorphismes
commutent (le diagramme du pentagone ci-dessus en est un exemple). Pour $n = 5$,
cela signifie que les expressions
$$
\begin{matrix}
(((X_1 \otimes X_2) \otimes X_3) \otimes X_4) \otimes X_5 &
((X_1 \otimes (X_2 \otimes X_3)) \otimes X_4) \otimes X_5 \\
((X_1 \otimes X_2) \otimes (X_3 \otimes X_4)) \otimes X_5 &
(X_1 \otimes ((X_2 \otimes X_3) \otimes X_4)) \otimes X_5 \\
(X_1 \otimes (X_2 \otimes (X_3 \otimes X_4))) \otimes X_5 &
((X_1 \otimes X_2) \otimes X_3) \otimes (X_4 \otimes X_5) \\
(X_1 \otimes (X_2 \otimes X_3)) \otimes (X_4 \otimes X_5) &
(X_1 \otimes X_2) \otimes ((X_3 \otimes X_4) \otimes X_5) \\
(X_1 \otimes X_2) \otimes (X_3 \otimes (X_4 \otimes X_5)) &
X_1 \otimes (((X_2 \otimes X_3) \otimes X_4) \otimes X_5) \\
X_1 \otimes ((X_2 \otimes (X_3 \otimes X_4)) \otimes X_5) &
X_1 \otimes ((X_2 \otimes X_3) \otimes (X_4 \otimes X_5)) \\
X_1 \otimes (X_2 \otimes ((X_3 \otimes X_4) \otimes X_5)) &
X_1 \otimes (X_2 \otimes (X_3 \otimes (X_4 \otimes X_5)))
\end{matrix}
$$
sont toutes fonctoriellement isomorphes à
$X_1 \otimes \ldots \otimes X_5$ et que ces isomorphismes sont tous
compatibles avec tous les isomorphismes possibles que l'on obtient entre ces
foncteurs en appliquant $\phi$ à trois positions consécutives (lorsque cela est
possible). Nous utiliserons ce fait sans autre mention dans la suite et
n'écrirons plus les parenthèses dans les itérations de $\otimes$.

\medskip\noindent
Une {\it unité} pour un triplet $(\mathcal{C}, \otimes, \phi)$ comme ci-dessus
est un objet $\mathbf{1}$ de $\mathcal{C}$ muni d'isomorphismes fonctoriels
$$
l : \mathbf{1} \otimes X \to X
\quad\text{et}\quad
r : X \otimes \mathbf{1} \to X
$$
tels que, pour tous objets $X, Y$, le diagramme
\begin{equation}
\label{categories-equation-triangle}
\vcenter{
\xymatrix{
X \otimes (\mathbf{1} \otimes Y) \ar[rr]_\phi \ar[rd]_{\text{id} \otimes l} & &
(X \otimes \mathbf{1}) \otimes Y \ar[ld]^{r \otimes \text{id}} \\
& X \otimes Y
}
}
\end{equation}
soit commutatif. Nous considérerons souvent les unités comme des paires
$(\mathbf{1}, 1)$, comme dans le lemme suivant.

\begin{lemma}
\label{categories-lemma-monoidal-unit}
Soit $(\mathcal{C}, \otimes, \phi)$ comme ci-dessus. Il existe une
correspondance bijective entre les unités $(\mathbf{1}, l, r)$ de $\mathcal{C}$
et les paires $(\mathbf{1}, 1)$ où $\mathbf{1}$ est un objet de $\mathcal{C}$ et
$1 : \mathbf{1} \otimes \mathbf{1} \to \mathbf{1}$
est un isomorphisme tel que les foncteurs
$L : X \mapsto \mathbf{1} \otimes X$
et $R : X \mapsto X \otimes \mathbf{1}$ soient des équivalences.
\end{lemma}

\begin{proof}
Étant donnée une unité $(\mathbf{1}, l, r)$, on obtient un isomorphisme
$r : \mathbf{1} \otimes \mathbf{1} \to \mathbf{1}$, et $L$ et $R$ sont
des équivalences puisqu'ils sont isomorphes au foncteur identité.
Réciproquement, supposons donnée $(\mathbf{1}, 1)$ telle que
$L$ et $R$ soient des équivalences. On obtient des isomorphismes fonctoriels
$l_X : \mathbf{1} \otimes X \to X$ et $r_X : X \otimes \mathbf{1} \to X$
caractérisés par $L(l_X) = 1 \otimes \text{id}_X$ et
$R(r_X) = \text{id}_X \otimes 1$. Il nous faut alors montrer que les deux
flèches $\text{id}_X \otimes l_Y$ et $r_X \otimes \text{id}_Y$ de
$X \otimes \mathbf{1} \otimes Y$ vers $X \otimes Y$ sont égales
pour tous $X$ et $Y$. Cette propriété ne dépend que des classes d'isomorphisme
de $X$ et $Y$.
Puisque $R$ et $L$ sont des équivalences, il suffit de
le vérifier pour $X = Z \otimes \mathbf{1}$ et $Y = \mathbf{1} \otimes W$
pour certains objets $Z$ et $W$.
Autrement dit, nous devons montrer que
$$
\text{id}_Z \otimes \text{id}_\mathbf{1} \otimes l_{\mathbf{1} \otimes W}
=
r_{Z \otimes \mathbf{1}} \otimes \text{id}_\mathbf{1} \otimes \text{id}_W
$$
Par construction, ces morphismes sont respectivement égaux à
$\text{id}_Z \otimes 1 \otimes \text{id}_\mathbf{1} \otimes \text{id}_W$ et
$\text{id}_Z \otimes \text{id}_\mathbf{1} \otimes 1 \otimes \text{id}_W$.
Il suffit donc de montrer que
$1 \otimes \text{id}_\mathbf{1} = \text{id}_\mathbf{1} \otimes 1$.

\medskip\noindent
Nous avons
$\text{id}_\mathbf{1} \otimes 1 = r_\mathbf{1} \otimes \text{id}_\mathbf{1}$
et
$l_\mathbf{1} \otimes \text{id}_\mathbf{1} = \text{id}_\mathbf{1} \otimes 1$.
On peut écrire $r_\mathbf{1} = a \circ 1$ et $l_\mathbf{1} = b \circ 1$
pour certains automorphismes $a, b$ de $\mathbf{1}$. Ainsi, on a
$\text{id}_\mathbf{1} \otimes 1 =
(a \otimes \text{id}_\mathbf{1}) \circ (1 \otimes \text{id}_\mathbf{1})$
et
$1 \otimes \text{id}_\mathbf{1} =
(\text{id}_\mathbf{1} \otimes b) \circ (\text{id}_\mathbf{1} \otimes 1)$.
On peut alors écrire
\begin{align*}
(1 \otimes \text{id}_\mathbf{1}) \circ
(\text{id}_\mathbf{1} \otimes 1 \otimes \text{id}_\mathbf{1})
& =
(1 \otimes \text{id}_\mathbf{1}) \circ
(\text{id}_\mathbf{1} \otimes \text{id}_\mathbf{1} \otimes b)
\circ
(\text{id}_\mathbf{1} \otimes \text{id}_\mathbf{1} \otimes 1) \\
& =
(\text{id}_\mathbf{1} \otimes b) \circ
(1 \otimes \text{id}_\mathbf{1}) \circ
(\text{id}_\mathbf{1} \otimes \text{id}_\mathbf{1} \otimes 1) \\
& =
(\text{id}_\mathbf{1} \otimes b) \circ
(1 \otimes 1)
\end{align*}
et l'on a aussi
\begin{align*}
(1 \otimes \text{id}_\mathbf{1}) \circ
(\text{id}_\mathbf{1} \otimes 1 \otimes \text{id}_\mathbf{1})
& =
(\text{id}_\mathbf{1} \otimes b) \circ
(\text{id}_\mathbf{1} \otimes 1) \circ
(a \otimes \text{id}_\mathbf{1} \otimes \text{id}_\mathbf{1}) \circ
(1 \otimes \text{id}_\mathbf{1} \otimes \text{id}_\mathbf{1}) \\
& =
(\text{id}_\mathbf{1} \otimes b) \circ
(a \otimes \text{id}_\mathbf{1}) \circ
(1 \otimes 1)
\end{align*}
Cela montre que $a \otimes \text{id}_\mathbf{1}$ est l'identité, et donc
que $a$ est l'identité, comme voulu.

\medskip\noindent
Pour achever la démonstration, notons que les règles ci-dessus déterminent
des équivalences inverses de catégories entre la catégorie des unités
(convenablement définie) et la catégorie des paires $(\mathbf{1}, 1)$.
\end{proof}

\begin{lemma}
\label{categories-lemma-monoidal-unit-unique}
Soit $(\mathcal{C}, \otimes, \phi)$ comme ci-dessus.
Soit $(\mathbf{1}, 1)$ une unité (voir le lemme \ref{categories-lemma-monoidal-unit}).
Alors
\begin{enumerate}
\item $1 \otimes \text{id}_\mathbf{1} = \text{id}_\mathbf{1} \otimes 1$
\item $\Gamma = \text{Mor}(\mathbf{1}, \mathbf{1})$ est un monoïde commutatif,
\item  $a = 1 \circ (a \otimes \text{id}_\mathbf{1}) \circ 1^{-1} =
1 \circ (\text{id}_\mathbf{1} \otimes a) \circ 1^{-1}$ pour tout
$a \in \Gamma$,
\item toute autre unité est isomorphe à $(\mathbf{1}, 1)$
par un unique isomorphisme.
\end{enumerate}
\end{lemma}

\begin{proof}
La partie (1) a été démontrée dans la preuve du
lemme \ref{categories-lemma-monoidal-unit}.
Pour $a \in \Gamma$, on a
\begin{align*}
(1 \circ (a \otimes \text{id}_\mathbf{1})) \otimes \text{id}_\mathbf{1}
& =
(1 \otimes \text{id}_\mathbf{1}) \circ
(a \otimes \text{id}_\mathbf{1} \otimes \text{id}_\mathbf{1}) \\
& =
(\text{id}_\mathbf{1} \otimes 1) \circ
(a \otimes \text{id}_\mathbf{1} \otimes \text{id}_\mathbf{1}) \\
& =
(a \otimes \text{id}_\mathbf{1}) \circ
(\text{id}_\mathbf{1} \otimes 1) \\
& =
(a \otimes \text{id}_\mathbf{1}) \circ
(1 \otimes \text{id}_\mathbf{1}) \\
& =
(a \circ 1) \otimes \text{id}_\mathbf{1}
\end{align*}
Ainsi, $1 \circ (a \otimes \text{id}_\mathbf{1}) = a \circ 1$,
ce qui donne une moitié de (3). L'autre moitié s'obtient
de la même manière. Pour voir (2), observons que, pour $a, b \in \Gamma$,
on a
$$
a \circ b =
1 \circ (a \otimes \text{id}_\mathbf{1}) \circ 1^{-1}
\circ
1 \circ (\text{id}_\mathbf{1} \otimes b) \circ 1^{-1}
=
1 \circ (a \otimes b) \circ 1^{-1}
$$
et la même expression donne $b \circ a$. Ainsi, (2) est satisfaite.
Pour voir (4), supposons que $(\mathbf{1}', 1')$ soit une seconde unité.
En utilisant $r$ et $l'$, on obtient des isomorphismes
$\mathbf{1}' \otimes \mathbf{1} \to \mathbf{1}'$ et
$\mathbf{1}' \otimes \mathbf{1} \to \mathbf{1}$. Il existe donc
un isomorphisme $t : \mathbf{1}' \to \mathbf{1}$.
Alors le diagramme
$$
\xymatrix{
\mathbf{1}' \otimes \mathbf{1}' \ar[r]_{t \otimes t} \ar[d]_{1'} &
\mathbf{1} \otimes \mathbf{1} \ar[d]^1 \\
\mathbf{1}' \ar[r]^t & \mathbf{1}
}
$$
commute à un automorphisme $a$ de $\mathbf{1}$ près. Après avoir remplacé
$t$ par $a \circ t$, le diagramme est commutatif (indication : utiliser (3) pour voir
que $1 \circ (a \otimes a) = a^2 \circ 1$).
\end{proof}

\begin{definition}
\label{categories-definition-monoidal-category}
Un triplet $(\mathcal{C}, \otimes, \phi)$ où $\mathcal{C}$ est une catégorie,
$\otimes : \mathcal{C} \times \mathcal{C} \to \mathcal{C}$ est un foncteur,
et $\phi$ est une contrainte d'associativité, est appelé
{\it catégorie monoïdale} s'il existe une unité $\mathbf{1}$.
\end{definition}

\noindent
Nous écrivons toujours $\mathbf{1}$ pour désigner une unité d'une catégorie
monoïdale et nous notons $1 : \mathbf{1} \otimes \mathbf{1} \to \mathbf{1}$
un isomorphisme choisi ; puisque la paire $(\mathbf{1}, 1)$ est déterminée
à un unique isomorphisme près (lemme \ref{categories-lemma-monoidal-unit-unique}),
on peut en choisir une sans inconvénient.

\begin{lemma}
\label{categories-lemma-monoidal-coherence}
Dans une catégorie monoïdale $\mathcal{C}, \otimes, \phi, \mathbf{1}, 1$,
et avec les notations de la démonstration du lemme \ref{categories-lemma-monoidal-unit},
on a
\begin{enumerate}
\item les flèches $1, r_\mathbf{1}, l_\mathbf{1} :
\mathbf{1} \otimes \mathbf{1} \to \mathbf{1}$ coïncident ;
\item les flèches
$l_X \otimes \text{id}_Y, l_{X \otimes Y} :
\mathbf{1} \otimes X \otimes Y \to X \otimes Y$ coïncident ; et
\item les flèches
$\text{id}_X \otimes r_Y , r_{X \otimes Y} :
X \otimes Y \otimes \mathbf{1} \to X \otimes Y$ coïncident.
\end{enumerate}
Une catégorie monoïdale satisfait les hypothèses de
\cite[Théorème 5.2]{associativity}.
\end{lemma}

\begin{proof}
Nous avons établi (1) dans la démonstration du
lemme \ref{categories-lemma-monoidal-unit}.
Nous avons vu dans la démonstration du lemme \ref{categories-lemma-monoidal-unit}
que $l_X$ et $l_{X \otimes Y}$ sont
les uniques morphismes tels que
$\text{id}_\mathbf{1} \otimes l_{X \otimes Y} =
1 \otimes \text{id}_{X \otimes Y}$ et
$\text{id}_\mathbf{1} \otimes l_X = 1 \otimes \text{id}_X$.
La partie (2) en découle immédiatement. La partie (3) se démontre de manière
analogue. Avec la commutativité de (\ref{categories-equation-pentagon})
et de (\ref{categories-equation-triangle}), cela donne la dernière assertion du lemme.
\end{proof}

\noindent
Dans \cite[Théorème 5.2]{associativity}, il est expliqué en quel sens un
quintuplet $(\mathcal{C}, \otimes, \phi, \mathbf{1}, 1)$ comme dans le lemme
ci-dessus est cohérent. Une reformulation est que, dans une catégorie monoïdale
(au sens où nous l'entendons), pour tous $n \geq i \geq 0$, on dispose
d'isomorphismes fonctoriels
$$
X_1 \otimes \ldots \otimes X_i \otimes \mathbf{1}
\otimes X_{i + 1} \otimes \ldots \otimes X_n
\to
X_1 \otimes \ldots \otimes X_n
$$
tels que tous les diagrammes possibles formés à partir de ces isomorphismes
commutent. Ainsi, par exemple, si l'on part de deux insertions de $\mathbf{1}$
et que l'on utilise les isomorphismes dans des ordres différents, on obtient le
même morphisme. Outre la convention consistant à supprimer les parenthèses dans
les itérations de $\otimes$, nous identifions désormais $X \otimes \mathbf{1}$
et $\mathbf{1} \otimes X$ à $X$ sans autre mention.
En outre, nous dirons « soit $\mathcal{C}$ une catégorie monoïdale »,
les données $\otimes, \phi, \mathbf{1}$ étant sous-entendues.

\begin{definition}
\label{categories-definition-functor-monoidal-categories}
Soient $\mathcal{C}$ et $\mathcal{C}'$ des catégories monoïdales.
Un {\it foncteur monoïdal} $F : \mathcal{C} \to \mathcal{C}'$
est la donnée du foncteur $F$ indiqué et d'un isomorphisme
$$
F(X) \otimes F(Y) \to F(X \otimes Y)
$$
fonctoriel en $X$ et $Y$, tel que, pour tous objets $X$, $Y$ et $Z$,
le diagramme
$$
\xymatrix{
F(X) \otimes (F(Y) \otimes F(Z)) \ar[r] \ar[d] &
F(X) \otimes F(Y \otimes Z) \ar[r] &
F(X \otimes (Y \otimes Z)) \ar[d] \\
(F(X) \otimes F(Y)) \otimes F(Z) \ar[r] &
F(X \otimes Y) \otimes F(Z) \ar[r] &
F((X \otimes Y) \otimes Z)
}
$$
est commutatif et tel que $F(\mathbf{1})$ soit une unité de $\mathcal{C}'$.
\end{definition}

\noindent
Selon nos conventions sur les unités, on peut toujours supposer que
$F(\mathbf{1}) = \mathbf{1}$ si $F$ est un foncteur monoïdal.
À titre d'exemple, si $A \to B$ est un homomorphisme d'anneaux, alors
le foncteur $M \mapsto M \otimes_A B$ est un foncteur monoïdal de
$\text{Mod}_A$ vers $\text{Mod}_B$.

\begin{lemma}
\label{categories-lemma-invertible}
Soit $\mathcal{C}$ une catégorie monoïdale. Soit $X$ un objet de
$\mathcal{C}$. Les conditions suivantes sont équivalentes :
\begin{enumerate}
\item le foncteur $L : Y \mapsto X \otimes Y$ est une équivalence ;
\item le foncteur $R : Y \mapsto Y \otimes X$ est une équivalence ;
\item il existe un objet $X'$ tel que
$X \otimes X' \cong X' \otimes X \cong \mathbf{1}$.
\end{enumerate}
\end{lemma}

\begin{proof}
Supposons (1). Choisissons $X'$ tel que $L(X') = \mathbf{1}$, c'est-à-dire
$X \otimes X' \cong \mathbf{1}$. Notons $L'$ et $R'$ les foncteurs
correspondant à $X'$. L'égalité $X \otimes X' \cong \mathbf{1}$
implique $L \circ L' \cong \text{id}$. Ainsi, $L'$ est nécessairement le
quasi-inverse de $L$ (qui existe par hypothèse). Par conséquent,
$L' \circ L \cong \text{id}$.
Il s'ensuit que $X' \otimes X \cong \mathbf{1}$. Ainsi, (3) est satisfaite.

\medskip\noindent
La démonstration de (2) $\Rightarrow$ (3) est duale de celle qui précède.

\medskip\noindent
Supposons (3). Il est alors clair que $L'$ et $L$ sont quasi-inverses
l'un de l'autre, et que $R'$ et $R$ le sont également. Ainsi, (1) et (2)
sont satisfaites.
\end{proof}

\begin{definition}
\label{categories-definition-invertible}
Soit $\mathcal{C}$ une catégorie monoïdale. Un objet $X$ de $\mathcal{C}$
est dit {\it inversible} si l'une quelconque (donc toutes) des conditions
équivalentes du lemme \ref{categories-lemma-invertible} est satisfaite.
\end{definition}

\noindent
Observons que, si $F : \mathcal{C} \to \mathcal{C}'$ est un foncteur
monoïdal, alors $F$ envoie les objets inversibles sur des objets
inversibles.

\begin{definition}
\label{categories-definition-dual}
Étant donnés une catégorie monoïdale $(\mathcal{C}, \otimes, \phi)$
et un objet $X$, un {\it dual à gauche} de celui-ci est un objet $Y$ muni de
morphismes $\eta : \mathbf{1} \to X \otimes Y$ et
$\epsilon : Y \otimes X \to \mathbf{1}$
tels que les diagrammes
$$
\vcenter{
\xymatrix{
X \ar[rd]_1 \ar[r]_-{\eta \otimes 1} &
X \otimes Y \otimes X \ar[d]^{1 \otimes \epsilon} \\
& X
}
}
\quad\text{et}\quad
\vcenter{
\xymatrix{
Y \ar[rd]_1 \ar[r]_-{1 \otimes \eta} &
Y \otimes X \otimes Y \ar[d]^{\epsilon \otimes 1} \\
& Y
}
}
$$
commutent. Dans cette situation, on dit que $X$ est un
{\it dual à droite} de $Y$.
\end{definition}

\noindent
Observons que, si $F : \mathcal{C} \to \mathcal{C}'$ est un foncteur
monoïdal, alors $F(Y)$ est un dual à gauche de $F(X)$ si
$Y$ est un dual à gauche de $X$.

\begin{lemma}
\label{categories-lemma-left-dual}
Soit $\mathcal{C}$ une catégorie monoïdale. Si $Y$ est un dual à gauche de $X$,
alors
$$
\Mor(Z' \otimes X, Z) = \Mor(Z', Z \otimes Y)
\quad\text{et}\quad
\Mor(Y \otimes Z', Z) = \Mor(Z', X \otimes Z)
$$
fonctoriellement en $Z$ et $Z'$.
\end{lemma}

\begin{proof}
Considérons les applications
$$
\Mor(Z' \otimes X, Z) \to
\Mor(Z' \otimes X \otimes Y, Z \otimes Y) \to
\Mor(Z', Z \otimes Y)
$$
où l'on utilise $\eta$ dans la seconde flèche,
ainsi que la suite d'applications
$$
\Mor(Z', Z \otimes Y) \to
\Mor(Z' \otimes X, Z \otimes Y \otimes X) \to
\Mor(Z' \otimes X, Z)
$$
où l'on utilise $\epsilon$ dans la seconde flèche.
Pour montrer que ces flèches sont inverses l'une de l'autre, considérons une
morphisme $a : Z' \to Z \otimes Y$. Nous devons montrer que
$$
Z' \xrightarrow{\text{id}_{Z'} \otimes \eta}
Z' \otimes X \otimes Y \xrightarrow{a \otimes \text{id}_{X \otimes Y}}
Z \otimes Y \otimes X \otimes Y
\xrightarrow{\text{id}_Z \otimes \epsilon \otimes \text{id}_Y}
Z \otimes Y
$$
est égale à $a$. La composée des deux premières flèches est égale à
$(\text{id}_{Z \otimes Y} \otimes \eta) \circ a :
Z' \to Z \otimes Y \to Z \otimes Y \otimes X \otimes Y$.
Puis la composée de $\text{id}_{Z \otimes Y} \otimes \eta$
et de $\text{id}_Z \otimes \epsilon \otimes \text{id}_Y$ est
l'identité par définition du dual. Il en va de même pour l'autre composée.
Nous omettons la démonstration de la seconde égalité.
\end{proof}

\begin{remark}
\label{categories-remark-left-dual-adjoint}
Le lemme \ref{categories-lemma-left-dual} affirme en particulier que
$Z \mapsto Z \otimes Y$ est l'adjoint à droite de
$Z' \mapsto Z' \otimes X$. En particulier, l'unicité des foncteurs adjoints
garantit qu'un dual à gauche de $X$, s'il existe, est unique à un unique
isomorphisme près.
Réciproquement, supposons que le foncteur $Z \mapsto Z \otimes Y$ soit un
adjoint à droite du foncteur $Z' \mapsto Z' \otimes X$, c'est-à-dire que l'on
se donne une bijection
$$
\Mor(Z' \otimes X, Z) \longrightarrow \Mor(Z', Z \otimes Y)
$$
fonctorielle à la fois en $Z$ et en $Z'$. L'unité de l'adjonction fournit
des morphismes
$$
\eta_Z : Z \to Z \otimes X \otimes Y
$$
fonctoriels en $Z$, et la counité de l'adjonction fournit des morphismes
$$
\epsilon_{Z'} : Z' \otimes Y \otimes X \to Z'
$$
fonctoriels en $Z'$. En particulier, on obtient
$\eta = \eta_\mathbf{1} : \mathbf{1} \to X \otimes Y$ et
$\epsilon = \epsilon_\mathbf{1} : Y \otimes X \to \mathbf{1}$.
En guise d'exercice sur la relation entre les unités, les counités et l'isomorphisme
d'adjonction, le lecteur peut montrer que l'on a
$$
(\epsilon \otimes \text{id}_Y) \circ \eta_Y = \text{id}_Y
\quad\text{et}\quad
\epsilon_X \circ (\eta \otimes \text{id}_X) = \text{id}_X
$$
Cela ne suffit toutefois pas à montrer que
$(\epsilon \otimes \text{id}_Y) \circ (\text{id}_Y \otimes \eta) =
\text{id}_Y$ et
$(\text{id}_X \otimes \epsilon) \circ (\eta \otimes \text{id}_X) =
\text{id}_X$, car on ne sait pas en général que
$\eta_Y = \text{id}_Y \otimes \eta$, ni que
$\epsilon_X = \epsilon \otimes \text{id}_X$. Pour cela, il suffirait
de savoir que notre isomorphisme d'adjonction possède la propriété suivante :
pour tous $W, Z, Z'$, le diagramme
$$
\xymatrix{
\Mor(Z' \otimes X, Z) \ar[r] \ar[d]_{\text{id}_W \otimes -} &
\Mor(Z', Z \otimes Y) \ar[d]^{\text{id}_W \otimes -} \\
\Mor(W \otimes Z' \otimes X, W \otimes Z) \ar[r] &
\Mor(W \otimes Z', W \otimes Z \otimes Y)
}
$$
est commutatif.
Si tel est le cas, nous dirons que {\it l'adjonction est compatible avec
la structure tensorielle donnée}. Ainsi, exiger que
$Z \mapsto Z \otimes Y$ soit l'adjoint à droite de $Z' \mapsto Z' \otimes X$
compatible avec la structure tensorielle donnée constitue une formulation
équivalente de la propriété d'être un dual à gauche.
\end{remark}

\begin{lemma}
\label{categories-lemma-tensor-dual}
Soit $\mathcal{C}$ une catégorie monoïdale. Si $Y_i$, $i = 1, 2$,
sont des duaux à gauche de $X_i$, $i = 1, 2$, alors $Y_2 \otimes Y_1$ est
un dual à gauche de $X_1 \otimes X_2$.
\end{lemma}

\begin{proof}
Cela découle de l'unicité des adjoints et de la
remarque \ref{categories-remark-left-dual-adjoint}.
\end{proof}

\noindent
Une {\it contrainte de commutativité} pour $(\mathcal{C}, \otimes)$ est un
isomorphisme fonctoriel
$$
\psi : X \otimes Y \longrightarrow Y \otimes X
$$
tel que la composée
$$
X \otimes Y \xrightarrow{\psi} Y \otimes X \xrightarrow{\psi} X \otimes Y
$$
soit l'identité. On dit que $\psi$ est {\it compatible} avec une contrainte
d'associativité donnée $\phi$ si, pour tous objets $X, Y, Z$, le diagramme
\begin{equation}
\label{categories-equation-hexagon}
\vcenter{
\xymatrix{
X \otimes (Y \otimes Z) \ar[r]_\phi \ar[d]^\psi &
(X \otimes Y) \otimes Z \ar[r]_\psi &
Z \otimes (X \otimes Y) \ar[d]^\phi \\
X \otimes (Z \otimes Y) \ar[r]^\phi &
(X \otimes Z) \otimes Y \ar[r]^\psi &
(Z \otimes X) \otimes Y
}
}
\end{equation}
est commutatif.

\begin{definition}
\label{categories-definition-symmetric-monoidal-category}
Un quadruplet $(\mathcal{C}, \otimes, \phi, \psi)$ où
$\mathcal{C}$ est une catégorie,
$\otimes : \mathcal{C} \times \mathcal{C} \to \mathcal{C}$ est un foncteur,
$\phi$ est une contrainte d'associativité et
$\psi$ est une contrainte de commutativité compatible avec $\phi$,
est appelé {\it catégorie monoïdale symétrique} s'il existe
une unité.
\end{definition}

\noindent
Bien entendu, si $(\mathcal{C}, \otimes, \phi, \psi)$ est une
catégorie monoïdale symétrique, alors $(\mathcal{C}, \otimes, \phi)$
est une catégorie monoïdale et nous pouvons employer le langage et les
notations introduits ci-dessus.

\begin{lemma}
\label{categories-lemma-symmetric-monoidal-coherence}
Dans une catégorie monoïdale symétrique
$\mathcal{C}, \otimes, \phi, \psi, \mathbf{1}, 1$,
on a
\begin{enumerate}
\item les flèches $1 \circ \psi, 1 :
\mathbf{1} \otimes \mathbf{1} \to \mathbf{1}$ coïncident ;
\item les flèches $\text{id}_X \otimes l_Y,
(l_X \otimes \text{id}) \circ (\psi \otimes \text{id}_Y):
X \otimes \mathbf{1} \otimes Y \to X \otimes Y$ coïncident,
\end{enumerate}
Une catégorie monoïdale symétrique satisfait les hypothèses de
\cite[Théorème 5.1]{associativity}.
\end{lemma}

\begin{proof}
On peut écrire $\psi = a \otimes \text{id}_\mathbf{1}$ pour un unique
isomorphisme $a : \mathbf{1} \to \mathbf{1}$.
Le lemme \ref{categories-lemma-monoidal-unit-unique} implique que
$a \otimes \text{id}_\mathbf{1} = \text{id}_\mathbf{1} \otimes a$.
La fonctorialité de $\psi$ affirme que le diagramme
$$
\xymatrix{
(\mathbf{1} \otimes \mathbf{1}) \otimes \mathbf{1}
\ar[r]_\psi
\ar[d]_{1 \otimes \text{id}_\mathbf{1}} &
\mathbf{1} \otimes (\mathbf{1} \otimes \mathbf{1})
\ar[d]^{\text{id}_\mathbf{1} \otimes 1} \\
\mathbf{1} \otimes \mathbf{1}  \ar[r]^\psi &
\mathbf{1} \otimes \mathbf{1}
}
$$
est commutatif. Ainsi, la flèche du haut est égale à
$a \otimes \text{id}_\mathbf{1} \otimes \text{id}_\mathbf{1}$.
Par conséquent, (\ref{categories-equation-hexagon}) pour
$X = Y = Z = \mathbf{1}$ affirme que
$a \otimes \text{id}_\mathbf{1} \otimes \text{id}_\mathbf{1}$
est égal à son carré. Ainsi, $a = \text{id}_\mathbf{1}$.
Cela démontre (1).

\medskip\noindent
La partie (2) affirme que $\psi : X \otimes \mathbf{1} \to \mathbf{1} \otimes X$
est l'identité lorsque l'on identifie la source et le but à $X$
dans notre catégorie monoïdale. Cela résulte du fait que, pour
(\ref{categories-equation-hexagon}) et $\mathbf{1}, X, \mathbf{1}$, le diagramme
$$
\xymatrix{
\mathbf{1} \otimes (X \otimes \mathbf{1}) \ar[r]_\phi \ar[d]^\psi &
(\mathbf{1} \otimes X) \otimes \mathbf{1} \ar[r]_\psi &
\mathbf{1} \otimes (\mathbf{1} \otimes X) \ar[d]^\phi \\
\mathbf{1} \otimes (\mathbf{1} \otimes X) \ar[r]^\phi &
(\mathbf{1} \otimes \mathbf{1}) \otimes X \ar[r]^\psi &
(\mathbf{1} \otimes \mathbf{1}) \otimes X
}
$$
est commutatif et que tous les morphismes $\psi$ sauf un
sont égaux à l'identité.

\medskip\noindent
Outre (1) et (2), la commutativité des diagrammes
(\ref{categories-equation-pentagon}), (\ref{categories-equation-triangle}), (\ref{categories-equation-hexagon})
et les résultats du lemme \ref{categories-lemma-monoidal-coherence}
impliquent la dernière assertion du lemme.
\end{proof}

\noindent
Dans \cite[Théorème 5.1]{associativity}, il est expliqué en quel sens un
sextuplet $(\mathcal{C}, \otimes, \phi, \psi, \mathbf{1}, 1)$ comme dans le
lemme ci-dessus est cohérent. Une reformulation est que, dans une catégorie
monoïdale symétrique (au sens où nous l'entendons), pour tout $n \geq 1$ et
toute permutation $\sigma$ de $\{1, \ldots, n\}$, on dispose d'isomorphismes
fonctoriels
$$
X_1 \otimes \ldots \otimes X_n
\to
X_{\sigma(1)} \otimes \ldots \otimes X_{\sigma(n)}
$$
tels que tous les diagrammes possibles formés à partir de ces isomorphismes
commutent et que ces isomorphismes soient compatibles avec la structure de
catégorie monoïdale (notamment avec les isomorphismes obtenus lorsque l'on insère
un $\mathbf{1}$ dans une position).

\begin{lemma}
\label{categories-lemma-dual-symmetric}
Soit $(\mathcal{C}, \otimes, \phi, \psi)$ une catégorie monoïdale symétrique.
Soit $X$ un objet de $\mathcal{C}$, et supposons que $Y$,
$\eta : \mathbf{1} \to X \otimes Y$, et
$\epsilon : Y \otimes X \to \mathbf{1}$
définissent un dual à gauche de $X$ comme dans la définition \ref{categories-definition-dual}.
Alors $\eta' = \psi \circ \eta : \mathbf{1} \to Y \otimes X$
et $\epsilon' = \epsilon \circ \psi : X \otimes Y \to \mathbf{1}$
font de $X$ un dual à gauche de $Y$.
\end{lemma}

\begin{proof}
Omis. Indication : un exercice agréable sur les définitions.
\end{proof}

\begin{definition}
\label{categories-definition-functor-symmetric-monoidal-categories}
Soient $\mathcal{C}$ et $\mathcal{C}'$ des catégories monoïdales symétriques.
Un {\it foncteur monoïdal symétrique}
$F : \mathcal{C} \to \mathcal{C}'$
est la donnée du foncteur $F$ indiqué et d'un isomorphisme
$$
F(X) \otimes F(Y) \to F(X \otimes Y)
$$
fonctoriel en $X$ et $Y$, tel que $F$ soit un foncteur monoïdal et tel que,
pour tous objets $X$ et $Y$, le diagramme
$$
\xymatrix{
F(X) \otimes F(Y) \ar[r] \ar[d] &
F(X \otimes Y) \ar[d] \\
F(Y) \otimes F(X) \ar[r] &
F(Y \otimes X)
}
$$
est commutatif.
\end{definition}

\begin{remark}
\label{categories-remark-internal-hom-monoidal}
Soit $\mathcal{C}$ une catégorie monoïdale. Nous disons que $\mathcal{C}$
possède un {\it Hom interne} si, pour toute paire d'objets $X, Y$ de
$\mathcal{C}$, il existe un objet $hom(X, Y)$ de $\mathcal{C}$ tel que l'on ait
$$
\Mor(X, hom(Y, Z)) = \Mor(X \otimes Y, Z)
$$
fonctoriellement en $X, Y, Z$. Par le lemme de Yoneda, le bifoncteur
$(X, Y) \mapsto hom(X, Y)$ est déterminé à un unique isomorphisme près,
s'il existe. Étant donné un Hom interne, on obtient des morphismes canoniques
\begin{enumerate}
\item $hom(X, Y) \otimes X \to Y$,
\item $hom(Y, Z) \otimes hom(X, Y) \to hom(X, Z)$,
\item $Z \otimes hom(X, Y) \to hom(X, Z \otimes Y)$,
\item $Y \to hom(X, Y \otimes X)$, et
\item $hom(Y, Z) \otimes X \to hom(hom(X, Y), Z)$ lorsque
$\mathcal{C}$ est monoïdale symétrique.
\end{enumerate}
En effet, le morphisme de (1) est l'image de $\text{id}_{hom(X, Y)}$
par $\Mor(hom(X, Y), hom(X, Y)) \to \Mor(hom(X, Y) \otimes X, Y)$.
Pour construire le morphisme de (2), la propriété caractéristique de
$hom(X, Z)$ ramène à construire un morphisme
$$
hom(Y, Z) \otimes hom(X, Y) \otimes X \longrightarrow Z
$$
et un tel morphisme existe puisque, d'après (1), on dispose des
morphismes $hom(X, Y) \otimes X \to Y$ et $hom(Y, Z) \otimes Y \to Z$.
Pour construire le morphisme de (3), la propriété caractéristique de
$hom(X, Z \otimes Y)$ ramène à construire un morphisme
$$
Z \otimes hom(X, Y) \otimes X \to Z \otimes Y
$$
pour laquelle on utilise $\text{id}_Z \otimes a$, où
$a$ est le morphisme de (1). Pour construire le morphisme de (4),
notons que nous disposons déjà du morphisme
$Y \otimes hom(X, X) \to hom(X, Y \otimes X)$ donnée par (3).
Il suffit donc de construire un morphisme $\mathbf{1} \to hom(X, X)$ ;
pour cela, on prend l'élément de $\Mor(\mathbf{1}, hom(X, X))$
qui correspond à l'isomorphisme canonique $\mathbf{1} \otimes X \to X$
dans $\Mor(\mathbf{1} \otimes X, X)$.
Enfin, considérons (5). Par la propriété universelle de
$hom(hom(X, Y), Z)$, il suffit de construire un morphisme
$$
hom(Y, Z) \otimes X \otimes hom(X, Y) \longrightarrow Z
$$
Pour ce faire, on permute les deux derniers facteurs tensoriels au moyen de la
contrainte de commutativité, puis on utilise les morphismes
$hom(X, Y) \otimes X \to Y$ et $hom(Y, Z) \otimes Y \to Z$.
\end{remark}





\section{Catégories de flèches pointillées}
\label{categories-section-dotted-arrows}

\noindent
Nous étudions certaines « catégories de flèches pointillées » dans les
$(2,1)$-catégories. Elles apparaîtront dans la formulation de divers critères de
relèvement pour les champs algébriques ; voir, par exemple, Morphismes de champs,
section \ref{stacks-morphisms-section-valuative} et
Compléments sur les morphismes de champs, section
\ref{stacks-more-morphisms-section-formally-smooth}.

\begin{definition}
\label{categories-definition-dotted-arrows}
Soit $\mathcal{C}$ une $(2,1)$-catégorie. Considérons un diagramme
$2$-commutatif formé par les flèches pleines
\begin{equation}
\label{categories-equation-dotted-arrows}
\vcenter{
\xymatrix{
S \ar[r]_-x \ar[d]_j & X \ar[d]^f \\
T \ar[r]^-y \ar@{..>}[ru] & Y
}
}
\end{equation}
dans $\mathcal{C}$. Fixons un $2$-isomorphisme
$$
\gamma : y \circ j \rightarrow f \circ x
$$
attestant la $2$-commutativité du diagramme.
Étant donnés (\ref{categories-equation-dotted-arrows}) et $\gamma$, une
\emph{flèche pointillée} est un triplet $(a, \alpha, \beta)$ constitué d'un
morphisme $a \colon T \to X$ et de $2$-isomorphismes
$\alpha : a \circ j \to x$, $\beta : y \to f \circ a$
tels que
$\gamma = (\text{id}_f \star \alpha) \circ (\beta \star \text{id}_j)$,
autrement dit tels que
$$
\xymatrix{
& f \circ a \circ j \ar[rd]^{\text{id}_f \star \alpha} \\
y \circ j \ar[ru]^{\beta \star \text{id}_j} \ar[rr]^\gamma & &
f \circ x
}
$$
soit commutatif. Un {\it morphisme de flèches pointillées}
$(a, \alpha, \beta) \to (a', \alpha', \beta')$ est une
$2$-flèche $\theta : a \to a'$ telle que
$\alpha = \alpha' \circ (\theta \star \text{id}_j)$ et
$\beta' = (\text{id}_f \star \theta) \circ \beta$.
\end{definition}

\noindent
Dans la situation de la définition \ref{categories-definition-dotted-arrows}, il existe
une \emph{catégorie de flèches pointillées} associée.
Cette catégorie est un groupoïde. Elle peut en général dépendre de $\gamma$.
Les deux lemmes suivants affirment que les catégories de flèches pointillées
sont compatibles avec le changement de base et avec la composition du
morphisme $f$.

\begin{lemma}
\label{categories-lemma-cat-dotted-arrows-base-change}
Soit $\mathcal{C}$ une $(2,1)$-catégorie. Supposons donné un diagramme
$2$-commutatif
$$
\xymatrix{
S \ar[r]_-{x'} \ar[d]_j &
X' \ar[d]^p \ar[r]_q &
X \ar[d]^f \\
T \ar[r]^-{y'} &
Y' \ar[r]^g &
Y
}
$$
dans $\mathcal{C}$, où le carré de droite est $2$-cartésien
relativement à un $2$-isomorphisme $\phi \colon g \circ p \to f \circ q$.
Choisissons une $2$-flèche
$\gamma' : y' \circ j \to p \circ x'$. Posons
$x = q \circ x'$, $y = g \circ y'$ et soit
$\gamma : y \circ j \to f \circ x$ le $2$-isomorphisme
$\gamma = (\phi \star \text{id}_{x'}) \circ (\text{id}_g \star \gamma')$.
Alors la catégorie $\mathcal{D}'$ des flèches pointillées
pour le carré de gauche et $\gamma'$ est équivalente à la catégorie
$\mathcal{D}$ des flèches pointillées
pour le rectangle extérieur et $\gamma$.
\end{lemma}

\begin{proof}
Il existe un foncteur $\mathcal{D}' \to \mathcal{D}$ qui est donné sur les
objets par $(a, \alpha, \beta) \mapsto (q \circ a, \text{id}_q \star \alpha,
(\phi \star \text{id}_a) \circ (\text{id}_g \star \beta))$
et sur les flèches par $\theta \mapsto \text{id}_q \star \theta$.
Le fait que ce foncteur $\mathcal{D}' \to \mathcal{D}$ soit une équivalence
découle formellement de la propriété universelle des $2$-produits fibrés,
comme dans la section \ref{categories-section-2-fibre-products}. Les détails sont omis.
\end{proof}

\begin{lemma}
\label{categories-lemma-cat-dotted-arrows-composition}
Soit $\mathcal{C}$ une $(2,1)$-catégorie. Supposons donné un diagramme
$2$-commutatif aux flèches pleines
$$
\xymatrix{
S \ar[r]_-x \ar[dd]_j & X \ar[d]^f \\
& Y \ar[d]^g \\
T \ar[r]^-z \ar@{..>}[ruu] & Z
}
$$
dans $\mathcal{C}$.
Choisissons un $2$-isomorphisme
$\gamma \colon z \circ j \to g \circ f \circ x$.
Soit $\mathcal{D}$ la catégorie des flèches pointillées pour
le rectangle extérieur et $\gamma$. Soit $\mathcal{D}'$
la catégorie des flèches pointillées pour le carré aux flèches pleines
$$
\xymatrix{
S \ar[r]_-{f \circ x} \ar[d]_j & Y \ar[d]^g \\
T \ar[r]^-z \ar@{..>}[ru] & Z
}
$$
et $\gamma$. Alors $\mathcal{D}$ est équivalente à
une catégorie $\mathcal{D}''$ qui possède la propriété suivante :
il existe un foncteur $\mathcal{D}'' \to \mathcal{D}'$ qui fait de
$\mathcal{D}''$ une catégorie fibrée en groupoïdes au-dessus de
$\mathcal{D}'$ et dont les catégories fibres sont isomorphes à des catégories
de flèches pointillées pour certains carrés aux flèches pleines de la forme
$$
\xymatrix{
S \ar[r]_-x \ar[d]_j & X \ar[d]^f \\
T \ar[r]^-y \ar@{..>}[ru] & Y
}
$$
et certains choix de $2$-isomorphisme $y \circ j \to f \circ x$.
\end{lemma}

\begin{proof}
Construisons la catégorie $\mathcal{D}''$ dont les objets sont les quintuplets
$(a,\alpha,\beta,b,\eta)$,
où $(a,\alpha,\beta)$ est un objet de $\mathcal{D}$,
$b \colon T \rightarrow Y$ est un $1$-morphisme
et $\eta \colon b \rightarrow f \circ a$ est un $2$-isomorphisme.
Les morphismes
$(a,\alpha,\beta,b,\eta) \rightarrow (a',\alpha',\beta',b',\eta')$
dans $\mathcal{D}''$
sont les paires $(\theta_1,\theta_2)$, où
$\theta_1 \colon a \rightarrow a'$ définit une flèche
$(a, \alpha, \beta) \rightarrow (a', \alpha', \beta')$
dans $\mathcal{D}$ et où $\theta_2 \colon b \rightarrow b'$ est un
$2$-isomorphisme satisfaisant à la condition de compatibilité
$\eta' \circ \theta_2 = (\text{id}_f \star \theta_1) \circ \eta$.

\medskip\noindent
Il existe un foncteur $\mathcal{D}'' \rightarrow \mathcal{D}'$ qui est donné
sur les objets par
$(a, \alpha, \beta, b, \eta) \mapsto
(b, (\text{id}_f \star \alpha) \circ (\eta \star \text{id}_j),
(\text{id}_g \star \eta^{-1}) \circ \beta)$
et sur les flèches par $(\theta_1,\theta_2) \mapsto \theta_2$.
Alors $\mathcal{D}'' \rightarrow \mathcal{D}'$ est fibrée en groupoïdes.

\medskip\noindent
Si $(y, \delta, \epsilon)$ est un objet de $\mathcal{D}'$, notons
$\mathcal{D}_{y,\delta}$
la catégorie des flèches pointillées pour le dernier diagramme affiché, avec
$y \circ j \rightarrow f \circ x$ donné par $\delta$.
Il existe un foncteur $\mathcal{D}_{y,\delta} \rightarrow \mathcal{D}''$ donné
sur les objets par
$(a, \alpha, \eta) \mapsto
(a, \alpha, (\text{id}_g \star \eta) \circ \epsilon, y, \eta)$
et sur les flèches par $\theta \mapsto (\theta, \text{id}_y)$.
Ce foncteur donne un isomorphisme de $\mathcal{D}_{y,\delta}$ vers
la catégorie fibre de $\mathcal{D}'' \rightarrow \mathcal{D}'$ au-dessus de
$(y,\delta,\epsilon)$.

\medskip\noindent
Il existe aussi un foncteur $\mathcal{D} \rightarrow \mathcal{D}''$ qui est
donné sur les objets par
$(a,\alpha,\beta) \mapsto (a,\alpha,\beta,f \circ a, \text{id}_{f \circ a})$
et sur les flèches par
$\theta \mapsto (\theta, \text{id}_f \star \theta)$.
Ce foncteur est pleinement fidèle et essentiellement surjectif ; c'est donc une
équivalence. Les détails sont omis.
\end{proof}












% OK, start here.
%

\chapter{Topologie}



\phantomsection
\label{topology-section-phantom}


\section{Introduction}
\label{topology-section-introduction}

\noindent
Nous exposerons dans ce chapitre les éléments de topologie générale dont nous
aurons besoin. On pourra consulter \cite{Engelking} comme ouvrage de référence.

\section{Notions de base}
\label{topology-section-topology-basic}

\noindent
Voici une liste de notions fondamentales de topologie. Certaines seront
étudiées plus en détail dans la suite, certaines sont définies dans la liste,
mais les autres sont tenues pour acquises et ne seront pas définies. Si la
plupart des notions en italique ne vous sont pas familières, nous vous
conseillons de consulter un ouvrage d'introduction à la topologie avant de
poursuivre.

\begin{enumerate}
\item
\label{topology-item-space}
$X$ est un {\it espace topologique},
\item
\label{topology-item-point}
$x\in X$ est un {\it point},
\item
\label{topology-item-locally-closed}
$E \subset X$ est une partie {\it localement fermée},
\item
\label{topology-item-closed-point}
$x\in X$ est un {\it point fermé},
\item
\label{topology-item-dense}
$E \subset X$ est une partie {\it dense},
\item
\label{topology-item-continuous}
$f : X_1 \to X_2$ est {\it continue},
\item
\label{topology-item-upper-semi-continuous}
une fonction à valeurs réelles étendues
$f : X \to \mathbf{R} \cup \{\infty, -\infty\}$
est {\it semi-continue supérieurement} si $\{x \in X \mid f(x) < a\}$ est
ouvert pour tout $a \in \mathbf{R}$,
\item
\label{topology-item-lower-semi-continuous}
une fonction à valeurs réelles étendues
$f : X \to \mathbf{R} \cup \{\infty, -\infty\}$
est {\it semi-continue inférieurement} si $\{x \in X \mid f(x) > a\}$ est
ouvert pour tout $a \in \mathbf{R}$,
\item une application continue d'espaces $f : X \to Y$ est {\it ouverte} si
$f(U)$ est ouvert dans $Y$ pour toute partie ouverte $U \subset X$,
\item une application continue d'espaces $f : X \to Y$ est {\it fermée} si
$f(Z)$ est fermé dans $Y$ pour toute partie fermée $Z \subset X$,
\item
\label{topology-item-neighbourhood}
un {\it voisinage de $x \in X$} est toute partie $E \subset X$ qui contient un
ouvert contenant $x$,
\item
\label{topology-item-induced-topology}
la {\it topologie induite} sur une partie $E \subset X$,
\item
\label{topology-item-covering}
$\mathcal{U} : U = \bigcup_{i \in I} U_i$ est un {\it recouvrement ouvert de}
$U$ (remarquons que chaque $U_i$ peut être vide et que, lorsque $U$ est vide,
nous admettons même que $I$ soit l'ensemble vide),
\item
\label{topology-item-subcover}
un {\it sous-recouvrement} d'un recouvrement comme en
(\ref{topology-item-covering}) est un recouvrement ouvert
$\mathcal{U}' : U = \bigcup_{i \in I'} U_i$, où $I' \subset I$,
\item
\label{topology-item-refinement}
le recouvrement ouvert $\mathcal{V}$ est un {\it raffinement} du recouvrement
ouvert $\mathcal{U}$ (si $\mathcal{V} : U = \bigcup_{j \in J} V_j$ et
$\mathcal{U} : U = \bigcup_{i \in I} U_i$, cela signifie que chaque $V_j$ est
contenu dans l'un des $U_i$),
\item
\label{topology-item-fundamental-system}
{\it $\{ E_i \}_{i \in I}$ est un système fondamental de voisinages de $x$
dans $X$},
\item
\label{topology-item-Hausdorff}
un espace topologique $X$ est dit {\it séparé}, ou {\it de Hausdorff}, si et
seulement si, pour toute paire de points distincts $x, y \in X$, il existe des
ouverts disjoints $U, V \subset X$ tels que $x \in U$, $y \in V$,
\item le {\it produit} de deux espaces topologiques,
\label{topology-item-product}
\item
\label{topology-item-fibre-product}
le {\it produit fibré $X \times_Y Z$} de deux applications continues
$f : X \to Y$ et $g : Z \to Y$,
\item
\label{topology-item-discrete-indiscrete}
la {\it topologie discrète} et la {\it topologie grossière} sur un ensemble,
\item etc.
\end{enumerate}



\section{Espaces séparés}
\label{topology-section-Hausdorff}

\noindent
La catégorie des espaces topologiques admet les produits finis.

\begin{lemma}
\label{topology-lemma-Hausdorff}
Soit $X$ un espace topologique. Les assertions suivantes sont équivalentes :
\begin{enumerate}
\item $X$ est séparé,
\item la diagonale $\Delta(X) \subset X \times X$ est fermée.
\end{enumerate}
\end{lemma}

\begin{proof}
Supposons que $X$ soit séparé. Soit $(x, y) \not \in \Delta(X)$,
c'est-à-dire $x \neq y$. Il existe deux ouverts disjoints $U$ et $V$ de $X$
tels que $x \in U$ et $y \in V$. Il s'ensuit que $U \times V \subset
(X \times X) \setminus \Delta(X) $. Cela montre que $ (X \times X) \setminus
\Delta(X)$ est un ouvert de $ X \times X$, ce qui revient à dire que
la diagonale $\Delta(X) \subset X \times X$ est fermée dans $X \times X$.
La réciproque est analogue :
le complémentaire $(X \times X) \setminus \Delta(X)$ est formé
précisément des $(x, y) \in X\times X$ tels que $ x \neq y$. Ainsi, si
$ \Delta(X) \subset X \times X$ est fermée, alors, par définition de la
topologie produit, pour tout tel couple $ (x, y)$, il existe des ouverts
$U, V \subset X$ tels que $(x, y) \in U \times V$ et
$(U \times V)\cap \Delta(X) = \emptyset$. Autrement dit, $x \in U$ et
$y \in V$, avec $U \cap V = \emptyset$.
\end{proof}

\begin{lemma}
\label{topology-lemma-graph-closed}
\begin{slogan}
Les graphes des applications à valeurs dans un espace séparé sont fermés.
\end{slogan}
Soit $f : X \to Y$ une application continue d'espaces topologiques.
Si $Y$ est séparé, alors le graphe de $f$ est fermé dans $X \times Y$.
\end{lemma}

\begin{proof}
Le graphe est l'image réciproque de la diagonale par l'application
$X \times Y \to Y \times Y$. Le lemme résulte donc du
lemme \ref{topology-lemma-Hausdorff}.
\end{proof}

\begin{lemma}
\label{topology-lemma-section-closed}
Soit $f : X \to Y$ une application continue d'espaces topologiques.
Soit $s : Y \to X$ une application continue telle que
$f \circ s = \text{id}_Y$.
Si $X$ est séparé, alors $s(Y)$ est fermé.
\end{lemma}

\begin{proof}
Cela résulte du lemme \ref{topology-lemma-Hausdorff}, puisque
$s(Y) = \{x \in X \mid x = s(f(x))\}$.
\end{proof}

\begin{lemma}
\label{topology-lemma-fibre-product-closed}
Soient $X \to Z$ et $Y \to Z$ des applications continues d'espaces
topologiques. Si $Z$ est séparé, alors $X \times_Z Y$ est fermé dans
$X \times Y$.
\end{lemma}

\begin{proof}
Cela résulte du lemme \ref{topology-lemma-Hausdorff}, car $X \times_Z Y$ est l'image
réciproque de $\Delta(Z)$ par $X \times Y \to Z \times Z$.
\end{proof}




\section{Applications séparées}
\label{topology-section-separated}

\noindent
Voici simplement la définition et quelques lemmes élémentaires.

\begin{definition}
\label{topology-definition-separated}
Une application continue $f : X \to Y$ d'espaces topologiques est dite
{\it séparée} si et seulement si la diagonale $\Delta : X \to X \times_Y X$
est une application fermée.
\end{definition}

\begin{lemma}
\label{topology-lemma-separated}
Soit $f : X \to Y$ une application continue d'espaces topologiques.
Les assertions suivantes sont équivalentes :
\begin{enumerate}
\item $f$ est séparée,
\item $\Delta(X) \subset X \times_Y X$ est une partie fermée,
\item étant donnés deux points distincts $x, x' \in X$ ayant la même image
dans $Y$, il existe des voisinages ouverts disjoints de $x$ et de $x'$.
\end{enumerate}
\end{lemma}

\begin{proof}
Si $f$ est séparée, la définition \ref{topology-definition-separated} dit que $\Delta$
est une application fermée. Le fait que $X$ soit fermé dans $X$ entraîne que
$\Delta(X)$ est fermé dans $ X\times_Y X $. Ainsi, (1) implique (2).

\medskip\noindent
Supposons que $\Delta(X) \subset X \times_Y X$ soit une partie fermée et
notons $U$ son complémentaire ouvert. Il existe donc un ouvert
$W \subset X \times X$ tel que $W \cap (X \times_Y X) = U$. Or, par
définition de la topologie produit, si
$(x, x') \in W \cap (X \times_Y X)$,
il existe des ouverts $V$ et $V'$ de $X$ tels que $x \in V$, $x' \in V'$ et
$V \times V' \subset W$. Si $V \cap V' \neq \emptyset$, il existerait
$z \in V \cap V'$. Mais $(z, z) \in X \times_Y X$, donc
$(z,z) \in (V \times V') \cap (X\times_Y X) \subset U$, ce qui est absurde.
Ainsi $V \cap V' = \emptyset $, et nous avons obtenu deux voisinages ouverts
disjoints de $x$ et $x'$. Cela montre que (2) implique (3).

\medskip\noindent
Supposons enfin que, pour tous points distincts $x, x' \in X$ ayant la même
image dans $Y$, il existe des voisinages ouverts disjoints de $x$ et de $x'$.
Soit $F$ un fermé de $X$. Montrons que $\Delta(F)$ est une partie fermée de
$X \times_Y X$. Soit $(x, x') \in X \times_Y X$ un point qui n'appartient
pas à $\Delta(F)$. Alors soit $x \not = x'$, soit $x \not \in F$.
Dans le premier cas, choisissons des voisinages ouverts disjoints
$V, V' \subset X$ de $x, x'$ ; alors $(V \times V') \cap X \times_Y X$
est un voisinage ouvert de $(x, x')$ qui ne rencontre pas $\Delta(F)$.
Dans le second cas, $((X \setminus F) \times X) \cap X \times_Y X$
est un voisinage ouvert de $(x, x')$ qui ne rencontre pas $\Delta(F)$.
Nous avons montré que (3) implique (1).
\end{proof}

\begin{lemma}
\label{topology-lemma-from-hausdorff}
Soit $f : X \to Y$ une application continue d'espaces topologiques.
Si $X$ est séparé, alors $f$ est séparée.
\end{lemma}

\begin{proof}
C'est immédiat d'après les lemmes \ref{topology-lemma-separated} et
\ref{topology-lemma-Hausdorff}, car le fait que $\Delta(X)$ soit fermé dans
$X \times X$ entraîne que $\Delta(X)$ est fermé dans $X \times_Y X$.
\end{proof}

\begin{lemma}
\label{topology-lemma-base-change-separated}
Soient $f : X \to Y$ et $Y' \to Y$ des applications continues d'espaces
topologiques. Si $f$ est séparée, alors
$f' : Y' \times_Y X \to Y'$ est séparée.
\end{lemma}

\begin{proof}
Cela résulte de la caractérisation (2) du lemme \ref{topology-lemma-separated}.
En effet, si $X' = Y' \times_Y X$, la diagonale $\Delta(X')$ dans le produit
fibré $X' \times_{Y'} X'$ est l'image réciproque de $\Delta(X)$ dans
$X \times_Y X$.
\end{proof}




\section{Bases}
\label{topology-section-bases}

\noindent
Rappels élémentaires sur les bases des espaces topologiques.

\begin{definition}
\label{topology-definition-base}
Soit $X$ un espace topologique. Une famille de parties $\mathcal{B}$ de $X$
est appelée une {\it base de la topologie sur $X$}, ou une {\it base de la
topologie de $X$}, si les conditions suivantes sont satisfaites :
\begin{enumerate}
\item Tout élément $B \in \mathcal{B}$ est ouvert dans $X$.
\item Pour tout ouvert $U \subset X$ et tout $x \in U$, il existe
$B \in \mathcal{B}$ tel que $x \in B \subset U$.
\end{enumerate}
\end{definition}

\noindent
Le lemme suivant sert parfois à définir une topologie.

\begin{lemma}
\label{topology-lemma-make-base}
Soit $X$ un ensemble et soit $\mathcal{B}$ une famille de parties.
Supposons que $X = \bigcup_{B \in \mathcal{B}} B$ et que, pour tous
$x \in B_1 \cap B_2$ avec $B_1, B_2 \in \mathcal{B}$, il existe
$B_3 \in \mathcal{B}$ tel que $x \in B_3 \subset B_1 \cap B_2$.
Il existe alors une unique topologie sur $X$ dont $\mathcal{B}$ est une base.
\end{lemma}

\begin{proof}
Notons $\sigma(\mathcal B)$ l'ensemble des parties de $X$ qui s'écrivent
comme réunions d'éléments de $\mathcal B$. Nous affirmons que
$\sigma(\mathcal{B})$ est une topologie. En effet, l'ensemble vide appartient
à $\sigma(\mathcal{B})$ (en tant que réunion vide), et $X$ appartient à
$\sigma(\mathcal{B})$ (en tant que réunion de tous les éléments de
$\mathcal{B}$). Il est clair que $\sigma(\mathcal{B})$ est stable par réunion.
Enfin, si $U, V \in \sigma(\mathcal{B})$, écrivons
$U = \bigcup_{i \in I} U_i$ et $V = \bigcup_{j \in J} V_j$, où
$U_i, V_j \in \mathcal{B}$ pour tout $i \in I$ et tout $j \in J$.
Alors
$$
U \cap V = \bigcup\nolimits_{i \in I, j \in J}
U_i \cap V_j
$$
L'hypothèse du lemme montre que $U_i \cap V_j \in \sigma(\mathcal{B})$,
et donc que $U \cap V$ y appartient également. Ainsi,
$\sigma(\mathcal{B})$ est une topologie.
Les propriétés (1) et (2) de la définition \ref{topology-definition-base} sont
immédiates pour cette topologie. Pour en montrer l'unicité, soit $\mathcal T$
une topologie sur $X$ dont $\mathcal B$ est une base. Tout élément de
$\mathcal{B}$ appartient évidemment à $\mathcal{T}$ d'après le point (1) de
la définition \ref{topology-definition-base}, et donc
$\sigma(\mathcal{B}) \subset \mathcal{T}$. Réciproquement, le point (2) de
la définition \ref{topology-definition-base} montre que tout élément de $\mathcal{T}$
est une réunion d'éléments de $\mathcal{B}$, c'est-à-dire que
$\mathcal{T} \subset \sigma(\mathcal{B})$. Cela achève la démonstration.
\end{proof}

\begin{lemma}
\label{topology-lemma-refine-covering-basis}
Soit $X$ un espace topologique.
Soit $\mathcal{B}$ une base de la topologie de $X$.
Soit $\mathcal{U} : U = \bigcup_i U_i$ un recouvrement ouvert de
$U \subset X$. Il existe un recouvrement ouvert $U = \bigcup V_j$ qui
raffine $\mathcal{U}$ et tel que chaque $V_j$ appartienne à la base
$\mathcal{B}$.
\end{lemma}

\begin{proof}
Si $ x \in U = \bigcup_{i\in I} U_i $, il existe $ i_x \in I $ tel que
$ x \in U_{i_x} $. Il existe donc $ B_{i_x} \in \mathcal{B}$ vérifiant
$ x \in B_{i_x} \subset U_{i_x}$. Posons
$J = \{i_x | x \in U\}$ et, pour $j = i_x \in J$, posons
$V_j = B_{i_x}$. On obtient ainsi le recouvrement ouvert voulu de $U$ par
$\{V_j\}_{j \in J}$.
\end{proof}

\begin{definition}
\label{topology-definition-subbase}
Soit $X$ un espace topologique. Une famille de parties $\mathcal{B}$ de $X$
est appelée une {\it sous-base de la topologie sur $X$}, ou une {\it sous-base
de la topologie de $X$}, si les intersections finies d'éléments de
$\mathcal{B}$ forment une base de la topologie de $X$.
\end{definition}

\noindent

En particulier, tout élément de $\mathcal{B}$ est ouvert.

\begin{lemma}
\label{topology-lemma-subbase}
Soit $X$ un ensemble. Pour toute famille $\mathcal{B}$ de parties de $X$,
il existe une unique topologie sur $X$ dont $\mathcal{B}$ est une sous-base.
\end{lemma}

\begin{proof}
Par convention, $\bigcap_\emptyset B = X $. Nous pouvons donc appliquer le
lemme \ref{topology-lemma-make-base} à l'ensemble des intersections finies d'éléments
de $\mathcal B$.
\end{proof}

\begin{lemma}
\label{topology-lemma-create-map-from-subcollection}
Soit $X$ un espace topologique. Soit $\mathcal{B}$ une famille d'ouverts de
$X$. Supposons que $X = \bigcup_{U \in \mathcal{B}} U$ et que, pour
$U, V \in \mathcal{B}$, on ait
$U \cap V = \bigcup_{W \in \mathcal{B}, W \subset U \cap V} W$.
Il existe alors une application continue $f : X \to Y$ d'espaces
topologiques telle que
\begin{enumerate}
\item pour $U \in \mathcal{B}$, l'image $f(U)$ soit ouverte,
\item pour $U \in \mathcal{B}$, on ait $f^{-1}(f(U)) = U$, et
\item les ouverts $f(U)$, $U \in \mathcal{B}$, forment une base de la
topologie de $Y$.
\end{enumerate}
\end{lemma}

\begin{proof}
Définissons une relation d'équivalence $\sim$ sur les points de $X$ par
$$
x \sim y \Leftrightarrow
(\forall U \in \mathcal{B} : x \in U \Leftrightarrow y \in U)
$$
Soit $Y$ l'ensemble des classes d'équivalence et soit $f : X \to Y$
l'application naturelle. Le point (2) est satisfait par construction.
Les hypothèses portant sur $\mathcal{B}$ sont exactement celles du
lemme \ref{topology-lemma-make-base} pour l'ensemble des parties $f(U)$,
$U \in \mathcal{B}$. Il existe donc une unique topologie sur $Y$ pour
laquelle (3) soit satisfaite. Le point (1) est alors clair lui aussi.
\end{proof}






\section{Applications submersives}
\label{topology-section-submersive}

\noindent
Si $X$ est un espace topologique et si $E \subset X$ est une partie, nous
munissons habituellement $E$ de la {\it topologie induite}.

\begin{lemma}
\label{topology-lemma-induced}
Soit $X$ un espace topologique. Soit $Y$ un ensemble et soit
$f : Y \to X$ une application injective d'ensembles. La topologie induite sur
$Y$ est la topologie caractérisée par chacune des assertions suivantes :
\begin{enumerate}
\item c'est la topologie la moins fine sur $Y$ pour laquelle $f$ soit continue,
\item les ouverts de $Y$ sont les $f^{-1}(U)$, où $U \subset X$ est ouvert,
\item les fermés de $Y$ sont les ensembles $f^{-1}(Z)$, où $Z \subset X$
est fermé.
\end{enumerate}
\end{lemma}

\begin{proof}
L'ensemble $\mathcal{T} = \{f^{-1}(U) | U \subset X \text{ ouvert}\}$ est
une topologie sur $Y$. Tout d'abord,
$\emptyset = f^{-1}(\emptyset)$ et $f^{-1}(X) = Y$.
Ainsi, $\mathcal{T}$ contient $\emptyset$ et $Y$.

\medskip\noindent
Soit maintenant $\{V_i\}_{i \in I}$ une famille de parties ouvertes avec
$V_i \in \mathcal{T}$, et écrivons $V_i = f^{-1}(U_i)$, où $U_i$ est une
partie ouverte de $X$. Alors
$$
\bigcup\nolimits_{i\in I} V_i =
\bigcup\nolimits_{i\in I} f^{-1}(U_i) =
f^{-1}\left(\bigcup\nolimits_{i\in I} U_i\right)
$$
Ainsi, $\bigcup_{i\in I} V_i \in \mathcal{T}$, puisque
$\bigcup_{i\in I} U_i$ est ouvert dans $X$.
Soient maintenant $V_1, V_2 \in \mathcal{T}$. Il existe des ouverts
$U_1, U_2$ de $X$ tels que $V_1 = f^{-1}(U_1)$ et
$V_2 = f^{-1}(U_2)$. Alors
$$
V_1 \cap V_2 = f^{-1}(U_1) \cap f^{-1}(U_2) = f^{-1}(U_1 \cap U_2)
$$
Ainsi, $V_1 \cap V_2 \in \mathcal{T}$, car $U_1 \cap U_2$ est ouvert dans
$X$.

\medskip\noindent
Toute topologie sur $Y$ pour laquelle $f$ est continue contient
$\mathcal{T}$, par définition d'une application continue. Ainsi,
$\mathcal{T}$ est bien la topologie la moins fine sur $Y$ pour laquelle $f$
soit continue. Cela prouve l'équivalence de (1) et (2).

\medskip\noindent
L'équivalence de (2) et (3) résulte de l'égalité
$f^{-1}(X \setminus E) = Y \setminus f^{-1}(E)$ pour toute partie
$E \subset X$.
\end{proof}

\noindent
Duellement, si $X$ est un espace topologique et si $X \to Y$ est une
surjection d'ensembles, on peut munir $Y$ de la {\it topologie quotient}.

\begin{lemma}
\label{topology-lemma-quotient}
Soit $X$ un espace topologique. Soit $Y$ un ensemble et soit $f : X \to Y$
une application surjective d'ensembles. La topologie quotient sur $Y$ est la
topologie caractérisée par chacune des assertions suivantes :
\begin{enumerate}
\item c'est la topologie la plus fine sur $Y$ pour laquelle $f$ soit continue,
\item une partie $V$ de $Y$ est ouverte si et seulement si $f^{-1}(V)$ est
ouverte,
\item une partie $Z$ de $Y$ est fermée si et seulement si $f^{-1}(Z)$ est
fermée.
\end{enumerate}
\end{lemma}

\begin{proof}
L'ensemble $\mathcal{T} = \{V\subset Y | f^{-1}(V) \text{ est ouvert}\}$ est
une topologie sur $Y$. Tout d'abord,
$\emptyset = f^{-1}(\emptyset)$ et $f^{-1}(Y) = X$.
Ainsi, $\mathcal{T}$ contient $\emptyset$ et $Y$.

\medskip\noindent
Soit $(V_i)_{i \in I}$ une famille d'éléments $V_i \in \mathcal{T}$. Alors
$$
\bigcup\nolimits_{i\in I} f^{-1}(V_i) =
f^{-1}\left(\bigcup\nolimits_{i\in I} V_i\right)
$$
Ainsi, $\bigcup_{i\in I} V_i \in \mathcal{T}$, puisque
$\bigcup_{i\in I}f^{-1}(V_i)$ est ouvert dans $X$.
De plus, si $V_1, V_2 \in \mathcal{T}$, alors
$$
f^{-1}(V_1) \cap f^{-1}(V_2) = f^{-1}(V_1 \cap V_2)
$$
Ainsi, $V_1 \cap V_2 \in \mathcal{T}$, car
$f^{-1}(V_1) \cap f^{-1}(V_2)$ est ouvert dans $X$.

\medskip\noindent
Enfin, toute topologie sur $Y$ pour laquelle $f$ est continue est contenue
dans $\mathcal{T}$, par définition d'une application continue. Par conséquent,
$\mathcal{T}$ est la topologie la plus fine sur $Y$ pour laquelle $f$ soit
continue. Cela prouve l'équivalence de (1) et (2).

\medskip\noindent
Enfin, l'équivalence de (2) et (3) résulte de
$f^{-1}(X\setminus E) = Y \setminus f^{-1}(E)$ pour toute partie
$E \subset X$.
\end{proof}

\noindent
Soit $f : X \to Y$ une application continue d'espaces topologiques.
Nous obtenons alors une factorisation $X \to f(X) \to Y$ d'applications
d'ensembles. Nous pouvons munir $f(X)$ soit de la topologie quotient provenant
de la surjection $X \to f(X)$, soit de la topologie induite provenant de
l'injection $f(X) \to Y$. L'application
$$
(f(X), \text{topologie quotient})
\longrightarrow
(f(X), \text{topologie induite})
$$
est continue.

\begin{definition}
\label{topology-definition-submersive}
Soit $f : X \to Y$ une application continue d'espaces topologiques.
\begin{enumerate}
\item On dit que $f$ est une {\it application stricte d'espaces topologiques}
si la topologie induite et la topologie quotient sur $f(X)$ coïncident
(voir la discussion ci-dessus).
\item On dit que $f$ est {\it submersive}\footnote{Cette notion est très
différente de celle de submersion entre variétés différentielles ! Il est sans
doute préférable d'employer « stricte et surjective » plutôt que
« submersive ».} si $f$ est surjective et stricte.
\end{enumerate}
\end{definition}

\noindent
Ainsi, une application continue $f : X \to Y$ est submersive si $f$ est une
surjection et si, pour toute partie $T \subset Y$, la partie $T$ est ouverte
(respectivement fermée) si et seulement si $f^{-1}(T)$ est ouverte
(respectivement fermée). Autrement dit, $Y$ porte la
topologie quotient relative à la surjection $X \to Y$.

\begin{lemma}
\label{topology-lemma-open-morphism-quotient-topology}
Soit $f : X \to Y$ une application surjective, ouverte et continue d'espaces
topologiques. Soit $T \subset Y$ une partie. Alors
\begin{enumerate}
\item $f^{-1}(\overline{T}) = \overline{f^{-1}(T)}$,
\item $T \subset Y$ est fermé si et seulement si $f^{-1}(T)$ est fermé,
\item $T \subset Y$ est ouvert si et seulement si $f^{-1}(T)$ est ouvert, et
\item $T \subset Y$ est localement fermé si et seulement si $f^{-1}(T)$
est localement fermé.
\end{enumerate}
En particulier, $f$ est submersive.
\end{lemma}

\begin{proof}
Il est clair que $\overline{f^{-1}(T)} \subset f^{-1}(\overline{T})$.
Si $x \in X$ et $x \not \in \overline{f^{-1}(T)}$, il existe un voisinage
ouvert $x \in U \subset X$ tel que $U \cap f^{-1}(T) = \emptyset$.
Comme $f$ est ouverte, $f(U)$ est un voisinage ouvert de $f(x)$ qui ne
rencontre pas $T$. Ainsi, $x \not \in f^{-1}(\overline{T})$. Cela prouve (1).
Le point (2) résulte immédiatement de (1).
Le point (3) est évident puisque $f$ est ouverte et surjective.
Pour (4), si $f^{-1}(T)$ est localement fermé, alors
$f^{-1}(T) \subset \overline{f^{-1}(T)} = f^{-1}(\overline{T})$
est ouvert. En appliquant alors (3) à l'application
$f^{-1}(\overline{T}) \to \overline{T}$, nous voyons que $T$ est ouvert
dans $\overline{T}$, c'est-à-dire que $T$ est localement fermé.
\end{proof}

\begin{lemma}
\label{topology-lemma-closed-morphism-quotient-topology}
Soit $f : X \to Y$ une application surjective, fermée et continue d'espaces
topologiques. Soit $T \subset Y$ une partie. Alors
\begin{enumerate}
\item $\overline{T} = f(\overline{f^{-1}(T)})$,
\item $T \subset Y$ est fermé si et seulement si $f^{-1}(T)$ est fermé,
\item $T \subset Y$ est ouvert si et seulement si $f^{-1}(T)$ est ouvert, et
\item $T \subset Y$ est localement fermé si et seulement si
$f^{-1}(T)$ est localement fermé.
\end{enumerate}
En particulier, $f$ est submersive.
\end{lemma}

\begin{proof}
Il est clair que $\overline{f^{-1}(T)} \subset f^{-1}(\overline{T})$.
Alors $T \subset f(\overline{f^{-1}(T)}) \subset \overline{T}$ est une
partie fermée, d'où (1). Le point (2) est évident puisque $f$ est fermée et
surjective. Le point (3) résulte de (2) appliqué au complémentaire de $T$.
Pour (4), si $f^{-1}(T)$ est localement fermé, alors
$f^{-1}(T) \subset \overline{f^{-1}(T)}$ est ouvert.
Puisque l'application $\overline{f^{-1}(T)} \to \overline{T}$ est surjective
d'après (1), nous pouvons appliquer le point (3) à l'application
$\overline{f^{-1}(T)} \to \overline{T}$ induite par $f$ et conclure que $T$
est ouvert dans $\overline{T}$, c'est-à-dire que $T$ est localement fermé.
\end{proof}








\section{Composantes connexes}
\label{topology-section-connected-components}

\begin{definition}
\label{topology-definition-connected-components}
Soit $X$ un espace topologique.
\begin{enumerate}
\item On dit que $X$ est {\it connexe} si $X$ n'est pas vide et si, chaque
fois que $X = T_1 \amalg T_2$ avec $T_i \subset X$ ouvert et fermé, on a
$T_1 = \emptyset$ ou $T_2 = \emptyset$.
\item On dit que $T \subset X$ est une {\it composante connexe} de $X$ si
$T$ est une partie connexe maximale de $X$.
\end{enumerate}
\end{definition}

\noindent
L'espace vide n'est pas connexe.

\begin{lemma}
\label{topology-lemma-image-connected-space}
Soit $f : X \to Y$ une application continue d'espaces topologiques.
Si $E \subset X$ est une partie connexe, alors $f(E) \subset Y$ est également
connexe.
\end{lemma}

\begin{proof}
Soit $A \subset f(E)$ une partie ouverte et fermée de $f(E)$. Comme $f$ est
continue, $f^{-1}(A)$ est une partie ouverte et fermée de $E$. Puisque $E$ est
connexe, $f^{-1}(A) = \emptyset$ ou $f^{-1}(A) = E$. Or $A\subset f(E)$
entraîne que $A = f(f^{-1}(A))$. En effet, si $x\in f(f^{-1}(A))$, il existe
$y\in f^{-1}(A)$ tel que $f(y) = x$ et, puisque $y\in f^{-1}(A)$, nous avons
$f(y) \in A$, c'est-à-dire $x\in A$. Réciproquement, si $x\in A$, l'inclusion
$A\subset f(E)$ implique qu'il existe $y\in E$ tel que $f(y) = x$. Alors
$y\in f^{-1}(A)$, puis $x\in f(f^{-1}(A))$. Ainsi, $A = \emptyset$ ou
$A = f(E)$, ce qui prouve que $f(E)$ est connexe.
\end{proof}

\begin{lemma}
\label{topology-lemma-connected-components}
Soit $X$ un espace topologique.
\begin{enumerate}
\item Si $T \subset X$ est connexe, son adhérence l'est aussi.
\item Toute composante connexe de $X$ est fermée (mais pas nécessairement
ouverte).
\item Toute partie connexe de $X$ est contenue dans une unique composante
connexe de $X$.
\item Tout point de $X$ appartient à une unique composante connexe ; autrement
dit, $X$ est la réunion disjointe de ses composantes connexes.
\end{enumerate}
\end{lemma}

\begin{proof}
Soit $\overline{T}$ l'adhérence de la partie connexe $T$.
Supposons $\overline{T} = T_1 \amalg T_2$, où
$T_i \subset \overline{T}$ est ouvert et fermé. Alors
$T = (T\cap T_1) \amalg (T \cap T_2)$. Par conséquent, $T$ est égal à l'une
des deux parties, disons $T = T_1 \cap T$. Ainsi,
$\overline{T} \subset T_1$. Cela entraîne (1) et (2).

\medskip\noindent
Soit $A$ un ensemble non vide de parties connexes de $X$ tel que
$\Omega = \bigcap_{T \in A} T$ soit non vide. Nous affirmons que
$E = \bigcup_{T \in A} T$ est connexe. En effet, $E$ est non vide puisqu'il
contient $\Omega$. Écrivons $E = E_1 \amalg E_2$, où $E_i$ est fermé dans
$E$. Nous pouvons supposer que $E_1$ rencontre $\Omega$ (quitte à renuméroter).
Alors tout $T \in A$ rencontre $E_1$ et doit donc être contenu dans $E_1$,
puisque $T$ est connexe. Ainsi $E \subset E_1$, ce qui prouve l'affirmation.

\medskip\noindent
Soit $W \subset X$ une partie connexe non vide. Appliquons le résultat du
paragraphe précédent à l'ensemble de toutes les parties connexes de $X$ qui
contiennent $W$. Nous voyons alors que $E$ est une composante connexe de $X$.
Cela prouve l'existence et l'unicité dans (3).

\medskip\noindent
Soit $x \in X$. En prenant $W = \{x\}$ dans le paragraphe précédent, nous
voyons que $x$ appartient à une unique composante connexe de $X$.
Deux composantes connexes distinctes sont nécessairement disjointes, d'après
le résultat du deuxième paragraphe.

\medskip\noindent
Pour obtenir un exemple où les composantes connexes ne sont pas ouvertes, il
suffit de considérer un produit infini
$\prod_{n \in \mathbf{N}} \{0, 1\}$ muni de la topologie produit. Ses
composantes connexes sont les singletons, qui ne sont pas ouverts.
\end{proof}

\begin{remark}
\label{topology-remark-quasi-components}
\begin{reference}
\cite[exemple 6.1.24]{Engelking}
\end{reference}
Soit $X$ un espace topologique et soit $x \in X$. Soit $Z \subset X$ la
composante connexe de $X$ passant par $x$. Considérons l'intersection $E$ de
toutes les parties ouvertes et fermées de $X$ qui contiennent $x$. Il est clair
que $Z \subset E$. En général, $Z \not = E$. Par exemple, soit
$X = \{x, y, z_1, z_2, \ldots\}$ muni de la topologie ayant pour base
d'ouverts $\{z_n\}$, $\{x, z_n, z_{n + 1}, \ldots\}$ et
$\{y, z_n, z_{n + 1}, \ldots\}$ pour tout $n$. Alors $Z = \{x\}$ et
$E = \{x, y\}$. Nous omettons les détails.
\end{remark}

\begin{lemma}
\label{topology-lemma-connected-fibres-quotient-topology-connected-components}
Soit $f : X \to Y$ une application continue d'espaces topologiques.
Supposons que
\begin{enumerate}
\item toutes les fibres de $f$ soient connexes, et
\item une partie $T \subset Y$ soit fermée si et seulement si $f^{-1}(T)$
est fermée.
\end{enumerate}
Alors $f$ induit une bijection entre les ensembles de composantes connexes de
$X$ et de $Y$.
\end{lemma}

\begin{proof}
Soit $T \subset Y$ une composante connexe. Remarquons que $T$ est fermée ;
voir le lemme \ref{topology-lemma-connected-components}. Le lemme résulte du fait que
$f^{-1}(T)$ est connexe, car toute partie connexe de $X$ s'envoie dans une
composante connexe de $Y$ d'après le lemme
\ref{topology-lemma-image-connected-space}.
Supposons que $f^{-1}(T) = Z_1 \amalg Z_2$, avec $Z_1$, $Z_2$ fermés.
Pour tout $t \in T$, nous avons
$f^{-1}(\{t\}) = Z_1 \cap f^{-1}(\{t\}) \amalg Z_2 \cap f^{-1}(\{t\})$.
D'après (1), $f^{-1}(\{t\})$ est connexe ; nous en concluons que soit
$f^{-1}(\{t\}) \subset Z_1$, soit $f^{-1}(\{t\}) \subset Z_2$.
Autrement dit, $T = T_1 \amalg T_2$, avec $f^{-1}(T_i) = Z_i$.
D'après (2), $T_i$ est fermé dans $Y$. Ainsi, $T_1 = \emptyset$ ou
$T_2 = \emptyset$, comme voulu.
\end{proof}

\begin{lemma}
\label{topology-lemma-connected-fibres-connected-components}
Soit $f : X \to Y$ une application continue d'espaces topologiques.
Supposons que
(a) $f$ soit ouverte,
(b) toutes les fibres de $f$ soient connexes.
Alors $f$ induit une bijection entre les ensembles de composantes connexes de
$X$ et de $Y$.
\end{lemma}

\begin{proof}
C'est un cas particulier du
lemme \ref{topology-lemma-connected-fibres-quotient-topology-connected-components}.
\end{proof}

\begin{lemma}
\label{topology-lemma-finite-fibre-connected-components}
Soit $f : X \to Y$ une application continue d'espaces topologiques non vides.
Supposons que
(a) $Y$ soit connexe,
(b) $f$ soit ouverte et fermée, et
(c) il existe un point $y\in Y$ tel que la fibre $f^{-1}(y)$ soit un ensemble
fini. Alors $X$ possède au plus $|f^{-1}(y)|$ composantes connexes. Par
conséquent, toute composante connexe $T$ de $X$ est ouverte et fermée, et
$f(T)$ est une partie ouverte et fermée non vide de $Y$, donc égale à $Y$.
\end{lemma}

\begin{proof}
Si l'espace topologique $X$ possède au moins $N$ composantes connexes pour un
certain $N \in \mathbf{N}$, nous obtenons par récurrence une décomposition
$X = X_1 \amalg \ldots \amalg X_N$ de $X$ en réunion disjointe de $N$ parties
ouvertes et fermées non vides $X_1, \ldots , X_N$ de $X$. Puisque $f$ est
ouverte et fermée, chaque $f(X_i)$ est une partie ouverte et fermée non vide de
$Y$, donc égale à $Y$. En particulier, l'intersection
$X_i \cap f^{-1}(y)$ est non vide pour tout $1 \leq i \leq N$. Ainsi,
$f^{-1}(y)$ possède au moins $N$ éléments.
\end{proof}

\begin{definition}
\label{topology-definition-totally-disconnected}
Un espace topologique est {\it totalement discontinu} si toutes ses composantes
connexes sont des singletons.
\end{definition}

\noindent
Un espace discret est totalement discontinu.
Un espace totalement discontinu n'est pas nécessairement discret ; par
exemple, $\mathbf{Q} \subset \mathbf{R}$ est totalement discontinu mais non
discret.

\begin{lemma}
\label{topology-lemma-space-connected-components}
Soit $X$ un espace topologique. Soit $\pi_0(X)$ l'ensemble des composantes
connexes de $X$. Soit $X \to \pi_0(X)$ l'application qui envoie $x \in X$
sur la composante connexe de $X$ passant par $x$. Munissons $\pi_0(X)$ de la
topologie quotient. Alors $\pi_0(X)$ est un espace totalement discontinu, et
toute application continue $X \to Y$ de $X$ vers un espace totalement
discontinu $Y$ se factorise par $\pi_0(X)$.
\end{lemma}

\begin{proof}
D'après le lemme
\ref{topology-lemma-connected-fibres-quotient-topology-connected-components},
les composantes connexes de $\pi_0(X)$ sont les singletons.
Nous omettons la démonstration de la seconde assertion.
\end{proof}

\begin{definition}
\label{topology-definition-locally-connected}
Un espace topologique $X$ est dit {\it localement connexe} si tout point
$x \in X$ possède un système fondamental de voisinages connexes.
\end{definition}

\begin{lemma}
\label{topology-lemma-locally-connected}
Soit $X$ un espace topologique. Si $X$ est localement connexe, alors
\begin{enumerate}
\item toute partie ouverte de $X$ est localement connexe, et
\item les composantes connexes de $X$ sont ouvertes.
\end{enumerate}
Les composantes connexes des parties ouvertes de $X$ sont donc elles aussi
ouvertes. En particulier, tout point possède un système fondamental de
voisinages ouverts connexes.
\end{lemma}

\begin{proof}
Pour tout $x\in X$, notons $\mathcal{N}(x)$ le système fondamental de
voisinages connexes de $x$, et soit $U\subset X$ une partie ouverte de $X$.
Pour tout $x\in U$, $U$ est un voisinage de $x$, donc l'ensemble
$\{V \in \mathcal{N}(x) | V\subset U\}$ n'est pas vide et forme un système
fondamental de voisinages connexes de $x$ dans $U$. Ainsi, $U$ est localement
connexe, ce qui prouve (1).

\medskip\noindent
Soit $x \in \mathcal{C} \subset X$, où $\mathcal{C}$ est la composante
connexe de $x$. Comme $X$ est localement connexe, il existe un voisinage
connexe $\mathcal{N}$ de $x$. Par définition d'une composante connexe, nous
avons donc $\mathcal{N} \subset \mathcal{C}$ ; ainsi, $\mathcal{C}$ est un
voisinage de $x$. Cela montre que $\mathcal{C}$ est un voisinage de chacun de
ses points ; autrement dit, $\mathcal{C}$ est ouverte, ce qui prouve (2).
\end{proof}




\section{Composantes irréductibles}
\label{topology-section-irreducible-components}

\begin{definition}
\label{topology-definition-irreducible-components}
Soit $X$ un espace topologique.
\begin{enumerate}
\item On dit que $X$ est {\it irréductible} si $X$ n'est pas vide et si,
chaque fois que $X = Z_1 \cup Z_2$ avec $Z_i$ fermé, on a $X = Z_1$ ou
$X = Z_2$.
\item On dit que $Z \subset X$ est une {\it composante irréductible} de $X$
si $Z$ est une partie irréductible maximale de $X$.
\end{enumerate}
\end{definition}

\noindent
Un espace irréductible est évidemment connexe.

\begin{lemma}
\label{topology-lemma-image-irreducible-space}
Soit $f : X \to Y$ une application continue d'espaces topologiques.
Si $E \subset X$ est une partie irréductible, alors $f(E) \subset Y$ est
également irréductible.
\end{lemma}

\begin{proof}
Nous pouvons évidemment supposer $E = X$ (c'est-à-dire $X$ irréductible) et
$f(E) = Y$ (c'est-à-dire $f$ surjective). Tout d'abord,
$Y \not = \emptyset$, puisque $X \not = \emptyset$. Supposons ensuite
$Y = Y_1 \cup Y_2$, où $Y_1$, $Y_2$ sont fermés. Alors
$X = X_1 \cup X_2$, où $X_i = f^{-1}(Y_i)$ est fermé dans $X$. Par hypothèse
sur $X$, on a nécessairement $X = X_1$ ou $X = X_2$, et donc $Y = Y_1$ ou
$Y = Y_2$, puisque $f$ est surjective.
\end{proof}

\begin{lemma}
\label{topology-lemma-irreducible}
Soit $X$ un espace topologique.
\begin{enumerate}
\item Si $T \subset X$ est irréductible, son adhérence dans $X$ l'est aussi.
\item Toute composante irréductible de $X$ est fermée.
\item Toute partie irréductible de $X$ est contenue dans une composante
irréductible de $X$.
\item Tout point de $X$ appartient à une composante irréductible de $X$ ;
autrement dit, $X$ est la réunion de ses composantes irréductibles.
\end{enumerate}
\end{lemma}

\begin{proof}
Soit $\overline{T}$ l'adhérence de la partie irréductible $T$.
Si $\overline{T} = Z_1 \cup Z_2$, avec $Z_i \subset \overline{T}$ fermé,
alors $T = (T\cap Z_1) \cup (T \cap Z_2)$ ; ainsi, $T$ est égal à l'une des
deux parties, disons $T = Z_1 \cap T$. Il est donc clair que
$\overline{T} \subset Z_1$. Cela prouve (1). Le point (2) résulte aussitôt de
(1) et de la définition des composantes irréductibles.

\medskip\noindent
Soit $T \subset X$ irréductible. Considérons l'ensemble $A$ des parties
irréductibles $T \subset T' \subset X$. Remarquons que $A$ est non vide,
puisque $T \in A$. L'inclusion définit sur $A$ un ordre partiel :
$T' \leq T'' \Leftrightarrow T' \subset T''$.
Choisissons une partie totalement ordonnée maximale $A' \subset A$ et posons
$E = \bigcup_{T' \in A'} T'$. Nous affirmons que $E$ est irréductible.
Supposons en effet que $E =  Z_1 \cup Z_2$ soit la réunion de deux parties
fermées de $E$. Pour tout $T' \in A'$, on a soit $T' \subset Z_1$, soit
$T' \subset Z_2$, par l'irréductibilité de $T'$. Supposons qu'il existe
$T'_0 \in A'$ tel que $T'_0 \not\subset Z_1$ (sinon, nous avons déjà fini).
Puisque $A'$ est totalement ordonné, nous voyons immédiatement que
$T' \subset Z_2$ pour tout $T' \in A'$. Ainsi $E = Z_2$.
Cela prouve (3). Le point (4) résulte immédiatement de (3), puisqu'un espace
réduit à un singleton est irréductible.
\end{proof}

\begin{lemma}
\label{topology-lemma-pick-irreducible-components}
Soit $X$ un espace topologique et supposons
$X = \bigcup_{i = 1, \ldots, n} X_i$, où chaque $X_i$ est une partie fermée
irréductible de $X$ et aucun $X_i$ n'est contenu dans la réunion des autres.
Alors chaque $X_i$ est une composante irréductible de $X$, et toute composante
irréductible de $X$ est l'un des $X_i$.
\end{lemma}

\begin{proof}
Soit $Y$ une composante irréductible de $X$. Écrivons
$Y = \bigcup_{i = 1, \ldots, n} (Y \cap X_i)$ et remarquons que chaque
$Y \cap X_i$ est fermé dans $Y$, puisque $X_i$ est fermé dans $X$.
L'irréductibilité de $Y$ montre que $Y$ est égal à l'un des $Y \cap X_i$,
c'est-à-dire que $Y \subset X_i$. Par maximalité des composantes
irréductibles, nous obtenons $Y = X_i$.

\medskip\noindent
Réciproquement, choisissons l'un des $X_i$. D'après le lemme
\ref{topology-lemma-irreducible}, il est contenu dans une composante irréductible $Y$
de $X$. Le paragraphe précédent donne $Y = X_j$ pour un certain $j$.
Ainsi $X_i \subset X_j$ et, comme la réunion initiale ne contient aucun terme
superflu, $X_i = X_j$, qui est une composante irréductible.
\end{proof}

\begin{lemma}
\label{topology-lemma-surjective-continuous-irreducible-components}
Soit $f : X \to Y$ une application continue et surjective d'espaces
topologiques. Si $X$ possède un nombre fini, disons $n$, de composantes
irréductibles, alors $Y$ possède un nombre de composantes irréductibles qui
est $\leq n$.
\end{lemma}

\begin{proof}
Soient $X_1, \ldots, X_n$ les composantes irréductibles de $X$.
D'après les lemmes \ref{topology-lemma-image-irreducible-space} et
\ref{topology-lemma-irreducible}, l'adhérence $Y_i \subset Y$ de $f(X_i)$ est
irréductible. Comme $f$ est surjective, $Y$ est la réunion des $Y_i$.
Nous pouvons choisir une partie minimale $I \subset \{1, \ldots, n\}$ telle
que $Y = \bigcup_{i \in I} Y_i$. Le lemme
\ref{topology-lemma-pick-irreducible-components} montre alors que les $Y_i$, pour
$i \in I$, sont les composantes irréductibles de $Y$.
\end{proof}

\noindent
Un singleton est irréductible. Ainsi, si $x \in X$ est un point, l'adhérence
$\overline{\{x\}}$ est une partie fermée irréductible de $X$.

\begin{definition}
\label{topology-definition-generic-point}
Soit $X$ un espace topologique.
\begin{enumerate}
\item Soit $Z \subset X$ une partie fermée irréductible. Un
{\it point générique} de $Z$ est un point $\xi \in Z$ tel que
$Z = \overline{\{\xi\}}$.
\item L'espace $X$ est dit {\it de Kolmogorov} si, pour tous
$x, x' \in X$ tels que $x \not = x'$, il existe une partie fermée de $X$ qui
contient exactement l'un des deux points.
\item L'espace $X$ est dit {\it quasi-sobre} si toute partie fermée
irréductible possède un point générique.
\item L'espace $X$ est dit {\it sobre} si toute partie fermée irréductible
possède un unique point générique.
\end{enumerate}
\end{definition}

\noindent
Un espace topologique $X$ est de Kolmogorov, quasi-sobre, resp.\ sobre si et
seulement si l'application $x\mapsto\overline{\{x\}}$ de $X$ vers l'ensemble
des parties fermées irréductibles de $X$ est injective, surjective, resp.\
bijective. Nous voyons donc qu'un espace topologique est sobre si et seulement
s'il est quasi-sobre et de Kolmogorov.

\begin{lemma}
\label{topology-lemma-sober-subspace}
Soit $X$ un espace topologique et soit $Y\subset X$.
\begin{enumerate}
\item Si $X$ est de Kolmogorov, alors $Y$ l'est aussi.
\item Supposons que $Y$ soit localement fermé dans $X$. Si $X$ est
quasi-sobre, alors $Y$ l'est aussi.
\item Supposons que $Y$ soit localement fermé dans $X$. Si $X$ est sobre,
alors $Y$ l'est aussi.
\end{enumerate}
\end{lemma}

\begin{proof}
Démonstration de (1). Supposons que $X$ soit de Kolmogorov. Soient
$x,y\in Y$ avec $x\neq y$. Alors
$\overline{\overline{\{x\}}\cap Y}=\overline{\{x\}}\neq\overline{\{y\}}=
\overline{\overline{\{y\}}\cap Y}$. Par conséquent,
$\overline{\{x\}}\cap Y\neq\overline{\{y\}}\cap Y$. Cela montre que $Y$ est
de Kolmogorov.

\medskip\noindent
Démonstration de (2). Supposons que $X$ soit quasi-sobre. Il suffit de traiter
les cas où $Y$ est fermé et où $Y$ est ouvert. Supposons d'abord $Y$ fermé.
Soit $Z$ une partie fermée irréductible de $Y$. Alors $Z$ est une partie fermée
irréductible de $X$. Il existe donc $x \in Z$ tel que
$\overline{\{x\}} = Z$. Il s'ensuit que $\overline{\{x\}} \cap Y = Z$.
Ainsi, $Y$ est quasi-sobre. Supposons ensuite $Y$ ouvert. Soit $Z$ une partie
fermée irréductible de $Y$. Alors $\overline{Z}$ est une partie fermée
irréductible de $X$. Il existe donc $x \in \overline{Z}$ tel que
$\overline{\{x\}}=\overline{Z}$. Si $x\notin Y$, nous obtenons la
contradiction
$Z=Z\cap Y\subset\overline{Z}\cap Y=\overline{\{x\}}\cap Y=\emptyset$.
Par conséquent, $x\in Y$. Il s'ensuit que
$Z=\overline{Z}\cap Y=\overline{\{x\}}\cap Y$. Ainsi, $Y$ est quasi-sobre.

\medskip\noindent
Démonstration de (3). Elle résulte immédiatement de (1) et (2).
\end{proof}

\begin{lemma}
\label{topology-lemma-sober-local}
Soit $X$ un espace topologique et soit $(X_i)_{i\in I}$ un recouvrement de
$X$.
\begin{enumerate}
\item Supposons $X_i$ localement fermé dans $X$ pour tout $i\in I$. Alors $X$
est de Kolmogorov si et seulement si $X_i$ l'est pour tout $i\in I$.
\item Supposons $X_i$ ouvert dans $X$ pour tout $i\in I$. Alors $X$ est
quasi-sobre si et seulement si $X_i$ l'est pour tout $i\in I$.
\item Supposons $X_i$ ouvert dans $X$ pour tout $i\in I$. Alors $X$ est sobre
si et seulement si $X_i$ l'est pour tout $i\in I$.
\end{enumerate}
\end{lemma}

\begin{proof}
Démonstration de (1). Si $X$ est de Kolmogorov, alors $X_i$ l'est pour tout
$i\in I$ d'après le lemme \ref{topology-lemma-sober-subspace}. Supposons $X_i$ de
Kolmogorov pour tout $i\in I$. Soient $x,y\in X$ tels que
$\overline{\{x\}}=\overline{\{y\}}$. Il existe $i\in I$ tel que $x\in X_i$.
Il existe un ouvert $U\subset X$ tel que $X_i$ soit fermé dans $U$.
Si $y\notin U$, nous obtenons la contradiction
$x\in\overline{\{x\}}\cap U=\overline{\{y\}}\cap U=\emptyset$. Ainsi,
$y\in U$. Il s'ensuit que
$y\in\overline{\{y\}}\cap U=\overline{\{x\}}\cap U\subset X_i$.
Cela montre que $y\in X_i$. Par conséquent,
$\overline{\{x\}}\cap X_i=\overline{\{y\}}\cap X_i$. Comme $X_i$ est de
Kolmogorov, nous obtenons $x=y$. Cela montre que $X$ est de Kolmogorov.

\medskip\noindent
Démonstration de (2). Si $X$ est quasi-sobre, alors $X_i$ l'est pour tout
$i\in I$ d'après le lemme \ref{topology-lemma-sober-subspace}. Supposons $X_i$
quasi-sobre pour tout $i\in I$. Soit $Y$ une partie fermée irréductible de
$X$. Comme $Y\neq\emptyset$, il existe $i\in I$ tel que
$X_i\cap Y\neq\emptyset$. Puisque $X_i$ est ouvert dans $X$, la partie
$X_i\cap Y$ est non vide et ouverte dans $Y$, donc irréductible et dense dans
$Y$. Ainsi, $X_i\cap Y$ est une partie fermée irréductible de $X_i$. Comme
$X_i$ est quasi-sobre, il existe $x\in X_i\cap Y$ tel que
$X_i\cap Y=\overline{\{x\}}\cap X_i\subset\overline{\{x\}}$. Puisque
$X_i\cap Y$ est dense dans $Y$ et que $Y$ est fermé dans $X$, il s'ensuit que
$Y=\overline{X_i\cap Y}\cap Y\subset\overline{X_i\cap Y}\subset
\overline{\{x\}}\subset Y$. Par conséquent,
$Y=\overline{\{x\}}$. Cela montre que $X$ est quasi-sobre.

\medskip\noindent
Démonstration de (3). Elle résulte immédiatement de (1) et (2).
\end{proof}

\begin{example}
\label{topology-example-quasi-sober-not-kolmogorov}
Soit $X$ un espace indiscret de cardinalité au moins $2$. Alors $X$ est
quasi-sobre mais non de Kolmogorov. De plus, la famille de ses singletons est
un recouvrement de $X$ par des espaces discrets, donc de Kolmogorov.
\end{example}

\begin{example}
\label{topology-example-kolmogorov-not-quasi-sober}
Soit $Y$ un ensemble infini, muni de la topologie dont les fermés sont $Y$ et
les parties finies de $Y$. Alors $Y$ est de Kolmogorov mais non quasi-sobre.
Cependant, la famille de ses singletons est un recouvrement par des espaces
discrets, donc sobres.
\end{example}

\begin{example}
\label{topology-example-not-kolmogorov-not-quasi-sober}
Soient $X$ et $Y$ comme dans l'exemple
\ref{topology-example-quasi-sober-not-kolmogorov} et l'exemple
\ref{topology-example-kolmogorov-not-quasi-sober}. Alors $X\amalg Y$ n'est ni de
Kolmogorov ni quasi-sobre.
\end{example}

\begin{example}
\label{topology-example-sober-subspace}
Soit $Z$ un ensemble infini et soit $z\in Z$. Munissons $Z$ de la topologie
dont les fermés sont $Z$ et les parties finies de $Z\setminus\{z\}$. Alors $Z$
est sobre, mais son sous-espace $Z\setminus\{z\}$ n'est pas quasi-sobre.
\end{example}

\begin{example}
\label{topology-example-Hausdorff}
Rappelons qu'un espace topologique $X$ est séparé si et seulement si, pour
toute paire de points distincts $x, y \in X$, il existe des ouverts disjoints
$U, V \subset X$ tels que $x \in U$, $y \in V$. Dans ce cas, $X$ est
irréductible si et seulement si $X$ est réduit à un singleton. De même, toute
partie de $X$ est irréductible si et seulement si elle est réduite à un
singleton. Ainsi, tout espace séparé est sobre.
\end{example}

\begin{lemma}
\label{topology-lemma-irreducible-on-top}
Soit $f : X \to Y$ une application continue d'espaces topologiques.
Supposons que
(a) $Y$ soit irréductible,
(b) $f$ soit ouverte, et
(c) il existe une famille dense de points $y \in Y$ tels que
$f^{-1}(y)$ soit irréductible.
Alors $X$ est irréductible.
\end{lemma}

\begin{proof}
Supposons $X = Z_1 \cup Z_2$, avec $Z_i$ fermé.
Considérons les ouverts $U_1 = Z_1 \setminus Z_2 = X \setminus Z_2$ et
$U_2 = Z_2 \setminus Z_1 = X \setminus Z_1$. Pour obtenir une contradiction,
supposons que $U_1$ et $U_2$ soient tous deux non vides. D'après (b),
$f(U_i)$ est ouvert. D'après (a), $Y$ est irréductible, et donc
$f(U_1) \cap f(U_2) \not = \emptyset$. D'après (c), cet ouvert non vide
contient un point $y$ tel que la fibre
$X_y = f^{-1}(y)$ soit irréductible. Alors $X_y \cap U_1$ et
$X_y \cap U_2$ sont des parties ouvertes non vides et disjointes de $X_y$,
ce qui est contradictoire.
\end{proof}

\begin{lemma}
\label{topology-lemma-irreducible-fibres-irreducible-components}
Soit $f : X \to Y$ une application continue d'espaces topologiques.
Supposons que (a) $f$ soit ouverte et que
(b) pour tout $y \in Y$, la fibre $f^{-1}(y)$ soit irréductible.
Alors $f$ induit une bijection entre les composantes irréductibles.
\end{lemma}

\begin{proof}
Remarquons que l'hypothèse (b) implique que $f$ est surjective (voir la
définition \ref{topology-definition-irreducible-components}).
Soit $T \subset Y$ une composante irréductible.
Remarquons que $T$ est fermée ; voir le lemme \ref{topology-lemma-irreducible}.
Le lemme résulte du fait que $f^{-1}(T)$ est irréductible, car toute partie
irréductible de $X$ s'envoie dans une composante irréductible de $Y$ d'après
le lemme \ref{topology-lemma-image-irreducible-space}.
Remarquons que $f^{-1}(T) \to T$ satisfait aux hypothèses du
lemme \ref{topology-lemma-irreducible-on-top}. Nous avons donc terminé.
\end{proof}

\noindent
La construction du lemme suivant est parfois appelée la « sobrification ».

\begin{lemma}
\label{topology-lemma-make-sober}
\begin{reference}
\cite[p. 68 et suiv.]{EGA1-second}
\end{reference}
Soit $X$ un espace topologique. Il existe une application continue canonique
$$
c : X \longrightarrow X'
$$
de $X$ vers un espace topologique sobre $X'$, universelle parmi les
applications continues de $X$ vers les espaces topologiques sobres.
De plus, l'application $U' \mapsto c^{-1}(U')$ est une bijection entre les
ouverts de $X'$ et ceux de $X$, qui commute aux intersections finies et aux
réunions quelconques. L'image $c(X)$ est un espace topologique de Kolmogorov,
et l'application $c : X \to c(X)$ est universelle parmi les applications de
$X$ vers les espaces de Kolmogorov.
\end{lemma}

\begin{proof}
Soit $X'$ l'ensemble des parties fermées irréductibles de $X$ et soit
$$
c : X \to X', \quad x \mapsto \overline{\{x\}}.
$$
Pour tout ouvert $U \subset X$, notons $U' \subset X'$ l'ensemble des
parties fermées irréductibles de $X$ qui rencontrent $U$.
Alors $c^{-1}(U') = U$. En particulier, si $U_1 \not = U_2$ sont ouverts dans
$X$, alors $U'_1 \not = U_2'$. Ainsi, $c$ induit une bijection entre les
parties de $X'$ de la forme $U'$ et les ouverts de $X$.

\medskip\noindent
Soient $U_1, U_2$ des ouverts de $X$. Supposons que $Z \in U'_1$ et
$Z \in U'_2$. Alors $Z \cap U_1$ et $Z \cap U_2$ sont des parties ouvertes
non vides de l'espace irréductible $Z$ ; par conséquent,
$Z \cap U_1 \cap U_2$ est non vide. Ainsi,
$(U_1 \cap U_2)' = U'_1 \cap U'_2$.
La règle $U \mapsto U'$ est également compatible avec les réunions
quelconques (nous omettons les détails). Il est donc clair que la famille des
$U'$ forme une topologie sur $X'$ et que l'on obtient la bijection annoncée
dans le lemme.

\medskip\noindent
Montrons ensuite que $X'$ est sobre. Soit $T \subset X'$ une partie fermée
irréductible. Soit $U \subset X$ l'ouvert tel que $X' \setminus T = U'$.
Alors $Z = X \setminus U$ est irréductible, en vertu des propriétés de la
bijection du lemme. Nous affirmons que $Z \in T$ est l'unique point générique.
En effet, tout ouvert de la forme $V' \subset X'$ qui ne contient pas $Z$
provient nécessairement d'un ouvert $V \subset X$ qui ne rencontre pas $Z$,
c'est-à-dire qui est contenu dans $U$.

\medskip\noindent
Vérifions enfin la propriété universelle. Soit $f : X \to Y$ une application
continue à valeurs dans un espace topologique sobre. Définissons alors
$f' : X' \to Y$ comme l'application qui envoie la partie fermée irréductible
$Z \subset X$ sur l'unique point générique de $\overline{f(Z)}$.
Il est immédiat que $f' \circ c = f$ comme applications d'ensembles, et les
propriétés de $c$ entraînent que $f'$ est continue. Nous omettons la
vérification de l'unicité de l'application continue $f'$. Nous omettons aussi
la démonstration des assertions concernant les espaces de Kolmogorov.
\end{proof}

\begin{lemma}
\label{topology-lemma-remove-irreducible-connected}
Soit $X$ un espace topologique connexe possédant un nombre fini de composantes
irréductibles $X_1, \ldots, X_n$. Si $n > 1$, il existe
$1 \leq j \leq n$ tel que $X' = \bigcup_{i \not = j} X_i$ soit connexe.
\end{lemma}

\begin{proof}
C'est un problème de théorie des graphes. Soit $\Gamma$ le graphe dont les
sommets sont $V = \{1, \ldots, n\}$ et dans lequel $i$ et $j$ sont reliés par
une arête si et seulement si $X_i \cap X_j$ est non vide. La connexité de $X$
signifie que $\Gamma$ est connexe. Il s'agit de trouver
$1 \leq j \leq n$ tel que $\Gamma \setminus \{j\}$ soit encore connexe.
On peut le faire en choisissant $j, j' \in E$ à distance maximale ; alors $j$
convient (choisir une feuille !). Nous omettons les détails.
\end{proof}

\section{Espaces topologiques noethériens}
\label{topology-section-noetherian}

\begin{definition}
\label{topology-definition-noetherian}
Un espace topologique est dit {\it noethérien} si les parties fermées de $X$
satisfont à la condition de chaîne descendante. Un espace topologique est dit
{\it localement noethérien} si tout point possède un voisinage noethérien.
\end{definition}

\begin{lemma}
\label{topology-lemma-Noetherian}
Soit $X$ un espace topologique noethérien.
\begin{enumerate}
\item Toute partie de $X$, munie de la topologie induite, est noethérienne.
\item L'espace $X$ possède un nombre fini de composantes irréductibles.
\item Chaque composante irréductible de $X$ contient un ouvert non vide de
$X$.
\end{enumerate}
\end{lemma}

\begin{proof}
Soit $T \subset X$ une partie de $X$.
Soit $T_1 \supset T_2 \supset \ldots$ une chaîne descendante de parties
fermées de $T$. Écrivons $T_i =  T \cap Z_i$, avec $Z_i \subset X$ fermé.
Considérons la chaîne descendante de parties fermées
$Z_1 \supset Z_1\cap Z_2 \supset Z_1 \cap Z_2 \cap Z_3 \ldots$
Elle est stationnaire par hypothèse, et la suite initiale des $T_i$ l'est donc
elle aussi. Ainsi $T$ est noethérien.

\medskip\noindent
Soit $A$ l'ensemble des parties fermées de $X$ qui n'ont pas un nombre fini de
composantes irréductibles. Supposons que $A$ soit non vide afin d'aboutir à une
contradiction. L'inclusion définit sur $A$ un ordre partiel :
$\alpha \leq \alpha' \Leftrightarrow Z_{\alpha} \subset Z_{\alpha'}$.
La condition de chaîne descendante permet de trouver un plus petit élément de
$A$, disons $Z$. Comme $Z$ n'est pas une réunion finie de composantes
irréductibles, il n'est pas irréductible. Nous pouvons donc écrire
$Z = Z' \cup Z''$, où les deux termes sont des parties fermées strictement plus
petites. Par construction, $Z' = \bigcup Z'_i$ et
$Z'' = \bigcup Z''_j$ sont des réunions finies de leurs composantes
irréductibles (lemme \ref{topology-lemma-irreducible}). Ainsi,
$Z = \bigcup Z'_i \cup \bigcup Z''_j$ est une réunion finie de parties fermées
irréductibles. Après suppression des termes superflus de cette expression, on
obtient la décomposition de $Z$ en ses composantes irréductibles (lemme
\ref{topology-lemma-pick-irreducible-components}), ce qui est contradictoire.

\medskip\noindent
Soit $Z \subset X$ une composante irréductible de $X$.
Soient $Z_1, \ldots, Z_n$ les autres composantes irréductibles de $X$.
Considérons $U = Z \setminus (Z_1\cup\ldots\cup Z_n)$.
Cette partie n'est pas vide, sinon l'espace irréductible $Z$ serait contenu
dans l'un des autres $Z_i$. Comme
$X = Z \cup Z_1 \cup \ldots Z_n$ (voir le lemme \ref{topology-lemma-irreducible}), on a
également $U = X \setminus (Z_1\cup\ldots\cup Z_n)$, qui est donc ouvert dans
$X$. Ainsi, $Z$ contient un ouvert non vide de $X$.
\end{proof}

\begin{lemma}
\label{topology-lemma-image-Noetherian}
Soit $f : X \to Y$ une application continue d'espaces topologiques.
\begin{enumerate}
\item Si $X$ est noethérien, alors $f(X)$ est noethérien.
\item Si $X$ est localement noethérien et si $f$ est ouverte, alors $f(X)$ est
localement noethérien.
\end{enumerate}
\end{lemma}

\begin{proof}
Dans le cas (1), supposons que
$Z_1 \supset Z_2 \supset Z_3 \supset \ldots$ soit une chaîne descendante de
parties fermées de $f(X)$ (muni, comme d'habitude, de la topologie induite en
tant que partie de $Y$). Alors
$f^{-1}(Z_1) \supset f^{-1}(Z_2) \supset f^{-1}(Z_3) \supset \ldots$ est une
chaîne descendante de parties fermées de $X$. Cette chaîne est donc
stationnaire. Puisque $f(f^{-1}(Z_i)) = Z_i$, nous en concluons que
$Z_1 \supset Z_2 \supset Z_3 \supset \ldots$ est elle aussi stationnaire.
Dans le cas (2), soit $y \in f(X)$. Choisissons $x \in X$ tel que
$f(x) = y$. Par hypothèse, il existe un voisinage $E \subset X$ de $x$ qui est
noethérien. Alors $f(E) \subset f(X)$ est un voisinage noethérien d'après le
point (1).
\end{proof}

\begin{lemma}
\label{topology-lemma-finite-union-Noetherian}
Soit $X$ un espace topologique.
Soient $X_i \subset X$, $i = 1, \ldots, n$, une famille finie de parties.
Si chaque $X_i$ est noethérien (pour la topologie induite), alors
$\bigcup_{i = 1, \ldots, n}  X_i$ est noethérien (pour la topologie induite).
\end{lemma}

\begin{proof}
Soit $\{F_m\}_{m \in \mathbf{N}}$ une suite décroissante de parties fermées de
$X' = \bigcup_{i = 1, \ldots, n}  X_i$ muni de la topologie induite.
On peut trouver une suite décroissante $\{G_m\}_{m \in \mathbf{N}}$ de parties
fermées de $X$ vérifiant $F_m = G_m \cap X'$ pour tout $m$ (nous omettons ce
petit détail). Comme $X_i$ est noethérien et que
$\{G_m \cap X_i\}_{m \in \mathbf{N}}$ est une suite décroissante de parties
fermées de $X_i$, il existe $m_i \in \mathbf{N}$ tel que, pour tout
$m \geq m_i$, on ait $G_m \cap X_i = G_{m_i} \cap X_i$.
Posons $m_0 = \max_{i = 1, \ldots, n} m_i$. Alors, manifestement,
$$
F_m = G_m \cap X' = G_m \cap (X_1 \cup \ldots \cup X_n) =
(G_m \cap X_1) \cup \ldots (G_m \cap X_n)
$$
est stationnaire pour $m \geq m_0$, ce qui achève la démonstration.
\end{proof}

\begin{example}
\label{topology-example-locally-Noetherian-no-closed-point}
Tout espace topologique noethérien non vide et de Kolmogorov possède un point
fermé (combiner les lemmes \ref{topology-lemma-quasi-compact-closed-point} et
\ref{topology-lemma-Noetherian-quasi-compact}).
Soit $X = \{1, 2, 3, \ldots \}$. Définissons sur $X$ la topologie dont les
ouverts sont $\emptyset$, $\{1, 2, \ldots, n\}$, $n \geq 1$, et $X$.
Alors $X$ est un espace topologique localement noethérien sans aucun point
fermé. Cet espace ne peut pas être l'espace topologique sous-jacent à un schéma
localement noethérien ; voir Propriétés, lemme
\ref{properties-lemma-locally-Noetherian-closed-point}.
\end{example}

\begin{lemma}
\label{topology-lemma-locally-Noetherian-locally-connected}
Soit $X$ un espace topologique localement noethérien.
Alors $X$ est localement connexe.
\end{lemma}

\begin{proof}
Soit $x \in X$. Soit $E$ un voisinage de $x$.
Nous devons trouver un voisinage connexe de $x$ contenu dans $E$.
Par hypothèse, il existe un voisinage $E'$ de $x$ qui est noethérien.
Alors $E \cap E'$ est noethérien ; voir le lemme \ref{topology-lemma-Noetherian}.
Soit $E \cap E' = Y_1 \cup \ldots \cup Y_n$ sa décomposition en composantes
irréductibles ; voir le lemme \ref{topology-lemma-Noetherian}.
Soit $E'' = \bigcup_{x \in Y_i} Y_i$. C'est une partie connexe de
$E \cap E'$ qui contient $x$. Elle contient l'ouvert
$E \cap E' \setminus (\bigcup_{x \not \in Y_i} Y_i)$ de $E \cap E'$ ; c'est
donc un voisinage de $x$ dans $X$. Cela prouve le lemme.
\end{proof}




\section{Dimension de Krull}
\label{topology-section-krull-dimension}

\begin{definition}
\label{topology-definition-Krull}
Soit $X$ un espace topologique.
\begin{enumerate}
\item Une {\it chaîne de parties fermées irréductibles} de $X$ est une suite
$Z_0 \subset Z_1 \subset \ldots \subset Z_n \subset X$, où $Z_i$ est fermé
irréductible et $Z_i \not = Z_{i + 1}$ pour $i = 0, \ldots, n - 1$.
\item La {\it longueur} d'une chaîne
$Z_0 \subset Z_1 \subset \ldots \subset Z_n \subset X$ de parties fermées
irréductibles de $X$ est l'entier $n$.
\item La {\it dimension}, ou plus précisément la {\it dimension de Krull},
$\dim(X)$ de $X$, est l'élément de
$\{-\infty, 0, 1, 2, 3, \ldots, \infty\}$ défini par la formule
$$
\dim(X) =
\sup \{\text{longueurs des chaînes de parties fermées irréductibles}\}
$$
Ainsi, $\dim(X) = -\infty$ si et seulement si $X$ est l'espace vide.
\item Soit $x \in X$. La {\it dimension de Krull de $X$ en $x$} est définie
par
$$
\dim_x(X) = \min \{\dim(U), x\in U\subset X\text{ ouvert}\}
$$
c'est-à-dire comme le minimum des $\dim(U)$ lorsque $U$ parcourt les
voisinages ouverts de $x$ dans $X$.
\end{enumerate}
\end{definition}

\noindent
Remarquons que, si $U' \subset U \subset X$ sont ouverts, alors
$\dim(U') \leq \dim(U)$. Ainsi, si $\dim_x(X) = d$, le point $x$ possède un
système fondamental de voisinages ouverts $U$ tels que
$\dim(U) = \dim_x(X)$.

\begin{lemma}
\label{topology-lemma-dimension-supremum-local-dimensions}
Soit $X$ un espace topologique. Alors $\dim(X) = \sup \dim_x(X)$, où la borne
supérieure est prise sur les points $x$ de $X$.
\end{lemma}

\begin{proof}
Il est clair que $\dim(X) \geq \dim_x(X)$ pour tout $x \in X$ (voir la
discussion qui suit la définition \ref{topology-definition-Krull}). Nous obtenons ainsi
l'une des inégalités. Pour la réciproque, soit $n \geq 0$ et supposons que
$\dim(X) \geq n$. On peut alors trouver une chaîne de parties fermées
irréductibles $Z_0 \subset Z_1 \subset \ldots \subset Z_n \subset X$.
Choisissons $x \in Z_0$. Pour tout voisinage ouvert $U$ de $x$, nous obtenons
une chaîne de parties fermées irréductibles
$$
Z_0 \cap U \subset Z_1 \cap U \subset \ldots \subset Z_n \cap U
$$
dans $U$. En effet, les ensembles $U \cap Z_i$ sont fermés irréductibles dans
$U$ et
les inclusions sont strictes (détails omis ; indication : l'adhérence de
$U \cap Z_i$ est $Z_i$).
Nous voyons ainsi que $\dim_x(X) \geq n$, ce qui prouve l'autre inégalité.
\end{proof}

\begin{example}
\label{topology-example-Krull-Rn}
La dimension de Krull de l'espace euclidien usuel $\mathbf{R}^n$ est $0$.
\end{example}

\begin{example}
\label{topology-example-krull-2set}
Soit $X = \{s, \eta\}$, dont les ouverts sont
$\{\emptyset, \{\eta\}, \{s, \eta\}\}$.
Dans ce cas, une chaîne maximale de parties fermées irréductibles est
$\{s\} \subset \{s, \eta\}$. Ainsi, $\dim(X) = 1$. Il est facile de
généraliser cet exemple pour obtenir un espace topologique à $(n + 1)$
éléments et de dimension de Krull $n$.
\end{example}

\begin{definition}
\label{topology-definition-equidimensional}
Soit $X$ un espace topologique.
On dit que $X$ est {\it équidimensionnel} si toutes les composantes
irréductibles de $X$ ont la même dimension.
\end{definition}





\section{Codimension et espaces caténaires}
\label{topology-section-catenary-spaces}

\noindent
Nous ne définissons que la codimension des parties fermées irréductibles.

\begin{definition}
\label{topology-definition-codimension}
Soit $X$ un espace topologique.
Soit $Y \subset X$ une partie fermée irréductible.
La {\it codimension} de $Y$ dans $X$ est la borne supérieure des longueurs
$e$ des chaînes
$$
Y = Y_0 \subset Y_1 \subset \ldots \subset Y_e \subset X
$$
de parties fermées irréductibles de $X$ qui commencent par $Y$.
Nous la noterons $\text{codim}(Y, X)$.
\end{definition}

\noindent
La codimension est un élément de $\{0, 1, 2, \ldots\} \cup \{\infty\}$.
Si $\text{codim}(Y, X) < \infty$, toute chaîne peut être prolongée en une
chaîne maximale (mais celles-ci n'ont pas nécessairement toutes la même
longueur).

\begin{lemma}
\label{topology-lemma-codimension-at-generic-point}
Soit $X$ un espace topologique.
Soit $Y \subset X$ une partie fermée irréductible.
Soit $U \subset X$ un ouvert tel que $Y \cap U$ soit non vide. Alors
$$
\text{codim}(Y, X) = \text{codim}(Y \cap U, U)
$$
\end{lemma}

\begin{proof}
La règle $T \mapsto \overline{T}$ définit une application bijective qui
préserve les inclusions entre les parties fermées irréductibles de $U$ et les
parties fermées irréductibles de $X$ qui rencontrent $U$.
Le lemme en résulte facilement. Nous omettons les détails.
\end{proof}

\begin{example}
\label{topology-example-Noetherian-infinite-codimension}
Soit $X = [0, 1]$ l'intervalle unité muni de la topologie suivante : les
ensembles $[0, 1]$, $(1 - 1/n, 1]$ pour $n \in \mathbf{N}$, et $\emptyset$
sont ouverts. Les fermés sont donc $\emptyset$, $\{0\}$,
$[0, 1 - 1/n]$ pour $n > 1$, et $[0, 1]$.
C'est manifestement un espace topologique noethérien. Mais la partie fermée
irréductible $Y = \{0\}$ est de codimension infinie :
$\text{codim}(Y, X) = \infty$.
Pour le voir, il suffit de remarquer que tous les fermés $[0, 1 - 1/n]$ sont
irréductibles.
\end{example}

\begin{definition}
\label{topology-definition-catenary}
Soit $X$ un espace topologique. On dit que $X$ est {\it caténaire} si, pour
toute paire de parties fermées irréductibles $T \subset T'$, on a
$\text{codim}(T, T') < \infty$ et si toute chaîne maximale de parties fermées
irréductibles
$$
T = T_0 \subset T_1 \subset \ldots \subset T_e = T'
$$
a la même longueur (égale à la codimension).
\end{definition}

\begin{lemma}
\label{topology-lemma-catenary}
Soit $X$ un espace topologique.
Les assertions suivantes sont équivalentes :
\begin{enumerate}
\item $X$ est caténaire,
\item $X$ possède un recouvrement ouvert par des espaces caténaires.
\end{enumerate}
De plus, dans ce cas, tout sous-espace localement fermé de $X$ est caténaire.
\end{lemma}

\begin{proof}
Supposons que $X$ soit caténaire et que $U \subset X$ soit ouvert.
La règle $T \mapsto \overline{T}$ définit une application bijective qui
préserve les inclusions entre les parties fermées irréductibles de $U$ et les
parties fermées irréductibles de $X$ qui rencontrent $U$.
Le lemme en résulte facilement. Nous omettons les détails.
\end{proof}

\begin{lemma}
\label{topology-lemma-catenary-in-codimension}
Soit $X$ un espace topologique. Les assertions suivantes sont équivalentes :
\begin{enumerate}
\item $X$ est caténaire, et
\item pour toute paire de parties fermées irréductibles $Y \subset Y'$, on a
$\text{codim}(Y, Y') < \infty$, et pour tout triplet
$Y \subset Y' \subset Y''$ de parties fermées irréductibles, on a
$$
\text{codim}(Y, Y'') = \text{codim}(Y, Y') + \text{codim}(Y', Y'').
$$
\end{enumerate}
\end{lemma}

\begin{proof}
Supposons que $X$ soit caténaire. D'après la définition
\ref{topology-definition-catenary}, pour toute paire de parties fermées irréductibles
$Y \subset Y'$, on a $\text{codim}(Y,Y') < \infty$.
Soit $Y \subset Y' \subset Y''$ un triplet de parties fermées irréductibles de
$X$. Soit
$$
Y = Y_0 \subset Y_1 \subset ... \subset Y_{e_1} = Y'
$$
une chaîne maximale de parties fermées irréductibles entre $Y$ et $Y'$, où
$e_1 = \text{codim}(Y,Y')$. Soit aussi
$$
Y' = Y_{e_1} \subset Y_{e_1 + 1}\subset ... \subset Y_{e_1 + e_2} = Y''
$$
une chaîne maximale de parties fermées irréductibles entre $Y'$ et $Y''$, où
$e_2 = \text{codim}(Y',Y'')$.
Comme les deux chaînes sont maximales, leur concaténation
$$
Y = Y_0\subset Y_1 \subset ... \subset Y_{e_1} = Y' =
Y_{e_1} \subset Y_{e_1+1}\subset ... \subset Y_{e_1+e_2}=Y''
$$
est elle aussi maximale (entre $Y$ et $Y''$), et sa longueur est
$e_1 + e_2$. Puisque $X$ est caténaire, toute chaîne maximale a la même
longueur, égale à la codimension. Le point (2), à savoir
$\text{codim}(Y,Y'') = e_1 + e_2 = \text{codim}(Y,Y') + \text{codim}(Y',Y'')$,
est donc vérifié.

\medskip\noindent
Réciproquement, une récurrence montre que, si
$Y = Y_1 \subset ... \subset Y_n = Y'$, alors
$ \text{codim}(Y,Y') = \text{codim}(Y_1,Y_2) + ... +
\text{codim}(Y_{n-1},Y_n)$. Cela force les chaînes maximales à avoir la même
longueur.
\end{proof}








\section{Espaces et applications quasi-compacts}
\label{topology-section-quasi-compact}

\noindent
Nous réserverons le mot « compact » aux espaces topologiques séparés. Or de
nombreux espaces qui interviennent en géométrie algébrique ne sont pas séparés.

\begin{definition}
\label{topology-definition-quasi-compact}
Quasi-compacité.
\begin{enumerate}
\item On dit qu'un espace topologique $X$ est {\it quasi-compact} si tout
recouvrement ouvert de $X$ admet un sous-recouvrement fini.
\item On dit qu'une application continue $f : X \to Y$ est
{\it quasi-compacte} si l'image réciproque $f^{-1}(V)$ de tout ouvert
quasi-compact $V \subset Y$ est quasi-compacte.
\item On dit qu'une partie $Z \subset X$ est {\it rétrocompacte} si
l'application d'inclusion $Z \to X$ est quasi-compacte.
\end{enumerate}
\end{definition}

\noindent
Dans de nombreux textes de topologie, un espace est dit {\it compact} s'il est
quasi-compact et séparé ; dans d'autres, la condition de séparation est omise.
Pour éviter toute ambiguïté en géométrie algébrique, nous employons le terme
quasi-compact. La quasi-compacité d'une application est très différente de la
notion d'« application propre » : pour celle-ci, on exige (outre le caractère
fermé et séparé) que l'image réciproque de toute partie quasi-compacte du but
soit quasi-compacte, tandis que la définition ci-dessus ne considère que les
ensembles {\it ouverts} quasi-compacts.

\begin{lemma}
\label{topology-lemma-composition-quasi-compact}
Une composée d'applications quasi-compactes est quasi-compacte.
\end{lemma}

\begin{proof}
C'est immédiat d'après la définition.
\end{proof}

\begin{lemma}
\label{topology-lemma-closed-in-quasi-compact}
Une partie fermée d'un espace topologique quasi-compact est quasi-compacte.
\end{lemma}

\begin{proof}
Soit $E \subset X$ une partie fermée de l'espace quasi-compact $X$.
Soit $E = \bigcup V_j$ un recouvrement ouvert. Choisissons des ouverts
$U_j \subset X$ tels que $V_j = E \cap U_j$. Alors
$X = (X \setminus E) \cup \bigcup U_j$ est un recouvrement ouvert de $X$.
On a donc $X = (X \setminus E) \cup U_{j_1} \cup \ldots \cup U_{j_n}$ pour
un certain $n$ et des indices $j_i$.
Ainsi $E = V_{j_1} \cup \ldots \cup V_{j_n}$, comme voulu.
\end{proof}

\begin{lemma}
\label{topology-lemma-quasi-compact-in-Hausdorff}
Soit $X$ un espace topologique séparé.
\begin{enumerate}
\item Si $E \subset X$ est quasi-compact, alors il est fermé.
\item Si $E_1, E_2 \subset X$ sont deux parties quasi-compactes disjointes,
il existe des ouverts $E_i \subset U_i$ tels que
$U_1 \cap U_2 = \emptyset$.
\end{enumerate}
\end{lemma}

\begin{proof}
Démonstration de (1). Soit $x \in X$, avec $x \not \in E$.
Pour tout $e \in E$, on peut trouver des ouverts disjoints $V_e$ et $U_e$
tels que $e \in V_e$ et $x \in U_e$. Comme $E \subset \bigcup V_e$, il
existe un nombre fini d'éléments $e_1, \ldots, e_n$ tels que
$E \subset V_{e_1} \cup \ldots \cup V_{e_n}$. Alors
$U = U_{e_1} \cap \ldots \cap U_{e_n}$ est un voisinage ouvert de $x$ qui ne
rencontre pas $V_{e_1} \cup \ldots \cup V_{e_n}$. En particulier, il ne
rencontre pas $E$. Ainsi $E$ est fermé.

\medskip\noindent
Démonstration de (2). Dans la démonstration de (1), nous avons vu que, pour
$x \in E_1$, on peut trouver un voisinage ouvert $x \in U_x$ et un ouvert
$E_2 \subset V_x$ tels que $U_x \cap V_x = \emptyset$. Comme $E_1$ est
quasi-compact, il existe un nombre fini de points $x_i \in E_1$ tels que
$E_1 \subset U = U_{x_1} \cup \ldots \cup U_{x_n}$.
Il suffit de prendre $V = V_{x_1} \cap \ldots \cap V_{x_n}$.
\end{proof}

\begin{lemma}
\label{topology-lemma-closed-in-compact}
Soit $X$ un espace séparé quasi-compact. Soit $E \subset X$.
Les assertions suivantes sont équivalentes : (a) $E$ est fermé dans $X$,
(b) $E$ est quasi-compact.
\end{lemma}

\begin{proof}
L'implication (a) $\Rightarrow$ (b) est le
lemme \ref{topology-lemma-closed-in-quasi-compact}.
L'implication (b) $\Rightarrow$ (a) est le
lemme \ref{topology-lemma-quasi-compact-in-Hausdorff}.
\end{proof}

\noindent
L'énoncé suivant n'est qu'une reformulation de la quasi-compacité.

\begin{lemma}
\label{topology-lemma-intersection-closed-in-quasi-compact}
Soit $X$ un espace topologique quasi-compact.
Si $\{Z_\alpha\}_{\alpha \in A}$ est une famille de parties fermées telle que
l'intersection de toute sous-famille finie soit non vide, alors
$\bigcap_{\alpha \in A} Z_\alpha$ est non vide.
\end{lemma}

\begin{proof}
Supposons que $\bigcap_{\alpha \in A} Z_{\alpha} = \emptyset$.
En passant aux complémentaires, on a
$\bigcup_{\alpha \in A} (X\setminus Z_{\alpha})=X$.
Comme les parties $Z_\alpha$ sont fermées,
$\bigcup_{\alpha \in A} (X \setminus Z_{\alpha})$ est un recouvrement ouvert
de l'espace quasi-compact $X$. Il existe donc une partie finie $J \subset A$
telle que $X = \bigcup_{\alpha \in J} (X\setminus Z_{\alpha})$.
Son complémentaire est alors vide, ce qui signifie que
$\bigcap_{\alpha \in J} Z_{\alpha} = \emptyset$. Il existe donc une
sous-famille finie de $\{Z_{\alpha}\}_{\alpha \in J}$ vérifiant
$\bigcap_{\alpha \in J} Z_{\alpha}=\emptyset$, ce qui conclut par
contraposition.
\end{proof}

\begin{lemma}
\label{topology-lemma-image-quasi-compact}
Soit $f : X \to Y$ une application continue d'espaces topologiques.
\begin{enumerate}
\item Si $X$ est quasi-compact, alors $f(X)$ est quasi-compact.
\item Si $f$ est quasi-compacte, alors $f(X)$ est rétrocompact.
\end{enumerate}
\end{lemma}

\begin{proof}
Si $f(X) = \bigcup V_i$ est un recouvrement ouvert, alors
$X = \bigcup f^{-1}(V_i)$ est un recouvrement ouvert. Par conséquent, si $X$
est quasi-compact, on a
$X = f^{-1}(V_{i_1}) \cup \ldots \cup f^{-1}(V_{i_n})$ pour certains
$i_1, \ldots, i_n \in I$, et donc
$f(X) = V_{i_1} \cup \ldots \cup V_{i_n}$. Cela prouve (1).
Supposons $f$ quasi-compacte et soit $V \subset Y$ un ouvert quasi-compact.
Alors $f^{-1}(V)$ est quasi-compact ; d'après (1),
$f(f^{-1}(V)) = f(X) \cap V$ est donc quasi-compact. Ainsi, $f(X)$ est
rétrocompact.
\end{proof}

\begin{lemma}
\label{topology-lemma-quasi-compact-closed-point}
Soit $X$ un espace topologique. Supposons que
\begin{enumerate}
\item $X$ soit non vide,
\item $X$ soit quasi-compact, et
\item $X$ soit de Kolmogorov.
\end{enumerate}
Alors $X$ possède un point fermé.
\end{lemma}

\begin{proof}
Considérons l'ensemble
$$
\mathcal{T} =
\{Z \subset X \mid Z = \overline{\{x\}} \text{ pour un certain }x \in X\}
$$
de toutes les adhérences de singletons de $X$. Il est non vide, puisque $X$
est non vide. Munissons $\mathcal{T}$ de l'ordre partiel défini par
l'inclusion. Supposons que $Z_\alpha$, $\alpha \in A$, soit une partie
totalement ordonnée de $\mathcal{T}$.
D'après le lemme \ref{topology-lemma-intersection-closed-in-quasi-compact},
$\bigcap_{\alpha \in A} Z_\alpha \not = \emptyset$. Il existe donc
$x \in \bigcap_{\alpha \in A} Z_\alpha$, et
$Z = \overline{\{x\}}\in \mathcal{T}$ est un minorant de la famille.
Le lemme de Zorn assure l'existence d'un élément minimal
$Z \in \mathcal{T}$. Comme $X$ est de Kolmogorov, nous en concluons que
$Z = \{x\}$ pour un certain $x$, et $x \in X$ est un point fermé.
\end{proof}

\begin{lemma}
\label{topology-lemma-closed-points-quasi-compact}
Soit $X$ un espace de Kolmogorov quasi-compact. Alors l'ensemble $X_0$ des
points fermés de $X$ est quasi-compact.
\end{lemma}

\begin{proof}
Soit $X_0 = \bigcup U_{i, 0}$ un recouvrement ouvert.
Écrivons $U_{i, 0} = X_0 \cap U_i$ pour un ouvert $U_i \subset X$.
Considérons le complémentaire $Z$ de $\bigcup U_i$. C'est une partie fermée de
$X$, donc quasi-compacte (lemme \ref{topology-lemma-closed-in-quasi-compact})
et de Kolmogorov. D'après le lemme \ref{topology-lemma-quasi-compact-closed-point},
si $Z$ n'était pas vide, il aurait un point
fermé, ce qui contredit le fait que $X_0 \subset \bigcup U_i$.
Ainsi $Z = \emptyset$ et $X = \bigcup U_i$. Comme $X$ est quasi-compact,
ce recouvrement admet un sous-recouvrement fini, ce qui permet de conclure.
\end{proof}

\begin{lemma}
\label{topology-lemma-connected-component-intersection}
Soit $X$ un espace topologique.
Supposons que
\begin{enumerate}
\item $X$ soit quasi-compact,
\item $X$ possède une base de la topologie formée d'ouverts quasi-compacts, et
\item l'intersection de deux ouverts quasi-compacts soit quasi-compacte.
\end{enumerate}
Pour tout $x \in X$, la composante connexe de $X$ contenant
$x$ est l'intersection de toutes les parties à la fois ouvertes et fermées
de $X$ contenant $x$.
\end{lemma}

\begin{proof}
Soit $T$ la composante connexe contenant $x$.
Soit $S = \bigcap_{\alpha \in A} Z_\alpha$ l'intersection de toutes les
parties $Z_\alpha$ de $X$, à la fois ouvertes et fermées, contenant $x$.
Remarquons que $S$ est fermé dans $X$.
Remarquons que toute intersection finie des $Z_\alpha$ est l'un des $Z_\alpha$.
Puisque $T$ est connexe et que $x \in T$, nous avons $T \subset S$.
Il suffit de montrer que $S$ est connexe.
Sinon, il existe une décomposition en réunion disjointe
$S = B \amalg C$ où $B$ et $C$ sont à la fois ouverts et fermés dans $S$.
En particulier, $B$ et $C$ sont fermés dans $X$, et donc quasi-compacts d'après le
lemme \ref{topology-lemma-closed-in-quasi-compact} et l'hypothèse (1).
D'après l'hypothèse (2), il existe des ouverts quasi-compacts
$U, V \subset X$ tels que $B = S \cap U$ et $C = S \cap V$ (détails omis).
Alors $U \cap V \cap S = \emptyset$.
Ainsi $\bigcap_\alpha U \cap V \cap Z_\alpha = \emptyset$.
D'après l'hypothèse (3), l'intersection $U \cap V$ est quasi-compacte.
D'après le lemme \ref{topology-lemma-intersection-closed-in-quasi-compact},
pour un certain $\alpha' \in A$, nous avons $U \cap V \cap Z_{\alpha'} = \emptyset$.
Comme $X \setminus (U \cup V)$ est disjoint de $S$
et fermé dans $X$, donc quasi-compact, nous pouvons appliquer le même lemme
pour voir que $Z_{\alpha''} \subset U \cup V$ pour un certain $\alpha'' \in A$.
Alors $Z_\alpha = Z_{\alpha'} \cap Z_{\alpha''}$ est contenu
dans $U \cup V$ et disjoint de $U \cap V$.
Ainsi $Z_\alpha = U \cap Z_\alpha \amalg V \cap Z_\alpha$
est une décomposition en deux parties ouvertes,
donc $U \cap Z_\alpha$ et $V \cap Z_\alpha$ sont ouverts et fermés dans $X$.
Ainsi, si $x \in B$ par exemple, alors nous voyons que $S \subset U \cap Z_\alpha$,
et nous concluons que $C = \emptyset$.
\end{proof}

\begin{lemma}
\label{topology-lemma-connected-component-intersection-compact-Hausdorff}
Soit $X$ un espace topologique. Supposons que $X$ soit quasi-compact et séparé.
Pour tout $x \in X$, la composante connexe de $X$ contenant
$x$ est l'intersection de toutes les parties à la fois ouvertes et fermées
de $X$ contenant $x$.
\end{lemma}

\begin{proof}
Soit $T$ la composante connexe contenant $x$.
Soit $S = \bigcap_{\alpha \in A} Z_\alpha$ l'intersection de toutes les
parties $Z_\alpha$ de $X$, à la fois ouvertes et fermées, contenant $x$.
Remarquons que $S$ est fermé dans $X$.
Remarquons que toute intersection finie des $Z_\alpha$ est l'un des $Z_\alpha$.
Puisque $T$ est connexe et que $x \in T$, nous avons $T \subset S$.
Il suffit de montrer que $S$ est connexe.
Sinon, il existe une décomposition en réunion disjointe
$S = B \amalg C$ où $B$ et $C$ sont à la fois ouverts et fermés dans $S$.
En particulier, $B$ et $C$ sont fermés dans $X$, et donc quasi-compacts d'après le
lemme \ref{topology-lemma-closed-in-quasi-compact}.
D'après le lemme \ref{topology-lemma-quasi-compact-in-Hausdorff},
il existe des ouverts disjoints $U, V \subset X$ tels que $B \subset U$ et
$C \subset V$. Alors $X \setminus U \cup V$ est fermé dans $X$,
donc quasi-compact (lemme \ref{topology-lemma-closed-in-quasi-compact}).
Il s'ensuit que $(X \setminus U \cup V) \cap Z_\alpha = \emptyset$
pour un certain $\alpha$, d'après le lemme \ref{topology-lemma-intersection-closed-in-quasi-compact}.
Autrement dit, $Z_\alpha \subset U \cup V$. Ainsi,
$Z_\alpha = Z_\alpha \cap V \amalg Z_\alpha \cap U$
est une décomposition en deux parties ouvertes,
donc $U \cap Z_\alpha$ et $V \cap Z_\alpha$ sont ouverts et fermés dans $X$.
Ainsi, si $x \in B$ par exemple, alors nous voyons que $S \subset U \cap Z_\alpha$,
et nous concluons que $C = \emptyset$.
\end{proof}

\begin{lemma}
\label{topology-lemma-closed-union-connected-components}
Soit $X$ un espace topologique.
Supposons que
\begin{enumerate}
\item $X$ soit quasi-compact,
\item $X$ possède une base de la topologie formée d'ouverts quasi-compacts, et
\item l'intersection de deux ouverts quasi-compacts soit quasi-compacte.
\end{enumerate}
Pour une partie $T \subset X$, les assertions suivantes sont équivalentes :
\begin{enumerate}
\item[(a)] $T$ est une intersection de parties à la fois ouvertes et fermées de $X$, et
\item[(b)] $T$ est fermée dans $X$ et est une réunion de composantes connexes de $X$.
\end{enumerate}
\end{lemma}

\begin{proof}
Il est clair que (a) implique (b).
Supposons (b). Soit $x \in X$, $x \not \in T$. Soit $x \in C \subset X$
la composante connexe de $X$ contenant $x$. D'après le
lemme \ref{topology-lemma-connected-component-intersection},
nous voyons que $C = \bigcap V_\alpha$ est l'intersection de toutes les parties à la fois ouvertes et
fermées $V_\alpha$ de $X$ qui contiennent $C$.
En particulier, toute intersection $V_\alpha \cap V_\beta$ de deux de ces parties
est elle-même un $V_\alpha$.
Comme $T$ est une réunion de composantes connexes
de $X$, nous voyons que $C \cap T = \emptyset$. Ainsi,
$T \cap \bigcap V_\alpha = \emptyset$. Puisque $T$ est quasi-compacte en tant que
partie fermée d'un espace quasi-compact (voir le
lemme \ref{topology-lemma-closed-in-quasi-compact}),
nous en déduisons que $T \cap V_\alpha = \emptyset$ pour un certain $\alpha$, voir le
lemme \ref{topology-lemma-intersection-closed-in-quasi-compact}.
Pour cet $\alpha$, nous voyons que $U_\alpha = X \setminus V_\alpha$
est une partie à la fois ouverte et fermée de $X$ qui contient $T$ mais pas $x$.
Le lemme en découle.
\end{proof}

\begin{lemma}
\label{topology-lemma-Noetherian-quasi-compact}
Soit $X$ un espace topologique noethérien.
\begin{enumerate}
\item L'espace $X$ est quasi-compact.
\item Toute partie de $X$ est rétrocompacte.
\end{enumerate}
\end{lemma}

\begin{proof}
Supposons que $X = \bigcup U_i$ soit un recouvrement ouvert de $X$ indexé
par l'ensemble $I$ et n'admettant pas de raffinement par un
recouvrement ouvert fini. Choisissons $i_1, i_2, \ldots $ dans $I$ par récurrence
de la manière suivante : choisissons $i_{n + 1}$ tel que $U_{i_{n + 1}}$
ne soit pas contenu dans $U_{i_1} \cup \ldots \cup U_{i_n}$. Ainsi, nous voyons que
$X \supset (X \setminus U_{i_1}) \supset
(X \setminus U_{i_1} \cup U_{i_2}) \supset \ldots$ est une suite infinie strictement
décroissante de parties fermées. Cela contredit
le fait que $X$ est noethérien. Cela démontre la première assertion.
La seconde assertion est maintenant claire, puisque toute partie de $X$ est noethérienne d'après le
lemme \ref{topology-lemma-Noetherian}.
\end{proof}

\begin{lemma}
\label{topology-lemma-quasi-compact-locally-Noetherian-Noetherian}
Un espace quasi-compact et localement noethérien est noethérien.
\end{lemma}

\begin{proof}
Les hypothèses impliquent immédiatement que $X$ admet un recouvrement fini par
des parties noethériennes, et est donc noethérien d'après le
lemme \ref{topology-lemma-finite-union-Noetherian}.
\end{proof}

\begin{lemma}[Théorème de la sous-base d'Alexander]
\label{topology-lemma-subbase-theorem}
Soit $X$ un espace topologique. Soit $\mathcal{B}$ une sous-base de $X$.
Si tout recouvrement de $X$ par des éléments de $\mathcal{B}$ admet un
raffinement fini, alors $X$ est quasi-compact.
\end{lemma}

\begin{proof}
Supposons qu'il existe un recouvrement ouvert de $X$ qui n'admette aucun
raffinement fini. À l'aide du lemme de Zorn, nous pouvons choisir un recouvrement ouvert maximal
$X = \bigcup_{i \in I} U_i$ qui n'admette aucun raffinement fini
(détails omis).
Autrement dit, si $U \subset X$ est un ouvert quelconque qui ne figure pas
parmi les $U_i$, alors le recouvrement $X = U \cup \bigcup_{i \in I} U_i$
admet un raffinement fini. Soit $I' \subset I$ l'ensemble des indices
tels que $U_i \in \mathcal{B}$. Alors $\bigcup_{i \in I'} U_i \not = X$,
car sinon nous obtiendrions un raffinement fini recouvrant $X$ en vertu de notre
hypothèse sur $\mathcal{B}$. Choisissons $x \in X$,
$x \not \in \bigcup_{i \in I'} U_i$. Choisissons $i \in I$ tel que $x \in U_i$.
Choisissons $V_1, \ldots, V_n \in \mathcal{B}$ tels que
$x \in V_1 \cap \ldots \cap V_n \subset U_i$. Cela est
possible puisque $\mathcal{B}$ est une sous-base de $X$. Remarquons que
$V_j$ ne figure pas parmi les $U_i$. Par maximalité du
recouvrement choisi, nous voyons que pour tout $j$, il existe
$i_{j, 1}, \ldots, i_{j, n_j} \in I$ tels que
$X = V_j \cup U_{i_{j, 1}} \cup \ldots \cup U_{i_{j, n_j}}$.
Puisque $V_1 \cap \ldots \cap V_n \subset U_i$, nous concluons que
$X = U_i \cup \bigcup U_{i_{j, l}}$, ce qui est une contradiction.
\end{proof}







\section{Espaces localement quasi-compacts}
\label{topology-section-locally-quasi-compact}

\noindent
Rappelons qu'un voisinage d'un point n'est pas nécessairement ouvert.

\begin{definition}
\label{topology-definition-locally-quasi-compact}
Un espace topologique $X$ est dit {\it localement
quasi-compact}\footnote{Cette terminologie n'est peut-être pas standard.
On rencontre dans la littérature les deux autres notions suivantes : (1) tout
point possède un voisinage quasi-compact ; (2) tout point possède un voisinage
fermé quasi-compact. Un schéma a la propriété que tout point possède un système
fondamental de voisinages ouverts quasi-compacts.} si tout point possède un
système fondamental de voisinages quasi-compacts.
\end{definition}

\noindent
Dans la littérature, l'expression {\it espace localement compact} désigne
souvent un espace comme dans le lemme suivant.

\begin{lemma}
\label{topology-lemma-locally-quasi-compact-Hausdorff}
Un espace séparé est localement quasi-compact si et seulement si tout point
possède un voisinage quasi-compact.
\end{lemma}

\begin{proof}
Soit $X$ un espace séparé. Soit $x \in X$ et soit
$x \in E \subset X$ un voisinage quasi-compact. Alors $E$ est fermé d'après le
lemme \ref{topology-lemma-quasi-compact-in-Hausdorff}.
Supposons que $x \in U \subset X$ soit un voisinage ouvert de $x$.
Alors $Z = E \setminus U$ est une partie fermée de $E$ qui ne contient pas
$x$. Il existe donc deux ouverts disjoints $W, V \subset E$ de $E$ tels que
$x \in V$ et $Z \subset W$ ; voir le
lemme \ref{topology-lemma-quasi-compact-in-Hausdorff}.
Il s'ensuit que $\overline{V} \subset E$ est un voisinage fermé de $x$ contenu
dans $E \cap U$. De plus, $\overline{V}$ est quasi-compact en tant que partie
fermée de $E$ (lemme \ref{topology-lemma-closed-in-quasi-compact}).
Nous obtenons ainsi un système fondamental de voisinages quasi-compacts de $x$.
\end{proof}

\begin{lemma}[Théorème de Baire]
\label{topology-lemma-baire-category-locally-compact}
Soit $X$ un espace séparé localement quasi-compact.
Soient $U_n \subset X$, $n \geq 1$, des ouverts denses. Alors
$\bigcap_{n \geq 1} U_n$ est dense dans $X$.
\end{lemma}

\begin{proof}
En remplaçant $U_n$ par $\bigcap_{i = 1, \ldots, n} U_i$, nous pouvons
supposer que $U_1 \supset U_2 \supset \ldots$.
Soit $x \in X$. Montrons que $x$ appartient à l'adhérence de
$\bigcap_{n \geq 1} U_n$. Soit donc $E$ un voisinage de $x$.
Pour montrer que $E \cap \bigcap_{n \geq 1} U_n$ est non vide, nous pouvons
remplacer $E$ par un voisinage plus petit. Après avoir remplacé $E$ par un
voisinage plus petit, nous pouvons supposer que $E$ est quasi-compact.

\medskip\noindent
Posons $x_0 = x$ et $E_0 = E$. Nous choisirons par récurrence un point
$x_i \in E_{i - 1} \cap U_i$ et un voisinage quasi-compact $E_i$ de $x_i$ tel
que $E_i \subset E_{i - 1} \cap U_i$.
Comme $X$ est séparé, les parties $E_i \subset X$ sont fermées
(lemme \ref{topology-lemma-quasi-compact-in-Hausdorff}).
Puisque les $E_i$ sont aussi non vides, nous concluons que
$\bigcap_{i \geq 1} E_i$ est non vide
(lemme \ref{topology-lemma-intersection-closed-in-quasi-compact}).
Comme $\bigcap_{i \geq 1} E_i \subset E \cap \bigcap_{n \geq 1} U_n$,
le lemme est démontré.

\medskip\noindent
Le cas initial $i = 0$ a été traité ci-dessus. Passons à l'étape de
récurrence. Puisque $E_{i - 1}$ est un voisinage de $x_{i - 1}$, il existe un
ouvert $x_{i - 1} \in W \subset E_{i - 1}$.
Comme $U_i$ est dense dans $X$, l'intersection $W \cap U_i$ est non vide.
Choisissons un point $x_i \in W \cap U_i$ quelconque.
Par définition des espaces localement quasi-compacts, il existe un voisinage
quasi-compact $E_i$ de $x_i$ contenu dans $W \cap U_i$.
Alors $E_i \subset E_{i - 1} \cap U_i$, comme voulu.
\end{proof}

\begin{lemma}
\label{topology-lemma-relatively-compact-refinement}
Soit $X$ un espace séparé et quasi-compact.
Soit $X = \bigcup_{i \in I} U_i$ un recouvrement ouvert.
Il existe alors un recouvrement ouvert $X = \bigcup_{i \in I} V_i$ tel que
$\overline{V_i} \subset U_i$ pour tout $i$.
\end{lemma}

\begin{proof}
Soit $x \in X$. Choisissons $i(x) \in I$ tel que $x \in U_{i(x)}$.
Comme $X \setminus U_{i(x)}$ et $\{x\}$ sont deux parties fermées disjointes
de $X$, les lemmes \ref{topology-lemma-closed-in-quasi-compact} et
\ref{topology-lemma-quasi-compact-in-Hausdorff} fournissent un voisinage ouvert $U_x$
de $x$ dont l'adhérence ne rencontre pas $X \setminus U_{i(x)}$.
Ainsi, $\overline{U_x} \subset U_{i(x)}$. Comme $X$ est quasi-compact, il
existe une liste finie de points $x_1, \ldots, x_m$ telle que
$X = U_{x_1} \cup \ldots \cup U_{x_m}$. En posant
$V_i = \bigcup_{i = i(x_j)} U_{x_j}$, on achève la démonstration.
\end{proof}

\begin{lemma}
\label{topology-lemma-refine-covering}
Soit $X$ un espace séparé et quasi-compact.
Soit $X = \bigcup_{i \in I} U_i$ un recouvrement ouvert.
Supposons donnés un entier $p \geq 0$ et, pour tout $(p + 1)$-uplet
$i_0, \ldots, i_p$ d'éléments de $I$, un recouvrement ouvert
$U_{i_0} \cap \ldots \cap U_{i_p} = \bigcup W_{i_0 \ldots i_p, k}$.
Il existe alors un recouvrement ouvert $X = \bigcup_{j \in J} V_j$ et une
application $\alpha : J \to I$ tels que $\overline{V_j} \subset U_{\alpha(j)}$
et que chaque $V_{j_0} \cap \ldots \cap V_{j_p}$ soit contenu dans
$W_{\alpha(j_0) \ldots \alpha(j_p), k}$ pour un certain $k$.
\end{lemma}

\begin{proof}
Comme $X$ est quasi-compact, on se ramène au cas où $I$ est fini (détails
omis). Nous démontrons le résultat pour $I$ fini par récurrence sur $p$.
Le cas initial $p = 0$ est immédiat : il suffit de prendre un recouvrement
comme dans le lemme \ref{topology-lemma-relatively-compact-refinement}, qui raffine le
recouvrement ouvert $X = \bigcup W_{i_0, k}$.

\medskip\noindent
Passons à l'étape de récurrence. Supposons le lemme démontré pour $p - 1$.
Pour tout $p$-uplet $i'_0, \ldots, i'_{p - 1}$ d'éléments de $I$, choisissons
$U_{i'_0} \cap \ldots \cap U_{i'_{p - 1}} =
\bigcup W_{i'_0 \ldots i'_{p - 1}, k}$ comme raffinement commun des
recouvrements
$U_{i_0} \cap \ldots \cap U_{i_p} = \bigcup W_{i_0 \ldots i_p, k}$ associés
aux $(p + 1)$-uplets tels que
$\{i'_0, \ldots, i'_{p - 1}\} = \{i_0, \ldots, i_p\}$ (égalité d'ensembles).
(Il n'y en a qu'un nombre fini, puisque $I$ est fini.)
Par récurrence, il existe une solution pour ces ouverts, disons
$X = \bigcup V_j$ et $\alpha : J \to I$.
À ce stade, le recouvrement $X = \bigcup_{j \in J} V_j$ et $\alpha$ vérifient
$\overline{V_j} \subset U_{\alpha(j)}$, et chaque
$V_{j_0} \cap \ldots \cap V_{j_p}$ est contenu dans
$W_{\alpha(j_0) \ldots \alpha(j_p), k}$ pour un certain $k$ dès qu'il y a une
répétition parmi $\alpha(j_0), \ldots, \alpha(j_p)$. Bien entendu, nous pouvons
et allons supposer que $J$ est fini.

\medskip\noindent
Fixons $i_0, \ldots, i_p \in I$ deux à deux distincts. Considérons les
$(p + 1)$-uplets $j_0, \ldots, j_p \in J$ tels que
$i_0 = \alpha(j_0), \ldots, i_p = \alpha(j_p)$ et pour lesquels
$V_{j_0} \cap \ldots \cap V_{j_p}$ n'est {\bf pas} contenu dans
$W_{\alpha(j_0) \ldots \alpha(j_p), k}$, quel que soit $k$.
Soit $N$ le nombre de ces $(p + 1)$-uplets. Montrons comment diminuer $N$.
Comme
$$
\overline{V_{j_0}} \cap \ldots \cap \overline{V_{j_p}} \subset
U_{i_0} \cap \ldots \cap U_{i_p} = \bigcup W_{i_0 \ldots i_p, k}
$$
il existe un ensemble fini $K = \{k_1, \ldots, k_t\}$ tel que le membre de
gauche soit contenu dans $\bigcup_{k \in K} W_{i_0 \ldots i_p, k}$.
Considérons alors le recouvrement ouvert
$$
V_{j_0} =
(V_{j_0} \setminus (\overline{V_{j_1}} \cap \ldots \cap \overline{V_{j_p}}))
\cup (\bigcup\nolimits_{k \in K} V_{j_0} \cap W_{i_0 \ldots i_p, k})
$$
Le premier ouvert du membre de droite a une intersection vide avec
$V_{j_1} \cap \ldots \cap V_{j_p}$, tandis que les
autres ouverts $V_{j_0, k}$ du membre de droite satisfont
$V_{j_0, k} \cap V_{j_1} \ldots \cap V_{j_p} \subset
W_{\alpha(j_0) \ldots \alpha(j_p), k}$.
Posons $J' = J \amalg K$. Pour $j \in J$, posons $V'_j = V_j$ si
$j \not = j_0$, et posons
$V'_{j_0} =
V_{j_0} \setminus (\overline{V_{j_1}} \cap \ldots \cap \overline{V_{j_p}})$.
Pour $k \in K$, posons $V'_k = V_{j_0, k}$. Enfin, l'application
$\alpha' : J' \to I$ est égale à $\alpha$ sur $J$ et envoie tout élément de
$K$ sur $i_0$. Une vérification immédiate montre que ce remplacement diminue
$N$ d'une unité. En répétant cette procédure $N$ fois, nous arrivons à la
situation où $N = 0$.

\medskip\noindent
Pour achever la démonstration, raisonnons par récurrence sur le nombre $M$ de
$(p + 1)$-uplets $i_0, \ldots, i_p \in I$ à entrées deux à deux distinctes
pour lesquels il existe un $(p + 1)$-uplet $j_0, \ldots, j_p \in J$ tel que
$i_0 = \alpha(j_0), \ldots, i_p = \alpha(j_p)$ et tel que
$V_{j_0} \cap \ldots \cap V_{j_p}$ ne soit {\bf pas} contenu dans
$W_{\alpha(j_0) \ldots \alpha(j_p), k}$, quel que soit $k$. À cette fin, nous
affirmons que l'opération du paragraphe précédent n'augmente pas $M$.
Cela résulte formellement du fait que l'application $\alpha' : J' \to I$ se
factorise par une application $\beta : J' \to J$ telle que
$V'_{j'} \subset V_{\beta(j')}$.
\end{proof}

\begin{lemma}
\label{topology-lemma-lift-covering-of-a-closed}
Soit $X$ un espace séparé et localement quasi-compact.
Soit $Z \subset X$ une partie quasi-compacte (donc fermée).
Supposons donnés un entier $p \geq 0$, un ensemble $I$, pour tout $i \in I$
un ouvert $U_i \subset X$, et pour tout $(p + 1)$-uplet
$i_0, \ldots, i_p$ d'éléments de $I$ un ouvert
$W_{i_0 \ldots i_p} \subset U_{i_0} \cap \ldots \cap U_{i_p}$
tel que
\begin{enumerate}
\item $Z \subset \bigcup U_i$, et
\item pour tous $i_0, \ldots, i_p$, on ait
$W_{i_0 \ldots i_p} \cap Z = U_{i_0} \cap \ldots \cap U_{i_p} \cap Z$.
\end{enumerate}
Il existe alors des ouverts $V_i$ de $X$ tels que
$Z \subset \bigcup V_i$,
pour tout $i$, on ait $\overline{V_i} \subset U_i$, et
$V_{i_0} \cap \ldots \cap V_{i_p} \subset W_{i_0 \ldots i_p}$
pour tout $(p + 1)$-uplet $i_0, \ldots, i_p$.
\end{lemma}

\begin{proof}
Puisque $Z$ est quasi-compact, on se ramène au cas où $I$ est fini (détails
omis). Comme $X$ est localement quasi-compact et $Z$ quasi-compact, nous
pouvons trouver un voisinage $Z \subset E$ qui soit quasi-compact,
c'est-à-dire que $E$ est quasi-compact et contient un voisinage ouvert de
$Z$ dans $X$. Si nous démontrons le résultat après avoir remplacé $X$ par
$E$, le résultat s'ensuit. Nous pouvons donc supposer $X$ quasi-compact.

\medskip\noindent
Nous démontrons le résultat lorsque $I$ est fini et $X$ quasi-compact, par
récurrence sur $p$. Le cas initial est $p = 0$. Dans ce cas,
$X = (X \setminus Z) \cup \bigcup W_i$. D'après le
lemme \ref{topology-lemma-relatively-compact-refinement}, il existe un recouvrement
$X = V \cup \bigcup V_i$ par des ouverts $V_i \subset W_i$ et
$V \subset X \setminus Z$, avec $\overline{V_i} \subset W_i$ pour tout $i$.
On voit alors que l'on obtient une solution du problème posé dans le lemme.

\medskip\noindent
Passons à l'étape de récurrence. Supposons le lemme démontré pour $p - 1$.
Posons $W_{j_0 \ldots j_{p - 1}}$ égal à l'intersection de tous les
$W_{i_0 \ldots i_p}$ tels que
$\{j_0, \ldots, j_{p - 1}\} = \{i_0, \ldots, i_p\}$ (égalité d'ensembles).
Par récurrence, il existe une solution pour ces ouverts, disons
$V_i \subset U_i$.
Il résulte de notre choix de $W_{j_0 \ldots j_{p - 1}}$ que
$V_{i_0} \cap \ldots \cap V_{i_p} \subset W_{i_0 \ldots i_p}$
pour tout $(p + 1)$-uplet $i_0, \ldots, i_p$ tel que $i_a = i_b$ pour
certains $0 \leq a < b \leq p$.
Il ne reste donc à modifier notre choix des $V_i$ que si
$V_{i_0} \cap \ldots \cap V_{i_p} \not \subset W_{i_0 \ldots i_p}$
pour un $(p + 1)$-uplet $i_0, \ldots, i_p$ dont les éléments sont deux à
deux distincts. Dans ce cas,
$$
T =
\overline{V_{i_0} \cap \ldots \cap V_{i_p} \setminus W_{i_0 \ldots i_p}}
\subset 
\overline{V_{i_0}} \cap \ldots \cap \overline{V_{i_p}} \setminus
W_{i_0 \ldots i_p}
$$
est un fermé de $X$, contenu dans $U_{i_0} \cap \ldots \cap U_{i_p}$ et
disjoint de $Z$. Nous pouvons donc remplacer $V_{i_0}$ par
$V_{i_0} \setminus T$ pour « corriger » le problème. Après avoir répété cette
opération un nombre fini de fois pour chacun des uplets problématiques, le
lemme est démontré.
\end{proof}

\begin{lemma}
\label{topology-lemma-lift-covering-of-quasi-compact-hausdorff-subset}
Soit $X$ un espace topologique. Soit $Z \subset X$ une partie quasi-compacte
telle que deux points quelconques de $Z$ admettent dans $X$ des voisinages
ouverts disjoints. Supposons donnés un entier $p \geq 0$, un ensemble $I$,
pour tout $i \in I$ un ouvert $U_i \subset X$, et pour tout $(p + 1)$-uplet
$i_0, \ldots, i_p$ d'éléments de $I$ un ouvert
$W_{i_0 \ldots i_p} \subset U_{i_0} \cap \ldots \cap U_{i_p}$
tel que
\begin{enumerate}
\item $Z \subset \bigcup U_i$, et
\item pour tous $i_0, \ldots, i_p$, on ait
$W_{i_0 \ldots i_p} \cap Z = U_{i_0} \cap \ldots \cap U_{i_p} \cap Z$.
\end{enumerate}
Il existe alors des ouverts $V_i$ de $X$ tels que
\begin{enumerate}
\item $Z \subset \bigcup V_i$,
\item $V_i \subset U_i$ pour tout $i$,
\item $\overline{V_i} \cap Z \subset U_i$ pour tout $i$, et
\item $V_{i_0} \cap \ldots \cap V_{i_p} \subset W_{i_0 \ldots i_p}$
pour tout $(p + 1)$-uplet $i_0, \ldots, i_p$.
\end{enumerate}
\end{lemma}

\begin{proof}
Puisque $Z$ est quasi-compact, on se ramène au cas où $I$ est fini (détails
omis). Nous démontrons le résultat lorsque $I$ est fini, par récurrence sur
$p$.

\medskip\noindent
Le cas initial est $p = 0$.
Pour $z \in Z \cap U_i$ et $z' \in Z \setminus U_i$, il existe des ouverts
disjoints $z \in V_{z, z'}$ et $z' \in W_{z, z'}$ de $X$.
Comme $Z \setminus U_i$ est quasi-compact
(lemme \ref{topology-lemma-closed-in-quasi-compact}), nous pouvons choisir un nombre
fini de points $z'_1, \ldots, z'_r$ tels que
$Z \setminus U_i \subset W_{z, z'_1} \cup \ldots \cup W_{z, z'_r}$.
On voit alors que
$V_z = V_{z, z'_1} \cap \ldots \cap V_{z, z'_r} \cap U_i$
est un voisinage ouvert de $z$ contenu dans $U_i$ et vérifiant
$\overline{V_z} \cap Z \subset U_i$.
Comme $z$ et $i$ étaient arbitraires et que $Z$ est quasi-compact, nous
pouvons trouver une liste finie $z_1, i_1, \ldots, z_t, i_t$ et des ouverts
$V_{z_j} \subset U_{i_j}$ tels que
$\overline{V_{z_j}} \cap Z \subset U_{i_j}$
et $Z \subset \bigcup V_{z_j}$.
Nous pouvons alors poser $V_i = W_i \cap (\bigcup_{j : i = i_j} V_{z_j})$
pour résoudre le problème lorsque $p = 0$.

\medskip\noindent
Passons à l'étape de récurrence. Supposons le lemme démontré pour $p - 1$.
Posons $W_{j_0 \ldots j_{p - 1}}$ égal à l'intersection de tous les
$W_{i_0 \ldots i_p}$ tels que
$\{j_0, \ldots, j_{p - 1}\} = \{i_0, \ldots, i_p\}$ (égalité d'ensembles).
Par récurrence, il existe une solution pour ces ouverts, disons
$V_i \subset U_i$.
Il résulte de notre choix de $W_{j_0 \ldots j_{p - 1}}$ que
$V_{i_0} \cap \ldots \cap V_{i_p} \subset W_{i_0 \ldots i_p}$
pour tout $(p + 1)$-uplet $i_0, \ldots, i_p$ tel que $i_a = i_b$ pour
certains $0 \leq a < b \leq p$.
Il ne reste donc à modifier notre choix des $V_i$ que si
$V_{i_0} \cap \ldots \cap V_{i_p} \not \subset W_{i_0 \ldots i_p}$
pour un $(p + 1)$-uplet $i_0, \ldots, i_p$ dont les éléments sont deux à
deux distincts. Dans ce cas,
$$
T =
\overline{V_{i_0} \cap \ldots \cap V_{i_p} \setminus W_{i_0 \ldots i_p}}
\subset
\overline{V_{i_0}} \cap \ldots \cap \overline{V_{i_p}} \setminus
W_{i_0 \ldots i_p}
$$
est un fermé de $X$ disjoint de $Z$, d'après la propriété (3) des ouverts
$V_i$. Nous pouvons donc remplacer $V_{i_0}$ par $V_{i_0} \setminus T$ pour
« corriger » le problème. Après avoir répété cette opération un nombre fini
de fois pour chacun des uplets problématiques, le lemme est démontré.
\end{proof}







\section{Limites d'espaces}
\label{topology-section-limits}

\noindent
La catégorie des espaces topologiques admet les produits. Plus précisément,
si $I$ est un ensemble et si, pour $i \in I$, on se donne un espace
topologique $X_i$, on munit $\prod_{i \in I} X_i$ de la {\it topologie
produit}. On prend pour base de cette topologie les ensembles de la forme
$\prod U_i$, où $U_i \subset X_i$ est ouvert et $U_i = X_i$ pour presque
tout $i$.

\medskip\noindent
La catégorie des espaces topologiques admet les égalisateurs. Plus
précisément, si $a, b : X \to Y$ sont des morphismes d'espaces topologiques,
l'égalisateur de $a$ et $b$ est la partie
$\{x \in X \mid a(x) = b(x)\} \subset X$, munie de la topologie induite.

\begin{lemma}
\label{topology-lemma-limits}
La catégorie des espaces topologiques admet les limites, et le foncteur
d'oubli vers les ensembles commute à celles-ci.
\end{lemma}

\begin{proof}
Cela résulte de la discussion précédente et de Catégories,
lemme \ref{categories-lemma-limits-products-equalizers}.
La description précédente montre que le foncteur d'oubli commute aux
limites. On peut aussi le voir en appliquant Catégories,
lemme \ref{categories-lemma-adjoint-exact}, et en observant que le foncteur
d'oubli admet un adjoint à gauche, à savoir le foncteur qui associe à un
ensemble l'espace topologique discret correspondant.
\end{proof}

\begin{lemma}
\label{topology-lemma-describe-limits}
Soit $\mathcal{I}$ une catégorie cofiltrante. Soit $i \mapsto X_i$ un
diagramme d'espaces topologiques indexé par $\mathcal{I}$. Soit
$X = \lim X_i$ sa limite, avec applications de projection $f_i : X \to X_i$.
\begin{enumerate}
\item Tout ouvert de $X$ est de la forme $\bigcup_{j \in J} f_j^{-1}(U_j)$
pour une partie $J \subset I$ et des ouverts $U_j \subset X_j$.
\item Tout ouvert quasi-compact de $X$ est de la forme
$f_i^{-1}(U_i)$ pour un certain $i$ et un ouvert $U_i \subset X_i$.
\end{enumerate}
\end{lemma}

\begin{proof}
La construction de la limite donnée ci-dessus montre que
$X \subset \prod X_i$, avec la topologie induite. Une base de la topologie de
$\prod X_i$ est donnée par les ouverts $\prod U_i$, où $U_i \subset X_i$ est
ouvert et $U_i = X_i$ pour presque tout $i$. Soient
$i_1, \ldots, i_n \in \Ob(\mathcal{I})$ les objets tels que
$U_{i_j} \not = X_{i_j}$. Alors
$$
X \cap \prod U_i = f_{i_1}^{-1}(U_{i_1}) \cap \ldots \cap f_{i_n}^{-1}(U_{i_n})
$$
Dans le cas d'une limite quelconque d'espaces topologiques, ces ouverts
forment une base de la topologie de $X$. Toutefois, si $\mathcal{I}$ est
cofiltrante comme dans l'énoncé du lemme, nous pouvons choisir un
$j \in \Ob(\mathcal{I})$ et des morphismes $j \to i_l$,
$l = 1, \ldots, n$. Posons
$$
U_j =
(X_j \to X_{i_1})^{-1}(U_{i_1}) \cap \ldots \cap
(X_j \to X_{i_n})^{-1}(U_{i_n})
$$
Il est alors clair que $X \cap \prod U_i = f_j^{-1}(U_j)$. Ainsi, pour tout
ouvert $W$ de $X$, il existe un ensemble $A$, une application
$\alpha : A \to \Ob(\mathcal{I})$ et des ouverts
$U_a \subset X_{\alpha(a)}$ tels que
$W = \bigcup f_{\alpha(a)}^{-1}(U_a)$. Posons $J = \Im(\alpha)$ et, pour
$j \in J$, posons $U_j = \bigcup_{\alpha(a) = j} U_a$ ; on obtient alors
$W = \bigcup_{j \in J} f_j^{-1}(U_j)$.
Ceci démontre (1).

\medskip\noindent
Pour démontrer (2), supposons que
$\bigcup_{j \in J} f_j^{-1}(U_j)$ soit quasi-compact. Cet ouvert est alors
égal à $f_{j_1}^{-1}(U_{j_1}) \cup \ldots \cup f_{j_m}^{-1}(U_{j_m})$
pour certains $j_1, \ldots, j_m \in J$. Puisque $\mathcal{I}$ est
cofiltrante, nous pouvons choisir un $i \in \Ob(\mathcal{I})$ et des
morphismes $i \to j_l$, $l = 1, \ldots, m$. Posons
$$
U_i =
(X_i \to X_{j_1})^{-1}(U_{j_1}) \cup \ldots \cup
(X_i \to X_{j_m})^{-1}(U_{j_m})
$$
Notre ouvert est alors égal à $f_i^{-1}(U_i)$, comme voulu.
\end{proof}

\begin{lemma}
\label{topology-lemma-characterize-limit}
Soit $\mathcal{I}$ une catégorie cofiltrante. Soit $i \mapsto X_i$ un
diagramme d'espaces topologiques indexé par $\mathcal{I}$. Soit $X$ un espace
topologique tel que
\begin{enumerate}
\item $X = \lim X_i$ comme ensemble (on note $f_i$ les applications de
projection),
\item les ensembles $f_i^{-1}(U_i)$, pour $i \in \Ob(\mathcal{I})$ et
$U_i \subset X_i$ ouvert, forment une base de la topologie de $X$.
\end{enumerate}
Alors $X$ est la limite des $X_i$ dans la catégorie des espaces topologiques.
\end{lemma}

\begin{proof}
Cela résulte de la description de la topologie limite donnée dans le
lemme \ref{topology-lemma-describe-limits}.
\end{proof}

\begin{theorem}[Tychonov]
\label{topology-theorem-tychonov}
Tout produit d'espaces quasi-compacts est quasi-compact.
\end{theorem}

\begin{proof}
Soit $I$ un ensemble et, pour $i \in I$, soit $X_i$ un espace topologique
quasi-compact. Posons $X = \prod X_i$. Soit $\mathcal{B}$ l'ensemble des
parties de $X$ de la forme
$U_i \times \prod_{i' \in I, i' \not = i} X_{i'}$, où
$U_i \subset X_i$ est ouvert. Par construction, cette famille est une
sous-base de la topologie de $X$. D'après le
lemme \ref{topology-lemma-subbase-theorem}, il suffit de montrer que tout recouvrement
$X = \bigcup_{j \in J} B_j$ par des éléments $B_j$ de $\mathcal{B}$ admet un
sous-recouvrement fini. Nous pouvons décomposer $J = \coprod J_i$ de telle
sorte que, si $j \in J_i$, alors
$B_j = U_j \times \prod_{i' \not = i} X_{i'}$, avec $U_j \subset X_i$
ouvert. Si $X_i = \bigcup_{j \in J_i} U_j$, il existe un sous-recouvrement
fini et nous concluons que $X = \bigcup_{j \in J} B_j$ admet un
sous-recouvrement fini. Si ce n'est pas le cas, nous pouvons choisir, pour
tout $i$, un point $x_i \in X_i$ qui n'appartient pas à
$\bigcup_{j \in J_i} U_j$. Mais alors le point $x = (x_i)_{i \in I}$ est un
élément de $X$ qui n'appartient pas à $\bigcup_{j \in J} B_j$, contradiction.
\end{proof}

\noindent
Le lemme suivant devient faux si l'on omet l'hypothèse que les espaces $X_i$
sont séparés ; voir Exemples,
section \ref{examples-section-lim-not-quasi-compact}.

\begin{lemma}
\label{topology-lemma-inverse-limit-quasi-compact}
Soit $\mathcal{I}$ une catégorie et soit $i \mapsto X_i$ un diagramme indexé
par $\mathcal{I}$ dans la catégorie des espaces topologiques. Si chaque $X_i$
est quasi-compact et séparé, alors $\lim X_i$ est quasi-compact.
\end{lemma}

\begin{proof}
Rappelons que $\lim X_i$ est un sous-espace de $\prod X_i$. D'après le
théorème \ref{topology-theorem-tychonov}, ce produit est quasi-compact. Il suffit donc
de montrer que $\lim X_i$ est un sous-espace fermé de $\prod X_i$
(lemme \ref{topology-lemma-closed-in-quasi-compact}).
Si $\varphi : j \to k$ est un morphisme de $\mathcal{I}$, notons
$\Gamma_\varphi \subset X_j \times X_k$ le graphe de l'application continue
correspondante $X_j \to X_k$. D'après le lemme \ref{topology-lemma-graph-closed}, ce
graphe est fermé. Il est clair que $\lim X_i$ est l'intersection des fermés
$$
\Gamma_\varphi \times \prod\nolimits_{l \not = j, k} X_l
\subset \prod X_i
$$
Le résultat s'ensuit.
\end{proof}

\noindent
Le lemme suivant généralise Catégories,
lemme \ref{categories-lemma-nonempty-limit}, et généralise partiellement le
lemme \ref{topology-lemma-intersection-closed-in-quasi-compact}.

\begin{lemma}
\label{topology-lemma-nonempty-limit}
Soit $\mathcal{I}$ une catégorie cofiltrante et soit $i \mapsto X_i$ un
diagramme indexé par $\mathcal{I}$ dans la catégorie des espaces
topologiques. Si chaque $X_i$ est quasi-compact, séparé et non vide, alors
$\lim X_i$ est non vide.
\end{lemma}

\begin{proof}
Dans la démonstration du lemme \ref{topology-lemma-inverse-limit-quasi-compact}, nous
avons vu que $X = \lim X_i$ est l'intersection des fermés
$$
Z_\varphi = \Gamma_\varphi \times \prod\nolimits_{l \not = j, k} X_l
$$
dans l'espace quasi-compact $\prod X_i$, où $\varphi : j \to k$ est un
morphisme de $\mathcal{I}$ et où $\Gamma_\varphi \subset X_j \times X_k$
est le graphe du morphisme correspondant $X_j \to X_k$. D'après le
lemme \ref{topology-lemma-intersection-closed-in-quasi-compact}, il suffit donc de
montrer que toute intersection finie de ces parties est non vide.
Soit $\varphi_t : j_t \to k_t$, $t = 1, \ldots, n$, une famille finie de
morphismes de $\mathcal{I}$. Puisque $\mathcal{I}$ est cofiltrante, nous
pouvons choisir un objet $j$ et un morphisme $\psi_t : j \to j_t$ pour tout
$t$. Pour chaque paire $t, t'$ telle que (a) $j_t = j_{t'}$, ou
(b) $j_t = k_{t'}$, ou (c) $k_t = k_{t'}$, nous obtenons deux morphismes
$j \to l$, avec $l = j_t$ dans les cas (a) et (b), ou $l = k_t$ dans le cas
(c). Comme $\mathcal{I}$ est cofiltrante et qu'il n'y a qu'un nombre fini de
paires $(t, t')$, nous pouvons choisir une application $j' \to j$ qui égalise
ces deux morphismes pour toutes ces paires $(t, t')$. Choisissons un élément
$x \in X_{j'}$ et, pour tout $t$, notons $x_{j_t}$, respectivement
$x_{k_t}$, l'image de $x$ par le morphisme $X_{j'} \to X_j \to X_{j_t}$,
respectivement $X_{j'} \to X_j \to X_{j_t} \to X_{k_t}$.
Pour tout indice $l \in \Ob(\mathcal{I})$ qui n'est égal ni à $j_t$ ni à
$k_t$ pour aucun $t$, choisissons un élément arbitraire $x_l \in X_l$ (en
utilisant l'axiome du choix). Alors $(x_i)_{i \in \Ob(\mathcal{I})}$ appartient, par
construction, à l'intersection
$$
Z_{\varphi_1} \cap \ldots \cap Z_{\varphi_n}
$$
et la démonstration est achevée.
\end{proof}


\section{Parties constructibles}
\label{topology-section-constructible}

\begin{definition}
\label{topology-definition-constructible}
Soit $X$ un espace topologique. Soit $E \subset X$ une partie de $X$.
\begin{enumerate}
\item On dit que $E$ est {\it constructible}\footnote{Dans la deuxième édition
d'EGA I \cite{EGA1-second}, cette partie était appelée « globalement
constructible » et la terminologie « constructible » était employée pour ce
que nous appelons une partie localement constructible.}
dans $X$ si $E$ est une réunion finie
de parties de la forme $U \cap V^c$, où $U, V \subset X$ sont des ouverts
rétrocompacts dans $X$.
\item On dit que $E$ est {\it localement constructible} dans $X$ s'il existe un
recouvrement ouvert $X = \bigcup V_i$ tel que chaque $E \cap V_i$ soit
constructible dans $V_i$.
\end{enumerate}
\end{definition}

\begin{lemma}
\label{topology-lemma-constructible}
La collection des parties constructibles est stable par
intersections finies, unions finies et passage au complémentaire.
\end{lemma}

\begin{proof}
Remarquons que si $U_1$, $U_2$ sont des ouverts rétrocompacts dans $X$,
alors $U_1 \cup U_2$ l'est aussi, car la réunion de deux parties
quasi-compactes de $X$ est quasi-compacte. Il est également vrai que
$U_1 \cap U_2$ est rétrocompact. En effet, supposons que $U \subset X$
soit un ouvert quasi-compact ; alors $U_2 \cap U$ est quasi-compact puisque
$U_2$ est rétrocompact dans $X$, et nous en déduisons que
$U_1 \cap (U_2 \cap U)$ est quasi-compact puisque $U_1$ est
rétrocompact dans $X$. À partir de là, il est formel de montrer que
le complémentaire d'une partie constructible est constructible,
que les unions finies de parties constructibles sont constructibles et
que les intersections finies de parties constructibles sont constructibles.
\end{proof}

\begin{lemma}
\label{topology-lemma-inverse-images-constructibles}
Soit $f : X \to Y$ une application continue d'espaces topologiques.
Si l'image réciproque de tout ouvert rétrocompact de $Y$
est rétrocompacte dans $X$, alors les images réciproques des parties
constructibles sont constructibles.
\end{lemma}

\begin{proof}
Cette assertion résulte de $f^{-1}(U \cap V^c) = f^{-1}(U) \cap f^{-1}(V)^c$
et de la définition des parties constructibles.
\end{proof}

\begin{lemma}
\label{topology-lemma-open-immersion-constructible-inverse-image}
Soit $U \subset X$ un ouvert. Pour une partie constructible
$E \subset X$, l'intersection $E \cap U$ est constructible
dans $U$.
\end{lemma}

\begin{proof}
Supposons que $V \subset X$ soit un ouvert rétrocompact dans $X$.
Il suffit de montrer que $V \cap U$ est rétrocompact dans $U$
d'après le lemme \ref{topology-lemma-inverse-images-constructibles}. Pour cela,
soit $W \subset U$ un ouvert quasi-compact. Alors $W$
est ouvert et quasi-compact dans $X$. Par conséquent, $V \cap W = V \cap U \cap W$
est quasi-compact puisque $V$ est rétrocompact dans $X$.
\end{proof}

\begin{lemma}
\label{topology-lemma-quasi-compact-open-immersion-constructible-image}
Soit $U \subset X$ un ouvert rétrocompact. Soit $E \subset U$.
Si $E$ est constructible dans $U$, alors $E$ est constructible dans $X$.
\end{lemma}

\begin{proof}
Supposons que $V, W \subset U$ soient des ouverts rétrocompacts dans $U$.
Alors $V, W$ sont des ouverts rétrocompacts dans $X$
(lemme \ref{topology-lemma-composition-quasi-compact}).
Ainsi, $V \cap (U \setminus W) = V \cap (X \setminus W)$
est constructible dans $X$. On conclut puisque toute partie constructible de $U$
est une réunion finie de parties de la forme $V \cap (U \setminus W)$.
\end{proof}

\begin{lemma}
\label{topology-lemma-collate-constructible}
Soit $X$ un espace topologique. Soit $E \subset X$ une partie.
Soit $X = V_1 \cup \ldots \cup V_m$ un recouvrement fini par des
ouverts rétrocompacts.
Alors $E$ est constructible dans $X$ si et seulement si $E \cap V_j$
est constructible dans $V_j$ pour chaque $j = 1, \ldots, m$.
\end{lemma}

\begin{proof}
Si $E$ est constructible dans $X$, alors le
lemme \ref{topology-lemma-open-immersion-constructible-inverse-image}
montre que $E \cap V_j$ est constructible dans $V_j$ pour tout $j$.
Réciproquement, supposons que $E \cap V_j$
soit constructible dans $V_j$ pour chaque $j = 1, \ldots, m$.
Alors $E = \bigcup E \cap V_j$ est une réunion finie de
parties constructibles d'après le
lemme \ref{topology-lemma-quasi-compact-open-immersion-constructible-image},
et est donc constructible.
\end{proof}

\begin{lemma}
\label{topology-lemma-intersect-constructible-with-closed}
Soit $X$ un espace topologique. Soit $Z \subset X$ une partie fermée
telle que $X \setminus Z$ soit quasi-compact.
Alors, pour toute partie constructible $E \subset X$, l'intersection
$E \cap Z$ est constructible dans $Z$.
\end{lemma}

\begin{proof}
Supposons que $V \subset X$ soit un ouvert rétrocompact dans $X$.
Il suffit de montrer que $V \cap Z$ est rétrocompact dans $Z$
d'après le lemme \ref{topology-lemma-inverse-images-constructibles}. Pour cela,
soit $W \subset Z$ un ouvert quasi-compact. La partie
$W' = W \cup (X \setminus Z)$ est quasi-compacte, ouverte, et $W = Z \cap W'$.
Par conséquent, $V \cap Z \cap W = V \cap Z \cap W'$
est une partie fermée de l'ouvert quasi-compact $V \cap W'$
puisque $V$ est rétrocompact dans $X$. Ainsi, $V \cap Z \cap W$ est quasi-compact
d'après le lemme \ref{topology-lemma-closed-in-quasi-compact}.
\end{proof}

\begin{lemma}
\label{topology-lemma-intersect-constructible-with-retrocompact}
Soit $X$ un espace topologique. Soit $T \subset X$ une partie. Supposons que
\begin{enumerate}
\item $T$ soit rétrocompacte dans $X$,
\item les ouverts quasi-compacts forment une base de la topologie de $X$.
\end{enumerate}
Alors, pour toute partie constructible $E \subset X$, l'intersection $E \cap T$ est
constructible dans $T$.
\end{lemma}

\begin{proof}
Supposons que $V \subset X$ soit un ouvert rétrocompact dans $X$.
Il suffit de montrer que $V \cap T$ est rétrocompact dans $T$
d'après le lemme \ref{topology-lemma-inverse-images-constructibles}. Pour cela,
soit $W \subset T$ un ouvert quasi-compact. Par l'hypothèse (2),
nous pouvons trouver un ouvert quasi-compact $W' \subset X$
tel que $W = T \cap W'$ (détails omis).
Ainsi, $V \cap T \cap W = V \cap T \cap W'$
est l'intersection de $T$ avec  l'ouvert quasi-compact $V \cap W'$
puisque $V$ est rétrocompact dans $X$. Ainsi, $V \cap T \cap W$ est quasi-compact.
\end{proof}

\begin{lemma}
\label{topology-lemma-closed-constructible-image}
Soit $Z \subset X$ une partie fermée dont le complémentaire est un ouvert rétrocompact.
Soit $E \subset Z$. Si $E$ est constructible dans $Z$, alors $E$ est constructible
dans $X$.
\end{lemma}

\begin{proof}
Supposons que $V \subset Z$ soit un ouvert rétrocompact dans $Z$. Considérons l'ouvert
$\tilde V = V \cup (X \setminus Z)$ de $X$. Soit $W \subset X$ un
ouvert quasi-compact. Alors
$$
W \cap \tilde V =
\left(V \cap W\right) \cup \left((X \setminus Z) \cap W\right).
$$
La première partie est quasi-compacte, car $V \cap W = V \cap (Z \cap W)$ et
$(Z \cap W)$ est un ouvert quasi-compact dans $Z$
(lemme \ref{topology-lemma-closed-in-quasi-compact}), et $V$ est rétrocompact dans $Z$.
La seconde partie est quasi-compacte puisque $(X \setminus Z)$ est rétrocompact dans $X$.
Nous voyons ainsi que $\tilde V$ est rétrocompact dans $X$.
Ainsi, si $V_1, V_2 \subset Z$ sont des ouverts rétrocompacts, alors
$$
V_1 \cap (Z \setminus V_2) = \tilde V_1 \cap (X \setminus \tilde V_2)
$$
est constructible dans $X$. On conclut puisque toute partie constructible de $Z$
est une réunion finie de parties de la forme $V_1 \cap (Z \setminus V_2)$.
\end{proof}

\begin{lemma}
\label{topology-lemma-constructible-is-retrocompact}
Soit $X$ un espace topologique. Toute partie constructible
de $X$ est rétrocompacte.
\end{lemma}

\begin{proof}
Soit $E = \bigcup_{i = 1, \ldots, n} U_i \cap V_i^c$, où $U_i, V_i$ sont des
ouverts rétrocompacts dans $X$. Soit $W \subset X$ un ouvert quasi-compact.
Alors $E \cap W = \bigcup_{i = 1, \ldots, n} U_i \cap V_i^c \cap W$.
Il suffit donc de montrer que $U \cap V^c \cap W$ est quasi-compact
si $U, V$ sont des ouverts rétrocompacts et $W$ un ouvert
quasi-compact. Cela est vrai parce que $U \cap V^c \cap W$ est une partie fermée
de la partie quasi-compacte $U \cap W$, si bien que le
lemme \ref{topology-lemma-closed-in-quasi-compact}
s'applique.
\end{proof}

\noindent
Question : le lemme suivant reste-t-il valable si l'on suppose que $X$ est un
espace topologique quasi-compact ? Comparer avec le
lemme \ref{topology-lemma-intersect-constructible-with-closed}.

\begin{lemma}
\label{topology-lemma-intersect-constructible-with-constructible}
Soit $X$ un espace topologique. Supposons que
$X$ possède une base constituée d'ouverts quasi-compacts.
Pour des parties $E, E'$ constructibles dans $X$, l'intersection
$E \cap E'$ est constructible dans $E$.
\end{lemma}

\begin{proof}
Combiner les lemmes \ref{topology-lemma-intersect-constructible-with-retrocompact} et
\ref{topology-lemma-constructible-is-retrocompact}.
\end{proof}

\begin{lemma}
\label{topology-lemma-constructible-in-constructible}
Soit $X$ un espace topologique. Supposons que
$X$ possède une base constituée d'ouverts quasi-compacts.
Soit $E$ une partie constructible dans $X$ et $F \subset E$ une partie constructible dans $E$.
Alors $F$ est constructible dans $X$.
\end{lemma}

\begin{proof}
Observons que toute partie rétrocompacte $T$ de $X$ possède, pour la topologie
induite, une base constituée d'ouverts quasi-compacts. Cela vaut en particulier
pour toute partie constructible
(lemme \ref{topology-lemma-constructible-is-retrocompact}).
Écrivons $E = E_1 \cup \ldots \cup E_n$, avec $E_i = U_i \cap V_i^c$,
où $U_i, V_i \subset X$ sont des ouverts rétrocompacts.
Notons que $E_i = E \cap E_i$ est constructible dans $E$ d'après le
lemme \ref{topology-lemma-intersect-constructible-with-constructible}.
Ainsi, $F \cap E_i$ est constructible dans $E_i$ d'après le
lemme \ref{topology-lemma-intersect-constructible-with-constructible}.
Il suffit donc de démontrer le lemme dans le cas où $E = U \cap V^c$,
où $U, V \subset X$ sont des ouverts rétrocompacts.
Dans ce cas, l'inclusion $E \subset X$ est la composée
$$
E = U \cap V^c \to U \to X
$$
On peut alors appliquer le lemme \ref{topology-lemma-closed-constructible-image}
à la première inclusion et le
lemme \ref{topology-lemma-quasi-compact-open-immersion-constructible-image}
à la seconde.
\end{proof}

\begin{lemma}
\label{topology-lemma-locally-closed-constructible-image}
Soit $X$ un espace topologique quasi-compact possédant une base constituée
d'ouverts quasi-compacts telle que l'intersection de deux
ouverts quasi-compacts quelconques soit quasi-compacte.
Soit $T \subset X$ une partie localement fermée
telle que $T$ soit quasi-compacte et que $T^c$ soit rétrocompacte dans $X$.
Alors $T$ est constructible dans $X$.
\end{lemma}

\begin{proof}
Notons que $T$ est quasi-compacte et ouverte dans $\overline{T}$.
En utilisant notre base d'ouverts quasi-compacts, nous pouvons écrire
$T = U \cap \overline{T}$, où $U$ est un ouvert quasi-compact de $X$.
Alors $V = U \setminus T = U \cap T^c$ est rétrocompact dans $U$, puisque $T^c$
est rétrocompact dans $X$. Par conséquent, $V$ est quasi-compact. Puisque
l'intersection de deux ouverts quasi-compacts quelconques est quasi-compacte,
tout ouvert quasi-compact de $X$ est rétrocompact. Ainsi, $T = U \cap V^c$,
où $U$ et $V = U \setminus T$ sont des ouverts rétrocompacts de $X$.
A fortiori, $T$ est constructible dans $X$.
\end{proof}

\begin{lemma}
\label{topology-lemma-collate-constructible-from-constructible}
Soit $X$ un espace topologique dont la topologie possède une base
constituée d'ouverts quasi-compacts. Soit $E \subset X$ une partie.
Soit $X = E_1 \cup \ldots \cup E_m$ un recouvrement fini par des parties
constructibles. Alors $E$ est constructible dans $X$ si et seulement si $E \cap E_j$
est constructible dans $E_j$ pour chaque $j = 1, \ldots, m$.
\end{lemma}

\begin{proof}
Combiner les
lemmes \ref{topology-lemma-intersect-constructible-with-constructible} et
\ref{topology-lemma-constructible-in-constructible}.
\end{proof}

\begin{lemma}
\label{topology-lemma-generic-point-in-constructible}
Soit $X$ un espace topologique. Supposons que
$Z \subset X$ soit irréductible. Soit $E \subset X$
une réunion finie de parties localement fermées (par exemple,\ $E$
est constructible). Les assertions suivantes sont équivalentes :
\begin{enumerate}
\item L'intersection $E \cap Z$ contient une partie ouverte
dense de $Z$.
\item L'intersection $E \cap Z$ est dense dans $Z$.
\end{enumerate}
Si $Z$ possède un point générique $\xi$, elles sont
également équivalentes à
\begin{enumerate}
\item[(3)] On a $\xi \in E$.
\end{enumerate}
\end{lemma}

\begin{proof}
L'implication (1) $\Rightarrow$ (2) est claire. Supposons (2).
Notons que $E \cap Z$ est une réunion finie de parties localement fermées
$Z_i$ de $Z$. Puisque $Z$ est irréductible, l'une des $Z_i$ doit
être dense dans $Z$. Cette partie $Z_i$ est alors ouverte et dense dans $Z$, car elle est
ouverte dans son adhérence. Par conséquent, (1) est vérifiée.

\medskip\noindent
Supposons que $\xi \in Z$ soit un point générique. Si les conditions équivalentes
(1) et (2) sont vérifiées, alors $\xi \in E$. Réciproquement, si $\xi \in E$, alors
$\xi \in E \cap Z$ et, par conséquent, $E \cap Z$ est dense dans $Z$.
\end{proof}






 

\section{Parties constructibles et espaces noethériens}
\label{topology-section-constructible-Noetherian}

\begin{lemma}
\label{topology-lemma-constructible-Noetherian-space}
Soit $X$ un espace topologique noethérien.
Les parties constructibles de $X$ sont exactement les réunions finies de
parties localement fermées de $X$.
\end{lemma}

\begin{proof}
Cela résulte immédiatement du lemme \ref{topology-lemma-Noetherian-quasi-compact}.
\end{proof}

\begin{lemma}
\label{topology-lemma-constructible-map-Noetherian}
Soit $f : X \to Y$ une application continue entre espaces topologiques
noethériens. Si $E \subset Y$ est constructible dans $Y$, alors $f^{-1}(E)$
est constructible dans $X$.
\end{lemma}

\begin{proof}
Cela résulte immédiatement du
lemme \ref{topology-lemma-constructible-Noetherian-space} et de la définition d'une
application continue.
\end{proof}

\begin{lemma}
\label{topology-lemma-characterize-constructible-Noetherian}
Soit $X$ un espace topologique noethérien.
Soit $E \subset X$ une partie.
Les propriétés suivantes sont équivalentes :
\begin{enumerate}
\item $E$ est constructible dans $X$, et
\item pour tout fermé irréductible $Z \subset X$, l'intersection $E \cap Z$
contient un ouvert non vide de $Z$ ou n'est pas dense dans $Z$.
\end{enumerate}
\end{lemma}

\begin{proof}
Supposons $E$ constructible et $Z \subset X$ fermé irréductible.
Alors $E \cap Z$ est constructible dans $Z$ d'après le
lemme \ref{topology-lemma-constructible-map-Noetherian}.
Ainsi, $E \cap Z$ est une réunion finie de parties localement fermées non
vides $T_i$ de $Z$. Il est clair que, si aucun des $T_i$ n'est ouvert dans
$Z$, alors $E \cap Z$ n'est pas dense dans $Z$. On voit ainsi que (1)
implique (2).

\medskip\noindent
Réciproquement, supposons (2) vérifiée. Considérons l'ensemble $\mathcal{S}$
des fermés $Y$ de $X$ tels que $E \cap Y$ ne soit pas constructible dans $Y$.
Si $\mathcal{S} \not = \emptyset$, il admet un plus petit élément $Y$, car
$X$ est noethérien.
Soit $Y = Y_1 \cup \ldots \cup Y_r$ la décomposition de $Y$ en ses
composantes irréductibles ; voir le lemme \ref{topology-lemma-Noetherian}.
Si $r > 1$, alors chaque $Y_i \cap E$ est constructible dans $Y_i$, donc est
une réunion finie de parties localement fermées de $Y_i$. Par conséquent,
$E \cap Y$ est lui aussi une réunion finie de parties localement fermées de
$Y$, et nous concluons que $E \cap Y$ est constructible dans $Y$ d'après le
lemme \ref{topology-lemma-constructible-Noetherian-space}.
C'est une contradiction ; ainsi $r = 1$. Si $r = 1$, alors $Y$ est
irréductible et, par l'hypothèse (2), on voit que $E \cap Y$ se trouve dans
l'un des deux cas suivants : (a) il contient un ouvert non vide $V$ de $Y$ ; ou
(b) il n'est pas dense dans $Y$.
Dans le cas (a), la minimalité de $Y$ montre que $E \cap (Y \setminus V)$
est une réunion finie de parties localement fermées de $Y \setminus V$.
Ainsi, $E \cap Y$ est une réunion finie de parties localement fermées de $Y$
et est constructible d'après le
lemme \ref{topology-lemma-constructible-Noetherian-space}.
C'est une contradiction ; nous sommes donc nécessairement dans le cas (b).
Dans le cas (b), on a $E \cap Y = E \cap Y'$ pour un fermé propre
$Y' \subset Y$. La minimalité de $Y$ montre que $E \cap Y'$ est une réunion
finie de parties localement fermées de $Y'$, puis que
$E \cap Y' = E \cap Y$ est une réunion finie de parties localement fermées
de $Y$ et est constructible d'après le
lemme \ref{topology-lemma-constructible-Noetherian-space}.
Cette contradiction achève la démonstration du lemme.
\end{proof}

\begin{lemma}
\label{topology-lemma-constructible-neighbourhood-Noetherian}
Soit $X$ un espace topologique noethérien.
Soit $x \in X$.
Soit $E \subset X$ constructible dans $X$.
Les propriétés suivantes sont équivalentes :
\begin{enumerate}
\item $E$ est un voisinage de $x$, et
\item pour tout fermé irréductible $Y$ de $X$ qui contient $x$,
l'intersection $E \cap Y$ est dense dans $Y$.
\end{enumerate}
\end{lemma}

\begin{proof}
Il est clair que (1) implique (2). Supposons (2).
Considérons l'ensemble $\mathcal{S}$ des fermés $Y$ de $X$ qui contiennent
$x$ et tels que $E \cap Y$ ne soit pas un voisinage de $x$ dans $Y$.
Si $\mathcal{S} \not = \emptyset$, il admet un élément minimal $Y$, car $X$
est noethérien. Supposons $Y = Y_1 \cup Y_2$, où $Y_1$, $Y_2$ sont deux
fermés non vides strictement plus petits. Si $x \in Y_i$ pour $i = 1, 2$,
alors $Y_i \cap E$ est un voisinage de $x$ dans $Y_i$, et nous concluons que
$Y \cap E$ est un voisinage de $x$ dans $Y$, contradiction. Si $x \in Y_1$
mais $x \not\in Y_2$ (disons), alors $Y_1 \cap E$ est un voisinage de $x$
dans $Y_1$, donc aussi dans $Y$, ce qui est encore une contradiction.
Nous concluons que $Y$ est un fermé irréductible. Par l'hypothèse (2),
$E \cap Y$ est dense dans $Y$. Ainsi, $E \cap Y$ contient un ouvert $V$ de
$Y$ ; voir le lemme \ref{topology-lemma-characterize-constructible-Noetherian}.
Si $x \in V$, alors $E \cap Y$ est un voisinage de $x$ dans $Y$, ce qui est
une contradiction. Si $x \not \in V$, alors $Y' = Y \setminus V$ est un
fermé propre de $Y$ qui contient $x$. La minimalité de $Y$ montre que
$E \cap Y'$ contient un voisinage ouvert $V' \subset Y'$ de $x$ dans $Y'$.
Mais alors $V' \cup V$ est un voisinage ouvert de $x$ dans $Y$ contenu dans
$E$, contradiction.
Cette contradiction achève la démonstration du lemme.
\end{proof}

\begin{lemma}
\label{topology-lemma-characterize-open-Noetherian}
Soit $X$ un espace topologique noethérien.
Soit $E \subset X$ une partie.
Les propriétés suivantes sont équivalentes :
\begin{enumerate}
\item $E$ est ouvert dans $X$, et
\item pour tout fermé irréductible $Y$ de $X$, l'intersection $E \cap Y$ est
vide ou contient un ouvert non vide de $Y$.
\end{enumerate}
\end{lemma}

\begin{proof}
Cela résulte formellement des
lemmes \ref{topology-lemma-characterize-constructible-Noetherian} et
\ref{topology-lemma-constructible-neighbourhood-Noetherian}.
\end{proof}





\section{Caractérisation des applications propres}
\label{topology-section-proper}

\noindent
Nous consacrons une section à la notion d'application propre en topologie
usuelle. Nous définissons une application continue d'espaces topologiques
comme étant {\it propre} si elle est universellement fermée et séparée.
Bien que cette définition s'accorde avec celle d'un morphisme propre en
géométrie algébrique, elle diffère de la définition de Bourbaki.
Avec notre définition d'une application propre d'espaces topologiques, le
théorème de changement de base propre (Cohomologie,
théorème \ref{cohomology-theorem-proper-base-change}) est valable sans aucune
hypothèse supplémentaire. De plus, pour un morphisme $f : X \to Y$ de schémas
de type fini sur $\mathbf{C}$, on a l'équivalence suivante : $f$ est propre
comme morphisme de schémas si et seulement si l'application continue
$f : X(\mathbf{C}) \to Y(\mathbf{C})$ sur les points à valeurs dans
$\mathbf{C}$, munis de la topologie classique, est propre.
Ceci est expliqué dans \cite[Exp. XII, Prop.\ 3.2(v)]{SGA1}, où une note
signale aussi que la notion topologique de propreté employée est celle de
Bourbaki, à laquelle on ajoute l'hypothèse de séparation.

\medskip\noindent
Il nous paraît utile de nommer trois notions distinctes : être séparée, être
universellement fermée, et posséder ces deux propriétés à la fois
(c'est-à-dire être propre). Pour une application continue $f : X \to Y$
entre espaces localement compacts et séparés, le mot « propre » désigne depuis
longtemps la propriété « $f^{-1}($compact$) =$ compact », laquelle équivaut à
être universellement fermée pour des espaces aussi réguliers.
En fait, nous verrons que, dans le cas général, la condition sur les images
réciproques, formulée par souci de clarté à l'aide du mot « quasi-compact »,
équivaut à être universellement fermée si l'on suppose en outre l'application
fermée. Voir aussi \cite[exercices 22 à 26 du chapitre II]{LangReal}, mais
attention : Lang emploie « propre » comme synonyme d'« universellement
fermée », à l'instar de Bourbaki.

\begin{lemma}[Lemme du tube]
\label{topology-lemma-tube}
Soient $X$ et $Y$ des espaces topologiques.
Soient $A \subset X$ et $B \subset Y$ des parties quasi-compactes.
Soit $A \times B \subset W \subset X \times Y$, où $W$ est ouvert dans
$X \times Y$. Il existe alors des ouverts $A \subset U \subset X$ et
$B \subset V \subset Y$ tels que $U \times V \subset W$.
\end{lemma}

\begin{proof}
Pour tous $a \in A$ et $b \in B$, il existe des ouverts $U_{(a, b)}$ de $X$
et $V_{(a, b)}$ de $Y$ tels que
$(a, b) \in U_{(a, b)} \times V_{(a, b)} \subset W$.
Fixons $b$ ; il existe alors un nombre fini de points $a_1, \ldots, a_n$ tels
que $A \subset U_{(a_1, b)} \cup \ldots \cup U_{(a_n, b)}$.
Par conséquent,
$$
A \times \{b\} \subset
(U_{(a_1, b)} \cup \ldots \cup U_{(a_n, b)}) \times
(V_{(a_1, b)} \cap \ldots \cap V_{(a_n, b)}) \subset W.
$$
Ainsi, pour tout $b \in B$, il existe des ouverts $U_b \subset X$ et
$V_b \subset Y$ tels que $A \times \{b\} \subset U_b \times V_b \subset W$.
Comme précédemment, il existe un nombre fini de points $b_1, \ldots, b_m$
tels que $B \subset V_{b_1} \cup \ldots \cup V_{b_m}$.
Le résultat en découle, car
$A \times B \subset
(U_{b_1} \cap \ldots \cap U_{b_m}) \times
(V_{b_1} \cup \ldots \cup V_{b_m})$.
\end{proof}

\noindent
La terminologie de la définition suivante peut différer légèrement de celle
dont le lecteur a l'habitude.

\begin{definition}
\label{topology-definition-proper-map}
Soit $f : X\to Y$ une application continue entre espaces topologiques.
\begin{enumerate}
\item On dit que l'application $f$ est {\it fermée} si l'image de toute partie
fermée est fermée.
\item On dit que l'application $f$ est {\it propre au sens de
Bourbaki}\footnote{C'est la terminologie employée dans \cite{Bourbaki}.
Cette propriété est parfois appelée « universellement fermée » dans la
littérature.} si l'application $Z \times X\to Z \times Y$ est fermée pour tout
espace topologique $Z$.
\item On dit que l'application $f$ est {\it quasi-propre} si l'image réciproque
$f^{-1}(V)$ de toute partie quasi-compacte $V \subset Y$ est quasi-compacte.
\item On dit que $f$ est {\it universellement fermée} si l'application
$f': Z \times_Y X \to Z$ est fermée pour toute application continue
$g: Z \to Y$.
\item On dit que $f$ est {\it propre} si $f$ est séparée et universellement
fermée.
\end{enumerate}
\end{definition}

\noindent
Le lemme suivant sera utile plus loin.

\begin{lemma}
\label{topology-lemma-characterize-quasi-compact}
\begin{reference}
Combinaison de
\cite[I, p. 75, Lemme 1]{Bourbaki} et de
\cite[I, p. 76, Corollaire 1]{Bourbaki}.
\end{reference}
Un espace topologique $X$ est quasi-compact si et seulement si l'application
de projection $Z \times X \to Z$ est fermée pour tout espace topologique $Z$.
\end{lemma}

\begin{proof}
(Voir aussi la remarque ci-dessous.)
Si $X$ n'est pas quasi-compact, il existe un recouvrement ouvert
$X = \bigcup_{i \in I} U_i$ tel qu'aucun nombre fini des $U_i$ ne recouvre
$X$.
Soit $Z$ la partie de l'ensemble des parties $\mathcal{P}(I)$ de $I$
constituée de $I$ et de toutes les parties finies non vides de $I$.
Définissons sur $Z$ une topologie ayant pour base les ensembles suivants :
\begin{enumerate}
\item Toutes les parties de $Z\setminus\{I\}$.
\item Pour toute partie finie $K$ de $I$, l'ensemble
$U_K := \{J\subset I \mid J \in Z, \ K\subset J \})$.
\end{enumerate}
Nous laissons au lecteur le soin de vérifier qu'il s'agit bien d'une base de
topologie. Considérons la partie de $Z \times X$ définie par la formule
$$
M = \{(J, x) \mid J \in Z, \ x \in \bigcap\nolimits_{i \in J} U_i^c)\}
$$
Si $(J, x) \not \in M$, alors $x \in U_i$ pour un certain $i \in J$.
Ainsi, $U_{\{i\}} \times U_i \subset Z \times X$ est un ouvert contenant
$(J, x)$ et ne rencontrant pas $M$. Par conséquent, $M$ est fermé. La
projection de $M$ sur $Z$ est $Z-\{I\}$, qui n'est pas fermée. L'application
$Z \times X \to Z$ n'est donc pas fermée.

\medskip\noindent
Supposons $X$ quasi-compact. Soit $Z$ un espace topologique.
Soit $M \subset  Z \times X$ fermé. Soit $z \in Z$ un point qui n'appartient
pas à $\text{pr}_1(M)$. D'après le lemme du tube \ref{topology-lemma-tube}, il existe un
ouvert $U \subset Z$ tel que $U \times X$ soit contenu dans le complémentaire
de $M$. Ainsi, $\text{pr}_1(M)$ est fermé.
\end{proof}

\begin{remark}
\label{topology-remark-lemma-literature}
Le lemme \ref{topology-lemma-characterize-quasi-compact} combine
\cite[I, p. 75, Lemme 1]{Bourbaki} et
\cite[I, p. 76, Corollaire 1]{Bourbaki}.
\end{remark}

\begin{theorem}
\label{topology-theorem-characterize-proper}
\begin{reference}
Dans \cite[I, p. 75, théorème 1]{Bourbaki}, on trouve :
(2) $\Leftrightarrow$ (4).
Dans \cite[I, p. 77, Proposition 6]{Bourbaki}, on trouve :
(2) $\Rightarrow$ (1).
\end{reference}
Soit $f: X\to Y$ une application continue entre espaces topologiques.
Les conditions suivantes sont équivalentes :
\begin{enumerate}
\item L'application $f$ est quasi-propre et fermée.
\item L'application $f$ est propre au sens de Bourbaki.
\item L'application $f$ est universellement fermée.
\item L'application $f$ est fermée et $f^{-1}(y)$ est quasi-compact pour tout
$y\in Y$.
\end{enumerate}
\end{theorem}

\begin{proof}
(Voir aussi la remarque ci-dessous.)
Si l'application $f$ vérifie (1), elle vérifie automatiquement (4), car tout
singleton est quasi-compact.

\medskip\noindent
Supposons que l'application $f$ vérifie (4).
Nous allons montrer qu'elle est universellement fermée, c'est-à-dire que (3)
est vérifiée. Soit $g : Z \to Y$ une application continue d'espaces
topologiques et considérons le diagramme
$$
\xymatrix{
Z \times_Y X \ar[r]_{g'} \ar[d]_{f'} & X \ar[d]^f \\
Z \ar[r]^g & Y
}
$$
Dans la démonstration, nous utiliserons le fait que
$Z \times_Y X \to Z \times X$ est un homéomorphisme sur son image,
c'est-à-dire que nous pouvons identifier $Z \times_Y X$ à la partie
correspondante de $Z \times X$, munie de la topologie induite.
L'image de $f' : Z \times_Y X \to Z$ est
$\Im(f') = \{z : g(z) \in f(X)\}$.
Comme $f(X)$ est fermé, on voit que $\Im(f')$ est un sous-espace fermé de $Z$.
Considérons un fermé $P \subset Z \times_Y X$.
Soit $z \in Z$, $z \not \in f'(P)$.
Si $z \not \in \Im(f')$, alors $Z \setminus \Im(f')$ est un voisinage
ouvert qui ne rencontre pas $f'(P)$.
Si $z$ appartient à $\Im(f')$, alors
$(f')^{-1}\{z\} = \{z\} \times f^{-1}\{g(z)\}$, et
$f^{-1}\{g(z)\}$ est quasi-compact par hypothèse. Comme $P$ est un fermé de
$Z \times_Y X$, il existe un fermé $P'$ de $Z \times X$ tel que
$P = P' \cap Z \times_Y X$.
Puisque $(f')^{-1}\{z\}$ est contenu dans
$P^c = P'^c \cup (Z \times_Y X)^c$ et puisque $(f')^{-1}\{z\}$ est disjoint
de $(Z \times_Y X)^c$, on voit que $(f')^{-1}\{z\}$ est contenu dans $P'^c$.
Nous pouvons appliquer le lemme du tube \ref{topology-lemma-tube} à
$(f')^{-1}\{z\} = \{z\} \times f^{-1}\{g(z)\}
\subset (P')^c \subset Z \times X$.
Nous obtenons ainsi $V \times U$ contenant $(f')^{-1}\{z\}$, où $U$ et $V$
sont des ouverts respectivement de $X$ et de $Z$, et tel que
$V \times U$ ne rencontre pas $P'$. Alors l'ensemble
$V \cap g^{-1}(Y-f(U^c))$ est ouvert dans $Z$, puisque $f$ est fermée ; il
contient $z$ et ne rencontre pas l'image de $P$. Ainsi, $f'(P)$ est fermé.
Autrement dit, l'application $f$ est universellement fermée.

\medskip\noindent
L'implication (3) $\Rightarrow$ (2) est immédiate.
En effet, étant donné un espace topologique $Z$, considérons la projection
$g : Z \times Y \to Y$. Il est alors facile de voir que $f'$ est
l'application $Z \times X \to Z \times Y$, autrement dit que
$(Z \times Y) \times_Y X = Z \times X$. (Cette identification est une
propriété purement catégorique, qui n'a rien de spécifiquement topologique.)

\medskip\noindent
Supposons que $f$ vérifie (2). Nous allons montrer qu'elle vérifie (1).
Notons que $f$ est fermée, car $f$ s'identifie à l'application
$\{pt\} \times X \to \{pt\} \times Y$, qui est fermée par hypothèse.
Choisissons une partie quasi-compacte quelconque $K \subset Y$.
Soit $Z$ un espace topologique quelconque.
Comme $Z \times X \to Z \times Y$ est fermée, on voit que l'application
$Z \times f^{-1}(K) \to Z \times K$ est fermée (si $T$ est fermé dans
$Z \times f^{-1}(K)$, écrivons
$T = Z \times f^{-1}(K) \cap T'$ pour un certain fermé
$T' \subset Z \times X$). Comme $K$ est quasi-compact,
$K \times Z\to Z$ est fermée d'après le
lemme \ref{topology-lemma-characterize-quasi-compact}.
Par conséquent, la composée $Z \times f^{-1}(K)\to Z \times K \to Z$ est
fermée ; le lemme \ref{topology-lemma-characterize-quasi-compact} montre alors, une
nouvelle fois, que $f^{-1}(K)$ est quasi-compact.
\end{proof}

\begin{remark}
\label{topology-remark-proof-literature}
Voici quelques références bibliographiques.
Dans \cite[I, p. 75, théorème 1]{Bourbaki}, on trouve :
(2) $\Leftrightarrow$ (4).
Dans \cite[I, p. 77, Proposition 6]{Bourbaki}, on trouve :
(2) $\Rightarrow$ (1).
Bien entendu, on a trivialement (1) $\Rightarrow$ (4).
Ainsi, (1), (2) et (4) sont équivalentes.
L'équivalence de (3) et (4) est \cite[chapitre II, exercice 25]{LangReal}.
\end{remark}

\begin{lemma}
\label{topology-lemma-closed-map}
\begin{slogan}
Une application d'un espace compact vers un espace séparé est
universellement fermée.
\end{slogan}
Soit $f : X \to Y$ une application continue d'espaces topologiques.
Si $X$ est quasi-compact et $Y$ séparé, alors $f$ est universellement fermée.
\end{lemma}

\begin{proof}
Comme tout point de $Y$ est fermé, le
lemme \ref{topology-lemma-closed-in-quasi-compact} montre que le fermé $f^{-1}(y)$ de
$X$ est quasi-compact pour tout $y \in Y$.
Ainsi, d'après le théorème \ref{topology-theorem-characterize-proper}, il suffit de
montrer que $f$ est fermée.
Si $E \subset X$ est fermé, alors il est quasi-compact
(lemme \ref{topology-lemma-closed-in-quasi-compact}), donc $f(E) \subset Y$ est
quasi-compact (lemme \ref{topology-lemma-image-quasi-compact}), donc $f(E)$ est fermé
dans $Y$ (lemme \ref{topology-lemma-quasi-compact-in-Hausdorff}).
\end{proof}

\begin{lemma}
\label{topology-lemma-bijective-map}
Soit $f : X \to Y$ une application continue d'espaces topologiques.
Si $f$ est bijective, si $X$ est quasi-compact et si $Y$ est séparé, alors
$f$ est un homéomorphisme.
\end{lemma}

\begin{proof}
Il suffit de montrer que $f$ est fermée, car il en résulte que $f^{-1}$ est
continue. Si $T \subset X$ est fermé, alors $T$ est quasi-compact d'après le
lemme \ref{topology-lemma-closed-in-quasi-compact}, donc $f(T)$ est quasi-compact
d'après le lemme \ref{topology-lemma-image-quasi-compact}, donc $f(T)$ est fermé
d'après le lemme \ref{topology-lemma-quasi-compact-in-Hausdorff}.
\end{proof}












\section{Espaces de Jacobson}
\label{topology-section-space-jacobson}

\begin{definition}
\label{topology-definition-space-jacobson}
Soit $X$ un espace topologique.
Soit $X_0$ l'ensemble des points fermés de $X$.
On dit que $X$ est {\it de Jacobson} si toute
partie fermée $Z \subset X$ est l'adhérence
de $Z \cap X_0$.
\end{definition}

\noindent
Notons qu'un espace topologique $X$ est de Jacobson si et seulement si
toute partie localement fermée non vide de $X$
possède un point fermé dans $X$.

\medskip\noindent
Soit $X$ un espace de Jacobson et soit $X_0$ l'ensemble
des points fermés de $X$, muni de la topologie induite.
Il est clair que la définition implique que le morphisme
$X_0 \to X$ induit une bijection entre les parties fermées
de $X_0$ et les parties fermées de $X$.
Ainsi, nombre de propriétés de $X$ passent à $X_0$.
Par exemple, les dimensions de Krull de $X$ et de $X_0$
sont égales.

\begin{lemma}
\label{topology-lemma-jacobson-check-irreducible-closed}
Soit $X$ un espace topologique. Soit $X_0$ l'ensemble
des points fermés de $X$.
Supposons que, pour tout point $x\in X$,
l'intersection $X_0 \cap \overline{\{x\}}$ soit dense dans $\overline{\{x\}}$.
Alors $X$ est de Jacobson.
\end{lemma}

\begin{proof}
Soient $Z$ une partie fermée de $X$
et $U$ un ouvert de $X$
tels que $U\cap Z$ soit non vide.
Alors, pour $x\in U\cap Z$, l'ensemble $\overline{\{x\}}\cap U$ est une partie
non vide de $Z\cap U$
et, par hypothèse, il contient un point fermé dans $X$, comme voulu.
\end{proof}

\begin{lemma}
\label{topology-lemma-non-jacobson-Noetherian-characterize}
Soit $X$ un espace topologique de Kolmogorov muni d'une base d'ouverts
quasi-compacts.
Si $X$ n'est pas de Jacobson, alors il existe un point non fermé
$x \in X$ tel que $\{x\}$ soit localement fermé.
\end{lemma}

\begin{proof}
Puisque $X$ n'est pas de Jacobson, il existe une partie fermée $Z$ et un ouvert $U$
de $X$ tels que $Z \cap U$ soit non vide et ne contienne aucun point fermé
dans $X$. Comme $X$ possède une base d'ouverts quasi-compacts, nous pouvons
remplacer $U$
par un voisinage ouvert quasi-compact d'un point de $Z\cap U$ et ainsi supposer
que $U$ est un ouvert quasi-compact. D'après le
lemme \ref{topology-lemma-quasi-compact-closed-point}, il existe un point
$x \in Z \cap U$ fermé dans $Z \cap U$ ;
ainsi, $\{x\}$ est localement fermé mais n'est pas fermé dans $X$.
\end{proof}

\begin{lemma}
\label{topology-lemma-jacobson-local}
Soit $X$ un espace topologique.
Soit $X = \bigcup U_i$ un recouvrement ouvert.
Alors $X$ est de Jacobson si et seulement si chaque $U_i$ est de Jacobson.
De plus, dans ce cas, $X_0 = \bigcup U_{i, 0}$.
\end{lemma}

\begin{proof}
Soit $X$ un espace topologique.
Soit $X_0$ l'ensemble des points fermés de $X$.
Soit $U_{i, 0}$ l'ensemble des points fermés de
$U_i$. Alors $X_0 \cap U_i \subset U_{i, 0}$,
mais l'égalité n'est pas toujours vérifiée.

\medskip\noindent
Supposons d'abord que chaque $U_i$ soit de Jacobson.
Nous affirmons qu'alors $X_0 \cap U_i = U_{i, 0}$.
En effet, supposons que $x \in U_{i, 0}$, c'est-à-dire que $x$ soit fermé dans
$U_i$. Soit $\overline{\{x\}}$ l'adhérence
dans $X$. Considérons $\overline{\{x\}} \cap U_j$.
Si $x \not \in U_j$, alors $\overline{\{x\}} \cap U_j = \emptyset$.
Si $x \in U_j$, alors $U_i \cap U_j \subset U_j$
est un ouvert de $U_j$ contenant $x$.
Soient $T' = U_j \setminus U_i \cap U_j$ et
$T = \{x\} \amalg T'$. Comme $x$ est fermé dans $U_i \cap U_j$
(puisque $x$ est fermé dans $U_i$), on voit que $T$ et $T'$
sont des parties fermées de $U_j$ et que $T$ contient
$x$. Comme $U_j$ est de Jacobson, les points fermés de
$U_j$ sont denses dans $T$. Puisque $T = \{x\} \amalg T'$,
cela n'est possible que si $x$ est fermé dans $U_j$.
Par conséquent, $\overline{\{x\}} \cap U_j = \{x\}$. Nous concluons
que $\overline{\{x\}} = \{ x \}$, comme voulu.

\medskip\noindent
Soit $Z \subset X$ une partie fermée (toujours
en supposant que chaque $U_i$ soit de Jacobson).
Puisque nous savons maintenant que $X_0 \cap Z  \cap U_i
= U_{i, 0} \cap Z$ est dense dans $Z \cap U_i$,
il s'ensuit immédiatement que $X_0 \cap Z$ est
dense dans $Z$.

\medskip\noindent
Réciproquement, supposons que $X$ soit de Jacobson.
Soit $Z \subset U_i$ une partie fermée. Alors
$X_0 \cap \overline{Z}$ est dense dans $\overline{Z}$.
Par conséquent, $X_0 \cap Z$ est également dense dans $Z$, car
$\overline{Z} \setminus Z$ est fermé. Comme $X_0 \cap U_i
\subset U_{i, 0}$, on voit que
$U_{i, 0} \cap Z$ est dense dans $Z$.
Ainsi, $U_i$ est de Jacobson, comme voulu.
\end{proof}

\begin{lemma}
\label{topology-lemma-jacobson-inherited}
Soit $X$ un espace de Jacobson. Les types suivants de parties $T \subset X$
sont de Jacobson :
\begin{enumerate}
\item Les sous-espaces ouverts.
\item Les sous-espaces fermés.
\item Les sous-espaces localement fermés.
\item Les réunions de sous-espaces localement fermés.
\item Les parties constructibles.
\item Toute partie $T \subset X$ qui, localement sur $X$,
est une réunion de parties localement fermées.
\end{enumerate}
Dans chacun de ces cas, les points fermés de $T$ sont
fermés dans $X$.
\end{lemma}

\begin{proof}
Soit $X_0$ l'ensemble des points fermés de $X$. Pour toute partie
$T \subset X$, notons $(*)$ la propriété suivante :
\begin{itemize}
\item[(*)] Toute partie localement fermée non vide de $T$ possède un point
fermé dans $X$.
\end{itemize}
Notons que l'on a toujours $X_0 \cap T \subset T_0$. La propriété $(*)$
implique donc que $T$ est de Jacobson. En outre, elle implique clairement
que tout point fermé de $T$ est fermé dans $X$.

\medskip\noindent
Supposons que $T=\bigcup_i T_i$, où $T_i$ est localement fermé dans $X$.
Prenons une partie $A\subset T$ localement fermée et non vide dans $T$ ;
on peut alors choisir $T_i$ de telle sorte que $A\cap T_i$ soit non vide ; cette intersection est
localement fermée dans $T_i$, donc dans $X$.
Comme $X$ est de Jacobson, $A$ possède un point fermé dans $X$.
\end{proof}

\begin{lemma}
\label{topology-lemma-finite-jacobson}
Un espace de Jacobson fini est discret.
Un espace de Jacobson qui ne possède qu'un nombre fini de points fermés est discret.
\end{lemma}

\begin{proof}
Si $X$ est fini, alors l'ensemble $X_0 \subset X$ de ses points fermés est fini.
Supposons que $X_0$ soit fini et que $X$ soit de Jacobson.
Alors $\overline{X_0} = X$ par la propriété de Jacobson.
Or $X_0 =\{x_1, \ldots, x_n\} = \bigcup_{i = 1, \ldots, n} \{x_i\}$
est une réunion finie de fermés, donc est fermé ; ainsi,
$X = \overline{X_0} = X_0$. Tout point est fermé et, par
finitude, tout point est ouvert.
\end{proof}

\begin{lemma}
\label{topology-lemma-jacobson-equivalent-locally-closed}
\begin{slogan}
Pour les espaces de Jacobson, les points fermés déterminent entièrement la topologie.
\end{slogan}
Supposons que $X$ soit un espace topologique de Jacobson.
Soit $X_0$ l'ensemble des points fermés de $X$.
Il existe une correspondance bijective, qui préserve les inclusions,
$$
\{\text{réunions finies de parties loc. fermées de } X\}
\leftrightarrow
\{\text{réunions finies de parties loc. fermées de } X_0\}
$$
définie par $E \mapsto E \cap X_0$. Cette correspondance préserve
les parties localement fermées, les parties ouvertes et les parties fermées.
\end{lemma}

\begin{proof}
Nous nous contentons de démontrer que la correspondance $E \mapsto E \cap X_0$ est injective.
En effet, si $E\neq E'$, alors, sans perte de généralité, $E\setminus E'$ est
non vide, et c'est une réunion finie de parties localement fermées (détails omis).
Comme $X$ est de Jacobson, on voit que
$(E \setminus E') \cap X_0 = E \cap X_0 \setminus E' \cap X_0$ n'est pas vide.
\end{proof}

\begin{lemma}
\label{topology-lemma-jacobson-equivalent-constructible}
Supposons que $X$ soit un espace topologique de Jacobson.
Soit $X_0$ l'ensemble des points fermés de $X$.
Il existe une correspondance bijective, qui préserve les inclusions,
$$
\{\text{parties constructibles de } X\}
\leftrightarrow
\{\text{parties constructibles de } X_0\}
$$
définie par $E \mapsto E \cap X_0$. Cette correspondance préserve
la sous-collection des ouverts rétrocompacts, ainsi que leurs complémentaires.
\end{lemma}

\begin{proof}
D'après le lemme \ref{topology-lemma-jacobson-equivalent-locally-closed} ci-dessus,
il suffit de voir que, si $U$ est ouvert dans $X$, alors $U\cap X_0$ est
rétrocompact dans $X_0$ si et seulement si $U$ est rétrocompact dans $X$.
Pour cela, il suffit de démontrer que, si $U$ est ouvert dans $X$, alors $U\cap X_0$ est
quasi-compact si et seulement si $U$ est quasi-compact.
D'après le lemme \ref{topology-lemma-jacobson-inherited}, nous pouvons remplacer
$X$ par $U$ et supposer que $U = X$.
Enfin, observons qu'une famille d'ouverts $\mathcal{U}$ de $X$ recouvre
$X$ si et seulement si elle recouvre $X_0$ : cela découle de la propriété
de Jacobson de $X$, appliquée au fermé $X\setminus \bigcup \mathcal{U}$ :
si ce fermé était non vide, il contiendrait un point de $X_0$.
\end{proof}



















\section{Spécialisation}
\label{topology-section-specialization}

\begin{definition}
\label{topology-definition-specialization}
Soit $X$ un espace topologique.
\begin{enumerate}
\item Si $x, x' \in X$, on dit que $x$ est une {\it spécialisation} de $x'$,
ou que $x'$ est une {\it générisation} de $x$, si
$x \in \overline{\{x'\}}$. Notation : $x' \leadsto x$.
\item Une partie $T \subset X$ est {\it stable par spécialisation} si, pour
tout $x' \in T$ et toute spécialisation $x' \leadsto x$, on a $x \in T$.
\item Une partie $T \subset X$ est {\it stable par générisation} si, pour tout
$x \in T$ et toute générisation $x' \leadsto x$, on a $x' \in T$.
\end{enumerate}
\end{definition}

\begin{lemma}
\label{topology-lemma-open-closed-specialization}
Soit $X$ un espace topologique.
\begin{enumerate}
\item Toute partie fermée de $X$ est stable par spécialisation.
\item Toute partie ouverte de $X$ est stable par générisation.
\item Une partie $T \subset X$ est stable par spécialisation si et seulement
si son complémentaire $T^c$ est stable par générisation.
\end{enumerate}
\end{lemma}

\begin{proof}
Soit $F$ une partie fermée de $X$. Si $y\in F$, alors $\{y\} \subset F$, donc
$\overline{\{y\}} \subset \overline{F} = F$, puisque $F$ est fermé. Ainsi,
pour toute spécialisation $x$ de $y$, on a $x\in F$.

\medskip\noindent
Soient $x, y\in X$ tels que $x\in \overline{\{y\}}$, et soit $T$ une partie
de $X$. Dire que $T$ est stable par spécialisation signifie que $y\in T$
implique $x\in T$ ; réciproquement, dire que $T$ est stable par générisation
signifie que $x\in T$ implique $y\in T$. L'assertion (3) résulte donc de la
contraposition.

\medskip\noindent
La deuxième assertion découle de (1) et (3) en passant au complémentaire.
\end{proof}

\begin{lemma}
\label{topology-lemma-stable-specialization}
Soit $T \subset X$ une partie d'un espace topologique $X$.
Les propriétés suivantes sont équivalentes :
\begin{enumerate}
\item $T$ est stable par spécialisation, et
\item $T$ est une réunion (filtrante) de parties fermées de $X$.
\end{enumerate}
\end{lemma}

\begin{proof}
Supposons $T$ stable par spécialisation. Alors, pour tout $y\in T$, on a
$\overline{\{y\}} \subset T$. Ainsi,
$T = \bigcup_{y\in T} \overline{\{y\}}$, qui est une réunion de parties
fermées de $X$. Réciproquement, supposons que
$T = \bigcup_{i\in I}F_i$, où les $F_i$ sont des parties fermées de $X$.
Si $y\in T$, il existe $i\in I$ tel que $y\in F_i$. Comme $F_i$ est fermé,
on a $\overline{\{y\}} \subset F_i \subset T$, ce qui prouve que $T$ est
stable par spécialisation.
\end{proof}

\begin{definition}
\label{topology-definition-lift-specializations}
Soit $f : X \to Y$ une application continue d'espaces topologiques.
\begin{enumerate}
\item On dit que {\it les spécialisations se relèvent le long de $f$}, ou que
$f$ est {\it spécialisante}, si, étant donnés $y' \leadsto y$ dans $Y$ et
$x'\in X$ tel que $f(x') = y'$, il existe une spécialisation
$x' \leadsto x$ de $x'$ dans $X$ telle que $f(x) = y$.
\item On dit que {\it les générisations se relèvent le long de $f$}, ou que
$f$ est {\it générisante}, si, étant donnés $y' \leadsto y$ dans $Y$ et
$x\in X$ tel que $f(x) = y$, il existe une générisation $x' \leadsto x$ de
$x$ dans $X$ telle que $f(x') = y'$.
\end{enumerate}
\end{definition}

\begin{lemma}
\label{topology-lemma-lift-specialization-composition}
Supposons que $f : X \to Y$ et $g : Y \to Z$ soient des applications
continues d'espaces topologiques. Si les spécialisations se relèvent le long
de $f$ et de $g$, elles se relèvent le long de $g \circ f$. Il en va de même
pour le relèvement des générisations.
\end{lemma}

\begin{proof}
Soit $z'\leadsto z$ une spécialisation dans $Z$ et soit $x' \in X$ tel que
$g\circ f (x') = z'$. Puisque les spécialisations se relèvent le long de $g$,
il existe une spécialisation $f(x') \leadsto y$ de $f(x')$ dans $Y$ telle que
$g(y) = z$. De même, puisqu'elles se relèvent le long de $f$, il existe une
spécialisation $x' \leadsto x$ de $x'$ dans $X$ telle que $f(x) = y$. Cela
fournit une spécialisation $x' \leadsto x$ de $x'$ dans $X$ telle que
$g\circ f(x) = z$. Autrement dit, les spécialisations se relèvent le long de
$g\circ f$.
\end{proof}

\begin{lemma}
\label{topology-lemma-lift-specializations-images}
Soit $f : X \to Y$ une application continue d'espaces topologiques.
\begin{enumerate}
\item Si les spécialisations se relèvent le long de $f$ et si $T \subset X$
est stable par spécialisation, alors $f(T) \subset Y$ est stable par
spécialisation.
\item Si les générisations se relèvent le long de $f$ et si $T \subset X$ est
stable par générisation, alors $f(T) \subset Y$ est stable par générisation.
\end{enumerate}
\end{lemma}

\begin{proof}
Soit $y' \leadsto y$ une spécialisation dans $Y$, où $y'\in f(T)$, et soit
$x'\in T$ tel que $f(x') = y'$. Puisque les spécialisations se relèvent le
long de $f$, il existe une spécialisation $x'\leadsto x$ de $x'$ dans $X$
telle que $f(x) = y$. Or $T$ est stable par spécialisation, donc $x\in T$, et
alors $y \in f(T)$. Par conséquent, $f(T)$ est stable par spécialisation.

\medskip\noindent
La démonstration de (2) est identique, en utilisant le relèvement des
générisations le long de $f$.
\end{proof}

\begin{lemma}
\label{topology-lemma-closed-open-map-specialization}
Soit $f : X \to Y$ une application continue d'espaces topologiques.
\begin{enumerate}
\item Si $f$ est fermée, alors les spécialisations se relèvent le long de $f$.
\item Si $f$ est ouverte, si $X$ est un espace topologique noethérien, si tout
fermé irréductible de $X$ possède un point générique et si $Y$ est un espace
de Kolmogorov, alors les générisations se relèvent le long de $f$.
\end{enumerate}
\end{lemma}

\begin{proof}
Supposons $f$ fermée. Soient $y' \leadsto y$ dans $Y$ et $x'\in X$ tel que
$f(x') = y'$. Considérons le fermé $T = \overline{\{x'\}}$ de $X$.
Alors $f(T) \subset Y$ est fermé et $y' \in f(T)$. Par conséquent,
$y \in f(T)$. Ainsi, $y = f(x)$ avec $x \in T$, c'est-à-dire
$x' \leadsto x$.

\medskip\noindent
Supposons $f$ ouverte, $X$ noethérien, tout fermé irréductible de $X$ pourvu
d'un point générique, et $Y$ un espace de Kolmogorov.
Soient $y' \leadsto y$ dans $Y$ et $x \in X$ tel que $f(x) = y$.
Considérons $T = f^{-1}(\{y'\}) \subset X$.
Prenons un voisinage ouvert $x \in U \subset X$ de $x$.
Alors $f(U) \subset Y$ est ouvert et $y \in f(U)$. Par conséquent,
$y' \in f(U)$. Autrement dit, $T \cap U \not = \emptyset$. Ceci prouve que
$x \in \overline{T}$. Puisque $X$ est noethérien, $T$ est noethérien
(lemme \ref{topology-lemma-Noetherian}).
Il admet donc une décomposition $T = T_1 \cup \ldots \cup T_n$ en composantes
irréductibles. On a alors, de même,
$\overline{T} = \overline{T_1} \cup \ldots \cup \overline{T_n}$.
D'après ce qui précède, $x \in \overline{T_i}$ pour un certain $i$. Par
hypothèse, il existe un point générique $x' \in \overline{T_i}$, et l'on voit
que $x' \leadsto x$. Comme $x' \in \overline{T}$, on a
$f(x') \in \overline{\{y'\}}$. Remarquons que
$f(\overline{T_i}) = f(\overline{\{x'\}}) \subset \overline{\{f(x')\}}$.
Si $f(x') \not = y'$, alors, puisque $Y$ est un espace de Kolmogorov, $f(x')$ n'est pas
un point générique du fermé irréductible $\overline{\{y'\}}$, et l'inclusion
$\overline{\{f(x')\}} \subset \overline{\{y'\}}$ est stricte, c'est-à-dire
que $y' \not \in f(\overline{T_i})$.
Ceci contredit l'égalité $f(T_i) = \{y'\}$.
Ainsi, $f(x') = y'$, ce qui conclut.
\end{proof}

\begin{lemma}
\label{topology-lemma-quotient-kolmogorov}
Supposons que $s, t : R \to U$ et $\pi : U \to X$ soient des applications
continues d'espaces topologiques telles que
\begin{enumerate}
\item $\pi$ est ouverte,
\item $U$ est sobre,
\item $s, t$ ont des fibres finies,
\item les générisations se relèvent le long de $s, t$,
\item $(t, s)(R) \subset U \times U$ est une relation d'équivalence sur $U$ et
$X$ est le quotient de $U$ par cette relation d'équivalence (comme ensemble).
\end{enumerate}
Alors $X$ est un espace de Kolmogorov.
\end{lemma}

\begin{proof}
Les propriétés (3) et (5) impliquent qu'un point $x$ correspond à une classe
d'équivalence finie $\{u_1, \ldots, u_n\} \subset U$. Supposons que
$x' \in X$ soit un second point correspondant à la classe d'équivalence
$\{u'_1, \ldots, u'_m\} \subset U$.
Supposons que $u_i \leadsto u'_j$ pour certains $i, j$. Alors, pour tout
$r' \in R$ tel que $s(r') = u'_j$, la propriété (4) fournit
$r \leadsto r'$ tel que $s(r) = u_i$. Par conséquent,
$t(r) \leadsto t(r')$. Puisque
$\{u'_1, \ldots, u'_m\} = t(s^{-1}(\{u'_j\}))$, nous concluons que tout
élément de $\{u'_1, \ldots, u'_m\}$ est une spécialisation d'un élément de
$\{u_1, \ldots, u_n\}$.
Ainsi, $\overline{\{u_1\}} \cup \ldots \cup \overline{\{u_n\}}$ est une
réunion de classes d'équivalence, donc est de la forme $\pi^{-1}(Z)$ pour une
partie $Z \subset X$. La propriété (1) montre que $Z$ est fermé dans $X$ et,
en fait, que $Z = \overline{\{x\}}$, car
$\pi(\overline{\{u_i\}}) \subset \overline{\{x\}}$ pour tout $i$.
Autrement dit, $x \leadsto x'$ si et seulement si un relèvement de $x$ dans
$U$ se spécialise en un relèvement de $x'$ dans $U$, si et seulement si tout
relèvement de $x'$ dans $U$ est une spécialisation d'un relèvement de $x$
dans $U$.

\medskip\noindent
Supposons que l'on ait à la fois $x \leadsto x'$ et $x' \leadsto x$.
Disons que $x$ correspond à $\{u_1, \ldots, u_n\}$ et $x'$ à
$\{u'_1, \ldots, u'_m\}$, comme ci-dessus. D'après les résultats du paragraphe
précédent, nous pouvons alors trouver une suite
$$
\ldots \leadsto u'_{j_3} \leadsto u_{i_3} \leadsto u'_{j_2} \leadsto
u_{i_2} \leadsto u'_{j_1} \leadsto u_{i_1}
$$
qui doit comporter une répétition ; la propriété (2) donne donc
$\{u_1, \ldots, u_n\} = \{u'_1, \ldots, u'_m\}$, c'est-à-dire $x = x'$.
Ainsi, $X$ est un espace de Kolmogorov.
\end{proof}


\begin{lemma}
\label{topology-lemma-dimension-specializations-lift}
Soit $f : X \to Y$ un morphisme d'espaces topologiques.
Supposons que $Y$ soit un espace topologique sobre et que $f$ soit surjectif.
Si les spécialisations ou les générisations se relèvent le long de $f$, alors
$\dim(X) \geq \dim(Y)$.
\end{lemma}

\begin{proof}
Supposons que les spécialisations se relèvent le long de $f$.
Soit $Z_0 \subset Z_1 \subset \ldots Z_e \subset Y$ une chaîne de fermés
irréductibles de $X$. Soit $\xi_e \in X$ un point s'envoyant sur le point
générique de $Z_e$. Par hypothèse, il existe une spécialisation
$\xi_e \leadsto \xi_{e - 1}$ dans $X$ telle que $\xi_{e - 1}$ s'envoie sur le
point générique de $Z_{e - 1}$. En poursuivant ainsi, on obtient une suite de
spécialisations
$$
\xi_e \leadsto \xi_{e - 1} \leadsto \ldots \leadsto \xi_0
$$
où $\xi_i$ s'envoie sur le point générique de $Z_i$.
Cela implique clairement que la suite de fermés irréductibles
$$
\overline{\{\xi_0\}} \subset
\overline{\{\xi_1\}} \subset \ldots
\overline{\{\xi_e\}}
$$
est une chaîne de longueur $e$ dans $X$.
Le cas où les générisations se relèvent le long de $f$ est analogue.
\end{proof}

\begin{lemma}
\label{topology-lemma-characterize-closed-Noetherian}
Soit $X$ un espace topologique noethérien et sobre.
Soit $E \subset X$ une partie de $X$.
\begin{enumerate}
\item Si $E$ est constructible et stable par spécialisation, alors $E$ est
fermé.
\item Si $E$ est constructible et stable par générisation, alors $E$ est
ouvert.
\end{enumerate}
\end{lemma}

\begin{proof}
Soit $E$ constructible et stable par générisation.
Soit $Y \subset X$ un fermé irréductible de point générique $\xi \in Y$.
Si $E \cap Y$ est non vide, alors il contient $\xi$ (par stabilité par
générisation), et est donc dense dans $Y$ ; il contient par conséquent un
ouvert non vide de $Y$ ; voir le
lemme \ref{topology-lemma-characterize-constructible-Noetherian}.
Ainsi, $E$ est ouvert d'après le
lemme \ref{topology-lemma-characterize-open-Noetherian}.
Ceci démontre (2). Pour démontrer (1), appliquons (2) au complémentaire de $E$
dans $X$.
\end{proof}









\section{Fonctions de dimension}
\label{topology-section-dimension-function}

\noindent
Il n'est guère pertinent de considérer des fonctions de dimension à moins que
l'espace considéré ne soit sobre (Définition \ref{topology-definition-generic-point}).
On pourrait donc améliorer la définition ci-dessous en considérant l'espace
topologique sobre associé à $X$. Puisque l'espace topologique sous-jacent à un
schéma est sobre, nous ne nous attarderons pas sur cette amélioration.

\begin{definition}
\label{topology-definition-dimension-function}
Soit $X$ un espace topologique.
\begin{enumerate}
\item Soient $x, y \in X$, $x \not = y$. Supposons que $x \leadsto y$,
c'est-à-dire que $y$ soit une spécialisation de $x$.
On dit que $y$ est une {\it spécialisation immédiate}
de $x$ s'il n'existe aucun
$z \in X \setminus \{x, y\}$ tel que $x \leadsto z$ et $z \leadsto y$.
\item Une application $\delta : X \to \mathbf{Z}$ est appelée une
{\it fonction de dimension}\footnote{Il s'agit vraisemblablement d'une
notation non standard. Cette notion n'est habituellement introduite que
pour les schémas (localement) noethériens, auquel cas la condition (a) résulte
de (b).} si
\begin{enumerate}
\item dès que $x \leadsto y$ et $x \not = y$,
on a $\delta(x) > \delta(y)$, et
\item pour toute spécialisation immédiate $x \leadsto y$ dans $X$,
on a $\delta(x) = \delta(y) + 1$.
\end{enumerate}
\end{enumerate}
\end{definition}

\noindent
Il est clair que si $\delta$ est une fonction de dimension, alors
$\delta + t$ l'est aussi pour tout $t \in \mathbf{Z}$. Voici un petit lemme.

\begin{lemma}
\label{topology-lemma-dimension-function-catenary}
Soit $X$ un espace topologique. Si $X$ est sobre et possède une fonction de
dimension, alors $X$ est caténaire. De plus, pour toute relation de
spécialisation $x \leadsto y$,
on a
$$
\delta(x) - \delta(y) =
\text{codim}\left(\overline{\{y\}}, \ \overline{\{x\}}\right).
$$
\end{lemma}

\begin{proof}
Supposons que $Y \subset Y' \subset X$ soient des fermés irréductibles.
Soient $\xi \in Y$, $\xi' \in Y'$ leurs points génériques.
Les définitions donnent immédiatement
$\text{codim}(Y, Y') \leq \delta(\xi) - \delta(\xi') < \infty$.
En fait, la première inégalité est une égalité. En effet, supposons que
$$
Y = Y_0 \subset Y_1 \subset \ldots \subset Y_e = Y'
$$
soit une chaîne maximale quelconque de fermés irréductibles. Notons
$\xi_i \in Y_i$ le point générique. Alors
$\xi_i \leadsto \xi_{i + 1}$ est une spécialisation immédiate.
Ainsi, $e = \delta(\xi) - \delta(\xi')$, comme voulu.
Cela démontre également la dernière assertion du lemme.
\end{proof}

\begin{lemma}
\label{topology-lemma-dimension-function-unique}
Soit $X$ un espace topologique.
Soient $\delta$, $\delta'$ deux fonctions de dimension sur $X$.
Si $X$ est localement noethérien et sobre, alors $\delta - \delta'$ est
localement constante sur $X$.
\end{lemma}

\begin{proof}
Soit $x \in X$ un point. Nous allons montrer que $\delta - \delta'$ est
constante dans un voisinage de $x$.
Nous pouvons remplacer $X$ par un voisinage ouvert
noethérien de $x$ dans $X$. Nous pouvons donc supposer que $X$ est
noethérien et sobre.
Soient $Z_1, \ldots, Z_r$ les composantes irréductibles de $X$ qui
contiennent $x$. (Il n'y en a qu'un nombre fini puisque
$X$ est noethérien, voir le Lemme \ref{topology-lemma-Noetherian}.)
Soit $\xi_i \in Z_i$ le point générique.
Notons que $Z_1 \cup \ldots \cup Z_r$ est un voisinage de $x$ dans $X$
(pas nécessairement fermé). Nous affirmons que $\delta - \delta'$ est
constante sur $Z_1 \cup \ldots \cup Z_r$. En effet, si $y \in Z_i$,
alors
$$
\delta(x) - \delta(y) = \delta(x) - \delta(\xi_i) + \delta(\xi_i) - \delta(y)
= - \text{codim}(\overline{\{x\}}, Z_i)
+ \text{codim}(\overline{\{y\}}, Z_i)
$$
d'après le Lemme \ref{topology-lemma-dimension-function-catenary}.
Il en va de même pour $\delta'$, d'où le résultat.
\end{proof}

\begin{lemma}
\label{topology-lemma-locally-dimension-function}
Soit $X$ un espace localement noethérien, sobre et caténaire.
Alors tout point possède un voisinage ouvert
$U \subset X$ qui admet une fonction de dimension.
\end{lemma}

\begin{proof}
Nous utiliserons à plusieurs reprises
le fait qu'un sous-espace ouvert d'un espace caténaire est caténaire, voir le
Lemme \ref{topology-lemma-catenary}, et qu'un espace topologique noethérien
n'a qu'un nombre fini de composantes irréductibles, voir le Lemme
\ref{topology-lemma-Noetherian}.
Dans la démonstration du Lemme \ref{topology-lemma-dimension-function-unique}, nous avons
vu comment construire une telle fonction. Tout d'abord, nous remplaçons $X$
par un voisinage ouvert noethérien de $x$. Ensuite, soient
$Z_1, \ldots, Z_r \subset X$ les composantes irréductibles de $X$. Soit
$$
Z_i \cap Z_j = \bigcup Z_{ijk}
$$
la décomposition en composantes irréductibles. Nous remplaçons
$X$ par
$$
X \setminus \left(
\bigcup\nolimits_{x \not \in Z_i} Z_i
\cup
\bigcup\nolimits_{x \not \in Z_{ijk}} Z_{ijk}
\right)
$$
de sorte que nous puissions supposer que $x \in Z_i$ pour tout $i$ et
$x \in Z_{ijk}$ pour tous $i, j, k$. Pour $y \in X$, choisissons un
$i$ quelconque tel que $y \in Z_i$ et posons
$$
\delta(y) = - \text{codim}(\overline{\{y\}}, Z_i)
+ \text{codim}(\overline{\{x\}}, Z_i).
$$
Nous affirmons qu'il s'agit d'une fonction de dimension. Montrons d'abord
qu'elle est bien définie, c'est-à-dire qu'elle ne dépend pas du choix de $i$.
Supposons en effet que $y \in Z_{ijk}$ pour certains $i, j, k$.
Alors nous avons (en utilisant le Lemme
\ref{topology-lemma-catenary-in-codimension})
\begin{align*}
\delta(y) & =
- \text{codim}(\overline{\{y\}}, Z_i)
+ \text{codim}(\overline{\{x\}}, Z_i) \\
& =
- \text{codim}(\overline{\{y\}}, Z_{ijk})
- \text{codim}(Z_{ijk}, Z_i)
+ \text{codim}(\overline{\{x\}}, Z_{ijk})
+ \text{codim}(Z_{ijk}, Z_i) \\
& =
- \text{codim}(\overline{\{y\}}, Z_{ijk})
+ \text{codim}(\overline{\{x\}}, Z_{ijk})
\end{align*}
ce qui est symétrique en $i$ et $j$.
Nous omettons la démonstration du fait qu'il s'agit d'une fonction de dimension.
\end{proof}

\begin{remark}
\label{topology-remark-obstruction-to-dimension-function}
En combinant les Lemmes \ref{topology-lemma-dimension-function-unique} et
\ref{topology-lemma-locally-dimension-function}, nous voyons que, sur un espace
topologique caténaire, localement noethérien et sobre, l'obstruction à
l'existence d'une fonction de dimension est un élément de
$H^1(X, \mathbf{Z})$.
\end{remark}



\section{Ensembles nulle part denses}
\label{topology-section-nowhere-dense}

\begin{definition}
\label{topology-definition-nowhere-dense}
Soit $X$ un espace topologique.
\begin{enumerate}
\item Étant donnée une partie $T \subset X$, l'{\it intérieur} de $T$ est le
plus grand ouvert de $X$ contenu dans $T$.
\item Une partie $T \subset X$ est dite {\it nulle part dense} si l'adhérence
de $T$ a un intérieur vide.
\end{enumerate}
\end{definition}

\begin{lemma}
\label{topology-lemma-nowhere-dense}
Soit $X$ un espace topologique. La réunion d'un nombre fini d'ensembles nulle
part denses est un ensemble nulle part dense.
\end{lemma}

\begin{proof}
Il suffit de démontrer le lemme pour deux ensembles nulle part denses, le cas
général en résultant par récurrence. Soient $A,B \subset X$ deux parties nulle
part denses. On a $\overline{A \cup B} = \overline{A} \cup \overline{B}$.
Ainsi, si $U \subset \overline{A \cup B}$ est un ouvert de $X$, alors
$U \setminus U \cap \overline{B}$ est un ouvert de $U$ et de $X$,
contenu dans $\overline{A}$, et est donc vide. De même,
$U \setminus U \cap \overline{A}$ est vide. Ainsi $U = \emptyset$,
comme voulu.
\end{proof}

\begin{lemma}
\label{topology-lemma-image-nowhere-dense-open}
Soit $X$ un espace topologique.
Soit $U \subset X$ un ouvert.
Soit $T \subset U$ une partie.
Si $T$ est nulle part dense dans $U$, alors $T$ est nulle part dense dans $X$.
\end{lemma}

\begin{proof}
Supposons que $T$ soit nulle part dense dans $U$.
Supposons que $x \in X$ soit un point intérieur à l'adhérence
$\overline{T}$ de $T$ dans $X$. Il existe alors
$x \in V \subset \overline{T}$, où $V \subset X$ est ouvert dans $X$.
Notons que $\overline{T} \cap U$ est
l'adhérence de $T$ dans $U$. Puisque l'intérieur de
$\overline{T} \cap U$ est vide, on a $V \cap U = \emptyset$. Par conséquent,
$x$ ne peut appartenir à l'adhérence de $U$, et a fortiori ne peut appartenir
à l'adhérence de $T$, ce qui est une contradiction.
\end{proof}

\begin{lemma}
\label{topology-lemma-nowhere-dense-local}
Soit $X$ un espace topologique.
Soit $X = \bigcup U_i$ un recouvrement ouvert.
Soit $T \subset X$ une partie.
Si $T \cap U_i$ est nulle part dense dans $U_i$ pour tout $i$,
alors $T$ est nulle part dense dans $X$.
\end{lemma}

\begin{proof}
Notons $\overline{T}_i$ l'adhérence de $T \cap U_i$ dans $U_i$.
On a $\overline{T} \cap U_i = \overline{T}_i$.
La prise de l'intérieur commute avec l'intersection par un ouvert, donc
$$
(\text{intérieur de }\overline{T}) \cap U_i =
\text{intérieur de }(\overline{T} \cap U_i) =
\text{intérieur dans }U_i\text{ de }\overline{T}_i
$$
Par hypothèse, le dernier de ces ensembles est vide. Ainsi,
$T$ est nulle part dense dans $X$.
\end{proof}

\begin{lemma}
\label{topology-lemma-closed-image-nowhere-dense}
Soit $f : X \to Y$ une application continue d'espaces topologiques.
Soit $T \subset X$ une partie.
Si $f$ est un homéomorphisme de $X$ sur une partie fermée de $Y$
et si $T$ est nulle part dense dans $X$, alors $f(T)$ est également nulle
part dense dans $Y$.
\end{lemma}

\begin{proof}
Puisque $f(X)$ est fermé dans $Y$ et que $f$ est un homéomorphisme de $X$
sur $f(X)$, l'adhérence de $f(T)$ dans $Y$ est égale à $f(\overline{T})$.
Ainsi, si $V \subset Y$ est ouvert et contenu dans l'adhérence de $f(T)$,
alors $U = f^{-1}(V)$ est ouvert et contenu dans $\overline{T}$.
Par conséquent, $U = \emptyset$, ce qui entraîne à son tour que
$V = \emptyset$, comme voulu.
\end{proof}

\begin{lemma}
\label{topology-lemma-open-inverse-image-closed-nowhere-dense}
Soit $f : X \to Y$ une application continue d'espaces topologiques.
Soit $T \subset Y$ une partie.
Si $f$ est ouverte et si $T$ est une partie fermée nulle part dense de $Y$,
alors $f^{-1}(T)$ est également une partie fermée nulle part dense de $X$.
Si $f$ est surjective et ouverte, alors
$T$ est fermée et nulle part dense si et seulement
si $f^{-1}(T)$ est fermée et nulle part dense.
\end{lemma}

\begin{proof}
Démonstration omise. (Indication : dans le premier cas, l'intérieur de
$f^{-1}(T)$ s'envoie dans l'intérieur de $T$, tandis que, dans le second cas,
l'intérieur de $f^{-1}(T)$ s'envoie sur l'intérieur de $T$.)
\end{proof}

\begin{lemma}
\label{topology-lemma-dense-in-constructible}
Soit $X$ un espace topologique. Soit $E \subset X$ une réunion finie de
parties localement fermées (par exemple, $E$ est constructible). Alors $E$
contient un ouvert dense de son adhérence.
\end{lemma}

\begin{proof}
Après avoir remplacé $X$ par l'adhérence de $E$, nous pouvons supposer que
$E$ est dense dans $X$ et devons montrer que $E$ contient un ouvert dense.
Le complémentaire $T$ de $E$ est alors une réunion finie de parties localement
fermées, disons $T = T_1 \cup \ldots \cup T_n$, et $T$ est nulle part dense
dans $X$. Puisque $T_i$ est un ouvert dense de son adhérence $Z_i$, nous voyons
que $Z_i$ est également nulle part dense dans $X$. L'adhérence
$Z = Z_1 \cup \ldots \cup Z_n$ de $T$ est donc nulle part dense dans $X$
d'après le Lemme \ref{topology-lemma-nowhere-dense}.
Nous en concluons que $E$ contient l'ouvert dense $X \setminus Z$.
\end{proof}





\section{Espaces profinis}
\label{topology-section-profinite}

\noindent
Voici la définition.

\begin{definition}
\label{topology-definition-profinite}
Un espace topologique est dit {\it profini} s'il est homéomorphe à la limite
d'un diagramme d'espaces discrets finis.
\end{definition}

\noindent
Ce n'est pas la caractérisation la plus commode d'un espace profini.

\begin{lemma}
\label{topology-lemma-profinite}
Soit $X$ un espace topologique.
Les assertions suivantes sont équivalentes :
\begin{enumerate}
\item $X$ est un espace profini, et
\item $X$ est séparé, quasi-compact et totalement discontinu.
\end{enumerate}
Si ces conditions sont satisfaites, alors $X$ est une limite cofiltrante
d'espaces discrets finis.
\end{lemma}

\begin{proof}
Supposons (1). Choisissons un diagramme $i \mapsto X_i$ d'espaces discrets finis
tel que $X = \lim X_i$. Comme chaque $X_i$ est séparé et quasi-compact, le
lemme \ref{topology-lemma-inverse-limit-quasi-compact} montre que $X$ est quasi-compact.
Si $x, x' \in X$ sont des points distincts, alors $x$ et $x'$ ont des images
distinctes dans un certain $X_i$. Par conséquent, $x$ et $x'$ possèdent des
voisinages ouverts disjoints, c'est-à-dire que $X$ est séparé. Exactement de la
même manière, nous voyons que $X$ est totalement discontinu.

\medskip\noindent
Supposons (2). Soit $\mathcal{I}$ l'ensemble des décompositions en réunion
disjointe finie $X = \coprod_{i \in I} U_i$, où $U_i$ est un ouvert non vide
(et fermé) pour tout $i \in I$.
Pour chaque $I \in \mathcal{I}$, il existe une application continue
$X \to I$ qui envoie tout point de $U_i$ sur $i$. Nous définissons un ordre
partiel : $I \leq I'$ pour $I, I' \in \mathcal{I}$ si et seulement si le
recouvrement correspondant à $I'$ raffine le recouvrement correspondant à
$I$. Dans ce cas, nous obtenons une application canonique $I' \to I$. En
d'autres termes, nous obtenons un système projectif d'espaces discrets finis
indexé par $\mathcal{I}$.
Les applications $X \to I$ sont compatibles et fournissent une application
continue
$$
X \longrightarrow \lim_{I \in \mathcal{I}} I
$$
Nous affirmons que cette application est un homéomorphisme, ce qui achève la
démonstration. (L'assertion finale en résulte également, car l'ensemble
partiellement ordonné $\mathcal{I}$ est dirigé : étant données deux
décompositions de $X$ en réunions disjointes, nous pouvons en trouver une
troisième qui raffine les deux.)
En effet, l'application est injective puisque $X$ est totalement discontinu
et que, par conséquent, $\{x\}$ est l'intersection de toutes les parties
ouvertes et fermées de $X$ qui contiennent $x$
(lemme \ref{topology-lemma-connected-component-intersection-compact-Hausdorff}) ;
l'application est surjective d'après le
lemme \ref{topology-lemma-intersection-closed-in-quasi-compact}.
D'après le lemme \ref{topology-lemma-bijective-map}, l'application est un homéomorphisme.
\end{proof}

\begin{lemma}
\label{topology-lemma-directed-inverse-limit-profinite}
Toute limite d'espaces profinis est un espace profini.
\end{lemma}

\begin{proof}
Soit $i \mapsto X_i$ un diagramme d'espaces profinis
indexé par la catégorie $\mathcal{I}$.
Utilisons la caractérisation des espaces profinis donnée par le
lemme \ref{topology-lemma-profinite}. En particulier, chaque $X_i$ est
séparé, quasi-compact et totalement discontinu.
D'après le lemme \ref{topology-lemma-limits}, la limite $X = \lim X_i$ existe.
D'après le lemme \ref{topology-lemma-inverse-limit-quasi-compact},
la limite $X$ est quasi-compacte. Soient $x, x' \in X$ des points distincts.
Il existe alors un $i$ tel que $x$ et $x'$ aient des images distinctes
$x_i$ et $x'_i$ dans $X_i$ par la projection $X \to X_i$.
Les points $x_i$ et $x'_i$ possèdent alors des voisinages ouverts disjoints
dans $X_i$. En prenant les images réciproques de ces ouverts, nous concluons
que $X$ est séparé.
De même, $x_i$ et $x'_i$ appartiennent à des composantes connexes distinctes
de $X_i$ ; nécessairement, $x$ et $x'$ appartiennent donc à des composantes
connexes distinctes de $X$. Ainsi, $X$ est totalement discontinu.
Cela achève la démonstration.
\end{proof}

\begin{lemma}
\label{topology-lemma-profinite-refine-open-covering}
Soit $X$ un espace profini. Tout recouvrement ouvert de $X$ admet un
raffinement par un recouvrement fini $X = \coprod U_i$ dont les $U_i$ sont
ouverts et fermés.
\end{lemma}

\begin{proof}
Écrivons $X = \lim X_i$ comme limite d'un système projectif d'espaces discrets
finis indexé par un ensemble dirigé $I$ (lemme \ref{topology-lemma-profinite}).
Notons $f_i : X \to X_i$ la projection.
Pour tout point $x = (x_i) \in X$, la famille des $f_i^{-1}(\{x_i\})$ forme
un système fondamental de voisinages ouverts. Ainsi, puisque $X$ est
quasi-compact, nous pouvons supposer donné un recouvrement ouvert
$$
X = f_{i_1}^{-1}(\{x_{i_1}\}) \cup \ldots \cup f_{i_n}^{-1}(\{x_{i_n}\})
$$
Choisissons $i \in I$ tel que $i \geq i_j$ pour $j = 1, \ldots, n$ (c'est
possible puisque $I$ est un ensemble dirigé). Nous voyons alors que le
recouvrement
$$
X = \coprod\nolimits_{t \in X_i} f_i^{-1}(\{t\})
$$
raffine le recouvrement donné et possède la forme voulue.
\end{proof}

\begin{lemma}
\label{topology-lemma-pi0-profinite}
Soit $X$ un espace topologique. Si $X$ est quasi-compact et si toute composante
connexe de $X$ est l'intersection des parties ouvertes et fermées qui la
contiennent, alors $\pi_0(X)$ est un espace profini.
\end{lemma}

\begin{proof}
Nous allons appliquer le lemme \ref{topology-lemma-profinite}.
Puisque $\pi_0(X)$ est l'image d'un espace quasi-compact, il est quasi-compact
(lemme \ref{topology-lemma-image-quasi-compact}).
Il est totalement discontinu par construction
(lemme \ref{topology-lemma-space-connected-components}).
Soient $C, D \subset X$ deux composantes connexes distinctes de $X$.
Écrivons $C = \bigcap U_\alpha$ comme intersection des parties ouvertes et
fermées de $X$ qui contiennent $C$. Toute intersection finie de $U_\alpha$
est encore une telle partie. Puisque $\bigcap U_\alpha \cap D = \emptyset$,
nous en concluons que $U_\alpha \cap D = \emptyset$ pour un certain $\alpha$
(utiliser les lemmes \ref{topology-lemma-connected-components},
\ref{topology-lemma-closed-in-quasi-compact} et
\ref{topology-lemma-intersection-closed-in-quasi-compact}).
Comme $U_\alpha$ est ouvert et fermé, il est la réunion des composantes
connexes qu'il contient ; autrement dit, $U_\alpha$ est l'image réciproque
d'une certaine partie ouverte et fermée $V_\alpha \subset \pi_0(X)$.
Cela montre que les points correspondant à $C$ et $D$ sont contenus dans des
ouverts disjoints, c'est-à-dire que $\pi_0(X)$ est séparé.
\end{proof}





\section{Espaces spectraux}
\label{topology-section-spectral}

\noindent
Le contenu de cette section est tiré de \cite{Hochster} et de
\cite{Hochster-thesis}. Dans sa thèse, Hochster démontre (entre autres)
que les espaces spectraux sont exactement les espaces topologiques qui
sont des spectres d'anneaux.

\begin{definition}
\label{topology-definition-spectral-space}
Un espace topologique $X$ est dit {\it spectral} s'il est sobre et
quasi-compact, si l'intersection de deux ouverts quasi-compacts est
quasi-compacte et si l'ensemble des ouverts quasi-compacts forme une
base de la topologie. Une application continue $f : X \to Y$ entre espaces
spectraux est dite {\it spectrale} si l'image réciproque de tout ouvert
quasi-compact est quasi-compacte.
\end{definition}

\noindent
Autrement dit, une application continue entre espaces spectraux est spectrale
si et seulement si elle est quasi-compacte (Définition
\ref{topology-definition-quasi-compact}).

\medskip\noindent
Soit $X$ un espace spectral. La {\it topologie constructible} sur $X$
est la topologie dont une sous-base d'ouverts est formée des ensembles $U$
et $U^c$, où $U$ est un ouvert quasi-compact de $X$. Puisque $X$ est spectral,
un ouvert $U \subset X$ est rétrocompact si et seulement si $U$ est
quasi-compact. La topologie constructible peut donc aussi être caractérisée
comme la topologie la moins fine pour laquelle toute partie constructible de
$X$ est à la fois ouverte et fermée (voir la section
\ref{topology-section-constructible} pour les définitions et les propriétés des parties
constructibles).
Il s'ensuit qu'une partie de $X$ est ouverte, resp.\ fermée, pour la topologie
constructible si et seulement si elle est une réunion, resp.\ une intersection,
de parties constructibles. Comme l'ensemble des ouverts quasi-compacts est une
base de la topologie de $X$, la topologie constructible est plus fine que la
topologie donnée sur $X$.

\begin{lemma}
\label{topology-lemma-constructible-hausdorff-quasi-compact}
Soit $X$ un espace spectral. La topologie constructible est séparée,
totalement discontinue et quasi-compacte.
\end{lemma}

\begin{proof}
Soient $x, y \in X$ avec $x \not = y$. Comme $X$ est sobre, il existe
un ouvert $U$ qui contient exactement l'un des deux points $x, y$.
Supposons que $x \in U$. On peut remplacer $U$ par un voisinage ouvert
quasi-compact de $x$ contenu dans $U$. Alors $U$ et $U^c$ sont ouverts et
fermés pour la topologie constructible. Ainsi, $X$ est séparé pour la
topologie constructible, car $x \in U$ et $y \in U^c$ appartiennent à des
ouverts disjoints de cette topologie. L'existence de $U$ implique aussi que
$x$ et $y$ appartiennent à des composantes connexes distinctes pour la
topologie constructible ; ainsi, $X$ est totalement discontinu pour cette
topologie.

\medskip\noindent
Soit $\mathcal{B}$ l'ensemble des parties $B \subset X$ telles que $B$ soit
un ouvert quasi-compact ou un fermé dont le complémentaire est
quasi-compact. Si $B \in \mathcal{B}$, alors $B^c \in \mathcal{B}$.
Il suffit de montrer que tout recouvrement $X = \bigcup_{i \in I} B_i$
avec $B_i \in \mathcal{B}$ admet un raffinement fini, voir le
Lemme \ref{topology-lemma-subbase-theorem}.
En passant aux complémentaires, on voit qu'il faut montrer que toute famille
$\{B_i\}_{i \in I}$ d'éléments de $\mathcal{B}$ telle que
$B_{i_1} \cap \ldots \cap B_{i_n} \not = \emptyset$
pour tout $n$ et tous $i_1, \ldots, i_n \in I$
possède un point commun. On peut supposer, et l'on suppose,
que $B_i \not = B_{i'}$ si $i \not = i'$.

\medskip\noindent
Raisonnons par l'absurde et supposons que $\{B_i\}_{i \in I}$ soit une famille
d'éléments de $\mathcal{B}$ ayant la propriété d'intersection finie, mais dont
l'intersection est vide. Une application du lemme de Zorn montre que l'on
peut supposer cette famille maximale (les détails sont omis).
Soit $I' \subset I$ l'ensemble des indices pour lesquels $B_i$ est fermé,
et posons $Z = \bigcap_{i \in I'} B_i$.
C'est une partie fermée non vide de $X$, d'après le Lemme
\ref{topology-lemma-intersection-closed-in-quasi-compact}.
Si $Z$ est réductible, on peut écrire $Z = Z' \cup Z''$
comme réunion de deux parties fermées, toutes deux distinctes de $Z$. On peut
donc, en particulier, trouver un ouvert quasi-compact $U' \subset X$ qui
rencontre $Z'$ mais pas $Z''$. De même, on peut trouver un ouvert quasi-compact
$U'' \subset X$ qui rencontre $Z''$ mais pas $Z'$. Posons
$B' = X \setminus U'$ et $B'' = X \setminus U''$. Remarquons que
$Z'' \subset B'$ et $Z' \subset B''$.
S'il existe un nombre fini d'indices $i_1, \ldots, i_n \in I$ tels que
$B' \cap B_{i_1} \cap \ldots \cap B_{i_n} = \emptyset$
ainsi qu'un nombre fini d'indices $j_1, \ldots, j_m \in I$ tels que
$B'' \cap B_{j_1} \cap \ldots \cap B_{j_m} = \emptyset$,
alors on obtient
$Z \cap B_{i_1} \cap \ldots \cap B_{i_n} \cap B_{j_1} \cap \ldots \cap B_{j_m}
= \emptyset$.
Cependant, l'ensemble
$B_{i_1} \cap \ldots \cap B_{i_n} \cap B_{j_1} \cap \ldots \cap B_{j_m}$
est quasi-compact ; il existerait donc un nombre fini d'indices
$i'_1, \ldots, i'_l \in I'$ tels que
$B_{i_1} \cap \ldots \cap B_{i_n} \cap B_{j_1} \cap \ldots \cap
B_{j_m} \cap B_{i'_1} \cap \ldots \cap B_{i'_l} = \emptyset$, ce qui est
contradictoire. On voit donc que l'on peut adjoindre soit $B'$, soit $B''$,
à la famille donnée, contrairement à sa maximalité. On en conclut que $Z$
est irréductible. Mais cela conduit encore à une contradiction, car tout
ouvert non vide $Z \cap B_i$, pour $i \in I \setminus I'$, contient alors
(par le même argument que ci-dessus) l'unique point générique de $Z$.
Cette contradiction démontre le lemme.
\end{proof}

\begin{lemma}
\label{topology-lemma-fibres-spectral-map-quasi-compact}
Soit $f : X \to Y$ une application spectrale entre espaces spectraux. Alors
\begin{enumerate}
\item $f$ est continue pour les topologies constructibles,
\item les fibres de $f$ sont quasi-compactes, et
\item l'image est fermée pour la topologie constructible.
\end{enumerate}
\end{lemma}

\begin{proof}
Notons $X'$ et $Y'$ les espaces $X$ et $Y$ munis de leur topologie
constructible ; ce sont des espaces quasi-compacts séparés d'après le
Lemme \ref{topology-lemma-constructible-hausdorff-quasi-compact}.
L'assertion (1) affirme que $X' \to Y'$ est continue et résulte immédiatement
des définitions. L'assertion (3) résulte du fait que $f(X')$ est une partie
quasi-compacte de l'espace séparé $Y'$, voir le
Lemme \ref{topology-lemma-quasi-compact-in-Hausdorff}.
On a un diagramme commutatif
$$
\xymatrix{
X' \ar[r] \ar[d] & X \ar[d] \\
Y' \ar[r] & Y
}
$$
d'applications continues entre espaces topologiques. Comme $Y'$ est séparé,
les fibres $X'_y$ sont fermées dans $X'$. Puisque $X'$ est quasi-compact,
$X'_y$ est quasi-compact
(Lemme \ref{topology-lemma-closed-in-quasi-compact}).
Comme $X'_y \to X_y$ est une application continue surjective, on en conclut
que $X_y$ est quasi-compact (Lemme \ref{topology-lemma-image-quasi-compact}).
\end{proof}

\begin{lemma}
\label{topology-lemma-spectral-if-continuous-wrt-constructible-top}
Soient $X$ et $Y$ des espaces spectraux et $f : X \to Y$ une application
continue. Alors $f$ est spectrale si et seulement si $f$ est continue pour
les topologies constructibles.
\end{lemma}

\begin{proof}
L'implication directe est le
Lemme \ref{topology-lemma-fibres-spectral-map-quasi-compact}.
Supposons $f$ continue pour les topologies constructibles.
Soit $V \subset Y$ un ouvert quasi-compact.
Alors $V$ est ouvert et fermé pour la topologie constructible.
Ainsi, $f^{-1}(V)$ est ouvert et fermé pour la topologie constructible.
Ainsi, $f^{-1}(V)$ est quasi-compact pour la topologie constructible, puisque
$X$ est quasi-compact pour cette topologie d'après le
Lemme \ref{topology-lemma-constructible-hausdorff-quasi-compact}.
L'application identité $f^{-1}(V) \to f^{-1}(V)$ étant continue et surjective
de la topologie constructible vers la topologie usuelle, on conclut que
$f^{-1}(V)$ est quasi-compact pour la topologie de $X$, d'après le
Lemme \ref{topology-lemma-image-quasi-compact}.
Ceci achève la démonstration.
\end{proof}

\begin{lemma}
\label{topology-lemma-spectral-sub}
Soit $X$ un espace spectral. Soit $E \subset X$ une partie fermée pour la
topologie constructible (par exemple une partie constructible ou fermée).
Alors $E$, muni de la topologie induite, est un espace spectral.
\end{lemma}

\begin{proof}
Soit $Z \subset E$ une partie fermée irréductible. Soit $\eta$ le point
générique de l'adhérence $\overline{Z}$ de $Z$ dans $X$. Pour démontrer que
$E$ est sobre, montrons que $\eta \in E$. Sinon, puisque $E$ est fermé pour
la topologie constructible, il existe une partie constructible $F \subset X$
telle que $\eta \in F$ et $F \cap E = \emptyset$.
D'après le Lemme \ref{topology-lemma-generic-point-in-constructible}, cela implique que
$F \cap \overline{Z}$ contient un ouvert non vide de $\overline{Z}$.
Mais c'est impossible puisque $\overline{Z}$ est l'adhérence de $Z$ et que
$Z \cap F = \emptyset$.

\medskip\noindent
Puisque $E$ est fermé pour la topologie constructible, il est quasi-compact
pour cette topologie
(Lemmes \ref{topology-lemma-closed-in-quasi-compact} et
\ref{topology-lemma-constructible-hausdorff-quasi-compact}). Il est donc, à plus forte
raison, quasi-compact pour la topologie induite par celle de $X$. Si
$U \subset X$ est un ouvert quasi-compact, alors $E \cap U$ est fermé pour la
topologie constructible, donc quasi-compact (comme on vient de le voir). Il en
résulte que les ouverts quasi-compacts de $E$ sont les intersections $E \cap U$
où $U$ est un ouvert quasi-compact de $X$. Ils forment une base de la topologie.
Enfin, étant donnés deux ouverts quasi-compacts $U, U' \subset X$, on a
$(E \cap U) \cap (E \cap U') = E \cap (U \cap U')$, et $U \cap U'$ est
quasi-compact puisque $X$ est spectral. Ceci achève la démonstration.
\end{proof}

\begin{lemma}
\label{topology-lemma-constructible-stable-specialization-closed}
Soit $X$ un espace spectral. Soit $E \subset X$ une partie fermée pour la
topologie constructible (par exemple une partie constructible).
\begin{enumerate}
\item Si $x \in \overline{E}$, alors $x$ est une spécialisation d'un point de
$E$.
\item Si $E$ est stable par spécialisation, alors $E$ est fermée.
\item Si $E' \subset X$ est ouverte pour la topologie constructible
(par exemple constructible) et stable par générisation, alors $E'$ est ouverte.
\end{enumerate}
\end{lemma}

\begin{proof}
Démonstration de (1). Soit $x \in \overline{E}$. Soit $\{U_i\}$ l'ensemble
des voisinages ouverts quasi-compacts de $x$. Toute intersection finie de
$U_i$ est encore un tel voisinage. L'intersection $U_i \cap E$ est non vide
pour tout $i$. Comme les parties $U_i \cap E$ sont fermées pour la topologie
constructible, l'intersection $\bigcap (U_i \cap E)$ est non vide d'après le
Lemme \ref{topology-lemma-constructible-hausdorff-quasi-compact} et le
Lemme \ref{topology-lemma-intersection-closed-in-quasi-compact}.
Puisque $\{U_i\}$ est un système fondamental de voisinages ouverts de $x$,
l'intersection $\bigcap U_i$ est l'ensemble des générisations de $x$.
Ainsi, $x$ est une spécialisation d'un point de $E$.

\medskip\noindent
L'assertion (2) résulte immédiatement de (1).

\medskip\noindent
Démonstration de (3). Supposons que $E'$ vérifie les hypothèses de (3).
Le complémentaire de $E'$ est fermé pour la topologie constructible
(Lemme \ref{topology-lemma-constructible}) et stable par spécialisation
(Lemme \ref{topology-lemma-open-closed-specialization}).
Ce complémentaire est donc fermé d'après (2), c'est-à-dire que $E'$ est ouverte.
\end{proof}

\begin{lemma}
\label{topology-lemma-two-points}
Soit $X$ un espace spectral et soient $x, y \in X$. Alors, ou bien il existe
un troisième point dont $x$ et $y$ sont des spécialisations, ou bien il existe
des voisinages ouverts disjoints de $x$ et de $y$.
\end{lemma}

\begin{proof}
Soit $\{U_i\}$ l'ensemble des voisinages ouverts quasi-compacts de $x$.
Toute intersection finie de $U_i$ est encore un tel voisinage.
Soit $\{V_j\}$ l'ensemble des voisinages ouverts quasi-compacts de $y$.
Toute intersection finie de $V_j$ est encore un tel voisinage.
Si $U_i \cap V_j$ est vide pour certains $i, j$, la conclusion est acquise.
Sinon, $U_i \cap V_j$ est non vide pour tous $i$ et $j$. Les ensembles
$U_i \cap V_j$ sont fermés pour la topologie constructible de $X$. D'après le
Lemme \ref{topology-lemma-constructible-hausdorff-quasi-compact},
l'intersection $\bigcap (U_i \cap V_j)$ est non vide
(Lemme \ref{topology-lemma-intersection-closed-in-quasi-compact}).
Comme $X$ est sobre et que $\{U_i\}$ est un système fondamental de voisinages
ouverts de $x$, l'intersection $\bigcap U_i$ est l'ensemble des générisations
de $x$. De même, $\bigcap V_j$ est l'ensemble des générisations de $y$.
Ainsi, tout élément de $\bigcap (U_i \cap V_j)$ a pour spécialisations à la
fois $x$ et $y$.
\end{proof}

\begin{lemma}
\label{topology-lemma-characterize-profinite-spectral}
Soit $X$ un espace spectral. Les propriétés suivantes sont équivalentes :
\begin{enumerate}
\item $X$ est profini,
\item $X$ est séparé,
\item $X$ est totalement discontinu,
\item tout ouvert quasi-compact est fermé,
\item il n'existe aucune spécialisation non triviale entre points,
\item tout point de $X$ est fermé,
\item tout point de $X$ est le point générique d'une composante irréductible
de $X$,
\item la topologie constructible est égale à la topologie donnée sur $X$, et
\item ajouter quelque chose ici.
\end{enumerate}
\end{lemma}

\begin{proof}
Le Lemme \ref{topology-lemma-profinite} donne l'implication (1) $\Rightarrow$ (3).
Les composantes irréductibles sont fermées ; si $X$ est totalement discontinu,
tout point est donc fermé. Ainsi, (3) implique (6). L'équivalence de (6) et
(5) est immédiate, et (6) $\Leftrightarrow$ (7) parce que $X$ est sobre.
Supposons (5). Alors toutes les parties constructibles de $X$ sont fermées
(Lemme \ref{topology-lemma-constructible-stable-specialization-closed}), et en
particulier tous les ouverts quasi-compacts sont fermés. Donc (5) implique (4).
Comme $X$ est sobre, pour deux points quelconques il existe un ouvert
quasi-compact qui contient exactement l'un des deux ; par conséquent, (4)
implique (2). Les assertions (4) et (8) sont équivalentes par définition de la
topologie constructible.
Il reste à démontrer que (2) implique (1). Supposons $X$ séparé. Tout ouvert
quasi-compact est alors fermé
(Lemme \ref{topology-lemma-quasi-compact-in-Hausdorff}). Cela implique que $X$ est
totalement discontinu. Il est donc profini d'après le
Lemme \ref{topology-lemma-profinite}.
\end{proof}

\begin{lemma}
\label{topology-lemma-spectral-pi0}
Si $X$ est un espace spectral, alors $\pi_0(X)$ est un espace profini.
\end{lemma}

\begin{proof}
Combiner les Lemmes \ref{topology-lemma-connected-component-intersection} et
\ref{topology-lemma-pi0-profinite}.
\end{proof}

\begin{lemma}
\label{topology-lemma-product-spectral-spaces}
Le produit de deux espaces spectraux est spectral.
\end{lemma}

\begin{proof}
Soient $X$ et $Y$ des espaces spectraux. Notons $p : X \times Y \to X$ et
$q : X \times Y \to Y$ les projections. Soit $Z \subset X \times Y$ une
partie fermée irréductible. Alors $p(Z) \subset X$ et $q(Z) \subset Y$ sont
irréductibles. Soit $x \in X$ le point générique de l'adhérence de $p(X)$ et
soit $y \in Y$ le point générique de l'adhérence de $q(Y)$. Si
$(x, y) \not \in Z$, il existe des ouverts $x \in U \subset X$ et
$y \in V \subset Y$ tels que $Z \cap U \times V = \emptyset$. Ainsi, $Z$
est contenu dans
$(X \setminus U) \times Y \cup X \times (Y \setminus V)$.
Comme $Z$ est irréductible, on a soit
$Z \subset (X \setminus U) \times Y$, soit
$Z \subset X \times (Y \setminus V)$.
Dans le premier cas, $p(Z) \subset (X \setminus U)$ et, dans le second,
$q(Z) \subset (Y \setminus V)$. Les deux cas sont absurdes, puisque $x$
appartient à l'adhérence de $p(Z)$ et $y$ à celle de $q(Z)$. On en conclut que
$(x, y) \in Z$, ce qui signifie que $(x, y)$ est le point générique de $Z$.

\medskip\noindent
Une base de la topologie de $X \times Y$ est formée des ouverts de la forme
$U \times V$, où $U \subset X$ et $V \subset Y$ sont des ouverts
quasi-compacts (on utilise ici le fait que $X$ et $Y$ sont spectraux).
Alors $U \times V$ est quasi-compact, car le produit de deux espaces
quasi-compacts est quasi-compact. De plus, tout ouvert quasi-compact de
$X \times Y$ est réunion finie de tels rectangles quasi-compacts $U \times V$.
Il s'ensuit que l'intersection de deux tels ouverts est encore quasi-compacte
(puisque $X$ et $Y$ sont spectraux). Ceci achève la démonstration.
\end{proof}

\begin{lemma}
\label{topology-lemma-spectral-bijective}
Soit $f : X \to Y$ une application continue entre espaces topologiques.
Supposons que
\begin{enumerate}
\item $X$ et $Y$ soient spectraux,
\item $f$ soit spectrale et bijective, et
\item les générisations (resp.\ les spécialisations) se relèvent le long de
$f$.
\end{enumerate}
Alors $f$ est un homéomorphisme.
\end{lemma}

\begin{proof}
Puisque $f$ est spectrale, elle définit une application continue entre $X$ et
$Y$ munis de leurs topologies constructibles. D'après les
Lemmes \ref{topology-lemma-constructible-hausdorff-quasi-compact} et
\ref{topology-lemma-bijective-map}, il en résulte que $X \to Y$ est un homéomorphisme
pour les topologies constructibles. Soit $U \subset X$ un ouvert
quasi-compact. Alors $f(U)$ est constructible dans $Y$. Soit $y \in Y$ une
spécialisation d'un point de $f(U)$. La dernière hypothèse montre que
$f^{-1}(y)$ est une spécialisation d'un point de $U$. Ainsi
$f^{-1}(y) \in U$, et donc $y \in f(U)$.
Il en résulte que $f(U)$ est ouvert, voir le
Lemme \ref{topology-lemma-constructible-stable-specialization-closed}.
Par conséquent, $f$ est un homéomorphisme.
Pour démontrer le lemme lorsque les spécialisations se relèvent le long de
$f$, on montre plutôt que $f(Z)$ est fermé si $X \setminus Z$ est un ouvert
quasi-compact de $X$.
\end{proof}

\begin{lemma}
\label{topology-lemma-directed-inverse-limit-finite-sober-spectral-spaces}
La limite projective d'un système projectif filtrant d'espaces topologiques
finis et sobres est un espace topologique spectral.
\end{lemma}

\begin{proof}
Soit $I$ un ensemble préordonné filtrant. Soit $X_i$ un système projectif
d'espaces finis et sobres indexé par $I$. Soit $X = \lim X_i$, qui existe
d'après le Lemme \ref{topology-lemma-limits}. En tant qu'ensemble, $X = \lim X_i$.
Notons $p_i : X \to X_i$ la projection.
Comme $I$ est filtrant, on peut appliquer le
Lemme \ref{topology-lemma-describe-limits}.
Les ouverts $p_i^{-1}(U_i)$, où $U_i \subset X_i$ est ouvert, forment une base
de la topologie. Comme tout recouvrement ouvert de $p_i^{-1}(U_i)$ est en
particulier un recouvrement ouvert pour la topologie profinie, on en conclut
que $p_i^{-1}(U_i)$ est quasi-compact. Si $U_i \subset X_i$ et
$U_j \subset X_j$, alors
$p_i^{-1}(U_i) \cap p_j^{-1}(U_j) = p_k^{-1}(U_k)$
pour un certain $k \geq i, j$ et un ouvert $U_k \subset X_k$. Enfin, si
$Z \subset X$ est fermée et irréductible, alors $p_i(Z) \subset X_i$ est
irréductible et possède donc un unique point générique $\xi_i$ (puisque $X_i$
est un espace topologique fini et sobre). Alors $\xi = \lim \xi_i$ est un
point générique de $Z$ (c'est un point de $Z$ puisque $Z$ est fermé).
Ceci achève la démonstration.
\end{proof}

\begin{lemma}
\label{topology-lemma-spectral-closed-in-product-two-point-space}
Soit $W$ l'espace topologique à deux points dont l'un est fermé et l'autre ne
l'est pas. Un espace topologique est spectral si et seulement s'il est
homéomorphe à un sous-espace d'un produit de copies de $W$ qui est fermé pour
la topologie constructible.
\end{lemma}

\begin{proof}
Écrivons $W = \{0, 1\}$, où $0$ est une spécialisation de $1$, mais non
réciproquement. Soit $I$ un ensemble. L'espace $\prod_{i \in I} W$ est spectral
d'après le
Lemme \ref{topology-lemma-directed-inverse-limit-finite-sober-spectral-spaces}.
On voit donc qu'un sous-espace de $\prod_{i \in I} W$ fermé pour la topologie
constructible est un espace spectral d'après le
Lemme \ref{topology-lemma-spectral-sub}.

\medskip\noindent
Réciproquement, soit $X$ un espace spectral. Soit $U \subset X$ un ouvert
quasi-compact. Considérons l'application continue
$$
f_U : X \longrightarrow W
$$
qui envoie tout point de $U$ sur $1$ et tout point de $X \setminus U$ sur $0$.
En prenant le produit de ces applications, on obtient une application continue
$$
f = \prod f_U : X \longrightarrow \prod\nolimits_U W
$$
Par construction, l'application $f : X \to Y$ est spectrale. D'après le
Lemme \ref{topology-lemma-fibres-spectral-map-quasi-compact}, l'image de $f$ est fermée
pour la topologie constructible. Si $x', x \in X$ sont distincts, alors, puisque
$X$ est sobre, ou bien $x'$ n'est pas une spécialisation de $x$, ou bien
ce dernier n'est pas une spécialisation du premier. Dans les deux cas (les ouverts
quasi-compacts formant une
base de la topologie de $X$), il existe un ouvert quasi-compact $U \subset X$
tel que $f_U(x') \not = f_U(x)$. Ainsi, $f$ est injective. Soit $Y = f(X)$,
muni de la topologie induite. Soit $y' \leadsto y$ une spécialisation dans $Y$,
et écrivons $f(x') = y'$ et $f(x) = y$. Le même raisonnement montre que
$x' \leadsto x$, car, sinon, il existerait un $U$ tel que $x \in U$ et
$x' \not \in U$, ce qui impliquerait
$f_U(x') \not \leadsto f_U(x)$.
On conclut que $f : X \to Y$ est un homéomorphisme d'après le
Lemme \ref{topology-lemma-spectral-bijective}.
\end{proof}

\begin{lemma}
\label{topology-lemma-spectral-inverse-limit-finite-sober-spaces}
Un espace topologique est spectral si et seulement s'il est limite projective
filtrante d'espaces topologiques finis et sobres.
\end{lemma}

\begin{proof}
Une implication est donnée par le
Lemme \ref{topology-lemma-directed-inverse-limit-finite-sober-spectral-spaces}.
Réciproquement, supposons $X$ spectral. On peut alors supposer que
$X \subset \prod_{i \in I} W$ est une partie fermée pour la topologie
constructible, où $W = \{0, 1\}$ est comme dans le
Lemme \ref{topology-lemma-spectral-closed-in-product-two-point-space}.
On peut écrire
$$
\prod\nolimits_{i \in I} W =
\lim_{J \subset I\text{ fini }} \prod\nolimits_{j \in J} W
$$
comme limite cofiltrante.
Pour chaque $J$, soit $X_J \subset \prod_{j \in J} W$ l'image de $X$.
On voit alors que $X = \lim X_J$ en tant qu'ensembles parce que $X$ est fermé
dans le produit muni de la topologie constructible (détail omis).
Un argument formel (omis) sur les limites montre que $X = \lim X_J$ comme
espaces topologiques.
\end{proof}

\begin{lemma}
\label{topology-lemma-Noetherian-goes-to-spectral}
Soit $X$ un espace topologique et soit $c : X \to X'$ l'application universelle
de $X$ vers un espace topologique sobre, voir le
Lemme \ref{topology-lemma-make-sober}.
\begin{enumerate}
\item Si $X$ est quasi-compact, alors $X'$ l'est aussi.
\item Si $X$ est quasi-compact, possède une base d'ouverts quasi-compacts et si
l'intersection de deux ouverts quasi-compacts est quasi-compacte, alors $X'$ est
spectral.
\item Si $X$ est noethérien, alors $X'$ est un espace spectral noethérien.
\end{enumerate}
\end{lemma}

\begin{proof}
Soit $U \subset X$ un ouvert et soit $U' \subset X'$ l'ouvert correspondant,
c'est-à-dire l'ouvert tel que $c^{-1}(U') = U$.
Alors $U$ est quasi-compact si et seulement si $U'$ est quasi-compact, car
l'image réciproque par $c$ établit une bijection entre les ouverts de $X$ et
ceux de $X'$ qui commute aux réunions. Cela démontre en particulier (1).

\medskip\noindent
Démonstration de (2). Il résulte de ce qui précède que $X'$ possède une base
d'ouverts quasi-compacts. Comme $c^{-1}$ commute aussi aux intersections de
paires d'ouverts, on voit que l'intersection de deux ouverts quasi-compacts de
$X'$ est quasi-compacte. Enfin, $X'$ est quasi-compact d'après (1) et sobre par
construction. Ainsi, $X'$ est spectral.

\medskip\noindent
Démonstration de (3). Il est immédiat que $X'$ est noethérien, puisque cette
propriété se définit par la condition de chaîne ascendante sur les ouverts,
laquelle est satisfaite pour $X$. On a déjà vu en (2) que $X'$ est spectral.
\end{proof}

\section{Limites d'espaces spectraux}
\label{topology-section-limits-spectral}

\noindent
Le Lemme \ref{topology-lemma-spectral-inverse-limit-finite-sober-spaces} nous dit
que tout espace spectral est une limite cofiltrante d'espaces finis
sobres. Tout espace fini sobre est un espace spectral, et toute application
continue entre espaces finis sobres est une application spectrale entre espaces
spectraux. Dans cette section, nous démontrons quelques lemmes concernant les
limites de systèmes d'espaces topologiques spectraux dont les applications de
transition sont spectrales.

\begin{lemma}
\label{topology-lemma-inverse-limit-spectral-spaces-quasi-compact}
Soit $\mathcal{I}$ une catégorie. Soit $i \mapsto X_i$ un diagramme
d'espaces spectraux tel que, pour $a : j \to i$ dans $\mathcal{I}$,
l'application correspondante $f_a : X_j \to X_i$ soit spectrale.
\begin{enumerate}
\item Si l'on se donne des parties $Z_i \subset X_i$ fermées pour la topologie
constructible telles que $f_a(Z_j) \subset Z_i$ pour tout $a : j \to i$ dans
$\mathcal{I}$, alors $\lim Z_i$ est quasi-compact.
\item L'espace $X = \lim X_i$ est quasi-compact.
\end{enumerate}
\end{lemma}

\begin{proof}
La limite $Z = \lim Z_i$ existe d'après le Lemme \ref{topology-lemma-limits}.
Notons $X'_i$ l'espace $X_i$ muni de la topologie constructible
et $Z'_i$ le sous-espace correspondant de $X'_i$.
Soit $a : j \to i$ un morphisme de $\mathcal{I}$. Comme $f_a$ est spectrale,
elle définit une application continue $f_a : X'_j \to X'_i$. Ainsi,
$f_a|_{Z_j} : Z'_j \to Z'_i$ est une application continue entre espaces
quasi-compacts séparés (d'après les
Lemmes \ref{topology-lemma-constructible-hausdorff-quasi-compact} et
\ref{topology-lemma-closed-in-quasi-compact}).
Ainsi, $Z' = \lim Z_i$ est quasi-compact d'après le
Lemme \ref{topology-lemma-inverse-limit-quasi-compact}.
Les applications $Z'_i \to Z_i$ sont continues ; par conséquent,
$Z' \to Z$ est continue et induit une bijection entre les ensembles
sous-jacents. Ainsi,
$Z$ est quasi-compact comme image de l'application continue surjective
$Z' \to Z$ (Lemme \ref{topology-lemma-image-quasi-compact}).
\end{proof}

\begin{lemma}
\label{topology-lemma-inverse-limit-spectral-spaces-nonempty}
Soit $\mathcal{I}$ une catégorie cofiltrante. Soit $i \mapsto X_i$ un diagramme
d'espaces spectraux tel que, pour $a : j \to i$ dans $\mathcal{I}$,
l'application correspondante $f_a : X_j \to X_i$ soit spectrale.
\begin{enumerate}
\item Si l'on se donne des parties non vides $Z_i \subset X_i$ fermées pour
la topologie constructible telles que $f_a(Z_j) \subset Z_i$ pour tout
$a : j \to i$ dans $\mathcal{I}$, alors $\lim Z_i$ est non vide.
\item Si chaque $X_i$ est non vide, alors $X = \lim X_i$ est non vide.
\end{enumerate}
\end{lemma}

\begin{proof}
Notons $X'_i$ l'espace $X_i$ muni de la topologie constructible
et $Z'_i$ le sous-espace correspondant de $X'_i$.
Soit $a : j \to i$ un morphisme de $\mathcal{I}$. Comme $f_a$ est spectrale,
elle définit une application continue $f_a : X'_j \to X'_i$. Ainsi,
$f_a|_{Z_j} : Z'_j \to Z'_i$ est une application continue entre espaces
quasi-compacts séparés (d'après les
Lemmes \ref{topology-lemma-constructible-hausdorff-quasi-compact} et
\ref{topology-lemma-closed-in-quasi-compact}). D'après le
Lemme \ref{topology-lemma-nonempty-limit}, l'espace $\lim Z'_i$ est non vide.
Puisque $\lim Z'_i = \lim Z_i$ comme ensembles, on conclut.
\end{proof}

\begin{lemma}
\label{topology-lemma-inverse-limit-spectral-spaces-equal}
Soit $\mathcal{I}$ une catégorie cofiltrante. Soit $i \mapsto X_i$ un diagramme
d'espaces spectraux tel que, pour $a : j \to i$ dans $\mathcal{I}$,
l'application correspondante $f_a : X_j \to X_i$ soit spectrale. Soit
$X = \lim X_i$, avec pour projections $p_i : X \to X_i$. Soient
$i \in \Ob(\mathcal{I})$ et $E, F \subset X_i$ des parties telles que $E$ soit
fermée et $F$ ouverte pour la topologie constructible.
Alors $p_i^{-1}(E) \subset p_i^{-1}(F)$ si et seulement s'il existe un morphisme
$a : j \to i$ dans $\mathcal{I}$ tel que
$f_a^{-1}(E) \subset f_a^{-1}(F)$.
\end{lemma}

\begin{proof}
Observons que
$$
p_i^{-1}(E) \setminus p_i^{-1}(F) =
\lim_{a : j \to i} f_a^{-1}(E) \setminus f_a^{-1}(F)
$$
Puisque $f_a$ est une application spectrale, elle est continue pour la
topologie constructible ; la partie
$f_a^{-1}(E) \setminus f_a^{-1}(F)$ est donc fermée pour la topologie
constructible. Le
Lemme \ref{topology-lemma-inverse-limit-spectral-spaces-nonempty}
montre alors que le membre de gauche est non vide si et seulement si chacun
des espaces du membre de droite est non vide.
\end{proof}

\begin{lemma}
\label{topology-lemma-inverse-limit-spectral-spaces-constructible}
Soit $\mathcal{I}$ une catégorie cofiltrante. Soit $i \mapsto X_i$ un diagramme
d'espaces spectraux tel que, pour $a : j \to i$ dans $\mathcal{I}$,
l'application correspondante $f_a : X_j \to X_i$ soit spectrale. Soit
$X = \lim X_i$, avec pour projections $p_i : X \to X_i$. Soit
$E \subset X$ une partie constructible. Il existe alors un
$i \in \Ob(\mathcal{I})$ et une partie constructible $E_i \subset X_i$ tels que
$p_i^{-1}(E_i) = E$. Si $E$ est ouverte, resp.\ fermée, on peut choisir $E_i$
ouverte, resp.\ fermée.
\end{lemma}

\begin{proof}
Supposons que $E$ soit un ouvert quasi-compact de $X$. D'après le
Lemme \ref{topology-lemma-describe-limits}, on peut écrire $E = p_i^{-1}(U_i)$
pour un certain $i$ et un certain ouvert $U_i \subset X_i$.
Écrivons $U_i = \bigcup U_{i, \alpha}$ comme réunion d'ouverts quasi-compacts.
Comme $E$ est quasi-compact, on peut trouver $\alpha_1, \ldots, \alpha_n$
tels que $E = p_i^{-1}(U_{i, \alpha_1} \cup \ldots \cup U_{i, \alpha_n})$.
Ainsi, $E_i = U_{i, \alpha_1} \cup \ldots \cup U_{i, \alpha_n}$ convient.

\medskip\noindent
Supposons que $E$ soit une partie constructible fermée. Alors $E^c$ est un
ouvert quasi-compact. Ainsi, $E^c = p_i^{-1}(F_i)$ pour un certain $i$ et
un ouvert quasi-compact $F_i \subset X_i$, d'après le résultat du paragraphe
précédent. Alors $E = p_i^{-1}(F_i^c)$, comme voulu.

\medskip\noindent
Si $E$ est quelconque, on peut écrire
$E = \bigcup_{l = 1, \ldots, n} U_l \cap Z_l$, où $U_l$ est un ouvert
constructible et $Z_l$ un fermé constructible. D'après les paragraphes
précédents, on peut écrire $U_l = p_{i_l}^{-1}(U_{l, i_l})$
et $Z_l = p_{j_l}^{-1}(Z_{l, j_l})$,
où $U_{l, i_l} \subset X_{i_l}$ est un ouvert constructible et
$Z_{l, j_l} \subset X_{j_l}$ un fermé constructible. Comme $\mathcal{I}$ est
cofiltrante, on peut choisir un objet $k$ de $\mathcal{I}$ et des morphismes
$a_l : k \to i_l$ et $b_l : k \to j_l$. En prenant alors
$E_k = \bigcup_{l = 1, \ldots, n}
f_{a_l}^{-1}(U_{l, i_l}) \cap f_{b_l}^{-1}(Z_{l, j_l})$,
on obtient une partie constructible de $X_k$ dont l'image réciproque dans $X$
est $E$.
\end{proof}

\begin{lemma}
\label{topology-lemma-directed-inverse-limit-spectral-spaces}
Soit $\mathcal{I}$ une catégorie d'indices cofiltrante.
Soit $i \mapsto X_i$ un diagramme d'espaces spectraux tel
que, pour $a : j \to i$ dans $\mathcal{I}$, l'application correspondante
$f_a : X_j \to X_i$ soit spectrale. Alors la limite projective
$X = \lim X_i$ est un espace topologique spectral, et les applications de
projection $p_i : X \to X_i$ sont spectrales.
\end{lemma}

\begin{proof}
La limite $X = \lim X_i$ existe (Lemme \ref{topology-lemma-limits})
et est quasi-compacte d'après le
Lemme \ref{topology-lemma-inverse-limit-spectral-spaces-quasi-compact}.

\medskip\noindent
Notons $p_i : X \to X_i$ la projection.
Comme $\mathcal{I}$ est cofiltrante, on peut appliquer le
Lemme \ref{topology-lemma-describe-limits}.
Une base de la topologie de $X$ est donc donnée par les ouverts
$p_i^{-1}(U_i)$, où $U_i \subset X_i$ est ouvert. Comme une base de
la topologie de $X_i$ est donnée par les ouverts quasi-compacts, on conclut
qu'une base de la topologie de $X$ est donnée par les $p_i^{-1}(U_i)$,
où $U_i \subset X_i$ est un ouvert quasi-compact. Un argument formel
montre que
$$
p_i^{-1}(U_i) = \lim_{a : j \to i} f_a^{-1}(U_i)
$$
comme espaces topologiques. Puisque chaque $f_a$ est spectrale, les parties
$f_a^{-1}(U_i)$ sont fermées pour la topologie constructible de $X_j$ ;
par conséquent, $p_i^{-1}(U_i)$ est quasi-compact d'après le
Lemme \ref{topology-lemma-inverse-limit-spectral-spaces-quasi-compact}.
Ainsi, $X$ possède une base de sa topologie formée d'ouverts quasi-compacts.

\medskip\noindent
Tout ouvert quasi-compact $U$ de $X$ est de la forme $U = p_i^{-1}(U_i)$
pour un certain $i$ et un certain ouvert quasi-compact $U_i \subset X_i$
(voir le Lemme \ref{topology-lemma-inverse-limit-spectral-spaces-constructible}).
Étant donnés des ouverts quasi-compacts $U_i \subset X_i$ et
$U_j \subset X_j$, on a
$p_i^{-1}(U_i) \cap p_j^{-1}(U_j) = p_k^{-1}(U_k)$ pour un certain $k$
et un certain ouvert quasi-compact $U_k \subset X_k$. En effet, choisissons
$k$ et des morphismes $k \to i$ et $k \to j$, et prenons pour $U_k$
l'intersection des images réciproques de $U_i$ et de $U_j$ dans $X_k$.
On voit ainsi que l'intersection de deux ouverts quasi-compacts de $X$ est
un ouvert quasi-compact.

\medskip\noindent
Enfin, soit $Z \subset X$ une partie fermée irréductible. Alors
$p_i(Z) \subset X_i$ est irréductible, et par conséquent
$Z_i = \overline{p_i(Z)}$ possède un unique point générique $\xi_i$
(puisque $X_i$ est un espace spectral). On a alors
$f_a(\xi_j) = \xi_i$ pour $a : j \to i$ dans $\mathcal{I}$, car
$\overline{f_a(Z_j)} = Z_i$. Ainsi, $\xi = \lim \xi_i$ est un point de $X$.
Nous affirmons que $\xi \in Z$. Sinon, on pourrait trouver un ouvert
quasi-compact contenant $\xi$ et disjoint de $Z$. Il serait de la forme
$p_i^{-1}(U_i)$ pour un certain $i$ et un certain ouvert quasi-compact
$U_i \subset X_i$. On aurait alors $\xi_i \in U_i$, mais
$p_i(Z) \cap U_i = \emptyset$, ce qui contredit
$\xi_i \in \overline{p_i(Z)}$. Ainsi, $\xi \in Z$, et donc
$\overline{\{\xi\}} \subset Z$. Réciproquement, tout $z \in Z$ appartient
à l'adhérence de $\xi$. En effet, étant donné un voisinage ouvert quasi-compact
$U$ de $z$, écrivons $U = p_i^{-1}(U_i)$ pour un certain $i$
et un certain ouvert quasi-compact $U_i \subset X_i$. On voit que
$p_i(z) \in U_i$, donc que $\xi_i \in U_i$, et donc que $\xi \in U$. Ainsi,
$\xi$ est un point générique de $Z$. Nous omettons de démontrer que $\xi$ est
l'unique point générique de $Z$ (indication : montrer qu'un second point
générique est nécessairement égal à $\xi$ en montrant qu'il a pour image
$\xi_i$ dans $X_i$, puisque $X_i$ est spectral et que l'irréductible $Z_i$
possède un unique point générique). Cela achève la démonstration.
\end{proof}

\begin{lemma}
\label{topology-lemma-descend-opens}
Soit $\mathcal{I}$ une catégorie d'indices cofiltrante.
Soit $i \mapsto X_i$ un diagramme d'espaces spectraux tel
que, pour $a : j \to i$ dans $\mathcal{I}$, l'application correspondante
$f_a : X_j \to X_i$ soit spectrale. Posons $X = \lim X_i$ et notons
$p_i : X \to X_i$ la projection.
\begin{enumerate}
\item Pour tout ouvert quasi-compact $U \subset X$,
il existe un $i \in \Ob(\mathcal{I})$ et un ouvert quasi-compact
$U_i \subset X_i$ tels que $p_i^{-1}(U_i) = U$.
\item Étant donnés des ouverts quasi-compacts $U_i \subset X_i$ et
$U_j \subset X_j$ tels que $p_i^{-1}(U_i) \subset p_j^{-1}(U_j)$,
il existe $k \in \Ob(\mathcal{I})$ et des morphismes
$a : k \to i$ et $b : k \to j$ tels que
$f_a^{-1}(U_i) \subset f_b^{-1}(U_j)$.
\item Si $U_i, U_{1, i}, \ldots, U_{n, i} \subset X_i$ sont des ouverts
quasi-compacts et si
$p_i^{-1}(U_i) = p_i^{-1}(U_{1, i}) \cup \ldots \cup p_i^{-1}(U_{n, i})$,
alors
$f_a^{-1}(U_i) = f_a^{-1}(U_{1, i}) \cup \ldots \cup f_a^{-1}(U_{n, i})$
pour un certain morphisme $a : j \to i$ dans $\mathcal{I}$.
\item Même énoncé qu'en (3), mais pour les intersections.
\end{enumerate}
\end{lemma}

\begin{proof}
L'assertion (1) est un cas particulier du
Lemme \ref{topology-lemma-inverse-limit-spectral-spaces-constructible}.
L'assertion (2) est un cas particulier du
Lemme \ref{topology-lemma-inverse-limit-spectral-spaces-equal},
car les ouverts quasi-compacts sont à la fois ouverts et fermés pour la
topologie constructible. Les assertions (3) et (4) résultent formellement de
(1) et (2), ainsi que du fait que la formation des images réciproques
commute aux réunions et aux intersections.
\end{proof}

\begin{lemma}
\label{topology-lemma-make-spectral-space}
Soit $W$ une partie d'un espace spectral $X$. Les conditions suivantes sont
équivalentes :
\begin{enumerate}
\item $W$ est une intersection de parties constructibles et est stable par
générisation,
\item $W$ est quasi-compacte et stable par générisation,
\item il existe une partie quasi-compacte $E \subset X$ telle que
$W$ soit l'ensemble des points qui se spécialisent en un point de $E$,
\item $W$ est une intersection d'ouverts quasi-compacts,
\item
\label{topology-item-intersection-quasi-compact-open}
il existe un ensemble non vide $I$ et des ouverts quasi-compacts
$U_i \subset X$, $i \in I$,
tels que $W = \bigcap U_i$ et que, pour tous $i, j \in I$, il existe un
$k \in I$ tel que $U_k \subset U_i \cap U_j$.
\end{enumerate}
Dans ce cas, (a) $W$ est un espace spectral, (b) $W = \lim U_i$
comme espaces topologiques et (c) pour tout ouvert $U$ contenant $W$,
il existe un $i$ tel que $U_i \subset U$.
\end{lemma}

\begin{proof}
Soit $W \subset X$ vérifiant (1). Alors $W$ est fermée pour la topologie
constructible, donc quasi-compacte pour la topologie constructible (d'après les
Lemmes \ref{topology-lemma-constructible-hausdorff-quasi-compact} et
\ref{topology-lemma-closed-in-quasi-compact}), donc quasi-compacte pour la topologie
de $X$ (car les ouverts de $X$ sont ouverts pour la topologie constructible).
Ainsi, (2) est satisfaite.

\medskip\noindent
Il est clair que (2) implique (3) en prenant $E = W$.

\medskip\noindent
Soit $X$ un espace spectral et soit $E \subset W$ comme en (3).
Puisque tout point de $W$ se spécialise en un point de $E$, on voit qu'un
ouvert de $W$ qui contient $E$ est égal à $W$. Par conséquent, puisque $E$
est quasi-compacte, $W$ l'est également.
Si $x \in X$, $x \not \in W$, alors $Z = \overline{\{x\}}$ est
disjoint de $W$. Comme $W$ est quasi-compacte, on peut trouver un
ouvert quasi-compact $U$ tel que $W \subset U$ et $U \cap Z = \emptyset$.
On conclut que (4) est satisfaite.

\medskip\noindent
Si $W = \bigcap_{j \in J} U_j$, alors, en prenant pour $I$ l'ensemble des
parties finies de $J$ et $U_i = U_{j_1} \cap \ldots \cap U_{j_r}$
pour $i = \{j_1, \ldots, j_r\}$, on voit que (4) implique (5). Il est immédiat
que (5) implique (1).

\medskip\noindent
Soient $I$ et $U_i$ comme en (5).
Puisque $W = \bigcap U_i$, on a $W = \lim U_i$ par la propriété universelle
des limites. Alors $W$ est un espace spectral d'après le
Lemme \ref{topology-lemma-directed-inverse-limit-spectral-spaces}.
Soit $U \subset X$ un voisinage ouvert de $W$.
Alors les $E_i = U_i \cap (X \setminus U)$ forment une famille de parties
constructibles de l'espace spectral $Z = X \setminus U$
d'intersection vide. En utilisant le fait que la topologie spectrale sur $Z$
est quasi-compacte (Lemme \ref{topology-lemma-constructible-hausdorff-quasi-compact}),
on déduit du
Lemme \ref{topology-lemma-intersection-closed-in-quasi-compact}
que $E_i = \emptyset$ pour un certain $i$.
\end{proof}

\begin{lemma}
\label{topology-lemma-make-spectral-space-minus}
Soit $X$ un espace spectral. Soit $E \subset X$ une partie constructible.
Soit $W \subset X$ l'ensemble des points de $X$ qui se spécialisent
en un point de $E$. Alors $W \setminus E$ est un espace spectral.
Si $W = \bigcap U_i$, avec les $U_i$ comme dans le
Lemme \ref{topology-lemma-make-spectral-space}
(\ref{topology-item-intersection-quasi-compact-open}),
alors $W \setminus E = \lim (U_i \setminus E)$.
\end{lemma}

\begin{proof}
Puisque $E$ est constructible, elle est quasi-compacte, et le
Lemme \ref{topology-lemma-make-spectral-space} s'applique donc à $W$.
Si $E$ est constructible, alors $E$ est constructible dans $U_i$
pour tout $i \in I$. Ainsi, $U_i \setminus E$
est spectral d'après le Lemme \ref{topology-lemma-spectral-sub}.
Puisque $W \setminus E = \bigcap (U_i \setminus E)$, on a
$W \setminus E = \lim U_i \setminus E$ par la propriété universelle des limites.
Alors $W \setminus E$ est un espace spectral d'après le
Lemme \ref{topology-lemma-directed-inverse-limit-spectral-spaces}.
\end{proof}






\section{Compactification de Stone-{\v C}ech}
\label{topology-section-stone-cech}

\noindent
La compactification de Stone-{\v C}ech d'un espace topologique $X$
est une application $X \to \beta(X)$ de $X$ vers un espace séparé
quasi-compact $\beta(X)$, universelle parmi de telles applications. Nous
démontrons son existence par un argument standard utilisant le lemme élémentaire
suivant.

\begin{lemma}
\label{topology-lemma-dense-image}
Soit $f : X \to Y$ une application continue entre espaces topologiques. Supposons
que $f(X)$ soit dense dans $Y$ et que $Y$ soit séparé. Alors le cardinal de
$Y$ est inférieur ou égal à celui de $P(P(X))$, où $P$ désigne l'opération
qui associe à un ensemble l'ensemble de ses parties.
\end{lemma}

\begin{proof}
Soit $S = f(X) \subset Y$. Soit $\mathcal{D}$ l'ensemble de tous les
fermés réguliers de $Y$, c'est-à-dire des parties $D \subset Y$ qui
sont égales à l'adhérence de leur intérieur. Notons que l'adhérence d'une
partie ouverte de $Y$ est un fermé régulier. Pour $y \in Y$, considérons
l'ensemble
$$
I_y = \{T \subset S \mid \text{ il existe }
D \in \mathcal{D}\text{ tel que }T = S \cap D\text{ et }y \in D\}.
$$
Comme $S$ est dense dans $Y$, pour tout fermé régulier $D$, on voit
que $S \cap D$ est dense dans $D$. Par conséquent, si
$D \cap S = D' \cap S$ pour $D, D' \in \mathcal{D}$, alors
$D = D'$. Ainsi, $I_y = I_{y'}$ implique $y = y'$, car la condition de
séparation assure que l'on peut trouver un fermé régulier contenant $y$
mais pas $y'$. Le résultat s'ensuit.
\end{proof}

\noindent
Soit $X$ un espace topologique. D'après le Lemme \ref{topology-lemma-dense-image}, il
existe un ensemble $I$
de classes d'isomorphisme d'applications continues $f : X \to Y$ à image
dense, où $Y$ est séparé et quasi-compact. Pour $i \in I$, choisissons
un représentant $f_i : X \to Y_i$. Considérons l'application
$$
\prod f_i : X \longrightarrow \prod\nolimits_{i \in I} Y_i
$$
et notons $\beta(X)$ l'adhérence de son image. Puisque chaque $Y_i$ est
séparé, $\beta(X)$ l'est aussi. Puisque chaque $Y_i$ est quasi-compact,
$\beta(X)$ l'est aussi (utiliser le Théorème \ref{topology-theorem-tychonov} et le
Lemme \ref{topology-lemma-closed-in-quasi-compact}).

\medskip\noindent
Montrons que l'application canonique $X \to \beta(X)$ satisfait à la propriété
universelle relativement aux applications vers des espaces séparés
quasi-compacts. Soit donc $f : X \to Y$ un tel morphisme. Soit
$Z \subset Y$ l'adhérence de $f(X)$. Alors $X \to Z$ est isomorphe à l'une des
applications $f_i : X \to Y_i$, disons $f_{i_0} : X \to Y_{i_0}$. Ainsi, $f$
se factorise sous la forme
$X \to \beta(X) \to \prod Y_i \to Y_{i_0} \cong Z \to Y$, comme voulu.

\begin{lemma}
\label{topology-lemma-one-point-compactification}
Soit $X$ un espace séparé et localement quasi-compact.
Il existe une application $X \to X^*$ qui identifie $X$ à un sous-espace
ouvert d'un espace séparé quasi-compact $X^*$ tel que
$X^* \setminus X$ soit un singleton (compactification par adjonction d'un point).
En particulier, l'application $X \to \beta(X)$ identifie $X$
à un sous-espace ouvert de $\beta(X)$.
\end{lemma}

\begin{proof}
Posons $X^* = X \amalg \{\infty\}$. Déclarons qu'une partie $V$ de $X^*$ est
ouverte si $V \subset X$ est ouverte dans $X$, ou bien si $\infty \in V$ et si
$U = V \cap X$ est un ouvert de $X$ tel que $X \setminus U$ soit quasi-compact.
Nous omettons de vérifier que cela définit une topologie. Il est clair
que $X \to X^*$ identifie $X$ à un sous-espace ouvert de $X$.

\medskip\noindent
Puisque $X$ est localement quasi-compact, tout point $x \in X$ possède un
voisinage quasi-compact $x \in E \subset X$. Alors $E$
est fermé (Lemme \ref{topology-lemma-quasi-compact-in-Hausdorff}, partie (1)), et
$V = (X \setminus E) \amalg \{\infty\}$ est un voisinage ouvert
de $\infty$ disjoint de l'intérieur de $E$. Ainsi, $X^*$ est séparé.

\medskip\noindent
Soit $X^* = \bigcup V_i$ un recouvrement ouvert. Alors, pour un certain $i$,
disons $i_0$, on a $\infty \in V_{i_0}$. Par construction,
$Z = X^* \setminus V_{i_0}$ est quasi-compact. Le recouvrement
$Z \subset \bigcup_{i \not = i_0} Z \cap V_i$ admet donc un raffinement fini,
ce qui implique que le recouvrement donné de $X^*$ admet un raffinement fini.
Ainsi, $X^*$ est quasi-compact.

\medskip\noindent
L'application $X \to X^*$ se factorise sous la forme
$X \to \beta(X) \to X^*$ par la propriété universelle de la compactification
de Stone-{\v C}ech. Soit
$\varphi : \beta(X) \to X^*$ cette factorisation.
Alors $X \to \varphi^{-1}(X)$ est une section de
$\varphi^{-1}(X) \to X$, et possède donc une image fermée
(Lemme \ref{topology-lemma-section-closed}).
Comme l'image de $X \to \beta(X)$ est dense, on conclut que
$X = \varphi^{-1}(X)$.
\end{proof}








\section{Espaces extrêmement discontinus}
\label{topology-section-extremally-disconnected}

\noindent
Le contenu de cette section est tiré de \cite{Gleason}
(avec une légère modification comme dans \cite{Rainwater}).
Dans l'article de Gleason, il est démontré que, dans la catégorie des espaces
séparés quasi-compacts, les « objets projectifs » sont exactement les espaces
extrêmement discontinus.

\begin{definition}
\label{topology-definition-extremally-disconnected}
On dit qu'un espace topologique $X$ est {\it extrêmement discontinu}
si l'adhérence de toute partie ouverte de $X$ est ouverte.
\end{definition}

\noindent
Si $X$ est séparé et extrêmement discontinu, alors $X$ est totalement
discontinu (cela est faux en général). Si $X$ est quasi-compact, séparé et
extrêmement discontinu, alors $X$ est profini d'après le
lemme \ref{topology-lemma-profinite}, mais la réciproque est fausse en général.
Par exemple, l'anneau des entiers $p$-adiques
$\mathbf{Z}_p = \lim \mathbf{Z}/p^n\mathbf{Z}$ est un espace profini
qui n'est pas extrêmement discontinu. En effet, si $U \subset \mathbf{Z}_p$
est l'ensemble des éléments non nuls dont la valuation est paire, alors $U$ est
ouvert, mais son adhérence est $U \cup \{0\}$, qui n'est pas ouvert.

\begin{lemma}
\label{topology-lemma-image-open-technical}
Soit $f : X \to Y$ une application continue entre espaces topologiques.
Supposons que $f$ soit surjective et que $f(E) \not = Y$ pour toute partie
fermée propre $E \subset X$. Alors, pour tout ouvert $U \subset X$, la partie
$f(U)$ est contenue dans l'adhérence de $Y \setminus f(X \setminus U)$.
\end{lemma}

\begin{proof}
Choisissons $y \in f(U)$ et un voisinage ouvert quelconque $V \subset Y$ de
$y$. Nous allons montrer que $V$ rencontre $Y \setminus f(X \setminus U)$.
Remarquons que $W = U \cap f^{-1}(V)$ est un ouvert non vide de $X$ ; ainsi,
$f(X \setminus W) \not = Y$. Prenons $y' \in Y$ tel que
$y' \not \in f(X \setminus W)$. Il est élémentaire de vérifier que $y' \in V$
et que $y' \in Y \setminus f(X \setminus U)$.
\end{proof}

\begin{lemma}
\label{topology-lemma-intersection-empty}
Soit $X$ un espace extrêmement discontinu.
Si $U, V \subset X$ sont des ouverts disjoints, alors
$\overline{U}$ et $\overline{V}$ sont eux aussi disjoints.
\end{lemma}

\begin{proof}
Par hypothèse, $\overline{U}$ est ouvert ; ainsi,
$V \cap \overline{U}$ est ouvert et disjoint de $U$, donc vide, puisque
$\overline{U}$ est l'intersection de toutes les parties fermées de $X$ qui
contiennent $U$. Cela signifie que l'ouvert
$\overline{V} \cap \overline{U}$ ne rencontre pas $V$ ; il est donc vide par
le même argument.
\end{proof}

\begin{lemma}
\label{topology-lemma-isomorphism}
Soit $f : X \to Y$ une application continue entre espaces topologiques
séparés quasi-compacts. Si $Y$ est extrêmement discontinu, si $f$ est
surjective et si $f(Z) \not = Y$ pour toute partie fermée propre $Z$ de $X$,
alors $f$ est un homéomorphisme.
\end{lemma}

\begin{proof}
D'après le lemme \ref{topology-lemma-bijective-map}, il suffit de montrer que $f$ est
injective. Supposons que $x, x' \in X$ soient des points distincts tels que
$y = f(x) = f(x')$. Choisissons des voisinages ouverts disjoints
$U, U' \subset X$ respectifs de $x, x'$. Remarquons que $f$ est fermée
(lemme \ref{topology-lemma-closed-map}) ; ainsi, $T = f(X \setminus U)$ et
$T' = f(X \setminus U')$ sont fermés dans $Y$. Puisque $X$ est la réunion de
$X \setminus U$ et de $X \setminus U'$, on voit que $Y = T \cup T'$.
D'après le lemme \ref{topology-lemma-image-open-technical}, on voit que $y$ appartient
à l'adhérence de $Y \setminus T$ et à celle de $Y \setminus T'$. D'autre part,
d'après le lemme \ref{topology-lemma-intersection-empty}, leur intersection est vide.
Nous obtenons ainsi la contradiction souhaitée.
\end{proof}

\begin{lemma}
\label{topology-lemma-find-compact-subset}
Soit $f : X \to Y$ une application continue surjective entre espaces
topologiques séparés quasi-compacts. Il existe une partie quasi-compacte
$E \subset X$ telle que $f(E) = Y$, mais $f(E') \not = Y$ pour toute partie
fermée propre $E' \subset E$.
\end{lemma}

\begin{proof}
Nous utiliserons sans plus le mentionner le fait que les parties quasi-compactes
de $X$ sont exactement les parties fermées
(lemme \ref{topology-lemma-closed-in-compact}).
Considérons la collection $\mathcal{E}$ de toutes les parties quasi-compactes
$E \subset X$ telles que $f(E) = Y$, ordonnée par inclusion. Nous allons
utiliser le lemme de Zorn pour montrer que $\mathcal{E}$ possède un élément
minimal. Pour cela, il suffit de montrer que, pour toute famille totalement
ordonnée $E_\lambda$ d'éléments de $\mathcal{E}$, l'intersection
$\bigcap E_\lambda$ appartient à $\mathcal{E}$. Elle est quasi-compacte car
elle est fermée. Pour tout $y \in Y$, les parties
$E_\lambda \cap f^{-1}(\{y\})$ sont non vides et fermées ; ainsi, l'intersection
$\bigcap E_\lambda \cap f^{-1}(\{y\}) = \bigcap (E_\lambda \cap f^{-1}(\{y\}))$
est non vide d'après le
lemme \ref{topology-lemma-intersection-closed-in-quasi-compact}.
Cela achève la démonstration.
\end{proof}

\begin{proposition}
\label{topology-proposition-projective-in-category-hausdorff-qc}
Soit $X$ un espace topologique séparé et quasi-compact.
Les assertions suivantes sont équivalentes :
\begin{enumerate}
\item $X$ est extrêmement discontinu,
\item pour toute application continue surjective $f : Y \to X$, où $Y$ est
séparé et quasi-compact, il existe une section continue, et
\item pour tout diagramme commutatif en traits pleins
$$
\xymatrix{
& Y \ar[d] \\
X \ar@{..>}[ru] \ar[r] & Z
}
$$
formé d'applications continues entre espaces séparés quasi-compacts, dans lequel
$Y \to Z$ est surjective, il existe dans la catégorie des espaces topologiques
une flèche pointillée qui rend le diagramme commutatif.
\end{enumerate}
\end{proposition}

\begin{proof}
Il est clair que (3) implique (2). D'autre part, si (2) est vérifiée et si
$X \to Z$ et $Y \to Z$ sont comme en (3), alors (2) assure l'existence d'une
section de la projection $X \times_Z Y \to X$, d'où l'existence d'une flèche
pointillée convenable (détails omis). Ainsi, (3) est équivalente à (2).

\medskip\noindent
Supposons $X$ extrêmement discontinu et soit $f : Y \to X$ comme en (2).
D'après le lemme \ref{topology-lemma-find-compact-subset}, il existe une partie
quasi-compacte $E \subset Y$ telle que $f(E) = X$, mais
$f(E') \not = X$ pour toute partie fermée propre $E' \subset E$. D'après le
lemme \ref{topology-lemma-isomorphism}, $f|_E : E \to X$ est un homéomorphisme, dont
l'inverse donne la section souhaitée.

\medskip\noindent
Supposons (2). Soit $U \subset X$ un ouvert de complémentaire $Z$.
Considérons la surjection continue $f : \overline{U} \amalg Z \to X$.
Soit $\sigma$ une section. Alors $\overline{U} = \sigma^{-1}(\overline{U})$
est ouverte. Ainsi, $X$ est extrêmement discontinu.
\end{proof}

\begin{lemma}
\label{topology-lemma-rainwater}
Soit $f : X \to X$ une application continue surjective d'un espace
topologique séparé dans lui-même. Si $f$ n'est pas $\text{id}_X$, alors il
existe une partie fermée propre $E \subset X$ telle que $X = E \cup f(E)$.
\end{lemma}

\begin{proof}
Choisissons $p \in X$ tel que $f(p) \not = p$. Choisissons des voisinages
ouverts disjoints $p \in U$, $f(p) \in V$ et posons
$E = X \setminus U \cap f^{-1}(V)$. Alors $p \not \in E$ ; ainsi, $E$ est une
partie fermée propre. Si $x \in X$, alors soit $x \in E$, soit, dans le cas
contraire, $x \in U \cap f^{-1}(V)$. Écrivons $x = f(y)$, ce qui est possible
puisque $f$ est surjective. Si $y \in U \cap f^{-1}(V)$, alors on aurait
$x = f(y) \in V$, ce qui contredit $x \in U$. Ainsi, $y \in E$ et
$x \in f(E)$.
\end{proof}

\begin{example}
\label{topology-example-stone-Cech-discrete}
La proposition \ref{topology-proposition-projective-in-category-hausdorff-qc} permet de
voir que la compactification de Stone-{\v C}ech $\beta(X)$ d'un espace discret
$X$ est extrêmement discontinue. En effet, soit $f : Y \to \beta(X)$ une
surjection continue, où $Y$ est quasi-compact et séparé. Nous pouvons alors
relever l'application $X \to \beta(X)$ en une application continue (!)
$X \to Y$, puisque $X$ est discret. Par la propriété universelle de la
compactification de Stone-{\v C}ech, nous obtenons une factorisation
$X \to \beta(X) \to Y$. Puisque $\beta(X) \to Y \to \beta(X)$ est égale à
l'identité sur la partie dense $X$, nous concluons que nous avons obtenu une
section.
En particulier, nous concluons que la compactification de Stone-{\v C}ech d'un
espace discret est totalement discontinue, donc profinie
(voir la discussion qui suit la
définition \ref{topology-definition-extremally-disconnected} et le
lemme \ref{topology-lemma-profinite}).
\end{example}

\noindent
Grâce aux espaces extrêmement discontinus fournis par
l'exemple \ref{topology-example-stone-Cech-discrete}, nous pouvons démontrer que tout
espace séparé quasi-compact possède une « enveloppe projective » dans la
catégorie des espaces séparés quasi-compacts.

\begin{lemma}
\label{topology-lemma-existence-projective-cover}
\begin{slogan}
Tout espace séparé quasi-compact possède une enveloppe
extrêmement discontinue canonique
\end{slogan}
Soit $X$ un espace séparé quasi-compact.
Il existe une surjection continue $X' \to X$ telle que $X'$ soit
quasi-compact, séparé et extrêmement discontinu.
Si l'on exige qu'aucune partie fermée propre de $X'$ n'ait pour image $X$,
alors $X'$ est unique à isomorphisme près.
\end{lemma}

\begin{proof}
Soit $Y = X$, mais muni de la topologie discrète. Soit $X' = \beta(Y)$.
L'application continue $Y \to X$ se factorise en $Y \to X' \to X$. Cela
démontre la première assertion du lemme d'après
l'exemple \ref{topology-example-stone-Cech-discrete}.

\medskip\noindent
D'après le lemme \ref{topology-lemma-find-compact-subset}, nous pouvons trouver une
partie quasi-compacte $E \subset X'$ dont l'image est $X$ et telle qu'aucune
partie fermée propre de $E$ n'ait pour image $X$.
Puisque $X'$ est extrêmement discontinu, il existe une application continue
$f : X' \to E$ au-dessus de $X$
(proposition \ref{topology-proposition-projective-in-category-hausdorff-qc}).
En composant $f$ avec l'application $E \to X'$, on obtient l'endomorphisme
continu $f|_E : E \to E$. Remarquons que $f|_E$ doit être surjectif, car sinon
son image serait une partie fermée propre d'image $X$.
Ainsi, $f|_E$ doit être $\text{id}_E$, car sinon le
lemme \ref{topology-lemma-rainwater} montrerait que $E$ n'est pas minimal.
Par conséquent, $\text{id}_E$ se factorise par l'espace extrêmement discontinu
$X'$. Un argument catégorique formel (utilisant la caractérisation de la
proposition \ref{topology-proposition-projective-in-category-hausdorff-qc}) montre que
$E$ est extrêmement discontinu.

\medskip\noindent
Pour démontrer l'unicité, supposons donnée une seconde enveloppe minimale
$X'' \to X$. Par la propriété de relèvement démontrée dans la
proposition \ref{topology-proposition-projective-in-category-hausdorff-qc}, nous pouvons
trouver une application continue $g : X' \to X''$ au-dessus de $X$.
Remarquons que $g$ est une application fermée
(lemme \ref{topology-lemma-closed-map}). Ainsi, $g(X') \subset X''$ est une partie
fermée d'image $X$, et nous concluons que
$g(X') = X''$ par minimalité de $X''$. D'autre part, si $E \subset X'$ est une
partie fermée
propre, alors $g(E) \not = X''$, car $E$ n'a pas pour image $X$ par minimalité
de $X'$. D'après le lemme \ref{topology-lemma-isomorphism}, $g$ est un
isomorphisme.
\end{proof}

\begin{remark}
\label{topology-remark-size-projective-cover}
Soit $X$ un espace séparé quasi-compact. Soit $\kappa$ un cardinal infini
supérieur ou égal au cardinal de $X$. Alors le cardinal de l'enveloppe minimale
séparée, quasi-compacte et extrêmement discontinue
$X' \to X$ (lemme \ref{topology-lemma-existence-projective-cover})
est au plus $2^{2^\kappa}$. En effet, choisissons une partie $S \subset X'$
qui s'envoie bijectivement sur $X$. Par minimalité de $X'$, la partie
$S$ est dense dans $X'$. Ainsi, $|X'| \leq 2^{2^\kappa}$ d'après le
lemme \ref{topology-lemma-dense-image}.
\end{remark}






\section{Divers}
\label{topology-section-miscellany}


\noindent
Le lemme suivant s'applique à l'espace topologique sous-jacent à un schéma
quasi-séparé.

\begin{lemma}
\label{topology-lemma-topology-quasi-separated-scheme}
Soit $X$ un espace topologique qui
\begin{enumerate}
\item possède une base de la topologie formée d'ouverts quasi-compacts, et
\item est tel que l'intersection de deux ouverts quasi-compacts
quelconques est quasi-compacte.
\end{enumerate}
Alors
\begin{enumerate}
\item $X$ est localement quasi-compact,
\item tout ouvert quasi-compact $U \subset X$ est rétrocompact,
\item tout ouvert quasi-compact $U \subset X$ possède un système cofinal de
recouvrements ouverts $\mathcal{U} : U = \bigcup_{j\in J} U_j$ où $J$ est fini
et où tous les $U_j$ et $U_j \cap U_{j'}$ sont quasi-compacts,
\item en ajouter ici.
\end{enumerate}
\end{lemma}

\begin{proof}
La démonstration est omise.
\end{proof}

\begin{definition}
\label{topology-definition-isolated-point}
Soit $X$ un espace topologique. On dit que $x \in X$ est un
{\it point isolé} de $X$ si $\{x\}$ est ouvert dans $X$.
\end{definition}

\section{Partitions et stratifications}
\label{topology-section-stratifications}

\noindent
Les stratifications peuvent être définies de bien des manières. Toute remarque
sur le choix des définitions adoptées dans cette section est la bienvenue.

\begin{definition}
\label{topology-definition-paritition}
Soit $X$ un espace topologique. Une {\it partition} de $X$ est une
décomposition $X = \coprod X_i$ en parties localement fermées $X_i$.
Les $X_i$ sont appelées les {\it parties} de la partition.
Étant donné deux partitions de $X$, nous disons que l'une {\it raffine}
l'autre si les parties de l'une sont des réunions de parties de l'autre.
\end{definition}

\noindent
Tout espace topologique $X$ admet une partition en composantes connexes.
Si $X$ possède un nombre fini de composantes irréductibles
$Z_1, \ldots, Z_r$, il existe une partition qui raffine la partition en
composantes connexes et dont les parties
$X_I = \bigcap_{i \in I} Z_i \setminus (\bigcup_{i \not \in I} Z_i)$
sont indexées par les sous-ensembles $I \subset \{1, \ldots, r\}$.

\begin{definition}
\label{topology-definition-good-stratification}
Soit $X$ un espace topologique. Une {\it bonne stratification}
de $X$ est une partition $X = \coprod X_i$ telle que, pour tous
$i, j \in I$, on ait
$$
X_i \cap \overline{X_j} \not = \emptyset
\Rightarrow
X_i \subset \overline{X_j}.
$$
\end{definition}

\noindent
Une bonne stratification $X = \coprod_{i \in I} X_i$ définit
un ordre partiel sur $I$ en posant $i \leq j$ si et seulement si
$X_i \subset \overline{X_j}$. On voit alors que
$$
\overline{X_j} = \bigcup\nolimits_{i \leq j} X_i
$$
Cependant, en géométrie algébrique, il arrive souvent que l'on sache seulement
que le membre de gauche est contenu dans le membre de droite dans la dernière
formule affichée. Cela conduit à la définition suivante.

\begin{definition}
\label{topology-definition-stratification}
Soit $X$ un espace topologique. Une {\it stratification} de $X$ est
la donnée d'une partition $X = \coprod_{i \in I} X_i$ et d'un ordre partiel
sur $I$ tels que, pour tout $j \in I$, on ait
$$
\overline{X_j} \subset \bigcup\nolimits_{i \leq j} X_i
$$
Les parties $X_i$ sont appelées les {\it strates} de la stratification.
\end{definition}

\noindent
Nous imposons souvent des conditions supplémentaires à la stratification.
Par exemple, les stratifications sont particulièrement agréables lorsqu'elles
sont {\it localement finies}, ce qui signifie que tout point possède un
voisinage ne rencontrant qu'un nombre fini de strates. Plus généralement,
nous introduisons la définition suivante.

\begin{definition}
\label{topology-definition-locally-finite}
Soit $X$ un espace topologique. Soit $I$ un ensemble et, pour $i \in I$,
soit $E_i \subset X$ une partie. Nous disons que la famille
$\{E_i\}_{i \in I}$ est {\it localement finie} si, pour tout $x \in X$,
il existe un voisinage ouvert $U$ de $x$ tel que
$\{i \in I | E_i \cap U \not = \emptyset\}$ soit fini.
\end{definition}


\begin{remark}
\label{topology-remark-locally-finite-stratification}
Étant donné une stratification localement finie $X = \coprod X_i$ d'un
espace topologique $X$, on obtient une famille de parties fermées
$Z_i = \bigcup_{j \leq i} X_j$ de $X$, indexée par $I$, telle que
$$
Z_i \cap Z_j = \bigcup\nolimits_{k \leq i, j} Z_k
$$
Réciproquement, supposons données des parties fermées $Z_i \subset X$
indexées par un ensemble ordonné $I$, telles que $X = \bigcup Z_i$, que tout
point possède un voisinage ne rencontrant qu'un nombre fini de $Z_i$ et que la
formule affichée soit vérifiée. On obtient alors une stratification localement
finie de $X$ en posant $X_i = Z_i \setminus \bigcup_{j < i} Z_j$.
\end{remark}

\begin{lemma}
\label{topology-lemma-partition-refined-by-stratification}
Soit $X$ un espace topologique et soit $X = \coprod X_i$ une partition finie
de $X$. Il existe alors une stratification finie de $X$ qui la raffine.
\end{lemma}

\begin{proof}
Posons $T_i = \overline{X_i}$ et $\Delta_i = T_i \setminus X_i$.
Soit $S$ l'ensemble de toutes les intersections de $T_i$ et de $\Delta_i$.
(Par exemple, $T_1 \cap T_2 \cap \Delta_4$ est un élément de $S$.)
Alors $S = \{Z_s\}$ est une famille finie de parties fermées de $X$ telle que
$Z_s \cap Z_{s'} \in S$ pour tous $s, s' \in S$. Munissons $S$ de l'ordre
partiel donné par l'inclusion. En posant ensuite
$Y_s = Z_s \setminus \bigcup_{s' < s} Z_{s'}$, on obtient la stratification
cherchée.
\end{proof}

\begin{lemma}
\label{topology-lemma-constructible-partition-refined-by-stratification}
Soit $X$ un espace topologique. Supposons que
$X = T_1 \cup \ldots \cup T_n$ soit écrit comme réunion de parties
constructibles. Il existe une stratification finie
$X = \coprod X_i$, dont chaque $X_i$ est constructible, telle que chaque
$T_k$ soit une réunion de strates.
\end{lemma}

\begin{proof}
Par définition des parties constructibles, chaque $T_i$ s'écrit comme réunion
finie de parties $U \cap V^c$, où $U, V \subset X$ sont des ouverts
rétrocompacts. Nous pouvons donc supposer que $T_i = U_i \cap V_i^c$,
où $U_i, V_i \subset X$ sont des ouverts rétrocompacts. Soit $S$ l'ensemble
fini de parties fermées de $X$ constitué de
$\emptyset, X, U_i^c, V_i^c$ et de leurs intersections finies.
Si $Z \in S$, alors $Z$ est constructible dans $X$
(Lemme \ref{topology-lemma-constructible}).
De plus, $Z \cap Z' \in S$ pour tous $Z, Z' \in S$.
Munissons $S$ de l'ordre partiel donné par l'inclusion. Pour $Z \in S$, posons
$X_Z = Z \setminus \bigcup_{Z' < Z,\ Z' \in S} Z'$.
On obtient ainsi une stratification $X = \coprod_{Z \in S} X_Z$
qui possède les propriétés énoncées dans le lemme.
\end{proof}

\begin{lemma}
\label{topology-lemma-noetherian-partition-refined-by-stratification}
Soit $X$ un espace topologique noethérien. Toute partition finie
de $X$ peut être raffinée par une bonne stratification finie.
\end{lemma}

\begin{proof}
Soit $X = \coprod X_i$ une partition finie de $X$.
Soit $Z$ une composante irréductible de $X$. Puisque
$X = \bigcup \overline{X_i}$ et que l'ensemble d'indices est fini, il existe
un $i$ tel que $Z \subset \overline{X_i}$. Comme $X_i$ est localement fermé,
il s'ensuit que $Z \cap X_i$ contient un ouvert de $Z$. Ainsi,
$Z \cap X_i$ contient un ouvert $U$ de $X$ (Lemme \ref{topology-lemma-Noetherian}).
Écrivons $X_i = U \amalg X_i^1 \amalg X_i^2$, où
$X_i^1 = (X_i \setminus U) \cap \overline{U}$ et
$X_i^2 = (X_i \setminus U) \cap \overline{U}^c$.
Pour $i' \not = i$, posons
$X_{i'}^1 = X_{i'} \cap \overline{U}$ et
$X_{i'}^2 = X_{i'} \cap \overline{U}^c$.
Alors
$$
X \setminus U = \coprod X^k_l
$$
est une partition telle que $\overline{U} \setminus U = \bigcup X_l^1$.
Remarquons que $X \setminus U$ est une partie fermée strictement contenue
dans $X$.
Par récurrence noethérienne, nous pouvons raffiner cette partition
par une bonne stratification finie
$X \setminus U = \coprod_{\alpha \in A} T_\alpha$.
Alors $X = U \amalg \coprod_{\alpha \in A} T_\alpha$ est une bonne
stratification finie de $X$ qui raffine la partition initiale.
\end{proof}







\section{Colimites d'espaces}
\label{topology-section-colimits}

\noindent
La catégorie des espaces topologiques possède des coproduits. En effet, si
$I$ est un ensemble et si, pour $i \in I$, on se donne un espace topologique
$X_i$, on munit l'ensemble $\coprod_{i \in I} X_i$ de la {\it topologie
somme}. Une base de cette topologie est formée des parties $U_i$, où
$U_i \subset X_i$ est ouvert.

\medskip\noindent
La catégorie des espaces topologiques possède des coégalisateurs. En effet,
si $a, b : X \to Y$ sont des morphismes d'espaces topologiques, alors le
coégalisateur de $a$ et $b$ est le coégalisateur $Y/\sim$ dans la catégorie
des ensembles, muni de la topologie quotient (section \ref{topology-section-submersive}).

\begin{lemma}
\label{topology-lemma-colimits}
La catégorie des espaces topologiques possède des colimites, et le foncteur
d'oubli vers la catégorie des ensembles commute aux colimites.
\end{lemma}

\begin{proof}
Cela résulte de la discussion précédente et de Catégories, Lemme
\ref{categories-lemma-colimits-coproducts-coequalizers}. Une autre preuve de
l'existence des colimites est esquissée dans Catégories, Remarque
\ref{categories-remark-how-to-use-it}. Il résulte de ce qui précède que le
foncteur d'oubli commute aux colimites. On peut également le voir en utilisant
Catégories, Lemme \ref{categories-lemma-adjoint-exact}, et le fait que le
foncteur d'oubli possède un adjoint à droite, à savoir le foncteur qui associe
à un ensemble l'espace topologique chaotique (ou indiscret) correspondant.
\end{proof}





\section{Groupes, anneaux et modules topologiques}
\label{topology-section-topological-groups}

\noindent
Cette courte section ne contient que des définitions et des propriétés
élémentaires.

\begin{definition}
\label{topology-definition-topological-group}
Un {\it groupe topologique} est un groupe $G$ muni d'une topologie telle que
la multiplication $G \times G \to G$, $(x, y) \mapsto xy$, et l'inversion
$G \to G$, $x \mapsto x^{-1}$, soient continues.
Un {\it homomorphisme de groupes topologiques} est un homomorphisme de groupes
qui est continu.
\end{definition}

\noindent
Si $G$ est un groupe topologique et $H \subset G$ un sous-groupe, alors $H$,
muni de la topologie induite, est un groupe topologique. Si $G$ est un groupe
topologique et $G \to H$ une surjection de groupes, alors $H$, muni de la
topologie quotient, est un groupe topologique.

\begin{example}
\label{topology-example-automorphisms-of-a-set}
Soit $E$ un ensemble. On peut munir l'ensemble de ses applications dans
lui-même, noté $\text{Map}(E, E)$, de la topologie compacte-ouverte,
c'est-à-dire de la topologie pour laquelle, étant donné $f : E \to E$, un
système fondamental de voisinages de $f$ est formé des ensembles
$U_S(f) = \{f' : E \to E \mid f'|_S = f|_S\}$, où $S \subset E$ est fini.
Pour cette topologie, l'action
$$
\text{Map}(E, E) \times E \longrightarrow E,\quad
(f, e) \longmapsto f(e)
$$
est continue lorsque $E$ est muni de la topologie discrète. Si $X$ est un
espace topologique et si $X \times E \to E$ est une application continue,
alors l'application $X \to \text{Map}(E, E)$ est continue. Autrement dit, la
topologie compacte-ouverte est la topologie la moins fine pour laquelle
l'application d'« action » affichée ci-dessus est continue. La composition
$$
\text{Map}(E, E) \times \text{Map}(E, E) \to \text{Map}(E, E)
$$
est elle aussi continue (comme on le vérifie aisément à l'aide de la
description précédente des voisinages). Enfin, si
$\text{Aut}(E) \subset \text{Map}(E, E)$ est la partie des applications
bijectives, alors l'inversion $i : \text{Aut}(E) \to \text{Aut}(E)$,
$f \mapsto f^{-1}$, est également continue. En effet, si $S \subset E$ est
fini, alors $i^{-1}(U_S(f^{-1})) = U_{f^{-1}(S)}(f)$. Ainsi,
$\text{Aut}(E)$ est un groupe topologique au sens de la
Définition \ref{topology-definition-topological-group}.
\end{example}

\begin{lemma}
\label{topology-lemma-topological-group-limits}
La catégorie des groupes topologiques possède des limites, et les foncteurs
d'oubli vers (a) la catégorie des espaces topologiques et (b) la catégorie des
groupes commutent aux limites.
\end{lemma}

\begin{proof}
Il suffit d'établir ces assertions pour les produits et les égalisateurs ;
voir Catégories, Lemme
\ref{categories-lemma-limits-products-equalizers}. Soient $G_i$, $i \in I$,
une famille de groupes topologiques. Munissons le produit usuel
$G = \prod G_i$ de la topologie produit. Puisque
$G \times G = \prod (G_i \times G_i)$ comme espace topologique (les produits
commutent entre eux dans toute catégorie), la multiplication sur $G$ est
continue. Il en va de même de l'inversion. Soient
$a, b : G \to H$ deux homomorphismes de groupes topologiques. Pour
l'égalisateur, il suffit alors de prendre l'égalisateur de $a$ et $b$ dans la
catégorie des espaces topologiques : il s'agit aussi de leur égalisateur dans
la catégorie des groupes, muni de la topologie induite.
\end{proof}

\begin{lemma}
\label{topology-lemma-profinite-group}
Soit $G$ un groupe topologique. Les conditions suivantes sont équivalentes :
\begin{enumerate}
\item $G$, considéré comme espace topologique, est profini,
\item $G$ est la limite d'un diagramme de groupes topologiques discrets finis,
\item $G$ est une limite cofiltrante de groupes topologiques discrets finis.
\end{enumerate}
\end{lemma}

\begin{proof}
Nous disposons du résultat correspondant pour les espaces topologiques ; voir
le Lemme \ref{topology-lemma-profinite}. Combiné au
Lemme \ref{topology-lemma-topological-group-limits}, ce résultat montre qu'il suffit de
prouver que (1) implique (3).

\medskip\noindent
Montrons d'abord que tout voisinage $E$ de l'élément neutre $e$ contient un
sous-groupe ouvert. En effet, puisque $G$ est une limite cofiltrante d'espaces
topologiques discrets finis (Lemme \ref{topology-lemma-profinite}), on peut choisir une
application continue $f : G \to T$ vers un espace discret fini $T$ telle que
$f^{-1}(f(\{e\})) \subset E$. Considérons
$$
H = \{g \in G \mid f(gg') = f(g')\text{ pour tout }g' \in G\}
$$
Il s'agit d'un sous-groupe de $G$ contenu dans $E$. Il suffit donc de montrer
que $H$ est ouvert. Fixons $t \in T$ et posons $W = f^{-1}(\{t\})$.
La partie $W \subset G$ est ouverte et fermée, donc en particulier
quasi-compacte. Pour tout $w \in W$, il existe des voisinages ouverts
$e \in U_w \subset G$ et $w \in U'_w \subset W$ tels que
$U_wU'_w \subset W$. Par quasi-compacité, on peut trouver
$w_1, \ldots, w_n$ tels que $W = \bigcup U'_{w_i}$. Alors
$U_t = U_{w_1} \cap \ldots \cap U_{w_n}$ est un voisinage ouvert de $e$ tel
que $f(gw) = t$ pour tout $w \in W$. Comme $T$ est fini,
$\bigcap_{t \in T} U_t \subset H$ est un voisinage ouvert de $e$. Puisque
$H \subset G$ est un sous-groupe, il s'ensuit que $H$ est ouvert.

\medskip\noindent
Supposons que $H \subset G$ soit un sous-groupe ouvert. Puisque $G$ est
quasi-compact, l'indice de $H$ dans $G$ est fini. Écrivons
$G = Hg_1 \cup \ldots \cup Hg_n$. Alors
$N = \bigcap_{i = 1, \ldots, n} g_iHg_i^{-1}$ est un sous-groupe normal
ouvert contenu dans $H$. Comme $N$ est lui aussi
d'indice fini, $G \to G/N$ est une surjection vers un groupe topologique
discret fini.

\medskip\noindent
Considérons l'application
$$
G \longrightarrow \lim_{N \subset G\text{ ouvert et normal}} G/N
$$
Nous affirmons que cette application est un isomorphisme de groupes
topologiques. Cela achève la preuve du lemme, car le système qui figure à
droite est cofiltrant (l'intersection de deux sous-groupes normaux ouverts est
ouverte et normale). L'application est continue puisque chaque $G \to G/N$
l'est.
Elle est injective parce que $G$ est séparé et que tout voisinage de $e$
contient un tel $N$, d'après les arguments précédents. Elle est surjective par
le Lemme \ref{topology-lemma-intersection-closed-in-quasi-compact}. Enfin, le
Lemme \ref{topology-lemma-bijective-map} montre que c'est un homéomorphisme.
\end{proof}

\begin{definition}
\label{topology-definition-profinite-group}
Un groupe topologique est appelé un {\it groupe profini} s'il satisfait aux
conditions équivalentes du Lemme \ref{topology-lemma-profinite-group}.
\end{definition}

\noindent
Si $G_1 \to G_2 \to G_3 \to \ldots$ est un système de groupes topologiques,
alors sa colimite $G = \colim G_n$ dans la catégorie des groupes topologiques
(Lemme \ref{topology-lemma-topological-group-colimits}) diffère en général de sa
colimite dans la catégorie des espaces topologiques
(Lemme \ref{topology-lemma-colimits}), bien que ces deux colimites aient le même
ensemble sous-jacent. Voir Exemples,
section \ref{examples-section-colimit-topology}.

\begin{lemma}
\label{topology-lemma-topological-group-colimits}
La catégorie des groupes topologiques possède des colimites, et le foncteur
d'oubli vers la catégorie des groupes commute aux colimites.
\end{lemma}

\begin{proof}
Pour démontrer l'existence des colimites, nous utiliserons l'argument de
Catégories, Remarque \ref{categories-remark-how-to-use-it}. Supposons donc que
$\mathcal{I} \to \textit{Top}$, $i \mapsto G_i$, soit un foncteur à valeurs
dans la catégorie $\textit{TopGroup}$ des groupes topologiques. Considérons
$$
F : \textit{TopGroup} \longrightarrow \textit{Ensembles},\quad
H \longmapsto \lim_\mathcal{I} \Mor_{\textit{TopGroup}}(G_i, H)
$$
Ce foncteur commute aux limites. De plus, étant donné un groupe topologique
$H$ et un élément $(\varphi_i : G_i \to H)$ de $F(H)$, il existe un
sous-groupe $H' \subset H$ de cardinal au plus $|\coprod G_i|$ (le coproduit
étant pris dans la catégorie des groupes, c'est-à-dire le produit libre des
$G_i$) tel que les morphismes $\varphi_i$ soient à valeurs dans $H'$. On peut
en effet choisir le sous-groupe engendré par les images des $\varphi_i$, muni
de la topologie induite. Les hypothèses de Catégories, Lemme
\ref{categories-lemma-a-version-of-brown}, sont donc satisfaites, et l'on
obtient un groupe topologique $G$ qui représente le foncteur $F$, ce qui
signifie précisément que $G$ est la colimite du diagramme $i \mapsto G_i$.

\medskip\noindent
Pour montrer que le foncteur d'oubli vers les groupes commute aux colimites,
nous utilisons Catégories, Lemme \ref{categories-lemma-adjoint-exact}. En
effet, le foncteur d'oubli possède un adjoint à droite : le foncteur qui associe
à un groupe le groupe topologique chaotique (ou indiscret) correspondant.
\end{proof}

\begin{definition}
\label{topology-definition-topological-ring}
Un {\it anneau topologique} est un anneau $R$ muni d'une topologie telle que
l'addition $R \times R \to R$, $(x, y) \mapsto x + y$, et la multiplication
$R \times R \to R$, $(x, y) \mapsto xy$, soient continues.
Un {\it homomorphisme d'anneaux topologiques} est un homomorphisme d'anneaux
qui est continu.
\end{definition}

\noindent
Dans le projet Champs, les anneaux sont commutatifs et possèdent une unité,
notée $1$. Si $R$ est un anneau topologique, alors $(R, +)$ est un groupe
topologique, puisque
$x \mapsto -x$ est continue. Si $R$ est un anneau topologique et
$R' \subset R$ un sous-anneau, alors $R'$, muni de la topologie induite, est
un anneau topologique. Si $R$ est un anneau topologique et $R \to R'$ une
surjection d'anneaux, alors $R'$, muni de la topologie quotient, est un anneau
topologique.

\begin{lemma}
\label{topology-lemma-topological-ring-limits}
La catégorie des anneaux topologiques possède des limites, et les foncteurs
d'oubli vers (a) la catégorie des espaces topologiques et (b) la catégorie des
anneaux commutent aux limites.
\end{lemma}

\begin{proof}
Il suffit d'établir ces assertions pour les produits et les égalisateurs ;
voir Catégories, Lemme
\ref{categories-lemma-limits-products-equalizers}. Soient $R_i$, $i \in I$,
une famille d'anneaux topologiques. Munissons le produit usuel
$R = \prod R_i$ de la topologie produit. Puisque
$R \times R = \prod (R_i \times R_i)$ comme espace topologique (les produits
commutent entre eux dans toute catégorie), l'addition et la multiplication
sur $R$ sont continues. Soient $a, b : R \to R'$ deux homomorphismes
d'anneaux topologiques. Pour l'égalisateur, il suffit alors de prendre
l'égalisateur de $a$ et $b$ dans la catégorie des espaces topologiques : il
s'agit aussi de leur égalisateur dans la catégorie des anneaux, muni de la
topologie induite.
\end{proof}

\begin{lemma}
\label{topology-lemma-topological-ring-colimits}
La catégorie des anneaux topologiques possède des colimites, et le foncteur
d'oubli vers la catégorie des anneaux commute aux colimites.
\end{lemma}

\begin{proof}
Le même argument que dans la preuve du
Lemme \ref{topology-lemma-topological-group-colimits} démontre l'existence des
colimites. Pour montrer que le foncteur d'oubli vers les anneaux commute aux
colimites, nous utilisons Catégories, Lemme
\ref{categories-lemma-adjoint-exact}. En effet, le foncteur d'oubli possède un
adjoint à droite : le foncteur qui associe à un anneau l'anneau topologique
chaotique (ou indiscret) correspondant.
\end{proof}

\begin{definition}
\label{topology-definition-topological-module}
Soit $R$ un anneau topologique. Un {\it module topologique} est un $R$-module
$M$ muni d'une topologie telle que l'addition $M \times M \to M$ et la
multiplication par les scalaires $R \times M \to M$ soient continues.
Un {\it homomorphisme de modules topologiques} est un homomorphisme de modules
qui est continu.
\end{definition}

\noindent
Si $R$ est un anneau topologique et $M$ un module topologique, alors
$(M, +)$ est un groupe topologique, puisque $x \mapsto -x$ est continue.
Si $R$ est un anneau topologique, $M$ un module topologique et
$M' \subset M$ un sous-module, alors $M'$, muni de la topologie induite, est
un module topologique. Si $R$ est un anneau topologique, $M$ un module
topologique et $M \to M'$ une surjection de modules, alors $M'$, muni de la
topologie quotient, est un module topologique.

\begin{lemma}
\label{topology-lemma-topological-module-limits}
Soit $R$ un anneau topologique. La catégorie des modules topologiques sur $R$
possède des limites, et les foncteurs d'oubli vers (a) la catégorie des espaces
topologiques et (b) la catégorie des $R$-modules commutent aux limites.
\end{lemma}

\begin{proof}
Il suffit d'établir ces assertions pour les produits et les égalisateurs ;
voir Catégories, Lemme
\ref{categories-lemma-limits-products-equalizers}. Soient $M_i$, $i \in I$,
une famille de modules topologiques sur $R$. Munissons le produit usuel
$M = \prod M_i$ de la topologie produit. Puisque
$M \times M = \prod (M_i \times M_i)$ comme espace topologique (les produits
commutent entre eux dans toute catégorie), l'addition sur $M$ est continue. Il
en va de même de la multiplication par les scalaires $R \times M \to M$.
Soient $a, b : M \to M'$ deux homomorphismes de modules topologiques sur $R$.
Pour l'égalisateur, il suffit alors de prendre l'égalisateur de $a$ et $b$
dans la catégorie des espaces topologiques : il s'agit aussi de leur
égalisateur dans la catégorie des modules, muni de la topologie induite.
\end{proof}

\begin{lemma}
\label{topology-lemma-topological-module-colimits}
Soit $R$ un anneau topologique. La catégorie des modules topologiques sur $R$
possède des colimites, et le foncteur d'oubli vers la catégorie des modules sur
$R$ commute aux colimites.
\end{lemma}

\begin{proof}
Le même argument que dans la preuve du
Lemme \ref{topology-lemma-topological-group-colimits} démontre l'existence des
colimites. Pour montrer que le foncteur d'oubli vers les $R$-modules commute
aux colimites, nous utilisons Catégories, Lemme
\ref{categories-lemma-adjoint-exact}. En effet, le foncteur d'oubli possède un
adjoint à droite : le foncteur qui associe à un module le module topologique
chaotique (ou indiscret) correspondant.
\end{proof}









% OK, start here.
%

\chapter{Faisceaux sur les espaces}



\phantomsection
\label{sheaves-section-phantom}


\section{Introduction}
\label{sheaves-section-introduction}

\noindent
Les propriétés fondamentales des faisceaux sur les espaces topologiques
seront exposées dans ce chapitre. On pourra consulter \cite{Godement}.

\medskip\noindent
La discussion des faisceaux sur les sites donnée plus loin dans ces documents
remplacera cet exposé. Mais il est peut-être utile de définir brièvement ici
certaines des notions en jeu.











\section{Notions de base}
\label{sheaves-section-sheaves-basic}

\noindent
Voici une liste de notions fondamentales de topologie.

\begin{enumerate}
\item Soit $X$ un espace topologique. L'expression : « Soit
$U = \bigcup_{i \in I} U_i$ un recouvrement ouvert » signifie ce qui
suit : $I$ est un ensemble et, pour chaque $i \in I$, on se donne
une partie ouverte $U_i \subset X$ telle que $U$ soit la
réunion des $U_i$. On autorise $I = \emptyset$ ;
dans ce cas, il n'y a aucun $U_i$ et $U = \emptyset$.
Lorsque $I \not = \emptyset$, il est également permis que
certains, voire tous, les $U_i$ soient vides.
\item etc., etc.
\end{enumerate}












\section{Préfaisceaux}
\label{sheaves-section-presheaves}


\begin{definition}
\label{sheaves-definition-presheaf}
Soit $X$ un espace topologique.
\begin{enumerate}
\item Un {\it préfaisceau d'ensembles $\mathcal{F}$ sur $X$} est une règle qui
associe à chaque ouvert $U \subset X$ un ensemble $\mathcal{F}(U)$ et
à chaque inclusion $V \subset U$ une application
$\rho^U_V : \mathcal{F}(U) \to \mathcal{F}(V)$ telle que
$\rho^U_U = \text{id}_{\mathcal{F}(U)}$ et que,
si $W \subset V \subset U$, on ait
$\rho^U_W = \rho^V_W \circ \rho ^U_V$.
\item Un {\it morphisme $\varphi : \mathcal{F} \to \mathcal{G}$
de préfaisceaux d'ensembles sur $X$} est une règle qui associe à chaque
ouvert $U \subset X$ une application d'ensembles $\varphi : \mathcal{F}(U)
\to \mathcal{G}(U)$ compatible avec les applications de restriction,
c'est-à-dire que, lorsque $V \subset U \subset X$ sont ouverts, le
diagramme
$$
\xymatrix{
\mathcal{F}(U) \ar[r]^\varphi \ar[d]^{\rho^U_V} &
\mathcal{G}(U) \ar[d]^{\rho^U_V} \\
\mathcal{F}(V) \ar[r]^\varphi & \mathcal{G}(V)
}
$$
est commutatif.
\item La catégorie des préfaisceaux d'ensembles sur $X$ sera notée
$\textit{PSh}(X)$.
\end{enumerate}
\end{definition}

\noindent
Les éléments de l'ensemble $\mathcal{F}(U)$ sont appelés
les {\it sections} de $\mathcal{F}$ sur $U$.
Pour tout $V \subset U$, l'application
$\rho^U_V : \mathcal{F}(U) \to \mathcal{F}(V)$
est appelée l'{\it application de restriction}. Nous emploierons la
notation $s|_V := \rho^U_V(s)$ si $s\in \mathcal{F}(U)$.
Cette notation est compatible avec la notion topologique de restriction
des fonctions, car si $W \subset V \subset U$
et si $s$ est une section de $\mathcal{F}$ sur $U$, alors
$s|_W = (s|_V)|_W$, par la propriété des applications de restriction
exprimée dans la définition ci-dessus.

\medskip\noindent
Une autre notation souvent employée consiste à désigner les sections
sur un ouvert $U$ par le symbole $\Gamma(U, -)$ ou par
$H^0(U, -)$. Autrement dit, les égalités suivantes
sont tautologiques :
$$
\Gamma(U, \mathcal{F}) = \mathcal{F}(U) = H^0(U, \mathcal{F}).
$$
Nous n'emploierons pas cette notation dans ce chapitre, mais nous
l'emploierons dans d'autres chapitres.

\begin{definition}
\label{sheaves-definition-constant-presheaf}
Soit $X$ un espace topologique. Soit $A$ un ensemble.
Le {\it préfaisceau constant de valeur $A$} est le
préfaisceau qui associe l'ensemble $A$ à tout ouvert
$U \subset X$ et dont toutes les applications de restriction
sont $\text{id}_A$.
\end{definition}

\section{Préfaisceaux abéliens}
\label{sheaves-section-abelian-presheaves}

\noindent
Dans cette section, nous indiquons brièvement quelques propriétés de la
catégorie des préfaisceaux qui permettent de définir les préfaisceaux
de groupes abéliens.

\begin{example}
\label{sheaves-example-singleton-presheaf}
Soit $X$ un espace topologique. Considérons une règle $\mathcal{F}$ qui
associe un singleton à toute partie ouverte de $X$. Comme tout ensemble admet
une unique application vers un singleton, les applications de restriction
$\rho^U_V$ existent et sont uniques. La structure ainsi obtenue est un
préfaisceau d'ensembles sur $X$. C'est un objet final de la catégorie des
préfaisceaux d'ensembles sur $X$, en vertu de la propriété des singletons
mentionnée ci-dessus. Il est donc
également unique à isomorphisme unique près. Nous désignerons parfois ce
préfaisceau par $*$.
\end{example}

\begin{lemma}
\label{sheaves-lemma-product-presheaves}
Soit $X$ un espace topologique. La catégorie des préfaisceaux d'ensembles
sur $X$ admet des produits (voir
Catégories, Définition \ref{categories-definition-product}).
De plus, l'ensemble des
sections du produit $\mathcal{F} \times \mathcal{G}$
sur un ouvert $U$ est le produit des ensembles de sections de
$\mathcal{F}$ et $\mathcal{G}$ sur $U$.
\end{lemma}

\begin{proof}
En effet, supposons que $\mathcal{F}$ et $\mathcal{G}$ soient des
préfaisceaux d'ensembles sur l'espace topologique $X$.
Considérons la règle $U \mapsto \mathcal{F}(U) \times \mathcal{G}(U)$,
notée $\mathcal{F} \times \mathcal{G}$. Si $V \subset U \subset X$
sont ouverts, définissons l'application de restriction
$$
(\mathcal{F} \times \mathcal{G})(U)
\longrightarrow
(\mathcal{F} \times \mathcal{G})(V)
$$
par la formule $(s, t) \mapsto (s|_V, t|_V)$. Il est alors immédiatement
clair que $\mathcal{F} \times \mathcal{G}$ est un préfaisceau.
Il existe en outre des morphismes de projection
$p : \mathcal{F} \times \mathcal{G} \to \mathcal{F}$
et
$q : \mathcal{F} \times \mathcal{G} \to \mathcal{G}$.
Nous laissons au lecteur le soin de montrer que,
pour tout autre préfaisceau $\mathcal{H}$, on a
$\Mor(\mathcal{H}, \mathcal{F} \times \mathcal{G})
= \Mor(\mathcal{H}, \mathcal{F}) \times
\Mor(\mathcal{H}, \mathcal{G})$.
\end{proof}

\noindent
Rappelons que si $(A, + : A \times A \to A, - : A \to A, 0\in A)$
est un groupe abélien, alors le zéro et l'application de passage à l'opposé
sont déterminés de manière unique par la loi d'addition. Autrement dit, il
est légitime de dire : « Soit $(A, +)$ un groupe abélien. »

\begin{lemma}
\label{sheaves-lemma-abelian-presheaves}
Soit $X$ un espace topologique.
Soit $\mathcal{F}$ un préfaisceau d'ensembles.
Considérons les types de structures suivants sur $\mathcal{F}$ :
\begin{enumerate}
\item Pour tout ouvert $U$, une structure de groupe abélien
sur $\mathcal{F}(U)$ telle que toutes les applications de restriction soient des
homomorphismes de groupes abéliens.
\item Un morphisme de préfaisceaux
$+ : \mathcal{F} \times \mathcal{F} \to \mathcal{F}$,
un morphisme de préfaisceaux $- : \mathcal{F} \to \mathcal{F}$
et un morphisme $0 : * \to \mathcal{F}$
(voir l'Exemple \ref{sheaves-example-singleton-presheaf})
satisfaisant tous les axiomes vérifiés par $+, -, 0$ dans un
groupe abélien usuel.
\item Un morphisme de préfaisceaux
$+ : \mathcal{F} \times \mathcal{F} \to \mathcal{F}$,
un morphisme de préfaisceaux $- : \mathcal{F} \to \mathcal{F}$
et un morphisme $0 : * \to \mathcal{F}$
tels que, pour chaque ouvert $U \subset X$, le quadruplet
$(\mathcal{F}(U), +, -, 0)$ soit un groupe abélien,
\item Un morphisme de préfaisceaux $+ : \mathcal{F} \times \mathcal{F}
\to \mathcal{F}$ tel que, pour tout ouvert $U \subset X$,
l'application $+ : \mathcal{F}(U) \times \mathcal{F}(U) \to \mathcal{F}(U)$
définisse une structure de groupe abélien.
\end{enumerate}
Il existe des bijections naturelles entre les collections de
données des types (1) à (4) ci-dessus.
\end{lemma}

\begin{proof}
Omis.
\end{proof}

\noindent
Le lemme affirme que munir $\mathcal{F}$ d'une structure d'objet en groupes
abéliens dans la catégorie des préfaisceaux revient à se donner un préfaisceau
d'ensembles
$\mathcal{F}$ tel que tous les ensembles $\mathcal{F}(U)$ soient munis
d'une structure de groupe abélien et que toutes les applications de restriction
soient des homomorphismes de groupes. Pour la plupart des structures algébriques,
nous procéderons de même pour les (pré)faisceaux de tels objets : autrement dit,
nous définirons un (pré)faisceau de tels objets comme un (pré)faisceau
$\mathcal{F}$ d'ensembles dont tous les ensembles de sections $\mathcal{F}(U)$
sont munis de cette structure d'une manière compatible avec les applications
de restriction.

\begin{definition}
\label{sheaves-definition-abelian-presheaves}
Soit $X$ un espace topologique.
\begin{enumerate}
\item Un {\it préfaisceau de groupes abéliens sur $X$}, ou un
{\it préfaisceau abélien sur $X$},
est un préfaisceau d'ensembles $\mathcal{F}$ tel que, pour chaque ouvert
$U \subset X$, l'ensemble $\mathcal{F}(U)$ soit muni
d'une structure de groupe abélien et que toutes les applications de restriction
$\rho^U_V$ soient des homomorphismes de groupes abéliens ; voir le
Lemme \ref{sheaves-lemma-abelian-presheaves} ci-dessus.
\item Un {\it morphisme de préfaisceaux abéliens sur $X$}
$\varphi : \mathcal{F} \to \mathcal{G}$ est un morphisme de préfaisceaux
d'ensembles qui induit
un homomorphisme de groupes abéliens $\mathcal{F}(U) \to \mathcal{G}(U)$
pour tout ouvert $U \subset X$.
\item La catégorie des préfaisceaux de groupes abéliens sur $X$ est notée
$\textit{PAb}(X)$.
\end{enumerate}
\end{definition}

\begin{example}
\label{sheaves-example-direct-sum-points}
Soit $X$ un espace topologique. Pour chaque $x \in X$, soit
$M_x$ un groupe abélien. Pour $U \subset X$ ouvert,
posons
$$
\mathcal{F}(U) = \bigoplus\nolimits_{x \in U} M_x.
$$
Un élément de ce groupe abélien sera noté sous la forme
$\sum_{i = 1}^n m_{x_i}$, où $x_i \in U$ et $m_{x_i} \in M_{x_i}$.
(Bien entendu, nous pouvons toujours choisir cette écriture de telle sorte que
$x_1, \ldots, x_n$ soient deux à deux distincts.)
Pour des ouverts $V \subset U \subset X$, nous définissons une application de
restriction $\mathcal{F}(U) \to \mathcal{F}(V)$ qui
envoie un élément $s = \sum_{i = 1}^n m_{x_i}$
sur l'élément $s|_V = \sum_{x_i \in V} m_{x_i}$.
Nous laissons au lecteur le soin de vérifier qu'il s'agit d'un
préfaisceau de groupes abéliens.
\end{example}

\section{Préfaisceaux de structures algébriques}
\label{sheaves-section-presheaves-structures}

\noindent
Précisons la définition
des préfaisceaux de structures algébriques.
Supposons que $\mathcal{C}$ soit une catégorie et
que $F : \mathcal{C} \to \textit{Ensembles}$ soit
un foncteur fidèle. En général, $F$ est un « foncteur
d'oubli ». Pour un objet $M \in \Ob(\mathcal{C})$,
nous appelons souvent $F(M)$ l'{\it ensemble sous-jacent} de
l'objet $M$. Si $M \to M'$ est un morphisme de $\mathcal{C}$,
nous appelons $F(M) \to F(M')$ l'{\it application ensembliste sous-jacente}.
En fait, nous ne distinguerons souvent pas un objet
de son ensemble sous-jacent, et de même pour les morphismes.
Nous dirons donc qu'une application d'ensembles $F(M) \to F(M')$
est un {\it morphisme de structures algébriques} si elle est
égale à $F(f)$ pour un certain morphisme $f : M \to M'$
de $\mathcal{C}$.

\medskip\noindent
Par analogie avec la définition \ref{sheaves-definition-abelian-presheaves}
ci-dessus, un « préfaisceau d'objets de $\mathcal{C}$ » pourrait être
défini par les données suivantes :
\begin{enumerate}
\item un préfaisceau d'ensembles $\mathcal{F}$, et
\item pour tout ouvert $U \subset X$, le choix
d'un objet $A(U) \in \Ob(\mathcal{C})$,
\end{enumerate}
soumises aux conditions suivantes (avec la terminologie ci-dessus) :
\begin{enumerate}
\item pour tout ouvert $U \subset X$, l'ensemble $\mathcal{F}(U)$
est l'ensemble sous-jacent de $A(U)$, et
\item pour tous ouverts $V \subset U \subset X$,
l'application d'ensembles $\rho_V^U: \mathcal{F}(U) \to \mathcal{F}(V)$
est un morphisme de structures algébriques.
\end{enumerate}
Autrement dit, pour tous ouverts $V \subset U$ de $X$,
l'application de restriction $\rho^U_V$ est l'image
$F(\alpha^U_V)$ d'un unique morphisme
$\alpha^U_V : A(U) \to A(V)$ de la catégorie $\mathcal{C}$.
L'unicité découle de ce que $F$ est
fidèle ; elle implique également que
$\alpha^U_W = \alpha^V_W \circ \alpha^U_V$
lorsque $W \subset V \subset U$ sont des ouverts de $X$.
Le système $(A(-), \alpha^U_V)$ est ce que nous définirons comme un
préfaisceau à valeurs dans $\mathcal{C}$ sur $X$ ; comparer
Sites, définition \ref{sites-definition-presheaf}.
Nous retrouvons notre préfaisceau d'ensembles $(\mathcal{F}, \rho_V^U)$
au moyen des règles $\mathcal{F}(U) = F(A(U))$ et
$\rho_V^U = F(\alpha_V^U)$.

\begin{definition}
\label{sheaves-definition-presheaf-values-in-category}
Soit $X$ un espace topologique.
Soit $\mathcal{C}$ une catégorie.
\begin{enumerate}
\item Un {\it préfaisceau $\mathcal{F}$ sur $X$ à valeurs dans $\mathcal{C}$}
est la donnée d'une règle qui associe à tout ouvert $U \subset X$
un objet $\mathcal{F}(U)$ de $\mathcal{C}$
et à toute inclusion $V \subset U$
un morphisme $\rho_V^U : \mathcal{F}(U) \to \mathcal{F}(V)$
de $\mathcal{C}$ tel que, lorsque $W \subset V \subset U$,
on ait $\rho_W^U = \rho_W^V \circ \rho_V^U$.
\item Un {\it morphisme $\varphi : \mathcal{F} \to \mathcal{G}$
de préfaisceaux à valeurs dans $\mathcal{C}$} est donné par un
morphisme $\varphi : \mathcal{F}(U) \to \mathcal{G}(U)$
de $\mathcal{C}$ compatible avec les morphismes de restriction.
\end{enumerate}
\end{definition}

\begin{definition}
\label{sheaves-definition-underlying-presheaf-sets}
Soit $X$ un espace topologique. Soit $\mathcal{C}$ une catégorie.
Soit $F : \mathcal{C} \to \textit{Ensembles}$ un foncteur fidèle.
Soit $\mathcal{F}$ un préfaisceau sur $X$ à valeurs dans $\mathcal{C}$.
Le préfaisceau d'ensembles $U \mapsto F(\mathcal{F}(U))$
est appelé le {\it préfaisceau d'ensembles sous-jacent à $\mathcal{F}$}.
\end{definition}

\noindent
Il est d'usage d'employer la même lettre $\mathcal{F}$ pour désigner
le préfaisceau d'ensembles sous-jacent, et cela est
justifié par la discussion qui précède la
définition \ref{sheaves-definition-presheaf-values-in-category}.
En particulier, l'expression « soit $s \in \mathcal{F}(U)$ »
ou « soit $s$ une section de $\mathcal{F}$ au-dessus de $U$ » signifie
que $s \in F(\mathcal{F}(U))$.

\medskip\noindent
Cette notation et ces définitions s'appliquent en particulier aux
{\it préfaisceaux de groupes (non nécessairement abéliens), d'anneaux, de modules
sur un anneau fixé, d'espaces vectoriels sur un corps fixé, } etc., ainsi qu'aux
{\it morphismes entre ceux-ci}.

\section{Préfaisceaux de modules}
\label{sheaves-section-presheaves-modules}

\noindent
Supposons que $\mathcal{O}$ soit un préfaisceau d'anneaux sur $X$.
Nous voudrions définir la notion de préfaisceau de
$\mathcal{O}$-modules sur $X$. Par analogie avec la définition
\ref{sheaves-definition-abelian-presheaves}, nous serions tentés de le définir
comme un préfaisceau d'ensembles $\mathcal{F}$ tel que, pour tout ouvert
$U \subset X$, l'ensemble $\mathcal{F}(U)$ soit muni d'une structure
de $\mathcal{O}(U)$-module compatible avec les applications de restriction
(de $\mathcal{F}$ et de $\mathcal{O}$). Il est cependant d'usage
(et équivalent) d'adopter la définition suivante.

\begin{definition}
\label{sheaves-definition-presheaf-modules}
Soit $X$ un espace topologique, et soit $\mathcal{O}$
un préfaisceau d'anneaux sur $X$.
\begin{enumerate}
\item Un {\it préfaisceau de $\mathcal{O}$-modules}
est la donnée d'un préfaisceau abélien $\mathcal{F}$ et d'une
application de préfaisceaux d'ensembles
$$
\mathcal{O} \times \mathcal{F} \longrightarrow \mathcal{F}
$$
telle que, pour tout ouvert $U \subset X$, l'application
$\mathcal{O}(U) \times \mathcal{F}(U) \to \mathcal{F}(U)$
définisse une structure de $\mathcal{O}(U)$-module
sur le groupe abélien $\mathcal{F}(U)$.
\item Un {\it morphisme $\varphi : \mathcal{F} \to \mathcal{G}$
de préfaisceaux de $\mathcal{O}$-modules} est un morphisme de préfaisceaux abéliens
$\varphi : \mathcal{F} \to \mathcal{G}$ tel que
le diagramme
$$
\xymatrix{
\mathcal{O} \times \mathcal{F} \ar[r] \ar[d]_{\text{id} \times \varphi} &
\mathcal{F} \ar[d]^{\varphi} \\
\mathcal{O} \times \mathcal{G} \ar[r] &
\mathcal{G}
}
$$
soit commutatif.
\item L'ensemble des morphismes de $\mathcal{O}$-modules ci-dessus est
noté $\Hom_\mathcal{O}(\mathcal{F}, \mathcal{G})$.
\item La catégorie des préfaisceaux de $\mathcal{O}$-modules
est notée $\textit{PMod}(\mathcal{O})$.
\end{enumerate}
\end{definition}

\noindent
Supposons que $\mathcal{O}_1 \to \mathcal{O}_2$ soit un
morphisme de préfaisceaux d'anneaux sur $X$. Dans ce cas,
si $\mathcal{F}$ est un préfaisceau de $\mathcal{O}_2$-modules,
nous pouvons considérer $\mathcal{F}$ comme un préfaisceau de
$\mathcal{O}_1$-modules en utilisant la composée
$$
\mathcal{O}_1 \times \mathcal{F}
\to
\mathcal{O}_2 \times \mathcal{F}
\to
\mathcal{F}.
$$
Nous notons parfois ce préfaisceau $\mathcal{F}_{\mathcal{O}_1}$
pour indiquer la restriction des scalaires. Nous appelons ceci
la {\it restriction de $\mathcal{F}$}. Nous obtenons le
foncteur de restriction
$$
\textit{PMod}(\mathcal{O}_2)
\longrightarrow
\textit{PMod}(\mathcal{O}_1)
$$

\medskip\noindent
D'autre part, étant donné un préfaisceau de $\mathcal{O}_1$-modules
$\mathcal{G}$,
nous pouvons construire un préfaisceau de $\mathcal{O}_2$-modules
$\mathcal{O}_2 \otimes_{p, \mathcal{O}_1} \mathcal{G}$
par la règle
$$
\left(\mathcal{O}_2 \otimes_{p, \mathcal{O}_1} \mathcal{G}\right)(U)
=
\mathcal{O}_2(U) \otimes_{\mathcal{O}_1(U)} \mathcal{G}(U)
$$
L'indice $p$ signifie « préfaisceau » et non « point ».
Ce préfaisceau est appelé le {\it préfaisceau produit tensoriel}. Nous obtenons
le foncteur d'{\it extension des scalaires}
$$
\textit{PMod}(\mathcal{O}_1)
\longrightarrow
\textit{PMod}(\mathcal{O}_2)
$$

\begin{lemma}
\label{sheaves-lemma-adjointness-tensor-restrict-presheaves}
Avec $X$, $\mathcal{O}_1$, $\mathcal{O}_2$, $\mathcal{F}$ et
$\mathcal{G}$ comme ci-dessus, il existe une bijection canonique
$$
\Hom_{\mathcal{O}_1}(\mathcal{G}, \mathcal{F}_{\mathcal{O}_1})
=
\Hom_{\mathcal{O}_2}(
\mathcal{O}_2 \otimes_{p, \mathcal{O}_1} \mathcal{G},
\mathcal{F}
)
$$
Autrement dit, les foncteurs de restriction et d'extension des scalaires
sont adjoints.
\end{lemma}

\begin{proof}
Cela résulte du fait que, pour un homomorphisme d'anneaux
$A \to B$, le foncteur de restriction et le foncteur
d'extension des scalaires sont adjoints.
\end{proof}





\section{Faisceaux}
\label{sheaves-section-sheaves}

\noindent
Dans cette section, nous expliquons la condition de faisceau.

\begin{definition}
\label{sheaves-definition-sheaf}
Soit $X$ un espace topologique.
\begin{enumerate}
\item Un {\it faisceau $\mathcal{F}$ d'ensembles sur $X$} est un préfaisceau
d'ensembles qui satisfait la propriété supplémentaire suivante : étant donnés
un recouvrement ouvert quelconque $U = \bigcup_{i \in I} U_i$ et une famille
de sections $s_i \in \mathcal{F}(U_i)$, $i \in I$, telle que
$\forall i, j\in I$
$$
s_i|_{U_i \cap U_j} = s_j|_{U_i \cap U_j}
$$
il existe une unique section $s \in \mathcal{F}(U)$ telle que
$s_i = s|_{U_i}$ pour tout $i \in I$.
\item Un {\it morphisme de faisceaux d'ensembles} est simplement un
morphisme de préfaisceaux d'ensembles.
\item La catégorie des faisceaux d'ensembles sur $X$ est notée
$\Sh(X)$.
\end{enumerate}
\end{definition}

\begin{remark}
\label{sheaves-remark-confusion}
Il subsiste toujours une certaine confusion quant à la
nécessité de dire quelque chose sur l'ensemble des sections d'un
faisceau au-dessus de l'ensemble vide $\emptyset \subset X$.
Cela est nécessaire, et nous l'avons déjà fait si vous lisez
correctement la définition. En effet, remarquez que l'ensemble vide est
recouvert par le recouvrement ouvert vide, si bien que la « famille
de sections $s_i$ » de la définition ci-dessus forme en réalité
un élément du produit vide, qui est l'objet final
de la catégorie dans laquelle le faisceau prend ses valeurs. Autrement dit,
si vous lisez correctement la définition, vous en déduisez automatiquement
que $\mathcal{F}(\emptyset) = \textit{un objet final}$,
ce qui, dans le cas d'un faisceau d'ensembles, est un singleton.
Si cet argument ne vous plaît pas, vous pouvez simplement exiger
que $\mathcal{F}(\emptyset) = \{*\}$.

\medskip\noindent
En particulier, cette condition garantira alors que, si
$U, V \subset X$ sont des ouverts {\it disjoints}, alors
$$
\mathcal{F}(U \cup V) = \mathcal{F}(U) \times \mathcal{F}(V).
$$
(En effet, le produit fibré au-dessus d'un objet final est un produit.)
\end{remark}

\begin{example}
\label{sheaves-example-basic-continuous-maps}
Soient $X$, $Y$ des espaces topologiques.
Considérons la règle $\mathcal{F}$ qui associe à
l'ouvert $U \subset X$ l'ensemble
$$
\mathcal{F}(U) = \{ f : U \to Y \mid f \text{ est continue}\}
$$
muni des applications de restriction évidentes. Nous affirmons que
$\mathcal{F}$ est un faisceau. Pour le voir, supposons que
$U = \bigcup_{i\in I} U_i$ soit un recouvrement ouvert, et que
l'on se donne $f_i \in \mathcal{F}(U_i)$, $i\in I$, avec
$f_i |_{U_i \cap U_j} = f_j|_{U_i \cap U_j}$ pour tous $i, j \in I$.
Dans ce cas, définissons $f : U \to Y$ en prenant $f(u)$
égal à la valeur de $f_i(u)$ pour n'importe quel $i \in I$ tel que
$u \in U_i$. Cette application est bien définie par hypothèse. De plus,
$f : U \to Y$ est une application dont la restriction à $U_i$
coïncide avec l'application continue $f_i$. Par conséquent, $f$ est
manifestement continue !
\end{example}

\noindent
Nous pouvons utiliser le résultat de l'exemple pour définir les
faisceaux constants. Soit donc $A$ un ensemble. Munissons $A$ de
la topologie discrète. Soit $U \subset X$ un ouvert.
Nous avons alors
$$
\{ f : U \to A \mid f\text{ continue}\}
=
\{ f : U \to A \mid f\text{ localement constante}\}.
$$
Ainsi, la règle qui associe à un ouvert l'ensemble des applications
localement constantes à valeurs dans $A$ est un faisceau.

\begin{definition}
\label{sheaves-definition-constant-sheaf}
Soit $X$ un espace topologique. Soit $A$ un ensemble.
Le {\it faisceau constant de valeur $A$}, noté $\underline{A}$ ou
$\underline{A}_X$, est le faisceau qui associe à un ouvert $U \subset X$
l'ensemble de toutes les applications localement constantes $U \to A$, les applications
de restriction étant données par les restrictions de fonctions.
\end{definition}

\begin{example}
\label{sheaves-example-sheaf-product-pointwise}
Soit $X$ un espace topologique. Soit $(A_x)_{x \in X}$
une famille d'ensembles $A_x$ indexée par les points $x \in X$. Nous
allons construire à partir de ces données un faisceau d'ensembles $\Pi$.
Pour tout ouvert $U \subset X$,
$$
\Pi(U) = \prod\nolimits_{x \in U} A_x.
$$
Pour tous ouverts $V \subset U \subset X$, définissons
une application de restriction par la règle suivante : un élément
$s = (a_x)_{x\in U} \in \Pi(U)$ a pour restriction
$s|_V = (a_x)_{x \in V}$. Il est clair que cela
définit un préfaisceau d'ensembles. Nous affirmons que ce préfaisceau est un faisceau.
Soit en effet $U = \bigcup U_i$ un recouvrement ouvert.
Supposons que les $s_i \in \Pi(U_i)$
soient tels que $s_i$ et $s_j$ coïncident sur $U_i \cap U_j$. Écrivons
$s_i = (a_{i, x})_{x\in U_i}$. La condition de compatibilité implique que
$a_{i, x} = a_{j, x}$ dans l'ensemble $A_x$ lorsque $x \in U_i \cap U_j$.
Il existe donc un unique élément $s = (a_x)_{x\in U}$
de $\Pi(U) = \prod_{x\in U} A_x$ tel que
$a_x = a_{i, x}$ dès que $x \in U_i$ pour un certain $i$. Bien entendu, cet
élément $s$ vérifie $s|_{U_i} = s_i$ pour tout $i$.
\end{example}

\begin{example}
\label{sheaves-example-direct-sum-points-not-sheaf}
Soit $X$ un espace topologique.
Supposons que, pour tout $x\in X$, un groupe abélien $M_x$ soit donné.
Considérons le préfaisceau $\mathcal{F} : U \mapsto \bigoplus_{x \in U} M_x$
défini dans l'exemple \ref{sheaves-example-direct-sum-points}. Ce
n'est pas un faisceau en général. Par exemple, si $X$ est un ensemble infini
muni de la topologie discrète, alors la condition de faisceau
impliquerait que $\mathcal{F}(X) = \prod_{x\in X} \mathcal{F}(\{x\})$,
mais, par définition, nous avons $\mathcal{F}(X)
= \bigoplus_{x \in X} M_x = \bigoplus_{x \in X} \mathcal{F}(\{x\})$.
Or une somme directe infinie diffère en général d'un
produit direct infini.

\medskip\noindent
Cependant, si $X$ est un espace topologique tel que tout ouvert
de $X$ soit quasi-compact, alors $\mathcal{F}$ {\it est} un faisceau.
Ceci est laissé en exercice au lecteur.
\end{example}

\section{Faisceaux abéliens}
\label{sheaves-section-abelian-sheaves}

\begin{definition}
\label{sheaves-definition-abelian-sheaf}
Soit $X$ un espace topologique.
\begin{enumerate}
\item Un {\it faisceau abélien sur $X$} ou un
{\it faisceau de groupes abéliens sur $X$}
est un préfaisceau abélien sur $X$ tel que le préfaisceau
d'ensembles sous-jacent soit un faisceau.
\item La catégorie des faisceaux de groupes abéliens
est notée $\textit{Ab}(X)$.
\end{enumerate}
\end{definition}

\noindent
Soit $X$ un espace topologique.
Dans le cas d'un préfaisceau abélien $\mathcal{F}$, la condition de
faisceau relative à un recouvrement ouvert $U = \bigcup U_i$
s'exprime souvent en disant que le complexe de groupes abéliens
$$
0 \to \mathcal{F}(U)
\to \prod\nolimits_i \mathcal{F}(U_i)
\to \prod\nolimits_{(i_0, i_1)} \mathcal{F}(U_{i_0} \cap U_{i_1})
$$
est exact. La première application est l'application usuelle, tandis que la seconde
envoie l'élément $(s_i)_{i \in I}$ sur l'élément
$$
(
s_{i_0}|_{U_{i_0} \cap U_{i_1}} -
s_{i_1}|_{U_{i_0} \cap U_{i_1}}
)_{(i_0, i_1)}
\in \prod\nolimits_{(i_0, i_1)} \mathcal{F}(U_{i_0} \cap U_{i_1})
$$

\section{Faisceaux de structures algébriques}
\label{sheaves-section-sheaves-structures}

\noindent
Précisons la définition des faisceaux de certaines espèces de structures.
Commençons par reformuler la condition de faisceau. Supposons donc que
$\mathcal{F}$ soit un préfaisceau d'ensembles sur l'espace topologique $X$.
La condition de faisceau peut se reformuler comme suit. Soit
$U = \bigcup_{i\in I} U_i$ un recouvrement ouvert. Considérons le
diagramme
$$
\xymatrix{
\mathcal{F}(U) \ar[r]
&
\prod\nolimits_{i\in I}
\mathcal{F}(U_i)
\ar@<1ex>[r] \ar@<-1ex>[r]
&
\prod\nolimits_{(i_0, i_1) \in I \times I}
\mathcal{F}(U_{i_0} \cap U_{i_1})
}
$$
Ici, l'application de gauche est définie par la règle
$s \mapsto \prod_{i \in I} s|_{U_i}$. Les deux applications
de droite sont les applications
$$
\prod\nolimits_i s_i
\mapsto
\prod\nolimits_{(i_0, i_1)} s_{i_0}|_{U_{i_0} \cap U_{i_1}}
\text{ resp. }
\prod\nolimits_i s_i
\mapsto
\prod\nolimits_{(i_0, i_1)} s_{i_1}|_{U_{i_0} \cap U_{i_1}}.
$$
La condition de faisceau signifie exactement que la flèche de gauche
est l'égalisateur des deux flèches de droite. Cette reformulation s'étend
immédiatement aux préfaisceaux à valeurs dans une
catégorie, pourvu que celle-ci admette des produits.

\begin{definition}
\label{sheaves-definition-sheaf-values-in-category}
Soit $X$ un espace topologique. Soit $\mathcal{C}$ une
catégorie admettant des produits. Un préfaisceau $\mathcal{F}$ à
valeurs dans $\mathcal{C}$ sur $X$ est un {\it faisceau}
si, pour tout recouvrement ouvert, le diagramme
$$
\xymatrix{
\mathcal{F}(U) \ar[r]
&
\prod\nolimits_{i\in I}
\mathcal{F}(U_i)
\ar@<1ex>[r] \ar@<-1ex>[r]
&
\prod\nolimits_{(i_0, i_1) \in I \times I}
\mathcal{F}(U_{i_0} \cap U_{i_1})
}
$$
est un diagramme égalisateur dans la catégorie $\mathcal{C}$.
\end{definition}

\noindent
Supposons que $\mathcal{C}$ soit une catégorie et que
$F : \mathcal{C} \to \textit{Ensembles}$ soit un foncteur fidèle.
Un bon exemple à garder à l'esprit est le cas où $\mathcal{C}$
est la catégorie des groupes abéliens et $F$ le foncteur d'oubli.
Considérons un préfaisceau $\mathcal{F}$ à valeurs dans $\mathcal{C}$ sur $X$.
Nous voudrions reformuler la condition précédente en termes
du préfaisceau d'ensembles sous-jacent
(Définition \ref{sheaves-definition-underlying-presheaf-sets}).
Remarquons que le
préfaisceau d'ensembles sous-jacent est un faisceau d'ensembles si et seulement si tous les
diagrammes
$$
\xymatrix{
F(\mathcal{F}(U)) \ar[r]
&
\prod\nolimits_{i\in I}
F(\mathcal{F}(U_i))
\ar@<1ex>[r] \ar@<-1ex>[r]
&
\prod\nolimits_{(i_0, i_1) \in I \times I}
F(\mathcal{F}(U_{i_0} \cap U_{i_1}))
}
$$
d'ensembles -- après application du foncteur d'oubli $F$ -- sont des
diagrammes égalisateurs ! Nous voudrions donc que $\mathcal{C}$ admette
des produits et des égalisateurs et que $F$ commute avec
ces constructions. Cela équivaut à demander que $\mathcal{C}$
admette des limites et que $F$ commute avec elles ; voir
Catégories, Lemme \ref{categories-lemma-limits-products-equalizers}.
Mais cela ne suffit pas encore
(voir l'Exemple \ref{sheaves-example-sheaves-topological-spaces}) ;
il faut aussi que $F$ {\it reflète les isomorphismes}.
Cette propriété signifie que, pour tout morphisme
$f : A \to A'$ de $\mathcal{C}$, $f$ est
un isomorphisme si (et seulement si) $F(f)$ est une bijection.

\begin{lemma}
\label{sheaves-lemma-sheaves-structure}
Supposons que la catégorie $\mathcal{C}$ et
le foncteur $F : \mathcal{C} \to \textit{Ensembles}$
possèdent les propriétés suivantes :
\begin{enumerate}
\item $F$ est fidèle,
\item $\mathcal{C}$ admet des limites et $F$ commute avec elles, et
\item le foncteur $F$ reflète les isomorphismes.
\end{enumerate}
Soit $X$ un espace topologique. Soit $\mathcal{F}$
un préfaisceau à valeurs dans $\mathcal{C}$.
Alors $\mathcal{F}$ est un faisceau si et seulement si le
préfaisceau d'ensembles sous-jacent est un faisceau.
\end{lemma}

\begin{proof}
Supposons que $\mathcal{F}$ soit un faisceau. Alors
$\mathcal{F}(U)$ est l'égalisateur du diagramme
ci-dessus et, par hypothèse, $F(\mathcal{F}(U))$
est l'égalisateur du diagramme correspondant
d'ensembles. Ainsi, $F(\mathcal{F})$ est un faisceau d'ensembles.

\medskip\noindent
Supposons que $F(\mathcal{F})$ soit un faisceau.
Soit $E \in \Ob(\mathcal{C})$
l'égalisateur des deux flèches parallèles de la
Définition \ref{sheaves-definition-sheaf-values-in-category}.
Nous obtenons un morphisme canonique $\mathcal{F}(U) \to E$,
simplement parce que $\mathcal{F}$ est un préfaisceau.
Par hypothèse, l'application induite $F(\mathcal{F}(U)) \to F(E)$
est un isomorphisme, puisque $F(E)$ est l'égalisateur
du diagramme correspondant d'ensembles. Nous voyons donc que
$\mathcal{F}(U) \to E$ est un isomorphisme
par la condition (3) du lemme.
\end{proof}

\noindent
Le lemme s'applique en particulier aux
{\it faisceaux de groupes, d'anneaux, d'algèbres sur un anneau fixé, de modules
sur un anneau fixé, d'espaces vectoriels sur un corps fixé, } etc.
Autrement dit, ce sont les préfaisceaux de groupes, d'anneaux,
de modules sur un anneau fixé, d'espaces vectoriels sur un corps fixé, etc.,
dont le préfaisceau d'ensembles sous-jacent est un faisceau.

\begin{example}
\label{sheaves-example-C0-sheaf-rings}
Soit $X$ un espace topologique. Pour tout ouvert $U \subset X$, considérons
la $\mathbf{R}$-algèbre
$\mathcal{C}^{0}(U) = \{ f : U \to \mathbf{R} \mid f\text{ est continue}\}$.
Les applications de restriction évidentes en font un
préfaisceau de $\mathbf{R}$-algèbres sur $X$.
D'après l'Exemple \ref{sheaves-example-basic-continuous-maps}, c'est un faisceau d'ensembles.
Le Lemme \ref{sheaves-lemma-sheaves-structure} montre donc que c'est un faisceau de
$\mathbf{R}$-algèbres sur $X$.
\end{example}

\begin{example}
\label{sheaves-example-sheaves-topological-spaces}
Considérons la catégorie des espaces topologiques $\textit{Top}$.
Il existe un foncteur fidèle naturel $\textit{Top} \to \textit{Ensembles}$
qui commute aux produits et aux égalisateurs. Mais il ne
reflète pas les isomorphismes. En fait, il s'avère que
l'analogue du Lemme \ref{sheaves-lemma-sheaves-structure} est faux.
En effet, supposons $X = \mathbf{N}$ muni de la topologie
discrète. Donnons-nous des espaces topologiques discrets $A_i$, pour
$i \in \mathbf{N}$. Pour toute partie $U \subset \mathbf{N}$,
définissons $\mathcal{F}(U) = \prod_{i\in U} A_i$, muni de la
topologie discrète. On obtient alors un préfaisceau d'espaces
topologiques dont le préfaisceau d'ensembles sous-jacent est un faisceau ; voir
l'Exemple \ref{sheaves-example-sheaf-product-pointwise}.
Cependant, si chaque $A_i$ possède au moins deux éléments, ce
préfaisceau n'est pas un faisceau d'espaces topologiques
au sens de la Définition \ref{sheaves-definition-sheaf-values-in-category}.
Le lecteur pourra vérifier qu'en munissant chaque
$\mathcal{F}(U) = \prod_{i\in U} A_i$ de la {\it topologie produit}, on
obtient bien un faisceau
d'espaces topologiques sur $X$.
\end{example}


\section{Faisceaux de modules}
\label{sheaves-section-sheaves-modules}

\begin{definition}
\label{sheaves-definition-sheaf-modules}
Soit $X$ un espace topologique.
Soit $\mathcal{O}$ un faisceau d'anneaux sur $X$.
\begin{enumerate}
\item Un {\it faisceau de $\mathcal{O}$-modules} est un préfaisceau
de $\mathcal{O}$-modules $\mathcal{F}$,
voir la Définition \ref{sheaves-definition-presheaf-modules},
tel que le préfaisceau de groupes abéliens sous-jacent $\mathcal{F}$
soit un faisceau.
\item Un {\it morphisme de faisceaux de $\mathcal{O}$-modules}
est un morphisme de préfaisceaux de $\mathcal{O}$-modules.
\item Étant donnés des faisceaux de $\mathcal{O}$-modules
$\mathcal{F}$ et $\mathcal{G}$, nous notons
$\Hom_\mathcal{O}(\mathcal{F}, \mathcal{G})$
l'ensemble des morphismes de faisceaux de $\mathcal{O}$-modules.
\item La catégorie des faisceaux de $\mathcal{O}$-modules
est notée $\textit{Mod}(\mathcal{O})$.
\end{enumerate}
\end{definition}

\noindent
Cette définition garde un certain sens même si $\mathcal{O}$ n'est qu'un
préfaisceau d'anneaux, bien que nous ne connaissions aucun exemple où
cela soit utile ; nous éviterons d'employer la terminologie
« faisceaux de $\mathcal{O}$-modules »
lorsque $\mathcal{O}$ n'est pas un faisceau d'anneaux.






\section{Fibres}
\label{sheaves-section-stalks}

\noindent
Soit $X$ un espace topologique. Soit $x \in X$ un point.
Soit $\mathcal{F}$ un préfaisceau d'ensembles sur $X$.
La {\it fibre de $\mathcal{F}$ en $x$} est l'ensemble
$$
\mathcal{F}_x
=
\colim_{x\in U} \mathcal{F}(U)
$$
où la limite inductive est prise sur l'ensemble des voisinages ouverts
$U$ de $x$ dans $X$. L'ensemble des voisinages ouverts est
partiellement ordonné par l'inclusion inverse :
nous posons $U \geq U' \Leftrightarrow U \subset U'$.
Les morphismes de transition du système sont
donnés par les applications de restriction de $\mathcal{F}$.
Voir Catégories, section \ref{categories-section-posets-limits}
pour les notations et la terminologie relatives aux (co)limites de systèmes.
Remarquons qu'il s'agit d'une limite inductive filtrante.
Il est donc facile de décrire $\mathcal{F}_x$. Plus précisément,
$$
\mathcal{F}_x
=
\{
(U, s)
\mid
x\in U, s\in \mathcal{F}(U)
\}/\sim
$$
où la relation d'équivalence est donnée par $(U, s) \sim (U', s')$ si et seulement s'il
existe un ouvert $U'' \subset U \cap U'$ tel que $x \in U''$ et
$s|_{U''} = s'|_{U''}$. Étant donné un couple $(U, s)$, nous notons parfois
$s_x$ l'élément de $\mathcal{F}_x$ correspondant à la classe d'équivalence
de $(U, s)$. Nous employons parfois l'expression
« image de $s$ dans $\mathcal{F}_x$ » pour désigner $s_x$.
Par exemple, étant donnés deux couples $(U, s)$ et $(U', s')$, nous disons parfois
« $s$ est égal à $s'$ dans $\mathcal{F}_x$ » pour indiquer
que $s_x = s'_x$. D'autres auteurs emploient la terminologie
« germe de $s$ en $x$ ».

\medskip\noindent
Une conséquence évidente de cette définition est que,
pour tout ouvert $U \subset X$, il existe une application canonique
$$
\mathcal{F}(U)
\longrightarrow
\prod\nolimits_{x \in U} \mathcal{F}_x
$$
définie par $s \mapsto \prod_{x \in U} (U, s)$. À méditer !

\begin{lemma}
\label{sheaves-lemma-sheaf-subset-stalks}
Soit $\mathcal{F}$ un faisceau d'ensembles sur l'espace topologique $X$.
Pour tout ouvert $U \subset X$, l'application
$$
\mathcal{F}(U)
\longrightarrow
\prod\nolimits_{x \in U} \mathcal{F}_x
$$
est injective.
\end{lemma}

\begin{proof}
Supposons que $s, s' \in \mathcal{F}(U)$ aient la même image
dans chaque fibre $\mathcal{F}_x$ pour tout $x \in U$. Cela signifie que,
pour tout $x \in U$, il existe un ouvert $V^x \subset U$,
$x \in V^x$, tel que $s|_{V^x} = s'|_{V^x}$. Mais alors
$U = \bigcup_{x \in U} V^x$ est un recouvrement ouvert. L'unicité
dans la condition de faisceau entraîne donc que $s = s'$.
\end{proof}

\begin{definition}
\label{sheaves-definition-separated}
Soit $X$ un espace topologique.
Un préfaisceau d'ensembles $\mathcal{F}$ sur $X$ est {\it séparé}
si, pour tout ouvert $U \subset X$, l'application
$\mathcal{F}(U) \to \prod_{x \in U} \mathcal{F}_x$ est
injective.
\end{definition}

\noindent
Observons aussi que la construction de la fibre
$\mathcal{F}_x$ est fonctorielle par rapport au préfaisceau $\mathcal{F}$.
Autrement dit, elle définit un foncteur
$$
\textit{PSh}(X) \longrightarrow \textit{Ensembles},
\ \mathcal{F} \longmapsto \mathcal{F}_x.
$$
Ce foncteur est appelé le {\it foncteur fibre}.
Plus précisément, si $\varphi : \mathcal{F} \to \mathcal{G}$ est
un morphisme de préfaisceaux, nous définissons
$\varphi_x : \mathcal{F}_x \to \mathcal{G}_x$
par la règle $(U, s) \mapsto (U, \varphi(s))$.
Pour vérifier que cette définition a un sens, il faut montrer que,
si $(U, s) = (U', s')$ dans $\mathcal{F}_x$, alors
$(U, \varphi(s)) = (U', \varphi(s'))$ dans $\mathcal{G}_x$.
C'est clair, puisque $\varphi$ est compatible avec les
applications de restriction.

\begin{example}
\label{sheaves-example-stalk-constant-presheaf}
Soit $X$ un espace topologique. Soit $A$ un ensemble.
Notons provisoirement $A_p$ le préfaisceau constant
de valeur $A$ ($p$ pour préfaisceau -- non pour point).
Il existe un morphisme canonique de préfaisceaux
$A_p \to \underline{A}$ vers le faisceau constant de
valeur $A$. En tout point, nous avons des
bijections canoniques $A = (A_p)_x = \underline{A}_x$, où
la seconde application est induite, par fonctorialité, par
l'application $A_p \to \underline{A}$.
\end{example}



\begin{example}
\label{sheaves-example-germs-functions}
Supposons que $X = \mathbf{R}^n$ soit muni de la topologie euclidienne.
Considérons le préfaisceau des fonctions $\mathcal{C}^\infty$
sur $X$, noté $\mathcal{C}^\infty_{\mathbf{R}^n}$.
Autrement dit, $\mathcal{C}^\infty_{\mathbf{R}^n}(U)$ est l'ensemble
des fonctions de classe $\mathcal{C}^\infty$, $f : U \to \mathbf{R}$.
Comme dans l'Exemple \ref{sheaves-example-basic-continuous-maps},
il est facile de montrer que c'est un faisceau. C'est même
un faisceau de $\mathbf{R}$-espaces vectoriels.

\medskip\noindent
Soit ensuite $x \in X = \mathbf{R}^n$ un point. Comment
se représenter un élément de la fibre $\mathcal{C}^\infty_{\mathbf{R}^n, x}$ ?
Un tel élément est représenté par une fonction $\mathcal{C}^\infty$
$f$ dont le domaine contient $x$. Deux telles
fonctions $f$, $g$ déterminent
le même élément de la fibre si elles coïncident sur un voisinage
de $x$. Autrement dit, un élément de $\mathcal{C}^\infty_{\mathbf{R}^n, x}$
n'est autre que ce qu'on appelle parfois
un {\it germe de fonction $\mathcal{C}^\infty$ en $x$}.
\end{example}

\begin{example}
\label{sheaves-example-sheaf-product-pointwise-stalk}
Soit $X$ un espace topologique. Soit $A_x$ un ensemble pour tout $x \in X$.
Considérons le faisceau $\mathcal{F} : U \mapsto \prod_{x\in U} A_x$ de l'Exemple
\ref{sheaves-example-sheaf-product-pointwise}. Nous voulons simplement signaler
ici que la fibre $\mathcal{F}_x$ de $\mathcal{F}$ en $x$
n'est en général {\it pas} égale à l'ensemble $A_x$. Il existe bien entendu
une application $\mathcal{F}_x \to A_x$, mais on ne peut en général rien dire
de plus. Supposons par exemple que $x = \lim x_n$ avec $x_n \not = x_m$ pour
tous $n \not = m$ et que $A_y = \{0, 1\}$ pour tout $y \in X$. Alors
$\mathcal{F}_x$ admet une surjection sur l'ensemble (infini) des queues de suites
de $0$ et de $1$. En effet, tout voisinage ouvert de $x$ contient presque
tous les $x_n$. En revanche, si tout voisinage de $x$ contient
un point $y$ tel que $A_y = \emptyset$, alors $\mathcal{F}_x = \emptyset$.
\end{example}



\section{Fibres des préfaisceaux abéliens}
\label{sheaves-section-stalks-abelian-presheaves}

\noindent
Nous traitons d'abord le cas des groupes abéliens, qui
servira de modèle au cas général.

\begin{lemma}
\label{sheaves-lemma-stalk-abelian-presheaf}
Soit $X$ un espace topologique. Soit $\mathcal{F}$ un préfaisceau
de groupes abéliens sur $X$. Il existe une unique structure de
groupe abélien sur $\mathcal{F}_x$ telle que, pour tout
$U \subset X$ ouvert avec $x\in U$, l'application $\mathcal{F}(U) \to \mathcal{F}_x$
soit un homomorphisme de groupes. De plus,
$$
\mathcal{F}_x
=
\colim_{x\in U} \mathcal{F}(U)
$$
dans la catégorie des groupes abéliens.
\end{lemma}

\begin{proof}
Nous définissons la somme des deux éléments
$(U, s)$ et $(V, t)$ comme le couple $(U \cap V, s|_{U\cap V} +
t|_{U \cap V})$. Le reste se vérifie aisément.
\end{proof}

\noindent
Le point crucial de la démonstration précédente est que l'ensemble partiellement ordonné des
voisinages ouverts est un ensemble dirigé (Catégories, Définition
\ref{categories-definition-directed-set}). En effet, le coproduit
de deux groupes abéliens $A, B$ est la somme directe $A \oplus B$, tandis que
le coproduit dans la catégorie des ensembles est la réunion
disjointe $A \amalg B$, ce qui montre que les limites inductives dans la catégorie
des groupes abéliens ne coïncident pas en général avec les limites inductives dans la
catégorie des ensembles.


\section{Fibres des préfaisceaux de structures algébriques}
\label{sheaves-section-stalks-presheaves-structures}

\noindent
La démonstration du Lemme \ref{sheaves-lemma-stalk-abelian-presheaf} vaut
pour toute espèce de structure algébrique pour laquelle le foncteur
d'oubli commute aux limites inductives filtrantes.

\begin{lemma}
\label{sheaves-lemma-stalk-presheaf-values-in-category}
Soit $\mathcal{C}$ une catégorie. Soit $F : \mathcal{C} \to \textit{Ensembles}$
un foncteur. Supposons que
\begin{enumerate}
\item $F$ soit fidèle, et
\item les limites inductives filtrantes existent dans $\mathcal{C}$ et que $F$ commute avec
elles.
\end{enumerate}
Soit $X$ un espace topologique. Soit $x \in X$. Soit $\mathcal{F}$
un préfaisceau à valeurs dans $\mathcal{C}$.
Alors
$$
\mathcal{F}_x = \colim_{x\in U} \mathcal{F}(U)
$$
existe dans $\mathcal{C}$. Son ensemble sous-jacent est égal à la
fibre du préfaisceau d'ensembles sous-jacent à $\mathcal{F}$.
En outre, la construction $\mathcal{F} \mapsto \mathcal{F}_x$
est un foncteur de la catégorie des préfaisceaux à valeurs dans
$\mathcal{C}$ vers $\mathcal{C}$.
\end{lemma}

\begin{proof}
Omis.
\end{proof}

\noindent
Par définition même, tous les morphismes $\mathcal{F}(U)
\to \mathcal{F}_x$ sont des morphismes de la catégorie $\mathcal{C}$
qui, après application du foncteur d'oubli $F$, deviennent
les applications correspondantes du faisceau d'ensembles sous-jacent.
Comme d'habitude, nous ne distinguerons pas le morphisme
de $\mathcal{C}$ de l'application d'ensembles sous-jacente, ce qui
est permis puisque $F$ est fidèle.

\medskip\noindent
Ce lemme s'applique en particulier aux
{\it préfaisceaux de groupes (non nécessairement abéliens), d'anneaux, de modules
sur un anneau fixé et d'espaces vectoriels sur un corps fixé}.

\section{Fibres des préfaisceaux de modules}
\label{sheaves-section-stalk-presheaves-modules}

\begin{lemma}
\label{sheaves-lemma-stalk-module}
Soit $X$ un espace topologique.
Soit $\mathcal{O}$ un préfaisceau d'anneaux sur $X$.
Soit $\mathcal{F}$ un préfaisceau de $\mathcal{O}$-modules.
Soit $x \in X$.
L'application canonique $\mathcal{O}_x \times \mathcal{F}_x
\to \mathcal{F}_x$ provenant de l'application de multiplication
$\mathcal{O} \times \mathcal{F} \to \mathcal{F}$ définit
une structure de $\mathcal{O}_x$-module sur le groupe abélien
$\mathcal{F}_x$.
\end{lemma}

\begin{proof}
Omis.
\end{proof}

\begin{lemma}
\label{sheaves-lemma-stalk-tensor-presheaf-modules}
Soit $X$ un espace topologique.
Soit $\mathcal{O} \to \mathcal{O}'$ un morphisme de
préfaisceaux d'anneaux sur $X$.
Soit $\mathcal{F}$ un préfaisceau de $\mathcal{O}$-modules.
Soit $x \in X$. Nous avons
$$
\mathcal{F}_x \otimes_{\mathcal{O}_x} \mathcal{O}'_x
=
(\mathcal{F} \otimes_{p, \mathcal{O}} \mathcal{O}')_x
$$
en tant que $\mathcal{O}'_x$-modules.
\end{lemma}

\begin{proof}
Omis.
\end{proof}
\section{Structures algébriques}
\label{sheaves-section-algebraic-structures}

\noindent
Dans cette section, nous formalisons légèrement les notions rencontrées dans
les sections précédentes.

\begin{definition}
\label{sheaves-definition-algebraic-structure}
Une {\it espèce de structure algébrique} est la donnée d'une catégorie
$\mathcal{C}$ et d'un foncteur $F : \mathcal{C} \to \textit{Ensembles}$ ayant les
propriétés suivantes :
\begin{enumerate}
\item $F$ est fidèle,
\item $\mathcal{C}$ admet des limites et $F$ commute aux limites,
\item $\mathcal{C}$ admet des limites inductives filtrantes et $F$ commute avec
celles-ci, et
\item $F$ reflète les isomorphismes.
\end{enumerate}
\end{definition}

\noindent
Nous introduisons cette définition afin de mettre en évidence les propriétés
que nous emploierons dans plusieurs arguments ci-dessous. Nous n'étudierons
cependant pas cette notion en détail, car nos conventions nous interdisent
d'étudier les catégories « grandes », à l'exception de celles énumérées dans
Catégories, remarque \ref{categories-remark-big-categories}.
Parmi celles-ci, les catégories suivantes possèdent les propriétés requises.

\begin{lemma}
\label{sheaves-lemma-list-algebraic-structures}
Les catégories suivantes, munies du foncteur d'oubli évident, définissent des
espèces de structures algébriques :
\begin{enumerate}
\item La catégorie des ensembles pointés.
\item La catégorie des groupes abéliens.
\item La catégorie des groupes.
\item La catégorie des monoïdes.
\item La catégorie des anneaux.
\item La catégorie des $R$-modules, pour un anneau fixé $R$.
\item La catégorie des algèbres de Lie sur un corps fixé.
\end{enumerate}
\end{lemma}

\begin{proof}
Omis.
\end{proof}

\noindent
Désormais, nous considérerons les (pré)faisceaux de structures algébriques et
leurs fibres au moyen des (pré)faisceaux d'ensembles sous-jacents. Cela est
permis par les lemmes \ref{sheaves-lemma-sheaves-structure} et
\ref{sheaves-lemma-stalk-presheaf-values-in-category}.

\medskip\noindent
Dans le reste de cette section, nous signalons quelques résultats sur les
structures algébriques qui seront utiles par la suite.

\begin{lemma}
\label{sheaves-lemma-properties-algebraic-structures}
Soit $(\mathcal{C}, F)$ une espèce de structure algébrique.
\begin{enumerate}
\item $\mathcal{C}$ admet un objet final $0$ et $F(0) = \{ * \}$.
\item $\mathcal{C}$ admet des produits et $F(\prod A_i) = \prod F(A_i)$.
\item $\mathcal{C}$ admet des produits fibrés et
$F(A \times_B C) = F(A)\times_{F(B)}F(C)$.
\item $\mathcal{C}$ admet des égalisateurs et, si $E \to A$ est
l'égalisateur de $a, b : A \to B$, alors $F(E) \to F(A)$ est l'égalisateur
de $F(a), F(b) : F(A) \to F(B)$.
\item $A \to B$ est un monomorphisme si et seulement si
$F(A) \to F(B)$ est injective.
\item si $F(a) : F(A) \to F(B)$ est surjective, alors $a$ est un
épimorphisme.
\item étant donnée une suite $A_1 \to A_2 \to A_3 \to \ldots$, la limite
inductive $\colim A_i$ existe et $F(\colim A_i) = \colim F(A_i)$ ; plus
généralement, il en va de même pour toute limite inductive filtrante.
\end{enumerate}
\end{lemma}

\begin{proof}
Omis. La seule assertion qui mérite commentaire est (5). Elle résulte du fait
que $A \to B$ est un monomorphisme si et seulement si
$A \to A \times_B A$ est un isomorphisme, puis de ce que $F$ reflète les
isomorphismes.
\end{proof}

\begin{lemma}
\label{sheaves-lemma-image-contained-in}
Soit $(\mathcal{C}, F)$ une espèce de structure algébrique.
Supposons que $A, B, C \in \Ob(\mathcal{C})$.
Soient $f : A \to B$ et $g : C \to B$ des morphismes de $\mathcal{C}$.
Si $F(g)$ est injective et si $\Im(F(f)) \subset \Im(F(g))$, alors $f$ se
factorise sous la forme $f = g \circ t$ pour un certain morphisme
$t : A \to C$.
\end{lemma}

\begin{proof}
Considérons $A \times_B C$. Les hypothèses impliquent que
$F(A \times_B C) = F(A) \times_{F(B)} F(C) = F(A)$.
Ainsi, $A = A \times_B C$, puisque $F$ reflète les isomorphismes.
Le résultat s'ensuit.
\end{proof}

\begin{example}
\label{sheaves-example-application-lemma-image-contained-in}
Nous appliquerons souvent le lemme dans la situation suivante.
Supposons donné un diagramme
$$
\xymatrix{
A \ar[r] & B \ar[d] \\
C \ar[r] & D
}
$$
dans $\mathcal{C}$. Supposons que $C \to D$ soit injective sur les ensembles
sous-jacents et que l'image, sur les ensembles sous-jacents, de la composée
$A \to B \to D$ soit contenue dans l'image de $C \to D$.
On obtient alors un diagramme commutatif
$$
\xymatrix{
A \ar[r] \ar[d] & B \ar[d] \\
C \ar[r] & D
}
$$
dans $\mathcal{C}$.
\end{example}

\begin{example}
\label{sheaves-example-sheaf-product-pointwise-algebraic-structure}
Soit $F : \mathcal{C} \to \textit{Ensembles}$ une espèce de structure algébrique.
Soit $X$ un espace topologique. Supposons que, pour tout $x \in X$, un objet
$A_x \in \Ob(\mathcal{C})$ soit donné. Considérons le préfaisceau $\Pi$ à
valeurs dans $\mathcal{C}$ sur $X$, défini par la règle
$\Pi(U) = \prod_{x \in U} A_x$ (avec les applications de restriction
évidentes). Remarquons que le préfaisceau d'ensembles associé
$U \mapsto F(\Pi(U)) = \prod_{x \in U} F(A_x)$ est un faisceau d'après
l'exemple \ref{sheaves-example-sheaf-product-pointwise}.
Ainsi, $\Pi$ est un faisceau de structures algébriques de type
$(\mathcal{C} , F)$. On obtient de cette manière de nombreux exemples de
faisceaux de groupes abéliens, de groupes, d'anneaux, etc.
\end{example}












\section{Exactitude et points}
\label{sheaves-section-exactness-points}

\noindent
Dans toute catégorie, on dispose des notions d'épimorphisme, de
monomorphisme, d'isomorphisme, etc.

\begin{lemma}
\label{sheaves-lemma-points-exactness}
Soit $X$ un espace topologique. Soit $\varphi : \mathcal{F} \to \mathcal{G}$
un morphisme de faisceaux d'ensembles sur $X$.
\begin{enumerate}
\item L'application $\varphi$ est un monomorphisme dans la catégorie des
faisceaux si et seulement si, pour tout $x \in X$, l'application
$\varphi_x : \mathcal{F}_x \to \mathcal{G}_x$ est injective.
\item L'application $\varphi$ est un épimorphisme dans la catégorie des
faisceaux si et seulement si, pour tout $x \in X$, l'application
$\varphi_x : \mathcal{F}_x \to \mathcal{G}_x$ est surjective.
\item L'application $\varphi$ est un isomorphisme dans la catégorie des
faisceaux si et seulement si, pour tout $x \in X$, l'application
$\varphi_x : \mathcal{F}_x \to \mathcal{G}_x$ est bijective.
\end{enumerate}
\end{lemma}

\begin{proof}
Omis.
\end{proof}

\noindent
Il en résulte que, dans la catégorie des faisceaux d'ensembles, les notions
d'épimorphisme et de monomorphisme se décrivent comme suit.

\begin{definition}
\label{sheaves-definition-injective-surjective}
Soit $X$ un espace topologique.
\begin{enumerate}
\item Un préfaisceau $\mathcal{F}$ est appelé un {\it sous-préfaisceau} d'un
préfaisceau $\mathcal{G}$ si
$\mathcal{F}(U) \subset \mathcal{G}(U)$ pour tout ouvert $U \subset X$ et
si les applications de restriction de $\mathcal{G}$ induisent celles de
$\mathcal{F}$. Si $\mathcal{F}$ et $\mathcal{G}$ sont des faisceaux, on dit
que $\mathcal{F}$ est un {\it sous-faisceau} de $\mathcal{G}$. Nous
l'indiquons parfois par la notation $\mathcal{F} \subset \mathcal{G}$.
\item Un morphisme de préfaisceaux d'ensembles
$\varphi : \mathcal{F} \to \mathcal{G}$ sur $X$ est dit {\it injectif} si
et seulement si $\mathcal{F}(U) \to \mathcal{G}(U)$ est injective pour tout
ouvert $U$ de $X$.
\item Un morphisme de préfaisceaux d'ensembles
$\varphi : \mathcal{F} \to \mathcal{G}$ sur $X$ est dit {\it surjectif} si
et seulement si $\mathcal{F}(U) \to \mathcal{G}(U)$ est surjective pour tout
ouvert $U$ de $X$.
\item Un morphisme de faisceaux d'ensembles
$\varphi : \mathcal{F} \to \mathcal{G}$ sur $X$ est dit {\it injectif} si
et seulement si $\mathcal{F}(U) \to \mathcal{G}(U)$ est injective pour tout
ouvert $U$ de $X$.
\item Un morphisme de faisceaux d'ensembles
$\varphi : \mathcal{F} \to \mathcal{G}$ sur $X$ est dit {\it surjectif} si
et seulement si, pour tout ouvert $U$ de $X$ et toute section $s$ de
$\mathcal{G}(U)$, il existe un recouvrement ouvert
$U = \bigcup U_i$ tel que $s|_{U_i}$ appartienne à l'image de
$\mathcal{F}(U_i) \to \mathcal{G}(U_i)$ pour tout $i$.
\end{enumerate}
\end{definition}


\begin{lemma}
\label{sheaves-lemma-characterize-epi-mono}
Soit $X$ un espace topologique.
\begin{enumerate}
\item Les épimorphismes (resp.\ les monomorphismes) dans la catégorie des
préfaisceaux sont exactement les morphismes surjectifs (resp.\ injectifs) de
préfaisceaux.
\item Les épimorphismes (resp.\ les monomorphismes) dans la catégorie des
faisceaux sont exactement les morphismes surjectifs (resp.\ injectifs) de
faisceaux ; ce sont également exactement les morphismes surjectifs
(resp.\ injectifs) sur toutes les fibres.
\end{enumerate}
\end{lemma}

\begin{proof}
Omis.
\end{proof}


\begin{lemma}
\label{sheaves-lemma-check-homomorphism-stalks}
Soit $X$ un espace topologique.
Soit $(\mathcal{C}, F)$ une espèce de structure algébrique.
Supposons que $\mathcal{F}$ et $\mathcal{G}$ soient des faisceaux sur $X$ à
valeurs dans $\mathcal{C}$.
Soit $\varphi : \mathcal{F} \to \mathcal{G}$ une application entre les
faisceaux d'ensembles sous-jacents.
Si, pour tout point $x \in X$, l'application
$\mathcal{F}_x \to \mathcal{G}_x$ est un morphisme de structures
algébriques, alors $\varphi$ est un morphisme de faisceaux de structures
algébriques.
\end{lemma}

\begin{proof}
Soit $U$ un ouvert de $X$. Considérons le diagramme d'ensembles sous-jacents
$$
\xymatrix{
\mathcal{F}(U) \ar[r] \ar[d] &
\prod_{x \in U} \mathcal{F}_x \ar[d] \\
\mathcal{G}(U) \ar[r] &
\prod_{x \in U} \mathcal{G}_x
}
$$
Par hypothèse et d'après les résultats précédents, toutes les flèches sauf la
flèche verticale de gauche sont des morphismes de structures algébriques. De
plus, la flèche horizontale inférieure est injective ; voir le
lemme \ref{sheaves-lemma-sheaf-subset-stalks}.
La conclusion résulte donc du lemme \ref{sheaves-lemma-image-contained-in} ; voir
également l'exemple \ref{sheaves-example-application-lemma-image-contained-in}.
\end{proof}

\noindent
Les suites exactes courtes de faisceaux abéliens, entre autres, seront étudiées
dans le chapitre consacré aux faisceaux de modules. Voir
Modules, section \ref{modules-section-kernels}.



\section{Faisceau associé}
\label{sheaves-section-sheafification}

\noindent
Dans cette section, nous expliquons comment construire le faisceau associé
à un préfaisceau sur un espace topologique. Nous utiliserons les fibres
pour décrire ici le faisceau associé. Cette construction diffère
du procédé général décrit dans Sites, section
\ref{sites-section-sheafification}, et elle est peut-être un peu
plus facile à comprendre.

\medskip\noindent
La construction fondamentale est la suivante. Soit $\mathcal{F}$ un préfaisceau
d'ensembles sur un espace topologique $X$.
Pour tout ouvert $U \subset X$, nous posons
$$
\mathcal{F}^{\#}(U)
=
\{
(s_u) \in \prod\nolimits_{u \in U} \mathcal{F}_u
\text{ tels que }(*)
\}
$$
où $(*)$ désigne la propriété suivante :
\begin{itemize}
\item[(*)] Pour tout $u \in U$, il existe un voisinage ouvert
$u \in V \subset U$ et une section $\sigma \in \mathcal{F}(V)$
tels que, pour tout $v \in V$, on ait $s_v = (V, \sigma)$
dans $\mathcal{F}_v$.
\end{itemize}
Remarquons que $(*)$ est une condition imposée en chaque $u \in U$
et que, $u \in U$ étant fixé, sa validité
ne dépend que des valeurs $s_v$ pour $v$ dans un voisinage ouvert quelconque
de $u$. Il est donc clair que,
si $V \subset U \subset X$ sont ouverts, les applications de projection
$$
\prod\nolimits_{u \in U} \mathcal{F}_u
\longrightarrow
\prod\nolimits_{v \in V} \mathcal{F}_v
$$
envoient les éléments de $\mathcal{F}^{\#}(U)$ dans $\mathcal{F}^{\#}(V)$.
En prenant ces applications pour applications de restriction, on munit $\mathcal{F}^\#$
d'une structure de préfaisceau d'ensembles sur $X$.

\medskip\noindent
De plus, l'application $\mathcal{F}(U) \to \prod_{u \in U} \mathcal{F}_u$
décrite dans la section \ref{sheaves-section-stalks} est manifestement à valeurs
dans $\mathcal{F}^{\#}(U)$. Par ailleurs, si $V \subset U \subset X$ sont
ouverts, on a le diagramme commutatif suivant :
$$
\xymatrix{
\mathcal{F}(U) \ar[r] \ar[d] &
\mathcal{F}^{\#}(U) \ar[r] \ar[d] &
\prod_{u\in U} \mathcal{F}_u \ar[d] \\
\mathcal{F}(V) \ar[r] &
\mathcal{F}^{\#}(V) \ar[r] &
\prod_{v\in V} \mathcal{F}_v
}
$$
où les flèches verticales sont induites par les
applications de restriction. On voit donc qu'il existe
un morphisme canonique de préfaisceaux
$\mathcal{F} \to \mathcal{F}^{\#}$.

\medskip\noindent
Dans l'Exemple \ref{sheaves-example-sheaf-product-pointwise}, nous avons vu
que la règle $\Pi(\mathcal{F}) : U \mapsto \prod_{u\in U} \mathcal{F}_u$
définit un faisceau, muni des applications de restriction évidentes. Par construction,
$\mathcal{F}^{\#}$ en est un sous-préfaisceau. Autrement dit, on
a des morphismes de préfaisceaux
$$
\mathcal{F} \to \mathcal{F}^\# \to \Pi(\mathcal{F}).
$$
En outre, la règle qui associe à $\mathcal{F}$
la suite ci-dessus est manifestement fonctorielle en le préfaisceau $\mathcal{F}$.
Cette notation sera
utilisée dans les démonstrations des lemmes qui suivent.

\begin{lemma}
\label{sheaves-lemma-sheafification-sheaf}
Le préfaisceau $\mathcal{F}^{\#}$ est un faisceau.
\end{lemma}

\begin{proof}
Il vaut sans doute mieux que le lecteur trouve lui-même une explication
plutôt que de lire la démonstration ci-dessous. En fait, le lemme est vrai
pour la même raison que le préfaisceau des fonctions
continues est un faisceau, voir l'Exemple \ref{sheaves-example-basic-continuous-maps}
(et cette analogie peut être précisée au moyen de l'« espace \'etal\'e »).

\medskip\noindent
Quoi qu'il en soit, soit $U = \bigcup U_i$ un recouvrement ouvert.
Supposons donnés des $s_i = (s_{i, u})_{u \in U_i} \in \mathcal{F}^{\#}(U_i)$
tels que $s_i$ et $s_j$ coïncident sur $U_i \cap U_j$.
Comme $\Pi(\mathcal{F})$ est un faisceau,
il existe un élément $s = (s_u)_{u\in U}$ de $\prod_{u\in U} \mathcal{F}_u$
qui se restreint en $s_i$ sur $U_i$. Il reste à vérifier la propriété $(*)$.
Soit $u \in U$. Alors $u \in U_i$ pour un certain $i$. La propriété $(*)$ appliquée à $s_i$
fournit donc un ouvert $V$, avec $u \in V \subset U_i$,
et un $\sigma \in \mathcal{F}(V)$
tels que $s_{i, v} = (V, \sigma)$ dans $\mathcal{F}_v$
pour tout $v \in V$. Puisque $s_{i, v} = s_v$, la propriété $(*)$ vaut pour $s$.
\end{proof}

\begin{lemma}
\label{sheaves-lemma-stalk-sheafification}
Soit $X$ un espace topologique.
Soit $\mathcal{F}$ un préfaisceau d'ensembles sur $X$.
Soit $x \in X$. Alors $\mathcal{F}_x = \mathcal{F}^\#_x$.
\end{lemma}

\begin{proof}
L'application $\mathcal{F}_x \to \mathcal{F}^\#_x$
est injective, puisque l'application
$\mathcal{F}_x \to \Pi(\mathcal{F})_x$ l'est déjà.
En effet, il existe une application canonique $\Pi(\mathcal{F})_x \to \mathcal{F}_x$
qui est un inverse à gauche de l'application $\mathcal{F}_x \to \Pi(\mathcal{F})_x$,
voir l'Exemple \ref{sheaves-example-sheaf-product-pointwise-stalk}.
Pour montrer qu'elle est surjective, soit
$\overline{s} \in \mathcal{F}^\#_x$.
On peut trouver un voisinage ouvert $U$ de $x$ tel que
$\overline{s}$ soit la classe d'équivalence de $(U, s)$,
où $s \in \mathcal{F}^\#(U)$.
Par définition, cela signifie qu'il existe un voisinage ouvert
$V \subset U$ de $x$ et une section $\sigma \in \mathcal{F}(V)$
tels que $s|_V$ soit l'image de $\sigma$ dans $\Pi(\mathcal{F})(V)$.
La classe de $(V, \sigma)$ définit manifestement un élément de
$\mathcal{F}_x$ qui s'envoie sur $\overline{s}$.
\end{proof}

\begin{lemma}
\label{sheaves-lemma-sheafify-universal}
Soit $\mathcal{F}$ un préfaisceau d'ensembles sur $X$.
Tout morphisme $\mathcal{F} \to \mathcal{G}$ vers un faisceau d'ensembles
se factorise de manière unique sous la forme
$\mathcal{F} \to \mathcal{F}^\# \to \mathcal{G}$.
\end{lemma}

\begin{proof}
On a manifestement un diagramme commutatif
$$
\xymatrix{
\mathcal{F} \ar[r] \ar[d] &
\mathcal{F}^\# \ar[r] \ar[d] &
\Pi(\mathcal{F}) \ar[d] \\
\mathcal{G} \ar[r] &
\mathcal{G}^\# \ar[r] &
\Pi(\mathcal{G}) \\
}
$$
Il suffit donc de montrer que $\mathcal{G} = \mathcal{G}^\#$.
D'après le Lemme \ref{sheaves-lemma-points-exactness}, il suffit pour cela de montrer,
pour tout point $x \in X$, que l'application
$\mathcal{G}_x \to \mathcal{G}^\#_x$ est bijective. C'est précisément
le Lemme \ref{sheaves-lemma-stalk-sheafification} ci-dessus.
\end{proof}

\noindent
Ce lemme affirme en réalité que les foncteurs
$i : \Sh(X) \to \textit{PSh}(X)$
(inclusion) et $\# : \textit{PSh}(X) \to \Sh(X)$
(faisceau associé) forment un couple de foncteurs adjoints. La formule est
$$
\Mor_{\textit{PSh}(X)}(\mathcal{F}, i(\mathcal{G}))
=
\Mor_{\Sh(X)}(\mathcal{F}^\#, \mathcal{G})
$$
ce qui signifie que le foncteur faisceau associé est adjoint à gauche du
foncteur d'inclusion. Voir Catégories, section
\ref{categories-section-adjoint}.

\begin{example}
\label{sheaves-example-sheafify-constant}
Pour les notations, voir l'Exemple \ref{sheaves-example-stalk-constant-presheaf}.
L'application $A_p \to \underline{A}$ induit une application
$A_p^\# \to \underline{A}$. On voit aisément que celle-ci
est un isomorphisme. Autrement dit, le faisceau associé
au préfaisceau constant de valeur $A$ est le
faisceau constant de valeur $A$.
\end{example}

\begin{lemma}
\label{sheaves-lemma-separated-presheaf-into-sheaf}
Soit $X$ un espace topologique.
Un préfaisceau $\mathcal{F}$ est séparé (voir la
Définition \ref{sheaves-definition-separated}) si et seulement si
l'application canonique $\mathcal{F} \to \mathcal{F}^\#$ est injective.
\end{lemma}

\begin{proof}
Cela résulte immédiatement de la construction de $\mathcal{F}^\#$ donnée dans cette
section.
\end{proof}

\begin{lemma}
\label{sheaves-lemma-sheafify-epi-mono}
Soit $X$ un espace topologique. Le morphisme de faisceaux associés à un morphisme surjectif
(resp.\ injectif) de préfaisceaux d'ensembles est surjectif
(resp.\ injectif).
\end{lemma}

\begin{proof}
Omis.
\end{proof}


\section{Faisceau associé à un préfaisceau abélien}
\label{sheaves-section-sheafify-abelian-presheaves}

\noindent
Le lemme suivant, d'aspect quelque peu étrange, est vraisemblablement superflu, mais
très commode pour traiter le passage au faisceau associé des préfaisceaux
de structures algébriques.

\begin{lemma}
\label{sheaves-lemma-diagram-fibre-product}
Soit $X$ un espace topologique. Soit $\mathcal{F}$
un préfaisceau d'ensembles sur $X$. Soit $U \subset X$ un ouvert.
On a un diagramme cartésien canonique
$$
\xymatrix{
\mathcal{F}^\#(U) \ar[d] \ar[r] &
\Pi(\mathcal{F})(U) \ar[d] \\
\prod_{x \in U} \mathcal{F}_x
\ar[r] &
\prod_{x \in U} \Pi(\mathcal{F})_x
}
$$
dont les flèches sont les suivantes :
\begin{enumerate}
\item Les composantes de la flèche verticale de gauche sont les applications
$\mathcal{F}^\#(U) \to \mathcal{F}^\#_x = \mathcal{F}_x$
où l'égalité résulte du Lemme \ref{sheaves-lemma-stalk-sheafification}.
\item La flèche horizontale supérieure provient du
morphisme de préfaisceaux $\mathcal{F} \to \Pi(\mathcal{F})$ décrit
dans la section \ref{sheaves-section-sheafification}.
\item Les composantes évidentes de la flèche verticale de droite sont les
applications $\Pi(\mathcal{F})(U) \to \Pi(\mathcal{F})_x$.
\item Les composantes de la flèche horizontale inférieure sont les applications
$\mathcal{F}_x \to \Pi(\mathcal{F})_x$
qui proviennent du morphisme de préfaisceaux
$\mathcal{F} \to \Pi(\mathcal{F})$ décrit
dans la section \ref{sheaves-section-sheafification}.
\end{enumerate}
\end{lemma}

\begin{proof}
Le diagramme est manifestement commutatif. Il faut montrer
qu'il est cartésien. La flèche horizontale inférieure
est injective puisque toutes les applications $\mathcal{F}_x \to \Pi(\mathcal{F})_x$
sont injectives (voir le début de la démonstration du
Lemme \ref{sheaves-lemma-stalk-sheafification}).
Une section $s \in \Pi(\mathcal{F})(U)$ appartient à $\mathcal{F}^\#$ si et
seulement si elle vérifie $(*)$. Or $(*)$ affirme qu'au voisinage de tout point
la section $s$ provient d'une section de $\mathcal{F}$. Par définition
des foncteurs fibres, cela équivaut à dire que
la valeur de $s$ dans toute fibre $\Pi(\mathcal{F})_x$ provient
d'un élément de la fibre $\mathcal{F}_x$. D'où le lemme.
\end{proof}


\begin{lemma}
\label{sheaves-lemma-sheafify-abelian-presheaf}
Soit $X$ un espace topologique.
Soit $\mathcal{F}$ un préfaisceau abélien sur $X$.
Il existe alors une unique structure de
faisceau abélien sur $\mathcal{F}^\#$ telle que
$\mathcal{F} \to \mathcal{F}^\#$ soit un morphisme
de préfaisceaux abéliens. En outre, on a la propriété
d'adjonction suivante :
$$
\Mor_{\textit{PAb}(X)}(\mathcal{F}, i(\mathcal{G}))
=
\Mor_{\textit{Ab}(X)}(\mathcal{F}^\#, \mathcal{G}).
$$
\end{lemma}

\begin{proof}
Rappelons le faisceau d'ensembles $\Pi(\mathcal{F})$ défini dans la
section \ref{sheaves-section-sheafification}. Toutes les fibres
$\mathcal{F}_x$ sont des groupes abéliens, voir le
Lemme \ref{sheaves-lemma-stalk-abelian-presheaf}.
Par conséquent, $\Pi(\mathcal{F})$ est un faisceau de groupes abéliens d'après
l'Exemple \ref{sheaves-example-sheaf-product-pointwise-algebraic-structure}.
Il est également clair que l'application $\mathcal{F} \to \Pi(\mathcal{F})$
est un morphisme de préfaisceaux abéliens. Si l'on montre que
la condition $(*)$ de la section \ref{sheaves-section-sheafification} définit un sous-groupe
de $\Pi(\mathcal{F})(U)$ pour tout ouvert $U \subset X$,
alors $\mathcal{F}^\#$ hérite canoniquement d'une structure de faisceau abélien.
Cela se vérifie aisément directement, et nous laissons au
lecteur le soin d'en trouver un argument simple. Voici l'argument que nous employons,
qui se généralise aux préfaisceaux de structures algébriques :
le Lemme \ref{sheaves-lemma-diagram-fibre-product} montre que
$\mathcal{F}^\#(U)$ est le produit fibré d'un diagramme
de groupes abéliens. Ainsi $\mathcal{F}^\#$ est bien le
sous-groupe abélien voulu.

\medskip\noindent
Remarquons qu'à ce stade $\mathcal{F}^\#_x$ est un groupe
abélien d'après le Lemme \ref{sheaves-lemma-stalk-abelian-presheaf},
et que $\mathcal{F}_x \to \mathcal{F}^\#_x$ est à la fois une
bijection (Lemme \ref{sheaves-lemma-stalk-sheafification})
et un homomorphisme de groupes abéliens. Par conséquent,
$\mathcal{F}_x \to \mathcal{F}^\#_x$ est un isomorphisme
de groupes abéliens. Nous l'utiliserons ci-dessous sans autre mention.

\medskip\noindent
Pour démontrer la propriété d'adjonction, nous utilisons celle du
foncteur faisceau associé pour les préfaisceaux d'ensembles. Par exemple,
si $\psi : \mathcal{F} \to i(\mathcal{G})$ est un morphisme de préfaisceaux,
on obtient un morphisme de faisceaux
$\psi' : \mathcal{F}^\# \to \mathcal{G}$. Il reste à vérifier
qu'il s'agit d'un morphisme de faisceaux abéliens. On peut le faire,
par exemple, en constatant que c'est vrai sur les fibres,
d'après le Lemme \ref{sheaves-lemma-stalk-sheafification}, puis en appliquant le
Lemme \ref{sheaves-lemma-check-homomorphism-stalks} ci-dessus.
\end{proof}



\section{Faisceau associé à un préfaisceau de structures algébriques}
\label{sheaves-section-sheafification-presheaves-structures}


\begin{lemma}
\label{sheaves-lemma-sheafify-presheaf-structures}
Soit $X$ un espace topologique.
Soit $(\mathcal{C}, F)$ une espèce de structure algébrique.
Soit $\mathcal{F}$ un préfaisceau à valeurs dans $\mathcal{C}$
sur $X$. Il existe alors un faisceau $\mathcal{F}^\#$ à valeurs
dans $\mathcal{C}$ et un morphisme $\mathcal{F} \to \mathcal{F}^\#$
de préfaisceaux à valeurs dans $\mathcal{C}$ possédant les
propriétés suivantes :
\begin{enumerate}
\item L'application $\mathcal{F} \to \mathcal{F}^\#$ identifie
le faisceau d'ensembles sous-jacent à $\mathcal{F}^\#$ au
faisceau associé au préfaisceau d'ensembles sous-jacent à $\mathcal{F}$.
\item Pour tout morphisme $\mathcal{F} \to \mathcal{G}$, où
$\mathcal{G}$ est un faisceau à valeurs dans $\mathcal{C}$, il existe
une unique factorisation $\mathcal{F} \to \mathcal{F}^\# \to \mathcal{G}$.
\end{enumerate}
\end{lemma}

\begin{proof}
La démonstration est la même que celle du
Lemme \ref{sheaves-lemma-sheafify-abelian-presheaf},
en appliquant à plusieurs reprises le
Lemme \ref{sheaves-lemma-image-contained-in} (voir aussi
l'Exemple \ref{sheaves-example-application-lemma-image-contained-in}).
L'idée principale consiste toutefois à définir $\mathcal{F}^\#(U)$
comme le produit fibré dans $\mathcal{C}$ du diagramme
$$
\xymatrix{
 &
\Pi(\mathcal{F})(U) \ar[d] \\
\prod_{x \in U} \mathcal{F}_x
\ar[r] &
\prod_{x \in U} \Pi(\mathcal{F})_x
}
$$
à comparer au Lemme \ref{sheaves-lemma-diagram-fibre-product}.
\end{proof}

\section{Faisceau associé à un préfaisceau de modules}
\label{sheaves-section-sheafification-presheaves-modules}

\begin{lemma}
\label{sheaves-lemma-sheafification-presheaf-modules}
Soit $X$ un espace topologique.
Soit $\mathcal{O}$ un préfaisceau d'anneaux sur $X$.
Soit $\mathcal{F}$ un préfaisceau de $\mathcal{O}$-modules.
Soit $\mathcal{O}^\#$ le faisceau associé à $\mathcal{O}$.
Soit $\mathcal{F}^\#$ le faisceau associé à $\mathcal{F}$
considéré comme préfaisceau de groupes abéliens. Il existe un morphisme de
faisceaux d'ensembles
$$
\mathcal{O}^\# \times \mathcal{F}^\#
\longrightarrow
\mathcal{F}^\#
$$
qui rend commutatif le diagramme
$$
\xymatrix{
\mathcal{O} \times \mathcal{F} \ar[r] \ar[d] &
\mathcal{F} \ar[d] \\
\mathcal{O}^\# \times \mathcal{F}^\# \ar[r] &
\mathcal{F}^\#
}
$$
et qui munit $\mathcal{F}^\#$ d'une structure de faisceau
de $\mathcal{O}^\#$-modules. En outre, si $\mathcal{G}$
est un faisceau de $\mathcal{O}^\#$-modules, tout morphisme
de préfaisceaux de $\mathcal{O}$-modules $\mathcal{F} \to \mathcal{G}$
(où $\mathcal{G}$ est considéré, par restriction des scalaires, comme $\mathcal{O}$-module)
se factorise de manière unique sous la forme $\mathcal{F} \to \mathcal{F}^\# \to \mathcal{G}$,
où $\mathcal{F}^\# \to \mathcal{G}$ est un morphisme de
$\mathcal{O}^\#$-modules.
\end{lemma}

\begin{proof}
Omis.
\end{proof}

\noindent
Cela signifie en fait que le foncteur
$i : \textit{Mod}(\mathcal{O}^\#) \to \textit{PMod}(\mathcal{O})$
(qui combine la restriction des scalaires et l'inclusion des faisceaux dans les préfaisceaux)
et le foncteur faisceau associé du lemme
${}^\# : \textit{PMod}(\mathcal{O}) \to \textit{Mod}(\mathcal{O}^\#)$
sont adjoints. En formule :
$$
\Mor_{\textit{PMod}(\mathcal{O})}(\mathcal{F}, i\mathcal{G})
=
\Mor_{\textit{Mod}(\mathcal{O}^\#)}(\mathcal{F}^\#, \mathcal{G})
$$

\medskip\noindent
Soit $X$ un espace topologique.
Soit $\mathcal{O}_1 \to \mathcal{O}_2$ un
morphisme de faisceaux d'anneaux sur $X$.
Dans la section \ref{sheaves-section-presheaves-modules},
nous avons défini un foncteur de restriction des scalaires
et un foncteur d'extension des scalaires sur les préfaisceaux de modules
associés à cette situation.

\medskip\noindent
Si $\mathcal{F}$ est un faisceau de $\mathcal{O}_2$-modules,
alors la restriction des scalaires $\mathcal{F}_{\mathcal{O}_1}$
de $\mathcal{F}$ est manifestement un faisceau
de $\mathcal{O}_1$-modules. On obtient le foncteur de restriction des scalaires
$$
\textit{Mod}(\mathcal{O}_2)
\longrightarrow
\textit{Mod}(\mathcal{O}_1)
$$

\medskip\noindent
Inversement, étant donné un faisceau de $\mathcal{O}_1$-modules
$\mathcal{G}$, le préfaisceau de $\mathcal{O}_2$-modules
$\mathcal{O}_2 \otimes_{p, \mathcal{O}_1} \mathcal{G}$
ne constitue pas en général un faisceau. Nous définissons donc le
{\it faisceau produit tensoriel}
$\mathcal{O}_2 \otimes_{\mathcal{O}_1} \mathcal{G}$
par la formule
$$
\mathcal{O}_2 \otimes_{\mathcal{O}_1} \mathcal{G}
=
(\mathcal{O}_2 \otimes_{p, \mathcal{O}_1} \mathcal{G})^\#
$$
comme le faisceau associé à la construction faite pour les préfaisceaux.
On obtient ainsi le foncteur {\it d'extension des scalaires}
$$
\textit{Mod}(\mathcal{O}_1)
\longrightarrow
\textit{Mod}(\mathcal{O}_2)
$$

\begin{lemma}
\label{sheaves-lemma-adjointness-tensor-restrict}
Avec $X$, $\mathcal{O}_1$, $\mathcal{O}_2$, $\mathcal{F}$ et
$\mathcal{G}$ comme ci-dessus, on a une bijection canonique
$$
\Hom_{\mathcal{O}_1}(\mathcal{G}, \mathcal{F}_{\mathcal{O}_1})
=
\Hom_{\mathcal{O}_2}(
\mathcal{O}_2 \otimes_{\mathcal{O}_1} \mathcal{G},
\mathcal{F}
)
$$
Autrement dit, les foncteurs de restriction et d'extension des scalaires
sont adjoints l'un de l'autre.
\end{lemma}

\begin{proof}
Cela résulte du
Lemme \ref{sheaves-lemma-adjointness-tensor-restrict-presheaves}
et du fait que
$\Hom_{\mathcal{O}_2}(
\mathcal{O}_2 \otimes_{\mathcal{O}_1} \mathcal{G},
\mathcal{F}
)
=
\Hom_{\mathcal{O}_2}(
\mathcal{O}_2 \otimes_{p, \mathcal{O}_1} \mathcal{G},
\mathcal{F}
)$
puisque $\mathcal{F}$ est un faisceau.
\end{proof}

\begin{lemma}
\label{sheaves-lemma-stalk-tensor-sheaf-modules}
Soit $X$ un espace topologique.
Soit $\mathcal{O} \to \mathcal{O}'$ un morphisme de
faisceaux d'anneaux sur $X$.
Soit $\mathcal{F}$ un faisceau de $\mathcal{O}$-modules.
Soit $x \in X$. On a
$$
\mathcal{F}_x \otimes_{\mathcal{O}_x} \mathcal{O}'_x
=
(\mathcal{F} \otimes_\mathcal{O} \mathcal{O}')_x
$$
en tant que $\mathcal{O}'_x$-modules.
\end{lemma}

\begin{proof}
Cela résulte directement du Lemme \ref{sheaves-lemma-stalk-tensor-presheaf-modules}
et du fait que le passage aux fibres commute au passage au faisceau associé.
\end{proof}













\section{Applications continues et faisceaux}
\label{sheaves-section-presheaves-functorial}

\noindent
Soit $f : X \to Y$ une application continue d'espaces topologiques.
Nous allons définir les foncteurs image directe et image inverse pour les
préfaisceaux et les faisceaux.

\medskip\noindent
Soit $\mathcal{F}$ un préfaisceau d'ensembles sur $X$. Nous définissons
l'{\it image directe} de $\mathcal{F}$ par la règle
$$
f_*\mathcal{F}(V) = \mathcal{F}(f^{-1}(V))
$$
pour tout ouvert $V \subset Y$.
Étant donnés des ouverts $V_1 \subset V_2 \subset Y$, l'application de restriction
est définie par la commutativité du diagramme
$$
\xymatrix{
f_*\mathcal{F}(V_2) \ar[d] \ar@{=}[r] &
\mathcal{F}(f^{-1}(V_2)) \ar[d]^{\text{restriction de }\mathcal{F}} \\
f_*\mathcal{F}(V_1) \ar@{=}[r] &
\mathcal{F}(f^{-1}(V_1))
}
$$
Il est clair que ceci définit un préfaisceau d'ensembles. La construction
est manifestement fonctorielle en $\mathcal{F}$ ; nous
obtenons donc un foncteur
$$
f_* : \textit{PSh}(X) \longrightarrow \textit{PSh}(Y).
$$

\begin{lemma}
\label{sheaves-lemma-pushforward-sheaf}
Soit $f : X \to Y$ une application continue.
Soit $\mathcal{F}$ un faisceau d'ensembles sur $X$.
Alors $f_*\mathcal{F}$ est un faisceau sur $Y$.
\end{lemma}

\begin{proof}
Cela résulte immédiatement du fait que,
si $V = \bigcup V_j$ est un recouvrement ouvert de $Y$,
alors $f^{-1}(V) = \bigcup f^{-1}(V_j)$ est un recouvrement ouvert de $X$.
\end{proof}

\noindent
Nous obtenons par conséquent un foncteur
$$
f_* : \Sh(X) \longrightarrow \Sh(Y).
$$
Celui-ci est compatible avec la composition au sens
fort suivant.

\begin{lemma}
\label{sheaves-lemma-pushforward-composition}
Soient $f : X \to Y$ et $g : Y \to Z$ des applications continues
d'espaces topologiques. Les foncteurs $(g \circ f)_*$
et $g_* \circ f_*$ sont égaux (tant sur les préfaisceaux
que sur les faisceaux d'ensembles).
\end{lemma}

\begin{proof}
Cela vient de ce que $(g \circ f)_*\mathcal{F}(W) =
\mathcal{F}((g \circ f)^{-1}W)$ et
$(g_* \circ f_*)\mathcal{F}(W) = \mathcal{F}(f^{-1} g^{-1} W)$,
et que $(g \circ f)^{-1}W = f^{-1} g^{-1} W$.
\end{proof}

\noindent
Soit $\mathcal{G}$ un préfaisceau d'ensembles sur $Y$.
Le {\it préfaisceau image inverse} $f_p\mathcal{G}$
d'un préfaisceau donné $\mathcal{G}$ est défini comme l'adjoint à gauche
de l'image directe $f_*$ sur les préfaisceaux. Autrement dit, ce
devrait être un préfaisceau $f_p \mathcal{G}$ sur $X$ tel que
$$
\Mor_{\textit{PSh}(X)}(f_p\mathcal{G}, \mathcal{F})
=
\Mor_{\textit{PSh}(Y)}(\mathcal{G}, f_*\mathcal{F}).
$$
Par le lemme de Yoneda, cette propriété détermine l'image inverse de manière unique.
Il se trouve qu'elle existe effectivement.

\begin{lemma}
\label{sheaves-lemma-pullback-presheaves}
Soit $f : X \to Y$ une application continue.
Il existe un foncteur
$f_p : \textit{PSh}(Y) \to \textit{PSh}(X)$
qui est adjoint à gauche de $f_*$. Pour un préfaisceau
$\mathcal{G}$, il est déterminé par la règle
$$
f_p\mathcal{G}(U) = \colim_{f(U) \subset V} \mathcal{G}(V)
$$
où la limite inductive porte sur la collection des voisinages ouverts
$V$ de $f(U)$ dans $Y$. Les limites inductives sont prises sur des
ensembles partiellement ordonnés filtrants.
(Les applications de restriction de $f_p\mathcal{G}$ sont expliquées dans la preuve.)
\end{lemma}

\begin{proof}
La limite inductive porte sur l'ensemble partiellement ordonné formé des ouverts
$V \subset Y$ qui contiennent $f(U)$, ordonné par l'inclusion opposée.
Cet ensemble partiellement ordonné est filtrant puisque, si $V, V'$ lui appartiennent,
il en va de même de $V \cap V'$. En outre, si
$U_1 \subset U_2$, alors tout voisinage ouvert de $f(U_2)$
est un voisinage ouvert de $f(U_1)$. Le système qui définit
$f_p\mathcal{G}(U_2)$ est donc un sous-système de celui qui définit
$f_p\mathcal{G}(U_1)$, et nous obtenons une application de restriction (par
exemple en appliquant les généralités de Catégories,
lemme \ref{categories-lemma-functorial-colimit}).

\medskip\noindent
Remarquons que la construction de la limite inductive est manifestement fonctorielle
en $\mathcal{G}$, de même que les applications de restriction.
Nous avons donc défini $f_p$ comme un foncteur.

\medskip\noindent
Une petite remarque utile est qu'il existe
une application canonique $\mathcal{G}(U) \to f_p\mathcal{G}(f^{-1}(U))$,
car le système des voisinages ouverts
de $f(f^{-1}(U))$ contient l'élément $U$. Celle-ci est compatible
avec les applications de restriction. Autrement dit, il existe un
morphisme canonique $i_\mathcal{G} : \mathcal{G} \to f_* f_p \mathcal{G}$.

\medskip\noindent
Soit $\mathcal{F}$ un préfaisceau d'ensembles sur $X$.
Supposons que $\psi : f_p\mathcal{G} \to \mathcal{F}$
soit un morphisme de préfaisceaux d'ensembles. Le morphisme correspondant
$\mathcal{G} \to f_*\mathcal{F}$ est le morphisme
$f_*\psi \circ i_\mathcal{G} :
\mathcal{G} \to f_* f_p \mathcal{G} \to f_* \mathcal{F}$.

\medskip\noindent
Une autre petite remarque utile est qu'il existe un
morphisme canonique $c_\mathcal{F} : f_p f_* \mathcal{F} \to \mathcal{F}$.
Soit en effet $U \subset X$ un ouvert.
Pour tout voisinage ouvert $V \supset f(U)$ de $Y$,
il existe une application
$f_*\mathcal{F}(V) = \mathcal{F}(f^{-1}(V))\to \mathcal{F}(U)$,
à savoir l'application de restriction de $\mathcal{F}$. Celle-ci est compatible
avec les applications de restriction entre les valeurs de $\mathcal{F}$
sur les images par $f^{-1}$ des différents ouverts contenant $f(U)$. Nous obtenons donc
une application canonique $f_p f_* \mathcal{F}(U) \to \mathcal{F}(U)$.
Une autre vérification immédiate montre que ces applications sont compatibles
avec les applications de restriction et définissent un morphisme $c_\mathcal{F}$
de préfaisceaux d'ensembles.

\medskip\noindent
Supposons que $\varphi : \mathcal{G} \to f_*\mathcal{F}$
soit un morphisme de préfaisceaux d'ensembles. Considérons $f_p\varphi :
f_p \mathcal{G} \to f_p f_* \mathcal{F}$.
En composant ensuite avec $c_\mathcal{F}$, nous obtenons le morphisme cherché
$c_\mathcal{F} \circ f_p\varphi : f_p\mathcal{G} \to \mathcal{F}$.
Nous omettons la vérification que cette construction est inverse
de la construction dans l'autre sens donnée ci-dessus.
\end{proof}

\begin{lemma}
\label{sheaves-lemma-stalk-pullback-presheaf}
Soit $f : X \to Y$ une application continue.
Soit $x \in X$. Soit $\mathcal{G}$ un préfaisceau d'ensembles sur $Y$.
Il existe une bijection canonique entre les fibres
$(f_p\mathcal{G})_x = \mathcal{G}_{f(x)}$.
\end{lemma}

\begin{proof}
On peut le voir comme suit :
\begin{eqnarray*}
(f_p\mathcal{G})_x
& = &
\colim_{x \in U} f_p\mathcal{G}(U) \\
& = &
\colim_{x \in U} \colim_{f(U) \subset V} \mathcal{G}(V) \\
& = &
\colim_{f(x) \in V} \mathcal{G}(V) \\
& = &
\mathcal{G}_{f(x)}
\end{eqnarray*}
Nous avons utilisé ici
Catégories, lemme \ref{categories-lemma-colimits-commute},
ainsi que le fait que tout ouvert $V$ de $Y$ contenant $f(x)$
intervient dans la troisième description ci-dessus. Nous omettons les détails.
\end{proof}

\noindent
Soit $\mathcal{G}$ un faisceau d'ensembles sur $Y$.
Le {\it faisceau image inverse} $f^{-1}\mathcal{G}$ est défini
par la formule
$$
f^{-1}\mathcal{G} = (f_p\mathcal{G})^\# .
$$
Le foncteur image inverse $f^{-1}$ est adjoint à gauche du
foncteur image directe sur les faisceaux. Autrement dit,
$$
\Mor_{\Sh(X)}(f^{-1}\mathcal{G}, \mathcal{F})
=
\Mor_{\Sh(Y)}(\mathcal{G}, f_*\mathcal{F}).
$$
En effet, nous avons
\begin{eqnarray*}
\Mor_{\Sh(X)}(f^{-1}\mathcal{G}, \mathcal{F})
& = &
\Mor_{\textit{PSh}(X)}(f_p\mathcal{G}, \mathcal{F}) \\
& = &
\Mor_{\textit{PSh}(Y)}(\mathcal{G}, f_*\mathcal{F}) \\
& = &
\Mor_{\Sh(Y)}(\mathcal{G}, f_*\mathcal{F})
\end{eqnarray*}
Pour la première égalité, nous utilisons que le foncteur « faisceau associé » est l'adjoint à gauche
du foncteur d'inclusion des faisceaux dans les préfaisceaux. Pour la deuxième égalité,
nous utilisons que $f_p$ est adjoint à gauche de $f_*$ sur les préfaisceaux. Nous reviendrons
sur cet énoncé dans la preuve du lemme \ref{sheaves-lemma-f-map}.

\begin{lemma}
\label{sheaves-lemma-stalk-pullback}
Soit $x \in X$. Soit $\mathcal{G}$ un faisceau d'ensembles sur $Y$.
Il existe une bijection canonique entre les fibres
$(f^{-1}\mathcal{G})_x = \mathcal{G}_{f(x)}$.
\end{lemma}

\begin{proof}
C'est une combinaison des lemmes \ref{sheaves-lemma-stalk-sheafification}
et \ref{sheaves-lemma-stalk-pullback-presheaf}.
\end{proof}

\begin{lemma}
\label{sheaves-lemma-pullback-composition}
Soient $f : X \to Y$ et $g : Y \to Z$ des applications continues
d'espaces topologiques. Les foncteurs $(g \circ f)^{-1}$
et $f^{-1} \circ g^{-1}$ sont canoniquement isomorphes.
De même, $(g \circ f)_p \cong f_p \circ g_p$ sur les
préfaisceaux.
\end{lemma}

\begin{proof}
Pour le voir,
utilisons le fait que les foncteurs adjoints sont uniques à isomorphisme unique près,
ainsi que le lemme \ref{sheaves-lemma-pushforward-composition}.
\end{proof}

\begin{definition}
\label{sheaves-definition-f-map}
Soit $f : X \to Y$ une application continue.
Soit $\mathcal{F}$ un faisceau d'ensembles sur $X$ et
soit $\mathcal{G}$ un faisceau d'ensembles sur $Y$.
Un {\it $f$-morphisme $\xi : \mathcal{G} \to \mathcal{F}$}
est une famille d'applications
$\xi_V : \mathcal{G}(V) \to \mathcal{F}(f^{-1}(V))$
indexée par les ouverts $V \subset Y$ telle que
$$
\xymatrix{
\mathcal{G}(V) \ar[r]_{\xi_V} \ar[d]_{\text{restriction de }\mathcal{G}} &
\mathcal{F}(f^{-1}V) \ar[d]^{\text{restriction de }\mathcal{F}} \\
\mathcal{G}(V') \ar[r]^{\xi_{V'}} &
\mathcal{F}(f^{-1}V')
}
$$
soit commutatif pour tous ouverts $V' \subset V \subset Y$.
\end{definition}

\noindent
Dans la littérature, on rencontre parfois une autre définition,
comme dans la partie (4) du lemme \ref{sheaves-lemma-f-map}, mais le lemme montre qu'il
n'y a en réalité aucune différence.

\begin{lemma}
\label{sheaves-lemma-f-map}
Soit $f : X \to Y$ une application continue.
Il existe des bijections entre les quatre ensembles suivants :
\begin{enumerate}
\item l'ensemble des morphismes $\mathcal{G} \to f_*\mathcal{F}$,
\item l'ensemble des morphismes $f^{-1}\mathcal{G} \to \mathcal{F}$,
\item l'ensemble des $f$-morphismes $\xi : \mathcal{G} \to \mathcal{F}$, et
\item l'ensemble de toutes les familles d'applications
$\xi_{U, V} : \mathcal{G}(V) \to \mathcal{F}(U)$ pour tous
ouverts $U \subset X$ et $V \subset Y$ tels que $f(U) \subset V$,
compatibles avec toutes les applications de restriction,
\end{enumerate}
fonctoriellement en $\mathcal{F} \in \Sh(X)$ et $\mathcal{G} \in \Sh(Y)$.
\end{lemma}

\begin{proof}
Un morphisme de faisceaux $a : \mathcal{G} \to f_*\mathcal{F}$
est par définition une règle qui associe à chaque ouvert $V$ de $Y$
une application $a_V : \mathcal{G}(V) \to f_*\mathcal{F}(V)$,
et nous avons $f_*\mathcal{F}(V) = \mathcal{F}(f^{-1}(V))$.
Ainsi, pour ce qui est des « données », on obtient exactement ce qu'il faut
pour un $f$-morphisme $\xi$ de $\mathcal{G}$ vers $\mathcal{F}$.
Pour montrer que les ensembles (1) et (3) sont en bijection, nous observons que
$a$ est un morphisme de faisceaux si et seulement si la famille d'applications $a_V$
correspondante satisfait la condition de la définition \ref{sheaves-definition-f-map}.

\medskip\noindent
Rappelons que $f^{-1}\mathcal{G}$ est le faisceau associé au préfaisceau $f_p\mathcal{G}$.
Par la propriété universelle du faisceau associé, un morphisme de faisceaux
$b : f^{-1}\mathcal{G} \to \mathcal{F}$ équivaut à un morphisme de
préfaisceaux $b_p : f_p\mathcal{G} \to \mathcal{F}$, où $f_p$ est le
foncteur défini plus haut dans cette section. Pour donner un tel morphisme $b_p$,
il faut spécifier, pour chaque ouvert $U$ de $X$, une application
$$
b_{p, U} :
\colim_{f(U) \subset V} \mathcal{G}(V)
\longrightarrow
\mathcal{F}(U)
$$
compatible avec les applications de restriction. Nous regardons $b_{p, U}$
comme une famille d'applications $b_{p, U, V} : \mathcal{G}(V) \to \mathcal{F}(U)$
pour tout ouvert $V$ de $Y$ tel que $f(U) \subset V$. Ces applications doivent être
compatibles avec toutes les applications de restriction possibles.
Autrement dit, nous voyons que $b_p$ correspond à une famille d'applications
comme en (4). Réciproquement, une telle famille définit bien sûr un morphisme
$b_p$, puis un morphisme $b : f^{-1}\mathcal{G} \to \mathcal{F}$.

\medskip\noindent
Pour achever la preuve du lemme, il faut montrer que, par
« oubli de structure », la règle qui, à une famille $\xi_{U, V}$
comme en (4), associe le $f$-morphisme $\xi$ défini par $\xi_V = \xi_{f^{-1}(V), V}$
est bijective. Pour cela, si $\xi$ est un $f$-morphisme usuel, nous
définissons simplement $\tilde \xi_{U, V}$ comme la composée de
$\xi_V : \mathcal{G}(V) \to \mathcal{F}(f^{-1}(V))$
avec l'application de restriction $\mathcal{F}(f^{-1}(V)) \to \mathcal{F}(U)$,
qui a un sens précisément parce que $f(U) \subset V$, c'est-à-dire
$U \subset f^{-1}(V)$. Ceci achève la preuve.
\end{proof}

\noindent
Il est parfois commode de considérer les $f$-morphismes
plutôt que les morphismes entre faisceaux, soit sur $X$, soit sur $Y$.
Nous définissons comme suit la composition des $f$-morphismes.

\begin{definition}
\label{sheaves-definition-composition-f-maps}
Supposons que $f : X \to Y$ et $g : Y \to Z$ soient des applications
continues d'espaces topologiques. Supposons que $\mathcal{F}$ soit
un faisceau sur $X$, que $\mathcal{G}$ soit un faisceau sur $Y$ et que
$\mathcal{H}$ soit un faisceau sur $Z$.
Soit $\varphi : \mathcal{G} \to \mathcal{F}$ un $f$-morphisme.
Soit $\psi : \mathcal{H} \to \mathcal{G}$ un $g$-morphisme.
La {\it composée de $\varphi$ et $\psi$} est le
$(g \circ f)$-morphisme $\varphi \circ \psi$ défini
par la commutativité du diagramme
$$
\xymatrix{
\mathcal{H}(W) \ar[rr]_{(\varphi \circ \psi)_W}
\ar[rd]_{\psi_W} & &
\mathcal{F}(f^{-1}g^{-1}W) \\
&
\mathcal{G}(g^{-1}W)
\ar[ru]_{\varphi_{g^{-1}W}}
}
$$
\end{definition}

\noindent
Nous laissons au lecteur le soin de vérifier que cela fonctionne.
Une autre manière de le voir consiste à considérer
$\varphi \circ \psi$ comme la composée
$$
\mathcal{H}
\xrightarrow{\psi}
g_*\mathcal{G}
\xrightarrow{g_*\varphi}
g_* f_* \mathcal{F} = (g \circ f)_* \mathcal{F}
$$
N'avez-vous pas maintenant l'impression que considérer les $f$-morphismes est
en quelque sorte plus simple ?

\medskip\noindent
Enfin, étant donnés une application continue $f : X \to Y$ et un
$f$-morphisme $\varphi : \mathcal{G} \to \mathcal{F}$, il existe
une application naturelle entre les fibres
$$
\varphi_x : \mathcal{G}_{f(x)} \longrightarrow \mathcal{F}_x
$$
pour tout $x \in X$. Un représentant $(V, s)$
d'un élément de $\mathcal{G}_{f(x)}$ est envoyé sur
l'élément de $\mathcal{F}_x$ ayant pour représentant $(f^{-1}V,
\varphi_V(s))$. Nous laissons au lecteur le soin de vérifier que cette application
est bien définie. Une autre formulation consiste à dire qu'il s'agit de
l'unique application pour laquelle tous les diagrammes
$$
\xymatrix{
\mathcal{F}(f^{-1}V) \ar[r] &
\mathcal{F}_x \\
\mathcal{G}(V) \ar[r] \ar[u]^{\varphi_V} &
\mathcal{G}_{f(x)} \ar[u]^{\varphi_x}
}
$$
(pour tout ouvert $f(x) \in V \subset Y$) sont commutatifs.

\begin{lemma}
\label{sheaves-lemma-compose-f-maps-stalks}
Supposons que $f : X \to Y$ et $g : Y \to Z$ soient des applications
continues d'espaces topologiques. Supposons que $\mathcal{F}$ soit
un faisceau sur $X$, que $\mathcal{G}$ soit un faisceau sur $Y$ et que
$\mathcal{H}$ soit un faisceau sur $Z$.
Soit $\varphi : \mathcal{G} \to \mathcal{F}$ un $f$-morphisme.
Soit $\psi : \mathcal{H} \to \mathcal{G}$ un $g$-morphisme.
Soit $x \in X$ un point. L'application induite sur les fibres
$(\varphi \circ \psi)_x : \mathcal{H}_{g(f(x))}
\to \mathcal{F}_x$ est la composée
$$
\mathcal{H}_{g(f(x))}
\xrightarrow{\psi_{f(x)}}
\mathcal{G}_{f(x)}
\xrightarrow{\varphi_x}
\mathcal{F}_x
$$
\end{lemma}

\begin{proof}
Cela résulte immédiatement de la définition \ref{sheaves-definition-composition-f-maps}
et de la définition de l'application sur les fibres ci-dessus.
\end{proof}
\section{Applications continues et faisceaux abéliens}
\label{sheaves-section-abelian-presheaves-functorial}

\noindent
Soit $f : X \to Y$ une application continue.
Nous affirmons qu'il existe des foncteurs
\begin{eqnarray*}
f_* : \textit{PAb}(X) & \longrightarrow & \textit{PAb}(Y) \\
f_* : \textit{Ab}(X) & \longrightarrow & \textit{Ab}(Y) \\
f_p : \textit{PAb}(Y) & \longrightarrow & \textit{PAb}(X) \\
f^{-1} : \textit{Ab}(Y) & \longrightarrow & \textit{Ab}(X)
\end{eqnarray*}
qui possèdent des propriétés analogues à celles de leurs homologues de la
section \ref{sheaves-section-presheaves-functorial}.
Pour le voir, nous procédons comme suit.

\medskip\noindent
Chacun de ces foncteurs sera construit de la même
manière que le foncteur correspondant de la
section \ref{sheaves-section-presheaves-functorial}.
Cela fonctionne parce que toutes les limites inductives de cette section
sont filtrantes (mais nous allons donner les détails ci-dessous).

\medskip\noindent
Tout d'abord, étant donnés un préfaisceau abélien $\mathcal{F}$ sur $X$ et
un préfaisceau abélien $\mathcal{G}$ sur $Y$, nous définissons
\begin{eqnarray*}
f_*\mathcal{F}(V) & = & \mathcal{F}(f^{-1}(V)) \\
f_p\mathcal{G}(U) & = & \colim_{f(U) \subset V} \mathcal{G}(V)
\end{eqnarray*}
comme groupes abéliens. Les applications de restriction sont les mêmes que
pour les préfaisceaux d'ensembles (et ce sont
toutes des homomorphismes de groupes abéliens).

\medskip\noindent
Les affectations $\mathcal{F} \mapsto f_*\mathcal{F}$ et
$\mathcal{G} \to f_p\mathcal{G}$ sont des foncteurs sur
les catégories de préfaisceaux de groupes abéliens.
C'est clair : par exemple, un morphisme de préfaisceaux abéliens
$\mathcal{G}_1 \to \mathcal{G}_2$ induit un morphisme de
systèmes filtrants
$\{\mathcal{G}_1(V)\}_{f(U) \subset V} \to
\{\mathcal{G}_2(V)\}_{f(U) \subset V}$
dont toutes les composantes sont des homomorphismes ;
il induit donc un homomorphisme de groupes abéliens
$f_p\mathcal{G}_1(U) \to f_p\mathcal{G}_2(U)$.

\medskip\noindent
Les foncteurs $f_*$ et $f_p$ sont adjoints sur la
catégorie des préfaisceaux de groupes abéliens, c'est-à-dire que
$$
\Mor_{\textit{PAb}(X)}(f_p\mathcal{G}, \mathcal{F})
=
\Mor_{\textit{PAb}(Y)}(\mathcal{G}, f_*\mathcal{F}).
$$
Pour le démontrer, remarquons que le morphisme
$i_\mathcal{G} : \mathcal{G} \to f_* f_p\mathcal{G}$ de la preuve
du Lemme \ref{sheaves-lemma-pullback-presheaves}
est un morphisme de préfaisceaux abéliens. Ainsi, si
$\psi : f_p\mathcal{G} \to \mathcal{F}$
est un morphisme de préfaisceaux abéliens, alors le morphisme correspondant
$\mathcal{G} \to f_*\mathcal{F}$, à savoir
$f_*\psi \circ i_\mathcal{G} :
\mathcal{G} \to f_* f_p \mathcal{G} \to f_* \mathcal{F}$,
est lui aussi un morphisme de préfaisceaux abéliens. Dans l'autre sens,
remarquons que le morphisme
$c_\mathcal{F} : f_p f_* \mathcal{F} \to \mathcal{F}$
de la preuve du Lemme \ref{sheaves-lemma-pullback-presheaves} est également un
morphisme de préfaisceaux abéliens (car il est construit à partir des applications
de restriction de $\mathcal{F}$, qui sont toutes des homomorphismes). Ainsi,
étant donné un morphisme de préfaisceaux abéliens $\varphi : \mathcal{G} \to f_*\mathcal{F}$,
le morphisme
$c_\mathcal{F} \circ f_p\varphi : f_p\mathcal{G} \to \mathcal{F}$
est lui aussi un morphisme de préfaisceaux abéliens. Puisque ces constructions
$\psi \mapsto f_*\psi$ et $\varphi \mapsto c_\mathcal{F} \circ f_p\varphi$
sont inverses l'une de l'autre sur les morphismes de préfaisceaux d'ensembles,
elles le sont aussi sur les morphismes de préfaisceaux abéliens.

\medskip\noindent
Si $\mathcal{F}$ est un faisceau abélien sur $Y$, alors $f_*\mathcal{F}$
est un faisceau abélien sur $X$. Cela résulte de la définition
d'un faisceau abélien et du résultat correspondant pour les faisceaux d'ensembles ;
voir le Lemme \ref{sheaves-lemma-pushforward-sheaf}. On obtient ainsi le foncteur
$f_*$ sur la catégorie des faisceaux abéliens.

\medskip\noindent
Nous définissons $f^{-1}\mathcal{G} = (f_p\mathcal{G})^\#$ comme précédemment.
L'adjonction de $f_*$ et $f^{-1}$ s'obtient formellement comme dans
le cas des préfaisceaux d'ensembles. Voici l'argument :
\begin{eqnarray*}
\Mor_{\textit{Ab}(X)}(f^{-1}\mathcal{G}, \mathcal{F})
& = &
\Mor_{\textit{PAb}(X)}(f_p\mathcal{G}, \mathcal{F}) \\
& = &
\Mor_{\textit{PAb}(Y)}(\mathcal{G}, f_*\mathcal{F}) \\
& = &
\Mor_{\textit{Ab}(Y)}(\mathcal{G}, f_*\mathcal{F})
\end{eqnarray*}

\begin{lemma}
\label{sheaves-lemma-pullback-abelian-stalk}
Soit $f : X \to Y$ une application continue.
\begin{enumerate}
\item Soit $\mathcal{G}$ un préfaisceau abélien sur $Y$.
Soit $x \in X$. La bijection
$\mathcal{G}_{f(x)} \to (f_p\mathcal{G})_x$ du
Lemme \ref{sheaves-lemma-stalk-pullback-presheaf} est un isomorphisme de groupes abéliens.
\item Soit $\mathcal{G}$ un faisceau abélien sur $Y$.
Soit $x \in X$. La bijection
$\mathcal{G}_{f(x)} \to (f^{-1}\mathcal{G})_x$ du
Lemme \ref{sheaves-lemma-stalk-pullback} est un isomorphisme de groupes abéliens.
\end{enumerate}
\end{lemma}

\begin{proof}
Omis.
\end{proof}

\noindent
Étant donnés une application continue $f : X \to Y$ et des faisceaux de groupes
abéliens $\mathcal{F}$ sur $X$, $\mathcal{G}$ sur $Y$, la notion
d'{\it $f$-morphisme $\mathcal{G} \to \mathcal{F}$
de faisceaux de groupes abéliens} a un sens. Nous pouvons la définir
exactement comme dans la Définition \ref{sheaves-definition-f-map} (en remplaçant les applications
d'ensembles par des homomorphismes de groupes abéliens), ou
dire simplement qu'il s'agit d'un morphisme de faisceaux abéliens
$\mathcal{G} \to f_*\mathcal{F}$. Nous emploierons librement cette
notion dans la suite. Le groupe des $f$-morphismes entre
$\mathcal{G}$ et $\mathcal{F}$ sera canoniquement en bijection
avec les groupes
$\Mor_{\textit{Ab}(X)}(f^{-1}\mathcal{G}, \mathcal{F})$
et
$\Mor_{\textit{Ab}(Y)}(\mathcal{G}, f_*\mathcal{F})$.

\medskip\noindent
La composition des $f$-morphismes se définit exactement de la
même manière que pour les $f$-morphismes de faisceaux
d'ensembles. De plus, étant donné un $f$-morphisme $\mathcal{G} \to \mathcal{F}$
comme ci-dessus, les applications induites sur les fibres
$$
\varphi_x : \mathcal{G}_{f(x)} \longrightarrow \mathcal{F}_x
$$
sont des homomorphismes de groupes abéliens.


\section{Applications continues et faisceaux de structures algébriques}
\label{sheaves-section-presheaves-structures-functorial}

\noindent
Soit $(\mathcal{C}, F)$ une espèce de structure algébrique.
Pour un espace topologique $X$, introduisons les notations suivantes :
\begin{enumerate}
\item $\textit{PSh}(X, \mathcal{C})$ désignera la catégorie
des préfaisceaux à valeurs dans $\mathcal{C}$.
\item $\Sh(X, \mathcal{C})$ désignera la catégorie
des faisceaux à valeurs dans $\mathcal{C}$.
\end{enumerate}
Soit $f : X \to Y$ une application continue d'espaces topologiques.
Les mêmes arguments que dans la section précédente montrent
qu'il existe des foncteurs
\begin{eqnarray*}
f_* : \textit{PSh}(X, \mathcal{C}) &
\longrightarrow &
\textit{PSh}(Y, \mathcal{C}) \\
f_* : \Sh(X, \mathcal{C}) &
\longrightarrow &
\Sh(Y, \mathcal{C}) \\
f_p : \textit{PSh}(Y, \mathcal{C}) &
\longrightarrow &
\textit{PSh}(X, \mathcal{C}) \\
f^{-1} : \Sh(Y, \mathcal{C}) &
\longrightarrow &
\Sh(X, \mathcal{C})
\end{eqnarray*}
construits de la même manière et possédant les mêmes
propriétés que les foncteurs construits pour les
(pré)faisceaux abéliens. En particulier, on a des diagrammes commutatifs
$$
\xymatrix{
\textit{PSh}(X, \mathcal{C}) \ar[r]^{f_*} \ar[d]^F &
\textit{PSh}(Y, \mathcal{C}) \ar[d]^F &
\Sh(X, \mathcal{C}) \ar[r]^{f_*} \ar[d]^F &
\Sh(Y, \mathcal{C}) \ar[d]^F
\\
\textit{PSh}(X) \ar[r]^{f_*} &
\textit{PSh}(Y) &
\Sh(X) \ar[r]^{f_*} &
\Sh(Y)
\\
\textit{PSh}(Y, \mathcal{C}) \ar[r]^{f_p} \ar[d]^F &
\textit{PSh}(X, \mathcal{C}) \ar[d]^F &
\Sh(Y, \mathcal{C}) \ar[r]^{f^{-1}} \ar[d]^F &
\Sh(X, \mathcal{C}) \ar[d]^F
\\
\textit{PSh}(Y) \ar[r]^{f_p} &
\textit{PSh}(X) &
\Sh(Y) \ar[r]^{f^{-1}} &
\Sh(X)
}
$$

\medskip\noindent
Les principales formules à retenir sont les suivantes :
\begin{eqnarray*}
f_*\mathcal{F}(V) & = & \mathcal{F}(f^{-1}(V)) \\
f_p\mathcal{G}(U) & = & \colim_{f(U) \subset V} \mathcal{G}(V) \\
f^{-1}\mathcal{G} & = & (f_p\mathcal{G})^\# \\
(f_p\mathcal{G})_x & = & \mathcal{G}_{f(x)} \\
(f^{-1}\mathcal{G})_x & = & \mathcal{G}_{f(x)}
\end{eqnarray*}
Chacune de ces formules est valable dans la
catégorie $\mathcal{C}$ et redonne, après passage aux ensembles
sous-jacents, la formule correspondante pour les préfaisceaux d'ensembles.
Nous avons en outre les propriétés d'adjonction
\begin{eqnarray*}
\Mor_{\textit{PSh}(X, \mathcal{C})}(f_p\mathcal{G}, \mathcal{F})
& = &
\Mor_{\textit{PSh}(Y, \mathcal{C})}(\mathcal{G}, f_*\mathcal{F}) \\
\Mor_{\Sh(X, \mathcal{C})}(f^{-1}\mathcal{G}, \mathcal{F})
& = &
\Mor_{\Sh(Y, \mathcal{C})}(\mathcal{G}, f_*\mathcal{F}).
\end{eqnarray*}
Pour les démontrer, l'étape principale consiste à construire les morphismes
$$
i_\mathcal{G} : \mathcal{G} \longrightarrow f_*f_p\mathcal{G}
$$
et
$$
c_\mathcal{F} : f_p f_* \mathcal{F} \longrightarrow \mathcal{F}
$$
qui interviennent dans la preuve du Lemme \ref{sheaves-lemma-pullback-presheaves} comme
morphismes de préfaisceaux à valeurs dans $\mathcal{C}$. Nous pouvons
laisser cette vérification au lecteur, car les constructions sont exactement les mêmes
que pour les préfaisceaux d'ensembles.

\medskip\noindent
Étant donnés une application continue $f : X \to Y$ et des faisceaux de structures
algébriques $\mathcal{F}$ sur $X$, $\mathcal{G}$ sur $Y$, la notion
d'{\it $f$-morphisme $\mathcal{G} \to \mathcal{F}$
de faisceaux de structures algébriques} a un sens. Nous pouvons la définir
exactement comme dans la Définition \ref{sheaves-definition-f-map} (en remplaçant les applications
d'ensembles par des morphismes de $\mathcal{C}$), ou
dire simplement qu'il s'agit d'un morphisme de faisceaux de structures
algébriques $\mathcal{G} \to f_*\mathcal{F}$. Nous emploierons librement cette
notion dans la suite. L'ensemble des $f$-morphismes entre
$\mathcal{G}$ et $\mathcal{F}$ sera canoniquement en bijection
avec les ensembles
$\Mor_{\Sh(X, \mathcal{C})}(f^{-1}\mathcal{G}, \mathcal{F})$
et
$\Mor_{\Sh(Y, \mathcal{C})}(\mathcal{G}, f_*\mathcal{F})$.

\medskip\noindent
La composition des $f$-morphismes se définit exactement de la
même manière que pour les $f$-morphismes de faisceaux
d'ensembles. De plus, étant donné un $f$-morphisme $\mathcal{G} \to \mathcal{F}$
comme ci-dessus, les applications induites sur les fibres
$$
\varphi_x : \mathcal{G}_{f(x)} \longrightarrow \mathcal{F}_x
$$
sont des homomorphismes de structures algébriques.


\begin{lemma}
\label{sheaves-lemma-f-map-sets-algebraic-structures}
Soit $f : X \to Y$ une application continue d'espaces topologiques.
Supposons donnés des faisceaux de structures algébriques
$\mathcal{F}$ sur $X$, $\mathcal{G}$ sur $Y$. Soit
$\varphi : \mathcal{G} \to \mathcal{F}$ un $f$-morphisme
des faisceaux d'ensembles sous-jacents. Si, pour tout ouvert $V \subset Y$,
l'application d'ensembles $\varphi_V : \mathcal{G}(V) \to \mathcal{F}(f^{-1}V)$
provient, sur les ensembles sous-jacents, d'un morphisme de $\mathcal{C}$,
alors $\varphi$ provient d'un unique $f$-morphisme entre
les faisceaux de structures algébriques.
\end{lemma}

\begin{proof}
Omis.
\end{proof}




\section{Applications continues et faisceaux de modules}
\label{sheaves-section-presheaves-modules-functorial}

\noindent
Le cas des faisceaux de modules est plus compliqué.
En effet, le cadre naturel pour définir
les foncteurs image inverse et image directe est celui
des espaces annelés, que nous définirons plus bas. Nous
énonçons d'abord quelques lemmes évidents.

\begin{lemma}
\label{sheaves-lemma-pushforward-presheaf-module}
Soit $f : X \to Y$ une application continue d'espaces topologiques.
Soit $\mathcal{O}$ un préfaisceau d'anneaux sur $X$. Soit
$\mathcal{F}$ un préfaisceau de $\mathcal{O}$-modules.
Il existe un morphisme naturel de préfaisceaux d'ensembles sous-jacents
$$
f_*\mathcal{O} \times f_*\mathcal{F}
\longrightarrow
f_*\mathcal{F}
$$
qui munit $f_*\mathcal{F}$ d'une structure de préfaisceau de
$f_*\mathcal{O}$-modules. Cette construction est
fonctorielle en $\mathcal{F}$.
\end{lemma}

\begin{proof}
Soit $V \subset Y$ un ouvert. Nous définissons l'application du lemme
comme l'application
$$
f_*\mathcal{O}(V) \times f_*\mathcal{F}(V)
=
\mathcal{O}(f^{-1}V) \times \mathcal{F}(f^{-1}V)
\to
\mathcal{F}(f^{-1}V)
=
f_*\mathcal{F}(V).
$$
Ici, la flèche du milieu est l'application de multiplication sur $X$.
Nous laissons au lecteur le soin de vérifier qu'elle est compatible avec les
applications de restriction et définit une structure de
$f_*\mathcal{O}$-module sur $f_*\mathcal{F}$.
\end{proof}

\begin{lemma}
\label{sheaves-lemma-pullback-presheaf-module}
Soit $f : X \to Y$ une application continue d'espaces topologiques.
Soit $\mathcal{O}$ un préfaisceau d'anneaux sur $Y$. Soit
$\mathcal{G}$ un préfaisceau de $\mathcal{O}$-modules.
Il existe un morphisme naturel de préfaisceaux d'ensembles sous-jacents
$$
f_p\mathcal{O} \times f_p\mathcal{G}
\longrightarrow
f_p\mathcal{G}
$$
qui munit $f_p\mathcal{G}$ d'une structure de préfaisceau de $f_p\mathcal{O}$-modules.
Cette construction est fonctorielle en $\mathcal{G}$.
\end{lemma}

\begin{proof}
Soit $U \subset X$ un ouvert. Nous définissons l'application du lemme
comme l'application
\begin{eqnarray*}
f_p\mathcal{O}(U) \times f_p\mathcal{G}(U)
& = &
\colim_{f(U) \subset V} \mathcal{O}(V)
\times
\colim_{f(U) \subset V} \mathcal{G}(V) \\
& = &
\colim_{f(U) \subset V} (\mathcal{O}(V)\times \mathcal{G}(V)) \\
& \to &
\colim_{f(U) \subset V} \mathcal{G}(V) \\
& = &
f_p\mathcal{G}(U).
\end{eqnarray*}
Ici, la flèche du milieu est l'application de multiplication sur $Y$.
La seconde égalité vaut parce que les limites inductives filtrantes commutent
aux limites finies ; voir
Catégories, Lemme \ref{categories-lemma-directed-commutes}.
Nous laissons au lecteur le soin de vérifier que cette application est compatible avec les
applications de restriction et définit une structure de
$f_p\mathcal{O}$-module sur $f_p\mathcal{G}$.
\end{proof}

\noindent
Soit $f : X \to Y$ une application continue.
Soit $\mathcal{O}_X$ un préfaisceau d'anneaux sur $X$ et
soit $\mathcal{O}_Y$ un préfaisceau d'anneaux sur $Y$.
À ce stade, nous avons donc défini des foncteurs
\begin{eqnarray*}
f_* : \textit{PMod}(\mathcal{O}_X) &
\longrightarrow &
\textit{PMod}(f_*\mathcal{O}_X) \\
f_p : \textit{PMod}(\mathcal{O}_Y) &
\longrightarrow &
\textit{PMod}(f_p\mathcal{O}_Y)
\end{eqnarray*}
Ils vérifient les compatibilités suivantes.

\begin{lemma}
\label{sheaves-lemma-adjoint-push-pull-presheaves-modules}
Soit $f : X \to Y$ une application continue d'espaces topologiques.
Soit $\mathcal{O}$ un préfaisceau d'anneaux sur $Y$.
Soit $\mathcal{G}$ un préfaisceau de $\mathcal{O}$-modules.
Soit $\mathcal{F}$ un préfaisceau de $f_p\mathcal{O}$-modules.
Alors
$$
\Mor_{\textit{PMod}(f_p\mathcal{O})}(f_p\mathcal{G}, \mathcal{F})
=
\Mor_{\textit{PMod}(\mathcal{O})}(\mathcal{G}, f_*\mathcal{F}).
$$
Nous utilisons ici
les Lemmes \ref{sheaves-lemma-pullback-presheaf-module}
et \ref{sheaves-lemma-pushforward-presheaf-module}, et considérons
$f_*\mathcal{F}$ comme un $\mathcal{O}$-module par restriction des scalaires le long de
$i_\mathcal{O} : \mathcal{O} \to f_*f_p\mathcal{O}$
(défini pour la première fois dans la preuve du Lemme \ref{sheaves-lemma-pullback-presheaves}).
\end{lemma}

\begin{proof}
Remarquons que
$$
\Mor_{\textit{PAb}(X)}(f_p\mathcal{G}, \mathcal{F})
=
\Mor_{\textit{PAb}(Y)}(\mathcal{G}, f_*\mathcal{F}).
$$
Cette égalité vaut d'après la section \ref{sheaves-section-abelian-presheaves-functorial}.
Il reste donc à montrer que, sous cette correspondance, les
sous-ensembles formés des morphismes de modules se correspondent. De plus, la correspondance
est déterminée par la règle
$$
(\psi : f_p\mathcal{G} \to \mathcal{F})
\longmapsto
(f_*\psi \circ i_\mathcal{G} :
\mathcal{G} \to f_* \mathcal{F})
$$
et, dans l'autre sens, par la règle
$$
(\varphi : \mathcal{G} \to f_* \mathcal{F})
\longmapsto
(c_\mathcal{F} \circ f_p\varphi : f_p\mathcal{G} \to \mathcal{F})
$$
où $i_\mathcal{G}$ et $c_\mathcal{F}$ sont ceux de la
section \ref{sheaves-section-abelian-presheaves-functorial}.
Ainsi, par fonctorialité de $f_*$ et $f_p$, nous voyons
qu'il suffit de vérifier que les morphismes
$i_\mathcal{G} : \mathcal{G} \to f_* f_p \mathcal{G}$ et
$c_\mathcal{F} : f_p f_* \mathcal{F} \to \mathcal{F}$
sont compatibles avec les structures de modules, ce que nous laissons au lecteur.
\end{proof}

\begin{lemma}
\label{sheaves-lemma-adjoint-pull-push-presheaves-modules}
Soit $f : X \to Y$ une application continue d'espaces topologiques.
Soit $\mathcal{O}$ un préfaisceau d'anneaux sur $X$.
Soit $\mathcal{F}$ un préfaisceau de $\mathcal{O}$-modules.
Soit $\mathcal{G}$ un préfaisceau de $f_*\mathcal{O}$-modules.
Alors
$$
\Mor_{\textit{PMod}(\mathcal{O})}(
\mathcal{O} \otimes_{p, f_pf_*\mathcal{O}} f_p\mathcal{G}, \mathcal{F})
=
\Mor_{\textit{PMod}(f_*\mathcal{O})}(\mathcal{G}, f_*\mathcal{F}).
$$
Nous utilisons ici
les Lemmes \ref{sheaves-lemma-pullback-presheaf-module}
et \ref{sheaves-lemma-pushforward-presheaf-module}, ainsi que
le morphisme $c_\mathcal{O} : f_pf_*\mathcal{O} \to \mathcal{O}$
dans la définition du produit tensoriel.
\end{lemma}

\begin{proof}
Cela résulte des égalités
\begin{eqnarray*}
\Mor_{\textit{PMod}(\mathcal{O})}(
\mathcal{O} \otimes_{p, f_pf_*\mathcal{O}} f_p\mathcal{G}, \mathcal{F})
& = &
\Mor_{\textit{PMod}(f_pf_*\mathcal{O})}(
f_p\mathcal{G}, \mathcal{F}_{f_pf_*\mathcal{O}}) \\
& = &
\Mor_{\textit{PMod}(f_*\mathcal{O})}(\mathcal{G},
f_*(\mathcal{F}_{f_pf_*\mathcal{O}})) \\
& = &
\Mor_{\textit{PMod}(f_*\mathcal{O})}(\mathcal{G}, f_*\mathcal{F}).
\end{eqnarray*}
La première égalité est le
Lemme \ref{sheaves-lemma-adjointness-tensor-restrict-presheaves}.
La deuxième égalité est le
Lemme \ref{sheaves-lemma-adjoint-push-pull-presheaves-modules}.
La troisième égalité est donnée par l'égalité
$f_*(\mathcal{F}_{f_pf_*\mathcal{O}}) = f_*\mathcal{F}$
de faisceaux abéliens qui est $f_*\mathcal{O}$-linéaire. En effet,
$\text{id}_{f_*\mathcal{O}}$ correspond à $c_\mathcal{O}$
par l'adjonction décrite dans la preuve du
Lemme \ref{sheaves-lemma-pullback-presheaves} ;
ainsi
$\text{id}_{f_*\mathcal{O}} = f_*c_\mathcal{O} \circ i_{f_*\mathcal{O}}$.
\end{proof}

\begin{lemma}
\label{sheaves-lemma-pushforward-module}
Soit $f : X \to Y$ une application continue d'espaces topologiques.
Soit $\mathcal{O}$ un faisceau d'anneaux sur $X$. Soit
$\mathcal{F}$ un faisceau de $\mathcal{O}$-modules.
L'image directe $f_*\mathcal{F}$, telle que définie dans le
Lemme \ref{sheaves-lemma-pushforward-presheaf-module},
est un faisceau de $f_*\mathcal{O}$-modules.
\end{lemma}

\begin{proof}
C'est évident d'après la définition et le Lemme \ref{sheaves-lemma-pushforward-sheaf}.
\end{proof}

\begin{lemma}
\label{sheaves-lemma-pullback-module}
Soit $f : X \to Y$ une application continue d'espaces topologiques.
Soit $\mathcal{O}$ un faisceau d'anneaux sur $Y$. Soit
$\mathcal{G}$ un faisceau de $\mathcal{O}$-modules.
Il existe un morphisme naturel de préfaisceaux d'ensembles sous-jacents
$$
f^{-1}\mathcal{O} \times f^{-1}\mathcal{G}
\longrightarrow
f^{-1}\mathcal{G}
$$
qui munit $f^{-1}\mathcal{G}$ d'une structure de
faisceau de $f^{-1}\mathcal{O}$-modules.
\end{lemma}

\begin{proof}
Rappelons que $f^{-1}$ est défini comme la composée du
foncteur $f_p$ et du foncteur faisceau associé. Le lemme
résulte donc du Lemme \ref{sheaves-lemma-pullback-presheaf-module}
et du Lemme \ref{sheaves-lemma-sheafification-presheaf-modules}.
\end{proof}

\noindent
Soit $f : X \to Y$ une application continue.
Soit $\mathcal{O}_X$ un faisceau d'anneaux sur $X$ et
soit $\mathcal{O}_Y$ un faisceau d'anneaux sur $Y$.
Nous avons donc maintenant défini des foncteurs
\begin{eqnarray*}
f_* : \textit{Mod}(\mathcal{O}_X) &
\longrightarrow &
\textit{Mod}(f_*\mathcal{O}_X) \\
f^{-1} : \textit{Mod}(\mathcal{O}_Y) &
\longrightarrow &
\textit{Mod}(f^{-1}\mathcal{O}_Y)
\end{eqnarray*}
Ils vérifient les compatibilités suivantes.

\begin{lemma}
\label{sheaves-lemma-adjoint-push-pull-modules}
Soit $f : X \to Y$ une application continue d'espaces topologiques.
Soit $\mathcal{O}$ un faisceau d'anneaux sur $Y$.
Soit $\mathcal{G}$ un faisceau de $\mathcal{O}$-modules.
Soit $\mathcal{F}$ un faisceau de $f^{-1}\mathcal{O}$-modules.
Alors
$$
\Mor_{\textit{Mod}(f^{-1}\mathcal{O})}(f^{-1}\mathcal{G}, \mathcal{F})
=
\Mor_{\textit{Mod}(\mathcal{O})}(\mathcal{G}, f_*\mathcal{F}).
$$
Nous utilisons ici
les Lemmes \ref{sheaves-lemma-pullback-module}
et \ref{sheaves-lemma-pushforward-module}, et considérons
$f_*\mathcal{F}$ comme un $\mathcal{O}$-module par restriction des scalaires le long de
$\mathcal{O} \to f_*f^{-1}\mathcal{O}$.
\end{lemma}

\begin{proof}
Le résultat découle des égalités
\begin{eqnarray*}
\Mor_{\textit{Mod}(f^{-1}\mathcal{O})}(f^{-1}\mathcal{G}, \mathcal{F})
& = &
\Mor_{\textit{Mod}(f_p\mathcal{O})}(f_p\mathcal{G}, \mathcal{F}) \\
& = &
\Mor_{\textit{Mod}(\mathcal{O})}(\mathcal{G}, f_*\mathcal{F}).
\end{eqnarray*}
La seconde égalité provient du
Lemme \ref{sheaves-lemma-adjoint-push-pull-presheaves-modules}
et la première du Lemme \ref{sheaves-lemma-sheafification-presheaf-modules}.
\end{proof}

\begin{lemma}
\label{sheaves-lemma-adjoint-pull-push-modules}
Soit $f : X \to Y$ une application continue d'espaces topologiques.
Soit $\mathcal{O}$ un faisceau d'anneaux sur $X$.
Soit $\mathcal{F}$ un faisceau de $\mathcal{O}$-modules.
Soit $\mathcal{G}$ un faisceau de $f_*\mathcal{O}$-modules.
Alors
$$
\Mor_{\textit{Mod}(\mathcal{O})}(
\mathcal{O} \otimes_{f^{-1}f_*\mathcal{O}} f^{-1}\mathcal{G}, \mathcal{F})
=
\Mor_{\textit{Mod}(f_*\mathcal{O})}(\mathcal{G}, f_*\mathcal{F}).
$$
Nous utilisons ici
les Lemmes \ref{sheaves-lemma-pullback-module}
et \ref{sheaves-lemma-pushforward-module}, ainsi que
le morphisme canonique $f^{-1}f_*\mathcal{O} \to \mathcal{O}$
dans la définition du produit tensoriel.
\end{lemma}

\begin{proof}
Cela résulte des égalités
\begin{eqnarray*}
\Mor_{\textit{Mod}(\mathcal{O})}(
\mathcal{O} \otimes_{f^{-1}f_*\mathcal{O}} f^{-1}\mathcal{G}, \mathcal{F})
& = &
\Mor_{\textit{Mod}(f^{-1}f_*\mathcal{O})}(
f^{-1}\mathcal{G}, \mathcal{F}_{f^{-1}f_*\mathcal{O}}) \\
& = &
\Mor_{\textit{Mod}(f_*\mathcal{O})}(\mathcal{G}, f_*\mathcal{F}).
\end{eqnarray*}
Ces égalités combinent le
Lemme \ref{sheaves-lemma-adjointness-tensor-restrict}
et le Lemme \ref{sheaves-lemma-adjoint-push-pull-modules}.
\end{proof}

\noindent
Soit $f : X \to Y$ une application continue.
Soit $\mathcal{O}_X$ un (pré)faisceau d'anneaux sur $X$ et
soit $\mathcal{O}_Y$ un (pré)faisceau d'anneaux sur $Y$.
À ce stade, nous avons donc défini les foncteurs
\begin{eqnarray*}
f_* : \textit{PMod}(\mathcal{O}_X) &
\longrightarrow &
\textit{PMod}(f_*\mathcal{O}_X) \\
f_* : \textit{Mod}(\mathcal{O}_X) &
\longrightarrow &
\textit{Mod}(f_*\mathcal{O}_X) \\
f_p : \textit{PMod}(\mathcal{O}_Y) &
\longrightarrow &
\textit{PMod}(f_p\mathcal{O}_Y) \\
f^{-1} : \textit{Mod}(\mathcal{O}_Y) &
\longrightarrow &
\textit{Mod}(f^{-1}\mathcal{O}_Y)
\end{eqnarray*}
Manifestement, en général, les deux foncteurs $(f_*, f^{-1})$
sur les faisceaux de modules ne sont pas adjoints, car leurs catégories d'arrivée
ne coïncident pas. En effet, comme nous l'avons vu plus haut, cela ne fonctionne que si, par miracle, les
faisceaux d'anneaux $\mathcal{O}_X, \mathcal{O}_Y$ vérifient les
relations $\mathcal{O}_X = f^{-1}\mathcal{O}_Y$ et
$\mathcal{O}_Y = f_*\mathcal{O}_X$. Ce n'est presque jamais
le cas en pratique. Nous interrompons la discussion pour définir
la notion correcte de morphisme pour laquelle il existe un couple
de foncteurs adjoints approprié sur les faisceaux de modules.

\section{Espaces annelés}
\label{sheaves-section-ringed-spaces}

\noindent
Soit $X$ un espace topologique et soit $\mathcal{O}_X$
un faisceau d'anneaux sur $X$. Nous devons
considérer le faisceau d'anneaux $\mathcal{O}_X$ comme un faisceau de fonctions
sur $X$. Et si $f : X \to Y$ est une application « convenable », alors, par composition,
une fonction sur $Y$ devient une fonction sur $X$.
Il devrait donc exister un
$f$-morphisme naturel de $\mathcal{O}_Y$ vers $\mathcal{O}_X$ ;
voir la définition \ref{sheaves-definition-f-map} et le
lemme \ref{sheaves-lemma-f-map}.
Pour un exemple précis, voir
l'exemple \ref{sheaves-example-continuous-map-ringed} ci-dessous. Voici la
définition abstraite pertinente.

\begin{definition}
\label{sheaves-definition-ringed-space}
Un {\it espace annelé} est un couple $(X, \mathcal{O}_X)$ formé
d'un espace topologique $X$ et d'un faisceau d'anneaux $\mathcal{O}_X$
sur $X$. Un {\it morphisme d'espaces annelés}
$(X, \mathcal{O}_X) \to (Y, \mathcal{O}_Y)$ est un couple
formé d'une application continue $f : X \to Y$ et d'un
$f$-morphisme de faisceaux d'anneaux
$f^\sharp : \mathcal{O}_Y \to \mathcal{O}_X$.
\end{definition}

\begin{example}
\label{sheaves-example-continuous-map-ringed}
Soit $f : X \to Y$ une application continue d'espaces topologiques.
Considérons les faisceaux de fonctions continues à valeurs réelles
$\mathcal{C}^0_X$ sur $X$ et $\mathcal{C}^0_Y$ sur $Y$ ;
voir l'exemple \ref{sheaves-example-C0-sheaf-rings}. Nous affirmons qu'il
existe un $f$-morphisme naturel $f^\sharp : \mathcal{C}^0_Y \to \mathcal{C}^0_X$
associé à $f$. Nous le définissons simplement par la
règle
\begin{eqnarray*}
\mathcal{C}^0_Y(V) & \longrightarrow & \mathcal{C}^0_X(f^{-1}V) \\
h & \longmapsto & h \circ f
\end{eqnarray*}
En toute rigueur, nous devrions écrire $f^\sharp(h) = h \circ f|_{f^{-1}(V)}$.
Il est clair qu'il s'agit d'une famille d'applications comme dans la
définition \ref{sheaves-definition-f-map}, compatible avec les structures de $\mathbf{R}$-algèbres.
C'est donc un $f$-morphisme de faisceaux de $\mathbf{R}$-algèbres ; voir le
lemme \ref{sheaves-lemma-f-map-sets-algebraic-structures}.

\medskip\noindent
Il existe bien sûr beaucoup d'autres situations dans lesquelles un
morphisme canonique d'espaces annelés est associé à un type
géométrique de morphisme. Par exemple, si $M$, $N$ sont des variétés $\mathcal{C}^\infty$
et si $f : M \to N$ est une application indéfiniment différentiable, alors
$f$ induit un morphisme canonique d'espaces annelés
$(M, \mathcal{C}_M^\infty) \to (N, \mathcal{C}^\infty_N)$.
La construction (qui est identique à celle ci-dessus) est laissée
au lecteur.
\end{example}

\noindent
La manière de composer les morphismes d'espaces
annelés n'étant peut-être pas tout à fait évidente, nous l'explicitons ici.

\begin{definition}
\label{sheaves-definition-composition-maps-ringed-spaces}
Soient
$(f, f^\sharp) : (X, \mathcal{O}_X) \to (Y, \mathcal{O}_Y)$ et
$(g, g^\sharp) : (Y, \mathcal{O}_Y) \to (Z, \mathcal{O}_Z)$
des morphismes d'espaces annelés. Nous définissons alors
la {\it composée de morphismes d'espaces annelés}
par la règle
$$
(g, g^\sharp) \circ (f, f^\sharp) = (g \circ f, f^\sharp \circ g^\sharp).
$$
Nous utilisons ici la composition des $f$-morphismes définie dans la
définition \ref{sheaves-definition-composition-f-maps}.
\end{definition}


\section{Morphismes d'espaces annelés et modules}
\label{sheaves-section-ringed-spaces-functoriality-modules}

\noindent
Nous avons maintenant introduit assez de notations pour pouvoir
définir l'image inverse et l'image directe des modules le long d'un morphisme
d'espaces annelés.

\begin{definition}
\label{sheaves-definition-pushforward}
Soit $(f, f^\sharp) : (X, \mathcal{O}_X) \to (Y, \mathcal{O}_Y)$
un morphisme d'espaces annelés.
\begin{enumerate}
\item Soit $\mathcal{F}$ un faisceau de $\mathcal{O}_X$-modules.
Nous définissons l'{\it image directe} de $\mathcal{F}$ comme le
faisceau de $\mathcal{O}_Y$-modules qui, comme faisceau
de groupes abéliens, est égal à $f_*\mathcal{F}$ et dont la
structure de module est obtenue par restriction des scalaires
via $f^\sharp : \mathcal{O}_Y \to f_*\mathcal{O}_X$
de la structure de module donnée
dans le lemme \ref{sheaves-lemma-pushforward-module}.
\item Soit $\mathcal{G}$ un faisceau de $\mathcal{O}_Y$-modules.
Nous définissons l'{\it image inverse} $f^*\mathcal{G}$ comme le
faisceau de $\mathcal{O}_X$-modules défini par la formule
$$
f^*\mathcal{G}
=
\mathcal{O}_X \otimes_{f^{-1}\mathcal{O}_Y} f^{-1}\mathcal{G}
$$
où l'homomorphisme d'anneaux $f^{-1}\mathcal{O}_Y \to \mathcal{O}_X$
est l'homomorphisme correspondant à $f^\sharp$, et où la structure de module
est donnée par le lemme \ref{sheaves-lemma-pullback-module}.
\end{enumerate}
\end{definition}

\noindent
Nous avons ainsi défini des foncteurs
\begin{eqnarray*}
f_* : \textit{Mod}(\mathcal{O}_X)
& \longrightarrow &
\textit{Mod}(\mathcal{O}_Y) \\
f^* : \textit{Mod}(\mathcal{O}_Y)
& \longrightarrow &
\textit{Mod}(\mathcal{O}_X)
\end{eqnarray*}
Le dernier résultat concernant ces foncteurs affirme qu'ils sont bien
adjoints, comme prévu.

\begin{lemma}
\label{sheaves-lemma-adjoint-pullback-pushforward-modules}
Soit $(f, f^\sharp) : (X, \mathcal{O}_X) \to (Y, \mathcal{O}_Y)$
un morphisme d'espaces annelés.
Soit $\mathcal{F}$ un faisceau de $\mathcal{O}_X$-modules.
Soit $\mathcal{G}$ un faisceau de $\mathcal{O}_Y$-modules.
Il existe une bijection canonique
$$
\Hom_{\mathcal{O}_X}(f^*\mathcal{G}, \mathcal{F})
=
\Hom_{\mathcal{O}_Y}(\mathcal{G}, f_*\mathcal{F}).
$$
Autrement dit, le foncteur $f^*$ est adjoint à gauche de
$f_*$.
\end{lemma}

\begin{proof}
Cela résulte du travail effectué précédemment :
\begin{eqnarray*}
\Hom_{\mathcal{O}_X}(f^*\mathcal{G}, \mathcal{F})
& = &
\Mor_{\textit{Mod}(\mathcal{O}_X)}(
\mathcal{O}_X \otimes_{f^{-1}\mathcal{O}_Y} f^{-1}\mathcal{G}, \mathcal{F}) \\
& = &
\Mor_{\textit{Mod}(f^{-1}\mathcal{O}_Y)}(
f^{-1}\mathcal{G}, \mathcal{F}_{f^{-1}\mathcal{O}_Y}) \\
& = &
\Hom_{\mathcal{O}_Y}(\mathcal{G}, f_*\mathcal{F}).
\end{eqnarray*}
Nous utilisons ici les lemmes \ref{sheaves-lemma-adjointness-tensor-restrict}
et \ref{sheaves-lemma-adjoint-push-pull-modules}.
\end{proof}

\begin{lemma}
\label{sheaves-lemma-push-pull-composition-modules}
Soient $f : X \to Y$ et $g : Y \to Z$ des morphismes d'espaces annelés.
Les foncteurs $(g \circ f)_*$ et $g_* \circ f_*$ sont égaux.
Il existe un isomorphisme canonique de foncteurs
$(g \circ f)^* \cong f^* \circ g^*$.
\end{lemma}

\begin{proof}
Le résultat sur les images directes est une conséquence du lemme
\ref{sheaves-lemma-pushforward-composition} et de nos définitions.
Le résultat sur les images inverses s'en déduit par le même
argument que dans la preuve du lemme \ref{sheaves-lemma-pullback-composition}.
\end{proof}


\noindent
Étant donnés un morphisme d'espaces annelés
$(f, f^\sharp) : (X, \mathcal{O}_X) \to (Y, \mathcal{O}_Y)$,
un faisceau de $\mathcal{O}_X$-modules $\mathcal{F}$
et un faisceau de $\mathcal{O}_Y$-modules $\mathcal{G}$ sur $Y$,
la notion d'{\it $f$-morphisme $\varphi : \mathcal{G} \to \mathcal{F}$
de faisceaux de modules} a un sens. Nous pouvons simplement le définir
comme un $f$-morphisme $\varphi : \mathcal{G} \to \mathcal{F}$
de faisceaux abéliens (voir la définition \ref{sheaves-definition-f-map} et le
lemme \ref{sheaves-lemma-f-map}) tel que, pour tout ouvert $V \subset Y$, l'application
$$
\mathcal{G}(V) \longrightarrow \mathcal{F}(f^{-1}V)
$$
soit un morphisme de $\mathcal{O}_Y(V)$-modules. Nous considérons ici
$\mathcal{F}(f^{-1}V)$ comme un $\mathcal{O}_Y(V)$-module
au moyen de l'homomorphisme $f^\sharp_V : \mathcal{O}_Y(V) \to \mathcal{O}_X(f^{-1}V)$.
L'ensemble des $f$-morphismes entre
$\mathcal{G}$ et $\mathcal{F}$ est en bijection canonique
avec les ensembles
$\Mor_{\textit{Mod}(\mathcal{O}_X)}(f^*\mathcal{G}, \mathcal{F})$
et
$\Mor_{\textit{Mod}(\mathcal{O}_Y)}(\mathcal{G}, f_*\mathcal{F})$.
Voir ci-dessus.

\medskip\noindent
La composition des $f$-morphismes est définie exactement de la
même manière que dans le cas des $f$-morphismes de faisceaux
d'ensembles. En outre, étant donnés un $f$-morphisme $\mathcal{G} \to \mathcal{F}$
comme ci-dessus et $x \in X$, l'application induite sur les fibres
$$
\varphi_x : \mathcal{G}_{f(x)} \longrightarrow \mathcal{F}_x
$$
est un morphisme de $\mathcal{O}_{Y, f(x)}$-modules, où la
structure de $\mathcal{O}_{Y, f(x)}$-module sur $\mathcal{F}_x$
provient de la structure de $\mathcal{O}_{X, x}$-module au moyen de
l'application $f^\sharp_x : \mathcal{O}_{Y, f(x)} \to \mathcal{O}_{X, x}$.
Voici un lemme apparenté.

\begin{lemma}
\label{sheaves-lemma-stalk-pullback-modules}
Soit $(f, f^\sharp) : (X, \mathcal{O}_X) \to (Y, \mathcal{O}_Y)$
un morphisme d'espaces annelés.
Soit $\mathcal{G}$ un faisceau de $\mathcal{O}_Y$-modules.
Soit $x \in X$. Alors
$$
(f^*\mathcal{G})_x =
\mathcal{G}_{f(x)}
\otimes_{\mathcal{O}_{Y, f(x)}}
\mathcal{O}_{X, x}
$$
comme $\mathcal{O}_{X, x}$-modules, le produit tensoriel de droite
utilisant $f^\sharp_x : \mathcal{O}_{Y, f(x)} \to \mathcal{O}_{X, x}$.
\end{lemma}

\begin{proof}
Cela résulte du lemme \ref{sheaves-lemma-stalk-tensor-sheaf-modules}
et de l'identification des fibres des faisceaux images inverses
en $x$ avec les fibres correspondantes en $f(x)$. Voir par
exemple les formules de la section \ref{sheaves-section-presheaves-structures-functorial}.
\end{proof}



\section{Faisceaux gratte-ciel et fibres}
\label{sheaves-section-skyscraper-sheaves}

\begin{definition}
\label{sheaves-definition-skyscraper-sheaf}
Soit $X$ un espace topologique.
\begin{enumerate}
\item Soit $x \in X$ un point. Notons $i_x : \{x\} \to X$ l'application d'inclusion.
Soit $A$ un ensemble et considérons $A$ comme un faisceau sur l'espace à un point $\{x\}$.
Nous appelons $i_{x, *}A$ le {\it faisceau gratte-ciel en $x$ de valeur $A$}.
\item Si, dans (1) ci-dessus, $A$ est un groupe abélien, nous considérons
$i_{x, *}A$ comme un faisceau de groupes abéliens sur $X$.
\item Si, dans (1) ci-dessus, $A$ est une structure algébrique, nous considérons
$i_{x, *}A$ comme un faisceau de structures algébriques.
\item Si $(X, \mathcal{O}_X)$ est un espace annelé, nous considérons
$i_x : \{x\} \to X$ comme un morphisme d'espaces annelés
$(\{x\}, \mathcal{O}_{X, x}) \to (X, \mathcal{O}_X)$
et, si $A$ est un $\mathcal{O}_{X, x}$-module, nous considérons
$i_{x, *}A$ comme un faisceau de $\mathcal{O}_X$-modules.
\item On dit qu'un faisceau d'ensembles $\mathcal{F}$ est un {\it faisceau gratte-ciel}
s'il existe un point $x$ de $X$ et un ensemble $A$ tels
que $\mathcal{F} \cong i_{x, *}A$.
\item On dit qu'un faisceau de groupes abéliens $\mathcal{F}$ est un
{\it faisceau gratte-ciel} s'il existe un point $x$ de $X$
et un groupe abélien $A$ tels que $\mathcal{F} \cong i_{x, *}A$
comme faisceaux de groupes abéliens.
\item On dit qu'un faisceau de structures algébriques $\mathcal{F}$ est un
{\it faisceau gratte-ciel} s'il existe un point $x$ de $X$
et une structure algébrique $A$ tels que $\mathcal{F} \cong i_{x, *}A$
comme faisceaux de structures algébriques.
\item Si $(X, \mathcal{O}_X)$ est un espace annelé et
si $\mathcal{F}$ est un faisceau de $\mathcal{O}_X$-modules,
on dit que $\mathcal{F}$ est un {\it faisceau gratte-ciel} s'il
existe un point $x \in X$ et un $\mathcal{O}_{X, x}$-module
$A$ tels que $\mathcal{F} \cong i_{x, *}A$
comme faisceaux de $\mathcal{O}_X$-modules.
\end{enumerate}
\end{definition}

\begin{lemma}
\label{sheaves-lemma-skyscraper-stalks}
Soient $X$ un espace topologique, $x \in X$ un point et
$A$ un ensemble. Pour tout point $x' \in X$, la fibre du
faisceau gratte-ciel en $x$ de valeur $A$ au point $x'$ est
$$
(i_{x, *}A)_{x'} =
\left\{
\begin{matrix}
A & \text{si} & x' \in \overline{\{x\}} \\
\{*\} & \text{si} & x' \not\in \overline{\{x\}}
\end{matrix}
\right.
$$
Une description analogue vaut dans le cas des
groupes abéliens, des structures algébriques et des
faisceaux de modules.
\end{lemma}

\begin{proof}
Démonstration omise.
\end{proof}

\begin{lemma}
\label{sheaves-lemma-stalk-skyscraper-adjoint}
Soit $X$ un espace topologique et soit $x \in X$ un point.
Les foncteurs $\mathcal{F} \mapsto \mathcal{F}_x$ et
$A \mapsto i_{x, *}A$ sont adjoints. Plus précisément,
$$
\Mor_{\textit{Ensembles}}(\mathcal{F}_x, A)
=
\Mor_{\Sh(X)}(\mathcal{F}, i_{x, *}A).
$$
Un énoncé analogue vaut dans le cas des
groupes abéliens et des structures algébriques. Dans le cas des
faisceaux de modules, nous avons
$$
\Hom_{\mathcal{O}_{X, x}}(\mathcal{F}_x, A)
=
\Hom_{\mathcal{O}_X}(\mathcal{F}, i_{x, *}A).
$$
\end{lemma}

\begin{proof}
Démonstration omise. Indication : le foncteur fibre peut être
considéré comme le foncteur image inverse pour le morphisme $i_x : \{x\} \to X$.
L'adjonction résulte alors de celle de
$i_x^{-1}$ et $i_{x, *}$ (resp.\ $i_x^*$ et $i_{x, *}$
dans le cas des faisceaux de modules).
\end{proof}






\section{Limites projectives et inductives de préfaisceaux}
\label{sheaves-section-limits-presheaves}

\noindent
Soit $X$ un espace topologique.
Soit $\mathcal{I} \to \textit{PSh}(X)$, $i \mapsto \mathcal{F}_i$,
un diagramme.
\begin{enumerate}
\item Les objets $\lim_i \mathcal{F}_i$ et $\colim_i \mathcal{F}_i$
existent tous deux.
\item Pour tout ouvert $U \subset X$, on a
$$
(\lim_i \mathcal{F}_i)(U) =
\lim_i \mathcal{F}_i(U)
$$
et
$$
(\colim_i \mathcal{F}_i)(U) =
\colim_i \mathcal{F}_i(U).
$$
\item Soit $x \in X$ un point. En général, la fibre de
$\lim_i \mathcal{F}_i$ en $x$ n'est pas égale à
la limite projective des fibres. C'est toutefois le cas si la catégorie d'indices est finie.
Autrement dit, le foncteur fibre est
exact à gauche (voir Catégories, définition \ref{categories-definition-exact}).
\item Soit $x \in X$. On a toujours
$$
(\colim_i \mathcal{F}_i)_x =
\colim_i \mathcal{F}_{i, x}.
$$
\end{enumerate}
Toutes les démonstrations sont faciles.

\section{Limites projectives et inductives de faisceaux}
\label{sheaves-section-limits-sheaves}

\noindent
Soit $X$ un espace topologique.
Soit $\mathcal{I} \to \Sh(X)$, $i \mapsto \mathcal{F}_i$,
un diagramme.
\begin{enumerate}
\item Les objets $\lim_i \mathcal{F}_i$ et $\colim_i \mathcal{F}_i$
existent tous deux.
\item Le foncteur d'inclusion $i : \Sh(X) \to \textit{PSh}(X)$
commute aux limites projectives. Autrement dit, on peut calculer la limite projective
dans la catégorie des faisceaux comme la limite projective dans la catégorie des
préfaisceaux. En particulier, pour tout ouvert $U \subset X$, on a
$$
(\lim_i \mathcal{F}_i)(U) = \lim_i \mathcal{F}_i(U).
$$
\item Le foncteur d'inclusion $i : \Sh(X) \to \textit{PSh}(X)$
ne commute pas en général aux limites inductives (pas même
aux limites inductives finies --- penser aux surjections). La limite inductive
s'obtient en prenant le faisceau associé à la limite inductive dans la
catégorie des préfaisceaux :
$$
\colim_i \mathcal{F}_i =
\Big(U \mapsto \colim_i \mathcal{F}_i(U)\Big)^\#.
$$
\item Soit $x \in X$ un point. En général, la fibre de
$\lim_i \mathcal{F}_i$ en $x$ n'est pas égale à
la limite projective des fibres. C'est toutefois le cas si la catégorie d'indices est finie.
Autrement dit, le foncteur fibre est
exact à gauche.
\item Soit $x \in X$. On a toujours
$$
(\colim_i \mathcal{F}_i)_x = \colim_i \mathcal{F}_{i, x}.
$$
\item Le foncteur faisceau associé
${}^\# : \textit{PSh}(X) \to \Sh(X)$ commute à toutes les
limites inductives et aux limites projectives finies. Il ne commute cependant pas
à toutes les limites projectives.
\end{enumerate}
Toutes les démonstrations sont faciles. Voici un exemple de ce qui vaut pour les limites
inductives filtrantes de faisceaux.

\begin{lemma}
\label{sheaves-lemma-directed-colimits-sections}
Soit $X$ un espace topologique. Soit $I$ un ensemble ordonné filtrant.
Soit $(\mathcal{F}_i, \varphi_{ii'})$ un système inductif de faisceaux d'ensembles
indexé par $I$, voir
Catégories, section \ref{categories-section-posets-limits}.
Soit $U \subset X$ un ouvert.
Considérons l'application canonique
$$
\Psi :
\colim_i \mathcal{F}_i(U)
\longrightarrow
\left(\colim_i \mathcal{F}_i\right)(U)
$$
\begin{enumerate}
\item Si toutes les applications de transition sont injectives, alors
$\Psi$ est injective pour tout ouvert $U$.
\item Si $U$ est quasi-compact, alors $\Psi$ est injective.
\item Si $U$ est quasi-compact et si toutes les applications de transition sont injectives,
alors $\Psi$ est un isomorphisme.
\item Si $U$ admet un système cofinal de recouvrements ouverts
$\mathcal{U} : U = \bigcup_{j\in J} U_j$, où
$J$ étant fini, tels que $U_j \cap U_{j'}$ soit quasi-compact
pour tous $j, j' \in J$, alors $\Psi$ est bijective.
\end{enumerate}
\end{lemma}

\begin{proof}
Supposons toutes les applications de transition injectives. Dans ce cas, le préfaisceau
$\mathcal{F}' : V \mapsto \colim_i \mathcal{F}_i(V)$ est
séparé (voir la Définition \ref{sheaves-definition-separated}).
D'après ce qui précède, on a
$(\mathcal{F}')^\# = \colim_i \mathcal{F}_i$.
Le Lemme \ref{sheaves-lemma-separated-presheaf-into-sheaf} montre que
$\mathcal{F}' \to (\mathcal{F}')^\#$ est injective. Cela démontre (1).

\medskip\noindent
Supposons $U$ quasi-compact.
Supposons que $s \in \mathcal{F}_i(U)$ et
$s' \in \mathcal{F}_{i'}(U)$ déterminent au membre de gauche
des éléments ayant la même image par $\Psi$.
Comme $U$ est quasi-compact, cela signifie qu'il existe
un recouvrement ouvert fini $U = \bigcup_{j = 1, \ldots, m} U_j$
et, pour chaque $j$, un indice $i_j \in I$, avec $i_j \geq i$, $i_j \geq i'$,
tels que $\varphi_{ii_j}(s)$ et $\varphi_{i'i_j}(s')$ aient la
même restriction à $U_j$.
Soit $i''\in I$ supérieur ou égal ($\geq$) à chacun des $i_j$.
On en conclut que $\varphi_{ii''}(s)$ et $\varphi_{i'i''}(s')$
coïncident sur les ouverts $U_j$ pour tout $j$, et donc que
$\varphi_{ii''}(s) = \varphi_{i'i''}(s')$. Cela démontre (2).

\medskip\noindent
Supposons $U$ quasi-compact et toutes les applications de transition injectives.
Soit $s$ un élément du but de $\Psi$.
Comme $U$ est quasi-compact,
il existe un recouvrement ouvert fini $U = \bigcup_{j = 1, \ldots, m} U_j$,
ainsi que, pour chaque $j$, un indice $i_j \in I$ et un $s_j \in \mathcal{F}_{i_j}(U_j)$,
tels que $s|_{U_j}$ provienne de $s_j$ pour tout $j$.
Choisissons $i \in I$ supérieur ou égal ($\geq$) à chacun des $i_j$.
D'après (1), les sections $\varphi_{i_ji}(s_j)$ coïncident sur les
intersections $U_j \cap U_{j'}$. Elles se recollent donc en une section
$s' \in \mathcal{F}_i(U)$ qui s'envoie sur $s$ par $\Psi$.
Cela démontre (3).

\medskip\noindent
Supposons l'hypothèse de (4). En particulier,
$U$ est quasi-compact et (2) assure donc l'injectivité de $\Psi$.
Soit $s$ un élément du but de $\Psi$.
Par hypothèse, il existe un recouvrement ouvert fini
$U = \bigcup_{j = 1, \ldots, m} U_j$, où $U_j \cap U_{j'}$ est
quasi-compact pour tous $j, j' \in J$, ainsi que,
pour chaque $j$, un indice $i_j \in I$ et un $s_j \in \mathcal{F}_{i_j}(U_j)$
tels que $s|_{U_j}$ soit l'image de $s_j$ pour tout $j$.
Puisque $U_j \cap U_{j'}$ est quasi-compact, on peut appliquer (2)
et voir qu'il existe un $i_{jj'} \in I$, avec
$i_{jj'} \geq i_j$, $i_{jj'} \geq i_{j'}$, tel que
$\varphi_{i_ji_{jj'}}(s_j)$ et $\varphi_{i_{j'}i_{jj'}}(s_{j'})$
coïncident sur $U_j \cap U_{j'}$. Choisissons un indice $i \in I$
majorant tous les $i_{jj'}$. On voit alors que
les sections $\varphi_{i_ji}(s_j)$ de $\mathcal{F}_i$ se recollent
en une section de $\mathcal{F}_i$ sur $U$. Cette section s'envoie
sur l'élément $s$, comme voulu.
\end{proof}

\begin{example}
\label{sheaves-example-conditions-needed-colimit}
Soit $X = \{s_1, s_2, \xi_1, \xi_2, \xi_3, \ldots\}$ comme ensemble.
Déclarons qu'une partie $U \subset X$ est ouverte
si $s_1 \in U$ ou $s_2 \in U$ entraîne que $U$ contient tous les
$\xi_i$. Soit $U_n = \{\xi_n, \xi_{n + 1}, \ldots\}$, et
soit $j_n : U_n \to X$ l'application d'inclusion.
Posons $\mathcal{F}_n = j_{n, *}\underline{\mathbf{Z}}$.
On a des applications de transition $\mathcal{F}_n \to \mathcal{F}_{n + 1}$.
Soit $\mathcal{F} = \colim \mathcal{F}_n$.
Remarquons que $\mathcal{F}_{n, \xi_m} = 0$ si $m < n$,
car $\{\xi_m\}$ est une partie ouverte de $X$ disjointe de $U_n$.
On voit donc que $\mathcal{F}_{\xi_n} = 0$ pour tout $n$.
En revanche, la fibre $\mathcal{F}_{s_i}$, pour $i = 1, 2$,
est la limite inductive
$$
M = \colim_n \prod\nolimits_{m \geq n} \mathbf{Z}
$$
qui n'est pas nulle. On en déduit que le faisceau
$\mathcal{F}$ est la somme directe des
faisceaux gratte-ciel de valeur $M$ aux points fermés
$s_1$ et $s_2$. Ainsi, $\Gamma(X, \mathcal{F}) = M \oplus M$.
D'autre part, le lecteur peut vérifier que
$\colim_n \Gamma(X, \mathcal{F}_n) = M$.
Une condition est donc nécessaire dans l'assertion (4) du
Lemme \ref{sheaves-lemma-directed-colimits-sections} ci-dessus.
\end{example}

\noindent
Il existe une version du lemme précédent pour les faisceaux sur un diagramme
d'espaces spectraux. Introduisons les notations nécessaires à son énoncé.
Soit $\mathcal{I}$ une catégorie d'indices cofiltrante.
Soit $i \mapsto X_i$ un diagramme d'espaces spectraux indexé par $\mathcal{I}$,
tel que, pour $a : j \to i$ dans $\mathcal{I}$, l'application correspondante
$f_a : X_j \to X_i$ soit spectrale. Posons $X = \lim X_i$ et notons
$p_i : X \to X_i$ la projection.

\begin{lemma}
\label{sheaves-lemma-compute-pullback-to-limit}
Dans la situation décrite ci-dessus, soit $i \in \Ob(\mathcal{I})$ et soit
$\mathcal{G}$ un faisceau sur $X_i$. Pour tout ouvert quasi-compact $U_i \subset X_i$,
on a
$$
p_i^{-1}\mathcal{G}(p_i^{-1}(U_i)) =
\colim_{a : j \to i} f_a^{-1}\mathcal{G}(f_a^{-1}(U_i))
$$
\end{lemma}

\begin{proof}
Montrons que l'application canonique
$\colim_{a : j \to i} f_a^{-1}\mathcal{G}(f_a^{-1}(U_i)) \to
p_i^{-1}\mathcal{G}(p_i^{-1}(U_i))$ est injective.
Soient $s, s'$ des sections de $f_a^{-1}\mathcal{G}$
sur $f_a^{-1}(U_i)$ pour un certain $a : j \to i$. Pour
$b : k \to j$, soit $Z_k \subset f_{a \circ b}^{-1}(U_i)$
la partie fermée formée des points $x$ tels que les images de
$s$ et $s'$ dans la fibre $(f_{a \circ b}^{-1}\mathcal{G})_x$
soient distinctes. Si $Z_k$ est non vide pour tout $b : k \to j$, alors le
lemme \ref{topology-lemma-inverse-limit-spectral-spaces-nonempty} de Topologie
montre que $\lim_{b : k \to j} Z_k$ est également non vide.
Pour $x \in \lim_{b : k \to j} Z_k \subset X$
(remarquons que $\mathcal{I}/j \to \mathcal{I}$ est initial),
les images de $s$ et $s'$ dans la fibre de
$p_i^{-1}\mathcal{G}$ en $x$ sont alors également distinctes, puisque
$(p_i^{-1}\mathcal{G})_x = (f_{b \circ a}^{-1}\mathcal{G})_{p_k(x)}$
pour tout $b : k \to j$ comme ci-dessus. Ainsi, si les images de $s$ et $s'$
dans $p_i^{-1}\mathcal{G}(p_i^{-1}(U_i))$ sont égales, alors $Z_k$
est vide pour un certain $b : k \to j$. Cela démontre l'injectivité.

\medskip\noindent
Surjectivité. Soit $s$ une section de $p_i^{-1}\mathcal{G}$ sur
$p_i^{-1}(U_i)$. D'après le
lemme \ref{topology-lemma-directed-inverse-limit-spectral-spaces} de Topologie,
la partie $p_i^{-1}(U_i)$ est un ouvert quasi-compact de l'espace spectral $X$.
Par construction du faisceau image inverse, on peut trouver un recouvrement ouvert
$p_i^{-1}(U_i) = \bigcup_{l \in L} W_l$,
des ouverts $V_{l, i} \subset X_i$ et
des sections $s_{l, i} \in \mathcal{G}(V_{l, i})$
tels que $p_i(W_l) \subset V_{l, i}$ et $p_i^{-1}s_{l, i}|_{W_l} = s|_{W_l}$.
Comme $X$ et $X_i$ sont spectraux et que $p_i^{-1}(U_i)$ est un ouvert quasi-compact,
on peut supposer $L$ fini et les $W_l$ et $V_{l, i}$ ouverts quasi-compacts
pour tout $l$. On peut alors appliquer le
lemme \ref{topology-lemma-descend-opens} de Topologie
pour trouver un $a : j \to i$ et un recouvrement ouvert
$f_a^{-1}(U_i) = \bigcup_{l \in L} W_{l, j}$ par des ouverts
quasi-compacts dont l'image inverse sur $X$ est le recouvrement
$p_i^{-1}(U_i) = \bigcup_{l \in L} W_l$
et tels qu'en outre $W_{l, j} \subset f_a^{-1}(V_{l, i})$.
Notons $s_{l, j}$ la restriction de l'image inverse de $s_{l, i}$
par $f_a$ à $W_{l, j}$. On voit alors que $s_{l, j}$ et $s_{l', j}$
se restreignent en des éléments de $(f_a^{-1}\mathcal{G})(W_{l, j} \cap W_{l', j})$
dont les images inverses donnent le même élément $(p_i^{-1}\mathcal{G})(W_l \cap W_{l'})$,
à savoir la restriction de $s$. L'injectivité permet donc de trouver
$b : k \to j$ tel que les sections $f_b^{-1}s_{l, j}$
se recollent en une section sur $f_{a \circ b}^{-1}(U_i)$, comme voulu.
\end{proof}

\noindent
Supposons maintenant donnés, en plus du système cofiltrant $X_i$ d'espaces spectraux,
les objets suivants :
\begin{enumerate}
\item un faisceau $\mathcal{F}_i$ sur $X_i$ pour tout $i \in \Ob(\mathcal{I})$ ;
\item pour $a : j \to i$, un $f_a$-morphisme
$\varphi_a : \mathcal{F}_i \to \mathcal{F}_j$
\end{enumerate}
tels que $\varphi_c = \varphi_b \circ \varphi_a$
dès que $c = a \circ b$. Posons $\mathcal{F} = \colim p_i^{-1}\mathcal{F}_i$
sur $X$.

\begin{lemma}
\label{sheaves-lemma-descend-opens}
Dans la situation décrite ci-dessus, soit $i \in \Ob(\mathcal{I})$ et soit
$U_i \subset X_i$ un ouvert quasi-compact. Alors
$$
\colim_{a : j \to i} \mathcal{F}_j(f_a^{-1}(U_i)) = \mathcal{F}(p_i^{-1}(U_i))
$$
\end{lemma}

\begin{proof}
Rappelons que $p_i^{-1}(U_i)$ est un ouvert quasi-compact de l'espace spectral
$X$, voir le
lemme \ref{topology-lemma-directed-inverse-limit-spectral-spaces} de Topologie.
Le Lemme \ref{sheaves-lemma-directed-colimits-sections} s'applique donc et donne
$$
\mathcal{F}(p_i^{-1}(U_i)) =
\colim_{a : j \to i} p_j^{-1}\mathcal{F}_j(p_i^{-1}(U_i)).
$$
Un argument formel montre que
$$
\colim_{a : j \to i} \mathcal{F}_j(f_a^{-1}(U_i)) =
\colim_{a : j \to i} \colim_{b : k \to j}
f_b^{-1}\mathcal{F}_j(f_{a \circ b}^{-1}(U_i))
$$
Il suffit donc de montrer que
$$
p_j^{-1}\mathcal{F}_j(p_i^{-1}(U_i)) =
\colim_{b : k \to j} f_b^{-1}\mathcal{F}_j(f_{a \circ b}^{-1}(U_i))
$$
C'est le Lemme \ref{sheaves-lemma-compute-pullback-to-limit}
appliqué à $\mathcal{F}_j$ et à l'ouvert quasi-compact $f_a^{-1}(U_i)$.
\end{proof}









\section{Bases et faisceaux}
\label{sheaves-section-bases}

\noindent
Il existe parfois une base de la topologie
formée d'ouverts avec lesquels il est plus facile de travailler
qu'avec des ouverts quelconques. Pour plus de commodité, nous donnons ici
quelques définitions et lemmes simples afin de
faciliter le travail avec les (pré)faisceaux dans une telle situation.

\begin{definition}
\label{sheaves-definition-presheaf-basis}
Soit $X$ un espace topologique. Soit $\mathcal{B}$ une
base de la topologie de $X$.
\begin{enumerate}
\item Un {\it préfaisceau d'ensembles $\mathcal{F}$ sur $\mathcal{B}$}
est une règle qui associe à chaque $U \in \mathcal{B}$ un ensemble
$\mathcal{F}(U)$ et à chaque inclusion $V \subset U$
d'éléments de $\mathcal{B}$ une application
$\rho^U_V : \mathcal{F}(U) \to \mathcal{F}(V)$ telle que
$\rho^U_U = \text{id}_{\mathcal{F}(U)}$ pour tout $U \in \mathcal{B}$ et que,
si $W \subset V \subset U$ dans $\mathcal{B}$, on ait
$\rho^U_W = \rho^V_W \circ \rho ^U_V$.
\item Un {\it morphisme $\varphi : \mathcal{F} \to \mathcal{G}$
de préfaisceaux d'ensembles sur $\mathcal{B}$} est une règle qui associe à chaque
élément $U \in \mathcal{B}$ une application d'ensembles $\varphi : \mathcal{F}(U)
\to \mathcal{G}(U)$ compatible avec les applications de restriction.
\end{enumerate}
\end{definition}

\noindent
Comme pour les préfaisceaux usuels, nous employons la terminologie des sections,
des restrictions de sections, etc. En particulier, nous pouvons définir la
{\it fibre} de $\mathcal{F}$ en un point $x \in X$ comme la
limite inductive
$$
\mathcal{F}_x = \colim_{U\in \mathcal{B}, x\in U} \mathcal{F}(U).
$$
Comme pour la fibre d'un préfaisceau sur $X$, cette limite inductive est
filtrante. En effet, la famille des $U\in \mathcal{B}$,
$x \in U$, est un système fondamental de voisinages ouverts de $x$.

\medskip\noindent
Il est facile de construire des exemples montrant que la notion de préfaisceau
sur $X$ est très différente de celle de préfaisceau sur une base
de la topologie de $X$. Ce phénomène ne se produit pas pour les
faisceaux. La notion suivante est donc beaucoup plus utile.

\begin{definition}
\label{sheaves-definition-sheaf-basis}
Soit $X$ un espace topologique. Soit $\mathcal{B}$ une
base de la topologie de $X$.
\begin{enumerate}
\item Un {\it faisceau d'ensembles $\mathcal{F}$ sur $\mathcal{B}$} est un préfaisceau
d'ensembles sur $\mathcal{B}$ qui satisfait la propriété supplémentaire
suivante : étant donnés $U \in \mathcal{B}$, un recouvrement
$U = \bigcup_{i \in I} U_i$ avec $U_i \in \mathcal{B}$, et
des recouvrements $U_i \cap U_j = \bigcup_{k \in I_{ij}} U_{ijk}$ avec
$U_{ijk} \in \mathcal{B}$, la condition de faisceau suivante est satisfaite :
\begin{itemize}
\item[(**)] Pour toute famille
de sections $s_i \in \mathcal{F}(U_i)$, $i \in I$, telle que
$\forall i, j\in I$, $\forall k\in I_{ij}$,
$$
s_i|_{U_{ijk}} = s_j|_{U_{ijk}}
$$
il existe une unique section $s \in \mathcal{F}(U)$ telle que
$s_i = s|_{U_i}$ pour tout $i \in I$.
\end{itemize}
\item Un {\it morphisme de faisceaux d'ensembles sur $\mathcal{B}$} est simplement un
morphisme de préfaisceaux d'ensembles.
\end{enumerate}
\end{definition}

\noindent
Nous expliquons d'abord qu'il suffit de
vérifier la condition de faisceau $(**)$
sur un système cofinal de recouvrements.
Dans la situation de la définition, supposons que
$U \in \mathcal{B}$. Notons temporairement
$\text{Cov}_\mathcal{B}(U)$ l'ensemble de tous les recouvrements
de $U$ par des éléments de $\mathcal{B}$.
Remarquons que $\text{Cov}_\mathcal{B}(U)$ est préordonné par le raffinement.
Une partie $C \subset \text{Cov}_\mathcal{B}(U)$ est un
système cofinal si, pour tout $\mathcal{U} \in \text{Cov}_\mathcal{B}(U)$,
il existe un recouvrement $\mathcal{V} \in C$ qui raffine $\mathcal{U}$.

\begin{lemma}
\label{sheaves-lemma-cofinal-systems-coverings}
Conservons les notations ci-dessus.
Pour chaque $U \in \mathcal{B}$, soit $C(U) \subset \text{Cov}_\mathcal{B}(U)$
un système cofinal. Pour chaque $U \in \mathcal{B}$ et chaque
$\mathcal{U} : U = \bigcup U_i$ appartenant à $C(U)$, supposons donnés des recouvrements
$\mathcal{U}_{ij} : U_i \cap U_j = \bigcup U_{ijk}$,
$U_{ijk} \in \mathcal{B}$.
Soit $\mathcal{F}$ un préfaisceau d'ensembles sur $\mathcal{B}$.
Les conditions suivantes sont équivalentes :
\begin{enumerate}
\item Le préfaisceau $\mathcal{F}$ est un faisceau sur $\mathcal{B}$.
\item Pour tout $U \in \mathcal{B}$ et tout recouvrement
$\mathcal{U} : U = \bigcup U_i$ appartenant à $C(U)$, la condition de faisceau
$(**)$ est satisfaite (pour les recouvrements $\mathcal{U}_{ij}$ donnés).
\end{enumerate}
\end{lemma}

\begin{proof}
Nous devons montrer que (2) implique (1).
Supposons que $U \in \mathcal{B}$ et que
$\mathcal{U} : U = \bigcup_{i\in I} U_i$ soit un recouvrement quelconque
par des éléments de $\mathcal{B}$. Comme le système $C(U)$ est cofinal,
nous pouvons trouver un élément $\mathcal{V} : U = \bigcup_{j \in J} V_j$
de $C(U)$ qui raffine $\mathcal{U}$. Cela signifie qu'il existe
une application $\alpha : J \to I$ telle que $V_j \subset U_{\alpha(j)}$.

\medskip\noindent
Remarquons que, si $s, s' \in \mathcal{F}(U)$ sont des sections telles
que $s|_{U_i} = s'|_{U_i}$, alors
$$
s|_{V_j}
= (s|_{U_{\alpha(j)}})|_{V_j}
= (s'|_{U_{\alpha(j)}})|_{V_j}
= s'|_{V_j}
$$
pour tout $j$. L'unicité dans $(**)$
pour le recouvrement $\mathcal{V}$ entraîne donc que $s = s'$.
Nous avons ainsi démontré la partie unicité de $(**)$
pour notre recouvrement quelconque $\mathcal{U}$.

\medskip\noindent
Supposons en outre que $U_i \cap U_{i'} = \bigcup_{k \in I_{ii'}} U_{ii'k}$
soient des recouvrements quelconques par des $U_{ii'k} \in \mathcal{B}$.
Essayons de démontrer la partie existence de $(**)$ pour le système
$(\mathcal{U}, \mathcal{U}_{ij})$. Soient donc $s_i \in \mathcal{F}(U_i)$
et supposons que
$$
s_i|_{U_{ii'k}} = s_{i'}|_{U_{ii'k}}
$$
pour tous $i, i', k$. Posons $t_j = s_{\alpha(j)}|_{V_j}$, où $\mathcal{V}$
et $\alpha$ sont comme ci-dessus.

\medskip\noindent
L'argument présente ici une petite difficulté. Soit en effet
$\mathcal{V}_{jj'} : V_j \cap V_{j'} = \bigcup_{l \in J_{jj'}} V_{jj'l}$
le recouvrement fourni par l'énoncé du lemme.
Il n'est pas clair a priori que
$$
t_j|_{V_{jj'l}} = t_{j'}|_{V_{jj'l}}
$$
pour tous $j, j', l$. Pour le voir, remarquons que nous avons bien
$$
t_j|_W = t_{j'}|_W \text{ pour tout } W \in \mathcal{B},
W \subset V_{jj'l} \cap U_{\alpha(j)\alpha(j')k}
$$
pour tout $k \in I_{\alpha(j)\alpha(j')}$, par notre hypothèse sur
la famille d'éléments $s_i$. Et puisque
$V_j \cap V_{j'} \subset U_{\alpha(j)} \cap U_{\alpha(j')}$,
nous voyons que $t_j|_{V_{jj'l}}$ et $t_{j'}|_{V_{jj'l}}$
coïncident sur les membres d'un recouvrement de $V_{jj'l}$ par des
éléments de $\mathcal{B}$. La partie unicité démontrée ci-dessus
donne donc enfin l'égalité voulue entre
$t_j|_{V_{jj'l}}$ et $t_{j'}|_{V_{jj'l}}$.
Nous obtenons alors l'existence d'un élément $t \in \mathcal{F}(U)$
par la propriété $(**)$ pour $(\mathcal{V}, \mathcal{V}_{jj'})$.

\medskip\noindent
Une petite difficulté subsiste. Nous savons que $t$ se restreint en $t_j$ sur $V_j$,
mais nous ne savons pas encore que $t$ se restreint en $s_i$ sur $U_i$. Pour conclure,
remarquons que les ensembles $U_i \cap V_j$, $j \in J$, recouvrent $U_i$. Par conséquent,
les ensembles $U_{i \alpha(j) k} \cap V_j$, $j\in J$, $k \in I_{i\alpha(j)}$,
recouvrent eux aussi $U_i$. Nous laissons au lecteur le soin de voir que $t$ et $s_i$ induisent
la même section de $\mathcal{F}$ sur tout $W \in \mathcal{B}$
contenu dans l'un des ouverts
$U_{i \alpha(j) k} \cap V_j$, $j\in J$, $k \in I_{i\alpha(j)}$.
La partie unicité établie ci-dessus permet donc de conclure.
\end{proof}

\begin{lemma}
\label{sheaves-lemma-cofinal-systems-coverings-standard-case}
Soit $X$ un espace topologique. Soit $\mathcal{B}$ une base de la
topologie de $X$. Supposons que, pour tout triplet $U, U', U'' \in \mathcal{B}$
tel que $U' \subset U$ et $U'' \subset U$, on ait $U' \cap U'' \in \mathcal{B}$.
Pour chaque $U \in \mathcal{B}$, soit $C(U) \subset \text{Cov}_\mathcal{B}(U)$
un système cofinal.
Soit $\mathcal{F}$ un préfaisceau d'ensembles sur $\mathcal{B}$.
Les conditions suivantes sont équivalentes :
\begin{enumerate}
\item Le préfaisceau $\mathcal{F}$ est un faisceau sur $\mathcal{B}$.
\item Pour tout $U \in \mathcal{B}$, tout recouvrement
$\mathcal{U} : U = \bigcup U_i$ appartenant à $C(U)$ et toute
famille de sections $s_i \in \mathcal{F}(U_i)$ telle
que $s_i|_{U_i \cap U_j} = s_j|_{U_i \cap U_j}$, il
existe une unique section $s \in \mathcal{F}(U)$ qui
se restreint en $s_i$ sur $U_i$.
\end{enumerate}
\end{lemma}

\begin{proof}
C'est une reformulation du
Lemme \ref{sheaves-lemma-cofinal-systems-coverings} ci-dessus
dans le cas particulier où chacun des recouvrements $\mathcal{U}_{ij}$
est formé d'un seul élément. Ce cas est d'ailleurs beaucoup
plus simple et constitue un exercice facile que l'on peut résoudre directement.
\end{proof}

\begin{lemma}
\label{sheaves-lemma-condition-star-sections}
Soit $X$ un espace topologique.
Soit $\mathcal{B}$ une base de la topologie de $X$.
Soit $U \in \mathcal{B}$.
Soit $\mathcal{F}$ un faisceau d'ensembles sur $\mathcal{B}$.
L'application
$$
\mathcal{F}(U) \to \prod\nolimits_{x \in U} \mathcal{F}_x
$$
identifie $\mathcal{F}(U)$ aux éléments $(s_x)_{x\in U}$
qui vérifient la propriété
\begin{itemize}
\item[(*)] Pour tout $x \in U$, il existe $V \in \mathcal{B}$
tel que $x \in V \subset U$, ainsi qu'une section $\sigma \in \mathcal{F}(V)$,
de sorte que, pour tout $y \in V$, on ait $s_y = (V, \sigma)$ dans $\mathcal{F}_y$.
\end{itemize}
\end{lemma}

\begin{proof}
Remarquons d'abord que l'application
$\mathcal{F}(U) \to \prod\nolimits_{x \in U} \mathcal{F}_x$
est injective par l'unicité dans la condition de faisceau
de la Définition \ref{sheaves-definition-sheaf-basis}. Soit $(s_x)$
un élément quelconque du membre de droite qui vérifie $(*)$.
Cela signifie clairement que l'on peut trouver un recouvrement $U = \bigcup U_i$,
$U_i \in \mathcal{B}$, tel que $(s_x)_{x \in U_i}$ provienne
de certains $\sigma_i \in \mathcal{F}(U_i)$. Pour tout $y \in U_i \cap U_j$,
les sections $\sigma_i$ et $\sigma_j$ coïncident dans la fibre
$\mathcal{F}_y$. Il existe donc un élément $V_{ijy} \in \mathcal{B}$,
$y \in V_{ijy}$, tel que $\sigma_i|_{V_{ijy}} = \sigma_j|_{V_{ijy}}$.
La condition de faisceau $(**)$ de la Définition \ref{sheaves-definition-sheaf-basis}
s'applique donc au système des $\sigma_i$, et nous obtenons une section
$s \in \mathcal{F}(U)$ qui possède la propriété voulue.
\end{proof}

\noindent
Soit $X$ un espace topologique.
Soit $\mathcal{B}$ une base de la topologie de $X$.
Il existe un foncteur de restriction naturel de la catégorie
des faisceaux d'ensembles sur $X$ vers la catégorie des faisceaux
d'ensembles sur $\mathcal{B}$. Ce foncteur est en fait une équivalence
de catégories. Concrètement, cela signifie ce qui suit.

\begin{lemma}
\label{sheaves-lemma-extend-off-basis}
Soit $X$ un espace topologique.
Soit $\mathcal{B}$ une base de la topologie de $X$.
Soit $\mathcal{F}$ un faisceau d'ensembles sur $\mathcal{B}$.
Il existe un unique faisceau d'ensembles $\mathcal{F}^{ext}$
sur $X$ tel que $\mathcal{F}^{ext}(U) = \mathcal{F}(U)$
pour tout $U \in \mathcal{B}$, de manière compatible avec les applications
de restriction.
\end{lemma}

\begin{proof}
Nous construisons d'abord un préfaisceau $\mathcal{F}^{ext}$ possédant la
propriété voulue. Pour un ouvert quelconque $U \subset X$, nous
définissons $\mathcal{F}^{ext}(U)$ comme l'ensemble des éléments
$(s_x)_{x \in U}$ qui vérifient la condition $(*)$ du
Lemme \ref{sheaves-lemma-condition-star-sections}.
Il est clair qu'il existe des applications de restriction
qui font de $\mathcal{F}^{ext}$ un préfaisceau d'ensembles.
En outre, le Lemme \ref{sheaves-lemma-condition-star-sections} montre que
$\mathcal{F}(U) = \mathcal{F}^{ext}(U)$
dès que $U$ est un élément de la base $\mathcal{B}$.
Pour voir que $\mathcal{F}^{ext}$ est un faisceau, on peut
raisonner comme dans la preuve du Lemme \ref{sheaves-lemma-sheafification-sheaf}.
\end{proof}

\noindent
Remarquons que l'on a
$$
\mathcal{F}_x = \mathcal{F}_x^{ext}
$$
dans la situation du lemme. En effet, la
famille des éléments de $\mathcal{B}$ qui contiennent
$x$ forme un système fondamental de voisinages ouverts de $x$.

\begin{lemma}
\label{sheaves-lemma-restrict-basis-equivalence}
Soit $X$ un espace topologique.
Soit $\mathcal{B}$ une base de la topologie de $X$.
Notons $\Sh(\mathcal{B})$ la catégorie des
faisceaux sur $\mathcal{B}$.
Il existe une équivalence de catégories
$$
\Sh(X) \longrightarrow \Sh(\mathcal{B})
$$
qui associe à un faisceau sur $X$ sa restriction aux
membres de $\mathcal{B}$.
\end{lemma}

\begin{proof}
Le foncteur inverse est donné par le Lemme \ref{sheaves-lemma-extend-off-basis} ci-dessus.
La vérification des fonctorialités évidentes est laissée au
lecteur.
\end{proof}

\noindent
Ceci termine l'étude des faisceaux d'ensembles sur
une base $\mathcal{B}$. Soit $(\mathcal{C}, F)$ une
espèce de structure algébrique. Pour terminer cette section,
nous souhaitons signaler que les constructions
ci-dessus valent pour les faisceaux à valeurs dans $\mathcal{C}$.
Définissons brièvement les notions pertinentes.

\begin{definition}
\label{sheaves-definition-sheaf-structures-basis}
Soit $X$ un espace topologique. Soit $\mathcal{B}$ une
base de la topologie de $X$. Soit $(\mathcal{C}, F)$
une espèce de structure algébrique.
\begin{enumerate}
\item Un {\it préfaisceau $\mathcal{F}$ à valeurs dans $\mathcal{C}$
sur $\mathcal{B}$} est une règle qui associe à chaque
$U \in \mathcal{B}$ un objet
$\mathcal{F}(U)$ de $\mathcal{C}$ et à chaque inclusion $V \subset U$
d'éléments de $\mathcal{B}$ un morphisme
$\rho^U_V : \mathcal{F}(U) \to \mathcal{F}(V)$ dans $\mathcal{C}$ tel que
$\rho^U_U = \text{id}_{\mathcal{F}(U)}$ pour tout $U \in \mathcal{B}$ et que,
si $W \subset V \subset U$ dans $\mathcal{B}$, on ait
$\rho^U_W = \rho^V_W \circ \rho ^U_V$.
\item Un {\it morphisme $\varphi : \mathcal{F} \to \mathcal{G}$
de préfaisceaux à valeurs dans $\mathcal{C}$
sur $\mathcal{B}$} est une règle qui associe à chaque
élément $U \in \mathcal{B}$ un morphisme de
structures algébriques $\varphi : \mathcal{F}(U) \to \mathcal{G}(U)$
compatible avec les applications de restriction.
\item Étant donné un préfaisceau $\mathcal{F}$ à valeurs dans $\mathcal{C}$
sur $\mathcal{B}$, nous appelons $U \mapsto F(\mathcal{F}(U))$ le
préfaisceau d'ensembles sous-jacent.
\item Un {\it faisceau $\mathcal{F}$ à valeurs dans $\mathcal{C}$
sur $\mathcal{B}$} est un préfaisceau à valeurs dans $\mathcal{C}$
sur $\mathcal{B}$ dont le préfaisceau d'ensembles sous-jacent est un faisceau.
\end{enumerate}
\end{definition}

\noindent
Nous pouvons maintenant définir la {\it fibre} en $x \in X$
d'un préfaisceau à valeurs dans $\mathcal{C}$ sur $\mathcal{B}$
comme la limite inductive filtrante
$$
\mathcal{F}_x = \colim_{U\in \mathcal{B}, x\in U} \mathcal{F}(U).
$$
Elle existe comme objet de $\mathcal{C}$
en vertu de nos hypothèses sur $\mathcal{C}$.
De plus, l'ensemble sous-jacent à $\mathcal{F}_x$
est la fibre du préfaisceau d'ensembles sous-jacent sur $\mathcal{B}$.

\medskip\noindent
Remarquons que les Lemmes \ref{sheaves-lemma-cofinal-systems-coverings},
\ref{sheaves-lemma-cofinal-systems-coverings-standard-case} et
\ref{sheaves-lemma-condition-star-sections} portent sur la propriété de faisceau,
que nous avons définie en termes du préfaisceau
d'ensembles associé. Ils se généralisent donc sans changement à la notion
de préfaisceau à valeurs dans $\mathcal{C}$. L'analogue du
Lemme \ref{sheaves-lemma-extend-off-basis} demande quelques précautions. Le voici.

\begin{lemma}
\label{sheaves-lemma-extend-off-basis-structures}
Soit $X$ un espace topologique. Soit $(\mathcal{C}, F)$
une espèce de structure algébrique.
Soit $\mathcal{B}$ une base de la topologie de $X$.
Soit $\mathcal{F}$ un faisceau à valeurs dans $\mathcal{C}$
sur $\mathcal{B}$.
Il existe un unique faisceau $\mathcal{F}^{ext}$ à valeurs dans $\mathcal{C}$
sur $X$ tel que $\mathcal{F}^{ext}(U) = \mathcal{F}(U)$
pour tout $U \in \mathcal{B}$, de manière compatible avec les applications
de restriction.
\end{lemma}

\begin{proof}
En vertu des conditions imposées au couple $(\mathcal{C}, F)$, il
suffit de construire un préfaisceau $\mathcal{F}^{ext}$
qui possède la propriété voulue au niveau des préfaisceaux
d'ensembles sous-jacents. Notre première tâche consiste donc à construire
un objet convenable $\mathcal{F}^{ext}(U)$ pour tout ouvert $U \subset X$.
Nous pourrions le faire en imitant le
Lemme \ref{sheaves-lemma-diagram-fibre-product} dans le cadre
des préfaisceaux sur $\mathcal{B}$. Voici toutefois une méthode légèrement différente
(mais essentiellement équivalente) :
définissons-le comme la limite inductive filtrante
$$
\mathcal{F}^{ext}(U)
:=
\colim_\mathcal{U} FIB(\mathcal{U})
$$
sur tous les recouvrements
$\mathcal{U} : U = \bigcup_{i\in I} U_i$ par des $U_i \in \mathcal{B}$
du produit fibré
$$
\xymatrix{
FIB(\mathcal{U}) \ar[r] \ar[d] &
\prod\nolimits_{x\in U} \mathcal{F}_x \ar[d] \\
\prod\nolimits_{i\in I} \mathcal{F}(U_i) \ar[r] &
\prod\nolimits_{i \in I} \prod\nolimits_{x\in U_i} \mathcal{F}_x
}
$$
Par les arguments usuels, voir le Lemme \ref{sheaves-lemma-image-contained-in}
et l'Exemple \ref{sheaves-example-application-lemma-image-contained-in},
il suffit de montrer que, sur les ensembles sous-jacents, cette construction
coïncide avec la définition ci-dessus au moyen de $(**)$.
Les détails sont laissés au lecteur.
\end{proof}

\noindent
Remarquons que l'on a
$$
\mathcal{F}_x = \mathcal{F}_x^{ext}
$$
comme objets de $\mathcal{C}$
dans la situation du lemme. En effet, la
famille des éléments de $\mathcal{B}$ qui contiennent
$x$ forme un système fondamental de voisinages ouverts de $x$.

\begin{lemma}
\label{sheaves-lemma-restrict-basis-equivalence-structures}
Soit $X$ un espace topologique.
Soit $\mathcal{B}$ une base de la topologie de $X$.
Soit $(\mathcal{C}, F)$ une espèce de structure algébrique.
Notons $\Sh(\mathcal{B}, \mathcal{C})$ la catégorie des
faisceaux à valeurs dans $\mathcal{C}$ sur $\mathcal{B}$.
Il existe une équivalence de catégories
$$
\Sh(X, \mathcal{C})
\longrightarrow
\Sh(\mathcal{B}, \mathcal{C})
$$
qui associe à un faisceau sur $X$ sa restriction aux
membres de $\mathcal{B}$.
\end{lemma}

\begin{proof}
Le foncteur inverse est donné par le
Lemme \ref{sheaves-lemma-extend-off-basis-structures} ci-dessus.
La vérification des fonctorialités évidentes est laissée au
lecteur.
\end{proof}

\noindent
Enfin, nous abordons le cas des (pré)faisceaux de modules
sur une base. Nous utiliserons le fait élémentaire que la catégorie
des préfaisceaux d'ensembles sur une base admet des produits et que
ceux-ci s'obtiennent en prenant les produits des valeurs sur les
éléments de la base.

\begin{definition}
\label{sheaves-definition-sheaf-modules-basis}
Soit $X$ un espace topologique. Soit $\mathcal{B}$ une
base de la topologie de $X$. Soit $\mathcal{O}$
un préfaisceau d'anneaux sur $\mathcal{B}$.
\begin{enumerate}
\item Un {\it préfaisceau de $\mathcal{O}$-modules $\mathcal{F}$
sur $\mathcal{B}$} est un préfaisceau de groupes abéliens sur
$\mathcal{B}$ muni d'un morphisme de préfaisceaux
d'ensembles $\mathcal{O} \times \mathcal{F} \to \mathcal{F}$
tel que, pour tout $U \in \mathcal{B}$, l'application
$\mathcal{O}(U) \times \mathcal{F}(U) \to \mathcal{F}(U)$
munisse le groupe $\mathcal{F}(U)$ d'une structure de $\mathcal{O}(U)$-module.
\item Un {\it morphisme $\varphi : \mathcal{F} \to \mathcal{G}$
de préfaisceaux de $\mathcal{O}$-modules sur $\mathcal{B}$}
est un morphisme de préfaisceaux abéliens sur $\mathcal{B}$
qui induit un homomorphisme de $\mathcal{O}(U)$-modules
$\mathcal{F}(U) \to \mathcal{G}(U)$ pour tout $U \in \mathcal{B}$.
\item Supposons que $\mathcal{O}$ soit un faisceau d'anneaux
sur $\mathcal{B}$. Un {\it faisceau $\mathcal{F}$ de $\mathcal{O}$-modules
sur $\mathcal{B}$} est un préfaisceau de $\mathcal{O}$-modules
sur $\mathcal{B}$ dont le préfaisceau de groupes abéliens sous-jacent
est un faisceau.
\end{enumerate}
\end{definition}

\noindent
Nous pouvons définir la {\it fibre} en $x \in X$
d'un préfaisceau de $\mathcal{O}$-modules sur $\mathcal{B}$
comme la limite inductive filtrante
$$
\mathcal{F}_x = \colim_{U\in \mathcal{B}, x\in U} \mathcal{F}(U).
$$
C'est un $\mathcal{O}_x$-module.

\medskip\noindent
Remarquons que les Lemmes \ref{sheaves-lemma-cofinal-systems-coverings},
\ref{sheaves-lemma-cofinal-systems-coverings-standard-case} et
\ref{sheaves-lemma-condition-star-sections} portent sur la propriété de faisceau,
que nous avons définie en termes du préfaisceau
d'ensembles associé. Ils se généralisent donc sans changement à la notion
de préfaisceau de $\mathcal{O}$-modules. L'analogue du
Lemme \ref{sheaves-lemma-extend-off-basis} est le suivant.

\begin{lemma}
\label{sheaves-lemma-extend-off-basis-module}
Soit $X$ un espace topologique.
Soit $\mathcal{B}$ une base de la topologie de $X$.
Soit $\mathcal{O}$ un faisceau d'anneaux sur $\mathcal{B}$.
Soit $\mathcal{F}$ un faisceau de $\mathcal{O}$-modules
sur $\mathcal{B}$. Soit $\mathcal{O}^{ext}$ le faisceau
d'anneaux sur $X$ qui prolonge $\mathcal{O}$, et soit
$\mathcal{F}^{ext}$ le faisceau abélien sur $X$ qui prolonge
$\mathcal{F}$ ; voir le Lemme \ref{sheaves-lemma-extend-off-basis-structures}.
Il existe un morphisme canonique
$$
\mathcal{O}^{ext} \times \mathcal{F}^{ext}
\longrightarrow
\mathcal{F}^{ext}
$$
qui coïncide avec le morphisme donné sur les éléments de $\mathcal{B}$
et qui munit $\mathcal{F}^{ext}$ d'une structure
de $\mathcal{O}^{ext}$-module.
\end{lemma}

\begin{proof}
Il suffit de construire le morphisme de multiplication
au niveau des préfaisceaux d'ensembles. La manière la plus simple
de le voir est peut-être de démontrer directement que, si
$(f_x)_{x \in U}$, $f_x \in \mathcal{O}_x$
et
$(m_x)_{x \in U}$, $m_x \in \mathcal{F}_x$
vérifient $(*)$, alors l'élément
$(f_xm_x)_{x \in U}$ vérifie lui aussi $(*)$.
On obtient alors le résultat voulu, car, dans la preuve
du Lemme \ref{sheaves-lemma-extend-off-basis}, nous construisons le prolongement
à l'aide de familles d'éléments de fibres qui vérifient $(*)$.
\end{proof}

\noindent
Remarquons que l'on a
$$
\mathcal{F}_x = \mathcal{F}_x^{ext}
$$
comme $\mathcal{O}_x$-modules dans la situation du lemme.
Cela tient au fait que la famille des éléments de $\mathcal{B}$ qui contiennent
$x$ forme un système fondamental de voisinages ouverts de $x$, ou
simplement au fait que l'assertion est vraie sur les ensembles sous-jacents.

\begin{lemma}
\label{sheaves-lemma-restrict-basis-equivalence-modules}
Soit $X$ un espace topologique.
Soit $\mathcal{B}$ une base de la topologie de $X$.
Soit $\mathcal{O}$ un faisceau d'anneaux sur $X$.
Notons $\textit{Mod}(\mathcal{O}|_\mathcal{B})$ la catégorie des
faisceaux de $\mathcal{O}|_\mathcal{B}$-modules sur $\mathcal{B}$.
Il existe une équivalence de catégories
$$
\textit{Mod}(\mathcal{O})
\longrightarrow
\textit{Mod}(\mathcal{O}|_\mathcal{B})
$$
qui associe à un faisceau de $\mathcal{O}$-modules sur $X$ sa restriction aux
membres de $\mathcal{B}$.
\end{lemma}

\begin{proof}
Le foncteur inverse est donné par le
Lemme \ref{sheaves-lemma-extend-off-basis-module} ci-dessus.
La vérification des fonctorialités évidentes est laissée au lecteur.
\end{proof}

\noindent
Enfin, nous étudions le lien de cette théorie avec les
applications continues. C'est maintenant très facile grâce au travail
accompli ci-dessus. Traitons d'abord le cas où une base de l'espace
d'arrivée est donnée.

\begin{lemma}
\label{sheaves-lemma-f-map-basis-below-structures}
Soit $f : X \to Y$ une application continue d'espaces topologiques.
Soit $(\mathcal{C}, F)$ une espèce de structure algébrique.
Soit $\mathcal{F}$ un faisceau à valeurs dans $\mathcal{C}$ sur $X$.
Soit $\mathcal{G}$ un faisceau à valeurs dans $\mathcal{C}$ sur $Y$.
Soit $\mathcal{B}$ une base de la topologie de $Y$.
Supposons donné, pour tout $V \in \mathcal{B}$, un morphisme
$$
\varphi_V :
\mathcal{G}(V)
\longrightarrow
\mathcal{F}(f^{-1}V)
$$
de $\mathcal{C}$ compatible avec les applications de restriction.
Il existe alors un unique $f$-morphisme (voir la Définition \ref{sheaves-definition-f-map}
et la discussion des $f$-morphismes dans la
section \ref{sheaves-section-presheaves-structures-functorial})
$\varphi : \mathcal{G} \to \mathcal{F}$
qui redonne $\varphi_V$ pour $V \in \mathcal{B}$.
\end{lemma}

\begin{proof}
C'est immédiat, car la famille de morphismes
revient à un morphisme entre les restrictions
de $\mathcal{G}$ et de $f_*\mathcal{F}$ à $\mathcal{B}$.
Par le Lemme \ref{sheaves-lemma-restrict-basis-equivalence-structures},
cela revient à donner un morphisme de $\mathcal{G}$
vers $f_*\mathcal{F}$, ce qui, par le Lemme \ref{sheaves-lemma-f-map},
équivaut à un $f$-morphisme. Voir aussi le
Lemme \ref{sheaves-lemma-f-map-sets-algebraic-structures}
et la discussion qui le précède
pour le traitement du cas des faisceaux de structures algébriques.
\end{proof}

\noindent
Voici l'analogue pour les espaces annelés.

\begin{lemma}
\label{sheaves-lemma-f-map-basis-below-modules}
Soit $(f, f^\sharp) : (X, \mathcal{O}_X) \to (Y, \mathcal{O}_Y)$
un morphisme d'espaces annelés.
Soit $\mathcal{F}$ un faisceau de $\mathcal{O}_X$-modules.
Soit $\mathcal{G}$ un faisceau de $\mathcal{O}_Y$-modules.
Soit $\mathcal{B}$ une base de la topologie de $Y$.
Supposons donné, pour tout $V \in \mathcal{B}$, un
morphisme de $\mathcal{O}_Y(V)$-modules
$$
\varphi_V :
\mathcal{G}(V)
\longrightarrow
\mathcal{F}(f^{-1}V)
$$
(où $\mathcal{F}(f^{-1}V)$ est muni d'une structure de module par restriction des scalaires le long de
$f^\sharp_V : \mathcal{O}_Y(V) \to \mathcal{O}_X(f^{-1}V)$)
compatible avec les applications de restriction.
Il existe alors un unique $f$-morphisme (voir la discussion des $f$-morphismes dans la
section \ref{sheaves-section-ringed-spaces-functoriality-modules})
$\varphi : \mathcal{G} \to \mathcal{F}$
qui redonne $\varphi_V$ pour $V \in \mathcal{B}$.
\end{lemma}

\begin{proof}
La preuve est la même que celle du lemme correspondant
pour les faisceaux de structures algébriques ci-dessus.
\end{proof}

\begin{lemma}
\label{sheaves-lemma-f-map-basis-above-and-below-structures}
Soit $f : X \to Y$ une application continue d'espaces topologiques.
Soit $(\mathcal{C}, F)$ une espèce de structure algébrique.
Soit $\mathcal{F}$ un faisceau à valeurs dans $\mathcal{C}$ sur $X$.
Soit $\mathcal{G}$ un faisceau à valeurs dans $\mathcal{C}$ sur $Y$.
Soit $\mathcal{B}_Y$ une base de la topologie de $Y$.
Soit $\mathcal{B}_X$ une base de la topologie de $X$.
Supposons donné, pour tous $V \in \mathcal{B}_Y$ et
$U \in \mathcal{B}_X$ tels que $f(U) \subset V$, un morphisme
$$
\varphi_V^U :
\mathcal{G}(V)
\longrightarrow
\mathcal{F}(U)
$$
de $\mathcal{C}$ compatible avec les applications de restriction.
Il existe alors un unique $f$-morphisme (voir
la Définition \ref{sheaves-definition-f-map} et la discussion
des $f$-morphismes dans la section \ref{sheaves-section-presheaves-structures-functorial})
$\varphi : \mathcal{G} \to \mathcal{F}$
qui redonne $\varphi_V^U$ comme la composée
$$
\mathcal{G}(V) \xrightarrow{\varphi_V}
\mathcal{F}(f^{-1}(V)) \xrightarrow{\text{rest.}}
\mathcal{F}(U)
$$
pour toute paire $(U, V)$ comme ci-dessus.
\end{lemma}

\begin{proof}
Démontrons d'abord ce résultat pour les faisceaux d'ensembles.
Fixons un ouvert $V \subset Y$. Choisissons $s \in \mathcal{G}(V)$.
Nous allons construire un élément
$\varphi_V(s) \in \mathcal{F}(f^{-1}V)$.
Nous pouvons définir une valeur $\varphi(s)_x$ dans la fibre $\mathcal{F}_x$
pour tout $x \in f^{-1}V$ en choisissant un $U \in \mathcal{B}_X$
tel que $x \in U \subset f^{-1}V$ et en prenant pour $\varphi(s)_x$
la classe d'équivalence de $(U, \varphi_V^U(s))$
dans la fibre. Il est clair que la famille $(\varphi(s)_x)_{x \in f^{-1}V}$
vérifie la condition $(*)$, car les morphismes $\varphi_V^U$,
lorsque $U$ varie,
sont compatibles avec les restrictions dans le faisceau $\mathcal{F}$.
Ainsi, par la preuve du Lemme \ref{sheaves-lemma-extend-off-basis},
nous voyons que $(\varphi(s)_x)_{x \in f^{-1}V}$ correspond
à un unique élément $\varphi_V(s)$ de $\mathcal{F}(f^{-1}V)$.
Nous avons donc défini une application d'ensembles
$\varphi_V : \mathcal{G}(V) \to \mathcal{F}(f^{-1}V)$.
La compatibilité entre $\varphi_V$ et $\varphi_V^U$
résulte du Lemme \ref{sheaves-lemma-condition-star-sections}.

\medskip\noindent
Nous laissons au lecteur le soin de montrer que la construction
de $\varphi_V$ est compatible avec les applications de restriction lorsque
$V \in \mathcal{B}_Y$ varie. Nous pouvons donc appliquer le Lemme
\ref{sheaves-lemma-f-map-basis-below-structures} ci-dessus pour les
``recoller'' et obtenir le $f$-morphisme voulu.

\medskip\noindent
Enfin, remarquons que le morphisme de faisceaux d'ensembles ainsi construit
possède la propriété suivante : le morphisme sur les fibres
$$
\mathcal{G}_{f(x)} \longrightarrow \mathcal{F}_x
$$
est la limite inductive du système de morphismes $\varphi_V^U$ lorsque
$V \in \mathcal{B}_Y$ parcourt les éléments qui
contiennent $f(x)$ et $U \in \mathcal{B}_X$ parcourt ceux qui
contiennent $x$. En particulier, si $\mathcal{G}$ et $\mathcal{F}$
sont les faisceaux d'ensembles sous-jacents à des faisceaux de structures algébriques,
nous voyons que le morphisme sur les fibres est un morphisme de structures
algébriques. Le morphisme correspondant de
faisceaux d'ensembles sous-jacents $f^{-1}\mathcal{G} \to \mathcal{F}$
satisfait donc les hypothèses du
Lemme \ref{sheaves-lemma-f-map-sets-algebraic-structures}.
Nous en concluons que $f^{-1}\mathcal{G} \to \mathcal{F}$
est un morphisme de faisceaux à valeurs dans $\mathcal{C}$.
Par adjonction, cela signifie que $\varphi$ est
un $f$-morphisme de faisceaux de structures algébriques.
\end{proof}

\begin{lemma}
\label{sheaves-lemma-f-map-basis-above-and-below-modules}
Soit $(f, f^\sharp) : (X, \mathcal{O}_X) \to (Y, \mathcal{O}_Y)$
un morphisme d'espaces annelés.
Soit $\mathcal{F}$ un faisceau de $\mathcal{O}_X$-modules.
Soit $\mathcal{G}$ un faisceau de $\mathcal{O}_Y$-modules.
Soit $\mathcal{B}_Y$ une base de la topologie de $Y$.
Soit $\mathcal{B}_X$ une base de la topologie de $X$.
Supposons donné, pour tous $V \in \mathcal{B}_Y$ et
$U \in \mathcal{B}_X$ tels que $f(U) \subset V$, un
morphisme de $\mathcal{O}_Y(V)$-modules
$$
\varphi_V^U :
\mathcal{G}(V)
\longrightarrow
\mathcal{F}(U)
$$
compatible avec les applications de restriction. Ici, la
structure de $\mathcal{O}_Y(V)$-module sur $\mathcal{F}(U)$
s'obtient par restriction des scalaires à partir de la structure de $\mathcal{O}_X(U)$-module
le long du morphisme $f^\sharp_V : \mathcal{O}_Y(V)
\to \mathcal{O}_X(f^{-1}V) \to \mathcal{O}_X(U)$.
Il existe alors un unique $f$-morphisme de faisceaux de modules (voir
la Définition \ref{sheaves-definition-f-map} et la discussion
des $f$-morphismes dans la section \ref{sheaves-section-ringed-spaces-functoriality-modules})
$\varphi : \mathcal{G} \to \mathcal{F}$
qui redonne $\varphi_V^U$ comme la composée
$$
\mathcal{G}(V) \xrightarrow{\varphi_V}
\mathcal{F}(f^{-1}(V)) \xrightarrow{\text{rest.}}
\mathcal{F}(U)
$$
pour toute paire $(U, V)$ comme ci-dessus.
\end{lemma}

\begin{proof}
La preuve est analogue à la précédente et est omise.
\end{proof}

\section{Immersions ouvertes et (pré)faisceaux}
\label{sheaves-section-open-immersions}

\noindent
Soit $X$ un espace topologique.
Soit $j : U \to X$ l'inclusion d'un sous-ensemble ouvert $U$ dans $X$.
Dans la section \ref{sheaves-section-presheaves-functorial}, nous avons défini
des foncteurs $j_*$ et $j^{-1}$ tels que $j_*$ soit adjoint à droite de
$j^{-1}$. Il se trouve que, pour une immersion ouverte, $j^{-1}$ possède
un adjoint à gauche, que nous noterons $j_!$. Commençons par remarquer que
$j^{-1}$ admet une description particulièrement simple dans le cas d'une
immersion ouverte.

\begin{lemma}
\label{sheaves-lemma-j-pullback}
Soit $X$ un espace topologique.
Soit $j : U \to X$ l'inclusion d'un sous-ensemble ouvert $U$ dans $X$.
\begin{enumerate}
\item Soit $\mathcal{G}$ un préfaisceau d'ensembles sur $X$.
Le préfaisceau $j_p\mathcal{G}$
(voir la section \ref{sheaves-section-presheaves-functorial}) est donné par la règle
$V \mapsto \mathcal{G}(V)$ pour tout ouvert $V \subset U$.
\item Soit $\mathcal{G}$ un faisceau d'ensembles sur $X$.
Le faisceau $j^{-1}\mathcal{G}$ est donné par la règle
$V \mapsto \mathcal{G}(V)$ pour tout ouvert $V \subset U$.
\item Pour tout point $u \in U$ et tout faisceau $\mathcal{G}$ sur $X$,
on dispose d'une identification canonique des fibres
$$
j^{-1}\mathcal{G}_u = (\mathcal{G}|_U)_u = \mathcal{G}_u.
$$
\item Dans la catégorie des préfaisceaux sur $U$, nous avons $j_pj_* = \text{id}$.
\item Dans la catégorie des faisceaux sur $U$, nous avons $j^{-1}j_* = \text{id}$.
\end{enumerate}
La même description vaut pour les (pré)faisceaux de groupes abéliens,
les (pré)faisceaux de structures algébriques et les (pré)faisceaux de modules.
\end{lemma}

\begin{proof}
La limite inductive dans la définition de $j_p\mathcal{G}(V)$
porte sur la collection de tous les ouverts $W \subset X$ tels que $V \subset W$,
ordonnée par l'inclusion opposée.
Elle possède donc un plus grand élément, à savoir $V$. Cela prouve (1).
Et (2) en résulte, car la règle $V \mapsto \mathcal{G}(V)$
pour tout ouvert $V \subset U$ est manifestement un faisceau si $\mathcal{G}$ est un
faisceau. L'assertion (3) résulte de (2), car la collection
des voisinages ouverts de $u$ contenus dans $U$ est cofinale
dans la collection de tous les voisinages ouverts de $u$ dans $X$.
Les parties (4) et (5) résultent du calcul
$j^{-1}j_*\mathcal{F}(V) = j_*\mathcal{F}(V) = \mathcal{F}(V)$.

\medskip\noindent
Les mêmes arguments s'appliquent mot pour mot aux (pré)faisceaux de groupes abéliens
et aux (pré)faisceaux de structures algébriques.
\end{proof}

\begin{definition}
\label{sheaves-definition-restriction}
Soit $X$ un espace topologique.
Soit $j : U \to X$ l'inclusion d'un sous-ensemble ouvert.
\begin{enumerate}
\item Soit $\mathcal{G}$ un préfaisceau d'ensembles, de groupes abéliens ou de
structures algébriques sur $X$. Le préfaisceau $j_p\mathcal{G}$ décrit
dans le lemme \ref{sheaves-lemma-j-pullback} est appelé
la {\it restriction de $\mathcal{G}$ à $U$} et noté $\mathcal{G}|_U$.
\item Soit $\mathcal{G}$ un faisceau d'ensembles sur $X$, ou un faisceau de groupes
abéliens ou de structures algébriques sur $X$. Le faisceau $j^{-1}\mathcal{G}$ est appelé
la {\it restriction de $\mathcal{G}$ à $U$} et noté $\mathcal{G}|_U$.
\item Si $(X, \mathcal{O})$ est un espace annelé, le couple
$(U, \mathcal{O}|_U)$ est appelé le
{\it sous-espace ouvert de $(X, \mathcal{O})$ associé à $U$}.
\item Si $\mathcal{G}$ est un préfaisceau de $\mathcal{O}$-modules,
alors $\mathcal{G}|_U$, muni de l'application de multiplication
$\mathcal{O}|_U \times \mathcal{G}|_U \to \mathcal{G}|_U$
(voir le lemme \ref{sheaves-lemma-pullback-module}),
est appelé la {\it restriction de $\mathcal{G}$ à $U$}.
\end{enumerate}
\end{definition}

\noindent
Nous laissons au lecteur la définition de la restriction des préfaisceaux
de modules. Dans cette section, nous allons donc
étudier un adjoint à gauche du foncteur de restriction.
Voici la définition dans le cas des (pré)faisceaux
d'ensembles.

\begin{definition}
\label{sheaves-definition-j-shriek}
Soit $X$ un espace topologique.
Soit $j : U \to X$ l'inclusion d'un sous-ensemble ouvert.
\begin{enumerate}
\item Soit $\mathcal{F}$ un préfaisceau d'ensembles sur $U$. Nous définissons
le {\it prolongement de $\mathcal{F}$ par l'ensemble vide $j_{p!}\mathcal{F}$}
comme le préfaisceau d'ensembles sur $X$ défini par la règle
$$
j_{p!}\mathcal{F}(V) =
\left\{
\begin{matrix}
\emptyset & \text{si} & V \not \subset U \\
\mathcal{F}(V) & \text{si} & V \subset U
\end{matrix}
\right.
$$
avec les applications de restriction évidentes.
\item Soit $\mathcal{F}$ un faisceau d'ensembles sur $U$. Nous définissons
le {\it prolongement de $\mathcal{F}$ par l'ensemble vide $j_!\mathcal{F}$}
comme le faisceau associé au préfaisceau $j_{p!}\mathcal{F}$.
\end{enumerate}
\end{definition}

\begin{lemma}
\label{sheaves-lemma-j-shriek}
Soit $X$ un espace topologique.
Soit $j : U \to X$ l'inclusion d'un sous-ensemble ouvert.
\begin{enumerate}
\item Le foncteur $j_{p!}$ est adjoint à gauche du
foncteur de restriction $j_p$ (voir le lemme \ref{sheaves-lemma-j-pullback}).
\item Le foncteur $j_!$ est adjoint à gauche de la restriction ;
plus précisément,
$$
\Mor_{\Sh(X)}(j_!\mathcal{F}, \mathcal{G})
=
\Mor_{\Sh(U)}(\mathcal{F}, j^{-1}\mathcal{G})
=
\Mor_{\Sh(U)}(\mathcal{F}, \mathcal{G}|_U)
$$
bifonctoriellement en $\mathcal{F}$ et $\mathcal{G}$.
\item Soit $\mathcal{F}$ un faisceau d'ensembles sur $U$.
Les fibres du faisceau $j_!\mathcal{F}$ sont décrites
comme suit :
$$
j_{!}\mathcal{F}_x =
\left\{
\begin{matrix}
\emptyset & \text{si} & x \not \in U \\
\mathcal{F}_x & \text{si} & x \in U
\end{matrix}
\right.
$$
\item Dans la catégorie des préfaisceaux sur $U$, nous avons $j_pj_{p!} = \text{id}$.
\item Dans la catégorie des faisceaux sur $U$, nous avons $j^{-1}j_! = \text{id}$.
\end{enumerate}
\end{lemma}

\begin{proof}
Pour définir un morphisme de $j_{p!}\mathcal{F}$ vers $\mathcal{G}$,
il suffit de définir des applications $\mathcal{F}(V) \to \mathcal{G}(V)$
lorsque $V \subset U$, compatibles avec les applications
de restriction. Et, d'après le lemme \ref{sheaves-lemma-j-pullback},
la même description vaut pour les morphismes
$\mathcal{F} \to \mathcal{G}|_U$.
L'adjonction entre $j_!$ et la restriction résulte
de ceci et des propriétés du faisceau associé.
L'identification des fibres est évidente d'après la
définition du prolongement par l'ensemble vide
et la définition d'une fibre.
Les assertions (4) et (5) résultent du calcul de la
valeur du faisceau sur tout ouvert de $U$.
\end{proof}

\noindent
Remarquons que, si $\mathcal{F}$ est un faisceau
de groupes abéliens sur $U$, alors $j_!\mathcal{F}$, tel qu'il est
défini ci-dessus, n'est en général pas un faisceau de groupes abéliens, par exemple
parce que certaines de ses fibres sont vides (et ne sont donc certainement pas
des groupes abéliens). Nous devons donc adapter la définition de
$j_!$ au type de faisceaux considéré.
La raison du choix de l'ensemble vide dans la définition du
prolongement par l'ensemble vide est qu'il s'agit de l'objet initial
de la catégorie des ensembles. Dans le cas des groupes abéliens,
nous utilisons donc $0$ (et, plus généralement, pour les faisceaux à valeurs dans
une catégorie abélienne quelconque).

\begin{definition}
\label{sheaves-definition-j-shriek-structures}
Soit $X$ un espace topologique.
Soit $j : U \to X$ l'inclusion d'un sous-ensemble ouvert.
\begin{enumerate}
\item Soit $\mathcal{F}$ un préfaisceau abélien sur $U$.
Nous définissons le {\it prolongement $j_{p!}\mathcal{F}$ de $\mathcal{F}$ par $0$}
comme le préfaisceau abélien sur $X$ défini par la règle
$$
j_{p!}\mathcal{F}(V) =
\left\{
\begin{matrix}
0 & \text{si} & V \not \subset U \\
\mathcal{F}(V) & \text{si} & V \subset U
\end{matrix}
\right.
$$
avec les applications de restriction évidentes.
\item Soit $\mathcal{F}$ un faisceau abélien sur $U$. Nous définissons
le {\it prolongement $j_!\mathcal{F}$ de $\mathcal{F}$ par $0$}
comme le faisceau associé au préfaisceau abélien $j_{p!}\mathcal{F}$.
\item Soit $\mathcal{C}$ une catégorie possédant un objet initial $e$.
Soit $\mathcal{F}$ un préfaisceau sur $U$ à valeurs dans $\mathcal{C}$.
Nous définissons le {\it prolongement $j_{p!}\mathcal{F}$ de $\mathcal{F}$ par $e$}
comme le préfaisceau sur $X$ à valeurs dans $\mathcal{C}$ défini par la
règle
$$
j_{p!}\mathcal{F}(V) =
\left\{
\begin{matrix}
e & \text{si} & V \not \subset U \\
\mathcal{F}(V) & \text{si} & V \subset U
\end{matrix}
\right.
$$
avec les applications de restriction évidentes.
\item Soit $(\mathcal{C}, F)$ un type de structure algébrique
tel que $\mathcal{C}$ possède un objet initial $e$.
Soit $\mathcal{F}$ un faisceau de structures algébriques sur $U$
(du type donné). Nous définissons le
{\it prolongement $j_!\mathcal{F}$ de $\mathcal{F}$ par $e$}
comme le faisceau associé au préfaisceau $j_{p!}\mathcal{F}$
défini ci-dessus.
\item Soit $\mathcal{O}$ un préfaisceau d'anneaux sur $X$.
Soit $\mathcal{F}$ un préfaisceau de $\mathcal{O}|_U$-modules.
Dans ce cas, nous définissons le {\it prolongement par $0$}
comme le préfaisceau de $\mathcal{O}$-modules qui est égal à
$j_{p!}\mathcal{F}$ comme préfaisceau abélien muni de
l'application de multiplication
$\mathcal{O} \times j_{p!}\mathcal{F} \to j_{p!}\mathcal{F}$.
\item Soit $\mathcal{O}$ un faisceau d'anneaux sur $X$.
Soit $\mathcal{F}$ un faisceau de $\mathcal{O}|_U$-modules.
Dans ce cas, nous définissons le {\it prolongement par $0$}
comme le $\mathcal{O}$-module qui est égal à
$j_!\mathcal{F}$ comme faisceau abélien muni de
l'application de multiplication $\mathcal{O} \times j_!\mathcal{F} \to j_!\mathcal{F}$.
\end{enumerate}
\end{definition}

\noindent
Il est vrai que l'on peut définir $j_!$ dans le cadre des faisceaux
de structures algébriques (voir ci-dessous). Toutefois, l'objet
obtenu dépend du type de structures algébriques considéré.
Par exemple, si $\mathcal{O}$ est un faisceau d'anneaux
sur $U$, alors $j_{!, rings}\mathcal{O} \not = j_{!, abelian}\mathcal{O}$,
car l'objet initial de la catégorie des anneaux
est $\mathbf{Z}$, tandis que l'objet initial de la catégorie
des groupes abéliens est $0$. En particulier, le foncteur $j_!$
{\it ne commute pas avec le passage aux faisceaux d'ensembles sous-jacents},
contrairement à ce que nous avons vu jusqu'ici ! Comme d'habitude, nous traitons séparément
les (pré)faisceaux de groupes abéliens, les (pré)faisceaux de structures algébriques
et les (pré)faisceaux de modules.

\begin{lemma}
\label{sheaves-lemma-j-shriek-abelian}
Soit $X$ un espace topologique.
Soit $j : U \to X$ l'inclusion d'un sous-ensemble ouvert.
Considérons les foncteurs de restriction et de prolongement
par $0$ pour les (pré)faisceaux abéliens.
\begin{enumerate}
\item Le foncteur $j_{p!}$ est adjoint à gauche du
foncteur de restriction $j_p$ (voir le lemme \ref{sheaves-lemma-j-pullback}).
\item Le foncteur $j_!$ est adjoint à gauche de la restriction ;
plus précisément,
$$
\Mor_{\textit{Ab}(X)}(j_!\mathcal{F}, \mathcal{G})
=
\Mor_{\textit{Ab}(U)}(\mathcal{F}, j^{-1}\mathcal{G})
=
\Mor_{\textit{Ab}(U)}(\mathcal{F}, \mathcal{G}|_U)
$$
bifonctoriellement en $\mathcal{F}$ et $\mathcal{G}$.
\item Soit $\mathcal{F}$ un faisceau abélien sur $U$.
Les fibres du faisceau $j_!\mathcal{F}$ sont décrites
comme suit :
$$
j_{!}\mathcal{F}_x =
\left\{
\begin{matrix}
0 & \text{si} & x \not \in U \\
\mathcal{F}_x & \text{si} & x \in U
\end{matrix}
\right.
$$
\item Dans la catégorie des préfaisceaux abéliens sur $U$,
nous avons $j_pj_{p!} = \text{id}$.
\item Dans la catégorie des faisceaux abéliens sur $U$,
nous avons $j^{-1}j_! = \text{id}$.
\end{enumerate}
\end{lemma}

\begin{proof}
Démonstration omise.
\end{proof}

\begin{lemma}
\label{sheaves-lemma-j-shriek-structures}
Soit $X$ un espace topologique.
Soit $j : U \to X$ l'inclusion d'un sous-ensemble ouvert.
Soit $(\mathcal{C}, F)$ un type de structure algébrique
tel que $\mathcal{C}$ possède un objet initial $e$.
Considérons les foncteurs de restriction et de prolongement
par $e$ des (pré)faisceaux de structures algébriques définis ci-dessus.
\begin{enumerate}
\item Le foncteur $j_{p!}$ est adjoint à gauche du
foncteur de restriction $j_p$ (voir le lemme \ref{sheaves-lemma-j-pullback}).
\item Le foncteur $j_!$ est adjoint à gauche de la restriction ;
plus précisément,
$$
\Mor_{\Sh(X, \mathcal{C})}(j_!\mathcal{F}, \mathcal{G})
=
\Mor_{\Sh(U, \mathcal{C})}(\mathcal{F}, j^{-1}\mathcal{G})
=
\Mor_{\Sh(U, \mathcal{C})}(\mathcal{F}, \mathcal{G}|_U)
$$
bifonctoriellement en $\mathcal{F}$ et $\mathcal{G}$.
\item Soit $\mathcal{F}$ un faisceau sur $U$.
Les fibres du faisceau $j_!\mathcal{F}$ sont décrites
comme suit :
$$
j_{!}\mathcal{F}_x =
\left\{
\begin{matrix}
e & \text{si} & x \not \in U \\
\mathcal{F}_x & \text{si} & x \in U
\end{matrix}
\right.
$$
\item Dans la catégorie des préfaisceaux de structures algébriques sur $U$,
nous avons $j_pj_{p!} = \text{id}$.
\item Dans la catégorie des faisceaux de structures algébriques sur $U$,
nous avons $j^{-1}j_! = \text{id}$.
\end{enumerate}
\end{lemma}

\begin{proof}
Démonstration omise.
\end{proof}

\begin{lemma}
\label{sheaves-lemma-j-shriek-modules}
Soit $(X, \mathcal{O})$ un espace annelé.
Soit $j : (U, \mathcal{O}|_U) \to (X, \mathcal{O})$
un sous-espace ouvert.
Considérons les foncteurs de restriction et de prolongement
par $0$ des (pré)faisceaux de modules définis ci-dessus.
\begin{enumerate}
\item Le foncteur $j_{p!}$ est adjoint à gauche de la restriction ;
plus précisément,
$$
\Mor_{\textit{PMod}(\mathcal{O})}(j_{p!}\mathcal{F}, \mathcal{G})
=
\Mor_{\textit{PMod}(\mathcal{O}|_U)}(\mathcal{F}, \mathcal{G}|_U)
$$
bifonctoriellement en $\mathcal{F}$ et $\mathcal{G}$.
\item Le foncteur $j_!$ est adjoint à gauche de la restriction ;
plus précisément,
$$
\Mor_{\textit{Mod}(\mathcal{O})}(j_!\mathcal{F}, \mathcal{G})
=
\Mor_{\textit{Mod}(\mathcal{O}|_U)}(\mathcal{F}, \mathcal{G}|_U)
$$
bifonctoriellement en $\mathcal{F}$ et $\mathcal{G}$.
\item Soit $\mathcal{F}$ un faisceau de $\mathcal{O}$-modules sur $U$.
Les fibres du faisceau $j_!\mathcal{F}$ sont décrites
comme suit :
$$
j_{!}\mathcal{F}_x =
\left\{
\begin{matrix}
0 & \text{si} & x \not \in U \\
\mathcal{F}_x & \text{si} & x \in U
\end{matrix}
\right.
$$
\item Dans la catégorie des faisceaux de $\mathcal{O}|_U$-modules sur $U$,
nous avons $j^{-1}j_! = \text{id}$.
\end{enumerate}
\end{lemma}

\begin{proof}
Démonstration omise.
\end{proof}

\noindent
Remarquons que, d'après les lemmes ci-dessus, les deux foncteurs
$j_*$ et $j_!$ sont des plongements pleinement fidèles de
la catégorie des faisceaux sur $U$ dans la catégorie
des faisceaux sur $X$. Ce n'est que pour le foncteur
$j_!$ que l'on peut décrire aisément l'image
essentielle.

\begin{lemma}
\label{sheaves-lemma-equivalence-categories-open}
Soit $X$ un espace topologique.
Soit $j : U \to X$ l'inclusion d'un sous-ensemble ouvert.
Le foncteur
$$
j_! : \Sh(U) \longrightarrow \Sh(X)
$$
est pleinement fidèle. Son image essentielle est formée exactement
des faisceaux $\mathcal{G}$ tels que
$\mathcal{G}_x = \emptyset$ pour tout $x \in X \setminus U$.
\end{lemma}

\begin{proof}
La pleine fidélité résulte formellement de $j^{-1} j_! = \text{id}$.
Nous avons vu que tout faisceau dans l'image du foncteur possède
la propriété sur les fibres mentionnée dans le lemme. Réciproquement, supposons
que $\mathcal{G}$ possède la propriété indiquée.
Il est alors facile de vérifier que
$$
j_! j^{-1} \mathcal{G} \to \mathcal{G}
$$
est un isomorphisme sur toutes les fibres, donc un isomorphisme.
\end{proof}

\begin{lemma}
\label{sheaves-lemma-equivalence-categories-open-abelian}
Soit $X$ un espace topologique.
Soit $j : U \to X$ l'inclusion d'un sous-ensemble ouvert.
Le foncteur
$$
j_! : \textit{Ab}(U) \longrightarrow \textit{Ab}(X)
$$
est pleinement fidèle. Son image essentielle est formée exactement
des faisceaux $\mathcal{G}$ tels que
$\mathcal{G}_x = 0$ pour tout $x \in X \setminus U$.
\end{lemma}

\begin{proof}
Démonstration omise.
\end{proof}

\begin{lemma}
\label{sheaves-lemma-equivalence-categories-open-structures}
Soit $X$ un espace topologique.
Soit $j : U \to X$ l'inclusion d'un sous-ensemble ouvert.
Soit $(\mathcal{C}, F)$ un type de structure algébrique
tel que $\mathcal{C}$ possède un objet initial $e$.
Le foncteur
$$
j_! : \Sh(U, \mathcal{C}) \longrightarrow \Sh(X, \mathcal{C})
$$
est pleinement fidèle. Son image essentielle est formée exactement
des faisceaux $\mathcal{G}$ tels que
$\mathcal{G}_x = e$ pour tout $x \in X \setminus U$.
\end{lemma}

\begin{proof}
Démonstration omise.
\end{proof}


\begin{lemma}
\label{sheaves-lemma-equivalence-categories-open-modules}
Soit $(X, \mathcal{O})$ un espace annelé.
Soit $j : (U, \mathcal{O}|_U) \to (X, \mathcal{O})$
un sous-espace ouvert.
Le foncteur
$$
j_! : \textit{Mod}(\mathcal{O}|_U) \longrightarrow \textit{Mod}(\mathcal{O})
$$
est pleinement fidèle. Son image essentielle est formée exactement
des faisceaux $\mathcal{G}$ tels que
$\mathcal{G}_x = 0$ pour tout $x \in X \setminus U$.
\end{lemma}

\begin{proof}
Démonstration omise.
\end{proof}

\begin{remark}
\label{sheaves-remark-j-shriek-not-exact}
Soit $j : U \to X$ une immersion ouverte d'espaces topologiques comme ci-dessus.
Soit $x \in X$, $x \not \in U$. Soit $\mathcal{F}$ un faisceau d'ensembles
sur $U$. Alors $j_!\mathcal{F}_x = \emptyset$ d'après le lemme \ref{sheaves-lemma-j-shriek}.
Ainsi, $j_!$ ne transforme pas un objet final de $\Sh(U)$
en un objet final de $\Sh(X)$, sauf si $U = X$.
D'après nos conventions de
Catégories, section \ref{categories-section-exact-functor},
cela signifie que le foncteur $j_!$ n'est pas exact à gauche
comme foncteur entre les catégories de faisceaux d'ensembles.
Nous montrerons plus loin que $j_!$ est exact sur les faisceaux abéliens ;
voir Modules, lemme \ref{modules-lemma-j-shriek-exact}.
\end{remark}







\section{Immersions fermées et (pré)faisceaux}
\label{sheaves-section-closed-immersions}

\noindent
Soit $X$ un espace topologique.
Soit $i : Z \to X$ l'inclusion d'un sous-ensemble fermé $Z$ dans $X$.
Dans la section \ref{sheaves-section-presheaves-functorial}, nous avons défini
des foncteurs $i_*$ et $i^{-1}$ tels que $i_*$ soit adjoint à droite de
$i^{-1}$.

\begin{lemma}
\label{sheaves-lemma-stalks-closed-pushforward}
Soit $X$ un espace topologique.
Soit $i : Z \to X$ l'inclusion d'un sous-ensemble fermé $Z$ dans $X$.
Soit $\mathcal{F}$ un faisceau d'ensembles sur $Z$.
Les fibres de $i_*\mathcal{F}$ se décrivent comme suit :
$$
i_*\mathcal{F}_x =
\left\{
\begin{matrix}
\{*\} & \text{si} & x \not \in Z \\
\mathcal{F}_x & \text{si} & x \in Z
\end{matrix}
\right.
$$
où $\{*\}$ désigne un ensemble à un élément. De plus,
$i^{-1}i_* = \text{id}$ sur la catégorie des faisceaux
d'ensembles sur $Z$. Le même énoncé vaut pour les faisceaux
abéliens sur $Z$, respectivement pour les faisceaux de structures
algébriques sur $Z$, en remplaçant $\{*\}$ par $0$, respectivement
par un objet final de la catégorie des structures algébriques.
\end{lemma}

\begin{proof}
Si $x \not \in Z$, il existe des voisinages ouverts arbitrairement petits
$U$ de $x$ qui ne rencontrent pas $Z$.
Comme $\mathcal{F}$ est un faisceau, on a
$\mathcal{F}(i^{-1}(U)) = \{*\}$ pour tout tel $U$,
voir la remarque \ref{sheaves-remark-confusion}. Cela démontre le premier cas.
Le second résulte du fait que, pour $z \in Z$, tout voisinage ouvert
de $z$ est de la forme $Z \cap U$ pour un ouvert $U$ de $X$.
Pour établir l'égalité $i^{-1}i_* = \text{id}$, considérons l'application
canonique $i^{-1}i_*\mathcal{F} \to \mathcal{F}$. C'est un isomorphisme
sur les fibres (d'après ce qui précède), donc un isomorphisme.

\medskip\noindent
On raisonne de même pour les faisceaux de groupes abéliens et les
faisceaux de structures algébriques.
\end{proof}


\begin{lemma}
\label{sheaves-lemma-equivalence-categories-closed}
Soit $X$ un espace topologique.
Soit $i : Z \to X$ l'inclusion d'un sous-ensemble fermé.
Le foncteur
$$
i_* : \Sh(Z) \longrightarrow \Sh(X)
$$
est pleinement fidèle. Son image essentielle est exactement constituée
des faisceaux $\mathcal{G}$ tels que
$\mathcal{G}_x = \{*\}$ pour tout $x \in X \setminus Z$.
\end{lemma}

\begin{proof}
La pleine fidélité résulte formellement de $i^{-1} i_* = \text{id}$.
Nous avons vu que tout faisceau dans l'image du foncteur possède
la propriété énoncée sur les fibres. Réciproquement, supposons
que $\mathcal{G}$ possède cette propriété.
On vérifie alors aisément que
$$
\mathcal{G} \to i_* i^{-1} \mathcal{G}
$$
est un isomorphisme sur toutes les fibres, donc un isomorphisme.
\end{proof}

\begin{lemma}
\label{sheaves-lemma-equivalence-categories-closed-abelian}
Soit $X$ un espace topologique.
Soit $i : Z \to X$ l'inclusion d'un sous-ensemble fermé.
Le foncteur
$$
i_* : \textit{Ab}(Z) \longrightarrow \textit{Ab}(X)
$$
est pleinement fidèle. Son image essentielle est exactement constituée
des faisceaux $\mathcal{G}$ tels que
$\mathcal{G}_x = 0$ pour tout $x \in X \setminus Z$.
\end{lemma}

\begin{proof}
Omis.
\end{proof}

\begin{lemma}
\label{sheaves-lemma-equivalence-categories-closed-structures}
Soit $X$ un espace topologique.
Soit $i : Z \to X$ l'inclusion d'un sous-ensemble fermé.
Soit $(\mathcal{C}, F)$ une espèce de structure algébrique
ayant un objet final $0$. Le foncteur
$$
i_* : \Sh(Z, \mathcal{C}) \longrightarrow \Sh(X, \mathcal{C})
$$
est pleinement fidèle. Son image essentielle est exactement constituée
des faisceaux $\mathcal{G}$ tels que
$\mathcal{G}_x = 0$ pour tout $x \in X \setminus Z$.
\end{lemma}

\begin{proof}
Omis.
\end{proof}

\begin{remark}
\label{sheaves-remark-i-star-not-exact}
Soit $i : Z \to X$ une immersion fermée d'espaces topologiques comme
ci-dessus. Soit $x \in X$, $x \not \in Z$. Soit $\mathcal{F}$ un faisceau
d'ensembles sur $Z$. Alors $(i_*\mathcal{F})_x = \{ * \}$
d'après le lemme \ref{sheaves-lemma-stalks-closed-pushforward}.
Ainsi, si $\mathcal{F} = * \amalg *$, où
$*$ est le faisceau à un élément, alors
$i_*\mathcal{F}_x = \{*\} \not = i_*(*)_x \amalg i_*(*)_x$,
car le membre de droite est un ensemble à deux éléments.
Selon nos conventions dans
Catégories, section \ref{categories-section-exact-functor},
cela signifie que le foncteur $i_*$ n'est pas exact à droite
comme foncteur entre les catégories de faisceaux d'ensembles.
En particulier, il ne peut pas avoir d'adjoint à droite, voir
Catégories, lemme \ref{categories-lemma-exact-adjoint}.

\medskip\noindent
En revanche, nous verrons plus loin (voir
Modules, lemme \ref{modules-lemma-i-star-right-adjoint})
que $i_*$ est exact sur les faisceaux abéliens et possède bien un adjoint
à droite, à savoir le foncteur qui associe à un faisceau abélien sur $X$
le faisceau de ses sections à support dans $Z$.
\end{remark}

\begin{remark}
\label{sheaves-remark-closed-immersion-spaces}
Nous n'avons pas étudié le lien entre les immersions fermées et les
espaces annelés. En effet, il est préférable d'aborder la notion
d'immersion fermée d'espaces annelés dans le cadre des faisceaux
quasi-cohérents ; voir Modules, section
\ref{modules-section-closed-immersion}.
\end{remark}

\section{Recollement de faisceaux}
\label{sheaves-section-glueing-sheaves}

\noindent
Dans cette section, nous recollons des faisceaux définis sur les membres
d'un recouvrement de $X$. Nous commençons par les applications.

\begin{lemma}
\label{sheaves-lemma-glue-maps}
Soit $X$ un espace topologique.
Soit $X = \bigcup U_i$ un recouvrement ouvert.
Soient $\mathcal{F}$ et $\mathcal{G}$ des faisceaux d'ensembles sur $X$.
Étant donnée une famille
$$
\varphi_i :
\mathcal{F}|_{U_i}
\longrightarrow
\mathcal{G}|_{U_i}
$$
d'applications de faisceaux telle que, pour tous $i, j \in I$, les
applications $\varphi_i, \varphi_j$ induisent par restriction la même application
$\mathcal{F}|_{U_i \cap U_j} \to \mathcal{G}|_{U_i \cap U_j}$,
il existe une unique application de faisceaux
$$
\varphi :
\mathcal{F}
\longrightarrow
\mathcal{G}
$$
dont la restriction à chaque $U_i$ coïncide avec $\varphi_i$.
\end{lemma}

\begin{proof}
Pour tout ouvert $U \subset X$, définissons
$$
\varphi_U : \mathcal{F}(U) \to \mathcal{G}(U), \quad
s \mapsto \varphi_U(s)
$$
où $\varphi_U(s)$ est l'unique section vérifiant
$$
(\varphi_U(s))|_{U \cap U_i}
= (\varphi_i)_{U \cap U_i}(s|_{U \cap U_i}).
$$
L'existence et l'unicité d'une telle section résultent des axiomes
des faisceaux, puisque
\begin{align*}
((\varphi_i)_{U \cap U_i}(s|_{U \cap U_i}))|_{U \cap U_i \cap U_j}
&= (\varphi_i)_{U \cap U_i \cap U_j}(s|_{U \cap U_i \cap U_j})\\
&= (\varphi_j)_{U \cap U_i \cap U_j}(s|_{U \cap U_i \cap U_j})\\
&= ((\varphi_j)_{U \cap U_j}(s|_{U \cap U_j}))|_{U \cap U_i \cap U_j}.
\end{align*}
Cette famille d'applications définit bien une application de faisceaux.
En effet, si $V \subset U \subset X$ sont des ouverts, alors
$$
(\varphi_U(s))|_V
= \varphi_V(s|_V)
$$
car, pour tout $i \in I$, on a
\begin{align*}
(\varphi_U(s))|_{V \cap U_i}
&= ((\varphi_U(s))|_{U \cap U_i})|_{V \cap U_i}\\
&= ((\varphi_i)_{U \cap U_i}(s|_{U \cap U_i}))|_{V \cap U_i}\\
&= (\varphi_i)_{V \cap U_i}(s|_{V \cap U_i})\\
&= \varphi_V(s_{V})|_{V \cap U_i}.
\end{align*}
En outre, sa restriction à chaque $U_i$ coïncide avec $\varphi_i$.
En effet, si $U \subset X$ est ouvert et
$s \in \mathcal{F}(U \cap U_i)$, alors
\begin{align*}
\varphi_{U \cap U_i}(s)
&= \varphi_{U \cap U_i}(s)|_{U \cap U_i}\\
&= (\varphi_i)_{U \cap U_i}(s|_{U \cap U_i})\\
&= (\varphi_i)_{U \cap U_i}(s).
\end{align*}
\end{proof}

\noindent
Le lemme précédent implique que, pour deux faisceaux $\mathcal{F}$ et
$\mathcal{G}$ sur l'espace topologique $X$, la règle
$$
U \longmapsto
\Mor_{\Sh(U)}(
\mathcal{F}|_U, \mathcal{G}|_U)
$$
définit un faisceau. Il s'agit d'une sorte de {\it faisceau Hom interne}.
On l'emploie rarement pour les faisceaux d'ensembles, et plus souvent
pour les faisceaux de modules ; voir Modules, section
\ref{modules-section-internal-hom}.

\medskip\noindent
Soit $X$ un espace topologique.
Soit $X = \bigcup_{i\in I} U_i$ un recouvrement ouvert.
Pour tout $i \in I$, soit $\mathcal{F}_i$ un faisceau d'ensembles sur $U_i$.
Pour toute paire $i, j \in I$, soit
$$
\varphi_{ij} :
\mathcal{F}_i|_{U_i \cap U_j}
\longrightarrow
\mathcal{F}_j|_{U_i \cap U_j}
$$
un isomorphisme de faisceaux d'ensembles. Supposons en outre que, pour
tout triplet d'indices $i, j, k \in I$, le diagramme suivant soit
commutatif :
$$
\xymatrix{
\mathcal{F}_i|_{U_i \cap U_j \cap U_k}
\ar[rr]_{\varphi_{ik}}
\ar[rd]_{\varphi_{ij}} & &
\mathcal{F}_k|_{U_i \cap U_j \cap U_k} \\
&
\mathcal{F}_j|_{U_i \cap U_j \cap U_k}
\ar[ru]_{\varphi_{jk}}
}
$$
Nous appellerons une telle famille de données
$(\mathcal{F}_i, \varphi_{ij})$
une {\it donnée de recollement de faisceaux d'ensembles relativement
au recouvrement $X = \bigcup U_i$}.

\begin{lemma}
\label{sheaves-lemma-glue-sheaves}
Soit $X$ un espace topologique.
Soit $X = \bigcup_{i\in I} U_i$ un recouvrement ouvert.
Pour toute donnée de recollement $(\mathcal{F}_i, \varphi_{ij})$
de faisceaux d'ensembles relativement au recouvrement
$X = \bigcup U_i$, il existe un faisceau d'ensembles $\mathcal{F}$ sur
$X$, muni d'isomorphismes
$$
\varphi_i : \mathcal{F}|_{U_i} \to \mathcal{F}_i
$$
tels que les diagrammes
$$
\xymatrix{
\mathcal{F}|_{U_i \cap U_j} \ar[r]_{\varphi_i} \ar[d]_{\text{id}} &
\mathcal{F}_i|_{U_i \cap U_j} \ar[d]^{\varphi_{ij}} \\
\mathcal{F}|_{U_i \cap U_j} \ar[r]^{\varphi_j} &
\mathcal{F}_j|_{U_i \cap U_j}
}
$$
soient commutatifs.
\end{lemma}

\begin{proof}
Première démonstration. Nous donnons ici une formule pour l'ensemble
des sections de $\mathcal{F}$ sur un ouvert $W \subset X$. Posons
$$
\mathcal{F}(W) =
\{
(s_i)_{i \in I} \mid
s_i \in \mathcal{F}_i(W \cap U_i),
\varphi_{ij}(s_i|_{W \cap U_i \cap U_j}) = s_j|_{W \cap U_i \cap U_j}
\}.
$$
Pour $W' \subset W$, les applications de restriction sont définies en
restreignant chacun des $s_i$ à $W' \cap U_i$. La condition de faisceau
pour $\mathcal{F}$ résulte immédiatement de celle de chacun des
$\mathcal{F}_i$.

\medskip\noindent
Il reste à démontrer que $\mathcal{F}|_{U_i}$ s'envoie
isomorphiquement sur $\mathcal{F}_i$. Soit $W \subset U_i$.
Dans ce cas, la condition de la définition de $\mathcal{F}(W)$ entraîne
$s_j = \varphi_{ij}(s_i|_{W \cap U_j})$.
La commutativité des diagrammes intervenant dans la définition d'une
donnée de recollement garantit alors que l'on peut partir d'une section
{\it quelconque} $s \in \mathcal{F}_i(W)$ et obtenir une famille compatible
en posant $s_i = s$ et
$s_j = \varphi_{ij}(s_i|_{W \cap U_j})$.

\medskip\noindent
Deuxième démonstration (esquisse). Soit $\mathcal{B}$ l'ensemble des ouverts
$U \subset X$ tels que $U \subset U_i$ pour un certain $i \in I$.
Alors $\mathcal{B}$ est une base de la topologie de $X$. Pour
$U \in \mathcal{B}$, choisissons $i \in I$ tel que $U \subset U_i$ et
posons $\mathcal{F}(U) = \mathcal{F}_i(U)$.
Les isomorphismes $\varphi_{ij}$ montrent que cette prescription est
«~indépendante du choix de $i$~». À l'aide des applications de restriction
des $\mathcal{F}_i$, on constate que $\mathcal{F}$ est un faisceau sur
$\mathcal{B}$. Enfin, le lemme \ref{sheaves-lemma-extend-off-basis} permet
d'étendre $\mathcal{F}$ en un unique faisceau $\mathcal{F}$ sur $X$.
\end{proof}

\begin{lemma}
\label{sheaves-lemma-glue-sheaves-structures}
Soit $X$ un espace topologique.
Soit $X = \bigcup U_i$ un recouvrement ouvert.
Soit $(\mathcal{F}_i, \varphi_{ij})$ une donnée de recollement
de faisceaux de groupes abéliens, respectivement de faisceaux de
structures algébriques, respectivement de faisceaux de
$\mathcal{O}$-modules pour un faisceau d'anneaux $\mathcal{O}$ sur $X$.
Alors la construction donnée dans la démonstration du
lemme \ref{sheaves-lemma-glue-sheaves} conduit respectivement à un faisceau de
groupes abéliens, à un faisceau de structures algébriques ou à un
faisceau de $\mathcal{O}$-modules.
\end{lemma}

\begin{proof}
Cela résulte du fait que, dans cette construction, l'ensemble des sections
$\mathcal{F}(W)$ sur un ouvert $W$ est donné par l'égalisateur des
applications
$$
\xymatrix{
\prod_{i \in I} \mathcal{F}_i(W \cap U_i)
\ar@<1ex>[r]
\ar@<-1ex>[r]
&
\prod_{i, j\in I} \mathcal{F}_i(W \cap U_i \cap U_j)
}
$$
Dans chacun des cas envisagés, cet égalisateur définit un objet de la
catégorie considérée dont l'ensemble sous-jacent est celui du lemme cité.
\end{proof}

\begin{lemma}
\label{sheaves-lemma-mapping-property-glue}
Soit $X$ un espace topologique.
Soit $X = \bigcup_{i\in I} U_i$ un recouvrement ouvert.
Le foncteur qui associe à un faisceau d'ensembles $\mathcal{F}$
la donnée de recollement suivante
$$
(\mathcal{F}|_{U_i},
(\mathcal{F}|_{U_i})|_{U_i \cap U_j}
\to
(\mathcal{F}|_{U_j})|_{U_i \cap U_j}
)
$$
relativement au recouvrement $X = \bigcup U_i$ définit une équivalence
de catégories entre $\Sh(X)$ et la catégorie des données de recollement.
Un énoncé analogue vaut pour les faisceaux abéliens, respectivement
les faisceaux de structures algébriques, respectivement les faisceaux
de $\mathcal{O}$-modules.
\end{lemma}

\begin{proof}
Le foncteur est pleinement fidèle d'après le
lemme \ref{sheaves-lemma-glue-maps}, et essentiellement surjectif (par un
quasi-inverse donné explicitement) d'après le
lemme \ref{sheaves-lemma-glue-sheaves}.
\end{proof}

\noindent
Ce lemme signifie que, si le faisceau $\mathcal{F}$ a été construit
à partir de la donnée de recollement $(\mathcal{F}_i, \varphi_{ij})$
et si $\mathcal{G}$ est un faisceau sur $X$, alors un morphisme
$f : \mathcal{F} \to \mathcal{G}$ équivaut à la donnée d'une famille
de morphismes de faisceaux
$$
f_i : \mathcal{F}_i \longrightarrow \mathcal{G}|_{U_i}
$$
compatible avec les applications de recollement $\varphi_{ij}$.
De même, se donner un morphisme de faisceaux
$g : \mathcal{G} \to \mathcal{F}$ revient à se donner une famille
de morphismes de faisceaux
$$
g_i : \mathcal{G}|_{U_i} \longrightarrow \mathcal{F}_i
$$
compatible avec les applications de recollement $\varphi_{ij}$.








% OK, start here.
%

\chapter{Sites et faisceaux}



\phantomsection
\label{sites-section-phantom}


\section{Introduction}
\label{sites-section-introduction}

\noindent
La notion de site a été introduite par Grothendieck afin de pouvoir étudier
les faisceaux pour la topologie \'etale des schémas. La référence fondamentale
pour cette notion est peut-être \cite{SGA4}. Notre notion de site diffère de
celle de \cite{SGA4} ; ce que nous appelons un site y est appelé une catégorie munie
d'une prétopologie dans \cite[Expos\'e II, D\'efinition 1.3]{SGA4}.
Nous faisons ce choix parce qu'en géométrie algébrique il est souvent commode de
travailler avec une classe donnée de recouvrements, par exemple lorsqu'on définit
ce que signifie, pour une propriété des schémas, être locale pour une topologie
donnée ; voir Descente, section \ref{descent-section-descending-properties}.
Notre exposé suivra de près \cite{ArtinTopologies}.
Nous n'utiliserons pas d'univers.












\section{Préfaisceaux}
\label{sites-section-presheaves}

\noindent
Soit $\mathcal{C}$ une catégorie.
Un {\it préfaisceau d'ensembles} est un foncteur contravariant $\mathcal{F}$
de $\mathcal{C}$ vers $\textit{Ensembles}$ (voir Catégories, Remarque
\ref{categories-remark-functor-into-sets}).
Ainsi, pour tout objet $U$ de $\mathcal{C}$, nous avons un ensemble
$\mathcal{F}(U)$. Les éléments de cet ensemble sont appelés
les {\it sections} de $\mathcal{F}$ sur $U$. Pour tout morphisme
$f : V \to U$, l'application $\mathcal{F}(f) : \mathcal{F}(U) \to \mathcal{F}(V)$
est appelée l'{\it application de restriction} et est souvent notée
$f^\ast : \mathcal{F}(U) \to \mathcal{F}(V)$. Une autre façon
d'exprimer ceci consiste à dire que $f^*(s)$ est l'{\it image inverse}
de $s$ par $f$. La fonctorialité signifie que $g^* f^* (s) = (f \circ g)^*(s)$.
Nous employons parfois la notation $s|_V := f^\ast(s)$.
Cette notation est compatible avec la notion de restriction
des fonctions en topologie, car si $W \to V \to U$
sont des morphismes dans $\mathcal{C}$ et si $s$ est une section de
$\mathcal{F}$ sur $U$, alors $s|_W = (s|_V)|_W$ en vertu de la
nature fonctorielle de $\mathcal{F}$. Il faut bien entendu être
prudent, puisqu'il peut fort bien arriver
qu'il existe plus d'un morphisme $V \to U$, et il n'est
certainement pas vrai alors que les applications d'image inverse
correspondantes soient égales.

\begin{definition}
\label{sites-definition-presheaves-sets}
Un {\it préfaisceau d'ensembles} sur $\mathcal{C}$ est un foncteur
contravariant de $\mathcal{C}$ vers $\textit{Ensembles}$. Les {\it morphismes
de préfaisceaux} sont les transformations naturelles de foncteurs. La catégorie
des préfaisceaux d'ensembles est notée $\textit{PSh}(\mathcal{C})$.
\end{definition}

\noindent
Remarquons que, pour tout objet $U$ de $\mathcal{C}$, le foncteur des
points $h_U$, voir Catégories, Exemple \ref{categories-example-hom-functor},
est un préfaisceau. Ceux-ci sont appelés les {\it préfaisceaux représentables}.
Ces préfaisceaux possèdent la propriété agréable suivante : pour tout
préfaisceau $\mathcal{F}$, nous avons
\begin{equation}
\label{sites-equation-map-representable-into-presheaf}
\Mor_{\textit{PSh}(\mathcal{C})}(h_U, \mathcal{F})
=
\mathcal{F}(U).
\end{equation}
C'est le lemme de Yoneda (Catégories, Lemme \ref{categories-lemma-yoneda}).

\medskip\noindent
De même, nous pouvons définir la notion de préfaisceau de groupes abéliens,
d'anneaux, etc. Plus généralement, nous pouvons définir un préfaisceau à valeurs
dans une catégorie.

\begin{definition}
\label{sites-definition-presheaf}
Soient $\mathcal{C}$, $\mathcal{A}$ des catégories.
Un {\it préfaisceau} $\mathcal{F}$ sur $\mathcal{C}$
à valeurs dans $\mathcal{A}$ est un foncteur
contravariant de $\mathcal{C}$ vers $\mathcal{A}$,
c'est-à-dire $\mathcal{F} : \mathcal{C}^{opp} \to \mathcal{A}$.
Un {\it morphisme} de préfaisceaux $\mathcal{F} \to \mathcal{G}$
sur $\mathcal{C}$ à valeurs dans $\mathcal{A}$ est une transformation
naturelle de foncteurs de $\mathcal{F}$ vers $\mathcal{G}$.
\end{definition}

\noindent
Ce sont les objets et les morphismes de la catégorie des préfaisceaux
sur $\mathcal{C}$ à valeurs dans $\mathcal{A}$.

\begin{remark}
\label{sites-remark-big-presheaves}
Comme nous l'avons déjà signalé, nous pouvons considérer la catégorie des
préfaisceaux à valeurs dans l'une quelconque des « grandes » catégories
énumérées dans Catégories, Remarque \ref{categories-remark-big-categories}.
Ce seront elles aussi de « grandes » catégories, et elles seront
ajoutées à la remarque précitée au fur et à mesure.
\end{remark}
















\section{Morphismes injectifs et surjectifs de préfaisceaux}
\label{sites-section-injective-surjective}

\begin{definition}
\label{sites-definition-presheaves-injective-surjective}
Soit $\mathcal{C}$ une catégorie, et soit $\varphi : \mathcal{F}
\to \mathcal{G}$ un morphisme de préfaisceaux d'ensembles.
\begin{enumerate}
\item Nous disons que $\varphi$ est {\it injectif} si, pour tout objet
$U$ de $\mathcal{C}$, l'application $\varphi_U : \mathcal{F}(U)
\to \mathcal{G}(U)$ est injective.
\item Nous disons que $\varphi$ est {\it surjectif} si, pour tout objet
$U$ de $\mathcal{C}$, l'application $\varphi_U : \mathcal{F}(U)
\to \mathcal{G}(U)$ est surjective.
\end{enumerate}
\end{definition}

\begin{lemma}
\label{sites-lemma-mono-epi}
Les morphismes injectifs (resp.\ surjectifs) définis ci-dessus
sont exactement les monomorphismes (resp.\ épimorphismes) de
$\textit{PSh}(\mathcal{C})$. Un morphisme est un isomorphisme
si et seulement s'il est à la fois injectif et surjectif.
\end{lemma}

\begin{proof}
Montrons que $\varphi : \mathcal{F} \to
\mathcal{G}$ est injectif si et seulement si c'est un monomorphisme
de $\textit{PSh}(\mathcal{C})$. Le sens direct
est immédiat ; montrons donc le sens réciproque.
Supposons que $\varphi$ soit un monomorphisme. Soit
$U \in \Ob(\mathcal{C})$ ; nous devons montrer que $\varphi_U$ est
injective. Soient donc $a, b \in \mathcal{F}(U)$ tels que
$\varphi_U (a) = \varphi_U (b)$ ; nous devons vérifier que $a = b$.
Sous l'isomorphisme
(\ref{sites-equation-map-representable-into-presheaf}), les éléments
$a$ et $b$ de $\mathcal{F}(U)$ correspondent à deux transformations
naturelles
$a', b' \in \Mor_{\textit{PSh}(\mathcal{C})}(h_U, \mathcal{F})$.
De même, sous l'isomorphisme analogue
$\Mor_{\textit{PSh}(\mathcal{C})}(h_U, \mathcal{G})
= \mathcal{G}(U)$,
les deux éléments égaux $\varphi_U (a)$ et $\varphi_U (b)$ de
$\mathcal{G}(U)$ correspondent aux deux transformations naturelles
$\varphi \circ a', \varphi \circ b'
\in \Mor_{\textit{PSh}(\mathcal{C})}(h_U, \mathcal{G})$,
qui doivent donc elles aussi être égales. Ainsi
$\varphi \circ a' = \varphi \circ b'$, et par conséquent $a' = b'$
(puisque $\varphi$ est un monomorphisme), d'où $a = b$. Ceci achève (1).

\medskip\noindent
Montrons que $\varphi : \mathcal{F} \to
\mathcal{G}$ est surjectif si et seulement si c'est un épimorphisme
de $\textit{PSh}(\mathcal{C})$. Le sens direct
est immédiat ; montrons donc le sens réciproque.
Supposons que $\varphi$ soit un épimorphisme.

\medskip\noindent
Pour deux morphismes quelconques $f : A \to B$ et $g : A \to C$ de la
catégorie $\textit{Ensembles}$, notons $\text{inl}_{f,g}$ et
$\text{inr}_{f,g}$ les deux applications canoniques de
$B$ et $C$ vers $B \coprod_A C$. (Ici, la somme amalgamée est
calculée dans $\textit{Ensembles}$.)

\medskip\noindent
Définissons maintenant un préfaisceau d'ensembles $\mathcal{H}$ sur $\mathcal{C}$
en posant $\mathcal{H}(U)
= \mathcal{G}(U) \coprod_{\mathcal{F}(U)} \mathcal{G}(U)$ (où
la somme amalgamée est calculée dans $\textit{Ensembles}$ et induite par
l'application $\varphi_U : \mathcal{F}(U) \to \mathcal{G}(U)$) pour
tout $U \in \Ob(\mathcal{C})$ ; son action sur les morphismes est
définie de manière évidente (par la fonctorialité de la somme amalgamée).
Il existe alors deux transformations
naturelles $i_1 : \mathcal{G} \to \mathcal{H}$ et
$i_2 : \mathcal{G} \to \mathcal{H}$ dont les composantes en un objet
$U \in \Ob(\mathcal{C})$ sont données respectivement par les applications
$\text{inl}_{\varphi_U, \varphi_U}$ et
$\text{inr}_{\varphi_U, \varphi_U}$. La
définition d'une somme amalgamée montre que $i_1 \circ \varphi
= i_2 \circ \varphi$, d'où $i_1 = i_2$ (puisque $\varphi$ est un
épimorphisme). Ainsi, pour tout $U \in \Ob(\mathcal{C})$, nous avons
$\text{inl}_{\varphi_U, \varphi_U}
= \text{inr}_{\varphi_U, \varphi_U}$. Par conséquent, $\varphi_U$
doit être surjective (car un argument combinatoire élémentaire montre
que si $f : A \to B$ est un morphisme dans $\textit{Ensembles}$, alors
$\text{inl}_{f,f} = \text{inr}_{f,f}$ si et
seulement si $f$ est surjective). Autrement dit, $\varphi$ est
surjectif, ce qui prouve (2).

\medskip\noindent
Montrons que $\varphi : \mathcal{F} \to
\mathcal{G}$ est à la fois injectif et surjectif si et seulement si c'est
un isomorphisme de $\textit{PSh}(\mathcal{C})$. Cette fois,
le sens réciproque est immédiat. Pour prouver le sens direct,
il suffit d'observer que, si $\varphi$ est à la fois
injectif et surjectif, alors $\varphi_U$ est une application inversible
pour tout $U \in \Ob(\mathcal{C})$, et que les inverses de ces
applications, pour tous les $U$, se combinent en une transformation naturelle
$\mathcal{G} \to \mathcal{F}$ qui est inverse de $\varphi$.
\end{proof}

\begin{definition}
\label{sites-definition-sub-presheaf}
Nous disons que $\mathcal{F}$ est un {\it sous-préfaisceau} de $\mathcal{G}$
si, pour tout objet $U \in \Ob(\mathcal{C})$, l'ensemble
$\mathcal{F}(U)$ est un sous-ensemble de $\mathcal{G}(U)$, de façon compatible
avec les applications de restriction.
\end{definition}

\noindent
Autrement dit, les applications
d'inclusion $\mathcal{F}(U) \to \mathcal{G}(U)$
se recollent pour donner un morphisme (injectif) de
préfaisceaux $\mathcal{F} \to \mathcal{G}$.

\begin{lemma}
\label{sites-lemma-image}
Soit $\mathcal{C}$ une catégorie.
Supposons que $\varphi : \mathcal{F} \to \mathcal{G}$ soit un
morphisme de préfaisceaux d'ensembles sur $\mathcal{C}$.
Il existe un unique sous-préfaisceau $\mathcal{G}' \subset \mathcal{G}$
tel que $\varphi$ se factorise en
$\mathcal{F} \to \mathcal{G}' \to \mathcal{G}$
et que le premier morphisme soit surjectif.
\end{lemma}

\begin{proof}
Pour établir l'existence, il suffit de poser
$\mathcal{G}'(U) = \varphi_U \left(\mathcal{F}(U)\right)$
pour tout $U \in \Ob(C)$ (et d'hériter de l'action sur les morphismes
de $\mathcal{G}$), puis de vérifier que cela définit un
sous-préfaisceau de $\mathcal{G}$ et que $\varphi$ se factorise en
$\mathcal{F} \to \mathcal{G}' \to \mathcal{G}$, le
premier morphisme étant surjectif. L'unicité est immédiate.
\end{proof}

\begin{definition}
\label{sites-definition-image}
Avec les notations du Lemme \ref{sites-lemma-image}, nous
disons que $\mathcal{G}'$ est l'{\it image de $\varphi$}.
\end{definition}
















\section{Limites projectives et limites inductives de préfaisceaux}
\label{sites-section-limits-colimits-PSh}

\noindent
Soit $\mathcal{C}$ une catégorie.
Les limites projectives et les limites inductives existent dans la catégorie
$\textit{PSh}(\mathcal{C})$. De plus, pour tout
$U \in \Ob(\mathcal{C})$, le foncteur
$$
\textit{PSh}(\mathcal{C})
\longrightarrow
\textit{Ensembles}, \quad
\mathcal{F}
\longmapsto
\mathcal{F}(U)
$$
commute aux limites projectives et aux limites inductives. La manière la plus simple
de prouver ces assertions est peut-être la suivante. Étant donné un diagramme
$
\mathcal{F} :
\mathcal{I}
\to
\textit{PSh}(\mathcal{C})
$
définissons les préfaisceaux
$$
\mathcal{F}_{\lim} :
U
\longmapsto
\lim_{i \in \mathcal{I}} \mathcal{F}_i(U)
\text{  et  }
\mathcal{F}_{\colim} :
U
\longmapsto
\colim_{i \in \mathcal{I}} \mathcal{F}_i(U)
$$
Il existe manifestement des morphismes de projection $\mathcal{F}_{\lim} \to \mathcal{F}_i$
et des morphismes canoniques $\mathcal{F}_i \to \mathcal{F}_{\colim}$. Ces
morphismes satisfont les conditions imposées aux morphismes d'une limite projective
(resp.\ d'une limite inductive) de Catégories, Définition \ref{categories-definition-limit}
(resp.\ Catégories, Définition \ref{categories-definition-colimit}).
En effet, ils forment manifestement un cône, resp. un cocône, sur $\mathcal{F}$.
En outre, si $(\mathcal{G}, q_i : \mathcal{G} \to \mathcal{F}_i)$
est un autre
système (comme dans la définition d'une limite projective), alors, pour tout
$U$, nous obtenons un système de morphismes $\mathcal{G}(U) \to \mathcal{F}_i(U)$
satisfaisant les conditions de fonctorialité appropriées. Il existe donc un unique
morphisme $\mathcal{G}(U) \to \mathcal{F}_{\lim}(U)$. On vérifie facilement
que ces morphismes sont compatibles lorsque $U$ varie et proviennent du
morphisme recherché $\mathcal{G} \to \mathcal{F}_{\lim}$.
Un argument analogue vaut dans le cas de la limite inductive.





















\section{Fonctorialité des catégories de préfaisceaux}
\label{sites-section-functoriality-PSh}

\noindent
Soit $u : \mathcal{C} \to \mathcal{D}$ un foncteur entre catégories.
Dans ce cas, nous notons
$$
u^p :
\textit{PSh}(\mathcal{D})
\longrightarrow
\textit{PSh}(\mathcal{C})
$$
le foncteur qui associe à $\mathcal{G}$ sur $\mathcal{D}$ le préfaisceau
$u^p\mathcal{G} = \mathcal{G} \circ u$. Remarquons, d'après la section précédente,
que ce foncteur commute à toutes les limites projectives et inductives.

\medskip\noindent
Pour $V \in \Ob(\mathcal{D})$, notons $\mathcal{I}^u_V$
la catégorie définie par
\begin{equation}
\label{sites-equation-colim-category}
\begin{matrix}
\Ob(\mathcal{I}^u_V)
&
=
&
\{
(U, \phi)
\mid
U \in \Ob(\mathcal{C}),
\phi : V \to u(U)
\}
\\
\Mor_{\mathcal{I}^u_V}((U, \phi), (U', \phi'))
&
=
&
\{
f : U \to U' \text{ dans }\mathcal{C}
\mid
u(f) \circ \phi = \phi'
\}
\end{matrix}
\end{equation}
Nous omettons parfois l'exposant ${}^u$ de la notation et écrivons simplement
$\mathcal{I}_V$.
Nous utiliserons ces catégories pour définir un adjoint à gauche du foncteur $u^p$.
Avant cela, démontrons quelques lemmes techniques.

\begin{lemma}
\label{sites-lemma-almost-directed}
Soit $u : \mathcal{C} \to \mathcal{D}$ un foncteur entre catégories.
Supposons que $\mathcal{C}$ admette les produits fibrés et les égalisateurs, et que
$u$ commute à ceux-ci. Alors les catégories $(\mathcal{I}_V)^{opp}$
satisfont les hypothèses du
lemme de Catégories \ref{categories-lemma-split-into-directed}.
\end{lemma}

\begin{proof}
Il y a deux conditions à vérifier.

\medskip\noindent
Premièrement, supposons donnés trois objets
$\phi : V \to u(U)$, $\phi' : V \to u(U')$ et $\phi'' : V \to u(U'')$,
ainsi que des morphismes $a : U' \to U$, $b : U'' \to U$ tels que
$u(a) \circ \phi' = \phi$ et $u(b) \circ \phi'' = \phi$.
Nous devons montrer qu'il existe un autre objet $\phi''' : V \to u(U''')$
et des morphismes $c : U''' \to U'$ et $d : U''' \to U''$ tels que
$u(c) \circ \phi''' = \phi'$, $u(d) \circ \phi''' = \phi''$ et
$a \circ c = b \circ d$. Prenons $U''' = U' \times_U U''$,
$c$ et $d$ étant les morphismes de projection. Cela convient puisque $u$ commute
aux produits fibrés ; nous omettons la vérification.

\medskip\noindent
Deuxièmement, supposons donnés deux objets
$\phi : V \to u(U)$ et $\phi' : V \to u(U')$,
ainsi que des morphismes $a, b : (U, \phi) \to (U', \phi')$.
Nous devons trouver un morphisme $c : (U'', \phi'') \to (U, \phi)$
qui égalise $a$ et $b$. Soit $c : U'' \to U$ l'égalisateur de
$a$ et $b$ dans la catégorie $\mathcal{C}$. Puisque $u$ commute
aux égalisateurs et que $u(a) \circ \phi = u(b) \circ \phi = \phi'$,
nous obtenons un morphisme $\phi'' : V \to u(U'')$.
\end{proof}

\begin{lemma}
\label{sites-lemma-directed}
Soit $u : \mathcal{C} \to \mathcal{D}$ un foncteur entre catégories.
Supposons que
\begin{enumerate}
\item la catégorie $\mathcal{C}$ possède un objet final $X$ et que
$u(X)$ soit un objet final de $\mathcal{D}$, et
\item la catégorie $\mathcal{C}$ admette les produits fibrés et que
$u$ commute à ceux-ci.
\end{enumerate}
Alors les catégories d'indices $(\mathcal{I}^u_V)^{opp}$ sont filtrantes (voir
Catégories, Définition \ref{categories-definition-directed}).
\end{lemma}

\begin{proof}
Les hypothèses impliquent que celles du
Lemme \ref{sites-lemma-almost-directed}
sont satisfaites (voir la discussion de
Catégories, section \ref{categories-section-finite-limits}).
D'après le
lemme de Catégories \ref{categories-lemma-split-into-directed},
$\mathcal{I}_V$ est une réunion disjointe (éventuellement vide) de
catégories filtrantes.
Il suffit donc de montrer que $\mathcal{I}_V$ est connexe.

\medskip\noindent
Montrons d'abord que $\mathcal{I}_V$ est non vide.
Soit en effet $X$ l'objet final de $\mathcal{C}$,
qui existe par hypothèse.
Soit $V \to u(X)$ l'unique morphisme donné par le fait
que $u(X)$ est final dans $\mathcal{D}$ par hypothèse.
Cela fournit un objet de $\mathcal{I}_V$.

\medskip\noindent
Montrons ensuite que $\mathcal{I}_V$ est connexe.
Soient $\phi_1 : V \to u(U_1)$ et $\phi_2 : V \to u(U_2)$
des objets de $\Ob(\mathcal{I}_V)$. Par hypothèse, $U_1\times U_2$
existe et $u(U_1\times U_2) = u(U_1)\times u(U_2)$.
Considérons le morphisme $\phi : V \to u(U_1\times U_2)$
correspondant à $(\phi_1, \phi_2)$ par la propriété universelle
des produits. L'objet $\phi : V \to u(U_1\times U_2)$
admet manifestement des morphismes vers $\phi_1 : V \to u(U_1)$ et $\phi_2 : V \to u(U_2)$.
\end{proof}

\noindent
Étant donné $g : V' \to V$ dans $\mathcal{D}$, nous obtenons un foncteur
$\overline{g} : \mathcal{I}_V \to \mathcal{I}_{V'}$
en posant $\overline{g}(U, \phi) = (U, \phi \circ g)$
sur les objets. Étant donné un préfaisceau $\mathcal{F}$ sur $\mathcal{C}$,
nous obtenons un foncteur
$$
\mathcal{F}_V :
\mathcal{I}_V^{opp}
\longrightarrow
\textit{Ensembles}, \quad
(U, \phi)
\longmapsto
\mathcal{F}(U).
$$
Autrement dit, $\mathcal{F}_V$ est un préfaisceau d'ensembles sur $\mathcal{I}_V$.
Remarquons que $\mathcal{F}_{V'} \circ \overline{g} = \mathcal{F}_V$.
Nous définissons
$$
u_p\mathcal{F}(V) =
\colim_{\mathcal{I}_V^{opp}} \mathcal{F}_V
$$
La limite inductive fournit, pour chaque $(U, \phi) \in \Ob(\mathcal{I}_V)$,
une application canonique $\mathcal{F}(U)\xrightarrow{c(\phi)}u_p\mathcal{F}(V)$.
Pour $g : V' \to V$ comme ci-dessus, il existe une
application canonique de restriction
$g^* : u_p\mathcal{F}(V) \to u_p\mathcal{F}(V')$
compatible avec
$\mathcal{F}_{V'} \circ \overline{g} = \mathcal{F}_V$
d'après Catégories, Lemme \ref{categories-lemma-functorial-colimit}.
C'est l'unique application telle que, pour tout $(U, \phi) \in \Ob(\mathcal{I}_V)$,
le diagramme
$$
\xymatrix{
\mathcal{F}(U) \ar[r]^{c(\phi)} \ar[d]_{\text{id}}
&
u_p\mathcal{F}(V) \ar[d]^{g^*}
\\
\mathcal{F}(U) \ar[r]^{c(\phi \circ g)}
&
u_p\mathcal{F}(V')
}
$$
soit commutatif. L'unicité de ces applications entraîne que nous obtenons un
préfaisceau. Ce préfaisceau sera noté $u_p\mathcal{F}$.

\begin{lemma}
\label{sites-lemma-recover}
Il existe une application canonique
$\mathcal{F}(U) \to u_p\mathcal{F}(u(U))$,
compatible avec les applications de restriction
(de $\mathcal{F}$ et de $u_p\mathcal{F}$).
\end{lemma}

\begin{proof}
C'est simplement l'application $c(\text{id}_{u(U)})$ introduite ci-dessus.
\end{proof}

\noindent
Remarquons que tout morphisme de préfaisceaux $\mathcal{F} \to \mathcal{F}'$
donne naissance à des systèmes compatibles de morphismes entre foncteurs
$\mathcal{F}_V \to \mathcal{F}'_V$, et donc à un morphisme
de préfaisceaux $u_p\mathcal{F} \to u_p\mathcal{F}'$. Autrement
dit, nous avons défini un foncteur
$$
u_p :
\textit{PSh}(\mathcal{C})
\longrightarrow
\textit{PSh}(\mathcal{D})
$$

\begin{lemma}
\label{sites-lemma-adjoints-u}
Le foncteur $u_p$ est adjoint à gauche du foncteur $u^p$.
Autrement dit, la formule
$$
\Mor_{\textit{PSh}(\mathcal{C})}(\mathcal{F}, u^p\mathcal{G})
=
\Mor_{\textit{PSh}(\mathcal{D})}(u_p\mathcal{F}, \mathcal{G})
$$
est bifonctorielle en $\mathcal{F}$ et $\mathcal{G}$.
\end{lemma}

\begin{proof}
Soit $\mathcal{G}$ un préfaisceau sur $\mathcal{D}$ et soit
$\mathcal{F}$ un préfaisceau sur $\mathcal{C}$.
Nous allons démontrer la formule affichée
en construisant des applications dans les deux sens. Nous laissons
au lecteur le soin de vérifier qu'elles sont inverses l'une de l'autre.

\medskip\noindent
Étant donné un morphisme $\alpha : u_p \mathcal{F} \to \mathcal{G}$,
nous obtenons $u^p\alpha : u^p u_p \mathcal{F} \to u^p \mathcal{G}$.
Le Lemme \ref{sites-lemma-recover} affirme qu'il existe un
morphisme $\mathcal{F} \to u^p u_p \mathcal{F}$. La composée
des deux donne le morphisme recherché. (L'intérêt de cette construction
est qu'elle est manifestement fonctorielle en toutes les données en présence.)

\medskip\noindent
Réciproquement, étant donné un morphisme $\beta : \mathcal{F} \to u^p\mathcal{G}$,
nous obtenons un morphisme $u_p\beta : u_p\mathcal{F} \to u_p u^p\mathcal{G}$.
Soit $V \in \Ob(\mathcal{D})$.
Nous affirmons que le foncteur $u^p\mathcal{G}_V$ sur $\mathcal{I}_V$
admet un morphisme canonique vers le foncteur constant de valeur $\mathcal{G}(V)$.
En effet, pour tout objet $(U, \phi)$ de $\mathcal{I}_V$,
la valeur de $u^p\mathcal{G}_V$ en cet objet est $\mathcal{G}(u(U))$,
qui s'envoie dans $\mathcal{G}(V)$ par $\mathcal{G}(\phi) = \phi^*$.
C'est une transformation naturelle de foncteurs parce que $\mathcal{G}$ est lui-même
un foncteur. Cela fournit un morphisme $u_p u^p \mathcal{G}(V) \to \mathcal{G}(V)$.
Une autre vérification immédiate montre que celui-ci est fonctoriel en $V$,
ce qui donne un morphisme de préfaisceaux $u_p u^p \mathcal{G} \to \mathcal{G}$.
La composée $u_p\mathcal{F} \to u_p u^p\mathcal{G} \to
\mathcal{G}$ est le morphisme recherché.
\end{proof}

\begin{remark}
\label{sites-remark-functoriality-presheaves-values}
Supposons que $\mathcal{A}$ soit une catégorie telle que
tout diagramme $\mathcal{I}_Y \to \mathcal{A}$ admette une
limite inductive dans $\mathcal{A}$. Il est alors clair
qu'il existe des foncteurs $u^p$ et $u_p$, définis
exactement comme ci-dessus, entre les catégories
de préfaisceaux à valeurs dans $\mathcal{A}$.
De plus, l'adjonction du couple
$u^p$ et $u_p$ subsiste dans ce cadre.
\end{remark}

\begin{lemma}
\label{sites-lemma-pullback-representable-presheaf}
Soit $u : \mathcal{C} \to \mathcal{D}$ un foncteur entre catégories.
Pour tout objet $U$ de $\mathcal{C}$, nous avons $u_ph_U = h_{u(U)}$.
\end{lemma}

\begin{proof}
Par l'adjonction de $u_p$ et $u^p$, nous avons
$$
\Mor_{\textit{PSh}(\mathcal{D})}(u_ph_U, \mathcal{G})
=
\Mor_{\textit{PSh}(\mathcal{C})}(h_U, u^p\mathcal{G})
=
u^p\mathcal{G}(U) =
\mathcal{G}(u(U))
$$
et le lemme de Yoneda montre donc que $u_ph_U = h_{u(U)}$ en tant que
préfaisceaux.
\end{proof}
   
    
 
 
 
  











\section{Sites}
\label{sites-section-sites-definitions}

\noindent
Notre notion de site fait appel au type de structures suivant.

\begin{definition}
\label{sites-definition-family-morphisms-fixed-target}
Soit $\mathcal{C}$ une catégorie ; voir Conventions,
section \ref{conventions-section-categories}.
Une {\it famille de morphismes à cible fixée} dans $\mathcal{C}$ est donnée
par un objet $U \in \Ob(\mathcal{C})$, un ensemble $I$ et, pour chaque
$i\in I$, un morphisme $U_i \to U$ de $\mathcal{C}$ de cible $U$.
Nous la noterons $\{U_i \to U\}_{i\in I}$.
\end{definition}

\noindent
Il se peut que l'ensemble $I$ soit vide ! Cette notation vise à évoquer
un recouvrement ouvert, comme en topologie.

\begin{definition}
\label{sites-definition-site}
Un {\it site}\footnote{Cette notation diffère de celle de \cite{SGA4},
comme il est expliqué dans l'introduction.} est la donnée d'une catégorie
$\mathcal{C}$ et d'un ensemble $\text{Cov}(\mathcal{C})$ de familles de
morphismes à cible fixée $\{U_i \to U\}_{i \in I}$, appelées
{\it recouvrements de $\mathcal{C}$}, qui satisfont aux axiomes suivants :
\begin{enumerate}
\item Si $V \to U$ est un isomorphisme, alors $\{V \to U\} \in
\text{Cov}(\mathcal{C})$.
\item Si $\{U_i \to U\}_{i\in I} \in \text{Cov}(\mathcal{C})$ et si, pour
chaque $i$, on a $\{V_{ij} \to U_i\}_{j\in J_i} \in
\text{Cov}(\mathcal{C})$, alors
$\{V_{ij} \to U\}_{i \in I, j\in J_i} \in \text{Cov}(\mathcal{C})$.
\item Si $\{U_i \to U\}_{i\in I}\in \text{Cov}(\mathcal{C})$
et si $V \to U$ est un morphisme de $\mathcal{C}$, alors
$U_i \times_U V$ existe pour tout $i$ et
$\{U_i \times_U V \to V \}_{i\in I} \in \text{Cov}(\mathcal{C})$.
\end{enumerate}
\end{definition}

\noindent
Précisions. Dans l'axiome (1), on exige qu'il existe un recouvrement
$\{U_i \to U\}_{i \in I}$ de $\mathcal{C}$ tel que $I = \{i\}$ soit un
singleton et que le morphisme $U_i \to U$ soit égal au morphisme
$V \to U$ donné en (1). Dans la suite, nous désignerons souvent par
$\{V \to U\}$ une famille de morphismes à cible fixée dont l'ensemble
d'indices est un singleton. Dans l'axiome (3), on exige l'existence du
recouvrement pour un certain choix des produits fibrés
$U_i \times_U V$, $i \in I$.

\begin{remark}
\label{sites-remark-no-big-sites}
(À propos des questions ensemblistes ; à omettre en première lecture.)
La raison principale pour laquelle nous introduisons les sites est l'étude
de la catégorie des faisceaux sur un site, car elle généralise la catégorie
des faisceaux sur un espace topologique, qui joue un rôle si important en
géométrie algébrique. Afin de ne pas avoir à considérer des objets tels que
des « classes de classes », et ainsi de suite, nous n'autoriserons pas les
sites à être des catégories « grandes », contrairement à ce que nous faisons
pour les catégories et les $2$-catégories.

\medskip\noindent
Supposons que $\mathcal{C}$ soit une catégorie et que
$\text{Cov}(\mathcal{C})$ soit une classe propre de recouvrements
satisfaisant (1), (2) et (3) ci-dessus. Nous ne l'autoriserons pas non plus
comme site, principalement parce que nous prendrons des limites sur les
recouvrements. Il existe toutefois plusieurs façons naturelles de remplacer
$\text{Cov}(\mathcal{C})$ par un ensemble de recouvrements, ou par une
structure légèrement différente, qui donnent lieu à la même catégorie de
faisceaux. Par exemple :
\begin{enumerate}
\item Dans Ensembles, section \ref{sets-section-coverings-site}, nous montrons
comment choisir un ensemble approprié de recouvrements qui donne la même
catégorie de faisceaux.
\item On peut aussi prendre la topologie associée (voir la
définition \ref{sites-definition-topology-associated-site}). La topologie ainsi
obtenue sur $\mathcal{C}$ possède la même catégorie de faisceaux. Deux
topologies ont les mêmes catégories de faisceaux si et seulement si elles
sont égales ; voir le théorème \ref{sites-theorem-topology-and-topos}. Une topologie
sur une catégorie est donnée par un choix de cribles sur les objets. La
collection de tous les cribles possibles, et même celle de toutes les
topologies possibles sur $\mathcal{C}$, est un ensemble.
\item On pourrait également modifier légèrement la notion de site ; voir la
remarque \ref{sites-remark-shrink-coverings} ci-dessous. On obtiendrait alors un
ensemble canonique de recouvrements.
\end{enumerate}
Chacune de ces solutions présente un inconvénient mineur. Pour la première,
il faut vérifier que les constructions ultérieures ne dépendent pas du choix
de l'ensemble de recouvrements. Pour la deuxième, il faut se familiariser avec
les topologies et reprendre de nombreux arguments concernant les sites. Pour
la troisième, voir la dernière phrase de la
remarque \ref{sites-remark-shrink-coverings}.

\medskip\noindent
Notre approche consistera à travailler avec les sites de la
définition \ref{sites-definition-site} ci-dessus. Étant donnée une catégorie
$\mathcal{C}$ munie d'une classe propre de recouvrements comme ci-dessus,
nous remplacerons celle-ci par un ensemble de recouvrements donnant un site,
au moyen d'Ensembles, lemme \ref{sets-lemma-coverings-site}. Le
lemme \ref{sites-lemma-choice-set-coverings-immaterial} ci-dessous montre que la
catégorie de faisceaux obtenue (le topos) est indépendante de ce choix. Nous
laissons au lecteur le soin d'employer, s'il le souhaite, l'une des deux
autres stratégies pour traiter ces questions.
\end{remark}

\begin{example}
\label{sites-example-site-topological}
Soit $X$ un espace topologique. Soit $X_{Zar}$ la catégorie dont les objets
sont tous les ouverts $U$ de $X$ et dont les morphismes sont simplement les
applications d'inclusion. Autrement dit, il existe au plus un morphisme entre
deux objets quelconques de $X_{Zar}$. Posons maintenant
$\{U_i \to U\}_{i \in I}\in \text{Cov}(X_{Zar})$ si et seulement si
$\bigcup U_i = U$. Les conditions (1) et (2) ci-dessus sont claires, et (3)
l'est également dès que l'on remarque que, dans $X_{Zar}$, on a
$U \times V = U \cap V$. Notons en particulier que l'ensemble vide doit être
un élément de $X_{Zar}$, faute de quoi cela ne fonctionnerait pas en général.
De plus, comme nous le verrons plus tard, il est tout aussi important
d'autoriser {\it le recouvrement vide de l'ensemble vide comme recouvrement} !
Nous faisons de $X_{Zar}$ un site en choisissant un ensemble approprié de
recouvrements $\text{Cov}(X_{Zar})_{\kappa, \alpha}$ comme dans Ensembles,
lemme \ref{sets-lemma-coverings-site}. Les préfaisceaux et les faisceaux
(au sens défini ci-dessous) sur le site $X_{Zar}$ coïncident exactement avec
les notions usuelles de préfaisceau et de faisceau sur un espace topologique,
définies dans Faisceaux, section \ref{sheaves-section-introduction}.
\end{example}

\begin{example}
\label{sites-example-site-on-group}
Soit $G$ un groupe. Considérons la catégorie $G\textit{-Ensembles}$ dont les objets
sont les ensembles $X$ munis d'une action à gauche de $G$ et dont les
morphismes sont les applications $G$-équivariantes. Un exemple important est
${}_GG$, le $G$-ensemble dont l'ensemble sous-jacent est $G$ et dont l'action
est donnée par la multiplication à gauche. Cette catégorie possède des
produits fibrés ; voir Catégories,
section \ref{categories-section-example-fibre-products}.
Nous déclarons que $\{\varphi_i : U_i \to U\}_{i\in I}$ est un recouvrement si
$\bigcup_{i\in I} \varphi_i(U_i) = U$. On obtient ainsi une classe de
recouvrements sur $G\textit{-Ensembles}$ dont on vérifie aisément qu'elle satisfait
les conditions (1), (2) et (3) de la définition \ref{sites-definition-site}. Le
résultat n'est pas un site, car tant la collection des objets de la catégorie
sous-jacente que la collection des recouvrements forment une classe propre.
Nous remplaçons d'abord $G\textit{-Ensembles}$ par une sous-catégorie pleine
$G\textit{-Ensembles}_\alpha$ comme dans Ensembles,
lemme \ref{sets-lemma-sets-with-group-action}. Le site
$(G\textit{-Ensembles}_\alpha,
\text{Cov}_{\kappa, \alpha'}(G\textit{-Ensembles}_\alpha))$
obtenu ensuite en restreignant convenablement la collection des recouvrements,
comme dans Ensembles, lemme \ref{sets-lemma-coverings-site}, sera noté
$\mathcal{T}_G$.

\medskip\noindent
Dans le cas particulier où le groupe $G$ est dénombrable, on peut prendre
pour $\mathcal{T}_G$ la catégorie des $G$-ensembles dénombrables et pour
recouvrements les familles conjointement surjectives de morphismes
$\{\varphi_i : U_i \to U\}_{i \in I}$ telles que $I$ soit dénombrable.
\end{example}

\begin{example}
\label{sites-example-indiscrete}
Soit $\mathcal{C}$ une catégorie. Il existe une manière canonique d'en faire
un site dans lequel les $\{f : V \to U \mid f\text{ est un isomorphisme}\}$
sont les recouvrements de $U$. Les faisceaux sur ce site sont les préfaisceaux
sur $\mathcal{C}$. La topologie correspondante est appelée {\it topologie
chaotique} ou {\it topologie grossière}.
\end{example}













\section{Faisceaux}
\label{sites-section-sheaves}

\noindent
Soit $\mathcal{C}$ un site. Avant d'introduire la notion de faisceau à valeurs
dans une catégorie, expliquons ce que signifie, pour un préfaisceau
d'ensembles, être un faisceau. Soit $\mathcal{F}$ un préfaisceau d'ensembles
sur $\mathcal{C}$ et soit $\{U_i \to U\}_{i\in I}$ un élément de
$\text{Cov}(\mathcal{C})$. Par hypothèse, tous les produits fibrés
$U_i \times_U U_j$ existent dans $\mathcal{C}$. Il existe deux applications
naturelles
$$
\xymatrix{
\prod\nolimits_{i\in I}
\mathcal{F}(U_i)
\ar@<1ex>[r]^-{\text{pr}_0^*} \ar@<-1ex>[r]_-{\text{pr}_1^*}
&
\prod\nolimits_{(i_0, i_1) \in I \times I}
\mathcal{F}(U_{i_0} \times_U U_{i_1})
}
$$
que nous noterons $\text{pr}^*_i$, $i = 0, 1$, comme indiqué dans la formule
ci-dessus. En effet, un élément du membre de gauche correspond à une famille
$(s_i)_{i\in I}$, où chaque $s_i$ est une section de $\mathcal{F}$ sur $U_i$.
Pour chaque couple $(i_0, i_1) \in I \times I$, on dispose des morphismes de
projection
$$
\text{pr}^{(i_0, i_1)}_{i_0} :
U_{i_0} \times_U U_{i_1}
\longrightarrow
U_{i_0}
\text{ et }
\text{pr}^{(i_0, i_1)}_{i_1} :
U_{i_0} \times_U U_{i_1}
\longrightarrow
U_{i_1}.
$$
On peut donc prendre soit l'image inverse de la section $s_{i_0}$ par la
première de ces applications, soit celle de la section $s_{i_1}$ par la
seconde. Explicitement, les applications mentionnées ci-dessus sont
$$
\text{pr}_0^* :
(s_i)_{i\in I}
\longmapsto
\Big(
\text{pr}^{(i_0, i_1), *}_{i_0}(s_{i_0})
\Big)_{(i_0, i_1) \in I \times I}
$$
et
$$
\text{pr}_1^* :
(s_i)_{i\in I}
\longmapsto
\Big(
\text{pr}^{(i_0, i_1), *}_{i_1}(s_{i_1})
\Big)_{(i_0, i_1) \in I \times I}.
$$
Considérons enfin l'application naturelle
$$
\mathcal{F}(U)
\longrightarrow
\prod\nolimits_{i\in I}
\mathcal{F}(U_i), \quad
s
\longmapsto
(s|_{U_i})_{i \in I}
$$
où la notation $s|_{U_i}$ désigne l'image inverse de $s$ par l'application
$U_i \to U$. La fonctorialité de $\mathcal{F}$ et la commutativité des
diagrammes de produits fibrés montrent clairement que
$\text{pr}_0^*( (s|_{U_i})_{i \in I} ) =
\text{pr}_1^*( (s|_{U_i})_{i \in I} )$.

\begin{definition}
\label{sites-definition-sheaf-sets}
Soit $\mathcal{C}$ un site et soit $\mathcal{F}$ un préfaisceau d'ensembles
sur $\mathcal{C}$. On dit que $\mathcal{F}$ est un {\it faisceau} si, pour
tout recouvrement $\{U_i \to U\}_{i \in I} \in
\text{Cov}(\mathcal{C})$, le diagramme
\begin{equation}
\label{sites-equation-sheaf-condition}
\xymatrix{
\mathcal{F}(U) \ar[r]
&
\prod\nolimits_{i\in I}
\mathcal{F}(U_i)
\ar@<1ex>[r]^-{\text{pr}_0^*} \ar@<-1ex>[r]_-{\text{pr}_1^*}
&
\prod\nolimits_{(i_0, i_1) \in I \times I}
\mathcal{F}(U_{i_0} \times_U U_{i_1})
}
\end{equation}
présente la première flèche comme l'égalisateur de $\text{pr}_0^*$ et
$\text{pr}_1^*$.
\end{definition}

\noindent
Autrement dit, si l'on se donne des sections
$s_i \in \mathcal{F}(U_i)$ telles que
$$
s_i|_{U_i \times_U U_j} = s_j|_{U_i \times_U U_j}
$$
dans $\mathcal{F}(U_i \times_U U_j)$ pour tout couple
$(i, j) \in I \times I$, alors il existe un unique
$s \in \mathcal{F}(U)$ tel que $s_i = s|_{U_i}$.

\begin{remark}
\label{sites-remark-sheaf-condition-empty-covering}
Si le recouvrement $\{U_i \to U\}_{i \in I}$ est la famille vide
(c'est-à-dire si $I = \emptyset$), la condition de faisceau signifie que
$\mathcal{F}(U) = \{*\}$ est un singleton. En effet, dans
(\ref{sites-equation-sheaf-condition}), les deuxième et troisième ensembles sont
des produits vides dans la catégorie des ensembles ; ce sont des objets
finaux de cette catégorie, donc des singletons.
\end{remark}

\begin{example}
\label{sites-example-sheaves-topological}
Soit $X$ un espace topologique. Soit $X_{Zar}$ le site construit dans
l'exemple \ref{sites-example-site-topological}. La notion de faisceau sur $X_{Zar}$
coïncide avec celle de faisceau sur $X$ introduite dans Faisceaux,
définition \ref{sheaves-definition-sheaf}.
\end{example}

\begin{example}
\label{sites-example-topological-wrong}
Soit $X$ un espace topologique. Considérons le site $X'_{Zar}$, identique au
site $X_{Zar}$ de l'exemple \ref{sites-example-site-topological}, à ceci près que
nous n'autorisons pas le recouvrement vide de l'ensemble vide. Autrement dit,
nous autorisons le recouvrement $\{\emptyset \to \emptyset\}$, mais non le
recouvrement dont l'ensemble d'indices est vide. On vérifie aisément que cela
définit encore un site. Nous affirmons toutefois que les faisceaux sur
$X'_{Zar}$ diffèrent de ceux sur $X_{Zar}$. Comme cas extrême, considérons par
exemple la situation où $X = \{p\}$ est un singleton. Les objets de
$X'_{Zar}$ sont alors $\emptyset, X$, tout recouvrement de $\emptyset$ admet
$\{\emptyset \to \emptyset\}$ comme raffinement et tout recouvrement de $X$
admet $\{X \to X\}$ comme raffinement. Un faisceau sur ce site est clairement
donné par un choix arbitraire d'un ensemble $\mathcal{F}(\emptyset)$ et d'un
ensemble $\mathcal{F}(X)$, ainsi que d'une application de restriction
$\mathcal{F}(X) \to \mathcal{F}(\emptyset)$. Ainsi, les faisceaux sur
$X'_{Zar}$ sont les faisceaux usuels sur l'espace à deux points
$\{\eta, p\}$ dont les ouverts sont
$\{\emptyset, \{\eta\}, \{p, \eta\}\}$. En général, les faisceaux sur
$X'_{Zar}$ sont les faisceaux sur l'espace $X \amalg \{\eta\}$ dont les
ouverts sont l'ensemble vide et tous les ensembles de la forme
$U \cup \{\eta\}$, où $U \subset X$ est ouvert.
\end{example}


\begin{definition}
\label{sites-definition-category-sheaves-sets}
La catégorie {\it $\Sh(\mathcal{C})$} des faisceaux d'ensembles est la
sous-catégorie pleine de la catégorie $\textit{PSh}(\mathcal{C})$ dont les
objets sont les faisceaux d'ensembles.
\end{definition}

\noindent
Soit $\mathcal{A}$ une catégorie. Si les produits indexés par $I$ et par
$I \times I$ existent dans $\mathcal{A}$ pour tout $I$ qui intervient comme
ensemble d'indices d'une famille couvrante, alors la
définition \ref{sites-definition-sheaf-sets} ci-dessus a un sens et définit une
notion de faisceau sur $\mathcal{C}$ à valeurs dans $\mathcal{A}$. Remarquons
que le diagramme dans $\mathcal{A}$
$$
\xymatrix{
\mathcal{F}(U) \ar[r]
&
\prod\nolimits_{i\in I}
\mathcal{F}(U_i)
\ar@<1ex>[r]^-{\text{pr}_0^*} \ar@<-1ex>[r]_-{\text{pr}_1^*}
&
\prod\nolimits_{(i_0, i_1) \in I \times I}
\mathcal{F}(U_{i_0} \times_U U_{i_1})
}
$$
est un diagramme égalisateur si et seulement si, pour tout objet $X$ de
$\mathcal{A}$, le diagramme d'ensembles
$$
\xymatrix{
\Mor_\mathcal{A}(X, \mathcal{F}(U)) \ar[r]
&
\prod
\Mor_\mathcal{A}(X, \mathcal{F}(U_i))
\ar@<1ex>[r]^-{\text{pr}_0^*} \ar@<-1ex>[r]_-{\text{pr}_1^*}
&
\prod
\Mor_\mathcal{A}(X, \mathcal{F}(U_{i_0} \times_U U_{i_1}))
}
$$
est un diagramme égalisateur.

\medskip\noindent
Supposons $\mathcal{A}$ arbitraire. Soit $\mathcal{F}$ un préfaisceau à
valeurs dans $\mathcal{A}$. Choisissons un objet quelconque
$X\in \Ob(\mathcal{A})$. On obtient alors un préfaisceau d'ensembles
$\mathcal{F}_X$ défini par la règle
$$
\mathcal{F}_X(U) = \Mor_\mathcal{A}(X, \mathcal{F}(U)).
$$
Il résulte de ce qui précède qu'une bonne définition consiste à exiger que
tous les préfaisceaux $\mathcal{F}_X$ soient des faisceaux d'ensembles.

\begin{definition}
\label{sites-definition-sheaf}
Soient $\mathcal{C}$ un site, $\mathcal{A}$ une catégorie et $\mathcal{F}$ un
préfaisceau sur $\mathcal{C}$ à valeurs dans $\mathcal{A}$. On dit que
$\mathcal{F}$ est un {\it faisceau} si, pour tout objet $X$ de $\mathcal{A}$,
le préfaisceau d'ensembles $\mathcal{F}_X$ (défini ci-dessus) est un faisceau.
\end{definition}










\section{Familles de morphismes à cible fixée}
\label{sites-section-refinements}

\noindent
Cette section introduit quelques notions relatives aux familles de morphismes
ayant la même cible.

\begin{definition}
\label{sites-definition-morphism-coverings}
Soit $\mathcal{C}$ une catégorie. Soit
$\mathcal{U} = \{U_i \to U\}_{i\in I}$ une famille de morphismes de
$\mathcal{C}$ à cible fixée. Soit
$\mathcal{V} = \{V_j \to V\}_{j\in J}$ une autre telle famille.
\begin{enumerate}
\item
Un {\it morphisme de familles de morphismes à cible fixée de $\mathcal{C}$ de
$\mathcal{U}$ vers $\mathcal{V}$}, ou simplement un {\it morphisme de
$\mathcal{U}$ vers $\mathcal{V}$}, est donné par un morphisme $U \to V$, une
application d'ensembles $\alpha : I \to J$ et, pour chaque $i\in I$, un
morphisme $U_i \to V_{\alpha(i)}$ tel que le diagramme
$$
\xymatrix{
U_i \ar[r] \ar[d]
&
V_{\alpha(i)} \ar[d]
\\
U \ar[r]
&
V
}
$$
soit commutatif.
\item Dans le cas particulier où $U = V$ et où $U \to V$ est l'identité,
nous disons que $\mathcal{U}$ est un {\it raffinement} de la famille
$\mathcal{V}$.
\end{enumerate}
\end{definition}

\noindent
Une remarque triviale mais importante est que, si
$\mathcal{V} = \{V_j \to V\}_{j \in J}$ est la {\it famille vide de
morphismes}, c'est-à-dire si $J = \emptyset$, aucune famille
$\mathcal{U} = \{U_i \to V\}_{i \in I}$ avec $I \not = \emptyset$ ne peut
raffiner $\mathcal{V}$ !

\begin{definition}
\label{sites-definition-combinatorial-tautological}
Soit $\mathcal{C}$ une catégorie. Soient
$\mathcal{U} = \{\varphi_i : U_i \to U\}_{i\in I}$ et
$\mathcal{V} = \{\psi_j : V_j \to U\}_{j\in J}$ deux familles de morphismes
à cible fixée.
\begin{enumerate}
\item Nous disons que $\mathcal{U}$ et $\mathcal{V}$ sont
{\it combinatoirement équivalentes} s'il existe des applications
$\alpha : I \to J$ et $\beta : J\to I$ telles que
$\varphi_i = \psi_{\alpha(i)}$ et $\psi_j = \varphi_{\beta(j)}$.
\item Nous disons que $\mathcal{U}$ et $\mathcal{V}$ sont
{\it tautologiquement équivalentes} s'il existe des applications
$\alpha : I \to J$ et $\beta : J\to I$ ainsi que, pour tout $i\in I$ et tout
$j \in J$, des diagrammes commutatifs
$$
\xymatrix{
U_i \ar[rd] \ar[rr] & &
V_{\alpha(i)} \ar[ld] & &
V_j \ar[rd] \ar[rr] & &
U_{\beta(j)} \ar[ld] \\
&
U & & & &
U &
}
$$
dont les flèches horizontales sont des isomorphismes.
\end{enumerate}
\end{definition}

\begin{lemma}
\label{sites-lemma-tautological-combinatorial}
Soit $\mathcal{C}$ une catégorie. Soient
$\mathcal{U} = \{\varphi_i : U_i \to U\}_{i\in I}$ et
$\mathcal{V} = \{\psi_j : V_j \to U\}_{j\in J}$ deux familles de morphismes
ayant la même cible fixée.
\begin{enumerate}
\item Si $\mathcal{U}$ et $\mathcal{V}$ sont combinatoirement équivalentes,
alors elles sont tautologiquement équivalentes.
\item Si $\mathcal{U}$ et $\mathcal{V}$ sont tautologiquement équivalentes,
alors $\mathcal{U}$ est un raffinement de $\mathcal{V}$ et $\mathcal{V}$ est
un raffinement de $\mathcal{U}$.
\item La relation « être combinatoirement équivalentes » est une relation
d'équivalence sur l'ensemble des familles de morphismes à cible fixée.
\item La relation « être tautologiquement équivalentes » est une relation
d'équivalence sur l'ensemble des familles de morphismes à cible fixée.
\item La relation « $\mathcal{U}$ raffine $\mathcal{V}$ et $\mathcal{V}$
raffine $\mathcal{U}$ » est une relation d'équivalence sur l'ensemble des
familles de morphismes à cible fixée.
\end{enumerate}
\end{lemma}

\begin{proof}
Omis.
\end{proof}

\noindent
Dans le lemme suivant, étant donnés une catégorie $\mathcal{C}$, un
préfaisceau $\mathcal{F}$ sur $\mathcal{C}$ et une famille
$\mathcal{U} = \{U_i \to U\}_{i\in I}$ telle que tous les produits fibrés
$U_i \times_U U_{i'}$ existent, nous dirons que {\it la condition de faisceau
pour $\mathcal{F}$ relativement à $\mathcal{U}$} est satisfaite si le
diagramme (\ref{sites-equation-sheaf-condition}) est un diagramme égalisateur.

\begin{lemma}
\label{sites-lemma-tautological-same-sheaf}
Soit $\mathcal{C}$ une catégorie. Soient
$\mathcal{U} = \{\varphi_i : U_i \to U\}_{i\in I}$ et
$\mathcal{V} = \{\psi_j : V_j \to U\}_{j\in J}$ deux familles de morphismes
ayant la même cible fixée. Supposons que les produits fibrés
$U_i \times_U U_{i'}$ et $V_j \times_U V_{j'}$ existent. Si $\mathcal{U}$ et
$\mathcal{V}$ sont tautologiquement équivalentes, alors, pour tout
préfaisceau $\mathcal{F}$ sur $\mathcal{C}$, la condition de faisceau pour
$\mathcal{F}$ relativement à $\mathcal{U}$ équivaut à la condition de
faisceau pour $\mathcal{F}$ relativement à $\mathcal{V}$.
\end{lemma}

\begin{proof}
Remarquons d'abord que, si $\varphi : A \to B$ est un isomorphisme dans la
catégorie $\mathcal{C}$, alors
$\varphi^* : \mathcal{F}(B) \to \mathcal{F}(A)$ est un isomorphisme. Soit
$\beta : J \to I$ une application et soient
$\chi_j : V_j \to U_{\beta(j)}$ des isomorphismes au-dessus de $U$, dont
l'existence résulte de l'hypothèse. Montrons que la condition de faisceau pour
$\mathcal{V}$ implique celle pour $\mathcal{U}$. Supposons données des
sections $s_i \in \mathcal{F}(U_i)$ telles que
$$
s_i|_{U_i \times_U U_{i'}} = s_{i'}|_{U_i \times_U U_{i'}}
$$
dans $\mathcal{F}(U_i \times_U U_{i'})$ pour tout couple
$(i, i') \in I \times I$. On peut alors poser
$s_j = \chi_j^*s_{\beta(j)}$. Pour tout couple
$(j, j') \in J \times J$, le morphisme
$\chi_j \times_{\text{id}_U} \chi_{j'} : V_j \times_U V_{j'} \to
U_{\beta(j)} \times_U U_{\beta(j')}$ est lui aussi un isomorphisme. Par
transport de structure, on voit donc que
$$
s_j|_{V_j \times_U V_{j'}} = s_{j'}|_{V_j \times_U V_{j'}}
$$
également. La condition de faisceau relativement à $\mathcal{V}$ implique
qu'il existe un unique $s$ tel que $s|_{V_j} = s_j$ pour tout $j \in J$.
D'après la première observation de la démonstration, cela implique aussi que
$s|_{U_i} = s_i$ pour tout $i \in \Im(\beta)$. Supposons que $i \in I$ et
$i \not \in \Im(\beta)$. Pour un tel $i$, on dispose d'isomorphismes
$U_i \to V_{\alpha(i)} \to U_{\beta(\alpha(i))}$ au-dessus de $U$. Cela donne
un morphisme $U_i \to U_i \times_U U_{\beta(\alpha(i))}$ qui est une section
de la projection. Puisque $s_i$ et $s_{\beta(\alpha(i))}$ ont la même
restriction au produit fibré, on en conclut que
$s_{\beta(\alpha(i))}$ a pour image inverse $s_i$ par le morphisme
$U_i \to U_{\beta(\alpha(i))}$. Ainsi,
on a également $s_i = s|_{U_i}$, comme voulu.
\end{proof}

\begin{lemma}
\label{sites-lemma-compare-separated-presheaf-condition}
Soit $\mathcal{C}$ une catégorie. Soit
$\mathcal{V} = \{V_j \to U\}_{j \in J} \to
\mathcal{U} = \{U_i \to U\}_{i \in I}$ un morphisme de familles de
morphismes à cible fixée de $\mathcal{C}$ donné par
$\text{id} : U \to U$, $\alpha : J \to I$ et
$f_j : V_j \to U_{\alpha(j)}$. Soit $\mathcal{F}$ un préfaisceau sur
$\mathcal{C}$. Si
$\mathcal{F}(U) \to \prod_{j \in J} \mathcal{F}(V_j)$ est injective, alors
$\mathcal{F}(U) \to \prod_{i \in I} \mathcal{F}(U_i)$ est injective.
\end{lemma}

\begin{proof}
Omis.
\end{proof}

\begin{lemma}
\label{sites-lemma-compare-sheaf-condition}
Soit $\mathcal{C}$ une catégorie. Soit
$\mathcal{V} = \{V_j \to U\}_{j \in J} \to
\mathcal{U} = \{U_i \to U\}_{i \in I}$ un morphisme de familles de
morphismes à cible fixée de $\mathcal{C}$ donné par
$\text{id} : U \to U$, $\alpha : J \to I$ et
$f_j : V_j \to U_{\alpha(j)}$. Soit $\mathcal{F}$ un préfaisceau sur
$\mathcal{C}$. Si
\begin{enumerate}
\item les produits fibrés $U_i \times_U U_{i'}$, $U_i \times_U V_j$ et
$V_j \times_U V_{j'}$ existent,
\item $\mathcal{F}$ satisfait à la condition de faisceau relativement à
$\mathcal{V}$, et
\item pour tout $i \in I$, l'application
$\mathcal{F}(U_i) \to \prod_{j \in J} \mathcal{F}(V_j \times_U U_i)$ est
injective,
\end{enumerate}
alors $\mathcal{F}$ satisfait à la condition de faisceau relativement à
$\mathcal{U}$.
\end{lemma}

\begin{proof}
D'après le lemme \ref{sites-lemma-compare-separated-presheaf-condition},
l'application $\mathcal{F}(U) \to \prod \mathcal{F}(U_i)$ est injective.
Supposons donnés des $s_i \in \mathcal{F}(U_i)$ tels que
$s_i|_{U_i \times_U U_{i'}} = s_{i'}|_{U_i \times_U U_{i'}}$ pour tous
$i, i' \in I$. Posons
$s_j = f_j^*(s_{\alpha(j)}) \in \mathcal{F}(V_j)$. Puisque les morphismes
$f_j$ sont des morphismes au-dessus de $U$, on obtient des morphismes induits
$f_{jj'} : V_j \times_U V_{j'} \to
U_{\alpha(i)} \times_U U_{\alpha(i')}$ compatibles avec $f_j, f_{j'}$ par
les morphismes de projection. Il s'ensuit que
$$
s_j|_{V_j \times_U V_{j'}}
= f_{jj'}^*(s_{\alpha(j)}|_{U_{\alpha(j)} \times_U U_{\alpha(j')}})
= f_{jj'}^*(s_{\alpha(j')}|_{U_{\alpha(j)} \times_U U_{\alpha(j')}})
= s_{j'}|_{V_j \times_U V_{j'}}
$$
pour tous $j, j' \in J$. La condition de faisceau pour $\mathcal{F}$
relativement à $\mathcal{V}$ fournit donc une section
$s \in \mathcal{F}(U)$ dont la restriction est $s_j$ sur chaque $V_j$. Il
suffit de montrer que $s$ se restreint en $s_i$ sur $U_i$, quel que soit
$i \in I$. Comme $\mathcal{F}$ vérifie (3), il suffit de montrer que $s$ et
$s_i$ ont la même restriction à $U_i \times_U V_j$ pour tout $j \in J$.
Pour cela, on utilise
$$
s|_{U_i \times_U V_j} = s_j|_{U_i \times_U V_j} =
(\text{id} \times f_j)^*s_{\alpha(j)}|_{U_i \times_U U_{\alpha(j)}} =
(\text{id} \times f_j)^*s_i|_{U_i \times_U U_{\alpha(j)}} =
s_i|_{U_i \times_U V_j}
$$
comme voulu.
\end{proof}

\begin{lemma}
\label{sites-lemma-refine-same-topology}
Soit $\mathcal{C}$ une catégorie. Soient $\text{Cov}_i$, $i = 1, 2$, deux
ensembles de familles de morphismes à cible fixée qui définissent chacun une
structure de site sur $\mathcal{C}$.
\begin{enumerate}
\item Si tout $\mathcal{U} \in \text{Cov}_1$ est tautologiquement équivalent
à un certain $\mathcal{V} \in \text{Cov}_2$, alors
$\Sh(\mathcal{C}, \text{Cov}_2) \subset
\Sh(\mathcal{C}, \text{Cov}_1)$.
Si, de plus, tout $\mathcal{U} \in \text{Cov}_2$ est tautologiquement
équivalent à un certain $\mathcal{V} \in \text{Cov}_1$, alors les catégories
de faisceaux sont égales.
\item Supposons que, pour tout $\mathcal{U} \in \text{Cov}_1$, il existe un
$\mathcal{V} \in \text{Cov}_2$ tel que $\mathcal{V}$ raffine $\mathcal{U}$.
Dans ce cas,
$\Sh(\mathcal{C}, \text{Cov}_2) \subset
\Sh(\mathcal{C}, \text{Cov}_1)$.
Si, de plus, pour tout $\mathcal{U} \in \text{Cov}_2$, il existe un
$\mathcal{V} \in \text{Cov}_1$ tel que $\mathcal{V}$ raffine $\mathcal{U}$,
alors les catégories de faisceaux sont égales.
\end{enumerate}
\end{lemma}

\begin{proof}
La partie (1) résulte directement du
lemme \ref{sites-lemma-tautological-same-sheaf} et des définitions.

\medskip\noindent
Démonstration de (2). Soit $\mathcal{F}$ un faisceau d'ensembles pour le site
$(\mathcal{C}, \text{Cov}_2)$. Soit $\mathcal{U} \in \text{Cov}_1$, disons
$\mathcal{U} = \{U_i \to U\}_{i \in I}$. Par hypothèse, on peut choisir un
raffinement $\mathcal{V} \in \text{Cov}_2$ de $\mathcal{U}$, disons
$\mathcal{V} = \{V_j \to U\}_{j \in J}$, le raffinement étant donné par
$\alpha : J \to I$ et $f_j : V_j \to U_{\alpha(j)}$. Observons que
$\mathcal{F}$ satisfait à la condition de faisceau pour $\mathcal{V}$ et pour
les recouvrements $\{V_j \times_U U_i \to U_i\}_{j \in J}$, puisque ceux-ci
appartiennent à $\text{Cov}_2$. Le lemme
\ref{sites-lemma-compare-sheaf-condition} montre donc que $\mathcal{F}$ satisfait à
la condition de faisceau pour $\mathcal{U}$.
\end{proof}

\begin{lemma}
\label{sites-lemma-choice-set-coverings-immaterial}
Soit $\mathcal{C}$ une catégorie. Soit $\text{Cov}(\mathcal{C})$ une classe
propre de recouvrements satisfaisant les conditions (1), (2) et (3) de la
définition \ref{sites-definition-site}. Soient
$\text{Cov}_1, \text{Cov}_2 \subset \text{Cov}(\mathcal{C})$ deux parties de
$\text{Cov}(\mathcal{C})$ qui munissent $\mathcal{C}$ d'une structure de site.
Si tout recouvrement $\mathcal{U} \in \text{Cov}(\mathcal{C})$ est
combinatoirement équivalent à un recouvrement de $\text{Cov}_1$ et
combinatoirement équivalent à un recouvrement de $\text{Cov}_2$, alors
$\Sh(\mathcal{C}, \text{Cov}_1) =
\Sh(\mathcal{C}, \text{Cov}_2)$.
\end{lemma}

\begin{proof}
C'est clair d'après les lemmes \ref{sites-lemma-refine-same-topology} et
\ref{sites-lemma-tautological-combinatorial} ci-dessus, puisque l'hypothèse implique
que tout recouvrement
$\mathcal{U} \in \text{Cov}_1 \subset \text{Cov}(\mathcal{C})$ est
combinatoirement équivalent à un élément de $\text{Cov}_2$, et de même après
échange des rôles de $\text{Cov}_1$ et de $\text{Cov}_2$.
\end{proof}












\section{L'exemple des G-ensembles}
\label{sites-section-example-sheaf-G-sets}

\noindent
Comme exemple, considérons le site $\mathcal{T}_G$ de
l'Exemple \ref{sites-example-site-on-group}. Nous allons décrire la
catégorie des faisceaux sur $\mathcal{T}_G$. La réponse se révélera
indépendante des choix faits dans la définition de $\mathcal{T}_G$.
En fait, au cours de la démonstration, nous n'aurons besoin que des
propriétés suivantes du site $\mathcal{T}_G$ :
\begin{enumerate}
\item[(a)] $\mathcal{T}_G$ est une sous-catégorie pleine de $G\textit{-Ensembles}$,
\item[(b)] $\mathcal{T}_G$ contient le $G$-ensemble ${}_GG$,
\item[(c)] $\mathcal{T}_G$ admet des produits fibrés et ceux-ci sont les mêmes que
dans $G\textit{-Ensembles}$,
\item[(d)] étant donné $U \in \Ob(\mathcal{T}_G)$ et un sous-ensemble
$G$-invariant $O \subset U$, il existe un objet de $\mathcal{T}_G$ isomorphe
à $O$, et
\item[(e)] toute famille surjective d'applications $\{U_i \to U\}_{i \in I}$, avec
$U, U_i \in \Ob(\mathcal{T}_G)$, est combinatoirement équivalente à un
recouvrement de $\mathcal{T}_G$.
\end{enumerate}
Ces propriétés découlent de Ensembles, Lemmes \ref{sets-lemma-what-is-in-it-G-sets}
et \ref{sets-lemma-coverings-site}.

\medskip\noindent
Remarquons que l'application
$$
\Hom_G({}_GG, {}_GG)
\longrightarrow
G^{opp},
\varphi
\longmapsto
\varphi(1)
$$
est un isomorphisme de groupes. L'application réciproque envoie $g \in G$
sur l'application $R_g : s \mapsto sg$ (c.-à-d.\ la multiplication à droite).
Notons que $R_{g_1g_2} = R_{g_2} \circ R_{g_1}$, de sorte que le passage au groupe
opposé est nécessaire.

\medskip\noindent
Il en résulte que, pour tout préfaisceau $\mathcal{F}$ sur $\mathcal{T}_G$,
l'ensemble $\mathcal{F}({}_GG)$ hérite d'une structure de $G$-ensemble
comme suit : $g \cdot s$, pour $g \in G$ et $s \in \mathcal{F}({}_GG)$,
est défini par $\mathcal{F}(R_g)(s)$. Il s'agit d'une action à gauche
car
$$
(g_1g_2) \cdot s  = \mathcal{F}(R_{g_1g_2})(s) =
\mathcal{F}(R_{g_2}\circ R_{g_1})(s) =
\mathcal{F}(R_{g_1})( \mathcal{F}(R_{g_2})(s)) =
g_1 \cdot (g_2 \cdot s).
$$
Nous avons utilisé ici le fait que $\mathcal{F}$
est contravariant. Notons que, si $\mathcal{F} \to \mathcal{G}$
est un morphisme de préfaisceaux d'ensembles sur $\mathcal{T}_G$,
nous obtenons alors une application $\mathcal{F}({}_GG) \to \mathcal{G}({}_GG)$
compatible avec les actions de $G$ que nous venons de définir.
En définitive, nous avons construit un foncteur
$$
\textit{PSh}(\mathcal{T}_G)
\longrightarrow
G\textit{-Ensembles}, \quad
\mathcal{F}
\longmapsto
\mathcal{F}({}_GG).
$$
Nous laissons au lecteur le soin de vérifier que cette construction
possède l'agréable propriété d'envoyer le préfaisceau représentable
$h_U$ sur un objet canoniquement isomorphe à $U$.
En formule, $h_U({}_GG) = \Hom_G({}_GG, U) \cong U$.

\medskip\noindent
Supposons que $S$ soit un $G$-ensemble. Nous définissons un préfaisceau
$\mathcal{F}_S$ par la formule\footnote{Il peut sembler qu'il s'agisse du
préfaisceau représentable défini par $S$. Ce n'est pas nécessairement le cas,
car $S$ peut ne pas être un objet de $\mathcal{T}_G$, lequel a été choisi
comme un ensemble suffisamment grand de $G$-ensembles.}
$$
\mathcal{F}_S(U)
=
\Mor_{G\textit{-Ensembles}}(U, S).
$$
C'est manifestement un préfaisceau. D'autre part, supposons que
$\{U_i \to U\}_{i\in I}$ soit un recouvrement de $\mathcal{T}_G$.
Il en résulte que $\coprod_i U_i \to U$ est surjective. Il est donc
clair que l'application
$$
\mathcal{F}_S(U)
=
\Mor_{G\textit{-Ensembles}}(U, S)
\longrightarrow
\prod \mathcal{F}_S(U_i)
=
\prod \Mor_{G\textit{-Ensembles}}(U_i, S)
$$
est injective. Et, étant donnée une famille d'applications
$G$-équivariantes $s_i : U_i \to S$, telle que tous les diagrammes
$$
\xymatrix{
U_i \times_U U_j \ar[d] \ar[r]
&
U_j \ar[d]^{s_j}
\\
U_i \ar[r]^{s_i}
&
S
}
$$
commutent, il existe une unique application $G$-équivariante
$s : U \to S$ telle que $s_i$ soit la composée
$U_i \to U \to S$. En effet, il suffit de définir $s(u) = s_i(u_i)$,
où $i\in I$ est un indice quelconque tel qu'il existe un
$u_i \in U_i$ s'envoyant sur $u$ par l'application $U_i \to U$.
La commutativité des diagrammes ci-dessus signifie exactement
que cette construction est bien définie. En définitive, nous avons
construit un foncteur
$$
G\textit{-Ensembles}
\longrightarrow
\Sh(\mathcal{T}_G), \quad
S
\longmapsto
\mathcal{F}_S
.
$$

\medskip\noindent
Nous avons maintenant le diagramme suivant de catégories et de foncteurs
$$
\xymatrix{
\textit{PSh}(\mathcal{T}_G) \ar[rr]^{\mathcal{F} \mapsto \mathcal{F}({}_GG)}
&
&
G\textit{-Ensembles} \ar[ld]_{S \mapsto \mathcal{F}_S}
\\
&
\Sh(\mathcal{T}_G) \ar[lu]
&
}
$$
Il découle immédiatement des définitions que $\mathcal{F}_S({}_GG)
= \Mor_G({}_GG, S) = S$, la dernière égalité résultant de l'évaluation en $1$.
Ceci démontre presque la proposition suivante.

\begin{proposition}
\label{sites-proposition-sheaves-on-group}
Les foncteurs $\mathcal{F} \mapsto \mathcal{F}({}_GG)$
et $S \mapsto \mathcal{F}_S$ définissent des équivalences
quasi-inverses entre $\Sh(\mathcal{T}_G)$
et $G\textit{-Ensembles}$.
\end{proposition}

\begin{proof}
Nous avons déjà vu que la composée des foncteurs dans un sens
est isomorphe au foncteur identité.
Dans l'autre sens, pour tout faisceau $\mathcal{H}$, il existe un
morphisme naturel de faisceaux
$$
can :
\mathcal{H}
\longrightarrow
\mathcal{F}_{\mathcal{H}({}_GG)}.
$$
En effet, pour tout objet $U$ de $\mathcal{T}_G$, nous définissons $can_U$
comme l'application
$$
\begin{matrix}
\mathcal{H}(U)
&
\longrightarrow
&
\mathcal{F}_{\mathcal{H}({}_GG)}(U)
=
\Mor_G(U, \mathcal{H}({}_GG))
\\
s
&
\longmapsto
&
(u \mapsto \alpha_u^*s).
\end{matrix}
$$
Ici, $\alpha_u : {}_GG \to U$ est l'application
$\alpha_u(g) = gu$ et $\alpha_u^* : \mathcal{H}(U)
\to \mathcal{H}({}_GG)$ est l'application de restriction. Une vérification
élémentaire mais déroutante montre qu'il s'agit bien d'un morphisme
de préfaisceaux. Nous devons montrer que $can$ est un isomorphisme.
Pour cela, nous montrons que $can_U$ est un isomorphisme pour tout
$U \in \Ob(\mathcal{T}_G)$. Nous laissons au lecteur le cas (important mais
facile) où $U = {}_GG$.
Un objet général $U$ de $\mathcal{T}_G$ est une réunion disjointe de
$G$-orbites : $U = \coprod_{i\in I} O_i$. La famille d'applications
$\{O_i \to U\}_{i \in I}$ est tautologiquement équivalente
à un recouvrement de $\mathcal{T}_G$ (par les propriétés de $\mathcal{T}_G$
énumérées au début de cette section). Par conséquent, d'après le Lemme
\ref{sites-lemma-tautological-same-sheaf}, le faisceau $\mathcal{H}$
satisfait la condition de faisceau relativement à
$\{O_i \to U\}_{i \in I}$. La condition de faisceau pour ce recouvrement
implique $\mathcal{H}(U) = \prod_i \mathcal{H}(O_i)$.
Il suffit donc de montrer que $can_U$ est un
isomorphisme lorsque $U$ consiste en une seule $G$-orbite. Soit $u \in U$
et soit $H \subset G$ son stabilisateur. Manifestement,
$\Mor_G(U, \mathcal{H}({}_GG)) = \mathcal{H}({}_GG)^H$
est le sous-ensemble des éléments invariants par $H$. Considérons d'autre part
le recouvrement $\{{}_GG \to U\}$ donné par $g \mapsto
gu$ (une fois encore, il est seulement combinatoirement équivalent à un recouvrement
de $\mathcal{T}_G$, ce qui, là encore, n'a aucune importance).
Notons que le produit fibré $({}_GG)\times_U ({}_GG)$
est égal à $\{(g, gh), g\in G, h\in H\} \cong \coprod_{h\in H} {}_GG$.
La condition de faisceau pour ce recouvrement s'écrit donc
$$
\xymatrix{
\mathcal{H}(U) \ar[r]
&
\mathcal{H}({}_GG)
\ar@<1ex>[r]^-{\text{pr}_0^*} \ar@<-1ex>[r]_-{\text{pr}_1^*}
&
\prod_{h \in H}
\mathcal{H}({}_GG).
}
$$
Les deux applications $\text{pr}_i^*$ à valeurs dans le facteur
$\mathcal{H}({}_GG)$ diffèrent par la multiplication par $h$.
Le résultat découle maintenant de ce fait et de ce que $can$
est un isomorphisme pour $U = {}_GG$.
\end{proof}





















\section{Passage au faisceau associé}
\label{sites-section-sheafification}

\noindent
Afin de définir le passage au faisceau associé, nous étudions le groupe de
cohomologie de degré zéro de {\v C}ech d'un recouvrement et ses propriétés
de fonctorialité.

\medskip\noindent
Soit $\mathcal{F}$ un préfaisceau d'ensembles sur $\mathcal{C}$, et soit
$\mathcal{U} = \{U_i \to U\}_{i \in I}$ un recouvrement de $\mathcal{C}$.
Nous employons la notation $\mathcal{F}(\mathcal{U})$ pour désigner l'égalisateur
$$
H^0(\mathcal{U}, \mathcal{F})
=
\{
(s_i)_{i\in I} \in \prod\nolimits_i \mathcal{F}(U_i)
\mid
s_i|_{U_i \times_U U_j} = s_j|_{U_i \times_U U_j}
\ \forall i, j \in I
\}.
$$
Comme nous le verrons plus loin, il s'agit de la cohomologie de degré zéro de
{\v C}ech de $\mathcal{F}$ sur $U$ relativement au recouvrement $\mathcal{U}$.
Une petite remarque est que nous pouvons définir $H^0(\mathcal{U}, \mathcal{F})$
dès que tous les morphismes $U_i \to U$ sont représentables, c.-à-d.\ que
$\mathcal{U}$ n'a pas besoin d'être un recouvrement du site.
Il existe une application canonique $\mathcal{F}(U) \to H^0(\mathcal{U}, \mathcal{F})$.
Il est clair qu'un morphisme de recouvrements $\mathcal{U} \to \mathcal{V}$
induit des diagrammes commutatifs
$$
\xymatrix{
& U_i \ar[rr] & & V_{\alpha(i)} \\
U_i \times_U U_j \ar[rr] \ar[ur] \ar[dr] & &
V_{\alpha(i)} \times_V V_{\alpha(j)} \ar[ur] \ar[dr] & \\
& U_j \ar[rr] & & V_{\alpha(j)}
}.
$$
Ceux-ci produisent à leur tour une application $H^0(\mathcal{V}, \mathcal{F}) \to
H^0(\mathcal{U}, \mathcal{F})$, compatible à l'application $\mathcal{F}(V)
\to \mathcal{F}(U)$.

\medskip\noindent
Par construction, un préfaisceau $\mathcal{F}$ est un faisceau si et seulement si,
pour tout recouvrement $\mathcal{U}$ de $\mathcal{C}$, l'application naturelle
$\mathcal{F}(U) \to H^0(\mathcal{U}, \mathcal{F})$ est bijective.
Nous emploierons cette notion pour démontrer le lemme simple suivant
sur les limites de faisceaux.

\begin{lemma}
\label{sites-lemma-limit-sheaf}
\begin{slogan}
La limite d'un diagramme de faisceaux existe et coïncide avec la limite dans les préfaisceaux
\end{slogan}
Soit $\mathcal{F} : \mathcal{I} \to \Sh(\mathcal{C})$
un diagramme. Alors $\lim_\mathcal{I} \mathcal{F}$ existe
et est égale à la limite dans la catégorie des préfaisceaux.
\end{lemma}

\begin{proof}
Soit $\lim_i \mathcal{F}_i$ la limite calculée comme préfaisceau.
Nous allons montrer qu'il s'agit d'un faisceau ; il en résultera alors trivialement
que c'est une limite dans la catégorie des faisceaux. Pour vérifier la condition
de faisceau, soit $\mathcal{V} = \{V_j \to V\}_{j\in J}$ un recouvrement.
Soit $(s_j)_{j\in J}$ un élément de $H^0(\mathcal{V}, \lim_i \mathcal{F}_i)$.
En utilisant les morphismes de projection, nous obtenons des éléments $(s_{j, i})_{j\in J}$
de $H^0(\mathcal{V}, \mathcal{F}_i)$. Par la condition de faisceau pour
$\mathcal{F}_i$, nous voyons qu'il existe un unique $s_i \in \mathcal{F}_i(V)$
tel que $s_{j, i} = s_i|_{V_j}$. Soit $\phi : i \to i'$ un morphisme
de la catégorie d'indices. Nous voulons montrer que
$\mathcal{F}(\phi) : \mathcal{F}_i \to \mathcal{F}_{i'}$
envoie $s_i$ sur $s_{i'}$. Nous savons que c'est vrai pour les sections
$s_{i, j}$ et $s_{i', j}$ pour tout $j$, et donc la condition de faisceau
pour $\mathcal{F}_{i'}$ montre que c'est vrai. Nous avons alors un
élément $s = (s_i)_{i \in \Ob(\mathcal{I})}$ de
$(\lim_i \mathcal{F}_i)(V)$. Nous laissons au lecteur le soin de vérifier
que cet élément possède la propriété requise $s_j = s|_{V_j}$.
\end{proof}

\begin{example}
\label{sites-example-singleton-sheaf}
Un exemple particulier est la limite du diagramme vide.
Elle donne l'objet final de la catégorie des (pré)faisceaux.
Il s'agit du préfaisceau qui associe
à tout objet $U$ de $\mathcal{C}$ un ensemble réduit à un élément, avec d'uniques
applications de restriction ; de plus, ce préfaisceau est un faisceau.
Nous notons souvent ce faisceau par $*$.
\end{example}

\noindent
Soit $\mathcal{J}_U$ la catégorie de tous les recouvrements de $U$.
Autrement dit, les objets de $\mathcal{J}_U$ sont les recouvrements
de $U$ dans $\mathcal{C}$, et les morphismes sont les raffinements.
Compte tenu de nos conventions sur les sites, il s'agit bien d'une catégorie, c.-à-d.\ que
la collection des objets et des morphismes forme un ensemble.
Notons que $\Ob(\mathcal{J}_U)$ n'est pas vide puisque
$\{\text{id}_U\}$ en est un objet. D'après les remarques
ci-dessus, la construction $\mathcal{U} \mapsto H^0(\mathcal{U}, \mathcal{F})$
est un foncteur contravariant sur $\mathcal{J}_U$.
Nous définissons
$$
\mathcal{F}^{+}(U)
=
\colim_{\mathcal{J}_U^{opp}}
H^0(\mathcal{U}, \mathcal{F})
$$
Voir Catégories, section \ref{categories-section-limits}, pour
une discussion des limites projectives et inductives. Signalons que nous
verrons plus loin que $\mathcal{F}^{+}(U)$ est la cohomologie de degré zéro de
{\v C}ech de $\mathcal{F}$ sur $U$.

\medskip\noindent
Avant d'en dire davantage sur la structure de la limite inductive, faisons
de la collection d'ensembles
$\mathcal{F}^{+}(U)$, $U \in \Ob(\mathcal{C})$,
un préfaisceau. Soit donc $V \to U$ un morphisme de $\mathcal{C}$.
Les axiomes d'un site fournissent un foncteur\footnote{Cette construction
fait en réalité intervenir un choix des produits fibrés $U_i \times_U V$
et donc l'axiome du choix. L'application obtenue ne dépend pas
des choix effectués, voir ci-dessous.}
$$
\mathcal{J}_U
\longrightarrow
\mathcal{J}_V, \quad
\{U_i \to U\}
\longmapsto
\{U_i \times_U V \to V\}.
$$
Notons que les projections fournissent un
morphisme fonctoriel de recouvrements $\{U_i \times_U V \to V\} \to \{U_i \to U\}$
et donc, par la construction ci-dessus, une application fonctorielle d'ensembles
$H^0(\{U_i \to U\}, \mathcal{F}) \to
H^0(\{U_i \times_U V \to V\}, \mathcal{F})$.
Autrement dit, il existe une transformation de foncteurs
de $H^0(-, \mathcal{F}) : \mathcal{J}_U^{opp} \to \textit{Ensembles}$
vers la composée
$\mathcal{J}_U^{opp} \to \mathcal{J}_V^{opp}
\xrightarrow{H^0(-, \mathcal{F})} \textit{Ensembles}$. Par les
généralités sur les limites inductives, nous obtenons donc une application canonique
$\mathcal{F}^+(U) \to \mathcal{F}^+(V)$. En termes de la description
ci-dessus de l'ensemble $\mathcal{F}^+(U)$, elle envoie simplement l'élément
associé à $s = (s_i) \in H^0(\{U_i \to U\}, \mathcal{F})$ sur l'élément
associé à $(s_i|_{V \times_U U_i})
\in H^0(\{U_i \times_U V \to V\}, \mathcal{F})$.

\begin{lemma}
\label{sites-lemma-plus-presheaf}
Les constructions ci-dessus définissent un préfaisceau
$\mathcal{F}^+$ muni d'un morphisme canonique
de préfaisceaux $\mathcal{F} \to \mathcal{F}^+$.
\end{lemma}

\begin{proof}
Il suffit de montrer qu'étant donnés des morphismes
$W \to V \to U$, la composée $\mathcal{F}^+(U)
\to \mathcal{F}^+(V) \to \mathcal{F}^+(W)$
est égale à l'application $\mathcal{F}^+(U) \to \mathcal{F}^+(W)$.
On peut le montrer directement en vérifiant que, étant donnés
un recouvrement $\{U_i \to U\}$ et
$s = (s_i) \in H^0(\{U_i \to U\}, \mathcal{F})$,
nous avons canoniquement
$W \times_U U_i \cong W \times_V (V \times_U U_i)$,
et que
$s_i|_{W \times_U U_i}$
correspond à
$(s_i|_{V \times_U U_i})|_{W \times_V (V \times_U U_i)}$
par cet isomorphisme.
\end{proof}

\noindent
Plus indirectement, le résultat du
Lemme \ref{sites-lemma-independent-refinement} montre que
nous pouvons tirer en arrière un élément $s$ comme ci-dessus par tout morphisme
d'un recouvrement quelconque de $W$ vers $\{U_i \to U\}$,
et que nous aboutirons toujours au même élément de
$\mathcal{F}^+(W)$.

\begin{lemma}
\label{sites-lemma-plus-functorial}
L'association $\mathcal{F} \mapsto
(\mathcal{F} \to \mathcal{F}^+)$
est un foncteur.
\end{lemma}

\begin{proof}
Au lieu de le démontrer, énonçons exactement ce qu'il faut prouver.
Soit $\mathcal{F} \to \mathcal{G}$ un morphisme de préfaisceaux. Il faut prouver
la commutativité de :
$$
\xymatrix{
\mathcal{F} \ar[r] \ar[d]
&
\mathcal{F}^{+} \ar[d]
\\
\mathcal{G} \ar[r]
&
\mathcal{G}^{+}
}
$$
\end{proof}

\noindent
Les deux lemmes suivants impliquent que les limites inductives ci-dessus sont indexées
par un ensemble ordonné filtrant.

\begin{lemma}
\label{sites-lemma-common-refinement}
Étant donnés deux recouvrements $\{U_i \to U\}$
et $\{V_j \to U\}$ d'un objet donné $U$ du site
$\mathcal{C}$, il existe un recouvrement qui en est un
raffinement commun.
\end{lemma}

\begin{proof}
Puisque $\mathcal{C}$ est un site, pour tout $i$, la
famille $\{V_j \times_U U_i \to U_i\}_j$ est un recouvrement.
Un autre axiome implique alors que $\{V_j \times_U U_i \to U\}_{i, j}$
est un recouvrement de $U$. Manifestement, ce recouvrement raffine les deux
recouvrements donnés.
\end{proof}

\begin{lemma}
\label{sites-lemma-independent-refinement}
Deux morphismes quelconques $f, g: \mathcal{U} \to \mathcal{V}$ de recouvrements
induisant le même morphisme $U \to V$ induisent la même
application $H^0(\mathcal{V}, \mathcal{F}) \to  H^0(\mathcal{U}, \mathcal{F})$.
\end{lemma}

\begin{proof}
Soient $\mathcal{U} = \{U_i \to U\}_{i\in I}$ et
$\mathcal{V} = \{V_j \to V\}_{j\in J}$.
Le morphisme $f$ consiste en une application $U\to V$, une application $\alpha : I \to J$ et
des morphismes $f_i : U_i \to V_{\alpha(i)}$.
De même, $g$~détermine une application $\beta : I \to J$ et des morphismes
$g_i : U_i \to V_{\beta(i)}$.
Comme $f$ et $g$ induisent la même application $U\to V$, le diagramme
$$
\xymatrix{
&
V_{\alpha(i)} \ar[dr]
\\
U_i \ar[ur]^{f_i} \ar[dr]_{g_i}
&
&
V
\\
&
V_{\beta(i)} \ar[ur]
}
$$
est commutatif pour tout $i\in I$. Par conséquent, $f$ et $g$ se factorisent par
le produit fibré
$$
\xymatrix{
&
V_{\alpha(i)}
\\
U_i \ar[r]^-\varphi \ar[ur]^{f_i} \ar[dr]_{g_i}
&
V_{\alpha(i)} \times_V V_{\beta(i)} \ar[u]_{\text{pr}_1} \ar[d]^{\text{pr}_2}
\\
&
V_{\beta(i)}.
}
$$
Soit maintenant $s = (s_j)_j \in H^0(\mathcal{V}, \mathcal{F})$.
Alors, pour tout $i\in I$ :
$$
(f^*s)_i
=
f_i^*(s_{\alpha(i)})
=
\varphi^*\text{pr}_1^*(s_{\alpha(i)})
=
\varphi^*\text{pr}_2^*(s_{\beta(i)})
=
g_i^*(s_{\beta(i)})
=
(g^*s)_i,
$$
où l'égalité du milieu résulte de la définition
de $H^0(\mathcal{V}, \mathcal{F})$.
Ceci montre que les applications
$H^0(\mathcal{V}, \mathcal{F}) \to H^0(\mathcal{U}, \mathcal{F})$
induites par $f$ et $g$ sont égales.
\end{proof}

\begin{remark}
\label{sites-remark-both-refine-same-H0}
En particulier, ce lemme montre que, si $\mathcal{U}$ est
un raffinement de $\mathcal{V}$ et si $\mathcal{V}$ est un
raffinement de $\mathcal{U}$, alors il existe une identification
canonique $H^0(\mathcal{U}, \mathcal{F}) =
H^0(\mathcal{V}, \mathcal{F})$.
\end{remark}

\noindent
Ces deux lemmes et le fait que $\mathcal{J}_U$ est non vide
impliquent que le diagramme $H^0(-, \mathcal{F}) : \mathcal{J}_U^{opp}
\to \textit{Ensembles}$ est filtrant, voir
Catégories, Définition \ref{categories-definition-directed}.
Par conséquent, d'après Catégories,
section \ref{categories-section-directed-colimits},
la limite inductive $\mathcal{F}^{+}(U)$ peut être décrite
de la manière élémentaire suivante. Tout élément de l'ensemble
$\mathcal{F}^{+}(U)$ provient d'un élément
$s \in H^0(\mathcal{U}, \mathcal{F})$ pour un certain recouvrement
$\mathcal{U}$ de $U$. Étant donné un second élément $s' \in
H^0(\mathcal{U}', \mathcal{F})$, les éléments $s$ et $s'$ déterminent
le même élément de la limite inductive si et seulement s'il existe un recouvrement
$\mathcal{V}$ de $U$ et des raffinements $f : \mathcal{V} \to \mathcal{U}$ et
$f' : \mathcal{V} \to \mathcal{U}'$ tels que $f^*s = (f')^*s'$
dans $H^0(\mathcal{V}, \mathcal{F})$. Comme le recouvrement trivial
$\{\text{id}_U\}$ est un objet de $\mathcal{J}_U$, nous obtenons
une application canonique $\mathcal{F}(U) \to \mathcal{F}^+(U)$.

\begin{lemma}
\label{sites-lemma-plus-surjective}
Le morphisme $\theta : \mathcal{F} \to \mathcal{F}^+$ possède la propriété
suivante : pour tout objet $U$ de $\mathcal{C}$ et toute section
$s \in \mathcal{F}^+(U)$, il existe un recouvrement $\{U_i \to U\}$
tel que $s|_{U_i}$ appartienne à l'image de $\theta : \mathcal{F}(U_i)
\to \mathcal{F}^{+}(U_i)$.
\end{lemma}

\begin{proof}
Soit en effet $\{U_i \to U\}$ un recouvrement tel que $s$ provienne
de l'élément $(s_i) \in H^0(\{U_i \to U\}, \mathcal{F})$.
D'après le Lemme \ref{sites-lemma-independent-refinement}, nous pouvons
considérer le recouvrement $\{U_i \to U_i\}$ et le morphisme (évident)
de recouvrements $\{U_i \to U_i\} \to \{U_i \to U\}$ pour calculer la
restriction de $s$ dans $\mathcal{F}^+(U_i)$. Et en effet,
ce recouvrement donne exactement $\theta(s_i)$ comme restriction
de $s$ à $U_i$.
\end{proof}

\begin{definition}
\label{sites-definition-separated}
Nous disons qu'un préfaisceau d'ensembles $\mathcal{F}$ sur un site
$\mathcal{C}$ est {\it séparé} si, pour tout recouvrement
$\{U_i \rightarrow U\}$, l'application
$\mathcal{F}(U) \to \prod \mathcal{F}(U_i)$ est injective.
\end{definition}

\begin{theorem}
\label{sites-theorem-plus}
Avec $\mathcal{F}$ comme ci-dessus,
\begin{enumerate}
\item
\label{sites-item-sep}
Le préfaisceau $\mathcal{F}^+$ est séparé.
\item
\label{sites-item-sheaf}
Si $\mathcal{F}$ est séparé, alors $\mathcal{F}^+$ est un faisceau
et le morphisme de préfaisceaux $\mathcal{F} \to \mathcal{F}^+$ est injectif.
\item
\label{sites-item-plus-iso}
Si $\mathcal{F}$ est un faisceau, alors $\mathcal{F} \to \mathcal{F}^+$
est un isomorphisme.
\item
\label{sites-item-plusplus}
Le préfaisceau $\mathcal{F}^{++}$ est toujours un faisceau.
\end{enumerate}
\end{theorem}

\begin{proof}
Démonstration de (\ref{sites-item-sep}).
Supposons que $s, s' \in \mathcal{F}^+(U)$ et qu'il
existe un recouvrement $\{U_i \to U\}$ tel que
$s|_{U_i} = s'|_{U_i}$ pour tout $i$. Nous avons alors trois recouvrements
de $U$ : le recouvrement $\{U_i \to U\}$ ci-dessus, un recouvrement $\mathcal{U}$
pour $s$ comme dans le Lemme \ref{sites-lemma-plus-surjective},
et un recouvrement analogue $\mathcal{U}'$ pour $s'$. D'après le Lemme
\ref{sites-lemma-common-refinement}, nous pouvons trouver un raffinement commun,
disons $\{W_j \to U\}$. Cela signifie qu'il existe $s_j, s'_j \in \mathcal{F}(W_j)$
tels que $s|_{W_j} = \theta(s_j)$, et de même pour $s'|_{W_j}$, et
tels que $\theta(s_j) = \theta(s'_j)$. Cette dernière égalité signifie
qu'il existe un recouvrement $\{W_{jk} \to W_j\}$ tel que
$s_j|_{W_{jk}} = s'_j|_{W_{jk}}$. Puisque $\{W_{jk} \to U\}$
est un recouvrement, nous voyons alors que $s, s'$ ont la même image dans
$H^0(\{W_{jk} \to U\}, \mathcal{F})$, comme voulu.

\medskip\noindent
Démonstration de (\ref{sites-item-sheaf}). Il est clair que $\mathcal{F} \to
\mathcal{F}^+$ est injectif, car toutes les applications
$\mathcal{F}(U) \to H^0(\mathcal{U}, \mathcal{F})$
sont injectives. Il est également clair que, si $\mathcal{U} \to
\mathcal{U}'$ est un raffinement, alors $H^0(\mathcal{U}', \mathcal{F})
\to H^0(\mathcal{U}, \mathcal{F})$ est injective. Maintenant,
supposons que $\{U_i \to U\}$ soit un recouvrement, et soit
$(s_i)$ une famille d'éléments de $\mathcal{F}^+(U_i)$
satisfaisant la condition de faisceau
$s_i|_{U_i \times_U U_{i'}} = s_{i'}|_{U_i \times_U U_{i'}}$
pour tous $i, i' \in I$. Choisissons des recouvrements (comme dans le
Lemme \ref{sites-lemma-plus-surjective}) $\{U_{ij} \to U_i\}$
tels que $s_i|_{U_{ij}}$ soit l'image de l'unique
élément $s_{ij} \in \mathcal{F}(U_{ij})$. La condition de faisceau
implique que $s_{ij}$ et $s_{i'j'}$ coïncident sur
$U_{ij} \times_U U_{i'j'}$, car ce produit s'envoie dans
$U_i \times_U U_{i'}$ et que nous y avons l'égalité.
Ainsi, $(s_{ij}) \in H^0(\{U_{ij} \to U\}, \mathcal{F})$
donne naissance à un élément $s \in \mathcal{F}^+(U)$. Nous laissons
au lecteur le soin de vérifier que $s|_{U_i} = s_i$.

\medskip\noindent
Démonstration de (\ref{sites-item-plus-iso}). Cela résulte immédiatement des définitions,
car la condition de faisceau dit exactement que toute application
$\mathcal{F} \to H^0(\mathcal{U}, \mathcal{F})$ est bijective
(pour tout recouvrement $\mathcal{U}$ de $U$).

\medskip\noindent
L'assertion (\ref{sites-item-plusplus}) est maintenant évidente.
\end{proof}

\begin{definition}
\label{sites-definition-associated-sheaf}
Soit $\mathcal{C}$ un site et soit $\mathcal{F}$ un préfaisceau
d'ensembles sur $\mathcal{C}$. Le faisceau $\mathcal{F}^\# := \mathcal{F}^{++}$,
muni du morphisme canonique $\mathcal{F} \to \mathcal{F}^\#$,
est appelé le {\it faisceau associé à $\mathcal{F}$}.
\end{definition}

\begin{proposition}
\label{sites-proposition-sheafification-adjoint}
Le morphisme canonique $\mathcal{F} \to \mathcal{F}^\#$ possède la
propriété universelle suivante : pour tout morphisme $\mathcal{F} \to \mathcal{G}$,
où $\mathcal{G}$ est un faisceau d'ensembles, il existe un unique morphisme
$\mathcal{F}^\# \to \mathcal{G}$ tel que $\mathcal{F} \to \mathcal{F}^\#
\to \mathcal{G}$ soit égal au morphisme donné.
\end{proposition}

\begin{proof}
Le Lemme \ref{sites-lemma-plus-functorial} donne un diagramme commutatif
$$
\xymatrix{
\mathcal{F} \ar[r] \ar[d]
&
\mathcal{F}^{+} \ar[r] \ar[d]
&
\mathcal{F}^{++} \ar[d]
\\
\mathcal{G} \ar[r]
&
\mathcal{G}^{+} \ar[r]
&
\mathcal{G}^{++}
}
$$
et, d'après le Théorème \ref{sites-theorem-plus}, les morphismes horizontaux inférieurs
sont des isomorphismes. L'unicité découle du Lemme
\ref{sites-lemma-plus-surjective}, qui affirme que toute section de
$\mathcal{F}^\#$ provient localement de sections de $\mathcal{F}$.
\end{proof}

\noindent
Il résulte clairement de cette proposition que le foncteur $\mathcal{F}
\mapsto (\mathcal{F} \to \mathcal{F}^\#)$ est unique
à un unique isomorphisme de foncteurs près. En fait, notons temporairement
$i : \Sh(\mathcal{C}) \to \textit{PSh}(\mathcal{C})$
le foncteur d'inclusion. Le résultat ci-dessus affirme précisément que
$$
\Mor_{\textit{PSh}(\mathcal{C})}(\mathcal{F}, i(\mathcal{G}))
=
\Mor_{\Sh(\mathcal{C})}(\mathcal{F}^\#, \mathcal{G}).
$$
Autrement dit, le foncteur de passage au faisceau associé est l'adjoint à gauche
du foncteur d'inclusion $i$.

\begin{example}
\label{sites-example-constant-sheaf}
Soit $\mathcal{C}$ un site. Soit $E$ un ensemble.
Le {\it faisceau constant $\underline{E}$ de valeur $E$}
est le faisceau associé au préfaisceau donné par le foncteur constant
de valeur $E$. Il est caractérisé par la règle
$\Mor_{\Sh(\mathcal{C})}(\underline{E}, \mathcal{F}) =
\Mor_{\textit{Ensembles}}(E, \Gamma(\mathcal{C}, \mathcal{F}))$
pour tout faisceau $\mathcal{F}$ sur $\mathcal{C}$, où la
notation est celle de la Définition \ref{sites-definition-global-sections}.
\end{example}

\noindent
Nous terminons cette section par quelques lemmes.

\begin{lemma}
\label{sites-lemma-colimit-sheaves}
\begin{slogan}
La limite inductive dans les faisceaux est le faisceau associé à la limite inductive dans les
préfaisceaux.
\end{slogan}
Soit $\mathcal{F} : \mathcal{I} \to \Sh(\mathcal{C})$
un diagramme. Alors $\colim_\mathcal{I} \mathcal{F}$ existe
et est le faisceau associé à la limite inductive dans la catégorie des préfaisceaux.
\end{lemma}

\begin{proof}
Puisque le foncteur de passage au faisceau associé est un adjoint à gauche, il commute
à la formation de toutes les limites inductives, voir Catégories,
Lemme \ref{categories-lemma-adjoint-exact}.
Ainsi, puisque $\textit{PSh}(\mathcal{C})$ admet des limites inductives, nous en déduisons
que $\Sh(\mathcal{C})$ admet des limites inductives (qui sont les
faisceaux associés aux limites inductives de préfaisceaux).
\end{proof}

\begin{lemma}
\label{sites-lemma-sheafification-exact}
Le foncteur $\textit{PSh}(\mathcal{C}) \to \Sh(\mathcal{C})$,
$\mathcal{F} \mapsto \mathcal{F}^\#$, est exact.
\end{lemma}

\begin{proof}
Puisqu'il est adjoint à gauche, il est exact à droite, voir
Catégories, Lemme \ref{categories-lemma-exact-adjoint}.
D'autre part, d'après les Lemmes \ref{sites-lemma-common-refinement}
et \ref{sites-lemma-independent-refinement}, les limites inductives
intervenant dans la construction de $\mathcal{F}^+$ sont en réalité indexées par
l'ensemble ordonné filtrant $\Ob(\mathcal{J}_U)$, où
$\mathcal{U} \geq \mathcal{U}'$ si et seulement si
$\mathcal{U}$ est un raffinement de $\mathcal{U}'$. Il résulte alors de
Catégories, Lemme \ref{categories-lemma-directed-commutes},
que $\mathcal{F} \to \mathcal{F}^+$ commute
aux limites finies (comme foncteur des préfaisceaux vers les
préfaisceaux). Nous concluons alors en utilisant le
Lemme \ref{sites-lemma-limit-sheaf}.
\end{proof}

\begin{lemma}
\label{sites-lemma-sections-sheafification}
Soit $\mathcal{C}$ un site.
Soit $\mathcal{F}$ un préfaisceau d'ensembles sur $\mathcal{C}$.
Notons $\theta^2 : \mathcal{F} \to \mathcal{F}^\#$ le morphisme canonique
de $\mathcal{F}$ vers son faisceau associé.
Soit $U$ un objet de $\mathcal{C}$.
Soit $s \in \mathcal{F}^\#(U)$. Il existe
un recouvrement $\{U_i \to U\}$ et des sections
$s_i \in \mathcal{F}(U_i)$ tels que
\begin{enumerate}
\item $s|_{U_i} = \theta^2(s_i)$, et
\item pour tous $i, j$, il existe un recouvrement
$\{U_{ijk} \to U_i \times_U U_j\}$ de $\mathcal{C}$ tel que
les restrictions de $s_i$ et $s_j$ à chaque $U_{ijk}$ coïncident.
\end{enumerate}
Réciproquement, étant donnés un recouvrement $\{U_i \to U\}$ et des éléments
$s_i \in \mathcal{F}(U_i)$ tels que (2) soit satisfaite, il
existe une unique section $s \in \mathcal{F}^\#(U)$ telle
que (1) soit satisfaite.
\end{lemma}

\begin{proof}
Omis.
\end{proof}

\begin{lemma}
\label{sites-lemma-isomorphism-sheafifications}
Soit $\mathcal{C}$ un site. Soit $\mathcal{F} \to \mathcal{G}$ un morphisme
de préfaisceaux d'ensembles sur $\mathcal{C}$. Notons $\mathcal{B}$ l'ensemble des
$U \in \Ob(\mathcal{C})$ tels que $\mathcal{F}(U) \to \mathcal{G}(U)$
soit bijective. Si tout objet de $\mathcal{C}$ admet un recouvrement par des éléments
de $\mathcal{B}$, alors $\mathcal{F}^\# \to \mathcal{G}^\#$ est un isomorphisme.
\end{lemma}

\begin{proof}
Soit $U \in \Ob(\mathcal{C})$.
Démontrons que $\mathcal{F}^\#(U) \to \mathcal{G}^\#(U)$ est
surjective. Pour ce faire, nous utiliserons le Lemme \ref{sites-lemma-sections-sheafification}
sans autre mention. Pour tout $s \in \mathcal{G}^\#(U)$, il existe
un recouvrement $\{U_i \to U\}$ et des sections $s_i \in \mathcal{G}(U_i)$ tels que
\begin{enumerate}
\item $s|_{U_i}$ soit l'image de $s_i$ par $\mathcal{G} \to \mathcal{G}^\#$,
et
\item pour tous $i, j$, il existe un recouvrement
$\{U_{ijk} \to U_i \times_U U_j\}$ de $\mathcal{D}$ tel que
les restrictions de $s_i$ et $s_j$ à chaque $U_{ijk}$ coïncident.
\end{enumerate}
Par hypothèse, pour tout $i$, nous pouvons choisir un recouvrement $\{U_{i, a} \to U_i\}$
avec $U_{i, a} \in \mathcal{B}$. Alors $\{U_{i, a} \to U\}$ est un recouvrement.
Notant $s_{i, a}$ l'image de $s_i$ dans $\mathcal{G}(U_{i, a})$,
nous voyons que les restrictions de $s_{i, a}$ et $s_{j, b}$ aux membres
du recouvrement
$$
\{(U_{i, a} \times_U U_{j, b}) \times_{U_i \times_U U_j} U_{ijk}
\to U_{i, a} \times_U U_{j, b}\}
$$
coïncident. Nous pouvons donc supposer que $U_i \in \mathcal{B}$. En répétant
l'argument, nous pouvons aussi supposer que $U_{ijk} \in \mathcal{B}$ pour
tous $i, j, k$ (détails omis).
Alors, puisque $\mathcal{F}(U_i) \to \mathcal{G}(U_i)$
est bijective par notre définition de $\mathcal{B}$ dans l'énoncé
du lemme, nous obtenons un unique $t_i \in \mathcal{F}(U_i)$
qui s'envoie sur $s_i$. Les restrictions de $t_i$ et $t_j$
à chaque $U_{ijk}$ coïncident dans $\mathcal{F}(U_{ijk})$, car
$U_{ijk} \in \mathcal{B}$ et parce que nous avons l'égalité
pour $t_i$ et $t_j$.
Le lemme nous dit alors qu'il existe un unique
$t \in \mathcal{F}^\#(U)$ tel que $t|_{U_i}$ soit l'image de $t_i$.
Il en résulte que $t$ s'envoie sur $s$.
Nous omettons la démonstration de l'injectivité de
$\mathcal{F}^\#(U) \to \mathcal{G}^\#(U)$.
\end{proof}








\section{Applications injectives et surjectives de faisceaux}
\label{sites-section-sheaves-injective}

\begin{definition}
\label{sites-definition-sheaves-injective-surjective}
Soit $\mathcal{C}$ un site, et soit $\varphi : \mathcal{F}
\to \mathcal{G}$ un morphisme de faisceaux d'ensembles.
\begin{enumerate}
\item Nous disons que $\varphi$ est {\it injectif} si, pour tout objet
$U$ de $\mathcal{C}$, l'application $\varphi : \mathcal{F}(U)
\to \mathcal{G}(U)$ est injective.
\item Nous disons que $\varphi$ est {\it surjectif} si, pour tout objet
$U$ de $\mathcal{C}$ et toute section $s\in \mathcal{G}(U)$,
il existe un recouvrement $\{U_i \to U\}$ tel que, pour tout
$i$, la restriction $s|_{U_i}$ appartienne à l'image de
$\varphi : \mathcal{F}(U_i) \to \mathcal{G}(U_i)$.
\end{enumerate}
\end{definition}

\begin{lemma}
\label{sites-lemma-mono-epi-sheaves}
Les morphismes injectifs (resp.\ surjectifs) définis ci-dessus
sont exactement les monomorphismes (resp.\ épimorphismes) de
la catégorie $\Sh(\mathcal{C})$. Un morphisme de faisceaux
est un isomorphisme si et seulement s'il est à la fois injectif
et surjectif.
\end{lemma}

\begin{proof}
Omis.
\end{proof}

\begin{lemma}
\label{sites-lemma-coequalizer-surjection}
Soit $\mathcal{C}$ un site. Soit $\mathcal{F} \to \mathcal{G}$
un morphisme surjectif de faisceaux d'ensembles. Alors le diagramme
$$
\xymatrix{
\mathcal{F} \times_\mathcal{G} \mathcal{F}
\ar@<1ex>[r] \ar@<-1ex>[r]
&
\mathcal{F} \ar[r]
&
\mathcal{G}}
$$
présente $\mathcal{G}$ comme un coégalisateur.
\end{lemma}

\begin{proof}
Soit $\mathcal{H}$ un faisceau d'ensembles et soit
$\varphi : \mathcal{F} \to \mathcal{H}$ un morphisme de faisceaux qui égalise
les deux morphismes $\mathcal{F} \times_\mathcal{G} \mathcal{F} \to \mathcal{F}$.
Soit $\mathcal{G}' \subset \mathcal{G}$ l'image préfaisceautique du
morphisme $\mathcal{F} \to \mathcal{G}$. Comme le produit
$\mathcal{F} \times_\mathcal{G} \mathcal{F}$ peut être calculé dans la
catégorie des préfaisceaux, il est égal au produit préfaisceautique
$\mathcal{F} \times_{\mathcal{G}'} \mathcal{F}$. Ainsi, $\varphi$
induit un unique morphisme de préfaisceaux $\psi' : \mathcal{G}' \to \mathcal{H}$.
Puisque $\mathcal{G}$ est le faisceau associé à $\mathcal{G}'$ d'après
le Lemme \ref{sites-lemma-mono-epi-sheaves},
nous en déduisons que $\psi'$ se prolonge de manière unique en un morphisme de faisceaux
$\psi : \mathcal{G} \to \mathcal{H}$. Nous omettons de vérifier que
$\varphi$ est égal à la composée de $\psi$ et du morphisme donné.
\end{proof}














\section{Faisceaux représentables}
\label{sites-section-representable-sheaves}

\noindent
Soit $\mathcal{C}$ une catégorie. La topologie canonique est
la topologie la plus fine pour laquelle tous les préfaisceaux représentables
sont des faisceaux (elle est formellement définie dans la
Définition \ref{sites-definition-canonical-topology}, mais nous n'en aurons pas
besoin).
Cette topologie n'est pas toujours celle qui est associée à la
structure de site sur $\mathcal{C}$.
Nous allons donner une famille de recouvrements qui engendre cette topologie
lorsque $\mathcal{C}$ possède des produits fibrés. Donnons d'abord
la définition générale suivante.

\begin{definition}
\label{sites-definition-universal-effective-epimorphisms}
Soit $\mathcal{C}$ une catégorie. Nous disons qu'une famille $\{U_i \to U\}_{i \in I}$
est une {\it famille épimorphique effective} si tous les morphismes $U_i \to U$ sont
représentables (voir
Catégories, Définition \ref{categories-definition-representable-morphism}),
et si, pour tout $X\in \Ob(\mathcal{C})$, la suite
$$
\xymatrix{
\Mor_\mathcal{C}(U, X) \ar[r]
&
\prod\nolimits_{i \in I} \Mor_\mathcal{C}(U_i, X)
\ar@<1ex>[r] \ar@<-1ex>[r]
&
\prod\nolimits_{(i, j) \in I^2} \Mor_\mathcal{C}(U_i \times_U U_j, X)
}
$$
est un diagramme égalisateur. Nous disons qu'une famille $\{U_i \to U\}$ est une
{\it famille épimorphique effective universelle} si, pour tout morphisme $V \to U$,
le changement de base $\{U_i \times_U V \to V\}$ est une famille épimorphique effective.
\end{definition}

\noindent
La classe des familles épimorphiques effectives universelles
satisfait les axiomes de la Définition \ref{sites-definition-site}.
Si $\mathcal{C}$ possède des produits fibrés, la topologie associée est
la topologie canonique. (Dans ce cas, pour obtenir un site, on raisonne comme dans Ensembles,
Lemme \ref{sets-lemma-coverings-site}.)

\medskip\noindent
Réciproquement, supposons que $\mathcal{C}$ soit un site tel que
tous les préfaisceaux représentables soient des faisceaux. Alors, manifestement, tous
les recouvrements sont des familles épimorphiques effectives universelles.
La définition suivante est donc la définition « correcte » dans le
cadre des sites.

\begin{definition}
\label{sites-definition-weaker-than-canonical}
Nous disons que la topologie d'un site $\mathcal{C}$ est
{\it moins fine que la topologie canonique}, ou qu'elle est
{\it sous-canonique}, si tous les recouvrements
de $\mathcal{C}$ sont des familles épimorphiques effectives universelles.
\end{definition}

\noindent
Un faisceau représentable est un préfaisceau représentable qui est aussi un
faisceau. Comme il vaut peut-être mieux éviter cette terminologie lorsque la
topologie n'est pas sous-canonique, nous ne la définissons formellement que dans ce cas.

\begin{definition}
\label{sites-definition-representable-sheaf}
Soit $\mathcal{C}$ un site dont la topologie est sous-canonique.
Le plongement de Yoneda $h$ (voir
Catégories, section \ref{categories-section-opposite})
présente $\mathcal{C}$ comme une sous-catégorie pleine de la
catégorie des faisceaux sur $\mathcal{C}$. Dans ce cas,
nous appelons les faisceaux de la forme $h_U$, avec $U \in \Ob(\mathcal{C})$,
des {\it faisceaux représentables} sur $\mathcal{C}$.
Notation : le faisceau représentable $h_U$ associé à $U$ est parfois
noté {\it $\underline{U}$}.
\end{definition}

\noindent
Remarquons que, dans la situation de la définition,
$$
\Mor_{\Sh(\mathcal{C})}(h_U, \mathcal{F}) = \mathcal{F}(U)
$$
pour tout faisceau $\mathcal{F}$, puisque cette égalité vaut pour les préfaisceaux ; voir
(\ref{sites-equation-map-representable-into-presheaf}). En général, les
préfaisceaux $h_U$ ne sont pas des faisceaux et il faut leur associer un
faisceau. Dans ce cas, nous avons encore
\begin{equation}
\label{sites-equation-map-representable-into-sheaf}
\Mor_{\Sh(\mathcal{C})}(h_U^\#, \mathcal{F}) =
\Mor_{\textit{PSh}(\mathcal{C})}(h_U, \mathcal{F}) =
\mathcal{F}(U)
\end{equation}
pour tout faisceau $\mathcal{F}$. En effet, la première égalité résulte
de la propriété d'adjonction de $\#$ et la seconde est
(\ref{sites-equation-map-representable-into-presheaf}).

\begin{lemma}
\label{sites-lemma-covering-surjective-after-sheafification}
\begin{slogan}
Les recouvrements deviennent surjectifs après passage au faisceau associé.
\end{slogan}
Soit $\mathcal{C}$ un site. Si
$\{U_i \to U\}_{i \in I}$ est un recouvrement du site
$\mathcal{C}$, alors le morphisme de préfaisceaux d'ensembles
$$
\coprod\nolimits_{i \in I} h_{U_i} \to h_U
$$
devient surjectif après passage aux faisceaux associés.
\end{lemma}

\begin{proof}
D'après le Lemme \ref{sites-lemma-mono-epi-sheaves} ci-dessus, il suffit de montrer que
$\coprod\nolimits_{i \in I} h_{U_i}^\# \to h_U^\#$
est un épimorphisme. Soit $\mathcal{F}$ un faisceau d'ensembles.
Un morphisme $h_U^\# \to \mathcal{F}$
correspond à une section $s \in \mathcal{F}(U)$.
L'injectivité de $\Mor(h_U^\#, \mathcal{F})
\to \prod_i \Mor(h_{U_i}^\#, \mathcal{F})$ découle donc
directement de la propriété de faisceau de $\mathcal{F}$.
\end{proof}

\noindent
Le lemme suivant affirme, lorsque la topologie est moins fine que la
topologie canonique, que tout faisceau est en quelque sorte construit à partir de
faisceaux représentables.

\begin{lemma}
\label{sites-lemma-sheaf-coequalizer-representable}
Soit $\mathcal{C}$ un site. Soit $E \subset \Ob(\mathcal{C})$ un
sous-ensemble tel que tout objet de $\mathcal{C}$ admette un recouvrement par
des éléments de $E$. Soit $\mathcal{F}$ un faisceau d'ensembles. Il existe un
diagramme de faisceaux d'ensembles
$$
\xymatrix{
\mathcal{F}_1 \ar@<1ex>[r] \ar@<-1ex>[r] &
\mathcal{F}_0 \ar[r] &
\mathcal{F}
}
$$
qui présente $\mathcal{F}$ comme un coégalisateur,
tel que les $\mathcal{F}_i$, $i = 0, 1$, soient des sommes disjointes
de faisceaux de la forme $h_U^\#$ avec $U \in E$.
\end{lemma}

\begin{proof}
Montrons d'abord qu'il existe un épimorphisme $\mathcal{F}_0 \to \mathcal{F}$
du type voulu. Il suffit de prendre
$$
\mathcal{F}_0 =
\coprod\nolimits_{U \in E, s \in \mathcal{F}(U)}
(h_U)^\# \longrightarrow \mathcal{F}
$$
Ici, la restriction de la flèche à la composante correspondant à $(U, s)$ envoie
l'élément $\text{id}_U \in h_U^\#(U)$ sur la section $s \in \mathcal{F}(U)$.
Il s'agit d'un épimorphisme d'après le Lemme \ref{sites-lemma-mono-epi-sheaves} et
notre condition sur $E$. Pour construire $\mathcal{F}_1$, posons d'abord
$\mathcal{G} = \mathcal{F}_0 \times_\mathcal{F} \mathcal{F}_0$, puis
construisons comme ci-dessus un épimorphisme $\mathcal{F}_1 \to \mathcal{G}$.
Voir le Lemme \ref{sites-lemma-coequalizer-surjection}.
\end{proof}

\begin{lemma}
\label{sites-lemma-colimit-representable}
Soit $\mathcal{C}$ un site. Soit $\mathcal{F}$ un faisceau d'ensembles sur
$\mathcal{C}$. Il existe alors un diagramme $\mathcal{I} \to \mathcal{C}$,
$i \mapsto U_i$, tel que
$$
\mathcal{F} = \colim_{i \in \mathcal{I}} h_{U_i}^\#
$$
De plus, si $E \subset \Ob(\mathcal{C})$ est un sous-ensemble tel que tout
objet de $\mathcal{C}$ admette un recouvrement par des éléments de $E$, nous pouvons
supposer que $U_i$ est un élément de $E$ pour tout $i \in \Ob(\mathcal{I})$.
\end{lemma}

\begin{proof}
Soit $\mathcal{I}$ la catégorie dont les objets sont les couples
$(U, s)$, avec $U \in \Ob(\mathcal{C})$ et $s \in \mathcal{F}(U)$,
et dont les morphismes $(U, s) \to (U', s')$ sont les morphismes $f : U \to U'$
de $\mathcal{C}$ tels que $f^*s' = s$. Pour tout objet $(U, s)$ de $\mathcal{I}$,
l'élément $s$ définit, par le lemme de Yoneda, un morphisme $s : h_U \to \mathcal{F}$
de préfaisceaux, et donc, par la propriété universelle du faisceau associé, un morphisme
$h_U^\# \to \mathcal{F}$. On voit immédiatement que ces morphismes sont compatibles
avec les morphismes de la catégorie $\mathcal{I}$. Nous obtenons ainsi
un morphisme $\colim_{(U, s)} h_U \to \mathcal{F}$ de préfaisceaux (où la
limite inductive est prise dans la catégorie des préfaisceaux) et un morphisme
$\colim_{(U, s)} (h_U)^\# \to \mathcal{F}$ de faisceaux (où la
limite inductive est prise dans la catégorie des faisceaux).
Puisque le passage au faisceau associé est adjoint à gauche du foncteur d'inclusion
$\Sh(\mathcal{C}) \to \textit{PSh}(\mathcal{C})$
(Proposition \ref{sites-proposition-sheafification-adjoint}), nous avons
$\colim (h_U)^\# = (\colim h_U)^\#$ d'après
Catégories, Lemme \ref{categories-lemma-adjoint-exact}.
Il suffit donc de montrer que $\colim_{(U, s)} h_U \to \mathcal{F}$
est un isomorphisme de préfaisceaux. Pour cela, montrons que, pour tout objet
$X$ de $\mathcal{C}$, l'application $\colim_{(U, s)} h_U(X) \to \mathcal{F}(X)$
est bijective. Elle admet en effet pour application inverse celle qui envoie l'élément
$t \in \mathcal{F}(X)$ sur l'élément de l'ensemble $\colim_{(U, s)} h_U(X)$
correspondant à $(X, t)$ et à $\text{id}_X \in h_X(X)$.

\medskip\noindent
Nous omettons la démonstration de la dernière assertion.
\end{proof}




\section{Foncteurs continus}
\label{sites-section-continuous-functors}

\begin{definition}
\label{sites-definition-continuous}
Soient $\mathcal{C}$ et $\mathcal{D}$ des sites.
Un foncteur $u : \mathcal{C} \to \mathcal{D}$ est dit
{\it continu} si, pour tout
$\{V_i \to V\}_{i\in I} \in \text{Cov}(\mathcal{C})$,
les conditions suivantes sont satisfaites :
\begin{enumerate}
\item $\{u(V_i) \to u(V)\}_{i\in I}$ appartient à $\text{Cov}(\mathcal{D})$, et
\item pour tout morphisme $T \to V$ dans $\mathcal{C}$, le morphisme
$u(T \times_V V_i) \to u(T) \times_{u(V)} u(V_i)$ est un isomorphisme.
\end{enumerate}
\end{definition}

\noindent
Rappelons qu'étant donné un foncteur $u$ comme ci-dessus et un préfaisceau d'ensembles
$\mathcal{F}$ sur $\mathcal{D}$, nous avons défini
$u^p\mathcal{F}$ comme le préfaisceau
$\mathcal{F} \circ u$, autrement dit
$$
u^p\mathcal{F} (V) = \mathcal{F}(u(V))
$$
pour tout objet $V$ de $\mathcal{C}$.

\begin{lemma}
\label{sites-lemma-pushforward-sheaf}
Soient $\mathcal{C}$ et $\mathcal{D}$ des sites.
Soit $u : \mathcal{C} \to \mathcal{D}$ un foncteur continu.
Si $\mathcal{F}$ est un faisceau sur $\mathcal{D}$, alors
$u^p\mathcal{F}$ est également un faisceau.
\end{lemma}

\begin{proof}
Soit $\{V_i \to V\}$ un recouvrement.
Par hypothèse, $\{u(V_i) \to u(V)\}$ est un recouvrement
dans $\mathcal{D}$ et $u(V_i \times_V V_j) =
u(V_i)\times_{u(V)}u(V_j)$. La condition de faisceau pour
$u^p\mathcal{F}$ et le recouvrement $\{V_i \to V\}$
est donc précisément la même que la condition de faisceau pour $\mathcal{F}$
et le recouvrement $\{u(V_i) \to u(V)\}$.
\end{proof}

\noindent
Pour éviter toute confusion, nous notons parfois
$$
u^s :
\Sh(\mathcal{D})
\longrightarrow
\Sh(\mathcal{C})
$$
la restriction du foncteur $u^p$ à la sous-catégorie des faisceaux d'ensembles.
Rappelons que $u^p$ admet un adjoint à gauche
$u_p : \textit{PSh}(\mathcal{C}) \to \textit{PSh}(\mathcal{D})$ ; voir
la section \ref{sites-section-functoriality-PSh}.

\begin{lemma}
\label{sites-lemma-adjoint-sheaves}
Dans la situation du Lemme \ref{sites-lemma-pushforward-sheaf},
le foncteur $u_s : \mathcal{G} \mapsto (u_p \mathcal{G})^\#$
est adjoint à gauche de $u^s$.
\end{lemma}

\begin{proof}
Cela résulte directement du Lemme \ref{sites-lemma-adjoints-u} et de la
Proposition \ref{sites-proposition-sheafification-adjoint}.
\end{proof}

\noindent
Voici un lemme technique.

\begin{lemma}
\label{sites-lemma-technical-up}
Dans la situation du Lemme \ref{sites-lemma-pushforward-sheaf},
pour tout préfaisceau $\mathcal{G}$ sur $\mathcal{C}$,
nous avons $(u_p\mathcal{G})^\# = (u_p(\mathcal{G}^\#))^\#$.
\end{lemma}

\begin{proof}
Pour tout faisceau $\mathcal{F}$ sur $\mathcal{D}$, nous avons
\begin{eqnarray*}
\Mor_{\Sh(\mathcal{D})}(u_s(\mathcal{G}^\#), \mathcal{F})
& = &
\Mor_{\Sh(\mathcal{C})}(\mathcal{G}^\#, u^s\mathcal{F}) \\
& = &
\Mor_{\textit{PSh}(\mathcal{C})}(\mathcal{G}^\#, u^p\mathcal{F}) \\
& = &
\Mor_{\textit{PSh}(\mathcal{C})}(\mathcal{G}, u^p\mathcal{F}) \\
& = &
\Mor_{\textit{PSh}(\mathcal{D})}(u_p\mathcal{G}, \mathcal{F}) \\
& = &
\Mor_{\Sh(\mathcal{D})}((u_p\mathcal{G})^\#, \mathcal{F})
\end{eqnarray*}
et le résultat découle du lemme de Yoneda.
\end{proof}

\begin{lemma}
\label{sites-lemma-pullback-representable-sheaf}
Soit $u : \mathcal{C} \to \mathcal{D}$ un foncteur continu
entre sites.
Pour tout objet $U$ de $\mathcal{C}$, nous avons $u_sh_U^\# = h_{u(U)}^\#$.
\end{lemma}

\begin{proof}
Cela résulte des
Lemmes \ref{sites-lemma-pullback-representable-presheaf}
et \ref{sites-lemma-technical-up}.
\end{proof}


\begin{remark}
\label{sites-remark-quasi-continuous}
(À omettre en première lecture.)
Soient $\mathcal{C}$ et $\mathcal{D}$ des sites. Utilisons
la définition des familles de morphismes tautologiquement équivalentes,
voir la Définition \ref{sites-definition-combinatorial-tautological},
pour affaiblir (légèrement) les conditions qui définissent la continuité.
Soit $u : \mathcal{C} \to \mathcal{D}$ un foncteur.
Disons que $u$ est {\it quasi-continu} si, pour tout
$\mathcal{V} = \{V_i \to V\}_{i\in I} \in \text{Cov}(\mathcal{C})$,
les conditions suivantes sont satisfaites :
\begin{enumerate}
\item[(1')] la famille de morphismes
$\{u(V_i) \to u(V)\}_{i\in I}$ est tautologiquement équivalente
à un élément de $\text{Cov}(\mathcal{D})$, et
\item[(2)] pour tout morphisme $T \to V$ dans $\mathcal{C}$, le morphisme
$u(T \times_V V_i) \to u(T) \times_{u(V)} u(V_i)$ est un isomorphisme.
\end{enumerate}
Nous allons voir que les Lemmes \ref{sites-lemma-pushforward-sheaf}
et \ref{sites-lemma-adjoint-sheaves} restent valables lorsque
$u$ est quasi-continu.

\medskip\noindent
Remarquons d'abord que les morphismes $u(V_i) \to u(V)$ sont représentables, car
ils sont isomorphes à des morphismes représentables (d'après la première condition).
En particulier, la famille $u(\mathcal{V}) = \{u(V_i) \to u(V)\}_{i\in I}$
donne naissance à un groupe de cohomologie de {\v C}ech de degré zéro
$H^0(u(\mathcal{V}), \mathcal{F})$ pour tout préfaisceau $\mathcal{F}$ sur
$\mathcal{D}$.
Soit $\mathcal{U} = \{U_j \to u(V)\}_{j \in J}$ un élément
de $\text{Cov}(\mathcal{D})$ tautologiquement
équivalent à $\{u(V_i) \to u(V)\}_{i \in I}$. Remarquons que $u(\mathcal{V})$
est un raffinement de $\mathcal{U}$ et réciproquement. La Remarque
\ref{sites-remark-both-refine-same-H0} donne donc
$H^0(u(\mathcal{V}), \mathcal{F}) = H^0(\mathcal{U}, \mathcal{F})$.
En particulier, si $\mathcal{F}$ est un faisceau, alors
$\mathcal{F}(u(V)) = H^0(u(\mathcal{V}), \mathcal{F})$ en vertu de
la propriété de faisceau exprimée au moyen des groupes de cohomologie de
{\v C}ech de degré zéro. Nous en déduisons que $u^p\mathcal{F}$ est un faisceau
si $\mathcal{F}$ est un faisceau, puisque $H^0(\mathcal{V}, u^p\mathcal{F}) =
H^0(u(\mathcal{V}), \mathcal{F})$, lequel est, comme nous venons de l'observer,
égal à $\mathcal{F}(u(V)) = u^p\mathcal{F}(V)$. Le Lemme
\ref{sites-lemma-pushforward-sheaf} est donc valable. Le Lemme
\ref{sites-lemma-adjoint-sheaves} en découle immédiatement.
\end{remark}






\section{Morphismes de sites}
\label{sites-section-morphism-sites}

\begin{definition}
\label{sites-definition-morphism-sites}
Soient $\mathcal{C}$ et $\mathcal{D}$ des sites.
Un {\it morphisme de sites} $f : \mathcal{D} \to \mathcal{C}$
est la donnée d'un foncteur continu $u : \mathcal{C} \to \mathcal{D}$
tel que le foncteur $u_s$ soit exact.
\end{definition}

\noindent
Notons que le foncteur $u$ va dans le sens {\it opposé}
au morphisme $f$. Si $f \leftrightarrow u$ est un morphisme de sites,
nous employons les notations $f^{-1} = u_s$ et $f_* = u^s$.
Le foncteur $f^{-1}$ est appelé {\it foncteur image inverse} et
le foncteur $f_*$ {\it foncteur image directe}.
Comme en topologie, nous avons la propriété d'adjonction suivante :
$$
\Mor_{\Sh(\mathcal{D})}(f^{-1}\mathcal{G}, \mathcal{F})
=
\Mor_{\Sh(\mathcal{C})}(\mathcal{G}, f_*\mathcal{F})
$$
La motivation de cette définition vient de l'exemple suivant.

\begin{example}
\label{sites-example-continuous-map}
Soit $f : X  \to Y$ une application continue d'espaces topologiques.
Rappelons que nous disposons des sites $X_{Zar}$ et $Y_{Zar}$ ;
voir l'Exemple \ref{sites-example-site-topological}. Considérons le foncteur
$u : Y_{Zar} \to X_{Zar}$, $V \mapsto f^{-1}(V)$.
Ce foncteur est manifestement continu puisque les images inverses de
recouvrements ouverts sont des recouvrements ouverts. (En réalité, cela dépend de la
manière dont on a choisi les ensembles de recouvrements de $X_{Zar}$ et de $Y_{Zar}$.
Mais, dans tous les cas, le foncteur est quasi-continu ; voir la Remarque
\ref{sites-remark-quasi-continuous}.)
Il est facile de vérifier que le foncteur $u^s$ est égal au foncteur image directe
usuel $f_*$ de la topologie. Ainsi, puisque $u_s$ est un adjoint et que
le foncteur image inverse topologique usuel $f^{-1}$ est lui aussi un adjoint,
nous obtenons un isomorphisme canonique $f^{-1} \cong u_s$. Comme $f^{-1}$
est exact, nous en déduisons que $u_s$ est exact. Le foncteur $u$ définit donc
un morphisme de sites $f : X_{Zar} \to Y_{Zar}$, que nous pouvons également noter
$f$, puisque nous avons déjà vu que les foncteurs $u_s, u^s$ coïncident
avec les notions usuelles correspondantes.
\end{example}

\begin{example}
\label{sites-example-finer-topology}
Soit $\mathcal{C}$ une catégorie. Soient
$$
\text{Cov}(\mathcal{C}) \supset \text{Cov}'(\mathcal{C})
$$
deux ensembles de familles de morphismes de même but
qui munissent $\mathcal{C}$ d'une structure de site. Notons $\mathcal{C}_\tau$
le site correspondant à $\text{Cov}(\mathcal{C})$ et
$\mathcal{C}_{\tau'}$ le site correspondant à $\text{Cov}'(\mathcal{C})$.
Nous affirmons que le foncteur identité de $\mathcal{C}$ définit un morphisme de sites
$$
\epsilon : \mathcal{C}_\tau \longrightarrow \mathcal{C}_{\tau'}
$$
En effet, $\text{id} : \mathcal{C}_{\tau'} \to \mathcal{C}_\tau$
est continu puisque tout recouvrement pour $\tau'$ est un recouvrement pour $\tau$.
Le foncteur $\epsilon_* = \text{id}^s$ est donc le foncteur identité sur les
préfaisceaux sous-jacents. Par conséquent, l'adjoint à gauche $\epsilon^{-1}$
de $\epsilon_*$ envoie un $\tau'$-faisceau $\mathcal{F}$ sur le
$\tau$-faisceau associé à $\mathcal{F}$ (d'après la
propriété universelle du faisceau associé).
Les limites finies de $\tau'$-faisceaux coïncident avec les limites finies
de préfaisceaux (Lemme \ref{sites-lemma-limit-sheaf}) et le passage au
$\tau$-faisceau associé commute aux limites finies
(Lemme \ref{sites-lemma-sheafification-exact}).
Ainsi, $\epsilon^{-1}$ est exact à gauche. Puisque $\epsilon^{-1}$
est adjoint à gauche, il est aussi exact à droite
(Catégories, Lemme \ref{categories-lemma-exact-adjoint}).
Par conséquent, $\epsilon^{-1}$ est exact et toutes les
conditions de la Définition \ref{sites-definition-morphism-sites} sont vérifiées.
\end{example}

\begin{lemma}
\label{sites-lemma-composition-morphisms-sites}
\begin{slogan}
La composition des foncteurs de sites préserve leur continuité.
\end{slogan}
Soient $\mathcal{C}_i$, $i = 1, 2, 3$, des sites. Soient
$u : \mathcal{C}_2 \to \mathcal{C}_1$ et
$v : \mathcal{C}_3 \to \mathcal{C}_2$ des foncteurs continus
qui induisent des morphismes de sites. Alors le foncteur
$u \circ v : \mathcal{C}_3 \to \mathcal{C}_1$ est continu et
définit un morphisme de sites $\mathcal{C}_1 \to \mathcal{C}_3$.
\end{lemma}

\begin{proof}
Il résulte immédiatement des définitions que $u \circ v$ est un foncteur continu.
De plus, nous avons manifestement $(u \circ v)^p = v^p \circ u^p$, et donc
$(u \circ v)^s = v^s \circ u^s$. Les foncteurs $(u \circ v)_s$ et
$u_s \circ v_s$ sont donc tous deux adjoints à gauche de $(u \circ v)^s$. Par conséquent,
$(u \circ v)_s \cong u_s \circ v_s$, et nous concluons que $(u \circ v)_s$
est exact en tant que composé de foncteurs exacts.
\end{proof}

\begin{definition}
\label{sites-definition-composition-morphisms-sites}
Soient $\mathcal{C}_i$, $i = 1, 2, 3$, des sites. Soient
$f : \mathcal{C}_1 \to \mathcal{C}_2$ et
$g : \mathcal{C}_2 \to \mathcal{C}_3$ des morphismes de sites
donnés par les foncteurs continus $u : \mathcal{C}_2 \to \mathcal{C}_1$
et $v : \mathcal{C}_3 \to \mathcal{C}_2$. Le {\it composé}
$g \circ f$ est le morphisme de sites correspondant au
foncteur $u \circ v$.
\end{definition}

\noindent
Dans cette situation, nous avons $(g \circ f)_* = g_* \circ f_*$ et
$(g \circ f)^{-1} = f^{-1} \circ g^{-1}$ (voir la démonstration du
Lemme \ref{sites-lemma-composition-morphisms-sites}).

\begin{lemma}
\label{sites-lemma-directed-morphism}
Soient $\mathcal{C}$ et $\mathcal{D}$ des sites. Soit
$u : \mathcal{C} \to \mathcal{D}$ un foncteur continu.
Supposons que toutes les catégories $(\mathcal{I}_V^u)^{opp}$ de la
section \ref{sites-section-functoriality-PSh}
soient filtrantes. Alors $u$ définit un morphisme de sites $\mathcal{D} \to
\mathcal{C}$ ; autrement dit, $u_s$ est exact.
\end{lemma}

\begin{proof}
Comme $u_s$ est adjoint à gauche de $u^s$, nous voyons que $u_s$ est exact à droite ;
voir Catégories, Lemme \ref{categories-lemma-exact-adjoint}.
Il suffit donc de montrer que $u_s$ est exact à gauche. Autrement dit,
nous devons montrer que $u_s$ commute aux limites finies.
Puisque les catégories $\mathcal{I}_Y^{opp}$ sont filtrantes,
$u_p$ commute aux limites finies ; voir
Catégories, Lemme \ref{categories-lemma-directed-commutes}
(on utilise également ici la description des limites dans $\textit{PSh}$ ;
voir la section \ref{sites-section-limits-colimits-PSh}).
Comme le passage au faisceau associé commute lui aussi aux limites finies
(Lemme \ref{sites-lemma-sheafification-exact}), nous concluons en utilisant
$u_s = \# \circ u_p$.
\end{proof}

\begin{proposition}
\label{sites-proposition-get-morphism}
Soient $\mathcal{C}$ et $\mathcal{D}$ des sites. Soit
$u : \mathcal{C} \to \mathcal{D}$ un foncteur continu.
Supposons en outre que :
\begin{enumerate}
\item la catégorie $\mathcal{C}$ possède un objet final $X$ et
$u(X)$ est un objet final de $\mathcal{D}$, et
\item la catégorie $\mathcal{C}$ possède des produits fibrés et
$u$ commute à ceux-ci.
\end{enumerate}
Alors $u$ définit un morphisme de sites $\mathcal{D} \to
\mathcal{C}$ ; autrement dit, $u_s$ est exact.
\end{proposition}

\begin{proof}
Cela résulte des Lemmes \ref{sites-lemma-directed} et
\ref{sites-lemma-directed-morphism}.
\end{proof}

\begin{remark}
\label{sites-remark-explain-left-exact}
Les conditions de la Proposition \ref{sites-proposition-get-morphism} ci-dessus
équivalent à dire que $u$ est exact à gauche, c'est-à-dire qu'il commute
aux limites finies. Voir
Catégories, Lemmes
\ref{categories-lemma-finite-limits-exist} et
\ref{categories-lemma-characterize-left-exact}.
Il paraît plus naturel de les formuler au moyen des objets finaux
et des produits fibrés, car cette formulation semble avoir un sens plus géométrique
dans les exemples.
\end{remark}

\noindent
Le Lemme \ref{sites-lemma-continuous-with-continuous-left-adjoint}
fournira une autre manière de démontrer qu'un foncteur continu
donne naissance à un morphisme de sites.

\begin{remark}
\label{sites-remark-quasi-continuous-morphism-sites}
(À omettre en première lecture.)
Soient $\mathcal{C}$ et $\mathcal{D}$ des sites. Par analogie avec la
Définition \ref{sites-definition-morphism-sites}, nous disons qu'un
{\it quasi-morphisme de sites $f : \mathcal{D} \to \mathcal{C}$}
est donné par un foncteur quasi-continu $u : \mathcal{C} \to \mathcal{D}$
(voir la Remarque \ref{sites-remark-quasi-continuous}) tel que $u_s$ soit exact.
L'analogue de la Proposition \ref{sites-proposition-get-morphism} dans ce
cadre s'obtient en remplaçant le mot « continu »
par « quasi-continu » et le mot
« morphisme » par « quasi-morphisme ». La démonstration est littéralement
la même.
\end{remark}

\noindent
Dans la Définition \ref{sites-definition-morphism-sites}, la condition d'exactitude de
$u_s$ ne peut être omise. Par exemple, la conclusion du lemme suivant
peut être fausse si l'on suppose seulement que $u$ est continu.

\begin{lemma}
\label{sites-lemma-morphism-of-sites-covering}
Soit $f : \mathcal{D} \to \mathcal{C}$ un morphisme de sites donné par le
foncteur $u : \mathcal{C} \to \mathcal{D}$. Pour tout objet $V$ de
$\mathcal{D}$, il existe un recouvrement $\{V_j \to V\}$ tel que, pour tout
$j$, il existe un morphisme $V_j \to u(U_j)$ pour un certain objet $U_j$
de $\mathcal{C}$.
\end{lemma}

\begin{proof}
Puisque $f^{-1} = u_s$ est exact, nous avons $f^{-1}* = *$, où $*$ désigne
l'objet final de la catégorie des faisceaux
(Exemple \ref{sites-example-singleton-sheaf}).
Comme $f^{-1}* = u_s*$ est le faisceau associé à $u_p*$, il existe
un recouvrement $\{V_j \to V\}$ tel que $(u_p*)(V_j)$ soit non vide.
Or $(u_p*)(V_j)$ est une limite inductive indexée par la catégorie
$\mathcal{I}^u_{V_j}$, dont les objets sont les morphismes $V_j \to u(U)$ ;
le lemme en découle.
\end{proof}

























\section{Topos}
\label{sites-section-topoi}

\noindent
Voici une définition du topos adaptée à nos besoins.
Un topos est la catégorie des faisceaux sur un site. Pour spécifier
un topos, il suffit donc de spécifier le site. La véritable différence entre un topos
et un site réside dans la définition des morphismes. En effet, il existe
de nombreux morphismes de topos qui ne proviennent pas de morphismes
des sites sous-jacents.

\begin{definition}[Topos]
\label{sites-definition-topos}
Un {\it topos} est la catégorie $\Sh(\mathcal{C})$ des faisceaux
sur un site $\mathcal{C}$.
\begin{enumerate}
\item Soient $\mathcal{C}$, $\mathcal{D}$ des sites.
Un {\it morphisme de topos} $f$ de $\Sh(\mathcal{D})$
vers $\Sh(\mathcal{C})$ est la donnée de deux foncteurs
$f_* : \Sh(\mathcal{D}) \to \Sh(\mathcal{C})$
et
$f^{-1} : \Sh(\mathcal{C}) \to \Sh(\mathcal{D})$
tels que :
\begin{enumerate}
\item nous ayons
$$
\Mor_{\Sh(\mathcal{D})}(f^{-1}\mathcal{G}, \mathcal{F})
=
\Mor_{\Sh(\mathcal{C})}(\mathcal{G}, f_*\mathcal{F})
$$
bifonctoriellement, et
\item le foncteur $f^{-1}$ commute aux limites finies, c'est-à-dire
qu'il est exact à gauche.
\end{enumerate}
\item Soient $\mathcal{C}$, $\mathcal{D}$, $\mathcal{E}$ des sites.
Étant donnés des morphismes de topos
$f :\Sh(\mathcal{D}) \to \Sh(\mathcal{C})$ et
$g :\Sh(\mathcal{E}) \to \Sh(\mathcal{D})$, le
{\it composé $f\circ g$} est le morphisme de topos défini
par les foncteurs
$(f \circ g)_* = f_* \circ g_*$ et
$(f \circ g)^{-1} = g^{-1} \circ f^{-1}$.
\end{enumerate}
\end{definition}

\noindent
Supposons que $\alpha : \mathcal{S}_1 \to \mathcal{S}_2$ soit une équivalence de
catégories (éventuellement « grandes »). Si $\mathcal{S}_1$, $\mathcal{S}_2$
sont des topos, poser $f_* = \alpha$ et prendre pour $f^{-1}$ un quasi-inverse
de $\alpha$ donne un morphisme de topos $f : \mathcal{S}_1 \to \mathcal{S}_2$.
Ce morphisme est en outre une équivalence dans la $2$-catégorie des topos (voir
la section \ref{sites-section-2-category}).
Il est donc légitime de dire que « $\mathcal{S}$ est un topos » si $\mathcal{S}$
est équivalente à la catégorie des faisceaux sur un site (sans être nécessairement
égale à la catégorie des faisceaux sur un site). Nous emploierons parfois
cet abus de notation.

\medskip\noindent
Le {\it topos vide} est le topos
des faisceaux sur le site $\mathcal{C}$, où $\mathcal{C}$
est la catégorie vide. Nous noterons parfois ce site $\emptyset$.
Ce site possède un unique faisceau (puisque $\emptyset$ n'a aucun
objet). Ainsi, $\Sh(\emptyset)$ est équivalente à la catégorie
qui possède un seul objet et un seul morphisme.

\medskip\noindent
Le {\it topos ponctuel} est le topos des faisceaux sur le site
$\mathcal{C}$ qui possède un unique objet $pt$, un unique morphisme
$\text{id}_{pt}$, et dont le seul recouvrement est
$\{\text{id}_{pt}\}$. Nous noterons simplement ce site $pt$.
Il est clair que la catégorie des faisceaux $ = $ la catégorie des
préfaisceaux $ = $ la catégorie des ensembles.
En formule, $\Sh(pt) = \textit{Ensembles}$.

\medskip\noindent
Soient $\mathcal{C}$ et $\mathcal{D}$ des sites. Soit
$f : \Sh(\mathcal{D}) \to \Sh(\mathcal{C})$ un morphisme de topos.
Remarquons que $f_*$ commute à toutes les limites et que
$f^{-1}$ commute à toutes les limites inductives ; voir Catégories,
Lemme \ref{categories-lemma-adjoint-exact}.
En particulier, la condition imposée à $f^{-1}$ dans la définition ci-dessus
garantit que $f^{-1}$ est exact. Les morphismes de topos sont souvent construits
à l'aide soit du Lemme \ref{sites-lemma-cocontinuous-morphism-topoi},
soit du lemme suivant.

\begin{lemma}
\label{sites-lemma-morphism-sites-topoi}
Étant donné un morphisme de sites $f : \mathcal{D} \to \mathcal{C}$
correspondant au foncteur $u : \mathcal{C} \to \mathcal{D}$,
le couple de foncteurs $(f^{-1} = u_s, f_* = u^s)$ est un morphisme de topos
$\Sh(\mathcal{D}) \to \Sh(\mathcal{C})$.
\end{lemma}

\begin{proof}
Cela résulte immédiatement de la Définition \ref{sites-definition-morphism-sites}.
\end{proof}

\begin{remark}
\label{sites-remark-pt-topos}
De nombreux sites donnent naissance au topos $\Sh(pt)$.
Voici un exemple utile. Supposons que $S$ soit un ensemble (d'ensembles)
qui contienne au moins un élément non vide. Soit $\mathcal{S}$ la
catégorie dont les objets sont les éléments de $S$ et dont les morphismes sont
les applications ensemblistes arbitraires. Supposons que $\mathcal{S}$ possède des produits fibrés.
Ce sera par exemple le cas si $S = \mathcal{P}(\text{ensemble infini})$
est l'ensemble des parties d'un ensemble infini quelconque (exercice de théorie des ensembles).
Munissons $\mathcal{S}$ d'une structure de site en déclarant que
les familles surjectives d'applications sont les recouvrements (et choisissons
un ensemble suffisamment grand convenable de familles de recouvrement, comme dans
Ensembles, section \ref{sets-section-coverings-site}).
Nous affirmons que $\Sh(\mathcal{S})$ est équivalente à la catégorie des
ensembles.

\medskip\noindent
Démontrons-le d'abord lorsque $S$ contient un singleton $e \in S$.
Dans ce cas, il existe une équivalence de topos
$i : \Sh(pt) \to \Sh(\mathcal{S})$ donnée par
les foncteurs
\begin{equation}
\label{sites-equation-sheaves-pt-sets}
i^{-1}\mathcal{F} = \mathcal{F}(e), \quad
i_*E = (U \mapsto \Mor_{\textit{Ensembles}}(U, E))
\end{equation}
En effet, supposons que $\mathcal{F}$ soit un faisceau sur $\mathcal{S}$.
Pour tout $U \in \Ob(\mathcal{S}) = S$, nous pouvons trouver un recouvrement
$\{\varphi_u : e \to U\}_{u \in U}$, où $\varphi_u$
envoie l'unique élément de $e$ sur $u \in U$. Dans ce cas, la condition de faisceau
implique que
$\mathcal{F}(U) = \prod_{u \in U} \mathcal{F}(e)$.
Autrement dit,
$\mathcal{F}(U) = \Mor_{\textit{Ensembles}}(U, \mathcal{F}(e))$.
Cette règle est en outre compatible avec les applications de restriction. Le foncteur
$$
i_* :
\textit{Ensembles} = \Sh(pt)
\longrightarrow
\Sh(\mathcal{S}), \quad
E \longmapsto (U \mapsto \Mor_{\textit{Ensembles}}(U, E))
$$
est donc une équivalence de catégories, et son inverse est le foncteur
$i^{-1}$ donné ci-dessus.

\medskip\noindent
Si $\mathcal{S}$ ne contient aucun singleton, le foncteur
$i_*$ défini ci-dessus a encore un sens. Pour montrer qu'il reste
une équivalence dans ce cas, choisissons un élément non vide $\tilde e \in S$
et une application $\varphi : \tilde e \to \tilde e$ dont l'image est un singleton.
Pour tout faisceau $\mathcal{F}$, posons
$$
\mathcal{F}(e) :=
\Im(
\mathcal{F}(\varphi) :
\mathcal{F}(\tilde e)
\longrightarrow
\mathcal{F}(\tilde e)
)
$$
et montrons qu'il s'agit d'un quasi-inverse de $i_*$. Les détails sont omis.
\end{remark}

\begin{remark}
\label{sites-remark-morphism-topoi-big}
(Questions ensemblistes relatives aux morphismes de topos. À omettre
en première lecture.)
Un morphisme de topos tel que défini ci-dessus n'est pas un ensemble, mais une classe.
Autrement dit, il est donné par une formule mathématique plutôt que par un
objet mathématique. Bien que nous puissions considérer
la classe de tous les morphismes entre deux topos donnés,
il n'est pas judicieux de l'introduire comme objet mathématique.
D'autre part, supposons que $\mathcal{C}$ et $\mathcal{D}$ soient
des sites donnés. Considérons un foncteur
$\Phi : \mathcal{C} \to \Sh(\mathcal{D})$.
Une telle donnée est un ensemble, autrement dit un objet mathématique.
Nous pouvons poser successivement les questions suivantes à propos de $\Phi$.
\begin{enumerate}
\item Est-il vrai que, pour tout faisceau $\mathcal{F}$ sur $\mathcal{D}$,
la règle
$U \mapsto \Mor_{\Sh(\mathcal{D})}(\Phi(U), \mathcal{F})$
définit un faisceau sur $\mathcal{C}$ ? Si oui, elle définit un foncteur
$\Phi_* : \Sh(\mathcal{D}) \to \Sh(\mathcal{C})$.
\item Est-il vrai que $\Phi_*$ admette un adjoint à gauche ? Si oui,
notons $\Phi^{-1}$ cet adjoint à gauche.
\item Est-il vrai que $\Phi^{-1}$ soit exact ?
\end{enumerate}
Si la réponse à la dernière question est encore « oui », nous obtenons
un morphisme de topos $(\Phi_*, \Phi^{-1})$. En outre, étant donné un
morphisme de topos $(f_*, f^{-1})$, nous pouvons poser
$\Phi(U) = f^{-1}(h_U^\#)$ et obtenir un foncteur $\Phi$ comme ci-dessus,
avec $f_* \cong \Phi_*$ et $f^{-1} \cong \Phi^{-1}$ (de manière compatible
avec la propriété d'adjonction).
En définitive, en travaillant avec la classe des $\Phi$
plutôt qu'avec les morphismes de topos, nous avons (a) remplacé la notion de
morphisme de topos par un objet mathématique et (b)
la classe des $\Phi$ est une classe (et non une classe
de classes). On peut bien sûr aller plus loin, par exemple en précisant
la signification des conditions (2) et (3) ci-dessus ;
nous le faisons pour les points de topos dans la section \ref{sites-section-points}.
\end{remark}

\begin{remark}
\label{sites-remark-quasi-continuous-morphism-topoi}
(À omettre en première lecture.)
Soient $\mathcal{C}$ et $\mathcal{D}$ des sites.
Un quasi-morphisme de sites $f : \mathcal{D} \to \mathcal{C}$
(voir la Remarque \ref{sites-remark-quasi-continuous-morphism-sites})
donne naissance, exactement comme dans le
Lemme \ref{sites-lemma-morphism-sites-topoi}, à un morphisme de topos $f$ de
$\Sh(\mathcal{D})$ vers $\Sh(\mathcal{C})$.
\end{remark}










\section{G-ensembles et morphismes}
\label{sites-section-G-sets-morphisms}

\noindent
Soit $\varphi : G \to H$ un homomorphisme de groupes.
Choisissons des sites (convenables) $\mathcal{T}_G$ et $\mathcal{T}_H$ comme dans
l'Exemple \ref{sites-example-site-on-group} et la
section \ref{sites-section-example-sheaf-G-sets}.
Soit $u : \mathcal{T}_H \to \mathcal{T}_G$ le foncteur qui associe
à un $H$-ensemble $U$ le $G$-ensemble $U_\varphi$, de même ensemble sous-jacent,
mais où l'action de $G$ est définie par $g \cdot \xi = \varphi(g)\xi$ pour
$g \in G$ et $\xi \in U$.
Il est clair que $u$ commute aux limites finies et qu'il est
continu\footnote{Remarque ensembliste : choisissons d'abord $\mathcal{T}_H$.
Choisissons ensuite $\mathcal{T}_G$ contenant $u(\mathcal{T}_H)$ et de sorte que
tout recouvrement de $\mathcal{T}_H$ corresponde à un recouvrement de
$\mathcal{T}_G$. Cela est possible d'après
Ensembles, Lemmes
\ref{sets-lemma-sets-with-group-action},
\ref{sets-lemma-what-is-in-it-G-sets} et
\ref{sets-lemma-coverings-site}.}.
En appliquant
la Proposition \ref{sites-proposition-get-morphism}
et le
Lemme \ref{sites-lemma-morphism-sites-topoi},
nous obtenons un morphisme de topos
$$
f : \Sh(\mathcal{T}_G) \longrightarrow \Sh(\mathcal{T}_H)
$$
associé à $\varphi$. La
Proposition \ref{sites-proposition-sheaves-on-group}
montre que nous obtenons un couple de foncteurs adjoints
$$
f_* : G\textit{-Ensembles} \longrightarrow H\textit{-Ensembles}, \quad
f^{-1} : H\textit{-Ensembles} \longrightarrow G\textit{-Ensembles}.
$$
Déterminons explicitement ces foncteurs dans ce cas.

\medskip\noindent
Commençons par établir une formule pour $f_*$.
Rappelons qu'au $G$-ensemble $S$ correspond le faisceau
$\mathcal{F}_S$ sur $\mathcal{T}_G$ donné par la règle
$\mathcal{F}_S(U) = \Mor_G(U, S)$. D'autre part, au
faisceau $\mathcal{G}$ sur $\mathcal{T}_H$ correspond le $H$-ensemble
donné par la règle $\mathcal{G}({}_HH)$. Nous obtenons donc
$$
f_*S = \Mor_{G\textit{-Ensembles}}(({}_HH)_\varphi, S).
$$
En développant un peu cette formule, nous obtenons
$$
f_*S = \{ a : H \to S \mid a(gh) = ga(h) \}
$$
avec l'action à gauche de $H$ donnée par
$(h \cdot a)(h') = a(h'h)$
pour tout élément $a \in f_*S$.

\medskip\noindent
Calculons ensuite explicitement $f^{-1}$. Comme la topologie
de $\mathcal{T}_G$ et celle de $\mathcal{T}_H$ sont sous-canoniques, tous les
préfaisceaux représentables sont des faisceaux. De plus, pour tout objet $V$
de $\mathcal{T}_H$, nous avons $f^{-1}h_V$ égal à $h_{u(V)}$ (voir le
Lemme \ref{sites-lemma-pullback-representable-sheaf}). Ainsi,
$f^{-1}S = S_\varphi$ pour les faisceaux représentables. Comme tout faisceau sur
$\mathcal{T}_H$ est une somme disjointe de faisceaux représentables, nous concluons que
cette formule vaut en général. Par conséquent, pour tout $H$-ensemble $T$,
$$
f^{-1}T = T_\varphi.
$$
L'adjonction entre $f^{-1}$ et $f_*$ est exprimée par la formule
$$
\Mor_{G\textit{-Ensembles}}(T_\varphi, S) =
\Mor_{H\textit{-Ensembles}}(T, f_*S)
$$
où $f_*S$ est défini ci-dessus. On peut le démontrer directement. Il est alors clair,
en outre, que $(f^{-1}, f_*)$ forme un couple adjoint et que $f^{-1}$ est exact.
On peut donc, au lieu de procéder comme ci-dessus, définir directement de cette
manière le morphisme de topos
$f : G\textit{-Ensembles} \to H\textit{-Ensembles}$.












\section{Objets quasi-compacts et limites inductives}
\label{sites-section-quasi-compact}

\noindent
Afin de pouvoir employer le même langage que pour les espaces topologiques,
nous introduisons la terminologie suivante.

\begin{definition}
\label{sites-definition-quasi-compact}
Soit $\mathcal{C}$ un site. Un objet $U$ de $\mathcal{C}$ est
{\it quasi-compact} si, pour tout recouvrement $\mathcal{U} = \{U_i \to U\}_{i \in I}$
de $\mathcal{C}$, il existe un autre recouvrement
$\mathcal{V} = \{V_j \to U\}_{j \in J}$ et un morphisme
$\mathcal{V} \to \mathcal{U}$ de familles de morphismes de même but
donné par $\text{id} : U \to U$, $\alpha : J \to I$ et $V_j \to U_{\alpha(j)}$
(voir la Définition \ref{sites-definition-morphism-coverings}),
tels que l'image de $\alpha$ soit une partie finie de $I$.
\end{definition}

\noindent
Bien entendu, la notion usuelle suffit pour conclure que
$U$ est quasi-compact.

\begin{lemma}
\label{sites-lemma-conclude-quasi-compact}
Soit $\mathcal{C}$ un site. Soit $U$ un objet de $\mathcal{C}$.
Considérons les conditions suivantes :
\begin{enumerate}
\item $U$ est quasi-compact,
\item pour tout recouvrement $\{U_i \to U\}_{i \in I}$ de $\mathcal{C}$,
il existe un recouvrement fini $\{V_j \to U\}_{j = 1, \ldots, m}$
de $\mathcal{C}$ qui raffine $\mathcal{U}$, et
\item pour tout recouvrement $\{U_i \to U\}_{i \in I}$ de $\mathcal{C}$,
il existe une partie finie $I' \subset I$ telle que
$\{U_i \to U\}_{i \in I'}$ soit un recouvrement de $\mathcal{C}$.
\end{enumerate}
Nous avons toujours (3) $\Rightarrow$ (2) $\Rightarrow$ (1),
mais les implications réciproques sont fausses en général.
\end{lemma}

\begin{proof}
Les implications résultent immédiatement des définitions.
Soit $X = [0, 1] \subset \mathbf{R}$
muni de sa topologie usuelle (la topologie usuelle $\epsilon$-$\delta$).
Soit $\mathcal{C}$ la catégorie des sous-espaces ouverts de $X$, les
morphismes étant les inclusions, munie des recouvrements ouverts usuels (comparer à
l'Exemple \ref{sites-example-site-topological}). Modifions cependant la notion
de recouvrement de $\mathcal{C}$ en excluant tous les recouvrements finis, sauf
ceux de la forme $\{U \to U\}$. On vérifie facilement que l'on obtient ainsi
un site au sens de la Définition \ref{sites-definition-site}.
Dans ce site, l'objet $X = U = [0, 1]$ est quasi-compact au sens de la
Définition \ref{sites-definition-quasi-compact}, mais $U$ ne satisfait pas (2).
Nous laissons au lecteur le soin de construire un exemple où (2) est vraie et (3) fausse.
\end{proof}

\noindent
Voici la signification, en théorie des topos, d'un objet quasi-compact.

\begin{lemma}
\label{sites-lemma-quasi-compact}
Soit $\mathcal{C}$ un site. Soit $U$ un objet de $\mathcal{C}$.
Les conditions suivantes sont équivalentes :
\begin{enumerate}
\item $U$ est quasi-compact, et
\item pour tout morphisme surjectif de faisceaux
$\coprod_{i \in I} \mathcal{F}_i \to h_U^\#$,
il existe une partie finie $J \subset I$ telle que
$\coprod_{i \in J} \mathcal{F}_i \to h_U^\#$ soit surjectif.
\end{enumerate}
\end{lemma}

\begin{proof}
Supposons (1) et soit $\coprod_{i \in I} \mathcal{F}_i \to h_U^\#$
un morphisme surjectif. Alors $\text{id}_U$ est une section de
$h_U^\#$ sur $U$. Il existe donc un recouvrement
$\{U_a \to U\}_{a \in A}$ et, pour tout $a \in A$,
une section $s_a$ de $\coprod_{i \in I} \mathcal{F}_i$
sur $U_a$ qui s'envoie sur $\text{id}_U$. Par la construction des sommes disjointes
comme faisceaux associés aux sommes disjointes de préfaisceaux
(Lemme \ref{sites-lemma-colimit-sheaves}), il existe, pour tout $a$,
un recouvrement $\{U_{ab} \to U_a\}_{b \in B_a}$ et,
pour tout $b \in B_a$, un indice $\iota(b) \in I$ et une section
$s_{b}$ de $\mathcal{F}_{\iota(b)}$ sur $U_{ab}$
qui s'envoie sur $\text{id}_U|_{U_{ab}}$. Ainsi, après avoir remplacé
le recouvrement $\{U_a \to U\}_{a \in A}$ par
$\{U_{ab} \to U\}_{a \in A, b \in B_a}$,
nous pouvons supposer que nous disposons d'une application $\iota : A \to I$
et, pour tout $a \in A$, d'une section $s_a$ de $\mathcal{F}_{\iota(a)}$
sur $U_a$ qui s'envoie sur $\text{id}_U$.
Puisque $U$ est quasi-compact, il existe un recouvrement
$\{V_c \to U\}_{c \in C}$, une application $\alpha : C \to A$
d'image finie et des morphismes $V_c \to U_{\alpha(c)}$ au-dessus de $U$.
Alors $J = \Im(\iota \circ \alpha) \subset I$ convient,
car $\coprod_{c \in C} h_{V_c}^\# \to h_U^\#$ est surjectif
(Lemme \ref{sites-lemma-covering-surjective-after-sheafification})
et se factorise par $\coprod_{i \in J} \mathcal{F}_i \to h_U^\#$.
(Nous utilisons ici le fait que le composé
$h_{V_c}^\# \to h_{U_{\alpha(c)}}
\xrightarrow{s_{\alpha(c)}} \mathcal{F}_{\iota(\alpha(c))} \to h_U^\#$
est le morphisme $h_{V_c}^\# \to h_U^\#$ provenant du morphisme
$V_c \to U$, puisque $s_{\alpha(c)}$ s'envoie sur
$\text{id}_U|_{U_{\alpha(c)}}$.)

\medskip\noindent
Supposons (2). Soit $\{U_i \to U\}_{i \in I}$ un recouvrement.
D'après le Lemme \ref{sites-lemma-covering-surjective-after-sheafification},
$\coprod_{i \in I} h_{U_i}^\# \to h_U^\#$ est surjectif.
Il existe donc une partie finie $J \subset I$ telle que
$\coprod_{j \in J} h_{U_j}^\# \to h_U^\#$ soit surjectif.
En raisonnant comme ci-dessus, nous trouvons un recouvrement
$\{V_c \to U\}_{c \in C}$ de $U$ dans $\mathcal{C}$
et une application $\iota : C \to J$
tels que $\text{id}_U$ se relève en une section
$s_c$ de $h_{U_{\iota(c)}}^\#$ sur $V_c$.
En raffinant encore le recouvrement, nous pouvons supposer que
$s_c \in h_{U_{\iota(c)}}(V_c)$ s'envoie sur $\text{id}_U$.
Alors $s_c : V_c \to U_{\iota(c)}$ est un morphisme au-dessus de $U$,
ce qui permet de conclure.
\end{proof}

\noindent
Le lemme précédent motive la définition suivante.

\begin{definition}
\label{sites-definition-quasi-compact-topos}
Un objet $\mathcal{F}$ d'un topos $\Sh(\mathcal{C})$ est {\it quasi-compact}
si, pour tout morphisme surjectif $\coprod_{i \in I} \mathcal{F}_i \to \mathcal{F}$
de $\Sh(\mathcal{C})$, il existe une partie finie $J \subset I$ telle
que $\coprod_{i \in J} \mathcal{F}_i \to \mathcal{F}$ soit surjectif.
Un topos $\Sh(\mathcal{C})$ est dit {\it quasi-compact}
si son objet final $*$ est un objet quasi-compact.
\end{definition}

\noindent
D'après le Lemme \ref{sites-lemma-quasi-compact},
si le site $\mathcal{C}$ possède un objet final $X$, alors
$\Sh(\mathcal{C})$ est quasi-compact si et seulement si $X$ est quasi-compact.

\begin{lemma}
\label{sites-lemma-image-of-quasi-compact}
\begin{slogan}
Les morphismes surjectifs de faisceaux transmettent la quasi-compacité.
\end{slogan}
Soit $\mathcal{C}$ un site.
\begin{enumerate}
\item Si $U \to V$ est un morphisme de $\mathcal{C}$ tel que
$h_U^\# \to h_V^\#$ soit surjectif et si $U$ est quasi-compact, alors
$V$ est quasi-compact.
\item Si $\mathcal{F} \to \mathcal{G}$ est un morphisme surjectif de faisceaux
d'ensembles et si $\mathcal{F}$ est quasi-compact, alors $\mathcal{G}$
est quasi-compact.
\end{enumerate}
\end{lemma}

\begin{proof}
Omis.
\end{proof}

\begin{lemma}
\label{sites-lemma-coproduct-quasi-compact}
\begin{slogan}
Les sommes disjointes finies de faisceaux préservent la quasi-compacité.
\end{slogan}
Soit $\mathcal{C}$ un site. Si $n \geq 1$ et si
$\mathcal{F}_1, \ldots, \mathcal{F}_n$ sont des faisceaux quasi-compacts
sur $\mathcal{C}$, alors $\coprod_{i = 1, \ldots, n} \mathcal{F}_i$
est quasi-compact.
\end{lemma}

\begin{proof}
Omis.
\end{proof}

\noindent
Les deux lemmes suivants sont les analogues, pour les sites, du
Lemme \ref{sheaves-lemma-directed-colimits-sections} de Faisceaux.

\begin{lemma}
\label{sites-lemma-directed-colimits-sections}
Soit $\mathcal{C}$ un site. Soit
$\mathcal{I} \to \Sh(\mathcal{C})$, $i \mapsto \mathcal{F}_i$,
un diagramme filtrant de faisceaux d'ensembles.
Soit $U \in \Ob(\mathcal{C})$.
Considérons l'application canonique
$$
\Psi :
\colim_i \mathcal{F}_i(U)
\longrightarrow
\left(\colim_i \mathcal{F}_i\right)(U)
$$
Avec la terminologie introduite ci-dessus :
\begin{enumerate}
\item Si tous les morphismes de transition sont injectifs, alors
$\Psi$ est injective pour tout $U$.
\item Si $U$ est quasi-compact, alors $\Psi$ est injective.
\item Si $U$ est quasi-compact et si tous les morphismes de transition sont injectifs,
alors $\Psi$ est un isomorphisme.
\item Si $U$ admet un système cofinal de recouvrements
$\{U_j \to U\}_{j \in J}$ où
$J$ est fini et $U_j \times_U U_{j'}$ est quasi-compact
pour tous $j, j' \in J$, alors $\Psi$ est bijective.
\end{enumerate}
\end{lemma}

\begin{proof}
Supposons tous les morphismes de transition injectifs. Dans ce cas, le préfaisceau
$\mathcal{F}' : V \mapsto \colim_i \mathcal{F}_i(V)$ est
séparé (voir la Définition \ref{sites-definition-separated}).
Par le Lemme \ref{sites-lemma-colimit-sheaves}, nous avons
$(\mathcal{F}')^\# = \colim_i \mathcal{F}_i$.
Le Théorème \ref{sites-theorem-plus}
montre que $\mathcal{F}' \to (\mathcal{F}')^\#$ est injectif.
Cela démontre (1).

\medskip\noindent
Supposons $U$ quasi-compact. Supposons que $s \in \mathcal{F}_i(U)$ et
$s' \in \mathcal{F}_{i'}(U)$ définissent des éléments du
membre de gauche ayant la même image par $\Psi$.
Cela signifie que nous pouvons choisir un recouvrement $\{U_a \to U\}_{a \in A}$
et, pour tout $a \in A$, un indice $i_a \in I$, $i_a \geq i$, $i_a \geq i'$,
tel que $\varphi_{ii_a}(s) = \varphi_{i'i_a}(s')$.
Comme $U$ est quasi-compact, nous pouvons choisir un recouvrement
$\{V_b \to U\}_{b \in B}$, une application $\alpha : B \to A$ d'image finie,
et des morphismes $V_b \to U_{\alpha(b)}$ au-dessus de $U$.
Choisissons $i''\in I$, supérieur ou égal ($\geq$) à tous les $i_{\alpha(b)}$,
ce qui est possible puisque l'image de $\alpha$ est finie.
Nous en déduisons que $\varphi_{ii''}(s)$ et $\varphi_{i'i''}(s)$
coïncident sur $V_b$ pour tout $b \in B$, et donc que
$\varphi_{ii''}(s) = \varphi_{i'i''}(s)$. Cela démontre (2).

\medskip\noindent
Supposons $U$ quasi-compact et tous les morphismes de transition injectifs.
Soit $s$ un élément du but de $\Psi$. Il existe un recouvrement
$\{U_a \to U\}_{a \in A}$ et, pour tout $a \in A$, un indice $i_a \in I$
et une section $s_a \in \mathcal{F}_{i_a}(U_a)$
tels que $s|_{U_a}$ provienne de $s_a$ pour tout $a \in A$.
Comme $U$ est quasi-compact, nous pouvons choisir un recouvrement
$\{V_b \to U\}_{b \in B}$, une application $\alpha : B \to A$ d'image finie,
et des morphismes $V_b \to U_{\alpha(b)}$ au-dessus de $U$.
Choisissons $i \in I$, supérieur ou égal ($\geq$) à tous les $i_{\alpha(b)}$,
ce qui est possible puisque l'image de $\alpha$ est finie.
D'après (1), les sections
$s_b = \varphi_{i_{\alpha(b)} i}(s_{\alpha(b)})|_{V_b}$
coïncident sur $V_b \times_U V_{b'}$.
Elles se recollent donc en une section
$s' \in \mathcal{F}_i(U)$ qui s'envoie sur $s$ par $\Psi$.
Cela démontre (3).

\medskip\noindent
Supposons vérifiée l'hypothèse de (4). D'après le Lemme
\ref{sites-lemma-conclude-quasi-compact}, l'objet $U$ est quasi-compact ;
ainsi, $\Psi$ est injective par (2).
Pour démontrer la surjectivité, soit $s$ un élément du but de $\Psi$.
Par hypothèse, il existe un recouvrement fini
$\{U_j \to U\}_{j = 1, \ldots, m}$, tel que $U_j \times_U U_{j'}$
soit quasi-compact pour tous $1 \leq j, j' \leq m$, ainsi que,
pour tout $j$, un indice $i_j \in I$ et une section
$s_j \in \mathcal{F}_{i_j}(U_j)$ tels que $s|_{U_j}$ soit l'image de $s_j$
pour tout $j$.
Puisque $U_j \times_U U_{j'}$ est quasi-compact, nous pouvons appliquer (2) :
il existe $i_{jj'} \in I$, avec $i_{jj'} \geq i_j$ et
$i_{jj'} \geq i_{j'}$, tel que
$\varphi_{i_ji_{jj'}}(s_j)$ et $\varphi_{i_{j'}i_{jj'}}(s_{j'})$
coïncident sur $U_j \times_U U_{j'}$. Choisissons un indice $i \in I$
supérieur ou égal à tous les $i_{jj'}$. Les sections
$\varphi_{i_ji}(s_j)$ de $\mathcal{F}_i$ se recollent alors
en une section de $\mathcal{F}_i$ sur $U$. L'image de cette section
est bien l'élément $s$ voulu.
\end{proof}

\begin{lemma}
\label{sites-lemma-directed-colimits-global-sections}
Soit $\mathcal{C}$ un site. Soit
$\mathcal{I} \to \Sh(\mathcal{C})$, $i \mapsto \mathcal{F}_i$,
un diagramme filtrant de faisceaux d'ensembles.
Considérons l'application canonique
$$
\Psi :
\colim_i \Gamma(\mathcal{C}, \mathcal{F}_i)
\longrightarrow
\Gamma(\mathcal{C}, \colim_i \mathcal{F}_i)
$$
On a les assertions suivantes :
\begin{enumerate}
\item Si tous les morphismes de transition sont injectifs, alors $\Psi$ est injective.
\item Si $\Sh(\mathcal{C})$ est quasi-compact, alors $\Psi$ est injective.
\item Si $\Sh(\mathcal{C})$ est quasi-compact et si tous les morphismes
de transition sont injectifs, alors $\Psi$ est un isomorphisme.
\item Supposons qu'il existe un ensemble $S \subset \Ob(\Sh(\mathcal{C}))$
ayant les propriétés suivantes :
\begin{enumerate}
\item pour toute surjection $\mathcal{F} \to *$, il existe
$\mathcal{K} \in S$ et un morphisme $\mathcal{K} \to \mathcal{F}$
tels que $\mathcal{K} \to *$ soit surjectif ;
\item pour $\mathcal{K} \in S$, le produit
$\mathcal{K} \times \mathcal{K}$ est quasi-compact.
\end{enumerate}
Alors $\Psi$ est bijective.
\end{enumerate}
\end{lemma}

\begin{proof}
Démonstration de (1). Supposons tous les morphismes de transition injectifs.
Dans ce cas, le préfaisceau
$\mathcal{F}' : V \mapsto \colim_i \mathcal{F}_i(V)$ est
séparé (voir la Définition \ref{sites-definition-separated}).
D'après le Lemme \ref{sites-lemma-colimit-sheaves}, on a
$(\mathcal{F}')^\# = \colim_i \mathcal{F}_i$.
Le Théorème \ref{sites-theorem-plus} montre que
$\mathcal{F}' \to (\mathcal{F}')^\#$ est injectif.
Cela démontre (1).

\medskip\noindent
Démonstration de (2). Supposons $\Sh(\mathcal{C})$ quasi-compact.
Rappelons que $\Gamma(\mathcal{C}, \mathcal{F}) = \Mor(*, \mathcal{F})$
pour tout $\mathcal{F}$ dans $\Sh(\mathcal{C})$.
Soient $a_i, b_i : * \to \mathcal{F}_i$ et, pour
$i' \geq i$, notons $a_{i'}, b_{i'} : * \to \mathcal{F}_{i'}$
leurs composés avec les morphismes de transition du système.
Posons $a = \colim_{i' \geq i} a_{i'}$, et de même pour $b$.
Pour $i' \geq i$, notons
$$
E_{i'} = \text{Égalisateur}(a_{i'}, b_{i'}) \subset *
\quad\text{et}\quad
E = \text{Égalisateur}(a, b) \subset *
$$
D'après Catégories, Lemme \ref{categories-lemma-directed-commutes}, on a
$E = \colim_{i' \geq i} E_{i'}$. Il s'ensuit que
$\coprod_{i' \geq i} E_{i'} \to E$
est un morphisme surjectif de faisceaux. Par conséquent, si $E = *$,
c'est-à-dire si $a = b$, alors, puisque $*$ est quasi-compact,
on a $E_{i'} = *$ pour un certain $i' \geq i$, d'où
$a_{i'} = b_{i'}$ pour un certain $i' \geq i$.
Cela démontre (2).

\medskip\noindent
Démonstration de (3). Supposons $\Sh(\mathcal{C})$ quasi-compact et tous
les morphismes de transition injectifs. Soit
$a : * \to \colim \mathcal{F}_i$ un morphisme. Alors
$E_i = a^{-1}(\mathcal{F}_i) \subset *$ est un sous-faisceau et
$\colim E_i = *$ (d'après la référence ci-dessus). Il existe donc
un indice $i$ tel que $E_i = *$, et l'image de $a$ est alors
contenue dans $\mathcal{F}_i$, comme voulu.

\medskip\noindent
Démonstration de (4). Soit $S \subset \Ob(\Sh(\mathcal{C}))$ satisfaisant
(4)(a) et (b). En appliquant (4)(a) à $\text{id} : * \to *$, nous trouvons
$\mathcal{K} \in S$ tel que $\mathcal{K} \to *$ soit surjectif.
Les morphismes $\mathcal{K} \times \mathcal{K} \to \mathcal{K} \to *$
sont surjectifs.
D'après (4)(b) et le Lemme \ref{sites-lemma-image-of-quasi-compact},
$\mathcal{K}$ et $\Sh(\mathcal{C})$ sont quasi-compacts.
Ainsi, $\Psi$ est injective par (2).
Posons $\mathcal{F} = \colim \mathcal{F}_i$.
Soit $s : * \to \mathcal{F}$ une section globale de la limite inductive.
Comme $\coprod \mathcal{F}_i \to \mathcal{F}$ est surjectif, la projection
$$
\coprod\nolimits_{i \in I} * \times_{s, \mathcal{F}} \mathcal{F}_i \to *
$$
est surjective. D'après (4)(a), il existe $\mathcal{K} \in S$ et un morphisme
$\mathcal{K} \to \coprod_{i \in I} * \times_{s, \mathcal{F}} \mathcal{F}_i$
tels que $\mathcal{K} \to *$ soit surjectif.
Puisque $\mathcal{K}$ est quasi-compact, ce morphisme se factorise par
$\mathcal{K} \to \coprod_{i' \in I'} * \times_{s, \mathcal{F}} \mathcal{F}_{i'}$
pour une partie finie $I' \subset I$. Soit $i \in I$ un majorant
de la partie finie $I'$. Les morphismes de transition définissent un morphisme
$\coprod_{i' \in I'} \mathcal{F}_{i'} \to \mathcal{F}_i$.
Celui-ci fournit à son tour un morphisme
$\mathcal{K} \to * \times_{s, \mathcal{F}} \mathcal{F}_i$.
Autrement dit, nous obtenons $\mathcal{K} \in S$, avec
$\mathcal{K} \to *$ surjectif, et un diagramme commutatif
$$
\xymatrix{
\mathcal{K} \times \mathcal{K} \ar@<1ex>[r] \ar@<-1ex>[r] &
\mathcal{K} \ar[d] \ar[r] & {*} \ar[d]^s \\
& \mathcal{F}_i \ar[r] & \mathcal{F} \ar@{=}[r] & \colim \mathcal{F}_i
}
$$
La ligne supérieure de ce diagramme est un coégalisateur.
Il suffit donc de montrer qu'après avoir augmenté $i$, les deux morphismes induits
$a_i, b_i : \mathcal{K} \times \mathcal{K} \to \mathcal{F}_i$
sont égaux. C'est ce que démontre le paragraphe suivant,
par exactement le même argument que dans la démonstration de (2) ;
nous invitons le lecteur à omettre le reste de la démonstration.

\medskip\noindent
Pour $i' \geq i$, notons
$a_{i'}, b_{i'} : \mathcal{K} \times \mathcal{K} \to \mathcal{F}_{i'}$
les composés de $a_i, b_i$ avec les morphismes de transition du système.
Posons $a = \colim_{i' \geq i} a_{i'} :
\mathcal{K} \times \mathcal{K} \to \mathcal{F}$, et de même pour $b$.
La commutativité du diagramme ci-dessus donne $a = b$.
Pour $i' \geq i$, notons
$$
E_{i'} = \text{Égalisateur}(a_{i'}, b_{i'}) \subset
\mathcal{K} \times \mathcal{K}
\quad\text{et}\quad
E = \text{Égalisateur}(a, b) \subset
\mathcal{K} \times \mathcal{K}
$$
D'après Catégories, Lemme \ref{categories-lemma-directed-commutes}, on a
$E = \colim_{i' \geq i} E_{i'}$. Il s'ensuit que
$\coprod_{i' \geq i} E_{i'} \to E$
est un morphisme surjectif de faisceaux. Comme $a = b$, on a
$E = \mathcal{K} \times \mathcal{K}$.
Puisque $\mathcal{K} \times \mathcal{K}$ est quasi-compact par (4)(b),
on a $E_{i'} = \mathcal{K} \times \mathcal{K}$ pour un certain
$i' \geq i$, et donc $a_{i'} = b_{i'}$ pour un certain $i' \geq i$.
\end{proof}

\begin{remark}
\label{sites-remark-stronger-conditions}
Soit $\mathcal{C}$ un site. Plusieurs conditions permettent d'assurer
les hypothèses de la partie (4) du
Lemme \ref{sites-lemma-directed-colimits-global-sections}.
En voici quelques-unes.
\begin{enumerate}
\item Supposons qu'il existe un ensemble $\mathcal{B} \subset \Ob(\mathcal{C})$
ayant les propriétés suivantes :
\begin{enumerate}
\item pour toute surjection $\mathcal{F} \to *$, il existe
$m \geq 0$ et $U_1, \ldots, U_m \in \mathcal{B}$ tels que
$\mathcal{F}(U_j)$ soit non vide et $\coprod h_{U_j}^\# \to *$ surjectif ;
\item pour $U, U' \in \mathcal{B}$, le faisceau
$h_U^\# \times h_{U'}^\#$ est quasi-compact.
\end{enumerate}
\item Supposons qu'il existe un ensemble $\mathcal{B} \subset \Ob(\mathcal{C})$
ayant les propriétés suivantes :
\begin{enumerate}
\item il existe $m \geq 0$ et $U_1, \ldots, U_m \in \mathcal{B}$ tels que
$\coprod h_{U_j}^\# \to *$ soit surjectif ;
\item pour $U \in \mathcal{B}$, tout recouvrement de $U$ peut être
raffiné par un recouvrement fini $\{U_j \to U\}_{j = 1, \ldots, m}$
avec $U_j \in \mathcal{B}$ ;
\item pour $U, U' \in \mathcal{B}$, il existe
$m \geq 0$, des objets $U_1, \ldots, U_m \in \mathcal{B}$ et
des morphismes $U_j \to U$ et $U_j \to U'$ tels que
$\coprod h_{U_j}^\# \to h_U^\# \times h_{U'}^\#$ soit surjectif.
\end{enumerate}
\item Supposons que :
\begin{enumerate}
\item $\Sh(\mathcal{C})$ soit quasi-compact ;
\item tout objet de $\mathcal{C}$ admette un recouvrement
dont les membres sont des objets quasi-compacts ;
\item si $U$ et $U'$ sont quasi-compacts, alors le faisceau
$h_U^\# \times h_{U'}^\#$ soit quasi-compact.
\end{enumerate}
\end{enumerate}
Dans les cas (1) et (2), nous prenons pour $S \subset \Ob(\Sh(\mathcal{C}))$
l'ensemble des coproduits finis des faisceaux $h_U^\#$, pour
$U \in \mathcal{B}$.
Dans le cas (3), nous prenons pour $S \subset \Ob(\Sh(\mathcal{C}))$
l'ensemble des coproduits finis des faisceaux $h_U^\#$, pour
$U \in \Ob(\mathcal{C})$ quasi-compact.
\end{remark}

\noindent
Nous aurons plus loin besoin d'une borne sur le comportement possible
des limites inductives, sous la forme suivante.

\begin{lemma}
\label{sites-lemma-colimit-over-ordinal-sections}
Soit $\mathcal{C}$ un site. Soit $\beta$ un ordinal.
Soit $\beta \to \Sh(\mathcal{C})$, $\alpha \mapsto \mathcal{F}_\alpha$,
un système de faisceaux indexé par $\beta$. Pour $U \in \Ob(\mathcal{C})$,
considérons l'application canonique
$$
\colim_{\alpha < \beta} \mathcal{F}_\alpha(U)
\longrightarrow
\left(\colim_{\alpha < \beta} \mathcal{F}_\alpha \right)(U)
$$
Si la cofinalité de $\beta$ est assez grande, cette application
est bijective pour tout $U$.
\end{lemma}

\begin{proof}
Le membre de gauche est la valeur en $U$ de la limite inductive
$\mathcal{F}_{\colim}$ prise dans la catégorie des préfaisceaux ; voir
la section \ref{sites-section-limits-colimits-PSh}.
Rappelons que $\colim_{\alpha < \beta} \mathcal{F}_\alpha$ est
le faisceau associé $\mathcal{F}_{\colim}^\#$ à $\mathcal{F}_{\colim}$ ;
voir le Lemme \ref{sites-lemma-colimit-sheaves}.
Soit $\mathcal{U} = \{U_i \to U\}_{i \in I}$ un élément de
l'ensemble $\text{Cov}(\mathcal{C})$ des recouvrements de $\mathcal{C}$.
Si la cofinalité de $\beta$ est strictement supérieure au cardinal de $I$,
nous affirmons que
$$
H^0(\mathcal{U}, \mathcal{F}_{\colim}) =
\colim H^0(\mathcal{U}, \mathcal{F}_\alpha) =
\colim \mathcal{F}_\alpha(U) = \mathcal{F}_{\colim}(U)
$$
Les deuxième et troisième égalités sont claires. Pour la première, écrivons
$s = (s_i) \in H^0(\mathcal{U}, \mathcal{F}_{\colim})$.
Pour tout $i$, l'élément $s_i$ provient alors d'un élément
$s_{i, \alpha_i} \in \mathcal{F}_{\alpha_i}(U_i)$ pour un certain
$\alpha_i < \beta$. Grâce à l'hypothèse sur la cofinalité, nous pouvons
choisir $\alpha_i = \alpha$ indépendamment de $i$. Les éléments $s_i$ et $s_j$
ont alors la même image dans
$\mathcal{F}_{\alpha_{i, j}}(U_i \times_U U_j)$
pour un certain $\alpha_{i, j} < \beta$. Comme le cardinal de
$I \times I$ est lui aussi strictement inférieur à la cofinalité de $\beta$,
nous pouvons, quitte à augmenter $\alpha$, supposer
$\alpha_{i, j} = \alpha$ pour tous $i, j$. Cela démontre la surjectivité
de l'application naturelle
$\colim H^0(\mathcal{U}, \mathcal{F}_\alpha)
\to H^0(\mathcal{U}, \mathcal{F}_{\colim})$.
Un argument tout à fait semblable montre qu'elle est injective.
En particulier, $\mathcal{F}_{\colim}$ satisfait la condition de faisceau
pour $\mathcal{U}$.
Ainsi, si la cofinalité de $\beta$ est strictement supérieure au supremum
des cardinaux des ensembles d'indices $I$ des recouvrements, la conclusion suit.
\end{proof}

\section{Limites inductives de sites}
\label{sites-section-colimit-sites}

\noindent
Il nous faut un analogue du Lemme \ref{sites-lemma-directed-colimits-sections}
lorsque le site est la limite d'un système projectif de sites.
Par souci de simplicité, nous n'expliquons la construction que dans le cas où
les ensembles d'indices des recouvrements sont finis.

\begin{situation}
\label{sites-situation-inverse-limit-sites}
On se donne ici :
\begin{enumerate}
\item une catégorie d'indices cofiltrante $\mathcal{I}$ ;
\item pour $i \in \Ob(\mathcal{I})$, un site $\mathcal{C}_i$ tel que tout
recouvrement de $\mathcal{C}_i$ possède un ensemble d'indices fini ;
\item pour un morphisme $a : i \to j$ de $\mathcal{I}$, un morphisme de sites
$f_a : \mathcal{C}_i \to \mathcal{C}_j$ défini par un foncteur continu
$u_a : \mathcal{C}_j \to \mathcal{C}_i$,
\end{enumerate}
tels que $f_a \circ f_b = f_c$ dès que $c = a \circ b$ dans $\mathcal{I}$.
\end{situation}

\begin{lemma}
\label{sites-lemma-colimit-sites}
Dans la Situation \ref{sites-situation-inverse-limit-sites}, nous pouvons construire
un site $(\mathcal{C}, \text{Cov}(\mathcal{C}))$ comme suit :
\begin{enumerate}
\item comme catégorie, $\mathcal{C} = \colim \mathcal{C}_i$ ;
\item $\text{Cov}(\mathcal{C})$ est la réunion des images
de $\text{Cov}(\mathcal{C}_i)$ par $u_i : \mathcal{C}_i \to \mathcal{C}$.
\end{enumerate}
\end{lemma}

\begin{proof}
Notre définition de la composition des morphismes de sites implique que
$u_b \circ u_a = u_c$ dès que $c = a \circ b$ dans $\mathcal{I}$.
La formule $\mathcal{C} = \colim \mathcal{C}_i$ signifie que
$\Ob(\mathcal{C}) = \colim \Ob(\mathcal{C}_i)$ et
$\text{Flèches}(\mathcal{C}) = \colim \text{Flèches}(\mathcal{C}_i)$.
La source, le but et la composition sont alors hérités de la source,
du but et de la composition sur $\text{Flèches}(\mathcal{C}_i)$.
Nous obtenons ainsi une catégorie.
Notons $u_i : \mathcal{C}_i \to \mathcal{C}$ le foncteur évident.
Remarquons que, pour tout diagramme fini dans $\mathcal{C}$,
il existe un $i$ tel que ce diagramme soit l'image d'un diagramme
dans $\mathcal{C}_i$.

\medskip\noindent
Soit $\{U^t \to U\}$ un recouvrement de $\mathcal{C}$. Montrons d'abord que,
si $V \to U$ est un morphisme de $\mathcal{C}$, alors $U^t \times_U V$ existe.
D'après la remarque précédente et notre définition des recouvrements,
nous pouvons trouver un $i$, un recouvrement $\{U_i^t \to U_i\}$ de
$\mathcal{C}_i$ et un morphisme $V_i \to U_i$ dont l'image par $u_i$
est la donnée considérée. Nous affirmons que $U^t \times_U V$ est l'image de
$U^t_i \times_{U_i} V_i$ par $u_i$.
En effet, pour tout $a : j \to i$ dans $\mathcal{I}$, le foncteur $u_a$
est continu, et donc
$u_a(U^t_i \times_{U_i} V_i) = u_a(U^t_i) \times_{u_a(U_i)} u_a(V_i)$.
En particulier, nous pouvons remplacer $i$ par $j$ si nous le souhaitons.
Ainsi, si $W$ est un autre objet de $\mathcal{C}$, nous pouvons supposer
$W = u_i(W_i)$, et nous obtenons
\begin{align*}
& \Mor_\mathcal{C}(W, u_i(U^t_i \times_{U_i} V_i)) \\
& =
\colim_{a : j \to i}
\Mor_{\mathcal{C}_j}(u_a(W_i), u_a(U^t_i \times_{U_i} V_i)) \\
& =
\colim_{a : j \to i}
\Mor_{\mathcal{C}_j}(u_a(W_i), u_a(U^t_i))
\times_{\Mor_{\mathcal{C}_j}(u_a(W_i), u_a(U_i))}
\Mor_{\mathcal{C}_j}(u_a(W_i), u_a(V_i)) \\
& =
\Mor_\mathcal{C}(W, U^t)
\times_{\Mor_\mathcal{C}(W, U)}
\Mor_\mathcal{C}(W, V)
\end{align*}
puisque les limites inductives filtrantes commutent aux limites finies
(Catégories, Lemme \ref{categories-lemma-directed-commutes}).
Il en résulte aussi que
$\{U^t \times_U V \to V\}$ est un recouvrement dans $\mathcal{C}$.
Nous voyons ainsi que l'axiome (3) de la Définition \ref{sites-definition-site}
est satisfait.

\medskip\noindent
Pour vérifier l'axiome (2) de la Définition \ref{sites-definition-site},
soit $\{U^t \to U\}_{t \in T}$ un recouvrement de $\mathcal{C}$ et,
pour tout $t$, soit $\{U^{ts} \to U^t\}$ un recouvrement de
$\mathcal{C}$. Nous pouvons alors trouver un $i$ et un recouvrement
$\{U^t_i \to U_i\}_{t \in T}$ de $\mathcal{C}_i$ dont l'image par $u_i$ est
$\{U^t \to U\}$. Comme $T$ est {\bf fini}, nous pouvons choisir
$a : j \to i$ dans $\mathcal{I}$ et des recouvrements
$\{U^{ts}_j \to u_a(U^t_i)\}$ de $\mathcal{C}_j$ dont l'image par $u_j$
donne $\{U^{ts} \to U^t\}$.
Nous en déduisons que $\{U^{ts} \to U\}$ est un recouvrement de $\mathcal{C}$
en appliquant l'axiome (2) au site $\mathcal{C}_j$.

\medskip\noindent
Nous omettons la démonstration de l'axiome (1) de la
Définition \ref{sites-definition-site}.
\end{proof}

\begin{lemma}
\label{sites-lemma-compute-pullback-to-limit}
Dans la Situation \ref{sites-situation-inverse-limit-sites}, soit
$u_i : \mathcal{C}_i \to \mathcal{C}$ le foncteur construit au
Lemme \ref{sites-lemma-colimit-sites}. Alors $u_i$ définit un morphisme
de sites $f_i : \mathcal{C} \to \mathcal{C}_i$. Pour
$U_i \in \Ob(\mathcal{C}_i)$ et un faisceau $\mathcal{F}$ sur $\mathcal{C}_i$,
on a
\begin{equation}
\label{sites-equation-compute-pullback-to-limit}
f_i^{-1}\mathcal{F}(u_i(U_i)) =
\colim_{a : j \to i} f_a^{-1}\mathcal{F}(u_a(U_i))
\end{equation}
\end{lemma}

\begin{proof}
Les arguments de la démonstration du Lemme \ref{sites-lemma-colimit-sites}
montrent immédiatement que les foncteurs $u_i$ sont continus.
Pour achever la démonstration, il faut montrer que
$f_i^{-1} := u_{i, s}$ est un foncteur exact
$\Sh(\mathcal{C}_i) \to \Sh(\mathcal{C})$.
Il suffit en fait de montrer que $f_i^{-1}$ est exact à gauche,
car il est exact à droite en tant qu'adjoint à gauche
(Catégories, Lemme \ref{categories-lemma-exact-adjoint}).
Nous démontrons d'abord (\ref{sites-equation-compute-pullback-to-limit}),
puis nous en déduisons l'exactitude.

\medskip\noindent
Pour un objet arbitraire $V$ de $\mathcal{C}$, nous pouvons choisir
un $a : j \to i$ et un objet $V_j \in \Ob(\mathcal{C})$ tels que
$V = u_j(V_j)$. Nous pouvons alors poser
$$
\mathcal{G}(V) = \colim_{b : k \to j} f_{a \circ b}^{-1}\mathcal{F}(u_b(V_j))
$$
La valeur $\mathcal{G}(V)$ de la limite inductive ne dépend ni du choix
de $b : j \to i$, ni de l'objet $V_j$ tel que $u_j(V_j) = V$ ;
nous omettons la vérification. En outre, si
$\alpha : V \to V'$ est un morphisme de $\mathcal{C}$, nous pouvons choisir
$b : j \to i$ et un morphisme $\alpha_j : V_j \to V'_j$ tels que
$u_j(\alpha_j) = \alpha$. Celui-ci induit une application
$\mathcal{G}(V') \to \mathcal{G}(V)$ à l'aide des restrictions le long des
morphismes $u_b(\alpha_j) : u_b(V_j) \to u_b(V'_j)$. Une vérification montre
que $\mathcal{G}$ est un préfaisceau (nous l'omettons).
En fait, $\mathcal{G}$ satisfait la condition de faisceau. En effet,
tout recouvrement $\mathcal{U} = \{U^t \to U\}$ de $\mathcal{C}$
provient d'un niveau fini. Supposons que
$\mathcal{U}_j = \{U^t_j \to U_j\}$ soit envoyé sur $\mathcal{U}$ par $u_j$
pour un certain $a : j \to i$ de $\mathcal{I}$. Alors
$$
H^0(\mathcal{U}, \mathcal{G}) =
\colim_{b : k \to j} H^0(u_b(\mathcal{U}_j), f_{b \circ a}^{-1}\mathcal{F}) =
\colim_{b : k \to j} f_{b \circ a}^{-1}\mathcal{F}(u_b(U_j)) =
\mathcal{G}(U)
$$
comme voulu. La première égalité est vraie parce que les limites inductives
filtrantes commutent aux limites finies
(Catégories, Lemme \ref{categories-lemma-directed-commutes}).
Par construction, $\mathcal{G}(U)$ est donné par le membre de droite de
(\ref{sites-equation-compute-pullback-to-limit}).
Ainsi, (\ref{sites-equation-compute-pullback-to-limit}) est vraie si nous pouvons
montrer que $\mathcal{G}$ est égal à $f_i^{-1}\mathcal{F}$.

\medskip\noindent
Vérifions dans ce paragraphe que $\mathcal{G}$ est canoniquement isomorphe à
$f_i^{-1}\mathcal{F}$. Nous recommandons vivement au lecteur de sauter ce paragraphe.
Pour cette vérification, nous devons établir une bijection
$\Mor_{\Sh(\mathcal{C})}(\mathcal{G}, \mathcal{H}) =
\Mor_{\Sh(\mathcal{C}_i)}(\mathcal{F}, f_{i, *}\mathcal{H})$
fonctorielle en le faisceau $\mathcal{H}$ sur $\mathcal{C}$,
où $f_{i, *} = u_i^p$. Un morphisme
$\mathcal{G} \to \mathcal{H}$ équivaut à un système compatible de morphismes
$$
\varphi_{a, b, V_j} :
f_{a \circ b}^{-1}\mathcal{F}(u_b(V_j))
\longrightarrow
\mathcal{H}(u_j(V_j))
$$
pour tous $a : j \to i$, $b : k \to j$ et
$V_j \in \Ob(\mathcal{C}_j)$.
Les compatibilités imposent que les morphismes $\varphi_{a, b, V_j}$ soient
égaux à $\varphi_{a \circ b, \text{id}, u_b(V_j)}$. Pour $a : j \to i$ fixé,
la famille de morphismes $\varphi_{a, \text{id}, V_j}$ correspond à un morphisme
de faisceaux $\varphi_a : f_a^{-1}\mathcal{F} \to f_{j, *}\mathcal{H}$.
Les compatibilités entre les $\varphi_{a, \text{id}, u_a(V_i)}$ et les
$\varphi_{\text{id}, \text{id}, V_i}$ impliquent que $\varphi_a$
est l'adjoint du morphisme $\varphi_{id}$ par les égalités
$$
\Mor_{\Sh(\mathcal{C}_j)}(f_a^{-1}\mathcal{F}, f_{j, *}\mathcal{H}) =
\Mor_{\Sh(\mathcal{C}_i)}(\mathcal{F}, f_{a, *}f_{j, *}\mathcal{H}) =
\Mor_{\Sh(\mathcal{C}_i)}(\mathcal{F}, f_{i, *}\mathcal{H})
$$
Nous voyons finalement que tout le système de morphismes
$\varphi_{a, b, V_j}$ est déterminé par le morphisme
$\varphi_{id} : \mathcal{F} \to f_{i, *}\mathcal{H}$.
Réciproquement, étant donné un tel morphisme
$\psi : \mathcal{F} \to f_{i, *}\mathcal{H}$, nous pouvons remonter
l'argument précédent pour construire la famille de morphismes
$\varphi_{a, b, V_j}$. Cela achève la démonstration de l'égalité
$\mathcal{G} = f_i^{-1}\mathcal{F}$.

\medskip\noindent
Supposons (\ref{sites-equation-compute-pullback-to-limit}) vérifiée. Alors le foncteur
$\mathcal{F} \mapsto f_i^{-1}\mathcal{F}(U)$ commute aux limites finies,
car les limites finies de faisceaux se calculent dans la catégorie des préfaisceaux
(Lemme \ref{sites-lemma-limit-sheaf}), les foncteurs $f_a^{-1}$ commutent aux limites
finies, et les limites inductives filtrantes commutent aux limites finies.
Pour voir que $\mathcal{F} \mapsto f_i^{-1}\mathcal{F}(V)$ commute aux limites
finies pour un objet général $V$ de $\mathcal{C}$, nous pouvons employer le même
argument avec la formule donnée ci-dessus pour
$f_i^{-1}\mathcal{F}(V) = \mathcal{G}(V)$.
Ainsi, $f_i^{-1}$ est exact à gauche, ce qui achève la démonstration du lemme.
\end{proof}

\begin{lemma}
\label{sites-lemma-colimit}
Dans la Situation \ref{sites-situation-inverse-limit-sites}, supposons donnés :
\begin{enumerate}
\item un faisceau $\mathcal{F}_i$ sur $\mathcal{C}_i$ pour tout
$i \in \Ob(\mathcal{I})$ ;
\item pour $a : j \to i$, un morphisme
$\varphi_a : f_a^{-1}\mathcal{F}_i \to \mathcal{F}_j$
de faisceaux sur $\mathcal{C}_j$.
\end{enumerate}
Supposons $\varphi_c = \varphi_b \circ f_b^{-1}\varphi_a$
dès que $c = a \circ b$. Posons
$\mathcal{F} = \colim f_i^{-1}\mathcal{F}_i$
sur le site $\mathcal{C}$ du Lemme \ref{sites-lemma-colimit-sites}.
Soient $i \in \Ob(\mathcal{I})$ et $X_i \in \text{Ob}(\mathcal{C}_i)$. Alors
$$
\colim_{a : j \to i} \mathcal{F}_j(u_a(X_i)) = \mathcal{F}(u_i(X_i))
$$
\end{lemma}

\begin{proof}
Un argument formel montre que
$$
\colim_{a : j \to i} \mathcal{F}_i(u_a(X_i)) =
\colim_{a : j \to i} \colim_{b : k \to j}
f_b^{-1}\mathcal{F}_j(u_{a \circ b}(X_i))
$$
D'après (\ref{sites-equation-compute-pullback-to-limit}), la limite inductive
intérieure est égale à $f_j^{-1}\mathcal{F}_j(u_i(X_i))$ ;
on conclut donc par le Lemme \ref{sites-lemma-directed-colimits-sections}.
\end{proof}

\begin{lemma}
\label{sites-lemma-colimit-push-pull}
Dans la Situation \ref{sites-situation-inverse-limit-sites}, supposons donné
un faisceau $\mathcal{F}$ sur $\mathcal{C}$. Alors
$$
\mathcal{F} = \colim f_i^{-1}f_{i, *}\mathcal{F}
$$
où les morphismes de transition sont les $f_j^{-1}\varphi_a$
pour $a : j \to i$, avec
$\varphi_a : f_a^{-1}f_{i, *}\mathcal{F} \to f_{j, *}\mathcal{F}$
un morphisme canonique satisfaisant une condition de cocycle comme au
Lemme \ref{sites-lemma-colimit}.
\end{lemma}

\begin{proof}
Comme morphisme
$$
\varphi_a : f_a^{-1}f_{i, *}\mathcal{F} \to f_{j, *}\mathcal{F}
$$
nous choisissons l'adjoint du morphisme identité
$$
f_{i, *}\mathcal{F} \to f_{a, *}f_{j, *}\mathcal{F}
$$
Ainsi, $\varphi_a$ est la co-unité de l'adjonction
$(f_a^{-1}, f_{a, *})$. Nous devons montrer que, pour tous
$a : j \to i$ et $b : k \to i$ de composé $c = a \circ b$, on a
$\varphi_c = \varphi_b \circ f_b^{-1}\varphi_a$.
Cela résulte de Catégories, Lemme \ref{categories-lemma-compose-counits}.
Enfin, il faut montrer que le morphisme fourni par adjonction
$$
\colim_{i \in I} f_i^{-1}f_{i, *}\mathcal{F}
\longrightarrow
\mathcal{F}
$$
est un isomorphisme. Pour un objet $U$ de $\mathcal{C}$,
nous devons montrer que l'application
$$
(\colim_{i \in I} f_i^{-1}\mathcal{F}_i)(U) \to \mathcal{F}(U)
$$
est bijective. Choisissons un $i$ et un objet $U_i$ de $\mathcal{C}_i$
tels que $u_i(U_i) = U$. Le membre de gauche est alors égal à
$$
(\colim_{i \in I} f_i^{-1}\mathcal{F}_i)(U) =
\colim_{a : j \to i} f_{j, *}\mathcal{F}(u_a(U_i))
$$
d'après le Lemme \ref{sites-lemma-colimit}. Comme $u_j(u_a(U_i)) = U$,
on a $f_{j, *}\mathcal{F}(u_a(U_i)) = \mathcal{F}(U)$
pour tout $a : j \to i$, par définition. La valeur de la limite inductive
est donc $\mathcal{F}(U)$, ce qui achève la démonstration.
\end{proof}











\section{Fonctorialité supplémentaire des préfaisceaux}
\label{sites-section-more-functoriality-PSh}

\noindent
Dans cette section, nous revenons sur le contenu de la
section \ref{sites-section-functoriality-PSh}.
Soit $u : \mathcal{C} \to \mathcal{D}$ un foncteur entre catégories.
Rappelons que
$$
u^p :
\textit{PSh}(\mathcal{D})
\longrightarrow
\textit{PSh}(\mathcal{C})
$$
est le foncteur qui associe à $\mathcal{G}$ sur $\mathcal{D}$ le préfaisceau
$u^p\mathcal{G} = \mathcal{G} \circ u$. Ce foncteur possède non seulement
un adjoint à gauche, à savoir $u_p$, mais aussi un adjoint à droite.

\medskip\noindent
Plus précisément, pour tout $V \in \Ob(\mathcal{D})$,
nous définissons une catégorie ${}_V\mathcal{I} = {}_V^u\mathcal{I}$.
Ses objets sont les couples $(U, \psi : u(U) \to V)$.
Remarquons que la flèche est orientée dans le sens opposé à celle utilisée
pour définir la catégorie $\mathcal{I}_V^u$ dans la
section \ref{sites-section-functoriality-PSh}.
Un morphisme $(U, \psi) \to (U', \psi')$ est donné par un morphisme
$\alpha : U \to U'$ tel que $\psi = \psi' \circ u(\alpha)$.
De plus, pour tout préfaisceau d'ensembles $\mathcal{F}$ sur $\mathcal{C}$,
nous introduisons le foncteur
${}_V\mathcal{F} : {}_V\mathcal{I}^{opp} \to \textit{Ensembles}$,
défini par la règle ${}_V\mathcal{F}(U, \psi) = \mathcal{F}(U)$.
Nous définissons
$$
{}_pu(\mathcal{F})(V) := \lim_{{}_V\mathcal{I}^{opp}} {}_V\mathcal{F}
$$
En tant que limite, cet objet possède des morphismes de projection
$c(\psi) : {}_pu(\mathcal{F})(V) \to \mathcal{F}(U)$
pour tout objet $(U, \psi)$ de ${}_V\mathcal{I}$.
En fait,
$$
{}_pu(\mathcal{F})(V)
=
\left\{
\begin{matrix}
\text{familles }
s_{(U, \psi)} \in \mathcal{F}(U) \\
\forall \beta : (U_1, \psi_1) \to (U_2, \psi_2)
\text{ dans }{}_V\mathcal{I} \\
\text{ telles que } \beta^*s_{(U_2, \psi_2)} = s_{(U_1, \psi_1)}
\end{matrix}
\right\}
$$
où la correspondance est donnée par
$s \mapsto s_{(U, \psi)} = c(\psi)(s)$.
Nous laissons au lecteur le soin de définir les applications de restriction
${}_pu(\mathcal{F})(V) \to {}_pu(\mathcal{F})(V')$
associées à tout morphisme $V' \to V$ de $\mathcal{D}$.
Le préfaisceau ainsi obtenu sera noté ${}_pu\mathcal{F}$.

\begin{lemma}
\label{sites-lemma-recover-pu}
Il existe une application canonique
${}_pu\mathcal{F}(u(U)) \to \mathcal{F}(U)$,
compatible aux applications de restriction.
\end{lemma}

\begin{proof}
C'est simplement le morphisme de projection $c(\text{id}_{u(U)})$ ci-dessus.
\end{proof}

\noindent
Remarquons que tout morphisme de préfaisceaux
$\mathcal{F} \to \mathcal{F}'$ donne naissance à des systèmes compatibles
de morphismes entre les foncteurs
${}_V\mathcal{F} \to {}_V\mathcal{F}'$, et donc à un morphisme de préfaisceaux
${}_pu\mathcal{F} \to {}_pu\mathcal{F}'$.
Autrement dit, nous avons défini un foncteur
$$
{}_pu :
\textit{PSh}(\mathcal{C})
\longrightarrow
\textit{PSh}(\mathcal{D})
$$

\begin{lemma}
\label{sites-lemma-adjoints-pu}
Le foncteur ${}_pu$ est adjoint à droite du foncteur $u^p$.
Autrement dit, la formule
$$
\Mor_{\textit{PSh}(\mathcal{C})}(u^p\mathcal{G}, \mathcal{F})
=
\Mor_{\textit{PSh}(\mathcal{D})}(\mathcal{G}, {}_pu\mathcal{F})
$$
est bifonctorielle en $\mathcal{F}$ et $\mathcal{G}$.
\end{lemma}

\begin{proof}
La démonstration est exactement la même que celle du
Lemme \ref{sites-lemma-adjoints-u}. Remarquons que le morphisme
$u^p{}_pu \mathcal{F} \to \mathcal{F}$ du
Lemme \ref{sites-lemma-recover-pu} est celui qui permet de passer
du membre de droite au membre de gauche.

\medskip\noindent
On peut aussi considérer un préfaisceau d'ensembles $\mathcal{F}$ sur
$\mathcal{C}$ comme un préfaisceau $\mathcal{F}'$ sur
$\mathcal{C}^{opp}$ à valeurs dans $\textit{Ensembles}^{opp}$,
et procéder de même sur $\mathcal{D}$.
Vérifions que $({}_pu \mathcal{F})' = u_p(\mathcal{F}')$ et que
$(u^p\mathcal{G})' = u^p(\mathcal{G}')$.
D'après la Remarque \ref{sites-remark-functoriality-presheaves-values},
$u_p$ et $u^p$ sont adjoints pour les préfaisceaux à valeurs dans
$\textit{Ensembles}^{opp}$. Le résultat s'en déduit formellement.
\end{proof}

\noindent
Ainsi, un foncteur $u : \mathcal{C} \to \mathcal{D}$ de catégories
fournit une suite de foncteurs
$$
u_p, u^p, {}_pu
$$
entre les catégories de préfaisceaux, dans laquelle, pour chaque couple
consécutif, le premier foncteur est adjoint à gauche du second.

\begin{lemma}
\label{sites-lemma-adjoint-functors}
Soient $u : \mathcal{C} \to \mathcal{D}$ et
$v : \mathcal{D} \to \mathcal{C}$ des foncteurs de catégories.
Supposons que $v$ soit adjoint à droite de $u$. Alors :
\begin{enumerate}
\item $u^ph_V = h_{v(V)}$ pour tout $V$ dans $\mathcal{D}$ ;
\item la catégorie $\mathcal{I}^v_U$ possède un objet initial ;
\item la catégorie ${}_V^u\mathcal{I}$ possède un objet final ;
\item ${}_pu = v^p$ ;
\item $u^p = v_p$.
\end{enumerate}
\end{lemma}

\begin{proof}
Démonstration de (1). Soit $V$ un objet de $\mathcal{D}$. On a
$u^ph_V = h_{v(V)}$ parce que
$u^ph_V(U) = \Mor_\mathcal{D}(u(U), V) = \Mor_\mathcal{C}(U, v(V))$
par hypothèse.

\medskip\noindent
Démonstration de (2). Soit $U$ un objet de $\mathcal{C}$. Soit
$\eta : U \to v(u(U))$ le morphisme adjoint du morphisme
$\text{id} : u(U) \to u(U)$. Nous affirmons que $(u(U), \eta)$
est un objet initial de $\mathcal{I}_U^v$. En effet, étant donné
un objet $(V, \phi : U \to v(V))$ de $\mathcal{I}_U^v$,
le morphisme $\phi$ est adjoint à un morphisme $\psi : u(U) \to V$,
qui définit alors un morphisme $(u(U), \eta) \to (V, \phi)$.

\medskip\noindent
Démonstration de (3). Soit $V$ un objet de $\mathcal{D}$. Soit
$\xi : u(v(V)) \to V$ le morphisme adjoint du morphisme
$\text{id} : v(V) \to v(V)$. Nous affirmons que $(v(V), \xi)$
est un objet final de ${}_V^u\mathcal{I}$. En effet, étant donné
un objet $(U, \psi : u(U) \to V)$ de ${}_V^u\mathcal{I}$,
le morphisme $\psi$ est adjoint à un morphisme $\phi : U \to v(V)$,
qui définit alors un morphisme $(U, \psi) \to (v(V), \xi)$.

\medskip\noindent
Par conséquent, pour tout préfaisceau $\mathcal{F}$ sur $\mathcal{C}$,
on a
\begin{eqnarray*}
v^p\mathcal{F}(V)
& = &
\mathcal{F}(v(V)) \\
& = &
\Mor_{\textit{PSh}(\mathcal{C})}(h_{v(V)}, \mathcal{F}) \\
& = &
\Mor_{\textit{PSh}(\mathcal{C})}(u^ph_V, \mathcal{F}) \\
& = &
\Mor_{\textit{PSh}(\mathcal{D})}(h_V, {}_pu\mathcal{F}) \\
& = &
{}_pu\mathcal{F}(V)
\end{eqnarray*}
ce qui démontre la partie (4). La partie (5) résulte de l'unicité
des foncteurs adjoints.
\end{proof}





\section{Foncteurs cocontinus}
\label{sites-section-cocontinuous-functors}

\noindent
Il existe une autre manière de construire des morphismes de topos.
Elle fait intervenir des foncteurs cocontinus entre sites, définis comme suit.

\begin{definition}
\label{sites-definition-cocontinuous}
Soient $\mathcal{C}$ et $\mathcal{D}$ des sites.
Soit $u : \mathcal{C} \to \mathcal{D}$ un foncteur.
Le foncteur $u$ est dit {\it cocontinu} si, pour tout
$U \in \Ob(\mathcal{C})$ et tout recouvrement
$\{V_j \to u(U)\}_{j \in J}$ de $\mathcal{D}$, il existe un recouvrement
$\{U_i \to U\}_{i\in I}$ de $\mathcal{C}$ tel que la famille de morphismes
$\{u(U_i) \to u(U)\}_{i \in I}$ raffine le recouvrement
$\{V_j \to u(U)\}_{j \in J}$.
\end{definition}

\noindent
Remarquons que $\{u(U_i) \to u(U)\}_{i \in I}$ n'est en général {\it pas}
un recouvrement du site $\mathcal{D}$.

\begin{lemma}
\label{sites-lemma-pu-sheaf}
Soient $\mathcal{C}$ et $\mathcal{D}$ des sites.
Soit $u : \mathcal{C} \to \mathcal{D}$ cocontinu.
Soit $\mathcal{F}$ un faisceau sur $\mathcal{C}$.
Alors ${}_pu\mathcal{F}$ est un faisceau sur $\mathcal{D}$,
que nous noterons ${}_su\mathcal{F}$.
\end{lemma}

\begin{proof}
Soit $\{V_j \to V\}_{j \in J}$ un recouvrement du site $\mathcal{D}$.
Nous devons montrer que
$$
\xymatrix{
&\ \phantom{}_pu\mathcal{F}(V) \ar[r] &
\prod {}_pu\mathcal{F}(V_j) \ar@<1ex>[r] \ar@<-1ex>[r] &
\prod {}_pu\mathcal{F}(V_j \times_V V_{j'}) &
}
$$
est un diagramme égalisateur. Comme ${}_pu$ est adjoint à droite de $u^p$,
nous avons
$$
{}_pu\mathcal{F}(V) =
\Mor_{\textit{PSh}(\mathcal{D})}(h_V, {}_pu\mathcal{F}) =
\Mor_{\textit{PSh}(\mathcal{C})}(u^ph_V, \mathcal{F}) =
\Mor_{\Sh(\mathcal{C})}((u^ph_V)^\#, \mathcal{F})
$$
Il suffit donc de montrer que
\begin{equation}
\label{sites-equation-coequalizer}
\xymatrix{
\coprod u^p h_{V_j \times_V V_{j'}} \ar@<1ex>[r] \ar@<-1ex>[r] &
\coprod u^p h_{V_j} \ar[r] &
u^p h_V
}
\end{equation}
devient un diagramme coégalisateur après passage aux faisceaux associés.
(Rappelons qu'un coproduit dans la catégorie des faisceaux est le faisceau
associé au coproduit dans la catégorie des préfaisceaux ; voir le
Lemme \ref{sites-lemma-colimit-sheaves}.)

\medskip\noindent
Montrons d'abord que la deuxième flèche de
(\ref{sites-equation-coequalizer}) devient surjective après passage aux faisceaux
associés. Nous utilisons pour cela le Lemme \ref{sites-lemma-mono-epi-sheaves}.
Il suffit donc de montrer qu'une section $s$ de $u^ph_V$ sur $U$ se relève
en une section de $\coprod u^p h_{V_j}$ sur les membres d'un recouvrement de $U$.
Remarquons que $s$ est un morphisme $s : u(U) \to V$. Alors
$\{V_j \times_{V, s} u(U) \to u(U)\}$ est un recouvrement de $\mathcal{D}$.
Comme $u$ est cocontinu, il existe donc un recouvrement $\{U_i \to U\}$
tel que $\{u(U_i) \to u(U)\}$ raffine
$\{V_j \times_{V, s} u(U) \to u(U)\}$.
Cela signifie que chaque restriction $s|_{U_i} : u(U_i) \to V$ se factorise
par un morphisme $s_i : u(U_i) \to V_j$ pour un certain $j$ ; autrement dit,
$s|_{U_i}$ appartient à l'image de
$u^ph_{V_j}(U_i) \to u^ph_V(U_i)$, comme voulu.

\medskip\noindent
Soient $s, s' \in (\coprod u^ph_{V_j})^\#(U)$ ayant la même image
dans $(u^ph_V)^\#(U)$. Pour achever la démonstration du lemme, montrons
qu'après avoir remplacé $U$ par les membres d'un recouvrement, $s, s'$ sont
l'image de la même section de $\coprod u^p h_{V_j \times_V V_{j'}}$
par les deux morphismes de (\ref{sites-equation-coequalizer}). Nous pouvons d'abord
remplacer $U$ par les membres d'un recouvrement et supposer
$s \in u^ph_{V_j}(U)$ et $s' \in u^ph_{V_{j'}}(U)$.
Un second remplacement de ce type assure que $s$ et $s'$ ont la même image
dans $u^ph_V(U)$, et non seulement dans le faisceau associé.
Ainsi, $s : u(U) \to V_j$ et $s' : u(U) \to V_{j'}$ sont des morphismes
de $\mathcal{D}$ tels que
$$
\xymatrix{
u(U) \ar[r]_{s'} \ar[d]_s & V_{j'} \ar[d] \\
V_j \ar[r] & V
}
$$
soit commutatif. Nous obtenons donc
$t = (s, s') : u(U) \to V_j \times_V V_{j'}$,
c'est-à-dire une section $t \in u^ph_{V_j \times_V V_{j'}}(U)$
qui s'envoie sur $s, s'$, comme voulu.
\end{proof}

\begin{lemma}
\label{sites-lemma-exact-cocontinuous}
Soient $\mathcal{C}$ et $\mathcal{D}$ des sites.
Soit $u : \mathcal{C} \to \mathcal{D}$ cocontinu.
Le foncteur
$\Sh(\mathcal{D}) \to \Sh(\mathcal{C})$,
$\mathcal{G} \mapsto (u^p\mathcal{G})^\#$
est adjoint à gauche du foncteur ${}_su$ introduit au
Lemme \ref{sites-lemma-pu-sheaf}. De plus, il est exact.
\end{lemma}

\begin{proof}
Démontrons la propriété d'adjonction par les égalités suivantes :
\begin{eqnarray*}
\Mor_{\Sh(\mathcal{C})}
((u^p\mathcal{G})^\#, \mathcal{F})
& = &
\Mor_{\textit{PSh}(\mathcal{C})}
(u^p\mathcal{G}, \mathcal{F}) \\
& = &
\Mor_{\textit{PSh}(\mathcal{D})}
(\mathcal{G}, {}_pu\mathcal{F}) \\
& = &
\Mor_{\Sh(\mathcal{D})}
(\mathcal{G}, {}_su\mathcal{F}).
\end{eqnarray*}
Il est donc adjoint à gauche, et par conséquent exact à droite ; voir
Catégories, Lemme \ref{categories-lemma-exact-adjoint}.
Nous avons vu que le passage au faisceau associé est exact à gauche ; voir le
Lemme \ref{sites-lemma-sheafification-exact}.
En outre, l'inclusion
$i : \Sh(\mathcal{D}) \to \textit{PSh}(\mathcal{D})$
est exacte à gauche d'après le Lemme \ref{sites-lemma-limit-sheaf}.
Enfin, le foncteur $u^p$ est exact à gauche parce qu'il est adjoint à droite,
à savoir de $u_p$. Le foncteur considéré est donc le composé
${}^\# \circ u^p \circ i$ de foncteurs exacts à gauche ;
il est par conséquent exact à gauche.
\end{proof}

\noindent
Terminons cette section par un lemme technique.

\begin{lemma}
\label{sites-lemma-technical-pu}
Dans la situation du Lemme \ref{sites-lemma-exact-cocontinuous},
pour tout préfaisceau $\mathcal{G}$ sur $\mathcal{D}$, on a
$(u^p\mathcal{G})^\# = (u^p(\mathcal{G}^\#))^\#$.
\end{lemma}

\begin{proof}
Pour tout faisceau $\mathcal{F}$ sur $\mathcal{C}$, on a
\begin{eqnarray*}
\Mor_{\Sh(\mathcal{C})}((u^p(\mathcal{G}^\#))^\#, \mathcal{F})
& = &
\Mor_{\Sh(\mathcal{D})}(\mathcal{G}^\#, {}_su\mathcal{F}) \\
& = &
\Mor_{\Sh(\mathcal{D})}(\mathcal{G}^\#, {}_pu\mathcal{F}) \\
& = &
\Mor_{\textit{PSh}(\mathcal{D})}(\mathcal{G}, {}_pu\mathcal{F}) \\
& = &
\Mor_{\textit{PSh}(\mathcal{C})}(u^p\mathcal{G}, \mathcal{F}) \\
& = &
\Mor_{\Sh(\mathcal{C})}((u^p\mathcal{G})^\#, \mathcal{F})
\end{eqnarray*}
et le résultat résulte du lemme de Yoneda.
\end{proof}

\begin{remark}
\label{sites-remark-cartesian-cocontinuous}
Soit $u : \mathcal{C} \to \mathcal{D}$ un foncteur entre catégories.
Étant donnés des morphismes $g : u(U) \to V$ et $f : W \to V$
de $\mathcal{D}$, nous pouvons considérer le foncteur
$$
\mathcal{C}^{opp} \longrightarrow \textit{Ensembles},\quad
T \longmapsto
\Mor_\mathcal{C}(T, U)
\times_{\Mor_\mathcal{D}(u(T), V)}
\Mor_\mathcal{D}(u(T), W)
$$
Si ce foncteur est représentable, notons
$U \times_{g, V, f} W$ l'objet correspondant de $\mathcal{C}$.
Supposons que $\mathcal{C}$ et $\mathcal{D}$ soient des sites.
Considérons la propriété $P$ suivante : pour tout recouvrement
$\{f_j : V_j \to V\}$ de $\mathcal{D}$ et tout morphisme
$g : u(U) \to V$, on a :
\begin{enumerate}
\item $U \times_{g, V, f_i} V_i$ existe pour tout $i$ ;
\item $\{U \times_{g, V, f_i} V_i \to U\}$ est un recouvrement de $\mathcal{C}$.
\end{enumerate}
Remarquons la similitude avec la définition des foncteurs continus.
Si $u$ possède la propriété $P$, alors $u$ est cocontinu (détails omis).
De nombreux foncteurs cocontinus que nous rencontrerons possèdent $P$.
\end{remark}

\section{Foncteurs cocontinus et morphismes de topos}
\label{sites-section-cocontinuous-morphism-topoi}

\noindent
Il ressort clairement de ce qui précède qu'un foncteur cocontinu $u$
définit un morphisme de topos dans le même sens que $u$.
Ce sens est donc opposé à celui du morphisme de topos associé
(sous certaines conditions) à un foncteur continu $u$, comme dans la
Définition \ref{sites-definition-morphism-sites}, la
Proposition \ref{sites-proposition-get-morphism} et le
Lemme \ref{sites-lemma-morphism-sites-topoi}.

\begin{lemma}
\label{sites-lemma-cocontinuous-morphism-topoi}
Soient $\mathcal{C}$ et $\mathcal{D}$ des sites.
Soit $u : \mathcal{C} \to \mathcal{D}$ cocontinu.
Les foncteurs $g_* = {}_su$ et $g^{-1} = (u^p\ )^\#$
définissent un morphisme de topos $g$ de $\Sh(\mathcal{C})$ vers
$\Sh(\mathcal{D})$.
\end{lemma}

\begin{proof}
C'est exactement le contenu du Lemme \ref{sites-lemma-exact-cocontinuous}.
\end{proof}

\begin{lemma}
\label{sites-lemma-composition-cocontinuous}
\begin{slogan}
La composition des foncteurs de sites préserve leur cocontinuité.
\end{slogan}
Soient $u : \mathcal{C} \to \mathcal{D}$ et
$v : \mathcal{D} \to \mathcal{E}$ des foncteurs cocontinus. Alors
$v \circ u$ est cocontinu et l'on a $h = g \circ f$, où
$f : \Sh(\mathcal{C}) \to \Sh(\mathcal{D})$, respectivement
$g : \Sh(\mathcal{D}) \to \Sh(\mathcal{E})$, respectivement
$h : \Sh(\mathcal{C}) \to \Sh(\mathcal{E})$, est le morphisme de topos
associé à $u$, respectivement à $v$, respectivement à $v \circ u$.
\end{lemma}

\begin{proof}
Soit $U \in \Ob(\mathcal{C})$.
Soit $\{E_i \to v(u(U))\}$ un recouvrement de $U$ dans $\mathcal{E}$.
Par hypothèse, il existe un recouvrement $\{D_j \to u(U)\}$ dans
$\mathcal{D}$ tel que $\{v(D_j) \to v(u(U))\}$ raffine
$\{E_i \to v(u(U))\}$. Toujours par hypothèse, il existe un recouvrement
$\{C_l \to U\}$ dans $\mathcal{C}$ tel que
$\{u(C_l) \to u(U)\}$ raffine $\{D_j \to u(U)\}$. Alors
$\{v(u(C_l)) \to v(u(U))\}$ raffine bien le recouvrement
$\{E_i \to v(u(U))\}$. Cela montre que $v \circ u$ est cocontinu.
Pour démontrer la dernière assertion, il suffit de montrer que
${}_sv \circ {}_su = {}_s(v \circ u)$. Il suffit de démontrer que
${}_pv \circ {}_pu = {}_p(v \circ u)$ ; voir le
Lemme \ref{sites-lemma-pu-sheaf}. Comme ${}_pu$, respectivement ${}_pv$,
respectivement ${}_p(v \circ u)$, est adjoint à droite de $u^p$,
respectivement de $v^p$, respectivement de $(v \circ u)^p$, il suffit de
montrer que $u^p \circ v^p = (v \circ u)^p$. Cela résulte directement des
définitions.
\end{proof}

\begin{example}
\label{sites-example-open-immersion-cocontinuous}
Soit $X$ un espace topologique.
Soit $j : U  \to X$ l'inclusion d'un sous-espace ouvert.
Rappelons que nous disposons des sites $X_{Zar}$ et $U_{Zar}$ ; voir
l'Exemple \ref{sites-example-site-topological}.
Rappelons également que le foncteur $u : X_{Zar} \to U_{Zar}$ associé à $j$
est continu et donne lieu à un morphisme de sites
$U_{Zar} \to X_{Zar}$ ; voir l'Exemple \ref{sites-example-continuous-map}.
Cela donne aussi un morphisme de topos $(j_*, j^{-1})$.
Considérons ensuite le foncteur
$v : U_{Zar} \to X_{Zar}$, $V \mapsto v(V) = V$
(le même ouvert, mais considéré maintenant comme un objet de $X_{Zar}$).
Ce foncteur est cocontinu. En effet, si
$v(V) = \bigcup_{j \in J} W_j$ est un recouvrement ouvert dans $X$, alors
chaque $W_j$ est nécessairement contenu dans $U$, donc est de la forme
$v(V_j)$, et $V = \bigcup_{j \in J} V_j$ est trivialement un recouvrement
ouvert dans $U$. Le Lemme \ref{sites-lemma-cocontinuous-morphism-topoi} ci-dessus
montre donc qu'il existe un morphisme de topos associé à $v$
$$
\Sh(U) \longrightarrow \Sh(X)
$$
donné par ${}_sv$ et $(v^p\ )^\#$. Nous affirmons qu'en fait
$(v^p\ )^\# = j^{-1}$ et que ${}_sv = j_*$ ; autrement dit, c'est le même
morphisme de topos que celui donné ci-dessus. La manière la plus simple de le
voir est peut-être de remarquer que, pour tout faisceau $\mathcal{G}$ sur $X$,
on a $v^p\mathcal{G}(V) = \mathcal{G}(V)$, ce qui, d'après Faisceaux,
Lemme \ref{sheaves-lemma-j-pullback}, décrit $j^{-1}\mathcal{G}$ (le passage au
faisceau associé est donc superflu dans ce cas). L'égalité de ${}_sv$ et de
$j_*$ résulte de l'unicité des foncteurs adjoints (mais on peut également la
calculer directement).
\end{example}

\begin{example}
\label{sites-example-open-map-cocontinuous}
Cet exemple généralise légèrement
l'Exemple \ref{sites-example-open-immersion-cocontinuous}.
Soit $f : X \to Y$ une application continue d'espaces topologiques.
Supposons que $f$ soit ouverte.
Rappelons que nous disposons des sites $X_{Zar}$ et $Y_{Zar}$ ; voir
l'Exemple \ref{sites-example-site-topological}.
Rappelons également que le foncteur $u : Y_{Zar} \to X_{Zar}$ associé à $f$
est continu et donne lieu à un morphisme de sites
$X_{Zar} \to Y_{Zar}$ ; voir l'Exemple \ref{sites-example-continuous-map}.
Cela donne aussi un morphisme de topos $(f_*, f^{-1})$.
Considérons ensuite le foncteur
$v : X_{Zar} \to Y_{Zar}$, $U \mapsto v(U) = f(U)$.
Ce foncteur est cocontinu. En effet, si
$f(U) = \bigcup_{j \in J} V_j$ est un recouvrement ouvert dans $Y$, alors, en
posant $U_j = f^{-1}(V_j) \cap U$, on obtient un recouvrement ouvert
$U = \bigcup U_j$ tel que $f(U) = \bigcup f(U_j)$ raffine
$f(U) = \bigcup V_j$.
Le Lemme \ref{sites-lemma-cocontinuous-morphism-topoi} ci-dessus montre donc qu'il
existe un morphisme de topos associé à $v$
$$
\Sh(X) \longrightarrow \Sh(Y)
$$
donné par ${}_sv$ et $(v^p\ )^\#$. Nous affirmons qu'en fait
$(v^p\ )^\# = f^{-1}$ et que ${}_sv = f_*$ ; autrement dit, c'est le même
morphisme de topos que celui donné ci-dessus. Pour tout faisceau
$\mathcal{G}$ sur $Y$, on a
$v^p\mathcal{G}(U) = \mathcal{G}(f(U))$.
D'autre part, on peut calculer
$u_p\mathcal{G}(U) = \colim_{f(U) \subset V} \mathcal{G}(V)
= \mathcal{G}(f(U))$, car $(f(U), U \subset f^{-1}(f(U)))$ est manifestement
un objet initial de la catégorie $\mathcal{I}_U^u$ de la
section \ref{sites-section-functoriality-PSh}.
Ainsi $u_p = v^p$ et l'on en conclut que
$f^{-1} = u_s = (v^p\ )^\#$.
L'égalité de ${}_sv$ et de $f_*$ résulte de l'unicité des foncteurs adjoints
(mais on peut également la calculer directement).
\end{example}

\noindent
Dans le premier Exemple \ref{sites-example-open-immersion-cocontinuous}, le
foncteur $v$ est également continu. Mais, dans le second
Exemple \ref{sites-example-open-map-cocontinuous}, il n'est généralement pas
continu, car la condition (2) de la Définition \ref{sites-definition-continuous}
peut ne pas être satisfaite. Le lemme suivant s'applique donc au premier
exemple, mais non au second.

\begin{lemma}
\label{sites-lemma-when-shriek}
\begin{slogan}
Le passage au faisceau associé est superflu dans les morphismes de topos
associés à des foncteurs de sites à la fois continus et cocontinus.
\end{slogan}
Soient $\mathcal{C}$ et $\mathcal{D}$ des sites.
Soit $u : \mathcal{C} \to \mathcal{D}$ un foncteur.
Supposons que
\begin{enumerate}
\item[(a)] $u$ soit cocontinu, et
\item[(b)] $u$ soit continu.
\end{enumerate}
Soit $g : \Sh(\mathcal{C}) \to \Sh(\mathcal{D})$
le morphisme de topos associé. Alors
\begin{enumerate}
\item le passage au faisceau associé dans la formule
$g^{-1} = (u^p\ )^\#$ est inutile ; autrement dit,
$g^{-1}(\mathcal{G})(U) = \mathcal{G}(u(U))$,
\item $g^{-1}$ possède un adjoint à gauche $g_{!} = (u_p\ )^\#$, et
\item $g^{-1}$ commute à toutes les limites projectives et inductives.
\end{enumerate}
\end{lemma}

\begin{proof}
D'après le Lemme \ref{sites-lemma-pushforward-sheaf}, pour tout faisceau
$\mathcal{G}$ sur $\mathcal{D}$, le préfaisceau $u^p\mathcal{G}$ est un
faisceau sur $\mathcal{C}$ (et est noté $u^s\mathcal{G}$). Le foncteur
$\mathcal{F} \mapsto (u_p\mathcal{F})^\#$ est adjoint à gauche de $u^s$
d'après le Lemme \ref{sites-lemma-adjoint-sheaves}. L'assertion relative aux limites
projectives et inductives résulte de la discussion de Catégories,
section \ref{categories-section-adjoint}.
\end{proof}

\noindent
Dans la situation du Lemme \ref{sites-lemma-when-shriek} ci-dessus, on voit que
l'on dispose d'une suite de foncteurs adjoints
$$
g_{!}, \ g^{-1}, \ g_*.
$$
Le foncteur $g_!$ n'est {\it pas} exact en général, car il ne transforme pas,
en général, un objet final de $\Sh(\mathcal{C})$ en un objet final de
$\Sh(\mathcal{D})$.
Voir Faisceaux, Remarque \ref{sheaves-remark-j-shriek-not-exact}.
En revanche, dans le cadre topologique de
l'Exemple \ref{sites-example-open-immersion-cocontinuous}, le foncteur $j_!$ est
exact sur les faisceaux abéliens ; voir Modules,
Lemme \ref{modules-lemma-j-shriek-exact}. Le lemme suivant donne la
généralisation au cas des sites.

\begin{lemma}
\label{sites-lemma-preserve-equalizers}
Soient $\mathcal{C}$ et $\mathcal{D}$ des sites.
Soit $u : \mathcal{C} \to \mathcal{D}$ un foncteur.
Supposons que
\begin{enumerate}
\item[(a)] $u$ soit cocontinu,
\item[(b)] $u$ soit continu, et
\item[(c)] les produits fibrés et les égalisateurs existent dans
$\mathcal{C}$ et que $u$ commute à ceux-ci.
\end{enumerate}
Dans ce cas, le foncteur $g_!$ ci-dessus commute aux produits fibrés et aux
égalisateurs (et, plus généralement, aux limites projectives finies connexes).
\end{lemma}

\begin{proof}
Supposons (a), (b) et (c).
On a $g_! = (u_p\ )^\#$. Rappelons (Lemme \ref{sites-lemma-limit-sheaf}) que les
limites projectives de faisceaux sont égales aux limites projectives
correspondantes calculées comme préfaisceaux. De plus, le passage au faisceau
associé commute aux limites projectives finies
(Lemme \ref{sites-lemma-sheafification-exact}). Il suffit donc de montrer que $u_p$
commute aux produits fibrés et aux égalisateurs. Pour cela, il suffit que les
limites inductives indexées par les catégories
$(\mathcal{I}_V^u)^{opp}$ de la
section \ref{sites-section-functoriality-PSh}
commutent aux produits fibrés et aux égalisateurs. Cela résulte du
Lemme \ref{sites-lemma-almost-directed} et de Catégories,
Lemme \ref{categories-lemma-almost-directed-commutes-equalizers}.
\end{proof}

\noindent
Le lemme suivant traite d'un cas qui ressemble encore davantage au morphisme
associé à une immersion ouverte d'espaces topologiques.

\begin{lemma}
\label{sites-lemma-back-and-forth}
Soient $\mathcal{C}$ et $\mathcal{D}$ des sites.
Soit $u : \mathcal{C} \to \mathcal{D}$ un foncteur.
Supposons que
\begin{enumerate}
\item[(a)] $u$ soit cocontinu,
\item[(b)] $u$ soit continu, et
\item[(c)] $u$ soit pleinement fidèle.
\end{enumerate}
Pour $g_!, g^{-1}, g_*$ comme ci-dessus, les morphismes canoniques
$\mathcal{F} \to g^{-1}g_!\mathcal{F}$ et
$g^{-1}g_*\mathcal{F} \to \mathcal{F}$ sont des isomorphismes pour tout
faisceau $\mathcal{F}$ sur $\mathcal{C}$.
\end{lemma}

\begin{proof}
Soit $X$ un objet de $\mathcal{C}$.
Nous avons vu dans les Lemmes \ref{sites-lemma-pu-sheaf} et
\ref{sites-lemma-when-shriek} que le passage au faisceau associé n'est pas
nécessaire pour les foncteurs
$g^{-1} = (u^p\ )^\#$ et $g_{*} = ({}_pu\ )^\#$.
On peut calculer
$(g^{-1}g_{*}\mathcal{F})(X) = g_{*}\mathcal{F}(u(X))
= \lim \mathcal{F}(Y)$. Ici, la limite est prise sur la catégorie des couples
$(Y, u(Y) \to u(X))$, où l'on n'exige pas que les morphismes
$u(Y) \to u(X)$ soient de la forme $u(\alpha)$, avec $\alpha$ un morphisme de
$\mathcal{C}$. D'après l'hypothèse (c), ils proviennent automatiquement de
morphismes de $\mathcal{C}$, et l'on en déduit que la limite est la valeur en
$(X, u(\text{id}_X))$, c'est-à-dire $\mathcal{F}(X)$.
Cela montre que $g^{-1}g_{*}\mathcal{F} = \mathcal{F}$.

\medskip\noindent
D'autre part, $(g^{-1}g_{!}\mathcal{F})(X) =
g_{!}\mathcal{F}(u(X)) = (u_p\mathcal{F})^\#(u(X))$, et
$u_p\mathcal{F}(u(X)) = \colim \mathcal{F}(Y)$.
Ici, la limite inductive est prise sur la catégorie des couples
$(Y, u(X) \to u(Y))$, où l'on n'exige pas que les morphismes
$u(X) \to u(Y)$ soient de la forme $u(\alpha)$, avec $\alpha$ un morphisme de
$\mathcal{C}$. D'après l'hypothèse (c), ils proviennent automatiquement de
morphismes de $\mathcal{C}$, et l'on en déduit que la limite inductive est la
valeur en $(X, u(\text{id}_X))$, c'est-à-dire $\mathcal{F}(X)$. Ainsi, pour
tout $X \in \Ob(\mathcal{C})$, on a
$u_p\mathcal{F}(u(X)) = \mathcal{F}(X)$.
Comme $u$ est cocontinu et continu, tout recouvrement de $u(X)$ dans
$\mathcal{D}$ peut être raffiné par un recouvrement (!)
$\{u(X_i) \to u(X)\}$ de $\mathcal{D}$, où $\{X_i \to X\}$ est un
recouvrement dans $\mathcal{C}$. Cela implique également
$(u_p\mathcal{F})^+(u(X)) = \mathcal{F}(X)$, car, dans la limite inductive qui
définit la valeur de $(u_p\mathcal{F})^+$ en $u(X)$, on peut se restreindre au
système cofinal des recouvrements $\{u(X_i) \to u(X)\}$ ci-dessus. On voit
donc que $(u_p\mathcal{F})^+(u(X)) = \mathcal{F}(X)$ pour tout objet $X$ de
$\mathcal{C}$. En répétant encore une fois cet argument, on obtient l'égalité
$(u_p\mathcal{F})^\#(u(X)) = \mathcal{F}(X)$ pour tout objet $X$ de
$\mathcal{C}$. Cela donne l'égalité voulue
$g^{-1}g_!\mathcal{F} = \mathcal{F}$.
\end{proof}

\noindent
Enfin, voici un cas qui ne possède aucun exemple topologique correspondant.
Nous utiliserons ce lemme pour voir ce qui se produit lorsque l'on agrandit
un « univers partiel » de schémas tout en conservant la même topologie. Dans
la situation du lemme, le morphisme de topos
$g : \Sh(\mathcal{C}) \to \Sh(\mathcal{D})$ identifie
$\Sh(\mathcal{C})$ à un sous-topos de $\Sh(\mathcal{D})$
(section \ref{sites-section-subtopoi}) et, de plus, le plongement donné admet une
rétraction.

\begin{lemma}
\label{sites-lemma-bigger-site}
Soient $\mathcal{C}$ et $\mathcal{D}$ des sites.
Soit $u : \mathcal{C} \to \mathcal{D}$ un foncteur.
Supposons que
\begin{enumerate}
\item[(a)] $u$ soit cocontinu,
\item[(b)] $u$ soit continu,
\item[(c)] $u$ soit pleinement fidèle,
\item[(d)] les produits fibrés existent dans $\mathcal{C}$ et que $u$ commute
à ceux-ci, et
\item[(e)] il existe des objets finaux
$e_\mathcal{C} \in \Ob(\mathcal{C})$,
$e_\mathcal{D} \in \Ob(\mathcal{D})$ tels que
$u(e_\mathcal{C}) = e_\mathcal{D}$.
\end{enumerate}
Soient $g_!, g^{-1}, g_*$ comme ci-dessus. Alors $u$ définit un morphisme de
sites $f : \mathcal{D} \to \mathcal{C}$ avec $f_* = g^{-1}$ et
$f^{-1} = g_!$. Le composé
$$
\xymatrix{
\Sh(\mathcal{C}) \ar[r]^g &
\Sh(\mathcal{D}) \ar[r]^f &
\Sh(\mathcal{C})
}
$$
est isomorphe au morphisme identité du topos $\Sh(\mathcal{C})$. De plus, le
foncteur $f^{-1}$ est pleinement fidèle.
\end{lemma}

\begin{proof}
Par hypothèse, le foncteur $u$ satisfait aux hypothèses de la
Proposition \ref{sites-proposition-get-morphism}. Ainsi, $u$ définit un morphisme de
sites, donc un morphisme de topos $f$ comme dans le
Lemme \ref{sites-lemma-morphism-sites-topoi}. Les formules $f_* = g^{-1}$ et
$f^{-1} = g_!$ résultent clairement du lemme cité et du
Lemme \ref{sites-lemma-when-shriek}.
On a
$f_* \circ g_* = g^{-1} \circ g_* \cong \text{id}$, et
$g^{-1} \circ f^{-1} = g^{-1} \circ g_! \cong \text{id}$
d'après le Lemme \ref{sites-lemma-back-and-forth}.

\medskip\noindent
Il reste à montrer que $f^{-1}$ est pleinement fidèle.
Soient $\mathcal{F}, \mathcal{G} \in \Ob(\Sh(\mathcal{C}))$.
Nous devons montrer que l'application
$$
\Mor_{\Sh(\mathcal{C})}(\mathcal{F}, \mathcal{G})
\longrightarrow
\Mor_{\Sh(\mathcal{D})}(f^{-1}\mathcal{F}, f^{-1}\mathcal{G})
$$
est bijective. Mais le membre de droite est égal à
\begin{align*}
\Mor_{\Sh(\mathcal{D})}(f^{-1}\mathcal{F}, f^{-1}\mathcal{G})
& =
\Mor_{\Sh(\mathcal{C})}(\mathcal{F}, f_*f^{-1}\mathcal{G}) \\
& =
\Mor_{\Sh(\mathcal{C})}(\mathcal{F}, g^{-1}f^{-1}\mathcal{G}) \\
& =
\Mor_{\Sh(\mathcal{C})}(\mathcal{F}, \mathcal{G})
\end{align*}
(la première égalité résultant de l'adjonction), ce qui démontre l'assertion.
\end{proof}

\begin{example}
\label{sites-example-closed-map-cocontinuous-false}
Soit $X$ un espace topologique.
Soit $i : Z  \to X$ l'inclusion d'une partie (munie de la topologie induite).
Considérons le foncteur $u : X_{Zar} \to Z_{Zar}$,
$U \mapsto u(U) = Z \cap U$.
À première vue, il pourrait sembler que ce foncteur soit également
cocontinu. Après tout, puisque $Z$ est muni de la topologie induite, tout
recouvrement de $U\cap Z$ ne devrait-il pas provenir d'un recouvrement de
$U$ dans $X$ ? Il n'en est rien ! En effet, que se passe-t-il si
$U \cap Z = \emptyset$ ? Dans ce cas, le recouvrement vide est un recouvrement
de $U \cap Z$, et le recouvrement vide ne peut être raffiné que par le
recouvrement vide. On en conclut donc que
$u$ cocontinu $\Rightarrow$ tout ouvert non vide $U$ de $X$ rencontre $Z$.
Mais cela ne suffit pas. Par exemple, si $X = \mathbf{R}$ est la droite réelle
munie de la topologie usuelle et si $Z = \mathbf{R} \setminus \{0\}$, alors
il existe un recouvrement ouvert de $Z$, à savoir
$Z = \{x < 0\} \cup \bigcup_n \{1/n < x\}$, qui ne peut être raffiné par la
restriction d'aucun recouvrement ouvert de $X$.
\end{example}






\section{Foncteurs cocontinus possédant un adjoint à droite}
\label{sites-section-cocontinuous-adjoint}

\noindent
Soient $\mathcal{C}$ et $\mathcal{D}$ des sites. Soient
$u : \mathcal{C} \to \mathcal{D}$ et $v : \mathcal{D} \to \mathcal{C}$ des
foncteurs des catégories sous-jacentes tels que $v$ soit adjoint à droite de
$u$. Dans ce cas, si $v$ est continu, alors $u$ est cocontinu
(Lemme \ref{sites-lemma-continuous-with-continuous-left-adjoint}).
Si $u$ est cocontinu, alors il arrive souvent (mais pas toujours ; voir
l'Exemple \ref{sites-example-not-equivalent}) que $v$ soit continu ; dans ce cas,
$v$ définit un morphisme de sites dont le morphisme de topos associé est le
même que celui défini par $u$.

\begin{lemma}
\label{sites-lemma-have-functor-other-way}
Soient $\mathcal{C}$ et $\mathcal{D}$ des sites. Soient
$u : \mathcal{C} \to \mathcal{D}$ et $v : \mathcal{D} \to \mathcal{C}$ des
foncteurs. Supposons que $u$ soit cocontinu et que $v$ soit adjoint à droite de
$u$. Soit $g : \Sh(\mathcal{C}) \to \Sh(\mathcal{D})$ le morphisme de topos
associé à $u$ ; voir le Lemme \ref{sites-lemma-cocontinuous-morphism-topoi}.
Alors
\begin{enumerate}
\item pour un faisceau $\mathcal{F}$ sur $\mathcal{C}$, le faisceau
$g_*\mathcal{F}$ est égal au préfaisceau $v^p\mathcal{F}$ ; autrement dit,
$(g_*\mathcal{F})(V) = \mathcal{F}(v(V))$, et
\item pour un faisceau $\mathcal{G}$ sur $\mathcal{D}$, on a
$g^{-1}\mathcal{G} = (v_p\mathcal{G})^\#$.
\end{enumerate}
\end{lemma}

\begin{proof}
Pour $\mathcal{F}$ comme en (1), on a
$$
g_*\mathcal{F} =
{}_su\mathcal{F} =
{}_pu\mathcal{F} =
v^p\mathcal{F} =
\mathcal{F} \circ v
$$
La première égalité résulte du Lemme \ref{sites-lemma-cocontinuous-morphism-topoi}.
La deuxième égalité résulte du Lemme \ref{sites-lemma-pu-sheaf}.
La troisième égalité résulte du Lemme \ref{sites-lemma-adjoint-functors}.
La dernière égalité est la définition de $v^p$ dans la
section \ref{sites-section-functoriality-PSh}. Cela démontre (1).
Pour $\mathcal{G}$ comme en (2), on a
$$
g^{-1}\mathcal{G} = (u^p\mathcal{G})^\# = (v_p\mathcal{G})^\#
$$
La première égalité résulte du Lemme \ref{sites-lemma-cocontinuous-morphism-topoi}.
La deuxième égalité résulte du Lemme \ref{sites-lemma-adjoint-functors}.
\end{proof}

\begin{lemma}
\label{sites-lemma-have-functor-other-way-morphism}
Les notations et les hypothèses sont celles du
Lemme \ref{sites-lemma-have-functor-other-way}.
Si, de plus, $v$ est continu, alors $v$ définit un morphisme de sites
$f : \mathcal{C} \to \mathcal{D}$ dont le morphisme de topos associé est égal
à $g$.
\end{lemma}

\begin{proof}
Nous utiliserons sans plus les mentionner les résultats du
Lemme \ref{sites-lemma-have-functor-other-way}. Pour montrer que $v$ définit un
morphisme de sites $f$ comme dans l'énoncé du lemme, nous devons montrer que
$v_s$ est un foncteur exact (voir la
Définition \ref{sites-definition-morphism-sites}). Comme
$v_s\mathcal{G} = (v_p\mathcal{G})^\# = g^{-1}\mathcal{G}$, cela résulte du
fait que $g$ est un morphisme de topos. On voit alors que
$f^{-1} = v_s = g^{-1}$ et l'on obtient $f = g$ comme morphismes de topos.
\end{proof}

\begin{example}
\label{sites-example-finer-topology-bis}
Cet exemple poursuit la discussion de l'Exemple \ref{sites-example-finer-topology},
auquel nous empruntons les notations $\mathcal{C}, \tau, \tau', \epsilon$.
Observons que le foncteur identité
$v : \mathcal{C}_{\tau'} \to \mathcal{C}_\tau$ est continu et que le foncteur
identité $u : \mathcal{C}_\tau \to \mathcal{C}_{\tau'}$ est cocontinu. De
plus, $u$ est adjoint à gauche de $v$. Les résultats des
Lemmes \ref{sites-lemma-have-functor-other-way} et
\ref{sites-lemma-have-functor-other-way-morphism} s'appliquent donc, et l'on conclut
que $v$ définit un morphisme de sites, à savoir
$$
\epsilon : \mathcal{C}_\tau \longrightarrow \mathcal{C}_{\tau'}
$$
dont le morphisme de topos correspondant est le même que le morphisme de topos
associé au foncteur cocontinu $u$.
\end{example}

\begin{lemma}
\label{sites-lemma-continuous-with-continuous-left-adjoint}
Soient $\mathcal{C}$ et $\mathcal{D}$ des sites et soit
$v : \mathcal{D} \to \mathcal{C}$ un foncteur continu.
Supposons que $v$ possède un adjoint à gauche
$u : \mathcal{C} \to \mathcal{D}$. Alors
\begin{enumerate}
\item $u$ est cocontinu,
\item les conclusions des Lemmes \ref{sites-lemma-have-functor-other-way} et
\ref{sites-lemma-have-functor-other-way-morphism} sont valables.
\end{enumerate}
En particulier, $v$ définit un morphisme de sites
$f : \mathcal{C} \to \mathcal{D}$.
\end{lemma}

\begin{proof}
Soit $U$ un objet de $\mathcal{C}$ et soit
$\{V_j \to u(U)\}_{j \in J}$ un recouvrement dans $\mathcal{D}$. Alors
$\{v(V_j) \to v(u(U))\}_{i \in I}$ est un recouvrement dans $\mathcal{C}$.
Par changement de base au moyen du morphisme d'adjonction
$U \to v(u(U))$, on obtient un recouvrement
$\{W_j \to U\}_{j \in J}$, où
$W_j = v(V_j) \times_{v(u(U))} U$. En notant
$p_j : W_j \to v(V_j)$ la première projection, on obtient des morphismes
$$
u(W_j) \xrightarrow{u(p_j)} u(v(V_j)) \longrightarrow V_j
$$
où la deuxième flèche est le morphisme d'adjonction. Cela détermine un
morphisme $\{u(W_j) \to u(U)\} \to \{V_j \to u(U)\}$ de familles de
morphismes à cible fixée, ce qui montre que $u$ est bien cocontinu. Les autres
assertions résultent immédiatement des
Lemmes \ref{sites-lemma-have-functor-other-way} et
\ref{sites-lemma-have-functor-other-way-morphism}.
\end{proof}

\begin{example}
\label{sites-example-not-equivalent}
Soient $\mathcal{C}$ et $\mathcal{D}$ des sites. Soient
$u : \mathcal{C} \to \mathcal{D}$ et $v : \mathcal{D} \to \mathcal{C}$ des
foncteurs des catégories sous-jacentes tels que $v$ soit adjoint à droite de
$u$. Le Lemme \ref{sites-lemma-continuous-with-continuous-left-adjoint} montre que,
si $v$ est continu, alors $u$ est cocontinu. Réciproquement, si $u$ est
cocontinu, on ne peut pas conclure que $v$ est continu. Nous donnerons un
exemple de ce phénomène à l'aide des grands sites \'etale et lisse d'un
schéma, mais il existe vraisemblablement aussi un exemple élémentaire.
Considérons en effet un schéma $S$ et les sites $(\Sch/S)_\etale$ et
$(\Sch/S)_{smooth}$. On peut supposer que ces sites ont la même catégorie
sous-jacente ; voir Topologies,
Remarque \ref{topologies-remark-choice-sites}.
Soit $u = v = \text{id}$. Alors $u$, considéré comme foncteur de
$(\Sch/S)_\etale$ vers $(\Sch/S)_{smooth}$, est cocontinu, car tout
recouvrement lisse d'un schéma peut être raffiné par un recouvrement \'etale ;
voir Compléments sur les morphismes,
Lemme \ref{more-morphisms-lemma-etale-dominates-smooth}.
Réciproquement, le foncteur $v$ de $(\Sch/S)_{smooth}$ vers
$(\Sch/S)_\etale$ n'est pas continu, car un recouvrement lisse n'est pas, en
général, un recouvrement \'etale.
\end{example}




\section{Foncteurs cocontinus possédant un adjoint à gauche}
\label{sites-section-cocontinuous-left-adjoint}

\noindent
Il peut arriver qu'un foncteur cocontinu $u$ possède un adjoint à gauche $w$.

\begin{lemma}
\label{sites-lemma-have-left-adjoint}
Soient $\mathcal{C}$ et $\mathcal{D}$ des sites. Soit
$g : \Sh(\mathcal{C}) \to \Sh(\mathcal{D})$ le morphisme de topos
associé à un foncteur continu et cocontinu
$u : \mathcal{C} \to \mathcal{D}$ ; voir les
Lemmes \ref{sites-lemma-cocontinuous-morphism-topoi} et
\ref{sites-lemma-when-shriek}.
\begin{enumerate}
\item Si $w : \mathcal{D} \to \mathcal{C}$ est adjoint à gauche de $u$, alors :
\begin{enumerate}
\item $g_!\mathcal{F}$ est le faisceau associé au préfaisceau
$w^p\mathcal{F}$ ;
\item $g_!$ est exact.
\end{enumerate}
\item Si $w$ est un adjoint à gauche continu, alors $g_!$
possède un adjoint à gauche.
\item Si $w$ est un adjoint à gauche cocontinu, alors $g_! = h^{-1}$ et
$g^{-1} = h_*$, où $h : \Sh(\mathcal{D}) \to \Sh(\mathcal{C})$ est
le morphisme de topos associé à $w$.
\end{enumerate}
\end{lemma}

\begin{proof}
Rappelons que $g_!\mathcal{F}$ est le faisceau associé à $u_p\mathcal{F}$.
Ainsi, (1)(a) résulte de l'égalité $u_p = w^p$ donnée par le
Lemme \ref{sites-lemma-adjoint-functors}.

\medskip\noindent
Pour établir (1)(b), remarquons que $g_!$ commute à toutes les limites
inductives, puisque $g_!$ est adjoint à gauche (Catégories, Lemme
\ref{categories-lemma-adjoint-exact}).
Soit $i \mapsto \mathcal{F}_i$ un diagramme fini dans $\Sh(\mathcal{C})$.
Alors $\lim \mathcal{F}_i$ se calcule dans la catégorie des préfaisceaux
(Lemme \ref{sites-lemma-limit-sheaf}). Comme $w^p$ est adjoint à droite
(Lemme \ref{sites-lemma-adjoints-u}), on a
$w^p \lim \mathcal{F}_i = \lim w^p\mathcal{F}_i$.
Puisque le passage au faisceau associé est exact
(Lemme \ref{sites-lemma-sheafification-exact}), on conclut par (1)(a).

\medskip\noindent
Supposons $w$ continu. Alors $g_! = (w^p\ )^\# = w^s$, mais le passage
au faisceau associé n'est pas nécessaire, et $w_s$ est un adjoint à gauche ;
voir les Lemmes \ref{sites-lemma-pushforward-sheaf} et
\ref{sites-lemma-adjoint-sheaves}.

\medskip\noindent
Supposons $w$ cocontinu. L'égalité $g_! = h^{-1}$ résulte de (1)(a)
et des définitions. L'égalité $g^{-1} = h_*$ résulte de
$g_! = h^{-1}$ et de l'unicité du foncteur adjoint.
On peut aussi la déduire du Lemme \ref{sites-lemma-have-functor-other-way}.
\end{proof}









\section{Existence de l'image directe extraordinaire}
\label{sites-section-lower-shriek}

\noindent
Dans cette section, nous étudions certains cas de morphismes de topos $f$
pour lesquels $f^{-1}$ possède un adjoint à gauche $f_!$.

\begin{lemma}
\label{sites-lemma-existence-lower-shriek}
Soient $\mathcal{C}$ et $\mathcal{D}$ deux sites.
Soit $f : \Sh(\mathcal{D}) \to \Sh(\mathcal{C})$ un morphisme de topos.
Soit $E \subset \Ob(\mathcal{D})$ une partie telle que :
\begin{enumerate}
\item pour $V \in E$, il existe un faisceau $\mathcal{G}$ sur
$\mathcal{C}$ tel que $f^{-1}\mathcal{F}(V) =
\Mor_{\Sh(\mathcal{C})}(\mathcal{G}, \mathcal{F})$ fonctoriellement
en $\mathcal{F}$ dans $\Sh(\mathcal{C})$ ;
\item tout objet de $\mathcal{D}$ admet un recouvrement par des objets de $E$.
\end{enumerate}
Alors $f^{-1}$ possède un adjoint à gauche $f_!$.
\end{lemma}

\begin{proof}
D'après le lemme de Yoneda (Catégories, Lemme
\ref{categories-lemma-yoneda}), le faisceau $\mathcal{G}_V$
correspondant à $V \in E$ est défini à isomorphisme unique près par la formule
$f^{-1}\mathcal{F}(V) =
\Mor_{\Sh(\mathcal{C})}(\mathcal{G}_V, \mathcal{F})$.
Rappelons que
$f^{-1}\mathcal{F}(V) =
\Mor_{\Sh(\mathcal{D})}(h_V^\#, f^{-1}\mathcal{F})$.
Notons $i_V : h_V^\# \to f^{-1}\mathcal{G}_V$ le morphisme correspondant à
$\text{id}$ dans $\Mor(\mathcal{G}_V, \mathcal{G}_V)$.
La fonctorialité de (1) implique que la bijection est donnée par
$$
\Mor_{\Sh(\mathcal{C})}(\mathcal{G}_V, \mathcal{F}) \to
\Mor_{\Sh(\mathcal{D})}(h_V^\#, f^{-1}\mathcal{F}),\quad
\varphi \mapsto f^{-1}\varphi \circ i_V
$$
Pour tous $V_1, V_2 \in E$, il existe une application canonique
$$
\Mor_{\Sh(\mathcal{D})}(h^\#_{V_2}, h^\#_{V_1})
\to
\Hom_{\Sh(\mathcal{C})}(\mathcal{G}_{V_2}, \mathcal{G}_{V_1}),\quad
\varphi \mapsto f_!(\varphi)
$$
caractérisée par
$f^{-1}(f_!(\varphi)) \circ i_{V_2} = i_{V_1} \circ \varphi$.
Remarquons que $\varphi \mapsto f_!(\varphi)$ est compatible à la composition ;
cela se voit directement à partir de la caractérisation.
Ainsi, $h_V^\# \mapsto \mathcal{G}_V$ et
$\varphi \mapsto f_!\varphi$ définissent un foncteur sur la sous-catégorie
pleine de $\Sh(\mathcal{D})$ dont les objets sont les $h_V^\#$.

\medskip\noindent
Soit $J$ un ensemble et soit $J \to E$, $j \mapsto V_j$, une application.
Nous avons alors une bijection fonctorielle
$$
\Mor_{\Sh(\mathcal{C})}(\coprod \mathcal{G}_{V_j}, \mathcal{F})
\longrightarrow
\Mor_{\Sh(\mathcal{D})}(\coprod h_{V_j}^\#, f^{-1}\mathcal{F})
$$
obtenue comme produit des bijections précédentes. Nous pouvons donc prolonger
le foncteur $f_!$ à la sous-catégorie pleine de $\Sh(\mathcal{D})$
dont les objets sont les coproduits de $h_V^\#$ avec $V \in E$.

\medskip\noindent
Étant donné un faisceau arbitraire $\mathcal{H}$ sur $\mathcal{D}$,
choisissons un diagramme coégalisateur
$$
\xymatrix{
\mathcal{H}_1 \ar@<1ex>[r] \ar@<-1ex>[r] &
\mathcal{H}_0 \ar[r] &
\mathcal{H}
}
$$
où $\mathcal{H}_i = \coprod h_{V_{i, j}}^\#$
est un coproduit avec $V_{i, j} \in E$.
Cela est possible par l'hypothèse (2) ; voir le
Lemme \ref{sites-lemma-sheaf-coequalizer-representable}
(pour qui s'inquiéterait des questions ensemblistes, remarquons que
la construction donnée dans le
Lemme \ref{sites-lemma-sheaf-coequalizer-representable} est canonique).
Définissons $f_!(\mathcal{H})$ comme le faisceau sur $\mathcal{C}$ qui fait de
$$
\xymatrix{
f_!\mathcal{H}_1 \ar@<1ex>[r] \ar@<-1ex>[r] &
f_!\mathcal{H}_0 \ar[r] &
f_!\mathcal{H}
}
$$
un diagramme coégalisateur. Alors
\begin{align*}
\Mor(f_!\mathcal{H}, \mathcal{F})
& =
\text{Égalisateur}(
\xymatrix{
\Mor(f_!\mathcal{H}_0, \mathcal{F}) \ar@<1ex>[r] \ar@<-1ex>[r] &
\Mor(f_!\mathcal{H}_1, \mathcal{F})
}
) \\
& =
\text{Égalisateur}(
\xymatrix{
\Mor(\mathcal{H}_0, f^{-1}\mathcal{F}) \ar@<1ex>[r] \ar@<-1ex>[r] &
\Mor(\mathcal{H}_1, f^{-1}\mathcal{F})
}
) \\
& =
\Hom(\mathcal{H}, f^{-1}\mathcal{F})
\end{align*}
Nous voyons ainsi que $f_!$ se prolonge à toute la catégorie des faisceaux
sur $\mathcal{D}$.
\end{proof}








\section{Localisation}
\label{sites-section-localize}

\noindent
Soit $\mathcal{C}$ un site.
Soit $U \in \Ob(\mathcal{C})$.
Voir Catégories, Exemple \ref{categories-example-category-over-X},
pour la définition de la catégorie $\mathcal{C}/U$ des objets au-dessus de $U$.
Nous munissons $\mathcal{C}/U$ d'une structure de site en déclarant qu'une
famille de morphismes $\{V_j \to V\}$ d'objets au-dessus de $U$ est un
recouvrement de $\mathcal{C}/U$ si et seulement si elle est un recouvrement
dans $\mathcal{C}$. Considérons le foncteur d'oubli
$$
j_U : \mathcal{C}/U \longrightarrow \mathcal{C}.
$$
Il est manifestement cocontinu et continu. Par le
Lemme \ref{sites-lemma-when-shriek}, nous obtenons donc un morphisme de topos
$$
j_U : \Sh(\mathcal{C}/U) \longrightarrow \Sh(\mathcal{C})
$$
donné par $j_U^{-1}$ et $j_{U*}$, ainsi qu'un foncteur $j_{U!}$
adjoint à gauche de $j_U^{-1}$.

\begin{definition}
\label{sites-definition-localize}
Soit $\mathcal{C}$ un site.
Soit $U \in \Ob(\mathcal{C})$.
\begin{enumerate}
\item Le site $\mathcal{C}/U$ est appelé la {\it localisation du site
$\mathcal{C}$ en l'objet $U$}.
\item Le morphisme de topos
$j_U : \Sh(\mathcal{C}/U) \to \Sh(\mathcal{C})$
est appelé le {\it morphisme de localisation}.
\item Le foncteur $j_{U*}$ est appelé le {\it foncteur image directe}.
\item Pour un faisceau $\mathcal{F}$ sur $\mathcal{C}$, le faisceau
$j_U^{-1}\mathcal{F}$ est appelé la {\it restriction de $\mathcal{F}$
à $\mathcal{C}/U$}.
\item Pour un faisceau $\mathcal{G}$ sur $\mathcal{C}/U$, le faisceau
$j_{U!}\mathcal{G}$ est appelé le
{\it prolongement de $\mathcal{G}$ par l'ensemble vide}.
\end{enumerate}
\end{definition}

\noindent
La restriction $j_U^{-1}\mathcal{F}$ est le faisceau défini par la règle
$j_U^{-1}\mathcal{F}(X/U) = \mathcal{F}(X)$, comme prévu.
Le prolongement par l'ensemble vide admet lui aussi une description très simple
dans ce cas ; la voici.

\begin{lemma}
\label{sites-lemma-describe-j-shriek}
Soit $\mathcal{C}$ un site.
Soit $U \in \Ob(\mathcal{C})$.
Soit $\mathcal{G}$ un préfaisceau sur $\mathcal{C}/U$.
Alors $j_{U!}(\mathcal{G}^\#)$ est le faisceau associé au préfaisceau
$$
V
\longmapsto
\coprod\nolimits_{\varphi \in \Mor_\mathcal{C}(V, U)}
\mathcal{G}(V \xrightarrow{\varphi} U)
$$
muni des applications de restriction évidentes.
\end{lemma}

\begin{proof}
D'après le Lemme \ref{sites-lemma-when-shriek}, on a
$j_{U!}(\mathcal{G}^\#) = ((j_U)_p\mathcal{G}^\#)^\#$.
D'après le Lemme \ref{sites-lemma-technical-up}, ceci est égal à
$((j_U)_p\mathcal{G})^\#$.
Il suffit donc de démontrer que $(j_U)_p$ est donné par la formule ci-dessus
pour tout préfaisceau $\mathcal{G}$ sur $\mathcal{C}/U$.
D'après la définition de la
section \ref{sites-section-functoriality-PSh}, on a
$$
(j_U)_p\mathcal{G}(V)
=
\colim_{(W/U, V \to W)} \mathcal{G}(W)
$$
Il est maintenant clair que la catégorie des couples $(W/U, V \to W)$
possède un objet $O_\varphi = (\varphi : V \to U, \text{id} : V \to V)$
pour tout $\varphi : V \to U$ et que, de plus, tout objet reçoit un unique
morphisme depuis l'un des $O_\varphi$. Le résultat s'ensuit.
\end{proof}

\begin{lemma}
\label{sites-lemma-describe-j-shriek-representable}
Soit $\mathcal{C}$ un site.
Soit $U \in \Ob(\mathcal{C})$.
Soit $X/U$ un objet de $\mathcal{C}/U$.
Alors $j_{U!}(h_{X/U}^\#) = h_X^\#$.
\end{lemma}

\begin{proof}
Notons $p : X \to U$ le morphisme structural de $X$.
D'après le Lemme \ref{sites-lemma-describe-j-shriek},
$j_{U!}(h_{X/U}^\#)$ est le faisceau associé au préfaisceau
$$
V
\longmapsto
\coprod\nolimits_{\varphi \in \Mor_\mathcal{C}(V, U)}
\{\psi : V \to X \mid p \circ \psi = \varphi\}
$$
Celui-ci n'est autre que $\Mor_\mathcal{C}(V, X)$.
Le lemme en résulte.
\end{proof}

\noindent
On a $j_{U!}(*) = h_U^\#$ par chacun des deux lemmes précédents.
Ainsi, pour tout faisceau $\mathcal{G}$ sur $\mathcal{C}/U$, il existe
un morphisme canonique de faisceaux $j_{U!}\mathcal{G} \to h_U^\#$.
Cela caractérise les faisceaux appartenant à l'image essentielle de $j_{U!}$.

\begin{lemma}
\label{sites-lemma-essential-image-j-shriek}
Soit $\mathcal{C}$ un site.
Soit $U \in \Ob(\mathcal{C})$.
Le foncteur $j_{U!}$ définit une équivalence de catégories
$$
\Sh(\mathcal{C}/U)
\longrightarrow
\Sh(\mathcal{C})/h_U^\#
$$
\end{lemma}

\begin{proof}
Notons les objets de $\mathcal{C}/U$ sous la forme de couples $(X, a)$,
où $X$ est un objet de $\mathcal{C}$ et $a : X \to U$ un morphisme de
$\mathcal{C}$. De même, les objets de $\Sh(\mathcal{C})/h_U^\#$ sont
des couples $(\mathcal{F}, \varphi)$.
Le foncteur $\Sh(\mathcal{C}/U) \to \Sh(\mathcal{C})/h_U^\#$
envoie $\mathcal{G}$ sur le couple $(j_{U!}\mathcal{G}, \gamma)$,
où $\gamma$ est le composé de $j_{U!}\mathcal{G} \to j_{U!}*$
avec l'identification $j_{U!}* = h_U^\#$.

\medskip\noindent
Construisons un foncteur de $\Sh(\mathcal{C})/h_U^\#$ vers
$\Sh(\mathcal{C}/U)$.
Supposons donné $(\mathcal{F}, \varphi)$.
Pour un objet $(X, a)$ de $\mathcal{C}/U$, considérons l'ensemble
$\mathcal{F}_\varphi(X, a)$ des éléments $s \in \mathcal{F}(X)$ dont l'image
par $\varphi$ est l'image de
$a \in \Mor_\mathcal{C}(X, U) = h_U(X)$ dans $h_U^\#(X)$.
Il est immédiat que $(X, a) \mapsto \mathcal{F}_\varphi(X, a)$ est un faisceau
sur $\mathcal{C}/U$. Il est clair que la règle
$(\mathcal{F}, \varphi) \mapsto \mathcal{F}_\varphi$
définit un foncteur $\Sh(\mathcal{C})/h_U^\# \to \Sh(\mathcal{C}/U)$.

\medskip\noindent
Considérons aussi le foncteur
$\textit{PSh}(\mathcal{C})/h_U \to \textit{PSh}(\mathcal{C}/U)$,
$(\mathcal{F}, \varphi) \mapsto \mathcal{F}_\varphi$,
où $\mathcal{F}_\varphi(X, a)$ est défini comme l'ensemble des éléments
de $\mathcal{F}(X)$ qui s'envoient sur $a \in h_U(X)$.
Nous affirmons que le diagramme
$$
\xymatrix{
\textit{PSh}(\mathcal{C})/h_U \ar[r] \ar[d] &
\textit{PSh}(\mathcal{C}/U) \ar[d] \\
\Sh(\mathcal{C})/h_U^\# \ar[r] &
\Sh(\mathcal{C}/U)
}
$$
est commutatif, où les flèches verticales sont données par le passage aux
faisceaux associés.
Pour le voir\footnote{Une autre méthode consiste à décrire
$\mathcal{F}_\varphi$ par le diagramme cartésien
$$
\vcenter{
\xymatrix{
\mathcal{F}_\varphi \ar[r] \ar[d] & {*} \ar[d] \\
\mathcal{F}|_{\mathcal{C}/U} \ar[r] & h_U|_{\mathcal{C}/U}
}
}
\quad\text{pour les préfaisceaux et}\quad
\vcenter{
\xymatrix{
\mathcal{F}_\varphi \ar[r] \ar[d] & {*} \ar[d] \\
\mathcal{F}|_{\mathcal{C}/U} \ar[r] & h_U^\#|_{\mathcal{C}/U}
}
}
$$
pour les faisceaux, puis à utiliser le fait que la restriction à
$\mathcal{C}/U$ commute au passage au faisceau associé.}, il suffit de
montrer que la construction commute au foncteur
$\mathcal{F} \mapsto \mathcal{F}^+$ des
Lemmes \ref{sites-lemma-plus-presheaf} et \ref{sites-lemma-plus-functorial},
ainsi que du Théorème \ref{sites-theorem-plus}.
La commutation avec $\mathcal{F} \mapsto \mathcal{F}^+$ résulte du fait que,
pour $(X, a)$ donné, la catégorie des recouvrements de $(X, a)$ dans
$\mathcal{C}/U$ s'identifie canoniquement à celle des recouvrements de $X$
dans $\mathcal{C}$.

\medskip\noindent
Ensuite, faisons envoyer par
$\textit{PSh}(\mathcal{C}/U) \to \textit{PSh}(\mathcal{C})/h_U$
le préfaisceau $\mathcal{G}$ sur le couple
$(j_{U!}^{PSh}\mathcal{G}, \gamma)$, où
$j_{U!}^{PSh}\mathcal{G}$ est le préfaisceau défini par la formule du
Lemme \ref{sites-lemma-describe-j-shriek}, et où $\gamma$ est le composé de
$j_{U!}^{PSh}\mathcal{G} \to j_{U!}*$ avec l'identification
$j_{U!}^{PSh}* = h_U$ (évidente d'après la formule).
Il est alors immédiatement clair que le diagramme
$$
\xymatrix{
\textit{PSh}(\mathcal{C}/U) \ar[r] \ar[d] &
\textit{PSh}(\mathcal{C})/h_U \ar[d] \\
\Sh(\mathcal{C}/U) \ar[r] &
\Sh(\mathcal{C})/h_U^\#
}
$$
est commutatif, où les flèches verticales sont les passages aux faisceaux
associés. En réunissant ces résultats, il suffit de montrer qu'il existe
des isomorphismes fonctoriels
$(j_{U!}^{PSh}\mathcal{G})_\gamma = \mathcal{G}$
pour $\mathcal{G}$ dans $\textit{PSh}(\mathcal{C}/U)$, et
$j_{U!}^{PSh}\mathcal{F}_\varphi = \mathcal{F}$
pour $(\mathcal{F}, \varphi)$ dans $\textit{PSh}(\mathcal{C})/h_U$.
La valeur du préfaisceau $(j_{U!}^{PSh}\mathcal{G})_\gamma$
en $(X, a)$ est la fibre de l'application
$$
\coprod\nolimits_{a' : X \to U} \mathcal{G}(X, a') \to \Mor_\mathcal{C}(X, U)
$$
au-dessus de $a$, laquelle est $\mathcal{G}(X, a)$.
Cela démontre la première égalité.
La valeur du préfaisceau $j_{U!}^{PSh}\mathcal{F}_\varphi$ en $X$ est
$$
\coprod\nolimits_{a : X \to U} \mathcal{F}_\varphi(X, a) =
\mathcal{F}(X)
$$
car, pour toute application d'ensembles $S \to S'$, l'ensemble $S$ est
la réunion disjointe de ses fibres.
\end{proof}

\noindent
Le Lemme \ref{sites-lemma-essential-image-j-shriek} affirme que le foncteur
$j_{U!}$ est le composé
$$
\Sh(\mathcal{C}/U) \rightarrow
\Sh(\mathcal{C})/h_U^\# \rightarrow
\Sh(\mathcal{C})
$$
où la première flèche est une équivalence.

\begin{lemma}
\label{sites-lemma-j-shriek-commutes-equalizers-fibre-products}
Soit $\mathcal{C}$ un site. Soit $U \in \Ob(\mathcal{C})$.
Le foncteur $j_{U!}$ commute aux produits fibrés et aux égalisateurs
(et, plus généralement, aux limites finies connexes). En particulier, si
$\mathcal{F} \subset \mathcal{F}'$ dans $\Sh(\mathcal{C}/U)$, alors
$j_{U!}\mathcal{F} \subset j_{U!}\mathcal{F}'$.
\end{lemma}

\begin{proof}
Grâce au Lemme \ref{sites-lemma-essential-image-j-shriek} et au fait qu'une
équivalence de catégories commute à toutes les limites, cela se ramène
au fait que le foncteur
$\Sh(\mathcal{C})/h_U^\# \rightarrow \Sh(\mathcal{C})$
commute aux produits fibrés et aux égalisateurs.
On peut également le démontrer directement à partir de la description de
$j_{U!}$ donnée au Lemme \ref{sites-lemma-describe-j-shriek}, en utilisant
l'exactitude du passage au faisceau associé. (De même, lorsque $\mathcal{C}$
possède des produits fibrés et des égalisateurs, le résultat découle du
Lemme \ref{sites-lemma-preserve-equalizers}.)
\end{proof}

\begin{lemma}
\label{sites-lemma-j-shriek-reflects-injections-surjections}
Soit $\mathcal{C}$ un site. Soit $U \in \Ob(\mathcal{C})$.
Le foncteur $j_{U!}$ reflète les injections et les surjections.
\end{lemma}

\begin{proof}
Nous devons montrer que $j_{U!}$ reflète les monomorphismes et les
épimorphismes ; voir le Lemme \ref{sites-lemma-mono-epi-sheaves}.
Grâce au Lemme \ref{sites-lemma-essential-image-j-shriek}, cela se ramène au fait
que le foncteur $\Sh(\mathcal{C})/h_U^\# \to \Sh(\mathcal{C})$
reflète les monomorphismes et les épimorphismes.
\end{proof}

\begin{lemma}
\label{sites-lemma-compute-j-shriek-restrict}
Soit $\mathcal{C}$ un site. Soit $U \in \Ob(\mathcal{C})$.
Pour tout faisceau $\mathcal{F}$ sur $\mathcal{C}$, on a
$j_{U!}j_U^{-1}\mathcal{F} = \mathcal{F} \times h_U^\#$.
\end{lemma}

\begin{proof}
Cela résulte immédiatement de la description de $j_{U!}$ donnée au
Lemme \ref{sites-lemma-describe-j-shriek}.
\end{proof}

\begin{lemma}
\label{sites-lemma-relocalize}
Soit $\mathcal{C}$ un site.
Soit $f : V \to U$ un morphisme de $\mathcal{C}$.
Il existe alors un diagramme commutatif
$$
\xymatrix{
\mathcal{C}/V \ar[rd]_{j_V} \ar[rr]_j & &
\mathcal{C}/U \ar[ld]^{j_U} \\
& \mathcal{C} &
}
$$
de foncteurs continus et cocontinus.
Le foncteur $j : \mathcal{C}/V \to \mathcal{C}/U$,
$(a : W \to V) \mapsto (f \circ a : W \to U)$,
s'identifie au foncteur
$j_{V/U} : (\mathcal{C}/U)/(V/U) \to \mathcal{C}/U$
par l'identification $(\mathcal{C}/U)/(V/U) = \mathcal{C}/V$.
En outre, on a $j_{V!} = j_{U!} \circ j_!$,
$j_V^{-1} = j^{-1} \circ j_U^{-1}$ et
$j_{V*} = j_{U*} \circ j_*$.
\end{lemma}

\begin{proof}
La commutativité du diagramme est immédiate.
L'identification de $j$ à $j_{V/U}$ résulte des définitions.
D'après le Lemme \ref{sites-lemma-composition-cocontinuous}, le diagramme suivant
de morphismes de topos
\begin{equation}
\label{sites-equation-relocalize}
\vcenter{
\xymatrix{
\Sh(\mathcal{C}/V) \ar[rd]_{j_V} \ar[rr]_j & &
\Sh(\mathcal{C}/U) \ar[ld]^{j_U} \\
& \Sh(\mathcal{C}) &
}
}
\end{equation}
est commutatif. Cela démontre les égalités
$j_V^{-1} = j^{-1} \circ j_U^{-1}$ et
$j_{V*} = j_{U*} \circ j_*$.
L'égalité $j_{V!} = j_{U!} \circ j_!$ résulte formellement des propriétés
d'adjonction.
\end{proof}

\begin{lemma}
\label{sites-lemma-relocalize-explicit}
Notons $\mathcal{C}$, $f : V \to U$, $j_U$, $j_V$ et $j$ comme au
Lemme \ref{sites-lemma-relocalize}. Par les identifications
$\Sh(\mathcal{C}/V) = \Sh(\mathcal{C})/h_V^\#$ et
$\Sh(\mathcal{C}/U) = \Sh(\mathcal{C})/h_U^\#$ du
Lemme \ref{sites-lemma-essential-image-j-shriek}, on a :
\begin{enumerate}
\item le foncteur $j^{-1}$ admet la description suivante :
$$
j^{-1}(\mathcal{H} \xrightarrow{\varphi} h_U^\#)
=
(\mathcal{H} \times_{\varphi, h_U^\#, f} h_V^\# \to h_V^\#).
$$
\item le foncteur $j_!$ admet la description suivante :
$$
j_!(\mathcal{H} \xrightarrow{\varphi} h_V^\#) =
(\mathcal{H} \xrightarrow{h_f \circ \varphi} h_U^\#)
$$
\end{enumerate}
\end{lemma}

\begin{proof}
Démonstration de (2). Rappelons que l'identification
$\Sh(\mathcal{C}/V) \to \Sh(\mathcal{C})/h_V^\#$
envoie $\mathcal{G}$ sur
$j_{V!}\mathcal{G} \to j_{V!}(*) = h_V^\#$,
et qu'il en va de même pour
$\Sh(\mathcal{C}/U) \to \Sh(\mathcal{C})/h_U^\#$.
Ainsi, $j_!\mathcal{G}$ s'envoie sur
$j_{U!}(j_!\mathcal{G}) \to j_{U!}(*) = h_U^\#$,
et (2) résulte de l'égalité $j_{U!}j_! = j_{V!}$ donnée par le
Lemme \ref{sites-lemma-relocalize}.

\medskip\noindent
Le lecteur peut maintenant démontrer (1) en utilisant le fait que $j^{-1}$
est adjoint à droite de $j_!$, et que la règle de (1) est adjointe à droite
de celle de (2). Voici une démonstration directe.
Supposons que $\varphi : \mathcal{H} \to h_U^\#$ soit un objet de
$\Sh(\mathcal{C})/h_U^\#$. D'après la démonstration du
Lemme \ref{sites-lemma-essential-image-j-shriek}, il correspond au faisceau
$\mathcal{H}_\varphi$ sur $\mathcal{C}/U$ défini par la règle
$$
(a : W \to U)
\longmapsto
\{ s \in \mathcal{H}(W) \mid \varphi(s) = a\}
$$
sur $\mathcal{C}/U$. L'image inverse $j^{-1}\mathcal{H}_\varphi$ sur
$\mathcal{C}/V$ est donnée par la règle
$$
(a : W \to V)
\longmapsto
\{ s \in \mathcal{H}(W) \mid \varphi(s) = f \circ a\}
$$
d'après la description de $j^{-1} = j_{U/V}^{-1}$ comme restriction de
$\mathcal{H}_\varphi$ à $\mathcal{C}/V$.
D'autre part, en appliquant la règle à l'objet
$$
\xymatrix{
\mathcal{H}' = \mathcal{H} \times_{\varphi, h_U^\#, f} h_V^\#
\ar[rr]^-{\varphi'} & & h_V^\#
}
$$
de $\Sh(\mathcal{C})/h_V^\#$, nous obtenons
$\mathcal{H}'_{\varphi'}$ donné par
\begin{align*}
(a : W \to V)
\longmapsto
&  \{ s' \in \mathcal{H}'(W) \mid \varphi'(s') = a\} \\
= &
\{ (s, a') \in \mathcal{H}(W) \times h_V^\#(W) \mid
a' = a \text{ et } \varphi(s) = f \circ a'\}
\end{align*}
qui est exactement la règle décrivant $j^{-1}\mathcal{H}_\varphi$ ci-dessus.
\end{proof}

\begin{remark}
\label{sites-remark-localize-presheaves}
Localisation et préfaisceaux. Soit $\mathcal{C}$ une catégorie.
Soit $U$ un objet de $\mathcal{C}$. À strictement parler, les foncteurs
$j_U^{-1}$, $j_{U*}$ et $j_{U!}$ n'ont pas été définis pour les préfaisceaux.
Nous pouvons cependant considérer un préfaisceau comme un faisceau pour la
topologie chaotique sur $\mathcal{C}$ (voir l'Exemple
\ref{sites-example-indiscrete}). Nous obtenons donc aussi un foncteur
$$
j_U^{-1} :
\textit{PSh}(\mathcal{C})
\longrightarrow
\textit{PSh}(\mathcal{C}/U)
$$
et des foncteurs
$$
j_{U*}, j_{U!} :
\textit{PSh}(\mathcal{C}/U)
\longrightarrow
\textit{PSh}(\mathcal{C})
$$
qui sont respectivement adjoints à droite et à gauche de $j_U^{-1}$.
D'après le Lemme \ref{sites-lemma-describe-j-shriek}, le préfaisceau
$j_{U!}\mathcal{G}$ est donné par
$$
V \longmapsto
\coprod\nolimits_{\varphi \in \Mor_\mathcal{C}(V, U)}
\mathcal{G}(V \xrightarrow{\varphi} U)
$$
De plus, le foncteur $j_{U!}$ commute aux produits fibrés et aux égalisateurs.
\end{remark}

\begin{remark}
\label{sites-remark-localization-cartesian-cocontinuous}
Soit $\mathcal{C}$ un site. Soit $U \to V$ un morphisme de $\mathcal{C}$.
Les foncteurs cocontinus $\mathcal{C}/U \to \mathcal{C}$ et
$j : \mathcal{C}/U \to \mathcal{C}/V$ (Lemme \ref{sites-lemma-relocalize})
satisfont la propriété $P$ de la Remarque
\ref{sites-remark-cartesian-cocontinuous}.
Par exemple, si l'on se donne des objets $(X/U)$, $(W/V)$, un morphisme
$g : j(X/U) \to (W/V)$ et un recouvrement
$\{f_i  : (W_i/V) \to (W/V)\}$, alors
$(X \times_W W_i/U)$ est une incarnation de
$(X/U) \times_{g, (W/V), f_i} (W_i/V)$, et la famille
$\{(X \times_W W_i/U) \to (X/U)\}$ est un recouvrement de $\mathcal{C}/U$.
\end{remark}



\section{Recollement de faisceaux}
\label{sites-section-glueing-sheaves}

\noindent
Cette section est l'analogue de
Faisceaux, section \ref{sheaves-section-glueing-sheaves}.

\begin{lemma}
\label{sites-lemma-glue-maps}
\begin{slogan}
Les morphismes de faisceaux se recollent.
\end{slogan}
Soit $\mathcal{C}$ un site.
Soit $\{U_i \to U\}$ un recouvrement du site $\mathcal{C}$.
Soient $\mathcal{F}$, $\mathcal{G}$ des faisceaux sur $\mathcal{C}$.
Étant donnée une famille
$$
\varphi_i :
\mathcal{F}|_{\mathcal{C}/U_i}
\longrightarrow
\mathcal{G}|_{\mathcal{C}/U_i}
$$
de morphismes de faisceaux telle que, pour tous $i, j \in I$, les morphismes
$\varphi_i, \varphi_j$ induisent par restriction le même morphisme
$\varphi_{ij} : \mathcal{F}|_{\mathcal{C}/U_i \times_U U_j} \to
\mathcal{G}|_{\mathcal{C}/U_i \times_U U_j}$,
il existe un unique morphisme de faisceaux
$$
\varphi :
\mathcal{F}|_{\mathcal{C}/U}
\longrightarrow
\mathcal{G}|_{\mathcal{C}/U}
$$
dont la restriction à chaque $\mathcal{C}/U_i$ coïncide avec $\varphi_i$.
\end{lemma}

\begin{proof}
Les restrictions employées dans le lemme sont celles du
Lemme \ref{sites-lemma-relocalize}.
Soit $V/U$ un objet de $\mathcal{C}/U$.
Posons $V_i = U_i \times_U V$ et notons $\mathcal{V} = \{V_i \to V\}$.
Remarquons que $(U_i \times_U U_j) \times_U V = V_i \times_V V_j$.
Nous avons alors
$\mathcal{F}|_{\mathcal{C}/U_i}(V_i/U_i) = \mathcal{F}(V_i)$
et
$\mathcal{F}|_{\mathcal{C}/U_i \times_U U_j}(V_i \times_V V_j/U_i \times_U U_j)
= \mathcal{F}(V_i \times_V V_j)$,
et de même pour $\mathcal{G}$.
Nous pouvons donc définir $\varphi$ sur les sections au-dessus de $V$
par les flèches pointillées du diagramme
$$
\xymatrix{
\mathcal{F}(V) \ar@{=}[r] &
H^0(\mathcal{V}, \mathcal{F}) \ar@{..>}[d] \ar[r] &
\prod \mathcal{F}(V_i)
\ar[d]_{\prod \varphi_i}
\ar@<1ex>[r] \ar@<-1ex>[r] &
\prod \mathcal{F}(V_i \times_V V_j) \ar[d]_{\prod \varphi_{ij}} \\
\mathcal{G}(V) \ar@{=}[r] &
H^0(\mathcal{V}, \mathcal{G}) \ar[r] &
\prod \mathcal{G}(V_i)
\ar@<1ex>[r] \ar@<-1ex>[r] &
\prod \mathcal{G}(V_i \times_V V_j)
}
$$
Les signes d'égalité proviennent de la condition de faisceau ; voir la
section \ref{sites-section-sheafification} pour la notation
$H^0(\mathcal{V}, -)$.
Nous omettons la vérification que ces morphismes sont compatibles
avec les applications de restriction.
\end{proof}

\noindent
Le lemme précédent implique qu'étant donnés deux faisceaux $\mathcal{F}$,
$\mathcal{G}$ sur un site $\mathcal{C}$, la règle
$$
U \longmapsto
\Mor_{\Sh(\mathcal{C}/U)}(
\mathcal{F}|_{\mathcal{C}/U}, \mathcal{G}|_{\mathcal{C}/U})
$$
définit un faisceau $\SheafHom(\mathcal{F},\mathcal{G})$.
Il s'agit d'une sorte de {\it faisceau Hom interne}. On l'emploie
rarement dans le cadre des faisceaux d'ensembles, et plus souvent
dans celui des faisceaux de modules ; voir
Modules sur les sites, section \ref{sites-modules-section-internal-hom}.

\begin{lemma}
\label{sites-lemma-internal-hom-sheaf}
\begin{slogan}
La catégorie des faisceaux sur un site est cartésienne fermée
\end{slogan}
Soit $\mathcal{C}$ un site. Soient $\mathcal{F}$, $\mathcal{G}$ et
$\mathcal{H}$ des faisceaux sur $\mathcal{C}$. Il existe une bijection canonique
$$
\Mor_{\Sh(\mathcal{C})}(\mathcal{F}\times\mathcal{G},\mathcal{H}) =
\Mor_{\Sh(\mathcal{C})}(\mathcal{F},\SheafHom(\mathcal{G},\mathcal{H}))
$$
qui est fonctorielle en ses trois arguments.
\end{lemma}

\begin{proof}
Le lemme affirme que les foncteurs $-\times\mathcal{G}$ et
$\SheafHom(\mathcal{G},-)$ sont adjoints l'un de l'autre.
Pour le montrer, nous utilisons les notions d'unité et de counité ; voir
Catégories, section \ref{categories-section-adjoint}.
L'unité
$$
\eta_\mathcal{F} :
\mathcal{F}
\longrightarrow
\SheafHom(\mathcal{G},\mathcal{F}\times\mathcal{G})
$$
envoie $s \in \mathcal{F}(U)$ sur le morphisme
$\mathcal{G}|_{\mathcal{C}/U} \to
\mathcal{F}|_{\mathcal{C}/U}\times\mathcal{G}|_{\mathcal{C}/U}$
qui, au-dessus de $V/U$, est donné par
$$
\mathcal{G}(V) \longrightarrow \mathcal{F}(V)\times \mathcal{G}(V), \quad
t \longmapsto (s|_{V},t).
$$
La counité
$$
\epsilon_{\mathcal{H}} :
\SheafHom(\mathcal{G}, \mathcal{H}) \times \mathcal{G}
\longrightarrow
\mathcal{H}
$$
est le morphisme d'évaluation. Il est donné par la règle
$$
\Mor_{\Sh(\mathcal{C}/U)}(
\mathcal{G}|_{\mathcal{C}/U}, \mathcal{H}|_{\mathcal{C}/U})
\times \mathcal{G}(U)
\longrightarrow
\mathcal{H}(U),\quad
(\varphi, s) \longmapsto \varphi(s).
$$
Alors, pour tout $\varphi : \mathcal{F} \times \mathcal{G} \to \mathcal{H}$,
le morphisme correspondant
$\mathcal{F} \to \SheafHom(\mathcal{G},\mathcal{H})$
est défini en envoyant chaque section $s \in \mathcal{F}(U)$
sur le morphisme de faisceaux sur $\mathcal{C}/U$ qui, sur les
sections au-dessus de $V/U$, est donné par
$$
\mathcal{G}(V) \longrightarrow \mathcal{H}(V),\quad
t \longmapsto \varphi(s|_V, t).
$$
Réciproquement, pour tout
$\psi : \mathcal{F} \to \SheafHom(\mathcal{G}, \mathcal{H})$,
le morphisme correspondant
$\mathcal{F} \times \mathcal{G} \to \mathcal{H}$ est donné par
$$
\mathcal{F}(U) \times \mathcal{G}(U) \longrightarrow \mathcal{H}(U),\quad
(s, t) \longmapsto \psi(s)(t)
$$
sur les sections au-dessus d'un objet $U$. Nous omettons les détails montrant
que ces constructions sont inverses l'une de l'autre.
\end{proof}

\begin{lemma}
\label{sites-lemma-hom-sheaf-hU}
Soit $\mathcal{C}$ un site et soit $U \in \Ob(\mathcal{C})$.
Alors $\SheafHom(h_U^\#, \mathcal{F}) = j_*(\mathcal{F}|_{\mathcal{C}/U})$
pour $\mathcal{F}$ dans $\Sh(\mathcal{C})$.
\end{lemma}

\begin{proof}
On peut le montrer en construisant directement un isomorphisme
de faisceaux. Nous raisonnons plutôt comme suit.
Soit $\mathcal{G}$ un faisceau sur $\mathcal{C}$.
Alors
\begin{align*}
\Mor(\mathcal{G}, j_*(\mathcal{F}|_{\mathcal{C}/U}))
& =
\Mor(\mathcal{G}|_{\mathcal{C}/U}, \mathcal{F}|_{\mathcal{C}/U}) \\
& =
\Mor(j_!(\mathcal{G}|_{\mathcal{C}/U}), \mathcal{F}) \\
& =
\Mor(\mathcal{G} \times h_U^\#, \mathcal{F}) \\
& =
\Mor(\mathcal{G}, \SheafHom(h_U^\#, \mathcal{F}))
\end{align*}
et nous concluons par le lemme de Yoneda. Nous avons utilisé ici les
Lemmes \ref{sites-lemma-internal-hom-sheaf} et
\ref{sites-lemma-compute-j-shriek-restrict}.
\end{proof}

\noindent
Soit $\mathcal{C}$ un site.
Soit $\{U_i \to U\}_{i \in I}$ un recouvrement du site $\mathcal{C}$.
Pour tout $i \in I$, soit $\mathcal{F}_i$ un faisceau d'ensembles sur
$\mathcal{C}/U_i$.
Pour toute paire $i, j \in I$, soit
$$
\varphi_{ij} :
\mathcal{F}_i|_{\mathcal{C}/U_i \times_U U_j}
\longrightarrow
\mathcal{F}_j|_{\mathcal{C}/U_i \times_U U_j}
$$
un isomorphisme de faisceaux d'ensembles. Supposons en outre
que, pour tout triplet d'indices $i, j, k \in I$, le
diagramme suivant soit commutatif :
$$
\xymatrix{
\mathcal{F}_i|_{\mathcal{C}/U_i \times_U U_j \times_U U_k}
\ar[rr]_{\varphi_{ik}}
\ar[rd]_{\varphi_{ij}} & &
\mathcal{F}_k|_{\mathcal{C}/U_i \times_U U_j \times_U U_k} \\
&
\mathcal{F}_j|_{\mathcal{C}/U_i \times_U U_j \times_U U_k}
\ar[ru]_{\varphi_{jk}}
}
$$
Nous appellerons une telle famille de données
$(\mathcal{F}_i, \varphi_{ij})$
une {\it donnée de recollement de faisceaux d'ensembles relativement
au recouvrement $\{U_i \to U\}_{i \in I}$}.

\begin{lemma}
\label{sites-lemma-glue-sheaves}
Soit $\mathcal{C}$ un site.
Soit $\{U_i \to U\}_{i \in I}$ un recouvrement du site $\mathcal{C}$.
Pour toute donnée de recollement $(\mathcal{F}_i, \varphi_{ij})$
de faisceaux d'ensembles relativement au recouvrement $\{U_i \to U\}_{i \in I}$,
il existe un faisceau d'ensembles $\mathcal{F}$ sur $\mathcal{C}/U$
muni d'isomorphismes
$$
\varphi_i : \mathcal{F}|_{\mathcal{C}/U_i} \to \mathcal{F}_i
$$
tels que les diagrammes
$$
\xymatrix{
\mathcal{F}|_{\mathcal{C}/U_i \times_U U_j}
\ar[d]_{\text{id}} \ar[r]_{\varphi_i} &
\mathcal{F}_i|_{\mathcal{C}/U_i \times_U U_j} \ar[d]^{\varphi_{ij}} \\
\mathcal{F}|_{\mathcal{C}/U_i \times_U U_j} \ar[r]^{\varphi_j} &
\mathcal{F}_j|_{\mathcal{C}/U_i \times_U U_j}
}
$$
soient commutatifs.
\end{lemma}

\begin{proof}
Décrivons la construction du faisceau $\mathcal{F}$ sur
$\mathcal{C}/U$. Soit $a : V \to U$ un objet de $\mathcal{C}/U$.
Alors
$$
\mathcal{F}(V/U) = \{
(s_i)_{i \in I} \in \prod_{i \in I} \mathcal{F}_i(U_i \times_U V/U_i)
\mid
\varphi_{ij}(s_i|_{U_i \times_U U_j \times_U V})
=
s_j|_{U_i \times_U U_j \times_U V}
\}
$$
Nous omettons la construction des applications de restriction.
Nous omettons la vérification qu'il s'agit d'un faisceau.
Nous omettons la construction des isomorphismes $\varphi_i$,
ainsi que la preuve de la commutativité des diagrammes
du lemme.
\end{proof}

\noindent
Soit $\mathcal{C}$ un site.
Soit $\{U_i \to U\}_{i \in I}$ un recouvrement du site $\mathcal{C}$.
Soit $\mathcal{F}$ un faisceau sur $\mathcal{C}/U$.
À $\mathcal{F}$ est associée sa
{\it donnée de recollement canonique}, constituée des restrictions
$\mathcal{F}|_{\mathcal{C}/U_i}$ et des isomorphismes canoniques
$$
\left(\mathcal{F}|_{\mathcal{C}/U_i}\right)|_{\mathcal{C}/U_i \times_U U_j}
=
\left(\mathcal{F}|_{\mathcal{C}/U_j}\right)|_{\mathcal{C}/U_i \times_U U_j}
$$
provenant du fait que les composés des foncteurs
$\mathcal{C}/U_i \times_U U_j \to \mathcal{C}/U_i \to \mathcal{C}/U$
et
$\mathcal{C}/U_i \times_U U_j \to \mathcal{C}/U_j \to \mathcal{C}/U$
sont égaux.

\begin{lemma}
\label{sites-lemma-mapping-property-glue}
Soit $\mathcal{C}$ un site.
Soit $\{U_i \to U\}_{i \in I}$ un recouvrement du site $\mathcal{C}$.
La catégorie $\Sh(\mathcal{C}/U)$ est équivalente
à la catégorie des données de recollement par le foncteur qui associe
à $\mathcal{F}$ sur $\mathcal{C}/U$ sa donnée de recollement canonique.
\end{lemma}

\begin{proof}
Dans le
Lemme \ref{sites-lemma-glue-maps},
nous avons vu que le foncteur est pleinement fidèle, et dans le
Lemme \ref{sites-lemma-glue-sheaves},
nous avons démontré qu'il est essentiellement surjectif (en construisant
explicitement un foncteur quasi-inverse).
\end{proof}

\noindent
Soit $\mathcal{C}$ un site. Nous allons étudier
une version du recollement de faisceaux sur le site $\mathcal{C}$ tout entier.
Pour tout objet $U$ de $\mathcal{C}$, soit $\mathcal{F}_U$
un faisceau sur $\mathcal{C}/U$. Rappelons qu'il existe un foncteur
$j_f : \mathcal{C}/V \rightarrow \mathcal{C}/U$ associé à chaque
morphisme $f : V \rightarrow U$ de $\mathcal{C}$, défini par
$(a : W \to V) \mapsto (f \circ a : W \to U)$. Pour chaque tel $f$, soit
$$
c_f : j_f^{-1} \mathcal{F}_U \to \mathcal{F}_V
$$
un isomorphisme de faisceaux. Supposons que, pour deux flèches quelconques
$f : V \to U$ et $g : W \to V$ de $\mathcal{C}$, le composé
$c_g \circ j_{g}^{-1} c_f$ soit égal à $c_{f \circ g}$. Nous appellerons
une telle famille de données $(\mathcal{F}_{U}, c_f)$ une
{\it donnée de recollement absolue de faisceaux d'ensembles sur $\mathcal{C}$}.
Un morphisme de données de recollement absolues
$(\mathcal{F}_{U}, c_f) \to (\mathcal{G}_{U}, c'_f)$
est donné par une famille $(\varphi_U)$
de morphismes de faisceaux $\varphi_U : \mathcal{F}_U \to \mathcal{G}_U$,
telle que
$$
\xymatrix{
j_f^{-1} \mathcal{F}_U \ar[r]_{c_f} \ar[d]_{j_f^{-1} \varphi_U} &
\mathcal{F}_V \ar[d]^{\varphi_V} \\
j_f^{-1} \mathcal{G}_U \ar[r]^{c'_f} &
\mathcal{G}_V
}
$$
soit commutatif pour tout morphisme $f : V \to U$ de $\mathcal{C}$.

\medskip\noindent
À tout faisceau $\mathcal{F}$ sur $\mathcal{C}$ est associée sa
{\it donnée canonique de recollement absolue} $(\mathcal{F}|_{\mathcal{C}/U},c_f)$,
où les isomorphismes canoniques
$c_f : j_f^{-1} \mathcal{F}|_{\mathcal{C}/U} \to
\mathcal{F}|_{\mathcal{C}/V}$
pour $f : V \to U$ proviennent de la relation
$j_V = j_U \circ j_f$ comme dans le Lemme \ref{sites-lemma-relocalize}.
Tout morphisme $\varphi : \mathcal{F} \to \mathcal{G}$
de faisceaux sur $\mathcal{C}$ induit un morphisme
$(\varphi|_{\mathcal{C}/U})$ de données canoniques de recollement absolues.

\begin{lemma}
\label{sites-lemma-glue-sheaves-absolute}
Soit $\mathcal{C}$ un site. La catégorie
$\Sh(\mathcal{C})$ est équivalente à la catégorie
des données de recollement absolues par le foncteur qui
associe à $\mathcal{F}$ sur $\mathcal{C}$ sa
donnée canonique de recollement absolue.
\end{lemma}

\begin{proof}
Étant donnée une donnée de recollement absolue $(\mathcal{F}_U, c_f)$,
nous construisons un faisceau $\mathcal{F}$ sur $\mathcal{C}$ en
posant $\mathcal{F}(U) = \mathcal{F}_U(U)$, la restriction
le long de $f : V \to U$ étant donnée par le diagramme commutatif
$$
\xymatrix{
\mathcal{F}_U(U) \ar[r] \ar@{=}[d] &
\mathcal{F}_U(V) \ar[r]^{c_f} &
\mathcal{F}_V(V) \ar@{=}[d] \\
\mathcal{F}(U) \ar[rr] &  &
\mathcal{F}(V)
}
$$
La condition de compatibilité $c_g \circ j_{g}^{-1} c_f = c_{f \circ g}$
garantit que $\mathcal{F}$ est un préfaisceau, ainsi que le fait que les morphismes
$c_f : \mathcal{F}_U(V) \to \mathcal{F}(V)$ définissent un isomorphisme
$\mathcal{F}_U \to \mathcal{F}|_{\mathcal{C}/U}$ pour tout $U$.
Puisque chaque $\mathcal{F}_U$ est un faisceau, il en résulte que $\mathcal{F}$
est lui aussi un faisceau. Le foncteur $(\mathcal{F}_U, c_f) \mapsto \mathcal{F}$
que nous venons de construire est quasi-inverse du foncteur qui envoie un faisceau
sur $\mathcal{C}$ sur sa donnée de recollement canonique. Nous omettons les détails.
\end{proof}

\begin{remark}
\label{sites-remark-variant-glue-sheaves-absolute}
Il existe une variante du Lemme \ref{sites-lemma-glue-sheaves-absolute}
qui intervient en géométrie algébrique. Supposons en effet que $\mathcal{C}$
soit un site admettant tous les produits fibrés et que, pour tout $U \in \Ob(\mathcal{C})$,
on se donne une sous-catégorie pleine $U_\tau \subset \mathcal{C}/U$ possédant
les propriétés suivantes :
\begin{enumerate}
\item $U/U$ appartient à $U_\tau$,
\item si $V/U$ appartient à $U_\tau$ et si $\{V_j \to V\}$ est un recouvrement
du site $\mathcal{C}$, alors $V_j/U$ appartient à $U_\tau$, et
\item pour un morphisme $U' \to U$ de $\mathcal{C}$ et $V/U$ dans $U_\tau$,
le changement de base $V' = V \times_U U'$ appartient à $U'_\tau$.
\end{enumerate}
Dans ce cadre, $U_\tau$ est un site pour tout $U$ de $\mathcal{C}$,
et le foncteur de changement de base $U_\tau \to U'_\tau$ définit un morphisme
$f_\tau : U_\tau \to U'_\tau$ de sites pour tout morphisme $f : U' \to U$
de $\mathcal{C}$. L'énoncé de recollement obtenu est alors le suivant :
Un faisceau $\mathcal{F}$ sur $\mathcal{C}$ est donné par les données suivantes :
\begin{enumerate}
\item pour tout $U \in \Ob(\mathcal{C})$, un faisceau
$\mathcal{F}_U$ sur $U_\tau$,
\item pour tout $f : U' \to U$ de $\mathcal{C}$, un morphisme
$c_f : f_\tau^{-1}\mathcal{F}_U \to \mathcal{F}_{U'}$.
\end{enumerate}
Ces données sont soumises aux conditions suivantes :
\begin{enumerate}
\item[(a)] étant donnés $f : U' \to U$ et $g : U'' \to U'$ dans
$\mathcal{C}$, le composé
$c_g \circ g_\tau^{-1}c_f$ est égal à $c_{f \circ g}$, et
\item[(b)] si $f : U' \to U$ appartient à $U_\tau$, alors $c_f$ est un isomorphisme.
\end{enumerate}
Si nous en avons un jour besoin, nous en donnerons ici un énoncé précis et une preuve.
(Remarquons que ce résultat diffère légèrement des énoncés
ci-dessus, puisque nous ne demandons pas que tous les morphismes $c_f$ soient des isomorphismes !)
\end{remark}



\section{Compléments sur la localisation}
\label{sites-section-more-localize}

\noindent
Dans cette section, nous démontrons quelques lemmes de localisation en
imposant des hypothèses supplémentaires sur le site ou sur l'objet
en lequel nous localisons.

\begin{lemma}
\label{sites-lemma-describe-j-shriek-good-site}
Soit $\mathcal{C}$ un site.
Soit $U \in \Ob(\mathcal{C})$.
Si la topologie de $\mathcal{C}$ est sous-canonique, voir la
Définition \ref{sites-definition-weaker-than-canonical},
et si $\mathcal{G}$ est un faisceau sur $\mathcal{C}/U$, alors
$$
j_{U!}(\mathcal{G})(V)
=
\coprod\nolimits_{\varphi \in \Mor_\mathcal{C}(V, U)}
\mathcal{G}(V \xrightarrow{\varphi} U),
$$
autrement dit, le passage au faisceau associé n'est pas nécessaire dans le
Lemme \ref{sites-lemma-describe-j-shriek}.
\end{lemma}

\begin{proof}
Soit $\mathcal{V} = \{V_i \to V\}_{i \in I}$ un recouvrement de $V$
dans le site $\mathcal{C}$.
Nous allons vérifier directement la condition de faisceau pour le préfaisceau $\mathcal{H}$
du Lemme \ref{sites-lemma-describe-j-shriek}.
Soit $(s_i, \varphi_i)_{i \in I} \in \prod_i \mathcal{H}(V_i)$.
Cela signifie que $\varphi_i : V_i \to U$ est un morphisme de $\mathcal{C}$ et que
$s_i \in \mathcal{G}(V_i \xrightarrow{\varphi_i} U)$.
La restriction du couple $(s_i, \varphi_i)$ à
$V_i \times_V V_j$ est le couple
$(s_i|_{V_i \times_V V_j/U}, \text{pr}_1 \circ \varphi_i)$, et
de même, la restriction du couple $(s_j, \varphi_j)$ à
$V_i \times_V V_j$ est le couple
$(s_j|_{V_i \times_V V_j/U}, \text{pr}_2 \circ \varphi_j)$.
Par conséquent, si la famille $(s_i, \varphi_i)$ appartient à
$\check{H}^0(\mathcal{V}, \mathcal{H})$, alors
$\text{pr}_1 \circ \varphi_i = \text{pr}_2 \circ \varphi_j$.
Le fait que la topologie de $\mathcal{C}$ soit moins fine que la topologie canonique
implique alors qu'il existe un unique morphisme
$\varphi : V \to U$ tel que $\varphi_i$ soit le composé
de $V_i \to V$ avec $\varphi$. La condition de faisceau pour
$\mathcal{G}$ garantit alors que les sections $s_i$ se recollent en une unique
section $s \in \mathcal{G}(V \xrightarrow{\varphi} U)$.
Ainsi $(s, \varphi) \in \mathcal{H}(V)$, comme voulu.
\end{proof}

\begin{lemma}
\label{sites-lemma-localize-given-products}
Soit $\mathcal{C}$ un site.
Soit $U \in \Ob(\mathcal{C})$.
Supposons que $\mathcal{C}$ admette les produits de deux objets.
Alors
\begin{enumerate}
\item le foncteur $j_U$ admet un adjoint à droite continu,
à savoir le foncteur $v(X) = X \times U / U$,
\item le foncteur $v$ définit un morphisme de sites
$\mathcal{C}/U \to \mathcal{C}$ dont le morphisme de topos associé est égal à
$j_U : \Sh(\mathcal{C}/U) \to \Sh(\mathcal{C})$, et
\item nous avons $j_{U*}\mathcal{F}(X) = \mathcal{F}(X \times U/U)$.
\end{enumerate}
\end{lemma}

\begin{proof}
Dire que le foncteur $v$ est adjoint à droite de $j_U$ signifie que, pour $Y/U$ et $X$,
nous avons
$$
\Mor_\mathcal{C}(Y, X)
=
\Mor_{\mathcal{C}/U}(Y/U, X \times U/U)
$$
ce qui est clair. Pour vérifier que $v$ est continu, soit $\{X_i \to X\}$
un recouvrement du site $\mathcal{C}$. Par le troisième axiome d'un site
(Définition \ref{sites-definition-site}),
nous voyons que
$$
\{X_i \times_X (X \times U) \to X \times_X (X \times U)\}
=
\{X_i \times U \to X \times U\}
$$
est encore un recouvrement du site $\mathcal{C}$. Ainsi $v$ est continu. Les autres
assertions du lemme résultent des Lemmes \ref{sites-lemma-have-functor-other-way}
et \ref{sites-lemma-have-functor-other-way-morphism}.
\end{proof}

\begin{lemma}
\label{sites-lemma-relocalize-given-fibre-products}
Soit $\mathcal{C}$ un site. Soit $U \to V$ un morphisme de $\mathcal{C}$.
Supposons que $\mathcal{C}$ admette les produits fibrés. Soit $j$ comme dans le
Lemme \ref{sites-lemma-relocalize}. Alors
\begin{enumerate}
\item le foncteur $j : \mathcal{C}/U \to \mathcal{C}/V$
admet un adjoint à droite continu, à savoir le foncteur
$v : (X/V) \mapsto (X \times_V U/U)$,
\item le foncteur $v$ définit un morphisme de sites
$\mathcal{C}/U \to \mathcal{C}/V$ dont le morphisme de topos associé est égal à
$j$, et
\item nous avons $j_*\mathcal{F}(X/V) = \mathcal{F}(X \times_V U/U)$.
\end{enumerate}
\end{lemma}

\begin{proof}
Cela résulte du Lemme \ref{sites-lemma-localize-given-products}, puisque $j$ peut être vu
comme un foncteur de localisation par le Lemme \ref{sites-lemma-relocalize}.
\end{proof}

\noindent
Une propriété fondamentale d'une immersion ouverte est
que la restriction de l'image directe et la restriction
du prolongement par l'ensemble vide redonnent le faisceau initial.
Ce n'est pas toujours vrai pour les foncteurs associés à $j_U$
ci-dessus. Cela est vrai lorsque $U$ est un « sous-objet de l'objet final ».

\begin{lemma}
\label{sites-lemma-restrict-back}
Soit $\mathcal{C}$ un site.
Soit $U \in \Ob(\mathcal{C})$.
Supposons que tout $X$ de $\mathcal{C}$ admette au plus
un morphisme vers $U$. Soit $\mathcal{F}$ un faisceau sur $\mathcal{C}/U$.
Les morphismes canoniques $\mathcal{F} \to j_U^{-1}j_{U!}\mathcal{F}$
et $j_U^{-1}j_{U*}\mathcal{F} \to \mathcal{F}$ sont des
isomorphismes.
\end{lemma}

\begin{proof}
C'est un cas particulier du Lemme \ref{sites-lemma-back-and-forth},
car l'hypothèse sur $U$ équivaut à la pleine fidélité
du foncteur de localisation $\mathcal{C}/U \to \mathcal{C}$.
\end{proof}

\begin{lemma}
\label{sites-lemma-localize-cartesian-square}
Soit $\mathcal{C}$ un site. Soit
$$
\xymatrix{
U' \ar[d] \ar[r] & U \ar[d] \\
V' \ar[r] & V
}
$$
un diagramme commutatif de $\mathcal{C}$. Les
morphismes du Lemme \ref{sites-lemma-relocalize}
fournissent les diagrammes commutatifs
$$
\vcenter{
\xymatrix{
\mathcal{C}/U' \ar[d]_{j_{U'/V'}} \ar[r]_{j_{U'/U}} &
\mathcal{C}/U \ar[d]^{j_{U/V}} \\
\mathcal{C}/V' \ar[r]^{j_{V'/V}} & \mathcal{C}/V
}
}
\quad\text{et}\quad
\vcenter{
\xymatrix{
\Sh(\mathcal{C}/U') \ar[d]_{j_{U'/V'}} \ar[r]_{j_{U'/U}} &
\Sh(\mathcal{C}/U) \ar[d]^{j_{U/V}} \\
\Sh(\mathcal{C}/V') \ar[r]^{j_{V'/V}} &
\Sh(\mathcal{C}/V)
}
}
$$
de foncteurs continus et cocontinus, et de topos.
De plus, si le diagramme initial de $\mathcal{C}$ est cartésien,
alors nous avons
$j_{V'/V}^{-1} \circ j_{U/V, *} = j_{U'/V', *} \circ j_{U'/U}^{-1}$.
\end{lemma}

\begin{proof}
La commutativité du carré de gauche dans la première assertion du lemme
résulte immédiatement des définitions. Elle implique la commutativité
du diagramme de topos par le Lemme \ref{sites-lemma-composition-cocontinuous}.
Supposons le diagramme cartésien.
Par l'unicité des foncteurs adjoints, montrer
$j_{V'/V}^{-1} \circ j_{U/V, *} = j_{U'/V', *} \circ j_{U'/U}^{-1}$
équivaut à montrer
$j_{U/V}^{-1} \circ j_{V'/V!} = j_{U'/U!} \circ j_{U'/V'}^{-1}$.
Par les identifications du Lemme \ref{sites-lemma-essential-image-j-shriek},
nous pouvons identifier notre diagramme de topos à
$$
\xymatrix{
\Sh(\mathcal{C})/h_{U'}^\# \ar[d] \ar[r] &
\Sh(\mathcal{C})/h_U^\# \ar[d] \\
\Sh(\mathcal{C})/h_{V'}^\# \ar[r] &
\Sh(\mathcal{C})/h_V^\#
}
$$
et nous savons interpréter les foncteurs $j^{-1}$ et $j_!$
grâce au Lemme \ref{sites-lemma-relocalize-explicit}. Nous devons donc montrer,
étant donné $\mathcal{F} \to h_{V'}^\#$, que
$$
\mathcal{F} \times_{h_{V'}^\#} h_{U'}^\# =
\mathcal{F} \times_{h_V^\#} h_U^\#
$$
en tant que faisceaux munis d'un morphisme vers $h_U^\#$.
C'est vrai parce que $h_{U'} = h_{V'} \times_{h_V} h_U$
et donc aussi
$$
h_{U'}^\# = h_{V'}^\# \times_{h_V^\#} h_U^\#
$$
puisque le passage au faisceau associé est exact.
\end{proof}











\section{Localisation et morphismes}
\label{sites-section-localize-morphisms}

\noindent
Le lemme suivant est important pour comprendre la relation
entre la localisation et les morphismes de sites et de topos.

\begin{lemma}
\label{sites-lemma-localize-morphism}
Soit $f : \mathcal{C} \to \mathcal{D}$ un morphisme de sites
correspondant au foncteur continu $u : \mathcal{D} \to \mathcal{C}$.
Soit $V \in \Ob(\mathcal{D})$ et posons $U = u(V)$.
Alors le foncteur $u' : \mathcal{D}/V \to \mathcal{C}/U$,
$V'/V \mapsto u(V')/U$, détermine un morphisme de sites
$f' : \mathcal{C}/U \to \mathcal{D}/V$.
Le morphisme $f'$ s'insère dans un diagramme commutatif de topos
$$
\xymatrix{
\Sh(\mathcal{C}/U) \ar[r]_{j_U} \ar[d]_{f'} &
\Sh(\mathcal{C}) \ar[d]^f \\
\Sh(\mathcal{D}/V) \ar[r]^{j_V} &
\Sh(\mathcal{D}).
}
$$
En utilisant les identifications
$\Sh(\mathcal{C}/U) = \Sh(\mathcal{C})/h_U^\#$ et
$\Sh(\mathcal{D}/V) = \Sh(\mathcal{D})/h_V^\#$ du
Lemme \ref{sites-lemma-essential-image-j-shriek},
le foncteur $(f')^{-1}$ est décrit par la règle
$$
(f')^{-1}(\mathcal{H} \xrightarrow{\varphi} h_V^\#)
=
(f^{-1}\mathcal{H} \xrightarrow{f^{-1}\varphi} h_U^\#).
$$
Enfin, nous avons $f'_*j_U^{-1} = j_V^{-1}f_*$.
\end{lemma}

\begin{proof}
Il est clair que $u'$ est continu, et nous obtenons donc les foncteurs
$f'_* = (u')^s = (u')^p$ (voir les
sections \ref{sites-section-functoriality-PSh}
et \ref{sites-section-continuous-functors})
et un adjoint $(f')^{-1} = (u')_s = ((u')_p\ )^\#$. L'assertion
$f'_*j_U^{-1} = j_V^{-1}f_*$ résulte de
$$
(j_V^{-1}f_*\mathcal{F})(V'/V)
= f_*\mathcal{F}(V') = \mathcal{F}(u(V'))
= (j_U^{-1}\mathcal{F})(u(V')/U)
= (f'_*j_U^{-1}\mathcal{F})(V'/V)
$$
ce qui vaut même pour les préfaisceaux. Ce qui n'est pas clair a priori, c'est
que $(f')^{-1}$ est exact, que le diagramme commute et que
la description de $(f')^{-1}$ est valable.

\medskip\noindent
Soit $\mathcal{H}$ un faisceau sur $\mathcal{D}/V$.
Calculons $j_{U!}(f')^{-1}\mathcal{H}$. Nous avons
\begin{eqnarray*}
j_{U!}(f')^{-1}\mathcal{H}
& =
((j_U)_p(u'_p\mathcal{H})^\#)^\# \\
& =
((j_U)_pu'_p\mathcal{H})^\# \\
& =
(u_p(j_V)_p\mathcal{H})^\# \\
& =
f^{-1}j_{V!}\mathcal{H}
\end{eqnarray*}
La première égalité résulte du développement des définitions.
La deuxième résulte du Lemme \ref{sites-lemma-technical-up}.
La troisième vient de ce que $u \circ j_V = j_U \circ u'$.
La quatrième résulte à nouveau du Lemme \ref{sites-lemma-technical-up}.
Toutes les égalités ci-dessus sont des isomorphismes de foncteurs,
et nous pouvons donc les interpréter en disant que le diagramme
suivant de catégories et de foncteurs est commutatif :
$$
\xymatrix{
\Sh(\mathcal{C}/U) \ar[r]_{j_{U!}} &
\Sh(\mathcal{C})/h_U^\# \ar[r] &
\Sh(\mathcal{C}) \\
\Sh(\mathcal{D}/V) \ar[r]^{j_{V!}} \ar[u]^{(f')^{-1}} &
\Sh(\mathcal{D})/h_V^\# \ar[r] \ar[u]^{f^{-1}} &
\Sh(\mathcal{D}) \ar[u]^{f^{-1}}
}
$$
La flèche médiane a un sens puisque $f^{-1}h_V^\# = (h_{u(V)})^\# = h_U^\#$, voir le
Lemme \ref{sites-lemma-pullback-representable-sheaf}.
Ceci démontre en particulier la description de $(f')^{-1}$ donnée
dans l'énoncé du lemme.
Puisque, d'après le
Lemme \ref{sites-lemma-essential-image-j-shriek},
les flèches horizontales de gauche sont des équivalences,
et puisque $f^{-1}$ est exact par hypothèse, nous concluons que
$(f')^{-1} = u'_s$ est exact. En effet, comme c'est un adjoint à gauche,
il est déjà exact à droite
(Catégories, Lemme \ref{categories-lemma-adjoint-exact}).
Il nous suffit donc de montrer qu'il
transforme un objet final en un objet final et commute
aux produits fibrés
(Catégories, Lemme \ref{categories-lemma-characterize-left-exact}).
Ces deux propriétés sont claires pour le foncteur induit
$f^{-1} : \Sh(\mathcal{D})/h_V^\# \to \Sh(\mathcal{C})/h_U^\#$.
Ceci démontre que $f'$ est un morphisme de sites.

\medskip\noindent
Il nous reste à vérifier que $(f')^{-1}j_V^{-1} = j_U^{-1}f^{-1}$.
Pour le voir, combinons la formule ci-dessus avec la description
du Lemme \ref{sites-lemma-compute-j-shriek-restrict}. Nous obtenons ainsi,
pour tout faisceau $\mathcal{G}$ sur $\mathcal{D}$,
$$
j_{U!}(f')^{-1}j_V^{-1}\mathcal{G}
=
f^{-1}j_{V!}j_V^{-1}\mathcal{G}
=
f^{-1}(\mathcal{G} \times h_V^\#)
=
f^{-1}\mathcal{G} \times h_U^\#
=
j_{U!}j_U^{-1}f^{-1}\mathcal{G}.
$$
Puisque le foncteur $j_{U!}$ induit une équivalence
$\Sh(\mathcal{C}/U) \to \Sh(\mathcal{C})/h_U^\#$,
nous concluons.
\end{proof}

\noindent
Le lemme suivant est un cas particulier du
Lemme \ref{sites-lemma-localize-morphism}
ci-dessus, plus général.

\begin{lemma}
\label{sites-lemma-localize-morphism-strong}
Soient $\mathcal{C}$ et $\mathcal{D}$ des sites.
Soit $u : \mathcal{D} \to \mathcal{C}$ un foncteur.
Soit $V \in \Ob(\mathcal{D})$. Posons $U = u(V)$.
Supposons que
\begin{enumerate}
\item $\mathcal{C}$ et $\mathcal{D}$ admettent
toutes les limites finies,
\item $u$ soit continu, et
\item $u$ commute aux limites finies.
\end{enumerate}
Il existe un diagramme commutatif de morphismes de sites
$$
\xymatrix{
\mathcal{C}/U \ar[r]_{j_U} \ar[d]_{f'} & \mathcal{C} \ar[d]^f \\
\mathcal{D}/V \ar[r]^{j_V} & \mathcal{D}
}
$$
où la flèche verticale de droite correspond à $u$,
la flèche verticale de gauche correspond au
foncteur $u' : \mathcal{D}/V \to \mathcal{C}/U$, $V'/V \mapsto u(V')/u(V)$,
et les flèches horizontales correspondent aux foncteurs
$\mathcal{C} \to \mathcal{C}/U$, $X \mapsto X \times U$,
et $\mathcal{D} \to \mathcal{D}/V$, $Y \mapsto Y \times V$,
comme dans le Lemme \ref{sites-lemma-localize-given-products}.
De plus, le diagramme associé de morphismes de topos est
égal au diagramme du
Lemme \ref{sites-lemma-localize-morphism}.
En particulier, nous avons $f'_*j_U^{-1} = j_V^{-1}f_*$.
\end{lemma}

\begin{proof}
Notons que $u$ satisfait aux hypothèses de la
Proposition \ref{sites-proposition-get-morphism} et induit donc
un morphisme de sites $f : \mathcal{C} \to \mathcal{D}$ par cette proposition.
Il est clair que $u$ induit un foncteur $u'$ comme indiqué.
Il est clair que ce foncteur satisfait également aux hypothèses de la
Proposition \ref{sites-proposition-get-morphism}.
Nous obtenons donc un morphisme de sites $f' : \mathcal{C}/U \to \mathcal{D}/V$.
Le diagramme commute par notre définition de la composition des morphismes de
sites (voir la Définition \ref{sites-definition-composition-morphisms-sites})
et par l'égalité
$$
u(Y \times V) = u(Y) \times u(V) = u(Y) \times U
$$
qui montre que le diagramme de catégories et de foncteurs opposé
à celui du lemme est commutatif.
\end{proof}

\noindent
À ce stade, nous pouvons localiser un site, nous savons le relocaliser,
et nous pouvons localiser un morphisme de sites en un objet du site d'arrivée.
En combinant ces opérations, nous obtenons le type de diagramme suivant.

\begin{lemma}
\label{sites-lemma-relocalize-morphism}
Soit $f : \mathcal{C} \to \mathcal{D}$ un morphisme de sites
correspondant au foncteur continu $u : \mathcal{D} \to \mathcal{C}$.
Soient $V \in \Ob(\mathcal{D})$, $U \in \Ob(\mathcal{C})$
et $c : U \to u(V)$ un morphisme de $\mathcal{C}$.
Il existe un diagramme commutatif de topos
$$
\xymatrix{
\Sh(\mathcal{C}/U) \ar[r]_{j_U} \ar[d]_{f_c} &
\Sh(\mathcal{C}) \ar[d]^f \\
\Sh(\mathcal{D}/V) \ar[r]^{j_V} &
\Sh(\mathcal{D}).
}
$$
Nous avons $f_c = f' \circ j_{U/u(V)}$, où
$f' : \Sh(\mathcal{C}/u(V)) \to \Sh(\mathcal{D}/V)$
est celui du
Lemme \ref{sites-lemma-localize-morphism},
et où
$j_{U/u(V)} : \Sh(\mathcal{C}/U) \to \Sh(\mathcal{C}/u(V))$
est celui du
Lemme \ref{sites-lemma-relocalize}.
En utilisant les identifications
$\Sh(\mathcal{C}/U) = \Sh(\mathcal{C})/h_U^\#$ et
$\Sh(\mathcal{D}/V) = \Sh(\mathcal{D})/h_V^\#$ du
Lemme \ref{sites-lemma-essential-image-j-shriek},
le foncteur $(f_c)^{-1}$ est décrit par la règle
$$
(f_c)^{-1}(\mathcal{H} \xrightarrow{\varphi} h_V^\#)
=
(f^{-1}\mathcal{H} \times_{f^{-1}\varphi, h_{u(V)}^\#, c} h_U^\#
\rightarrow h_U^\#).
$$
Enfin, étant donnés des morphismes $b : V' \to V$, $a : U' \to U$ et
$c' : U' \to u(V')$ tels que
$$
\xymatrix{
U' \ar[r]_-{c'} \ar[d]_a & u(V') \ar[d]^{u(b)} \\
U \ar[r]^-c & u(V)
}
$$
commute, le diagramme
$$
\xymatrix{
\Sh(\mathcal{C}/U') \ar[r]_{j_{U'/U}} \ar[d]_{f_{c'}} &
\Sh(\mathcal{C}/U) \ar[d]^{f_c} \\
\Sh(\mathcal{D}/V') \ar[r]^{j_{V'/V}} &
\Sh(\mathcal{D}/V).
}
$$
est commutatif.
\end{lemma}

\begin{proof}
Ce lemme se démontre de lui-même et constitue plutôt une synthèse des faits connus
à ce stade du développement de la théorie. Par exemple, la commutativité
du premier carré résulte de la commutativité du
Diagramme (\ref{sites-equation-relocalize})
et de celle du diagramme du
Lemme \ref{sites-lemma-localize-morphism}.
La description de $f_c^{-1}$ résulte de la combinaison du
Lemme \ref{sites-lemma-relocalize-explicit}
avec le
Lemme \ref{sites-lemma-localize-morphism}.
La commutativité du dernier carré résulte alors de
l'égalité
$$
f^{-1}\mathcal{H} \times_{h_{u(V)}^\#, c} h_U^\# \times_{h_U^\#} h_{U'}^\#
=
f^{-1}(\mathcal{H} \times_{h_V^\#} h_{V'}^\#)
\times_{h_{u(V'), c'}^\#} h_{U'}^\#
$$
qui est formelle puisque $f^{-1}h_V^\# = h_{u(V)}^\#$ et
$f^{-1}h_{V'}^\# = h_{u(V')}^\#$, voir le
Lemme \ref{sites-lemma-pullback-representable-sheaf}.
\end{proof}

\noindent
Dans le lemme suivant, nous rencontrons une autre forme de fonctorialité de la
localisation, lorsque le morphisme de topos provient d'un
foncteur cocontinu. Ce type de diagramme diffère
du diagramme du
Lemme \ref{sites-lemma-localize-morphism},
et, en particulier, l'égalité $f'_*j_U^{-1} = j_V^{-1}f_*$ vue dans le
Lemme \ref{sites-lemma-localize-morphism}
n'est en général pas valable dans la situation du lemme suivant.

\begin{lemma}
\label{sites-lemma-localize-cocontinuous}
Soient $\mathcal{C}$ et $\mathcal{D}$ des sites.
Soit $u : \mathcal{C} \to \mathcal{D}$ un foncteur cocontinu.
Soit $U$ un objet de $\mathcal{C}$, et posons $V = u(U)$.
Nous avons un diagramme commutatif
$$
\xymatrix{
\mathcal{C}/U \ar[r]_{j_U} \ar[d]_{u'} & \mathcal{C} \ar[d]^u \\
\mathcal{D}/V \ar[r]^-{j_V} & \mathcal{D}
}
$$
où la flèche verticale de gauche est
$u' : \mathcal{C}/U \to \mathcal{D}/V$, $U'/U \mapsto V'/V$.
Alors $u'$ est également cocontinu et nous obtenons un diagramme commutatif de topos
$$
\xymatrix{
\Sh(\mathcal{C}/U) \ar[r]_{j_U} \ar[d]_{f'} &
\Sh(\mathcal{C}) \ar[d]^f \\
\Sh(\mathcal{D}/V) \ar[r]^-{j_V} &
\Sh(\mathcal{D})
}
$$
où $f$ (resp.\ $f'$) correspond à $u$ (resp.\ $u'$).
\end{lemma}

\begin{proof}
La commutativité du premier diagramme est claire.
Elle implique celle du second diagramme pourvu que nous
montrions que $u'$ est cocontinu.

\medskip\noindent
Soit $U'/U$ un objet de $\mathcal{C}/U$.
Soit $\{V_j/V \to u(U')/V\}_{j \in J}$ un recouvrement de $u(U')/V$
dans $\mathcal{D}/V$. Puisque $u$ est cocontinu, il existe un
recouvrement $\{U_i' \to U'\}_{i \in I}$ tel que la famille
$\{u(U_i') \to u(U')\}$ raffine le recouvrement
$\{V_j \to u(U')\}$ dans $\mathcal{D}$. Autrement dit, il existe
une application $\alpha : I \to J$ et des morphismes
$\phi_i : u(U_i') \to V_{\alpha(i)}$ au-dessus de $U'$.
Considérons $U_i'$ comme un objet au-dessus de
$U$ par la composée $U'_i \to U' \to U$. Alors
$\{U'_i/U \to U'/U\}$ est un recouvrement de $\mathcal{C}/U$ tel que
$\{u(U_i')/V \to u(U')/V\}$ raffine $\{V_j/V \to u(U')/V\}$
(en utilisant les mêmes $\alpha$ et $\phi_i$). Ainsi,
$u' : \mathcal{C}/U \to \mathcal{D}/V$ est cocontinu.
\end{proof}

\begin{lemma}
\label{sites-lemma-localize-cocontinuous-downstairs}
Soient $\mathcal{C}$ et $\mathcal{D}$ des sites.
Soit $u : \mathcal{C} \to \mathcal{D}$ un foncteur cocontinu.
Soit $V$ un objet de $\mathcal{D}$. Soit
${}^u_V\mathcal{I}$ la catégorie introduite dans la
section \ref{sites-section-more-functoriality-PSh}.
Nous avons un diagramme commutatif
$$
\vcenter{
\xymatrix{
\,_V^u\mathcal{I} \ar[r]_j \ar[d]_{u'} &
\mathcal{C} \ar[d]^u \\
\mathcal{D}/V \ar[r]^-{j_V} &
\mathcal{D}
}
}
\quad\text{où}\quad
\begin{matrix}
j : (U, \psi) \mapsto U \\
u' : (U, \psi) \mapsto (\psi : u(U) \to V)
\end{matrix}
$$
Déclarons qu'une famille de morphismes $\{(U_i, \psi_i) \to (U, \psi)\}$
de ${}^u_V\mathcal{I}$ est un recouvrement si et seulement si
$\{U_i \to U\}$ est un recouvrement de $\mathcal{C}$.
Alors
\begin{enumerate}
\item ${}^u_V\mathcal{I}$ est un site,
\item $j$ est continu et cocontinu,
\item $u'$ est cocontinu,
\item nous obtenons un diagramme commutatif de topos
$$
\xymatrix{
\Sh({}^u_V\mathcal{I}) \ar[r]_j \ar[d]_{f'} &
\Sh(\mathcal{C}) \ar[d]^f \\
\Sh(\mathcal{D}/V) \ar[r]^-{j_V} &
\Sh(\mathcal{D})
}
$$
où $f$ (resp.\ $f'$) correspond à $u$ (resp.\ $u'$), et
\item nous avons $f'_*j^{-1} = j_V^{-1}f_*$.
\end{enumerate}
\end{lemma}

\begin{proof}
Les parties (1), (2), (3) et (4) sont des conséquences immédiates des
définitions et du fait que le foncteur $j$ commute aux produits fibrés.
Nous omettons les détails. Pour voir (5), rappelons que $f_*$ est donné par
${}_su = {}_pu$. Ainsi, la valeur de $j_V^{-1}f_*\mathcal{F}$ en
$V'/V$ est la valeur de ${}_pu\mathcal{F}$ en $V'$, laquelle est la
limite des valeurs de $\mathcal{F}$ sur la catégorie
${}^u_{V'}\mathcal{I}$. Il existe manifestement une équivalence
de catégories
$$
{}^u_{V'}\mathcal{I} \to {}^{u'}_{V'/V}\mathcal{I}
$$
Puisque la valeur de $f'_*j^{-1}\mathcal{F}$ en $V'/V$ est
donnée par la limite des valeurs de $j^{-1}\mathcal{F}$
sur la catégorie ${}^{u'}_{V'/V}\mathcal{I}$, et puisque
les valeurs de $j^{-1}\mathcal{F}$ sur les objets de
${}^u_V\mathcal{I}$ sont simplement les valeurs de $\mathcal{F}$
(d'après le Lemme \ref{sites-lemma-when-shriek}, puisque $j$ est continu et
cocontinu),
nous voyons que (5) est vraie.
\end{proof}

\noindent
Les deux résultats suivants sont d'une nature légèrement différente.

\begin{lemma}
\label{sites-lemma-special-square-cocontinuous}
Supposons donnés des sites $\mathcal{C}', \mathcal{C}, \mathcal{D}', \mathcal{D}$
et des foncteurs
$$
\xymatrix{
\mathcal{C}' \ar[r]_{v'} \ar[d]_{u'} &
\mathcal{C} \ar[d]^u \\
\mathcal{D}' \ar[r]^v &
\mathcal{D}
}
$$
Supposons que
\begin{enumerate}
\item $u$, $u'$, $v$ et $v'$ soient cocontinus et donnent naissance aux morphismes de
topos $f$, $f'$, $g$ et $g'$ par le Lemme \ref{sites-lemma-cocontinuous-morphism-topoi},
\item $v \circ u' = u \circ v'$,
\item $v$ et $v'$ soient continus aussi bien que cocontinus, et
\item pour tout objet $V'$ de $\mathcal{D}'$, le foncteur
${}^{u'}_{V'}\mathcal{I} \to {}^{\ \ \ u}_{v(V')}\mathcal{I}$
donné par $v$ soit cofinal.
\end{enumerate}
Alors $f'_* \circ (g')^{-1} = g^{-1} \circ f_*$ et
$g'_! \circ (f')^{-1} = f^{-1} \circ g_!$.
\end{lemma}

\begin{proof}
Les catégories ${}^{u'}_{V'}\mathcal{I}$ et
${}^{\ \ \ u}_{v(V')}\mathcal{I}$ sont définies dans la
section \ref{sites-section-more-functoriality-PSh}.
Le foncteur de la condition (4) envoie l'objet
$\psi : u'(U') \to V'$ de ${}^{u'}_{V'}\mathcal{I}$ sur l'objet
$v(\psi) : uv'(U') = vu'(U') \to v(V')$ de ${}^{\ \ \ u}_{v(V')}\mathcal{I}$.
Rappelons que $g^{-1}$ est donné par $v^p$ (Lemme \ref{sites-lemma-when-shriek}) et
que $f_*$ est donné par ${}_su = {}_pu$. La valeur de
$g^{-1}f_*\mathcal{F}$ en $V'$ est donc la valeur de ${}_pu\mathcal{F}$
en $v(V')$, c'est-à-dire la limite
$$
\lim_{u(U) \to v(V') \in \Ob({}^{\ \ \ u}_{v(V')}\mathcal{I}^{opp})}
\mathcal{F}(U)
$$
Par le même raisonnement, la valeur de $f'_*(g')^{-1}\mathcal{F}$
en $V'$ est donnée par la limite
$$
\lim_{u'(U') \to V' \in \Ob({}^{u'}_{V'}\mathcal{I}^{opp})} \mathcal{F}(v'(U'))
$$
Ainsi, l'hypothèse (4) et Catégories, Lemme \ref{categories-lemma-initial},
montrent que ces limites coïncident, ce qui démontre la première égalité du
lemme. La seconde égalité découle de la première par unicité des
adjoints.
\end{proof}

\begin{lemma}
\label{sites-lemma-special-square-continuous}
Supposons donnés des sites $\mathcal{C}', \mathcal{C}, \mathcal{D}', \mathcal{D}$
et des foncteurs
$$
\xymatrix{
\mathcal{C}' \ar[r]_{v'} &
\mathcal{C} \\
\mathcal{D}' \ar[r]^v \ar[u]^{u'} &
\mathcal{D} \ar[u]_u
}
$$
Avec les notations des sections \ref{sites-section-morphism-sites}
et \ref{sites-section-cocontinuous-morphism-topoi}, supposons que
\begin{enumerate}
\item $u$ et $u'$ soient continus et donnent naissance aux morphismes de
sites $f$ et $f'$,
\item $v$ et $v'$ soient cocontinus et donnent naissance aux morphismes
de topos $g$ et $g'$,
\item $u \circ v = v' \circ u'$, et
\item $v$ et $v'$ soient continus aussi bien que cocontinus.
\end{enumerate}
Alors\footnote{Dans cette généralité,
nous ne savons pas si $f \circ g'$ est égal à $g \circ f'$
comme morphismes de topos (il existe une $2$-flèche canonique du
premier vers le second, qui peut ne pas être un isomorphisme).}
$f'_* \circ (g')^{-1} = g^{-1} \circ f_*$ et
$g'_! \circ (f')^{-1} = f^{-1} \circ g_!$.
\end{lemma}

\begin{proof}
En effet, nous avons
$$
f'_*(g')^{-1}\mathcal{F} = (u')^p((v')^p\mathcal{F})^\# =
(u')^p(v')^p\mathcal{F}
$$
La première égalité résulte de la définition, et la
seconde du Lemme \ref{sites-lemma-when-shriek}. Nous avons
$$
g^{-1}f_*\mathcal{F} = (v^pu^p\mathcal{F})^\# =
((u')^p(v')^p\mathcal{F})^\# = (u')^p(v')^p\mathcal{F}
$$
La première égalité résulte de la définition, la deuxième de ce que
$u \circ v = v' \circ u'$, et la troisième de ce que nous avons déjà vu que
$(u')^p(v')^p\mathcal{F}$ est un faisceau. Cela démontre
$f'_* \circ (g')^{-1} = g^{-1} \circ f_*$ et l'égalité
$g'_! \circ (f')^{-1} = f^{-1} \circ g_!$ résulte de
l'unicité des adjoints à gauche.
\end{proof}
















\section{Morphismes de topos}
\label{sites-section-morphisms-topoi}

\noindent
Dans cette section, nous montrons que tout morphisme de topos est équivalent
à un morphisme de topos qui provient d'un morphisme de sites.
Comparer avec \cite[Expos\'e III, Proposition 4.9.4]{SGA4}.

\begin{lemma}
\label{sites-lemma-equivalence}
Soient $\mathcal{C}$, $\mathcal{D}$ des sites.
Soit $u : \mathcal{C} \to \mathcal{D}$ un foncteur.
Supposons que
\begin{enumerate}
\item $u$ soit cocontinu,
\item $u$ soit continu,
\item étant donnés $a, b : U' \to U$ dans $\mathcal{C}$ tels que
$u(a) = u(b)$, il existe un recouvrement $\{f_i : U'_i \to U'\}$
dans $\mathcal{C}$ tel que $a \circ f_i = b \circ f_i$,
\item étant donnés $U', U \in \Ob(\mathcal{C})$ et
un morphisme $c : u(U') \to u(U)$ dans $\mathcal{D}$, il existe
un recouvrement $\{f_i : U_i' \to U'\}$ dans $\mathcal{C}$
et des morphismes $c_i : U_i' \to U$ tels que $u(c_i) = c \circ u(f_i)$, et
\item pour tout $V \in \Ob(\mathcal{D})$, il existe un recouvrement
de $V$ dans $\mathcal{D}$ de la forme $\{u(U_i) \to V\}_{i \in I}$.
\end{enumerate}
Alors le morphisme de topos
$$
g : \Sh(\mathcal{C}) \longrightarrow \Sh(\mathcal{D})
$$
associé au foncteur cocontinu $u$ par le
Lemme \ref{sites-lemma-cocontinuous-morphism-topoi}
est une équivalence.
\end{lemma}

\begin{proof}
Supposons que $u$ satisfasse les propriétés (1) -- (5). Nous allons montrer que
les morphismes d'adjonction
$$
\mathcal{G} \longrightarrow g_*g^{-1}\mathcal{G}
\quad\text{et}\quad
g^{-1}g_*\mathcal{F} \longrightarrow \mathcal{F}
$$
sont des isomorphismes.

\medskip\noindent
Remarquons que le Lemme \ref{sites-lemma-when-shriek} s'applique et que nous avons
$g^{-1}\mathcal{G}(U) = \mathcal{G}(u(U))$ pour tout faisceau $\mathcal{G}$
sur $\mathcal{D}$. Ensuite, soit $\mathcal{F}$ un faisceau sur $\mathcal{C}$,
et soit $V$ un objet de $\mathcal{D}$. Par définition, nous avons
$g_*\mathcal{F}(V) = \lim_{u(U) \to V} \mathcal{F}(U)$.
Par conséquent
$$
g^{-1}g_*\mathcal{F}(U) = \lim_{U', u(U') \to u(U)} \mathcal{F}(U')
$$
où les morphismes $\psi : u(U') \to u(U)$ ne sont pas nécessairement de la forme
$u(\alpha)$. La catégorie de ces couples $(U', \psi)$ possède un objet final,
à savoir $(U, \text{id})$, qui donne le morphisme de
la limite vers $\mathcal{F}(U)$. Soit $(s_{(U', \psi)})$ un élément
de la limite. Nous voulons montrer que $s_{(U', \psi)}$ est uniquement déterminé
par la valeur $s_{(U, \text{id})} \in \mathcal{F}(U)$. D'après la propriété (4), pour
tout $(U', \psi)$, il existe un recouvrement $\{U'_i \to U'\}$ tel que les
composés $u(U'_i) \to u(U') \to u(U)$ soient de la forme $u(c_i)$
pour certains $c_i : U'_i \to U$ dans $\mathcal{C}$. Par conséquent
$$
s_{(U', \psi)}|_{U'_i} = c_i^*(s_{(U, \text{id})}).
$$
Comme $\mathcal{F}$ est un faisceau, il s'ensuit que $s_{(U', \psi)}$
est bien déterminé par $s_{(U, \text{id})}$. Cela démontre l'unicité.
Pour l'existence, supposons donnés
$s \in \mathcal{F}(U)$, $\psi : u(U') \to u(U)$, $\{f_i : U_i' \to U'\}$
et $c_i : U_i' \to U$ tels que $\psi \circ u(f_i) = u(c_i)$ comme ci-dessus.
Nous affirmons qu'il existe un (unique) élément
$s_{(U', \psi)} \in \mathcal{F}(U')$ tel que
$$
s_{(U', \psi)}|_{U'_i} = c_i^*(s).
$$
En effet, a priori, il n'est pas clair que les éléments
$c_i^*(s)|_{U_i' \times_{U'} U_j'}$
et $c_j^*(s)|_{U_i' \times_{U'} U_j'}$ coïncident, car
le diagramme
$$
\xymatrix{
U_i' \times_{U'} U_j' \ar[r]_-{\text{pr}_2} \ar[d]_{\text{pr}_1} &
U_j' \ar[d]^{c_j} \\
U_i' \ar[r]^{c_i} & U}
$$
ne commute pas nécessairement. Mais la condition (3) du lemme garantit qu'il
existe des recouvrements
$\{f_{ijk} : U'_{ijk} \to U_i' \times_{U'} U_j'\}_{k \in K_{ij}}$ tels que
$c_i \circ \text{pr}_1 \circ f_{ijk} = c_j \circ \text{pr}_2 \circ f_{ijk}$.
Par conséquent
$$
f_{ijk}^* \left(c_i^*s|_{U_i' \times_{U'} U_j'}\right)
=
f_{ijk}^* \left(c_j^*s|_{U_i' \times_{U'} U_j'}\right)
$$
Par conséquent $c_i^*(s)|_{U_i' \times_{U'} U_j'} = c_j^*(s)|_{U_i' \times_{U'} U_j'}$
par la condition de faisceau pour $\mathcal{F}$, et l'existence de
$s_{(U', \psi)}$ en résulte également par la condition de faisceau pour $\mathcal{F}$. L'unicité
garantit que la famille $(s_{(U', \psi)})$ ainsi obtenue est un élément
de la limite avec $s_{(U, \psi)} = s$. Cela démontre
que $g^{-1}g_*\mathcal{F} \to \mathcal{F}$ est un isomorphisme.

\medskip\noindent
Soit $\mathcal{G}$ un faisceau sur $\mathcal{D}$. Soit $V$ un
objet de $\mathcal{D}$. Alors nous voyons que
$$
g_*g^{-1}\mathcal{G}(V) =
\lim_{U, \psi : u(U) \to V} \mathcal{G}(u(U))
$$
D'après le paragraphe précédent, nous voyons que la valeur du faisceau
$g_*g^{-1}\mathcal{G}$ sur un objet $V$ de la forme $V = u(U)$
est égale à $\mathcal{G}(u(U))$. (Formellement, cela découle de
$g^{-1}g_*g^{-1} \cong g^{-1}$ et de la description
de $g^{-1}$ donnée au début de la démonstration ; informellement, il suffit de
comparer les limites ici et ci-dessus.)
Ainsi, le morphisme d'adjonction $\mathcal{G} \to g_*g^{-1}\mathcal{G}$
a la propriété d'être une bijection sur les sections au-dessus de tout objet de
la forme $u(U)$. Comme, par l'axiome (5), il
existe un recouvrement de $V$ par des objets de la forme $u(U)$, nous voyons
facilement que le morphisme d'adjonction est un isomorphisme.
\end{proof}

\noindent
Il sera commode de donner un nom aux foncteurs cocontinus comme dans le
Lemme \ref{sites-lemma-equivalence}.

\begin{definition}
\label{sites-definition-special-cocontinuous-functor}
Soient $\mathcal{C}$, $\mathcal{D}$ des sites.
Un {\it foncteur cocontinu spécial $u$ de $\mathcal{C}$ vers $\mathcal{D}$}
est un foncteur cocontinu $u : \mathcal{C} \to \mathcal{D}$ satisfaisant
aux hypothèses et aux conclusions du Lemme \ref{sites-lemma-equivalence}.
\end{definition}

\begin{lemma}
\label{sites-lemma-localize-special-cocontinuous}
Soient $\mathcal{C}$, $\mathcal{D}$ des sites.
Soit $u : \mathcal{C} \to \mathcal{D}$ un foncteur cocontinu spécial.
Pour tout objet $U$ de $\mathcal{C}$, nous avons un diagramme commutatif
$$
\xymatrix{
\mathcal{C}/U \ar[r]_{j_U} \ar[d] & \mathcal{C} \ar[d]^u \\
\mathcal{D}/u(U) \ar[r]^-{j_{u(U)}} & \mathcal{D}
}
$$
comme dans le Lemme \ref{sites-lemma-localize-cocontinuous}.
La flèche verticale de gauche est un foncteur cocontinu spécial.
Ainsi, dans le diagramme commutatif de topos
$$
\xymatrix{
\Sh(\mathcal{C}/U) \ar[r]_{j_U} \ar[d] &
\Sh(\mathcal{C}) \ar[d]^u \\
\Sh(\mathcal{D}/u(U)) \ar[r]^-{j_{u(U)}} &
\Sh(\mathcal{D})
}
$$
les flèches verticales sont des équivalences.
\end{lemma}

\begin{proof}
Nous avons établi l'existence et la commutativité des diagrammes dans le
Lemme \ref{sites-lemma-localize-cocontinuous}. Il reste à vérifier
les hypothèses (1) -- (5) du Lemme \ref{sites-lemma-equivalence} pour le
foncteur induit $u : \mathcal{C}/U \to \mathcal{D}/u(U)$.
Cette vérification est entièrement formelle.

\medskip\noindent
Propriété (1). C'est le Lemme \ref{sites-lemma-localize-cocontinuous}.

\medskip\noindent
Propriété (2). Soit $\{U_i'/U \to U'/U\}_{i \in I}$ un recouvrement
de $U'/U$ dans $\mathcal{C}/U$. Comme $u$ est continu, nous voyons que
$\{u(U_i')/u(U) \to u(U')/u(U)\}_{i \in I}$ est un recouvrement
de $u(U')/u(U)$ dans $\mathcal{D}/u(U)$. Ainsi, (2) est vérifiée
pour $u : \mathcal{C}/U \to \mathcal{D}/u(U)$.

\medskip\noindent
Propriété (3). Soient $a, b : U''/U \to U'/U$ dans $\mathcal{C}/U$
des morphismes tels que $u(a) = u(b)$ dans $\mathcal{D}/u(U)$.
Comme $u$ satisfait (3), nous voyons qu'il existe un recouvrement
$\{f_i : U''_i \to U''\}$ dans $\mathcal{C}$ tel que
$a \circ f_i = b \circ f_i$. Cela donne un recouvrement
$\{f_i : U''_i/U \to U''/U\}$ dans $\mathcal{C}/U$ tel que
$a \circ f_i = b \circ f_i$. Ainsi, (3) est vérifiée
pour $u : \mathcal{C}/U \to \mathcal{D}/u(U)$.

\medskip\noindent
Propriété (4). Soient $U''/U, U'/U \in \Ob(\mathcal{C}/U)$ et
soit donné un morphisme $c : u(U'')/u(U) \to u(U')/u(U)$ dans $\mathcal{D}/u(U)$.
Comme $u$ satisfait la propriété (4), il existe
un recouvrement $\{f_i : U_i'' \to U''\}$ dans $\mathcal{C}$
et des morphismes $c_i : U_i'' \to U'$ tels que $u(c_i) = c \circ u(f_i)$.
Nous considérons $U_i''$ comme un objet au-dessus de $U$ par la composée
$U_i'' \to U'' \to U$.
Il n'est pas nécessairement vrai que $c_i$ soit un morphisme au-dessus de $U$ !
Mais puisque $u(c_i)$ est un morphisme au-dessus de $u(U)$, nous pouvons appliquer
la propriété (3) à $u$ et trouver des recouvrements $\{f_{ik} : U''_{ik} \to U''_i\}$
tels que $c_{ik} = c_i \circ f_{ik} : U''_{ik} \to U'$ soient des morphismes au-dessus de $U$.
Ainsi, $\{f_i \circ f_{ik} : U''_{ik}/U \to U''/U\}$ est un recouvrement
dans $\mathcal{C}/U$ tel que $u(c_{ik}) = c \circ u(f_{ik})$.
Ainsi, (4) est vérifiée
pour $u : \mathcal{C}/U \to \mathcal{D}/u(U)$.

\medskip\noindent
Propriété (5). Soit $h : V \to u(U)$ un objet de $\mathcal{D}/u(U)$.
Comme $u$ satisfait la propriété (5), il existe un recouvrement
$\{c_i : u(U_i) \to V\}$ dans $\mathcal{D}$. Par la propriété (4), nous pouvons trouver
des recouvrements $\{f_{ij} : U_{ij} \to U_i\}$ et des morphismes
$c_{ij} : U_{ij} \to U$ tels que $u(c_{ij}) = h \circ c_i \circ u(f_{ij})$.
Ainsi, $\{u(U_{ij})/u(U) \to V/u(U)\}$ est un recouvrement dans
$\mathcal{D}/u(U)$ de la forme voulue, et nous concluons que
(5) est vérifiée pour $u : \mathcal{C}/U \to \mathcal{D}/u(U)$.
\end{proof}

\begin{lemma}
\label{sites-lemma-special-equivalence}
Soit $\mathcal{C}$ un site. Soit
$\mathcal{C}' \subset \Sh(\mathcal{C})$
une sous-catégorie pleine (possédant un ensemble d'objets) telle que
\begin{enumerate}
\item $h_U^\# \in \Ob(\mathcal{C}')$ pour tout
$U \in \Ob(\mathcal{C})$, et
\item $\mathcal{C}'$ soit stable par produits fibrés dans
$\Sh(\mathcal{C})$.
\end{enumerate}
Déclarons qu'un recouvrement de $\mathcal{C}'$ est toute famille
$\{\mathcal{F}_i \to \mathcal{F}\}_{i \in I}$ de morphismes telle que
$\coprod_{i \in I} \mathcal{F}_i \to \mathcal{F}$ soit un morphisme surjectif
de faisceaux. Alors
\begin{enumerate}
\item $\mathcal{C}'$ est un site (après avoir
choisi un ensemble de recouvrements, voir Ensembles, Lemme \ref{sets-lemma-coverings-site}),
\item les préfaisceaux représentables sur $\mathcal{C}'$ sont des faisceaux
(c'est-à-dire que la topologie sur $\mathcal{C}'$ est sous-canonique, voir
Définition \ref{sites-definition-weaker-than-canonical}),
\item le foncteur $v : \mathcal{C} \to \mathcal{C}'$,
$U \mapsto h_U^\#$ est un foncteur cocontinu spécial, et induit donc une
équivalence $g : \Sh(\mathcal{C}) \to \Sh(\mathcal{C}')$,
\item pour tout $\mathcal{F} \in \Ob(\mathcal{C}')$, nous avons
$g^{-1}h_\mathcal{F} = \mathcal{F}$, et
\item pour tout $U \in \Ob(\mathcal{C})$, nous avons
$g_*h_U^\# = h_{v(U)} = h_{h_U^\#}$.
\end{enumerate}
\end{lemma}

\begin{proof}
Avertissement : certains des énoncés ci-dessus peuvent sembler un peu déroutants au premier abord ;
cela vient de ce que les objets de $\mathcal{C}'$ peuvent aussi être considérés comme des faisceaux sur
$\mathcal{C}$ ! Nous omettons de vérifier que les recouvrements de $\mathcal{C}'$
décrits dans le lemme satisfont les conditions de la
Définition \ref{sites-definition-site}.

\medskip\noindent
Supposons que $\{\mathcal{F}_i \to \mathcal{F}\}$ soit une famille surjective
de morphismes de faisceaux. Soit $\mathcal{G}$ un autre faisceau.
La partie (2) du lemme affirme que l'égalisateur de
$$
\xymatrix{
\Mor_{\Sh(\mathcal{C})}(
\coprod_{i \in I} \mathcal{F}_i, \mathcal{G})
\ar@<1ex>[r] \ar@<-1ex>[r]
&
\Mor_{\Sh(\mathcal{C})}(
\coprod_{(i_0, i_1) \in I \times I}
\mathcal{F}_{i_0} \times_\mathcal{F} \mathcal{F}_{i_1}, \mathcal{G})
}
$$
est $\Mor_{\Sh(\mathcal{C})}(\mathcal{F}, \mathcal{G}).$
C'est clair (utiliser par exemple le Lemme \ref{sites-lemma-coequalizer-surjection}).

\medskip\noindent
Pour démontrer (3), nous
devons vérifier les conditions (1) -- (5) du Lemme \ref{sites-lemma-equivalence}.
Le fait que $v$ soit cocontinu est
équivalent à la description des morphismes surjectifs de faisceaux dans le
Lemme \ref{sites-lemma-mono-epi-sheaves}.
Le foncteur $v$ est continu parce que
$U \mapsto h_U^\#$ commute aux produits fibrés,
et transforme les recouvrements en recouvrements (voir le
Lemme \ref{sites-lemma-sheafification-exact}, et le
Lemme \ref{sites-lemma-covering-surjective-after-sheafification}).
Les propriétés (3), (4) du Lemme \ref{sites-lemma-equivalence}
sont des énoncés portant sur les morphismes $f : h_{U'}^\# \to h_U^\#$.
Un tel morphisme revient à un élément de $h_U^\#(U')$.
Ainsi, (3) et (4) découlent immédiatement de la construction du faisceau associé.
La propriété (5) du Lemme \ref{sites-lemma-equivalence} est le
Lemme \ref{sites-lemma-sheaf-coequalizer-representable}.
Soit $g : \Sh(\mathcal{C}) \to \Sh(\mathcal{C}')$
l'équivalence de topos associée à $v$ par le Lemme \ref{sites-lemma-equivalence}.

\medskip\noindent
Soit $\mathcal{F}$ comme dans la partie (4) du lemme.
Pour tout $U \in \Ob(\mathcal{C})$, nous avons
$$
g^{-1}h_\mathcal{F}(U) = h_\mathcal{F}(v(U))
= \Mor_{\Sh(\mathcal{C})}(h_U^\#, \mathcal{F})
= \mathcal{F}(U)
$$
La première égalité résulte du
Lemme \ref{sites-lemma-when-shriek}.
La partie (4) est donc démontrée.

\medskip\noindent
Soit $\mathcal{F} \in \Ob(\mathcal{C}')$.
Soit $U \in \Ob(\mathcal{C})$.
Alors
\begin{align*}
g_*h_U^\#(\mathcal{F})
& =
\Mor_{\Sh(\mathcal{C}')}(h_\mathcal{F}, g_*h_U^\#) \\
& =
\Mor_{\Sh(\mathcal{C})}(g^{-1}h_\mathcal{F}, h_U^\#) \\
& =
\Mor_{\Sh(\mathcal{C})}(\mathcal{F}, h_U^\#) \\
& =
\Mor_{\mathcal{C}'}(\mathcal{F}, h_U^\#)
\end{align*}
comme voulu (où la troisième égalité a été démontrée ci-dessus).
\end{proof}

\noindent
Cela nous permet de réaliser tout topos sur un site possédant
toutes les limites finies.

\begin{lemma}
\label{sites-lemma-topos-good-site}
Soit $\Sh(\mathcal{C})$ un topos. Soit $\{\mathcal{F}_i\}_{i \in I}$
un ensemble de faisceaux sur $\mathcal{C}$. Il existe une équivalence de topos
$g : \Sh(\mathcal{C}) \to \Sh(\mathcal{C}')$ induite par un foncteur
cocontinu spécial $u : \mathcal{C} \to \mathcal{C}'$ de sites
telle que
\begin{enumerate}
\item $\mathcal{C}'$ possède une topologie sous-canonique,
\item une famille $\{V_j \to V\}$ de morphismes de $\mathcal{C}'$
est (combinatoirement équivalente à) un recouvrement de $\mathcal{C}'$
si et seulement si $\coprod h_{V_j} \to h_V$ est surjectif,
\item $\mathcal{C}'$ possède des produits fibrés et un objet final
(c'est-à-dire que $\mathcal{C}'$ possède toutes les limites finies),
\item tout sous-faisceau d'un faisceau représentable sur $\mathcal{C}'$
est représentable, et
\item chaque $g_*\mathcal{F}_i$ est un faisceau représentable.
\end{enumerate}
\end{lemma}

\begin{proof}
Considérons la sous-catégorie pleine
$\mathcal{C}_1 \subset \Sh(\mathcal{C})$ formée de tous les
$h_U^\#$ pour tout $U \in \Ob(\mathcal{C})$, des faisceaux donnés
$\mathcal{F}_i$ et du faisceau final $*$ (voir
Exemple \ref{sites-example-singleton-sheaf}). Nous allons définir par récurrence
des sous-catégories pleines
$$
\mathcal{C}_1 \subset \mathcal{C}_2 \subset \mathcal{C}_2 \subset \ldots
\subset \Sh(\mathcal{C})
$$
Plus précisément, $\mathcal{C}_n$ étant donnée, soit $\mathcal{C}_{n + 1}$ la sous-catégorie pleine
formée de tous les produits fibrés et sous-faisceaux d'objets
de $\mathcal{C}_n$. (Remarquons que $\mathcal{C}_{n + 1}$ possède un ensemble
d'objets.) Posons
$\mathcal{C}' = \bigcup_{n \geq 1} \mathcal{C}_n$.
Un recouvrement de $\mathcal{C}'$ est toute famille
$\{\mathcal{G}_j \to \mathcal{G}\}_{j \in J}$ de morphismes d'objets
de $\mathcal{C}'$ telle que
$\coprod \mathcal{G}_j \to \mathcal{G}$ soit surjectif
comme morphisme de faisceaux sur $\mathcal{C}$.
Le foncteur $v : \mathcal{C} \to \mathcal{C'}$ est donné par
$U \mapsto h_U^\#$. Appliquons le
Lemme \ref{sites-lemma-special-equivalence}.
\end{proof}

\noindent
Voici le but de la présente section.

\begin{lemma}
\label{sites-lemma-morphism-topoi-comes-from-morphism-sites}
\begin{reference}
Cet énoncé est étroitement lié à
\cite[Proposition 4.9.4. Expos\'e IV]{SGA4}.
Pour obtenir le résultat complet, il faut aussi utiliser
\cite[Remarque 4.7.4, Expos\'e IV]{SGA4}.
\end{reference}
Soient $\mathcal{C}$, $\mathcal{D}$ des sites.
Soit $f : \Sh(\mathcal{C}) \to \Sh(\mathcal{D})$ un
morphisme de topos.
Alors il existe un site $\mathcal{C}'$ et un diagramme de foncteurs
$$
\xymatrix{
\mathcal{C} \ar[r]_v & \mathcal{C}' & \mathcal{D} \ar[l]^u
}
$$
tels que
\begin{enumerate}
\item le foncteur $v$ soit un foncteur cocontinu spécial,
\item le foncteur $u$ commute aux produits fibrés, soit
continu et définisse un morphisme de sites
$\mathcal{C}' \to \mathcal{D}$, et
\item le morphisme de topos $f$ coïncide avec le composé
des morphismes de topos
$$
\Sh(\mathcal{C}) \longrightarrow
\Sh(\mathcal{C}') \longrightarrow
\Sh(\mathcal{D})
$$
où la première flèche provient de $v$ par le Lemme \ref{sites-lemma-equivalence}
et la seconde de $u$ par le Lemme \ref{sites-lemma-morphism-sites-topoi}.
\end{enumerate}
\end{lemma}

\begin{proof}
Considérons la sous-catégorie pleine
$\mathcal{C}_1 \subset \Sh(\mathcal{C})$ formée de tous les
$h_U^\#$ et de tous les $f^{-1}h_V^\#$ pour tout
$U \in \Ob(\mathcal{C})$ et tout $V \in \Ob(\mathcal{D})$.
Soit $\mathcal{C}_{n + 1}$ la sous-catégorie pleine formée de tous les
produits fibrés d'objets de $\mathcal{C}_n$. Posons
$\mathcal{C}' = \bigcup_{n \geq 1} \mathcal{C}_n$.
Un recouvrement de $\mathcal{C}'$ est toute famille
$\{\mathcal{F}_i \to \mathcal{F}\}_{i \in I}$ telle que
$\coprod_{i \in I} \mathcal{F}_i \to \mathcal{F}$ soit surjectif
comme morphisme de faisceaux sur $\mathcal{C}$.
Le foncteur $v : \mathcal{C} \to \mathcal{C'}$ est donné par
$U \mapsto h_U^\#$.
Le foncteur $u : \mathcal{D} \to \mathcal{C'}$ est donné par
$V \mapsto f^{-1}h_V^\#$.

\medskip\noindent
La partie (1) résulte du Lemme \ref{sites-lemma-special-equivalence}.

\medskip\noindent
Démonstration de (2) et (3) du lemme. Le foncteur $u$ commute aux produits
fibrés, puisque $V \mapsto h_V^\#$ et $f^{-1}$ le font tous deux. En outre,
comme $f^{-1}$ est exact et commute aux limites inductives arbitraires, nous voyons
qu'il transforme un recouvrement en une famille surjective de morphismes de
faisceaux. Ainsi, $u$ est continu. Pour voir qu'il définit un morphisme de
sites, il reste à vérifier que $u_s$ est exact. À cette fin,
nous allons montrer que $g^{-1} \circ u_s = f^{-1}$. En effet, puisque $g^{-1}$
est une équivalence et que $f^{-1}$ est exact, nous pourrons conclure.
Comme $g^{-1}$ est adjoint à $g_*$, que $u_s$ est adjoint à
$u^s$, et que $f^{-1}$ est adjoint à $f_*$, il suffit aussi de démontrer que
$u^s \circ g_* = f_*$. Soit $\mathcal{F}$ un faisceau sur $\mathcal{C}$
et soit $V$ un objet de $\mathcal{D}$. Alors
\begin{align*}
(u^s g_{\ast}\mathcal{F})(V)
& =
(g_{\ast}\mathcal{F})(f^{-1}h_V^\#) \\
& =
\Mor_{Sh(\mathcal{C}')}(h_{f^{-1}h_V^\#},g_{\ast}\mathcal{F}) \\
& =
\Mor_{Sh(\mathcal{C})}(g^{-1}h_{f^{-1}h_V^\#},\mathcal{F}) \\
& =
\Mor_{Sh(\mathcal{C})}(f^{-1}h_V^\#,\mathcal{F}) \\
& =
\Mor_{Sh(\mathcal{D})}(h_V^\#,f_{\ast}\mathcal{F}) \\
& =
f_{\ast}\mathcal{F}(V)
\end{align*}
La première égalité vient de ce que $u^s = u^p$.
La deuxième égalité est le lemme de Yoneda.
La troisième égalité résulte de l'adjonction entre $g^{-1}$ et $g_*$.
La quatrième égalité résulte du Lemme \ref{sites-lemma-special-equivalence} (4).
La cinquième égalité résulte de l'adjonction entre $f^{-1}$ et $f_*$.
La sixième égalité est le lemme de Yoneda.
Ainsi, $u^s g_*\mathcal{F} = f_*\mathcal{F}$ et
cela achève la démonstration du lemme.
\end{proof}

\begin{remark}
\label{sites-remark-morphism-topoi-comes-from-morphism-sites}
Notations et hypothèses
comme dans le Lemme \ref{sites-lemma-morphism-topoi-comes-from-morphism-sites}.
Si le site $\mathcal{D}$ possède un objet final et des produits fibrés,
alors le foncteur $u : \mathcal{D} \to \mathcal{C}'$ satisfait
toutes les hypothèses de la Proposition \ref{sites-proposition-get-morphism}.
En effet, outre les propriétés mentionnées dans le lemme, $u$
transforme aussi l'objet final de $\mathcal{D}$ en l'objet final
de $\mathcal{C}'$. Cela ressort clairement de la construction de $u$.
Ainsi, si nous appliquons d'abord le
Lemme \ref{sites-lemma-topos-good-site}
à $\mathcal{D}$
puis le
Lemme \ref{sites-lemma-morphism-topoi-comes-from-morphism-sites}
au morphisme de topos qui en résulte
$\Sh(\mathcal{C}) \to \Sh(\mathcal{D}')$,
nous obtenons l'énoncé suivant :
Tout morphisme de topos
$f : \Sh(\mathcal{C}) \to \Sh(\mathcal{D})$
s'insère dans un diagramme commutatif
$$
\xymatrix{
\Sh(\mathcal{C}) \ar[d]_g \ar[r]_f &
\Sh(\mathcal{D}) \ar[d]^e \\
\Sh(\mathcal{C}') \ar[r]^{f'} &
\Sh(\mathcal{D}')
}
$$
où les propriétés suivantes sont satisfaites :
\begin{enumerate}
\item les morphismes $e$ et $g$ sont des équivalences données par des
foncteurs cocontinus spéciaux $\mathcal{C} \to \mathcal{C}'$ et
$\mathcal{D} \to \mathcal{D}'$,
\item les sites $\mathcal{C}'$ et $\mathcal{D}'$ possèdent des produits fibrés et un objet
final et sont munis de topologies sous-canoniques,
\item le morphisme $f' : \mathcal{C}' \to \mathcal{D}'$ provient d'un
morphisme de sites correspondant à un foncteur
$u : \mathcal{D}' \to \mathcal{C}'$ auquel la
Proposition \ref{sites-proposition-get-morphism}
s'applique, et
\item étant donné un ensemble quelconque de faisceaux $\mathcal{F}_i$ (resp.\ $\mathcal{G}_j$)
sur $\mathcal{C}$ (resp.\ $\mathcal{D}$), nous pouvons supposer que chacun d'eux est
un faisceau représentable sur $\mathcal{C}'$ (resp.\ $\mathcal{D}'$).
\end{enumerate}
Il est souvent utile de remplacer $\mathcal{C}$ et $\mathcal{D}$ par
$\mathcal{C}'$ et $\mathcal{D}'$.
\end{remark}

\begin{remark}
\label{sites-remark-equivalence-topoi-comes-from-morphism-sites}
Notations et hypothèses
comme dans le Lemme \ref{sites-lemma-morphism-topoi-comes-from-morphism-sites}.
Supposons en outre que le morphisme de topos d'origine
$\Sh(\mathcal{C}) \to \Sh(\mathcal{D})$ soit une équivalence.
Alors la construction de la démonstration du
Lemme \ref{sites-lemma-morphism-topoi-comes-from-morphism-sites}
donne deux foncteurs
$$
\mathcal{C} \rightarrow \mathcal{C}' \leftarrow \mathcal{D}
$$
qui sont tous deux des foncteurs cocontinus spéciaux.
Ainsi, dans ce cas, nous pouvons effectivement
factoriser le morphisme de topos comme un composé
$$
\Sh(\mathcal{C}) \rightarrow
\Sh(\mathcal{C}') =
\Sh(\mathcal{D}') \leftarrow
\Sh(\mathcal{D})
$$
comme dans la Remarque \ref{sites-remark-morphism-topoi-comes-from-morphism-sites}, mais
avec le morphisme du milieu égal à l'identité.
\end{remark}








\section{Localisation des topos}
\label{sites-section-localize-topoi}

\noindent
Nous reprenons certains résultats sur la localisation dans le cadre, en
apparence plus général, des topos. En réalité, ce cadre n'est pas plus général,
car nous pouvons toujours agrandir les sites sous-jacents pour supposer que
nous localisons en des objets du site.

\begin{lemma}
\label{sites-lemma-localize-topos}
Soit $\mathcal{C}$ un site.
Soit $\mathcal{F}$ un faisceau sur $\mathcal{C}$.
Alors la catégorie $\Sh(\mathcal{C})/\mathcal{F}$ est un topos. Il existe un
morphisme canonique de topos
$$
j_\mathcal{F} :
\Sh(\mathcal{C})/\mathcal{F}
\longrightarrow
\Sh(\mathcal{C})
$$
qui est une localisation au sens de la section \ref{sites-section-localize} et tel
que
\begin{enumerate}
\item le foncteur $j_\mathcal{F}^{-1}$ soit le foncteur
$\mathcal{H} \mapsto \mathcal{H} \times \mathcal{F}/\mathcal{F}$, et
\item le foncteur $j_{\mathcal{F}!}$ soit le foncteur d'oubli
$\mathcal{G}/\mathcal{F} \mapsto \mathcal{G}$.
\end{enumerate}
\end{lemma}

\begin{proof}
Appliquons le Lemme \ref{sites-lemma-topos-good-site}.
Cela signifie que nous pouvons supposer que $\mathcal{C}$ est un site muni
d'une topologie sous-canonique et que $\mathcal{F} = h_U = h_U^\#$ pour un
certain $U \in \Ob(\mathcal{C})$.
Les résultats de la section \ref{sites-section-localize} s'appliquent donc. En
particulier, il existe une équivalence
$\Sh(\mathcal{C}/U) = \Sh(\mathcal{C})/h_U^\#$ telle que le composé
$$
\Sh(\mathcal{C}/U)
\to
\Sh(\mathcal{C})/h_U^\#
\to \Sh(\mathcal{C})
$$
soit égal à $j_{U!}$ ; voir le
Lemme \ref{sites-lemma-essential-image-j-shriek}.
Notons $a : \Sh(\mathcal{C})/h_U^\# \to \Sh(\mathcal{C}/U)$ le foncteur
inverse, de sorte que $j_{\mathcal{F}!} = j_{U!} \circ a$,
$j_\mathcal{F}^{-1} = a^{-1} \circ j_U^{-1}$ et
$j_{\mathcal{F}, *} = j_{U, *} \circ a$. La description de
$j_{\mathcal{F}!}$ résulte de ce qui précède. Celle de
$j_\mathcal{F}^{-1}$ résulte du
Lemme \ref{sites-lemma-compute-j-shriek-restrict}.
\end{proof}

\begin{lemma}
\label{sites-lemma-localize-pushforward}
Dans la situation du Lemme \ref{sites-lemma-localize-topos}, le foncteur
$j_{\mathcal{F}, *}$ associe à
$\varphi : \mathcal{G} \to \mathcal{F}$ le faisceau
$$
U
\longmapsto
\{\alpha : \mathcal{F}|_U \to \mathcal{G}|_U
\text{ telle que } \alpha \text{ est un inverse à droite de }\varphi|_U \}.
$$
\end{lemma}

\begin{proof}
Pour tout $\varphi : \mathcal{G} \to \mathcal{F}$, convenons de noter
$\mathcal{G}/\mathcal{F}$ l'objet correspondant de
$\Sh(\mathcal{C})/\mathcal{F}$. On a
$$
(j_{\mathcal{F}, *}(\mathcal{G}/\mathcal{F}))(U) =
\Mor_{\Sh(\mathcal{C})}(h_U^\#, j_{\mathcal{F}, *}(\mathcal{G}/\mathcal{F})) =
\Mor_{\Sh(\mathcal{C})/\mathcal{F}}(j_{\mathcal{F}}^{-1}h_U^{\#},
(\mathcal{G}/\mathcal{F})).
$$
D'après le Lemme \ref{sites-lemma-localize-topos}, cet ensemble est la fibre
au-dessus de l'élément $h_U^\# \times \mathcal{F} \to \mathcal{F}$ pour
l'application d'ensembles
$$
\Mor_{\Sh(\mathcal{C})}(h_U^\# \times \mathcal{F}, \mathcal{G})
\xrightarrow{\varphi \circ}
\Mor_{\Sh(\mathcal{C})}(h_U^\# \times \mathcal{F}, \mathcal{F}).
$$
Par l'adjonction du Lemme \ref{sites-lemma-internal-hom-sheaf}, on a
\begin{align*}
\Mor_{\Sh(\mathcal{C})}(h_U^{\#}\times\mathcal{F}, \mathcal{G})
& =
\Mor_{\Sh(\mathcal{C})}(h_U^{\#},\SheafHom(\mathcal{F}, \mathcal{G})) \\
& =
\Mor_{\Sh(\mathcal{C}/U)}(\mathcal{F}|_{\mathcal{C}/U},
\mathcal{G}|_{\mathcal{C}/U}), \\
\Mor_{\Sh(\mathcal{C})}(h_U^{\#} \times \mathcal{F}, \mathcal{F})
& =
\Mor_{\Sh(\mathcal{C})}(h_U^{\#},\SheafHom(\mathcal{F},\mathcal{F})) \\
& =
\Mor_{\Sh(\mathcal{C}/U)}(\mathcal{F}|_{\mathcal{C}/U},
\mathcal{F}|_{\mathcal{C}/U}),
\end{align*}
et, sous cette adjonction, l'application $\varphi\circ$ devient
$$
\Mor_{\Sh(\mathcal{C}/U)}(\mathcal{F}|_{\mathcal{C}/U},
\mathcal{G}|_{\mathcal{C}/U})
\longrightarrow
\Mor_{\Sh(\mathcal{C}/U)}(\mathcal{F}|_{\mathcal{C}/U},
\mathcal{F}|_{\mathcal{C}/U}),\quad
\psi \longmapsto \varphi|_{\mathcal{C}/U} \circ \psi,
$$
tandis que l'élément $h_U^\# \times \mathcal{F} \to \mathcal{F}$ devient
$\text{id}_{\mathcal{F}|_{\mathcal{C}/U}}$.
Par conséquent, $(j_{\mathcal{F}, *}\mathcal{G}/\mathcal{F})(U)$ est
isomorphe à la fibre au-dessus de $\text{id}_{\mathcal{F}|_{\mathcal{C}/U}}$
pour l'application
$$
\Mor_{\Sh(\mathcal{C}/U)}(\mathcal{F}|_{\mathcal{C}/U},
\mathcal{G}|_{\mathcal{C}/U})
\xrightarrow{\varphi|_{\mathcal{C}/U}\circ}
\Mor_{\Sh(\mathcal{C}/U)}(\mathcal{F}|_{\mathcal{C}/U},
\mathcal{F}|_{\mathcal{C}/U}),
$$
laquelle est $\{\alpha : \mathcal{F}|_U \to \mathcal{G}|_U
\text{ telle que } \alpha \text{ est un inverse à droite de }\varphi|_U \}$,
comme voulu.
\end{proof}

\begin{lemma}
\label{sites-lemma-localize-topos-site}
Soit $\mathcal{C}$ un site. Soit $\mathcal{F}$ un faisceau sur $\mathcal{C}$.
Soit $\mathcal{C}/\mathcal{F}$ la catégorie des couples $(U, s)$, où
$U \in \Ob(\mathcal{C})$ et $s \in \mathcal{F}(U)$. Déclarons qu'un
recouvrement de $\mathcal{C}/\mathcal{F}$ est une famille
$\{(U_i, s_i) \to (U, s)\}$ telle que $\{U_i \to U\}$ soit un recouvrement de
$\mathcal{C}$. Alors $j : \mathcal{C}/\mathcal{F} \to \mathcal{C}$ est un
foncteur de sites continu et cocontinu, qui induit un morphisme de topos
$j : \Sh(\mathcal{C}/\mathcal{F}) \to \Sh(\mathcal{C})$. En fait, il existe
une équivalence $\Sh(\mathcal{C}/\mathcal{F}) =
\Sh(\mathcal{C})/\mathcal{F}$ qui transforme $j$ en $j_\mathcal{F}$.
\end{lemma}

\begin{proof}
Nous omettons de vérifier que $\mathcal{C}/\mathcal{F}$ est un site et que
$j$ est continu et cocontinu. D'après le Lemme \ref{sites-lemma-when-shriek}, il
existe un morphisme de topos $j$ comme indiqué, avec
$j^{-1}\mathcal{G}(U, s) = \mathcal{G}(U)$, et il existe un adjoint à gauche
$j_!$ de $j^{-1}$. Un morphisme $\varphi : * \to j^{-1}\mathcal{G}$ sur
$\mathcal{C}/\mathcal{F}$ revient à se donner une règle qui associe à chaque
couple $(U, s)$ une section $\varphi(s) \in \mathcal{G}(U)$ compatible avec les
applications de restriction. Cela revient donc à se donner un morphisme
$\varphi : \mathcal{F} \to \mathcal{G}$ sur $\mathcal{C}$.
On en conclut que $j_!* = \mathcal{F}$. En particulier, pour tout
$\mathcal{H} \in \Sh(\mathcal{C}/\mathcal{F})$, il existe un morphisme
canonique
$$
j_!\mathcal{H} \to j_!* = \mathcal{F}
$$
c'est-à-dire que l'on obtient un foncteur
$j'_! : \Sh(\mathcal{C}/\mathcal{F}) \to
\Sh(\mathcal{C})/\mathcal{F}$.
Un inverse de ce foncteur est la règle qui associe à un objet
$\varphi : \mathcal{G} \to \mathcal{F}$ de
$\Sh(\mathcal{C})/\mathcal{F}$ le faisceau
$$
a(\mathcal{G}/\mathcal{F}) :
(U, s) \longmapsto \{t \in \mathcal{G}(U) \mid \varphi(t) = s\}
$$
Nous omettons de vérifier que $a(\mathcal{G}/\mathcal{F})$ est un faisceau et
que $a$ est inverse de $j'_!$.
\end{proof}

\begin{definition}
\label{sites-definition-localize-topos}
Soit $\mathcal{C}$ un site.
Soit $\mathcal{F}$ un faisceau sur $\mathcal{C}$.
\begin{enumerate}
\item Le topos $\Sh(\mathcal{C})/\mathcal{F}$ est appelé la
{\it localisation du topos $\Sh(\mathcal{C})$ en $\mathcal{F}$}.
\item Le morphisme de topos
$j_\mathcal{F} :
\Sh(\mathcal{C})/\mathcal{F}
\to
\Sh(\mathcal{C})$ du
Lemme \ref{sites-lemma-localize-topos}
est appelé le {\it morphisme de localisation}.
\end{enumerate}
\end{definition}

\noindent
Nous allons montrer que, chaque fois que le faisceau $\mathcal{F}$ est égal à
$h_U^\#$ pour un certain objet $U$ du site, la localisation du topos est égale
à la catégorie des faisceaux sur la localisation du site en $U$. Nous allons
en outre vérifier que toutes les fonctorialités sont compatibles avec cette
identification.

\begin{lemma}
\label{sites-lemma-localize-compare}
Soit $\mathcal{C}$ un site. Soit $\mathcal{F} = h_U^\#$ pour un certain objet
$U$ de $\mathcal{C}$. Alors le morphisme
$j_\mathcal{F} : \Sh(\mathcal{C})/\mathcal{F}
\to \Sh(\mathcal{C})$ construit dans le
Lemme \ref{sites-lemma-localize-topos}
coïncide avec le morphisme de topos
$j_U : \Sh(\mathcal{C}/U) \to \Sh(\mathcal{C})$
construit dans la section \ref{sites-section-localize} par l'identification
$\Sh(\mathcal{C}/U) = \Sh(\mathcal{C})/h_U^\#$ du
Lemme \ref{sites-lemma-essential-image-j-shriek}.
\end{lemma}

\begin{proof}
Nous avons vu dans le Lemme \ref{sites-lemma-essential-image-j-shriek} que le
composé
$\Sh(\mathcal{C}/U) \to \Sh(\mathcal{C})/h_U^\#
\to \Sh(\mathcal{C})$
est $j_{U!}$. Le foncteur
$\Sh(\mathcal{C})/h_U^\# \to \Sh(\mathcal{C})$
est $j_{\mathcal{F}!}$ d'après le Lemme \ref{sites-lemma-localize-topos}.
Ainsi, $j_{\mathcal{F}!} = j_{U!}$ par l'identification.
Donc $j_\mathcal{F}^{-1} = j_U^{-1}$ (par adjonction), puis
$j_{\mathcal{F}, *} = j_{U, *}$ (encore par adjonction).
\end{proof}

\begin{lemma}
\label{sites-lemma-relocalize-topos}
Soit $\mathcal{C}$ un site.
Si $s : \mathcal{G} \to \mathcal{F}$ est un morphisme de faisceaux sur
$\mathcal{C}$, alors il existe un diagramme naturel commutatif de morphismes de
topos
$$
\xymatrix{
\Sh(\mathcal{C})/\mathcal{G} \ar[rd]_{j_\mathcal{G}} \ar[rr]_j & &
\Sh(\mathcal{C})/\mathcal{F} \ar[ld]^{j_\mathcal{F}} \\
& \Sh(\mathcal{C}) &
}
$$
où $j = j_{\mathcal{G}/\mathcal{F}}$ est la localisation du topos
$\Sh(\mathcal{C})/\mathcal{F}$ en l'objet $\mathcal{G}/\mathcal{F}$. En
particulier, on a
$$
j^{-1}(\mathcal{H} \to \mathcal{F}) =
(\mathcal{H} \times_\mathcal{F} \mathcal{G} \to \mathcal{G})
$$
et
$$
j_!(\mathcal{E} \xrightarrow{e} \mathcal{G}) =
(\mathcal{E} \xrightarrow{s \circ e} \mathcal{F}).
$$
\end{lemma}

\begin{proof}
La description de $j^{-1}$ et de $j_!$ résulte de la description de ces
foncteurs dans le Lemme \ref{sites-lemma-localize-topos}.
L'égalité de foncteurs $j_{\mathcal{G}!} = j_{\mathcal{F}!} \circ j_!$ résulte
clairement de leur description comme foncteurs d'oubli. Par adjonction, on
obtient également les égalités
$j_\mathcal{G}^{-1} = j^{-1} \circ j_\mathcal{F}^{-1}$ et
$j_{\mathcal{G}, *} = j_{\mathcal{F}, *} \circ j_*$.
\end{proof}

\begin{lemma}
\label{sites-lemma-relocalize-compare}
Supposons que $\mathcal{C}$ et $s : \mathcal{G} \to \mathcal{F}$ soient comme
dans le Lemme \ref{sites-lemma-relocalize-topos}.
Si $\mathcal{G} = h_V^\#$, $\mathcal{F} = h_U^\#$ et si
$s : \mathcal{G} \to \mathcal{F}$ provient d'un morphisme $V \to U$ de
$\mathcal{C}$, alors le diagramme du Lemme \ref{sites-lemma-relocalize-topos}
s'identifie au diagramme (\ref{sites-equation-relocalize}) par les identifications
$\Sh(\mathcal{C}/V) = \Sh(\mathcal{C})/h_V^\#$ et
$\Sh(\mathcal{C}/U) = \Sh(\mathcal{C})/h_U^\#$ du
Lemme \ref{sites-lemma-essential-image-j-shriek}.
\end{lemma}

\begin{proof}
C'est vrai parce que les descriptions de $j^{-1}$ coïncident. Voir le
Lemme \ref{sites-lemma-relocalize-explicit} et le
Lemme \ref{sites-lemma-relocalize-topos}.
\end{proof}






\section{Localisation et morphismes de topos}
\label{sites-section-localize-morphisms-topoi}

\noindent
Cette section est l'analogue de la section \ref{sites-section-localize-morphisms}
pour les morphismes de topos.

\begin{lemma}
\label{sites-lemma-localize-morphism-topoi}
Soit $f : \Sh(\mathcal{C}) \to \Sh(\mathcal{D})$ un morphisme de topos. Soit
$\mathcal{G}$ un faisceau sur $\mathcal{D}$. Posons
$\mathcal{F} = f^{-1}\mathcal{G}$. Il existe alors un diagramme commutatif de
topos
$$
\xymatrix{
\Sh(\mathcal{C})/\mathcal{F} \ar[r]_{j_\mathcal{F}} \ar[d]_{f'} &
\Sh(\mathcal{C}) \ar[d]^f \\
\Sh(\mathcal{D})/\mathcal{G} \ar[r]^{j_\mathcal{G}} &
\Sh(\mathcal{D}).
}
$$
Le morphisme $f'$ est caractérisé par la propriété
$$
(f')^{-1}(\mathcal{H} \xrightarrow{\varphi} \mathcal{G})
=
(f^{-1}\mathcal{H} \xrightarrow{f^{-1}\varphi} \mathcal{F})
$$
et l'on a $f'_*j_\mathcal{F}^{-1} = j_\mathcal{G}^{-1}f_*$.
\end{lemma}

\begin{proof}
Comme l'énoncé porte sur des topos et ne fait pas référence aux sites
sous-jacents, nous pouvons changer ceux-ci à volonté. D'après la discussion de
la Remarque \ref{sites-remark-morphism-topoi-comes-from-morphism-sites}, nous pouvons
donc supposer que $f$ est donné par un foncteur continu
$u : \mathcal{D} \to \mathcal{C}$ satisfaisant aux hypothèses de la
Proposition \ref{sites-proposition-get-morphism}, entre des sites possédant toutes
les limites finies et des topologies sous-canoniques, et tel que
$\mathcal{G} = h_V$ pour un certain objet $V$ de $\mathcal{D}$. Alors
$\mathcal{F} = f^{-1}h_V = h_{u(V)}$ d'après le
Lemme \ref{sites-lemma-pullback-representable-sheaf}.
D'après le Lemme \ref{sites-lemma-localize-morphism}, on obtient un diagramme
commutatif de morphismes de topos
$$
\xymatrix{
\Sh(\mathcal{C}/U) \ar[r]_{j_U} \ar[d]_{f'} &
\Sh(\mathcal{C}) \ar[d]^f \\
\Sh(\mathcal{D}/V) \ar[r]^{j_V} &
\Sh(\mathcal{D}),
}
$$
et l'on a $f'_*j_U^{-1} = j_V^{-1}f_*$. D'après le
Lemme \ref{sites-lemma-localize-compare}, on peut identifier $j_\mathcal{F}$ à
$j_U$, et $j_\mathcal{G}$ à $j_V$. La description de $(f')^{-1}$ est donnée
dans le Lemme \ref{sites-lemma-localize-morphism}.
\end{proof}

\begin{lemma}
\label{sites-lemma-localize-morphism-compare}
Soit $f : \mathcal{C} \to \mathcal{D}$ un morphisme de sites donné par le
foncteur continu $u : \mathcal{D} \to \mathcal{C}$.
Soit $V$ un objet de $\mathcal{D}$. Posons $U = u(V)$.
Posons $\mathcal{G} = h_V^\#$ et
$\mathcal{F} = h_U^\# = f^{-1}h_V^\#$ (voir le
Lemme \ref{sites-lemma-pullback-representable-sheaf}).
Alors le diagramme de morphismes de topos du
Lemme \ref{sites-lemma-localize-morphism-topoi} coïncide avec le diagramme de
morphismes de topos du Lemme \ref{sites-lemma-localize-morphism}, par les
identifications $j_\mathcal{F}= j_U$ et $j_\mathcal{G} = j_V$ du
Lemme \ref{sites-lemma-localize-compare}.
\end{lemma}

\begin{proof}
Ce n'est pas tout à fait trivial, car le morphisme de sites donnant naissance
à $f$ choisi dans la démonstration du
Lemme \ref{sites-lemma-localize-morphism-topoi} peut différer du morphisme de sites
qui nous est donné dans le lemme. Mais, dans les deux cas, le foncteur
$(f')^{-1}$ est décrit par la même règle. Les deux foncteurs coïncident donc,
ainsi que les morphismes de topos associés. Nous omettons certains détails.
\end{proof}

\begin{lemma}
\label{sites-lemma-relocalize-morphism-topoi}
Soit $f : \Sh(\mathcal{C}) \to \Sh(\mathcal{D})$ un morphisme de topos.
Soient $\mathcal{G} \in \Sh(\mathcal{D})$,
$\mathcal{F} \in \Sh(\mathcal{C})$ et
$s : \mathcal{F} \to f^{-1}\mathcal{G}$ un morphisme de faisceaux.
Il existe un diagramme commutatif de topos
$$
\xymatrix{
\Sh(\mathcal{C})/\mathcal{F} \ar[r]_{j_\mathcal{F}} \ar[d]_{f_s} &
\Sh(\mathcal{C}) \ar[d]^f \\
\Sh(\mathcal{D})/\mathcal{G} \ar[r]^{j_\mathcal{G}} &
\Sh(\mathcal{D}).
}
$$
On a $f_s = f' \circ j_{\mathcal{F}/f^{-1}\mathcal{G}}$, où
$f' :
\Sh(\mathcal{C})/f^{-1}\mathcal{G}
\to
\Sh(\mathcal{D})/\mathcal{F}$
est comme dans le Lemme \ref{sites-lemma-localize-morphism-topoi}, et
$j_{\mathcal{F}/f^{-1}\mathcal{G}} :
\Sh(\mathcal{C})/\mathcal{F}
\to
\Sh(\mathcal{C})/f^{-1}\mathcal{G}$
est comme dans le Lemme \ref{sites-lemma-relocalize-topos}.
Le foncteur $(f_s)^{-1}$ est décrit par la règle
$$
(f_s)^{-1}(\mathcal{H} \xrightarrow{\varphi} \mathcal{G})
=
(f^{-1}\mathcal{H} \times_{f^{-1}\varphi, f^{-1}\mathcal{G}, s} \mathcal{F}
\rightarrow \mathcal{F}).
$$
Enfin, étant donnés des morphismes quelconques
$b : \mathcal{G}' \to \mathcal{G}$,
$a : \mathcal{F}' \to \mathcal{F}$ et
$s' : \mathcal{F}' \to f^{-1}\mathcal{G}'$ tels que
$$
\xymatrix{
\mathcal{F}' \ar[r]_-{s'} \ar[d]_a & f^{-1}\mathcal{G}' \ar[d]^{f^{-1}b} \\
\mathcal{F} \ar[r]^-s & f^{-1}\mathcal{G}
}
$$
soit commutatif, le diagramme
$$
\xymatrix{
\Sh(\mathcal{C})/\mathcal{F}'
\ar[r]_{j_{\mathcal{F}'/\mathcal{F}}} \ar[d]_{f_{s'}} &
\Sh(\mathcal{C})/\mathcal{F} \ar[d]^{f_s} \\
\Sh(\mathcal{D})/\mathcal{G}' \ar[r]^{j_{\mathcal{G}'/\mathcal{G}}} &
\Sh(\mathcal{D})/\mathcal{G}.
}
$$
est commutatif.
\end{lemma}

\begin{proof}
La commutativité du premier carré résulte de celle du diagramme du
Lemme \ref{sites-lemma-relocalize-topos} et de celle du diagramme du
Lemme \ref{sites-lemma-localize-morphism-topoi}. La description de $f_s^{-1}$
résulte de la combinaison de la description de $(f')^{-1}$ dans le
Lemme \ref{sites-lemma-localize-morphism-topoi} et de la description de
$(j_{\mathcal{F}/f^{-1}\mathcal{G}})^{-1}$ dans le
Lemme \ref{sites-lemma-relocalize-topos}. La commutativité du dernier carré résulte
alors de l'égalité
$$
f^{-1}\mathcal{H} \times_{f^{-1}\mathcal{G}, s} \mathcal{F}
\times_\mathcal{F} \mathcal{F}'
=
f^{-1}(\mathcal{H} \times_\mathcal{G} \mathcal{G}')
\times_{f^{-1}\mathcal{G}', s'} \mathcal{F}'
$$
qui est formelle.
\end{proof}

\begin{lemma}
\label{sites-lemma-relocalize-morphism-compare}
Soit $f : \mathcal{C} \to \mathcal{D}$ un morphisme de sites donné par le
foncteur continu $u : \mathcal{D} \to \mathcal{C}$.
Soit $V$ un objet de $\mathcal{D}$. Soit $c : U \to u(V)$ un morphisme.
Posons $\mathcal{G} = h_V^\#$ et
$\mathcal{F} = h_U^\# = f^{-1}h_V^\#$.
Soit $s : \mathcal{F} \to f^{-1}\mathcal{G}$ l'application induite par $c$.
Alors le diagramme de morphismes de topos du
Lemme \ref{sites-lemma-relocalize-morphism} coïncide avec le diagramme de morphismes
de topos du Lemme \ref{sites-lemma-relocalize-morphism-topoi} par les identifications
$j_\mathcal{F} = j_U$ et $j_\mathcal{G} = j_V$ du
Lemme \ref{sites-lemma-localize-compare}.
\end{lemma}

\begin{proof}
Cela résulte de la combinaison des Lemmes \ref{sites-lemma-relocalize-compare} et
\ref{sites-lemma-localize-morphism-compare}.
\end{proof}















\section{Points}
\label{sites-section-points}

\begin{definition}
\label{sites-definition-point-topos}
Soit $\mathcal{C}$ un site.
Un {\it point du topos $\Sh(\mathcal{C})$} est un morphisme de topos $p$ de
$\Sh(pt)$ vers $\Sh(\mathcal{C})$.
\end{definition}

\noindent
Nous définirons un point d'un site au moyen d'un foncteur
$u : \mathcal{C} \to \textit{Ensembles}$.
Nous verrons plus loin que $u$ définit un morphisme de sites qui donne
naissance à un point du topos associé à $\mathcal{C}$ ; voir le
Lemme \ref{sites-lemma-site-point-morphism}.

\medskip\noindent
Soit $\mathcal{C}$ un site. Soit $p = u$ un foncteur
$u : \mathcal{C} \to \textit{Ensembles}$.
Nous introduisons cette curieuse terminologie parce qu'il paraît étrange de
parler des voisinages de foncteurs ; nous considérons donc un « point » $p$
comme un objet géométrique donné par une donnée catégorique, à savoir le
foncteur $u$. Le fait que $p$ soit effectivement égal à $u$ importe peu.
Un {\it voisinage} de $p$ est un couple $(U, x)$ avec
$U \in \Ob(\mathcal{C})$ et $x \in u(U)$.
Un {\it morphisme de voisinages} $(V, y) \to (U, x)$ est donné par un
morphisme $\alpha :V \to U$ de $\mathcal{C}$ tel que
$u(\alpha)(y) = x$. Remarquons que la catégorie des voisinages n'est pas une
catégorie « grande ».

\medskip\noindent
Nous définissons la {\it fibre} d'un préfaisceau $\mathcal{F}$ en $p$ par
\begin{equation}
\label{sites-equation-stalk}
\mathcal{F}_p = \colim_{\{(U, x)\}^{opp}} \mathcal{F}(U).
\end{equation}
La limite inductive est prise sur l'opposée de la catégorie des
voisinages de $p$. Autrement dit, un élément de $\mathcal{F}_p$ est donné par
un triplet $(U, x, s)$, où $(U, x)$ est un voisinage de $p$ et
$s \in \mathcal{F}(U)$. L'égalité des triplets est la relation d'équivalence
engendrée par $(U, x, s) \sim (V, y, \alpha^*s)$ lorsque $\alpha$ est comme
ci-dessus.

\medskip\noindent
Remarquons que, si $\varphi : \mathcal{F} \to \mathcal{G}$ est un morphisme
de préfaisceaux d'ensembles, on obtient une application canonique entre les
fibres $\varphi_p : \mathcal{F}_p \to \mathcal{G}_p$. On obtient ainsi un
{\it foncteur fibre}
$$
\textit{PSh}(\mathcal{C}) \longrightarrow \textit{Ensembles}, \quad
\mathcal{F} \longmapsto \mathcal{F}_p.
$$
Nous avons défini le foncteur fibre à partir d'un foncteur quelconque
$p = u : \mathcal{C} \to \textit{Ensembles}$. Aucune condition n'est nécessaire
pour que la définition ait un sens\footnote{Il convient d'éviter le cas où
$u(U) = \emptyset$ pour tout $U$.}. En revanche, il vaut probablement mieux
ne pas employer cette notion à moins que $p$ ne soit effectivement un point
(voir la définition ci-dessous), car le foncteur fibre ne possède pas, en
général, de bonnes propriétés.

\begin{definition}
\label{sites-definition-point}
Soit $\mathcal{C}$ un site. Un {\it point $p$ du site $\mathcal{C}$} est donné
par un foncteur $u : \mathcal{C} \to \textit{Ensembles}$ tel que
\begin{enumerate}
\item Pour tout recouvrement $\{U_i \to U\}$ de $\mathcal{C}$, l'application
$\coprod u(U_i) \to u(U)$ est surjective.
\item Pour tout recouvrement $\{U_i \to U\}$ de $\mathcal{C}$ et tout
morphisme $V \to U$, les applications
$u(U_i \times_U V) \to u(U_i) \times_{u(U)} u(V)$ sont bijectives.
\item Le foncteur fibre $\Sh(\mathcal{C}) \to \textit{Ensembles}$,
$\mathcal{F} \mapsto \mathcal{F}_p$, est exact à gauche.
\end{enumerate}
\end{definition}

\noindent
Ces conditions devraient être familières, car elles sont calquées sur celles
des Définitions \ref{sites-definition-continuous} et
\ref{sites-definition-morphism-sites}.
Remarquons que (3) implique que $*_p = \{*\}$ ; voir
l'Exemple \ref{sites-example-singleton-sheaf}. Ainsi,
$u(U) \not = \emptyset$ pour au moins un $U$ (car une limite inductive vide
donne l'ensemble vide). Nous montrerons plus bas
(Lemme \ref{sites-lemma-point-site-topos}) que cela donne bien naissance à un point
du topos $\Sh(\mathcal{C})$.
Avant cela, démontrons quelques lemmes pour des foncteurs $u$ généraux.

\begin{lemma}
\label{sites-lemma-points-recover}
Soit $\mathcal{C}$ un site.
Soit $p = u : \mathcal{C} \to \textit{Ensembles}$ un foncteur.
Il existe des isomorphismes fonctoriels $(h_U)_p = u(U)$ pour
$U \in \Ob(\mathcal{C})$.
\end{lemma}

\begin{proof}
Un élément de $(h_U)_p$ est donné par un triplet $(V, y, f)$, où
$V \in \Ob(\mathcal{C})$, $y\in u(V)$ et
$f \in h_U(V) = \Mor_\mathcal{C}(V, U)$.
Deux triplets de ce type $(V, y, f)$, $(V', y', f')$ déterminent le même
élément s'il existe un morphisme $\phi : V \to V'$ tel que
$u(\phi)(y) = y'$ et $f' \circ \phi = f$ ; en général, il faut prendre des
chaînes d'identifications de ce type pour obtenir la bonne relation
d'équivalence. Quoi qu'il en soit, tout $(V, y, f)$ est équivalent à l'élément
$(U, u(f)(y), \text{id}_U)$. Si $\phi$ existe comme ci-dessus, alors les
triplets $(V, y, f)$, $(V', y', f')$ déterminent le même triplet
$(U, u(f)(y), \text{id}_U) = (U, u(f')(y'), \text{id}_U)$.
Cela montre que l'application
$u(U) \to (h_U)_p$, $x \mapsto \text{classe de }(U, x, \text{id}_U)$
est bijective.
\end{proof}

\noindent
Soit $\mathcal{C}$ un site. Soit
$p = u : \mathcal{C} \to \textit{Ensembles}$ un foncteur. Par analogie avec les
constructions de la section \ref{sites-section-functoriality-PSh}, étant donné un
ensemble $E$, nous définissons un préfaisceau $u^pE$ par la règle
\begin{equation}
\label{sites-equation-skyscraper}
U
\longmapsto
u^pE(U) = \Mor_{\textit{Ensembles}}(u(U), E) = \text{Map}(u(U), E).
\end{equation}
Cela définit un foncteur
$u^p : \textit{Ensembles} \to \textit{PSh}(\mathcal{C})$, $E \mapsto u^pE$.

\begin{lemma}
\label{sites-lemma-adjoint-point-push-stalk}
Pour tout foncteur $u : \mathcal{C} \to \textit{Ensembles}$, le foncteur $u^p$ est
adjoint à droite du foncteur fibre sur les préfaisceaux.
\end{lemma}

\begin{proof}
Soit $\mathcal{F}$ un préfaisceau sur $\mathcal{C}$.
Soit $E$ un ensemble. Un morphisme $\mathcal{F} \to u^pE$ est donné par un
système compatible d'applications
$\mathcal{F}(U) \to \text{Map}(u(U), E)$, c'est-à-dire par un système
compatible d'applications $\mathcal{F}(U) \times u(U) \to E$.
Et, par définition de $\mathcal{F}_p$, une application
$\mathcal{F}_p \to E$ est donnée par une règle qui associe à chaque triplet
$(U, x, \sigma)$ un élément de $E$, de telle sorte que des triplets équivalents
aient la même image ; voir la discussion entourant
l'Équation (\ref{sites-equation-stalk}).
C'est encore la donnée d'un système compatible d'applications
$\mathcal{F}(U) \times u(U) \to E$.
\end{proof}

\noindent
Par analogie avec la section \ref{sites-section-continuous-functors}, on a le lemme
suivant.

\begin{lemma}
\label{sites-lemma-point-pushforward-sheaf}
Soit $\mathcal{C}$ un site. Soit
$p = u : \mathcal{C} \to \textit{Ensembles}$ un foncteur. Supposons que, pour tout
recouvrement $\{U_i \to U\}$ de $\mathcal{C}$,
\begin{enumerate}
\item l'application $\coprod u(U_i) \to u(U)$ soit surjective, et
\item les applications
$u(U_i \times_U U_j) \to u(U_i) \times_{u(U)} u(U_j)$ soient surjectives.
\end{enumerate}
Alors
\begin{enumerate}
\item le préfaisceau $u^pE$ est un faisceau pour tout ensemble $E$ ;
notons-le $u^sE$,
\item le foncteur fibre $\Sh(\mathcal{C}) \to \textit{Ensembles}$ et le foncteur
$u^s: \textit{Ensembles} \to \Sh(\mathcal{C})$ sont adjoints, et
\item on a $\mathcal{F}_p = \mathcal{F}^\#_p$ pour tout préfaisceau
d'ensembles $\mathcal{F}$.
\end{enumerate}
\end{lemma}

\begin{proof}
La première assertion résulte immédiatement de la définition d'un faisceau,
des hypothèses (1) et (2), et de la définition de $u^pE$. La deuxième est une
reformulation de l'adjonction entre $u^p$ et le foncteur fibre
(Lemme \ref{sites-lemma-adjoint-point-push-stalk}) restreinte aux faisceaux.
La troisième assertion résulte de ce que, pour tout ensemble $E$, on a
$$
\text{Map}(\mathcal{F}_p, E) =
\Mor_{\textit{PSh}(\mathcal{C})}(\mathcal{F}, u^pE) =
\Mor_{\Sh(\mathcal{C})}(\mathcal{F}^\#, u^sE) =
\text{Map}(\mathcal{F}^\#_p, E)
$$
par la propriété d'adjonction du passage au faisceau associé.
\end{proof}

\noindent
En particulier, le Lemme \ref{sites-lemma-point-pushforward-sheaf} s'applique
lorsque $p = u$ est un point. Dans ce cas, nous considérons le faisceau $u^sE$
comme le faisceau « gratte-ciel » de valeur $E$ en $p$.

\begin{definition}
\label{sites-definition-pushforward-point}
Soit $p$ un point du site $\mathcal{C}$ donné par le foncteur $u$.
Pour un ensemble $E$, nous définissons $p_*E = u^sE$ comme le faisceau décrit
dans le Lemme \ref{sites-lemma-point-pushforward-sheaf} ci-dessus.
Nous l'appelons parfois un {\it faisceau gratte-ciel}.
\end{definition}

\noindent
On a en particulier la propriété d'adjonction suivante pour les faisceaux
gratte-ciel et les fibres :
$$
\Mor_{\Sh(\mathcal{C})}(\mathcal{F}, p_*E)
=
\text{Map}(\mathcal{F}_p, E)
$$
Cela motive la notation $p^{-1}\mathcal{F} = \mathcal{F}_p$ que nous
emploierons parfois.

\begin{lemma}
\label{sites-lemma-point-site-topos}
Soit $\mathcal{C}$ un site.
\begin{enumerate}
\item Soit $p$ un point du site $\mathcal{C}$.
Alors le couple de foncteurs $(p_*, p^{-1})$ introduit ci-dessus définit un
morphisme de topos $\Sh(pt) \to \Sh(\mathcal{C})$.
\item Soit $p = (p_*, p^{-1})$ un point du topos $\Sh(\mathcal{C})$.
Alors le foncteur $u : U \mapsto p^{-1}(h_U^\#)$ donne naissance à un point
$p'$ du site $\mathcal{C}$ dont le morphisme de topos associé
$(p'_*, (p')^{-1})$ est égal à $p$.
\end{enumerate}
\end{lemma}

\begin{proof}
Démonstration de (1). D'après ce qui précède, les foncteurs $p_*$ et $p^{-1}$
sont adjoints. Le foncteur $p^{-1}$ est exact par la
Définition \ref{sites-definition-point}.
Toutes les conditions imposées dans la Définition \ref{sites-definition-topos} sont
donc satisfaites, ce qui démontre (1).

\medskip\noindent
Démonstration de (2). Soit $\{U_i \to U\}$ un recouvrement de $\mathcal{C}$.
Alors $\coprod (h_{U_i})^\# \to h_U^\#$ est surjective ; voir le
Lemme \ref{sites-lemma-covering-surjective-after-sheafification}.
Comme $p^{-1}$ est exact (par définition d'un morphisme de topos), on en
conclut que $\coprod u(U_i) \to u(U)$ est surjective.
Cela démontre la partie (1) de la Définition \ref{sites-definition-point}.
Le passage au faisceau associé est exact ; voir le
Lemme \ref{sites-lemma-sheafification-exact}.
Ainsi, si $U \times_V W$ existe dans $\mathcal{C}$, alors
$$
h_{U \times_V W}^\# =  h_U^\# \times_{h_V^\#} h_W^\#
$$
et l'on voit que $u(U \times_V W) = u(U) \times_{u(V)} u(W)$, puisque
$p^{-1}$ est exact. Cela démontre la partie (2) de la
Définition \ref{sites-definition-point}.
Posons $p' = u$ et soit $\mathcal{F}_{p'}$ le foncteur fibre défini par
l'Équation (\ref{sites-equation-stalk}) au moyen de $u$. Il existe une application
canonique de comparaison
$c : \mathcal{F}_{p'} \to \mathcal{F}_p = p^{-1}\mathcal{F}$.
En effet, étant donné un triplet $(U, x, \sigma)$ représentant un élément
$\xi$ de $\mathcal{F}_{p'}$, nous considérons $\sigma$ comme une application
$\sigma : h_U^\# \to \mathcal{F}$ et nous pouvons poser
$c(\xi) = p^{-1}(\sigma)(x)$, puisque
$x \in u(U) = p^{-1}(h_U^\#)$. D'après le
Lemme \ref{sites-lemma-points-recover}, on a $(h_U)_{p'} = u(U)$. Puisque les
conditions (1) et (2) de la Définition \ref{sites-definition-point} sont satisfaites
pour $p'$, on a également $(h_U^\#)_{p'} = (h_U)_{p'}$ d'après le
Lemme \ref{sites-lemma-point-pushforward-sheaf}. Ainsi,
$$
(h_U^\#)_{p'} = (h_U)_{p'} = u(U) = p^{-1}(h_U^\#)
$$
Nous affirmons que cette bijection est égale à l'application de comparaison
$c : (h_U^\#)_{p'} \to p^{-1}(h_U^\#)$ (vérification omise).
Tout faisceau sur $\mathcal{C}$ est un coégalisateur d'applications entre des
coproduits de faisceaux de la forme $h_U^\#$ ; voir le
Lemme \ref{sites-lemma-sheaf-coequalizer-representable}.
Le foncteur fibre $\mathcal{F} \mapsto \mathcal{F}_{p'}$ et le foncteur
$p^{-1}$ commutent aux limites inductives arbitraires (car ils sont tous deux
adjoints à gauche).
On en conclut que $c$ est un isomorphisme pour tout faisceau $\mathcal{F}$.
Ainsi, le foncteur fibre $\mathcal{F} \mapsto \mathcal{F}_{p'}$ est isomorphe
à $p^{-1}$ et, en particulier, il est exact. Cela démontre que la condition
(3) de la Définition \ref{sites-definition-point} est satisfaite et que $p'$ est un
point. La dernière assertion a déjà été démontrée ci-dessus, puisque nous
avons vu que $p^{-1} = (p')^{-1}$.
\end{proof}

\noindent
En fait, un point correspond toujours à un morphisme de sites, comme le montre
le lemme suivant.

\begin{lemma}
\label{sites-lemma-site-point-morphism}
Soit $\mathcal{C}$ un site. Soit $p$ un point de $\mathcal{C}$ donné par
$u : \mathcal{C} \to \textit{Ensembles}$. Soit $S_0$ un ensemble infini tel que
$u(U) \subset S_0$ pour tout $U \in \Ob(\mathcal{C})$. Soit $\mathcal{S}$ le
site construit à partir de l'ensemble des parties
$S = \mathcal{P}(S_0)$ dans la Remarque \ref{sites-remark-pt-topos}.
Alors
\begin{enumerate}
\item il existe une équivalence $i : \Sh(pt) \to \Sh(\mathcal{S})$,
\item le foncteur $u : \mathcal{C} \to \mathcal{S}$ induit un morphisme de
sites $f : \mathcal{S} \to \mathcal{C}$, et
\item le composé
$$
\Sh(pt) \to
\Sh(\mathcal{S}) \to
\Sh(\mathcal{C})
$$
est le morphisme de topos $(p_*, p^{-1})$ du
Lemme \ref{sites-lemma-point-site-topos}.
\end{enumerate}
\end{lemma}

\begin{proof}
La partie (1) a été établie dans la Remarque \ref{sites-remark-pt-topos}.
Rappelons en outre que l'équivalence associe à l'ensemble $E$ le faisceau
$i_*E$ sur $\mathcal{S}$ défini par la règle
$V \mapsto \Mor_{\textit{Ensembles}}(V, E)$.
La partie (2) résulte clairement de la définition d'un point de $\mathcal{C}$
(Définition \ref{sites-definition-point}) et de la définition d'un morphisme de
sites (Définition \ref{sites-definition-morphism-sites}).
Enfin, considérons $f_*i_*E$. Par construction, on a
$$
f_*i_*E(U) = i_*E(u(U)) = \Mor_{\textit{Ensembles}}(u(U), E)
$$
ce qui est égal à $p_*E(U)$ ; voir
l'Équation (\ref{sites-equation-skyscraper}).
Cela démontre (3).
\end{proof}

\noindent
Contrairement à ce qui se passe dans le cas topologique, la fibre du faisceau
gratte-ciel de valeur $E$ n'est pas toujours égale à $E$. Voici ce qui est
toujours vrai.

\begin{lemma}
\label{sites-lemma-stalk-skyscraper}
Soit $\mathcal{C}$ un site. Soit
$p : \Sh(pt) \to \Sh(\mathcal{C})$ un point du topos associé à
$\mathcal{C}$. Pour tout ensemble $E$, il existe des applications canoniques
$$
E \longrightarrow (p_*E)_p \longrightarrow E
$$
dont le composé est $\text{id}_E$.
\end{lemma}

\begin{proof}
Il existe toujours un morphisme d'adjonction
$(p_*E)_p = p^{-1}p_*E \to E$.
Ce morphisme est un isomorphisme lorsque $E = \{*\}$, car $p_*$ et $p^{-1}$
sont tous deux exacts à gauche et transforment donc l'objet final en l'objet
final. Étant donné $e \in E$, on peut donc considérer l'application
$i_e : \{*\} \to E$, qui donne
$$
\xymatrix{
p^{-1}p_*\{*\} \ar[rr]_{p^{-1}p_*i_e} \ar[d]_{\cong} & & p^{-1}p_*E \ar[d] \\
\{*\} \ar[rr]^{i_e} & & E
}
$$
d'où l'application $E \to (p_*E)_p = p^{-1}p_*E$ voulue.
\end{proof}

\begin{lemma}
\label{sites-lemma-skyscraper-functor-exact}
Soit $\mathcal{C}$ un site. Soit
$p : \Sh(pt) \to \Sh(\mathcal{C})$ un point du topos associé à
$\mathcal{C}$. Le foncteur $p_* : \textit{Ensembles} \to \Sh(\mathcal{C})$
possède les propriétés suivantes : il commute aux limites projectives
arbitraires, il est exact à gauche, il est fidèle, il transforme les
surjections en surjections, il commute aux coégalisateurs, il reflète les
injections, il reflète les surjections et il reflète les isomorphismes.
\end{lemma}

\begin{proof}
Comme $p_*$ est adjoint à droite, il commute aux limites projectives
arbitraires et est exact à gauche. Le fait que
$p^{-1}p_*E \to E$ soit une surjection canoniquement scindée implique que
$p_*$ est fidèle, reflète les injections, reflète les surjections et reflète
les isomorphismes. D'après le Lemme \ref{sites-lemma-point-site-topos}, on peut
supposer que $p$ provient d'un point
$u : \mathcal{C} \to \textit{Ensembles}$ du site sous-jacent $\mathcal{C}$.
Dans ce cas, le faisceau $p_*E$ est donné par
$$
p_*E(U) = \Mor_{\textit{Ensembles}}(u(U), E)
$$
voir l'Équation (\ref{sites-equation-skyscraper}) et la
Définition \ref{sites-definition-pushforward-point}.
Il résulte immédiatement de cette formule que $p_*$ transforme les
surjections en surjections et les coégalisateurs en coégalisateurs.
\end{proof}



\section{Construction de points}
\label{sites-section-construct-points}

\noindent
Dans cette section, nous donnons des critères pour qu'un foncteur
d'un site vers la catégorie des ensembles définisse un point de ce site.

\begin{lemma}
\label{sites-lemma-neighbourhoods-cofiltered}
Soit $\mathcal{C}$ un site. Soit $p = u : \mathcal{C} \to \textit{Ensembles}$
un foncteur. Si la catégorie des voisinages de $p$ est
cofiltrante, alors le foncteur fibre (\ref{sites-equation-stalk})
est exact à gauche.
\end{lemma}

\begin{proof}
Soit $\mathcal{I} \to \Sh(\mathcal{C})$,
$i \mapsto \mathcal{F}_i$ un diagramme fini de faisceaux.
Nous devons montrer que la fibre de la limite de ce
système coïncide avec la limite des fibres.
Soit $\mathcal{F}$ la limite du système en tant que {\it préfaisceau}.
D'après le Lemme \ref{sites-lemma-limit-sheaf}, c'est un faisceau et
c'est la limite dans la catégorie des faisceaux.
Il nous faut donc montrer que
$\mathcal{F}_p = \lim_\mathcal{I} \mathcal{F}_{i, p}$.
Rappelons aussi que $\mathcal{F}$ admet une description simple ; voir la
section \ref{sites-section-limits-colimits-PSh}. Il nous faut ainsi montrer que
$$
\lim_i \colim_{\{(U, x)\}^{opp}} \mathcal{F}_i(U)
=
\colim_{\{(U, x)\}^{opp}} \lim_i \mathcal{F}_i(U).
$$
Cela résulte de Catégories, Lemme \ref{categories-lemma-directed-commutes},
car l'opposée de la catégorie des voisinages est filtrante
par hypothèse.
\end{proof}

\begin{lemma}
\label{sites-lemma-neighbourhoods-directed}
Soit $\mathcal{C}$ un site. Supposons que $\mathcal{C}$ possède
un objet final $X$ et des produits fibrés.
Soit $p = u : \mathcal{C} \to \textit{Ensembles}$ un foncteur tel que
\begin{enumerate}
\item $u(X)$ soit un singleton, et
\item pour toute paire de morphismes $U \to W$ et $V \to W$ de
même but, l'application
$u(U \times_W V) \to u(U) \times_{u(W)} u(V)$ soit bijective.
\end{enumerate}
Alors la catégorie des voisinages de $p$ est cofiltrante
et, par conséquent, le foncteur fibre $\Sh(\mathcal{C}) \to \textit{Ensembles}$,
$\mathcal{F} \to \mathcal{F}_p$ commute aux limites finies.
\end{lemma}

\begin{proof}
Remarquons l'analogie avec le Lemme \ref{sites-lemma-directed}.
Les hypothèses sur $\mathcal{C}$ impliquent que $\mathcal{C}$ admet les limites finies.
Voir Catégories, Lemme \ref{categories-lemma-finite-limits-exist}.
L'hypothèse (1) implique que la catégorie des voisinages
n'est pas vide. Soient $(U, x)$ et $(V, y)$ deux voisinages.
Alors
$u(U \times V) = u(U \times_X V) =
u(U) \times_{u(X)} u(V) = u(U) \times u(V)$ d'après (2).
Il existe donc un voisinage $(U \times_X V, z)$ se projetant
sur $(U, x)$ et sur $(V, y)$.
Soient $a, b : (V, y) \to (U, x)$ deux morphismes
de la catégorie des voisinages. Soit $W$ l'égalisateur de
$a, b : V \to U$. Comme dans la démonstration de
Catégories, Lemme \ref{categories-lemma-finite-limits-exist},
nous pouvons écrire $W$ au moyen de produits fibrés :
$$
W = (V \times_{a, U, b} V) \times_{(pr_1, pr_2), V \times V, \Delta} V
$$
La bijectivité de (2) garantit l'existence d'un élément $z \in u(W)$
dont l'image est $((y, y), y)$.
Alors $(W, z) \to (V, y)$ égalise $a, b$, comme voulu.
La conclusion « par conséquent » est le Lemme \ref{sites-lemma-neighbourhoods-cofiltered}.
\end{proof}

\begin{proposition}
\label{sites-proposition-point-limits}
Soit $\mathcal{C}$ un site. Supposons que les limites finies existent
dans $\mathcal{C}$. (Autrement dit, $\mathcal{C}$ possède des produits fibrés et un
objet final.) Un point $p$ d'un tel site $\mathcal{C}$
est donné par un foncteur $u : \mathcal{C} \to \textit{Ensembles}$ tel que
\begin{enumerate}
\item $u$ commute aux limites finies, et
\item si $\{U_i \to U\}$ est un recouvrement, alors
$\coprod_i u(U_i) \to u(U)$ est surjective.
\end{enumerate}
\end{proposition}

\begin{proof}
Supposons d'abord que $p$ soit un point (Définition \ref{sites-definition-point})
donné par un foncteur $u$. La condition (2) résulte directement de
la définition d'un point. D'après le Lemme \ref{sites-lemma-points-recover},
nous avons $(h_U)_p = u(U)$. D'après le Lemme \ref{sites-lemma-point-pushforward-sheaf},
nous avons $(h_U^\#)_p = (h_U)_p$. Ainsi, $u$
est égal au composé des foncteurs
$$
\mathcal{C} \xrightarrow{h}
\textit{PSh}(\mathcal{C}) \xrightarrow{{}^\#}
\Sh(\mathcal{C}) \xrightarrow{()_p}
\textit{Ensembles}
$$
Chacun de ces foncteurs est exact à gauche ; ainsi, $u$ satisfait (1).

\medskip\noindent
Réciproquement, supposons que $u$ satisfasse (1) et (2).
Nous voyons aussitôt que $u$ satisfait les deux premières
conditions de la Définition \ref{sites-definition-point}. De plus, son
foncteur fibre est exact, car c'est un adjoint à gauche d'après le
Lemme \ref{sites-lemma-point-pushforward-sheaf}, et il commute
aux limites finies d'après le Lemme \ref{sites-lemma-neighbourhoods-directed}.
\end{proof}

\begin{remark}
\label{sites-remark-improve-proposition-points-limits}
En fait, soit $\mathcal{C}$ un site. Supposons que $\mathcal{C}$ possède un objet final
$X$ et des produits fibrés. Soit $p = u: \mathcal{C} \to \textit{Ensembles}$ un
foncteur tel que
\begin{enumerate}
\item $u(X) = \{*\}$ soit un singleton, et
\item pour toute paire de morphismes $U \to W$ et $V \to W$ de
même but, l'application
$u(U \times_W V) \to u(U) \times_{u(W)} u(V)$ soit surjective.
\item pour tout recouvrement $\{U_i \to U\}$, l'application
$\coprod u(U_i) \to u(U)$ soit surjective.
\end{enumerate}
Alors, en général, $p$ n'est {\bf pas} un point de $\mathcal{C}$.
Un exemple est la catégorie $\mathcal{C}$ à deux objets $\{U, X\}$
et une seule flèche non identique, à savoir $U \to X$. Munissons $\mathcal{C}$
de la topologie triviale, c'est-à-dire que les seuls recouvrements sont $\{U \to U\}$ et
$\{X \to X\}$. Un faisceau $\mathcal{F}$ est alors la même chose qu'un préfaisceau et
consiste en un triplet $(A, B, A \to B)$ : à savoir $A = \mathcal{F}(X)$,
$B = \mathcal{F}(U)$, où $A \to B$ est l'application de restriction correspondant
à $U \to X$. Remarquons que $U \times_X U = U$ ; les produits fibrés existent donc.
Considérons le foncteur $u = p$ défini par $u(X) = \{*\}$ et $u(U) = \{*_1, *_2\}$.
Il satisfait (1), (2) et (3), mais le foncteur fibre correspondant
(\ref{sites-equation-stalk}) est le foncteur
$$
(A, B, A \to B) \longmapsto B \amalg_A B
$$
qui n'est pas exact. En effet, considérons
$(\emptyset, \{1\}, \emptyset \to \{1\}) \to (\{1\}, \{1\}, \{1\} \to \{1\})$,
qui est une application injective de faisceaux mais qui est transformée en l'application
non injective d'ensembles
$$
\{1\} \amalg \{1\} \longrightarrow \{1\} \amalg_{\{1\}} \{1\}
$$
par le foncteur fibre.
\end{remark}

\begin{example}
\label{sites-example-point-topological}
Soit $X$ un espace topologique. Soit $X_{Zar}$ le site de
l'Exemple \ref{sites-example-site-topological}.
Soit $x \in X$ un point. Considérons le foncteur
$$
u : X_{Zar} \longrightarrow \textit{Ensembles}, \quad
U \mapsto
\left\{
\begin{matrix}
\emptyset & \text{si} & x \not \in U \\
\{*\} & \text{si} & x \in U
\end{matrix}
\right.
$$
Ce foncteur commute aux produits et aux produits fibrés,
et transforme les recouvrements en familles surjectives d'applications.
Nous obtenons donc un point $p$ du site $X_{Zar}$.
On vérifie immédiatement que le foncteur fibre
coïncide avec la fibre en $x$ définie dans
Faisceaux, section
\ref{sheaves-section-stalks}.
\end{example}

\begin{example}
\label{sites-example-point-topology}
Soit $X$ un espace topologique. Quels sont les points du topos
$\Sh(X)$ ? Pour le déterminer, soit $X_{Zar}$ le site de
l'Exemple \ref{sites-example-site-topological}.
D'après le
Lemme \ref{sites-lemma-point-site-topos},
un point de $\Sh(X)$ correspond à un point de ce site.
Soit $p$ un point du site $X_{Zar}$ donné par le foncteur
$u : X_{Zar} \to \textit{Ensembles}$. Nous allons utiliser la
caractérisation de tels $u$ donnée dans la
Proposition \ref{sites-proposition-point-limits}.
Elle implique immédiatement que $u(\emptyset) = \emptyset$ et
$u(U \cap V) = u(U) \times u(V)$. En particulier,
$u(U) = u(U) \times u(U)$ par l'application diagonale, ce qui implique que $u(U)$
est soit un singleton, soit vide. En outre, si $U = \bigcup U_i$ est un
recouvrement ouvert, alors
$$
u(U) = \emptyset \Rightarrow \forall i, \ u(U_i) = \emptyset
\quad\text{et}\quad
u(U) \not = \emptyset \Rightarrow \exists i, \ u(U_i) \not = \emptyset.
$$
Nous en concluons qu'il existe un unique plus grand ouvert $W \subset X$ tel que
$u(W) = \emptyset$, à savoir la réunion de tous les ouverts $U$ tels que
$u(U) = \emptyset$. Posons $Z = X \setminus W$. Si $Z = Z_1 \cup Z_2$ avec
$Z_i \subset Z$ fermés, alors $W = (X \setminus Z_1) \cap (X \setminus Z_2)$ ;
ainsi $\emptyset = u(W) = u(X \setminus Z_1) \times u(X \setminus Z_2)$,
et nous concluons que $u(X \setminus Z_1) = \emptyset$ ou que
$u(X \setminus Z_2) = \emptyset$. Cela signifie que $X \setminus Z_1 = W$
ou que $X \setminus Z_2 = W$. Autrement dit, $Z$ est irréductible.
Nous voyons maintenant que $u$ est décrit par la règle
$$
u : X_{Zar} \longrightarrow \textit{Ensembles}, \quad
U \mapsto
\left\{
\begin{matrix}
\emptyset & \text{si} & Z \cap U = \emptyset \\
\{*\} & \text{si} & Z \cap U \not = \emptyset
\end{matrix}
\right.
$$
Remarquons que, pour tout fermé irréductible $Z \subset X$, ce
foncteur satisfait les hypothèses (1), (2) de la
Proposition \ref{sites-proposition-point-limits}
et définit donc un point. Autrement dit, les points du
site $X_{Zar}$ sont en correspondance bijective avec les
sous-ensembles fermés irréductibles de $X$. En particulier, si $X$ est
un espace topologique sobre, alors les points de $X_{Zar}$ et les
points de $X$ sont en correspondance bijective ; voir
l'Exemple \ref{sites-example-point-topological}.
\end{example}

\begin{example}
\label{sites-example-point-G-sets}
Considérons le site $\mathcal{T}_G$ décrit dans
l'Exemple \ref{sites-example-site-on-group} et la
section \ref{sites-section-example-sheaf-G-sets}.
Le foncteur d'oubli $u : \mathcal{T}_G \to \textit{Ensembles}$
commute aux produits et aux produits fibrés, et transforme
les recouvrements en familles surjectives. Il définit donc un point
de $\mathcal{T}_G$. Identifions $\Sh(\mathcal{T}_G)$
et $G\textit{-Ensembles}$. Le foncteur fibre
$$
p^{-1} :
\Sh(\mathcal{T}_G) = G\textit{-Ensembles}
\longrightarrow
\textit{Ensembles}
$$
est le foncteur d'oubli. L'image directe $p_*$ est le
foncteur
$$
\textit{Ensembles}
\longrightarrow
\Sh(\mathcal{T}_G) = G\textit{-Ensembles}
$$
qui envoie un ensemble $S$ sur le $G$-ensemble $\text{Map}(G, S)$ muni de
l'action $g \cdot \psi = \psi \circ R_g$, où $R_g$ est la multiplication
à droite. En particulier,
$p^{-1}p_*S = \text{Map}(G, S)$ en tant qu'ensemble, et les applications
$S \to \text{Map}(G, S) \to S$ du
Lemme \ref{sites-lemma-stalk-skyscraper}
sont les applications évidentes.
\end{example}

\begin{example}
\label{sites-example-indiscrete-points}
Soit $\mathcal{C}$ une catégorie munie de la topologie grossière
(Exemple \ref{sites-example-indiscrete}). Pour tout objet $U_0$ de $\mathcal{C}$,
le foncteur $u : U \mapsto \Mor_\mathcal{C}(U_0, U)$ définit un point
$p$ de $\mathcal{C}$. En effet, les conditions (1) et (2) de la
Définition \ref{sites-definition-point} sont immédiates puisque les seuls recouvrements
sont donnés par les applications identiques. La condition (3) est satisfaite car
$\mathcal{F}_p = \mathcal{F}(U_0)$ et, la topologie étant grossière,
prendre les sections sur $U_0$ est un foncteur exact.
\end{example}





\section{Points et morphismes de topos}
\label{sites-section-functorial-points}

\noindent
Dans cette section, nous formulons quelques remarques sur les points et les morphismes
de topos.

\begin{lemma}
\label{sites-lemma-point-functor}
Soit $u : \mathcal{C} \to \mathcal{D}$ un foncteur. Soit
$v : \mathcal{D} \to \textit{Ensembles}$ un foncteur, et posons
$w = v \circ u$. Notons $q$, respectivement $p$, le foncteur fibre
(\ref{sites-equation-stalk}) associé à $v$, respectivement à $w$.
Alors $(u_p\mathcal{F})_q = \mathcal{F}_p$, fonctoriellement par rapport au
préfaisceau $\mathcal{F}$ sur $\mathcal{C}$.
\end{lemma}

\begin{proof}
C'est un fait catégorique simple. Nous avons
\begin{align*}
(u_p\mathcal{F})_q
& =
\colim_{(V, y)} \colim_{U, \phi : V \to u(U)} \mathcal{F}(U) \\
& = \colim_{(V, y, U, \phi : V \to u(U))} \mathcal{F}(U) \\
& = \colim_{(U, x)} \mathcal{F}(U) \\
& = \mathcal{F}_p
\end{align*}
La première égalité résulte de la définition de $u_p$ et de la
définition du foncteur fibre. Observons que $y \in v(V)$.
Dans la deuxième égalité, nous réunissons simplement les limites inductives.
Pour établir la troisième égalité, appliquons
Catégories, Lemme \ref{categories-lemma-colimit-constant-connected-fibers},
au foncteur $F$ des catégories de diagrammes défini par la règle
$$
F((V, y, U, \phi : V \to u(U))) = (U, v(\phi)(y)).
$$
Cela a un sens car $w(U) = v(u(U))$. Vérifions les hypothèses de
Catégories, Lemme \ref{categories-lemma-colimit-constant-connected-fibers}.
Observons que $F$ possède un inverse à droite, à savoir
$(U, x) \mapsto (u(U), x, U, \text{id} : u(U) \to u(U))$.
Cela a encore un sens car $x \in w(U) = v(u(U))$. D'autre part,
il existe toujours un morphisme
$$
(V, y, U, \phi : V \to u(U))
\longrightarrow
(u(U), v(\phi)(y), U, \text{id} : u(U) \to u(U))
$$
dans la catégorie fibre au-dessus de $(U, x)$, ce qui montre que les catégories fibres
sont connexes. La quatrième et dernière égalité est claire.
\end{proof}

\begin{lemma}
\label{sites-lemma-point-morphism-sites}
\begin{slogan}
Un morphisme de sites définit un morphisme sur les points, et l'image inverse respecte les fibres.
\end{slogan}
Soit $f : \mathcal{D} \to \mathcal{C}$ un morphisme de sites
donné par un foncteur continu $u : \mathcal{C} \to \mathcal{D}$.
Soit $q$ un point de $\mathcal{D}$ donné par le foncteur
$v : \mathcal{D} \to \textit{Ensembles}$ ; voir la
Définition \ref{sites-definition-point}.
Alors le foncteur $v \circ u : \mathcal{C} \to \textit{Ensembles}$
définit un point $p$ de $\mathcal{C}$ et il existe en outre
une identification canonique
$$
(f^{-1}\mathcal{F})_q = \mathcal{F}_p
$$
pour tout faisceau $\mathcal{F}$ sur $\mathcal{C}$.
\end{lemma}

\begin{proof}[Première démonstration du Lemme \ref{sites-lemma-point-morphism-sites}]
Puisque $u$ est continu et que $v$ définit un point,
il est immédiat que $v \circ u$ satisfait les conditions (1) et (2) de la
Définition \ref{sites-definition-point}. Démontrons l'égalité affichée.
Soit $\mathcal{F}$ un faisceau sur $\mathcal{C}$. Alors
$$
(f^{-1}\mathcal{F})_q = (u_s\mathcal{F})_q =
(u_p \mathcal{F})_q = \mathcal{F}_p
$$
La première égalité tient à $f^{-1} = u_s$, la deuxième au
Lemme \ref{sites-lemma-point-pushforward-sheaf}, et la troisième
au Lemme \ref{sites-lemma-point-functor}.
Nous voyons alors que $p$ satisfait aussi la condition (3) de la
Définition \ref{sites-definition-point},
car son foncteur fibre est un composé de foncteurs exacts. Cela achève la démonstration.
\end{proof}

\begin{proof}[Seconde démonstration du Lemme \ref{sites-lemma-point-morphism-sites}]
D'après le
Lemme \ref{sites-lemma-site-point-morphism},
nous pouvons factoriser $(q_*, q^{-1})$ sous la forme
$$
\Sh(pt) \xrightarrow{i} \Sh(\mathcal{S})
\xrightarrow{h} \Sh(\mathcal{D})
$$
où le second morphisme de topos provient d'un morphisme de sites
$h : \mathcal{S} \to \mathcal{D}$ induit par le foncteur
$v : \mathcal{D} \to \mathcal{S}$ (ce qui a un sens puisque
$\mathcal{S} \subset \textit{Ensembles}$ est une sous-catégorie pleine contenant
tout objet de l'image de $v$). D'après le
Lemme \ref{sites-lemma-composition-morphisms-sites},
le composé $v \circ u : \mathcal{C} \to \mathcal{S}$
définit un morphisme de sites $g : \mathcal{S} \to \mathcal{C}$.
En particulier, le foncteur
$v \circ u : \mathcal{C} \to \mathcal{S}$
est continu, ce qui, par la définition des recouvrements
de $\mathcal{S}$ (voir la
Remarque \ref{sites-remark-pt-topos}),
signifie que $v \circ u$ satisfait les conditions (1) et (2) de la
Définition \ref{sites-definition-point}.
D'autre part, nous voyons que
$$
g_*i_*E(U) = i_*E(v(u(U)) = \Mor_{\textit{Ensembles}}(v(u(U)), E)
$$
par la construction de $i$ dans la
Remarque \ref{sites-remark-pt-topos}.
Remarquons qu'il s'agit précisément de la formule pour
$(v \circ u)^pE$ ; voir l'Équation (\ref{sites-equation-skyscraper}).
D'après le
Lemme \ref{sites-lemma-point-pushforward-sheaf},
le foncteur $g_*i_* = (v \circ u)^p = (v \circ u)^s$
est adjoint à droite du foncteur fibre
$\mathcal{F} \mapsto \mathcal{F}_q$.
Nous voyons donc que le foncteur fibre $p^{-1}$ est canoniquement
isomorphe à $i^{-1} \circ g^{-1}$. Il est donc exact, et
nous concluons que $p$ est un point. Enfin, puisque
$g = f \circ h$ par construction, nous voyons que
$p^{-1} = i^{-1} \circ h^{-1} \circ f^{-1} = q^{-1} \circ f^{-1}$,
c'est-à-dire que nous avons la formule affichée du lemme.
\end{proof}

\begin{lemma}
\label{sites-lemma-point-morphism-topoi}
Soit $f : \Sh(\mathcal{D}) \to \Sh(\mathcal{C})$
un morphisme de topos. Soit $q : \Sh(pt) \to \Sh(\mathcal{D})$
un point. Alors $p = f \circ q$ est un point du topos
$\Sh(\mathcal{C})$ et nous avons
une identification canonique
$$
(f^{-1}\mathcal{F})_q = \mathcal{F}_p
$$
pour tout faisceau $\mathcal{F}$ sur $\mathcal{C}$.
\end{lemma}

\begin{proof}
Cela résulte immédiatement des définitions et du fait que nous pouvons
composer les morphismes de topos.
\end{proof}





\section{Localisation et points}
\label{sites-section-localize-points}

\noindent
Dans cette section, nous montrons que les points d'une localisation $\mathcal{C}/U$
se construisent simplement à partir des points de $\mathcal{C}$.

\begin{lemma}
\label{sites-lemma-point-localize}
Soit $\mathcal{C}$ un site. Soit $p$ un point de $\mathcal{C}$ donné par
$u : \mathcal{C} \to \textit{Ensembles}$. Soit $U$ un objet de $\mathcal{C}$
et soit $x \in u(U)$. Le foncteur
$$
v : \mathcal{C}/U \longrightarrow \textit{Ensembles}, \quad
(\varphi : V \to U) \longmapsto \{y \in u(V) \mid u(\varphi)(y) = x\}
$$
définit un point $q$ du site $\mathcal{C}/U$ tel que le diagramme
$$
\xymatrix{
& \Sh(pt) \ar[d]^p \ar[ld]_q \\
\Sh(\mathcal{C}/U) \ar[r]^{j_U} &
\Sh(\mathcal{C})
}
$$
soit commutatif. Autrement dit,
$\mathcal{F}_p = (j_U^{-1}\mathcal{F})_q$ pour tout
faisceau sur $\mathcal{C}$.
\end{lemma}

\begin{proof}
Choisissons $S$ et $\mathcal{S}$ comme dans le
Lemme \ref{sites-lemma-site-point-morphism}.
Nous pouvons identifier $\Sh(pt) = \Sh(\mathcal{S})$
comme dans ce lemme, et écrire
$p = f : \Sh(\mathcal{S}) \to \Sh(\mathcal{C})$
pour le morphisme de topos induit par $u$.
D'après le
Lemme \ref{sites-lemma-localize-morphism},
nous obtenons un diagramme commutatif de topos
$$
\xymatrix{
\Sh(\mathcal{S}/u(U)) \ar[r]_-{j_{u(U)}} \ar[d]_{p'} &
\Sh(\mathcal{S}) \ar[d]^p \\
\Sh(\mathcal{C}/U) \ar[r]^{j_U} &
\Sh(\mathcal{C}),
}
$$
où $p'$ est donné par le foncteur $u' : \mathcal{C}/U \to \mathcal{S}/u(U)$,
$V/U \mapsto u(V)/u(U)$.
Considérons le foncteur $j_x : \mathcal{S} \cong \mathcal{S}/x$ obtenu
en associant à un ensemble $E$ l'ensemble $E$ muni de l'application constante
$E \to u(U)$ de valeur $x$.
Alors $j_x$ est un foncteur cocontinu pleinement fidèle qui possède un
adjoint à droite continu
$v_x : (\psi : E \to u(U)) \mapsto \psi^{-1}(\{x\})$.
Remarquons que $j_{u(U)} \circ j_x = \text{id}_\mathcal{S}$ et que
$v_x \circ u' = v$.
Ces observations impliquent le diagramme commutatif
de topos suivant :
$$
\xymatrix{
\Sh(\mathcal{S}) \ar[rd]^a \ar[rdd]_q
\ar `r[rrr] `d[dd]^p [rrdd] & & & \\
& \Sh(\mathcal{S}/u(U)) \ar[r]_-{j_{u(U)}} \ar[d]^{p'} &
\Sh(\mathcal{S}) \ar[d]^p & \\
& \Sh(\mathcal{C}/U) \ar[r]^{j_U} &
\Sh(\mathcal{C}) &
}
$$
En effet :
\begin{enumerate}
\item Le morphisme
$a : \Sh(\mathcal{S}) \to \Sh(\mathcal{S}/u(U))$
est le morphisme de topos associé au foncteur cocontinu
$j_x$, lequel est égal au morphisme associé au foncteur continu
$v_x$ ; voir le
Lemme \ref{sites-lemma-cocontinuous-morphism-topoi}
et la
section \ref{sites-section-cocontinuous-adjoint}.
\item Le composé $p \circ j_{u(U)} \circ a = p$ puisque
$j_{u(U)} \circ j_x = \text{id}_\mathcal{S}$.
\item Le composé $p' \circ a$ donne un morphisme de topos.
De plus, c'est le morphisme de topos associé au foncteur continu
$v_x \circ u' = v$. Ainsi, $v$ définit bien un point $q$ de
$\mathcal{C}/U$, qui s'insère par construction dans le diagramme ci-dessus.
\end{enumerate}
Cela achève la démonstration du lemme.
\end{proof}

\begin{lemma}
\label{sites-lemma-points-above-point}
Soient $\mathcal{C}$, $p$, $u$, $U$ comme dans le
Lemme \ref{sites-lemma-point-localize}.
La construction du
Lemme \ref{sites-lemma-point-localize}
donne une correspondance bijective entre les points $q$
de $\mathcal{C}/U$ situés au-dessus de $p$ et les éléments $x$ de $u(U)$.
\end{lemma}

\begin{proof}
Soit $q$ un point de $\mathcal{C}/U$ donné par le foncteur
$v : \mathcal{C}/U \to \textit{Ensembles}$ tel que $j_U \circ q = p$
comme morphismes de topos. Rappelons que $u(V) = p^{-1}(h_V^\#)$ pour tout
objet $V$ de $\mathcal{C}$ ; voir le
Lemme \ref{sites-lemma-point-site-topos}.
De même, $v(V/U) = q^{-1}(h_{V/U}^\#)$ pour tout objet
$V/U$ de $\mathcal{C}/U$.
Considérons les deux diagrammes suivants :
$$
\vcenter{
\xymatrix{
\Mor_{\mathcal{C}/U}(W/U, V/U) \ar[r] \ar[d] &
\Mor_\mathcal{C}(W, V) \ar[d] \\
\Mor_{\mathcal{C}/U}(W/U, U/U) \ar[r] &
\Mor_\mathcal{C}(W, U)
}
}
\quad
\vcenter{
\xymatrix{
h_{V/U}^\# \ar[r] \ar[d] & j_U^{-1}(h_V^\#) \ar[d] \\
h_{U/U}^\# \ar[r] & j_U^{-1}(h_U^\#)
}
}
$$
Le diagramme de droite est le faisceau associé au diagramme des
préfaisceaux sur $\mathcal{C}/U$ qui envoie $W/U$ sur le diagramme de gauche
d'ensembles. (Il y a ici un petit point technique :
nous avons $(j_U^{-1}h_V)^\# = j_U^{-1}(h_V^\#)$, et de même pour $h_U$ ; voir le
Lemme \ref{sites-lemma-technical-pu}.)
Remarquons que le diagramme de gauche d'ensembles est cartésien.
Comme le passage au faisceau associé est exact
(Lemme \ref{sites-lemma-sheafification-exact}),
nous en concluons que le diagramme de droite est cartésien.

\medskip\noindent
Appliquons le foncteur exact $q^{-1}$ au diagramme de droite
pour obtenir un diagramme cartésien
$$
\xymatrix{
v(V/U) \ar[r] \ar[d] & u(V) \ar[d] \\
v(U/U) \ar[r] & u(U)
}
$$
d'ensembles. Nous avons utilisé ici
$q^{-1} \circ j^{-1} = p^{-1}$. Puisque $U/U$ est un objet final de
$\mathcal{C}/U$, nous voyons que $v(U/U)$ est un singleton. L'image de
$v(U/U)$ dans $u(U)$ est donc un élément $x$, et l'application horizontale
supérieure donne une bijection
$v(V/U) \to  \{y \in u(V) \mid y \mapsto x\text{ dans }u(U)\}$,
comme voulu.
\end{proof}

\begin{lemma}
\label{sites-lemma-stalk-j-shriek}
Soit $\mathcal{C}$ un site. Soit $p$ un point de $\mathcal{C}$ donné par
$u : \mathcal{C} \to \textit{Ensembles}$. Soit $U$ un objet de $\mathcal{C}$.
Pour tout faisceau $\mathcal{G}$ sur $\mathcal{C}/U$, nous avons
$$
(j_{U!}\mathcal{G})_p =
\coprod\nolimits_q \mathcal{G}_q
$$
où la somme disjointe porte sur les points $q$ de $\mathcal{C}/U$
associés aux éléments $x \in u(U)$ comme dans le
Lemme \ref{sites-lemma-point-localize}.
\end{lemma}

\begin{proof}
Nous utilisons la description de $j_{U!}\mathcal{G}$ comme le faisceau associé
au préfaisceau
$V \mapsto
\coprod\nolimits_{\varphi \in \Mor_\mathcal{C}(V, U)}
\mathcal{G}(V/_\varphi U)$
du
Lemme \ref{sites-lemma-describe-j-shriek}.
De plus, la fibre de $j_{U!}\mathcal{G}$ en $p$ est égale à la
fibre de ce préfaisceau ; voir le
Lemme \ref{sites-lemma-point-pushforward-sheaf}.
Nous voyons donc que
$$
(j_{U!}\mathcal{G})_p =
\colim_{(V, y)} \coprod\nolimits_{\varphi : V \to U} \mathcal{G}(V/_\varphi U)
$$
À chaque élément $(V, y, \varphi, s)$ de cette limite inductive, nous pouvons associer
$x = u(\varphi)(y) \in u(U)$. Nous obtenons donc
$$
(j_{U!}\mathcal{G})_p =
\coprod\nolimits_{x \in u(U)}
\colim_{(\varphi : V \to U, y), \ u(\varphi)(y) = x} \mathcal{G}(V/_\varphi U).
$$
Cette expression est égale à celle du lemme par notre
construction des points $q$.
\end{proof}

\begin{remark}
\label{sites-remark-not-pushforward}
Attention : le résultat du
Lemme \ref{sites-lemma-stalk-j-shriek}
n'a pas d'analogue pour $j_{U, *}$.
\end{remark}




\section{2-flèches de topos}
\label{sites-section-2-category}

\noindent
Voici une brève section consacrée à la notion de $2$-flèche
de topos.

\begin{definition}
\label{sites-definition-2-morphism-topoi}
Soient $f, g : \Sh(\mathcal{C}) \to \Sh(\mathcal{D})$
deux morphismes de topos. Une {\it 2-flèche de $f$ vers $g$}
est donné par une transformation de foncteurs $t : f_* \to g_*$.
\end{definition}

\noindent
Nous représentons parfois $t$ par le diagramme suivant :
$$
\xymatrix{
\Sh(\mathcal{C})
\rrtwocell^f_g{t}
&
&
\Sh(\mathcal{D})
}
$$
Puisque $f^{-1}$ est adjoint à gauche de $f_*$ et que
$g^{-1}$ est adjoint à gauche de $g_*$, nous voyons que $t$ induit aussi
une transformation de foncteurs
$t : g^{-1} \to f^{-1}$ (notée généralement du même symbole),
caractérisée de manière unique par la commutativité du diagramme
$$
\xymatrix{
\Mor_{\Sh(\mathcal{D})}(\mathcal{G}, f_*\mathcal{F})
\ar[d]_{t \circ -} \ar@{=}[r] &
\Mor_{\Sh(\mathcal{C})}(f^{-1}\mathcal{G}, \mathcal{F})
\ar[d]^{- \circ t}
\\
\Mor_{\Sh(\mathcal{D})}(\mathcal{G}, g_*\mathcal{F})
\ar@{=}[r] &
\Mor_{\Sh(\mathcal{C})}(g^{-1}\mathcal{G}, \mathcal{F})
}
$$
En raison de difficultés ensemblistes (voir la
Remarque \ref{sites-remark-morphism-topoi-big}),
nous n'obtenons pas une 2-catégorie de topos. Nous pouvons néanmoins définir
les compositions horizontale et verticale, et montrer que les axiomes d'une
2-catégorie stricte énumérés dans
Catégories, section \ref{categories-section-2-categories},
sont satisfaits.
En effet, la composition verticale des 2-flèches est claire (il suffit de composer
les transformations de foncteurs), la composition des 1-morphismes a été
définie dans la
Définition \ref{sites-definition-topos},
et la composition horizontale de
$$
\xymatrix{
\Sh(\mathcal{C})
\rtwocell^f_g{t}
&
\Sh(\mathcal{D})
\rtwocell^{f'}_{g'}{s}
&
\Sh(\mathcal{E})
}
$$
est définie par la transformation de foncteurs $s \star t$
introduite dans
Catégories, Définition \ref{categories-definition-horizontal-composition}.
Explicitement, $s \star t$ est donnée par
$$
\xymatrix{
f'_*f_*\mathcal{F}
\ar[r]^{f'_*t} &
f'_*g_*\mathcal{F}
\ar[r]^s &
g'_*g_*\mathcal{F}
}
\quad\text{ou}\quad
\xymatrix{
f'_*f_*\mathcal{F}
\ar[r]^s &
g'_*f_*\mathcal{F}
\ar[r]^{g'_*t} &
g'_*g_*\mathcal{F}
}
$$
(ces applications sont égales). Comme ces définitions coïncident avec celles de
Catégories, section \ref{categories-section-formal-cat-cat},
il résulte de
Catégories, Lemme \ref{categories-lemma-properties-2-cat-cats},
que les axiomes d'une 2-catégorie stricte sont satisfaits avec ces définitions.




\section{Morphismes entre points}
\label{sites-section-morphisms-points}

\begin{lemma}
\label{sites-lemma-maps-u-points}
Soit $\mathcal{C}$ un site.
Soient $u, u' : \mathcal{C} \to \textit{Ensembles}$ deux
foncteurs, et soit $t : u' \to u$ une transformation de foncteurs.
Nous obtenons alors une transformation canonique de foncteurs fibres
$t_{stalk} : \mathcal{F}_{p'} \to \mathcal{F}_p$
qui coïncide avec $t$ par les identifications du
Lemme \ref{sites-lemma-points-recover}.
\end{lemma}

\begin{proof}
Omis.
\end{proof}

\begin{definition}
\label{sites-definition-morphism-points}
Soit $\mathcal{C}$ un site. Soient $p, p'$ des points de $\mathcal{C}$
donnés par des foncteurs $u, u' : \mathcal{C} \to \textit{Ensembles}$.
Un {\it morphisme $f : p \to p'$} est donné par une transformation de
foncteurs
$$
f_u : u' \to u.
$$
\end{definition}

\noindent
Remarquons que la transformation de foncteurs va dans le sens opposé.
Cela s'explique, comme nous le verrons plus loin, en considérant
le morphisme $f$ comme une sorte de $2$-flèche, représentée par le diagramme
suivant :
$$
\xymatrix{
\textit{Ensembles}
=
\Sh(pt)
\rrtwocell^p_{p'}{f}
&
&
\Sh(\mathcal{C})
}
$$
En effet, nous verrons plus loin que $f_u$ induit une transformation
canonique de foncteurs $p_* \to p'_*$ entre les
constructions de faisceaux gratte-ciel.

\medskip\noindent
Cette notion est assez importante et mérite ici un traitement plus complet.
Liste des compléments souhaitables :
\begin{enumerate}
\item Décrire les automorphismes du point de $\mathcal{T}_G$
décrit dans l'Exemple \ref{sites-example-point-G-sets}.
\item Décrire $\Mor(p, p')$ au moyen de $\Mor(p_*, p'_*)$.
\item Spécialisation des points dans les espaces topologiques.
Montrer que, si $x' \in \overline{\{x\}}$ dans l'espace topologique
$X$, alors il existe un morphisme $p \to p'$, où $p$ (respectivement $p'$)
est le point de $X_{Zar}$ associé à $x$ (respectivement $x'$).
\end{enumerate}




\section{Sites ayant assez de points}
\label{sites-section-sites-enough-points}

\begin{definition}
\label{sites-definition-enough-points}
Soit $\mathcal{C}$ un site.
\begin{enumerate}
\item Une famille de points $\{p_i\}_{i\in I}$ est dite {\it conservative}
si toute application de faisceaux $\phi : \mathcal{F} \to \mathcal{G}$
qui est un isomorphisme sur toutes les fibres $\mathcal{F}_{p_i}
\to \mathcal{G}_{p_i}$ est un isomorphisme.
\item On dit que $\mathcal{C}$ {\it a assez de points}
s'il existe une famille conservative de points.
\end{enumerate}
\end{definition}

\noindent
On peut alors vérifier l'« exactitude » sur les fibres.

\begin{lemma}
\label{sites-lemma-exactness-stalks}
Soit $\mathcal{C}$ un site, et soit $\{p_i\}_{i\in I}$ une famille conservative
de points. Alors :
\begin{enumerate}
\item Pour toute application de faisceaux $\varphi : \mathcal{F} \to \mathcal{G}$,
l'assertion « $\forall i, \varphi_{p_i}$ est injective » implique que $\varphi$ est injective.
\item Pour toute application de faisceaux $\varphi : \mathcal{F} \to \mathcal{G}$,
l'assertion « $\forall i, \varphi_{p_i}$ est surjective » implique que $\varphi$ est surjective.
\item Pour toute paire d'applications de faisceaux
$\varphi_1, \varphi_2 : \mathcal{F} \to \mathcal{G}$,
l'égalité $\forall i, \varphi_{1, p_i} = \varphi_{2, p_i}$
implique $\varphi_1 = \varphi_2$.
\item Soient un diagramme fini $\mathcal{G} : \mathcal{J}
\to \Sh(\mathcal{C})$, un faisceau $\mathcal{F}$ et des morphismes
$q_j : \mathcal{F} \to \mathcal{G}_j$. Alors $(\mathcal{F}, q_j)$
est une limite du diagramme si et seulement si, pour tout $i$, la fibre
$(\mathcal{F}_{p_i}, (q_j)_{p_i})$ en est une.
\item Soient un diagramme fini $\mathcal{F} : \mathcal{J}
\to \Sh(\mathcal{C})$, un faisceau $\mathcal{G}$ et des morphismes
$e_j : \mathcal{F}_j \to \mathcal{G}$. Alors $(\mathcal{G}, e_j)$
est une limite inductive du diagramme si et seulement si, pour tout $i$, la fibre
$(\mathcal{G}_{p_i}, (e_j)_{p_i})$ en est une.
\end{enumerate}
\end{lemma}

\begin{proof}
Nous utiliserons à plusieurs reprises le fait que tous les foncteurs fibres commutent
aux limites et limites inductives finies, donc aux produits, produits fibrés,
etc. Nous utiliserons aussi le fait que les applications injectives sont les
monomorphismes, et que les applications surjectives sont les épimorphismes.
Une application de faisceaux $\varphi : \mathcal{F} \to \mathcal{G}$
est injective si et seulement si
$\mathcal{F} \to \mathcal{F} \times_\mathcal{G} \mathcal{F}$
est un isomorphisme. D'où (1).
De même, $\varphi : \mathcal{F} \to \mathcal{G}$
est surjective si et seulement si
$\mathcal{G} \amalg_\mathcal{F} \mathcal{G} \to \mathcal{G}$
est un isomorphisme. D'où (2).
Les applications $a, b : \mathcal{F} \to \mathcal{G}$
sont égales si et seulement si
$\text{Égalisateur}(a, b : \mathcal{F} \to \mathcal{G}) \to \mathcal{F}$
est un isomorphisme\footnote{L'égalisateur s'exprime
au moyen de produits et de produits fibrés ; voir la démonstration de
Catégories, Lemme \ref{categories-lemma-finite-limits-exist}.}. D'où (3).
Les assertions (4) et (5) résultent immédiatement des définitions
et des remarques faites au début de cette démonstration.
\end{proof}

\begin{lemma}
\label{sites-lemma-enough}
Soit $\mathcal{C}$ un site, et soit $\{(p_i, u_i)\}_{i\in I}$ une
famille de points. Cette famille est conservative si et seulement si, pour tout
faisceau $\mathcal{F}$, tout $U\in \Ob(\mathcal{C})$ et toute
paire de sections distinctes $s, s' \in \mathcal{F}(U)$, $s \not = s'$, il
existe $i$ et $x\in u_i(U)$ tels que les triplets
$(U, x, s)$ et $(U, x, s')$ définissent des éléments distincts de
$\mathcal{F}_{p_i}$.
\end{lemma}

\begin{proof}
Supposons la famille conservative, et soient $\mathcal{F}$, $U$ et
$s, s'$ comme dans le lemme. Les sections $s$, $s'$ définissent des applications
$a, a' : (h_U)^\# \to \mathcal{F}$ qui sont distinctes. D'après le Lemme
\ref{sites-lemma-exactness-stalks}, il existe donc $i$ tel que $a_{p_i}
\not = a'_{p_i}$. Rappelons que $(h_U)^\#_{p_i} = u_i(U)$, d'après les
Lemmes \ref{sites-lemma-points-recover} et
\ref{sites-lemma-point-pushforward-sheaf}.
Il existe donc $x \in u_i(U)$ tel que $a_{p_i}(x)
\not = a'_{p_i}(x)$ dans $\mathcal{F}_{p_i}$.
En déroulant les définitions, on voit que $(U, x, s)$
et $(U, x, s')$ sont comme dans l'énoncé du lemme.

\medskip\noindent
Pour démontrer la réciproque, supposons satisfaite la condition d'existence de
points donnée dans le lemme. Soit $\phi : \mathcal{F} \to \mathcal{G}$
une application de faisceaux qui est un isomorphisme sur toutes les fibres.
Nous devons montrer que $\phi$ est à la fois
injective et surjective ; voir le Lemme \ref{sites-lemma-mono-epi-sheaves}.
L'injectivité résulte immédiatement de l'hypothèse.
Pour montrer la surjectivité, il nous faut montrer que
$\mathcal{G} \amalg_\mathcal{F} \mathcal{G} \to \mathcal{G}$
est un isomorphisme
(Catégories, Lemme \ref{categories-lemma-characterize-mono-epi}).
Comme cette application est manifestement surjective, il suffit d'en vérifier l'injectivité,
qui résulte du fait que $\mathcal{G} \amalg_\mathcal{F} \mathcal{G} \to \mathcal{G}$
est injective sur toutes les fibres par hypothèse.
\end{proof}

\noindent
Dans le lemme suivant, les points $q_{i, x}$ sont exactement tous les points
de $\mathcal{C}/U$ situés au-dessus du point $p_i$, d'après le
Lemme \ref{sites-lemma-points-above-point}.

\begin{lemma}
\label{sites-lemma-localize-enough}
Soit $\mathcal{C}$ un site. Soit $U$ un objet de $\mathcal{C}$.
Soit $\{(p_i, u_i)\}_{i\in I}$ une famille de points de $\mathcal{C}$.
Pour $x \in u_i(U)$, soit $q_{i, x}$ le point de $\mathcal{C}/U$
construit dans le
Lemme \ref{sites-lemma-point-localize}.
Si $\{p_i\}$ est une famille conservative de points, alors
$\{q_{i, x}\}_{i \in I, x \in u_i(U)}$ est une famille conservative
de points de $\mathcal{C}/U$.
En particulier, si $\mathcal{C}$ a assez de points, il en va de même
de toute localisation $\mathcal{C}/U$.
\end{lemma}

\begin{proof}
Nous savons que $j_{U!}$ induit une équivalence
$j_{U!} : \Sh(\mathcal{C}/U) \to \Sh(\mathcal{C})/h_U^\#$ ; voir le
Lemme \ref{sites-lemma-essential-image-j-shriek}.
De plus, nous savons que
$(j_{U!}\mathcal{G})_{p_i} = \coprod_x \mathcal{G}_{q_{i, x}}$ ; voir le
Lemme \ref{sites-lemma-stalk-j-shriek}.
Le résultat en découle donc formellement.
\end{proof}

\noindent
Le lemme suivant affirme que l'existence de points peut être vérifiée
localement sur le site.

\begin{lemma}
\label{sites-lemma-enough-points-local}
Soit $\mathcal{C}$ un site. Soit $\{U_i\}_{i \in I}$ une famille
d'objets de $\mathcal{C}$. Supposons que
\begin{enumerate}
\item $\coprod h_{U_i}^\# \to *$ soit une application surjective de faisceaux, et
\item chaque localisation $\mathcal{C}/U_i$ ait assez de points.
\end{enumerate}
Alors $\mathcal{C}$ a assez de points.
\end{lemma}

\begin{proof}
Pour tout $i \in I$, soit $\{p_j\}_{j \in J_i}$ une famille conservative
de points de $\mathcal{C}/U_i$. Pour $j \in J_i$, notons
$q_j : \Sh(pt) \to \Sh(\mathcal{C})$ le composé
de $p_j$ avec le morphisme de localisation
$\Sh(\mathcal{C}/U_i) \to \Sh(\mathcal{C})$.
Alors $q_j$ est un point ; voir le
Lemme \ref{sites-lemma-point-morphism-topoi}.
Nous affirmons que la famille de points $\{q_j\}_{j \in \coprod J_i}$
est conservative.
Soit en effet $\mathcal{F} \to \mathcal{G}$ une application de faisceaux
sur $\mathcal{C}$ telle que $\mathcal{F}_{q_j} \to \mathcal{G}_{q_j}$
soit un isomorphisme pour tout $j \in \coprod J_i$.
Soit $W$ un objet de $\mathcal{C}$.
D'après l'hypothèse (1), il existe un recouvrement $\{W_a \to W\}$ et
des morphismes $W_a \to U_{i(a)}$.
Comme $(\mathcal{F}|_{\mathcal{C}/U_{i(a)}})_{p_j} = \mathcal{F}_{q_j}$
et $(\mathcal{G}|_{\mathcal{C}/U_{i(a)}})_{p_j} = \mathcal{G}_{q_j}$ d'après le
Lemme \ref{sites-lemma-point-morphism-topoi},
nous voyons que
$\mathcal{F}|_{U_{i(a)}} \to \mathcal{G}|_{U_{i(a)}}$ est un isomorphisme,
puisque la famille de points $\{p_j\}_{j \in J_{i(a)}}$ est conservative.
Ainsi, $\mathcal{F}(W_a) \to \mathcal{G}(W_a)$ est bijective pour tout $a$.
De même, $\mathcal{F}(W_a \times_W W_b) \to \mathcal{G}(W_a \times_W W_b)$
est bijective pour tous $a, b$.
La condition de faisceau montre alors que
$\mathcal{F}(W) \to \mathcal{G}(W)$ est bijective, c'est-à-dire que
$\mathcal{F} \to \mathcal{G}$ est un isomorphisme.
\end{proof}

\begin{lemma}
\label{sites-lemma-check-morphism-sites}
Soit $u : \mathcal{C} \to \mathcal{D}$ un foncteur continu de sites.
Soit $\{(q_i, v_i)\}_{i\in I}$ une famille conservative de points de
$\mathcal{D}$. Si chaque foncteur $u_i = v_i \circ u$ définit
un point de $\mathcal{C}$,
alors $u$ définit un morphisme de sites $f : \mathcal{D} \to \mathcal{C}$.
\end{lemma}

\begin{proof}
Notons $p_i$ le foncteur fibre (\ref{sites-equation-stalk}) sur
$\textit{PSh}(\mathcal{C})$ correspondant au foncteur $u_i$. Nous avons
$$
(f^{-1}\mathcal{F})_{q_i} =
(u_s\mathcal{F})_{q_i} =
(u_p\mathcal{F})_{q_i} =
\mathcal{F}_{p_i}
$$
La première égalité tient à $f^{-1} = u_s$, la deuxième au
Lemme \ref{sites-lemma-point-pushforward-sheaf}, et la troisième
au Lemme \ref{sites-lemma-point-functor}.
Ainsi, si $p_i$ est un point, alors prendre l'image inverse par $f$
puis les fibres en $q_i$ est un foncteur exact.
Comme la famille de points $\{q_i\}$ est conservative, cela
implique que $f^{-1}$ est un foncteur exact ; nous voyons donc
que $f$ est un morphisme de sites par la Définition \ref{sites-definition-morphism-sites}.
\end{proof}





\section{Critère d'existence de points}
\label{sites-section-criterion-points}

\noindent
Cette section correspond à l'appendice de Deligne à \cite[Expos\'e VI]{SGA4}.
Elle lui est en fait presque littéralement identique.

\medskip\noindent
Soit $\mathcal{C}$ un site.
Supposons que $(I, \geq)$ soit un ensemble ordonné filtrant,
et que $(U_i, f_{ii'})$ soit un système projectif sur $I$ ; voir
Catégories, Définition \ref{categories-definition-system-over-poset}.
Aux données $(I, \geq, U_i, f_{ii'})$, nous associons le foncteur
$$
u : \mathcal{C} \longrightarrow \textit{Ensembles}, \quad
u(V) = \colim_i \Mor_\mathcal{C}(U_i , V)
$$
Soit $\mathcal{F} \mapsto \mathcal{F}_p$ le foncteur fibre
associé à $u$ comme dans la section \ref{sites-section-points}.
Il résulte directement de la définition que
$$
\mathcal{F}_p = \colim_i \mathcal{F}(U_i)
$$
dans ce cas particulier.
Remarquons que $u$ commute à toutes les limites finies (c'est-à-dire à celles qui
sont représentables dans $\mathcal{C}$), car chacun des foncteurs
$V \mapsto \Mor_\mathcal{C}(U_i , V)$ le fait ; voir
Catégories, Lemme \ref{categories-lemma-directed-commutes}.

\medskip\noindent
On dit qu'un système $(I, \geq, U_i, f_{ii'})$
est un {\it raffinement} de $(J, \geq, V_j, g_{jj'})$ si
$J \subset I$, si l'ordre sur $J$ est celui induit par l'ordre de $I$,
et si $V_j = U_j$, $g_{jj'} = f_{jj'}$ (autrement dit, si le système projectif
sur $J$ est induit par celui sur $I$). Soit $u$ le foncteur
associé à $(I, \geq, U_i, f_{ii'})$, et soit $u'$ le
foncteur associé à $(J, \geq, V_j, g_{jj'})$.
Cela induit une transformation de foncteurs
$$
u' \longrightarrow u
$$
simplement parce que les limites inductives définissant $u'$ portent sur un sous-système
des systèmes intervenant dans les limites inductives définissant $u$.
En particulier, nous obtenons une transformation associée des
foncteurs fibres $\mathcal{F}_{p'} \to \mathcal{F}_p$ ;
voir le Lemme \ref{sites-lemma-maps-u-points}.

\begin{lemma}
\label{sites-lemma-refine}
Soit $\mathcal{C}$ un site.
Soit $(J, \geq, V_j, g_{jj'})$ un système comme ci-dessus, muni de la
paire de foncteurs associée $(u', p')$.
Soit $\mathcal{F}$ un faisceau sur $\mathcal{C}$.
Soient $s, s' \in \mathcal{F}_{p'}$ des éléments distincts.
Soit $\{W_k \to W\}$ un recouvrement fini de $\mathcal{C}$.
Soit $f \in u'(W)$.
Il existe un raffinement $(I, \geq, U_i, f_{ii'})$
de $(J, \geq, V_j, g_{jj'})$ tel que les images de $s, s'$ dans
$\mathcal{F}_p$ soient distinctes et que
l'image de $f$ dans $u(W)$ appartienne à l'image de l'un des
$u(W_k)$.
\end{lemma}

\begin{proof}
Il existe $j_0 \in J$ tel que $f$ soit défini par $f' : V_{j_0} \to W$.
Pour $j \geq j_0$, posons $V_{j, k} = V_j \times_{f'\circ f_{j j_0}, W} W_k$.
Alors $\{V_{j, k} \to V_j\}$ est un recouvrement fini du site
$\mathcal{C}$. Par conséquent,
$\mathcal{F}(V_j) \subset \prod_k \mathcal{F}(V_{j, k})$.
En appliquant encore Catégories, Lemme \ref{categories-lemma-directed-commutes},
nous voyons que
$$
\mathcal{F}_{p'} =
\colim_j \mathcal{F}(V_j)
\longrightarrow
\prod\nolimits_k \colim_j \mathcal{F}(V_{j, k})
$$
est injective. Il existe donc $k$ tel que $s$ et $s'$ aient des images
distinctes dans $\colim_j \mathcal{F}(V_{j, k})$.
Posons $J_0 = \{j \in J, j \geq j_0\}$ et $I = J \amalg J_0$.
Ordonnons $I$ de sorte qu'aucun élément du second terme
ne soit inférieur à un élément du premier terme, et utilisons par ailleurs
l'ordre sur $J$. Si $j \in I$ appartient au premier
terme, nous prenons $V_j$ ; si $j \in I$ appartient au second terme,
nous prenons $V_{j, k}$. Nous omettons la définition
des morphismes de transition du système projectif. Il résulte de ce qui précède
que les images de $s, s'$ dans $\mathcal{F}_p$ sont distinctes.
En outre, la restriction de $f'$ à $V_{j, k}$ se factorise
par $W_k$ par construction.
\end{proof}

\begin{lemma}
\label{sites-lemma-refine-all-at-once}
Soit $\mathcal{C}$ un site.
Soit $(J, \geq, V_j, g_{jj'})$ un système comme ci-dessus, muni de la
paire de foncteurs associée $(u', p')$.
Soit $\mathcal{F}$ un faisceau sur $\mathcal{C}$.
Soient $s, s' \in \mathcal{F}_{p'}$ des éléments distincts.
Il existe un raffinement $(I, \geq, U_i, f_{ii'})$
de $(J, \geq, V_j, g_{jj'})$ tel que les images de $s, s'$ dans
$\mathcal{F}_p$ soient distinctes et que,
pour tout recouvrement fini $\{W_k \to W\}$ du site
$\mathcal{C}$ et tout $f \in u'(W)$, l'image de $f$ dans $u(W)$
appartienne à l'image de l'un des $u(W_k)$.
\end{lemma}

\begin{proof}
Soit $E$ l'ensemble des couples $(\{W_k \to W\}, f\in u'(W))$.
Considérons les couples $(E' \subset E, (I, \geq, U_i, f_{ii'}))$
tels que
\begin{enumerate}
\item $(I, \geq, U_i, g_{ii'})$ soit un raffinement de $(J, \geq, V_j, g_{jj'})$,
\item les images de $s, s'$ dans $\mathcal{F}_p$ soient distinctes, et
\item pour tout couple $(\{W_k \to W\}, f\in u'(W)) \in E'$, l'image de
$f$ dans $u(W)$ appartienne à l'image de l'un des $u(W_k)$.
\end{enumerate}
Ordonnons ces couples par inclusion sur le premier facteur et
par raffinement sur le second. Notons $\mathcal{S}$ la classe
de tous les couples $(E' \subset E, (I, \geq, U_i, f_{ii'}))$ ci-dessus.
Nous affirmons que l'hypothèse du lemme de Zorn est satisfaite pour $\mathcal{S}$.
Supposons en effet que $(E'_a, (I_a, \geq, U_i, f_{ii'}))_{a \in A}$
soit une partie totalement ordonnée de $\mathcal{S}$. Nous pouvons définir
$E' = \bigcup_{a \in A} E'_a$ et poser $I = \bigcup_{a \in A} I_a$.
Nous affirmons que le couple correspondant
$(E' , (I, \geq, U_i, f_{ii'}))$ est un élément de $\mathcal{S}$.
Les conditions (1) et (3) sont claires. Pour la condition (2), remarquons
que
$$
u = \colim_{a \in A} u_a
\text{ et, par conséquent, }
\mathcal{F}_p = \colim_{a \in A} \mathcal{F}_{p_a}
$$
Le fait que les images de $s, s'$ dans cette fibre soient distinctes résulte
de la description d'une limite inductive filtrante d'ensembles ; voir
Catégories, section \ref{categories-section-directed-colimits}.
Nous écrirons simplement
$(E', (I, \ldots)) = \bigcup_{a \in A}(E'_a, (I_a, \ldots))$
dans cette situation.

\medskip\noindent
Le lemme de Zorn s'appliquerait donc si $\mathcal{S}$ était un ensemble,
et, combiné au Lemme \ref{sites-lemma-refine} ci-dessus, il démontrerait aisément
le lemme. Mais $\mathcal{S}$ est une classe. Pour contourner cette difficulté,
choisissons un bon ordre sur $E$.
Pour $e \in E$, posons $E'_e = \{e' \in E \mid e' \leq e\}$.
Par récurrence transfinie, construisons des couples
$(E'_e, (I_e, \ldots)) \in \mathcal{S}$ tels que
$e_1 \leq e_2 \Rightarrow (E'_{e_1}, (I_{e_1}, \ldots))
\leq (E'_{e_2}, (I_{e_2}, \ldots))$.
Soit $e \in E$, disons $e = (\{W_k \to W\}, f\in u'(W))$.
Si $e$ possède un prédécesseur $e - 1$, prenons
$(I_e, \ldots)$ comme un raffinement de $(I_{e - 1}, \ldots)$
comme dans le Lemme \ref{sites-lemma-refine}, relativement au système
$e = (\{W_k \to W\}, f\in u'(W))$.
Si $e$ ne possède pas de prédécesseur, prenons
$(I_e, \ldots)$ comme un raffinement de $\bigcup_{e' < e} (I_{e'}, \ldots)$
relativement au système
$e = (\{W_k \to W\}, f\in u'(W))$.
Enfin, la réunion $\bigcup_{e \in E} I_e$ est une solution au
problème posé dans le lemme.
\end{proof}

\begin{proposition}
\label{sites-proposition-criterion-points}
\begin{reference}
\cite[Expos\'e VI, Appendice de Deligne, Proposition 9.0]{SGA4}
\end{reference}
Soit $\mathcal{C}$ un site. Supposons que
\begin{enumerate}
\item les limites finies existent dans $\mathcal{C}$, et que
\item tout recouvrement $\{U_i \to U\}_{i \in I}$
admette un raffinement par un recouvrement fini de $\mathcal{C}$.
\end{enumerate}
Alors $\mathcal{C}$ a assez de points.
\end{proposition}

\begin{proof}
Nous devons montrer que, pour tout faisceau
$\mathcal{F}$ sur $\mathcal{C}$, tout $U \in \Ob(\mathcal{C})$
et toutes sections distinctes $s, s' \in \mathcal{F}(U)$, il existe
un point $p$ tel que les images de $s, s'$ dans
$\mathcal{F}_p$ soient distinctes. Voir le Lemme \ref{sites-lemma-enough}.
Considérons le système $(J, \geq, V_j, g_{jj'})$
défini par $J = \{1\}$, $V_1 = U$, $g_{11} = \text{id}_U$.
Appliquons le Lemme \ref{sites-lemma-refine-all-at-once}.
Ce lemme fournit un système
$(I, \geq, U_i, f_{ii'})$ raffinant le nôtre de sorte que
$s_p \not = s'_p$ et que, de plus, pour tout
recouvrement fini $\{W_k \to W\}$ du site $\mathcal{C}$, l'application
$\coprod_k u(W_k) \to u(W)$ soit surjective.
Comme tout recouvrement de $\mathcal{C}$ admet un raffinement par
un recouvrement fini, nous en concluons que
$\coprod_k u(W_k) \to u(W)$ est surjective pour {\it tout}
recouvrement $\{W_k \to W\}$ du site $\mathcal{C}$.
Cela implique que $u = p$ est un point ; voir la
Proposition \ref{sites-proposition-point-limits} (ainsi que la discussion
au début de cette section, qui garantit que $u$
commute aux limites finies).
\end{proof}

\begin{lemma}
\label{sites-lemma-criterion-points}
Soit $\mathcal{C}$ un site. Soit $I$ un ensemble et, pour
$i \in I$, soit $U_i$ un objet de $\mathcal{C}$ tel que
\begin{enumerate}
\item $\coprod h_{U_i}$ se projette surjectivement sur
l'objet final de $\Sh(\mathcal{C})$, et
\item $\mathcal{C}/U_i$ satisfasse les hypothèses de la
Proposition \ref{sites-proposition-criterion-points}.
\end{enumerate}
Alors $\mathcal{C}$ a assez de points.
\end{lemma}

\begin{proof}
D'après l'hypothèse (2) et la proposition, $\mathcal{C}/U_i$ a assez de points.
Les points de $\mathcal{C}/U_i$ donnent des points de $\mathcal{C}$
par le procédé du Lemme \ref{sites-lemma-point-morphism-sites}.
Il suffit donc de montrer ceci : si $\phi : \mathcal{F} \to \mathcal{G}$
est une application de faisceaux sur $\mathcal{C}$ telle que $\phi|_{\mathcal{C}/U_i}$
soit un isomorphisme pour tout $i$, alors $\phi$ est un isomorphisme.
D'après l'hypothèse (1), pour tout objet $W$ de $\mathcal{C}$,
il existe un recouvrement $\{W_j \to W\}_{j \in J}$
tel que, pour $j \in J$, il existe $i \in I$ et un morphisme
$f_j : W_j \to U_i$. Les applications
$\mathcal{F}(W_j) \to \mathcal{G}(W_j)$
sont alors bijectives, de même que
$\mathcal{F}(W_j \times_W W_{j'}) \to \mathcal{G}(W_j \times_W W_{j'})$.
La condition de faisceau nous dit que $\mathcal{F}(W) \to \mathcal{G}(W)$
est bijective, comme voulu.
\end{proof}




\section{Objets faiblement contractiles}
\label{sites-section-w-contractible}

\noindent
Un objet d'un site est dit {\it faiblement contractile} s'il satisfait
aux conditions équivalentes du lemme suivant.

\begin{lemma}
\label{sites-lemma-w-contractible}
Soit $\mathcal{C}$ un site. Soit $U$ un objet de $\mathcal{C}$.
Les conditions suivantes sont équivalentes :
\begin{enumerate}
\item Pour toute famille couvrante $\{U_i \to U\}$, il existe un morphisme de
faisceaux $h_U^\# \to \coprod h_{U_i}^\#$ inverse à droite du morphisme de
faisceaux associé à $\coprod h_{U_i} \to h_U$.
\item Pour toute surjection de faisceaux d'ensembles $\mathcal{F} \to \mathcal{G}$,
l'application $\mathcal{F}(U) \to \mathcal{G}(U)$ est surjective.
\end{enumerate}
\end{lemma}

\begin{proof}
Supposons (1) et soit $\mathcal{F} \to \mathcal{G}$ un morphisme surjectif de
faisceaux d'ensembles. Pour $s \in \mathcal{G}(U)$, il existe une famille couvrante
$\{U_i \to U\}$ et des $t_i \in \mathcal{F}(U_i)$ dont l'image est
$s|_{U_i}$, voir la définition \ref{sites-definition-sheaves-injective-surjective}.
Considérons $t_i$ comme un morphisme $t_i : h_{U_i}^\# \to \mathcal{F}$ grâce à
(\ref{sites-equation-map-representable-into-sheaf}).
En précomposant alors $\coprod t_i : \coprod h_{U_i}^\# \to \mathcal{F}$
par le morphisme $h_U^\# \to \coprod h_{U_i}^\#$ fourni par (1),
nous obtenons une section $t \in \mathcal{F}(U)$ dont l'image est $s$.
Ainsi (2) est vérifiée.

\medskip\noindent
Supposons (2). Soit $\{U_i \to U\}$ une famille couvrante.
Alors $\coprod h_{U_i}^\# \to h_U^\#$ est surjectif
(lemme \ref{sites-lemma-covering-surjective-after-sheafification}).
Il existe donc par (2) une section $s$ de $\coprod h_{U_i}^\#$ dont l'image
est la section $\text{id}_U$ de $h_U^\#$. Cette section correspond à un morphisme
$h_U^\# \to \coprod h_{U_i}^\#$ qui est inverse à droite du morphisme de
faisceaux associé à $\coprod h_{U_i} \to h_U$, ce qui prouve (1).
\end{proof}

\begin{definition}
\label{sites-definition-w-contractible}
Soit $\mathcal{C}$ un site.
\begin{enumerate}
\item On dit qu'un objet $U$ de $\mathcal{C}$ est {\it faiblement contractile}
si les conditions équivalentes du lemme \ref{sites-lemma-w-contractible} sont vérifiées.
\item On dit qu'un site possède {\it suffisamment d'objets faiblement contractiles}
si tout objet $U$ de $\mathcal{C}$ admet une famille couvrante $\{U_i \to U\}$
telle que $U_i$ soit faiblement contractile pour tout $i$.
\item Plus généralement, si $P$ est une propriété des objets de $\mathcal{C}$,
on dit que $\mathcal{C}$ possède {\it suffisamment d'objets vérifiant $P$} si tout objet $U$ de
$\mathcal{C}$ admet une famille couvrante $\{U_i \to U\}$ telle que $U_i$ vérifie $P$
pour tout $i$.
\end{enumerate}
\end{definition}

\noindent
Le petit site \'etale de $\mathbf{A}^1_\mathbf{C}$ ne possède aucun
objet faiblement contractile. En revanche, le petit site
pro-\'etale de tout schéma possède suffisamment d'objets contractiles.







\section{Propriétés d'exactitude de l'image directe}
\label{sites-section-pushforward}

\noindent
Soit $f$ un morphisme de topos. En général, le foncteur $f_*$ n'est
qu'exact à gauche. On peut imposer à ce foncteur de nombreuses conditions
supplémentaires afin de distinguer certaines classes de morphismes de topos.
Nous les rassemblons ici et relevons quelques-unes de leurs implications logiques.
Certaines parties du lemme suivant sont de nature purement catégorique (c'est-à-dire
qu'elles ne dépendent pas de l'existence d'un morphisme de topos, l'existence d'une paire
de foncteurs adjoints suffit).

\begin{lemma}
\label{sites-lemma-exactness-properties}
Soit $f : \Sh(\mathcal{C}) \to \Sh(\mathcal{D})$ un
morphisme de topos. Considérons les propriétés suivantes (pour les faisceaux
d'ensembles) :
\begin{enumerate}
\item $f_*$ est fidèle,
\item $f_*$ est pleinement fidèle,
\item $f^{-1}f_*\mathcal{F} \to \mathcal{F}$ est surjectif pour
tout $\mathcal{F}$ de $\Sh(\mathcal{C})$,
\item $f_*$ transforme les surjections en surjections,
\item $f_*$ commute aux coégalisateurs,
\item $f_*$ commute aux sommes amalgamées,
\item $f^{-1}f_*\mathcal{F} \to \mathcal{F}$ est un isomorphisme pour
tout $\mathcal{F}$ de $\Sh(\mathcal{C})$,
\item $f_*$ reflète les injections,
\item $f_*$ reflète les surjections,
\item $f_*$ reflète les bijections, et
\item pour toute surjection $\mathcal{F} \to f^{-1}\mathcal{G}$, il
existe une surjection $\mathcal{G}' \to \mathcal{G}$ telle que
$f^{-1}\mathcal{G}' \to f^{-1}\mathcal{G}$ se factorise par
$\mathcal{F} \to f^{-1}\mathcal{G}$.
\end{enumerate}
On a alors les implications suivantes :
\begin{enumerate}
\item[(a)] (2) $\Rightarrow$ (1),
\item[(b)] (3) $\Rightarrow$ (1),
\item[(c)] (7) $\Rightarrow$ (1), (2), (3), (8), (9), (10).
\item[(d)] (3) $\Leftrightarrow$ (9),
\item[(e)] (6) $\Rightarrow$ (4) et (5) $\Rightarrow$ (4),
\item[(f)] (4) $\Leftrightarrow$ (11),
\item[(g)] (9) $\Rightarrow$ (8), (10), et
\item[(h)] (2) $\Leftrightarrow$ (7).
\end{enumerate}
Le diagramme est le suivant :
$$
\xymatrix{
(6) \ar@{=>}[rd] & & & & & (9) \ar@{=>}[r] \ar@{=>}[rd] & (8) \\
& (4) \ar@{<=>}[r] & (11) &
(2) \ar@{<=>}[r] &
(7) \ar@{=>}[ru] \ar@{=>}[rd] & & (10) \\
(5) \ar@{=>}[ur] & & & & & (3) \ar@{=>}[r] & (1)
}
$$
\end{lemma}

\begin{proof}
Démonstration de (a) : cela résulte immédiatement des définitions.

\medskip\noindent
Démonstration de (b). Supposons que $a, b : \mathcal{F} \to \mathcal{F}'$ soient des
morphismes de faisceaux sur $\mathcal{C}$. Si $f_*a = f_*b$, alors
$f^{-1}f_*a = f^{-1}f_*b$. Considérons le diagramme commutatif
$$
\xymatrix{
\mathcal{F} \ar@<-1ex>[r] \ar@<1ex>[r] & \mathcal{F}' \\
f^{-1}f_*\mathcal{F} \ar@<-1ex>[r] \ar@<1ex>[r] \ar[u] &
f^{-1}f_*\mathcal{F}' \ar[u]
}
$$
Si les deux flèches inférieures sont égales et les flèches verticales surjectives,
alors les deux flèches supérieures sont égales. D'où (b).

\medskip\noindent
Démonstration de (c). Soit $a : \mathcal{F} \to \mathcal{F}'$ un
morphisme de faisceaux sur $\mathcal{C}$. Considérons le diagramme commutatif
$$
\xymatrix{
\mathcal{F} \ar[r] & \mathcal{F}' \\
f^{-1}f_*\mathcal{F} \ar[r] \ar[u] &
f^{-1}f_*\mathcal{F}' \ar[u]
}
$$
Si (7) est vérifiée, les flèches verticales sont des isomorphismes.
Ainsi, si $f_*a$ est injectif (resp.\ surjectif, resp.\ bijectif),
alors la flèche inférieure est injective (resp.\ surjective, resp.\ bijective), et
il en va donc de même pour la flèche supérieure.
On voit ainsi que (7) implique (8), (9), (10). Il est clair que (7) implique (3).
Les implications (7) $\Rightarrow$ (2), (1) résultent de (a) et (h), que
nous verrons ci-dessous.

\medskip\noindent
Démonstration de (d). Supposons (3). Soit $a : \mathcal{F} \to \mathcal{F}'$
un morphisme de faisceaux sur $\mathcal{C}$ tel que $f_*a$ soit surjectif.
Comme $f^{-1}$ est exact, il en résulte que
$f^{-1}f_*a : f^{-1}f_*\mathcal{F} \to f^{-1}f_*\mathcal{F}'$
est surjectif. Avec (3), cela implique que $a$ est surjectif.
Ainsi (9) est vérifiée.
Supposons (9). Soit $\mathcal{F}$ un faisceau sur $\mathcal{C}$.
Il faut montrer que le morphisme $f^{-1}f_*\mathcal{F} \to \mathcal{F}$ est
surjectif. Il suffit de montrer que
$f_*f^{-1}f_*\mathcal{F} \to f_*\mathcal{F}$ est surjectif.
Or cela est vrai puisqu'il existe un morphisme canonique
$f_*\mathcal{F} \to f_*f^{-1}f_*\mathcal{F}$ qui est un inverse d'un côté de ce morphisme.

\medskip\noindent
Démonstration de (e). Nous utilisons
Catégories, lemme \ref{categories-lemma-characterize-mono-epi}
sans autre mention.
Si $\mathcal{F} \to \mathcal{F}'$ est surjectif, alors
$\mathcal{F}' \amalg_\mathcal{F} \mathcal{F}' \to \mathcal{F}'$
est un isomorphisme. Par conséquent, (6) implique que
$$
f_*\mathcal{F}' \amalg_{f_*\mathcal{F}} f_*\mathcal{F}' =
f_*(\mathcal{F}' \amalg_\mathcal{F} \mathcal{F}')
\longrightarrow
f_*\mathcal{F}'
$$
est également un isomorphisme. Cela implique à son tour que
$f_*\mathcal{F} \to f_*\mathcal{F}'$ est surjectif.
On voit donc que (6) implique (4). Si $\mathcal{F} \to \mathcal{F}'$
est surjectif, alors $\mathcal{F}'$ est le coégalisateur des deux
projections $\mathcal{F} \times_{\mathcal{F}'} \mathcal{F} \to \mathcal{F}$
d'après le lemme \ref{sites-lemma-coequalizer-surjection}.
Par conséquent, si (5) est vérifiée, alors $f_*\mathcal{F}'$ est le coégalisateur
des deux projections
$$
f_*(\mathcal{F} \times_{\mathcal{F}'} \mathcal{F}) =
f_*\mathcal{F} \times_{f_*\mathcal{F}'} f_*\mathcal{F}
\longrightarrow
f_*\mathcal{F}
$$
ce qui signifie clairement que $f_*\mathcal{F} \to f_*\mathcal{F}'$ est
surjectif. Ainsi (5) implique également (4).

\medskip\noindent
Démonstration de (f). Supposons (4). Soit $\mathcal{F} \to f^{-1}\mathcal{G}$
un morphisme surjectif de faisceaux sur $\mathcal{C}$. Par (4), on voit que
$f_*\mathcal{F} \to f_*f^{-1}\mathcal{G}$ est surjectif. Soit
$\mathcal{G}'$ le produit fibré
$$
\xymatrix{
f_*\mathcal{F} \ar[r] & f_*f^{-1}\mathcal{G} \\
\mathcal{G}' \ar[u] \ar[r] & \mathcal{G} \ar[u]
}
$$
de sorte que $\mathcal{G}' \to \mathcal{G}$ est lui aussi surjectif. Considérons
le diagramme commutatif
$$
\xymatrix{
\mathcal{F} \ar[r] & f^{-1}\mathcal{G} \\
f^{-1}f_*\mathcal{F} \ar[r] \ar[u] & f^{-1}f_*f^{-1}\mathcal{G} \ar[u] \\
f^{-1}\mathcal{G}' \ar[u] \ar[r] & f^{-1}\mathcal{G} \ar[u]
}
$$
qui donne le résultat voulu. Réciproquement, supposons (11).
Soit $a : \mathcal{F} \to \mathcal{F}'$ un morphisme surjectif de faisceaux
sur $\mathcal{C}$. Considérons le diagramme cartésien
$$
\xymatrix{
\mathcal{F} \ar[r] & \mathcal{F}' \\
\mathcal{F}'' \ar[u] \ar[r] & f^{-1}f_*\mathcal{F}' \ar[u]
}
$$
Comme la flèche horizontale inférieure est surjective, (11) permet de
trouver une surjection $\gamma : \mathcal{G}' \to f_*\mathcal{F}'$ telle
que $f^{-1}\gamma$ se factorise par $\mathcal{F}'' \to f^{-1}f_*\mathcal{F}'$ :
$$
\xymatrix{
& \mathcal{F} \ar[r] & \mathcal{F}' \\
f^{-1}\mathcal{G}' \ar[r] & \mathcal{F}'' \ar[u] \ar[r] &
f^{-1}f_*\mathcal{F}' \ar[u]
}
$$
En appliquant $f_*$, on obtient un diagramme commutatif
$$
\xymatrix{
& f_*\mathcal{F} \ar[r] & f_*\mathcal{F}' \\
f_*f^{-1}\mathcal{G}' \ar[r] & f_*\mathcal{F}'' \ar[u] \ar[r] &
f_*f^{-1}f_*\mathcal{F}' \ar[u] \\
\mathcal{G}' \ar[u] \ar[rr] & & f_*\mathcal{F}' \ar[u]
}
$$
ce qui prouve (4).

\medskip\noindent
Démonstration de (g). Supposons (9). Nous utilisons
Catégories, lemme \ref{categories-lemma-characterize-mono-epi}
sans autre mention. Soit $a : \mathcal{F} \to \mathcal{F}'$ un morphisme
de faisceaux sur $\mathcal{C}$ tel que $f_*a$ soit injectif. Cela signifie que
$f_*\mathcal{F} \to f_*\mathcal{F} \times_{f_*\mathcal{F}'} f_*\mathcal{F} =
f_*(\mathcal{F} \times_{\mathcal{F}'} \mathcal{F})$ est un isomorphisme.
Ainsi, par (9), on voit que
$\mathcal{F} \to \mathcal{F} \times_{\mathcal{F}'} \mathcal{F}$ est
surjectif, c'est-à-dire un isomorphisme. Ainsi $a$ est injectif, c'est-à-dire que (8) est vérifiée.
Puisque (10) est trivialement équivalente à (8) $+$ (9), cela achève la démonstration de (g).

\medskip\noindent
Démonstration de (h). C'est
Catégories, lemme \ref{categories-lemma-adjoint-fully-faithful}.
\end{proof}

\noindent
Voici une condition portant sur un morphisme de sites qui garantit que
le foncteur $f_*$ transforme les morphismes surjectifs en morphismes surjectifs.

\begin{lemma}
\label{sites-lemma-weaker}
Soit $f : \mathcal{D} \to \mathcal{C}$ un morphisme de sites associé au
foncteur continu $u : \mathcal{C} \to \mathcal{D}$.
Supposons que, pour tout objet $U$ de $\mathcal{C}$ et toute famille couvrante
$\{V_j \to u(U)\}$ de $\mathcal{D}$, il existe une famille couvrante $\{U_i \to U\}$
de $\mathcal{C}$ telle que le morphisme de faisceaux
$$
\coprod h_{u(U_i)}^\# \to h_{u(U)}^\#
$$
se factorise par le morphisme de faisceaux
$$
\coprod h_{V_j}^\# \to h_{u(U)}^\#.
$$
Alors $f_*$ transforme les morphismes surjectifs de faisceaux en morphismes surjectifs de
faisceaux.
\end{lemma}

\begin{proof}
Soit $a : \mathcal{F} \to \mathcal{G}$ un morphisme surjectif de faisceaux sur
$\mathcal{D}$. Soit $U$ un objet de $\mathcal{C}$ et soit
$s \in f_*\mathcal{G}(U) = \mathcal{G}(u(U))$. Par hypothèse, il existe
une famille couvrante $\{V_j \to u(U)\}$ et des sections $s_j \in \mathcal{F}(V_j)$
telles que $a(s_j) = s|_{V_j}$. On peut maintenant considérer les sections $s$,
$s_j$ et $a$ comme définissant un diagramme commutatif de morphismes de faisceaux
$$
\xymatrix{
\coprod h_{V_j}^\# \ar[r]_-{\coprod s_j} \ar[d] & \mathcal{F} \ar[d]^a \\
h_{u(U)}^\# \ar[r]^s & \mathcal{G}
}
$$
Par hypothèse, il existe une famille couvrante $\{U_i \to U\}$ qui permet
d'agrandir comme suit le diagramme commutatif ci-dessus :
$$
\xymatrix{
& \coprod h_{V_j}^\# \ar[r]_-{\coprod s_j} \ar[d] & \mathcal{F} \ar[d]^a \\
\coprod h_{u(U_i)}^\# \ar[r] \ar[ur] &
h_{u(U)}^\# \ar[r]^s & \mathcal{G}
}
$$
Comme $\mathcal{F}$ est un faisceau, le morphisme du coin inférieur gauche vers
le coin supérieur droit correspond à une famille de sections
$s_i \in \mathcal{F}(u(U_i))$, c'est-à-dire à des sections $s_i \in f_*\mathcal{F}(U_i)$.
La commutativité du diagramme implique que $a(s_i)$ est égal à la
restriction de $s$ à $U_i$. Autrement dit, nous avons montré que $f_*a$ est un
morphisme surjectif de faisceaux.
\end{proof}

\begin{example}
\label{sites-example-weaker-not-coequalizer}
Supposons que $f : \mathcal{D} \to \mathcal{C}$ satisfasse aux hypothèses du
lemme \ref{sites-lemma-weaker}.
En général, $f_*$ ne commute alors ni aux coégalisateurs
ni aux sommes amalgamées. En effet, supposons que $f$ soit le morphisme de sites associé
au morphisme d'espaces topologiques $X = \{1, 2\} \to Y = \{*\}$ (voir
l'exemple \ref{sites-example-continuous-map}),
où $Y$ est un espace réduit à un point et $X = \{1, 2\}$ est un espace discret
à deux points. Un faisceau $\mathcal{F}$ sur $X$ est donné par une paire
$(A_1, A_2)$ d'ensembles. Alors $f_*\mathcal{F}$ correspond à
l'ensemble $A_1 \times A_2$. Par conséquent, si
$a = (a_1, a_2), b = (b_1, b_2) : (A_1, A_2) \to (B_1, B_2)$
sont des morphismes de faisceaux sur $X$, alors le
coégalisateur de $a, b$ est $(C_1, C_2)$, où $C_i$ est le coégalisateur de
$a_i, b_i$, tandis que le coégalisateur de $f_*a, f_*b$ est le coégalisateur de
$$
a_1 \times a_2, b_1 \times b_2 :
A_1 \times A_2 \longrightarrow B_1 \times B_2
$$
qui est en général différent de $C_1 \times C_2$. En effet, si
$A_2 = \emptyset$, alors $A_1 \times A_2 = \emptyset$, et par conséquent le
coégalisateur des flèches affichées est $B_1 \times B_2$, mais en général
$C_1 \not = B_1$. Un exemple analogue vaut pour les sommes amalgamées.
\end{example}

\noindent
Le lemme suivant donne un critère assurant que, pour un morphisme de sites,
le foncteur $f_*$ reflète les injections et les surjections.
Remarquons que cela implique aussi que $f_*$ est fidèle et que le
morphisme $f^{-1}f_*\mathcal{F} \to \mathcal{F}$ est toujours surjectif.

\begin{lemma}
\label{sites-lemma-cover-from-below}
Soit $f : \mathcal{D} \to \mathcal{C}$ un morphisme de sites défini
par le foncteur $u : \mathcal{C} \to \mathcal{D}$.
Supposons que, pour tout objet $V$ de $\mathcal{D}$, il existe des objets
$U_i$ de $\mathcal{C}$ et des morphismes $u(U_i) \to V$ tels que
$\{u(U_i) \to V\}$ soit une famille couvrante de $\mathcal{D}$. Dans ce cas, le foncteur
$f_* : \Sh(\mathcal{D}) \to \Sh(\mathcal{C})$ reflète les
injections et les surjections.
\end{lemma}

\begin{proof}
Soit $\alpha : \mathcal{F} \to \mathcal{G}$ un morphisme de faisceaux
sur $\mathcal{D}$. Par hypothèse, pour tout objet $V$ de $\mathcal{D}$,
on a $\mathcal{F}(V) \subset \prod \mathcal{F}(u(U_i)) =
\prod f_*\mathcal{F}(U_i)$ en vertu de la condition de faisceau,
pour certains $U_i \in \Ob(\mathcal{C})$, et de même pour $\mathcal{G}$.
Il est donc clair que, si $f_*\alpha$ est injectif, alors
$\alpha$ est injectif. Autrement dit, $f_*$ reflète les injections.

\medskip\noindent
Supposons que $f_*\alpha$ soit surjectif. Pour $V, U_i, u(U_i) \to V$
comme ci-dessus et une section $s \in \mathcal{G}(V)$, il existe alors des familles couvrantes
$\{U_{ij} \to U_i\}$ telles que $s|_{u(U_{ij})}$ appartienne à l'image
de $\mathcal{F}(u(U_{ij}))$. Comme $\{u(U_{ij}) \to V\}$ est une famille couvrante
(puisque $u$ est continu et d'après les axiomes d'un site), on conclut que
$s$ appartient localement à l'image. Ainsi $\alpha$ est surjectif. Autrement
dit, $f_*$ reflète les surjections.
\end{proof}

\begin{example}
\label{sites-example-cover-from-below}
Construisons un morphisme $f : \mathcal{D} \to \mathcal{C}$ qui satisfait
aux hypothèses du lemme \ref{sites-lemma-cover-from-below}. Soit donc
$\varphi : G \to H$ un morphisme de groupes finis. Considérons les sites
$\mathcal{D} = \mathcal{T}_G$ et $\mathcal{C} = \mathcal{T}_H$ des $G$-ensembles et
$H$-ensembles dénombrables, dont les familles couvrantes sont les familles dénombrables de
morphismes conjointement surjectifs (exemple \ref{sites-example-site-on-group}).
Soit $u : \mathcal{T}_H \to \mathcal{T}_G$ le foncteur décrit
dans la section \ref{sites-section-G-sets-morphisms},
et $f : \mathcal{T}_G \to \mathcal{T}_H$ le
morphisme de sites correspondant. Si $\varphi$ est injectif, alors
tout $G$-ensemble dénombrable est, en tant que $G$-ensemble, le quotient d'un
$H$-ensemble dénombrable (cela est faux si $\varphi$ n'est pas injectif). Ainsi $f$ satisfait
à l'hypothèse du lemme \ref{sites-lemma-cover-from-below}.
Si le faisceau $\mathcal{F}$ sur $\mathcal{T}_G$ correspond au $G$-ensemble
$S$, alors le morphisme canonique
$$
f^{-1}f_*\mathcal{F} \longrightarrow \mathcal{F}
$$
correspond à l'application
$$
\text{Map}_G(H, S) \longrightarrow S,\quad
a \longmapsto a(1_H)
$$
Si $\varphi$ est injectif mais non surjectif, cette application est surjective
(comme il se doit d'après le lemme \ref{sites-lemma-cover-from-below}),
mais elle n'est pas injective en général
(par exemple, prendre $G = \{1\}$, $H = \{1, \sigma\}$ et $S = \{1, 2\}$).
De plus, le foncteur $f_*$ ne commute ni aux coégalisateurs ni aux
sommes amalgamées (pour $G = \{1\}$ et $H = \{1, \sigma\}$).
\end{example}





\section{Foncteurs presque cocontinus}
\label{sites-section-almost-cocontinuous}

\noindent
Soit $\mathcal{C}$ un site. La catégorie $\textit{PSh}(\mathcal{C})$ possède
un objet initial, à savoir le préfaisceau qui associe l'ensemble vide
à chaque objet de $\mathcal{C}$. Notons ce préfaisceau
$\emptyset$. Il résulte des propriétés du passage au faisceau associé que le
faisceau $\emptyset^\#$ associé à $\emptyset$ est un objet initial
de la catégorie $\Sh(\mathcal{C})$ des faisceaux sur $\mathcal{C}$.

\begin{definition}
\label{sites-definition-empty}
Soit $\mathcal{C}$ un site. On dit qu'un objet $U$ de $\mathcal{C}$
est {\it vide au sens des faisceaux} si $\emptyset^\# \to h_U^\#$
est un isomorphisme de faisceaux.
\end{definition}

\noindent
Le lemme suivant explicite cette notion.

\begin{lemma}
\label{sites-lemma-characterize-empty}
Soit $\mathcal{C}$ un site. Soit $U$ un objet de $\mathcal{C}$.
Les conditions suivantes sont équivalentes :
\begin{enumerate}
\item $U$ est vide au sens des faisceaux,
\item $\mathcal{F}(U)$ est un ensemble à un élément pour tout faisceau $\mathcal{F}$,
\item $\emptyset^\#(U)$ est un ensemble à un élément,
\item $\emptyset^\#(U)$ est non vide, et
\item la famille vide est une famille couvrante de $U$ dans $\mathcal{C}$.
\end{enumerate}
De plus, si $U$ est vide au sens des faisceaux, alors, pour tout morphisme
$U' \to U$ de $\mathcal{C}$, l'objet $U'$ est vide au sens des faisceaux.
\end{lemma}

\begin{proof}
Pour tout faisceau $\mathcal{F}$, on a
$\mathcal{F}(U) = \Mor_{\Sh(\mathcal{C})}(h_U^\#, \mathcal{F})$.
On voit donc que (1) et (2) sont équivalentes.
Il est clair que (2) implique (3), qui implique (4).
Si toute famille couvrante de $U$ est une famille non vide,
alors $\emptyset^+(U)$ est vide par définition de la construction plus.
Remarquons que $\emptyset^+ = \emptyset^\#$, puisque $\emptyset$ est un
préfaisceau séparé, voir le
théorème \ref{sites-theorem-plus}.
On voit ainsi que (4) implique (5). Si (5) est vérifiée, alors
$\mathcal{F}(U)$ est un ensemble à un élément pour tout faisceau $\mathcal{F}$
en vertu de la condition de faisceau pour $\mathcal{F}$, voir la
remarque \ref{sites-remark-sheaf-condition-empty-covering}.
Ainsi (5) implique (2), et (1) -- (5) sont équivalentes. La dernière
assertion du lemme résulte de l'axiome (3) de la
définition \ref{sites-definition-site}
appliqué à la famille couvrante vide de $U$.
\end{proof}

\begin{definition}
\label{sites-definition-almost-cocontinuous}
Soient $\mathcal{C}$ et $\mathcal{D}$ des sites.
Soit $u : \mathcal{C} \to \mathcal{D}$ un foncteur.
On dit que $u$ est {\it presque cocontinu} si, pour tout
objet $U$ de $\mathcal{C}$ et toute famille couvrante
$\{V_j \to u(U)\}_{j \in J}$, il existe une famille couvrante
$\{U_i \to U\}_{i \in I}$ dans $\mathcal{C}$ telle que,
pour tout $i$ dans $I$, l'une au moins des deux conditions suivantes soit vérifiée :
\begin{enumerate}
\item $u(U_i)$ est vide au sens des faisceaux, ou
\item le morphisme $u(U_i) \to u(U)$ se factorise par $V_j$ pour un certain $j \in J$.
\end{enumerate}
\end{definition}

\noindent
La motivation de cette définition provient d'une immersion fermée
$i : Z \to X$ d'espaces topologiques. Comme on l'a vu dans
l'exemple \ref{sites-example-closed-map-cocontinuous-false},
le foncteur continu
$X_{Zar} \to Z_{Zar}$, $U \mapsto Z \cap U$, n'est
pas cocontinu. Mais il est presque cocontinu au sens défini ci-dessus.
Nous savons que $i_*$, bien qu'il ne soit pas exact sur les faisceaux d'ensembles, est
exact sur les faisceaux de groupes abéliens, voir
Faisceaux, remarque \ref{sheaves-remark-i-star-not-exact}.
Et cela vaut en général pour les foncteurs continus et presque cocontinus.

\begin{lemma}
\label{sites-lemma-almost-cocontinuous-sheafification}
Soient $\mathcal{C}$ et $\mathcal{D}$ des sites.
Soit $u : \mathcal{C} \to \mathcal{D}$ un foncteur.
Supposons que $u$ soit continu et presque cocontinu.
Soit $\mathcal{G}$ un préfaisceau sur $\mathcal{D}$ tel que $\mathcal{G}(V)$ soit
un ensemble à un élément dès que $V$ est vide au sens des faisceaux.
Alors $(u^p\mathcal{G})^\# = u^p(\mathcal{G}^\#)$.
\end{lemma}

\begin{proof}
Soit $U \in \Ob(\mathcal{C})$. Il faut montrer que
$(u^p\mathcal{G})^\#(U) = u^p(\mathcal{G}^\#)(U)$.
Il suffit de montrer que
$(u^p\mathcal{G})^+(U) = u^p(\mathcal{G}^+)(U)$,
puisque $\mathcal{G}^+$ est un autre préfaisceau qui satisfait à
l'hypothèse du lemme. On a
$$
u^p(\mathcal{G}^+)(U) =
\mathcal{G}^+(u(U)) =
\colim_\mathcal{V} \check H^0(\mathcal{V}, \mathcal{G})
$$
où la limite inductive porte sur les familles couvrantes $\mathcal{V}$ de $u(U)$ dans
$\mathcal{D}$. D'autre part, on a
$$
u^p(\mathcal{G})^+(U) =
\colim_\mathcal{U} \check H^0(u(\mathcal{U}), \mathcal{G})
$$
où la limite inductive porte sur la catégorie des familles couvrantes
$\mathcal{U} = \{U_i \to U\}_{i \in I}$ de $U$ dans $\mathcal{C}$, et
$u(\mathcal{U}) = \{u(U_i) \to u(U)\}_{i \in I}$. La continuité
de $u$ signifie que chaque $u(\mathcal{U})$ est une famille couvrante.
Écrivons $I = I_1 \amalg I_2$, où
$$
I_2 = \{i \in I \mid u(U_i)\text{ est vide au sens des faisceaux}\}
$$
Alors $u(\mathcal{U})' = \{u(U_i) \to u(U)\}_{i \in I_1}$ est encore une famille couvrante de
car chacune des autres composantes peut être recouverte par la famille vide
et peut donc être omise en vertu de l'axiome (2) de la
définition \ref{sites-definition-site}.
En outre,
$\check H^0(u(\mathcal{U}), \mathcal{G}) =
\check H^0(u(\mathcal{U})', \mathcal{G})$
par notre hypothèse sur $\mathcal{G}$. Enfin, la presque cocontinuité de $u$
implique que, pour toute famille couvrante $\mathcal{V}$
de $u(U)$, il existe une famille couvrante $\mathcal{U}$ de $U$ telle que
$u(\mathcal{U})'$ raffine $\mathcal{V}$. Il s'ensuit que les deux limites inductives
affichées ci-dessus ont la même valeur, comme souhaité.
\end{proof}

\begin{lemma}
\label{sites-lemma-continuous-almost-cocontinuous}
Soient $\mathcal{C}$ et $\mathcal{D}$ des sites.
Soit $u : \mathcal{C} \to \mathcal{D}$ un foncteur.
Supposons que $u$ soit continu et presque cocontinu.
Alors $u^s = u^p : \Sh(\mathcal{D}) \to \Sh(\mathcal{C})$
commute aux sommes amalgamées et aux coégalisateurs (et, plus généralement,
aux limites inductives connexes finies).
\end{lemma}

\begin{proof}
Soit $\mathcal{I}$ une catégorie d'indices finie et connexe.
Soit $\mathcal{I} \to \Sh(\mathcal{D})$,
$i \mapsto \mathcal{G}_i$ un diagramme. On sait que la limite inductive de
ce diagramme est le faisceau associé à la limite inductive dans la catégorie des
préfaisceaux, voir le
lemme \ref{sites-lemma-colimit-sheaves}.
Notons $\colim^{Psh}$ la limite inductive dans la catégorie
des préfaisceaux. Puisque $\mathcal{I}$ est finie et connexe,
on voit que $\colim^{Psh}_i \mathcal{G}_i$
est un préfaisceau satisfaisant aux hypothèses du
lemme \ref{sites-lemma-almost-cocontinuous-sheafification}
(car une limite inductive finie connexe d'ensembles à un élément est un
ensemble à un élément). Ce lemme donne donc
\begin{align*}
u^s(\colim_i \mathcal{G}_i) & =
u^s((\colim^{Psh}_i \mathcal{G}_i)^\#) \\
& = (u^p(\colim^{Psh}_i \mathcal{G}_i))^\# \\
& = (\colim^{PSh}_i u^p(\mathcal{G}_i))^\# \\
& = \colim_i u^s(\mathcal{G}_i)
\end{align*}
comme voulu.
\end{proof}

\begin{lemma}
\label{sites-lemma-morphism-of-sites-almost-cocontinuous}
Soit $f : \mathcal{D} \to \mathcal{C}$ un morphisme de sites
associé au foncteur continu $u : \mathcal{C} \to \mathcal{D}$.
Si $u$ est presque cocontinu, alors $f_*$ commute aux
sommes amalgamées et aux coégalisateurs (et, plus généralement, aux limites inductives connexes finies).
\end{lemma}

\begin{proof}
C'est un cas particulier du lemme \ref{sites-lemma-continuous-almost-cocontinuous}.
\end{proof}








\section{Sous-topos}
\label{sites-section-subtopoi}

\noindent
Voici la définition.

\begin{definition}
\label{sites-definition-embedding}
Soient $\mathcal{C}$ et $\mathcal{D}$ des sites.
Un morphisme de topos $f : \Sh(\mathcal{D}) \to \Sh(\mathcal{C})$
est appelé un {\it plongement} si $f_*$ est pleinement fidèle.
\end{definition}

\noindent
D'après le lemme \ref{sites-lemma-exactness-properties},
cela équivaut à demander que le
morphisme d'adjonction $f^{-1}f_*\mathcal{F} \to \mathcal{F}$
soit un isomorphisme pour tout faisceau $\mathcal{F}$ sur $\mathcal{D}$.

\begin{definition}
\label{sites-definition-subtopos}
Soit $\mathcal{C}$ un site. Une sous-catégorie strictement pleine
$E \subset \Sh(\mathcal{C})$ est un {\it sous-topos} s'il
existe un plongement de topos $f : \Sh(\mathcal{D}) \to \Sh(\mathcal{C})$
tel que $E$ soit égale à l'image essentielle du foncteur $f_*$.
\end{definition}

\noindent
Les sous-topos construits dans le lemme suivant seront qualifiés
d'« ouverts » dans la définition donnée plus loin.

\begin{lemma}
\label{sites-lemma-open-subtopos}
Soit $\mathcal{C}$ un site. Soit $\mathcal{F}$ un faisceau sur
$\mathcal{C}$. Les assertions suivantes sont équivalentes :
\begin{enumerate}
\item $\mathcal{F}$ est un sous-objet de l'objet final de
$\Sh(\mathcal{C})$, et
\item le topos $\Sh(\mathcal{C})/\mathcal{F}$ est un sous-topos de
$\Sh(\mathcal{C})$.
\end{enumerate}
\end{lemma}

\begin{proof}
Nous avons vu dans le lemme \ref{sites-lemma-localize-topos} que
$\Sh(\mathcal{C})/\mathcal{F}$ est un topos. Rappelons en fait la
démonstration. Nous appliquons d'abord le lemme \ref{sites-lemma-topos-good-site}
pour voir que nous pouvons supposer que $\mathcal{C}$ est un site à topologie
sous-canonique, possédant les produits fibrés, un objet final $X$ et un objet $U$ tel que
$\mathcal{F} = h_U$. La démonstration du
lemme \ref{sites-lemma-localize-topos}
montre que le morphisme de topos
$j_\mathcal{F} : \Sh(\mathcal{C})/\mathcal{F} \to \Sh(\mathcal{C})$
est égal (à certaines identifications près) au morphisme de localisation
$j_U : \Sh(\mathcal{C}/U) \to \Sh(\mathcal{C})$.

\medskip\noindent
Supposons (2). Cela signifie que $j_U^{-1}j_{U, *}\mathcal{G} \to \mathcal{G}$
est un isomorphisme pour tout faisceau $\mathcal{G}$ sur $\mathcal{C}/U$.
Pour tout objet $Z/U$ de $\mathcal{C}/U$, nous avons
$$
(j_{U, *}h_{Z/U})(U) = \Mor_{\mathcal{C}/U}(U \times_X U/U, Z/U)
$$
d'après le lemme \ref{sites-lemma-localize-given-products}.
En posant $\mathcal{G} = h_{Z/U}$ dans l'égalité ci-dessus, nous obtenons
$$
\Mor_{\mathcal{C}/U}(U \times_X U/U, Z/U) = \Mor_{\mathcal{C}/U}(U, Z/U)
$$
pour tout $Z/U$. D'après le lemme de Yoneda
(Catégories, Lemme \ref{categories-lemma-yoneda}),
cela implique $U \times_X U = U$. D'après
Catégories, Lemme \ref{categories-lemma-characterize-mono-epi},
$U \to X$ est un monomorphisme, autrement dit (1) est satisfaite.

\medskip\noindent
Supposons (1). Alors $j_U^{-1} j_{U, *} = \text{id}$ d'après le
lemme \ref{sites-lemma-restrict-back}.
\end{proof}

\begin{definition}
\label{sites-definition-open-subtopos}
Soit $\mathcal{C}$ un site. Une sous-catégorie strictement pleine
$E \subset \Sh(\mathcal{C})$ est un {\it sous-topos ouvert}
s'il existe un sous-faisceau $\mathcal{F}$ de l'objet final
de $\Sh(\mathcal{C})$ tel que $E$ soit le sous-topos
$\Sh(\mathcal{C})/\mathcal{F}$ décrit dans le lemme \ref{sites-lemma-open-subtopos}.
\end{definition}

\noindent
Cela signifie qu'il existe une bijection entre la collection des sous-topos ouverts
de $\Sh(\mathcal{C})$ et l'ensemble des sous-objets de l'objet final de
$\Sh(\mathcal{C})$. À tout sous-topos ouvert correspond un complémentaire « fermé ».

\begin{lemma}
\label{sites-lemma-closed-subtopos}
Soit $\mathcal{C}$ un site. Soit $\mathcal{F}$ un sous-faisceau de l'objet final
$*$ de $\Sh(\mathcal{C})$. La sous-catégorie pleine des faisceaux
$\mathcal{G}$ tels que $\mathcal{F} \times \mathcal{G} \to \mathcal{F}$
soit un isomorphisme est un sous-topos de $\Sh(\mathcal{C})$.
\end{lemma}

\begin{proof}
Nous appliquons le lemme \ref{sites-lemma-topos-good-site} pour voir que nous pouvons supposer
que $\mathcal{C}$ est un site possédant les propriétés énumérées dans ce lemme.
En particulier, $\mathcal{C}$ possède un objet final $X$ (de sorte que
$* = h_X$) et un objet $U$ tel que $\mathcal{F} = h_U$.

\medskip\noindent
Soit $\mathcal{D} = \mathcal{C}$ en tant que catégorie, mais déclarons qu'un recouvrement
est une famille $\{V_j \to V\}$ de morphismes telle que
$\{V_i \to V\} \cup \{U \times_X V \to V\}$ soit un recouvrement.
Par notre choix de $\mathcal{C}$, cela signifie exactement que
$$
h_{U \times_X V} \amalg \coprod h_{V_i} \longrightarrow h_V
$$
est surjectif. Nous affirmons que $\mathcal{D}$ est un site, c'est-à-dire que les recouvrements
satisfont les conditions (1), (2), (3) de la définition \ref{sites-definition-site}.
La condition (1) est satisfaite. Pour la condition (2), supposons que
$\{V_i \to V\}$ et $\{V_{ij} \to V_i\}$ soient des recouvrements de $\mathcal{D}$.
Alors la composée
$$
\coprod \left(
h_{U \times_X V_i} \amalg \coprod h_{V_{ij}}
\right) \longrightarrow
h_{U \times_X V} \amalg \coprod h_{V_i} \longrightarrow h_V
$$
est surjective. Comme chacun des morphismes $U \times_X V_i \to V$
se factorise par $U \times_X V$, nous voyons que
$$
h_{U \times_X V} \amalg \coprod h_{V_{ij}} \longrightarrow h_V
$$
est surjectif, c'est-à-dire que $\{V_{ij} \to V\}$ est un recouvrement de $V$ dans
$\mathcal{D}$. La condition (3) s'en déduit de même, puisque tout changement de base
d'un morphisme surjectif de faisceaux est surjectif.

\medskip\noindent
Notons que le foncteur (identité) $u : \mathcal{C} \to \mathcal{D}$ est
continu et commute aux produits fibrés et aux objets finaux. Nous obtenons donc
un morphisme $f : \mathcal{D} \to \mathcal{C}$ de sites
(proposition \ref{sites-proposition-get-morphism}).
Remarquons que $f_*$ est le foncteur identité sur les
préfaisceaux sous-jacents, donc qu'il est pleinement fidèle. Pour achever la démonstration, il nous faut
montrer que l'image essentielle de $f_*$ est la sous-catégorie pleine
$E \subset \Sh(\mathcal{C})$ distinguée dans le lemme. Pour cela, notons
que $\mathcal{G} \in \Ob(\Sh(\mathcal{C}))$ appartient à $E$ si et seulement si
$\mathcal{G}(U \times_X V)$ est un singleton pour tout objet
$V$ de $\mathcal{C}$. Un tel faisceau satisfait donc la
propriété de faisceau pour tous les recouvrements de $\mathcal{D}$ (argument omis).
Réciproquement, si $\mathcal{G}$ satisfait la propriété de faisceau
pour tous les recouvrements de $\mathcal{D}$, alors $\mathcal{G}(U \times_X V)$
est un singleton, puisque, dans $\mathcal{D}$, l'objet $U \times_X V$ est
recouvert par le recouvrement vide.
\end{proof}

\begin{definition}
\label{sites-definition-closed-subtopos}
Soit $\mathcal{C}$ un site. Une sous-catégorie strictement pleine
$E \subset \Sh(\mathcal{C})$ est un {\it sous-topos fermé}
s'il existe un sous-faisceau $\mathcal{F}$ de l'objet final
de $\Sh(\mathcal{C})$ tel que $E$ soit le sous-topos
décrit dans le lemme \ref{sites-lemma-closed-subtopos}.
\end{definition}

\noindent
Nous pouvons maintenant définir ce que sont une immersion fermée
et une immersion ouverte de topos.

\begin{definition}
\label{sites-definition-immersion-topoi}
Soit $f : \Sh(\mathcal{D}) \to \Sh(\mathcal{C})$ un morphisme de topos.
\begin{enumerate}
\item Nous disons que $f$ est une {\it immersion ouverte} si $f$ est un plongement
et si l'image essentielle de $f_*$ est un sous-topos ouvert.
\item Nous disons que $f$ est une {\it immersion fermée} si $f$ est un plongement
et si l'image essentielle de $f_*$ est un sous-topos fermé.
\end{enumerate}
\end{definition}

\begin{lemma}
\label{sites-lemma-closed-immersion}
Soit $i : \Sh(\mathcal{D}) \to \Sh(\mathcal{C})$ une immersion fermée de
topos. Alors $i_*$ est pleinement fidèle, transforme les surjections en surjections,
commute aux coégalisateurs, commute aux sommes amalgamées, reflète les injections,
reflète les surjections et reflète les bijections.
\end{lemma}

\begin{proof}
Soit $\mathcal{F}$ un sous-faisceau de l'objet final $*$ de $\Sh(\mathcal{C})$
et soit $E \subset \Sh(\mathcal{C})$ la sous-catégorie pleine formée
des $\mathcal{G}$ tels que
$\mathcal{F} \times \mathcal{G} \to \mathcal{F}$ soit un isomorphisme.
D'après le lemme \ref{sites-lemma-closed-subtopos},
le foncteur $i_*$ est isomorphe au
foncteur d'inclusion $\iota : E \to \Sh(\mathcal{C})$.

\medskip\noindent
Soit $j_{\mathcal{F}} : \Sh(\mathcal{C})/\mathcal{F} \to \Sh(\mathcal{C})$
le foncteur de localisation (lemme \ref{sites-lemma-localize-topos}).
Notons que $E$ peut également être décrite comme
la collection des faisceaux $\mathcal{G}$ tels que
$j_\mathcal{F}^{-1}\mathcal{G} = *$.

\medskip\noindent
Soient $a, b : \mathcal{G}_1 \to \mathcal{G}_2$ deux morphismes de $E$.
Pour prouver que $\iota$ commute aux coégalisateurs, il suffit de montrer que
le coégalisateur de $a$, $b$ dans $\Sh(\mathcal{C})$ appartient à $E$.
C'est clair, car
le coégalisateur de deux morphismes $* \to *$ est $*$ et parce que
$j_\mathcal{F}^{-1}$ est exact. Il en va de même pour les sommes amalgamées.

\medskip\noindent
Ainsi $i_*$ satisfait les propriétés (5), (6) et (7) du
lemme \ref{sites-lemma-exactness-properties} et, par suite,
le morphisme $i$ satisfait toutes les propriétés mentionnées dans ce
lemme, en particulier celles qui sont mentionnées dans le présent lemme.
\end{proof}








\section{Faisceaux de structures algébriques}
\label{sites-section-sheaves-algebraic-structures}

\noindent
Dans Faisceaux, section \ref{sheaves-section-algebraic-structures},
nous avons défini un type de structure algébrique comme une paire
$(\mathcal{A}, s)$, où $\mathcal{A}$ est une catégorie
et $s : \mathcal{A} \to \textit{Ensembles}$ est un foncteur tel
que
\begin{enumerate}
\item $s$ soit fidèle,
\item $\mathcal{A}$ possède les limites et $s$ commute aux limites,
\item $\mathcal{A}$ possède les colimites filtrantes et $s$ commute à celles-ci, et
\item $s$ reflète les isomorphismes.
\end{enumerate}
Pour un tel type de structure algébrique, nous avons vu qu'un préfaisceau
$\mathcal{F}$ à valeurs dans $\mathcal{A}$ sur un espace $X$ est un faisceau si et
seulement si le préfaisceau d'ensembles associé est un faisceau. De plus,
nous avons élaboré la notion de fibre et, étant donnée une application continue
$f : X \to Y$, nous avons défini les foncteurs adjoints image directe et image inverse
sur les faisceaux de structures algébriques, lesquels coïncident avec l'image directe
et l'image inverse sur les faisceaux d'ensembles sous-jacents. En outre, le prolongement
d'un faisceau de structures algébriques d'une base à tous les ouverts
d'un espace se comporte comme prévu.

\medskip\noindent
Une partie de ces résultats reste valable dans le cadre des sites et des faisceaux.
Soit $(\mathcal{A}, s)$ un type de structure algébrique.
Soit $\mathcal{C}$ un site. Notons
$\textit{PSh}(\mathcal{C}, \mathcal{A})$,
resp.\ $\Sh(\mathcal{C}, \mathcal{A})$ la catégorie
des préfaisceaux, resp.\ des faisceaux à valeurs dans $\mathcal{A}$ sur $\mathcal{C}$.
\begin{itemize}
\item[] ($\alpha$) Un préfaisceau à valeurs dans $\mathcal{A}$ est
un faisceau si et seulement si son préfaisceau d'ensembles sous-jacent est un faisceau.
Voir la démonstration de Faisceaux, Lemme \ref{sheaves-lemma-sheaves-structure}.
\item[] ($\beta$) Étant donné un préfaisceau $\mathcal{F}$ à valeurs dans
$\mathcal{A}$, le préfaisceau ${\mathcal{F}}^\# = (\mathcal{F}^+)^+$
est un faisceau. Cela est vrai parce que les colimites intervenant dans la faisceautisation
sont filtrantes, et même des colimites sur des ensembles ordonnés filtrants (voir
la section \ref{sites-section-sheafification}, notamment la démonstration du
lemme \ref{sites-lemma-sheafification-exact}),
et parce que $s$ commute aux colimites filtrantes.
\item[] ($\gamma$) Nous obtenons le diagramme commutatif suivant
$$
\xymatrix{
\Sh(\mathcal{C}, \mathcal{A}) \ar@<1ex>[r] \ar[d]_s &
\textit{PSh}(\mathcal{C}, \mathcal{A}) \ar@<1ex>[l]^{{}^\#} \ar[d]^s\\
\Sh(\mathcal{C}) \ar@<1ex>[r] &
\textit{PSh}(\mathcal{C}) \ar@<1ex>[l]
}
$$
\item[] ($\delta$) Nous avons $\mathcal{F} = \mathcal{F}^\#$ si et seulement si
$\mathcal{F}$ est un faisceau de structures algébriques.
\item[] ($\epsilon$) Le foncteur ${}^\#$ est adjoint au foncteur d'inclusion :
$$
\Mor_{\textit{PSh}(\mathcal{C}, \mathcal{A})}(\mathcal{G}, \mathcal{F})
=
\Mor_{\Sh(\mathcal{C}, \mathcal{A})}(\mathcal{G}^\#, \mathcal{F})
$$
La démonstration est la même que celle de la
proposition \ref{sites-proposition-sheafification-adjoint}.
\item[] ($\zeta$) Le foncteur
$\mathcal{F} \mapsto \mathcal{F}^\#$ est exact à gauche.
La démonstration est la même que celle du lemme \ref{sites-lemma-sheafification-exact}.
\end{itemize}

\begin{definition}
\label{sites-definition-pushforward-algebraic-structures}
Soit $f : \mathcal{D} \to \mathcal{C}$ un morphisme de sites
donné par un foncteur $u : \mathcal{C} \to \mathcal{D}$.
Nous définissons le foncteur {\it image directe} pour les préfaisceaux de structures algébriques
par la règle $u^p\mathcal{F}(U) = \mathcal{F}(uU)$,
et pour les faisceaux de structures algébriques par la même règle, à savoir
$f_*\mathcal{F}(U) = \mathcal{F}(uU)$.
\end{definition}

\noindent
Le problème survient lorsqu'on cherche à définir l'image inverse.
La raison en est que les colimites définissant
le foncteur $u_p$ dans la section \ref{sites-section-functoriality-PSh}
ne sont pas nécessairement filtrantes. Les axiomes ci-dessus ne suffisent donc pas
en général pour définir l'image inverse d'un (pré)faisceau de structures
algébriques. Néanmoins, dans presque tous les cas, le lemme
suivant suffit pour définir l'image directe et l'image inverse
des (pré)faisceaux de structures algébriques.

\begin{lemma}
\label{sites-lemma-push-pull-good-case}
Supposons que le foncteur $u : \mathcal{C} \to \mathcal{D}$ satisfasse
aux hypothèses de la proposition \ref{sites-proposition-get-morphism}
et donne donc naissance à un morphisme de sites
$f : \mathcal{D} \to \mathcal{C}$. Dans ce cas,
le foncteur image inverse $f^{-1}$ (resp.\ $u_p$) et le foncteur image directe
$f_*$ (resp. $u^p$) s'étendent en une paire de foncteurs adjoints sur
les catégories de faisceaux (resp.\ de préfaisceaux) de structures algébriques.
De plus, ces foncteurs commutent au passage au
faisceau (resp.\ au préfaisceau) d'ensembles sous-jacent.
\end{lemma}

\begin{proof}
Nous avons défini $f_* = u^p$ ci-dessus.
Au cours de la démonstration de la proposition \ref{sites-proposition-get-morphism},
nous avons vu que toutes les colimites utilisées pour définir $u_p$ sont
filtrantes sous les hypothèses de la proposition.
Nous déduisons donc de la définition d'un type de
structure algébrique que nous pouvons définir $u_p$, en tant que foncteur sur
les préfaisceaux de structures algébriques, au moyen exactement des mêmes colimites.
L'adjonction de $u_p$ et $u^p$ se démontre exactement de la
même manière que le lemme \ref{sites-lemma-adjoints-u}.
La discussion ci-dessus de la faisceautisation des préfaisceaux de
structures algébriques implique alors que nous pouvons définir
$f^{-1}(\mathcal{F}) = (u_p\mathcal{F})^\#$.
\end{proof}

\noindent
Nous discutons brièvement une méthode pour traiter l'image inverse et
l'image directe pour un morphisme général de sites et, plus généralement,
pour tout morphisme de topos.

\medskip\noindent
Soit $\mathcal{C}$ un site.
Dans le cas $\mathcal{A} = \textit{Ab}$,
nous pouvons considérer un (pré)faisceau abélien sur $\mathcal{C}$
comme un quadruplet $(\mathcal{F}, +, 0, i)$.
Les données sont les suivantes :
\begin{enumerate}
\item[(D1)] $\mathcal{F}$ est un faisceau d'ensembles,
\item[(D2)] $+ : \mathcal{F} \times \mathcal{F} \to \mathcal{F}$ est
un morphisme de faisceaux d'ensembles,
\item[(D3)] $0 : * \to \mathcal{F}$ est un morphisme depuis le
faisceau singleton (voir l'exemple \ref{sites-example-singleton-sheaf})
vers $\mathcal{F}$, et
\item[(D4)] $i : \mathcal{F} \to \mathcal{F}$ est un morphisme de faisceaux
d'ensembles.
\end{enumerate}
Ces données doivent satisfaire les axiomes suivants :
\begin{enumerate}
\item[(A1)] $+$ est associative et commutative,
\item[(A2)] $0$ est un élément neutre pour $+$, et
\item[(A3)] $+ \circ (1, i) = 0 \circ (\mathcal{F} \to *)$.
\end{enumerate}
Comparer avec Faisceaux, Lemme \ref{sheaves-lemma-abelian-presheaves}.
Soit $f : \mathcal{D} \to \mathcal{C}$ un morphisme de sites.
Notons que, puisque $f^{-1}$ est exact, nous avons
$f^{-1}* = *$ et
$f^{-1}(\mathcal{F} \times \mathcal{F}) =
f^{-1}\mathcal{F} \times f^{-1}\mathcal{F}$.
Nous pouvons donc définir $f^{-1}\mathcal{F}$ simplement comme le quadruplet
$(f^{-1}\mathcal{F}, f^{-1}+, f^{-1}0, f^{-1}i)$. Les axiomes
seront préservés parce que $f^{-1}$ est un foncteur
qui commute aux limites finies. Enfin, il n'est pas difficile
de vérifier que $f_*$ et $f^{-1}$ sont adjoints comme à l'ordinaire.

\medskip\noindent
Dans \cite{SGA4}, cette méthode est utilisée. On y introduit ce
qui est appelé une « {\it esp\`ece de structure alg\'ebrique $\ll$d\'efinie
par limites projectives finies$\gg$} ». Pour une telle esp\`ece, on
peut utiliser la méthode décrite ci-dessus pour définir une paire de foncteurs adjoints
$f^{-1}$ et $f_*$ comme ci-dessus. Cela fonctionne manifestement pour la plupart
des structures algébriques que l'on rencontre en pratique.
Au lieu de formaliser cette construction, nous énumérons simplement les
structures algébriques pour lesquelles cette méthode fonctionne (ce qui doit être
vérifié cas par cas). En fait, cette méthode fonctionne pour tout
morphisme de topos.

\begin{proposition}
\label{sites-proposition-functoriality-algebraic-structures-topoi}
\begin{slogan}
Les morphismes de topos préservent les structures algébriques.
\end{slogan}
Soient $\mathcal{C}$ et $\mathcal{D}$ des sites.
Soit $f = (f^{-1}, f_*)$ un morphisme de topos
de $\Sh(\mathcal{D}) \to \Sh(\mathcal{C})$.
La méthode introduite ci-dessus donne naissance à une
paire de foncteurs adjoints $(f^{-1}, f_*)$ sur les faisceaux de structures algébriques,
compatible avec le passage aux faisceaux d'ensembles sous-jacents,
pour les types suivants de structures algébriques :
\begin{enumerate}
\item ensembles pointés,
\item groupes abéliens,
\item groupes,
\item monoïdes,
\item anneaux,
\item modules sur un anneau fixé, et
\item algèbres de Lie sur un corps fixé.
\end{enumerate}
De plus, dans chacun de ces cas, les résultats ci-dessus étiquetés ($\alpha$),
($\beta$), ($\gamma$), ($\delta$), ($\epsilon$) et ($\zeta$) sont valables.
\end{proposition}

\begin{proof}
La dernière assertion de la proposition est satisfaite simplement parce que chacune des
catégories énumérées, munie du foncteur d'oubli évident, est bien un type de
structure algébrique au sens expliqué au début de cette section.
Voir Faisceaux, Lemme \ref{sheaves-lemma-list-algebraic-structures}.

\medskip\noindent
Démonstration de (2). Nous considérons un faisceau de groupes abéliens comme
un quadruplet $(\mathcal{F}, +, 0, i)$, ainsi qu'il a été expliqué dans la discussion précédant
la proposition.
Si $(\mathcal{F}, +, 0, i)$ est défini sur $\mathcal{C}$, alors son image inverse
est définie comme $(f^{-1}\mathcal{F}, f^{-1}+, f^{-1}0, f^{-1}i)$.
Si $(\mathcal{G}, +, 0, i)$ est défini sur $\mathcal{D}$, alors son image directe
est définie comme $(f_*\mathcal{G}, f_*+, f_*0, f_*i)$. Cette définition convient puisque
$f_*(\mathcal{G} \times \mathcal{G}) = f_*\mathcal{G} \times f_*\mathcal{G}$.
L'adjonction résulte de l'adjonction des foncteurs sur les ensembles,
puisque
$$
\Mor_{\textit{Ab}(\mathcal{C})}
((\mathcal{F}_1, +, 0, i), (\mathcal{F}_2, +, 0, i))
=
\left\{
\begin{matrix}
\varphi \in \Mor_{\Sh(\mathcal{C})}
(\mathcal{F}_1, \mathcal{F}_2) \\
\varphi \text{ est compatible avec }+, 0, i
\end{matrix}
\right\}
$$
Les détails sont laissés au lecteur.

\medskip\noindent
Cette méthode fonctionne également pour les faisceaux d'anneaux lorsque l'on considère
un faisceau d'anneaux (unitaires) comme un sextuplet
$(\mathcal{O}, + , 0, i, \cdot, 1)$ satisfaisant une liste
d'axiomes que l'on peut trouver dans tout ouvrage élémentaire
d'algèbre.

\medskip\noindent
Un faisceau d'ensembles pointés est une paire $(\mathcal{F}, p)$, où
$\mathcal{F}$ est un faisceau d'ensembles et $p : * \to \mathcal{F}$
est un morphisme de faisceaux d'ensembles.

\medskip\noindent
Un faisceau de groupes est donné par un quadruplet $(\mathcal{F}, \cdot, 1, i)$
muni des axiomes appropriés.

\medskip\noindent
Un faisceau de monoïdes est donné par une paire $(\mathcal{F}, \cdot)$
munie de l'axiome approprié.

\medskip\noindent
Soit $R$ un anneau. Un faisceau de $R$-modules est donné par
un quintuplet $(\mathcal{F}, +, 0, i, \{\lambda_r\}_{r \in R})$,
où le quadruplet $(\mathcal{F}, +, 0, i)$ est un faisceau de
groupes abéliens comme ci-dessus et $\lambda_r : \mathcal{F} \to \mathcal{F}$
est une famille de morphismes de faisceaux d'ensembles
telle que
$\lambda_r \circ 0 = 0$,
$\lambda_r \circ + = + \circ (\lambda_r, \lambda_r)$,
$\lambda_{r + r'} =
+ \circ \lambda_r \times \lambda_{r'} \circ (\text{id}, \text{id})$,
$\lambda_{rr'} = \lambda_r \circ \lambda_{r'}$,
$\lambda_1 = \text{id}$, $\lambda_0 = 0 \circ (\mathcal{F} \to *)$.
\end{proof}

\noindent
Nous étudierons la catégorie des faisceaux de modules sur un faisceau d'anneaux dans
Modules sur les sites, section \ref{sites-modules-section-sheaves-modules}.

\begin{remark}
\label{sites-remark-no-pullback-presheaves}
Soient $\mathcal{C}$ et $\mathcal{D}$ des sites.
Soit $u : \mathcal{D} \to \mathcal{C}$ un foncteur continu
qui donne naissance à un morphisme de sites $\mathcal{C} \to \mathcal{D}$.
Notons que, même dans le cas des groupes abéliens, nous n'avons pas défini
de foncteur image inverse pour les préfaisceaux de groupes abéliens.
Puisque toutes les colimites sont représentables dans
la catégorie des groupes abéliens, nous pouvons assurément définir
un foncteur $u_p^{ab}$ sur les préfaisceaux abéliens au moyen des mêmes colimites
que celles que nous avons utilisées pour définir $u_p$ sur les préfaisceaux d'ensembles.
Il sera également vrai que $u_p^{ab}$ est adjoint à
$u^p$ sur les catégories de préfaisceaux abéliens.
Cependant, $u_p^{ab}$ ne coïncidera pas toujours avec
$u_p$ sur les préfaisceaux d'ensembles sous-jacents.
\end{remark}












\section{Applications d'image inverse}
\label{sites-section-pullback}

\noindent
Il arrive parfois qu'un site $\mathcal{C}$ ne possède
pas d'objet final. Dans ce cas, nous définissons le
foncteur des sections globales comme suit.

\begin{definition}
\label{sites-definition-global-sections}
L'{\it ensemble des sections globales} d'un préfaisceau d'ensembles $\mathcal{F}$ sur un
site $\mathcal{C}$ est l'ensemble
$$
\Gamma(\mathcal{C}, \mathcal{F}) =
\Mor_{\textit{PSh}(\mathcal{C})}(*, \mathcal{F})
$$
où $*$ est l'objet final de la catégorie des préfaisceaux sur
$\mathcal{C}$, c'est-à-dire le préfaisceau qui associe à tout objet
un singleton.
\end{definition}

\noindent
Bien entendu, la même définition s'applique aussi aux faisceaux.
Voici une manière de calculer les sections globales.

\begin{lemma}
\label{sites-lemma-compute-global-sections}
Soit $\mathcal{C}$ un site. Soient $a, b : V \to U$ des objets de $\mathcal{C}$
tels que
$$
\xymatrix{
h_V^\# \ar@<1ex>[r] \ar@<-1ex>[r] & h_U^\# \ar[r] & {*}
}
$$
soit un coégalisateur dans $\Sh(\mathcal{C})$. Alors
$\Gamma(\mathcal{C}, \mathcal{F})$ est l'égalisateur de
$a^*, b^* : \mathcal{F}(U) \to \mathcal{F}(V)$.
\end{lemma}

\begin{proof}
Puisque $\Mor_{\Sh(\mathcal{C})}(h_U^\#, \mathcal{F}) = \mathcal{F}(U)$,
cela résulte immédiatement des définitions.
\end{proof}

\noindent
Soit maintenant $f : \Sh(\mathcal{D}) \to \Sh(\mathcal{C})$ un morphisme de topos.
Alors, pour tout faisceau $\mathcal{F}$ sur $\mathcal{C}$, il existe une application d'image inverse
$$
f^{-1} :
\Gamma(\mathcal{C}, \mathcal{F})
\longrightarrow
\Gamma(\mathcal{D}, f^{-1}\mathcal{F})
$$
En effet, comme $f^{-1}$ est exact, il transforme $*$ en $*$. Ainsi, une section
globale $s$ de $\mathcal{F}$ sur $\mathcal{C}$, qui est un morphisme
de faisceaux $s : * \to \mathcal{F}$, a pour image inverse le morphisme
$f^{-1}s : * = f^{-1}* \to f^{-1}\mathcal{F}$.

\medskip\noindent
Nous pouvons généraliser légèrement cette construction en considérant une paire de faisceaux
$\mathcal{F}$, $\mathcal{G}$ sur $\mathcal{C}$, $\mathcal{D}$,
munie d'un morphisme $f^{-1}\mathcal{F} \to \mathcal{G}$. Nous composons alors
la construction ci-dessus avec le morphisme évident
$\Gamma(\mathcal{D}, f^{-1}\mathcal{F}) \to \Gamma(\mathcal{D}, \mathcal{G})$
pour obtenir un morphisme
$$
\Gamma(\mathcal{C}, \mathcal{F})
\longrightarrow
\Gamma(\mathcal{D}, \mathcal{G})
$$
Ce morphisme est parfois lui aussi appelé application d'image inverse.

\medskip\noindent
Une construction un peu plus générale, qui intervient fréquemment de façon naturelle,
est la suivante. Supposons donné un diagramme commutatif de
morphismes de topos
$$
\xymatrix{
\Sh(\mathcal{D}) \ar[rd]_h \ar[rr]_f & &
\Sh(\mathcal{C}) \ar[ld]^g \\
& \Sh(\mathcal{B})
}
$$
Supposons ensuite donné un faisceau $\mathcal{F}$ sur $\mathcal{C}$.
Il existe alors une {\it application d'image inverse}
$$
f^{-1} : g_*\mathcal{F} \longrightarrow h_*f^{-1}\mathcal{F}
$$
En effet, il s'agit simplement du morphisme provenant de l'identification
$g_*f_*f^{-1}\mathcal{F} = h_*f^{-1}\mathcal{F}$ et de
l'application de $g_*$ au morphisme canonique $\mathcal{F} \to f_*f^{-1}\mathcal{F}$.
Si $g$ est l'identité, alors ce morphisme sur les sections globales coïncide
avec l'application d'image inverse ci-dessus.

\medskip\noindent
Dans la situation du paragraphe précédent, supposons que
nous ayons une paire de faisceaux $\mathcal{F}$, $\mathcal{G}$
sur $\mathcal{C}$, $\mathcal{D}$, munie d'un morphisme
$f^{-1}\mathcal{F} \to \mathcal{G}$ ; nous composons alors
l'application d'image inverse ci-dessus avec l'application de $h_*$ au morphisme
$f^{-1}\mathcal{F} \to \mathcal{G}$
pour obtenir un morphisme
$$
g_*\mathcal{F} \longrightarrow h_*\mathcal{G}
$$
En se restreignant aux sections sur un objet de $\mathcal{B}$, on retrouve
l'« application d'image inverse » sur les sections globales discutée ci-dessus (pour des
choix convenables de sites).

\medskip\noindent
Une situation encore plus générale est celle où nous avons un diagramme commutatif de topos
$$
\xymatrix{
\Sh(\mathcal{D}) \ar[d]_h \ar[r]_f &
\Sh(\mathcal{C}) \ar[d]^g \\
\Sh(\mathcal{B}) \ar[r]^e & \Sh(\mathcal{A})
}
$$
et un faisceau $\mathcal{G}$ sur $\mathcal{C}$. Il existe alors un
{\it morphisme de changement de base}
$$
e^{-1}g_*\mathcal{G} \longrightarrow h_*f^{-1}\mathcal{G}.
$$
En effet, ce morphisme est adjoint à un morphisme
$g_*\mathcal{G} \to e_*h_*f^{-1}\mathcal{G} = (e \circ h)_*f^{-1}\mathcal{G}$,
qui est l'application d'image inverse que nous venons de décrire.

\begin{remark}
\label{sites-remark-compose-base-change}
Considérons un diagramme commutatif
$$
\xymatrix{
\Sh(\mathcal{B}') \ar[r]_k \ar[d]_{f'} &
\Sh(\mathcal{B}) \ar[d]^f \\
\Sh(\mathcal{C}') \ar[r]^l \ar[d]_{g'} &
\Sh(\mathcal{C}) \ar[d]^g \\
\Sh(\mathcal{D}') \ar[r]^m &
\Sh(\mathcal{D})
}
$$
de topos. Alors les morphismes de changement de base des deux carrés se composent pour donner le
morphisme de changement de base du rectangle extérieur. Plus précisément, la composée
\begin{align*}
m^{-1} \circ (g \circ f)_*
& =
m^{-1} \circ g_* \circ f_* \\
& \to g'_* \circ l^{-1} \circ f_* \\
& \to g'_* \circ f'_* \circ k^{-1} \\
& = (g' \circ f')_* \circ k^{-1}
\end{align*}
est le morphisme de changement de base du rectangle. Nous omettons la vérification.
\end{remark}

\begin{remark}
\label{sites-remark-compose-base-change-horizontal}
Considérons un diagramme commutatif
$$
\xymatrix{
\Sh(\mathcal{C}'') \ar[r]_{g'} \ar[d]_{f''} &
\Sh(\mathcal{C}') \ar[r]_g \ar[d]_{f'} &
\Sh(\mathcal{C}) \ar[d]^f \\
\Sh(\mathcal{D}'') \ar[r]^{h'} &
\Sh(\mathcal{D}') \ar[r]^h &
\Sh(\mathcal{D})
}
$$
de topos annelés. Alors les morphismes de changement de base
des deux carrés se composent pour donner le morphisme de
changement de base du rectangle extérieur. Plus précisément,
la composée
\begin{align*}
(h \circ h')^{-1} \circ f_*
& =
(h')^{-1} \circ h^{-1} \circ f_* \\
& \to (h')^{-1} \circ f'_* \circ g^{-1} \\
& \to f''_* \circ (g')^{-1} \circ g^{-1} \\
& = f''_* \circ (g \circ g')^{-1}
\end{align*}
est le morphisme de changement de base du rectangle. Nous omettons la vérification.
\end{remark}






\section{Comparaison avec SGA4}
\label{sites-section-notation-comparison}

\noindent
Nos notations pour les foncteurs
$u^p$ et $u_p$ de la section \ref{sites-section-functoriality-PSh} ainsi que
$u^s$ et $u_s$ de la section \ref{sites-section-continuous-functors}
sont empruntées à \cite[pages 14 et 42]{ArtinTopologies}. Ces choix étant faits,
la notation pour le foncteur
${}_pu$ de la section \ref{sites-section-more-functoriality-PSh} et
${}_su$ de la section \ref{sites-section-cocontinuous-functors}
paraît raisonnable.
Dans cette section, nous comparons nos notations à celles de SGA4.

\medskip\noindent
{\bf Préfaisceaux :} Soit $u : \mathcal{C} \to \mathcal{D}$ un foncteur
entre catégories.
Le foncteur $u^p$ est noté $u^*$ dans \cite[Exposé I, section 5]{SGA4}.
Le foncteur $u_p$ est noté $u_!$ dans \cite[Exposé I, proposition 5.1]{SGA4}.
Le foncteur ${}_pu$ est noté $u_*$ dans
\cite[Exposé I, proposition 5.1]{SGA4}.
Autrement dit, nous avons
$$
u_p, u^p, {}_pu\quad(SP)
\quad\text{contre}\quad
u_!, u^*, u_*\quad(SGA4)
$$
Nous attirons l'attention du lecteur sur le fait que des notations différentes
sont employées pour ces foncteurs dans différentes parties de SGA4.

\medskip\noindent
{\bf Faisceaux et foncteurs continus :}
Supposons que $\mathcal{C}$ et $\mathcal{D}$ soient des sites et que
$u : \mathcal{C} \to \mathcal{D}$ soit
un foncteur continu (définition \ref{sites-definition-continuous}).
Le foncteur $u^s$ est noté $u_s$ dans \cite[Exposé III, 1.11]{SGA4}.
Le foncteur $u_s$ est noté $u^s$ dans
\cite[Exposé III, proposition 1.2]{SGA4}.
Autrement dit, nous avons
$$
u_s, u^s\quad(SP)
\quad\text{contre}\quad
u^s, u_s\quad(SGA4)
$$
Lorsque $u$ définit un morphisme de sites
$f : \mathcal{D} \to \mathcal{C}$
(définition \ref{sites-definition-morphism-sites}),
nous voyons que le morphisme de topos associé
(lemme \ref{sites-lemma-morphism-sites-topoi}) est le même que celui de
\cite[Exposé IV, (4.9.1.1)]{SGA4}.

\medskip\noindent
{\bf Faisceaux et foncteurs cocontinus :}
Supposons que $\mathcal{C}$ et $\mathcal{D}$ soient des sites et que
$u : \mathcal{C} \to \mathcal{D}$ soit
un foncteur cocontinu (définition \ref{sites-definition-cocontinuous}).
Le foncteur ${}_su$ (lemme \ref{sites-lemma-pu-sheaf}) est noté $u_*$ dans
\cite[Exposé III, proposition 2.3]{SGA4}.
Le foncteur $(u^p\ )^\#$ est noté $u^*$ dans
\cite[Exposé III, proposition 2.3]{SGA4}.
Autrement dit, nous avons
$$
(u^p\ )^\#, {}_su\quad(SP)
\quad\text{contre}\quad
u^*, u_*\quad(SGA4)
$$
Ainsi, le morphisme de topos associé à $u$ dans le
lemme \ref{sites-lemma-cocontinuous-morphism-topoi}
est le même que celui de \cite[Exposé IV, 4.7]{SGA4}.

\medskip\noindent
{\bf Morphismes de topos :} Si $f$ est un morphisme de topos donné
par les foncteurs $(f^{-1}, f_*)$, alors le foncteur $f^{-1}$
est noté $f^*$ dans \cite[Exposé IV, définition 3.1]{SGA4}.
Nous emploierons $f^{-1}$ pour désigner l'image inverse des faisceaux d'ensembles
ou, plus généralement, des faisceaux de structures algébriques
(section \ref{sites-section-sheaves-algebraic-structures}).
Nous emploierons $f^*$ pour désigner l'image inverse des faisceaux de modules
pour un morphisme de topos annelés
(Modules sur les sites, définition \ref{sites-modules-definition-pushforward}).



\section{Topologies}
\label{sites-section-topologies}

\noindent
Dans cette section, nous définissons ce qu'est une topologie sur une catégorie,
au sens de \cite{SGA4}. On peut développer dans ce langage toute la théorie des
faisceaux et des topos. On trouvera un exposé moderne de cette matière
dans \cite{KS}. Toutefois, le cas qui intéresse le plus la géométrie algébrique
est celui de la topologie définie par un site sur sa catégorie sous-jacente.
Nous conseillons donc vivement au lecteur qui aborde ce texte pour la première fois de
{\bf sauter cette section, ainsi que toutes les autres sections de ce chapitre} !

\begin{definition}
\label{sites-definition-sieve}
Soit $\mathcal{C}$ une catégorie. Soit $U \in \Ob(\mathcal{C})$.
Un {\it crible $S$ sur $U$} est un sous-préfaisceau $S \subset h_U$.
\end{definition}

\noindent
Autrement dit, un crible sur $U$ détermine, pour chaque objet
$T \in \Ob(\mathcal{C})$, un sous-ensemble $S(T)$ de l'ensemble
de tous les morphismes $T \to U$. En fait, l'unique condition
imposée à la famille de sous-ensembles
$S(T) \subset h_U(T) = \Mor_\mathcal{C}(T, U)$
est la règle suivante :
\begin{equation}
\label{sites-equation-property-sieve}
\left.
\begin{matrix}
(\alpha : T \to U) \in S(T) \\
g : T' \to T
\end{matrix}
\right\} \Rightarrow
(\alpha \circ g : T' \to U) \in S(T')
\end{equation}
Une bonne représentation mentale consiste à considérer
l'application $S \to h_U$ comme un ``morphisme de $S$ vers $U$''.

\begin{lemma}
\label{sites-lemma-sieves-set}
Soit $\mathcal{C}$ une catégorie. Soit $U \in \Ob(\mathcal{C})$.
\begin{enumerate}
\item La collection des cribles sur $U$ est un ensemble.
\item L'inclusion définit un ordre partiel sur cet ensemble.
\item Les réunions et les intersections de cribles sont des cribles.
\item
\label{sites-item-sieve-generated}
Étant donnée une famille de morphismes $\{U_i \to U\}_{i\in I}$
de $\mathcal{C}$ de but $U$,
il existe un unique plus petit crible $S$ sur $U$ tel que
chaque $U_i \to U$ appartienne à $S(U_i)$.
\item Le crible $S = h_U$ est le crible maximal.
\item Le sous-préfaisceau vide est le crible minimal.
\end{enumerate}
\end{lemma}

\begin{proof}
D'après notre définition d'un sous-préfaisceau, la collection de
tous les sous-préfaisceaux d'un préfaisceau $\mathcal{F}$ est une partie de
$\prod_{U \in \Ob(\mathcal{C})} \mathcal{P}(\mathcal{F}(U))$.
Or ceci est un ensemble. (Ici $\mathcal{P}(A)$ désigne
l'ensemble des parties de $A$.) Par conséquent, la collection des cribles sur $U$
est un ensemble.

\medskip\noindent
L'ordre partiel est défini par : $S \leq S'$ si et seulement si
$S(T) \subset S'(T)$ pour tout $T \to U$. Notation : $S \subset S'$.

\medskip\noindent
Étant donnée une famille de cribles $S_i$, $i \in I$, sur $U$, nous pouvons
définir $\bigcup S_i$ comme le crible prenant les valeurs
$(\bigcup S_i)(T) = \bigcup S_i(T)$ pour tout
$T \in \Ob(\mathcal{C})$.
Nous définissons l'intersection $\bigcap S_i$ de la même manière.

\medskip\noindent
Étant donnée $\{U_i \to U\}_{i\in I}$ comme dans l'énoncé, considérons
les morphismes de préfaisceaux $h_{U_i} \to h_U$. Nous définissons simplement
$S$ comme la réunion des images (définition \ref{sites-definition-image})
de ces applications de préfaisceaux.

\medskip\noindent
Les deux dernières assertions du lemme sont évidentes.
\end{proof}

\begin{definition}
\label{sites-definition-sieve-generated}
Soit $\mathcal{C}$ une catégorie.
Étant donnée une famille de morphismes $\{f_i : U_i \to U\}_{i\in I}$
de $\mathcal{C}$ de but $U$, nous disons que le crible
$S$ sur $U$ décrit dans le lemme \ref{sites-lemma-sieves-set},
partie (\ref{sites-item-sieve-generated}), est le {\it crible sur $U$
engendré par les morphismes $f_i$}.
\end{definition}

\begin{definition}
\label{sites-definition-pullback-sieve}
Soit $\mathcal{C}$ une catégorie.
Soit $f : V \to U$ un morphisme de $\mathcal{C}$.
Soit $S \subset h_U$ un crible. Nous définissons
{\it l'image réciproque de $S$ par $f$} comme le crible
$S \times_U V$ de $V$ défini par la règle
$$
(\alpha : T \to V) \in (S \times_U V)(T)
\Leftrightarrow
(f \circ \alpha : T \to U) \in S(T)
$$
\end{definition}

\noindent
Nous laissons au lecteur le soin de vérifier qu'il s'agit bien d'un crible
(indication : utiliser l'équation \ref{sites-equation-property-sieve}).
Nous appelons aussi parfois $S \times_U V$ le {\it changement de base}
de $S$ par $f : V \to U$.

\begin{lemma}
\label{sites-lemma-pullback-sieve-section}
Soit $\mathcal{C}$ une catégorie.
Soit $U \in \Ob(\mathcal{C})$.
Soit $S$ un crible sur $U$.
Si $f : V \to U$ appartient à $S$, alors
$S \times_U V = h_V$ est maximal.
\end{lemma}

\begin{proof}
Cela résulte immédiatement des définitions.
\end{proof}

\begin{definition}
\label{sites-definition-topology}
Soit $\mathcal{C}$ une catégorie. Une {\it topologie sur $\mathcal{C}$} consiste
en une règle qui associe à tout $U \in \Ob(\mathcal{C})$
une partie $J(U)$ de l'ensemble de tous les cribles sur $U$, satisfaisant
aux conditions suivantes :
\begin{enumerate}
\item Pour tout morphisme $f : V \to U$ de $\mathcal{C}$ et
tout élément $S \in J(U)$, l'image réciproque $S \times_U V$
est un élément de $J(V)$.
\item Si $S$ et $S'$ sont des cribles sur $U \in \Ob(\mathcal{C})$,
si $S \in J(U)$, et si, pour tout $f \in S(V)$, l'image réciproque
$S' \times_U V$ appartient à $J(V)$, alors $S'$ appartient à $J(U)$.
\item Pour tout $U \in \Ob(\mathcal{C})$, le
crible maximal $S = h_U$ appartient à $J(U)$.
\end{enumerate}
\end{definition}

\noindent
Dans ce cas, les cribles appartenant à $J(U)$ sont appelés
les {\it cribles couvrants}.

\begin{lemma}
\label{sites-lemma-topology-basic}
Soit $\mathcal{C}$ une catégorie.
Soit $J$ une topologie sur $\mathcal{C}$.
Soit $U \in \Ob(\mathcal{C})$.
\begin{enumerate}
\item Les intersections finies d'éléments de $J(U)$ appartiennent à $J(U)$.
\item Si $S \in J(U)$ et $S' \supset S$, alors $S' \in J(U)$.
\end{enumerate}
\end{lemma}

\begin{proof}
Démonstration de (1). Le cas de l'intersection de la famille vide résulte de la condition
(3) de la définition \ref{sites-definition-topology}.
Soient $S, S' \in J(U)$. Considérons $S'' = S \cap S'$. Pour tout
$V \to U$ appartenant à $S(V)$, nous avons
$$
S' \times_U V = S'' \times_U V
$$
simplement parce que $V \to U$ appartient déjà à $S$. Le deuxième
axiome de la définition montre donc que $S'' \in J(U)$.

\medskip\noindent
Démonstration de (2). Soient $S \in J(U)$ et $S' \supset S$. Pour tout
$V \to U$ appartenant à $S(V)$, nous avons $S' \times_U V = h_V$ d'après le
lemme \ref{sites-lemma-pullback-sieve-section}. Ainsi,
$S' \times_U V \in J(V)$ par le troisième axiome. Par conséquent,
$S' \in J(U)$ par le deuxième axiome.
\end{proof}

\begin{definition}
\label{sites-definition-finer}
Soit $\mathcal{C}$ une catégorie. Soient $J$, $J'$ deux
topologies sur $\mathcal{C}$. Nous disons que $J$ est
{\it plus fine} ou {\it plus forte} que $J'$ si et seulement si, pour tout objet
$U$ de $\mathcal{C}$, nous avons $J'(U) \subset J(U)$.
Dans ce cas, nous disons aussi que $J'$ est
{\it moins fine} ou {\it plus faible} que $J$.
\end{definition}

\noindent
Autrement dit, tout crible couvrant pour $J'$ est un
crible couvrant pour $J$. Il existe une topologie la plus fine
sur $\mathcal{C}$, à savoir celle pour laquelle tout crible
est un crible couvrant. On l'appelle la
{\it topologie discrète} de $\mathcal{C}$.
Il existe également une topologie la moins fine.
Il s'agit de la topologie pour laquelle $J(U) = \{h_U\}$
pour tout objet $U$. On l'appelle la
{\it topologie chaotique} ou {\it topologie indiscrète}.

\begin{lemma}
\label{sites-lemma-play-with-topologies}
Soit $\mathcal{C}$ une catégorie.
Soit $\{J_i\}_{i\in I}$ un ensemble de topologies.
\begin{enumerate}
\item La règle $J(U) = \bigcap J_i(U)$ définit
une topologie sur $\mathcal{C}$.
\item Il existe une topologie la moins fine parmi celles qui sont plus fines que
toutes les topologies $J_i$.
\end{enumerate}
\end{lemma}

\begin{proof}
La première assertion résulte directement des définitions.
La seconde s'obtient en prenant l'intersection
de toutes les topologies plus fines que chacun des $J_i$.
\end{proof}

\noindent
Nous pouvons à présent définir,
sans aucune motivation, ce qu'est un faisceau.

\begin{definition}
\label{sites-definition-sheaf-sets-topology}
Soit $\mathcal{C}$ une catégorie munie d'une
topologie $J$. Soit $\mathcal{F}$ un préfaisceau d'ensembles
sur $\mathcal{C}$.
Nous disons que $\mathcal{F}$ est un
{\it faisceau} sur $\mathcal{C}$
si, pour tout $U \in \Ob(\mathcal{C})$ et pour
tout crible couvrant $S$ sur $U$, l'application canonique
$$
\Mor_{\textit{PSh}(\mathcal{C})}(h_U, \mathcal{F})
\longrightarrow
\Mor_{\textit{PSh}(\mathcal{C})}(S, \mathcal{F})
$$
est bijective.
\end{definition}

\noindent
Rappelons que le membre de gauche de la formule
affichée est égal à $\mathcal{F}(U)$. Autrement dit, $\mathcal{F}$
est un faisceau si et seulement si une section de $\mathcal{F}$
sur $U$ revient au même qu'une famille compatible de sections
$s_{T, \alpha} \in \mathcal{F}(T)$ paramétrée par
$(\alpha : T \to U) \in S(T)$, et ce pour tout crible couvrant $S$
sur $U$.

\begin{lemma}
\label{sites-lemma-topology-presheaves-sheaves}
Soit $\mathcal{C}$ une catégorie. Soit $\{ \mathcal{F}_i \}_{i\in I}$ une
famille de préfaisceaux d'ensembles sur $\mathcal{C}$. Pour chaque
$U \in \Ob(\mathcal{C})$, notons
$J(U)$ l'ensemble des cribles $S$ sur $U$ possédant la propriété suivante :
Pour tout morphisme $V \to U$, les applications
$$
\Mor_{\textit{PSh}(\mathcal{C})}(h_V, \mathcal{F}_i)
\longrightarrow
\Mor_{\textit{PSh}(\mathcal{C})}(S \times_U V, \mathcal{F}_i)
$$
sont bijectives pour tout $i \in I$. Alors $J$ définit une
topologie sur $\mathcal{C}$. Cette topologie est la plus fine
parmi celles pour lesquelles tous les $\mathcal{F}_i$ sont des faisceaux.
\end{lemma}

\begin{proof}
Remarquons d'abord que, si $J$ est une topologie, la dernière assertion du
lemme en résulte immédiatement.

\medskip\noindent
Pour tout $i \in I$ et tout $U \in \Ob(\mathcal{C})$, soit $J_i(U)$
l'ensemble des cribles $S$ sur $U$ possédant la propriété suivante :
Pour tout morphisme $V \to U$, l'application
$$
\Mor_{\textit{PSh}(\mathcal{C})}(h_V, \mathcal{F}_i)
\longrightarrow
\Mor_{\textit{PSh}(\mathcal{C})}(S \times_U V, \mathcal{F}_i)
$$
est bijective. Remarquons que $J(U) = \bigcap J_i(U)$. Ainsi, si nous montrons
que $J_i$ définit une topologie, nous concluons par le
lemme \ref{sites-lemma-play-with-topologies}.
Nous sommes donc ramenés au cas étudié dans le paragraphe suivant.

\medskip\noindent
Montrons que $J$ est une topologie lorsque notre famille est réduite à un seul
préfaisceau $\mathcal{F}$. Les premier et troisième axiomes d'une
topologie sont immédiatement vérifiés. Supposons donc que nous ayons
un objet $U$ et des cribles $S, S'$ sur $U$
tels que $S \in J(U)$ et que, pour tout $V \to U$ appartenant à $S(V)$,
nous ayons $S' \times_U V \in J(V)$. Nous devons montrer que
$S' \in J(U)$. Autrement dit, nous devons montrer que, pour
tout $f : W \to U$, l'application
$$
\mathcal{F}(W) =
\Mor_{\textit{PSh}(\mathcal{C})}(h_W, \mathcal{F})
\longrightarrow
\Mor_{\textit{PSh}(\mathcal{C})}(S' \times_U W, \mathcal{F})
$$
est bijective. Choisissons un élément $\varphi \in
\Mor_{\textit{PSh}(\mathcal{C})}(S' \times_U W, \mathcal{F})$.
Nous allons construire une section $s \in \mathcal{F}(W)$
qui s'envoie sur $\varphi$.

\medskip\noindent
Supposons que $\alpha : V \to W$ soit un élément de $S \times_U W$.
D'après la définition des images réciproques, nous voyons que
la composée $f \circ \alpha : V \to W  \to U$ appartient à $S$. Ainsi,
$S' \times_U V$ appartient à $J(V)$ par l'hypothèse faite sur le couple
de cribles $S, S'$. Nous avons alors un diagramme commutatif
de préfaisceaux
$$
\xymatrix{
S' \times_U V \ar[r] \ar[d] & h_V \ar[d] \\
S' \times_U W \ar[r] & h_W
}
$$
La restriction de $\varphi$ à $S' \times_U V$
correspond à un élément $s_{V, \alpha} \in \mathcal{F}(V)$.
Cela résulte de la définition de $J$
et du fait que $S' \times_U V$ appartient à $J(V)$.
Nous laissons au lecteur le soin de vérifier
que la règle $(V, \alpha) \mapsto s_{V, \alpha}$ définit
un élément
$\psi \in
\Mor_{\textit{PSh}(\mathcal{C})}(S \times_U W, \mathcal{F})$.
Puisque $S \in J(U)$, la définition de $J$ montre immédiatement
que $\psi$ correspond à un élément $s$ de $\mathcal{F}(W)$.

\medskip\noindent
Nous laissons au lecteur le soin de vérifier que la construction
$\varphi \mapsto s$ est inverse de l'application naturelle affichée ci-dessus.
\end{proof}

\begin{definition}
\label{sites-definition-canonical-topology}
Soit $\mathcal{C}$ une catégorie.
La topologie la plus fine sur $\mathcal{C}$ pour laquelle
tous les préfaisceaux représentables sont des faisceaux, voir le
lemme \ref{sites-lemma-topology-presheaves-sheaves},
est appelée la {\it topologie canonique} de $\mathcal{C}$.
\end{definition}














\section{La topologie définie par un site}
\label{sites-section-topology-site}

\noindent
Supposons que $\mathcal{C}$ soit une catégorie et que
$\text{Cov}_1(\mathcal{C})$ et $\text{Cov}_2(\mathcal{C})$
soient des ensembles de recouvrements définissant chacun une structure de site
sur $\mathcal{C}$. Dans cette situation, il peut arriver que
les catégories de faisceaux (d'ensembles) pour $\text{Cov}_1(\mathcal{C})$
et $\text{Cov}_2(\mathcal{C})$ soient les mêmes ; voir, par exemple, le
lemme \ref{sites-lemma-refine-same-topology}.

\medskip\noindent
Il se trouve que la catégorie des faisceaux sur $\mathcal{C}$
relativement à une topologie $J$
détermine la topologie $J$ et est déterminée par elle.
C'est un énoncé non trivial que nous traiterons plus loin ;
voir le théorème \ref{sites-theorem-topology-and-topos}.

\medskip\noindent
En l'admettant pour le moment, il est naturel d'étudier la
topologie déterminée par un site.

\begin{lemma}
\label{sites-lemma-site-gives-topology}
Soit $\mathcal{C}$ un site dont les recouvrements sont $\text{Cov}(\mathcal{C})$.
Pour tout objet $U$ de $\mathcal{C}$, notons $J(U)$
l'ensemble des cribles $S$ sur $U$ possédant la propriété suivante :
il existe un recouvrement
$\{f_i : U_i \to U\}_{i\in I} \in \text{Cov}(\mathcal{C})$
tel que le crible $S'$ engendré par les $f_i$ (voir la définition
\ref{sites-definition-sieve-generated}) soit contenu dans $S$.
\begin{enumerate}
\item Ce $J$ est une topologie sur $\mathcal{C}$.
\item Un préfaisceau $\mathcal{F}$ est un faisceau pour cette topologie
(voir la définition \ref{sites-definition-sheaf-sets-topology})
si et seulement si c'est un faisceau sur le site (voir la
définition \ref{sites-definition-sheaf-sets}).
\end{enumerate}
\end{lemma}

\begin{proof}
Pour démontrer la première assertion, il suffit de remarquer que les axiomes
(1), (2) et (3) de la définition d'un site
(définition \ref{sites-definition-site})
entraînent directement les axiomes
(3), (2) et (1) de la définition d'une topologie
(définition \ref{sites-definition-topology}). À titre d'exemple, nous
montrons que $J$ possède la propriété (2). Soit donc $U$ un objet
de $\mathcal{C}$, et soient $S, S'$ des cribles sur $U$ tels que
$S \in J(U)$ et tels que, pour tout $V \to U$ appartenant à $S(V)$,
nous ayons $S' \times_U V \in J(V)$. Par définition de $J(U)$,
nous pouvons trouver un recouvrement $\{f_i : U_i \to U\}$ du site
tel que $S$ l'image de $h_{U_i} \to h_U$ soit contenue
dans $S$. Comme chaque $S'\times_U U_i$ appartient à $J(U_i)$, nous
voyons qu'il existe des recouvrements $\{U_{ij} \to U_i\}$ du
site tels que $h_{U_{ij}} \to h_{U_i}$ soit contenu
dans $S' \times_U U_i$. Par définition du changement de base,
cela signifie que $h_{U_{ij}} \to h_U$ est contenu
dans le sous-préfaisceau $S' \subset h_U$. Par l'axiome (2) pour les
sites, nous voyons que $\{U_{ij} \to U\}$ est un recouvrement de
$U$, et nous concluons que $S' \in J(U)$ par définition de $J$.

\medskip\noindent
Soit $\mathcal{F}$ un préfaisceau. Supposons que $\mathcal{F}$
soit un faisceau pour la topologie $J$. Nous allons montrer que $\mathcal{F}$
est aussi un faisceau sur le site. Soit $\{f_i : U_i \to U\}_{i\in I}$
un recouvrement du site. Soit $s_i \in \mathcal{F}(U_i)$ une
famille de sections telle que
$s_i|_{U_i \times_U U_j} = s_j|_{U_i \times_U U_j}$ pour tous
$i, j$. Nous devons montrer qu'il existe une unique section
$s \in \mathcal{F}(U)$ qui se restreint en $s_i$ sur $U_i$.
Soit $S \subset h_U$ le crible engendré par les $f_i$.
Remarquons que $S \in J(U)$ par définition. Au lieu de construire
$s$, il suffit, d'après la condition de faisceau pour la topologie,
de construire un élément
$$
\varphi \in \Mor_{\textit{PSh}(\mathcal{C})}(S, \mathcal{F}).
$$
Prenons $\alpha \in S(T)$ pour un objet $T \in \mathcal{U}$.
Cela signifie exactement que $\alpha : T \to U$ est un morphisme
qui se factorise par $f_i$ pour un certain $i\in I$ (et peut-être plus de $1$).
Choisissons un tel indice $i$ et une factorisation $\alpha = f_i \circ \alpha_i$.
Posons $\varphi(\alpha) = \alpha_i^* s_i$. Si $i'$ et $\alpha
= f_i \circ \alpha_{i'}'$ donnent un second choix, alors
$\alpha_i^* s_i = (\alpha_{i'}')^* s_{i'}$ précisément en vertu de notre
condition $s_i|_{U_i \times_U U_j} = s_j|_{U_i \times_U U_j}$ pour tous
$i, j$. Ainsi, $\varphi(\alpha)$ est bien défini. Nous laissons au lecteur
le soin de vérifier que $\varphi$, qui à son tour détermine $s$, convient,
au sens où $s$ se restreint en $s_i$.

\medskip\noindent
Soit $\mathcal{F}$ un préfaisceau. Supposons que $\mathcal{F}$
soit un faisceau sur le site $(\mathcal{C}, \text{Cov}(\mathcal{C}))$.
Nous allons montrer que $\mathcal{F}$ est aussi un faisceau pour la topologie $J$.
Soit $U$ un objet de $\mathcal{C}$. Soit
$S$ un crible couvrant sur $U$ relativement à la topologie $J$.
Soit
$$
\varphi \in \Mor_{\textit{PSh}(\mathcal{C})}(S, \mathcal{F}).
$$
Nous devons montrer qu'il existe un unique élément de
$\mathcal{F}(U) = \Mor_{\textit{PSh}(\mathcal{C})}(h_U, \mathcal{F})$
dont la restriction est $\varphi$. Par définition, il existe
un recouvrement $\{f_i : U_i \to U\}_{i\in I} \in \text{Cov}(\mathcal{C})$
tel que $f_i : U_i \to U$ appartienne à $S(U_i)$. Ainsi,
nous pouvons poser $s_i = \varphi(f_i) \in \mathcal{F}(U_i)$.
C'est alors un exercice agréable de vérifier que
$s_i|_{U_i \times_U U_j} = s_j|_{U_i \times_U U_j}$ pour tous
$i, j$. Nous obtenons donc la section voulue $s$ par la condition de
faisceau pour $\mathcal{F}$ sur le site
$(\mathcal{C}, \text{Cov}(\mathcal{C}))$.
Les détails sont laissés au lecteur.
\end{proof}

\begin{definition}
\label{sites-definition-topology-associated-site}
Soit $\mathcal{C}$ un site dont les recouvrements sont $\text{Cov}(\mathcal{C})$.
La {\it topologie associée à $\mathcal{C}$} est la topologie
$J$ construite dans le lemme \ref{sites-lemma-site-gives-topology} ci-dessus.
\end{definition}

\noindent
Soit $\mathcal{C}$ une catégorie.
Soient $\text{Cov}_1(\mathcal{C})$ et $\text{Cov}_2(\mathcal{C})$
deux recouvrements définissant chacun une structure de site sur $\mathcal{C}$.
Il peut fort bien arriver que les topologies ainsi définies
soient les mêmes. Lorsque cela se produit,
nous disons que $\text{Cov}_1(\mathcal{C})$ et
$\text{Cov}_2(\mathcal{C})$ {\it définissent la même topologie} sur
$\mathcal{C}$. Dans ce cas, les catégories de faisceaux
sont les mêmes, par le lemme \ref{sites-lemma-site-gives-topology}.

\medskip\noindent
En général, nous ne nous intéressons qu'à la topologie définie
par une collection de recouvrements, et nous considérons la possibilité de choisir
différents ensembles de recouvrements comme un outil pour étudier la topologie.

\begin{remark}
\label{sites-remark-enlarge-coverings}
Élargissement de la classe des recouvrements.
Il est clair que, si $\text{Cov}(\mathcal{C})$
définit une structure de site sur $\mathcal{C}$, nous pouvons
ajouter à $\mathcal{C}$ n'importe quel ensemble de familles de morphismes de but fixé
tautologiquement équivalentes
(voir la définition \ref{sites-definition-combinatorial-tautological})
à des éléments de $\text{Cov}(\mathcal{C})$ sans changer la topologie.
\end{remark}

\begin{remark}
\label{sites-remark-shrink-coverings}
Réduction de la classe des recouvrements. Soit $\mathcal{C}$ un site.
Considérons l'ensemble
$$
\mathcal{S} = P(\text{Arrows}(\mathcal{C})) \times \Ob(\mathcal{C})
$$
où $P(\text{Arrows}(\mathcal{C}))$ est l'ensemble des parties de l'ensemble des morphismes,
c'est-à-dire l'ensemble de tous les ensembles de morphismes.
Soit $\mathcal{S}_\tau \subset \mathcal{S}$
la partie formée des $(T, U) \in \mathcal{S}$ tels que
(a) tout $\varphi \in T$ ait pour but $U$,
(b) la famille $\{\varphi\}_{\varphi \in T}$ soit tautologiquement
équivalente (voir la définition \ref{sites-definition-combinatorial-tautological})
à un recouvrement de $\text{Cov}(\mathcal{C})$.
Il est clair qu'en considérant les éléments de $\mathcal{S}_\tau$ comme
les recouvrements, nous n'obtenons pas exactement la notion de site
telle qu'elle est définie dans la définition \ref{sites-definition-site}.
La structure $(\mathcal{C}, \mathcal{S}_\tau)$
ainsi obtenue satisfait des conditions légèrement modifiées. Ces
conditions modifiées sont les suivantes :
\begin{enumerate}
\item[(0')] $\text{Cov}(\mathcal{C}) \subset
P(\text{Arrows}(\mathcal{C})) \times \Ob(\mathcal{C})$,
\item[(1')] Si $V \to U$ est un isomorphisme, alors
$(\{V \to U\}, U) \in \text{Cov}(\mathcal{C})$.
\item[(2')] Si $(T, U) \in \text{Cov}(\mathcal{C})$
et si, pour $f : U' \to U$ dans $T$, on s'est donné
$(T_f, U') \in \text{Cov}(\mathcal{C})$,
alors, en posant $T' = \{f \circ f' \mid f \in T,\ f' \in T_f\}$,
nous obtenons $(T', U) \in \text{Cov}(\mathcal{C})$.
\item[(3')] Si $(T, U) \in \text{Cov}(\mathcal{C})$ et si $g : V \to U$
est un morphisme de $\mathcal{C}$, alors
\begin{enumerate}
\item $U' \times_{f, U, g} V$ existe pour $f : U' \to U$ dans $T$, et
\item en posant
$T' = \{\text{pr}_2 : U' \times_{f, U, g} V \to V \mid f : U' \to U \in T\}$
pour un certain choix de produits fibrés, nous obtenons $(T', V) \in \text{Cov}(\mathcal{C})$.
\end{enumerate}
\end{enumerate}
Il est alors facile de vérifier que, partant d'une structure satisfaisant
aux conditions (0') -- (3') ci-dessus, puis en élargissant convenablement
$\text{Cov}(\mathcal{C})$ (comparer avec Ensembles,
section \ref{sets-section-coverings-site}), nous obtenons un site.
Manifestement, il y a peu de différence entre cette notion et la
notion usuelle de site, du moins du point de vue de la
topologie. Elle présente deux avantages :
en raison de la condition (0') ci-dessus, les recouvrements forment automatiquement
un ensemble, et, toujours en raison de (0'), la totalité des structures
de ce type forme elle aussi un ensemble.
Le prix à payer est qu'il faut écrire partout
``tautologiquement équivalente''.
\end{remark}

















\section{Faisceautisation pour une topologie}
\label{sites-section-sheafification-topology}

\noindent
Dans cette section, nous expliquons l'analogue de la
construction de faisceautisation pour une topologie.

\medskip\noindent
Soit $\mathcal{C}$ une catégorie.
Soit $J$ une topologie sur $\mathcal{C}$.
Soit $\mathcal{F}$ un préfaisceau d'ensembles.
Pour tout $U \in \Ob(\mathcal{C})$, nous
définissons
$$
L\mathcal{F}(U)
=
\colim_{S \in J(U)^{opp}}
\Mor_{\textit{PSh}(\mathcal{C})}(S, \mathcal{F})
$$
comme une colimite. Nous considérons ici $J(U)$ comme un ensemble
partiellement ordonné par inclusion, voir le lemme \ref{sites-lemma-sieves-set}. Les
applications de transition du système sont définies comme suit.
Si $S \subset S'$ appartiennent à $J(U)$, alors $S \to S'$ est
un morphisme de préfaisceaux. Il existe donc une application naturelle
de restriction
$$
\Mor_{\textit{PSh}(\mathcal{C})}(S', \mathcal{F})
\longrightarrow
\Mor_{\textit{PSh}(\mathcal{C})}(S, \mathcal{F}).
$$
Ainsi, nous voyons que
$S \mapsto \Mor_{\textit{PSh}(\mathcal{C})}(S, \mathcal{F})$
est un système inductif au sens de Catégories,
définition \ref{categories-definition-system-over-poset},
à condition de renverser l'ordre
sur $J(U)$ (ce qu'indique l'exposant ${}^{opp}$).
En particulier, puisque $h_U \in J(U)$,
il existe une application canonique
$$
\ell : \mathcal{F}(U) \longrightarrow L\mathcal{F}(U)
$$
provenant de l'identification
$\mathcal{F}(U) = \Mor_{\textit{PSh}(\mathcal{C})}(h_U, \mathcal{F})$.
En outre, la colimite définissant $L\mathcal{F}(U)$ est une colimite filtrante,
car, pour toute paire de cribles couvrants $S, S'$ sur $U$, le
crible $S \cap S'$ est encore couvrant, voir le lemme \ref{sites-lemma-sieves-set}.

\medskip\noindent
Soit $f : V \to U$ un morphisme de $\mathcal{C}$.
Soit $S \in J(U)$. Il existe un diagramme commutatif
$$
\xymatrix{
S \times_U V \ar[r] \ar[d] & h_V \ar[d] \\
S \ar[r] & h_U
}
$$
Nous pouvons utiliser la flèche verticale de gauche pour obtenir des applications canoniques de restriction
$$
\Mor_{\textit{PSh}(\mathcal{C})}(S, \mathcal{F})
\to \Mor_{\textit{PSh}(\mathcal{C})}(S \times_U V, \mathcal{F}).
$$
Le changement de base $S \mapsto S \times_U V$ induit une application
croissante $J(U) \to J(V)$. Les applications de restriction
définissent une transformation de foncteurs comme dans
Catégories, lemme \ref{categories-lemma-functorial-colimit}.
Nous obtenons donc une application naturelle de restriction
$$
L\mathcal{F}(U) \longrightarrow L\mathcal{F}(V).
$$

\begin{lemma}
\label{sites-lemma-L-presheaf}
Dans la situation ci-dessus :
\begin{enumerate}
\item La correspondance $U \mapsto L\mathcal{F}(U)$, munie des
applications de restriction définies ci-dessus, est un préfaisceau.
\item Les applications $\ell$ se recollent en un morphisme de préfaisceaux
$\ell : \mathcal{F} \to L\mathcal{F}$.
\item La règle $\mathcal{F} \mapsto (\mathcal{F} \xrightarrow{\ell}
L\mathcal{F})$ est un foncteur.
\item Si $\mathcal{F}$ est un sous-préfaisceau de $\mathcal{G}$, alors
$L\mathcal{F}$ est un sous-préfaisceau de $L\mathcal{G}$.
\item L'application $\ell : \mathcal{F} \to L\mathcal{F}$ possède la
propriété suivante : pour toute section $s \in L\mathcal{F}(U)$,
il existe un crible couvrant $S$ sur $U$ et un élément
$\varphi \in \Mor_{\textit{PSh}(\mathcal{C})}(S, \mathcal{F})$
tels que $\ell(\varphi)$ soit égal à la restriction de
$s$ à $S$.
\end{enumerate}
\end{lemma}

\begin{proof}
La démonstration est omise.
\end{proof}

\begin{definition}
\label{sites-definition-presheaf-separated-topology}
Soit $\mathcal{C}$ une catégorie.
Soit $J$ une topologie sur $\mathcal{C}$.
Nous disons qu'un préfaisceau d'ensembles $\mathcal{F}$
est {\it séparé} si, pour tout objet $U$ et
tout crible couvrant $S$ sur $U$, l'application canonique
$\mathcal{F}(U) \to \Mor_{\textit{PSh}(\mathcal{C})}(S, \mathcal{F})$
est injective.
\end{definition}

\begin{theorem}
\label{sites-theorem-L-topology}
Soit $\mathcal{C}$ une catégorie.
Soit $J$ une topologie sur $\mathcal{C}$.
Soit $\mathcal{F}$ un préfaisceau d'ensembles.
\begin{enumerate}
\item Le préfaisceau $L\mathcal{F}$ est séparé.
\item Si $\mathcal{F}$ est séparé, alors $L\mathcal{F}$ est un faisceau
et le morphisme de préfaisceaux $\mathcal{F} \to L\mathcal{F}$ est injectif.
\item Si $\mathcal{F}$ est un faisceau, alors $\mathcal{F} \to L\mathcal{F}$
est un isomorphisme.
\item Le préfaisceau $LL\mathcal{F}$ est toujours un faisceau.
\end{enumerate}
\end{theorem}

\begin{proof}
L'assertion (3) résulte immédiatement de la définition de $L$ et
de celle d'un faisceau (définition \ref{sites-definition-sheaf-sets-topology}).
L'assertion (4) résulte formellement des autres.

\medskip\noindent
Esquissons la démonstration de (1). Soit $S$ un crible couvrant
sur l'objet $U$. Supposons que $\varphi_i \in L\mathcal{F}(U)$,
$i = 1, 2$, aient la même image dans
$\Mor_{\textit{PSh}(\mathcal{C})}(S, L\mathcal{F})$.
Nous pouvons trouver un même crible couvrant $S'$ sur $U$
tel que les deux $\varphi_i$ soient représentés par des éléments
$\varphi_i \in \Mor_{\textit{PSh}(\mathcal{C})}(S', \mathcal{F})$.
Nous pouvons supposer que $S' = S$ en remplaçant $S$ et $S'$ par
$S' \cap S$, qui est encore un crible couvrant, voir le lemme \ref{sites-lemma-sieves-set}.
Soient $V\in \Ob(\mathcal{C})$ et
$\alpha : V \to U$ dans $S(V)$.
Alors $S \times_U V = h_V$,
voir le lemme \ref{sites-lemma-pullback-sieve-section}. Ainsi, les restrictions
des $\varphi_i$ par $V \to U$ correspondent à des sections $s_{i, V, \alpha}$
de $\mathcal{F}$ sur $V$. L'hypothèse affirme qu'il existe
un crible couvrant $S_{V, \alpha}$ sur $V$ tel que
les $s_{i, V, \alpha}$ se restreignent au même élément de
$\Mor_{\textit{PSh}(\mathcal{C})}(S_{V, \alpha}, \mathcal{F})$.
Considérons le crible $S''$ sur $U$ défini par la règle
\begin{eqnarray}
\label{sites-equation-S-prime-prime}
(f : T \to U) \in S''(T)
& \Leftrightarrow &
\exists\ V , \ \alpha : V \to U, \ \alpha \in S(V), \nonumber \\
& &
\exists\ g : T \to V, \ g \in S_{V, \alpha}(T), \\
& &
f = \alpha \circ g \nonumber
\end{eqnarray}
L'axiome (2) d'une topologie montre que $S''$ est un crible
couvrant sur $U$. Par construction, $\varphi_1$
et $\varphi_2$ se restreignent au même élément de
$\Mor_{\textit{PSh}(\mathcal{C})}(S'', L\mathcal{F})$,
comme souhaité.

\medskip\noindent
Esquissons la démonstration de (2). Supposons que $\mathcal{F}$ soit un
préfaisceau d'ensembles séparé sur $\mathcal{C}$ relativement à
la topologie $J$.
Soit $S$ un crible couvrant sur l'objet $U$ de $\mathcal{C}$.
Supposons que $\varphi \in \Mor_\mathcal{C}(S, L\mathcal{F})$.
Nous devons trouver un élément $s \in L\mathcal{F}(U)$ dont la restriction
soit $\varphi$. Soient $V\in \Ob(\mathcal{C})$ et
$\alpha : V \to U$ dans $S(V)$. La valeur $\varphi(\alpha)
\in L\mathcal{F}(V)$ est donnée par un crible couvrant
$S_{V, \alpha}$ sur $V$ et un morphisme de préfaisceaux
$\varphi_{V, \alpha} : S_{V, \alpha} \to \mathcal{F}$.
Comme dans la démonstration précédente, définissons un crible couvrant $S''$ sur $U$ par
l'équation (\ref{sites-equation-S-prime-prime}). Nous définissons
$$
\varphi'' : S'' \longrightarrow \mathcal{F}
$$
par la règle simple suivante : pour tout $f : T \to U$,
$f \in S''(T)$, choisissons $V, \alpha, g$ comme dans
l'équation (\ref{sites-equation-S-prime-prime}). Posons alors
$$
\varphi''(f) = \varphi_{V, \alpha}(g).
$$
Nous affirmons que ceci est indépendant du
choix de $V, \alpha, g$.
Considérons un second choix$ V', \alpha', g'$.
Les restrictions de $\varphi_{V, \alpha}$ et de
$\varphi_{V', \alpha'}$ à l'intersection
des cribles couvrants suivants sur $T$
$$
(S_{V, \alpha} \times_{V, g} T) \cap (S_{V', \alpha'} \times_{V', g'} T)
$$
coïncident. En effet, ces deux restrictions correspondent à la
restriction de $\varphi$ à $T$ (par $f$), et l'égalité voulue
résulte de ce que $\mathcal{F}$ est séparé.
Notons $\psi$ la restriction commune.
L'indépendance du choix résulte de
$\varphi_{V, \alpha}(g) = \psi(\text{id}_T) =
\varphi_{V', \alpha'}(g')$. Ainsi, $\varphi''$
définit un élément $s \in L\mathcal{F}(U)$. Nous laissons au
lecteur le soin de vérifier que la restriction de $s$ est $\varphi$.
\end{proof}

\begin{definition}
\label{sites-definition-associated-sheaf-topology}
Soit $\mathcal{C}$ une catégorie munie d'une topologie $J$.
Soit $\mathcal{F}$ un préfaisceau d'ensembles sur $\mathcal{C}$.
Le faisceau $\mathcal{F}^\# := LL\mathcal{F}$,
muni du morphisme canonique $\mathcal{F} \to \mathcal{F}^\#$,
est appelé le {\it faisceau associé à $\mathcal{F}$}.
\end{definition}

\begin{proposition}
\label{sites-proposition-sheafification-adjoint-topology}
Soit $\mathcal{C}$ une catégorie munie d'une topologie.
Soit $\mathcal{F}$ un préfaisceau d'ensembles sur $\mathcal{C}$.
Le morphisme canonique $\mathcal{F} \to \mathcal{F}^\#$ possède la
propriété universelle suivante : pour tout morphisme
$\mathcal{F} \to \mathcal{G}$,
où $\mathcal{G}$ est un faisceau d'ensembles, il existe un unique morphisme
$\mathcal{F}^\# \to \mathcal{G}$ tel que $\mathcal{F} \to \mathcal{F}^\#
\to \mathcal{G}$ soit le morphisme donné.
\end{proposition}

\begin{proof}
La démonstration est la même que celle de la proposition \ref{sites-proposition-sheafification-adjoint}.
\end{proof}














\section{Topologies et faisceaux}
\label{sites-section-topology-and-sheaves}

\begin{lemma}
\label{sites-lemma-sieve-sheafification}
Soit $\mathcal{C}$ une catégorie munie d'une topologie $J$.
Soit $U$ un objet de $\mathcal{C}$.
Soit $S$ un crible sur $U$. Les propriétés suivantes sont équivalentes :
\begin{enumerate}
\item Le crible $S$ est couvrant.
\item La faisceautisation $S^\# \to h_U^\#$
du morphisme $S \to h_U$ est un isomorphisme.
\end{enumerate}
\end{lemma}

\begin{proof}
Commençons par quelques remarques générales.
Nous utiliserons les égalités $S^\# = LLS$ et $h_U^\# = LLh_U$.
En particulier, le lemme \ref{sites-lemma-L-presheaf} montre que
$S^\# \to h_U^\#$ est injective. Remarquons que
$\text{id}_U \in h_U(U)$. Cet élément donne donc naissance à
des sections de $Lh_U$ et de $h_U^\# = LLh_U$ sur $U$ que
nous noterons encore $\text{id}_U$.

\medskip\noindent
Supposons que $S$ soit un crible couvrant. Il suffit manifestement de
trouver un morphisme $h_U \to S^\#$ tel que le composé
$h_U \to h_U^\#$ soit le morphisme canonique. Pour trouver un tel morphisme,
il suffit de trouver une section $s \in S^\#(U)$ dont la restriction
soit $\text{id}_U$. Or, puisque $S$ est un
crible couvrant, l'élément
$\text{id}_S \in \Mor_{\textit{PSh}(\mathcal{C})}(S, S)$
donne naissance à une section de $LS$ sur $U$ dont la restriction
est $\text{id}_U$ dans $Lh_U$. Ceci conclut.

\medskip\noindent
Supposons que $S^\# \to h_U^\#$ soit un isomorphisme.
Soit $1 \in S^\#(U)$ l'élément correspondant à
$\text{id}_U$ dans $h_U^\#(U)$. Puisque $S^\# = LLS$,
il existe un crible couvrant $S'$ sur $U$ tel que
$1$ provienne d'un élément
$$
\varphi \in \Mor_{\textit{PSh}(\mathcal{C})}(S', LS).
$$
Cela signifie à son tour que, pour tout $\alpha : V \to U$,
$\alpha\in S'(V)$, il existe un crible couvrant $S_{V, \alpha}$
sur $V$ tel que $\varphi(\alpha)$ corresponde à
un morphisme de préfaisceaux $S_{V, \alpha} \to S$. Autrement dit,
$S_{V, \alpha}$ est contenu dans $S \times_U V$. Le second
axiome d'une topologie montre que $S$ est un crible couvrant.
\end{proof}

\begin{theorem}
\label{sites-theorem-topology-and-topos}
Soit $\mathcal{C}$ une catégorie.
Soient $J$ et $J'$ des topologies sur $\mathcal{C}$.
Les propriétés suivantes sont équivalentes :
\begin{enumerate}
\item $J = J'$,
\item les faisceaux pour la topologie $J$ sont les mêmes que
les faisceaux pour la topologie $J'$.
\end{enumerate}
\end{theorem}

\begin{proof}
Il est tautologique que, si $J = J'$, les notions de faisceau
coïncident. Réciproquement, le lemme \ref{sites-lemma-sieve-sheafification}
caractérise les cribles couvrants au moyen du foncteur de faisceautisation.
Or le foncteur de faisceautisation
$\textit{PSh}(\mathcal{C}) \to \Sh(\mathcal{C}, J)$
est l'adjoint à gauche du foncteur d'inclusion
$\Sh(\mathcal{C}, J) \to \textit{PSh}(\mathcal{C})$.
Par conséquent, si les sous-catégories
$\Sh(\mathcal{C}, J)$ et
$\Sh(\mathcal{C}, J')$ sont les mêmes, alors les foncteurs de faisceautisation
sont les mêmes, et les familles de cribles couvrants sont donc
les mêmes.
\end{proof}

\begin{lemma}
\label{sites-lemma-finer-topology}
Hypothèses et notations comme dans le théorème \ref{sites-theorem-topology-and-topos}.
Alors $J \subset J'$ si et seulement si tout faisceau pour la
topologie $J'$ est un faisceau pour la topologie $J$.
\end{lemma}

\begin{proof}
Une implication est claire. Pour la réciproque, supposons que
$\Sh(\mathcal{C}, J') \subset \Sh(\mathcal{C}, J)$.
Par les propriétés formelles, cela implique
que, si $\mathcal{F}$ est un préfaisceau d'ensembles,
et si $\mathcal{F} \to \mathcal{F}^\#$,
resp.\ $\mathcal{F} \to \mathcal{F}^{\#, \prime}$,
est la faisceautisation relativement à $J$, resp.\ $J'$, alors il existe
un morphisme canonique $\mathcal{F}^\# \to \mathcal{F}^{\#, \prime}$
tel que
$\mathcal{F} \to \mathcal{F}^\# \to \mathcal{F}^{\#, \prime}$
soit le morphisme canonique $\mathcal{F} \to \mathcal{F}^{\#, \prime}$.
Bien entendu, $\mathcal{F}^\# \to \mathcal{F}^{\#, \prime}$
identifie le second faisceau à la faisceautisation du premier
relativement à la topologie $J'$.
Appliquons ceci au morphisme $S \to h_U$ du
lemme \ref{sites-lemma-sieve-sheafification}. Nous obtenons un diagramme
commutatif
$$
\xymatrix{
S \ar[r] \ar[d] &
S^\# \ar[r] \ar[d] &
S^{\#, \prime} \ar[d] \\
h_U \ar[r] &
h_U^\# \ar[r] &
h_U^{\#, \prime}
}
$$
De plus, si $S$ est un crible couvrant pour la topologie $J$,
alors la flèche verticale médiane est un isomorphisme (par le lemme),
et nous en déduisons que la flèche verticale de droite est un isomorphisme,
car elle est la faisceautisation de celle du milieu relativement à $J'$.
Le lemme montre de nouveau que $S$ est aussi un crible couvrant
pour $J'$.
\end{proof}




\section{Topologies et foncteurs continus}
\label{sites-section-topologies-continuous-functors}

\noindent
Expliquer comment un foncteur continu donne une paire adjointe
de foncteurs sur les faisceaux.





\section{Points et topologies}
\label{sites-section-points-topologies}

\noindent
Rappelons, d'après la section \ref{sites-section-points}, qu'étant donné un foncteur
$p = u : \mathcal{C} \to \textit{Ensembles}$, nous pouvons définir
un foncteur fibre
$$
\textit{PSh}(\mathcal{C}) \longrightarrow \textit{Ensembles},
\mathcal{F} \longmapsto \mathcal{F}_p.
$$

\begin{definition}
\label{sites-definition-point-topology}
Soit $\mathcal{C}$ une catégorie.
Soit $J$ une topologie sur $\mathcal{C}$.
Un {\it point $p$} de la topologie est la donnée d'un foncteur
$u : \mathcal{C} \to \textit{Ensembles}$ tel que
\begin{enumerate}
\item Pour tout crible couvrant $S$ sur $U$, l'application
$S_p \to (h_U)_p$ est surjective.
\item Le foncteur fibre $\Sh(\mathcal{C}) \to \textit{Ensembles}$,
$\mathcal{F} \to \mathcal{F}_p$, est exact.
\end{enumerate}
\end{definition}
















% OK, start here.
%

\chapter{Champs}



\phantomsection
\label{stacks-section-phantom}


\section{Introduction}
\label{stacks-section-introduction}

\noindent
Dans ce très court chapitre, nous introduisons les champs et les
champs en groupoïdes. Voir \cite{DM} et \cite{Vis2}.


\section{Préfaisceaux de morphismes associés aux catégories fibrées}
\label{stacks-section-morphisms}

\noindent
Soit $\mathcal{C}$ une catégorie.
Soit $p : \mathcal{S} \to \mathcal{C}$ une catégorie fibrée,
voir Catégories, section \ref{categories-section-fibred-categories}.
Supposons que $x, y\in \Ob(\mathcal{S}_U)$ soient des
objets de la catégorie fibre au-dessus de $U$. Nous allons définir
un foncteur
$$
\mathit{Mor}(x, y) : (\mathcal{C}/U)^{opp} \longrightarrow \textit{Ensembles}.
$$
Autrement dit, ce sera un préfaisceau sur $\mathcal{C}/U$, voir
Sites, définition \ref{sites-definition-presheaf}.
Faisons un choix d'images inverses comme dans
Catégories,
définition \ref{categories-definition-pullback-functor-fibred-category}.
Alors, pour $f : V \to U$, nous posons
$$
\mathit{Mor}(x, y)(f : V \to U) =
\Mor_{\mathcal{S}_V}(f^\ast x, f^\ast y).
$$
Soit $f' : V' \to U$ un second objet de $\mathcal{C}/U$.
Nous devons aussi définir l'application de restriction correspondant à un
morphisme $g : V'/U  \to V/U$ dans $\mathcal{C}/U$,
autrement dit $g : V' \to V$ et $f' = f \circ g$.
Ce sera une application
$$
\Mor_{\mathcal{S}_V}(f^\ast x, f^\ast y)
\longrightarrow
\Mor_{\mathcal{S}_{V'}}({f'}^\ast x, {f'}^\ast y), \quad
\phi \longmapsto \phi|_{V'}
$$
Cette application sera essentiellement $g^\ast$, à ceci près qu'elle transforme
un élément $\phi$ du membre de gauche en un élément
$g^\ast \phi$
de $\Mor_{\mathcal{S}_{V'}}(g^\ast f^\ast x, g^\ast f^\ast y)$.
À ce stade, nous utilisons la transformation $\alpha_{g, f}$ du
lemme \ref{categories-lemma-fibred} de Catégories.
En formule, l'application de restriction est décrite par
$$
\phi|_{V'} =
(\alpha_{g, f})_y^{-1} \circ
g^\ast \phi \circ
(\alpha_{g, f})_x.
$$
Bien entendu, personne ne conçoit cette application de restriction de cette
manière. Nous ne le ferons qu'une fois, afin de vérifier le lemme suivant.

\begin{lemma}
\label{stacks-lemma-painful}
Cela définit effectivement un préfaisceau.
\end{lemma}

\begin{proof}
Soient $g : V'/U \to V/U$ comme ci-dessus et, de même,
$g' : V''/U \to V'/U$ des morphismes de $\mathcal{C}/U$.
Ainsi $f' = f \circ g$ et $f'' = f' \circ g' = f \circ g \circ g'$.
Soit $\phi \in \Mor_{\mathcal{S}_V}(f^\ast x, f^\ast y)$.
Nous avons alors
\begin{eqnarray*}
& &
(\alpha_{g \circ g', f})_y^{-1} \circ
(g \circ g')^\ast \phi \circ
(\alpha_{g \circ g', f})_x
\\
& = &
(\alpha_{g \circ g', f})_y^{-1} \circ
(\alpha_{g', g})_{f^*y}^{-1} \circ
(g')^*g^\ast \phi \circ
(\alpha_{g', g})_{f^*x} \circ
(\alpha_{g \circ g', f})_x
\\
& = &
(\alpha_{g', f'})_y^{-1} \circ
(g')^*(\alpha_{g, f})_y^{-1} \circ
(g')^* g^\ast \phi \circ
(g')^*(\alpha_{g, f})_x
\circ
(\alpha_{g', f'})_x
\\
& = &
(\alpha_{g', f'})_y^{-1} \circ
(g')^*\Big(
(\alpha_{g, f})_y^{-1} \circ
g^\ast \phi \circ
(\alpha_{g, f})_x
\Big) \circ
(\alpha_{g', f'})_x
\end{eqnarray*}
ce qui est précisément l'égalité voulue,
$\phi|_{V''} = (\phi|_{V'})|_{V''}$.
La première égalité vaut parce que
$\alpha_{g', g}$ est une transformation de foncteurs, de sorte que
$$
\xymatrix{
(g \circ g')^*f^*x
\ar[rr]_{(g \circ g')^\ast \phi}
\ar[d]_{(\alpha_{g', g})_{f^*x}} & &
(g \circ g')^*f^*y
\ar[d]^{(\alpha_{g', g})_{f^*y}} \\
(g')^*g^*f^*x
\ar[rr]^{(g')^*g^\ast \phi} & &
(g')^*g^*f^*y
}
$$
est commutatif. La deuxième égalité résulte de la propriété (d) d'un
pseudo-foncteur, puisque $f' = f \circ g$ (voir
Catégories, définition \ref{categories-definition-functor-into-2-category}).
La dernière égalité découle du fait que $(g')^*$ est un foncteur.
\end{proof}

\noindent
Désormais, nous omettrons souvent de mentionner les transformations
$\alpha_{g, f}$ et identifierons simplement les foncteurs
$g^* \circ f^*$ et $(f \circ g)^*$. En particulier,
pour $g : V'/U \to V/U$, les applications de restriction
du préfaisceau $\mathit{Mor}(x, y)$ seront parfois notées
$\phi \mapsto g^*\phi$. Nous formalisons cette construction dans une définition.

\begin{definition}
\label{stacks-definition-mor-presheaf}
Soit $\mathcal{C}$ une catégorie.
Soit $p : \mathcal{S} \to \mathcal{C}$ une catégorie fibrée,
voir Catégories, section \ref{categories-section-fibred-categories}.
Étant donnés un objet $U$ de $\mathcal{C}$ et des objets
$x$, $y$ de la catégorie fibre, le {\it préfaisceau
des morphismes de $x$ vers $y$} est le préfaisceau
$$
(f : V \to U) \longmapsto \Mor_{\mathcal{S}_V}(f^*x, f^*y)
$$
décrit ci-dessus. On le note $\mathit{Mor}(x, y)$.
Le sous-préfaisceau $\mathit{Isom}(x, y)$ dont la valeur
au-dessus de $V$ est l'ensemble des isomorphismes
$f^*x \to f^*y$ dans la catégorie fibre $\mathcal{S}_V$
est appelé le {\it préfaisceau des isomorphismes de $x$ vers $y$}.
\end{definition}

\noindent
Si $\mathcal{S}$ est fibrée en groupoïdes, alors bien entendu
$\mathit{Isom}(x, y) = \mathit{Mor}(x, y)$, et il est
d'usage d'employer la notation $\mathit{Isom}$.

\begin{lemma}
\label{stacks-lemma-presheaf-mor-map-fibred-categories}
Soit $F : \mathcal{S}_1 \to \mathcal{S}_2$ un $1$-morphisme de catégories
fibrées au-dessus de la catégorie $\mathcal{C}$. Soient
$U \in \Ob(\mathcal{C})$ et $x, y\in \Ob((\mathcal{S}_1)_U)$.
Alors $F$ définit un morphisme canonique de préfaisceaux
$$
\mathit{Mor}_{\mathcal{S}_1}(x, y)
\longrightarrow
\mathit{Mor}_{\mathcal{S}_2}(F(x), F(y))
$$
sur $\mathcal{C}/U$.
\end{lemma}

\begin{proof}
D'après la définition \ref{categories-definition-fibred-categories-over-C}
de Catégories, le foncteur $F$ envoie les morphismes fortement cartésiens sur
des morphismes fortement cartésiens. Ainsi, si $f : V \to U$ est un morphisme
de $\mathcal{C}$, il existe des isomorphismes canoniques
$\alpha_V : f^*F(x) \to F(f^*x)$,
$\beta_V : f^*F(y) \to F(f^*y)$ tels que $f^*F(x) \to F(f^*x) \to F(x)$
soit le morphisme canonique $f^*F(x) \to F(x)$, et de même pour $\beta_V$.
Nous pouvons donc définir
$$
\xymatrix{
\mathit{Mor}_{\mathcal{S}_1}(x, y)(f : V \to U) \ar@{=}[r] &
\Mor_{\mathcal{S}_{1, V}}(f^\ast x, f^\ast y) \ar[d] \\
\mathit{Mor}_{\mathcal{S}_2}(F(x), F(y))(f : V \to U) \ar@{=}[r] &
\Mor_{\mathcal{S}_{2, V}}(f^\ast F(x), f^\ast F(y))
}
$$
par $\phi \mapsto \beta_V^{-1} \circ F(\phi) \circ \alpha_V$.
Nous omettons la vérification de la compatibilité avec les applications de
restriction.
\end{proof}

\begin{remark}
\label{stacks-remark-alternative}
Supposons que $p : \mathcal{S} \to \mathcal{C}$ soit fibrée en groupoïdes.
Dans ce cas, nous pouvons démontrer le lemme \ref{stacks-lemma-painful}
à l'aide du lemme \ref{categories-lemma-fibred-strict} de Catégories,
qui dit que $\mathcal{S} \to \mathcal{C}$ est équivalente à la catégorie
associée à un foncteur contravariant
$F : \mathcal{C} \to \textit{Groupoïdes}$.
Dans le cas de la catégorie fibrée associée à $F$,
nous avons exactement $g^* \circ f^* = (f \circ g)^*$,
et il n'est pas nécessaire d'utiliser les applications $\alpha_{g, f}$.
Dans ce cas, le lemme est (encore plus) trivial. Bien entendu, on utilise alors
le fait que le préfaisceau $\mathit{Mor}(x, y)$ reste inchangé lorsque l'on
passe à une catégorie fibrée équivalente, ce qui découle du
lemme \ref{stacks-lemma-presheaf-mor-map-fibred-categories}.
\end{remark}

\begin{lemma}
\label{stacks-lemma-isom-as-2-fibre-product}
Soit $\mathcal{C}$ une catégorie.
Soit $p : \mathcal{S} \to \mathcal{C}$ une catégorie fibrée,
voir Catégories, section \ref{categories-section-fibred-categories}.
Soit $U \in \Ob(\mathcal{C})$ et soient $x, y \in \Ob(\mathcal{S}_U)$.
Notons encore $x, y : \mathcal{C}/U \to \mathcal{S}$ les
$1$-morphismes correspondants, voir
Catégories, lemme \ref{categories-lemma-yoneda-2category}.
Alors
\begin{enumerate}
\item le $2$-produit fibré
$\mathcal{S} \times_{\mathcal{S} \times \mathcal{S}, (x, y)} \mathcal{C}/U$
est fibré en setoïdes au-dessus de $\mathcal{C}/U$, et
\item $\mathit{Isom}(x, y)$ est le préfaisceau d'ensembles correspondant
à cette catégorie fibrée en setoïdes, voir
Catégories, lemme \ref{categories-lemma-2-category-fibred-setoids}.
\end{enumerate}
\end{lemma}

\begin{proof}
Omis. Indication : les objets du $2$-produit fibré sont
$(a : V \to U, z, (\alpha, \beta))$, où
$\alpha : z \to a^*x$ et $\beta : z \to a^*y$ sont des isomorphismes
dans $\mathcal{S}_V$. Le lien avec $\mathit{Isom}(x, y)$
s'obtient donc en associant à un tel objet l'isomorphisme
$\beta \circ \alpha^{-1}$.
\end{proof}







\section{Données de descente dans les catégories fibrées}
\label{stacks-section-descent-data}

\noindent
Dans cette section, nous définissons la notion de donnée de descente
dans le cadre abstrait d'une catégorie fibrée. Signalons auparavant qu'elle
est entièrement analogue aux données de descente pour les faisceaux
quasi-cohérents (Descente, section \ref{descent-section-equivalence})
et aux données de descente pour les schémas au-dessus de schémas
(Descente, section \ref{descent-section-descent-datum}).

\medskip\noindent
Nous adopterons la convention selon laquelle les projections
$\text{pr}_i : X \times \ldots \times X \to X$
sont numérotées à partir de $i = 0$. Ainsi, nous avons
$\text{pr}_0, \text{pr}_1 : X \times X  \to X$,
$\text{pr}_0, \text{pr}_1, \text{pr}_2 : X \times X \times X  \to X$,
etc.

\begin{definition}
\label{stacks-definition-descent-data}
Soit $\mathcal{C}$ une catégorie.
Soit $p : \mathcal{S} \to \mathcal{C}$ une catégorie fibrée.
Faisons un choix d'images inverses comme dans Catégories,
définition \ref{categories-definition-pullback-functor-fibred-category}.
Soit $\mathcal{U} = \{f_i : U_i \to U\}_{i \in I}$
une famille de morphismes de $\mathcal{C}$. Supposons que tous les produits
fibrés $U_i \times_U U_j$ et $U_i \times_U U_j \times_U U_k$ existent.
\begin{enumerate}
\item Une {\it donnée de descente $(X_i, \varphi_{ij})$ dans $\mathcal{S}$
relative à la famille $\{f_i : U_i \to U\}$} consiste en un objet $X_i$
de $\mathcal{S}_{U_i}$ pour chaque $i \in I$, et en un isomorphisme
$\varphi_{ij} : \text{pr}_0^*X_i \to \text{pr}_1^*X_j$
dans $\mathcal{S}_{U_i \times_U U_j}$ pour chaque paire $(i, j) \in I^2$,
tels que, pour tout triplet d'indices $(i, j, k) \in I^3$, le diagramme
$$
\xymatrix{
\text{pr}_0^*X_i \ar[rd]_{\text{pr}_{01}^*\varphi_{ij}}
\ar[rr]_{\text{pr}_{02}^*\varphi_{ik}} & &
\text{pr}_2^*X_k \\
& \text{pr}_1^*X_j \ar[ru]_{\text{pr}_{12}^*\varphi_{jk}} &
}
$$
dans la catégorie $\mathcal{S}_{U_i \times_U U_j \times_U U_k}$
soit commutatif. C'est ce qu'on appelle la {\it condition de cocycle}.
\item Un {\it morphisme $\psi : (X_i, \varphi_{ij}) \to
(X'_i, \varphi'_{ij})$ de données de descente} consiste en une famille
$\psi = (\psi_i)_{i\in I}$ de morphismes
$\psi_i : X_i \to X'_i$ dans $\mathcal{S}_{U_i}$
telle que tous les diagrammes
$$
\xymatrix{
\text{pr}_0^*X_i \ar[r]_{\varphi_{ij}} \ar[d]_{\text{pr}_0^*\psi_i}
& \text{pr}_1^*X_j \ar[d]^{\text{pr}_1^*\psi_j} \\
\text{pr}_0^*X'_i \ar[r]^{\varphi'_{ij}} &
\text{pr}_1^*X'_j \\
}
$$
dans les catégories $\mathcal{S}_{U_i \times_U U_j}$ soient commutatifs.
\item La catégorie des données de descente relatives à
$\mathcal{U}$ est notée $DD(\mathcal{U})$.
\end{enumerate}
\end{definition}

\noindent
Les produits fibrés $U_i \times_U U_j$ et
$U_i \times_U U_j \times_U U_k$ existent si chacun des morphismes
$f_i : U_i \to U$ est {\it représentable}, voir
Catégories, définition \ref{categories-definition-representable-morphism}.
Rappelons que, dans un site, l'une des conditions imposées à un recouvrement
$\{U_i \to U\}$ est que chacun des morphismes soit représentable, voir
Sites, définition \ref{sites-definition-site}, partie (3).
En fait, la définition ci-dessus présente surtout un intérêt lorsque
$\mathcal{C}$ est un site et que $\{U_i \to U\}$ est un recouvrement de
$\mathcal{C}$. Une donnée de descente n'est toutefois qu'une construction
abstraite que l'on peut définir comme ci-dessus. C'est utile : par exemple,
étant donnée une catégorie fibrée au-dessus de $\mathcal{C}$, on peut considérer
l'ensemble des familles relativement auxquelles les données de descente sont
effectives et tenter de les employer comme famille de recouvrements d'un site.

\begin{remarks}
\label{stacks-remarks-definition-descent-datum}
Deux remarques au sujet de la définition \ref{stacks-definition-descent-data}
s'imposent. Soit $p : \mathcal{S} \to \mathcal{C}$ une catégorie fibrée.
Soient $\{f_i : U_i \to U\}_{i \in I}$ et $(X_i, \varphi_{ij})$
comme dans la définition \ref{stacks-definition-descent-data}.
\begin{enumerate}
\item Il existe un morphisme diagonal $\Delta : U_i \to U_i \times_U U_i$.
Nous pouvons tirer $\varphi_{ii}$ en arrière par ce morphisme et obtenir un
automorphisme $\Delta^\ast \varphi_{ii} \in \text{Aut}_{U_i}(X_i)$.
En tirant la condition de cocycle pour le triplet $(i, i, i)$ en arrière par
$\Delta_{123} : U_i \to U_i \times_U U_i \times_U U_i$, nous déduisons que
$\Delta^\ast \varphi_{ii} \circ \Delta^\ast \varphi_{ii} =
\Delta^\ast \varphi_{ii}$ ; ainsi $\Delta^\ast \varphi_{ii} =
\text{id}_{X_i}$.
\item Il existe un morphisme
$\Delta_{13}: U_i \times_U U_j \to U_i \times_U U_j \times_U U_i$,
et nous pouvons tirer en arrière la condition de cocycle pour le triplet
$(i, j, i)$ afin d'obtenir l'identité
$(\sigma^\ast \varphi_{ji}) \circ \varphi_{ij} =
\text{id}_{\text{pr}_0^\ast X_i}$, où
$\sigma : U_i \times_U U_j \to U_j \times_U U_i$ est le morphisme de permutation.
\end{enumerate}
\end{remarks}

\begin{lemma}
\label{stacks-lemma-pullback}
(Changement de base des données de descente.)
Soit $\mathcal{C}$ une catégorie.
Soit $p : \mathcal{S} \to \mathcal{C}$ une catégorie fibrée.
Faisons un choix d'images inverses comme dans Catégories,
définition \ref{categories-definition-pullback-functor-fibred-category}.
Soient $\mathcal{U} = \{f_i : U_i \to U\}_{i \in I}$ et
$\mathcal{V} = \{V_j \to V\}_{j \in J}$
des familles de morphismes de $\mathcal{C}$ ayant chacune un but fixé.
Supposons que tous les produits fibrés
$U_i \times_U U_{i'}$, $U_i \times_U U_{i'} \times_U U_{i''}$,
$V_j \times_V V_{j'}$ et $V_j \times_V V_{j'} \times_V V_{j''}$ existent.
Les données $\alpha : I \to J$, $h : U \to V$ et
$g_i : U_i \to V_{\alpha(i)}$ définissent un morphisme de familles
d'applications à but fixé, voir
Sites, définition \ref{sites-definition-morphism-coverings}.
\begin{enumerate}
\item Soit $(Y_j, \varphi_{jj'})$ une donnée de descente relative à la
famille $\{V_j \to V\}$. Le système
$$
\left(
g_i^*Y_{\alpha(i)},
(g_i \times g_{i'})^*\varphi_{\alpha(i)\alpha(i')}
\right)
$$
est une donnée de descente relative à $\mathcal{U}$.
\item Cette construction définit un foncteur des données de descente relatives
à $\mathcal{V}$ vers les données de descente relatives à $\mathcal{U}$.
\item Étant donné un second morphisme de familles d'applications à but fixé
formé de $\alpha' : I \to J$, $h' : U \to V$ et
$g'_i : U_i \to V_{\alpha'(i)}$, si $h = h'$, alors les deux foncteurs
obtenus entre catégories de données de descente sont canoniquement isomorphes.
\end{enumerate}
\end{lemma}

\begin{proof}
Omis.
\end{proof}

\begin{definition}
\label{stacks-definition-pullback-functor}
Avec $\mathcal{U} = \{U_i \to U\}_{i \in I}$,
$\mathcal{V} = \{V_j \to V\}_{j \in J}$,
$\alpha : I \to J$, $h : U \to V$ et
$g_i : U_i \to V_{\alpha(i)}$ comme dans le lemme \ref{stacks-lemma-pullback},
le foncteur
$$
(Y_j, \varphi_{jj'}) \longmapsto
(g_i^*Y_{\alpha(i)}, (g_i \times g_{i'})^*\varphi_{\alpha(i)\alpha(i')})
$$
construit dans ce lemme est appelé le {\it foncteur de changement de base}
sur les données de descente.
\end{definition}

\noindent
Étant donné $h : U \to V$, s'il existe un morphisme
$\tilde h : \mathcal{U} \to \mathcal{V}$ au-dessus de $h$,
alors $\tilde h^*$ ne dépend pas du choix de $\tilde h$, comme nous l'avons vu
dans le lemme \ref{stacks-lemma-pullback}. Nous écrirons donc parfois simplement
$h^*$ pour désigner le foncteur de changement de base.

\begin{definition}
\label{stacks-definition-effective-descent-datum}
Soit $\mathcal{C}$ une catégorie.
Soit $p : \mathcal{S} \to \mathcal{C}$ une catégorie fibrée.
Faisons un choix d'images inverses comme dans Catégories,
définition \ref{categories-definition-pullback-functor-fibred-category}.
Soit $\mathcal{U} = \{f_i : U_i \to U\}_{i \in I}$ une famille de morphismes
de but $U$. Supposons que tous les produits fibrés
$U_i \times_U U_j$ et $U_i \times_U U_j \times_U U_k$ existent.
\begin{enumerate}
\item Étant donné un objet $X$ de $\mathcal{S}_U$, la {\it donnée de descente
triviale} est la donnée de descente $(X, \text{id}_X)$ relative à la famille
$\{\text{id}_U : U \to U\}$.
\item Étant donné un objet $X$ de $\mathcal{S}_U$, nous obtenons une
{\it donnée de descente canonique} sur la famille des objets $f_i^*X$ en
tirant en arrière la donnée de descente triviale $(X, \text{id}_X)$ par
l'application évidente
$\{f_i : U_i \to U\} \to \{\text{id}_U : U \to U\}$.
Nous notons cette donnée de descente $(f_i^*X, can)$.
\item Une donnée de descente $(X_i, \varphi_{ij})$
relative à $\{f_i : U_i \to U\}$ est dite {\it effective}
s'il existe un objet $X$ de $\mathcal{S}_U$ tel que
$(X_i, \varphi_{ij})$ soit isomorphe à $(f_i^*X, can)$.
\end{enumerate}
\end{definition}

\noindent
Remarquons que la règle qui associe à $X \in \mathcal{S}_U$ sa donnée de
descente canonique relative à $\mathcal{U}$ définit un foncteur
$$
\mathcal{S}_U \longrightarrow DD(\mathcal{U}).
$$
Une donnée de descente est effective si et seulement si elle appartient à
l'image essentielle de ce foncteur. Explicitons maintenant la donnée de
descente canonique.

\begin{lemma}
\label{stacks-lemma-trivial-cocycle}
Dans la situation de la partie (2) de la
définition \ref{stacks-definition-effective-descent-datum}, les applications
$can_{ij} : \text{pr}_0^*f_i^*X \to \text{pr}_1^*f_j^*X$ sont égales à
$(\alpha_{\text{pr}_1, f_j})_X \circ (\alpha_{\text{pr}_0, f_i})_X^{-1}$,
où $\alpha_{\cdot, \cdot}$ est comme dans le
lemme \ref{categories-lemma-fibred} de Catégories et où nous utilisons
l'égalité $f_i \circ \text{pr}_0 = f_j \circ \text{pr}_1$
des applications $U_i \times_U U_j \to U$.
\end{lemma}

\begin{proof}
Omis.
\end{proof}

\begin{lemma}
\label{stacks-lemma-compare-descent-condition}
Soit $\mathcal{C}$ une catégorie. Soit
$\mathcal{V} = \{V_j \to U\}_{j \in J} \to
\mathcal{U} = \{U_i \to U\}_{i \in I}$
un morphisme de familles d'applications à but fixé de $\mathcal{C}$, donné par
$\text{id} : U \to U$, $\alpha : J \to I$ et
$f_j : V_j \to U_{\alpha(j)}$. Soit
$p : \mathcal{S} \to \mathcal{C}$ une catégorie fibrée. Si
\begin{enumerate}
\item pour $0 \leq p \leq 3$ et $0 \leq q \leq 3$, avec $p + q \geq 2$,
et pour $i_1, \ldots, i_p \in I$ et $j_1, \ldots, j_q \in J$,
les produits fibrés $U_{i_1} \times_U \ldots \times_U U_{i_p} \times_U
V_{j_1} \times_U \ldots \times_U V_{j_q}$ existent,
\item le foncteur $\mathcal{S}_U \to DD(\mathcal{V})$
est une équivalence,
\item pour tout $i \in I$, le foncteur
$\mathcal{S}_{U_i} \to DD(\mathcal{V}_i)$
est pleinement fidèle, et
\item pour tous $i, i' \in I$, le foncteur
$\mathcal{S}_{U_i \times_U U_{i'}} \to DD(\mathcal{V}_{ii'})$
est fidèle.
\end{enumerate}
Ici $\mathcal{V}_i = \{U_i \times_U V_j \to U_i\}_{j \in J}$ et
$\mathcal{V}_{ii'} =
\{U_i \times_U U_{i'} \times_U V_j \to U_i \times_U U_{i'}\}_{j \in J}$.
Alors $\mathcal{S}_U \to DD(\mathcal{U})$ est une équivalence.
\end{lemma}

\begin{proof}
La condition (1) garantit l'existence d'un nombre suffisant de produits fibrés
pour que l'énoncé ait un sens. Nous allons montrer que le foncteur
$\mathcal{S}_U \to DD(\mathcal{U})$ est essentiellement surjectif.
Supposons donnée une donnée de descente $(X_i, \varphi_{ii'})$ relative à
$\mathcal{U}$. D'après le lemme \ref{stacks-lemma-pullback}, nous pouvons la tirer en
arrière en une donnée de descente $(X_j, \varphi_{jj'})$ pour $\mathcal{V}$.
D'après l'hypothèse (2), cette donnée de descente est effective ; nous obtenons
donc un objet $X$ de $\mathcal{S}_U$ tel que l'image inverse de la donnée de
descente triviale $(X, \text{id}_X)$ par le morphisme
$\mathcal{V} \to \{U \to U\}$ soit isomorphe à
$(X_j, \varphi_{jj'})$. Observons ensuite que nous avons un diagramme
$$
\xymatrix{
\mathcal{V}_i \ar[r] \ar[d] &
\mathcal{V} \ar[r] &
\mathcal{U} \ar[d] \\
\{U_i \to U_i\} \ar[rr] \ar[rru] & &
\{U \to U\}
}
$$
de morphismes de familles d'applications à but fixé de $\mathcal{C}$.
Ce diagramme n'est pas commutatif, mais, d'après le
lemme \ref{stacks-lemma-pullback}, les foncteurs de changement de base ainsi obtenus
sur les données de descente sont canoniquement isomorphes. Par conséquent,
$(X, \text{id}_X)$ et $(X_i, \text{id}_{X_i})$ ont des images inverses
isomorphes dans $DD(\mathcal{V}_i)$. L'hypothèse (3) fournit donc un
isomorphisme $(U_i \to U)^*X \to X_i$ dans la catégorie
$\mathcal{S}_{U_i}$. Nous omettons de vérifier que ces flèches sont compatibles
avec les morphismes $\varphi_{ii'}$ ; indication : utiliser la fidélité des
foncteurs de la condition (4). Nous omettons également de vérifier que le
foncteur $\mathcal{S}_U \to DD(\mathcal{U})$ est pleinement fidèle.
\end{proof}











\section{Champs}
\label{stacks-section-definition}

\noindent
Voici la définition d'un champ. Elle combine la notion de catégorie
fibrée et celle de descente.

\begin{definition}
\label{stacks-definition-stack}
Soit $\mathcal{C}$ un site. Un {\it champ} sur $\mathcal{C}$
est une catégorie $p : \mathcal{S} \to \mathcal{C}$ au-dessus de $\mathcal{C}$ qui
satisfait aux conditions suivantes :
\begin{enumerate}
\item $p : \mathcal{S} \to \mathcal{C}$ est une catégorie fibrée, voir
Catégories, définition \ref{categories-definition-fibred-category},
\item pour tout $U \in \Ob(\mathcal{C})$ et tous $x, y \in \mathcal{S}_U$,
le préfaisceau $\mathit{Mor}(x, y)$ (voir
Définition \ref{stacks-definition-mor-presheaf}) est un faisceau sur
le site $\mathcal{C}/U$, et
\item pour tout recouvrement $\mathcal{U} = \{f_i : U_i \to U\}_{i \in I}$
du site $\mathcal{C}$, toute donnée de descente dans $\mathcal{S}$
relativement à $\mathcal{U}$ est effective.
\end{enumerate}
\end{definition}

\noindent
La formulation ci-dessus nous paraît être la façon la plus commode de concevoir
un champ. En effet, étant donné une catégorie au-dessus de $\mathcal{C}$, pour vérifier
qu'il s'agit d'un champ, on vérifie les propriétés (1), (2) et
(3), dans cet ordre. Bien entendu, les propriétés (2) et (3) n'ont pas de sens
si la catégorie n'est pas fibrée. Sans (2), nous ne pouvons pas prouver que la
descente de (3) est unique à un unique isomorphisme près et fonctorielle.

\medskip\noindent
Le lemme suivant fournit une autre définition.

\begin{lemma}
\label{stacks-lemma-stack-equivalences}
Soit $\mathcal{C}$ un site.
Soit $p : \mathcal{S} \to \mathcal{C}$ une catégorie fibrée
au-dessus de $\mathcal{C}$. Les assertions suivantes sont équivalentes :
\begin{enumerate}
\item $\mathcal{S}$ est un champ sur $\mathcal{C}$, et
\item pour tout recouvrement $\mathcal{U} = \{f_i : U_i \to U\}_{i \in I}$
du site $\mathcal{C}$, le foncteur
$$
\mathcal{S}_U \longrightarrow DD(\mathcal{U})
$$
qui associe à un objet sa donnée de descente canonique est une équivalence.
\end{enumerate}
\end{lemma}

\begin{proof}
Démonstration omise.
\end{proof}

\begin{lemma}
\label{stacks-lemma-substack}
Soit $p : \mathcal{S} \to \mathcal{C}$ un champ sur le site $\mathcal{C}$.
Soit $\mathcal{S}'$ une sous-catégorie de $\mathcal{S}$.
Supposons que
\begin{enumerate}
\item si $\varphi : y \to x$ est un morphisme fortement cartésien
de $\mathcal{S}$ et si
$x$ est un objet de $\mathcal{S}'$, alors $y$ est isomorphe à un
objet de $\mathcal{S}'$,
\item $\mathcal{S}'$ est une sous-catégorie pleine de $\mathcal{S}$, et
\item si $\{f_i : U_i \to U\}$ est un recouvrement de $\mathcal{C}$
et si $x$ est un objet de $\mathcal{S}$ au-dessus de $U$ tel que $f_i^*x$
soit isomorphe à un objet de $\mathcal{S}'$ pour tout $i$,
alors $x$ est isomorphe à un objet de $\mathcal{S}'$.
\end{enumerate}
Alors $\mathcal{S}' \to \mathcal{C}$ est un champ.
\end{lemma}

\begin{proof}
Démonstration omise. Indications :
la première condition garantit que $\mathcal{S}'$ est une catégorie fibrée.
La deuxième garantit que les préfaisceaux $\mathit{Isom}$
de $\mathcal{S}'$ sont des faisceaux (puisqu'ils sont identiques à leurs homologues
dans $\mathcal{S}$). La troisième garantit que la condition de descente
est satisfaite dans $\mathcal{S}'$, car nous pouvons d'abord effectuer la descente dans
$\mathcal{S}$, puis (3) implique que l'objet obtenu est isomorphe à un objet de
$\mathcal{S}'$.
\end{proof}

\begin{lemma}
\label{stacks-lemma-stack-equivalent}
Soit $\mathcal{C}$ un site.
Soient $\mathcal{S}_1$, $\mathcal{S}_2$ des catégories au-dessus de $\mathcal{C}$.
Supposons que $\mathcal{S}_1$ et $\mathcal{S}_2$ soient équivalentes
comme catégories au-dessus de $\mathcal{C}$.
Alors $\mathcal{S}_1$ est un champ sur $\mathcal{C}$ si et seulement si
$\mathcal{S}_2$ est un champ sur $\mathcal{C}$.
\end{lemma}

\begin{proof}
Soient $F : \mathcal{S}_1 \to \mathcal{S}_2$,
$G : \mathcal{S}_2 \to \mathcal{S}_1$ des foncteurs au-dessus de $\mathcal{C}$, et soient
$i : F \circ G \to \text{id}_{\mathcal{S}_2}$,
$j : G \circ F \to \text{id}_{\mathcal{S}_1}$ des isomorphismes de
foncteurs au-dessus de $\mathcal{C}$. D'après
Catégories, lemme \ref{categories-lemma-fibred-equivalent},
nous voyons que $\mathcal{S}_1$ est fibrée si et seulement si $\mathcal{S}_2$
est fibrée au-dessus de $\mathcal{C}$. Nous pouvons donc supposer que
$\mathcal{S}_1$ et $\mathcal{S}_2$ sont toutes deux fibrées. De plus, la démonstration de
Catégories, lemme \ref{categories-lemma-fibred-equivalent},
montre que $F$ et $G$ envoient les morphismes fortement cartésiens sur des morphismes
fortement cartésiens, c'est-à-dire que $F$ et $G$ sont des $1$-morphismes de catégories
fibrées au-dessus de $\mathcal{C}$. Cela signifie que, étant donnés
$U \in \Ob(\mathcal{C})$ et $x, y \in \mathcal{S}_{1, U}$, les préfaisceaux
$$
\mathit{Mor}_{\mathcal{S}_1}(x, y),
\mathit{Mor}_{\mathcal{S}_1}(F(x), F(y)) :
(\mathcal{C}/U)^{opp} \longrightarrow \textit{Ensembles}.
$$
s'identifient, voir
lemme \ref{stacks-lemma-presheaf-mor-map-fibred-categories}. Par conséquent,
le premier est un faisceau si et seulement si le second est un faisceau.
Enfin, nous devons montrer que, si toute donnée de descente dans $\mathcal{S}_1$
est effective, il en va de même de toute donnée de descente dans $\mathcal{S}_2$.
Pour cela, soit $(X_i, \varphi_{ii'})$ une donnée de descente
dans $\mathcal{S}_2$ relativement au recouvrement $\{U_i \to U\}$ du site
$\mathcal{C}$. Alors $(G(X_i), G(\varphi_{ii'}))$ est une donnée de descente
dans $\mathcal{S}_1$ relativement au recouvrement $\{U_i \to U\}$.
Soit $X$ un objet de $\mathcal{S}_{1, U}$ tel que la
donnée de descente $(f_i^*X, can)$ soit isomorphe à
$(G(X_i), G(\varphi_{ii'}))$. Alors $F(X)$ est un objet de $\mathcal{S}_{2, U}$
tel que la donnée de descente $(f_i^*F(X), can)$ soit isomorphe à
$(F(G(X_i)), F(G(\varphi_{ii'})))$, laquelle est à son tour isomorphe à
la donnée de descente initiale $(X_i, \varphi_{ii'})$ au moyen de $i$.
\end{proof}

\noindent
La $2$-catégorie des champs sur $\mathcal{C}$
est définie comme suit.

\begin{definition}
\label{stacks-definition-stacks-over-C}
Soit $\mathcal{C}$ un site.
La {\it $2$-catégorie des champs sur $\mathcal{C}$}
est la sous-$2$-catégorie de la $2$-catégorie des catégories fibrées
au-dessus de $\mathcal{C}$ (voir
Catégories, définition \ref{categories-definition-fibred-categories-over-C})
définie comme suit :
\begin{enumerate}
\item Ses objets sont les champs $p : \mathcal{S} \to \mathcal{C}$.
\item Ses $1$-morphismes $(\mathcal{S}, p) \to (\mathcal{S}', p')$
sont les foncteurs $G : \mathcal{S} \to \mathcal{S}'$ tels que
$p' \circ G = p$ et tels que $G$ envoie les morphismes fortement cartésiens
sur des morphismes fortement cartésiens.
\item Ses $2$-morphismes $t : G \to H$, pour
$G, H : (\mathcal{S}, p) \to (\mathcal{S}', p')$,
sont les morphismes de foncteurs
tels que $p'(t_x) = \text{id}_{p(x)}$
pour tout $x \in \Ob(\mathcal{S})$.
\end{enumerate}
\end{definition}

\begin{lemma}
\label{stacks-lemma-2-product-stacks}
Soit $\mathcal{C}$ un site.
La $(2, 1)$-catégorie des champs sur $\mathcal{C}$
admet des 2-produits fibrés, et ceux-ci sont décrits comme dans
Catégories, lemme \ref{categories-lemma-2-product-categories-over-C}.
\end{lemma}

\begin{proof}
Soient $f : \mathcal{X} \to \mathcal{S}$ et
$g : \mathcal{Y} \to \mathcal{S}$ des
$1$-morphismes de champs sur $\mathcal{C}$
au sens de la définition ci-dessus. La catégorie
$\mathcal{X} \times_\mathcal{S} \mathcal{Y}$
décrite dans
Catégories, lemme \ref{categories-lemma-2-product-categories-over-C}, est une
catégorie fibrée d'après
Catégories, lemme \ref{categories-lemma-2-product-fibred-categories-over-C}.
(C'est ici que nous utilisons le fait que $f$ et $g$ préservent les morphismes
fortement cartésiens.) Il reste à montrer que les préfaisceaux de morphismes sont des faisceaux
et que la descente relativement aux recouvrements de $\mathcal{C}$ est effective.

\medskip\noindent
Rappelons qu'un objet de $\mathcal{X} \times_\mathcal{S} \mathcal{Y}$
est donné par un quadruplet $(U, x, y, \phi)$.
Il se trouve au-dessus de l'objet
$U$ de $\mathcal{C}$. Soit ensuite $(U, x', y', \phi')$ un second
objet au-dessus de $U$.
Rappelons que $\phi : f(x) \to g(y)$ et $\phi' : f(x') \to g(y')$
sont des isomorphismes dans la catégorie $\mathcal{S}_U$. Utilisons ces
isomorphismes pour identifier $z = f(x) = g(y)$ et
$z' = f(x') = g(y')$. Avec ces identifications,
il est clair que
$$
\mathit{Mor}((U, x, y, \phi), (U, x', y', \phi'))
=
\mathit{Mor}(x, x')
\times_{\mathit{Mor}(z, z')}
\mathit{Mor}(y, y')
$$
comme préfaisceaux. Or, puisque le produit fibré dans la catégorie des
préfaisceaux préserve les faisceaux (Sites, lemme \ref{sites-lemma-limit-sheaf}),
nous voyons que c'est un faisceau.

\medskip\noindent
Soit $\mathcal{U} = \{f_i : U_i \to U\}_{i \in I}$ un recouvrement du site
$\mathcal{C}$. Soit $(X_i, \chi_{ij})$ une donnée de descente
dans $\mathcal{X} \times_\mathcal{S} \mathcal{Y}$ relativement à $\mathcal{U}$.
Écrivons $X_i = (U_i, x_i, y_i, \phi_i)$ comme ci-dessus. Écrivons
$\chi_{ij} = (\varphi_{ij}, \psi_{ij})$ comme dans la définition de
la catégorie $\mathcal{X} \times_\mathcal{S} \mathcal{Y}$ (voir
Catégories, lemme \ref{categories-lemma-2-product-categories-over-C}).
Il est clair que $(x_i, \varphi_{ij})$ est une donnée de descente dans
$\mathcal{X}$ et que $(y_i, \psi_{ij})$ est une donnée de descente dans
$\mathcal{Y}$. Puisque $\mathcal{X}$ et $\mathcal{Y}$ sont des champs, ces
données de descente sont effectives. Nous obtenons donc
$x \in \Ob(\mathcal{X}_U)$ et $y \in \Ob(\mathcal{Y}_U)$
avec $x_i = x|_{U_i}$ et $y_i = y|_{U_i}$, de façon compatible avec les données de descente.
Posons $z = f(x)$ et $z' = g(y)$, qui sont tous deux des objets de $\mathcal{S}_U$.
Les morphismes $\phi_i$ sont des éléments de
$\mathit{Isom}(z, z')(U_i)$ vérifiant
$\phi_i|_{U_i \times_U U_j} = \phi_j|_{U_i \times_U U_j}$.
Par la propriété de faisceau de $\mathit{Isom}(z, z')$,
nous obtenons donc un isomorphisme $\phi : z = f(x) \to z' = g(y)$.
Nous omettons la vérification que la donnée de descente canonique associée à
l'objet $(U, x, y, \phi)$ de
$(\mathcal{X} \times_\mathcal{S} \mathcal{Y})_U$ est isomorphe
à la donnée de descente dont nous sommes partis.
\end{proof}

\begin{lemma}
\label{stacks-lemma-characterize-ff}
Soit $\mathcal{C}$ un site.
Soient $\mathcal{S}_1$, $\mathcal{S}_2$ des champs sur $\mathcal{C}$.
Soit $F : \mathcal{S}_1 \to \mathcal{S}_2$ un $1$-morphisme.
Alors les assertions suivantes sont équivalentes :
\begin{enumerate}
\item $F$ est pleinement fidèle,
\item pour tout $U \in \Ob(\mathcal{C})$ et tous
$x, y \in \Ob(\mathcal{S}_{1, U})$, l'application
$$
F :
\mathit{Mor}_{\mathcal{S}_1}(x, y)
\longrightarrow
\mathit{Mor}_{\mathcal{S}_2}(F(x), F(y))
$$
est un isomorphisme de faisceaux sur $\mathcal{C}/U$.
\end{enumerate}
\end{lemma}

\begin{proof}
Supposons (1). Pour $U, x, y$ comme dans (2), l'application affichée $F$ s'évalue en
$F : \Mor_{\mathcal{S}_{1, V}}(x|_V, y|_V) \to
\Mor_{\mathcal{S}_{2, V}}(F(x|_V), F(y|_V))$
sur un objet $V$ de $\mathcal{C}$ au-dessus de $U$.
Or, puisque $F$ est pleinement fidèle, l'application correspondante
$\Mor_{\mathcal{S}_1}(x|_V, y|_V) \to \Mor_{\mathcal{S}_2}(F(x|_V), F(y|_V))$
est une bijection. Les morphismes de la catégorie fibre $\mathcal{S}_{1, V}$ sont
exactement les morphismes entre $x|_V$ et $y|_V$ dans $\mathcal{S}_1$ qui sont
au-dessus de $\text{id}_V$. De même, les morphismes de la catégorie fibre
$\mathcal{S}_{2, V}$ sont exactement les morphismes entre $F(x|_V)$ et
$F(y|_V)$ dans $\mathcal{S}_2$ qui sont au-dessus de $\text{id}_V$. Nous constatons donc que $F$
induit également une bijection entre ceux-ci. Par conséquent, (2) est satisfaite.

\medskip\noindent
Supposons (2). Soient $U$, $V$ des objets de $\mathcal{C}$, ainsi que
$x \in \Ob(\mathcal{S}_{1, U})$ et
$y \in \Ob(\mathcal{S}_{1, V})$. Pour montrer que $F$ est pleinement fidèle,
il suffit de prouver qu'il induit une bijection sur les
morphismes au-dessus d'un $f : U \to V$ fixé. Choisissons un morphisme fortement cartésien
$f^*y \to y$ dans $\mathcal{S}_1$ au-dessus de $f$. Nous obtenons ainsi une
bijection entre l'ensemble des morphismes $x \to y$ dans $\mathcal{S}_1$ au-dessus
de $f$ et $\Mor_{\mathcal{S}_{1, U}}(x, f^*y)$. Puisque $F$ préserve
les morphismes fortement cartésiens en tant que $1$-morphisme dans la $2$-catégorie
des champs sur $\mathcal{C}$, nous obtenons également une bijection
entre l'ensemble des morphismes $F(x) \to F(y)$ dans $\mathcal{S}_2$ au-dessus
de $f$ et $\Mor_{\mathcal{S}_{2, U}}(F(x), F(f^*y))$.
Comme $F$ induit une bijection
$\Mor_{\mathcal{S}_{1, U}}(x, f^*y) \to
\Mor_{\mathcal{S}_{2, U}}(F(x), F(f^*y))$,
nous concluons que (1) est satisfaite.
\end{proof}

\begin{lemma}
\label{stacks-lemma-characterize-essentially-surjective-when-ff}
Soit $\mathcal{C}$ un site.
Soient $\mathcal{S}_1$, $\mathcal{S}_2$ des champs sur $\mathcal{C}$.
Soit $F : \mathcal{S}_1 \to \mathcal{S}_2$ un $1$-morphisme
pleinement fidèle. Alors les assertions suivantes sont équivalentes :
\begin{enumerate}
\item $F$ est une équivalence,
\item pour tout $U \in \Ob(\mathcal{C})$ et tout
$x \in \Ob(\mathcal{S}_{2, U})$, il existe un recouvrement
$\{f_i : U_i \to U\}$ tel que $f_i^*x$ appartienne à l'image essentielle
du foncteur $F : \mathcal{S}_{1, U_i} \to \mathcal{S}_{2, U_i}$.
\end{enumerate}
\end{lemma}

\begin{proof}
L'implication (1) $\Rightarrow$ (2) est immédiate.
Pour voir que (2) implique (1), nous devons montrer que tout
$x$ comme dans (2) appartient à l'image essentielle du foncteur $F$.
Pour cela, choisissons un recouvrement comme dans (2),
$x_i \in \Ob(\mathcal{S}_{1, U_i})$, et des
isomorphismes $\varphi_i : F(x_i) \to f_i^*x$. Nous obtenons alors une donnée de
descente pour $\mathcal{S}_1$ relativement à $\{f_i : U_i \to U\}$
en prenant pour
$$
\varphi_{ij} :
x_i|_{U_i \times_U U_j}
\longrightarrow
x_j|_{U_i \times_U U_j}
$$
la flèche telle que $F(\varphi_{ij}) = \varphi_j^{-1} \circ \varphi_i$.
Cette donnée de descente est effective par les axiomes d'un champ ; nous
obtenons donc un objet $x_1$ de $\mathcal{S}_1$ au-dessus de $U$. Nous omettons la
vérification que $F(x_1)$ est isomorphe à $x$ au-dessus de $U$.
\end{proof}

\begin{remark}
\label{stacks-remark-stack-make-small}
(Réduction d'un « grand » champ pour obtenir un champ.)
Soit $\mathcal{C}$ un site. Supposons que $p : \mathcal{S} \to \mathcal{C}$
soit un foncteur d'une « grande » catégorie vers $\mathcal{C}$, c'est-à-dire supposons
que la collection des objets de $\mathcal{S}$ forme une classe propre.
Supposons enfin que $p : \mathcal{S} \to \mathcal{C}$ satisfasse
aux conditions (1), (2), (3) de la
Définition \ref{stacks-definition-stack}.
En général, il n'existe aucun moyen de remplacer $p : \mathcal{S} \to \mathcal{C}$
par une catégorie équivalente de façon à obtenir un champ. La raison en est qu'il
peut arriver qu'une catégorie fibre $\mathcal{S}_U$ ait une classe propre
de classes d'isomorphisme d'objets.
Supposons en revanche que
\begin{enumerate}
\item[(4)] pour tout $U \in \Ob(\mathcal{C})$, il existe un ensemble
$S_U \subset \Ob(\mathcal{S}_U)$ tel que tout objet de
$\mathcal{S}_U$ soit isomorphe dans $\mathcal{S}_U$ à un élément de $S_U$.
\end{enumerate}
Dans ce cas, nous pouvons trouver une sous-catégorie pleine $\mathcal{S}_{small}$
de $\mathcal{S}$ telle qu'en posant $p_{small} = p|_{\mathcal{S}_{small}}$,
nous ayons
\begin{enumerate}
\item[(a)] le foncteur $p_{small} : \mathcal{S}_{small} \to \mathcal{C}$
définit un champ, et
\item[(b)] l'inclusion $\mathcal{S}_{small} \to \mathcal{S}$
est pleinement fidèle et essentiellement surjective.
\end{enumerate}
(Indication : pour tout $U \in \Ob(\mathcal{C})$,
notons $\alpha(U)$ le plus petit ordinal tel que
$\Ob(\mathcal{S}_U) \cap V_{\alpha(U)}$ s'envoie surjectivement sur l'ensemble
des classes d'isomorphisme de $\mathcal{S}_U$, et posons
$\alpha = \sup_{U \in \Ob(\mathcal{C})} \alpha(U)$.
Prenons alors
$\Ob(\mathcal{S}_{small}) = \Ob(\mathcal{S}) \cap V_\alpha$.
Pour les notations utilisées, voir Ensembles, section \ref{sets-section-sets-hierarchy}.)
\end{remark}







\section{Champs en groupoïdes}
\label{stacks-section-stacks-in-groupoids}

\noindent
Parmi les champs, ceux qui sont fibrés en groupoïdes sont un peu plus faciles
à appréhender. Nous les redéfinissons comme suit.

\begin{definition}
\label{stacks-definition-stack-in-groupoids}
Un {\it champ en groupoïdes} sur un site $\mathcal{C}$ est une
catégorie $p : \mathcal{S} \to \mathcal{C}$ au-dessus de $\mathcal{C}$
telle que
\begin{enumerate}
\item $p : \mathcal{S} \to \mathcal{C}$ est fibrée
en groupoïdes au-dessus de $\mathcal{C}$ (voir
Catégories, définition \ref{categories-definition-fibred-groupoids}),
\item pour tout $U \in \Ob(\mathcal{C})$ et
tous $x, y\in \Ob(\mathcal{S}_U)$, le préfaisceau
$\mathit{Isom}(x, y)$ est un faisceau sur le site $\mathcal{C}/U$, et
\item pour tout recouvrement $\mathcal{U} = \{U_i \to U\}$ dans $\mathcal{C}$,
toute donnée de descente $(x_i, \phi_{ij})$ relative à $\mathcal{U}$ est effective.
\end{enumerate}
\end{definition}

\noindent
La troisième condition est généralement la plus difficile à vérifier.
Voici le lemme qui compare cette notion à celle de champ.

\begin{lemma}
\label{stacks-lemma-stack-in-groupoids-stack}
Soit $\mathcal{C}$ un site.
Soit $p : \mathcal{S} \to \mathcal{C}$ une catégorie au-dessus de $\mathcal{C}$.
Les assertions suivantes sont équivalentes :
\begin{enumerate}
\item $\mathcal{S}$ est un champ en groupoïdes sur $\mathcal{C}$,
\item $\mathcal{S}$ est un champ sur $\mathcal{C}$ et toutes ses
catégories fibres sont des groupoïdes, et
\item $\mathcal{S}$ est fibrée en groupoïdes au-dessus de $\mathcal{C}$
et est un champ sur $\mathcal{C}$.
\end{enumerate}
\end{lemma}

\begin{proof}
Démonstration omise, mais voir Catégories, lemme \ref{categories-lemma-fibred-groupoids}.
\end{proof}

\begin{lemma}
\label{stacks-lemma-stack-gives-stack-groupoids}
Soit $\mathcal{C}$ un site.
Soit $p : \mathcal{S} \to \mathcal{C}$ un champ.
Soit $p' : \mathcal{S}' \to \mathcal{C}$
la catégorie fibrée en groupoïdes associée à $\mathcal{S}$,
construite dans
Catégories, lemme \ref{categories-lemma-fibred-gives-fibred-groupoids}.
Alors $p' : \mathcal{S}' \to \mathcal{C}$ est un champ en groupoïdes.
\end{lemma}

\begin{proof}
Rappelons que les morphismes de $\mathcal{S}'$ sont exactement les
morphismes fortement cartésiens de $\mathcal{S}$ et que tout isomorphisme de
$\mathcal{S}$ est un tel morphisme. Ainsi, les données de descente dans $\mathcal{S}'$
sont exactement les mêmes que les données de descente dans $\mathcal{S}$. Appliquons maintenant le
lemme \ref{stacks-lemma-stack-equivalences}. Certains détails sont omis.
\end{proof}

\begin{lemma}
\label{stacks-lemma-stack-in-groupoids-equivalent}
Soit $\mathcal{C}$ un site.
Soient $\mathcal{S}_1$, $\mathcal{S}_2$ des catégories au-dessus de $\mathcal{C}$.
Supposons que $\mathcal{S}_1$ et $\mathcal{S}_2$ soient équivalentes
comme catégories au-dessus de $\mathcal{C}$.
Alors $\mathcal{S}_1$ est un champ en groupoïdes sur $\mathcal{C}$ si et seulement si
$\mathcal{S}_2$ est un champ en groupoïdes sur $\mathcal{C}$.
\end{lemma}

\begin{proof}
Cela résulte de la combinaison des
lemmes \ref{stacks-lemma-stack-in-groupoids-stack} et \ref{stacks-lemma-stack-equivalent}.
\end{proof}

\noindent
La $2$-catégorie des champs en groupoïdes sur $\mathcal{C}$
est définie comme suit.

\begin{definition}
\label{stacks-definition-stacks-in-groupoids-over-C}
Soit $\mathcal{C}$ un site.
La {\it $2$-catégorie des champs en groupoïdes sur $\mathcal{C}$}
est la sous-$2$-catégorie de la $2$-catégorie des champs
sur $\mathcal{C}$ (voir Définition \ref{stacks-definition-stacks-over-C})
définie comme suit :
\begin{enumerate}
\item Ses objets sont les champs en groupoïdes
$p : \mathcal{S} \to \mathcal{C}$.
\item Ses $1$-morphismes $(\mathcal{S}, p) \to (\mathcal{S}', p')$
sont les foncteurs $G : \mathcal{S} \to \mathcal{S}'$ tels que
$p' \circ G = p$. (Puisque tout morphisme est fortement cartésien,
tout foncteur les préserve.)
\item Ses $2$-morphismes $t : G \to H$, pour
$G, H : (\mathcal{S}, p) \to (\mathcal{S}', p')$,
sont les morphismes de foncteurs
tels que $p'(t_x) = \text{id}_{p(x)}$
pour tout $x \in \Ob(\mathcal{S})$.
\end{enumerate}
\end{definition}

\noindent
Remarquons que tout $2$-morphisme est automatiquement un isomorphisme, de sorte
que la $2$-catégorie des champs en groupoïdes sur $\mathcal{C}$
est en fait une $(2, 1)$-catégorie (stricte).

\begin{lemma}
\label{stacks-lemma-2-product-stacks-in-groupoids}
Soit $\mathcal{C}$ une catégorie.
La $2$-catégorie des champs en groupoïdes sur $\mathcal{C}$
admet des 2-produits fibrés, et ceux-ci sont décrits comme dans
Catégories, lemme \ref{categories-lemma-2-product-categories-over-C}.
\end{lemma}

\begin{proof}
Cela résulte immédiatement de
Catégories, lemme \ref{categories-lemma-2-product-fibred-categories},
et des lemmes \ref{stacks-lemma-stack-in-groupoids-stack}
et \ref{stacks-lemma-2-product-stacks}.
\end{proof}




\section{Champs en setoïdes}
\label{stacks-section-stacks-in-setoids}

\noindent
Cette brève section se contente d'indiquer qu'un champ en ensembles
est la même chose qu'un faisceau d'ensembles. Pour les notations, consulter
Catégories, section \ref{categories-section-fibred-in-setoids}.

\begin{definition}
\label{stacks-definition-stack-in-sets}
Soit $\mathcal{C}$ un site.
\begin{enumerate}
\item Un {\it champ en setoïdes} sur $\mathcal{C}$
est un champ sur $\mathcal{C}$ dont toutes les catégories fibres sont des
setoïdes.
\item Un {\it champ en ensembles}, ou un {\it champ en catégories discrètes},
est un champ sur $\mathcal{C}$ dont toutes les catégories fibres sont discrètes.
\end{enumerate}
\end{definition}

\noindent
D'après la discussion de la
section \ref{stacks-section-stacks-in-groupoids},
c'est la même chose qu'un champ en groupoïdes dont les catégories fibres
sont des setoïdes (resp.\ discrètes). De plus, c'est aussi la même chose
qu'une catégorie fibrée en setoïdes (resp.\ en ensembles) qui est un champ.

\begin{lemma}
\label{stacks-lemma-when-stack-in-sets}
Soit $\mathcal{C}$ un site. Sous l'équivalence
$$
\left\{
\begin{matrix}
\text{la catégorie des préfaisceaux}\\
\text{d'ensembles sur }\mathcal{C}
\end{matrix}
\right\}
\leftrightarrow
\left\{
\begin{matrix}
\text{la catégorie des catégories}\\
\text{fibrées en ensembles au-dessus de }\mathcal{C}
\end{matrix}
\right\}
$$
de
Catégories, lemme \ref{categories-lemma-2-category-fibred-sets},
les champs en ensembles correspondent précisément aux faisceaux.
\end{lemma}

\begin{proof}
Démonstration omise. Indication : montrer que l'effectivité de la descente correspond exactement à
la condition de faisceau.
\end{proof}

\begin{lemma}
\label{stacks-lemma-stack-in-setoids-characterize}
Soit $\mathcal{C}$ un site.
Soit $\mathcal{S}$ une catégorie fibrée en setoïdes au-dessus de $\mathcal{C}$.
Alors $\mathcal{S}$ est un champ en setoïdes si et seulement si l'unique
catégorie $\mathcal{S}'$ fibrée en ensembles qui lui est équivalente (voir
Catégories, lemme \ref{categories-lemma-setoid-fibres})
est un champ en ensembles. Autrement dit, si et seulement si le préfaisceau
$$
U \longmapsto \Ob(\mathcal{S}_U)/\!\!\cong
$$
est un faisceau.
\end{lemma}

\begin{proof}
Démonstration omise.
\end{proof}

\begin{lemma}
\label{stacks-lemma-stack-in-setoids-equivalent}
Soit $\mathcal{C}$ un site.
Soient $\mathcal{S}_1$, $\mathcal{S}_2$ des catégories au-dessus de $\mathcal{C}$.
Supposons que $\mathcal{S}_1$ et $\mathcal{S}_2$ soient équivalentes
comme catégories au-dessus de $\mathcal{C}$.
Alors $\mathcal{S}_1$ est un champ en setoïdes sur $\mathcal{C}$ si et seulement si
$\mathcal{S}_2$ est un champ en setoïdes sur $\mathcal{C}$.
\end{lemma}

\begin{proof}
D'après
Catégories, lemme \ref{categories-lemma-setoid-fibres},
nous voyons qu'une catégorie $\mathcal{S}$ au-dessus de $\mathcal{C}$ est fibrée en setoïdes
au-dessus de $\mathcal{C}$ si et seulement si elle est équivalente au-dessus de $\mathcal{C}$
à une catégorie fibrée en ensembles.
Nous voyons donc que $\mathcal{S}_1$ est fibrée en setoïdes au-dessus de $\mathcal{C}$
si et seulement si $\mathcal{S}_2$ est fibrée en setoïdes au-dessus de $\mathcal{C}$.
Le lemme résulte alors du
lemme \ref{stacks-lemma-stack-in-setoids-characterize}.
\end{proof}

\noindent
La $2$-catégorie des champs en setoïdes sur $\mathcal{C}$
est définie comme suit.

\begin{definition}
\label{stacks-definition-stacks-in-setoids-over-C}
Soit $\mathcal{C}$ un site.
La {\it $2$-catégorie des champs en setoïdes sur $\mathcal{C}$}
est la sous-$2$-catégorie de la $2$-catégorie des champs
sur $\mathcal{C}$ (voir Définition \ref{stacks-definition-stacks-over-C})
définie comme suit :
\begin{enumerate}
\item Ses objets sont les champs en setoïdes
$p : \mathcal{S} \to \mathcal{C}$.
\item Ses $1$-morphismes $(\mathcal{S}, p) \to (\mathcal{S}', p')$
sont les foncteurs $G : \mathcal{S} \to \mathcal{S}'$ tels que
$p' \circ G = p$. (Puisque tout morphisme est fortement cartésien,
tout foncteur les préserve.)
\item Ses $2$-morphismes $t : G \to H$, pour
$G, H : (\mathcal{S}, p) \to (\mathcal{S}', p')$,
sont les morphismes de foncteurs
tels que $p'(t_x) = \text{id}_{p(x)}$
pour tout $x \in \Ob(\mathcal{S})$.
\end{enumerate}
\end{definition}

\noindent
Remarquons que tout $2$-morphisme est automatiquement un isomorphisme, de sorte
que la $2$-catégorie des champs en setoïdes sur $\mathcal{C}$
est en fait une $(2, 1)$-catégorie (stricte).

\begin{lemma}
\label{stacks-lemma-2-product-stacks-in-setoids}
Soit $\mathcal{C}$ un site.
La $2$-catégorie des champs en setoïdes sur $\mathcal{C}$
admet des 2-produits fibrés, et ceux-ci sont décrits comme dans
Catégories, lemme \ref{categories-lemma-2-product-categories-over-C}.
\end{lemma}

\begin{proof}
Cela résulte immédiatement de
Catégories, lemmes \ref{categories-lemma-2-product-fibred-categories} et
\ref{categories-lemma-2-product-categories-fibred-setoids},
ainsi que des
lemmes \ref{stacks-lemma-stack-in-groupoids-stack} et
\ref{stacks-lemma-2-product-stacks}.
\end{proof}

\begin{lemma}
\label{stacks-lemma-2-fibre-product-gives-stack-in-setoids}
Soit $\mathcal{C}$ un site.
Soient $\mathcal{S}, \mathcal{T}$ des champs en groupoïdes sur $\mathcal{C}$
et soit $\mathcal{R}$ un champ en setoïdes sur $\mathcal{C}$.
Soient $f : \mathcal{T} \to \mathcal{S}$ et $g : \mathcal{R} \to \mathcal{S}$
des $1$-morphismes. Si $f$ est fidèle, alors le $2$-produit fibré
$$
\mathcal{T} \times_{f, \mathcal{S}, g} \mathcal{R}
$$
est un champ en setoïdes sur $\mathcal{C}$.
\end{lemma}

\begin{proof}
Cela résulte immédiatement de la description explicite du $2$-produit fibré dans
Catégories, lemme \ref{categories-lemma-2-product-categories-over-C}.
\end{proof}

\begin{lemma}
\label{stacks-lemma-2-fibre-product-stacks-in-setoids-over-stack-in-groupoids}
Soit $\mathcal{C}$ un site.
Soit $\mathcal{S}$ un champ en groupoïdes sur $\mathcal{C}$ et
soient $\mathcal{S}_i$, $i = 1, 2$, des champs en setoïdes sur $\mathcal{C}$.
Soient $f_i : \mathcal{S}_i \to \mathcal{S}$ des $1$-morphismes.
Alors le $2$-produit fibré
$$
\mathcal{S}_1 \times_{f_1, \mathcal{S}, f_2} \mathcal{S}_2
$$
est un champ en setoïdes sur $\mathcal{C}$.
\end{lemma}

\begin{proof}
C'est un cas particulier du
lemme \ref{stacks-lemma-2-fibre-product-gives-stack-in-setoids},
car $f_2$ est fidèle.
\end{proof}

\begin{lemma}
\label{stacks-lemma-faithful-descent}
Soit $\mathcal{C}$ un site. Soit
$$
\xymatrix{
\mathcal{T}_2 \ar[r] \ar[d]_{G'} & \mathcal{T}_1 \ar[d]^G \\
\mathcal{S}_2 \ar[r]^F & \mathcal{S}_1
}
$$
un diagramme $2$-cartésien de champs en groupoïdes sur $\mathcal{C}$.
Supposons que
\begin{enumerate}
\item pour tout $U \in \Ob(\mathcal{C})$ et tout
$x \in \Ob((\mathcal{S}_1)_U)$, il existe un recouvrement
$\{U_i \to U\}$ tel que $x|_{U_i}$ appartienne à l'image
essentielle de $F : (\mathcal{S}_2)_{U_i} \to (\mathcal{S}_1)_{U_i}$, et
\item $G'$ soit fidèle,
\end{enumerate}
alors $G$ est fidèle.
\end{lemma}

\begin{proof}
Nous pouvons supposer que $\mathcal{T}_2$ est la catégorie
$\mathcal{S}_2 \times_{\mathcal{S}_1} \mathcal{T}_1$
décrite dans
Catégories, lemme \ref{categories-lemma-2-product-categories-over-C}.
D'après
Catégories, lemme \ref{categories-lemma-equivalence-fibred-categories},
la fidélité de $G, G'$ peut être vérifiée sur les catégories fibres.
Supposons que $y, y'$ soient des objets de $\mathcal{T}_1$ au-dessus de l'objet $U$
de $\mathcal{C}$. Soient $\alpha, \beta : y \to y'$ des morphismes de
$(\mathcal{T}_1)_U$ tels que $G(\alpha) = G(\beta)$. Notre but est de
montrer que $\alpha = \beta$. En considérant plutôt
$\gamma = \alpha^{-1} \circ \beta$, nous voyons que $G(\gamma) = \text{id}_{G(y)}$
et nous devons montrer que $\gamma = \text{id}_y$. Par hypothèse, nous
pouvons trouver un recouvrement $\{U_i \to U\}$ tel que $G(y)|_{U_i}$ appartienne à
l'image essentielle de $F :(\mathcal{S}_2)_{U_i} \to (\mathcal{S}_1)_{U_i}$.
Puisqu'il suffit de montrer que $\gamma|_{U_i} = \text{id}$ pour tout $i$,
nous pouvons donc supposer qu'il existe
$f : F(x) \to G(y)$ pour un objet $x$ de $\mathcal{S}_2$ au-dessus de $U$,
la flèche $f$ étant un morphisme de $(\mathcal{S}_1)_U$. Dans ce cas, nous obtenons
un morphisme
$$
(1, \gamma) : (U, x, y, f) \longrightarrow (U, x, y, f)
$$
dans la catégorie fibre de $\mathcal{S}_2 \times_{\mathcal{S}_1} \mathcal{T}_1$
au-dessus de $U$, dont l'image par $G'$ dans $\mathcal{S}_1$ est $\text{id}_x$.
Comme $G'$ est fidèle, nous concluons que $\gamma = \text{id}_y$, ce qui termine la preuve.
\end{proof}

\begin{lemma}
\label{stacks-lemma-stack-in-setoids-descent}
Soit $\mathcal{C}$ un site.
Soit
$$
\xymatrix{
\mathcal{T}_2 \ar[r] \ar[d] & \mathcal{T}_1 \ar[d]^G \\
\mathcal{S}_2 \ar[r]^F & \mathcal{S}_1
}
$$
un diagramme $2$-cartésien de champs en groupoïdes sur $\mathcal{C}$.
Si
\begin{enumerate}
\item $F : \mathcal{S}_2 \to \mathcal{S}_1$ est pleinement fidèle,
\item pour tout $U \in \Ob(\mathcal{C})$ et tout
$x \in \Ob((\mathcal{S}_1)_U)$, il existe un recouvrement
$\{U_i \to U\}$ tel que $x|_{U_i}$ appartienne à l'image essentielle
de $F : (\mathcal{S}_2)_{U_i} \to (\mathcal{S}_1)_{U_i}$, et
\item $\mathcal{T}_2$ est un champ en setoïdes,
\end{enumerate}
alors $\mathcal{T}_1$ est un champ en setoïdes.
\end{lemma}

\begin{proof}
Nous pouvons supposer que $\mathcal{T}_2$ est la catégorie
$\mathcal{S}_2 \times_{\mathcal{S}_1} \mathcal{T}_1$
décrite dans
Catégories, lemme \ref{categories-lemma-2-product-categories-over-C}.
Choisissons $U \in \Ob(\mathcal{C})$ et
$y \in \Ob((\mathcal{T}_1)_U)$.
Nous devons montrer que le faisceau $\mathit{Aut}(y)$ sur $\mathcal{C}/U$
est trivial. Pour cela, nous pouvons remplacer $U$ par les membres d'un
recouvrement de $U$. Par l'hypothèse (2), nous pouvons donc supposer
qu'il existe un objet $x \in \Ob((\mathcal{S}_2)_U)$
et un isomorphisme $f : F(x) \to G(y)$.
Alors $y' = (U, x, y, f)$ est un objet de $\mathcal{T}_2$ au-dessus de $U$
qui est envoyé sur $y$ par la projection $\mathcal{T}_2 \to \mathcal{T}_1$.
Puisque $F$ est pleinement fidèle d'après (1), l'application
$\mathit{Aut}(y') \to \mathit{Aut}(y)$ est surjective, comme le montre la description
explicite des morphismes de $\mathcal{T}_2$ donnée dans
Catégories, lemme \ref{categories-lemma-2-product-categories-over-C}.
Puisque, par (3), le faisceau $\mathit{Aut}(y')$ est trivial,
nous obtenons le résultat du lemme.
\end{proof}

\begin{lemma}
\label{stacks-lemma-relative-sheaf-over-stack-is-stack}
Soit $\mathcal{C}$ un site. Soit $F : \mathcal{S} \to \mathcal{T}$
un $1$-morphisme de catégories fibrées en groupoïdes au-dessus de $\mathcal{C}$.
Supposons que
\begin{enumerate}
\item $\mathcal{T}$ soit un champ en groupoïdes sur $\mathcal{C}$,
\item pour tout $U \in \Ob(\mathcal{C})$, le foncteur
$\mathcal{S}_U \to \mathcal{T}_U$ entre catégories fibres soit fidèle,
\item pour tout $U$ et tout $y \in \Ob(\mathcal{T}_U)$, le préfaisceau
$$
(h : V \to U)
\longmapsto
\{(x, f) \mid x \in \Ob(\mathcal{S}_V), f : F(x) \to f^*y\text{ au-dessus de }V\}/\cong
$$
soit un faisceau sur $\mathcal{C}/U$.
\end{enumerate}
Alors $\mathcal{S}$ est un champ en groupoïdes sur $\mathcal{C}$.
\end{lemma}

\begin{proof}
Nous devons démontrer la descente des morphismes et la descente des objets.

\medskip\noindent
Descente des morphismes. Soit $\{U_i \to U\}$ un recouvrement de $\mathcal{C}$.
Soient $x, x'$ des objets de $\mathcal{S}$ au-dessus de $U$. Pour tout $i$,
soit $\alpha_i : x|_{U_i} \to x'|_{U_i}$ un morphisme au-dessus de $U_i$
tel que les restrictions de $\alpha_i$ et de $\alpha_j$ coïncident comme morphismes
$x|_{U_i \times_U U_j} \to x'|_{U_i \times_U U_j}$.
Puisque $\mathcal{T}$ est un champ en groupoïdes, il existe un morphisme
$\beta : F(x) \to F(x')$ au-dessus de $U$ dont la restriction à $U_i$ est $F(\alpha_i)$.
Nous pouvons alors considérer $\xi = (x, \beta)$ et $\xi' = (x', \text{id}_{F(x')})$
comme des sections du préfaisceau associé à $y = F(x')$ au-dessus de $U$
dans l'hypothèse (3). D'autre part, les restrictions de
$\xi$ et $\xi'$ à $U_i$ sont $(x|_{U_i}, F(\alpha_i))$
et $(x'|_{U_i}, \text{id}_{F(x'|_{U_i})})$.
Elles sont isomorphes l'une à l'autre par le morphisme $\alpha_i$.
Ainsi, $\xi$ et $\xi'$ sont isomorphes d'après l'hypothèse (3). Cela signifie qu'il existe un
morphisme $\alpha : x \to x'$ au-dessus de $U$ tel que $F(\alpha) = \beta$.
Puisque $F$ est fidèle sur les catégories fibres, nous obtenons
$\alpha|_{U_i} = \alpha_i$.

\medskip\noindent
Descente des objets. Soit $\{U_i \to U\}$ un recouvrement de $\mathcal{C}$.
Soit $(x_i, \varphi_{ij})$ une donnée de descente pour $\mathcal{S}$
relativement au recouvrement donné. Puisque $\mathcal{T}$ est un champ en groupoïdes,
il existe un objet $y$ de $\mathcal{T}_U$ et des isomorphismes
$\beta_i : F(x_i) \to y|_{U_i}$ tels que
$F(\varphi_{ij}) = \beta_j|_{U_i \times_U U_j} \circ
(\beta_i|_{U_i \times_U U_j})^{-1}$.
Alors les $(x_i, \beta_i)$ sont des sections du préfaisceau associé à
$y$ au-dessus de $U$ défini dans l'hypothèse (3).
De plus, $\varphi_{ij}$ définit un isomorphisme
de la paire $(x_i, \beta_i)|_{U_i \times_U U_j}$ vers
la paire $(x_j, \beta_j)|_{U_i \times_U U_j}$.
Ainsi, d'après l'hypothèse (3), il existe une paire $(x, \beta)$ au-dessus de $U$
dont la restriction à $U_i$ est isomorphe à $(x_i, \beta_i)$.
Cela signifie qu'il existe des morphismes $\alpha_i : x_i \to x|_{U_i}$
tels que $\beta_i = \beta|_{U_i} \circ F(\alpha_i)$.
Puisque $F$ est fidèle sur les catégories fibres, un calcul montre
que $\varphi_{ij} = \alpha_j|_{U_i \times_U U_j} \circ
(\alpha_i|_{U_i \times_U U_j})^{-1}$. Cela achève la démonstration.
\end{proof}










\section{Le champ d'inertie}
\label{stacks-section-the-inertia-stack}

\noindent
Soient
$p : \mathcal{S} \to \mathcal{C}$ et
$p' : \mathcal{S}' \to \mathcal{C}$
des catégories fibrées au-dessus de la catégorie $\mathcal{C}$.
Soit $F : \mathcal{S} \to \mathcal{S}'$ un $1$-morphisme de
catégories fibrées au-dessus de $\mathcal{C}$.
Rappelons que nous avons défini dans
Catégories, définition \ref{categories-definition-inertia-fibred-category},
une {\it catégorie fibrée d'inertie relative}
$\mathcal{I}_{\mathcal{S}/\mathcal{S}'} \to \mathcal{C}$ comme la
catégorie dont les objets sont les paires $(x , \alpha)$ où
$x \in \Ob(\mathcal{S})$ et $\alpha : x \to x$ vérifie
$F(\alpha) = \text{id}_{F(x)}$. Il existe aussi une version absolue,
à savoir l'{\it inertie} $\mathcal{I}_\mathcal{S}$ de $\mathcal{S}$.
Ces catégories d'inertie sont en fait des champs sur $\mathcal{C}$ dès que
$\mathcal{S}$ et $\mathcal{S}'$ sont des champs.

\begin{lemma}
\label{stacks-lemma-inertia}
Soit $\mathcal{C}$ un site. Soient
$p : \mathcal{S} \to \mathcal{C}$ et
$p' : \mathcal{S}' \to \mathcal{C}$
des champs sur le site $\mathcal{C}$.
Soit $F : \mathcal{S} \to \mathcal{S}'$ un $1$-morphisme de
champs sur $\mathcal{C}$.
\begin{enumerate}
\item Les inerties $\mathcal{I}_{\mathcal{S}/\mathcal{S}'}$ et
$\mathcal{I}_\mathcal{S}$ sont des champs sur $\mathcal{C}$.
\item Si $\mathcal{S}, \mathcal{S}'$ sont des champs en groupoïdes sur
$\mathcal{C}$, alors $\mathcal{I}_{\mathcal{S}/\mathcal{S}'}$ et
$\mathcal{I}_\mathcal{S}$ le sont également.
\item Si $\mathcal{S}, \mathcal{S}'$ sont des champs en setoïdes sur $\mathcal{C}$,
alors $\mathcal{I}_{\mathcal{S}/\mathcal{S}'}$ et
$\mathcal{I}_\mathcal{S}$ le sont également.
\end{enumerate}
\end{lemma}

\begin{proof}
Les trois assertions résultent des
lemmes \ref{stacks-lemma-2-product-stacks},
\ref{stacks-lemma-2-product-stacks-in-groupoids} et
\ref{stacks-lemma-2-product-stacks-in-setoids},
ainsi que de l'équivalence de
Catégories, lemme \ref{categories-lemma-inertia-fibred-category}, partie (1).
\end{proof}

\begin{lemma}
\label{stacks-lemma-characterize-stack-in-setoids}
Soit $\mathcal{C}$ un site.
Si $\mathcal{S}$ est un champ en groupoïdes, alors le
$1$-morphisme canonique $\mathcal{I}_\mathcal{S} \to \mathcal{S}$
est une équivalence si et seulement si $\mathcal{S}$ est un champ en setoïdes.
\end{lemma}

\begin{proof}
Cela résulte directement de
Catégories, lemme \ref{categories-lemma-characterize-fibred-setoids-inertia}.
\end{proof}

\section{Champification des catégories fibrées}
\label{stacks-section-stackify}

\noindent
Le résultat de la procédure du lemme suivant sera appelé la
{\it champification} d'une catégorie fibrée sur un site.

\begin{lemma}
\label{stacks-lemma-stackify}
Soit $\mathcal{C}$ un site.
Soit $p : \mathcal{S} \to \mathcal{C}$ une catégorie fibrée sur $\mathcal{C}$.
Il existe un champ $p' : \mathcal{S}' \to \mathcal{C}$ et un
$1$-morphisme $G : \mathcal{S} \to \mathcal{S}'$
de catégories fibrées sur $\mathcal{C}$ (voir
Catégories, Définition \ref{categories-definition-fibred-categories-over-C})
tels que
\begin{enumerate}
\item pour tout $U \in \Ob(\mathcal{C})$ et tous
$x, y \in \Ob(\mathcal{S}_U)$, l'application
$$
\mathit{Mor}(x, y) \longrightarrow \mathit{Mor}(G(x), G(y))
$$
induite par $G$ identifie le membre de droite à la faisceautisation
du membre de gauche, et
\item pour tout $U \in \Ob(\mathcal{C})$ et tout
$x' \in \Ob(\mathcal{S}'_U)$, il existe un recouvrement
$\{U_i \to U\}_{i \in I}$ tel que, pour tout $i \in I$,
l'objet $x'|_{U_i}$ appartienne à l'image essentielle du
foncteur $G : \mathcal{S}_{U_i} \to \mathcal{S}'_{U_i}$.
\end{enumerate}
De plus, ces conditions déterminent le champ $\mathcal{S}'$ à un unique
$2$-isomorphisme près.
\end{lemma}

\begin{proof}[Démonstration par la méthode naïve]
Dans cette méthode de démonstration, nous procédons par étapes :

\medskip\noindent
Tout d'abord, étant donnés $x$ au-dessus de $U$ et un objet quelconque $y$ de
$\mathcal{S}$, on dit que deux morphismes
$a, b : x \to y$ de $\mathcal{S}$
au-dessus d'une même flèche de $\mathcal{C}$
sont {\it localement égaux}
s'il existe un recouvrement $\{f_i : U_i \to U\}$ de $\mathcal{C}$
tel que les composées
$$
f_i^*x \to x \xrightarrow{a} y,
\quad
f_i^*x \to x \xrightarrow{b} y
$$
soient égales. Cela définit une relation d'équivalence $\sim$
sur les flèches de $\mathcal{S}$. Si $b \sim b'$, alors
$a \circ b \circ c \sim a \circ b' \circ c$ (vérification omise).
On peut donc quotienter par cette relation d'équivalence pour
obtenir une nouvelle catégorie $\mathcal{S}^1$ sur $\mathcal{C}$,
ainsi qu'un morphisme $G^1 : \mathcal{S} \to \mathcal{S}^1$.

\medskip\noindent
On vérifie que $G^1$ préserve les morphismes fortement cartésiens
et que $\mathcal{S}^1$ est une catégorie fibrée sur $\mathcal{C}$.
Les vérifications sont omises. On se ramène ainsi au cas où des
morphismes localement égaux sont égaux.

\medskip\noindent
Ensuite, on ajoute des morphismes comme suit. Étant donnés
$x$ au-dessus de $U$ et un objet quelconque $y$ au-dessus de $V$,
un {\it morphisme de $x$ vers $y$ défini localement} consiste en
\begin{enumerate}
\item un morphisme $f : U \to V$,
\item un recouvrement $\{f_i : U_i \to U\}$ de $U$, et
\item des morphismes $a_i : f_i^*x \to y$ avec $p(a_i) = f \circ f_i$
\end{enumerate}
ayant la propriété que les composées
$$
(f_i \times f_j)^*x \to f_i^*x \xrightarrow{a_i} y,
\quad
(f_i \times f_j)^*x \to f_j^*x \xrightarrow{a_j} y
$$
sont égales. Remarquons qu'un morphisme usuel $a : x \to y$ fournit le
morphisme défini localement $(p(a) : U \to V, \{\text{id}_U\}, a)$.
On dit que deux morphismes définis localement
$(f, \{f_i : U_i \to U\}, a_i)$ et $(g, \{g_j : U'_j \to U\}, b_j)$
sont {\it égaux} si $f = g$ et si les composées
$$
(f_i \times g_j)^*x \to f_i^*x \xrightarrow{a_i} y,
\quad
(f_i \times g_j)^*x \to g_j^*x \xrightarrow{b_j} y
$$
sont égales (c'est la bonne condition puisque nous sommes dans le
cas où des morphismes localement égaux sont égaux).
Pour composer les morphismes définis localement
$(f, \{f_i : U_i \to U\}, a_i)$ de $x$ vers $y$ et
$(g, \{g_j : V_j \to V\}, b_j)$ de $y$ vers $z$ au-dessus de $W$,
il suffit de prendre $g \circ f : U \to W$, le recouvrement
$\{U_i \times_V V_j \to U\}$ et, comme applications, les composées
$$
x|_{U_i \times_V V_j}
\xrightarrow{\text{pr}_0^*a_i}
y|_{V_j}
\xrightarrow{b_j}
z
$$
Nous omettons la vérification qu'il s'agit d'un morphisme défini localement.

\medskip\noindent
On vérifie que $\mathcal{S}^2$, qui a les mêmes objets que
$\mathcal{S}$ et pour morphismes les morphismes définis localement,
est une catégorie sur $\mathcal{C}$, qu'il existe un foncteur
$G^2 : \mathcal{S} \to \mathcal{S}^2$ sur $\mathcal{C}$,
que ce foncteur préserve les objets fortement cartésiens
et que $\mathcal{S}^2$ est une catégorie fibrée sur $\mathcal{C}$.
Les vérifications sont omises. On se ramène ainsi au cas où tous les
préfaisceaux de morphismes de $\mathcal{S}$ sont des faisceaux, en
vérifiant que l'emploi de morphismes définis localement a pour effet de
faisceautiser les préfaisceaux (séparés) de morphismes.

\medskip\noindent
Enfin, lorsque les préfaisceaux de morphismes sont tous des faisceaux,
il faut ajouter des objets afin que les données de descente soient
effectives dans le résultat final. La façon la plus simple consiste à
considérer la catégorie $\mathcal{S}'$ dont les objets sont les
paires $(\mathcal{U}, \xi)$ où
$\mathcal{U} = \{U_i \to U\}$ est un recouvrement de $\mathcal{C}$ et
$\xi = (X_i, \varphi_{ii'})$ est une donnée de descente relativement à $\mathcal{U}$.
Supposons données deux telles données
$(\mathcal{U}, \xi) = (\{f_i : U_i \to U\},  x_i, \varphi_{ii'})$ et
$(\mathcal{V}, \eta) = (\{g_j : V_j \to V\},  y_j, \psi_{jj'})$.
On définit
$$
\Mor_{\mathcal{S}'}((\mathcal{U}, \xi), (\mathcal{V}, \eta))
$$
comme l'ensemble des $(f, a_{ij})$, où $f : U \to V$ et
$$
a_{ij} :
x_i|_{U_i \times_V V_j}
\longrightarrow
y_j
$$
sont des morphismes de $\mathcal{S}$ au-dessus de $U_i \times_V V_j \to V_j$.
Ils doivent satisfaire la condition suivante : pour tous
$i, i' \in I$ et $j, j' \in J$, posons
$W = (U_i \times_U U_{i'}) \times_V (V_j \times_V V_{j'})$. Alors
$$
\xymatrix{
x_i|_W \ar[r]_{a_{ij}|_W} \ar[d]_{\varphi_{ii'}|_W} &
y_j|_W \ar[d]^{\psi_{jj'}|_W} \\
x_{i'}|_W \ar[r]^{a_{i'j'}|_W} &
y_{j'}|_W
}
$$
est commutatif. Il faut alors vérifier les points suivants :
\begin{enumerate}
\item les morphismes ci-dessus admettent une composition bien définie,
\item cela fait de $\mathcal{S}'$ une catégorie sur $\mathcal{C}$,
\item il existe un foncteur $G : \mathcal{S} \to \mathcal{S}'$ sur $\mathcal{C}$,
\item si $x, y$ sont des objets de $\mathcal{S}$, alors
$\Mor_\mathcal{S}(x, y) = \Mor_{\mathcal{S}'}(G(x), G(y))$,
\item tout objet de $\mathcal{S}'$ provient localement d'un objet de
$\mathcal{S}$, autrement dit l'assertion (2) du lemme est satisfaite,
\item $G$ préserve les morphismes fortement cartésiens,
\item $\mathcal{S}'$ est une catégorie fibrée sur $\mathcal{C}$, et
\item $\mathcal{S}'$ est un champ sur $\mathcal{C}$.
\end{enumerate}
Rien de tout cela n'est difficile, mais les vérifications sont nombreuses. Détails omis.
\end{proof}

\begin{proof}[Démonstration moins naïve]
Voici une démonstration moins naïve.
D'après Catégories, Lemme \ref{categories-lemma-fibred-strict},
il existe une équivalence de
catégories fibrées $\mathcal{S} \to \mathcal{S}'$ telle que $\mathcal{S}'$
soit une catégorie fibrée scindée, c'est-à-dire une catégorie dans laquelle les foncteurs
d'image inverse se composent au sens strict. Le lemme pour $\mathcal{S}'$
implique évidemment le lemme pour $\mathcal{S}$. On peut donc considérer $\mathcal{S}$
comme un préfaisceau en catégories.

\medskip\noindent
Considérons temporairement la $2$-catégorie $\textit{Cat}$ comme une
catégorie, en oubliant les $2$-morphismes.
Considérons une catégorie comme un quintuplet
$(\text{Ob}, \text{Arrows}, s, t, \circ)$ comme dans
Catégories, section \ref{categories-section-definition-categories}.
Considérons le foncteur d'oubli
$$
forget : \textit{Cat} \to \textit{Ensembles} \times \textit{Ensembles}, \quad
(\text{Ob}, \text{Arrows}, s, t, \circ)
\mapsto
(\text{Ob}, \text{Arrows}).
$$
Alors $forget$ est fidèle, $\textit{Cat}$ admet des limites et
$forget$ commute à celles-ci, $\textit{Cat}$ admet des colimites filtrantes et
$forget$ commute à celles-ci, et $forget$ reflète les isomorphismes.
On peut faisceautiser les préfaisceaux à valeurs dans $\textit{Cat}$ et,
par un argument analogue à celui de la première partie de
Sites, section \ref{sites-section-sheaves-algebraic-structures},
le résultat commute avec $forget$. En appliquant ceci à
$\mathcal{S}$, on obtient une faisceautisation $\mathcal{S}^\#$
qui possède un faisceau d'objets et un faisceau de morphismes,
tous deux faisceautisations des préfaisceaux correspondants
de $\mathcal{S}$. Dans ce cas, il est assez
facile de voir que l'application $\mathcal{S} \to \mathcal{S}^\#$
possède les propriétés (1) et (2) du lemme.

\medskip\noindent
Cependant, la catégorie $\mathcal{S}^\#$ n'est pas nécessairement encore un
champ car, bien que le préfaisceau d'objets soit un faisceau,
la condition de descente peut ne pas être satisfaite.
Pour y remédier, il faut ajouter davantage d'objets. Mais l'argument
ci-dessus nous ramène au cas où $\mathcal{S} = \mathcal{S}_F$
pour un faisceau (!) de catégories
$F : \mathcal{C}^{opp} \to \textit{Cat}$. Dans ce cas, considérons le foncteur
$F' : \mathcal{C}^{opp} \to \textit{Cat}$ défini comme suit :
\begin{enumerate}
\item L'ensemble $\Ob(F'(U))$ est l'ensemble des paires
$(\mathcal{U}, \xi)$ où $\mathcal{U} = \{U_i \to U\}$
est un recouvrement de $U$ et $\xi = (x_i, \varphi_{ii'})$ est
une donnée de descente relativement à $\mathcal{U}$.
\item Un morphisme dans $F'(U)$ de
$(\mathcal{U}, \xi)$ vers $(\mathcal{V}, \eta)$
est un élément de
$$
\colim \Mor_{DD(\mathcal{W})}(a^*\xi, b^*\eta)
$$
où la colimite porte sur tous les raffinements communs
$a : \mathcal{W} \to \mathcal{U}$, $b : \mathcal{W} \to \mathcal{V}$.
Cette colimite est filtrante (vérification omise).
La composition des morphismes dans $F(U)$ est donc définie en
choisissant un raffinement commun puis en composant dans $DD(\mathcal{W})$.
\item Étant donnés $h : V \to U$ et un objet
$(\mathcal{U}, \xi)$ de $F'(U)$, on définit $F'(h)(\mathcal{U}, \xi)$
comme $(V \times_U \mathcal{U}, \text{pr}_1^*\xi)$.
Plus précisément, si $\mathcal{U} = \{U_i \to U\}$
et $\xi = (x_i, \varphi_{ii'})$, alors
$V \times_U \mathcal{U} = \{V \times_U U_i \to V\}$,
muni d'un morphisme canonique
$\text{pr}_1 : V \times_U \mathcal{U} \to \mathcal{U}$, et
$\text{pr}_1^*\xi$ est l'image inverse de $\xi$ par
ce morphisme (voir la Définition \ref{stacks-definition-pullback-functor}).
\item Soient $h : V \to U$, des objets $(\mathcal{U}, \xi)$
et $(\mathcal{V}, \eta)$, ainsi qu'un morphisme entre eux représenté par
$a : \mathcal{W} \to \mathcal{U}$, $b : \mathcal{W} \to \mathcal{V}$
et $\alpha : a^*\xi \to b^*\eta$ ; alors $F'(h)(\alpha)$ est
représenté par
$a' : V \times_U\mathcal{W} \to V \times_U\mathcal{U}$,
$b' : V \times_U\mathcal{W} \to V \times_U\mathcal{V}$
et l'image inverse $\alpha'$ du morphisme $\alpha$ par
l'application $V \times_U \mathcal{W} \to \mathcal{W}$. Cela fonctionne
car les images inverses dans $\mathcal{S}_F$ commutent au sens strict.
\end{enumerate}
Il existe une application $F \to F'$ qui associe à
un objet $x$ de $F(U)$ l'objet $(\{U \to U\}, (x, triv))$ de
$F'(U)$. Il reste alors à vérifier que le foncteur correspondant
$\mathcal{S}_F \to \mathcal{S}_{F'}$ possède les propriétés (1)
et (2) du lemme, et enfin que $\mathcal{S}_{F'}$ est
un champ. Détails omis.
\end{proof}

\begin{lemma}
\label{stacks-lemma-stackify-universal-property}
Soit $\mathcal{C}$ un site.
Soit $p : \mathcal{S} \to \mathcal{C}$ une catégorie fibrée sur $\mathcal{C}$.
Soient $p' : \mathcal{S}' \to \mathcal{C}$ et $G : \mathcal{S} \to \mathcal{S}'$
respectivement le champ et le $1$-morphisme construits dans le Lemme \ref{stacks-lemma-stackify}.
Cette construction possède la propriété universelle suivante : étant donnés un champ
$q : \mathcal{X} \to \mathcal{C}$ et un $1$-morphisme
$F : \mathcal{S} \to \mathcal{X}$ de catégories fibrées sur $\mathcal{C}$,
il existe un $1$-morphisme $H : \mathcal{S}' \to \mathcal{X}$
tel que le diagramme
$$
\xymatrix{
\mathcal{S} \ar[rr]_F \ar[rd]_G & & \mathcal{X} \\
& \mathcal{S}' \ar[ru]_H
}
$$
soit $2$-commutatif.
\end{lemma}

\begin{proof}
Omis. Indication : supposons que $x' \in \Ob(\mathcal{S}'_U)$.
D'après le Lemme \ref{stacks-lemma-stackify},
il existe un recouvrement $\{U_i \to U\}_{i \in I}$
tel que $x'|_{U_i} = G(x_i)$ pour un certain $x_i \in \Ob(\mathcal{S}_{U_i})$.
De plus, il existe des recouvrements $\{U_{ijk} \to U_i \times_U U_j\}$
et des isomorphismes $\alpha_{ijk} : x_i|_{U_{ijk}} \to x_j|_{U_{ijk}}$
tels que $G(\alpha_{ijk}) = \text{id}_{x'|_{U_{ijk}}}$. Posons $y_i = F(x_i)$.
On peut alors vérifier que les morphismes
$$
F(\alpha_{ijk}) : y_i|_{U_{ijk}} \to y_j|_{U_{ijk}}
$$
coïncident sur les intersections et définissent donc (puisque $\mathcal{X}$ est un champ)
un morphisme $\beta_{ij} : y_i|_{U_i \times_U U_j} \to y_j|_{U_i \times_U U_j}$.
On vérifie ensuite que les $\beta_{ij}$ définissent une donnée de descente. Comme
$\mathcal{X}$ est un champ, ces données de descente sont effectives et l'on trouve
un objet $y$ de $\mathcal{X}_U$ coïncidant avec $G(x_i)$ au-dessus de $U_i$.
L'indication consiste à poser $H(x') = y$.
\end{proof}

\begin{lemma}
\label{stacks-lemma-stackify-universal-property-more}
Les notations et les hypothèses sont celles du
Lemme \ref{stacks-lemma-stackify-universal-property}.
Il existe une équivalence canonique de catégories
$$
\Mor_{\textit{Fib}/\mathcal{C}}(\mathcal{S}, \mathcal{X})
=
\Mor_{\textit{Champs}/\mathcal{C}}(\mathcal{S}', \mathcal{X})
$$
donnée par les constructions de la démonstration du lemme précité.
\end{lemma}

\begin{proof}
Omis.
\end{proof}

\begin{lemma}
\label{stacks-lemma-stackification-fibre-product-fibred-categories}
Soit $\mathcal{C}$ un site.
Soient $f : \mathcal{X} \to \mathcal{Y}$ et $g : \mathcal{Z} \to \mathcal{Y}$
des morphismes de catégories fibrées sur $\mathcal{C}$.
Alors la champification du $2$-produit fibré est le $2$-produit
fibré des champifications.
\end{lemma}

\begin{proof}
Notons $\mathcal{X}', \mathcal{Y}', \mathcal{Z}'$ les champifications
et $\mathcal{W}$ la champification de
$\mathcal{X} \times_\mathcal{Y} \mathcal{Z}$. La construction des $2$-produits
fibrés fournit un $1$-morphisme canonique
$\mathcal{X} \times_\mathcal{Y} \mathcal{Z} \to
\mathcal{X}' \times_{\mathcal{Y}'} \mathcal{Z}'$.
Comme le second $2$-produit fibré est un champ (voir le
Lemme \ref{stacks-lemma-2-product-stacks}),
ce $1$-morphisme induit un $1$-morphisme
$h : \mathcal{W} \to \mathcal{X}' \times_{\mathcal{Y}'} \mathcal{Z}'$
par la propriété universelle de la champification, voir le
Lemme \ref{stacks-lemma-stackify-universal-property}.
Or $h$ est un morphisme de champs, et l'on peut vérifier que c'est une
équivalence à l'aide des
Lemmes \ref{stacks-lemma-characterize-ff} et
\ref{stacks-lemma-characterize-essentially-surjective-when-ff}.

\medskip\noindent
Montrons donc d'abord que $h$ induit des isomorphismes de faisceaux $\mathit{Mor}$.
Soient $\xi, \xi'$ des objets de $\mathcal{W}$ au-dessus de
$U \in \Ob(\mathcal{C})$. Nous voulons montrer que
$$
h : \mathit{Mor}(\xi, \xi') \longrightarrow \mathit{Mor}(h(\xi), h(\xi'))
$$
est un isomorphisme. Pour cela, on peut travailler localement sur $U$ (voir
Sites, section \ref{sites-section-glueing-sheaves}).
Par la construction de $\mathcal{W}$ (voir le
Lemme \ref{stacks-lemma-stackify}),
on peut donc supposer que $\xi, \xi'$ proviennent effectivement
des objets $(x, z, \alpha)$ et $(x', z', \alpha')$
de $\mathcal{X} \times_\mathcal{Y} \mathcal{Z}$ au-dessus de $U$.
Le même lemme montre encore que, dans ce cas,
$\mathit{Mor}(\xi, \xi')$ est la faisceautisation de
$$
V/U \longmapsto
\Mor_{\mathcal{X}_V}(x|_V, x'|_V)
\times_{\Mor_{\mathcal{Y}_V}(f(x)|_V, f(x')|_V)}
\Mor_{\mathcal{Z}_V}(z|_V, z'|_V)
$$
et que $\mathit{Mor}(h(\xi), h(\xi'))$ est égal au produit fibré
$$
\mathit{Mor}(i(x), i(x'))
\times_{\mathit{Mor}(j(f(x)), j(f(x'))}
\mathit{Mor}(k(z), k(z'))
$$
où $i : \mathcal{X} \to \mathcal{X}'$,
$j : \mathcal{Y} \to \mathcal{Y}'$ et
$k : \mathcal{Z} \to \mathcal{Z}'$ sont les foncteurs canoniques.
La première application affichée de ce paragraphe est donc un isomorphisme, car
la faisceautisation est exacte (et la faisceautisation d'un produit fibré
de préfaisceaux est par conséquent le produit fibré des faisceautisations).

\medskip\noindent
Enfin, il faut vérifier que tout objet de
$\mathcal{X}' \times_{\mathcal{Y}'} \mathcal{Z}'$
au-dessus de $U$ appartient localement sur $U$ à l'image essentielle de $h$.
Écrivons un tel objet sous la forme d'un triplet $(x', z', \alpha)$.
Alors $x'$ provient localement d'un objet de $\mathcal{X}$,
$z'$ provient localement d'un objet de $\mathcal{Z}$ et,
après avoir effectué des remplacements convenables de $x'$ et $z'$, le morphisme
$\alpha$ de $\mathcal{Y}'_U$ provient localement d'un morphisme de
$\mathcal{Y}$. Autrement dit, nous avons montré que tout objet de
$\mathcal{X}' \times_{\mathcal{Y}'} \mathcal{Z}'$
au-dessus de $U$ appartient localement sur $U$ à l'image essentielle de
$\mathcal{X} \times_\mathcal{Y} \mathcal{Z} \to
\mathcal{X}' \times_{\mathcal{Y}'} \mathcal{Z}'$ ; il appartient donc
a fortiori localement à l'image essentielle de $h$.
\end{proof}

\begin{lemma}
\label{stacks-lemma-stackification-inertia}
Soit $\mathcal{C}$ un site.
Soit $\mathcal{X}$ une catégorie fibrée sur $\mathcal{C}$.
La champification de la catégorie fibrée d'inertie $\mathcal{I}_\mathcal{X}$
est l'inertie de la champification de $\mathcal{X}$.
\end{lemma}

\begin{proof}
Cela résulte du fait que la champification est compatible
avec les $2$-produits fibrés d'après le
Lemme \ref{stacks-lemma-stackification-fibre-product-fibred-categories},
et du fait que l'inertie s'exprime en termes de
$2$-produits fibrés de catégories sur $\mathcal{C}$, voir
Catégories, Lemme \ref{categories-lemma-inertia-fibred-category}.
\end{proof}






\section{Champification des catégories fibrées en groupoïdes}
\label{stacks-section-stackify-groupoids}

\noindent
Voici le résultat.

\begin{lemma}
\label{stacks-lemma-stackify-groupoids}
Soit $\mathcal{C}$ un site.
Soit $p : \mathcal{S} \to \mathcal{C}$ une catégorie
fibrée en groupoïdes sur $\mathcal{C}$.
Il existe un champ en groupoïdes
$p' : \mathcal{S}' \to \mathcal{C}$ et un
$1$-morphisme $G : \mathcal{S} \to \mathcal{S}'$
de catégories fibrées en groupoïdes sur $\mathcal{C}$ (voir
Catégories, Définition
\ref{categories-definition-categories-fibred-in-groupoids-over-C})
tels que
\begin{enumerate}
\item pour tout $U \in \Ob(\mathcal{C})$ et tous
$x, y \in \Ob(\mathcal{S}_U)$, l'application
$$
\mathit{Mor}(x, y) \longrightarrow \mathit{Mor}(G(x), G(y))
$$
induite par $G$ identifie le membre de droite à la faisceautisation
du membre de gauche, et
\item pour tout $U \in \Ob(\mathcal{C})$ et tout
$x' \in \Ob(\mathcal{S}'_U)$, il existe un recouvrement
$\{U_i \to U\}_{i \in I}$ tel que, pour tout $i \in I$,
l'objet $x'|_{U_i}$ appartienne à l'image essentielle du
foncteur $G : \mathcal{S}_{U_i} \to \mathcal{S}'_{U_i}$.
\end{enumerate}
De plus, ces conditions déterminent le champ en groupoïdes $\mathcal{S}'$
à un unique $2$-isomorphisme près.
\end{lemma}

\begin{proof}
Appliquons le Lemme \ref{stacks-lemma-stackify}. Le résultat est un
champ en groupoïdes par le Lemme \ref{stacks-lemma-stack-in-groupoids-stack}.
\end{proof}

\begin{lemma}
\label{stacks-lemma-stackify-groupoids-universal-property}
Soit $\mathcal{C}$ un site.
Soit $p : \mathcal{S} \to \mathcal{C}$ une catégorie fibrée en groupoïdes
sur $\mathcal{C}$. Soient $p' : \mathcal{S}' \to \mathcal{C}$ et
$G : \mathcal{S} \to \mathcal{S}'$
respectivement le champ en groupoïdes et le $1$-morphisme construits dans le
Lemme \ref{stacks-lemma-stackify-groupoids}.
Cette construction possède la propriété universelle suivante : étant donnés un champ
en groupoïdes $q : \mathcal{X} \to \mathcal{C}$ et un $1$-morphisme
$F : \mathcal{S} \to \mathcal{X}$ de catégories sur $\mathcal{C}$,
il existe un $1$-morphisme $H : \mathcal{S}' \to \mathcal{X}$
tel que le diagramme
$$
\xymatrix{
\mathcal{S} \ar[rr]_F \ar[rd]_G & & \mathcal{X} \\
& \mathcal{S}' \ar[ru]_H
}
$$
soit $2$-commutatif.
\end{lemma}

\begin{proof}
C'est un cas particulier du
Lemme \ref{stacks-lemma-stackify-universal-property}.
\end{proof}

\begin{lemma}
\label{stacks-lemma-stackification-fibre-product-categories-fibred-in-groupoids}
Soit $\mathcal{C}$ un site.
Soient $f : \mathcal{X} \to \mathcal{Y}$ et $g : \mathcal{Z} \to \mathcal{Y}$
des morphismes de catégories fibrées en groupoïdes sur $\mathcal{C}$.
Alors la champification du $2$-produit fibré est le $2$-produit
fibré des champifications.
\end{lemma}

\begin{proof}
C'est un cas particulier du
Lemme \ref{stacks-lemma-stackification-fibre-product-fibred-categories}.
\end{proof}





\section{Topologies héritées}
\label{stacks-section-topology}

\noindent
Il s'avère qu'une catégorie fibrée sur un site hérite d'une topologie
canonique du site sous-jacent.

\begin{lemma}
\label{stacks-lemma-topology-inherited}
Soit $\mathcal{C}$ un site. Soit $p : \mathcal{S} \to \mathcal{C}$
une catégorie fibrée. Soit $\text{Cov}(\mathcal{S})$
l'ensemble des familles $\{x_i \to x\}_{i \in I}$ de morphismes de $\mathcal{S}$
de but fixé telles que (a) chaque $x_i \to x$ soit fortement cartésien
et (b) $\{p(x_i) \to p(x)\}_{i \in I}$ soit un recouvrement de $\mathcal{C}$.
Alors $(\mathcal{S}, \text{Cov}(\mathcal{S}))$ est un site.
\end{lemma}

\begin{proof}
Il faut vérifier les trois conditions de
Sites, Définition \ref{sites-definition-site}.
\begin{enumerate}
\item Si $x \to y$ est un isomorphisme de $\mathcal{S}$, alors
il est fortement cartésien d'après
Catégories, Lemme \ref{categories-lemma-composition-cartesian},
et $p(x) \to p(y)$ est un isomorphisme de $\mathcal{C}$. Par conséquent,
$\{p(x) \to p(y)\}$ est un recouvrement de $\mathcal{C}$, d'où
$\{x \to y\} \in \text{Cov}(\mathcal{S})$.
\item Si $\{x_i \to x\}_{i\in I} \in \text{Cov}(\mathcal{S})$ et si, pour tout
$i$, on a $\{y_{ij} \to x_i\}_{j\in J_i} \in \text{Cov}(\mathcal{S})$, alors
chaque composée $y_{ij} \to x$ est fortement cartésienne d'après
Catégories, Lemme \ref{categories-lemma-composition-cartesian},
et $\{p(y_{ij}) \to p(x)\}_{i \in I, j\in J_i} \in \text{Cov}(\mathcal{C})$.
On a donc aussi $\{y_{ij} \to x\}_{i \in I, j\in J_i} \in \text{Cov}(\mathcal{S})$.
\item Supposons $\{x_i \to x\}_{i\in I}\in \text{Cov}(\mathcal{S})$
et soit $y \to x$ un morphisme de $\mathcal{S}$. Puisque $\{p(x_i) \to p(x)\}$
est un recouvrement de $\mathcal{C}$, le produit fibré
$p(x_i) \times_{p(x)} p(y)$ existe. Par conséquent,
Catégories, Lemme \ref{categories-lemma-fibred-category-representable-goes-up}
implique que $x_i \times_x y$ existe, que $p(x_i \times_x y) =
p(x_i) \times_{p(x)} p(y)$ et que $x_i \times_x y \to y$ est
fortement cartésien. Comme, de plus,
$\{p(x_i) \times_{p(x)} p(y) \to p(y) \}_{i\in I} \in \text{Cov}(\mathcal{C})$,
on en conclut que
$\{x_i \times_x y \to y \}_{i\in I} \in \text{Cov}(\mathcal{S})$.
\end{enumerate}
Cela achève la démonstration.
\end{proof}

\noindent
Notons que, si $p : \mathcal{S} \to \mathcal{C}$ est fibré en groupoïdes, alors
les recouvrements du site $\mathcal{S}$ du
Lemme \ref{stacks-lemma-topology-inherited}
se caractérisent par
$$
\{x_i \to x\} \in \text{Cov}(\mathcal{S})
\Leftrightarrow
\{p(x_i) \to p(x)\} \in \text{Cov}(\mathcal{C})
$$
car tout morphisme de $\mathcal{S}$ est fortement cartésien.

\begin{definition}
\label{stacks-definition-topology-inherited}
Soit $\mathcal{C}$ un site. Soit $p : \mathcal{S} \to \mathcal{C}$ une
catégorie fibrée. On dit que $(\mathcal{S}, \text{Cov}(\mathcal{S}))$, comme dans le
Lemme \ref{stacks-lemma-topology-inherited},
est la {\it structure de site sur $\mathcal{S}$ héritée de $\mathcal{C}$}.
On l'exprime parfois en disant que
{\it $\mathcal{S}$ est munie de la topologie héritée de $\mathcal{C}$}.
\end{definition}

\noindent
On obtient en particulier, dans cette situation, un topos de faisceaux
$\Sh(\mathcal{S})$. Il s'avère que ce topos est fonctoriel par rapport
aux $1$-morphismes de catégories fibrées.

\begin{lemma}
\label{stacks-lemma-topology-inherited-functorial}
Soit $\mathcal{C}$ un site. Soit $F : \mathcal{X} \to \mathcal{Y}$
un $1$-morphisme de catégories fibrées sur $\mathcal{C}$.
Alors $F$ est un foncteur continu et cocontinu entre les structures
de site héritées de $\mathcal{C}$. Ainsi, $F$ induit un morphisme de topos
$f : \Sh(\mathcal{X}) \to \Sh(\mathcal{Y})$ avec
$f_* = {}_sF = {}_pF$ et $f^{-1} = F^s = F^p$. En particulier,
$f^{-1}(\mathcal{G})(x) = \mathcal{G}(F(x))$
pour tout faisceau $\mathcal{G}$ sur $\mathcal{Y}$ et tout objet $x$ de $\mathcal{X}$.
\end{lemma}

\begin{proof}
Montrons d'abord que $F$ est continu.
Soit $\{x_i \to x\}_{i \in I}$ un recouvrement de $\mathcal{X}$. D'après
Catégories, Définition \ref{categories-definition-fibred-categories-over-C},
le foncteur $F$ transforme les morphismes fortement cartésiens en morphismes
fortement cartésiens ; ainsi, $\{F(x_i) \to F(x)\}_{i \in I}$ est un recouvrement
de $\mathcal{Y}$. Cela prouve la partie (1) de
Sites, Définition \ref{sites-definition-continuous}.
De plus, soit $x' \to x$ un morphisme de $\mathcal{X}$. D'après
Catégories, Lemme \ref{categories-lemma-fibred-category-representable-goes-up},
le produit fibré $x_i \times_x x'$ existe et $x_i \times_x x' \to x'$
est fortement cartésien. Ainsi, $F(x_i \times_x x') \to F(x')$ est fortement
cartésien. D'après
Catégories, Lemme \ref{categories-lemma-fibred-category-representable-goes-up},
appliqué à $\mathcal{Y}$, cela signifie que
$F(x_i \times_x x') = F(x_i) \times_{F(x)} F(x')$.
Cela prouve la partie (2) de
Sites, Définition \ref{sites-definition-continuous},
et l'on en conclut que $F$ est continu.

\medskip\noindent
Montrons ensuite que $F$ est cocontinu.
Soit $x \in \Ob(\mathcal{X})$ et soit $\{y_i \to F(x)\}_{i \in I}$
un recouvrement de $\mathcal{Y}$. Notons $\{U_i \to U\}_{i \in I}$ le
recouvrement correspondant de $\mathcal{C}$. Pour chaque $i$, choisissons un morphisme
fortement cartésien $x_i \to x$ de $\mathcal{X}$ au-dessus de $U_i \to U$.
Alors $F(x_i) \to F(x)$ et $y_i \to F(x)$ sont tous deux des morphismes
fortement cartésiens de $\mathcal{Y}$ au-dessus de $U_i \to U$. Il existe donc
un unique isomorphisme $F(x_i) \to y_i$ dans $\mathcal{Y}_{U_i}$ compatible
avec les applications vers $F(x)$. Ainsi, $\{x_i \to x\}_{i \in I}$ est un recouvrement de
$\mathcal{X}$ tel que $\{F(x_i) \to F(x)\}_{i \in I}$ soit isomorphe à
$\{y_i \to F(x)\}_{i \in I}$. Le foncteur $F$ est donc cocontinu, voir
Sites, Définition \ref{sites-definition-cocontinuous}.

\medskip\noindent
L'assertion finale découle des deux premières, voir
Sites, Lemmes
\ref{sites-lemma-cocontinuous-morphism-topoi},
\ref{sites-lemma-pu-sheaf} et
\ref{sites-lemma-when-shriek}.
\end{proof}

\begin{lemma}
\label{stacks-lemma-localizing}
Soit $\mathcal{C}$ un site. Soit $p : \mathcal{X} \to \mathcal{C}$
une catégorie fibrée en groupoïdes. Soit $x \in \Ob(\mathcal{X})$
au-dessus de $U = p(x)$. Le foncteur $p$ induit une équivalence de sites
$\mathcal{X}/x \to \mathcal{C}/U$, où $\mathcal{X}$ est munie de
la topologie héritée de $\mathcal{C}$.
\end{lemma}

\begin{proof}
Ici, $\mathcal{C}/U$ est la localisation du site
$\mathcal{C}$ en l'objet $U$, et de même pour $\mathcal{X}/x$.
Il résulte de
Catégories, Définition \ref{categories-definition-fibred-groupoids},
que la règle $x'/x \mapsto p(x')/p(x)$ définit une équivalence de
catégories $\mathcal{X}/x \to \mathcal{C}/U$. Il résulte alors de la
Définition \ref{stacks-definition-topology-inherited}
que les recouvrements de $x'$ dans $\mathcal{X}/x$ correspondent bijectivement
aux recouvrements de $p(x')$ dans $\mathcal{C}/U$.
\end{proof}

\begin{lemma}
\label{stacks-lemma-stack-in-groupoids-over-stack-in-groupoids}
Soit $\mathcal{C}$ un site. Soient $p : \mathcal{X} \to \mathcal{C}$
et $q : \mathcal{Y} \to \mathcal{C}$
des champs en groupoïdes. Soit $F : \mathcal{X} \to \mathcal{Y}$
un $1$-morphisme de catégories sur $\mathcal{C}$. Si $F$ fait de
$\mathcal{X}$ une catégorie fibrée en groupoïdes sur $\mathcal{Y}$,
alors $\mathcal{X}$ est un champ en groupoïdes sur $\mathcal{Y}$ (pour la
topologie héritée de $\mathcal{C}$).
\end{lemma}

\begin{proof}
Montrons la descente des objets. Soit $\{y_i \to y\}$ un recouvrement
de $\mathcal{Y}$. Soit $(x_i, \varphi_{ij})$ une donnée de descente dans
$\mathcal{X}$ relativement à ce recouvrement. Alors $(x_i, \varphi_{ij})$ est aussi une donnée
de descente relativement au recouvrement $\{q(y_i) \to q(y)\}$ de
$\mathcal{C}$. Comme $\mathcal{X}$ est un champ en groupoïdes, on obtient un
objet $x$ au-dessus de $q(y)$ et des isomorphismes
$\psi_i : x|_{q(y_i)} \to x_i$
au-dessus de $q(y_i)$, compatibles avec les $\varphi_{ij}$, c'est-à-dire tels que
$$
\varphi_{ij} = \psi_j|_{q(y_i) \times_{q(y)} q(y_j)}
\circ \psi_i^{-1}|_{q(y_i) \times_{q(y)} q(y_j)}.
$$
Considérons le faisceau $\mathit{I} = \mathit{Isom}_\mathcal{Y}(F(x), y)$
sur $\mathcal{C}/p(x)$. Notons que $s_i = F(\psi_i) \in \mathit{I}(q(x_i))$
car $F(x_i) = y_i$. Puisque $F(\varphi_{ij}) = \text{id}$ (car nous sommes partis
d'une donnée de descente sur $\{y_i \to y\}$), la formule affichée montre
que $s_i|_{q(y_i) \times_{q(y)} q(y_j)} = s_j|_{q(y_i) \times_{q(y)} q(y_j)}$.
Les sections locales $s_i$ se recollent donc en $s : F(x) \to y$. Comme $F$ est fibré en
groupoïdes, on voit que $x$ est isomorphe à un objet $x'$ tel que $F(x') = y$.
Nous omettons de vérifier que $x'$, dans la catégorie fibre de $\mathcal{X}$
au-dessus de $y$, résout le problème de descente posé par la donnée de descente
$(x_i, \varphi_{ij})$. Nous omettons aussi la démonstration de la propriété de faisceau
des préfaisceaux $\mathit{Isom}$ de $\mathcal{X}/\mathcal{Y}$.
\end{proof}

\begin{lemma}
\label{stacks-lemma-stack-over-stack}
Soit $\mathcal{C}$ un site. Soit $p : \mathcal{X} \to \mathcal{C}$
un champ. Munissons $\mathcal{X}$ de la topologie héritée de
$\mathcal{C}$ et soit $q : \mathcal{Y} \to \mathcal{X}$ un champ.
Alors $\mathcal{Y}$ est un champ sur $\mathcal{C}$.
Si $p$ et $q$ définissent des champs en groupoïdes, alors
$\mathcal{Y}$ est un champ en groupoïdes sur $\mathcal{C}$.
\end{lemma}

\begin{proof}
Vérifions les trois conditions de la Définition \ref{stacks-definition-stack}
pour montrer que $\mathcal{Y}$ est un champ sur $\mathcal{C}$.
D'après Catégories, Lemme \ref{categories-lemma-fibred-over-fibred},
$\mathcal{Y}$ est une catégorie fibrée sur $\mathcal{C}$.
La condition (1) est donc satisfaite.

\medskip\noindent
Soit $U$ un objet de $\mathcal{C}$ et soient $y_1, y_2$ des objets
de $\mathcal{Y}$ au-dessus de $U$. Notons $x_i = q(y_i)$ dans $\mathcal{X}$.
Considérons l'application de préfaisceaux
$$
q : \mathit{Mor}_{\mathcal{Y}/\mathcal{C}}(y_1, y_2)
\longrightarrow
\mathit{Mor}_{\mathcal{X}/\mathcal{C}}(x_1, x_2)
$$
sur $\mathcal{C}/U$, voir le Lemme \ref{stacks-lemma-presheaf-mor-map-fibred-categories}.
Soit $\{U_i \to U\}$ un recouvrement et soient les $\varphi_i$ des sections
du préfaisceau de gauche sur les $U_i$, telles que $\varphi_i$ et
$\varphi_j$ aient la même restriction sur $U_i \times_U U_j$.
Il faut trouver un morphisme $\varphi : y_1 \to y_2$ dont la restriction soit $\varphi_i$.
Notons que $q(\varphi_i) = \psi|_{U_i}$ pour un certain morphisme
$\psi : x_1 \to x_2$ au-dessus de $U$, car le second préfaisceau est un faisceau
(par hypothèse). Soit $y_{12} \to y_2$ le morphisme fortement
$\mathcal{X}$-cartésien de $\mathcal{Y}$ au-dessus de $\psi$. Alors $\varphi_i$ correspond
à un morphisme $\varphi'_i : y_1|_{U_i} \to y_{12}|_{U_i}$ au-dessus de $x_1|_{U_i}$.
Autrement dit, les $\varphi'_i$ définissent maintenant des sections locales du préfaisceau
$$
\mathit{Mor}_{\mathcal{Y}/\mathcal{X}}(y_1, y_{12})
$$
sur les membres du recouvrement $\{x_1|_{U_i} \to x_1\}$. Par hypothèse, celles-ci
se recollent en un unique morphisme $y_1 \to y_{12}$ qui, composé avec le morphisme
donné $y_{12} \to y_2$, fournit le morphisme recherché $y_1 \to y_2$.

\medskip\noindent
Enfin, montrons que les données de descente sont effectives. Soit $\{f_i : U_i \to U\}$
un recouvrement de $\mathcal{C}$ et soit $(y_i, \varphi_{ij})$ une donnée de descente
relativement à ce recouvrement (Définition \ref{stacks-definition-descent-data}).
En posant $x_i = q(y_i)$ et $\psi_{ij} = q(\varphi_{ij})$,
on obtient une donnée de descente $(x_i, \psi_{ij})$ pour le recouvrement dans $\mathcal{X}$.
Par l'hypothèse sur $\mathcal{X}$, on peut supposer que $x_i = x|_{U_i}$
et que les $\psi_{ij}$ sont égaux à la donnée de descente canonique
(Définition \ref{stacks-definition-effective-descent-datum}).
Dans ce cas, $\{x|_{U_i} \to x\}$ est un recouvrement et l'on peut considérer
$(y_i, \varphi_{ij})$ comme une donnée de descente relativement à ce recouvrement.
Notre hypothèse selon laquelle $\mathcal{Y}$ est un champ sur $\mathcal{C}$
montre qu'elle est effective, ce qui achève la démonstration de la condition (3).

\medskip\noindent
L'assertion finale résulte de ce que $\mathcal{Y}$ est un champ sur
$\mathcal{C}$ et est fibrée en groupoïdes d'après
Catégories, Lemme
\ref{categories-lemma-fibred-in-groupoids-over-fibred-in-groupoids}.
\end{proof}





\section{Gerbes}
\label{stacks-section-gerbes}

\noindent
Les gerbes sont une classe particulière de champs en groupoïdes.

\begin{definition}
\label{stacks-definition-gerbe}
Une {\it gerbe} sur un site $\mathcal{C}$ est une catégorie
$p : \mathcal{S} \to \mathcal{C}$ sur $\mathcal{C}$ telle que
\begin{enumerate}
\item $p : \mathcal{S} \to \mathcal{C}$ soit un champ
en groupoïdes sur $\mathcal{C}$ (voir la
Définition \ref{stacks-definition-stack-in-groupoids}),
\item pour tout $U \in \Ob(\mathcal{C})$, il existe
un recouvrement $\{U_i \to U\}$ dans $\mathcal{C}$ tel que
$\mathcal{S}_{U_i}$ soit non vide, et
\item pour tout $U \in \Ob(\mathcal{C})$ et tous
$x, y \in \Ob(\mathcal{S}_U)$, il existe
un recouvrement $\{U_i \to U\}$ dans $\mathcal{C}$ tel que
$x|_{U_i} \cong y|_{U_i}$ dans $\mathcal{S}_{U_i}$.
\end{enumerate}
\end{definition}

\noindent
Autrement dit, une gerbe est un champ en groupoïdes dans lequel deux objets
quelconques sont localement isomorphes et où des objets existent localement.

\begin{lemma}
\label{stacks-lemma-gerbe-equivalent}
Soit $\mathcal{C}$ un site.
Soient $\mathcal{S}_1$ et $\mathcal{S}_2$ des catégories sur $\mathcal{C}$.
Supposons que $\mathcal{S}_1$ et $\mathcal{S}_2$ soient équivalentes
en tant que catégories sur $\mathcal{C}$.
Alors $\mathcal{S}_1$ est une gerbe sur $\mathcal{C}$ si et seulement si
$\mathcal{S}_2$ est une gerbe sur $\mathcal{C}$.
\end{lemma}

\begin{proof}
Supposons que $\mathcal{S}_1$ soit une gerbe sur $\mathcal{C}$. D'après le
Lemme \ref{stacks-lemma-stack-in-groupoids-equivalent},
$\mathcal{S}_2$ est un champ en groupoïdes sur $\mathcal{C}$.
Soient $F : \mathcal{S}_1 \to \mathcal{S}_2$ et
$G : \mathcal{S}_2 \to \mathcal{S}_1$ des équivalences de catégories sur
$\mathcal{C}$. Étant donné $U \in \Ob(\mathcal{C})$, il existe
un recouvrement $\{U_i \to U\}$ tel que $(\mathcal{S}_1)_{U_i}$ soit
non vide. En appliquant $F$, on voit que $(\mathcal{S}_2)_{U_i}$ est non vide.
Étant donnés $U \in \Ob(\mathcal{C})$ et
$x, y \in \Ob((\mathcal{S}_2)_U)$, il existe
un recouvrement $\{U_i \to U\}$ dans $\mathcal{C}$ tel que
$G(x)|_{U_i} \cong G(y)|_{U_i}$ dans $(\mathcal{S}_1)_{U_i}$. D'après
Catégories, Lemme \ref{categories-lemma-equivalence-fibred-categories},
cela implique que $x|_{U_i} \cong y|_{U_i}$ dans $(\mathcal{S}_2)_{U_i}$.
\end{proof}

\noindent
Nous voulons généraliser quelque peu la définition des gerbes. Soit donc
$F : \mathcal{X} \to \mathcal{Y}$ un $1$-morphisme de champs en groupoïdes
sur un site $\mathcal{C}$. Nous voulons préciser ce que signifie
pour $\mathcal{X}$ d'être une gerbe sur $\mathcal{Y}$. D'après la
section \ref{stacks-section-topology},
la catégorie $\mathcal{Y}$ hérite de $\mathcal{C}$ d'une structure de site.
Une idée naïve serait d'exiger simplement que
$\mathcal{X} \to \mathcal{Y}$ {\it soit} une gerbe au sens ci-dessus. Toutefois,
la notion ainsi obtenue n'est pas invariante lorsque l'on remplace $\mathcal{X}$
par un champ en groupoïdes équivalent sur $\mathcal{C}$ ; c'est même le
cas de la propriété d'être fibrée en groupoïdes sur $\mathcal{Y}$.
Il s'avère cependant que l'on peut remplacer $\mathcal{X}$ par un champ
en groupoïdes équivalent sur $\mathcal{C}$ qui soit fibré en groupoïdes sur
$\mathcal{Y}$, et que la propriété d'être une gerbe sur $\mathcal{Y}$
est alors indépendante de ce choix. Voici l'énoncé précis.

\begin{lemma}
\label{stacks-lemma-when-gerbe}
Soit $\mathcal{C}$ un site. Soient $p : \mathcal{X} \to \mathcal{C}$
et $q : \mathcal{Y} \to \mathcal{C}$ des champs en groupoïdes.
Soit $F : \mathcal{X} \to \mathcal{Y}$ un $1$-morphisme de catégories
sur $\mathcal{C}$. Les assertions suivantes sont équivalentes :
\begin{enumerate}
\item Pour une certaine factorisation (de manière équivalente, pour toute factorisation)
$F = F' \circ a$ où
$a : \mathcal{X} \to \mathcal{X}'$ est une équivalence de catégories sur
$\mathcal{C}$ et où $F'$ est fibré en groupoïdes, l'application
$F' : \mathcal{X}' \to \mathcal{Y}$ est une gerbe (pour la topologie
sur $\mathcal{Y}$ héritée de $\mathcal{C}$).
\item Les deux conditions suivantes sont satisfaites :
\begin{enumerate}
\item pour tout $y \in \Ob(\mathcal{Y})$ au-dessus de
$U \in \Ob(\mathcal{C})$, il existe un recouvrement
$\{U_i \to U\}$ dans $\mathcal{C}$ et des objets $x_i$ de $\mathcal{X}$
au-dessus de $U_i$ tels que $F(x_i) \cong y|_{U_i}$ dans $\mathcal{Y}_{U_i}$, et
\item pour $U \in \Ob(\mathcal{C})$,
$x, x' \in \Ob(\mathcal{X}_U)$ et $b : F(x) \to F(x')$ dans
$\mathcal{Y}_U$, il existe
un recouvrement $\{U_i \to U\}$ dans $\mathcal{C}$ et des morphismes
$a_i : x|_{U_i} \to x'|_{U_i}$ dans $\mathcal{X}_{U_i}$ tels que
$F(a_i) = b|_{U_i}$.
\end{enumerate}
\end{enumerate}
\end{lemma}

\begin{proof}
D'après
Catégories, Lemme
\ref{categories-lemma-ameliorate-morphism-categories-fibred-groupoids},
il existe une factorisation $F = F' \circ a$ où
$a : \mathcal{X} \to \mathcal{X}'$ est une équivalence de catégories sur
$\mathcal{C}$ et où $F'$ est fibré en groupoïdes. D'après
Catégories, Lemme \ref{categories-lemma-amelioration-unique},
étant données deux telles factorisations $F = F' \circ a = F'' \circ b$,
on a une équivalence de $\mathcal{X}'$ et $\mathcal{X}''$ en tant que
catégories sur $\mathcal{Y}$. Par conséquent, le
Lemme \ref{stacks-lemma-gerbe-equivalent}
garantit que la condition (1) ne dépend pas du choix de la
factorisation. De plus, cela signifie que l'on peut supposer
$\mathcal{X}' = \mathcal{X} \times_{F, \mathcal{Y}, \text{id}} \mathcal{Y}$
comme dans la démonstration de
Catégories, Lemme
\ref{categories-lemma-ameliorate-morphism-categories-fibred-groupoids}.

\medskip\noindent
Montrons que (a) et (b) impliquent que $\mathcal{X}' \to \mathcal{Y}$
est une gerbe. Tout d'abord, d'après le
Lemme \ref{stacks-lemma-stack-in-groupoids-over-stack-in-groupoids},
$\mathcal{X}' \to \mathcal{Y}$ est un champ en groupoïdes.
Ensuite, soit $y$ un objet de $\mathcal{Y}$ au-dessus de
$U \in \Ob(\mathcal{C})$. D'après (a), on peut trouver un recouvrement
$\{U_i \to U\}$ dans $\mathcal{C}$, des objets $x_i$ de $\mathcal{X}$
au-dessus de $U_i$ et des isomorphismes $f_i : F(x_i) \to y|_{U_i}$ dans
$\mathcal{Y}_{U_i}$. Alors les $(U_i, x_i, y|_{U_i}, f_i)$ sont des objets
de $\mathcal{X}'_{U_i}$ ; la deuxième condition de la
Définition \ref{stacks-definition-gerbe}
est donc satisfaite. Enfin, soient $(U, x, y, f)$ et $(U, x', y, f')$ des objets
de $\mathcal{X}'$ au-dessus du même objet $y \in \Ob(\mathcal{Y})$.
Posons $b = (f')^{-1} \circ f$. D'après la condition (b), on peut trouver un recouvrement
$\{U_i \to U\}$ et des isomorphismes $a_i : x|_{U_i} \to x'|_{U_i}$
dans $\mathcal{X}_{U_i}$ tels que $F(a_i) = b|_{U_i}$. Alors
$$
(a_i, \text{id}) : (U, x, y, f)|_{U_i} \to (U, x', y, f')|_{U_i}
$$
est un morphisme dans $\mathcal{X}'_{U_i}$ comme voulu. Cela prouve que
(2) implique (1).

\medskip\noindent
Pour montrer que (1) implique (2), il suffit de lire à rebours les arguments
du paragraphe précédent. Détails omis.
\end{proof}

\begin{definition}
\label{stacks-definition-gerbe-over-stack-in-groupoids}
Soit $\mathcal{C}$ un site. Soient $\mathcal{X}$
et $\mathcal{Y}$ des champs en groupoïdes sur $\mathcal{C}$.
Soit $F : \mathcal{X} \to \mathcal{Y}$ un $1$-morphisme de catégories
sur $\mathcal{C}$. On dit que $\mathcal{X}$ est une {\it gerbe sur} $\mathcal{Y}$
si les conditions équivalentes du
Lemme \ref{stacks-lemma-when-gerbe}
sont satisfaites.
\end{definition}

\noindent
Cette définition n'entre pas en conflit avec la
Définition \ref{stacks-definition-gerbe}
lorsque $\mathcal{Y} = \mathcal{C}$, car on peut alors prendre
$\mathcal{X}' = \mathcal{X}$ dans la partie (1) du
Lemme \ref{stacks-lemma-when-gerbe}.
Notons que les conditions (2)(a) et (2)(b) du
Lemme \ref{stacks-lemma-when-gerbe}
sont très proches, dans leur esprit, des conditions (2) et (3) de la
Définition \ref{stacks-definition-gerbe}.
En effet, (2)(a) affirme que l'application des préfaisceaux de classes
d'isomorphisme d'objets devient une surjection après faisceautisation.
De plus, (2)(b) affirme que
$$
\mathit{Isom}_\mathcal{X}(x, x')
\longrightarrow
\mathit{Isom}_\mathcal{Y}(F(x), F(x'))
$$
est une surjection de faisceaux sur $\mathcal{C}/U$, pour tous $U$ et
$x, x' \in \Ob(\mathcal{X}_U)$.

\begin{lemma}
\label{stacks-lemma-base-change-gerbe}
Soit $\mathcal{C}$ un site. Soit
$$
\xymatrix{
\mathcal{X}' \ar[r]_{G'} \ar[d]_{F'} & \mathcal{X} \ar[d]^F \\
\mathcal{Y}' \ar[r]^G & \mathcal{Y}
}
$$
un $2$-produit fibré de champs en groupoïdes sur $\mathcal{C}$.
Si $\mathcal{X}$ est une gerbe sur $\mathcal{Y}$, alors
$\mathcal{X}'$ est une gerbe sur $\mathcal{Y}'$.
\end{lemma}

\begin{proof}
Par la propriété d'unicité d'un $2$-produit fibré, on peut supposer que
$\mathcal{X}' = \mathcal{Y}' \times_\mathcal{Y} \mathcal{X}$
comme dans
Catégories, Lemme \ref{categories-lemma-2-product-categories-over-C}.
Montrons les propriétés (2)(a) et (2)(b) du
Lemme \ref{stacks-lemma-when-gerbe}
pour $\mathcal{Y}' \times_\mathcal{Y} \mathcal{X} \to \mathcal{Y}'$.

\medskip\noindent
Soit $y'$ un objet de $\mathcal{Y}'$ au-dessus de l'objet $U$
de $\mathcal{C}$. Par hypothèse, il existe
un recouvrement $\{U_i \to U\}$ de $U$ et des objets
$x_i \in \mathcal{X}_{U_i}$ munis d'isomorphismes
$\alpha_i : G(y')|_{U_i} \to F(x_i)$.
Alors $(U_i, y'|_{U_i}, x_i, \alpha_i)$ est un objet de
$\mathcal{Y}' \times_\mathcal{Y} \mathcal{X}$ au-dessus de $U_i$
dont l'image dans $\mathcal{Y}'$ est $y'|_{U_i}$. La condition (2)(a) est donc satisfaite.

\medskip\noindent
Soit $U \in \Ob(\mathcal{C})$, soient $x'_1, x'_2$ des objets de
$\mathcal{Y}' \times_\mathcal{Y} \mathcal{X}$ au-dessus de $U$, et soit
$b' : F'(x'_1) \to F'(x'_2)$ un morphisme dans $\mathcal{Y}'_U$.
Écrivons $x'_i = (U, y'_i, x_i, \alpha_i)$. Notons que $F'(x'_i) = x_i$
et $G'(x'_i) = y'_i$. Par hypothèse, il existe
un recouvrement $\{U_i \to U\}$ dans $\mathcal{C}$ et des morphismes
$a_i : x_1|_{U_i} \to x_2|_{U_i}$ dans $\mathcal{X}_{U_i}$ tels que
$F(a_i) = G(b')|_{U_i}$. Alors $(b'|_{U_i}, a_i)$ est un morphisme
$x'_1|_{U_i} \to x'_2|_{U_i}$ comme l'exige (2)(b).
\end{proof}

\begin{lemma}
\label{stacks-lemma-composition-gerbe}
Soit $\mathcal{C}$ un site. Soient
$F : \mathcal{X} \to \mathcal{Y}$ et $G : \mathcal{Y} \to \mathcal{Z}$
des $1$-morphismes de champs en groupoïdes sur $\mathcal{C}$.
Si $\mathcal{X}$ est une gerbe sur $\mathcal{Y}$ et si
$\mathcal{Y}$ est une gerbe sur $\mathcal{Z}$, alors
$\mathcal{X}$ est une gerbe sur $\mathcal{Z}$.
\end{lemma}

\begin{proof}
Montrons les propriétés (2)(a) et (2)(b) du
Lemme \ref{stacks-lemma-when-gerbe}
pour $\mathcal{X} \to \mathcal{Z}$.

\medskip\noindent
Soit $z$ un objet de $\mathcal{Z}$ au-dessus de l'objet $U$
de $\mathcal{C}$. Par l'hypothèse sur $G$, il existe
un recouvrement $\{U_i \to U\}$ de $U$ et des objets
$y_i \in \mathcal{Y}_{U_i}$ tels que $G(y_i) \cong z|_{U_i}$.
Par l'hypothèse sur $F$, il existe des recouvrements $\{U_{ij} \to U_i\}$
et des objets $x_{ij} \in \mathcal{X}_{U_{ij}}$ tels que
$F(x_{ij}) \cong y_i|_{U_{ij}}$.
Alors $\{U_{ij} \to U\}$ est un recouvrement de $\mathcal{C}$
et $(G \circ F)(x_{ij}) \cong z|_{U_{ij}}$. La condition (2)(a) est donc satisfaite.

\medskip\noindent
Soit $U \in \Ob(\mathcal{C})$, soient $x_1, x_2$ des objets de
$\mathcal{X}$ au-dessus de $U$, et soit
$c : (G \circ F)(x_1) \to (G \circ F)(x_2)$ un morphisme dans
$\mathcal{Z}_U$. Par l'hypothèse sur $G$, il existe
un recouvrement $\{U_i \to U\}$ de $U$ et des morphismes
$b_i : F(x_1)|_{U_i} \to F(x_2)|_{U_i}$ dans $\mathcal{Y}_{U_i}$
tels que $G(b_i) = c|_{U_i}$.
Par l'hypothèse sur $F$, il existe des recouvrements $\{U_{ij} \to U_i\}$
et des morphismes $a_{ij} : x_1|_{U_{ij}} \to x_2|_{U_{ij}}$ dans
$\mathcal{X}_{U_{ij}}$ tels que $F(a_{ij}) = b_i|_{U_{ij}}$.
Alors $\{U_{ij} \to U\}$ est un recouvrement de $\mathcal{C}$
et $(G \circ F)(a_{ij}) = c|_{U_{ij}}$, comme l'exige (2)(b).
\end{proof}

\begin{lemma}
\label{stacks-lemma-gerbe-descent}
Soit $\mathcal{C}$ un site. Soit
$$
\xymatrix{
\mathcal{X}' \ar[r]_{G'} \ar[d]_{F'} & \mathcal{X} \ar[d]^F \\
\mathcal{Y}' \ar[r]^G & \mathcal{Y}
}
$$
un diagramme $2$-cartésien de champs en groupoïdes sur $\mathcal{C}$.
Si, pour tout $U \in \Ob(\mathcal{C})$ et tout
$x \in \Ob(\mathcal{Y}_U)$, il existe un recouvrement
$\{U_i \to U\}$ tel que $x|_{U_i}$ appartienne à l'image
essentielle de $G : \mathcal{Y}'_{U_i} \to \mathcal{Y}_{U_i}$ et si
$\mathcal{X}'$ est une gerbe sur $\mathcal{Y}'$, alors
$\mathcal{X}$ est une gerbe sur $\mathcal{Y}$.
\end{lemma}

\begin{proof}
Par la propriété d'unicité d'un $2$-produit fibré, on peut supposer que
$\mathcal{X}' = \mathcal{Y}' \times_\mathcal{Y} \mathcal{X}$
comme dans
Catégories, Lemme \ref{categories-lemma-2-product-categories-over-C}.
Montrons les propriétés (2)(a) et (2)(b) du
Lemme \ref{stacks-lemma-when-gerbe}
pour $\mathcal{X} \to \mathcal{Y}$.

\medskip\noindent
Soit $y$ un objet de $\mathcal{Y}$ au-dessus de l'objet $U$
de $\mathcal{C}$. Par hypothèse, il existe
un recouvrement $\{U_i \to U\}$ de $U$ et des objets
$y'_i \in \mathcal{Y}'_{U_i}$ tels que $G(y'_i) \cong y|_{U_i}$.
D'après (2)(a) pour $\mathcal{X}' \to \mathcal{Y}'$, il existe des
recouvrements $\{U_{ij} \to U_i\}$ et des objets $x'_{ij}$ de
$\mathcal{X}'$ au-dessus de $U_{ij}$ tels que $F'(x'_{ij})$ soit isomorphe
à la restriction de $y'_i$ à $U_{ij}$. Alors
$\{U_{ij} \to U\}$ est un recouvrement de $\mathcal{C}$ et les
$G'(x'_{ij})$ sont des objets de $\mathcal{X}$ au-dessus de $U_{ij}$
dont les images dans $\mathcal{Y}$ sont isomorphes aux restrictions
$y|_{U_{ij}}$. Cela prouve (2)(a) pour $\mathcal{X} \to \mathcal{Y}$.

\medskip\noindent
Soit $U \in \Ob(\mathcal{C})$, soient $x_1, x_2$ des objets de
$\mathcal{X}$ au-dessus de $U$, et soit $b : F(x_1) \to F(x_2)$ un morphisme
dans $\mathcal{Y}_U$. Par hypothèse, on peut choisir un recouvrement
$\{U_i \to U\}$ et des objets $y'_i$ de $\mathcal{Y}'$
au-dessus de $U_i$ tels qu'il existe des isomorphismes
$\alpha_i : G(y'_i) \to F(x_1)|_{U_i}$.
On obtient alors des objets
$$
x'_{1i} = (U_i, y'_i, x_1|_{U_i}, \alpha_i)
\quad\text{et}\quad
x'_{2i} = (U_i, y'_i, x_2|_{U_i}, b|_{U_i} \circ \alpha_i)
$$
de $\mathcal{X}'$ au-dessus de $U_i$. Le morphisme identité de $y'_i$ est un
morphisme $F'(x'_{1i}) \to F'(x'_{2i})$. D'après (2)(b) pour
$\mathcal{X}' \to \mathcal{Y}'$, il existe des recouvrements
$\{U_{ij} \to U_i\}$ et des morphismes
$a'_{ij} : x'_{1i}|_{U_{ij}} \to x'_{2i}|_{U_{ij}}$
tels que $F'(a'_{ij}) = \text{id}_{y'_i}|_{U_{ij}}$. En explicitant la définition
des morphismes de $\mathcal{Y}' \times_\mathcal{Y} \mathcal{X}$,
on voit que les $G'(a'_{ij}) : x_1|_{U_{ij}} \to x_2|_{U_{ij}}$ sont
les morphismes recherchés, c'est-à-dire que (2)(b) est satisfaite pour
$\mathcal{X} \to \mathcal{Y}$.
\end{proof}

\noindent
Les gerbes dont tous les faisceaux d'automorphismes sont abéliens jouent un rôle
important en géométrie algébrique.

\begin{lemma}
\label{stacks-lemma-gerbe-abelian-auts}
Soit $p : \mathcal{S} \to \mathcal{C}$ une gerbe sur un site $\mathcal{C}$.
Supposons que, pour tous $U \in \Ob(\mathcal{C})$ et $x \in \Ob(\mathcal{S}_U)$,
le faisceau de groupes $\mathit{Aut}(x) = \mathit{Isom}(x, x)$ sur $\mathcal{C}/U$
soit abélien. Alors il existe
\begin{enumerate}
\item un faisceau $\mathcal{G}$ de groupes abéliens sur $\mathcal{C}$,
\item pour tout $U \in \Ob(\mathcal{C})$ et tout $x \in \Ob(\mathcal{S}_U)$,
un isomorphisme $\mathcal{G}|_U \to \mathit{Aut}(x)$
\end{enumerate}
tel que, pour tout $U$ et tout morphisme $\varphi : x \to y$
dans $\mathcal{S}_U$, le diagramme
$$
\xymatrix{
\mathcal{G}|_U \ar[d] \ar@{=}[rr] & & \mathcal{G}|_U \ar[d] \\
\mathit{Aut}(x)
\ar[rr]^{\alpha \mapsto \varphi \circ \alpha \circ \varphi^{-1}} & &
\mathit{Aut}(y)
}
$$
soit commutatif.
\end{lemma}

\begin{proof}
Soient $x, y$ deux objets de $\mathcal{S}$ avec $U = p(x) = p(y)$.

\medskip\noindent
S'il existe un morphisme $\varphi : x \to y$ au-dessus de $U$, c'est un
isomorphisme et l'on obtient bien un isomorphisme
$\mathit{Aut}(x) \to \mathit{Aut}(y)$ qui envoie
$\alpha$ sur $\varphi \circ \alpha \circ \varphi^{-1}$.
De plus, comme on suppose $\mathit{Aut}(x)$
commutatif, cet isomorphisme ne dépend pas du choix
de $\varphi$, par un calcul immédiat : en effet, si $\psi$ est
une seconde application de ce type, alors
$$
\varphi \circ \alpha \circ \varphi^{-1} =
\psi \circ \psi^{-1} \circ \varphi \circ \alpha \circ \varphi^{-1} =
\psi \circ \alpha \circ \psi^{-1} \circ \varphi \circ \varphi^{-1} =
\psi \circ \alpha \circ \psi^{-1}
$$
Il en résulte un isomorphisme canonique de faisceaux
$\mathit{Aut}(x) \to \mathit{Aut}(y)$. En outre, s'il existe
un troisième objet $z$ et un morphisme $y \to z$ (et donc
aussi un morphisme $x \to z$), alors les
isomorphismes canoniques $\mathit{Aut}(x) \to \mathit{Aut}(y)$,
$\mathit{Aut}(y) \to \mathit{Aut}(z)$ et
$\mathit{Aut}(x) \to \mathit{Aut}(z)$ sont compatibles au sens
où le diagramme
$$
\xymatrix{
\mathit{Aut}(x) \ar[rd] \ar[rr] & &  \mathit{Aut}(z) \\
& \mathit{Aut}(y) \ar[ru]
}
$$
est commutatif.

\medskip\noindent
S'il n'existe aucun morphisme de $x$ vers $y$ au-dessus de $U$, on peut
choisir un recouvrement $\{U_i \to U\}$ tel qu'il existe des morphismes
$x|_{U_i} \to y|_{U_i}$. On obtient ainsi des isomorphismes canoniques
$$
\mathit{Aut}(x)|_{U_i} \longrightarrow \mathit{Aut}(y)|_{U_i}
$$
qui coïncident sur $U_i \times_U U_j$ (par canonicité). Par recollement des
faisceaux (Sites, Lemme \ref{sites-lemma-glue-maps}), on obtient un unique
isomorphisme $\mathit{Aut}(x) \to \mathit{Aut}(y)$ dont la restriction
à chaque $U_i$ est l'isomorphisme canonique du paragraphe précédent.
Comme ci-dessus, ces isomorphismes canoniques satisfont une
compatibilité en présence d'un troisième objet au-dessus de $U$.

\medskip\noindent
Que se passe-t-il si la catégorie fibre de $\mathcal{S}$ au-dessus de $U$ est vide ?
Dans ce cas, on peut trouver un recouvrement $\{U_i \to U\}$
et des objets $x_i$ de $\mathcal{S}$ au-dessus de $U_i$. Posons alors
$\mathcal{G}_i = \mathit{Aut}(x_i)$. D'après ce qui précède, on obtient
des isomorphismes canoniques
$$
\varphi_{ij} :
\mathcal{G}_i|_{U_i \times_U U_j}
\longrightarrow
\mathcal{G}_j|_{U_i \times_U U_j}
$$
dont les restrictions à $U_i \times_U U_j \times_U U_k$ satisfont
la condition de cocycle décrite dans Sites, section
\ref{sites-section-glueing-sheaves}.
D'après Sites, Lemme \ref{sites-lemma-glue-sheaves},
on obtient un faisceau $\mathcal{G}$ sur $U$ dont la
restriction à $U_i$ donne $\mathcal{G}_i$
d'une manière compatible avec les applications de recollement $\varphi_{ij}$.

\medskip\noindent
Si $\mathcal{C}$ possède un objet final $U$, cela achève la démonstration,
car on peut prendre pour $\mathcal{G}$ le faisceau que l'on vient de construire.
Dans le cas général, il faut vérifier que les faisceaux $\mathcal{G}$
construits sur les différents $U$ sont canoniquement compatibles.
Cette vérification est omise.
\end{proof}




\section{Fonctorialité des champs}
\label{stacks-section-inverse-image}

\noindent
Dans cette section, nous étudions ce qui se passe lorsque l'on veut changer
le site de base d'un champ. Cette section peut être omise en première lecture.

\medskip\noindent
Soit $u : \mathcal{C} \to \mathcal{D}$ un foncteur entre catégories.
Soit $p : \mathcal{S} \to \mathcal{D}$ une catégorie au-dessus de $\mathcal{D}$.
Dans cette situation, nous notons $u^p\mathcal{S}$ la catégorie au-dessus de $\mathcal{C}$
définie comme suit :
\begin{enumerate}
\item Un objet de $u^p\mathcal{S}$ est un couple $(U, y)$ formé
d'un objet $U$ de $\mathcal{C}$ et d'un objet $y$ de $\mathcal{S}_{u(U)}$.
\item Un morphisme $(a, \beta) : (U, y) \to (U', y')$ est donné par
un morphisme $a : U \to U'$ de $\mathcal{C}$ et un morphisme $\beta : y \to y'$
de $\mathcal{S}$ tels que $p(\beta) = u(a)$.
\end{enumerate}
Remarquons qu'avec ces définitions, la catégorie fibre de
$u^p\mathcal{S}$ au-dessus de $U$ est égale à la catégorie fibre de
$\mathcal{S}$ au-dessus de $u(U)$.

\begin{lemma}
\label{stacks-lemma-fibred-category-pushforward}
Dans la situation ci-dessus, si $\mathcal{S}$ est une catégorie fibrée au-dessus de
$\mathcal{D}$, alors $u^p\mathcal{S}$ est une catégorie fibrée au-dessus de $\mathcal{C}$.
\end{lemma}

\begin{proof}
Avant de lire cette démonstration, prière de se reporter à la discussion qui entoure
Catégories, Définitions \ref{categories-definition-cartesian-over-C} et
\ref{categories-definition-fibred-category}.
Soit $(a, \beta) : (U, y) \to (U', y')$ un morphisme de $u^p\mathcal{S}$.
Nous affirmons que $(a, \beta)$ est fortement cartésien si et seulement si
$\beta$ est fortement cartésien. Supposons d'abord que $\beta$ soit fortement cartésien.
Considérons un autre morphisme quelconque
$(a_1, \beta_1) : (U_1, y_1) \to (U', y')$ de $u^p\mathcal{S}$. Alors
\begin{align*}
& \Mor_{u^p\mathcal{S}}((U_1, y_1), (U, y)) \\
& =
\Mor_\mathcal{C}(U_1, U)
\times_{\Mor_\mathcal{D}(u(U_1), u(U))}
\Mor_\mathcal{S}(y_1, y) \\
& =
\Mor_\mathcal{C}(U_1, U)
\times_{\Mor_\mathcal{D}(u(U_1), u(U))}
\Mor_\mathcal{S}(y_1, y')
\times_{\Mor_\mathcal{D}(u(U_1), u(U'))}
\Mor_\mathcal{D}(u(U_1), u(U)) \\
& =
\Mor_\mathcal{S}(y_1, y')
\times_{\Mor_\mathcal{D}(u(U_1), u(U'))}
\Mor_\mathcal{C}(U_1, U) \\
& =
\Mor_{u^p\mathcal{S}}((U_1, y_1), (U', y'))
\times_{\Mor_\mathcal{C}(U_1, U')}
\Mor_\mathcal{C}(U_1, U)
\end{align*}
la deuxième égalité résultant de ce que $\beta$ est fortement cartésien. Nous voyons donc
que $(a, \beta)$ est bien fortement cartésien. Réciproquement, supposons que
$(a, \beta)$ soit fortement cartésien. Choisissons un morphisme fortement cartésien
$\beta' : y'' \to y'$ dans $\mathcal{S}$ tel que $p(\beta') = u(a)$.
Alors les deux morphismes $(a, \beta) : (U, y) \to (U, y')$ et
$(a, \beta') : (U, y'') \to (U, y)$ sont fortement cartésiens et
relèvent $a$. Par l'unicité des morphismes fortement cartésiens
(voir la discussion dans
Catégories, section \ref{categories-section-fibred-categories}),
il existe donc un isomorphisme $\iota : y \to y''$ dans $\mathcal{S}_{u(U)}$
tel que $\beta = \beta' \circ \iota$, ce qui implique que
$\beta$ est fortement cartésien dans $\mathcal{S}$ d'après
Catégories, Lemme \ref{categories-lemma-composition-cartesian}.

\medskip\noindent
Enfin, il nous faut montrer
qu'étant donnés $(U', y')$ et $U \to U'$, on peut trouver un morphisme
fortement cartésien $(U, y) \to (U', y')$ dans $u^p\mathcal{S}$
relevant le morphisme $U \to U'$. Cela résulte de ce qui précède puisque, par hypothèse,
on peut trouver un morphisme fortement cartésien $y \to y'$ relevant le
morphisme $u(U) \to u(U')$.
\end{proof}

\begin{lemma}
\label{stacks-lemma-stack-pushforward}
Soit $u : \mathcal{C} \to \mathcal{D}$ un foncteur continu de sites.
Soit $p : \mathcal{S} \to \mathcal{D}$ un champ sur $\mathcal{D}$.
Alors $u^p\mathcal{S}$ est un champ sur $\mathcal{C}$.
\end{lemma}

\begin{proof}
Nous avons vu dans le
Lemme \ref{stacks-lemma-fibred-category-pushforward}
que $u^p\mathcal{S}$ est une catégorie fibrée au-dessus de $\mathcal{C}$.
De plus, dans la démonstration de ce lemme, nous avons vu qu'un morphisme
$(a, \beta)$ de $u^p\mathcal{S}$ est fortement cartésien si et seulement si
$\beta$ est fortement cartésien dans $\mathcal{S}$. Ainsi, étant donné
un morphisme $a : U \to U'$ de $\mathcal{C}$, non seulement
avons-nous les égalités
$(u^p\mathcal{S})_U = \mathcal{S}_U$ et
$(u^p\mathcal{S})_{U'} = \mathcal{S}_{U'}$, mais, par ces
égalités, les foncteurs de changement de base coïncident ; en formule,
$a^*(U', y') = (U, u(a)^*y')$.

\medskip\noindent
Ceci étant dit, soit $\mathcal{U} = \{U_i \to U\}$ un recouvrement
de $\mathcal{C}$. Comme $u$ est continu, nous voyons que
$\mathcal{V} = \{u(U_i) \to u(U)\}$ est un recouvrement de $\mathcal{D}$,
que $u(U_i \times_U U_j) = u(U_i) \times_{u(U)} u(U_j)$, et de même
pour les produits fibrés triples $U_i \times_U U_j \times_U U_k$.
Compte tenu des identifications des catégories fibres et des changements de base,
les données de descente relatives à $\mathcal{U}$ sont identiques
aux données de descente relatives à $\mathcal{V}$. Puisque, par hypothèse, la descente
est effective dans $\mathcal{S}$, nous concluons qu'il en est de même dans
$u^p\mathcal{S}$.
\end{proof}

\begin{lemma}
\label{stacks-lemma-stack-in-groupoids-pushforward}
Soit $u : \mathcal{C} \to \mathcal{D}$ un foncteur continu de sites.
Soit $p : \mathcal{S} \to \mathcal{D}$ un champ en groupoïdes
sur $\mathcal{D}$. Alors $u^p\mathcal{S}$ est un champ en groupoïdes
sur $\mathcal{C}$.
\end{lemma}

\begin{proof}
Cela résulte immédiatement du
Lemme \ref{stacks-lemma-stack-pushforward}
et du fait que toutes les catégories fibres sont des groupoïdes.
\end{proof}

\begin{definition}
\label{stacks-definition-pushforward-stack}
Soit $f : \mathcal{D} \to \mathcal{C}$ un morphisme de sites
donné par le foncteur continu $u : \mathcal{C} \to \mathcal{D}$.
Soit $\mathcal{S}$ une catégorie fibrée au-dessus de $\mathcal{D}$.
Dans ce cadre, nous notons {\it $f_*\mathcal{S}$} la catégorie
fibrée $u^p\mathcal{S}$ définie ci-dessus. Nous disons que
$f_*\mathcal{S}$ est l'{\it image directe de $\mathcal{S}$ par $f$}.
\end{definition}

\noindent
D'après les résultats précédents, nous savons que $f_*\mathcal{S}$ est un champ (en groupoïdes)
si $\mathcal{S}$ est un champ (en groupoïdes).
Il est plus difficile de définir l'image inverse d'un champ (et notre construction
particulière nécessitera des hypothèses supplémentaires -- le lecteur peut rédiger et
soumettre une construction plus générale). Nous procédons en plusieurs étapes.

\medskip\noindent
Soit $u : \mathcal{C} \to \mathcal{D}$ un foncteur entre catégories.
Soit $p : \mathcal{S} \to \mathcal{C}$ une catégorie au-dessus de $\mathcal{C}$.
Dans ce cadre, nous définissons une catégorie $u_{pp}\mathcal{S}$ comme suit :
\begin{enumerate}
\item Un objet de $u_{pp}\mathcal{S}$ est un triplet
$(U, \phi : V \to u(U), x)$ où $U \in \Ob(\mathcal{C})$,
l'application $\phi : V \to u(U)$ est un morphisme de $\mathcal{D}$,
et $x \in \Ob(\mathcal{S}_U)$.
\item Un morphisme
$$
(U_1, \phi_1 : V_1 \to u(U_1), x_1)
\longrightarrow
(U_2, \phi_2 : V_2 \to u(U_2), x_2)
$$
de $u_{pp}\mathcal{S}$ est donné par un
triplet $(a, b, \alpha)$ où $a : U_1 \to U_2$ est un morphisme de $\mathcal{C}$,
$b : V_1 \to V_2$ est un morphisme de $\mathcal{D}$, et
$\alpha : x_1 \to x_2$ est un morphisme de $\mathcal{S}$,
tels que $p(\alpha) = a$ et que le diagramme
$$
\xymatrix{
V_1 \ar[d]_{\phi_1} \ar[r]_b & V_2 \ar[d]^{\phi_2} \\
u(U_1) \ar[r]^{u(a)} & u(U_2)
}
$$
soit commutatif dans $\mathcal{D}$.
\end{enumerate}
Nous considérons $u_{pp}\mathcal{S}$ comme une catégorie au-dessus de $\mathcal{D}$ au moyen de
$$
p_{pp} : u_{pp}\mathcal{S} \longrightarrow \mathcal{D},
\quad
(U, \phi : V \to u(U), x) \longmapsto V.
$$
La catégorie fibre de $u_{pp}\mathcal{S}$ au-dessus d'un objet $V$ de $\mathcal{D}$
n'admet pas de description simple.

\begin{lemma}
\label{stacks-lemma-right-multiplicative-system}
Dans la situation ci-dessus, supposons que
\begin{enumerate}
\item $p : \mathcal{S} \to \mathcal{C}$ soit une catégorie fibrée,
\item $\mathcal{C}$ possède les limites finies non vides, et que
\item $u : \mathcal{C} \to \mathcal{D}$ commute aux limites finies non vides.
\end{enumerate}
Considérons l'ensemble $R \subset \text{Flèches}(u_{pp}\mathcal{S})$ des morphismes
de la forme
$$
(a, \text{id}_V, \alpha) :
(U', \phi' : V \to u(U'), x')
\longrightarrow
(U, \phi : V \to u(U), x)
$$
où $\alpha$ est fortement cartésien. Alors $R$ est un système multiplicatif à droite.
\end{lemma}

\begin{proof}
D'après
Catégories, Définition \ref{categories-definition-multiplicative-system},
nous devons vérifier RMS1, RMS2 et RMS3.
La condition RMS1 est satisfaite, car une composée de morphismes fortement cartésiens est
fortement cartésienne ; voir
Catégories, Lemme \ref{categories-lemma-composition-cartesian}.

\medskip\noindent
Pour vérifier RMS2, supposons donnés un morphisme
$$
(a, b, \alpha) :
(U_1, \phi_1 : V_1 \to u(U_1), x_1)
\longrightarrow
(U, \phi : V \to u(U), x)
$$
de $u_{pp}\mathcal{S}$ et un morphisme
$$
(c, \text{id}_V, \gamma) :
(U', \phi' : V \to u(U'), x')
\longrightarrow
(U, \phi : V \to u(U), x)
$$
où $\gamma$ est fortement cartésien, appartenant à $R$. Dans cette situation, posons
$U'_1 = U_1 \times_U U'$, et notons $a' : U'_1 \to U'$ et
$c' : U'_1 \to U_1$ les projections.
Puisque $u(U'_1) = u(U_1) \times_{u(U)} u(U')$,
$\phi'_1 = (\phi_1, \phi') : V_1 \to u(U'_1)$ est
un morphisme de $\mathcal{D}$. Soit $\gamma_1 : x_1' \to x_1$
un morphisme fortement cartésien de $\mathcal{S}$ tel que
$p(\gamma_1) = \phi'_1$ (il existe puisque $\mathcal{S}$ est une
catégorie fibrée au-dessus de $\mathcal{C}$). Comme $\gamma : x' \to x$ est
fortement cartésien, il existe alors un unique morphisme
$\alpha' : x'_1 \to x'$ tel que $p(\alpha') = a'$.
À ce stade, nous voyons que
$$
(a', b, \alpha') :
(U_1, \phi_1 : V_1 \to u(U'_1), x'_1)
\longrightarrow
(U, \phi : V \to u(U'), x')
$$
est un morphisme et que
$$
(c', \text{id}_{V_1}, \gamma_1) :
(U'_1, \phi'_1 : V_1 \to u(U'_1), x'_1)
\longrightarrow
(U_1, \phi : V_1 \to u(U_1), x_1)
$$
est un élément de $R$ ; ensemble, ces morphismes donnent une solution au problème
d'existence posé par RMS2.

\medskip\noindent
Enfin, supposons que
$$
(a, b, \alpha), (a', b', \alpha') :
(U_1, \phi_1 : V_1 \to u(U_1), x_1)
\longrightarrow
(U, \phi : V \to u(U), x)
$$
soient deux morphismes de $u_{pp}\mathcal{S}$ et supposons que
$$
(c, \text{id}_V, \gamma) :
(U, \phi : V \to u(U), x)
\longrightarrow
(U', \phi : V \to u(U'), x')
$$
soit un élément de $R$ qui égalise les morphismes
$(a, b, \alpha)$ et $(a', b', \alpha')$. Cela implique en particulier
que $b = b'$. Soit $d : U_2 \to U_1$ l'égalisateur de $a, a'$, lequel
existe (voir
Catégories, Lemme \ref{categories-lemma-almost-finite-limits-exist}).
En outre, $u(d) : u(U_2) \to u(U_1)$ est l'égalisateur de $u(a), u(a')$ ;
il existe donc (puisque $b = b'$) un morphisme $\phi_2 : V_1 \to u(U_2)$ tel que
$\phi_1 = u(d) \circ \phi_1$. Soit $\delta : x_2 \to x_1$ un morphisme fortement
cartésien de $\mathcal{S}$ tel que $p(\delta) = u(d)$.
Nous affirmons maintenant que $\alpha \circ \delta = \alpha' \circ \delta$.
C'est vrai parce que
$\gamma$ est fortement cartésien,
$\gamma \circ \alpha \circ \delta = \gamma \circ \alpha' \circ \delta$, et
$p(\alpha \circ \delta) = p(\alpha' \circ \delta)$.
Par conséquent, la flèche
$$
(d, \text{id}_{V_1}, \delta) :
(U_2, \phi_2 : V_1 \to u(U_2), x_2)
\longrightarrow
(U_1, \phi_1 : V_1 \to u(U_1), x_1)
$$
est un élément de $R$ et égalise $(a, b, \alpha)$ et $(a', b', \alpha')$.
Ainsi, $R$ satisfait également RMS3.
\end{proof}

\begin{lemma}
\label{stacks-lemma-fibred-category-pullback}
Avec les notations et les hypothèses du
Lemme \ref{stacks-lemma-right-multiplicative-system},
posons $u_p\mathcal{S} = R^{-1}u_{pp}\mathcal{S}$ ; voir
Catégories, section \ref{categories-section-localization}.
Alors $u_p\mathcal{S}$ est une catégorie fibrée au-dessus de $\mathcal{D}$.
\end{lemma}

\begin{proof}
Nous utilisons la description de $u_p\mathcal{S}$ donnée juste avant
Catégories, Lemme \ref{categories-lemma-right-localization}.
Remarquons que le foncteur $p_{pp} : u_{pp}\mathcal{S} \to \mathcal{D}$
transforme tout élément de $R$ en un morphisme identité.
Ainsi, d'après
Catégories, Lemme \ref{categories-lemma-properties-right-localization},
nous obtenons un foncteur canonique $p_p : u_p\mathcal{S} \to \mathcal{D}$
qui prolonge le foncteur donné. C'est ainsi que nous considérons
$u_p\mathcal{S}$ comme une catégorie au-dessus de $\mathcal{D}$.

\medskip\noindent
Nous voulons d'abord caractériser les morphismes fortement cartésiens au-dessus de $\mathcal{D}$
dans $u_p\mathcal{S}$.
Un morphisme $f : X \to Y$ de $u_p\mathcal{S}$ est la classe d'équivalence
d'un couple $(f' : X' \to Y, r : X' \to X)$ avec $r \in R$.
En effet, dans $u_p\mathcal{S}$, nous avons $f = (f', 1) \circ (r, 1)^{-1}$
avec les notations évidentes.
Remarquons qu'un isomorphisme est toujours fortement cartésien, de même qu'une
composée de morphismes fortement cartésiens ; voir
Catégories, Lemme \ref{categories-lemma-composition-cartesian}.
Ainsi, $f$ est fortement cartésien si et seulement si $(f', 1)$ l'est.
L'affirmation suivante caractérise donc complètement les morphismes fortement cartésiens.
Affirmation : un morphisme
$$
(a, b, \alpha) :
X_1 = (U_1, \phi_1 : V_1 \to u(U_1), x_1)
\longrightarrow
(U_2, \phi_2 : V_2 \to u(U_2), x_2) = X_2
$$
de $u_{pp}\mathcal{S}$ a une image $f = ((a, b, \alpha), 1)$
fortement cartésienne dans $u_p\mathcal{S}$ si et seulement si $\alpha$
est un morphisme fortement cartésien de $\mathcal{S}$.

\medskip\noindent
Supposons $\alpha$ fortement cartésien.
Soit $X = (U, \phi : V \to u(U), x)$ un autre objet, et soit
$f_2 : X \to X_2$ un morphisme de $u_p\mathcal{S}$
tel que $p_p(f_2) = b \circ b_1$ pour un certain $b_1 : U \to U_1$.
Pour montrer que $f$ est fortement cartésien, nous devons montrer
qu'il existe un unique morphisme $f_1 : X \to X_1$ dans $u_p\mathcal{S}$
tel que $p_p(f_1) = b_1$ et $f_2 = f \circ f_1$ dans $u_p\mathcal{S}$.
Écrivons $f_2 = (f'_2 : X' \to X_2, r : X' \to X)$. Là encore, nous pouvons écrire
$f_2 = (f'_2, 1) \circ (r, 1)^{-1}$ dans $u_p\mathcal{S}$. Puisque $(r, 1)$
est un isomorphisme dont l'image dans $\mathcal{D}$ est une identité, nous
voyons que trouver un morphisme $f_1 : X \to X_1$ possédant les propriétés
requises revient à trouver un morphisme $f'_1 : X' \to X_1$
dans $u_p\mathcal{S}$ tel que $p(f'_1) = b_1$ et $f'_2 = f \circ f'_1$.
Nous pouvons donc supposer que $f_2$ est de la forme
$f_2 = ((a_2, b_2, \alpha_2), 1)$ avec $b_2 = b \circ b_1$. Voici
un diagramme :
$$
\xymatrix{
& & (U_1, V_1 \to u(U_1), x_1) \ar[d]^{(a, b, \alpha)} \\
(U, V \to u(U), x) \ar[rr]^{(a_2, b_2, \alpha_2)} & &
(U_2, V_2 \to u(U_2), x_2)
}
$$
La construction du morphisme $f_1$ est maintenant claire.
Posons en effet $U' = U \times_{U_2} U_1$, avec pour projections
$c : U' \to U$ et $a_1 : U' \to U_1$.
Choisissons un morphisme fortement cartésien $\gamma : x' \to x$
relevant le morphisme $c$. Puisque $b_2 = b \circ b_1$
et que $u(U') = u(U) \times_{u(U_2)} u(U_1)$, nous avons
$\phi' = (\phi, \phi_1 \circ b_1) : V \to u(U')$. Comme
$\alpha$ est fortement cartésien et que
$a \circ a_1 = a_2 \circ c = p(\alpha_2 \circ \gamma)$,
il existe un morphisme
$\alpha_1 : x' \to x_1$ relevant $a_1$ tel que
$\alpha \circ \alpha_1 = \alpha_2 \circ \gamma$.
Posons $X' = (U', \phi' : V \to u(U'), x')$.
Nous voyons ainsi que
$$
f_1 =
((a_1, b_1, \alpha_1) : X' \to X_1, (c, \text{id}_V, \gamma) : X' \to X) :
X \longrightarrow X_1
$$
convient ; en effet, le diagramme
$$
\xymatrix{
(U', \phi' : V \to u(U'), x')
\ar[d]_{(c, \text{id}_V, \gamma)}
\ar[rr]_{(a_1, b_1, \alpha_1)}
& & (U_1, V_1 \to u(U_1), x_1) \ar[d]^{(a, b, \alpha)} \\
(U, V \to u(U), x) \ar[rr]^{(a_2, b_2, \alpha_2)} & &
(U_2, V_2 \to u(U_2), x_2)
}
$$
est commutatif par construction. Cela démontre l'existence.

\medskip\noindent
Démontrons ensuite l'unicité, toujours dans le cas particulier
$f = ((a, b, \alpha), 1)$ et $f_2 = ((a_2, b_2, \alpha_2), 1)$.
Nous recommandons vivement au lecteur de passer cette partie.
Supposons que $g_1, g'_1 : X \to X_1$ soient deux morphismes de
$u_p\mathcal{S}$ tels que $p_p(g_1) = p_p(g'_1) = b_1$ et
$f_2 = f \circ g_1 = f \circ g'_1$. Notre but est de montrer que
$g_1 = g'_1$. D'après
Catégories, Lemme \ref{categories-lemma-morphisms-right-localization},
nous pouvons représenter $g_1$ et $g'_1$ par les classes d'équivalence de
$(f_1 : X' \to X_1, r : X' \to X)$ et
$(f'_1 : X' \to X_1, r : X' \to X)$ pour un certain $r \in R$. D'après
Catégories, Lemme \ref{categories-lemma-equality-morphisms-right-localization},
l'égalité $f_2 = f \circ g_1 = f \circ g'_1$ signifie qu'il
existe un morphisme $r' : X'' \to X'$ dans $u_{pp}\mathcal{S}$ tel que
$r' \circ r \in R$ et
$$
(a, b, \alpha) \circ f_1 \circ r' =
(a, b, \alpha) \circ f'_1 \circ r' =
(a_2, b_2, \alpha_2) \circ r'
$$
dans $u_{pp}\mathcal{S}$. Remarquons que
$g_1$ est maintenant représenté par $(f_1 \circ r', r \circ r')$ et
de même pour $g'_1$. Nous pouvons donc supposer que
$$
(a, b, \alpha) \circ f_1 =
(a, b, \alpha) \circ f'_1 =
(a_2, b_2, \alpha_2).
$$
Écrivons $r = (c, \text{id}_V, \gamma) : (U', \phi' : V \to u(U'), x')$,
$f_1 = (a_1, b_1, \alpha_1)$ et $f'_1 = (a'_1, b_1, \alpha'_1)$.
Nous avons utilisé ici la condition $p_p(g_1) = p_p(g'_1)$.
Les égalités précédentes équivalent maintenant à
$a \circ a_1 = a \circ a'_1 = a_2 \circ c$ et
$\alpha \circ \alpha_1 = \alpha \circ \alpha'_1 = \alpha_2 \circ \gamma$.
Dans cette situation, il n'est pas nécessaire que $a_1 = a'_1$.
Nous devons donc précomposer par un morphisme supplémentaire appartenant à $R$.
Soit en effet $U'' = \text{Eq}(a_1, a'_1)$ l'égalisateur de
$a_1$ et $a'_1$, qui est un sous-objet de $U'$. Notons
$c' : U'' \to U'$ le monomorphisme canonique.
En raison des relations entre les morphismes
ci-dessus, nous voyons que $V \to u(U')$ se factorise par
$u(U'') = u(\text{Eq}(a_1, a'_1)) = \text{Eq}(u(a_1), u(a'_1))$.
Nous obtenons donc un nouvel objet $(U'', \phi'' : V \to u(U''), x'')$, où
$\gamma' : x'' \to x'$ est un morphisme fortement cartésien relevant $\gamma$.
Nous voyons alors que nous pouvons précomposer $f_1$ et $f'_1$ par l'élément
$(c', \text{id}_V, \gamma')$ de $R$. Après cette opération, c'est-à-dire après avoir remplacé
$(U', \phi' : V \to u(U'), x')$ par
$(U'', \phi'' : V \to u(U''), x'')$, nous nous retrouvons dans la situation
précédente, avec en outre $a_1 = a'_1$.
Dans ce cas, il résulte formellement du fait que $\alpha$
est fortement cartésien (!) que $\alpha_1 = \alpha'_1$.
Cela montre que $g_1 = g'_1$, comme souhaité.

\medskip\noindent
Nous omettons la démonstration du fait que, pour tout morphisme fortement cartésien
de $u_p\mathcal{S}$ de la forme $((a, b, \alpha), 1)$, le morphisme
$\alpha$ est fortement cartésien dans $\mathcal{S}$.
(Nous n'avons pas besoin de cette caractérisation des morphismes fortement cartésiens
dans la suite de la démonstration, bien que nous l'utilisions plus loin dans cette section.)

\medskip\noindent
Soit $(U, \phi : V \to u(U), x)$ un objet de $u_p\mathcal{S}$.
Soit $b : V' \to V$ un morphisme de $\mathcal{D}$. Alors le morphisme
$$
(\text{id}_U, b, \text{id}_x) :
(U, \phi \circ b : V' \to u(U), x)
\longrightarrow
(U, \phi : V \to u(U), x)
$$
est fortement cartésien d'après le résultat des paragraphes précédents, ce qui conclut la démonstration.
\end{proof}

\begin{lemma}
\label{stacks-lemma-fibred-groupoids-category-pullback}
Avec les notations et les hypothèses du
Lemme \ref{stacks-lemma-fibred-category-pullback},
si $\mathcal{S}$ est fibrée en groupoïdes, alors $u_p\mathcal{S}$ est fibrée
en groupoïdes.
\end{lemma}

\begin{proof}
D'après le
Lemme \ref{stacks-lemma-fibred-category-pullback},
nous savons que $u_p\mathcal{S}$ est une catégorie fibrée.
Soit $f : X \to Y$ un morphisme de $u_p\mathcal{S}$ tel que
$p_p(f) = \text{id}_V$. Il suffit de montrer que $f$ est inversible ; voir
Catégories, Lemme \ref{categories-lemma-fibred-groupoids}.
Écrivons $f$ comme la classe d'équivalence
d'un couple $((a, b, \alpha), r)$ avec $r \in R$. Alors $p_p(r) = \text{id}_V$,
donc $p_{pp}((a, b, \alpha)) = \text{id}_V$. Par conséquent, $b = \text{id}_V$.
Or tout morphisme de $\mathcal{S}$ est fortement cartésien ; voir
Catégories, Lemme \ref{categories-lemma-fibred-groupoids}.
Ainsi, $(a, b, \alpha) \in R$ est inversible
dans $u_p\mathcal{S}$, comme souhaité.
\end{proof}

\begin{lemma}
\label{stacks-lemma-adjointness-pullback-pushforward}
Soit $u : \mathcal{C} \to \mathcal{D}$ un foncteur.
Soient $p : \mathcal{S} \to \mathcal{C}$ et $q : \mathcal{T} \to \mathcal{D}$
des catégories au-dessus de $\mathcal{C}$ et de $\mathcal{D}$. Supposons que
\begin{enumerate}
\item $p : \mathcal{S} \to \mathcal{C}$ soit une catégorie fibrée,
\item $q : \mathcal{T} \to \mathcal{D}$ soit une catégorie fibrée,
\item $\mathcal{C}$ possède les limites finies non vides, et que
\item $u : \mathcal{C} \to \mathcal{D}$ commute aux limites finies non vides.
\end{enumerate}
Alors nous avons une équivalence canonique de catégories
$$
\Mor_{\textit{Fib}/\mathcal{C}}(\mathcal{S}, u^p\mathcal{T})
=
\Mor_{\textit{Fib}/\mathcal{D}}(u_p\mathcal{S}, \mathcal{T})
$$
entre les catégories de morphismes.
\end{lemma}

\begin{proof}
Dans cette démonstration, nous employons la notation $x/U$ pour désigner un objet
$x$ de $\mathcal{S}$ situé au-dessus de $U$ dans $\mathcal{C}$.
De même, $y/V$ désigne un objet $y$ de $\mathcal{T}$
situé au-dessus de $V$ dans $\mathcal{D}$. Dans le même esprit,
$\alpha/a : x/U \to x'/U'$ désigne le morphisme
$\alpha : x \to x'$ dont l'image est $a : U \to U'$ dans $\mathcal{C}$.

\medskip\noindent
Soit $G : u_p\mathcal{S} \to \mathcal{T}$ un $1$-morphisme de catégories fibrées
au-dessus de $\mathcal{D}$. Notons $G' : u_{pp}\mathcal{S} \to \mathcal{T}$
la composée de $G$ avec le foncteur canonique (de localisation)
$u_{pp}\mathcal{S} \to u_p\mathcal{S}$. Considérons alors le foncteur
$H : \mathcal{S} \to u^p\mathcal{T}$ donné par
$$
H(x/U) = (U, G'(U, \text{id}_{u(U)} : u(U) \to u(U), x))
$$
sur les objets et par
$$
H((\alpha, a) : x/U \to x'/U') = G'(a, u(a), \alpha)
$$
sur les morphismes. Comme $G$ transforme les morphismes fortement cartésiens
en morphismes fortement cartésiens, nous voyons que, si $\alpha$ est fortement
cartésien, alors $H(\alpha)$ est fortement cartésien.
En effet, nous avons vu dans la démonstration du
Lemme \ref{stacks-lemma-fibred-category-pullback}
que, dans ce cas, l'application $(a, u(a), \alpha)$ devient
fortement cartésienne dans $u_p\mathcal{S}$. Cette construction est manifestement
fonctorielle en $G$, et nous obtenons un foncteur
$$
A :
\Mor_{\textit{Fib}/\mathcal{D}}(u_p\mathcal{S}, \mathcal{T})
\longrightarrow
\Mor_{\textit{Fib}/\mathcal{C}}(\mathcal{S}, u^p\mathcal{T})
$$

\medskip\noindent
Réciproquement, soit $H : \mathcal{S} \to u^p\mathcal{T}$ un $1$-morphisme
de catégories fibrées. Rappelons qu'un objet de $u^p\mathcal{T}$ est
un couple $(U, y)$ avec $y \in \Ob(\mathcal{T}_{u(U)})$. Notons
$\text{pr} : u^p\mathcal{T} \to \mathcal{T}$ le foncteur $(U, y) \mapsto y$.
Dans ce cas, nous définissons un foncteur
$G' : u_{pp}\mathcal{S} \to \mathcal{T}$ par les règles
$$
G'(U, \phi : V \to u(U), x) = \phi^*\text{pr}(H(x))
$$
sur les objets et, sur les morphismes, par
$$
G'((a, b, \alpha) : (U, \phi : V \to u(U), x)
\to (U', \phi' : V' \to u(U'), x'))
=
\beta
$$
où $\beta : \phi^*\text{pr}(H(x)) \to (\phi')^*\text{pr}(H(x'))$
est l'unique morphisme
tel que $q(\beta) = b$ et que le diagramme
$$
\xymatrix{
\phi^*\text{pr}(H(x)) \ar[d] \ar[r]_-{\beta} &
(\phi')^*\text{pr}(H(x')) \ar[d] \\
\text{pr}(H(x)) \ar[r]^{\text{pr}(H(a, \alpha))} & \text{pr}(H(x'))
}
$$
soit commutatif. Un tel morphisme existe et est unique parce que $\mathcal{T}$ est une catégorie
fibrée.

\medskip\noindent
Vérifions que $G'(r)$ est un isomorphisme si $r \in R$.
En effet, si
$$
(a, \text{id}_V, \alpha) :
(U', \phi' : V \to u(U'), x')
\longrightarrow
(U, \phi : V \to u(U), x)
$$
avec $\alpha$ fortement cartésien est un élément du système multiplicatif à
droite $R$ du
Lemme \ref{stacks-lemma-right-multiplicative-system},
alors $H(\alpha)$ est fortement cartésien, et
$\text{pr}(H(\alpha))$ est fortement cartésien ; voir la démonstration du
Lemme \ref{stacks-lemma-fibred-category-pushforward}.
Dans ce cas, le morphisme $\beta$ vérifie donc $q(\beta) = \text{id}_V$
et est fortement cartésien. Ainsi, $\beta$ est un isomorphisme d'après
Catégories, Lemme \ref{categories-lemma-composition-cartesian}.
Par conséquent, d'après
Catégories, Lemme \ref{categories-lemma-properties-right-localization},
nous obtenons un prolongement canonique $G : u_p\mathcal{S} \to \mathcal{T}$.

\medskip\noindent
Montrons ensuite que $G$ transforme les morphismes fortement cartésiens
en morphismes fortement cartésiens.
Supposons que $f : X \to Y$ soit fortement cartésien. D'après la caractérisation
des morphismes fortement cartésiens dans $u_p\mathcal{S}$, nous pouvons écrire $f$ sous la forme
$((a, b, \alpha) : X' \to Y, r : X' \to Y)$, où $r \in R$ et où
$\alpha$ est fortement cartésien dans $\mathcal{S}$. D'après ce qui précède, il suffit
de montrer que $G'(a, b \alpha)$ est fortement cartésien. Comme précédemment,
la condition que $\alpha$ soit fortement cartésien implique que
$\text{pr}(H(a, \alpha)) : \text{pr}(H(x)) \to \text{pr}(H(x'))$
est fortement cartésien dans $\mathcal{T}$. Puisque, dans le carré commutatif
ci-dessus, toutes les flèches sauf peut-être $\beta$ sont maintenant fortement cartésiennes,
il s'ensuit que $\beta$ est lui aussi fortement cartésien, comme souhaité.
La construction $H \mapsto G$ est manifestement fonctorielle en $H$, et nous
obtenons un foncteur
$$
B :
\Mor_{\textit{Fib}/\mathcal{C}}(\mathcal{S}, u^p\mathcal{T})
\longrightarrow
\Mor_{\textit{Fib}/\mathcal{D}}(u_p\mathcal{S}, \mathcal{T})
$$
Pour achever la démonstration du lemme, il faut montrer que les foncteurs
$A$ et $B$ sont quasi-inverses l'un de l'autre. Nous omettons les vérifications.
\end{proof}

\begin{definition}
\label{stacks-definition-pullback-stack}
Soit $f : \mathcal{D} \to \mathcal{C}$ un morphisme de sites
donné par un foncteur continu $u : \mathcal{C} \to \mathcal{D}$
qui satisfait les hypothèses et les conclusions de
Sites, Proposition \ref{sites-proposition-get-morphism}.
Soit $\mathcal{S}$ un champ sur $\mathcal{C}$.
Dans ce cadre, nous notons {\it $f^{-1}\mathcal{S}$} la champification
de la catégorie fibrée $u_p\mathcal{S}$ au-dessus de $\mathcal{D}$ construite
ci-dessus. Nous disons que $f^{-1}\mathcal{S}$ est
l'{\it image inverse de $\mathcal{S}$ par $f$}.
\end{definition}

\noindent
Bien entendu, si $\mathcal{S}$ est un champ en groupoïdes, alors $f^{-1}\mathcal{S}$
est un champ en groupoïdes d'après les
Lemmes \ref{stacks-lemma-stackify-groupoids} et
\ref{stacks-lemma-fibred-groupoids-category-pullback}.

\begin{lemma}
\label{stacks-lemma-adjointness-pullback-pushforward-stacks}
Soit $f : \mathcal{D} \to \mathcal{C}$ un morphisme de sites
donné par un foncteur continu $u : \mathcal{C} \to \mathcal{D}$
qui satisfait les hypothèses et les conclusions de
Sites, Proposition \ref{sites-proposition-get-morphism}.
Soient $p : \mathcal{S} \to \mathcal{C}$ et
$q : \mathcal{T} \to \mathcal{D}$ des champs.
Alors nous avons une équivalence canonique de catégories
$$
\Mor_{\textit{Champs}/\mathcal{C}}(\mathcal{S}, f_*\mathcal{T})
=
\Mor_{\textit{Champs}/\mathcal{D}}(f^{-1}\mathcal{S}, \mathcal{T})
$$
entre les catégories de morphismes.
\end{lemma}

\begin{proof}
Pour $i = 1, 2$, un $i$-morphisme de champs est la même chose qu'un
$i$-morphisme de catégories fibrées ; voir la
Définition \ref{stacks-definition-stacks-over-C}.
D'après le
Lemme \ref{stacks-lemma-adjointness-pullback-pushforward},
nous avons déjà
$$
\Mor_{\textit{Fib}/\mathcal{C}}(\mathcal{S}, u^p\mathcal{T})
=
\Mor_{\textit{Fib}/\mathcal{D}}(u_p\mathcal{S}, \mathcal{T})
$$
Le résultat découle donc du
Lemme \ref{stacks-lemma-stackify-universal-property-more},
puisque $u^p\mathcal{T} = f_*\mathcal{T}$ et que $f^{-1}\mathcal{S}$ est la
champification de $u_p\mathcal{S}$.
\end{proof}

\begin{lemma}
\label{stacks-lemma-technical-up}
Soit $f : \mathcal{D} \to \mathcal{C}$ un morphisme de sites
donné par un foncteur continu $u : \mathcal{C} \to \mathcal{D}$
qui satisfait les hypothèses et les conclusions de
Sites, Proposition \ref{sites-proposition-get-morphism}.
Soit $\mathcal{S} \to \mathcal{C}$ une catégorie fibrée, et
soit $\mathcal{S} \to \mathcal{S}'$ la champification de $\mathcal{S}$.
Alors $f^{-1}\mathcal{S}'$ est la champification de
$u_p\mathcal{S}$.
\end{lemma}

\begin{proof}
Omis. Indication :
C'est l'analogue de Sites, Lemme \ref{sites-lemma-technical-up}.
\end{proof}

\noindent
Le lemme suivant nous dit que la $2$-catégorie des champs
sur $\Sch_{fppf}$ est une ``2-sous-catégorie pleine'' de la
$2$-catégorie des champs sur $\Sch'_{fppf}$ dès que
$\Sch'_{fppf}$ contient $\Sch_{fppf}$ (voir
Topologies, section \ref{topologies-section-change-alpha}).

\begin{lemma}
\label{stacks-lemma-bigger-site}
Soient $\mathcal{C}$ et $\mathcal{D}$ des sites.
Soit $u : \mathcal{C} \to \mathcal{D}$ un foncteur satisfaisant aux
hypothèses de
Sites, Lemme \ref{sites-lemma-bigger-site}.
Soit $f : \mathcal{D} \to \mathcal{C}$ le morphisme de sites
correspondant. Alors
\begin{enumerate}
\item pour tout champ $p : \mathcal{S} \to \mathcal{C}$, le
foncteur canonique $\mathcal{S} \to f_*f^{-1}\mathcal{S}$ est
une équivalence de champs ;
\item étant donnés des champs $\mathcal{S}, \mathcal{S}'$ sur $\mathcal{C}$,
la construction $f^{-1}$
induit une équivalence
$$
\Mor_{\textit{Champs}/\mathcal{C}}(\mathcal{S}, \mathcal{S}')
\longrightarrow
\Mor_{\textit{Champs}/\mathcal{D}}(f^{-1}\mathcal{S}, f^{-1}\mathcal{S}')
$$
entre les catégories de morphismes.
\end{enumerate}
\end{lemma}

\begin{proof}
Remarquons que, d'après le
Lemme \ref{stacks-lemma-adjointness-pullback-pushforward-stacks},
nous avons une équivalence de catégories
$$
\Mor_{\textit{Champs}/\mathcal{D}}(f^{-1}\mathcal{S}, f^{-1}\mathcal{S}')
=
\Mor_{\textit{Champs}/\mathcal{C}}(\mathcal{S}, f_*f^{-1}\mathcal{S}')
$$
Ainsi, (2) découle de (1).

\medskip\noindent
Pour démontrer (1), nous allons utiliser le
Lemme \ref{stacks-lemma-characterize-essentially-surjective-when-ff}.
Ce lemme nous dit qu'il faut montrer que
$can : \mathcal{S} \to f_*f^{-1}\mathcal{S}$ est
pleinement fidèle et que tous les objets de $f_*f^{-1}\mathcal{S}$
sont localement dans l'image essentielle.

\medskip\noindent
Décrivons rapidement le foncteur $can$ ; voir la démonstration du
Lemme \ref{stacks-lemma-adjointness-pullback-pushforward}.
Pour ce faire, introduisons le foncteur
$c'' : \mathcal{S} \to u_{pp}\mathcal{S}$ défini par
$c''(x/U) = (U, \text{id} : u(U) \to u(U), x)$ et
$c''(\alpha/a) = (a, u(a), \alpha)$. Nous définissons
$c' : \mathcal{S} \to u_p\mathcal{S}$ comme la composée
de $c''$ et du foncteur canonique $u_{pp}\mathcal{S} \to u_p\mathcal{S}$.
Nous définissons $c : \mathcal{S} \to f^{-1}\mathcal{S}$ comme la composée
de $c'$ et du foncteur canonique $u_p\mathcal{S} \to f^{-1}\mathcal{S}$.
Alors $can : \mathcal{S} \to f_*f^{-1}\mathcal{S}$ est le foncteur
qui associe à $x/U$ le couple $(U, c(x))$ et à
$\alpha/a$ le morphisme $(a, c(\alpha))$.

\medskip\noindent
Pleine fidélité. Pour la démontrer, nous allons utiliser le
Lemme \ref{stacks-lemma-characterize-ff}.
Soit $U \in \Ob(\mathcal{C})$.
Soient $x, y \in \mathcal{S}_U$.
Tout d'abord, comme $u$ est pleinement fidèle, nous avons
$$
\Mor_{(f_*f^{-1}\mathcal{S})_U}(can(x), can(y))
=
\Mor_{(f^{-1}\mathcal{S})_{u(U)}}(c(x), c(y))
$$
directement par la définition de $f_*$. Il en va de même après
changement de base à tout $U'/U$. Puisque $f^{-1}\mathcal{S}$ est la
champification de $u_p\mathcal{S}$, et puisque $u$ est continu
et cocontinu, le préfaisceau
$$
U'/U \longmapsto
\Mor_{(f^{-1}\mathcal{S})_{u(U')}}(c(x|_{U'}), c(y|_{U'}))
$$
est le faisceau associé au préfaisceau
$$
U'/U \longmapsto
\Mor_{(u_p\mathcal{S})_{u(U')}}(c'(x|_{U'}), c'(y|_{U'}))
$$
Ainsi, pour achever la démonstration de la pleine fidélité, il suffit
de montrer que, quels que soient $U$ et $x, y$, l'application
$$
\Mor_{\mathcal{S}_U}(x, y)
\longrightarrow
\Mor_{(u_p\mathcal{S})_U}(c'(x), c'(y))
$$
est bijective. Un morphisme $f : x \to y$ dans $u_p\mathcal{S}$ au-dessus de $u(U)$
est donné par une classe d'équivalence de diagrammes
$$
\xymatrix{
(U', \phi : u(U) \to u(U'), x')
\ar[d]_{(c, \text{id}_{u(U)}, \gamma)}
\ar[r]_{(a, b, \alpha)} &
(U, \text{id} : u(U) \to u(U), y)
\\
(U, \text{id} : u(U) \to u(U), x)
}
$$
avec $\gamma$ fortement cartésien et $b = \text{id}_{u(U)}$.
Mais, puisque $u$ est pleinement fidèle, nous pouvons écrire $\phi = u(c')$ pour un
morphisme $c' : U \to U'$ ;
nous voyons alors que $a \circ c' = \text{id}_U$ et
$c \circ c' = \text{id}_{U'}$. Comme $\gamma$ est fortement cartésien,
nous pouvons trouver un morphisme $\gamma' : x \to x'$ relevant $c'$ tel que
$\gamma \circ \gamma' = \text{id}_x$. Par définition des classes
d'équivalence qui définissent les morphismes dans $u_p\mathcal{S}$, il s'ensuit que
le morphisme
$$
\xymatrix{
(U, \text{id} : u(U) \to u(U), x)
\ar[rr]_{(\text{id}, \text{id}, \alpha \circ \gamma')} & &
(U, \text{id} : u(U) \to u(U), y)
}
$$
de $u_{pp}\mathcal{S}$ induit le morphisme $f$ dans $u_p\mathcal{S}$.
Cela prouve que l'application est surjective. Nous omettons la démonstration de son
injectivité.

\medskip\noindent
Enfin, nous devons montrer que tout objet de $f_*f^{-1}\mathcal{S}$
provient localement d'un objet de $\mathcal{S}$. Cela ressort des
constructions (détails omis).
\end{proof}




\section{Champs et localisation}
\label{stacks-section-localize}

\noindent
Soit $\mathcal{C}$ un site.
Soit $U$ un objet de $\mathcal{C}$.
Nous voulons interpréter les champs sur $\mathcal{C}/U$ comme des champs
sur $\mathcal{C}$ munis d'un morphisme vers $U$.
Le lemme suivant explique pourquoi cela est plus facile
lorsque le préfaisceau $h_U$ est un faisceau.

\begin{lemma}
\label{stacks-lemma-when-localization-stack}
Soit $\mathcal{C}$ un site. Soit $U \in \Ob(\mathcal{C})$.
Alors $j_U : \mathcal{C}/U \to \mathcal{C}$ est un champ sur $\mathcal{C}$
si et seulement si $h_U$ est un faisceau.
\end{lemma}

\begin{proof}
Combiner le
Lemme \ref{stacks-lemma-stack-in-setoids-characterize}
avec
Catégories,
Exemple \ref{categories-example-fibred-category-from-functor-of-points}.
\end{proof}

\noindent
Supposons que $\mathcal{C}$ soit un site
et que $U$ soit un objet de $\mathcal{C}$ dont le préfaisceau représentable
associé soit un faisceau. Nous notons $j : \mathcal{C}/U \to \mathcal{C}$
le foncteur de localisation.

\medskip\noindent
{\bf Construction A.} Soit $p : \mathcal{S} \to \mathcal{C}/U$ un champ
sur le site $\mathcal{C}/U$. Nous définissons un champ
$j_!p : j_!\mathcal{S} \to \mathcal{C}$ comme suit :
\begin{enumerate}
\item En tant que catégorie, $j_!\mathcal{S} = \mathcal{S}$, et
\item le foncteur $j_!p : j_!\mathcal{S} \to \mathcal{C}$
est simplement la composée $j \circ p$.
\end{enumerate}
Nous omettons la vérification qu'il s'agit d'un champ (indication : utiliser le fait que
$h_U$ est un faisceau pour recoller les morphismes vers $U$). Il existe un foncteur
canonique
$$
j_!\mathcal{S} \longrightarrow \mathcal{C}/U
$$
à savoir le foncteur $p$, qui est un $1$-morphisme de champs sur $\mathcal{C}$.

\medskip\noindent
{\bf Construction B.}
Soit $q : \mathcal{T} \to \mathcal{C}$ un champ sur $\mathcal{C}$
muni d'un morphisme de champs $p : \mathcal{T} \to \mathcal{C}/U$
sur $\mathcal{C}$. Alors $p : \mathcal{T} \to \mathcal{C}/U$
est automatiquement un champ sur $\mathcal{C}/U$.

\begin{lemma}
\label{stacks-lemma-localize-stacks}
Supposons que $\mathcal{C}$ soit un site
et que $U$ soit un objet de $\mathcal{C}$ dont le préfaisceau représentable
associé soit un faisceau. Les constructions A et B ci-dessus définissent
des foncteurs de $2$-catégories mutuellement inverses (!)
$$
\left\{
\begin{matrix}
2\text{-catégorie des}\\
\text{champs sur }\mathcal{C}/U
\end{matrix}
\right\}
\leftrightarrow
\left\{
\begin{matrix}
2\text{-catégorie des couples }(\mathcal{T}, p)
\text{ constitués} \\
\text{d'un champ }\mathcal{T}\text{ sur }\mathcal{C}\text{ et d'un morphisme} \\
p : \mathcal{T} \to \mathcal{C}/U\text{ de champs sur }\mathcal{C}
\end{matrix}
\right\}
$$
\end{lemma}

\begin{proof}
C'est clair.
\end{proof}








% OK, start here.
%

\chapter{Corps}



\phantomsection
\label{fields-section-phantom}



\section{Introduction}
\label{fields-section-introduction}

\noindent
Dans ce chapitre, nous étudierons la théorie des corps. Rappelons qu'un
{\it corps} est un anneau non nul dans lequel tout élément non nul est inversible.
De manière équivalente, les deux seuls idéaux d'un corps sont $(0)$ et $(1)$,
puisque tout élément non nul est une unité. Les corps constitueront donc les
cas les plus simples d'une grande partie de la théorie développée ultérieurement.

\medskip\noindent
La théorie des extensions de corps présente un aspect différent de l'algèbre
commutative classique puisque, par exemple, tout morphisme de corps est injectif.
Il s'avère néanmoins que des questions portant sur les anneaux peuvent souvent
se ramener à des questions sur les corps. Par exemple, tout anneau intègre se plonge dans un corps
(son corps des fractions), et à tout {\it anneau local} (c'est-à-dire un anneau possédant un unique
idéal maximal ; nous n'avons pas encore défini ce terme) est associé son
corps résiduel (c'est-à-dire son quotient par l'idéal maximal).
La connaissance des extensions de corps sera donc utile.




\section{Définitions de base}
\label{fields-section-definitions}

\noindent
Comme nous avons placé ce chapitre avant celui consacré à l'algèbre
commutative, nous devons introduire ici certaines définitions élémentaires
avant de les étudier plus en détail dans les chapitres d'algèbre.

\begin{definition}
\label{fields-definition-field}
Un {\it corps} est un anneau non nul dont tout élément non nul est inversible.
Étant donné un corps, un {\it sous-corps} est un sous-anneau qui est lui-même un corps.
\end{definition}

\noindent
Pour un corps $k$, nous notons $k^*$ le sous-ensemble $k \setminus \{0\}$.
Cette notation généralise la notation usuelle $R^*$ qui désigne le groupe des
éléments inversibles d'un anneau $R$.

\begin{definition}
\label{fields-definition-domain}
Un {\it anneau intègre}, ou un {\it domaine d'intégrité}, est un anneau non nul dans lequel $0$
est le seul diviseur de zéro.
\end{definition}



\section{Exemples de corps}
\label{fields-section-examples}

\noindent
Pour commencer, donnons plusieurs exemples de corps. Le lecteur se rappellera
que, si $R$ est un anneau et $I \subset R$ un idéal, alors $R/I$ est un corps
si et seulement si $I$ est un idéal maximal.

\begin{example}[Nombres rationnels]
\label{fields-example-rational-numbers}
Les nombres rationnels forment un corps. On l'appelle le
{\it corps des nombres rationnels} et on le note $\mathbf{Q}$.
\end{example}

\begin{example}[Corps premiers]
\label{fields-example-prime-field}
Si $p$ est un nombre premier, alors $\mathbf{Z}/(p)$ est un corps, noté
$\mathbf{F}_p$. En effet, $(p)$ est un
idéal maximal de $\mathbf{Z}$. Ainsi, un corps peut être fini : $\mathbf{F}_p$
contient $p$ éléments.
\end{example}

\begin{example}
\label{fields-example-quotient-polymial-ring}
Dans un anneau principal, tout idéal engendré par un élément irréductible
est maximal. Or, si $k$ est un corps, l'anneau de polynômes $k[x]$ est un
anneau principal. Il s'ensuit que, si $P \in k[x]$ est un polynôme irréductible
(c'est-à-dire un polynôme non constant
qui n'admet pas de factorisation en facteurs de degrés plus petits), alors
$k[x]/(P)$ est un corps. Il contient naturellement une copie de $k$.
C'est une méthode très générale pour construire des corps. Par exemple, les
nombres complexes $\mathbf{C}$
peuvent être construits sous la forme $\mathbf{R}[x]/(x^2 + 1)$.
\end{example}

\begin{example}[Corps des fractions]
\label{fields-example-quotient-field}
Rappelons qu'étant donné un anneau intègre $A$, il existe un plongement $A \to F$ dans un
corps $F$ construit à partir de $A$ exactement de la même manière que
$\mathbf{Q}$ est construit à partir de $\mathbf{Z}$. Formellement, les éléments
de $F$ sont des fractions $a/b$ (ou, plus précisément, leurs classes d'équivalence),
$a, b \in A$, $b \not = 0$. Comme d'habitude, $a/b = a'/b'$ si et seulement si $ab' = ba'$.
Le corps $F$ est appelé le {\it corps des fractions} de $A$ ; les calques
{\it corps quotient} et {\it corps de fractions} ne seront pas employés.
Le corps des fractions possède la propriété universelle suivante : pour tout
morphisme injectif d'anneaux $\varphi : A \to K$ vers un corps $K$, il existe un unique
morphisme $\psi : F \to K$ qui rende le diagramme
$$
\xymatrix{
F \ar[r]_\psi & K \\
A \ar[u] \ar[ru]_\varphi
}
$$
commutatif. En effet, la façon de définir un tel morphisme est claire : on pose
$\psi(a/b) = \varphi(a)\varphi(b)^{-1}$, l'injectivité de $\varphi$
assurant que $\varphi(b) \not = 0$ si $ b \not = 0$.
\end{example}

\begin{example}[Corps des fonctions rationnelles]
\label{fields-example-field-of-rational-functions}
Si $k$ est un corps, on peut considérer le corps $k(x)$ des fonctions
rationnelles sur $k$. C'est le corps des fractions de l'anneau de polynômes
$k[x]$. Autrement dit, c'est l'ensemble des quotients $F/G$ avec
$F, G \in k[x]$, $G \not = 0$, muni de la relation d'équivalence évidente.
\end{example}

\begin{example}
\label{fields-example-field-of-meromorphic-functions}
Soit $X$ une surface de Riemann. Notons $\mathbf{C}(X)$
l'ensemble des fonctions méromorphes sur $X$. Alors $\mathbf{C}(X)$ est un anneau pour
la multiplication et l'addition des fonctions. Il s'avère qu'en fait
$\mathbf{C}(X)$ est un corps. En effet, si une fonction non nulle $f(z)$ est
méromorphe, il en va de même de $1/f(z)$. Par exemple, soit $S^2$ la sphère
de Riemann ; l'analyse complexe nous apprend que l'anneau des fonctions méromorphes
$\mathbf{C}(S^2)$ est le corps des fonctions rationnelles $\mathbf{C}(z)$.
\end{example}



\section{Espaces vectoriels}
\label{fields-section-vector-spaces}

\noindent
L'une des raisons pour lesquelles les corps sont si commodes est que la théorie des modules sur un corps
(c'est-à-dire des espaces vectoriels) est très simple.

\begin{lemma}
\label{fields-lemma-vector-space-is-free}
Si $k$ est un corps, tout $k$-module est libre.
\end{lemma}

\begin{proof}
En effet, l'algèbre linéaire nous apprend qu'un $k$-module (c'est-à-dire un espace vectoriel)
$V$ possède une {\it base} $\mathcal{B} \subset V$, laquelle définit un isomorphisme
de l'espace vectoriel libre sur $\mathcal{B}$ vers $V$.
\end{proof}

\begin{lemma}
\label{fields-lemma-field-semi-simple}
Toute suite exacte de modules sur un corps est scindée.
\end{lemma}

\begin{proof}
Cela résulte du Lemme \ref{fields-lemma-vector-space-is-free}, puisque tout espace
vectoriel est un module projectif.
\end{proof}

\noindent
C'est une autre raison pour laquelle une grande partie de la théorie des chapitres ultérieurs dira
peu de choses sur les corps, puisque les modules s'y comportent de manière si simple.
Notons que le Lemme \ref{fields-lemma-field-semi-simple} est un énoncé portant sur la
{\it catégorie} des $k$-modules (où $k$ est un corps), car la notion
d'exactitude est intrinsèquement de nature morphique, c'est-à-dire qu'elle ne fait intervenir que des notions
catégoriques, et peut en fait se formuler dans ce que l'on appelle une {\it catégorie abélienne}.

\medskip\noindent
Désormais, puisque l'étude des modules sur un corps relève de l'algèbre linéaire et
que la théorie des idéaux d'un corps présente peu d'intérêt, nous étudierons le véritable
objet de ce chapitre : les {\it extensions} de corps.


\section{La caractéristique d'un corps}
\label{fields-section-more-fields}

\noindent
Dans la catégorie des anneaux, il existe un {\it objet initial} $\mathbf{Z}$ : tout
anneau $R$ reçoit un morphisme de $\mathbf{Z}$ d'une manière et d'une seule. Pour les corps,
il n'existe pas d'objet initial.
Il existe néanmoins une famille d'objets telle que tout corps reçoive d'une manière
unique un morphisme provenant d'exactement l'un d'entre eux, et aucun morphisme provenant des autres.

\medskip\noindent
Soit $F$ un corps. Considérons $F$ comme un anneau afin d'obtenir un morphisme d'anneaux
$f : \mathbf{Z} \to F$. L'image de ce morphisme d'anneaux est un anneau intègre
(en tant que sous-anneau d'un corps) ; le noyau de $f$ est donc un idéal premier
de $\mathbf{Z}$. Par conséquent, le noyau de $f$ est soit $(0)$, soit $(p)$ pour
un certain nombre premier $p$.

\medskip\noindent
Dans le premier cas, $f$ est injectif, et nous considérons alors
$\mathbf{Z}$ comme un sous-anneau de $F$. De plus, puisque tout
élément non nul de $F$ est inversible, il est légitime de
parler de $p/q \in F$ pour $p, q \in \mathbf{Z}$ avec $q \not = 0$.
Dans ce cas, nous pouvons donc considérer, et nous considérons, $\mathbf{Q}$ comme un sous-anneau de $F$.
On vérifie aisément que c'est alors le plus petit sous-corps de $F$.

\medskip\noindent
Dans le second cas, c'est-à-dire lorsque $\Ker(f) = (p)$, on voit que
$\mathbf{Z}/(p) = \mathbf{F}_p$ est un sous-anneau de $F$. C'est manifestement le
plus petit sous-corps de $F$.

\medskip\noindent
Ce raisonnement montre que tout corps contient un plus petit sous-corps,
qui est soit $\mathbf{Q}$, soit le corps fini $\mathbf{F}_p$ pour un certain
nombre premier $p$.

\begin{definition}
\label{fields-definition-characteristic}
La {\it caractéristique} d'un corps $F$ est $0$ si
$\mathbf{Z} \subset F$, et est un nombre premier $p$ si $p = 0$ dans $F$.
Le {\it sous-corps premier de $F$} est le plus petit sous-corps de $F$ ;
il s'agit soit de $\mathbf{Q} \subset F$ si la caractéristique est nulle, soit de
$\mathbf{F}_p \subset F$ si la caractéristique est $p > 0$.
\end{definition}

\noindent
Il est facile de voir que, si $E \subset F$ est un sous-corps, la
caractéristique de $E$ est la même que celle de $F$.

\begin{example}
\label{fields-example-characteristic}
La caractéristique de $\mathbf{F}_p$ est $p$, et celle de $\mathbf{Q}$ est $0$.
\end{example}


\section{Extensions de corps}
\label{fields-section-extensions}

\noindent
En général, cependant, ce ne sont pas tant les corps eux-mêmes qui nous intéressent que leurs
{\it extensions}. Cette démarche est peut-être analogue à celle qui consiste à étudier non les anneaux,
mais les {\it algèbres} sur un anneau fixé.
Dans le cas des corps, la notion de ``corps sur un autre corps''
redonne simplement celle d'extension de corps, comme le montre le résultat suivant.

\begin{lemma}
\label{fields-lemma-field-maps-injective}
Si $F$ est un corps et $R$ un anneau non nul, alors tout homomorphisme d'anneaux
$\varphi : F \to R$ est injectif.
\end{lemma}

\begin{proof}
En effet, soit $a \in \Ker(\varphi)$ un élément non nul. On a alors
$\varphi(1) = \varphi(a^{-1} a) = \varphi(a^{-1}) \varphi(a) = 0$.
Ainsi $1 = \varphi(1) = 0$, et $R$ est l'anneau nul.
\end{proof}

\begin{definition}
\label{fields-definition-extension}
Si $F$ est un corps contenu dans un corps $E$, on dit que $E$ est
une {\it extension de corps} de $F$. Nous écrirons $E/F$ pour indiquer
que $E$ est une extension de $F$.
\end{definition}

\noindent
Ainsi, si $F, F'$ sont des corps et $F \to F'$ un homomorphisme d'anneaux quelconque, le
Lemme \ref{fields-lemma-field-maps-injective} montre qu'il est injectif, et $F'$ peut être
considéré comme une extension de $F$, par un léger abus de langage. De façon équivalente,
une extension de corps de $F$ n'est autre qu'une $F$-algèbre qui est un corps.
La situation est entièrement différente pour les anneaux généraux, puisqu'un
homomorphisme d'anneaux n'est pas nécessairement injectif.

\medskip\noindent
Soit $k$ un corps. Il existe une {\it catégorie} des extensions de corps de $k$.
Un objet de cette catégorie est une extension $E/k$, c'est-à-dire un
morphisme de corps (nécessairement injectif)
$$
k \to E,
$$
tandis qu'un morphisme entre les extensions $E/k$ et $E'/k$ est un morphisme de
$k$-algèbres $E \to E'$ ; de façon équivalente, c'est un diagramme commutatif
$$
\xymatrix{
E \ar[rr] & & E' \\
& k \ar[ru] \ar[lu] &
}
$$
L'ensemble des morphismes $E \to E'$ dans la catégorie des extensions de $k$
sera noté $\Mor_k(E, E')$.

\begin{definition}
\label{fields-definition-tower}
Une {\it tour} de corps $E_n/E_{n - 1}/\ldots/E_0$ consiste en une suite
d'extensions de corps
$E_n/E_{n - 1}$, $E_{n - 1}/E_{n - 2}$, $\ldots$, $E_1/E_0$.
\end{definition}

\noindent
Donnons quelques exemples d'extensions de corps.

\begin{example}
\label{fields-example-monogenic-extension}
Soient $k$ un corps et $P \in k[x]$ un polynôme irréductible. Nous avons
vu que $k[x]/(P)$ est un corps (Exemple \ref{fields-example-quotient-polymial-ring}).
Comme c'est aussi une $k$-algèbre de façon évidente, il s'agit d'une extension de $k$.
\end{example}

\begin{example}
\label{fields-example-field-of-meromorphic-functions-extension-C}
Si $X$ est une surface de Riemann, le corps des fonctions méromorphes
$\mathbf{C}(X)$ (Exemple \ref{fields-example-field-of-meromorphic-functions})
est une extension de corps de $\mathbf{C}$, car tout élément de $\mathbf{C}$
définit sur $X$ une fonction constante méromorphe --- et même holomorphe.
\end{example}

\noindent
Soit $F/k$ une extension de corps. Soit $S \subset F$ un sous-ensemble quelconque.
Il existe alors une {\it plus petite} sous-extension de $F$ (c'est-à-dire un sous-corps de
$F$ contenant $k$) qui contienne $S$. Pour le voir, considérons la famille des
sous-corps de $F $ contenant $S$ et $k$, puis prenons leur intersection ; on
vérifie qu'il s'agit d'un corps. Un argument classique montre en fait que
c'est l'ensemble des éléments de $F$ que l'on peut obtenir au moyen d'un nombre fini
d'opérations algébriques élémentaires (addition, multiplication, soustraction
et division) portant sur des éléments de $k$ et de $S$.

\begin{definition}
\label{fields-definition-generated-by}
Soit $k$ un corps. Si $F/k$ est une extension de corps et
$S \subset F$, nous notons $k(S)$ le plus petit sous-corps de $F$
contenant $k$ et $S$. Nous dirons que $S$ {\it engendre
l'extension de corps} $k(S)/k$. Si $S = \{\alpha\}$ est un singleton, nous
écrirons $k(\alpha)$ au lieu de $k(\{\alpha\})$. On dit que $F/k$ est une
{\it extension de corps de type fini} s'il existe un
sous-ensemble fini $S \subset F$ tel que $F = k(S)$.
\end{definition}

\noindent
Par exemple, $\mathbf{C}$ est engendré par $i$ sur $\mathbf{R}$.

\begin{exercise}
\label{fields-exercise-C-not-countably-generated}
Montrer que $\mathbf{C}$ ne possède pas de famille dénombrable de générateurs sur
$\mathbf{Q}$.
\end{exercise}

\noindent
Classifions maintenant les extensions engendrées par un élément.

\begin{lemma}[Classification des extensions simples]
\label{fields-lemma-field-extension-generated-by-one-element}
Si une extension de corps $F/k$ est engendrée par un élément, alors elle est
$k$-isomorphe soit au corps des fonctions rationnelles $k(t)/k$, soit à l'une
des extensions $k[t]/(P)$, où $P \in k[t]$ est irréductible.
\end{lemma}

\noindent
Nous verrons que bon nombre des cas les plus importants d'extensions de corps sont
engendrés par un élément ; ce résultat est donc effectivement utile.

\begin{proof}
Soit $\alpha \in F$ tel que $F = k(\alpha)$ ; par hypothèse, un tel
$\alpha$ existe. Il existe un morphisme d'anneaux
$$
k[t] \to F
$$
qui envoie l'indéterminée $t$ sur $\alpha$. Son image est un anneau intègre, donc son
noyau est un idéal premier. Celui-ci est donc soit $(0)$, soit $(P)$ pour un $P \in k[t]$
irréductible.

\medskip\noindent
Si le noyau est $(P)$ pour un $P \in k[t]$ irréductible, le morphisme se factorise
par $k[t]/(P)$ et induit un morphisme de corps $k[t]/(P) \to F$. Puisque
l'image contient $\alpha$, on voit aisément que ce morphisme est surjectif, donc
est un isomorphisme. Dans ce cas, $k[t]/(P) \simeq F$.

\medskip\noindent
Si le noyau est trivial, on obtient une injection $k[t] \to F$.
On peut donc définir un morphisme du corps des fractions $k(t)$ dans $F$ ; à un
quotient $R(t)/Q(t)$ avec $R(t), Q(t) \in k[t]$, on associe
$R(\alpha)/Q(\alpha)$. L'hypothèse selon laquelle $k[t] \to F$ est injectif implique
que $Q(\alpha) \neq 0$, sauf si $Q$ est le polynôme nul.
Le corps des fractions de $k[t]$ est le corps des fonctions rationnelles $k(t)$ ; on obtient donc
un morphisme $k(t) \to F$
dont l'image contient $\alpha$. Il est par conséquent surjectif, donc est un isomorphisme.
\end{proof}

\begin{lemma}
\label{fields-lemma-common-extension-field}
Soient $k$ un corps et $E/k$, $F/k$ des extensions de corps.
Il existe alors une extension de corps commune $M/k$, c'est-à-dire une extension
de corps telle qu'il existe des morphismes $E \to M$ et $F \to M$ d'extensions de $k$.
\end{lemma}

\begin{proof}
Nous ne démontrons ce résultat que lorsque $E$ est une extension de corps de type fini
de $k$ ; le cas général s'en déduit par un argument de type lemme de Zorn
(les détails sont omis).

\medskip\noindent
Supposons d'abord que $E$ soit une extension simple de $k$. D'après le
Lemme \ref{fields-lemma-field-extension-generated-by-one-element},
cela signifie soit que $E = k(t)$ est le corps des fonctions rationnelles,
soit que $E = k[t]/(P)$ pour un certain polynôme irréductible $P \in k[t]$.
Dans le premier cas, prenons pour $M = F(t)$ le corps des fonctions rationnelles,
muni des morphismes évidents $E \to M$ et $F \to M$. Dans le second cas, nous
choisissons un facteur irréductible $Q$ de l'image de $P$ dans $F[t]$
et prenons $M = F[t]/(Q)$, muni des morphismes évidents $E \to M$ et $F \to M$.

\medskip\noindent
Si $E = k(\alpha_1, \ldots, \alpha_n)$, une récurrence sur $n$
permet de trouver une extension $M/k$ et des morphismes $F \to M$ et
$k(\alpha_1, \ldots, \alpha_{n - 1}) \to M$. En appliquant le cas simple
étudié au paragraphe précédent,
on trouve une extension $M'/k(\alpha_1, \ldots, \alpha_{n - 1})$
et des morphismes $M \to M'$ et $k(\alpha_1, \ldots, \alpha_n) \to M'$.
Alors $M'$, considéré comme une extension de $k$, convient.
\end{proof}




\section{Extensions finies}
\label{fields-section-finite-extensions}

\noindent
Si $F/E$ est une extension de corps, alors $F$ est aussi, de manière évidente, un espace vectoriel
sur $E$ (l'action scalaire n'est autre que la multiplication dans $F$).

\begin{definition}
\label{fields-definition-degree}
Soit $F/E$ une extension de corps. La dimension de $F$, considéré comme un
$E$-espace vectoriel, est appelée le {\it degré} de l'extension et est
notée $[F : E]$. Si $[F : E] < \infty$, on dit que $F$ est une extension
{\it finie} de $E$.
\end{definition}

\begin{example}
\label{fields-example-C-over-R}
Le corps $\mathbf{C}$ est un espace vectoriel de dimension deux sur $\mathbf{R}$,
de base $1, i$. Ainsi $\mathbf{C}$ est une extension finie de $\mathbf{R}$
de degré 2.
\end{example}

\begin{lemma}
\label{fields-lemma-finite-goes-up}
Soit $K/E/F$ une tour d'extensions algébriques de corps.
Si $K$ est fini sur $F$, alors $K$ est fini sur $E$.
\end{lemma}

\begin{proof}
Cela résulte directement de la définition.
\end{proof}

\noindent
Calculons maintenant le degré dans l'exemple particulier le plus important, celui
fourni par le Lemme \ref{fields-lemma-field-extension-generated-by-one-element}, au moyen
des deux exemples suivants.

\begin{example}[Degré d'un corps de fonctions rationnelles]
\label{fields-example-degree-rational-function-field}
Si $k$ est un corps quelconque, le corps des fonctions rationnelles $k(t)$ n'est
{\it pas} une extension finie. Par exemple, les éléments
$\left\{t^n, n \in \mathbf{Z}\right\}$ sont linéairement indépendants sur $k$.

\medskip\noindent
En fait, si $k$ est non dénombrable, alors $k(t)$ est de dimension {\it non dénombrable}
comme $k$-espace vectoriel. Pour le montrer, affirmons que la famille d'éléments
$\{1/(t- \alpha), \alpha \in k\} \subset k(t)$ est linéairement indépendante sur
$k$. Une relation non triviale entre eux conduirait à une contradiction : par
exemple, si l'on travaille sur $\mathbf{C}$, cela résulte du fait que
$\frac{1}{t-\alpha}$, considéré comme une fonction méromorphe sur
$\mathbf{C}$, possède un pôle en $\alpha$ et nulle part ailleurs.
Par conséquent, toute somme $\sum c_i \frac{1}{t - \alpha_i}$, où les $c_i \in k^*$
et les $\alpha_i \in k$ sont distincts, aurait un pôle en chacun des $\alpha_i$.
Elle ne pourrait notamment pas être nulle.

\medskip\noindent
Fait amusant, ceci donne une démonstration rapide du Nullstellensatz de Hilbert sur
les nombres complexes. Pour un résultat légèrement plus général, voir
Algèbre, Théorème \ref{algebra-theorem-uncountable-nullstellensatz}.
\end{example}

\begin{lemma}
\label{fields-lemma-finite-finitely-generated}
Toute extension finie de corps est une extension de corps de type fini.
La réciproque est fausse.
\end{lemma}

\begin{proof}
Soit $F/E$ une extension finie de corps. Soient $\alpha_1, \ldots, \alpha_n$
une base de $F$ comme espace vectoriel sur $E$. Alors
$F = E(\alpha_1, \ldots, \alpha_n)$ ; ainsi $F/E$ est une extension de corps
de type fini. La réciproque est fausse d'après
l'Exemple \ref{fields-example-degree-rational-function-field}.
\end{proof}

\begin{example}[Degré d'une extension algébrique simple]
\label{fields-example-degree-simple-algebraic-extension}
Considérons une extension de corps monogène $E/k$ de la forme étudiée dans
l'Exemple \ref{fields-example-monogenic-extension}.
Autrement dit, $E = k[t]/(P)$, où $P \in k[t]$ est un polynôme irréductible.
Alors le degré $[E : k]$ est simplement le degré $d = \deg(P)$ du
polynôme $P$. Écrivons en effet
\begin{equation}
\label{fields-equation-P}
P = a_d t^d + a_{d - 1} t^{d - 1} + \ldots + a_0.
\end{equation}
avec $a_d \not = 0$. Les images de $1, t, \ldots, t^{d - 1}$ dans
$k[t]/(P)$ sont alors linéairement indépendantes sur $k$, car toute relation entre
elles serait de degré strictement inférieur à celui de $P$, alors que $P$ est
l'élément de plus petit degré de l'idéal $(P)$.

\medskip\noindent
Réciproquement, l'ensemble $S = \{1, t, \ldots, t^{d - 1}\}$ (ou, plus
précisément, l'ensemble de leurs images) engendre $k[t]/(P)$ comme espace vectoriel.
En effet, il résulte de (\ref{fields-equation-P}) que $a_d t^d$ appartient au sous-espace engendré par $S$.
Comme $a_d$ est inversible, on voit que $t^d$ appartient au sous-espace engendré par $S$.
De même, la relation $t P(t) = 0$ montre que l'image de $t^{d + 1}$
appartient au sous-espace engendré par $\{1, t, \ldots, t^d\}$ --- donc, d'après ce qui précède,
au sous-espace engendré par $S$. Une récurrence ascendante montre alors
que l'image de $t^n$, pour $n \geq d$, appartient au sous-espace engendré par $S$.
\end{example}

\noindent
Ceci confirme, par exemple, que $[\mathbf{C}: \mathbf{R}] = 2$. Plus généralement,
si $k$ est un corps et $\alpha \in k$ n'est pas un carré, le polynôme irréductible
$x^2 - \alpha \in k[x]$ permet de construire une extension $k[x]/(x^2 - \alpha)$
de degré deux.
Nous noterons celle-ci $k(\sqrt{\alpha})$. De telles extensions seront dites
{\it quadratiques,} pour des raisons évidentes.

\medskip\noindent
Le fait fondamental concernant le degré est qu'il est {\it multiplicatif dans les tours.}

\begin{lemma}[Multiplicativité]
\label{fields-lemma-multiplicativity-degrees}
Supposons donnée une tour de corps $F/E/k$. Alors
$$
[F:k] = [F:E][E:k]
$$
\end{lemma}

\begin{proof}
Soient $\alpha_1, \ldots, \alpha_n \in F$ une $E$-base de $F$, et
$\beta_1, \ldots, \beta_m \in E$ une $k$-base de $E$. L'affirmation est
que l'ensemble des produits
$\{\alpha_i \beta_j, 1 \leq i \leq n, 1 \leq j \leq m\}$
est une $k$-base de $F$. Vérifions d'abord qu'ils engendrent $F$ sur $k$.

\medskip\noindent
Par hypothèse, les $\{\alpha_i\}$ engendrent $F$ sur $E$. Ainsi, si
$f \in F$, il existe des $a_i \in E$ tels que
$$
f = \sum\nolimits_i a_i \alpha_i,
$$
et, pour chaque $i$, on peut écrire $a_i = \sum b_{ij} \beta_j$ avec des
$b_{ij} \in k$. En regroupant ces égalités, on obtient
$$
f = \sum\nolimits_{i,j} b_{ij} \alpha_i \beta_j,
$$
ce qui prouve que les $\{\alpha_i \beta_j\}$ engendrent $F$ sur $k$.

\medskip\noindent
Supposons maintenant qu'il existe une relation non triviale
$$
\sum\nolimits_{i,j} c_{ij} \alpha_i \beta_j = 0
$$
avec des $c_{ij} \in k$. On aurait alors
$$
\sum\nolimits_i \alpha_i \left( \sum\nolimits_j c_{ij} \beta_j \right) = 0,
$$
et les termes entre parenthèses appartiennent à $E$, puisque les $\beta_j$ y appartiennent.
L'indépendance $E$-linéaire des $\{\alpha_i\}$ montre que toutes les sommes intérieures sont nulles.
Puis l'indépendance $k$-linéaire des $\{\beta_j\}$ montre que tous les
$c_{ij}$ sont nuls.
\end{proof}

\noindent
Ouvrons une parenthèse pour donner une définition quelque peu tangente à notre propos.

\begin{definition}
\label{fields-definition-number-field}
On dit qu'un corps $K$ est un {\it corps de nombres} s'il est de caractéristique
$0$ et si l'extension $K/\mathbf{Q}$ est finie.
\end{definition}

\noindent
Les corps de nombres sont les objets fondamentaux de la théorie algébrique des nombres. Nous verrons
plus loin que,
pour l'analogue des entiers $\mathbf{Z}$ dans un corps de nombres, une propriété
assez proche de la factorisation unique subsiste (bien qu'il n'y ait généralement
pas de factorisation unique au sens strict !).


\section{Extensions algébriques}
\label{fields-section-algebraic-extensions}

\noindent
Une classe importante d'extensions est formée de celles dans lesquelles chaque élément engendre
une extension finie.

\begin{definition}
\label{fields-definition-algebraic}
Considérons une extension de corps $F/E$. Un élément $\alpha \in F$ est dit
{\it algébrique} sur $E$ si $\alpha$ est racine d'un polynôme non nul
à coefficients dans $E$. Si tous les éléments de $F$ sont algébriques, on dit que $F$ est
une {\it extension algébrique} de $E$.
\end{definition}

\noindent
D'après le Lemme \ref{fields-lemma-field-extension-generated-by-one-element}, la
sous-extension $E(\alpha)$ est isomorphe soit au corps des fonctions rationnelles
$E(t)$, soit à un anneau quotient $E[t]/(P)$, où $P \in E[t]$ est un
polynôme irréductible. Dans le second cas, $\alpha$ est algébrique sur
$E$ (en fait, la preuve du
Lemme \ref{fields-lemma-field-extension-generated-by-one-element}
montre que l'on peut choisir $P$ de sorte que $\alpha$ soit une racine de $P$) ;
dans le premier cas, il ne l'est pas.

\begin{example}
\label{fields-example-C-algebraic-over-R}
Le corps $\mathbf{C}$ est algébrique sur $\mathbf{R}$. En effet, si
$\alpha = a + ib$ dans $\mathbf{C}$, alors $\alpha^2 - 2a\alpha + a^2 + b^2 = 0$
est une équation polynomiale satisfaite par $\alpha$ sur $\mathbf{R}$.
\end{example}

\begin{example}
\label{fields-example-compact-riemann-surface-is-finite-over-P1}
Soient $X$ une surface de Riemann compacte et
$f \in \mathbf{C}(X) - \mathbf{C}$ une fonction méromorphe non constante quelconque
sur $X$ (voir l'Exemple \ref{fields-example-field-of-meromorphic-functions}). On sait alors
que $\mathbf{C}(X)$ est algébrique sur la sous-extension
$\mathbf{C}(f)$ engendrée par $f$. Nous ne le démontrerons pas.
\end{example}

\begin{lemma}
\label{fields-lemma-algebraic-goes-up}
Soit $K/E/F$ une tour d'extensions de corps.
\begin{enumerate}
\item Si $\alpha \in K$ est algébrique sur $F$, alors $\alpha$ est algébrique
sur $E$.
\item Si $K$ est algébrique sur $F$, alors $K$ est algébrique sur $E$.
\end{enumerate}
\end{lemma}

\begin{proof}
Cela résulte immédiatement des définitions.
\end{proof}

\noindent
Montrons à présent qu'il existe un lien étroit entre la finitude et le caractère
algébrique.

\begin{lemma}
\label{fields-lemma-finite-is-algebraic}
Toute extension finie est algébrique. En fait, une extension $E/k$ est algébrique
si et seulement si toute sous-extension $k(\alpha)/k$ engendrée par un
$\alpha \in E$ est finie.
\end{lemma}

\noindent
En général, il est tout à fait faux qu'une extension algébrique soit finie.

\begin{proof}
Soit $E/k$ finie, de degré $n$. Choisissons $\alpha \in E$. Les
éléments $1, \alpha, \ldots, \alpha^n$ sont linéairement
dépendants sur $k$, sans quoi on aurait nécessairement $[E : k] > n$. Une relation de
dépendance linéaire fournit alors le polynôme recherché que $\alpha$ doit annuler.

\medskip\noindent
Pour la dernière assertion, remarquons qu'une extension monogène $k(\alpha)/k$ est
finie si et seulement si $\alpha$ est algébrique sur $k$, d'après les
Exemples \ref{fields-example-degree-rational-function-field} et
\ref{fields-example-degree-simple-algebraic-extension}.
Ainsi, si $E/k$ est algébrique, chaque $k(\alpha)/k$, pour $\alpha \in E$, est une extension
finie, et réciproquement.
\end{proof}

\noindent
Nous pouvons extraire un lemme de la preuve précédente (en réalité, des
Exemples \ref{fields-example-degree-rational-function-field} et
\ref{fields-example-degree-simple-algebraic-extension}) :
une extension monogène est finie si et seulement si elle est algébrique.
Nous utiliserons cette observation dans le résultat suivant.

\begin{lemma}
\label{fields-lemma-algebraic-finitely-generated}
\begin{slogan}
Une extension algébrique de type fini est finie.
\end{slogan}
Soient $k$ un corps et $\alpha_1, \alpha_2, \ldots, \alpha_n$ des éléments
d'un corps d'extension tels que chaque $\alpha_i$ soit algébrique sur $k$. Alors
l'extension $k(\alpha_1, \ldots, \alpha_n)/k$ est finie.
Autrement dit, une extension algébrique de type fini est finie.
\end{lemma}

\begin{proof}
En effet, chaque extension
$k(\alpha_{1}, \ldots, \alpha_{i+1})/k(\alpha_1, \ldots, \alpha_{i})$
est engendrée par un élément et algébrique, donc finie.
La multiplicativité du degré (Lemme \ref{fields-lemma-multiplicativity-degrees})
donne le résultat.
\end{proof}

\noindent
L'ensemble des nombres complexes qui sont algébriques sur $\mathbf{Q}$ est simplement
appelé l'ensemble des {\it nombres algébriques.} Par exemple, $\sqrt{2}$ est algébrique,
$i$ est algébrique, mais $\pi$ ne l'est pas.
C'est un fait fondamental que les nombres algébriques forment un corps, bien qu'il ne soit pas
évident de le démontrer à partir de la définition selon laquelle un nombre est algébrique
exactement lorsqu'il satisfait une équation polynomiale non nulle à coefficients
rationnels (par exemple, en manipulant des équations polynomiales).

\begin{lemma}
\label{fields-lemma-algebraic-elements}
Soit $E/k$ une extension de corps. Les éléments de $E$ algébriques sur $k$
forment une sous-extension de $E/k$.
\end{lemma}

\begin{proof}
Soient $\alpha, \beta \in E$ algébriques sur $k$. Alors $k(\alpha, \beta)/k$
est une extension finie d'après le Lemme \ref{fields-lemma-algebraic-finitely-generated}.
Il s'ensuit que $k(\alpha + \beta) \subset k(\alpha, \beta)$ est une extension
finie, ce qui implique que $\alpha + \beta$ est algébrique d'après le
Lemme \ref{fields-lemma-finite-is-algebraic}. Il en va de même pour la différence,
le produit et, lorsque $\beta \neq 0$, le quotient $\alpha/\beta$.
\end{proof}

\noindent
De nombreuses propriétés remarquables des extensions de corps, tout comme celles des anneaux, sont
préservées par les tours et les corps composés.

\begin{lemma}
\label{fields-lemma-algebraic-permanence}
Soient $E/k$ et $F/E$ des extensions algébriques de corps. Alors $F/k$ est une
extension algébrique de corps.
\end{lemma}

\begin{proof}
Choisissons $\alpha \in F$. Alors $\alpha$ est algébrique sur $E$.
L'observation essentielle est que $\alpha$ est algébrique sur une
sous-extension de type fini de $k$.
Autrement dit, il existe un ensemble fini $S \subset E$ tel que $\alpha $ soit algébrique
sur $k(S)$ : cela est clair, puisque le fait d'être algébrique signifie qu'il existe un
polynôme de $E[x]$ annulé par $\alpha$, et l'on peut prendre pour $S$ l'ensemble des
coefficients de ce polynôme. Il en résulte que $\alpha$ est algébrique sur
$k(S)$. En particulier, l'extension $k(S, \alpha)/ k(S)$ est finie.
Comme $S$ est un ensemble fini et que $k(S)/k$ est algébrique, le
Lemme \ref{fields-lemma-algebraic-finitely-generated} montre que
$k(S)/k$ est finie. En utilisant la multiplicativité
(Lemme \ref{fields-lemma-multiplicativity-degrees}),
on trouve que $k(S,\alpha)/k$ est finie ; ainsi $\alpha$ est algébrique sur $k$.
\end{proof}

\noindent
La méthode de la preuve précédente --- le fait que la propriété d'être algébrique
sur $E$ {\it descendait} à une sous-extension de type fini
de $E$ --- est une idée récurrente en algèbre.
Elle permet souvent, par exemple, de ramener des questions générales d'algèbre commutative
au cas noethérien.

\begin{lemma}
\label{fields-lemma-size-algebraic-extension}
Soit $E/F$ une extension algébrique de corps. Alors le cardinal $|E|$
de $E$ est au plus $\max(\aleph_0, |F|)$.
\end{lemma}

\begin{proof}
Soit $S$ l'ensemble des polynômes non constants à coefficients dans $F$.
Pour tout $P \in S$, l'ensemble des racines
$r(P, E) = \{\alpha \in E \mid P(\alpha) = 0\}$
est fini (les détails sont omis). De plus, le fait que $E$ soit algébrique
sur $F$ implique que $E = \bigcup_{P \in S} r(P, E)$.
Il est clair que le cardinal de $S$ est majoré par $\max(\aleph_0, |F|)$,
car cet ensemble est une réunion dénombrable de produits finis de copies de $F$.
Il en va donc de même pour $E$.
\end{proof}

\begin{lemma}
\label{fields-lemma-subalgebra-algebraic-extension-field}
Soit $E/F$ une extension finie, ou plus généralement algébrique, de corps.
Tout sous-anneau $F \subset R \subset E$ est un corps.
\end{lemma}

\begin{proof}
Soit $\alpha \in R$ non nul. Alors $1, \alpha, \alpha^2, \ldots$
appartiennent à $R$. D'après le Lemme \ref{fields-lemma-finite-is-algebraic},
il existe une relation non triviale
$a_0 + a_1 \alpha + \ldots + a_d \alpha^d = 0$, avec $a_i \in F$.
Nous pouvons supposer $a_0 \not = 0$, car sinon nous pouvons diviser la relation
par $\alpha$ afin de diminuer $d$. On obtient alors
$$
a_0 = \alpha (- a_1  - \ldots - a_d \alpha^{d - 1})
$$
ce qui prouve que l'inverse de $\alpha$ est l'élément
$a_0^{-1} (- a_1  - \ldots - a_d \alpha^{d - 1})$
de $R$.
\end{proof}

\begin{lemma}
\label{fields-lemma-algebraic-extension-self-map}
Soit $E/F$ une extension algébrique de corps. Tout morphisme de $F$-algèbres
$f : E \to E$ est un automorphisme.
\end{lemma}

\begin{proof}
Si $E/F$ est finie, alors $f : E \to E$ est une application $F$-linéaire
injective (Lemme \ref{fields-lemma-field-maps-injective})
entre espaces vectoriels de dimension finie, donc est bijective.
Dans le cas général, on voit encore que $f$ est injective.
Soit $\alpha \in E$, et soit $P \in F[x]$ un
polynôme tel que $P(\alpha) = 0$.
Soit $E' \subset E$ le sous-corps de $E$ engendré
par les racines $\alpha = \alpha_1, \ldots, \alpha_n$ de $P$ dans $E$.
Alors $E'$ est fini sur $F$ d'après le Lemme \ref{fields-lemma-algebraic-finitely-generated}.
Comme $f$ préserve l'ensemble des racines, on a
$f|_{E'} : E' \to E'$. Ainsi $f|_{E'}$ est un isomorphisme
d'après la première partie de la preuve, et l'on conclut que $\alpha$
appartient à l'image de $f$.
\end{proof}








\section{Polynômes minimaux}
\label{fields-section-minimal-polynomials}

\noindent
Soient $E/k$ une extension de corps et $\alpha \in E$ algébrique sur $k$.
Alors $\alpha$ satisfait une équation polynomiale (non triviale) dans $k[x]$.
Considérons l'ensemble des polynômes $P \in k[x]$ tels que $P(\alpha) = 0$ ; par
hypothèse, cet ensemble ne contient pas seulement le polynôme nul.
On vérifie facilement que cet ensemble est un {\it idéal.} C'est en effet le noyau
du morphisme
$$
k[x] \to E, \quad x \mapsto \alpha
$$
Puisque $k[x]$ est un anneau principal, il existe un {\it générateur} $P \in k[x]$ de cet
idéal. Si l'on suppose, sans perte de généralité, que $P$ est unitaire, alors $P$ est
déterminé de manière unique.

\begin{definition}
\label{fields-definition-minimal-polynomial}
Le polynôme $P$ ci-dessus est appelé le {\it polynôme minimal}
de $\alpha$ sur $k$.
\end{definition}

\noindent
Le polynôme minimal possède la caractérisation suivante : c'est le polynôme unitaire,
de plus petit degré, qui annule $\alpha$. Tout multiple non constant de $P$
sera de degré plus grand, et seuls les multiples de $P$ peuvent
annuler $\alpha$. Cela explique le qualificatif {\it minimal}.

\medskip\noindent
Le polynôme minimal est manifestement {\it irréductible}. Cela équivaut à affirmer
que l'idéal de $k[x]$ formé des polynômes qui annulent
$\alpha$ est premier. Ce fait résulte de ce que le morphisme
$k[x] \to E, x \mapsto \alpha$ prend ses valeurs dans un anneau intègre (et même dans un corps), de sorte que son
noyau est un idéal premier.

\begin{lemma}
\label{fields-lemma-degree-minimal-polynomial}
Le degré du polynôme minimal est $[k(\alpha) : k]$.
\end{lemma}

\begin{proof}
Il s'agit simplement d'une reformulation de l'argument donné dans le
Lemme \ref{fields-lemma-field-extension-generated-by-one-element} : l'observation
est que, si $P$ est le polynôme minimal de $\alpha$, le morphisme
$$
k[x]/(P) \to k(\alpha), \quad x \mapsto \alpha
$$
est un isomorphisme comme dans la preuve précitée, et nous avons calculé le
degré d'une telle extension (voir
l'Exemple \ref{fields-example-degree-simple-algebraic-extension}).
\end{proof}

\noindent
L'observation de la preuve ci-dessus est donc que, si $\alpha \in E$ est algébrique,
alors $k(\alpha) \subset E$ est isomorphe à $k[x]/(P)$.


\section{Clôture algébrique}
\label{fields-section-algebraic-closure}

\noindent
Le ``théorème fondamental de l'algèbre'' affirme que $\mathbf{C}$ est
algébriquement clos. Une belle démonstration de ce résultat utilise
le théorème de Liouville en analyse complexe ; nous en donnerons une autre
(voir le Lemme \ref{fields-lemma-C-algebraically-closed}).

\begin{definition}
\label{fields-definition-algebraically-closed}
On dit qu'un corps $F$ est {\it algébriquement clos} si toute extension algébrique
$E/F$ est triviale, c'est-à-dire si $E = F$.
\end{definition}

\noindent
Ce n'est peut-être pas la définition retenue dans tous les ouvrages. Le lemme suivant
la compare à l'autre définition usuelle.

\begin{lemma}
\label{fields-lemma-algebraically-closed}
Soit $F$ un corps. Les assertions suivantes sont équivalentes :
\begin{enumerate}
\item $F$ est algébriquement clos,
\item tout polynôme irréductible sur $F$ est de degré un,
\item tout polynôme non constant sur $F$ possède une racine,
\item tout polynôme non constant sur $F$ est un produit de facteurs de degré un.
\end{enumerate}
\end{lemma}

\begin{proof}
Si $F$ est algébriquement clos, alors tout polynôme irréductible est de degré un.
En effet, s'il existe un polynôme irréductible de degré $> 1$, celui-ci
engendre une extension de corps finie non triviale (donc algébrique) ; voir
l'Exemple \ref{fields-example-degree-simple-algebraic-extension}.
Ainsi (1) implique (2). Si tout polynôme irréductible
est de degré un, alors tout polynôme irréductible possède une racine, et par conséquent tout
polynôme non constant possède une racine. Ainsi (2) implique (3).

\medskip\noindent
Supposons que tout polynôme non constant possède une racine. Soit $P \in F[x]$
non constant. Si $P(\alpha) = 0$ avec $\alpha \in F$, on voit alors
que $P = (x - \alpha)Q$ pour un certain $Q \in F[x]$ (par division euclidienne).
On peut donc raisonner par récurrence sur le degré pour écrire tout polynôme non constant
sous la forme d'un produit $c \prod (x - \alpha_i)$.

\medskip\noindent
Supposons enfin que tout polynôme non constant sur $F$ soit un produit de
facteurs de degré un. Soit $E/F$ une extension algébrique. Toutes les sous-extensions
simples $F(\alpha)/F$ de $E$ sont alors nécessairement triviales (car, par
hypothèse, les seuls polynômes irréductibles sont de degré un). Ainsi $E = F$.
Nous voyons que (4) implique (1), ce qui achève la preuve.
\end{proof}

\noindent
Nous voulons maintenant définir une extension algébrique ``universelle'' d'un corps.
En réalité, il convient d'être prudent : la clôture algébrique n'est {\it pas} un
objet universel. En effet, la clôture algébrique n'est pas unique à
{\it unique} isomorphisme près : elle n'est unique qu'à isomorphisme près. Elle
n'en sera pas moins très utile, bien que non fonctorielle.

\begin{definition}
\label{fields-definition-algebraic-closure}
Soit $F$ un corps. Une {\it clôture algébrique} de $F$ est un corps
$\overline{F}$ contenant $F$ tel que :
\begin{enumerate}
\item $\overline{F}$ est algébrique sur $F$.
\item $\overline{F}$ est algébriquement clos.
\end{enumerate}
\end{definition}

\noindent
Si $F$ est algébriquement clos, alors $F$ est sa propre clôture algébrique.
Démontrons maintenant le résultat fondamental d'existence.

\begin{theorem}
\label{fields-theorem-existence-algebraic-closure}
Tout corps possède une clôture algébrique.
\end{theorem}

\noindent
La preuve sera pour l'essentiel sans rapport avec le reste du chapitre. Nous voudrons toutefois
savoir qu'il est {\it possible} de plonger un corps dans un
corps algébriquement clos, et nous supposerons souvent ce plongement effectué.

\begin{proof}
Soit $F$ un corps. D'après le Lemme \ref{fields-lemma-size-algebraic-extension}, le
cardinal d'une extension algébrique de $F$ est majoré par
$\max(\aleph_0, |F|)$. Choisissons un ensemble $S$ contenant $F$ tel que
$|S| > \max(\aleph_0, |F|)$. Considérons les triplets
$(E, \sigma_E, \mu_E)$ où
\begin{enumerate}
\item $E$ est un ensemble tel que $F \subset E \subset S$, et
\item $\sigma_E : E \times E \to E$ et $\mu_E : E \times E \to E$
sont des applications d'ensembles telles que $(E, \sigma_E, \mu_E)$ définisse la structure
d'une extension de corps de $F$ (en particulier, $\sigma_E(a, b) = a +_F b$
pour $a, b \in F$, et de même pour $\mu_E$), et
\item $E/F$ est une extension algébrique de corps.
\end{enumerate}
La collection de tous les triplets $(E, \sigma_E, \mu_E)$ forme un ensemble $I$.
Pour $i \in I$, nous noterons $E_i = (E_i, \sigma_i, \mu_i)$ l'extension de corps
de $F$ correspondante. Nous définissons un ordre partiel sur
$I$ en déclarant que $i \leq i'$ si et seulement si $E_i \subset E_{i'}$
(cela a un sens puisque $E_i$ et $E_{i'}$ sont des sous-ensembles du même ensemble $S$)
et si $\sigma_i = \sigma_{i'}|_{E_i \times E_i}$ et
$\mu_i = \mu_{i'}|_{E_i \times E_i}$ ; autrement dit, $E_{i'}$ est une extension de corps
de $E_i$.

\medskip\noindent
Soit $T \subset I$ un sous-ensemble totalement ordonné. Il est alors clair que
$E_T = \bigcup_{i \in T} E_i$, muni des applications induites $\sigma_T = \bigcup \sigma_i$
et $\mu_T = \bigcup \mu_i$, est un autre élément de $I$. Autrement dit,
tout sous-ensemble totalement ordonné de $I$ possède un majorant dans $I$. D'après le lemme de Zorn,
il existe un élément maximal $(E, \sigma_E, \mu_E)$ de $I$. Nous affirmons que
$E$ est une clôture algébrique. Puisque, par définition de $I$, l'extension
$E/F$ est algébrique, il suffit de montrer que $E$ est algébriquement clos.

\medskip\noindent
Pour le voir, raisonnons par l'absurde. Supposons donc que $E$ ne soit pas
algébriquement clos. Il existe alors un polynôme irréductible
$P$ sur $E$ de degré $> 1$ ; voir le Lemme \ref{fields-lemma-algebraically-closed}.
D'après le Lemme \ref{fields-lemma-finite-is-algebraic}, on obtient une extension finie
non triviale $E' = E[x]/(P)$. Remarquons que $E'/F$ est algébrique d'après le
Lemme \ref{fields-lemma-algebraic-permanence}.
Le cardinal de $E'$ est donc $\leq \max(\aleph_0, |F|)$.
La théorie élémentaire des ensembles permet de prolonger l'injection donnée
$E \subset S$ en une injection $E' \to S$. Autrement dit, nous pouvons
considérer $E'$ comme un élément de notre ensemble $I$, ce qui contredit la
maximalité de $E$. Cette contradiction achève la preuve.
\end{proof}

\begin{lemma}
\label{fields-lemma-map-into-algebraic-closure}
Soit $F$ un corps. Soit $\overline{F}$ une clôture algébrique de $F$.
Soit $M/F$ une extension algébrique. Il existe alors un morphisme
d'extensions de $F$, $M \to \overline{F}$.
\end{lemma}

\begin{proof}
Considérons l'ensemble $I$ des couples $(E, \varphi)$, où $F \subset E \subset M$
est une sous-extension et $\varphi : E \to \overline{F}$ un morphisme
d'extensions de $F$. Munissons $I$ de l'ordre partiel défini en déclarant que
$(E, \varphi) \leq (E', \varphi')$ si et seulement si $E \subset E'$ et
$\varphi'|_E = \varphi$. Si $T = \{(E_t,  \varphi_t)\} \subset I$
est un sous-ensemble totalement ordonné, alors
$\bigcup \varphi_t : \bigcup E_t \to \overline{F}$ est un élément de $I$.
Ainsi, tout sous-ensemble totalement ordonné de $I$ possède un majorant.
Le lemme de Zorn assure l'existence d'un élément maximal $(E, \varphi)$ de $I$.
Nous affirmons que $E = M$, ce qui achèvera la preuve. Sinon,
choisissons $\alpha \in M$, $\alpha \not \in E$. Alors $\alpha$ est algébrique
sur $E$ ; voir le Lemme \ref{fields-lemma-algebraic-goes-up}.
Soit $P$ le polynôme minimal de $\alpha$ sur $E$.
Soit $P^\varphi$ l'image de $P$ par $\varphi$ dans $\overline{F}[x]$.
Comme $\overline{F}$ est algébriquement clos, il existe une racine $\beta$
de $P^\varphi$ dans $\overline{F}$. On peut alors prolonger $\varphi$ en
$\varphi' : E(\alpha) = E[x]/(P) \to \overline{F}$ en envoyant
$x$ sur $\beta$. Cela contredit la maximalité de $(E, \varphi)$,
comme voulu.
\end{proof}

\begin{lemma}
\label{fields-lemma-algebraic-closures-isomorphic}
Deux clôtures algébriques quelconques d'un corps sont isomorphes.
\end{lemma}

\begin{proof}
Soit $F$ un corps. Si $M$ et $\overline{F}$ sont des clôtures algébriques de
$F$, il existe alors un morphisme d'extensions de $F$
$\varphi : M \to \overline{F}$ d'après le
Lemme \ref{fields-lemma-map-into-algebraic-closure}.
Or l'image $\varphi(M)$ est algébriquement close.
D'autre part, l'extension $\varphi(M) \subset \overline{F}$
est algébrique d'après le Lemme \ref{fields-lemma-algebraic-goes-up}.
Ainsi $\varphi(M) = \overline{F}$.
\end{proof}





\section{Polynômes premiers entre eux}
\label{fields-section-relatively-prime}

\noindent
Soit $K$ un corps algébriquement clos. L'anneau $K[x]$ possède alors une structure
d'idéaux très simple, comme nous l'avons vu au Lemme \ref{fields-lemma-algebraically-closed}.
En particulier, tout polynôme $P \in K[x]$ peut s'écrire
$$
P = c(x - \alpha_1) \ldots (x - \alpha_n),
$$
où $c$ est le coefficient dominant et où les $\alpha_1, \ldots, \alpha_n \in K$
sont les racines de $P$ (comptées avec multiplicité). Manifestement, les seuls polynômes
irréductibles de $K[x]$ sont les polynômes de degré un $c(x - \alpha)$,
$c, \alpha \in K$ (avec $c \neq 0$).

\begin{definition}
\label{fields-definition-relatively-prime}
Si $k$ est un corps quelconque, nous disons que deux polynômes de $k[x]$ sont
{\it premiers entre eux} s'ils engendrent l'idéal unité de $k[x]$.
\end{definition}

\noindent
Poursuivons la discussion précédente. Si $K$ est un corps algébriquement clos,
deux polynômes de $K[x]$ sont premiers entre eux si et seulement s'ils n'ont aucune
racine commune. Cela résulte du fait que les idéaux maximaux de $K[x]$ sont de la forme
$(x - \alpha)$, $\alpha \in K$. Ainsi, si $F, G \in K[x]$ n'ont aucune racine commune,
alors $(F, G)$ ne peut être contenu dans aucun $(x - \alpha)$ (car ils auraient alors
une racine commune en $\alpha$).

\medskip\noindent
Si $k$ n'est {\it pas} algébriquement clos, ceci fournit encore des
informations permettant de déterminer quand deux polynômes de $k[x]$ engendrent l'idéal unité.

\begin{lemma}
\label{fields-lemma-relatively-prime-polynomials}
Deux polynômes de $k[x]$ sont premiers entre eux si et seulement s'ils
n'ont aucune racine commune dans une clôture algébrique $\overline{k}$ de $k$.
\end{lemma}

\begin{proof}
L'affirmation est que deux polynômes quelconques $P, Q$ engendrent $(1)$ dans $k[x]$ si et
seulement s'ils engendrent $(1)$ dans $\overline{k}[x]$. C'est une question
d'algèbre linéaire : un système d'équations linéaires à coefficients dans $k$ possède
une solution si et seulement s'il possède une solution dans une extension quelconque de $k$.
On peut par conséquent se ramener au cas d'un corps algébriquement clos ; dans
ce cas, le résultat découle clairement de ce qui précède.
\end{proof}





\section{Extensions algébriques séparables}
\label{fields-section-separable-extensions}

\noindent
En caractéristique $p$, un phénomène curieux se produit pour les polynômes irréductibles
sur un corps. Le lemme suivant l'explique.

\begin{lemma}
\label{fields-lemma-irreducible-polynomials}
Soit $F$ un corps. Soit $P \in F[x]$ un polynôme irréductible sur $F$.
Soit $P' = \text{d}P/\text{d}x$ le polynôme dérivé de $P$ par rapport
à $x$. L'un des deux cas suivants se produit :
\begin{enumerate}
\item $P$ et $P'$ sont premiers entre eux, ou
\item $P'$ est le polynôme nul.
\end{enumerate}
Le second cas ne peut se produire que si $F$ est de caractéristique $p > 0$.
Dans ce cas, $P(x) = Q(x^q)$, où $q = p^f$ est une puissance de $p$ et
$Q \in F[x]$ un polynôme irréductible tel que $Q$ et $Q'$ soient
premiers entre eux.
\end{lemma}

\begin{proof}
Remarquons que $P'$ est de degré $< \deg(P)$. Par conséquent, si $P$ et $P'$ ne sont pas premiers
entre eux, alors $(P, P') = (R)$, où $R$ est un polynôme de degré $< \deg(P)$,
ce qui contredit l'irréductibilité de $P$. Cela établit la dichotomie
entre (1) et (2).

\medskip\noindent
Supposons que nous soyons dans le cas (2) et que $P = a_d x^d + \ldots + a_0$. Alors
$P' = da_d x^{d - 1} + \ldots + a_1$. En caractéristique $0$, on voit
que cela impose $a_d, \ldots, a_1 = 0$, ce qui signifierait que $P$ est constant,
une contradiction. On conclut donc que la caractéristique $p$ est positive.
Dans ce cas, la condition $P' = 0$ impose $a_i = 0$ chaque fois que $p$ ne
divise pas $i$.
Autrement dit, $P(x) = P_1(x^p)$ pour un certain polynôme non constant $P_1$.
Manifestement, $P_1$ est lui aussi irréductible. Par récurrence sur le degré, on
voit que $P_1(x) = Q(x^q)$ comme dans l'énoncé du lemme ; ainsi
$P(x) = Q(x^{pq})$, et le lemme est démontré.
\end{proof}

\begin{definition}
\label{fields-definition-separable}
Soit $F$ un corps. Soit $K/F$ une extension de corps.
\begin{enumerate}
\item On dit qu'un polynôme irréductible $P$ sur $F$ est {\it séparable}
s'il est premier avec son polynôme dérivé.
\item Étant donné $\alpha \in K$ algébrique sur $F$, on dit que $\alpha$ est
{\it séparable} sur $F$ si son polynôme minimal est séparable sur $F$.
\item Si $K$ est une extension algébrique de $F$, on dit que $K$ est
{\it séparable}\footnote{Pour les extensions non algébriques,
cette définition n'a pas de sens et n'est pas la bonne. Nous renvoyons
le lecteur à Algèbre, sections \ref{algebra-section-separability} et
\ref{algebra-section-separability-continued}.}
sur $F$ si tout élément de $K$ est séparable sur $F$.
\end{enumerate}
\end{definition}

\noindent
D'après le Lemme \ref{fields-lemma-irreducible-polynomials}, en caractéristique $0$, tout
polynôme irréductible est séparable, tout élément algébrique d'une extension
est séparable et toute extension algébrique est séparable.

\begin{lemma}
\label{fields-lemma-separable-goes-up}
Soit $K/E/F$ une tour d'extensions algébriques de corps.
\begin{enumerate}
\item Si $\alpha \in K$ est séparable sur $F$, alors $\alpha$ est séparable
sur $E$.
\item Si $K$ est séparable sur $F$, alors $K$ est séparable sur $E$.
\end{enumerate}
\end{lemma}

\begin{proof}
Nous utiliserons le Lemme \ref{fields-lemma-irreducible-polynomials} sans le mentionner davantage.
Soit $P$ le polynôme minimal de $\alpha$ sur $F$.
Soit $Q$ le polynôme minimal de $\alpha$ sur $E$.
Alors $Q$ divise $P$ dans l'anneau de polynômes $E[x]$, disons $P = QR$.
Alors $P' = Q'R + QR'$. Ainsi, si $Q' = 0$, alors $Q$ divise $P$ et $P'$,
donc $P' = 0$ d'après le lemme. Ceci démontre (1). L'assertion (2)
découle immédiatement de (1) et des définitions.
\end{proof}

\begin{lemma}
\label{fields-lemma-recognize-separable}
Soit $F$ un corps. Un polynôme irréductible $P$ sur $F$
est séparable si et seulement si $P$ possède des racines deux à deux distinctes dans une
clôture algébrique de $F$.
\end{lemma}

\begin{proof}
Supposons que $\alpha \in \overline{F}$ soit une racine à la fois de $P$ et de $P'$.
Alors $P = (x - \alpha)Q$ pour un certain polynôme $Q$. En dérivant,
nous obtenons $P' = Q + (x - \alpha)Q'$. Ainsi $\alpha$ est une racine de $Q$.
On voit donc que si $P$ et $P'$ ont une racine commune, alors $P$
ne possède pas de racines deux à deux distinctes. Réciproquement, si $P$ possède
une racine multiple, c'est-à-dire si $(x - \alpha)^2$ divise $P$, alors $\alpha$
est une racine à la fois de $P$ et de $P'$. Conjointement avec le
Lemme \ref{fields-lemma-relatively-prime-polynomials}, ceci démontre le lemme.
\end{proof}

\begin{lemma}
\label{fields-lemma-nr-roots-unchanged}
Soit $F$ un corps et soit $\overline{F}$ une clôture algébrique de $F$.
Soit $p > 0$ la caractéristique de $F$. Soit $P$ un polynôme
sur $F$. Alors les ensembles des racines de $P$ et de $P(x^p)$ dans $\overline{F}$
ont le même cardinal (sans compter les multiplicités).
\end{lemma}

\begin{proof}
Manifestement, $\alpha$ est une racine de $P(x^p)$ si et seulement si $\alpha^p$ est une
racine de $P$. Autrement dit, les racines de $P(x^p)$ sont les racines de
$x^p - \beta$, où $\beta$ est une racine de $P$. Il suffit donc de montrer
que l'application $\overline{F} \to \overline{F}$, $\alpha \mapsto \alpha^p$
est bijective. Elle est surjective, puisque $\overline{F}$ est algébriquement clos,
ce qui signifie que tout élément possède une racine $p$-ième. Elle est injective car
$\alpha^p = \beta^p$ implique $(\alpha - \beta)^p = 0$, puisque
la caractéristique est $p$. Et bien sûr, dans un corps, $x^p = 0$ implique
$x = 0$.
\end{proof}

\noindent
Soit $F$ un corps et soit $P$ un polynôme irréductible sur $F$.
Nous savons alors que $P = Q(x^q)$ pour un certain polynôme irréductible séparable $Q$
(Lemme \ref{fields-lemma-irreducible-polynomials}), où $q$ est une puissance de
la caractéristique $p$ (et, si la caractéristique est nulle, alors
$q = 1$\footnote{Une bonne convention pour ce chapitre consiste à poser $0^0 = 1$.}
et $Q = P$). D'après le Lemme \ref{fields-lemma-nr-roots-unchanged}, le nombre de
racines de $P$ et de $Q$ dans toute clôture algébrique de $F$ est le même.
D'après le Lemme \ref{fields-lemma-recognize-separable}, ce nombre est égal au degré
de $Q$.

\begin{definition}
\label{fields-definition-separable-degree}
Soit $F$ un corps. Soit $P$ un polynôme irréductible sur $F$.
Le {\it degré séparable} de $P$ est le cardinal de
l'ensemble des racines de $P$ dans toute clôture algébrique de $F$ (voir la discussion
ci-dessus). Notation $\deg_s(P)$.
\end{definition}

\noindent
Le degré séparable de $P$ divise toujours son degré, et le quotient
est une puissance de la caractéristique. Si la caractéristique est nulle, alors
$\deg_s(P) = \deg(P)$.

\begin{situation}
\label{fields-situation-finitely-generated}
Ici, $F$ est un corps et $K/F$ est une extension finie engendrée par des éléments
$\alpha_1, \ldots, \alpha_n \in K$. Nous posons $K_0 = F$ et
$$
K_i = F(\alpha_1, \ldots, \alpha_i)
$$
pour obtenir une tour d'extensions finies
$K = K_n / K_{n - 1} / \ldots / K_0 = F$.
Notons $P_i$ le polynôme minimal de $\alpha_i$ sur $K_{i - 1}$.
Enfin, fixons une clôture algébrique $\overline{F}$ de $F$.
\end{situation}

\noindent
Soient $F$, $K$, les $\alpha_i$ et $\overline{F}$ comme dans la
Situation \ref{fields-situation-finitely-generated}.
Supposons que $\varphi : K \to \overline{F}$ soit un morphisme d'extensions
de $F$. Nous obtenons alors des applications $\varphi_i : K_i \to \overline{F}$.
L'image de $P_i \in K_{i - 1}[x]$ par $\varphi_{i - 1}$
est un polynôme $P_i^{\varphi_{i - 1}} \in \overline{F}[x]$
et $\varphi(\alpha_i)$ est une racine de ce polynôme.

\begin{lemma}
\label{fields-lemma-count-embeddings}
Dans la Situation \ref{fields-situation-finitely-generated}, la correspondance
$$
\Mor_F(K, \overline{F})
\longrightarrow
\{(\beta_1, \ldots, \beta_n)\text{ comme ci-dessous}\},
\quad
\varphi \longmapsto (\varphi(\alpha_1), \ldots, \varphi(\alpha_n))
$$
est une bijection. Ici, le membre de droite est l'ensemble des $n$-uplets
$(\beta_1, \ldots, \beta_n)$ d'éléments de $\overline{F}$ possédant
la propriété suivante :
\begin{enumerate}
\item $\beta_1 \in \overline{F}$ est une racine de $P_1$ ;
soit $\varphi_1 : K_1 \to \overline{F}$ l'homomorphisme
de $F$-algèbres qui envoie $\alpha_1$ sur $\beta_1$,
\item $\beta_2 \in \overline{F}$ est une racine de $P_2^{\varphi_1}$ ;
soit $\varphi_2 : K_2 \to \overline{F}$ l'homomorphisme
qui prolonge $\varphi_1$ et envoie $\alpha_2$ sur $\beta_2$,
\item et ainsi de suite jusqu'au dernier cas,
\item $\beta_n \in \overline{F}$ est une racine de $P_n^{\varphi_{n - 1}}$.
\end{enumerate}
À chaque étape, l'homomorphisme $\varphi_i$ existe et est unique, car
$K_i = K_{i - 1}[x]/(P_i)$ et $\beta_i$ est une racine de $P_i^{\varphi_{i - 1}}$.
\end{lemma}

\begin{proof}
L'application de gauche à droite a été décrite avant le lemme.
Soit $(\beta_1, \ldots, \beta_n)$ un élément du membre de droite.
Nous définissons alors $\varphi : K = K_n \to \overline{F}$ comme l'unique homomorphisme
qui prolonge $\varphi_{n - 1}$ et envoie $\alpha_n$ sur $\beta_n$.
L'unicité implique que les deux constructions sont inverses l'une de l'autre.
\end{proof}

\begin{lemma}
\label{fields-lemma-count-embeddings-explicitly}
Dans la Situation \ref{fields-situation-finitely-generated}, nous avons
$|\Mor_F(K, \overline{F})| = \prod_{i = 1}^n \deg_s(P_i)$.
\end{lemma}

\begin{proof}
Ceci découle immédiatement du Lemme \ref{fields-lemma-count-embeddings}.
Remarquons qu'un ingrédient essentiel que nous utilisons ici implicitement est le
caractère bien défini du degré séparable d'un polynôme irréductible,
qui a été constaté juste avant la
Définition \ref{fields-definition-separable-degree}.
\end{proof}

\noindent
Nous utilisons maintenant le résultat précédent pour caractériser les extensions de corps séparables.

\begin{lemma}
\label{fields-lemma-separably-generated-separable}
Hypothèses et notations comme dans la Situation \ref{fields-situation-finitely-generated}.
Si chaque $P_i$ est séparable, c'est-à-dire si chaque $\alpha_i$ est séparable sur
$K_{i - 1}$, alors
$$
|\Mor_F(K, \overline{F})| = [K : F]
$$
et l'extension de corps $K/F$ est séparable. Si l'un des $\alpha_i$ n'est
pas séparable sur $K_{i - 1}$, alors
$|\Mor_F(K, \overline{F})| < [K : F]$.
\end{lemma}

\begin{proof}
Si $\alpha_i$ est séparable sur $K_{i - 1}$, alors
$\deg_s(P_i) = \deg(P_i) = [K_i : K_{i - 1}]$
(la dernière égalité résultant du Lemme \ref{fields-lemma-degree-minimal-polynomial}).
Par multiplicativité (Lemme \ref{fields-lemma-multiplicativity-degrees}), nous avons
$$
[K : F] = \prod [K_i : K_{i - 1}] = \prod \deg(P_i) =
\prod \deg_s(P_i) = |\Mor_F(K, \overline{F})|
$$
où la dernière égalité est le Lemme \ref{fields-lemma-count-embeddings-explicitly}.
Par exactement le même raisonnement, nous obtenons l'inégalité stricte
$|\Mor_F(K, \overline{F})| < [K : F]$ si l'un des $\alpha_i$ n'est
pas séparable sur $K_{i - 1}$.

\medskip\noindent
Supposons enfin, à nouveau, que chaque $\alpha_i$ soit séparable sur $K_{i - 1}$.
Nous allons montrer que $K/F$ est séparable.
Soit $\gamma = \gamma_1 \in K$ quelconque. Nous pouvons alors trouver des éléments supplémentaires
$\gamma_2, \ldots, \gamma_m$ tels que
$K = F(\gamma_1, \ldots, \gamma_m)$ (nous pourrions par exemple prendre
$\gamma_2 = \alpha_1, \ldots, \gamma_{n + 1} = \alpha_n$).
La dernière partie du lemme (déjà démontrée ci-dessus) montre alors
que si $\gamma$ n'était pas séparable sur $F$, nous aurions l'inégalité
stricte $|\Mor_F(K, \overline{F})| < [K : F]$,
en contradiction avec la toute première partie du lemme (elle aussi déjà démontrée
ci-dessus).
\end{proof}

\begin{lemma}
\label{fields-lemma-separable-equality}
Soit $K/F$ une extension finie de corps. Soit $\overline{F}$ une
clôture algébrique de $F$. Nous avons alors
$$
|\Mor_F(K, \overline{F})| \leq [K : F]
$$
avec égalité si et seulement si $K$ est séparable sur $F$.
\end{lemma}

\begin{proof}
C'est un corollaire du Lemme \ref{fields-lemma-separably-generated-separable}.
En effet, puisque $K/F$ est finie, nous pouvons trouver un nombre fini d'éléments
$\alpha_1, \ldots, \alpha_n \in K$ qui engendrent $K$ sur $F$ (nous pouvons par exemple
choisir les $\alpha_i$ de manière à former une base de $K$ sur $F$).
Si $K/F$ est séparable, alors chaque $\alpha_i$ est séparable sur
$F(\alpha_1, \ldots, \alpha_{i - 1})$ d'après le Lemme \ref{fields-lemma-separable-goes-up},
et le Lemme \ref{fields-lemma-separably-generated-separable} donne l'égalité.
Réciproquement, si nous avons l'égalité, alors, quelle que soit la manière de choisir
$\alpha_1, \ldots, \alpha_n$, le Lemme \ref{fields-lemma-separably-generated-separable}
montre que $\alpha_1$ est séparable sur $F$. Comme nous pouvons
commencer la suite par un élément arbitraire de $K$, il s'ensuit
que $K$ est séparable sur $F$.
\end{proof}

\begin{lemma}
\label{fields-lemma-separable-permanence}
Soient $E/k$ et $F/E$ des extensions algébriques séparables de corps. Alors $F/k$
est une extension séparable de corps.
\end{lemma}

\begin{proof}
Choisissons $\alpha \in F$. Alors $\alpha$ est algébrique séparable sur $E$.
Soit $P = x^d + \sum_{i < d} a_i x^i$ le polynôme minimal de
$\alpha$ sur $E$. Chaque $a_i$ est algébrique séparable sur $k$.
Considérons la tour de corps
$$
k \subset k(a_0) \subset k(a_0, a_1) \subset \ldots \subset
k(a_0, \ldots, a_{d - 1}) \subset k(a_0, \ldots, a_{d - 1}, \alpha)
$$
Puisque $a_i$ est algébrique séparable sur $k$, il est algébrique séparable
sur $k(a_0, \ldots, a_{i - 1})$ d'après le Lemme \ref{fields-lemma-separable-goes-up}.
Enfin, $\alpha$ est algébrique séparable sur $k(a_0, \ldots, a_{d - 1})$
car c'est une racine de $P$, qui est irréductible
(puisqu'il est irréductible sur le corps éventuellement plus grand $E$)
et séparable (puisqu'il est séparable sur $E$).
Ainsi $k(a_0, \ldots, a_{d - 1}, \alpha)$ est séparable sur $k$
d'après le Lemme \ref{fields-lemma-separably-generated-separable},
et nous concluons que $\alpha$ est séparable sur $k$, comme voulu.
\end{proof}

\begin{lemma}
\label{fields-lemma-separable-elements}
Soit $E/k$ une extension de corps. Alors les éléments de $E$ séparables
sur $k$ forment une sous-extension de $E/k$.
\end{lemma}

\begin{proof}

Soient $\alpha, \beta \in E$ séparables sur $k$. Alors $\beta$ est séparable
sur $k(\alpha)$ d'après le Lemme \ref{fields-lemma-separable-goes-up}.
Le Lemme \ref{fields-lemma-separably-generated-separable} (appliqué avec $n = 2$,
$\alpha_1 = \alpha$ et $\alpha_2 = \beta$) montre que
$k(\alpha, \beta)$ est séparable sur $k$.

\end{proof}
\section{Indépendance linéaire des caractères}
\label{fields-section-independence-characters}

\noindent
Voici l'énoncé.

\begin{lemma}
\label{fields-lemma-independence-characters}
Soit $L$ un corps. Soit $G$ un monoïde, par exemple un groupe.
Soient
$\chi_1, \ldots, \chi_n : G \to L$ des homomorphismes de monoïdes deux à deux
distincts, où $L$ est considéré comme un monoïde pour la multiplication.
Alors $\chi_1, \ldots, \chi_n$ sont linéairement indépendants sur $L$ :
si $\lambda_1, \ldots, \lambda_n \in L$ ne sont pas tous nuls, alors
$\sum \lambda_i\chi_i(g) \not = 0$ pour un certain $g \in G$.
\end{lemma}

\begin{proof}
Si $n = 1$, cela résulte de $\chi_1(e) = 1$, où $e \in G$ est l'élément
neutre (identité). Nous démontrons le résultat par récurrence pour $n > 1$.
Supposons que $\lambda_1, \ldots, \lambda_n \in L$ ne soient pas tous nuls.
Si $\lambda_i = 0$ pour un certain $i$, le résultat découle de la récurrence sur $n$.
Puisque nous voulons montrer que $\sum \lambda_i\chi_i(g) \not = 0$
pour un certain $g \in G$, nous pouvons, après division par $-\lambda_n$,
supposer que $\lambda_n = -1$. La seule difficulté possible
surviendrait si
$$
\chi_n(g) = \sum\nolimits_{i = 1, \ldots, n - 1} \lambda_i\chi_i(g)
$$
pour tout $g \in G$. Fixons $h \in G$. Nous aurions alors également
\begin{align*}
\chi_n(h)\chi_n(g) & = \chi_n(hg) \\
& = \sum\nolimits_{i = 1, \ldots, n - 1} \lambda_i\chi_i(hg) \\
& = \sum\nolimits_{i = 1, \ldots, n - 1} \lambda_i\chi_i(h) \chi_i(g)
\end{align*}
En multipliant la relation précédente par $\chi_n(h)$ puis en soustrayant, on obtient
$$
0 = \sum\nolimits_{i = 1, \ldots, n - 1}
\lambda_i (\chi_n(h) - \chi_i(h)) \chi_i(g)
$$
pour tout $g \in G$. Comme $\lambda_i \not = 0$, la récurrence donne
$\chi_n(h) = \chi_i(h)$ pour tout $i$.
Le choix de $h$ étant arbitraire, nous en déduisons
que $\chi_i = \chi_n$ pour $i \leq n - 1$, ce qui contredit
l'hypothèse selon laquelle les caractères $\chi_i$ sont deux à deux distincts.
\end{proof}

\begin{lemma}
\label{fields-lemma-sums-of-powers}
Soient $L$ un corps, $n \geq 1$ et $\alpha_1, \ldots, \alpha_n \in L$
des éléments deux à deux distincts de $L$. Alors il existe un
$e \geq 0$ tel que $\sum_{i = 1, \ldots, n} \alpha_i^e \not = 0$.
\end{lemma}

\begin{proof}
Appliquons l'indépendance linéaire des caractères
(Lemme \ref{fields-lemma-independence-characters})
aux homomorphismes de monoïdes $\mathbf{Z}_{\geq 0} \to L$,
$e \mapsto \alpha_i^e$.
\end{proof}

\begin{lemma}
\label{fields-lemma-independence-embeddings}
Soient $K/F$ et $L/F$ des extensions de corps. Soient
$\sigma_1, \ldots, \sigma_n : K \to L$ des morphismes de $F$-extensions
deux à deux distincts. Alors $\sigma_1, \ldots, \sigma_n$ sont linéairement
indépendants sur $L$ : si $\lambda_1, \ldots, \lambda_n \in L$ ne sont pas
tous nuls, alors $\sum \lambda_i\sigma_i(\alpha) \not = 0$
pour un certain $\alpha \in K$.
\end{lemma}

\begin{proof}
Appliquons le Lemme \ref{fields-lemma-independence-characters} aux restrictions
des $\sigma_i$ aux groupes des unités.
\end{proof}

\begin{lemma}
\label{fields-lemma-finite-separable-tensor-alg-closed}
Soient $K/F$ et $L/F$ des extensions de corps, avec
$K/F$ finie séparable et $L$ algébriquement clos.
Alors l'application
$$
K \otimes_F L
\longrightarrow
\prod\nolimits_{\sigma \in \Hom_F(K, L)} L,\quad
\alpha \otimes \beta \mapsto (\sigma(\alpha)\beta)_\sigma
$$
est un isomorphisme de $L$-algèbres.
\end{lemma}

\begin{proof}
Choisissons une base $\alpha_1, \ldots, \alpha_n$ de $K$ comme espace vectoriel sur $F$.
D'après le Lemme \ref{fields-lemma-separable-equality} (et un petit argument omis), l'ensemble
$\Hom_F(K, L)$ possède $n$ éléments, disons $\sigma_1, \ldots, \sigma_n$.
En particulier, les deux membres ont la même dimension $n$ comme espaces
vectoriels sur $L$. Ainsi, si l'application n'est pas un isomorphisme, son
noyau est non nul. Autrement dit, il existerait des
$\mu_j \in L$, $j = 1, \ldots, n$, non tous nuls,
tels que $\sum \alpha_j \otimes \mu_j$ appartienne au noyau.
Autrement dit, $\sum \sigma_i(\alpha_j)\mu_j = 0$ pour tout $i$.
Cela signifierait que la matrice $n \times n$ de coefficients
$\sigma_i(\alpha_j)$ n'est pas inversible. Nous pouvons donc trouver
$\lambda_1, \ldots, \lambda_n \in L$, non tous nuls,
tels que $\sum \lambda_i\sigma_i(\alpha_j) = 0$ pour tout $j$.
Or tout élément $\alpha \in K$ peut s'écrire
$\alpha = \sum \beta_j \alpha_j$ avec $\beta_j \in F$, et nous obtiendrions
$$
\sum \lambda_i\sigma_i(\alpha) =
\sum \lambda_i\sigma_i(\sum \beta_j \alpha_j) =
\sum \beta_j \sum \lambda_i\sigma_i(\alpha_j) = 0
$$
ce qui contredit le Lemme \ref{fields-lemma-independence-embeddings}.
\end{proof}







\section{Extensions purement inséparables}
\label{fields-section-purely-inseparable}

\noindent
Les extensions purement inséparables sont à l'opposé des extensions
séparables définies à la section précédente. Elles n'apparaissent
qu'en caractéristique positive.

\begin{definition}
\label{fields-definition-purely-inseparable}
Soit $F$ un corps de caractéristique $p > 0$. Soit $K/F$ une extension.
\begin{enumerate}
\item Un élément $\alpha \in K$ est dit {\it purement inséparable} sur $F$
s'il existe une puissance $q$ de $p$ telle que $\alpha^q \in F$.
\item L'extension $K/F$ est dite {\it purement inséparable}
si et seulement si tout élément de $K$ est purement inséparable sur $F$.
\end{enumerate}
Si $L/M$ est une extension de corps (sans condition sur la caractéristique
de $M$), nous dirons que l'extension est {\it purement inséparable}
si $L = M$, ou si la caractéristique de $M$ est un nombre premier $p$
et que $L/M$ est purement inséparable au sens défini ci-dessus.
\end{definition}

\noindent
Observons qu'une extension purement inséparable est nécessairement algébrique.
Soit $F$ un corps de caractéristique $p > 0$.
On obtient un exemple d'extension purement inséparable en adjoignant
la racine $p$-ième d'un élément $t \in F$ qui n'en possède pas encore. En effet,
le lemme ci-dessous montre que $P = x^p - t$ est irréductible, et donc
$$
K = F[x]/(P) = F[t^{1/p}]
$$
est un corps. De plus, $K$ est purement inséparable sur $F$, car tout élément
$$
a_0 + a_1t^{1/p} + \ldots + a_{p - 1}t^{(p - 1)/p},\quad a_i \in F
$$
de $K$ a pour puissance $p$-ième
$$
(a_0 + a_1t^{1/p} + \ldots + a_{p - 1}t^{(p - 1)/p})^p =
a_0^p + a_1^p t + \ldots + a_{p - 1}^pt^{p - 1} \in F
$$
Cette situation se présente pour le corps
$\mathbf{F}_p(t)$ des fonctions rationnelles sur $\mathbf{F}_p$.

\begin{lemma}
\label{fields-lemma-take-pth-root}
Soit $p$ un nombre premier. Soit $F$ un corps de caractéristique $p$.
Soit $t \in F$ un élément qui n'a pas de racine $p$-ième dans $F$.
Alors le polynôme $x^p - t$ est irréductible sur $F$.
\end{lemma}

\begin{proof}
Pour le voir, supposons donnée une factorisation
$x^p - t = f g$. En dérivant, on obtient $f' g + f g' = 0$.
Notons que ni $f' = 0$ ni $g' = 0$, puisque les degrés de $f$ et de $g$
sont strictement inférieurs à $p$. En outre, $\deg(f') < \deg(f)$ et $\deg(g') < \deg(g)$.
Nous en déduisons que $f$ et $g$ ont un facteur commun. Ainsi, si $x^p - t$
est réductible, il est de la forme $x^p - t = c f^n$ pour un certain polynôme
irréductible $f$, un certain $c \in F^*$ et un certain $n > 1$. Comme $p$ est premier,
cela entraîne $n = p$ et $f$ linéaire, ce qui impliquerait que $x^p - t$ possède
une racine dans $F$. Contradiction.
\end{proof}

\noindent
Nous verrons que l'extraction de racines $p$-ièmes est une opération très importante
en caractéristique $p$.

\begin{lemma}
\label{fields-lemma-purely-inseparable-permanence}
Soient $E/k$ et $F/E$ des extensions de corps purement inséparables. Alors $F/k$
est une extension de corps purement inséparable.
\end{lemma}

\begin{proof}
Si la caractéristique est nulle, les deux extensions sont triviales et il n'y a
rien à démontrer. Supposons donc que la caractéristique de $k$ soit un nombre
premier $p$. Choisissons $\alpha \in F$. Alors
$\alpha^q \in E$ pour une certaine puissance de $p$, notée $q$. Dès lors, $(\alpha^q)^{q'} \in k$
pour une certaine puissance de $p$, notée $q'$. Par conséquent, $\alpha^{qq'} \in k$.
\end{proof}

\begin{lemma}
\label{fields-lemma-purely-inseparable-elements}
Soit $E/k$ une extension de corps. Alors les éléments de $E$ purement inséparables
sur $k$ forment une sous-extension de $E/k$.
\end{lemma}

\begin{proof}
Soit $p$ la caractéristique de $k$.
Si celle-ci est nulle, l'assertion est immédiate, car tout élément purement
inséparable appartient alors au corps de base. Supposons donc cette caractéristique
positive.
Soient $\alpha, \beta \in E$ purement inséparables sur $k$. Écrivons
$\alpha^q \in k$ et $\beta^{q'} \in k$ pour certaines puissances de $p$, notées $q, q'$.
Si $q''$ est une puissance de $p$, alors
$(\alpha + \beta)^{q''} = \alpha^{q''} + \beta^{q''}$.
Par conséquent, si $q'' \geq q, q'$, nous en déduisons que $\alpha + \beta$
est purement inséparable sur $k$. Il en va de même de la différence,
du produit et du quotient de $\alpha$ et $\beta$.
\end{proof}

\begin{lemma}
\label{fields-lemma-finite-purely-inseparable}
Soit $E/F$ une extension de corps finie purement inséparable de
caractéristique $p > 0$. Alors il existe une suite d'éléments
$\alpha_1, \ldots, \alpha_n \in E$ qui donne une tour
de corps
$$
E = F(\alpha_1, \ldots, \alpha_n) \supset
F(\alpha_1, \ldots, \alpha_{n - 1}) \supset
\ldots
\supset F(\alpha_1) \supset F
$$
telle que chaque extension intermédiaire soit de degré $p$ et s'obtienne
par adjonction d'une racine $p$-ième. Plus précisément,
$\alpha_i^p \in F(\alpha_1, \ldots, \alpha_{i - 1})$
est un élément qui ne possède pas de racine $p$-ième dans
$F(\alpha_1, \ldots, \alpha_{i - 1})$, pour $i = 1, \ldots, n$.
\end{lemma}

\begin{proof}
Raisonnons par récurrence sur le degré de $E/F$. Si ce degré vaut $1$,
le résultat est clair (avec $n = 0$). Sinon, choisissons
$\alpha \in E$, $\alpha \not \in F$. Écrivons $\alpha^{p^r} \in F$ pour un certain
$r > 0$. Prenons $r$ minimal et remplaçons $\alpha$ par $\alpha^{p^{r - 1}}$.
Alors $\alpha \not \in F$, mais $\alpha^p \in F$. Ainsi $t = \alpha^p$ n'est pas
une puissance $p$-ième dans $F$ (car cela entraînerait $\alpha \in F$, voir le
Lemme \ref{fields-lemma-nr-roots-unchanged} ou sa démonstration).
Ainsi, $F \subset F(\alpha)$ est une sous-extension de degré $p$
(Lemme \ref{fields-lemma-take-pth-root}). Par récurrence, nous trouvons
$\alpha_1, \ldots, \alpha_n \in E$ qui engendrent $E/F(\alpha)$
et satisfont aux conclusions du lemme.
La suite $\alpha, \alpha_1, \ldots, \alpha_n$ convient alors
pour l'extension $E/F$.
\end{proof}

\begin{lemma}
\label{fields-lemma-separable-first}
\begin{slogan}
Toute extension algébrique de corps est, de manière unique, une extension
séparable suivie d'une extension purement inséparable.
\end{slogan}
Soit $E/F$ une extension algébrique de corps. Il existe une unique sous-extension
$E/E_{sep}/F$ telle que $E_{sep}/F$ soit séparable et que $E/E_{sep}$ soit
purement inséparable.
\end{lemma}

\begin{proof}
Si la caractéristique est nulle, posons $E_{sep} = E$. Supposons désormais la caractéristique
égale à $p > 0$. Soit $E_{sep}$ l'ensemble des éléments de $E$ qui sont séparables
sur $F$. C'est une sous-extension d'après le Lemme \ref{fields-lemma-separable-elements},
et bien entendu $E_{sep}$ est séparable sur $F$. Étant donné $\alpha$ dans $E$,
il existe une puissance de $p$, notée $q$, telle que $\alpha^q$ soit séparable sur $F$.
En effet, $q$ est la puissance de $p$ telle que le polynôme minimal de
$\alpha$ soit de la forme $P(x^q)$, où $P$ est un polynôme séparable, voir le
Lemme \ref{fields-lemma-irreducible-polynomials}. Ainsi $E/E_{sep}$ est purement
inséparable. L'unicité est claire.
\end{proof}

\begin{definition}
\label{fields-definition-insep-degree}
Soit $E/F$ une extension algébrique de corps. Soit $E_{sep}$ la sous-extension
obtenue au Lemme \ref{fields-lemma-separable-first}.
\begin{enumerate}
\item Le degré $[E_{sep} : F]$ est appelé le {\it degré séparable}
de l'extension. Notation : $[E : F]_s$.
\item Le degré $[E : E_{sep}]$ est appelé le {\it degré inséparable},
ou le {\it degré d'inséparabilité} de l'extension.
Notation : $[E : F]_i$.
\end{enumerate}
\end{definition}

\noindent
Bien entendu, en caractéristique $0$, on a $[E : F] = [E : F]_s$ et
$[E : F]_i = 1$. Par multiplicativité
(Lemme \ref{fields-lemma-multiplicativity-degrees}), on a
$$
[E : F] = [E : F]_s [E : F]_i
$$
même lorsque certains de ces degrés sont infinis. En fait, le degré séparable
et le degré inséparable sont eux aussi multiplicatifs (voir le
Lemme \ref{fields-lemma-multiplicativity-all-degrees}).

\begin{lemma}
\label{fields-lemma-separable-degree}
Soit $K/F$ une extension finie. Soit $\overline{F}$ une clôture algébrique
de $F$. Alors $[K : F]_s = |\Mor_F(K, \overline{F})|$.
\end{lemma}

\begin{proof}
Démontrons d'abord ce résultat lorsque $K/F$ est purement inséparable. Nous affirmons
que, dans ce cas, il existe un unique $F$-plongement $K \to \overline{F}$. On peut le
voir en choisissant une suite d'éléments $\alpha_1, \ldots, \alpha_n \in K$
comme au Lemme \ref{fields-lemma-finite-purely-inseparable}. Le polynôme irréductible
de $\alpha_i$ sur $F(\alpha_1, \ldots, \alpha_{i - 1})$ est $x^p - \alpha_i^p$.
En appliquant le Lemme \ref{fields-lemma-count-embeddings-explicitly}, on voit que
$|\Mor_F(K, \overline{F})| = 1$. D'autre part, $[K : F]_s = 1$
dans ce cas, d'où l'égalité.

\medskip\noindent
Revenons à une extension finie générale $K/F$. Dans ce cas,
choisissons $F \subset K_s \subset K$ comme au Lemme \ref{fields-lemma-separable-first}.
D'après le Lemme \ref{fields-lemma-separable-equality}, on a
$|\Mor_F(K_s, \overline{F})| = [K_s : F] = [K : F]_s$.
D'autre part, tout morphisme de corps $\sigma' : K_s \to \overline{F}$
se prolonge de manière unique en un morphisme de corps $\sigma : K \to \overline{F}$, d'après le
résultat du paragraphe précédent. Autrement dit,
$|\Mor_F(K, \overline{F})| = |\Mor_F(K_s, \overline{F})|$,
ce qui achève la démonstration.
\end{proof}

\begin{lemma}[Multiplicativité]
\label{fields-lemma-multiplicativity-all-degrees}
Supposons donnée une tour d'extensions algébriques de corps $K/E/F$. Alors
$$
[K : F]_s = [K : E]_s [E : F]_s
\quad\text{et}\quad
[K : F]_i = [K : E]_i [E : F]_i
$$
\end{lemma}

\begin{proof}
Démontrons d'abord le résultat lorsque $K$ est fini sur $F$. Comme la
multiplicativité vaut pour le degré usuel (d'après le
Lemme \ref{fields-lemma-multiplicativity-degrees}), il suffit de démontrer
l'une des deux formules. D'après le Lemme \ref{fields-lemma-separable-degree}, on a
$[K : F]_s = |\Mor_F(K, \overline{F})|$. Par le même lemme,
pour tout $\sigma \in \Mor_F(E, \overline{F})$, le nombre de prolongements
de $\sigma$ en une application $\tau : K \to \overline{F}$ est $[K : E]_s$.
En effet, via $E \cong \sigma(E) \subset \overline{F}$, nous pouvons considérer
$\overline{F}$ comme une clôture algébrique de $E$. En combinant ceci avec le
fait qu'il existe $[E : F]_s = |\Mor_F(E, \overline{F})|$ choix
pour $\sigma$, on obtient le résultat.

\medskip\noindent
Nous omettons la démonstration lorsque les extensions sont infinies.
\end{proof}









\section{Extensions normales}
\label{fields-section-normal}

\noindent
Soit $P \in F[x]$ un polynôme non constant sur un corps $F$. Nous dirons que $P$
{\it se décompose complètement en facteurs linéaires sur $F$} ou
{\it est complètement décomposé sur $F$} s'il existe
$c \in F^*$, $n \geq 1$, $\alpha_1, \ldots, \alpha_n \in F$ tels que
$$
P = c(x - \alpha_1) \ldots (x - \alpha_n)
$$
dans $F[x]$. Les extensions normales sont définies comme suit.

\begin{definition}
\label{fields-definition-normal}
Soit $E/F$ une extension algébrique de corps. Nous dirons que $E$ est {\it normale}
sur $F$ si, pour tout $\alpha \in E$, le polynôme minimal $P$
de $\alpha$ sur $F$ se décompose complètement en facteurs linéaires sur $E$.
\end{definition}

\noindent
Comme pour les extensions séparables, l'établissement des propriétés
fondamentales de cette notion demande un peu de travail.

\begin{lemma}
\label{fields-lemma-normal-goes-up}
Soit $K/E/F$ une tour d'extensions algébriques de corps.
Si $K$ est normale sur $F$, alors $K$ est normale sur $E$.
\end{lemma}

\begin{proof}
Soit $\alpha \in K$. Soit $P$ le polynôme minimal de $\alpha$ sur $F$.
Soit $Q$ le polynôme minimal de $\alpha$ sur $E$.
Alors $Q$ divise $P$ dans l'anneau de polynômes $E[x]$, disons $P = QR$.
Par conséquent, si $P$ est complètement décomposé sur $K$, il en va de même de $Q$.
\end{proof}

\begin{lemma}
\label{fields-lemma-intersect-normal}
Soit $F$ un corps. Soit $M/F$ une extension algébrique. Soient
$M/E_i/F$, $i \in I$, une famille non vide de sous-extensions telles que
$E_i/F$ soit normale. Alors $\bigcap E_i$ est normale sur $F$.
\end{lemma}

\begin{proof}
Cela résulte directement des définitions.
\end{proof}

\begin{lemma}
\label{fields-lemma-separable-first-normal}
Soit $E/F$ une extension algébrique de corps normale. Alors la sous-extension
$E/E_{sep}/F$ du Lemme \ref{fields-lemma-separable-first} est normale.
\end{lemma}

\begin{proof}
Si la caractéristique est nulle, alors $E_{sep} = E$, et le résultat
est clair. Si la caractéristique vaut $p > 0$, alors $E_{sep}$
est l'ensemble des éléments de $E$ qui sont séparables sur $F$.
Alors, si $\alpha \in E_{sep}$ a pour polynôme minimal $P$,
écrivons $P = c(x - \alpha)(x - \alpha_2) \ldots (x - \alpha_d)$
avec $\alpha_2, \ldots, \alpha_d \in E$. Puisque
$P$ est un polynôme séparable et que $\alpha_i$
est une racine de $P$, nous en déduisons $\alpha_i \in E_{sep}$, comme voulu.
\end{proof}

\begin{lemma}
\label{fields-lemma-characterize-normal}
Soit $E/F$ une extension algébrique de corps. Soit $\overline{F}$ une
clôture algébrique de $F$. Les assertions suivantes sont équivalentes :
\begin{enumerate}
\item $E$ est normale sur $F$ ;
\item pour toute paire $\sigma, \sigma' \in \Mor_F(E, \overline{F})$, on
a $\sigma(E) = \sigma'(E)$.
\end{enumerate}
\end{lemma}

\begin{proof}
Soit $\mathcal{P}$ l'ensemble de tous les polynômes minimaux sur $F$
de tous les éléments de $E$. Posons
$$
T =
\{\beta \in \overline{F} \mid P(\beta) = 0\text{ pour un certain }P \in \mathcal{P}\}
$$
Il est clair que, si $E$ est normale sur $F$, alors $\sigma(E) = T$
pour tout $\sigma \in \Mor_F(E, \overline{F})$. Nous voyons donc que (1)
implique (2).

\medskip\noindent
Réciproquement, supposons (2). Choisissons $\beta \in T$.
Nous pouvons trouver $\alpha \in E$ dont le polynôme minimal
$P \in \mathcal{P}$ annule $\beta$. Puisque $F(\alpha) = F[x]/(P)$,
nous pouvons trouver un élément $\sigma_0 \in \Mor_F(F(\alpha), \overline{F})$ qui envoie
$\alpha$ sur $\beta$. D'après le Lemme \ref{fields-lemma-map-into-algebraic-closure},
nous pouvons prolonger $\sigma_0$ en un morphisme $\sigma \in \Mor_F(E, \overline{F})$.
Nous voyons ainsi que $\beta$ appartient à l'image commune de tous les plongements
$\sigma : E \to \overline{F}$. Il s'ensuit que $\sigma(E) = T$
pour tout $\sigma$. Fixons un $\sigma$. Soit maintenant $P \in \mathcal{P}$. Alors nous
pouvons écrire
$$
P = (x - \beta_1) \ldots (x - \beta_n)
$$
pour un certain $n$ et des $\beta_i \in \overline{F}$, d'après le
Lemme \ref{fields-lemma-algebraically-closed}. Observons que $\beta_i \in T$.
Ainsi $\beta_i = \sigma(\alpha_i)$ pour un certain $\alpha_i \in E$. Par conséquent,
$P = (x - \alpha_1) \ldots (x - \alpha_n)$ se décompose complètement sur $E$.
Cela achève la démonstration.
\end{proof}

\begin{lemma}
\label{fields-lemma-normally-generated}
Soit $E/F$ une extension algébrique de corps.
Si $E$ est engendrée par les éléments $\alpha_i \in E$, $i \in I$,
sur $F$, et si, pour tout $i$, le polynôme minimal
de $\alpha_i$ sur $F$ se décompose complètement dans $E$, alors
$E/F$ est normale.
\end{lemma}

\begin{proof}
Soit $P_i$ le polynôme minimal de $\alpha_i$ sur $F$.
Soient $\alpha_i = \alpha_{i, 1}, \alpha_{i, 2}, \ldots, \alpha_{i, d_i}$
les racines de $P_i$ sur $E$. Étant donnés deux plongements
$\sigma, \sigma' : E \to \overline{F}$ au-dessus de $F$, on voit que
$$
\{\sigma(\alpha_{i, 1}), \ldots, \sigma(\alpha_{i, d_i})\} =
\{\sigma'(\alpha_{i, 1}), \ldots, \sigma'(\alpha_{i, d_i})\}
$$
car les deux membres sont égaux à l'ensemble des racines de $P_i$
dans $\overline{F}$. Les éléments $\alpha_{i, j}$
engendrent $E$ sur $F$, et l'on obtient $\sigma(E) = \sigma'(E)$.
Ainsi $E/F$ est normale d'après le Lemme \ref{fields-lemma-characterize-normal}.
\end{proof}

\begin{lemma}
\label{fields-lemma-lift-maps}
Soit $L/M/K$ une tour d'extensions algébriques.
\begin{enumerate}
\item Si $M/K$ est normale, alors tout automorphisme $\tau$ de $L/K$
induit un automorphisme $\tau|_M : M \to M$.
\item Si $L/K$ est normale, alors tout morphisme de $K$-algèbres $\sigma : M \to L$
se prolonge en un automorphisme de $L$.
\end{enumerate}
\end{lemma}

\begin{proof}
Choisissons une clôture algébrique $\overline{L}$ de $L$
(Théorème \ref{fields-theorem-existence-algebraic-closure}).

\medskip\noindent
Soit $\tau$ comme en (1). Alors $\tau(M) = M$ comme sous-corps de $\overline{L}$
d'après le Lemme \ref{fields-lemma-characterize-normal}, et donc
$\tau|_M : M \to M$ est un automorphisme.

\medskip\noindent
Soit $\sigma : M \to L$ comme en (2).
D'après le Lemme \ref{fields-lemma-map-into-algebraic-closure},
nous pouvons prolonger $\sigma$ en une application
$\tau : L \to \overline{L}$, c'est-à-dire de telle sorte que le diagramme
$$
\xymatrix{
L \ar[r]_\tau & \overline{L} \\
M \ar[u] \ar[ru]_\sigma & K \ar[l] \ar[u]
}
$$
soit commutatif. D'après le Lemme \ref{fields-lemma-characterize-normal}, on voit que
$\tau(L) = L$. Ainsi $\tau : L \to L$ est un automorphisme qui
prolonge $\sigma$.
\end{proof}

\begin{definition}
\label{fields-definition-automorphisms}
Soit $E/F$ une extension de corps. Alors $\text{Aut}(E/F)$ ou
$\text{Aut}_F(E)$ désigne le groupe des automorphismes de $E$ comme objet
de la catégorie des $F$-extensions. Les éléments de $\text{Aut}(E/F)$
sont appelés {\it automorphismes de $E$ sur $F$} ou
{\it automorphismes de $E/F$}.
\end{definition}

\noindent
Voici une caractérisation des extensions normales en termes d'automorphismes.

\begin{lemma}
\label{fields-lemma-normal-and-automorphisms}
Soit $E/F$ une extension finie. On a
$$
|\text{Aut}(E/F)| \leq [E : F]_s
$$
avec égalité si et seulement si $E$ est normale sur $F$.
\end{lemma}

\begin{proof}
Choisissons une clôture algébrique $\overline{F}$ de $F$. Rappelons que
$[E : F]_s = |\Mor_F(E, \overline{F})|$. Choisissons un élément
$\sigma_0 \in \Mor_F(E, \overline{F})$. Alors l'application
$$
\text{Aut}(E/F) \longrightarrow \Mor_F(E, \overline{F}),\quad
\tau \longmapsto \sigma_0 \circ \tau
$$
est injective. D'où l'inégalité. S'il y a égalité, tout
$\sigma \in \Mor_F(E, \overline{F})$ s'obtient en précomposant
$\sigma_0$ par un automorphisme. Ainsi $\sigma(E) = \sigma_0(E)$.
Donc $E$ est normale sur $F$ d'après le Lemme \ref{fields-lemma-characterize-normal}.

\medskip\noindent
Réciproquement, supposons que $E/F$ est normale. Alors, d'après le
Lemme \ref{fields-lemma-characterize-normal}, on a $\sigma(E) = \sigma_0(E)$
pour tout $\sigma \in \Mor_F(E, \overline{F})$.
On obtient donc un automorphisme de $E$ sur $F$ en posant
$\tau = \sigma_0^{-1} \circ \sigma$. Par conséquent, l'application ci-dessus
est surjective.
\end{proof}

\begin{lemma}
\label{fields-lemma-normal-embeddings-differ-by-aut}
Soit $L/K$ une extension algébrique normale de corps.
Soit $E/K$ une extension de corps. Alors, soit
il n'existe aucun $K$-plongement de $L$ dans $E$, soit
il en existe un $\tau : L \to E$ et tout autre
est de la forme $\tau \circ \sigma$, où $\sigma \in \text{Aut}(L/K)$.
\end{lemma}

\begin{proof}
Étant donné $\tau$, remplaçons $L$ par $\tau(L) \subset E$ et appliquons le
Lemme \ref{fields-lemma-lift-maps}.
\end{proof}





\section{Corps de décomposition}
\label{fields-section-splitting-fieds}

\noindent
Le lemme suivant est un outil utile pour construire des extensions normales de corps.

\begin{lemma}
\label{fields-lemma-splitting-field}
Soit $F$ un corps. Soit $P \in F[x]$ un polynôme non constant.
Il existe une plus petite extension de corps $E/F$ telle que $P$
soit scindé sur $E$. De plus, l'extension de corps $E/F$ est normale
et unique à isomorphisme près, cet isomorphisme n'étant pas nécessairement unique.
\end{lemma}

\begin{proof}
Choisissons une clôture algébrique $\overline{F}$. Nous pouvons alors écrire
$P = c (x - \beta_1) \ldots (x - \beta_n)$ dans $\overline{F}[x]$, voir le
Lemme \ref{fields-lemma-algebraically-closed}. Notons que $c \in F^*$. Posons
$E = F(\beta_1, \ldots, \beta_n)$. Il est alors clair que $E$ est
minimale sous la condition que $P$ soit scindé sur $E$.

\medskip\noindent
Soit ensuite $E'$ une autre extension minimale de $F$ telle que
$P$ soit scindé sur $E'$. Écrivons
$P = c (x - \alpha_1) \ldots (x - \alpha_n)$ avec $c \in F$ et
$\alpha_i \in E'$. La minimalité entraîne de nouveau que
$E' = F(\alpha_1, \ldots, \alpha_n)$. De plus, si nous choisissons
un $\sigma : E' \to \overline{F}$ quelconque
(Lemme \ref{fields-lemma-map-into-algebraic-closure}),
nous voyons immédiatement que $\sigma(\alpha_i) = \beta_{\tau(i)}$
pour une certaine permutation $\tau : \{1, \ldots, n\} \to \{1, \ldots, n\}$.
Ainsi $\sigma(E') = E$. Cela implique que $E'$ est une extension normale
de $F$ d'après le Lemme \ref{fields-lemma-characterize-normal},
et que $E \cong E'$ comme extensions de $F$, ce qui achève la démonstration.
\end{proof}

\begin{definition}
\label{fields-definition-splitting-field}
Soit $F$ un corps. Soit $P \in F[x]$ un polynôme non constant.
L'extension de corps $E/F$ construite au Lemme \ref{fields-lemma-splitting-field}
est appelée le {\it corps de décomposition de $P$ sur $F$}.
\end{definition}

\begin{lemma}
\label{fields-lemma-normal-closure}
\begin{slogan}
Existence de la clôture normale des extensions finies de corps.
\end{slogan}
Soit $E/F$ une extension finie de corps. Il existe une unique
plus petite extension finie $K/E$ telle que $K$ soit normale sur $F$.
\end{lemma}

\begin{proof}
Choisissons des générateurs $\alpha_1, \ldots, \alpha_n$ de $E$ sur $F$.
Soient $P_1, \ldots, P_n$ les polynômes minimaux de
$\alpha_1, \ldots, \alpha_n$ sur $F$. Posons $P = P_1 \ldots P_n$.
Observons que $(x - \alpha_1) \ldots (x - \alpha_n)$ divise $P$, puisque
chaque $(x - \alpha_i)$ divise $P_i$. Écrivons
$P = (x - \alpha_1) \ldots (x - \alpha_n)Q$.
Soit $K/E$ le corps de décomposition de $P$ sur $E$.
Nous affirmons que $K$ est également le corps de décomposition de $P$ sur $F$
(ce qui implique que $K$ est normale sur $F$).
Cela est clair, car $K/E$ est engendrée par les racines de
$Q$ sur $E$, tandis que $E$ est engendrée par les racines de
$(x - \alpha_1) \ldots (x - \alpha_n)$ sur $F$ ; ainsi,
$K$ est engendrée par les racines de $P$ sur $F$.

\medskip\noindent
Unicité. Supposons que $K'/E$ soit une seconde extension minimale telle que
$K'/F$ soit normale. Choisissons une clôture algébrique $\overline{F}$ et un
plongement $\sigma_0 : E \to \overline{F}$. D'après le
Lemme \ref{fields-lemma-map-into-algebraic-closure},
nous pouvons prolonger $\sigma_0$ en $\sigma : K \to \overline{F}$ et en
$\sigma' : K' \to \overline{F}$.
D'après le Lemme \ref{fields-lemma-intersect-normal}, on voit que
$\sigma(K) \cap \sigma'(K')$ est normale sur $F$.
La minimalité donne $\sigma(K) = \sigma'(K')$.
Ainsi $\sigma^{-1} \circ \sigma' : K' \to K$ fournit un isomorphisme
d'extensions de $E$.
\end{proof}

\begin{definition}
\label{fields-definition-normal-closure}
Soit $E/F$ une extension finie de corps. L'extension de corps $K/E$
construite au Lemme \ref{fields-lemma-normal-closure}
est appelée la {\it clôture normale de $E$ sur $F$}.
\end{definition}

\noindent
On peut construire la clôture normale à l'intérieur de toute extension normale donnée.

\begin{lemma}
\label{fields-lemma-normal-closure-inside-normal}
Soit $L/K$ une extension algébrique normale.
\begin{enumerate}
\item Si $L/M/K$ est une sous-extension avec $M/K$ finie, alors il existe
une tour $L/M'/M/K$ telle que $M'/K$ soit finie et normale.
\item Si $L/M'/M/K$ est une tour telle que $M/K$ soit normale et $M'/M$ finie,
alors il existe une tour $L/M''/M'/M/K$ telle que $M''/M$ soit
finie et $M''/K$ normale.
\end{enumerate}
\end{lemma}

\begin{proof}
Démonstration de (1). Soit $M'$ la plus petite sous-extension de $L/K$ contenant $M$
qui soit normale sur $K$. D'après le Lemme \ref{fields-lemma-normal-closure},
c'est la clôture normale de $M/K$, et elle est finie sur $K$.

\medskip\noindent
Démonstration de (2). Soient $\alpha_1, \ldots, \alpha_n \in M'$ des générateurs de $M'$ sur $M$.
Soient $P_1, \ldots, P_n$ les polynômes minimaux de
$\alpha_1, \ldots, \alpha_n$ sur $K$. Soient $\alpha_{i, j}$ les racines
de $P_i$ dans $L$. Posons $M''  = M(\alpha_{i, j})$. Il résulte du
Lemme \ref{fields-lemma-normally-generated}
(appliqué à l'ensemble de générateurs $M \cup \{\alpha_{i, j}\}$)
que $M''$ est normale sur $K$.
\end{proof}

\noindent
Le lemme suivant peut parfois servir à démontrer des propriétés
de la clôture normale.

\begin{lemma}
\label{fields-lemma-normal-closure-tensor-product}
Soit $L/K$ une extension finie. Soit $M/L$ la clôture
normale de $L$ sur $K$. Alors il existe une application surjective
$$
L \otimes_K L \otimes_K \ldots \otimes_K L \longrightarrow M
$$
de $K$-algèbres, où le nombre de facteurs tensoriels peut être pris égal à
$[L : K]_s \leq [L : K]$.
\end{lemma}

\begin{proof}
Choisissons une clôture algébrique $\overline{K}$ de $K$.
Posons $n = [L : K]_s = |\Mor_K(L, \overline{K})|$, l'égalité résultant du
Lemme \ref{fields-lemma-separable-degree}.
Écrivons $\Mor_K(L, \overline{K}) = \{\sigma_1, \ldots, \sigma_n\}$.
Soit $M' \subset \overline{K}$ la $K$-sous-algèbre
engendrée par $\sigma_i(L)$, $i = 1, \ldots, n$. Alors $M'$ est
un corps, car toute $K$-sous-algèbre de $\overline{K}$ est un corps.
Tout morphisme de $K$-algèbres de $M'$ vers $\overline{K}$ permute les $\sigma_i$,
et envoie donc le sous-corps $M'$ sur $M'$. Par construction, le corps $M'$
est engendré par les conjugués d'éléments de $\sigma_1(L)$. Cela étant,
il résulte du Lemme \ref{fields-lemma-characterize-normal}
que $M'$ est normale sur $K$ et qu'elle est la
plus petite sous-extension normale de $\overline{K}$ contenant
$\sigma_1(L)$. Par unicité de la clôture normale, on a $M \cong M'$.
Enfin, il existe une application surjective
$$
L \otimes_K L \otimes_K \ldots \otimes_K L \longrightarrow M',
\quad
\lambda_1 \otimes \ldots \otimes \lambda_n \longmapsto
\sigma_1(\lambda_1) \ldots \sigma_n(\lambda_n)
$$
et notons que $n \leq [L : K]$ par définition.
\end{proof}












\section{Racines de l'unité}
\label{fields-section-roots-of-1}

\noindent
Soit $F$ un corps. Pour un entier $n \geq 1$, posons
$$
\mu_n(F) = \{\zeta \in F \mid \zeta^n = 1\}
$$
On appelle cet ensemble le {\it groupe des racines $n$-ièmes de l'unité} ; ses
éléments sont les {\it racines $n$-ièmes de $1$}. C'est un groupe abélien pour la multiplication,
dont l'élément neutre est $1$.
Observons que, dans un corps, le nombre de racines d'un polynôme de degré $d$
est toujours au plus $d$. Ainsi $|\mu_n(F)| \leq n$,
puisque cet ensemble est défini par une équation polynomiale de degré $n$.
Bien entendu, tout élément de $\mu_n(F)$ est d'ordre divisant $n$.
En outre, les sous-groupes
$$
\mu_d(F) \subset \mu_n(F),\quad d | n
$$
ont chacun au plus $d$ éléments. Cela implique que $\mu_n(F)$ est cyclique.

\begin{lemma}
\label{fields-lemma-cyclic}
Soit $A$ un groupe abélien d'exposant divisant $n$ tel que
$\{x \in A \mid dx = 0\}$ soit de cardinal au plus $d$ pour tout $d | n$.
Alors $A$ est cyclique d'ordre divisant $n$.
\end{lemma}

\begin{proof}
Si $A$ est trivial, il n'y a rien à démontrer. Supposons désormais que $A$
ne soit pas trivial.
Les conditions impliquent que $|A| \leq n$ ; en particulier, $A$ est fini.
La structure des groupes abéliens finis montre que
$A = \mathbf{Z}/e_1\mathbf{Z} \oplus \ldots \oplus \mathbf{Z}/e_r\mathbf{Z}$
pour des entiers $1 < e_1 | e_2 | \ldots | e_r$. Cela impliquerait
que $\{x \in A \mid e_1 x = 0\}$ est de cardinal $e_1^r$. Par conséquent,
$r = 1$.
\end{proof}

\noindent
Appliqué au corps $\mathbf{F}_p$, ceci donne le célèbre résultat
selon lequel le groupe $(\mathbf{Z}/p\mathbf{Z})^*$ est cyclique. Nous reviendrons sur
cette question dans la section consacrée aux corps finis.

\medskip\noindent
Une autre observation est souvent utile : si $F$ est de caractéristique
$p > 0$, alors $\mu_{p^n}(F) = \{1\}$. Cela tient à ce que l'élévation
à la puissance $p$ est une application injective sur les corps de caractéristique $p$,
comme nous l'avons vu dans la démonstration du Lemme \ref{fields-lemma-nr-roots-unchanged}.
(Bien entendu, cela résulte aussi de l'énoncé même de ce lemme.)






\section{Corps finis}
\label{fields-section-finite}

\noindent
Soit $F$ un corps fini. Il est clair que $F$ est de caractéristique positive,
car il ne peut exister d'injection $\mathbf{Q} \to F$. Supposons que la caractéristique
de $F$ soit $p$. L'extension $\mathbf{F}_p \subset F$ est finie.
Nous voyons ainsi que $F$ possède $q = p^f$ éléments pour un certain $f \geq 1$.

\medskip\noindent
Considérons le groupe des unités $F^*$. C'est un groupe abélien
fini, il possède donc un certain exposant $e$. Alors $F^* = \mu_e(F)$, et la
discussion de la section \ref{fields-section-roots-of-1} montre que $F^*$
est un groupe cyclique d'ordre $q - 1$. (A posteriori, il en résulte aussi que
$e = q - 1$.) En particulier, si $\alpha \in F^*$ est un générateur,
alors il est clair que
$$
F = \mathbf{F}_p(\alpha)
$$
Autrement dit, l'extension $F/\mathbf{F}_p$ est engendrée par un seul
élément. Bien entendu, il en va de même de toute extension de corps finis
$E/F$ (car $E$ est déjà engendré par un seul élément sur
le corps premier).





\section{Éléments primitifs}
\label{fields-section-primitive-element}

\noindent
Soit $E/F$ une extension finie de corps. Un élément $\alpha \in E$
est appelé {\it élément primitif de $E$ sur $F$} si $E = F(\alpha)$.

\begin{lemma}[Élément primitif]
\label{fields-lemma-primitive-element}
Soit $E/F$ une extension finie de corps. Les assertions suivantes sont équivalentes :
\begin{enumerate}
\item il existe un élément primitif de $E$ sur $F$ ;
\item il existe un nombre fini de sous-extensions $E/K/F$.
\end{enumerate}
De plus, (1) et (2) sont vérifiées si $E/F$ est séparable.
\end{lemma}

\begin{proof}
Soit $\alpha \in E$ un élément primitif. Soit $P$ le polynôme minimal
de $\alpha$ sur $F$. Soit ensuite $E/K/F$ une sous-extension.
Soit $Q$ le polynôme minimal de $\alpha$ sur $K$. Observons que
$\deg(Q) = [E : K]$. En écrivant $Q = x^d + \sum_{i < d} a_i x^i$, nous affirmons que
$K$ est égal à $L = F(a_0, \ldots, a_{d - 1})$. En effet, $\alpha$ est de degré
$d$ sur $L$ et $L \subset K$. Ainsi $[E : L] = [E : K]$, d'où
$[K : L] = 1$, c'est-à-dire $K = L$.
Il suffit donc de montrer qu'il n'existe qu'un nombre fini de possibilités
pour le polynôme $Q$. Cela est clair, puisque nous avons une factorisation
$P = QR$  dans $K[x]$, et en particulier dans $E[x]$. Comme la factorisation
est unique dans $E[x]$, le polynôme $P$ n'a qu'un nombre fini de facteurs
unitaires dans $E[x]$.

\medskip\noindent
Si $F$ est un corps fini (ce qui équivaut à dire que $E$ est un corps fini), alors
$E/F$ possède un élément primitif d'après la discussion de la
section \ref{fields-section-finite}.
Supposons ensuite $F$ infini et qu'il n'existe qu'un nombre fini de sous-corps
propres $E/K/F$. Énumérons-les, disons $K_1, \ldots, K_N$. Alors
chaque $K_i \subset E$ est un sous-espace vectoriel propre sur $F$. Comme $F$ est infini,
nous pouvons trouver un vecteur $\alpha \in E$ tel que $\alpha \not \in K_i$ pour tout $i$
(un espace vectoriel ne peut jamais être la réunion d'un nombre fini de sous-espaces propres ;
nous omettons les détails). Alors $\alpha$ est un élément primitif de $E$ sur $F$.

\medskip\noindent
Après avoir établi l'équivalence de (1) et (2), démontrons maintenant
la dernière assertion du lemme. Choisissons une clôture algébrique
$\overline{F}$ de $F$. Énumérons les éléments
$\sigma_1, \ldots, \sigma_n \in \Mor_F(E, \overline{F})$.
Puisque $E/F$ est séparable, on a $n = [E : F]$ d'après le
Lemme \ref{fields-lemma-separable-equality}.
Notons que, si $i \not = j$, alors
$$
V_{ij} = \Ker(\sigma_i - \sigma_j : E \longrightarrow \overline{F})
$$
n'est pas égal à $E$. En raisonnant comme au paragraphe précédent,
nous pouvons donc trouver $\alpha \in E$ tel que $\alpha \not \in V_{ij}$ pour tous
$i \not = j$. Il s'ensuit que $|\Mor_F(F(\alpha), \overline{F})| \geq n$.
D'autre part, $[F(\alpha) : F] \leq [E : F]$. Il y a donc égalité
d'après le Lemme \ref{fields-lemma-separable-equality},
et nous concluons que $E = F(\alpha)$.
\end{proof}






\section{Trace et norme}
\label{fields-section-trace-pairing}

\noindent
Soit $L/K$ une extension finie de corps. D'après le
Lemme \ref{fields-lemma-vector-space-is-free},
nous pouvons choisir un isomorphisme $L \cong K^{\oplus n}$ de $K$-modules.
Bien entendu, $n = [L : K]$ est le degré de l'extension de corps.
À l'aide de cet isomorphisme, nous obtenons un morphisme de $K$-algèbres
$$
L \longrightarrow \text{Mat}(n \times n, K),\quad
\alpha \longmapsto \text{matrice de la multiplication par }\alpha
$$
Ainsi, étant donné $\alpha \in L$, nous pouvons prendre la trace et le déterminant
de la matrice correspondante. Bien entendu, ces quantités ne dépendent pas
du choix de la base effectué ci-dessus. Plus canoniquement, en considérant simplement
$L$ comme un espace vectoriel de dimension finie sur $K$, nous avons
$\text{Trace}_K(\alpha : L \to L)$ et le déterminant
$\det_K(\alpha : L \to L)$.

\begin{definition}
\label{fields-definition-trace-norm}
Soit $L/K$ une extension finie de corps. Pour $\alpha \in L$, nous définissons
la {\it trace}
$\text{Trace}_{L/K}(\alpha) = \text{Trace}_K(\alpha : L \to L)$
et la {\it norme}
$\text{Norm}_{L/K}(\alpha) = \det_K(\alpha : L \to L)$.
\end{definition}

\noindent
Il résulte clairement de la définition que
$\text{Trace}_{L/K}$ est $K$-linéaire et vérifie
$\text{Trace}_{L/K}(\alpha) = [L : K]\alpha$ pour $\alpha \in K$.
De même, $\text{Norm}_{L/K}$ est multiplicative et
$\text{Norm}_{L/K}(\alpha) = \alpha^{[L : K]}$ pour $\alpha \in K$.
C'est un cas particulier de la construction plus générale examinée
dans les Exercices \ref{exercises-exercise-trace-det} et
\ref{exercises-exercise-trace-det-rings}.

\begin{lemma}
\label{fields-lemma-characteristic-vs-minimal-polynomial}
Soit $L/K$ une extension finie de corps. Soit $\alpha \in L$, et soit $P$
le polynôme minimal de $\alpha$ sur $K$. Alors le polynôme caractéristique
de l'application $K$-linéaire $\alpha : L \to L$ est égal à
$P^e$, avec $e \deg(P) = [L : K]$.
\end{lemma}

\begin{proof}
Choisissons une base $\beta_1, \ldots, \beta_e$ de $L$ sur $K(\alpha)$.
Alors $e$ vérifie $e \deg(P) = [L : K]$ d'après les
Lemmes \ref{fields-lemma-degree-minimal-polynomial} et
\ref{fields-lemma-multiplicativity-degrees}.
Nous voyons alors que $L = \bigoplus K(\alpha) \beta_i$ est une
décomposition en somme directe de sous-espaces invariants par $\alpha$ ;
le polynôme caractéristique de $\alpha : L \to L$
est donc égal au polynôme caractéristique de
$\alpha : K(\alpha) \to K(\alpha)$ élevé à la puissance $e$.

\medskip\noindent
Pour achever la démonstration, nous pouvons supposer que $L = K(\alpha)$.
Dans ce cas, le théorème de Cayley-Hamilton montre que $\alpha$
est une racine du polynôme caractéristique. Et puisque le
polynôme caractéristique a le même degré que le polynôme
minimal, nous obtenons l'égalité.
\end{proof}

\begin{lemma}
\label{fields-lemma-trace-and-norm-from-minimal-polynomial}
Soit $L/K$ une extension finie de corps. Soit $\alpha \in L$, et soit
$P = x^d + a_1 x^{d - 1} + \ldots + a_d$
le polynôme minimal de $\alpha$ sur $K$. Alors
$$
\text{Norm}_{L/K}(\alpha) = (-1)^{[L : K]} a_d^e
\quad\text{et}\quad
\text{Trace}_{L/K}(\alpha) = - e a_1
$$
où $e d = [L : K]$.
\end{lemma}

\begin{proof}
Cela résulte immédiatement du Lemme \ref{fields-lemma-characteristic-vs-minimal-polynomial}
et des définitions.
\end{proof}

\begin{lemma}
\label{fields-lemma-trace-and-norm-linear}
Soit $L/K$ une extension finie de corps. Soit $V$ un espace vectoriel de dimension
finie sur $L$. Soit $\varphi : V \to V$ une application $L$-linéaire.
Alors
$$
\text{Trace}_K(\varphi : V \to V) =
\text{Trace}_{L/K}(\text{Trace}_L(\varphi : V \to V))
$$
et
$$
\det\nolimits_K(\varphi : V \to V) =
\text{Norm}_{L/K}(\det\nolimits_L(\varphi : V \to V))
$$
\end{lemma}

\begin{proof}
Choisissons un isomorphisme $V = L^{\oplus n}$ de sorte que $\varphi$ corresponde
à une matrice $n \times n$. Dans le cas des traces, les deux membres de la formule
sont additifs en $\varphi$. Nous pouvons donc supposer que $\varphi$
corresponde à la matrice dont l'unique coefficient non nul est en position $(i, j)$.
Dans ce cas, un calcul direct montre que les deux membres sont égaux.

\medskip\noindent
Dans le cas des normes, les deux membres sont nuls si $\varphi$ possède un noyau non nul.
Nous pouvons donc supposer que $\varphi$ corresponde à un élément de
$\text{GL}_n(L)$. Les deux membres de la formule sont multiplicatifs en $\varphi$.
Puisque tout élément de $\text{GL}_n(L)$ est un produit de matrices
élémentaires, nous pouvons supposer que $\varphi$ soit de la forme
$$
E_{12}(\lambda) =
\left(
\begin{matrix}
1 & \lambda & \ldots \\
0 & 1 & \ldots \\
\ldots & \ldots & \ldots
\end{matrix}
\right)
\quad\text{ou}\quad
E_1(a) =
\left(
\begin{matrix}
a & 0 & \ldots \\
0 & 1 & \ldots \\
\ldots & \ldots & \ldots
\end{matrix}
\right)
$$
(car nous pouvons également permuter les éléments de la base à notre convenance).
Dans les deux cas, la formule se vérifie aisément par un calcul direct.
\end{proof}

\begin{lemma}
\label{fields-lemma-trace-and-norm-tower}
Soit $M/L/K$ une tour d'extensions finies de corps. Alors
$$
\text{Trace}_{M/K} = \text{Trace}_{L/K} \circ \text{Trace}_{M/L}
\quad\text{et}\quad
\text{Norm}_{M/K} = \text{Norm}_{L/K} \circ \text{Norm}_{M/L}
$$
\end{lemma}

\begin{proof}
Considérons $M$ comme un espace vectoriel sur $L$ et appliquons le
Lemme \ref{fields-lemma-trace-and-norm-linear}.
\end{proof}

\noindent
La forme trace se définit au moyen de la trace.

\begin{definition}
\label{fields-definition-trace-pairing}
Soit $L/K$ une extension finie de corps. La {\it forme trace}
de $L/K$ est la forme $K$-bilinéaire symétrique
$$
Q_{L/K} : L \times L \longrightarrow K,\quad
(\alpha, \beta) \longmapsto \text{Trace}_{L/K}(\alpha\beta)
$$
\end{definition}

\noindent
Il se trouve qu'une extension finie de corps est séparable si et seulement
si la forme trace est non dégénérée.

\begin{lemma}
\label{fields-lemma-separable-trace-pairing}
Soit $L/K$ une extension finie de corps. Les assertions suivantes sont équivalentes :
\begin{enumerate}
\item $L/K$ est séparable,
\item $\text{Trace}_{L/K}$ n'est pas identiquement nulle, et
\item la forme trace $Q_{L/K}$ est non dégénérée.
\end{enumerate}
\end{lemma}

\begin{proof}
Il est clair que (3) implique (2). Si (2) est vérifiée, choisissons $\gamma \in L$
tel que $\text{Trace}_{L/K}(\gamma) \not = 0$. Alors, si $\alpha \in L$
est non nul, nous voyons que $Q_{L/K}(\alpha, \gamma/\alpha) \not = 0$.
Ainsi, $Q_{L/K}$ est non dégénérée. Cela démontre l'équivalence de
(2) et (3).

\medskip\noindent
Supposons que $K$ soit de caractéristique $p$ et que $L = K(\alpha)$ avec
$\alpha \not \in K$ et $\alpha^p \in K$. Alors $\text{Trace}_{L/K}(1) = p = 0$.
Pour $i = 1, \ldots, p - 1$, nous voyons que $x^p - \alpha^{pi}$ est le polynôme
minimal de $\alpha^i$ sur $K$, et nous obtenons
$\text{Trace}_{L/K}(\alpha^i) = 0$ d'après le
Lemme \ref{fields-lemma-trace-and-norm-from-minimal-polynomial}.
Ainsi, pour ce type d'extension purement inséparable de degré $p$,
nous voyons que $\text{Trace}_{L/K}$ est identiquement nulle.

\medskip\noindent
Supposons que $L/K$ ne soit pas séparable. Il existe alors un corps intermédiaire
$L/K'/K$ tel que $L/K'$ soit une extension purement inséparable de degré $p$
comme au paragraphe précédent ; voir les
Lemmes \ref{fields-lemma-separable-first} et \ref{fields-lemma-finite-purely-inseparable}.
Ainsi, d'après le Lemme \ref{fields-lemma-trace-and-norm-tower},
nous voyons que $\text{Trace}_{L/K}$ est identiquement nulle.

\medskip\noindent
Supposons au contraire que $L/K$ soit séparable.
Par récurrence sur le degré, nous allons montrer que
$\text{Trace}_{L/K}$ n'est pas identiquement nulle.
Ainsi, d'après le Lemme \ref{fields-lemma-trace-and-norm-tower}, nous pouvons supposer que
$L/K$ est engendrée par un seul élément $\alpha$ (utiliser le fait que, si
la trace est non nulle, alors elle est surjective).
Nous devons montrer que $\text{Trace}_{L/K}(\alpha^e)$
est non nulle pour un certain $e \geq 0$.
Soit $P = x^d + a_1 x^{d - 1} + \ldots + a_d$ le polynôme
minimal de $\alpha$ sur $K$.
Alors $P$ est aussi le polynôme caractéristique de l'application linéaire
$\alpha : L \to L$ ; voir le
Lemme \ref{fields-lemma-characteristic-vs-minimal-polynomial}.
Puisque $L/K$ est séparable, le Lemme \ref{fields-lemma-recognize-separable} montre
que $P$ a $d$ racines deux à deux distinctes $\alpha_1, \ldots, \alpha_d$
dans une clôture algébrique $\overline{K}$ de $K$. Ce sont donc les
valeurs propres de $\alpha : L \to L$.
D'après l'algèbre linéaire, la trace de $\alpha^e$ est
égale à $\alpha_1^e + \ldots + \alpha_d^e$.
Nous concluons donc par le Lemme \ref{fields-lemma-sums-of-powers}.
\end{proof}

\noindent
Soit $K$ un corps et soit $Q : V \times V \to K$ une forme bilinéaire
sur un espace vectoriel de dimension finie sur $K$. Posons $\dim_K(V) = n$.
Alors $Q$ définit une application linéaire $Q : V \to V^*$, $v \mapsto Q(v, -)$,
où $V^* = \Hom_K(V, K)$ est l'espace vectoriel dual. D'où une application linéaire
$$
\det(Q) : \wedge^n(V) \longrightarrow \wedge^n(V)^*
$$
Si nous choisissons un vecteur de base $\omega \in \wedge^n(V)$, nous pouvons
écrire $\det(Q)(\omega) = \lambda \omega^*$, où $\omega^*$
est le vecteur de base dual dans $\wedge^n(V)^*$. Si nous remplaçons notre
choix de $\omega$ par $c \omega$ pour un certain $c \in K^*$, alors
$\omega^*$ devient $c^{-1} \omega^*$ et, par conséquent,
$\lambda$ devient $c^2 \lambda$. Ainsi, la classe de
$\lambda$ dans $K/(K^*)^2$ est bien définie et s'appelle le
{\it discriminant de $Q$}. En développant les définitions, nous voyons que
$$
\lambda = \det(Q(v_i, v_j)_{1 \leq i, j \leq n})
$$
si $\{v_1, \ldots, v_n\}$ est une base de $V$ sur $K$. Notons que
le discriminant est non nul si et seulement si $Q$ est non dégénérée.

\begin{definition}
\label{fields-definition-discriminant}
Soit $L/K$ une extension finie de corps. Le
{\it discriminant de $L/K$} est le discriminant de
la forme trace $Q_{L/K}$.
\end{definition}

\noindent
D'après la discussion ci-dessus et le Lemme \ref{fields-lemma-separable-trace-pairing},
nous voyons que le discriminant est non nul si et seulement
si $L/K$ est séparable. Pour $a \in K$, nous disons souvent
« le discriminant vaut $a$ » alors qu'il serait plus exact
de dire que le discriminant est la classe de $a$ dans $K/(K^*)^2$.

\begin{exercise}
\label{fields-exercise-quadratic-discriminant}
Soit $L/K$ une extension de degré $2$. Montrer qu'il se produit exactement
l'un des cas suivants :
\begin{enumerate}
\item le discriminant vaut $0$, la caractéristique de $K$ vaut $2$,
et $L/K$ est purement inséparable et s'obtient en adjoignant une racine carrée
d'un élément de $K$,
\item le discriminant vaut $1$, la caractéristique de $K$ vaut $2$, et
$L/K$ est séparable de degré $2$,
\item le discriminant n'est pas un carré, la caractéristique de $K$
n'est pas $2$, et $L$ s'obtient à partir de $K$ en adjoignant la racine carrée
du discriminant.
\end{enumerate}
\end{exercise}







\section{Théorie de Galois}
\label{fields-section-galois-theory}

\noindent
Voici la définition.

\begin{definition}
\label{fields-definition-galois}
Une extension de corps $E/F$ est dite {\it galoisienne} si elle est algébrique,
séparable et normale.
\end{definition}

\noindent
On verra qu'une extension finie est galoisienne si et seulement si elle possède
le nombre « correct » d'automorphismes.

\begin{lemma}
\label{fields-lemma-finite-Galois}
Soit $E/F$ une extension finie de corps. Alors $E$ est galoisienne sur $F$
si et seulement si $|\text{Aut}(E/F)| = [E : F]$.
\end{lemma}

\begin{proof}
Supposons que $|\text{Aut}(E/F)| = [E : F]$. D'après le
Lemme \ref{fields-lemma-normal-and-automorphisms}, cela implique
que $E/F$ est séparable et normale, donc galoisienne.
Réciproquement, si $E/F$ est séparable, alors $[E : F] = [E : F]_s$, et
si $E/F$ est de plus normale, le
Lemme \ref{fields-lemma-normal-and-automorphisms} implique que
$|\text{Aut}(E/F)| = [E : F]$.
\end{proof}

\noindent
Motivés par le lemme précédent, nous introduisons le groupe de Galois comme suit.

\begin{definition}
\label{fields-definition-galois-group}
Si $E/F$ est une extension galoisienne, le groupe $\text{Aut}(E/F)$ est
appelé le {\it groupe de Galois} et est noté $\text{Gal}(E/F)$.
\end{definition}

\noindent
Si $L/K$ est une extension galoisienne infinie, il faut considérer le
groupe de Galois comme un groupe topologique. Nous y reviendrons à la
section \ref{fields-section-infinite-galois}.

\begin{lemma}
\label{fields-lemma-galois-goes-up}
Soit $K/E/F$ une tour d'extensions algébriques de corps.
Si $K$ est galoisienne sur $F$, alors $K$ est galoisienne sur $E$.
\end{lemma}

\begin{proof}
Il suffit de combiner les Lemmes \ref{fields-lemma-normal-goes-up} et \ref{fields-lemma-separable-goes-up}.
\end{proof}

\begin{lemma}
\label{fields-lemma-normal-closure-galois}
Soit $L/K$ une extension finie séparable de corps.
Soit $M$ la clôture normale de $L$ sur $K$
(Définition \ref{fields-definition-normal-closure}).
Alors $M/K$ est galoisienne.
\end{lemma}

\begin{proof}
L'extension intermédiaire $M/M_{sep}/K$ du Lemme \ref{fields-lemma-separable-first}
est normale d'après le Lemme \ref{fields-lemma-separable-first-normal}.
Puisque $L/K$ est séparable, nous avons $L \subset M_{sep}$.
Par minimalité, $M = M_{sep}$, ce qui achève la démonstration.
\end{proof}

\noindent
Soit $G$ un groupe agissant sur un corps $K$ (par automorphismes de corps).
Nous emploierons souvent la notation
$$
K^G = \{x \in K \mid \sigma(x) = x \ \forall \sigma \in G\}
$$
et appellerons ce corps le {\it corps des invariants} de l'action de $G$ sur $K$.

\begin{lemma}
\label{fields-lemma-galois-over-fixed-field}
Soit $K$ un corps. Soit $G$ un groupe fini agissant fidèlement sur $K$.
Alors l'extension $K/K^G$ est galoisienne, on a $[K : K^G] = |G|$,
et le groupe de Galois de l'extension est $G$.
\end{lemma}

\begin{proof}
Étant donné $\alpha \in K$, considérons l'orbite $G \cdot \alpha \subset K$
de $\alpha$ sous l'action du groupe. Considérons le polynôme
$$
P = \prod\nolimits_{\beta \in G \cdot \alpha} (x - \beta) \in K[x]
$$
Le point clé de tout le lemme est que ce polynôme est invariant
sous l'action de $G$ et possède donc ses coefficients dans $K^G$.
En effet, pour $\tau \in G$, nous avons
$$
P^\tau = \prod\nolimits_{\beta \in G \cdot \alpha} (x - \tau(\beta)) =
\prod\nolimits_{\beta \in G \cdot \alpha} (x - \beta) = P
$$
car l'application $\beta \mapsto \tau(\beta)$ est une permutation de
l'orbite $G \cdot \alpha$. Ainsi $P \in K^G[x]$. De plus,
$P(\alpha) = 0$, puisque $\alpha$ appartient à son orbite ;
nous en concluons que l'extension $K/K^G$ est algébrique. En outre,
le polynôme minimal $Q$ de $\alpha$ sur $K^G$ divise le
polynôme $P$ que nous venons de construire. Par conséquent, $Q$ est séparable
(par exemple d'après le Lemme \ref{fields-lemma-recognize-separable}) et
nous en concluons que $K/K^G$ est séparable. Ainsi $K/K^G$ est galoisienne.
Pour achever la démonstration, il suffit de montrer que $[K : K^G] = |G|$,
car alors $G$ sera le groupe de Galois d'après le
Lemme \ref{fields-lemma-finite-Galois}.

\medskip\noindent
Choisissons un nombre fini d'éléments $\alpha_i \in K$, $i = 1, \ldots, n$,
de sorte que les égalités $\sigma(\alpha_i) = \alpha_i$ pour $i = 1, \ldots, n$ impliquent
que $\sigma$ est l'élément neutre de $G$. Posons
$$
L = K^G(\{\sigma(\alpha_i); 1 \leq i \leq n, \sigma \in G\}) \subset K
$$
et observons que l'action de $G$ sur $K$ induit une action de $G$ sur $L$.
Nous allons montrer que $L$ est de degré $|G|$ sur $K^G$. Cela achèvera la
démonstration : en effet, si $L \subset K$ était stricte, nous pourrions ajouter un élément
$\alpha \in K$, $\alpha \not \in L$, à notre liste d'éléments
$\alpha_1, \ldots, \alpha_n$ sans agrandir $L$, ce qui est absurde.
Nous sommes ainsi ramenés au cas où $K/K^G$ est finie, lequel est
traité dans le paragraphe suivant.

\medskip\noindent
Supposons $K/K^G$ finie. D'après le Lemme \ref{fields-lemma-primitive-element},
nous pouvons trouver $\alpha \in K$ tel que $K = K^G(\alpha)$.
La construction du premier paragraphe de cette démonstration montre
que $\alpha$ est de degré au plus $|G|$ sur $K^G$. Toutefois, ce
degré ne peut être inférieur à $|G|$, puisque $G$ agit fidèlement sur
$K^G(\alpha) = K$ par construction et en vertu de l'inégalité du
Lemme \ref{fields-lemma-normal-and-automorphisms}.
\end{proof}

\begin{theorem}[Théorème fondamental de la théorie de Galois]
\label{fields-theorem-galois-theory}
Soit $L/K$ une extension galoisienne finie de groupe de Galois $G$.
Alors $K = L^G$ et l'application
$$
\{\text{sous-groupes de }G\}
\longrightarrow
\{\text{extensions intermédiaires }L/M/K\},\quad
H \longmapsto L^H
$$
est une bijection dont l'inverse envoie $M$ sur $\text{Gal}(L/M)$.
Les sous-groupes normaux $H$ de $G$ correspondent exactement aux
extensions intermédiaires $M$ telles que $M/K$ soit galoisienne.
\end{theorem}

\begin{proof}
D'après le Lemme \ref{fields-lemma-galois-goes-up}, pour toute extension intermédiaire $L/M/K$,
l'extension $L/M$ est galoisienne. Bien entendu, $L/M$ est aussi finie
(Lemme \ref{fields-lemma-finite-goes-up}). Ainsi $|\text{Gal}(L/M)| = [L : M]$
d'après le Lemme \ref{fields-lemma-finite-Galois}.
Réciproquement, si $H \subset G$ est un sous-groupe fini, alors
$[L : L^H] = |H|$ d'après le Lemme \ref{fields-lemma-galois-over-fixed-field}.
Il résulte formellement de ces deux observations que nous obtenons
la correspondance bijective annoncée dans le théorème.

\medskip\noindent
Si $H \subset G$ est normal, alors $L^H$ est stable sous l'action de
$G$ et nous obtenons une application canonique $G/H \to \text{Aut}(L^H/K)$.
Cette application est nécessairement injective puisque $\text{Gal}(L/L^H) = H$. Ainsi
$|G/H| = [L^H : K]$ et $L^H$ est galoisienne d'après le
Lemme \ref{fields-lemma-finite-Galois}.

\medskip\noindent
Réciproquement, supposons que $K \subset M \subset L$ et que $M/K$ soit galoisienne.
D'après le Lemme \ref{fields-lemma-lift-maps}, tout
élément $\tau \in \text{Gal}(L/K)$ induit un élément
$\tau|_M \in \text{Gal}(M/K)$. On obtient ainsi un homomorphisme
de groupes de Galois $\text{Gal}(L/K) \to \text{Gal}(M/K)$ dont le
noyau est $H$. Ainsi $H$ est un sous-groupe normal.
\end{proof}

\begin{lemma}
\label{fields-lemma-ses-galois}
Soit $L/M/K$ une tour de corps. Supposons $L/K$ et $M/K$ finies galoisiennes.
Alors nous obtenons une suite exacte courte
$$
1 \to \text{Gal}(L/M) \to \text{Gal}(L/K) \to \text{Gal}(M/K) \to 1
$$
de groupes finis.
\end{lemma}

\begin{proof}
En effet, d'après le Lemme \ref{fields-lemma-lift-maps},
nous voyons que tout élément $\tau \in \text{Gal}(L/K)$ induit un élément
$\tau|_M \in \text{Gal}(M/K)$, ce qui fournit l'homomorphisme de
droite. L'application de gauche identifie le groupe de gauche au noyau
de la flèche de droite. La suite est exacte parce que les ordres
des groupes coïncident comme il convient, par multiplicativité des degrés
dans les tours d'extensions finies
(Lemme \ref{fields-lemma-multiplicativity-degrees}).
On peut aussi appliquer directement le Lemme \ref{fields-lemma-lift-maps}
pour voir que l'application de droite est surjective.
\end{proof}






\section{Théorie de Galois infinie}
\label{fields-section-infinite-galois}

\noindent
Le groupe de Galois est muni d'une topologie canonique.

\begin{lemma}
\label{fields-lemma-galois-profinite}
Soit $E/F$ une extension galoisienne. Munissons $\text{Gal}(E/F)$ de la topologie
la moins fine telle que
$$
\text{Gal}(E/F) \times E \longrightarrow E
$$
soit continue lorsque $E$ est muni de la topologie discrète. Alors
\begin{enumerate}
\item pour tout espace topologique $X$ et toute application $X \to \text{Gal}(E/F)$
tels que l'action $X \times E \to E$ soit continue, l'application induite
$X \to \text{Gal}(E/F)$ est continue,
\item cette topologie fait de $\text{Gal}(E/F)$
un groupe topologique profini.
\end{enumerate}
\end{lemma}

\begin{proof}
Dans toute cette démonstration, nous considérons $E$ comme un espace topologique discret.
Rappelons que la topologie compacte-ouverte sur l'ensemble des applications d'un ensemble dans lui-même
$\text{Map}(E, E)$ est la topologie universelle telle que l'action
$\text{Map}(E, E) \times E \to E$ soit continue. Voir
Topologie, Exemple \ref{topology-example-automorphisms-of-a-set}
pour un énoncé précis. La topologie du lemme sur
$\text{Gal}(E/F)$ est la topologie induite par l'application injective
$\text{Gal}(E/F) \to \text{Map}(E, E)$.
La propriété universelle (1) découle donc de la propriété universelle
correspondante de la topologie compacte-ouverte.
Comme l'ensemble des applications bijectives de l'espace discret considéré, $\text{Aut}(E)$,
muni de la topologie compacte-ouverte, forme un groupe topologique, voir
Topologie, Exemple \ref{topology-example-automorphisms-of-a-set},
et comme $\text{Gal}(E/F) = \text{Aut}(E/F) \to \text{Map}(E, E)$
se factorise par $\text{Aut}(E)$, nous obtenons un groupe topologique.
Autrement dit, nous utilisons l'injection
$$
\text{Gal}(E/F) \subset \text{Aut}(E)
$$
pour munir $\text{Gal}(E/F)$ de la structure induite de groupe topologique
(voir Topologie, section \ref{topology-section-topological-groups})
et, par construction, il s'agit de la structure de groupe topologique
la moins fine pour laquelle l'action $\text{Gal}(E/F) \times E \to E$ est continue.

\medskip\noindent
Pour montrer que $\text{Gal}(E/F)$ est profini, nous raisonnons comme suit
(notre argument est nécessairement non standard, car nous avons défini
la topologie avant de montrer que le groupe de Galois est une limite
projective de groupes finis).
D'après Topologie, Lemme \ref{topology-lemma-profinite-group},
il suffit de montrer que l'espace topologique
sous-jacent à $\text{Gal}(E/F)$ est profini.
Pour toute partie $S \subset E$, considérons l'ensemble
$$
G(S) = \{ f : S \to E \mid
\begin{matrix}
f(\alpha)\text{ est une racine du polynôme minimal}\\
\text{de }\alpha\text{ sur }F\text{ pour tout }\alpha \in S
\end{matrix}
\}
$$
Comme un polynôme n'a qu'un nombre fini de racines, nous voyons que
$G(S)$ est fini pour toute partie finie $S \subset E$. Si $S \subset S'$,
la restriction fournit une application $G(S') \to G(S)$. Observons également
que si $\alpha \in S \cap F$ et $f \in G(S)$, alors $f(\alpha) = \alpha$
car le polynôme minimal est linéaire dans ce cas.
Considérons l'espace topologique profini
$$
G = \lim_{S \subset E\text{ fini}} G(S)
$$
Considérons l'application canonique
$$
c : \text{Gal}(E/F) \longrightarrow G,\quad
\sigma \longmapsto (\sigma|_S : S \to E)_S
$$
Elle est injective et, en explicitant les définitions, le lecteur constate
que la topologie sur $\text{Gal}(E/F)$ définie ci-dessus
est la topologie induite par $G$. Un élément $(f_S) \in G$ appartient à
l'image de $c$ exactement lorsque
(A) $f_S(\alpha) + f_S(\beta) = f_S(\alpha + \beta)$ et
(M) $f_S(\alpha)f_S(\beta) = f_S(\alpha\beta)$ chaque fois que
ces expressions ont un sens (c'est-à-dire lorsque
$\alpha, \beta, \alpha + \beta, \alpha\beta \in S$).
En effet, cela signifie que
$\lim f_S : E \to E$ sera un morphisme de $F$-algèbres,
et donc un automorphisme d'après le
Lemme \ref{fields-lemma-algebraic-extension-self-map}.
Les conditions (A) et (M) pour un triplet donné
$(S, \alpha, \beta)$ définissent un fermé de
$G$ ; ainsi $\text{Gal}(E/F)$ est homéomorphe
à un fermé d'un espace profini, et est donc lui-même profini.
\end{proof}

\begin{lemma}
\label{fields-lemma-galois-infinite}
Soit $L/M/K$ une tour de corps. Supposons que $L/K$ et
$M/K$ soient toutes deux galoisiennes. Alors il existe un homomorphisme canonique
continu surjectif $c : \text{Gal}(L/K) \to \text{Gal}(M/K)$.
\end{lemma}

\begin{proof}
D'après le Lemme \ref{fields-lemma-lift-maps}, pour $\tau : L \to L$ appartenant à
$\text{Gal}(L/K)$, la restriction $\tau|_M : M \to M$
est un élément de $\text{Gal}(M/K)$. Cela définit l'homomorphisme $c$.
La continuité découle de la propriété universelle de la topologie :
l'action
$$
\text{Gal}(L/K) \times M \longrightarrow M,\quad
(\tau, x) \longmapsto \tau(x) = c(\tau)(x)
$$
est continue puisque $M \subset L$ et que l'action
$\text{Gal}(L/K) \times L \to L$ est continue.
La continuité de $c$ résulte donc de la partie (1) du
Lemme \ref{fields-lemma-galois-profinite}.
Le Lemme \ref{fields-lemma-lift-maps} montre également
que l'application est surjective.
\end{proof}

\noindent
Voici une manière plus classique de considérer
le groupe de Galois d'une extension galoisienne infinie.

\begin{lemma}
\label{fields-lemma-infinite-galois-limit}
Soit $L/K$ une extension galoisienne de groupe de Galois $G$.
Soit $\Lambda$ l'ensemble des extensions intermédiaires finies galoisiennes,
c'est-à-dire que $\lambda \in \Lambda$ correspond à $L/L_\lambda/K$,
où $L_\lambda/K$ est finie galoisienne de groupe de Galois $G_\lambda$.
Définissons un ordre partiel sur $\Lambda$ par la règle
$\lambda \geq \lambda'$ si et seulement si
$L_\lambda \supset L_{\lambda'}$. Alors
\begin{enumerate}
\item $\Lambda$ est un ensemble ordonné filtrant,
\item les $L_\lambda$ forment un système d'extensions de $K$ indexé par $\Lambda$
et $L = \colim L_\lambda$,
\item les $G_\lambda$ forment un système projectif de groupes finis indexé par $\Lambda$,
les morphismes de transition sont surjectifs, et
$$
G = \lim_{\lambda \in \Lambda} G_\lambda
$$
comme groupe profini, et
\item chacune des projections $G \to G_\lambda$ est continue et surjective.
\end{enumerate}
\end{lemma}

\begin{proof}
Tout sous-corps de $L$ contenant $K$ est séparable sur $K$
(cela résulte immédiatement de la définition). Soit $S \subset L$
une partie finie. Alors $K(S)/K$ est finie et il existe
une tour $L/E/K(S)/K$ telle que $E/K$ soit finie galoisienne, voir le
Lemme \ref{fields-lemma-normal-closure-inside-normal}.
Ainsi $E = L_\lambda$ pour un certain $\lambda \in \Lambda$.
Cela implique en particulier que l'ensemble $\Lambda$ n'est pas vide.
De plus, étant donnés $\lambda_1, \lambda_2 \in \Lambda$, nous pouvons
écrire $L_{\lambda_i} = K(S_i)$ pour des ensembles finis
$S_1, S_2 \subset L$ (Lemme \ref{fields-lemma-finite-finitely-generated}).
Il existe alors un $\lambda \in \Lambda$ tel que
$K(S_1 \cup S_2) \subset L_\lambda$. Par conséquent,
$\lambda \geq \lambda_1, \lambda_2$ et
$\Lambda$ est filtrant (Catégories, Définition
\ref{categories-definition-directed-system}).
Enfin, puisque tout élément de $L$ appartient à
$L_\lambda$ pour un certain $\lambda \in \Lambda$, il résulte
de la description des limites inductives filtrantes dans
Catégories, section \ref{categories-section-directed-colimits},
que $\colim L_\lambda = L$.

\medskip\noindent
Si $\lambda \geq \lambda'$ dans $\Lambda$, nous obtenons une
application canonique surjective $G_\lambda \to G_{\lambda'}$,
$\sigma \mapsto \sigma|_{L_{\lambda'}}$,
d'après le Lemme \ref{fields-lemma-ses-galois}. Nous obtenons ainsi un système
projectif de groupes finis dont les morphismes de transition sont surjectifs.

\medskip\noindent
Rappelons que $G = \text{Aut}(L/K)$. D'après le Lemme \ref{fields-lemma-galois-infinite},
la restriction $\sigma|_{L_\lambda}$ d'un $\sigma \in G$ à $L_\lambda$
est un élément de $G_\lambda$. De plus, ce procédé fournit une surjection
continue $G \to G_\lambda$. Puisque les morphismes de transition
du système projectif des $G_\lambda$ sont eux aussi donnés par restriction,
il est clair que nous obtenons une application continue canonique
$$
G \longrightarrow \lim_{\lambda \in \Lambda} G_\lambda
$$
La continuité résulte de la définition des limites dans la catégorie des groupes
topologiques ; rappelons que ces limites commutent avec le foncteur d'oubli
vers les catégories des ensembles et des espaces topologiques d'après
Topologie, Lemme \ref{topology-lemma-topological-group-limits}.
D'autre part, puisque $L = \colim L_\lambda$, il est clair
que tout élément de la limite projective (considérée comme un ensemble) définit un
automorphisme de $L$. L'application est donc bijective. Comme la topologie
des deux membres est profinie et qu'une application continue bijective
entre espaces profinis est un homéomorphisme
(Topologie, Lemme \ref{topology-lemma-bijective-map}), la démonstration est achevée.
\end{proof}

\begin{theorem}[Théorème fondamental de la théorie de Galois infinie]
\label{fields-theorem-inifinite-galois-theory}
Soit $L/K$ une extension galoisienne. Soit $G = \text{Gal}(L/K)$
le groupe de Galois considéré comme un groupe topologique profini
(Lemme \ref{fields-lemma-galois-profinite}). Alors $K = L^G$ et l'application
$$
\{\text{sous-groupes fermés de }G\}
\longrightarrow
\{\text{extensions intermédiaires }L/M/K\},\quad
H \longmapsto L^H
$$
est une bijection dont l'inverse envoie $M$ sur $\text{Gal}(L/M)$.
Les extensions intermédiaires finies $M$ correspondent exactement aux sous-groupes
ouverts $H \subset G$. Les sous-groupes fermés normaux $H$ de $G$
correspondent exactement aux extensions intermédiaires $M$ galoisiennes sur $K$.
\end{theorem}

\begin{proof}
Nous utiliserons le résultat de la théorie de Galois finie
(Théorème \ref{fields-theorem-galois-theory})
sans autre mention.
Soit $S \subset L$ une partie finie. Il existe une tour
$L/E/K$ telle que $K(S) \subset E$ et que
$E/K$ soit finie galoisienne, voir le Lemme \ref{fields-lemma-normal-closure-inside-normal}.
Autrement dit, $L/K$ est la réunion de ses extensions intermédiaires
finies galoisiennes.
Pour un tel $E$, d'après le Lemme \ref{fields-lemma-galois-infinite},
l'application $\text{Gal}(L/K) \to \text{Gal}(E/K)$ est surjective
et continue, c'est-à-dire que son noyau est ouvert puisque la topologie
sur $\text{Gal}(E/K)$ est discrète.
En particulier, aucun élément de $L \setminus K$ n'est fixé par
$\text{Gal}(L/K)$, car $E^{\text{Gal}(E/K)} = K$.
Cela prouve que $L^G = K$.

\medskip\noindent
D'après le Lemme \ref{fields-lemma-galois-goes-up}, pour toute extension intermédiaire $L/M/K$,
l'extension $L/M$ est galoisienne. Il résulte immédiatement de la définition
de la topologie sur $G$ que le sous-groupe $\text{Gal}(L/M)$ est fermé.
En appliquant ce qui précède à $L/M$, nous voyons que $L^{\text{Gal}(L/M)} = M$.

\medskip\noindent
Réciproquement, soit $H \subset G$ un sous-groupe fermé. Nous affirmons que
$H = \text{Gal}(L/L^H)$. L'inclusion $H \subset \text{Gal}(L/L^H)$
est claire. Supposons que $g \in \text{Gal}(L/L^H)$. Soit $S \subset L$
une partie finie. Nous allons montrer que le voisinage ouvert
$U_S(g) = \{g' \in G \mid g'(s) = g(s)\text{ pour tout }s \in S\}$ de $g$ rencontre $H$.
Cela implique que $g \in H$, puisque $H$ est fermé.
Soit $L/E/K$ une extension intermédiaire finie galoisienne contenant $K(S)$,
comme dans le premier paragraphe de la démonstration, et considérons l'homomorphisme
$c : \text{Gal}(L/K) \to \text{Gal}(E/K)$.
Alors $L^H \cap E = E^{c(H)}$. Puisque $g$ fixe $L^H$, il fixe
$E^{c(H)}$ et donc $c(g) \in c(H)$ par la théorie de Galois finie.
Choisissons $h \in H$ tel que $c(h) = c(g)$. Alors $h \in U_S(g)$, comme voulu.

\medskip\noindent
À ce stade, nous avons établi la correspondance entre sous-groupes fermés
et extensions intermédiaires.

\medskip\noindent
Supposons $H \subset G$ ouvert. En raisonnant comme ci-dessus, nous voyons que
$H$ contient $\text{Gal}(L/E)$ pour une extension intermédiaire finie
galoisienne $E$ assez grande, et que $L^H$ est contenu
dans $E$, donc fini sur $K$. Réciproquement, si $M$ est une extension
intermédiaire finie, alors $M$ est engendré par une partie finie $S$,
et le sous-groupe correspondant est l'ouvert $U_S(e)$,
où $e \in G$ est l'élément neutre.

\medskip\noindent
Supposons que $K \subset M \subset L$ et que $M/K$ soit galoisienne.
D'après le Lemme \ref{fields-lemma-galois-infinite}, il existe un homomorphisme surjectif
continu de groupes de Galois
$\text{Gal}(L/K) \to \text{Gal}(M/K)$ dont le
noyau est $\text{Gal}(L/M)$. Ainsi $\text{Gal}(L/M)$ est un sous-groupe
fermé normal.

\medskip\noindent
Enfin, supposons $N \subset G$ normal et fermé. Pour toute
extension intermédiaire $L/E/K$ comme dans le premier paragraphe de la démonstration, l'image
$c(N) \subset \text{Gal}(E/K)$ est un sous-groupe normal.
Par conséquent, $L^N = \bigcup E^{c(N)}$ est une réunion d'extensions galoisiennes
de $K$ (par la théorie de Galois finie), et est donc galoisienne sur $K$.
\end{proof}

\begin{lemma}
\label{fields-lemma-ses-infinite-galois}
Soit $L/M/K$ une tour de corps. Supposons $L/K$ et $M/K$ galoisiennes.
Alors nous obtenons une suite exacte courte
$$
1 \to \text{Gal}(L/M) \to \text{Gal}(L/K) \to \text{Gal}(M/K) \to 1
$$
de groupes topologiques profinis.
\end{lemma}

\begin{proof}
C'est une reformulation du Lemme \ref{fields-lemma-galois-infinite}.
\end{proof}





\section{Les nombres complexes}
\label{fields-section-complex-numbers}

\noindent
Le théorème fondamental de l'algèbre affirme que le corps des nombres complexes
est algébriquement clos. Nous en discutons brièvement dans cette section.

\medskip\noindent
La première remarque que nous souhaitons faire est qu'il faut utiliser quelques
résultats d'analyse pour le démontrer. Nous utiliserons le fait intuitivement
clair que tout polynôme de degré impair à coefficients réels possède une racine réelle.
En effet, soit
$P(x) = a_{2k + 1} x^{2k + 1} + \ldots + a_0 \in \mathbf{R}[x]$
pour un certain $k \geq 0$ et avec $a_{2k + 1} \not = 0$.
Nous pouvons supposer, et supposons, que $a_{2k + 1} > 0$. Alors, pour
$x \in \mathbf{R}$ positif et très grand, nous avons $P(x) > 0$, car le terme
$a_{2k + 1} x^{2k + 1}$ domine tous les autres termes. De même,
si $x \ll 0$, alors $P(x) < 0$ pour la même raison (et c'est ici
que nous utilisons que le degré est impair). Le théorème des valeurs
intermédiaires donne donc un $x \in \mathbf{R}$ tel que $P(x) = 0$.

\medskip\noindent
Une conséquence de ce qui précède est que $\mathbf{R}$ ne possède
aucune extension de corps non triviale de degré impair, car les éléments de telles extensions
auraient des polynômes minimaux de degré impair.

\medskip\noindent
Soit ensuite $K/\mathbf{R}$ une extension finie galoisienne de groupe de Galois $G$.
Soit $P \subset G$ un $2$-sous-groupe de Sylow. Alors $K^P/\mathbf{R}$ est une
extension de degré impair, donc $K^P = \mathbf{R}$ d'après ce qui précède,
ce qui implique à son tour
$G = P$. (Tous ces arguments reposent bien entendu sur la théorie de Galois.)
Ainsi $G$ est un $2$-groupe. Si $G$ est non trivial, nous voyons alors que
$\mathbf{C} \subset K$, puisque $\mathbf{C}$ est (à isomorphisme près) l'unique
extension de degré $2$ de $\mathbf{R}$. Si $G$ avait plus de $2$ éléments,
nous obtiendrions une extension quadratique de $\mathbf{C}$.
C'est absurde, car tout nombre complexe possède une racine carrée.

\medskip\noindent
Conclusion : $\mathbf{C}$ est algébriquement clos. En effet, sinon
nous obtiendrions une extension finie non triviale $K/\mathbf{C}$
que nous pourrions supposer normale (donc galoisienne) sur $\mathbf{R}$ d'après le
Lemme \ref{fields-lemma-normal-closure}. Mais nous avons vu ci-dessus qu'alors
$K = \mathbf{C}$.

\begin{lemma}[Théorème fondamental de l'algèbre]
\label{fields-lemma-C-algebraically-closed}
Le corps $\mathbf{C}$ est algébriquement clos.
\end{lemma}

\begin{proof}
Voir la discussion ci-dessus.
\end{proof}





\section{Extensions de Kummer}
\label{fields-section-Kummer}

\noindent
Soit $K$ un corps. Soit $n \geq 2$ un entier tel que $K$ contienne
une racine primitive $n$-ième de $1$. Soit $a \in K^*$. Soit $L$ une extension
de $K$ obtenue en adjoignant une racine $b$ de l'équation $x^n = a$.
Alors $L/K$ est galoisienne. Si $G = \text{Gal}(L/K)$ est le groupe de Galois,
l'application
$$
G \longrightarrow \mu_n(K),\quad \sigma \longmapsto \sigma(b)/b
$$
est un homomorphisme injectif de groupes. En particulier, $G$ est cyclique
et son ordre divise $n$, puisqu'il s'identifie à un sous-groupe du groupe cyclique $\mu_n(K)$.
La théorie de Kummer fournit une réciproque.

\begin{lemma}[Extensions de Kummer]
\label{fields-lemma-Kummer}
Soit $L/K$ une extension galoisienne de corps dont le groupe de Galois est
$\mathbf{Z}/n\mathbf{Z}$. Supposons de plus que la caractéristique de $K$
soit première à $n$ et que $K$ contienne une racine primitive $n$-ième de $1$.
Alors $L = K[z]$ avec $z^n \in K$.
\end{lemma}

\begin{proof}
Soit $\zeta \in K$ une racine primitive $n$-ième de $1$.
Soit $\sigma$ un générateur de $\text{Gal}(L/K)$.
Considérons $\sigma : L \to L$ comme un opérateur $K$-linéaire.
Notons que $\sigma^n - 1 = 0$ comme opérateur linéaire.
En appliquant l'indépendance linéaire des caractères
(Lemme \ref{fields-lemma-independence-characters}), nous voyons
qu'aucun polynôme sur $K$ de degré $< n$ ne peut
annuler $\sigma$. Le polynôme minimal de $\sigma$
comme opérateur linéaire est donc $x^n - 1$. 
Puisque $\zeta$ est une racine de $x^n - 1$, l'algèbre linéaire
donne un $0 \neq z \in L$ tel que $\sigma(z) = \zeta z$.
Cet élément $z$ vérifie $z^n \in K$, car
$\sigma(z^n) = (\zeta z)^n = z^n$. De plus, les éléments
$z, \sigma(z), \ldots, \sigma^{n - 1}(z) =
z, \zeta z, \ldots \zeta^{n - 1} z$ sont deux à deux distincts,
ce qui garantit que $z$ engendre $L$ sur $K$.
Ainsi $L = K[z]$, ce qu'il fallait démontrer.
\end{proof}

\begin{lemma}
\label{fields-lemma-adjoint-pth-root-unity}
Soit $K$ un corps, muni d'une clôture algébrique $\overline{K}$.
Soit $p$ un nombre premier distinct de la caractéristique de $K$.
Soit $\zeta \in \overline{K}$ une racine primitive $p$-ième
de $1$. Alors $K(\zeta)/K$ est une extension galoisienne dont le degré divise $p - 1$.
\end{lemma}

\begin{proof}
Le polynôme $x^p - 1$ se décompose entièrement sur
$K(\zeta)$, ses racines étant $1, \zeta, \zeta^2, \ldots, \zeta^{p - 1}$.
Ainsi, dans l'extension $K(\zeta)/K$, le corps supérieur est un corps de
décomposition ; cette extension est donc normale.
L'extension est séparable, car $x^p - 1$ est un polynôme séparable.
Elle est donc galoisienne. Tout automorphisme de $K(\zeta)$ sur $K$
envoie $\zeta$ sur $\zeta^i$ pour un certain $1 \leq i \leq p - 1$.
Le groupe de Galois est donc un sous-groupe de $(\mathbf{Z}/p\mathbf{Z})^*$.
\end{proof}

\begin{lemma}
\label{fields-lemma-subfields-kummer}
\begin{reference}
\cite[Théorème 5.2]{Radical}
\end{reference}
Soit $K$ un corps. Soit $L/K$ une extension finie de degré $e$
engendrée par un élément $\alpha$ tel que $a = \alpha^e \in K$.
Si toute racine $e$-ième de l'unité appartenant à $L$ est contenue dans $K$, alors 
toute extension intermédiaire $L/L'/K$ est engendrée par $\alpha^d$ pour un certain $d | e$.
\end{lemma}

\begin{proof}
Observons que, pour $d | e$, le sous-corps $K(\alpha^d)$ vérifie
$[K(\alpha^d) : K] = e/d$ et $[L : K(\alpha^d)] = d$.
Soit $L/L'/K$ une extension intermédiaire. Posons $d = [L : L']$. Si $\alpha^d \in L'$,
alors $L' = K(\alpha^d)$ pour des raisons de degré.
Soit $P \in L'[x]$ le polynôme minimal de $\alpha$ sur $L'$.
Alors $P$ divise $x^e - a$ et $P$ est de degré $d$.
Écrivons
$$
x^e - a = \prod\nolimits_{i = 1, \ldots, e} (x - \zeta_i \alpha)
$$
dans un corps de décomposition de $x^e - a$ sur $L$. Les $\zeta_i$ sont des racines
$e$-ièmes de l'unité et, après renumérotation, nous avons
$$
P = \prod\nolimits_{i = 1, \ldots, d} (x - \zeta_i \alpha)
$$
Le terme constant de $P$ est égal à
$$
c = (-1)^d(\prod\nolimits_{i = 1, \ldots, d} \zeta_i) \alpha^d
$$
et appartient à $L' \subset L$. Comme $\alpha \in L$, cela implique que
$\zeta = \prod_{i = 1, \ldots, d} \zeta_i$ appartient à $L$, et donc à $K$ par
notre hypothèse. Ainsi $\alpha^d = (-1)^d\zeta^{-1}c \in L'$, ce qui conclut.
\end{proof}




\section{Extensions d'Artin–Schreier}
\label{fields-section-Artin-Schreier}

\noindent
Soit $K$ un corps de caractéristique $p > 0$. Soit $a \in K$. Soit $L$ une
extension de $K$ obtenue en adjoignant une racine $b$ de l'équation
$x^p - x = a$. Alors $L/K$ est galoisienne. Si $G = \text{Gal}(L/K)$ est le groupe
de Galois, alors l'application
$$
G \longrightarrow \mathbf{Z}/p\mathbf{Z},\quad
\sigma \longmapsto \sigma(b) - b
$$
est un homomorphisme injectif de groupes. En particulier, $G$ est cyclique
d'ordre divisant $p$ en tant que sous-groupe de $\mathbf{Z}/p\mathbf{Z}$.
La théorie des extensions d'Artin–Schreier fournit une réciproque.

\begin{lemma}[Extensions d'Artin–Schreier]
\label{fields-lemma-Artin-Schreier}
Soit $L/K$ une extension galoisienne de corps de caractéristique $p > 0$
de groupe de Galois $\mathbf{Z}/p\mathbf{Z}$. Alors $L = K[z]$ avec
$z^p - z \in K$.
\end{lemma}

\begin{proof}
Soit $\sigma$ un générateur de $\text{Gal}(L/K)$.
Considérons $\sigma : L \to L$ comme un opérateur $K$-linéaire.
Observons que $\sigma^p - 1 = 0$ comme opérateur linéaire.
En appliquant l'indépendance linéaire des caractères
(Lemme \ref{fields-lemma-independence-characters}),
aucun polynôme de degré $< p$ ne peut
annuler $\sigma$. Nous en concluons que le polynôme minimal
de $\sigma$ est $x^p - 1 = (x - 1)^p$.
Cela implique qu'il existe $w \in L$ tel que
$(\sigma - 1)^{p - 1}(w) = y$ soit non nul. Alors
$\sigma(y) = y$, c'est-à-dire $y \in K$. Ainsi
$z = y^{-1}(\sigma - 1)^{p - 2}(w)$ vérifie
$\sigma(z) = z + 1$. Comme $z \not \in K$, nous avons $L = K[z]$.
De plus, puisque $\sigma(z^p - z) = (z + 1)^p - (z + 1) = z^p - z$,
nous voyons que $z^p - z \in K$, ce qui achève la démonstration.
\end{proof}






\section{Transcendance}
\label{fields-section-transcendence}

\noindent
Rappelons les définitions usuelles.

\begin{definition}
\label{fields-definition-transcendence}
Soit $K/k$ une extension de corps.
\begin{enumerate}
\item Une famille d'éléments $\{x_i\}_{i \in I}$ de $K$ est dite
{\it algébriquement indépendante} sur $k$ si l'application
$$
k[X_i; i\in I] \longrightarrow K
$$
qui envoie $X_i$ sur $x_i$ est injective.
\item Le corps des fractions d'un anneau de polynômes
$k[x_i; i \in I]$ est noté $k(x_i; i\in I)$.
\item Une {\it extension transcendante pure} de $k$ est toute
extension de corps $K/k$ isomorphe au corps des
fractions d'un anneau de polynômes sur $k$.
\item Une {\it base de transcendance} de $K/k$ est une
famille d'éléments $\{x_i\}_{i \in I}$ qui sont
algébriquement indépendants sur $k$ et tels que
l'extension $K/k(x_i; i\in I)$ soit algébrique.
\end{enumerate}
\end{definition}

\begin{example}
\label{fields-example-pi-transcendental}
Le corps $\mathbf{Q}(\pi)$ est une extension transcendante pure, car
$\pi$ n'est racine d'aucun polynôme non nul à coefficients rationnels.
En particulier, $\mathbf{Q}(\pi) \cong \mathbf{Q}(x)$.
\end{example}

\begin{lemma}
\label{fields-lemma-transcendence-degree}
Soit $E/F$ une extension de corps. Il existe une base de transcendance de $E$ sur $F$.
Deux bases de transcendance quelconques ont même cardinal.
\end{lemma}

\begin{proof}
Soit $A$ une partie algébriquement indépendante de $E$. Soit $G$ une partie
de $E$ contenant $A$ qui engendre $E/F$. Nous affirmons qu'il existe une
base de transcendance $B$ telle que $A \subset B \subset G$.
Pour le démontrer, considérons l'ensemble $\mathcal{B}$ des parties
algébriquement indépendantes qui sont des parties de $G$ contenant $A$.
Ordonnons partiellement $\mathcal{B}$ par inclusion.
Alors $\mathcal{B}$ contient au moins l'élément $A$.
La réunion des éléments d'une partie totalement ordonnée $T$ de $\mathcal{B}$
est une partie de $E$ algébriquement indépendante sur $F$, car toute relation
de dépendance algébrique apparaîtrait déjà dans l'un des éléments de $T$
(puisque les polynômes ne font intervenir qu'un nombre fini de variables). Cette réunion
contient aussi $A$ et est contenue dans $G$. D'après le lemme de Zorn, il existe un
élément maximal $B \in \mathcal{B}$. Nous affirmons maintenant que $E$ est algébrique sur $F(B)$.
En effet, sinon il existerait un élément
$f \in G$ transcendant sur $F(B)$, puisque $F(G) = E$. Alors
$B \cup\{f\}$ serait algébriquement indépendante, contrairement à la
maximalité de $B$. Ainsi $B$ est la base de transcendance recherchée.

\medskip\noindent
Soient $B$ et $B'$ deux bases de transcendance. Sans perte de généralité, nous
pouvons supposer que $|B'| \leq |B|$. Divisons la démonstration en deux cas :
le premier est celui où $B$ est un ensemble infini. Alors, pour chaque $\alpha \in B'$,
il existe un ensemble fini $B_{\alpha} \subset B$ tel que $\alpha$ soit algébrique
sur
$F(B_{\alpha})$, puisque toute relation de dépendance algébrique n'utilise qu'un nombre fini
d'indéterminées. Posons alors $B^* = \bigcup_{\alpha\in B'} B_{\alpha}$.
Par construction, $B^* \subset B$, mais nous affirmons qu'en fait ces deux ensembles sont
égaux. Pour le voir, supposons qu'ils ne le soient pas et qu'il existe, par exemple, un élément
$\beta \in B \setminus B^*$. Nous savons que $\beta$ est algébrique sur $F(B')$, qui
est algébrique sur $F(B^*)$. Par conséquent, $\beta$ est algébrique sur $F(B^*)$, ce qui est
une contradiction. Ainsi $|B| \leq |\bigcup_{\alpha \in B'} B_{\alpha}|$.
Si $B'$ est fini, alors $B$ l'est également ; nous pouvons donc supposer $B'$ infini.
Il vient
$$
|B| \leq |\bigcup\nolimits_{\alpha \in B'} B_{\alpha}| = |B'|
$$
car chaque $B_\alpha$ est fini et $B'$ est infini. Par conséquent, dans le
cas infini, $|B| = |B'|$.

\medskip\noindent
Il reste à étudier le cas où $B$ est fini.
Dans ce cas, $B'$ est également fini ; écrivons
$B = \{\alpha_1, \ldots, \alpha_n\}$ et
$B' = \{\beta_1, \ldots, \beta_m\}$ avec $m \leq n$.
Raisonnons par récurrence sur $m$ : si $m = 0$, alors $E/F$ est algébrique, donc
$B = \emptyset$ et $n = 0$. Si $m > 0$, il existe un polynôme irréductible
$f \in F[x, y_1, \ldots, y_n]$ tel que
$f(\beta_1, \alpha_1, \ldots, \alpha_n) = 0$ et tel que $x$ intervienne dans $f$.
Puisque $\beta_1$ n'est pas algébrique sur $F$, $f$ doit faire intervenir un certain $y_i$ ;
sans perte de généralité, supposons que $f$ fasse intervenir $y_1$.
Posons $B^* = \{\beta_1, \alpha_2, \ldots, \alpha_n\}$.
Nous affirmons que $B^*$ est une base de transcendance de l'extension $E/F$. Pour le montrer, remarquons que
nous avons une tour d'extensions algébriques
$$
E/ F(B^*, \alpha_1) / F(B^*)
$$
puisque $\alpha_1$ est algébrique sur $F(B^*)$.
Nous affirmons maintenant que $B^*$ (en comptant les éléments avec multiplicité) est
algébriquement indépendante sur $F$ : en effet, sinon, il existerait un
polynôme irréductible $g\in F[x, y_2, \ldots, y_n]$ tel que
$g(\beta_1, \alpha_2, \ldots, \alpha_n) = 0$,
dans lequel $x$ devrait intervenir, rendant $\beta_1$
algébrique sur $F(\alpha_2, \ldots, \alpha_n)$, ce qui rendrait $\alpha_1$
algébrique sur $F(\alpha_2, \ldots, \alpha_n)$, chose impossible.
Cela signifie donc que $\{\alpha_2, \ldots, \alpha_n\}$ et
$\{\beta_2, \ldots, \beta_m\}$ sont des bases de transcendance de $E$ sur $F(\beta_1)$ ;
par récurrence, $m = n$.
\end{proof}

\begin{definition}
\label{fields-definition-transcendence-degree}
Soit $K/k$ une extension de corps.
Le {\it degré de transcendance} de $K$ sur $k$ est
le cardinal d'une base de transcendance de $K$ sur $k$.
Il est noté $\text{trdeg}_k(K)$.
\end{definition}

\begin{lemma}
\label{fields-lemma-transcendence-degree-tower}
Soient $L/K/k$ des extensions de corps.
Alors
$$
\text{trdeg}_k(L) =
\text{trdeg}_K(L) +
\text{trdeg}_k(K).
$$
\end{lemma}

\begin{proof}
Choisissons une base de transcendance $A \subset K$ de $K$ sur $k$.
Choisissons une base de transcendance $B \subset L$ de $L$ sur $K$.
Il est alors immédiat que $A \cup B$ est une base de transcendance
de $L$ sur $k$.
\end{proof}

\begin{example}
\label{fields-example-pi-e-transcendental}
Considérons l'extension de corps $\mathbf{Q}(e, \pi)$ obtenue en
adjoignant les nombres $e$ et $\pi$. Cette extension de corps a un degré de transcendance
au moins égal à $1$, puisque $e$ et $\pi$ sont tous deux transcendants sur les
rationnels. Toutefois, cette extension de corps pourrait avoir un degré de transcendance
égal à $2$ si $e$ et $\pi$ sont algébriquement indépendants. On ignore si
c'est le cas ; le problème de la détermination de
$\text{trdeg}(\mathbf{Q}(e, \pi))$ est donc ouvert.
\end{example}

\begin{example}
\label{fields-example-function-field-not-unique-transcendence-basis}
Soit $F$ un corps et soit $E = F(t)$. Alors $\{t\}$ est une
base de transcendance puisque $E = F(t)$. Toutefois, $\{t^2\}$
est aussi une base de transcendance puisque $F(t)/F(t^2)$ est algébrique.
Cela montre que, bien que nous puissions toujours décomposer une extension
$E/F$ en une extension algébrique $E/F'$ et une
extension transcendante pure $F'/F$, cette décomposition n'est pas unique et
dépend du choix d'une base de transcendance.
\end{example}

\begin{example}
\label{fields-example-riemann-surface-transcendence}
Soit $X$ une surface de Riemann compacte. Alors le corps des fonctions
$\mathbf{C}(X)$ (voir Exemple \ref{fields-example-field-of-meromorphic-functions})
est de degré de transcendance un sur $\mathbf{C}$. En fait, {\it toute}
extension de type fini de $\mathbf{C}$ de degré de transcendance
un provient d'une surface de Riemann. Il existe même une équivalence de
catégories entre la catégorie des surfaces de Riemann compactes et des
applications holomorphes (non constantes), et l'opposée de la catégorie des extensions
de type fini de $\mathbf{C}$ de degré de transcendance $1$
et des morphismes de $\mathbf{C}$-algèbres. Voir \cite{Forster}.

\medskip\noindent
Il existe également une version algébrique de l'énoncé précédent. Étant donnée une
courbe algébrique (irréductible) dans l'espace projectif sur un corps algébriquement
clos $k$ (par exemple le corps des nombres complexes), on peut considérer son
« corps des fonctions rationnelles » : il s'agit essentiellement des fonctions qui se présentent comme
des quotients de polynômes, dont le dénominateur ne s'annule pas identiquement
sur la courbe. Il existe une anti-équivalence analogue de catégories
(Courbes algébriques, Théorème \ref{curves-theorem-curves-rational-maps})
entre les courbes projectives lisses et
les morphismes non constants de courbes, et les extensions de type fini de $k$ de
degré de transcendance un. Voir \cite{H}.
\end{example}

\begin{definition}
\label{fields-definition-algebraically-closed-in}
Soit $K/k$ une extension de corps.
\begin{enumerate}
\item La {\it clôture algébrique de $k$ dans $K$} est le sous-corps
$k'$ de $K$ formé des éléments de $K$ qui sont algébriques sur $k$.
\item Nous disons que $k$ est {\it algébriquement clos dans $K$} si
tout élément de $K$ qui est algébrique sur $k$
appartient à $k$.
\end{enumerate}
\end{definition}

\begin{lemma}
\label{fields-lemma-purely-transcendental-degree}
Soit $k'/k$ une extension finie de corps. Soit
$k'(x_1, \ldots, x_r)/k(x_1, \ldots, x_r)$
l'extension induite des extensions transcendantes pures.
Alors $[k'(x_1, \ldots, x_r) : k(x_1, \ldots, x_r)] = [k' : k] < \infty$.
\end{lemma}

\begin{proof}
Par multiplicativité des degrés des extensions
(Lemme \ref{fields-lemma-multiplicativity-degrees}),
il suffit de le démontrer lorsque $k'$ est engendré par un unique élément
$\alpha \in k'$ sur $k$. Soit $f \in k[T]$ le polynôme minimal
de $\alpha$ sur $k$. Alors $k'(x_1, \ldots, x_r)$ est engendré
par $\alpha, x_1, \ldots, x_r$ sur $k$, et donc $k'(x_1, \ldots, x_r)$
est engendré par $\alpha$ sur $k(x_1, \ldots, x_r)$.
Il suffit donc de montrer que $f$ reste irréductible comme
élément de $k(x_1, \ldots, x_r)[T]$. Nous ne faisons qu'esquisser la démonstration.
Il est clair que $f$ est irréductible comme élément de
$k[x_1, \ldots, x_r, T]$, par exemple parce que $f$ est unitaire
comme polynôme en $T$ et que toute factorisation éventuelle dans
$k[x_1, \ldots, x_r, T]$ conduirait à une factorisation dans $k[T]$
en posant les $x_i$ égaux à $0$. Le lemme de Gauss permet de conclure.
\end{proof}

\begin{lemma}
\label{fields-lemma-algebraic-closure-in-finitely-generated}
Soit $K/k$ une extension de corps de type fini.
La clôture algébrique de $k$ dans $K$ est finie sur $k$.
\end{lemma}

\begin{proof}
Supposons que $x_1, \ldots, x_r \in K$ forment une base de transcendance de $K$
sur $k$. Alors $n = [K : k(x_1, \ldots, x_r)] < \infty$.
Supposons que $k \subset k' \subset K$ et que $k'/k$ soit finie.
Dans ce cas,
$[k'(x_1, \ldots, x_r) : k(x_1, \ldots, x_r)] = [k' : k] < \infty$, voir le
Lemme \ref{fields-lemma-purely-transcendental-degree}.
Par conséquent,
$$
[k' : k] = [k'(x_1, \ldots, x_r) : k(x_1, \ldots, x_r)] \leq
[K : k(x_1, \ldots, x_r)] = n.
$$
Autrement dit, les degrés des extensions intermédiaires finies sont bornés,
et le lemme en résulte.
\end{proof}






\section{Extensions linéairement disjointes}
\label{fields-section-linearly-disjoint}

\noindent
Soit $k$ un corps, et soient $K$ et $L$ des extensions de corps de $k$.
Supposons aussi que $K$ et $L$ soient plongés dans un corps plus grand $\Omega$.

\begin{definition}
\label{fields-definition-compositum}
Considérons un diagramme
\begin{equation}
\label{fields-equation-inside-omega}
\vcenter{
\xymatrix{
L \ar[r] & \Omega \\
k \ar[r] \ar[u] & K \ar[u]
}
}
\end{equation}
d'extensions de corps. Le {\it composé de $K$ et $L$ dans $\Omega$},
noté $KL$, est le plus petit sous-corps de $\Omega$ contenant à la fois
$L$ et $K$.
\end{definition}

\noindent
Il est clair que $KL$ est engendré par l'ensemble $K \cup L$ sur $k$,
par l'ensemble $K$ sur $L$, et par l'ensemble $L$ sur $K$.

\medskip\noindent
{\bf Avertissement :} La classe d'isomorphisme du corps composé dépend du
choix des plongements de $K$ et $L$ dans $\Omega$. Considérons par exemple
les corps de nombres $K = \mathbf{Q}(2^{1/8}) \subset \mathbf{R}$ et
$L = \mathbf{Q}(2^{1/12}) \subset \mathbf{R}$. Leur composé dans
$\mathbf{R}$ est le corps $\mathbf{Q}(2^{1/24})$, de degré $24$ sur
$\mathbf{Q}$. Toutefois, si nous plongeons $K = \mathbf{Q}[x]/(x^8 - 2)$ dans
$\mathbf{C}$ en envoyant $x$ sur $2^{1/8}e^{2\pi i/8}$, alors le composé
$\mathbf{Q}(2^{1/12}, 2^{1/8}e^{2\pi i/8})$ contient $i = e^{2\pi i/4}$ et est de
degré $48$ sur $\mathbf{Q}$ (nous omettons de montrer que le degré vaut $48$, mais
l'existence de $i$ prouve assurément que les deux composés ne sont pas isomorphes).

\begin{definition}
\label{fields-definition-linearly-disjoint}
Considérons un diagramme de corps comme dans (\ref{fields-equation-inside-omega}).
Nous disons que $K$ et $L$ sont {\it linéairement disjoints sur $k$ dans $\Omega$}
si l'application
$$
K \otimes_k L \longrightarrow KL,\quad
\sum x_i \otimes y_i \longmapsto \sum x_i y_i
$$
est injective.
\end{definition}

\noindent
Le lemme suivant ne semble trouver sa place nulle part ailleurs.

\begin{lemma}
\label{fields-lemma-normal-case}
Soit $E/F$ une extension algébrique normale de corps. Il existe des extensions intermédiaires
$E / E_{sep} /F$ et $E / E_{insep} / F$ telles que
\begin{enumerate}
\item $F \subset E_{sep}$ soit galoisienne et $E_{sep} \subset E$
soit purement inséparable,
\item $F \subset E_{insep}$ soit purement inséparable et $E_{insep} \subset E$
soit galoisienne,
\item l'application de multiplication $E_{sep} \otimes_F E_{insep} \longrightarrow E$
est un isomorphisme.
\end{enumerate}
\end{lemma}

\begin{proof}
Nous avons trouvé le sous-corps $E_{sep}$ dans le Lemme \ref{fields-lemma-separable-first}.
Nous posons $E_{insep} = E^{\text{Aut}(E/F)}$. Les détails sont omis.
\end{proof}









\section{Rappels}
\label{fields-section-algebraic}

\noindent
Dans cette section, nous récapitulons brièvement les résultats obtenus ci-dessus.

\medskip\noindent
Soit $K/k$ une extension de corps. Soit $\alpha \in K$. Les
possibilités suivantes se présentent :
\begin{enumerate}
\item L'élément $\alpha$ est transcendant sur $k$.
\item L'élément $\alpha$ est algébrique sur $k$. Notons
$P(T) \in k[T]$ son {\it polynôme minimal}. C'est un polynôme unitaire
$P(T) = T^d + a_1 T^{d - 1} + \ldots + a_d$ à coefficients dans
$k$. Il est irréductible et $P(\alpha) = 0$. Ces propriétés
déterminent $P$ de manière unique, et l'entier $d$ est appelé le
{\it degré de $\alpha$ sur $k$}. Il existe deux sous-cas :
\begin{enumerate}
\item Le polynôme $\text{d}P/\text{d}T$ n'est pas identiquement nul.
Cela équivaut à la condition
$P(T) = \prod_{i = 1, \ldots, d} (T - \alpha_i)$ pour des
éléments deux à deux distincts $\alpha_1, \ldots, \alpha_d$
de la clôture algébrique de $k$.
Dans ce cas, nous disons que $\alpha$ est {\it séparable} sur $k$.
\item Le $\text{d}P/\text{d}T$ est identiquement nul. Dans ce cas, la
caractéristique $p$ de $k$ est $ > 0$, et $P$ est en fait un polynôme
en $T^p$. Il existe manifestement une plus grande puissance $q = p^e$ telle que $P$ soit
un polynôme en $T^q$. Alors l'élément $\alpha^q$ est séparable sur $k$.
\end{enumerate}
\end{enumerate}

\begin{definition}
\label{fields-definition-separable-algebraic}
Extensions algébriques de corps.
\begin{enumerate}
\item Une extension de corps $K/k$ est dite {\it algébrique}
si tout élément de $K$ est algébrique sur $k$.
\item Une extension algébrique $k'/k$ est dite {\it séparable}
si tout $\alpha \in k'$ est séparable sur $k$.
\item Une extension algébrique
$k'/k$ est dite {\it purement inséparable} si
la caractéristique de $k$ est $p > 0$ et si, pour tout élément
$\alpha \in k'$, il existe une puissance $q$ de $p$ telle que
$\alpha^q \in k$.
\item Une extension algébrique $k'/k$ est dite {\it normale}
si, pour tout $\alpha \in k'$, le polynôme minimal $P(T) \in k[T]$
de $\alpha$ sur $k$ se décompose entièrement en facteurs linéaires sur $k'$.
\item Une extension algébrique $k'/k$ est dite {\it galoisienne}
si elle est séparable et normale.
\end{enumerate}
\end{definition}

\noindent
Le lemme suivant ne semble trouver sa place nulle part ailleurs.

\begin{lemma}
\label{fields-lemma-pth-root}
Soit $K$ un corps de caractéristique $p > 0$. Soit $L/K$ une extension
algébrique séparable. Soit $\alpha \in L$.
\begin{enumerate}
\item Si les coefficients du polynôme minimal de $\alpha$
sur $K$ sont des puissances $p$-ièmes dans $K$, alors $\alpha$ est une puissance
$p$-ième dans $L$.
\item Plus généralement, si $P \in K[T]$ est un polynôme tel que (a) $\alpha$
soit une racine de $P$, (b) $P$ possède des racines deux à deux distinctes dans une clôture algébrique,
et (c) tous les coefficients de $P$ soient des puissances $p$-ièmes, alors $\alpha$ est une
puissance $p$-ième dans $L$.
\end{enumerate}
\end{lemma}

\begin{proof}
Il résulte des définitions que (2) implique (1). Supposons que $P$ satisfasse aux hypothèses de (2).
Écrivons $P(T) = \sum\nolimits_{i = 0}^d a_i T^{d - i}$ et $a_i = b_i^p$.
Le polynôme $Q(T) = \sum\nolimits_{i = 0}^d b_i T^{d - i}$ possède lui aussi des
racines distinctes dans une clôture algébrique, car les racines de $Q$
sont les racines $p$-ièmes des racines de $P$. Si $\alpha$ n'est pas une puissance $p$-ième,
alors $T^p - \alpha$ est un polynôme irréductible sur $L$
(Lemme \ref{fields-lemma-take-pth-root}).
De plus, $Q$ et $T^p - \alpha$ ont une racine commune dans
une clôture algébrique $\overline{L}$.
Ainsi $Q$ et $T^p - \alpha$ ne sont pas premiers entre eux, ce qui
implique que $T^p - \alpha | Q$ dans $L[T]$. Cela contredit le fait que
les racines de $Q$ sont deux à deux distinctes.
\end{proof}















% OK, start here.
%

\chapter{Algèbre commutative}



\phantomsection
\label{algebra-section-phantom}






\section{Introduction}
\label{algebra-section-introduction}

\noindent
Les notions de base d'algèbre commutative seront exposées dans ce document.
Une référence est \cite{MatCA}.






\section{Conventions}
\label{algebra-section-conventions}

\noindent
Un anneau est commutatif et possède un élément unité $1$. L'anneau nul est un anneau. En fait, c'est
le seul anneau qui ne possède pas d'idéal premier. Nous utiliserons le symbole
de Kronecker $\delta_{ij}$. Si $R \to S$ est un morphisme d'anneaux et si
$\mathfrak q$ est un idéal premier de $S$, nous employons alors la notation
``$\mathfrak p = R \cap \mathfrak q$''
pour désigner l'idéal premier qui est l'image réciproque de $\mathfrak q$ par
$R \to S$, même si $R$ n'est pas un sous-anneau de $S$ et même si $R \to S$
n'est pas injectif.






\section{Notions de base}
\label{algebra-section-rings-basic}

\noindent
Voici une liste de notions de base d'algèbre commutative. Certaines de ces
notions sont étudiées plus en détail dans la suite du texte et certaines sont
définies dans la liste, mais d'autres sont considérées comme élémentaires et ne seront pas définies.
Si la plupart des notions en italique ne vous sont pas familières, nous vous conseillons
de consulter un texte d'introduction à l'algèbre avant de poursuivre.

\begin{enumerate}
\item $R$ est un {\it anneau},
\label{algebra-item-ring}
\item $x\in R$ est {\it nilpotent},
\label{algebra-item-ring-element-nilpotent}
\item $x\in R$ est un {\it diviseur de zéro},
\label{algebra-item-ring-element-zerodivisor}
\item $x\in R$ est une {\it unité},
\label{algebra-item-ring-element-unit}
\item $e \in R$ est un {\it idempotent},
\label{algebra-item-ring-element-idempotent}
\item un idempotent $e \in R$ est dit {\it trivial} si $e = 1$ ou $e = 0$,
\label{algebra-item-idempotent-trivial}
\item $\varphi : R_1 \to R_2$ est un {\it homomorphisme d'anneaux},
\label{algebra-item-ring-homomorphism}
\item
\label{algebra-item-ring-homomorphism-finite-presentation}
$\varphi : R_1 \to R_2$ est {\it de présentation finie}, ou
{\it $R_2$ est une $R_1$-algèbre de présentation finie},
voir la Définition \ref{algebra-definition-finite-type},
\item
\label{algebra-item-ring-homomorphism-finite-type}
$\varphi : R_1 \to R_2$ est {\it de type fini}, ou
{\it $R_2$ est une $R_1$-algèbre de type fini},
voir la Définition \ref{algebra-definition-finite-type},
\item
\label{algebra-item-ring-homomorphism-finite}
$\varphi : R_1 \to R_2$ est {\it fini}, ou
{\it $R_2$ est une $R_1$-algèbre finie},
\item $R$ est un {\it anneau intègre},
\label{algebra-item-ring-domain}
\item $R$ est {\it réduit},
\label{algebra-item-ring-reduced}
\item $R$ est {\it noethérien},
\label{algebra-item-ring-Noetherian}
\item $R$ est un {\it anneau principal} ou un {\it PID},
\label{algebra-item-ring-PID}
\item $R$ est un {\it anneau euclidien},
\label{algebra-item-ring-Euclidean}
\item $R$ est un {\it anneau factoriel} ou un {\it UFD},
\label{algebra-item-ring-UFD}
\item $R$ est un {\it anneau de valuation discrète} ou un {\it dvr},
\label{algebra-item-ring-dvr}
\item $K$ est un {\it corps},
\label{algebra-item-field}
\item $L/K$ est une {\it extension de corps},
\label{algebra-item-field-extension}
\item $L/K$ est une {\it extension algébrique de corps},
\label{algebra-item-field-extension-algebraic}
\item $\{t_i\}_{i\in I}$ est une {\it base de transcendance} de $L$ sur $K$,
\label{algebra-item-transcendence-basis}
\item le {\it degré de transcendance} $\text{trdeg}(L/K)$ de $L$ sur $K$,
\label{algebra-item-transcendence-degree}
\item le corps $k$ est {\it algébriquement clos},
\label{algebra-item-algebraically-closed}
\item
\label{algebra-item-extend-into-algebraically-closed}
si $L/K$ est algébrique et si $\Omega/K$ est une extension avec $\Omega$
algébriquement clos,
alors il existe un morphisme d'anneaux $L \to \Omega$ qui prolonge le morphisme sur $K$,
\item $I \subset R$ est un {\it idéal},
\label{algebra-item-ideal}
\item $I \subset R$ est {\it radical},
\label{algebra-item-ideal-radical}
\item si $I$ est un idéal, nous disposons de son {\it radical} $\sqrt{I}$,
\label{algebra-item-radical-ideal}
\item
\label{algebra-item-ideal-nilpotent}
$I \subset R$ est {\it nilpotent} signifie que $I^n = 0$ pour
un certain $n \in \mathbf{N}$,
\item
\label{algebra-item-ideal-locally-nilpotent}
$I \subset R$ est {\it localement nilpotent} signifie que tout
élément de $I$ est nilpotent,
\item $\mathfrak p \subset R$ est un {\it idéal premier},
\label{algebra-item-prime-ideal}
\item
\label{algebra-item-prime-product-ideals}
si $\mathfrak p \subset R$ est premier, si $I, J \subset R$
sont des idéaux et si $IJ\subset \mathfrak p$, alors
$I \subset \mathfrak p$ ou $J \subset \mathfrak p$.
\item $\mathfrak m \subset R$ est un {\it idéal maximal},
\label{algebra-item-maximal-ideal}
\item tout anneau non nul possède un idéal maximal,
\label{algebra-item-exists-maximal-ideal}
\item
\label{algebra-item-jacobson-radical}
le {\it radical de Jacobson} de $R$ est $\text{rad}(R) =
\bigcap_{\mathfrak m \subset R} \mathfrak m$, l'intersection
de tous les idéaux maximaux de $R$,
\item l'idéal $(T)$ {\it engendré} par un sous-ensemble $T \subset R$,
\label{algebra-item-ideal-generated-by}
\item l'{\it anneau quotient} $R/I$,
\label{algebra-item-quotient-ring}
\item un idéal $I$ de l'anneau $R$ est premier si et seulement si $R/I$ est intègre,
\label{algebra-item-characterize-prime-ideal}
\item
\label{algebra-item-characterize-maximal-ideal}
un idéal $I$ de l'anneau $R$ est maximal si et seulement si
l'anneau $R/I$ est un corps,
\item
\label{algebra-item-inverse-image-ideal}
si $\varphi : R_1 \to R_2$ est un homomorphisme d'anneaux et si
$I \subset R_2$ est un idéal, alors $\varphi^{-1}(I)$ est un
idéal de $R_1$,
\item
\label{algebra-item-image-ideal}
si $\varphi : R_1 \to R_2$ est un homomorphisme d'anneaux et si
$I \subset R_1$ est un idéal, alors $\varphi(I) \cdot R_2$ (parfois
noté $I \cdot R_2$ ou $IR_2$) est l'idéal de $R_2$ engendré
par $\varphi(I)$,
\item
\label{algebra-item-inverse-image-prime}
si $\varphi : R_1 \to R_2$ est un homomorphisme d'anneaux et si
$\mathfrak p \subset R_2$ est un idéal premier, alors
$\varphi^{-1}(\mathfrak p)$ est un idéal premier de $R_1$,
\item $M$ est un {\it $R$-module},
\label{algebra-item-module}
\item
\label{algebra-item-annihilator}
pour $m \in M$, l'{\it annulateur}
$I = \{f \in R \mid fm = 0\}$ de $m$ dans $R$,
\item $N \subset M$ est un {\it $R$-sous-module},
\label{algebra-item-submodule}
\item $M$ est un {\it $R$-module noethérien},
\label{algebra-item-Noetherian-module}
\item $M$ est un {\it $R$-module de type fini},
\label{algebra-item-finite-module}
\item $M$ est un {\it $R$-module engendré par un nombre fini d'éléments},
\label{algebra-item-finitely-generated-module}
\item $M$ est un {\it $R$-module de présentation finie},
\label{algebra-item-finitely-presented-module}
\item $M$ est un {\it $R$-module libre},
\label{algebra-item-free-module}
\item
\label{algebra-item-extension-free}
si $0 \to K \to L \to M \to 0$ est une suite exacte courte
de $R$-modules et si $K$, $M$ sont libres, alors $L$ est libre,
\item si $N \subset M \subset L$ sont des $R$-modules, alors $L/M = (L/N)/(M/N)$,
\label{algebra-item-isomorphism-theorem}
\item $S$ est une {\it partie multiplicative de $R$},
\label{algebra-item-multiplicative-subset}
\item la {\it localisation} $R \to S^{-1}R$ de $R$,
\label{algebra-item-localization-ring}
\item
\label{algebra-item-localization-zero}
si $R$ est un anneau et si $S$ est une partie multiplicative
de $R$, alors $S^{-1}R$ est l'anneau nul si et seulement si $S$ contient $0$,
\item
\label{algebra-item-localize-nonzerodivisors}
si $R$ est un anneau et si la partie multiplicative $S$
est entièrement constituée d'éléments non diviseurs de zéro, alors $R \to S^{-1}R$
est injectif,
\item si $\varphi : R_1 \to R_2$ est un homomorphisme d'anneaux et si
$S$ est une partie multiplicative de $R_1$, alors $\varphi(S)$ est
une partie multiplicative de $R_2$,
\item
\label{algebra-item-products-multiplicative-subsets}
si $S$, $S'$ sont des parties multiplicatives de $R$,
et si $SS'$ désigne l'ensemble des produits $SS' =
\{r \in R \mid \exists s\in S, \exists s' \in S', r = ss'\}$,
alors $SS'$ est une partie multiplicative de $R$,
\item
\label{algebra-item-localization-localization}
si $S$, $S'$ sont des parties multiplicatives de $R$,
et si $\overline{S}$ désigne l'image de $S$ dans $(S')^{-1}R$,
alors $(SS')^{-1}R = \overline{S}^{-1}((S')^{-1}R)$,
\item la {\it localisation} $S^{-1}M$ du $R$-module $M$,
\label{algebra-item-localization-module}
\item
\label{algebra-item-localization-exact}
le foncteur $M \mapsto S^{-1}M$ préserve les applications injectives,
les applications surjectives et l'exactitude,
\item
\label{algebra-item-localization-localization-module}
si $S$, $S'$ sont des parties multiplicatives de $R$,
et si $M$ est un $R$-module, alors
$(SS')^{-1}M = S^{-1}((S')^{-1}M)$,
\item
\label{algebra-item-localize-ideal}
si $R$ est un anneau, si $I$ est un idéal de $R$ et si $S$ est une partie multiplicative
de $R$, alors $S^{-1}I$ est un idéal de $S^{-1}R$, et nous avons
$S^{-1}R/S^{-1}I = \overline{S}^{-1}(R/I)$, où $\overline{S}$
est l'image de $S$ dans $R/I$,
\item
\label{algebra-item-ideal-in-localization}
si $R$ est un anneau et si $S$ est une partie multiplicative
de $R$, alors tout idéal $I'$ de $S^{-1}R$ est
de la forme $S^{-1}I$, où l'on peut prendre pour $I$
l'image réciproque de $I'$ dans $R$,
\item
\label{algebra-item-submodule-in-localization}
si $R$ est un anneau, si $M$ est un $R$-module et si $S$ est une partie multiplicative
de $R$, alors tout sous-module $N'$ de $S^{-1}M$ est de la forme
$S^{-1}N$ pour un certain sous-module $N \subset M$, où
l'on peut prendre pour $N$ l'image réciproque de $N'$ dans $M$,
\item si $S = \{1, f, f^2, \ldots\}$, alors $R_f = S^{-1}R$ et $M_f = S^{-1}M$,
\label{algebra-item-localize-f}
\item
\label{algebra-item-localize-p}
si $S = R \setminus \mathfrak p = \{x\in R \mid x\not\in \mathfrak p\}$
pour un certain idéal premier $\mathfrak p$,
alors on note usuellement $R_{\mathfrak p} = S^{-1}R$
et $M_{\mathfrak p} = S^{-1}M$,
\item un {\it anneau local} est un anneau qui possède exactement un idéal maximal,
\label{algebra-item-local-ring}
\item un {\it anneau semi-local} est un anneau qui possède un nombre fini d'idéaux maximaux,
\label{algebra-item-semi-local-ring}
\item
\label{algebra-item-localize-p-local-ring}
si $\mathfrak p$ est un idéal premier de $R$, alors $R_{\mathfrak p}$ est
un anneau local d'idéal maximal $\mathfrak p R_{\mathfrak p}$,
\item
\label{algebra-item-residue-field}
le {\it corps résiduel}, noté $\kappa(\mathfrak p)$,
de l'idéal premier $\mathfrak p$ de l'anneau $R$ est le
corps des fractions de l'anneau intègre $R/\mathfrak p$ ;
il est égal à $R_\mathfrak p/\mathfrak pR_\mathfrak p
= (R \setminus \mathfrak p)^{-1}R/\mathfrak p$,
\item étant donnés $R$ et $M_1$, $M_2$, le {\it produit tensoriel} $M_1 \otimes_R M_2$,
\label{algebra-item-tensor-product}
\item
\label{algebra-item-cauchy-binet}
étant données des matrices $A$ et $B$ dans un anneau $R$, de tailles $m \times n$ et
$n \times m$, nous avons $\det(AB) = \sum \det(A_S)\det({}_SB)$ dans $R$, où
la somme porte sur les sous-ensembles $S \subset \{1, \ldots, n\}$ de cardinal $m$,
et $A_S$ est la sous-matrice $m \times m$ de $A$ formée des colonnes
correspondant à $S$, tandis que ${}_SB$ est la sous-matrice $m \times m$ de $B$
formée des lignes correspondant à $S$,
\item etc.
\end{enumerate}




\section{Lemme du serpent}
\label{algebra-section-snake}

\noindent
Le lemme du serpent et ses variantes sont étudiés dans le cadre des
catégories abéliennes dans
Homologie, section \ref{homology-section-abelian-categories}.

\begin{lemma}
\label{algebra-lemma-snake}
\begin{reference}
\cite[III, Lemma 3.3]{Cartan-Eilenberg}
\end{reference}
Étant donné un diagramme commutatif
$$
\xymatrix{
& X \ar[r] \ar[d]^\alpha &
Y \ar[r] \ar[d]^\beta &
Z \ar[r] \ar[d]^\gamma &
0 \\
0 \ar[r] & U \ar[r] & V \ar[r] & W
}
$$
de groupes abéliens dont les lignes sont exactes, il existe une suite exacte canonique
$$
\Ker(\alpha) \to \Ker(\beta) \to \Ker(\gamma)
\to
\Coker(\alpha) \to \Coker(\beta) \to \Coker(\gamma)
$$
De plus, si $X \to Y$ est injectif, la première application est
injective ; si $V \to W$ est surjectif, la dernière
application est surjective.
\end{lemma}

\begin{proof}
L'application $\partial : \Ker(\gamma) \to \Coker(\alpha)$ est définie
comme suit. Soit $z \in \Ker(\gamma)$. Choisissons $y \in Y$ d'image $z$.
Alors $\beta(y) \in V$ s'envoie sur zéro dans $W$. Ainsi, $\beta(y)$ est l'image d'un
élément $u \in U$. Posons $\partial z = \overline{u}$, la classe de $u$ dans le
conoyau de $\alpha$. La démonstration de l'exactitude est omise.
\end{proof}

\section{Modules de type fini et modules de présentation finie}
\label{algebra-section-module-finite-type}

\noindent
Voici simplement quelques notations et lemmes élémentaires.

\begin{definition}
\label{algebra-definition-module-finite-type}
Soit $R$ un anneau. Soit $M$ un $R$-module.
\begin{enumerate}
\item Nous disons que $M$ est un {\it $R$-module de type fini}, ou un {\it
$R$-module engendré par un nombre fini d'éléments}, s'il existe $n \in \mathbf{N}$ et $x_1, \ldots, x_n \in M$
tels que tout élément de $M$ soit une combinaison $R$-linéaire des $x_i$.
De manière équivalente, cela signifie qu'il existe une surjection
$R^{\oplus n} \to M$ pour un certain $n \in \mathbf{N}$.
\item Nous disons que $M$ est un {\it $R$-module de présentation finie} ou un
{\it $R$-module de présentation finie} s'il existe des entiers
$n, m \in \mathbf{N}$ et une suite exacte
$$
R^{\oplus m} \longrightarrow R^{\oplus n} \longrightarrow M \longrightarrow 0
$$
\end{enumerate}
\end{definition}

\noindent
De manière informelle, $M$ est un $R$-module de présentation finie si et seulement s'il
est de type fini et si le module des relations entre ces
générateurs est lui aussi de type fini.
Le choix d'une suite exacte comme dans la définition s'appelle une
{\it présentation} de $M$.

\begin{lemma}
\label{algebra-lemma-lift-map}
Soit $R$ un anneau. Soient $\alpha : R^{\oplus n} \to M$ et $\beta : N \to M$ des
homomorphismes de modules. Si $\Im(\alpha) \subset \Im(\beta)$, il existe alors
un homomorphisme de $R$-modules $\gamma : R^{\oplus n} \to N$ tel que
$\alpha = \beta \circ \gamma$.
\end{lemma}

\begin{proof}
Soit $e_i = (0, \ldots, 0, 1, 0, \ldots, 0)$ le $i$-ème vecteur de base
de $R^{\oplus n}$. Soit $x_i \in N$ un élément tel que
$\alpha(e_i) = \beta(x_i)$ ; il existe par hypothèse. Posons
$\gamma(a_1, \ldots, a_n) = \sum a_i x_i$. Par construction,
$\alpha = \beta \circ \gamma$.
\end{proof}

\begin{lemma}
\label{algebra-lemma-extension}
Soit $R$ un anneau.
Soit
$$
0 \to M_1 \to M_2 \to M_3 \to 0
$$
une suite exacte courte de $R$-modules.
\begin{enumerate}
\item Si $M_1$ et $M_3$ sont des $R$-modules de type fini, alors $M_2$ est un
$R$-module de type fini.
\item Si $M_1$ et $M_3$ sont des $R$-modules de présentation finie, alors $M_2$
est un $R$-module de présentation finie.
\item Si $M_2$ est un $R$-module de type fini, alors $M_3$ est un $R$-module de type fini.
\item Si $M_2$ est un $R$-module de présentation finie et si $M_1$ est un
$R$-module de type fini, alors $M_3$ est un $R$-module de présentation finie.
\item Si $M_3$ est un $R$-module de présentation finie et si $M_2$ est un
$R$-module de type fini, alors $M_1$ est un $R$-module de type fini.
\end{enumerate}
\end{lemma}

\begin{proof}
Démonstration de (1). Si $x_1, \ldots, x_n$ sont des générateurs de $M_1$ et si
$y_1, \ldots, y_m \in M_2$ sont des éléments dont les images dans $M_3$ sont des
générateurs de $M_3$, alors $x_1, \ldots, x_n, y_1, \ldots, y_m$
engendrent $M_2$.

\medskip\noindent
L'assertion (3) découle immédiatement de la définition.

\medskip\noindent
Démonstration de (5). Supposons $M_3$ de présentation finie et $M_2$ de type fini.
Choisissons une présentation
$$
R^{\oplus m} \to R^{\oplus n} \to M_3 \to 0
$$
D'après le Lemme \ref{algebra-lemma-lift-map}, il existe une application
$R^{\oplus n} \to M_2$ telle que
le diagramme en traits pleins
$$
\xymatrix{
& R^{\oplus m} \ar[r] \ar@{..>}[d] & R^{\oplus n} \ar[r] \ar[d] &
M_3 \ar[r] \ar[d]^{\text{id}} & 0 \\
0 \ar[r] & M_1 \ar[r] & M_2 \ar[r] & M_3 \ar[r] & 0
}
$$
soit commutatif. On obtient ainsi la flèche en pointillé. Par le lemme du serpent
(Lemme \ref{algebra-lemma-snake}), on voit que l'on obtient un isomorphisme
$$
\Coker(R^{\oplus m} \to M_1)
\cong
\Coker(R^{\oplus n} \to M_2)
$$
En particulier, nous en déduisons que $\Coker(R^{\oplus m} \to M_1)$
est un $R$-module de type fini. Puisque $\Im(R^{\oplus m} \to M_1)$
est de type fini d'après (3), on voit que $M_1$ est de type fini d'après l'assertion (1).

\medskip\noindent
Démonstration de (4). Supposons $M_2$ de présentation finie et $M_1$ de type fini.
Choisissons une présentation $R^{\oplus m} \to R^{\oplus n} \to M_2 \to 0$.
Choisissons une surjection $R^{\oplus k} \to M_1$. D'après le Lemme \ref{algebra-lemma-lift-map},
il existe une factorisation $R^{\oplus k} \to R^{\oplus n} \to M_2$
de la composée $R^{\oplus k} \to M_1 \to M_2$. Alors
$R^{\oplus k + m} \to R^{\oplus n} \to M_3 \to 0$
est une présentation.

\medskip\noindent
Démonstration de (2). Supposons $M_1$ et $M_3$ de présentation finie.
L'argument de la démonstration de l'assertion (1) fournit un diagramme commutatif
$$
\xymatrix{
0 \ar[r] & R^{\oplus n} \ar[d] \ar[r] & R^{\oplus n + m} \ar[d] \ar[r] &
R^{\oplus m} \ar[d] \ar[r] & 0 \\
0 \ar[r] & M_1 \ar[r] & M_2 \ar[r] & M_3 \ar[r] & 0
}
$$
dont les flèches verticales sont surjectives. Par le lemme du serpent, nous obtenons une suite
exacte courte
$$
0 \to \Ker(R^{\oplus n} \to M_1) \to
\Ker(R^{\oplus n + m} \to M_2) \to
\Ker(R^{\oplus m} \to M_3) \to 0
$$
D'après l'assertion (5), les deux modules extrêmes sont de type fini. Par conséquent, celui du
milieu l'est aussi. D'après (4), on voit que $M_2$ est de présentation finie.
\end{proof}

\begin{lemma}
\label{algebra-lemma-trivial-filter-finite-module}
\begin{slogan}
Les modules de type fini admettent des filtrations dont les quotients successifs sont des
modules cycliques.
\end{slogan}
Soit $R$ un anneau et soit $M$ un $R$-module de type fini.
Il existe une filtration par des $R$-sous-modules de type fini
$$
0 = M_0 \subset M_1 \subset \ldots \subset M_n = M
$$
telle que chaque quotient $M_i/M_{i - 1}$ soit isomorphe
à $R/I_i$ pour un certain idéal $I_i$ de $R$.
\end{lemma}

\begin{proof}
Nous raisonnons par récurrence sur le nombre de générateurs de $M$. Soient
$x_1, \ldots, x_r \in M$ des générateurs.
Posons $M' = Rx_1 \subset M$. Alors $M/M'$ possède $r - 1$ générateurs
et l'hypothèse de récurrence s'applique. De plus, il est clair que $M' \cong R/I_1$
avec $I_1 = \{f \in R \mid fx_1 = 0\}$.
\end{proof}

\begin{lemma}
\label{algebra-lemma-finite-over-subring}
Soit $R \to S$ un morphisme d'anneaux.
Soit $M$ un $S$-module.
Si $M$ est de type fini comme $R$-module, alors $M$ est de type fini comme $S$-module.
\end{lemma}

\begin{proof}
En effet, tout système générateur sur $R$ de $M$ est aussi un système générateur sur $S$ de
$M$, puisque la structure de $R$-module est induite par l'image de $R$ dans $S$.
\end{proof}




\section{Morphismes d'anneaux de type fini et de présentation finie}
\label{algebra-section-finite-type}

\begin{definition}
\label{algebra-definition-finite-type}
Soit $R \to S$ un morphisme d'anneaux.
\begin{enumerate}
\item Nous disons que $R \to S$ est de {\it type fini}, ou que {\it $S$ est une
$R$-algèbre de type fini}, s'il existe un $n \in \mathbf{N}$ et une surjection
de $R$-algèbres $R[x_1, \ldots, x_n] \to S$.
\item Nous disons que $R \to S$ est de {\it présentation finie} s'il
existe des entiers $n, m \in \mathbf{N}$, des polynômes
$f_1, \ldots, f_m \in R[x_1, \ldots, x_n]$
et un isomorphisme de $R$-algèbres
$R[x_1, \ldots, x_n]/(f_1, \ldots, f_m) \cong S$.
\end{enumerate}
\end{definition}

\noindent
De manière informelle, $R \to S$ est de présentation finie si et seulement si
$S$ est de type fini comme $R$-algèbre
et si l'idéal des relations entre les générateurs est de type fini.
Le choix d'une surjection $R[x_1, \ldots, x_n] \to S$ comme dans la définition
est parfois appelé une {\it présentation} de $S$.

\begin{lemma}
\label{algebra-lemma-compose-finite-type}
Les notions de type fini et de présentation finie possèdent les propriétés
de permanence suivantes.
\begin{enumerate}
\item Une composée de morphismes d'anneaux de type fini est de type fini.
\item Une composée de morphismes d'anneaux de présentation finie est de présentation
finie.
\item Étant donné $R \to S' \to S$ avec $R \to S$ de type fini,
le morphisme $S' \to S$ est alors de type fini.
\item Étant donné $R \to S' \to S$, avec $R \to S$ de présentation finie
et $R \to S'$ de type fini, le morphisme $S' \to S$ est alors de présentation finie.
\end{enumerate}
\end{lemma}

\begin{proof}
Nous ne démontrons que la dernière assertion.
Écrivons $S = R[x_1, \ldots, x_n]/(f_1, \ldots, f_m)$
et $S' = R[y_1, \ldots, y_a]/I$. Supposons que la classe
$\bar y_i$ de $y_i$ s'envoie
sur $h_i \bmod (f_1, \ldots, f_m)$ dans $S$.
Il est alors clair que
$S = S'[x_1, \ldots, x_n]/(f_1, \ldots, f_m,
h_1 - \bar y_1, \ldots, h_a - \bar y_a)$.
\end{proof}

\begin{lemma}
\label{algebra-lemma-finite-presentation-independent}
Soit $R \to S$ un morphisme d'anneaux de présentation finie.
Pour toute surjection $\alpha : R[x_1, \ldots, x_n] \to S$, le
noyau de $\alpha$ est un idéal de type fini de $R[x_1, \ldots, x_n]$.
\end{lemma}

\begin{proof}
Écrivons $S = R[y_1, \ldots, y_m]/(f_1, \ldots, f_k)$.
Choisissons des $g_i \in R[y_1, \ldots, y_m]$ qui relèvent
les $\alpha(x_i)$. On voit alors que $S = R[x_i, y_j]/(f_l, x_i - g_i)$.
Choisissons $h_j \in R[x_1, \ldots, x_n]$ tel que $\alpha(h_j)$
corresponde à $y_j \bmod (f_1, \ldots, f_k)$. Considérons
l'application $\psi : R[x_i, y_j] \to R[x_i]$, $x_i \mapsto x_i$,
$y_j \mapsto h_j$. Le noyau de $\alpha$
est alors l'image de $(f_l, x_i - g_i)$ par $\psi$, ce qui conclut.
\end{proof}

\begin{lemma}
\label{algebra-lemma-finitely-presented-over-subring}
Soit $R \to S$ un morphisme d'anneaux.
Soit $M$ un $S$-module.
Supposons $R \to S$ de type fini et
$M$ de présentation finie comme $R$-module.
Alors $M$ est de présentation finie comme $S$-module.
\end{lemma}

\begin{proof}
La démonstration est semblable à celle de l'assertion (4) du
Lemme \ref{algebra-lemma-compose-finite-type}.
Nous pouvons supposer $S = R[x_1, \ldots, x_n]/J$.
Choisissons $y_1, \ldots, y_m \in M$ qui engendrent $M$ comme $R$-module
et des relations $\sum a_{ij} y_j = 0$, $i = 1, \ldots, t$, qui
engendrent le noyau de $R^{\oplus m} \to M$. Pour tous
$i = 1, \ldots, n$ et $j = 1, \ldots, m$, écrivons
$$
x_i y_j = \sum a_{ijk} y_k
$$
pour certains $a_{ijk} \in R$. Considérons le $S$-module $N$ engendré par
$y_1, \ldots, y_m$ et soumis aux relations
$\sum a_{ij} y_j = 0$, $i = 1, \ldots, t$, ainsi que
$x_i y_j = \sum a_{ijk} y_k$, $i = 1, \ldots, n$ et $j = 1, \ldots, m$.
Alors $N$ possède une présentation
$$
S^{\oplus nm + t} \longrightarrow S^{\oplus m} \longrightarrow N
\longrightarrow 0
$$
Par construction, il existe une application surjective $\varphi : N \to M$.
Pour achever la démonstration, montrons que $\varphi$ est injective.
Supposons $z = \sum b_j y_j \in N$ pour certains $b_j \in S$.
Nous pouvons considérer les $b_j$ comme des polynômes en $x_1, \ldots, x_n$
à coefficients dans $R$.
En appliquant les relations de la forme $x_i y_j = \sum a_{ijk} y_k$,
nous pouvons abaisser par récurrence le degré des polynômes.
Nous voyons donc que $z = \sum c_j y_j$ pour certains $c_j \in R$.
Par conséquent, si $\varphi(z) = 0$, le vecteur $(c_1, \ldots, c_m)$
est une combinaison $R$-linéaire des vecteurs $(a_{i1}, \ldots, a_{im})$,
et nous concluons que $z = 0$, comme souhaité.
\end{proof}




\section{Morphismes d'anneaux finis}
\label{algebra-section-finite}

\noindent
Voici la définition.

\begin{definition}
\label{algebra-definition-finite-ring-map}
Soit $\varphi : R \to S$ un morphisme d'anneaux. Nous disons que $\varphi : R \to S$ est
{\it fini} si $S$ est de type fini comme $R$-module.
\end{definition}

\begin{lemma}
\label{algebra-lemma-finite-module-over-finite-extension}
Soit $R \to S$ un morphisme d'anneaux fini.
Soit $M$ un $S$-module.
Alors $M$ est de type fini comme $R$-module si et seulement si $M$ est de type fini
comme $S$-module.
\end{lemma}

\begin{proof}
L'une des implications résulte du
Lemme \ref{algebra-lemma-finite-over-subring}.
Pour voir l'autre, supposons que $M$ soit de type fini comme $S$-module.
Choisissons $x_1, \ldots, x_n \in S$ qui engendrent $S$ comme $R$-module.
Choisissons $y_1, \ldots, y_m \in M$ qui engendrent $M$ comme $S$-module.
Alors les $x_i y_j$ engendrent $M$ comme $R$-module.
\end{proof}

\begin{lemma}
\label{algebra-lemma-finite-transitive}
Supposons que $R \to S$ et $S \to T$ soient des morphismes d'anneaux finis.
Alors $R \to T$ est fini.
\end{lemma}

\begin{proof}
Si les $t_i$ engendrent $T$ comme $S$-module et si les $s_j$ engendrent $S$ comme
$R$-module, alors les $t_i s_j$ engendrent $T$ comme $R$-module.
(Cela résulte également du
Lemme \ref{algebra-lemma-finite-module-over-finite-extension}.)
\end{proof}

\begin{lemma}
\label{algebra-lemma-finite-finite-type}
Soit $\varphi : R \to S$ un morphisme d'anneaux.
\begin{enumerate}
\item Si $\varphi$ est fini, alors $\varphi$ est de type fini.
\item Si $S$ est de présentation finie comme $R$-module, alors
$\varphi$ est de présentation finie.
\end{enumerate}
\end{lemma}

\begin{proof}
Pour (1), si $x_1, \ldots, x_n \in S$ engendrent $S$ comme $R$-module,
alors $x_1, \ldots, x_n$ engendrent $S$ comme $R$-algèbre. Pour (2),
supposons que $\sum r_j^ix_i = 0$, $j = 1, \ldots, m$, soit un système
de générateurs des relations entre les $x_i$ considérés comme des
générateurs du $R$-module $S$. Écrivons en outre
$1 = \sum r_ix_i$ pour certains $r_i \in R$, et
$x_ix_j = \sum r_{ij}^k x_k$ pour certains $r_{ij}^k \in R$.
Alors
$$
S = R[t_1, \ldots, t_n]/
(\sum r_j^it_i,\ 1 - \sum r_it_i,\ t_it_j - \sum r_{ij}^k t_k)
$$
comme $R$-algèbre, ce qui démontre (2).
\end{proof}

\noindent
Pour plus de renseignements sur les morphismes d'anneaux finis, voir
la section \ref{algebra-section-finite-ring-extensions}.


\section{Colimites}
\label{algebra-section-colimits}

\noindent
Une partie du contenu de cette section recoupe l'étude générale
des colimites dans
Catégories, sections \ref{categories-section-limits} --
\ref{categories-section-posets-limits}.
La notion d'ensemble préordonné est définie dans
Catégories, Définition \ref{categories-definition-directed-set}.
Elle est légèrement plus faible que celle d'ensemble partiellement ordonné.

\begin{definition}
\label{algebra-definition-directed-system}
Soit $(I, \leq)$ un ensemble préordonné.
Un {\it système $(M_i, \mu_{ij})$ de $R$-modules sur $I$}
est constitué d'une famille de $R$-modules $\{M_i\}_{i\in I}$ indexée
par $I$ et d'une famille d'homomorphismes de $R$-modules $\{\mu_{ij} : M_i \to M_j\}_{i \leq j}$
telles que, pour tous $i \leq j \leq k$,
$$
\mu_{ii} = \text{id}_{M_i}\quad
\mu_{ik} = \mu_{jk}\circ \mu_{ij}
$$
Nous disons que $(M_i, \mu_{ij})$ est un {\it système filtrant} si $I$ est un ensemble filtrant.
\end{definition}

\noindent
C'est la même notion que celle définie dans Catégories,
Définition \ref{categories-definition-system-over-poset}
et section \ref{categories-section-posets-limits}.
Nous renvoyons à Catégories, Définition \ref{categories-definition-colimit},
pour la définition d'une colimite d'un diagramme ou d'un système dans une
catégorie quelconque.

\begin{lemma}
\label{algebra-lemma-colimit}
Soit $(M_i, \mu_{ij})$ un système de $R$-modules sur l'ensemble préordonné $I$.
La colimite du système $(M_i, \mu_{ij})$ est le $R$-module quotient
$(\bigoplus_{i\in I} M_i) /Q$, où $Q$ est le
$R$-sous-module engendré par tous les éléments
$$
\iota_i(x_i) - \iota_j(\mu_{ij}(x_i))
$$
où $\iota_i : M_i \to \bigoplus_{i\in I} M_i$
est l'inclusion canonique. Nous notons la colimite
$M = \colim_i M_i$. Nous notons
$\pi : \bigoplus_{i\in I} M_i \to M$ la
projection et
$\phi_i = \pi \circ \iota_i : M_i \to M$.
\end{lemma}

\begin{proof}
Ce lemme est un cas particulier de
Catégories, Lemme \ref{categories-lemma-colimits-coproducts-coequalizers},
mais nous allons aussi le démontrer directement dans ce cas.
Remarquons en effet que $\phi_i = \phi_j\circ \mu_{ij}$ dans la
construction ci-dessus. Pour montrer que la paire $(M, \phi_i)$ est la colimite, nous devons
montrer qu'elle satisfait la propriété universelle : pour toute autre paire de ce type
$(Y, \psi_i)$ avec $\psi_i : M_i \to
Y$, $\psi_i = \psi_j\circ \mu_{ij}$, il existe un unique homomorphisme de
$R$-modules $g : M \to Y$ tel que le
diagramme suivant soit commutatif :
$$
\xymatrix{
M_i \ar[rr]^{\mu_{ij}} \ar[dr]^{\phi_i} \ar[ddr]_{\psi_i} & &
M_j\ar[dl]_{\phi_j} \ar[ddl]^{\psi_j} \\
& M \ar[d]^{g}\\
& Y
}
$$
Cela est clair, car nous pouvons définir $g$ en prenant
l'application $\psi_i$ sur le facteur $M_i$ de la somme directe
$\bigoplus M_i$.
\end{proof}

\begin{lemma}
\label{algebra-lemma-directed-colimit}
Soit $(M_i, \mu_{ij})$ un système de $R$-modules sur
l'ensemble préordonné $I$. Supposons que $I$ soit filtrant.
La colimite du système $(M_i, \mu_{ij})$ est canoniquement
isomorphe au module $M$ défini comme suit :
\begin{enumerate}
\item comme ensemble, posons
$$
M = \left(\coprod\nolimits_{i \in I} M_i\right)/\sim
$$
où, pour $m \in M_i$ et $m' \in M_{i'}$, nous avons
$$
m \sim m' \Leftrightarrow
\mu_{ij}(m) = \mu_{i'j}(m')\text{ pour un certain }j \geq i, i'
$$
\item comme groupe abélien, pour $m \in M_i$ et $m' \in M_{i'}$,
nous définissons la somme des classes de $m$ et $m'$ dans $M$
comme étant la classe de $\mu_{ij}(m) + \mu_{i'j}(m')$, où
$j \in I$ est un indice quelconque tel que $i \leq j$ et $i' \leq j$, et
\item comme $R$-module, pour $m \in M_i$ et $x \in R$, définissons
le produit de $x$ par la classe de $m$ dans $M$ comme étant la
classe de $xm$ dans $M$.
\end{enumerate}
Les applications canoniques $\phi_i : M_i \to M$ sont induites par les applications
canoniques $M_i \to \coprod_{i \in I} M_i$.
\end{lemma}

\begin{proof}
Démonstration omise. Comparer à
Catégories, section \ref{categories-section-directed-colimits}.
\end{proof}

\begin{lemma}
\label{algebra-lemma-zero-directed-limit}
Soit $(M_i, \mu_{ij})$ un système filtrant.
Soit $M = \colim M_i$ avec $\mu_i : M_i \to M$.
Alors $\mu_i(x_i) = 0$ pour $x_i \in M_i$ si et seulement s'il
existe $j \geq i$ tel que $\mu_{ij}(x_i) = 0$.
\end{lemma}

\begin{proof}
Cela résulte clairement de la description de la colimite filtrante
du Lemme \ref{algebra-lemma-directed-colimit}.
\end{proof}

\begin{example}
\label{algebra-example-zero-colimit-different}
Considérons l'ensemble partiellement ordonné $I = \{a, b, c\}$ avec
$a < b$ et $a < c$, et aucune autre inégalité stricte.
Un système $(M_a, M_b, M_c, \mu_{ab}, \mu_{ac})$
sur $I$ est constitué de trois $R$-modules $M_a, M_b, M_c$
et de deux homomorphismes de $R$-modules $\mu_{ab} : M_a \to M_b$ et
$\mu_{ac} : M_a \to M_c$.
La colimite du système est simplement
$$
M := \colim_{i \in I} M_i = \Coker(M_a \to M_b \oplus M_c)
$$
où l'application est $\mu_{ab} \oplus -\mu_{ac}$. Le noyau de
l'application canonique $M_a \to M$ est donc $\Ker(\mu_{ab}) + \Ker(\mu_{ac})$.
Et le noyau de l'application canonique $M_b \to M$ est l'image de
$\Ker(\mu_{ac})$ par l'application $\mu_{ab}$. Par conséquent, il est clair que
le résultat du Lemme \ref{algebra-lemma-zero-directed-limit} est faux pour les
systèmes généraux.
\end{example}

\begin{definition}
\label{algebra-definition-homomorphism-directed-systems}
Soient $(M_i, \mu_{ij})$, $(N_i, \nu_{ij})$ des
systèmes de $R$-modules sur le même ensemble préordonné $I$.
Un {\it homomorphisme de systèmes} $\Phi$ de $(M_i, \mu_{ij})$ vers
$(N_i, \nu_{ij})$ est, par définition, une famille d'homomorphismes de $R$-modules
$\phi_i : M_i \to N_i$
tels que $\phi_j \circ \mu_{ij} = \nu_{ij} \circ \phi_i$
pour tous $i \leq j$.
\end{definition}

\noindent
C'est la même notion que celle de transformation naturelle de foncteurs
entre les diagrammes associés $M : I \to \text{Mod}_R$
et $N : I \to \text{Mod}_R$, dans le langage des
catégories.
Le lemme suivant est un cas particulier de
Catégories, Lemme \ref{categories-lemma-functorial-colimit}.

\begin{lemma}
\label{algebra-lemma-homomorphism-limit}
Soient $(M_i, \mu_{ij})$, $(N_i, \nu_{ij})$ des
systèmes de $R$-modules sur le même ensemble préordonné.
Un morphisme de systèmes $\Phi = (\phi_i)$ de $(M_i, \mu_{ij})$ vers
$(N_i, \nu_{ij})$ induit un unique homomorphisme
$$
\colim \phi_i : \colim M_i \longrightarrow \colim N_i
$$
tel que
$$
\xymatrix{
M_i \ar[r] \ar[d]_{\phi_i} & \colim M_i \ar[d]^{\colim \phi_i} \\
N_i \ar[r] & \colim N_i
}
$$
soit commutatif pour tout $i \in I$.
\end{lemma}

\begin{proof}
Écrivons $M = \colim M_i$, $N = \colim N_i$ et $\phi = \colim \phi_i$
(ce dernier restant à construire). Nous utiliserons sans autre mention la description explicite de $M$
et de $N$ donnée au Lemme \ref{algebra-lemma-colimit}.
La condition du lemme équivaut à ce que
$$
\xymatrix{
\bigoplus_{i\in I} M_i \ar[r] \ar[d]_{\bigoplus\phi_i} & M \ar[d]^\phi \\
\bigoplus_{i\in I} N_i \ar[r] & N
}
$$
soit commutatif. Il est donc clair que, si $\phi$ existe, il est unique.
Pour voir que $\phi$ existe, il suffit de montrer que le noyau de la
flèche horizontale supérieure est envoyé par $\bigoplus \phi_i$ dans le noyau
de la flèche horizontale inférieure. Pour cela, soient $j \leq k$ et
$x_j \in M_j$. Alors
$$
(\bigoplus \phi_i)(x_j - \mu_{jk}(x_j))
=
\phi_j(x_j) - \phi_k(\mu_{jk}(x_j))
=
\phi_j(x_j) - \nu_{jk}(\phi_j(x_j))
$$
appartient au noyau de la flèche horizontale inférieure, comme voulu.
\end{proof}

\begin{lemma}
\label{algebra-lemma-directed-colimit-exact}
\begin{slogan}
Les colimites filtrantes sont exactes. Les colimites sur un ensemble filtrant sont exactes.
\end{slogan}
Soit $I$ un ensemble filtrant.
Soient $(L_i, \lambda_{ij})$, $(M_i, \mu_{ij})$ et
$(N_i, \nu_{ij})$ des systèmes de $R$-modules sur $I$.
Soient $\varphi_i : L_i \to M_i$ et $\psi_i : M_i \to N_i$ des
morphismes de systèmes sur $I$. Supposons que, pour tout $i \in I$, la
suite de $R$-modules
$$
\xymatrix{
L_i \ar[r]^{\varphi_i} &
M_i \ar[r]^{\psi_i} &
N_i
}
$$
soit un complexe d'homologie $H_i$.
Alors les $R$-modules $H_i$ forment un système sur $I$,
la suite de $R$-modules
$$
\xymatrix{
\colim_i L_i \ar[r]^\varphi &
\colim_i M_i \ar[r]^\psi &
\colim_i N_i
}
$$
est elle aussi un complexe et, en notant $H$ son homologie, nous avons
$$
H = \colim_i H_i.
$$
\end{lemma}

\begin{proof}
Il est clair que
$
\xymatrix{
\colim_i L_i \ar[r]^\varphi &
\colim_i M_i \ar[r]^\psi &
\colim_i N_i
}
$
est un complexe. Pour chaque $i \in I$, il existe un morphisme canonique de
$R$-modules $H_i \to H$ (qui envoie chaque
$[m] \in H_i = \Ker(\psi_i) / \Im(\varphi_i)$ sur la
classe dans $H = \Ker(\psi) / \Im(\varphi)$
de l'image de $m$ dans $\colim_i M_i$). Ces morphismes donnent lieu
à un morphisme $\colim_i H_i \to H$. Il reste à
montrer que ce morphisme est surjectif et injectif.

\medskip\noindent
Nous allons utiliser à plusieurs reprises, sans autre mention, la description des colimites sur $I$
donnée au Lemme \ref{algebra-lemma-directed-colimit}.
Soit $h \in H$.
Puisque $H = \Ker(\psi)/\Im(\varphi)$, on voit que
$h$ est la classe modulo $\Im(\varphi)$ d'un élément $[m]$
de $\Ker(\psi) \subset \colim_i M_i$. Choisissons un
$i$ tel que $[m]$ provienne d'un élément $m \in M_i$. Choisissons
$j \geq i$ tel que $\nu_{ij}(\psi_i(m)) = 0$, ce qui est possible
puisque $[m] \in \Ker(\psi)$. Après avoir remplacé $i$ par $j$ et
$m$ par $\mu_{ij}(m)$, nous voyons que nous pouvons supposer $m \in \Ker(\psi_i)$.
Cela montre que l'application $\colim_i H_i \to H$ est surjective.

\medskip\noindent
Supposons qu'un $h_i \in H_i$ ait une image nulle dans $H$. Puisque
$H_i = \Ker(\psi_i)/\Im(\varphi_i)$, nous pouvons représenter
$h_i$ par un élément $m \in \Ker(\psi_i) \subset M_i$.
L'hypothèse d'annulation de $h_i$ dans $H$ signifie que
la classe de $m$ dans $\colim_i M_i$ appartient à l'image
de $\varphi$. Il existe donc un $j \geq i$ et un $l \in L_j$
tels que $\varphi_j(l) = \mu_{ij}(m)$. Cela montre clairement que
l'image de $h_i$ dans $H_j$ est nulle. Ceci démontre
l'injectivité de $\colim_i H_i \to H$.
\end{proof}

\begin{example}
\label{algebra-example-colimit-not-exact}
La prise des colimites n'est pas exacte en général.
Considérons l'ensemble partiellement ordonné $I = \{a, b, c\}$ avec
$a < b$ et $a < c$, et aucune autre inégalité stricte,
comme dans l'Exemple \ref{algebra-example-zero-colimit-different}.
Considérons l'application de systèmes
$(0, \mathbf{Z}, \mathbf{Z}, 0, 0) \to
(\mathbf{Z}, \mathbf{Z}, \mathbf{Z}, 1, 1)$.
D'après la description de la colimite dans
l'Exemple \ref{algebra-example-zero-colimit-different},
nous voyons que l'application associée entre les colimites n'est pas injective,
bien que l'application de systèmes soit injective sur chaque objet.
Par conséquent, le résultat du Lemme \ref{algebra-lemma-directed-colimit-exact}
est faux pour les systèmes généraux.
\end{example}

\begin{lemma}
\label{algebra-lemma-almost-directed-colimit-exact}
Soit $\mathcal{I}$ une catégorie d'indices satisfaisant aux hypothèses de
Catégories, Lemme \ref{categories-lemma-split-into-directed}.
Alors la prise des colimites de diagrammes de groupes abéliens sur $\mathcal{I}$
est exacte (c'est-à-dire que l'analogue du
Lemme \ref{algebra-lemma-directed-colimit-exact}
vaut dans cette situation).
\end{lemma}

\begin{proof}
D'après
Catégories, Lemme \ref{categories-lemma-split-into-directed},
nous pouvons écrire $\mathcal{I} = \coprod_{j \in J} \mathcal{I}_j$, où chaque
$\mathcal{I}_j$ est une catégorie filtrante et où $J$ peut être vide. D'après
Catégories, Lemme \ref{categories-lemma-directed-category-system},
prendre les colimites sur les catégories d'indices $\mathcal{I}_j$ revient à
prendre la colimite sur un certain ensemble filtrant. Le
Lemme \ref{algebra-lemma-directed-colimit-exact}
s'applique donc à ces colimites. Le problème se ramène ainsi à montrer que
les coproduits dans la catégorie des $R$-modules indexés par l'ensemble $J$ sont exacts.
Autrement dit, étant données des suites exactes
$L_j \to M_j \to N_j$ de $R$-modules, nous devons montrer que
$$
\bigoplus\nolimits_{j \in J} L_j
\longrightarrow
\bigoplus\nolimits_{j \in J} M_j
\longrightarrow
\bigoplus\nolimits_{j \in J} N_j
$$
est exacte. Cela se vérifie directement et reste vrai même si $J$ est vide.
\end{proof}






\section{Localisation}
\label{algebra-section-localization}

\begin{definition}
\label{algebra-definition-multiplicative-subset}
Soit $R$ un anneau et soit $S$ un sous-ensemble de $R$.
Nous disons que $S$ est une {\it partie multiplicative de $R$} si
$1\in S$ et si $S$ est stable
par multiplication, c'est-à-dire si $s, s' \in S \Rightarrow ss' \in S$.
\end{definition}

\noindent
Étant donnés un anneau $A$ et une partie multiplicative $S$, nous
définissons une relation sur $A \times S$ comme suit :
$$
(x, s) \sim (y, t)
\Leftrightarrow
\exists u \in S \text{ tel que } (xt-ys)u = 0
$$
On vérifie aisément qu'il s'agit d'une relation d'équivalence.
Notons $x/s$ (ou $\frac{x}{s}$) la classe d'équivalence de $(x, s)$ et
$S^{-1}A$ l'ensemble de toutes les classes d'équivalence. Définissons l'addition
et la multiplication dans $S^{-1}A$ comme suit :
$$
x/s + y/t = (xt + ys)/st, \quad
x/s \cdot y/t = xy/st
$$
On vérifie que ces opérations font de $S^{-1}A$ un anneau.

\begin{definition}
\label{algebra-definition-localization}
Cet anneau est appelé la {\it localisation de $A$ relativement à $S$}.
\end{definition}

\noindent
Il existe un morphisme naturel d'anneaux de $A$ vers sa localisation $S^{-1}A$,
$$
A \longrightarrow S^{-1}A, \quad x \longmapsto x/1
$$
parfois appelé le {\it morphisme de localisation}. En général, le
morphisme de localisation n'est pas injectif, à moins que $S$ ne contienne aucun diviseur de zéro.
En effet, si $x/1 = 0$, il existe un $u\in S$ tel que $xu = 0$ dans $A$, et
donc $x = 0$ puisqu'il n'y a pas de diviseur de zéro dans $S$.
La localisation d'un anneau possède la propriété universelle suivante.

\begin{proposition}
\label{algebra-proposition-universal-property-localization}
Soit $f : A \to B$ un morphisme d'anneaux qui envoie chaque élément de $S$ sur une unité
de $B$. Il existe alors un unique homomorphisme $g : S^{-1}A \to B$ tel
que le diagramme suivant soit commutatif.
$$
\xymatrix{
A \ar[rr]^{f} \ar[dr] & & B \\
& S^{-1}A \ar[ur]_g
}
$$
\end{proposition}

\begin{proof}
Existence. Définissons une application $g$ comme suit. Pour $x/s\in S^{-1}A$, posons
$g(x/s) = f(x)f(s)^{-1}\in B$. La définition permet de vérifier aisément
qu'il s'agit d'un morphisme d'anneaux bien défini. Il est également clair que
ce morphisme rend le diagramme commutatif.

\medskip\noindent
Unicité. Montrons maintenant que, si $g' : S^{-1}A \to B$
satisfait $g'(x/1) = f(x)$, alors $g = g'$. La commutativité du diagramme donne donc $f(s) = g'(s/1)$ pour
$s \in S$. Mais alors $g'(1/s)f(s) = 1$
dans $B$, ce qui implique $g'(1/s) = f(s)^{-1}$, puis
$g'(x/s) = g'(x/1)g'(1/s) = f(x)f(s)^{-1} = g(x/s)$.
\end{proof}

\begin{lemma}
\label{algebra-lemma-localization-zero}
La localisation $S^{-1}A$ est l'anneau nul si et seulement si $0\in S$.
\end{lemma}

\begin{proof}
Si $0\in S$, toute paire $(a, s)\sim (0, 1)$ par définition.
Si $0\not \in S$, alors il est clair que $1/1 \neq 0/1$ dans $S^{-1}A$.
\end{proof}

\begin{lemma}
\label{algebra-lemma-localization-and-modules}
Soit $R$ un anneau. Soit $S \subset R$ une partie multiplicative.
La catégorie des $S^{-1}R$-modules est équivalente à la catégorie
des $R$-modules $N$ ayant la propriété que tout $s \in S$ agit sur
$N$ par un automorphisme.
\end{lemma}

\begin{proof}
Le foncteur qui définit l'équivalence associe à un $S^{-1}R$-module
$M$ le même module, considéré maintenant comme un $R$-module au moyen du morphisme de localisation
$R \to S^{-1}R$. Réciproquement, si $N$ est un $R$-module tel que tout
$s \in S$ agisse par un automorphisme $s_N$, nous pouvons considérer $N$ comme un
$S^{-1}R$-module en faisant agir $x/s$ par $x_N  \circ s_N^{-1}$.
Nous omettons de vérifier que ces deux foncteurs sont quasi-inverses l'un de
l'autre.
\end{proof}

\noindent
La notion de localisation d'un anneau se généralise à la
localisation d'un module. Soient $A$ un anneau, $S$ une partie multiplicative
de $A$ et $M$ un $A$-module. Nous définissons une relation sur
$M \times S$ comme suit :
$$
(m, s) \sim (n, t)
\Leftrightarrow
\exists u\in S \text{ tel que } (mt-ns)u = 0
$$
Il s'agit clairement d'une relation d'équivalence. Notons $m/s$ (ou
$\frac{m}{s}$) la classe d'équivalence de $(m, s)$ et $S^{-1}M$
l'ensemble de toutes les classes d'équivalence. Définissons l'addition et la multiplication
scalaire comme suit :
$$
m/s + n/t = (mt + ns)/st,\quad
m/s\cdot n/t = mn/st
$$
Il est clair que cela fait de $S^{-1}M$ un $S^{-1}A$-module.

\begin{definition}
\label{algebra-definition-localization-module}
Le $S^{-1}A$-module $S^{-1}M$ est appelé la {\it localisation} de $M$
relativement à $S$.
\end{definition}

\noindent
Remarquons qu'il existe un homomorphisme de $A$-modules $M \to S^{-1}M$, $m \mapsto m/1$,
parfois
appelé le {\it morphisme de localisation}. Il satisfait la propriété
universelle suivante.

\begin{lemma}
\label{algebra-lemma-universal-property-localization-module}
Soit $R$ un anneau. Soit $S \subset R$ une partie multiplicative. Soient $M$, $N$
des $R$-modules. Supposons que tous les éléments de $S$ agissent sur $N$ par des automorphismes.
Alors l'application canonique
$$
\Hom_R(S^{-1}M, N) \longrightarrow \Hom_R(M, N)
$$
induite par le morphisme de localisation est un isomorphisme.
\end{lemma}

\begin{proof}
Il est clair que l'application est bien définie et $R$-linéaire.
Injectivité : soit $\alpha \in \Hom_R(S^{-1}M, N)$ et prenons un
élément arbitraire $m/s \in S^{-1}M$. Puisque $s \cdot \alpha(m/s) = \alpha(m/1)$,
nous avons $ \alpha(m/s) =s^{-1}(\alpha (m/1))$ ; ainsi, $\alpha$ est entièrement
déterminé par ses valeurs sur l'image de $M$ dans $S^{-1}M$.
Surjectivité : soit $\beta : M \rightarrow N$ une application R-linéaire donnée.
Nous devons montrer qu'elle peut être ``prolongée'' à $S^{-1}M$. Définissons une application d'ensembles
$$
M \times S \rightarrow N,\quad
(m,s) \mapsto s^{-1}\beta(m)
$$
Cette application respecte clairement la relation d'équivalence ci-dessus ; elle
passe donc au quotient en une application bien définie $\alpha : S^{-1}M \rightarrow N$.
Il reste à montrer que cette application est $R$-linéaire ; prenons donc
$r, r' \in R$, ainsi que $s, s' \in S$ et
$m, m' \in M$. Alors
\begin{align*}
\alpha(r \cdot m/s + r' \cdot m' /s')
& =
\alpha((r \cdot s' \cdot m + r' \cdot s \cdot m') /(ss')) \\
& =
(ss')^{-1}\beta(r \cdot s' \cdot m + r' \cdot s \cdot m') \\
& =
(ss')^{-1} (r \cdot s' \beta (m) + r' \cdot s \beta (m')) \\
& =
r \alpha (m/s) + r' \alpha (m' /s')
\end{align*}
ce qui conclut.
\end{proof}

\begin{example}
\label{algebra-example-localize-at-prime}
Soit $A$ un anneau et soit $M$ un $A$-module.
Voici quelques exemples importants de localisations.
\begin{enumerate}
\item Étant donné un idéal premier $\mathfrak p$ de $A$, considérons
$S = A\setminus\mathfrak p$. On vérifie
immédiatement que $S$ est une partie multiplicative. Dans ce cas,
nous notons $A_\mathfrak p$ et $M_\mathfrak p$ les localisations de
$A$ et de $M$ relativement à $S$, respectivement. Elles sont
appelées la {\it localisation de $A$, resp.\ de $M$, en $\mathfrak p$}.
\item Soit $f\in A$. Considérons $S = \{1, f, f^2, \ldots\}$.
C'est clairement une partie multiplicative de $A$.
Dans ce cas, nous notons $A_f$
(resp. $M_f$) la localisation $S^{-1}A$ (resp. $S^{-1}M$).
Elle est appelée la {\it localisation de $A$, resp.\ de $M$, relativement
à $f$}.
Remarquons que $A_f = 0$ si et seulement si $f$ est nilpotent dans $A$.
\item Soit $S = \{f \in A \mid f \text{ n'est pas un diviseur de zéro dans }A\}$.
C'est une partie multiplicative de $A$. Dans ce cas, l'anneau
$Q(A) = S^{-1}A$ est appelé soit l'{\it anneau total des quotients},
soit l'{\it anneau total des fractions}
de $A$.
\item Si $A$ est un anneau intègre, alors l'anneau total des quotients $Q(A)$ est
le corps des fractions de $A$. Voir
Corps, Exemple \ref{fields-example-quotient-field}.
\end{enumerate}
\end{example}

\begin{lemma}
\label{algebra-lemma-localization-colimit}
Soit $R$ un anneau.
Soit $S \subset R$ une partie multiplicative.
Soit $M$ un $R$-module.
Alors
$$
S^{-1}M = \colim_{f \in S} M_f
$$
où le préordre sur $S$ est donné par
$f \geq f' \Leftrightarrow f = f'f''$ pour un certain $f'' \in R$,
auquel cas l'application $M_{f'} \to M_f$ est donnée
par $m/(f')^e \mapsto m(f'')^e/f^e$.
\end{lemma}

\begin{proof}
Démonstration omise. Indication : utiliser la propriété universelle du
Lemme \ref{algebra-lemma-universal-property-localization-module}.
\end{proof}

\noindent
Dans le paragraphe suivant,
$A$ désigne un anneau,
et $M, N$ désignent des modules sur $A$.

\medskip\noindent
Si $S$ et $S'$ sont des parties multiplicatives de $A$, il est alors
clair que
$$
SS' = \{ss' : s\in S, \ s'\in S'\}
$$
est également une partie multiplicative de $A$. On a alors le résultat suivant.

\begin{proposition}
\label{algebra-proposition-localize-twice}
Soit $\overline{S}$ l'image de $S$ dans $S'^{-1}A$ ; alors
$(SS')^{-1}A$ est isomorphe à $\overline{S}^{-1}(S'^{-1}A)$.
\end{proposition}

\begin{proof}
L'application qui envoie $x\in A$ sur $x/1\in (SS')^{-1}A$ induit, par propriété
universelle, une application qui envoie $x/s\in S'^{-1}A$ sur $x/s \in (SS')^{-1}A$.
Les images des éléments de $\overline{S}$ sont inversibles
dans $(SS')^{-1}A$. La propriété universelle fournit une application
$f : \overline{S}^{-1}(S'^{-1}A) \to (SS')^{-1}A$ qui envoie
$(x/s')/(s/1)$ sur $x/ss'$.

\medskip\noindent
D'autre part, l'application de $A$ vers $\overline{S}^{-1}(S'^{-1}A)$
qui envoie $x\in A$ sur $(x/1)/(1/1)$ induit elle aussi, par la propriété universelle,
une application $g : (SS')^{-1}A \to \overline{S}^{-1}(S'^{-1}A)$ qui envoie $x/ss'$
sur $(x/s')/(s/1)$. On vérifie
immédiatement que $f$ et $g$ sont inverses l'un de l'autre ;
ce sont donc tous deux des isomorphismes.
\end{proof}

\noindent
Pour le module $M$, nous avons

\begin{proposition}
\label{algebra-proposition-localize-twice-module}
Considérons $S'^{-1}M$ comme un $A$-module ; alors $S^{-1}(S'^{-1}M)$ est
isomorphe à $(SS')^{-1}M$.
\end{proposition}

\begin{proof}
Remarquons qu'étant donné un $A$-module M, nous n'avons démontré aucune
propriété universelle pour $S^{-1}M$. Nous ne pouvons donc pas raisonner
comme dans la démonstration précédente ; nous devons construire explicitement l'isomorphisme.

\medskip\noindent
Définissons les applications comme suit :
\begin{align*}
& f : S^{-1}(S'^{-1}M) \longrightarrow (SS')^{-1}M, \quad \frac{x/s'}{s}\mapsto
x/ss'\\
& g : (SS')^{-1}M \longrightarrow S^{-1}(S'^{-1}M), \quad x/t\mapsto
\frac{x/s'}{s}\ \text{pour certains }s\in S, s'\in S', \text{ et }
t = ss'
\end{align*}
Nous devons vérifier que ces homomorphismes sont bien définis, c'est-à-dire
indépendants du choix de la fraction. Cela se vérifie aisément et il est également
immédiat qu'ils sont inverses l'un de l'autre.
\end{proof}

\noindent
Si $u : M \to N$ est un homomorphisme de $A$-modules, la localisation
induit bien un homomorphisme de $S^{-1}A$-modules $S^{-1}u : S^{-1}M \to
S^{-1}N$ qui envoie $x/s$ sur $u(x)/s$. On vérifie immédiatement que
cette construction est fonctorielle, si bien que $S^{-1}$
est effectivement un foncteur de la catégorie des $A$-modules vers la
catégorie des $S^{-1}A$-modules. De plus, ce foncteur est exact,
comme le montre la proposition suivante.

\begin{proposition}
\label{algebra-proposition-localization-exact}
\begin{slogan}
La localisation est exacte.
\end{slogan}
Soit $L\xrightarrow{u} M\xrightarrow{v} N$ une suite exacte
de $R$-modules. Alors
$S^{-1}L \to S^{-1}M \to S^{-1}N$ est également exacte.
\end{proposition}

\begin{proof}
Tout d'abord, il est clair que $S^{-1}L \to S^{-1}M \to S^{-1}N$ est un complexe,
puisque la localisation est un foncteur. Supposons ensuite que $x/s$ s'envoie sur zéro
dans $S^{-1}N$ pour un certain $x/s \in S^{-1}M$. Il existe alors, par définition, un
$t\in S$ tel que $v(xt) = v(x)t = 0$ dans $N$, ce qui signifie que
$xt \in \Ker(v)$. Par l'exactitude de $L \to M \to N$, nous avons
$xt = u(y)$ pour un certain $y$ dans $L$. Ainsi, $x/s$ est l'image de $y/st$.
Cela démontre l'exactitude.
\end{proof}

\begin{lemma}
\label{algebra-lemma-localize-quotient-modules}
La localisation respecte les quotients : si $N$ est un sous-module de
$M$, alors $S^{-1}(M/N)\simeq (S^{-1}M)/(S^{-1}N)$.
\end{lemma}

\begin{proof}
La suite exacte
$$
0 \longrightarrow N \longrightarrow M \longrightarrow M/N \longrightarrow 0
$$
donne
$$
0 \longrightarrow S^{-1}N \longrightarrow S^{-1}M
\longrightarrow S^{-1}(M/N) \longrightarrow 0
$$
L'énoncé en résulte.
\end{proof}

\noindent
Si, dans le lemme précédent, nous prenons $N = I$ et $M = A$ pour un idéal $I$ de
$A$, nous voyons que $S^{-1}A/S^{-1}I \simeq S^{-1}(A/I)$ comme $A$-modules. La
proposition suivante montre qu'ils sont isomorphes comme anneaux.

\begin{proposition}
\label{algebra-proposition-localize-quotient}
Soient $I$ un idéal de $A$ et $S$ une partie multiplicative de $A$. Alors
$S^{-1}I$ est un idéal de $S^{-1}A$ et $\overline{S}^{-1}(A/I)$ est
isomorphe à $S^{-1}A/S^{-1}I$, où $\overline{S}$ est
l'image de $S$ dans $A/I$.
\end{proposition}

\begin{proof}
Le fait que $S^{-1}I$ soit un idéal est clair puisque $I$ est lui-même un
idéal. Définissons
$$
f : S^{-1}A\longrightarrow \overline{S}^{-1}(A/I), \quad x/s\mapsto
\overline{x}/\overline{s}
$$
où $\overline{x}$ et $\overline{s}$ sont les images de $x$ et de
$s$ dans $A/I$. Nous conserverons des notations semblables dans cette démonstration.
Cette application est bien définie par la propriété universelle de
$S^{-1}A$, et $S^{-1}I$ est contenu dans son noyau ;
elle induit donc une application
$$
\overline{f} : S^{-1}A/S^{-1}I \longrightarrow \overline{S}^{-1}(A/I), \quad
\overline{x/s}\mapsto \overline{x}/\overline{s}
$$

\medskip\noindent
D'autre part, l'application $A \to S^{-1}A/S^{-1}I$ qui envoie $x$ sur
$\overline{x/1}$ induit une application $A/I \to S^{-1}A/S^{-1}I$ qui envoie
$\overline{x}$ sur $\overline{x/1}$. L'image de $\overline{S}$ est
inversible dans $S^{-1}A/S^{-1}I$ ; elle induit donc une application
$$
g : \overline{S}^{-1}(A/I) \longrightarrow S^{-1}A/S^{-1}I, \quad
\frac{\overline{x}}{\overline{s}}\mapsto \overline{x/s}
$$
par la propriété universelle. Il est alors clair que $\overline{f}$ et $g$
sont inverses l'un de l'autre ; ce sont donc tous deux des isomorphismes.
\end{proof}

\noindent
Étudions maintenant le comportement des sous-modules par localisation.

\begin{lemma}
\label{algebra-lemma-submodule-localization}
Tout sous-module $N'$ de $S^{-1}M$ est de la forme $S^{-1}N$ pour un certain
$N\subset M$. On peut en effet prendre pour $N$ l'image réciproque de
$N'$ dans $M$.
\end{lemma}

\begin{proof}
Soit $N$ l'image réciproque de $N'$ dans $M$. On voit alors que
$S^{-1}N\supset N'$. Pour montrer qu'ils sont égaux, prenons $x/s$ dans
$S^{-1}N$, où $s\in S$ et $x\in N$. Il s'ensuit que $x/1\in
N'$. Puisque $N'$ est un $S^{-1}R$-sous-module, nous avons
$x/s = x/1\cdot 1/s\in N'$. Cela achève la démonstration.
\end{proof}

\noindent
En prenant $M = A$ et $N = I$, où I est un idéal de $A$, nous obtenons le
corollaire suivant, que l'on peut considérer comme une réciproque de la première partie de la
Proposition \ref{algebra-proposition-localize-quotient}.

\begin{lemma}
\label{algebra-lemma-ideal-in-localization}
\begin{slogan}
Les idéaux de la localisation d'un anneau sont les localisés d'idéaux.
\end{slogan}
Tout idéal $I'$ de $S^{-1}A$ est de la forme $S^{-1}I$, où l'on peut
prendre pour $I$ l'image réciproque de $I'$ dans $A$.
\end{lemma}

\begin{proof}
Cela résulte immédiatement du Lemme \ref{algebra-lemma-submodule-localization}.
\end{proof}









\section{Hom interne}
\label{algebra-section-hom}

\noindent
Si $R$ est un anneau et si $M$, $N$ sont des $R$-modules, alors
$$
\Hom_R(M, N) = \{ \varphi : M \to N\}
$$
est l'ensemble des applications $R$-linéaires de $M$ vers $N$. Cet ensemble est muni
d'une structure de groupe abélien en posant, comme d'habitude,
$(\varphi + \psi)(m) = \varphi(m) + \psi(m)$.
En fait, $\Hom_R(M, N)$ est aussi un $R$-module par la règle
$(x \varphi)(m) = x \varphi(m) = \varphi(xm)$.

\medskip\noindent
Étant donnés des homomorphismes $a : M \to M'$ et $b : N \to N'$ de $R$-modules, nous pouvons
précomposer et postcomposer les homomorphismes par $a$ et $b$. On obtient
le diagramme commutatif suivant :
$$
\xymatrix{
\Hom_R(M', N) \ar[d]_{- \circ a} \ar[r]_{b \circ -} &
\Hom_R(M', N') \ar[d]^{- \circ a} \\
\Hom_R(M, N) \ar[r]^{b \circ -} &
\Hom_R(M, N')
}
$$
En fait, les applications de ce diagramme sont des homomorphismes de $R$-modules.
Ainsi, $\Hom_R$ définit un foncteur additif
$$
\text{Mod}_R^{opp} \times \text{Mod}_R \longrightarrow \text{Mod}_R, \quad
(M, N) \longmapsto \Hom_R(M, N)
$$

\begin{lemma}
\label{algebra-lemma-hom-exact}
Exactitude et $\Hom_R$. Soit $R$ un anneau. Soient $M_1$, $M_2$, $M_3$
des $R$-modules. Soient $M_1 \to M_2$ et $M_2 \to M_3$ des homomorphismes de $R$-modules.
\begin{enumerate}
\item $M_1 \to M_2 \to M_3 \to 0$ est exacte si et seulement si
$0 \to \Hom_R(M_3, N) \to \Hom_R(M_2, N) \to \Hom_R(M_1, N)$
est exacte pour tout $R$-module $N$.
\item $0 \to M_1 \to M_2 \to M_3$ est exacte si et seulement si
$0 \to \Hom_R(N, M_1) \to \Hom_R(N, M_2) \to \Hom_R(N, M_3)$
est exacte pour tout $R$-module $N$.
\end{enumerate}
\end{lemma}

\begin{proof}
Démonstration omise.
\end{proof}

\begin{lemma}
\label{algebra-lemma-hom-from-finitely-presented}
Soit $R$ un anneau. Soit $M$ un $R$-module de présentation finie.
Soit $N$ un $R$-module.
\begin{enumerate}
\item Pour $f \in R$, nous avons
$\Hom_R(M, N)_f = \Hom_{R_f}(M_f, N_f) = \Hom_R(M_f, N_f)$,
\item pour une partie multiplicative $S$ de $R$, nous avons
$$
S^{-1}\Hom_R(M, N) = \Hom_{S^{-1}R}(S^{-1}M, S^{-1}N) =
\Hom_R(S^{-1}M, S^{-1}N).
$$
\end{enumerate}
\end{lemma}

\begin{proof}
L'assertion (1) est un cas particulier de l'assertion (2).
La seconde égalité dans (2) résulte du
Lemme \ref{algebra-lemma-localization-and-modules}.
Choisissons une présentation
$$
\bigoplus\nolimits_{j = 1, \ldots, m} R
\longrightarrow
\bigoplus\nolimits_{i = 1, \ldots, n} R
\to M \to 0.
$$
D'après le
Lemme \ref{algebra-lemma-hom-exact},
nous obtenons une suite exacte
$$
0 \to
\Hom_R(M, N) \to
\bigoplus\nolimits_{i = 1, \ldots, n} N
\longrightarrow
\bigoplus\nolimits_{j = 1, \ldots, m} N.
$$
En inversant $S$ et en utilisant la Proposition \ref{algebra-proposition-localization-exact},
nous obtenons une suite exacte
$$
0 \to
S^{-1}\Hom_R(M, N) \to
\bigoplus\nolimits_{i = 1, \ldots, n} S^{-1}N
\longrightarrow
\bigoplus\nolimits_{j = 1, \ldots, m} S^{-1}N
$$
et le résultat en découle puisque $S^{-1}M$ s'insère dans
une suite exacte
$$
\bigoplus\nolimits_{j = 1, \ldots, m} S^{-1}R
\longrightarrow
\bigoplus\nolimits_{i = 1, \ldots, n} S^{-1}R \to S^{-1}M \to 0
$$
qui induit (d'après le Lemme \ref{algebra-lemma-hom-exact})
la suite exacte
$$
0 \to
\Hom_{S^{-1}R}(S^{-1}M, S^{-1}N) \to
\bigoplus\nolimits_{i = 1, \ldots, n} S^{-1}N
\longrightarrow
\bigoplus\nolimits_{j = 1, \ldots, m} S^{-1}N
$$
qui est la même que ci-dessus.
\end{proof}




\section{Caractérisation des modules de type fini et de présentation finie}
\label{algebra-section-colim-and-hom}

\noindent
Étant donné un module $N$ sur un anneau $R$, on peut caractériser si
$N$ est ou non un module de type fini ou un module de présentation finie
à l'aide du foncteur $\Hom_R(N, -)$.

\begin{lemma}
\label{algebra-lemma-characterize-finite-module-hom}
Soit $R$ un anneau. Soit $N$ un $R$-module. Les assertions suivantes sont équivalentes :
\begin{enumerate}
\item $N$ est un $R$-module de type fini,
\item pour toute colimite filtrante $M = \colim M_i$ de $R$-modules, l'application
$\colim \Hom_R(N, M_i) \to \Hom_R(N, M)$ est injective.
\end{enumerate}
\end{lemma}

\begin{proof}
Supposons (1) et choisissons des générateurs $x_1, \ldots, x_m$ de $N$.
Si $N \to M_i$ est un homomorphisme de modules et si la composée
$N \to M_i \to M$ est nulle, alors, puisque $M = \colim_{i' \geq i} M_{i'}$,
pour tout $j \in \{1, \ldots, m\}$, nous pouvons trouver un $i' \geq i$ tel que
$x_j$ s'envoie sur zéro dans $M_{i'}$. Comme les $x_j$ sont en nombre fini,
nous pouvons trouver un seul $i'$ qui convienne à tous.
La composée $N \to M_i \to M_{i'}$ est alors nulle, et nous en concluons que
l'application est injective, c'est-à-dire que l'assertion (2) vaut.

\medskip\noindent
Supposons (2). Pour un sous-ensemble fini $E \subset N$, notons $N_E \subset N$
le $R$-sous-module engendré par les éléments de $E$. Alors
$0 = \colim N/N_E$ est une colimite filtrante. Nous voyons donc que
$\text{id} : N \to N$ se factorise par $N_E$ pour un certain $E$, c'est-à-dire que $N$
est de type fini.
\end{proof}

\noindent
Pour référence, nous définissons ce que signifie l'existence d'une relation
entre des éléments d'un module.

\begin{definition}
\label{algebra-definition-relation}
Soit $R$ un anneau. Soit $M$ un $R$-module.
Soient $n \geq 0$ et $x_i \in M$ pour $i = 1, \ldots, n$.
Une {\it relation} entre $x_1, \ldots, x_n$ dans $M$ est une
suite d'éléments $f_1, \ldots, f_n \in R$ telle que
$\sum_{i = 1, \ldots, n} f_i x_i = 0$.
\end{definition}

\begin{lemma}
\label{algebra-lemma-module-colimit-fp}
Soit $R$ un anneau et soit $M$ un $R$-module.
Alors $M$ est la colimite d'un système filtrant
$(M_i, \mu_{ij})$ de $R$-modules
dont tous les $M_i$ sont des $R$-modules de présentation finie.
\end{lemma}

\begin{proof}
Considérons un sous-ensemble fini quelconque $S \subset M$ et une famille finie
de relations $E$ entre les éléments
de $S$. Ainsi, chaque $s \in S$ correspond à un $x_s \in M$ et
chaque $e \in E$ est constitué d'un vecteur
d'éléments $f_{e, s} \in R$ tel que $\sum f_{e, s} x_s = 0$.
Soit $M_{S, E}$ le conoyau de l'application
$$
R^{\# E} \longrightarrow R^{\# S}, \quad
(g_e)_{e\in E} \longmapsto (\sum g_e f_{e, s})_{s\in S}.
$$
Il existe des applications canoniques $M_{S, E} \to M$.
Si $S \subset S'$ et si les éléments de
$E$ correspondent, par cette application, à des relations
de $E'$, il existe alors une application évidente
$M_{S, E} \to M_{S', E'}$ qui commute aux
applications vers $M$. Soit $I$ l'ensemble des paires
$(S, E)$, ordonné par l'inclusion décrite ci-dessus.
Il est clair que la colimite de ce système filtrant est $M$.
\end{proof}

\begin{lemma}
\label{algebra-lemma-characterize-finitely-presented-module-hom}
Soit $R$ un anneau. Soit $N$ un $R$-module. Les assertions suivantes sont équivalentes :
\begin{enumerate}
\item $N$ est un $R$-module de présentation finie,
\item pour toute colimite filtrante $M = \colim M_i$ de $R$-modules, l'application
$\colim \Hom_R(N, M_i) \to \Hom_R(N, M)$ est bijective.
\end{enumerate}
\end{lemma}

\begin{proof}
Supposons (1) et choisissons une suite exacte $F_{-1} \to F_0 \to N \to 0$
où les $F_i$ sont libres de type fini. Nous avons alors une suite exacte
$$
0 \to \Hom_R(N, M) \to \Hom_R(F_0, M) \to \Hom_R(F_{-1}, M)
$$
fonctorielle en le $R$-module $M$. Les foncteurs $\Hom_R(F_i, M)$ commutent
aux colimites filtrantes puisque $\Hom_R(R^{\oplus n}, M) = M^{\oplus n}$.
Comme les colimites filtrantes sont exactes
(Lemme \ref{algebra-lemma-directed-colimit-exact}),
nous voyons que (2) vaut.

\medskip\noindent
Supposons (2). D'après le Lemme \ref{algebra-lemma-module-colimit-fp},
nous pouvons écrire $N = \colim N_i$ comme une colimite
filtrante telle que $N_i$ soit de présentation finie pour tout $i$.
Ainsi, $\text{id}_N$ se factorise par $N_i$ pour un certain $i$.
Cela signifie que $N$ est facteur direct d'un module de
présentation finie sur $R$ (à savoir $N_i$), et est donc de présentation finie.
\end{proof}












\section{Produits tensoriels}
\label{algebra-section-tensor-product}

\begin{definition}
\label{algebra-definition-bilinear}
Soient $R$ un anneau et $M, N, P$ trois $R$-modules.
Une application $f : M \times N \to P$ (où $M \times N$
est considéré uniquement comme le produit cartésien de deux $R$-modules) est dite
{\it $R$-bilinéaire} si, pour tout $x \in M$,
l'application $y\mapsto f(x, y)$ de $N$ dans $P$ est $R$-linéaire et si, pour tout
$y\in N$, l'application $x\mapsto f(x, y)$ est également $R$-linéaire.
\end{definition}

\begin{lemma}
\label{algebra-lemma-tensor-product}
Soient $M, N$ des $R$-modules. Il existe alors une paire $(T, g)$,
où $T$ est un $R$-module et
$g : M \times N \to T$ une application $R$-bilinéaire,
qui possède la propriété universelle suivante :
pour tout $R$-module $P$ et toute application $R$-bilinéaire
$f : M \times N \to P$, il
existe une unique application $R$-linéaire
$\tilde{f} : T \to P$ telle que $f = \tilde{f} \circ g$.
Autrement dit, le diagramme suivant est commutatif :
$$
\xymatrix{
M \times N \ar[rr]^f \ar[dr]_g & & P\\
& T \ar[ur]_{\tilde f}
}
$$
De plus, si $(T, g)$ et $(T', g')$
sont deux paires possédant cette propriété, il
existe un unique isomorphisme
$j : T \to T'$ tel que $j\circ g = g'$.
\end{lemma}

\noindent
Le $R$-module $T$ qui satisfait à la propriété universelle ci-dessus est appelé
le {\it produit tensoriel} des $R$-modules $M$ et $N$, et est noté
$M \otimes_R N$.

\begin{proof}
Démontrons d'abord l'existence d'un tel $R$-module $T$.
Soient $M, N$ des $R$-modules.
Soit $T$ le module quotient
$P/Q$, où $P$ est le $R$-module libre $R^{(M \times N)}$ et où $Q$ est le
$R$-module engendré par tous les éléments
des types suivants : ($x\in M, y\in N$)
\begin{align*}
(x + x', y) - (x, y) - (x', y), \\
(x, y + y') - (x, y) - (x, y'), \\
(ax, y) - a(x, y), \\
(x, ay) - a(x, y)
\end{align*}
Notons $\pi : M \times N \to T$ l'application canonique.
Cette application est $R$-bilinéaire, comme
le montrent les relations ci-dessus
lorsque l'on vérifie les conditions de bilinéarité. Notons l'image
$\pi(x, y) = x \otimes
y$ ; ces éléments engendrent alors
$T$. Soit maintenant $f : M \times N \to P$ une application $R$-bilinéaire ;
nous pouvons définir
$f' : T \to P$ en prolongeant l'application
$f'(x \otimes y) = f(x, y)$. Il est clair que $f = f'\circ \pi$. De plus, $f'$ est
uniquement déterminée par ses valeurs sur le
système générateur $\{x \otimes y : x\in M, y\in N\}$.
Supposons qu'il existe une autre paire $(T', g')$ satisfaisant aux mêmes propriétés.
Il existe alors un unique $j : T \to T'$, ainsi
qu'un $j' : T' \to T$, tels que $g' = j\circ g$, $g = j'\circ g'$.
Mais les deux applications $(j\circ j') \circ g$ et $g$
satisfont alors aux propriétés universelles ; par unicité, elles sont donc égales,
et par conséquent $j'\circ j$ est l'identité sur $T$.
De même, $(j'\circ j) \circ g' = g'$ et $j\circ j'$ est l'identité sur $T'$.
Ainsi, $j$ est un isomorphisme.
\end{proof}

\begin{lemma}
\label{algebra-lemma-flip-tensor-product}
Soient $M, N, P$ des $R$-modules ; les applications bilinéaires
\begin{align*}
(x, y) & \mapsto y \otimes x\\
(x + y, z) & \mapsto x \otimes z + y \otimes z\\
(r, x) & \mapsto rx
\end{align*}
induisent des isomorphismes uniques
\begin{align*}
M \otimes_R N & \to N \otimes_R M, \\
(M\oplus N)\otimes_R P & \to (M \otimes_R P)\oplus(N \otimes_R P),  \\
R \otimes_R M & \to M
\end{align*}
\end{lemma}

\begin{proof}
Démonstration omise.
\end{proof}

\noindent
Nous pouvons généraliser le produit tensoriel de deux $R$-modules à un nombre fini de
$R$-modules et établir une
correspondance entre le produit multitensoriel et les applications multilinéaires.
Une construction presque identique
permet de démontrer le résultat suivant :

\begin{lemma}
\label{algebra-lemma-multilinear}
Soient $M_1, \ldots, M_r$ des $R$-modules. Il existe alors une paire $(T, g)$
constituée d'un $R$-module T et d'une application $R$-multilinéaire
$g : M_1\times \ldots \times M_r \to T$ possédant la propriété
universelle suivante : pour toute application $R$-multilinéaire
$f : M_1\times \ldots \times M_r \to P$, il existe un unique homomorphisme de $R$-modules
$f' : T \to P$ tel que $f'\circ g = f$.
Un tel module $T$ est unique à isomorphisme unique près. Nous le notons
$M_1\otimes_R \ldots \otimes_R M_r$ et nous notons l'application
multilinéaire universelle $(m_1, \ldots, m_r) \mapsto m_1 \otimes \ldots \otimes m_r$.
\end{lemma}

\begin{proof}
Démonstration omise.
\end{proof}

\begin{lemma}
\label{algebra-lemma-transitive}
Les homomorphismes
$$
(M \otimes_R N)\otimes_R P \to
M \otimes_R N \otimes_R P \to
M \otimes_R (N \otimes_R P)
$$
tels que
$f((x \otimes y)\otimes z) = x \otimes y \otimes z$
et $g(x \otimes y \otimes z) = x \otimes (y \otimes z)$,
$x\in M, y\in N, z\in P$, sont bien définis et sont des isomorphismes.
\end{lemma}

\begin{proof}
Nous allons montrer que $f$ est bien défini et est un isomorphisme ; la démonstration
s'applique de manière analogue à $g$. Fixons un
$z\in P$ quelconque. L'application $(x, y)\mapsto x \otimes y \otimes z$,
$x\in M, y\in N$, est alors $R$-bilinéaire en $x$ et $y$,
et induit donc un homomorphisme $f_z : M \otimes N \to M \otimes N \otimes P$
qui vérifie
$f_z(x \otimes y) = x \otimes y \otimes z$.
Considérons ensuite $(M \otimes N)\times P \to M \otimes N \otimes P$ donnée par
$(w, z)\mapsto f_z(w)$. Cette application est
$R$-bilinéaire et induit donc
$f : (M \otimes_R N)\otimes_R P \to M \otimes_R N \otimes_R P$,
avec $f((x \otimes y)\otimes z) = x \otimes y \otimes z$.
Pour construire l'inverse, remarquons que l'application
$\pi : M \times N \times P \to (M \otimes N)\otimes P$ est
$R$-trilinéaire.
Elle induit donc une application $R$-linéaire
$h : M \otimes N \otimes P \to (M \otimes N)\otimes P$ qui
s'accorde avec la propriété universelle. Nous voyons ici que
$h(x \otimes y \otimes z) = (x \otimes y)\otimes z$.
D'après les expressions explicites de $f$ et de $h$, $f\circ h$ et $h\circ f$ sont
les applications identiques de $M \otimes N \otimes
P$ et de $(M \otimes N)\otimes P$, respectivement ; ainsi, $f$ est bien
l'isomorphisme recherché.
\end{proof}

\noindent
En raisonnant par récurrence, nous voyons que cela s'étend aux produits multitensoriels. En combinant
ce résultat avec le Lemme \ref{algebra-lemma-flip-tensor-product}, nous voyons que
l'opération de produit tensoriel sur la catégorie des $R$-modules est associative,
commutative et distributive.

\begin{definition}
\label{algebra-definition-bimodule}
Un groupe abélien $N$ est appelé un {\it $(A, B)$-bimodule} s'il est à la fois un
$A$-module et un $B$-module et si, pour tous $a \in A$ et $b \in B$, les
multiplications par $a$ et par $b$ commutent, de sorte que $b(an) = a(bn)$ pour tout $n \in N$.
Dans cette situation, nous écrivons habituellement l'action de $B$ à droite : ainsi,
pour $b \in B$ et $n \in N$, le résultat de la multiplication de $n$ par $b$
est noté $nb$. Avec cette convention, la compatibilité ci-dessus s'écrit
$(ax)b = a(xb)$ pour tous $a\in A, b\in B, x\in N$.
La notation abrégée $_AN_B$ désigne un $(A, B)$-bimodule $N$.
\end{definition}

\begin{lemma}
\label{algebra-lemma-tensor-with-bimodule}
Pour un $A$-module $M$, un $B$-module $P$ et un $(A, B)$-bimodule $N$, les modules
$(M \otimes_A N)\otimes_B P$ et $M \otimes_A(N \otimes_B P)$ peuvent tous deux être
munis d'une structure de $(A, B)$-bimodule,
et de plus
$$
(M \otimes_A N)\otimes_B P \cong M \otimes_A(N \otimes_B P).
$$
\end{lemma}

\begin{proof}
A priori, $M \otimes_A N$ est un $A$-module, mais nous pouvons le munir d'une
structure de $B$-module en posant
$$
(x \otimes y)b = x \otimes yb, \quad x\in M, y\in N, b\in B
$$
Ainsi, $M \otimes_A N$ devient un $(A, B)$-bimodule. Il en va de même pour
$N \otimes_B P$, et donc pour
$(M \otimes_A N)\otimes_B P$ et $M \otimes_A(N \otimes_B P)$. D'après le
Lemme \ref{algebra-lemma-transitive}, ces deux
modules sont isomorphes à la fois comme $A$-modules et comme $B$-modules par la même
application.
\end{proof}

\begin{lemma}
\label{algebra-lemma-hom-from-tensor-product}
Pour trois $R$-modules quelconques $M, N, P$,
$$
\Hom_R(M \otimes_R N, P) \cong \Hom_R(M, \Hom_R(N, P))
$$
\end{lemma}

\begin{proof}
Une application $R$-linéaire $\hat{f}\in \Hom_R(M \otimes_R N, P)$ correspond à une
application $R$-bilinéaire $f : M \times N \to P$. Pour
tout $x\in M$, l'application $y\mapsto f(x, y)$ est $R$-linéaire par la propriété
universelle. Ainsi, $f$ correspond à une
application $\phi_f : M \to \Hom_R(N, P)$. Cette application est $R$-linéaire puisque
$$
\phi_f(ax + y)(z) =
f(ax + y, z) = af(x, z)+f(y, z) =
(a\phi_f(x)+\phi_f(y))(z),
$$
pour tous $a \in R$, $x \in M$, $y \in M$ et
$z \in N$. Réciproquement, tout
$f \in \Hom_R(M, \Hom_R(N, P))$ définit une application $R$-bilinéaire
$M \times N \to P$, à savoir $(x, y)\mapsto f(x)(y)$.
Nous obtenons donc une bijection naturelle entre les
deux modules
$\Hom_R(M \otimes_R N, P)$ et $\Hom_R(M, \Hom_R(N, P))$.
\end{proof}

\begin{lemma}[Les produits tensoriels commutent aux colimites]
\label{algebra-lemma-tensor-products-commute-with-limits}
Soit $(M_i, \mu_{ij})$ un système sur l'ensemble préordonné $I$.
Soit $N$ un $R$-module. Alors
$$
\colim (M_i \otimes N) \cong (\colim M_i)\otimes N.
$$
De plus, l'isomorphisme est induit par les homomorphismes
$\mu_i \otimes 1: M_i \otimes N \to M \otimes N$,
où $M = \colim_i M_i$ et où $\mu_i : M_i \to M$ sont les applications canoniques.
\end{lemma}

\begin{proof}
Première démonstration. Le foncteur $M' \mapsto M' \otimes_R N$ est adjoint à gauche
du foncteur $N' \mapsto \Hom_R(N, N')$ d'après le
Lemme \ref{algebra-lemma-hom-from-tensor-product}. Ainsi, $M' \mapsto M' \otimes_R N$
commute à toutes les colimites ; voir
Catégories, Lemme \ref{categories-lemma-adjoint-exact}.

\medskip\noindent
Seconde démonstration directe. Soit $P = \colim (M_i \otimes N)$, muni des coprojections
$\lambda_i : M_i \otimes N \to P$. Soit $M = \colim M_i$, muni des coprojections
$\mu_i : M_i \to M$. Alors, pour tous $i\leq j$, le diagramme suivant est commutatif :
$$
\xymatrix{
M_i \otimes N \ar[r]_{\mu_i \otimes 1} \ar[d]_{\mu_{ij} \otimes 1} &
M \otimes N \ar[d]^{\text{id}} \\
M_j \otimes N \ar[r]^{\mu_j \otimes 1} &
M \otimes N
}
$$
D'après le Lemme \ref{algebra-lemma-homomorphism-limit}, ces applications induisent un unique homomorphisme
$\psi : P \to M \otimes N$ tel que
$\mu_i \otimes 1 = \psi \circ \lambda_i$.

\medskip\noindent
Pour construire l'application inverse, il existe, pour tout $i\in I$, l'application canonique
$R$-bilinéaire $g_i : M_i \times N \to
M_i \otimes N$. Celle-ci induit une unique application
$\widehat{\phi} : M \times N \to P$
telle que $\widehat{\phi} \circ (\mu_i \times 1) = \lambda_i \circ g_i$.
Elle est $R$-bilinéaire et induit donc une
application $R$-linéaire $\phi : M \otimes N \to P$.
Le diagramme commutatif ci-dessous :
$$
\xymatrix{
M_i \times N \ar[r]^{g_i} \ar[d]^{\mu_i \times \text{id}} &
M_i \otimes N\ar[r]_{\text{id}} \ar[d]_{\lambda_i} &
M_i \otimes N \ar[d]_{\mu_i \otimes \text{id}} \ar[rd]^{\lambda_i} \\
M \times N \ar[r]^{\widehat{\phi}} &
P \ar[r]^{\psi} & M \otimes N \ar[r]^{\phi} & P
}
$$
montre que $\psi\circ\widehat{\phi} = g$, où g est l'application $R$-bilinéaire canonique
$g : M \times N \to M \otimes N$. Ainsi,
$\psi\circ\phi$ est l'identité de $M \otimes N$. Le carré et le
triangle de droite montrent que $\phi\circ\psi$ est également
l'identité de $P$.
\end{proof}

\begin{lemma}
\label{algebra-lemma-tensor-product-exact}
Soit
\begin{align*}
M_1\xrightarrow{f} M_2\xrightarrow{g} M_3 \to 0
\end{align*}
une suite exacte de $R$-modules et d'homomorphismes, et soit $N$ un
$R$-module quelconque. Alors la suite
\begin{equation}
\label{algebra-equation-2ndex}
M_1\otimes N\xrightarrow{f \otimes 1} M_2\otimes N \xrightarrow{g \otimes 1}
M_3\otimes N \to 0
\end{equation}
est exacte. Autrement dit, le foncteur $- \otimes_R N$ est
{\it exact à droite}, au sens où tensoriser
chaque terme de la suite exacte à droite initiale préserve l'exactitude.
\end{lemma}

\begin{proof}
Pour tout $R$-module $P$,
appliquons le foncteur $\Hom(-, \Hom(N, P))$ à la première suite
exacte. Nous obtenons
$$
0 \to
\Hom(M_3, \Hom(N, P)) \to
\Hom(M_2, \Hom(N, P)) \to
\Hom(M_1, \Hom(N, P))
$$
qui est exacte d'après le Lemme \ref{algebra-lemma-hom-exact} (1).
D'après le Lemme \ref{algebra-lemma-hom-from-tensor-product}, cette suite devient
$$
0 \to \Hom(M_3 \otimes N, P) \to
\Hom(M_2 \otimes N, P) \to \Hom(M_1 \otimes N, P)
$$
et est donc elle aussi exacte. En appliquant de nouveau le
Lemme \ref{algebra-lemma-hom-exact} (1), nous obtenons la suite exacte recherchée.
\end{proof}

\begin{remark}
\label{algebra-remark-tensor-product-not-exact}
Cependant, le produit tensoriel NE préserve PAS les suites exactes en général.
Autrement dit, si $M_1 \to M_2 \to M_3$ est
exacte, il n'est pas nécessaire que
$M_1 \otimes N \to M_2 \otimes N \to M_3 \otimes N$
soit exacte pour un $R$-module $N$ arbitraire.
\end{remark}

\begin{example}
\label{algebra-example-tensor-product-not-exact}
Considérons l'application injective $2 : \mathbf{Z}\to \mathbf{Z}$
comme une application de $\mathbf{Z}$-modules.
Soit $N = \mathbf{Z}/2$. L'application induite
$\mathbf{Z} \otimes \mathbf{Z}/2 \to \mathbf{Z} \otimes \mathbf{Z}/2$
n'est PAS injective. En effet, pour
$x \otimes y\in \mathbf{Z} \otimes \mathbf{Z}/2$,
$$
(2 \otimes 1)(x \otimes y) = 2x \otimes y = x \otimes 2y = x \otimes 0 = 0
$$
L'application induite est donc l'application nulle, tandis que $\mathbf{Z} \otimes N\neq 0$.
\end{example}

\begin{remark}
\label{algebra-remark-flat-module}
Pour un $R$-module $N$, si le
foncteur $-\otimes_R N$ est exact, c'est-à-dire si la tensorisation
par $N$ préserve toutes les suites
exactes, alors $N$ est appelé un $R$-module {\it plat}.
Nous étudierons cette notion plus loin dans la section \ref{algebra-section-flat}.
\end{remark}

\begin{lemma}
\label{algebra-lemma-tensor-finiteness}
Soit $R$ un anneau. Soient $M$ et $N$ des $R$-modules.
\begin{enumerate}
\item Si $N$ et $M$ sont de type fini, alors $M \otimes_R N$ l'est aussi.
\item Si $N$ et $M$ sont de présentation finie, alors $M \otimes_R N$ l'est aussi.
\end{enumerate}
\end{lemma}

\begin{proof}
Supposons $M$ de type fini. Choisissons alors une présentation
$0 \to K \to R^{\oplus n} \to M \to 0$. Elle donne une suite exacte
$K \otimes_R N \to N^{\oplus n} \to M \otimes_R N \to 0$ d'après le
Lemme \ref{algebra-lemma-tensor-product-exact}.
Nous en concluons que, si $N$ est également de type fini, alors $M \otimes_R N$
est un quotient d'un module de type fini, donc est de type fini ; voir le
Lemme \ref{algebra-lemma-extension}.
De même, si $N$ et $M$ sont tous deux de présentation finie,
nous voyons que $K$ est de type fini et que $M \otimes_R N$
est le quotient du module de présentation finie $N^{\oplus n}$ par
un module de type fini, à savoir $K \otimes_R N$, et est donc de présentation finie ; voir le
Lemme \ref{algebra-lemma-extension}.
\end{proof}

\begin{lemma}
\label{algebra-lemma-tensor-localization}
Soit $M$ un $R$-module. Alors les $S^{-1}R$-modules $S^{-1}M$
et $S^{-1}R \otimes_R M$ sont canoniquement isomorphes, et
l'isomorphisme canonique $f : S^{-1}R \otimes_R M \to S^{-1}M$
est donné par
$$
f((a/s) \otimes m) = am/s, \forall a \in R, m \in M, s \in S
$$
\end{lemma}

\begin{proof}
Il est évident que l'application
$f' : S^{-1}R \times M \to S^{-1}M$ donnée par $f'(a/s, m) = am/s$ est
bilinéaire ; par la
propriété universelle, elle induit donc un unique homomorphisme de $S^{-1}R$-modules
$f : S^{-1}R \otimes_R M \to S^{-1}M$ comme dans l'énoncé du lemme.
En fait, tout élément de $S^{-1}M$ est de la forme $m/s$, $m\in M, s\in S$, et
tout élément de
$S^{-1}R \otimes_R M$ est de la forme $1/s \otimes m$. Pour voir ce dernier fait,
écrivons un élément de
$S^{-1}R \otimes_R M$ sous la forme
$$
\sum_k \frac{a_k}{s_k} \otimes m_k =
\sum_k \frac{a_k t_k}{s} \otimes m_k =
\frac{1}{s} \otimes \sum_k {a_k t_k}m_k = \frac{1}{s} \otimes m
$$
où $m = \sum_k {a_k t_k}m_k$. Il est alors évident que $f$ est surjectif,
et si $f(\frac{1}{s} \otimes m) = m/s = 0$, il existe $t \in S$ tel que
$tm = 0$ dans $M$. Nous avons alors
$$
\frac{1}{s} \otimes m = \frac{1}{st} \otimes tm = \frac{1}{st} \otimes 0 = 0
$$
Par conséquent, $f$ est injectif.
\end{proof}

\begin{lemma}
\label{algebra-lemma-tensor-product-localization}
Soient $M, N$ des $R$-modules. Il existe alors un
isomorphisme canonique de $S^{-1}R$-modules
$f : S^{-1}M \otimes_{S^{-1}R}S^{-1}N \to S^{-1}(M \otimes_R N)$,
donné par
$$
f((m/s)\otimes(n/t)) = (m \otimes n)/st
$$
\end{lemma}

\begin{proof}
Nous pouvons utiliser à plusieurs reprises le Lemme \ref{algebra-lemma-tensor-with-bimodule}
et le Lemme \ref{algebra-lemma-tensor-localization} pour
voir que ces deux
$S^{-1}R$-modules sont isomorphes, en remarquant que $S^{-1}R$ est un
$(R, S^{-1}R)$-bimodule :
\begin{align*}
S^{-1}(M \otimes_R N) & \cong S^{-1}R \otimes_R (M \otimes_R N)\\
 & \cong S^{-1}M \otimes_R N\\
 & \cong (S^{-1}M \otimes_{S^{-1}R}S^{-1}R)\otimes_R N\\
 & \cong S^{-1}M \otimes_{S^{-1}R}(S^{-1}R \otimes_R N)\\
 & \cong S^{-1}M \otimes_{S^{-1}R}S^{-1}N
\end{align*}
On vérifie aisément que cet isomorphisme est bien celui énoncé dans le lemme.
\end{proof}

\section{Algèbre tensorielle}
\label{algebra-section-tensor-algebra}

\noindent
Soit $R$ un anneau. Soit $M$ un $R$-module.
Nous définissons l'{\it algèbre tensorielle de $M$ sur $R$}
comme la $R$-algèbre non commutative
$$
\text{T}(M) = \text{T}_R(M) =
\bigoplus\nolimits_{n \geq 0} \text{T}^n(M)
$$
où
$\text{T}^0(M) = R$,
$\text{T}^1(M) = M$,
$\text{T}^2(M) = M \otimes_R M$,
$\text{T}^3(M) = M \otimes_R M \otimes_R M$, et ainsi de suite.
La multiplication est définie par la règle suivante sur les tenseurs purs :
$$
(x_1 \otimes x_2 \otimes \ldots \otimes x_n)
\cdot
(y_1 \otimes y_2 \otimes \ldots \otimes y_m)
=
x_1 \otimes x_2 \otimes \ldots \otimes x_n \otimes
y_1 \otimes y_2 \otimes \ldots \otimes y_m
$$
et nous l'étendons par linéarité.

\medskip\noindent
Nous définissons l'{\it algèbre extérieure $\wedge(M)$ de $M$ sur $R$}
comme le quotient de $\text{T}(M)$ par l'idéal bilatère
engendré par les éléments $x \otimes x \in \text{T}^2(M)$.
L'image d'un tenseur pur $x_1 \otimes \ldots \otimes x_n$
dans $\wedge^n(M)$ est notée $x_1 \wedge \ldots \wedge x_n$.
Ces éléments engendrent $\wedge^n(M)$, ils dépendent $R$-linéairement
de chaque $x_i$ et ils sont nuls lorsque deux des $x_i$ sont égaux
(autrement dit, ils sont alternés en
$x_1, x_2, \ldots, x_n$). La multiplication sur $\wedge(M)$ est
commutative au sens gradué, c'est-à-dire que tous $x \in M$ et $y \in M$
vérifient $x \wedge y = - y \wedge x$.

\medskip\noindent
Considérons le cas où $M = Rx_1 \oplus \ldots \oplus Rx_n$
est un module libre de type fini. Alors $\wedge(M)$ est libre sur
$R$, de base formée des éléments
$$
x_{i_1} \wedge \ldots \wedge x_{i_r}
$$
où $0 \leq r \leq n$ et $1 \leq i_1 < i_2 < \ldots < i_r \leq n$.

\medskip\noindent
Nous définissons l'{\it algèbre symétrique $\text{Sym}(M)$ de $M$ sur $R$}
comme le quotient de $\text{T}(M)$ par l'idéal bilatère
engendré par les éléments $x \otimes y - y \otimes x \in \text{T}^2(M)$.
L'image d'un tenseur pur $x_1 \otimes \ldots \otimes x_n$
dans $\text{Sym}^n(M)$ est simplement notée $x_1 \ldots x_n$.
Ces éléments engendrent $\text{Sym}^n(M)$, ils dépendent $R$-linéairement
de chaque $x_i$, et $x_1 \ldots x_n = x_1' \ldots x_n'$ si la
suite d'éléments $x_1, \ldots, x_n$ est une permutation de la
suite $x_1', \ldots, x_n'$. Ainsi, $\text{Sym}(M)$
est commutative.

\medskip\noindent
Considérons le cas où $M = Rx_1 \oplus \ldots \oplus Rx_n$
est un module libre de type fini. Alors
$\text{Sym}(M) = R[x_1, \ldots, x_n]$ est une algèbre polynomiale.

\begin{lemma}
\label{algebra-lemma-free-tensor-algebra}
Soit $R$ un anneau. Soit $M$ un $R$-module.
Si $M$ est un $R$-module libre, chacune de ses puissances symétriques
et extérieures l'est aussi.
\end{lemma}

\begin{proof}
La démonstration est omise ; voir toutefois ci-dessus le cas libre de type fini.
\end{proof}

\begin{lemma}
\label{algebra-lemma-presentation-sym-exterior}
Soit $R$ un anneau.
Soit $M_2 \to M_1 \to M \to 0$ une suite exacte de $R$-modules.
Il existe des suites exactes
$$
M_2 \otimes_R \text{Sym}^{n - 1}(M_1)
\to
\text{Sym}^n(M_1)
\to
\text{Sym}^n(M)
\to
0
$$
et, de même,
$$
M_2 \otimes_R \wedge^{n - 1}(M_1)
\to
\wedge^n(M_1)
\to
\wedge^n(M)
\to
0
$$
\end{lemma}

\begin{proof}
La démonstration est omise.
\end{proof}

\begin{lemma}
\label{algebra-lemma-present-sym-wedge}
Soit $R$ un anneau.
Soit $M$ un $R$-module.
Soit $x_i$, $i \in I$, une famille donnée de générateurs de
$M$ comme $R$-module. Soit $n \geq 2$.
Il existe une suite exacte canonique
$$
\bigoplus_{1 \leq j_1 < j_2 \leq n}
\bigoplus_{i_1, i_2 \in I}
\text{T}^{n - 2}(M)
\oplus
\bigoplus_{1 \leq j_1 < j_2 \leq n}
\bigoplus_{i \in I}
\text{T}^{n - 2}(M)
\to
\text{T}^n(M)
\to
\wedge^n(M)
\to
0
$$
où le tenseur pur $m_1 \otimes \ldots \otimes m_{n - 2}$ du premier
facteur direct est envoyé sur
\begin{align*}
\underbrace{
m_1 \otimes \ldots \otimes x_{i_1} \otimes \ldots
\otimes x_{i_2} \otimes \ldots \otimes m_{n - 2}
}_{\text{avec } x_{i_1} \text{ et } x_{i_2}
\text{ occupant les places } j_1 \text{ et } j_2
\text{ dans le tenseur}} \\
+
\underbrace{
m_1 \otimes \ldots \otimes x_{i_2} \otimes \ldots
\otimes x_{i_1} \otimes \ldots \otimes m_{n - 2}
}_{\text{avec } x_{i_2} \text{ et } x_{i_1}
\text{ occupant les places } j_1 \text{ et } j_2
\text{ dans le tenseur}}
\end{align*}
et $m_1 \otimes \ldots \otimes m_{n - 2}$ du second
facteur direct est envoyé sur
$$
\underbrace{
m_1 \otimes \ldots \otimes x_i \otimes \ldots
\otimes x_i \otimes \ldots \otimes m_{n - 2}
}_{\text{avec } x_{i} \text{ et } x_{i}
\text{ occupant les places } j_1 \text{ et } j_2
\text{ dans le tenseur}}
$$
Il existe aussi une suite exacte canonique
$$
\bigoplus_{1 \leq j_1 < j_2 \leq n}
\bigoplus_{i_1, i_2 \in I}
\text{T}^{n - 2}(M)
\to
\text{T}^n(M)
\to
\text{Sym}^n(M)
\to
0
$$
où le tenseur pur $m_1 \otimes \ldots \otimes m_{n - 2}$ est envoyé sur
\begin{align*}
\underbrace{
m_1 \otimes \ldots \otimes x_{i_1} \otimes \ldots
\otimes x_{i_2} \otimes \ldots \otimes m_{n - 2}
}_{\text{avec } x_{i_1} \text{ et } x_{i_2}
\text{ occupant les places } j_1 \text{ et } j_2
\text{ dans le tenseur}} \\
-
\underbrace{
m_1 \otimes \ldots \otimes x_{i_2} \otimes \ldots
\otimes x_{i_1} \otimes \ldots \otimes m_{n - 2}
}_{\text{avec } x_{i_2} \text{ et } x_{i_1}
\text{ occupant les places } j_1 \text{ et } j_2
\text{ dans le tenseur}}
\end{align*}
\end{lemma}

\begin{proof}
La démonstration est omise.
\end{proof}

\begin{lemma}
\label{algebra-lemma-present-wedge}
Soit $A \to B$ un morphisme d'anneaux. Soit $M$ un $B$-module.
Soit $n > 1$. Le noyau de l'application $A$-linéaire
$M \otimes_A \ldots \otimes_A M \to \wedge^n_B(M)$
est engendré comme $A$-module par les éléments
$m_1 \otimes \ldots \otimes m_n$ tels que $m_i = m_j$
pour $i \not = j$, avec $m_1, \ldots, m_n \in M$, et par les éléments
$m_1 \otimes \ldots \otimes bm_i \otimes \ldots \otimes m_n -
m_1 \otimes \ldots \otimes bm_j \otimes \ldots \otimes m_n$
pour $i \not = j$, $m_1, \ldots, m_n \in M$ et $b \in B$.
\end{lemma}

\begin{proof}
La démonstration est omise.
\end{proof}

\begin{lemma}
\label{algebra-lemma-colimit-tensor-algebra}
\begin{slogan}
La formation des algèbres tensorielles commute aux colimites filtrantes.
\end{slogan}
Soit $R$ un anneau. Supposons que les $M_i$ forment un système filtrant de
$R$-modules. Alors
$\colim_i \text{T}(M_i) = \text{T}(\colim_i M_i)$,
et il en va de même pour les algèbres symétriques et extérieures.
\end{lemma}

\begin{proof}
La démonstration est omise. Indication : appliquer le lemme \ref{algebra-lemma-tensor-products-commute-with-limits}.
\end{proof}

\begin{lemma}
\label{algebra-lemma-tensor-algebra-localization}
Soit $R$ un anneau et soit $S \subset R$ une partie multiplicative.
Alors $S^{-1}T_R(M) = T_{S^{-1}R}(S^{-1}M)$ pour tout $R$-module $M$.
Il en va de même pour les algèbres symétriques et extérieures.
\end{lemma}

\begin{proof}
La démonstration est omise. Indication : appliquer le lemme \ref{algebra-lemma-tensor-product-localization}.
\end{proof}

\section{Changement de base}
\label{algebra-section-base-change}

\noindent
Nous introduisons formellement le changement de base en algèbre comme suit.

\begin{definition}
\label{algebra-definition-base-change}
Soit $\varphi : R \to S$ un morphisme d'anneaux. Soit $M$ un $S$-module.
Soit $R \to R'$ un morphisme d'anneaux quelconque. Le {\it changement de base}
de $\varphi$ par $R \to R'$ est le morphisme d'anneaux $R' \to S \otimes_R R'$.
Dans cette situation, nous écrivons souvent $S' = S \otimes_R R'$.
Le {\it changement de base} du $S$-module $M$ est le $S'$-module
$M \otimes_R R'$.
\end{definition}

\noindent
Si $S = R[x_i]/(f_j)$ pour une famille de variables $x_i$, $i \in I$,
et une famille de polynômes $f_j \in R[x_i]$, $j \in J$, alors
$S \otimes_R R' = R'[x_i]/(f'_j)$, où $f'_j \in R'[x_i]$ est l'image
de $f_j$ par le morphisme $R[x_i] \to R'[x_i]$ induit par $R \to R'$.
Cette remarque simple est la clé de la compréhension du changement de base.

\begin{lemma}
\label{algebra-lemma-base-change-finiteness}
Soit $R \to S$ un morphisme d'anneaux. Soit $M$ un $S$-module.
Soit $R \to R'$ un morphisme d'anneaux, et soient
$S' = S \otimes_R R'$ et $M' = M \otimes_R R'$ les changements de base.
\begin{enumerate}
\item Si $M$ est un $S$-module de type fini, alors son changement de base
$M'$ est un $S'$-module de type fini.
\item Si $M$ est un $S$-module de présentation finie, alors son
changement de base $M'$ est un $S'$-module de présentation finie.
\item Si $R \to S$ est de type fini, alors son changement de base
$R' \to S'$ est de type fini.
\item Si $R \to S$ est de présentation finie, alors son
changement de base $R' \to S'$ est de présentation finie.
\end{enumerate}
\end{lemma}

\begin{proof}
Démonstration de (1). Choisissons une application $S$-linéaire surjective
$S^{\oplus n} \to M \to 0$.
Par les lemmes \ref{algebra-lemma-flip-tensor-product} et \ref{algebra-lemma-tensor-product-exact},
sa tensorisation par $R^\prime$ donne une surjection
${S^\prime}^{\oplus n} \to M^\prime \rightarrow 0$,
de sorte que $M^\prime$ est un $S^\prime$-module de type fini.
Démonstration de (2). Choisissons une présentation
$S^{\oplus m} \to S^{\oplus n} \to M \to 0$.
Par les lemmes \ref{algebra-lemma-flip-tensor-product} et \ref{algebra-lemma-tensor-product-exact},
sa tensorisation par $R^\prime$ donne une présentation finie
${S^\prime}^{\oplus m} \to {S^\prime}^{\oplus n} \to M^\prime \to 0$
du $S^\prime$-module $M^\prime$. Démonstration de (3). Cela résulte
de la remarque précédant le lemme, puisque nous pouvons supposer $I$ fini.
Démonstration de (4). Cela résulte de la remarque précédant le lemme,
puisque nous pouvons supposer $I$ et $J$ finis.
\end{proof}

\noindent
Soit $\varphi : R \to S$ un morphisme d'anneaux. Étant donné un $S$-module
$N$, nous obtenons un $R$-module $N_R$ par la règle
$r \cdot n = \varphi(r)n$.
On l'appelle parfois la {\it restriction des scalaires} de $N$ à $R$.

\begin{lemma}
\label{algebra-lemma-adjoint-tensor-restrict}
Soit $R \to S$ un morphisme d'anneaux. Les foncteurs
$\text{Mod}_S \to \text{Mod}_R$, $N \mapsto N_R$ (restriction des scalaires),
et $\text{Mod}_R \to \text{Mod}_S$, $M \mapsto M \otimes_R S$
(changement de base) sont adjoints. Plus précisément,
$$
\Hom_R(M, N_R) = \Hom_S(M \otimes_R S, N)
$$
\end{lemma}

\begin{proof}
Si $\alpha : M \to N_R$ est un morphisme de $R$-modules, nous définissons
$\alpha' : M \otimes_R S \to N$ par la règle
$\alpha'(m \otimes s) = s\alpha(m)$. Si $\beta : M \otimes_R S \to N$
est un morphisme de $S$-modules, nous définissons
$\beta' : M \to N_R$ par la règle
$\beta'(m) = \beta(m \otimes 1)$.
Nous omettons la vérification que ces constructions sont inverses l'une de l'autre.
\end{proof}

\noindent
Le lemme précédent nous dit que la restriction des scalaires admet un adjoint
à gauche, à savoir le changement de base. Elle admet aussi un adjoint à droite.

\begin{lemma}
\label{algebra-lemma-adjoint-hom-restrict}
Soit $R \to S$ un morphisme d'anneaux. Les foncteurs
$\text{Mod}_S \to \text{Mod}_R$, $N \mapsto N_R$ (restriction des scalaires),
et $\text{Mod}_R \to \text{Mod}_S$, $M \mapsto \Hom_R(S, M)$
sont adjoints. Plus précisément,
$$
\Hom_R(N_R, M) = \Hom_S(N, \Hom_R(S, M))
$$
\end{lemma}

\begin{proof}
Si $\alpha : N_R \to M$ est un morphisme de $R$-modules, nous définissons
$\alpha' : N \to \Hom_R(S, M)$ par la règle
$\alpha'(n) = (s \mapsto \alpha(sn))$. Si $\beta : N \to \Hom_R(S, M)$
est un morphisme de $S$-modules, nous définissons
$\beta' : N_R \to M$ par la règle
$\beta'(n) = \beta(n)(1)$.
Nous omettons la vérification que ces constructions sont inverses l'une de l'autre.
\end{proof}

\begin{lemma}
\label{algebra-lemma-hom-from-tensor-product-variant}
Soit $R \to S$ un morphisme d'anneaux. Étant donnés des $S$-modules $M, N$
et un $R$-module $P$, nous avons
$$
\Hom_R(M \otimes_S N, P) = \Hom_S(M, \Hom_R(N, P))
$$
\end{lemma}

\begin{proof}
On peut le démontrer directement, mais c'est aussi une conséquence
des lemmes \ref{algebra-lemma-adjoint-hom-restrict} et \ref{algebra-lemma-hom-from-tensor-product}.
En effet, nous avons
\begin{align*}
\Hom_R(M \otimes_S N, P)
& =
\Hom_S(M \otimes_S N, \Hom_R(S, P)) \\
& =
\Hom_S(M, \Hom_S(N, \Hom_R(S, P))) \\
& =
\Hom_S(M, \Hom_R(N, P))
\end{align*}
comme voulu.
\end{proof}

\section{Divers}
\label{algebra-section-miscellany}

\noindent
Les démonstrations de cette section ne doivent faire appel à aucun résultat
autre que ceux de la section sur les notions de base, la section \ref{algebra-section-rings-basic}.

\begin{lemma}
\label{algebra-lemma-product-ideals-in-prime}
Soient $R$ un anneau, $I$ et $J$ deux idéaux, et $\mathfrak p$ un idéal premier
contenant le produit $IJ$. Alors $\mathfrak{p}$ contient $I$ ou $J$.
\end{lemma}

\begin{proof}
Supposons le contraire et choisissons $x \in I \setminus \mathfrak p$ et
$y \in J \setminus \mathfrak p$. Leur produit appartient à
$IJ \subset \mathfrak p$, ce qui contredit le fait que
$\mathfrak p$ est premier.
\end{proof}

\begin{lemma}[Évitement des idéaux premiers]
\label{algebra-lemma-silly}
\begin{slogan}
1. Dans un schéma affine, tout ensemble fini de points contenu dans un
ouvert est contenu dans un ouvert principal plus petit.
2. Les ouverts affines sont cofinaux parmi les voisinages d'une partie finie
donnée d'un schéma affine.
\end{slogan}
Soit $R$ un anneau. Soient $I_i \subset R$, $i = 1, \ldots, r$,
et $J \subset R$ des idéaux. Supposons que
\begin{enumerate}
\item $J \not\subset I_i$ pour $i = 1, \ldots, r$, et
\item tous les $I_i$ sauf deux sont des idéaux premiers.
\end{enumerate}
Alors il existe $x \in J$ tel que $x\not\in I_i$ pour tout $i$.
\end{lemma}

\begin{proof}
Le résultat est vrai pour $r = 1$. Si $r = 2$, choisissons $x, y \in J$ tels que
$x \not \in I_1$ et $y \not \in I_2$. Nous avons conclu, sauf si $x \in I_2$
et $y \in I_1$. Mais alors l'élément $x + y$ ne peut appartenir à $I_1$
(car cela entraînerait $x + y - y \in I_1$), ni à $I_2$.

\medskip\noindent
Pour $r \geq 3$, supposons le résultat vrai pour $r - 1$. Après renumérotation,
nous pouvons supposer que $I_r$ est premier. Nous pouvons aussi supposer
qu'il n'existe aucune inclusion entre les $I_i$. Choisissons $x \in J$ tel que
$x \not \in I_i$ pour tout $i = 1, \ldots, r - 1$. Si $x \not\in I_r$,
nous avons conclu. Supposons donc $x \in I_r$. Si
$J I_1 \ldots I_{r - 1} \subset I_r$, alors
$J \subset I_r$ (par le lemme \ref{algebra-lemma-product-ideals-in-prime}), contradiction.
Choisissons $y \in J I_1 \ldots I_{r - 1}$ tel que $y \not \in I_r$.
Alors $x + y$ convient.
\end{proof}

\begin{lemma}
\label{algebra-lemma-silly-silly}
Soit $R$ un anneau. Soient $x \in R$, $I \subset R$ un idéal, et
$\mathfrak p_i$, $i = 1, \ldots, r$, des idéaux premiers.
Supposons que $x + I \not \subset \mathfrak p_i$ pour
$i = 1, \ldots, r$. Alors il existe $y \in I$
tel que $x + y \not \in \mathfrak p_i$ pour tout $i$.
\end{lemma}

\begin{proof}
Nous pouvons supposer qu'il n'existe aucune inclusion entre les $\mathfrak p_i$.
Après réordonnement, nous pouvons supposer $x \not \in \mathfrak p_i$ pour $i < s$
et $x \in \mathfrak p_i$ pour $i \geq s$. Si $s = r + 1$, nous avons conclu.
Sinon, nous pouvons trouver $y \in I$ tel que $y \not \in \mathfrak p_s$.
Choisissons $f \in \bigcap_{i < s} \mathfrak p_i$ tel que
$f \not \in \mathfrak p_s$.
Alors $x + fy$ n'appartient à aucun de $\mathfrak p_1, \ldots, \mathfrak p_s$.
Nous concluons donc par récurrence sur $s$.
\end{proof}

\begin{lemma}[Théorème des restes chinois]
\label{algebra-lemma-chinese-remainder}
Soit $R$ un anneau.
\begin{enumerate}
\item Si $I_1, \ldots, I_r$ sont des idéaux tels que $I_a + I_b = R$
lorsque $a \not = b$, alors $I_1 \cap \ldots \cap I_r =
I_1I_2\ldots I_r$ et $R/(I_1I_2\ldots I_r)
\cong R/I_1 \times \ldots \times R/I_r$.
\item Si $\mathfrak m_1, \ldots, \mathfrak m_r$ sont des idéaux maximaux
deux à deux distincts, alors $\mathfrak m_a + \mathfrak m_b = R$ pour
$a \not = b$, et ce qui précède s'applique.
\end{enumerate}
\end{lemma}

\begin{proof}
Démontrons d'abord $I_1 \cap \ldots \cap I_r = I_1 \ldots I_r$,
car cela entraînera aussi l'injectivité du morphisme d'anneaux induit
$R/(I_1 \ldots I_r) \rightarrow R/I_1 \times \ldots \times R/I_r$.
L'inclusion $I_1 \cap \ldots \cap I_r \supset I_1 \ldots I_r$ est toujours
vraie, puisque les idéaux sont stables par multiplication par des éléments
quelconques de l'anneau. Pour démontrer l'autre inclusion, affirmons que les idéaux
$$
I_1 \ldots \hat I_i \ldots I_r,\quad i = 1, \ldots, r
$$
engendrent l'anneau $R$. Nous le démontrons par récurrence sur $r$.
C'est vrai pour $r = 2$. Si $r > 2$, nous voyons que $R$ est la somme des idéaux
$I_1 \ldots \hat I_i \ldots I_{r - 1}$, $i = 1, \ldots, r - 1$.
Par conséquent, $I_r$ est la somme des idéaux
$I_1 \ldots \hat I_i \ldots I_r$, $i = 1, \ldots, r - 1$.
En appliquant le même argument aux idéaux pris dans l'ordre inverse,
nous voyons que $I_1$ est la somme des idéaux
$I_1 \ldots \hat I_i \ldots I_r$, $i = 2, \ldots, r$.
Puisque $R = I_1 + I_r$ par hypothèse, nous voyons que $R$ est la somme
des idéaux affichés ci-dessus. Nous pouvons donc trouver des éléments
$a_i \in I_1 \ldots \hat I_i \ldots I_r$
dont la somme vaut un. En multipliant cette égalité par un élément
de $I_1 \cap \ldots \cap I_r$, nous obtenons l'autre inclusion.
Il reste à montrer que l'application canonique
$R/(I_1 \ldots I_r) \rightarrow R/I_1 \times \ldots \times R/I_r$
est surjective. Pour cela, considérons son action sur l'égalité
$1 = \sum_{i=1}^r a_i$ obtenue ci-dessus. D'une part, un
morphisme d'anneaux envoie 1 sur 1 ; d'autre part, l'image de tout
$a_i$ est nulle dans $R/I_j$ pour $j \neq i$. L'image de $a_i$
dans $R/I_i$ est donc l'élément unité. Ainsi, pour tout élément
$(\bar{b_1}, \ldots, \bar{b_r}) \in R/I_1 \times \ldots \times R/I_r$,
l'élément $\sum_{i=1}^r a_i \cdot b_i$ en est un antécédent dans $R$.

\medskip\noindent
Pour voir (2), la définition même d'idéaux maximaux distincts donne
$\mathfrak{m}_a + \mathfrak{m}_b = R$ pour $a \neq b$,
et ce qui précède s'applique.
\end{proof}

\begin{lemma}
\label{algebra-lemma-matrix-left-inverse}
Soit $R$ un anneau. Soit $n \geq m$. Soit $A$ une matrice
$n \times m$ à coefficients dans $R$. Soit $J \subset R$
l'idéal engendré par les mineurs $m \times m$ de $A$.
\begin{enumerate}
\item Pour tout $f \in J$, il existe une matrice $m \times n$, $B$,
telle que $BA = f 1_{m \times m}$.
\item Si $f \in R$ et $BA = f 1_{m \times m}$ pour une matrice
$m \times n$, $B$, alors $f^m \in J$.
\end{enumerate}
\end{lemma}

\begin{proof}
Pour $I \subset \{1, \ldots, n\}$ avec $|I| = m$, notons
$E_I$ la matrice $m \times n$ de la projection
$$
R^{\oplus n} = \bigoplus\nolimits_{i \in \{1, \ldots, n\}} R
\longrightarrow \bigoplus\nolimits_{i \in I} R
$$
et posons $A_I = E_I A$ ; ainsi, $A_I$ est la matrice $m \times m$
dont les lignes sont celles de $A$ d'indices appartenant à $I$.
Soit $B_I$ la matrice adjointe (transposée de la matrice
des cofacteurs) de $A_I$, c'est-à-dire telle que
$A_I B_I = B_I A_I = \det(A_I) 1_{m \times m}$.
Les mineurs $m \times m$ de $A$ sont les déterminants $\det A_I$
pour tous les $I \subset \{1, \ldots, n\}$ tels que $|I| = m$.
Si $f \in J$, nous pouvons écrire $f = \sum c_I \det(A_I)$ avec
$c_I \in R$. Posons $B = \sum c_I B_I E_I$ ; cela démontre (1).

\medskip\noindent
Si $f 1_{m \times m} = BA$, la formule de Cauchy-Binet
(\ref{algebra-item-cauchy-binet}) donne
$f^m = \sum b_I \det(A_I)$, où $b_I$ est le déterminant
de la matrice $m \times m$ dont les colonnes sont celles de $B$
d'indices appartenant à $I$.
\end{proof}

\begin{lemma}
\label{algebra-lemma-matrix-right-inverse}
Soit $R$ un anneau. Soit $n \geq m$. Soit $A = (a_{ij})$ une matrice
$n \times m$ à coefficients dans $R$, écrite par blocs sous la forme
$$
A =
\left(
\begin{matrix}
A_1 \\
A_2
\end{matrix}
\right)
$$
où $A_1$ est de taille $m \times m$. Soit $B$ la matrice adjointe
(transposée de la matrice des cofacteurs) de $A_1$. Alors
$$
AB = 
\left(
\begin{matrix}
f 1_{m \times m} \\
C
\end{matrix}
\right)
$$
où $f = \det(A_1)$ et où $c_{ij}$ est, au signe près, le déterminant du
mineur $m \times m$ de $A$ correspondant aux lignes
$1, \ldots, \hat j, \ldots, m, i$.
\end{lemma}

\begin{proof}
Puisque la matrice adjointe vérifie $A_1B = B A_1 = f$, le premier bloc
de l'expression de $AB$ est correct. Remarquons que
$$
c_{ij} = \sum\nolimits_k a_{ik}b_{kj} = \sum (-1)^{j + k}a_{ik} \det(A_1^{jk})
$$
où $A_1^{ij}$ désigne la matrice $A_1$ privée de sa $j$-ième ligne
et de sa $k$-ième colonne. Cette dernière expression est le développement
suivant une ligne du déterminant de la matrice de l'énoncé du lemme.
\end{proof}

\begin{lemma}
\label{algebra-lemma-map-cannot-be-injective}
\begin{slogan}
Un morphisme de modules libres de type fini ne peut être injectif si la source
est de rang strictement supérieur à celui du but.
\end{slogan}
Soit $R$ un anneau non nul. Soit $n \geq 1$. Soit $M$ un $R$-module engendré
par un nombre d'éléments $< n$. Alors tout morphisme de $R$-modules
$f : R^{\oplus n} \to M$ a un noyau non nul.
\end{lemma}

\begin{proof}
Choisissons une surjection $R^{\oplus n - 1} \to M$.
Nous pouvons relever le morphisme $f$ en un morphisme
$f' : R^{\oplus n} \to R^{\oplus n - 1}$
(lemme \ref{algebra-lemma-lift-map}).
Il suffit de démontrer que $f'$ a un noyau non nul.
Le morphisme $f' : R^{\oplus n} \to R^{\oplus n - 1}$ est donné par une
matrice $A = (a_{ij})$. Si l'un des $a_{ij}$ n'est pas nilpotent, disons
que $a = a_{ij}$ ne l'est pas, nous pouvons remplacer $R$ par la localisation
$R_a$ et supposer que $a_{ij}$ est une unité. En effet, si nous trouvons un
noyau non nul après localisation, il existait déjà un noyau non nul,
car la localisation est exacte ; voir la proposition
\ref{algebra-proposition-localization-exact}.
Dans ce cas, nous pouvons effectuer un changement de base à la fois dans
$R^{\oplus n}$ et dans $R^{\oplus n - 1}$ et nous ramener au cas où
$$
A =
\left(
\begin{matrix}
1 & 0 & 0 & \ldots \\
0 & a_{22} & a_{23} & \ldots \\
0 & a_{32} & \ldots \\
\ldots & \ldots
\end{matrix}
\right)
$$
Nous concluons alors par récurrence sur $n$. Sinon, chaque
$a_{ij}$ est nilpotent. Posons $I = (a_{ij}) \subset R$. Remarquons que
$I^{m + 1} = 0$ pour un certain $m \geq 0$. Soit $m$ le plus grand entier
tel que $I^m \not = 0$. Alors $(I^m)^{\oplus n}$ est
contenu dans le noyau du morphisme, ce qui conclut.
\end{proof}

\begin{lemma}
\label{algebra-lemma-rank}
\begin{slogan}
Le rang d'un module libre de type fini est bien défini.
\end{slogan}
Soit $R$ un anneau non nul. Soient $n, m \geq 0$ des entiers.
Si $R^{\oplus n}$ est isomorphe à $R^{\oplus m}$ comme
$R$-modules, alors $n = m$.
\end{lemma}

\begin{proof}
C'est immédiat d'après le lemme \ref{algebra-lemma-map-cannot-be-injective}.
\end{proof}

\section{Cayley-Hamilton}
\label{algebra-section-cayley-hamilton}

\begin{lemma}
\label{algebra-lemma-charpoly}
Soit $R$ un anneau. Soit $A = (a_{ij})$ une matrice $n \times n$
à coefficients dans $R$. Soit $P(x) \in R[x]$
le polynôme caractéristique de $A$ (défini
comme $\det(x\text{id}_{n \times n} - A)$).
Alors $P(A) = 0$ dans $\text{Mat}(n \times n, R)$.
\end{lemma}

\begin{proof}
Nous ramenons la question au théorème de Cayley-Hamilton bien connu
d'algèbre linéaire, en plusieurs étapes :
\begin{enumerate}
\item Si $\phi :S \rightarrow R$ est un morphisme d'anneaux et si les $b_{ij}$
sont des antécédents des $a_{ij}$ par ce morphisme, il suffit de démontrer
l'énoncé pour $S$ et $(b_{ij})$, puisque $\phi$ est un morphisme d'anneaux.
\item Si $\psi :R \hookrightarrow S$ est un morphisme injectif d'anneaux,
il suffit clairement de démontrer le résultat pour $S$ et pour les $a_{ij}$
considérés comme éléments de $S$.
\item Nous pouvons donc d'abord nous ramener au cas $R = \mathbf{Z}[X_{ij}]$,
$a_{ij} = X_{ij}$, d'un anneau de polynômes, puis au
cas $R = \mathbf{Q}(X_{ij})$, où nous pouvons enfin appliquer Cayley-Hamilton.
\end{enumerate}
\end{proof}

\begin{lemma}
\label{algebra-lemma-charpoly-module}
Soit $R$ un anneau.
Soit $M$ un $R$-module de type fini.
Soit $\varphi : M \to M$ un endomorphisme.
Alors il existe un polynôme unitaire $P \in R[T]$ tel que
$P(\varphi) = 0$ comme endomorphisme de $M$.
\end{lemma}

\begin{proof}
Choisissons un morphisme surjectif de $R$-modules $R^{\oplus n} \to M$, donné par
$(a_1, \ldots, a_n) \mapsto \sum a_ix_i$, pour des générateurs $x_i \in M$.
Choisissons $(a_{i1}, \ldots, a_{in}) \in R^{\oplus n}$ tels que
$\varphi(x_i) = \sum a_{ij} x_j$. Autrement dit, le diagramme
$$
\xymatrix{
R^{\oplus n} \ar[d]_A \ar[r] & M \ar[d]^\varphi \\
R^{\oplus n} \ar[r] & M
}
$$
est commutatif, où $A = (a_{ij})$. D'après
le lemme \ref{algebra-lemma-charpoly},
il existe un polynôme unitaire $P$ tel que $P(A) = 0$.
Il s'ensuit que $P(\varphi) = 0$.
\end{proof}

\begin{lemma}
\label{algebra-lemma-charpoly-module-ideal}
Soit $R$ un anneau. Soit $I \subset R$ un idéal.
Soit $M$ un $R$-module de type fini.
Soit $\varphi : M \to M$ un endomorphisme tel
que $\varphi(M) \subset IM$.
Alors il existe un polynôme unitaire
$P = t^n + a_1 t^{n - 1} + \ldots + a_n \in R[T]$
tel que $a_j \in I^j$ et que $P(\varphi) = 0$ comme endomorphisme de $M$.
\end{lemma}

\begin{proof}
Choisissons un morphisme surjectif de $R$-modules $R^{\oplus n} \to M$, donné par
$(a_1, \ldots, a_n) \mapsto \sum a_ix_i$, pour des générateurs $x_i \in M$.
Choisissons $(a_{i1}, \ldots, a_{in}) \in I^{\oplus n}$ tels que
$\varphi(x_i) = \sum a_{ij} x_j$. Autrement dit, le diagramme
$$
\xymatrix{
R^{\oplus n} \ar[d]_A \ar[r] & M \ar[d]^\varphi \\
I^{\oplus n} \ar[r] & M
}
$$
est commutatif, où $A = (a_{ij})$. D'après
le lemme \ref{algebra-lemma-charpoly},
le polynôme
$P(t) = \det(t\text{id}_{n \times n} - A)$
possède toutes les propriétés voulues.
\end{proof}

\noindent
À titre d'application plaisante, démontrons le lemme surprenant suivant.

\begin{lemma}
\label{algebra-lemma-fun}
Soit $R$ un anneau.
Soit $M$ un $R$-module de type fini.
Soit $\varphi : M \to M$ un morphisme surjectif de $R$-modules.
Alors $\varphi$ est un isomorphisme.
\end{lemma}

\begin{proof}[Première démonstration]
Écrivons $R' = R[x]$ et considérons $M$ comme un $R'$-module de type fini,
$x$ agissant par $\varphi$. Posons $I = (x) \subset R'$. L'hypothèse de
surjectivité de $\varphi$ donne $IM = M$. Nous pouvons donc appliquer
le lemme \ref{algebra-lemma-charpoly-module-ideal}
à $M$ considéré comme un $R'$-module, à l'idéal $I$ et à l'endomorphisme
$\text{id}_M$.
Nous en déduisons que
$(1 + a_1 + \ldots + a_n)\text{id}_M = 0$ avec $a_j \in I$.
Écrivons $a_j = b_j(x)x$ pour un certain $b_j(x) \in R[x]$.
En revenant à $\varphi$, nous voyons que
$\text{id}_M = -(\sum_{j = 1, \ldots, n} b_j(\varphi)) \varphi$,
et donc que $\varphi$ est inversible.
\end{proof}

\begin{proof}[Seconde démonstration]
Nous raisonnons par récurrence sur le nombre de générateurs de $M$ sur $R$. Si
$M$ est engendré par un élément, alors $M \cong R/I$ pour un certain idéal
$I \subset R$. Nous pouvons alors remplacer $R$ par $R/I$, de sorte que $M = R$.
Dans ce cas, $\varphi : R \to R$ est donné par la multiplication sur $M$ par un
élément $r \in R$. La surjectivité de $\varphi$ force $r$ à être inversible,
puisque $\varphi$ doit atteindre $1$ ; il s'ensuit que $\varphi$ est
inversible.

\medskip\noindent
Supposons maintenant le lemme démontré pour les modules
engendrés par $n - 1$ éléments, et considérons un module $M$ engendré
par $n$ éléments. Notons $A$ l'anneau $R[t]$ et regardons le module
$M$ comme un $A$-module en faisant agir $t$ par $\varphi$ ; comme $M$ est
de type fini sur $R$, il l'est aussi sur $R[t]$. Puisque nous cherchons
à démontrer que $\varphi$ est injectif, ce qui est une propriété ensembliste,
nous pouvons tout aussi bien démontrer que l'endomorphisme
$t : M \to M$ sur $A$ est injectif. Nous avons ramené le problème au cas
où notre endomorphisme est la multiplication par un élément de l'anneau de base.
Notons $M' \subset M$ le sous-$A$-module engendré par les $n - 1$ premiers
générateurs de $M$, et considérons le diagramme
$$
\xymatrix{
0 \ar[r] & M' \ar[r]\ar[d]^{\varphi\mid_{M'}} & M\ar[d]^\varphi \ar[r] &
M/M' \ar[d]^{\varphi \bmod M'} \ar[r] & 0 \\
0 \ar[r] & M' \ar[r]                                  & M \ar[r]
           & M/M' \ar[r]                                  & 0,
}
$$
où la restriction de $\varphi$ à $M'$ et l'application induite par $\varphi$
sur le quotient $M/M'$ sont bien définies, puisque $\varphi$ est la
multiplication par un élément de la base, et que $M'$ et $M/M'$ sont
eux-mêmes des $A$-modules. Par le cas $n = 1$, l'application
$M/M' \to M/M'$ est un isomorphisme. Une chasse au diagramme montre que
$\varphi|_{M'}$ est surjective ; par récurrence, $\varphi|_{M'}$ est donc
un isomorphisme. Le lemme du serpent impose alors à la colonne du milieu
d'être un isomorphisme.
\end{proof}

\section{Le spectre d'un anneau}
\label{algebra-section-spectrum-ring}

\noindent
Nous décidons arbitrairement que le spectre d'un anneau, considéré comme espace
topologique, relève du chapitre d'algèbre, tandis qu'un schéma affine relève du
chapitre sur les schémas.

\begin{definition}
\label{algebra-definition-spectrum-ring}
Soit $R$ un anneau.
\begin{enumerate}
\item Le {\it spectre} de $R$ est l'ensemble des idéaux premiers de $R$.
On le note habituellement $\Spec(R)$.
\item Étant donnée une partie $T \subset R$, nous notons
$V(T) \subset \Spec(R)$ l'ensemble des idéaux premiers contenant $T$,
c'est-à-dire $V(T) = \{ \mathfrak p \in
\Spec(R) \mid \forall f\in T, f\in \mathfrak p\}$.
\item Étant donné un élément $f \in R$, nous notons $D(f) \subset \Spec(R)$
l'ensemble des idéaux premiers ne contenant pas $f$.
\end{enumerate}
\end{definition}

\begin{lemma}
\label{algebra-lemma-Zariski-topology}
Soit $R$ un anneau.
\begin{enumerate}
\item Le spectre d'un anneau $R$ est vide si et seulement si $R$
est l'anneau nul.
\item Tout anneau non nul possède un idéal maximal.
\item Tout anneau non nul possède un idéal premier minimal.
\item Étant donnés un idéal $I \subset R$ et un idéal premier vérifiant
$I \subset \mathfrak p$, il existe un idéal premier
$I \subset \mathfrak q \subset \mathfrak p$ tel que
$\mathfrak q$ soit minimal au-dessus de $I$.
\item Si $T \subset R$ et si $(T)$ est l'idéal engendré par
$T$ dans $R$, alors $V((T)) = V(T)$.
\item Si $I$ est un idéal et si $\sqrt{I}$ est son radical,
voir la notion de base (\ref{algebra-item-radical-ideal}), alors $V(I) = V(\sqrt{I})$.
\item Pour tout idéal $I$ de $R$, nous avons $\sqrt{I} =
\bigcap_{I \subset \mathfrak p} \mathfrak p$.
\item Si $I$ est un idéal, alors $V(I) = \emptyset$ si et seulement
si $I$ est l'idéal unité.
\item Si $I$ et $J$ sont des idéaux de $R$, alors $V(I) \cup V(J) =
V(I \cap J)$.
\item Si $(I_a)_{a\in A}$ est une famille d'idéaux de $R$, alors
$\bigcap_{a\in A} V(I_a) = V(\bigcup_{a\in A} I_a)$.
\item Si $f \in R$, alors $D(f) \amalg V(f) = \Spec(R)$.
\item Si $f \in R$, alors $D(f) = \emptyset$ si et seulement si $f$
est nilpotent.
\item Si $f = u f'$ pour une unité $u \in R$, alors $D(f) = D(f')$.
\item Si $I \subset R$ est un idéal et si $\mathfrak p$ est un idéal premier de
$R$ tel que $\mathfrak p \not\in V(I)$, alors il existe $f \in R$
tel que $\mathfrak p \in D(f)$ et $D(f) \cap V(I) = \emptyset$.
\item Si $f, g \in R$, alors $D(fg) = D(f) \cap D(g)$.
\item Si $f_i \in R$ pour $i \in I$, alors
$\bigcup_{i\in I} D(f_i)$ est le complémentaire de
$V(\{f_i \}_{i\in I})$ dans $\Spec(R)$.
\item Si $f \in R$ et $D(f) = \Spec(R)$, alors $f$ est une unité.
\end{enumerate}
\end{lemma}

\begin{proof}
Nous traitons chaque assertion dans l'item correspondant ci-dessous.
\begin{enumerate}
\item C'est une conséquence directe de (2) ou de (3).
\item Soit $\mathfrak{A}$ l'ensemble de tous les idéaux propres de $R$.
Cet ensemble est ordonné par inclusion et n'est pas vide, puisque
$(0) \in \mathfrak{A}$ est un idéal propre.
Soit $A$ une partie totalement ordonnée de $\mathfrak A$.
Alors $\bigcup_{I \in A} I$ est en fait un idéal. Comme 1 $\notin I$
pour tout $I \in A$, la réunion ne contient pas 1 et est donc propre.
Ainsi, $\bigcup_{I \in A} I$ appartient à $\mathfrak{A}$ et majore
la partie $A$. Le lemme de Zorn assure donc à $\mathfrak{A}$
l'existence d'un élément maximal, qui est l'idéal maximal recherché.
\item Puisque $R$ est non nul, il contient un idéal maximal, donc premier.
Ainsi, l'ensemble $\mathfrak{A}$ de tous les idéaux premiers de $R$
n'est pas vide. $\mathfrak{A}$ est ordonné par l'inclusion opposée.
Soit $A$ une partie totalement ordonnée de $\mathfrak{A}$. Il est clair que
$J = \bigcap_{I \in A} I$ est un idéal. Il est moins clair qu'il soit premier.
Soit $xy \in J$. Alors $xy \in I$ pour tout $I \in A$.
Posons $B = \{I \in A | y \in I\}$ et $K
= \bigcap_{I \in B} I$. Puisque $A$ est totalement ordonné, ou bien
$K = J$ (et nous avons conclu, car alors $y \in J$), ou bien $K \supset J$
et, pour tout $I \in A$ tel que $I$ soit strictement contenu dans $K$,
nous avons $y \notin I$. Mais cela signifie que, pour tous ces idéaux
$I, x \in I$, puisqu'ils sont premiers. Ainsi $x \in J$. Dans les deux cas,
$J$ est premier, comme voulu. Le lemme de Zorn fournit donc un élément
maximal, qui est ici un idéal premier minimal.
\item C'est exactement le même argument que dans (3), en ne considérant que
les idéaux premiers contenus dans $\mathfrak{p}$ et contenant $I$.
\item $(T)$ est le plus petit idéal contenant $T$. Par conséquent, si
$T \subset I$ pour un certain idéal, alors $(T) \subset I$ également.
Ainsi, si $I \in V(T)$, alors $I \in V((T))$ également.
L'autre inclusion est évidente.
\item Puisque $I \subset \sqrt{I}, V(\sqrt{I}) \subset V(I)$, prenons
$\mathfrak{p} \in V(I)$. Soit $x \in \sqrt{I}$. Alors $x^n \in I$
pour un certain $n$. Donc $x^n \in \mathfrak{p}$. Comme $\mathfrak{p}$
est premier, une récurrence immédiate donne $x \in \mathfrak{p}$.
Ainsi $\sqrt{I} \subset \mathfrak{p}$ et
$\mathfrak{p} \in V(\sqrt{I})$.
\item Soit $f \in R \setminus \sqrt{I}$. Alors $f^n \notin I$
pour tout $n$. Par conséquent,
$S = \{1, f, f^2, \ldots\}$ est une partie multiplicative ne contenant pas $0$.
Choisissons un idéal premier
$\bar{\mathfrak{p}} \subset S^{-1}R$ contenant $S^{-1}I$.
Alors $\mathfrak{p}$, image réciproque dans $R$ de $\bar{\mathfrak{p}}$,
est un idéal premier contenant $I$ et ne rencontrant pas $S$.
Cela montre que $\bigcap_{I \subset
\mathfrak p} \mathfrak p \subset \sqrt{I}$. Maintenant, si
$a \in \sqrt{I}$, alors $a^n \in I$ pour un certain $n$.
Ainsi, si $I \subset \mathfrak{p}$, alors $a^n \in
\mathfrak{p}$. Comme $\mathfrak{p}$ est premier, nous avons
$a \in \mathfrak{p}$. L'égalité est donc démontrée.
\item $I$ n'est pas l'idéal unité si et seulement si $I$
est contenu dans un idéal maximal (pour le voir, appliquer (2) à l'anneau
$R/I$), lequel est donc premier.
\item Si $\mathfrak{p} \in V(I) \cup V(J)$, alors $I \subset \mathfrak{p}$
ou $J \subset \mathfrak{p}$, ce qui signifie que
$I \cap J \subset \mathfrak{p}$. Réciproquement, si
$I \cap J \subset \mathfrak{p}$, alors $IJ \subset \mathfrak{p}$ ;
ainsi, soit $I$, soit $J$, est contenu dans $\mathfrak{p}$,
puisque $\mathfrak{p}$ est premier.
\item $\mathfrak{p} \in \bigcap_{a \in A} V(I_a) \Leftrightarrow
I_a \subset \mathfrak{p}, \forall a \in A \Leftrightarrow
\mathfrak{p} \in V(\bigcup_{a\in A} I_a)$
\item Si $\mathfrak{p}$ est un idéal premier et $f \in R$, alors ou bien
$f \in \mathfrak{p}$, ou bien $f \notin \mathfrak{p}$, de façon exclusive ;
c'est précisément ce qu'exprime la réunion disjointe.
\item Si $a \in R$ est nilpotent, alors $a^n = 0$ pour un certain $n$.
Ainsi $a^n \in \mathfrak{p}$ pour tout idéal premier. Une récurrence donne
donc $a \in \mathfrak{p}$ et $D(a) = \emptyset$. Réciproquement, comme
nous l'avons montré en (7), si $a \in R$ n'est pas nilpotent, il existe
un idéal premier qui ne le contient pas.
\item $f \in \mathfrak{p} \Leftrightarrow uf \in \mathfrak{p}$,
puisque $u$ est inversible.
\item Si $\mathfrak{p} \notin V(I)$, alors
$\exists f \in I \setminus \mathfrak{p}$. Ainsi
$f \notin \mathfrak{p}$, donc $\mathfrak{p} \in D(f)$. En outre, si
$\mathfrak{q} \in D(f)$, alors $f \notin \mathfrak{q}$, si bien que $I$
n'est pas contenu dans $\mathfrak{q}$. Ainsi
$D(f) \cap V(I) = \emptyset$.
\item Si $fg \in \mathfrak{p}$, alors $f \in \mathfrak{p}$ ou
$g \in \mathfrak{p}$. Par conséquent, si $f \notin \mathfrak{p}$ et
$g \notin \mathfrak{p}$, alors $fg \notin \mathfrak{p}$.
Puisque $\mathfrak{p}$ est un idéal, si $fg \notin
\mathfrak{p}$, alors $f \notin \mathfrak{p}$ et $g \notin \mathfrak{p}$.
\item $\mathfrak{p} \in \bigcup_{i \in I} D(f_i) \Leftrightarrow \exists i \in
I, f_i \notin \mathfrak{p} \Leftrightarrow \mathfrak{p} \in \Spec(R)
\setminus V(\{f_i\}_{i \in I})$
\item Si $D(f) = \Spec(R)$, alors $V(f) = \emptyset$ et donc
$fR = R$, de sorte que $f$ est une unité.
\end{enumerate}
\end{proof}

\noindent
Le lemme entraîne que les parties $V(T)$ de la
définition \ref{algebra-definition-spectrum-ring} sont les fermés
d'une topologie sur $\Spec(R)$. Il montre aussi que
les ensembles $D(f)$ sont ouverts et forment une base de cette
topologie.

\begin{definition}
\label{algebra-definition-Zariski-topology}
Soit $R$ un anneau.
La topologie sur $\Spec(R)$ dont les fermés sont les
ensembles $V(T)$ est appelée topologie {\it de Zariski}. Les ouverts
$D(f)$ sont appelés les {\it ouverts principaux} de $\Spec(R)$.
\end{definition}

\noindent
Le contexte indiquera clairement si nous considérons $\Spec(R)$
seulement comme un ensemble ou comme un espace topologique.

\begin{lemma}
\label{algebra-lemma-spec-functorial}
\begin{slogan}
Fonctorialité du spectre
\end{slogan}
Supposons que $\varphi : R \to R'$ soit un morphisme d'anneaux.
L'application induite
$$
\Spec(\varphi) : \Spec(R') \longrightarrow \Spec(R),
\quad
\mathfrak p' \longmapsto \varphi^{-1}(\mathfrak p')
$$
est continue pour les topologies de Zariski. En fait, pour tout
élément $f \in R$, nous avons
$\Spec(\varphi)^{-1}(D(f)) = D(\varphi(f))$.
\end{lemma}

\begin{proof}
La notion de base (\ref{algebra-item-inverse-image-prime}) assure que
$\mathfrak p := \varphi^{-1}(\mathfrak p')$
est bien un idéal premier de $R$. La dernière assertion
du lemme résulte directement des définitions
et entraîne la première.
\end{proof}

\noindent
Si $\varphi' : R' \to R''$ est un second morphisme d'anneaux,
alors la composée
$$
\Spec(R'')
\longrightarrow
\Spec(R')
\longrightarrow
\Spec(R)
$$
est égale à $\Spec(\varphi' \circ \varphi)$. Autrement dit,
$\Spec$ est un foncteur contravariant de la catégorie des anneaux
vers la catégorie des espaces topologiques.

\begin{lemma}
\label{algebra-lemma-spec-localization}
Soit $R$ un anneau. Soit $S \subset R$ une partie multiplicative.
Le morphisme $R \to S^{-1}R$ induit, par fonctorialité de $\Spec$,
un homéomorphisme
$$
\Spec(S^{-1}R)
\longrightarrow
\{\mathfrak p \in \Spec(R) \mid S \cap \mathfrak p = \emptyset \}
$$
où le membre de droite est muni de la topologie induite par la
topologie de Zariski sur $\Spec(R)$. L'application inverse est donnée
par $\mathfrak p \mapsto S^{-1}\mathfrak p = \mathfrak p(S^{-1}R)$.
\end{lemma}

\begin{proof}
Notons $D$ le membre de droite de la flèche du lemme.
Choisissons un idéal premier $\mathfrak p' \subset S^{-1}R$ et soit
$\mathfrak p$ l'image réciproque de $\mathfrak p'$ dans $R$.
Puisque $\mathfrak p'$ ne contient pas $1$, nous voyons que $\mathfrak p$
ne contient aucun élément de $S$. Ainsi $\mathfrak p \in D$, et nous voyons que
l'image est contenue dans $D$. Soit $\mathfrak p \in D$.
Par hypothèse, l'image $\overline{S}$ ne contient pas $0$.
D'après la notion de base (\ref{algebra-item-localization-zero}),
$\overline{S}^{-1}(R/\mathfrak p)$ n'est pas l'anneau nul.
D'après la notion de base (\ref{algebra-item-localize-ideal}), nous voyons que
$S^{-1}R / S^{-1}\mathfrak p = \overline{S}^{-1}(R/\mathfrak p)$
est un anneau intègre, et donc que $S^{-1}\mathfrak p$ est premier.
L'égalité d'anneaux montre aussi que l'image réciproque de
$S^{-1}\mathfrak p$ dans $R$ est égale à $\mathfrak p$,
car $R/\mathfrak p \to \overline{S}^{-1}(R/\mathfrak p)$
est injective d'après la notion de base
(\ref{algebra-item-localize-nonzerodivisors}).
Cela démontre que l'application $\Spec(S^{-1}R) \to \Spec(R)$
est une bijection sur $D$, dont l'inverse est celui indiqué.
Elle est continue d'après le lemme \ref{algebra-lemma-spec-functorial}.
Enfin, soit $D(g) \subset \Spec(S^{-1}R)$ un ouvert principal.
Écrivons $g = h/s$ pour certains $h\in R$ et $s\in S$.
Puisque $g$ s'obtient à partir de $h/1$ par multiplication par une unité, nous avons
$D(g) = D(h/1)$ dans $\Spec(S^{-1}R)$.
Le lemme \ref{algebra-lemma-spec-functorial} et la bijectivité ci-dessus montrent donc
que l'image de $D(g) = D(h/1)$ est $D \cap D(h)$.
Cela prouve que l'application est également ouverte.
\end{proof}

\begin{lemma}
\label{algebra-lemma-standard-open}
Soit $R$ un anneau. Soit $f \in R$.
Le morphisme $R \to R_f$ induit, par fonctorialité de
$\Spec$, un homéomorphisme
$$
\Spec(R_f) \longrightarrow D(f) \subset \Spec(R).
$$
L'application inverse est donnée par
$\mathfrak p \mapsto \mathfrak p \cdot R_f$.
\end{lemma}

\begin{proof}
C'est un cas particulier du lemme \ref{algebra-lemma-spec-localization}.
\end{proof}

\noindent
Tout « ouvert affine » d'un spectre n'est pas nécessairement
un ouvert principal. Voir l'exemple
\ref{algebra-example-affine-open-not-standard}.

\begin{lemma}
\label{algebra-lemma-spec-closed}
Soit $R$ un anneau. Soit $I \subset R$ un idéal.
Le morphisme $R \to R/I$ induit, par fonctorialité de
$\Spec$, un homéomorphisme
$$
\Spec(R/I) \longrightarrow V(I) \subset \Spec(R).
$$
L'application inverse est donnée par
$\mathfrak p \mapsto \mathfrak p / I$.
\end{lemma}

\begin{proof}
Il est immédiat que l'image est contenue dans $V(I)$.
D'autre part, si $\mathfrak p \in V(I)$, alors
$\mathfrak p \supset I$, et nous pouvons considérer
l'idéal $\mathfrak p /I \subset R/I$. La notion de base
(\ref{algebra-item-isomorphism-theorem}) montre que
$(R/I)/(\mathfrak p/I) = R/\mathfrak p$ est un anneau intègre,
et donc que $\mathfrak p/I$ est un idéal premier. Il est alors
immédiat que l'image de $D(f + I)$ est
$D(f) \cap V(I)$, et donc que l'application est un homéomorphisme.
\end{proof}



\begin{lemma}
\label{algebra-lemma-quasi-compact}
\begin{slogan}
Le spectre d'un anneau est quasi-compact
\end{slogan}
Soit $R$ un anneau. L'espace $\Spec(R)$ est quasi-compact.
\end{lemma}

\begin{proof}
Il suffit de démontrer que de tout recouvrement de $\Spec(R)$
par des ouverts principaux on peut extraire un recouvrement fini.
Supposons donc que $\Spec(R) = \cup D(f_i)$
pour une famille d'éléments $\{f_i\}_{i\in I}$ de $R$. Cela signifie que
$\cap V(f_i) = \emptyset$. D'après le lemme
\ref{algebra-lemma-Zariski-topology}, cela signifie que
$V(\{f_i \}) = \emptyset$. D'après le
même lemme, cela signifie que l'idéal engendré
par les $f_i$ est l'idéal unité de $R$. Nous pouvons donc
écrire $1$ comme une somme {\it finie} :
$1 = \sum_{i \in J} r_i f_i$ avec $J \subset I$ fini.
Il s'ensuit alors que $\Spec(R)
= \cup_{i \in J} D(f_i)$.
\end{proof}

\begin{lemma}
\label{algebra-lemma-topology-spec}
Soit $R$ un anneau.
La topologie sur $X = \Spec(R)$ possède les propriétés suivantes :
\begin{enumerate}
\item $X$ est quasi-compact,
\item $X$ possède une base de la topologie formée d'ouverts quasi-compacts, et
\item l'intersection de deux ouverts quasi-compacts quelconques est quasi-compacte.
\end{enumerate}
\end{lemma}

\begin{proof}
Le spectre d'un anneau est quasi-compact ; voir le
lemme \ref{algebra-lemma-quasi-compact}.
Il possède une base de la topologie formée des ouverts principaux
$D(f) = \Spec(R_f)$
(lemme \ref{algebra-lemma-standard-open}),
qui sont quasi-compacts d'après la première assertion.
L'intersection de deux ouverts principaux est quasi-compacte,
car $D(f) \cap D(g) = D(fg)$. Étant donnés deux ouverts quasi-compacts
$U, V \subset X$, nous pouvons écrire $U = D(f_1) \cup \ldots \cup D(f_n)$
et $V = D(g_1) \cup \ldots \cup D(g_m)$. Alors
$U \cap V = \bigcup D(f_ig_j)$, qui est quasi-compact.
\end{proof}

\section{Anneaux locaux}
\label{algebra-section-local-rings}

\noindent
Les anneaux locaux sont au cœur de la géométrie algébrique.

\begin{definition}
\label{algebra-definition-local-ring}
Un {\it anneau local} est un anneau qui possède exactement un idéal maximal.
Si $R$ est un anneau local, on note souvent son idéal maximal
$\mathfrak m_R$, et le corps $R/\mathfrak m_R$ est appelé le
{\it corps résiduel} de l'anneau local $R$.
Nous disons souvent « soit $(R, \mathfrak m)$ un anneau local »
ou « soit $(R, \mathfrak m, \kappa)$ un anneau local »
pour indiquer que $R$ est local, que $\mathfrak m$ est son unique
idéal maximal et que $\kappa = R/\mathfrak m$ est son corps résiduel.
Un {\it morphisme local d'anneaux locaux} est un morphisme d'anneaux
$\varphi : R \to S$ tel que $R$ et $S$ soient des anneaux locaux et
que $\varphi(\mathfrak m_R) \subset \mathfrak m_S$.
S'il est entendu que $R$ et $S$ sont des anneaux locaux, l'expression
« {\it morphisme local d'anneaux $\varphi : R \to S$} » signifie que
$\varphi$ est un morphisme local d'anneaux locaux.
\end{definition}

\noindent
Un corps est un anneau local. Tout morphisme d'anneaux entre corps est
un morphisme local d'anneaux locaux.

\medskip\noindent
La localisation $R_\mathfrak p$ d'un anneau $R$ en un idéal premier
$\mathfrak p$ est un anneau local d'idéal maximal
$\mathfrak p R_\mathfrak p$. En effet, d'après le
lemme \ref{algebra-lemma-spec-localization}, tout idéal premier de $R_\mathfrak p$
est contenu dans l'idéal premier $\mathfrak p R_\mathfrak p$
(qui est donc un idéal maximal, et l'unique idéal maximal de $R_\mathfrak p$).
Le corps résiduel de $R_\mathfrak p$ est noté $\kappa(\mathfrak p)$ ;
nous l'appelons le {\it corps résiduel de $\mathfrak p$}. D'après la
proposition \ref{algebra-proposition-localize-quotient}, nous pouvons identifier
$\kappa(\mathfrak p)$ au corps des fractions de l'anneau intègre
$R/\mathfrak p$. Par la composée
$$
\Spec(\kappa(\mathfrak p)) \to \Spec(R_\mathfrak p) \to \Spec(R)
$$
l'unique point de la source est envoyé sur le point $\mathfrak p$ du but.

\medskip\noindent
Soit $\varphi : R \to S$ un morphisme d'anneaux. Soit
$\mathfrak q \subset S$ un idéal premier, et considérons l'idéal premier
$\mathfrak p = \varphi^{-1}(\mathfrak q)$ de $R$.
Puisque $\varphi(\mathfrak p) \subset \mathfrak q$, le morphisme d'anneaux induit
$$
R_\mathfrak p \to S_\mathfrak q,\quad r/g \mapsto \varphi(r)/\varphi(g)
$$
est un morphisme local d'anneaux locaux, et nous obtenons un morphisme induit
entre les corps résiduels $\kappa(\mathfrak p) \to \kappa(\mathfrak q)$.


\begin{example}
\label{algebra-example-not-local}
Si $R$ est un anneau local et si $\mathfrak p \subset R$ est un idéal premier
non maximal, alors $R \to R_\mathfrak p$ n'est pas un morphisme local.
\end{example}

\begin{lemma}
\label{algebra-lemma-characterize-local-ring}
Soit $R$ un anneau. Les conditions suivantes sont équivalentes :
\begin{enumerate}
\item $R$ est un anneau local,
\item $\Spec(R)$ possède exactement un point fermé,
\item $R$ possède un idéal maximal $\mathfrak m$
et tout élément de $R \setminus \mathfrak m$
est une unité, et
\item $R$ n'est pas l'anneau nul et, pour tout $x \in R$, ou bien $x$,
ou bien $1 - x$, ou bien les deux sont inversibles.
\end{enumerate}
\end{lemma}

\begin{proof}
Soient $R$ un anneau et $\mathfrak m$ un idéal maximal.
Si $x \in R \setminus \mathfrak m$ et si $x$ n'est pas une unité,
alors il existe un idéal maximal $\mathfrak m'$ contenant $x$.
Ainsi $R$ possède au moins deux idéaux maximaux. Réciproquement,
si $\mathfrak m'$ est un autre idéal maximal, choisissons
$x \in \mathfrak m'$ tel que $x \not \in \mathfrak m$. Clairement,
$x$ n'est pas une unité. Cela démontre l'équivalence de (1) et de (3).
L'équivalence de (1) et de (2) est tautologique.
Si $R$ est local, alors (4) est vraie, puisque $x$ appartient ou non à
$\mathfrak m$. Si (4) est vraie et si $\mathfrak m$, $\mathfrak m'$
sont des idéaux maximaux distincts, nous pouvons choisir $x \in R$ tel que
$x \bmod \mathfrak m' = 0$ et $x \bmod \mathfrak m = 1$
par le théorème des restes chinois
(lemme \ref{algebra-lemma-chinese-remainder}).
Cet élément $x$ n'est pas inversible, et $1 - x$ ne l'est pas davantage,
ce qui est une contradiction. Ainsi (4) et (1) sont équivalentes.
\end{proof}

\begin{lemma}
\label{algebra-lemma-characterize-local-ring-map}
Soit $\varphi : R \to S$ un morphisme d'anneaux. Supposons que $R$ et $S$
soient des anneaux locaux. Les conditions suivantes sont équivalentes :
\begin{enumerate}
\item $\varphi$ est un morphisme local d'anneaux,
\item $\varphi(\mathfrak m_R) \subset \mathfrak m_S$,
\item $\varphi^{-1}(\mathfrak m_S) = \mathfrak m_R$, et
\item pour tout $x \in R$, si $\varphi(x)$ est inversible dans $S$, alors $x$
est inversible dans $R$.
\end{enumerate}
\end{lemma}

\begin{proof}
Les conditions (1) et (2) sont équivalentes par définition.
Si (3) est vraie, alors (2) l'est. Réciproquement, si (2) est vraie,
$\varphi^{-1}(\mathfrak m_S)$ est un idéal premier contenant
l'idéal maximal $\mathfrak m_R$, donc
$\varphi^{-1}(\mathfrak m_S) = \mathfrak m_R$. Enfin, (4) est la
contraposée de (2) d'après le lemme \ref{algebra-lemma-characterize-local-ring}.
\end{proof}

\begin{remark}
\label{algebra-remark-fundamental-diagram}
Un diagramme commutatif fondamental associé à un morphisme d'anneaux
$\varphi : R \to S$ et à un idéal premier $\mathfrak p \subset R$ est le suivant :
$$
\xymatrix{
\kappa(\mathfrak p) \otimes_R S =
S_{\mathfrak p}/{\mathfrak p}S_{\mathfrak p}
&
S_{\mathfrak p} \ar[l]
&
S \ar[r] \ar[l]
&
S/\mathfrak pS \ar[r]
&
(R \setminus \mathfrak p)^{-1}S/\mathfrak pS
\\
\kappa(\mathfrak p) =
R_{\mathfrak p}/{\mathfrak p}R_{\mathfrak p} \ar[u]
&
R_{\mathfrak p} \ar[u] \ar[l]
&
R \ar[u] \ar[r] \ar[l]
&
R/\mathfrak p \ar[u] \ar[r]
&
\kappa(\mathfrak p) \ar[u]
}
$$
Dans ce diagramme, les colonnes extrêmes gauche et droite sont identiques.
Sur les spectres, les morphismes horizontaux induisent des homéomorphismes
sur leurs images, et les carrés induisent des carrés cartésiens d'espaces
topologiques (voir les lemmes \ref{algebra-lemma-spec-localization}
et \ref{algebra-lemma-spec-closed}).
Cela montre que $\mathfrak p$ appartient à l'image de l'application
entre spectres si et seulement si
$S \otimes_R \kappa(\mathfrak p)$ n'est pas l'anneau nul.
S'il existe un idéal premier $\mathfrak q \subset S$ au-dessus de
$\mathfrak p$, c'est-à-dire tel que
$\mathfrak p = \varphi^{-1}(\mathfrak q)$,
alors nous pouvons compléter le diagramme en le diagramme suivant :
$$
\xymatrix{
\kappa(\mathfrak q) = S_{\mathfrak q}/{\mathfrak q}S_{\mathfrak q}
&
S_{\mathfrak q} \ar[l]
&
S \ar[r] \ar[l]
&
S/\mathfrak q \ar[r]
&
\kappa(\mathfrak q)
\\
\kappa(\mathfrak p) \otimes_R S =
S_{\mathfrak p}/{\mathfrak p}S_{\mathfrak p} \ar[u]
&
S_{\mathfrak p} \ar[u] \ar[l]
&
S \ar[u] \ar[r] \ar[l]
&
S/\mathfrak pS \ar[u] \ar[r]
&
(R \setminus \mathfrak p)^{-1}S/\mathfrak pS \ar[u]
\\
\kappa(\mathfrak p) =
R_{\mathfrak p}/{\mathfrak p}R_{\mathfrak p} \ar[u]
&
R_{\mathfrak p} \ar[u] \ar[l]
&
R \ar[u] \ar[r] \ar[l]
&
R/\mathfrak p \ar[u] \ar[r]
&
\kappa(\mathfrak p) \ar[u]
}
$$
Dans ce diagramme encore, les colonnes extrêmes gauche et droite sont
identiques et, sur les spectres, les morphismes horizontaux induisent
des homéomorphismes sur leurs images.
\end{remark}

\begin{lemma}
\label{algebra-lemma-in-image}
Soit $\varphi : R \to S$ un morphisme d'anneaux. Soit $\mathfrak p$
un idéal premier de $R$. Les conditions suivantes sont équivalentes :
\begin{enumerate}
\item $\mathfrak p$ appartient à l'image de
$\Spec(S) \to \Spec(R)$,
\item $S \otimes_R \kappa(\mathfrak p) \not = 0$,
\item $S_{\mathfrak p}/\mathfrak p S_{\mathfrak p} \not = 0$,
\item $(S/\mathfrak pS)_{\mathfrak p} \not = 0$, et
\item $\mathfrak p = \varphi^{-1}(\mathfrak pS)$.
\end{enumerate}
\end{lemma}

\begin{proof}
Nous avons déjà vu l'équivalence des deux premières conditions
dans la remarque \ref{algebra-remark-fundamental-diagram}. Les autres
n'en sont que des reformulations.
\end{proof}

\begin{remark}
\label{algebra-remark-local-ring-fibre}
Soit $R \to S$ un morphisme d'anneaux. Soit $\mathfrak q$ un idéal premier
de $S$ au-dessus de l'idéal premier $\mathfrak p$ de $R$.
D'après la remarque \ref{algebra-remark-fundamental-diagram},
l'idéal premier $\mathfrak q$ correspond à un unique idéal premier
$\overline{\mathfrak q}$ de l'anneau de la fibre
$F = S \otimes_R \kappa(\mathfrak p)$.
Alors nous avons
$$
F_{\overline{\mathfrak q}}
\cong S_\mathfrak q \otimes_{R_\mathfrak p} \kappa(\mathfrak p)
\cong S_\mathfrak q/\mathfrak p S_\mathfrak q
$$
En effet, il existe un morphisme d'anneaux évident
$F \to S_\mathfrak q \otimes_{R_\mathfrak p} \kappa(\mathfrak p)$
qui s'identifie aisément à $F \to F_{\overline{\mathfrak q}}$.
Le second isomorphisme résulte du fait que $\kappa(\mathfrak p)$
est le quotient de $R_\mathfrak p$ par $\mathfrak pR_\mathfrak p$.
\end{remark}

\section{Le radical de Jacobson d'un anneau}
\label{algebra-section-radical}

\noindent
Rappelons que le {\it radical de Jacobson} $\text{rad}(R)$ d'un anneau $R$
est l'intersection de tous les idéaux maximaux de $R$. Si $R$ est local,
alors $\text{rad}(R)$ est l'idéal maximal de $R$.

\begin{lemma}
\label{algebra-lemma-contained-in-radical}
Soit $R$ un anneau de radical de Jacobson $\text{rad}(R)$.
Soit $I \subset R$ un idéal. Les conditions suivantes sont
équivalentes :
\begin{enumerate}
\item $I \subset \text{rad}(R)$, et
\item tout élément de $1 + I$ est une unité de $R$.
\end{enumerate}
Dans ce cas, tout élément de $R$ dont l'image est une unité de $R/I$
est une unité.
\end{lemma}

\begin{proof}
Si $f \in \text{rad}(R)$, alors $f \in \mathfrak m$ pour tout
idéal maximal $\mathfrak m$ de $R$. Ainsi $1 + f \not \in \mathfrak m$
pour tout idéal maximal $\mathfrak m$ de $R$. Le fermé
$V(1 + f)$ de $\Spec(R)$ est donc vide. Cela entraîne
que $1 + f$ est une unité ; voir le lemme \ref{algebra-lemma-Zariski-topology}.

\medskip\noindent
Réciproquement, supposons que $1 + f$ soit une unité pour tout $f \in I$.
Si $\mathfrak m$ est un idéal maximal et si $I \not \subset \mathfrak m$,
alors $I + \mathfrak m = R$. Par conséquent,
$1 = f + g$ pour un certain $g \in \mathfrak m$ et un certain
$f \in I$. Mais $g = 1 + (-f)$ n'est pas une unité, contradiction.

\medskip\noindent
Pour la dernière assertion, soit $f \in R$ dont l'image est une unité de $R/I$.
Nous pouvons trouver $g \in R$ dont l'image est l'inverse multiplicatif
de $f \bmod I$. Alors $fg = 1 \bmod I$. Ainsi $fg$ est une unité de $R$
d'après (2), ce qui entraîne que $f$ est une unité.
\end{proof}

\begin{lemma}
\label{algebra-lemma-surjective-on-spec-units}
Soit $\varphi : R \to S$ un morphisme d'anneaux tel que l'application induite
$\Spec(S) \to \Spec(R)$ soit surjective. Alors un élément $x \in R$
est une unité si et seulement si $\varphi(x) \in S$ est une unité.
\end{lemma}

\begin{proof}
Si $x$ est une unité, alors $\varphi(x)$ l'est aussi. Réciproquement, si
$\varphi(x)$ est une unité, alors
$\varphi(x) \not \in \mathfrak q$ pour tout
$\mathfrak q \in \Spec(S)$. Par conséquent,
$x \not \in \varphi^{-1}(\mathfrak q) = \Spec(\varphi)(\mathfrak q)$
pour tout $\mathfrak q \in \Spec(S)$. Puisque $\Spec(\varphi)$ est surjective,
nous concluons que $x$ est une unité d'après
la partie (17) du lemme \ref{algebra-lemma-Zariski-topology}.
\end{proof}












\section{Le lemme de Nakayama}
\label{algebra-section-nakayama}

\noindent
Nous citons \cite{MatCA} : « Ce lemme simple mais
important est dû à T.~Nakayama, G.~Azumaya et W.~Krull. La question de la priorité
est obscure et, bien qu'il soit habituellement appelé lemme de Nakayama, feu
le professeur Nakayama n'aimait pas ce nom. »

\begin{lemma}[Lemme de Nakayama]
\label{algebra-lemma-NAK}
\begin{reference}
\cite[1.M Lemme (NAK), page 11]{MatCA}
\end{reference}
\begin{history}
Nous citons \cite{MatCA} : « Ce lemme simple mais
important est dû à T.~Nakayama, G.~Azumaya et W.~Krull. La question de la priorité
est obscure et, bien qu'il soit habituellement appelé lemme de Nakayama, feu
le professeur Nakayama n'aimait pas ce nom. »
\end{history}
Soit $R$ un anneau de radical de Jacobson $\text{rad}(R)$.
Soit $M$ un $R$-module. Soit $I \subset R$
un idéal.
\begin{enumerate}
\item
\label{algebra-item-nakayama}
Si $IM = M$ et si $M$ est de type fini, alors il existe $f \in 1 + I$ tel que
$fM = 0$.
\item Si $IM = M$, si $M$ est de type fini et si $I \subset \text{rad}(R)$, alors $M = 0$.
\item Si $N, N' \subset M$, si $M = N + IN'$ et si $N'$ est de type fini,
alors il existe $f \in 1 + I$ tel que $fM \subset N$ et $M_f = N_f$.
\item Si $N, N' \subset M$, si $M = N + IN'$, si $N'$ est de type fini et si
$I \subset \text{rad}(R)$, alors $M = N$.
\item Si $N \to M$ est un morphisme de modules, si $N/IN \to M/IM$ est
surjectif et si $M$ est de type fini, alors il existe $f \in 1 + I$
tel que $N_f \to M_f$ soit surjectif.
\item Si $N \to M$ est un morphisme de modules, si $N/IN \to M/IM$ est
surjectif, si $M$ est de type fini et si $I \subset \text{rad}(R)$,
alors $N \to M$ est surjectif.
\item Si $x_1, \ldots, x_n \in M$ engendrent $M/IM$ et si $M$ est de type fini,
alors il existe $f \in 1 + I$ tel que $x_1, \ldots, x_n$
engendrent $M_f$ comme $R_f$-module.
\item Si $x_1, \ldots, x_n \in M$ engendrent $M/IM$, si $M$ est de type fini et si
$I \subset \text{rad}(R)$, alors $M$ est engendré par $x_1, \ldots, x_n$.
\item Si $IM = M$ et si $I$ est nilpotent, alors $M = 0$.
\item Si $N, N' \subset M$, si $M = N + IN'$ et si $I$ est nilpotent, alors $M = N$.
\item Si $N \to M$ est un morphisme de modules, si $I$ est nilpotent et si $N/IN \to M/IM$
est surjectif, alors $N \to M$ est surjectif.
\item Si $\{x_\alpha\}_{\alpha \in A}$ est un ensemble d'éléments de $M$
qui engendrent $M/IM$ et si $I$ est nilpotent, alors $M$ est engendré
par les $x_\alpha$.
\end{enumerate}
\end{lemma}

\begin{proof}
Démonstration de (\ref{algebra-item-nakayama}). Choisissons des générateurs $y_1, \ldots, y_m$ de $M$
comme $R$-module. Pour chaque $i$, nous pouvons écrire $y_i = \sum z_{ij} y_j$ avec
$z_{ij} \in I$ (puisque $M = IM$).
Autrement dit, $\sum_j (\delta_{ij} - z_{ij})y_j = 0$.
Soit $f$ le déterminant de la matrice $m \times m$
$A = (\delta_{ij} - z_{ij})$. Notons que $f \in 1 + I$
(puisque la matrice $A$ est congrue terme à terme à la
matrice identité $m \times m$ modulo $I$).
D'après le lemme \ref{algebra-lemma-matrix-left-inverse} (1),
il existe une matrice $m \times m$
$B$ telle que $BA = f 1_{m \times m}$. En développant, nous voyons que
$\sum_{i} b_{hi} a_{ij} = f \delta_{hj}$ pour tous
$h$ et $j$ ; ainsi, $\sum_{i, j} b_{hi} a_{ij} y_j
= \sum_{j} f \delta_{hj} y_j = f y_h$ pour tout $h$.
Autrement dit, $0 = f y_h$ pour tout $h$ (puisque chaque
$i$ vérifie $\sum_j a_{ij} y_j = 0$).
Cela entraîne que $f$ annule $M$.

\medskip\noindent
D'après le lemme \ref{algebra-lemma-contained-in-radical}, tout élément de $1 + \text{rad}(R)$ est
un élément inversible de $R$. Nous voyons donc que (\ref{algebra-item-nakayama}) entraîne
(2). Nous obtenons (3) en appliquant (1) à $M/N$, qui est de type fini puisque $N'$ l'est.
Nous obtenons (4) en appliquant (2) à $M/N$, qui est de type fini puisque $N'$ l'est.
Nous obtenons (5) en appliquant (3) à $M$ et aux sous-modules $\Im(N \to M)$
et $M$. Nous obtenons (6) en appliquant (4) à $M$ et aux sous-modules
$\Im(N \to M)$ et $M$.
Nous obtenons (7) en appliquant (5) au morphisme $R^{\oplus n} \to M$,
$(a_1, \ldots, a_n) \mapsto a_1x_1 + \ldots + a_nx_n$.
Nous obtenons (8) en appliquant (6) au morphisme $R^{\oplus n} \to M$,
$(a_1, \ldots, a_n) \mapsto a_1x_1 + \ldots + a_nx_n$.

\medskip\noindent
L'assertion (9) est vraie car, si $M = IM$, alors $M = I^nM$ pour tout $n \geq 0$,
et la nilpotence de $I$ signifie que $I^n = 0$ pour un certain $n \gg 0$. Les assertions
(10), (11) et (12) résultent de (9) par les arguments utilisés ci-dessus.
\end{proof}

\begin{lemma}
\label{algebra-lemma-NAK-localization}
Soit $R$ un anneau, soit $S \subset R$ une partie multiplicative,
soit $I \subset R$ un idéal et soit $M$ un $R$-module de type fini.
Si $x_1, \ldots, x_r \in M$ engendrent $S^{-1}(M/IM)$
comme $S^{-1}(R/I)$-module, alors il existe $f \in S + I$
tel que $x_1, \ldots, x_r$ engendrent $M_f$ comme
$R_f$-module.\footnote{Cas particuliers : (I) $I = 0$. Le lemme affirme que,
si $x_1, \ldots, x_r$ engendrent $S^{-1}M$, alors $x_1, \ldots, x_r$
engendrent $M_f$ pour un certain $f \in S$. (II) $I = \mathfrak p$ est
un idéal premier et $S = R \setminus \mathfrak p$. Le lemme affirme que, si
$x_1, \ldots, x_r$ engendrent $M \otimes_R \kappa(\mathfrak p)$,
alors $x_1, \ldots, x_r$ engendrent $M_f$ pour un certain
$f \in R$, $f \not \in \mathfrak p$.}
\end{lemma}

\begin{proof}
Cas particulier $I = 0$. Soient $y_1, \ldots, y_s$ des générateurs de $M$ comme $R$-module.
Puisque $S^{-1}M$ est engendré par $x_1, \ldots, x_r$, pour chaque $i$,
nous pouvons écrire $y_i = \sum (a_{ij}/s_{ij})x_j$ dans $S^{-1}M$
pour certains $a_{ij} \in R$ et $s_{ij} \in S$. En multipliant par le produit
$s \in S$ des $s_{ij}$, nous voyons que $sy_i = \sum a'_{ij}x_j$
dans $S^{-1}M$ pour certains $a'_{ij} \in R$.
Cela signifie à son tour qu'il existe des $t_i \in S$ tels que
$t_isy_i = \sum t_ia'_{ij}x_j$ dans $M$. Ainsi, si $t \in S$
est le produit des $t_i$, alors nous
voyons que $y_i$ appartient au sous-$R_{st}$-module engendré par $x_1, \ldots, x_r$
dans $M_{st}$. Par conséquent, $x_1, \ldots, x_r$ engendrent $M_{st}$.

\medskip\noindent
Cas général. D'après le cas particulier, nous pouvons trouver $s \in S$
tel que $x_1, \ldots, x_r$ engendrent $(M/IM)_s$ comme $(R/I)_s$-module.
D'après le lemme \ref{algebra-lemma-NAK}, nous pouvons trouver $g \in 1 + I_s \subset R_s$
tel que $x_1, \ldots, x_r$ engendrent $(M_s)_g$ comme $(R_s)_g$-module.
Écrivons $g = 1 + i/s'$. Alors $f = ss' + is$ convient ; nous omettons les détails.
\end{proof}

\begin{lemma}
\label{algebra-lemma-when-surjective-local}
Soit $A \to B$ un morphisme local d'anneaux locaux.
Supposons que
\begin{enumerate}
\item $B$ soit de type fini comme $A$-module,
\item $\mathfrak m_B$ soit un idéal de type fini,
\item $A \to B$ induise un isomorphisme sur les corps résiduels, et
\item $\mathfrak m_A/\mathfrak m_A^2 \to \mathfrak m_B/\mathfrak m_B^2$
soit surjectif.
\end{enumerate}
Alors $A \to B$ est surjectif.
\end{lemma}

\begin{proof}
Pour montrer que $A \to B$ est surjectif, considérons-le comme un morphisme de $A$-modules
et appliquons le lemme \ref{algebra-lemma-NAK} (6). Nous en concluons qu'il suffit
de montrer que $A/\mathfrak m_A \to B/\mathfrak m_AB$ est surjectif.
Comme $A/\mathfrak m_A = B/\mathfrak m_B$, il suffit de montrer que
$\mathfrak m_AB \to \mathfrak m_B$ est surjectif. Considérons
$\mathfrak m_AB \to \mathfrak m_B$ comme un morphisme de $B$-modules et appliquons
le lemme \ref{algebra-lemma-NAK} (6). Nous en concluons qu'il suffit de voir que
$\mathfrak m_AB/\mathfrak m_A\mathfrak m_B \to \mathfrak m_B/\mathfrak m_B^2$
est surjectif. Cela résulte de l'hypothèse (4).
\end{proof}







\section{Sous-ensembles ouverts et fermés des spectres}
\label{algebra-section-open-and-closed}

\noindent
Il se trouve que les sous-ensembles à la fois ouverts et fermés d'un spectre
correspondent aux idempotents de l'anneau.

\begin{lemma}
\label{algebra-lemma-idempotent-spec}
Soit $R$ un anneau. Soit $e \in R$ un idempotent.
Dans ce cas,
$$
\Spec(R) = D(e) \amalg D(1-e).
$$
\end{lemma}

\begin{proof}
Notons qu'un idempotent $e$ d'un anneau intègre est égal à $1$ ou à $0$.
Nous voyons donc que
\begin{eqnarray*}
D(e)
& = &
\{ \mathfrak p \in \Spec(R)
\mid
e \not\in \mathfrak p \} \\
& = &
\{ \mathfrak p \in \Spec(R)
\mid
e \not = 0\text{ dans }\kappa(\mathfrak p) \} \\
& = &
\{ \mathfrak p \in \Spec(R)
\mid
e = 1\text{ dans }\kappa(\mathfrak p) \}
\end{eqnarray*}
De même, nous avons
\begin{eqnarray*}
D(1-e)
& = &
\{ \mathfrak p \in \Spec(R)
\mid
1 - e \not\in \mathfrak p \} \\
& = &
\{ \mathfrak p \in \Spec(R)
\mid
e \not = 1\text{ dans }\kappa(\mathfrak p) \} \\
& = &
\{ \mathfrak p \in \Spec(R)
\mid
e = 0\text{ dans }\kappa(\mathfrak p) \}
\end{eqnarray*}
Puisque l'image de $e$ dans tout corps résiduel est égale à $1$ ou à $0$,
nous en déduisons que $D(e)$ et $D(1-e)$ recouvrent tout $\Spec(R)$.
\end{proof}

\begin{lemma}
\label{algebra-lemma-spec-product}
Soient $R_1$ et $R_2$ des anneaux.
Soit $R = R_1 \times R_2$.
Les morphismes $R \to R_1$, $(x, y) \mapsto x$ et $R \to R_2$,
$(x, y) \mapsto y$
induisent des applications continues $\Spec(R_1) \to \Spec(R)$ et
$\Spec(R_2) \to \Spec(R)$.
L'application ainsi obtenue
$$
\Spec(R_1) \amalg \Spec(R_2)
\longrightarrow
\Spec(R)
$$
est un homéomorphisme. Autrement dit,
le spectre de $R = R_1\times R_2$ est la
réunion disjointe du spectre de $R_1$ et du
spectre de $R_2$.
\end{lemma}

\begin{proof}
Écrivons $1 = e_1 + e_2$ avec $e_1 = (1, 0)$ et $e_2 = (0, 1)$.
Notons que $e_1$ et $e_2 = 1 - e_1$ sont des idempotents.
Nous laissons au lecteur le soin de montrer que
$R_1 = R_{e_1}$ est la localisation de $R$ en $e_1$.
Il en va de même pour $e_2$.
Ainsi, l'assertion du lemme résulte du lemme
\ref{algebra-lemma-idempotent-spec} combiné au lemme
\ref{algebra-lemma-standard-open}.
\end{proof}

\noindent
Nous redémontrerons le lemme suivant plus loin, après avoir introduit
un lemme de recollement pour les fonctions. Voir la section
\ref{algebra-section-tilde-module-sheaf}.

\begin{lemma}
\label{algebra-lemma-disjoint-decomposition}
Soit $R$ un anneau. Pour tout $U \subset \Spec(R)$
qui est à la fois ouvert et fermé,
il existe un unique idempotent $e \in R$ tel que
$U = D(e)$. Cela induit une correspondance bijective entre les
sous-ensembles à la fois ouverts et fermés $U \subset \Spec(R)$ et les
idempotents $e \in R$.
\end{lemma}

\begin{proof}
Soit $U \subset \Spec(R)$ un sous-ensemble ouvert et fermé.
Puisque $U$ est fermé, il est quasi-compact d'après le
lemme \ref{algebra-lemma-quasi-compact}, et il en va de même pour
son complémentaire.
Écrivons $U = \bigcup_{i = 1}^n D(f_i)$ comme réunion finie d'ouverts principaux.
De même, écrivons $\Spec(R) \setminus U = \bigcup_{j = 1}^m D(g_j)$
comme réunion finie d'ouverts principaux. Puisque $\emptyset =
D(f_i) \cap D(g_j) = D(f_i g_j)$, nous voyons que $f_i g_j$ est
nilpotent d'après le lemme \ref{algebra-lemma-Zariski-topology}.
Soit $I = (f_1, \ldots, f_n) \subset R$ et soit
$J = (g_1, \ldots, g_m) \subset R$.
Notons que $V(J)$ est égal à $U$, que $V(I)$
est égal au complémentaire de $U$, de sorte que $\Spec(R) = V(I) \amalg V(J)$.
D'après la remarque ci-dessus concernant la nilpotence,
nous voyons que $(IJ)^N = (0)$ pour un entier $N$ suffisamment grand.
Puisque $\bigcup D(f_i) \cup \bigcup D(g_j) = \Spec(R)$,
nous voyons que $I + J = R$ ; voir le lemme \ref{algebra-lemma-Zariski-topology}.
En élevant cette égalité à la puissance $2N$, nous en concluons que
$I^N + J^N = R$. Écrivons $1 = x + y$ avec $x \in I^N$ et $y \in J^N$.
Alors $0 = xy = x(1 - x)$ puisque $I^N J^N = (0)$. Ainsi $x = x^2$
est idempotent et appartient
à $I^N \subset I$. L'idempotent $y = 1 - x$ appartient à $J^N \subset J$.
Cela montre que l'idempotent $x$ a pour image $1$ dans tout corps résiduel
$\kappa(\mathfrak p)$ pour $\mathfrak p \in V(J)$ et que $x$ a pour image $0$
dans $\kappa(\mathfrak p)$ pour tout $\mathfrak p \in V(I)$.

\medskip\noindent
Pour voir l'unicité, supposons que $e_1, e_2$ soient des
idempotents distincts de $R$. Nous devons montrer qu'il
existe un idéal premier $\mathfrak p$ tel que $e_1 \in \mathfrak p$
et $e_2 \not \in \mathfrak p$, ou l'inverse.
Écrivons $e_i' = 1 - e_i$. Si $e_1 \not = e_2$, alors
$0 \not = e_1 - e_2  = e_1(e_2 + e_2') - (e_1 + e_1')e_2
= e_1 e_2' - e_1' e_2$. Par conséquent, soit l'idempotent
$e_1 e_2' \not = 0$, soit $e_1' e_2 \not = 0$. Un idempotent non nul
n'est pas nilpotent ; nous trouvons donc un idéal premier
$\mathfrak p$ tel que soit $e_1e_2' \not \in \mathfrak p$,
soit $e_1'e_2 \not \in \mathfrak p$, d'après le lemme \ref{algebra-lemma-Zariski-topology}.
Il est facile de voir que cela fournit l'idéal premier recherché.
\end{proof}

\begin{lemma}
\label{algebra-lemma-characterize-spec-connected}
Soit $R$ un anneau non nul. Alors $\Spec(R)$ est
connexe si et seulement si $R$ ne possède aucun idempotent
non trivial.
\end{lemma}

\begin{proof}
Cela résulte immédiatement du lemme \ref{algebra-lemma-disjoint-decomposition}
et de la définition d'un espace topologique connexe.
\end{proof}

\begin{lemma}
\label{algebra-lemma-ideal-is-squared-union-connected}
Soit $I \subset R$ un idéal de type fini d'un anneau $R$
tel que $I = I^2$. Alors
\begin{enumerate}
\item il existe un idempotent $e \in R$ tel que $I = (e)$,
\item $R/I \cong R_{e'}$ pour l'idempotent $e' = 1 - e \in R$, et
\item $V(I)$ est ouvert et fermé dans $\Spec(R)$.
\end{enumerate}
\end{lemma}

\begin{proof}
D'après le lemme de Nakayama \ref{algebra-lemma-NAK}, il existe un élément
$f = 1 + i$, $i \in I$ tel que $fI = 0$. Alors $f^2 = f + fi = f$
est un idempotent. Considérons l'idempotent $e = 1 - f = -i \in I$.
Pour $j \in I$, nous avons $ej = j - fj = j$, donc $I = (e)$.
Cela démontre (1).

\medskip\noindent
Les parties (2) et (3) résultent de (1). En effet, nous avons
$V(I) = V(e) = \Spec(R) \setminus D(e)$, qui est ouvert et
fermé d'après le lemme \ref{algebra-lemma-idempotent-spec} ou le
lemme \ref{algebra-lemma-disjoint-decomposition}. Cela démontre (3).
Pour (2), observons que le morphisme $R \to R_{e'}$ est surjectif
puisque $x/(e')^n = x/e' = xe'/(e')^2 = xe'/e' = x/1$ dans $R_{e'}$.
Le noyau du morphisme $R \to R_{e'}$ est l'ensemble des éléments de
$R$ annulés par une puissance strictement positive de $e'$. Puisque $e'$ est
idempotent, c'est l'idéal des éléments annulés par $e'$,
qui est l'idéal $I = (e)$ puisque $e + e' = 1$ est une paire
d'idempotents orthogonaux. Cela démontre (2).
\end{proof}






\section{Composantes connexes des spectres}
\label{algebra-section-connected-components}

\noindent
Les composantes connexes des spectres ne sont pas aussi faciles à comprendre
qu'on pourrait le penser de prime abord. Cela vient de ce que nous sommes habitués à la topologie des
espaces localement connexes, alors que le spectre d'un anneau n'est en général
pas localement connexe.

\begin{lemma}
\label{algebra-lemma-closed-union-connected-components}
Soit $R$ un anneau. Soit $T \subset \Spec(R)$ un sous-ensemble du spectre.
Les conditions suivantes sont équivalentes :
\begin{enumerate}
\item $T$ est fermé et est une réunion de composantes connexes de
$\Spec(R)$,
\item $T$ est une intersection de sous-ensembles ouverts et fermés de
$\Spec(R)$, et
\item $T = V(I)$, où $I \subset R$ est un idéal engendré par des idempotents.
\end{enumerate}
De plus, l'idéal de (3), s'il existe, est unique.
\end{lemma}

\begin{proof}
D'après
le lemme \ref{algebra-lemma-topology-spec}
et
Topologie, lemme \ref{topology-lemma-closed-union-connected-components},
nous voyons que (1) et (2) sont équivalentes.
Supposons (2) et écrivons $T = \bigcap U_\alpha$ avec
$U_\alpha \subset \Spec(R)$ ouvert et fermé.
Alors $U_\alpha = D(e_\alpha)$ pour un certain idempotent $e_\alpha \in R$ d'après le
lemme \ref{algebra-lemma-disjoint-decomposition}.
En posant alors $I = (1 - e_\alpha)$, nous voyons que $T = V(I)$, c'est-à-dire que (3) est satisfaite.
Enfin, supposons (3). Écrivons $T = V(I)$ et $I = (e_\alpha)$ pour une certaine
famille d'idempotents $e_\alpha$. Il est alors clair que
$T = \bigcap V(e_\alpha) = \bigcap D(1 - e_\alpha)$.

\medskip\noindent
Supposons que $I$ soit un idéal engendré par des idempotents.
Soit $e \in R$ un idempotent tel que $V(I) \subset V(e)$. Alors, d'après le
lemme \ref{algebra-lemma-Zariski-topology},
nous voyons que $e^n \in I$ pour un certain $n \geq 1$. Comme $e$ est un idempotent,
cela signifie que $e \in I$. Nous voyons donc que $I$ est engendré par
exactement les idempotents $e$ tels que $T \subset V(e)$.
Autrement dit, l'idéal $I$ est entièrement déterminé par le
sous-ensemble fermé $T$, ce qui démontre l'unicité.
\end{proof}

\begin{lemma}
\label{algebra-lemma-connected-component}
Soit $R$ un anneau.
Une composante connexe de
$\Spec(R)$ est de la forme $V(I)$,
où $I$ est un idéal engendré par des idempotents
tel que tout idempotent de $R$ ait pour image
$0$ ou $1$ dans $R/I$.
\end{lemma}

\begin{proof}
Soit $\mathfrak p$ un idéal premier de $R$. D'après le
lemme \ref{algebra-lemma-topology-spec},
nous voyons que les hypothèses de
Topologie, lemme \ref{topology-lemma-connected-component-intersection},
sont satisfaites pour l'espace topologique $\Spec(R)$.
Par conséquent, la composante connexe de $\mathfrak p$ dans $\Spec(R)$
est l'intersection des sous-ensembles ouverts et fermés de $\Spec(R)$
qui contiennent $\mathfrak p$. Elle est donc égale à $V(I)$, où
$I$ est engendré par les idempotents $e \in R$, tels que $e$ ait pour image $0$
dans $\kappa(\mathfrak p)$ ; voir le
lemme \ref{algebra-lemma-disjoint-decomposition}.
Tout idempotent $e$ qui n'appartient pas à cette famille a manifestement pour image $1$
dans $R/I$.
\end{proof}









\section{Propriétés de recollement}
\label{algebra-section-more-glueing}

\noindent
Dans cette section, nous rassemblons plusieurs résultats classiques de la
forme suivante : si une propriété est vraie sur tous les membres d'un
recouvrement par des ouverts principaux, alors elle est vraie. En fait, il suffit souvent de vérifier
les propriétés au niveau des anneaux locaux, comme dans le lemme suivant.

\begin{lemma}
\label{algebra-lemma-characterize-zero-local}
Soit $R$ un anneau.
\begin{enumerate}
\item Pour un élément $x$ d'un $R$-module $M$, les conditions suivantes sont équivalentes :
\begin{enumerate}
\item $x = 0$,
\item l'image de $x$ dans $M_\mathfrak p$ est nulle pour tout $\mathfrak p \in \Spec(R)$,
\item l'image de $x$ dans $M_{\mathfrak m}$ est nulle pour tout idéal maximal
$\mathfrak m$ de $R$.
\end{enumerate}
Autrement dit, le morphisme $M \to \prod_{\mathfrak m} M_{\mathfrak m}$
est injectif.
\item Pour un $R$-module $M$, les conditions suivantes sont équivalentes :
\begin{enumerate}
\item $M$ est nul,
\item $M_{\mathfrak p}$ est nul pour tout $\mathfrak p \in \Spec(R)$,
\item $M_{\mathfrak m}$ est nul pour tout idéal maximal $\mathfrak m$ de $R$.
\end{enumerate}
\item Pour un complexe $M_1 \to M_2 \to M_3$
de $R$-modules, les conditions suivantes sont équivalentes :
\begin{enumerate}
\item $M_1 \to M_2 \to M_3$ est exact,
\item pour tout idéal premier $\mathfrak p$ de $R$, le complexe localisé
$M_{1, \mathfrak p} \to M_{2, \mathfrak p} \to M_{3, \mathfrak p}$
est exact,
\item pour tout idéal maximal $\mathfrak m$ de $R$, le complexe localisé
$M_{1, \mathfrak m} \to M_{2, \mathfrak m} \to M_{3, \mathfrak m}$
est exact.
\end{enumerate}
\item Pour un morphisme $f : M \to M'$ de $R$-modules, les conditions suivantes sont équivalentes :
\begin{enumerate}
\item $f$ est injectif,
\item $f_{\mathfrak p} : M_\mathfrak p \to M'_\mathfrak p$ est injectif
pour tout idéal premier $\mathfrak p$ de $R$,
\item $f_{\mathfrak m} : M_\mathfrak m \to M'_\mathfrak m$ est injectif
pour tout idéal maximal $\mathfrak m$ de $R$.
\end{enumerate}
\item Pour un morphisme $f : M \to M'$ de $R$-modules, les conditions suivantes sont équivalentes :
\begin{enumerate}
\item $f$ est surjectif,
\item $f_{\mathfrak p} : M_\mathfrak p \to M'_\mathfrak p$ est surjectif
pour tout idéal premier $\mathfrak p$ de $R$,
\item $f_{\mathfrak m} : M_\mathfrak m \to M'_\mathfrak m$ est surjectif
pour tout idéal maximal $\mathfrak m$ de $R$.
\end{enumerate}
\item Pour un morphisme $f : M \to M'$ de $R$-modules, les conditions suivantes sont équivalentes :
\begin{enumerate}
\item $f$ est bijectif,
\item $f_{\mathfrak p} : M_\mathfrak p \to M'_\mathfrak p$ est bijectif
pour tout idéal premier $\mathfrak p$ de $R$,
\item $f_{\mathfrak m} : M_\mathfrak m \to M'_\mathfrak m$ est bijectif
pour tout idéal maximal $\mathfrak m$ de $R$.
\end{enumerate}
\end{enumerate}
\end{lemma}

\begin{proof}
Soit $x \in M$ comme dans (1). Soit $I = \{f \in R \mid fx = 0\}$.
Il est facile de voir que $I$ est un idéal (c'est
l'annulateur de $x$). La condition (1)(c) signifie que, pour
tout idéal maximal $\mathfrak m$, il existe
$f \in R \setminus \mathfrak m$ tel que $fx =0$.
Autrement dit, $V(I)$ ne contient aucun point fermé.
D'après le lemme \ref{algebra-lemma-Zariski-topology}, nous voyons que $I$ est l'idéal unité.
Ainsi $x$ est nul, c'est-à-dire que (1)(a) est satisfaite. Cela démontre (1).

\medskip\noindent
La partie (2) résulte de l'application simultanée de (1) à tous les éléments de $M$.

\medskip\noindent
Démonstration de (3). Soit $H$ l'homologie de la suite, c'est-à-dire
$H = \Ker(M_2 \to M_3)/\Im(M_1 \to M_2)$. D'après la
proposition \ref{algebra-proposition-localization-exact},
$H_\mathfrak p$ est l'homologie de la suite
$M_{1, \mathfrak p} \to M_{2, \mathfrak p} \to M_{3, \mathfrak p}$.
Ainsi (3) résulte de (2).

\medskip\noindent
Les parties (4) et (5) sont des cas particuliers de (3). La partie (6) résulte
formellement de la combinaison de (4) et (5).
\end{proof}

\begin{lemma}
\label{algebra-lemma-cover}
\begin{slogan}
Propriétés locales pour Zariski des modules et des algèbres
\end{slogan}
Soit $R$ un anneau. Soit $M$ un $R$-module. Soit $S$ une $R$-algèbre.
Supposons que $f_1, \ldots, f_n$ soit une liste finie
d'éléments de $R$ telle que $\bigcup D(f_i) = \Spec(R)$,
autrement dit $(f_1, \ldots, f_n) = R$.
\begin{enumerate}
\item Si chaque $M_{f_i} = 0$, alors $M = 0$.
\item Si chaque $M_{f_i}$ est un $R_{f_i}$-module de type fini,
alors $M$ est un $R$-module de type fini.
\item Si chaque $M_{f_i}$ est un $R_{f_i}$-module de présentation finie,
alors $M$ est un $R$-module de présentation finie.
\item Soit $M \to N$ un morphisme de $R$-modules. Si $M_{f_i} \to N_{f_i}$
est un isomorphisme pour tout $i$, alors $M \to N$ est un isomorphisme.
\item Soit $0 \to M'' \to M \to M' \to 0$ un complexe de $R$-modules.
Si $0 \to M''_{f_i} \to M_{f_i} \to M'_{f_i} \to 0$ est exact pour tout $i$,
alors $0 \to M'' \to M \to M' \to 0$ est exact.
\item Si chaque $R_{f_i}$ est noethérien, alors $R$ est noethérien.
\item Si chaque $S_{f_i}$ est une $R_{f_i}$-algèbre de type fini, alors
$S$ est une $R$-algèbre de type fini.
\item Si chaque $S_{f_i}$ est de présentation finie sur $R_{f_i}$, alors
$S$ est une $R$-algèbre de présentation finie.
\end{enumerate}
\end{lemma}

\begin{proof}
Démontrons successivement chacune des parties.
\begin{enumerate}
\item D'après la proposition \ref{algebra-proposition-localize-twice},
cela entraîne que $M_\mathfrak p = 0$ pour tout $\mathfrak p \in \Spec(R)$ ;
nous concluons donc par le lemme \ref{algebra-lemma-characterize-zero-local}.
\item Pour chaque $i$, choisissons un ensemble fini de générateurs $X_i$ de $M_{f_i}$.
Nous pouvons supposer sans perte de généralité que les éléments de $X_i$
appartiennent à l'image du morphisme de localisation $M \rightarrow M_{f_i}$ ; prenons donc
un ensemble fini $Y_i$ d'antécédents des éléments de $X_i$ dans $M$. Soit $Y$
la réunion de ces ensembles. C'est encore un ensemble fini.
Considérons le morphisme $R$-linéaire évident $R^Y \rightarrow M$ qui envoie l'élément
de base $e_y$ sur $y$. D'après l'hypothèse, ce morphisme est surjectif après localisation
en tout idéal premier $\mathfrak p$ de $R$ ; il est donc surjectif d'après le
lemme \ref{algebra-lemma-characterize-zero-local},
et $M$ est de type fini.
\item D'après (2), nous avons une suite exacte courte
$$
0 \rightarrow K \rightarrow R^n \rightarrow M \rightarrow 0
$$
Puisque la localisation est un foncteur exact et que $M_{f_i}$ est de présentation
finie, nous voyons que $K_{f_i}$ est de type fini pour tout
$1 \leq i \leq n$ d'après le lemme \ref{algebra-lemma-extension}.
D'après (2), cela entraîne que $K$ est un $R$-module de type fini et donc que
$M$ est de présentation finie.
\item D'après la proposition \ref{algebra-proposition-localize-twice},
l'hypothèse entraîne que le morphisme induit
sur les localisés en tous les idéaux premiers est un isomorphisme ; nous concluons donc
par le lemme \ref{algebra-lemma-characterize-zero-local}.
\item D'après la proposition \ref{algebra-proposition-localize-twice}, l'hypothèse
entraîne que la
suite des localisés en tous les idéaux premiers est exacte courte ; nous
concluons donc par le lemme \ref{algebra-lemma-characterize-zero-local}.
\item Nous allons montrer que tout idéal de $R$ possède un ensemble fini de générateurs.
Pour cela, soit $I \subset R$ un idéal arbitraire. D'après la
proposition \ref{algebra-proposition-localization-exact},
chaque $I_{f_i} \subset R_{f_i}$ est un idéal. Ils sont tous
de type fini par hypothèse, et nous concluons donc par (2).
\item Pour chaque $i$, choisissons un ensemble fini de générateurs $X_i$ de $S_{f_i}$.
Nous pouvons supposer sans perte de généralité que les éléments de $X_i$
appartiennent à l'image du morphisme de localisation $S \rightarrow S_{f_i}$ ; prenons donc
un ensemble fini $Y_i$ d'antécédents des éléments de $X_i$ dans
$S$. Soit $Y$ la réunion de ces ensembles. C'est encore un ensemble fini.
Considérons le morphisme d'algèbres $R[X_y]_{y \in Y} \rightarrow S$
induit par $Y$. Puisqu'il s'agit d'un morphisme d'algèbres, l'image $T$
est un $R$-sous-module du $R$-module $S$ ; nous pouvons donc considérer le
module quotient $S/T$. D'après l'hypothèse, celui-ci est nul après localisation
en chacun des $f_i$ ; il est donc nul d'après (1), et $S$ est par conséquent une
$R$-algèbre de type fini.
\item D'après le point précédent, il existe un morphisme surjectif de $R$-algèbres
$R[X_1, \ldots, X_n] \rightarrow S$. Soit $K$ le noyau
de ce morphisme. C'est un idéal de $R[X_1, \ldots, X_n]$, de type fini
après chacune des localisations en $f_i$. Puisque les $f_i$ engendrent l'idéal unité
dans $R$, ils engendrent également l'idéal unité dans $R[X_1, \ldots, X_n]$ ; une
application de (2) achève donc la démonstration.
\end{enumerate}
\end{proof}

\begin{lemma}
\label{algebra-lemma-cover-upstairs}
Soit $R \to S$ un morphisme d'anneaux.
Supposons que $g_1, \ldots, g_n$ soit une liste finie
d'éléments de $S$ telle que $\bigcup D(g_i) = \Spec(S)$,
autrement dit $(g_1, \ldots, g_n) = S$.
\begin{enumerate}
\item Si chaque $S_{g_i}$ est de type fini sur $R$, alors $S$ est
de type fini sur $R$.
\item Si chaque $S_{g_i}$ est de présentation finie sur $R$,
alors $S$ est de présentation finie sur $R$.
\end{enumerate}
\end{lemma}

\begin{proof}
Choisissons $h_1, \ldots, h_n \in S$ tels que $\sum h_i g_i = 1$.

\medskip\noindent
Démonstration de (1). Pour chaque $i$, choisissons une liste finie d'éléments
$x_{i, j} \in S_{g_i}$, $j = 1, \ldots, m_i$,
qui engendrent $S_{g_i}$ comme $R$-algèbre.
Écrivons $x_{i, j} = y_{i, j}/g_i^{n_{i, j}}$ pour un certain $y_{i, j} \in S$ et
un certain $n_{i, j} \ge 0$. Considérons la $R$-sous-algèbre $S' \subset S$
engendrée par $g_1, \ldots, g_n$, $h_1, \ldots, h_n$ et
les $y_{i, j}$, $i = 1, \ldots, n$, $j = 1, \ldots, m_i$.
Puisque la localisation est exacte (proposition \ref{algebra-proposition-localization-exact}),
nous voyons que $S'_{g_i} \to S_{g_i}$ est injectif.
D'autre part, il est surjectif par notre choix des $y_{i, j}$.
Les éléments $g_1, \ldots, g_n$ engendrent l'idéal unité dans $S'$
puisque $h_1, \ldots, h_n \in S'$.
Ainsi, $S' \to S$, considéré comme un morphisme de $S'$-modules, est un isomorphisme
d'après le lemme \ref{algebra-lemma-cover}.

\medskip\noindent
Démonstration de (2). Nous savons déjà que $S$ est de type fini.
Écrivons $S = R[x_1, \ldots, x_m]/J$ pour un certain idéal $J$.
Pour chaque $i$, choisissons un relèvement $g'_i \in R[x_1, \ldots, x_m]$ de $g_i$
et choisissons un relèvement $h'_i \in R[x_1, \ldots, x_m]$ de $h_i$.
Nous voyons alors que
$$
S_{g_i} = R[x_1, \ldots, x_m, y_i]/(J_i + (1 - y_ig'_i))
$$
où $J_i$ est l'idéal de $R[x_1, \ldots, x_m, y_i]$
engendré par $J$. Nous omettons un petit détail. D'après le
lemme \ref{algebra-lemma-finite-presentation-independent},
nous pouvons choisir une liste finie d'éléments
$f_{i, j} \in J$, $j = 1, \ldots, m_i$,
telle que les images des $f_{i, j}$ dans $J_i$ ainsi que $1 - y_ig'_i$
engendrent l'idéal $J_i + (1 - y_ig'_i)$.
Posons
$$
S' = R[x_1, \ldots, x_m]/\left(\sum h'_ig'_i - 1, f_{i, j}; 
i = 1, \ldots, n, j = 1, \ldots, m_i\right)
$$
Il existe un morphisme surjectif de $R$-algèbres $S' \to S$.
Les classes des éléments $g'_1, \ldots, g'_n$ dans $S'$
engendrent l'idéal unité et, par construction, les morphismes
$S'_{g'_i} \to S_{g_i}$ sont injectifs.
Nous concluons donc comme dans la partie (1).
\end{proof}





\section{Recollement de fonctions}
\label{algebra-section-tilde-module-sheaf}

\noindent
Dans cette section, nous montrons qu'étant donné un recouvrement ouvert
$$
\Spec(R) = \bigcup\nolimits_{i = 1}^n D(f_i)
$$
par des ouverts principaux, et étant donné un élément $h_i \in R_{f_i}$
pour chaque $i$, tel que $h_i = h_j$ comme éléments de $R_{f_i f_j}$,
il existe un unique $h \in R$ tel que l'image de
$h$ dans $R_{f_i}$ soit $h_i$. Ce résultat peut s'interpréter
de deux manières :
\begin{enumerate}
\item La règle $D(f) \mapsto R_f$ est un faisceau d'anneaux
sur les ouverts principaux ; voir Faisceaux, section \ref{sheaves-section-bases}.
\item Si nous considérons les éléments de $R_f$ comme les fonctions « algébriques »
ou « régulières » sur $D(f)$, alors celles-ci se recollent
comme le feraient les fonctions continues, resp.\ différentiables,
sur une variété topologique, resp.\ différentiable.
\end{enumerate}

\begin{lemma}
\label{algebra-lemma-cover-module}
Soit $R$ un anneau. Soient $f_1, \ldots, f_n$ des éléments de $R$
qui engendrent l'idéal unité. Soit $M$ un $R$-module.
La suite
$$
0 \to
M \xrightarrow{\alpha}
\bigoplus\nolimits_{i = 1}^n M_{f_i} \xrightarrow{\beta}
\bigoplus\nolimits_{i, j = 1}^n M_{f_i f_j}
$$
est exacte, où $\alpha(m) = (m/1, \ldots, m/1)$
et $\beta(m_1/f_1^{e_1}, \ldots, m_n/f_n^{e_n})
= (m_i/f_i^{e_i} - m_j/f_j^{e_j})_{(i, j)}$.
\end{lemma}

\begin{proof}
Il suffit de montrer que la localisation de la suite en
tout idéal maximal $\mathfrak m$ est exacte ; voir le
lemme \ref{algebra-lemma-characterize-zero-local}.
Puisque $f_1, \ldots, f_n$ engendrent l'idéal unité,
il existe $i$ tel que $f_i \not \in \mathfrak m$.
Après renumérotation, nous pouvons supposer que $i = 1$.
Notons que $(M_{f_i})_\mathfrak m = (M_\mathfrak m)_{f_i}$
et que $(M_{f_if_j})_\mathfrak m = (M_\mathfrak m)_{f_if_j}$ ; voir la
proposition \ref{algebra-proposition-localize-twice-module}.
En particulier, $(M_{f_1})_\mathfrak m = M_\mathfrak m$ et
$(M_{f_1 f_i})_\mathfrak m = (M_\mathfrak m)_{f_i}$, car
$f_1$ est une unité.
Notons que les morphismes de la suite sont les morphismes canoniques
provenant du
lemme \ref{algebra-lemma-universal-property-localization-module}
et du morphisme identité de $M$.
Cela étant dit, après avoir remplacé $R$ par $R_\mathfrak m$,
$M$ par $M_\mathfrak m$, les $f_i$ par leurs images dans $R_\mathfrak m$
et $f_1$ par $1 \in R_\mathfrak m$,
nous nous ramenons au cas où $f_1 = 1$.

\medskip\noindent
Supposons $f_1 = 1$. L'injectivité de $\alpha$ est alors immédiate. Soit
$m = (m_i) \in \bigoplus_{i = 1}^n M_{f_i}$ un élément du noyau de $\beta$.
Alors $m_1 \in M_{f_1} = M$. De plus, $\beta(m) = 0$
entraîne que $m_1$ et $m_i$ ont la même image dans
$M_{f_1f_i} = M_{f_i}$. Ainsi $\alpha(m_1) = m$, ce qui
achève la démonstration.
\end{proof}

\begin{lemma}
\label{algebra-lemma-standard-covering}
Soit $R$ un anneau, et soient $f_1, f_2, \ldots f_n\in R$ des éléments qui engendrent
l'idéal unité de $R$.
Alors la suite suivante est exacte :
$$
0 \longrightarrow
R \longrightarrow
\bigoplus\nolimits_i R_{f_i} \longrightarrow
\bigoplus\nolimits_{i, j}R_{f_if_j}
$$
où les morphismes $\alpha : R \longrightarrow \bigoplus_i R_{f_i}$
et $\beta : \bigoplus_i R_{f_i} \longrightarrow \bigoplus_{i, j} R_{f_if_j}$
sont définis par
$$
\alpha(x) = \left(\frac{x}{1}, \ldots, \frac{x}{1}\right)
\text{ et }
\beta\left(\frac{x_1}{f_1^{r_1}}, \ldots, \frac{x_n}{f_n^{r_n}}\right)
=
\left(\frac{x_i}{f_i^{r_i}}-\frac{x_j}{f_j^{r_j}}~\text{dans}~R_{f_if_j}\right).
$$
\end{lemma}

\begin{proof}
C'est un cas particulier du lemme \ref{algebra-lemma-cover-module}.
\end{proof}

\noindent
Nous avons déjà vu le résultat suivant ci-dessus, mais nous l'énonçons explicitement ici
pour plus de commodité.

\begin{lemma}
\label{algebra-lemma-disjoint-implies-product}
Soit $R$ un anneau.
Si $\Spec(R) = U \amalg V$, où $U$ et $V$ sont tous deux ouverts,
alors $R \cong R_1 \times R_2$, avec $U \cong \Spec(R_1)$
et $V \cong \Spec(R_2)$ via les morphismes du lemme \ref{algebra-lemma-spec-product}.
De plus, $R_1$ et $R_2$ sont tous deux des localisations ainsi que des quotients
de l'anneau $R$.
\end{lemma}

\begin{proof}
D'après le lemme \ref{algebra-lemma-disjoint-decomposition}, nous avons
$U = D(e)$ et $V = D(1-e)$ pour un certain idempotent $e$.
D'après le lemme \ref{algebra-lemma-standard-covering}, nous voyons que
$R \cong R_e \times R_{1 - e}$ (puisque, manifestement, $R_{e(1-e)} = 0$,
de sorte que la condition de recollement est triviale ; bien entendu, il est
également immédiat de démontrer directement la décomposition en produit dans ce
cas). Le lemme en résulte.
\end{proof}

\begin{lemma}
\label{algebra-lemma-when-injective-covering}
Soit $R$ un anneau.
Soient $f_1, \ldots, f_n \in R$.
Soit $M$ un $R$-module.
Alors $M \to \bigoplus M_{f_i}$ est injectif si et seulement si
$$
M \longrightarrow \bigoplus\nolimits_{i = 1, \ldots, n} M, \quad
m \longmapsto (f_1m, \ldots, f_nm)
$$
est injectif.
\end{lemma}

\begin{proof}
Le morphisme $M \to \bigoplus M_{f_i}$ est injectif si et seulement si,
pour tous $m \in M$ et $e_1, \ldots, e_n \geq 1$ tels que
$f_i^{e_i}m = 0$, $i = 1, \ldots, n$, nous avons $m = 0$.
Cela entraîne manifestement que le morphisme affiché est injectif.
Réciproquement, supposons le morphisme affiché injectif et que
$m \in M$ et $e_1, \ldots, e_n \geq 1$ soient tels que
$f_i^{e_i}m = 0$, $i = 1, \ldots, n$. Si $e_i = 1$ pour tout $i$,
alors l'injectivité du morphisme affiché donne immédiatement $m = 0$.
Démontrons ensuite que cela vaut pour toutes données de ce type
par récurrence sur $e = \sum e_i$. Le cas initial est $e = n$, et nous venons
de le traiter. Si un certain $e_i > 1$, posons $m' = f_im$.
Par récurrence, nous voyons que $m' = 0$. Ainsi $f_i m = 0$,
c'est-à-dire que nous pouvons prendre $e_i = 1$, ce qui diminue $e$ et achève la récurrence.
\end{proof}

\noindent
Le lemme suivant s'énonce et se démontre mieux dans le cadre plus général
de la descente plate. Il est néanmoins naturel de l'énoncer ici,
car il s'accorde bien avec ce qui précède.

\begin{lemma}
\label{algebra-lemma-glue-modules}
Soit $R$ un anneau. Soient $f_1, \ldots, f_n \in R$. Supposons données
les données suivantes :
\begin{enumerate}
\item Pour chaque $i$, un $R_{f_i}$-module $M_i$.
\item Pour chaque paire $i, j$, un isomorphisme de $R_{f_if_j}$-modules
$\psi_{ij} : (M_i)_{f_j} \to (M_j)_{f_i}$.
\end{enumerate}
qui satisfont la « condition de cocycle » selon laquelle tous les diagrammes
$$
\xymatrix{
(M_i)_{f_jf_k}
\ar[rd]_{\psi_{ij}}
\ar[rr]^{\psi_{ik}}
& &
(M_k)_{f_if_j} \\
&
(M_j)_{f_if_k} \ar[ru]_{\psi_{jk}}
}
$$
sont commutatifs (pour tous les triplets $i, j, k$). À partir de ces données, définissons
$$
M = \Ker\left(
\bigoplus\nolimits_{1 \leq i \leq n} M_i
\longrightarrow
\bigoplus\nolimits_{1 \leq i, j \leq n} (M_i)_{f_j}
\right)
$$
où $(m_1, \ldots, m_n)$ est envoyé sur l'élément dont
la composante d'indice $(i, j)$ est $m_i/1 - \psi_{ji}(m_j/1)$.
Alors le morphisme naturel $M \to M_i$ induit un isomorphisme
$M_{f_i} \to M_i$. De plus, $\psi_{ij}(m/1) = m/1$
pour tout $m \in M$ (avec les notations évidentes).
\end{lemma}

\begin{proof}
Pour montrer que $M_{f_1} \to M_1$ est un isomorphisme, il suffit
de montrer que sa localisation en tout idéal premier $\mathfrak p'$
de $R_{f_1}$ est un isomorphisme ; voir le
lemme \ref{algebra-lemma-characterize-zero-local}.
Écrivons $\mathfrak p' = \mathfrak p R_{f_1}$
pour un certain idéal premier $\mathfrak p \subset R$, avec $f_1 \not \in \mathfrak p$ ; voir le
lemme \ref{algebra-lemma-standard-open}.
Puisque la localisation est exacte
(proposition \ref{algebra-proposition-localization-exact}),
nous voyons que
\begin{align*}
(M_{f_1})_{\mathfrak p'} & =
M_\mathfrak p \\
& =
\Ker\left(
\bigoplus\nolimits_{1 \leq i \leq n} M_{i, \mathfrak p}
\longrightarrow
\bigoplus\nolimits_{1 \leq i, j \leq n} ((M_i)_{f_j})_\mathfrak p
\right) \\
& =
\Ker\left(
\bigoplus\nolimits_{1 \leq i \leq n} M_{i, \mathfrak p}
\longrightarrow
\bigoplus\nolimits_{1 \leq i, j \leq n} (M_{i, \mathfrak p})_{f_j}
\right)
\end{align*}
Nous avons également utilisé ici la proposition \ref{algebra-proposition-localize-twice-module}.
Puisque $f_1$ est une unité de $R_\mathfrak p$, nous nous ramenons au cas
où $f_1 = 1$ en remplaçant $R$ par $R_\mathfrak p$, les $f_i$ par les
images des $f_i$ dans $R_\mathfrak p$, $M$ par $M_\mathfrak p$ et
$f_1$ par $1$.

\medskip\noindent
Supposons $f_1 = 1$. Alors $\psi_{1j} : (M_1)_{f_j} \to M_j$
est un isomorphisme pour $j = 2, \ldots, n$. Si nous utilisons ces
isomorphismes pour identifier $M_j = (M_1)_{f_j}$, nous voyons alors
que $\psi_{ij} : (M_1)_{f_if_j} \to (M_1)_{f_if_j}$ est
l'identification canonique. Ainsi, le complexe
$$
0 \to M_1 \to
\bigoplus\nolimits_{1 \leq i \leq n} (M_1)_{f_i}
\longrightarrow
\bigoplus\nolimits_{1 \leq i, j \leq n}
(M_1)_{f_if_j}
$$
est exact d'après le lemme \ref{algebra-lemma-cover-module}.
Ainsi, le premier morphisme identifie $M_1$ à $M$ dans ce cas,
et tout est clair.
\end{proof}












\section{Diviseurs de zéro et anneaux totaux des fractions}
\label{algebra-section-total-quotient-ring}

\noindent
L'anneau local en un idéal premier minimal possède les propriétés suivantes.

\begin{lemma}
\label{algebra-lemma-minimal-prime-reduced-ring}
Soit $\mathfrak p$ un idéal premier minimal d'un anneau $R$.
Tout élément de l'idéal maximal de $R_{\mathfrak p}$
est nilpotent. Si $R$ est réduit, alors $R_{\mathfrak p}$
est un corps.
\end{lemma}

\begin{proof}
Si un élément $x$ de ${\mathfrak p}R_{\mathfrak p}$
n'est pas nilpotent, alors $D(x) \not = \emptyset$ ; voir le
lemme \ref{algebra-lemma-Zariski-topology}. Cela contredit
la minimalité de $\mathfrak p$. Si $R$ est réduit,
alors ${\mathfrak p}R_{\mathfrak p} = 0$ et
l'anneau local est donc un corps.
\end{proof}

\begin{lemma}
\label{algebra-lemma-reduced-ring-sub-product-fields}
Soit $R$ un anneau réduit. Alors
\begin{enumerate}
\item $R$ est un sous-anneau d'un produit de corps,
\item $R \to \prod_{\mathfrak p\text{ minimal}} R_{\mathfrak p}$
est un plongement dans un produit de corps,
\item $\bigcup_{\mathfrak p\text{ minimal}} \mathfrak p$ est l'ensemble
des diviseurs de zéro de $R$.
\end{enumerate}
\end{lemma}

\begin{proof}
D'après le lemme \ref{algebra-lemma-minimal-prime-reduced-ring}, chacun des anneaux
$R_\mathfrak p$ est un corps. En particulier, le noyau du morphisme d'anneaux
$R \to R_\mathfrak p$ est $\mathfrak p$.
D'après le lemme \ref{algebra-lemma-Zariski-topology},
nous avons $\bigcap_{\mathfrak p} \mathfrak p = (0)$.
Ainsi (2) et (1) sont vraies. Si $x y = 0$ et $y \not = 0$, alors
$y \not \in \mathfrak p$ pour un certain idéal premier minimal $\mathfrak p$.
Par conséquent, $x \in \mathfrak p$. Ainsi, tout diviseur de zéro de $R$ appartient
à $\bigcup_{\mathfrak p\text{ minimal}} \mathfrak p$.
Réciproquement, supposons que $x \in \mathfrak p$ pour un certain idéal
premier minimal $\mathfrak p$. Alors l'image de $x$ dans $R_\mathfrak p$ est nulle ;
il existe donc $y \in R$, $y \not \in \mathfrak p$ tel que
$xy = 0$. Autrement dit, $x$ est un diviseur de zéro. Cela achève la
démonstration de (3) et du lemme.
\end{proof}

\noindent
L'anneau total des fractions $Q(R)$ d'un anneau $R$ a été introduit dans
l'exemple \ref{algebra-example-localize-at-prime}.

\begin{lemma}
\label{algebra-lemma-total-ring-fractions}
Soit $R$ un anneau.
Soit $S \subset R$ une partie multiplicative formée d'éléments non diviseurs de zéro.
Alors $Q(R) \cong Q(S^{-1}R)$.
En particulier, $Q(R) \cong Q(Q(R))$.
\end{lemma}

\begin{proof}
Si $x \in S^{-1}R$ est un élément non diviseur de zéro et si
$x = r/f$ pour certains $r \in R$, $f \in S$, alors
$r$ est un élément non diviseur de zéro de $R$. Le lemme en résulte.
\end{proof}

\noindent
Nous pouvons appliquer les résultats de recollement pour démontrer une propriété des
anneaux totaux des fractions $Q(R)$, que nous avons introduits dans
l'exemple \ref{algebra-example-localize-at-prime}.

\begin{lemma}
\label{algebra-lemma-total-ring-fractions-no-embedded-points}
Soit $R$ un anneau.
Supposons que $R$ possède un nombre fini d'idéaux premiers minimaux
$\mathfrak q_1, \ldots, \mathfrak q_t$, et que
$\mathfrak q_1 \cup \ldots \cup \mathfrak q_t$ soit l'ensemble
des diviseurs de zéro de $R$.
Alors l'anneau total des fractions $Q(R)$ est égal à
$R_{\mathfrak q_1} \times \ldots \times R_{\mathfrak q_t}$.
\end{lemma}

\begin{proof}
Il existe des morphismes naturels $Q(R) \to R_{\mathfrak q_i}$ puisque
tout élément non diviseur de zéro appartient à $R \setminus \mathfrak q_i$.
Il existe donc un morphisme naturel
$Q(R) \to R_{\mathfrak q_1} \times \ldots \times R_{\mathfrak q_t}$.
Pour tout idéal premier non minimal $\mathfrak p \subset R$, nous voyons que
$\mathfrak p \not \subset \mathfrak q_1 \cup \ldots \cup \mathfrak q_t$
d'après le lemme \ref{algebra-lemma-silly}. Par conséquent,
$\Spec(Q(R)) = \{\mathfrak q_1, \ldots, \mathfrak q_t\}$
(comme sous-ensembles de $\Spec(R)$ ; voir le lemme \ref{algebra-lemma-spec-localization}).
Ainsi $\Spec(Q(R))$ est un ensemble discret fini et
il s'ensuit que $Q(R) = A_1 \times \ldots \times A_t$,
avec $\Spec(A_i) = \{\mathfrak{q}_i\}$ ; voir le
lemme \ref{algebra-lemma-disjoint-implies-product}.
De plus, $A_i$ est un anneau local qui est une localisation
de $R$. Ainsi $A_i \cong R_{\mathfrak q_i}$.
\end{proof}

















\section{Composantes irréductibles des spectres}
\label{algebra-section-irreducible}

\noindent
Nous montrons que les composantes irréductibles du
spectre d'un anneau correspondent aux
idéaux premiers minimaux de l'anneau.

\begin{lemma}
\label{algebra-lemma-irreducible}
Soit $R$ un anneau.
\begin{enumerate}
\item Pour un idéal premier $\mathfrak p \subset R$, l'adhérence
de $\{\mathfrak p\}$ dans la topologie de Zariski est $V(\mathfrak p)$.
Sous forme de formule, $\overline{\{\mathfrak p\}} = V(\mathfrak p)$.
\item Les sous-ensembles fermés irréductibles de $\Spec(R)$ sont
exactement les sous-ensembles $V(\mathfrak p)$, où $\mathfrak p \subset R$
est un idéal premier.
\item Les composantes irréductibles (voir Topologie,
définition \ref{topology-definition-irreducible-components})
de $\Spec(R)$ sont exactement les sous-ensembles $V(\mathfrak p)$,
où $\mathfrak p \subset R$ est un idéal premier minimal.
\end{enumerate}
\end{lemma}

\begin{proof}
Notons que, si $ \mathfrak p \in V(I)$, alors
$I \subset \mathfrak p$. Par conséquent,
il est clair que $\overline{\{\mathfrak p\}} = V(\mathfrak p)$.
En particulier, $V(\mathfrak p)$ est l'adhérence d'un
singleton et est donc irréductible.
La deuxième assertion entraîne la troisième.
Pour démontrer la deuxième, soit
$V(I) \subset \Spec(R)$, où $I$ est un idéal radical.
Si $I$ n'est pas premier, choisissons $a, b\in R$, $a, b\not \in I$
avec $ab\in I$. Dans ce cas, $V(I, a) \cup V(I, b) = V(I)$,
mais ni $V(I, b) = V(I)$ ni $V(I, a) = V(I)$, d'après le
lemme \ref{algebra-lemma-Zariski-topology}. Ainsi $V(I)$ n'est pas
irréductible.
\end{proof}

\noindent
Autrement dit, ce lemme montre que tout sous-ensemble fermé irréductible
de $\Spec(R)$ est de la forme $V(\mathfrak p)$ pour
un certain idéal premier $\mathfrak p$. Puisque $V(\mathfrak p) = \overline{\{\mathfrak p\}}$,
nous voyons que tout sous-ensemble fermé irréductible possède un unique point générique ;
voir Topologie, définition \ref{topology-definition-generic-point}.
En particulier, $\Spec(R)$ est un espace topologique sobre.
Consignons ce fait dans le lemme suivant.

\begin{lemma}
\label{algebra-lemma-spec-spectral}
Le spectre d'un anneau est un espace spectral ; voir Topologie, définition
\ref{topology-definition-spectral-space}.
\end{lemma}

\begin{proof}
Formellement, cela résulte du lemme \ref{algebra-lemma-irreducible} et du
lemme \ref{algebra-lemma-topology-spec}. Voir aussi la discussion ci-dessus.
\end{proof}

\begin{lemma}
\label{algebra-lemma-irreducible-components-containing-x}
Soit $R$ un anneau. Soit $\mathfrak p \subset R$ un idéal premier.
\begin{enumerate}
\item L'ensemble des sous-ensembles fermés irréductibles de $\Spec(R)$
passant par $\mathfrak p$ est en correspondance bijective avec les
idéaux premiers $\mathfrak q \subset R_{\mathfrak p}$.
\item L'ensemble des composantes irréductibles de $\Spec(R)$ passant par
$\mathfrak p$ est en correspondance bijective avec les idéaux
premiers minimaux $\mathfrak q \subset R_{\mathfrak p}$.
\end{enumerate}
\end{lemma}

\begin{proof}
Cela résulte du lemme \ref{algebra-lemma-irreducible}
et de la description de $\Spec(R_\mathfrak p)$ donnée dans le
lemme \ref{algebra-lemma-spec-localization}, qui montre que
$\Spec(R_\mathfrak p)$ correspond aux idéaux premiers $\mathfrak q$ de $R$
tels que $\mathfrak q \subset \mathfrak p$.
\end{proof}

\begin{lemma}
\label{algebra-lemma-standard-open-containing-maximal-point}
Soit $R$ un anneau.
Soit $\mathfrak p$ un idéal premier minimal de $R$.
Soit $W \subset \Spec(R)$ un ouvert quasi-compact
qui ne contient pas le point $\mathfrak p$. Alors il
existe $f \in R$, $f \not \in \mathfrak p$, tel
que $D(f) \cap W = \emptyset$.
\end{lemma}

\begin{proof}
Puisque $W$ est quasi-compact, nous pouvons l'écrire comme réunion finie
d'ouverts affines principaux $D(g_i)$, $i = 1, \ldots, n$.
Puisque $\mathfrak p \not \in W$, nous avons $g_i \in \mathfrak p$ pour
tout $i$. D'après le lemme \ref{algebra-lemma-minimal-prime-reduced-ring},
chaque $g_i$ est nilpotent dans $R_{\mathfrak p}$. Nous pouvons donc trouver
$f \in R$, $f \not \in \mathfrak p$, tel que, pour tout $i$, nous ayons
$f g_i^{n_i} = 0$ pour un certain $n_i > 0$. Alors $D(f)$ convient.
\end{proof}

\begin{lemma}
\label{algebra-lemma-ring-with-only-minimal-primes}
Soit $R$ un anneau. Soit $X = \Spec(R)$ comme espace topologique.
Les conditions suivantes sont équivalentes :
\begin{enumerate}
\item $X$ est profini,
\item $X$ est séparé,
\item $X$ est totalement discontinu.
\item tout ouvert quasi-compact de $X$ est fermé,
\item il n'existe aucune inclusion non triviale entre ses idéaux premiers,
\item tout idéal premier est un idéal maximal,
\item tout idéal premier est minimal,
\item tout ouvert principal $D(f) \subset X$ est fermé, et
\item ajouter d'autres conditions ici.
\end{enumerate}
\end{lemma}

\begin{proof}
Première démonstration. Il est clair que (5), (6) et (7) sont équivalentes.
Il est clair que (4) et (8) sont équivalentes, puisque tout ouvert quasi-compact
est une réunion finie d'ouverts principaux.
L'implication (7) $\Rightarrow$ (4) résulte du
lemme \ref{algebra-lemma-standard-open-containing-maximal-point}.
Supposons (4). Soient $\mathfrak p, \mathfrak p'$ deux idéaux premiers
distincts de $R$. Choisissons $f \in \mathfrak p'$, $f \not \in \mathfrak p$
(au besoin, échangeons $\mathfrak p$ et $\mathfrak p'$).
Alors $\mathfrak p' \not \in D(f)$ et $\mathfrak p \in D(f)$.
D'après (4), l'ouvert $D(f)$ est également fermé.
Ainsi $\mathfrak p$ et $\mathfrak p'$ appartiennent à des voisinages ouverts
disjoints dont la réunion est $X$. Par conséquent, $X$ est séparé et totalement
discontinu. Ainsi (4) $\Rightarrow$ (2) et (3).
Si (3) est satisfaite, il ne peut exister aucune spécialisation
entre les points de $\Spec(R)$, et nous voyons que (5) est satisfaite.
Si $X$ est séparé, tout point est fermé, donc (2) entraîne (6).
Ainsi (2), (3), (4), (5), (6), (7) et (8) sont équivalentes.
Tout espace profini est séparé, donc (1) entraîne (2).
Si $X$ satisfait (2) et (3), alors $X$ (qui est quasi-compact d'après le
lemme \ref{algebra-lemma-quasi-compact}) est profini d'après
Topologie, lemme \ref{topology-lemma-profinite}.

\medskip\noindent
Deuxième démonstration. Outre l'équivalence de (4) et (8), cela résulte
du lemme \ref{algebra-lemma-spec-spectral} et de faits purement topologiques ; voir
Topologie, lemme \ref{topology-lemma-characterize-profinite-spectral}.
\end{proof}













\section{Exemples de spectres d'anneaux}
\label{algebra-section-examples-spectra}

\noindent
Dans cette section, nous donnons quelques exemples de spectres.

\begin{example}
\label{algebra-example-spec-Zxmodx2minus4}
Dans cet exemple, nous décrivons $X = \Spec(\mathbf{Z}[x]/(x^2 - 4))$.
Soit $\mathfrak{p}$ un idéal premier arbitraire de $X$.
Soit $\phi : \mathbf{Z} \to \mathbf{Z}[x]/(x^2 - 4)$ le morphisme d'anneaux naturel.
Alors $ \phi^{-1}(\mathfrak p)$ est un idéal premier de $\mathbf{Z}$.
Si $ \phi^{-1}(\mathfrak p) = (2)$, alors, puisque $\mathfrak p$ contient $2$,
il correspond à un idéal premier de
$\mathbf{Z}[x]/(x^2 - 4, 2) \cong (\mathbf{Z}/2\mathbf{Z})[x]/(x^2)$
via le morphisme $ \mathbf{Z}[x]/(x^2 - 4) \to  \mathbf{Z}[x]/(x^2 - 4, 2)$.
Tout idéal premier de $(\mathbf{Z}/2\mathbf{Z})[x]/(x^2)$ correspond à un idéal premier
de $(\mathbf{Z}/2\mathbf{Z})[x]$ contenant $(x^2)$. Un tel idéal premier
contient alors $x$. Puisque
$(\mathbf{Z}/2\mathbf{Z}) \cong (\mathbf{Z}/2\mathbf{Z})[x]/(x)$ est un corps,
$(x)$ est un idéal maximal. Puisque tout idéal premier contient $(x)$ et que $(x)$ est
maximal, l'anneau ne possède qu'un seul idéal premier, $(x)$. Ainsi, dans ce cas,
$\mathfrak p = (2, x)$. Maintenant, si $ \phi^{-1}(\mathfrak p) = (q)$ pour
$q > 2$, alors, puisque $\mathfrak p$ contient $q$, il correspond à un
idéal premier de
$\mathbf{Z}[x]/(x^2 - 4, q) \cong (\mathbf{Z}/q\mathbf{Z})[x]/(x^2 - 4)$
via le morphisme $ \mathbf{Z}[x]/(x^2 - 4) \to  \mathbf{Z}[x]/(x^2 - 4, q)$.
Tout idéal premier de $(\mathbf{Z}/q\mathbf{Z})[x]/(x^2 - 4)$ correspond à un
idéal premier de $(\mathbf{Z}/q\mathbf{Z})[x]$ contenant $(x^2 - 4) = (x -2)(x + 2)$.
Par conséquent, ces idéaux premiers doivent contenir soit $x -2$, soit $x + 2$. Puisque
$(\mathbf{Z}/q\mathbf{Z})[x]$ est un anneau principal, tous les idéaux premiers
non nuls sont maximaux, et il existe donc
exactement deux idéaux premiers de $(\mathbf{Z}/q\mathbf{Z})[x]$ contenant
$(x-2)(x + 2)$, à savoir $(x-2)$ et $(x + 2)$. En conclusion, il existe deux
idéaux premiers $(q, x-2)$ et $(q, x + 2)$ puisque $2 \neq -2 \in \mathbf{Z}/(q)$.
Enfin, traitons le cas où $\phi^{-1}(\mathfrak p) = (0)$. Remarquons
que $\mathfrak p$ correspond à un idéal premier de $\mathbf{Z}[x]$ qui
contient $(x^2 - 4) = (x -2)(x + 2)$. Ainsi, $\mathfrak p$ contient soit
$(x-2)$, soit $(x + 2)$. Par conséquent, $\mathfrak p$ correspond à un idéal premier de
$\mathbf{Z}[x]/(x - 2)$ ou de $\mathbf{Z}[x]/(x + 2)$ dont l'intersection
avec $\mathbf{Z}$ est réduite à $0$, par hypothèse. Puisque
$\mathbf{Z}[x]/(x - 2) \cong \mathbf{Z}$ et
$\mathbf{Z}[x]/(x + 2) \cong \mathbf{Z}$, cela signifie que $\mathfrak p$
doit correspondre à $0$ dans l'un de ces anneaux. Ainsi,
$\mathfrak p = (x - 2)$ ou $\mathfrak p = (x + 2)$ dans l'anneau initial.
\end{example}

\begin{example}
\label{algebra-example-spec-Zx}
Dans cet exemple, nous décrivons $X = \Spec(\mathbf{Z}[x])$.
Fixons $\mathfrak p \in X$.
Soit $\phi : \mathbf{Z} \to \mathbf{Z}[x]$ et remarquons
que $\phi^{-1}(\mathfrak p) \in \Spec(\mathbf{Z})$.
Si $\phi^{-1}(\mathfrak p) = (q)$ pour $q$ un nombre premier tel que $q > 0$,
alors $\mathfrak p$ correspond à un idéal premier de $(\mathbf{Z}/(q))[x]$,
qui doit être engendré par un polynôme irréductible dans
$(\mathbf{Z}/(q))[x]$. Si nous choisissons un représentant de ce polynôme
de degré minimal, il sera également irréductible dans $\mathbf{Z}[x]$.
Ainsi, dans ce cas, $\mathfrak p = (q, f_q)$, où $f_q$ est un polynôme
irréductible de $\mathbf{Z}[x]$ qui reste irréductible lorsqu'on le considère
dans $(\mathbf{Z}/(q) [x])$. Supposons maintenant que $\phi^{-1}(\mathfrak p) = (0)$.
Dans ce cas, si $\mathfrak p$ est non nul, il doit être engendré par des polynômes non constants
qui, puisque $\mathfrak p$ est premier, peuvent être supposés irréductibles dans
$\mathbf{Z}[x]$. D'après le lemme de Gauss, ces polynômes sont aussi irréductibles
dans $\mathbf{Q}[x]$. Puisque $\mathbf{Q}[x]$ est un anneau euclidien, s'il
existe au moins deux polynômes irréductibles distincts $f, g$ qui engendrent $\mathfrak p$,
alors $1 = af + bg$ pour certains $a, b \in \mathbf{Q}[x]$. En multipliant par
un dénominateur commun, nous voyons que $m = \bar{a}f + \bar{b} g$ pour
certains $\bar{a}, \bar{b} \in \mathbf{Z}[x]$ et un certain $m \in \mathbf{Z}$ non nul.
C'est une contradiction. Ainsi, lorsqu'il est non nul, $\mathfrak p$ est engendré par un
polynôme irréductible de $\mathbf{Z}[x]$.
\end{example}

\begin{example}
\label{algebra-example-spec-kxy}
Dans cet exemple, nous décrivons $X = \Spec(k[x, y])$
lorsque $k$ est un corps arbitraire.
Il est clair que $(0)$ est premier, et tout idéal principal engendré par un
polynôme irréductible est également premier puisque $k[x, y]$ est un
anneau factoriel. Supposons maintenant que $\mathfrak p$ soit un
élément de $X$ qui n'est pas principal. Puisque $k[x, y]$ est un
anneau factoriel noethérien, l'idéal premier $\mathfrak p$ peut être engendré
par un nombre fini de polynômes irréductibles $(f_1, \ldots, f_n)$.
Affirmons maintenant que, si $f, g$ sont des polynômes irréductibles de $k[x, y]$
qui ne sont pas associés, alors $(f, g) \cap k[x] \neq 0$. Pour le montrer,
il suffit d'établir que $f$ et $g$ sont premiers entre eux lorsqu'on les
considère dans $k(x)[y]$. Dans ce cas, $k(x)[y]$ est un anneau euclidien ;
en appliquant l'algorithme d'Euclide et en chassant les dénominateurs, nous
obtenons $p = af + bg$ pour $p, a, b \in k[x]$. Supposons donc que ce ne soit pas
le cas, c'est-à-dire qu'un certain élément non inversible $h \in k(x)[y]$ divise à la fois
$f$ et $g$. Alors, d'après le lemme de Gauss, pour certains $a, b \in k(x)$, nous
avons $ah | f$ et $bh | g$ avec $ah, bh \in k[x]$. Par
irréductibilité, $ah = f$ et
$bh = g$ (puisque $h \notin k(x)$). Ainsi, de retour dans $k(x)[y]$, $f, g $
sont associés, car $\frac{a}{b} g = f$. Puisque
$k(x)$ est le corps des fractions de $k[x]$, nous pouvons écrire $g = \frac{r}{s} f $
pour des éléments $r , s \in k[x]$ sans facteur commun. Cela
entraîne que $sg = rf$ dans $k[x, y]$, et donc que $s$ doit diviser $f$,
puisque $k[x, y]$ est un anneau factoriel. Ainsi, $s = 1$ ou $s = f$. Si $s = f$,
alors $r = g$, ce qui entraîne $f, g \in k[x]$ ; ils doivent donc être des unités de
$k(x)$ et être premiers entre eux dans $k(x)[y]$, en contradiction avec notre
hypothèse. Si $s = 1$, alors $g = rf$, autre contradiction.
Ainsi, $f, g$ doivent être premiers entre eux dans $k(x)[y]$, qui est un
anneau euclidien. Nous nous sommes donc ramenés au cas où $\mathfrak p$
contient un polynôme irréductible $p \in k[x] \subset k[x, y]$.
D'après ce qui précède, $\mathfrak p$ correspond à un idéal premier de l'anneau
$k[x, y]/(p) = k(\alpha)[y]$, où $\alpha$ est un élément
algébrique sur $k$ de polynôme minimal $p$. Cet anneau est
principal, et tout idéal premier correspond donc à $(0)$ ou à un
polynôme irréductible de $k(\alpha)[y]$. Ainsi, $\mathfrak p$
est de la forme $(p)$ ou $(p, f)$, où $f$ est un
polynôme de $k[x, y]$ qui est irréductible dans le quotient
$k[x, y]/(p)$.
\end{example}

\begin{example}
\label{algebra-example-affine-open-not-standard}
Considérons l'anneau
$$
R = \{ f \in \mathbf{Q}[z]\text{ vérifiant }f(0) = f(1) \}.
$$
Considérons le morphisme
$$
\varphi : \mathbf{Q}[A, B] \to R
$$
défini par $\varphi(A) = z^2-z$ et $\varphi(B) = z^3-z^2$. On vérifie
aisément que $(A^3 - B^2 + AB) \subset \Ker(\varphi)$ et que
$A^3 - B^2 + AB$ est irréductible. Supposons que $\varphi$ soit surjectif ;
alors, puisque $R$ est un anneau intègre (c'est un sous-anneau d'un anneau
intègre), $\Ker(\varphi)$ doit être un idéal premier de $\mathbf{Q}[A, B]$.
Les idéaux premiers qui contiennent $(A^3-B^2 + AB)$ sont $(A^3-B^2 + AB)$
lui-même et tout idéal maximal $(f, g)$, avec $f, g\in\mathbf{Q}[A, B]$,
tel que $f$ soit irréductible modulo $g$. Mais $R$ n'est pas un corps, donc le
noyau doit être $(A^3-B^2 + AB)$ ; ainsi, $\varphi$ donne un isomorphisme
$R \to \mathbf{Q}[A, B]/(A^3-B^2 + AB)$.

\medskip\noindent
Pour voir que $\varphi$ est surjectif, nous devons exprimer tout
$f\in R$ comme un polynôme à coefficients dans $\mathbf{Q}$ en $A(z) = z^2-z$
et $B(z) = z^3-z^2$. Notons la relation $zA(z) = B(z)$. Soit
$a = f(0) = f(1)$. Alors $z(z-1)$ doit diviser $f(z)-a$, et nous pouvons donc écrire
$f(z) = z(z-1)g(z)+a = A(z)g(z)+a$. Si $\deg(g) < 2$, alors
$h(z) = c_1z + c_0$ et $f(z) = A(z)(c_1z + c_0)+a = c_1B(z)+c_0A(z)+a$, et nous
avons terminé. Si $\deg(g)\geq 2$, alors, d'après l'algorithme de division
euclidienne des polynômes, nous pouvons écrire $g(z) = A(z)h(z)+b_1z + b_0$
($\deg(h)\leq\deg(g)-2$), de sorte que $f(z) = A(z)^2h(z)+b_1B(z)+b_0A(z)$.
En appliquant la division à $h(z)$ et en itérant, nous obtenons une expression
de $f(z)$ comme polynôme en $A(z)$ et $B(z)$ ; ainsi, $\varphi$ est
surjectif.

\medskip\noindent
Soit maintenant $a \in \mathbf{Q}$, $a \neq 0, \frac{1}{2}, 1$, et
considérons
$$
R_a = \{ f \in \mathbf{Q}[z, \frac{1}{z-a}]\text{ vérifiant }f(0) = f(1)
\}.
$$
C'est encore une $\mathbf{Q}$-algèbre de type fini : on vérifie
aisément que les fonctions $z^2-z$, $z^3-z$ et
$\frac{a^2-a}{z-a}+z$ engendrent $R_a$ comme $\mathbf{Q}$-algèbre. Nous
avons les inclusions suivantes :
$$
R\subset R_a\subset\mathbf{Q}[z, \frac{1}{z-a}], \quad
R\subset\mathbf{Q}[z]\subset\mathbf{Q}[z, \frac{1}{z-a}].
$$
Rappelons (lemme \ref{algebra-lemma-spec-localization}) que, pour un anneau T et une
partie multiplicative $S\subset T$, le morphisme d'anneaux $T \to S^{-1}T$
induit un morphisme sur les spectres $\Spec(S^{-1}T) \to \Spec(T)$
qui est un homéomorphisme sur le sous-ensemble
$$
\{\mathfrak p \in \Spec(T) \mid S \cap \mathfrak p = \emptyset\}
\subset \Spec(T).
$$
Lorsque $S = \{ 1, f, f^2, \ldots\}$ pour un certain $f\in T$, il s'agit
de l'ouvert $D(f)\subset T$. Vérifions maintenant une propriété analogue
pour le morphisme d'anneaux $R \to R_a$ : nous allons montrer que le morphisme
$\theta : \Spec(R_a) \to \Spec(R)$ induit par l'inclusion
$R\subset R_a$ est un homéomorphisme sur un sous-ensemble ouvert de
$\Spec(R)$ en vérifiant que $\theta$ est un homéomorphisme local
injectif. Nous le faisons relativement à un recouvrement ouvert de
$\Spec(R_a)$ par deux ouverts principaux, que nous décrivons maintenant.
Pour tout $r\in\mathbf{Q}$, soit $\text{ev}_r : R \to \mathbf{Q}$ le
morphisme donné par l'évaluation en $r$. Notons que, pour $r = 0$ et
$r = 1-a$, celui-ci se prolonge en un morphisme
$\text{ev}_r' : R_a \to \mathbf{Q}$ (dans le second cas, parce que $\frac{1}{z-a}$
est bien défini en $z = 1-a$, puisque $a\neq\frac{1}{2}$). Toutefois,
$\text{ev}_a$ ne se prolonge pas à $R_a$. Posons
$\mathfrak{m}_r = \Ker(\text{ev}_r)$. Nous avons
$$
\mathfrak{m}_0 = (z^2-z, z^3-z),
$$
$$
\mathfrak{m}_a = ((z-1 + a)(z-a), (z^2-1 + a)(z-a)), \text{ et}
$$
$$
\mathfrak{m}_{1-a} = ((z-1 + a)(z-a), (z-1 + a)(z^2-a)).
$$
Pour le vérifier, notons que les membres de droite sont clairement contenus dans
les membres de gauche. Vérifions ensuite que les membres de droite sont des
idéaux maximaux en écrivant les générateurs en fonction de $A$ et $B$,
et en considérant $R$ comme $\mathbf{Q}[A, B]/(A^3-B^2 + AB)$. Notons que
$\mathfrak{m}_a$ n'appartient pas à l'image de $\theta$ : nous avons
$$
(z^2 - z)^2(z - a)\left(\frac{a^2 - a}{z - a} + z\right) =
(z^2 - z)^2(a^2 - a) + (z^2 - z)^2(z - a)z
$$
Le membre de gauche appartient à $\mathfrak m_a R_a$ parce que
$(z^2 - z)(z - a)$ appartient à $\mathfrak m_a$ et que
$(z^2 - z)(\frac{a^2 - a}{z - a} + z)$ appartient à $R_a$. De même,
l'élément $(z^2 - z)^2(z - a)z$ appartient à $\mathfrak m_a R_a$
parce que $(z^2 - z)$ appartient à $R_a$ et que $(z^2 - z)(z - a)$ appartient à $\mathfrak m_a$.
Comme $a \not \in \{0, 1\}$, nous en concluons que
$(z^2 - z)^2 \in \mathfrak m_a R_a$. Par conséquent,
aucun idéal $I$ de $R_a$ ne peut vérifier $I \cap R = \mathfrak m_a$, car un tel
$I$ devrait contenir $(z^2 - z)^2$, qui appartient à $R$ mais pas à
$\mathfrak m_a$. L'ouvert principal
$D((z-1 + a)(z-a))\subset\Spec(R)$ est égal au complémentaire du
fermé $\{\mathfrak{m}_a, \mathfrak{m}_{1-a}\}$.
Vérifions alors que $R_{(z-1 + a)(z-a)} = (R_a)_{(z-1 + a)(z-a)}$ ; en
notant cet anneau localisé $R'$, il s'ensuit que le morphisme $R \to R'$
se factorise en $R \to R_a \to R'$. D'après le lemme
\ref{algebra-lemma-spec-localization}, ces morphismes présentent alors
$\Spec(R') \subset \Spec(R_a)$ et
$\Spec(R') \subset \Spec(R)$ comme des sous-ensembles ouverts ; ainsi,
$\theta : \Spec(R_a) \to \Spec(R)$, restreint à
$D((z-1 + a)(z-a))$, est un homéomorphisme sur un sous-ensemble ouvert.
De même, la restriction de $\theta$ à
$D((z^2 + z + 2a-2)(z-a)) \subset \Spec(R_a)$ est un homéomorphisme
sur le sous-ensemble ouvert $D((z^2 + z + 2a-2)(z-a)) \subset \Spec(R)$.
Selon que $z^2 + z + 2a-2$ est irréductible ou non sur
$\mathbf{Q}$, le complémentaire de ce dernier ouvert principal
est constitué d'un ou deux points fermés, ainsi que du point fermé
$\mathfrak{m}_a$. De plus, l'idéal de $R_a$ engendré par les
éléments $(z^2 + z + 2a-a)(z-a)$ et $(z-1 + a)(z-a)$ est égal à $R_a$ tout entier ; ainsi,
ces deux ouverts principaux recouvrent $\Spec(R_a)$. Par conséquent, pour
montrer que $\theta$ est un homéomorphisme sur
$\Spec(R)-\{\mathfrak{m}_a\}$, il suffit de montrer
que ces un ou deux points ne peuvent jamais être égaux à $\mathfrak{m}_{1-a}$.
Et c'est bien le cas, puisque $1-a$ est une racine de $z^2 + z + 2a-2$
si et seulement si $a = 0$ ou $a = 1$, deux cas qui ne se produisent pas.

\medskip\noindent
Malgré cet homéomorphisme, qui imite le comportement d'une
localisation en un élément de $R$, tandis que
$\mathbf{Q}[z, \frac{1}{z-a}]$ est la localisation de $\mathbf{Q}[z]$
relativement à l'élément $(z-a)$, l'anneau $R_a$ n'est {\it pas} une
localisation de $R$ : toute localisation $S^{-1}R$ produit davantage
d'unités que l'anneau initial $R$. Les unités de $R$ sont
$\mathbf{Q}^\times$, les unités de $\mathbf{Q}$. En fait, il est
facile de voir que les unités de $R_a$ sont $\mathbf{Q}^*$.
En effet, les unités de $\mathbf{Q}[z, \frac{1}{z - a}]$ sont
$c (z - a)^n$ pour $c \in \mathbf{Q}^*$ et $n \in \mathbf{Z}$,
et il est clair que celles-ci n'appartiennent à $R_a$ que si $n = 0$.
Ainsi, $R_a$ ne possède pas davantage d'unités que
$R$ et ne peut donc pas être une localisation de $R$.

\medskip\noindent
Nous avons utilisé le fait que $a\neq 0, 1$ pour garantir que
$\frac{1}{z-a}$ est défini en $z = 0, 1$. Nous avons utilisé le fait que
$a\neq 1/2$ en plusieurs endroits : (1) pour pouvoir parler du
noyau de $\text{ev}_{1-a}$ sur $R_a$, ce qui garantit que
$\mathfrak{m}_{1-a}$ est un point de $R_a$ (c'est-à-dire que $R_a$
ne possède qu'un point de moins que $R$). (2) À la fin, pour conclure
que $(z-a)^{k + \ell}$ ne peut appartenir à $R$ que pour $k = \ell = 0$ ; en effet, si
$a = 1/2$, cet élément appartient à $R$ dès que $k + \ell$ est pair. Il
y aurait alors effectivement davantage d'unités dans $R_a$ que dans $R$, et $R_a$
pourrait éventuellement être une localisation de $R$.
\end{example}





\section{Une méta-observation sur les idéaux premiers}
\label{algebra-section-oka-families}

\noindent
Cette section est tirée du projet CRing. Soit $R$ un anneau et
soit $S \subset R$ une partie multiplicative.
Une conséquence du
lemme \ref{algebra-lemma-spec-localization}
est qu'un idéal $I \subset R$ maximal parmi ceux qui ne
rencontrent pas $S$ est premier. La raison en est que $I = R \cap \mathfrak m$
pour un certain idéal maximal $\mathfrak m$ de l'anneau $S^{-1}R$.
Il s'avère que, pour de nombreuses propriétés d'idéaux, les éléments maximaux
sont premiers. Une méthode générale permettant de le voir a été développée dans \cite{Lam-Reyes}.
Dans cette section, nous faisons une digression pour expliquer ce phénomène.

\medskip\noindent
Soit $R$ un anneau. Si $I$ est un idéal de $R$ et $a \in R$, nous
définissons
$$
(I : a) = \left\{ x \in R \mid xa \in I\right\}.
$$
Plus généralement, si $J \subset R$ est un idéal, nous définissons
$$
(I : J) = \left\{ x \in R \mid xJ \subset I\right\}.
$$

\begin{lemma}
\label{algebra-lemma-colon}
Soit $R$ un anneau. Pour un idéal principal $J \subset R$ et pour tout idéal
$I \subset J$, nous avons $I = J (I : J)$.
\end{lemma}

\begin{proof}
Écrivons $J = (a)$. Alors $(I : J) = (I : a)$.
Puisque $I \subset J$, nous voyons que tout $y \in I$ est de la forme
$y = xa$ pour un certain $x \in (I : a)$. Ainsi, $I \subset J (I : J)$.
Réciproquement, si $x \in (I : a)$, alors $xJ = (xa) \subset I$, ce qui
démontre l'autre inclusion.
\end{proof}

\noindent
Soit $\mathcal{F}$ une collection d'idéaux de $R$. Nous cherchons des
conditions garantissant que les éléments maximaux du complémentaire
de $\mathcal{F}$ sont premiers.

\begin{definition}
\label{algebra-definition-oka-family}
Soit $R$ un anneau. Soit $\mathcal{F}$ un ensemble d'idéaux de $R$. Nous disons
que $\mathcal{F}$ est une {\it famille d'Oka} si $R \in \mathcal{F}$ et si,
chaque fois que $I \subset R$ est un idéal et que $(I : a), (I, a) \in \mathcal{F}$
pour un certain $a \in R$, alors $I \in \mathcal{F}$.
\end{definition}

\noindent
Donnons quelques exemples de familles d'Oka. Le premier est l'exemple
fondamental présenté dans l'introduction de cette section.

\begin{example}
\label{algebra-example-oka-family-not-meet-multiplicative-set}
Soit $R$ un anneau et soit $S$ une partie multiplicative de $R$.
Nous affirmons que $\mathcal{F} = \{I \subset R \mid I \cap S \not = \emptyset\}$
est une famille d'Oka. En effet, supposons que $(I : a), (I, a) \in \mathcal{F}$
pour un certain $a \in R$. Choisissons alors $s \in (I, a) \cap S$ et
$s' \in (I : a) \cap S$. Alors $ss' \in I \cap S$, et donc
$I \in \mathcal{F}$. Ainsi, $\mathcal{F}$ est une famille d'Oka.
\end{example}

\begin{example}
\label{algebra-example-oka-family-finitely-generated}
Soient $R$ un anneau, $I \subset R$ un idéal et $a \in R$.
Si $(I : a)$ est engendré par $a_1, \ldots, a_n$ et si $(I, a)$
est engendré par $a, b_1, \ldots, b_m$, avec
$b_1, \ldots, b_m \in I$, alors $I$ est engendré par
$aa_1, \ldots, aa_n, b_1, \ldots, b_m$.
Pour le voir, notons que, si $x \in I$, alors $x \in (I, a)$
est une combinaison linéaire de $a, b_1, \ldots, b_m$, mais le
coefficient de $a$ doit appartenir à $(I:a)$.
Nous en déduisons que
la famille des idéaux de type fini est une famille d'Oka.
\end{example}

\begin{example}
\label{algebra-example-oka-family-principal}
Montrons que la famille des idéaux principaux d'un anneau $R$ est une famille d'Oka.
En effet, supposons que $I \subset R$ soit un idéal, que $a \in R$ et que $(I, a)$ et
$(I : a)$ soient principaux. Notons que $(I : a) = (I : (I, a))$.
En posant $J = (I, a)$, nous constatons que $J$ est principal et que $(I : J)$ l'est aussi. D'après
le lemme \ref{algebra-lemma-colon},
nous avons $I = J (I : J)$.
Ainsi, dans notre situation, puisque $J = (I, a)$ et $(I : J)$
sont principaux, $I$ est principal.
\end{example}

\begin{example}
\label{algebra-example-oka-family-bound-cardinality}
Soit $R$ un anneau.
Soit $\kappa$ un cardinal infini.
La famille des idéaux qui peuvent être engendrés par au plus $\kappa$ éléments
est une famille d'Oka. L'argument est analogue à celui de
l'exemple \ref{algebra-example-oka-family-finitely-generated}
et est omis.
\end{example}

\begin{example}
\label{algebra-example-oka-family-property-modules}
Soient $A$ un anneau, $I \subset A$ un idéal et $a \in A$ un élément.
Il existe une suite exacte courte $0 \to A/(I : a) \to A/I \to A/(I, a) \to 0$
dont la première flèche est donnée par la multiplication par $a$. Ainsi, si $P$
est une propriété des $A$-modules qui est stable par extensions et qui est vérifiée
par $0$, alors la famille des idéaux $I$ tels que $A/I$ possède $P$
est une famille d'Oka.
\end{example}

\begin{proposition}
\label{algebra-proposition-oka}
Si $\mathcal{F}$ est une famille d'Oka d'idéaux, alors tout élément maximal du
complémentaire de $\mathcal{F}$ est premier.
\end{proposition}

\begin{proof}
Supposons que $I \not \in \mathcal{F}$ soit maximal parmi les idéaux qui n'appartiennent pas à
$\mathcal{F}$, mais que $I$ ne soit pas premier. Notons que $I \not = R$, car
$R \in \mathcal{F}$. Puisque $I$ n'est pas premier, nous pouvons trouver $a, b \in R - I$
tels que $ab \in I$. Il s'ensuit que $(I, a) \neq I$ et que $(I : a)$ contient
$b \not \in I$, donc que $(I : a) \neq I$ également. Ainsi, $(I : a), (I, a)$ contiennent tous deux
strictement $I$ et doivent donc appartenir à $\mathcal{F}$.
Par la condition d'Oka, nous avons $I \in \mathcal{F}$, ce qui est une contradiction.
\end{proof}

\noindent
Nous pouvons maintenant convertir la plupart des exemples ci-dessus en
un lemme sur les idéaux premiers d'un anneau.

\begin{lemma}
\label{algebra-lemma-simple}
Soit $R$ un anneau. Soit $S$ une partie multiplicative de $R$.
Un idéal $I \subset R$ maximal parmi ceux qui vérifient
$I \cap S = \emptyset$ est premier.
\end{lemma}

\begin{proof}
C'est l'exemple présenté dans l'introduction de cette section.
Pour une autre démonstration, combiner
l'exemple \ref{algebra-example-oka-family-not-meet-multiplicative-set}
avec la
proposition \ref{algebra-proposition-oka}.
\end{proof}

\begin{lemma}
\label{algebra-lemma-cohen}
Soit $R$ un anneau.
\begin{enumerate}
\item Un idéal $I \subset R$ maximal parmi ceux qui ne sont pas
de type fini est premier.
\item Si tout idéal premier de $R$ est
de type fini, alors
tout idéal de $R$ est de type fini\footnote{Nous dirons plus loin
que $R$ est noethérien.}.
\end{enumerate}
\end{lemma}

\begin{proof}
La première assertion est une conséquence immédiate de
l'exemple \ref{algebra-example-oka-family-finitely-generated} et de la
proposition \ref{algebra-proposition-oka}. Pour la seconde,
supposons qu'il existe un idéal $I \subset R$ qui ne soit pas de type
fini. La réunion d'une chaîne totalement ordonnée $\left\{I_\alpha\right\}$
d'idéaux qui ne sont pas de type fini n'est pas de type fini ;
en effet, si $I = \bigcup I_\alpha$ était engendré par
$a_1, \ldots, a_n$, alors tous les générateurs appartiendraient à un certain
$I_\alpha $ et l'engendreraient par conséquent.
D'après le lemme de Zorn, il existe un idéal maximal parmi ceux qui ne sont pas de type
fini. D'après la première partie, cet idéal est premier.
\end{proof}

\begin{lemma}
\label{algebra-lemma-primes-principal}
Soit $R$ un anneau.
\begin{enumerate}
\item Un idéal $I \subset R$ maximal parmi ceux qui ne sont pas
principaux est premier.
\item Si tout idéal premier de $R$ est principal, alors
tout idéal de $R$ est principal.
\end{enumerate}
\end{lemma}

\begin{proof}
La première partie résulte de
l'exemple \ref{algebra-example-oka-family-principal} et de la
proposition \ref{algebra-proposition-oka}.
Pour la seconde, supposons qu'il existe un idéal $I \subset R$
qui ne soit pas principal. La réunion d'une chaîne totalement ordonnée
$\left\{I_\alpha\right\}$ d'idéaux qui ne sont pas principaux n'est pas principale ;
en effet, si $I = \bigcup I_\alpha$ était engendré par
$a$, alors $a$ appartiendrait à un certain $I_\alpha $ et $a$ l'engendrerait.
D'après le lemme de Zorn, il existe un idéal maximal parmi ceux qui ne sont pas
principaux. Cet idéal est nécessairement premier d'après la première partie.
\end{proof}

\begin{lemma}
\label{algebra-lemma-characterize-domain}
Soit $R$ un anneau.
\begin{enumerate}
\item Un idéal maximal parmi les idéaux qui ne contiennent aucun
élément non diviseur de zéro est premier.
\item Si $R$ est non nul et si tout idéal premier non nul de $R$
contient un élément non diviseur de zéro, alors $R$ est un anneau intègre.
\end{enumerate}
\end{lemma}

\begin{proof}
Considérons l'ensemble $S$ des éléments non diviseurs de zéro. C'est une partie
multiplicative de $R$. Ainsi, tout idéal maximal parmi ceux qui ne rencontrent pas
$S$ est premier, voir le
lemme \ref{algebra-lemma-simple}.
Par conséquent, si tout idéal premier non nul contient un élément non diviseur de zéro, alors
$(0)$ est premier, c'est-à-dire que $R$ est un anneau intègre.
\end{proof}

\begin{remark}
\label{algebra-remark-cohen-bound-cardinality}
Soit $R$ un anneau. Soit $\kappa$ un cardinal infini.
En appliquant
l'exemple \ref{algebra-example-oka-family-bound-cardinality} et la
proposition \ref{algebra-proposition-oka},
nous voyons que tout idéal maximal parmi ceux qui ne peuvent pas être
engendrés par $\kappa$ éléments est premier. Ce résultat n'est pas très
utile, car il existe un anneau $R$ dont tout idéal premier
peut être engendré par $\aleph_0$ éléments, mais dont un certain
idéal ne le peut pas. En effet, soit $k$ un corps, soit $T$ un ensemble dont
la cardinalité est strictement supérieure à $\aleph_0$, et posons
$$
R = k[\{x_n\}_{n \geq 1}, \{z_{t, n}\}_{t \in T, n \geq 0}]/
(x_n^2, z_{t, n}^2, x_n z_{t, n} - z_{t, n - 1})
$$
C'est un anneau local dont l'unique idéal premier est
$\mathfrak m = (x_n)$. Mais l'idéal $(z_{t, n})$ ne peut pas
être engendré par une famille dénombrable d'éléments.
\end{remark}

\begin{example}
\label{algebra-example-noetherian-topology-on-spec}
\begin{reference}
Commentaire de Lukas Heger du 12 novembre 2020.
\end{reference}
Soient $R$ un anneau et $X = \Spec(R)$. Puisque les fermés de $X$ correspondent aux
idéaux radicaux de $R$ (lemme \ref{algebra-lemma-Zariski-topology}), nous voyons que $X$ est un
espace topologique noethérien si et seulement si les idéaux radicaux vérifient la condition de chaîne ascendante.
Cela équivaut à dire que tout idéal radical est le radical d'un idéal de type
fini (détails omis). Posons
$$
\mathcal{F} = \{I \subset R \mid
\sqrt{I} = \sqrt{(f_1, \ldots, f_n)}\text{ pour un certain }n
\text{ et }f_1, \ldots, f_n \in R\}.
$$
Le lecteur peut montrer que $\mathcal{F}$ est une famille d'Oka en utilisant l'identité
$$
\sqrt{I} = \sqrt{(I, a)(I : a)}
$$
qui vaut pour tout idéal $I \subset R$ et tout élément $a \in R$.
D'autre part, si nous avons une chaîne totalement ordonnée d'idéaux
$\{I_\alpha\}$ dont aucun n'appartient à $\mathcal{F}$, alors la réunion
$I = \bigcup I_\alpha$ ne peut pas non plus appartenir à $\mathcal{F}$. Sinon,
$\sqrt{I} = \sqrt{(f_1, \ldots, f_n)}$, puis $f_i^e \in I$ pour un certain $e$,
puis $f_i^e \in I_\alpha$ pour un certain $\alpha$ indépendant de $i$, et enfin
$\sqrt{I_\alpha} = \sqrt{(f_1, \ldots, f_n)}$, contradiction.
Ainsi, si l'ensemble des idéaux qui n'appartiennent pas à $\mathcal{F}$ est non vide, il
possède des éléments maximaux et,
exactement comme dans le lemme \ref{algebra-lemma-cohen}, nous concluons que $X$ est un
espace topologique noethérien si et seulement si tout idéal premier de $R$
est égal à $\sqrt{(f_1, \ldots, f_n)}$ pour certains $f_1, \ldots, f_n \in R$.
Si ce résultat nous est un jour nécessaire, nous l'énoncerons et le démontrerons
ici avec soin.
\end{example}








\section{Images des morphismes d'anneaux de présentation finie}
\label{algebra-section-images-finite-presentation}

\noindent
Dans cette section, nous démontrons quelques résultats sur la
topologie des applications $\Spec(S) \to \Spec(R)$
induites par des morphismes d'anneaux $R \to S$, principalement le théorème de Chevalley.
Pour ce faire, nous utiliserons les notions de parties constructibles,
de parties quasi-compactes, de parties rétrocompactes, etc.,
qui sont définies dans Topologie, section \ref{topology-section-constructible}.

\begin{lemma}
\label{algebra-lemma-qc-open}
Soit $U \subset \Spec(R)$ un ouvert. Les conditions suivantes
sont équivalentes :
\begin{enumerate}
\item $U$ est rétrocompact dans $\Spec(R)$,
\item $U$ est quasi-compact,
\item $U$ est une réunion finie d'ouverts principaux, et
\item il existe un idéal de type fini $I \subset R$ tel
que $X \setminus V(I) = U$.
\end{enumerate}
\end{lemma}

\begin{proof}
Nous avons (1) $\Rightarrow$ (2), car $\Spec(R)$ est quasi-compact, voir le
lemme \ref{algebra-lemma-quasi-compact}. Nous avons (2) $\Rightarrow$ (3), car les
ouverts principaux forment une base de la topologie. Démontrons (3) $\Rightarrow$ (1).
Soit $U = \bigcup_{i = 1\ldots n} D(f_i)$. Pour montrer que $U$ est rétrocompact
dans $\Spec(R)$, il suffit de montrer que $U \cap V$ est quasi-compact pour tout
ouvert quasi-compact $V$ de $\Spec(R)$. Écrivons
$V = \bigcup_{j = 1\ldots m} D(g_j)$, ce qui est possible par (2) $\Rightarrow$
(3). Tout ouvert principal est homéomorphe au spectre d'un anneau et est donc
quasi-compact, voir les lemmes \ref{algebra-lemma-standard-open} et
\ref{algebra-lemma-quasi-compact}. Ainsi,
$U \cap V =
(\bigcup_{i = 1\ldots n} D(f_i)) \cap (\bigcup_{j = 1\ldots m} D(g_j))
= \bigcup_{i, j} D(f_i g_j)$ est une réunion finie d'ouverts quasi-compacts,
donc est quasi-compact. Pour achever la démonstration, notons
que (4) équivaut à (3) d'après le
lemme \ref{algebra-lemma-Zariski-topology}.
\end{proof}

\begin{lemma}
\label{algebra-lemma-affine-map-quasi-compact}
Soit $\varphi : R \to S$ un morphisme d'anneaux.
L'application continue induite $f : \Spec(S) \to \Spec(R)$
est quasi-compacte. Pour toute partie constructible $E \subset \Spec(R)$,
l'image réciproque $f^{-1}(E)$ est constructible dans $\Spec(S)$.
\end{lemma}

\begin{proof}
Montrons d'abord que l'image réciproque de tout ouvert quasi-compact
$U \subset \Spec(R)$ est quasi-compacte. D'après le
lemme \ref{algebra-lemma-qc-open}, nous pouvons écrire $U$ comme une réunion
finie d'ouverts principaux. Le lemme \ref{algebra-lemma-spec-functorial}
montre donc que $f^{-1}(U)$ est une réunion finie d'ouverts principaux.
Ainsi, $f^{-1}(U)$ est quasi-compact d'après le lemme \ref{algebra-lemma-qc-open}.
La seconde assertion résulte alors de Topologie, lemme
\ref{topology-lemma-inverse-images-constructibles}.
\end{proof}

\begin{lemma}
\label{algebra-lemma-constructible}
Soit $R$ un anneau. Une partie de $\Spec(R)$ est constructible si et seulement
si elle peut s'écrire comme une réunion finie de parties de la forme
$D(f) \cap V(g_1, \ldots, g_m)$, où $f, g_1, \ldots, g_m \in R$.
\end{lemma}

\begin{proof}
D'après le lemme \ref{algebra-lemma-qc-open}, la partie $D(f)$ et le complémentaire de
$V(g_1, \ldots, g_m)$ sont des ouverts rétrocompacts. Ainsi,
$D(f) \cap V(g_1, \ldots, g_m)$ est une partie constructible, de même que
toute réunion finie de telles parties. Réciproquement, soit $T \subset \Spec(R)$ une
partie constructible. D'après Topologie, définition \ref{topology-definition-constructible},
nous pouvons supposer que $T = U \cap V^c$, où $U, V \subset \Spec(R)$
sont des ouverts rétrocompacts. D'après le lemme \ref{algebra-lemma-qc-open}, nous pouvons écrire
$U = \bigcup_{i = 1, \ldots, n} D(f_i)$ et
$V = \bigcup_{j = 1, \ldots, m} D(g_j)$. Alors
$T = \bigcup_{i = 1, \ldots, n} \big(D(f_i) \cap V(g_1, \ldots, g_m)\big)$.
\end{proof}

\begin{lemma}
\label{algebra-lemma-constructible-is-image}
Soit $R$ un anneau et soit $T \subset \Spec(R)$
une partie constructible. Il existe alors un morphisme d'anneaux $R \to S$ de
présentation finie tel que $T$ soit l'image de
$\Spec(S)$ dans $\Spec(R)$.
\end{lemma}

\begin{proof}
Le spectre d'un produit fini d'anneaux
est la réunion disjointe des spectres, voir le
lemme \ref{algebra-lemma-spec-product}. Par conséquent, si $T = T_1 \cup T_2$
et si le résultat vaut pour $T_1$ et $T_2$, alors il
vaut pour $T$.
D'après le lemme \ref{algebra-lemma-constructible}, nous pouvons supposer
que $T = D(f) \cap V(g_1, \ldots, g_m)$.
Dans ce cas, $T$ est l'image du morphisme
$\Spec((R/(g_1, \ldots, g_m))_f) \to \Spec(R)$, voir les lemmes
\ref{algebra-lemma-standard-open} et \ref{algebra-lemma-spec-closed}.
\end{proof}

\begin{lemma}
\label{algebra-lemma-open-fp}
Soit $R$ un anneau.
Soit $f$ un élément de $R$.
Soit $S = R_f$.
Alors l'image d'une partie constructible de $\Spec(S)$
est constructible dans $\Spec(R)$.
\end{lemma}

\begin{proof}
Nous utilisons à plusieurs reprises le lemme \ref{algebra-lemma-qc-open} sans le mentionner.
Soient $U, V$ des ouverts quasi-compacts de $\Spec(S)$.
Nous allons montrer que l'image de $U \cap V^c$ est constructible.
Sous l'identification
$\Spec(S) = D(f)$ du lemme \ref{algebra-lemma-standard-open},
les parties $U, V$ correspondent à des ouverts quasi-compacts
$U', V'$ de $\Spec(R)$.
Il suffit donc de montrer que $U' \cap (V')^c$
est constructible dans $\Spec(R)$, ce qui est clair.
\end{proof}

\begin{lemma}
\label{algebra-lemma-closed-fp}
Soit $R$ un anneau.
Soit $I$ un idéal de type fini de $R$.
Soit $S = R/I$.
Alors l'image d'une partie constructible de $\Spec(S)$
est constructible dans $\Spec(R)$.
\end{lemma}

\begin{proof}
Si $I = (f_1, \ldots, f_m)$, alors nous voyons que
$V(I)$ est le complémentaire de $\bigcup D(f_i)$,
voir le lemme \ref{algebra-lemma-Zariski-topology}.
Il est donc constructible, d'après le lemme \ref{algebra-lemma-qc-open}.
Notons le morphisme $R \to S$ par $f \mapsto \overline{f}$.
Nous devons montrer que, si $\overline{U}, \overline{V}$
sont des ouverts rétrocompacts de $\Spec(S)$, alors l'image de
$\overline{U} \cap \overline{V}^c$
dans $\Spec(R)$ est constructible.
D'après le lemme \ref{algebra-lemma-qc-open}, nous pouvons écrire
$\overline{U} = \bigcup D(\overline{g_i})$.
En posant $U = \bigcup D({g_i})$, nous voyons que $\overline{U}$
a pour image $U \cap V(I)$, qui est constructible dans
$\Spec(R)$. De même, l'image de $\overline{V}$ est égale à
$V \cap V(I)$ pour un certain ouvert rétrocompact $V$ de $\Spec(R)$.
Ainsi, l'image de $\overline{U} \cap \overline{V}^c$
est égale à $U \cap V(I) \cap V^c$, comme voulu.
\end{proof}

\begin{lemma}
\label{algebra-lemma-affineline-open}
Soit $R$ un anneau. Le morphisme $\Spec(R[x]) \to \Spec(R)$
est ouvert, et l'image de tout ouvert principal est un ouvert
quasi-compact.
\end{lemma}

\begin{proof}
Il suffit de montrer que l'image d'un ouvert principal
$D(f)$, $f\in R[x]$, est un ouvert quasi-compact.
L'image de $D(f)$ est celle du morphisme
$\Spec(R[x]_f) \to \Spec(R)$.
Soit $\mathfrak p \subset R$ un idéal premier.
Soit $\overline{f}$ l'image de $f$ dans
$\kappa(\mathfrak p)[x]$.
Rappelons, voir le lemme \ref{algebra-lemma-in-image},
que $\mathfrak p$ appartient à l'image
si et seulement si $R[x]_f \otimes_R \kappa(\mathfrak p) =
\kappa(\mathfrak p)[x]_{\overline{f}}$ n'est pas l'anneau
nul. C'est exactement la condition que $f$ ne s'envoie pas
sur zéro dans $\kappa(\mathfrak p)[x]$, autrement dit qu'un
coefficient de $f$ n'appartienne pas à $\mathfrak p$.
Ainsi, si $f = a_d x^d + \ldots + a_0$, alors
l'image de $D(f)$ est $D(a_d) \cup \ldots \cup D(a_0)$.
\end{proof}

\noindent
Nous démontrons une propriété des polynômes caractéristiques qui
sera utilisée ci-dessous.

\begin{lemma}
\label{algebra-lemma-characteristic-polynomial-prime}
Soit $R \to A$ un morphisme d'anneaux.
Supposons que $A \cong R^{\oplus n}$ comme $R$-module.
Soit $f \in A$. L'application de multiplication $m_f: A
\to A$ est $R$-linéaire et possède donc
un polynôme caractéristique
$P(T) = T^n + r_{n-1}T^{n-1} + \ldots + r_0 \in R[T]$.
Pour tout idéal premier
$\mathfrak{p} \in \Spec(R)$, $f$ agit de manière nilpotente sur $A
\otimes_R \kappa(\mathfrak{p})$ si et seulement si $\mathfrak p \in
V(r_0, \ldots, r_{n-1})$.
\end{lemma}

\begin{proof}
Cela résulte assez facilement du fait que le polynôme
caractéristique $\bar P(T) \in \kappa(\mathfrak p)[T]$ de l'application
de multiplication $m_{\bar f}: A \otimes_R \kappa(\mathfrak p) \to
A \otimes_R \kappa(\mathfrak p)$, qui multiplie les éléments de $A
\otimes_R \kappa(\mathfrak p)$ par $\bar f$, l'image de $f$ dans cette algèbre sur
$\kappa(\mathfrak p)$, est précisément l'image de $P(T)$ dans
$\kappa(\mathfrak p)[T]$. Soit $(a_{ij})$ la matrice de l'application
$m_f$, à coefficients dans $R$, relativement à une base $e_1, \ldots, e_n$
de $A$ comme $R$-module.
Alors $A \otimes_R \kappa(\mathfrak p) \cong (R \otimes_R
\kappa(\mathfrak p))^{\oplus n} = \kappa(\mathfrak p)^n$, qui est
un espace vectoriel de dimension $n$ sur $\kappa(\mathfrak p)$, de
base $e_1 \otimes 1, \ldots, e_n \otimes 1$. L'image est $\bar f = f
\otimes 1$, et l'application de multiplication $m_{\bar f}$ a donc pour matrice
$(a_{ij} \otimes 1)$. Ainsi, le polynôme caractéristique est
précisément l'image de $P(T)$.

\medskip\noindent
Nous savons par l'algèbre linéaire qu'un endomorphisme agit
de manière nilpotente sur un espace vectoriel de dimension $n$ si et seulement si son
polynôme caractéristique est $T^n$ (car le polynôme caractéristique
divise une puissance du polynôme minimal). Ainsi,
$f$ agit de manière nilpotente sur $A \otimes_R \kappa(\mathfrak p)$ si et
seulement si $\bar P(T) = T^n$. Cela se produit si et seulement si $r_i \in
\mathfrak p$ pour tout $0 \leq i \leq n - 1$, c'est-à-dire lorsque $\mathfrak p \in
V(r_0, \ldots, r_{n - 1}).$
\end{proof}

\begin{lemma}
\label{algebra-lemma-affineline-special}
Soit $R$ un anneau. Soient $f, g \in R[x]$ des polynômes.
Supposons que le coefficient dominant de $g$ soit une unité de $R$.
Il existe des éléments $r_i\in R$, $i = 1\ldots, n$, tels que
l'image de $D(f) \cap V(g)$ dans $\Spec(R)$ soit
$\bigcup_{i = 1, \ldots, n} D(r_i)$.
\end{lemma}

\begin{proof}
Écrivons $g = ux^d + a_{d-1}x^{d-1} + \ldots + a_0$, où
$d$ est le degré de $g$, et donc $u \in R^*$.
Considérons l'anneau $A = R[x]/(g)$.
Comme $R$-module, il est libre de type fini, de base les images
de $1, x, \ldots, x^{d-1}$. Considérons la multiplication
par (l'image de) $f$ sur $A$. C'est un morphisme de $R$-modules.
Nous pouvons donc noter $P(T) \in R[T]$ le polynôme caractéristique
de ce morphisme. Écrivons $P(T) = T^d + r_{d-1} T^{d-1} + \ldots + r_0$.
Nous affirmons que $r_0, \ldots, r_{d-1}$ ont la propriété voulue.
Nous utiliserons ci-dessous la propriété des polynômes caractéristiques
selon laquelle
$$
\mathfrak p \in V(r_0, \ldots, r_{d-1})
\Leftrightarrow
\text{la multiplication par }f\text{ est nilpotente sur }
A \otimes_R \kappa(\mathfrak p).
$$
Elle a été démontrée dans le lemme \ref{algebra-lemma-characteristic-polynomial-prime}.

\medskip\noindent
Supposons que $\mathfrak q\in D(f) \cap V(g)$, et soit
$\mathfrak p = \mathfrak q \cap R$. Il existe alors un morphisme non nul
$A \otimes_R \kappa(\mathfrak p) \to \kappa(\mathfrak q)$ qui
est compatible avec la multiplication par $f$.
Et $f$ agit comme une unité sur $\kappa(\mathfrak q)$.
Nous en concluons que $\mathfrak p \not \in  V(r_0, \ldots, r_{d-1})$.

\medskip\noindent
D'autre part, supposons que $r_i \not\in \mathfrak p$ pour un certain
idéal premier $\mathfrak p$ de $R$ et un certain $0 \leq i \leq d - 1$.
Alors la multiplication par $f$ n'est pas nilpotente sur l'algèbre
$A \otimes_R \kappa(\mathfrak p)$.
Il existe donc un idéal premier $\overline{\mathfrak q} \subset
A \otimes_R \kappa(\mathfrak p)$ qui ne contient pas l'image de $f$.
L'image réciproque de $\overline{\mathfrak q}$ dans $R[x]$
est un élément de $D(f) \cap V(g)$ qui s'envoie sur $\mathfrak p$.
\end{proof}

\begin{theorem}[Théorème de Chevalley]
\label{algebra-theorem-chevalley}
Supposons que $R \to S$ soit de présentation finie.
L'image d'une partie constructible de
$\Spec(S)$ dans $\Spec(R)$ est constructible.
\end{theorem}

\begin{proof}
Écrivons $S = R[x_1, \ldots, x_n]/(f_1, \ldots, f_m)$.
Nous pouvons factoriser $R \to S$ sous la forme $R \to R[x_1] \to R[x_1, x_2]
\to \ldots \to R[x_1, \ldots, x_{n-1}] \to S$. Nous pouvons donc
supposer que $S = R[x]/(f_1, \ldots, f_m)$.
Dans ce cas, nous factorisons le morphisme sous la forme $R \to R[x] \to S$
et, grâce au lemme \ref{algebra-lemma-closed-fp}, nous nous ramenons au
cas $S = R[x]$. D'après le lemme \ref{algebra-lemma-qc-open}, il suffit
de montrer que, si
$T = (\bigcup_{i = 1\ldots n} D(f_i)) \cap V(g_1, \ldots, g_m)$
pour $f_i , g_j \in R[x]$, alors l'image dans $\Spec(R)$ est
constructible. Puisque les réunions finies de parties constructibles
sont constructibles, il suffit de traiter le cas $n = 1$,
c'est-à-dire celui où $T = D(f) \cap V(g_1, \ldots, g_m)$.

\medskip\noindent
Notons que, si $c \in R$, alors
$$
\Spec(R) =
V(c) \amalg D(c) =
\Spec(R/(c)) \amalg \Spec(R_c),
$$
et, de façon correspondante, $\Spec(R[x]) =
V(c) \amalg D(c) = \Spec(R/(c)[x]) \amalg
\Spec(R_c[x])$. L'intersection de $T = D(f) \cap V(g_1, \ldots, g_m)$
avec chaque partie conserve la même forme, où $f$, $g_i$ sont remplacés
par leurs images dans $R/(c)[x]$ et $R_c[x]$, respectivement.
Notons que l'image de $T$
dans $\Spec(R)$ est la réunion des images de
$T \cap V(c)$ et de $T \cap D(c)$. D'après les lemmes \ref{algebra-lemma-open-fp}
et \ref{algebra-lemma-closed-fp}, il suffit de montrer que les images des deux
parties sont constructibles dans $\Spec(R/(c))$ et
$\Spec(R_c)$, respectivement.

\medskip\noindent
Supposons que $T = D(f) \cap V(g_1, \ldots, g_m)$
comme ci-dessus, avec $\deg(g_1) \leq \deg(g_2) \leq \ldots \leq \deg(g_m)$.
Nous allons raisonner par récurrence sur $m$ et sur les
degrés des $g_i$. Soit $d_1 = \deg(g_1)$, c'est-à-dire $g_1 = c x^{d_1} + l.o.t$
avec $c \in R$ non nul. En décomposant $R$ en les deux morceaux
$R/(c)$ et $R_c$, soit nous abaissons le degré de $g_1$ (et ce
cas est couvert par récurrence),
soit nous nous ramenons au cas où $c$ est inversible.
Si $c$ est inversible et si $m > 1$, écrivons
$g_2 = c' x^{d_2} + l.o.t$. Dans ce cas, considérons
$g_2' = g_2 - (c'/c) x^{d_2 - d_1} g_1$. Puisque les idéaux
$(g_1, g_2, \ldots, g_m)$ et $(g_1, g_2', g_3, \ldots, g_m)$
sont égaux, nous voyons que $T = D(f) \cap V(g_1, g_2', g_3\ldots, g_m)$.
Mais ici, le degré de $g_2'$ est strictement inférieur à celui
de $g_2$, et ce cas est donc couvert par récurrence.

\medskip\noindent
Les cas de base de la récurrence ci-dessus sont les cas
(a) $T = D(f) \cap V(g)$ où le coefficient dominant
de $g$ est inversible, et (b) $T = D(f)$. Ces deux cas
sont traités par les lemmes \ref{algebra-lemma-affineline-special}
et \ref{algebra-lemma-affineline-open}.
\end{proof}











\section{Compléments sur les images}
\label{algebra-section-more-images}

\noindent
Dans cette section, nous rassemblons quelques lemmes supplémentaires concernant l'image
au niveau de $\Spec$ pour les morphismes d'anneaux. Voir aussi, par exemple, la section
\ref{algebra-section-going-up}.

\begin{lemma}
\label{algebra-lemma-generic-finite-presentation}
Soit $R \subset S$ une inclusion d'anneaux intègres.
Supposons que $R \to S$ soit de type fini.
Il existe un élément non nul $f \in R$ et un élément non nul $g \in S$
tels que $R_f \to S_{fg}$ soit de présentation finie.
\end{lemma}

\begin{proof}
Nous raisonnons par récurrence sur le nombre de générateurs de $S$ sur $R$.
Au cours de la démonstration, nous pouvons remplacer $R$ par $R_f$ et $S$ par $S_f$
pour un certain élément non nul $f \in R$.

\medskip\noindent
Supposons que $S$ soit engendré par un seul élément
sur $R$. Alors $S = R[x]/\mathfrak q$ pour un certain
idéal premier $\mathfrak q \subset R[x]$. Si $\mathfrak q = (0)$,
il n'y a rien à démontrer. Si $\mathfrak q \not = (0)$,
soit $h \in \mathfrak q$ un élément non nul de degré minimal
en $x$. Écrivons $h = f x^d + a_{d - 1} x^{d - 1} + \ldots + a_0$,
avec $a_i \in R$ et $f \not = 0$. Après avoir inversé $f$
dans $R$ et $S$, nous pouvons supposer que $h$ est unitaire. Nous obtenons
un morphisme surjectif de $R$-algèbres $R[x]/(h) \to S$.
Nous avons $R[x]/(h) = R \oplus Rx \oplus \ldots \oplus Rx^{d - 1}$
comme $R$-module et, par minimalité de $d$, nous voyons que
$R[x]/(h)$ s'injecte dans $S$. Ainsi, $R[x]/(h) \cong S$
est de présentation finie sur $R$.

\medskip\noindent
Supposons que $S$ soit engendré par $n > 1$ éléments sur $R$.
Disons que $x_1, \ldots, x_n \in S$ engendrent $S$. Notons $S' \subset S$
la sous-$R$-algèbre engendrée par $x_1, \ldots, x_{n-1}$. Par l'hypothèse
de récurrence, nous voyons qu'il existe des éléments non nuls $f\in R$ et $g \in S'$
tels que $R_f \to S'_{fg}$ soit de présentation finie.
Nous appliquons ensuite l'hypothèse de récurrence à $S'_{fg} \to S_{fg}$
pour voir qu'il existe $f' \in S'_{fg}$ et
$g' \in S_{fg}$ tels que $S'_{fgf'} \to S_{fgf'g'}$
soit de présentation finie. Nous laissons au lecteur le soin de conclure.
\end{proof}

\begin{lemma}
\label{algebra-lemma-characterize-image-finite-type}
Soit $R \to S$ un morphisme d'anneaux de type fini.
Notons $X = \Spec(R)$ et $Y = \Spec(S)$.
Notons $f : Y \to X$ le morphisme induit
sur les spectres. Soit $E \subset Y = \Spec(S)$ une
partie constructible.
Si un point $\xi \in X$ appartient à $f(E)$, alors
$\overline{\{\xi\}} \cap f(E)$ contient un sous-ensemble ouvert
dense de $\overline{\{\xi\}}$.
\end{lemma}

\begin{proof}
Soit $\xi \in X$ un point de $f(E)$. Choisissons un point $\eta \in E$
qui s'envoie sur $\xi$. Soit $\mathfrak p \subset R$ l'idéal premier
correspondant à $\xi$, et soit $\mathfrak q \subset S$ l'idéal premier
correspondant à $\eta$. Considérons le diagramme
$$
\xymatrix{
\eta \ar[r] \ar@{|->}[d] & E \cap Y' \ar[r] \ar[d] &
Y' = \Spec(S/\mathfrak q) \ar[r] \ar[d] &
Y \ar[d] \\
\xi \ar[r] & f(E) \cap X' \ar[r] &
X' = \Spec(R/\mathfrak p) \ar[r] &
X
}
$$
D'après le lemme \ref{algebra-lemma-affine-map-quasi-compact}, la partie $E \cap Y'$
est constructible dans $Y'$.
Il s'ensuit que nous pouvons remplacer $X$ par $X'$ et
$Y$ par $Y'$. Nous pouvons donc supposer que $R \subset S$ est une
inclusion d'anneaux intègres, que $\xi$ est le point générique
de $X$ et que $\eta$ est le point générique de $Y$.
En combinant le lemme \ref{algebra-lemma-generic-finite-presentation}
avec le théorème de Chevalley
(théorème \ref{algebra-theorem-chevalley}),
nous voyons qu'il existe des ouverts denses $U \subset X$ et
$V \subset Y$ tels que $f(V) \subset U$ et
que $f : V \to U$ envoie les parties constructibles
sur des parties constructibles. Notons que $E \cap V$ est
constructible dans $V$, voir Topologie,
lemme \ref{topology-lemma-open-immersion-constructible-inverse-image}.
Ainsi, $f(E \cap V)$ est constructible dans $U$ et contient $\xi$.
D'après Topologie, lemme \ref{topology-lemma-generic-point-in-constructible},
nous voyons que $f(E \cap V)$ contient un ouvert dense $U' \subset U$.
\end{proof}

\noindent
À la fin de cette section, nous présentons quelques résultats supplémentaires sur les
images des applications entre spectres qui n'ont rien à voir avec les parties
constructibles.

\begin{lemma}
\label{algebra-lemma-surjective-spec-radical-ideal}
Soit $\varphi : R \to S$ un morphisme d'anneaux.
Les conditions suivantes sont équivalentes :
\begin{enumerate}
\item L'application $\Spec(S) \to \Spec(R)$ est surjective.
\item Pour tout idéal $I \subset R$,
l'image réciproque de $\sqrt{IS}$ dans $R$ est égale à $\sqrt{I}$.
\item Pour tout idéal radical $I \subset R$, l'image réciproque
de $IS$ dans $R$ est égale à $I$.
\item Pour tout idéal premier $\mathfrak p$ de $R$, l'image
réciproque de $\mathfrak p S$ dans $R$ est $\mathfrak p$.
\end{enumerate}
Dans ce cas, il en va de même après tout changement de base : étant donné un morphisme d'anneaux
$R \to R'$, le morphisme d'anneaux $R' \to R' \otimes_R S$ possède lui aussi les
propriétés équivalentes (1), (2), (3).
\end{lemma}

\begin{proof}
Si $J \subset S$ est un idéal, alors
$\sqrt{\varphi^{-1}(J)} = \varphi^{-1}(\sqrt{J})$. Cela montre que (2)
et (3) sont équivalentes.
L'implication (3) $\Rightarrow$ (4) est immédiate.
Si $I \subset R$ est un idéal radical, alors le
lemme \ref{algebra-lemma-Zariski-topology}
garantit que $I = \bigcap_{I \subset \mathfrak p} \mathfrak p$.
Ainsi, (4) $\Rightarrow$ (2). D'après le
lemme \ref{algebra-lemma-in-image},
nous avons $\mathfrak p = \varphi^{-1}(\mathfrak p S)$ si et seulement si
$\mathfrak p$ appartient à l'image. Ainsi, (1) $\Leftrightarrow$ (4).
Les conditions (1), (2), (3) et (4) sont donc équivalentes.

\medskip\noindent
Supposons que (1) soit vérifiée. Soit $R \to R'$ un morphisme d'anneaux. Soit
$\mathfrak p' \subset R'$ un idéal premier au-dessus de l'idéal premier
$\mathfrak p$ de $R$. Pour voir que $\mathfrak p'$ appartient à l'image
de $\Spec(R' \otimes_R S) \to \Spec(R')$, nous devons montrer
que $(R' \otimes_R S) \otimes_{R'} \kappa(\mathfrak p')$ n'est pas nul, voir le
lemme \ref{algebra-lemma-in-image}.
Mais nous avons
$$
(R' \otimes_R S) \otimes_{R'} \kappa(\mathfrak p') =
S \otimes_R \kappa(\mathfrak p)
\otimes_{\kappa(\mathfrak p)} \kappa(\mathfrak p')
$$
qui n'est pas nul, car $S \otimes_R \kappa(\mathfrak p)$ n'est pas nul
par hypothèse et $\kappa(\mathfrak p) \to \kappa(\mathfrak p')$ est
une extension de corps.
\end{proof}

\begin{lemma}
\label{algebra-lemma-domain-image-dense-set-points-generic-point}
Soit $R$ un anneau intègre. Soit $\varphi : R \to S$ un morphisme d'anneaux.
Les conditions suivantes sont équivalentes :
\begin{enumerate}
\item Le morphisme d'anneaux $R \to S$ est injectif.
\item L'image de $\Spec(S) \to \Spec(R)$
contient un ensemble dense de points.
\item Il existe un idéal premier $\mathfrak q \subset S$
dont l'image réciproque dans $R$ est $(0)$.
\end{enumerate}
\end{lemma}

\begin{proof}
Soit $K$ le corps des fractions de l'anneau intègre $R$.
Supposons que $R \to S$ soit injectif. Puisque la localisation
est exacte, nous voyons que $K \to S \otimes_R K$ est injectif.
Il existe donc un idéal premier s'envoyant sur $(0)$ d'après le
lemme \ref{algebra-lemma-in-image}.

\medskip\noindent
Notons que $(0)$ est dense dans $\Spec(R)$, de sorte que la
dernière condition implique la seconde.

\medskip\noindent
Supposons que la seconde condition soit vérifiée. Soit $f \in R$,
$f \not = 0$. Puisque $R$ est un anneau intègre, nous voyons que $V(f)$
est un fermé propre du spectre de $R$. Par hypothèse,
il existe un idéal premier $\mathfrak q$
de $S$ tel que $\varphi(f) \not \in \mathfrak q$.
Ainsi, $\varphi(f) \not = 0$.
Le morphisme $R \to S$ est donc injectif.
\end{proof}

\begin{lemma}
\label{algebra-lemma-injective-minimal-primes-in-image}
Soit $R \subset S$ un morphisme d'anneaux injectif.
Alors $\Spec(S) \to \Spec(R)$
atteint tous les idéaux premiers minimaux.
\end{lemma}

\begin{proof}
Soit $\mathfrak p \subset R$ un idéal premier minimal.
Dans ce cas, $R_{\mathfrak p}$ possède un unique idéal premier.
Il suffit donc de montrer que $S_{\mathfrak p}$ n'est pas nul.
Cela résulte du fait que la localisation est exacte,
voir la proposition \ref{algebra-proposition-localization-exact}.
\end{proof}

\begin{lemma}
\label{algebra-lemma-image-dense-generic-points}
Soit $R \to S$ un morphisme d'anneaux. Les conditions suivantes sont équivalentes :
\begin{enumerate}
\item Le noyau de $R \to S$ est constitué d'éléments nilpotents.
\item Les idéaux premiers minimaux de $R$ appartiennent à l'image de
$\Spec(S) \to \Spec(R)$.
\item L'image de $\Spec(S) \to \Spec(R)$ est dense
dans $\Spec(R)$.
\end{enumerate}
\end{lemma}

\begin{proof}
Soit $I = \Ker(R \to S)$. Notons que
$\sqrt{(0)} = \bigcap_{\mathfrak q \subset S} \mathfrak q$, voir le
lemme \ref{algebra-lemma-Zariski-topology}.
Ainsi, $\sqrt{I} = \bigcap_{\mathfrak q \subset S} R \cap \mathfrak q$.
Donc $V(I) = V(\sqrt{I})$ est l'adhérence de l'image de
$\Spec(S) \to \Spec(R)$.
Cela montre que (1) équivaut à (3). Il est clair que
(2) implique (3). Enfin, supposons (1). Nous pouvons remplacer
$R$ par $R/I$ et $S$ par $S/IS$ sans changer la topologie
des spectres ni l'application. Ainsi, l'implication (1) $\Rightarrow$ (2)
résulte du lemme \ref{algebra-lemma-injective-minimal-primes-in-image}.
\end{proof}

\begin{lemma}
\label{algebra-lemma-minimal-prime-image-minimal-prime}
Soit $R \to S$ un morphisme d'anneaux. Si un idéal premier minimal $\mathfrak p \subset R$
appartient à l'image de $\Spec(S) \to \Spec(R)$, alors il est l'image
d'un idéal premier minimal.
\end{lemma}

\begin{proof}
Écrivons $\mathfrak p = \mathfrak q \cap R$. Choisissons alors un
idéal premier minimal $\mathfrak r \subset S$ tel que $\mathfrak r \subset \mathfrak q$, voir le
lemme \ref{algebra-lemma-Zariski-topology}.
Par minimalité de $\mathfrak p$, nous voyons que
$\mathfrak p = \mathfrak r \cap R$.
\end{proof}

\begin{lemma}
\label{algebra-lemma-intersection-not-zero}
Soit $A \subset B$ une inclusion d'anneaux intègres induisant une extension algébrique
des corps des fractions. Si $J \subset B$ est un idéal non nul, alors $A \cap J$
est également non nul. Ainsi, l'image d'un fermé propre de $\Spec(B)$
n'est pas dense dans $\Spec(A)$.
\end{lemma}

\begin{proof}
Soit $x \in J$ un élément non nul. Puisque $x$ est algébrique sur le corps des
fractions de $A$, il existe un $d \geq 1$ et des $a_0, \ldots, a_d \in A$
tels que $a_0, a_d \not = 0$ et que
$a_d x^d + a_{d - 1} x^{d - 1} + \ldots + a_0 = 0$ dans $B$.
Alors $a_0 \in A \cap J$.
\end{proof}




\section{Anneaux noethériens}
\label{algebra-section-Noetherian}

\noindent
Un anneau $R$ est {\it noethérien} si tout idéal de $R$ est de type fini.
Cela équivaut clairement à la condition de chaîne ascendante pour les idéaux de $R$.
D'après le
lemme \ref{algebra-lemma-cohen},
il suffit de vérifier que tout idéal premier de $R$ est de type fini.

\begin{lemma}
\label{algebra-lemma-Noetherian-permanence}
\begin{slogan}
La propriété noethérienne est stable par extension de type fini
et par localisation.
\end{slogan}
Tout anneau de type fini sur un anneau noethérien
est noethérien. Toute localisation d'un anneau noethérien
est noethérienne.
\end{lemma}

\begin{proof}
L'assertion sur les localisations résulte du fait
que tout idéal $J \subset S^{-1}R$ est de la forme
$I \cdot S^{-1}R$. Tout quotient $R/I$ d'un anneau noethérien
$R$ est noethérien, car tout idéal $\overline{J} \subset R/I$
est de la forme $J/I$ pour un certain idéal $I \subset J \subset R$.
Il suffit donc de montrer que, si $R$ est noethérien, alors
$R[X]$ l'est aussi. Supposons que $J_1 \subset J_2 \subset \ldots$ soit une
chaîne ascendante d'idéaux de $R[X]$. Considérons les idéaux $I_{i, d}$
définis comme les idéaux constitués des éléments de $R$ qui apparaissent comme coefficients
dominants de polynômes de degré $d$ appartenant à $J_i$.
Il est clair que $I_{i, d} \subset I_{i', d'}$ dès que
$i \leq i'$ et $d \leq d'$. Par la condition de chaîne ascendante
dans $R$, il n'y a qu'un nombre fini d'idéaux distincts parmi tous les
$I_{i, d}$.
(Indication : tout ensemble infini d'éléments de
$\mathbf{N} \times \mathbf{N}$ contient une suite infinie
croissante.)
Prenons $i_0$ assez grand pour que $I_{i, d} = I_{i_0, d}$
pour tous $i \geq i_0$ et tout $d$. Supposons que $f \in J_i$ pour un certain $i \geq i_0$.
Par récurrence sur le degré $d = \deg(f)$, montrons que $f \in J_{i_0}$.
En effet, il existe un $g\in J_{i_0}$ de degré $d$ ayant
le même coefficient dominant que $f$. Par récurrence,
$f - g \in J_{i_0}$, et le résultat s'ensuit.
\end{proof}

\begin{lemma}
\label{algebra-lemma-Noetherian-power-series}
Si $R$ est un anneau noethérien, alors l'anneau de séries
formelles $R[[x_1, \ldots, x_n]]$ l'est également.
\end{lemma}

\begin{proof}
Puisque $R[[x_1, \ldots, x_{n + 1}]] \cong R[[x_1, \ldots, x_n]][[x_{n + 1}]]$,
il suffit de montrer que $R[[x]]$ est noethérien lorsque
$R$ l'est. Soit $I \subset R[[x]]$ un idéal.
Nous devons montrer que $I$ est un idéal de type fini.
Pour tout entier naturel
$d$, notons $I_d = \{a \in R \mid ax^d + \text{termes d'ordre supérieur} \in I\}$.
Alors nous voyons que $I_0 \subset I_1 \subset \ldots$ se stabilise, puisque $R$
est noethérien. Choisissons $d_0$ tel que $I_{d_0} = I_{d_0 + 1} = \ldots$.
Pour tout $d \leq d_0$, choisissons des éléments $f_{d, j} \in I \cap (x^d)$,
$j = 1, \ldots, n_d$, tels que, si nous écrivons
$f_{d, j} = a_{d, j}x^d + \text{termes d'ordre supérieur}$, alors $I_d = (a_{d, j})$.
Notons $I' = (\{f_{d, j}\}_{d = 0, \ldots, d_0, j = 1, \ldots, n_d})$.
Il est alors clair que $I' \subset I$. Choisissons $f \in I$.
Nous pouvons d'abord choisir des $c_{d, i} \in R$ tels que
$$
f - \sum c_{d, i} f_{d, i} \in (x^{d_0 + 1}) \cap I.
$$
Ensuite, nous pouvons choisir $c_{i, 1} \in R$, $i = 1, \ldots, n_{d_0}$, tels que
$$
f - \sum c_{d, i} f_{d, i} - \sum c_{i, 1}xf_{d_0, i} \in (x^{d_0 + 2}) \cap I.
$$
Ensuite, nous pouvons choisir $c_{i, 2} \in R$, $i = 1, \ldots, n_{d_0}$, tels que
$$
f - \sum c_{d, i} f_{d, i} - \sum c_{i, 1}xf_{d_0, i}
- \sum c_{i, 2}x^2f_{d_0, i}
\in (x^{d_0 + 3}) \cap I.
$$
Et ainsi de suite. Finalement, nous voyons que
$$
f = \sum c_{d, i} f_{d, i} +
\sum\nolimits_i (\sum\nolimits_e c_{i, e} x^e)f_{d_0, i}
$$
appartient à $I'$, comme voulu.
\end{proof}

\noindent
Le lemme suivant, quoique facile, est utile, car les
$\mathbf{Z}$-algèbres de type fini apparaissent assez souvent dans
une technique appelée « réduction noethérienne absolue ».

\begin{lemma}
\label{algebra-lemma-obvious-Noetherian}
Toute algèbre de type fini sur un corps est noethérienne.
Toute algèbre de type fini sur $\mathbf{Z}$ est noethérienne.
\end{lemma}

\begin{proof}
Cela résulte immédiatement du lemme \ref{algebra-lemma-Noetherian-permanence}
et du fait que les corps sont des anneaux noethériens et que
$\mathbf{Z}$ est un anneau noethérien (car c'est un
anneau principal).
\end{proof}

\begin{lemma}
\label{algebra-lemma-Noetherian-finite-type-is-finite-presentation}
Soit $R$ un anneau noethérien.
\begin{enumerate}
\item Tout $R$-module de type fini est de présentation finie.
\item Tout sous-module d'un $R$-module de type fini est de type fini.
\item Toute $R$-algèbre de type fini est de présentation finie sur $R$.
\end{enumerate}
\end{lemma}

\begin{proof}
Soit $M$ un $R$-module de type fini. D'après le
lemme \ref{algebra-lemma-trivial-filter-finite-module},
nous pouvons trouver une filtration finie de $M$ dont les quotients successifs sont
de la forme $R/I$. Puisque tout idéal est de type fini, chacun des
quotients $R/I$ est de présentation finie. Ainsi, $M$ est de présentation
finie d'après le
lemme \ref{algebra-lemma-extension}.
Cela démontre (1).

\medskip\noindent
Soit $N \subset M$ un sous-module. Puisque $M$ est de type fini, le quotient
$M/N$ est de type fini. Ainsi, $M/N$ est de présentation finie d'après la partie (1).
Nous voyons donc que $N$ est de type fini d'après le lemme \ref{algebra-lemma-extension}, partie (5).
Cela démontre la partie (2).

\medskip\noindent
Pour voir (3), notons que tout idéal de
$R[x_1, \ldots, x_n]$ est de type fini d'après le
lemme \ref{algebra-lemma-Noetherian-permanence}.
\end{proof}

\begin{lemma}
\label{algebra-lemma-Noetherian-topology}
Si $R$ est un anneau noethérien, alors $\Spec(R)$
est un espace topologique noethérien, voir Topologie,
définition \ref{topology-definition-noetherian}.
\end{lemma}

\begin{proof}
Cela vient du fait que tout fermé de $\Spec(R)$
s'écrit de manière unique sous la forme $V(I)$, où $I$ est un idéal radical,
voir le lemme \ref{algebra-lemma-Zariski-topology}.
Et cette correspondance renverse les inclusions.
Le résultat découle donc des définitions.
\end{proof}

\begin{lemma}
\label{algebra-lemma-Noetherian-irreducible-components}
\begin{slogan}
Un schéma affine noethérien possède un nombre fini de points génériques.
\end{slogan}
Si $R$ est un anneau noethérien, alors $\Spec(R)$
possède un nombre fini de composantes irréductibles. Autrement dit,
$R$ possède un nombre fini d'idéaux premiers minimaux.
\end{lemma}

\begin{proof}
D'après le lemme \ref{algebra-lemma-Noetherian-topology} et
Topologie, lemme \ref{topology-lemma-Noetherian},
nous voyons qu'il n'y a qu'un nombre fini de composantes irréductibles.
D'après le lemme \ref{algebra-lemma-irreducible}, celles-ci correspondent aux
idéaux premiers minimaux de $R$.
\end{proof}

\begin{lemma}
\label{algebra-lemma-Noetherian-base-change-finite-type}
Soit $R \to S$ un morphisme d'anneaux. Supposons que $R \to R'$ soit de type fini.
Si $S$ est noethérien, alors le changement de base $S' = R' \otimes_R S$
est noethérien.
\end{lemma}

\begin{proof}
D'après le lemme \ref{algebra-lemma-base-change-finiteness}, le caractère de type fini est stable par
changement de base. Ainsi, $S \to S'$ est de type fini. Puisque $S$ est noethérien, nous
pouvons appliquer le lemme \ref{algebra-lemma-Noetherian-permanence}.
\end{proof}

\begin{lemma}
\label{algebra-lemma-Noetherian-field-extension}
Soit $k$ un corps et soit $R$ une $k$-algèbre noethérienne.
Si $K/k$ est une extension de corps de type fini, alors
$K \otimes_k R$ est noethérien.
\end{lemma}

\begin{proof}
Puisque $K/k$ est une extension de corps de type fini, il existe
une $k$-algèbre de type fini $B \subset K$ telle que $K$ soit
le corps des fractions de $B$. Autrement dit, $K = S^{-1}B$,
où $S = B \setminus \{0\}$. Alors $K \otimes_k R = S^{-1}(B \otimes_k R)$.
Puis, $B \otimes_k R$ est noethérien d'après le
lemme \ref{algebra-lemma-Noetherian-base-change-finite-type}.
Enfin, $K \otimes_k R = S^{-1}(B \otimes_k R)$ est noethérien d'après le
lemme \ref{algebra-lemma-Noetherian-permanence}.
\end{proof}

\noindent
Voici quelques lemmes amusants qui sont parfois utiles.

\begin{lemma}
\label{algebra-lemma-subring-of-local-ring}
Soient $R$ un anneau et $\mathfrak p \subset R$ un idéal premier.
Il existe un $f \in R$, $f \not \in \mathfrak p$, tel
que $R_f \to R_\mathfrak p$ soit injectif dans chacun des
cas suivants :
\begin{enumerate}
\item $R$ est un anneau intègre,
\item $R$ est noethérien, ou
\item $R$ est réduit et possède un nombre fini d'idéaux premiers minimaux.
\end{enumerate}
\end{lemma}

\begin{proof}
Si $R$ est un anneau intègre, alors $R \subset R_\mathfrak p$, et donc $f = 1$ convient.
Si $R$ est noethérien, alors le noyau $I$ de $R \to R_\mathfrak p$
est un idéal de type fini, et nous pouvons trouver
$f \in R$, $f \not \in \mathfrak p$, tel que $IR_f = 0$.
Pour ce $f$, le morphisme $R_f \to R_\mathfrak p$ est injectif,
et $f$ convient. Si $R$ est réduit et possède un nombre
fini d'idéaux premiers minimaux $\mathfrak p_1, \ldots, \mathfrak p_n$,
alors nous pouvons choisir
$f \in \bigcap_{\mathfrak p_i \not \subset \mathfrak p} \mathfrak p_i$,
$f \not \in \mathfrak p$. En effet, si $\mathfrak{p}_i\not\subset
\mathfrak{p}$, alors il existe $f_i \in \mathfrak{p}_i$,
$f_i \not\in \mathfrak{p}$, et $f = \prod f_i$ convient.
Pour ce $f$, nous avons $R_f \subset R_\mathfrak p$, car les idéaux premiers
minimaux de $R_f$ correspondent aux idéaux premiers minimaux de $R_\mathfrak p$
et nous pouvons appliquer le lemme \ref{algebra-lemma-reduced-ring-sub-product-fields}
(certains détails sont omis).
\end{proof}

\begin{lemma}
\label{algebra-lemma-surjective-endo-noetherian-ring-is-iso}
Tout endomorphisme surjectif d'un anneau noethérien est un isomorphisme.
\end{lemma}

\begin{proof}
Si $f : R \to R$ était un tel endomorphisme sans être injectif, alors
$$
\Ker(f) \subset \Ker(f \circ f) \subset
\Ker(f \circ f \circ f) \subset \ldots
$$
serait une chaîne strictement croissante d'idéaux.
\end{proof}







\section{Idéaux localement nilpotents}
\label{algebra-section-locally-nilpotent}

\noindent
Voici la définition.

\begin{definition}
\label{algebra-definition-locally-nilpotent-ideal}
Soit $R$ un anneau. Soit $I \subset R$ un idéal.
Nous disons que $I$ est {\it localement nilpotent} si, pour tout
$x \in I$, il existe un $n \in \mathbf{N}$ tel
que $x^n = 0$. Nous disons que $I$ est {\it nilpotent} s'il
existe un $n \in \mathbf{N}$ tel que $I^n = 0$.
\end{definition}

\begin{example}
\label{algebra-example-locally-nilpotent-not-nilpotent}
Soit $R = k[x_n | n \in \mathbf{N}]$ l'anneau de polynômes en une infinité
de variables sur un corps $k$. Soit $I$ l'idéal engendré par
les éléments $x_n^n$ pour $n \in \mathbf{N}$, et soit $S = R/I$. Alors l'idéal
$J \subset S$ engendré par les images des $x_n$, $n \in \mathbf{N}$,
est localement nilpotent, mais n'est pas nilpotent. En effet, puisque les combinaisons
$S$-linéaires d'éléments nilpotents sont nilpotentes, pour prouver que $J$ est localement
nilpotent, il suffit d'observer que tous ses générateurs sont nilpotents
(ce qu'ils sont manifestement). D'autre part, pour tout $n \in \mathbf{N}$,
nous avons $x_{n + 1}^n \not \in I$, de sorte que $J^n \not = 0$.
Il s'ensuit que $J$ n'est pas nilpotent.
\end{example}

\begin{lemma}
\label{algebra-lemma-locally-nilpotent}
Soit $R \to R'$ un morphisme d'anneaux et soit $I \subset R$ un idéal localement
nilpotent. Alors $IR'$ est un idéal localement nilpotent de $R'$.
\end{lemma}

\begin{proof}
Cela résulte du fait que, si $x, y \in R'$ sont nilpotents, alors
$x + y$ l'est aussi. En effet, si $x^n = 0$ et $y^m = 0$, alors
$(x + y)^{n + m - 1} = 0$.
\end{proof}

\begin{lemma}
\label{algebra-lemma-locally-nilpotent-unit}
Soit $R$ un anneau et soit $I \subset R$ un idéal localement
nilpotent.
Un élément $x$ de $R$ est une unité si et seulement si l'image de $x$
dans $R/I$ est une unité.
\end{lemma}

\begin{proof}
Si $x$ est une unité de $R$, alors son image est clairement une unité de $R/I$.
Il reste à démontrer la réciproque.
Supposons que l'image de $y \in R$ dans $R/I$ soit l'inverse de l'image de $x$.
Alors $xy = 1 - z$ pour un certain $z \in I$.
Cela signifie que $1\equiv z$ modulo $xR$.
Puisque $z$ appartient à l'idéal localement nilpotent
$I$, nous avons $z^N = 0$ pour un certain $N$ assez grand.
Il s'ensuit que $1 = 1^N \equiv z^N = 0$ modulo $xR$.
Autrement dit, $x$ divise $1$ et est donc une unité.
\end{proof}

\begin{lemma}
\label{algebra-lemma-Noetherian-power}
\begin{slogan}
Un idéal d'un anneau noethérien est nilpotent si chacun de ses éléments
est nilpotent.
\end{slogan}
Soit $R$ un anneau noethérien. Soient $I, J$ des idéaux de $R$.
Supposons que $J \subset \sqrt{I}$. Alors $J^n \subset I$ pour un certain $n$.
En particulier, dans un anneau noethérien, les notions d'
« idéal localement nilpotent »
et d'« idéal nilpotent » coïncident.
\end{lemma}

\begin{proof}
Écrivons $J = (f_1, \ldots, f_s)$.
Par hypothèse, $f_i^{d_i} \in I$.
Prenons $n = d_1 + d_2 + \ldots + d_s + 1$.
\end{proof}

\begin{lemma}
\label{algebra-lemma-lift-idempotents}
Soit $R$ un anneau. Soit $I \subset R$ un idéal localement nilpotent.
Alors $R \to R/I$ induit une bijection sur les idempotents.
\end{lemma}

\begin{proof}[Première démonstration du lemme \ref{algebra-lemma-lift-idempotents}]
Puisque $I$ est localement nilpotent, il est contenu dans tout idéal premier.
Ainsi, $\Spec(R/I) = V(I) = \Spec(R)$. Le
lemme résulte donc du lemme \ref{algebra-lemma-disjoint-decomposition}.
\end{proof}

\begin{proof}[Seconde démonstration du lemme \ref{algebra-lemma-lift-idempotents}]
Supposons que $\overline{e} \in R/I$ soit un idempotent.
Nous devons relever $\overline{e}$ en un idempotent de $R$.

\medskip\noindent
Choisissons d'abord un relèvement quelconque $f \in R$ de $\overline{e}$ et posons
$x = f^2 - f$. Alors $x \in I$, donc $x$ est nilpotent (puisque $I$
est localement nilpotent). Soit maintenant $J$ l'idéal de $R$ engendré
par $x$. Alors $J$ est nilpotent (et pas seulement localement nilpotent),
puisqu'il est engendré par l'élément nilpotent $x$.

\medskip\noindent
Supposons maintenant que nous ayons trouvé un relèvement $e \in R$ de $\overline{e}$
tel que $e^2 - e \in J^k$ pour un certain $k \geq 1$.
Soit $e' = e - (2e - 1)(e^2 - e) = 3e^2 - 2e^3$, qui est un autre
relèvement de $\overline{e}$ (car l'idempotence de $\overline{e}$
donne $e^2 - e \in I$). Alors
$$
(e')^2 - e' = (4e^2 - 4e - 3)(e^2 - e)^2 \in J^{2k}
$$
par un calcul direct.

\medskip\noindent
Nous sommes donc partis d'un relèvement $e$ de $\overline{e}$ tel
que $e^2 - e \in J^k$, et nous avons obtenu un relèvement $e'$ de
$\overline{e}$ tel que $(e')^2 - e' \in J^{2k}$.
Nous pouvons ainsi améliorer successivement l'approximation
(en commençant par $e = f$, qui convient pour $k = 1$).
Nous finissons par atteindre un stade où $J^k = 0$, et à ce
stade nous disposons d'un relèvement $e$ de $\overline{e}$ tel que
$e^2 - e \in J^k = 0$, c'est-à-dire que cet $e$ est idempotent.

\medskip\noindent
Nous avons ainsi vu que, si $\overline{e} \in R/I$ est un
idempotent quelconque, il existe un relèvement de $\overline{e}$
qui est un idempotent de $R$.
Il reste à montrer que ce relèvement est unique. En effet, soient
$e_1$ et $e_2$ deux tels relèvements. Nous
devons montrer que $e_1 = e_2$.

\medskip\noindent
Par définition de $e_1$ et de $e_2$, nous avons $e_1 \equiv e_2
\mod I$, et $e_1$ et $e_2$ sont tous deux idempotents. De
$e_1 \equiv e_2 \mod I$, nous déduisons que $e_1 - e_2 \in I$,
de sorte que $e_1 - e_2$ est nilpotent (puisque $I$ est localement nilpotent).
Un calcul direct
(utilisant l'idempotence de $e_1$ et de $e_2$)
montre que $(e_1 - e_2)^3 = e_1 - e_2$. En utilisant cela et
une récurrence, nous obtenons $(e_1 - e_2)^k = e_1 - e_2$ pour tout
entier impair positif $k$. Comme tout $k$ assez grand vérifie
$(e_1 - e_2)^k = 0$ (puisque $e_1 - e_2$ est nilpotent),
cela montre que $e_1 - e_2 = 0$, et donc que $e_1 = e_2$, ce qui
achève la démonstration.
\end{proof}

\begin{lemma}
\label{algebra-lemma-lift-idempotents-noncommutative}
Soit $A$ une algèbre éventuellement non commutative.
Soit $e \in A$ un élément tel que $x = e^2 - e$ soit nilpotent.
Il existe alors un idempotent de la forme
$e' = e + x(\sum a_{i, j}e^ix^j) \in A$,
avec $a_{i, j} \in \mathbf{Z}$.
\end{lemma}

\begin{proof}
Considérons l'anneau $R_n = \mathbf{Z}[e]/((e^2 - e)^n)$. Il est clair que,
si nous pouvons démontrer le résultat pour chaque $R_n$, le lemme en résulte.
Dans $R_n$, considérons l'idéal $I = (e^2 - e)$ et appliquons le
lemme \ref{algebra-lemma-lift-idempotents}.
\end{proof}

\begin{lemma}
\label{algebra-lemma-lift-nth-roots}
Soit $R$ un anneau. Soit $I \subset R$ un idéal localement nilpotent.
Soit $n \geq 1$ un entier inversible dans $R/I$. Alors
\begin{enumerate}
\item l'application d'élévation à la puissance $n$-ième $1 + I \to 1 + I$, $1 + x \mapsto (1 + x)^n$
est une bijection,
\item une unité de $R$ est une puissance $n$-ième si et seulement si son image dans $R/I$
est une puissance $n$-ième.
\end{enumerate}
\end{lemma}

\begin{proof}
Soit $a \in R$ une unité dont l'image dans $R/I$ est égale à l'image
de $b^n$, avec $b \in R$. Alors $b$ est une unité
(lemme \ref{algebra-lemma-locally-nilpotent-unit}) et
$ab^{-n} = 1 + x$ pour un certain $x \in I$. Ainsi, $ab^{-n} = c^n$ d'après
la partie (1). La partie (2) résulte donc de la partie (1).

\medskip\noindent
Démontrons (1). Cela est vrai parce que l'application
$1 + x \mapsto (1 + x)^n$ possède une inverse. Nous pouvons en effet considérer l'application
qui envoie $1 + x$ sur
\begin{align*}
(1 + x)^{1/n}
& =
1 + {1/n \choose 1}x +
{1/n \choose 2}x^2 +
{1/n \choose 3}x^3 + \ldots \\
& =
1 + \frac{1}{n} x + \frac{1 - n}{2n^2}x^2 +
\frac{(1 - n)(1 - 2n)}{6n^3}x^3 + \ldots
\end{align*}
comme en calcul élémentaire. Cette expression a un sens, car la série est finie,
puisque $x^k = 0$ pour tout $k \gg 0$, et que chaque coefficient
$ {1/n \choose k} \in \mathbf{Z}[1/n]$ (détails omis ; observons que
$n$ est inversible dans $R$ d'après le lemme \ref{algebra-lemma-locally-nilpotent-unit}).
\end{proof}






\section{Curiosité}
\label{algebra-section-curiosity}

\noindent
Le lemme \ref{algebra-lemma-disjoint-implies-product}
explique ce qui se passe lorsque $V(I)$ est ouvert pour un certain idéal $I \subset R$.
Mais qu'en est-il si $\Spec(S^{-1}R)$ est fermé dans $\Spec(R)$ ?
Les deux lemmes suivants donnent une réponse partielle. Pour plus d'informations, voir la
section \ref{algebra-section-pure-ideals}.

\begin{lemma}
\label{algebra-lemma-invert-closed-quotient}
Soit $R$ un anneau. Soit $S \subset R$ une partie multiplicative.
Supposons que l'image du morphisme $\Spec(S^{-1}R) \to \Spec(R)$
soit fermée. Alors $S^{-1}R \cong R/I$ pour un certain idéal $I \subset R$.
\end{lemma}

\begin{proof}
Soit $I = \Ker(R \to S^{-1}R)$, de sorte que $V(I)$ contient l'image.
Disons que l'image est le fermé $V(I') \subset \Spec(R)$ pour
un certain idéal $I' \subset R$. Ainsi, $V(I') \subset V(I)$.
Pour $f \in I'$, nous voyons que $f/1 \in S^{-1}R$
appartient à tout idéal premier. Ainsi, $f^n$ s'envoie sur zéro dans $S^{-1}R$
pour un certain $n \geq 1$ (lemme \ref{algebra-lemma-Zariski-topology}).
Par conséquent, $V(I') = V(I)$.
Cela implique que tout $g \in S$ est inversible modulo $I$.
Nous obtenons donc des morphismes d'anneaux $R/I \to S^{-1}R$ et $S^{-1}R \to R/I$.
Le premier morphisme est injectif par le choix de $I$.
Le second est le morphisme $S^{-1}R \to S^{-1}(R/I) = R/I$, qui
a pour noyau $S^{-1}I$, car la localisation est exacte.
Puisque $S^{-1}I = 0$, nous voyons que le second morphisme est lui aussi injectif.
Ainsi, $S^{-1}R \cong R/I$.
\end{proof}

\begin{lemma}
\label{algebra-lemma-invert-closed-split}
Soit $R$ un anneau. Soit $S \subset R$ une partie multiplicative.
Supposons que l'image du morphisme $\Spec(S^{-1}R) \to \Spec(R)$
soit fermée. Si $R$ est noethérien, ou si $\Spec(R)$ est un
espace topologique noethérien, ou si $S$ est de type fini comme monoïde,
alors $R \cong S^{-1}R \times R'$ pour un certain anneau $R'$.
\end{lemma}

\begin{proof}
D'après le lemme \ref{algebra-lemma-invert-closed-quotient}, nous avons $S^{-1}R \cong R/I$
pour un certain idéal $I \subset R$. D'après le lemme \ref{algebra-lemma-disjoint-implies-product},
il suffit de montrer que $V(I)$ est ouvert.
Si $R$ est noethérien, alors $\Spec(R)$ est un espace topologique
noethérien, voir le lemme \ref{algebra-lemma-Noetherian-topology}.
Si $\Spec(R)$ est un espace topologique noethérien,
alors le complémentaire $\Spec(R) \setminus V(I)$ est quasi-compact, voir
Topologie, lemme \ref{topology-lemma-Noetherian-quasi-compact}.
Il existe donc un nombre fini d'éléments $f_1, \ldots, f_n \in I$ tels
que $V(I) = V(f_1, \ldots, f_n)$.
Puisque chaque $f_i$ s'envoie sur zéro dans $S^{-1}R$,
il existe un $g \in S$ tel que $gf_i = 0$ pour
$i = 1, \ldots, n$. Ainsi, $D(g) = V(I)$, comme voulu.
Dans le cas où $S$ est de type fini comme monoïde, disons que $S$ est engendré
par $g_1, \ldots, g_m$, nous avons $S^{-1}R \cong R_{g_1 \ldots g_m}$,
et nous concluons que $V(I) = D(g_1 \ldots g_m)$.
\end{proof}














\section{Nullstellensatz de Hilbert}
\label{algebra-section-nullstellensatz}

\begin{theorem}[Nullstellensatz de Hilbert]
\label{algebra-theorem-nullstellensatz}
Soit $k$ un corps.
\begin{enumerate}
\item
\label{algebra-item-finite-kappa}
Pour tout idéal maximal $\mathfrak m \subset k[x_1, \ldots, x_n]$,
l'extension de corps $\kappa(\mathfrak m)/k$ est finie.
\item
\label{algebra-item-polynomial-ring-Jacobson}
Tout idéal radical $I \subset k[x_1, \ldots, x_n]$
est l'intersection des idéaux maximaux qui le contiennent.
\end{enumerate}
Il en va de même dans toute $k$-algèbre de type fini.
\end{theorem}

\begin{proof}
Il suffit de démontrer la partie (\ref{algebra-item-finite-kappa}) du
théorème dans le cas d'une algèbre polynomiale
$k[x_1, \ldots, x_n]$, car toute $k$-algèbre de type fini
est un quotient d'une telle algèbre polynomiale.
Nous le démontrons par récurrence sur $n$. Le cas $n = 0$ est clair.
Supposons que $\mathfrak m$ soit un idéal maximal de $k[x_1, \ldots, x_n]$.
Soit $\mathfrak p \subset k[x_n]$ l'intersection
de $\mathfrak m$ avec $k[x_n]$.

\medskip\noindent
Si $\mathfrak p \not = (0)$,
alors $\mathfrak p$ est maximal et engendré par un polynôme
irréductible unitaire $P$ (grâce à l'algorithme d'Euclide
dans $k[x_n]$). Alors
$k' = k[x_n]/\mathfrak p$ est une extension finie de $k$
et est contenu dans $\kappa(\mathfrak m)$. Dans ce cas,
nous obtenons une surjection
$$
k'[x_1, \ldots, x_{n-1}]
\to
k'[x_1, \ldots, x_n] =
k' \otimes_k k[x_1, \ldots, x_n]
\longrightarrow
\kappa(\mathfrak m)
$$
et nous voyons donc que $\kappa(\mathfrak m)$ est une extension finie
de $k'$ par l'hypothèse de récurrence. Ainsi, $\kappa(\mathfrak m)$
est également finie sur $k$.

\medskip\noindent
Si $\mathfrak p = (0)$, nous considérons l'extension d'anneaux
$k[x_n] \subset k[x_1, \ldots, x_n]/\mathfrak m$.
Cette extension d'anneaux est de type fini, donc
de présentation finie d'après les
lemmes \ref{algebra-lemma-obvious-Noetherian} et
\ref{algebra-lemma-Noetherian-finite-type-is-finite-presentation}.
Ainsi, l'image de $\Spec(k[x_1, \ldots, x_n]/\mathfrak m)$
dans $\Spec(k[x_n])$ est constructible d'après le
théorème \ref{algebra-theorem-chevalley}. Puisque l'image
contient $(0)$, nous concluons qu'elle contient un ouvert
standard $D(f)$ pour un certain $f\in k[x_n]$ non nul. Comme
$D(f)$ est manifestement infini, nous obtenons une contradiction avec
l'hypothèse que $k[x_1, \ldots, x_n]/\mathfrak m$ est
un corps (et a donc un spectre constitué d'un seul point).

\medskip\noindent
Démonstration de (\ref{algebra-item-polynomial-ring-Jacobson}). Soit
$I \subset R$ un idéal radical, où $R$ est de type fini sur $k$.
Soit $f \in R$, $f \not \in I$. Nous devons trouver un idéal maximal
$\mathfrak m \subset R$ tel que $I \subset \mathfrak m$ et
$f \not \in \mathfrak m$. L'anneau $(R/I)_f$ est non nul, car
$1 = 0$ dans cet anneau signifierait que $f^n \in I$ et, puisque $I$ est
radical, cela signifierait que $f \in I$, contrairement à notre hypothèse sur $f$.
Nous pouvons donc choisir un idéal maximal $\mathfrak m'$
de $(R/I)_f$, voir le lemme \ref{algebra-lemma-Zariski-topology}.
Soit $\mathfrak m \subset R$
l'image réciproque de $\mathfrak m'$ dans $R$. Nous voyons que
$I \subset \mathfrak m$
et $f \not \in \mathfrak m$. Si nous montrons que $\mathfrak m$ est un idéal
maximal de $R$, alors nous avons terminé. Nous avons manifestement
$$
k \subset R/\mathfrak m \subset \kappa(\mathfrak m').
$$
D'après la partie (\ref{algebra-item-finite-kappa}), l'extension de corps
$\kappa(\mathfrak m')/k$ est finie. Ainsi,
$R/\mathfrak m$ est un corps d'après Corps, lemme
\ref{fields-lemma-subalgebra-algebraic-extension-field}.
Par conséquent, $\mathfrak m$ est maximal, ce qui achève la démonstration.
\end{proof}

\begin{lemma}
\label{algebra-lemma-field-finite-type-over-domain}
Soit $R$ un anneau. Soit $K$ un corps.
Si $R \subset K$ et si $K$ est de type fini sur $R$,
alors il existe un $f \in R$ tel que $R_f$ soit un corps
et que $K/R_f$ soit une extension finie de corps.
\end{lemma}

\begin{proof}
D'après le lemme \ref{algebra-lemma-characterize-image-finite-type}, il
existe un ouvert non vide $U \subset \Spec(R)$
contenu dans l'image $\{(0)\}$ de $\Spec(K) \to \Spec(R)$.
Choisissons $f \in R$, $f \not = 0$, tel que $D(f) \subset U$, c'est-à-dire
$D(f) = \{(0)\}$. Alors $R_f$ est un anneau intègre dont le spectre a exactement un
point, et $R_f$ est un corps. Ensuite, $K$ est une algèbre de type fini
sur le corps $R_f$, donc une extension finie de
$R_f$ d'après le Nullstellensatz de Hilbert (théorème \ref{algebra-theorem-nullstellensatz}).
\end{proof}





















\section{Anneaux de Jacobson}
\label{algebra-section-ring-jacobson}

\noindent
Soit $R$ un anneau. Les points fermés de $\Spec(R)$ sont les
idéaux maximaux de $R$. Souvent, les anneaux qui apparaissent naturellement en
géométrie algébrique ont beaucoup d'idéaux maximaux. C'est par exemple le cas des algèbres
de type fini sur un corps ou sur $\mathbf{Z}$. Nous montrerons que celles-ci
sont des exemples d'anneaux de Jacobson.

\begin{definition}
\label{algebra-definition-ring-jacobson}
Soit $R$ un anneau. Nous disons que $R$ est un
{\it anneau de Jacobson} si tout idéal
radical $I$ est l'intersection des
idéaux maximaux qui le contiennent.
\end{definition}

\begin{lemma}
\label{algebra-lemma-finite-type-field-Jacobson}
Toute algèbre de type fini sur un corps est de Jacobson.
\end{lemma}

\begin{proof}
Cela résulte du théorème \ref{algebra-theorem-nullstellensatz}
et de la définition \ref{algebra-definition-ring-jacobson}.
\end{proof}

\begin{lemma}
\label{algebra-lemma-jacobson-prime}
Soit $R$ un anneau. Si tout idéal premier de $R$ est
l'intersection des idéaux maximaux qui le contiennent,
alors $R$ est de Jacobson.
\end{lemma}

\begin{proof}
C'est immédiat du fait que
chaque idéal radical $I \subset R$ est
l'intersection des idéaux premiers qui le contiennent.
Voir le lemme \ref{algebra-lemma-Zariski-topology}.
\end{proof}

\begin{lemma}
\label{algebra-lemma-jacobson}
Un anneau $R$ est de Jacobson si et seulement si $\Spec(R)$
est un espace de Jacobson, voir Topologie,
définition \ref{topology-definition-space-jacobson}.
\end{lemma}

\begin{proof}
Supposons que $R$ soit de Jacobson. Soit $Z \subset \Spec(R)$
un fermé. Nous devons montrer que l'ensemble des points
fermés de $Z$ est dense dans $Z$. Soit $U \subset \Spec(R)$
un ouvert tel que $U \cap Z$ soit non vide.
Nous devons montrer que $Z \cap U$ contient un point fermé
de $\Spec(R)$. Nous pouvons
supposer que $U = D(f)$, car les ouverts standards forment une base de la
topologie de $\Spec(R)$. D'après le
lemme \ref{algebra-lemma-Zariski-topology}, nous pouvons supposer que
$Z = V(I)$, où $I$ est un idéal radical. Nous voyons également
que $f \not \in I$. Par hypothèse, il existe un
idéal maximal $\mathfrak m \subset R$ tel que
$I \subset \mathfrak m$ mais $f \not\in \mathfrak m$.
Ainsi, $\mathfrak m \in D(f) \cap V(I) = U \cap Z$, comme voulu.

\medskip\noindent
Réciproquement, supposons que $\Spec(R)$ soit un espace de Jacobson.
Soit $I \subset R$ un idéal radical. Soit
$J = \cap_{I \subset \mathfrak m} \mathfrak m$
l'intersection des idéaux maximaux qui contiennent $I$.
Manifestement, $J$ est un idéal radical, $V(J) \subset V(I)$, et
$V(J)$ est le plus petit fermé de $V(I)$ qui contient
tous les points fermés de $V(I)$. Par hypothèse, nous voyons que
$V(J) = V(I)$. Mais le lemme \ref{algebra-lemma-Zariski-topology}
montre qu'il existe une bijection entre les fermés de Zariski
et les idéaux radicaux ; ainsi, $I = J$, comme voulu.
\end{proof}

\begin{lemma}
\label{algebra-lemma-characterize-jacobson}
Soit $R$ un anneau. Si $R$ n'est pas de Jacobson, il existe
un idéal premier $\mathfrak p \subset R$ et un élément $f \in R$
tels que les propriétés suivantes soient satisfaites :
\begin{enumerate}
\item $\mathfrak p$ n'est pas un idéal maximal,
\item $f \not \in \mathfrak p$,
\item $V(\mathfrak p) \cap D(f) = \{\mathfrak p\}$, et
\item $(R/\mathfrak p)_f$ est un corps.
\end{enumerate}
En revanche, si $R$ est de Jacobson, alors, pour toute paire $(\mathfrak p, f)$
qui satisfait (1) et (2), l'ensemble $V(\mathfrak p) \cap D(f)$ est
infini.
\end{lemma}

\begin{proof}
Supposons que $R$ ne soit pas de Jacobson.
D'après le lemme \ref{algebra-lemma-jacobson}, cela signifie qu'il existe un
fermé $T \subset \Spec(R)$
dont l'ensemble $T_0 \subset T$ des points fermés n'est pas dense dans $T$.
Choisissons un $f \in R$ tel que $T_0 \subset V(f)$ mais
$T \not \subset V(f)$. Notons que $T \cap D(f)$
est homéomorphe à $\Spec((R/I)_f)$ si $T = V(I)$, voir les
lemmes \ref{algebra-lemma-spec-closed} et \ref{algebra-lemma-standard-open}.
Comme tout anneau possède un idéal maximal
(lemme \ref{algebra-lemma-Zariski-topology}), nous pouvons choisir un point fermé $t$ de
l'espace $T \cap D(f)$. Alors $t$ correspond à un idéal premier
$\mathfrak p \subset R$ qui n'est pas maximal (car $t \not \in T_0$).
Ainsi, (1) est satisfaite. Par construction, $f \not \in \mathfrak p$, d'où (2).
Comme $t$ est un point fermé de $T \cap D(f)$, nous voyons que
$V(\mathfrak p) \cap D(f) = \{\mathfrak p\}$, c'est-à-dire que (3) est satisfaite. Nous
concluons donc que $(R/\mathfrak p)_f$ est un anneau intègre dont le
spectre a un seul point, d'où (4)
(on peut par exemple combiner les lemmes \ref{algebra-lemma-characterize-local-ring} et
\ref{algebra-lemma-minimal-prime-reduced-ring}).

\medskip\noindent
Réciproquement, supposons que $R$ soit de Jacobson et que $(\mathfrak p, f)$
satisfasse (1) et (2). Si
$V(\mathfrak p) \cap D(f) =
\{\mathfrak p, \mathfrak q_1, \ldots, \mathfrak q_t\}$,
alors $\mathfrak p \not = \mathfrak q_i$
implique qu'il existe un élément $g \in R$ tel que $g \not \in \mathfrak p$
mais $g \in \mathfrak q_i$ pour tout $i$. Ainsi,
$V(\mathfrak p) \cap D(fg) = \{\mathfrak p\}$, ce qui
est impossible puisque tout sous-ensemble localement fermé de $\Spec(R)$
contient au moins un point fermé, car $\Spec(R)$ est
un espace topologique de Jacobson.
\end{proof}

\begin{lemma}
\label{algebra-lemma-pid-jacobson}
L'anneau $\mathbf{Z}$ est un anneau de Jacobson.
Plus généralement, soit $R$ un anneau tel que
\begin{enumerate}
\item $R$ soit un anneau intègre,
\item $R$ soit noethérien,
\item tout idéal premier non nul soit un idéal maximal, et
\item $R$ possède une infinité d'idéaux maximaux.
\end{enumerate}
Alors $R$ est un anneau de Jacobson.
\end{lemma}

\begin{proof}
Soit $R$ satisfaisant (1), (2), (3) et (4). L'assertion
signifie que $(0) = \bigcap_{\mathfrak m \subset R} \mathfrak m$.
Puisque $R$ possède une infinité d'idéaux maximaux, il suffit de
montrer que tout $x \in R$ non nul appartient à seulement
un nombre fini d'idéaux maximaux, autrement dit que $V(x)$ est fini.
D'après le lemme \ref{algebra-lemma-spec-closed},
nous voyons que $V(x)$ est homéomorphe à $\Spec(R/xR)$.
Par l'hypothèse (3), tout idéal premier de $R/xR$ est minimal et
correspond donc à une composante irréductible de $\Spec(R/xR)$
(lemme \ref{algebra-lemma-irreducible}).
Comme $R/xR$ est noethérien, l'espace topologique $\Spec(R/xR)$
est noethérien (lemme \ref{algebra-lemma-Noetherian-topology})
et possède un nombre fini de composantes irréductibles
(Topologie, lemme \ref{topology-lemma-Noetherian}).
Ainsi, $V(x)$ est fini, comme voulu.
\end{proof}

\begin{example}
\label{algebra-example-infinite-product-fields-jacobson}
Soit $A$ un ensemble infini.
Pour chaque $\alpha \in A$, soit $k_\alpha$ un corps.
Nous affirmons que $R = \prod_{\alpha\in A} k_\alpha$ est de Jacobson.
Remarquons d'abord que tout élément $f \in R$ est de la forme
$f = ue$, où $u \in R$ est une unité et $e\in R$ un idempotent
(laissé au lecteur).
Ainsi, $D(f) = D(e)$, et $R_f = R_e = R/(1-e)$ est un quotient de $R$.
En fait, tout anneau possédant cette propriété est de Jacobson.
En effet, soit $\mathfrak p \subset R$ un idéal premier
et $f \in R$, $f \not \in \mathfrak p$. Nous devons trouver
un idéal maximal $\mathfrak m$ de $R$ tel que
$\mathfrak p \subset \mathfrak m$ et $f \not\in \mathfrak m$.
Puisque $R_f$ est un quotient de $R$, nous voyons que tout idéal
maximal de $R_f$ correspond à un idéal maximal de $R$
ne contenant pas $f$. Le résultat s'obtient donc
en choisissant un idéal maximal de $R_f$ contenant $\mathfrak p R_f$.
\end{example}

\begin{example}
\label{algebra-example-not-jacobson}
Un anneau intègre $R$ possédant un nombre fini d'idéaux maximaux
$\mathfrak m_i$, $i = 1, \ldots, n$, n'est pas un
anneau de Jacobson, sauf s'il s'agit d'un corps.
En effet, dans ce cas, $(0)$ n'est pas l'intersection
des idéaux maximaux : $(0) \not =
\mathfrak m_1 \cap \mathfrak m_2 \cap \ldots \cap \mathfrak m_n
\supset \mathfrak m_1 \cdot \mathfrak m_2 \cdot \ldots
\cdot \mathfrak m_n \not = 0$.
En particulier, un anneau de valuation discrète, ou tout anneau local ayant
au moins deux idéaux premiers, n'est pas un anneau
de Jacobson.
\end{example}

\begin{lemma}
\label{algebra-lemma-finite-residue-extension-closed}
Soit $R \to S$ un morphisme d'anneaux.
Soit $\mathfrak m \subset R$ un idéal maximal.
Soit $\mathfrak q \subset S$ un idéal premier
au-dessus de $\mathfrak m$ tel que $\kappa(\mathfrak q)/\kappa(\mathfrak m)$
soit une extension algébrique de corps.
Alors $\mathfrak q$ est un idéal maximal de $S$.
\end{lemma}

\begin{proof}
Considérons le diagramme
$$
\xymatrix{
S \ar[r] & S/\mathfrak q \ar[r] & \kappa(\mathfrak q) \\
R \ar[r] \ar[u] & R/\mathfrak m \ar[u]
}
$$
Nous voyons que $\kappa(\mathfrak m) \subset S/\mathfrak q \subset
\kappa(\mathfrak q)$. Puisque l'extension de corps
$\kappa(\mathfrak m) \subset \kappa(\mathfrak q)$
est algébrique, tout anneau compris entre $\kappa(\mathfrak m)$
et $\kappa(\mathfrak q)$ est un corps
(Corps, lemme \ref{fields-lemma-subalgebra-algebraic-extension-field}).
Ainsi, $S/\mathfrak q$ est un corps, et est donc égal
à $\kappa(\mathfrak q)$.
\end{proof}

\begin{lemma}
\label{algebra-lemma-dimension}
Supposons que $k$ soit un corps et que $V$ soit un espace vectoriel non nul
sur $k$. Supposons que la dimension de $V$ (qui est un nombre cardinal)
soit strictement inférieure au cardinal de $k$. Alors, pour tout endomorphisme
$T : V \to V$, il existe un polynôme unitaire $P(t) \in k[t]$ tel que
$P(T)$ ne soit pas inversible.
\end{lemma}

\begin{proof}
Sinon, $V$ hérite d'une structure d'espace vectoriel sur
le corps $k(t)$. Or la dimension de $k(t)$ sur $k$ est
au moins égale au cardinal de $k$, par exemple parce que les éléments
$\frac{1}{t - \lambda}$ sont linéairement indépendants sur $k$.
\end{proof}

\noindent
Voici une autre version du Nullstellensatz de Hilbert.

\begin{theorem}
\label{algebra-theorem-uncountable-nullstellensatz}
Soit $k$ un corps. Soit $S$ une $k$-algèbre engendrée sur $k$
par les éléments $\{x_i\}_{i \in I}$. Supposons que le cardinal de $I$
soit strictement inférieur au cardinal de $k$. Alors
\begin{enumerate}
\item pour tout idéal maximal $\mathfrak m \subset S$, l'extension de corps
$\kappa(\mathfrak m)/k$
est algébrique, et
\item $S$ est un anneau de Jacobson.
\end{enumerate}
\end{theorem}

\begin{proof}
Si $I$ est fini, le résultat découle du Nullstellensatz de Hilbert,
théorème \ref{algebra-theorem-nullstellensatz}. Dans la suite de la démonstration, nous supposons
$I$ infini. Il suffit de démontrer le résultat lorsque
$\mathfrak m \subset k[\{x_i\}_{i \in I}]$ est maximal dans l'anneau
de polynômes en les variables $x_i$, puisque $S$ en est un quotient.
Puisque $I$ est infini, l'ensemble
des monômes $x_{i_1}^{e_1} \ldots x_{i_r}^{e_r}$, $i_1, \ldots, i_r \in I$
et $e_1, \ldots, e_r \geq 0$, est de cardinal inférieur ou égal au
cardinal de $I$. En effet, le cardinal de $I \times \ldots \times I$
est le cardinal de $I$, et le cardinal de
$\bigcup_{n \geq 0} I^n$ est lui aussi le même.
(Si $I$ est fini, cela n'est plus vrai et,
dans ce cas, cette démonstration ne fonctionne que si $k$ est non dénombrable.)

\medskip\noindent
Pour obtenir une contradiction, choisissons $T \in \kappa(\mathfrak m)$ transcendant
sur $k$. Notons que l'application $k$-linéaire
$T : \kappa(\mathfrak m) \to \kappa(\mathfrak m)$
donnée par la multiplication par $T$ possède la propriété que $P(T)$ est inversible
pour tout polynôme unitaire $P(t) \in k[t]$.
De plus, $\kappa(\mathfrak m)$ est de dimension au plus égale au cardinal de $I$
sur $k$, puisqu'il s'agit d'un quotient de l'espace vectoriel
$k[\{x_i\}_{i \in I}]$ sur $k$ (dont la dimension vaut $\# I$, comme vu ci-dessus).
C'est impossible d'après le lemme \ref{algebra-lemma-dimension}.

\medskip\noindent
Pour montrer que $S$ est de Jacobson, raisonnons comme suit. Sinon,
il existe un idéal premier $\mathfrak q \subset S$ et un élément $f \in S$,
$f \not \in \mathfrak q$, tels que $\mathfrak q$ ne soit pas maximal
et que $(S/\mathfrak q)_f$ soit un corps, voir le
lemme \ref{algebra-lemma-characterize-jacobson}.
Mais remarquons que $(S/\mathfrak q)_f$ est engendré par au plus $\# I + 1$ éléments.
Ainsi, l'extension de corps $(S/\mathfrak q)_f/k$ est algébrique
(d'après la première partie de la démonstration).
Cela implique que $\kappa(\mathfrak q)$ est une extension algébrique de $k$ ;
donc $\mathfrak q$ est maximal d'après le
lemme \ref{algebra-lemma-finite-residue-extension-closed}. Cette contradiction
achève la démonstration.
\end{proof}

\begin{lemma}
\label{algebra-lemma-base-change-Jacobson}
Soit $k$ un corps. Soit $S$ une $k$-algèbre.
Pour toute extension de corps $K/k$ dont le cardinal est strictement supérieur
au cardinal de $S$, nous avons :
\begin{enumerate}
\item pour tout idéal maximal $\mathfrak m$ de $S_K$, le corps
$\kappa(\mathfrak m)$ est algébrique sur $K$, et
\item $S_K$ est un anneau de Jacobson.
\end{enumerate}
\end{lemma}

\begin{proof}
Considérons $k \subset K$ tel que le cardinal de $K$ soit supérieur
au cardinal de $S$. Puisque les éléments de $S$ engendrent
la $K$-algèbre $S_K$, nous voyons que le
théorème \ref{algebra-theorem-uncountable-nullstellensatz}
s'applique.
\end{proof}

\begin{example}
\label{algebra-example-countable-trick-does-not-work}
L'astuce employée dans la démonstration du
théorème \ref{algebra-theorem-uncountable-nullstellensatz}
ne fonctionne réellement pas si $k$ est un corps dénombrable et $I$ est lui aussi dénombrable.
Soit $k$ un corps dénombrable. Soit $x$ une indéterminée,
et soit $k(x)$ le corps des fractions rationnelles en $x$.
Considérons l'algèbre polynomiale $R = k[x, \{x_f\}_{f \in k[x]-\{0\}}]$.
Soit $I = (\{fx_f - 1\}_{f\in k[x] - \{0\}})$. Remarquons que
$I$ est un idéal propre de $R$.
Choisissons un idéal maximal $I \subset \mathfrak m$.
Alors $k \subset R/\mathfrak m$ est isomorphe à
$k(x)$, et n'est pas algébrique sur $k$.
\end{example}

\begin{lemma}
\label{algebra-lemma-Jacobson-invert-element}
Soit $R$ un anneau de Jacobson. Soit $f \in R$. L'anneau $R_f$ est de Jacobson, et
les idéaux maximaux de $R_f$ correspondent aux idéaux maximaux de $R$ qui ne contiennent pas $f$.
\end{lemma}

\begin{proof}
D'après Topologie, lemme \ref{topology-lemma-jacobson-inherited},
nous voyons que $D(f) = \Spec(R_f)$ est de Jacobson et
que les points fermés de $D(f)$
correspondent aux points fermés de $\Spec(R)$
qui appartiennent à $D(f)$. Ainsi, $R_f$ est de Jacobson d'après le
lemme \ref{algebra-lemma-jacobson}.
\end{proof}

\begin{example}
\label{algebra-example-localize-not-preserve-closed-points}
Voici un exemple simple qui montre que le lemme \ref{algebra-lemma-Jacobson-invert-element}
est faux si $R$ n'est pas de Jacobson.
Considérons l'anneau $R = \mathbf{Z}_{(2)}$, c'est-à-dire la localisation
de $\mathbf{Z}$ en l'idéal premier $(2)$. La localisation de $R$ par
l'élément $2$ est isomorphe à $\mathbf{Q}$ ; sous forme d'une formule :
$R_2 \cong \mathbf{Q}$. Manifestement, le morphisme $R \to R_2$ envoie le
point fermé de $\Spec(\mathbf{Q})$ sur le point générique
de $\Spec(R)$.
\end{example}

\begin{example}
\label{algebra-example-infinite-localize-not-preserve-closed-points}
Voici un exemple simple qui montre que le
lemme \ref{algebra-lemma-Jacobson-invert-element}
est faux si $R$ est de Jacobson mais que nous localisons par une infinité
d'éléments.
En effet, soit $R = \mathbf{Z}$ et considérons la localisation
$(R \setminus \{0\})^{-1}R \cong \mathbf{Q}$
de $R$ par l'ensemble de tous les éléments non nuls.
Manifestement, le morphisme $\mathbf{Z} \to \mathbf{Q}$ envoie le
point fermé de $\Spec(\mathbf{Q})$ sur le point générique
de $\Spec(\mathbf{Z})$.
\end{example}

\begin{lemma}
\label{algebra-lemma-Jacobson-mod-ideal}
Soit $R$ un anneau de Jacobson. Soit $I \subset R$ un idéal.
L'anneau $R/I$ est de Jacobson, et les idéaux maximaux
de $R/I$ correspondent aux idéaux maximaux de $R$ qui contiennent $I$.
\end{lemma}

\begin{proof}
La démonstration est la même que celle du
lemme \ref{algebra-lemma-Jacobson-invert-element}.
\end{proof}

\begin{lemma}
\label{algebra-lemma-silly-jacobson}
Soit $R$ un anneau de Jacobson. Soit $K$ un corps. Supposons que $R \subset K$ et
que $K$ soit de type fini sur $R$. Alors $R$ est un corps et $K/R$
est une extension finie de corps.
\end{lemma}

\begin{proof}
Remarquons d'abord que $R$ est un anneau intègre.
D'après le lemme \ref{algebra-lemma-field-finite-type-over-domain},
nous voyons que $R_f$ est un corps et que $K/R_f$ est une extension finie de corps
pour un certain $f \in R$ non nul. Ainsi, $(0)$ est un idéal maximal de $R_f$
et, d'après le
lemme \ref{algebra-lemma-Jacobson-invert-element},
nous concluons que $(0)$ est un idéal maximal de $R$.
\end{proof}

\begin{proposition}
\label{algebra-proposition-Jacobson-permanence}
Soit $R$ un anneau de Jacobson. Soit $R \to S$ un
morphisme d'anneaux de type fini. Alors :
\begin{enumerate}
\item L'anneau $S$ est de Jacobson.
\item Le morphisme $\Spec(S) \to \Spec(R)$ envoie
les points fermés sur des points fermés.
\item Pour $\mathfrak m' \subset S$ maximal au-dessus de $\mathfrak m \subset R$,
l'extension de corps $\kappa(\mathfrak m')/\kappa(\mathfrak m)$
est finie.
\end{enumerate}
\end{proposition}

\begin{proof}
Soit $\mathfrak m' \subset S$ un idéal maximal, et
$R \cap \mathfrak m' = \mathfrak m$.
Alors $R/\mathfrak m \to S/\mathfrak m'$ satisfait
les conditions du lemme \ref{algebra-lemma-silly-jacobson}
d'après le lemme \ref{algebra-lemma-Jacobson-mod-ideal}.
Ainsi, $R/\mathfrak m$ est un corps,
$\mathfrak m$ est un idéal maximal et l'extension induite
des corps résiduels est finie. Cela démontre (2) et (3).  

\medskip\noindent
Si $S$ n'est pas de Jacobson, alors, d'après le lemme \ref{algebra-lemma-characterize-jacobson}, il
existe un idéal premier non maximal $\mathfrak q$ de $S$ et un
$g \in S$, $g \not\in \mathfrak q$, tels que $(S/\mathfrak q)_g$ soit un corps.
Pour obtenir une contradiction, nous montrons que $\mathfrak q$ est un idéal maximal.
Soit $\mathfrak p = \mathfrak q \cap R$. Alors
$R/\mathfrak p \to (S/\mathfrak q)_g$ satisfait les conditions du
lemme \ref{algebra-lemma-silly-jacobson} d'après le
lemme \ref{algebra-lemma-Jacobson-mod-ideal}.
Ainsi, $R/\mathfrak p$ est un corps et l'extension de corps
$\kappa(\mathfrak p) \to (S/\mathfrak q)_g = \kappa(\mathfrak q)$ est
finie, donc algébrique. Alors $\mathfrak q$ est un idéal maximal de $S$ d'après le
lemme \ref{algebra-lemma-finite-residue-extension-closed}. Contradiction.
\end{proof}

\begin{lemma}
\label{algebra-lemma-corollary-jacobson}
Toute algèbre de type fini sur $\mathbf{Z}$ est de Jacobson.
\end{lemma}

\begin{proof}
Combiner le lemme \ref{algebra-lemma-pid-jacobson} et la
proposition \ref{algebra-proposition-Jacobson-permanence}.
\end{proof}

\begin{lemma}
\label{algebra-lemma-image-finite-type-map-Jacobson-rings}
Soit $R \to S$ un morphisme de type fini entre anneaux de Jacobson.
Posons $X = \Spec(R)$ et $Y = \Spec(S)$.
Notons $f : Y \to X$ le morphisme induit
sur les spectres. Soit $E \subset Y = \Spec(S)$ un ensemble
constructible. Notons par un indice ${}_0$ l'ensemble
des points fermés d'un espace topologique.
\begin{enumerate}
\item Nous avons $f(E)_0 = f(E_0) = X_0 \cap f(E)$.
\item Un point $\xi \in X$ appartient à $f(E)$ si et seulement si
$\overline{\{\xi\}} \cap f(E_0)$ est dense dans $\overline{\{\xi\}}$.
\end{enumerate}
\end{lemma}

\begin{proof}
Nous avons un diagramme commutatif d'applications continues
$$
\xymatrix{
E \ar[r] \ar[d] & Y \ar[d] \\
f(E) \ar[r] & X
}
$$
Supposons que $x \in f(E)$ soit fermé dans $f(E)$. Alors $f^{-1}(\{x\})\cap E$
est non vide et fermé dans $E$. En appliquant le
lemme \ref{topology-lemma-jacobson-inherited} de Topologie
aux deux inclusions
$$
f^{-1}(\{x\}) \cap E \subset E \subset Y
$$
nous trouvons qu'il existe un point $y \in f^{-1}(\{x\}) \cap E$ qui est
fermé dans $Y$. Autrement dit, il existe $y \in Y_0$ et $y \in E_0$
qui s'envoie sur $x$. Ainsi, $x \in f(E_0)$.
Cela démontre que $f(E)_0 \subset f(E_0)$.
La proposition \ref{algebra-proposition-Jacobson-permanence} implique que
$f(E_0) \subset X_0 \cap f(E)$. L'inclusion
$X_0 \cap f(E) \subset f(E)_0$ est triviale. Cela démontre la
première assertion.

\medskip\noindent
Supposons que $\xi \in f(E)$. D'après le
lemme \ref{algebra-lemma-characterize-image-finite-type},
l'ensemble $f(E) \cap \overline{\{\xi\}}$ contient un ouvert
dense de $\overline{\{\xi\}}$. Puisque $X$ est de Jacobson,
nous concluons que $f(E) \cap \overline{\{\xi\}}$ contient un
ensemble dense de points fermés, voir Topologie,
lemme \ref{topology-lemma-jacobson-inherited}.
Nous concluons par la partie (1) du lemme.

\medskip\noindent
Réciproquement, supposons que $\overline{\{\xi\}} \cap f(E_0)$
soit dense dans $\overline{\{\xi\}}$. D'après le
lemme \ref{algebra-lemma-constructible-is-image},
il existe un morphisme d'anneaux $S \to S'$ de présentation finie
tel que $E$ soit l'image de $Y' := \Spec(S') \to Y$.
Alors $E_0$ est l'image de $Y'_0$ d'après la première partie du
lemme appliquée au morphisme d'anneaux $S \to S'$. Nous pouvons donc supposer que
$E = Y$ en remplaçant $S$ par $S'$. Supposons que $\xi$ corresponde
à $\mathfrak p \subset R$. Considérons le diagramme
$$
\xymatrix{
S \ar[r] & S/\mathfrak p S \\
R \ar[r] \ar[u] & R/\mathfrak p \ar[u]
}
$$
Ce diagramme et la densité de $f(Y_0) \cap V(\mathfrak p)$
dans $V(\mathfrak p)$
montrent que le morphisme $R/\mathfrak p \to S/\mathfrak p S$
satisfait la condition (2) du
lemme \ref{algebra-lemma-domain-image-dense-set-points-generic-point}.
Nous en concluons
qu'il existe un idéal premier $\overline{\mathfrak q} \subset S/\mathfrak pS$
qui s'envoie sur $(0)$. Autrement dit, l'image réciproque $\mathfrak q$
de $\overline{\mathfrak q}$ dans $S$ s'envoie sur $\mathfrak p$, comme voulu.
\end{proof}

\noindent
La conclusion du lemme ci-dessus est que nous pouvons déterminer
l'image de $f$ à partir de l'ensemble des points fermés de l'image.
C'est un peu plus commode lorsque le morphisme est de présentation finie,
car nous savons alors que l'image d'une partie constructible est constructible.
Avant de l'énoncer, introduisons quelques notations.
Notons $\text{Constr}(X)$ l'ensemble des parties constructibles.
Soit $R \to S$ un morphisme d'anneaux.
Posons $X = \Spec(R)$ et $Y = \Spec(S)$.
Notons $f : Y \to X$ le morphisme induit sur les spectres.
Notons par un indice ${}_0$ l'ensemble
des points fermés d'un espace topologique.


\begin{lemma}
\label{algebra-lemma-conclude-jacobson-Noetherian}
Avec les notations ci-dessus. Supposons que $R$ soit un anneau de Jacobson noethérien.
Supposons en outre que $R \to S$ soit de type fini.
Il existe un diagramme commutatif
$$
\xymatrix{
\text{Constr}(Y) \ar[r]^{E \mapsto E_0} \ar[d]^{E \mapsto f(E)} &
\text{Constr}(Y_0) \ar[d]^{E \mapsto f(E)} \\
\text{Constr}(X) \ar[r]^{E \mapsto E_0} &
\text{Constr}(X_0)
}
$$
où les flèches horizontales sont les bijections du
lemme \ref{topology-lemma-jacobson-equivalent-constructible} de Topologie.
\end{lemma}

\begin{proof}
Puisque $R \to S$ est de type fini, il est de présentation finie,
voir le lemme \ref{algebra-lemma-Noetherian-finite-type-is-finite-presentation}.
Ainsi, l'image d'une partie constructible dans $X$ est constructible
dans $Y$ d'après le théorème de Chevalley
(théorème \ref{algebra-theorem-chevalley}).
Combiné au
lemme \ref{algebra-lemma-image-finite-type-map-Jacobson-rings},
cela démontre le lemme.
\end{proof}

\noindent
Pour illustrer l'emploi des anneaux de Jacobson, donnons les deux exemples suivants.

\begin{example}
\label{algebra-example-product-matrices-zero}
Soit $k$ un corps. L'espace $\Spec(k[x, y]/(xy))$
possède deux composantes irréductibles : l'axe des $x$ et
l'axe des $y$. Comme généralisation, soit
$$
R = k[x_{11}, x_{12}, x_{21}, x_{22}, y_{11}, y_{12}, y_{21}, y_{22}]/
\mathfrak a,
$$
où $\mathfrak a$ est l'idéal de
$k[x_{11}, x_{12}, x_{21}, x_{22}, y_{11}, y_{12}, y_{21}, y_{22}]$
engendré par les coefficients du produit des matrices $2 \times 2$
$$
\left(
\begin{matrix}
x_{11} & x_{12}\\
x_{21} & x_{22}
\end{matrix}
\right)
 \left(
\begin{matrix}
y_{11} & y_{12}\\
y_{21} & y_{22}
\end{matrix}
\right).
$$
Dans cet exemple, nous allons décrire $\Spec(R)$.

\medskip\noindent
Pour démontrer l'assertion concernant $\Spec(k[x, y]/(xy))$, raisonnons comme suit.
Si $\mathfrak p \subset k[x, y]$ est un idéal premier contenant $xy$, alors
soit $x$, soit $y$ appartient à $\mathfrak p$. Les idéaux premiers minimaux
ayant cette propriété sont donc simplement $(x)$ et $(y)$. Lorsque $k$ est
algébriquement clos, le $\text{max-Spec}$ de ces composantes
se visualise alors comme les ensembles de points des axes des $y$ et des $x$.

\medskip\noindent
Pour la généralisation, remarquons que nous pouvons identifier les points
fermés du spectre de
$k[x_{11}, x_{12}, x_{21}, x_{22}, y_{11}, y_{12}, y_{21}, y_{22}])$
à l'espace des matrices
$$
\left\{ (X, Y) \in \text{Mat}(2, k)\times \text{Mat}(2, k) \mid
X = \left(
\begin{matrix}
x_{11} & x_{12}\\
x_{21} & x_{22}
\end{matrix}
\right),
Y= \left(
\begin{matrix}
y_{11} & y_{12}\\
y_{21} & y_{22}
\end{matrix}
\right)
\right\}
$$
du moins si $k$ est algébriquement clos.
Définissons maintenant une action du groupe
$\text{GL}(2, k)\times \text{GL}(2, k)\times \text{GL}(2, k)$
sur l'espace des matrices $\{(X, Y)\}$ par
$$
(g_1, g_2, g_3) \times (X, Y) \mapsto ((g_1Xg_2^{-1}, g_2Yg_3^{-1})).
$$
Observons également que l'ensemble algébrique
$$
\text{GL}(2, k)\times \text{GL}(2, k)\times \text{GL}(2, k) \subset
\text{Mat}(2, k)\times \text{Mat}(2, k) \times \text{Mat}(2, k)
$$
est irréductible, puisqu'il est le spectre maximal de l'anneau intègre
$$
k[x_{11}, x_{12}, \ldots, z_{21}, z_{22}, (x_{11}x_{22}-x_{12}x_{21})^{-1}
, (y_{11}y_{22}-y_{12}y_{21})^{-1}, (z_{11}z_{22}-z_{12}z_{21})^{-1}].
$$
Puisque l'image d'un ensemble algébrique irréductible est encore
irréductible, il suffit de classifier les orbites de l'ensemble
$\{(X, Y)\in \text{Mat}(2, k)\times \text{Mat}(2, k)|XY = 0\}$ et de prendre leurs
adhérences. L'algèbre linéaire usuelle nous ramène aux
trois cas suivants :
\begin{enumerate}
\item $\exists (g_1, g_2)$ tel que $g_1Xg_2^{-1} = I_{2\times 2}$.
Alors $Y$ vaut nécessairement $0$ ; cet ensemble algébrique est
invariant sous l'action du groupe. Il s'ensuit que cette orbite est
contenue dans l'ensemble algébrique irréductible défini par l'idéal premier
$(y_{11}, y_{12}, y_{21}, y_{22})$. En prenant l'adhérence, nous voyons
que $(y_{11}, y_{12}, y_{21}, y_{22})$ est effectivement une composante.
\item $\exists (g_1, g_2)$ tel que
$$
g_1Xg_2^{-1} = \left(
\begin{matrix}
1 & 0 \\
0 & 0
\end{matrix}
\right).
$$
Ce cas se produit si et seulement si $X$ est une matrice de rang 1,
et, de plus, $Y$ est annulée par une telle matrice $X$ si et seulement si
$$
x_{11}y_{11}+x_{12}y_{21} = 0; \quad x_{11}y_{12}+x_{12}y_{22} = 0;
$$
$$
x_{21}y_{11}+x_{22}y_{21} = 0; \quad x_{21}y_{12}+x_{22}y_{22} = 0.
$$
Fixons une matrice $X$ de rang 1 ; les matrices $Y$ non nulles qui satisfont les
équations ci-dessus forment un ensemble algébrique irréductible pour la
raison suivante ($Y = 0$ relève du cas précédent) :
$0 = g_1Xg_2^{-1}g_2Y$ implique que
$$
g_2Y = \left(
\begin{matrix}
0 & 0 \\
y_{21}' & y_{22}'
\end{matrix}
\right).
$$
À l'aide d'une autre action de $\text{GL}(2, k)$ à droite par $g_3$,
$g_2Y$ peut être amenée à la forme
$$
g_2Yg_3^{-1} = \left(
\begin{matrix}
0 & 0 \\
0 & 1
\end{matrix}
\right),
$$
et ces matrices $Y$ forment donc un ensemble algébrique irréductible
isomorphe à l'image de $\text{GL}(2, k)$ sous cette action. Enfin,
remarquons que la condition « de rang 1 » sur les matrices $X$ définit un ouvert dense
de l'ensemble algébrique irréductible
$\det X = x_{11}x_{22} - x_{12}x_{21} = 0$.
Il s'ensuit maintenant que les cinq équations définissent une composante irréductible
$(x_{11}y_{11}+x_{12}y_{21}, x_{11}y_{12}+x_{12}y_{22}, x_{21}y_{11}
+x_{22}y_{21}, x_{21}y_{12}+x_{22}y_{22}, x_{11}x_{22}-x_{12}x_{21})$
dans l'ouvert de l'espace des paires de matrices non nulles.
On peut montrer que les deux équations
$\det X = 0$, $\det Y = 0$ découpent $\Spec(R)$ suivant une composante irréductible
dont le lieu ci-dessus est un ouvert dense.
\item $\exists (g_1, g_2)$ tel que $g_1Xg_2^{-1} = 0$, ou,
de manière équivalente, $X = 0$. Alors $Y$ peut être arbitraire, et cette composante
est donc définie par $(x_{11}, x_{12}, x_{21}, x_{22})$.
\end{enumerate}
\end{example}


\begin{example}
\label{algebra-example-idempotent-matrices}
Pour un autre exemple, considérons
$R = k[\{t_{ij}\}_{i, j = 1}^{n}]/\mathfrak a$, où $\mathfrak a$ est
l'idéal engendré par les coefficients de la matrice produit $T^2-T$,
$T = (t_{ij})$. L'algèbre linéaire nous apprend que, sous
l'action de $GL(n, k)$ définie par $g, T \mapsto gTg^{-1}$, $T$ est
classifiée par son rang, et chaque $T$ est conjuguée à une matrice
$\text{diag}(1, \ldots, 1, 0, \ldots, 0)$, qui possède $r$ coefficients égaux à 1 et $n-r$ coefficients égaux à 0.
Ainsi, chaque orbite d'une telle matrice $\text{diag}(1, \ldots, 1, 0, \ldots, 0)$ sous
l'action du groupe forme une composante irréductible, et toute matrice
idempotente appartient à l'une de ces orbites. Montrons maintenant que deux
orbites distinctes sont nécessairement disjointes. Pour cela, il
suffit de construire des fonctions polynomiales qui prennent des
valeurs différentes sur des orbites différentes. En caractéristique 0, une telle
fonction peut être
$f(t_{ij}) = trace(T) = \sum_{i = 1}^nt_{ii}$. En caractéristique
positive, les choses sont un peu plus délicates, car nous
pouvons avoir $trace(T) = 0$ même si $T \neq 0$. Par exemple, $char = 3$
$$
trace\left(
\begin{matrix}
1 & & \\
& 1 & \\
& & 1
\end{matrix}
\right) = 3 = 0
$$
Quoi qu'il en soit, ces composantes peuvent être séparées à l'aide d'autres fonctions.
Par exemple, en caractéristique 3, $tr(\wedge^3T)$ prend
la valeur 1 sur les composantes correspondant à $diag(1, 1, 1)$ et 0 sur
les autres composantes.
\end{example}





































\section{Extensions d'anneaux finies et entières}
\label{algebra-section-finite-ring-extensions}

\noindent
Lemmes élémentaires concernant les morphismes d'anneaux finis et entiers.
Rappelons la définition.

\begin{definition}
\label{algebra-definition-integral-ring-map}
Soit $\varphi : R \to S$ un morphisme d'anneaux.
\begin{enumerate}
\item Un élément $s \in S$
est {\it entier sur $R$} s'il existe un polynôme
unitaire $P(x) \in R[x]$ tel que
$P^\varphi(s) = 0$, où $P^\varphi(x) \in S[x]$
est l'image de $P$ par $\varphi : R[x] \to S[x]$.
\item Le morphisme d'anneaux $\varphi$ est {\it entier}
si tout $s \in S$ est entier sur $R$.
\end{enumerate}
\end{definition}

\begin{lemma}
\label{algebra-lemma-characterize-integral-element}
Soit $\varphi : R \to S$ un morphisme d'anneaux. Soit $y \in S$. S'il existe un
sous-$R$-module de type fini $M$ de $S$ tel que $1 \in M$ et $yM \subset M$,
alors $y$ est entier sur $R$.
\end{lemma}

\begin{proof}
Considérons l'application $\varphi : M \to M$, $x \mapsto y \cdot x$.
D'après le lemme \ref{algebra-lemma-charpoly-module}, il existe un polynôme unitaire
$P \in R[T]$ tel que $P(\varphi) = 0$. Dans l'anneau $S$, nous obtenons
$P(y) = P(y) \cdot 1 = P(\varphi)(1) = 0$.
\end{proof}

\begin{lemma}
\label{algebra-lemma-finite-is-integral}
Un morphisme d'anneaux fini est entier.
\end{lemma}

\begin{proof}
Soit $R \to S$ fini. Soit $y \in S$. Appliquons le
lemme \ref{algebra-lemma-characterize-integral-element}
à $M = S$ pour voir que $y$ est entier sur $R$.
\end{proof}

\begin{lemma}
\label{algebra-lemma-characterize-integral}
Soit $\varphi : R \to S$ un morphisme d'anneaux. Soient $s_1, \ldots, s_n$
un ensemble fini d'éléments de $S$.
Alors les $s_i$ sont entiers sur $R$ pour tout $i = 1, \ldots, n$
si et seulement s'il
existe une $R$-sous-algèbre $S' \subset S$ finie sur $R$
qui contient tous les $s_i$.
\end{lemma}

\begin{proof}
Si chaque $s_i$ est entier, alors la sous-algèbre
engendrée par $\varphi(R)$ et les $s_i$ est finie
sur $R$. En effet, si $s_i$ satisfait une équation unitaire
de degré $d_i$ sur $R$, alors cette sous-algèbre est engendrée comme
$R$-module par les éléments $s_1^{e_1} \ldots s_n^{e_n}$
avec $0 \leq e_i \leq d_i - 1$.
Réciproquement, supposons donnée une $R$-sous-algèbre finie
$S'$ qui contient tous les $s_i$. Alors tous les
$s_i$ sont entiers d'après le lemme \ref{algebra-lemma-finite-is-integral}.
\end{proof}

\begin{lemma}
\label{algebra-lemma-characterize-finite-in-terms-of-integral}
Soit $R \to S$ un morphisme d'anneaux. Les propriétés suivantes sont équivalentes :
\begin{enumerate}
\item $R \to S$ est fini,
\item $R \to S$ est entier et de type fini, et
\item il existe $x_1, \ldots, x_n \in S$ qui engendrent $S$ comme
algèbre sur $R$ et tels que chaque $x_i$ soit entier sur $R$.
\end{enumerate}
\end{lemma}

\begin{proof}
C'est clair d'après le lemme \ref{algebra-lemma-characterize-integral}.
\end{proof}

\begin{lemma}
\label{algebra-lemma-integral-transitive}
\begin{slogan}
Une composée de morphismes d'anneaux entiers est entière
\end{slogan}
Supposons que $R \to S$ et $S \to T$ soient des
morphismes d'anneaux entiers. Alors $R \to T$ est entier.
\end{lemma}

\begin{proof}
Soit $t \in T$. Soit $P(x) \in S[x]$ un
polynôme unitaire tel que $P(t) = 0$.
Appliquons le lemme \ref{algebra-lemma-characterize-integral}
à l'ensemble fini des coefficients de $P$.
Ainsi, $t$ est entier sur une certaine sous-algèbre
$S' \subset S$ finie sur $R$. Appliquons à nouveau le lemme
\ref{algebra-lemma-characterize-integral} pour trouver
une sous-algèbre $T' \subset T$ finie sur $S'$ et
contenant $t$. Le lemme \ref{algebra-lemma-finite-transitive}
appliqué à $R \to S' \to T'$ montre que $T'$ est finie
sur $R$. L'intégralité de $t$ sur $R$
résulte alors du lemme \ref{algebra-lemma-finite-is-integral}.
\end{proof}

\begin{lemma}
\label{algebra-lemma-integral-closure-is-ring}
Soit $R \to S$ un morphisme d'anneaux.
L'ensemble
$$
S' = \{s \in S \mid s\text{ est entier sur }R\}
$$
est une $R$-sous-algèbre de $S$.
\end{lemma}

\begin{proof}
C'est clair d'après les lemmes \ref{algebra-lemma-characterize-integral}
et \ref{algebra-lemma-finite-is-integral}.
\end{proof}

\begin{lemma}
\label{algebra-lemma-finite-product-integral}
Soient $R_i\to S_i$ des morphismes d'anneaux, $i = 1, \ldots, n$.
Notons respectivement $R$ et $S$ les produits des $R_i$ et des $S_i$.
Alors un élément $s = (s_1, \ldots, s_n) \in S$ est entier sur $R$
si et seulement si chaque $s_i$ est entier sur $R_i$.
\end{lemma}

\begin{proof}
Omis.
\end{proof}

\begin{definition}
\label{algebra-definition-integral-closure}
Soit $R \to S$ un morphisme d'anneaux.
L'anneau $S' \subset S$ des éléments entiers sur
$R$, voir le lemme \ref{algebra-lemma-integral-closure-is-ring},
est appelé la {\it clôture intégrale} de $R$
dans $S$. Si $R \subset S$, nous disons que $R$ est
{\it intégralement clos} dans $S$ si $R = S'$.
\end{definition}

\noindent
En particulier, nous voyons que $R \to S$ est entier si et seulement
si la clôture intégrale de $R$ dans $S$ est $S$ tout entier.

\begin{lemma}
\label{algebra-lemma-finite-product-integral-closure}
Soient $R_i\to S_i$ des morphismes d'anneaux, $i = 1, \ldots, n$.
Notons la clôture intégrale de $R_i$ dans $S_i$ par $S'_i$.
Notons en outre respectivement $R$ et $S$ les produits des $R_i$ et des $S_i$.
Alors la clôture intégrale de $R$ dans $S$
est le produit des $S'_i$. En particulier, $R \to S$ est
intégralement clos si et seulement si chaque $R_i \to S_i$ est intégralement clos.
\end{lemma}

\begin{proof}
Cela résulte immédiatement du lemme \ref{algebra-lemma-finite-product-integral}.
\end{proof}

\begin{lemma}
\label{algebra-lemma-integral-closure-localize}
La clôture intégrale commute à la localisation : si $A \to B$ est un
morphisme d'anneaux et $S \subset A$ une partie multiplicative, alors la clôture
intégrale de $S^{-1}A$ dans $S^{-1}B$ est $S^{-1}B'$, où $B' \subset B$
est la clôture intégrale de $A$ dans $B$.
\end{lemma}

\begin{proof}
Puisque la localisation est exacte, nous voyons que $S^{-1}B' \subset S^{-1}B$.
Supposons que $x \in B'$ et $f \in S$. Alors
$x^d + \sum_{i = 1, \ldots, d} a_i x^{d - i} = 0$
dans $B$ pour certains $a_i \in A$. Ainsi,
$$
(x/f)^d + \sum\nolimits_{i = 1, \ldots, d} a_i/f^i (x/f)^{d - i} = 0
$$
dans $S^{-1}B$. Nous voyons ainsi que $S^{-1}B'$ est contenu dans
la clôture intégrale de $S^{-1}A$ dans $S^{-1}B$. Réciproquement, supposons
que $x/f \in S^{-1}B$ soit entier sur $S^{-1}A$. Alors nous avons
$$
(x/f)^d + \sum\nolimits_{i = 1, \ldots, d} (a_i/f_i) (x/f)^{d - i} = 0
$$
dans $S^{-1}B$ pour certains $a_i \in A$ et $f_i \in S$. Cela signifie que
$$
(f'f_1 \ldots f_d x)^d +
\sum\nolimits_{i = 1, \ldots, d}
f^i(f')^if_1^i \ldots f_i^{i - 1} \ldots f_d^i a_i
(f'f_1 \ldots f_dx)^{d - i} = 0
$$
pour un $f' \in S$ convenable. Ainsi, $f'f_1\ldots f_dx \in B'$ et donc
$x/f \in S^{-1}B'$, comme voulu.
\end{proof}

\begin{lemma}
\label{algebra-lemma-integral-closure-stalks}
\begin{slogan}
Un élément d'une algèbre sur un anneau est entier sur cet anneau
si et seulement s'il est localement entier en tout idéal premier de l'anneau.
\end{slogan}
Soit $\varphi : R \to S$ un morphisme d'anneaux.
Soit $x \in S$. Les propriétés suivantes sont équivalentes :
\begin{enumerate}
\item $x$ est entier sur $R$, et
\item pour tout idéal premier $\mathfrak p \subset R$, l'élément
$x \in S_{\mathfrak p}$ est entier sur $R_{\mathfrak p}$.
\end{enumerate}
\end{lemma}

\begin{proof}
Il est clair que (1) implique (2). Supposons (2). Considérons la $R$-algèbre
$S' \subset S$ engendrée par $\varphi(R)$ et $x$. Soit $\mathfrak p$ un
idéal premier de $R$. Nous savons alors que
$x^d + \sum_{i = 1, \ldots, d} \varphi(a_i) x^{d - i} = 0$
dans $S_{\mathfrak p}$ pour certains $a_i \in R_{\mathfrak p}$. En examinant
les dénominateurs qui interviennent, nous voyons donc que,
pour un certain $f \in R$, $f \not \in \mathfrak p$, nous avons
$a_i \in R_f$ et
$x^d + \sum_{i = 1, \ldots, d} \varphi(a_i) x^{d - i} = 0$
dans $S_f$. Cela implique que $S'_f$ est fini sur $R_f$.
Puisque $\mathfrak p$ était arbitraire et que $\Spec(R)$ est quasi-compact
(lemme \ref{algebra-lemma-quasi-compact}), nous pouvons trouver un nombre fini d'éléments
$f_1, \ldots, f_n \in R$
qui engendrent l'idéal unité de $R$ tels que $S'_{f_i}$ soit fini
sur $R_{f_i}$. Nous concluons donc du lemme \ref{algebra-lemma-cover} que
$S'$ est fini sur $R$. Ainsi, $x$ est entier sur $R$ d'après le
lemme \ref{algebra-lemma-characterize-integral}.
\end{proof}

\begin{lemma}
\label{algebra-lemma-base-change-integral}
\begin{slogan}
L'intégralité et la finitude sont préservées par changement de base.
\end{slogan}
Soient $R \to S$ et $R \to R'$ des morphismes d'anneaux.
Posons $S' = R' \otimes_R S$.
\begin{enumerate}
\item Si $R \to S$ est entier, alors $R' \to S'$ l'est aussi.
\item Si $R \to S$ est fini, alors $R' \to S'$ l'est aussi.
\end{enumerate}
\end{lemma}

\begin{proof}
Nous démontrons (1).
Soient $s_i \in S$ des générateurs de $S$ sur $R$.
Chacun d'eux satisfait une équation polynomiale unitaire $P_i$
sur $R$. Ainsi, les éléments $1 \otimes s_i \in S'$ engendrent
$S'$ sur $R'$ et satisfont le polynôme correspondant
$P_i'$ sur $R'$. Puisque ces éléments engendrent $S'$ sur $R'$,
nous voyons que $S'$ est entier sur $R'$.
Démonstration de (2) omise.
\end{proof}

\begin{lemma}
\label{algebra-lemma-integral-local}
Soit $R \to S$ un morphisme d'anneaux.
Soient $f_1, \ldots, f_n \in R$ engendrant l'idéal unité.
\begin{enumerate}
\item Si chaque $R_{f_i} \to S_{f_i}$ est entier, alors $R \to S$ l'est aussi.
\item Si chaque $R_{f_i} \to S_{f_i}$ est fini, alors $R \to S$ l'est aussi.
\end{enumerate}
\end{lemma}

\begin{proof}
Démonstration de (1).
Soit $s \in S$. Considérons l'idéal $I \subset R[x]$ des
polynômes $P$ tels que $P(s) = 0$. Notons $J \subset R$
l'idéal (!) des coefficients dominants des éléments de $I$.
Par hypothèse et en chassant les dénominateurs,
nous voyons que $f_i^{n_i} \in J$ pour tout $i$
et certains $n_i \geq 0$. Ainsi, $J$ contient $1$, et nous voyons que
$s$ est entier sur $R$. Démonstration de (2) omise.
\end{proof}

\begin{lemma}
\label{algebra-lemma-integral-permanence}
Soient $A \to B \to C$ des morphismes d'anneaux.
\begin{enumerate}
\item Si $A \to C$ est entier, alors $B \to C$ l'est aussi.
\item Si $A \to C$ est fini, alors $B \to C$ l'est aussi.
\end{enumerate}
\end{lemma}

\begin{proof}
Omis.
\end{proof}

\begin{lemma}
\label{algebra-lemma-integral-closure-transitive}
Soient $A \to B \to C$ des morphismes d'anneaux.
Soit $B'$ la clôture intégrale de $A$ dans $B$,
et soit $C'$ la clôture intégrale de $B'$ dans $C$. Alors
$C'$ est la clôture intégrale de $A$ dans $C$.
\end{lemma}

\begin{proof}
Omis.
\end{proof}

\begin{lemma}
\label{algebra-lemma-integral-overring-surjective}
Supposons que $R \to S$ soit une extension
d'anneaux entière avec $R \subset S$.
Alors $\varphi : \Spec(S) \to \Spec(R)$
est surjective.
\end{lemma}

\begin{proof}
Soit $\mathfrak p \subset R$ un idéal premier.
Nous devons montrer que $\mathfrak pS_{\mathfrak p} \not = S_{\mathfrak p}$, voir le
lemme \ref{algebra-lemma-in-image}.
La localisation $R_{\mathfrak p} \to S_{\mathfrak p}$ est injective
(car la localisation est exacte) et entière d'après le
lemme \ref{algebra-lemma-integral-closure-localize} ou
\ref{algebra-lemma-base-change-integral}.
Nous pouvons donc remplacer $R$, $S$ par $R_{\mathfrak p}$, $S_{\mathfrak p}$ et
supposer que $R$ est local d'idéal maximal $\mathfrak m$ ; il
suffit alors de montrer que $\mathfrak mS \not = S$.
Supposons que $1 = \sum f_i s_i$, avec $f_i \in \mathfrak m$
et $s_i \in S$, afin d'obtenir une contradiction.
Soit $R \subset S' \subset S$
tel que $R \to S'$ soit fini et que $s_i \in S'$, voir le
lemme \ref{algebra-lemma-characterize-integral}.
L'équation $1 = \sum f_i s_i$ implique que
le $R$-module de type fini $S'$ satisfait $S' = \mathfrak m S'$. Le
lemme de Nakayama \ref{algebra-lemma-NAK} donne donc
$S' = 0$. Contradiction.
\end{proof}

\begin{lemma}
\label{algebra-lemma-integral-under-field}
Soit $R$ un anneau. Soit $K$ un corps.
Si $R \subset K$ et si $K$ est entier sur $R$,
alors $R$ est un corps et $K$ est une extension algébrique.
Si $R \subset K$ et si $K$ est fini sur $R$,
alors $R$ est un corps et $K$ est une extension algébrique finie.
\end{lemma}

\begin{proof}
Supposons que $R \subset K$ soit entier.
D'après le lemme \ref{algebra-lemma-integral-overring-surjective}, nous voyons que
$\Spec(R)$ possède $1$ point. Comme $R$ est manifestement un anneau intègre, nous voyons
que $R = R_{(0)}$ est un corps (lemme \ref{algebra-lemma-minimal-prime-reduced-ring}).
Les autres assertions en résultent immédiatement.
\end{proof}

\begin{lemma}
\label{algebra-lemma-integral-over-field}
Soit $k$ un corps. Soit $S$ une $k$-algèbre sur $k$.
\begin{enumerate}
\item Si $S$ est un anneau intègre de dimension finie sur $k$,
alors $S$ est un corps.
\item Si $S$ est entier sur $k$ et est un anneau intègre,
alors $S$ est un corps.
\item Si $S$ est entier sur $k$, alors tout idéal premier de
$S$ est un idéal maximal (voir le
lemme \ref{algebra-lemma-ring-with-only-minimal-primes}
pour d'autres conséquences).
\end{enumerate}
\end{lemma}

\begin{proof}
L'assertion sur les idéaux premiers résulte de l'assertion
« entier $+$ intègre $\Rightarrow$ corps ».
Soit $S$ entier sur $k$ et supposons que $S$ soit un anneau intègre.
Prenons $s \in S$. D'après le lemme
\ref{algebra-lemma-characterize-integral}, nous pouvons trouver une
$k$-sous-algèbre de dimension finie $k \subset S' \subset S$
qui contient $s$. Ainsi, $S$ est un corps si nous pouvons démontrer la
première assertion. Supposons $S$ de dimension finie
sur $k$ et intègre. Choisissons $s\in S$ non nul.
Puisque $S$ est un anneau intègre, l'application de multiplication
$s : S \to S$ est surjective pour des raisons de
dimension. Il existe donc un élément $s_1 \in S$
tel que $ss_1 = 1$. Ainsi, $S$ est un corps.
\end{proof}

\begin{lemma}
\label{algebra-lemma-integral-no-inclusion}
Supposons que $R \to S$ soit entier.
Soient $\mathfrak q, \mathfrak q' \in \Spec(S)$
deux idéaux premiers distincts
ayant la même image dans $\Spec(R)$.
Alors ni $\mathfrak q \subset \mathfrak q'$
ni $\mathfrak q' \subset \mathfrak q$.
\end{lemma}

\begin{proof}
Soit $\mathfrak p \subset R$ l'image.
D'après la remarque \ref{algebra-remark-fundamental-diagram},
les idéaux premiers $\mathfrak q, \mathfrak q'$
correspondent à des idéaux de
$S \otimes_R \kappa(\mathfrak p)$.
Le lemme résulte donc du lemme \ref{algebra-lemma-integral-over-field}.
\end{proof}

\begin{lemma}
\label{algebra-lemma-finite-finite-fibres}
Supposons que $R \to S$ soit fini.
Alors les fibres de $\Spec(S) \to \Spec(R)$ sont finies.
\end{lemma}

\begin{proof}
D'après la discussion de la
remarque \ref{algebra-remark-fundamental-diagram},
les fibres sont les spectres des anneaux $S \otimes_R \kappa(\mathfrak p)$.
Comme $R \to S$ est fini, ces anneaux fibres sont finis sur
$\kappa(\mathfrak p)$, donc noethériens d'après le
lemme \ref{algebra-lemma-Noetherian-permanence}.
D'après le
lemme \ref{algebra-lemma-integral-no-inclusion},
tout idéal premier de $S \otimes_R \kappa(\mathfrak p)$ est
minimal. Ainsi, d'après le
lemme \ref{algebra-lemma-Noetherian-irreducible-components},
il n'y en a qu'un nombre fini.
\end{proof}

\begin{lemma}
\label{algebra-lemma-integral-going-up}
Soit $R \to S$ un morphisme d'anneaux tel que
$S$ soit entier sur $R$.
Soient $\mathfrak p \subset \mathfrak p' \subset R$
des idéaux premiers. Soit $\mathfrak q$ un idéal premier de $S$ qui s'envoie
sur $\mathfrak p$. Alors il existe un idéal premier $\mathfrak q'$
tel que $\mathfrak q \subset \mathfrak q'$
et qui s'envoie sur $\mathfrak p'$.
\end{lemma}

\begin{proof}
Nous pouvons remplacer $R$ par $R/\mathfrak p$ et $S$ par $S/\mathfrak q$.
Cela nous ramène à une extension entière d'anneaux
intègres $R \subset S$ et à un idéal premier $\mathfrak p' \subset R$.
Le lemme \ref{algebra-lemma-integral-overring-surjective} permet de conclure.
\end{proof}

\noindent
La propriété exprimée dans le lemme ci-dessus est appelée
la « propriété de montée » pour le morphisme d'anneaux $R \to S$,
voir la définition \ref{algebra-definition-going-up-down}.

\begin{lemma}
\label{algebra-lemma-finite-finitely-presented-extension}
Soit $R \to S$ un morphisme d'anneaux fini et de présentation finie.
Soit $M$ un $S$-module.
Alors $M$ est de présentation finie comme $R$-module si et seulement si
$M$ est de présentation finie comme $S$-module.
\end{lemma}

\begin{proof}
L'une des implications résulte du
lemme \ref{algebra-lemma-finitely-presented-over-subring}.
Pour voir l'autre, supposons que $M$ soit de présentation finie comme $S$-module.
Choisissons une présentation
$$
S^{\oplus m} \longrightarrow
S^{\oplus n} \longrightarrow
M \longrightarrow 0
$$
Comme $S$ est fini comme $R$-module, le noyau de
$S^{\oplus n} \to M$ est un $R$-module de type fini. Ainsi, le
lemme \ref{algebra-lemma-extension}
montre qu'il suffit de démontrer que $S$ est de présentation finie comme
$R$-module.

\medskip\noindent
Choisissons $y_1, \ldots, y_n \in S$ tels que $y_1, \ldots, y_n$ engendrent $S$
comme $R$-module. D'après le lemme \ref{algebra-lemma-characterize-integral-element},
chaque $y_i$ est entier sur $R$. Choisissons des polynômes unitaires
$P_i(x) \in R[x]$ tels que $P_i(y_i) = 0$. Considérons l'anneau
$$
S' = R[x_1, \ldots, x_n]/(P_1(x_1), \ldots, P_n(x_n))
$$
Nous voyons alors que $S$ est de présentation finie comme $S'$-algèbre
d'après le lemme \ref{algebra-lemma-compose-finite-type}. Puisque $S' \to S$ est surjectif,
le noyau $J = \Ker(S' \to S)$ est engendré par un nombre fini d'éléments comme idéal d'après le
lemme \ref{algebra-lemma-finite-presentation-independent}. Ainsi, $J$ est un
$S'$-module de type fini (cela découle immédiatement des définitions).
Par conséquent, $S = \Coker(J \to S')$ est de présentation finie comme $S'$-module
d'après le lemme \ref{algebra-lemma-extension}.
Ainsi, en raisonnant comme dans le premier paragraphe, il suffit de montrer que $S'$ est
de présentation finie comme $R$-module. En fait, $S'$ est libre comme
$R$-module, de base les monômes $x_1^{e_1} \ldots x_n^{e_n}$
pour $0 \leq e_i < \deg(P_i)$. En effet, écrivons $R \to S'$ comme la composée
$$
R \to R[x_1]/(P_1(x_1)) \to R[x_1, x_2]/(P_1(x_1), P_2(x_2)) \to
\ldots \to S'
$$
Cela montre que le $i$-ème anneau de cette suite est libre comme module sur le
$(i - 1)$-ième, de base $1, x_i, \ldots, x_i^{\deg(P_i) - 1}$. Le résultat
s'en déduit aisément par récurrence. Certains détails sont omis.
\end{proof}

\begin{lemma}
\label{algebra-lemma-silly-normal}
Soit $R$ un anneau. Soient $x, y \in R$ des non-diviseurs de zéro.
Soit $R[x/y] \subset R_{xy}$ la $R$-sous-algèbre engendrée
par $x/y$, et de même pour les sous-algèbres $R[y/x]$ et $R[x/y, y/x]$.
Si $R$ est intégralement clos dans $R_x$ ou $R_y$, alors la suite
$$
0 \to R \xrightarrow{(-1, 1)} R[x/y] \oplus R[y/x] \xrightarrow{(1, 1)}
R[x/y, y/x] \to 0
$$
est une suite exacte courte de $R$-modules.
\end{lemma}

\begin{proof}
Puisque $x/y \cdot y/x = 1$, il est clair que l'application
$R[x/y] \oplus R[y/x] \to R[x/y, y/x]$ est surjective.
Soit $\alpha \in R[x/y] \cap R[y/x]$. Pour montrer l'exactitude au milieu,
nous devons prouver que $\alpha \in R$. Par hypothèse, nous pouvons écrire
$$
\alpha = a_0 + a_1 x/y + \ldots + a_n (x/y)^n =
b_0 + b_1 y/x + \ldots + b_m(y/x)^m
$$
pour certains $n, m \geq 0$ et $a_i, b_j \in R$.
Choisissons un $N > \max(n, m)$.
Considérons le $R$-sous-module de type fini $M$ de $R_{xy}$ engendré par les éléments
$$
(x/y)^N, (x/y)^{N - 1}, \ldots, x/y, 1, y/x, \ldots, (y/x)^{N - 1}, (y/x)^N
$$
Nous affirmons que $\alpha M \subset M$. En effet, il est clair que
$(x/y)^i (b_0 + b_1 y/x + \ldots + b_m(y/x)^m) \in M$ pour
$0 \leq i \leq N$ et que
$(y/x)^i (a_0 + a_1 x/y + \ldots + a_n(x/y)^n) \in M$ pour
$0 \leq i \leq N$. Ainsi, $\alpha$ est entier sur $R$ d'après le
lemme \ref{algebra-lemma-characterize-integral-element}. Remarquons que
$\alpha \in R_x$ ; si $R$ est intégralement clos dans $R_x$,
alors $\alpha \in R$, comme voulu.
\end{proof}







\section{Anneaux normaux}
\label{algebra-section-normal-rings}

\noindent
Introduisons d'abord la notion d'anneau intègre normal, puis
la notion (très générale) d'anneau normal.

\begin{definition}
\label{algebra-definition-domain-normal}
Un anneau intègre $R$ est dit {\it normal} s'il est intégralement
clos dans son corps des fractions.
\end{definition}

\begin{lemma}
\label{algebra-lemma-integral-closure-in-normal}
Soit $R \to S$ un morphisme d'anneaux.
Si $S$ est un anneau intègre normal, alors la clôture intégrale de $R$
dans $S$ est un anneau intègre normal.
\end{lemma}

\begin{proof}
Omis.
\end{proof}

\noindent
La notion suivante est parfois utile dans l'étude
de la normalité.

\begin{definition}
\label{algebra-definition-almost-integral}
Soit $R$ un anneau intègre.
\begin{enumerate}
\item Un élément $g$ du corps des
fractions de $R$ est dit {\it presque entier sur $R$}
s'il existe un élément $r \in R$, $r\not = 0$,
tel que $rg^n \in R$ pour tout $n \geq 0$.
\item L'anneau intègre $R$ est dit {\it complètement normal} si tout
élément presque entier du corps des fractions de $R$
appartient à $R$.
\end{enumerate}
\end{definition}

\noindent
Le lemme suivant montre qu'un anneau intègre noethérien est normal
si et seulement s'il est complètement normal.

\begin{lemma}
\label{algebra-lemma-almost-integral}
Soit $R$ un anneau intègre de corps des fractions $K$.
Si $u, v \in K$ sont presque entiers sur $R$, alors
$u + v$ et $uv$ le sont aussi. Tout élément $g \in K$ qui est entier sur $R$
est presque entier sur $R$. Si $R$ est noethérien,
alors la réciproque est également vraie.
\end{lemma}

\begin{proof}
Si $ru^n \in R$ pour tout $n \geq 0$ et
$v^nr' \in R$ pour tout $n \geq 0$, alors
$(uv)^nrr'$ et $(u + v)^nrr'$ appartiennent à $R$ pour
tout $n \geq 0$. D'où la première assertion.
Supposons que $g \in K$ soit entier sur $R$.
Dans ce cas, il existe un $d > 0$ tel que
l'anneau $R[g]$ soit engendré par $1, g, \ldots, g^d$ comme $R$-module.
Soit $r \in R$ un dénominateur commun des éléments
$1, g, \ldots, g^d \in K$. Il s'ensuit que $rR[g] \subset R$,
et donc que $g$ est presque entier sur $R$.

\medskip\noindent
Supposons que $R$ soit noethérien et que $g \in K$ soit presque entier sur $R$.
Soit $r \in R$, $r\not = 0$, comme dans la définition.
Alors $R[g] \subset \frac{1}{r}R$ comme $R$-module.
Puisque $R$ est noethérien, cela implique que $R[g]$ est
fini sur $R$. Ainsi, $g$ est entier sur $R$, voir le
lemme \ref{algebra-lemma-finite-is-integral}.
\end{proof}

\begin{lemma}
\label{algebra-lemma-localize-normal-domain}
Toute localisation d'un anneau intègre normal est normale.
\end{lemma}

\begin{proof}
Soit $R$ un anneau intègre normal, et soit $S \subset R$ une
partie multiplicative. Supposons que $g$ soit un élément
du corps des fractions de $R$ qui soit entier sur $S^{-1}R$.
Soit $P = x^d + \sum_{j < d} a_j x^j$ un polynôme
avec $a_i \in S^{-1}R$ tel que $P(g) = 0$.
Choisissons $s \in S$ tel que $sa_i \in R$ pour tout $i$.
Alors $sg$ satisfait le polynôme unitaire
$x^d + \sum_{j < d} s^{d-j}a_j x^j$, dont les coefficients
$s^{d-j}a_j$ appartiennent à $R$. Ainsi, $sg \in R$ puisque $R$ est normal.
Donc $g \in S^{-1}R$.
\end{proof}

\begin{lemma}
\label{algebra-lemma-PID-normal}
Un anneau principal est normal.
\end{lemma}

\begin{proof}
Soit $R$ un anneau principal.
Soit $g = a/b$ un élément du corps des fractions
de $R$ qui soit entier sur $R$. Puisque $R$ est un anneau principal,
nous pouvons diviser $a$ et $b$ par un facteur commun
et supposer que $(a, b) = R$. Dans ce cas, toute équation
$(a/b)^n + r_{n-1} (a/b)^{n-1} + \ldots + r_0 = 0$
avec $r_i \in R$ impliquerait que $a^n \in (b)$. Cela
contredit $(a, b) = R$, sauf si $b$ est une unité de $R$.
\end{proof}

\begin{lemma}
\label{algebra-lemma-prepare-polynomial-ring-normal}
Soit $R$ un anneau intègre de corps des fractions $K$.
Supposons que $f = \sum \alpha_i x^i$ soit un
élément de $K[x]$.
\begin{enumerate}
\item Si $f$ est entier sur $R[x]$,
alors tous les $\alpha_i$ sont entiers sur $R$, et
\item Si $f$ est presque entier sur $R[x]$,
alors tous les $\alpha_i$ sont presque entiers sur $R$.
\end{enumerate}
\end{lemma}

\begin{proof}
Démontrons d'abord la seconde assertion.
Écrivons $f = \alpha_0 + \alpha_1 x + \ldots + \alpha_r x^r$
avec $\alpha_r \not = 0$.
Par hypothèse, il existe $h = b_0 + b_1 x + \ldots + b_s x^s \in R[x]$,
$b_s \not = 0$, tel que $f^n h \in R[x]$ pour tout
$n \geq 0$. Cela implique que $b_s \alpha_r^n \in R$
pour tout $n \geq 0$. Ainsi, $\alpha_r$ est presque
entier sur $R$. Puisque l'ensemble des éléments presque
entiers forme un sous-anneau (lemme \ref{algebra-lemma-almost-integral}), nous en déduisons que
$f - \alpha_r x^r = \alpha_0 + \alpha_1 x + \ldots + \alpha_{r - 1} x^{r - 1}$
est presque entier sur $R[x]$. Une récurrence sur $r$ permet de conclure.

\medskip\noindent
Pour démontrer la première assertion, nous utiliserons une réduction
noethérienne absolue. Écrivons $\alpha_i = a_i / b_i$ et
soit $P(t) = t^d + \sum_{j < d} f_j t^j$ un polynôme
à coefficients $f_j \in R[x]$ tel que $P(f) = 0$.
Soit $f_j = \sum f_{ji}x^i$. Considérons le sous-anneau
$R_0 \subset R$ engendré par la liste finie des éléments
$a_i, b_i, f_{ji}$ de $R$. C'est un anneau intègre ; soit
$K_0$ son corps des fractions. Puisque $R_0$ est une
$\mathbf{Z}$-algèbre de type fini, il est noethérien, voir le
lemme \ref{algebra-lemma-obvious-Noetherian}. Il est toujours
vrai que $f \in K_0[x]$ est entier sur $R_0[x]$,
car toutes les identités dans $R$
entre les éléments $a_i, b_i, f_{ji}$ sont également valables dans $R_0$.
D'après le lemme \ref{algebra-lemma-almost-integral}, l'élément
$f$ est presque entier sur $R_0[x]$. D'après la seconde assertion du
lemme, les éléments $\alpha_i$ sont presque entiers
sur $R_0$. Et puisque $R_0$ est noethérien, ils sont
entiers sur $R_0$, voir le lemme \ref{algebra-lemma-almost-integral}.
Ils sont alors bien sûr entiers sur $R$.
\end{proof}

\begin{lemma}
\label{algebra-lemma-polynomial-domain-normal}
Soit $R$ un anneau intègre normal.
Alors $R[x]$ est un anneau intègre normal.
\end{lemma}

\begin{proof}
Le résultat est vrai si $R$ est un corps $K$, car
$K[x]$ est un anneau euclidien, donc un anneau principal,
et donc normal d'après le lemme \ref{algebra-lemma-PID-normal}.
Soit $g$ un élément du corps des fractions de
$R[x]$ qui soit entier sur $R[x]$. Puisque $g$
est entier sur $K[x]$, où $K$ est le corps des
fractions de $R$, nous pouvons écrire $g = \alpha_d x^d + \alpha_{d-1}x^{d-1} +
\ldots + \alpha_0$ avec $\alpha_i \in K$.
D'après le lemme \ref{algebra-lemma-prepare-polynomial-ring-normal},
les éléments $\alpha_i$ sont entiers sur $R$ et
appartiennent donc à $R$.
\end{proof}

\begin{lemma}
\label{algebra-lemma-power-series-over-Noetherian-normal-domain}
Soit $R$ un anneau intègre noethérien normal. Alors $R[[x]]$ est
un anneau intègre noethérien normal.
\end{lemma}

\begin{proof}
L'anneau de séries formelles est noethérien d'après le
lemme \ref{algebra-lemma-Noetherian-power-series}.
Soient $f, g \in R[[x]]$ des éléments non nuls tels que
$w = f/g$ soit entier sur $R[[x]]$.
Soit $K$ le corps des fractions de $R$. Puisque l'anneau des séries de Laurent
$K((x)) = K[[x]][1/x]$ est un corps, nous pouvons écrire
$w = a_n x^n + a_{n + 1} x^{n + 1} + \ldots$
pour certains $n \in \mathbf{Z}$, $a_i \in K$, et $a_n \not = 0$.
D'après le lemme \ref{algebra-lemma-almost-integral}, nous voyons qu'il existe un
élément non nul $h = b_m x^m + b_{m + 1} x^{m + 1} + \ldots$
de $R[[x]]$, avec $b_m \not = 0$, tel que
$w^e h \in R[[x]]$ pour tout $e \geq 1$. Nous concluons que $n \geq 0$ et que
$b_m a_n^e \in R$ pour tout $e \geq 1$.
Puisque $R$ est noethérien, le même lemme implique que $a_n \in R$.
Maintenant, si $a_n, a_{n + 1}, \ldots, a_{N - 1} \in R$,
nous pouvons appliquer le même argument à
$w - a_n x^n - \ldots - a_{N - 1} x^{N - 1} = a_N x^N + \ldots$.
Nous voyons ainsi que tous les $a_i \in R$, ce qui démontre le lemme.
\end{proof}

\begin{lemma}
\label{algebra-lemma-normality-is-local}
Soit $R$ un anneau intègre. Les propriétés suivantes sont équivalentes :
\begin{enumerate}
\item L'anneau intègre $R$ est normal,
\item pour tout idéal premier $\mathfrak p \subset R$, l'anneau local
$R_{\mathfrak p}$ est un anneau intègre normal, et
\item pour tout idéal maximal $\mathfrak m$, l'anneau $R_{\mathfrak m}$
est un anneau intègre normal.
\end{enumerate}
\end{lemma}

\begin{proof}
Nous déduisons (1) $\Rightarrow$ (2) du lemme \ref{algebra-lemma-localize-normal-domain}.
L'implication (2) $\Rightarrow$ (3) est immédiate. L'implication
(3) $\Rightarrow$ (1) résulte du fait que, pour tout anneau intègre $R$, nous avons
$$
R = \bigcap\nolimits_{\mathfrak m} R_{\mathfrak m}
$$
dans le corps des fractions de $R$. En effet, si $g$ est un élément du
membre de droite, alors l'idéal $I = \{x \in R \mid xg \in R\}$
n'est contenu dans aucun idéal maximal $\mathfrak m$, d'où $I = R$.
\end{proof}

\noindent
Le lemme \ref{algebra-lemma-normality-is-local} montre que la définition suivante
est compatible avec la définition \ref{algebra-definition-domain-normal}. (Il s'agit de la
définition d'EGA -- voir \cite[IV, 5.13.5 et 0, 4.1.4]{EGA}.)

\begin{definition}
\label{algebra-definition-ring-normal}
Un anneau $R$ est dit {\it normal} si, pour tout idéal premier
$\mathfrak p \subset R$, la localisation $R_{\mathfrak p}$ est
un anneau intègre normal (voir la définition \ref{algebra-definition-domain-normal}).
\end{definition}

\noindent
Remarquons qu'un anneau normal est réduit, puisque $R$ est un sous-anneau du produit
de ses localisations en tous les idéaux premiers (voir par exemple le
lemme \ref{algebra-lemma-characterize-zero-local}).

\begin{lemma}
\label{algebra-lemma-normal-ring-integrally-closed}
Un anneau normal est intégralement clos dans son anneau total des fractions.
\end{lemma}

\begin{proof}
Soit $R$ un anneau normal. Soit $x \in Q(R)$ un élément de l'anneau total
des fractions de $R$ entier sur $R$. Posons $I = \{f \in R, fx \in R\}$. Soit
$\mathfrak p \subset R$ un idéal premier. Comme $R \to R_{\mathfrak p}$ est
plat, nous voyons que $R_{\mathfrak p} \subset Q(R) \otimes_R R_{\mathfrak p}$. Comme
$R_{\mathfrak p}$ est un anneau intègre normal, nous voyons que $x \otimes 1$ est un élément de
$R_{\mathfrak p}$. Nous pouvons donc trouver $a, f \in R$, $f \not \in \mathfrak p$,
tels que $x \otimes 1 = a \otimes 1/f$. Cela signifie que $fx - a$ s'envoie sur
zéro dans $Q(R) \otimes_R R_{\mathfrak p} = Q(R)_{\mathfrak p}$, ce qui
signifie à son tour qu'il existe un $f' \in R$, $f' \not \in \mathfrak p$,
tel que $f'fx = f'a$ dans $R$. Autrement dit, $ff' \in I$. Ainsi, $I$
est un idéal qui n'est contenu dans aucun idéal premier de $R$, c'est-à-dire
$I = R$ et $x \in R$.
\end{proof}

\begin{lemma}
\label{algebra-lemma-localization-normal-ring}
Une localisation d'un anneau normal est un anneau normal.
\end{lemma}

\begin{proof}
Omis.
\end{proof}

\begin{lemma}
\label{algebra-lemma-polynomial-ring-normal}
Soit $R$ un anneau normal. Alors $R[x]$ est un anneau normal.
\end{lemma}

\begin{proof}
Soit $\mathfrak q$ un idéal premier de $R[x]$. Posons $\mathfrak p = R \cap \mathfrak q$.
Nous voyons alors que $R_{\mathfrak p}[x]$ est un anneau intègre normal d'après le
lemme \ref{algebra-lemma-polynomial-domain-normal}.
Ainsi, $(R[x])_{\mathfrak q}$ est un anneau intègre normal d'après le
lemme \ref{algebra-lemma-localize-normal-domain}.
\end{proof}

\begin{lemma}
\label{algebra-lemma-finite-product-normal}
Un produit fini d'anneaux normaux est normal.
\end{lemma}

\begin{proof}
Il suffit de montrer que le produit de deux anneaux normaux, disons $R$ et $S$, est
normal. D'après le lemme \ref{algebra-lemma-disjoint-decomposition}, les idéaux premiers de
$R\times S$ sont de la forme $\mathfrak{p}\times S$ et $R\times
\mathfrak{q}$, où $\mathfrak{p}$ et $\mathfrak{q}$ sont des idéaux premiers respectivement de $R$
et de $S$. La localisation donne
$(R\times S)_{\mathfrak{p}\times S}=R_{\mathfrak{p}}$, qui est un anneau intègre normal
par hypothèse. De même pour $S$.
\end{proof}

\begin{lemma}
\label{algebra-lemma-characterize-reduced-ring-normal}
Soit $R$ un anneau. Supposons que $R$ soit réduit et possède un nombre fini
d'idéaux premiers minimaux. Alors les propriétés suivantes sont équivalentes :
\begin{enumerate}
\item $R$ est un anneau normal,
\item $R$ est intégralement clos dans son anneau total des fractions, et
\item $R$ est un produit fini d'anneaux intègres normaux.
\end{enumerate}
\end{lemma}

\begin{proof}
Les implications (1) $\Rightarrow$ (2) et
(3) $\Rightarrow$ (1) sont vraies en général,
voir les lemmes \ref{algebra-lemma-normal-ring-integrally-closed} et
\ref{algebra-lemma-finite-product-normal}.

\medskip\noindent
Soient $\mathfrak p_1, \ldots, \mathfrak p_n$ les idéaux premiers minimaux de $R$.
D'après les lemmes \ref{algebra-lemma-reduced-ring-sub-product-fields} et
\ref{algebra-lemma-total-ring-fractions-no-embedded-points}, nous avons
$Q(R) = R_{\mathfrak p_1} \times \ldots \times R_{\mathfrak p_n}$, et,
d'après le lemme \ref{algebra-lemma-minimal-prime-reduced-ring}, chaque facteur est un corps.
Notons $e_i = (0, \ldots, 0, 1, 0, \ldots, 0)$ le $i$-ème idempotent
de $Q(R)$.

\medskip\noindent
Si $R$ est intégralement clos dans $Q(R)$, alors il contient en particulier
les idempotents $e_i$, et nous voyons que $R$ est un produit de $n$
anneaux intègres (voir les sections \ref{algebra-section-connected-components} et
\ref{algebra-section-tilde-module-sheaf}). Chaque facteur est de la forme
$R/\mathfrak p_i$, de corps des fractions $R_{\mathfrak p_i}$.
D'après le lemme \ref{algebra-lemma-finite-product-integral-closure}, chaque morphisme
$R/\mathfrak p_i \to R_{\mathfrak p_i}$ est intégralement clos.
Ainsi, $R$ est un produit fini d'anneaux intègres normaux.
\end{proof}

\begin{lemma}
\label{algebra-lemma-colimit-normal-ring}
Soit $(R_i, \varphi_{ii'})$ un système inductif
(Catégories, définition \ref{algebra-definition-directed-system})
d'anneaux. Si chaque $R_i$ est un anneau normal, alors
$R = \colim_i R_i$ l'est aussi.
\end{lemma}

\begin{proof}
Soit $\mathfrak p \subset R$ un idéal premier.
Posons $\mathfrak p_i = R_i \cap \mathfrak p$ (abus de notation usuel).
Nous voyons alors que
$R_{\mathfrak p} = \colim_i (R_i)_{\mathfrak p_i}$.
Puisque chaque $(R_i)_{\mathfrak p_i}$ est un anneau intègre normal, nous
nous ramenons à démontrer l'énoncé du lemme pour des anneaux intègres
normaux. Si $a, b \in R$ et si $a/b$ satisfait un polynôme unitaire
$P(T) \in R[T]$, alors nous pouvons trouver un $i \in I$ (suffisamment grand)
tel que $a, b$ proviennent d'éléments $a_i, b_i$ sur $R_i$, que $P$ provienne d'un
polynôme unitaire $P_i\in R_i[T]$ et que $P_i(a_i/b_i)=0$. Puisque $R_i$ est normal, nous
voyons que $a_i/b_i \in R_i$, et donc aussi que $a/b \in R$.
\end{proof}







\section{Descente pour une extension entière au-dessus d'un anneau normal}
\label{algebra-section-going-down-integral-over-normal}

\noindent
Commençons par manipuler un peu la notion d'éléments
entiers sur un idéal, puis démontrons le théorème mentionné
dans le titre de la section.

\begin{definition}
\label{algebra-definition-integral-over-ideal}
Soit $\varphi : R \to S$ un morphisme d'anneaux.
Soit $I \subset R$ un idéal.
Nous disons qu'un élément $g \in S$ est
{\it entier sur $I$} s'il
existe un polynôme unitaire
$P = x^d + \sum_{j < d} a_j x^j$
à coefficients $a_j \in I^{d-j}$ tel
que $P^\varphi(g) = 0$ dans $S$.
\end{definition}

\noindent
Ceci est surtout utilisé lorsque $\varphi = \text{id}_R : R \to R$.
Dans ce cas, l'ensemble $I'$ des éléments entiers sur $I$ est appelé
la {\it clôture intégrale de $I$}. Nous verrons que $I'$ est
un idéal de $R$ (et, bien sûr, $I \subset I'$).

\begin{lemma}
\label{algebra-lemma-characterize-integral-ideal}
Soit $\varphi : R \to S$ un morphisme d'anneaux.
Soit $I \subset R$ un idéal.
Soit $A = \sum I^nt^n \subset R[t]$ le
sous-anneau de l'anneau de polynômes
engendré par $R \oplus It \subset R[t]$.
Un élément $s \in S$ est entier sur $I$ si
et seulement si l'élément $st \in S[t]$
est entier sur $A$.
\end{lemma}

\begin{proof}
Supposons que $st$ soit entier sur $A$.
Soit $P = x^d + \sum_{j < d} a_j x^j$
un polynôme unitaire à coefficients dans $A$
tel que $P^\varphi(st) = 0$. Soit $a_j' \in A$
la partie de degré $d-j$ de $a_j$, autrement
dit $a_j' = a_j'' t^{d-j}$ avec $a_j'' \in I^{d-j}$.
Pour des raisons de degré, nous avons encore
$(st)^d + \sum_{j < d} \varphi(a_j'') t^{d-j} (st)^j = 0$.
Ainsi, $s^d + \sum_{j < d} \varphi(a_j'') s^j = 0$,
et nous voyons que $s$ est entier sur $I$.

\medskip\noindent
Supposons que $s$ soit entier sur $I$.
Écrivons $P = x^d + \sum_{j < d} a_j x^j$
avec $a_j \in I^{d-j}$. Nous trouvons alors immédiatement un
polynôme $Q = x^d + \sum_{j < d} (a_j t^{d-j}) x^j$
à coefficients dans $A$ qui démontre que
$st$ est entier sur $A$.
\end{proof}

\begin{lemma}
\label{algebra-lemma-integral-over-ideal-is-submodule}
Soit $\varphi : R \to S$ un morphisme d'anneaux.
Soit $I \subset R$ un idéal.
L'ensemble des éléments de $S$ qui sont entiers
sur $I$ forme un $R$-sous-module de $S$.
En outre, si $s \in S$ est entier sur
$R$ et si $s'$ est entier sur $I$, alors
$ss'$ est entier sur $I$.
\end{lemma}

\begin{proof}
Nous utiliserons le lemme \ref{algebra-lemma-integral-closure-is-ring} sans
autre mention. La stabilité par addition résulte clairement de la
caractérisation du lemme \ref{algebra-lemma-characterize-integral-ideal},
dont nous adoptons les notations.
Tout élément $s \in S$ qui est entier sur
$R$ correspond à l'élément de degré $0$, $s$, de $S[t]$,
qui est entier sur $A$ (car $R \subset A$).
Nous voyons donc que la multiplication par $s$ sur $S[t]$
préserve la propriété d'être entier sur $A$.
\end{proof}

\begin{lemma}
\label{algebra-lemma-integral-integral-over-ideal}
Supposons que $\varphi : R \to S$ soit entier.
Supposons que $I \subset R$ soit un idéal.
Alors tout élément de $IS$ est entier sur $I$.
\end{lemma}

\begin{proof}
C'est immédiat d'après le lemme \ref{algebra-lemma-integral-over-ideal-is-submodule}.
\end{proof}

\begin{lemma}
\label{algebra-lemma-polynomials-divide}
Soit $K$ un corps. Soient $n, m \in \mathbf{N}$ et
$a_0, \ldots, a_{n - 1}, b_0, \ldots, b_{m - 1} \in K$.
Si le polynôme $x^n + a_{n - 1}x^{n - 1} + \ldots + a_0$
divise le polynôme $x^m + b_{m - 1} x^{m - 1} + \ldots + b_0$
dans $K[x]$, alors :
\begin{enumerate}
\item $a_0, \ldots, a_{n - 1}$ sont entiers sur tout sous-anneau
$R_0$ de $K$ contenant les éléments $b_0, \ldots, b_{m - 1}$, et
\item chaque $a_i$ appartient à $\sqrt{(b_0, \ldots, b_{m-1})R}$
pour tout sous-anneau $R \subset K$ contenant les éléments
$a_0, \ldots, a_{n - 1}, b_0, \ldots, b_{m - 1}$.
\end{enumerate}
\end{lemma}

\begin{proof}
Soit $L/K$ une extension de corps telle que nous puissions écrire
$x^m + b_{m - 1} x^{m - 1} + \ldots + b_0 =
\prod_{i = 1}^m (x - \beta_i)$ avec $\beta_i \in L$.
Voir Corps, section \ref{fields-section-splitting-fieds}.
Chaque $\beta_i$ est entier sur $R_0$.
Puisque chaque $a_i$ est un polynôme homogène en $\beta_1, \ldots, \beta_m$,
nous en déduisons la même propriété pour les $a_i$
(utiliser le lemme \ref{algebra-lemma-integral-closure-is-ring}). Cela démontre (1).

\medskip\noindent
Soit $R$ comme dans (2). Choisissons $c_0, \ldots, c_{m - n - 1} \in K$ tels que
$$
\begin{matrix}
x^m + b_{m - 1} x^{m - 1} + \ldots + b_0 =  \\
(x^n + a_{n - 1}x^{n - 1} + \ldots + a_0)
(x^{m - n} + c_{m - n - 1}x^{m - n - 1}+ \ldots + c_0).
\end{matrix}
$$
Cette équation implique
\begin{align*}
c_{m - n - 1} & = b_{m - 1} - a_{n - 1}, \\
c_{m - n - 2} & = b_{m - 2} - a_{n - 2} - a_{n - 1}c_{m - n - 1}, \\
\ldots
\end{align*}
Ainsi, $c_j \in R$ pour tout $j$. En quotientant par le radical
$\sqrt{(b_0, \ldots, b_{m - 1})}$, nous obtenons un anneau réduit $\overline{R}$.
Nous devons montrer que les images $\overline{a}_i \in \overline{R}$
sont nulles. Et, dans
$\overline{R}[x]$, nous avons la relation
$$
\begin{matrix}
x^m = x^m + \overline{b}_{m - 1} x^{m - 1} + \ldots + \overline{b}_0 = \\
(x^n + \overline{a}_{n - 1}x^{n - 1} + \ldots + \overline{a}_0)
(x^{m - n} + \overline{c}_{m - n - 1}x^{m - n - 1}+ \ldots + \overline{c}_0).
\end{matrix}
$$
Il est facile de voir que cela implique $\overline{a}_i = 0$ pour tout $i$. En effet,
d'après le lemme \ref{algebra-lemma-minimal-prime-reduced-ring}, la localisation de
$\overline{R}$ en un idéal premier minimal $\mathfrak{p}$ est un corps et
$\overline{R}_{\mathfrak p}[x]$ est un anneau factoriel. Ainsi,
$f = x^n + \sum \overline{a}_i x^i$
est associé à $x^n$ et, puisque $f$ est unitaire, $f = x^n$
dans $\overline{R}_{\mathfrak p}[x]$.
Il existe alors un $s \in \overline{R}$, $s \not\in \mathfrak p$,
tel que $s(f - x^n) = 0$. Par conséquent, tous les $\overline{a}_i$ appartiennent
à $\mathfrak p$, et nous concluons par le
lemme \ref{algebra-lemma-reduced-ring-sub-product-fields}.
\end{proof}

\begin{lemma}
\label{algebra-lemma-minimal-polynomial-normal-domain}
Soit $R \subset S$ une inclusion d'anneaux intègres.
Supposons que $R$ soit normal. Soit $g \in S$ entier
sur $R$. Alors le polynôme minimal de $g$
est à coefficients dans $R$.
\end{lemma}

\begin{proof}
Soit $P = x^m + b_{m-1} x^{m-1} + \ldots + b_0$
un polynôme à coefficients dans $R$
tel que $P(g) = 0$. Soit $Q = x^n + a_{n-1}x^{n-1} + \ldots + a_0$
le polynôme minimal de $g$ sur le corps des fractions
$K$ de $R$. Alors $Q$ divise $P$ dans $K[x]$. D'après le lemme
\ref{algebra-lemma-polynomials-divide}, nous voyons que les $a_i$ sont
entiers sur $R$. Puisque $R$ est normal, cela
signifie qu'ils appartiennent à $R$.
\end{proof}

\begin{proposition}
\label{algebra-proposition-going-down-normal-integral}
Soit $R \subset S$ une inclusion d'anneaux intègres.
Supposons que $R$ soit normal et que $S$ soit entier sur $R$.
Soient $\mathfrak p \subset \mathfrak p' \subset R$
des idéaux premiers. Soit $\mathfrak q'$ un idéal premier de $S$
tel que $\mathfrak p' = R \cap \mathfrak q'$.
Alors il existe un idéal premier $\mathfrak q$
tel que $\mathfrak q \subset \mathfrak q'$
et $\mathfrak p = R \cap \mathfrak q$. Autrement dit,
la propriété de descente est satisfaite pour $R \to S$, voir la
définition \ref{algebra-definition-going-up-down}.
\end{proposition}

\begin{proof}
Soient $\mathfrak p$, $\mathfrak p'$ et $\mathfrak q'$
comme dans l'énoncé. Nous devons montrer qu'il existe un idéal premier
$\mathfrak q$, avec $\mathfrak q \subset \mathfrak q'$ et
$R \cap \mathfrak q = \mathfrak p$. Cela revient
à trouver un idéal premier de
$S_{\mathfrak q'}$ qui s'envoie sur $\mathfrak p$.
D'après le lemme \ref{algebra-lemma-in-image}, nous devons montrer
que $\mathfrak p S_{\mathfrak q'} \cap R
= \mathfrak p$. Choisissons $z \in \mathfrak p S_{\mathfrak q'} \cap R$.
Nous pouvons écrire $z = y/g$ avec $y \in \mathfrak pS$ et
$g \in S$, $g \not\in \mathfrak q'$. Autrement dit,
nous avons $zg = y$.

\medskip\noindent
D'après le lemme \ref{algebra-lemma-integral-integral-over-ideal},
il existe un polynôme unitaire
$P = x^m + b_{m-1} x^{m-1} + \ldots + b_0$
avec $b_i \in \mathfrak p$ tel que $P(y) = 0$.

\medskip\noindent
D'après le lemme \ref{algebra-lemma-minimal-polynomial-normal-domain},
le polynôme minimal de $g$ sur $K$ est à coefficients
dans $R$. Écrivons-le $Q = x^n + a_{n-1} x^{n-1} + \ldots
+ a_0$. Remarquons que les $a_i$, $i = n-1, \ldots, 0$,
n'appartiennent pas tous à $\mathfrak p$, car cela impliquerait que
$g^n = \sum_{j < n} a_j g^j \in \mathfrak pS
\subset \mathfrak p'S
\subset \mathfrak q'$,
ce qui est une contradiction.

\medskip\noindent
Puisque $y = zg$, ce qui précède montre immédiatement
que $Q' = x^n + za_{n-1} x^{n-1} + \ldots + z^{n}a_0$
est le polynôme minimal de $y$. Ainsi,
$Q'$ divise $P$ et, d'après le lemme \ref{algebra-lemma-polynomials-divide},
nous voyons que $z^ja_{n - j} \in \sqrt{(b_0, \ldots, b_{m-1})}
\subset \mathfrak p$, $j =  1, \ldots, n$.
Puisque les $a_i$, $i = n-1, \ldots, 0$,
n'appartiennent pas tous à $\mathfrak p$, nous concluons que $z \in \mathfrak p$,
comme voulu.
\end{proof}














\section{Modules plats et morphismes d'anneaux plats}
\label{algebra-section-flat}

\noindent
Un résultat souvent utilisé est le suivant : si $M = \colim_{i\in \mathcal{I}} M_i$
est une colimite de $R$-modules et si $N$ est un $R$-module, alors
$$
M \otimes N
=
\colim_{i\in \mathcal{I}} M_i \otimes_R N,
$$
voir le lemme \ref{algebra-lemma-tensor-products-commute-with-limits}.
On exprime habituellement cette propriété en disant
que {\it $\otimes$ commute aux colimites}.
Un autre résultat souvent utilisé est le suivant : si $0 \to N_1 \to N_2 \to N_3 \to 0$
est une suite exacte et si $M$ est un $R$-module quelconque, alors
$$
M \otimes_R N_1
\to
M \otimes_R N_2
\to
M \otimes_R N_3
\to
0
$$
est encore exacte, voir le lemme \ref{algebra-lemma-tensor-product-exact}.
Ces deux propriétés nous disent que le foncteur
$N \mapsto M \otimes_R N$ {\it est exact à droite}.
Voir Catégories, section \ref{categories-section-exact-functor},
et Homologie, section \ref{homology-section-functors}.
Un $R$-module $M$ est plat si $N \mapsto N \otimes_R M$ est aussi exact à gauche,
c'est-à-dire s'il est exact. Voici la définition précise.

\begin{definition}
\label{algebra-definition-flat}
Soit $R$ un anneau.
\begin{enumerate}
\item Un $R$-module $M$ est dit {\it plat} si, chaque fois que
$N_1 \to N_2 \to N_3$ est une suite exacte de $R$-modules,
la suite $M \otimes_R N_1 \to M \otimes_R N_2 \to M \otimes_R N_3$
est également exacte.
\item Un $R$-module $M$ est dit {\it fidèlement plat} si le
complexe de $R$-modules
$N_1 \to N_2 \to N_3$ est exact si et seulement si
la suite $M \otimes_R N_1 \to M \otimes_R N_2 \to M \otimes_R N_3$
est exacte.
\item Un morphisme d'anneaux $R \to S$ est dit {\it plat} si
$S$ est plat comme $R$-module.
\item Un morphisme d'anneaux $R \to S$ est dit {\it fidèlement plat} si
$S$ est fidèlement plat comme $R$-module.
\end{enumerate}
\end{definition}

\noindent
Voici un exemple d'utilisation de la condition de platitude.

\begin{lemma}
\label{algebra-lemma-flat-intersect-ideals}
Soit $R$ un anneau. Soient $I, J \subset R$ des idéaux. Soit $M$ un
$R$-module plat. Alors $IM \cap JM = (I \cap J)M$.
\end{lemma}

\begin{proof}
Considérons la suite exacte $0 \to I \cap J \to R \to R/I \oplus R/J$.
En tensorisant par le module plat $M$, nous obtenons une suite exacte
$$
0 \to (I \cap J) \otimes_R M \to M \to M/IM \oplus M/JM
$$
Puisque le noyau de $M \to M/IM \oplus M/JM$ est égal à
$IM \cap JM$, nous concluons.
\end{proof}

\begin{lemma}
\label{algebra-lemma-colimit-flat}
Soit $R$ un anneau. Soit $\{M_i, \varphi_{ii'}\}$ un système filtrant de
$R$-modules plats. Alors $\colim_i M_i$ est un $R$-module plat.
\end{lemma}

\begin{proof}
Cela résulte du fait que $\otimes$ commute aux colimites et que
les colimites filtrantes sont exactes, voir le
lemme \ref{algebra-lemma-directed-colimit-exact}.
\end{proof}

\begin{lemma}
\label{algebra-lemma-composition-flat}
Une composée de morphismes d'anneaux (fidèlement) plats est
(fidèlement) plate.
Si $R \to R'$ est (fidèlement) plat et si $M'$ est un
$R'$-module (fidèlement) plat, alors $M'$ est un
$R$-module (fidèlement) plat.
\end{lemma}

\begin{proof}
La première assertion du lemme est un cas particulier de la
seconde, de sorte qu'il suffit clairement de démontrer cette dernière. Soit
$R \to R'$ un morphisme d'anneaux plat et soit $M'$ un $R'$-module plat.
Nous devons démontrer que $M'$ est un $R$-module plat. Soit
$N_1 \to N_2 \to N_3$ un complexe exact de $R$-modules.
Alors le complexe $R' \otimes_R N_1 \to
R' \otimes_R N_2 \to R' \otimes_R N_3$ est exact (puisque $R'$
est plat comme $R$-module), et le complexe
$M' \otimes_{R'} \left(R' \otimes_R N_1\right)
\to M' \otimes_{R'} \left(R' \otimes_R N_2\right)
\to M' \otimes_{R'} \left(R' \otimes_R N_3\right)$ est donc
exact (puisque $M'$ est un $R'$-module plat). Comme
$M' \otimes_{R'} \left(R' \otimes_R N\right)
\cong \left(M' \otimes_{R'} R'\right) \otimes_R N
\cong M' \otimes_R N$ pour tout $R$-module $N$, fonctoriellement
(d'après les lemmes \ref{algebra-lemma-tensor-with-bimodule} et
\ref{algebra-lemma-flip-tensor-product}), ce complexe est isomorphe
au complexe
$M' \otimes_R N_1 \to M' \otimes_R N_2 \to M' \otimes_R N_3$,
qui est donc lui aussi exact. Cela montre que $M'$ est un
$R$-module plat. En remontant cet argument, nous pouvons montrer
que, si $R \to R'$ est fidèlement plat et si $M'$ est
fidèlement plat comme $R'$-module, alors $M'$ est fidèlement
plat comme $R$-module.
\end{proof}

\begin{lemma}
\label{algebra-lemma-flat}
Soit $M$ un $R$-module. Les assertions suivantes sont équivalentes :
\begin{enumerate}
\item
\label{algebra-item-flat}
$M$ est plat sur $R$.
\item
\label{algebra-item-injective}
pour toute injection de $R$-modules $N \subset N'$,
l'application $N \otimes_R M \to N'\otimes_R M$ est injective.
\item
\label{algebra-item-f-ideal}
pour tout idéal $I \subset R$, l'application
$I \otimes_R M \to R \otimes_R M = M$ est injective.
\item
\label{algebra-item-ffg-ideal}
pour tout idéal de type fini $I \subset R$,
l'application $I \otimes_R M \to R \otimes_R M = M$ est injective.
\end{enumerate}
\end{lemma}

\begin{proof}
Les implications successives de (\ref{algebra-item-flat}) à (\ref{algebra-item-injective}), puis
à (\ref{algebra-item-f-ideal}) et à (\ref{algebra-item-ffg-ideal}), sont immédiates. Nous démontrons
donc que (\ref{algebra-item-ffg-ideal}) implique (\ref{algebra-item-flat}).
Supposons que $N_1 \to N_2 \to N_3$ soit exacte.
Posons $K = \Ker(N_2 \to N_3)$ et $Q = \Im(N_2 \to N_3)$.
Nous obtenons alors des applications
$$
N_1 \otimes_R M \to
K \otimes_R M \to
N_2 \otimes_R M \to
Q \otimes_R M \to
N_3 \otimes_R M
$$
Remarquons que les première et troisième flèches sont surjectives. Ainsi, si nous montrons
que les deuxième et quatrième flèches sont injectives, nous aurons
conclu\footnote{Voici l'argument plus en détail :
supposons que nous sachions que les deuxième et quatrième flèches sont
injectives. Le lemme \ref{algebra-lemma-tensor-product-exact} (appliqué
à la suite exacte $K \to N_2 \to Q \to 0$) donne que
la suite $K \otimes_R M \to N_2 \otimes_R M \to
Q \otimes_R M \to 0$ est exacte. Par conséquent,
$\Ker \left(N_2 \otimes_R M \to Q \otimes_R M\right)
= \Im \left(K \otimes_R M \to N_2 \otimes_R M\right)$.
Comme
$\Im \left(K \otimes_R M \to N_2 \otimes_R M\right)
= \Im \left(N_1 \otimes_R M \to N_2 \otimes_R M\right)$
(en raison de la surjectivité de $N_1 \otimes_R M \to
K \otimes_R M$) et
$\Ker \left(N_2 \otimes_R M \to Q \otimes_R M\right)
= \Ker \left(N_2 \otimes_R M \to N_3 \otimes_R M\right)$
(en raison de l'injectivité de $Q \otimes_R M \to
N_3 \otimes_R M$), cela devient
$\Ker \left(N_2 \otimes_R M \to N_3 \otimes_R M\right)
= \Im \left(N_1 \otimes_R M \to N_2 \otimes_R M\right)$,
ce qui montre que le foncteur $- \otimes_R M$ est exact,
et donc que $M$ est plat.}.
Il suffit donc de montrer que $- \otimes_R M$ transforme
les applications injectives de $R$-modules en applications injectives de $R$-modules.

\medskip\noindent
Supposons que $K \to N$ soit une application injective de $R$-modules et
soit $x \in \Ker(K \otimes_R M \to N \otimes_R M)$.
Nous devons montrer que $x$ est nul.
Le $R$-module $K$ est la réunion de ses
$R$-sous-modules de type fini ; par conséquent, $K \otimes_R M$ est
la colimite des $R$-modules de la forme
$K_i \otimes_R M$, où $K_i$ parcourt tous les
$R$-sous-modules de type fini de $K$
(car le produit tensoriel commute aux colimites).
Ainsi, pour un certain $i$, notre $x$ provient d'un élément
$x_i \in K_i \otimes_R M$. Nous pouvons donc supposer que $K$
est un $R$-module de type fini. Faisons cette hypothèse. Nous considérons
l'injection $K \to N$ comme une inclusion, de sorte que
$K \subset N$.

\medskip\noindent
Le $R$-module $N$ est la réunion de ses
$R$-sous-modules de type fini qui contiennent $K$. Par conséquent, $N \otimes_R M$
est la colimite des $R$-modules de la forme
$N_i \otimes_R M$, où $N_i$ parcourt tous les
$R$-sous-modules de type fini de $N$ qui contiennent $K$
(de nouveau parce que le produit tensoriel commute aux colimites).
Remarquons qu'il s'agit d'une colimite sur un système filtrant
(puisque la somme de deux sous-modules de type fini de $N$ est
encore de type fini).
Par conséquent, (d'après le lemme \ref{algebra-lemma-zero-directed-limit}),
l'élément $x \in K \otimes_R M$ s'envoie sur
zéro dans au moins un de ces $R$-modules
$N_i \otimes_R M$ (puisque $x$ s'envoie sur zéro
dans $N \otimes_R M$).
Nous pouvons donc supposer que $N$ est un $R$-module de type fini.

\medskip\noindent
Supposons $N$ de type fini sur $R$. Écrivons $N = R^{\oplus n}/L$ et $K = L'/L$
pour certains $L \subset L' \subset R^{\oplus n}$.
Pour tout $R$-sous-module $G \subset R^{\oplus n}$,
nous disposons d'une application canonique $G \otimes_R M \to M^{\oplus n}$
obtenue en composant
$G \otimes_R M \to R^n \otimes_R M = M^{\oplus n}$.
Il suffit de prouver que $L \otimes_R M \to M^{\oplus n}$
et $L' \otimes_R M \to M^{\oplus n}$ sont injectives.
En effet, si tel est le cas, nous voyons que
$K \otimes_R M = L' \otimes_R M/L \otimes_R M \to M^{\oplus n}/L \otimes_R M$
est elle aussi injective\footnote{Cela devient évident si nous
identifions $L' \otimes_R M$ et $L \otimes_R M$ à des
sous-modules de $M^{\oplus n}$ (ce qui est légitime puisque
les applications $L \otimes_R M \to M^{\oplus n}$
et $L' \otimes_R M \to M^{\oplus n}$ sont injectives et
commutent avec l'application évidente $L' \otimes_R M \to L \otimes_R M$).}.

\medskip\noindent
Il suffit donc de montrer que $L \otimes_R M \to M^{\oplus n}$
est injective lorsque $L \subset R^{\oplus n}$ est un $R$-sous-module.
Nous le faisons par récurrence sur $n$. Nous traitons ci-dessous le cas initial $n = 1$.
Pour l'étape de récurrence, supposons $n > 1$ et posons
$L' = L \cap R \oplus 0^{\oplus n - 1}$. Alors $L'' = L/L'$ est un sous-module
de $R^{\oplus n - 1}$. Nous obtenons un diagramme
$$
\xymatrix{
&
L' \otimes_R M \ar[r] \ar[d] &
L \otimes_R M \ar[r] \ar[d] &
L'' \otimes_R M \ar[r] \ar[d] &
0 \\
0 \ar[r] &
M \ar[r] &
M^{\oplus n} \ar[r] &
M^{\oplus n - 1} \ar[r] & 0
}
$$
D'après l'hypothèse de récurrence et le cas initial, les flèches verticales
de gauche et de droite sont injectives. Les lignes sont exactes. Il s'ensuit que la flèche verticale
du milieu est elle aussi injective.

\medskip\noindent
Le cas initial de la récurrence ci-dessus est celui où $L \subset R$ est un idéal.
Autrement dit, nous devons montrer que $I \otimes_R M \to M$ est injective
pour tout idéal $I$ de $R$. Nous savons que cela est vrai lorsque $I$ est de type
fini. Toutefois, $I = \bigcup I_\alpha$ est la réunion des
idéaux de type fini $I_\alpha$ qu'il contient. Autrement dit,
$I = \colim I_\alpha$. Puisque $\otimes$ commute aux colimites, nous voyons que
$I \otimes_R M = \colim I_\alpha \otimes_R M$ et, puisque toutes
les applications $I_\alpha \otimes_R M \to M$ sont injectives par
hypothèse, il en va de même de $I \otimes_R M \to M$.
\end{proof}

\begin{lemma}
\label{algebra-lemma-colimit-rings-flat}
Soit $\{R_i, \varphi_{ii'}\}$ un système d'anneaux sur l'ensemble filtrant $I$.
Soit $R = \colim_i R_i$.
\begin{enumerate}
\item Si $M$ est un $R$-module tel que $M$ soit plat comme $R_i$-module
pour tout $i$, alors $M$ est plat comme $R$-module.
\item Pour $i \in I$, soit $M_i$ un $R_i$-module plat et,
pour $i' \geq i$, soit $f_{ii'} : M_i \to M_{i'}$ une application
$\varphi_{ii'}$-linéaire telle que $f_{i' i''} \circ f_{i i'} = f_{i i''}$.
Alors $M = \colim_{i \in I} M_i$ est un $R$-module plat.
\end{enumerate}
\end{lemma}

\begin{proof}
La partie (1) est un cas particulier de la partie (2), avec $M_i = M$ pour tout $i$
et $f_{i i'} = \text{id}_M$. Démontrons (2).
Soit $\mathfrak a \subset R$ un idéal de type fini. D'après le
lemme \ref{algebra-lemma-flat},
il suffit de montrer que $\mathfrak a \otimes_R M \to M$ est
injective. Nous pouvons trouver $i \in I$ et un idéal de type fini
$\mathfrak a' \subset R_i$ tels que $\mathfrak a = \mathfrak a'R$.
Alors $\mathfrak a = \colim_{i' \geq i} \mathfrak a'R_{i'}$.
Puisque $\otimes$ commute aux colimites, l'application
$\mathfrak a \otimes_R M \to M$ est la colimite des applications
$$
\mathfrak a'R_{i'} \otimes_{R_{i'}} M_{i'} \longrightarrow M_{i'}
$$
Toutes ces applications sont injectives par hypothèse. Puisque les colimites sur $I$
sont exactes d'après le lemme \ref{algebra-lemma-directed-colimit-exact}, nous concluons.
\end{proof}

\begin{lemma}
\label{algebra-lemma-flat-base-change}
Supposons que $M$ soit (fidèlement) plat sur $R$ et que $R \to R'$
soit un morphisme d'anneaux. Alors $M \otimes_R R'$ est (fidèlement) plat sur $R'$.
\end{lemma}

\begin{proof}
Pour tout $R'$-module $N$, nous avons un
isomorphisme canonique $N \otimes_{R'} (R'\otimes_R M)
= N \otimes_R M$. Les propriétés d'exactitude voulues du foncteur
$-\otimes_{R'}(R'\otimes_R M)$ résultent donc
des propriétés d'exactitude correspondantes du foncteur $-\otimes_R M$.
\end{proof}

\begin{lemma}
\label{algebra-lemma-flatness-descends}
Soit $R \to R'$ un morphisme d'anneaux fidèlement plat.
Soit $M$ un module sur $R$, et posons $M' = R' \otimes_R M$.
Alors $M$ est plat sur $R$ si et seulement si $M'$ est plat sur $R'$.
\end{lemma}

\begin{proof}
D'après le lemme \ref{algebra-lemma-flat-base-change}, si $M$ est plat,
alors $M'$ est plat. Réciproquement, supposons $M'$ plat.
Soit $N_1 \to N_2 \to N_3$ une suite exacte de $R$-modules.
Nous voulons montrer que $N_1 \otimes_R M \to N_2 \otimes_R M \to N_3 \otimes_R M$
est exacte. Nous savons que
$N_1 \otimes_R R' \to N_2 \otimes_R R' \to N_3 \otimes_R R'$ est
exacte, car $R \to R'$ est plat. La platitude de $M'$ implique que
$N_1 \otimes_R R' \otimes_{R'} M'
\to N_2 \otimes_R R' \otimes_{R'} M'
\to N_3 \otimes_R R' \otimes_{R'} M'$ est exacte.
Nous pouvons l'écrire
$N_1 \otimes_R M \otimes_R R'
\to N_2 \otimes_R M \otimes_R R'
\to N_3 \otimes_R M \otimes_R R'$.
Enfin, la fidèle platitude implique que
$N_1 \otimes_R M \to N_2 \otimes_R M \to N_3 \otimes_R M$
est exacte.
\end{proof}

\begin{lemma}
\label{algebra-lemma-flatness-descends-more-general}
Soit $R$ un anneau. Soit $S \to S'$ un morphisme plat de $R$-algèbres.
Soit $M$ un module sur $S$, et posons $M' = S' \otimes_S M$.
\begin{enumerate}
\item Si $M$ est plat sur $R$, alors $M'$ est plat sur $R$.
\item Si $S \to S'$ est fidèlement plat, alors $M$ est plat
sur $R$ si et seulement si $M'$ est plat sur $R$.
\end{enumerate}
\end{lemma}

\begin{proof}
Soit $N \to N'$ une injection de $R$-modules. Par la platitude
de $S \to S'$, nous avons
$$
\Ker(N \otimes_R M \to N' \otimes_R M) \otimes_S S'
=
\Ker(N \otimes_R M' \to N' \otimes_R M')
$$
Si $M$ est plat sur $R$, alors le membre de gauche est nul et
nous trouvons que $M'$ est plat sur $R$ par la deuxième caractérisation
de la platitude dans le lemme \ref{algebra-lemma-flat}.
Si $M'$ est plat sur $R$, le membre de droite est nul et, si
en outre $S \to S'$ est fidèlement plat, cela implique que
$\Ker(N \otimes_R M \to N' \otimes_R M)$ est nul, ce qui
montre à son tour que $M$ est plat sur $R$.
\end{proof}

\begin{lemma}
\label{algebra-lemma-flat-permanence}
Soit $R \to S$ un morphisme d'anneaux. Soit $M$ un $S$-module.
Si $M$ est plat comme $R$-module et fidèlement plat comme $S$-module,
alors $R \to S$ est plat.
\end{lemma}

\begin{proof}
Soit $N_1 \to N_2 \to N_3$ une suite exacte de $R$-modules.
Par hypothèse, $N_1 \otimes_R M \to N_2 \otimes_R M \to N_3 \otimes_R M$
est exacte. Nous pouvons l'écrire
$$
N_1 \otimes_R S \otimes_S M
\to
N_2 \otimes_R S \otimes_S M
\to
N_3 \otimes_R S \otimes_S M.
$$
Par fidèle platitude de $M$ sur $S$, nous concluons que
$N_1 \otimes_R S \to N_2 \otimes_R S \to N_3 \otimes_R S$ est exacte.
Ainsi $R \to S$ est plat.
\end{proof}

\noindent
Soit $R$ un anneau.
Soit $M$ un $R$-module.
Soit $\sum f_i x_i = 0$ une relation dans $M$.
Nous disons que la relation $\sum f_i x_i$
est {\it triviale} s'il existe un entier $m \geq 0$,
des éléments $y_j \in M$, $j = 1, \ldots, m$, et des éléments $a_{ij} \in R$,
$i = 1, \ldots, n$, $j = 1, \ldots, m$, tels que
$$
x_i = \sum\nolimits_j a_{ij} y_j, \forall i,
\quad\text{et}\quad
0 = \sum\nolimits_i f_ia_{ij}, \forall j.
$$

\begin{lemma}[Critère équationnel de platitude]
\label{algebra-lemma-flat-eq}
Un module $M$ sur $R$ est plat si et seulement si
toute relation dans $M$ est triviale.
\end{lemma}

\begin{proof}
Supposons $M$ plat et soit $\sum f_i x_i = 0$ une relation dans $M$.
Posons $I = (f_1, \ldots, f_n)$ et
$K = \Ker(R^n \to I, (a_1, \ldots, a_n) \mapsto \sum_i a_i f_i)$.
Nous avons donc la suite exacte courte
$0 \to K \to R^n \to I \to 0$. Alors $\sum f_i \otimes x_i$
est un élément de $I \otimes_R M$ qui s'envoie
sur zéro dans $R \otimes_R M = M$. Par platitude,
$\sum f_i \otimes x_i$ est nul dans $I \otimes_R M$.
Il existe donc un élément de $K \otimes_R M$ dont l'image
est $\sum e_i \otimes x_i \in R^n \otimes_R M$, où $e_i$
est le $i$-ième vecteur de base de $R^n$.
Écrivons cet élément sous la forme $\sum k_j \otimes y_j$,
puis l'image de $k_j$ dans $R^n$ sous la forme
$\sum a_{ij} e_i$, ce qui donne le résultat.

\medskip\noindent
Supposons toute relation triviale, soit $I$
un idéal de type fini, et soit $x = \sum f_i \otimes x_i$
un élément de $I \otimes_R M$ qui s'envoie sur zéro dans $R \otimes_R M = M$.
Cela signifie exactement que $\sum f_i x_i$ est une relation dans
$M$. Le fait qu'elle soit triviale implique aisément que
$x$ est nul, car
$$
x
=
\sum f_i \otimes x_i
=
\sum f_i \otimes \left(\sum a_{ij}y_j\right)
=
\sum \left(\sum f_i a_{ij}\right) \otimes y_j
=
0
$$
\end{proof}

\begin{lemma}
\label{algebra-lemma-flat-tor-zero}
Supposons que $R$ soit un anneau, que $0 \to M'' \to M' \to M \to 0$
soit une suite exacte courte et que $N$ soit un $R$-module. Si $M$ est plat,
alors $N \otimes_R M'' \to N \otimes_R M'$ est injective, c'est-à-dire que la
suite
$$
0 \to N \otimes_R M'' \to N \otimes_R M' \to N \otimes_R M \to 0
$$
est une suite exacte courte.
\end{lemma}

\begin{proof}
Soit $R^{(I)} \to N$ une surjection d'un module libre
sur $N$, de noyau $K$. Le résultat découle
du lemme du serpent appliqué au diagramme suivant
$$
\begin{matrix}
 & & 0 & & 0 & & 0 & & \\
 & & \uparrow & & \uparrow & & \uparrow & & \\
 & & M''\otimes_R N & \to & M' \otimes_R N & \to & M \otimes_R N & \to & 0 \\
 & & \uparrow & & \uparrow & & \uparrow & & \\
0 & \to & (M'')^{(I)} & \to & (M')^{(I)} & \to & M^{(I)} & \to & 0 \\
 & & \uparrow & & \uparrow & & \uparrow & & \\
 & & M''\otimes_R K & \to & M' \otimes_R K & \to & M \otimes_R K & \to & 0 \\
 & & & & & & \uparrow & & \\
 & & & & & & 0 & &
\end{matrix}
$$
dont les lignes et les colonnes sont exactes. La ligne du milieu est exacte car tensoriser
par le module libre $R^{(I)}$ est exact.
\end{proof}


\begin{lemma}
\label{algebra-lemma-flat-ses}
Supposons que $0 \to M' \to M \to M'' \to 0$ soit
une suite exacte courte de $R$-modules.
Si $M'$ et $M''$ sont plats, alors $M$ l'est aussi.
Si $M$ et $M''$ sont plats, alors $M'$ l'est aussi.
\end{lemma}

\begin{proof}
Nous utiliserons le critère suivant : un module $N$ est plat si, pour
tout idéal $I \subset R$, l'application $N \otimes_R I \to N$ est injective ;
voir le lemme \ref{algebra-lemma-flat}.
Considérons un idéal $I \subset R$.
Considérons le diagramme
$$
\begin{matrix}
0 & \to & M' & \to & M & \to & M'' & \to & 0 \\
& & \uparrow & & \uparrow & & \uparrow & & \\
& & M'\otimes_R I & \to & M \otimes_R I & \to & M''\otimes_R I & \to & 0
\end{matrix}
$$
dont les lignes sont exactes. Cela démontre immédiatement la première assertion.
La seconde résulte du fait que, si $M''$ est plat, alors la flèche horizontale
inférieure gauche est injective d'après le lemme \ref{algebra-lemma-flat-tor-zero}.
\end{proof}

\begin{lemma}
\label{algebra-lemma-easy-ff}
Soit $R$ un anneau.
Soit $M$ un $R$-module.
Les assertions suivantes sont équivalentes :
\begin{enumerate}
\item $M$ est fidèlement plat, et
\item $M$ est plat et, pour tous les homomorphismes de $R$-modules $\alpha : N \to N'$,
nous avons $\alpha = 0$ si et seulement si $\alpha \otimes \text{id}_M = 0$.
\end{enumerate}
\end{lemma}

\begin{proof}
Si $M$ est fidèlement plat, alors
$0 \to \Ker(\alpha) \to N \to N'$ est exacte si et seulement si elle le reste
après tensorisation par $M$. Cela démontre que (1) implique (2).
Réciproquement, supposons (2). Soit $N_1 \to N_2 \to N_3$
un complexe et supposons que le complexe
$N_1 \otimes_R M \to N_2 \otimes_R M \to N_3\otimes_R M$
soit exact. Prenons $x \in \Ker(N_2 \to N_3)$
et considérons l'application $\alpha : R \to N_2/\Im(N_1)$,
$r \mapsto rx + \Im(N_1)$. Par l'exactitude
du complexe $-\otimes_R M$, nous voyons que $\alpha \otimes
\text{id}_M$ est nulle. Par hypothèse, nous obtenons que $\alpha$ est
nulle. Ainsi $x $ appartient à l'image de $N_1 \to N_2$.
\end{proof}

\begin{lemma}
\label{algebra-lemma-ff}
\begin{slogan}
Un module plat est fidèlement plat si et seulement si ses fibres sont non nulles.
\end{slogan}
Soit $M$ un $R$-module plat.
Les assertions suivantes sont équivalentes :
\begin{enumerate}
\item $M$ est fidèlement plat,
\item pour tout $R$-module non nul $N$, le produit tensoriel $M \otimes_R N$
est non nul,
\item pour tout $\mathfrak p \in \Spec(R)$,
le produit tensoriel $M \otimes_R \kappa(\mathfrak p)$ est non nul, et
\item pour tout idéal maximal $\mathfrak m$ de $R$,
le produit tensoriel $M \otimes_R \kappa(\mathfrak m) = M/{\mathfrak m}M$
est non nul.
\end{enumerate}
\end{lemma}

\begin{proof}
Supposons $M$ fidèlement plat et $N \not = 0$. D'après le lemme \ref{algebra-lemma-easy-ff},
l'application non nulle $1 : N \to N$ induit une application non nulle
$M \otimes_R N \to M \otimes_R N$, donc $M \otimes_R N \not = 0$.
Ainsi (1) implique (2). Les implications (2) $\Rightarrow$ (3) $\Rightarrow$ (4)
sont immédiates.

\medskip\noindent
Supposons (4). Supposons que $N_1 \to N_2 \to N_3$ soit un complexe et
que $N_1 \otimes_R M \to N_2\otimes_R M \to
N_3\otimes_R M$ soit exact. Soit $H$ l'homologie du complexe,
donc $H = \Ker(N_2 \to N_3)/\Im(N_1 \to N_2)$. Pour achever la preuve,
nous allons montrer que $H = 0$. Par platitude, nous voyons que $H \otimes_R M = 0$.
Prenons $x \in H$ et soit $I = \{f \in R \mid fx = 0 \}$
son annulateur. Puisque $R/I \subset H$, nous obtenons
$M/IM \subset H \otimes_R M = 0$ par platitude de $M$.
Si $I \not =  R$, nous pouvons choisir
un idéal maximal $I \subset \mathfrak m \subset R$.
Cela donne immédiatement une contradiction.
\end{proof}

\begin{lemma}
\label{algebra-lemma-ff-rings}
Soit $R \to S$ un morphisme d'anneaux plat.
Les assertions suivantes sont équivalentes :
\begin{enumerate}
\item $R \to S$ est fidèlement plat,
\item l'application induite sur $\Spec$ est surjective, et
\item tout point fermé $x \in \Spec(R)$ appartient
à l'image de l'application $\Spec(S) \to \Spec(R)$.
\end{enumerate}
\end{lemma}

\begin{proof}
Cela résulte rapidement du lemme \ref{algebra-lemma-ff}, car nous
avons vu dans la remarque \ref{algebra-remark-fundamental-diagram}
que $\mathfrak p$ appartient à l'image
si et seulement si l'anneau $S \otimes_R \kappa(\mathfrak p)$
est non nul.
\end{proof}

\begin{lemma}
\label{algebra-lemma-local-flat-ff}
Un homomorphisme local plat d'anneaux locaux est fidèlement plat.
\end{lemma}

\begin{proof}
Cela résulte immédiatement du lemme \ref{algebra-lemma-ff-rings}.
\end{proof}

\noindent
La platitude se comporte bien avec la localisation.

\begin{lemma}
\label{algebra-lemma-flat-localization}
Soit $R$ un anneau. Soit $S \subset R$ une partie multiplicative.
\begin{enumerate}
\item La localisation $S^{-1}R$ est une $R$-algèbre plate.
\item Si $M$ est un $S^{-1}R$-module, alors $M$ est un $R$-module plat
si et seulement si $M$ est un $S^{-1}R$-module plat.
\item Supposons que $M$ soit un $R$-module. Alors
$M$ est un $R$-module plat si et seulement si $M_{\mathfrak p}$ est un
$R_{\mathfrak p}$-module plat pour tout idéal premier $\mathfrak p$ de $R$.
\item Supposons que $M$ soit un $R$-module. Alors $M$ est un $R$-module plat si
et seulement si $M_{\mathfrak m}$ est un
$R_{\mathfrak m}$-module plat pour tout idéal maximal $\mathfrak m$ de $R$.
\item Supposons que $R \to A$ soit un morphisme d'anneaux, que $M$ soit un $A$-module,
et que $g_1, \ldots, g_m \in A$ soient des éléments qui engendrent l'idéal
unité de $A$. Alors $M$ est plat sur $R$ si et seulement si chaque localisation
$M_{g_i}$ est plate sur $R$.
\item Supposons que $R \to A$ soit un morphisme d'anneaux et que $M$ soit un $A$-module.
Alors $M$ est un $R$-module plat si et seulement si la localisation
$M_{\mathfrak q}$ est un $R_{\mathfrak p}$-module plat
(où $\mathfrak p$ est l'idéal premier de $R$ situé sous $\mathfrak q$)
pour tout idéal premier $\mathfrak q$ de $A$.
\item Supposons que $R \to A$ soit un morphisme d'anneaux et que $M$ soit un $A$-module.
Alors $M$ est un $R$-module plat si et seulement si la localisation
$M_{\mathfrak m}$ est un $R_{\mathfrak p}$-module plat
(où $\mathfrak p = R \cap \mathfrak m$)
pour tout idéal maximal $\mathfrak m$ de $A$.
\end{enumerate}
\end{lemma}

\begin{proof}
Démontrons la dernière assertion du lemme.
Dans la preuve, nous utiliserons à plusieurs reprises le fait que la localisation est exacte
et commute au produit tensoriel, voir les sections \ref{algebra-section-localization}
et \ref{algebra-section-tensor-product}.

\medskip\noindent
Supposons que $R \to A$ soit un morphisme d'anneaux et que $M$ soit un $A$-module.
Supposons que $M_{\mathfrak m}$ soit un $R_{\mathfrak p}$-module plat
pour tout idéal maximal $\mathfrak m$ de $A$ (avec
$\mathfrak p = R \cap \mathfrak m$). Soit $I \subset R$ un idéal.
Nous devons montrer que l'application $I \otimes_R M \to M$ est injective.
Nous pouvons la considérer comme une application de $A$-modules.
Par hypothèse, la localisation
$(I \otimes_R M)_{\mathfrak m} \to M_{\mathfrak m}$ est injective
car
$(I \otimes_R M)_{\mathfrak m} =
I_{\mathfrak p} \otimes_{R_{\mathfrak p}} M_{\mathfrak m}$.
Ainsi, le noyau de $I \otimes_R M \to M$ est nul d'après le
lemme \ref{algebra-lemma-characterize-zero-local}.
Par conséquent, $M$ est plat sur $R$.

\medskip\noindent
Réciproquement, supposons $M$ plat sur $R$. Prenons un idéal premier
$\mathfrak q$ de $A$ situé au-dessus de l'idéal premier $\mathfrak p$ de $R$.
Supposons que $I \subset R_{\mathfrak p}$ soit un idéal. Nous devons montrer que
$I \otimes_{R_{\mathfrak p}} M_{\mathfrak q} \to M_{\mathfrak q}$
est injective. Nous pouvons écrire $I = J_{\mathfrak p}$ pour un certain
idéal $J \subset R$. Alors l'application
$I \otimes_{R_{\mathfrak p}} M_{\mathfrak q} \to M_{\mathfrak q}$
est simplement la localisation (en $\mathfrak q$) de l'application
$J \otimes_R M \to M$, qui est injective. Puisque la localisation est exacte,
nous voyons que $M_{\mathfrak q}$ est un $R_{\mathfrak p}$-module plat.

\medskip\noindent
Cela démontre (7) et (6). Les autres assertions résultent de manière immédiate
de la dernière assertion (preuves omises).
\end{proof}

\begin{lemma}
\label{algebra-lemma-flat-going-down}
Soit $R \to S$ plat. Soient $\mathfrak p \subset \mathfrak p'$
des idéaux premiers de $R$. Soit $\mathfrak q' \subset S$ un idéal premier de $S$
qui s'envoie sur $\mathfrak p'$. Alors il existe un idéal premier
$\mathfrak q \subset \mathfrak q'$ qui s'envoie sur $\mathfrak p$.
\end{lemma}

\begin{proof}
D'après le lemme \ref{algebra-lemma-flat-localization}, le morphisme d'anneaux locaux
$R_{\mathfrak p'} \to S_{\mathfrak q'}$ est plat.
D'après le lemme \ref{algebra-lemma-local-flat-ff}, ce morphisme d'anneaux locaux est fidèlement
plat. D'après le lemme \ref{algebra-lemma-ff-rings}, il existe un idéal premier qui s'envoie sur
$\mathfrak p R_{\mathfrak p'}$. L'image réciproque de cet
idéal premier dans $S$ convient.
\end{proof}

\noindent
La propriété de $R \to S$ décrite dans le lemme est appelée
« propriété de descente ». Voir la définition \ref{algebra-definition-going-up-down}.

\begin{lemma}
\label{algebra-lemma-colimit-faithfully-flat}
Soit $R$ un anneau. Soit $\{S_i, \varphi_{ii'}\}$ un système filtrant de
$R$-algèbres fidèlement plates. Alors $S = \colim_i S_i$ est une
$R$-algèbre fidèlement plate.
\end{lemma}

\begin{proof}
D'après le lemme \ref{algebra-lemma-colimit-flat}, nous voyons que $S$ est plate.
Soit $\mathfrak m \subset R$ un idéal maximal. D'après le
lemme \ref{algebra-lemma-ff-rings},
aucun des anneaux $S_i/\mathfrak m S_i$ n'est nul.
Ainsi $S/\mathfrak mS = \colim S_i/\mathfrak mS_i$ est lui aussi non nul,
car $1$ n'est pas égal à zéro. L'image de
$\Spec(S) \to \Spec(R)$ contient donc $\mathfrak m$, et nous voyons que
$R \to S$ est fidèlement plat d'après le lemme \ref{algebra-lemma-ff-rings}.
\end{proof}





\section{Supports et annulateurs}
\label{algebra-section-supp-and-ann}

\noindent
Quelques définitions et lemmes très élémentaires.

\begin{definition}
\label{algebra-definition-support-module}
Soit $R$ un anneau et soit $M$ un $R$-module.
Le {\it support de $M$} est l'ensemble
$$
\text{Supp}(M)
=
\{
\mathfrak p \in \Spec(R)
\mid
M_{\mathfrak p} \not = 0
\}
$$
\end{definition}

\begin{lemma}
\label{algebra-lemma-support-zero}
\begin{slogan}
Un module sur un anneau est de support vide si et seulement s'il est nul.
\end{slogan}
Soit $R$ un anneau. Soit $M$ un $R$-module. Alors
$$
M = (0) \Leftrightarrow \text{Supp}(M) = \emptyset.
$$
\end{lemma}

\begin{proof}
En fait, le
lemme \ref{algebra-lemma-characterize-zero-local}
montre même que $\text{Supp}(M)$ contient toujours un idéal maximal
si $M$ n'est pas nul.
\end{proof}

\begin{definition}
\label{algebra-definition-annihilator}
Soit $R$ un anneau. Soit $M$ un $R$-module.
\begin{enumerate}
\item Étant donné un élément $m \in M$, l'{\it annulateur de $m$}
est l'idéal
$$
\text{Ann}_R(m) = \text{Ann}(m) = \{f \in R \mid fm = 0\}.
$$
\item L'{\it annulateur de $M$}
est l'idéal
$$
\text{Ann}_R(M) = \text{Ann}(M) = \{f \in R \mid fm = 0\ \forall m \in M\}.
$$
\end{enumerate}
\end{definition}

\begin{lemma}
\label{algebra-lemma-annihilator-flat-base-change}
Soit $R \to S$ un morphisme d'anneaux plat. Soit $M$ un $R$-module et
$m \in M$. Alors $\text{Ann}_R(m) S = \text{Ann}_S(m \otimes 1)$.
Si $M$ est un $R$-module de type fini, alors
$\text{Ann}_R(M) S = \text{Ann}_S(M \otimes_R S)$.
\end{lemma}

\begin{proof}
Posons $I = \text{Ann}_R(m)$. Par définition, il existe une suite exacte
$0 \to I \to R \to M$, où l'application $R \to M$ envoie $f$ sur $fm$. En utilisant
la platitude, nous obtenons une suite exacte
$0 \to I \otimes_R S \to S \to M \otimes_R S$, ce qui démontre la première
assertion. Si $m_1, \ldots, m_n$ est un système de générateurs de $M$,
alors $\text{Ann}_R(M) = \bigcap \text{Ann}_R(m_i)$. De même,
$\text{Ann}_S(M \otimes_R S) = \bigcap \text{Ann}_S(m_i \otimes 1)$.
Posons $I_i = \text{Ann}_R(m_i)$. Il suffit alors de montrer que
$\bigcap_{i = 1, \ldots, n} (I_i S) = (\bigcap_{i = 1, \ldots, n} I_i)S$.
C'est le lemme \ref{algebra-lemma-flat-intersect-ideals}.
\end{proof}

\begin{lemma}
\label{algebra-lemma-support-closed}
Soit $R$ un anneau et soit $M$ un $R$-module.
Si $M$ est de type fini, alors $\text{Supp}(M)$ est fermé.
Plus précisément, si $I = \text{Ann}(M)$ est l'annulateur de $M$, alors
$V(I) = \text{Supp}(M)$.
\end{lemma}

\begin{proof}
Nous allons montrer que $V(I) = \text{Supp}(M)$.

\medskip\noindent
Supposons $\mathfrak p \in \text{Supp}(M)$. Alors $M_{\mathfrak p} \not = 0$.
Choisissons un élément $m \in M$ dont l'image dans $M_\mathfrak p$ est non nulle.
L'annulateur de $m$ est alors contenu dans $\mathfrak p$ par construction
de la localisation $M_\mathfrak p$. Par conséquent,
à plus forte raison, $I = \text{Ann}(M)$ doit être contenu dans $\mathfrak p$.

\medskip\noindent
Réciproquement, supposons que $\mathfrak p \not \in \text{Supp}(M)$.
Alors $M_{\mathfrak p} = 0$.
Soient $x_1, \ldots, x_r \in M$ des générateurs.
D'après le lemme \ref{algebra-lemma-localization-colimit}, il existe
$f \in R$, $f\not\in \mathfrak p$, tel que
$x_i/1 = 0$ dans $M_f$. Ainsi $f^{n_i} x_i = 0$ pour un certain $n_i \geq 1$.
Par conséquent $f^nM = 0$ pour $n = \max\{n_i\}$, comme voulu.
\end{proof}

\begin{lemma}
\label{algebra-lemma-support-base-change}
Soit $R \to R'$ un morphisme d'anneaux et soit $M$ un $R$-module de type fini.
Alors $\text{Supp}(M \otimes_R R')$ est l'image réciproque de
$\text{Supp}(M)$.
\end{lemma}

\begin{proof}
Soit $\mathfrak p \in \text{Supp}(M)$. D'après le lemme de Nakayama
(lemme \ref{algebra-lemma-NAK}), nous voyons que
$$
M \otimes_R \kappa(\mathfrak p) = M_\mathfrak p/\mathfrak p M_\mathfrak p
$$
est un espace vectoriel non nul sur $\kappa(\mathfrak p)$.
Ainsi, pour tout idéal premier $\mathfrak p' \subset R'$ situé
au-dessus de $\mathfrak p$, nous voyons que
$$
(M \otimes_R R')_{\mathfrak p'}/\mathfrak p' (M \otimes_R R')_{\mathfrak p'} =
(M \otimes_R R') \otimes_{R'} \kappa(\mathfrak p') =
M \otimes_R \kappa(\mathfrak p) \otimes_{\kappa(\mathfrak p)}
\kappa(\mathfrak p')
$$
est non nul. Cela implique $\mathfrak p' \in \text{Supp}(M \otimes_R R')$.
Réciproquement, si $\mathfrak p' \subset R'$ est un idéal premier situé
au-dessus d'un idéal premier quelconque $\mathfrak p \subset R$, alors
$$
(M \otimes_R R')_{\mathfrak p'} =
M_\mathfrak p \otimes_{R_\mathfrak p} R'_{\mathfrak p'}.
$$
Par conséquent, si $\mathfrak p' \in \text{Supp}(M \otimes_R R')$
est situé au-dessus de l'idéal premier $\mathfrak p \subset R$, alors
$\mathfrak p \in \text{Supp}(M)$.
\end{proof}

\begin{lemma}
\label{algebra-lemma-support-element}
Soit $R$ un anneau, soit $M$ un $R$-module et soit $m \in M$.
Alors $\mathfrak p \in V(\text{Ann}(m))$ si et seulement si
$m$ ne s'envoie pas sur zéro dans $M_\mathfrak p$.
\end{lemma}

\begin{proof}
Nous pouvons remplacer $M$ par $Rm \subset M$. Alors (1) $\text{Ann}(m) = \text{Ann}(M)$
et (2) $m$ ne s'envoie pas sur zéro dans $M_\mathfrak p$ si et seulement si
$\mathfrak p \in \text{Supp}(M)$.
Le résultat découle maintenant du lemme \ref{algebra-lemma-support-closed}.
\end{proof}

\begin{lemma}
\label{algebra-lemma-support-finite-presentation-constructible}
Soit $R$ un anneau et soit $M$ un $R$-module.
Si $M$ est un $R$-module de présentation finie, alors $\text{Supp}(M)$ est une
partie fermée de $\Spec(R)$ dont le complémentaire est quasi-compact.
\end{lemma}

\begin{proof}
Choisissons une présentation
$$
R^{\oplus m} \longrightarrow R^{\oplus n} \longrightarrow M \to 0
$$
Soit $A \in \text{Mat}(n \times m, R)$ la matrice de la première
application. D'après le lemme de Nakayama \ref{algebra-lemma-NAK}, nous voyons que
$$
M_{\mathfrak p} \not = 0 \Leftrightarrow
M \otimes \kappa(\mathfrak p) \not = 0 \Leftrightarrow
\text{rang}(A \bmod \mathfrak p) < n.
$$
Ainsi, si $I$ est l'idéal de $R$ engendré par les mineurs $n \times n$
de $A$, alors $\text{Supp}(M) = V(I)$. Comme $I$
est de type fini, disons $I = (f_1, \ldots, f_t)$,
nous voyons que $\Spec(R) \setminus V(I)$ est
une réunion finie des ouverts standards $D(f_i)$, donc est quasi-compact.
\end{proof}

\begin{lemma}
\label{algebra-lemma-support-quotient}
Soit $R$ un anneau et soit $M$ un $R$-module.
\begin{enumerate}
\item Si $M$ est de type fini, alors le support
de $M/IM$ est $\text{Supp}(M) \cap V(I)$.
\item Si $N \subset M$, alors $\text{Supp}(N) \subset
\text{Supp}(M)$.
\item Si $Q$ est un module quotient de $M$, alors $\text{Supp}(Q) \subset
\text{Supp}(M)$.
\item Si $0 \to N \to M \to Q \to 0$ est une suite exacte courte,
alors $\text{Supp}(M) = \text{Supp}(Q) \cup \text{Supp}(N)$.
\end{enumerate}
\end{lemma}

\begin{proof}
Les foncteurs $M \mapsto M_{\mathfrak p}$ sont exacts. Cela implique immédiatement
toutes les assertions sauf la première. Pour la première assertion,
nous devons montrer que $M_\mathfrak p \not = 0$ et
$I \subset \mathfrak p$ impliquent $(M/IM)_{\mathfrak p}
= M_\mathfrak p/IM_\mathfrak p \not = 0$. Cela résulte
du lemme de Nakayama \ref{algebra-lemma-NAK}.
\end{proof}





\section{Montée et descente}
\label{algebra-section-going-up}

\noindent
Supposons que $\mathfrak p$, $\mathfrak p'$ soient des idéaux premiers
de l'anneau $R$. Soit $X = \Spec(R)$ muni de la topologie de
Zariski. Notons $x \in X$ le point correspondant
à $\mathfrak p$ et $x' \in X$ le point correspondant
à $\mathfrak p'$. Nous avons alors :
$$
x' \leadsto x \Leftrightarrow \mathfrak p' \subset \mathfrak p.
$$
Autrement dit, $x$ est une spécialisation de $x'$ si et
seulement si $\mathfrak p' \subset \mathfrak p$.
Voir Topologie, section \ref{topology-section-specialization},
pour la terminologie et les notations.

\begin{definition}
\label{algebra-definition-going-up-down}
Soit $\varphi : R \to S$ un morphisme d'anneaux.
\begin{enumerate}
\item Nous disons que $\varphi : R \to S$ vérifie la {\it propriété de montée} si,
étant donnés des idéaux premiers $\mathfrak p \subset \mathfrak p'$ de $R$
et un idéal premier $\mathfrak q$ de $S$ situé au-dessus de $\mathfrak p$,
il existe un idéal premier $\mathfrak q'$ de $S$ tel que
(a) $\mathfrak q \subset \mathfrak q'$ et (b)
$\mathfrak q'$ soit situé au-dessus de $\mathfrak p'$.
\item Nous disons que $\varphi : R \to S$ vérifie la {\it propriété de descente} si,
étant donnés des idéaux premiers $\mathfrak p \subset \mathfrak p'$ de $R$
et un idéal premier $\mathfrak q'$ de $S$ situé au-dessus de $\mathfrak p'$,
il existe un idéal premier $\mathfrak q$ de $S$ tel que
(a) $\mathfrak q \subset \mathfrak q'$ et (b)
$\mathfrak q$ soit situé au-dessus de $\mathfrak p$.
\end{enumerate}
\end{definition}

\noindent
Nous avons jusqu'ici rencontré les cas suivants :
\begin{enumerate}
\item Un morphisme d'anneaux entier vérifie la propriété de montée, voir le
lemme \ref{algebra-lemma-integral-going-up}.
\item En particulier, les morphismes d'anneaux finis vérifient la propriété de montée.
\item En particulier, les morphismes quotients $R \to R/I$ vérifient la propriété de montée.
\item Un morphisme d'anneaux plat vérifie la propriété de descente, voir le
lemme \ref{algebra-lemma-flat-going-down}.
\item En particulier, toute localisation vérifie la propriété de descente.
\item Une extension $R \subset S$ d'anneaux intègres, où $R$ est normal
et $S$ est entier sur $R$, vérifie la propriété de descente, voir la
proposition \ref{algebra-proposition-going-down-normal-integral}.
\end{enumerate}
Voici un autre cas où la propriété de descente est vérifiée.

\begin{lemma}
\label{algebra-lemma-open-going-down}
Soit $R \to S$ un morphisme d'anneaux. Si l'application induite
$\varphi : \Spec(S) \to \Spec(R)$ est ouverte, alors
$R \to S$ vérifie la propriété de descente.
\end{lemma}

\begin{proof}
Supposons que $\mathfrak p \subset \mathfrak p' \subset R$ et qu'un
$\mathfrak q' \subset S$ soit situé au-dessus de $\mathfrak p'$. Comme $\varphi$ est ouverte,
pour tout $g \in S$, $g \not \in \mathfrak q'$, nous voyons que $\mathfrak p$
appartient à l'image de $D(g) \subset \Spec(S)$. Autrement dit,
$S_g \otimes_R \kappa(\mathfrak p)$ n'est pas nul. Puisque $S_{\mathfrak q'}$
est la colimite filtrante de ces $S_g$, cela implique
que $S_{\mathfrak q'} \otimes_R \kappa(\mathfrak p)$ n'est pas
nul, voir les
lemmes \ref{algebra-lemma-localization-colimit} et
\ref{algebra-lemma-tensor-products-commute-with-limits}.
Par conséquent, $\mathfrak p$ appartient à l'image de
$\Spec(S_{\mathfrak q'}) \to \Spec(R)$, comme voulu.
\end{proof}

\begin{lemma}
\label{algebra-lemma-going-up-down-specialization}
Soit $R \to S$ un morphisme d'anneaux.
\begin{enumerate}
\item $R \to S$ vérifie la propriété de descente si et seulement si
les généralisations se relèvent le long de l'application $\Spec(S) \to \Spec(R)$,
voir Topologie, définition \ref{topology-definition-lift-specializations}.
\item $R \to S$ vérifie la propriété de montée si et seulement si
les spécialisations se relèvent le long de l'application $\Spec(S) \to \Spec(R)$,
voir Topologie, définition \ref{topology-definition-lift-specializations}.
\end{enumerate}
\end{lemma}

\begin{proof}
Omis.
\end{proof}

\begin{lemma}
\label{algebra-lemma-going-up-down-composition}
Supposons que $R \to S$ et $S \to T$ soient des morphismes d'anneaux vérifiant
la propriété de descente. Alors $R \to T$ la vérifie aussi. Il en va de même pour la propriété de montée.
\end{lemma}

\begin{proof}
D'après le lemme \ref{algebra-lemma-going-up-down-specialization},
cela résulte du
lemme de Topologie \ref{topology-lemma-lift-specialization-composition}.
\end{proof}

\begin{lemma}
\label{algebra-lemma-image-stable-specialization-closed}
Soit $R \to S$ un morphisme d'anneaux. Soit $T \subset \Spec(R)$
l'image de $\Spec(S)$. Si $T$ est stable par spécialisation,
alors $T$ est fermé.
\end{lemma}

\begin{proof}
Nous donnons deux preuves.

\medskip\noindent
Première preuve. Soit $\mathfrak p \subset R$ un idéal premier tel que
le point correspondant de $\Spec(R)$ appartienne à l'adhérence
de $T$. Cela signifie que, pour tout $f \in R$, $f \not \in \mathfrak p$,
nous avons $D(f) \cap T \not = \emptyset$. Remarquons que $D(f) \cap T$
est l'image de $\Spec(S_f)$ dans $\Spec(R)$. Nous
concluons donc que $S_f \not = 0$. Autrement dit, $1 \not = 0$ dans
l'anneau $S_f$. Puisque $S_{\mathfrak p}$ est la colimite filtrante
des anneaux $S_f$, nous concluons que $1 \not = 0$ dans
$S_{\mathfrak p}$. Autrement dit, $S_{\mathfrak p} \not = 0$ et,
en considérant l'image de $\Spec(S_{\mathfrak p})
\to \Spec(S) \to \Spec(R)$, nous voyons qu'il existe
$\mathfrak p' \in T$ avec $\mathfrak p' \subset \mathfrak p$.
Comme nous avons supposé $T$ stable par spécialisation, nous concluons que
$\mathfrak p$ est un point de $T$, comme voulu.

\medskip\noindent
Deuxième preuve. Posons $I = \Ker(R \to S)$. Nous pouvons remplacer $R$ par $R/I$.
Dans ce cas, le morphisme d'anneaux $R \to S$ est injectif.
D'après le lemme \ref{algebra-lemma-injective-minimal-primes-in-image},
tous les idéaux premiers minimaux de $R$ appartiennent à l'image $T$. Par conséquent,
si $T$ est stable par spécialisation, alors il contient tous les idéaux premiers.
\end{proof}

\begin{lemma}
\label{algebra-lemma-going-up-closed}
Soit $R \to S$ un morphisme d'anneaux. Les assertions suivantes sont équivalentes :
\begin{enumerate}
\item La propriété de montée est vérifiée pour $R \to S$, et
\item l'application $\Spec(S) \to \Spec(R)$ est fermée.
\end{enumerate}
\end{lemma}

\begin{proof}
Les spécialisations se relèvent le long de toute
application fermée d'espaces topologiques ; voir le
lemme de Topologie \ref{topology-lemma-closed-open-map-specialization}.
La seconde condition implique donc la première.

\medskip\noindent
Supposons que la propriété de montée soit vérifiée pour $R \to S$.
Soit $V(I) \subset \Spec(S)$ une partie fermée.
Nous voulons montrer que l'image de $V(I)$ dans $\Spec(R)$ est fermée.
Le morphisme d'anneaux $S \to S/I$ vérifie évidemment la propriété de montée.
Ainsi $R \to S \to S/I$ vérifie la propriété de montée,
d'après le lemme \ref{algebra-lemma-going-up-down-composition}.
En remplaçant $S$ par $S/I$, il suffit de montrer que l'image $T$
de $\Spec(S)$ dans $\Spec(R)$ est fermée.
D'après les lemmes de Topologie \ref{topology-lemma-open-closed-specialization}
et \ref{topology-lemma-lift-specializations-images}, cette
image est stable par spécialisation. Le résultat découle donc
du lemme \ref{algebra-lemma-image-stable-specialization-closed}.
\end{proof}

\begin{lemma}
\label{algebra-lemma-constructible-stable-specialization-closed}
Soit $R$ un anneau. Soit $E \subset \Spec(R)$ une partie constructible.
\begin{enumerate}
\item Si $E$ est stable par spécialisation, alors $E$ est fermée.
\item Si $E$ est stable par généralisation, alors $E$ est ouverte.
\end{enumerate}
\end{lemma}

\begin{proof}
Première preuve. La première assertion
résulte du lemme \ref{algebra-lemma-image-stable-specialization-closed}
combiné au lemme \ref{algebra-lemma-constructible-is-image}.
La seconde résulte du fait que le complémentaire d'une partie constructible
est constructible
(voir Topologie, lemme \ref{topology-lemma-constructible}),
de la première partie du lemme et du
lemme de Topologie \ref{topology-lemma-open-closed-specialization}.

\medskip\noindent
Deuxième preuve. Puisque $\Spec(R)$ est un espace spectral d'après le
lemme \ref{algebra-lemma-spec-spectral}, il s'agit d'un cas particulier du
lemme de Topologie
\ref{topology-lemma-constructible-stable-specialization-closed}.
\end{proof}

\begin{proposition}
\label{algebra-proposition-fppf-open}
Soit $R \to S$ plat et de présentation finie.
Alors $\Spec(S) \to \Spec(R)$ est ouverte.
Plus généralement, cela vaut pour tout morphisme d'anneaux $R \to S$ de
présentation finie qui vérifie la propriété de descente.
\end{proposition}

\begin{proof}
Si $R \to S$ est plat, alors $R \to S$ vérifie la propriété de descente d'après le
lemme \ref{algebra-lemma-flat-going-down}. Pour démontrer la proposition, nous pouvons donc
supposer que $R \to S$ est de présentation finie et vérifie la propriété de descente.

\medskip\noindent
Puisque les ouverts standards $D(g) \subset \Spec(S)$, $g \in S$,
forment une base de la topologie, il suffit de prouver que
l'image de $D(g)$ est ouverte. Rappelons que $\Spec(S_g) \to \Spec(S)$
est un homéomorphisme de $\Spec(S_g)$ sur $D(g)$
(lemme \ref{algebra-lemma-standard-open}).
Puisque $S \to S_g$ vérifie la propriété de descente (voir ci-dessus), nous voyons que
$R \to S_g$ vérifie la propriété de descente d'après le
lemme \ref{algebra-lemma-going-up-down-composition}. Ainsi, après avoir remplacé
$S$ par $S_g$, nous voyons qu'il suffit de prouver que l'image est ouverte.
D'après le théorème de Chevalley (théorème \ref{algebra-theorem-chevalley}),
l'image est une partie constructible $E$. Et $E$ est stable
par généralisation parce que $R \to S$ vérifie la propriété de descente ;
voir les lemmes de Topologie \ref{topology-lemma-open-closed-specialization}
et \ref{topology-lemma-lift-specializations-images}.
Par conséquent, $E$ est ouverte d'après le
lemme \ref{algebra-lemma-constructible-stable-specialization-closed}.
\end{proof}

\begin{lemma}
\label{algebra-lemma-same-image}
Soit $k$ un corps, et soient $R$, $S$ des $k$-algèbres.
Soit $S' \subset S$ une sous-$k$-algèbre, et soit $f \in S' \otimes_k R$.
Dans le diagramme commutatif
$$
\xymatrix{
\Spec((S \otimes_k R)_f) \ar[rd] \ar[rr] & &
\Spec((S' \otimes_k R)_f) \ar[ld] \\
& \Spec(R) &
}
$$
les images des flèches diagonales sont les mêmes.
\end{lemma}

\begin{proof}
Soit $\mathfrak p \subset R$ dans l'image de la flèche
sud-ouest. Cela signifie (lemme \ref{algebra-lemma-in-image}) que
$$
(S' \otimes_k R)_f \otimes_R \kappa(\mathfrak p)
=
(S' \otimes_k \kappa(\mathfrak p))_f
$$
n'est pas l'anneau nul, c'est-à-dire que $S' \otimes_k \kappa(\mathfrak p)$
n'est pas l'anneau nul et que l'image de $f$ dans cet anneau n'est pas nilpotente.
Le morphisme d'anneaux
$S' \otimes_k \kappa(\mathfrak p) \to S \otimes_k \kappa(\mathfrak p)$
est injectif. Ainsi, $S \otimes_k \kappa(\mathfrak p)$
n'est pas non plus l'anneau nul et l'image de $f$ n'y est pas nilpotente.
Par conséquent, $(S \otimes_k R)_f \otimes_R \kappa(\mathfrak p)$
n'est pas l'anneau nul. Ainsi (lemme \ref{algebra-lemma-in-image}),
nous voyons que $\mathfrak p$ appartient à l'image de la flèche sud-est,
comme voulu.
\end{proof}

\begin{lemma}
\label{algebra-lemma-map-into-tensor-algebra-open}
Soit $k$ un corps.
Soient $R$ et $S$ des $k$-algèbres.
L'application $\Spec(S \otimes_k R) \to \Spec(R)$
est ouverte.
\end{lemma}

\begin{proof}
Soit $f \in S \otimes_k R$.
Il suffit de prouver que l'image de l'ouvert standard $D(f)$ est ouverte.
Soit $S' \subset S$ une sous-$k$-algèbre de type fini telle que
$f \in S' \otimes_k R$. L'application $R \to S' \otimes_k R$ est plate
et de présentation finie ; l'image $U$ de
$\Spec((S' \otimes_k R)_f) \to \Spec(R)$ est donc ouverte
d'après la proposition \ref{algebra-proposition-fppf-open}.
D'après le lemme \ref{algebra-lemma-same-image}, c'est aussi l'image de $D(f)$, et nous concluons.
\end{proof}

\noindent
Voici un lemme délicat qui est parfois utile.

\begin{lemma}
\label{algebra-lemma-unique-prime-over-localize-below}
Soit $R \to S$ un morphisme d'anneaux.
Soit $\mathfrak p \subset R$ un idéal premier.
Supposons que
\begin{enumerate}
\item il existe un unique idéal premier $\mathfrak q \subset S$ situé au-dessus de
$\mathfrak p$, et
\item soit
\begin{enumerate}
\item la propriété de montée est vérifiée pour $R \to S$, soit
\item la propriété de descente est vérifiée pour $R \to S$ et il y a au plus un idéal premier
de $S$ au-dessus de chaque idéal premier de $R$.
\end{enumerate}
\end{enumerate}
Alors $S_{\mathfrak p} = S_{\mathfrak q}$.
\end{lemma}

\begin{proof}
Considérons un idéal premier quelconque $\mathfrak q' \subset S$ qui correspond à
un point de $\Spec(S_{\mathfrak p})$. Cela signifie que
$\mathfrak p' = R \cap \mathfrak q'$ est contenu dans $\mathfrak p$.
Voici un diagramme
$$
\xymatrix{
\mathfrak q' \ar@{-}[d] \ar@{-}[r] & ? \ar@{-}[r] \ar@{-}[d] & S \ar@{-}[d] \\
\mathfrak p' \ar@{-}[r] & \mathfrak p \ar@{-}[r] & R
}
$$
Supposons (1) et (2)(a).
Par montée, il existe un idéal premier $\mathfrak q'' \subset S$
tel que $\mathfrak q' \subset \mathfrak q''$ et que $\mathfrak q''$
soit situé au-dessus de $\mathfrak p$. Par l'unicité de $\mathfrak q$, nous
concluons que $\mathfrak q'' = \mathfrak q$. Autrement dit,
$\mathfrak q'$ définit un point de $\Spec(S_{\mathfrak q})$.

\medskip\noindent
Supposons (1) et (2)(b).
Par descente, il existe un idéal premier $\mathfrak q'' \subset \mathfrak q$
situé au-dessus de $\mathfrak p'$. Par l'unicité des idéaux premiers situés au-dessus de
$\mathfrak p'$, nous voyons que $\mathfrak q' = \mathfrak q''$. Autrement dit,
$\mathfrak q'$ définit un point de $\Spec(S_{\mathfrak q})$.

\medskip\noindent
Dans les deux cas, nous concluons que l'application
$\Spec(S_{\mathfrak q}) \to \Spec(S_{\mathfrak p})$
est bijective. Cela signifie clairement que tous les éléments de $S - \mathfrak q$
sont inversibles dans $S_{\mathfrak p}$, autrement dit
$S_{\mathfrak p} = S_{\mathfrak q}$.
\end{proof}

\noindent
Le lemme suivant est une généralisation de la propriété de descente pour les
morphismes d'anneaux plats.

\begin{lemma}
\label{algebra-lemma-going-down-flat-module}
Soit $R \to S$ un morphisme d'anneaux. Soit $N$ un $S$-module de type fini plat sur $R$.
Munissons $\text{Supp}(N) \subset \Spec(S)$ de la topologie induite.
Alors les généralisations se relèvent le long de $\text{Supp}(N) \to \Spec(R)$.
\end{lemma}

\begin{proof}
L'énoncé signifie ce qui suit. Soient
$\mathfrak p \subset \mathfrak p' \subset R$ des idéaux premiers. Soit
$\mathfrak q' \subset S$ un idéal premier tel que $\mathfrak q' \in \text{Supp}(N)$.
Il existe alors un idéal premier $\mathfrak q \subset \mathfrak q'$,
$\mathfrak q \in \text{Supp}(N)$, situé au-dessus de $\mathfrak p$.
Comme $N$ est plat sur $R$, nous voyons que $N_{\mathfrak q'}$ est plat
sur $R_{\mathfrak p'}$, voir le lemme \ref{algebra-lemma-flat-localization}.
Comme $N_{\mathfrak q'}$ est de type fini sur $S_{\mathfrak q'}$
et non nul puisque $\mathfrak q' \in \text{Supp}(N)$, nous voyons
que $N_{\mathfrak q'} \otimes_{S_{\mathfrak q'}} \kappa(\mathfrak q')$
est non nul d'après le lemme de Nakayama \ref{algebra-lemma-NAK}.
Ainsi $N_{\mathfrak q'} \otimes_{R_{\mathfrak p'}} \kappa(\mathfrak p')$
n'est pas nul non plus. Nous déduisons du lemme \ref{algebra-lemma-ff}
que $N_{\mathfrak q'} \otimes_{R_{\mathfrak p'}} \kappa(\mathfrak p)$
est non nul. Soit
$J \subset S_{\mathfrak q'} \otimes_{R_{\mathfrak p'}} \kappa(\mathfrak p)$
l'annulateur du module non nul de type fini
$N_{\mathfrak q'} \otimes_{R_{\mathfrak p'}} \kappa(\mathfrak p)$.
Puisque $J$ est un idéal propre, nous pouvons choisir un idéal premier
$\mathfrak q \subset S$ qui correspond à un idéal premier de
$S_{\mathfrak q'} \otimes_{R_{\mathfrak p'}} \kappa(\mathfrak p)/J$.
Cet idéal premier appartient au support de $N$, est situé au-dessus de
$\mathfrak p$ et est contenu dans $\mathfrak q'$, comme voulu.
\end{proof}

\section{Extensions séparables}
\label{algebra-section-separability}

\noindent
Dans cette section, nous étudions la séparabilité pour les extensions de
corps non algébriques. Elle est étroitement liée à la notion d'algèbres
géométriquement réduites, voir la
Définition \ref{algebra-definition-geometrically-reduced}.

\begin{definition}
\label{algebra-definition-separable-field-extension}
Soit $K/k$ une extension de corps.
\begin{enumerate}
\item On dit que $K$ est {\it séparablement engendré sur $k$} s'il existe
une base de transcendance $\{x_i; i \in I\}$ de $K/k$ telle que l'extension
$K/k(x_i; i \in I)$ soit une extension algébrique séparable.
\item On dit que $K$ est {\it séparable sur $k$} si, pour toute sous-extension
$k \subset K' \subset K$ avec $K'$ de type fini
sur $k$, l'extension $K'/k$ est séparablement engendrée.
\end{enumerate}
\end{definition}

\noindent
Avec cette définition quelque peu maladroite, il n'est pas clair qu'une
extension de corps séparablement engendrée soit elle-même séparable.
Nous verrons que c'est bien le cas, voir le
Lemme \ref{algebra-lemma-separably-generated-separable}.

\begin{lemma}
\label{algebra-lemma-subextensions-are-separable}
Soit $K/k$ une extension de corps séparable.
Pour toute sous-extension $K/K'/k$, l'extension de corps
$K'/k$ est séparable.
\end{lemma}

\begin{proof}
C'est immédiat d'après la définition.
\end{proof}

\begin{lemma}
\label{algebra-lemma-generating-finitely-generated-separable-field-extensions}
Soit $K/k$ une extension de corps séparablement engendrée et de type fini.
Posons $r = \text{trdeg}_k(K)$. Alors il existe des éléments
$x_1, \ldots, x_{r + 1}$ de $K$ tels que
\begin{enumerate}
\item $x_1, \ldots, x_r$ soit une base de transcendance de $K$ sur $k$,
\item $K = k(x_1, \ldots, x_{r + 1})$, et
\item $x_{r + 1}$ soit séparable sur $k(x_1, \ldots, x_r)$.
\end{enumerate}
\end{lemma}

\begin{proof}
Il suffit de combiner la définition avec le
lemme \ref{fields-lemma-primitive-element} de Corps.
\end{proof}

\begin{lemma}
\label{algebra-lemma-make-separably-generated}
Soit $K/k$ une extension de corps de type fini.
Il existe un diagramme
$$
\xymatrix{
K \ar[r] & K' \\
k \ar[u] \ar[r] & k' \ar[u]
}
$$
où $k'/k$ et $K'/K$ sont des extensions de corps finies purement
inséparables, telles que $K'/k'$ soit une extension de corps
séparablement engendrée.
\end{lemma}

\begin{proof}
Ce lemme n'est intéressant que lorsque la caractéristique de $k$ est
$p > 0$. Choisissons $x_1, \ldots, x_r$ comme base de transcendance
de $K$ sur $k$. Puisque $K$ est de type fini sur $k$, l'extension
$k(x_1, \ldots, x_r) \subset K$ est finie.
Soit $K/K_{sep}/k(x_1, \ldots, x_r)$ la sous-extension trouvée dans le
lemme \ref{fields-lemma-separable-first} de Corps.
Si $K = K_{sep}$, c'est terminé.
Nous raisonnerons par récurrence sur $d = [K : K_{sep}]$.

\medskip\noindent
Supposons que $d > 1$. Choisissons un $\beta \in K$ tel que
$\alpha = \beta^p \in K_{sep}$ et $\beta \not \in K_{sep}$.
Soit $P = T^n + a_1T^{n - 1} + \ldots + a_n$
le polynôme minimal de $\alpha$ sur $k(x_1, \ldots, x_r)$.
Soit $k'/k$ une extension finie purement inséparable
obtenue en adjoignant des racines $p$-ièmes, de sorte que chaque $a_i$
soit une puissance $p$-ième dans $k'(x_1^{1/p}, \ldots, x_r^{1/p})$.
Une telle extension existe ; les détails sont omis.
Soit $L$ un corps s'insérant dans le diagramme
$$
\xymatrix{
K \ar[r] & L \\
k(x_1, \ldots, x_r) \ar[u] \ar[r] & k'(x_1^{1/p}, \ldots, x_r^{1/p}) \ar[u]
}
$$
Nous pouvons et allons supposer que $L$ est le compositum de $K$ et
de $k'(x_1^{1/p}, \ldots, x_r^{1/p})$.
Soit $L/L_{sep}/k'(x_1^{1/p}, \ldots, x_r^{1/p})$
la sous-extension trouvée dans le
lemme \ref{fields-lemma-separable-first} de Corps.
Alors $L_{sep}$ est le compositum de
$K_{sep}$ et de $k'(x_1^{1/p}, \ldots, x_r^{1/p})$.
L'élément $\alpha \in L_{sep}$ est une racine du polynôme
$P$, dont tous les coefficients sont des puissances $p$-ièmes dans
$k'(x_1^{1/p}, \ldots, x_r^{1/p})$ et dont les racines sont
deux à deux distinctes. D'après le
lemme \ref{fields-lemma-pth-root} de Corps,
nous voyons que $\alpha = (\alpha')^p$ pour un certain $\alpha' \in L_{sep}$.
Cela signifie clairement que $\beta$ s'envoie sur $\alpha' \in L_{sep}$.
Autrement dit, nous obtenons la tour de corps
$$
\xymatrix{
K \ar[r] & L \\
K_{sep}(\beta) \ar[r] \ar[u] & L_{sep} \ar[u] \\
K_{sep} \ar[r] \ar[u] & L_{sep} \ar@{=}[u] \\
k(x_1, \ldots, x_r) \ar[u] \ar[r] & k'(x_1^{1/p}, \ldots, x_r^{1/p}) \ar[u] \\
k \ar[r] \ar[u] & k' \ar[u]
}
$$
Ainsi cette construction conduit à une nouvelle situation avec
$[L : L_{sep}] < [K : K_{sep}]$. Par récurrence, nous pouvons trouver
$k' \subset k''$ et $L \subset L'$ comme dans le lemme pour
l'extension $L/k'$. Alors les extensions $k''/k$ et $L'/K$
conviennent pour l'extension $K/k$.
Ceci démontre le lemme.
\end{proof}

\section{Algèbres géométriquement réduites}
\label{algebra-section-geometrically-reduced}

\noindent
Le résultat principal sur les algèbres géométriquement réduites est le
Lemme \ref{algebra-lemma-geometrically-reduced-finite-purely-inseparable-extension}.
Nous suggérons au lecteur de passer directement à ce lemme après avoir lu la définition.

\begin{definition}
\label{algebra-definition-geometrically-reduced}
Soit $k$ un corps. Soit $S$ une $k$-algèbre.
On dit que $S$ est {\it géométriquement réduite sur $k$}
si, pour toute extension de corps $K/k$, la
$K$-algèbre $K \otimes_k S$ est réduite.
\end{definition}

\noindent
Soient $k$ un corps et $S$ une $k$-algèbre réduite.
Pour vérifier que $S$ est géométriquement réduite, il suffit
de vérifier que $\overline{k} \otimes_k S$ est réduite (où
$\overline{k}$ désigne la clôture algébrique de $k$).
En fait, il suffit de le vérifier pour les extensions de corps finies
purement inséparables $k'/k$. Voir le
Lemme \ref{algebra-lemma-geometrically-reduced-finite-purely-inseparable-extension}.

\begin{lemma}
\label{algebra-lemma-subalgebra-separable}
Propriétés élémentaires des algèbres géométriquement réduites.
Soit $k$ un corps. Soit $S$ une $k$-algèbre.
\begin{enumerate}
\item Si $S$ est géométriquement réduite sur $k$, toute
$k$-sous-algèbre.
\item Si toutes les $k$-sous-algèbres de type fini de $S$ sont
géométriquement réduites, alors $S$ est géométriquement réduite.
\item Une colimite filtrante de $k$-algèbres géométriquement réduites
est géométriquement réduite.
\item Si $S$ est géométriquement réduite sur $k$, toute localisation
de $S$ est géométriquement réduite sur $k$.
\end{enumerate}
\end{lemma}

\begin{proof}
Omis. Les deuxième et troisième propriétés résultent du fait que
le produit tensoriel commute aux colimites.
\end{proof}

\begin{lemma}
\label{algebra-lemma-geometrically-reduced-permanence}
Soit $k$ un corps.
Si $R$ est géométriquement réduite sur $k$
et si $S \subset R$ est une partie multiplicative, alors la localisation
$S^{-1}R$ est géométriquement réduite sur $k$.
Si $R$ est géométriquement réduite sur $k$, alors $R[x]$ est
géométriquement réduite sur $k$.
\end{lemma}

\begin{proof}
Omis. Indications : une localisation d'un anneau réduit est réduite,
et la localisation commute aux produits tensoriels.
\end{proof}

\noindent
Dans les preuves des lemmes suivants, nous utiliserons à plusieurs reprises
l'observation suivante : supposons que $R' \subset R$ et
$S' \subset S$ soient des inclusions de $k$-algèbres.
Alors l'application $R' \otimes_k S' \to R \otimes_k S$
est injective.

\begin{lemma}
\label{algebra-lemma-limit-argument}
Soit $k$ un corps. Soient $R$ et $S$ des $k$-algèbres.
\begin{enumerate}
\item Si $R \otimes_k S$ n'est pas réduit, il existe des
sous-algèbres de type fini $R' \subset R$ et
$S' \subset S$ telles que $R' \otimes_k S'$ ne soit pas réduit.
\item Si $R \otimes_k S$ contient un diviseur de zéro non nul, il existe des
sous-algèbres de type fini $R' \subset R$ et
$S' \subset S$ telles que $R' \otimes_k S'$ contienne un diviseur de zéro non nul.
\item Si $R \otimes_k S$ contient un idempotent non trivial, il existe des
sous-algèbres de type fini $R' \subset R$ et
$S' \subset S$ telles que $R' \otimes_k S'$ contienne un idempotent non trivial.
\end{enumerate}
\end{lemma}

\begin{proof}
Supposons que $z \in R \otimes_k S$ soit nilpotent. On peut écrire
$z = \sum_{i = 1, \ldots, n} x_i \otimes y_i$.
On peut donc prendre pour $R'$ la $k$-sous-algèbre engendrée par
les $x_i$, et pour $S'$ la $k$-sous-algèbre engendrée par les $y_i$.
Les deuxième et troisième assertions se démontrent de la même manière.
\end{proof}

\begin{lemma}
\label{algebra-lemma-geometrically-reduced-any-reduced-base-change}
Soit $k$ un corps.
Soit $S$ une $k$-algèbre géométriquement réduite.
Soit $R$ une $k$-algèbre réduite quelconque.
Alors $R \otimes_k S$ est réduite.
\end{lemma}

\begin{proof}
D'après le Lemme \ref{algebra-lemma-limit-argument},
nous pouvons supposer que $R$ est de type fini sur $k$.
Alors, comme anneau réduit noethérien, $R$ se plonge dans un produit
fini de corps
(voir les Lemmes \ref{algebra-lemma-total-ring-fractions-no-embedded-points},
\ref{algebra-lemma-Noetherian-irreducible-components} et
\ref{algebra-lemma-minimal-prime-reduced-ring}).
Nous pouvons donc supposer que $R$ est un produit fini de corps.
Dans ce cas, la
Définition \ref{algebra-definition-geometrically-reduced}
montre que $R \otimes_k S$ est réduite.
\end{proof}

\begin{lemma}
\label{algebra-lemma-separable-extension-preserves-reducedness}
Soit $k$ un corps.
Soit $S$ une $k$-algèbre réduite.
Soit $K/k$ une extension de corps séparable
ou une extension de corps séparablement engendrée.
Alors $K \otimes_k S$ est réduite.
\end{lemma}

\begin{proof}
Supposons que $k \subset K$ soit séparable.
D'après le Lemme \ref{algebra-lemma-limit-argument},
nous pouvons supposer que $S$ est de type fini sur $k$
et que $K$ est de type fini sur $k$.
Alors $S$ se plonge dans un produit fini de corps,
à savoir son anneau total des fractions (voir les
Lemmes \ref{algebra-lemma-minimal-prime-reduced-ring} et
\ref{algebra-lemma-total-ring-fractions-no-embedded-points}).
Nous pouvons donc en fait supposer que $S$ est un domaine.
Choisissons $x_1, \ldots, x_{r + 1} \in K$ comme dans le
Lemme \ref{algebra-lemma-generating-finitely-generated-separable-field-extensions}.
Soit $P \in k(x_1, \ldots, x_r)[T]$
le polynôme minimal de $x_{r + 1}$. C'est un polynôme séparable.
Il est facile de voir que
$k[x_1, \ldots, x_r] \otimes_k S = S[x_1, \ldots, x_r]$ est un domaine.
Il s'ensuit que $k(x_1, \ldots, x_r) \otimes_k S$ est un domaine,
car c'est une localisation de $S[x_1, \ldots, x_r]$.
L'extension d'anneaux
$k(x_1, \ldots, x_r) \otimes_k S \subset K \otimes_k S$
est engendrée par un seul élément $x_{r + 1}$ qui satisfait une seule
équation, à savoir $P$. Par conséquent, $K \otimes_k S$ se plonge dans
$F[T]/(P)$, où $F$ est le corps des fractions de
$k(x_1, \ldots, x_r) \otimes_k S$.
Comme $P$ est séparable, ceci est un produit fini de corps, et la conclusion s'ensuit.

\medskip\noindent
À ce stade, nous ne savons pas encore qu'une extension de corps
séparablement engendrée est séparable ; il faut donc aussi démontrer
le lemme dans ce cas.
Pour cela, supposons que $\{x_i\}_{i \in I}$ soit une base de
transcendance séparante de $K$ sur $k$. Pour tout ensemble fini
d'éléments $\lambda_j \in K$, il existe un sous-ensemble fini
$T \subset I$ tel que
$k(\{x_i\}_{i\in T}) \subset k(\{x_i\}_{i \in T} \cup \{\lambda_j\})$
soit finie et séparable. Il s'ensuit que $K$ est une colimite filtrante
d'extensions de $k$ de type fini et séparablement engendrées.
L'argument du paragraphe précédent s'applique donc également à ce cas.
\end{proof}

\begin{lemma}
\label{algebra-lemma-generic-points-geometrically-reduced}
Soit $k$ un corps et soit $S$ une $k$-algèbre. Supposons que
$S$ soit réduite et que $S_{\mathfrak p}$ soit géométriquement
réduite pour tout idéal premier minimal $\mathfrak p$ de $S$.
Alors $S$ est géométriquement réduite.
\end{lemma}

\begin{proof}
Comme $S$ est réduite, l'application
$S \to \prod_{\mathfrak p\text{ minimal}} S_{\mathfrak p}$
est injective, voir le
Lemme \ref{algebra-lemma-reduced-ring-sub-product-fields}.
Si $K/k$ est une extension de corps, alors les applications
$$
S \otimes_k K \to (\prod S_\mathfrak p) \otimes_k K \to
\prod S_\mathfrak p \otimes_k K
$$
sont injectives : la première parce que $k \to K$ est plate, et la seconde
par inspection puisque $K$ est un $k$-module libre. Comme $S_\mathfrak p$ est
géométriquement réduite, l'anneau de droite est réduit. Ainsi $S \otimes_k K$
est réduit comme sous-anneau d'un anneau réduit.
\end{proof}

\begin{lemma}
\label{algebra-lemma-separable-algebraic-diagonal}
Soit $k'/k$ une extension algébrique séparable.
Alors il existe une partie multiplicative $S \subset k' \otimes_k k'$
telle que l'application de multiplication $k' \otimes_k k' \to k'$
s'identifie à $k' \otimes_k k' \to S^{-1}(k' \otimes_k k')$.
\end{lemma}

\begin{proof}
Supposons d'abord que $k'/k$ soit finie et séparable. Alors
$k' = k(\alpha)$, voir le lemme \ref{fields-lemma-primitive-element} de Corps.
Soit $P \in k[x]$ le polynôme minimal de $\alpha$ sur $k$.
Alors $P$ est un polynôme irréductible, séparable et unitaire, voir
la section \ref{fields-section-separable-extensions} de Corps.
Alors $k'[x]/(P) \to k' \otimes_k k'$,
$\sum \alpha_i x^i \mapsto \alpha_i \otimes \alpha^i$ est un isomorphisme.
On peut factoriser $P = (x - \alpha) Q$ dans $k'[x]$ et, puisque $P$
est séparable, on a $Q(\alpha) \not = 0$.
Il est alors clair que la partie multiplicative $S'$ engendrée par
$Q$ dans $k'[x]/(P)$ convient, c'est-à-dire
$k' = (S')^{-1}(k'[x]/(P))$.
Par transport de structure, l'image $S$ de $S'$ dans $k' \otimes_k k'$
convient.

\medskip\noindent
Dans le cas général, écrivons $k' = \bigcup k_i$ comme réunion
de ses sous-extensions finies sur $k$. Pour chaque $i$, il existe une
partie multiplicative $S_i \subset k_i \otimes_k k_i$
telle que $k_i = S_i^{-1}(k_i \otimes_k k_i)$. Soit $S$ la clôture
multiplicative de $\bigcup S_i \subset k' \otimes_k k'$.
Il est clair que l'application de multiplication envoie chaque élément de $S$
sur un élément inversible de $k'$. Comme $k' \otimes_k k'$ est la réunion
des anneaux $k_i \otimes_k k_i$, tout élément du noyau de l'application
de multiplication est envoyé sur zéro dans $S^{-1}(k' \otimes_k k')$.
Par exactitude de la localisation, le résultat s'ensuit.
\end{proof}

\begin{lemma}
\label{algebra-lemma-geometrically-reduced-over-separable-algebraic}
Soit $k'/k$ une extension de corps algébrique séparable.
Soit $A$ une algèbre sur $k'$. Alors $A$ est géométriquement
réduite sur $k$ si et seulement si elle est géométriquement réduite sur $k'$.
\end{lemma}

\begin{proof}
Supposons que $A$ soit géométriquement réduite sur $k'$.
Soit $K/k$ une extension de corps. Alors $K \otimes_k k'$ est
un anneau réduit d'après le
Lemme \ref{algebra-lemma-separable-extension-preserves-reducedness}.
Ainsi, d'après le
Lemme \ref{algebra-lemma-geometrically-reduced-any-reduced-base-change},
nous obtenons que $K \otimes_k A = (K \otimes_k k') \otimes_{k'} A$ est réduit.

\medskip\noindent
Supposons que $A$ soit géométriquement réduite sur $k$. Soit $K/k'$ une extension
de corps. Alors
$$
K \otimes_{k'} A = (K \otimes_k A) \otimes_{(k' \otimes_k k')} k'
$$
Comme $k' \otimes_k k' \to k'$ est une localisation d'après le
Lemme \ref{algebra-lemma-separable-algebraic-diagonal}, nous voyons que
$K \otimes_{k'} A$ est une localisation d'une algèbre réduite, donc réduite.
\end{proof}

\section{Extensions séparables, suite}
\label{algebra-section-separability-continued}

\noindent
Dans cette section, nous poursuivons la discussion commencée dans la
section \ref{algebra-section-separability}.

\begin{lemma}
\label{algebra-lemma-mini-separability}
Soit $k$ un corps de caractéristique $p > 1$. Soit $K/k$
une extension de corps engendrée par $x_1, \ldots, x_{n + 1} \in K$ telle que
\begin{enumerate}
\item $\{x_1, \ldots, x_n\}$ soit une base de transcendance de $K/k$,
\item pour tout sous-ensemble $k$-linéairement indépendant
$\{a_1, \ldots, a_m\}$ de $K$, l'ensemble
$\{a^p_1, \ldots, a_m^p\}$ soit linéairement indépendant sur $k$.
\end{enumerate}
Alors il existe $1 \leq j \leq n+1$ tel que
$\{ x_1, \ldots, \widehat{x}_j, \ldots, x_{n+1}\}$
soit une base de transcendance séparante de $K / k$.
\end{lemma}

\begin{proof}
Par hypothèse, $x_{n + 1}$ est algébrique sur $k(x_1, \ldots, x_n)$,
et il existe donc un polynôme non nul $F \in k[X_1, \ldots, X_{n + 1}]$
tel que $F(x_1, \ldots, x_{n+1}) = 0$. Choisissons $F$ de degré total minimal.
Alors $F$ est irréductible, car au moins un facteur irréductible
doit également avoir cette propriété.

\medskip\noindent
Nous affirmons que, pour un certain $i$, toutes les puissances de $X_i$
qui apparaissent dans $F$ ne sont pas multiples de $p$. Supposons par
contradiction que tous les exposants apparaissant dans $F$ soient
des multiples de $p$ ; alors l'ensemble
$$
\{x_1^{\alpha_1} \ldots x^{\alpha_{n+1}}_{n+1} \mid \lambda_\alpha \neq 0\}
\subset
K
$$
est linéairement dépendant sur $k$, où les $\lambda_\alpha$ sont les
coefficients de $F$. Par l'hypothèse (2), nous concluons que l'ensemble
$$
\{x_1^{\alpha_1 / p} \ldots x^{\alpha_{n+1} / p}_{n+1}
\mid \lambda_\alpha \neq 0 \}
$$
est également linéairement dépendant sur $k$, ce qui contredit la minimalité
de $\deg(F)$.

\medskip\noindent
Choisissons $i$ tel qu'une puissance non-$p$-ième de $X_i$ apparaisse
dans $F$. Alors $x_i$ est algébrique sur
$L = k(x_1, \ldots, x_{i - 1}, x_{i + 1}, \ldots, x_{n+1})$.
D'après le lemme \ref{fields-lemma-transcendence-degree} de Corps,
nous voyons que $x_1, \ldots, x_{i - 1}, x_{i + 1}, \ldots, x_{n+1}$
est une base de transcendance de $K/k$. Ainsi $L$ est le corps des fractions
de l'anneau de polynômes sur $k$ en
$x_1, \ldots, x_{i - 1}, x_{i + 1}, \ldots, x_{n + 1}$.
Par le Lemme de Gauss, nous concluons que
$$
P(T) =
F(x_1, \ldots, x_{i - 1}, T, x_{i + 1}, \ldots, x_{n + 1}) \in L[T]
$$
est irréductible. Par construction, $P(T)$ n'est pas contenu dans $L[T^p]$.
Ainsi $K/L$ est séparable, comme requis.
\end{proof}

\noindent
Soit $p$ un nombre premier et soit $k$ un corps de caractéristique $p$.
Dans ce cas, nous écrivons $k^{1/p}$ pour l'extension de $k$ obtenue
en adjoignant les racines $p$-ièmes de tous les éléments de $k$
dans $k$.
(Autrement dit, c'est le sous-corps d'une clôture algébrique de
$k$ engendré par les racines $p$-ièmes des éléments de $k$.)

\begin{lemma}
\label{algebra-lemma-characterize-separable-field-extensions}
Soit $k$ un corps de caractéristique $p > 0$.
Soit $K/k$ une extension de corps.
Les assertions suivantes sont équivalentes :
\begin{enumerate}
\item $K$ est séparable sur $k$,
\item pour tout sous-ensemble $k$-linéairement indépendant
$\{a_1, \ldots, a_m\}$ de $K$, l'ensemble $\{a^p_1, \ldots, a_m^p\}$ est linéairement
indépendant sur $k$,
\item l'anneau $K \otimes_k k^{1/p}$ est réduit, et
\item $K$ est géométriquement réduit sur $k$.
\end{enumerate}
\end{lemma}

\begin{proof}
L'implication (1) $\Rightarrow$ (4) résulte du
Lemme \ref{algebra-lemma-separable-extension-preserves-reducedness}.
L'implication (4) $\Rightarrow$ (3) est immédiate.

\medskip\noindent
Supposons (3). Considérons l'homomorphisme d'anneaux
$m : K \otimes_k k^{1/p} \rightarrow K$ défini par
$$
\lambda \otimes \mu \rightarrow \lambda^p \mu^p
$$
Remarquons que $x^p = m(x) \otimes 1$ pour tout $x \in K \otimes_k k^{1/p}$.
Comme $K \otimes_k k^{1/p}$ est réduit, nous voyons que $m$ est injectif.
Si $\{a_1, \ldots, a_m\} \subset K$ est linéairement indépendant sur $k$,
alors $\{a_1 \otimes 1, \ldots, a_m \otimes 1\}$ est linéairement
indépendant sur $k^{1/p}$. Par injectivité de $m$, nous déduisons
qu'aucune combinaison $k$-linéaire non triviale de
$a_1^p, \ldots, a_m^p$ n'est nulle. Ainsi (3) implique (2).

\medskip\noindent
Supposons (2). Pour démontrer (1), nous pouvons supposer que $K$ est
de type fini sur $k$, et nous devons montrer que $K$ est séparablement
engendrée sur $k$.
Soit $\{x_1, \ldots, x_d\}$ une base de transcendance de $K/k$.
D'après le lemme \ref{fields-lemma-algebraic-finitely-generated} de Corps,
nous avons $[K : K'] < \infty$, où $K' = k(x_1, \ldots, x_d)$.
Choisissons la base de transcendance de sorte que le degré d'inséparabilité
$[K : K']_i$ soit minimal. Si $K / K'$ est séparable, la conclusion est acquise.
Supposons le contraire afin d'obtenir une contradiction. Alors il existe
$x_{d + 1} \in K$ qui n'est pas séparable sur $K'$ et, en particulier,
$[K'(x_{d+1}) : K']_i > 1$. Alors, d'après le
Lemme \ref{algebra-lemma-mini-separability}, il existe $1 \leq j \leq n + 1$
tel que $K'' = k(x_1, \ldots, \widehat{x}_j, \ldots, x_{d+1})$
satisfasse $[K'(x_{d+1}) : K'']_i = 1$. Par multiplicativité,
$[K : K'']_i < [K : K']_i$, ce qui donne la contradiction.
\end{proof}

\begin{lemma}
\label{algebra-lemma-separably-generated-separable}
Une extension de corps séparablement engendrée est séparable.
\end{lemma}

\begin{proof}
Combiner le Lemme \ref{algebra-lemma-separable-extension-preserves-reducedness}
avec le Lemme \ref{algebra-lemma-characterize-separable-field-extensions}.
\end{proof}

\noindent
Dans le lemme suivant, nous utiliserons la notion de clôture parfaite,
qui est définie dans la
Définition \ref{algebra-definition-perfection}.

\begin{lemma}
\label{algebra-lemma-geometrically-reduced-finite-purely-inseparable-extension}
Soit $k$ un corps. Soit $S$ une $k$-algèbre.
Les assertions suivantes sont équivalentes :
\begin{enumerate}
\item $k' \otimes_k S$ est réduit pour toute extension finie
purement inséparable $k'$ de $k$,
\item $k^{1/p} \otimes_k S$ est réduit,
\item $k^{perf} \otimes_k S$ est réduit, où $k^{perf}$ est la
clôture parfaite de $k$,
\item $\overline{k} \otimes_k S$ est réduit, où $\overline{k}$ est la
clôture algébrique de $k$, et
\item $S$ est géométriquement réduite sur $k$.
\end{enumerate}
\end{lemma}

\begin{proof}
Remarquons que toute extension finie purement inséparable $k'/k$ se plonge
dans $k^{perf}$. De plus, $k^{1/p}$ se plonge dans $k^{perf}$, qui se
plonge lui-même dans $\overline{k}$. Il est donc clair que
(5) $\Rightarrow$ (4) $\Rightarrow$ (3) $\Rightarrow$ (2)
et que (3) $\Rightarrow$ (1).

\medskip\noindent
Prouvons que (1) $\Rightarrow$ (5).
Supposons que $k' \otimes_k S$ soit réduit pour toute extension finie
purement inséparable $k'$ de $k$. Soit $K/k$ une extension de corps.
Nous devons montrer que $K \otimes_k S$ est réduit. D'après le
Lemme \ref{algebra-lemma-limit-argument}, nous nous ramenons au cas où
$K/k$ est une extension de corps de type fini. Choisissons un diagramme
$$
\xymatrix{
K \ar[r] & K' \\
k \ar[u] \ar[r] & k' \ar[u]
}
$$
comme dans le Lemme \ref{algebra-lemma-make-separably-generated}.
Par hypothèse, $k' \otimes_k S$ est réduit.
D'après le Lemme \ref{algebra-lemma-separable-extension-preserves-reducedness},
il s'ensuit que $K' \otimes_k S$ est réduit.
Nous concluons donc que $K \otimes_k S$ est réduit, comme voulu.

\medskip\noindent
Enfin, démontrons que (2) $\Rightarrow$ (5).
Supposons que $k^{1/p} \otimes_k S$ soit réduit. Alors $S$ est réduite.
De plus, pour chaque localisation $S_{\mathfrak p}$ en un idéal premier
minimal $\mathfrak p$, l'anneau $k^{1/p}\otimes_k S_{\mathfrak p}$
est une localisation de $k^{1/p} \otimes_k S$, donc est réduit.
Mais $S_{\mathfrak p}$ est un corps d'après le
Lemme \ref{algebra-lemma-minimal-prime-reduced-ring},
donc $S_{\mathfrak p}$ est géométriquement réduit d'après le
Lemme \ref{algebra-lemma-characterize-separable-field-extensions}.
Il s'ensuit du Lemme \ref{algebra-lemma-generic-points-geometrically-reduced}
que $S$ est géométriquement réduite.
\end{proof}

\section{Corps parfaits}
\label{algebra-section-perfect-fields}

\noindent
Voici la définition.

\begin{definition}
\label{algebra-definition-perfect}
Soit $k$ un corps. On dit que $k$ est {\it parfait}
si toute extension de corps de $k$ est séparable sur $k$.
\end{definition}

\begin{lemma}
\label{algebra-lemma-perfect}
Un corps $k$ est parfait si et seulement s'il est de caractéristique $0$
ou de caractéristique $p > 0$ et si tout élément possède une
racine $p$-ième.
\end{lemma}

\begin{proof}
Le cas de caractéristique nulle est clair.
Supposons que la caractéristique de $k$ soit $p > 0$.
Si $k$ est parfait, toutes les extensions de corps obtenues en adjoignant
une racine $p$-ième d'un élément de $k$ doivent être triviales ; tout
élément de $k$ possède donc une racine $p$-ième. Réciproquement, si tout
élément possède une racine $p$-ième, alors $k = k^{1/p}$ et toute extension
de corps de $k$ est séparable d'après le
Lemme \ref{algebra-lemma-characterize-separable-field-extensions}.
\end{proof}

\begin{lemma}
\label{algebra-lemma-make-separable}
Soit $K/k$ une extension de corps de type fini.
Il existe un diagramme
$$
\xymatrix{
K \ar[r] & K' \\
k \ar[u] \ar[r] & k' \ar[u]
}
$$
où $k'/k$ et $K'/K$ sont des extensions de corps finies
purement inséparables telles que $K'/k'$ soit une extension de corps séparable.
Dans cette situation, on peut supposer que $K' = k'K$ est le compositum,
et aussi que $K' = (k' \otimes_k K)_{red}$.
\end{lemma}

\begin{proof}
D'après le Lemme \ref{algebra-lemma-make-separably-generated},
nous pouvons trouver un tel diagramme avec $K'/k'$ séparablement engendrée.
D'après le
Lemme \ref{algebra-lemma-separably-generated-separable},
il s'ensuit que $K'$ est séparable sur $k'$.
Le compositum $k'K$ est une sous-extension de $K'/k'$ et donc
$k' \subset k'K$ est séparable d'après le
Lemme \ref{algebra-lemma-subextensions-are-separable}.
L'anneau $(k' \otimes_k K)_{red}$ est un domaine, car pour un certain
$n \gg 0$ l'application $x \mapsto x^{p^n}$ l'envoie dans $K$.
Il est donc un corps d'après le
Lemme \ref{algebra-lemma-integral-over-field}.
Ainsi $(k' \otimes_k K)_{red} \to K'$ l'envoie isomorphiquement sur $k'K$.
\end{proof}

\begin{lemma}
\label{algebra-lemma-perfection}
\begin{slogan}
Tout corps possède une clôture parfaite unique.
\end{slogan}
Pour tout corps $k$, il existe une extension purement inséparable
$k'/k$ telle que $k'$ soit parfaite. L'extension de corps
$k'/k$ est unique à isomorphisme unique près.
\end{lemma}

\begin{proof}
Si la caractéristique de $k$ est nulle, alors $k' = k$ est l'unique choix.
Supposons que la caractéristique de $k$ soit $p > 0$.
Pour tout $n > 0$, il existe une unique extension algébrique
$k \subset k^{1/p^n}$ telle que (a) tout élément $\lambda \in k$
possède une racine $p^n$-ième dans $k^{1/p^n}$ et (b) pour tout élément
$\mu \in k^{1/p^n}$, on ait $\mu^{p^n} \in k$.
C'est-à-dire, considérons l'application d'anneaux
$k \to k^{1/p^n} = k$, $x \mapsto x^{p^n}$.
Elle est injective et satisfait (a) et (b). Il est clair que
$k^{1/p^n} \subset k^{1/p^{n + 1}}$ comme extensions de $k$ via
l'application $y \mapsto y^p$. Nous pouvons donc prendre
$k' = \bigcup k^{1/p^n}$.
Quelques détails sont omis.
\end{proof}

\begin{definition}
\label{algebra-definition-perfection}
Soit $k$ un corps. L'extension de corps $k'/k$ du
Lemme \ref{algebra-lemma-perfection} s'appelle la {\it clôture parfaite} de $k$.
Notation : $k^{perf}/k$.
\end{definition}

\noindent
Remarquons que si $k'/k$ est une extension algébrique purement inséparable,
alors $k'$ est une sous-extension de $k^{perf}$, c'est-à-dire
$k^{perf}/k'/k$. En effet, $(k')^{perf}$ est isomorphe à $k^{perf}$
par l'unicité du
Lemme \ref{algebra-lemma-perfection}.

\begin{lemma}
\label{algebra-lemma-perfect-reduced}
Soit $k$ un corps parfait.
Toute $k$-algèbre réduite est géométriquement réduite sur $k$.
Soient $R$ et $S$ des $k$-algèbres.
Supposons que $R$ et $S$ soient toutes deux réduites.
Alors la $k$-algèbre $R \otimes_k S$ est réduite.
\end{lemma}

\begin{proof}
La première assertion résulte du
Lemme \ref{algebra-lemma-geometrically-reduced-finite-purely-inseparable-extension}.
Pour la seconde, utiliser la première assertion et le
Lemme \ref{algebra-lemma-geometrically-reduced-any-reduced-base-change}.
\end{proof}

\section{Homéomorphismes universels}
\label{algebra-section-universal-homeomorphism}

\noindent
Soit $k'/k$ une extension de corps algébrique purement inséparable.
Alors, pour toute $k$-algèbre $R$, l'application d'anneaux
$R \to k' \otimes_k R$ induit un homéomorphisme des spectres.
La raison en est le
Lemme \ref{algebra-lemma-p-ring-map} un peu plus général ci-dessous.

\begin{lemma}
\label{algebra-lemma-surjective-locally-nilpotent-kernel}
Soit $\varphi : R \to S$ un homomorphisme surjectif à noyau localement nilpotent.
Alors $\varphi$ induit un homéomorphisme des spectres et des isomorphismes
sur les corps résiduels. Pour toute application d'anneaux $R \to R'$, l'application
$R' \to R' \otimes_R S$ est surjective à noyau localement nilpotent.
\end{lemma}

\begin{proof}
D'après le Lemme \ref{algebra-lemma-spec-closed}, l'application
$\Spec(S) \to \Spec(R)$ est un homéomorphisme sur le sous-ensemble fermé
$V(\Ker(\varphi))$. Bien sûr, $V(\Ker(\varphi)) = \Spec(R)$, car tout idéal
premier de $R$ contient tout élément nilpotent de $R$. Cela implique également
l'assertion sur les corps résiduels. Par l'exactitude à droite du produit
tensoriel, nous voyons que $\Ker(\varphi)R'$ est le noyau de l'application
surjective $R' \to R' \otimes_R S$. La dernière assertion résulte donc du
Lemme \ref{algebra-lemma-locally-nilpotent}.
\end{proof}

\begin{lemma}
\label{algebra-lemma-powers-field}
\begin{reference}
\cite[Lemma 3.1.6]{Alper-adequate}
\end{reference}
Soit $k'/k$ une extension de corps. Les assertions suivantes sont équivalentes :
\begin{enumerate}
\item pour tout $x \in k'$, il existe un $n > 0$ tel que $x^n \in k$, et
\item $k' = k$, ou bien $k$ et $k'$ sont de caractéristique $p > 0$ et
soit $k'/k$ est une extension purement inséparable, soit
$k$ et $k'$ sont des extensions algébriques de $\mathbf{F}_p$.
\end{enumerate}
\end{lemma}

\begin{proof}
Remarquons que chacune des possibilités énumérées en (2) satisfait (1).
Supposons donc que $k'/k$ satisfasse (1), et montrons que l'un des cas de (2)
se produit. En écartant le cas $k = k'$, nous pouvons supposer que
$k' \not = k$. Il est clair que $k'/k$ est algébrique.
Nous pouvons donc supposer que $k'/k$ est une extension finie non triviale.
Soit $k'/k'_{sep}/k$ la sous-extension séparable
trouvée dans le lemme \ref{fields-lemma-separable-first} de Corps.
Nous devons montrer que $k = k'_{sep}$ ou que $k$ est algébrique sur
$\mathbf{F}_p$. Nous pouvons donc supposer que $k'/k$
est une extension finie séparable non triviale et nous devons montrer
que $k$ est algébrique sur $\mathbf{F}_p$.

\medskip\noindent
Choisissons $x \in k'$, $x \not \in k$. Choisissons $n, m > 0$ tels que
$x^n \in k$ et $(x + 1)^m \in k$. Soit $\overline{k}$ une clôture
algébrique de $k$. Nous pouvons choisir des plongements
$\sigma, \tau : k' \to \overline{k}$ tels que $\sigma(x) \not = \tau(x)$.
Cela résulte de la discussion de la
section \ref{fields-section-separable-extensions} de Corps
(plus précisément, après avoir remplacé $k'$ par l'extension de $k$
engendrée par $x$, cela résulte du
lemme \ref{fields-lemma-count-embeddings} de Corps).
Nous voyons alors que $\sigma(x) = \zeta \tau(x)$ pour une
racine $n$-ième de l'unité $\zeta$ dans $\overline{k}$.
De même, $\sigma(x + 1) = \zeta' \tau(x + 1)$
pour une racine $m$-ième de l'unité $\zeta' \in \overline{k}$.
Comme $\sigma(x + 1) \not = \tau(x + 1)$, nous voyons que $\zeta' \not = 1$.
Alors
$$
\zeta' (\tau(x) + 1) =
\zeta' \tau(x + 1) =
\sigma(x + 1) =
\sigma(x) + 1 =
\zeta \tau(x) + 1
$$
implique que
$$
\tau(x) (\zeta' - \zeta) = 1 - \zeta'
$$
d'où $\zeta' \not = \zeta$ et
$$
\tau(x) = (1 - \zeta')/(\zeta' - \zeta)
$$
Ainsi, tout élément de $k'$ qui n'appartient pas à $k$ est algébrique
sur le sous-corps premier. Comme $k'$ est engendré sur le sous-corps premier
par les éléments de $k'$ qui n'appartiennent pas à $k$, nous concluons que
$k'$ (et donc $k$) est algébrique sur le sous-corps premier.

\medskip\noindent
Enfin, si la caractéristique de $k$ est $0$, ce qui précède conduit à une
contradiction comme suit (nous encourageons le lecteur à trouver sa propre preuve).
Pour tout nombre rationnel $y$, nous obtenons de même une racine de l'unité
$\zeta_y$ telle que $\sigma(x + y) = \zeta_y\tau(x + y)$.
Nous trouvons alors
$$
\zeta \tau(x) + y = \zeta_y(\tau(x) + y)
$$
et, grâce à la formule ci-dessus pour $\tau(x)$, nous concluons que
$\zeta_y \in \mathbf{Q}(\zeta, \zeta')$. Comme le corps de nombres
$\mathbf{Q}(\zeta, \zeta')$ ne contient qu'un nombre fini de racines de
l'unité, il existe deux nombres rationnels distincts $y, y'$ tels que
$\zeta_y = \zeta_{y'}$. Nous concluons alors que
$$
y - y' =
\sigma(x + y) - \sigma(x + y') =
\zeta_y(\tau(x + y)) - \zeta_{y'}\tau(x + y') = \zeta_y(y - y')
$$
ce qui implique $\zeta_y = 1$, contradiction.
\end{proof}

\begin{lemma}
\label{algebra-lemma-powers}
Soit $\varphi : R \to S$ une application d'anneaux. Supposons que
\begin{enumerate}
\item pour tout $x \in S$, il existe $n > 0$ tel que
$x^n$ appartienne à l'image de $\varphi$, et
\item $\Ker(\varphi)$ soit localement nilpotent,
\end{enumerate}
alors $\varphi$ induit un homéomorphisme des spectres et induit des
extensions de corps résiduels satisfaisant aux conditions équivalentes du
Lemme \ref{algebra-lemma-powers-field}.
\end{lemma}

\begin{proof}
Supposons (1) et (2). Soient $\mathfrak q, \mathfrak q'$ des idéaux premiers
de $S$ au-dessus du même idéal premier $\mathfrak p$ de $R$. Supposons que
$x \in S$ vérifie $x \in \mathfrak q$, $x \not \in \mathfrak q'$. Alors
$x^n \in \mathfrak q$ et $x^n \not \in \mathfrak q'$ pour tout $n > 0$. Si
$x^n = \varphi(y)$ avec $y \in R$ pour un certain $n > 0$, alors
$$
x^n \in \mathfrak q \Rightarrow y \in \mathfrak p \Rightarrow
x^n \in \mathfrak q'
$$
ce qui est une contradiction. Il n'existe donc pas de $x$ comme ci-dessus,
et nous concluons que $\mathfrak q = \mathfrak q'$, c'est-à-dire que
l'application sur les spectres est injective. D'après (2), le noyau
$I = \Ker(\varphi)$ est contenu dans tout idéal premier, donc
$\Spec(R) = \Spec(R/I)$ comme espaces topologiques. Comme l'application
induite $R/I \to S$ est entière par (1), le
Lemme \ref{algebra-lemma-integral-overring-surjective} montre que
$\Spec(S) \to \Spec(R/I)$ est surjective. En combinant ce qui précède, nous
voyons que $\Spec(S) \to \Spec(R)$ est bijective.
Si $x \in S$ est quelconque et si nous choisissons $y \in R$ tel que
$\varphi(y) = x^n$ pour un certain $n > 0$, alors l'ouvert
$D(x) \subset \Spec(S)$ correspond à l'ouvert
$D(y) \subset \Spec(R)$ par la bijection ci-dessus. Ainsi
l'application $\Spec(S) \to \Spec(R)$ est un homéomorphisme.

\medskip\noindent
Pour l'assertion sur les corps résiduels, soit
$\mathfrak q \subset S$ un idéal premier au-dessus d'un idéal premier
$\mathfrak p \subset R$. Soit $x \in \kappa(\mathfrak q)$. Si nous considérons
$\kappa(\mathfrak q)$ comme le corps résiduel de l'anneau local
$S_\mathfrak q$, nous voyons que $x$ est l'image d'un certain
$y/z \in S_\mathfrak q$ avec $y \in S$, $z \in S$, $z \not \in \mathfrak q$.
Choisissons $n, m > 0$ tels que $y^n, z^m$ appartiennent à l'image de $\varphi$.
Alors $x^{nm}$ est le résidu de $(y/z)^{nm} = (y^n)^m/(z^m)^n$,
qui appartient à l'image de $R_\mathfrak p \to S_\mathfrak q$.
Ainsi $x^{nm}$ appartient à l'image de
$\kappa(\mathfrak p) \to \kappa(\mathfrak q)$.
\end{proof}

\begin{lemma}
\label{algebra-lemma-2-3-ring-map}
Soit $\varphi : R \to S$ une application d'anneaux. Supposons que
\begin{enumerate}
\item[(a)] $S$ soit engendrée comme $R$-algèbre par des éléments $x$ tels que
$x^2, x^3 \in \varphi(R)$, et
\item[(b)] $\Ker(\varphi)$ soit localement nilpotent,
\end{enumerate}
Alors $\varphi$ induit des isomorphismes sur les corps résiduels et
un homéomorphisme des spectres. Pour toute application d'anneaux $R \to R'$,
l'application $R' \to R' \otimes_R S$ satisfait également (a) et (b).
\end{lemma}

\begin{proof}
Supposons (a) et (b). L'application sur les spectres est fermée, car $S$ est
entière sur $R$, voir les Lemmes \ref{algebra-lemma-going-up-closed} et
\ref{algebra-lemma-integral-going-up}. L'image est dense d'après le
Lemme \ref{algebra-lemma-image-dense-generic-points}. Ainsi
$\Spec(S) \to \Spec(R)$ est surjective. Si
$\mathfrak q \subset S$ est un idéal premier au-dessus de
$\mathfrak p \subset R$, alors l'extension de corps
$\kappa(\mathfrak q)/\kappa(\mathfrak p)$ est engendrée par des éléments
$\alpha \in \kappa(\mathfrak q)$ dont le carré et le cube appartiennent à
$\kappa(\mathfrak p)$. Il est alors clair que
$\alpha \in \kappa(\mathfrak p)$ et nous obtenons
$\kappa(\mathfrak q) = \kappa(\mathfrak p)$.
Si $\mathfrak q, \mathfrak q'$ étaient deux idéaux premiers distincts
au-dessus de $\mathfrak p$, au moins l'un des générateurs $x$ de $S$ comme
dans (a) aurait des images distinctes dans
$\kappa(\mathfrak q) = \kappa(\mathfrak p)$ et
$\kappa(\mathfrak q') = \kappa(\mathfrak p)$.
Cela contredirait le fait que $x^2$ et $x^3$ ont bien la même image.
Cela montre que $\Spec(S) \to \Spec(R)$ est injective, donc un homéomorphisme
(d'après ce qui a déjà été démontré).

\medskip\noindent
Puisque $\varphi$ induit un homéomorphisme des spectres, elle est en particulier
surjective sur les spectres, propriété préservée par tout changement de base,
voir le Lemme \ref{algebra-lemma-surjective-spec-radical-ideal}.
Par conséquent, pour toute application $R \to R'$, le noyau de l'application
d'anneaux $R' \to R' \otimes_R S$ est constitué d'éléments nilpotents,
voir le Lemme \ref{algebra-lemma-image-dense-generic-points}; autrement dit, (b) est
satisfaite pour $R' \to R' \otimes_R S$.
Il est clair que (a) est préservée par changement de base.
\end{proof}

\begin{lemma}
\label{algebra-lemma-help-with-powers}
Soit $p$ un nombre premier. Soient $n, m > 0$ deux entiers. Il existe
un entier $a$ tel que
$(x + y)^{p^a}, p^a(x + y) \in \mathbf{Z}[x^{p^n}, p^nx, y^{p^m}, p^my]$.
\end{lemma}

\begin{proof}
C'est clair pour $p^a(x + y)$ dès que $a \geq n, m$.
En fait, choisissons $a \gg n, m$. Écrivons
$$
(x + y)^{p^a}  = \sum\nolimits_{i, j \geq 0, i + j = p^a}
{p^a \choose i, j} x^iy^j
$$
Pour tous $i, j \geq 0$ tels que $i + j = p^a$, écrivons
$i = q p^n + r$ avec $r \in \{0, \ldots, p^n - 1\}$ et
$j = q' p^m + r'$ avec $r' \in \{0, \ldots, p^m - 1\}$.
La condition $(x + y)^{p^a} \in \mathbf{Z}[x^{p^n}, p^nx, y^{p^m}, p^my]$
est satisfaite si
$$
p^{nr + mr'} \text{ divise } {p^a \choose i, j}
$$
Si $r = r' = 0$, la divisibilité est satisfaite. Si $r \not = 0$, nous
écrivons
$$
{p^a \choose i, j} = \frac{p^a}{i} {p^a - 1 \choose i - 1, j}
$$
Comme $r \not = 0$, le nombre rationnel $p^a/i$ a une valuation
$p$-adique au moins égale à $a - (n - 1)$ (car $i$ n'est pas divisible par
$p^n$). Ainsi ${p^a \choose i, j}$ est divisible par $p^{a - n + 1}$
dans ce cas. De même, si $r' \not = 0$, alors ${p^a \choose i, j}$ est
divisible par $p^{a - m + 1}$. Le choix
$a = np^n + mp^m + n + m$ convient.
\end{proof}

\begin{lemma}
\label{algebra-lemma-p-ring-map-field}
Soit $k'/k$ une extension de corps. Soit $p$ un nombre premier.
Les assertions suivantes sont équivalentes
\begin{enumerate}
\item $k'$ est engendré comme extension de corps de $k$ par des éléments
$x$ tels qu'il existe un $n > 0$ avec $x^{p^n} \in k$ et
$p^nx \in k$, et
\item $k = k'$ ou bien la caractéristique de $k$ et de $k'$ est $p$
et $k'/k$ est purement inséparable.
\end{enumerate}
\end{lemma}

\begin{proof}
Soit $x \in k'$. S'il existe un $n > 0$ tel que $x^{p^n} \in k$ et
$p^nx \in k$ et si la caractéristique n'est pas $p$, alors $x \in k$.
Si la caractéristique est $p$, alors $x^{p^n} \in k$ et donc $x$ est
purement inséparable sur $k$.
\end{proof}

\begin{lemma}
\label{algebra-lemma-p-ring-map}
Soit $\varphi : R \to S$ une application d'anneaux. Soit $p$ un nombre
premier. Supposons
\begin{enumerate}
\item[(a)] que $S$ soit engendrée comme $R$-algèbre par des éléments $x$
tels qu'il existe un $n > 0$ avec $x^{p^n} \in \varphi(R)$ et
$p^nx \in \varphi(R)$, et
\item[(b)] que $\Ker(\varphi)$ soit localement nilpotent,
\end{enumerate}
Alors $\varphi$ induit un homéomorphisme des spectres et des extensions
de corps résiduels satisfaisant aux conditions équivalentes du
Lemme \ref{algebra-lemma-p-ring-map-field}. Pour toute application d'anneaux
$R \to R'$, l'application d'anneaux $R' \to R' \otimes_R S$ satisfait
également (a) et (b).
\end{lemma}

\begin{proof}
Supposons (a) et (b). Remarquons que (b) équivaut à la condition (2)
du Lemme \ref{algebra-lemma-powers}. Soit $T \subset S$ l'ensemble des
éléments $x \in S$ tels qu'il existe un entier
$n > 0$ tel que $x^{p^n} , p^n x \in \varphi(R)$.
Nous affirmons que $T = S$. Cela montrera que la condition (1) du
Lemme \ref{algebra-lemma-powers} est satisfaite et donc que $\varphi$ induit
un homéomorphisme sur les spectres.
Par (a), il suffit de montrer que $T \subset S$ est une sous-algèbre sur
$R$ de l'algèbre ambiante. Si $x \in T$ et $y \in R$, il est clair que $yx \in T$.
Supposons $x, y \in T$ et $n, m > 0$ tels que
$x^{p^n}, y^{p^m}, p^n x, p^m y \in \varphi(R)$.
Alors $(xy)^{p^{n + m}}, p^{n + m}xy \in \varphi(R)$,
d'où $xy \in T$. Nous avons $x + y \in T$ par le
Lemme \ref{algebra-lemma-help-with-powers}, ce qui prouve l'affirmation.

\medskip\noindent
Puisque $\varphi$ induit un homéomorphisme sur les spectres, elle est en
particulier surjective sur les spectres, propriété préservée par tout
changement de base, voir le Lemme \ref{algebra-lemma-surjective-spec-radical-ideal}.
Par conséquent, pour toute application $R \to R'$, le noyau de
l'application d'anneaux $R' \to R' \otimes_R S$ est constitué d'éléments
nilpotents, voir le Lemme \ref{algebra-lemma-image-dense-generic-points};
autrement dit, (b) est satisfaite pour $R' \to R' \otimes_R S$.
Il est clair que (a) est préservée par changement de base.
Enfin, la condition sur les corps résiduels découle de (a), car les
générateurs de $S$ comme $R$-algèbre donnent des générateurs des
extensions de corps résiduels.
\end{proof}

\begin{lemma}
\label{algebra-lemma-radicial}
Soit $\varphi : R \to S$ une application d'anneaux. Supposons que
\begin{enumerate}
\item $\varphi$ induise une application injective sur les spectres,
\item $\varphi$ induise des extensions de corps résiduels purement
inséparables.
\end{enumerate}
Alors, pour toute application d'anneaux $R \to R'$, les propriétés (1) et
(2) sont vraies pour $R' \to R' \otimes_R S$.
\end{lemma}

\begin{proof}
Posons $S' = R' \otimes_R S$, de sorte que nous ayons un diagramme
commutatif d'applications continues de spectres d'anneaux
$$
\xymatrix{
\Spec(S') \ar[r] \ar[d] & \Spec(S) \ar[d] \\
\Spec(R') \ar[r] & \Spec(R)
}
$$
Soit $\mathfrak p' \subset R'$ un idéal premier au-dessus de
$\mathfrak p \subset R$. S'il n'existe aucun idéal premier de $S$
au-dessus de $\mathfrak p$, alors il n'existe aucun idéal premier de
$S'$ au-dessus de $\mathfrak p'$. Sinon, d'après la
Remarque \ref{algebra-remark-fundamental-diagram}, il existe un unique idéal
premier $\mathfrak r$ de $F = S \otimes_R \kappa(\mathfrak p)$ dont le
corps résiduel est une extension purement inséparable de
$\kappa(\mathfrak p)$. Considérons les applications d'anneaux
$$
\kappa(\mathfrak p) \to F \to \kappa(\mathfrak r)
$$
Par le Lemme \ref{algebra-lemma-minimal-prime-reduced-ring}, l'idéal
$\mathfrak r \subset F$ est localement nilpotent; nous pouvons donc
appliquer le Lemme \ref{algebra-lemma-surjective-locally-nilpotent-kernel}
à l'application d'anneaux $F \to \kappa(\mathfrak r)$.
Nous pouvons appliquer le Lemme \ref{algebra-lemma-p-ring-map}
à l'application $\kappa(\mathfrak p) \to \kappa(\mathfrak r)$.
Ainsi, la composée et la seconde flèche des applications
$$
\kappa(\mathfrak p') \to
\kappa(\mathfrak p') \otimes_{\kappa(\mathfrak p)} F \to
\kappa(\mathfrak p') \otimes_{\kappa(\mathfrak p)} \kappa(\mathfrak r)
$$
induisent des bijections sur les spectres et des extensions de corps
résiduels purement inséparables. Il en est donc de même pour la première
application. Comme
$$
\kappa(\mathfrak p') \otimes_{\kappa(\mathfrak p)} F =
\kappa(\mathfrak p') \otimes_{\kappa(\mathfrak p)}
\kappa(\mathfrak p) \otimes_R S =
\kappa(\mathfrak p') \otimes_R S =
\kappa(\mathfrak p') \otimes_{R'} R' \otimes_R S
$$
nous concluons par la discussion de la
Remarque \ref{algebra-remark-fundamental-diagram}.
\end{proof}

\begin{lemma}
\label{algebra-lemma-radicial-integral}
Soit $\varphi : R \to S$ une application d'anneaux. Supposons que
\begin{enumerate}
\item $\varphi$ soit entière,
\item $\varphi$ induise une application injective sur les spectres,
\item $\varphi$ induise des extensions de corps résiduels purement
inséparables.
\end{enumerate}
Alors $\varphi$ induit un homéomorphisme de $\Spec(S)$ sur un sous-ensemble
fermé de $\Spec(R)$ et, pour toute application d'anneaux
$R \to R'$, les propriétés (1), (2), (3) sont vraies pour
$R' \to R' \otimes_R S$.
\end{lemma}

\begin{proof}
L'application sur les spectres est fermée par les
Lemmes \ref{algebra-lemma-going-up-closed} et \ref{algebra-lemma-integral-going-up}.
Les propriétés sont préservées par changement de base, par les
Lemmes \ref{algebra-lemma-radicial} et \ref{algebra-lemma-base-change-integral}.
\end{proof}

\begin{lemma}
\label{algebra-lemma-radicial-integral-bijective}
Soit $\varphi : R \to S$ une application d'anneaux. Supposons que
\begin{enumerate}
\item $\varphi$ soit entière,
\item $\varphi$ induise une application bijective sur les spectres,
\item $\varphi$ induise des extensions de corps résiduels purement
inséparables.
\end{enumerate}
Alors $\varphi$ induit un homéomorphisme sur les spectres et, pour toute
application d'anneaux $R \to R'$, les propriétés (1), (2), (3) sont vraies
pour $R' \to R' \otimes_R S$.
\end{lemma}

\begin{proof}
Cela résulte des Lemmes \ref{algebra-lemma-radicial-integral} et
\ref{algebra-lemma-surjective-spec-radical-ideal}.
\end{proof}

\begin{lemma}
\label{algebra-lemma-universally-bijective}
Soit $\varphi : R \to S$ une application d'anneaux telle que
\begin{enumerate}
\item le noyau de $\varphi$ soit localement nilpotent, et
\item $S$ soit engendrée comme $R$-algèbre par des éléments $x$
tels qu'il existe un $n > 0$ et un polynôme $P(T) \in R[T]$
dont l'image dans $S[T]$ soit $(T - x)^n$.
\end{enumerate}
Alors $\Spec(S) \to \Spec(R)$ est un homéomorphisme et $R \to S$
induit des extensions de corps résiduels purement inséparables.
De plus, les conditions (1) et (2) restent vraies après tout changement
de base.
\end{lemma}

\begin{proof}
Nous pouvons remplacer $R$ par $R/\Ker(\varphi)$, voir le
Lemme \ref{algebra-lemma-surjective-locally-nilpotent-kernel}.
L'hypothèse (2) implique que $S$ est engendrée sur $R$ par des éléments
entiers sur $R$. Ainsi $R \subset S$ est une extension entière
(Lemme \ref{algebra-lemma-integral-closure-is-ring}).
En particulier, $\Spec(S) \to \Spec(R)$ est surjective et fermée
(Lemmes \ref{algebra-lemma-integral-overring-surjective},
\ref{algebra-lemma-going-up-closed} et \ref{algebra-lemma-integral-going-up}).

\medskip\noindent
Soit $x \in S$ l'un des générateurs de (2), c'est-à-dire qu'il existe un
$n > 0$ tel que $(T - x)^n \in R[T]$.
Soit $\mathfrak p \subset R$ un idéal premier.
L'anneau $\kappa(\mathfrak p) \otimes_R S$ est non nul d'après ce qui
précède et le Lemme \ref{algebra-lemma-in-image}.
Si la caractéristique de $\kappa(\mathfrak p)$ est nulle,
nous voyons que $nx \in R$ implique que $1 \otimes x$ appartient à l'image
de $\kappa(\mathfrak p) \to \kappa(\mathfrak p) \otimes_R S$.
Ainsi $\kappa(\mathfrak p) \to \kappa(\mathfrak p) \otimes_R S$
est un isomorphisme.
Si la caractéristique de $\kappa(\mathfrak p)$ est $p > 0$,
écrivons $n = p^k m$ avec $m$ premier à $p$.
Dans $\kappa(\mathfrak p) \otimes_R S[T]$, nous avons
$$
(T - 1 \otimes x)^n = ((T - 1 \otimes x)^{p^k})^m =
(T^{p^k} - 1 \otimes x^{p^k})^m
$$
et nous voyons que $mx^{p^k} \in R$. Cela implique que
$1 \otimes x^{p^k}$ appartient à l'image de
$\kappa(\mathfrak p) \to \kappa(\mathfrak p) \otimes_R S$.
Le Lemme \ref{algebra-lemma-p-ring-map} s'applique donc à
$\kappa(\mathfrak p) \to \kappa(\mathfrak p) \otimes_R S$.
Dans les deux cas, nous concluons que
$\kappa(\mathfrak p) \otimes_R S$ possède un unique idéal premier dont
le corps résiduel est une extension purement inséparable de
$\kappa(\mathfrak p)$. Par la Remarque \ref{algebra-remark-fundamental-diagram},
nous concluons que $\varphi$ est bijective sur les spectres.

\medskip\noindent
L'assertion sur le changement de base est immédiate.
\end{proof}

\section{Algèbres géométriquement irréductibles}
\label{algebra-section-algebras-over-fields}

\noindent
Une algèbre $S$ sur un corps $k$ est dite géométriquement irréductible si
l'algèbre $S \otimes_k k'$ possède un unique idéal premier minimal pour
toute extension de corps $k'/k$. Dans cette section, nous développons
quelques éléments de théorie pertinents pour cette notion.

\begin{lemma}
\label{algebra-lemma-flat-fibres-irreducible}
Soit $R \to S$ une application d'anneaux. Supposons que
\begin{enumerate}
\item[(a)] $\Spec(R)$ soit irréductible,
\item[(b)] $R \to S$ soit plate,
\item[(c)] $R \to S$ soit de présentation finie,
\item[(d)] les anneaux fibres $S \otimes_R \kappa(\mathfrak p)$ aient des
spectres irréductibles pour une famille dense d'idéaux premiers
$\mathfrak p$ de $R$.
\end{enumerate}
Alors $\Spec(S)$ est irréductible.
Cela reste vrai plus généralement si (b) $+$ (c) est remplacé par
« l'application $\Spec(S) \to \Spec(R)$ est ouverte ».
\end{lemma}

\begin{proof}
Les hypothèses (b) et (c) impliquent que l'application sur les spectres
est ouverte, voir la Proposition \ref{algebra-proposition-fppf-open}.
Le lemme résulte donc du
lemme \ref{topology-lemma-irreducible-on-top} de Topologie.
\end{proof}

\begin{lemma}
\label{algebra-lemma-separably-closed-irreducible}
Soit $k$ un corps séparablement clos.
Soient $R$ et $S$ des $k$-algèbres. Si $R$ et $S$ possèdent un unique
idéal premier minimal, il en est de même pour $R \otimes_k S$.
\end{lemma}

\begin{proof}
Soit $k \subset \overline{k}$ une clôture parfaite, voir la
Définition \ref{algebra-definition-perfection}.
D'après l'hypothèse, $\overline{k}$ est algébriquement clos.
Les applications d'anneaux $R \to R \otimes_k \overline{k}$ et
$S \to S \otimes_k \overline{k}$ ainsi que
$R \otimes_k S \to (R \otimes_k S) \otimes_k \overline{k}
= (R \otimes_k \overline{k}) \otimes_{\overline{k}} (S \otimes_k \overline{k})$
satisfont aux hypothèses du Lemme \ref{algebra-lemma-p-ring-map}.
Nous pouvons donc supposer que $k$ est algébriquement clos.

\medskip\noindent
Nous pouvons remplacer $R$ et $S$ par leurs réductions.
Nous pouvons donc supposer que $R$ et $S$ sont des domaines.
Par le Lemme \ref{algebra-lemma-perfect-reduced}, $R \otimes_k S$ est réduit.
Son spectre est donc réductible si et seulement s'il contient un diviseur
de zéro non nul. Par le Lemme \ref{algebra-lemma-limit-argument}, nous nous
ramenons au cas où $R$ et $S$ sont des domaines de type fini sur le corps
$k$ algébriquement clos.

\medskip\noindent
Remarquons que l'application d'anneaux $R \to R \otimes_k S$ est de
présentation finie et plate. De plus, pour tout idéal maximal
$\mathfrak m$ de $R$, nous avons
$(R \otimes_k S) \otimes_R R/\mathfrak m \cong S$ car
$k \cong R/\mathfrak m$ par le théorème des zéros de Hilbert,
Théorème \ref{algebra-theorem-nullstellensatz}. En outre, l'ensemble des idéaux
maximaux est dense dans le spectre de $R$ puisque $\Spec(R)$ est un
espace de Jacobson, voir le Lemme \ref{algebra-lemma-finite-type-field-Jacobson}.
Le Lemme \ref{algebra-lemma-flat-fibres-irreducible} s'applique donc à
l'application $R \to R \otimes_k S$, et nous concluons que le spectre de
$R \otimes_k S$ est irréductible, comme voulu.
\end{proof}

\begin{lemma}
\label{algebra-lemma-geometrically-irreducible}
Soit $k$ un corps.
Soit $R$ une $k$-algèbre.
Les assertions suivantes sont équivalentes
\begin{enumerate}
\item pour toute extension de corps $k'/k$, le spectre de
$R \otimes_k k'$ est irréductible,
\item pour toute extension de corps finie et séparable $k'/k$, le spectre de
$R \otimes_k k'$ est irréductible,
\item le spectre de $R \otimes_k \overline{k}$ est irréductible,
où $\overline{k}$ est la clôture algébrique séparable de $k$, et
\item le spectre de $R \otimes_k \overline{k}$ est irréductible,
où $\overline{k}$ est la clôture algébrique de $k$.
\end{enumerate}
\end{lemma}

\begin{proof}
Il est clair que (1) implique (2).

\medskip\noindent
Supposons (2) et soit $\overline{k}$ la clôture algébrique séparable de $k$.
Supposons que $\mathfrak q_i \subset R \otimes_k \overline{k}$, $i = 1, 2$,
soient deux idéaux premiers minimaux. Pour toute sous-extension finie
$\overline{k}/k'/k$, l'extension $k'/k$ est séparable et l'application
$R \otimes_k k' \to R \otimes_k \overline{k}$ est plate.
Ainsi $\mathfrak p_i = (R \otimes_k k') \cap \mathfrak q_i$ sont des
idéaux premiers minimaux (car nous avons la descente des idéaux premiers
pour les applications d'anneaux plates, par le
Lemme \ref{algebra-lemma-flat-going-down}). Nous avons donc
$\mathfrak p_1 = \mathfrak p_2$ par (2). Comme
$\overline{k} = \bigcup k'$, nous concluons que
$\mathfrak q_1 = \mathfrak q_2$. Ainsi $\Spec(R \otimes_k \overline{k})$
est irréductible.

\medskip\noindent
Supposons (3) et soit $\overline{k}$ la clôture algébrique de $k$.
Soit $\overline{k}/\overline{k}'/k$ la clôture algébrique séparable
correspondante de $k$. Alors $\overline{k}/\overline{k}'$ est purement
inséparable (en caractéristique positive) ou triviale.
Ainsi $R \otimes_k \overline{k}' \to R \otimes_k \overline{k}$
induit un homéomorphisme sur les spectres, par exemple par le
Lemme \ref{algebra-lemma-p-ring-map}. Nous obtenons donc (4).

\medskip\noindent
Supposons (4). Soit $k'/k$ une extension de corps quelconque et soit
$\overline{k}$ la clôture algébrique de $k$. Choisissons un corps $F$
tel que $k'$ et $\overline{k}$ soient tous deux isomorphes à des
sous-corps de $F$. Alors
$$
R \otimes_k F = (R \otimes_k \overline{k}) \otimes_{\overline{k}} F
$$
et, par le Lemme \ref{algebra-lemma-separably-closed-irreducible},
$R \otimes_k F$ possède un unique idéal premier minimal.
Enfin, l'application d'anneaux $R \otimes_k k' \to R \otimes_k F$ est
plate et injective, et tout idéal premier minimal de $R \otimes_k k'$ est
donc l'image d'un idéal premier minimal de $R \otimes_k F$ (par le
Lemme \ref{algebra-lemma-injective-minimal-primes-in-image} et la descente).
Nous concluons qu'il n'existe qu'un seul idéal premier minimal, ce qui
achève la démonstration.
\end{proof}

\begin{definition}
\label{algebra-definition-geometrically-irreducible}
Soit $k$ un corps.
Soit $S$ une $k$-algèbre.
Nous disons que $S$ est {\it géométriquement irréductible sur $k$}
si, pour toute extension de corps $k'/k$, le spectre de
$S \otimes_k k'$ est irréductible\footnote{Un espace irréductible est non vide.}.
\end{definition}

\noindent
Par le Lemme \ref{algebra-lemma-geometrically-irreducible}, il suffit de vérifier
cette propriété pour les extensions de corps finies et séparables $k'/k$
ou pour $k'$ égal à la clôture algébrique séparable de $k$.

\begin{lemma}
\label{algebra-lemma-separably-closed-irreducible-implies-geometric}
Soit $k$ un corps.
Soit $R$ une $k$-algèbre.
Si $k$ est séparablement algébriquement clos, alors $R$ est
géométriquement irréductible sur $k$ si et seulement si le spectre de
$R$ est irréductible.
\end{lemma}

\begin{proof}
Cela résulte immédiatement de la remarque qui suit la
Définition \ref{algebra-definition-geometrically-irreducible}.
\end{proof}

\begin{lemma}
\label{algebra-lemma-subalgebra-geometrically-irreducible}
Soit $k$ un corps. Soit $S$ une $k$-algèbre.
\begin{enumerate}
\item Si $S$ est géométriquement irréductible sur $k$, toute
$k$-sous-algèbre l'est également.
\item Si toutes les $k$-sous-algèbres de type fini de $S$ sont
géométriquement irréductibles, alors $S$ est géométriquement irréductible.
\item Une colimite filtrante de $k$-algèbres géométriquement
irréductibles est géométriquement irréductible.
\end{enumerate}
\end{lemma}

\begin{proof}
Soit $S' \subset S$ une sous-algèbre. Alors, pour toute extension $k'/k$,
l'application d'anneaux $S' \otimes_k k' \to S \otimes_k k'$ est également
injective. Ainsi (1) résulte du
Lemme \ref{algebra-lemma-injective-minimal-primes-in-image}
(et du fait que l'image d'un espace irréductible par une application
continue est irréductible). Les deuxième et troisième assertions résultent
du fait que le produit tensoriel commute aux colimites.
\end{proof}

\begin{lemma}
\label{algebra-lemma-geometrically-irreducible-any-base-change}
Soit $k$ un corps.
Soit $S$ une $k$-algèbre géométriquement irréductible.
Soit $R$ une $k$-algèbre quelconque.
L'application
$$
\Spec(R \otimes_k S) \longrightarrow \Spec(R)
$$
induit une bijection sur les composantes irréductibles.
\end{lemma}

\begin{proof}
Rappelons que les composantes irréductibles correspondent aux idéaux
premiers minimaux (Lemme \ref{algebra-lemma-irreducible}).
Comme $R \to R \otimes_k S$ est plate, la descente montre
(Lemme \ref{algebra-lemma-flat-going-down}) que tout idéal premier minimal de
$R \otimes_k S$ est au-dessus d'un idéal premier minimal de $R$.
Réciproquement, si $\mathfrak p \subset R$ est un idéal premier
(minimal), alors
$$
R \otimes_k S/\mathfrak p(R \otimes_k S)
=
(R/\mathfrak p) \otimes_k S
\subset
\kappa(\mathfrak p) \otimes_k S
$$
par platitude de $R \to R \otimes_k S$. L'anneau
$\kappa(\mathfrak p) \otimes_k S$ a un spectre irréductible
par hypothèse. Il s'ensuit que
$R \otimes_k S/\mathfrak p(R \otimes_k S)$ possède un unique idéal
premier minimal (Lemme \ref{algebra-lemma-injective-minimal-primes-in-image}).
Autrement dit, l'image inverse de l'ensemble irréductible $V(\mathfrak p)$
est irréductible. Le lemme en résulte.
\end{proof}

\noindent
Faisons maintenant quelques remarques sur la notion d'extensions de corps
géométriquement irréductibles.

\begin{lemma}
\label{algebra-lemma-field-extension-geometrically-irreducible}
Soit $K/k$ une extension de corps. Si $k$ est algébriquement clos dans $K$,
alors $K$ est géométriquement irréductible sur $k$.
\end{lemma}

\begin{proof}
Supposons que $k$ soit algébriquement clos dans $K$. Par la
Définition \ref{algebra-definition-geometrically-irreducible} et le
Lemme \ref{algebra-lemma-geometrically-irreducible}, il suffit de montrer que le
spectre de $K \otimes_k k'$ est irréductible pour toute extension séparable
finie $k'/k$. Disons que $k'$ est engendré par $\alpha \in k'$ sur $k$, voir
le lemme \ref{fields-lemma-primitive-element} de Corps. Soit
$P = T^d + a_1 T^{d - 1} + \ldots + a_d \in k[T]$ le polynôme minimal
de $\alpha$. Alors $K \otimes_k k' \cong K[T]/(P)$.
La seule façon pour que le spectre de $K[T]/(P)$ soit réductible est que
$P$ soit réductible dans $K[T]$. Supposons donc que $P = P_1 P_2$ soit une
factorisation non triviale dans $K[T]$, afin d'obtenir une contradiction.
Par le Lemme \ref{algebra-lemma-polynomials-divide}, les coefficients de $P_1$ et
$P_2$ sont algébriques sur $k$. Notre hypothèse implique que les
coefficients de $P_1$ et $P_2$ appartiennent à $k$, ce qui contredit le fait
que $P$ soit irréductible sur $k$.
\end{proof}

\begin{lemma}
\label{algebra-lemma-geometrically-irreducible-transitive}
Soit $K/k$ une extension de corps géométriquement irréductible.
Soit $S$ une $K$-algèbre géométriquement irréductible.
Alors $S$ est géométriquement irréductible sur $k$.
\end{lemma}

\begin{proof}
Par la Définition \ref{algebra-definition-geometrically-irreducible} et le
Lemme \ref{algebra-lemma-geometrically-irreducible}, il suffit de montrer que le
spectre de $S \otimes_k k'$ est irréductible pour toute extension séparable
finie $k'/k$. Comme $K$ est géométriquement irréductible sur $k$, nous
voyons que $K' = K \otimes_k k'$ est une extension de corps finie et
séparable de $K$. Par conséquent, le spectre de
$S \otimes_k k' = S \otimes_K K'$ est irréductible, puisque $S$ est,
par hypothèse, géométriquement irréductible sur $K$.
\end{proof}

\begin{lemma}
\label{algebra-lemma-geometrically-irreducible-base-change-transcendental}
Soit $K/k$ une extension de corps. Les assertions suivantes sont équivalentes
\begin{enumerate}
\item $K$ est géométriquement irréductible sur $k$, et
\item l'extension induite $K(t)/k(t)$ d'extensions purement transcendantes
est géométriquement irréductible.
\end{enumerate}
\end{lemma}

\begin{proof}
Supposons (1). Notons $\Omega$ une clôture algébrique de $k(t)$.
Par la Définition \ref{algebra-definition-geometrically-irreducible}, nous obtenons
que le spectre de
$$
K \otimes_k \Omega = K \otimes_k k(t) \otimes_{k(t)} \Omega
$$
est irréductible. Comme $K(t)$ est une localisation de
$K \otimes_k k(T)$, nous concluons que le spectre de
$K(t) \otimes_{k(t)} \Omega$ est irréductible. Ainsi, par le
Lemme \ref{algebra-lemma-geometrically-irreducible}, $K(t)/k(t)$ est
géométriquement irréductible.

\medskip\noindent
Supposons (2). Soit $k'/k$ une extension de corps.
Nous devons montrer que $K \otimes_k k'$ possède un unique idéal premier
minimal. Nous savons que le spectre de
$$
K(t) \otimes_{k(t)} k'(t)
$$
est irréductible, c'est-à-dire qu'il possède un unique idéal premier minimal.
Comme il existe une application injective
$K \otimes_k k' \to K(t) \otimes_{k(t)} k'(t)$ (détails omis), nous
concluons par les Lemmes \ref{algebra-lemma-injective-minimal-primes-in-image} et
\ref{algebra-lemma-minimal-prime-image-minimal-prime}.
\end{proof}

\begin{lemma}
\label{algebra-lemma-geometrically-irreducible-add-transcendental}
Soit $K/L/M$ une tour de corps telle que $L/M$ soit géométriquement
irréductible. Soit $x \in K$ transcendant sur $L$. Alors $L(x)/M(x)$ est
géométriquement irréductible.
\end{lemma}

\begin{proof}
Cela résulte du
Lemme \ref{algebra-lemma-geometrically-irreducible-base-change-transcendental},
car les corps $L(x)$ et $M(x)$ sont des extensions purement transcendantes
de $L$ et $M$.
\end{proof}

\begin{lemma}
\label{algebra-lemma-geometrically-irreducible-separable-elements}
Soit $K/k$ une extension de corps. Les assertions suivantes sont équivalentes
\begin{enumerate}
\item $K/k$ est géométriquement irréductible, et
\item tout élément $\alpha \in K$ qui est algébrique séparable sur $k$
appartient à $k$.
\end{enumerate}
\end{lemma}

\begin{proof}
Supposons (1) et soit $\alpha \in K$ algébrique séparable sur $k$.
Alors $k' = k(\alpha)$ est une extension finie et séparable de $k$ contenue
dans $K$. Par le Lemme \ref{algebra-lemma-subalgebra-geometrically-irreducible},
l'extension $k'/k$ est géométriquement irréductible.
En particulier, le spectre de $k' \otimes_k \overline{k}$ est
irréductible (et donc, s'il s'agit d'un produit de corps, il n'y a
exactement qu'un seul facteur). Par le
lemme \ref{fields-lemma-finite-separable-tensor-alg-closed} de Corps,
$\Hom_k(k', \overline{k})$ possède un seul élément, ce qui implique à son
tour que $k' = k$ par le
lemme \ref{fields-lemma-separable-equality} de Corps. L'assertion (2)
en résulte.

\medskip\noindent
Supposons (2). Soit $k' \subset K$ le sous-corps constitué des éléments
algébriques sur $k$. Par le
Lemme \ref{algebra-lemma-field-extension-geometrically-irreducible}, l'extension
$K/k'$ est géométriquement irréductible. Par hypothèse, $k'/k$ est une
extension purement inséparable. Par le Lemme \ref{algebra-lemma-p-ring-map},
l'extension $k'/k$ est géométriquement irréductible. Ainsi, par le
Lemme \ref{algebra-lemma-geometrically-irreducible-transitive}, $K/k$ est
géométriquement irréductible.
\end{proof}

\begin{lemma}
\label{algebra-lemma-make-geometrically-irreducible}
Soit $K/k$ une extension de corps. Considérons la sous-extension
$K/k'/k$ constituée des éléments algébriques séparables sur $k$.
Alors $K$ est géométriquement irréductible sur $k'$. Si $K/k$ est une
extension de corps de type fini, alors $[k' : k] < \infty$.
\end{lemma}

\begin{proof}
La première assertion résulte immédiatement du
Lemme \ref{algebra-lemma-geometrically-irreducible-separable-elements} et du fait
que les éléments algébriques séparables sur $k'$ appartiennent à $k'$ par
transitivité des extensions algébriques séparables, voir le
lemme \ref{fields-lemma-separable-permanence} de Corps.
Si $K/k$ est de type fini, alors $k'$ est une extension finie de $k$ par le
lemme \ref{fields-lemma-algebraic-closure-in-finitely-generated} de Corps.
\end{proof}

\begin{lemma}
\label{algebra-lemma-Galois-orbit}
Soit $K/k$ une extension de corps.
Soit $\overline{k}/k$ une clôture algébrique séparable.
Alors $\text{Gal}(\overline{k}/k)$ agit transitivement sur les idéaux
premiers de $\overline{k} \otimes_k K$.
\end{lemma}

\begin{proof}
Soit $K/k'/k$ la sous-extension obtenue dans le
Lemme \ref{algebra-lemma-make-geometrically-irreducible}.
Remarquons que, comme $k \subset \overline{k}$ est entière, tous les
idéaux premiers de $\overline{k} \otimes_k K$ et
$\overline{k} \otimes_k k'$ sont maximaux, voir le
Lemme \ref{algebra-lemma-integral-no-inclusion}.
Par le Lemme \ref{algebra-lemma-geometrically-irreducible-any-base-change},
l'application
$$
\Spec(\overline{k} \otimes_k K) \to \Spec(\overline{k} \otimes_k k')
$$
est bijective car (1) tous les idéaux premiers sont minimaux, (2)
$$
\overline{k} \otimes_k K = (\overline{k} \otimes_k k') \otimes_{k'} K
$$
et (3) $K$ est géométriquement irréductible sur $k'$.
Il suffit donc de démontrer le lemme pour l'action de
$\text{Gal}(\overline{k}/k)$ sur les idéaux premiers de
$\overline{k} \otimes_k k'$.

\medskip\noindent
Comme tout idéal premier de $\overline{k} \otimes_k k'$ est maximal, les
corps résiduels sont isomorphes à $\overline{k}$. Les idéaux premiers de
$\overline{k} \otimes_k k'$ correspondent donc bijectivement aux éléments
de $\Hom_k(k', \overline{k})$, où $\sigma \in \Hom_k(k', \overline{k})$
correspond au noyau $\mathfrak p_\sigma$ de
$1 \otimes \sigma : \overline{k} \otimes_k k' \to \overline{k}$.
En particulier, $\text{Gal}(\overline{k}/k)$ agit transitivement sur cet
ensemble, comme voulu.
\end{proof}

\section{Algèbres géométriquement connexes}
\label{algebra-section-geometrically-connected}

\begin{lemma}
\label{algebra-lemma-separably-closed-connected}
Soit $k$ un corps séparablement algébriquement clos.
Soient $R$ et $S$ des $k$-algèbres. Si $\Spec(R)$ et
$\Spec(S)$ sont connexes, il en est de même pour
$\Spec(R \otimes_k S)$.
\end{lemma}

\begin{proof}
Rappelons que $\Spec(R)$ est connexe si et seulement si
$R$ ne possède pas d'idempotent non trivial, voir le
Lemme \ref{algebra-lemma-characterize-spec-connected}.
Par conséquent, par le Lemme \ref{algebra-lemma-limit-argument}, nous pouvons
supposer que $R$ et $S$ sont de type fini sur $k$.
Dans ce cas, $R$ et $S$ sont noethériens et possèdent un nombre fini
d'idéaux premiers minimaux, voir le
Lemme \ref{algebra-lemma-Noetherian-irreducible-components}.
Nous pouvons donc raisonner par récurrence sur $n + m$, où $n$, resp.\ $m$,
est le nombre de composantes irréductibles de $\Spec(R)$, resp.\ $\Spec(S)$.
Le cas où $n$ ou $m$ est nul est bien entendu trivial. Si $n = m = 1$,
c'est-à-dire si $\Spec(R)$ et $\Spec(S)$ ont chacun une seule composante
irréductible, le résultat est donné par le
Lemme \ref{algebra-lemma-separably-closed-irreducible}.
Supposons $n > 1$. Soit $\mathfrak p \subset R$ un idéal premier minimal
correspondant au fermé irréductible $T \subset \Spec(R)$.
Soit $T' \subset \Spec(R)$ l'union des $n - 1$ autres composantes
irréductibles. Choisissons un idéal $I \subset R$ tel que
$T' = V(I) = \Spec(R/I)$ (Lemme \ref{algebra-lemma-spec-closed}).
En choisissant convenablement notre idéal premier minimal, nous pouvons
de plus faire en sorte que $T'$ soit connexe, voir le
lemme \ref{topology-lemma-remove-irreducible-connected} de Topologie.
Alors $T \cup T' = \Spec(R)$ et
$T \cap T' = V(\mathfrak p + I) = \Spec(R/(\mathfrak p + I))$
n'est pas vide puisque $\Spec(R)$ est supposé connexe.
L'image inverse de $T$ dans $\Spec(R \otimes_k S)$ est
$\Spec(R/\mathfrak p \otimes_k S)$, et l'image inverse de $T'$ dans
$\Spec(R \otimes_k S)$ est $\Spec(R/I \otimes_k S)$.
Par récurrence, ces deux espaces sont connexes. L'image inverse de
$T \cap T'$ est $\Spec(R/(\mathfrak p + I) \otimes_k S)$, qui n'est pas vide.
Ainsi $\Spec(R \otimes_k S)$ est connexe.
\end{proof}

\begin{lemma}
\label{algebra-lemma-geometrically-connected}
Soit $k$ un corps.
Soit $R$ une $k$-algèbre.
Les assertions suivantes sont équivalentes
\begin{enumerate}
\item pour toute extension de corps $k'/k$, le spectre de
$R \otimes_k k'$ est connexe, et
\item pour toute extension de corps finie et séparable $k'/k$, le spectre
de $R \otimes_k k'$ est connexe.
\end{enumerate}
\end{lemma}

\begin{proof}
Pour toute extension de corps $k'/k$, la connexité du spectre de
$R \otimes_k k'$ équivaut à l'absence d'idempotents non triviaux dans
$R \otimes_k k'$, voir le Lemme \ref{algebra-lemma-characterize-spec-connected}.
Supposons (2). Soit $k \subset \overline{k}$ une clôture algébrique
séparable de $k$. Par le Lemme \ref{algebra-lemma-limit-argument}, (2) équivaut
à l'absence d'idempotents non triviaux dans $R \otimes_k \overline{k}$.
Pour toute extension de corps $k'/k$, il existe une extension
$\overline{k}'/\overline{k}$ telle que $k' \subset \overline{k}'$.
Par le Lemme \ref{algebra-lemma-separably-closed-connected},
$R \otimes_k \overline{k}'$ ne possède pas d'idempotent non trivial.
Si $R \otimes_k k'$ possédait un idempotent non trivial, il en serait de
même dans $R \otimes_k \overline{k}'$, contradiction.
\end{proof}

\begin{definition}
\label{algebra-definition-geometrically-connected}
Soit $k$ un corps.
Soit $S$ une $k$-algèbre.
Nous disons que $S$ est {\it géométriquement connexe sur $k$}
si, pour toute extension de corps $k'/k$, le spectre
de $S \otimes_k k'$ est connexe.
\end{definition}

\noindent
Par le Lemme \ref{algebra-lemma-geometrically-connected}, il suffit de vérifier
cette propriété pour les extensions de corps finies et séparables $k'/k$.

\begin{lemma}
\label{algebra-lemma-separably-closed-connected-implies-geometric}
Soit $k$ un corps.
Soit $R$ une $k$-algèbre.
Si $k$ est séparablement algébriquement clos, alors $R$ est
géométriquement connexe sur $k$ si et seulement si le spectre de
$R$ est connexe.
\end{lemma}

\begin{proof}
Cela résulte immédiatement de la remarque qui suit la
Définition \ref{algebra-definition-geometrically-connected}.
\end{proof}

\begin{lemma}
\label{algebra-lemma-subalgebra-geometrically-connected}
Soit $k$ un corps. Soit $S$ une $k$-algèbre.
\begin{enumerate}
\item Si $S$ est géométriquement connexe sur $k$, toute
$k$-sous-algèbre l'est également.
\item Si toutes les $k$-sous-algèbres de type fini de $S$ sont
géométriquement connexes, alors $S$ est géométriquement connexe.
\item Une colimite filtrante de $k$-algèbres géométriquement connexes
est géométriquement connexe.
\end{enumerate}
\end{lemma}

\begin{proof}
Cela résulte de la caractérisation de la connexité par l'absence
d'idempotents non triviaux. Les deuxième et troisième assertions résultent
du fait que le produit tensoriel commute aux colimites.
\end{proof}

\noindent
Le lemme suivant sera remplacé par le résultat plus général du
lemme \ref{varieties-lemma-bijection-connected-components} de Variétés.

\begin{lemma}
\label{algebra-lemma-geometrically-connected-any-base-change}
Soit $k$ un corps.
Soit $S$ une $k$-algèbre géométriquement connexe.
Soit $R$ une $k$-algèbre quelconque.
L'application
$$
R \longrightarrow R \otimes_k S
$$
induit une bijection sur les idempotents, et l'application
$$
\Spec(R \otimes_k S) \longrightarrow \Spec(R)
$$
induit une bijection sur les composantes connexes.
\end{lemma}

\begin{proof}
La seconde assertion résulte de la première et du
Lemme \ref{algebra-lemma-connected-component}.
Par les Lemmes \ref{algebra-lemma-subalgebra-geometrically-connected} et
\ref{algebra-lemma-limit-argument}, nous pouvons supposer que $R$ et $S$ sont de
type fini sur $k$. Alors $R \otimes_k S$ est également de type fini sur $k$.
Remarquons que, dans ce cas, tous les anneaux sont noethériens et que leurs
spectres possèdent un nombre fini de composantes connexes (puisqu'ils ont
un nombre fini de composantes irréductibles, voir le
Lemme \ref{algebra-lemma-Noetherian-irreducible-components}).
En particulier, toutes les composantes connexes en question sont ouvertes !
Par le Lemme \ref{algebra-lemma-disjoint-implies-product}, la première assertion
du lemme est alors équivalente à la seconde. Démontrons cette dernière.
Comme l'algèbre $S$ est géométriquement connexe et non nulle, toutes les
fibres de $X = \Spec(R \otimes_k S) \to \Spec(R) = Y$ sont connexes
et non vides. De plus, comme $R \to R \otimes_k S$ est plate et de
présentation finie, l'application $X \to Y$ est ouverte
(Proposition \ref{algebra-proposition-fppf-open}).
Le Lemme \ref{topology-lemma-connected-fibres-connected-components} de
Topologie montre que $X \to Y$ induit une bijection sur les composantes
connexes.
\end{proof}

\section{Algèbres géométriquement intègres}
\label{algebra-section-geometrically-integral}

\noindent
Voici la définition.

\begin{definition}
\label{algebra-definition-geometrically-integral}
Soit $k$ un corps.
Soit $S$ une $k$-algèbre.
Nous disons que $S$ est {\it géométriquement intègre sur $k$}
si, pour toute extension de corps $k'/k$, l'anneau
$S \otimes_k k'$ est un domaine intègre.
\end{definition}

\noindent
Toute question concernant les algèbres géométriquement intègres peut être
traduite en une question concernant les algèbres géométriquement réduites
et irréductibles.

\begin{lemma}
\label{algebra-lemma-geometrically-integral}
Soit $k$ un corps.
Soit $S$ une $k$-algèbre.
Alors $S$ est géométriquement intègre sur $k$ si et seulement si
$S$ est à la fois géométriquement irréductible et géométriquement réduite
sur $k$.
\end{lemma}

\begin{proof}
Démonstration omise.
\end{proof}

\begin{lemma}
\label{algebra-lemma-characterize-geometrically-integral}
Soit $k$ un corps. Soit $S$ une $k$-algèbre.
Les assertions suivantes sont équivalentes
\begin{enumerate}
\item $S$ est géométriquement intègre sur $k$,
\item pour toute extension finie de corps $k'/k$, l'anneau
$S \otimes_k k'$ est un domaine intègre,
\item $S \otimes_k \overline{k}$ est un domaine intègre,
où $\overline{k}$ est la clôture algébrique de $k$.
\end{enumerate}
\end{lemma}

\begin{proof}
Cela résulte des Lemmes \ref{algebra-lemma-geometrically-integral},
\ref{algebra-lemma-geometrically-reduced-finite-purely-inseparable-extension} et
\ref{algebra-lemma-geometrically-irreducible}.
\end{proof}

\begin{lemma}
\label{algebra-lemma-geometrically-integral-any-integral-base-change}
Soit $k$ un corps. Soit $S$ une $k$-algèbre géométriquement intègre.
Soit $R$ une $k$-algèbre qui est un domaine intègre. Alors
$R \otimes_k S$ est un domaine intègre.
\end{lemma}

\begin{proof}
Par le Lemme \ref{algebra-lemma-geometrically-reduced-any-reduced-base-change},
l'anneau $R \otimes_k S$ est réduit et, par le
Lemme \ref{algebra-lemma-geometrically-irreducible-any-base-change}, l'anneau
$R \otimes_k S$ est irréductible (son spectre possède une seule composante
irréductible), donc $R \otimes_k S$ est un domaine intègre.
\end{proof}


\section{Anneaux de valuation}
\label{algebra-section-valuation-rings}

\noindent
Voici quelques définitions.

\begin{definition}
\label{algebra-definition-valuation-ring}
Anneaux de valuation.
\begin{enumerate}
\item Soit $K$ un corps. Soient $A$, $B$ des anneaux locaux contenus
dans $K$. Nous disons que $B$ {\it domine} $A$ si $A \subset B$
et si $\mathfrak m_A = A \cap \mathfrak m_B$.
\item Soit $A$ un anneau. Nous disons que $A$ est un {\it anneau de valuation}
si $A$ est un domaine local et si $A$ est maximal
pour la relation de domination parmi les anneaux locaux contenus dans
le corps des fractions de $A$.
\item Soit $A$ un anneau de valuation de corps des fractions $K$.
Si $R \subset K$ est un sous-anneau de $K$, nous disons que $A$
est {\it centré} sur $R$ si $R \subset A$.
\end{enumerate}
\end{definition}

\noindent
Avec cette définition, un corps est un anneau de valuation.

\begin{lemma}
\label{algebra-lemma-dominate}
Soit $K$ un corps. Soit $A \subset K$ un sous-anneau local.
Alors il existe un anneau de valuation de corps des fractions $K$
qui domine $A$.
\end{lemma}

\begin{proof}
Nous considérons l'ensemble des sous-anneaux locaux
de $K$ comme un ensemble partiellement ordonné par la relation de domination.
Supposons que $\{A_i\}_{i \in I}$ soit une famille totalement ordonnée
de sous-anneaux locaux de $K$. Alors $B = \bigcup A_i$
est un sous-anneau local qui domine tous les $A_i$.
Par conséquent, par le lemme de Zorn, il suffit de montrer que si $A \subset K$
est un anneau local dont le corps des fractions n'est pas $K$, alors il
existe un anneau local $B \subset K$, $B \not = A$, qui domine $A$.

\medskip\noindent
Choisissons $t \in K$ qui n'appartient pas au corps des fractions de $A$.
Si $t$ est transcendant sur $A$, alors $A[t] \subset K$
et donc $A[t]_{(t, \mathfrak m)} \subset K$ est un anneau local
distinct de $A$ qui domine $A$. Supposons que $t$ soit algébrique sur $A$.
Alors, pour un certain $a \in A$ non nul, l'élément $at$ est entier sur $A$.
Dans ce cas, le sous-anneau $A' \subset K$ engendré par $A$ et
$ta$ est fini sur $A$.
Par le Lemme \ref{algebra-lemma-integral-overring-surjective}, il existe
un idéal premier $\mathfrak m' \subset A'$ au-dessus de
$\mathfrak m$. Alors $A'_{\mathfrak m'}$ domine
$A$. Si $A = A'_{\mathfrak m'}$, alors $t$
appartient au corps des fractions de $A$, ce que nous avons supposé faux.
Ainsi $A \not = A'_{\mathfrak m'}$, comme voulu.
\end{proof}

\begin{lemma}
\label{algebra-lemma-valuation-ring-normal}
Soit $A$ un anneau de valuation.
Alors $A$ est un domaine normal.
\end{lemma}

\begin{proof}
Supposons que $x$ appartienne au corps des fractions de $A$ et soit entier sur $A$.
Soit $A'$ le sous-anneau de $K$ engendré par $A$ et $x$.
Comme $A\subset A'$ est une extension entière, le
Lemme \ref{algebra-lemma-integral-overring-surjective} montre qu'il existe
un idéal premier $\mathfrak m' \subset A'$ au-dessus de
$\mathfrak m$. Alors $A'_{\mathfrak m'}$ domine
$A$. Comme $A$ est un anneau de valuation, nous concluons que $A=A'_{\mathfrak m'}$
et donc que $x\in A$.
\end{proof}


\begin{lemma}
\label{algebra-lemma-valuation-ring-x-or-x-inverse}
Soit $A$ un anneau de valuation d'idéal maximal $\mathfrak m$ et de
corps des fractions $K$.
Soit $x \in K$. Alors $x \in A$ ou bien $x^{-1} \in A$, ou les deux.
\end{lemma}

\begin{proof}
Supposons que $x$ n'appartienne pas à $A$.
Soit $A'$ le sous-anneau de $K$ engendré par $A$ et $x$.
Comme $A$ est un anneau de valuation, aucun idéal premier de $A'$
n'est au-dessus de $\mathfrak m$. Comme $\mathfrak m$ est maximal,
nous avons $V(\mathfrak m A') = \emptyset$. Alors $\mathfrak m A' = A'$
par le Lemme \ref{algebra-lemma-Zariski-topology}.
Nous pouvons donc écrire
$1 = \sum_{i = 0}^d t_i x^i$ avec $t_i \in \mathfrak m$.
Cela implique que $(1 - t_0) (x^{-1})^d - \sum t_i (x^{-1})^{d - i} = 0$.
En particulier, $x^{-1}$ est entier sur $A$, et donc
$x^{-1} \in A$ par le
Lemme \ref{algebra-lemma-valuation-ring-normal}.
\end{proof}

\begin{lemma}
\label{algebra-lemma-x-or-x-inverse-valuation-ring}
Soit $A \subset K$ un sous-anneau d'un corps $K$ tel que, pour tout
$x \in K$, on ait $x \in A$ ou bien $x^{-1} \in A$, ou les deux.
Alors $A$ est un anneau de valuation de corps des fractions $K$.
\end{lemma}

\begin{proof}
Si $A$ n'est pas $K$, alors $A$ n'est pas un corps et possède un idéal
maximal non nul $\mathfrak m$.
Si $\mathfrak m'$ est un second idéal maximal, choisissons $x, y \in A$
tels que $x \in \mathfrak m$, $y \not \in \mathfrak m$,
$x \not \in \mathfrak m'$, et $y \in \mathfrak m'$.
Alors ni $x/y \in A$ ni $y/x \in A$,
ce qui contredit l'hypothèse du lemme. Ainsi $A$ est
un anneau local. Supposons que $A'$ soit un anneau local contenu dans $K$
qui domine $A$. Soit $x \in A'$. Nous devons montrer que $x \in A$. Sinon,
$x^{-1} \in A$, et bien sûr $x^{-1} \in \mathfrak m_A$. Mais alors
$x^{-1} \in \mathfrak m_{A'}$, ce qui contredit $x \in A'$.
\end{proof}


\begin{lemma}
\label{algebra-lemma-colimit-valuation-rings}
\begin{slogan}
Les anneaux de valuation sont stables par limites directes filtrantes
\end{slogan}
Soit $I$ un ensemble dirigé. Soit $(A_i, \varphi_{ij})$
un système d'anneaux de valuation indexé par $I$.
Alors $A = \colim A_i$ est un anneau de valuation.
\end{lemma}

\begin{proof}
Il est clair que $A$ est un domaine. Soient $a, b \in A$.
Le Lemme \ref{algebra-lemma-x-or-x-inverse-valuation-ring} nous dit qu'il faut
montrer que $a | b$ ou $b | a$ dans $A$. Choisissons $i$
assez grand pour qu'il existe $a_i, b_i \in A_i$ dont les images soient $a, b$.
Alors le Lemme \ref{algebra-lemma-valuation-ring-x-or-x-inverse}
appliqué à $a_i, b_i$ dans $A_i$ donne le résultat pour $a, b$ dans $A$.
\end{proof}

\begin{lemma}
\label{algebra-lemma-valuation-ring-cap-field}
Soit $L/K$ une extension de corps. Si $B \subset L$
est un anneau de valuation, alors $A = K \cap B$ est un anneau de valuation.
\end{lemma}

\begin{proof}
Nous pouvons remplacer $L$ par le corps des fractions $F$ de $B$ et $K$ par
$K \cap F$. Le lemme résulte alors d'une combinaison des
Lemmes \ref{algebra-lemma-valuation-ring-x-or-x-inverse} et
\ref{algebra-lemma-x-or-x-inverse-valuation-ring}.
\end{proof}

\begin{lemma}
\label{algebra-lemma-valuation-ring-cap-field-finite}
Soit $L/K$ une extension algébrique de corps. Si $B \subset L$
est un anneau de valuation de corps des fractions $L$ et n'est pas un corps, alors
$A = K \cap B$ est un anneau de valuation et n'est pas un corps.
\end{lemma}

\begin{proof}
Par le Lemme \ref{algebra-lemma-valuation-ring-cap-field}, l'anneau $A$ est un anneau de
valuation. Si $A$ est un corps, alors $A = K$. Alors $A = K \subset B$ est une extension
entière, donc il n'y a pas d'inclusions propres entre les idéaux premiers de $B$
(Lemme \ref{algebra-lemma-integral-no-inclusion}).
Cela contredit le fait que $B$ est un domaine local et n'est pas un corps.
\end{proof}

\begin{lemma}
\label{algebra-lemma-make-valuation-rings}
Soit $A$ un anneau de valuation. Pour tout idéal premier $\mathfrak p \subset A$, le
quotient $A/\mathfrak p$ est un anneau de valuation. Il en va de même pour la
localisation $A_\mathfrak p$ et, en fait, pour toute localisation de $A$.
\end{lemma}

\begin{proof}
Utiliser la caractérisation des anneaux de valuation donnée
dans le Lemme \ref{algebra-lemma-x-or-x-inverse-valuation-ring}.
\end{proof}

\begin{lemma}
\label{algebra-lemma-stack-valuation-rings}
Soit $A'$ un anneau de valuation de corps résiduel $K$.
Soit $A$ un anneau de valuation de corps des fractions $K$.
Alors
$C = \{\lambda \in A' \mid \lambda \bmod \mathfrak m_{A'} \in A\}$
est un anneau de valuation.
\end{lemma}

\begin{proof}
Remarquons que $\mathfrak m_{A'} \subset C$ et que $C/\mathfrak m_{A'} = A$.
En particulier, le corps des fractions de $C$ est égal au corps des fractions
de $A'$. Nous utiliserons le critère du
Lemme \ref{algebra-lemma-x-or-x-inverse-valuation-ring} pour démontrer le lemme.
Soit $x$ un élément du corps des fractions de $C$.
D'après le lemme, nous pouvons supposer que $x \in A'$. Si $x \in \mathfrak m_{A'}$,
alors $x \in C$. Sinon, $x$ est une unité de $A'$ et nous avons
également $x^{-1} \in A'$. Ainsi $x$ ou $x^{-1}$ s'envoie sur
un élément de $A$, à nouveau par le lemme.
\end{proof}

\begin{lemma}
\label{algebra-lemma-find-valuation-rings}
Soit $A$ un domaine normal de corps des fractions $K$.
\begin{enumerate}
\item Pour tout $x \in K$, $x \not \in A$, il existe un anneau de valuation
$A \subset V \subset K$ de corps des fractions $K$ tel que $x \not \in V$.
\item Si $A$ est local, nous pouvons de plus choisir $V$ qui domine $A$.
\end{enumerate}
Autrement dit, $A$ est l'intersection de tous les anneaux de valuation dans $K$
contenant $A$ et, si $A$ est local, $A$ est l'intersection de tous les anneaux de
valuation dans $K$ qui dominent $A$.
\end{lemma}

\begin{proof}
Supposons que $x \in K$, $x \not \in A$. Considérons $B = A[x^{-1}]$.
Alors $x \not \in B$. En effet, si $x = a_0 + a_1x^{-1} + \ldots + a_d x^{-d}$
alors $x^{d + 1} - a_0x^d - \ldots - a_d = 0$ et $x$ est entier
sur $A$, contradiction avec le fait que $A$ est normal.
Ainsi $x^{-1}$ n'est pas une unité de $B$. Donc $V(x^{-1}) \subset \Spec(B)$
n'est pas vide (Lemme \ref{algebra-lemma-Zariski-topology}), et nous pouvons choisir un idéal premier
$\mathfrak p \subset B$ tel que $x^{-1} \in \mathfrak p$.
Choisissons un anneau de valuation $V \subset K$ dominant $B_\mathfrak p$
(Lemme \ref{algebra-lemma-dominate}).
Alors $x \not \in V$ puisque $x^{-1} \in \mathfrak m_V$.

\medskip\noindent
Si $A$ est local, nous affirmons que $x^{-1} B + \mathfrak m_A B \not = B$.
En effet, si $1 = (a_0 + a_1x^{-1} + \ldots + a_d x^{-d})x^{-1} +
a'_0 + \ldots + a'_d x^{-d}$ avec $a_i \in A$ et $a'_i \in \mathfrak m_A$,
alors nous obtenons
$$
(1 - a'_0) x^{d + 1} - (a_0 + a'_1) x^d - \ldots - a_d = 0
$$
Comme $a'_0 \in \mathfrak m_A$, $1 - a'_0$ est une unité de $A$
et nous concluons que $x$ serait entier sur $A$, contradiction comme
précédemment. Choisissons alors l'idéal premier $\mathfrak p \supset x^{-1} B + \mathfrak m_A B$ ;
nous trouvons $V$ dominant $A$.
\end{proof}

\noindent
Un {\it groupe abélien totalement ordonné} est un couple $(\Gamma, \geq)$
constitué d'un groupe abélien $\Gamma$ muni d'un
ordre total $\geq$ tel que $\gamma \geq \gamma' \Rightarrow
\gamma + \gamma'' \geq \gamma' + \gamma''$ pour tous
$\gamma, \gamma', \gamma'' \in \Gamma$.

\begin{lemma}
\label{algebra-lemma-valuation-group}
Soit $A$ un anneau de valuation de corps des fractions $K$.
Posons $\Gamma = K^*/A^*$ (la loi de groupe étant écrite additivement).
Pour $\gamma, \gamma' \in \Gamma$,
définissons $\gamma \geq \gamma'$ si et seulement si
$\gamma - \gamma'$ appartient à l'image de $A - \{0\} \to \Gamma$.
Alors $(\Gamma, \geq)$ est un groupe abélien totalement ordonné.
\end{lemma}

\begin{proof}
Démonstration omise, mais le résultat découle facilement du
Lemme \ref{algebra-lemma-valuation-ring-x-or-x-inverse}.
Remarquons que dans le cas $A = K$, on obtient le groupe nul
$\Gamma = \{0\}$ muni de son unique ordre total.
\end{proof}

\begin{definition}
\label{algebra-definition-value-group}
Soit $A$ un anneau de valuation.
\begin{enumerate}
\item Le groupe abélien totalement ordonné $(\Gamma, \geq)$ du
Lemme \ref{algebra-lemma-valuation-group} est appelé le
{\it groupe des valeurs} de l'anneau de valuation $A$.
\item L'application $v : A - \{0\} \to \Gamma$ et aussi $v : K^* \to \Gamma$ est
appelée la {\it valuation} associée à $A$.
\item L'anneau de valuation $A$ est appelé un {\it anneau de valuation discrète}
si $\Gamma \cong \mathbf{Z}$.
\end{enumerate}
\end{definition}

\noindent
Remarquons que si $\Gamma \cong \mathbf{Z}$, il existe un unique isomorphisme de ce type
tel que $1 \geq 0$. Si l'isomorphisme est choisi de cette manière,
l'ordre devient l'ordre usuel des entiers.

\begin{lemma}
\label{algebra-lemma-properties-valuation}
Soit $A$ un anneau de valuation. La valuation $v : A -\{0\} \to \Gamma_{\geq 0}$
possède les propriétés suivantes :
\begin{enumerate}
\item $v(a) = 0 \Leftrightarrow a \in A^*$,
\item $v(ab) = v(a) + v(b)$,
\item $v(a + b) \geq \min(v(a), v(b))$ pourvu que $a + b \not = 0$.
\end{enumerate}
\end{lemma}

\begin{proof}
Démonstration omise.
\end{proof}

\begin{lemma}
\label{algebra-lemma-characterize-valuation-ring}
Soit $A$ un anneau. Les assertions suivantes sont équivalentes
\begin{enumerate}
\item $A$ est un anneau de valuation,
\item $A$ est un domaine local et tout idéal
de $A$ engendré par un nombre fini d'éléments est principal.
\end{enumerate}
\end{lemma}

\begin{proof}
Supposons que $A$ soit un anneau de valuation et soient $f_1, \ldots, f_n \in A$.
Choisissons $i$ tel que $v(f_i)$ soit minimal parmi les $v(f_j)$.
Alors $(f_i) = (f_1, \ldots, f_n)$. Réciproquement, supposons que $A$ soit
un domaine local et que tout idéal de $A$ engendré par un nombre fini d'éléments soit principal.
Prenons $f, g \in A$ et écrivons $(f, g) = (h)$. Alors $f = ah$ et $g = bh$
et $h = cf + dg$ pour certains $a, b, c, d \in A$. Ainsi $ac + bd = 1$
et nous voyons que $a$ ou $b$ est une unité, c'est-à-dire que $g/f$
ou $f/g$ appartient à $A$. Cela montre que $A$ est un anneau de valuation
par le Lemme \ref{algebra-lemma-x-or-x-inverse-valuation-ring}.
\end{proof}

\begin{lemma}
\label{algebra-lemma-valuation-valuation-ring}
Soit $(\Gamma, \geq)$ un groupe abélien totalement ordonné.
Soit $K$ un corps. Soit $v : K^* \to \Gamma$ un homomorphisme
de groupes abéliens tel que $v(a + b) \geq \min(v(a), v(b))$ pour
$a, b \in K$ avec $a, b, a + b$ non nuls. Alors
$$
A =
\{
x \in K \mid x = 0 \text{ ou } v(x) \geq 0
\}
$$
est un anneau de valuation de groupe des valeurs $\Im(v) \subset \Gamma$,
d'idéal maximal
$$
\mathfrak m =
\{
x \in K \mid x = 0 \text{ ou } v(x) > 0
\}
$$
et de groupe des unités
$$
A^* =
\{
x \in K^* \mid v(x) = 0
\}.
$$
\end{lemma}

\begin{proof}
Démonstration omise.
\end{proof}

\noindent
Soit $(\Gamma, \geq)$ un groupe abélien totalement ordonné.
Un {\it idéal de $\Gamma$} est un sous-ensemble $I \subset \Gamma$ tel
que tous les éléments de $I$ soient $\geq 0$ et que, si $\gamma \in I$,
$\gamma' \geq \gamma$ implique $\gamma' \in I$. Nous disons qu'un tel
idéal est {\it premier} si $0 \not \in I$ et si
$\gamma + \gamma' \in I, \gamma, \gamma' \geq 0
\Rightarrow \gamma \in I \text{ ou } \gamma' \in I$.

\begin{lemma}
\label{algebra-lemma-ideals-valuation-ring}
Soit $A$ un anneau de valuation.
Les idéaux de $A$ correspondent $1 - 1$ aux idéaux de $\Gamma$.
Cette bijection respecte l'inclusion et envoie les idéaux premiers
sur les idéaux premiers.
\end{lemma}

\begin{proof}
Démonstration omise.
\end{proof}

\begin{lemma}
\label{algebra-lemma-valuation-ring-Noetherian-discrete}
Un anneau de valuation est noethérien si et seulement si c'est
un anneau de valuation discrète ou un corps.
\end{lemma}

\begin{proof}
Supposons que $A$ soit un anneau de valuation discrète
de valuation $v : A \setminus \{0\} \to \mathbf{Z}$
normalisée de sorte que $\Im(v) = \mathbf{Z}_{\geq 0}$.
Par le Lemme \ref{algebra-lemma-ideals-valuation-ring}, les idéaux de $A$ sont les sous-ensembles
$I_n = \{0\} \cup v^{-1}(\mathbf{Z}_{\geq n})$. Il est clair
que tout élément $x \in A$ tel que $v(x) = n$ engendre $I_n$.
Ainsi $A$ est un PID, donc certainement noethérien.

\medskip\noindent
Supposons que $A$ soit un anneau de valuation noethérien de groupe des valeurs $\Gamma$.
Par le Lemme \ref{algebra-lemma-ideals-valuation-ring}, la condition de chaîne ascendante
est satisfaite pour les idéaux de $\Gamma$.
Nous pouvons supposer que $A$ n'est pas un corps, c'est-à-dire qu'il existe un $\gamma \in \Gamma$
tel que $\gamma > 0$. En appliquant la condition de chaîne ascendante aux sous-ensembles
$\gamma + \Gamma_{\geq 0}$ avec $\gamma > 0$, nous voyons
qu'il existe un plus petit élément $\gamma_0$ strictement supérieur à $0$.
Soit $\gamma \in \Gamma$ un élément tel que $\gamma > 0$. Considérons la suite
des éléments $\gamma$, $\gamma - \gamma_0$, $\gamma - 2\gamma_0$,
etc. Par la condition de chaîne ascendante, ils ne peuvent pas tous être $> 0$.
Soit $\gamma - n \gamma_0$ le dernier qui soit $\geq 0$. Par minimalité
de $\gamma_0$, nous voyons que $0 = \gamma - n \gamma_0$. Ainsi $\Gamma$
est un groupe cyclique, comme voulu.
\end{proof}


\section{Autres anneaux noethériens}
\label{algebra-section-Noetherian-again}


\begin{lemma}
\label{algebra-lemma-Noetherian-basic}
Soit $R$ un anneau noethérien.
Tout $R$-module fini est de présentation finie.
Tout sous-module d'un $R$-module fini est fini.
La condition de chaîne ascendante est satisfaite pour les sous-modules
$R$-d'un $R$-module fini.
\end{lemma}

\begin{proof}
Montrons d'abord que tout sous-module $N$ d'un $R$-module fini
$M$ est fini. Nous procédons par récurrence sur le nombre de
générateurs de $M$. Si ce nombre vaut $1$, alors $N = J/I \subset
M = R/I$ pour certains idéaux $I \subset J \subset R$. Ainsi la définition
de noethérien donne le résultat. Si le nombre de générateurs de
$M$ est supérieur à $1$, nous pouvons trouver une suite exacte courte
$0 \to M' \to M \to M'' \to 0$ où $M'$ et $M''$ ont moins de
générateurs. Remarquons qu'en posant $N' = M' \cap N$ et
$N'' = \Im(N \to M'')$, nous obtenons une suite exacte courte analogue pour $N$.
Le résultat découle donc de l'hypothèse de récurrence
puisque le nombre de générateurs de $N$ est au plus le nombre de
générateurs de $N'$ plus le nombre de générateurs de $N''$.

\medskip\noindent
Pour montrer que $M$ est de présentation finie, il suffit d'appliquer le résultat précédent
au noyau d'une présentation $R^n \to M$.

\medskip\noindent
Il est bien connu et facile de démontrer que la condition de chaîne ascendante pour les
sous-modules $R$ de $M$ est équivalente à la condition que tout sous-module
de $M$ soit un $R$-module fini. Nous omettons la démonstration.
\end{proof}

\begin{lemma}[Artin-Rees]
\label{algebra-lemma-Artin-Rees}
Supposons que $R$ soit noethérien et que $I \subset R$ soit un idéal.
Soient $N \subset M$ des $R$-modules finis.
Il existe une constante $c > 0$ telle que
$I^n M \cap N  =  I^{n-c}(I^cM \cap N)$ pour tout $n \geq c$.
\end{lemma}

\begin{proof}
Considérons l'anneau $S = R \oplus I \oplus I^2 \oplus \ldots
= \bigoplus_{n \geq 0} I^n$. Convention : $I^0 = R$.
La multiplication envoie $I^n \times I^m$
dans $I^{n + m}$ par la multiplication dans $R$.
Remarquons que si $I = (f_1, \ldots, f_t)$,
alors $S$ est un quotient de l'anneau noethérien $R[X_1, \ldots, X_t]$.
L'application envoie simplement le monôme $X_1^{e_1}\ldots X_t^{e_t}$
sur $f_1^{e_1}\ldots f_t^{e_t}$. Ainsi $S$ est noethérien.
De même, considérons le module $M \oplus IM \oplus I^2M \oplus \ldots
= \bigoplus_{n \geq 0} I^nM$. C'est un $S$-module de type fini.
En effet, si $x_1, \ldots, x_r$ engendrent $M$ sur $R$, ils engendrent aussi
$\bigoplus_{n \geq 0} I^nM$ sur $S$. Ensuite, considérons le
sous-module $\bigoplus_{n \geq 0} I^nM \cap N$.
Il s'agit d'un $S$-sous-module, comme on le vérifie facilement. Par le
Lemme \ref{algebra-lemma-Noetherian-basic}, il est de type fini comme
$S$-module,
disons engendré par $\xi_j \in \bigoplus_{n \geq 0} I^nM \cap N$, $j = 1, \ldots, s$.
Nous pouvons supposer, en décomposant chaque $\xi_j$ en ses composantes homogènes,
que $\xi_j \in I^{d_j}M \cap N$ pour un certain $d_j$.
Posons $c = \max\{d_j\}$. Alors, pour tout $n \geq c$, tout élément
de $I^nM \cap N$ est de la forme $\sum h_j \xi_j$ avec
$h_j \in I^{n - d_j}$. Le lemme résulte alors de cette observation et du fait trivial que
$I^{n-d_j}(I^{d_j}M \cap N) \subset I^{n-c}(I^cM \cap N)$.
\end{proof}

\begin{lemma}
\label{algebra-lemma-map-AR}
Supposons que $0 \to K \to M \xrightarrow{f} N$ soit une
suite exacte de modules de type fini
sur un anneau noethérien $R$. Soit $I \subset R$ un idéal.
Alors il existe un $c$ tel que
$$
f^{-1}(I^nN) = K + I^{n-c}f^{-1}(I^cN)
\quad\text{et}\quad
f(M) \cap I^nN \subset f(I^{n - c}M)
$$
pour tout $n \geq c$.
\end{lemma}

\begin{proof}
Appliquer le Lemme \ref{algebra-lemma-Artin-Rees} à
$\Im(f) \subset N$ et remarquer que
$f : I^{n-c}M \to I^{n-c}f(M)$ est surjectif.
\end{proof}

\begin{lemma}[théorème d'intersection de Krull]
\label{algebra-lemma-intersect-powers-ideal-module-zero}
Soit $R$ un anneau local noethérien. Soit $I \subset R$ un
idéal propre. Soit $M$ un $R$-module fini.
Alors $\bigcap_{n \geq 0} I^nM = 0$.
\end{lemma}

\begin{proof}
Soit $N = \bigcap_{n \geq 0} I^nM$.
Alors $N = I^nM \cap N$ pour tout $n \geq 0$.
Par le Lemme d'Artin-Rees \ref{algebra-lemma-Artin-Rees},
nous voyons que $N = I^nM \cap N \subset IN$
pour un certain $n$ suffisamment grand. Par le Lemme de Nakayama
\ref{algebra-lemma-NAK}, nous avons $N = 0$.
\end{proof}

\begin{lemma}
\label{algebra-lemma-intersection-powers-ideal-module}
Soit $R$ un anneau noethérien. Soit $I \subset R$ un idéal.
Soit $M$ un $R$-module fini. Soit $N = \bigcap_n I^n M$.
\begin{enumerate}
\item Pour tout idéal premier $\mathfrak p$, $I \subset \mathfrak p$, il
existe un $f \in R$, $f \not \in \mathfrak p$, tel que $N_f = 0$.
\item Si $I$ est contenu dans le radical de Jacobson
de $R$, alors $N = 0$.
\end{enumerate}
\end{lemma}

\begin{proof}
Démontrons (1). Soient $x_1, \ldots, x_n$ des générateurs du module $N$,
voir le Lemme \ref{algebra-lemma-Noetherian-basic}. Pour tout idéal premier
$\mathfrak p$, $I \subset \mathfrak p$, l'image de $N$ dans la localisation $M_{\mathfrak p}$
est nulle, par le Lemme \ref{algebra-lemma-intersect-powers-ideal-module-zero}.
Nous pouvons donc trouver $g_i \in R$, $g_i \not \in \mathfrak p$
tel que $x_i$ s'envoie sur zéro dans $N_{g_i}$. Ainsi
$N_{g_1g_2\ldots g_n} = 0$.

\medskip\noindent
La partie (2) découle de (1) et du Lemme \ref{algebra-lemma-characterize-zero-local}.
\end{proof}

\begin{remark}
\label{algebra-remark-intersection-powers-ideal}
Le Lemme \ref{algebra-lemma-intersect-powers-ideal-module-zero} implique en particulier
que $\bigcap_n I^n = (0)$ lorsque $I \subset R$ est un idéal non unité dans un
anneau local noethérien $R$. Plus généralement, soit $R$ un anneau noethérien et
$I \subset R$ un idéal. Supposons que $f \in \bigcap_{n \in \mathbf{N}} I^n$.
Alors le Lemme \ref{algebra-lemma-intersection-powers-ideal-module}
dit que, pour tout idéal premier $I \subset \mathfrak p$,
il existe un $g \in R$, $g \not \in \mathfrak p$,
tel que $f$ s'envoie sur zéro dans $R_g$. En géométrie algébrique, nous
exprimons cela en disant que « $f$ est nul dans un voisinage ouvert
de l'ensemble fermé $V(I)$ de $\Spec(R)$ ».
\end{remark}

\begin{lemma}[Artin-Tate]
\label{algebra-lemma-Artin-Tate}
Soit $R$ un anneau noethérien. Soit $S$ une $R$-algèbre
de type fini. Si $T \subset S$ est une $R$-sous-algèbre telle que
$S$ soit de type fini comme $T$-module, alors $T$ est de
type fini sur $R$.
\end{lemma}

\begin{proof}
Choisissons des éléments $x_1, \ldots, x_n \in S$ qui engendrent $S$ comme $R$-algèbre.
Choisissons $y_1, \ldots, y_m$ dans $S$ qui engendrent $S$ comme $T$-module.
Il existe donc des $a_{ij} \in T$ tels que
$x_i = \sum a_{ij} y_j$. Il existe également des $b_{ijk} \in T$ tels
que $y_i y_j = \sum b_{ijk} y_k$. Soit $T' \subset T$ la
$R$-sous-algèbre engendrée par les $a_{ij}$ et $b_{ijk}$. Il s'agit d'une $R$-algèbre
de type fini, donc noethérienne. Considérons l'algèbre
$$
S' = T'[Y_1, \ldots, Y_m]/(Y_i Y_j - \sum b_{ijk} Y_k).
$$
Remarquons que $S'$ est fini sur $T'$, car en tant que $T'$-module il est
engendré par les classes de $1, Y_1, \ldots, Y_m$.
Considérons l'homomorphisme de $T'$-algèbres $S' \to S$ qui envoie
$Y_i$ sur $y_i$. Comme $a_{ij} \in T'$, nous voyons que $x_j$ appartient à l'image de cette application. Ainsi $S' \to S$ est surjectif.
Par conséquent, $S$ est également fini sur $T'$. Comme $T'$ est noethérien,
nous concluons que $T \subset S$ est fini sur $T'$ et nous avons terminé.
\end{proof}


\section{Longueur}
\label{algebra-section-length}

\begin{definition}
\label{algebra-definition-length}
Soit $R$ un anneau. Pour tout $R$-module $M$,
nous définissons la {\it longueur} de $M$ sur $R$ par la
formule
$$
\text{length}_R(M)
=
\sup
\{
n
\mid
\exists\ 0 = M_0 \subset M_1 \subset \ldots \subset M_n = M,
\text{ }M_i \not = M_{i + 1}
\}.
$$
\end{definition}

\noindent
Autrement dit, c'est le supremum des longueurs des chaînes
de sous-modules. Il existe une notion évidente de chaîne
de sous-modules qui raffine une autre. Cela donne un
ordre partiel sur l'ensemble de toutes les chaînes de sous-modules,
la chaîne minimale ayant la forme $0 = M_0 \subset M_1 = M$
si $M$ n'est pas nul.
Remarquons le fait évident que si la longueur de
$M$ est finie, toute chaîne peut être raffinée en une
chaîne maximale. Mais il est moins évident que toutes les chaînes maximales
ont la même longueur (comme nous le verrons plus loin).

\begin{lemma}
\label{algebra-lemma-finite-length-finite}
\begin{slogan}
Les modules de longueur finie sont finis.
\end{slogan}
Soit $R$ un anneau.
Soit $M$ un $R$-module.
Si $\text{length}_R(M) < \infty$, alors $M$ est un $R$-module fini.
\end{lemma}

\begin{proof}
Démonstration omise.
\end{proof}

\begin{lemma}
\label{algebra-lemma-length-additive}
\begin{slogan}
La longueur est additive dans les suites exactes courtes.
\end{slogan}
Si $0 \to M' \to M \to M'' \to 0$
est une suite exacte courte de modules sur $R$, alors
la longueur de $M$ est la somme des
longueurs de $M'$ et $M''$.
\end{lemma}

\begin{proof}
Étant données des filtrations de $M'$ et $M''$ de longueurs $n', n''$,
il est facile de construire une filtration correspondante de $M$
de longueur $n' + n''$. Ainsi
$\text{length}_R M \geq \text{length}_R M' + \text{length}_R M''$.
Réciproquement, étant donnée une filtration
$M_0 \subset M_1 \subset \ldots \subset M_n$ de
$M$, considérons les filtrations induites
$M_i' = M_i \cap M'$ et $M_i'' = \Im(M_i \to M'')$.
Soit $n'$ (resp.\ $n''$) le nombre d'étapes de la filtration
$\{M'_i\}$ (resp.\ $\{M''_i\}$).
Si $M_i' = M_{i + 1}'$ et $M_i'' = M_{i + 1}''$, alors
$M_i = M_{i + 1}$. Nous concluons donc que $n' + n'' \geq n$.
Avec le résultat précédent, c'est démontré.
\end{proof}

\begin{lemma}
\label{algebra-lemma-length-infinite}
Soit $R$ un anneau local d'idéal maximal $\mathfrak m$.
Si $M$ est un $R$-module et $\mathfrak m^n M \not = 0$ pour tout $n \geq 0$,
alors $\text{length}_R(M) = \infty$. Autrement dit, si $M$
est de longueur finie, alors $\mathfrak m^nM = 0$ pour un certain $n$.
\end{lemma}

\begin{proof}
Supposons que $\mathfrak m^n M \not = 0$ pour tout $n\geq 0$.
Choisissons $x \in M$ et $f_1, \ldots, f_n \in \mathfrak m$
tels que $f_1f_2 \ldots f_n x \not = 0$.
Les $n$ premières étapes de la filtration
$$
0 \subset R f_1 \ldots f_n x \subset R f_1 \ldots f_{n - 1} x
\subset \ldots \subset R x \subset M
$$
sont distinctes. Par exemple, si
$R f_1 x = R f_1 f_2 x$, alors $f_1 x = g f_1 f_2 x$ pour un certain $g$,
d'où $(1 - gf_2) f_1 x = 0$, donc $f_1 x = 0$ puisque $1 - gf_2$ est une unité,
ce qui contredit le choix de $x$ et de $f_1, \ldots, f_n$.
Ainsi la longueur est infinie.
\end{proof}

\begin{lemma}
\label{algebra-lemma-length-independent}
Soit $R \to S$ un morphisme d'anneaux. Soit $M$ un $S$-module.
Nous avons toujours $\text{length}_R(M) \geq \text{length}_S(M)$.
Si $R \to S$ est surjectif, alors il y a égalité.
\end{lemma}

\begin{proof}
Une filtration de $M$ par des $S$-sous-modules donne une filtration
de $M$ par des $R$-sous-modules. Cela démontre l'inégalité.
Et si $R \to S$ est surjectif, tout $R$-sous-module
de $M$ est automatiquement un $S$-sous-module. Il y a donc égalité
dans ce cas.
\end{proof}

\begin{lemma}
\label{algebra-lemma-dimension-is-length}
Soit $R$ un anneau d'idéal maximal $\mathfrak m$.
Supposons que $M$ soit un $R$-module tel que
$\mathfrak m M  =  0$. Alors la longueur de $M$ comme
$R$-module coïncide avec la dimension de $M$ comme
espace vectoriel sur $R/\mathfrak m$.
La longueur est finie si et seulement si $M$ est un $R$-module fini.
\end{lemma}

\begin{proof}
La première assertion est un cas particulier du
Lemme \ref{algebra-lemma-length-independent}.
Ainsi la longueur est finie si et seulement si $M$ possède une base finie
comme espace vectoriel sur $R/\mathfrak m$, si et seulement si $M$ possède un ensemble fini
de générateurs comme $R$-module.
\end{proof}

\begin{lemma}
\label{algebra-lemma-length-localize}
Soit $R$ un anneau. Soit $M$ un $R$-module. Soit $S \subset R$ un
ensemble multiplicatif. Alors
$\text{length}_R(M) \geq \text{length}_{S^{-1}R}(S^{-1}M)$.
\end{lemma}

\begin{proof}
Tout sous-module $N' \subset S^{-1}M$ est de la forme
$S^{-1}N$ pour un certain sous-module $R$-module $N \subset M$, par le Lemme
\ref{algebra-lemma-submodule-localization}. Le lemme en découle.
\end{proof}

\begin{lemma}
\label{algebra-lemma-length-finite}
Soit $R$ un anneau dont l'idéal maximal $\mathfrak m$ est engendré par un nombre fini d'éléments.
(Par exemple, $R$ est noethérien.)
Supposons que $M$ soit un $R$-module fini tel que
$\mathfrak m^n M  =  0$ pour un certain $n$.
Alors $\text{length}_R(M) < \infty$.
\end{lemma}

\begin{proof}
Considérons la filtration
$0 = \mathfrak m^n M \subset
\mathfrak m^{n-1} M \subset
\ldots \subset \mathfrak m M \subset M$.
Tous les quotients successifs sont des $R$-modules de type fini
auxquels s'applique le Lemme \ref{algebra-lemma-dimension-is-length}. Nous concluons
par additivité, voir le Lemme \ref{algebra-lemma-length-additive}.
\end{proof}


\begin{definition}
\label{algebra-definition-simple-module}
Soit $R$ un anneau. Soit $M$ un $R$-module.
Nous disons que $M$ est {\it simple} si $M \not = 0$ et
tout sous-module de $M$ est soit égal à $M$, soit
à $0$.
\end{definition}

\begin{lemma}
\label{algebra-lemma-characterize-length-1}
Soit $R$ un anneau. Soit $M$ un $R$-module.
Les assertions suivantes sont équivalentes :
\begin{enumerate}
\item $M$ est simple,
\item $\text{length}_R(M) = 1$, et
\item $M \cong R/\mathfrak m$ pour un idéal maximal
$\mathfrak m \subset R$.
\end{enumerate}
\end{lemma}

\begin{proof}
Soit $\mathfrak m$ un idéal maximal de $R$.
Par le Lemme \ref{algebra-lemma-dimension-is-length}, le module
$R/\mathfrak m$ est de longueur $1$. L'équivalence des deux
premières assertions est tautologique.
Supposons que $M$ soit simple. Choisissons $x \in M$, $x \not = 0$.
Comme $M$ est simple, nous avons $M = R \cdot x$.
Soit $I \subset R$ l'annulateur de $x$, c'est-à-dire
$I = \{f \in R \mid fx = 0\}$. L'application $R/I \to M$,
$f \bmod I \mapsto fx$ est un isomorphisme, donc
$R/I$ est un $R$-module simple. Comme $R/I \not = 0$, nous voyons que $I \not = R$.
Soit $I \subset \mathfrak m$ un idéal maximal contenant $I$.
Si $I \not = \mathfrak m$, alors $\mathfrak m /I \subset R/I$
est un sous-module non trivial, ce qui contredit la simplicité
de $R/I$. Nous obtenons donc $I = \mathfrak m$, comme voulu.
\end{proof}

\begin{lemma}
\label{algebra-lemma-simple-pieces}
Soit $R$ un anneau. Soit $M$ un $R$-module de longueur finie.
Choisissons une chaîne maximale quelconque de sous-modules
$$
0 = M_0 \subset M_1 \subset M_2 \subset \ldots \subset M_n = M
$$
telle que $M_i \not = M_{i-1}$, $i = 1, \ldots, n$. Alors
\begin{enumerate}
\item $n = \text{length}_R(M)$,
\item chaque $M_i/M_{i-1}$ est simple,
\item chaque $M_i/M_{i-1}$ est de la forme
$R/\mathfrak m_i$ pour un certain idéal maximal $\mathfrak m_i$,
\item étant donné un idéal maximal $\mathfrak m \subset R$,
nous avons
$$
\# \{i \mid \mathfrak m_i = \mathfrak m\}
=
\text{length}_{R_{\mathfrak m}} (M_{\mathfrak m}).
$$
\end{enumerate}
\end{lemma}

\begin{proof}
Si $M_i/M_{i-1}$ n'est pas simple, nous pouvons raffiner la filtration
et la filtration n'est pas maximale. Ainsi $M_i/M_{i-1}$
est simple. Par le Lemme \ref{algebra-lemma-characterize-length-1}, les modules
$M_i/M_{i-1}$ sont de longueur $1$ et sont de la forme $R/\mathfrak m_i$
pour certains idéaux maximaux $\mathfrak m_i$. Par additivité de la longueur,
Lemme \ref{algebra-lemma-length-additive}, nous voyons que $n = \text{length}_R(M)$.
Comme la localisation est exacte, nous voyons que
$$
0 = (M_0)_{\mathfrak m}
\subset (M_1)_{\mathfrak m}
\subset (M_2)_{\mathfrak m}
\subset \ldots
\subset (M_n)_{\mathfrak m} = M_{\mathfrak m}
$$
est une filtration de $M_{\mathfrak m}$ dont les quotients successifs sont
$(M_i/M_{i-1})_{\mathfrak m}$. La dernière assertion découle alors
directement du fait que pour des idéaux maximaux $\mathfrak m$,
$\mathfrak m'$ de $R$, nous avons
$$
(R/\mathfrak m')_{\mathfrak m}
\cong
\left\{
\begin{matrix}
0 &
\text{si } \mathfrak m \not = \mathfrak m', \\
R_{\mathfrak m}/\mathfrak m R_{\mathfrak m} &
\text{si } \mathfrak m  = \mathfrak m'
\end{matrix}
\right.
$$
Nous laissons cela au lecteur.
\end{proof}

\begin{lemma}
\label{algebra-lemma-pushdown-module}
Soit $A$ un anneau local d'idéal maximal $\mathfrak m$.
Soit $B$ un anneau semi-local d'idéaux maximaux $\mathfrak m_i$,
$i = 1, \ldots, n$.
Supposons que $A \to B$ soit un homomorphisme tel que chaque $\mathfrak m_i$
soit au-dessus de $\mathfrak m$ et tel que
$$
[\kappa(\mathfrak m_i) : \kappa(\mathfrak m)] < \infty.
$$
Soit $M$ un $B$-module de longueur finie.
Alors
$$
\text{length}_A(M) = \sum\nolimits_{i = 1, \ldots, n}
[\kappa(\mathfrak m_i) : \kappa(\mathfrak m)]
\text{length}_{B_{\mathfrak m_i}}(M_{\mathfrak m_i}),
$$
en particulier $\text{length}_A(M) < \infty$.
\end{lemma}

\begin{proof}
Choisissons une chaîne maximale
$$
0 = M_0
\subset M_1
\subset M_2
\subset \ldots
\subset M_m = M
$$
par des $B$-sous-modules comme dans le Lemme \ref{algebra-lemma-simple-pieces}.
Alors chaque quotient $M_j/M_{j - 1}$ est isomorphe à
$\kappa(\mathfrak m_{i(j)})$ pour un certain $i(j) \in \{1, \ldots, n\}$.
De plus
$\text{length}_A(\kappa(\mathfrak m_i)) =
[\kappa(\mathfrak m_i) : \kappa(\mathfrak m)]$ par le
Lemme \ref{algebra-lemma-dimension-is-length}. Le lemme découle
de l'additivité des longueurs (Lemme \ref{algebra-lemma-length-additive}).
\end{proof}

\begin{lemma}
\label{algebra-lemma-pullback-module}
Soit $A \to B$ un homomorphisme local plat d'anneaux locaux.
Alors, pour tout $A$-module $M$, nous avons
$$
\text{length}_A(M) \text{length}_B(B/\mathfrak m_AB)
=
\text{length}_B(M \otimes_A B).
$$
En particulier, si $\text{length}_B(B/\mathfrak m_AB) < \infty$,
alors $M$ est de longueur finie si et seulement si $M \otimes_A B$ est de longueur finie.
\end{lemma}

\begin{proof}
Le morphisme d'anneaux $A \to B$ est fidèlement plat par le
Lemme \ref{algebra-lemma-local-flat-ff}.
Ainsi, si $0 = M_0 \subset M_1 \subset \ldots \subset M_n = M$
est une chaîne de longueur $n$ dans $M$, la chaîne correspondante
$0 = M_0 \otimes_A B \subset M_1 \otimes_A B \subset
\ldots \subset M_n \otimes_A B = M \otimes_A B$ est aussi de longueur $n$.
Cela démontre
$\text{length}_A(M) = \infty \Rightarrow
\text{length}_B(M \otimes_A B) = \infty$.
Supposons maintenant $\text{length}_A(M) < \infty$. Dans ce cas, $M$ possède une filtration de longueur $\ell = \text{length}_A(M)$
dont les quotients sont $A/\mathfrak m_A$. En raisonnant comme ci-dessus,
nous voyons que $M \otimes_A B$ possède une filtration de longueur $\ell$
dont les quotients sont isomorphes à
$B \otimes_A A/\mathfrak m_A =  B/\mathfrak m_AB$.
Le lemme en découle.
\end{proof}

\begin{lemma}
\label{algebra-lemma-pullback-transitive}
Soient $A \to B \to C$ des homomorphismes locaux plats d'anneaux locaux. Alors
$$
\text{length}_B(B/\mathfrak m_A B)
\text{length}_C(C/\mathfrak m_B C)
=
\text{length}_C(C/\mathfrak m_A C)
$$
\end{lemma}

\begin{proof}
Cela résulte du Lemme \ref{algebra-lemma-pullback-module} appliqué au morphisme d'anneaux
$B \to C$ et au $B$-module $M = B/\mathfrak m_A B$.
\end{proof}


\section{Anneaux artiniens}
\label{algebra-section-artinian}

\noindent
Les anneaux artiniens, et en particulier les anneaux artiniens locaux,
jouent un rôle important en géométrie algébrique, par exemple
dans la théorie des déformations.

\begin{definition}
\label{algebra-definition-artinian}
Un anneau $R$ est {\it artinien} s'il satisfait la
condition de chaîne descendante pour les idéaux.
\end{definition}

\begin{lemma}
\label{algebra-lemma-finite-dimensional-algebra}
Supposons que $R$ soit une algèbre de dimension finie sur un corps.
Alors $R$ est artinien.
\end{lemma}

\begin{proof}
La condition de chaîne descendante pour les idéaux est évidente.
\end{proof}

\begin{lemma}
\label{algebra-lemma-artinian-finite-nr-max}
Si $R$ est artinien, alors $R$ ne possède qu'un nombre fini d'idéaux maximaux.
\end{lemma}

\begin{proof}
Supposons que $\mathfrak m_i$, $i = 1, 2, 3, \ldots$, soient
des idéaux maximaux deux à deux distincts.
Alors $\mathfrak m_1 \supset \mathfrak m_1\cap \mathfrak m_2
\supset \mathfrak m_1 \cap \mathfrak m_2 \cap \mathfrak m_3 \supset \ldots$
est une suite descendante infinie (car, par le
théorème des restes chinois, toutes les applications
$R \to \oplus_{i = 1}^n R/\mathfrak m_i$ sont surjectives).
\end{proof}

\begin{lemma}
\label{algebra-lemma-artinian-radical-nilpotent}
Soit $R$ artinien. Le radical de Jacobson de $R$ est un idéal nilpotent.
\end{lemma}

\begin{proof}
Soit $I \subset R$ le radical de Jacobson.
Remarquons que $I \supset I^2 \supset I^3 \supset \ldots$ est une suite
descendante. Ainsi $I^n = I^{n + 1}$ pour un certain $n$.
Posons $J = \{ x\in R \mid xI^n = 0\}$. Nous devons montrer que $J = R$.
Sinon, choisissons un idéal $J' \not = J$, $J \subset J'$ minimal
(possible par la propriété artinienne). Alors $J'/J$ est un
$R$-module simple, donc isomorphe à $R/\mathfrak m$ pour un certain idéal maximal $\mathfrak m$,
voir le Lemme \ref{algebra-lemma-characterize-length-1}.
Alors $\mathfrak m I^n$ annule $J'$. Comme $I \subset \mathfrak m$,
nous concluons que $I^{n + 1} = I^n$ annule $J'$. Donc $J' = J$,
ce qui est une contradiction.
\end{proof}

\begin{lemma}
\label{algebra-lemma-product-local}
Tout anneau ayant un nombre fini d'idéaux maximaux et
un radical de Jacobson localement nilpotent est le produit de ses localisations
en ses idéaux maximaux. De plus, tous les idéaux premiers sont maximaux.
\end{lemma}

\begin{proof}
Soit $R$ un anneau ayant un nombre fini d'idéaux maximaux
$\mathfrak m_1, \ldots, \mathfrak m_n$.
Soit $I = \bigcap_{i = 1}^n \mathfrak m_i$
le radical de Jacobson de $R$. Supposons que $I$ soit localement nilpotent.
Soit $\mathfrak p$ un idéal premier de $R$.
Comme tout idéal premier contient tout élément
nilpotent de $R$, nous voyons
$ \mathfrak p \supset \mathfrak m_1 \cap \ldots \cap \mathfrak m_n$.
Comme $\mathfrak m_1 \cap \ldots \cap \mathfrak m_n \supset
\mathfrak m_1 \ldots \mathfrak m_n$,
nous concluons $\mathfrak p \supset \mathfrak m_1 \ldots \mathfrak m_n$.
Donc $\mathfrak p \supset \mathfrak m_i$ pour un certain $i$, et ainsi
$\mathfrak p = \mathfrak m_i$. Le spectre de $R$
est donc l'espace topologique discret $\{\mathfrak m_1, \ldots, \mathfrak m_n\}$.
Par le Lemme \ref{algebra-lemma-disjoint-implies-product}, appliqué $n - 1$ fois,
nous trouvons que $R = R_1 \times \ldots \times R_n$, où le
spectre de chaque $R_i$ est un singleton pour chaque $i$.
Ainsi $R_i$ est un anneau local et, comme il est une localisation de $R$
(par le lemme), il est l'un des anneaux locaux de $R$, comme voulu.
\end{proof}

\begin{lemma}
\label{algebra-lemma-artinian-finite-length}
Un anneau $R$ est artinien si et seulement s'il est de longueur finie
comme module sur lui-même. Un tel anneau $R$ est à la fois artinien et
noethérien, tout idéal premier de $R$ est maximal, et $R$ est égal
au produit (fini) de ses localisations en ses idéaux maximaux.
\end{lemma}

\begin{proof}
Si $R$ est de longueur finie sur lui-même, il satisfait à la fois
la condition de chaîne ascendante et la condition de chaîne descendante
pour les idéaux. Il est donc à la fois noethérien et artinien.
Tout anneau artinien est égal au produit de ses localisations
en ses idéaux maximaux par les Lemmes \ref{algebra-lemma-artinian-finite-nr-max},
\ref{algebra-lemma-artinian-radical-nilpotent} et \ref{algebra-lemma-product-local}.

\medskip\noindent
Supposons que $R$ soit artinien. Nous allons montrer que $R$ est de longueur finie
sur lui-même. Il suffit d'exhiber une chaîne de
sous-modules dont les quotients successifs sont de longueur finie.
D'après ce qui précède, nous pouvons supposer que $R$ est local, d'idéal maximal $\mathfrak m$.
Par le Lemme \ref{algebra-lemma-artinian-radical-nilpotent}, nous avons
$\mathfrak m^n =0$ pour un certain $n$.
Considérons la suite
$0 = \mathfrak m^n \subset \mathfrak m^{n-1} \subset
\ldots \subset \mathfrak m \subset R$. Par le Lemme
\ref{algebra-lemma-dimension-is-length}, la longueur de chaque quotient successif
$\mathfrak m^j/\mathfrak m^{j + 1}$ est sa dimension
comme espace vectoriel sur $\kappa(\mathfrak m)$. Celle-ci doit être
finie, car sinon nous aurions une chaîne descendante infinie
de sous-espaces vectoriels, qui correspondrait à une chaîne descendante infinie d'idéaux dans $R$.
\end{proof}

\section{Homomorphismes essentiellement de type fini}
\label{algebra-section-essentially-of-finite-type}

\noindent
Quelques remarques simples sur les localisations de morphismes d'anneaux de type fini.

\begin{definition}
\label{algebra-definition-essentially-finite-p-t}
Soit $R \to S$ un morphisme d'anneaux.
\begin{enumerate}
\item Nous disons que $R \to S$ est {\it essentiellement de type fini} si
$S$ est la localisation d'une $R$-algèbre de type fini.
\item Nous disons que $R \to S$ est {\it essentiellement de présentation finie} si
$S$ est la localisation d'une $R$-algèbre de présentation finie.
\end{enumerate}
\end{definition}

\begin{lemma}
\label{algebra-lemma-composition-essentially-of-finite-type}
La classe des morphismes d'anneaux qui sont essentiellement de type fini
est stable par composition. De même pour les morphismes essentiellement de
présentation finie.
\end{lemma}

\begin{proof}
Omis.
\end{proof}

\begin{lemma}
\label{algebra-lemma-base-change-essentially-of-finite-type}
La classe des morphismes d'anneaux qui sont essentiellement de type fini
est stable par changement de base. De même pour les morphismes essentiellement de
présentation finie.
\end{lemma}

\begin{proof}
Omis.
\end{proof}

\begin{lemma}
\label{algebra-lemma-essentially-of-finite-type-into-artinian-local}
Soit $R \to S$ un morphisme d'anneaux. Supposons que $S$ soit un anneau local
artinien d'idéal maximal $\mathfrak m$. Alors
\begin{enumerate}
\item $R \to S$ est fini si et seulement si $R \to S/\mathfrak m$ est fini,
\item $R \to S$ est de type fini si et seulement si $R \to S/\mathfrak m$
est de type fini.
\item $R \to S$ est essentiellement de type fini si et seulement si la
composition $R \to S/\mathfrak m$ est essentiellement de type fini.
\end{enumerate}
\end{lemma}

\begin{proof}
Si $R \to S$ est fini, alors $R \to S/\mathfrak m$
est fini par le Lemme \ref{algebra-lemma-finite-transitive}.
Réciproquement, supposons que $R \to S/\mathfrak m$ soit fini.
Comme $S$ est de longueur finie sur lui-même
(Lemme \ref{algebra-lemma-artinian-finite-length}), nous pouvons choisir une filtration
$$
0 \subset I_1 \subset \ldots \subset I_n = S
$$
par des idéaux telle que $I_i/I_{i - 1} \cong S/\mathfrak m$ comme
$S$-modules. Ainsi, $S$ possède une filtration par des $R$-sous-modules
$I_i$ dont chaque quotient successif est un $R$-module fini. Donc $S$ est
un $R$-module fini par le Lemme \ref{algebra-lemma-extension}.

\medskip\noindent
Si $R \to S$ est de type fini, alors $R \to S/\mathfrak m$
est de type fini par le Lemme \ref{algebra-lemma-compose-finite-type}.
Réciproquement, supposons que $R \to S/\mathfrak m$ soit de type fini.
Choisissons $f_1, \ldots, f_n \in S$ qui se projettent sur des générateurs
de $S/\mathfrak m$. Alors $A = R[x_1, \ldots, x_n] \to S$, $x_i \mapsto f_i$
est un morphisme d'anneaux tel que $A \to S/\mathfrak m$ soit surjectif
(en particulier fini). Ainsi $A \to S$ est fini par (1), et nous voyons que
$R \to S$ est de type fini par le Lemme \ref{algebra-lemma-compose-finite-type}.

\medskip\noindent
Si $R \to S$ est essentiellement de type fini, alors $R \to S/\mathfrak m$
est essentiellement de type fini par le
Lemme \ref{algebra-lemma-composition-essentially-of-finite-type}.
Réciproquement, supposons que $R \to S/\mathfrak m$ soit essentiellement
de type fini. Supposons que $S/\mathfrak m$ soit la localisation
de $R[x_1, \ldots, x_n]/I$. Choisissons $f_1, \ldots, f_n \in S$
dont les classes modulo $\mathfrak m$ correspondent aux classes de
$x_1, \ldots, x_n$ modulo $I$.
Considérons le morphisme $R[x_1, \ldots, x_n] \to S$, $x_i \mapsto f_i$
de noyau $J$. Posons $A = R[x_1, \ldots, x_n]/J \subset S$
et $\mathfrak p = A \cap \mathfrak m$. Remarquons que
$A/\mathfrak p \subset S/\mathfrak m$ est égal à l'image
de $R[x_1, \ldots, x_n]/I$ dans $S/\mathfrak m$. Ainsi
$\kappa(\mathfrak p) = S/\mathfrak m$. Donc $A_\mathfrak p \to S$
est fini par (1). Nous concluons que $S$ est essentiellement de type
fini par le Lemme \ref{algebra-lemma-composition-essentially-of-finite-type}.
\end{proof}

\noindent
Le lemme suivant peut être démontré en utilisant la propreté de l'espace
projectif au lieu de l'argument algébrique que nous donnons ici.

\begin{lemma}
\label{algebra-lemma-localization-at-closed-point-special-fibre}
Soit $\varphi : R \to S$ essentiellement de type fini, avec $R$ et $S$
locaux (mais $\varphi$ n'est pas nécessairement local). Alors il existe
un $n$ et un idéal maximal $\mathfrak m \subset R[x_1, \ldots, x_n]$
au-dessus de $\mathfrak m_R$ tels que $S$ soit une localisation d'un
quotient de $R[x_1, \ldots, x_n]_\mathfrak m$.
\end{lemma}

\begin{proof}
Nous pouvons écrire $S$ comme une localisation d'un quotient de
$R[x_1, \ldots, x_n]$. Il suffit donc de démontrer le lemme dans le cas où
$S = R[x_1, \ldots, x_n]_\mathfrak q$ pour un idéal premier
$\mathfrak q \subset R[x_1, \ldots, x_n]$.
Si $\mathfrak q + \mathfrak m_R R[x_1, \ldots, x_n] \not = R[x_1, \ldots, x_n]$
alors nous pouvons trouver un idéal maximal $\mathfrak m$ comme dans l'énoncé
du lemme avec $\mathfrak q \subset \mathfrak m$, et le résultat est clair.

\medskip\noindent
Choisissons un anneau de valuation $A \subset \kappa(\mathfrak q)$
qui domine l'image de $R \to \kappa(\mathfrak q)$
(Lemme \ref{algebra-lemma-dominate}). Si l'image $\lambda_i \in \kappa(\mathfrak q)$
de $x_i$ est contenue dans $A$, alors $\mathfrak q$ est contenue dans
l'image réciproque de $\mathfrak m_A$ par $R[x_1, \ldots, x_n] \to A$,
ce qui nous ramène au cas précédent.
Il existe donc un $i$ tel que $\lambda_i^{-1} \in A$ et tel que
$\lambda_j/\lambda_i \in A$ pour tout $j = 1, \ldots, n$
(car le groupe des valeurs de $A$ est totalement ordonné, voir le
Lemme \ref{algebra-lemma-valuation-group}). Nous considérons alors le morphisme
$$
R[y_0, y_1, \ldots, \hat{y_i}, \ldots, y_n]
\to R[x_1, \ldots, x_n]_\mathfrak q,\quad
y_0 \mapsto 1/x_i,\quad y_j \mapsto x_j/x_i
$$
Soit $\mathfrak q' \subset R[y_0, \ldots, \hat{y_i}, \ldots, y_n]$
l'image réciproque de $\mathfrak q$. Comme $y_0 \not \in \mathfrak q'$, il est
facile de voir que la flèche affichée définit un isomorphisme entre les
localisations. D'autre part, le résultat du premier paragraphe s'applique à
$R[y_0, \ldots, \hat{y_i}, \ldots, y_n]$
car $y_j$ s'envoie sur un élément de $A$.
Ceci achève la démonstration.
\end{proof}

\section{Groupes K}
\label{algebra-section-K-groups}

\noindent
Soit $R$ un anneau. Nous allons introduire deux groupes abéliens associés
à $R$. Le premier est noté $K'_0(R)$ et possède les propriétés suivantes
\footnote{La définition a un sens pour tout anneau, mais elle est rarement
utilisée lorsque $R$ n'est pas noethérien.} :
\begin{enumerate}
\item Pour tout $R$-module fini $M$, un élément $[M]$ de
$K'_0(R)$ est donné,
\item pour toute suite exacte courte $0 \to M' \to M \to M'' \to 0$
de $R$-modules finis, nous avons la relation $[M] = [M'] + [M'']$,
\item le groupe $K'_0(R)$ est engendré par les éléments $[M]$, et
\item toutes les relations dans $K'_0(R)$ entre les générateurs $[M]$
sont des combinaisons $\mathbf{Z}$-linéaires
des relations provenant des suites exactes comme ci-dessus.
\end{enumerate}
La construction effective est un peu plus délicate, car il faut prendre garde
à ce que la collection de tous les $R$-modules de type fini soit une classe
propre. Cependant, on peut surmonter ce problème en prenant comme ensemble
de générateurs du groupe $K'_0(R)$ les éléments $[R^n/K]$, où $n$ parcourt
les entiers et où $K$ parcourt tous les sous-modules $K \subset R^n$.
Les générateurs du sous-groupe des relations imposées sur ces éléments
seront les relations provenant des suites exactes courtes dont les termes
sont de la forme $R^n/K$. L'élément $[M]$ est défini en
choisissant $n$ et $K$ tels que $M \cong R^n/K$ et en posant
$[M] = [R^n/K]$. Les détails sont laissés au lecteur.

\begin{lemma}
\label{algebra-lemma-length-K0}
Si $R$ est un anneau local artinien, la fonction longueur
définit un homomorphisme naturel de groupes abéliens
$\text{length}_R : K'_0(R) \to \mathbf{Z}$.
\end{lemma}

\begin{proof}
La longueur de tout $R$-module fini est finie,
car il est quotient de $R^n$, qui est de longueur finie par
le Lemme \ref{algebra-lemma-artinian-finite-length}. Et la fonction longueur
est additive, voir le Lemme \ref{algebra-lemma-length-additive}.
\end{proof}

\noindent
Le second est noté $K_0(R)$ et possède les propriétés suivantes :
\begin{enumerate}
\item Pour tout $R$-module projectif fini $M$, un élément $[M]$ de
$K_0(R)$ est donné,
\item pour toute suite exacte courte $0 \to M' \to M \to M'' \to 0$
de $R$-modules projectifs finis, nous avons la relation $[M] = [M'] + [M'']$,
\item le groupe $K_0(R)$ est engendré par les éléments $[M]$, et
\item toutes les relations dans $K_0(R)$ sont des combinaisons
$\mathbf{Z}$-linéaires des relations provenant des suites exactes comme ci-dessus.
\end{enumerate}
La construction de ce groupe se fait comme ci-dessus.

\medskip\noindent
Nous notons qu'il existe un morphisme évident $K_0(R) \to K'_0(R)$
qui n'est pas un isomorphisme en général.

\begin{example}
\label{algebra-example-K0-field}
Notons que si $R = k$ est un corps, alors nous avons clairement
$K_0(k) = K'_0(k) \cong \mathbf{Z}$, l'isomorphisme
étant donné par la fonction dimension (qui est aussi la fonction longueur).
\end{example}

\begin{example}
\label{algebra-example-K0-PID}
Soit $R$ un anneau principal. Nous affirmons que $K_0(R) = K'_0(R) = \mathbf{Z}$.
En effet, tout $R$-module projectif fini est libre fini.
Un module libre fini possède un rang bien défini par le Lemme \ref{algebra-lemma-rank}.
Étant donnée une suite exacte courte de modules libres finis
$$
0 \to M' \to M \to M'' \to 0
$$
nous avons $\text{rank}(M) = \text{rank}(M') + \text{rank}(M'')$
car dans ce cas $M \cong M' \oplus M'$ (par exemple,
la suite se scinde par le Lemme \ref{algebra-lemma-lift-map}).
Nous en concluons que $K_0(R) = \mathbf{Z}$.

\medskip\noindent
Le théorème de structure des modules sur un anneau principal dit que
tout $R$-module de type fini est de la forme
$M = R^{\oplus r} \oplus R/(d_1) \oplus \ldots \oplus R/(d_k)$.
Considérons la suite exacte courte
$$
0 \to (d_i) \to R \to R/(d_i) \to 0
$$
Comme l'idéal $(d_i)$ est isomorphe à $R$ comme module
(il est libre de générateur $d_i$), dans $K'_0(R)$ nous avons
$[(d_i)] = [R]$. Alors $[R/(d_i)] = [(d_i)]-[R] = 0$. Il
s'ensuit qu'un module de torsion a une classe nulle dans $K'_0(R)$.
En utilisant le rang de la partie libre, on obtient une identification
$K'_0(R) = \mathbf{Z}$ et l'homomorphisme canonique
$K_0(R) \to K'_0(R)$ est un isomorphisme.
\end{example}

\begin{example}
\label{algebra-example-K0-polynomial-ring}
Soit $k$ un corps. Alors $K_0(k[x]) = K'_0(k[x]) = \mathbf{Z}$.
Ceci découle de l'Exemple \ref{algebra-example-K0-PID}, puisque
$R = k[x]$ est un anneau principal.
\end{example}

\begin{example}
\label{algebra-example-K0-node}
Soit $k$ un corps. Soit $R = \{f \in k[x] \mid f(0) = f(1)\}$,
voir l'Exemple \ref{algebra-example-affine-open-not-standard}.
Dans ce cas, $K_0(R) \cong k^* \oplus \mathbf{Z}$, mais
$K'_0(R) = \mathbf{Z}$.
\end{example}

\begin{lemma}
\label{algebra-lemma-K0-product}
Soit $R = R_1 \times R_2$. Alors $K_0(R) = K_0(R_1) \times K_0(R_2)$
et $K'_0(R) = K'_0(R_1) \times K'_0(R_2)$.
\end{lemma}

\begin{proof}
Omis.
\end{proof}

\begin{lemma}
\label{algebra-lemma-K0prime-Artinian}
Soit $R$ un anneau local artinien.
L'application $\text{length}_R : K'_0(R) \to \mathbf{Z}$
du Lemme \ref{algebra-lemma-length-K0} est un isomorphisme.
\end{lemma}

\begin{proof}
Omis.
\end{proof}

\begin{lemma}
\label{algebra-lemma-K0-local}
Soit $(R, \mathfrak m)$ un anneau local. Tout $R$-module projectif fini
est libre fini. L'application $\text{rank}_R : K_0(R) \to \mathbf{Z}$
donnée par $[M] \to \text{rank}_R(M)$ est bien définie
et est un isomorphisme.
\end{lemma}

\begin{proof}
Soit $P$ un $R$-module projectif fini. Choisissons des éléments
$x_1, \ldots, x_n \in P$ qui se projettent sur une base de $P/\mathfrak m P$.
Par le Lemme de Nakayama \ref{algebra-lemma-NAK}, ces éléments engendrent $P$.
La surjection correspondante $u : R^{\oplus n} \to P$ admet une section
car $P$ est projectif. Donc $R^{\oplus n} = P \oplus Q$ avec $Q = \Ker(u)$.
Il s'ensuit que $Q/\mathfrak m Q = 0$, donc $Q$ est nul par
le lemme de Nakayama. Nous voyons ainsi que tout $R$-module projectif
fini est libre fini.
Un module libre fini possède un rang bien défini par le Lemme \ref{algebra-lemma-rank}.
Étant donnée une suite exacte courte de $R$-modules libres finis
$$
0 \to M' \to M \to M'' \to 0
$$
nous avons $\text{rank}(M) = \text{rank}(M') + \text{rank}(M'')$
car dans ce cas $M \cong M' \oplus M'$ (par exemple,
la suite se scinde par le Lemme \ref{algebra-lemma-lift-map}).
Nous en concluons que $K_0(R) = \mathbf{Z}$.
\end{proof}

\begin{lemma}
\label{algebra-lemma-K0-and-K0prime-Artinian-local}
Soit $R$ un anneau artinien local. Il existe un diagramme
commutatif
$$
\xymatrix{
K_0(R) \ar[rr] \ar[d]_{\text{rank}_R} & &
K'_0(R) \ar[d]^{\text{length}_R} \\
\mathbf{Z} \ar[rr]^{\text{length}_R(R)} & & 
\mathbf{Z}
}
$$
où les applications verticales sont des isomorphismes par
les Lemmes \ref{algebra-lemma-K0prime-Artinian} et \ref{algebra-lemma-K0-local}.
\end{lemma}

\begin{proof}
Soit $P$ un $R$-module projectif fini. Nous devons montrer que
$\text{length}_R(P) = \text{rank}_R(P) \text{length}_R(R)$.
Par le Lemme \ref{algebra-lemma-K0-local}, le module $P$ est libre fini. Donc
$P \cong R^{\oplus n}$ pour un certain $n \geq 0$. Alors
$\text{rank}_R(P) = n$ et
$\text{length}_R(R^{\oplus n}) = n \text{length}_R(R)$
par additivité des longueurs (Lemme \ref{algebra-lemma-length-additive}).
Le résultat s'ensuit.
\end{proof}


\section{Anneaux gradués}
\label{algebra-section-graded}

\noindent
Un {\it anneau gradué} sera pour nous un anneau $S$ muni
d'une décomposition en somme directe $S = \bigoplus_{d \geq 0} S_d$
du groupe abélien sous-jacent, telle que $S_d \cdot S_e \subset S_{d + e}$.
Remarquons que nous n'autorisons pas d'éléments non nuls en degrés négatifs.
L'{\it idéal non pertinent} est l'idéal $S_{+} = \bigoplus_{d > 0} S_d$.
Un {\it module gradué}
sera un $S$-module $M$ muni d'une décomposition en somme directe
$M = \bigoplus_{n\in \mathbf{Z}} M_n$ du groupe abélien sous-jacent
telle que $S_d \cdot M_e \subset M_{d + e}$. Remarquons que pour les modules
nous autorisons des éléments non nuls en degrés négatifs.
Nous considérons $S$ comme un $S$-module gradué en posant $S_{-k} = (0)$
pour $k > 0$. Un élément $x$ (resp.\ $f$) de $M$ (resp.\ $S$) est appelé
{\it homogène}
si $x \in M_d$ (resp.\ $f \in S_d$) pour un certain $d$.
Un {\it morphisme de $S$-modules gradués} est un morphisme de $S$-modules
$\varphi : M \to M'$ tel que $\varphi(M_d) \subset M'_d$.
Nous n'autorisons pas les morphismes à décaler les degrés. Notons
$\text{GrHom}_0(M, N)$ le $S_0$-module des homomorphismes
de modules gradués de $M$ vers $N$.

\medskip\noindent
À ce stade interviennent les notions d'idéal gradué,
d'anneau quotient gradué, de sous-module gradué, de module quotient gradué,
de produit tensoriel gradué, etc. Nous laissons au lecteur le soin de trouver
les définitions et les lemmes pertinents. Par exemple : une suite exacte courte
de modules gradués est exacte courte à chaque degré.

\medskip\noindent
Étant donné un anneau gradué $S$, un $S$-module gradué $M$ et $n \in \mathbf{Z}$,
nous notons $M(n)$ le $S$-module gradué tel que $M(n)_d = M_{n + d}$.
C'est ce que l'on appelle le {\it tordu de $M$ par $n$}. En particulier, nous obtenons
les modules $S(n)$, $n \in \mathbf{Z}$, qui joueront un rôle important
dans l'étude des schémas projectifs. Il existe des isomorphismes
fonctoriels évidents tels que
$(M \oplus N)(n) = M(n) \oplus N(n)$,
$(M \otimes_S N)(n) = M \otimes_S N(n) = M(n) \otimes_S N$.
De plus, nous pouvons définir une structure de $S$-module gradué sur
le $S_0$-module
$$
\text{GrHom}(M, N) =
\bigoplus\nolimits_{n \in \mathbf{Z}} \text{GrHom}_n(M, N),
\quad
\text{GrHom}_n(M, N) = \text{GrHom}_0(M, N(n)).
$$
Nous omettons la définition de la multiplication.

\begin{lemma}
\label{algebra-lemma-graded-NAK}
Soit $S$ un anneau gradué. Soit $M$ un $S$-module gradué.
\begin{enumerate}
\item Si $S_+M = M$ et si $M$ est fini, alors $M = 0$.
\item Si $N, N' \subset M$ sont des sous-modules gradués,
$M = N + S_+N'$, et si $N'$ est fini, alors $M = N$.
\item Si $N \to M$ est un morphisme de modules gradués, si $N/S_+N \to M/S_+M$
est surjectif, et si $M$ est fini, alors $N \to M$ est surjectif.
\item Si $x_1, \ldots, x_n \in M$ sont homogènes, engendrent $M/S_+M$
et si $M$ est fini, alors $x_1, \ldots, x_n$ engendrent $M$.
\end{enumerate}
\end{lemma}

\begin{proof}
Preuve de (1). Choisissons des générateurs $y_1, \ldots, y_r$ de $M$ sur $S$.
Nous pouvons supposer que $y_i$ est homogène de degré $d_i$. Après
réindexation, nous pouvons supposer $d_r = \min(d_i)$. Alors la condition
$S_+M = M$ implique $y_r = 0$. Donc $M = 0$ par récurrence sur $r$.
La partie (2) s'obtient en appliquant (1) à $M/N$. La partie (3) s'obtient en
appliquant (2) aux sous-modules $\Im(N \to M)$ et $M$.
La partie (4) s'obtient en appliquant (3) au morphisme de modules
$\bigoplus S(-d_i) \to M$, $(s_1, \ldots, s_n) \mapsto \sum s_i x_i$.
\end{proof}

\noindent
Soit $S$ un anneau gradué. Soit $d \geq 1$ un entier.
Posons $S^{(d)} = \bigoplus_{n \geq 0} S_{nd}$. Nous considérons
$S^{(d)}$ comme un anneau gradué dont la composante de degré $n$ est
$(S^{(d)})_n = S_{nd}$. Étant donné un $S$-module gradué $M$, nous
pouvons de même considérer $M^{(d)} = \bigoplus_{n \in \mathbf{Z}} M_{nd}$,
qui est un $S^{(d)}$-module gradué.

\begin{lemma}
\label{algebra-lemma-uple-generated-degree-1}
Soit $S$ un anneau gradué, qui est de type fini sur $S_0$.
Alors, pour tout $d$ suffisamment divisible, l'algèbre
$S^{(d)}$ est engendrée en degré $1$ sur $S_0$.
\end{lemma}

\begin{proof}
Disons que $S$ est engendré par $f_1, \ldots, f_r \in S$ sur $S_0$.
Après remplacement des $f_i$ par leurs composantes homogènes, nous pouvons
supposer que $f_i$ est homogène de degré $d_i > 0$. Alors tout élément de
$S_n$ est une combinaison linéaire à coefficients dans $S_0$ de monômes
$f_1^{e_1} \ldots f_r^{e_r}$ tels que $\sum e_i d_i = n$.
Soit $m$ un multiple de $\text{lcm}(d_i)$. Pour tout $N \geq r$, si
$$
\sum e_i d_i = N m
$$
alors il existe un $i$ tel que $e_i \geq m/d_i$, par un argument élémentaire.
Ainsi tout monôme de degré $N m$ est le produit d'un monôme
de degré $m$, à savoir $f_i^{m/d_i}$, et d'un monôme de degré $(N - 1)m$.
Il s'ensuit que tout monôme de degré $nrm$, pour $n \geq 2$,
est un produit de monômes de degré $rm$. Ainsi $S^{(rm)}$ est engendrée
en degré $1$ sur $S_0$.
\end{proof}

\begin{lemma}
\label{algebra-lemma-integral-closure-graded}
\begin{reference}
\cite[Theorem 2.3.2]{Huneke-Swanson}
\end{reference}
Soit $R \to S$ un homomorphisme d'anneaux gradués.
Soit $S' \subset S$ la clôture intégrale de $R$ dans $S$.
Alors
$$
S' = \bigoplus\nolimits_{d \geq 0} S' \cap S_d,
$$
c'est-à-dire que $S'$ est une $R$-sous-algèbre graduée de $S$.
\end{lemma}

\begin{proof}
Nous devons montrer ce qui suit : si
$s = s_n + s_{n + 1} + \ldots + s_m \in S'$, alors chaque
composante homogène $s_j \in S'$. Nous le démontrerons par récurrence
sur $m - n$ pour tous les homomorphismes $R \to S$ d'anneaux gradués.
Tout d'abord, il est immédiat que $s_0$ est entier sur $R_0$ (et donc sur $R$),
car il existe un morphisme d'anneaux $S \to S_0$ compatible avec le morphisme
d'anneaux $R \to R_0$. Ainsi, après avoir remplacé $s$ par $s - s_0$, nous
pouvons supposer $n > 0$. Considérons l'extension d'anneaux gradués
$R[t, t^{-1}] \to S[t, t^{-1}]$ où $t$ est de degré $0$. Il existe un diagramme
commutatif
$$
\xymatrix{
S[t, t^{-1}] \ar[rr]_{s \mapsto t^{\deg(s)}s} & & S[t, t^{-1}] \\
R[t, t^{-1}] \ar[u] \ar[rr]^{r \mapsto t^{\deg(r)}r} & &  R[t, t^{-1}] \ar[u]
}
$$
où les morphismes horizontaux sont des automorphismes d'anneaux. Ainsi la clôture
intégrale $C$ de $S[t, t^{-1}]$ sur $R[t, t^{-1}]$ se transforme en elle-même.
Nous voyons donc que
$$
t^m(s_n + s_{n + 1} + \ldots + s_m) -
(t^ns_n + t^{n + 1}s_{n + 1} + \ldots + t^ms_m) \in C
$$
ce qui implique, par l'hypothèse de récurrence, que chaque $(t^m - t^i)s_i \in C$
pour $i = n, \ldots, m - 1$. Remarquons que pour tout anneau $A$ et
$m > i \geq n > 0$, nous avons
$A[t, t^{-1}]/(t^m - t^i - 1) \cong A[t]/(t^m - t^i - 1) \supset A$
car $t(t^{m - 1} - t^{i - 1}) = 1$ dans $A[t]/(t^m - t^i - 1)$.
Comme $t^m - t^i$ s'envoie sur $1$, nous voyons que l'image de $s_i$ dans
l'anneau $S[t]/(t^m - t^i - 1)$ est entière sur
$R[t]/(t^m - t^i - 1)$ pour $i = n, \ldots, m - 1$. Comme
$R \to R[t]/(t^m - t^i - 1)$ est fini, nous voyons que $s_i$ est entier
sur $R$ par transitivité, voir le Lemme \ref{algebra-lemma-integral-transitive}.
Enfin, nous concluons également que $s_m = s - \sum_{i = n, \ldots, m - 1} s_i$
est entier sur $R$.
\end{proof}

\section{Proj d'un anneau gradué}
\label{algebra-section-proj}

\noindent
Soit $S$ un anneau gradué.
Un {\it idéal homogène} est simplement un idéal
$I \subset S$ qui est aussi un sous-module gradué de $S$.
De manière équivalente, c'est un idéal engendré par des éléments homogènes.
De manière équivalente, si $f \in I$ et
$$
f = f_0 + f_1 + \ldots + f_n
$$
est la décomposition de $f$ en composantes homogènes dans $S$, alors $f_i \in I$
pour tout $i$. Pour vérifier qu'un idéal homogène $\mathfrak p$
est premier, il suffit de vérifier que si $ab \in \mathfrak p$
avec $a, b$ homogènes, alors soit $a \in \mathfrak p$, soit
$b \in \mathfrak p$.

\begin{definition}
\label{algebra-definition-proj}
Soit $S$ un anneau gradué.
Nous définissons $\text{Proj}(S)$ comme l'ensemble des idéaux
premiers homogènes $\mathfrak p$ de $S$ tels que
$S_{+} \not \subset \mathfrak p$.
L'ensemble $\text{Proj}(S)$ est un sous-ensemble de $\Spec(S)$
et nous le munissons de la topologie induite.
L'espace topologique $\text{Proj}(S)$ est appelé le
{\it spectre homogène} de l'anneau gradué $S$.
\end{definition}

\noindent
Remarquons que, par construction, il existe un morphisme continu
$$
\text{Proj}(S) \longrightarrow \Spec(S_0).
$$

\medskip\noindent
Soit $S = \oplus_{d \geq 0} S_d$ un anneau gradué.
Soit $f\in S_d$ et supposons que $d \geq 1$.
Nous définissons $S_{(f)}$ comme le sous-anneau de $S_f$
constitué des éléments de la forme $r/f^n$, avec $r$ homogène et
$\deg(r) = nd$. Si $M$ est un $S$-module gradué,
nous définissons le $S_{(f)}$-module $M_{(f)}$ comme le
sous-module de $M_f$ constitué des éléments de
la forme $x/f^n$, où $x$ est homogène de degré $nd$.

\begin{lemma}
\label{algebra-lemma-Z-graded}
Soit $S$ un anneau gradué par $\mathbf{Z}$ contenant un
élément inversible homogène de degré strictement positif. Alors l'ensemble
$G \subset \Spec(S)$ des idéaux premiers gradués par $\mathbf{Z}$ de $S$
(muni de la topologie induite) s'envoie homéomorphiquement sur $\Spec(S_0)$.
\end{lemma}

\begin{proof}
Montrons d'abord que le morphisme est une bijection en construisant un inverse.
Soit $f \in S_d$, $d > 0$, inversible dans $S$.
Si $\mathfrak p_0$ est un idéal premier de $S_0$, alors
$\mathfrak p_0S$ est un idéal gradué par $\mathbf{Z}$ de $S$ tel que
$\mathfrak p_0S \cap S_0 = \mathfrak p_0$. Et si $ab \in \mathfrak p_0S$
avec $a$, $b$ homogènes, alors
$a^db^d/f^{\deg(a) + \deg(b)} \in \mathfrak p_0$.
Ainsi, soit $a^d/f^{\deg(a)} \in \mathfrak p_0$, soit
$b^d/f^{\deg(b)} \in \mathfrak p_0$, autrement dit soit
$a^d \in \mathfrak p_0S$, soit $b^d \in \mathfrak p_0S$.
Il s'ensuit que $\sqrt{\mathfrak p_0S}$ est un idéal premier
gradué par $\mathbf{Z}$ de $S$ dont l'intersection avec $S_0$ est $\mathfrak p_0$.

\medskip\noindent
Pour montrer que le morphisme est un homéomorphisme, montrons que
l'image de $G \cap D(g)$ est ouverte. Si $g = \sum g_i$
avec $g_i \in S_i$, alors, d'après ce qui précède, $G \cap D(g)$
s'envoie sur l'ensemble $\bigcup D(g_i^d/f^i)$, qui est ouvert.
\end{proof}

\noindent
Pour $f \in S$ homogène de degré $> 0$, nous définissons
$$
D_{+}(f) = \{ \mathfrak p \in \text{Proj}(S) \mid f \not\in \mathfrak p \}.
$$
Enfin, pour un idéal homogène $I \subset S$, nous définissons
$$
V_{+}(I) = \{ \mathfrak p \in \text{Proj}(S) \mid I \subset \mathfrak p \}.
$$
Nous utiliserons plus généralement la notation $V_{+}(E)$ pour tout
ensemble $E$ d'éléments homogènes $E \subset S$.

\begin{lemma}[Topologie sur Proj]
\label{algebra-lemma-topology-proj}
Soit $S = \oplus_{d \geq 0} S_d$ un anneau gradué.
\begin{enumerate}
\item Les ensembles $D_{+}(f)$ sont ouverts dans $\text{Proj}(S)$.
\item Nous avons $D_{+}(ff') = D_{+}(f) \cap D_{+}(f')$.
\item Soit $g = g_0 + \ldots + g_m$ un élément
de $S$ avec $g_i \in S_i$. Alors
$$
D(g) \cap \text{Proj}(S) =
(D(g_0) \cap \text{Proj}(S))
\cup
\bigcup\nolimits_{i \geq 1} D_{+}(g_i).
$$
\item
Soit $g_0\in S_0$ un élément homogène de degré $0$. Alors
$$
D(g_0) \cap \text{Proj}(S)
=
\bigcup\nolimits_{f \in S_d, \ d\geq 1} D_{+}(g_0 f).
$$
\item Les ouverts $D_{+}(f)$ forment une
base de la topologie de $\text{Proj}(S)$.
\item Soit $f \in S$ homogène de degré positif.
L'anneau $S_f$ possède une graduation naturelle par $\mathbf{Z}$.
Les morphismes d'anneaux $S \to S_f \leftarrow S_{(f)}$ induisent
des homéomorphismes
$$
D_{+}(f)
\leftarrow
\{\mathbf{Z}\text{-idéaux premiers gradués de }S_f\}
\to
\Spec(S_{(f)}).
$$
\item Il existe un $S$ tel que $\text{Proj}(S)$ ne soit pas
quasi-compact.
\item Les ensembles $V_{+}(I)$ sont fermés.
\item Tout fermé $T \subset \text{Proj}(S)$ est de
la forme $V_{+}(I)$ pour un certain idéal homogène $I \subset S$.
\item Pour tout idéal gradué $I \subset S$, nous avons
$V_{+}(I) = \emptyset$ si et seulement si $S_{+} \subset \sqrt{I}$.
\end{enumerate}
\end{lemma}

\begin{proof}
Puisque $D_{+}(f) = \text{Proj}(S) \cap D(f)$, ces ensembles sont ouverts.
Ceci démontre (1). De même, (2) découle de $D(ff') = D(f) \cap D(f')$.
De manière analogue, les ensembles $V_{+}(I) = \text{Proj}(S) \cap V(I)$
sont fermés. Ceci démontre (8).

\medskip\noindent
Supposons que $T \subset \text{Proj}(S)$ soit fermé.
Alors nous pouvons écrire $T = \text{Proj}(S) \cap V(J)$ pour un
idéal $J \subset S$. Par définition d'un idéal homogène,
si $g \in J$, $g = g_0 + \ldots + g_m$
avec $g_d \in S_d$, alors $g_d \in \mathfrak p$ pour tout
$\mathfrak p \in T$. Ainsi, si $I \subset S$
est l'idéal engendré par les composantes homogènes des éléments
de $J$, nous avons $T = V_{+}(I)$. Ceci démontre (9).

\medskip\noindent
La formule pour $\text{Proj}(S) \cap D(g)$, avec $g \in S$, découle directement
des définitions. Ceci démontre (3).
Considérons la formule pour $\text{Proj}(S) \cap D(g_0)$.
L'inclusion du membre de droite dans le membre de gauche est
évidente. Pour l'autre inclusion, supposons $g_0 \not \in \mathfrak p$
avec $\mathfrak p \in \text{Proj}(S)$. Si $g_0f \in \mathfrak p$
pour tout $f$ homogène de degré positif, alors nous voyons que
$S_{+} \subset \mathfrak p$, ce qui est une contradiction. Cela donne
l'autre inclusion. Ceci démontre (4).

\medskip\noindent
La collection des ouverts $D(g) \cap \text{Proj}(S)$
forme une base de la topologie, puisque les ouverts standards
$D(g) \subset \Spec(S)$ forment une base de la topologie de
$\Spec(S)$. D'après les formules ci-dessus, nous pouvons exprimer
$D(g) \cap \text{Proj}(S)$ comme une réunion d'ouverts $D_{+}(f)$.
Ainsi la collection des ouverts $D_{+}(f)$ forme aussi une base de la topologie.
Ceci démontre (5).

\medskip\noindent
Preuve de (6). Notons d'abord que $D_{+}(f)$ peut être identifié
à un sous-ensemble (muni de la topologie induite) de $D(f) = \Spec(S_f)$
via le Lemme \ref{algebra-lemma-standard-open}. Remarquons que l'anneau
$S_f$ est gradué par $\mathbf{Z}$. Les éléments homogènes sont
de la forme $r/f^n$ avec $r \in S$ homogène et ont
pour degré $\deg(r/f^n) = \deg(r) - n\deg(f)$. Le sous-ensemble
$D_{+}(f)$ correspond exactement aux idéaux premiers
$\mathfrak p \subset S_f$ qui sont des idéaux gradués par $\mathbf{Z}$
(c'est-à-dire engendrés par des éléments homogènes). Il suffit donc de montrer que
l'ensemble des idéaux premiers gradués par $\mathbf{Z}$ de $S_f$ s'envoie
homéomorphiquement sur $\Spec(S_{(f)})$. Cela découle du Lemme \ref{algebra-lemma-Z-graded}.

\medskip\noindent
Soit $S = \mathbf{Z}[X_1, X_2, X_3, \ldots]$ avec une graduation telle que
chaque $X_i$ soit de degré $1$. Alors il est facile de voir que
$$
\text{Proj}(S) = \bigcup\nolimits_{i = 1}^\infty D_{+}(X_i)
$$
n'admet pas de sous-recouvrement fini. Ceci démontre (7).

\medskip\noindent
Soit $I \subset S$ un idéal gradué.
Si $\sqrt{I} \supset S_{+}$, alors $V_{+}(I) = \emptyset$, puisque
tout idéal premier $\mathfrak p \in \text{Proj}(S)$ ne contient pas
$S_{+}$ par définition. Réciproquement, supposons que
$S_{+} \not \subset \sqrt{I}$. Nous pouvons trouver un élément
$f \in S_{+}$ tel que $f$ ne soit pas nilpotent modulo $I$.
Cela signifie clairement que l'une des composantes homogènes de $f$
n'est pas nilpotente modulo $I$; autrement dit, nous pouvons (et allons)
supposer que $f$ est homogène. Cela implique que
$I S_f \not = S_f$, autrement dit que $(S/I)_f$ n'est pas
nul. Donc $(S/I)_{(f)} \not = 0$, puisque c'est un anneau qui
se plonge dans $(S/I)_f$. Choisissons un idéal premier
$\mathfrak q \subset (S/I)_{(f)}$. Celui-ci correspond à
un idéal premier gradué de $S/I$, qui ne contient pas l'idéal non pertinent
$(S/I)_{+}$. Cela correspond à son tour à un idéal premier
gradué $\mathfrak p$ de $S$, contenant $I$ mais ne contenant pas $S_{+}$,
comme voulu. Ceci démontre (10) et achève la démonstration.
\end{proof}

\begin{example}
\label{algebra-example-proj-polynomial-ring-1-variable}
Soit $R$ un anneau. Si $S = R[X]$ avec $\deg(X) = 1$, alors le morphisme naturel
$\text{Proj}(S) \to \Spec(R)$ est une bijection et même un homéomorphisme.
En effet, soit $\mathfrak p \in \text{Proj}(S)$. Puisque
$S_{+} \not \subset \mathfrak p$, nous voyons que $X \not \in \mathfrak p$.
Ainsi, si $aX^n \in \mathfrak p$ avec $a \in R$ et $n > 0$, alors
$a \in \mathfrak p$. Il s'ensuit que $\mathfrak p = \mathfrak p_0S$
avec $\mathfrak p_0 = \mathfrak p \cap R$.
\end{example}

\noindent
Si $\mathfrak p \in \text{Proj}(S)$, nous
définissons $S_{(\mathfrak p)}$ comme l'anneau dont les
éléments sont les fractions $r/f$, où $r, f \in S$ sont des
éléments homogènes de même degré tels que $f \not\in \mathfrak p$.
Comme d'habitude, nous disons que $r/f = r'/f'$ si et seulement s'il existe
un certain $f'' \in S$ homogène, $f'' \not \in \mathfrak p$, tel
que $f''(rf' - r'f) = 0$.
Étant donné un $S$-module gradué $M$, nous posons
$M_{(\mathfrak p)}$ égal au $S_{(\mathfrak p)}$-module
dont les éléments sont les fractions $x/f$, où $x \in M$
et $f \in S$ sont homogènes de même degré, avec
$f \not \in \mathfrak p$. Nous disons que $x/f = x'/f'$
si et seulement s'il existe un certain $f'' \in S$ homogène,
$f'' \not \in \mathfrak p$, tel que $f''(xf' - x'f) = 0$.

\begin{lemma}
\label{algebra-lemma-proj-prime}
Soit $S$ un anneau gradué. Soit $M$ un $S$-module gradué.
Soit $\mathfrak p$ un élément de $\text{Proj}(S)$.
Soit $f \in S$ un élément homogène de degré positif
tel que $f \not \in \mathfrak p$, c'est-à-dire $\mathfrak p \in D_{+}(f)$.
Soit $\mathfrak p' \subset S_{(f)}$ l'élément de
$\Spec(S_{(f)})$ correspondant à $\mathfrak p$ comme dans
le Lemme \ref{algebra-lemma-topology-proj}. Alors
$S_{(\mathfrak p)} = (S_{(f)})_{\mathfrak p'}$
et, de manière compatible,
$M_{(\mathfrak p)} = (M_{(f)})_{\mathfrak p'}$.
\end{lemma}

\begin{proof}
Nous définissons un morphisme $\psi : M_{(\mathfrak p)} \to (M_{(f)})_{\mathfrak p'}$.
Soit $x/g \in M_{(\mathfrak p)}$. Posons
$$
\psi(x/g) = (x g^{\deg(f) - 1}/f^{\deg(x)})/(g^{\deg(f)}/f^{\deg(g)}).
$$
Ceci a un sens puisque $\deg(x) = \deg(g)$ et puisque
$g^{\deg(f)}/f^{\deg(g)} \not \in \mathfrak p'$.
Nous omettons la vérification que $\psi$ est bien défini, est un morphisme
de modules et est un isomorphisme. Indication : l'inverse envoie
$(x/f^n)/(g/f^m)$ sur $(xf^m)/(g f^n)$.
\end{proof}

\noindent
Voici une variante graduée du Lemme \ref{algebra-lemma-silly}.

\begin{lemma}
\label{algebra-lemma-graded-silly}
Supposons que $S$ soit un anneau gradué, que les $\mathfrak p_i$, $i = 1, \ldots, r$,
soient des idéaux premiers homogènes et que $I \subset S_{+}$ soit un idéal gradué.
Supposons que $I \not\subset \mathfrak p_i$ pour tout $i$. Alors il
existe un élément homogène $x\in I$ de degré positif tel que
$x\not\in \mathfrak p_i$ pour tout $i$.
\end{lemma}

\begin{proof}
Nous pouvons supposer qu'il n'y a aucune inclusion entre les $\mathfrak p_i$.
Le résultat est vrai pour $r = 1$. Supposons le résultat vrai pour $r - 1$.
Choisissons $x \in I$ homogène de degré positif tel que
$x \not \in \mathfrak p_i$ pour tous les $i = 1, \ldots, r - 1$.
Si $x \not\in \mathfrak p_r$, c'est terminé. Supposons donc $x \in \mathfrak p_r$.
Si $I \mathfrak p_1 \ldots \mathfrak p_{r-1} \subset \mathfrak p_r$,
alors $I \subset \mathfrak p_r$, contradiction.
Choisissons $y \in I\mathfrak p_1 \ldots \mathfrak p_{r-1}$ homogène
et $y \not \in \mathfrak p_r$. Alors $x^{\deg(y)} + y^{\deg(x)}$ convient.
\end{proof}

\begin{lemma}
\label{algebra-lemma-smear-out}
Soit $S$ un anneau gradué.
Soit $\mathfrak p \subset S$ un idéal premier.
Soit $\mathfrak q$ l'idéal homogène de $S$ engendré par les
éléments homogènes de $\mathfrak p$. Alors $\mathfrak q$ est un
idéal premier de $S$.
\end{lemma}

\begin{proof}
Pour prouver que $\mathfrak q$ est premier, il suffit de vérifier que
si $f, g \in S$ sont homogènes et si $fg \in \mathfrak q$, alors
$f$ ou $g$ appartient à $\mathfrak q$. Alors $fg \in \mathfrak p$
car $\mathfrak p \subset \mathfrak q$. Comme $\mathfrak p$
est premier, nous voyons que soit $f \in \mathfrak p$, soit $g \in \mathfrak p$.
Comme $f$ et $g$ sont homogènes, il est alors clair que soit
$f \in \mathfrak q$, soit $g \in \mathfrak q$.
\end{proof}

\begin{lemma}
\label{algebra-lemma-graded-ring-minimal-prime}
Soit $S$ un anneau gradué.
\begin{enumerate}
\item Tout idéal premier minimal de $S$ est un idéal homogène de $S$.
\item Étant donné un idéal homogène $I \subset S$, tout idéal premier
minimal au-dessus de $I$ est homogène.
\end{enumerate}
\end{lemma}

\begin{proof}
La première assertion est vraie car l'idéal premier $\mathfrak q$ construit dans
le Lemme \ref{algebra-lemma-smear-out} vérifie $\mathfrak q \subset \mathfrak p$.
La seconde découle de ce que nous pouvons considérer $S/I$ et appliquer la première partie.
\end{proof}

\begin{lemma}
\label{algebra-lemma-dehomogenize-finite-type}
Soit $R$ un anneau. Soit $S$ une $R$-algèbre graduée. Soit $f \in S_{+}$
homogène. Supposons que $S$ soit de type fini sur $R$. Alors
\begin{enumerate}
\item l'anneau $S_{(f)}$ est de type fini sur $R$, et
\item pour tout $S$-module gradué fini $M$, le module $M_{(f)}$
est un $S_{(f)}$-module fini.
\end{enumerate}
\end{lemma}

\begin{proof}
Choisissons $f_1, \ldots, f_n \in S$ qui engendrent $S$ comme $R$-algèbre.
Nous pouvons supposer que chaque $f_i$ est homogène (en décomposant chaque $f_i$
en ses composantes homogènes). Un élément de $S_{(f)}$ est une somme
de la forme
$$
\sum\nolimits_{e\deg(f) =
\sum e_i\deg(f_i)} \lambda_{e_1 \ldots e_n} f_1^{e_1} \ldots f_n^{e_n}/f^e
$$
avec $\lambda_{e_1 \ldots e_n} \in R$. Ainsi $S_{(f)}$ est engendré
comme $R$-algèbre par les $f_1^{e_1} \ldots f_n^{e_n} /f^e$ vérifiant
$e\deg(f) = \sum e_i\deg(f_i)$. Si $e_i \geq \deg(f)$,
nous pouvons écrire ceci sous la forme
$$
f_1^{e_1} \ldots f_n^{e_n}/f^e =
f_i^{\deg(f)}/f^{\deg(f_i)} \cdot
f_1^{e_1} \ldots f_i^{e_i - \deg(f)} \ldots f_n^{e_n}/f^{e - \deg(f_i)}
$$
Ainsi, il ne nous faut que les éléments $f_i^{\deg(f)}/f^{\deg(f_i)}$, ainsi
que les éléments $f_1^{e_1} \ldots f_n^{e_n} /f^e$ tels que
$e \deg(f) = \sum e_i \deg(f_i)$ et $e_i < \deg(f)$.
Il s'agit d'une liste finie, et nous voyons que (1) est vraie.

\medskip\noindent
Pour voir (2), supposons que $M$ soit engendré par des éléments homogènes
$x_1, \ldots, x_m$. En raisonnant
comme ci-dessus, nous trouvons que $M_{(f)}$ est engendré comme $S_{(f)}$-module
par la liste finie des éléments de la forme
$f_1^{e_1} \ldots f_n^{e_n} x_j /f^e$
avec $e \deg(f) = \sum e_i \deg(f_i) + \deg(x_j)$ et
$e_i < \deg(f)$.
\end{proof}

\begin{lemma}
\label{algebra-lemma-homogenize}
Soit $R$ un anneau.
Soit $R'$ une $R$-algèbre de type fini, et soit $M$ un $R'$-module fini.
Il existe une $R$-algèbre graduée $S$, un $S$-module gradué $N$ et
un élément $f \in S$ homogène de degré $1$ tels que
\begin{enumerate}
\item $R' \cong S_{(f)}$ et $M \cong N_{(f)}$ (comme modules),
\item $S_0 = R$ et $S$ est engendré par un nombre fini d'éléments
de degré $1$ sur $R$, et
\item $N$ est un $S$-module fini.
\end{enumerate}
\end{lemma}

\begin{proof}
Nous pouvons écrire $R' = R[x_1, \ldots, x_n]/I$ pour un certain idéal $I$.
Pour un élément $g \in R[x_1, \ldots, x_n]$, notons
$\tilde g \in R[X_0, \ldots, X_n]$ l'élément homogène de degré minimal
tel que $g = \tilde g(1, x_1, \ldots, x_n)$.
Soit $\tilde I \subset R[X_0, \ldots, X_n]$ l'idéal engendré par tous les
éléments $\tilde g$, $g \in I$.
Posons $S = R[X_0, \ldots, X_n]/\tilde I$ et notons $f$ l'image
de $X_0$ dans $S$. Par construction, nous avons un isomorphisme
$$
S_{(f)} \longrightarrow R', \quad
X_i/X_0 \longmapsto x_i.
$$
Pour faire de même avec le module $M$, choisissons une présentation
$$
M = (R')^{\oplus r}/\sum\nolimits_{j \in J} R'k_j
$$
avec $k_j = (k_{1j}, \ldots, k_{rj})$. Posons $d_{ij} = \deg(\tilde k_{ij})$.
Posons $d_j = \max\{d_{ij}\}$. Posons $K_{ij} = X_0^{d_j - d_{ij}}\tilde k_{ij}$,
qui est homogène de degré $d_j$. Avec cette notation, posons
$$
N = \Coker\Big(
\bigoplus\nolimits_{j \in J} S(-d_j) \xrightarrow{(K_{ij})} S^{\oplus r}
\Big)
$$
ce qui convient. Certains détails sont omis.
\end{proof}

\section{Anneaux gradués noethériens}
\label{algebra-section-noetherian-graded}

\noindent
Quelques éléments de théorie des anneaux gradués noethériens,
notamment sur les polynômes de Hilbert.

\begin{lemma}
\label{algebra-lemma-S-plus-generated}
Soit $S$ un anneau gradué. Un ensemble d'éléments homogènes
$f_i \in S_{+}$ engendre $S$ comme algèbre sur $S_0$ si et seulement
si ces éléments engendrent $S_{+}$ comme idéal de $S$.
\end{lemma}

\begin{proof}
Si les $f_i$ engendrent $S$ comme algèbre sur $S_0$, tout élément
de $S_{+}$ est un polynôme sans terme constant en les $f_i$, et donc
$S_{+}$ est engendré par les $f_i$ comme idéal. Réciproquement, supposons que
$S_{+} = \sum Sf_i$. Nous montrerons que tout élément $f$ de $S$ s'écrit
comme un polynôme en les $f_i$ à coefficients dans $S_0$. Il suffit
de le faire pour les éléments homogènes. Supposons que $f$ soit de degré $d$.
Nous pouvons alors raisonner par récurrence sur $d$. Le cas $d = 0$ est immédiat.
Si $d > 0$, alors $f \in S_{+}$, donc nous pouvons écrire $f = \sum g_i f_i$
pour certains $g_i \in S$. Comme $S$ est gradué, nous pouvons remplacer
$g_i$ par sa composante homogène de degré $d - \deg(f_i)$. Par récurrence,
nous voyons que chaque $g_i$ est un polynôme en les $f_i$, et le résultat suit.
\end{proof}

\begin{lemma}
\label{algebra-lemma-graded-Noetherian}
Un anneau gradué $S$ est noethérien si et seulement si $S_0$ est
noethérien et si $S_{+}$ est de type fini comme idéal de $S$.
\end{lemma}

\begin{proof}
Il est clair que si $S$ est noethérien, alors $S_0 = S/S_{+}$ est noethérien
et $S_{+}$ est de type fini. Réciproquement, supposons $S_0$ noethérien
et $S_{+}$ de type fini comme idéal de $S$. Choisissons des générateurs
$S_{+} = (f_1, \ldots, f_n)$. En décomposant les $f_i$ en composantes homogènes,
nous pouvons supposer que chaque $f_i$ est homogène. Par le
Lemme \ref{algebra-lemma-S-plus-generated}, nous voyons que
$S_0[X_1, \ldots X_n] \to S$ qui envoie $X_i$ sur $f_i$
est surjectif. Ainsi $S$ est noethérien par le
Lemme \ref{algebra-lemma-Noetherian-permanence}.
\end{proof}

\begin{definition}
\label{algebra-definition-numerical-polynomial}
Soit $A$ un groupe abélien.
Nous disons qu'une fonction $f : n \mapsto f(n) \in A$
définie pour tous les entiers $n$ suffisamment grands est un
{\it polynôme numérique} s'il existe $r \geq 0$,
des éléments $a_0, \ldots, a_r\in A$ tels que
$$
f(n) = \sum\nolimits_{i = 0}^r \binom{n}{i} a_i
$$
pour tout $n \gg 0$.
\end{definition}

\noindent
La raison d'utiliser les coefficients binomiaux est le
fait élémentaire que tout polynôme $P \in \mathbf{Q}[T]$
dont toutes les valeurs aux points entiers sont entières est
égal à une somme $P(T) = \sum a_i \binom{T}{i}$ avec
$a_i \in \mathbf{Z}$. Remarquons qu'en particulier les
expressions $\binom{T + 1}{i + 1}$ sont de cette forme.

\begin{lemma}
\label{algebra-lemma-numerical-polynomial-functorial}
Si $A \to A'$ est un homomorphisme de groupes abéliens et si
$f : n \mapsto f(n) \in A$ est un polynôme numérique,
alors la composée l'est également.
\end{lemma}

\begin{proof}
Ceci est immédiat d'après les définitions.
\end{proof}

\begin{lemma}
\label{algebra-lemma-numerical-polynomial}
Supposons que $f: n \mapsto f(n) \in A$
soit définie pour tout $n$ suffisamment grand,
et supposons que $n \mapsto f(n) - f(n-1)$
soit un polynôme numérique. Alors $f$ est un
polynôme numérique.
\end{lemma}

\begin{proof}
Soit $f(n) - f(n-1) = \sum\nolimits_{i = 0}^r \binom{n}{i} a_i$
pour tout $n \gg 0$. Posons
$g(n) = f(n) - \sum\nolimits_{i = 0}^r \binom{n + 1}{i + 1} a_i$.
Alors $g(n) - g(n-1) = 0$ pour tout $n \gg 0$. Ainsi $g$ est
finalement constante, disons égale à $a_{-1}$. Nous laissons
au lecteur le soin de montrer que
$a_{-1} + \sum\nolimits_{i = 0}^r \binom{n + 1}{i + 1} a_i$
a la forme requise (voir la remarque au-dessus du lemme).
\end{proof}

\begin{lemma}
\label{algebra-lemma-graded-module-fg}
Si $M$ est un $S$-module gradué de type fini,
et si $S$ est de type fini sur $S_0$, alors
chaque $M_n$ est un $S_0$-module fini.
\end{lemma}

\begin{proof}
Supposons que les générateurs de $M$ soient les $m_i$ et les générateurs
de $S$ les $f_i$. En prenant les composantes homogènes, nous pouvons
supposer que les $m_i$ et les $f_i$ sont homogènes
et nous pouvons supposer $f_i \in S_{+}$. Dans ce cas, il est
clair que chaque $M_n$ est engendré sur $S_0$
par les « monômes » $\prod f_i^{e_i} m_j$ dont le
degré est $n$.
\end{proof}

\begin{proposition}
\label{algebra-proposition-graded-hilbert-polynomial}
Supposons que $S$ soit un anneau gradué noethérien
et $M$ un $S$-module gradué fini. Considérons la
fonction
$$
\mathbf{Z} \longrightarrow K'_0(S_0), \quad
n \longmapsto [M_n]
$$
voir le Lemme \ref{algebra-lemma-graded-module-fg}.
Si $S_{+}$ est engendré par des éléments de degré $1$,
alors cette fonction est un polynôme numérique.
\end{proposition}

\begin{proof}
Nous démontrons ceci par récurrence sur le nombre minimal de
générateurs de $S_1$. Si ce nombre est $0$, alors
$M_n = 0$ pour tout $n \gg 0$ et le résultat est vrai.
Pour l'étape de récurrence, soit $x\in S_1$
l'un des générateurs d'un ensemble minimal, tel que
l'hypothèse de récurrence s'applique à l'anneau gradué $S/(x)$.

\medskip\noindent
Montrons d'abord que le résultat est vrai si $x$ est nilpotent sur $M$.
Nous raisonnons par récurrence sur le plus petit entier $r$ tel que
$x^r M  = 0$. Si $r = 1$, alors $M$ est un module sur $S/xS$
et le résultat est vrai (par l'autre hypothèse de récurrence).
Si $r > 1$, nous pouvons trouver une suite exacte courte
$0 \to M' \to M \to M'' \to 0$ telle que les entiers
$r', r''$ soient strictement plus petits que $r$. Nous connaissons donc
le résultat pour $M''$ et $M'$. Ainsi
nous obtenons le résultat pour $M$ en raison de la relation
$
[M_d]  = [M'_d] + [M''_d]
$
dans $K'_0(S_0)$.

\medskip\noindent
Si $x$ n'est pas nilpotent sur $M$, soit $M' \subset M$ le
plus grand sous-module sur lequel $x$ est nilpotent.
Considérons la suite exacte $0 \to M' \to M \to M/M' \to 0$ :
nous voyons à nouveau qu'il suffit de prouver le résultat pour $M/M'$.
Autrement dit, nous pouvons supposer que la multiplication par $x$ est injective.

\medskip\noindent
Posons $\overline{M} = M/xM$. Remarquons que l'application $x : M \to M$
n'est {\it pas} une application de $S$-modules gradués, car elle
n'envoie pas $M_d$ dans $M_d$. En effet, pour chaque $d$, nous avons la
suite exacte courte suivante
$$
0 \to M_d \xrightarrow{x} M_{d + 1} \to \overline{M}_{d + 1} \to 0
$$
Ceci montre que $[M_{d + 1}] - [M_d] = [\overline{M}_{d + 1}]$.
Le résultat découle alors du Lemme \ref{algebra-lemma-numerical-polynomial}.
\end{proof}

\begin{remark}
\label{algebra-remark-period-polynomial}
Si $S$ est toujours noethérien mais $S$ n'est pas engendré en degré $1$,
alors la fonction associée à un $S$-module gradué est un
polynôme périodique (c'est-à-dire un polynôme numérique sur les
classes de congruence des entiers modulo $n$ pour un certain $n$).
\end{remark}

\begin{example}
\label{algebra-example-hilbert-function}
Supposons que $S = k[X_1, \ldots, X_d]$.
Par l'Exemple \ref{algebra-example-K0-field}, nous pouvons identifier
$K_0(k) = K'_0(k) = \mathbf{Z}$. Ainsi tout module gradué
de type fini sur $k[X_1, \ldots, X_d]$
donne lieu à un polynôme numérique
$n \mapsto \dim_k(M_n)$.
\end{example}

\begin{lemma}
\label{algebra-lemma-quotient-smaller-d}
Soit $k$ un corps. Supposons que $I \subset k[X_1, \ldots, X_d]$
soit un idéal gradué non nul. Soit $M = k[X_1, \ldots, X_d]/I$.
Alors le polynôme numérique $n \mapsto \dim_k(M_n)$ (voir
l'Exemple \ref{algebra-example-hilbert-function})
est de degré $ < d - 1$ (ou est nul si $d = 1$).
\end{lemma}

\begin{proof}
Le polynôme numérique associé au module gradué
$k[X_1, \ldots, X_d]$ est $n \mapsto \binom{n - 1 + d}{d - 1}$.
Pour tout élément homogène non nul $f \in I$ de degré $e$
et tout degré $n >> e$, nous avons $I_n \supset f \cdot k[X_1, \ldots, X_d]_{n-e}$
et donc $\dim_k(I_n) \geq \binom{n - e - 1 + d}{d - 1}$. Ainsi
$\dim_k(M_n) \leq \binom{n - 1 + d}{d - 1} - \binom{n - e - 1 + d}{d - 1}$.
Le résultat suit car la dernière expression
est de degré $ < d - 1$ (ou est nulle si $d = 1$).
\end{proof}

\section{Anneaux locaux noethériens}
\label{algebra-section-Noetherian-local}

\noindent
Dans toute cette section, $(R, \mathfrak m, \kappa)$ est un anneau local
noethérien. Nous y développons quelques éléments de théorie des fonctions
de Hilbert des modules. Soit $M$ un $R$-module de type fini. Nous définissons
la {\it fonction de Hilbert} de $M$ comme la fonction
$$
\varphi_M : n
\longmapsto
\text{longueur}_R(\mathfrak m^nM/{\mathfrak m}^{n + 1}M)
$$
définie pour tout entier $n \geq 0$. Un autre invariant important est la
fonction
$$
\chi_M : n
\longmapsto
\text{longueur}_R(M/{\mathfrak m}^{n + 1}M)
$$
définie pour tout entier $n \geq 0$.
D'après le Lemme \ref{algebra-lemma-length-additive}, nous avons
$$
\chi_M(n) = \sum\nolimits_{i = 0}^n \varphi_M(i).
$$
Il existe une variante de cette construction qui utilise un idéal de définition.

\begin{definition}
\label{algebra-definition-ideal-definition}
Soit $(R, \mathfrak m)$ un anneau local noethérien.
Un idéal $I \subset R$ tel que $\sqrt{I} = \mathfrak m$ est appelé
{\it idéal de définition de $R$}.
\end{definition}

\noindent
Soit $I \subset R$ un idéal de définition.
Comme $R$ est noethérien, cela implique que
$\mathfrak m^r \subset I$ pour un certain $r$; voir le Lemme
\ref{algebra-lemma-Noetherian-power}. Tout $R$-module de type fini
annulé par une puissance de $I$ est donc de longueur finie; voir le Lemme
\ref{algebra-lemma-length-finite}.
Il est ainsi légitime de définir
$$
\varphi_{I, M}(n) = \text{longueur}_R(I^nM/I^{n + 1}M)
\quad\text{et}\quad
\chi_{I, M}(n) = \text{longueur}_R(M/I^{n + 1}M)
$$
pour tout $n \geq 0$. Nous avons encore
$$
\chi_{I, M}(n) = \sum\nolimits_{i = 0}^n \varphi_{I, M}(i).
$$

\begin{lemma}
\label{algebra-lemma-differ-finite}
Supposons que $M' \subset M$ soient des $R$-modules de type fini
dont le quotient est de longueur finie. Il existe alors des
constantes $c_1, c_2$ telles que, pour tout $n \geq c_2$, nous ayons
$$
c_1 + \chi_{I, M'}(n - c_2) \leq \chi_{I, M}(n) \leq
c_1 + \chi_{I, M'}(n)
$$
\end{lemma}

\begin{proof}
Comme $M/M'$ est de longueur finie, il existe $c_2 \geq 0$ tel que
$I^{c_2}M \subset M'$. Posons $c_1 = \text{longueur}_R(M/M')$.
Pour $n \geq c_2$, nous avons
\begin{eqnarray*}
\chi_{I, M}(n)
& = &
\text{longueur}_R(M/I^{n + 1}M) \\
& = &
c_1 + \text{longueur}_R(M'/I^{n + 1}M) \\
& \leq &
c_1 + \text{longueur}_R(M'/I^{n + 1}M') \\
& = &
c_1 + \chi_{I, M'}(n)
\end{eqnarray*}
D'autre part, comme $I^{c_2}M \subset M'$,
nous avons $I^nM \subset I^{n - c_2}M'$ pour $n \geq c_2$.
Ainsi, pour $n \geq c_2$, nous obtenons
\begin{eqnarray*}
\chi_{I, M}(n)
& = &
\text{longueur}_R(M/I^{n + 1}M) \\
& = &
c_1 + \text{longueur}_R(M'/I^{n + 1}M) \\
& \geq &
c_1 + \text{longueur}_R(M'/I^{n + 1 - c_2}M') \\
& = &
c_1 + \chi_{I, M'}(n - c_2)
\end{eqnarray*}
ce qui achève la démonstration.
\end{proof}

\begin{lemma}
\label{algebra-lemma-hilbert-ses}
Supposons que $0 \to M' \to M \to M'' \to 0$
soit une suite exacte courte de $R$-modules de type fini.
Il existe alors un sous-module $N \subset M'$ dont le quotient est de
longueur finie $l$, et un entier $c \geq 0$, tels que
$$
\chi_{I, M}(n) = \chi_{I, M''}(n) + \chi_{I, N}(n - c) + l
$$
et
$$
\varphi_{I, M}(n) = \varphi_{I, M''}(n) + \varphi_{I, N}(n - c)
$$
pour tout $n \geq c$.
\end{lemma}

\begin{proof}
Remarquons que $M/I^nM \to M''/I^nM''$ est surjectif
et a pour noyau $M' / M' \cap I^nM$. D'après le lemme d'Artin--Rees
(Lemme \ref{algebra-lemma-Artin-Rees}), il existe une constante
$c$ telle que $M' \cap I^nM =
I^{n - c}(M' \cap I^cM)$. Posons $N = M' \cap I^cM$.
Remarquons que $I^c M' \subset N \subset M'$.
Ainsi $\text{longueur}_R(M' / M' \cap I^nM)
= \text{longueur}_R(M'/N) + \text{longueur}_R(N/I^{n - c}N)$ pour $n \geq c$.
La suite exacte courte
$$
0 \to M' / M' \cap I^nM \to M/I^nM \to M''/I^nM'' \to 0
$$
et l'additivité des longueurs (Lemme \ref{algebra-lemma-length-additive})
donnent l'égalité
$$
\chi_{I, M}(n - 1)
=
\chi_{I, M''}(n - 1)
+
\chi_{I, N}(n - c - 1)
+
\text{longueur}_R(M'/N)
$$
pour $n \geq c$. Nous avons
$\varphi_{I, M}(n) = \chi_{I, M}(n) - \chi_{I, M}(n - 1)$,
et de même pour les modules $M''$ et $N$. Nous obtenons donc
$\varphi_{I, M}(n) = \varphi_{I, M''}(n) + \varphi_{I, N}(n-c)$ pour
$n \geq c$.
\end{proof}

\begin{lemma}
\label{algebra-lemma-hilbert-change-I}
Supposons que $I$, $I'$ soient deux idéaux de définition
de l'anneau local noethérien $R$. Soit $M$ un
$R$-module de type fini. Il existe une constante $a$ telle que
$\chi_{I, M}(n) \leq \chi_{I', M}(an)$ pour $n \geq 1$.
\end{lemma}

\begin{proof}
Il existe un entier $c \geq 1$ tel que $(I')^c \subset I$.
Nous obtenons donc une surjection $M/(I')^{c(n + 1)}M \to M/I^{n + 1}M$.
Le résultat s'ensuit avec $a = 2c - 1$.
\end{proof}

\begin{proposition}
\label{algebra-proposition-hilbert-function-polynomial}
Soit $R$ un anneau local noethérien. Soit $M$ un $R$-module de type fini.
Soit $I \subset R$ un idéal de définition.
La fonction de Hilbert $\varphi_{I, M}$ et la fonction
$\chi_{I, M}$ sont des polynômes numériques.
\end{proposition}

\begin{proof}
Considérons l'anneau gradué $S = R/I \oplus I/I^2 \oplus I^2/I^3 \oplus
\ldots = \bigoplus_{d \geq 0} I^d/I^{d + 1}$. Considérons le
$S$-module gradué $N = M/IM \oplus IM/I^2M \oplus \ldots =
\bigoplus_{d \geq 0} I^dM/I^{d + 1}M$. La paire $(S, N)$ satisfait
aux hypothèses de la Proposition \ref{algebra-proposition-graded-hilbert-polynomial}.
Le résultat pour $\varphi_{I, M}$ découle donc de cette proposition et du
Lemme \ref{algebra-lemma-length-K0}. Le résultat pour $\chi_{I, M}$ découle
de ce qui précède et du Lemme \ref{algebra-lemma-numerical-polynomial}.
\end{proof}

\begin{definition}
\label{algebra-definition-hilbert-polynomial}
Soit $R$ un anneau local noethérien. Soit $M$ un $R$-module de type fini.
Le {\it polynôme de Hilbert} de $M$ sur $R$ est l'élément
$P(t) \in \mathbf{Q}[t]$ tel que $P(n) = \varphi_M(n)$ pour $n \gg 0$.
\end{definition}

\noindent
La Proposition \ref{algebra-proposition-hilbert-function-polynomial}
montre que le polynôme de Hilbert existe.

\begin{lemma}
\label{algebra-lemma-d-independent}
Soit $R$ un anneau local noethérien. Soit $M$ un $R$-module de type fini.
\begin{enumerate}
\item Le degré du polynôme numérique $\varphi_{I, M}$ est indépendant
de l'idéal de définition $I$.
\item Le degré du polynôme numérique $\chi_{I, M}$ est indépendant
de l'idéal de définition $I$.
\end{enumerate}
\end{lemma}

\begin{proof}
La partie (2) découle immédiatement du Lemme \ref{algebra-lemma-hilbert-change-I}.
La partie (1) découle de (2), car
$\varphi_{I, M}(n) = \chi_{I, M}(n) - \chi_{I, M}(n - 1)$
pour $n \geq 1$.
\end{proof}

\begin{definition}
\label{algebra-definition-d}
Soit $R$ un anneau local noethérien et $M$ un $R$-module de type fini.
On note {\it $d(M)$} l'élément de $\{-\infty, 0, 1, 2, \ldots \}$
défini comme suit :
\begin{enumerate}
\item si $M = 0$, nous posons $d(M) = -\infty$;
\item si $M \not = 0$, alors $d(M)$ est le degré du polynôme
numérique $\chi_M$.
\end{enumerate}
\end{definition}

\noindent
Si $\mathfrak m^nM \not = 0$ pour tout $n$, alors nous voyons que
$d(M)$ est le degré, augmenté d'une unité ($+1$), du polynôme de Hilbert
de $M$.

\begin{lemma}
\label{algebra-lemma-differ-finite-chi}
Soit $R$ un anneau local noethérien. Soit $I \subset R$ un idéal
de définition. Soit $M$ un $R$-module de type fini
qui n'est pas de longueur finie. Si $M' \subset M$ est un sous-module
dont le quotient est de longueur finie, alors $\chi_{I, M} - \chi_{I, M'}$
est un polynôme de degré strictement inférieur ($<$) à celui de
chacun des deux polynômes.
\end{lemma}

\begin{proof}
Cela découle du Lemme \ref{algebra-lemma-differ-finite} par un calcul élémentaire.
\end{proof}

\begin{lemma}
\label{algebra-lemma-hilbert-ses-chi}
Soit $R$ un anneau local noethérien. Soit $I \subset R$ un idéal de
définition. Soit $0 \to M' \to M \to M'' \to 0$ une suite exacte courte
de $R$-modules de type fini. Alors :
\begin{enumerate}
\item si $M'$ n'est pas de longueur finie, alors
$\chi_{I, M} - \chi_{I, M''} - \chi_{I, M'}$
est un polynôme numérique de degré strictement inférieur ($<$) à celui de
$\chi_{I, M'}$;
\item $\max\{ \deg(\chi_{I, M'}), \deg(\chi_{I, M''}) \} = \deg(\chi_{I, M})$;
et
\item $\max\{d(M'), d(M'')\} = d(M)$.
\end{enumerate}
\end{lemma}

\begin{proof}
Démontrons d'abord (1). Soit $N \subset M'$ comme dans le
Lemme \ref{algebra-lemma-hilbert-ses}.
D'après le Lemme \ref{algebra-lemma-differ-finite-chi}, le polynôme numérique
$\chi_{I, M'} - \chi_{I, N}$ est de degré strictement inférieur ($<$)
au degré commun de
$\chi_{I, M'}$ et $\chi_{I, N}$. D'après le Lemme \ref{algebra-lemma-hilbert-ses},
la différence
$$
\chi_{I, M}(n) - \chi_{I, M''}(n) - \chi_{I, N}(n - c)
$$
est constante pour $n \gg 0$. Par un calcul élémentaire, la différence
$\chi_{I, N}(n) - \chi_{I, N}(n - c)$ est de degré strictement
inférieur ($<$) au degré de
$\chi_{I, N}$, qui est strictement positif (voir ci-dessus). En réunissant
ces observations, nous obtenons (1).

\medskip\noindent
Remarquons que les coefficients dominants de $\chi_{I, M'}$ et
$\chi_{I, M''}$ sont non négatifs. Le degré de
$\chi_{I, M'} + \chi_{I, M''}$ est donc égal au maximum des degrés.
Ainsi, si $M'$ n'est pas de longueur finie, alors (2) découle de (1).
Si $M'$ est de longueur finie, alors
$I^nM \to I^nM''$ est un isomorphisme pour tout $n \gg 0$ d'après
Artin--Rees (Lemme \ref{algebra-lemma-Artin-Rees}). Ainsi
$M/I^nM \to M''/I^nM''$ est surjectif et a pour noyau $M'$ pour
$n \gg 0$, et nous voyons que
$\chi_{I, M}(n) - \chi_{I, M''}(n) = \text{longueur}(M')$
pour tout $n \gg 0$. La partie (2) est donc également vraie dans ce cas.

\medskip\noindent
Démonstration de (3). Cela découle de (2), sauf si l'un des modules
$M$, $M'$ ou $M''$ est nul. Nous omettons la démonstration dans ces cas particuliers.
\end{proof}

\section{Dimension}
\label{algebra-section-dimension}

\noindent
Comparer avec
Topologie, section \ref{topology-section-krull-dimension}.

\begin{definition}
\label{algebra-definition-chain-primes}
Soit $R$ un anneau. Une {\it chaîne d'idéaux premiers} est une suite
$\mathfrak p_0 \subset \mathfrak p_1 \subset \ldots \subset \mathfrak p_n$
d'idéaux premiers de $R$ telle que $\mathfrak p_i \not = \mathfrak p_{i + 1}$
pour $i = 0, \ldots, n - 1$. La {\it longueur} de cette chaîne d'idéaux
premiers est $n$.
\end{definition}

\noindent
Rappelons qu'il existe une bijection renversant les inclusions entre les
idéaux premiers d'un anneau $R$ et les sous-ensembles fermés irréductibles de
$\Spec(R)$; voir le Lemme \ref{algebra-lemma-irreducible}.

\begin{definition}
\label{algebra-definition-Krull}
La {\it dimension de Krull} de l'anneau $R$ est la dimension de Krull
de l'espace topologique $\Spec(R)$; voir
Topologie, Définition \ref{topology-definition-Krull}.
Autrement dit, c'est la borne supérieure des entiers $n\geq 0$
pour lesquels $R$ possède une chaîne d'idéaux premiers
$$
\mathfrak p_0
\subset
\mathfrak p_1
\subset
\ldots
\subset
\mathfrak p_n, \quad
\mathfrak p_i \not = \mathfrak p_{i + 1}.
$$
Cette chaîne est de longueur $n$.
\end{definition}

\begin{definition}
\label{algebra-definition-height}
La {\it hauteur} d'un idéal premier $\mathfrak p$ d'un anneau $R$
est la dimension de l'anneau local $R_{\mathfrak p}$.
\end{definition}

\begin{lemma}
\label{algebra-lemma-dimension-height}
La dimension de Krull de $R$ est la borne supérieure des hauteurs de ses
idéaux premiers (maximaux).
\end{lemma}

\begin{proof}
Cela résulte du fait que l'on peut toujours ajouter un idéal maximal à la fin
d'une chaîne d'idéaux premiers.
\end{proof}

\begin{lemma}
\label{algebra-lemma-Noetherian-dimension-0}
Un anneau noethérien de dimension $0$ est artinien.
Réciproquement, tout anneau artinien est noethérien de dimension zéro.
\end{lemma}

\begin{proof}
Supposons que $R$ soit un anneau noethérien de dimension $0$.
D'après le Lemme \ref{algebra-lemma-Noetherian-topology}, l'espace $\Spec(R)$
est noethérien. D'après Topologie, Lemme \ref{topology-lemma-Noetherian},
$\Spec(R)$ possède un nombre fini de composantes irréductibles, disons
$\Spec(R) = Z_1 \cup \ldots \cup Z_r$.
D'après le Lemme \ref{algebra-lemma-irreducible}, chaque $Z_i = V(\mathfrak p_i)$,
où $\mathfrak p_i$ est un idéal premier minimal. Comme la dimension est $0$,
ces $\mathfrak p_i$ sont également maximaux. Ainsi $\Spec(R)$ est l'espace
topologique discret dont les éléments sont les $\mathfrak p_i$.
Tout élément $f$ du radical de Jacobson $\bigcap \mathfrak p_i$ est nilpotent,
car sinon $R_f$ ne serait pas l'anneau nul et nous aurions un autre idéal premier.
D'après le Lemme \ref{algebra-lemma-product-local}, $R$ est égal à
$\prod R_{\mathfrak p_i}$. Comme $R_{\mathfrak p_i}$ est lui aussi noethérien
et de dimension $0$, les arguments précédents montrent que son radical
$\mathfrak p_iR_{\mathfrak p_i}$ est localement nilpotent.
Le Lemme \ref{algebra-lemma-Noetherian-power} donne
$\mathfrak p_i^nR_{\mathfrak p_i} = 0$ pour un certain $n \geq 1$.
D'après le Lemme \ref{algebra-lemma-length-finite}, nous en concluons que
$R_{\mathfrak p_i}$ est de longueur finie sur $R$. Le Lemme
\ref{algebra-lemma-artinian-finite-length} montre alors que $R$ est artinien.

\medskip\noindent
Si $R$ est un anneau artinien, alors le Lemme
\ref{algebra-lemma-artinian-finite-length} montre qu'il est noethérien. Tous ses
idéaux premiers sont maximaux par combinaison des Lemmes
\ref{algebra-lemma-artinian-finite-nr-max},
\ref{algebra-lemma-artinian-radical-nilpotent} et \ref{algebra-lemma-product-local}.
\end{proof}

\noindent
Dans la suite, nous utiliserons l'invariant $d(-)$ défini dans la
Définition \ref{algebra-definition-d}. Commençons par un lemme préparatoire.

\begin{lemma}
\label{algebra-lemma-dimension-0-d-0}
Soit $R$ un anneau local noethérien.
Alors $\dim(R) = 0 \Leftrightarrow d(R) = 0$.
\end{lemma}

\begin{proof}
En effet, $d(R) = 0$ si et seulement si $R$ est de longueur finie
comme $R$-module. Voir le Lemme \ref{algebra-lemma-artinian-finite-length}.
\end{proof}

\begin{proposition}
\label{algebra-proposition-dimension-zero-ring}
Soit $R$ un anneau. Les conditions suivantes sont équivalentes :
\begin{enumerate}
\item $R$ est artinien;
\item $R$ est noethérien et $\dim(R) = 0$;
\item $R$ est de longueur finie comme module sur lui-même;
\item $R$ est un produit fini d'anneaux locaux artiniens;
\item $R$ est noethérien et $\Spec(R)$ est un espace topologique discret fini;
\item $R$ est un produit fini d'anneaux locaux noethériens de dimension $0$;
\item $R$ est un produit fini d'anneaux locaux noethériens $R_i$
tels que $d(R_i) = 0$;
\item $R$ est un produit fini d'anneaux locaux noethériens $R_i$
dont les idéaux maximaux sont nilpotents;
\item $R$ est noethérien, possède un nombre fini d'idéaux maximaux
et son radical de Jacobson est nilpotent; et
\item $R$ est noethérien et il n'existe aucune inclusion stricte entre
ses idéaux premiers.
\end{enumerate}
\end{proposition}

\begin{proof}
C'est une combinaison des Lemmes
\ref{algebra-lemma-product-local},
\ref{algebra-lemma-artinian-finite-length},
\ref{algebra-lemma-Noetherian-dimension-0} et
\ref{algebra-lemma-dimension-0-d-0}.
\end{proof}

\begin{lemma}
\label{algebra-lemma-height-1}
Soit $R$ un anneau local noethérien.
Les conditions suivantes sont équivalentes :
\begin{enumerate}
\item
\label{algebra-item-dim-1}
$\dim(R) = 1$;
\item
\label{algebra-item-d-1}
$d(R) = 1$;
\item
\label{algebra-item-Vx}
il existe $x \in \mathfrak m$, avec $x$ non nilpotent, tel que
$V(x) = \{\mathfrak m\}$;
\item
\label{algebra-item-x}
il existe $x \in \mathfrak m$, avec $x$ non nilpotent, tel que
$\mathfrak m = \sqrt{(x)}$; et
\item
\label{algebra-item-ideal-1}
il existe un idéal de définition engendré par $1$ élément,
et aucun idéal de définition n'est engendré par $0$ élément.
\end{enumerate}
\end{lemma}

\begin{proof}
Supposons d'abord que $\dim(R) = 1$.
Soient $\mathfrak p_i$ les idéaux premiers minimaux de $R$.
Comme la dimension est $1$, le seul autre idéal premier de $R$ est
$\mathfrak m$.
D'après le Lemme \ref{algebra-lemma-Noetherian-irreducible-components},
ils sont en nombre fini. Nous pouvons donc trouver $x \in \mathfrak m$,
$x \not \in \mathfrak p_i$; voir le Lemme \ref{algebra-lemma-silly}.
Le seul idéal premier contenant $x$ est ainsi $\mathfrak m$, d'où
(\ref{algebra-item-Vx}).

\medskip\noindent
Si (\ref{algebra-item-Vx}) est satisfaite, alors
$\mathfrak m = \sqrt{(x)}$ d'après le Lemme
\ref{algebra-lemma-Zariski-topology}, d'où (\ref{algebra-item-x}).
La réciproque est également claire.
L'équivalence entre (\ref{algebra-item-x}) et (\ref{algebra-item-ideal-1}) découle
directement des définitions.

\medskip\noindent
Supposons (\ref{algebra-item-ideal-1}).
Soit $I = (x)$ un idéal de définition.
Remarquons que $I^n/I^{n + 1}$ est un quotient de $R/I$ par la multiplication
par $x^n$, et donc que $\text{longueur}_R(I^n/I^{n + 1})$ est bornée.
Ainsi $d(R) = 0$ ou $d(R) = 1$, mais $d(R) = 0$ est exclu par
l'hypothèse selon laquelle $0$ n'est pas un idéal de définition.

\medskip\noindent
Supposons (\ref{algebra-item-d-1}). Pour obtenir une contradiction, supposons qu'il
existe des idéaux premiers $\mathfrak p \subset \mathfrak q \subset \mathfrak m$,
les deux inclusions étant strictes. Choisissons un idéal de définition
$I \subset R$. Nous utiliserons plusieurs fois le
Lemme \ref{algebra-lemma-hilbert-ses-chi}. Tout d'abord, appliqué à la suite exacte
$0 \to \mathfrak p \to R \to R/\mathfrak p \to 0$,
il implique que $d(R/\mathfrak p) \leq 1$. Mais cette valeur ne peut
évidemment pas être nulle. Choisissons $x\in \mathfrak q$,
$x\not \in \mathfrak p$. Considérons la suite exacte courte
$$
0 \to R/\mathfrak p \to R/\mathfrak p \to R/(xR + \mathfrak p) \to 0.
$$
Cela implique que $\chi_{I, R/\mathfrak p} - \chi_{I, R/\mathfrak p}
- \chi_{I, R/(xR + \mathfrak p)} = - \chi_{I, R/(xR + \mathfrak p)}$
est de degré strictement inférieur ($ < 1$). Autrement dit,
$d(R/(xR + \mathfrak p)) = 0$, et donc
$\dim(R/(xR + \mathfrak p)) = 0$, d'après le
Lemme \ref{algebra-lemma-dimension-0-d-0}. Mais $R/(xR + \mathfrak p)$ possède les
idéaux premiers distincts $\mathfrak q/(xR + \mathfrak p)$ et
$\mathfrak m/(xR + \mathfrak p)$, ce qui donne la contradiction cherchée.
\end{proof}

\begin{proposition}
\label{algebra-proposition-dimension}
Soit $R$ un anneau local noethérien. Soit $d \geq 0$ un entier.
Les conditions suivantes sont équivalentes :
\begin{enumerate}
\item
\label{algebra-item-dim-d}
$\dim(R) = d$;
\item
\label{algebra-item-d-d}
$d(R) = d$;
\item
\label{algebra-item-ideal-d}
il existe un idéal de définition engendré par $d$ éléments,
et aucun idéal de définition n'est engendré par moins de $d$ éléments.
\end{enumerate}
\end{proposition}

\begin{proof}
Cette démonstration est essentiellement la même que celle du Lemme
\ref{algebra-lemma-height-1}. Nous procédons par récurrence sur $d$.
D'après les Lemmes \ref{algebra-lemma-dimension-0-d-0} et \ref{algebra-lemma-height-1},
nous pouvons supposer que $d > 1$. Notons le nombre minimal de
générateurs d'un idéal de définition de $R$ par $d'(R)$.
Nous allons démontrer les inégalités
$\dim(R) \geq d'(R) \geq d(R) \geq \dim(R)$,
ce qui montrera que ces nombres sont tous égaux.

\medskip\noindent
Supposons d'abord que $\dim(R) = d$.
Soient $\mathfrak p_i$ les idéaux premiers minimaux de $R$.
D'après le Lemme \ref{algebra-lemma-Noetherian-irreducible-components},
ils sont en nombre fini. Nous pouvons donc trouver $x \in \mathfrak m$,
$x \not \in \mathfrak p_i$; voir le Lemme \ref{algebra-lemma-silly}.
Toute chaîne maximale d'idéaux premiers commence par un certain
$\mathfrak p_i$; la dimension de $R/xR$ est donc au plus $d-1$.
Par récurrence, il existe $x_2, \ldots, x_d$ qui engendrent un idéal de
définition dans $R/xR$. Ainsi $R$ possède un idéal de définition engendré
par (au plus) $d$ éléments.

\medskip\noindent
Supposons $d'(R) = d$. Soit $I = (x_1, \ldots, x_d)$ un idéal de définition.
Remarquons que $I^n/I^{n + 1}$ est un quotient d'une somme directe de
$\binom{d + n - 1}{d - 1}$ copies de $R/I$, par multiplication par tous les
monômes de degré $n$ en $x_1, \ldots, x_d$.
Ainsi $\text{longueur}_R(I^n/I^{n + 1})$ est bornée par un polynôme
de degré $d-1$. Donc $d(R) \leq d$.

\medskip\noindent
Supposons $d(R) = d$. Considérons une chaîne d'idéaux premiers
$\mathfrak p \subset \mathfrak q \subset
\mathfrak q_2 \subset \ldots \subset \mathfrak q_e = \mathfrak m$,
où toutes les inclusions sont strictes et $e \geq 2$.
Choisissons un idéal de définition $I \subset R$.
Nous utiliserons plusieurs fois le Lemme \ref{algebra-lemma-hilbert-ses-chi}.
Tout d'abord, appliqué à la suite exacte
$0 \to \mathfrak p \to R \to R/\mathfrak p \to 0$,
il implique que $d(R/\mathfrak p) \leq d$. Mais cette valeur ne peut
évidemment pas être nulle. Choisissons $x\in \mathfrak q$,
$x\not \in \mathfrak p$. Considérons la suite exacte courte
$$
0 \to R/\mathfrak p \to R/\mathfrak p \to R/(xR + \mathfrak p) \to 0.
$$
Cela implique que $\chi_{I, R/\mathfrak p} - \chi_{I, R/\mathfrak p}
- \chi_{I, R/(xR + \mathfrak p)} = - \chi_{I, R/(xR + \mathfrak p)}$
est de degré strictement inférieur ($ < d$). Autrement dit,
$d(R/(xR + \mathfrak p)) \leq d - 1$, et donc
$\dim(R/(xR + \mathfrak p)) \leq d - 1$, par récurrence.
Or $R/(xR + \mathfrak p)$ possède la chaîne d'idéaux premiers
$\mathfrak q/(xR + \mathfrak p) \subset \mathfrak q_2/(xR + \mathfrak p)
\subset \ldots \subset \mathfrak q_e/(xR + \mathfrak p)$, ce qui donne
$e - 1 \leq d - 1$. Comme nous sommes partis d'une chaîne arbitraire
d'idéaux premiers, cela prouve que $\dim(R) \leq d(R)$.

\medskip\noindent
En relisant la démonstration, le lecteur constatera que nous avons obtenu
les inégalités cycliques voulues.
\end{proof}

\noindent
Soit $(R, \mathfrak m)$ un anneau local noethérien.
D'après ce qui précède, il est clair que $\mathfrak m$ ne peut être
engendré par moins de $\dim(R)$ éléments.
D'après le lemme de Nakayama \ref{algebra-lemma-NAK}, le nombre minimal de
générateurs de $\mathfrak m$ est égal à $\dim_{\kappa(\mathfrak m)}
\mathfrak m/\mathfrak m^2$. Nous avons donc l'inégalité fondamentale
$$
\dim(R) \leq \dim_{\kappa(\mathfrak m)} \mathfrak m/\mathfrak m^2.
$$
Il se trouve que les anneaux pour lesquels l'égalité est vérifiée
possèdent de nombreuses bonnes propriétés. On les appelle anneaux locaux
réguliers.

\begin{definition}
\label{algebra-definition-regular-local}
Soit $(R, \mathfrak m)$ un anneau local noethérien de dimension $d$.
\begin{enumerate}
\item Un {\it système de paramètres de $R$} est une suite d'éléments
$x_1, \ldots, x_d \in \mathfrak m$ qui engendre un idéal de définition
de $R$;
\item s'il existe $x_1, \ldots, x_d \in \mathfrak m$ tels que
$\mathfrak m = (x_1, \ldots, x_d)$, alors nous appelons $R$ un
{\it anneau local régulier}, et $x_1, \ldots, x_d$ un {\it système
régulier de paramètres}.
\end{enumerate}
\end{definition}

\noindent
Les lemmes suivants ressortent clairement des démonstrations des lemmes
et de la proposition ci-dessus, mais nous les énonçons explicitement afin
de disposer de références commodes.

\begin{lemma}
\label{algebra-lemma-minimal-over-1}
Soit $R$ un anneau noethérien. Soit $x \in R$.
\begin{enumerate}
\item Si $\mathfrak p$ est minimal au-dessus de $(x)$,
alors la hauteur de $\mathfrak p$ est $0$ ou $1$.
\item Si $\mathfrak p, \mathfrak q \in \Spec(R)$ et si $\mathfrak q$
est minimal au-dessus de $(\mathfrak p, x)$, alors il n'existe aucun idéal
premier strictement compris entre $\mathfrak p$ et $\mathfrak q$.
\end{enumerate}
\end{lemma}

\begin{proof}
Démonstration de (1). Si $\mathfrak p$ est minimal parmi les idéaux premiers
contenant $x$, alors le seul idéal premier de $R_\mathfrak p$ contenant $x$
est l'idéal maximal $\mathfrak p R_\mathfrak p$. En effet, les idéaux premiers
de $R_\mathfrak p$ correspondent de façon $1$ à $1$ aux idéaux premiers de
$R$ contenus dans $\mathfrak p$; voir le Lemme
\ref{algebra-lemma-spec-localization}.
Le Lemme \ref{algebra-lemma-height-1} montre donc que
$\dim(R_\mathfrak p) = 1$ si $x$ n'est pas nilpotent dans
$R_\mathfrak p$. Bien entendu, si $x$ est nilpotent dans
$R_\mathfrak p$, l'argument montre que $\mathfrak pR_\mathfrak p$ est le
seul idéal premier et que la hauteur est $0$.

\medskip\noindent
Démonstration de (2). D'après la partie (1),
$\mathfrak q/\mathfrak p$ est un idéal premier de hauteur $1$ ou $0$
dans $R/\mathfrak p$. Il en résulte immédiatement qu'il ne peut exister
d'idéal premier strictement compris entre $\mathfrak p$ et $\mathfrak q$.
\end{proof}

\begin{lemma}
\label{algebra-lemma-minimal-over-r}
Soit $R$ un anneau noethérien. Soient $f_1, \ldots, f_r \in R$.
\begin{enumerate}
\item Si $\mathfrak p$ est minimal au-dessus de $(f_1, \ldots, f_r)$,
alors la hauteur de $\mathfrak p$ est $\leq r$.
\item Si $\mathfrak p, \mathfrak q \in \Spec(R)$ et si
$\mathfrak q$ est minimal au-dessus de $(\mathfrak p, f_1, \ldots, f_r)$,
alors toute chaîne d'idéaux premiers entre $\mathfrak p$ et
$\mathfrak q$ est de longueur au plus $r$.
\end{enumerate}
\end{lemma}

\begin{proof}
Démonstration de (1). Si $\mathfrak p$ est minimal parmi les idéaux premiers
contenant $f_1, \ldots, f_r$, alors le seul idéal premier de
$R_\mathfrak p$ contenant $f_1, \ldots, f_r$ est l'idéal maximal
$\mathfrak p R_\mathfrak p$. En effet, les idéaux premiers de
$R_\mathfrak p$ correspondent de façon $1$ à $1$ aux idéaux premiers de
$R$ contenus dans $\mathfrak p$; voir le Lemme
\ref{algebra-lemma-spec-localization}. La Proposition \ref{algebra-proposition-dimension}
montre donc que $\dim(R_\mathfrak p) \leq r$.

\medskip\noindent
Démonstration de (2). D'après la partie (1),
$\mathfrak q/\mathfrak p$ est un idéal premier de hauteur $\leq r$.
Cela implique immédiatement l'assertion sur les chaînes d'idéaux premiers
entre $\mathfrak p$ et $\mathfrak q$.
\end{proof}

\begin{lemma}
\label{algebra-lemma-one-equation}
Supposons que $R$ soit un anneau local noethérien et que
$x\in \mathfrak m$ soit un élément de son idéal maximal. Alors
$\dim R \leq \dim R/xR + 1$.
Si $x$ n'appartient à aucun des idéaux premiers minimaux de $R$,
alors il y a égalité. (Par exemple, si $x$ est un non-diviseur de zéro.)
\end{lemma}

\begin{proof}
Si $x_1, \ldots, x_{\dim R/xR} \in R$ s'envoient sur des éléments de
$R/xR$ qui engendrent un idéal de définition de $R/xR$, alors
$x, x_1, \ldots, x_{\dim R/xR}$ engendrent un idéal de définition de $R$.
L'inégalité résulte donc de la Proposition \ref{algebra-proposition-dimension}.
D'autre part, si $x$ n'appartient à aucun idéal premier minimal de $R$,
alors toute chaîne d'idéaux premiers de $R/xR$ donne naissance à une
chaîne dans $R$ à laquelle il manque encore au moins une étape pour être
maximale.
\end{proof}

\begin{lemma}
\label{algebra-lemma-elements-generate-ideal-definition}
Soit $(R, \mathfrak m)$ un anneau local noethérien.
Supposons que $x_1, \ldots, x_d \in \mathfrak m$ engendrent un idéal
de définition et que $d = \dim(R)$. Alors
$\dim(R/(x_1, \ldots, x_i)) = d - i$ pour tout $i = 1, \ldots, d$.
\end{lemma}

\begin{proof}
Cela découle soit de la démonstration de la Proposition
\ref{algebra-proposition-dimension}, soit d'une récurrence sur $d$ utilisant le
Lemme \ref{algebra-lemma-one-equation}.
\end{proof}

\section{Applications de la théorie de la dimension}
\label{algebra-section-applications-dimension-theory}

\noindent
Nous pouvons utiliser les résultats sur la dimension pour montrer que certains
anneaux ont un spectre infini et pour construire de nouveaux anneaux de
Jacobson.

\begin{lemma}
\label{algebra-lemma-Noetherian-local-domain-dim-2-infinite-opens}
Soit $R$ un domaine local noethérien de dimension $\geq 2$.
Toute partie ouverte non vide $U \subset \Spec(R)$ est infinie.
\end{lemma}

\begin{proof}
Pour obtenir une contradiction, supposons que $U \subset \Spec(R)$ soit finie.
Dans ce cas, $(0) \in U$ et $\{(0)\}$ est une partie ouverte de $U$ (car le
complémentaire de $\{(0)\}$ est la réunion des adhérences des autres points).
Nous pouvons donc supposer que $U = \{(0)\}$.
Soit $\mathfrak m \subset R$ l'idéal maximal.
Nous pouvons trouver $x \in \mathfrak m$, $x \not = 0$, tel que
$V(x) \cup U = \Spec(R)$. Autrement dit,
$D(x) = \{(0)\}$. En particulier,
$\dim(R/xR) = \dim(R) - 1 \geq 1$, voir le Lemme \ref{algebra-lemma-one-equation}.
Soient $\overline{y}_2, \ldots, \overline{y}_{\dim(R)} \in R/xR$ des éléments
engendrant un idéal de définition de $R/xR$, voir la Proposition
\ref{algebra-proposition-dimension}. Choisissons des relèvements
$y_2, \ldots, y_{\dim(R)} \in R$, de sorte que
$x, y_2, \ldots, y_{\dim(R)}$ engendrent un idéal de définition de $R$.
Il s'ensuit que $\dim(R/(y_2)) = \dim(R) - 1$ et que
$\dim(R/(y_2, x)) = \dim(R) - 2$, voir le Lemme
\ref{algebra-lemma-elements-generate-ideal-definition}.
Il existe donc un idéal premier $\mathfrak p$ contenant $y_2$ mais pas $x$.
Cela contredit l'égalité $D(x) = \{(0)\}$.
\end{proof}

\noindent
Les anneaux $k[[t]]$, où $k$ est un corps, ainsi que l'anneau des nombres
$p$-adiques sont des anneaux noethériens de dimension $1$ qui n'ont qu'un
nombre fini d'idéaux premiers. C'est la plus grande dimension pour laquelle
cela puisse se produire.

\begin{lemma}
\label{algebra-lemma-Noetherian-finite-nr-primes}
Un anneau noethérien qui n'a qu'un nombre fini d'idéaux premiers est de
dimension $\leq 1$.
\end{lemma}

\begin{proof}
Soit $R$ un anneau noethérien qui n'a qu'un nombre fini d'idéaux premiers.
Si $R$ est un domaine local, le lemme découle du Lemme
\ref{algebra-lemma-Noetherian-local-domain-dim-2-infinite-opens}.
Si $R$ est un domaine, alors $R_\mathfrak m$ est de dimension $\leq 1$ pour
tout idéal maximal $\mathfrak m$, d'après le cas local.
Ainsi, $\dim(R) \leq 1$ d'après le Lemme \ref{algebra-lemma-dimension-height}.
Si $R$ est quelconque, alors $\dim(R/\mathfrak q) \leq 1$ pour tout idéal premier
minimal $\mathfrak q$ de $R$. Puisque tout idéal premier contient un idéal
premier minimal (Lemme \ref{algebra-lemma-Zariski-topology}), il s'ensuit que
$\dim(R) \leq 1$.
\end{proof}

\begin{lemma}
\label{algebra-lemma-finite-type-algebra-finite-nr-primes}
Soit $S$ une algèbre non nulle de type fini sur un corps $k$.
Les conditions suivantes sont équivalentes :
\begin{enumerate}
\item $\dim(S) = 0$,
\item $S$ n'a qu'un nombre fini d'idéaux premiers,
\item $S$ n'a qu'un nombre fini d'idéaux maximaux,
\item $\Spec(S)$ satisfait l'une des conditions équivalentes du
Lemme \ref{algebra-lemma-ring-with-only-minimal-primes},
\item $\dim_k(S) < \infty$,
\item $S$ est artinien,
\item $\Spec(S)$ est un espace topologique discret,
\item ajouter ici d'autres conditions.
\end{enumerate}
\end{lemma}

\begin{proof}
Il résulte immédiatement des définitions que (1) équivaut à (4), en
considérant l'assertion (5) du Lemme
\ref{algebra-lemma-ring-with-only-minimal-primes}.
Rappelons que $\Spec(S)$ est sobre, noethérien et de Jacobson, voir les Lemmes
\ref{algebra-lemma-spec-spectral}, \ref{algebra-lemma-Noetherian-topology},
\ref{algebra-lemma-finite-type-field-Jacobson} et \ref{algebra-lemma-jacobson}.
Si $S$ est de dimension $0$, chaque point définit une composante irréductible,
et il n'y a qu'un nombre fini de composantes irréductibles (Topologie,
Lemme \ref{topology-lemma-Noetherian}). Ainsi, (1) implique (2).
Il est clair que (2) implique (3). Si (3) est satisfaite, alors $\Spec(S)$ est
discret d'après Topologie, Lemme \ref{topology-lemma-finite-jacobson}, et la
dimension de $S$ est donc $0$.

\medskip\noindent
Nous savons à présent que (1) -- (4) sont équivalentes.
L'implication (5) $\Rightarrow$ (6) est le Lemme
\ref{algebra-lemma-finite-dimensional-algebra}.
L'implication (6) $\Rightarrow$ (7) résulte de la Proposition
\ref{algebra-proposition-dimension-zero-ring}.
L'implication (7) $\Rightarrow$ (4) est immédiate.
Réciproquement, si $S$ satisfait (1) -- (4), alors $S$ possède un nombre fini
d'idéaux premiers $\mathfrak m_1, \ldots, \mathfrak m_r$, tous maximaux.
Remarquons que $\kappa(\mathfrak m_i)$ est une extension finie de $k$ d'après
le Nullstellensatz de Hilbert (Théorème \ref{algebra-theorem-nullstellensatz}).
La Proposition \ref{algebra-proposition-dimension-zero-ring} montre également que
$S$ est artinien. Ensuite, le Lemme \ref{algebra-lemma-artinian-finite-length} donne
$\text{longueur}_S(S) < \infty$. Ainsi, $\dim_k(S) < \infty$ d'après le
Lemme \ref{algebra-lemma-pushdown-module}. Nous concluons que (1) -- (7) sont
équivalentes. (Remarque : une autre manière, plus classique, de démontrer que
$\dim(S) = 0 \Rightarrow \dim_k(S) < \infty$ consiste à utiliser la
normalisation de Noether, mais nous ne disposons pas encore de ce résultat.)
\end{proof}

\begin{lemma}
\label{algebra-lemma-noetherian-dim-1-Jacobson}
Anneaux de Jacobson noethériens.
\begin{enumerate}
\item Tout domaine noethérien $R$ de dimension $1$ qui possède une infinité
d'idéaux premiers est un anneau de Jacobson.
\item Tout anneau noethérien tel que chaque idéal premier $\mathfrak p$ soit
maximal ou soit contenu dans une infinité d'idéaux premiers est un anneau de
Jacobson.
\end{enumerate}
\end{lemma}

\begin{proof}
L'assertion (1) est une reformulation du Lemme \ref{algebra-lemma-pid-jacobson}.

\medskip\noindent
Soit $R$ un anneau noethérien tel que tout idéal premier non maximal
$\mathfrak p$ soit contenu dans une infinité d'idéaux premiers.
Supposons, pour obtenir une contradiction, que $\Spec(R)$ ne soit pas un
espace de Jacobson. Les Lemmes \ref{algebra-lemma-irreducible} et
\ref{algebra-lemma-Noetherian-topology} montrent que $\Spec(R)$ est un espace
topologique sobre et noethérien.
D'après Topologie, Lemme
\ref{topology-lemma-non-jacobson-Noetherian-characterize}, il existe un idéal
non maximal $\mathfrak p \subset R$ tel que $\{\mathfrak p\}$ soit une partie
localement fermée de $\Spec(R)$. Autrement dit, $\mathfrak p$ n'est pas
maximal et $\{\mathfrak p\}$ est une partie ouverte de $V(\mathfrak p)$.
Considérons un idéal premier $\mathfrak q \subset R$ tel que
$\mathfrak p \subset \mathfrak q$. Rappelons que la topologie du spectre de
$(R/\mathfrak p)_{\mathfrak q} = R_{\mathfrak q}/\mathfrak pR_{\mathfrak q}$
est induite par celle de $\Spec(R)$, voir les Lemmes
\ref{algebra-lemma-spec-localization} et \ref{algebra-lemma-spec-closed}.
Ainsi, $\{(0)\}$ est une partie localement fermée de
$\Spec((R/\mathfrak p)_{\mathfrak q})$. Le Lemme
\ref{algebra-lemma-Noetherian-local-domain-dim-2-infinite-opens} donne alors
$\dim((R/\mathfrak p)_{\mathfrak q}) = 1$.
Comme cela vaut pour tout $\mathfrak q \supset \mathfrak p$, nous concluons
que $\dim(R/\mathfrak p) = 1$. Utilisons maintenant l'hypothèse selon laquelle
$\mathfrak p$ est contenu dans une infinité d'idéaux premiers : elle montre
que $\Spec(R/\mathfrak p)$ est infini. L'assertion (1) du lemme implique donc
que $V(\mathfrak p) \cong \Spec(R/\mathfrak p)$ est l'adhérence de l'ensemble
de ses points fermés. C'est la contradiction recherchée, puisque cela signifie
que $\{\mathfrak p\} \subset V(\mathfrak p)$ ne peut pas être ouvert.
\end{proof}

\section{Support et dimension des modules}
\label{algebra-section-support}

\noindent
Quelques résultats élémentaires sur le support et la dimension des modules.

\begin{lemma}
\label{algebra-lemma-filter-Noetherian-module}
Soit $R$ un anneau noethérien et soit $M$ un $R$-module de type fini.
Il existe une filtration par des sous-$R$-modules
$$
0 = M_0 \subset M_1 \subset \ldots \subset M_n = M
$$
telle que chaque quotient $M_i/M_{i-1}$ soit isomorphe à
$R/\mathfrak p_i$ pour un certain idéal premier $\mathfrak p_i$ de $R$.
\end{lemma}

\begin{proof}[Première démonstration]
D'après le Lemme \ref{algebra-lemma-trivial-filter-finite-module}, il suffit de traiter
le cas $M = R/I$ pour un certain idéal $I$.
Considérons l'ensemble $S$ des idéaux $J$ tels que le lemme ne soit pas vrai
pour le module $R/J$, et ordonnons-le par inclusion. Pour obtenir une
contradiction, supposons que $S$ ne soit pas vide. Comme $R$ est noethérien,
$S$ possède un élément maximal $J$. Par définition de $S$, l'idéal $J$ ne
peut pas être premier. Choisissons $a, b\in R$ tels que $ab \in J$, mais que
ni $a \in J$ ni $b\in J$. Considérons la filtration
$0 \subset aR/(J \cap aR) \subset R/J$.
Remarquons que le sous-module $aR/(J \cap aR)$ et le module quotient
$(R/J)/(aR/(J \cap aR))$ sont tous deux cycliques ; écrivons-les respectivement
$R/J'$ et $R/J''$, de sorte que nous ayons une suite exacte courte
$0 \to R/J' \to R/J \to R/J'' \to 0$.
L'inclusion $J \subset J'$ est stricte puisque $b \in J'$, et l'inclusion
$J \subset J''$ est stricte puisque $a \in J''$.
Par maximalité de $J$, les modules $R/J'$ et $R/J''$ possèdent donc tous deux
une filtration comme ci-dessus, et il en va de même de $R/J$. Contradiction.
\end{proof}

\begin{proof}[Seconde démonstration]
Pour un $R$-module $M$, disons que $P(M)$ est satisfaite s'il existe une
filtration comme dans l'énoncé du lemme. Remarquons que $P$ est stable par
extensions et qu'elle est satisfaite pour $0$. D'après le Lemme
\ref{algebra-lemma-trivial-filter-finite-module}, il suffit de démontrer que $P(R/I)$
est satisfaite pour tout idéal $I$. Dans le cas contraire, comme $R$ est
noethérien, il existe un contre-exemple maximal $J$. D'après l'Exemple
\ref{algebra-example-oka-family-property-modules} et la Proposition
\ref{algebra-proposition-oka}, l'idéal $J$ est premier, ce qui est une contradiction.
\end{proof}

\begin{lemma}
\label{algebra-lemma-filter-primes-in-support}
Soient $R$, $M$, $M_i$ et $\mathfrak p_i$ comme dans le Lemme
\ref{algebra-lemma-filter-Noetherian-module}.
Alors $\text{Supp}(M) = \bigcup V(\mathfrak p_i)$ et, en particulier,
$\mathfrak p_i \in \text{Supp}(M)$.
\end{lemma}

\begin{proof}
Cela découle des Lemmes \ref{algebra-lemma-support-closed} et
\ref{algebra-lemma-support-quotient}.
\end{proof}

\begin{lemma}
\label{algebra-lemma-support-point}
Supposons que $R$ soit un anneau local noethérien d'idéal maximal
$\mathfrak m$. Soit $M$ un $R$-module non nul de type fini.
Alors $\text{Supp}(M) = \{ \mathfrak m\}$ si et seulement si $M$ est de
longueur finie sur $R$.
\end{lemma}

\begin{proof}
Supposons que $\text{Supp}(M) = \{ \mathfrak m\}$.
Il suffit de montrer que tous les idéaux premiers $\mathfrak p_i$ de la
filtration du Lemme \ref{algebra-lemma-filter-Noetherian-module} sont égaux à l'idéal
maximal. Cela résulte immédiatement du Lemme
\ref{algebra-lemma-filter-primes-in-support}.

\medskip\noindent
Supposons que $M$ soit de longueur finie sur $R$.
Alors $\mathfrak m^n M = 0$ d'après le Lemme \ref{algebra-lemma-length-infinite}.
Comme un élément de $\mathfrak m$ devient une unité dans $R_{\mathfrak p}$
pour tout idéal premier $\mathfrak p \not = \mathfrak m$ de $R$, nous voyons
que $M_{\mathfrak p} = 0$.
\end{proof}

\begin{lemma}
\label{algebra-lemma-Noetherian-power-ideal-kills-module}
Soit $R$ un anneau noethérien.
Soit $I \subset R$ un idéal.
Soit $M$ un $R$-module de type fini.
Alors $I^nM = 0$ pour un certain $n \geq 0$ si et seulement si
$\text{Supp}(M) \subset V(I)$.
\end{lemma}

\begin{proof}
En effet, $I^nM = 0$ équivaut à $I^n \subset \text{Ann}(M)$.
Puisque $R$ est noethérien, cela équivaut à
$I \subset \sqrt{\text{Ann}(M)}$, voir le Lemme
\ref{algebra-lemma-Noetherian-power}.
À son tour, cette condition équivaut à
$V(I) \supset V(\text{Ann}(M))$, voir le Lemme
\ref{algebra-lemma-Zariski-topology}.
D'après le Lemme \ref{algebra-lemma-support-closed}, elle équivaut à
$V(I) \supset \text{Supp}(M)$.
\end{proof}

\begin{lemma}
\label{algebra-lemma-filter-minimal-primes-in-support}
Soient $R$, $M$, $M_i$ et $\mathfrak p_i$ comme dans le Lemme
\ref{algebra-lemma-filter-Noetherian-module}.
Les éléments minimaux de l'ensemble $\{\mathfrak p_i\}$ sont les éléments
minimaux de $\text{Supp}(M)$. Le nombre de fois qu'un idéal premier minimal
$\mathfrak p$ apparaît est
$$
\#\{i \mid \mathfrak p_i = \mathfrak p\}
=
\text{longueur}_{R_\mathfrak p} M_{\mathfrak p}.
$$
\end{lemma}

\begin{proof}
La première assertion résulte de l'égalité
$\text{Supp}(M) = \bigcup V(\mathfrak p_i)$, voir le Lemme
\ref{algebra-lemma-filter-primes-in-support}.
Soit $\mathfrak p \in \text{Supp}(M)$ un élément minimal.
Le support de $M_{\mathfrak p}$ est l'ensemble réduit à l'idéal maximal
$\mathfrak p R_{\mathfrak p}$. D'après le Lemme \ref{algebra-lemma-support-point},
la longueur de $M_{\mathfrak p}$ est donc finie et $> 0$.
Remarquons ensuite que $M_{\mathfrak p}$ possède une filtration dont les
sous-quotients sont
$
(R/\mathfrak p_i)_{\mathfrak p}
=
R_{\mathfrak p}/{\mathfrak p_i}R_{\mathfrak p}
$.
Ils sont nuls si $\mathfrak p_i \not \subset \mathfrak p$, et ils sont égaux à
$\kappa(\mathfrak p)$ si $\mathfrak p_i \subset \mathfrak p$, car la
minimalité de $\mathfrak p$ donne alors $\mathfrak p_i = \mathfrak p$.
Le résultat en découle puisque $\kappa(\mathfrak p)$ est de longueur $1$.
\end{proof}

\begin{lemma}
\label{algebra-lemma-support-dimension-d}
Soit $R$ un anneau local noethérien.
Soit $M$ un $R$-module de type fini.
Alors $d(M) = \dim(\text{Supp}(M))$, où $d(M)$ est défini dans la
Définition \ref{algebra-definition-d}.
\end{lemma}

\begin{proof}
Soient $M_i, \mathfrak p_i$ comme dans le Lemme
\ref{algebra-lemma-filter-Noetherian-module}.
Le Lemme \ref{algebra-lemma-hilbert-ses-chi} donne l'égalité
$d(M) = \max \{ d(R/\mathfrak p_i) \}$. D'après la Proposition
\ref{algebra-proposition-dimension}, nous avons
$d(R/\mathfrak p_i) = \dim(R/\mathfrak p_i)$.
Il est clair que $\dim(R/\mathfrak p_i) = \dim V(\mathfrak p_i)$.
Comme tous les idéaux premiers minimaux de $\text{Supp}(M)$ figurent parmi les
$\mathfrak p_i$ (Lemme \ref{algebra-lemma-filter-minimal-primes-in-support}), le
résultat en découle.
\end{proof}

\begin{lemma}
\label{algebra-lemma-ses-dimension}
Soit $R$ un anneau noethérien. Soit
$0 \to M' \to M \to M'' \to 0$ une suite exacte courte de $R$-modules de
type fini. Alors
$\max\{\dim(\text{Supp}(M')), \dim(\text{Supp}(M''))\} =
\dim(\text{Supp}(M))$.
\end{lemma}

\begin{proof}
Si $R$ est local, cela découle immédiatement des Lemmes
\ref{algebra-lemma-support-dimension-d} et \ref{algebra-lemma-hilbert-ses-chi}.
Un argument plus élémentaire, qui vaut aussi lorsque $R$ n'est pas local,
consiste à utiliser le fait que $\text{Supp}(M')$, $\text{Supp}(M'')$ et
$\text{Supp}(M)$ sont fermés (Lemme \ref{algebra-lemma-support-closed}) et que
$\text{Supp}(M) = \text{Supp}(M') \cup \text{Supp}(M'')$
(Lemme \ref{algebra-lemma-support-quotient}).
\end{proof}

\section{Idéaux premiers associés}
\label{algebra-section-ass}

\noindent
Voici la définition usuelle. Pour les anneaux non noethériens et les modules
qui ne sont pas de type fini, il peut être plus approprié d'utiliser la
définition de la section \ref{algebra-section-weakly-ass}.

\begin{definition}
\label{algebra-definition-associated}
Soit $R$ un anneau. Soit $M$ un $R$-module.
Un idéal premier $\mathfrak p$ de $R$ est {\it associé} à $M$ s'il existe un
élément $m \in M$ dont l'annulateur est $\mathfrak p$.
L'ensemble de tous ces idéaux premiers se note $\text{Ass}_R(M)$ ou
$\text{Ass}(M)$.
\end{definition}

\begin{lemma}
\label{algebra-lemma-ass-support}
Soit $R$ un anneau. Soit $M$ un $R$-module.
Alors $\text{Ass}(M) \subset \text{Supp}(M)$.
\end{lemma}

\begin{proof}
Si $m \in M$ a pour annulateur $\mathfrak p$, alors aucun élément de
$R \setminus \mathfrak p$ n'annule $m$. Ainsi, $m$ est un élément non nul de
$M_{\mathfrak p}$, c'est-à-dire que $\mathfrak p \in \text{Supp}(M)$.
\end{proof}

\begin{lemma}
\label{algebra-lemma-ass}
Soit $R$ un anneau. Soit $0 \to M' \to M \to M'' \to 0$ une suite exacte
courte de $R$-modules. Alors $\text{Ass}(M') \subset \text{Ass}(M)$ et
$\text{Ass}(M) \subset \text{Ass}(M') \cup \text{Ass}(M'')$.
De plus, $\text{Ass}(M' \oplus M'') = \text{Ass}(M') \cup \text{Ass}(M'')$.
\end{lemma}

\begin{proof}
Si $m' \in M'$, l'annulateur de $m'$ considéré comme élément de $M'$ est le
même que l'annulateur de $m'$ considéré comme élément de $M$. On obtient donc
l'inclusion $\text{Ass}(M') \subset \text{Ass}(M)$. Soit $m \in M$ un élément
dont l'annulateur est un idéal premier $\mathfrak p$. S'il existe
$g \in R$, $g \not \in \mathfrak p$, tel que $m' = gm \in M'$, alors
l'annulateur de $m'$ est $\mathfrak p$. S'il n'existe aucun
$g \in R$, $g \not \in \mathfrak p$, tel que $gm \in M'$, alors l'annulateur
de l'image $m'' \in M''$ de $m$ est $\mathfrak p$. Cela démontre l'inclusion
$\text{Ass}(M) \subset \text{Ass}(M') \cup \text{Ass}(M'')$.
Nous omettons la démonstration de la dernière assertion.
\end{proof}

\begin{lemma}
\label{algebra-lemma-ass-filter}
Soit $R$ un anneau et soit $M$ un $R$-module.
Supposons qu'il existe une filtration par des sous-$R$-modules
$$
0 = M_0 \subset M_1 \subset \ldots \subset M_n = M
$$
telle que chaque quotient $M_i/M_{i-1}$ soit isomorphe à $R/\mathfrak p_i$
pour un certain idéal premier $\mathfrak p_i$ de $R$.
Alors $\text{Ass}(M) \subset \{\mathfrak p_1, \ldots, \mathfrak p_n\}$.
\end{lemma}

\begin{proof}
Raisonnons par récurrence sur la longueur $n$ de la filtration $\{ M_i \}$.
Choisissons $m \in M$ dont l'annulateur est un idéal premier $\mathfrak p$.
Si $m \in M_{n-1}$, le résultat découle de l'hypothèse de récurrence. Sinon,
$m$ s'envoie sur un élément non nul de
$M/M_{n-1} \cong R/\mathfrak p_n$. Nous avons donc
$\mathfrak p \subset \mathfrak p_n$.
S'il n'y a pas égalité, nous pouvons trouver $f \in \mathfrak p_n$,
$f \not\in \mathfrak p$. Dans ce cas, l'annulateur de $fm$ est encore
$\mathfrak p$ et $fm \in M_{n-1}$. L'hypothèse de récurrence conclut.
\end{proof}

\begin{lemma}
\label{algebra-lemma-finite-ass}
Soit $R$ un anneau noethérien.
Soit $M$ un $R$-module de type fini.
Alors $\text{Ass}(M)$ est fini.
\end{lemma}

\begin{proof}
Cela résulte immédiatement du Lemme \ref{algebra-lemma-ass-filter} et du Lemme
\ref{algebra-lemma-filter-Noetherian-module}.
\end{proof}

\begin{proposition}
\label{algebra-proposition-minimal-primes-associated-primes}
Soit $R$ un anneau noethérien.
Soit $M$ un $R$-module de type fini.
Les ensembles d'idéaux premiers suivants sont égaux :
\begin{enumerate}
\item les idéaux premiers minimaux du support de $M$ ;
\item les idéaux premiers minimaux de $\text{Ass}(M)$ ;
\item pour toute filtration $0 = M_0 \subset M_1 \subset \ldots
\subset M_{n-1} \subset M_n = M$ telle que
$M_i/M_{i-1} \cong R/\mathfrak p_i$, les idéaux premiers minimaux de
l'ensemble $\{\mathfrak p_i\}$.
\end{enumerate}
\end{proposition}

\begin{proof}
Choisissons une filtration comme en (3).
Le Lemme \ref{algebra-lemma-filter-minimal-primes-in-support} montre que les ensembles
de (1) et (3) sont égaux.

\medskip\noindent
Soit $\mathfrak p$ un élément minimal de l'ensemble $\{\mathfrak p_i\}$.
Choisissons $i$ minimal tel que $\mathfrak p = \mathfrak p_i$.
Prenons $m \in M_i$, $m \not \in M_{i-1}$. L'annulateur de $m$ est contenu
dans $\mathfrak p_i = \mathfrak p$ et contient
$\mathfrak p_1 \mathfrak p_2 \ldots \mathfrak p_i$. Par notre choix de
$i$ et de $\mathfrak p$, nous avons $\mathfrak p_j \not \subset \mathfrak p$
pour $j < i$, et donc
$\mathfrak p_1 \mathfrak p_2 \ldots \mathfrak p_{i - 1}
\not \subset \mathfrak p_i$. Choisissons
$f \in \mathfrak p_1 \mathfrak p_2 \ldots \mathfrak p_{i - 1}$,
$f \not \in \mathfrak p$. Alors $fm$ a pour annulateur $\mathfrak p$.
Ainsi, $\mathfrak p$ est un idéal premier associé à $M$.
Le Lemme \ref{algebra-lemma-ass-support} donne
$\text{Ass}(M) \subset \text{Supp}(M)$ ; par conséquent, $\mathfrak p$ est
minimal dans $\text{Ass}(M)$. L'ensemble des idéaux premiers de (1) est donc
contenu dans l'ensemble des idéaux premiers de (2).

\medskip\noindent
Soit $\mathfrak p$ un élément minimal de $\text{Ass}(M)$.
Puisque $\text{Ass}(M) \subset \text{Supp}(M)$, il existe un élément minimal
$\mathfrak q$ de $\text{Supp}(M)$ tel que
$\mathfrak q \subset \mathfrak p$. Nous venons de montrer que
$\mathfrak q \in \text{Ass}(M)$. Nous avons donc
$\mathfrak q = \mathfrak p$ par minimalité de $\mathfrak p$. Ainsi, l'ensemble des idéaux premiers de (2) est
contenu dans celui de (1).
\end{proof}

\begin{lemma}
\label{algebra-lemma-ass-zero}
\begin{slogan}
Sur un anneau noethérien, tout module non nul possède un idéal premier associé.
\end{slogan}
Soit $R$ un anneau noethérien. Soit $M$ un $R$-module.
Alors
$$
M = (0) \Leftrightarrow \text{Ass}(M) = \emptyset.
$$
\end{lemma}

\begin{proof}
Si $M = (0)$, alors $\text{Ass}(M) = \emptyset$ par définition.
Si $M \not = 0$, choisissons un sous-module non nul de type fini
$M' \subset M$, par exemple le sous-module engendré par un seul élément non
nul. Le Lemme \ref{algebra-lemma-support-zero} montre que $\text{Supp}(M')$ n'est pas
vide. D'après la Proposition
\ref{algebra-proposition-minimal-primes-associated-primes}, il en résulte que
$\text{Ass}(M')$ n'est pas vide. Le Lemme \ref{algebra-lemma-ass} implique alors que
$\text{Ass}(M) \not = \emptyset$.
\end{proof}

\begin{lemma}
\label{algebra-lemma-ass-minimal-prime-support}
Soit $R$ un anneau noethérien.
Soit $M$ un $R$-module.
Tout $\mathfrak p \in \text{Supp}(M)$ qui est minimal parmi les éléments de
$\text{Supp}(M)$ appartient à $\text{Ass}(M)$.
\end{lemma}

\begin{proof}
Si $M$ est un $R$-module de type fini, cela résulte de la Proposition
\ref{algebra-proposition-minimal-primes-associated-primes}.
Dans le cas général, écrivons $M = \bigcup M_\lambda$ comme réunion de ses
sous-modules de type fini et utilisons les égalités
$\text{Supp}(M) = \bigcup \text{Supp}(M_\lambda)$ et
$\text{Ass}(M) = \bigcup \text{Ass}(M_\lambda)$.
\end{proof}

\begin{lemma}
\label{algebra-lemma-ass-zero-divisors}
Soit $R$ un anneau noethérien.
Soit $M$ un $R$-module.
La réunion $\bigcup_{\mathfrak q \in \text{Ass}(M)} \mathfrak q$ est
l'ensemble des éléments de $R$ qui sont des diviseurs de zéro dans $M$.
\end{lemma}

\begin{proof}
Tout élément appartenant à un idéal premier associé est manifestement un
diviseur de zéro dans $M$. Réciproquement, supposons que $x \in R$ soit un
diviseur de zéro dans $M$. Considérons le sous-module
$N = \{m \in M \mid xm = 0\}$. Comme $N$ n'est pas nul, il possède un idéal
premier associé $\mathfrak q$ d'après le Lemme \ref{algebra-lemma-ass-zero}.
Alors $x \in \mathfrak q$, et $\mathfrak q$ est un idéal premier associé à
$M$ d'après le Lemme \ref{algebra-lemma-ass}.
\end{proof}

\begin{lemma}
\label{algebra-lemma-one-equation-module}
Soient $R$ un anneau local noethérien, $M$ un $R$-module de type fini et
$f \in \mathfrak m$ un élément de l'idéal maximal de $R$. Alors
$$
\dim(\text{Supp}(M/fM)) \leq
\dim(\text{Supp}(M)) \leq
\dim(\text{Supp}(M/fM)) + 1
$$
Si $f$ n'appartient à aucun des idéaux premiers minimaux du support de $M$
(par exemple, si $f$ est un non-diviseur de zéro dans $M$), alors l'inégalité
de droite est une égalité.
\end{lemma}

\begin{proof}
(L'assertion entre parenthèses résulte du Lemme
\ref{algebra-lemma-ass-zero-divisors}.) La première inégalité découle de
$\text{Supp}(M/fM) \subset \text{Supp}(M)$, voir le Lemme
\ref{algebra-lemma-support-quotient}. Pour la seconde, remarquons que
$\text{Supp}(M/fM) = \text{Supp}(M) \cap V(f)$, voir le Lemme
\ref{algebra-lemma-support-quotient}. Il en résulte, par exemple d'après le Lemme
\ref{algebra-lemma-filter-primes-in-support} et les propriétés élémentaires de la
dimension, qu'il suffit de montrer que
$\dim V(\mathfrak p) \leq \dim (V(\mathfrak p) \cap V(f)) + 1$
pour les idéaux premiers $\mathfrak p$ de $R$. C'est une conséquence du
Lemme \ref{algebra-lemma-one-equation}.
Enfin, si $f$ n'est contenu dans aucun idéal premier minimal du support de $M$,
alors toute chaîne d'idéaux premiers de $\text{Supp}(M/fM)$ donne naissance à
une chaîne dans $\text{Supp}(M)$ à laquelle il manque encore au moins une étape
pour être maximale.
\end{proof}

\begin{lemma}
\label{algebra-lemma-ass-functorial}
Soit $\varphi : R \to S$ un morphisme d'anneaux.
Soit $M$ un $S$-module.
Alors $\Spec(\varphi)(\text{Ass}_S(M)) \subset \text{Ass}_R(M)$.
\end{lemma}

\begin{proof}
Si $\mathfrak q \in \text{Ass}_S(M)$, il existe un élément $m$ de $M$ tel que
l'annulateur de $m$ dans $S$ soit $\mathfrak q$. L'annulateur de $m$ dans $R$
est alors $\mathfrak q \cap R$.
\end{proof}

\begin{remark}
\label{algebra-remark-ass-reverse-functorial}
Soit $\varphi : R \to S$ un morphisme d'anneaux.
Soit $M$ un $S$-module.
Il n'est pas toujours vrai que
$\Spec(\varphi)(\text{Ass}_S(M)) \supset \text{Ass}_R(M)$.
Par exemple, considérons le morphisme d'anneaux
$R = k \to S = k[x_1, x_2, x_3, \ldots]/(x_i^2)$ et $M = S$.
Alors $\text{Ass}_R(M)$ n'est pas vide, tandis que $\text{Ass}_S(S)$ est vide.
\end{remark}

\begin{lemma}
\label{algebra-lemma-ass-functorial-Noetherian}
Soit $\varphi : R \to S$ un morphisme d'anneaux.
Soit $M$ un $S$-module. Si $S$ est noethérien, alors
$\Spec(\varphi)(\text{Ass}_S(M)) = \text{Ass}_R(M)$.
\end{lemma}

\begin{proof}
Nous avons déjà vu dans le Lemme \ref{algebra-lemma-ass-functorial} que
$\Spec(\varphi)(\text{Ass}_S(M)) \subset \text{Ass}_R(M)$.
Pour l'inclusion réciproque, choisissons un idéal premier
$\mathfrak p \in \text{Ass}_R(M)$. Soit $m \in M$ un élément tel que
l'annulateur de $m$ dans $R$ est $\mathfrak p$. Soit
$I = \{g \in S \mid gm = 0\}$ l'annulateur de $m$ dans $S$.
Alors $R/\mathfrak p \subset S/I$ est injectif. En combinant les Lemmes
\ref{algebra-lemma-injective-minimal-primes-in-image} et
\ref{algebra-lemma-minimal-prime-image-minimal-prime}, nous voyons qu'il existe un
idéal premier $\mathfrak q \subset S$ minimal au-dessus de $I$ qui s'envoie
sur $\mathfrak p$. D'après la Proposition
\ref{algebra-proposition-minimal-primes-associated-primes}, $\mathfrak q$ est un idéal
premier associé à $S/I$ ; par le Lemme \ref{algebra-lemma-ass}, $\mathfrak q$ est donc
un idéal premier associé à $M$, ce qui conclut.
\end{proof}

\begin{lemma}
\label{algebra-lemma-ass-quotient-ring}
Soit $R$ un anneau.
Soit $I$ un idéal.
Soit $M$ un $R/I$-module.
Via l'injection canonique $\Spec(R/I) \to \Spec(R)$, nous avons
$\text{Ass}_{R/I}(M) = \text{Ass}_R(M)$.
\end{lemma}

\begin{proof}
Omis.
\end{proof}

\begin{lemma}
\label{algebra-lemma-associated-primes-localize}
Soit $R$ un anneau.
Soit $M$ un $R$-module.
Soit $\mathfrak p \subset R$ un idéal premier.
\begin{enumerate}
\item Si $\mathfrak p \in \text{Ass}(M)$, alors
$\mathfrak pR_{\mathfrak p} \in \text{Ass}(M_{\mathfrak p})$.
\item Si $\mathfrak p$ est de type fini, la réciproque est également vraie.
\end{enumerate}
\end{lemma}

\begin{proof}
Si $\mathfrak p \in \text{Ass}(M)$, il existe un élément $m \in M$ dont
l'annulateur est $\mathfrak p$. Comme la localisation est exacte
(Proposition \ref{algebra-proposition-localization-exact}), l'annulateur de $m/1$ dans
$M_{\mathfrak p}$ est $\mathfrak pR_{\mathfrak p}$ ; cela démontre (1).
Supposons $\mathfrak pR_{\mathfrak p} \in \text{Ass}(M_{\mathfrak p})$ et
$\mathfrak p = (f_1, \ldots, f_n)$. Soit $m/g$ un élément de
$M_{\mathfrak p}$ dont l'annulateur est $\mathfrak pR_{\mathfrak p}$.
Cela implique que l'annulateur de $m$ est contenu dans $\mathfrak p$.
Comme $f_i m/g = 0$ dans $M_{\mathfrak p}$, il existe
$g_i \in R$, $g_i \not \in \mathfrak p$, tel que $g_i f_i m = 0$ dans $M$.
En combinant ces faits, nous voyons que l'annulateur de $g_1\ldots g_nm$ est
$\mathfrak p$. Ainsi, $\mathfrak p \in \text{Ass}(M)$.
\end{proof}

\begin{lemma}
\label{algebra-lemma-localize-ass}
Soit $R$ un anneau. Soit $M$ un $R$-module.
Soit $S \subset R$ une partie multiplicative.
Via l'injection canonique $\Spec(S^{-1}R) \to \Spec(R)$, nous avons
\begin{enumerate}
\item $\text{Ass}_R(S^{-1}M) = \text{Ass}_{S^{-1}R}(S^{-1}M)$,
\item
$\text{Ass}_R(M) \cap \Spec(S^{-1}R) \subset \text{Ass}_R(S^{-1}M)$, et
\item si $R$ est noethérien, cette inclusion est une égalité.
\end{enumerate}
\end{lemma}

\begin{proof}
Pour $m \in S^{-1}M$, soient $I \subset R$ et $J \subset S^{-1}R$ les
annulateurs de $m$. Alors $I$ est l'image réciproque de $J$ par le morphisme
$R \to S^{-1}R$, et $J = S^{-1}I$. L'égalité de (1) résulte de la description
du morphisme $\Spec(S^{-1}R) \to \Spec(R)$ donnée dans le Lemme
\ref{algebra-lemma-spec-localization}.
Pour $m \in M$, soit $I \subset R$ l'annulateur de $m$ dans $R$, et soit
$J \subset S^{-1}R$ l'annulateur de $m/1 \in S^{-1}M$.
Nous avons $J = S^{-1}I$, ce qui implique (2). Dans le cas noethérien,
l'égalité résulte du Lemme \ref{algebra-lemma-associated-primes-localize}, car pour
$\mathfrak p \in R$, $S \cap \mathfrak p = \emptyset$, nous avons
$M_{\mathfrak p} = (S^{-1}M)_{S^{-1}\mathfrak p}$.
\end{proof}

\begin{lemma}
\label{algebra-lemma-localize-ass-nonzero-divisors}
Soit $R$ un anneau. Soit $M$ un $R$-module.
Soit $S \subset R$ une partie multiplicative.
Supposons que tout $s \in S$ soit un non-diviseur de zéro dans $M$.
Alors
$$
\text{Ass}_R(M) = \text{Ass}_R(S^{-1}M).
$$
\end{lemma}

\begin{proof}
Puisque $M \subset S^{-1}M$ par hypothèse, le Lemme \ref{algebra-lemma-ass} donne
l'inclusion $\text{Ass}(M) \subset \text{Ass}(S^{-1}M)$.
Réciproquement, supposons que $n/s \in S^{-1}M$ soit un élément dont
l'annulateur est un idéal premier $\mathfrak p$. Alors l'annulateur de
$n \in M$ est lui aussi $\mathfrak p$.
\end{proof}

\begin{lemma}
\label{algebra-lemma-ideal-nonzerodivisor}
Soit $R$ un anneau local noethérien d'idéal maximal $\mathfrak m$.
Soit $I \subset \mathfrak m$ un idéal. Soit $M$ un $R$-module de type fini.
Les conditions suivantes sont équivalentes :
\begin{enumerate}
\item Il existe $x \in I$ qui n'est pas un diviseur de zéro dans $M$.
\item Nous avons $I \not \subset \mathfrak q$ pour tout
$\mathfrak q \in \text{Ass}(M)$.
\end{enumerate}
\end{lemma}

\begin{proof}
S'il existe un non-diviseur de zéro $x$ dans $I$, alors $x$ ne peut
manifestement appartenir à aucun idéal premier associé à $M$.
Réciproquement, supposons $I \not \subset \mathfrak q$ pour tout
$\mathfrak q \in \text{Ass}(M)$. Nous pouvons alors choisir
$x \in I$, $x \not \in \mathfrak q$ pour tout
$\mathfrak q \in \text{Ass}(M)$, d'après les Lemmes
\ref{algebra-lemma-finite-ass} et \ref{algebra-lemma-silly}.
Le Lemme \ref{algebra-lemma-ass-zero-divisors} montre que l'élément $x$ n'est pas un
diviseur de zéro dans $M$.
\end{proof}

\begin{lemma}
\label{algebra-lemma-zero-at-ass-zero}
Soit $R$ un anneau. Soit $M$ un $R$-module. Si $R$ est noethérien,
l'application
$$
M
\longrightarrow
\prod\nolimits_{\mathfrak p \in \text{Ass}(M)} M_{\mathfrak p}
$$
est injective.
\end{lemma}

\begin{proof}
Soit $x \in M$ un élément du noyau de l'application.
Alors, si $\mathfrak p$ est un idéal premier associé à $Rx \subset M$, nous
voyons d'une part que $\mathfrak p \in \text{Ass}(M)$
(Lemme \ref{algebra-lemma-ass}) et d'autre part que
$(Rx)_{\mathfrak p} \subset M_{\mathfrak p}$ n'est pas nul.
Cette contradiction montre que $\text{Ass}(Rx) = \emptyset$.
Ainsi, $Rx = 0$ d'après le Lemme \ref{algebra-lemma-ass-zero}.
\end{proof}

\noindent
Ce lemme devrait probablement être placé ailleurs.

\begin{lemma}
\label{algebra-lemma-dim-not-zero-exists-nonzerodivisor-nonunit}
Soit $k$ un corps. Soit $S$ une $k$-algèbre de type fini.
Si $\dim(S) > 0$, alors il existe un élément $f \in S$ qui est un
non-diviseur de zéro et n'est pas inversible.
\end{lemma}

\begin{proof}
D'après le Lemme \ref{algebra-lemma-finite-ass}, l'anneau $S$ ne possède qu'un nombre
fini d'idéaux premiers associés. D'après le Lemme
\ref{algebra-lemma-finite-type-algebra-finite-nr-primes}, l'anneau $S$ possède une
infinité d'idéaux maximaux. Nous pouvons donc choisir un idéal maximal
$\mathfrak m \subset S$ qui ne soit pas un idéal premier associé à $S$.
Par évitement des idéaux premiers (Lemme \ref{algebra-lemma-silly}), nous pouvons
choisir un élément non nul $f \in \mathfrak m$ qui n'appartient à aucun des
idéaux premiers associés à $S$. D'après le Lemme
\ref{algebra-lemma-ass-zero-divisors}, l'élément $f$ est un non-diviseur de zéro ; et,
comme $f \in \mathfrak m$, nous voyons que $f$ n'est pas une unité.
\end{proof}

\section{Puissances symboliques}
\label{algebra-section-symbolic-power}

\noindent
Voici la définition.

\begin{definition}
\label{algebra-definition-symbolic-power}
Soit $R$ un anneau. Soit $\mathfrak p$ un idéal premier. Pour $n \geq 0$, la
$n$-ième {\it puissance symbolique} de $\mathfrak p$ est l'idéal
$\mathfrak p^{(n)} = \Ker(R \to R_\mathfrak p/\mathfrak p^nR_\mathfrak p)$.
\end{definition}

\noindent
Remarquons que $\mathfrak p^n \subset \mathfrak p^{(n)}$, mais qu'il n'y a
pas toujours égalité.

\begin{lemma}
\label{algebra-lemma-symbolic-power-associated}
Soit $R$ un anneau noethérien.
Soit $\mathfrak p$ un idéal premier.
Soit $n > 0$. Alors $\text{Ass}(R/\mathfrak p^{(n)}) = \{\mathfrak p\}$.
\end{lemma}

\begin{proof}
Si $\mathfrak q$ est un idéal premier associé à $R/\mathfrak p^{(n)}$, alors
$\mathfrak p \subset \mathfrak q$ manifestement.
D'autre part, tout élément $x \in R$, $x \not \in \mathfrak p$, est un
non-diviseur de zéro dans $R/\mathfrak p^{(n)}$.
En effet, si $y \in R$ et
$xy \in \mathfrak p^{(n)} = R \cap \mathfrak p^nR_{\mathfrak p}$, alors
$y \in \mathfrak p^nR_{\mathfrak p}$, d'où
$y \in \mathfrak p^{(n)}$. Le lemme en découle.
\end{proof}

\begin{lemma}
\label{algebra-lemma-symbolic-power-flat-extension}
Soit $R \to S$ un morphisme d'anneaux plat. Soit $\mathfrak p \subset R$ un
idéal premier tel que $\mathfrak q = \mathfrak p S$ soit un idéal premier de
$S$. Alors $\mathfrak p^{(n)} S = \mathfrak q^{(n)}$.
\end{lemma}

\begin{proof}
Puisque
$\mathfrak p^{(n)} = \Ker(R \to R_\mathfrak p/\mathfrak p^nR_\mathfrak p)$,
la platitude montre que $\mathfrak p^{(n)} S$ est le noyau du morphisme
$S \to S_\mathfrak p/\mathfrak p^nS_\mathfrak p$. D'autre part,
$\mathfrak q^{(n)}$ est le noyau du morphisme
$S \to S_\mathfrak q/\mathfrak q^nS_\mathfrak q =
S_\mathfrak q/\mathfrak p^nS_\mathfrak q$. Il suffit donc de montrer que
$$
S_\mathfrak p/\mathfrak p^nS_\mathfrak p
\longrightarrow
S_\mathfrak q/\mathfrak p^nS_\mathfrak q
$$
est injectif. Remarquons que le module de droite est le localisé du module de
gauche par les éléments $f \in S$, $f \not \in \mathfrak q$.
Il suffit donc de montrer que ces éléments sont des non-diviseurs de zéro dans
$S_\mathfrak p/\mathfrak p^nS_\mathfrak p$. Par platitude, le module
$S_\mathfrak p/\mathfrak p^nS_\mathfrak p$ possède une filtration finie dont
les sous-quotients sont
$$
\mathfrak p^iS_\mathfrak p/\mathfrak p^{i + 1}S_\mathfrak p
\cong \mathfrak p^iR_\mathfrak p/\mathfrak p^{i + 1}R_\mathfrak p
\otimes_{R_\mathfrak p} S_\mathfrak p \cong
V \otimes_{\kappa(\mathfrak p)} (S/\mathfrak q)_\mathfrak p
$$
où $V$ est un espace vectoriel sur $\kappa(\mathfrak p)$. Ainsi, $f$ agit de
manière inversible, comme voulu.
\end{proof}



\section{Assassin relatif}
\label{algebra-section-relative-assassin}

\noindent
Discussion des assassins relatifs. Soit $R \to S$ un morphisme d'anneaux.
Soit $N$ un $S$-module. Dans cette situation, nous pouvons introduire les
ensembles suivants d'idéaux premiers $\mathfrak q$ de $S$ :
\begin{enumerate}
\item $A$ : en posant $\mathfrak p = R \cap \mathfrak q$, on a
$\mathfrak q \in \text{Ass}_S(N \otimes_R \kappa(\mathfrak p))$,
\item $A'$ : en posant $\mathfrak p = R \cap \mathfrak q$, l'idéal
$\mathfrak q$ appartient à l'image de
$\text{Ass}_{S \otimes \kappa(\mathfrak p)}(N \otimes_R \kappa(\mathfrak p))$
par le morphisme canonique
$\Spec(S \otimes_R \kappa(\mathfrak p)) \to \Spec(S)$,
\item $A_{fin}$ : en posant $\mathfrak p = R \cap \mathfrak q$, on a
$\mathfrak q \in \text{Ass}_S(N/\mathfrak pN)$,
\item $A'_{fin}$ : il existe un idéal premier $\mathfrak p' \subset R$ tel que
$\mathfrak q \in \text{Ass}_S(N/\mathfrak p'N)$,
\item $B$ : il existe un $R$-module $M$ tel que
$\mathfrak q \in \text{Ass}_S(N \otimes_R M)$, et
\item $B_{fin}$ : il existe un $R$-module de type fini $M$ tel que
$\mathfrak q \in \text{Ass}_S(N \otimes_R M)$.
\end{enumerate}
Déterminons quelques-unes des relations entre ces ensembles.

\begin{lemma}
\label{algebra-lemma-compare-relative-assassins}
Soit $R \to S$ un morphisme d'anneaux. Soit $N$ un $S$-module.
Soient $A$, $A'$, $A_{fin}$, $B$ et $B_{fin}$ les sous-ensembles de
$\Spec(S)$ introduits ci-dessus.
\begin{enumerate}
\item Nous avons toujours $A = A'$.
\item Nous avons toujours $A_{fin} \subset A$,
$B_{fin} \subset B$, $A_{fin} \subset A'_{fin} \subset B_{fin}$
et $A \subset B$.
\item Si $S$ est noethérien, alors $A = A_{fin}$ et $B = B_{fin}$.
\item Si $N$ est plat sur $R$, alors $A = A_{fin} = A'_{fin}$ et $B = B_{fin}$.
\item Si $R$ est noethérien et si $N$ est plat sur $R$, alors tous ces
ensembles sont égaux, c'est-à-dire
$A = A' = A_{fin} = A'_{fin} = B = B_{fin}$.
\end{enumerate}
\end{lemma}

\begin{proof}
Certains arguments de la démonstration seront répétés dans les démonstrations
de lemmes ultérieurs, qui sont plus précis que celui-ci (car ils portent sur
un module $M$ donné ou sur un idéal premier $\mathfrak p$ donné, et non sur
l'ensemble de tous ces objets).

\medskip\noindent
Démonstration de (1). Soit $\mathfrak p$ un idéal premier de $R$. Nous avons
alors
$$
\text{Ass}_S(N \otimes_R \kappa(\mathfrak p)) =
\text{Ass}_{S/\mathfrak pS}(N \otimes_R \kappa(\mathfrak p)) =
\text{Ass}_{S \otimes_R \kappa(\mathfrak p)}(N \otimes_R \kappa(\mathfrak p))
$$
la première égalité d'après le
Lemme \ref{algebra-lemma-ass-quotient-ring}
et la seconde d'après le
Lemme \ref{algebra-lemma-localize-ass}, partie (1). Cela montre que $A = A'$.
L'inclusion $A_{fin} \subset A'_{fin}$ est claire.

\medskip\noindent
Démonstration de (2). Chacune des inclusions découle immédiatement des
définitions, à l'exception peut-être de $A_{fin} \subset A$, qui résulte du
Lemme \ref{algebra-lemma-localize-ass}
et du fait que nous imposons $\mathfrak p = R \cap \mathfrak q$ dans la
définition de $A_{fin}$.

\medskip\noindent
Démonstration de (3). L'égalité $A = A_{fin}$ découle du
Lemme \ref{algebra-lemma-localize-ass}, partie (3), si $S$ est noethérien.
Soit $\mathfrak q = (g_1, \ldots, g_m)$ un idéal premier de type fini de $S$.
Soit $z \in N \otimes_R M$ un élément dont l'annulateur est $\mathfrak q$.
Nous pouvons choisir un sous-module de type fini $M' \subset M$ tel que $z$
soit l'image d'un élément $z' \in N \otimes_R M'$. Alors
$\text{Ann}_S(z') \subset \mathfrak q = \text{Ann}_S(z)$.
Puisque $N \otimes_R -$ commute aux colimites et que $M$ est la colimite
filtrante de ses sous-$R$-modules de type fini, nous pouvons trouver
$M' \subset M'' \subset M$ tel que l'image
$z'' \in N \otimes_R M''$ soit annulée par $g_1, \ldots, g_m$. Ainsi,
$\text{Ann}_S(z'') = \mathfrak q$. Cela montre que $B = B_{fin}$ lorsque $S$
est noethérien.

\medskip\noindent
Démonstration de (4). Si $N$ est plat, le foncteur $N \otimes_R -$ est exact.
En particulier, si $M' \subset M$, alors
$N \otimes_R M' \subset N \otimes_R M$. Ainsi, si
$z \in N \otimes_R M$ est un élément dont l'annulateur
$\mathfrak q = \text{Ann}_S(z)$ est premier, nous pouvons choisir un
sous-$R$-module de type fini quelconque $M' \subset M$ tel que
$z \in N \otimes_R M'$ ; l'annulateur de $z$ comme élément de
$N \otimes_R M'$ est alors égal à $\mathfrak q$. Ainsi, $B = B_{fin}$.
Soit $\mathfrak p'$ un idéal premier de $R$, et soit $\mathfrak q$ un idéal
premier de $S$ associé à $N/\mathfrak p'N$. Cela implique que
$\mathfrak p'S \subset \mathfrak q$. Comme $N$ est plat sur $R$, le module
$N/\mathfrak p'N$ est plat sur l'anneau intègre $R/\mathfrak p'$.
Par conséquent, tout élément non nul de $R/\mathfrak p'$ est un
non-diviseur de zéro dans $N/\mathfrak p'$. L'image d'aucun de ces éléments
ne peut donc appartenir à $\mathfrak q$, et nous concluons que
$\mathfrak p' = R \cap \mathfrak q$. Ainsi, $A_{fin} = A'_{fin}$. Enfin, le
Lemme \ref{algebra-lemma-localize-ass-nonzero-divisors}
donne
$\text{Ass}_S(N/\mathfrak p'N) =
\text{Ass}_S(N \otimes_R \kappa(\mathfrak p'))$, c'est-à-dire
$A'_{fin} = A$.

\medskip\noindent
Démonstration de (5). Il suffit de prouver $A'_{fin} = B_{fin}$, puisque les
autres égalités ont été établies en (4). Pour cela, soit $M$ un $R$-module de
type fini. D'après le
Lemme \ref{algebra-lemma-filter-Noetherian-module},
il existe une filtration par des sous-$R$-modules
$$
0 = M_0 \subset M_1 \subset \ldots \subset M_n = M
$$
telle que chaque quotient $M_i/M_{i-1}$ soit isomorphe à
$R/\mathfrak p_i$ pour un certain idéal premier $\mathfrak p_i$ de $R$.
Puisque $N$ est plat, nous obtenons une filtration par des sous-$S$-modules
$$
0 = N \otimes_R M_0 \subset N \otimes_R M_1 \subset \ldots \subset
N \otimes_R M_n = N \otimes_R M
$$
dont chaque sous-quotient est isomorphe à $N/\mathfrak p_iN$. D'après le
Lemme \ref{algebra-lemma-ass},
nous en déduisons que
$\text{Ass}_S(N \otimes_R M) \subset \bigcup \text{Ass}_S(N/\mathfrak p_iN)$.
Ainsi, $B_{fin} \subset A'_{fin}$. L'inclusion réciproque faisant partie de
(2), l'égalité est démontrée.
\end{proof}

\noindent
Nous définissons l'assassin de $N$ relativement à $S/R$ comme
l'ensemble $A = A'$ ci-dessus. Pour motiver cette définition, remarquons
qu'il ne dépend que des modules fibres $N \otimes_R \kappa(\mathfrak p)$ sur
les anneaux fibres. Comme dans le cas de l'assassin d'un module, nous
avertissons le lecteur que cette notion est surtout pertinente lorsque les
anneaux fibres $S \otimes_R \kappa(\mathfrak p)$ sont noethériens, par exemple
lorsque $R \to S$ est de type fini.

\begin{definition}
\label{algebra-definition-relative-assassin}
Soit $R \to S$ un morphisme d'anneaux. Soit $N$ un $S$-module.
L'{\it assassin de $N$ relativement à $S/R$} est l'ensemble
$$
\text{Ass}_{S/R}(N)
=
\{ \mathfrak q \subset S \mid
\mathfrak q \in \text{Ass}_S(N \otimes_R \kappa(\mathfrak p))
\text{ avec }\mathfrak p = R \cap \mathfrak q\}.
$$
C'est l'ensemble noté $A$ dans le
Lemme \ref{algebra-lemma-compare-relative-assassins}.
\end{definition}

\noindent
L'esprit des quelques résultats suivants est qu'ils portent sur l'assassin
relatif, même si cela n'apparaît pas immédiatement.

\begin{lemma}
\label{algebra-lemma-bourbaki}
Soit $R \to S$ un morphisme d'anneaux.
Soit $M$ un $R$-module et soit $N$ un $S$-module.
Si $N$ est plat comme $R$-module, alors
$$
\text{Ass}_S(M \otimes_R N)
\supset
\bigcup\nolimits_{\mathfrak p \in \text{Ass}_R(M)} \text{Ass}_S(N/\mathfrak pN)
$$
et, si $R$ est noethérien, il y a égalité.
\end{lemma}

\begin{proof}
Si $\mathfrak p \in \text{Ass}_R(M)$, il existe une injection
$R/\mathfrak p \to M$. Comme $N$ est plat sur $R$, nous obtenons une
injection $R/\mathfrak p \otimes_R N \to M \otimes_R N$. Puisque
$R/\mathfrak p \otimes_R N = N/\mathfrak pN$, nous concluons que
$\text{Ass}_S(N/\mathfrak pN) \subset \text{Ass}_S(M \otimes_R N)$, voir le
Lemme \ref{algebra-lemma-ass}. Le membre de droite est donc contenu dans le membre de
gauche.

\medskip\noindent
Écrivons $M = \bigcup M_\lambda$ comme la réunion de ses sous-$R$-modules de
type fini. Alors $N \otimes_R M = \bigcup N \otimes_R M_\lambda$ également
(car $N$ est plat sur $R$). La définition des idéaux premiers associés donne
$\text{Ass}_S(N \otimes_R M) = \bigcup \text{Ass}_S(N \otimes_R M_\lambda)$
et $\text{Ass}_R(M) = \bigcup \text{Ass}(M_\lambda)$. Nous pouvons donc
supposer que $M$ est de type fini.

\medskip\noindent
Soit $\mathfrak q \in \text{Ass}_S(M \otimes_R N)$, et supposons que $R$ soit
noethérien et que $M$ soit un $R$-module de type fini. Pour achever la
démonstration, nous devons montrer que $\mathfrak q$ appartient au membre de
droite. Remarquons d'abord que
$\mathfrak qS_{\mathfrak q} \in
\text{Ass}_{S_{\mathfrak q}}((M \otimes_R N)_{\mathfrak q})$,
voir le Lemme \ref{algebra-lemma-associated-primes-localize}.
Soit $\mathfrak p$ l'idéal premier correspondant de $R$.
Remarquons que
$$
(M \otimes_R N)_{\mathfrak q} = M \otimes_R N_{\mathfrak q}
= M_{\mathfrak p} \otimes_{R_{\mathfrak p}} N_{\mathfrak q}
$$
Si
$\mathfrak pR_{\mathfrak p} \not \in
\text{Ass}_{R_{\mathfrak p}}(M_{\mathfrak p})$,
il existe un élément $x \in \mathfrak pR_{\mathfrak p}$ qui est un
non-diviseur de zéro dans $M_{\mathfrak p}$ (voir le
Lemme \ref{algebra-lemma-ideal-nonzerodivisor}). Comme
$N_{\mathfrak q}$ est plat sur $R_{\mathfrak p}$, l'image de $x$ dans
$\mathfrak qS_{\mathfrak q}$ est un non-diviseur de zéro dans
$(M \otimes_R N)_{\mathfrak q}$. Cela contredit l'hypothèse
$\mathfrak qS_{\mathfrak q} \in \text{Ass}_S((M \otimes_R N)_{\mathfrak q})$.
Nous concluons donc que $\mathfrak p$ est l'un des idéaux premiers associés à
$M$.

\medskip\noindent
Poursuivant l'argument, choisissons une filtration
$$
0 = M_0 \subset M_1 \subset \ldots \subset M_n = M
$$
telle que chaque quotient $M_i/M_{i-1}$ soit isomorphe à
$R/\mathfrak p_i$ pour un certain idéal premier $\mathfrak p_i$ de $R$, voir
le Lemme \ref{algebra-lemma-filter-Noetherian-module}.
(D'après le Lemme \ref{algebra-lemma-ass-filter}, nous avons
$\mathfrak p_i = \mathfrak p$ pour au moins un $i$.) Nous obtenons une
filtration
$$
0 = M_0 \otimes_R N \subset M_1 \otimes_R N \subset \ldots
\subset M_n \otimes_R N = M \otimes_R N
$$
dont les sous-quotients sont isomorphes à $N/\mathfrak p_iN$. Si
$\mathfrak p_i \not = \mathfrak p$, alors $\mathfrak q$ ne peut être associé
au module $N/\mathfrak p_iN$, d'après le résultat du paragraphe précédent
(car $\text{Ass}_R(R/\mathfrak p_i) = \{\mathfrak p_i\}$). Nous concluons
donc que $\mathfrak q$ est associé à $N/\mathfrak pN$, comme voulu.
\end{proof}

\begin{lemma}
\label{algebra-lemma-post-bourbaki}
Soit $R \to S$ un morphisme d'anneaux.
Soit $N$ un $S$-module.
Supposons que $N$ soit plat comme $R$-module et que $R$ soit un anneau
intègre de corps des fractions $K$. Alors
$$
\text{Ass}_S(N) =
\text{Ass}_S(N \otimes_R K) =
\text{Ass}_{S \otimes_R K}(N \otimes_R K)
$$
par l'inclusion canonique
$\Spec(S \otimes_R K) \subset \Spec(S)$.
\end{lemma}

\begin{proof}
Remarquons que $S \otimes_R K = (R \setminus \{0\})^{-1}S$ et
$N \otimes_R K = (R \setminus \{0\})^{-1}N$.
Pour tout élément non nul $x \in R$, la multiplication par $x$ dans $N$ est
injective, puisque $N$ est plat sur $R$. Le lemme résulte donc du
Lemme \ref{algebra-lemma-localize-ass-nonzero-divisors}
combiné avec le
Lemme \ref{algebra-lemma-localize-ass}, partie (1).
\end{proof}

\begin{lemma}
\label{algebra-lemma-bourbaki-fibres}
Soit $R \to S$ un morphisme d'anneaux.
Soit $M$ un $R$-module et soit $N$ un $S$-module.
Supposons que $N$ soit plat comme $R$-module. Alors
$$
\text{Ass}_S(M \otimes_R N)
\supset
\bigcup\nolimits_{\mathfrak p \in \text{Ass}_R(M)}
\text{Ass}_{S \otimes_R \kappa(\mathfrak p)}(N \otimes_R \kappa(\mathfrak p))
$$
où nous utilisons la
Remarque \ref{algebra-remark-fundamental-diagram}
pour considérer les spectres des anneaux fibres comme des sous-ensembles de
$\Spec(S)$. Si $R$ est noethérien, cette inclusion est une égalité.
\end{lemma}

\begin{proof}
Ceci équivaut au
Lemme \ref{algebra-lemma-bourbaki}
d'après les
Lemmes \ref{algebra-lemma-ass-quotient-ring},
\ref{algebra-lemma-flat-base-change} et
\ref{algebra-lemma-post-bourbaki}.
\end{proof}

\begin{remark}
\label{algebra-remark-bourbaki}
Soit $R \to S$ un morphisme d'anneaux. Soit $N$ un $S$-module.
Soit $\mathfrak p$ un idéal premier de $R$. Alors
$$
\text{Ass}_S(N \otimes_R \kappa(\mathfrak p)) =
\text{Ass}_{S/\mathfrak pS}(N \otimes_R \kappa(\mathfrak p)) =
\text{Ass}_{S \otimes_R \kappa(\mathfrak p)}(N \otimes_R \kappa(\mathfrak p)).
$$
La première égalité résulte du
Lemme \ref{algebra-lemma-ass-quotient-ring}
et la seconde du
Lemme \ref{algebra-lemma-localize-ass}, partie (1).
\end{remark}
















\section{Idéaux premiers faiblement associés}
\label{algebra-section-weakly-ass}

\noindent
Il s'agit d'une variante de la notion d'idéal premier associé qui est utile
pour les anneaux non noethériens et les modules qui ne sont pas de type fini.

\begin{definition}
\label{algebra-definition-weakly-associated}
Soit $R$ un anneau. Soit $M$ un $R$-module.
Un idéal premier $\mathfrak p$ de $R$ est {\it faiblement associé} à $M$
s'il existe un élément $m \in M$ tel que $\mathfrak p$ soit minimal
parmi les idéaux premiers contenant l'annulateur
$\text{Ann}(m) = \{f \in R \mid fm = 0\}$.
L'ensemble de tous ces idéaux premiers est noté $\text{WeakAss}_R(M)$
ou $\text{WeakAss}(M)$.
\end{definition}

\noindent
Ainsi, tout idéal premier associé est faiblement associé.
Voici une caractérisation à l'aide de la localisation en l'idéal premier.

\begin{lemma}
\label{algebra-lemma-weakly-ass-local}
Soit $R$ un anneau. Soit $M$ un $R$-module.
Soit $\mathfrak p$ un idéal premier de $R$.
Les assertions suivantes sont équivalentes :
\begin{enumerate}
\item $\mathfrak p$ est faiblement associé à $M$,
\item $\mathfrak pR_{\mathfrak p}$ est faiblement associé à $M_{\mathfrak p}$,
et
\item $M_{\mathfrak p}$ contient un élément dont
l'annulateur a pour radical $\mathfrak pR_{\mathfrak p}$.
\end{enumerate}
\end{lemma}

\begin{proof}
Supposons (1). Il existe alors un élément $m \in M$ tel que
$\mathfrak p$ soit minimal parmi les idéaux premiers contenant l'annulateur
$I = \{x \in R \mid xm = 0\}$ de $m$. Comme la localisation est exacte,
l'annulateur de $m$ dans $M_{\mathfrak p}$ est $I_{\mathfrak p}$.
Ainsi, $\mathfrak pR_{\mathfrak p}$ est un idéal premier minimal de
$R_{\mathfrak p}$ contenant l'annulateur $I_{\mathfrak p}$
de $m$ dans $M_{\mathfrak p}$. Ceci montre (2), ainsi que (3),
car $\sqrt{I_{\mathfrak p}} = \mathfrak pR_{\mathfrak p}$.

\medskip\noindent
En appliquant l'implication (1) $\Rightarrow$ (3) à $M_{\mathfrak p}$
sur $R_{\mathfrak p}$, on voit que (2) $\Rightarrow$ (3).

\medskip\noindent
Enfin, supposons (3). Il existe donc un élément
$m/f \in M_{\mathfrak p}$ dont l'annulateur a pour radical
$\mathfrak pR_{\mathfrak p}$. L'annulateur
$I = \{x \in R \mid xm = 0\}$ de $m$ dans $M$ vérifie alors
$\sqrt{I_{\mathfrak p}} = \mathfrak pR_{\mathfrak p}$. Cela signifie
clairement que $\mathfrak p$ contient $I$ et est minimal parmi les
idéaux premiers contenant $I$, c'est-à-dire que (1) est vérifiée.
\end{proof}

\begin{lemma}
\label{algebra-lemma-reduced-weakly-ass-minimal}
Pour un anneau réduit, les idéaux premiers faiblement associés à l'anneau sont
les idéaux premiers minimaux.
\end{lemma}

\begin{proof}
Soit $(R, \mathfrak m)$ un anneau local réduit.
Supposons que $x \in R$ soit un élément dont l'annulateur
a pour radical $\mathfrak m$. Si $\mathfrak m \not = 0$, alors $x$
ne peut pas être inversible, donc $x \in \mathfrak m$. En particulier, $x^{1 + n} = 0$
pour un certain $n \geq 0$. Ainsi $x = 0$, ce qui contredit l'hypothèse
selon laquelle l'annulateur de $x$ est contenu dans $\mathfrak m$.
On voit donc que $\mathfrak m = 0$, c'est-à-dire que $R$ est un corps.
Le Lemme \ref{algebra-lemma-weakly-ass-local} donne alors
l'assertion du lemme.
\end{proof}

\begin{lemma}
\label{algebra-lemma-weakly-ass}
Soit $R$ un anneau.
Soit $0 \to M' \to M \to M'' \to 0$ une suite exacte courte
de $R$-modules.
Alors $\text{WeakAss}(M') \subset \text{WeakAss}(M)$ et
$\text{WeakAss}(M) \subset \text{WeakAss}(M') \cup \text{WeakAss}(M'')$.
\end{lemma}

\begin{proof}
Nous utiliserons la caractérisation des idéaux premiers faiblement associés du
Lemme \ref{algebra-lemma-weakly-ass-local}.
Soit $\mathfrak p$ un idéal premier de $R$. L'exactitude de la localisation donne
la suite exacte courte
$0 \to M'_{\mathfrak p} \to M_{\mathfrak p} \to M''_{\mathfrak p} \to 0$.
Supposons que $m \in M_{\mathfrak p}$ soit un élément dont l'annulateur
a pour radical $\mathfrak pR_{\mathfrak p}$. Alors, ou bien l'image $\overline{m}$
de $m$ dans $M''_{\mathfrak p}$ est nulle et $m \in M'_{\mathfrak p}$, ou bien le
radical de l'annulateur de $\overline{m}$ est $\mathfrak pR_{\mathfrak p}$.
Ceci montre que
$\text{WeakAss}(M) \subset \text{WeakAss}(M') \cup \text{WeakAss}(M'')$.
L'inclusion $\text{WeakAss}(M') \subset \text{WeakAss}(M)$ résulte immédiatement
des définitions.
\end{proof}

\begin{lemma}
\label{algebra-lemma-weakly-ass-zero}
\begin{slogan}
Tout module non nul possède un idéal premier faiblement associé.
\end{slogan}
Soit $R$ un anneau. Soit $M$ un $R$-module. Alors
$$
M = (0) \Leftrightarrow \text{WeakAss}(M) = \emptyset
$$
\end{lemma}

\begin{proof}
Si $M = (0)$, alors $\text{WeakAss}(M) = \emptyset$ par définition.
Réciproquement, supposons $M \not = 0$. Choisissons un élément non nul $m \in M$.
Notons $I = \{x \in R \mid xm = 0\}$ l'annulateur de $m$.
Alors $R/I \subset M$. Ainsi $\text{WeakAss}(R/I) \subset \text{WeakAss}(M)$ par le
Lemme \ref{algebra-lemma-weakly-ass}.
Or, puisque $I \not = R$, l'ensemble $V(I) = \Spec(R/I)$ contient un idéal premier
minimal ; voir les
Lemmes \ref{algebra-lemma-Zariski-topology} et
\ref{algebra-lemma-spec-closed},
ce qui conclut.
\end{proof}

\begin{lemma}
\label{algebra-lemma-weakly-ass-support}
Soit $R$ un anneau. Soit $M$ un $R$-module. Alors
$$
\text{Ass}(M) \subset \text{WeakAss}(M) \subset \text{Supp}(M).
$$
\end{lemma}

\begin{proof}
La première inclusion résulte immédiatement des définitions.
Si $\mathfrak p \in \text{WeakAss}(M)$, alors, d'après le
Lemme \ref{algebra-lemma-weakly-ass-local},
on a $M_{\mathfrak p} \not = 0$, donc $\mathfrak p \in \text{Supp}(M)$.
\end{proof}

\begin{lemma}
\label{algebra-lemma-weakly-ass-zero-divisors}
Soit $R$ un anneau.
Soit $M$ un $R$-module.
La réunion $\bigcup_{\mathfrak q \in \text{WeakAss}(M)} \mathfrak q$
est l'ensemble des éléments de $R$ qui sont des diviseurs de zéro sur $M$.
\end{lemma}

\begin{proof}
Supposons $f \in \mathfrak q \in \text{WeakAss}(M)$.
Il existe un élément $m \in M$ tel que
$\mathfrak q$ soit minimal au-dessus de $I = \{x \in R \mid xm = 0\}$.
Il existe donc $g \in R$, avec $g \not \in \mathfrak q$, et $n > 0$
tels que $f^ngm = 0$. Remarquons que $gm \not = 0$ puisque $g \not \in I$.
Si l'on choisit $n$ minimal avec cette propriété, alors $f (f^{n - 1}gm) = 0$
et $f^{n - 1}gm \not = 0$ ; ainsi $f$ est un diviseur de zéro sur $M$.
Réciproquement, supposons que $f \in R$ soit un diviseur de zéro sur $M$.
Considérons le sous-module $N = \{m \in M \mid fm = 0\}$.
Comme $N$ n'est pas nul, il possède un idéal premier faiblement associé $\mathfrak q$ par le
Lemme \ref{algebra-lemma-weakly-ass-zero}.
On a clairement $f \in \mathfrak q$ et, d'après le
Lemme \ref{algebra-lemma-weakly-ass},
$\mathfrak q$ est un idéal premier faiblement associé à $M$.
\end{proof}

\begin{lemma}
\label{algebra-lemma-weakly-ass-minimal-prime-support}
Soit $R$ un anneau.
Soit $M$ un $R$-module.
Tout $\mathfrak p \in \text{Supp}(M)$ qui est minimal parmi les éléments
de $\text{Supp}(M)$ appartient à $\text{WeakAss}(M)$.
\end{lemma}

\begin{proof}
Remarquons que $\text{Supp}(M_{\mathfrak p}) = \{\mathfrak pR_{\mathfrak p}\}$
dans $\Spec(R_{\mathfrak p})$. En particulier, $M_{\mathfrak p}$
est non nul, et donc $\text{WeakAss}(M_{\mathfrak p}) \not = \emptyset$ d'après le
Lemme \ref{algebra-lemma-weakly-ass-zero}.
Puisque $\text{WeakAss}(M_{\mathfrak p}) \subset \text{Supp}(M_{\mathfrak p})$
d'après le
Lemme \ref{algebra-lemma-weakly-ass-support},
on en déduit que
$\text{WeakAss}(M_{\mathfrak p}) = \{\mathfrak pR_{\mathfrak p}\}$,
d'où $\mathfrak p \in \text{WeakAss}(M)$ par le
Lemme \ref{algebra-lemma-weakly-ass-local}.
\end{proof}

\begin{lemma}
\label{algebra-lemma-ass-weakly-ass}
Soit $R$ un anneau. Soit $M$ un $R$-module.
Soit $\mathfrak p$ un idéal premier de $R$ qui est de type fini.
Alors
$$
\mathfrak p \in \text{Ass}(M) \Leftrightarrow
\mathfrak p \in \text{WeakAss}(M).
$$
En particulier, si $R$ est noethérien, alors $\text{Ass}(M) = \text{WeakAss}(M)$.
\end{lemma}

\begin{proof}
Écrivons $\mathfrak p = (g_1, \ldots, g_n)$ avec $g_i \in R$.
Il suffit de démontrer l'implication « $\Leftarrow$ », puisque l'autre
implication est toujours vraie ; voir le
Lemme \ref{algebra-lemma-weakly-ass-support}.
Supposons $\mathfrak p \in \text{WeakAss}(M)$.
D'après le
Lemme \ref{algebra-lemma-weakly-ass-local},
il existe un élément $m \in M_{\mathfrak p}$ tel que
$I = \{x \in R_{\mathfrak p} \mid xm = 0\}$ ait pour radical
$\mathfrak pR_{\mathfrak p}$. Pour tout $i$, il existe donc
un plus petit $e_i > 0$ tel que $g_i^{e_i}m = 0$ dans $M_{\mathfrak p}$.
Si $e_i > 1$ pour un certain $i$, on peut remplacer $m$ par
$g_i^{e_i - 1} m \not = 0$ et diminuer $\sum e_i$.
On peut donc supposer que l'annulateur de $m \in M_{\mathfrak p}$ est
$(g_1, \ldots, g_n)R_{\mathfrak p} = \mathfrak p R_{\mathfrak p}$. Le
Lemme \ref{algebra-lemma-associated-primes-localize}
montre alors que $\mathfrak p \in \text{Ass}(M)$.
\end{proof}

\begin{remark}
\label{algebra-remark-weakly-ass-not-functorial}
Soit $\varphi : R \to S$ un homomorphisme d'anneaux. Soit $M$ un $S$-module.
On n'a pas toujours
$\Spec(\varphi)(\text{WeakAss}_S(M)) \subset \text{WeakAss}_R(M)$
contrairement au cas des idéaux premiers associés (voir le
Lemme \ref{algebra-lemma-ass-functorial}).
Un exemple est donné par l'homomorphisme d'anneaux
$$
R = k[x_1, x_2, x_3, \ldots] \to
S = k[x_1, x_2, x_3, \ldots, y_1, y_2, y_3, \ldots]/
(x_1y_1, x_2y_2, x_3y_3, \ldots)
$$
et par $M = S$. Dans ce cas, $\mathfrak q = \sum x_iS$ est un idéal premier minimal de
$S$, donc un idéal premier faiblement associé à $M = S$ (voir le
Lemme \ref{algebra-lemma-weakly-ass-minimal-prime-support}).
Mais, d'autre part, pour tout élément non nul de $S$, son annulateur dans $R$
est de type fini et n'a donc pas
pour radical $R \cap \mathfrak q = (x_1, x_2, x_3, \ldots)$
(détails omis).
\end{remark}

\begin{lemma}
\label{algebra-lemma-weakly-ass-reverse-functorial}
Soit $\varphi : R \to S$ un homomorphisme d'anneaux. Soit $M$ un $S$-module.
Alors on a
$\Spec(\varphi)(\text{WeakAss}_S(M)) \supset \text{WeakAss}_R(M)$.
\end{lemma}

\begin{proof}
Soit $\mathfrak p$ un élément de $\text{WeakAss}_R(M)$.
Il existe alors $m \in M_{\mathfrak p}$ dont l'annulateur
$I = \{x \in R_{\mathfrak p} \mid xm = 0\}$ a pour radical
$\mathfrak pR_{\mathfrak p}$. Considérons l'annulateur
$J = \{x \in S_{\mathfrak p} \mid xm = 0 \}$ de $m$ dans $S_{\mathfrak p}$.
Comme $IS_{\mathfrak p} \subset J$, tout idéal premier minimal
$\mathfrak q \subset S_{\mathfrak p}$ au-dessus de $J$ est au-dessus de $\mathfrak p$.
De plus, un tel $\mathfrak q$ correspond à un idéal premier faiblement associé
à $M$, par exemple d'après le
Lemme \ref{algebra-lemma-weakly-ass-local}.
\end{proof}

\begin{remark}
\label{algebra-remark-ass-functorial}
Soit $\varphi : R \to S$ un homomorphisme d'anneaux. Soit $M$ un $S$-module.
Notons $f : \Spec(S) \to \Spec(R)$ l'application induite sur les spectres.
Alors on a
$$
f(\text{Ass}_S(M)) \subset
\text{Ass}_R(M) \subset
\text{WeakAss}_R(M) \subset
f(\text{WeakAss}_S(M))
$$
voir les
Lemmes \ref{algebra-lemma-ass-functorial},
\ref{algebra-lemma-weakly-ass-reverse-functorial} et
\ref{algebra-lemma-weakly-ass-support}.
En général, toutes les inclusions peuvent être strictes ; voir les
Remarques \ref{algebra-remark-ass-reverse-functorial} et
\ref{algebra-remark-weakly-ass-not-functorial}.
Si $S$ est noethérien, alors toutes les inclusions sont des égalités puisque,
d'après le
Lemme \ref{algebra-lemma-ass-weakly-ass}, les deux extrêmes sont égaux.
\end{remark}

\begin{lemma}
\label{algebra-lemma-weakly-ass-finite-ring-map}
Soit $\varphi : R \to S$ un homomorphisme d'anneaux. Soit $M$ un $S$-module.
Notons $f : \Spec(S) \to \Spec(R)$ l'application induite sur les spectres.
Si $\varphi$ est un homomorphisme d'anneaux fini, alors
$$
\text{WeakAss}_R(M) = f(\text{WeakAss}_S(M)).
$$
\end{lemma}

\begin{proof}
L'une des inclusions a déjà été démontrée ; voir la
Remarque \ref{algebra-remark-ass-functorial}.
Pour démontrer l'autre, supposons $\mathfrak q \in \text{WeakAss}_S(M)$
et soit $\mathfrak p$ l'idéal premier correspondant de $R$. Soit $m \in M$
un élément tel que $\mathfrak q$ soit un idéal premier minimal au-dessus de
$J = \{g \in S \mid gm = 0\}$. Le radical de
$JS_{\mathfrak q}$ est donc $\mathfrak qS_{\mathfrak q}$.
Comme $R \to S$ est fini, il n'existe qu'un nombre
fini d'idéaux premiers
$\mathfrak q = \mathfrak q_1, \mathfrak q_2, \ldots, \mathfrak q_l$
au-dessus de $\mathfrak p$ ; voir le
Lemme \ref{algebra-lemma-finite-finite-fibres}.
Choisissons $x \in \mathfrak q$ tel que $x \not \in \mathfrak q_i$ pour $i > 1$ ; voir le
Lemme \ref{algebra-lemma-silly}.
D'après ce qui précède, il existe un élément $y \in S$, avec $y \not \in \mathfrak q$,
et un entier $t > 0$ tels que $y x^t m = 0$. Ainsi, l'élément
$ym \in M$ est annulé par $x^t$ ; ainsi $ym$ a une image nulle dans
$M_{\mathfrak q_i}$ pour $i = 2, \ldots, l$. En revanche, l'image de $ym$ n'est
pas nulle dans $S_{\mathfrak q}$.

\medskip\noindent
L'anneau $S_{\mathfrak p}$ est semi-local, d'idéaux maximaux
$\mathfrak q_i S_{\mathfrak p}$, par le théorème de montée pour les homomorphismes
d'anneaux finis ; voir le Lemme \ref{algebra-lemma-integral-going-up}.
Si $f \in \mathfrak pR_{\mathfrak p}$, une puissance de $f$ appartient
à $JS_{\mathfrak q}$ ; ainsi, pour un certain $t > 0$, l'image de
$f^t ym$ est nulle dans $M_{\mathfrak q}$. Comme $ym$ s'annule aux
autres idéaux maximaux de $S_{\mathfrak p}$, on en déduit que $f^t ym$ est nul
dans $M_{\mathfrak p}$ ; voir le
Lemme \ref{algebra-lemma-characterize-zero-local}.
On voit ainsi que $\mathfrak p$ est un idéal premier minimal au-dessus de
l'annulateur de $ym$ dans $R$, ce qui conclut.
\end{proof}

\begin{lemma}
\label{algebra-lemma-weakly-ass-quotient-ring}
Soit $R$ un anneau.
Soit $I$ un idéal.
Soit $M$ un $R/I$-module.
Par l'injection canonique
$\Spec(R/I) \to \Spec(R)$
on a $\text{WeakAss}_{R/I}(M) = \text{WeakAss}_R(M)$.
\end{lemma}

\begin{proof}
C'est un cas particulier du Lemme \ref{algebra-lemma-weakly-ass-finite-ring-map}.
\end{proof}

\begin{lemma}
\label{algebra-lemma-localize-weakly-ass}
Soit $R$ un anneau. Soit $M$ un $R$-module.
Soit $S \subset R$ une partie multiplicative.
Par l'injection canonique $\Spec(S^{-1}R) \to \Spec(R)$,
on a $\text{WeakAss}_R(S^{-1}M) = \text{WeakAss}_{S^{-1}R}(S^{-1}M)$
et
$$
\text{WeakAss}(M) \cap \Spec(S^{-1}R) = \text{WeakAss}(S^{-1}M).
$$
\end{lemma}

\begin{proof}
Supposons $m \in S^{-1}M$. Soient $I = \{x \in R \mid xm = 0\}$
et $I' = \{x' \in S^{-1}R \mid x'm = 0\}$. Alors $I' = S^{-1}I$
et $I \cap S = \emptyset$, sauf si $I = R$ (vérifications omises).
Ainsi, les idéaux premiers de $S^{-1}R$ minimaux au-dessus de $I'$ correspondent bijectivement
aux idéaux premiers de $R$ minimaux au-dessus de $I$ et disjoints de $S$. Ceci démontre
l'égalité $\text{WeakAss}_R(S^{-1}M) = \text{WeakAss}_{S^{-1}R}(S^{-1}M)$.
La seconde égalité résulte du
Lemme \ref{algebra-lemma-weakly-ass-local},
car, pour $\mathfrak p \in R$ tel que $S \cap \mathfrak p = \emptyset$, on a
$M_{\mathfrak p} = (S^{-1}M)_{S^{-1}\mathfrak p}$.
\end{proof}

\begin{lemma}
\label{algebra-lemma-localize-weakly-ass-nonzero-divisors}
Soit $R$ un anneau. Soit $M$ un $R$-module.
Soit $S \subset R$ une partie multiplicative.
Supposons que tout $s \in S$ soit un non-diviseur de zéro sur $M$.
Alors
$$
\text{WeakAss}(M) = \text{WeakAss}(S^{-1}M).
$$
\end{lemma}

\begin{proof}
Puisque, par hypothèse, $M \subset S^{-1}M$, le
Lemme \ref{algebra-lemma-weakly-ass} donne
$\text{WeakAss}(M) \subset \text{WeakAss}(S^{-1}M)$.
Réciproquement, soit $n/s \in S^{-1}M$ un élément d'annulateur
$I$, et soit $\mathfrak p$ un idéal premier minimal au-dessus de $I$.
Alors l'annulateur de $n \in M$ est $I$, et $\mathfrak p$ est un idéal premier
minimal au-dessus de $I$.
\end{proof}

\begin{lemma}
\label{algebra-lemma-zero-at-weakly-ass-zero}
Soit $R$ un anneau. Soit $M$ un $R$-module. L'application
$$
M
\longrightarrow
\prod\nolimits_{\mathfrak p \in \text{WeakAss}(M)} M_{\mathfrak p}
$$
est injective.
\end{lemma}

\begin{proof}
Soit $x \in M$ un élément du noyau de l'application. Posons
$N = Rx \subset M$. Si $\mathfrak p$ est un idéal premier faiblement associé à $N$,
on a d'une part $\mathfrak p \in \text{WeakAss}(M)$
(Lemme \ref{algebra-lemma-weakly-ass})
et, d'autre part, $N_{\mathfrak p} \subset M_{\mathfrak p}$
n'est pas nul. Cette contradiction montre que $\text{WeakAss}(N) = \emptyset$.
Ainsi $N = 0$, c'est-à-dire $x = 0$, d'après le
Lemme \ref{algebra-lemma-weakly-ass-zero}.
\end{proof}

\begin{lemma}
\label{algebra-lemma-weak-post-bourbaki}
Soit $R \to S$ un homomorphisme d'anneaux.
Soit $N$ un $S$-module.
Supposons que $N$ soit plat comme $R$-module et que
$R$ soit intègre, de corps des fractions $K$.
Alors
$$
\text{WeakAss}_S(N) = \text{WeakAss}_{S \otimes_R K}(N \otimes_R K)
$$
par l'inclusion canonique
$\Spec(S \otimes_R K) \subset \Spec(S)$.
\end{lemma}

\begin{proof}
Remarquons que $S \otimes_R K = (R \setminus \{0\})^{-1}S$ et
$N \otimes_R K = (R \setminus \{0\})^{-1}N$.
Pour tout $x \in R$ non nul, la multiplication par $x$ dans $N$ est injective, car
$N$ est plat sur $R$. Le lemme résulte donc du
Lemme \ref{algebra-lemma-localize-weakly-ass-nonzero-divisors}.
\end{proof}

\begin{lemma}
\label{algebra-lemma-weakly-ass-change-fields}
Soit $K/k$ une extension de corps. Soit $R$ une $k$-algèbre.
Soit $M$ un $R$-module. Soit $\mathfrak q \subset R \otimes_k K$
un idéal premier au-dessus de $\mathfrak p \subset R$. Si
$\mathfrak q$ est faiblement associé à $M \otimes_k K$,
alors $\mathfrak p$ est faiblement associé à $M$.
\end{lemma}

\begin{proof}
Soit $z \in M \otimes_k K$ un élément tel que $\mathfrak q$
soit minimal au-dessus de l'annulateur $J \subset R \otimes_k K$ de $z$.
Choisissons une sous-extension de type fini $K/L/k$ telle que
$z \in M \otimes_k L$. Comme $R \otimes_k L \to R \otimes_k K$
est plat, on a $J = I(R \otimes_k K)$, où $I \subset R \otimes_k L$
est l'annulateur de $z$ dans le plus petit anneau
(Lemme \ref{algebra-lemma-annihilator-flat-base-change}).
Ainsi, $\mathfrak q \cap (R \otimes_k L)$ est minimal au-dessus de $I$ par le
théorème de descente (Lemme \ref{algebra-lemma-flat-going-down}).
On se ramène ainsi au cas décrit dans le paragraphe suivant.

\medskip\noindent
Supposons que $K/k$ soit une extension de corps de type fini.
Soient $x_1, \ldots, x_r \in K$ une base de transcendance
de $K$ sur $k$ ; voir Corps, section \ref{fields-section-transcendence}.
Posons $L = k(x_1, \ldots, x_r)$. Écrivons $[K : L] = n$. Alors
$R \otimes_k L \to R \otimes_k K$ est un homomorphisme d'anneaux fini.
Le Lemme \ref{algebra-lemma-weakly-ass-finite-ring-map} montre donc que
$\mathfrak q \cap (R \otimes_k L)$
est un idéal premier faiblement associé à $M \otimes_k K$
considéré comme $R \otimes_k L$-module.
Comme $M \otimes_k K \cong (M \otimes_k L)^{\oplus n}$
en tant que $R \otimes_k L$-module, on voit que
$\mathfrak q \cap (R \otimes_k L)$
est un idéal premier faiblement associé à $M \otimes_k L$
(par exemple en utilisant le Lemme \ref{algebra-lemma-weakly-ass} et une récurrence).
On se ramène ainsi au cas examiné dans le paragraphe suivant.

\medskip\noindent
Supposons que $K = k(x_1, \ldots, x_r)$ soit une extension transcendante pure.
On peut remplacer $R$ par $R_\mathfrak p$, $M$ par $M_\mathfrak p$
et $\mathfrak q$ par $\mathfrak q(R_\mathfrak p \otimes_k K)$.
Voir le Lemme \ref{algebra-lemma-localize-weakly-ass}.
On se ramène ainsi au cas examiné dans le paragraphe suivant.

\medskip\noindent
Supposons que $K = k(x_1, \ldots, x_r)$ soit une extension transcendante pure
et que $R$ soit local d'idéal maximal $\mathfrak p$. Nous affirmons que tout
$f \in R \otimes_k K$, avec $f \not \in \mathfrak p(R \otimes_k K)$,
est un non-diviseur de zéro sur $M \otimes_k K$. En effet, soit
$z \in M \otimes_k K$ un élément.
Il existe un $R$-sous-module de type fini $M' \subset M$ tel que
$z \in M' \otimes_k K$ et tel que $M'$ soit minimal pour
cette propriété : choisissons une base $\{t_\alpha\}$ de $K$ comme
$k$-espace vectoriel, écrivons $z = \sum m_\alpha \otimes t_\alpha$ et prenons pour
$M'$ le $R$-sous-module engendré par les $m_\alpha$.
Si $z \in \mathfrak p(M' \otimes_k K) = \mathfrak p M' \otimes_k K$,
alors $\mathfrak pM' = M'$ et $M' = 0$ d'après le Lemme \ref{algebra-lemma-NAK},
ce qui est une contradiction.
Ainsi, $z$ a une image non nulle $\overline{z}$ dans $M'/\mathfrak p M' \otimes_k K$.
Or $R/\mathfrak p \otimes_k K$ est intègre, car c'est une localisation
de $\kappa(\mathfrak p)[x_1, \ldots, x_n]$, et
$M'/\mathfrak p M' \otimes_k K$ est un module libre ; par conséquent,
$f\overline{z} \not = 0$. Cela démontre l'affirmation.

\medskip\noindent
Enfin, choisissons $z \in M \otimes_k K$ tel que $\mathfrak q$
soit minimal au-dessus de l'annulateur $J \subset R \otimes_k K$ de $z$.
Pour $f \in \mathfrak p$, il existe $n \geq 1$ et
$g \in R \otimes_k K$, avec $g \not \in \mathfrak q$, tels que
$g f^n z \in J$, c'est-à-dire $g f^n z = 0$.
(Cela vient de ce que $\mathfrak q$ est au-dessus de $\mathfrak p$
et que $\mathfrak q$ est minimal au-dessus de $J$.)
On a vu plus haut que $g$ est un non-diviseur de zéro, donc $f^n z = 0$.
Ainsi, $\mathfrak p$ est un idéal premier faiblement associé
à $M \otimes_k K$ considéré comme $R$-module.
Comme $M \otimes_k K$ est une somme directe de copies de $M$,
on en déduit, comme précédemment, que $\mathfrak p$ est un idéal premier
faiblement associé à $M$.
\end{proof}





\section{Idéaux premiers immergés}
\label{algebra-section-embedded-primes}

\noindent
Voici la définition.

\begin{definition}
\label{algebra-definition-embedded-primes}
Soit $R$ un anneau.
Soit $M$ un $R$-module.
\begin{enumerate}
\item  Les idéaux premiers associés à $M$ qui ne sont
pas minimaux parmi les idéaux premiers associés à $M$ sont appelés les
{\it idéaux premiers associés immergés} de $M$.
\item Les {\it idéaux premiers immergés de $R$}
sont les idéaux premiers associés immergés de $R$ considéré comme $R$-module.
\end{enumerate}
\end{definition}

\noindent
Voici une manière de les éliminer.

\begin{lemma}
\label{algebra-lemma-remove-embedded-primes}
Soit $R$ un anneau noethérien.
Soit $M$ un $R$-module de type fini.
Considérons l'ensemble des $R$-sous-modules
$$
\{
K \subset M
\mid
\text{Supp}(K)
\text{ rare dans }
\text{Supp}(M)
\}.
$$
Cet ensemble possède un élément maximal $K$, et le quotient
$M' = M/K$ possède les propriétés suivantes :
\begin{enumerate}
\item $\text{Supp}(M) = \text{Supp}(M')$,
\item $M'$ n'a aucun idéal premier associé immergé,
\item pour tout $f \in R$ appartenant à tous les
idéaux premiers associés immergés de $M$, on a $M_f \cong M'_f$.
\end{enumerate}
\end{lemma}

\begin{proof}
Nous utiliserons le Lemme \ref{algebra-lemma-finite-ass} et la
Proposition \ref{algebra-proposition-minimal-primes-associated-primes}
sans les mentionner davantage.
Notons $\mathfrak q_1, \ldots, \mathfrak q_t$ les
idéaux premiers minimaux du support de $M$. Notons
$\mathfrak p_1, \ldots, \mathfrak p_s$ les
idéaux premiers associés immergés de $M$. Alors
$\text{Ass}(M) = \{\mathfrak q_j, \mathfrak p_i\}$.
Posons
$$
K = \{m \in M \mid \text{Supp}(Rm) \subset \bigcup V(\mathfrak p_i)\}
$$
On voit immédiatement que c'est un sous-module. Puisque $M$ est de type fini sur un
anneau noethérien, $K$ est lui aussi de type fini. Ainsi $\text{Supp}(K)$
est rare dans $\text{Supp}(M)$. Soit $K' \subset M$ un autre sous-module
dont le support est rare dans $\text{Supp}(M)$. Cela signifie
que $K_{\mathfrak q_j} = 0$. Ainsi, si $m \in K'$,
l'image de $m$ dans $M_{\mathfrak q_j}$ est nulle, ce qui
implique à son tour $(Rm)_{\mathfrak q_j} = 0$.
D'autre part, on a $\text{Ass}(Rm) \subset \text{Ass}(M)$.
Le support de $Rm$ est donc contenu dans $\bigcup V(\mathfrak p_i)$.
Ainsi $m \in K$, puis $K' \subset K$, puisque $m$ était arbitraire dans $K'$.

\medskip\noindent
Soit $M' = M/K$. Puisque $K_{\mathfrak q_j}=0$, on a
$M'_{\mathfrak q_j} = M_{\mathfrak q_j}$ pour tout $j$.
Ainsi $M$ et $M'$ ont le même support.

\medskip\noindent
Supposons $\mathfrak q = \text{Ann}(\overline{m}) \in \text{Ass}(M')$,
où $\overline{m} \in M'$ est l'image de $m \in M$.
Alors $m \not \in K$ ; le support de $Rm$ doit donc contenir l'un des
$\mathfrak q_j$. Puisque $M_{\mathfrak q_j} = M'_{\mathfrak q_j}$,
on sait que l'image de $\overline{m}$ dans $M'_{\mathfrak q_j}$ n'est pas nulle.
Ainsi $\mathfrak q \subset \mathfrak q_j$ (en fait, il y a égalité),
ce qui signifie qu'aucun des idéaux premiers associés à $M'$ n'est immergé.

\medskip\noindent
Soit $f$ un élément appartenant à tous les $\mathfrak p_i$.
Alors $D(f) \cap \text{supp}(K) = 0$. Ainsi $M_f = M'_f$
parce que $K_f = 0$.
\end{proof}

\begin{lemma}
\label{algebra-lemma-remove-embedded-primes-localize}
Soit $R$ un anneau noethérien.
Soit $M$ un $R$-module de type fini.
Pour tout $f \in R$, on a $(M')_f = (M_f)'$, où
$M \to M'$ et $M_f \to (M_f)'$ sont les quotients
construits dans le Lemme \ref{algebra-lemma-remove-embedded-primes}.
\end{lemma}

\begin{proof}
Omis.
\end{proof}

\begin{lemma}
\label{algebra-lemma-no-embedded-primes-endos}
Soit $R$ un anneau noethérien.
Soit $M$ un $R$-module de type fini sans idéal premier associé immergé.
Soit $I = \{x \in R \mid xM = 0\}$. Alors l'anneau $R/I$ n'a aucun
idéal premier immergé.
\end{lemma}

\begin{proof}
On peut remplacer $R$ par $R/I$.
On peut donc supposer que tout élément non nul
de $R$ agit non trivialement sur $M$.
D'après le Lemme \ref{algebra-lemma-support-closed}, cela implique que
$\Spec(R)$ est égal au support de $M$.
Supposons que $\mathfrak p$ soit un idéal premier immergé de $R$.
Soit $x \in R$ un élément dont l'annulateur est $\mathfrak p$.
Considérons le module non nul $N = xM \subset M$. Il est annulé
par $\mathfrak p$. Tout idéal premier associé $\mathfrak q$ à $N$
contient donc $\mathfrak p$ et est aussi un idéal premier associé à $M$.
Alors $\mathfrak q$ serait un idéal premier associé immergé de
$M$, ce qui contredit l'hypothèse du lemme.
\end{proof}

















\section{Suites régulières}
\label{algebra-section-regular-sequences}

\noindent
Dans cette section, nous établissons quelques propriétés élémentaires des suites régulières.

\begin{definition}
\label{algebra-definition-regular-sequence}
Soit $R$ un anneau. Soit $M$ un $R$-module. Une suite d'éléments
$f_1, \ldots, f_r$ de $R$ est appelée une {\it suite $M$-régulière}
si les conditions suivantes sont satisfaites :
\begin{enumerate}
\item $f_i$ est un non-diviseur de zéro sur
$M/(f_1, \ldots, f_{i - 1})M$
pour tout $i = 1, \ldots, r$, et
\item le module $M/(f_1, \ldots, f_r)M$ n'est pas nul.
\end{enumerate}
Si $I$ est un idéal de $R$ et si $f_1, \ldots, f_r \in I$,
nous appelons $f_1, \ldots, f_r$ une {\it suite $M$-régulière
dans $I$}. Si $M = R$, nous appelons simplement $f_1, \ldots, f_r$ une
{\it suite régulière} (dans $I$).
\end{definition}

\noindent
Notons bien que la définition dépend de l'ordre
des éléments $f_1, \ldots, f_r$ (voir les exemples ci-dessous). Certains articles ou livres
ne demandent pas que le module $M/(f_1, \ldots, f_r)M$ soit non nul.
Cela présente l'avantage que la propriété d'être une suite régulière est préservée par
localisation. Cependant, nous utiliserons surtout cette définition pour définir la
profondeur d'un module lorsque $R$ est local ; dans ce cas, les $f_i$ doivent
appartenir à l'idéal maximal, condition qui n'est pas préservée lorsqu'on passe
de $R$ à une localisation $R_\mathfrak p$.

\begin{example}
\label{algebra-example-global-regular}
Soit $k$ un corps. Dans l'anneau $k[x, y, z]$,
la suite $x, y(1-x), z(1-x)$ est régulière,
mais la suite $y(1-x), z(1-x), x$ ne l'est pas.
\end{example}

\begin{example}
\label{algebra-example-local-regular}
Soit $k$ un corps. Considérons l'anneau
$k[x, y, w_0, w_1, w_2, \ldots]/I$
où $I$ est engendré par les $yw_i$, $i = 0, 1, 2, \ldots$, et les
$w_i - xw_{i + 1}$, $i = 0, 1, 2, \ldots$.
La suite $x, y$ est régulière, mais $y$ est un diviseur de zéro.
On peut de plus localiser en l'idéal maximal
$(x, y, w_i)$ et obtenir encore un exemple.
\end{example}

\begin{lemma}
\label{algebra-lemma-permute-xi}
Soit $R$ un anneau local noethérien.
Soit $M$ un $R$-module de type fini.
Soit $x_1, \ldots, x_c$ une suite $M$-régulière.
Alors toute permutation des $x_i$ est elle aussi une suite
régulière.
\end{lemma}

\begin{proof}
Traitons d'abord le cas $c = 2$.
Considérons le noyau $K \subset M$ de $x_2 : M \to M$.
Pour tout $z \in K$, on sait que $z = x_1 z'$
pour un certain $z' \in M$, car
$x_2$ est un non-diviseur de zéro sur $M/x_1M$.
Comme $x_1$ est un non-diviseur de zéro sur $M$, on a aussi $x_2 z' = 0$.
Ainsi $x_1 : K \to K$ est surjectif.
Le Lemme de Nakayama \ref{algebra-lemma-NAK} donne donc $K = 0$.
Considérons ensuite la multiplication par $x_1$ sur $M/x_2M$.
Si $z \in M$ a pour image un élément $\overline{z} \in M/x_2M$
du noyau de cette application, alors $x_1 z = x_2 y$ pour un certain $y \in M$.
Mais, puisque $x_1, x_2$ est une suite régulière, on a alors
$y = x_1 y'$ pour un certain $y' \in M$. Ainsi $x_1 ( z - x_2 y' ) =0$,
puis $z = x_2 y'$, et donc $\overline{z} = 0$, comme voulu.

\medskip\noindent
Dans le cas général, remarquons que toute permutation est
une composée de transpositions d'indices adjacents.
Il suffit donc de démontrer que
$$
x_1, \ldots, x_{i-2}, x_i, x_{i-1}, x_{i + 1}, \ldots, x_c
$$
est une suite $M$-régulière. Cela résulte du cas que nous venons
de traiter, appliqué au module $M/(x_1, \ldots, x_{i-2})$
et à la suite régulière de longueur $2$ formée par $x_{i-1}, x_i$.
\end{proof}

\begin{lemma}
\label{algebra-lemma-flat-increases-depth}
\begin{slogan}
Les homomorphismes locaux plats d'anneaux préservent et reflètent les suites régulières.
\end{slogan}
Soient $R, S$ des anneaux locaux. Soit $R \to S$ un homomorphisme local plat d'anneaux.
Soit $x_1, \ldots, x_r$ une suite dans $R$. Soit $M$ un $R$-module.
Les assertions suivantes sont équivalentes :
\begin{enumerate}
\item $x_1, \ldots, x_r$ est une suite $M$-régulière dans $R$, et
\item les images de $x_1, \ldots, x_r$ dans $S$ forment une suite
$M \otimes_R S$-régulière.
\end{enumerate}
\end{lemma}

\begin{proof}
Cela vient de ce que $R \to S$ est fidèlement plat
d'après le Lemme \ref{algebra-lemma-local-flat-ff}.
\end{proof}

\begin{lemma}
\label{algebra-lemma-regular-sequence-in-neighbourhood}
Soit $R$ un anneau noethérien. Soit $M$ un $R$-module de type fini.
Soit $\mathfrak p$ un idéal premier. Soit $x_1, \ldots, x_r$ une suite
dans $R$ dont l'image dans $R_{\mathfrak p}$ forme une suite
$M_{\mathfrak p}$-régulière. Alors il existe $g \in R$, avec $g \not \in \mathfrak p$,
tel que l'image de $x_1, \ldots, x_r$ dans $R_g$ forme
une suite $M_g$-régulière.
\end{lemma}

\begin{proof}
Posons
$$
K_i = \Ker\left(x_i : M/(x_1, \ldots, x_{i - 1})M \to
M/(x_1, \ldots, x_{i - 1})M\right).
$$
C'est un $R$-module de type fini dont la localisation en $\mathfrak p$ est
nulle par hypothèse. Il existe donc $g \in R$, avec $g \not \in \mathfrak p$,
tel que $(K_i)_g = 0$ pour tout $i = 1, \ldots, r$. Ce $g$ convient.
\end{proof}

\begin{lemma}
\label{algebra-lemma-join-regular-sequences}
Soit $A$ un anneau. Soit $I$ un idéal engendré par une suite régulière
$f_1, \ldots, f_n$ dans $A$. Soient $g_1, \ldots, g_m \in A$ des
éléments dont les images $\overline{g}_1, \ldots, \overline{g}_m$ forment une
suite régulière dans $A/I$. Alors la suite $f_1, \ldots, f_n, g_1, \ldots, g_m$
est régulière dans $A$.
\end{lemma}

\begin{proof}
Cela résulte immédiatement des définitions.
\end{proof}

\begin{lemma}
\label{algebra-lemma-regular-sequence-short-exact-sequence}
Soit $R$ un anneau. Soit $0 \to M_1 \to M_2 \to M_3 \to 0$
une suite exacte courte de $R$-modules. Soit $f_1, \ldots, f_r \in R$.
Si $f_1, \ldots, f_r$ est $M_1$-régulière et $M_3$-régulière, alors
$f_1, \ldots, f_r$ est $M_2$-régulière.
\end{lemma}

\begin{proof}
D'après le Lemme \ref{algebra-lemma-snake}, si $f_1 : M_1 \to M_1$ et
$f_1 : M_3 \to M_3$ sont injectifs, alors $f_1 : M_2 \to M_2$ l'est aussi,
et l'on obtient une suite exacte courte
$$
0 \to M_1/f_1M_1 \to M_2/f_1M_2 \to M_3/f_1M_3 \to 0
$$
Le lemme en résulte par récurrence sur $r$. Certains détails sont omis.
\end{proof}

\begin{lemma}
\label{algebra-lemma-regular-sequence-powers}
Soit $R$ un anneau. Soit $M$ un $R$-module.
Soient $f_1, \ldots, f_r \in R$ et des entiers $e_1, \ldots, e_r > 0$.
Alors la suite $f_1, \ldots, f_r$ est $M$-régulière
si et seulement si $f_1^{e_1}, \ldots, f_r^{e_r}$
est une suite $M$-régulière.
\end{lemma}

\begin{proof}
Nous le démontrons par récurrence sur $r$. Si $r = 1$, cela résulte des
deux faits élémentaires suivants : (a) une puissance d'un non-diviseur de zéro sur $M$
est un non-diviseur de zéro sur $M$, et (b) un diviseur d'un non-diviseur de zéro sur $M$
est un non-diviseur de zéro sur $M$.
Si $r > 1$, la récurrence appliquée à $M/f_1M$ montre que
$f_1, f_2, \ldots, f_r$ est une suite $M$-régulière si et seulement si
$f_1, f_2^{e_2}, \ldots, f_r^{e_r}$ est une suite $M$-régulière.
Il suffit donc de montrer, étant donné $e > 0$, que $f_1^e, f_2, \ldots, f_r$
est une suite $M$-régulière si et seulement si $f_1, \ldots, f_r$
est une suite $M$-régulière. Nous le démontrons
par récurrence sur $e$. Le cas $e = 1$ est trivial. Puisque $f_1$ est un
non-diviseur de zéro sous les deux hypothèses (d'après le cas $r = 1$),
on a une suite exacte courte
$$
0 \to M/f_1M \xrightarrow{f_1^{e - 1}} M/f_1^eM \to M/f_1^{e - 1}M \to 0
$$
Supposons que $f_1, f_2, \ldots, f_r$ soit une suite $M$-régulière.
Alors, par récurrence, les éléments $f_2, \ldots, f_r$ sont des suites régulières sur $M/f_1M$ et
$M/f_1^{e - 1}M$. D'après le
Lemme \ref{algebra-lemma-regular-sequence-short-exact-sequence},
$f_2, \ldots, f_r$ est $M/f_1^eM$-régulière. Ainsi $f_1^e, f_2, \ldots, f_r$
est $M$-régulière. Réciproquement, supposons
que $f_1^e, f_2, \ldots, f_r$ soit une suite $M$-régulière. Alors
$f_2 : M/f_1^eM \to M/f_1^eM$ est injectif, donc
$f_2 : M/f_1M \to M/f_1M$ est injectif, puis, par récurrence (!),
$f_2 : M/f_1^{e - 1}M \to M/f_1^{e - 1}M$ est injectif, et donc
$$
0 \to
M/(f_1, f_2)M \xrightarrow{f_1^{e - 1}}
M/(f_1^e, f_2)M \to
M/(f_1^{e - 1}, f_2)M \to 0
$$
est une suite exacte courte d'après le Lemme \ref{algebra-lemma-snake}. Cela démontre la
réciproque pour $r = 2$. Si $r > 2$, alors
$f_3 : M/(f_1^e, f_2)M \to M/(f_1^e, f_2)M$ est injectif, donc
$f_3 : M/(f_1, f_2)M \to M/(f_1, f_2)M$ est injectif, et ainsi de suite.
Certains détails sont omis.
\end{proof}

\begin{lemma}
\label{algebra-lemma-regular-sequence-in-polynomial-ring}
Soit $R$ un anneau. Soient $f_1, \ldots, f_r \in R$ qui n'engendrent pas
l'idéal unité. Les assertions suivantes sont équivalentes :
\begin{enumerate}
\item toute permutation de $f_1, \ldots, f_r$ est une suite régulière,
\item toute sous-suite de $f_1, \ldots, f_r$ (dans l'ordre donné) est
une suite régulière, et
\item $f_1x_1, \ldots, f_rx_r$ est une suite régulière dans l'anneau de
polynômes $R[x_1, \ldots, x_r]$.
\end{enumerate}
\end{lemma}

\begin{proof}
Il est clair que (1) implique (2). Montrons par récurrence sur $r$
que (2) implique (1). Le cas $r = 1$ est trivial. Pour $r = 2$, il s'agit de dire que si
$a, b \in R$ est une suite régulière et si $b$ est un non-diviseur de zéro, alors
$b, a$ est une suite régulière. Cela est clair, car le noyau de
$a : R/(b) \to R/(b)$ est isomorphe au noyau de $b : R/(a) \to R/(a)$
si $a$ et $b$ sont tous deux des non-diviseurs de zéro. Supposons $r > 2$.
Supposons (2) vérifiée et cherchons à montrer que $f_{\sigma(1)}, \ldots, f_{\sigma(r)}$
est une suite régulière pour une permutation $\sigma$. Nous savons déjà
par récurrence que $f_{\sigma(1)}, \ldots, f_{\sigma(r - 1)}$ est une suite régulière.
Il suffit donc de montrer que $f_s$, où $s = \sigma(r)$,
est un non-diviseur de zéro modulo $f_1, \ldots, \hat f_s, \ldots, f_r$.
Si $s = r$, c'est terminé. Si $s < r$, remarquons que $f_s$ et $f_r$
sont tous deux des non-diviseurs de zéro dans l'anneau
$R/(f_1, \ldots, \hat f_s, \ldots, f_{r - 1})$
(encore par hypothèse de récurrence). Puisque $f_s, f_r$ est une
suite régulière dans cet anneau, le cas des suites de longueur
$2$ montre que $f_r, f_s$ l'est également.

\medskip\noindent
Remarquons que $R[x_1, \ldots, x_r]/(f_1x_1, \ldots, f_ix_i)$, comme $R$-module,
est une somme directe des modules
$$
R/I_E \cdot x_1^{e_1} \ldots x_r^{e_r}
$$
indexés par les multi-indices $E = (e_1, \ldots, e_r)$, où
$I_E$ est l'idéal engendré par les $f_j$ tels que $1 \leq j \leq i$
et $e_j > 0$. Ainsi $f_{i + 1}x_i$ est un non-diviseur de zéro sur ce module si
et seulement si $f_{i + 1}$ est un non-diviseur de zéro sur $R/I_E$ pour tout $E$.
En prenant $E$ dont toutes les composantes sont strictement positives, on voit que $f_{i + 1}$
est un non-diviseur de zéro sur $R/(f_1, \ldots, f_i)$. Ainsi (3) implique (2).
Réciproquement, si (2) est vérifiée, toute sous-suite de
$f_1, \ldots, f_i, f_{i + 1}$ est une suite régulière ;
en particulier, $f_{i + 1}$ est un non-diviseur de zéro sur tous les $R/I_E$.
On voit ainsi que (2) implique (3).
\end{proof}










\section{Suites quasi-régulières}
\label{algebra-section-quasi-regular}

\noindent
Nous introduisons la notion de suite quasi-régulière, légèrement plus faible
que celle de suite régulière et plus facile à utiliser.
Soit $R$ un anneau et soient $f_1, \ldots, f_c \in R$.
Posons $J = (f_1, \ldots, f_c)$. Soit $M$ un $R$-module.
Il existe alors une application canonique
\begin{equation}
\label{algebra-equation-quasi-regular}
M/JM \otimes_{R/J} R/J[X_1, \ldots, X_c]
\longrightarrow
\bigoplus\nolimits_{n \geq 0} J^nM/J^{n + 1}M
\end{equation}
de $R/J[X_1, \ldots, X_c]$-modules gradués définie par la règle
$$
\overline{m} \otimes X_1^{e_1} \ldots X_c^{e_c} \longmapsto
f_1^{e_1} \ldots f_c^{e_c} m \bmod J^{e_1 + \ldots + e_c + 1}M.
$$
Remarquons que (\ref{algebra-equation-quasi-regular}) est toujours surjective.

\begin{definition}
\label{algebra-definition-quasi-regular-sequence}
Soit $R$ un anneau.
Soit $M$ un $R$-module.
Une suite d'éléments $f_1, \ldots, f_c$ de $R$ est dite
{\it $M$-quasi-régulière} si (\ref{algebra-equation-quasi-regular})
est un isomorphisme. Si $M = R$, nous appelons simplement $f_1, \ldots, f_c$ une
{\it suite quasi-régulière}.
\end{definition}

\noindent
Ainsi, si $f_1, \ldots, f_c$ est une suite quasi-régulière, alors
$$
R/J[X_1, \ldots, X_c] = \bigoplus\nolimits_{n \geq 0} J^n/J^{n + 1}
$$
où $J = (f_1, \ldots, f_c)$. Il est clair que la propriété d'être une suite
quasi-régulière ne dépend pas de l'ordre de $f_1, \ldots, f_c$.

\begin{lemma}
\label{algebra-lemma-regular-quasi-regular}
Soit $R$ un anneau.
\begin{enumerate}
\item Une suite régulière $f_1, \ldots, f_c$ de $R$ est une suite
quasi-régulière.
\item Supposons que $M$ soit un $R$-module et que $f_1, \ldots, f_c$
soit une suite $M$-régulière. Alors $f_1, \ldots, f_c$ est une suite
$M$-quasi-régulière.
\end{enumerate}
\end{lemma}

\begin{proof}
Posons $J = (f_1, \ldots, f_c)$.
Démontrons la première assertion par récurrence sur $c$.
Il faut montrer que, pour toute relation
$\sum_{|I| = n} a_I f^I \in J^{n + 1}$ avec $a_I \in R$, on a en fait
$a_I \in J$ pour tout multi-indice $I$. Puisque
tout élément de $J^{n + 1}$ est de la forme $\sum_{|I| = n} b_I f^I$
avec $b_I \in J$, on peut supposer, après avoir remplacé $a_I$ par $a_I - b_I$,
que la relation s'écrit $\sum_{|I| = n} a_I f^I = 0$. On peut la réécrire
sous la forme
$$
\sum\nolimits_{e = 0}^n
\left(
\sum\nolimits_{|I'| = n - e}
a_{I', e} f^{I'}
\right)
f_c^e
=
0
$$
Ici et ci-dessous, les multi-indices « primés » $I'$ sont supposés être de la forme
$I' = (i_1, \ldots, i_{c - 1}, 0)$. Nous allons montrer par
récurrence sur $l \in \{0, \ldots, n\}$
que, si l'on a une relation
$$
\sum\nolimits_{e = 0}^l
\left(
\sum\nolimits_{|I'| = n - e}
a_{I', e} f^{I'}
\right)
f_c^e
=
0
$$
alors $a_{I', e} \in J$ pour tous $I', e$.
En effet, posons $J' = (f_1, \ldots, f_{c-1})$.
Remarquons que $\sum\nolimits_{|I'| = n - l} a_{I', l} f^{I'}$
est envoyé dans $(J')^{n - l + 1}$ par $f_c^{l}$.
L'hypothèse de récurrence (pour la récurrence sur $c$)
montre que $f_c^l a_{I', l} \in J'$.
Comme $f_c$ n'est pas un diviseur de zéro sur $R/J'$ (car $f_1, \ldots, f_c$
est une suite régulière), on en déduit que $a_{I', l} \in J'$.
Cela permet de réécrire le terme
$(\sum\nolimits_{|I'| = n - l} a_{I', l} f^{I'})f_c^l$
 sous la forme $(\sum\nolimits_{|I'| = n - l + 1} f_c b_{I', l - 1}
f^{I'})f_c^{l-1}$. On obtient ainsi une nouvelle relation de la forme
$$
\left(\sum\nolimits_{|I'| = n - l + 1}
(a_{I', l-1} + f_c b_{I', l - 1}) f^{I'}\right)f_c^{l-1}
+
\sum\nolimits_{e = 0}^{l - 2}
\left(
\sum\nolimits_{|I'| = n - e}
a_{I', e} f^{I'}
\right)
f_c^e
=
0
$$
L'hypothèse de récurrence (sur $l$ cette fois) montre maintenant que
tous les $a_{I', l-1} + f_c b_{I', l - 1} \in J$ et que
tous les $a_{I', e} \in J$ pour $e \leq l - 2$. Avec la relation
$a_{I', l} \in J' \subset J$ obtenue ci-dessus, cela achève la démonstration du
pas de récurrence.

\medskip\noindent
La seconde assertion signifie que, pour toute expression formelle
$F = \sum_{|I| = n} m_I X^I$, avec $m_I \in M$ et $\sum m_I f^I
\in J^{n + 1}M$, tous les coefficients $m_I$ appartiennent à $J$.
Cela se démontre exactement comme le résultat correspondant
de la première assertion ci-dessus.
\end{proof}

\begin{lemma}
\label{algebra-lemma-flat-base-change-quasi-regular}
Soit $R \to R'$ un homomorphisme plat d'anneaux. Soit $M$ un $R$-module.
Supposons que la suite $f_1, \ldots, f_r \in R$ soit $M$-quasi-régulière.
Alors les images de $f_1, \ldots, f_r$ dans
$R'$ forment une suite $M \otimes_R R'$-quasi-régulière.
\end{lemma}

\begin{proof}
Posons $J = (f_1, \ldots, f_r)$, $J' = JR'$ et $M' = M \otimes_R R'$.
Il faut montrer que l'application canonique
$\mu : R'/J'[X_1, \ldots X_r] \otimes_{R'/J'} M'/J'M' \to
\bigoplus (J')^nM'/(J')^{n + 1}M'$ est un isomorphisme.
Puisque $R \to R'$ est plat, les suites
$0 \to J^nM \to M$ et
$0 \to J^{n + 1}M \to J^nM \to J^nM/J^{n + 1}M \to 0$
restent exactes après tensorisation par $R'$. Cela implique d'abord
$J^nM \otimes_R R' = (J')^nM'$, puis
$(J')^nM'/(J')^{n + 1}M' = J^nM/J^{n + 1}M \otimes_R R'$.
Ainsi $\mu$ est le produit tensoriel de (\ref{algebra-equation-quasi-regular}),
qui est un isomorphisme par hypothèse,
avec $\text{id}_{R'}$, ce qui conclut.
\end{proof}

\begin{lemma}
\label{algebra-lemma-quasi-regular-sequence-in-neighbourhood}
Soit $R$ un anneau noethérien. Soit $M$ un $R$-module de type fini.
Soit $\mathfrak p$ un idéal premier. Soit $x_1, \ldots, x_c$ une suite
dans $R$ dont l'image dans $R_{\mathfrak p}$ forme une suite
$M_{\mathfrak p}$-quasi-régulière. Alors il existe
$g \in R$, avec $g \not \in \mathfrak p$,
tel que l'image de $x_1, \ldots, x_c$ dans $R_g$ forme
une suite $M_g$-quasi-régulière.
\end{lemma}

\begin{proof}
Considérons le noyau $K$ de l'application (\ref{algebra-equation-quasi-regular}).
Comme $M/JM \otimes_{R/J} R/J[X_1, \ldots, X_c]$ est un
$R/J[X_1, \ldots, X_c]$-module de type fini et que $R/J[X_1, \ldots, X_c]$ est
noethérien, $K$ est aussi un $R/J[X_1, \ldots, X_c]$-module de type fini.
Choisissons des générateurs homogènes $k_1, \ldots, k_t \in K$.
Par hypothèse, pour tout $i = 1, \ldots, t$, il existe $g_i \in R$,
avec $g_i \not \in \mathfrak p$, tel que $g_i k_i = 0$.
Alors $g = g_1 \ldots g_t$ convient.
\end{proof}

\begin{lemma}
\label{algebra-lemma-truncate-quasi-regular}
Soit $R$ un anneau. Soit $M$ un $R$-module.
Soit $f_1, \ldots, f_c \in R$ une suite $M$-quasi-régulière.
Pour tout $i$, la suite
$\overline{f}_{i + 1}, \ldots, \overline{f}_c$
de $\overline{R} = R/(f_1, \ldots, f_i)$ est une suite
$\overline{M} = M/(f_1, \ldots, f_i)M$-quasi-régulière.
\end{lemma}

\begin{proof}
Il suffit de le démontrer pour $i = 1$. Posons
$\overline{J} = (\overline{f}_2, \ldots, \overline{f}_c) \subset \overline{R}$.
Alors
\begin{align*}
\overline{J}^n\overline{M}/\overline{J}^{n + 1}\overline{M}
& =
(J^nM + f_1M)/(J^{n + 1}M + f_1M) \\
& = J^nM / (J^{n + 1}M + J^nM \cap f_1M).
\end{align*}
Ainsi, pour démontrer le lemme, il suffit de montrer que
$J^{n + 1}M + J^nM \cap f_1M = J^{n + 1}M + f_1J^{n - 1}M$
car cela montrera que
$\bigoplus_{n \geq 0}
\overline{J}^n\overline{M}/\overline{J}^{n + 1}\overline{M}$
est le quotient de
$\bigoplus_{n \geq 0} J^nM/J^{n + 1}M \cong M/JM[X_1, \ldots, X_c]$
par $X_1$. En fait, on a $J^nM \cap f_1M = f_1J^{n - 1}M$.
En effet, si $m \not \in J^{n - 1}M$, alors $f_1m \not \in J^nM$
car $\bigoplus J^nM/J^{n + 1}M$ est l'algèbre de polynômes
$M/J[X_1, \ldots, X_c]$ par hypothèse.
\end{proof}

\begin{lemma}
\label{algebra-lemma-quasi-regular-regular}
Soit $(R, \mathfrak m)$ un anneau local noethérien.
Soit $M$ un $R$-module de type fini non nul.
Soit $f_1, \ldots, f_c \in \mathfrak m$ une suite $M$-quasi-régulière.
Alors la suite $f_1, \ldots, f_c$ est $M$-régulière.
\end{lemma}

\begin{proof}
Posons $J = (f_1, \ldots, f_c)$.
Montrons que $f_1$ est un non-diviseur de zéro sur $M$.
Supposons $x \in M$ non nul.
D'après le théorème d'intersection de Krull, il existe un entier $r$
tel que $x \in J^rM$ mais $x \not \in J^{r + 1}M$ ; voir le
Lemme \ref{algebra-lemma-intersect-powers-ideal-module-zero}.
Alors $f_1 x \in J^{r + 1}M$ est un élément dont la classe dans
$J^{r + 1}M/J^{r + 2}M$ est non nulle, d'après la structure supposée de
$\bigoplus J^nM/J^{n + 1}M$. Par conséquent $f_1x \not = 0$.

\medskip\noindent
On peut maintenant achever la démonstration par récurrence sur $c$ en utilisant le
Lemme \ref{algebra-lemma-truncate-quasi-regular}.
\end{proof}

\begin{remark}[Autres types de suites régulières]
\label{algebra-remark-koszul-regular}
Dans l'article \cite{Kabele}, l'auteur étudie deux autres
conditions de régularité pour les suites $x_1, \ldots, x_r$ d'éléments
d'un anneau $R$. On dit précisément que la suite est {\it Koszul-régulière}
si $H_i(K_{\bullet}(R, x_{\bullet})) = 0$ pour $i \geq 1$, où
$K_{\bullet}(R, x_{\bullet})$ est le complexe de Koszul. La suite est
dite {\it $H_1$-régulière} si $H_1(K_{\bullet}(R, x_{\bullet})) = 0$.
On a les implications régulière $\Rightarrow$ Koszul-régulière $\Rightarrow$
$H_1$-régulière $\Rightarrow$ quasi-régulière. Par des exemples, l'auteur montre que
ces implications ne sont en général pas réversibles, même si $R$ est un anneau local
(non noethérien) et si la suite engendre l'idéal maximal de
$R$. Nous présentons ces notions plus en détail dans
Compléments d'algèbre, section \ref{more-algebra-section-koszul-regular}.
\end{remark}

\begin{remark}
\label{algebra-remark-join-quasi-regular-sequences}
Soit $k$ un corps. Considérons l'anneau
$$
A = k[x, y, w, z_0, z_1, z_2, \ldots]/
(y^2z_0 - wx, z_0 - yz_1, z_1 - yz_2, \ldots)
$$
Dans cet anneau, $x$ est un non-diviseur de zéro et l'image de $y$ dans
$A/xA$ donne une suite quasi-régulière. Mais il n'est pas vrai que la suite
$x, y$ soit quasi-régulière dans $A$, car $(x, y)/(x, y)^2$
n'est pas libre de rang deux sur $A/(x, y)$ : en effet,
$wx = 0$ dans $(x, y)/(x, y)^2$, tandis que $w$ n'est pas nul dans $A/(x, y)$.
Ainsi, l'analogue du
Lemme \ref{algebra-lemma-join-regular-sequences}
n'est pas vrai pour les suites quasi-régulières.
\end{remark}

\begin{lemma}
\label{algebra-lemma-quasi-regular-on-quotient}
Soit $R$ un anneau. Soit $J = (f_1, \ldots, f_r)$ un idéal de $R$.
Soit $M$ un $R$-module. Posons $\overline{R} = R/\bigcap_{n \geq 0} J^n$,
$\overline{M} = M/\bigcap_{n \geq 0} J^nM$, et notons
$\overline{f}_i$ l'image de $f_i$ dans $\overline{R}$.
Alors la suite $f_1, \ldots, f_r$ est $M$-quasi-régulière si et seulement si
la suite $\overline{f}_1, \ldots, \overline{f}_r$ est $\overline{M}$-quasi-régulière.
\end{lemma}

\begin{proof}
Cela est vrai parce que
$J^nM/J^{n + 1}M \cong
\overline{J}^n\overline{M}/\overline{J}^{n + 1}\overline{M}$.
\end{proof}





\section{Algèbres d'éclatement}
\label{algebra-section-blow-up}

\noindent
Dans cette section, nous présentons quelques observations élémentaires sur l'éclatement.

\begin{definition}
\label{algebra-definition-blow-up}
Soit $R$ un anneau.
Soit $I \subset R$ un idéal.
\begin{enumerate}
\item L'{\it algèbre d'éclatement}, ou {\it algèbre de Rees}, associée à
la paire $(R, I)$ est la $R$-algèbre graduée
$$
\text{Bl}_I(R) =
\bigoplus\nolimits_{n \geq 0} I^n =
R \oplus I \oplus I^2 \oplus \ldots
$$
où le facteur direct $I^n$ est placé en degré $n$.
\item Soit $a \in I$ un élément. Notons $a^{(1)}$ l'élément $a$
considéré comme un élément de degré $1$ de l'algèbre de Rees. L'
{\it algèbre d'éclatement affine} $R[\frac{I}{a}]$ est alors l'algèbre
$(\text{Bl}_I(R))_{(a^{(1)})}$ construite dans la section \ref{algebra-section-proj}.
\end{enumerate}
\end{definition}

\noindent
Autrement dit, un élément de $R[\frac{I}{a}]$ est représenté par
une expression de la forme $x/a^n$ avec $x \in I^n$. Deux représentants
$x/a^n$ et $y/a^m$ définissent le même élément si et seulement si
$a^k(a^mx - a^ny) = 0$ pour un certain $k \geq 0$.

\begin{lemma}
\label{algebra-lemma-affine-blowup}
Soit $R$ un anneau, $I \subset R$ un idéal et $a \in I$.
Soit $R' = R[\frac{I}{a}]$ l'algèbre d'éclatement affine. Alors
\begin{enumerate}
\item l'image de $a$ dans $R'$ est un non-diviseur de zéro,
\item $IR' = aR'$, et
\item $(R')_a = R_a$.
\end{enumerate}
\end{lemma}

\begin{proof}
Cela résulte immédiatement de la description de $R[\frac{I}{a}]$ ci-dessus.
\end{proof}

\begin{lemma}
\label{algebra-lemma-blowup-base-change}
Soit $R \to S$ un homomorphisme d'anneaux. Soit $I \subset R$ un idéal
et soit $a \in I$. Posons $J = IS$ et soit $b \in J$ l'image de $a$.
Alors $S[\frac{J}{b}]$ est le quotient de $S \otimes_R R[\frac{I}{a}]$
par l'idéal des éléments annulés par une puissance de $b$.
\end{lemma}

\begin{proof}
Soit $S'$ le quotient de $S \otimes_R R[\frac{I}{a}]$ par ses
éléments de torsion pour les puissances de $b$. L'homomorphisme d'anneaux
$$
S \otimes_R R[\textstyle{\frac{I}{a}}]
\longrightarrow
S[\textstyle{\frac{J}{b}}]
$$
est surjectif et annule la torsion pour les puissances de $a$, puisque $b$ est un non-diviseur de zéro
dans $S[\frac{J}{b}]$. On obtient donc une application surjective $S' \to S[\frac{J}{b}]$.
Pour voir que son noyau est trivial, construisons une application inverse. Soit
$z = y/b^n$ un élément de $S[\frac{J}{b}]$, c'est-à-dire $y \in J^n$.
Écrivons $y = \sum x_is_i$ avec $x_i \in I^n$ et $s_i \in S$.
On envoie $z$ sur la classe de $\sum s_i \otimes x_i/a^n$ dans
$S'$. Cette définition est bonne, car un élément du noyau de l'application
$S \otimes_R I^n \to J^n$ est annulé par $a^n$, donc a une image nulle dans $S'$.
\end{proof}

\begin{example}
\label{algebra-example-rees-algebra-polynomial}
Soit $R$ un anneau. Soit $P = R[t_1, \ldots, t_n]$ l'algèbre de polynômes.
Soit $I = (t_1, \ldots, t_n) \subset P$. Avec les notations de la
Définition \ref{algebra-definition-blow-up}, il existe un isomorphisme
$$
P[T_1, \ldots, T_n]/(t_iT_j - t_jT_i) \longrightarrow \text{Bl}_I(P)
$$
qui envoie $T_i$ sur $t_i^{(1)}$. Nous laissons au lecteur le soin de montrer que
cette application est bien définie. Puisque $I$ est engendré par $t_1, \ldots, t_n$,
on voit que notre application est surjective. Pour montrer qu'elle est injective,
il faut établir ceci : pour tout $e \geq 1$, le $P$-module $I^e$ est engendré
par les monômes $t^E = t_1^{e_1} \ldots x_n^{e_n}$, pour les multi-indices
$E = (e_1, \ldots, e_n)$ de degré $|E| = e$, soumis uniquement aux
relations $t_i t^E = t_j t^{E'}$ lorsque $|E| = |E'| = e$ et
$e_a + \delta_{a i} = e'_a + \delta_{a j},\ a = 1, \ldots, n$
(delta de Kronecker). Nous omettons les détails.
\end{example}

\begin{example}
\label{algebra-example-affine-blowup-algebra-polynomial}
Soit $R$ un anneau. Soit $P = R[t_1, \ldots, t_n]$ l'algèbre de polynômes.
Soit $I = (t_1, \ldots, t_n) \subset P$. Soit $a = t_1$. Avec les notations de la
Définition \ref{algebra-definition-blow-up}, il existe un isomorphisme
$$
P[x_2, \ldots, x_n]/(t_1x_2 - t_2, \ldots, t_1x_n - t_n)
\longrightarrow
\textstyle{P[\frac{I}{a}] = P[\frac{I}{t_1}]}
$$
qui envoie $x_i$ sur $t_i/t_1$. Nous laissons au lecteur le soin de montrer que
cette application est bien définie. Puisque $I$ est engendré par $t_1, \ldots, t_n$,
on voit que notre application est surjective. Pour montrer qu'elle est injective,
le lecteur peut observer que la source et le but de notre application sont
sans $t_1$-torsion et que l'application est un isomorphisme après inversion de
$t_1$ ; voir le Lemme \ref{algebra-lemma-affine-blowup}. Le lecteur peut aussi
utiliser la description de l'algèbre de Rees donnée dans
l'Exemple \ref{algebra-example-rees-algebra-polynomial}.
Nous omettons les détails.
\end{example}

\begin{lemma}
\label{algebra-lemma-affine-blowup-quotient-description}
Soit $R$ un anneau. Soit $I = (a_1, \ldots, a_n)$ un idéal de $R$.
Soit $a = a_1$. Il existe alors une surjection
$$
R[x_2, \ldots, x_n]/(a x_2 - a_2, \ldots, a x_n - a_n)
\longrightarrow
\textstyle{R[\frac{I}{a}]}
$$
dont le noyau est la torsion pour les puissances de $a$ dans la source.
\end{lemma}

\begin{proof}
Considérons l'homomorphisme d'anneaux $P = \mathbf{Z}[t_1, \ldots, t_n] \to R$
qui envoie $t_i$ sur $a_i$. Posons $J = (t_1, \ldots, t_n)$.
D'après l'Exemple \ref{algebra-example-affine-blowup-algebra-polynomial}, on a
$P[\frac{J}{t_1}] =
P[x_2, \ldots, x_n]/(t_1 x_2 - t_2, \ldots, t_1 x_n - t_n)$.
On conclut en appliquant le Lemme \ref{algebra-lemma-blowup-base-change} à l'application $P \to A$.
\end{proof}

\begin{lemma}
\label{algebra-lemma-blowup-in-principal}
Soit $R$ un anneau, $I \subset R$ un idéal et $a \in I$.
Posons $R' = R[\frac{I}{a}]$. Si $f \in R$ vérifie $V(f) = V(I)$,
alors l'image de $f$ dans $R'$ est un non-diviseur de zéro et $R'_f = R'_a = R_a$.
\end{lemma}

\begin{proof}
Nous utiliserons les résultats du Lemme \ref{algebra-lemma-affine-blowup}
sans les mentionner davantage.
L'hypothèse $V(f) = V(I)$ implique $V(fR') = V(IR') = V(aR')$.
Ainsi $a^n = fb$ et $f^m = ac$ pour certains $b, c \in R'$.
Le lemme en résulte.
\end{proof}

\begin{lemma}
\label{algebra-lemma-blowup-add-principal}
Soit $R$ un anneau, $I \subset R$ un idéal, $a \in I$ et $f \in R$.
Posons $R' = R[\frac{I}{a}]$ et $R'' = R[\frac{fI}{fa}]$. Alors
il existe un homomorphisme surjectif de $R$-algèbres $R' \to R''$ dont le noyau
est l'ensemble des éléments de $f$-torsion de $R'$.
\end{lemma}

\begin{proof}
L'application envoie $x/a^n$, pour $x \in I^n$, sur $f^nx/(fa)^n$.
On vérifie immédiatement que cette application est bien définie et surjective.
Puisque $af$ est un non-diviseur de zéro dans $R''$
(Lemme \ref{algebra-lemma-affine-blowup}), l'ensemble des éléments de torsion
pour les puissances de $f$ est envoyé sur zéro. Réciproquement, si $x \in R'$
et si $f^n x \not = 0$ pour tout $n > 0$, alors $(af)^n x \not = 0$
pour tout $n$, puisque $a$ est un non-diviseur de zéro dans $R'$. Il s'ensuit
que l'image de $x$ dans $R''$ n'est pas nulle, d'après la description de
$R''$ donnée après la Définition \ref{algebra-definition-blow-up}.
\end{proof}

\begin{lemma}
\label{algebra-lemma-blowup-reduced}
\begin{slogan}
La propriété d'être réduit est invariante par éclatement
\end{slogan}
Si $R$ est réduit, toute algèbre d'éclatement (affine) de $R$ est réduite.
\end{lemma}

\begin{proof}
Soit $I \subset R$ un idéal et soit $a \in I$. Supposons que $x/a^n$, avec
$x \in I^n$, soit un élément nilpotent de $R[\frac{I}{a}]$. Alors
$(x/a^n)^m = 0$. Ainsi $a^N x^m = 0$ dans $R$ pour un certain $N \geq 0$.
Après avoir augmenté $N$ si nécessaire, on peut supposer $N = me$ pour un certain
$e \geq 0$. Alors $(a^e x)^m = 0$ et, puisque $R$ est réduit, on obtient
$a^e x = 0$. Cela signifie que $x/a^n = 0$ dans $R[\frac{I}{a}]$.
\end{proof}

\begin{lemma}
\label{algebra-lemma-blowup-domain}
Soit $R$ un anneau intègre, $I \subset R$ un idéal et $a \in I$ un élément
non nul. Alors l'algèbre d'éclatement affine $R[\frac{I}{a}]$ est intègre.
\end{lemma}

\begin{proof}
Supposons que $x/a^n$, $y/a^m$, avec $x \in I^n$, $y \in I^m$,
soient des éléments de $R[\frac{I}{a}]$ dont le produit est nul.
Alors $a^N x y = 0$ dans $R$. Puisque $R$ est intègre, on en déduit
que $x = 0$ ou $y = 0$.
\end{proof}

\begin{lemma}
\label{algebra-lemma-blowup-dominant}
Soit $R$ un anneau. Soit $I \subset R$ un idéal. Soit $a \in I$.
Si $a$ n'appartient à aucun idéal premier minimal de $R$, alors
$\Spec(R[\frac{I}{a}]) \to \Spec(R)$ a une image dense.
\end{lemma}

\begin{proof}
Si $a^k x = 0$ pour $x \in R$, alors $x$ appartient à tous les
idéaux premiers minimaux de $R$ et est donc nilpotent ; voir le
Lemme \ref{algebra-lemma-Zariski-topology}.
Ainsi, le noyau de $R \to R[\frac{I}{a}]$ est constitué d'éléments
nilpotents. Le résultat découle donc du
Lemme \ref{algebra-lemma-image-dense-generic-points}.
\end{proof}

\begin{lemma}
\label{algebra-lemma-valuation-ring-colimit-affine-blowups}
Soit $(R, \mathfrak m)$ un anneau local intègre de corps des fractions $K$.
Soit $R \subset A \subset K$ un anneau de valuation qui domine $R$.
Alors
$$
A = \colim R[\textstyle{\frac{I}{a}}]
$$
est une colimite filtrante d'éclatements affines $R \to R[\frac{I}{a}]$ possédant
les propriétés suivantes :
\begin{enumerate}
\item $a \in I \subset \mathfrak m$,
\item $I$ est de type fini, et
\item l'anneau de la fibre de $R \to R[\frac{I}{a}]$ en $\mathfrak m$
n'est pas nul.
\end{enumerate}
\end{lemma}

\begin{proof}
Toute algèbre d'éclatement $R[\frac{I}{a}]$ est un anneau intègre contenu dans $K$ ; voir le
Lemme \ref{algebra-lemma-blowup-domain}. Le lemme dit simplement que $A$ est la
réunion filtrante de celles pour lesquelles $a \in I$ et les propriétés (1), (2), (3) sont satisfaites.
Si $R[\frac{I}{a}] \subset A$ et $R[\frac{J}{b}] \subset A$, alors
on a
$$
R[\textstyle{\frac{I}{a}}] \cup R[\textstyle{\frac{J}{b}}]  \subset
R[\textstyle{\frac{IJ}{ab}}] \subset A
$$
La première inclusion vient de ce que $x/a^n = b^nx/(ab)^n$, et la seconde
de ce que, si $z \in (IJ)^n$, alors $z = \sum x_iy_i$ avec $x_i \in I^n$
et $y_i \in J^n$, et donc $z/(ab)^n = \sum (x_i/a^n)(y_i/b^n)$
appartient à $A$.

\medskip\noindent
Considérons une partie finie $E \subset A$. Écrivons $E = \{e_1, \ldots, e_n\}$.
Choisissons $a \in R$ non nul tel que l'on puisse écrire $e_i = f_i/a$ pour
tout $i = 1, \ldots, n$. Posons $I = (f_1, \ldots, f_n, a)$.
Nous affirmons que $R[\frac{I}{a}] \subset A$. C'est clair, car un élément
de $R[\frac{I}{a}]$ peut être représenté par un polynôme en les éléments
$e_i$. Le lemme résulte immédiatement de cette observation.
\end{proof}










\section{Groupes Ext}
\label{algebra-section-ext}

\noindent
Dans cette section, nous faisons un peu d'algèbre homologique
afin d'établir quelques propriétés fondamentales de la
profondeur sur les anneaux locaux noethériens.

\begin{lemma}
\label{algebra-lemma-resolution-by-finite-free}
Soit $R$ un anneau. Soit $M$ un $R$-module.
\begin{enumerate}
\item Il existe un complexe exact
$$
\ldots \to F_2 \to F_1 \to F_0 \to M \to 0.
$$
où les $F_i$ sont des $R$-modules libres.
\item Si $R$ est noethérien et si $M$ est de type fini sur $R$, on
peut choisir le complexe de sorte que les $F_i$ soient libres de type fini.
Autrement dit, on peut trouver un complexe exact
$$
\ldots \to R^{\oplus n_2} \to R^{\oplus n_1} \to R^{\oplus n_0} \to M \to 0.
$$
\end{enumerate}
\end{lemma}

\begin{proof}
Expliquons seulement le cas noethérien.
On choisit d'abord une surjection $R^{n_0} \to M$.
Puis, après avoir construit un complexe exact de longueur
$e$, il suffit de choisir une surjection $R^{n_{e + 1}} \to
\Ker(R^{n_e} \to R^{n_{e-1}})$, ce qui est possible
puisque $R$ est noethérien.
\end{proof}

\begin{definition}
\label{algebra-definition-finite-free-resolution}
Soit $R$ un anneau. Soit $M$ un $R$-module.
\begin{enumerate}
\item Une {\it résolution} (à gauche) $F_\bullet \to M$ de $M$ est un complexe exact
$$
\ldots \to F_2 \to F_1 \to F_0 \to M \to 0
$$
de $R$-modules.
\item Une {\it résolution de $M$ par des $R$-modules libres} est une résolution
$F_\bullet \to M$ dans laquelle chaque $F_i$ est un $R$-module libre.
\item Une {\it résolution de $M$ par des $R$-modules libres de type fini} est une résolution
$F_\bullet \to M$ dans laquelle chaque $F_i$ est un $R$-module libre de type fini.
\end{enumerate}
\end{definition}

\noindent
Nous employons souvent la notation $F_{\bullet}$ pour désigner un complexe
de $R$-modules
$$
\ldots \to F_i \to F_{i-1} \to \ldots
$$
Dans ce cas, nous notons souvent $d_i$ ou $d_{F, i}$ l'application
$F_i \to F_{i-1}$. Dans cette section, nous supposerons toujours
que $F_0$ est le dernier terme non nul du complexe.
Le {\it $i$-ième groupe d'homologie du complexe} $F_{\bullet}$
est le groupe $H_i = \Ker(d_{F, i})/\Im(d_{F, i + 1})$.
Un {\it morphisme de complexes $\alpha : F_{\bullet} \to G_{\bullet}$}
est donné par des applications $\alpha_i : F_i \to G_i$ telles que
$\alpha_{i-1} \circ d_{F, i} = d_{G, i-1} \circ \alpha_i$.
Un tel morphisme induit en homologie une application $H_i(\alpha) :
H_i(F_{\bullet}) \to H_i(G_{\bullet})$. Si $\alpha, \beta
:  F_{\bullet} \to G_{\bullet}$ sont des morphismes de complexes, une
{\it homotopie} entre $\alpha$ et $\beta$ est donnée par
une famille d'applications $h_i : F_i \to G_{i + 1}$ telles que
$\alpha_i - \beta_i = d_{G, i + 1} \circ h_i +
h_{i-1} \circ d_{F, i}$.
Deux morphismes $\alpha, \beta : F_{\bullet} \to G_{\bullet}$ sont dits
{\it homotopes} s'il existe une homotopie entre $\alpha$
et $\beta$.

\medskip\noindent
Nous emploierons des notations tout à fait analogues pour les complexes
de la forme $F^{\bullet}$ qui s'écrivent
$$
\ldots \to F^i \xrightarrow{d^i} F^{i + 1} \to \ldots
$$
On dispose de morphismes de complexes, d'homotopies, etc.
Dans ce cas, nous posons $H^i(F^{\bullet}) =
\Ker(d^i)/\Im(d^{i - 1})$ et nous l'appelons le
{\it $i$-ième groupe de cohomologie}.

\begin{lemma}
\label{algebra-lemma-homotopic-equal-homology}
Deux morphismes de complexes homotopes quelconques induisent les mêmes applications sur les
groupes de (co)homologie.
\end{lemma}

\begin{proof}
Omis.
\end{proof}

\begin{lemma}
\label{algebra-lemma-compare-resolutions}
Soit $R$ un anneau. Soit $M \to N$ un homomorphisme de $R$-modules.
Soit $N_\bullet \to N$ une résolution quelconque.
Soit
$$
\ldots \to F_2 \to F_1 \to F_0 \to M
$$
un complexe de $R$-modules dans lequel chaque $F_i$ est un $R$-module libre. Alors
\begin{enumerate}
\item il existe un morphisme de complexes $F_\bullet \to N_\bullet$ tel que
$$
\xymatrix{
F_0 \ar[r] \ar[d] & M \ar[d] \\
N_0 \ar[r] & N
}
$$
soit commutatif, et
\item deux morphismes quelconques $\alpha, \beta : F_\bullet \to N_\bullet$ comme en (1)
sont homotopes.
\end{enumerate}
\end{lemma}

\begin{proof}
Démonstration de (1). Puisque $F_0$ est libre, on peut trouver une application $F_0 \to N_0$
qui relève l'application $F_0 \to M \to N$. On obtient une application induite
$F_1 \to F_0 \to N_0$ dont l'image est contenue dans l'image
de $N_1 \to N_0$. Puisque $F_1$ est libre, on peut la relever
en une application $F_1 \to N_1$. Celle-ci induit à son tour une application
$F_2 \to F_1 \to N_1$ dont l'image dans
$N_0$ est nulle. Comme $N_\bullet$ est exact, on voit que
l'image de cette application est contenue dans l'image
de $N_2 \to N_1$. On peut donc la relever en une application
$F_2 \to N_2$. On poursuit ainsi.

\medskip\noindent
Démonstration de (2). Pour montrer que $\alpha, \beta$ sont homotopes, il suffit
de montrer que la différence $\gamma = \alpha - \beta$ est homotope
à zéro. Remarquons que l'image de $\gamma_0 : F_0 \to N_0$
est contenue dans l'image de $N_1 \to N_0$. On peut donc relever
$\gamma_0$ en une application $h_0 : F_0 \to N_1$. Considérons l'application
$\gamma_1' = \gamma_1 - h_0 \circ d_{F, 1}$. Par le choix de $h_0$,
on voit que l'image de $\gamma_1'$ est contenue dans
le noyau de $N_1 \to N_0$. Comme $N_\bullet$ est exact,
on peut relever $\gamma_1'$ en une application $h_1 : F_1 \to N_2$.
À ce stade, on a $\gamma_1 = h_0 \circ d_{F, 1}
+ d_{N, 2} \circ h_1$. On poursuit ainsi.
\end{proof}

\noindent
Nous sommes maintenant en mesure de définir les groupes
$\Ext^i_R(M, N)$. Choisissons une résolution
$F_{\bullet}$ de $M$ par des $R$-modules libres, voir le Lemme
\ref{algebra-lemma-resolution-by-finite-free}. Considérons
le complexe (cohomologique)
$$
\Hom_R(F_\bullet, N) :
\Hom_R(F_0, N) \to
\Hom_R(F_1, N) \to
\Hom_R(F_2, N) \to \ldots
$$
On définit $\Ext^i_R(M, N)$, pour $i \geq 0$, comme le $i$-ième
groupe de cohomologie de ce complexe\footnote{À ce stade,
il conviendrait peut-être davantage de dire ``un'' plutôt
que ``le'' groupe Ext.}. Pour $i < 0$, on pose $\Ext^i_R(M, N) = 0$.
Avant de poursuivre, signalons que
$$
\Ext^0_R(M, N) = \Ker(\Hom_R(F_0, N) \to \Hom_R(F_1, N)) =
\Hom_R(M, N)
$$
car on peut appliquer la partie (1) du Lemme \ref{algebra-lemma-hom-exact} à
la suite exacte $F_1 \to F_0 \to M \to 0$.
Le lemme suivant précise
en quel sens cette définition est bien définie.

\begin{lemma}
\label{algebra-lemma-ext-welldefined}
Soit $R$ un anneau. Soient $M_1, M_2, N$ des $R$-modules.
Supposons que $F_{\bullet}$ soit une résolution libre du module $M_1$,
et que $G_{\bullet}$ soit une résolution libre du module $M_2$.
Soit $\varphi : M_1 \to M_2$ un homomorphisme de modules.
Soit $\alpha : F_{\bullet} \to G_{\bullet}$
un morphisme de complexes qui induit $\varphi$ sur
$M_1 = \Coker(d_{F, 1}) \to M_2 = \Coker(d_{G, 1})$,
voir le Lemme \ref{algebra-lemma-compare-resolutions}.
Alors les applications induites
$$
H^i(\alpha) :
H^i(\Hom_R(F_{\bullet}, N))
\longrightarrow
H^i(\Hom_R(G_{\bullet}, N))
$$
ne dépendent pas du choix de $\alpha$.
Si $\varphi$ est un isomorphisme, toutes les applications
$H^i(\alpha)$ sont des isomorphismes. Si $M_1 = M_2$, $F_\bullet = G_\bullet$, et
si $\varphi$ est l'identité, toutes les applications $H_i(\alpha)$ le sont aussi.
\end{lemma}

\begin{proof}
Un autre morphisme $\beta : F_{\bullet} \to G_{\bullet}$
qui induit $\varphi$ est homotope à $\alpha$ d'après le
Lemme \ref{algebra-lemma-compare-resolutions}. Les
applications $\Hom_R(F_\bullet, N) \to
\Hom_R(G_\bullet, N)$ sont donc homotopes.
L'indépendance résulte alors du
Lemme \ref{algebra-lemma-homotopic-equal-homology}.

\medskip\noindent
Supposons que $\varphi$ soit un isomorphisme.
Soit $\psi : M_2 \to M_1$ un inverse.
Choisissons $\beta : G_{\bullet} \to F_{\bullet}$
qui induit $\psi :
M_2 = \Coker(d_{G, 1}) \to M_1 = \Coker(d_{F, 1})$,
voir le Lemme \ref{algebra-lemma-compare-resolutions}.
Considérons maintenant l'application
$H^i(\alpha) \circ H^i(\beta) =
H^i(\alpha \circ \beta)$. D'après ce qui précède, l'application
$H^i(\alpha \circ \beta)$ est la {\it même}
que l'application $H^i(\text{id}_{G_{\bullet}}) = \text{id}$.
Il en va de même de la composée $H^i(\beta) \circ H^i(\alpha)$.
Ainsi $H^i(\alpha)$ et $H^i(\beta)$ sont inverses l'une de l'autre.
\end{proof}

\begin{lemma}
\label{algebra-lemma-long-exact-seq-ext}
Soit $R$ un anneau. Soit $M$ un $R$-module.
Soit $0 \to N' \to N \to N'' \to 0$ une
suite exacte courte. On obtient alors une suite exacte
longue
$$
\begin{matrix}
0
\to \Hom_R(M, N')
\to \Hom_R(M, N)
\to \Hom_R(M, N'')
\\
\phantom{0\ }
\to \Ext^1_R(M, N')
\to \Ext^1_R(M, N)
\to \Ext^1_R(M, N'')
\to \ldots
\end{matrix}
$$
\end{lemma}

\begin{proof}
Choisissons une résolution libre $F_{\bullet} \to M$.
Comme chacun des $F_i$ est libre, on obtient
une suite exacte courte de complexes
$$
0 \to
\Hom_R(F_{\bullet}, N') \to
\Hom_R(F_{\bullet}, N) \to
\Hom_R(F_{\bullet}, N'') \to
0
$$
La suite exacte longue résulte alors du
lemme du serpent appliqué à cette suite.
\end{proof}

\begin{lemma}
\label{algebra-lemma-reverse-long-exact-seq-ext}
Soit $R$ un anneau. Soit $N$ un $R$-module.
Soit $0 \to M' \to M \to M'' \to 0$ une
suite exacte courte. On obtient alors une suite exacte
longue
$$
\begin{matrix}
0
\to \Hom_R(M'', N)
\to \Hom_R(M, N)
\to \Hom_R(M', N)
\\
\phantom{0\ }
\to \Ext^1_R(M'', N)
\to \Ext^1_R(M, N)
\to \Ext^1_R(M', N)
\to \ldots
\end{matrix}
$$
\end{lemma}

\begin{proof}
Choisissons des systèmes de générateurs $\{m'_{i'}\}_{i' \in I'}$ et
$\{m''_{i''}\}_{i'' \in I''}$ de $M'$ et de $M''$.
Pour chaque $i'' \in I''$, choisissons un relèvement $\tilde m''_{i''} \in M$
de l'élément $m''_{i''} \in M''$. Posons $F' = \bigoplus_{i' \in I'} R$,
$F'' = \bigoplus_{i'' \in I''} R$ et $F = F' \oplus F''$.
En envoyant les générateurs de ces modules libres sur les générateurs choisis
correspondants, on obtient des homomorphismes surjectifs de $R$-modules $F' \to M'$,
$F'' \to M''$ et $F \to M$. On obtient un morphisme de suites exactes courtes
$$
\begin{matrix}
0 & \to & M' & \to & M & \to & M'' & \to & 0 \\
& & \uparrow & & \uparrow & & \uparrow \\
0 & \to & F' & \to & F & \to & F'' & \to & 0 \\
\end{matrix}
$$
Le lemme du serpent montre que la suite des noyaux
$0 \to K' \to K \to K'' \to 0$ est une suite exacte courte de $R$-modules.
On peut donc poursuivre indéfiniment ce procédé. Autrement dit,
on obtient une suite exacte courte de résolutions qui s'insère dans le diagramme
$$
\begin{matrix}
0 & \to & M' & \to & M & \to & M'' & \to & 0 \\
& & \uparrow & & \uparrow & & \uparrow \\
0 & \to & F_\bullet' & \to & F_\bullet & \to & F_\bullet'' & \to & 0 \\
\end{matrix}
$$
Comme chacune des suites $0 \to F'_n \to F_n \to F''_n \to 0$
est exacte scindée (par construction), on obtient une suite exacte courte de
complexes
$$
0 \to
\Hom_R(F''_{\bullet}, N) \to
\Hom_R(F_{\bullet}, N) \to
\Hom_R(F'_{\bullet}, N) \to
0
$$
en appliquant le foncteur $\Hom_R(-, N)$.
La suite exacte longue résulte alors du
lemme du serpent appliqué à cette suite.
\end{proof}

\begin{lemma}
\label{algebra-lemma-annihilate-ext}
Soit $R$ un anneau. Soient $M$, $N$ des $R$-modules.
Tout $x\in R$ tel que $xN = 0$ ou $xM = 0$
annule chacun des modules $\Ext^i_R(M, N)$.
\end{lemma}

\begin{proof}
Choisissons une résolution libre $F_{\bullet}$ de $M$.
Puisque $\Ext^i_R(M, N)$
est défini comme la cohomologie du complexe
$\Hom_R(F_{\bullet}, N)$, le lemme est
clair lorsque $xN = 0$. Si $xM = 0$, alors
la multiplication par $x$ sur $F_{\bullet}$
relève l'application nulle sur $M$. Ainsi, d'après le Lemme
\ref{algebra-lemma-ext-welldefined}, elle induit
la même application sur les groupes Ext que l'application
nulle.
\end{proof}

\begin{lemma}
\label{algebra-lemma-ext-noetherian}
Soit $R$ un anneau noethérien. Soient $M$, $N$ des $R$-modules de type fini.
Alors $\Ext^i_R(M, N)$ est un $R$-module de type fini pour tout $i$.
\end{lemma}

\begin{proof}
Cela résulte de ce que $\Ext^i_R(M, N)$ se calcule comme les
groupes de cohomologie d'un complexe $\Hom_R(F_\bullet, N)$
où chaque $F_n$ est un $R$-module libre de type fini ; voir le
Lemme \ref{algebra-lemma-resolution-by-finite-free}.
\end{proof}




\section{Profondeur}
\label{algebra-section-depth}

\noindent
Voici notre définition.

\begin{definition}
\label{algebra-definition-depth}
Soit $R$ un anneau, et soit $I \subset R$ un idéal. Soit $M$ un $R$-module de type fini.
La {\it profondeur relativement à $I$} de $M$, notée $\text{depth}_I(M)$, est définie comme suit :
\begin{enumerate}
\item si $IM \not = M$, alors $\text{depth}_I(M)$ est la borne supérieure dans
$\{0, 1, 2, \ldots, \infty\}$ des longueurs des suites $M$-régulières dans $I$,
\item si $IM = M$, on pose $\text{depth}_I(M) = \infty$.
\end{enumerate}
Si $(R, \mathfrak m)$ est local, on appelle simplement $\text{depth}_{\mathfrak m}(M)$ la
{\it profondeur} de $M$.
\end{definition}

\noindent
Explication. D'après la Définition \ref{algebra-definition-regular-sequence}, la suite vide
n'est pas une suite régulière sur le module nul, mais, en pratique,
il s'avère commode de fixer la profondeur du module $0$
égale à $+\infty$. Remarquons que, si $I = R$, alors $\text{depth}_I(M) = \infty$
pour tout $R$-module de type fini $M$. Si $I$ est contenu dans le radical de Jacobson
de $R$ (par exemple, si $R$ est local et si $I \subset \mathfrak m_R$), alors
$M \not = 0 \Rightarrow IM \not = M$ d'après le lemme de Nakayama.
Un module $M$ a une profondeur relativement à $I$ égale à $0$ si et seulement si $M$ est non nul et si $I$
ne contient aucun non-diviseur de zéro sur $M$.

\medskip\noindent
L'Exemple \ref{algebra-example-global-regular} montre que la profondeur ne se comporte pas
bien, même si l'anneau est noethérien, et l'Exemple
\ref{algebra-example-local-regular} montre qu'elle ne se comporte pas
bien si l'anneau est local mais non noethérien.
Nous verrons que la profondeur se comporte bien si l'anneau est local noethérien.

\begin{lemma}
\label{algebra-lemma-depth-weak-sequence}
Soit $R$ un anneau, $I \subset R$ un idéal, et $M$ un $R$-module de type fini.
Alors $\text{depth}_I(M)$ est égale à la borne supérieure des longueurs des
suites $f_1, \ldots, f_r \in I$ telles que $f_i$ soit un non-diviseur de zéro
sur $M/(f_1, \ldots, f_{i - 1})M$.
\end{lemma}

\begin{proof}
Supposons que $IM = M$. Alors le Lemme \ref{algebra-lemma-NAK} montre qu'il existe
$f \in I$ tel que $f : M \to M$ soit $\text{id}_M$. Ainsi
$f, 0, 0, 0, \ldots$ est une suite infinie de
non-diviseurs de zéro successifs, et les deux quantités coïncident dans ce cas.
Si $IM \not =  M$, une suite comme dans le lemme
est une suite $M$-régulière, et les deux quantités coïncident là encore.
\end{proof}

\begin{lemma}
\label{algebra-lemma-bound-depth}
Soit $(R, \mathfrak m)$ un anneau local noethérien.
Soit $M$ un $R$-module de type fini non nul.
Alors $\dim(\text{Supp}(M)) \geq \text{depth}(M)$.
\end{lemma}

\begin{proof}
La démonstration se fait par récurrence sur $\dim(\text{Supp}(M))$.
Si $\dim(\text{Supp}(M)) = 0$, alors
$\text{Supp}(M) = \{\mathfrak m\}$, d'où $\text{Ass}(M) = \{\mathfrak m\}$
(d'après les Lemmes \ref{algebra-lemma-ass-support} et \ref{algebra-lemma-ass-zero}), et donc
la profondeur de $M$ est nulle, par exemple d'après le
Lemme \ref{algebra-lemma-ideal-nonzerodivisor}.
Pour l'étape de récurrence, supposons $\dim(\text{Supp}(M)) > 0$.
Soit $f_1, \ldots, f_d$ une suite d'éléments de $\mathfrak m$
telle que $f_i$ soit un non-diviseur de zéro sur $M/(f_1, \ldots, f_{i - 1})M$.
D'après le Lemme \ref{algebra-lemma-depth-weak-sequence}, il suffit de montrer que
$\dim(\text{Supp}(M)) \geq d$. On peut supposer
$d > 0$, sans quoi le lemme est établi. D'après le
Lemme \ref{algebra-lemma-one-equation-module},
on a $\dim(\text{Supp}(M/f_1M)) = \dim(\text{Supp}(M)) - 1$.
Par récurrence, on conclut que $\dim(\text{Supp}(M/f_1M)) \geq d - 1$,
comme souhaité.
\end{proof}

\begin{lemma}
\label{algebra-lemma-depth-finite-noetherian}
Soit $R$ un anneau noethérien, $I \subset R$ un idéal, et $M$ un
$R$-module non nul de type fini tel que $IM \not = M$. Alors
$\text{depth}_I(M) < \infty$.
\end{lemma}

\begin{proof}
Puisque $M/IM$ est non nul, on peut choisir $\mathfrak p \in \text{Supp}(M/IM)$
d'après le Lemme \ref{algebra-lemma-support-zero}. Alors $(M/IM)_\mathfrak p \not = 0$,
ce qui implique $I \subset \mathfrak p$ et, de plus,
$M_\mathfrak p \not = IM_\mathfrak p$, puisque la localisation est exacte.
Soit $f_1, \ldots, f_r \in I$ une suite $M$-régulière.
Alors $M_\mathfrak p/(f_1, \ldots, f_r)M_\mathfrak p$ est
non nul, car $(f_1, \ldots, f_r) \subset I$. La localisation étant
plate, les images de $f_1, \ldots, f_r$ forment une
suite $M_\mathfrak p$-régulière dans $I_\mathfrak p$. Comme cela vaut
pour toute suite $M$-régulière dans $I$, on conclut que
$\text{depth}_I(M) \leq \text{depth}_{I_\mathfrak p}(M_\mathfrak p)$.
Ce dernier membre est $\leq \text{depth}(M_\mathfrak p)$, qui est
$< \infty$ d'après le Lemme \ref{algebra-lemma-bound-depth}.
\end{proof}

\begin{lemma}
\label{algebra-lemma-depth-ext}
Soit $R$ un anneau local noethérien d'idéal maximal $\mathfrak m$.
Soit $M$ un $R$-module de type fini non nul. Alors $\text{depth}(M)$
est égale au plus petit entier $i$ tel que
$\Ext^i_R(R/\mathfrak m, M)$ soit non nul.
\end{lemma}

\begin{proof}
Notons $\delta(M)$ la profondeur de $M$ et $i(M)$ le
plus petit entier $i$ tel que $\Ext^i_R(R/\mathfrak m, M)$
soit non nul. Nous verrons dans un instant que $i(M) < \infty$.
D'après le Lemme \ref{algebra-lemma-ideal-nonzerodivisor}, on a
$\delta(M) = 0$ si et seulement si $i(M) = 0$, car
$\mathfrak m \in \text{Ass}(M)$ signifie exactement
que $i(M) = 0$. Ainsi, si $\delta(M)$ ou $i(M)$ est $> 0$, on peut
choisir $x \in \mathfrak m$ tel que (a) $x$ soit un non-diviseur de zéro
sur $M$, et (b) $\text{depth}(M/xM) = \delta(M) - 1$.
Considérons la suite exacte longue
de groupes Ext associée à la suite exacte courte
$0 \to M \to M \to M/xM \to 0$ par le Lemme \ref{algebra-lemma-long-exact-seq-ext} :
$$
\begin{matrix}
0
\to \Hom_R(\kappa, M)
\to \Hom_R(\kappa, M)
\to \Hom_R(\kappa, M/xM)
\\
\phantom{0\ }
\to \Ext^1_R(\kappa, M)
\to \Ext^1_R(\kappa, M)
\to \Ext^1_R(\kappa, M/xM)
\to \ldots
\end{matrix}
$$
Puisque $x \in \mathfrak m$, toutes les applications $\Ext^i_R(\kappa, M)
\to \Ext^i_R(\kappa, M)$ sont nulles ; voir le
Lemme \ref{algebra-lemma-annihilate-ext}.
Il est donc clair que $i(M/xM) = i(M) - 1$. Une récurrence sur
$\delta(M)$ achève la démonstration.
\end{proof}

\begin{lemma}
\label{algebra-lemma-depth-in-ses}
Soit $R$ un anneau local noethérien. Soit $0 \to N' \to N \to N'' \to 0$
une suite exacte courte de $R$-modules non nuls de type fini.
\begin{enumerate}
\item
$\text{depth}(N) \geq \min\{\text{depth}(N'), \text{depth}(N'')\}$
\item
$\text{depth}(N'') \geq \min\{\text{depth}(N), \text{depth}(N') - 1\}$
\item
$\text{depth}(N') \geq \min\{\text{depth}(N), \text{depth}(N'') + 1\}$
\end{enumerate}
\end{lemma}

\begin{proof}
Utilisons la caractérisation de la profondeur par les groupes Ext
$\Ext^i(\kappa, N)$, voir le Lemme \ref{algebra-lemma-depth-ext},
ainsi que la suite exacte longue de cohomologie
$$
\begin{matrix}
0
\to \Hom_R(\kappa, N')
\to \Hom_R(\kappa, N)
\to \Hom_R(\kappa, N'')
\\
\phantom{0\ }
\to \Ext^1_R(\kappa, N')
\to \Ext^1_R(\kappa, N)
\to \Ext^1_R(\kappa, N'')
\to \ldots
\end{matrix}
$$
du Lemme \ref{algebra-lemma-long-exact-seq-ext}.
\end{proof}

\begin{lemma}
\label{algebra-lemma-depth-drops-by-one}
Soit $R$ un anneau local noethérien et $M$ un $R$-module de type fini non nul.
\begin{enumerate}
\item Si $x \in \mathfrak m$ est un non-diviseur de zéro sur $M$, alors
$\text{depth}(M/xM) = \text{depth}(M) - 1$.
\item Toute suite $M$-régulière $x_1, \ldots, x_r$ peut être prolongée en une
suite $M$-régulière de longueur $\text{depth}(M)$.
\end{enumerate}
\end{lemma}

\begin{proof}
La partie (2) est une conséquence formelle de la partie (1). Soit $x \in R$ comme en (1).
La suite exacte courte $0 \to M \to M \to M/xM \to 0$
et le Lemme \ref{algebra-lemma-depth-in-ses} montrent que la profondeur diminue d'au plus 1.
D'autre part, si $x_1, \ldots, x_r \in \mathfrak m$
est une suite régulière pour $M/xM$, alors $x, x_1, \ldots, x_r$
est une suite régulière pour $M$. On voit donc que la profondeur diminue
d'au moins 1.
\end{proof}

\begin{lemma}
\label{algebra-lemma-inherit-minimal-primes}
Soit $(R, \mathfrak m)$ un anneau local noethérien et $M$ un $R$-module de type fini.
Soient $x \in \mathfrak m$, $\mathfrak p \in \text{Ass}(M)$, et $\mathfrak q$
minimal au-dessus de $\mathfrak p + (x)$. Alors $\mathfrak q \in \text{Ass}(M/x^nM)$
pour un certain $n \geq 1$.
\end{lemma}

\begin{proof}
Choisissons un sous-module $N \subset M$ tel que $N \cong R/\mathfrak p$.
D'après le lemme d'Artin-Rees (Lemme \ref{algebra-lemma-Artin-Rees}),
on peut choisir $n > 0$ tel que $N \cap x^nM \subset xN$.
Soit $\overline{N} \subset M/x^nM$ l'image de $N \to M \to M/x^nM$.
D'après le Lemme \ref{algebra-lemma-ass}, il suffit de montrer que
$\mathfrak q \in \text{Ass}(\overline{N})$.
Par notre choix de $n$, il existe une surjection
$\overline{N} \to N/xN = R/\mathfrak p + (x)$,
et donc $\mathfrak q$ appartient au support de $\overline{N}$.
Puisque $\overline{N}$ est annulé par $x^n$ et $\mathfrak p$, on voit que
$\mathfrak q$ est minimal parmi les idéaux premiers du support de $\overline{N}$.
Ainsi $\mathfrak q$ est un idéal premier associé à $\overline{N}$ d'après le
Lemme \ref{algebra-lemma-ass-minimal-prime-support}.
\end{proof}

\begin{lemma}
\label{algebra-lemma-depth-dim-associated-primes}
Soit $(R, \mathfrak m)$ un anneau local noethérien et $M$ un $R$-module de type fini.
Pour $\mathfrak p \in \text{Ass}(M)$, on a
$\dim(R/\mathfrak p) \geq \text{depth}(M)$.
\end{lemma}

\begin{proof}
Si $\mathfrak m \in \text{Ass}(M)$, il existe un élément non nul
$x \in M$ annulé par tous les éléments de $\mathfrak m$.
Ainsi $\text{depth}(M) = 0$. En particulier, le lemme est vrai dans ce cas.

\medskip\noindent
Si $\text{depth}(M) = 1$, alors le premier paragraphe
montre que $\mathfrak m \not \in \text{Ass}(M)$.
Ainsi $\dim(R/\mathfrak p) \geq 1$ pour tout $\mathfrak p \in \text{Ass}(M)$,
et le lemme est encore vrai dans ce cas.

\medskip\noindent
Démontrons le lemme en général par récurrence sur $\text{depth}(M)$,
que l'on peut et va supposer $> 1$. Choisissons $x \in \mathfrak m$ qui soit
un non-diviseur de zéro sur $M$. Remarquons que $x \not \in \mathfrak p$
(Lemme \ref{algebra-lemma-ass-zero-divisors}).
D'après le Lemme \ref{algebra-lemma-one-equation}, on a
$\dim(R/\mathfrak p + (x)) = \dim(R/\mathfrak p) - 1$.
Il existe donc un idéal premier $\mathfrak q$ minimal au-dessus de $\mathfrak p + (x)$ tel que
$\dim(R/\mathfrak q) = \dim(R/\mathfrak p) - 1$ (petit argument omis ;
indication : la dimension d'un anneau local noethérien $A$ est le maximum
des dimensions de $A/\mathfrak r$ lorsque $\mathfrak r$ parcourt les idéaux premiers
minimaux de $A$). Choisissons $n$ comme dans le
Lemme \ref{algebra-lemma-inherit-minimal-primes}, de sorte que
$\mathfrak q$ soit un idéal premier associé à $M/x^nM$.
On peut appliquer l'hypothèse de récurrence à $M/x^nM$ et $\mathfrak q$,
car $\text{depth}(M/x^nM) = \text{depth}(M) - 1$ d'après le
Lemme \ref{algebra-lemma-depth-drops-by-one}. On obtient
$\dim(R/\mathfrak q) \geq \text{depth}(M/x^nM)$, ce qui conclut.
\end{proof}

\begin{lemma}
\label{algebra-lemma-depth-localization}
Soit $R$ un anneau local noethérien et $M$ un $R$-module de type fini.
Pour un idéal premier $\mathfrak p \subset R$, on a
$\text{depth}(M_\mathfrak p) + \dim(R/\mathfrak p) \geq \text{depth}(M)$.
\end{lemma}

\begin{proof}
Si $M_\mathfrak p = 0$, alors $\text{depth}(M_\mathfrak p) = \infty$, et
le lemme est établi.
Si $\text{depth}(M) \leq \dim(R/\mathfrak p)$, alors le lemme est vrai.
Si $\text{depth}(M) > \dim(R/\mathfrak p)$, alors $\mathfrak p$ n'est contenu
dans aucun idéal premier associé $\mathfrak q$ à $M$, d'après le
Lemme \ref{algebra-lemma-depth-dim-associated-primes}.
On peut donc trouver $x \in \mathfrak p$ qui n'appartient à aucun
idéal premier associé à $M$, d'après le Lemme \ref{algebra-lemma-silly} et le
Lemme \ref{algebra-lemma-finite-ass}. Alors $x$ est un non-diviseur de zéro
sur $M$ ; voir le Lemme \ref{algebra-lemma-ass-zero-divisors}.
Ainsi $\text{depth}(M/xM) = \text{depth}(M) - 1$ et
$\text{depth}(M_\mathfrak p / x M_\mathfrak p) =
\text{depth}(M_\mathfrak p) - 1$, pourvu que $M_\mathfrak p$ soit non nul ;
voir le Lemme \ref{algebra-lemma-depth-drops-by-one}.
On conclut donc par récurrence sur $\text{depth}(M)$.
\end{proof}

\begin{lemma}
\label{algebra-lemma-depth-goes-down-finite}
Soit $(R, \mathfrak m)$ un anneau local noethérien. Soit $R \to S$
un homomorphisme fini d'anneaux. Soient $\mathfrak m_1, \ldots, \mathfrak m_n$
les idéaux maximaux de $S$. Soit $N$ un $S$-module de type fini.
Alors
$$
\min\nolimits_{i = 1, \ldots, n} \text{depth}(N_{\mathfrak m_i}) =
\text{depth}_\mathfrak m(N)
$$
\end{lemma}

\begin{proof}
D'après les Lemmes \ref{algebra-lemma-integral-no-inclusion}, \ref{algebra-lemma-integral-going-up},
et le Lemme \ref{algebra-lemma-finite-finite-fibres}, les idéaux maximaux de
$S$ sont exactement les idéaux premiers de $S$ au-dessus de $\mathfrak m$, et
ils sont en nombre fini. L'énoncé du lemme a donc
un sens. Nous allons démontrer le lemme par récurrence sur
$k = \min\nolimits_{i = 1, \ldots, n} \text{depth}(N_{\mathfrak m_i})$.
Si $k = 0$, alors $\text{depth}(N_{\mathfrak m_i}) = 0$ pour un certain $i$.
D'après le Lemme \ref{algebra-lemma-depth-ext}, cela signifie que
$\mathfrak m_i S_{\mathfrak m_i}$ est un idéal premier associé
à $N_{\mathfrak m_i}$, et donc que $\mathfrak m_i$ est un
idéal premier associé à $N$ (Lemme \ref{algebra-lemma-localize-ass}).
D'après le Lemme \ref{algebra-lemma-ass-functorial-Noetherian}, on voit que
$\mathfrak m$ est un idéal premier associé à $N$ comme $R$-module.
Ainsi $\text{depth}_\mathfrak m(N) = 0$. Cela établit le cas initial.
Si $k > 0$, alors $\mathfrak m_i \not \in \text{Ass}_S(N)$.
Ainsi $\mathfrak m \not \in \text{Ass}_R(N)$, encore d'après le
Lemme \ref{algebra-lemma-ass-functorial-Noetherian}.
On peut donc trouver $f \in \mathfrak m$ qui n'est pas un diviseur de zéro sur
$N$ ; voir le Lemme \ref{algebra-lemma-ideal-nonzerodivisor}. D'après le
Lemme \ref{algebra-lemma-depth-drops-by-one},
toutes les profondeurs diminuent exactement de $1$ lorsque l'on passe de $N$ à
$N/fN$, et l'hypothèse de récurrence conclut.
\end{proof}




\section{Fonctorialités de Ext}
\label{algebra-section-functoriality-ext}

\noindent
Dans cette section, nous examinons brièvement la fonctorialité
de $\Ext$ par rapport au changement d'anneau, etc.
Voici une liste de points à développer.
\begin{enumerate}
\item Étant donnés $R \to R'$, un $R$-module
$M$ et un $R'$-module $N'$,
le $R$-module $\Ext^i_R(M, N')$
possède une structure naturelle de $R'$-module. De plus, il existe un
homomorphisme canonique $R'$-linéaire $\Ext^i_{R'}(M \otimes_R R', N') \to
\Ext^i_R(M, N')$.
\item Étant donnés $R \to R'$ et des $R$-modules $M$, $N$, il existe un
homomorphisme naturel de $R$-modules
$\Ext^i_R(M, N) \to \text{Ext}^i_R(M, N \otimes_R R')$.
\end{enumerate}

\begin{lemma}
\label{algebra-lemma-flat-base-change-ext}
Étant donnés un homomorphisme plat d'anneaux $R \to R'$, un $R$-module $M$ et un
$R'$-module $N'$, l'homomorphisme naturel
$$
\Ext^i_{R'}(M \otimes_R R', N') \to \text{Ext}^i_R(M, N')
$$
est un isomorphisme pour $i \geq 0$.
\end{lemma}

\begin{proof}
Choisissons une résolution libre $F_\bullet$ de $M$.
Comme $R \to R'$ est plat, $F_\bullet \otimes_R R'$ est
une résolution libre de $M \otimes_R R'$ sur $R'$.
L'énoncé affirme que l'homomorphisme
$$
\Hom_{R'}(F_\bullet \otimes_R R', N') \to
\Hom_R(F_\bullet, N')
$$
induit un isomorphisme sur les groupes d'homologie, ce qui est vrai puisqu'il
est un isomorphisme de complexes d'après le
Lemme \ref{algebra-lemma-adjoint-tensor-restrict}.
\end{proof}






\section{Une application des groupes Ext}
\label{algebra-section-ext-application}

\noindent
La voici.

\begin{lemma}
\label{algebra-lemma-split-injection-after-completion}
Soit $R$ un anneau noethérien. Soit $I \subset R$ un idéal
contenu dans le radical de Jacobson de $R$.
Soit $N \to M$ un homomorphisme de $R$-modules de type fini.
Supposons qu'il existe des entiers $n$ arbitrairement grands tels que
$N/I^nN \to M/I^nM$ soit une injection scindée.
Alors $N \to M$ est une injection scindée.
\end{lemma}

\begin{proof}
Supposons que $\varphi : N \to M$ satisfasse les hypothèses du lemme.
Remarquons que cela implique $\Ker(\varphi) \subset I^nN$
pour des entiers $n$ arbitrairement grands. Ainsi, d'après le
Lemme \ref{algebra-lemma-intersection-powers-ideal-module}, on voit que $\varphi$
est injectif. Posons $Q = M/N$, de sorte que l'on ait une suite exacte courte
$$
0 \to N \to M \to Q \to 0.
$$
Soit
$$
F_2 \xrightarrow{d_2} F_1 \xrightarrow{d_1} F_0 \to Q \to 0
$$
une résolution libre de type fini de $Q$. On peut choisir une application
$\alpha : F_0 \to M$ qui relève l'application $F_0 \to Q$. Elle induit une application
$\beta : F_1 \to N$ telle que $\beta \circ d_2 = 0$. L'extension
ci-dessus est scindée si et seulement s'il existe une application $\gamma : F_0 \to N$
telle que $\beta = \gamma \circ d_1$. Autrement dit, la classe de
$\beta$ dans $\Ext^1_R(Q, N)$ est l'obstruction au scindage
de la suite exacte courte ci-dessus.

\medskip\noindent
Supposons que $n$ soit un entier assez grand tel que $N/I^nN \to M/I^nM$ soit une
injection scindée. Cela implique que
$$
0 \to N/I^nN \to M/I^nM \to Q/I^nQ \to 0.
$$
est encore une suite exacte courte. De plus, la suite
$$
F_1/I^nF_1 \xrightarrow{d_1} F_0/I^nF_0 \to Q/I^nQ \to 0
$$
est encore exacte. En raisonnant comme ci-dessus, on voit que l'application
$\overline{\beta} : F_1/I^nF_1 \to N/I^nN$
induite par $\beta$ est égale à $\overline{\gamma_n} \circ d_1$ pour une certaine
application $\overline{\gamma_n} : F_0/I^nF_0 \to N/I^nN$.
Comme $F_0$ est libre, on peut relever $\overline{\gamma_n}$ en une application
$\gamma_n : F_0 \to N$, et l'on voit alors que
$\beta - \gamma_n \circ d_1$ est une application de $F_1$ dans $I^nN$.
Autrement dit, on conclut que
$$
\beta \in
\Im\Big(\Hom_R(F_0, N) \to \Hom_R(F_1, N)\Big) + I^n\Hom_R(F_1, N).
$$
pour cet entier $n$.

\medskip\noindent
Puisque, par hypothèse, cette propriété vaut pour des entiers $n$ arbitrairement grands,
on conclut que l'image de $\beta$ dans le conoyau de
$\Hom_R(F_0, N) \to \Hom_R(F_1, N)$ est nulle d'après le
Lemme \ref{algebra-lemma-intersection-powers-ideal-module}. Ainsi
$\beta$ appartient à l'image de l'application $\Hom_R(F_0, N) \to \Hom_R(F_1, N)$,
comme souhaité.
\end{proof}











\section{Groupes Tor et platitude}
\label{algebra-section-tor}

\noindent
Dans cette section, nous utilisons une partie de l'algèbre homologique
développée dans la section précédente pour expliquer ce que sont
les groupes Tor. Supposons donc que $R$ soit un anneau
et que $M$, $N$ soient deux $R$-modules. Choisissons
une résolution $F_\bullet$ de $M$ par des $R$-modules libres.
Voir le Lemme \ref{algebra-lemma-resolution-by-finite-free}.
Considérons le complexe homologique
$$
F_\bullet \otimes_R N
:
\ldots
\to F_2 \otimes_R N
\to F_1 \otimes_R N
\to F_0 \otimes_R N
$$
On définit $\text{Tor}^R_i(M, N)$ comme le $i$-ième groupe d'homologie
de ce complexe. Le lemme suivant précise en
quel sens cette définition est bien définie.

\begin{lemma}
\label{algebra-lemma-tor-welldefined}
Soit $R$ un anneau. Soient $M_1, M_2, N$ des $R$-modules.
Supposons que $F_\bullet$ soit une résolution libre du
module $M_1$ et que $G_\bullet$ soit une résolution libre du
module $M_2$. Soit $\varphi : M_1 \to M_2$
un homomorphisme de modules. Soit $\alpha : F_\bullet \to G_\bullet$
un morphisme de complexes qui induit $\varphi$ sur
$M_1 = \Coker(d_{F, 1}) \to M_2 = \Coker(d_{G, 1})$,
voir le Lemme \ref{algebra-lemma-compare-resolutions}.
Alors les applications induites
$$
H_i(\alpha) :
H_i(F_\bullet \otimes_R N)
\longrightarrow
H_i(G_\bullet \otimes_R N)
$$
ne dépendent pas du choix de $\alpha$. Si $\varphi$
est un isomorphisme, toutes les applications $H_i(\alpha)$ le sont aussi.
Si $M_1 = M_2$, $F_\bullet = G_\bullet$, et
si $\varphi$ est l'identité, toutes les applications $H_i(\alpha)$ le sont aussi.
\end{lemma}

\begin{proof}
La démonstration de ce lemme est identique à celle du Lemme
\ref{algebra-lemma-ext-welldefined}.
\end{proof}

\noindent
Ce lemme implique non seulement que les modules Tor sont bien définis,
mais il fournit aussi la fonctorialité des constructions
$(M, N) \mapsto \text{Tor}_i^R(M, N)$ en la première variable. Bien entendu,
la fonctorialité en la seconde variable est évidente. Nous laissons au
lecteur le soin de vérifier que chacun des $\text{Tor}_i^R$ est en fait
un foncteur
$$
\text{Mod}_R \times \text{Mod}_R \to \text{Mod}_R.
$$
Ici $\text{Mod}_R$ désigne la catégorie des $R$-modules et,
pour la définition de la catégorie produit,
voir Catégories, Définition \ref{categories-definition-product-category}.
Plus précisément, étant donnés des morphismes de $R$-modules $M_1 \to M_2$
et $N_1 \to N_2$, on obtient un diagramme commutatif
$$
\xymatrix{
\text{Tor}_i^R(M_1, N_1) \ar[r] \ar[d] &
\text{Tor}_i^R(M_1, N_2) \ar[d] \\
\text{Tor}_i^R(M_2, N_1) \ar[r] &
\text{Tor}_i^R(M_2, N_2) \\
}
$$

\begin{lemma}
\label{algebra-lemma-long-exact-sequence-tor}
Soit $R$ un anneau et soit $M$ un $R$-module.
Supposons que $0 \to N' \to N \to N'' \to 0$ soit une suite
exacte courte de $R$-modules. Il existe une suite exacte
longue
$$
\text{Tor}_1^R(M, N')
\to \text{Tor}_1^R(M, N)
\to \text{Tor}_1^R(M, N'')
\to
M \otimes_R N'
\to M \otimes_R N
\to M \otimes_R N''
\to 0
$$
\end{lemma}

\begin{proof}
La démonstration est la même que celle du
Lemme \ref{algebra-lemma-long-exact-seq-ext}.
\end{proof}

\noindent
Considérons un bicomplexe homologique de $R$-modules
$$
\xymatrix{
\ldots \ar[r]^d &
A_{2, 0} \ar[r]^d &
A_{1, 0} \ar[r]^d &
A_{0, 0} \\
\ldots \ar[r]^d &
A_{2, 1} \ar[r]^d \ar[u]^\delta &
A_{1, 1} \ar[r]^d \ar[u]^\delta &
A_{0, 1} \ar[u]^\delta \\
\ldots \ar[r]^d &
A_{2, 2} \ar[r]^d \ar[u]^\delta &
A_{1, 2} \ar[r]^d \ar[u]^\delta &
A_{0, 2} \ar[u]^\delta \\
&
\ldots \ar[u]^\delta &
\ldots \ar[u]^\delta &
\ldots \ar[u]^\delta \\
}
$$
Cela signifie que $d_{i, j} : A_{i, j} \to A_{i-1, j}$
et $\delta_{i, j} : A_{i, j} \to A_{i, j-1}$ possèdent les propriétés
suivantes :
\begin{enumerate}
\item Toute composée de deux $d_{i, j}$ est nulle. Autrement dit,
les lignes du bicomplexe sont des complexes.
\item Toute composée de deux $\delta_{i, j}$ est nulle. Autrement dit,
les colonnes du bicomplexe sont des complexes.
\item Pour toute paire $(i, j)$, on a $\delta_{i-1, j} \circ d_{i, j}
= d_{i, j-1} \circ \delta_{i, j}$. Autrement dit, tous les carrés
sont commutatifs.
\end{enumerate}
La bonne construction consiste à associer une suite spectrale à
chacun de ces bicomplexes. Cependant, pour le moment, nous pouvons nous contenter
d'une construction légèrement plus simple.

\medskip\noindent
Plus précisément, pour les seuls besoins de cette section, étant donné un bicomplexe
$(A_{\bullet, \bullet}, d, \delta)$, posons
$R(A)_j = \Coker(A_{1, j} \to A_{0, j})$ et
$U(A)_i = \Coker(A_{i, 1} \to A_{i, 0})$. (Les lettres
$R$ et $U$ évoquent les mots anglais pour « droite » et « haut ».)
On munit $R(A)_\bullet$ d'une structure de complexe
au moyen des applications $\delta$. De même, on munit $U(A)_\bullet$
d'une structure de complexe au moyen des applications $d$.
Autrement dit, on obtient l'immense diagramme commutatif suivant
$$
\xymatrix{
\ldots \ar[r]^d &
U(A)_2 \ar[r]^d &
U(A)_1 \ar[r]^d &
U(A)_0 &
\\
\ldots \ar[r]^d &
A_{2, 0} \ar[r]^d \ar[u] &
A_{1, 0} \ar[r]^d \ar[u] &
A_{0, 0} \ar[r] \ar[u] &
R(A)_0 \\
\ldots \ar[r]^d &
A_{2, 1} \ar[r]^d \ar[u]^\delta &
A_{1, 1} \ar[r]^d \ar[u]^\delta &
A_{0, 1} \ar[r] \ar[u]^\delta &
R(A)_1 \ar[u]^\delta \\
\ldots \ar[r]^d &
A_{2, 2} \ar[r]^d \ar[u]^\delta &
A_{1, 2} \ar[r]^d \ar[u]^\delta &
A_{0, 2} \ar[r] \ar[u]^\delta &
R(A)_2 \ar[u]^\delta \\
&
\ldots \ar[u]^\delta &
\ldots \ar[u]^\delta &
\ldots \ar[u]^\delta &
\ldots \ar[u]^\delta \\
}
$$
(Bien entendu, ce n'est plus un bicomplexe.)
La notion de morphisme $\Phi : (A_{\bullet, \bullet}, d, \delta)
\to (B_{\bullet, \bullet}, d, \delta)$ de bicomplexes
est claire, et il est clair qu'un tel morphisme induit des morphismes de complexes
$R(\Phi) : R(A)_\bullet \to R(B)_\bullet$ et
$U(\Phi) : U(A)_\bullet \to U(B)_\bullet$.

\begin{lemma}
\label{algebra-lemma-no-spectral-sequence}
Soit $(A_{\bullet, \bullet}, d, \delta)$ un bicomplexe tel
que
\begin{enumerate}
\item Chaque ligne $A_{\bullet, j}$ soit une résolution de $R(A)_j$.
\item Chaque colonne $A_{i, \bullet}$ soit une résolution de $U(A)_i$.
\end{enumerate}
Il existe alors des isomorphismes canoniques
$$
H_i(R(A)_\bullet)
\cong
H_i(U(A)_\bullet).
$$
Ces isomorphismes sont fonctoriels par rapport aux morphismes
de bicomplexes possédant les propriétés ci-dessus.
\end{lemma}

\begin{proof}
Nous allons montrer que $H_i(R(A)_\bullet))$
et $H_i(U(A)_\bullet)$ sont canoniquement
isomorphes à un troisième groupe, à savoir
$$
\mathbf{H}_i(A) :=
\frac{
\{
(a_{i, 0}, a_{i-1, 1}, \ldots, a_{0, i})
\mid
d(a_{i, 0}) = \delta(a_{i-1, 1}), \ldots,
d(a_{1, i-1}) = \delta(a_{0, i})
\}}
{
\{
d(a_{i + 1, 0}) + \delta(a_{i, 1}),
d(a_{i, 1}) + \delta(a_{i-1, 2}),
\ldots,
d(a_{1, i}) + \delta(a_{0, i + 1})
\}
}
$$
Nous adoptons ici la convention de notation selon laquelle $a_{i, j}$ désigne
un élément de $A_{i, j}$. Autrement dit, un élément de $\mathbf{H}_i$
est représenté par un zigzag, figuré comme suit lorsque $i = 2$
$$
\xymatrix{
a_{2, 0} \ar@{|->}[r] & d(a_{2, 0}) = \delta(a_{1, 1}) & \\
& a_{1, 1} \ar@{|->}[u] \ar@{|->}[r] & d(a_{1, 1}) = \delta(a_{0, 2}) \\
& & a_{0, 2} \ar@{|->}[u] \\
}
$$
Naturellement, on quotient par les zigzags « triviaux », c'est-à-dire par le sous-module
engendré par les éléments de la forme $(0, \ldots, 0, -\delta(a_{t + 1, t-i}),
d(a_{t + 1, t-i}), 0, \ldots, 0)$. Remarquons qu'il existe des
homomorphismes canoniques
$$
\mathbf{H}_i(A) \to H_i(R(A)_\bullet), \quad
(a_{i, 0}, a_{i-1, 1}, \ldots, a_{0, i}) \mapsto
\text{classe de l'image de }a_{0, i}
$$
et
$$
\mathbf{H}_i(A) \to H_i(U(A)_\bullet), \quad
(a_{i, 0}, a_{i-1, 1}, \ldots, a_{0, i}) \mapsto
\text{classe de l'image de }a_{i, 0}
$$

\medskip\noindent
Montrons d'abord que ces homomorphismes sont surjectifs.
Supposons que $\overline{r} \in H_i(R(A)_\bullet)$.
Soit $r \in R(A)_i$ un cycle qui représente la
classe de $\overline{r}$.
Soit $a_{0, i} \in A_{0, i}$ un élément dont
l'image est $r$. Puisque $\delta(r) = 0$,
on voit que $\delta(a_{0, i})$ appartient à
l'image de $d$. Il existe donc un élément
$a_{1, i-1} \in A_{1, i-1}$ tel que
$d(a_{1, i-1}) = \delta(a_{0, i})$. Cela implique à son tour
que $\delta(a_{1, i-1})$ appartient au noyau
de $d$ (car $d(\delta(a_{1, i-1})) = \delta(d(a_{1, i-1}))
= \delta(\delta(a_{0, i})) = 0$). Par exactitude des
lignes, on trouve un élément $a_{2, i-2}$ tel que
$d(a_{2, i-2}) = \delta(a_{1, i-1})$. Et l'on poursuit ainsi
jusqu'à obtenir un zigzag complet. Bien entendu, la surjectivité
de $\mathbf{H}_i \to H_i(U(A))$ se démontre de la même manière.

\medskip\noindent
Pour démontrer l'injectivité, on raisonne exactement de la même manière.
Supposons donc que l'on ait un zigzag
$(a_{i, 0}, a_{i-1, 1}, \ldots, a_{0, i})$
dont l'image dans $H_i(R(A)_\bullet)$ est nulle.
Cela signifie que l'image de $a_{0, i}$ est un élément
de $\Coker(A_{i, 1} \to A_{i, 0})$
qui appartient à l'image de
$\delta : \Coker(A_{i + 1, 1} \to A_{i + 1, 0}) \to
\Coker(A_{i, 1} \to A_{i, 0})$.
Autrement dit, $a_{0, i}$ appartient à l'image de
$\delta \oplus d : A_{0, i + 1} \oplus A_{1, i} \to A_{0, i}$.
La définition des zigzags triviaux montre que
l'on peut modifier notre zigzag par un zigzag trivial et
supposer que $a_{0, i} = 0$. Cela implique immédiatement
que $d(a_{1, i-1}) = 0$. Comme les lignes
sont exactes, il s'ensuit que $a_{1, i-1}$ appartient
à l'image de $d : A_{2, i-1} \to A_{1, i-1}$.
On peut donc modifier de nouveau notre zigzag par un
zigzag trivial et supposer qu'il est de la forme
$(a_{i, 0}, a_{i-1, 1}, \ldots, a_{2, i-2}, 0, 0)$.
En poursuivant ainsi, on obtient l'injectivité voulue.

\medskip\noindent
Si $\Phi : (A_{\bullet, \bullet}, d, \delta)
\to (B_{\bullet, \bullet}, d, \delta)$ est un morphisme
de bicomplexes satisfaisant tous deux aux conditions
du lemme, on obtient clairement un diagramme
commutatif
$$
\xymatrix{
H_i(U(A)_\bullet) \ar[d] &
\mathbf{H}_i(A) \ar[r] \ar[l] \ar[d] &
H_i(R(A)_\bullet) \ar[d] \\
H_i(U(B)_\bullet) &
\mathbf{H}_i(B) \ar[r] \ar[l] &
H_i(R(B)_\bullet) \\
}
$$
Cela démontre la fonctorialité.
\end{proof}

\begin{remark}
\label{algebra-remark-signs-double-complex}
L'isomorphisme construit ci-dessus n'est l'isomorphisme « correct » qu'aux signes près.
Une bonne partie de l'algèbre homologique consiste à choisir les signes de
diverses applications et à démontrer la commutativité de diagrammes en faisant intervenir
des signes appropriés. Pour le moment, nous utiliserons simplement l'isomorphisme
donné dans la démonstration ci-dessus et nous nous préoccuperons des signes plus tard.
\end{remark}

\begin{lemma}
\label{algebra-lemma-tor-left-right}
Soit $R$ un anneau. Pour tout $i \geq 0$, les foncteurs
$\text{Mod}_R \times \text{Mod}_R \to \text{Mod}_R$,
$(M, N) \mapsto \text{Tor}_i^R(M, N)$ et
$(M, N) \mapsto \text{Tor}_i^R(N, M)$ sont
canoniquement isomorphes.
\end{lemma}

\begin{proof}
Soit $F_\bullet$ une résolution libre du module $M$ et
soit $G_\bullet$ une résolution libre du module $N$.
Considérons le bicomplexe $(A_{i, j}, d, \delta)$ défini
comme suit :
\begin{enumerate}
\item posons $A_{i, j} = F_i \otimes_R G_j$,
\item définissons $d_{i, j} : F_i \otimes_R G_j \to F_{i-1} \otimes G_j$
comme $d_{F, i} \otimes \text{id}$, et
\item définissons $\delta_{i, j} : F_i \otimes_R G_j \to F_i \otimes G_{j-1}$
comme $\text{id} \otimes d_{G, j}$.
\end{enumerate}
Ce bicomplexe est généralement noté simplement $F_\bullet \otimes_R G_\bullet$.

\medskip\noindent
Puisque chaque $G_j$ est libre, donc plat, chaque
ligne du bicomplexe est exacte sauf en degré homologique
$0$. Puisque chaque $F_i$ est libre, donc plat, chaque
colonne du bicomplexe est exacte sauf en degré homologique
$0$. Le bicomplexe satisfait donc aux conditions
du Lemme \ref{algebra-lemma-no-spectral-sequence}.

\medskip\noindent
Pour comprendre ce que dit le lemme, calculons $R(A)_\bullet$ et $U(A)_\bullet$.
On a
\begin{eqnarray*}
R(A)_i & = & \Coker(A_{1, i} \to A_{0, i}) \\
& = & \Coker(F_1 \otimes_R G_i \to F_0 \otimes_R G_i) \\
& = & \Coker(F_1 \to F_0) \otimes_R G_i \\
& = & M \otimes_R G_i
\end{eqnarray*}
En fait, ces isomorphismes sont compatibles avec les différentielles
$\delta$, et l'on voit que $R(A)_\bullet = M \otimes_R G_\bullet$
comme complexes homologiques. De manière tout à fait analogue, on voit que
$U(A)_\bullet = F_\bullet \otimes_R N$. On obtient
\begin{eqnarray*}
\text{Tor}_i^R(M, N)
& = & H_i(F_\bullet \otimes_R N) \\
& = & H_i(U(A)_\bullet) \\
& = & H_i(R(A)_\bullet) \\
& = & H_i(M \otimes_R G_\bullet) \\
& = & H_i(G_\bullet \otimes_R M) \\
& = & \text{Tor}_i^R(N, M)
\end{eqnarray*}
Ici, la troisième égalité est le Lemme \ref{algebra-lemma-no-spectral-sequence}, et
la cinquième utilise l'isomorphisme $V \otimes W = W \otimes V$
du produit tensoriel.

\medskip\noindent
Fonctorialité. Supposons que l'on ait des $R$-modules $M_\nu$, $N_\nu$,
$\nu = 1, 2$. Soient $\varphi : M_1 \to M_2$ et $\psi : N_1 \to N_2$
des morphismes de $R$-modules.
Supposons que l'on ait des résolutions libres $F_{\nu, \bullet}$
de $M_\nu$ et des résolutions libres $G_{\nu, \bullet}$ de $N_\nu$.
D'après le Lemme \ref{algebra-lemma-compare-resolutions}, on peut choisir
des morphismes de complexes $\alpha : F_{1, \bullet} \to F_{2, \bullet}$
et $\beta : G_{1, \bullet} \to G_{2, \bullet}$ compatibles
avec $\varphi$ et $\psi$. Nous affirmons que
la paire $(\alpha, \beta)$ induit un morphisme de
bicomplexes
$$
\alpha \otimes \beta :
F_{1, \bullet} \otimes_R G_{1, \bullet}
\longrightarrow
F_{2, \bullet} \otimes_R G_{2, \bullet}
$$
Il s'agit d'une vérification tout à fait directe : l'application
$F_{1, i} \otimes_R G_{1, j} \to F_{2, i} \otimes_R G_{2, j}$
est donnée par $\alpha_i \otimes \beta_j$, où $\alpha_i$,
resp.\ $\beta_j$, est la composante de degré $i$, resp.\ $j$, de $\alpha$,
resp.\ $\beta$. Le lecteur vérifie aussi aisément que l'application induite
$R(F_{1, \bullet} \otimes_R G_{1, \bullet})_\bullet
\to R(F_{2, \bullet} \otimes_R G_{2, \bullet})_\bullet$
coïncide avec l'application $M_1 \otimes_R G_{1, \bullet}
\to M_2 \otimes_R G_{2, \bullet}$ induite par $\varphi \otimes \beta$.
Il en va de même de l'application induite sur les complexes $U(-)_\bullet$.
L'assertion de fonctorialité résulte donc de l'assertion
de fonctorialité du Lemme \ref{algebra-lemma-no-spectral-sequence}.
\end{proof}

\begin{remark}
\label{algebra-remark-curiosity-signs-swap}
Un cas intéressant se présente lorsque $M = N$ ci-dessus.
On obtient alors une application canonique $\text{Tor}_i^R(M, M)
\to \text{Tor}_i^R(M, M)$. Remarquons que cette application n'est pas
l'identité, car même lorsque $i = 0$, elle n'est pas
l'identité ! Par exemple, si $V$ est un espace vectoriel de dimension
$n$ sur un corps, alors l'application d'échange $V \otimes_k V \to V \otimes_k V$
possède $(n^2 + n)/2$ valeurs propres égales à $+1$ et $(n^2-n)/2$ valeurs propres
égales à $-1$. En caractéristique $2$, elle n'est même pas diagonalisable.
Remarquons que même changer le signe de l'application ne suffira pas à supprimer
ce phénomène.
\end{remark}

\begin{lemma}
\label{algebra-lemma-tor-noetherian}
Soit $R$ un anneau noethérien. Soient $M$, $N$ des $R$-modules de type fini.
Alors $\text{Tor}_p^R(M, N)$ est un $R$-module de type fini pour tout $p$.
\end{lemma}

\begin{proof}
Cela résulte de ce que $\text{Tor}_p^R(M, N)$ se calcule comme les
groupes d'homologie d'un complexe $F_\bullet \otimes_R N$
où chaque $F_n$ est un $R$-module libre de type fini ; voir le
Lemme \ref{algebra-lemma-resolution-by-finite-free}.
\end{proof}

\begin{lemma}
\label{algebra-lemma-characterize-flat}
Soit $R$ un anneau. Soit $M$ un $R$-module.
Les assertions suivantes sont équivalentes :
\begin{enumerate}
\item Le module $M$ est plat sur $R$.
\item Pour tout $i > 0$, le foncteur $\text{Tor}_i^R(M, -)$ est nul.
\item Le foncteur $\text{Tor}_1^R(M, -)$ est nul.
\item Pour tout idéal $I \subset R$, on a $\text{Tor}_1^R(M, R/I) = 0$.
\item Pour tout idéal de type fini $I \subset R$, on a
$\text{Tor}_1^R(M, R/I) = 0$.
\end{enumerate}
\end{lemma}

\begin{proof}
Supposons $M$ plat. Soit $N$ un $R$-module.
Soit $F_\bullet$ une résolution libre de $N$.
Alors $F_\bullet \otimes_R M$ est une résolution de $N \otimes_R M$,
par platitude de $M$. Tous les groupes Tor supérieurs s'annulent donc.

\medskip\noindent
Il suffit maintenant de montrer que la dernière condition implique que
$M$ est plat. Soit $I \subset R$ un idéal.
Considérons la suite exacte courte
$0 \to I \to R \to R/I \to 0$. Appliquons le
Lemme \ref{algebra-lemma-long-exact-sequence-tor}. On obtient une
suite exacte
$$
\text{Tor}_1^R(M, R/I) \to
M \otimes_R I \to
M \otimes_R R \to
M \otimes_R R/I \to
0
$$
Comme, de toute évidence, $M \otimes_R R = M$, on conclut que la
dernière hypothèse implique que $M \otimes_R I \to M$ est
injectif pour tout idéal de type fini $I$.
Ainsi $M$ est plat d'après le Lemme \ref{algebra-lemma-flat}.
\end{proof}

\begin{remark}
\label{algebra-remark-Tor-ring-mod-ideal}
La démonstration du Lemme \ref{algebra-lemma-characterize-flat} montre en fait
que
$$
\text{Tor}_1^R(M, R/I)
=
\Ker(I \otimes_R M \to M).
$$
\end{remark}














\section{Fonctorialités de Tor}
\label{algebra-section-functoriality-tor}

\noindent
Dans cette section, nous examinons brièvement la fonctorialité
de $\text{Tor}$ par rapport au changement d'anneau, etc.
Voici une liste de points à développer.
\begin{enumerate}
\item Étant donnés un homomorphisme d'anneaux $R \to R'$, un $R$-module
$M$ et un $R'$-module $N'$,
les $R$-modules $\text{Tor}_i^R(M, N')$ possèdent
une structure naturelle de $R'$-module.
\item Étant donnés un homomorphisme d'anneaux $R \to R'$ et des $R$-modules
$M$, $N$, il existe un homomorphisme naturel de $R$-modules
$\text{Tor}_i^R(M, N) \to \text{Tor}_i^{R'}(M \otimes_R R', N \otimes_R R')$.
\item Étant donnés un homomorphisme d'anneaux $R \to R'$, un $R$-module $M$ et
un $R'$-module $N'$, il existe un homomorphisme naturel
de $R'$-modules
$\text{Tor}_i^R(M, N') \to \text{Tor}_i^{R'}(M \otimes_R R', N')$.
\end{enumerate}

\begin{lemma}
\label{algebra-lemma-flat-base-change-tor}
Étant donnés un homomorphisme plat d'anneaux $R \to R'$ et des $R$-modules
$M$, $N$, l'homomorphisme naturel de $R$-modules
$\text{Tor}_i^R(M, N)\otimes_R R'
\to \text{Tor}_i^{R'}(M \otimes_R R', N \otimes_R R')$
est un isomorphisme pour tout $i$.
\end{lemma}

\begin{proof}
Omis. Cela résulte de ce qu'une résolution libre $F_\bullet$ de $M$ sur
$R$ reste exacte après tensorisation par $R'$ sur $R$, et donc que
$(F_\bullet \otimes_R N)\otimes_R R'$ calcule les groupes Tor
sur $R'$.
\end{proof}

\noindent
Le lemme suivant ne semble trouver sa place nulle part ailleurs.

\begin{lemma}
\label{algebra-lemma-tor-commutes-filtered-colimits}
Soit $R$ un anneau. Soit $M = \colim M_i$ une colimite filtrante de
$R$-modules. Soit $N$ un $R$-module. Alors
$\text{Tor}_n^R(M, N) = \colim \text{Tor}_n^R(M_i, N)$ pour tout $n$.
\end{lemma}

\begin{proof}
Choisissons une résolution libre $F_\bullet$ de $N$. Alors
$F_\bullet \otimes_R M = \colim F_\bullet \otimes_R M_i$
comme complexes, d'après le Lemme \ref{algebra-lemma-tensor-products-commute-with-limits}.
Le résultat découle donc du Lemme \ref{algebra-lemma-directed-colimit-exact}.
\end{proof}








\section{Modules projectifs}
\label{algebra-section-projective}

\noindent
Quelques lemmes sur les modules projectifs.

\begin{definition}
\label{algebra-definition-projective}
Soit $R$ un anneau. Un $R$-module $P$ est {\it projectif} si et seulement si
le foncteur $\Hom_R(P, -) : \text{Mod}_R \to \text{Mod}_R$ est
un foncteur exact.
\end{definition}

\noindent
Le foncteur $\Hom_R(M, - )$ est exact à gauche pour tout $R$-module $M$ ; voir le
Lemme \ref{algebra-lemma-hom-exact}.
La condition que $P$ soit projectif signifie donc précisément que, pour toute
surjection de $R$-modules $N \to N'$, l'application
$\Hom_R(P, N) \to \Hom_R(P, N')$ est surjective.

\begin{lemma}
\label{algebra-lemma-characterize-projective}
Soit $R$ un anneau. Soit $P$ un $R$-module.
Les assertions suivantes sont équivalentes :
\begin{enumerate}
\item $P$ est projectif,
\item $P$ est un facteur direct d'un $R$-module libre, et
\item $\Ext^1_R(P, M) = 0$ pour tout $R$-module $M$.
\end{enumerate}
\end{lemma}

\begin{proof}
Supposons $P$ projectif. Choisissons une surjection $\pi : F \to P$, où $F$
est un $R$-module libre. Comme $P$ est projectif, il existe
$i \in \Hom_R(P, F)$ tel que $\pi \circ i = \text{id}_P$.
Autrement dit, $F \cong \Ker(\pi) \oplus i(P)$, et l'on voit
que $P$ est un facteur direct de $F$.

\medskip\noindent
Réciproquement, supposons que $P \oplus Q = F$ soit un $R$-module libre.
Remarquons que le module libre $F = \bigoplus_{i \in I} R$ est projectif,
car $\Hom_R(F, M) = \prod_{i \in I} M$ et le foncteur
$M \mapsto \prod_{i \in I} M$ est exact.
Alors $\Hom_R(F, -) = \Hom_R(P, -) \times \Hom_R(Q, -)$
comme foncteurs, et donc $P$ et $Q$ sont tous deux projectifs.

\medskip\noindent
Supposons que $P \oplus Q = F$ soit un $R$-module libre. On dispose alors d'une
résolution libre $F_\bullet$ de la forme
$$
\ldots F \xrightarrow{a} F \xrightarrow{b} F \to P \to 0
$$
où les applications $a, b$ alternent et sont égales aux projecteurs sur
$P$ et $Q$. Le complexe $\Hom_R(F_\bullet, M)$ est donc exact scindé
en degrés $\geq 1$, d'où l'annulation de (3).

\medskip\noindent
Supposons que $\Ext^1_R(P, M) = 0$ pour tout $R$-module $M$.
Choisissons une résolution libre $F_\bullet \to P$. Posons
$M = \Im(F_1 \to F_0) = \Ker(F_0 \to P)$.
Considérons l'élément $\xi \in \Ext^1_R(P, M)$ donné par
la classe de l'application quotient $\pi : F_1 \to M$. Comme $\xi$ est nul,
il existe une application $s : F_0 \to M$ telle que $\pi = s \circ (F_1 \to F_0)$.
Cela signifie clairement que
$$
F_0 = \Ker(s) \oplus \Ker(F_0 \to P) =
P \oplus \Ker(F_0 \to P)
$$
et l'on conclut.
\end{proof}

\begin{lemma}
\label{algebra-lemma-characterize-finite-projective-noetherian}
Soit $R$ un anneau noethérien. Soit $P$ un $R$-module de type fini.
Si $\Ext^1_R(P, M) = 0$ pour tout $R$-module de type fini $M$, alors
$P$ est projectif.
\end{lemma}

\noindent
On peut renforcer ce lemme : il existe une version pour les
$R$-modules de présentation finie sans supposer $R$ noethérien. Dans le cas noethérien,
il existe une version où $M$ parcourt tous les modules de longueur finie.

\begin{proof}
Choisissons une surjection $R^{\oplus n} \to P$ de noyau $M$.
Puisque $\Ext^1_R(P, M) = 0$, cette surjection est scindée, et
on conclut par le Lemme \ref{algebra-lemma-characterize-projective}.
\end{proof}

\begin{lemma}
\label{algebra-lemma-direct-sum-projective}
Une somme directe de modules projectifs est projective.
\end{lemma}

\begin{proof}
Cela résulte de la caractérisation des modules projectifs comme facteurs
directs de modules libres dans le
Lemme \ref{algebra-lemma-characterize-projective}.
\end{proof}

\begin{lemma}
\label{algebra-lemma-lift-projective-module}
Soit $R$ un anneau. Soit $I \subset R$ un idéal nilpotent. Soit
$\overline{P}$ un $R/I$-module projectif. Il existe alors un
$R$-module projectif $P$ tel que $P/IP \cong \overline{P}$.
\end{lemma}

\begin{proof}
D'après le Lemme \ref{algebra-lemma-characterize-projective},
on peut choisir un ensemble $A$ et une décomposition en somme directe
$\bigoplus_{\alpha \in A} R/I = \overline{P} \oplus \overline{K}$
pour un certain $R/I$-module $\overline{K}$. Notons $F = \bigoplus_{\alpha \in A} R$
le $R$-module libre de base $A$. Choisissons un relèvement
$p : F \to F$ du projecteur $\overline{p}$ associé au
facteur direct $\overline{P}$ de $\bigoplus_{\alpha \in A} R/I$.
Remarquons que $p^2 - p \in \text{End}_R(F)$ est un endomorphisme
nilpotent de $F$ (puisque $I$ est nilpotent et que les coefficients matriciels de
$p^2 - p$ appartiennent à $I$ ; plus précisément, si $I^n = 0$, alors $(p^2 - p)^n = 0$).
Ainsi, d'après le Lemme \ref{algebra-lemma-lift-idempotents-noncommutative},
on peut modifier le choix de $p$ et supposer que $p$ est un projecteur.
Posons $P = \Im(p)$.
\end{proof}

\begin{lemma}
\label{algebra-lemma-lift-finite-projective-module}
Soit $R$ un anneau. Soit $I \subset R$ un idéal localement nilpotent. Soit
$\overline{P}$ un $R/I$-module projectif de type fini. Il existe alors un
$R$-module projectif de type fini $P$ tel que $P/IP \cong \overline{P}$.
\end{lemma}

\begin{proof}
Rappelons que $\overline{P}$ est un facteur direct d'un $R/I$-module libre
$\bigoplus_{\alpha \in A} R/I$ d'après le Lemme \ref{algebra-lemma-characterize-projective}.
Comme $\overline{P}$ est de type fini, il s'ensuit que $\overline{P}$ est contenu
dans $\bigoplus_{\alpha \in A'} R/I$ pour une partie finie $A' \subset A$.
On peut donc supposer que l'on a une décomposition en somme directe
$(R/I)^{\oplus n} = \overline{P} \oplus \overline{K}$
pour un certain $n$ et un certain $R/I$-module $\overline{K}$. Choisissons un relèvement
$p \in \text{Mat}(n \times n, R)$ du projecteur $\overline{p}$
associé au facteur direct $\overline{P}$ de $(R/I)^{\oplus n}$.
Remarquons que $p^2 - p \in \text{Mat}(n \times n, R)$ est nilpotent :
comme $I$ est localement nilpotent et que les coefficients matriciels $c_{ij}$ de
$p^2 - p$ appartiennent à $I$, on a $c_{ij}^t = 0$ pour un certain $t > 0$, et
alors $(p^2 - p)^{tn^2} = 0$ (comme on le voit sur les coefficients matriciels).
Ainsi, d'après le Lemme \ref{algebra-lemma-lift-idempotents-noncommutative},
on peut modifier le choix de $p$ et supposer que $p$ est un projecteur.
Posons $P = \Im(p)$.
\end{proof}

\begin{lemma}
\label{algebra-lemma-lift-projective}
Soit $R$ un anneau.
Soit $I \subset R$ un idéal.
Soit $M$ un $R$-module.
Supposons que
\begin{enumerate}
\item $I$ soit nilpotent,
\item $M/IM$ soit un $R/I$-module projectif,
\item $M$ soit un $R$-module plat.
\end{enumerate}
Alors $M$ est un $R$-module projectif.
\end{lemma}

\begin{proof}
D'après le Lemme \ref{algebra-lemma-lift-projective-module}, on peut trouver un
$R$-module projectif $P$ et un isomorphisme $P/IP \to M/IM$. Nous allons montrer
que $M$ est isomorphe à $P$, ce qui achèvera la démonstration. Comme $P$
est projectif, on peut relever l'application $P \to P/IP \to M/IM$ en un homomorphisme de $R$-modules
$P \to M$ qui est un isomorphisme modulo $I$. Puisque $I^n = 0$
pour un certain $n$, on peut utiliser les filtrations
\begin{align*}
0 = I^nM \subset I^{n - 1}M \subset \ldots \subset IM \subset M \\
0 = I^nP \subset I^{n - 1}P \subset \ldots \subset IP \subset P
\end{align*}
pour voir qu'il suffit de montrer que les applications induites
$I^aP/I^{a + 1}P \to I^aM/I^{a + 1}M$ sont bijectives. Comme $P$
et $M$ sont tous deux des $R$-modules plats, on peut identifier cette application à
$$
I^a/I^{a + 1} \otimes_{R/I} P/IP
\longrightarrow
I^a/I^{a + 1} \otimes_{R/I} M/IM
$$
induite par $P \to M$. Puisque l'on a choisi $P \to M$ de sorte que l'application induite
$P/IP \to M/IM$ soit un isomorphisme, on conclut.
\end{proof}

\begin{lemma}
\label{algebra-lemma-projective-modulo-two-ideals}
Soit $R$ un anneau. Soient $I, J \subset R$ des idéaux tels que $I \cap J = 0$.
Soit $P$ un $R$-module tel que $P/IP$ soit un $R/I$-module projectif et que
$P/JP$ soit un $R/J$-module projectif. Alors $P$ est un $R$-module projectif.
\end{lemma}

\begin{proof}
Choisissons une surjection $p : F \to P$, où $F$ est un $R$-module libre.
Puisque $P/IP$ est un $R/I$-module projectif, on peut choisir une application
$f : P/IP \to F/IF$ qui soit un inverse à droite de $p \bmod I$.
Considérons l'application $q : F/JF \to F/(I + J)F \times_{P/(I + J)P} P/JP$
induite par l'application quotient $F/JF \to F/(I + J)F$ et par $p$.
Remarquons que $q$ est un homomorphisme surjectif de $R/J$-modules (petit détail omis).
Considérons l'application $f' : P/JP \to F/(I + J)F \times_{P/(I + J)P} P/JP$
induite par $f$ et par l'application identité.
Puisque $P/JP$ est un $R/J$-module projectif et que $q$ est surjectif,
on peut choisir une application $g : P/JP \to F/JF$ telle que $q \circ g = f'$.
Alors $f \bmod I + J = g \bmod I + J$. Puisque $F = F/JF \times_{F/(I + J)F} F/IF$
car $I \cap J = 0$, on conclut que l'on obtient un homomorphisme bien défini
$h = (f, g) : P \to F$ de $R$-modules. L'application $a = p \circ h : P \to P$
se réduit à l'identité modulo $I$ et modulo $J$. Alors $a$ induit aussi
l'identité sur le sous-module $IP$ : si $x = \sum i_\alpha x_\alpha$
avec $i_\alpha \in I$ et $x_\alpha \in P$, alors
$a(x) = \sum i_\alpha a(x_\alpha) = \sum i_\alpha x_\alpha = x$,
car $a(x_\alpha) = x_\alpha + j_\alpha$ avec $j_\alpha \in J$
et $i_\alpha j_\alpha = 0$. Ainsi $a : P \to P$ agit comme l'identité
sur $IP$ et sur $P/IP$, et l'on conclut que $a$ est un automorphisme de $P$
(petit détail omis). Ainsi $P$ est un facteur direct de $F$, donc projectif.
\end{proof}










\section{Modules projectifs de type fini}
\label{algebra-section-finite-projective-modules}

\begin{definition}
\label{algebra-definition-locally-free}
Soit $R$ un anneau et $M$ un $R$-module.
\begin{enumerate}
\item On dit que $M$ est {\it localement libre} si l'on peut recouvrir $\Spec(R)$ par
des ouverts standards $D(f_i)$, $i \in I$, tels que $M_{f_i}$ soit un
$R_{f_i}$-module libre pour tout $i \in I$.
\item On dit que $M$ est {\it localement libre de type fini} si l'on peut choisir
le recouvrement de sorte que chaque $M_{f_i}$ soit libre de type fini.
\item On dit que $M$ est {\it localement libre de type fini de rang $r$}
si l'on peut choisir le recouvrement de sorte que chaque $M_{f_i}$ soit isomorphe
à $R_{f_i}^{\oplus r}$.
\end{enumerate}
\end{definition}

\noindent
Remarquons qu'un $R$-module localement libre de type fini est automatiquement
de présentation finie d'après le Lemme \ref{algebra-lemma-cover}. De plus, si $M$ est un
module localement libre de type fini de rang $r$ sur un anneau $R$ et si
$R$ est non nul, alors $r$ est déterminé de manière unique par le
Lemme \ref{algebra-lemma-rank} (car au moins l'un des anneaux
localisés $	R_{f_i}$ est non nul).

\begin{lemma}
\label{algebra-lemma-finite-projective}
Soit $R$ un anneau et soit $M$ un $R$-module.
Les assertions suivantes sont équivalentes :
\begin{enumerate}
\item $M$ est de présentation finie et plat sur $R$,
\item $M$ est projectif de type fini,
\item $M$ est un facteur direct d'un $R$-module libre de type fini,
\item $M$ est de présentation finie et,
pour tout $\mathfrak p \in \Spec(R)$, le
module localisé $M_{\mathfrak p}$ est libre,
\item $M$ est de présentation finie et,
pour tout idéal maximal $\mathfrak m \subset R$, le
module localisé $M_{\mathfrak m}$ est libre,
\item $M$ est de type fini et localement libre,
\item $M$ est localement libre de type fini, et
\item $M$ est de type fini, pour tout idéal premier $\mathfrak p$ le module
$M_{\mathfrak p}$ est libre, et la fonction
$$
\rho_M : \Spec(R) \to \mathbf{Z}, \quad
\mathfrak p
\longmapsto
\dim_{\kappa(\mathfrak p)} M \otimes_R \kappa(\mathfrak p)
$$
est localement constante pour la topologie de Zariski.
\end{enumerate}
\end{lemma}

\begin{proof}
Supposons d'abord $M$ projectif de type fini, c'est-à-dire que (2) est vérifiée.
Prenons une surjection $R^n \to M$ et notons $K$ son noyau.
Puisque $M$ est projectif,
$0 \to K \to R^n \to M \to 0$ est scindée.
Ainsi (2) $\Rightarrow$ (3).
L'implication (3) $\Rightarrow$ (2) résulte de ce qu'un
facteur direct d'un module projectif est projectif ; voir le
Lemme \ref{algebra-lemma-characterize-projective}.

\medskip\noindent
Supposons (3), de sorte que l'on puisse écrire $K \oplus M \cong R^{\oplus n}$.
Alors $K$ est un facteur direct de $R^n$, et donc de type fini.
Cela montre que $M = R^{\oplus n}/K$ est de présentation finie.
Autrement dit, (3) $\Rightarrow$ (1).

\medskip\noindent
Supposons que $M$ soit de présentation finie et plat, c'est-à-dire que (1) soit vérifiée.
Nous allons démontrer (7). Choisissons un idéal premier $\mathfrak p$ quelconque et
$x_1, \ldots, x_r \in M$ dont les images forment une base de
$M \otimes_R \kappa(\mathfrak p)$. D'après le
lemme de Nakayama (sous la forme du Lemme \ref{algebra-lemma-NAK-localization}),
ces éléments engendrent $M_g$ pour un certain $g \in R$, $g \not \in \mathfrak p$.
La surjection correspondante $\varphi : R_g^{\oplus r} \to M_g$
possède les deux propriétés suivantes : (a) $\Ker(\varphi)$ est un
$R_g$-module de type fini (voir le Lemme \ref{algebra-lemma-extension})
et (b) $\Ker(\varphi) \otimes \kappa(\mathfrak p) = 0$
par platitude de $M_g$ sur $R_g$ (voir le
Lemme \ref{algebra-lemma-flat-tor-zero}).
Ainsi, d'après le lemme de Nakayama encore une fois, il existe $g' \in R_g$ tel que
$\Ker(\varphi)_{g'} = 0$. Autrement dit, $M_{gg'}$ est libre.

\medskip\noindent
Un module localement libre de type fini est de type fini ; voir le
Lemme \ref{algebra-lemma-cover},
donc (7) $\Rightarrow$ (6).
Il est clair que (6) $\Rightarrow$ (7) et que (7) $\Rightarrow$ (8).

\medskip\noindent
Un module localement libre de type fini est de présentation finie ; voir le
Lemme \ref{algebra-lemma-cover},
donc (7) $\Rightarrow$ (4).
Bien entendu, (4) implique (5).
Puisque l'on peut vérifier la platitude localement (voir le
Lemme \ref{algebra-lemma-flat-localization}),
on conclut que (5) implique (1).
À ce stade, on a
$$
\xymatrix{
(2) \ar@{<=>}[r] & (3) \ar@{=>}[r] & (1) \ar@{=>}[r] &
(7)  \ar@{<=>}[r] \ar@{=>}[rd] \ar@{=>}[d] & (6) \\
& & (5) \ar@{=>}[u] & (4) \ar@{=>}[l] & (8)
}
$$

\medskip\noindent
Supposons que $M$ satisfasse (1), (4), (5), (6) et (7).
Nous allons démontrer (3). Il suffit
de montrer que $M$ est projectif. Il faut montrer que $\Hom_R(M, -)$
est exact. Soit $0 \to N'' \to N \to N'\to 0$ une suite exacte courte de
$R$-modules. Il faut montrer que
$0 \to \Hom_R(M, N'') \to \Hom_R(M, N) \to
\Hom_R(M, N') \to 0$ est exacte.
Comme $M$ est localement libre de type fini, il existe un recouvrement
$\Spec(R) = \bigcup D(f_i)$ tel que $M_{f_i}$ soit libre de type fini.
D'après le
Lemme \ref{algebra-lemma-hom-from-finitely-presented},
on voit que
$$
0 \to \Hom_R(M, N'')_{f_i} \to \Hom_R(M, N)_{f_i} \to
\Hom_R(M, N')_{f_i} \to 0
$$
est égale à
$0 \to \Hom_{R_{f_i}}(M_{f_i}, N''_{f_i}) \to
\Hom_{R_{f_i}}(M_{f_i}, N_{f_i}) \to
\Hom_{R_{f_i}}(M_{f_i}, N'_{f_i}) \to 0$,
qui est exacte puisque $M_{f_i}$ est libre et que la localisation
$0 \to N''_{f_i} \to N_{f_i} \to N'_{f_i} \to 0$
est exacte (car la localisation est exacte). On voit donc que
$0 \to \Hom_R(M, N'') \to \Hom_R(M, N) \to
\Hom_R(M, N') \to 0$ est exacte d'après le
Lemme \ref{algebra-lemma-cover}.

\medskip\noindent
Enfin, supposons que (8) soit vérifiée. Choisissons un idéal maximal $\mathfrak m \subset R$.
Choisissons $x_1, \ldots, x_r \in M$ dont les images forment une base sur $\kappa(\mathfrak m)$ de
$M \otimes_R \kappa(\mathfrak m) = M/\mathfrak mM$. En particulier,
$\rho_M(\mathfrak m) = r$. D'après le
Lemme de Nakayama \ref{algebra-lemma-NAK},
il existe $f \in R$, $f \not \in \mathfrak m$, tel que
$x_1, \ldots, x_r$ engendrent $M_f$ sur $R_f$. Comme, par hypothèse,
$\rho_M$ est localement constante, il existe $g \in R$, $g \not \in \mathfrak m$,
tel que $\rho_M$ soit constante de valeur $r$ sur $D(g)$. Nous affirmons que
$$
\Psi : R_{fg}^{\oplus r} \longrightarrow M_{fg}, \quad
(a_1, \ldots, a_r) \longmapsto \sum a_i x_i
$$
est un isomorphisme. Cela montrera que $M$ est localement libre de type
fini, c'est-à-dire que (7) est vérifiée. Pour établir l'affirmation,
il suffit de montrer que l'application induite sur les localisations
$\Psi_{\mathfrak p} : R_{\mathfrak p}^{\oplus r} \to M_{\mathfrak p}$
est un isomorphisme pour tout $\mathfrak p \in D(fg)$ ; voir le
Lemme \ref{algebra-lemma-characterize-zero-local}.
Par le choix de $f$, l'application $\Psi_{\mathfrak p}$
est surjective. D'après l'hypothèse (8), on a
$M_{\mathfrak p} \cong R_{\mathfrak p}^{\oplus \rho_M(\mathfrak p)}$,
et, par le choix de $g$, on a $\rho_M(\mathfrak p) = r$.
Ainsi $\Psi_{\mathfrak p}$ détermine une surjection
$R_{\mathfrak p}^{\oplus r} \to
M_{\mathfrak p} \cong R_{\mathfrak p}^{\oplus r}$,
qui est donc un isomorphisme d'après le
Lemme \ref{algebra-lemma-fun}.
(Bien entendu, ce dernier fait résulte aussi d'un simple argument matriciel.)
\end{proof}

\begin{lemma}
\label{algebra-lemma-finite-projective-reduced}
Soit $R$ un anneau réduit et soit $M$ un $R$-module. Alors les
conditions équivalentes du Lemme \ref{algebra-lemma-finite-projective}
sont aussi équivalentes à
\begin{enumerate}
\item[(9)] $M$ est de type fini et la fonction
$\rho_M : \Spec(R) \to \mathbf{Z}$, $\mathfrak p \mapsto
\dim_{\kappa(\mathfrak p)} M \otimes_R \kappa(\mathfrak p)$
est localement constante pour la topologie de Zariski.
\end{enumerate}
\end{lemma}

\begin{proof}
Choisissons un idéal maximal $\mathfrak m \subset R$.
Choisissons $x_1, \ldots, x_r \in M$ dont les images forment une base sur $\kappa(\mathfrak m)$ de
$M \otimes_R \kappa(\mathfrak m) = M/\mathfrak mM$. En particulier,
$\rho_M(\mathfrak m) = r$. D'après le
Lemme de Nakayama \ref{algebra-lemma-NAK},
il existe $f \in R$, $f \not \in \mathfrak m$, tel que
$x_1, \ldots, x_r$ engendrent $M_f$ sur $R_f$. Comme, par hypothèse,
$\rho_M$ est localement constante, il existe $g \in R$, $g \not \in \mathfrak m$,
tel que $\rho_M$ soit constante de valeur $r$ sur $D(g)$. Nous affirmons que
$$
\Psi : R_{fg}^{\oplus r} \longrightarrow M_{fg}, \quad
(a_1, \ldots, a_r) \longmapsto \sum a_i x_i
$$
est un isomorphisme. Cela montrera que $M$ est localement libre de type
fini, c'est-à-dire que (7) est vérifiée. Puisque $\Psi$ est surjective, il suffit de
montrer que $\Psi$ est injective. Puisque $R_{fg}$ est réduit,
il suffit de montrer que $\Psi$ est injective après localisation
en tous les idéaux premiers minimaux $\mathfrak p$ de $R_{fg}$ ; voir le
Lemme \ref{algebra-lemma-reduced-ring-sub-product-fields}.
Cependant, on sait que $R_\mathfrak p = \kappa(\mathfrak p)$
d'après le Lemme \ref{algebra-lemma-minimal-prime-reduced-ring}, et
$\rho_M(\mathfrak p) = r$ ; ainsi $\Psi_\mathfrak p :
R_\mathfrak p^{\oplus r} \to M \otimes_R \kappa(\mathfrak p)$
est un isomorphisme, puisqu'il s'agit d'une application surjective entre espaces vectoriels
de dimension finie et de même dimension.
\end{proof}

\begin{remark}
\label{algebra-remark-warning}
Il n'est pas vrai qu'un $R$-module de type fini qui soit
plat sur $R$ soit automatiquement projectif. Un contre-
exemple s'obtient en prenant $R = \mathcal{C}^\infty(\mathbf{R})$
égal à l'anneau des fonctions indéfiniment différentiables sur
$\mathbf{R}$, et $M = R_{\mathfrak m} = R/I$, où
$\mathfrak m = \{f \in R \mid f(0) = 0\}$ et
$I = \{f \in R \mid \exists \epsilon, \epsilon > 0 :
f(x) = 0\ \forall x, |x| < \epsilon\}$.
\end{remark}

\begin{lemma}
\label{algebra-lemma-finite-flat-local}
(Attention : voir la Remarque \ref{algebra-remark-warning}.)
Supposons que $R$ soit un anneau local et que $M$ soit un $R$-module
plat de type fini. Alors $M$ est libre de type fini.
\end{lemma}

\begin{proof}
Cela résulte du critère équationnel de platitude ; voir le
Lemme \ref{algebra-lemma-flat-eq}. Plus précisément, supposons que
$x_1, \ldots, x_r \in M$ aient pour images une base de
$M/\mathfrak mM$. D'après le Lemme de Nakayama \ref{algebra-lemma-NAK},
ces éléments engendrent $M$. Nous voulons montrer qu'il
n'existe aucune relation entre les $x_i$. Nous allons plutôt montrer
par récurrence sur $n$ que, si $x_1, \ldots, x_n \in M$
sont linéairement indépendants dans l'espace vectoriel
$M/\mathfrak mM$, alors ils sont indépendants sur $R$.

\medskip\noindent
Le cas initial de la récurrence est celui où l'on a
$x \in M$, $x \not\in \mathfrak mM$ et une relation
$fx = 0$. D'après le critère équationnel, il
existe $y_j \in M$ et $a_j \in R$ tels que
$x = \sum a_j y_j$ et $fa_j = 0$ pour tout $j$.
Puisque $x \not\in \mathfrak mM$, on voit qu'au moins
un $a_j$ est une unité, et donc que $f = 0 $.

\medskip\noindent
Supposons que $\sum f_i x_i$ soit une relation entre $x_1, \ldots, x_n$.
Par le choix des $x_i$, on a $f_i \in \mathfrak m$.
D'après le critère équationnel de platitude, il existe
$a_{ij} \in R$ et $y_j \in M$ tels que
$x_i = \sum a_{ij} y_j$ et $\sum f_i a_{ij} = 0$.
Puisque $x_n \not \in \mathfrak mM$, on voit que
$a_{nj}\not\in \mathfrak m$ pour au moins un $j$.
Puisque $\sum f_i a_{ij} = 0$, on obtient
$f_n = \sum_{i = 1}^{n-1} (-a_{ij}/a_{nj}) f_i$.
La relation $\sum f_i x_i = 0$ peut alors se réécrire
$\sum_{i = 1}^{n-1} f_i( x_i + (-a_{ij}/a_{nj}) x_n) = 0$.
Remarquons que les éléments $x_i + (-a_{ij}/a_{nj}) x_n$ ont pour images
$n-1$ éléments linéairement indépendants de $M/\mathfrak mM$.
Par hypothèse de récurrence, tous les $f_i$, $i \leq n-1$,
doivent être nuls, ainsi que $f_n = \sum_{i = 1}^{n-1} (-a_{ij}/a_{nj}) f_i$.
Cela établit l'étape de récurrence.
\end{proof}

\begin{lemma}
\label{algebra-lemma-finite-projective-descends}
Soit $R \to S$ un homomorphisme local plat d'anneaux locaux.
Soit $M$ un $R$-module de type fini. Alors $M$ est projectif de type fini
sur $R$ si et seulement si $M \otimes_R S$ est projectif de type fini
sur $S$.
\end{lemma}

\begin{proof}
D'après le Lemme \ref{algebra-lemma-finite-projective}, être projectif de type fini
sur un anneau local revient à être libre de type fini.
Supposons que $M \otimes_R S$ soit un $S$-module libre de type fini.
Choisissons $x_1, \ldots, x_r \in M$ dont les images dans $M/\mathfrak m_RM$
forment une base sur $\kappa(\mathfrak m)$. Alors
$x_1 \otimes 1, \ldots, x_r \otimes 1$
forment une base de $M \otimes_R S$. Cela implique que
l'application $R^{\oplus r} \to M, (a_i) \mapsto \sum a_i x_i$
devient un isomorphisme après tensorisation par $S$.
Par fidèle platitude de $R \to S$, voir le Lemme \ref{algebra-lemma-local-flat-ff},
on voit que c'est un isomorphisme.
\end{proof}

\begin{lemma}
\label{algebra-lemma-locally-free-semi-local-free}
Soit $R$ un anneau semi-local.
Soit $M$ un module localement libre de type fini.
Si $M$ est de rang constant, alors $M$ est libre. En particulier, si $R$ a
un spectre connexe, alors $M$ est libre.
\end{lemma}

\begin{proof}
Omis. Indications : montrons d'abord que $M/\mathfrak m_iM$ a la
même dimension $d$ pour tous les idéaux maximaux $\mathfrak m_1, \ldots, \mathfrak m_n$
de $R$, en utilisant la constance du rang.
Montrons ensuite qu'il existe des éléments $x_1, \ldots, x_d \in M$
qui forment une base de chaque $M/\mathfrak m_iM$, grâce au théorème
chinois des restes. Montrons enfin que $x_1, \ldots, x_d$ est une base de $M$.
\end{proof}

\noindent
Voici un lemme technique utilisé dans le chapitre sur les groupoïdes.

\begin{lemma}
\label{algebra-lemma-semi-local-module-basis-in-submodule}
Soit $R$ un anneau local d'idéal maximal $\mathfrak m$ et de
corps résiduel infini.
Soit $R \to S$ un homomorphisme d'anneaux.
Soit $M$ un $S$-module et soit $N \subset M$ un $R$-sous-module.
Supposons que
\begin{enumerate}
\item $S$ soit semi-local et que $\mathfrak mS$ soit contenu dans le
radical de Jacobson de $S$,
\item $M$ soit un $S$-module libre de type fini, et
\item $N$ engendre $M$ comme $S$-module.
\end{enumerate}
Alors $N$ contient une $S$-base de $M$.
\end{lemma}

\begin{proof}
Supposons $M$ libre de rang $n$. Soit $I \subset S$ le radical de Jacobson.
D'après le Lemme de Nakayama \ref{algebra-lemma-NAK}, une suite d'éléments
$m_1, \ldots, m_n$ est une base de $M$ si et seulement si
$\overline{m}_i \in M/IM$ engendrent $M/IM$. On peut donc remplacer
$M$ par $M/IM$, $N$ par $N/(N \cap IM)$, $R$ par $R/\mathfrak m$,
et $S$ par $S/IS$. Dans ce cas, on voit que $S$ est un produit fini
de corps $S = k_1 \times \ldots \times k_r$ et que
$M = k_1^{\oplus n} \times \ldots \times k_r^{\oplus n}$.
Le fait que $N \subset M$ engendre $M$ comme $S$-module
signifie qu'il existe $x_j \in N$ tels qu'une combinaison linéaire
$\sum a_j x_j$, avec $a_j \in S$, ait une composante non nulle dans chaque
facteur $k_i^{\oplus n}$.
Puisque $R = k$ est un corps infini, cela signifie aussi qu'une certaine
combinaison linéaire $y = \sum c_j x_j$, avec $c_j \in k$, a une
composante non nulle dans chaque facteur. Ainsi $y \in N$ engendre un
facteur direct libre $Sy \subset M$. Par récurrence sur $n$, le résultat
vaut pour $M/Sy$ et le sous-module $\overline{N} = N/(N \cap Sy)$.
Autrement dit, il existe $\overline{y}_2, \ldots, \overline{y}_n$
dans $\overline{N}$ qui engendrent librement $M/Sy$. Alors
$y, y_2, \ldots, y_n$ engendrent librement $M$, et l'on conclut.
\end{proof}

\begin{lemma}
\label{algebra-lemma-evaluation-map-iso-finite-projective}
Soit $R$ un anneau. Soient $L$, $M$, $N$ des $R$-modules.
L'application canonique
$$
\Hom_R(M, N) \otimes_R L \to \Hom_R(M, N \otimes_R L)
$$
est un isomorphisme si $M$ est projectif de type fini.
\end{lemma}

\begin{proof}
D'après le Lemme \ref{algebra-lemma-finite-projective}, on voit que $M$
est de présentation finie et localement libre de type fini.
D'après les Lemmes \ref{algebra-lemma-hom-from-finitely-presented} et
\ref{algebra-lemma-tensor-product-localization}, la formation du
membre de gauche et du membre de droite de la flèche commute à la
localisation. On peut vérifier que notre application est un isomorphisme
après localisation ; voir le Lemme \ref{algebra-lemma-cover}.
On peut donc supposer $M$ libre de type fini. Dans ce cas,
le lemme est immédiat.
\end{proof}


 


\section{Lieux ouverts définis par des homomorphismes de modules}
\label{algebra-section-loci-maps}

\noindent
L'ensemble des idéaux premiers où un homomorphisme de modules donné est surjectif ou est un isomorphisme
est parfois ouvert. Dans le cas de modules projectifs de type fini, on peut considérer
le rang de l'homomorphisme.

\begin{lemma}
\label{algebra-lemma-map-between-finite}
Soit $R$ un anneau. Soit $\varphi : M \to N$ un homomorphisme de $R$-modules,
où $N$ est un $R$-module de type fini. Alors on a l'égalité
\begin{align*}
U & = \{\mathfrak p \subset R \mid
\varphi_{\mathfrak p} : M_{\mathfrak p} \to N_{\mathfrak p}
\text{ est surjective}\} \\
& = \{\mathfrak p \subset R \mid
\varphi \otimes \kappa(\mathfrak p) :
M \otimes \kappa(\mathfrak p) \to N \otimes \kappa(\mathfrak p)
\text{ est surjective}\}
\end{align*}
et $U$ est un ouvert de $\Spec(R)$. De plus, pour tout $f \in R$
tel que $D(f) \subset U$, l'homomorphisme $M_f \to N_f$ est surjectif.
\end{lemma}

\begin{proof}
L'égalité de la formule affichée résulte du lemme de Nakayama.
Le lemme de Nakayama implique aussi que $U$ est ouvert. Voir le
Lemme \ref{algebra-lemma-NAK}, en particulier la partie (3). Si $D(f) \subset U$, alors
$M_f \to N_f$ est surjectif sur toutes les localisations aux idéaux premiers de
$R_f$, et il est donc surjectif d'après le Lemme \ref{algebra-lemma-characterize-zero-local}.
\end{proof}

\begin{lemma}
\label{algebra-lemma-map-between-finitely-presented}
Soit $R$ un anneau. Soit $\varphi : M \to N$ un homomorphisme de $R$-modules,
où $M$ est de type fini et $N$ de présentation finie. Alors
$$
U = \{\mathfrak p \subset R \mid
\varphi_{\mathfrak p} : M_{\mathfrak p} \to N_{\mathfrak p}
\text{ est un isomorphisme}\}
$$
est un ouvert de $\Spec(R)$.
\end{lemma}

\begin{proof}
Soit $\mathfrak p \in U$. Choisissons une présentation
$N = R^{\oplus n}/\sum_{j = 1, \ldots, m} R k_j$.
Notons $e_i$ l'image dans $N$ du $i$-ième vecteur de base de $R^{\oplus n}$.
Pour chaque $i \in \{1, \ldots, n\}$, choisissons un élément
$m_i \in M_{\mathfrak p}$ tel que $\varphi(m_i) = f_i e_i$
pour un certain $f_i \in R$, $f_i \not \in \mathfrak p$. C'est possible
car $\varphi_{\mathfrak p}$ est un isomorphisme. Posons $f = f_1 \ldots f_n$
et soit $\psi : R_f^{\oplus n} \to M_f$ l'homomorphisme qui envoie
le $i$-ième vecteur de base sur $m_i/f_i$. Remarquons que
$\varphi_f \circ \psi$ est la localisation
en $f$ de l'homomorphisme donné $R^{\oplus n} \to N$.
Comme $\varphi_{\mathfrak p}$ est un isomorphisme, on
voit que $\psi(k_j)$ est un élément de $M$ dont l'image est nulle
dans $M_{\mathfrak p}$. Il existe donc
$g_j \in R$, $g_j \not \in \mathfrak p$, tel que $g_j \psi(k_j) = 0$.
En posant $g = g_1 \ldots g_m$, on voit que $\psi_g$ se factorise par
$N_{fg}$ et donne un homomorphisme $\chi : N_{fg} \to M_{fg}$. Par construction,
$\chi$ est un inverse à droite de $\varphi_{fg}$. Il s'ensuit que
$\chi_\mathfrak p$ est un isomorphisme. D'après le Lemme \ref{algebra-lemma-map-between-finite},
il existe $h \in R$, $h \not \in \mathfrak p$, tel que
$\chi_h : N_{fgh} \to M_{fgh}$ soit surjectif.
Ainsi $\varphi_{fgh}$ et $\chi_h$ sont des homomorphismes inverses l'un de l'autre,
ce qui implique $D(fgh) \subset U$, comme souhaité.
\end{proof}

\begin{lemma}
\label{algebra-lemma-finitely-presented-localization-free}
Soit $R$ un anneau. Soit $\mathfrak p \subset R$ un idéal premier.
Soit $M$ un $R$-module de présentation finie. Si $M_\mathfrak p$
est libre, alors il existe $f \in R$, $f \not \in \mathfrak p$,
tel que $M_f$ soit un $R_f$-module libre.
\end{lemma}

\begin{proof}
Choisissons une base $x_1, \ldots, x_n \in M_\mathfrak p$. On peut choisir
$f \in R$, $f \not \in \mathfrak p$, tel que $x_i$ soit l'image
d'un certain $y_i \in M_f$. Après avoir remplacé $y_i$ par $f^m y_i$ avec $m \gg 0$,
on peut supposer $y_i \in M$. Cela revient à remplacer $x_1, \ldots, x_n$
par $f^mx_1, \ldots, f^mx_n$, qui est encore une base puisque l'image de $f$ est
une unité dans $R_\mathfrak p$. On obtient donc
un homomorphisme $\varphi = (y_1, \ldots, y_n) : R^{\oplus n} \to M$
de $R$-modules dont la localisation en $\mathfrak p$ est un isomorphisme.
D'après le Lemme \ref{algebra-lemma-map-between-finitely-presented},
on peut trouver $f \in R$, $f \not \in \mathfrak p$,
tel que $\varphi_\mathfrak q$ soit un isomorphisme pour tout idéal premier
$\mathfrak q \subset R$ vérifiant $f \not \in \mathfrak q$.
Il résulte alors du Lemme \ref{algebra-lemma-characterize-zero-local}
que $\varphi_f$ est un isomorphisme, ce qui achève la démonstration.
\end{proof}

\begin{lemma}
\label{algebra-lemma-cokernel-flat}
Soit $R$ un anneau. Soit $\varphi : P_1 \to P_2$ un homomorphisme de
modules projectifs de type fini. Alors
\begin{enumerate}
\item L'ensemble $U$ des idéaux premiers $\mathfrak p \in \Spec(R)$ tels que
$\varphi \otimes \kappa(\mathfrak p)$ soit injectif est ouvert et,
pour tout $f\in R$ tel que $D(f) \subset U$, on a
\begin{enumerate}
\item $P_{1, f} \to P_{2, f}$ est injectif, et
\item le module $\Coker(\varphi)_f$ est projectif de type fini sur $R_f$.
\end{enumerate}
\item L'ensemble $W$ des idéaux premiers $\mathfrak p \in \Spec(R)$ tels que
$\varphi \otimes \kappa(\mathfrak p)$ soit surjectif est ouvert et,
pour tout $f\in R$ tel que $D(f) \subset W$, on a
\begin{enumerate}
\item $P_{1, f} \to P_{2, f}$ est surjectif, et
\item le module $\Ker(\varphi)_f$ est projectif de type fini sur $R_f$.
\end{enumerate}
\item L'ensemble $V$ des idéaux premiers $\mathfrak p \in \Spec(R)$ tels que
$\varphi \otimes \kappa(\mathfrak p)$ soit un isomorphisme est ouvert et,
pour tout $f\in R$ tel que $D(f) \subset V$, l'homomorphisme
$\varphi : P_{1, f} \to P_{2, f}$ est un isomorphisme de modules sur $R_f$.
\end{enumerate}
\end{lemma}

\begin{proof}
Pour démontrer que l'ensemble $U$ est ouvert, on peut travailler localement sur
$\Spec(R)$. On peut donc remplacer $R$ par une localisation convenable
et supposer que $P_1 = R^{n_1}$ et $P_2 = R^{n_2}$ ; voir le Lemme
\ref{algebra-lemma-finite-projective}. Dans ce cas, l'injectivité de
$\varphi \otimes \kappa(\mathfrak p)$ équivaut à $n_1 \leq n_2$
et à ce qu'un certain mineur $n_1 \times n_1$, noté $f$, de la matrice de $\varphi$, soit
inversible dans $\kappa(\mathfrak p)$. Ainsi $D(f) \subset U$.
Cet argument montre aussi que $P_{1, \mathfrak p} \to P_{2, \mathfrak p}$
est injectif pour $\mathfrak p \in U$.

\medskip\noindent
Supposons maintenant $D(f) \subset U$. La remarque du paragraphe précédent
et le Lemme \ref{algebra-lemma-characterize-zero-local} montrent que
$P_{1, f} \to P_{2, f}$ est injectif, c'est-à-dire que (1)(a) est vérifiée.
D'après le Lemme \ref{algebra-lemma-finite-projective}, pour démontrer (1)(b),
il suffit de montrer que $\Coker(\varphi)$ est projectif de type fini
localement sur $D(f)$. Ainsi, comme on l'a vu ci-dessus, on peut
supposer que $P_1 = R^{n_1}$ et $P_2 = R^{n_2}$
et qu'un certain mineur de la matrice de $\varphi$ est inversible dans $R$.
Si le mineur considéré correspond aux $n_1$ premiers
vecteurs de base de $R^{n_2}$, alors les $n_2 - n_1$ derniers vecteurs de base
donnent un homomorphisme $R^{n_2 - n_1} \to R^{n_2} \to \Coker(\varphi)$
dont on vérifie facilement que c'est un isomorphisme.

\medskip\noindent
L'ouverture de $W$ et (2)(a) lorsque $D(f) \subset W$ résultent du
Lemme \ref{algebra-lemma-map-between-finite}. Puisque $P_{2, f}$ est projectif
sur $R_f$, on voit que $\varphi_f : P_{1, f} \to P_{2, f}$ possède une section,
et il s'ensuit que $\Ker(\varphi)_f$ est un facteur direct de $P_{2, f}$.
Par conséquent, $\Ker(\varphi)_f$ est projectif de type fini. Ainsi (2)(b) est également vérifiée.

\medskip\noindent
Il est clair que $V = U \cap W$ est ouvert, et l'autre assertion de (3)
résulte de (1)(a) et (2)(a).
\end{proof}






\section{Descente fidèlement plate de la projectivité des modules}
\label{algebra-section-ffdescent-projectivity-introduction}

\medskip\noindent
Dans les quelques sections suivantes, nous démontrons, à la suite de
Raynaud et Gruson \cite{GruRay}, que la
projectivité des modules descend le long des homomorphismes d'anneaux fidèlement plats. L'idée de la
démonstration est d'utiliser un dévissage à la Kaplansky \cite{Kaplansky} pour se ramener
au
cas des modules dénombrablement engendrés. Étant donnée une filtration convenable d'un
module $M$, le dévissage permet d'exprimer $M$ comme somme directe des quotients
successifs des sous-modules de la filtration (voir la
section \ref{algebra-section-transfinite-devissage}).
Grâce à cette technique, nous démontrons qu'un module projectif est une somme directe de
modules dénombrablement engendrés (Théorème \ref{algebra-theorem-projective-direct-sum}). Pour
démontrer la descente de la projectivité des modules dénombrablement engendrés, nous introduisons une
condition « de Mittag-Leffler » sur les modules, montrons qu'un module dénombrablement engendré
est projectif si et seulement s'il est plat et de Mittag-Leffler (Théorème
\ref{algebra-theorem-projectivity-characterization}), puis montrons que la propriété
d'être un module de Mittag-Leffler descend (Lemme
\ref{algebra-lemma-ffdescent-ML}). Enfin, étant donné un module quelconque $M$ dont le
changement de base par un homomorphisme fidèlement plat d'anneaux est projectif, nous filtrons $M$ par
des sous-modules dont les quotients successifs sont des modules projectifs
dénombrablement engendrés, puis le dévissage montre que $M$ est une somme directe de modules projectifs,
donc est lui-même projectif (Théorème \ref{algebra-theorem-ffdescent-projectivity}).

\medskip\noindent
Signalons qu'il y a une erreur dans la démonstration de la descente fidèlement plate
de la projectivité dans \cite{GruRay}. La descente
de la projectivité le long des homomorphismes d'anneaux fidèlement plats y est déduite de la descente de la
projectivité le long d'un type plus général d'homomorphisme d'anneaux
(\cite[exemple 3.1.4(1) de la partie II]{GruRay}).
Cependant, la démonstration de la descente le long de ce type plus général
d'homomorphisme est incorrecte. Dans \cite{G}, Gruson explique ce qui ne va pas,
mais ne fournit pas de correction pour le cas qui nous intéresse.
Combler cette lacune dans la démonstration de la descente fidèlement plate de la projectivité
revient à démontrer que la propriété d'être un module de Mittag-Leffler
descend le long des homomorphismes d'anneaux fidèlement plats. C'est ce que nous faisons dans le
Lemme \ref{algebra-lemma-ffdescent-ML}.





\section{Caractérisation de la platitude}
\label{algebra-section-characterize-flatness}

\noindent
Dans cette section, nous étudions des critères de platitude. Le résultat principal de
cette section est le théorème de Lazard (Théorème \ref{algebra-theorem-lazard}
ci-dessous), qui affirme qu'un module plat est la colimite d'un système filtrant de
modules libres de type fini.
Rappelons au lecteur le « critère équationnel de platitude » ; voir le
Lemme \ref{algebra-lemma-flat-eq}.
Il se trouve que l'on peut le transformer en une propriété apparemment bien plus forte.

\begin{lemma}
\label{algebra-lemma-flat-factors-free}
Soit $M$ un $R$-module. Les assertions suivantes sont équivalentes :
\begin{enumerate}
\item $M$ est plat.
\item Si $f: R^n \to M$ est un homomorphisme de modules et $x \in \Ker(f)$, alors il existe
des homomorphismes de modules $h: R^n \to R^m$ et $g: R^m \to M$ tels que
$f = g \circ h$ et $x \in \Ker(h)$.
\item Supposons que $f: R^n \to M$ soit un homomorphisme de modules, que $N \subset \Ker(f)$ soit un
sous-module quelconque et que $h: R^n \to R^{m}$ soit un homomorphisme tel que $N \subset \Ker(h)$
et que $f$ se factorise par $h$. Alors, pour tout $x \in \Ker(f)$, on peut trouver un homomorphisme
$h': R^n \to R^{m'}$ tel que $N + Rx \subset \Ker(h')$ et que $f$
se factorise par $h'$.
\item Si $f: R^n \to M$ est un homomorphisme de modules et si $N \subset \Ker(f)$ est un
sous-module de type fini, alors il existe des homomorphismes de modules $h: R^n \to
R^m$ et $g: R^m \to M$ tels que $f = g \circ h$ et $N \subset
\Ker(h)$.
\end{enumerate}
\end{lemma}

\begin{proof}
L'équivalence de (1) et (2) n'est qu'une reformulation du critère équationnel
de platitude\footnote{En fait, un homomorphisme de modules $f : R^n \to M$
correspond au choix d'éléments $x_1, x_2, \ldots, x_n$ de $M$
(à savoir les images des vecteurs de la base canonique $e_1, e_2, \ldots,
e_n$) ; de plus, un élément $x \in \Ker(f)$ correspond à une
relation entre ces $x_1, x_2, \ldots, x_n$ (à savoir la relation
$\sum_i f_i x_i = 0$, où les $f_i$ sont les coordonnées de $x$).
L'homomorphisme de modules $h$ (représenté par une matrice $m \times n$)
correspond à la matrice $(a_{ij})$ du Lemme \ref{algebra-lemma-flat-eq},
et les $y_j$ du Lemme \ref{algebra-lemma-flat-eq} sont les images des
vecteurs de la base canonique de $R^m$ par $g$.}. Pour montrer que
(2) implique (3), soit $g: R^m \to M$
l'homomorphisme tel que $f$ se factorise sous la forme $f = g \circ h$. D'après (2), choisissons $h'': R^m
\to R^{m'}$ tel que $h''$ annule $h(x)$ et que $g: R^m \to M$
se factorise par $h''$. Alors $h' = h'' \circ h$ convient. L'assertion (3) implique (4)
par récurrence sur le nombre de générateurs de $N \subset \Ker(f)$ dans (4).
Il est clair que (4) implique (2).
\end{proof}

\begin{lemma}
\label{algebra-lemma-flat-factors-fp}
Soit $M$ un $R$-module. Alors $M$ est plat si et seulement si la condition
suivante est satisfaite : si $P$ est un $R$-module de présentation finie et $f: P
\to M$ un homomorphisme de modules, alors il existe un $R$-module libre de type fini $F$ et des
homomorphismes de modules $h: P \to F$ et $g: F \to M$ tels que $f = g
\circ h$.
\end{lemma}

\begin{proof}
Ce n'est qu'une reformulation de la condition (4) du
Lemme \ref{algebra-lemma-flat-factors-free}.
\end{proof}

\begin{lemma}
\label{algebra-lemma-flat-surjective-hom}
Soit $M$ un $R$-module. Alors $M$ est plat si et seulement si la condition
suivante est satisfaite : pour tout $R$-module de présentation finie $P$, si $N \to
M$ est un homomorphisme surjectif de $R$-modules, alors l'homomorphisme induit $\Hom_R(P, N)
\to \Hom_R(P, M)$ est surjectif.
\end{lemma}

\begin{proof}
Supposons d'abord $M$ plat. Il faut montrer que, si $P$ est de présentation finie,
alors toute application $f: P \to M$ se factorise par l'application $N
\to M$. D'après le
Lemme \ref{algebra-lemma-flat-factors-fp},
l'application $f$ se factorise par une application $F \to M$, où $F$ est libre de type fini.
Puisque $F$ est libre, cette application se factorise par
$N \to M$. Ainsi $f$ se factorise par $N \to M$.

\medskip\noindent
Réciproquement, supposons que la condition du lemme soit satisfaite. Soit $f: P \to M$
une application issue d'un module de présentation finie $P$. Choisissons un module libre $N$ et une
surjection $N \to M$ sur $M$. Alors $f$ se factorise par $N \to
M$ et, puisque $P$ est de type fini, $f$ se factorise par un sous-module libre
de type fini de $N$. Ainsi $M$ satisfait à la condition du Lemme
\ref{algebra-lemma-flat-factors-fp}, et est donc plat.
\end{proof}

\begin{theorem}[Théorème de Lazard]
\label{algebra-theorem-lazard}
Soit $M$ un $R$-module. Alors $M$ est plat si et seulement s'il est la colimite
d'un système filtrant de $R$-modules libres de type fini.
\end{theorem}

\begin{proof}
Une colimite d'un système filtrant de modules plats est plate, puisque les colimites
filtrantes sont exactes et commutent au produit tensoriel. Ainsi, si $M$ est la colimite
d'un système filtrant de modules libres de type fini, alors $M$ est plat.

\medskip\noindent
Pour la réciproque, rappelons d'abord que tout module $M$ peut s'écrire comme la
colimite d'un système filtrant de modules de présentation finie, de la manière
suivante. Choisissons une surjection $f: R^I \to M$ pour un certain ensemble $I$, et soit $K$
son noyau. Soit $E$ l'ensemble des couples ordonnés $(J, N)$, où $J$ est une
partie finie de $I$ et $N$ un sous-module de type fini de $R^J \cap K$.
On fait de $E$ un ensemble ordonné filtrant en posant $(J, N) \leq
(J', N')$ si et seulement si $J \subset J'$ et $N \subset N'$. Définissons $M_e =
R^J/N$ pour $e = (J, N)$, et définissons $f_{ee'}: M_e \to M_{e'}$ comme
l'application naturelle pour $e \leq e'$. Alors $(M_e, f_{ee'})$ est un système filtrant, et
les applications naturelles $f_e: M_e \to M$ induisent un isomorphisme
$\colim_{e
\in E} M_e \xrightarrow{\cong} M$.

\medskip\noindent
Supposons maintenant $M$ plat. Soit $I = M \times \mathbf{Z}$, notons $(x_i)$ la
base canonique de $R^{I}$ et prenons, dans la discussion ci-dessus, $f: R^I
\to M$ égal à l'application qui envoie $x_i$ sur la projection de $i$ sur $M$.
Pour démontrer le théorème, il suffit de montrer que les $e \in E$ tels que $M_e$
soit libre forment une partie cofinale de $E$. Soit donc $e = (J, N) \in E$ quelconque.
D'après le Lemme \ref{algebra-lemma-flat-factors-fp}, il existe un module libre de type fini $F$ et des applications
$h: R^J/N \to F$ et $g: F \to M$ telles que l'application naturelle
$f_e: R^J/N \to M$ se factorise sous la forme $R^J/N \xrightarrow{h} F
\xrightarrow{g} M$. Nous allons réaliser $F$ comme $M_{e'}$ pour un certain $e' \geq
e$.

\medskip\noindent
Soit $\{ b_1, \ldots, b_n \}$ une base finie de $F$. Choisissons $n$ éléments distincts
$i_1, \ldots, i_n \in I$ tels que $i_{\ell} \notin J$ pour tout $\ell$,
et tels que l'image de $x_{i_{\ell}}$ par $f: R^I \to M$ soit égale
à l'image de $b_{\ell}$ par $g: F \to M$. C'est possible puisque
tout élément de $M$ peut s'écrire sous la forme $f(x_i)$ pour une infinité d'éléments
distincts $i \in I$ (par le choix de $I$). Posons maintenant
$J' = J \cup \{i_1, \ldots , i_n \}$, et définissons $R^{J'}
\to F$ par $x_i \mapsto h(x_i)$ pour $i \in J$ et $x_{i_{\ell}} \mapsto
b_{\ell}$ pour $\ell = 1, \ldots, n$. Soit $N' = \Ker(R^{J'} \to F)$.
Remarquons que :
\begin{enumerate}
\item Le carré
$$
\xymatrix{
R^{J'} \ar[r] \ar@{^{(}->}[d] & F \ar[d]^{g} \\
R^{I} \ar[r]_{f} & M
}
$$
est commutatif,
%$R^{J'} \to F$ factors $f: R^I \to M$,
donc $N' \subset K = \Ker(f)$ ;
\item $R^{J'} \to F$ est une surjection sur un module libre de type fini, elle
est donc scindée, et $N'$ est de type fini ;
\item $J \subset J'$ et $N \subset N'$.
\end{enumerate}
D'après (1) et (2), $e' = (J', N')$ appartient à $E$ ; d'après (3), $e' \geq e$ ; et, par
construction, $M_{e'} = R^{J'}/N' \cong F$ est libre.
\end{proof}

\section{Morphismes de modules universellement injectifs}
\label{algebra-section-universally-injective}

\noindent
Nous étudions maintenant les morphismes de modules universellement injectifs, qui
sont en un certain sens complémentaires des modules plats (voir le lemme
\ref{algebra-lemma-flat-universally-injective}). Nous suivons la thèse de Lazard
\cite{Autour} ; voir aussi \cite{Lam}.

\begin{definition}
\label{algebra-definition-universally-injective}
Soit $f: M \to N$ un morphisme de $R$-modules. On dit que $f$ est
{\it universellement injectif} si, pour tout $R$-module $Q$, le morphisme $f
\otimes_R \text{id}_Q: M \otimes_R Q \to N \otimes_R Q$
est injectif. Une suite $0 \to M_1 \to M_2 \to M_3
\to 0$ de $R$-modules est dite {\it universellement exacte} si elle est exacte
et si $M_1 \to M_2$ est universellement injectif.
\end{definition}

\begin{example}
\label{algebra-example-universally-exact}
Exemples de suites universellement exactes.
\begin{enumerate}
\item Une suite exacte courte scindée est universellement exacte, puisque le produit tensoriel
commute aux sommes directes.
\item La colimite d'un système inductif de suites universellement exactes est
universellement exacte. Cela résulte de l'exactitude des colimites filtrantes et
du fait que le produit tensoriel commute aux colimites. En particulier, la
colimite d'un système inductif de suites exactes scindées est universellement exacte. Nous
verrons plus bas que, réciproquement, toute suite universellement exacte s'obtient de
cette manière.
\end{enumerate}
\end{example}

\noindent
Nous donnons maintenant une liste de critères pour qu'une suite exacte courte soit universellement
exacte. Ils sont analogues aux critères de platitude donnés plus haut. Les assertions (3)--(6)
ci-dessous correspondent respectivement aux critères de platitude donnés dans les
lemmes \ref{algebra-lemma-flat-eq}, \ref{algebra-lemma-flat-factors-free},
\ref{algebra-lemma-flat-surjective-hom} et le
théorème \ref{algebra-theorem-lazard}.

\begin{theorem}
\label{algebra-theorem-universally-exact-criteria}
Soit
$$
0 \to M_1 \xrightarrow{f_1} M_2 \xrightarrow{f_2} M_3 \to 0
$$
une suite exacte de $R$-modules. Les assertions suivantes sont équivalentes :
\begin{enumerate}
\item La suite $0 \to M_1 \to M_2 \to M_3
\to 0$ est universellement exacte.
\item Pour tout $R$-module de présentation finie $Q$, la suite
$$
0 \to M_1 \otimes_R Q \to M_2 \otimes_R Q \to
M_3 \otimes_R Q \to 0
$$
est exacte.
\item Étant donnés des éléments $x_i \in M_1$ $(i = 1, \ldots, n)$, $y_j \in M_2$ $(j = 1,
\ldots, m)$ et $a_{ij} \in R$ $(i = 1, \ldots, n, j = 1, \ldots, m)$ tels que,
pour tout $i$,
$$
f_1(x_i) = \sum\nolimits_j a_{ij} y_j,
$$
il existe $z_j \in M_1$ $(j =1, \ldots, m)$ tels que, pour tout $i$,
$$
x_i = \sum\nolimits_j a_{ij} z_j .
$$
\item Étant donné un diagramme commutatif de morphismes de $R$-modules
$$
\xymatrix{
R^n \ar[r] \ar[d] &  R^m \ar[d] \\
M_1 \ar[r]^{f_1}        &  M_2
}
$$
où $m$ et $n$ sont des entiers, il existe un morphisme $R^m \to M_1$ rendant
commutatif le triangle supérieur.
\item Pour tout $R$-module de présentation finie $P$, le morphisme de $R$-modules
$\Hom_R(P, M_2) \to \Hom_R(P, M_3)$ est surjectif.
\item La suite $0 \to M_1 \to M_2 \to M_3
\to 0$ est la colimite d'un système inductif de suites exactes scindées de
la forme
$$
0 \to M_{1} \to M_{2, i} \to M_{3, i} \to 0
$$
où les $M_{3, i}$ sont de présentation finie.
\end{enumerate}
\end{theorem}

\begin{proof}
Il est clair que (1) implique (2).

\medskip\noindent
Montrons maintenant que (2) implique (3). Soient $f_1(x_i) = \sum_j a_{ij} y_j$ des relations
comme en (3). Soient $(d_j)$ une base de $R^m$, $(e_i)$ une base de $R^n$, et
$R^m \to R^n$ le morphisme donné par $d_j \mapsto \sum_i a_{ij} e_i$.
Soit $Q$ le conoyau de $R^m \to R^n$. En tensorisant
$R^m \to R^n \to Q \to 0$ par le morphisme $f_1: M_1 \to M_2$, on obtient un
diagramme commutatif
$$
\xymatrix{
M_1^{\oplus m} \ar[r] \ar[d] & M_1^{\oplus n} \ar[r] \ar[d] & M_1 \otimes_R Q
\ar[r] \ar[d] & 0 \\
M_2^{\oplus m} \ar[r] & M_2^{\oplus n} \ar[r] & M_2 \otimes_R Q \ar[r] & 0
}
$$
où $M_1^{\oplus m} \to M_1^{\oplus n}$ est donné par
$$
(z_1, \ldots, z_m) \mapsto
(\sum\nolimits_j a_{1j} z_j, \ldots, \sum\nolimits_j a_{nj} z_j),
$$
et où $M_2^{\oplus m} \to M_2^{\oplus n}$ est défini de même. Nous voulons
montrer que $x = (x_1, \ldots, x_n) \in M_1^{\oplus n}$ appartient à l'image de $M_1^{\oplus
m} \to M_1^{\oplus n}$. Par (2), le morphisme $M_1 \otimes Q \to M_2
\otimes Q$ est injectif ; par exactitude de la ligne supérieure, il suffit donc de montrer que
$x$ s'envoie sur $0$ dans $M_2 \otimes Q$ ; par exactitude de la ligne inférieure, il suffit ainsi
de montrer que l'image de $x$ dans $M_2^{\oplus n}$ appartient à l'image de
$M_2^{\oplus m} \to M_2^{\oplus n}$. Cela résulte de l'hypothèse.

\medskip\noindent
La condition (4) n'est que la traduction de (3) sous forme de diagramme.

\medskip\noindent
Montrons maintenant que (4) implique (5). Soit $\varphi : P \to M_3$ un morphisme issu d'un
$R$-module de présentation finie $P$. Nous devons montrer que $\varphi$ se relève en un morphisme
$P \to M_2$. Choisissons une présentation de $P$,
$$
R^n \xrightarrow{g_1} R^m \xrightarrow{g_2} P \to 0.
$$
La liberté de $R^n$ et de $R^m$ permet de construire $h_2: R^m \to M_2$
puis $h_1: R^n \to M_1$ de sorte que le diagramme suivant soit commutatif :
$$
\xymatrix{
         & R^n \ar[r]^{g_1} \ar[d]^{h_1} & R^m \ar[r]^{g_2} \ar[d]^{h_2} & P
\ar[r] \ar[d]^{\varphi} & 0 \\
0 \ar[r] & M_1 \ar[r]^{f_1} & M_2 \ar[r]^{f_2} & M_3 \ar[r] & 0 .
}
$$
Par (4), il existe un morphisme $k_1: R^m \to M_1$ tel que $k_1 \circ g_1 =
h_1$. Définissons maintenant $h'_2: R^m \to M_2$ par $h_2' = h_2 - f_1 \circ k_1$.
Alors
$$
h'_2 \circ g_1 = h_2 \circ g_1 - f_1 \circ k_1 \circ g_1 =
h_2 \circ g_1 - f_1 \circ h_1 = 0 .
$$
Ainsi, par passage au quotient, $h'_2$ définit un morphisme $\varphi': P \to
M_2$ tel que $\varphi' \circ g_2 = h_2'$. Sous forme de diagramme, on a
$$
\xymatrix{
R^m \ar[r]^{g_2} \ar[d]_{h'_2} & P \ar[d]^{\varphi} \ar[dl]_{\varphi'} \\
M_2 \ar[r]^{f_2} & M_3.
}
$$
où le triangle supérieur est commutatif. Nous affirmons que $\varphi'$ est le relèvement voulu,
c'est-à-dire que $f_2 \circ \varphi' = \varphi$. D'après les définitions, on a
$$
f_2 \circ \varphi' \circ g_2 = f_2 \circ h'_2 =
f_2 \circ h_2 - f_2 \circ f_1 \circ k_1 = f_2 \circ h_2 =
\varphi \circ g_2.
$$
Comme $g_2$ est surjectif, ceci achève la démonstration.

\medskip\noindent
Montrons maintenant que (5) implique (6). Écrivons $M_{3}$ comme colimite d'un système inductif
de modules de présentation finie $M_{3, i}$, voir le
lemme \ref{algebra-lemma-module-colimit-fp}. Soit $M_{2, i}$ le produit fibré
de $M_{3, i}$ et $M_{2}$ au-dessus de $M_{3}$ ; par définition, c'est le sous-module de
$M_2 \times M_{3, i}$ formé des éléments dont les deux projections sur $M_3$
sont égales. Soit $M_{1, i}$ le noyau de la projection
$M_{2, i} \to M_{3, i}$. On obtient alors un système inductif de suites exactes
$$
0 \to M_{1, i} \to M_{2, i} \to M_{3, i} \to 0,
$$
et, pour chaque $i$, un morphisme de suites exactes
$$
\xymatrix{
0  \ar[r] & M_{1, i} \ar[d] \ar[r]  &  M_{2, i}  \ar[r] \ar[d] & M_{3, i} \ar[d]
\ar[r] & 0 \\
0  \ar[r] & M_{1} \ar[r]  &  M_{2}  \ar[r] & M_{3} \ar[r] & 0
}
$$
compatible au système inductif. Il résulte de la définition du produit fibré
$M_{2, i}$ que le morphisme $M_{1, i} \to M_1$ est un isomorphisme.
 Par (5), il existe un morphisme $M_{3, i} \to M_{2}$ relevant $M_{3, i} \to
M_3$ et, par la propriété universelle du produit fibré, celui-ci donne lieu à une
section de $M_{2, i} \to M_{3, i}$. Ainsi, les suites
$$
0 \to M_{1, i} \to M_{2, i} \to M_{3, i} \to 0
$$
sont scindées. En passant à la colimite, on obtient un diagramme commutatif
$$
\xymatrix{
0  \ar[r] & \colim M_{1, i} \ar[d]^{\cong} \ar[r]  &  \colim M_{2, i}
 \ar[r]
\ar[d] & \colim M_{3, i} \ar[d]^{\cong} \ar[r] & 0 \\
0  \ar[r] & M_{1} \ar[r]  &  M_{2}  \ar[r] & M_{3} \ar[r] & 0
}
$$
à lignes exactes et dont les flèches verticales extérieures sont des isomorphismes. Ainsi, $\colim
M_{2, i}
\to M_2$ est également un isomorphisme, et (6) est satisfaite.

\medskip\noindent
La condition (6) implique (1) d'après
l'exemple \ref{algebra-example-universally-exact} (2).
\end{proof}

\noindent
Le théorème précédent montre qu'une suite universellement exacte est toujours une
colimite de suites exactes courtes scindées. Si le conoyau d'un morphisme universellement
injectif est de présentation finie, alors le morphisme lui-même est en fait scindé :

\begin{lemma}
\label{algebra-lemma-universally-exact-split}
Soit
$$
0 \to M_1 \to M_2 \to M_3 \to 0
$$
une suite exacte de $R$-modules. Supposons que $M_3$ soit de présentation finie.
Alors
$$
0 \to M_1 \to M_2 \to M_3 \to 0
$$
est universellement exacte si et seulement si elle est scindée.
\end{lemma}

\begin{proof}
Une suite exacte courte scindée est toujours universellement exacte, voir
l'exemple \ref{algebra-example-universally-exact}.
Réciproquement, si la suite est universellement exacte, alors le
théorème \ref{algebra-theorem-universally-exact-criteria} (5),
appliqué à $P = M_3$, montre que le morphisme $M_2 \to M_3$ admet une section.
\end{proof}

\noindent
Le lemme suivant montre en quel sens les morphismes universellement injectifs sont complémentaires des
modules plats.

\begin{lemma}
\label{algebra-lemma-flat-universally-injective}
Soit $M$ un $R$-module. Alors $M$ est plat si et seulement si toute suite exacte
de $R$-modules
$$
0 \to M_1 \to M_2 \to M \to 0
$$
est universellement exacte.
\end{lemma}

\begin{proof}
Cela résulte du lemme \ref{algebra-lemma-flat-surjective-hom} et du
théorème \ref{algebra-theorem-universally-exact-criteria} (5).
\end{proof}

\begin{example}
\label{algebra-example-universally-exact-non-split-non-flat}
Suites universellement exactes non scindées et non plates.
\begin{enumerate}
\item Malgré le
lemme \ref{algebra-lemma-universally-exact-split},
il peut exister une suite exacte courte de $R$-modules
$$
0 \to M_1 \to M_2 \to M_3 \to 0
$$
qui soit universellement exacte sans être scindée. Par exemple, prenons $R = \mathbf{Z}$,
$M_1 = \bigoplus_{n=1}^{\infty} \mathbf{Z}$, $M_{2} = \prod_{n =
1}^{\infty} \mathbf{Z}$, et pour $M_{3}$ le conoyau de l'inclusion $M_1
\to M_2$. Alors $M_1, M_2, M_3$ sont tous plats, car ils sont sans torsion
(Compléments d'algèbre, lemme \ref{more-algebra-lemma-dedekind-torsion-free-flat}) ;
par le lemme \ref{algebra-lemma-flat-universally-injective},
$$
0 \to M_1 \to M_2 \to M_3 \to 0
$$
est donc universellement exacte. Il ne peut toutefois exister de section $s: M_3 \to
M_2$. En effet, si $x$ est l'image de $(2, 2^2, 2^3, \ldots) \in M_2$ dans $M_3$,
alors tout morphisme de modules $s: M_3 \to M_2$ doit annuler $x$. En effet, $x
\in 2^n M_3$ pour tout $n \geq 1$ ; ainsi, $s(x)$ est divisible par $2^n$ pour tout $n
\geq 1$ et doit donc être $0$.
\item Malgré le lemme \ref{algebra-lemma-flat-universally-injective}, il peut exister
une suite exacte courte de $R$-modules
$$
0 \to M_1 \to M_2 \to M_3 \to 0
$$
qui soit universellement exacte alors que $M_1, M_2, M_3$ sont tous non plats. En effet, si $M$
est un module non plat quelconque, il suffit de prendre la suite exacte scindée
$$
0 \to M \to M \oplus M \to M \to 0.
$$
Par exemple, pour $R = \mathbf{Z}$, on peut prendre pour $M$ un module de torsion quelconque.
\item En prenant la somme directe d'une suite exacte comme en (1) et d'une autre comme en (2),
on obtient une suite exacte courte de $R$-modules
$$
0 \to M_1 \to M_2 \to M_3 \to 0
$$
qui est universellement exacte, non scindée, et telle que
$M_1, M_2, M_3$ sont tous non plats.
\end{enumerate}
\end{example}

\begin{lemma}
\label{algebra-lemma-ui-flat-domain}
Soit $0 \to M_1 \to M_2 \to M_3 \to 0$ une suite universellement exacte
de $R$-modules, et supposons $M_2$ plat.
Alors $M_1$ et $M_3$ sont plats.
\end{lemma}

\begin{proof}
Soit $0 \to N \to N' \to N'' \to 0$ une suite exacte courte de
$R$-modules. Considérons le diagramme commutatif
$$
\xymatrix{
M_1 \otimes_R N \ar[r] \ar[d] &
M_2 \otimes_R N \ar[r] \ar[d] &
M_3 \otimes_R N \ar[d] \\
M_1 \otimes_R N' \ar[r] \ar[d] &
M_2 \otimes_R N' \ar[r] \ar[d] &
M_3 \otimes_R N' \ar[d] \\
M_1 \otimes_R N'' \ar[r] &
M_2 \otimes_R N'' \ar[r] &
M_3 \otimes_R N''
}
$$
(nous avons omis les $0$ du bord).
Par hypothèse, les lignes sont des suites exactes courtes et la flèche
$M_2 \otimes N \to M_2 \otimes N'$ est injective. Cela implique clairement
que $M_1 \otimes N \to M_1 \otimes N'$ est injective, et l'on voit que $M_1$
est plat. En particulier, les colonnes de gauche et du milieu sont des suites
exactes courtes. Une chasse au diagramme montre alors que la flèche
$M_3 \otimes N \to M_3 \otimes N'$ est injective. Ainsi, $M_3$ est plat.
\end{proof}

\begin{lemma}
\label{algebra-lemma-universally-injective-tensor}
Soit $R$ un anneau.
Soit $M \to M'$ un morphisme universellement injectif de $R$-modules.
Alors, pour tout $R$-module $N$, le morphisme $M \otimes_R N \to M' \otimes_R N$
est universellement injectif.
\end{lemma}

\begin{proof}
Démonstration omise.
\end{proof}

\begin{lemma}
\label{algebra-lemma-composition-universally-injective}
Soit $R$ un anneau. Toute composée de morphismes universellement injectifs de
$R$-modules est universellement injective.
\end{lemma}

\begin{proof}
Démonstration omise.
\end{proof}

\begin{lemma}
\label{algebra-lemma-universally-injective-permanence}
Soit $R$ un anneau. Soient $M \to M'$ et $M' \to M''$ des morphismes de $R$-modules.
Si leur composée $M \to M''$ est universellement injective, alors
$M \to M'$ est universellement injectif.
\end{lemma}

\begin{proof}
Démonstration omise.
\end{proof}

\begin{lemma}
\label{algebra-lemma-universally-injective-check-stalks}
Soit $R \to S$ un homomorphisme d'anneaux.
Soit $M \to M'$ un morphisme de $S$-modules.
Les assertions suivantes sont équivalentes :
\begin{enumerate}
\item $M \to M'$ est universellement injectif comme morphisme de $R$-modules,
\item pour tout idéal premier $\mathfrak q$ de $S$, le morphisme
$M_{\mathfrak q} \to M'_{\mathfrak q}$ est universellement injectif
comme morphisme de $R$-modules,
\item pour tout idéal maximal $\mathfrak m$ de $S$, le morphisme
$M_{\mathfrak m} \to M'_{\mathfrak m}$ est universellement injectif
comme morphisme de $R$-modules,
\item pour tout idéal premier $\mathfrak q$ de $S$, le morphisme
$M_{\mathfrak q} \to M'_{\mathfrak q}$ est universellement injectif
comme morphisme de $R_{\mathfrak p}$-modules, où $\mathfrak p$ est l'image
réciproque de $\mathfrak q$ dans $R$, et
\item pour tout idéal maximal $\mathfrak m$ de $S$, le morphisme
$M_{\mathfrak m} \to M'_{\mathfrak m}$ est universellement injectif
comme morphisme de $R_{\mathfrak p}$-modules, où $\mathfrak p$ est l'image
réciproque de $\mathfrak m$ dans $R$.
\end{enumerate}
\end{lemma}

\begin{proof}
Soit $N$ un $R$-module. Soit $\mathfrak q$ un idéal premier de $S$ au-dessus de
l'idéal premier $\mathfrak p$ de $R$. On a alors
$$
(M \otimes_R N)_{\mathfrak q} =
M_{\mathfrak q} \otimes_R N =
M_{\mathfrak q} \otimes_{R_{\mathfrak p}} N_{\mathfrak p}.
$$
La même chose vaut en outre pour $M'$, et la localisation est exacte.
De plus, si $N$ est un $R_{\mathfrak p}$-module, alors $N_{\mathfrak p} = N$.
Ces remarques permettent de démontrer directement les équivalences.

\medskip\noindent
Supposons par exemple que (5) soit satisfaite. Soit
$K = \Ker(M \otimes_R N \to M' \otimes_R N)$. D'après les remarques
précédentes, $K_{\mathfrak m} = 0$ pour tout idéal maximal $\mathfrak m$
de $S$. Ainsi $K = 0$ d'après le
lemme \ref{algebra-lemma-characterize-zero-local}.
L'assertion (1) est donc satisfaite. Réciproquement, supposons (1) satisfaite. Prenons un
$\mathfrak q \subset S$ quelconque au-dessus de $\mathfrak p \subset R$.
Prenons un module $N$ quelconque sur $R_{\mathfrak p}$. Par hypothèse,
$\Ker(M \otimes_R N \to M' \otimes_R N) = 0$.
Les formules ci-dessus et le fait que $N = N_{\mathfrak p}$
montrent alors que
$\Ker(M_{\mathfrak q} \otimes_{R_{\mathfrak p}} N \to
M'_{\mathfrak q} \otimes_{R_{\mathfrak p}} N) = 0$. Autrement dit,
(4) est satisfaite. Bien sûr, (4) $\Rightarrow$ (5) est immédiat. Ainsi,
(1), (4) et (5) sont toutes équivalentes.
Nous omettons la démonstration des autres équivalences.
\end{proof}

\begin{lemma}
\label{algebra-lemma-universally-injective-localize}
Soit $\varphi : A \to B$ un homomorphisme d'anneaux. Soient $S \subset A$ et
$S' \subset B$ des parties multiplicatives telles que $\varphi(S) \subset S'$.
Soit $M \to M'$ un morphisme de $B$-modules.
\begin{enumerate}
\item Si $M \to M'$ est universellement injectif comme morphisme de $A$-modules,
alors $(S')^{-1}M \to (S')^{-1}M'$ est universellement injectif comme morphisme de
$A$-modules et comme morphisme de $S^{-1}A$-modules.
\item Si $M$ et $M'$ sont des $(S')^{-1}B$-modules, alors $M \to M'$
est universellement injectif comme morphisme de $A$-modules si et seulement s'il
est universellement injectif comme morphisme de $S^{-1}A$-modules.
\end{enumerate}
\end{lemma}

\begin{proof}
On peut démontrer ceci à l'aide du
lemme \ref{algebra-lemma-universally-injective-check-stalks},
mais on peut aussi le démontrer directement comme suit.
Supposons $M \to M'$ universellement injectif comme morphisme de $A$-modules.
Soit $Q$ un $A$-module. Alors $Q \otimes_A M \to Q \otimes_A M'$
est injectif. L'exactitude de la localisation montre que
$(S')^{-1}(Q \otimes_A M) \to (S')^{-1}(Q \otimes_A M')$ est injectif.
Comme $(S')^{-1}(Q \otimes_A M) = Q \otimes_A (S')^{-1}M$, et de même pour $M'$,
on voit que
$Q \otimes_A (S')^{-1}M \to Q \otimes_A (S')^{-1}M'$ est injectif ; ainsi,
$(S')^{-1}M \to (S')^{-1}M'$ est universellement injectif comme morphisme de
$A$-modules. Ceci démontre la première partie de (1).
Pour établir (2), on peut utiliser les deux faits suivants : (a) si $Q$ est un
$S^{-1}A$-module, alors $Q \otimes_A S^{-1}A = Q$, c'est-à-dire que tensoriser
par $Q$ sur $A$ revient à tensoriser par $Q$ sur $S^{-1}A$ ;
(b) si $M$ est un $A$-module quelconque sur lequel les éléments de $S$ sont inversibles,
alors $M \otimes_A Q = M \otimes_{S^{-1}A} S^{-1}Q$.
L'assertion (2) en résulte immédiatement.
\end{proof}

\begin{lemma}
\label{algebra-lemma-check-universally-injective-into-flat}
Soit $R$ un anneau et soit $M \to M'$ un morphisme de $R$-modules.
Si $M'$ est plat, alors $M \to M'$ est universellement injectif si
et seulement si $M/IM \to M'/IM'$ est injectif pour tout idéal de type
fini $I$ de $R$.
\end{lemma}

\begin{proof}
Il suffit de montrer que $M \otimes_R Q \to M' \otimes_R Q$ est
injectif pour tout $R$-module de type fini $Q$, voir le
théorème \ref{algebra-theorem-universally-exact-criteria}.
Alors $Q$ possède une filtration finie
$0 = Q_0 \subset Q_1 \subset \ldots \subset Q_n = Q$
par des sous-modules dont les sous-quotients
sont isomorphes à des modules cycliques $R/I_i$, voir le
lemme \ref{algebra-lemma-trivial-filter-finite-module}.
Comme $M'$ est plat, on obtient une filtration
$$
\xymatrix{
M \otimes Q_1 \ar[r] \ar[d] &
M \otimes Q_2 \ar[r] \ar[d] &
\ldots \ar[r] &
M \otimes Q \ar[d] \\
M' \otimes Q_1 \ar@{^{(}->}[r] &
M' \otimes Q_2 \ar@{^{(}->}[r] &
\ldots \ar@{^{(}->}[r] &
M' \otimes Q
}
$$
de $M' \otimes_R Q$ par les sous-modules $M' \otimes_R Q_i$, dont les quotients
successifs sont $M' \otimes_R R/I_i = M'/I_iM'$. Une récurrence simple
montre qu'il suffit de vérifier l'injectivité de $M/I_i M \to M'/I_i M'$.
Notons que l'ensemble des idéaux de type fini $I'_i \subset I_i$
est filtrant. Ainsi, $M/I_iM = \colim M/I'_iM$ est une colimite
filtrante, et de même pour $M'$ ; les morphismes $M/I'_iM \to M'/I'_i M'$ sont
injectifs par hypothèse et, puisque les colimites filtrantes sont exactes
(lemme \ref{algebra-lemma-directed-colimit-exact}), on conclut.
\end{proof}

\begin{lemma}
\label{algebra-lemma-base-change-along-universally-injective-ring-map}
Soit $R \to S$ un homomorphisme d'anneaux universellement injectif
comme morphisme de $R$-modules. Alors le foncteur $M \mapsto M \otimes_R S$
sur les $R$-modules reflète les morphismes injectifs, surjectifs et les isomorphismes.
\end{lemma}

\begin{proof}
Soit $M \to N$ un morphisme de $R$-modules, de noyau $K$ et de conoyau $Q$.
Si $M \otimes_R S \to N \otimes_R S$ est injectif, alors
l'image de $K \otimes_R S \to M \otimes_R S$ est nulle.
Comme $K \subset K \otimes_R S$ et $M \subset M \otimes_R S$,
on en déduit que $K = 0$.
D'autre part, si $M \otimes_R S \to N \otimes_R S$ est surjectif,
alors $Q \otimes_R S$ est nul (par exactitude à droite du produit tensoriel).
Ainsi, $Q \subset Q \otimes_R S$ est également nul.
\end{proof}

\begin{lemma}
\label{algebra-lemma-faithfully-flat-universally-injective}
Soit $R \to S$ un homomorphisme d'anneaux fidèlement plat.
Alors $R \to S$ est universellement injectif comme morphisme de $R$-modules.
En particulier, $R \cap IS = I$ pour tout idéal $I \subset R$.
\end{lemma}

\begin{proof}
Soit $N$ un $R$-module. Nous devons montrer que $N \to N \otimes_R S$ est
injectif. Comme $S$ est fidèlement plat comme $R$-module, il suffit de le démontrer
après tensorisation par $S$. Il suffit donc de montrer que
$N \otimes_R S \to N \otimes_R S \otimes_R S$,
$n \otimes s \mapsto n \otimes 1 \otimes s$ est injectif. C'est bien le cas
car il existe une rétraction, à savoir
$n \otimes s \otimes s' \mapsto n \otimes ss'$.
\end{proof}




\section{Descente pour les modules projectifs de type fini}
\label{algebra-section-finite-projective}

\noindent
Dans cette section, nous donnons une démonstration élémentaire du fait que la propriété
d'être un module projectif {\it de type fini} descend le long des homomorphismes
d'anneaux fidèlement plats. La démonstration ne s'applique plus sans l'hypothèse de finitude.
Toutefois, la méthode préfigure celle que nous utiliserons pour démontrer la descente
de la propriété d'être un module projectif {\it dénombrablement engendré} ; voir
les remarques à la fin de cette section.

\begin{lemma}
\label{algebra-lemma-finite-projective-again}
Soit $M$ un $R$-module. Alors $M$ est projectif de type fini si et seulement si $M$ est
de présentation finie et plat.
\end{lemma}

\begin{proof}
Ceci fait partie du
lemme \ref{algebra-lemma-finite-projective}.
À ce stade, nous pouvons toutefois donner une démonstration plus élégante de l'implication
(1) $\Rightarrow$ (2) de ce lemme, comme suit.
Si $M$ est de présentation finie et plat, choisissons une surjection
$R^n \to M$. D'après le
lemme \ref{algebra-lemma-flat-surjective-hom},
appliqué à $P = M$, le morphisme $R^n \to M$ admet une section.
Ainsi, $M$ est un facteur direct d'un module libre, et est donc projectif.
\end{proof}

\noindent
Voici quelques propriétés des modules qui descendent.

\begin{lemma}
\label{algebra-lemma-descend-properties-modules}
Soit $R \to S$ un homomorphisme d'anneaux fidèlement plat.
Soit $M$ un $R$-module. Alors
\begin{enumerate}
\item si le $S$-module $M \otimes_R S$ est de type fini, alors
$M$ est de type fini,
\item si le $S$-module $M \otimes_R S$ est de présentation finie, alors
$M$ est de présentation finie,
\item si le $S$-module $M \otimes_R S$ est plat, alors
$M$ est plat, et
\item ajouter ici d'autres propriétés au besoin.
\end{enumerate}
\end{lemma}

\begin{proof}
Supposons $M \otimes_R S$ de type fini. Soient $y_1, \ldots, y_m$ des générateurs
de $M \otimes_R S$ sur $S$. Écrivons $y_j = \sum x_i \otimes f_i$ pour certains
$x_1, \ldots, x_n \in M$. On voit alors que le morphisme
$\varphi : R^{\oplus n} \to M$
a la propriété que
$\varphi \otimes \text{id}_S : S^{\oplus n} \to M \otimes_R S$
est surjectif. Comme $R \to S$ est fidèlement plat, on en déduit que
$\varphi$ est surjectif et que $M$ est de type fini.

\medskip\noindent
Supposons $M \otimes_R S$ de présentation finie. D'après (1),
$M$ est de type fini. Choisissons une surjection $R^{\oplus n} \to M$ et
notons $K$ son noyau. Comme $R \to S$ est plat, $K \otimes_R S$
est le noyau du changement de base $S^{\oplus n} \to M \otimes_R S$.
Puisque $M \otimes_R S$ est de présentation finie, on en déduit que $K \otimes_R S$
est de type fini. D'après (1), $K$ est donc de type fini,
et $M$ est par conséquent de présentation finie.

\medskip\noindent
L'assertion (3) est le
lemme \ref{algebra-lemma-flatness-descends}.
\end{proof}

\begin{proposition}
\label{algebra-proposition-ffdescent-finite-projectivity}
Soit $R \to S$ un homomorphisme d'anneaux fidèlement plat. Soit $M$ un $R$-module.
Si le $S$-module $M \otimes_R S$ est projectif de type fini, alors $M$ est projectif
de type fini.
\end{proposition}

\begin{proof}
Cela résulte des
lemmes \ref{algebra-lemma-finite-projective-again} et
\ref{algebra-lemma-descend-properties-modules}.
\end{proof}

\noindent
Les quelques sections suivantes visent à supprimer l'hypothèse de finitude en recourant au
d\'evissage pour se ramener au cas dénombrablement engendré. Dans ce dernier
cas, la stratégie consiste à trouver une caractérisation des modules projectifs
dénombrablement engendrés analogue au
lemme \ref{algebra-lemma-finite-projective-again},
puis à démontrer directement que cette caractérisation descend. Pour cela, nous
introduisons la notion de module de Mittag-Leffler et démontrons qu'un module
$M$ dénombrablement engendré est projectif si et seulement s'il est plat et
de Mittag-Leffler (théorème \ref{algebra-theorem-projectivity-characterization}). Lorsque $M$
est de type fini, cet énoncé se réduit au lemme
\ref{algebra-lemma-finite-projective-again} (puisque, d'après
l'exemple \ref{algebra-example-ML} (1),
un module de type fini est de Mittag-Leffler si et seulement s'il est
de présentation finie).

\section{D\'evissage transfini des modules}
\label{algebra-section-transfinite-devissage}

\noindent
Dans cette section, nous introduisons une technique de d\'evissage permettant de décomposer un module
en somme directe. Le résultat principal affirme qu'un module projectif est une somme directe
de modules dénombrablement engendrés (théorème \ref{algebra-theorem-projective-direct-sum}
ci-dessous). Nous suivons \cite{Kaplansky}.


\begin{definition}
\label{algebra-definition-devissage}
Soit $M$ un $R$-module. Un {\it d\'evissage par sommes directes} de $M$ est une famille
de sous-modules $(M_{\alpha})_{\alpha \in S}$, indexée par un ordinal $S$ et
croissante (pour l'inclusion), telle que :
\begin{enumerate}
\item[(0)] $M_0 = 0$ ;
\item[(1)] $M = \bigcup_{\alpha} M_{\alpha}$ ;
\item[(2)] si $\alpha \in S$ est un ordinal limite, alors $M_{\alpha} =
\bigcup_{\beta < \alpha} M_{\beta}$ ;
\item[(3)] si $\alpha + 1 \in S$, alors $M_{\alpha}$ est un facteur direct de
$M_{\alpha + 1}$.
\end{enumerate}
Si de plus
\begin{enumerate}
\item[(4)] $M_{\alpha + 1}/M_{\alpha}$ est dénombrablement engendré pour
$\alpha + 1 \in S$,
\end{enumerate}
alors $(M_{\alpha})_{\alpha \in S}$ est appelé un {\it d\'evissage de Kaplansky}
de $M$.
\end{definition}

\noindent
La terminologie est justifiée par le lemme suivant.

\begin{lemma}
\label{algebra-lemma-direct-sum-devissage}
Soit $M$ un $R$-module. Si $(M_{\alpha})_{\alpha \in S}$ est un d\'evissage
par sommes directes de $M$, alors
$M \cong \bigoplus_{\alpha + 1 \in S} M_{\alpha + 1}/M_{\alpha}$.
\end{lemma}

\begin{proof}
Par la propriété (3) d'un d\'evissage par sommes directes, il existe une inclusion
$M_{\alpha + 1}/M_{\alpha} \to M$ pour tout $\alpha \in S$. Considérons le
morphisme
$$
f : \bigoplus\nolimits_{\alpha + 1\in S} M_{\alpha + 1}/M_{\alpha} \to M
$$
donné par la somme de ces inclusions.
Considérons en outre les restrictions
$$
f_{\beta} :
\bigoplus\nolimits_{\alpha + 1 \leq \beta} M_{\alpha + 1}/M_{\alpha}
\longrightarrow
M
$$
pour $\beta\in S$. Une récurrence transfinie sur $S$ montre que l'image de
$f_{\beta}$ est $M_{\beta}$. Pour $\beta=0$, cela résulte de $(0)$. Si $\beta+1$
est un ordinal successeur et si l'assertion vaut pour $\beta$, elle vaut alors pour
$\beta + 1$ par (3). Si enfin $\beta$ est un ordinal limite et que l'assertion vaut pour
$\alpha < \beta$, elle vaut pour $\beta$ par (2). Ainsi, $f$ est surjectif
par (1).

\medskip\noindent
Une récurrence transfinie sur $S$ montre aussi que les restrictions $f_{\beta}$
sont injectives. Pour $\beta = 0$, c'est vrai. Si $\beta+1$ est un
ordinal successeur et que $f_{\beta}$ est injectif, soit $x$ un élément du noyau ;
écrivons $x = (x_{\alpha + 1})_{\alpha + 1 \leq \beta + 1}$ selon ses
composantes $x_{\alpha + 1} \in M_{\alpha + 1}/M_{\alpha}$. Par la propriété (3) et
le fait que l'image de $f_{\beta}$ est $M_{\beta}$, tant
$(x_{\alpha + 1})_{\alpha + 1 \leq \beta}$ que $x_{\beta + 1}$ s'envoient sur $0$.
Ainsi, $x_{\beta+1} = 0$ et, puisque la restriction
$f_{\beta}$ est injective, on a aussi $x_{\alpha + 1} = 0$
pour tout $\alpha + 1 \leq \beta$. Ainsi $x = 0$ et $f_{\beta+1}$ est injectif.
Si $\beta$ est un ordinal limite, considérons un élément $x$ du noyau. Alors $x$
appartient déjà au domaine de $f_{\alpha}$ pour un certain $\alpha < \beta$.
Ainsi $x = 0$, ce qui achève la récurrence.
On conclut que $f$ est injectif, puisque $f_{\beta}$ l'est pour tout $\beta \in S$.
\end{proof}

\begin{lemma}
\label{algebra-lemma-Kaplansky-devissage}
Soit $M$ un $R$-module. Alors $M$ est une somme directe de $R$-modules
dénombrablement engendrés si et seulement s'il admet un d\'evissage de Kaplansky.
\end{lemma}

\begin{proof}
Le lemme précédent traite le sens « si ». Réciproquement, supposons $M =
\bigoplus_{i \in I} N_i$, où chaque $N_i$ est un $R$-module dénombrablement engendré.
Munissons $I$ d'un bon ordre afin de pouvoir l'identifier à un ordinal. Alors, en posant $M_i =
\bigoplus_{j < i} N_j$, on obtient un d\'evissage de Kaplansky $(M_i)_{i \in I}$ de
$M$.
\end{proof}

\begin{theorem}
\label{algebra-theorem-kaplansky-direct-sum}
Supposons que $M$ soit une somme directe de $R$-modules dénombrablement engendrés. Si $P$ est un
facteur direct de $M$, alors $P$ est lui aussi une somme directe de $R$-modules
dénombrablement engendrés.
\end{theorem}

\begin{proof}
Écrivons $M = P \oplus Q$. Nous allons construire un d\'evissage de Kaplansky
$(M_{\alpha})_{\alpha \in S}$ de $M$ qui, outre les propriétés de définition
(0)--(4), vérifie :
\begin{enumerate}
\item[(5)] Chaque $M_{\alpha}$ est un facteur direct de $M$ ;
\item[(6)] $M_{\alpha} = P_{\alpha} \oplus Q_{\alpha}$, où $P_{\alpha} =P
\cap M_{\alpha}$ et $Q_\alpha = Q \cap M_{\alpha}$.
\end{enumerate}
(Remarque : si les propriétés (0)--(2) sont satisfaites, alors la propriété (3) équivaut en fait à
la propriété (5).)

\medskip\noindent
Pour voir comment ceci implique le théorème, il suffit de montrer que
$(P_{\alpha})_{\alpha \in S}$ forme un d\'evissage de Kaplansky de $P$. Les propriétés
(0), (1) et (2) sont claires. D'après (5) et (6) pour $(M_{\alpha})$, chaque
$P_{\alpha}$ est un facteur direct de $M$. Comme $P_{\alpha} \subset P_{\alpha +
1}$, il s'ensuit que $P_{\alpha}$ est un facteur direct de $P_{\alpha + 1}$ ;
ainsi, (3) vaut pour $(P_{\alpha})$. Pour (4), remarquons que
$$
M_{\alpha + 1}/M_{\alpha} \cong P_{\alpha + 1}/P_{\alpha} \oplus
Q_{\alpha + 1}/Q_{\alpha},
$$
de sorte que $P_{\alpha + 1}/P_{\alpha}$ est dénombrablement engendré, puisque
$M_{\alpha + 1}/M_{\alpha}$ l'est.

\medskip\noindent
Il reste à construire les $M_{\alpha}$. Écrivons $M = \bigoplus_{i \in I} N_i$,
où chaque $N_i$ est un $R$-module dénombrablement engendré. Choisissons un bon ordre
sur $I$. Par récurrence transfinie, nous allons définir une famille croissante
de sous-modules $M_{\alpha}$ de $M$, un pour chaque ordinal $\alpha$, telle que
$M_{\alpha}$ soit la somme directe d'une certaine sous-famille des $N_i$.

\medskip\noindent
Pour $\alpha = 0$, posons $M_{0} = 0$. Si $\alpha$ est un ordinal limite et que
$M_{\beta}$ a été défini pour tout $\beta < \alpha$, posons $M_{\alpha}
= \bigcup_{\beta < \alpha} M_{\beta}$. Puisque chaque $M_{\beta}$, pour $\beta <
\alpha$, est la somme directe d'une sous-famille des $N_i$, il en va de même de
$M_{\alpha}$. Si $\alpha + 1$ est un ordinal successeur et que $M_{\alpha}$ a été
défini, définissons $M_{\alpha + 1}$ comme suit. Si $M_{\alpha} = M$, posons
$M_{\alpha + 1} = M$. Sinon, choisissons le plus petit $j \in I$ tel que $N_j$ ne soit
pas contenu dans $M_{\alpha}$. Nous allons construire une matrice infinie $(x_{mn}),
m, n = 1, 2, 3, \ldots$ telle que :
\begin{enumerate}
\item $N_j$ est contenu dans le sous-module de $M$ engendré par les entrées
$x_{mn}$ ;
\item si l'on écrit une entrée quelconque $x_{k\ell}$ selon ses composantes dans $P$ et
$Q$, $x_{k\ell} = y_{k\ell} + z_{k\ell}$, alors la matrice $(x_{mn})$
contient une famille de générateurs de chaque $N_i$ pour lequel $y_{k\ell}$ ou
$z_{k\ell}$ possède une composante non nulle.
\end{enumerate}
Nous définissons alors $M_{\alpha + 1}$ comme le sous-module de $M$ engendré par
$M_{\alpha}$ et tous les $x_{mn}$ ; d'après la propriété (2) de la matrice $(x_{mn})$,
$M_{\alpha + 1}$ sera la somme directe d'une certaine sous-famille des $N_i$.
Pour construire la matrice $(x_{mn})$, soit $x_{11}, x_{12}, x_{13}, \ldots$
une famille dénombrable de générateurs de $N_j$. Si ensuite
$x_{11} = y_{11} + z_{11}$ est la décomposition selon les composantes dans $P$ et
$Q$, soit $x_{21}, x_{22}, x_{23}, \ldots$ une famille dénombrable
de générateurs de la somme des $N_i$ pour lesquels $y_{11}$ ou $z_{11}$ possède une
composante non nulle. Répétons ce procédé avec $x_{12}$ pour obtenir $x_{31},
x_{32}, \ldots$, la troisième ligne de notre matrice. Répétons-le avec $x_{21}$ pour obtenir la
quatrième ligne, avec $x_{13}$ pour obtenir la cinquième, et ainsi de suite, en parcourant
successivement les antidiagonales comme indiqué ci-dessous :
$$
\left(
\vcenter{
\xymatrix@R=2mm@C=2mm{
x_{11} & x_{12} \ar[dl] & x_{13} \ar[dl] & x_{14} \ar[dl] & \ldots  \\
x_{21} & x_{22} \ar[dl] & x_{23} \ar[dl] & \ldots  \\
x_{31} & x_{32} \ar[dl] & \ldots \\
x_{41} & \ldots \\
\ldots
}
}
\right).
$$

\medskip\noindent
Une récurrence transfinie sur $I$ (utilisant le fait que nous avons construit
$M_{\alpha + 1}$
de façon qu'il contienne $N_j$ pour le plus petit $j$ tel que $N_j$ ne soit pas contenu dans
$M_{\alpha}$) montre que, pour tout $i \in I$, $N_i$ est contenu dans un certain
$M_{\alpha}$. Il existe donc un ordinal $S$ assez grand vérifiant : pour
tout $i \in I$, il existe $\alpha \in S$ tel que $N_i$ soit contenu dans
$M_{\alpha}$. Cela signifie que $(M_{\alpha})_{\alpha \in S}$ vérifie la propriété (1)
d'un d\'evissage de Kaplansky de $M$. La famille $(M_{\alpha})_{\alpha \in S}$
vérifie en outre les autres propriétés de définition, ainsi que (5) et (6) ci-dessus :
les propriétés (0), (2), (4) et (6) sont claires par construction ; la propriété (5) est
vraie parce que chaque $M_{\alpha}$ est par construction une somme directe de certains $N_i$ ;
et (3) résulte de (5) et du fait que $M_{\alpha} \subset M_{\alpha + 1}$.
\end{proof}

\noindent
On obtient comme corollaire le résultat sur les modules projectifs annoncé au début
de la section.

\begin{theorem}
\label{algebra-theorem-projective-direct-sum}
\begin{slogan}
Tout module projectif est une somme directe de modules projectifs
dénombrablement engendrés.
\end{slogan}
Si $P$ est un $R$-module projectif, alors $P$ est une somme directe de $R$-modules projectifs
dénombrablement engendrés.
\end{theorem}

\begin{proof}
Un module est projectif si et seulement s'il est un facteur direct d'un module libre ;
le résultat découle donc du théorème \ref{algebra-theorem-kaplansky-direct-sum}.
\end{proof}



\section{Modules projectifs sur un anneau local}
\label{algebra-section-projective-local-ring}

\noindent
Dans cette section, nous démontrons un très joli résultat :
un module projectif $M$ sur un anneau local est libre
(théorème \ref{algebra-theorem-projective-free-over-local-ring} ci-dessous).
Notons que, sous l'hypothèse supplémentaire que $M$ est de type fini, ce résultat est le
lemme \ref{algebra-lemma-finite-flat-local}.
En général, on a :

\begin{lemma}
\label{algebra-lemma-projective-free}
Soit $R$ un anneau. Alors tout $R$-module projectif est libre si et seulement si
tout $R$-module projectif dénombrablement engendré est libre.
\end{lemma}

\begin{proof}
Cela résulte immédiatement du
théorème \ref{algebra-theorem-projective-direct-sum}.
\end{proof}

\noindent
Voici un critère pour qu'un module dénombrablement engendré soit libre.

\begin{lemma}
\label{algebra-lemma-freeness-criteria}
Soit $M$ un $R$-module dénombrablement engendré possédant la propriété suivante :
si $M = N \oplus N'$, où $N'$ est un $R$-module libre de type fini, alors
tout élément de $N$ appartient à un facteur direct libre de $N$.
Alors $M$ est libre.
\end{lemma}

\begin{proof}
Soient $x_1, x_2, \ldots$ une famille dénombrable de générateurs de $M$.
Construisons par récurrence des facteurs directs libres de type fini
$F_1, F_2, \ldots$ de $M$ tels que, pour tout $n$,
$F_1 \oplus \ldots \oplus F_n$ soit un facteur direct de $M$
contenant $x_1, \ldots, x_n$. Plus précisément, étant donnés $F_1, \ldots, F_n$
ayant les propriétés voulues, écrivons
$$
M = F_1 \oplus \ldots \oplus F_n \oplus N
$$
et soit $x \in N$ l'image de $x_{n + 1}$.
L'hypothèse de l'énoncé du lemme fournit alors un facteur direct libre $F_{n + 1} \subset N$
contenant $x$.
On peut bien sûr remplacer $F_{n + 1}$ par un facteur direct libre de type fini
de $F_{n + 1}$, ce qui achève l'étape de récurrence.
Alors $M = \bigoplus_{i = 1}^{\infty} F_i$ est libre.
\end{proof}

\begin{lemma}
\label{algebra-lemma-projective-freeness-criteria}
Soit $P$ un module projectif sur un anneau local $R$. Alors tout élément de $P$
appartient à un facteur direct libre de $P$.
\end{lemma}

\begin{proof}
Comme $P$ est projectif, il est facteur direct d'un certain $R$-module libre $F$ ; écrivons
$F = P \oplus Q$. Soit $x \in P$ l'élément dont nous voulons montrer qu'il
appartient à un facteur direct libre de $P$. Soit $B$ une base de $F$ telle que
le nombre d'éléments de base nécessaires pour exprimer $x$ soit minimal ; écrivons $x
= \sum_{i=1}^n a_i e_i$ pour certains $e_i \in B$ et $a_i \in R$. Alors aucun $a_j$
ne peut s'écrire comme combinaison linéaire des autres $a_i$ ; en effet, si $a_j =
\sum_{i \neq  j} a_i b_i$ pour certains $b_i \in R$, le remplacement de $e_i$ par $e_i +
b_ie_j$ pour $i \neq j$, les autres éléments de $B$ restant inchangés, donne
une nouvelle base de $F$ dans laquelle $x$ possède une expression plus courte.

\medskip\noindent
Soit $e_i = y_i + z_i, y_i \in P, z_i \in Q$ la décomposition de $e_i$ selon
ses composantes dans $P$ et $Q$. Écrivons $y_i = \sum_{j=1}^{n} b_{ij} e_j + t_i$,
où $t_i$ est une combinaison linéaire d'éléments de $B$ autres que $e_1, \ldots,
e_n$. Pour achever la démonstration, il suffit de montrer que la matrice $(b_{ij})$ est
inversible. Alors, en effet, le morphisme $F \to F$ envoyant $e_i \mapsto y_i$ pour
$i=1, \ldots, n$ et fixant $B \setminus \{e_1, \ldots, e_n\}$ est un
isomorphisme ;
ainsi, $y_1, \ldots, y_n$ et $B \setminus \{e_1, \ldots, e_n\}$
forment une base de $F$. Le sous-module $N$ engendré par $y_1, \ldots, y_n$
est alors un sous-module libre de $P$ ; $N$ est un facteur direct de $P$, car $N \subset P$
et $N$ comme $P$ sont des facteurs directs de $F$ ; enfin $x \in N$, puisque $x \in P$
implique $x = \sum_{i=1}^n a_i e_i = \sum_{i=1}^n a_i y_i$.

\medskip\noindent
Démontrons maintenant que $(b_{ij})$ est inversible. En substituant $y_i = \sum_{j=1}^{n}
b_{ij} e_j + t_i$ dans $\sum_{i=1}^n a_i e_i = \sum_{i=1}^n a_i y_i$ et
en identifiant les coefficients de $e_j$, on obtient $a_j = \sum_{i=1}^n a_i b_{ij}$. Or,
comme observé ci-dessus, le choix de $B$ garantit qu'aucun $a_j$ ne peut s'écrire comme
combinaison linéaire des autres $a_i$. Ainsi, $b_{ij}$ est non inversible pour $i \neq
j$, et $1-b_{ii}$ est non inversible --- donc, en particulier, $b_{ii}$ est inversible --- pour
tout $i$. Une matrice sur un anneau local dont les coefficients diagonaux sont inversibles et
les autres non inversibles est inversible, car son déterminant est inversible.
\end{proof}

\begin{theorem}
\label{algebra-theorem-projective-free-over-local-ring}
\begin{slogan}
Les modules projectifs sur les anneaux locaux sont libres.
\end{slogan}
Si $P$ est un module projectif sur un anneau local $R$, alors $P$ est libre.
\end{theorem}

\begin{proof}
Cela résulte des lemmes \ref{algebra-lemma-projective-free}, \ref{algebra-lemma-freeness-criteria}
et \ref{algebra-lemma-projective-freeness-criteria}.
\end{proof}



\section{Systèmes de Mittag-Leffler}
\label{algebra-section-mittag-leffler}

\noindent
Cette section a pour but de définir les systèmes de Mittag-Leffler
et d'expliquer l'utilité de cette notion.

\medskip\noindent
Dans la suite, $I$ sera un ensemble préordonné filtrant, voir
Catégories, définition \ref{categories-definition-directed-set}.
Soit $(A_i, \varphi_{ji}: A_j \to A_i)$ un système projectif
d'ensembles ou de modules indexé par $I$, voir
Catégories, définition \ref{categories-definition-directed-system}.
Il s'agit d'un système projectif filtrant, puisque nous avons supposé $I$ filtrant
(Catégories, définition \ref{categories-definition-directed-system}).
Pour tout $i \in I$, les images $\varphi_{ji}(A_j) \subset A_i$, pour $j \geq i$,
forment une famille filtrante décroissante de parties (ou de sous-modules) de $A_i$. Posons
$A'_i = \bigcap_{j \geq i} \varphi_{ji}(A_j)$.
Alors $\varphi_{ji}(A'_j) \subset A'_i$ pour $j \geq i$ ; par restriction,
on obtient donc un système projectif filtrant $(A'_i, \varphi_{ji}|_{A'_j})$.
La construction de la limite d'un système projectif dans la catégorie
des ensembles ou des modules donne $\lim A_i = \lim A'_i$. La condition de Mittag-Leffler
sur $(A_i, \varphi_{ji})$ exige que $A'_i$ soit égal à
$\varphi_{ji}(A_j)$ pour un certain $j \geq i$ (et donc égal à
$\varphi_{ki}(A_k)$ pour tout $k \geq j$) :

\begin{definition}
\label{algebra-definition-ML-system}
Soit $(A_i, \varphi_{ji})$ un système projectif filtrant d'ensembles sur $I$. On
dit que $(A_i, \varphi_{ji})$ est {\it de Mittag-Leffler} si, pour
tout $i \in I$, la famille $\varphi_{ji}(A_j) \subset A_i$, pour
$j \geq i$, se stabilise. Explicitement, cela signifie que, pour tout $i \in I$, il
existe $j \geq i$ tel que, pour $k \geq j$, on ait $\varphi_{ki}(A_k) =
\varphi_{ji}( A_j)$. Si $(A_i, \varphi_{ji})$ est un système projectif filtrant
de modules sur un anneau $R$, on dit qu'il est de Mittag-Leffler si le système
projectif d'ensembles sous-jacent est de Mittag-Leffler.
\end{definition}

\begin{example}
\label{algebra-example-ML-surjective-maps}
Si $(A_i, \varphi_{ji})$ est un système projectif filtrant d'ensembles ou de modules et si
les morphismes $\varphi_{ji}$ sont surjectifs, alors le système est clairement
de Mittag-Leffler. Réciproquement, supposons $(A_i, \varphi_{ji})$ de Mittag-Leffler.
Soit $A'_i \subset A_i$ l'image stable des $\varphi_{ji}(A_j)$ pour $j \geq
i$. Alors $\varphi_{ji}|_{A'_j}: A'_j \to A'_i$ est surjectif pour
$j \geq i$ et $\lim A_i = \lim A'_i$. Ainsi, la limite du système de Mittag-Leffler
$(A_i, \varphi_{ji})$ peut aussi s'écrire comme la limite d'un système
projectif filtrant sur $I$ à morphismes de transition surjectifs.
\end{example}

\begin{lemma}
\label{algebra-lemma-ML-limit-nonempty}
Soit $(A_i, \varphi_{ji})$ un système projectif filtrant sur $I$. Supposons $I$
dénombrable. Si $(A_i, \varphi_{ji})$ est de Mittag-Leffler et que les $A_i$ sont
non vides, alors $\lim A_i$ est non vide.
\end{lemma}

\begin{proof}
Soit $i_1, i_2, i_3, \ldots$ une énumération des éléments de $I$. Définissons
par récurrence une suite d'éléments $j_n \in I$, pour $n = 1, 2, 3, \ldots$, par les
conditions suivantes : $j_1 = i_1$, et $j_n \geq i_n$ et $j_n \geq j_m$ pour $m < n$.
 La suite $j_n$ est alors croissante et forme une partie cofinale de $I$.
Nous pouvons donc supposer $I =\{1, 2, 3, \ldots \}$. L'exemple
\ref{algebra-example-ML-surjective-maps} nous ramène alors à montrer que la limite d'un
système projectif d'ensembles non vides, à morphismes surjectifs et indexé par les entiers
strictement positifs, est non vide. Cela résulte de l'axiome du choix.
\end{proof}

\noindent
La condition de Mittag-Leffler sera importante pour nous en raison de la propriété
d'exactitude suivante.

\begin{lemma}
\label{algebra-lemma-ML-exact-sequence}
Soit
$$
0 \to A_i \xrightarrow{f_i} B_i \xrightarrow{g_i} C_i \to 0
$$
une suite exacte de systèmes projectifs filtrants de groupes abéliens sur $I$.
Supposons $I$ dénombrable. Si $(A_i)$ est de Mittag-Leffler, alors
$$
0 \to \lim A_i \to \lim B_i \to \lim C_i\to 0
$$
est exacte.
\end{lemma}

\begin{proof}
Le passage à la limite des systèmes projectifs filtrants est exact à gauche ; il suffit donc de
démontrer la surjectivité de $\lim B_i \to \lim C_i$. Soit donc $(c_i) \in \lim
C_i$. Pour tout $i \in I$, posons $E_i = g_i^{-1}(c_i)$, qui est non vide puisque
$g_i: B_i \to C_i$ est surjectif. Le système de morphismes $\varphi_{ji}: B_j
\to B_i$ du système $(B_i)$ induit par restriction des morphismes $E_j \to E_i$ qui
font de $(E_i)$ un système projectif d'ensembles non vides. Il suffit de montrer
que $(E_i)$ est de Mittag-Leffler. Alors le lemme \ref{algebra-lemma-ML-limit-nonempty}
montrera que $\lim E_i$ est non vide, et le choix d'un élément de $\lim E_i$
fournira un élément de $\lim B_i$ s'envoyant sur $(c_i)$.

\medskip\noindent
Grâce à l'injection $f_i: A_i \to B_i$, nous considérerons $A_i$ comme une partie de
$B_i$. Puisque $(A_i)$ est de Mittag-Leffler, pour $i \in I$, il existe $j \geq
i$ tel que $\varphi_{ki}(A_k) = \varphi_{ji}(A_j)$ pour $k \geq j$. Nous affirmons
qu'on a aussi $\varphi_{ki}(E_k) = \varphi_{ji}(E_j)$ pour $k \geq j$. On a toujours
$\varphi_{ki}(E_k) \subset \varphi_{ji}(E_j)$ pour $k \geq j$. Pour l'inclusion
réciproque, soit $e_j \in E_j$ ; nous devons trouver $x_k \in E_k$ tel que
$\varphi_{ki}(x_k) = \varphi_{ji}(e_j)$. Soit $e'_k \in E_k$ un élément quelconque,
et posons $e'_j = \varphi_{kj}(e'_k)$. Alors $g_j(e_j - e'_j) = c_j - c_j = 0$ ;
ainsi $e_j - e'_j = a_j \in A_j$. Comme $\varphi_{ki}(A_k) =
\varphi_{ji}(A_j)$, il existe $a_k \in A_k$ tel que $\varphi_{ki}(a_k) =
\varphi_{ji}(a_j)$. Ainsi
$$
\varphi_{ki}(e'_k + a_k) = \varphi_{ji}(e'_j) + \varphi_{ji}(a_j) =
\varphi_{ji}(e_j),
$$
on peut donc prendre $x_k = e'_k + a_k$.
\end{proof}




\section{Systèmes projectifs}
\label{algebra-section-inverse-systems}

\noindent
Dans de nombreux articles (et dans cette section), le terme {\it système projectif} est
employé pour désigner un système projectif sur l'ensemble ordonné
$(\mathbf{N}, \geq)$. Nous étudions brièvement de tels systèmes dans cette section.
Ce matériel sera traité plus généralement dans
Homologie, section \ref{homology-section-inverse-systems}.
Supposons donnés un anneau $R$ et une suite de $R$-modules
$$
M_1 \xleftarrow{\varphi_2} M_2 \xleftarrow{\varphi_3} M_3 \leftarrow \ldots
$$
avec les morphismes indiqués. En composant les morphismes successifs, on obtient des morphismes
$\varphi_{ii'} : M_i \to M_{i'}$ pour $i \geq i'$ tels qu'en outre
$\varphi_{ii''} = \varphi_{i'i''} \circ \varphi_{i i'}$ dès que
$i \geq i' \geq i''$. Réciproquement, étant donné le système des morphismes $\varphi_{ii'}$,
on peut poser $\varphi_i = \varphi_{i(i-1)}$ et retrouver les morphismes représentés
ci-dessus. Dans ce cas,
$$
\lim M_i
=
\{(x_i) \in \prod M_i \mid \varphi_i(x_i) = x_{i - 1}, \ i = 2, 3, \ldots\}
$$
voir
Catégories, section \ref{categories-section-limit-sets}.
Comme expliqué dans
Homologie, section \ref{homology-section-inverse-systems},
il s'agit bien d'une limite dans la catégorie des $R$-modules, au sens défini dans
Catégories, section \ref{categories-section-limits}.

\begin{lemma}
\label{algebra-lemma-Mittag-Leffler}
Soit $R$ un anneau.
Soient $0 \to K_i \to L_i \to M_i \to 0$ des suites exactes courtes de
$R$-modules, $i \geq 1$, s'insérant dans des morphismes de suites exactes courtes
$$
\xymatrix{
0 \ar[r] &
K_i \ar[r] &
L_i \ar[r] &
M_i \ar[r] &
0 \\
0 \ar[r] &
K_{i + 1} \ar[r] \ar[u] &
L_{i + 1} \ar[r] \ar[u] &
M_{i + 1} \ar[r] \ar[u] &
0}
$$
Si, pour tout $i$, il existe $c = c(i) \geq i$ tel que
$\Im(K_c \to K_i) = \Im(K_j \to K_i)$
pour tout $j \geq c$, alors la suite
$$
0 \to \lim K_i \to \lim L_i \to \lim M_i \to 0
$$
est exacte.
\end{lemma}

\begin{proof}
C'est un cas particulier du lemme plus général
\ref{algebra-lemma-ML-exact-sequence}.
\end{proof}



\section{Modules de Mittag-Leffler}
\label{algebra-section-mittag-leffler-modules}

\noindent
Un module de Mittag-Leffler est (très grossièrement) un module qui peut s'écrire
comme une colimite filtrante dont le dual est un système de Mittag-Leffler. Afin de
donner une définition précise, il nous faut effectuer quelques préparatifs.

\begin{definition}
\label{algebra-definition-ML-inductive-system}
Soit $(M_i, f_{ij})$ un système inductif de $R$-modules. On dit que
$(M_i, f_{ij})$ est un {\it système inductif de Mittag-Leffler} si chaque
$M_i$ est un $R$-module de présentation finie et si, pour tout $R$-module $N$,
le système projectif
$$
(\Hom_R(M_i, N), \Hom_R(f_{ij}, N))
$$
est de Mittag-Leffler.
\end{definition}

\noindent
Nous allons caractériser les $R$-modules qui sont des colimites de
systèmes inductifs de Mittag-Leffler.

\begin{definition}
\label{algebra-definition-domination}
Soient $f: M \to N$ et $g: M \to M'$ des morphismes de $R$-modules.
On dit que $g$ {\it domine} $f$ si, pour tout $R$-module $Q$, on a $\Ker(f
\otimes_R \text{id}_Q) \subset \Ker(g \otimes_R \text{id}_Q)$.
\end{definition}

\noindent
Il suffit de vérifier cette condition sur les modules de présentation finie.

\begin{lemma}
\label{algebra-lemma-domination-fp}
Soient $f: M \to N$ et $g: M \to M'$ des morphismes de $R$-modules.
Alors $g$ domine $f$ si et seulement si, pour tout $R$-module de présentation finie
$Q$, on a $\Ker(f \otimes_R \text{id}_Q) \subset \Ker(g \otimes_R
\text{id}_Q)$.
\end{lemma}

\begin{proof}
Supposons $\Ker(f \otimes_R \text{id}_Q) \subset \Ker(g \otimes_R
\text{id}_Q)$ pour tout module de présentation finie $Q$. Si $Q$ est un
module quelconque, écrivons $Q = \colim_{i \in I} Q_i$ comme colimite d'un
système
inductif de modules de présentation finie $Q_i$. Alors $\Ker(f \otimes_R
\text{id}_{Q_i}) \subset \Ker(g \otimes_R \text{id}_{Q_i})$ pour
tout $i$. Puisque le passage aux colimites filtrantes est exact et commute au produit
tensoriel, il s'ensuit que $\Ker(f \otimes_R \text{id}_Q) \subset \Ker(g
\otimes_R \text{id}_Q)$.
\end{proof}

\begin{lemma}
\label{algebra-lemma-domination-universally-injective}
Soient $f : M \to N$ et $g : M \to M'$ des morphismes de $R$-modules.
Considérons la somme amalgamée de $f$ et $g$,
$$
\xymatrix{
M  \ar[r]_f \ar[d]_g & N \ar[d]^{g'} \\
M' \ar[r]^{f'} & N'
}
$$
Alors $g$ domine $f$ si et seulement si $f'$ est universellement injectif.
\end{lemma}

\begin{proof}
Rappelons que $N'$ est le quotient de $M' \oplus N$ par le sous-module formé des éléments
$(g(x), -f(x))$ pour $x \in M$.
La construction de $N'$ fournit une suite exacte courte
$$
0 \to \Ker(f) \cap \Ker(g) \to \Ker(f) \to \Ker(f')
\to 0.
$$
Comme le produit tensoriel commute aux sommes amalgamées, on a une telle suite exacte
courte
$$
0 \to \Ker(f \otimes \text{id}_Q ) \cap \Ker(g \otimes
\text{id}_Q) \to \Ker(f \otimes \text{id}_Q)
\to \Ker(f' \otimes \text{id}_Q) \to 0
$$
pour tout $R$-module $Q$. Ainsi, $f'$ est universellement injectif si et seulement si
$\Ker(f \otimes \text{id}_Q ) \subset \Ker(g \otimes
\text{id}_Q)$ pour tout $Q$, si et seulement si $g$ domine $f$.
\end{proof}

\noindent
La définition de la domination donnée ci-dessus est parfois reliée à la notion usuelle
de domination de morphismes, comme le montre le lemme suivant.

\begin{lemma}
\label{algebra-lemma-domination}
Soient $f: M \to N$ et $g: M \to M'$ des morphismes de $R$-modules.
Supposons $\Coker(f)$ de présentation finie. Alors $g$ domine $f$ si
et
seulement si $g$ se factorise par $f$, c'est-à-dire\ s'il existe un morphisme de modules $h: N
\to M'$ tel que $g = h \circ f$.
\end{lemma}

\begin{proof}
Considérons la somme amalgamée de $f$ et $g$ comme dans l'énoncé du
lemme \ref{algebra-lemma-domination-universally-injective}.
La construction de la somme amalgamée montre que
$\Coker(f') = \Coker(f)$ ; ainsi, $\Coker(f')$ est de
présentation finie. D'après le
lemme \ref{algebra-lemma-universally-exact-split}, $f'$ est universellement injectif si et
seulement si
$$
0 \to M' \xrightarrow{f'} N' \to \Coker(f') \to 0
$$
est scindée. C'est le cas si et seulement s'il existe un morphisme $h' : N' \to M'$
tel que $h' \circ f' = \text{id}_{M'}$. Par la propriété
universelle de la somme amalgamée, l'existence d'un tel $h'$ équivaut à ce que $g$
se factorise par $f$.
\end{proof}


\begin{proposition}
\label{algebra-proposition-ML-characterization}
Soit $M$ un $R$-module. Soit $(M_i, f_{ij})$ un système inductif de $R$-modules de
présentation finie, indexé par $I$, tel que $M = \colim M_i$. Soit
$f_i:
M_i \to M$ le morphisme canonique. Les assertions suivantes sont équivalentes :
\begin{enumerate}
\item Pour tout $R$-module de présentation finie $P$ et tout morphisme de modules $f: P
\to M$, il existe un $R$-module de présentation finie $Q$ et un morphisme de modules
$g: P \to Q$ tels que $g$ et $f$ se dominent mutuellement, c'est-à-dire
$\Ker(f \otimes_R \text{id}_N) = \Ker(g \otimes_R \text{id}_N)$
pour tout $R$-module $N$.
\item Pour tout $i \in I$, il existe $j \geq i$ tel que $f_{ij}: M_i
\to M_j$ domine $f_i: M_i \to M$.
\item Pour tout $i \in I$, il existe $j \geq i$ tel que $f_{ij}: M_i
\to M_j$ se factorise par $f_{ik}: M_i \to M_k$ pour tout $k \geq
i$.
\item Pour tout $R$-module $N$, le système projectif
$(\Hom_R(M_i, N), \Hom_R(f_{ij}, N))$ est de Mittag-Leffler.
\item Pour $N = \prod_{s \in I} M_s$, le système projectif
$(\Hom_R(M_i, N), \Hom_R(f_{ij}, N))$ est de Mittag-Leffler.
\end{enumerate}
\end{proposition}

\begin{proof}
Démontrons d'abord l'équivalence de (1) et (2). Supposons (1) satisfaite et soit $i
\in I$. Au morphisme $f_i: M_i \to M$, on peut associer un $g:
M_i \to Q$ comme en (1). Comme $M_i$ et $Q$ sont de présentation finie,
$\Coker(g)$ l'est aussi. D'après le lemme \ref{algebra-lemma-domination}, $f_i : M_i
\to M$ se factorise par $g: M_i \to Q$ ; écrivons $f_i = h \circ g$
pour un certain $h: Q \to M$. Comme $Q$ est de présentation finie, $h$
se factorise par $M_j \to M$ pour un certain $j \geq i$ ; écrivons $h = f_j \circ h'$
pour un certain $h': Q \to M_j$. On obtient au total un diagramme commutatif
$$
\xymatrix{
  &  M  & \\
M_i \ar[dr]_g \ar[ur]^{f_i} \ar[rr]^{f_{ij}} &
& M_j \ar[ul]_{f_j} \\
  & Q  \ar[ur]_{h'} &
}
$$
Ainsi, $f_{ij}$ domine $g$. Or $g$ domine $f_i$, donc $f_{ij}$ domine
$f_i$.

\medskip\noindent
Réciproquement, supposons (2) satisfaite. Soit $P$ de présentation finie et soit $f: P
\to M$ un morphisme de modules. Alors $f$ se factorise par $f_i: M_i \to M$
pour un certain $i \in I$ ; écrivons $f = f_i \circ g'$ pour un certain $g': P \to M_i$.
Choisissons, grâce à (2), un $j \geq i$ tel que $f_{ij}$ domine $f_i$. On a un
diagramme commutatif
$$
\xymatrix{
P \ar[d]_{g'} \ar[r]^{f}           & M  \\
M_i \ar[ur]^{f_i} \ar[r]_{f_{ij}} & M_j \ar[u]_{f_j}
}
$$
Le diagramme et le fait que $f_{ij}$ domine $f_i$ montrent que $f$
et $f_{ij} \circ g'$ se dominent mutuellement. Ainsi, on peut prendre $g = f_{ij} \circ g' :
P \to M_j$.

\medskip\noindent
Montrons ensuite que (2) équivaut à (3). Soit $i \in I$. Il est toujours vrai que
$f_i$ domine $f_{ik}$ pour $k \geq i$, puisque $f_i$ se factorise par
$f_{ik}$. Si (2) est satisfaite, choisissons $j \geq i$ tel que $f_{ij}$ domine
$f_i$. La domination étant transitive, $f_{ij}$ domine alors
$f_{ik}$ pour $k \geq i$. Tous les $M_i$ sont de présentation finie, donc
$\Coker(f_{ik})$ est de présentation finie pour $k \geq i$. D'après le lemme
\ref{algebra-lemma-domination}, $f_{ij}$ se factorise par $f_{ik}$ pour tout $k \geq i$.
Ainsi, (2) implique (3). D'autre part, si (3) est satisfaite, alors, pour tout $R$-module
$N$, $f_{ij} \otimes_R \text{id}_N$ se factorise par $f_{ik}
\otimes_R \text{id}_N$ pour $k \geq i$. Ainsi, $\Ker(f_{ik} \otimes_R
\text{id}_N) \subset \Ker(f_{ij} \otimes_R \text{id}_N)$ pour $k
\geq i$. Or $\Ker(f_i \otimes_R \text{id}_N: M_i \otimes_R N
\to M \otimes_R N)$ est la réunion des $\Ker(f_{ik} \otimes_R
\text{id}_N)$ pour $k \geq i$. Ainsi, $\Ker(f_i \otimes_R
\text{id}_N) \subset \Ker(f_{ij} \otimes_R \text{id}_N)$ pour
tout $R$-module $N$, ce qui signifie par définition que $f_{ij}$ domine $f_i$.

\medskip\noindent
Il est immédiat que (3) implique (4), qui implique (5). Montrons que (5) implique (3). Soit
$N = \prod_{s \in I} M_s$. Si (5) est satisfaite, étant donné $i \in I$, choisissons $j \geq i$
tel que
$$
\Im( \Hom(M_j, N) \to  \Hom(M_i, N)) =
\Im( \Hom(M_k, N) \to  \Hom(M_i, N))
$$
pour tout $k \geq j$. En faisant sortir le produit sur $s \in I$ des
$\Hom$
et en examinant les morphismes sur chaque composante du produit, ceci signifie
$$
\Im( \Hom(M_j, M_s) \to  \Hom(M_i, M_s)) =
\Im( \Hom(M_k, M_s) \to   \Hom(M_i, M_s))
$$
pour tous $k \geq j$ et $s \in I$. En prenant $s = j$, on obtient
$$
\Im( \Hom(M_j, M_j) \to  \Hom(M_i, M_j)) =
\Im( \Hom(M_k, M_j) \to   \Hom(M_i, M_j))
$$
pour tout $k \geq j$. Puisque $f_{ij}$ est l'image de
$\text{id} \in  \Hom(M_j, M_j)$ par
$\Hom(M_j, M_j) \to  \Hom(M_i, M_j)$,
cela montre que, pour tout $k \geq j$, il existe $h \in \Hom(M_k, M_j)$
tel que $f_{ij} = h \circ f_{ik}$. Si $j \geq k$, on peut prendre
$h = f_{kj}$. Ainsi, (3) est satisfaite.
\end{proof}

\begin{definition}
\label{algebra-definition-mittag-leffler-module}
Soit $M$ un $R$-module. On dit que $M$ est {\it de Mittag-Leffler} si les
conditions équivalentes de la
proposition \ref{algebra-proposition-ML-characterization}
sont satisfaites.
\end{definition}

\noindent
En particulier, tout module de présentation finie est de Mittag-Leffler.

\begin{remark}
\label{algebra-remark-flat-ML}
Soit $M$ un $R$-module plat. D'après le théorème de Lazard
(théorème \ref{algebra-theorem-lazard}),
on peut écrire $M = \colim M_i$ comme la colimite d'un
système inductif $(M_i, f_{ij})$ où les $M_i$ sont des $R$-modules libres
de type fini. Pour que $M$ soit de Mittag-Leffler, il suffit que le système projectif
des duaux $(\Hom_R(M_i, R), \Hom_R(f_{ij}, R))$ soit de
Mittag-Leffler. Cela résulte du critère (4) de la
proposition \ref{algebra-proposition-ML-characterization}
et du fait que, pour un $R$-module libre de type fini $F$,
il existe un isomorphisme fonctoriel
$\Hom_R(F, R) \otimes_R N \cong \Hom_R(F, N)$
pour tout $R$-module $N$.
\end{remark}

\begin{lemma}
\label{algebra-lemma-tensor-ML-modules}
Si $R$ est un anneau et si $M$, $N$ sont des $R$-modules de Mittag-Leffler,
alors $M \otimes_R N$ est un module de Mittag-Leffler.
\end{lemma}

\begin{proof}
Écrivons $M = \colim_{i \in I} M_i$ et $N = \colim_{j \in J} N_j$
comme des colimites filtrantes de $R$-modules de présentation finie.
Notons $f_{ii'} : M_i \to M_{i'}$ et $g_{jj'} : N_j \to N_{j'}$ les
morphismes de transition. Alors $M_i \otimes_R N_j$ est un
$R$-module de présentation finie (voir le
lemme \ref{algebra-lemma-tensor-finiteness}),
et $M \otimes_R N = \colim_{(i, j) \in I \times J} M_i \otimes_R M_j$.
Choisissons $(i, j) \in I \times J$. Par définition d'un module de Mittag-Leffler,
la
proposition \ref{algebra-proposition-ML-characterization} (3)
s'applique aux deux systèmes. Autrement dit, il existe $i' \geq i$ et $j' \geq j$
tels que, pour tous $i'' \geq i$ et $j'' \geq j$, il existe des
morphismes $a : M_{i''} \to M_{i'}$ et $b : M_{j''} \to M_{j'}$ tels que
$f_{ii'} = a \circ f_{ii''}$ et $g_{jj'} = b \circ g_{jj''}$.
Il est alors clair que
$a \otimes b : M_{i''} \otimes_R N_{j''} \to M_{i'} \otimes_R N_{j'}$
joue le même rôle pour le système
$(M_i \otimes_R N_j, f_{ii'} \otimes g_{jj'})$.
Ainsi, d'après la caractérisation
de la proposition \ref{algebra-proposition-ML-characterization} (3),
on conclut que $M \otimes_R N$ est de Mittag-Leffler.
\end{proof}

\begin{lemma}
\label{algebra-lemma-ML-also}
Soient $R$ un anneau et $M$ un $R$-module. Alors $M$ est de Mittag-Leffler si et
seulement si, pour tout $R$-module libre de type fini $F$ et tout morphisme de modules
$f: F \to M$, il existe un $R$-module de présentation finie $Q$
et un morphisme de modules $g : F \to Q$ tels que $g$ et $f$ se dominent mutuellement, c'est-à-dire
$\Ker(f \otimes_R \text{id}_N) = \Ker(g \otimes_R \text{id}_N)$
pour tout $R$-module $N$.
\end{lemma}

\begin{proof}
Comme cette condition est manifestement plus faible que la condition (1) de la
proposition \ref{algebra-proposition-ML-characterization},
on voit qu'un module de Mittag-Leffler la satisfait.
Réciproquement, supposons que $M$ satisfasse cette condition et que
$f : P \to M$ soit un morphisme de $R$-modules d'un $R$-module de présentation finie
$P$ vers $M$. Choisissons une surjection $F \to P$, où
$F$ est un $R$-module libre de type fini. Par hypothèse, on peut trouver un morphisme
$F \to Q$, où $Q$ est un $R$-module de présentation finie, tel que
$F \to Q$ et $F \to M$ se dominent mutuellement. En particulier, le noyau
de $F \to Q$ contient le noyau de $F \to P$ ; on obtient donc un
morphisme de $R$-modules $g : P \to Q$ tel que $F \to Q$ soit égal à
la composée $F \to P \to Q$. Soit $N$ un $R$-module quelconque et
considérons le diagramme commutatif
$$
\xymatrix{
F \otimes_R N \ar[d] \ar[r] & Q \otimes_R N \\
P \otimes_R N \ar[ru] \ar[r] & M \otimes_R N
}
$$
Par hypothèse, les noyaux de $F \otimes_R N \to Q \otimes_R N$
et de $F \otimes_R N \to M \otimes_R N$ sont égaux. Comme
$F \otimes_R N \to P \otimes_R N$ est surjectif, les noyaux
de $P \otimes_R N \to Q \otimes_R N$
et de $P \otimes_R N \to M \otimes_R N$ sont donc eux aussi égaux.
\end{proof}

\begin{lemma}
\label{algebra-lemma-restrict-ML-modules}
Soit $R \to S$ un homomorphisme d'anneaux fini et de présentation finie.
Soit $M$ un $S$-module.
Si $M$ est un module de Mittag-Leffler sur $S$, alors
$M$ est un module de Mittag-Leffler sur $R$.
\end{lemma}

\begin{proof}
Supposons que $M$ soit un module de Mittag-Leffler sur $S$.
Écrivons $M = \colim M_i$ comme colimite filtrante de
$S$-modules de présentation finie $M_i$. Comme $M$ est de Mittag-Leffler sur $S$, il existe, pour tout
$i$, un indice $j \geq i$ tel que, pour tout $k \geq j$, il existe une factorisation
$f_{ij} = h \circ f_{ik}$ (où $h$ dépend de $i$, du choix de $j$ et de
$k$). Notons que, d'après le
lemme \ref{algebra-lemma-finite-finitely-presented-extension},
les modules $M_i$ sont aussi de présentation finie comme $R$-modules. De plus, tous
les morphismes $f_{ij}, f_{ik}, h$ sont des morphismes de $R$-modules. Le
système $(M_i, f_{ij})$ satisfait donc la même condition lorsqu'on le considère comme un système
de $R$-modules. Ainsi, $M$ est de Mittag-Leffler comme $R$-module.
\end{proof}

\begin{lemma}
\label{algebra-lemma-mod-ideal-ML-modules}
Soit $R$ un anneau.
Soit $S = R/I$ pour un idéal de type fini $I$.
Soit $M$ un $S$-module.
Alors $M$ est un module de Mittag-Leffler sur $R$ si et seulement si
$M$ est un module de Mittag-Leffler sur $S$.
\end{lemma}

\begin{proof}
Une implication résulte du
lemme \ref{algebra-lemma-restrict-ML-modules}.
Pour démontrer l'autre, supposons $M$ de Mittag-Leffler comme $R$-module.
Écrivons $M = \colim M_i$ comme colimite filtrante de
$S$-modules de présentation finie. Comme $I$ est de type fini, l'anneau $S$ est fini et de présentation
finie comme $R$-algèbre ; les modules $M_i$ sont donc de présentation
finie comme $R$-modules, voir le
lemme \ref{algebra-lemma-finite-finitely-presented-extension}.
Soit ensuite $N$ un $S$-module quelconque. Notons que, pour tout $i$, on a
$\Hom_R(M_i, N) = \Hom_S(M_i, N)$, puisque $R \to S$ est surjectif.
Ainsi, la condition que le système projectif
$(\Hom_R(M_i, N))_i$ soit de Mittag-Leffler implique que le système
$(\Hom_S(M_i, N))_i$ est de Mittag-Leffler. Par définition, $M$ est donc
de Mittag-Leffler sur $S$.
\end{proof}

\begin{remark}
\label{algebra-remark-go-up-ML-modules}
Soit $R \to S$ un homomorphisme d'anneaux fini et de présentation finie.
Soit $M$ un $S$-module qui est de Mittag-Leffler comme $R$-module.
En général, il n'est alors pas vrai que $M$ soit de Mittag-Leffler comme
$S$-module. Supposons par exemple que $S$ soit l'anneau des nombres duaux
sur $R$, c'est-à-dire $S = R \oplus R\epsilon$ avec $\epsilon^2 = 0$. Un
$S$-module consiste alors en un $R$-module $M$ muni d'un endomorphisme
$R$-linéaire de carré nul $\epsilon : M \to M$. Supposons maintenant que $M_0$
soit un $R$-module qui ne soit pas de Mittag-Leffler. Choisissons une présentation
$F_1 \xrightarrow{u} F_0 \to M_0 \to 0$, où $F_1$ et $F_0$ sont des $R$-modules libres.
Posons $M = F_1 \oplus F_0$ et
$$
\epsilon =
\left(
\begin{matrix}
0 & 0 \\
u & 0
\end{matrix}
\right) : M \longrightarrow M.
$$
Alors $M/\epsilon M \cong F_1 \oplus M_0$ n'est pas de Mittag-Leffler sur
$R = S/\epsilon S$, et n'est donc pas de Mittag-Leffler sur $S$ (voir le
lemme \ref{algebra-lemma-mod-ideal-ML-modules}).
D'autre part, $M/\epsilon M = M \otimes_S S/\epsilon S$, qui serait
de Mittag-Leffler sur $S$ si $M$ l'était, voir le
lemme \ref{algebra-lemma-tensor-ML-modules}.
\end{remark}


\section{Commutation des produits directs avec le produit tensoriel}
\label{algebra-section-products-tensor}

\noindent
Soit $M$ un $R$-module et soit $(Q_{\alpha})_{\alpha \in A}$ une famille de
$R$-modules. Il existe alors un morphisme canonique $M \otimes_R \left( \prod_{\alpha
\in A} Q_{\alpha} \right) \to \prod_{\alpha \in A} ( M \otimes_R
Q_{\alpha})$ donné sur les tenseurs purs par $x \otimes (q_{\alpha}) \mapsto (x
\otimes q_{\alpha})$. Ce morphisme n'est pas nécessairement injectif ni surjectif, comme
le montre l'exemple suivant.

\begin{example}
\label{algebra-example-Q-not-ML}
Prenons $R = \mathbf{Z}$, $M = \mathbf{Q}$, et considérons la famille $Q_n =
\mathbf{Z}/n$ pour $n \geq 1$. Alors $\prod_n (M \otimes Q_n) = 0$. Toutefois,
il existe une injection $\mathbf{Q} \to M \otimes (\prod_n Q_n)$
obtenue en tensorisant l'injection $\mathbf{Z} \to \prod_n Q_n$ par
$M$ ; ainsi, $M \otimes (\prod_n Q_n)$ est non nul. Le morphisme $M \otimes (\prod_n
Q_n) \to \prod_n (M \otimes Q_n)$ n'est donc pas injectif.

\medskip\noindent
D'autre part, prenons encore $R = \mathbf{Z}$, $M = \mathbf{Q}$, et posons $Q_n
= \mathbf{Z}$ pour $n \geq 1$. L'image de $M \otimes (\prod_n Q_n)
\to \prod_n (M \otimes Q_n) = \prod_n M$ est formée exactement des
suites de la forme $(a_n/m)_{n \geq 1}$, où $a_n \in \mathbf{Z}$ et où $m$ est
un entier non nul. Le morphisme n'est donc pas surjectif.
\end{example}

\noindent
Nous déterminons ci-dessous les conditions précises à imposer à $M$ pour que le morphisme $M
\otimes_R \left( \prod_{\alpha} Q_{\alpha} \right) \to \prod_{\alpha}
(M \otimes_R Q_{\alpha})$ soit surjectif, bijectif ou injectif pour tout
choix de $(Q_{\alpha})_{\alpha \in A}$. Cela est pertinent, car les modules
pour lesquels il est injectif sont exactement les modules de Mittag-Leffler
(proposition \ref{algebra-proposition-ML-tensor}). Dans la suite, si $M$ est un
$R$-module et $A$ un ensemble, nous notons $M^A$ le produit $\prod_{\alpha \in A}
M$.

\begin{proposition}
\label{algebra-proposition-fg-tensor}
Soit $M$ un $R$-module. Les assertions suivantes sont équivalentes :
\begin{enumerate}
\item $M$ est de type fini.
\item Pour toute famille $(Q_{\alpha})_{\alpha \in A}$ de $R$-modules, le
morphisme canonique $M \otimes_R \left( \prod_{\alpha} Q_{\alpha} \right)
\to \prod_{\alpha} (M \otimes_R Q_{\alpha})$ est surjectif.
\item Pour tout $R$-module $Q$ et tout ensemble $A$, le morphisme canonique $M
\otimes_R Q^{A} \to (M \otimes_R Q)^{A}$ est surjectif.
\item Pour tout ensemble $A$, le morphisme canonique $M \otimes_R R^{A} \to
M^{A}$ est surjectif.
\end{enumerate}
\end{proposition}

\begin{proof}
Montrons d'abord que (1) implique (2). Choisissons une surjection $R^n \to M$ et
considérons le diagramme commutatif
$$
\xymatrix{
R^n \otimes_R (\prod_{\alpha} Q_{\alpha})  \ar[r]^{\cong} \ar[d] &
\prod_{\alpha} (R^n \otimes_R Q_{\alpha}) \ar[d] \\
M \otimes_R (\prod_{\alpha} Q_{\alpha})  \ar[r] & \prod_{\alpha} ( M
\otimes_R Q_{\alpha}).
}
$$
La flèche supérieure est un isomorphisme et les flèches verticales sont surjectives. On
en conclut que la flèche inférieure est surjective.

\medskip\noindent
Il est clair que (2) implique (3), qui implique (4) ; il reste donc à montrer que (4) implique (1).
En fait, pour que (1) soit satisfaite, il suffit que l'élément $d = (x)_{x \in M}$ de
$M^M$ appartienne à l'image du morphisme $f: M \otimes_R R^{M} \to M^M$. Dans
ce cas, $d = \sum_{i = 1}^{n} f(x_i \otimes a_i)$ pour certains $x_i \in M$ et
$a_i \in R^M$. Si, pour $x \in M$, on note $p_x: M^M \to M$ la
projection sur le facteur indexé par $x$, alors
$$
x = p_x(d) = \sum\nolimits_{i = 1}^{n} p_x(f(x_i \otimes a_i)) =
\sum\nolimits_{i=1}^{n} p_x(a_i) x_i.
$$
Ainsi, $x_1, \ldots, x_n$ engendrent $M$.
\end{proof}

\begin{proposition}
\label{algebra-proposition-fp-tensor}
Soit $M$ un $R$-module. Les assertions suivantes sont équivalentes :
\begin{enumerate}
\item $M$ est de présentation finie.
\item Pour toute famille $(Q_{\alpha})_{\alpha \in A}$ de $R$-modules, le
morphisme canonique $M \otimes_R \left( \prod_{\alpha} Q_{\alpha} \right)
\to \prod_{\alpha} (M \otimes_R Q_{\alpha})$ est bijectif.
\item Pour tout $R$-module $Q$ et tout ensemble $A$, le morphisme canonique $M
\otimes_R Q^{A} \to (M \otimes_R Q)^{A}$ est bijectif.
\item Pour tout ensemble $A$, le morphisme canonique $M \otimes_R R^{A} \to
M^{A}$ est bijectif.
\end{enumerate}
\end{proposition}

\begin{proof}
Montrons d'abord que (1) implique (2). Choisissons une présentation $R^m \to R^n
\to M$ et considérons le diagramme commutatif
$$
\xymatrix{
R^m \otimes_R (\prod_{\alpha} Q_{\alpha}) \ar[r] \ar[d]^{\cong} & R^n
\otimes_R (\prod_{\alpha} Q_{\alpha}) \ar[r] \ar[d]^{\cong} & M \otimes_R
(\prod_{\alpha} Q_{\alpha}) \ar[r] \ar[d] & 0 \\
\prod_{\alpha} (R^m \otimes_R Q_{\alpha}) \ar[r] & \prod_{\alpha} (R^n
\otimes_R Q_{\alpha}) \ar[r] & \prod_{\alpha} (M \otimes_R Q_{\alpha})
\ar[r] & 0.
}
$$
Les deux premières flèches verticales sont des isomorphismes et les lignes sont exactes. Cela
implique que le morphisme
$M \otimes_R (\prod_{\alpha} Q_{\alpha}) \to
\prod_{\alpha} ( M \otimes_R Q_{\alpha})$
est surjectif et, par une chasse au diagramme, également injectif. Ainsi, (2) est satisfaite.

\medskip\noindent
Il est clair que (2) implique (3), qui implique (4) ; il reste donc à montrer que (4) implique (1).
D'après la proposition \ref{algebra-proposition-fg-tensor}, si (4) est satisfaite, on sait déjà que
$M$ est de type fini. On peut donc choisir une surjection $F \to M$,
où $F$ est libre de type fini. Soit $K$ le noyau. Nous devons montrer que $K$ est
de type fini. Pour tout ensemble $A$, on a un diagramme commutatif
$$
\xymatrix{
& K \otimes_R R^A \ar[r] \ar[d]_{f_3} & F \otimes_R R^A \ar[r]
\ar[d]_{f_2}^{\cong} & M \otimes_R R^A \ar[r] \ar[d]_{f_1}^{\cong} & 0 \\
0 \ar[r] & K^A \ar[r] & F^A \ar[r] & M^A \ar[r] & 0 .
}
$$
Le morphisme $f_1$ est un isomorphisme par hypothèse, le morphisme $f_2$ est un isomorphisme
puisque $F$ est libre de type fini, et les lignes sont exactes. Une chasse au diagramme montre
que $f_3$ est surjectif ; la proposition \ref{algebra-proposition-fg-tensor} montre donc
que $K$ est de type fini.
\end{proof}

\noindent
Nous aurons besoin du lemme suivant pour la prochaine proposition.

\begin{lemma}
\label{algebra-lemma-kernel-tensored-fp}
Soient $M$ un $R$-module, $P$ un $R$-module de présentation finie et $f: P
\to M$ un morphisme. Soit $Q$ un $R$-module et supposons que $x \in \Ker(P
\otimes Q \to M \otimes Q)$. Alors il existe un $R$-module de présentation finie
$P'$ et un morphisme $f': P \to P'$ tels que $f$ se factorise par
$f'$ et que $x \in \Ker(P \otimes Q \to P' \otimes Q)$.
\end{lemma}

\begin{proof}
Écrivons $M$ comme une colimite $M = \colim_{i \in I} M_i$ d'un système filtrant de
modules de présentation finie $M_i$. Puisque $P$ est de présentation finie, le morphisme $f:
P \to M$ se factorise par $M_j \to M$ pour un certain $j \in I$. Après
tensorisation par $Q$, nous avons un diagramme commutatif
$$
\xymatrix{
& M_j \otimes Q \ar[dr] & \\
P \otimes Q \ar[ur] \ar[rr] & & M \otimes Q .
}
$$
L'image $y$ de $x$ dans $M_j \otimes Q$ appartient au noyau de $M_j \otimes Q
\to M \otimes Q$. Puisque $M \otimes Q = \colim_{i \in I} (M_i
\otimes
Q)$, cela signifie que $y$ s'envoie sur $0$ dans $M_{j'} \otimes Q$ pour un certain $j' \geq j$.
Nous pouvons donc prendre $P' = M_{j'}$ et pour $f'$ le composé $P \to M_j
\to M_{j'}$.
\end{proof}

\begin{proposition}
\label{algebra-proposition-ML-tensor}
Soit $M$ un $R$-module. Les assertions suivantes sont équivalentes :
\begin{enumerate}
\item $M$ est de Mittag-Leffler.
\item Pour toute famille $(Q_{\alpha})_{\alpha \in A}$ de $R$-modules, le
morphisme canonique $M \otimes_R \left( \prod_{\alpha} Q_{\alpha} \right)
\to \prod_{\alpha} (M \otimes_R Q_{\alpha})$ est injectif.
\end{enumerate}
\end{proposition}

\begin{proof}
Montrons d'abord que (1) implique (2). Supposons que $M$ soit de Mittag-Leffler et soit $x$
dans le noyau de $M \otimes_R (\prod_{\alpha} Q_{\alpha}) \to
\prod_{\alpha} (M \otimes_R Q_{\alpha})$. Écrivons $M$ comme une colimite $M =
\colim_{i \in I} M_i$ d'un système filtrant de modules de présentation finie
$M_i$.
Alors $M \otimes_R (\prod_{\alpha} Q_{\alpha})$ est la colimite des $M_i
\otimes_R (\prod_{\alpha} Q_{\alpha})$. Ainsi, $x$ est l'image d'un élément
$x_i \in M_i \otimes_R (\prod_{\alpha} Q_{\alpha})$. Nous devons montrer que $x_i$
s'envoie sur $0$ dans $M_j \otimes_R (\prod_{\alpha} Q_{\alpha})$ pour un certain $j \geq
i$. Puisque $M$ est de Mittag-Leffler, nous pouvons choisir $j \geq i$ tel que $M_i
\to M_j$ et $M_i \to M$ se dominent mutuellement. Considérons alors
le diagramme commutatif
$$
\xymatrix{
M \otimes_R (\prod_{\alpha} Q_{\alpha}) \ar[r] & \prod_{\alpha} (M
\otimes_R Q_{\alpha}) \\
M_i \otimes_R (\prod_{\alpha} Q_{\alpha}) \ar[r]^{\cong} \ar[d] \ar[u] &
\prod_{\alpha} (M_i \otimes_R Q_{\alpha}) \ar[d] \ar[u] \\
M_j \otimes_R (\prod_{\alpha} Q_{\alpha}) \ar[r]^{\cong} & \prod_{\alpha}
(M_j \otimes_R Q_{\alpha})
}
$$
dont les deux flèches horizontales inférieures sont des isomorphismes, d'après la proposition
\ref{algebra-proposition-fp-tensor}. Puisque $x_i$ s'envoie sur $0$ dans $\prod_{\alpha} (M
\otimes_R Q_{\alpha})$, son image dans $\prod_{\alpha} (M_i \otimes_R
Q_{\alpha})$ appartient au noyau du morphisme $\prod_{\alpha} (M_i \otimes_R
Q_{\alpha}) \to \prod_{\alpha} (M \otimes_R Q_{\alpha})$. Mais ce
noyau est égal au noyau de $\prod_{\alpha} (M_i \otimes_R Q_{\alpha})
\to \prod_{\alpha} (M_j \otimes_R Q_{\alpha})$ par le
choix de $j$. Ainsi, $x_i$ s'envoie sur $0$ dans $\prod_{\alpha} (M_j \otimes_R
Q_{\alpha})$, et donc sur $0$ dans $M_j \otimes_R (\prod_{\alpha} Q_{\alpha})$.

\medskip\noindent
Supposons maintenant que (2) soit satisfaite. Nous montrons que $M$ satisfait la formulation (1) de la propriété
de Mittag-Leffler donnée dans la proposition \ref{algebra-proposition-ML-characterization}. Soit $f:
P \to M$ un morphisme d'un module de présentation finie $P$ vers $M$. Choisissons
un ensemble $B$ de représentants des classes d'isomorphisme de $R$-modules de présentation finie.
Soit $A$ l'ensemble des couples $(Q, x)$ où $Q \in B$ et $x \in
\Ker(P \otimes Q \to M \otimes Q)$. Pour $\alpha = (Q, x) \in A$, nous
notons $Q_{\alpha}$ le module $Q$ et $x_{\alpha}$ l'élément $x$. Considérons le diagramme
commutatif
$$
\xymatrix{
M \otimes_R (\prod_{\alpha} Q_{\alpha}) \ar[r] &
\prod_{\alpha} (M \otimes_R Q_{\alpha}) \\
P \otimes_R (\prod_{\alpha} Q_{\alpha}) \ar[r]^{\cong} \ar[u] &
\prod_{\alpha} (P \otimes_R Q_{\alpha}) \ar[u] .
}
$$
La flèche supérieure est injective par hypothèse, et la flèche inférieure est un
isomorphisme par la proposition \ref{algebra-proposition-fp-tensor}. Soit $x \in P
\otimes_R (\prod_{\alpha} Q_{\alpha})$ l'élément correspondant à
$(x_{\alpha}) \in \prod_{\alpha} (P \otimes_R Q_{\alpha})$ par cet
isomorphisme. Alors $x \in \Ker( P \otimes_R (\prod_{\alpha} Q_{\alpha})
\to M \otimes_R (\prod_{\alpha} Q_{\alpha}))$, puisque la flèche supérieure du
diagramme est injective. Par le lemme \ref{algebra-lemma-kernel-tensored-fp}, nous obtenons un
module de présentation finie $P'$ et un morphisme $f': P \to P'$ tels que $f: P
\to M$ se factorise par $f'$ et que $x \in \Ker(P \otimes_R
(\prod_{\alpha} Q_{\alpha}) \to P' \otimes_R (\prod_{\alpha}
Q_{\alpha}))$. Nous avons un diagramme commutatif
$$
\xymatrix{
P' \otimes_R (\prod_{\alpha} Q_{\alpha}) \ar[r]^{\cong} &
\prod_{\alpha} (P' \otimes_R Q_{\alpha}) \\
P \otimes_R (\prod_{\alpha} Q_{\alpha}) \ar[r]^{\cong} \ar[u] &
\prod_{\alpha} (P \otimes_R Q_{\alpha}) \ar[u] .
}
$$
où les flèches supérieure et inférieure sont toutes deux des isomorphismes par la proposition
\ref{algebra-proposition-fp-tensor}. Ainsi, puisque $x$ appartient au noyau du morphisme vertical
de gauche, $(x_{\alpha})$ appartient au noyau du morphisme vertical de droite. Cela
signifie que $x_{\alpha} \in \Ker(P \otimes_R Q_{\alpha} \to P' \otimes_R
Q_{\alpha})$ pour tout $\alpha \in A$. Par définition de $A$, cela signifie que
$\Ker(P \otimes_R Q \to P' \otimes_R Q) \supset \Ker(P \otimes_R
Q \to M \otimes_R Q)$ pour tout $Q$ de présentation finie et, puisque $f: P
\to M$ se factorise par $f': P \to P'$, nous avons en fait l'égalité.
Par le lemme \ref{algebra-lemma-domination-fp}, $f$ et $f'$ se dominent mutuellement.
\end{proof}

\begin{lemma}
\label{algebra-lemma-minimal-contains}
Soit $M$ un module plat de Mittag-Leffler sur $R$. Soit $F$ un $R$-module
et soit $x \in F \otimes_R M$. Alors il existe un plus petit sous-module
$F' \subset F$ tel que $x \in F' \otimes_R M$.
De plus, $F'$ est un $R$-module de type fini.
\end{lemma}

\begin{proof}
Puisque $M$ est plat, nous avons $F' \otimes_R M \subset F \otimes_R M$
si $F' \subset F$ est un sous-module, de sorte que l'énoncé a un sens.
Soit $I = \{F' \subset F \mid x \in F' \otimes_R M\}$ et, pour
$i \in I$, notons $F_i \subset F$ le sous-module correspondant.
Alors $x$ s'envoie sur zéro par le morphisme
$$
F \otimes_R M \longrightarrow \prod (F/F_i \otimes_R M)
$$
d'où, par la proposition \ref{algebra-proposition-ML-tensor},
$x$ s'envoie sur zéro par le morphisme
$$
F \otimes_R M \longrightarrow \left(\prod F/F_i\right) \otimes_R M
$$
Puisque $M$ est plat, le noyau de cette flèche est
$(\bigcap F_i) \otimes_R M$, ce qui prouve que $F' = \bigcap F_i$ convient.
Pour voir que $F'$ est un module de type fini, supposons que
$x = \sum_{j = 1, \ldots, m} f_j \otimes m_j$ avec $f_j \in F'$
et $m_j \in M$. Alors $x \in F'' \otimes_R M$, où $F'' \subset F'$ est
le sous-module engendré par $f_1, \ldots, f_m$. Bien entendu,
$F'' = F'$, ce qui donne la dernière assertion.
\end{proof}

\begin{lemma}
\label{algebra-lemma-pure-submodule-ML}
Soit $0 \to M_1 \to M_2 \to M_3 \to 0$ une suite
universellement exacte de $R$-modules. Alors :
\begin{enumerate}
\item Si $M_2$ est de Mittag-Leffler, alors $M_1$ est de Mittag-Leffler.
\item Si $M_1$ et $M_3$ sont de Mittag-Leffler, alors $M_2$ est de Mittag-Leffler.
\end{enumerate}
\end{lemma}

\begin{proof}
Pour toute famille $(Q_{\alpha})_{\alpha \in A}$ de $R$-modules, nous avons un
diagramme commutatif
$$
\xymatrix{
0 \ar[r] & M_1 \otimes_R (\prod_{\alpha} Q_{\alpha}) \ar[r] \ar[d] & M_2
\otimes_R (\prod_{\alpha} Q_{\alpha}) \ar[r] \ar[d] & M_3 \otimes_R
(\prod_{\alpha} Q_{\alpha}) \ar[r] \ar[d] & 0 \\
0 \ar[r] & \prod_{\alpha}(M_1 \otimes Q_{\alpha}) \ar[r] & \prod_{\alpha}(M_2
\otimes Q_{\alpha}) \ar[r] & \prod_{\alpha}(M_3 \otimes Q_{\alpha})\ar[r] & 0
}
$$
dont les lignes sont exactes. Ainsi, (1) et (2) résultent de la
proposition \ref{algebra-proposition-ML-tensor}.
\end{proof}

\begin{lemma}
\label{algebra-lemma-quotient-module-ML}
Soit $M_1 \to M_2 \to M_3 \to 0$ une suite exacte de $R$-modules.
Si $M_1$ est de type fini et $M_2$ est de Mittag-Leffler, alors $M_3$
est de Mittag-Leffler.
\end{lemma}

\begin{proof}
Pour toute famille $(Q_{\alpha})_{\alpha \in A}$ de $R$-modules,
puisque le produit tensoriel est exact à droite, nous avons un diagramme commutatif
$$
\xymatrix{
M_1 \otimes_R (\prod_{\alpha} Q_{\alpha}) \ar[r] \ar[d] & M_2
\otimes_R (\prod_{\alpha} Q_{\alpha}) \ar[r] \ar[d] & M_3 \otimes_R
(\prod_{\alpha} Q_{\alpha}) \ar[r] \ar[d] & 0 \\
\prod_{\alpha}(M_1 \otimes Q_{\alpha}) \ar[r] & \prod_{\alpha}(M_2
\otimes Q_{\alpha}) \ar[r] & \prod_{\alpha}(M_3 \otimes Q_{\alpha})\ar[r] & 0
}
$$
dont les lignes sont exactes. Par la proposition \ref{algebra-proposition-fg-tensor},
la flèche verticale de gauche est surjective. Par la
proposition \ref{algebra-proposition-ML-tensor}, la flèche verticale du milieu
est injective. Une chasse au diagramme montre que la flèche verticale de droite
est injective. Ainsi, $M_3$ est de Mittag-Leffler par la
proposition \ref{algebra-proposition-ML-tensor}.
\end{proof}


\begin{lemma}
\label{algebra-lemma-colimit-universally-injective-ML}
Si $M = \colim M_i$ est la colimite d'un système filtrant de $R$-modules de Mittag-Leffler
$M_i$ dont les morphismes de transition sont universellement injectifs, alors $M$ est de
Mittag-Leffler.
\end{lemma}

\begin{proof}
Soit $(Q_{\alpha})_{\alpha \in A}$ une famille de $R$-modules. Nous devons
montrer que $M \otimes_R (\prod Q_\alpha) \to \prod M \otimes_R Q_\alpha$
est injectif, et nous savons que
$M_i \otimes_R (\prod Q_\alpha) \to \prod M_i \otimes_R Q_\alpha$
est injectif pour tout $i$, voir la proposition \ref{algebra-proposition-ML-tensor}.
Comme $\otimes$ commute aux colimites filtrantes, il suffit de montrer que
$\prod M_i \otimes_R Q_\alpha \to \prod M \otimes_R Q_\alpha$
est injectif. Cela est clair, car chacun des morphismes
$M_i \otimes_R Q_\alpha \to M \otimes_R Q_\alpha$ est injectif
par l'hypothèse que les morphismes de transition sont universellement injectifs.
\end{proof}

\begin{lemma}
\label{algebra-lemma-direct-sum-ML}
Si $M = \bigoplus_{i \in I} M_i$ est une somme directe de $R$-modules, alors $M$ est
de Mittag-Leffler si et seulement si chaque $M_i$ est de Mittag-Leffler.
\end{lemma}

\begin{proof}
Le sens « seulement si » résulte du lemme \ref{algebra-lemma-pure-submodule-ML} (1)
et du fait qu'une suite exacte courte scindée est universellement exacte. La
réciproque résulte du lemme \ref{algebra-lemma-colimit-universally-injective-ML},
mais nous pouvons aussi raisonner directement comme suit. Remarquons d'abord que si $I$ est
fini, cela résulte du lemme
\ref{algebra-lemma-pure-submodule-ML} (2). Pour $I$ quelconque, si tous les $M_i$ sont
de Mittag-Leffler, nous montrons qu'il en est de même de $M$ en vérifiant la condition (1) de la
proposition \ref{algebra-proposition-ML-characterization}.
Soit $f: P \to M$ un morphisme d'un module de présentation finie $P$.
Alors $f$ se factorise sous la forme
$P \xrightarrow{f'} \bigoplus_{i' \in I'} M_{i'} \hookrightarrow
\bigoplus_{i \in I} M_i$
pour une certaine partie finie $I'$ de $I$. Par le cas fini,
$\bigoplus_{i' \in I'} M_{i'}$ est de Mittag-Leffler, et il existe donc
un module de présentation finie $Q$ et un morphisme $g: P \to Q$ tels que $g$
et $f'$ se dominent mutuellement. Alors $g$ et $f$ se dominent également mutuellement.
\end{proof}

\begin{lemma}
\label{algebra-lemma-flat-ML-over-ML-ring}
Soit $R \to S$ un homomorphisme d'anneaux. Soit $M$ un $S$-module.
Si $S$ est de Mittag-Leffler comme $R$-module, et si $M$ est plat et de Mittag-Leffler
comme $S$-module, alors $M$ est de Mittag-Leffler comme $R$-module.
\end{lemma}

\begin{proof}
Nous déduisons ceci de la caractérisation de la
proposition \ref{algebra-proposition-ML-tensor}.
En effet, supposons que les $Q_\alpha$ forment une famille de $R$-modules.
Considérons le composé
$$
\xymatrix{
M \otimes_R \prod_\alpha Q_\alpha =
M \otimes_S S \otimes_R \prod_\alpha Q_\alpha \ar[d] \\
M \otimes_S \prod_\alpha (S \otimes_R Q_\alpha) \ar[d] \\
\prod_\alpha (M \otimes_S S \otimes_R Q_\alpha) =
\prod_\alpha (M \otimes_R Q_\alpha)
}
$$
La première flèche est injective puisque $M$ est plat sur $S$ et que $S$ est
de Mittag-Leffler sur $R$, et la seconde flèche est injective puisque
$M$ est de Mittag-Leffler sur $S$. Ainsi, $M$ est de Mittag-Leffler sur $R$.
\end{proof}


\section{Anneaux cohérents}
\label{algebra-section-coherent}

\noindent
Nous utilisons la discussion sur la commutation de $\prod$ et $\otimes$ pour déterminer
les anneaux pour lesquels les produits de modules plats sont plats. Il s'avère qu'il
s'agit des anneaux dits cohérents. La notion de module $\mathcal{O}_X$ cohérent
sur un espace annelé vous est peut-être plus familière, voir
Modules, section \ref{modules-section-coherent}.

\begin{definition}
\label{algebra-definition-coherent}
Soit $R$ un anneau. Soit $M$ un $R$-module.
\begin{enumerate}
\item Nous disons que $M$ est un {\it module cohérent} s'il est de type fini
et si tout sous-module de type fini de $M$ est de présentation finie sur
$R$.
\item Nous disons que $R$ est un {\it anneau cohérent} s'il est cohérent comme module
sur lui-même.
\end{enumerate}
\end{definition}

\noindent
Ainsi, un anneau est cohérent si et seulement si tout idéal de type fini est de
présentation finie comme module.

\begin{example}
\label{algebra-example-valuation-ring-coherent}
Un anneau de valuation est un anneau cohérent. En effet, tout idéal non nul de type fini
est principal (lemme \ref{algebra-lemma-characterize-valuation-ring}),
donc libre puisqu'un anneau de valuation est intègre, et donc de présentation finie.
\end{example}

\noindent
La catégorie des modules cohérents est abélienne.

\begin{lemma}
\label{algebra-lemma-coherent}
Soit $R$ un anneau.
\begin{enumerate}
\item Un sous-module de type fini d'un module cohérent est cohérent.
\item Soit $\varphi : N \to M$ un homomorphisme d'un module de type fini
vers un module cohérent. Alors $\Ker(\varphi)$ est de type fini,
$\Im(\varphi)$ est
cohérent et $\Coker(\varphi)$ est cohérent.
\item Soit $\varphi : N \to M$ un homomorphisme de modules cohérents.
Alors $\Ker(\varphi)$ et $\Coker(\varphi)$ sont des modules
cohérents.
\item Étant donnée une suite exacte courte de $R$-modules
$0 \to M_1 \to M_2 \to M_3 \to 0$, si deux d'entre eux sont cohérents,
le troisième l'est aussi.
\end{enumerate}
\end{lemma}

\begin{proof}
La première assertion résulte immédiatement de la définition.

\medskip\noindent
Soit $\varphi : N \to M$ satisfaisant aux hypothèses de (2). Tout d'abord,
$\Im(\varphi)$ est de type fini,
donc cohérent par (1). En particulier, $\Im(\varphi)$ est de présentation finie ;
en appliquant le lemme \ref{algebra-lemma-extension} à la suite exacte
$0 \to \Ker(\varphi) \to N \to \Im(\varphi) \to 0$, nous voyons que $\Ker(\varphi)$
est de type fini.
Pour montrer que $\Coker(\varphi)$ est cohérent, soit $E \subset \Coker(\varphi)$
un sous-module de type fini, et soit $E'$ son image réciproque dans $M$.
De la suite
exacte $0 \to \Im(\varphi) \to E' \to E \to 0$ et puisque $\Ker(\varphi)$
est de type fini, nous concluons par le lemme \ref{algebra-lemma-extension} que $E' \subset M$ est
de type fini,
donc de présentation finie puisque $M$ est cohérent. La même suite exacte
montre alors que $E$ est de présentation finie, d'où l'assertion.

\medskip\noindent
La partie (3) résulte immédiatement de (1) et (2).

\medskip\noindent
Soit $0 \to M_1 \xrightarrow{i} M_2 \xrightarrow{p} M_3 \to 0$
une suite exacte courte de $R$-modules comme dans (4). Il reste
à montrer que si $M_1$ et $M_3$ sont cohérents,
alors $M_2$ l'est. Par le lemme \ref{algebra-lemma-extension},
nous voyons que $M_2$ est de type fini. Soit $N_2 \subset M_2$ un sous-module
de type fini.
Posons $N_3 = p(N_2) \subset M_3$ et $N_1 = i^{-1}(N_2) \subset M_1$.
Nous avons une suite exacte $0 \to N_1 \to N_2 \to N_3 \to 0$. Il est clair que $N_3$
est de type fini (comme quotient de $N_2$), donc de présentation finie (comme sous-module
de type fini de $M_3$). Il résulte du lemme \ref{algebra-lemma-extension} (5) que
$N_1$ est de type fini,
donc de présentation finie (comme sous-module de type fini de $M_1$). Nous concluons par le
lemme \ref{algebra-lemma-extension} (2) que $N_2$ est de présentation finie.
\end{proof}

\begin{lemma}
\label{algebra-lemma-coherent-ring}
Soit $R$ un anneau. Si $R$ est cohérent, alors un module est cohérent
si et seulement s'il est de présentation finie.
\end{lemma}

\begin{proof}
Il est clair qu'un module cohérent est de présentation finie (sur tout anneau).
Réciproquement, si $R$ est cohérent, alors $R^{\oplus n}$ est cohérent, de même que
le conoyau de tout morphisme $R^{\oplus m} \to R^{\oplus n}$, voir le
lemme \ref{algebra-lemma-coherent}.
\end{proof}

\begin{lemma}
\label{algebra-lemma-Noetherian-coherent}
Un anneau noethérien est un anneau cohérent.
\end{lemma}

\begin{proof}
Par le
lemme \ref{algebra-lemma-Noetherian-finite-type-is-finite-presentation},
tout $R$-module de type fini est de présentation finie. En particulier, tout idéal de $R$
est de présentation finie.
\end{proof}

\begin{proposition}
\label{algebra-proposition-characterize-coherent}
\begin{reference}
C'est \cite[Theorem 2.1]{Chase}.
\end{reference}
Soit $R$ un anneau. Les assertions suivantes sont équivalentes :
\begin{enumerate}
\item $R$ est cohérent,
\item tout produit de $R$-modules plats est plat, et
\item pour tout ensemble $A$, le module $R^A$ est plat.
\end{enumerate}
\end{proposition}

\begin{proof}
Supposons $R$ cohérent, et soit $Q_\alpha$, $\alpha \in A$, un ensemble de
$R$-modules plats. Nous devons montrer que
$I \otimes_R \prod_\alpha Q_\alpha \to \prod Q_\alpha$ est injectif
pour tout idéal de type fini $I$ de $R$, voir le
lemme \ref{algebra-lemma-flat}.
Puisque $R$ est cohérent, $I$ est un $R$-module de présentation finie.
Ainsi, $I \otimes_R \prod_\alpha Q_\alpha = \prod I \otimes_R Q_\alpha$ par la
proposition \ref{algebra-proposition-fp-tensor}.
L'injectivité voulue résulte de ce que $I \otimes_R Q_\alpha \to Q_\alpha$
est injectif par platitude de $Q_\alpha$.

\medskip\noindent
L'implication (2) $\Rightarrow$ (3) est triviale.

\medskip\noindent
Supposons que le $R$-module $R^A$ soit plat pour tout ensemble $A$. Soit $I$
un idéal de type fini de $R$. Alors $I \otimes_R R^A \to R^A$
est injectif par hypothèse. Par la
proposition \ref{algebra-proposition-fg-tensor}
et la finitude de $I$, l'image est égale à $I^A$. Ainsi,
$I \otimes_R R^A = I^A$ pour tout ensemble $A$, et nous concluons que $I$
est de présentation finie par la
proposition \ref{algebra-proposition-fp-tensor}.
\end{proof}







\section{Exemples et contre-exemples de modules de Mittag-Leffler}
\label{algebra-section-examples}

\noindent
Nous terminons cette section par quelques exemples et contre-exemples de modules de
Mittag-Leffler.

\begin{example}
\label{algebra-example-ML}
Modules de Mittag-Leffler.
\begin{enumerate}
\item Tout module de présentation finie est de Mittag-Leffler. Cela résulte, par
exemple, de la proposition \ref{algebra-proposition-ML-characterization} (1). Plus
généralement, un module de type fini est de Mittag-Leffler si et
seulement s'il est de présentation finie. Cela résulte des propositions
\ref{algebra-proposition-fg-tensor}, \ref{algebra-proposition-fp-tensor} et
\ref{algebra-proposition-ML-tensor}.
\item Un module libre est de Mittag-Leffler puisqu'il satisfait la condition (1) de la
proposition \ref{algebra-proposition-ML-characterization}.
\item Par l'exemple précédent et le lemme \ref{algebra-lemma-direct-sum-ML},
les modules projectifs sont de Mittag-Leffler.
\end{enumerate}
\end{example}

\noindent
Nous voulons aussi ajouter à notre liste d'exemples les anneaux de séries formelles sur un
anneau noethérien $R$. Ce sera une conséquence du lemme suivant.

\begin{lemma}
\label{algebra-lemma-flat-ML-criterion}
Soit $M$ un $R$-module plat. Les assertions suivantes sont équivalentes :
\begin{enumerate}
\item $M$ est de Mittag-Leffler, et
\item si $F$ est un $R$-module libre de type fini et si
$x \in F \otimes_R M$, alors il existe un plus petit sous-module $F'$ de $F$
tel que $x \in F' \otimes_R M$.
\end{enumerate}
\end{lemma}

\begin{proof}
L'implication (1) $\Rightarrow$ (2) est un cas particulier du
lemme \ref{algebra-lemma-minimal-contains}. Supposons (2).
Par le théorème \ref{algebra-theorem-lazard}, nous pouvons écrire $M$ comme la colimite
$M = \colim_{i \in I} M_i$ d'un système filtrant $(M_i, f_{ij})$ de
$R$-modules libres de type fini.
Par la remarque \ref{algebra-remark-flat-ML}, il suffit de montrer que le système projectif
$(\Hom_R(M_i, R), \Hom_R(f_{ij}, R))$ est de Mittag-Leffler. Autrement
dit,
fixons $i \in I$ et, pour $j \geq i$, soit $Q_j$ l'image de
$\Hom_R(M_j, R)
\to \Hom_R(M_i, R)$ ; nous devons montrer que les $Q_j$ se stabilisent.

\medskip\noindent
Puisque $M_i$ est libre de type fini, nous pouvons faire l'identification
$\Hom_R(M_i, M_j) = \Hom_R(M_i, R) \otimes_R  M_j$ pour tout $j$.
En utilisant le
fait que les $M_j$ sont libres, il s'ensuit que, pour $j \geq i$, $Q_j$ est le
plus petit sous-module de $\Hom_R(M_i, R)$ tel que $f_{ij} \in Q_j
\otimes_R
M_j$. Sous l'identification $\Hom_R(M_i, M) = \Hom_R(M_i, R)
\otimes_R
M$, le morphisme canonique $f_i: M_i \to M$ appartient à $\Hom_R(M_i, R)
\otimes_R M$. Par l'hypothèse sur $M$, il existe un plus petit sous-module
$Q$ de $\Hom_R(M_i, R)$ tel que $f_i \in Q \otimes_R M$. Nous allons
montrer
que les $Q_j$ se stabilisent à $Q$.

\medskip\noindent
Pour $j \geq i$, nous avons un diagramme commutatif
$$
\xymatrix{
Q_j \otimes_R M_j \ar[r] \ar[d] & \Hom_R(M_i, R) \otimes_R M_j
\ar[d] \\
Q_j \otimes_R M \ar[r] & \Hom_R(M_i, R) \otimes_R M.
}
$$
Puisque $f_{ij} \in Q_j \otimes_R M_j$ s'envoie sur $f_i \in \Hom_R(M_i, R)
\otimes_R M$, il s'ensuit que $f_i \in Q_j \otimes_R M$. Ainsi, par le
choix de $Q$, nous avons $Q \subset Q_j$ pour tout $j \geq i$.

\medskip\noindent
Puisque les $Q_j$ sont décroissants et que $Q \subset Q_j$ pour tout $j \geq i$, pour montrer
que les $Q_j$ se stabilisent à $Q$, il suffit de trouver un $j \geq i$ tel que $Q_j
\subset Q$. Comme élément de
$$
\Hom_R(M_i, R) \otimes_R M = \colim_{j \in J}
(\Hom_R(M_i, R) \otimes_R
M_j),
$$
$f_i$ est la colimite des $f_{ij}$ pour $j \geq i$, et $f_i$ appartient également au
sous-module
$$
\colim_{j \in J} (Q \otimes_R M_j) \subset \colim_{j \in J}
(\Hom_R(M_i, R)
\otimes_R M_j).
$$
Il s'ensuit que, pour un certain $j \geq i$, $f_{ij}$ appartient à $Q \otimes_R M_j$.
Puisque $Q_j$ est le plus petit sous-module de $\Hom_R(M_i, R)$ tel que $f_{ij}
\in
Q_j \otimes_R M_j$, nous concluons que $Q_j\subset Q$.
\end{proof}

\begin{lemma}
\label{algebra-lemma-product-over-Noetherian-ring}
Soient $R$ un anneau noethérien et $A$ un ensemble.
Alors $M = R^A$ est un $R$-module plat et de Mittag-Leffler.
\end{lemma}

\begin{proof}
En combinant le
lemme \ref{algebra-lemma-Noetherian-coherent}
et la
proposition \ref{algebra-proposition-characterize-coherent},
nous voyons que $M$ est plat sur $R$. Nous montrons que $M$ satisfait la condition du
lemme \ref{algebra-lemma-flat-ML-criterion}.
Soit $F$ un $R$-module libre de type fini. Si $F'$ est un sous-module quelconque de $F$, alors il
est de présentation finie puisque $R$ est noethérien. Par la
proposition \ref{algebra-proposition-fp-tensor},
nous avons donc un diagramme commutatif
$$
\xymatrix{
F' \otimes_R M \ar[r] \ar[d]^{\cong} & F \otimes_R M \ar[d]^{\cong} \\
(F')^A \ar[r] & F^A
}
$$
par lequel nous pouvons identifier le morphisme $F' \otimes_R M \to F \otimes_R M$
à $(F')^A \to F^A$. Ainsi, si $x \in F \otimes_R M$ correspond à
$(x_\alpha) \in F^A$, alors le sous-module $F'$ de $F$ engendré par les
$x_\alpha$ est le plus petit sous-module de $F$ tel que $x \in F' \otimes_R M$.
\end{proof}

\begin{lemma}
\label{algebra-lemma-power-series-ML}
Soient $R$ un anneau noethérien et $n$ un entier positif. Alors le $R$-module
$M = R[[t_1, \ldots, t_n]]$ est plat et de Mittag-Leffler.
\end{lemma}

\begin{proof}
Comme $R$-module, nous avons $M = R^A$ pour un ensemble (dénombrable) $A$.
Ce lemme est donc un cas particulier du
lemme \ref{algebra-lemma-product-over-Noetherian-ring}.
\end{proof}

\begin{example}
\label{algebra-example-not-ML}
Modules qui ne sont pas de Mittag-Leffler.
\begin{enumerate}
\item Par l'exemple \ref{algebra-example-Q-not-ML} et la
proposition \ref{algebra-proposition-ML-tensor}, $\mathbf{Q}$ n'est pas un
$\mathbf{Z}$-module de Mittag-Leffler.
\item Nous démontrons ci-dessous (théorème \ref{algebra-theorem-projectivity-characterization}) que,
pour un module plat et dénombrablement engendré, la projectivité équivaut à être
de Mittag-Leffler. Ainsi, tout module $M$ plat, dénombrablement engendré et non projectif
est un exemple de module qui n'est pas de Mittag-Leffler. Pour un tel exemple, voir la
remarque \ref{algebra-remark-warning}.
\item Soit $k$ un corps. Soit $R = k[[x]]$. Le $R$-module
$M = \prod_{n \in \mathbf{N}} R/(x^n)$ n'est pas de Mittag-Leffler.
En effet, considérons l'élément $\xi = (\xi_1, \xi_2, \xi_3, \ldots)$
défini par $\xi_{2^m} = x^{2^{m - 1}}$ et $\xi_n = 0$ sinon, de sorte que
$$
\xi = (0, x, 0, x^2, 0, 0, 0, x^4, 0, 0, 0, 0, 0, 0, 0, x^8, \ldots)
$$
Alors l'annulateur de $\xi$ dans $M/x^{2^m}M$ est engendré par $x^{2^{m - 1}}$
pour $m \gg 0$. Mais si $M$ était de Mittag-Leffler, il existerait un
$R$-module de type fini $Q$ et un élément $\xi' \in Q$ tels que l'annulateur
de $\xi'$ dans $Q/x^l Q$ coïncide avec l'annulateur de $\xi$ dans $M/x^l M$
pour tout $l \geq 1$, voir la
proposition \ref{algebra-proposition-ML-characterization} (1).
Vous pouvez alors montrer qu'il existe un entier $a \geq 0$ tel que
l'annulateur de $\xi'$ dans $Q/x^l Q$ soit engendré
soit par $x^a$, soit par $x^{l - a}$ pour tout $l \gg 0$ (selon que
$\xi' \in Q$ est de torsion ou non). En combinant ce qui précède,
on obtiendrait, pour tout $l = 2^m >> 0$, l'égalité $a = l/2$ ou $l - a = l/2$,
ce qui est absurde.
\item Le même argument montre que le complété $(x)$-adique de
$\bigoplus_{n \in \mathbf{N}} R/(x^n)$ n'est pas de Mittag-Leffler sur
$R = k[[x]]$ (indication : $\xi$ est en fait un élément de ce complété).
\item Soit $R = k[a, b]/(a^2, ab, b^2)$. Soit $S$ la $R$-algèbre de présentation finie
présentée par $S = R[t]/(at - b)$. Alors, comme $R$-module,
$S$ est dénombrablement engendré et indécomposable (détails omis).
D'autre part, $R$ est local artinien, donc local complet,
et donc un anneau local hensélien, voir le
lemme \ref{algebra-lemma-complete-henselian}.
Si $S$ était de Mittag-Leffler comme $R$-module, alors ce serait une
somme directe de $R$-modules de type fini par le
lemme \ref{algebra-lemma-split-ML-henselian}.
Nous concluons donc que $S$ n'est pas de Mittag-Leffler comme $R$-module.
\end{enumerate}
\end{example}




\section{Modules de Mittag-Leffler dénombrablement engendrés}
\label{algebra-section-ML-countable}

\noindent
Il s'avère que les modules de Mittag-Leffler dénombrablement engendrés ont une
structure particulièrement simple.

\begin{lemma}
\label{algebra-lemma-ML-countable-colimit}
Soit $M$ un $R$-module. Écrivons $M = \colim_{i \in I} M_i$, où $(M_i,
f_{ij})$ est un système filtrant de $R$-modules de présentation finie. Si $M$ est
de Mittag-Leffler et dénombrablement engendré, alors il existe une partie filtrante dénombrable
$I' \subset I$ telle que $M \cong \colim_{i \in I'} M_i$.
\end{lemma}

\begin{proof}
Soit $x_1, x_2, \ldots$ une famille dénombrable de générateurs de $M$. Pour chaque $x_n$,
choisissons $i \in I$ tel que $x_n$ appartienne à l'image du morphisme canonique $f_i:
M_i \to M$ ; soit $I'_{0} \subset I$ l'ensemble de tous ces $i$. Puisque
$M$ est de Mittag-Leffler, pour chaque $i \in I'_{0}$, nous pouvons choisir $j \in I$
tel que $j \geq i$ et que $f_{ij}: M_i \to M_j$ se factorise par
$f_{ik}: M_i \to M_k$ pour tout $k \geq i$ (condition (3) de la proposition
\ref{algebra-proposition-ML-characterization}) ; soit $I'_1$ la réunion de $I'_0$ et de
tous ces $j$. Puisque $I'_1$ est dénombrable, nous pouvons l'agrandir en un
ensemble préordonné filtrant dénombrable $I'_{2} \subset I$. Nous pouvons maintenant appliquer à $I'_{2}$
le même procédé qu'à $I'_{0}$ pour obtenir un nouvel ensemble dénombrable $I'_{3} \subset
I$. Puis nous agrandissons $I'_{3}$ en un ensemble préordonné filtrant dénombrable $I'_{4}$. En continuant
ainsi --- en ajoutant un $j$ comme dans la proposition
\ref{algebra-proposition-ML-characterization} (3) pour chaque $ i \in I'_{\ell}$ si $\ell$
est pair, et en agrandissant $I'_{\ell}$ en un ensemble préordonné filtrant si $\ell$ est impair --- nous obtenons une
suite de parties $I'_{\ell} \subset I$ pour $\ell \geq 0$. La réunion $I' =
\bigcup I'_{\ell}$ satisfait :
\begin{enumerate}
\item $I'$ est dénombrable et filtrant ;
\item chaque $x_n$ appartient à l'image de $f_i: M_i \to M$ pour un certain $i
\in I'$ ;
\item si $i \in I'$, alors il existe $j \in I'$ tel que $j \geq i$ et que $f_{ij}:
M_i \to M_j$ se factorise par $f_{ik}: M_i \to M_k$ pour tout
$k \in I$ tel que $k \geq i$. En particulier, $\Ker(f_{ik}) \subset \Ker(f_{ij})$
pour $k \geq i$.
\end{enumerate}
Nous affirmons que le morphisme canonique $\colim_{i \in I'} M_i \to
\colim_{i
\in I} M_i = M$ est un isomorphisme. Par (2), il est surjectif. Pour l'injectivité,
supposons que $x \in \colim_{i \in I'} M_i$ s'envoie sur $0$ dans $\colim_{i \in
I} M_i$.
En représentant $x$ par un élément $\tilde{x} \in M_i$ pour un certain $i \in I'$, cela
signifie que $f_{ik}(\tilde{x}) = 0$ pour un certain $k \in I, k \geq i$. Mais, par
(3), il existe $j \in I', j \geq i,$ tel que $f_{ij}(\tilde{x}) = 0$. Ainsi, $x
= 0$ dans $\colim_{i \in I'} M_i$.
\end{proof}

\noindent
Le lemme \ref{algebra-lemma-ML-countable-colimit}
implique qu'un module de Mittag-Leffler $M$ dénombrablement engendré sur
$R$ est la colimite d'un système
$$
M_1 \to M_2 \to M_3 \to M_4 \to \ldots
$$
où chaque $M_n$ est un $R$-module de présentation finie. Pour le voir, raisonnons comme dans la
démonstration du
lemme \ref{algebra-lemma-ML-limit-nonempty}
pour constater qu'un ensemble préordonné filtrant dénombrable possède une partie cofinale
isomorphe à $(\mathbf{N}, \geq)$. Supposons que
$R = k[x_1, x_2, x_3, \ldots]$ et $M = R/(x_i)$. Alors $M$ est de type
fini, mais n'est pas de présentation finie, et n'est donc pas de Mittag-Leffler (voir
l'exemple \ref{algebra-example-ML}, partie (1)).
Bien entendu, on peut écrire $M = \colim_n M_n$ en prenant
$M_n = R/(x_1, \ldots, x_n)$ ; ainsi, la possibilité d'écrire
$M$ comme une telle colimite n'implique pas que $M$ soit de Mittag-Leffler.

\begin{lemma}
\label{algebra-lemma-ML-countable}
Soit $R$ un anneau.
Soit $M$ un $R$-module.
Supposons $M$ de Mittag-Leffler et dénombrablement engendré.
Pour tout morphisme de $R$-modules $f : P \to M$, où $P$ est de type fini, il
existe un endomorphisme $\alpha : M \to M$ tel que
\begin{enumerate}
\item $\alpha : M \to M$ se factorise par un $R$-module de présentation finie, et
\item $\alpha \circ f = f$.
\end{enumerate}
\end{lemma}

\begin{proof}
Écrivons $M = \colim_{i \in I} M_i$ comme colimite filtrante de $R$-modules de
présentation finie, avec $I$ dénombrable, voir le
lemme \ref{algebra-lemma-ML-countable-colimit}.
Les morphismes de transition sont notés $f_{ij}$ et nous utilisons $f_i : M_i \to M$
pour désigner les morphismes canoniques vers $M$. Posons $N = \prod_{s \in I} M_s$. Notons
$$
M_i^* = \Hom_R(M_i, N) = \prod\nolimits_{s \in I} \Hom_R(M_i, M_s)
$$
de sorte que $(M_i^*)$ est un système projectif de $R$-modules indexé par $I$.
Remarquons que $\Hom_R(M, N) = \lim M_i^*$.
Comme $M$ est de Mittag-Leffler, nous trouvons, pour tout
$i \in I$, un indice $k(i) \geq i$ tel que
$$
E_i := \bigcap\nolimits_{i' \geq i} \Im(M_{i'}^* \to M_i^*)
=
\Im(M_{k(i)}^* \to M_i^*)
$$
Choisissons et fixons $j \in I$ tel que $\Im(P \to M) \subset \Im(M_j \to M)$.
C'est possible puisque $P$ est de type fini. Posons $k = k(j)$.
Soit
$x = (0, \ldots, 0, \text{id}_{M_k}, 0, \ldots, 0) \in M_k^*$ et
remarquons qu'il s'envoie sur $y = (0, \ldots, 0, f_{jk}, 0, \ldots, 0) \in M_j^*$.
Par notre choix de $k$, nous voyons que $y \in E_j$. Par
l'exemple \ref{algebra-example-ML-surjective-maps},
les morphismes de transition $E_i \to E_j$ sont surjectifs pour tout $i \geq j$
et $\lim E_i = \lim M_i^* = \Hom_R(M, N)$. Ainsi, le
lemme \ref{algebra-lemma-ML-limit-nonempty}
garantit l'existence d'un élément $z \in \Hom_R(M, N)$
qui s'envoie sur $y$ dans $E_j \subset M_j^*$. Soit $z_k$ la $k$-ième composante
de $z$. Alors $z_k : M \to M_k$ est un homomorphisme tel que
$$
\xymatrix{
M \ar[r]_{z_k} & M_k \\
M_j \ar[ru]_{f_{jk}} \ar[u]^{f_j}
}
$$
commute. Soit $\alpha : M \to M$ le composé
$f_k \circ z_k : M \to M_k \to M$.
Alors $\alpha$ se factorise par un module de présentation finie par construction et
$\alpha \circ f_j = f_j$. Puisque l'image de $f$ est contenue dans l'image de
$f_j$, cela implique aussi que $\alpha \circ f = f$.
\end{proof}

\noindent
Nous verrons plus loin (voir le
lemme \ref{algebra-lemma-split-ML-henselian})
que le
lemme \ref{algebra-lemma-ML-countable}
signifie qu'un module de Mittag-Leffler dénombrablement engendré sur un anneau local
hensélien est une somme directe de modules de présentation finie.




\section{Caractérisation des modules projectifs}
\label{algebra-section-characterize-projective}

\noindent
L'objectif de cette section est de démontrer qu'un module est projectif si et seulement s'il
est plat, de Mittag-Leffler et somme directe de modules dénombrablement engendrés
(théorème \ref{algebra-theorem-projectivity-characterization} ci-dessous).

\begin{lemma}
\label{algebra-lemma-countgen-projective}
Soit $M$ un $R$-module. Si $M$ est plat, de Mittag-Leffler et
dénombrablement engendré, alors $M$ est projectif.
\end{lemma}

\begin{proof}
Par le théorème de Lazard (théorème \ref{algebra-theorem-lazard}), nous pouvons écrire $M =
\colim_{i
\in I} M_i$ pour un système filtrant de $R$-modules libres de type fini $(M_i, f_{ij})$
indexé par un ensemble $I$. Par le lemme \ref{algebra-lemma-ML-countable-colimit}, nous pouvons supposer
$I$ dénombrable. Soit maintenant
$$
0 \to N_1 \to N_2 \to N_3 \to 0
$$
une suite exacte de $R$-modules. Nous devons montrer que l'application de
$\Hom_R(M, -)$
préserve l'exactitude. Puisque $M_i$ est libre de type fini,
$$
0 \to \Hom_R(M_i, N_1) \to \Hom_R(M_i, N_2) \to
\Hom_R(M_i, N_3) \to 0
$$
est exacte pour tout $i$. Puisque $M$ est de Mittag-Leffler, $(\Hom_R(M_i,
N_{1}))$ est
un système projectif de Mittag-Leffler. Par le lemme \ref{algebra-lemma-ML-exact-sequence},
$$
0 \to \lim_{i \in I} \Hom_R(M_i, N_1) \to
\lim_{i \in I} \Hom_R(M_i, N_2) \to
\lim_{i \in I} \Hom_R(M_i, N_3) \to 0
$$
est donc exacte. Mais, pour tout $R$-module $N$, il existe un isomorphisme fonctoriel
$\Hom_R(M, N) \cong \lim_{i \in I} \Hom_R(M_i, N)$, de sorte que
$$
0 \to \Hom_R(M, N_1) \to \Hom_R(M, N_2) \to
\Hom_R(M, N_3) \to 0
$$
est exacte.
\end{proof}

\begin{remark}
\label{algebra-remark-characterize-projective}
Le lemme \ref{algebra-lemma-countgen-projective} n'est plus vrai sans l'hypothèse
d'engendrement dénombrable. Par exemple, le $\mathbf Z$-module $M =
\mathbf{Z}[[x]]$ est plat et de Mittag-Leffler, mais n'est pas projectif. Il est
de Mittag-Leffler par le lemme \ref{algebra-lemma-power-series-ML}. Les sous-groupes des groupes abéliens
libres sont libres ; un $\mathbf Z$-module projectif est donc en fait libre, et il en va de même
de ses sous-modules. Ainsi, pour montrer que $M$ n'est pas projectif, il suffit de construire
un sous-module non libre. Fixons un nombre premier $p$ et considérons le sous-module $N$
formé des séries formelles $f(x) = \sum a_i x^i$ telles que, pour tout entier $m
\geq 1$, $p^m$ divise $a_i$ pour tous les $i$ sauf un nombre fini. Alors $\sum a_i p^i
x^i$ appartient à $N$ pour tous $a_i \in \mathbf{Z}$, et $N$ n'est donc pas dénombrable. Ainsi, si
$N$ était libre, il serait de rang non dénombrable, et la dimension de $N/pN$ sur
$\mathbf{Z}/p$ ne serait pas dénombrable. Ce n'est pas le cas, car les éléments $x^i \in
N/pN$, pour $i \geq 0$, engendrent $N/pN$.
\end{remark}

\begin{theorem}
\label{algebra-theorem-projectivity-characterization}
Soit $M$ un $R$-module. Alors $M$ est projectif si et seulement si
\begin{enumerate}
\item $M$ est plat,
\item $M$ est de Mittag-Leffler,
\item $M$ est une somme directe de $R$-modules dénombrablement engendrés.
\end{enumerate}
\end{theorem}

\begin{proof}
Supposons d'abord que $M$ soit projectif. Alors $M$ est un facteur direct d'un module
libre, de sorte que $M$ est plat et de Mittag-Leffler, puisque ces propriétés passent aux facteurs
directs. Par le théorème de Kaplansky (théorème \ref{algebra-theorem-projective-direct-sum}),
$M$ satisfait (3).

\medskip\noindent
Réciproquement, supposons que $M$ satisfasse (1)-(3). Puisque les propriétés d'être plat et de Mittag-Leffler
passent aux facteurs directs, $M$ est une somme directe de $R$-modules plats, de Mittag-Leffler et
dénombrablement engendrés.
Le lemme \ref{algebra-lemma-countgen-projective}
implique que $M$ est une somme directe de modules projectifs. Ainsi, $M$ est projectif.
\end{proof}

\begin{lemma}
\label{algebra-lemma-ML-ui-descent}
Soit $f: M \to N$ un morphisme universellement injectif de $R$-modules. Supposons que
$M$ soit une somme directe de $R$-modules dénombrablement engendrés, et supposons que $N$ soit plat
et de Mittag-Leffler. Alors $M$ est projectif.
\end{lemma}

\begin{proof}
Par les
lemmes \ref{algebra-lemma-ui-flat-domain} et
\ref{algebra-lemma-pure-submodule-ML},
$M$ est plat et de Mittag-Leffler ; la conclusion résulte donc du théorème
\ref{algebra-theorem-projectivity-characterization}.
\end{proof}

\begin{lemma}
\label{algebra-lemma-universally-injective-submodule-powerseries}
Soient $R$ un anneau noethérien et $M$ un $R$-module. Supposons que $M$ soit une
somme directe de $R$-modules dénombrablement engendrés, et supposons qu'il existe un
morphisme universellement injectif $M \to R[[t_1, \ldots, t_n]]$ pour un certain $n$.
Alors $M$ est projectif.
\end{lemma}

\begin{proof}
Cela résulte des
lemmes \ref{algebra-lemma-ML-ui-descent} et
\ref{algebra-lemma-power-series-ML}.
\end{proof}



\section{Montée de propriétés des modules}
\label{algebra-section-ascent}

\noindent
Toutes les propriétés d'un module figurant dans le théorème
\ref{algebra-theorem-projectivity-characterization} montent le long d'homomorphismes d'anneaux quelconques :

\begin{lemma}
\label{algebra-lemma-ascend-properties-modules}
Soit $R \to S$ un homomorphisme d'anneaux. Soit $M$ un $R$-module. Alors :
\begin{enumerate}
\item Si $M$ est plat, alors le $S$-module $M \otimes_R S$ est plat.
\item Si $M$ est de Mittag-Leffler, alors le $S$-module $M \otimes_R S$ est
de Mittag-Leffler.
\item Si $M$ est une somme directe de $R$-modules dénombrablement engendrés, alors le
$S$-module $M \otimes_R S$ est une somme directe de $S$-modules dénombrablement engendrés.
\item Si $M$ est projectif, alors le $S$-module $M \otimes_R S$ est projectif.
\end{enumerate}
\end{lemma}

\begin{proof}
Toutes ces assertions sont évidentes, sauf (2). Pour celle-ci, utilisons la formulation (3) de la propriété
de Mittag-Leffler donnée dans la proposition \ref{algebra-proposition-ML-characterization} ainsi que le
fait que la tensorisation commute aux colimites. On peut aussi
utiliser la caractérisation de la proposition \ref{algebra-proposition-ML-tensor}.
\end{proof}



\section{Descente de propriétés des modules}
\label{algebra-section-descent}

\noindent
Nous abordons la descente fidèlement plate des propriétés du théorème
\ref{algebra-theorem-projectivity-characterization} qui caractérisent la projectivité.
En présence de la platitude, la propriété d'être un module de Mittag-Leffler
descend :

\begin{lemma}
\label{algebra-lemma-ffdescent-ML}
\begin{reference}
Courriel de Juan Pablo Acosta Lopez daté du 20 décembre 2014.
\end{reference}
Soit $R \to S$ un homomorphisme d'anneaux fidèlement plat. Soit $M$ un $R$-module. Si le
$S$-module $M \otimes_R S$ est de Mittag-Leffler, alors $M$ est de Mittag-Leffler.
\end{lemma}

\begin{proof}
Écrivons $M = \colim_{i\in I} M_i$ comme colimite filtrante
de $R$-modules de présentation finie $M_i$.
En utilisant la proposition \ref{algebra-proposition-ML-characterization}, nous voyons que nous
devons montrer que, pour tout $i \in I$, il existe $i \leq j$, $j\in I$,
tel que $M_i\rightarrow M_j$ domine $M_i\rightarrow M$.

\medskip\noindent
Soit $N$ la somme amalgamée
$$
\xymatrix{
M_i  \ar[r] \ar[d] & M_j \ar[d] \\
M \ar[r] & N
}
$$
Le lemme équivaut alors à l'existence de $j$ tel que
$M_j\rightarrow N$ soit universellement injectif, voir le
lemme \ref{algebra-lemma-domination-universally-injective}.
Observons que sa tensorisation par $S$,
$$
\xymatrix{
M_i\otimes_R S  \ar[r] \ar[d] & M_j\otimes_R S \ar[d] \\
M\otimes_R S \ar[r] & N\otimes_R S
}
$$
est un diagramme de somme amalgamée. Ainsi, puisque
$M \otimes_R S = \colim_{i\in I} M_i \otimes_R S$
exprime $M\otimes_R S$ comme une colimite de $S$-modules de
présentation finie, et puisque $M\otimes_R S$ est de Mittag-Leffler,
il existe $j \geq i$ tel que $M_j\otimes_R S\rightarrow N\otimes_R S$
soit universellement injectif. Comme $R\rightarrow S$ est fidèlement plat,
nous concluons que $M_j\rightarrow N$ est lui aussi universellement injectif.
\end{proof}

\begin{lemma}
\label{algebra-lemma-ffdescent-countable}
Soit $R \to S$ un homomorphisme d'anneaux fidèlement plat. Soit $M$ un $R$-module. Si
le $S$-module $M \otimes_R S$ est dénombrablement engendré, alors $M$
est dénombrablement engendré.
\end{lemma}

\begin{proof}
Supposons que $M \otimes_R S$ soit engendré par les éléments
$y_i$, $i = 1, 2, 3, \ldots$. Écrivons $y_i = \sum_{j = 1, \ldots, n_i}
x_{ij} \otimes s_{ij}$ pour certains $n_i \geq 0$, $x_{ij} \in M$
et $s_{ij} \in S$. Notons $M' \subset M$ le sous-module engendré
par la famille dénombrable des éléments $x_{ij}$. Alors
$M' \otimes_R S \to M \otimes_R S$ est surjectif puisque son
image contient les générateurs $y_i$. Comme $S$ est fidèlement plat
sur $R$, nous concluons que $M' = M$, comme voulu.
\end{proof}

\noindent
À ce stade, la descente fidèlement plate des modules projectifs
dénombrablement engendrés s'obtient aisément.

\begin{lemma}
\label{algebra-lemma-ffdescent-countable-projectivity}
Soit $R \to S$ un homomorphisme d'anneaux fidèlement plat. Soit $M$ un $R$-module.
Si le $S$-module $M \otimes_R S$ est dénombrablement engendré et projectif,
alors $M$ est dénombrablement engendré et projectif.
\end{lemma}

\begin{proof}
Cela résulte des lemmes \ref{algebra-lemma-descend-properties-modules},
\ref{algebra-lemma-ffdescent-ML} et \ref{algebra-lemma-ffdescent-countable}, et du
théorème \ref{algebra-theorem-projectivity-characterization}.
\end{proof}

\noindent
Il ne reste qu'à utiliser un dévissage pour ramener la descente de la projectivité dans le
cas général au cas dénombrablement engendré. Commençons par deux lemmes simples.

\begin{lemma}
\label{algebra-lemma-lift-countably-generated-submodule}
Soient $R \to S$ un homomorphisme d'anneaux, $M$ un $R$-module et $Q$ un
$S$-sous-module dénombrablement engendré de $M \otimes_R S$. Alors il existe un
$R$-sous-module dénombrablement engendré $P$ de $M$ tel que
$\Im(P \otimes_R S \to M \otimes_R S)$ contienne $Q$.
\end{lemma}

\begin{proof}
Soient $y_1, y_2, \ldots$ des générateurs de $Q$ et écrivons $y_j = \sum_k x_{jk}
\otimes s_{jk}$ pour certains $x_{jk} \in M$ et $s_{jk} \in S$. Prenons alors pour $P$
le sous-module de $M$ engendré par les $x_{jk}$.
\end{proof}

\begin{lemma}
\label{algebra-lemma-adapted-submodule}
Soient $R \to S$ un homomorphisme d'anneaux et $M$ un $R$-module. Supposons que $M
\otimes_R S = \bigoplus_{i \in I} Q_i$ soit une somme directe de
$S$-modules dénombrablement engendrés $Q_i$. Si $N$ est un sous-module dénombrablement engendré de $M$, alors
il existe un sous-module dénombrablement engendré $N'$ de $M$ tel que $N' \supset N$
et $\Im(N' \otimes_R S \to M \otimes_R S) =
\bigoplus_{i \in I'} Q_i$ pour une certaine partie $I' \subset I$.
\end{lemma}

\begin{proof}
Soit $N'_0 = N$. Nous construisons par récurrence une suite croissante de sous-modules
dénombrablement engendrés $N'_{\ell} \subset M$ pour $\ell = 0, 1, 2, \ldots$
telle que : si $I'_{\ell}$ est l'ensemble des $i \in I$ tels que la projection de
$\Im(N'_{\ell} \otimes_R S \to M \otimes_R S)$ sur $Q_i$ est
non nulle, alors $\Im(N'_{\ell + 1} \otimes_R S \to M \otimes_R
S)$ contient $Q_i$ pour tout $i \in I'_{\ell}$. Pour construire $N'_{\ell + 1}$
à partir de $N'_\ell$, soit $Q$ la somme des $Q_i$ (en nombre dénombrable) pour
$i \in I'_{\ell}$, choisissons $P$ comme dans le lemme
\ref{algebra-lemma-lift-countably-generated-submodule}, puis posons $N'_{\ell + 1} =
N'_{\ell} + P$. Une fois les $N'_{\ell}$ construits, prenons simplement $N' =
\bigcup_{\ell} N'_{\ell}$ et $I' = \bigcup_{\ell} I'_{\ell}$.
\end{proof}

\begin{theorem}
\label{algebra-theorem-ffdescent-projectivity}
Soit $R \to S$ un homomorphisme d'anneaux fidèlement plat. Soit $M$ un $R$-module.
Si le $S$-module $M \otimes_R S$ est projectif, alors $M$ est projectif.
\end{theorem}

\begin{proof}
Nous allons construire un dévissage de Kaplansky de $M$ pour montrer que c'est une
somme directe de modules projectifs, et donc qu'il est projectif. Par le théorème
\ref{algebra-theorem-projective-direct-sum}, nous pouvons écrire $M \otimes_R S =
\bigoplus_{i \in I} Q_i$ comme une somme directe de $S$-modules dénombrablement engendrés
$Q_i$. Choisissons un bon ordre sur $M$. Par récurrence transfinie, nous allons
définir une famille croissante de sous-modules $M_{\alpha}$ de $M$, un pour chaque
ordinal $\alpha$, telle que $M_{\alpha} \otimes_R S$ soit une somme directe d'une certaine
partie des $Q_i$.

\medskip\noindent
Pour $\alpha = 0$, posons $M_0 = 0$. Si $\alpha$ est un ordinal limite et si $M_{\beta}$
a été défini pour tout $\beta < \alpha$, définissons alors $M_\alpha =
\bigcup_{\beta < \alpha} M_{\beta}$. Puisque chaque $M_{\beta} \otimes_R S$, pour
$\beta < \alpha$, est une somme directe d'une partie des $Q_i$, il en va de même de
$M_{\alpha} \otimes_R S$. Si $\alpha + 1$ est un ordinal successeur et si
$M_{\alpha}$ a été défini, définissons $M_{\alpha + 1}$ comme suit. Si
$M_{\alpha} = M$, posons $M_{\alpha +1} = M$. Sinon, choisissons le plus petit
$x \in M$ (pour le bon ordre fixé) tel que $x \notin
M_{\alpha}$. Puisque $S$ est plat sur $R$, $(M/M_{\alpha}) \otimes_R S = M
\otimes_R S/M_{\alpha} \otimes_R S$ ; comme $M_{\alpha} \otimes_R S$ est
une somme directe de certains $Q_i$, il en va de même de $(M/M_{\alpha}) \otimes_R
S$. Par le lemme \ref{algebra-lemma-adapted-submodule}, nous pouvons trouver un
$R$-sous-module dénombrablement engendré $P$ de $M/M_{\alpha}$ contenant l'image de $x$ dans
$M/M_{\alpha}$ et tel que $P \otimes_R S$ (qui est égal à $\Im(P
\otimes_R S \to M \otimes_R S)$ puisque $S$ est plat sur $R$) soit une
somme directe de certains $Q_i$. Puisque $M \otimes_R S = \bigoplus_{i \in I} Q_i$
est projectif et que la projectivité passe aux facteurs directs, $P \otimes_R S$ est
également projectif. Ainsi, par le lemme \ref{algebra-lemma-ffdescent-countable-projectivity},
$P$ est projectif. Enfin, nous définissons $M_{\alpha + 1}$ comme l'image réciproque de $P$
dans $M$, de sorte que $M_{\alpha + 1}/M_{\alpha} = P$ est dénombrablement engendré et
projectif. En particulier, $M_{\alpha}$ est un facteur direct de $M_{\alpha + 1}$,
car la projectivité de $M_{\alpha + 1}/M_{\alpha}$ implique que la suite $0
\to M_{\alpha} \to M_{\alpha + 1} \to
M_{\alpha + 1}/M_{\alpha} \to 0$ est scindée.

\medskip\noindent
Une récurrence transfinie sur $M$ (en utilisant le fait que nous avons construit
$M_{\alpha + 1}$ de façon qu'il contienne le plus petit $x \in M$ non contenu dans
$M_{\alpha}$) montre que
chaque $x \in M$ est contenu dans un certain $M_{\alpha}$. Ainsi, il existe un ordinal $S$
assez grand satisfaisant : pour chaque $x \in M$, il existe $\alpha \in S$ tel
que $x \in M_{\alpha}$. Cela signifie que $(M_{\alpha})_{\alpha \in S}$ satisfait
la propriété (1) d'un dévissage de Kaplansky de $M$. Les autres propriétés sont claires
par construction. Nous concluons que
$M = \bigoplus_{\alpha + 1 \in S} M_{\alpha + 1}/M_{\alpha}$.
Puisque chaque $M_{\alpha + 1}/M_{\alpha}$ est projectif
par construction, $M$ est projectif.
\end{proof}












\section{Complétion}
\label{algebra-section-completion}

\noindent
Supposons que $R$ soit un anneau et que $I$ soit un idéal.
Nous définissons le {\it complété de $R$ relativement à $I$}
comme la limite
$$
R^\wedge = \lim_n R/I^n.
$$
Un élément de $R^\wedge$ est donné par une suite
d'éléments $f_n \in R/I^n$ tels que $f_n \equiv f_{n + 1} \bmod I^n$
pour tout $n$. Nous considérerons $R^\wedge$ comme une $R$-algèbre.
De même, si $M$ est un $R$-module, nous définissons le
{\it complété de $M$ relativement à $I$}
comme la limite
$$
M^\wedge = \lim_n M/I^nM.
$$
Un élément de $M^\wedge$ est donné par une suite
d'éléments $m_n \in M/I^nM$ tels que $m_n \equiv m_{n + 1} \bmod I^nM$
pour tout $n$. Nous considérerons $M^\wedge$ comme un $R^\wedge$-module.
Cette description montre clairement qu'il existe
toujours des morphismes canoniques
$$
M \longrightarrow M^\wedge
\quad\text{et}\quad
M \otimes_R R^\wedge \longrightarrow M^\wedge.
$$
De plus, étant donné un morphisme de modules $\varphi : M \to N$, nous obtenons sur les complétés un
morphisme induit $\varphi^\wedge : M^\wedge \to N^\wedge$ qui rend le
diagramme
$$
\xymatrix{
M \ar[r] \ar[d] & N \ar[d] \\
M^\wedge \ar[r] & N^\wedge
}
$$
commutatif. En général, la complétion n'est pas un foncteur exact, voir
Exemples, section \ref{examples-section-completion-not-exact}.
Voici quelques premiers résultats positifs.

\begin{lemma}
\label{algebra-lemma-completion-generalities}
Soit $R$ un anneau. Soit $I \subset R$ un idéal.
Soit $\varphi : M \to N$ un morphisme de $R$-modules.
\begin{enumerate}
\item Si $M/IM \to N/IN$ est surjectif, alors $M^\wedge \to N^\wedge$
est surjectif.
\item Si $M \to N$ est surjectif, alors $M^\wedge \to N^\wedge$ est surjectif.
\item Si $0 \to K \to M \to N \to 0$ est une suite exacte courte de
$R$-modules et si $N$ est plat, alors
$0 \to K^\wedge \to M^\wedge \to N^\wedge \to 0$ est une suite exacte courte.
\item Le morphisme $M \otimes_R R^\wedge \to M^\wedge$ est
surjectif pour tout $R$-module de type fini $M$.
\end{enumerate}
\end{lemma}

\begin{proof}
Supposons que $M/IM \to N/IN$ soit surjectif. Alors le morphisme $M/I^nM \to N/I^nN$
est surjectif pour tout $n \geq 1$ par le lemme de Nakayama. Plus précisément,
appliquons la partie (11) du lemme \ref{algebra-lemma-NAK} au
morphisme $M/I^nM \to N/I^nN$ sur l'anneau $R/I^n$ et à l'idéal nilpotent
$I/I^n$. Posons $K_n = \{x \in M \mid \varphi(x) \in I^nN\}$.
Nous obtenons ainsi des suites exactes courtes
$$
0 \to K_n/I^nM \to M/I^nM \to N/I^nN \to 0
$$
Nous affirmons que le morphisme canonique $K_{n + 1}/I^{n + 1}M \to K_n/I^nM$
est surjectif. En effet, si $x \in K_n$, écrivons
$\varphi(x) = \sum z_j n_j$ avec $z_j \in I^n$, $n_j \in N$.
Par hypothèse, nous pouvons écrire $n_j = \varphi(m_j) + \sum z_{jk}n_{jk}$
avec $m_j \in M$, $z_{jk} \in I$ et $n_{jk} \in N$. Par conséquent,
$$
\varphi(x - \sum z_j m_j) = \sum z_jz_{jk} n_{jk}.
$$
Cela signifie que $x' = x - \sum z_j m_j \in K_{n + 1}$ s'envoie
sur $x \bmod I^nM$, ce qui prouve l'assertion. Nous pouvons maintenant appliquer le
lemme \ref{algebra-lemma-Mittag-Leffler}
au système projectif de suites exactes courtes ci-dessus pour obtenir (1).
La partie (2) est un cas particulier de (1).
Si les hypothèses de (3) sont satisfaites, alors, pour tout $n$, la suite
$$
0 \to K/I^nK \to M/I^nM \to N/I^nN \to 0
$$
est exacte courte par le
lemme \ref{algebra-lemma-flat-tor-zero}.
Nous pouvons donc appliquer directement le
lemme \ref{algebra-lemma-Mittag-Leffler}
pour conclure que (3) est vraie.
Pour voir (4), choisissons des générateurs $x_i \in M$, $i = 1, \ldots, n$.
Alors le morphisme $R^{\oplus n} \to M$, $(a_1, \ldots, a_n) \mapsto \sum a_ix_i$
est surjectif. Par (2), nous voyons donc que
$(R^\wedge)^{\oplus n} \to M^\wedge$, $(a_1, \ldots, a_n) \mapsto \sum a_ix_i$
est surjectif. L'assertion (4) en résulte.
\end{proof}

\begin{definition}
\label{algebra-definition-complete}
Soit $R$ un anneau. Soit $I \subset R$ un idéal.
Soit $M$ un $R$-module. Nous disons que $M$ est {\it $I$-adiquement complet}
si le morphisme
$$
M \longrightarrow M^\wedge = \lim_n M/I^nM
$$
est un isomorphisme\footnote{Cela comprend la condition
$\bigcap I^nM = 0$.}. Nous disons que $R$ est {\it $I$-adiquement complet}
si $R$ est $I$-adiquement complet comme $R$-module.
\end{definition}

\noindent
Il n'est pas vrai que le complété d'un $R$-module $M$ relativement
à $I$ soit $I$-adiquement complet. Pour un exemple, voir
Exemples, section \ref{examples-section-noncomplete-completion}.
Si l'idéal est de type fini, alors le complété est complet.

\begin{lemma}
\label{algebra-lemma-hathat-finitely-generated}
\begin{reference}
\cite[Theorem 15]{Matlis}. La démonstration élégante donnée ici provient
d'un courriel de Bjorn Poonen daté du 5 novembre 2016.
\end{reference}
Soit $R$ un anneau. Soit $I$ un idéal de type fini de $R$.
Soit $M$ un $R$-module. Alors
\begin{enumerate}
\item le complété $M^\wedge$ est $I$-adiquement complet, et
\item $I^nM^\wedge = \Ker(M^\wedge \to M/I^nM) = (I^nM)^\wedge$ pour tout
$n \geq 1$.
\end{enumerate}
En particulier, $R^\wedge$ est $I$-adiquement complet,
$I^nR^\wedge = (I^n)^\wedge$, et
$R^\wedge/I^nR^\wedge = R/I^n$.
\end{lemma}

\begin{proof}
Puisque $I$ est de type fini,
$I^n$ est de type fini, disons engendré par $f_1, \ldots, f_r$.
L'application de la partie (2) du lemme \ref{algebra-lemma-completion-generalities}
à la surjection $(f_1, \ldots, f_r) : M^{\oplus r} \to I^n M$
donne une surjection
$$
(M^\wedge)^{\oplus r} \xrightarrow{(f_1, \ldots, f_r)} (I^n M)^\wedge =
\lim_{m \geq n} I^n M/I^m M = \Ker(M^\wedge \to M/I^n M).
$$
D'autre part, l'image de
$(f_1, \ldots, f_r) : (M^\wedge)^{\oplus r} \to M^\wedge$
est $I^n M^\wedge$.
Ainsi, $M^\wedge / I^n M^\wedge \simeq M/I^n M$.
En prenant les limites projectives, nous obtenons $(M^\wedge)^\wedge \simeq M^\wedge$ ;
autrement dit, $M^\wedge$ est $I$-adiquement complet.
\end{proof}

\begin{lemma}
\label{algebra-lemma-completion-differ-by-torsion}
Soit $R$ un anneau. Soit $I \subset R$ un idéal. Soit
$0 \to M \to N \to Q \to 0$ une suite exacte de
$R$-modules telle que $Q$ soit annulé par une puissance de $I$.
Alors la complétion fournit une suite exacte
$0 \to M^\wedge \to N^\wedge \to Q \to 0$.
\end{lemma}

\begin{proof}
Disons que $I^c Q = 0$. Alors $Q/I^nQ = Q$ pour $n \geq c$.
D'autre part, il est clair que
$I^nM \subset M \cap I^nN \subset I^{n - c}M$ pour $n \geq c$.
Ainsi, $M^\wedge = \lim M/(M \cap I^n N)$. Appliquons le lemme \ref{algebra-lemma-Mittag-Leffler}
au système de suites exactes
$$
0 \to M/(M \cap I^n N) \to N/I^n N \to Q \to 0
$$
pour $n \geq c$ afin de conclure.
\end{proof}

\begin{lemma}
\label{algebra-lemma-hathat}
\begin{reference}
Tiré d'une note non publiée de Lenstra et de Smit.
\end{reference}
Soit $R$ un anneau. Soit $I \subset R$ un idéal. Soit $M$ un $R$-module.
Notons $K_n = \Ker(M^\wedge \to M/I^nM)$. Alors $M^\wedge$ est $I$-adiquement
complet si et seulement si $K_n$ est égal à $I^nM^\wedge$ pour tout $n \geq 1$.
\end{lemma}

\begin{proof}
Le module $I^n M^\wedge$ est contenu dans $K_n$.
Ainsi, pour tout $n \geq 1$, il existe une suite exacte canonique
$$
0 \to K_n/I^nM^\wedge \to M^\wedge/I^nM^\wedge \to M/I^nM \to 0.
$$
Comme $I^nM^\wedge$ se surjecte sur $I^nM/I^{n + 1}M$, nous voyons que
$K_{n + 1} + I^n M^\wedge = K_n$. Ainsi, le système projectif
$\{K_n/I^n M^\wedge\}_{n \geq 1}$ a des morphismes de transition surjectifs.
Par le
lemme \ref{algebra-lemma-Mittag-Leffler},
nous voyons qu'il existe une suite exacte courte
$$
0 \to
\lim_n K_n/I^n M^\wedge \to
(M^\wedge)^\wedge \to
M^\wedge \to 0
$$
Ainsi, $M^\wedge$ est complet si et seulement si $K_n/I^n M^\wedge = 0$
pour tout $n \geq 1$.
\end{proof}

\begin{lemma}
\label{algebra-lemma-radical-completion}
Soit $R$ un anneau, soit $I \subset R$ un idéal et soit
$R^\wedge = \lim R/I^n$.
\begin{enumerate}
\item tout élément de $R^\wedge$ qui s'envoie sur une unité de $R/I$ est une unité,
\item tout élément de $1 + I$ s'envoie sur un élément inversible de $R^\wedge$,
\item tout élément de $1 + IR^\wedge$ est inversible dans $R^\wedge$, et
\item les idéaux $IR^\wedge$ et $\Ker(R^\wedge \to R/I)$ sont contenus
dans le radical de Jacobson de $R^\wedge$.
\end{enumerate}
\end{lemma}

\begin{proof}
Soit $x \in R^\wedge$ qui s'envoie sur une unité $x_1$ dans $R/I$.
Alors $x$ s'envoie sur une unité $x_n$ dans $R/I^n$ pour tout $n$, par le
lemme \ref{algebra-lemma-locally-nilpotent-unit}.
Ainsi, $y = (x_n^{-1}) \in \lim R/I^n = R^\wedge$ est un inverse de $x$.
Les parties (2) et (3) résultent immédiatement de (1).
La partie (4) résulte de (1) et du lemme \ref{algebra-lemma-contained-in-radical}.
\end{proof}

\begin{lemma}
\label{algebra-lemma-when-surjective-to-completion}
Soit $A$ un anneau. Soit $I = (f_1, \ldots, f_r)$ un idéal de type
fini. Si $M \to \lim M/f_i^nM$ est surjectif pour
chaque $i$, alors $M \to \lim M/I^nM$ est surjectif.
\end{lemma}

\begin{proof}
Remarquons que $\lim M/I^nM = \lim M/(f_1^n, \ldots, f_r^n)M$, puisque
$I^n \supset (f_1^n, \ldots, f_r^n) \supset I^{rn}$.
Un élément $\xi$ de $\lim M/(f_1^n, \ldots, f_r^n)M$ peut s'écrire formellement
sous la forme
$$
\xi = \sum\nolimits_{n \geq 0} \sum\nolimits_i f_i^n x_{n, i}
$$
avec $x_{n, i} \in M$. Si $M \to \lim M/f_i^nM$ est surjectif, alors il existe
un $x_i \in M$ qui s'envoie sur $\sum x_{n, i} f_i^n$ dans $\lim M/f_i^nM$.
Alors $x = \sum x_i$ s'envoie sur $\xi$ dans $\lim M/I^nM$.
\end{proof}

\begin{lemma}
\label{algebra-lemma-complete-by-sub}
Soit $A$ un anneau. Soient $I \subset J \subset A$ des idéaux.
Si $M$ est $J$-adiquement complet et si $I$ est de type fini, alors
$M$ est $I$-adiquement complet.
\end{lemma}

\begin{proof}
Supposons que $M$ soit $J$-adiquement complet et que $I$ soit de type fini.
Nous avons $\bigcap I^nM = 0$ parce que $\bigcap J^nM = 0$. Par le
lemme \ref{algebra-lemma-when-surjective-to-completion},
il suffit de démontrer la surjectivité de $M \to \lim M/I^nM$ dans le cas où
$I$ est engendré par un seul élément. Disons $I = (f)$.
Soient $x_n \in M$ tels que $x_{n + 1} - x_n \in f^nM$. Nous devons montrer qu'il existe
un $x \in M$ tel que $x_n - x \in f^nM$ pour tout $n$.
Comme $x_{n + 1} - x_n \in J^nM$ et comme $M$ est $J$-adiquement complet,
il existe un élément $x \in M$ tel que $x_n - x \in J^nM$.
En remplaçant $x_n$ par $x_n - x$, nous pouvons supposer que $x_n \in J^nM$.
Pour terminer la démonstration, nous allons montrer que cela implique $x_n \in I^nM$.
En effet, écrivons $x_n - x_{n + 1} = f^nz_n$.
Alors
$$
x_n = f^n(z_n + fz_{n + 1} + f^2z_{n + 2} + \ldots)
$$
La somme $z_n + fz_{n + 1} + f^2z_{n + 2} + \ldots$ converge dans $M$
puisque $f^c \in J^c$. La somme $f^n(z_n + fz_{n + 1} + f^2z_{n + 2} + \ldots)$
converge dans $M$ vers $x_n$, car
les sommes partielles sont égales à $x_n - x_{n + c}$ et $x_{n + c} \in J^{n + c}M$.
\end{proof}

\begin{lemma}
\label{algebra-lemma-change-ideal-completion}
Soit $R$ un anneau.
Soient $I$, $J$ des idéaux de $R$.
Supposons qu'il existe des entiers $c, d > 0$ tels que
$I^c \subset J$ et $J^d \subset I$.
Alors la complétion relativement à $I$ s'identifie à la complétion
relativement à $J$ pour tout $R$-module.
En particulier, un $R$-module $M$ est $I$-adiquement complet
si et seulement s'il est $J$-adiquement complet.
\end{lemma}

\begin{proof}
Considérons le système de morphismes
$M/I^nM \to M/J^{\lfloor n/c \rfloor}M$ et
le système de morphismes $M/J^mM \to M/I^{\lfloor m/d \rfloor}M$
pour obtenir des morphismes mutuellement inverses entre les complétés.
\end{proof}

\begin{lemma}
\label{algebra-lemma-quotient-complete}
Soit $R$ un anneau. Soit $I$ un idéal de $R$.
Soit $M$ un module $I$-adiquement complet sur $R$,
et soit $K \subset M$ un $R$-sous-module.
Les assertions suivantes sont équivalentes :
\begin{enumerate}
\item $K = \bigcap (K + I^nM)$ et
\item $M/K$ est $I$-adiquement complet.
\end{enumerate}
\end{lemma}

\begin{proof}
Posons $N = M/K$. Par le
lemme \ref{algebra-lemma-completion-generalities},
le morphisme $M = M^\wedge \to N^\wedge$ est surjectif.
Ainsi, $N \to N^\wedge$ est surjectif. Il est facile de voir que le
noyau de $N \to N^\wedge$ est le module $\bigcap (K + I^nM) / K$.
\end{proof}

\begin{lemma}
\label{algebra-lemma-when-finite-module-complete-over-complete-ring}
Soit $R$ un anneau. Soit $I$ un idéal de $R$.
Soit $M$ un $R$-module.
Si (a) $R$ est $I$-adiquement complet, (b) $M$ est un $R$-module de type fini,
et (c) $\bigcap I^nM = (0)$, alors $M$ est $I$-adiquement complet.
\end{lemma}

\begin{proof}
Par le lemme \ref{algebra-lemma-completion-generalities},
le morphisme $M = M \otimes_R R = M \otimes_R R^\wedge \to M^\wedge$
est surjectif. Le noyau de ce morphisme est $\bigcap I^nM$, donc nul
par hypothèse. Ainsi, $M \cong M^\wedge$ et $M$ est complet.
\end{proof}

\begin{lemma}
\label{algebra-lemma-finite-over-complete-ring}
Soit $R$ un anneau. Soit $I \subset R$ un idéal. Soit $M$ un $R$-module.
Supposons que
\begin{enumerate}
\item $R$ soit $I$-adiquement complet,
\item $\bigcap_{n \geq 1} I^nM = (0)$, et
\item $M/IM$ soit un $R/I$-module de type fini.
\end{enumerate}
Alors $M$ est un $R$-module de type fini.
\end{lemma}

\begin{proof}
Soient $x_1, \ldots, x_n \in M$ des éléments dont les images dans $M/IM$ engendrent
$M/IM$ comme $R/I$-module. Notons $M' \subset M$ le $R$-sous-module
engendré par $x_1, \ldots, x_n$. Par le lemme \ref{algebra-lemma-completion-generalities},
le morphisme $(M')^\wedge \to M^\wedge$ est surjectif.
Puisque $\bigcap I^nM = 0$, nous voyons en particulier que $\bigcap I^nM' = (0)$.
Par le lemme \ref{algebra-lemma-when-finite-module-complete-over-complete-ring},
nous voyons donc que $M'$ est complet, et nous concluons que $M' \to M^\wedge$
est surjectif. Enfin, le noyau de $M \to M^\wedge$ est
nul, puisqu'il est égal à $\bigcap I^nM = (0)$.
Nous concluons donc que $M \cong M' \cong M^\wedge$
est de type fini.
\end{proof}









\section{Complétion pour les anneaux noethériens}
\label{algebra-section-completion-noetherian}

\noindent
Dans cette section, nous étudions la complétion relativement à des idéaux dans les
anneaux noethériens.

\begin{lemma}
\label{algebra-lemma-completion-tensor}
Soit $I$ un idéal d'un anneau noethérien $R$.
Notons ${}^\wedge$ la complétion relativement à $I$.
\begin{enumerate}
\item Si $K \to N$ est un morphisme injectif de $R$-modules de type fini,
alors le morphisme induit sur les complétés $K^\wedge \to N^\wedge$ est injectif.
\item Si $0 \to K \to N \to M \to 0$ est une suite exacte courte
de $R$-modules de type fini, alors $0 \to K^\wedge \to N^\wedge \to M^\wedge \to 0$
est une suite exacte courte.
\item Si $M$ est un $R$-module de type fini, alors $M^\wedge = M \otimes_R R^\wedge$.
\end{enumerate}
\end{lemma}

\begin{proof}
En posant $M = N/K$, nous constatons que la partie (1) résulte de la partie (2).
Soit $0 \to K \to N \to M \to 0$ comme dans (2).
Pour tout $n$, nous obtenons la suite exacte courte
$$
0 \to K/(I^nN \cap K) \to N/I^nN \to M/I^nM \to 0.
$$
Par le lemme \ref{algebra-lemma-Mittag-Leffler},
nous obtenons la suite exacte
$$
0 \to \lim K/(I^nN \cap K) \to N^\wedge \to M^\wedge \to 0.
$$
Par le lemme d'Artin-Rees \ref{algebra-lemma-Artin-Rees}, nous pouvons choisir $c$ tel que
$I^nK \subset I^n N \cap K \subset I^{n-c} K$ pour $n \geq c$.
Ainsi, $K^\wedge = \lim K/I^nK = \lim K/(I^nN \cap K)$,
et nous concluons que (2) est vraie.

\medskip\noindent
Soit $M$ comme dans (3), et soit $0 \to K \to R^{\oplus t} \to M \to 0$
une présentation de $M$. Nous obtenons un diagramme commutatif
$$
\xymatrix{
&
K \otimes_R R^\wedge \ar[r] \ar[d] &
R^{\oplus t} \otimes_R R^\wedge \ar[r] \ar[d] &
M \otimes_R R^\wedge \ar[r] \ar[d] &
0 \\
0 \ar[r] &
K^\wedge \ar[r] &
(R^{\oplus t})^\wedge \ar[r] &
M^\wedge \ar[r] & 0
}
$$
La ligne supérieure est exacte, voir la section \ref{algebra-section-flat}.
La ligne inférieure est exacte par la partie (2).
Par le lemme \ref{algebra-lemma-completion-generalities},
les flèches verticales sont surjectives.
La flèche verticale du milieu est un isomorphisme.
Nous concluons que (3) est satisfaite par le lemme du serpent \ref{algebra-lemma-snake}.
\end{proof}

\begin{lemma}
\label{algebra-lemma-completion-flat}
Soit $I$ un idéal d'un anneau noethérien $R$.
Notons ${}^\wedge$ la complétion relativement à $I$.
\begin{enumerate}
\item L'homomorphisme d'anneaux $R \to R^\wedge$ est plat.
\item Le foncteur $M \mapsto M^\wedge$ est exact sur la catégorie des
$R$-modules de type fini.
\end{enumerate}
\end{lemma}

\begin{proof}
Considérons $J \otimes_R R^\wedge \to R \otimes_R R^\wedge = R^\wedge$,
où $J$ est un idéal quelconque de $R$.
D'après le lemme \ref{algebra-lemma-completion-tensor}, ce morphisme
s'identifie à $J^\wedge \to R^\wedge$, et $J^\wedge \to R^\wedge$
est injectif. La partie (1) résulte du lemme \ref{algebra-lemma-flat}.
La partie (2) est une reformulation de la
partie (2) du lemme \ref{algebra-lemma-completion-tensor}.
\end{proof}

\begin{lemma}
\label{algebra-lemma-completion-faithfully-flat}
Soit $I$ un idéal d'un anneau noethérien $R$. Notons $R^\wedge$
le complété de $R$ relativement à $I$. Si $I$ est contenu
dans le radical de Jacobson de $R$, alors l'homomorphisme d'anneaux $R \to R^\wedge$
est fidèlement plat. En particulier, si $(R, \mathfrak m)$ est un anneau local
noethérien, alors le complété $\lim_n R/\mathfrak m^n$ est fidèlement plat.
\end{lemma}

\begin{proof}
Il est plat par le lemme \ref{algebra-lemma-completion-flat}.
Le composé $R \to R^\wedge \to R/I$, où
le dernier morphisme est la projection $R^\wedge \to R/I$, montre que
tout idéal maximal de $R$ appartient à l'image de $\Spec(R^\wedge)
\to \Spec(R)$. Ainsi, le morphisme est fidèlement
plat par le lemme \ref{algebra-lemma-ff}.
\end{proof}

\begin{lemma}
\label{algebra-lemma-completion-complete}
Soit $R$ un anneau noethérien.
Soit $I$ un idéal de $R$.
Soit $M$ un $R$-module.
Alors le complété $M^\wedge$
de $M$ relativement à $I$ est $I$-adiquement complet,
$I^n M^\wedge = (I^nM)^\wedge$, et $M^\wedge/I^nM^\wedge = M/I^nM$.
\end{lemma}

\begin{proof}
C'est un cas particulier du
lemme \ref{algebra-lemma-hathat-finitely-generated},
car $I$ est un idéal de type fini.
\end{proof}

\begin{lemma}
\label{algebra-lemma-completion-Noetherian}
Soit $I$ un idéal d'un anneau $R$. Supposons que
\begin{enumerate}
\item $R/I$ soit un anneau noethérien,
\item $I$ soit de type fini.
\end{enumerate}
Alors le complété $R^\wedge$ de $R$ relativement à $I$
est un anneau noethérien complet relativement à $IR^\wedge$.
\end{lemma}

\begin{proof}
Par le
lemme \ref{algebra-lemma-hathat-finitely-generated},
nous voyons que $R^\wedge$ est $I$-adiquement complet. Il est donc aussi
$IR^\wedge$-adiquement complet. Puisque $R^\wedge/IR^\wedge = R/I$ est
noethérien, nous voyons qu'après avoir remplacé $R$ par $R^\wedge$, nous pouvons,
en plus des hypothèses (1) et (2), supposer que $R$ est lui aussi $I$-adiquement
complet.

\medskip\noindent
Soient $f_1, \ldots, f_t$ des générateurs de $I$.
Il existe alors une surjection d'anneaux
$R/I[T_1, \ldots, T_t] \to \bigoplus I^n/I^{n + 1}$
qui envoie $T_i$ sur l'élément $\overline{f}_i \in I/I^2$.
Ainsi, $\bigoplus I^n/I^{n + 1}$ est un anneau noethérien.
Soit $J \subset R$ un idéal. Considérons l'idéal
$$
\bigoplus J \cap I^n/J \cap I^{n + 1} \subset \bigoplus I^n/I^{n + 1}.
$$
Soient $\overline{g}_1, \ldots, \overline{g}_m$ des générateurs de cet
idéal. Nous pouvons choisir $\overline{g}_j$ homogène
de degré $d_j$ et choisir $g_j \in J \cap I^{d_j}$ qui s'envoie sur
$\overline{g}_j \in J \cap I^{d_j}/J \cap I^{d_j + 1}$. Nous affirmons
que $g_1, \ldots, g_m$ engendrent $J$.

\medskip\noindent
Soit $x \in J \cap I^n$. Il existe $a_j \in I^{\max(0, n - d_j)}$ tels que
$x - \sum a_j g_j \in J \cap I^{n + 1}$.
La raison en est que $J \cap I^n/J \cap I^{n + 1}$ est égal à
$\sum \overline{g}_j I^{n - d_j}/I^{n - d_j + 1}$ par notre choix
de $g_1, \ldots, g_m$. Ainsi, en partant de $x \in J$, nous pouvons trouver
une suite de vecteurs $(a_{1, n}, \ldots, a_{m, n})_{n \geq 0}$
avec $a_{j, n} \in I^{\max(0, n - d_j)}$ telle que
$$
x =
\sum\nolimits_{n = 0, \ldots, N}
\sum\nolimits_{j = 1, \ldots, m} a_{j, n} g_j \bmod I^{N + 1}
$$
En posant $A_j = \sum_{n \geq 0} a_{j, n}$, nous voyons que
$x = \sum A_j g_j$, puisque $R$ est complet. Ainsi, $J$ est de type fini, et
le résultat est démontré.
\end{proof}

\begin{lemma}
\label{algebra-lemma-completion-Noetherian-Noetherian}
Soit $R$ un anneau noethérien.
Soit $I$ un idéal de $R$.
Le complété $R^\wedge$ de $R$ relativement à $I$ est
noethérien.
\end{lemma}

\begin{proof}
C'est une conséquence du
lemme \ref{algebra-lemma-completion-Noetherian}.
On peut aussi le voir directement comme suit.
Choisissons des générateurs $f_1, \ldots, f_n$ de $I$.
Considérons le morphisme
$$
R[[x_1, \ldots, x_n]] \longrightarrow R^\wedge,
\quad
x_i \longmapsto f_i.
$$
C'est un homomorphisme d'anneaux bien défini et surjectif
(détails omis).
Puisque $R[[x_1, \ldots, x_n]]$ est noethérien (voir le
lemme \ref{algebra-lemma-Noetherian-power-series}), le résultat est démontré.
\end{proof}

\noindent
Supposons que $R \to S$ soit un homomorphisme local d'anneaux locaux $(R, \mathfrak m)$
et $(S, \mathfrak n)$. Soit $S^\wedge$ le complété
de $S$ relativement à $\mathfrak n$. En général, $S^\wedge$ n'est pas
le complété $\mathfrak m$-adique de $S$. Si
$\mathfrak n^t \subset \mathfrak mS$ pour un certain $t \geq 1$,
alors nous avons bien $S^\wedge = \lim S/\mathfrak m^nS$ par le
lemme \ref{algebra-lemma-change-ideal-completion}. Dans certains cas, cela
implique même que $S^\wedge$ est de type fini sur $R^\wedge$.

\begin{lemma}
\label{algebra-lemma-finite-after-completion}
Soit $R \to S$ un homomorphisme local d'anneaux locaux $(R, \mathfrak m)$
et $(S, \mathfrak n)$. Soient $R^\wedge$, resp.\ $S^\wedge$, le complété
de $R$, resp.\ de $S$, relativement à $\mathfrak m$, resp.\ à $\mathfrak n$.
Si $\mathfrak m$ et $\mathfrak n$ sont de type fini et si
$\dim_{\kappa(\mathfrak m)} S/\mathfrak mS < \infty$, alors
\begin{enumerate}
\item $S^\wedge$ est égal au complété $\mathfrak m$-adique de $S$, et
\item $S^\wedge$ est un $R^\wedge$-module de type fini.
\end{enumerate}
\end{lemma}

\begin{proof}
Nous avons $\mathfrak mS \subset \mathfrak n$, puisque $R \to S$ est un homomorphisme
local.
L'hypothèse $\dim_{\kappa(\mathfrak m)} S/\mathfrak mS < \infty$
implique que $S/\mathfrak mS$ est un anneau artinien, voir le
lemme \ref{algebra-lemma-finite-dimensional-algebra}.
Il est donc de dimension $0$, voir le
lemme \ref{algebra-lemma-Noetherian-dimension-0},
et par conséquent $\mathfrak n = \sqrt{\mathfrak mS}$.
Cela et le fait que $\mathfrak n$ est de type fini
impliquent que $\mathfrak n^t \subset \mathfrak mS$ pour
un certain $t \geq 1$. Par le
lemme \ref{algebra-lemma-change-ideal-completion},
nous voyons que $S^\wedge$ peut être identifié au complété
$\mathfrak m$-adique de $S$. Comme $\mathfrak m$ est de type fini, il résulte du
lemme \ref{algebra-lemma-hathat-finitely-generated}
que $S^\wedge$ et $R^\wedge$ sont $\mathfrak m$-adiquement complets.
Nous pouvons alors appliquer le
lemme \ref{algebra-lemma-finite-over-complete-ring}
à $S^\wedge$ comme $R^\wedge$-module pour conclure.
\end{proof}

\begin{lemma}
\label{algebra-lemma-completion-finite-extension}
Soit $R$ un anneau noethérien. Soit $R \to S$ un homomorphisme d'anneaux fini.
Soit $\mathfrak p \subset R$ un idéal premier et soient
$\mathfrak q_1, \ldots, \mathfrak q_m$ les idéaux premiers de $S$
au-dessus de $\mathfrak p$
(lemme \ref{algebra-lemma-finite-finite-fibres}).
Alors
$$
R_\mathfrak p^\wedge \otimes_R S =
(S_\mathfrak p)^\wedge =
S_{\mathfrak q_1}^\wedge \times \ldots \times S_{\mathfrak q_m}^\wedge
$$
où $(S_\mathfrak p)^\wedge$ est le complété relativement à
$\mathfrak p$, et où les anneaux locaux $R_\mathfrak p$ et
$S_{\mathfrak q_i}$ sont complétés relativement à leurs idéaux maximaux.
\end{lemma}

\begin{proof}
Nous pouvons remplacer $R$ par le localisé $R_\mathfrak p$ et
$S$ par $S_\mathfrak p = S \otimes_R R_\mathfrak p$.
Nous pouvons donc supposer que $R$ est un anneau local noethérien et
que $\mathfrak p = \mathfrak m$ est son idéal maximal.
Le complété $\mathfrak q_iS_{\mathfrak q_i}$-adique
$S_{\mathfrak q_i}^\wedge$ est égal au complété $\mathfrak m$-adique
par le lemme \ref{algebra-lemma-finite-after-completion}.
Pour tout $n \geq 1$, les idéaux premiers de $S/\mathfrak m^nS$ sont en
correspondance bijective avec les idéaux maximaux
$\mathfrak q_1, \ldots, \mathfrak q_m$ de $S$
(par montée pour $S$ sur $R$, voir le
lemme \ref{algebra-lemma-integral-going-up}).
Ainsi,
$S/\mathfrak m^nS = \prod S_{\mathfrak q_i}/\mathfrak m^nS_{\mathfrak q_i}$
par le lemme \ref{algebra-lemma-artinian-finite-length}
(en utilisant par exemple la
proposition \ref{algebra-proposition-dimension-zero-ring}
pour voir que $S/\mathfrak m^nS$ est artinien).
Ainsi, le complété $\mathfrak m$-adique $S^\wedge$ de $S$ est égal à
$\prod S_{\mathfrak q_i}^\wedge$. Enfin, nous avons
$R^\wedge \otimes_R S = S^\wedge$ par le lemme \ref{algebra-lemma-completion-tensor}.
\end{proof}

\begin{lemma}
\label{algebra-lemma-split-completed-sequence}
Soit $R$ un anneau. Soit $I \subset R$ un idéal.
Soit $0 \to K \to P \to M \to 0$ une suite exacte courte de
$R$-modules. Si $M$ est plat sur $R$ et si $M/IM$ est un
$R/I$-module projectif, alors la suite des complétés $I$-adiques
$$
0 \to K^\wedge \to P^\wedge \to M^\wedge \to 0
$$
est une suite exacte scindée.
\end{lemma}

\begin{proof}
Puisque $M$ est plat, chacune des suites
$$
0 \to K/I^nK \to P/I^nP \to M/I^nM \to 0
$$
est exacte courte, voir le
lemme \ref{algebra-lemma-flat-tor-zero},
et la suite $0 \to K^\wedge \to P^\wedge \to M^\wedge \to 0$
est une suite exacte courte, voir le
lemme \ref{algebra-lemma-completion-generalities}.
Il suffit de montrer que nous pouvons trouver des scindages
$s_n : M/I^nM \to P/I^nP$ tels que $s_{n + 1} \bmod I^n = s_n$.
Nous allons construire ces $s_n$ par récurrence sur $n$.
Choisissons un scindage quelconque $s_1$, qui existe puisque $M/IM$ est un $R/I$-module projectif.
Supposons $s_n$ donné pour un certain $n > 0$. Posons
$P_{n + 1} = \{x \in P \mid x \bmod I^nP \in \Im(s_n)\}$.
Le morphisme $\pi : P_{n + 1}/I^{n + 1}P_{n + 1} \to M/I^{n + 1}M$ est surjectif
(détails omis). Puisque $M/I^{n + 1}M$ est projectif comme $R/I^{n + 1}$-module
par le
lemme \ref{algebra-lemma-lift-projective},
nous pouvons choisir une section
$t : M/I^{n + 1}M \to P_{n + 1}/I^{n + 1}P_{n + 1}$
de $\pi$. En prenant pour $s_{n + 1}$ le composé de
$t$ et du morphisme canonique $P_{n + 1}/I^{n + 1}P_{n + 1} \to P/I^{n + 1}P$,
nous obtenons le résultat voulu.
\end{proof}

\begin{lemma}
\label{algebra-lemma-complete-modulo-nilpotent}
Soit $A$ un anneau noethérien. Soient $I, J \subset A$ des idéaux.
Si $A$ est $I$-adiquement complet et si $A/I$ est $J$-adiquement complet,
alors $A$ est $J$-adiquement complet.
\end{lemma}

\begin{proof}
Soit $B$ le complété $(I + J)$-adique de $A$. Par le
lemme \ref{algebra-lemma-completion-flat}, $B/IB$ est le complété $J$-adique de $A/I$,
donc est isomorphe à $A/I$ par hypothèse. De plus, $B$ est $I$-adiquement
complet par le lemme \ref{algebra-lemma-complete-by-sub}. Ainsi, $B$ est un
$A$-module de type fini par le lemme \ref{algebra-lemma-finite-over-complete-ring}.
Par le lemme de Nakayama (lemme \ref{algebra-lemma-NAK}, en utilisant que $I$ est contenu dans le
radical de Jacobson de $A$
par le lemme \ref{algebra-lemma-radical-completion}), nous trouvons que $A \to B$ est surjectif.
Le morphisme $A \to B$ est plat par le lemme \ref{algebra-lemma-completion-flat}.
L'image de $\Spec(B) \to \Spec(A)$ contient $V(I)$ et, comme $I$
est contenu dans le radical de Jacobson de $A$, nous trouvons que $A \to B$ est fidèlement plat
(lemme \ref{algebra-lemma-ff-rings}). Ainsi, $A \to B$ est injectif. Par conséquent, $A$ est
complet relativement à $I + J$, et donc a fortiori complet
relativement à $J$.
\end{proof}









\section{Passage à la limite de modules}
\label{algebra-section-limits}

\noindent
Dans cette section, nous étudions ce qui se produit lorsqu'on prend une limite de modules.

\begin{lemma}
\label{algebra-lemma-limit-complete-pre}
Soit $I \subset A$ un idéal de type fini d'un anneau.
Soit $(M_n)$ un système projectif de $A$-modules tel que $I^n M_n = 0$.
Alors $M = \lim M_n$ est $I$-adiquement complet.
\end{lemma}

\begin{proof}
Nous avons $M \to M/I^nM \to M_n$. En prenant la limite, nous obtenons
$M \to M^\wedge \to M$. Ainsi, $M$ est un facteur direct de $M^\wedge$.
Puisque $M^\wedge$ est $I$-adiquement complet par le
lemme \ref{algebra-lemma-hathat-finitely-generated}, il en va de même de $M$.
\end{proof}

\begin{lemma}
\label{algebra-lemma-limit-complete}
Soit $I \subset A$ un idéal de type fini d'un anneau.
Soit $(M_n)$ un système projectif de $A$-modules tel que
$M_n = M_{n + 1}/I^nM_{n + 1}$. Posons $M = \lim M_n$.
Alors $M/I^nM = M_n$ et $M$ est $I$-adiquement complet.
\end{lemma}

\begin{proof}
Par le lemme \ref{algebra-lemma-limit-complete-pre}, nous voyons que $M$ est $I$-adiquement
complet. Puisque les morphismes de transition sont surjectifs, les morphismes $M \to M_n$
sont surjectifs. Considérons le système projectif de suites exactes courtes
$$
0 \to N_n \to M \to M_n \to 0
$$
qui définit $N_n$. Puisque $M_n = M_{n + 1}/I^nM_{n + 1}$, le morphisme
$N_{n + 1} + I^nM \to N_n$ est surjectif. Ainsi,
$N_{n + 1}/(N_{n + 1} \cap I^{n + 1}M) \to N_n/(N_n \cap I^nM)$
est surjectif. En prenant la limite projective des suites exactes courtes
$$
0 \to N_n/(N_n \cap I^nM) \to M/I^nM \to M_n \to 0
$$
nous obtenons une suite exacte
$$
0 \to \lim N_n/(N_n \cap I^nM) \to M^\wedge \to M
$$
Puisque $M$ est $I$-adiquement complet, nous concluons que
$\lim N_n/(N_n \cap I^nM) = 0$ et donc, par la surjectivité
des morphismes de transition, que $N_n/(N_n \cap I^nM) = 0$ pour tout $n$.
Ainsi, $M_n = M/I^nM$, comme voulu.
\end{proof}

\begin{lemma}
\label{algebra-lemma-finiteness-graded}
Soit $A$ un anneau gradué noethérien. Soit $I \subset A_+$ un idéal
homogène. Soit $(N_n)$ un système projectif de $A$-modules gradués de type fini tel que
$N_n = N_{n + 1}/I^n N_{n + 1}$. Alors il existe un $A$-module gradué de type fini
$N$ tel que $N_n = N/I^nN$ comme modules gradués pour tout $n$.
\end{lemma}

\begin{proof}
Choisissons $r$ et des éléments homogènes $x_{1, 1}, \ldots, x_{1, r} \in N_1$ de
degrés $d_1, \ldots, d_r$ qui engendrent $N_1$. Puisque les morphismes de transition
sont surjectifs, nous pouvons choisir un système compatible d'éléments homogènes
$x_{n, i} \in N_n$ relevant $x_{1, i}$. Par le lemme de Nakayama gradué
(lemme \ref{algebra-lemma-graded-NAK}), nous voyons que
$N_n$ est engendré par les éléments $x_{n, 1}, \ldots, x_{n, r}$
placés dans les degrés $d_1, \ldots, d_r$.
Ainsi, pour $m \leq n$, nous voyons que $N_n \to N_n/I^m N_n$
est un isomorphisme en degrés $< \min(d_i) + m$ (puisque $I^mN_n$ est nul
en ces degrés). Le système projectif des composantes de degré $d$
$$
\ldots
= N_{2 + d - \min(d_i), d}
= N_{1 + d - \min(d_i), d}
= N_{d - \min(d_i), d} \to N_{-1 + d - \min(d_i), d} \to \ldots
$$
se stabilise donc comme indiqué. Soit $N$ le $A$-module gradué dont la
composante graduée de degré $d$ est cette valeur stationnaire. En particulier, nous avons les
éléments $x_i = \lim x_{n, i}$ dans $N$. Nous affirmons que les $x_i$ engendrent $N$ :
tout $x \in N_d$ est une combinaison linéaire de $x_1, \ldots, x_r$,
car nous pouvons le vérifier dans $N_{d - \min(d_i), d}$, où cela est vrai
puisque les $x_{d - \min(d_i), i}$ engendrent $N_{d - \min(d_i)}$.
Enfin, le lecteur vérifie que le morphisme surjectif
$N/I^nN \to N_n$ est un isomorphisme en examinant
ce qui se produit dans chaque degré comme ci-dessus. Les détails sont omis.
\end{proof}

\begin{lemma}
\label{algebra-lemma-daniel-litt}
Soit $A$ un anneau gradué. Soit $I \subset A_+$ un idéal homogène.
Notons $A' = \lim A/I^n$. Soit $(G_n)$ un système projectif
de $A$-modules gradués tel que $G_n$ soit annulé par $I^n$.
Soit $M$ un $A$-module gradué et soit
$\varphi_n : M \to G_n$ un système compatible de morphismes gradués de
$A$-modules. Si le morphisme induit
$$
\varphi : M \otimes_A A' \longrightarrow \lim G_n
$$
est un isomorphisme, alors $M_d \to \lim G_{n, d}$
est un isomorphisme pour tout $d \in \mathbf{Z}$.
\end{lemma}

\begin{proof}
Par convention, les anneaux gradués sont en degrés $\geq 0$, tandis que les modules gradués
peuvent avoir des composantes non nulles de tout degré, voir la section \ref{algebra-section-graded}.
Le morphisme $\varphi$ existe parce que $\lim G_n$ est
un module sur $A'$, puisque $G_n$ est annulé par $I^n$.
Une autre observation utile à garder à l'esprit est que nous avons
$$
\bigoplus\nolimits_{d \in \mathbf{Z}} \lim G_{n, d} \subset
\lim G_n \subset
\prod\nolimits_{d \in \mathbf{Z}} \lim G_{n, d}
$$
où un indice ${\ }_d$ désigne la composante graduée de degré $d$.

\medskip\noindent
Injectivité. Soit $x \in M_d$. Si $x \mapsto 0$ dans $\lim G_{n, d}$,
alors $x \otimes 1 = 0$ dans $M \otimes_A A'$. Nous pouvons alors trouver
un sous-module de type fini $M' \subset M$ tel que $x \in M'$
et que $x \otimes 1$ soit nul dans $M' \otimes_A A'$.
Disons que $M'$ est engendré par des éléments homogènes placés
en degrés $d_1, \ldots, d_r$. Soit $n = d - \min(d_i) + 1$.
Puisque $A'$ possède un morphisme vers $A/I^n$ et puisque
$A \to A/I^n$ est un isomorphisme en degrés $\leq n - 1$,
nous voyons que $M' \to M' \otimes_A A'$ est injectif en
degrés $\leq n - 1$. Ainsi, $x = 0$, comme voulu.

\medskip\noindent
Surjectivité. Soit $y \in \lim G_{n, d}$. Choisissons une somme
finie $\sum x_i \otimes f'_i$ dans $M \otimes_A A'$ qui s'envoie sur $y$.
Nous pouvons supposer $x_i$ homogène, disons de degré $d_i$.
Observons que, bien que $A'$ ne soit pas un anneau gradué, c'est une
limite des anneaux gradués $A/I^nA$ et que, de plus, dans tout
degré donné, les morphismes de transition deviennent finalement des isomorphismes
(voir ci-dessus). Cela donne
$$
A = \bigoplus\nolimits_{d \geq 0} A_d \subset A' \subset
\prod\nolimits_{d \geq 0} A_d
$$
Nous pouvons donc écrire
$$
f'_i = \sum\nolimits_{j = 0, \ldots, d - d_i - 1} f_{i, j} + f_i + g'_i
$$
avec $f_{i, j} \in A_j$, $f_i \in A_{d - d_i}$ et
$g'_i \in A'$ qui s'envoie sur zéro dans $\prod_{j \leq d - d_i} A_j$.
Si maintenant nous calculons $\varphi_n(\sum_{i, j} f_{i, j}x_i) \in G_n$,
nous obtenons une somme d'éléments homogènes de degré $< d$.
Ainsi, $\varphi(\sum x_i \otimes f_{i, j})$ s'envoie sur zéro dans
$\lim G_{n, d}$.
De même, un calcul montre que l'élément $\varphi(\sum x_i \otimes g'_i)$
s'envoie sur zéro dans $\prod_{d' \leq d} \lim G_{n, d'}$.
Puisque nous savons que $\varphi(\sum x_i \otimes f'_i)$
est $y$, nous concluons que $\sum f_ix_i \in M_d$ s'envoie sur $y$, comme voulu.
\end{proof}













\section{Critères de platitude}
\label{algebra-section-criteria-flatness}

\noindent
Dans cette section, nous démontrons quelques lemmes techniques importants dans le cas
noethérien. Nous les généraliserons (partiellement) au cas non noethérien
dans la section \ref{algebra-section-more-flatness-criteria}.

\begin{lemma}
\label{algebra-lemma-mod-injective}
Supposons que $R \to S$ soit un homomorphisme local d'anneaux locaux, avec $S$
noethérien.
Notons $\mathfrak m$ l'idéal maximal de $R$. Soient $M$ un $R$-module plat
et $N$ un $S$-module de type fini. Soit $u : N \to M$ un morphisme de $R$-modules.
Si $\overline{u} : N/\mathfrak m N \to M/\mathfrak m M$
est injectif, alors $u$ est injectif.
Dans ce cas, $M/u(N)$ est plat sur $R$.
\end{lemma}

\begin{proof}
Affirmons d'abord que $u_n : N/{\mathfrak m}^nN \to M/{\mathfrak m}^nM$
est injectif pour tout $n \geq 1$. Procédons par récurrence, le cas initial
étant que $\overline{u} = u_1$ est injectif. Comme $M$
est plat sur $R$, nous avons une suite exacte courte
$0 \to M \otimes_R {\mathfrak m}^n/{\mathfrak m}^{n + 1}
\to M/{\mathfrak m}^{n + 1}M \to M/{\mathfrak m}^n M \to 0$.
De plus, $M \otimes_R {\mathfrak m}^n/{\mathfrak m}^{n + 1}
= M/{\mathfrak m}M \otimes_{R/{\mathfrak m}}
{\mathfrak m}^n/{\mathfrak m}^{n + 1}$. Nous avons
une suite exacte analogue $N \otimes_R {\mathfrak m}^n/{\mathfrak m}^{n + 1}
\to N/{\mathfrak m}^{n + 1}N \to N/{\mathfrak m}^n N \to 0$
pour $N$, sauf que nous n'avons pas le zéro à gauche. Nous avons aussi
$N \otimes_R {\mathfrak m}^n/{\mathfrak m}^{n + 1}
= N/{\mathfrak m}N \otimes_{R/{\mathfrak m}}
{\mathfrak m}^n/{\mathfrak m}^{n + 1}$. Ainsi, le morphisme $u_{n + 1}$ est
injectif, puisque $u_n$ et le morphisme
$\overline{u} \otimes \text{id}_{{\mathfrak m}^n/{\mathfrak m}^{n + 1}}$ le sont.

\medskip\noindent
Par le théorème d'intersection de Krull
(lemme \ref{algebra-lemma-intersect-powers-ideal-module-zero})
appliqué à $N$ sur l'anneau $S$ et à l'idéal $\mathfrak mS$,
nous avons $\bigcap \mathfrak m^nN = 0$. L'injectivité
de $u_n$ pour tout $n$ implique donc que $u$ est injectif.

\medskip\noindent
Pour montrer que $M/u(N)$ est plat sur $R$, il suffit de montrer que
$\text{Tor}_1^R(M/u(N), R/I) = 0$ pour tout idéal $I \subset R$,
voir le lemme \ref{algebra-lemma-characterize-flat}. De la suite exacte courte
$$
0 \to N \xrightarrow{u} M \to M/u(N) \to 0
$$
et de la platitude de $M$, nous obtenons une suite exacte de groupes Tor
$$
0 \to \text{Tor}_1^R(M/u(N), R/I) \to N/IN \to M/IM
$$
Voir le lemme \ref{algebra-lemma-long-exact-sequence-tor}. Il suffit donc de montrer que
$N/IN$ s'injecte dans $M/IM$. Notons que $R/I \to S/IS$ est un homomorphisme
local d'anneaux locaux avec $S/IS$ noethérien, que $N/IN \to M/IM$ est un morphisme
de $R/I$-modules, que $N/IN$ est de type fini sur $S/IS$, que $M/IM$ est plat sur
$R/I$ et que $u \bmod I : N/IN \to M/IM$ est injectif modulo $\mathfrak m$.
Nous pouvons donc appliquer la première partie de la démonstration à $u \bmod I$, ce qui conclut.
\end{proof}

\begin{lemma}
\label{algebra-lemma-grothendieck}
Supposons que $R \to S$ soit un homomorphisme d'anneaux plat et local entre anneaux
locaux noethériens. Notons $\mathfrak m$ l'idéal maximal de $R$.
Supposons que $f \in S$ soit un non-diviseur de zéro dans $S/{\mathfrak m}S$.
Alors $S/fS$ est plat sur $R$, et $f$ est un non-diviseur de zéro dans $S$.
\end{lemma}

\begin{proof}
Cela découle directement du lemme \ref{algebra-lemma-mod-injective}.
\end{proof}

\begin{lemma}
\label{algebra-lemma-grothendieck-regular-sequence}
Supposons que $R \to S$ soit un homomorphisme d'anneaux plat et local entre anneaux
locaux noethériens. Notons $\mathfrak m$ l'idéal maximal de $R$.
Supposons que $f_1, \ldots, f_c$ soit une suite d'éléments de
$S$ telle que les images $\overline{f}_1, \ldots, \overline{f}_c$
forment une suite régulière dans $S/{\mathfrak m}S$.
Alors $f_1, \ldots, f_c$ est une suite régulière dans $S$ et chacun
des quotients $S/(f_1, \ldots, f_i)$ est plat sur $R$.
\end{lemma}

\begin{proof}
Par récurrence et par le lemme \ref{algebra-lemma-grothendieck}.
\end{proof}

\begin{lemma}
\label{algebra-lemma-free-fibre-flat-free}
Soit $R \to S$ un homomorphisme local d'anneaux
locaux noethériens. Soit $\mathfrak m$ l'idéal
maximal de $R$. Soit $M$ un $S$-module non nul de type fini.
Supposons que (a) $M/\mathfrak mM$
soit un $S/\mathfrak mS$-module libre, et que (b) $M$ soit plat sur $R$.
Alors $M$ est libre et $S$ est plat sur $R$.
\end{lemma}

\begin{proof}
Soit $\overline{x}_1, \ldots, \overline{x}_n$ une base
du module libre $M/\mathfrak mM$. Choisissons
$x_1, \ldots, x_n \in M$, avec $x_i$ s'envoyant sur $\overline{x}_i$. Soit
$u : S^{\oplus n} \to M$ le morphisme qui envoie le $i$-ème vecteur
de la base canonique sur $x_i$. Le lemme \ref{algebra-lemma-mod-injective}
montre que $u$ est injectif. D'autre part, par le
lemme de Nakayama \ref{algebra-lemma-NAK}, ce morphisme est surjectif. Le
résultat en découle.
\end{proof}

\begin{lemma}
\label{algebra-lemma-complex-exact-mod}
Soit $R \to S$ un homomorphisme local d'anneaux locaux
noethériens. Soit $\mathfrak m$ l'idéal maximal de $R$.
Soit $0 \to F_e \to F_{e-1} \to \ldots \to F_0$
un complexe fini de $S$-modules de type fini. Supposons que
chaque $F_i$ soit plat sur $R$, et que le complexe
$0 \to F_e/\mathfrak m F_e \to F_{e-1}/\mathfrak m F_{e-1}
\to \ldots \to F_0 / \mathfrak m F_0$ soit exact.
Alors $0 \to F_e \to F_{e-1} \to \ldots \to F_0$
est exact et, de plus, le module
$\Coker(F_1 \to F_0)$ est plat sur $R$.
\end{lemma}

\begin{proof}
Raisonnons par récurrence sur $e$. Si $e = 1$, c'est exactement le
lemme \ref{algebra-lemma-mod-injective}. Si $e > 1$, le
lemme \ref{algebra-lemma-mod-injective} montre que $F_e \to F_{e-1}$
est injectif et que $C = \Coker(F_e \to F_{e-1})$
est un $S$-module de type fini plat sur $R$. Nous pouvons donc
appliquer l'hypothèse de récurrence au complexe
$0 \to C \to F_{e-2} \to \ldots \to F_0$.
Nous en déduisons que $C \to F_{e-2}$ est injectif,
d'où l'exactitude du complexe ainsi que
la platitude du conoyau de $F_1 \to F_0$.
\end{proof}

\noindent
Dans la suite de cette section, nous démontrons deux versions de ce que l'on appelle le
« {\it critère local de platitude} ». Notons aussi l'intéressant
lemme \ref{algebra-lemma-CM-over-regular-flat} ci-dessous.

\begin{lemma}
\label{algebra-lemma-prepare-local-criterion-flatness}
Soit $R$ un anneau local d'idéal maximal $\mathfrak m$
et de corps résiduel $\kappa = R/\mathfrak m$.
Soit $M$ un $R$-module. Si $\text{Tor}_1^R(\kappa, M) = 0$,
alors, pour tout $R$-module $N$ de longueur finie, nous avons
$\text{Tor}_1^R(N, M) = 0$.
\end{lemma}

\begin{proof}
Raisonnons par récurrence sur la longueur de $N$.
Si la longueur de $N$ vaut $1$, alors $N \cong \kappa$
et le résultat est acquis. Si la longueur de $N$ est strictement supérieure à
$1$, nous pouvons insérer $N$ dans une suite exacte courte
$0 \to N' \to N \to N'' \to 0$, où $N'$ et $N''$ sont des
$R$-modules de longueur finie et de longueurs plus petites.
L'annulation de $\text{Tor}_1^R(N, M)$ découle
de celles de $\text{Tor}_1^R(N', M)$
et de $\text{Tor}_1^R(N'', M)$ (hypothèse de récurrence),
ainsi que de la suite exacte longue des groupes Tor, voir le lemme
\ref{algebra-lemma-long-exact-sequence-tor}.
\end{proof}

\begin{lemma}[Critère local de platitude]
\label{algebra-lemma-local-criterion-flatness}
Soit $R \to S$ un homomorphisme local d'anneaux locaux
noethériens. Soit $\mathfrak m$ l'idéal maximal de $R$,
et soit $\kappa = R/\mathfrak m$.
Soit $M$ un $S$-module de type fini. Si $\text{Tor}_1^R(\kappa, M) = 0$,
alors $M$ est plat sur $R$.
\end{lemma}

\begin{proof}
Soit $I \subset R$ un idéal. D'après le lemme \ref{algebra-lemma-flat}, il suffit
de montrer que $I \otimes_R M \to M$ est injectif. Par la remarque
\ref{algebra-remark-Tor-ring-mod-ideal}, nous voyons que son noyau est
égal à $\text{Tor}_1^R(M, R/I)$. Le
lemme \ref{algebra-lemma-prepare-local-criterion-flatness}
montre que $J \otimes_R M \to M$ est injectif pour tout idéal
de colongueur finie.

\medskip\noindent
Choisissons $n >> 0$ et considérons la suite exacte
courte suivante
$$
0
\to I \cap \mathfrak m^n
\to I \oplus \mathfrak m^n
\to I + \mathfrak m^n
\to 0
$$
C'est une sous-suite de la suite exacte courte
$0 \to R \to R^{\oplus 2} \to R \to 0$. Nous obtenons donc le diagramme
$$
\xymatrix{
(I\cap \mathfrak m^n) \otimes_R M \ar[r] \ar[d] &
I \otimes_R M \oplus \mathfrak m^n \otimes_R M \ar[r] \ar[d] &
(I + \mathfrak m^n) \otimes_R M \ar[d] \\
M \ar[r] &
M \oplus M \ar[r] &
M
}
$$
Notons que $I + \mathfrak m^n$ et $\mathfrak m^n$
sont des idéaux de colongueur finie.
Une chasse au diagramme montre donc que
$\Ker((I \cap \mathfrak m^n)\otimes_R M \to M)
\to \Ker(I \otimes_R M \to M)$
est surjectif. Nous en concluons en particulier que
$K = \Ker(I \otimes_R M \to M)$ est contenu
dans l'image de $(I \cap \mathfrak m^n) \otimes_R M$
dans $I \otimes_R M$. Par le lemme d'Artin-Rees \ref{algebra-lemma-Artin-Rees},
nous voyons que $K$ est contenu
dans $\mathfrak m^{n-c}(I \otimes_R M)$ pour un certain $c > 0$
et tout $n >> 0$. Comme $I \otimes_R M$ est un
$S$-module de type fini (!) et que $S$ est noethérien, nous voyons
que cela implique $K = 0$. En effet, ce qui précède implique que
$K$ s'envoie sur zéro dans la complétion $\mathfrak mS$-adique
de $I \otimes_R M$. Mais le morphisme de $S$
vers sa complétion $\mathfrak mS$-adique est fidèlement plat
par le lemme \ref{algebra-lemma-completion-faithfully-flat}.
Ainsi $K = 0$, comme souhaité.
\end{proof}

\noindent
Dans ce qui suit, nous rencontrons souvent les conditions
« $M/IM$ est plat sur $R/I$ et $\text{Tor}_1^R(R/I, M) = 0$ ».
Le lemme suivant donne quelques conséquences de ces conditions
(il s'agit d'une généralisation du
lemme \ref{algebra-lemma-prepare-local-criterion-flatness}).

\begin{lemma}
\label{algebra-lemma-what-does-it-mean}
Soit $R$ un anneau.
Soit $I \subset R$ un idéal.
Soit $M$ un $R$-module.
Si $M/IM$ est plat sur $R/I$ et $\text{Tor}_1^R(R/I, M) = 0$, alors
\begin{enumerate}
\item $M/I^nM$ est plat sur $R/I^n$ pour tout $n \geq 1$, et
\item pour tout module $N$ annulé par $I^m$ pour un certain $m \geq 0$,
nous avons $\text{Tor}_1^R(N, M) = 0$.
\end{enumerate}
En particulier, si $I$ est nilpotent, alors $M$ est plat sur $R$.
\end{lemma}

\begin{proof}
Supposons que $M/IM$ soit plat sur $R/I$ et que $\text{Tor}_1^R(R/I, M) = 0$.
Soit $N$ un $R/I$-module. Choisissons une suite exacte courte
$$
0 \to K \to \bigoplus\nolimits_{i \in I} R/I \to N \to 0
$$
Par la suite exacte longue de $\text{Tor}$ et l'annulation de
$\text{Tor}_1^R(R/I, M)$, nous obtenons
$$
0 \to \text{Tor}_1^R(N, M) \to K \otimes_R M \to
(\bigoplus\nolimits_{i \in I} R/I) \otimes_R M \to N \otimes_R M \to 0
$$
Mais comme $K$, $\bigoplus_{i \in I} R/I$ et $N$ sont tous annulés
par $I$, nous voyons que
\begin{align*}
K \otimes_R M & = K \otimes_{R/I} M/IM, \\
(\bigoplus\nolimits_{i \in I} R/I) \otimes_R M & =
(\bigoplus\nolimits_{i \in I} R/I) \otimes_{R/I} M/IM, \\
N \otimes_R M & = N \otimes_{R/I} M/IM.
\end{align*}
Puisque $M/IM$ est plat sur $R/I$, nous concluons que
$$
0 \to K \otimes_{R/I} M/IM \to
(\bigoplus\nolimits_{i \in I} R/I) \otimes_{R/I} M/IM \to
N \otimes_{R/} M/IM \to 0
$$
est exacte. En combinant ceci avec ce qui précède, nous concluons que
$\text{Tor}_1^R(N, M) = 0$ pour tout $R$-module $N$ annulé par $I$.

\medskip\noindent
Démontrons (2) par récurrence sur $m$. Le cas $m = 1$ a été
traité au paragraphe précédent. Pour $N$ annulé par $I^m$
avec $m > 1$, nous pouvons choisir une suite exacte
$0 \to N' \to N \to N'' \to 0$, où $N'$ et $N''$ sont annulés
par $I^{m - 1}$. Par exemple, on peut prendre $N' = IN$ et $N'' = N/IN$.
Alors la suite exacte
$$
\text{Tor}_1^R(N', M) \to
\text{Tor}_1^R(N, M) \to
\text{Tor}_1^R(N'', M)
$$
et la récurrence prouvent l'annulation voulue.

\medskip\noindent
Enfin, démontrons (1). Étant donné $n \geq 1$, nous devons montrer que $M/I^nM$
est plat sur $R/I^n$. Autrement dit, nous devons montrer que le foncteur
$N \mapsto N \otimes_{R/I^n} M/I^nM$ est exact sur la catégorie des $R$-modules
$N$ annulés par $I^n$. Or, pour un tel $N$, nous avons
$N \otimes_{R/I^n} M/I^nM = N \otimes_R M$. L'annulation de
$\text{Tor}_1$ dans (2) montre que le foncteur $N \mapsto N \otimes_R M$
est exact sur la catégorie des $N$ annulés par une puissance de $I$,
ce qui conclut.
\end{proof}

\begin{lemma}
\label{algebra-lemma-what-does-it-mean-again}
Soient $R$ un anneau, $I \subset R$ un idéal et
$M$ un $R$-module.
\begin{enumerate}
\item Si $M/IM$ est plat sur $R/I$ et si $M \otimes_R I/I^2 \to IM/I^2M$
est injectif, alors $M/I^2M$ est plat sur $R/I^2$.
\item Si $M/IM$ est plat sur $R/I$ et si $M \otimes_R I^n/I^{n + 1}
\to I^nM/I^{n + 1}M$ est injectif pour $n = 1, \ldots, k$,
alors $M/I^{k + 1}M$ est plat sur $R/I^{k + 1}$.
\end{enumerate}
\end{lemma}

\begin{proof}
La première assertion découle du
lemme \ref{algebra-lemma-what-does-it-mean} appliqué en remplaçant $R$ par $R/I^2$
et $M$ par $M/I^2M$, en utilisant l'égalité
$$
\text{Tor}_1^{R/I^2}(M/I^2M, R/I) =
\Ker(M \otimes_R I/I^2 \to IM/I^2M),
$$
voir la remarque \ref{algebra-remark-Tor-ring-mod-ideal}.
La seconde assertion s'obtient de la même manière en raisonnant par récurrence
sur $n$ pour montrer que $M/I^{n + 1}M$ est plat sur $R/I^{n + 1}$ pour
$n = 1, \ldots, k$. Nous utilisons ici l'égalité
$$
\text{Tor}_1^{R/I^{n + 1}}(M/I^{n + 1}M, R/I^n) =
\Ker(M \otimes_R I^n/I^{n + 1} \to I^nM/I^{n + 1}M)
$$
pour tout $n$.
\end{proof}

\begin{lemma}[Variante du critère local]
\label{algebra-lemma-variant-local-criterion-flatness}
Soit $R \to S$ un homomorphisme local d'anneaux locaux
noethériens. Soit $I \not = R$ un idéal de $R$.
Soit $M$ un $S$-module de type fini. Si $\text{Tor}_1^R(M, R/I) = 0$
et si $M/IM$ est plat sur $R/I$, alors $M$ est plat sur $R$.
\end{lemma}

\begin{proof}
Première démonstration : par le
lemme \ref{algebra-lemma-what-does-it-mean},
nous voyons que $\text{Tor}_1^R(\kappa, M)$ est nul, où $\kappa$
est le corps résiduel de $R$. Le module $M$ est donc
plat sur $R$ par le
lemme \ref{algebra-lemma-local-criterion-flatness}.

\medskip\noindent
Seconde démonstration : soit $\mathfrak m$ l'idéal maximal de $R$.
Nous allons montrer que $\mathfrak m \otimes_R M \to M$ est injectif,
puis appliquer le
lemme \ref{algebra-lemma-local-criterion-flatness}.
Supposons que $\sum f_i \otimes x_i \in \mathfrak m \otimes_R M$
et que $\sum f_i x_i = 0$ dans $M$. Par le critère équationnel
de platitude, lemme \ref{algebra-lemma-flat-eq}, appliqué à $M/IM$
sur $R/I$, nous voyons qu'il existe des $\overline{a}_{ij} \in R/I$
et des $\overline{y}_j \in M/IM$ tels que
$x_i \bmod IM = \sum_j \overline{a}_{ij} \overline{y}_j $
et $0 = \sum_i (f_i \bmod I) \overline{a}_{ij}$.
Soit $a_{ij} \in R$ un relèvement de $\overline{a}_{ij}$ et,
de même, soit $y_j \in M$ un relèvement de $\overline{y}_j$.
Nous voyons alors que
\begin{eqnarray*}
\sum f_i \otimes x_i
& = &
\sum f_i \otimes x_i +
\sum f_ia_{ij} \otimes y_j -
\sum f_i \otimes a_{ij} y_j
\\
& = &
\sum f_i \otimes (x_i - \sum a_{ij} y_j) +
\sum (\sum f_i a_{ij}) \otimes y_j
\end{eqnarray*}
Comme $x_i - \sum a_{ij} y_j \in IM$ et
$\sum f_i a_{ij} \in I$, nous voyons qu'il existe
un élément de $I \otimes_R M$ qui s'envoie sur notre
élément donné $\sum f_i \otimes x_i$ dans $\mathfrak m \otimes_R M$.
Mais $I \otimes_R M \to M$ est injectif par hypothèse (voir la
remarque \ref{algebra-remark-Tor-ring-mod-ideal}), ce qui conclut.
\end{proof}

\noindent
En particulier, dans la situation du
lemme \ref{algebra-lemma-variant-local-criterion-flatness}, supposons que
$I = (x)$ soit engendré par un seul élément $x$ qui soit
un non-diviseur de zéro dans $R$. Alors $\text{Tor}_1^R(M, R/(x)) = (0)$
si et seulement si $x$ est un non-diviseur de zéro sur $M$.

\begin{lemma}
\label{algebra-lemma-flat-module-powers}
Soit $R \to S$ un homomorphisme d'anneaux. Soit $I \subset R$ un idéal.
Soit $M$ un $S$-module. Supposons que
\begin{enumerate}
\item $R$ soit un anneau noethérien,
\item $S$ soit un anneau noethérien,
\item $M$ soit un $S$-module de type fini, et
\item pour tout $n \geq 1$, le module $M/I^n M$ soit plat sur
$R/I^n$.
\end{enumerate}
Alors, pour tout $\mathfrak q \in V(IS)$,
le localisé $M_{\mathfrak q}$ est plat sur $R$.
En particulier, si $S$ est local et si $IS$ est contenu
dans son idéal maximal, alors $M$ est plat sur $R$.
\end{lemma}

\begin{proof}
Nous allons utiliser le
lemme \ref{algebra-lemma-variant-local-criterion-flatness}.
Par hypothèse, $M/IM$ est plat sur $R/I$. Il suffit donc de vérifier
que $\text{Tor}_1^R(M, R/I)$ est nul après localisation en $\mathfrak q$. Par la
remarque \ref{algebra-remark-Tor-ring-mod-ideal},
ce groupe Tor est égal à $K = \Ker(I \otimes_R M \to M)$.
Nous savons que le noyau de
$I/I^n \otimes_{R/I^n} M/I^nM \to M/I^nM$ est nul pour tout $n \geq 1$.
Tout élément de $K$ s'envoie donc sur zéro dans $I/I^n \otimes_{R/I^n} M/I^nM$.
Comme
$$
I/I^n \otimes_{R/I^n} M/I^nM = I/I^n \otimes_R M =
(I \otimes_R M)/I^{n - 1}(I \otimes_R M)
$$
nous concluons que $K \subset I^{n - 1}(I \otimes_R M)$ pour tout $n \geq 1$.
Par le lemme d'Artin-Rees, et plus précisément par le
lemme \ref{algebra-lemma-intersection-powers-ideal-module},
nous concluons que $K_{\mathfrak q} = 0$, comme souhaité.
\end{proof}

\begin{lemma}
\label{algebra-lemma-surjective-on-tor-one}
Soient $R \to R' \to R''$ des homomorphismes d'anneaux.
Soit $M$ un $R$-module. Supposons que $M \otimes_R R'$
soit plat sur $R'$. Alors le morphisme naturel
$\text{Tor}_1^R(M, R') \otimes_{R'} R'' \to
\text{Tor}_1^R(M, R'')$ est surjectif.
\end{lemma}

\begin{proof}
Soit $F_\bullet$ une résolution libre de $M$ sur $R$.
Le complexe $F_2 \otimes_R R' \to F_1\otimes_R R' \to F_0 \otimes_R R'$
calcule $\text{Tor}_1^R(M, R')$.
Le complexe $F_2 \otimes_R R'' \to F_1\otimes_R R'' \to F_0 \otimes_R R''$
calcule $\text{Tor}_1^R(M, R'')$. Notons que
$F_i \otimes_R R' \otimes_{R'} R'' = F_i \otimes_R R''$. Posons
$K' = \Ker(F_1\otimes_R R' \to F_0 \otimes_R R')$ et,
de même, $K'' = \Ker(F_1\otimes_R R'' \to F_0 \otimes_R R'')$.
Nous avons donc une suite exacte
$$
0 \to K' \to F_1\otimes_R R' \to F_0 \otimes_R R' \to M \otimes_R R' \to 0.
$$
Puisque, par hypothèse, $M \otimes_R R'$ est plat sur $R'$,
la suite
$$
K' \otimes_{R'} R'' \to
F_1 \otimes_R R'' \to
F_0 \otimes_R R'' \to
M \otimes_R R'' \to 0
$$
reste exacte. Cela signifie que $K' \otimes_{R'} R'' \to K''$
est surjectif. Comme $\text{Tor}_1^R(M, R')$ est un quotient de $K'$ et que
$\text{Tor}_1^R(M, R'')$ est un quotient de $K''$, le résultat en découle.
\end{proof}

\begin{lemma}
\label{algebra-lemma-surjective-on-tor-one-trivial}
Soit $R \to R'$ un homomorphisme d'anneaux. Soient $I \subset R$
un idéal et $I' = IR'$. Soit $M$ un $R$-module
et posons $M' = M \otimes_R R'$. Le morphisme naturel
$\text{Tor}_1^R(R'/I', M) \to \text{Tor}_1^{R'}(R'/I', M')$
est surjectif.
\end{lemma}

\begin{proof}
Soit $F_2 \to F_1 \to F_0 \to M \to 0$ une résolution libre de
$M$ sur $R$. Posons $F_i' = F_i \otimes_R R'$. La suite
$F_2' \to F_1' \to F_0' \to M' \to 0$ peut ne plus être exacte
en $F_1'$. Une résolution libre de $M'$ sur $R'$ est donc de la forme
suivante :
$$
F_2' \oplus F_2'' \to F_1' \to F_0' \to M' \to 0
$$
pour un module libre convenable $F_2''$ sur $R'$. Ensuite, notons que
$F_i \otimes_R R'/I' = F_i' / IF_i' = F_i'/I'F_i'$.
Le complexe $F_2'/I'F_2' \to F_1'/I'F_1' \to F_0'/I'F_0'$
calcule donc $\text{Tor}_1^R(M, R'/I')$. D'autre part,
$F_i' \otimes_{R'} R'/I' = F_i'/I'F_i'$ et il en va de même
pour $F_2''$. Ainsi, le complexe
$F_2'/I'F_2' \oplus F_2''/I'F_2'' \to F_1'/I'F_1' \to F_0'/I'F_0'$
calcule $\text{Tor}_1^{R'}(M', R'/I')$. Comme le morphisme vertical
de complexes
$$
\xymatrix{
F_2'/I'F_2' \ar[r] \ar[d] &
F_1'/I'F_1' \ar[r] \ar[d] &
F_0'/I'F_0' \ar[d] \\
F_2'/I'F_2' \oplus F_2''/I'F_2'' \ar[r] &
F_1'/I'F_1' \ar[r] &
F_0'/I'F_0'
}
$$
induit manifestement une surjection en homologie, le résultat en découle.
\end{proof}

\begin{lemma}
\label{algebra-lemma-another-variant-local-criterion-flatness}
Soit
$$
\xymatrix{
S \ar[r] & S' \\
R \ar[r] \ar[u] & R' \ar[u]
}
$$
un diagramme commutatif d'homomorphismes locaux d'anneaux locaux noethériens.
Soit $I \subset R$ un idéal propre.
Soit $M$ un $S$-module de type fini.
Notons $I' = IR'$ et $M' = M \otimes_S S'$.
Supposons que
\begin{enumerate}
\item $S'$ soit un localisé du produit tensoriel
$S \otimes_R R'$,
\item $M/IM$ soit plat sur $R/I$,
\item $\text{Tor}_1^R(M, R/I) \to \text{Tor}_1^{R'}(M', R'/I')$
soit nul.
\end{enumerate}
Alors $M'$ est plat sur $R'$.
\end{lemma}

\begin{proof}
Puisque $S'$ est un localisé de $S \otimes_R R'$, nous voyons que
$M'$ est un localisé de $M \otimes_R R'$. Notons que,
par le lemme \ref{algebra-lemma-flat-base-change}, le module $M/IM \otimes_{R/I} R'/I'
= M \otimes_R R' /I'(M \otimes_R R')$ est plat sur $R'/I'$. Par conséquent,
$M'/I'M'$ est lui aussi plat sur $R'/I'$, puisque tout localisé d'un module plat
est plat. D'après le lemme \ref{algebra-lemma-variant-local-criterion-flatness},
il suffit de montrer que $\text{Tor}_1^{R'}(M', R'/I')$ est nul.
Comme $M'$ est un localisé de $M \otimes_R R'$, la dernière hypothèse
implique qu'il suffit de montrer que
$\text{Tor}_1^R(M, R/I) \otimes_R R'
\to
\text{Tor}_1^{R'}(M \otimes_R R', R'/I')$
est surjectif.

\medskip\noindent
Le lemme \ref{algebra-lemma-surjective-on-tor-one-trivial} montre que
$\text{Tor}_1^R(M, R'/I') \to \text{Tor}_1^{R'}(M \otimes_R R', R'/I')$
est surjectif. Il suffit donc maintenant de montrer que
$\text{Tor}_1^R(M, R/I) \otimes_R R'
\to
\text{Tor}_1^R(M, R'/I')$
est surjectif. Cela découle du lemme \ref{algebra-lemma-surjective-on-tor-one}
appliqué aux homomorphismes d'anneaux $R \to R/I \to R'/I'$ et au module $M$.
\end{proof}

\noindent
Veuillez comparer le lemme ci-dessous au
lemme \ref{algebra-lemma-criterion-flatness-fibre-nilpotent}
(le cas d'un idéal nilpotent) et au
lemme \ref{algebra-lemma-criterion-flatness-fibre}
(le cas des algèbres de présentation finie).

\begin{lemma}[Crit\`ere de platitude par fibres ; cas noethérien]
\label{algebra-lemma-criterion-flatness-fibre-Noetherian}
Soient $R$, $S$ et $S'$ des anneaux locaux noethériens, et soient $R \to S \to S'$
des homomorphismes locaux d'anneaux. Soit $\mathfrak m \subset R$
l'idéal maximal. Soit $M$ un $S'$-module. Supposons que
\begin{enumerate}
\item le module $M$ soit de type fini sur $S'$ ;
\item le module $M$ soit non nul ;
\item le module $M/\mathfrak m M$
soit un $S/\mathfrak m S$-module plat ;
\item le module $M$ soit un $R$-module plat.
\end{enumerate}
Alors $S$ est plat sur $R$ et $M$ est un $S$-module plat.
\end{lemma}

\begin{proof}
Posons $I = \mathfrak mS \subset S$. Nous voyons que $M/IM$ est un
$S/I$-module plat en vertu de (3). Comme
$\mathfrak m \otimes_R S' \to I \otimes_S S'$ est surjectif, nous voyons
que $\mathfrak m \otimes_R M \to I \otimes_S M$ est également surjectif.
Considérons
$$
\mathfrak m \otimes_R M \to I \otimes_S M \to M.
$$
Comme $M$ est plat sur $R$, la composée est injective,
et les deux flèches le sont donc.
En particulier, $\text{Tor}_1^S(S/I, M) = 0$, voir la
remarque \ref{algebra-remark-Tor-ring-mod-ideal}. Par le
lemme \ref{algebra-lemma-variant-local-criterion-flatness}, nous concluons
que $M$ est plat sur $S$. Notons que, puisque $M/\mathfrak m_{S'}M$
n'est pas nul par le lemme de Nakayama \ref{algebra-lemma-NAK},
nous voyons qu'en fait $M$ est fidèlement plat sur $S$ par le
lemme \ref{algebra-lemma-ff} (car cela impose $M/\mathfrak m_SM \not = 0$).

\medskip\noindent
Considérons la suite exacte
$0 \to \mathfrak m \to R \to \kappa \to 0$.
Elle donne une suite exacte
$0 \to \text{Tor}_1^R(\kappa, S) \to \mathfrak m \otimes_R S \to I \to 0$.
Puisque $M$ est plat sur $S$, celle-ci donne une suite exacte
$0 \to \text{Tor}_1^R(\kappa, S)\otimes_S M \to
\mathfrak m \otimes_R M \to I \otimes_S M \to 0$.
D'après ce qui précède, cela implique que $\text{Tor}_1^R(\kappa, S)\otimes_S M = 0$.
Puisque $M$ est fidèlement plat sur $S$, il s'ensuit que
$\text{Tor}_1^R(\kappa, S) = 0$, et nous concluons que
$S$ est plat sur $R$ par le lemme \ref{algebra-lemma-local-criterion-flatness}.
\end{proof}

\begin{lemma}
\label{algebra-lemma-flatness-scallop}
Soient $A$ un anneau, $M$ un $A$-module et $f \in A$. Si
\begin{enumerate}
\item $f$ est un non-diviseur de zéro sur $A$ et sur $M$,
\item $M_f$ est un $A_f$-module plat, et
\item $M/fM$ est un $A/fA$-module plat,
\end{enumerate}
alors $M$ est un $A$-module plat. Il en va de même en remplaçant « plat » par
« fidèlement plat ».
\end{lemma}

\begin{proof}
Comme $0 \to A \to A \to A/fA \to 0$ et $0 \to M \to M \to M/fM \to 0$ sont
exactes, nous obtenons $\text{Tor}_i^A(M, A/fA) = 0$ pour $i = 1$ et $i = 2$.
Par le lemme \ref{algebra-lemma-what-does-it-mean}, nous concluons que
$\text{Tor}_1^A(M, N) = 0$ pour tout $A$-module $N$ annulé par $f$
(ceci utilise la platitude de $M/fM$ sur $A/fA$).
Étant donné un $A$-module $N$ annulé par $f$, nous pouvons
choisir une suite exacte courte $0 \to N' \to F \to N \to 0$ de $A$-modules,
où $F$ est une somme directe de copies de $A/fA$. De la suite exacte
$$
\text{Tor}_2^A(M, F) \to
\text{Tor}_2^A(M, N) \to
\text{Tor}_1^A(M, N')
$$
nous concluons que $\text{Tor}_2^A(M, N) = 0$ pour tout $A$-module
$N$ annulé par $f$. Soit ensuite $K$ un $A$-module arbitraire.
Nous pouvons décomposer le morphisme $f : K \to K$ en deux suites exactes courtes
$$
0 \to K[f] \to K \to K' \to 0
\quad\text{et}\quad
0 \to K' \to K \to K/fK \to 0
$$
où $K' = K/K[f] \cong fK$. En appliquant les suites exactes de Tor, nous
obtenons les suites exactes
$$
\text{Tor}_1^A(K[f], M) \to
\text{Tor}_1^A(K, M) \to
\text{Tor}_1^A(K', M)
$$
et
$$
\text{Tor}_2^A(K/fK, M) \to
\text{Tor}_1^A(K', M) \to
\text{Tor}_1^A(K, M)
$$
En utilisant l'annulation de
$\text{Tor}_1^A(K[f], M)$ et de
$\text{Tor}_2^A(K/fK, M)$,
nous concluons que $f : K \to K$ induit un morphisme injectif sur
$\text{Tor}^A_1(M, K)$. Autrement dit, nous voyons que $\text{Tor}^A_1(M, K)$
est nul si et seulement si $\text{Tor}^A_1(M, K) \otimes_A A_f$ est nul.
Par le lemme \ref{algebra-lemma-flat-base-change-tor}, nous avons
$$
\text{Tor}^A_1(M, K) \otimes_A A_f = \text{Tor}^{A_f}_1(M_f, K_f)  = 0
$$
Cette annulation résulte de la platitude de $M_f$ sur $A_f$
(lemme \ref{algebra-lemma-characterize-flat}).
Nous concluons que $\text{Tor}_1^A(M, K) = 0$
pour tout $A$-module $K$. Par conséquent,
$M$ est plat sur $A$ par le lemme \ref{algebra-lemma-characterize-flat}.

\medskip\noindent
Nous omettons l'argument dans le cas des modules fidèlement plats.
\end{proof}

\begin{lemma}
\label{algebra-lemma-flatness-scallop-pre}
Soient $A$ un anneau et $M$ un $A$-module. Soit $I = (f_1, \ldots, f_r)$
un idéal de $A$ engendré par $r \geq 1$ éléments. Si
\begin{enumerate}
\item $M_{f_i}$ est un $A_{f_i}$-module plat pour $i = 1, \ldots, r$,
\item $M/IM$ est un $A/I$-module plat,
\item $\text{Tor}_i^A(M, A/I) = 0$ pour $i = 1, \ldots, r + 1$,
\end{enumerate}
alors $M$ est un $A$-module plat. Il en va de même en remplaçant « plat » par
« fidèlement plat ».
\end{lemma}

\begin{proof}
Le lemme \ref{algebra-lemma-what-does-it-mean} montre que
$\text{Tor}_1^A(M, K) = 0$ pour tout $A$-module $K$ annulé par $I$.
Étant donné un $A$-module $K$ annulé par $I$, nous pouvons choisir une suite
exacte courte $0 \to N \to F \to K \to 0$ de $A$-modules, où $F$
est une somme directe de copies de $A/I$. Nous obtenons une suite exacte
$$
\text{Tor}_i^A(M, F) \to
\text{Tor}_i^A(M, K) \to
\text{Tor}_{i - 1}^A(M, N)
$$
Ainsi, en utilisant l'hypothèse (3) et en raisonnant par récurrence sur $i$,
nous concluons que $\text{Tor}_i^A(M, K) = 0$
pour tout $A$-module $K$ annulé par $I$ et $i = 1, \ldots, r + 1$.

\medskip\noindent
Supposons que, pour un certain $1 \leq j \leq r$,
nous ayons montré que $\text{Tor}_i^A(M, K) = 0$
pour tout $A$-module $K$ annulé par $f_1, \ldots, f_j$
et $i = 1, \ldots, j + 1$.
Soit $K$ un $A$-module annulé par $f_1, \ldots, f_{j - 1}$.
Nous pouvons décomposer le morphisme $f_j : K \to K$ en deux suites exactes courtes
$$
0 \to K[f_j] \to K \to K' \to 0
\quad\text{et}\quad
0 \to K' \to K \to K/f_jK \to 0
$$
où $K' = K/K[f_j] \cong f_jK$. Soit $1 \leq i \leq j$.  En appliquant les
suites exactes de Tor, nous obtenons les suites exactes
$$
\text{Tor}_i^A(K[f_j], M) \to
\text{Tor}_i^A(K, M) \to
\text{Tor}_i^A(K', M)
$$
et
$$
\text{Tor}_{i + 1}^A(K/f_jK, M) \to
\text{Tor}_i^A(K', M) \to
\text{Tor}_i^A(K, M)
$$
En utilisant l'annulation de $\text{Tor}_i^A(K[f_j], M)$ et de
$\text{Tor}_{i + 1}^A(K/f_jK, M)$, nous concluons que $f_j : K \to K$
induit un morphisme injectif sur $\text{Tor}^A_i(M, K)$.
Autrement dit, nous voyons que $\text{Tor}^A_i(M, K)$
est nul si et seulement si $\text{Tor}^A_i(M, K) \otimes_A A_{f_j}$ est nul.
Par le lemme \ref{algebra-lemma-flat-base-change-tor}, nous avons
$$
\text{Tor}^A_i(M, K) \otimes_A A_{f_j} =
\text{Tor}^{A_{f_j}}_i(M_{f_j}, K_{f_j})  = 0
$$
Cette annulation résulte de la platitude de $M_{f_j}$ sur $A_{f_j}$
(lemme \ref{algebra-lemma-characterize-flat}).
Nous concluons que $\text{Tor}_i^A(M, K) = 0$
pour tout $A$-module $K$ annulé par $f_1, \ldots, f_{j - 1}$
et $i = 1, \ldots, j$.
Par récurrence descendante sur $j$, nous concluons que
ceci vaut pour $j = 0$, c'est-à-dire que
$\text{Tor}_1^A(M, K) = 0$
pour tout $A$-module $K$. Par conséquent,
$M$ est plat sur $A$ par le lemme \ref{algebra-lemma-characterize-flat}.

\medskip\noindent
Nous omettons l'argument dans le cas des modules fidèlement plats.
\end{proof}








\section{Changement de base et platitude}
\label{algebra-section-base-change-flat}

\noindent
Quelques lemmes décrivant ce qu'il advient de la platitude lors
d'un changement de base.

\begin{lemma}
\label{algebra-lemma-base-change-flat-up-down}
Soit
$$
\xymatrix{
S \ar[r] & S' \\
R \ar[r] \ar[u] & R' \ar[u]
}
$$
un diagramme commutatif d'homomorphismes locaux d'anneaux locaux.
Supposons que $S'$ soit un localisé du produit tensoriel $S \otimes_R R'$.
Soit $M$ un $S$-module et posons $M' = S' \otimes_S M$.
\begin{enumerate}
\item Si $M$ est plat sur $R$, alors $M'$ est plat sur $R'$.
\item Si $M'$ est plat sur $R'$ et si $R \to R'$ est plat, alors
$M$ est plat sur $R$.
\end{enumerate}
En particulier, nous avons
\begin{enumerate}
\item[(3)] Si $S$ est plat sur $R$, alors $S'$ est plat sur $R'$.
\item[(4)] Si $R' \to S'$ et $R \to R'$ sont plats, alors $S$ est plat sur $R$.
\end{enumerate}
\end{lemma}

\begin{proof}
Démonstration de (1). Si $M$ est plat sur $R$, alors $M \otimes_R R'$
est plat sur $R'$ par le
lemme \ref{algebra-lemma-flat-base-change}.
Si $W \subset S \otimes_R R'$ est la partie multiplicative telle que
$W^{-1}(S \otimes_R R') = S'$, alors $M' = W^{-1}(M \otimes_R R')$.
Ainsi, $M'$ est plat sur $R'$ en tant que localisé d'un module plat, voir la
partie (5) du lemme \ref{algebra-lemma-flat-localization}. Ceci démontre (1) et,
en particulier, montre que (3) est vraie.

\medskip\noindent
Démonstration de (2). Supposons que $M'$ soit plat sur $R'$ et que $R \to R'$ soit plat.
En appliquant (3) au diagramme réfléchi par rapport à la diagonale nord-ouest,
nous voyons que $S \to S'$ est plat. Ainsi, $S \to S'$ est fidèlement plat par le
lemme \ref{algebra-lemma-local-flat-ff}.
Nous allons utiliser le critère du
lemme \ref{algebra-lemma-flat} (\ref{algebra-item-f-ideal})
pour montrer que $M$ est plat.
Soit $I \subset R$ un idéal. Si $I \otimes_R M \to M$ a un noyau non nul,
il en va de même de $(I \otimes_R M) \otimes_S S' \to M \otimes_S S' = M'$.
Notons que $I \otimes_R R' = IR'$, puisque $R \to R'$ est plat, et que
$$
(I \otimes_R M) \otimes_S S' =
(I \otimes_R R') \otimes_{R'} (M \otimes_S S') =
IR' \otimes_{R'} M'.
$$
La platitude de $M'$ sur $R'$
implique que ce module s'injecte dans $M'$.
Ceci achève la démonstration de (2), et (4) est donc également vraie.
\end{proof}

\noindent
Voici encore une application du critère local de platitude.

\begin{lemma}
\label{algebra-lemma-yet-another-variant-local-criterion-flatness}
Considérons un diagramme commutatif d'anneaux locaux et d'homomorphismes locaux
$$
\xymatrix{
S \ar[r] & S' \\
R \ar[r] \ar[u] & R' \ar[u]
}
$$
Soit $M$ un $S$-module de type fini. Supposons que
\begin{enumerate}
\item les flèches horizontales soient des homomorphismes d'anneaux plats,
\item $M$ soit plat sur $R$,
\item $\mathfrak m_R R' = \mathfrak m_{R'}$,
\item $R'$ et $S'$ soient noethériens.
\end{enumerate}
Alors $M' = M \otimes_S S'$ est plat sur $R'$.
\end{lemma}

\begin{proof}
Comme $\mathfrak m_R \subset R$ et que $R \to R'$ est plat, nous obtenons
$\mathfrak m_R \otimes_R R' = \mathfrak m_R R' = \mathfrak m_{R'}$
par l'hypothèse (3). Observons que $M'$ est un $S'$-module de type fini
qui est plat sur $R$ par le lemme \ref{algebra-lemma-flatness-descends-more-general}.
Ainsi, $\mathfrak m_R \otimes_R M' \to M'$ est injectif.
Nous obtenons alors
$$
\mathfrak m_R \otimes_R M' =
\mathfrak m_R \otimes_R R' \otimes_{R'} M' =
\mathfrak m_{R'} \otimes_{R'} M'
$$
Ainsi, $\mathfrak m_{R'} \otimes_{R'} M' \to M'$ est injectif.
Ceci montre que $\text{Tor}_1^{R'}(\kappa_{R'}, M') = 0$
(remarque \ref{algebra-remark-Tor-ring-mod-ideal}).
Le module $M'$ est donc plat sur $R'$ par le
lemme \ref{algebra-lemma-local-criterion-flatness}.
\end{proof}







\section{Critères de platitude sur les anneaux artiniens}
\label{algebra-section-flatness-artinian}

\noindent
Nous examinons quelques critères de platitude pour les modules sur les anneaux artiniens.
Notons que l'idéal maximal d'un anneau local artinien est nilpotent,
de sorte que les deux lemmes suivants s'appliquent aux anneaux locaux artiniens.

\begin{lemma}
\label{algebra-lemma-local-artinian-basis-when-flat}
Soit $(R, \mathfrak m)$ un anneau local dont l'idéal maximal
$\mathfrak m$ est nilpotent. Soit $M$ un $R$-module plat.
Si $A$ est un ensemble et si $x_\alpha \in M$, $\alpha \in A$, est une famille
d'éléments de $M$, alors les assertions suivantes sont équivalentes :
\begin{enumerate}
\item $\{\overline{x}_\alpha\}_{\alpha \in A}$ est une base
de l'espace vectoriel $M/\mathfrak mM$ sur $R/\mathfrak m$, et
\item $\{x_\alpha\}_{\alpha \in A}$ est une base de $M$ sur $R$.
\end{enumerate}
\end{lemma}

\begin{proof}
L'implication (2) $\Rightarrow$ (1) est immédiate.
Supposons (1). Par le lemme de Nakayama \ref{algebra-lemma-NAK},
les éléments $x_\alpha$ engendrent $M$. On obtient alors une suite exacte
courte
$$
0 \to K \to \bigoplus\nolimits_{\alpha \in A} R \to M \to 0
$$
En tensorisant par $R/\mathfrak m$ et en utilisant le lemme \ref{algebra-lemma-flat-tor-zero},
nous obtenons $K/\mathfrak mK = 0$. Par le lemme de Nakayama \ref{algebra-lemma-NAK},
nous concluons que $K = 0$.
\end{proof}

\begin{lemma}
\label{algebra-lemma-local-artinian-characterize-flat}
Soit $R$ un anneau local d'idéal maximal nilpotent. Soit $M$ un $R$-module.
Les assertions suivantes sont équivalentes :
\begin{enumerate}
\item $M$ est plat sur $R$,
\item $M$ est un $R$-module libre, et
\item $M$ est un $R$-module projectif.
\end{enumerate}
\end{lemma}

\begin{proof}
Comme tout module projectif est plat (en tant que facteur direct d'un module libre)
et que tout module libre est projectif, il suffit de montrer qu'un module plat
est libre. Soit $M$ un module plat. Soient $A$ un ensemble et des éléments $x_\alpha \in M$,
$\alpha \in A$, tels que les
$\overline{x_\alpha} \in M/\mathfrak m M$ forment une base sur le corps
résiduel de $R$. Par le
lemme \ref{algebra-lemma-local-artinian-basis-when-flat},
les $x_\alpha$ sont une base de $M$ sur $R$, ce qui conclut.
\end{proof}

\begin{lemma}
\label{algebra-lemma-lift-basis}
Soit $R$ un anneau.
Soit $I \subset R$ un idéal.
Soit $M$ un $R$-module.
Soient $A$ un ensemble et $x_\alpha \in M$, $\alpha \in A$, une famille
d'éléments de $M$.
Supposons que
\begin{enumerate}
\item $I$ soit nilpotent,
\item $\{\overline{x}_\alpha\}_{\alpha \in A}$ soit une base de $M/IM$ sur
$R/I$, et
\item $\text{Tor}_1^R(R/I, M) = 0$.
\end{enumerate}
Alors $M$ est libre de base $\{x_\alpha\}_{\alpha \in A}$ sur $R$.
\end{lemma}

\begin{proof}
Soient $R$, $I$, $M$, $\{x_\alpha\}_{\alpha \in A}$ comme dans le lemme,
et supposons les hypothèses (1), (2) et (3) satisfaites. Par le
lemme de Nakayama \ref{algebra-lemma-NAK},
les éléments $x_\alpha$ engendrent $M$ sur $R$.
L'hypothèse $\text{Tor}_1^R(R/I, M) = 0$ implique que nous avons une suite
exacte courte
$$
0 \to I \otimes_R M \to M \to M/IM \to 0.
$$
Soit $\sum f_\alpha x_\alpha = 0$ une relation dans $M$.
Par le choix des $x_\alpha$, nous voyons que $f_\alpha \in I$.
Nous concluons donc que $\sum f_\alpha \otimes x_\alpha = 0$ dans
$I \otimes_R M$. Le morphisme $I \otimes_R M \to I/I^2 \otimes_{R/I} M/IM$
et le fait que $\{x_\alpha\}_{\alpha \in A}$ soit une base
de $M/IM$ impliquent que $f_\alpha \in I^2$ ! Nous concluons donc qu'il
n'existe aucune relation entre les images des $x_\alpha$ dans
$M/I^2M$. Autrement dit, nous voyons que $M/I^2M$ est libre, de base
les images des $x_\alpha$. En utilisant le morphisme
$I \otimes_R M \to I/I^3 \otimes_{R/I^2} M/I^2M$,
nous concluons ensuite que $f_\alpha \in I^3$ !
Et ainsi de suite. Comme $I^n = 0$ pour un certain $n$ par l'hypothèse (1), le résultat suit.
\end{proof}

\begin{lemma}
\label{algebra-lemma-prepare-lift-flatness}
Soit $\varphi : R \to R'$ un homomorphisme d'anneaux.
Soit $I \subset R$ un idéal.
Soit $M$ un $R$-module.
Supposons que
\begin{enumerate}
\item $M/IM$ soit plat sur $R/I$, et
\item $R' \otimes_R M$ soit plat sur $R'$.
\end{enumerate}
Posons $I_2 = \varphi^{-1}(\varphi(I^2)R')$.
Alors $M/I_2M$ est plat sur $R/I_2$.
\end{lemma}

\begin{proof}
Nous pouvons remplacer $R$, $M$ et $R'$ par $R/I_2$, $M/I_2M$ et
$R'/\varphi(I)^2R'$. Alors $I^2 = 0$ et $\varphi$ est injectif. Par le
lemme \ref{algebra-lemma-what-does-it-mean}
et puisque $I^2 = 0$, il suffit de montrer que
$\text{Tor}^R_1(R/I, M) = K = \Ker(I \otimes_R M \to M)$ est nul.
Posons $M' = M \otimes_R R'$ et $I' = IR'$.
Par hypothèse, le morphisme $I' \otimes_{R'} M' \to M'$ est injectif.
Ainsi, $K$ s'envoie sur zéro dans
$$
I' \otimes_{R'} M' = I' \otimes_R M = I' \otimes_{R/I} M/IM.
$$
Le morphisme $I \to I'$ est alors un morphisme injectif de $R/I$-modules.
Comme $M/IM$ est plat sur $R/I$, le morphisme
$$
I \otimes_{R/I} M/IM \longrightarrow I' \otimes_{R/I} M/IM
$$
est injectif. Ceci implique que $K$ est nul dans
$I \otimes_R M = I \otimes_{R/I} M/IM$, comme souhaité.
\end{proof}

\begin{lemma}
\label{algebra-lemma-lift-flatness}
Soit $\varphi : R \to R'$ un homomorphisme d'anneaux.
Soit $I \subset R$ un idéal.
Soit $M$ un $R$-module.
Supposons que
\begin{enumerate}
\item $I$ soit nilpotent,
\item $R \to R'$ soit injectif,
\item $M/IM$ soit plat sur $R/I$, et
\item $R' \otimes_R M$ soit plat sur $R'$.
\end{enumerate}
Alors $M$ est plat sur $R$.
\end{lemma}

\begin{proof}
Définissons par récurrence $I_1 = I$ et $I_{n + 1} = \varphi^{-1}(\varphi(I_n)^2R')$
pour $n \geq 1$. Notons que, par le
lemme \ref{algebra-lemma-prepare-lift-flatness},
nous obtenons que $M/I_nM$ est plat sur $R/I_n$ pour tout $n \geq 1$.
Il est clair que $\varphi(I_{n + 1}) \subset \varphi(I)^{2^n}R'$. Comme
$I$ est nilpotent, nous voyons que $\varphi(I_n) = 0$ pour un certain $n$. Puisque
$\varphi$ est injectif, nous concluons que $I_n = 0$ pour un certain $n$, ce qui
achève la démonstration.
\end{proof}

\noindent
Voici la version artinienne locale du critère local de platitude.

\begin{lemma}
\label{algebra-lemma-artinian-variant-local-criterion-flatness}
Soit $R$ un anneau local artinien. Soit $M$ un $R$-module.
Soit $I \subset R$ un idéal propre. Les assertions suivantes sont
équivalentes :
\begin{enumerate}
\item $M$ est plat sur $R$, et
\item $M/IM$ est plat sur $R/I$ et $\text{Tor}_1^R(R/I, M) = 0$.
\end{enumerate}
\end{lemma}

\begin{proof}
L'implication (1) $\Rightarrow$ (2) découle immédiatement des
définitions. Supposons que $M/IM$ soit plat sur $R/I$ et que
$\text{Tor}_1^R(R/I, M) = 0$. Par le
lemme \ref{algebra-lemma-local-artinian-characterize-flat},
ceci implique que $M/IM$ est libre sur $R/I$. Choisissons un ensemble $A$
et des éléments $x_\alpha \in M$ dont les images dans $M/IM$ forment
une base. Par le
lemme \ref{algebra-lemma-lift-basis},
nous concluons que $M$ est libre et, en particulier, plat.
\end{proof}

\noindent
Il se trouve que la platitude descend par tout homomorphisme injectif
dont la source est un anneau artinien.

\begin{lemma}
\label{algebra-lemma-descent-flatness-injective-map-artinian-rings}
Soit $R \to S$ un homomorphisme d'anneaux. Soit $M$ un $R$-module.
Supposons que
\begin{enumerate}
\item $R$ soit artinien,
\item $R \to S$ soit injectif, et
\item $M \otimes_R S$ soit un $S$-module plat.
\end{enumerate}
Alors $M$ est un $R$-module plat.
\end{lemma}

\begin{proof}
Première démonstration : soit $I \subset R$ le radical de Jacobson de $R$.
Alors $I$ est nilpotent et $M/IM$ est plat sur $R/I$, puisque $R/I$
est un produit de corps, voir la
section \ref{algebra-section-artinian}.
Par conséquent, $M$ est plat par une application du
lemme \ref{algebra-lemma-lift-flatness}.

\medskip\noindent
Seconde démonstration : par le
lemme \ref{algebra-lemma-artinian-finite-length},
nous pouvons écrire $R = \prod R_i$ comme produit fini d'anneaux locaux
artiniens. Ceci induit des décompositions en produits analogues pour $R$ et $S$.
Nous nous ramenons donc au cas où $R$ est local artinien (détails omis).

\medskip\noindent
Supposons que $R \to S$ et $M$ soient comme dans le lemme, satisfassent (1), (2) et (3),
et qu'en outre $R$ soit local d'idéal maximal $\mathfrak m$.
Soient $A$ un ensemble et $x_\alpha \in A$ des éléments tels que
$\overline{x}_\alpha$ forme une base de $M/\mathfrak mM$
sur $R/\mathfrak m$. Par le
lemme de Nakayama \ref{algebra-lemma-NAK},
nous voyons que les éléments $x_\alpha$ engendrent $M$ comme $R$-module.
Posons $N = S \otimes_R M$ et $I = \mathfrak mS$.
Alors $\{1 \otimes x_\alpha\}_{\alpha \in A}$ est une famille d'éléments
de $N$ qui forme une base de $N/IN$. De plus, puisque $N$ est plat sur
$S$, nous avons $\text{Tor}_1^S(S/I, N) = 0$. Le
lemme \ref{algebra-lemma-lift-basis}
permet donc de conclure que $N$ est libre de base $\{1 \otimes x_\alpha\}_{\alpha \in A}$.
L'injectivité de $R \to S$ garantit alors qu'il ne peut exister de
relation non triviale entre les $x_\alpha$ à coefficients dans $R$.
\end{proof}

\noindent
Veuillez comparer le lemme ci-dessous au
lemme \ref{algebra-lemma-criterion-flatness-fibre-Noetherian}
(le cas des anneaux locaux noethériens), au
lemme \ref{algebra-lemma-criterion-flatness-fibre}
(le cas des algèbres de présentation finie), et au
lemme \ref{algebra-lemma-criterion-flatness-fibre-locally-nilpotent}
(le cas des idéaux localement nilpotents).

\begin{lemma}[Crit\`ere de platitude par fibres : cas nilpotent]
\label{algebra-lemma-criterion-flatness-fibre-nilpotent}
Soit
$$
\xymatrix{
S \ar[rr] & & S' \\
& R \ar[lu] \ar[ru]
}
$$
un diagramme commutatif dans la catégorie des anneaux.
Soient $I \subset R$ un idéal nilpotent et $M$ un $S'$-module. Supposons que
\begin{enumerate}
\item le module $M/IM$ soit un $S/IS$-module plat ;
\item le module $M$ soit un $R$-module plat.
\end{enumerate}
Alors $M$ est un $S$-module plat et $S_{\mathfrak q}$ est plat sur $R$
pour tout $\mathfrak q \subset S$ tel que $M \otimes_S \kappa(\mathfrak q)$
soit non nul.
\end{lemma}

\begin{proof}
Puisque $M$ est plat sur $R$, la tensorisation par la suite exacte
courte $0 \to I \to R \to R/I \to 0$ donne une suite exacte courte
$$
0 \to I \otimes_R M \to M \to M/IM \to 0.
$$
Notons que $I \otimes_R M \to IS \otimes_S M$ est surjectif. Combiné avec
ce qui précède, cela signifie que les deux morphismes de
$$
I \otimes_R M \to IS \otimes_S M \to M
$$
sont injectifs. Ainsi, $\text{Tor}_1^S(IS, M) = 0$ (voir la
remarque \ref{algebra-remark-Tor-ring-mod-ideal}),
et nous concluons que $M$ est un $S$-module plat par le
lemme \ref{algebra-lemma-what-does-it-mean}.
Pour terminer, nous devons montrer que $S_{\mathfrak q}$ est plat sur
$R$ pour tout idéal premier $\mathfrak q \subset S$ tel que
$M \otimes_S \kappa(\mathfrak q)$ soit non nul. Ceci découle du
lemme \ref{algebra-lemma-ff} et du lemme \ref{algebra-lemma-flat-permanence}.
\end{proof}








\section{Qu'est-ce qui rend un complexe exact ?}
\label{algebra-section-complex-exact}

\noindent
Une partie de ce matériau se trouve dans l'article \cite{WhatExact}
de Buchsbaum et Eisenbud.

\begin{situation}
\label{algebra-situation-complex}
Ici, $R$ est un anneau, et nous avons un complexe
$$
0
\to
R^{n_e}
\xrightarrow{\varphi_e}
R^{n_{e-1}}
\xrightarrow{\varphi_{e-1}}
\ldots
\xrightarrow{\varphi_{i + 1}}
R^{n_i}
\xrightarrow{\varphi_i}
R^{n_{i-1}}
\xrightarrow{\varphi_{i-1}}
\ldots
\xrightarrow{\varphi_1}
R^{n_0}
$$
Autrement dit, nous imposons $\varphi_i \circ \varphi_{i + 1} = 0$
pour $i = 1, \ldots, e - 1$.
\end{situation}

\begin{lemma}
\label{algebra-lemma-add-trivial-complex}
Supposons que $R$ soit un anneau. Soit
$$
\ldots
\xrightarrow{\varphi_{i + 1}}
R^{n_i}
\xrightarrow{\varphi_i}
R^{n_{i-1}}
\xrightarrow{\varphi_{i-1}}
\ldots
$$
un complexe de $R$-modules libres de rang fini. Supposons que, pour un certain $i$,
un coefficient matriciel du morphisme $\varphi_i$ soit inversible.
Alors le complexe affiché est isomorphe à la somme directe d'un complexe
$$
\ldots \to
R^{n_{i + 2}} \xrightarrow{\varphi_{i + 2}}
R^{n_{i + 1}} \to
R^{n_i - 1} \to
R^{n_{i - 1} - 1} \to
R^{n_{i - 2}} \xrightarrow{\varphi_{i - 2}}
R^{n_{i - 3}} \to
\ldots
$$
et du complexe $\ldots \to 0 \to R \to R \to 0 \to \ldots$,
où le morphisme $R \to R$ est l'identité.
\end{lemma}

\begin{proof}
L'hypothèse signifie, après un changement de base de
$R^{n_i}$ et de $R^{n_{i-1}}$, que le premier vecteur de la base
de $R^{n_i}$ est envoyé par $\varphi_i$ sur le premier vecteur de la base
de $R^{n_{i-1}}$. Notons $e_j$ le
$j$-ème vecteur de la base de $R^{n_i}$ et $f_k$ le $k$-ème
vecteur de la base de $R^{n_{i-1}}$. Écrivons $\varphi_i(e_j)
= \sum a_{jk} f_k$. Ainsi, $a_{1k} = 0$ sauf si $k = 1$,
et $a_{11} = 1$. Changeons de nouveau la base de $R^{n_i}$
en posant $e'_j = e_j - a_{j1} e_1$ pour $j > 1$.
Après ce changement de coordonnées, nous avons $a_{j1} = 0$
pour $j > 1$. Notons que l'image
de $R^{n_{i + 1}} \to R^{n_i}$ est contenue dans le
sous-module engendré par les $e_j$, $j > 1$. Notons également
que $R^{n_{i-1}} \to R^{n_{i-2}}$ doit annuler
$f_1$, puisqu'il appartient à l'image. Ces conditions
et la forme de la matrice $(a_{jk})$ de $\varphi_i$
impliquent le lemme.
\end{proof}

\noindent
Dans la situation \ref{algebra-situation-complex}, nous disons qu'un complexe de la forme
$$
0 \to \ldots \to 0 \to R \xrightarrow{1} R \to 0 \to \ldots \to 0
$$
ou de la forme
$$
0 \to \ldots \to 0 \to R
$$
est {\it trivial}. Plus précisément, nous disons que
$0 \to R^{n_e} \to R^{n_{e-1}} \to \ldots \to R^{n_0}$
est trivial s'il existe soit un $e \geq i \geq 1$
tel que $n_i = n_{i - 1} = 1$, $\varphi_i = \text{id}_R$, et que
$n_j = 0$ pour $j \not \in \{i, i - 1\}$, soit
$n_0 = 1$ et $n_i = 0$ pour $i > 0$.
Le lemme précédent dit clairement que
tout complexe fini de modules libres de rang fini sur un anneau local est, à des
sommes directes de complexes triviaux près, identique à un complexe
dont tous les morphismes ont tous leurs coefficients matriciels dans
l'idéal maximal.

\begin{lemma}
\label{algebra-lemma-exact-depth-zero-local}
Dans la situation \ref{algebra-situation-complex}, supposons que $R$ soit
un anneau local noethérien d'idéal maximal $\mathfrak m$.
Supposons que $\mathfrak m \in \text{Ass}(R)$, autrement dit que
$R$ soit de profondeur $0$. Supposons que
$0 \to R^{n_e} \to R^{n_{e-1}} \to \ldots \to R^{n_0}$
soit exact en $R^{n_e}, \ldots, R^{n_1}$.
Alors le complexe est isomorphe à une somme directe de complexes
triviaux.
\end{lemma}

\begin{proof}
Choisissons $x \in R$, $x \not = 0$, tel que $\mathfrak m x = 0$.
Soit $i$ le plus grand indice tel que $n_i > 0$.
Si $i = 0$, alors l'assertion est vraie. Si
$i > 0$, notons $f_1$ le premier vecteur de la base de $R^{n_i}$.
Comme $xf_1$ ne s'envoie pas sur zéro en raison de
l'exactitude du complexe, nous en déduisons qu'un coefficient
matriciel du morphisme $R^{n_i} \to R^{n_{i - 1}}$
n'appartient pas à $\mathfrak m$.
Le lemme \ref{algebra-lemma-add-trivial-complex} nous permet alors
de diminuer $n_e + \ldots + n_1$. Une récurrence achève la démonstration.
\end{proof}

\begin{lemma}
\label{algebra-lemma-exact-artinian-local}
Dans la situation \ref{algebra-situation-complex}, soit $R$ un anneau local artinien.
Supposons que $0 \to R^{n_e} \to R^{n_{e-1}} \to \ldots \to R^{n_0}$
soit exact en $R^{n_e}, \ldots, R^{n_1}$. Alors le complexe est isomorphe
à une somme directe de complexes triviaux.
\end{lemma}

\begin{proof}
C'est un cas particulier du lemme \ref{algebra-lemma-exact-depth-zero-local},
car tout anneau local artinien est de profondeur $0$.
\end{proof}

\noindent
Nous définissons ci-dessous le rang d'un morphisme de modules libres de rang fini.
Ce n'est qu'une définition possible du rang. C'est
simplement celle qui convient dans cette section ; il en
existe d'autres qui peuvent être plus commodes dans d'autres situations.

\begin{definition}
\label{algebra-definition-rank}
Soit $R$ un anneau. Supposons que $\varphi : R^m \to R^n$ soit un morphisme
de modules libres de rang fini.
\begin{enumerate}
\item Le {\it rang} de $\varphi$ est le plus grand $r$ tel que
$\wedge^r \varphi : \wedge^r R^m \to \wedge^r R^n$ soit non nul.
\item Nous notons $I(\varphi) \subset R$ l'idéal engendré par
les mineurs $r \times r$ de la matrice de $\varphi$, où $r$
est le rang défini ci-dessus.
\end{enumerate}
\end{definition}

\noindent
Le rang de $\varphi : R^m \to R^n$ vaut $0$
si et seulement si $\varphi = 0$, et dans ce cas $I(\varphi) = R$.

\begin{lemma}
\label{algebra-lemma-trivial-case-exact}
Dans la situation \ref{algebra-situation-complex}, supposons que le complexe soit
isomorphe à une somme directe de complexes triviaux. Alors
nous avons
\begin{enumerate}
\item les morphismes $\varphi_i$ ont pour rang
$r_i = n_i - n_{i + 1} + \ldots + (-1)^{e-i-1} n_{e-1} + (-1)^{e-i} n_e$,
\item pour tout $i$, $1 \leq i \leq e - 1$, nous avons
$\text{rank}(\varphi_{i + 1}) + \text{rank}(\varphi_i) = n_i$,
\item chaque $I(\varphi_i) = R$.
\end{enumerate}
\end{lemma}

\begin{proof}
Nous pouvons supposer que le complexe est la somme directe de complexes
triviaux. Pour chaque $i$, nous pouvons alors répartir les vecteurs de la base
canonique de $R^{n_i}$ entre ceux qui s'envoient sur un vecteur de la base
de $R^{n_{i-1}}$ et ceux qui s'envoient sur zéro (et ces derniers
sont atteints par des vecteurs de la base de $R^{n_{i + 1}}$ si $i > 0$).
En raisonnant par
récurrence descendante à partir de $i = e$, il est facile de prouver qu'il y
a $r_{i + 1}$ vecteurs de la base de $R^{n_i}$ qui s'envoient
sur zéro, et $r_i$ qui s'envoient sur des vecteurs de la base de
$R^{n_{i-1}}$. Le résultat en découle.
\end{proof}

\begin{lemma}
\label{algebra-lemma-div-x-exact-one-less}
Dans la situation \ref{algebra-situation-complex}, supposons que $R$ soit
un anneau local d'idéal maximal $\mathfrak m$.
Supposons que $0 \to R^{n_e} \to R^{n_{e-1}} \to \ldots \to R^{n_0}$
soit exact en $R^{n_e}, \ldots, R^{n_1}$.
Soit $x \in \mathfrak m$ un non-diviseur de zéro. Le complexe
$0 \to (R/xR)^{n_e} \to \ldots \to (R/xR)^{n_1}$
est exact en $(R/xR)^{n_e}, \ldots, (R/xR)^{n_2}$.
\end{lemma}

\begin{proof}
Notons $F_\bullet$ le complexe de termes $F_i = R^{n_i}$
et dont la différentielle est donnée par $\varphi_i$. Nous avons alors une suite
exacte courte de complexes
$$
0 \to F_\bullet \xrightarrow{x} F_\bullet \to F_\bullet/xF_\bullet \to 0
$$
En appliquant le lemme du serpent, nous obtenons une suite exacte longue
$$
H_i(F_\bullet) \xrightarrow{x} H_i(F_\bullet) \to
H_i(F_\bullet/xF_\bullet) \to H_{i - 1}(F_\bullet)
\xrightarrow{x} H_{i - 1}(F_\bullet)
$$
Le lemme en découle.
\end{proof}

\begin{lemma}[Lemme d'acyclicité]
\label{algebra-lemma-acyclic}
\begin{reference}
\cite[Lemma 1.8]{Peskine-Szpiro}
\end{reference}
Soit $R$ un anneau local noethérien.
Soit $0 \to M_e \to M_{e-1} \to \ldots \to M_0$
un complexe de $R$-modules de type fini.
Supposons que $\text{depth}(M_i) \geq i$.
Soit $i$ le plus grand indice où le complexe n'est
pas exact en $M_i$. Si $i > 0$, alors
$\Ker(M_i \to M_{i-1})/\Im(M_{i + 1} \to M_i)$
est de profondeur $\geq 1$.
\end{lemma}

\begin{proof}
Soit $H = \Ker(M_i \to M_{i-1})/\Im(M_{i + 1} \to M_i)$ le
groupe d'homologie considéré.
Nous pouvons décomposer le complexe en suites exactes courtes
$0 \to M_e \to M_{e-1} \to K_{e-2} \to 0$,
$0 \to K_j \to M_j \to K_{j-1} \to 0$, pour $i + 2 \leq j \leq e-2 $,
$0 \to K_{i + 1} \to M_{i + 1} \to B_i \to 0$,
$0 \to K_i \to M_i \to M_{i-1}$, et
$0 \to B_i \to K_i \to H \to 0$.
Nous remontons ces complexes pour
prouver les assertions sur les profondeurs, en utilisant de façon répétée le
lemme \ref{algebra-lemma-depth-in-ses}.
Tout d'abord, puisque $\text{depth}(M_e) \geq e$
et que $\text{depth}(M_{e-1}) \geq e-1$, nous en déduisons
que $\text{depth}(K_{e-2}) \geq e - 1$. À ce stade, les
suites $0 \to K_j \to M_j \to K_{j-1} \to 0$ pour $i + 2 \leq j \leq e-2 $
impliquent de même que $\text{depth}(K_{j-1}) \geq j$ pour
$i + 2 \leq j \leq e-2$. La suite
$0 \to K_{i + 1} \to M_{i + 1} \to B_i \to 0$
montre alors que $\text{depth}(B_i) \geq i + 1$. La suite
$0 \to K_i \to M_i \to M_{i-1}$ montre que $\text{depth}(K_i) \geq 1$,
puisque $M_i$ est de profondeur $\geq i \geq 1$ par hypothèse.
La suite $0 \to B_i \to K_i \to H \to 0$
implique alors le résultat.
\end{proof}

\begin{proposition}
\label{algebra-proposition-what-exact}
\begin{reference}
\cite[Corollary 1]{WhatExact}
\end{reference}
Dans la situation \ref{algebra-situation-complex}, supposons que $R$ soit
un anneau local noethérien. Les assertions suivantes sont équivalentes :
\begin{enumerate}
\item $0 \to R^{n_e} \to R^{n_{e-1}} \to \ldots \to R^{n_0}$
est exact en $R^{n_e}, \ldots, R^{n_1}$, et
\item pour tout $i$, $1 \leq i \leq e$,
les deux conditions suivantes sont satisfaites :
\begin{enumerate}
\item $\text{rank}(\varphi_i) = r_i$, où
$r_i = n_i - n_{i + 1} + \ldots + (-1)^{e-i-1} n_{e-1} + (-1)^{e-i} n_e$,
\item $I(\varphi_i) = R$, ou bien $I(\varphi_i)$ contient une
suite régulière de longueur $i$.
\end{enumerate}
\end{enumerate}
\end{proposition}

\begin{proof}
Si, pour un certain $i$, un coefficient matriciel de $\varphi_i$
n'appartient pas à $\mathfrak m$, nous appliquons le lemme \ref{algebra-lemma-add-trivial-complex}.
Il est facile de voir que la proposition pour un complexe et
pour le même complexe auquel on a ajouté un complexe trivial
sont équivalentes. Nous pouvons donc supposer que tous les coefficients matriciels
de chaque $\varphi_i$ sont des éléments de l'idéal maximal.
Nous pouvons également supposer que $e \geq 1$.

\medskip\noindent
Supposons que le complexe soit exact en $R^{n_e}, \ldots, R^{n_1}$.
Soit $\mathfrak q \in \text{Ass}(R)$.
Notons que l'anneau $R_{\mathfrak q}$ est de profondeur $0$
et que le complexe reste exact après localisation en $\mathfrak q$.
Nous appliquons les lemmes \ref{algebra-lemma-exact-depth-zero-local} et
\ref{algebra-lemma-trivial-case-exact} au complexe localisé
sur $R_{\mathfrak q}$. Nous concluons que
$\varphi_{i, \mathfrak q}$ est de rang $r_i$ pour tout $i$.
Comme $R \to \bigoplus_{\mathfrak q \in \text{Ass}(R)} R_\mathfrak q$
est injectif (lemme \ref{algebra-lemma-zero-at-ass-zero}), nous concluons que
$\varphi_i$ est de rang $r_i$ sur $R$, par la définition du rang donnée
dans la définition \ref{algebra-definition-rank}. Nous voyons donc que
$I(\varphi_i)_\mathfrak q = I(\varphi_{i, \mathfrak q})$,
puisque les rangs ne changent pas. Comme tous les idéaux
$I(\varphi_i)_{\mathfrak q}$, $e \geq i \geq 1$,
sont égaux à $R_{\mathfrak q}$ (par les lemmes cités ci-dessus),
nous concluons qu'aucun des idéaux $I(\varphi_i)$ n'est contenu dans $\mathfrak q$.
Ceci implique que $I(\varphi_e)I(\varphi_{e-1})\ldots I(\varphi_1)$
n'est contenu dans aucun des idéaux premiers associés
à $R$. Par le lemme \ref{algebra-lemma-silly}, nous pouvons choisir
$x \in I(\varphi_e)I(\varphi_{e - 1})\ldots I(\varphi_1)$,
$x \not \in \mathfrak q$ pour tout $\mathfrak q \in \text{Ass}(R)$.
Observons que $x$ est un non-diviseur de zéro (lemme \ref{algebra-lemma-ass-zero-divisors}).
D'après le lemme \ref{algebra-lemma-div-x-exact-one-less},
le complexe $0 \to (R/xR)^{n_e} \to \ldots \to (R/xR)^{n_1}$ est exact
en $(R/xR)^{n_e}, \ldots, (R/xR)^{n_2}$. Par récurrence
sur $e$, tous les idéaux $I(\varphi_i)/xR$ possèdent une suite
régulière de longueur $i - 1$. Ceci prouve que $I(\varphi_i)$
contient une suite régulière de longueur $i$.

\medskip\noindent
Supposons que (2)(a) et (2)(b) soient satisfaites. Nous allons prouver (1) par
récurrence sur $\dim(R)$. Si $\dim(R) = 0$, alors nous devons avoir
$I(\varphi_i) = R$ pour $1 \leq i \leq e$ par (2)(b). Comme les
coefficients de $\varphi_i$ sont contenus dans l'idéal maximal,
cela ne peut se produire que si $r_i = 0$ pour tout $i$. Par (2)(a), nous
concluons que $e = 0$, et (1) est satisfaite. Supposons que $\dim(R) > 0$.
Affirmons que, pour tout idéal premier $\mathfrak p \subset R$, les conditions
(2)(a) et (2)(b) sont satisfaites pour le complexe
$0 \to R_\mathfrak p^{n_e} \to R_\mathfrak p^{n_{e - 1}} \to \ldots \to
R_\mathfrak p^{n_0}$, de morphismes $\varphi_{i, \mathfrak p}$,
sur $R_\mathfrak p$. En effet, comme $I(\varphi_i)$ contient un
non-diviseur de zéro, l'image de $I(\varphi_i)$ dans $R_\mathfrak p$
est non nulle. Ceci implique que le rang de $\varphi_{i, \mathfrak p}$
est le même que celui de $\varphi_i$ : le rang, tel que défini ci-dessus,
d'une matrice $\varphi$ sur un anneau $R$ ne peut que diminuer lors du passage
à une $R$-algèbre $R'$, et cela se produit si et seulement si $I(\varphi)$
s'envoie sur zéro dans $R'$. Ainsi, (2)(a) est satisfaite. Cela étant dit,
nous savons que $I(\varphi_{i, \mathfrak p}) = I(\varphi_i)_\mathfrak p$,
et nous voyons que (2)(b) est également préservée par localisation.
Par récurrence sur la dimension de $R$, nous pouvons supposer que le complexe
est exact après localisation en tout idéal premier non maximal $\mathfrak p$ de $R$.
Ainsi, $\Ker(\varphi_i)/\Im(\varphi_{i + 1})$ a son support contenu dans
$\{\mathfrak m\}$ et, s'il est non nul, est donc de profondeur $0$.
Comme $I(\varphi_i) \subset \mathfrak m$ pour tout $i$, en vertu
de ce qui a été dit dans le premier paragraphe de la démonstration, nous
voyons que (2)(b) implique $\text{depth}(R) \geq e$.
Par le lemme \ref{algebra-lemma-acyclic}, nous voyons
que le complexe est exact en $R^{n_e}, \ldots, R^{n_1}$,
ce qui achève la démonstration.
\end{proof}

\begin{remark}
\label{algebra-remark-what-exact}
Si, dans la proposition \ref{algebra-proposition-what-exact}, les conditions équivalentes (1)
et (2) sont satisfaites, alors il existe un $j$ tel que $I(\varphi_i) = R$
si et seulement si $i \geq j$. Comme dans la démonstration de la proposition, il suffit
de le voir lorsque toutes les matrices ont leurs coefficients dans l'idéal maximal
$\mathfrak m$ de $R$. Dans ce cas, nous voyons que $I(\varphi_j) = R$
si et seulement si $\varphi_j = 0$. Mais si $\varphi_j = 0$, alors nous obtenons
des complexes exacts arbitrairement longs
$0 \to R^{n_e} \to R^{n_{e-1}} \to \ldots
\to R^{n_j} \to 0 \to 0 \to \ldots \to 0$,
et la proposition montre donc que $I(\varphi_i)$, pour $i > j$,
doit être égal à $R$ (sinon, c'est un idéal propre d'un anneau local
noethérien contenant des suites régulières arbitrairement longues, ce qui est impossible).
\end{remark}











\section{Modules de Cohen-Macaulay}
\label{algebra-section-CM}

\noindent
Nous montrons ici que les modules de Cohen-Macaulay possèdent de bonnes propriétés. Nous remettons
à plus tard l'emploi des groupes Ext pour établir le lien avec la dualité, entre autres.

\begin{definition}
\label{algebra-definition-CM}
Soit $R$ un anneau local noethérien.
Soit $M$ un $R$-module de type fini.
Nous disons que $M$ est {\it de Cohen-Macaulay}
si $\dim(\text{Supp}(M)) = \text{depth}(M)$.
\end{definition}

\noindent
Un premier objectif sera d'établir la proposition \ref{algebra-proposition-CM-module}.
Nous y parvenons par une suite (peut-être non standard) de lemmes élémentaires
ne faisant presque appel à aucun des résultats antérieurs sur la profondeur. Introduisons
quelques notations.

\medskip\noindent
Soit $R$ un anneau local noethérien. Soient $M$ un module
de Cohen-Macaulay et $f_1, \ldots, f_d$
une suite $M$-régulière, avec $d = \dim(\text{Supp}(M))$.
Nous disons que $g \in \mathfrak m$ est {\it bon relativement à
$(M, f_1, \ldots, f_d)$} si, pour tout $i = 0, 1, \ldots, d-1$,
nous avons $\dim (\text{Supp}(M) \cap V(g, f_1, \ldots, f_i))
= d - i - 1$. Ceci équivaut à la condition
$\dim(\text{Supp}(M/(f_1, \ldots, f_i)M) \cap V(g)) =
d - i - 1$ pour $i = 0, 1, \ldots, d - 1$.

\begin{lemma}
\label{algebra-lemma-good-element}
Notations et hypothèses comme ci-dessus. Si $g$ est bon relativement à
$(M, f_1, \ldots, f_d)$, alors (a) $g$ est un non-diviseur de zéro sur $M$,
et (b) $M/gM$ est de Cohen-Macaulay, de suite régulière
maximale $f_1, \ldots, f_{d - 1}$.
\end{lemma}

\begin{proof}
Nous démontrons le lemme par récurrence sur $d$.
Si $d = 0$, alors $M$ est de type fini et il n'existe aucun cas
auquel le lemme s'applique.
Si $d = 1$, nous devons montrer que $g : M \to M$ est
injectif. Le noyau $K$ a pour support $\{\mathfrak m\}$,
car, par hypothèse, $\dim \text{Supp}(M) \cap V(g) = 0$.
Ainsi, $K$ est de longueur finie. L'injectivité de $f_1 : K \to K$
implique donc que son image a la même longueur que $K$, et par conséquent que
$f_1 K = K$, ce qui, par le lemme de Nakayama \ref{algebra-lemma-NAK}, implique $K = 0$.
De plus, $\dim \text{Supp}(M/gM) = 0$, si bien que $M/gM$ est de Cohen-Macaulay
et de profondeur $0$.

\medskip\noindent
Supposons que $d > 1$. Observons que $g$ est bon pour $(M/f_1M, f_2, \ldots, f_d)$,
comme on le voit aisément d'après la définition. Par récurrence, nous avons que
(a) $g$ est un non-diviseur de zéro sur $M/f_1M$ et que
(b) $M/(g, f_1)M$ est de Cohen-Macaulay, de suite régulière maximale
$f_2, \ldots, f_{d - 1}$. Par le
lemme \ref{algebra-lemma-permute-xi},
nous voyons que $g, f_1$ est une suite $M$-régulière.
Ainsi, $g$ est un non-diviseur de zéro sur $M$ et
$f_1, \ldots, f_{d - 1}$ est une suite $M/gM$-régulière.
\end{proof}

\begin{lemma}
\label{algebra-lemma-CM-one-g}
Soit $R$ un anneau local noethérien.
Soit $M$ un module de Cohen-Macaulay sur $R$.
Supposons que $g \in \mathfrak m$ vérifie $\dim(\text{Supp}(M) \cap V(g))
= \dim(\text{Supp}(M)) - 1$. Alors (a) $g$ est un non-diviseur de zéro sur $M$,
et (b) $M/gM$ est de Cohen-Macaulay, de profondeur diminuée d'une unité.
\end{lemma}

\begin{proof}
Choisissons une suite $M$-régulière $f_1, \ldots, f_d$, avec
$d = \dim(\text{Supp}(M))$. Si $g$ est bon relativement à
$(M, f_1, \ldots, f_d)$, le lemme \ref{algebra-lemma-good-element} conclut.
En particulier, le lemme est vrai si $d = 1$. (Le cas $d = 0$ ne
se produit pas.) Supposons que $d > 1$. Choisissons un élément $h \in R$ tel que
(\romannumeral1) $h$ soit bon relativement à $(M, f_1, \ldots, f_d)$,
et que (\romannumeral2) $\dim(\text{Supp}(M) \cap V(h, g)) = d - 2$.
Pour voir que $h$ existe, soit $\{\mathfrak q_j\}$ l'ensemble (fini) des
idéaux premiers minimaux des fermés $\text{Supp}(M)$,
$\text{Supp}(M)\cap V(f_1, \ldots, f_i)$, $i = 1, \ldots, d - 1$,
et $\text{Supp}(M) \cap V(g)$. Aucun de ces $\mathfrak q_j$
n'est égal à $\mathfrak m$ ; nous pouvons donc trouver $h \in \mathfrak m$,
$h \not \in \mathfrak q_j$, par le lemme \ref{algebra-lemma-silly}. Il est clair
que $h$ satisfait (\romannumeral1) et (\romannumeral2). Du
lemme \ref{algebra-lemma-good-element}, nous déduisons que
$M/hM$ est de Cohen-Macaulay. Par (\romannumeral2), nous voyons que la paire
$(M/hM, g)$ satisfait l'hypothèse de récurrence. Ainsi,
$M/(h, g)M$ est de Cohen-Macaulay et $g : M/hM \to M/hM$
est injectif. Par le lemme \ref{algebra-lemma-permute-xi}, nous voyons
que $g : M \to M$ et $h : M/gM \to M/gM$
sont injectifs. Combiné au fait que $M/(g, h)M$
est de Cohen-Macaulay, ceci achève la démonstration.
\end{proof}

\begin{proposition}
\label{algebra-proposition-CM-module}
Soit $R$ un anneau local noethérien, d'idéal maximal $\mathfrak m$.
Soit $M$ un module de Cohen-Macaulay sur $R$, dont le support est de dimension $d$.
Supposons que $g_1, \ldots, g_c$ soient des éléments de
$\mathfrak m$ tels que $\dim(\text{Supp}(M/(g_1, \ldots, g_c)M))
= d - c$. Alors $g_1, \ldots, g_c$ est une suite $M$-régulière,
et peut être prolongée en une suite $M$-régulière maximale.
\end{proposition}

\begin{proof}
Soit $Z = \text{Supp}(M) \subset \Spec(R)$.
Par le lemme \ref{algebra-lemma-one-equation}, dans la chaîne
$Z \supset Z \cap V(g_1) \supset \ldots \supset Z \cap V(g_1, \ldots, g_c)$,
chaque étape diminue la dimension d'au plus $1$. Par hypothèse,
chaque étape diminue donc la dimension d'exactement $1$. Nous pouvons ainsi
appliquer successivement le lemme \ref{algebra-lemma-CM-one-g} aux modules
$M/(g_1, \ldots, g_i)$ et à l'élément $g_{i + 1}$.

\medskip\noindent
Pour prolonger $g_1, \ldots, g_c$ d'un élément si $c < d$, il suffit de
choisir un élément $g_{c + 1} \in \mathfrak m$ qui n'appartient
à aucun des idéaux premiers minimaux, en nombre fini, de $Z \cap V(g_1, \ldots, g_c)$,
en utilisant le lemme \ref{algebra-lemma-silly}.
\end{proof}

\noindent
La proposition \ref{algebra-proposition-CM-module} étant démontrée, poursuivons le
développement de la théorie standard.

\begin{lemma}
\label{algebra-lemma-nonzerodivisor-on-CM}
Soit $R$ un anneau local noethérien d'idéal maximal $\mathfrak m$.
Soit $M$ un $R$-module de type fini. Soit $x \in \mathfrak m$ un
non-diviseur de zéro sur $M$. Alors $M$ est de Cohen-Macaulay si et seulement
si $M/xM$ est de Cohen-Macaulay.
\end{lemma}

\begin{proof}
Par le lemme \ref{algebra-lemma-depth-drops-by-one}, nous avons
$\text{depth}(M/xM) = \text{depth}(M)-1$.
Par le lemme \ref{algebra-lemma-one-equation-module},
nous avons $\dim(\text{Supp}(M/xM)) = \dim(\text{Supp}(M)) - 1$.
\end{proof}

\begin{lemma}
\label{algebra-lemma-CM-over-quotient}
Soit $R \to S$ un homomorphisme surjectif d'anneaux locaux noethériens.
Soit $N$ un $S$-module de type fini. Alors $N$ est de Cohen-Macaulay comme $S$-module
si et seulement si $N$ est de Cohen-Macaulay comme $R$-module.
\end{lemma}

\begin{proof}
Démonstration omise.
\end{proof}

\begin{lemma}
\label{algebra-lemma-CM-ass-minimal-support}
\begin{reference}
\cite[Chapter 0, Proposition 16.5.4]{EGA}
\end{reference}
Soit $R$ un anneau local noethérien. Soit $M$ un $R$-module de Cohen-Macaulay
de type fini. Si $\mathfrak p \in \text{Ass}(M)$, alors
$\dim(R/\mathfrak p) = \dim(\text{Supp}(M))$ et $\mathfrak p$
est un idéal premier minimal du support de $M$.
En particulier, $M$ n'a aucun idéal premier associé immergé.
\end{lemma}

\begin{proof}
Par le lemme \ref{algebra-lemma-depth-dim-associated-primes}, nous avons
$\text{depth}(M) \leq \dim(R/\mathfrak p)$.
Bien sûr, $\dim(R/\mathfrak p) \leq \dim(\text{Supp}(M))$,
car $\mathfrak p \in \text{Supp}(M)$ (lemme \ref{algebra-lemma-ass-support}).
Nous avons donc égalité dans les deux inégalités, puisque $M$ est de Cohen-Macaulay.
Alors $\mathfrak p$ doit être minimal dans $\text{Supp}(M)$, sans quoi
nous aurions $\dim(R/\mathfrak p) < \dim(\text{Supp}(M))$.
Enfin, les idéaux premiers minimaux du support de $M$ sont égaux
aux éléments minimaux de $\text{Ass}(M)$
(proposition \ref{algebra-proposition-minimal-primes-associated-primes}) ;
ainsi, $M$ n'a aucun idéal premier associé immergé
(définition \ref{algebra-definition-embedded-primes}).
\end{proof}

\begin{definition}
\label{algebra-definition-maximal-CM}
Soit $R$ un anneau local noethérien.
Un module $M$ de type fini sur $R$ est appelé {\it module de Cohen-Macaulay
maximal} si $\text{depth}(M) = \dim(R)$.
\end{definition}

\noindent
Autrement dit, un module de Cohen-Macaulay maximal sur un anneau local
noethérien est un module de type fini ayant la plus grande profondeur possible sur cet anneau.
De manière équivalente, un module de Cohen-Macaulay maximal sur un anneau local
noethérien $R$ est un module de Cohen-Macaulay dont la dimension est égale à celle
de l'anneau. En particulier, si $M$ est un $R$-module de Cohen-Macaulay tel que
$\Spec(R) = \text{Supp}(M)$, alors $M$ est de Cohen-Macaulay maximal.
Les deux lemmes suivants portent donc sur les modules de Cohen-Macaulay maximaux.

\begin{lemma}
\label{algebra-lemma-maximal-chain-maximal-CM}
\begin{slogan}
Dans un anneau local de Cohen-Macaulay, toute chaîne maximale d'idéaux premiers est de
longueur égale à la dimension.
\end{slogan}
Soit $R$ un anneau local noethérien. Supposons qu'il existe un
module de Cohen-Macaulay $M$ tel que $\Spec(R) = \text{Supp}(M)$.
Alors toute chaîne maximale d'idéaux premiers $\mathfrak p_0 \subset
\mathfrak p_1 \subset \ldots \subset \mathfrak p_n$
est de longueur $n = \dim(R)$.
\end{lemma}

\begin{proof}
Nous allons le démontrer par récurrence sur $\dim(R)$. Si $\dim(R) = 0$,
l'assertion est claire. Supposons que $\dim(R) > 0$. Alors $n > 0$.
Choisissons un élément $x \in \mathfrak p_1$, avec $x$ n'appartenant
à aucun des idéaux premiers minimaux de $R$, et en particulier
$x \not \in \mathfrak p_0$. (Voir le lemme \ref{algebra-lemma-silly}.)
Alors $\dim(R/xR) = \dim(R) - 1$ par le lemme \ref{algebra-lemma-one-equation}.
Le module $M/xM$ est de Cohen-Macaulay sur $R/xR$ par la
proposition \ref{algebra-proposition-CM-module} et le
lemme \ref{algebra-lemma-CM-over-quotient}.
Le support de $M/xM$ est $\Spec(R/xR)$ par le
lemme \ref{algebra-lemma-support-quotient}.
Après avoir remplacé $x$ par $x^n$ pour un certain $n$,
nous pouvons supposer que $\mathfrak p_1$ est un idéal premier associé à $M/xM$, voir le
lemme \ref{algebra-lemma-inherit-minimal-primes}.
Par le lemme \ref{algebra-lemma-CM-ass-minimal-support},
nous concluons que $\mathfrak p_1/(x)$ est un idéal premier minimal de $R/xR$.
Il s'ensuit que la chaîne
$\mathfrak p_1/(x) \subset \ldots \subset \mathfrak p_n/(x)$
est une chaîne maximale d'idéaux premiers de $R/xR$.
Par récurrence, nous trouvons que cette chaîne
est de longueur $\dim(R/xR) = \dim(R) - 1$, comme souhaité.
\end{proof}

\begin{lemma}
\label{algebra-lemma-dim-formula-maximal-CM}
Supposons que $R$ soit un anneau local noethérien. Supposons qu'il existe un
module de Cohen-Macaulay $M$ tel que $\Spec(R) = \text{Supp}(M)$. Alors, pour
un idéal premier $\mathfrak p \subset R$, nous avons
$$
\dim(R) = \dim(R_{\mathfrak p}) + \dim(R/\mathfrak p).
$$
\end{lemma}

\begin{proof}
Cela découle immédiatement du lemme \ref{algebra-lemma-maximal-chain-maximal-CM}.
\end{proof}

\begin{lemma}
\label{algebra-lemma-localize-CM-module}
Supposons que $R$ soit un anneau local noethérien. Soit $M$ un module de Cohen-Macaulay
sur $R$. Pour tout idéal premier $\mathfrak p \subset R$, le
module $M_{\mathfrak p}$ est de Cohen-Macaulay sur $R_\mathfrak p$.
\end{lemma}

\begin{proof}
Nous pouvons supposer, et supposons, que $\mathfrak p \not = \mathfrak m$ et que $M$ est non nul.
Choisissons une chaîne maximale d'idéaux premiers $\mathfrak p = \mathfrak p_c \subset
\mathfrak p_{c - 1} \subset \ldots \subset \mathfrak p_1 \subset \mathfrak m$.
Si nous prouvons le résultat pour $M_{\mathfrak p_1}$ sur $R_{\mathfrak p_1}$,
alors le lemme découlera d'une récurrence sur $c$. Nous pouvons donc supposer qu'il
n'existe aucun idéal premier strictement compris entre $\mathfrak p$ et $\mathfrak m$.
Notons que
$\dim(\text{Supp}(M_\mathfrak p)) \leq \dim(\text{Supp}(M)) - 1$,
car toute chaîne d'idéaux premiers du support de $M_\mathfrak p$
peut être prolongée d'un idéal premier supplémentaire (à savoir $\mathfrak m$) dans le
support de $M$. D'autre part, nous avons
$\text{depth}(M_\mathfrak p) \geq \text{depth}(M) - \dim(R/\mathfrak p) =
\text{depth}(M) - 1$ par le lemme \ref{algebra-lemma-depth-localization} et notre
choix de $\mathfrak p$.
Ainsi, $\text{depth}(M_\mathfrak p) \geq \dim(\text{Supp}(M_\mathfrak p))$,
comme souhaité (l'autre inégalité est le lemme \ref{algebra-lemma-bound-depth}).
\end{proof}

\begin{definition}
\label{algebra-definition-module-CM}
Soit $R$ un anneau noethérien. Soit $M$ un $R$-module de type fini.
Nous disons que $M$ est {\it de Cohen-Macaulay} si $M_\mathfrak p$ est un module
de Cohen-Macaulay sur $R_\mathfrak p$ pour tout idéal premier $\mathfrak p$ de $R$.
\end{definition}

\noindent
Par le lemme \ref{algebra-lemma-localize-CM-module}, il suffit de le vérifier
aux idéaux maximaux de $R$.

\begin{lemma}
\label{algebra-lemma-maximal-CM-polynomial-algebra}
Soit $R$ un anneau noethérien. Soit $M$ un module de Cohen-Macaulay
sur $R$. Alors $M \otimes_R R[x_1, \ldots, x_n]$ est un module
de Cohen-Macaulay sur $R[x_1, \ldots, x_n]$.
\end{lemma}

\begin{proof}
Par récurrence sur le nombre de variables, il suffit de le démontrer pour
$M[x] = M \otimes_R R[x]$ sur $R[x]$. Soit $\mathfrak m \subset R[x]$
un idéal maximal, et soit $\mathfrak p = R \cap \mathfrak m$.
Soit $f_1, \ldots, f_d$ une suite $M_\mathfrak p$-régulière dans l'idéal
maximal de $R_{\mathfrak p}$, de longueur $d = \dim(\text{Supp}(M_{\mathfrak p}))$.
Notons que, puisque $R[x]$ est plat sur $R$, le localisé
$R[x]_{\mathfrak m}$ est plat sur $R_{\mathfrak p}$.
Ainsi, par le lemme \ref{algebra-lemma-flat-increases-depth}, la suite
$f_1, \ldots, f_d$ est une suite $M[x]_{\mathfrak m}$-régulière de longueur $d$
dans $R[x]_{\mathfrak m}$. Le quotient
$$
Q = M[x]_{\mathfrak m}/(f_1, \ldots, f_d)M[x]_{\mathfrak m} =
M_{\mathfrak p}/(f_1, \ldots, f_d)M_{\mathfrak p}
\otimes_{R_\mathfrak p} R[x]_{\mathfrak m}
$$
a pour support les idéaux premiers situés au-dessus de $\mathfrak p$,
car $R_\mathfrak p \to R[x]_\mathfrak m$ est plat et
le support de $M_{\mathfrak p}/(f_1, \ldots, f_d)M_{\mathfrak p}$
est égal à $\{\mathfrak p\}$ (détails omis ; indication : cela découle des
lemmes \ref{algebra-lemma-annihilator-flat-base-change} et
\ref{algebra-lemma-support-closed}). La dimension vaut donc $1$.
Pour achever la démonstration, il suffit de trouver un $f \in \mathfrak m$
qui soit un non-diviseur de zéro sur $Q$. Comme $\mathfrak m$ est
un idéal maximal, l'extension de corps
$\kappa(\mathfrak m)/\kappa(\mathfrak p)$
est finie (théorème \ref{algebra-theorem-nullstellensatz}).
Nous pouvons donc trouver $f \in \mathfrak m$ qui,
vu comme un polynôme en $x$, a son coefficient
dominant hors de $\mathfrak p$. Un tel $f$ agit comme
un non-diviseur de zéro sur
$$
M_{\mathfrak p}/(f_1, \ldots, f_d)M_{\mathfrak p} \otimes_R R[x] =
\bigoplus\nolimits_{n \geq 0}
M_{\mathfrak p}/(f_1, \ldots, f_d)M_{\mathfrak p} \cdot x^n
$$
et agit donc comme un non-diviseur de zéro sur $Q$.
\end{proof}















\section{Anneaux de Cohen-Macaulay}
\label{algebra-section-CM-ring}

\noindent
La plupart des résultats de cette section sont des cas particuliers de ceux
de la section \ref{algebra-section-CM}.

\begin{definition}
\label{algebra-definition-local-ring-CM}
Un anneau local noethérien $R$ est dit {\it de Cohen-Macaulay}
s'il est de Cohen-Macaulay comme module sur lui-même.
\end{definition}

\noindent
Notons que cela équivaut à demander l'existence
d'une suite $R$-régulière $x_1, \ldots, x_d$ de l'idéal
maximal telle que $R/(x_1, \ldots, x_d)$ soit de dimension $0$.
Nous dirons généralement simplement « suite régulière », et non
« suite $R$-régulière ».

\begin{lemma}
\label{algebra-lemma-reformulate-CM}
\begin{slogan}
Dans les anneaux locaux de Cohen-Macaulay, les suites régulières sont caractérisées
par le fait de découper un objet de la bonne dimension.
\end{slogan}
Soit $R$ un anneau local noethérien de Cohen-Macaulay, d'idéal
maximal $\mathfrak m $. Soient $x_1, \ldots, x_c \in \mathfrak m$ des
éléments. Alors
$$
x_1, \ldots, x_c
\text{ est une suite régulière }
\Leftrightarrow
\dim(R/(x_1, \ldots, x_c)) = \dim(R) - c
$$
Dans ce cas,
$x_1, \ldots, x_c$ peut être prolongée en
une suite régulière de longueur $\dim(R)$, et chaque quotient
$R/(x_1, \ldots, x_i)$ est un anneau de Cohen-Macaulay de dimension
$\dim(R) - i$.
\end{lemma}

\begin{proof}
C'est un cas particulier de la proposition \ref{algebra-proposition-CM-module}.
\end{proof}

\begin{lemma}
\label{algebra-lemma-maximal-chain-CM}
Soit $R$ un anneau local noethérien.
Supposons que $R$ soit de Cohen-Macaulay de dimension $d$.
Toute chaîne maximale d'idéaux premiers $\mathfrak p_0 \subset
\mathfrak p_1 \subset \ldots \subset \mathfrak p_n$
est de longueur $n = d$.
\end{lemma}

\begin{proof}
C'est un cas particulier du lemme \ref{algebra-lemma-maximal-chain-maximal-CM}.
\end{proof}

\begin{lemma}
\label{algebra-lemma-CM-dim-formula}
Supposons que $R$ soit un anneau local noethérien de Cohen-Macaulay de dimension $d$.
Pour tout idéal premier $\mathfrak p \subset R$, nous avons
$$
\dim(R) = \dim(R_{\mathfrak p}) + \dim(R/\mathfrak p).
$$
\end{lemma}

\begin{proof}
Cela découle immédiatement du lemme \ref{algebra-lemma-maximal-chain-CM}.
(C'est aussi un cas particulier du lemme \ref{algebra-lemma-dim-formula-maximal-CM}.)
\end{proof}

\begin{lemma}
\label{algebra-lemma-localize-CM}
Supposons que $R$ soit un anneau local de Cohen-Macaulay.
Pour tout idéal premier $\mathfrak p \subset R$, l'anneau
$R_{\mathfrak p}$ est lui aussi de Cohen-Macaulay.
\end{lemma}

\begin{proof}
C'est un cas particulier du lemme \ref{algebra-lemma-localize-CM-module}.
\end{proof}

\begin{definition}
\label{algebra-definition-ring-CM}
Un anneau noethérien $R$ est dit {\it de Cohen-Macaulay} si tous
ses anneaux locaux sont de Cohen-Macaulay.
\end{definition}

\begin{lemma}
\label{algebra-lemma-CM-polynomial-algebra}
Supposons que $R$ soit un anneau noethérien de Cohen-Macaulay.
Toute algèbre polynomiale sur $R$ est de Cohen-Macaulay.
\end{lemma}

\begin{proof}
C'est un cas particulier du lemme \ref{algebra-lemma-maximal-CM-polynomial-algebra}.
\end{proof}

\begin{lemma}
\label{algebra-lemma-dimension-shift}
Soit $R$ un anneau local noethérien de Cohen-Macaulay de dimension $d$.
Soit $0 \to K \to R^{\oplus n} \to M \to 0$ une suite exacte de
$R$-modules. Alors soit $M = 0$, soit $\text{depth}(K) > \text{depth}(M)$, soit
$\text{depth}(K) = \text{depth}(M) = d$.
\end{lemma}

\begin{proof}
C'est un cas particulier du lemme \ref{algebra-lemma-depth-in-ses}.
\end{proof}

\begin{lemma}
\label{algebra-lemma-mcm-resolution}
Soit $R$ un anneau local noethérien de Cohen-Macaulay de dimension $d$.
Soit $M$ un $R$-module de type fini et de profondeur $e$.
Il existe un complexe exact
$$
0 \to K \to F_{d-e-1} \to \ldots \to F_0 \to M \to 0
$$
où chaque $F_i$ est libre de rang fini et $K$ est de Cohen-Macaulay maximal.
\end{lemma}

\begin{proof}
Cela découle immédiatement de la définition et du lemme \ref{algebra-lemma-dimension-shift}.
\end{proof}

\begin{lemma}
\label{algebra-lemma-find-sequence-image-regular}
Soit $\varphi : A \to B$ un homomorphisme d'anneaux locaux.
Supposons que $B$ soit noethérien et de Cohen-Macaulay, et que
$\mathfrak m_B = \sqrt{\varphi(\mathfrak m_A) B}$. Alors il existe
une suite d'éléments $f_1, \ldots, f_{\dim(B)}$ de $A$
telle que $\varphi(f_1), \ldots, \varphi(f_{\dim(B)})$ soit une
suite régulière dans $B$.
\end{lemma}

\begin{proof}
Par récurrence sur $\dim(B)$, il suffit de prouver ceci : si $\dim(B) \geq 1$, alors nous
pouvons trouver un élément $f$ de $A$ qui s'envoie sur un non-diviseur de zéro dans $B$.
Par le
lemme \ref{algebra-lemma-reformulate-CM},
il suffit de trouver $f \in A$ dont l'image dans $B$ n'est contenue dans aucun
des idéaux premiers minimaux, en nombre fini, $\mathfrak q_1, \ldots, \mathfrak q_r$
de $B$. L'hypothèse
$\mathfrak m_B = \sqrt{\varphi(\mathfrak m_A) B}$
montre que $\mathfrak m_A \not \subset \varphi^{-1}(\mathfrak q_i)$.
Nous pouvons donc trouver $f$ grâce au
lemme \ref{algebra-lemma-silly}.
\end{proof}








\section{Anneaux caténaires}
\label{algebra-section-catenary}

\noindent
Comparer avec Topologie, section \ref{topology-section-catenary-spaces}.

\begin{definition}
\label{algebra-definition-catenary}
On dit qu'un anneau $R$ est {\it caténaire} si, pour toute paire d'idéaux premiers
$\mathfrak p \subset \mathfrak q$, il existe un entier qui majore les
longueurs de toutes les chaînes finies d'idéaux premiers
$\mathfrak p = \mathfrak p_0 \subset \mathfrak p_1 \subset \ldots \subset
\mathfrak p_e = \mathfrak q$ et si toutes les chaînes maximales de ce type ont la même longueur.
\end{definition}

\begin{lemma}
\label{algebra-lemma-catenary}
Un anneau $R$ est caténaire si et seulement si l'espace topologique
$\Spec(R)$ est caténaire (voir
Topologie, définition \ref{topology-definition-catenary}).
\end{lemma}

\begin{proof}
Cela résulte immédiatement de la définition et de la caractérisation des
fermés irréductibles dans le lemme \ref{algebra-lemma-irreducible}.
\end{proof}

\noindent
En général, il n'est pas vrai qu'une $R$-algèbre de type fini
soit caténaire lorsque $R$ l'est. Nous posons donc la
définition suivante.

\begin{definition}
\label{algebra-definition-universally-catenary}
Un anneau noethérien $R$ est dit {\it universellement caténaire}
si toute $R$-algèbre de type fini est caténaire.
\end{definition}

\noindent
Nous nous restreignons aux anneaux noethériens, car il n'est pas clair que
cette définition soit la bonne pour les anneaux non noethériens.
D'après le lemme \ref{algebra-lemma-quotient-catenary}, pour vérifier qu'un anneau
noethérien $R$ est universellement caténaire, il suffit de vérifier que
toute algèbre polynomiale $R[x_1, \ldots, x_n]$ est caténaire.

\begin{lemma}
\label{algebra-lemma-localization-catenary}
Toute localisation d'un anneau caténaire est caténaire.
Toute localisation d'un anneau noethérien universellement caténaire
est universellement caténaire.
\end{lemma}

\begin{proof}
Soient $A$ un anneau et $S \subset A$ une partie multiplicative.
La description de $\Spec(S^{-1}A)$ dans le lemme \ref{algebra-lemma-spec-localization}
montre que, si $A$ est caténaire, alors $S^{-1}A$ l'est aussi. Si $S^{-1}A \to C$
est de type fini, alors $C = S^{-1}B$ pour un homomorphisme d'anneaux de type fini
$A \to B$. Ainsi, si $A$ est noethérien et universellement caténaire, alors
$B$ est caténaire et nous voyons que $C$ est également caténaire. Cela prouve le lemme.
\end{proof}

\begin{lemma}
\label{algebra-lemma-universally-catenary}
Soit $A$ un anneau noethérien universellement caténaire.
Toute $A$-algèbre essentiellement de type fini sur $A$
est universellement caténaire.
\end{lemma}

\begin{proof}
Si $B$ est une $A$-algèbre de type fini, alors $B$ est noethérien
d'après le lemme \ref{algebra-lemma-Noetherian-permanence}. Toute
$B$-algèbre de type fini est une $A$-algèbre de type fini et est donc caténaire
puisque, par hypothèse, $A$ est universellement caténaire. Ainsi, $B$
est universellement caténaire. Toute localisation de $B$ est universellement
caténaire d'après le lemme \ref{algebra-lemma-localization-catenary}, ce qui
achève la preuve.
\end{proof}

\begin{lemma}
\label{algebra-lemma-catenary-check-local}
Soit $R$ un anneau. Les assertions suivantes sont équivalentes :
\begin{enumerate}
\item $R$ est caténaire,
\item $R_\mathfrak p$ est caténaire pour tout idéal premier $\mathfrak p$,
\item $R_\mathfrak m$ est caténaire pour tout idéal maximal $\mathfrak m$.
\end{enumerate}
Supposons $R$ noethérien. Les assertions suivantes sont équivalentes :
\begin{enumerate}
\item $R$ est universellement caténaire,
\item $R_\mathfrak p$ est universellement caténaire pour tout idéal premier
$\mathfrak p$,
\item $R_\mathfrak m$ est universellement caténaire pour tout idéal maximal
$\mathfrak m$.
\end{enumerate}
\end{lemma}

\begin{proof}
L'implication (1) $\Rightarrow$ (2) découle du
lemme \ref{algebra-lemma-localization-catenary} dans les deux cas.
L'implication (2) $\Rightarrow$ (3) est immédiate dans les deux cas.
Supposons que $R_\mathfrak m$ soit caténaire pour tout idéal maximal
$\mathfrak m$ de $R$. Si $\mathfrak p \subset \mathfrak q$ sont des idéaux premiers
de $R$, choisissons un idéal maximal vérifiant $\mathfrak q \subset \mathfrak m$.
Les chaînes d'idéaux premiers entre $\mathfrak p$ et $\mathfrak q$ sont
en correspondance bijective avec les chaînes d'idéaux premiers entre
$\mathfrak pR_\mathfrak m$ et $\mathfrak qR_\mathfrak m$ ; nous voyons donc que
$R$ est caténaire. Supposons $R$ noethérien et $R_\mathfrak m$
universellement caténaire pour tout idéal maximal $\mathfrak m$ de $R$.
Soit $R \to S$ un homomorphisme d'anneaux de type fini. Soit $\mathfrak q$ un idéal premier
de $S$ au-dessus de l'idéal premier $\mathfrak p \subset R$.
Choisissons un idéal maximal vérifiant $\mathfrak p \subset \mathfrak m$ dans $R$.
Alors $R_\mathfrak p$ est une localisation de $R_\mathfrak m$, donc est
universellement caténaire d'après le lemme \ref{algebra-lemma-localization-catenary}.
Ainsi, $S_\mathfrak p$ est caténaire comme anneau de type fini sur $R_\mathfrak p$.
Par conséquent, $S_\mathfrak q$ est caténaire comme localisation. Ainsi, $S$ est caténaire
d'après le premier cas traité ci-dessus.
\end{proof}

\begin{lemma}
\label{algebra-lemma-quotient-catenary}
Tout quotient d'un anneau caténaire est caténaire.
Tout quotient d'un anneau noethérien universellement caténaire est
universellement caténaire.
\end{lemma}

\begin{proof}
Soient $A$ un anneau et $I \subset A$ un idéal.
La description de $\Spec(A/I)$ dans le lemme \ref{algebra-lemma-spec-closed}
montre que, si $A$ est caténaire, alors $A/I$ l'est aussi.
La seconde assertion est un cas particulier du
lemme \ref{algebra-lemma-universally-catenary}.
\end{proof}

\begin{lemma}
\label{algebra-lemma-catenary-check-irreducible}
Soit $R$ un anneau noethérien.
\begin{enumerate}
\item $R$ est caténaire si et seulement si $R/\mathfrak p$ est caténaire
pour tout idéal premier minimal $\mathfrak p$.
\item $R$ est universellement caténaire si et seulement si $R/\mathfrak p$ est
universellement caténaire pour tout idéal premier minimal $\mathfrak p$.
\end{enumerate}
\end{lemma}

\begin{proof}
Si $\mathfrak a \subset \mathfrak b$ est une inclusion d'idéaux premiers de $R$,
alors nous pouvons trouver un idéal premier minimal $\mathfrak p \subset \mathfrak a$,
et la première assertion est claire. Nous omettons la preuve de la seconde.
\end{proof}

\begin{lemma}
\label{algebra-lemma-CM-ring-catenary}
Un anneau de Cohen-Macaulay noethérien est universellement caténaire.
Plus généralement, si $R$ est un anneau noethérien et $M$ un
$R$-module de Cohen-Macaulay tel que $\text{Supp}(M) = \Spec(R)$,
alors $R$ est universellement caténaire.
\end{lemma}

\begin{proof}
Puisqu'une algèbre polynomiale sur $R$ est de Cohen-Macaulay
d'après le lemme \ref{algebra-lemma-CM-polynomial-algebra},
il suffit de montrer qu'un anneau de Cohen-Macaulay est
caténaire.
Soit $R$ un anneau de Cohen-Macaulay et soient $\mathfrak p \subset \mathfrak q$
des idéaux premiers de $R$. Par définition, $R_{\mathfrak q}$ et
$R_{\mathfrak p}$ sont de Cohen-Macaulay.
Prenons une chaîne maximale d'idéaux premiers
$\mathfrak p = \mathfrak p_0 \subset \mathfrak p_1 \subset
\ldots \subset \mathfrak p_n = \mathfrak q$.
Choisissons ensuite une chaîne maximale d'idéaux premiers
$\mathfrak q_0 \subset \mathfrak q_1 \subset \ldots \subset
\mathfrak q_m = \mathfrak p$. D'après le
lemme \ref{algebra-lemma-maximal-chain-CM},
nous avons $n + m = \dim(R_{\mathfrak q})$. Et nous avons
$m = \dim(R_{\mathfrak p})$ d'après le même lemme.
Ainsi, $n = \dim(R_{\mathfrak q}) - \dim(R_{\mathfrak p})$
est indépendant des choix.

\medskip\noindent
Pour prouver l'assertion plus générale, raisonnons exactement comme ci-dessus, mais
en utilisant les lemmes \ref{algebra-lemma-maximal-CM-polynomial-algebra}
et \ref{algebra-lemma-maximal-chain-maximal-CM}.
\end{proof}

\begin{lemma}
\label{algebra-lemma-catenary-Noetherian-local}
Soit $(A, \mathfrak m)$ un anneau local noethérien. Les assertions suivantes sont équivalentes :
\begin{enumerate}
\item $A$ est caténaire, et
\item $\mathfrak p \mapsto \dim(A/\mathfrak p)$ est une fonction de dimension
sur $\Spec(A)$.
\end{enumerate}
\end{lemma}

\begin{proof}
Si $A$ est caténaire, alors $\Spec(A)$ possède une fonction de dimension $\delta$ d'après
Topologie, lemme \ref{topology-lemma-locally-dimension-function}
(et le lemme \ref{algebra-lemma-catenary}). Nous pouvons supposer que $\delta(\mathfrak m) = 0$.
Nous voyons alors que
$$
\delta(\mathfrak p) = \text{codim}(V(\mathfrak m), V(\mathfrak p)) =
\dim(A/\mathfrak p)
$$
d'après Topologie, lemme \ref{topology-lemma-dimension-function-catenary}.
Nous voyons ainsi que (1) implique (2). L'implication réciproque
découle de
Topologie, lemme \ref{topology-lemma-dimension-function-catenary}
également.
\end{proof}














\section{Anneaux locaux réguliers}
\label{algebra-section-regular}

\noindent
Les anneaux locaux réguliers sont définis dans la définition \ref{algebra-definition-regular-local}.
Il n'est pas si facile de montrer que toutes les localisations premières d'un anneau local
régulier sont régulières. En fait, une bonne partie des résultats développés jusqu'ici
vise à prouver ce fait. Voir la
proposition \ref{algebra-proposition-finite-gl-dim-regular}, et
remonter la chaîne des références.

\begin{lemma}
\label{algebra-lemma-regular-graded}
Soit $(R, \mathfrak m, \kappa)$ un anneau local régulier de dimension $d$.
L'anneau gradué $\bigoplus \mathfrak m^n / \mathfrak m^{n + 1}$
est isomorphe à l'algèbre polynomiale graduée
$\kappa[X_1, \ldots, X_d]$.
\end{lemma}

\begin{proof}
Soit $x_1, \ldots, x_d$ un système minimal de générateurs
de l'idéal maximal $\mathfrak m$, voir la
définition \ref{algebra-definition-regular-local}.
Il existe une surjection $\kappa[X_1, \ldots, X_d]
\to \bigoplus \mathfrak m^n/\mathfrak m^{n + 1}$,
qui envoie $X_i$ sur la classe de $x_i$ dans $\mathfrak m/\mathfrak m^2$.
Puisque $d(R) = d$ d'après la
proposition \ref{algebra-proposition-dimension},
nous savons que le polynôme
numérique $n \mapsto \dim_\kappa \mathfrak m^n/\mathfrak m^{n + 1}$
est de degré $d - 1$. D'après le lemme \ref{algebra-lemma-quotient-smaller-d}, nous
concluons que la surjection $\kappa[X_1, \ldots, X_d]
\to \bigoplus \mathfrak m^n/\mathfrak m^{n + 1}$ est un isomorphisme.
\end{proof}

\begin{lemma}
\label{algebra-lemma-regular-domain}
Tout anneau local régulier est intègre.
\end{lemma}

\begin{proof}
Nous allons utiliser l'égalité $\bigcap \mathfrak m^n = 0$
donnée par le lemme \ref{algebra-lemma-intersect-powers-ideal-module-zero}.
Soient $f, g \in R$ tels que $fg = 0$.
Supposons que $f \in \mathfrak m^a$ et
$g \in \mathfrak m^b$, avec $a, b$ maximaux.
Puisque $fg = 0 \in \mathfrak m^{a + b + 1}$,
le résultat du lemme \ref{algebra-lemma-regular-graded} montre
que soit $f \in \mathfrak m^{a + 1}$, soit
$g \in \mathfrak m^{b + 1}$. Contradiction.
\end{proof}

\begin{lemma}
\label{algebra-lemma-regular-ring-CM}
Soit $R$ un anneau local régulier et soit
$x_1, \ldots, x_d$ un système minimal de générateurs
de l'idéal maximal $\mathfrak m$. Alors
$x_1, \ldots, x_d$ est une suite régulière, et
chaque $R/(x_1, \ldots, x_c)$ est un anneau local régulier
de dimension $d - c$. En particulier, $R$ est de Cohen-Macaulay.
\end{lemma}

\begin{proof}
Remarquons que $R/x_1R$ est un anneau local noethérien de dimension $\geq d - 1$
d'après le lemme \ref{algebra-lemma-one-equation}, tandis que $x_2, \ldots, x_d$
engendrent l'idéal maximal. Il s'agit donc d'un anneau local régulier par définition.
Puisque $R$ est intègre d'après le lemme \ref{algebra-lemma-regular-domain},
$x_1$ est un non-diviseur de zéro.
\end{proof}

\begin{lemma}
\label{algebra-lemma-regular-quotient-regular}
Soit $R$ un anneau local régulier. Soit $I \subset R$ un idéal
tel que $R/I$ soit lui aussi un anneau local régulier. Alors
il existe un système minimal de générateurs $x_1, \ldots, x_d$
de l'idéal maximal $\mathfrak m$ de $R$ tel que
$I = (x_1, \ldots, x_c)$ pour un certain $0 \leq c \leq d$.
\end{lemma}

\begin{proof}
Écrivons $\dim(R) = d$ et $\dim(R/I) = d - c$.
Notons $\overline{\mathfrak m} = \mathfrak m/I$ l'idéal
maximal de $R/I$. Soit $\kappa = R/\mathfrak m$. Nous avons
$$
\dim_\kappa((I + \mathfrak m^2)/\mathfrak m^2) =
\dim_\kappa(\mathfrak m/\mathfrak m^2)
- \dim(\overline{\mathfrak m}/\overline{\mathfrak m}^2) = d - (d - c) = c
$$
par définition d'un anneau local régulier. Nous pouvons donc choisir
$x_1, \ldots, x_c \in I$ dont les images dans $\mathfrak m/\mathfrak m^2$
sont linéairement indépendantes, et les compléter par
$x_{c + 1}, \ldots, x_d$ pour obtenir un système minimal de générateurs de
$\mathfrak m$. Le morphisme induit $R/(x_1, \ldots, x_c) \to R/I$ est une
surjection entre anneaux locaux réguliers de même dimension
(lemme \ref{algebra-lemma-regular-ring-CM}). Il s'ensuit que
le noyau est nul, c'est-à-dire que $I = (x_1, \ldots, x_c)$. En effet, sinon,
nous aurions $\dim(R/I) < \dim(R/(x_1, \ldots, x_c))$ d'après les
lemmes \ref{algebra-lemma-regular-domain} et \ref{algebra-lemma-one-equation}.
\end{proof}

\begin{lemma}
\label{algebra-lemma-free-mod-x}
Soit $R$ un anneau local noethérien.
Soit $x \in \mathfrak m$.
Soit $M$ un $R$-module fini tel que
$x$ soit un non-diviseur de zéro sur $M$ et que
$M/xM$ soit libre sur $R/xR$.
Alors $M$ est libre sur $R$.
\end{lemma}

\begin{proof}
Soient $m_1, \ldots, m_r$ des éléments de $M$ qui s'envoient sur
une $R/xR$-base de $M/xM$. D'après le lemme de Nakayama \ref{algebra-lemma-NAK},
$m_1, \ldots, m_r$ engendrent $M$. Si $\sum a_i m_i = 0$
est une relation, alors $a_i \in xR$ pour tout $i$. Ainsi,
$a_i = b_i x$ pour certains $b_i \in R$. Par conséquent,
le noyau $K$ de $R^r \to M$ vérifie $xK = K$
et est donc nul d'après le lemme de Nakayama.
\end{proof}

\begin{lemma}
\label{algebra-lemma-regular-mcm-free}
Soit $R$ un anneau local régulier.
Tout module de Cohen-Macaulay maximal sur $R$ est libre.
\end{lemma}

\begin{proof}
Soit $M$ un module de Cohen-Macaulay maximal sur $R$.
Soit $x \in \mathfrak m$ un élément d'une suite régulière
engendrant $\mathfrak m$. Alors $x$ est un non-diviseur de zéro
sur $M$ d'après la proposition \ref{algebra-proposition-CM-module}, et
$M/xM$ est un module de Cohen-Macaulay maximal sur $R/xR$.
Par récurrence sur $\dim(R)$, nous voyons que $M/xM$ est libre.
Le lemme \ref{algebra-lemma-free-mod-x} permet de conclure.
\end{proof}

\begin{lemma}
\label{algebra-lemma-regular-mod-x}
Supposons que $R$ soit un anneau local noethérien.
Soit $x \in \mathfrak m$ un non-diviseur de zéro
tel que $R/xR$ soit un anneau local régulier. Alors $R$ est un anneau local régulier.
Plus généralement, si $x_1, \ldots, x_r$ est une suite régulière dans $R$
telle que $R/(x_1, \ldots, x_r)$ soit un anneau local régulier, alors
$R$ est un anneau local régulier.
\end{lemma}

\begin{proof}
C'est vrai parce que $x$, avec les relèvements d'un système
minimal de générateurs de l'idéal maximal de $R/xR$, fournit
$\dim(R)$ générateurs de $\mathfrak m$.
Utiliser le lemme \ref{algebra-lemma-one-equation}.
La dernière assertion découle de la première et d'une récurrence.
\end{proof}

\begin{lemma}
\label{algebra-lemma-colimit-regular}
Soit $(R_i, \varphi_{ii'})$ un système inductif d'anneaux locaux dont les
morphismes de transition sont des homomorphismes locaux d'anneaux. Si chaque $R_i$ est un anneau local régulier et si
$R = \colim R_i$ est noethérien, alors $R$ est un anneau local régulier.
\end{lemma}

\begin{proof}
Soit $\mathfrak m \subset R$ l'idéal maximal ; c'est la colimite
des idéaux maximaux $\mathfrak m_i \subset R_i$.
Nous prouvons le lemme par récurrence sur $d = \dim \mathfrak m/\mathfrak m^2$.
Si $d = 0$, alors $R = R/\mathfrak m$ est un corps et $R$ est un anneau local régulier.
Si $d > 0$, choisissons $x \in \mathfrak m$, $x \not \in \mathfrak m^2$.
Pour un certain $i$, nous pouvons trouver $x_i \in \mathfrak m_i$ qui s'envoie sur $x$.
Remarquons que $R/xR = \colim_{i' \geq i} R_{i'}/x_iR_{i'}$ est un anneau
local noethérien. D'après le
lemme \ref{algebra-lemma-regular-ring-CM},
nous voyons que $R_{i'}/x_iR_{i'}$ est un anneau local régulier.
Ainsi, par récurrence, nous voyons
que $R/xR$ est un anneau local régulier. Puisque chaque $R_i$ est intègre
(lemme \ref{algebra-lemma-regular-graded}), nous voyons que $R$ est intègre.
Ainsi, $x$ est un non-diviseur de zéro et nous concluons que $R$ est
un anneau local régulier d'après le lemme \ref{algebra-lemma-regular-mod-x}.
\end{proof}





\section{Épimorphismes d'anneaux}
\label{algebra-section-epimorphism}

\noindent
Dans toute catégorie, il existe une notion d'{\it épimorphisme}.
Une partie de ces résultats provient de \cite{Autour} et de \cite{Mazet}.

\begin{lemma}
\label{algebra-lemma-epimorphism}
Soit $R \to S$ un homomorphisme d'anneaux. Les assertions suivantes sont équivalentes :
\begin{enumerate}
\item $R \to S$ est un épimorphisme,
\item les deux homomorphismes d'anneaux $S \to S \otimes_R S$ sont égaux,
\item l'un ou l'autre des homomorphismes d'anneaux $S \to S \otimes_R S$ est un isomorphisme, et
\item l'homomorphisme d'anneaux $S \otimes_R S \to S$ est un isomorphisme.
\end{enumerate}
\end{lemma}

\begin{proof}
Omis.
\end{proof}

\begin{lemma}
\label{algebra-lemma-composition-epimorphism}
La composée de deux épimorphismes d'anneaux est un épimorphisme.
\end{lemma}

\begin{proof}
Omis. Indication : c'est vrai dans toute catégorie.
\end{proof}

\begin{lemma}
\label{algebra-lemma-base-change-epimorphism}
Si $R \to S$ est un épimorphisme d'anneaux et $R \to R'$ un homomorphisme d'anneaux quelconque,
alors $R' \to R' \otimes_R S$ est un épimorphisme.
\end{lemma}

\begin{proof}
Omis. Indication : c'est vrai dans toute catégorie possédant des sommes amalgamées.
\end{proof}

\begin{lemma}
\label{algebra-lemma-permanence-epimorphism}
Si $A \to B \to C$ sont des homomorphismes d'anneaux et $A \to C$ un épimorphisme, alors
$B \to C$ l'est aussi.
\end{lemma}

\begin{proof}
Omis. Indication : c'est vrai dans toute catégorie.
\end{proof}

\noindent
Cela signifie en particulier que, si $R \to S$ est un épimorphisme d'image
$\overline{R} \subset S$, alors $\overline{R} \to S$ est un épimorphisme.
Ainsi, lorsque nous prouvons des résultats sur les épimorphismes, nous pouvons souvent supposer
que l'homomorphisme est injectif. Le lemme suivant signifie en particulier que
toute localisation est un épimorphisme.

\begin{lemma}
\label{algebra-lemma-epimorphism-local}
Soit $R \to S$ un homomorphisme d'anneaux. Les assertions suivantes sont équivalentes :
\begin{enumerate}
\item $R \to S$ est un épimorphisme, et
\item $R_{\mathfrak p} \to S_{\mathfrak p}$ est un épimorphisme pour
tout idéal premier $\mathfrak p$ de $R$.
\end{enumerate}
\end{lemma}

\begin{proof}
Puisque $S_{\mathfrak p} = R_{\mathfrak p} \otimes_R S$ (voir le
lemme \ref{algebra-lemma-tensor-localization}),
nous voyons que (1) implique (2) d'après le
lemme \ref{algebra-lemma-base-change-epimorphism}.
Réciproquement, supposons que (2) soit vérifiée. Soient $a, b : S \to A$ deux homomorphismes d'anneaux
de $S$ vers un anneau $A$ qui coïncident sur l'homomorphisme $R \to S$. Par hypothèse, nous voyons
que, pour tout idéal premier $\mathfrak p$ de $R$, les homomorphismes induits
$a_{\mathfrak p}, b_{\mathfrak p} : S_{\mathfrak p} \to A_{\mathfrak p}$ sont
égaux. Ainsi, $a = b$ puisque $A \subset \prod_{\mathfrak p} A_{\mathfrak p}$, voir le
lemme \ref{algebra-lemma-characterize-zero-local}.
\end{proof}

\begin{lemma}
\label{algebra-lemma-finite-epimorphism-surjective}
\begin{slogan}
Un homomorphisme d'anneaux est surjectif si et seulement si c'est un épimorphisme fini.
\end{slogan}
Soit $R \to S$ un homomorphisme d'anneaux. Les assertions suivantes sont équivalentes :
\begin{enumerate}
\item $R \to S$ est un épimorphisme fini, et
\item $R \to S$ est surjectif.
\end{enumerate}
\end{lemma}

\begin{proof}
(Ce lemme semble avoir été redémontré de nombreuses fois dans la littérature, et
possède de nombreuses preuves différentes.)
Il est clair qu'un homomorphisme d'anneaux surjectif est un épimorphisme.
Supposons que $R \to S$ soit un homomorphisme d'anneaux fini tel que
$S \otimes_R S \to S$ soit un isomorphisme. Notre but est de montrer que
$R \to S$ est surjectif. Supposons que $S/R$ ne soit pas nul.
La suite exacte $R \to S \to S/R \to 0$
donne une suite exacte
$$
R \otimes_R S \to S \otimes_R S \to S/R \otimes_R S \to 0.
$$
Notre hypothèse implique que la première flèche est un isomorphisme ; ainsi,
nous concluons que $S/R \otimes_R S = 0$. Donc aussi $S/R \otimes_R S/R = 0$. D'après le
lemme \ref{algebra-lemma-trivial-filter-finite-module},
il existe une surjection de $R$-modules $S/R \to R/I$ pour un certain idéal propre
$I \subset R$. Il existe donc une
surjection $S/R \otimes_R S/R \to R/I \otimes_R R/I = R/I \not = 0$,
ce qui est une contradiction.
\end{proof}

\begin{lemma}
\label{algebra-lemma-faithfully-flat-epimorphism}
Un épimorphisme fidèlement plat est un isomorphisme.
\end{lemma}

\begin{proof}
Cela résulte clairement de la partie (3) du
lemme \ref{algebra-lemma-epimorphism},
car l'homomorphisme $S \to S \otimes_R S$ est l'homomorphisme $R \to S$ tensorisé avec $S$.
\end{proof}

\begin{lemma}
\label{algebra-lemma-epimorphism-over-field}
Si $k \to S$ est un épimorphisme et $k$ un corps, alors $S = k$ ou $S = 0$.
\end{lemma}

\begin{proof}
Cela résulte clairement du
lemme \ref{algebra-lemma-faithfully-flat-epimorphism}
(puisque toute algèbre non nulle sur $k$ est fidèlement plate), ou
en montrant directement que $R \to R \otimes_k R$ ne peut pas être
surjectif sauf si $\dim_k(R) \leq 1$.
\end{proof}

\begin{lemma}
\label{algebra-lemma-epimorphism-injective-spec}
Soit $R \to S$ un épimorphisme d'anneaux. Alors
\begin{enumerate}
\item $\Spec(S) \to \Spec(R)$ est injectif, et
\item pour $\mathfrak q \subset S$ au-dessus de $\mathfrak p \subset R$,
nous avons $\kappa(\mathfrak p) = \kappa(\mathfrak q)$.
\end{enumerate}
\end{lemma}

\begin{proof}
Soit $\mathfrak p$ un idéal premier de $R$. La fibre de l'homomorphisme est le spectre
de l'anneau fibre $S \otimes_R \kappa(\mathfrak p)$. D'après le
lemme \ref{algebra-lemma-base-change-epimorphism},
l'homomorphisme $\kappa(\mathfrak p) \to S \otimes_R \kappa(\mathfrak p)$
est un épimorphisme ; ainsi, d'après le
lemme \ref{algebra-lemma-epimorphism-over-field},
nous avons soit $S \otimes_R \kappa(\mathfrak p) = 0$,
soit $S \otimes_R \kappa(\mathfrak p) = \kappa(\mathfrak p)$,
ce qui prouve (1) et (2).
\end{proof}

\begin{lemma}
\label{algebra-lemma-relations}
Soit $R$ un anneau.
Soient $M$, $N$ des $R$-modules.
Soit $\{x_i\}_{i \in I}$ un système de générateurs de $M$.
Soit $\{y_j\}_{j \in J}$ un système de générateurs de $N$.
Soit $\{m_j\}_{j \in J}$ une famille d'éléments de $M$ telle que $m_j = 0$
pour presque tout $j$.
Alors
$$
\sum\nolimits_{j \in J} m_j \otimes y_j = 0 \text{ dans } M \otimes_R N
$$
équivaut à ce qui suit :
il existe des $a_{i, j} \in R$, avec $a_{i, j} = 0$ pour presque toute
paire $(i, j)$, tels que
\begin{align*}
m_j & = \sum\nolimits_{i \in I} a_{i, j} x_i \quad\text{pour tout } j \in J, \\
0 & = \sum\nolimits_{j \in J} a_{i, j} y_j \quad\text{pour tout } i \in I.
\end{align*}
\end{lemma}

\begin{proof}
La suffisance est immédiate. Supposons que
$\sum_{j \in J} m_j \otimes y_j = 0$.
Considérons la suite exacte courte
$$
0 \to K \to \bigoplus\nolimits_{j \in J} R \to N \to 0
$$
dans laquelle le $j$-ième vecteur de base de $\bigoplus\nolimits_{j \in J} R$ s'envoie
sur $y_j$. Tensorisons-la avec $M$ pour obtenir la suite exacte
$$
K \otimes_R M \to \bigoplus\nolimits_{j \in J} M \to N \otimes_R M \to 0.
$$
L'hypothèse implique qu'il existe des éléments $k_i \in K$ tels que
$\sum k_i \otimes x_i$ s'envoie sur l'élément $(m_j)_{j \in J}$ du terme médian.
En écrivant $k_i = (a_{i, j})_{j \in J}$, nous obtenons le résultat voulu.
\end{proof}

\begin{lemma}
\label{algebra-lemma-kernel-difference-projections}
Soit $\varphi : R \to S$ un homomorphisme d'anneaux.
Soit $g \in S$. Les assertions suivantes sont équivalentes :
\begin{enumerate}
\item $g \otimes 1 = 1 \otimes g$ dans $S \otimes_R S$, et
\item il existe $n \geq 0$, des éléments $y_i, z_j \in S$
et des $x_{i, j} \in R$ pour $1 \leq i, j \leq n$ tels que
\begin{enumerate}
\item $g = \sum_{i, j \leq n} x_{i, j} y_i z_j$,
\item pour tout $j$, nous avons $\sum x_{i, j}y_i \in \varphi(R)$, et
\item pour tout $i$, nous avons $\sum x_{i, j}z_j \in \varphi(R)$.
\end{enumerate}
\end{enumerate}
\end{lemma}

\begin{proof}
Il est clair que (2) implique (1). Réciproquement, supposons que
$g \otimes 1 = 1 \otimes g$. Choisissons un système de générateurs $\{s_i\}_{i \in I}$
de $S$ comme $R$-module avec $0, 1 \in I$, $s_0 = 1$ et $s_1 = g$.
Appliquons le
lemme \ref{algebra-lemma-relations}
à la relation $g \otimes s_0 + (-1) \otimes s_1 = 0$.
Nous voyons qu'il existe des $a_{i, j} \in R$ tels que
$g = \sum_i a_{i, 0} s_i$, $-1 = \sum_i a_{i, 1} s_i$ et, pour
$j \not = 0, 1$, $0 = \sum_i a_{i, j} s_i$, et tels qu'en outre,
pour tout $i$, nous ayons $\sum_j a_{i, j}s_j = 0$.
Nous avons alors
$$
\sum\nolimits_{i, j \not = 0} a_{i, j} s_i s_j = -g + a_{0, 0}
$$
et, pour tout $j \not = 0$, nous avons
$\sum_{i \not = 0} a_{i, j}s_i \in R$. Cela prouve que $-g + a_{0, 0}$
peut s'écrire comme dans (2). Il s'ensuit que $g$ peut s'écrire
comme dans (2). Les détails sont omis.
Indication : montrer que l'ensemble des éléments de $S$ qui possèdent une
expression comme dans (2) forme une $R$-sous-algèbre de $S$.
\end{proof}

\begin{remark}
\label{algebra-remark-matrices-associated-to-elements-epicenter}
Soit $R \to S$ un homomorphisme d'anneaux. L'ensemble des éléments
$g \in S$ tels que $g \otimes 1 = 1 \otimes g$ est parfois appelé l'
{\it épicentre} de $S$. C'est une $R$-algèbre. La construction du
lemme \ref{algebra-lemma-kernel-difference-projections}
donne, pour chaque $g$ dans l'épicentre, une factorisation matricielle
$$
(g) = Y X Z
$$
avec $X \in \text{Mat}(n \times n, R)$,
$Y \in \text{Mat}(1 \times n, S)$ et
$Z \in \text{Mat}(n \times 1, S)$. En effet, soient $x_{i, j}, y_i, z_j$
comme dans la partie (2) du lemme. Posons $X = (x_{i, j})$, soit $y$ le
vecteur ligne dont les coefficients sont les $y_i$, et soit $z$ le vecteur colonne
dont les coefficients sont les $z_j$. Avec ces notations, les conditions (b) et (c) du
lemme \ref{algebra-lemma-kernel-difference-projections}
signifient exactement que $Y X \in \text{Mat}(1 \times n, R)$,
$X Z \in \text{Mat}(n \times 1, R)$.
Il s'avère très commode de considérer le triplet de
matrices $(X, YX, XZ)$. Étant donnés $n \in \mathbf{N}$ et un triplet
$(P, U, V)$, nous disons que $(P, U, V)$ est un {\it $n$-triplet associé à $g$}
s'il existe une factorisation matricielle comme ci-dessus telle que
$P = X$, $U = YX$ et $V = XZ$.
\end{remark}

\begin{lemma}
\label{algebra-lemma-epimorphism-cardinality}
Soit $R \to S$ un épimorphisme d'anneaux.
Alors le cardinal de $S$ est au plus celui de $R$.
En formule : $|S| \leq |R|$.
\end{lemma}

\begin{proof}
La condition que $R \to S$ soit un épimorphisme signifie que tout $g \in S$
vérifie $g \otimes 1 = 1 \otimes g$, voir le
lemme \ref{algebra-lemma-epimorphism}.
Nous allons utiliser les notations introduites dans la
remarque \ref{algebra-remark-matrices-associated-to-elements-epicenter}.
Supposons que $g, g' \in S$ et que $(P, U, V)$ soit un $n$-triplet
associé à la fois à $g$ et à $g'$. Nous affirmons alors que
$g = g'$. En effet, écrivons $(P, U, V) = (X, YX, XZ)$ pour une factorisation
matricielle $(g) = YXZ$ de $g$, et écrivons $(P, U, V) = (X', Y'X', X'Z')$
pour une factorisation matricielle $(g') = Y'X'Z'$ de $g'$.
Nous voyons alors que
$$
(g) = YXZ = UZ = Y'X'Z = Y'PZ = Y'XZ = Y'V = Y'X'Z' = (g')
$$
et donc que $g = g'$. Cela implique que le cardinal de $S$ est majoré
par le nombre de triplets possibles, dont le cardinal est au plus
$\sup_{n \in \mathbf{N}} |R|^n$. Si $R$ est infini, celui-ci est
au plus $|R|$, voir \cite[Ch. I, 10.13]{Kunen}.

\medskip\noindent
Si $R$ est un anneau fini, l'argument ci-dessus prouve seulement que $S$
est au plus dénombrable. En fait, dans ce cas, $R$ est artinien et
l'homomorphisme $R \to S$ est surjectif. Nous omettons la preuve de ce cas.
\end{proof}

\begin{lemma}
\label{algebra-lemma-epimorphism-modules}
Pour un homomorphisme d'anneaux $R \to S$, les assertions suivantes sont équivalentes :
\begin{enumerate}
\item $R \to S$ est un épimorphisme d'anneaux,
\item pour tous $S$-modules $N_1, N_2$, nous avons
$\Hom_S(N_1, N_2) = \Hom_R(N_1, N_2)$, et
\item le foncteur de restriction des scalaires $\text{Mod}_S \to \text{Mod}_R$
est pleinement fidèle.
\end{enumerate}
\end{lemma}

\begin{proof}
Remarquons que (2) et (3) sont équivalentes par définition.

\medskip\noindent
Supposons (1), soient $N_1, N_2$ des $S$-modules et soit $\varphi : N_1 \to N_2$
un homomorphisme $R$-linéaire. Pour tout $x \in N_1$, considérons l'homomorphisme
$S \otimes_R S \to N_2$ défini par la règle
$g \otimes g' \mapsto g\varphi(g'x)$. Puisque les deux homomorphismes $S \to S \otimes_R S$
sont des isomorphismes (lemme \ref{algebra-lemma-epimorphism}), nous concluons que
$g \varphi(g'x) = gg'\varphi(x) = \varphi(gg' x)$. Ainsi, $\varphi$
est $S$-linéaire.

\medskip\noindent
Supposons (2). Soit $N_1 = S \otimes_R S$ considéré comme un $S$-module
par l'action à gauche de $S$, et soit $N_2 = S \otimes_R S$
muni de l'action à droite de $S$. Puisque $N_1 = N_2$ comme $R$-modules,
(2) montre que les deux structures de $S$-module sur $S \otimes_R S$
coïncident. Cela implique (1) d'après le lemme \ref{algebra-lemma-epimorphism}.
\end{proof}



\section{Idéaux purs}
\label{algebra-section-pure-ideals}

\noindent
Les résultats de cette section sont abordés dans de nombreux articles, voir par exemple
\cite{Lazard}, \cite{Bkouche} et \cite{DeMarco}.

\begin{definition}
\label{algebra-definition-pure-ideal}
Soit $R$ un anneau. Nous disons qu'un idéal $I \subset R$ est {\it pur}
si l'anneau quotient $R/I$ est plat sur $R$.
\end{definition}

\begin{lemma}
\label{algebra-lemma-pure}
Soit $R$ un anneau.
Soit $I \subset R$ un idéal.
Les assertions suivantes sont équivalentes :
\begin{enumerate}
\item $I$ est pur,
\item pour tout idéal $J \subset R$, nous avons $J \cap I = IJ$,
\item pour tout idéal de type fini $J \subset R$, nous avons
$J \cap I = JI$,
\item pour tout $x \in R$, nous avons $(x) \cap I = xI$,
\item pour tout $x \in I$, nous avons $x = yx$ pour un certain $y \in I$,
\item pour tous $x_1, \ldots, x_n \in I$, il existe un
$y \in I$ tel que $x_i = yx_i$ pour tout $i = 1, \ldots, n$,
\item pour tout idéal premier $\mathfrak p$ de $R$, nous avons
$IR_{\mathfrak p} = 0$ ou $IR_{\mathfrak p} = R_{\mathfrak p}$,
\item $\text{Supp}(I) = \Spec(R) \setminus V(I)$,
\item $I$ est le noyau de l'homomorphisme $R \to (1 + I)^{-1}R$,
\item $R/I \cong S^{-1}R$ comme $R$-algèbres pour une certaine partie
multiplicative $S$ de $R$, et
\item $R/I \cong (1 + I)^{-1}R$ comme $R$-algèbres.
\end{enumerate}
\end{lemma}

\begin{proof}
Pour tout idéal $J$ de $R$, nous avons la suite exacte courte
$0 \to J \to R \to R/J \to 0$. En tensorisant avec $R/I$, nous obtenons
une suite exacte $J \otimes_R R/I \to R/I \to R/I + J \to 0$
et $J \otimes_R R/I = J/JI$. Ainsi, l'
équivalence de (1), (2) et (3) découle du
lemme \ref{algebra-lemma-flat}. En outre, ces assertions impliquent (4).

\medskip\noindent
L'implication (4) $\Rightarrow$ (5) est triviale.
Supposons (5) et soient $x_1, \ldots, x_n \in I$.
Choisissons $y_i \in I$ tel que $x_i = y_ix_i$.
Soit $y \in I$ l'élément tel que
$1 - y = \prod_{i = 1, \ldots, n} (1 - y_i)$.
Alors $x_i = yx_i$ pour tout $i = 1, \ldots, n$.
Ainsi, (6) est vérifiée, et il s'ensuit que (5) $\Leftrightarrow$ (6).

\medskip\noindent
Supposons (5). Soit $x \in I$. Alors $x = yx$ pour un certain $y \in I$.
Ainsi, $x(1 - y) = 0$, ce qui montre que $x$ s'envoie sur zéro dans
$(1 + I)^{-1}R$. Bien entendu, le noyau de l'homomorphisme
$R \to (1 + I)^{-1}R$ est toujours contenu dans $I$. Nous
voyons donc que (5) implique (9). Supposons (9). Alors, pour tout $x \in I$,
nous voyons que $x(1 - y) = 0$ pour un certain $y \in I$.
Autrement dit, $x = yx$. Nous concluons que (5) équivaut à (9).

\medskip\noindent
Supposons (5). Soit $\mathfrak p$ un idéal premier de $R$.
Si $\mathfrak p \not \in V(I)$, alors $IR_{\mathfrak p} = R_{\mathfrak p}$.
Si $\mathfrak p \in V(I)$, autrement dit, si $I \subset \mathfrak p$,
alors, pour $x \in I$, nous avons $x(1 - y) = 0$ pour un certain $y \in I$, ce qui implique que
$x$ s'envoie sur zéro dans $R_{\mathfrak p}$, c'est-à-dire que $IR_{\mathfrak p} = 0$.
Nous voyons donc que (7) est vérifiée.

\medskip\noindent
Supposons (7). Alors $(R/I)_{\mathfrak p}$ vaut soit $0$, soit $R_{\mathfrak p}$
pour tout idéal premier $\mathfrak p$ de $R$. Ainsi, d'après le
lemme \ref{algebra-lemma-flat-localization},
nous voyons que (1) est vérifiée. À ce stade, nous voyons que toutes les assertions
(1) -- (7) et (9) sont équivalentes.

\medskip\noindent
Comme $IR_{\mathfrak p} = I_{\mathfrak p}$, nous voyons que (7) implique (8).
Enfin, si (8) est vérifiée, cela signifie exactement que $I_{\mathfrak p}$
est le module nul si et seulement si $\mathfrak p \in V(I)$, ce qui
exprime clairement que (7) est vérifiée. Ainsi, (1) -- (9) sont équivalentes.

\medskip\noindent
Supposons que (1) -- (9) soient vérifiées. Alors $R/I \subset (1 + I)^{-1}R$ d'après (9), et
l'homomorphisme $R/I \to (1 + I)^{-1}R$ est également surjectif d'après la description
des localisations aux idéaux premiers fournie par (7). Ainsi, (11) est vérifiée.

\medskip\noindent
L'implication (11) $\Rightarrow$ (10) est triviale.
Et (10) implique que (1) est vérifiée, car une localisation de
$R$ est plate sur $R$, voir le
lemme \ref{algebra-lemma-flat-localization}.
\end{proof}

\begin{lemma}
\label{algebra-lemma-pure-ideal-determined-by-zero-set}
\begin{slogan}
Les idéaux purs sont déterminés par leur lieu d'annulation.
\end{slogan}
Soit $R$ un anneau.
Si $I, J \subset R$ sont des idéaux purs, alors $V(I) = V(J)$
implique $I = J$.
\end{lemma}

\begin{proof}
Par exemple, d'après la propriété (7) du
lemme \ref{algebra-lemma-pure},
nous voyons que
$I = \Ker(R \to \prod_{\mathfrak p \in V(I)} R_{\mathfrak p})$
peut être retrouvé à partir du fermé qui lui est associé.
\end{proof}

\begin{lemma}
\label{algebra-lemma-pure-open-closed-specializations}
Soit $R$ un anneau. La règle
$I \mapsto V(I)$
détermine une bijection
$$
\{I \subset R \text{ pur}\}
\leftrightarrow
\{Z \subset \Spec(R)\text{ fermé et stable par générisation}\}
$$
\end{lemma}

\begin{proof}
Soit $I$ un idéal pur. Puisque $R \to R/I$ est plat, le théorème de descente
montre que les générisations se relèvent le long de l'homomorphisme $\Spec(R/I) \to \Spec(R)$.
Ainsi, $V(I)$ est stable par générisation. Cela montre que l'application
est bien définie. D'après le
lemme \ref{algebra-lemma-pure-ideal-determined-by-zero-set},
l'application est injective. Supposons que
$Z \subset \Spec(R)$ soit fermé et stable par générisation.
Soit $J \subset R$ l'idéal radical tel que $Z = V(J)$.
Soit $I = \{x \in R : x \in xJ\}$. Remarquons que $I$ est un idéal :
si $x, y \in I$, alors il existe $f, g \in J$ tels que
$x = xf$ et $y = yg$. Alors
$$
x + y = (x + y)(f + g - fg)
$$
La vérification est laissée au lecteur.
Nous affirmons que $I$ est pur et que $V(I) = V(J)$.
Si cette affirmation est vraie, alors l'application du lemme est surjective et
le lemme est démontré.

\medskip\noindent
Remarquons que $I \subset J$, de sorte que $V(J) \subset V(I)$.
Soit $I \subset \mathfrak p$, ce dernier idéal étant premier. Considérons la partie multiplicative
$S = (R \setminus \mathfrak p)(1 + J)$. La définition de
$I$ et l'inclusion $I \subset \mathfrak p$ montrent que $0 \not \in S$.
Nous pouvons donc trouver un idéal premier $\mathfrak q$ de $R$ disjoint
de $S$, voir les
lemmes \ref{algebra-lemma-localization-zero} et
\ref{algebra-lemma-spec-localization}.
Ainsi, $\mathfrak q \subset \mathfrak p$ et
$\mathfrak q \cap (1 + J) = \emptyset$.
Cela implique que $\mathfrak q + J$ est un idéal propre de $R$.
Soit $\mathfrak m$ un idéal maximal contenant $\mathfrak q + J$.
Nous obtenons alors
$\mathfrak m \in V(J)$, et donc $\mathfrak q \in V(J) = Z$
puisque $Z$ est supposé stable par générisation.
Cela implique à son tour que $\mathfrak p \in V(J)$, car
$\mathfrak q \subset \mathfrak p$. Ainsi, nous voyons que $V(I) = V(J)$.

\medskip\noindent
Enfin, puisque $V(I) = V(J)$ (et que $J$ est radical), nous voyons que $J = \sqrt{I}$.
Choisissons $x \in I$ ; par définition, $x = xy$ pour un certain $y \in J$.
Alors $x = xy = xy^2 = \ldots = xy^n$. Puisque $y^n \in I$ pour un certain $n > 0$,
nous concluons que la propriété (5) du
lemme \ref{algebra-lemma-pure}
est vérifiée et nous voyons que $I$ est bien pur.
\end{proof}

\begin{lemma}
\label{algebra-lemma-finitely-generated-pure-ideal}
Soit $R$ un anneau. Soit $I \subset R$ un idéal.
Les assertions suivantes sont équivalentes :
\begin{enumerate}
\item $I$ est pur et de type fini,
\item $I$ est engendré par un idempotent,
\item $I$ est pur et $V(I)$ est ouvert, et
\item $R/I$ est un $R$-module projectif.
\end{enumerate}
\end{lemma}

\begin{proof}
Si (1) est vérifiée, alors $I = I \cap I = I^2$ d'après le
lemme \ref{algebra-lemma-pure}.
Ainsi, $I$ est engendré par un idempotent d'après le
lemme \ref{algebra-lemma-ideal-is-squared-union-connected}.
Ainsi, (1) $\Rightarrow$ (2).
Si (2) est vérifiée, alors $I = (e)$ et $R = (1 - e) \oplus (e)$ comme
$R$-module ; ainsi, $R/I$ est plat, $I$ est pur et $V(I) = D(1 - e)$
est ouvert. Donc (2) $\Rightarrow$ (1) $+$ (3).
Enfin, supposons (3). Alors $V(I)$ est ouvert et fermé, donc
$V(I) = D(1 - e)$ pour un certain idempotent $e$ de $R$, voir le
lemme \ref{algebra-lemma-disjoint-decomposition}.
L'idéal $J = (e)$ est un idéal pur tel que $V(J) = V(I)$ ; ainsi,
$I = J$ d'après le
lemme \ref{algebra-lemma-pure-ideal-determined-by-zero-set}.
Nous voyons ainsi que (3) $\Rightarrow$ (2). D'après le
lemme \ref{algebra-lemma-finite-projective},
nous voyons que (4) équivaut à l'assertion que $I$ est pur
et que $R/I$ est de présentation finie. En outre, $R/I$ est de présentation finie
si et seulement si $I$ est de type fini, voir le lemme \ref{algebra-lemma-extension}.
Ainsi, (4) équivaut à (1).
\end{proof}

\noindent
Nous pouvons utiliser ce qui précède pour caractériser les anneaux sur lesquels tout module
plat fini est de présentation finie.

\begin{lemma}
\label{algebra-lemma-finite-flat-module-finitely-presented}
Soit $R$ un anneau. Les assertions suivantes sont équivalentes :
\begin{enumerate}
\item tout $Z \subset \Spec(R)$ qui est fermé et stable par
générisation est également ouvert, et
\item tout $R$-module plat fini est localement libre de type fini.
\end{enumerate}
\end{lemma}

\begin{proof}
Si tout $R$-module plat fini est localement libre de type fini, alors le support
de $R/I$, où $I$ est un idéal pur, est ouvert. Ainsi, l'implication
(2) $\Rightarrow$ (1) découle du
lemme \ref{algebra-lemma-pure-ideal-determined-by-zero-set}.

\medskip\noindent
Pour la réciproque, supposons que $R$ vérifie (1).
Soit $M$ un $R$-module plat fini.
Le support $Z = \text{Supp}(M)$ de $M$ est fermé, voir le
lemme \ref{algebra-lemma-support-closed}.
D'autre part, si $\mathfrak p \subset \mathfrak p'$, alors, d'après le
lemme \ref{algebra-lemma-finite-flat-local},
le module $M_{\mathfrak p'}$ est libre, et
$M_{\mathfrak p} = M_{\mathfrak p'} \otimes_{R_{\mathfrak p'}} R_{\mathfrak p}$
Ainsi,
$\mathfrak p' \in \text{Supp}(M) \Rightarrow \mathfrak p \in \text{Supp}(M)$,
autrement dit, le support est stable par générisation.
Comme $R$ vérifie (1), nous voyons que le support de $M$ est ouvert et fermé.
Supposons que $M$ soit engendré par $r$ éléments $m_1, \ldots, m_r$.
Les modules $\wedge^i(M)$, $i = 1, \ldots, r$, sont eux aussi des $R$-modules
plats finis, car $\wedge^i(M)_{\mathfrak p} = \wedge^i(M_{\mathfrak p})$
est libre sur $R_{\mathfrak p}$. Remarquons que
$\text{Supp}(\wedge^{i + 1}(M)) \subset \text{Supp}(\wedge^i(M))$.
Nous voyons donc qu'il existe une décomposition
$$
\Spec(R) = U_0 \amalg U_1 \amalg \ldots \amalg U_r
$$
en parties ouvertes et fermées telle que le support de
$\wedge^i(M)$ soit $U_r \cup \ldots \cup U_i$ pour tout $i = 0, \ldots, r$.
Soit $\mathfrak p$ un idéal premier de $R$, et supposons que $\mathfrak p \in U_i$.
Remarquons que
$\wedge^i(M) \otimes_R \kappa(\mathfrak p) =
\wedge^i(M \otimes_R \kappa(\mathfrak p))$.
Ainsi, après avoir éventuellement renuméroté $m_1, \ldots, m_r$, nous pouvons supposer que
$m_1, \ldots, m_i$ engendrent $M \otimes_R \kappa(\mathfrak p)$. D'après le
lemme de Nakayama \ref{algebra-lemma-NAK},
nous obtenons une surjection
$$
R_f^{\oplus i} \longrightarrow M_f, \quad
(a_1, \ldots, a_i) \longmapsto \sum a_im_i
$$
pour un certain $f \in R$, $f \not \in \mathfrak p$. Nous pouvons aussi supposer que
$D(f) \subset U_i$. Cela signifie que $\wedge^i(M_f) = \wedge^i(M)_f$
est un $R_f$-module plat dont le support est tout $\Spec(R_f)$.
D'après ce qui précède, il est engendré par un seul élément, à savoir
$m_1 \wedge \ldots \wedge m_i$. Ainsi, $\wedge^i(M)_f \cong R_f/J$
pour un certain idéal pur $J \subset R_f$ tel que $V(J) = \Spec(R_f)$.
Cela signifie clairement que $J = (0)$, voir le
lemme \ref{algebra-lemma-pure-ideal-determined-by-zero-set}.
Ainsi, $m_1 \wedge \ldots \wedge m_i$ est une base de
$\wedge^i(M_f)$, et il s'ensuit que l'homomorphisme affiché est injectif
aussi bien que surjectif. Cela prouve que $M$ est localement libre de type fini,
comme voulu.
\end{proof}







\section{Anneaux de dimension globale finie}
\label{algebra-section-ring-finite-gl-dim}

\noindent
Le lemme suivant est souvent utilisé pour comparer différentes
résolutions projectives d'un module donné.

\begin{lemma}[Lemme de Schanuel]
\label{algebra-lemma-Schanuel}
Soit $R$ un anneau. Soit $M$ un $R$-module.
Supposons que
$$
0 \to K \xrightarrow{c_1} P_1 \xrightarrow{p_1} M \to 0
\quad\text{et}\quad
0 \to L \xrightarrow{c_2} P_2 \xrightarrow{p_2} M \to 0
$$
soient deux suites exactes courtes, avec $P_i$ projectif.
Alors $K \oplus P_2 \cong L \oplus P_1$. Plus précisément,
il existe un diagramme commutatif
$$
\xymatrix{
0 \ar[r] &
K \oplus P_2 \ar[r]_{(c_1, \text{id})} \ar[d] &
P_1 \oplus P_2 \ar[r]_{(p_1, 0)} \ar[d] &
M \ar[r] \ar@{=}[d] &
0 \\
0 \ar[r] &
P_1 \oplus L \ar[r]^{(\text{id}, c_2)} &
P_1 \oplus P_2 \ar[r]^{(0, p_2)} &
M \ar[r] &
0
}
$$
dont les flèches verticales sont des isomorphismes.
\end{lemma}

\begin{proof}
Considérons le module $N$ défini par la suite exacte courte
$0 \to N \to P_1 \oplus P_2 \to M \to 0$,
où le dernier homomorphisme est la somme des deux homomorphismes
$P_i \to M$. Il est facile de voir que la projection
$N \to P_1$ est surjective de noyau $L$, et que
$N \to P_2$ est surjective de noyau $K$.
Puisque les $P_i$ sont projectifs, nous avons $N \cong K \oplus P_2
\cong L \oplus P_1$. Cela prouve la première assertion.

\medskip\noindent
Pour prouver la seconde assertion (et prouver de nouveau la première), choisissons
$a : P_1 \to P_2$ et $b : P_2 \to P_1$ tels que
$p_1 = p_2 \circ a$ et $p_2 = p_1 \circ b$. C'est possible
parce que $P_1$ et $P_2$ sont projectifs. Nous obtenons alors un diagramme commutatif
$$
\xymatrix{
0 \ar[r] &
K \oplus P_2 \ar[r]_{(c_1, \text{id})} &
P_1 \oplus P_2 \ar[r]_{(p_1, 0)} &
M \ar[r] &
0 \\
0 \ar[r] &
N \ar[r] \ar[d] \ar[u] &
P_1 \oplus P_2 \ar[r]_{(p_1, p_2)}
\ar[d]_S \ar[u]^T &
M \ar[r] \ar@{=}[d] \ar@{=}[u] &
0 \\
0 \ar[r] &
P_1 \oplus L \ar[r]^{(\text{id}, c_2)} &
P_1 \oplus P_2 \ar[r]^{(0, p_2)} &
M \ar[r] &
0
}
$$
où $T$ et $S$ sont donnés par les matrices
$$
S = \left(
\begin{matrix}
\text{id} & 0 \\
a & \text{id}
\end{matrix}
\right)
\quad\text{et}\quad
T = \left(
\begin{matrix}
\text{id} & b \\
0 & \text{id}
\end{matrix}
\right)
$$
Alors $S$, $T$ et les homomorphismes $N \to P_1 \oplus L$
et $N \to K \oplus P_2$ sont des isomorphismes, comme voulu.
\end{proof}

\begin{definition}
\label{algebra-definition-finite-proj-dim}
Soit $R$ un anneau. Soit $M$ un $R$-module. Nous disons que $M$ est de
{\it dimension projective finie} s'il possède une résolution de longueur finie
par des $R$-modules projectifs. La longueur minimale d'une telle
résolution est appelée la {\it dimension projective}
de $M$.
\end{definition}

\noindent
Il est clair que la dimension projective de $M$ est $0$ si et
seulement si $M$ est un module projectif.
Le lemme suivant explique dans quelle mesure la dimension
projective est indépendante du choix d'une
résolution projective.

\begin{lemma}
\label{algebra-lemma-independent-resolution}
Soit $R$ un anneau. Supposons que $M$ soit un $R$-module de dimension
projective $d$. Supposons que $F_e \to F_{e-1} \to \ldots \to F_0 \to M \to 0$
soit exacte, avec les $F_i$ projectifs et $e \geq d - 1$.
Alors le noyau de $F_e \to F_{e-1}$ est projectif
(ou bien le noyau de $F_0 \to M$ est projectif dans le cas où
$e = 0$).
\end{lemma}

\begin{proof}
Nous le prouvons par récurrence sur $d$. Si $d = 0$, alors
$M$ est projectif. Dans ce cas, il existe une décomposition
$F_0 = \Ker(F_0 \to M) \oplus M$, et donc
$\Ker(F_0 \to M)$ est projectif. Cela achève
la preuve si $e = 0$ ; si $e > 0$, en remplaçant
$M$ par $\Ker(F_0 \to M)$, nous diminuons $e$.

\medskip\noindent
Supposons maintenant $d > 0$.
Soit $0 \to P_d \to P_{d-1} \to \ldots \to P_0 \to M \to 0$
une résolution de longueur minimale, avec les $P_i$ projectifs.
D'après le
lemme de Schanuel \ref{algebra-lemma-Schanuel},
nous avons
$P_0 \oplus \Ker(F_0 \to M) \cong F_0 \oplus \Ker(P_0 \to M)$.
Cela prouve le cas $d = 1$, $e = 0$, car le membre de
droite est alors $F_0 \oplus P_1$, qui est projectif. Nous pouvons donc désormais
supposer $e > 0$. Le module $F_0 \oplus \Ker(P_0 \to M)$ possède la
résolution projective finie
$$
0 \to P_d \to P_{d-1} \to \ldots \to
P_2 \to P_1 \oplus F_0 \to \Ker(P_0 \to M) \oplus F_0 \to 0
$$
de longueur $d - 1$. En appliquant l'hypothèse de récurrence à la suite exacte
$$
F_e \to F_{e-1} \to \ldots \to F_2 \to P_0 \oplus F_1 \to
P_0 \oplus \Ker(F_0 \to M) \to 0
$$
de longueur $e - 1$, nous concluons que $\Ker(F_e \to F_{e - 1})$
est projectif (si $e \geq 2$),
ou que $\Ker(F_1 \oplus P_0 \to F_0 \oplus P_0)$ est projectif.
Cela implique le lemme.
\end{proof}

\begin{lemma}
\label{algebra-lemma-what-kind-of-resolutions}
Soit $R$ un anneau. Soit $M$ un $R$-module. Soit $d \geq 0$.
Les assertions suivantes sont équivalentes :
\begin{enumerate}
\item $M$ est de dimension projective $\leq d$,
\item il existe une résolution
$0 \to P_d \to P_{d - 1} \to \ldots \to P_0 \to M \to 0$
avec les $P_i$ projectifs,
\item pour une certaine résolution
$\ldots \to P_2 \to P_1 \to P_0 \to M \to 0$ avec les
$P_i$ projectifs, $\Ker(P_{d - 1} \to P_{d - 2})$
est projectif si $d \geq 2$, ou $\Ker(P_0 \to M)$ est projectif si
$d = 1$, ou $M$ est projectif si $d = 0$,
\item pour toute résolution
$\ldots \to P_2 \to P_1 \to P_0 \to M \to 0$ avec les
$P_i$ projectifs, $\Ker(P_{d - 1} \to P_{d - 2})$
est projectif si $d \geq 2$, ou $\Ker(P_0 \to M)$ est projectif si
$d = 1$, ou $M$ est projectif si $d = 0$.
\end{enumerate}
\end{lemma}

\begin{proof}
L'équivalence de (1) et (2) est la définition de la dimension
projective, voir la définition \ref{algebra-definition-finite-proj-dim}.
Nous avons (2) $\Rightarrow$ (4) d'après le lemme \ref{algebra-lemma-independent-resolution}.
Les implications (4) $\Rightarrow$ (3) et (3) $\Rightarrow$ (2) sont
immédiates.
\end{proof}

\begin{lemma}
\label{algebra-lemma-what-kind-of-resolutions-local}
Soit $R$ un anneau local. Soit $M$ un $R$-module. Soit $d \geq 0$.
Les conditions équivalentes (1) -- (4) du
lemme \ref{algebra-lemma-what-kind-of-resolutions}
sont également équivalentes à
\begin{enumerate}
\item[(5)] il existe une résolution
$0 \to P_d \to P_{d - 1} \to \ldots \to P_0 \to M \to 0$
avec les $P_i$ libres.
\end{enumerate}
\end{lemma}

\begin{proof}
Cela découle du lemme \ref{algebra-lemma-what-kind-of-resolutions} et du
théorème \ref{algebra-theorem-projective-free-over-local-ring}.
\end{proof}

\begin{lemma}
\label{algebra-lemma-what-kind-of-resolutions-Noetherian}
Soit $R$ un anneau noethérien. Soit $M$ un $R$-module fini.
Soit $d \geq 0$. Les conditions équivalentes (1) -- (4) du
lemme \ref{algebra-lemma-what-kind-of-resolutions}
sont également équivalentes à
\begin{enumerate}
\item[(6)] il existe une résolution
$0 \to P_d \to P_{d - 1} \to \ldots \to P_0 \to M \to 0$
avec les $P_i$ projectifs finis.
\end{enumerate}
\end{lemma}

\begin{proof}
Choisissons une résolution $\ldots \to F_2 \to F_1 \to F_0 \to M \to 0$
avec les $F_i$ libres finis (lemme \ref{algebra-lemma-resolution-by-finite-free}).
D'après le lemme \ref{algebra-lemma-what-kind-of-resolutions}, nous voyons que
$P_d = \Ker(F_{d - 1} \to F_{d - 2})$ est projectif, du moins si $d \geq 2$.
Alors $P_d$ est un $R$-module fini puisque $R$ est noethérien et que
$P_d \subset F_{d - 1}$, qui est libre fini.
Ainsi, $0 \to P_d \to F_{d - 1} \to \ldots \to F_1 \to F_0 \to M \to 0$
est la résolution voulue.
\end{proof}

\begin{lemma}
\label{algebra-lemma-what-kind-of-resolutions-Noetherian-local}
Soit $R$ un anneau local noethérien. Soit $M$ un $R$-module fini.
Soit $d \geq 0$. Les conditions équivalentes (1) -- (4) du
lemme \ref{algebra-lemma-what-kind-of-resolutions},
la condition (5) du lemme \ref{algebra-lemma-what-kind-of-resolutions-local},
et la condition (6) du lemme \ref{algebra-lemma-what-kind-of-resolutions-Noetherian}
sont également équivalentes à
\begin{enumerate}
\item[(7)] il existe une résolution
$0 \to F_d \to F_{d - 1} \to \ldots \to F_0 \to M \to 0$
avec les $F_i$ libres finis.
\end{enumerate}
\end{lemma}

\begin{proof}
Cela découle des lemmes \ref{algebra-lemma-what-kind-of-resolutions},
\ref{algebra-lemma-what-kind-of-resolutions-local} et
\ref{algebra-lemma-what-kind-of-resolutions-Noetherian},
ainsi que du fait qu'un module projectif fini sur un anneau local
est libre fini, voir le lemme \ref{algebra-lemma-finite-projective}.
\end{proof}

\begin{lemma}
\label{algebra-lemma-projective-dimension-ext}
Soit $R$ un anneau. Soit $M$ un $R$-module. Soit $n \geq 0$.
Les assertions suivantes sont équivalentes :
\begin{enumerate}
\item $M$ est de dimension projective $\leq n$,
\item $\Ext^i_R(M, N) = 0$ pour tout $R$-module $N$ et tout
$i \geq n + 1$, et
\item $\Ext^{n + 1}_R(M, N) = 0$ pour tout $R$-module $N$.
\end{enumerate}
\end{lemma}

\begin{proof}
Supposons (1). Choisissons une résolution libre $F_\bullet \to M$ de $M$. Notons
$d_e : F_e \to F_{e - 1}$. D'après le
lemme \ref{algebra-lemma-independent-resolution},
nous voyons que $P_e = \Ker(d_e)$ est projectif pour $e \geq n - 1$.
Cela implique que $F_e \cong P_e \oplus P_{e - 1}$ pour $e \geq n$,
où $d_e$ envoie isomorphiquement le facteur $P_{e - 1}$ sur $P_{e - 1}$
dans $F_{e - 1}$. Ainsi, pour tout $R$-module $N$, le complexe
$\Hom_R(F_\bullet, N)$ est exact scindé en degrés $\geq n + 1$.
Donc (2) est vérifiée. L'implication (2) $\Rightarrow$ (3) est triviale.

\medskip\noindent
Supposons que (3) soit vérifiée. Si $n = 0$, alors $M$ est projectif d'après le
lemme \ref{algebra-lemma-characterize-projective},
et nous voyons que (1) est vérifiée. Si $n > 0$, choisissons un $R$-module libre $F$
et une surjection $F \to M$ de noyau $K$. D'après le
lemme \ref{algebra-lemma-reverse-long-exact-seq-ext}
et l'annulation de $\Ext_R^i(F, N)$ pour tout $i > 0$ donnée par la partie (1),
nous voyons que $\Ext_R^n(K, N) = 0$ pour tout $R$-module $N$.
Ainsi, par récurrence, nous voyons que $K$ est de dimension projective $\leq n - 1$.
Alors $M$ est de dimension projective $\leq n$, car toute résolution projective finie
de $K$ donne une résolution projective de longueur augmentée de un
de $M$ en ajoutant $F$ au début.
\end{proof}

\begin{lemma}
\label{algebra-lemma-exact-sequence-projective-dimension}
Soit $R$ un anneau. Soit $0 \to M' \to M \to M'' \to 0$ une suite
exacte courte de $R$-modules.
\begin{enumerate}
\item Si $M$ est de dimension projective $\leq n$ et $M''$
de dimension projective $\leq n + 1$, alors $M'$ est de dimension projective
$\leq n$.
\item Si $M'$ et $M''$ sont de dimension projective
$\leq n$, alors $M$ est de dimension projective $\leq n$.
\item Si $M'$ est de dimension projective $\leq n$ et
$M$ de dimension projective $\leq n + 1$, alors
$M''$ est de dimension projective $\leq n + 1$.
\end{enumerate}
\end{lemma}

\begin{proof}
Combiner la caractérisation de la dimension projective du
lemme \ref{algebra-lemma-projective-dimension-ext}
avec la suite exacte longue de groupes Ext du
lemme \ref{algebra-lemma-reverse-long-exact-seq-ext}.
\end{proof}

\begin{definition}
\label{algebra-definition-finite-gl-dim}
Soit $R$ un anneau. On dit que l'anneau
$R$ est de {\it dimension globale finie}
s'il existe un entier $n$ tel que
tout $R$-module possède une résolution par des
$R$-modules projectifs de longueur au plus $n$.
Le plus petit tel $n$ est alors appelé la {\it dimension globale}
de $R$.
\end{definition}

\noindent
L'argument de la preuve du lemme suivant se trouve
dans l'article \cite{Auslander} d'Auslander.

\begin{lemma}
\label{algebra-lemma-colimit-projective-dimension}
Soit $R$ un anneau. Supposons donné un module $M = \bigcup_{e \in E} M_e$
où les $M_e$ sont des sous-modules bien ordonnés par inclusion. Supposons que les quotients
$M_e/\bigcup\nolimits_{e' < e} M_{e'}$ soient de dimension projective $\leq n$.
Alors $M$ est de dimension projective $\leq n$.
\end{lemma}

\begin{proof}
Nous allons le prouver par récurrence sur $n$.

\medskip\noindent
Cas initial : $n = 0$. Alors $P_e = M_e/\bigcup_{e' < e} M_{e'}$ est projectif.
Nous pouvons donc choisir une section $P_e \to M_e$ de la projection $M_e \to P_e$.
Nous affirmons que l'homomorphisme induit $\psi : \bigoplus_{e \in E} P_e \to M$ est un
isomorphisme. En effet, si $x = \sum x_e \in \bigoplus P_e$ est non nul,
soit $e_{max}$ maximal tel que $x_{e_{max}}$ soit non nul ;
nous concluons que $y = \psi(x) = \psi(\sum x_e)$ est non nul parce que
$y \in M_{e_{max}}$ a pour image non nulle $x_{e_{max}}$ dans $P_{e_{max}}$.
D'autre part, soit $y \in M$. Alors $y \in M_e$ pour un certain $e$.
Nous montrons que $y \in \Im(\psi)$ par récurrence transfinie sur $e$.
Soit $x_e \in P_e$ l'image de $y$. Alors
$y - \psi(x_e) \in \bigcup_{e' < e} M_{e'}$.
Par hypothèse de récurrence, nous concluons que $y - \psi(x_e) \in \Im(\psi)$,
et donc que $y \in \Im(\psi)$. L'affirmation est ainsi vraie et
$\psi$ est un isomorphisme. Nous concluons que $M$ est projectif comme
somme directe de modules projectifs, voir le
lemme \ref{algebra-lemma-direct-sum-projective}.

\medskip\noindent
Si $n > 0$, alors, pour $e \in E$, nous notons $F_e$ le $R$-module libre
sur l'ensemble des éléments de $M_e$. Nous avons alors un système de
suites exactes courtes
$$
0 \to K_e \to F_e \to M_e \to 0
$$
indexé par l'ensemble bien ordonné $E$. Remarquons que les homomorphismes de transition
$F_{e'} \to F_e$ et $K_{e'} \to K_e$ sont eux aussi injectifs.
Posons $F = \bigcup F_e$ et $K = \bigcup K_e$. Alors
$$
0 \to
K_e/\bigcup\nolimits_{e' < e} K_{e'} \to
F_e/\bigcup\nolimits_{e' < e} F_{e'} \to
M_e/\bigcup\nolimits_{e' < e} M_{e'} \to 0
$$
est aussi une suite exacte courte de $R$-modules, et
$F_e/\bigcup_{e' < e} F_{e'}$ est le $R$-module libre sur l'
ensemble des éléments de $M_e$ qui ne sont pas contenus dans $\bigcup_{e' < e} M_{e'}$.
Ainsi, d'après le
lemme \ref{algebra-lemma-exact-sequence-projective-dimension},
nous voyons que la dimension projective de $K_e/\bigcup_{e' < e} K_{e'}$
est au plus $n - 1$. Par récurrence, nous concluons que $K$ est de dimension
projective au plus $n - 1$. Donc $M$ est de dimension projective au plus
$n$, ce qui conclut.
\end{proof}

\begin{lemma}
\label{algebra-lemma-finite-gl-dim}
Soit $R$ un anneau. Les assertions suivantes sont équivalentes :
\begin{enumerate}
\item $R$ est de dimension globale finie $\leq n$,
\item tout $R$-module fini est de dimension projective $\leq n$, et
\item tout $R$-module cyclique $R/I$ est de dimension projective $\leq n$.
\end{enumerate}
\end{lemma}

\begin{proof}
Il est clair que (1) $\Rightarrow$ (2) et (2) $\Rightarrow$ (3).
Supposons (3). Choisissons un ensemble $E \subset M$ de générateurs de $M$.
Choisissons un bon ordre sur $E$. Pour $e \in E$, notons
$M_e$ le sous-module de $M$ engendré par les éléments $e' \in E$
tels que $e' \leq e$. Alors $M = \bigcup_{e \in E} M_e$.
Remarquons que, pour chaque $e \in E$, le quotient
$$
M_e/\bigcup\nolimits_{e' < e} M_{e'}
$$
est soit nul, soit engendré par un élément ; il est donc de dimension
projective $\leq n$ d'après (3). D'après le lemme \ref{algebra-lemma-colimit-projective-dimension},
cela signifie que $M$ est de dimension projective $\leq n$.
\end{proof}

\begin{lemma}
\label{algebra-lemma-localize-finite-gl-dim}
Soit $R$ un anneau. Soit $M$ un $R$-module.
Soit $S \subset R$ une partie multiplicative.
\begin{enumerate}
\item Si $M$ est de dimension projective $\leq n$, alors $S^{-1}M$ est de
dimension projective $\leq n$ sur $S^{-1}R$.
\item Si $R$ est de dimension globale finie $\leq n$, alors
$S^{-1}R$ est de dimension globale finie $\leq n$.
\end{enumerate}
\end{lemma}

\begin{proof}
Soit $0 \to P_n \to P_{n - 1} \to \ldots \to P_0 \to M \to 0$
une résolution projective. Comme la localisation est exacte, voir la
proposition \ref{algebra-proposition-localization-exact},
et comme chaque $S^{-1}P_i$ est un $S^{-1}R$-module projectif, voir le
lemme \ref{algebra-lemma-ascend-properties-modules},
nous voyons que $0 \to S^{-1}P_n \to \ldots \to S^{-1}P_0 \to S^{-1}M \to 0$
est une résolution projective de $S^{-1}M$. Cela prouve (1).
Soit $M'$ un $S^{-1}R$-module.
Remarquons que $M' = S^{-1}M'$.
Nous voyons donc que (2) découle de (1).
\end{proof}










\section{Anneaux réguliers et dimension globale}
\label{algebra-section-regular-finite-gl-dim}

\noindent
Nous pouvons utiliser les résultats sur les anneaux de dimension globale finie
pour donner une autre caractérisation des anneaux locaux réguliers.

\begin{proposition}
\label{algebra-proposition-regular-finite-gl-dim}
Soit $R$ un anneau local régulier de dimension $d$.
Tout $R$-module fini $M$ de profondeur $e$ possède une résolution
libre finie
$$
0 \to F_{d-e} \to \ldots \to F_0 \to M \to 0.
$$
En particulier, un anneau local régulier est de dimension globale $\leq d$.
\end{proposition}

\begin{proof}
La première partie découle du lemme \ref{algebra-lemma-regular-mcm-free}
et du lemme \ref{algebra-lemma-mcm-resolution}. La dernière partie en découle
avec le lemme \ref{algebra-lemma-finite-gl-dim}.
\end{proof}

\begin{lemma}
\label{algebra-lemma-finite-gl-dim-primes}
Soit $R$ un anneau noethérien. Soit $n \geq 0$ un entier.
Alors $R$ est de dimension globale finie $\leq n$ si et
seulement si, pour tout idéal maximal $\mathfrak m$ de $R$,
l'anneau $R_{\mathfrak m}$ est de dimension globale $\leq n$.
\end{lemma}

\begin{proof}
Nous avons vu, dans le lemme \ref{algebra-lemma-localize-finite-gl-dim},
que, si $R$ est de dimension globale finie $n$,
alors toutes les localisations $R_{\mathfrak m}$
sont de dimension globale finie au plus $n$.
Réciproquement, supposons que tous les $R_{\mathfrak m}$
soient de dimension globale $\leq n$. Soit $M$ un
$R$-module fini. Soit
$0 \to K_n \to F_{n-1} \to \ldots \to F_0 \to M \to 0$
une résolution avec les $F_i$ libres finis.
Alors $K_n$ est un $R$-module fini.
D'après le
lemme \ref{algebra-lemma-independent-resolution}
et l'hypothèse, tous les modules $K_n \otimes_R R_{\mathfrak m}$
sont projectifs. Ainsi, d'après le
lemme \ref{algebra-lemma-finite-projective},
le module $K_n$ est projectif fini.
\end{proof}

\begin{lemma}
\label{algebra-lemma-length-resolution-residue-field}
Supposons que $R$ soit un anneau local noethérien
d'idéal maximal $\mathfrak m$ et de
corps résiduel $\kappa$. Alors
la dimension projective de $\kappa$ est
$\geq \dim_\kappa \mathfrak m / \mathfrak m^2$.
\end{lemma}

\begin{proof}
Soient $x_1 , \ldots, x_n$ des éléments de $\mathfrak m$
dont les images dans $\mathfrak m / \mathfrak m^2$ forment une base.
Considérons le {\it complexe de Koszul} sur $x_1, \ldots, x_n$.
C'est le complexe
$$
0 \to \wedge^n R^n \to \wedge^{n-1} R^n \to \wedge^{n-2} R^n \to
\ldots \to \wedge^i R^n \to \ldots \to R^n \to R
$$
dont les homomorphismes sont donnés par
$$
e_{j_1} \wedge \ldots \wedge e_{j_i}
\longmapsto
\sum_{a = 1}^i (-1)^{a + 1} x_{j_a} e_{j_1} \wedge \ldots
\wedge \hat e_{j_a} \wedge \ldots \wedge e_{j_i}
$$
Il est facile de voir qu'il s'agit d'un complexe $K_{\bullet}(R, x_{\bullet})$.
Remarquons que le conoyau du dernier homomorphisme de $K_{\bullet}(R, x_{\bullet})$
est $\kappa$ d'après la partie (8) du lemme \ref{algebra-lemma-NAK}.

\medskip\noindent
Si $\kappa$ est de dimension projective finie $d$, nous pouvons trouver
une résolution $F_{\bullet} \to \kappa$ par des $R$-modules libres finis
de longueur $d$
(lemme \ref{algebra-lemma-what-kind-of-resolutions-Noetherian-local}).
D'après le lemme \ref{algebra-lemma-add-trivial-complex},
nous pouvons supposer que tous les homomorphismes du complexe $F_{\bullet}$
ont la propriété que $\Im(F_i \to F_{i-1})
\subset \mathfrak m F_{i-1}$, car supprimer un facteur trivial
de la résolution peut tout au plus raccourcir celle-ci.
D'après le lemme \ref{algebra-lemma-compare-resolutions}, nous pouvons trouver un morphisme
de complexes $\alpha : K_{\bullet}(R, x_{\bullet}) \to F_{\bullet}$
induisant l'identité sur $\kappa$. Nous allons prouver par récurrence
que les homomorphismes $\alpha_i : \wedge^i R^n = K_i(R, x_{\bullet}) \to F_i$
ont la propriété que
$\alpha_i \otimes \kappa : \wedge^i \kappa^n \to F_i \otimes \kappa$
sont injectifs. Cela montre que $F_n \not = 0$ et donc que $d \geq n$,
comme voulu.

\medskip\noindent
Le résultat est clair pour $i = 0$, car la composée
$R \xrightarrow{\alpha_0} F_0 \to \kappa$ est non nulle.
Remarquons que $F_0$ doit être de rang $1$, puisque
sinon l'homomorphisme $F_1 \to F_0$, dont le conoyau est une seule
copie de $\kappa$, ne peut pas avoir son image contenue dans $\mathfrak m F_0$.

\medskip\noindent
Vérifions ensuite le cas $i = 1$, que nous jugeons instructif ;
le lecteur peut l'omettre, puisque l'étape de récurrence déduira le cas $i = 1$
du cas $i = 0$. Nous avons vu ci-dessus que
$F_0 = R$ et que l'homomorphisme $F_1 \to F_0 = R$ a pour image $\mathfrak m$.
Nous avons un diagramme commutatif
$$
\begin{matrix}
R^n & = & K_1(R, x_{\bullet}) & \to & K_0(R, x_{\bullet}) & = & R \\
& & \downarrow & & \downarrow & & \downarrow \\
& & F_1 & \to & F_0 & = & R
\end{matrix}
$$
où la flèche verticale la plus à droite est donnée par la multiplication
par une unité. Nous voyons donc que l'image de la composée
$R^n \to F_1 \to F_0 = R$ est elle aussi égale à $\mathfrak m$.
Ainsi, l'homomorphisme $R^n \otimes \kappa \to F_1 \otimes \kappa$
doit être injectif puisque $\dim_\kappa (\mathfrak m / \mathfrak m^2) = n$.

\medskip\noindent
Soit $i \geq 1$ et supposons l'injectivité de $\alpha_j \otimes \kappa$ prouvée
pour tout $j \leq i - 1$. Considérons le diagramme commutatif
$$
\begin{matrix}
\wedge^i R^n & = & K_i(R, x_{\bullet}) & \to & K_{i-1}(R, x_{\bullet})
& = & \wedge^{i-1} R^n \\
& & \downarrow & & \downarrow & & \\
& & F_i & \to & F_{i-1} & &
\end{matrix}
$$
Nous savons que $\wedge^{i-1} \kappa^n \to F_{i-1} \otimes \kappa$
est injectif. Cela prouve que
$\wedge^{i-1} \kappa^n \otimes_{\kappa} \mathfrak m/\mathfrak m^2
\to F_{i-1} \otimes \mathfrak m/\mathfrak m^2$ est injectif.
En outre, par notre choix du complexe, $F_i$ s'envoie dans
$\mathfrak mF_{i-1}$, et il en va de même pour le complexe de Koszul.
Nous obtenons donc un diagramme commutatif
$$
\begin{matrix}
\wedge^i \kappa^n & \to &
\wedge^{i-1} \kappa^n \otimes \mathfrak m/\mathfrak m^2 \\
\downarrow & & \downarrow \\
F_i \otimes \kappa & \to & F_{i-1} \otimes \mathfrak m/\mathfrak m^2
\end{matrix}
$$
À ce stade, il suffit de vérifier que l'homomorphisme
$\wedge^i \kappa^n \to
\wedge^{i-1} \kappa^n \otimes \mathfrak m/\mathfrak m^2$
est injectif, ce qui se vérifie directement.
\end{proof}

\begin{lemma}
\label{algebra-lemma-dim-gl-dim}
Soit $R$ un anneau local noethérien.
Supposons que le corps résiduel $\kappa$ ait une dimension
projective finie $n$ sur $R$.
Alors $\dim(R) \geq n$.
\end{lemma}

\begin{proof}
Soit $F_{\bullet}$ une résolution finie de $\kappa$ par des
$R$-modules libres finis (lemme \ref{algebra-lemma-what-kind-of-resolutions-Noetherian-local}).
D'après le lemme \ref{algebra-lemma-add-trivial-complex},
nous pouvons supposer que tous les homomorphismes du complexe $F_{\bullet}$
ont la propriété que $\Im(F_i \to F_{i-1})
\subset \mathfrak m F_{i-1}$, car supprimer un facteur trivial
de la résolution peut tout au plus raccourcir celle-ci.
Écrivons $F_n \not = 0$ et $F_i = 0$ pour $i > n$, de sorte que
la dimension projective de $\kappa$ soit $n$.
D'après la proposition \ref{algebra-proposition-what-exact}, nous voyons que
$\text{depth}_{I(\varphi_n)}(R) \geq n$, puisque $I(\varphi_n)$
ne peut pas être égal à $R$ par notre choix du complexe.
Ainsi, d'après le lemme \ref{algebra-lemma-bound-depth}, nous avons aussi $\dim(R) \geq n$.
\end{proof}

\begin{proposition}
\label{algebra-proposition-finite-gl-dim-regular}
Soit $(R, \mathfrak m, \kappa)$ un anneau local noethérien.
Les assertions suivantes sont équivalentes :
\begin{enumerate}
\item $\kappa$ est de dimension projective finie comme $R$-module,
\item $R$ est de dimension globale finie,
\item $R$ est un anneau local régulier.
\end{enumerate}
De plus, dans ce cas, la dimension globale de $R$ est égale à
$\dim(R) = \dim_\kappa(\mathfrak m/\mathfrak m^2)$.
\end{proposition}

\begin{proof}
Nous avons (3) $\Rightarrow$ (2) d'après la
proposition \ref{algebra-proposition-regular-finite-gl-dim}.
L'implication (2) $\Rightarrow$ (1) est triviale.
Supposons (1). D'après les lemmes \ref{algebra-lemma-length-resolution-residue-field}
et \ref{algebra-lemma-dim-gl-dim}, nous voyons que
$\dim(R) \geq \dim_\kappa(\mathfrak m /\mathfrak m^2)$.
Ainsi, $R$ est régulier, voir la définition
\ref{algebra-definition-regular-local} et la discussion qui la précède.
Supposons que les conditions équivalentes (1) -- (3) soient vérifiées.
D'après la proposition \ref{algebra-proposition-regular-finite-gl-dim},
la dimension globale de $R$ est au plus $\dim(R)$ et, d'après le
lemme \ref{algebra-lemma-length-resolution-residue-field},
elle est au moins $\dim_\kappa(\mathfrak m/\mathfrak m^2)$.
L'égalité annoncée est donc vérifiée.
\end{proof}

\begin{lemma}
\label{algebra-lemma-localization-of-regular-local-is-regular}
Un anneau local noethérien $R$ est régulier si et seulement si
il est de dimension globale finie. Dans ce cas,
$R_{\mathfrak p}$ est un anneau local régulier pour tout idéal premier $\mathfrak p$.
\end{lemma}

\begin{proof}
D'après les propositions \ref{algebra-proposition-finite-gl-dim-regular} et
\ref{algebra-proposition-regular-finite-gl-dim},
nous voyons qu'un anneau local noethérien est régulier si et seulement si
il est de dimension globale finie. En outre, toute localisation
$R_{\mathfrak p}$ est de dimension globale finie,
voir le lemme \ref{algebra-lemma-localize-finite-gl-dim},
et est donc un anneau local régulier.
\end{proof}

\noindent
D'après le lemme \ref{algebra-lemma-localization-of-regular-local-is-regular},
il est raisonnable de poser la définition suivante,
car elle n'entre pas en conflit avec la définition antérieure
d'un anneau local régulier.

\begin{definition}
\label{algebra-definition-regular}
Un anneau noethérien $R$ est dit {\it régulier}
si toutes ses localisations $R_{\mathfrak p}$ aux idéaux premiers sont des
anneaux locaux réguliers.
\end{definition}

\noindent
Il suffit d'exiger que les anneaux locaux aux idéaux maximaux soient réguliers.
Remarquons que cela ne revient pas à demander que $R$ soit de dimension
globale finie, même en supposant $R$ noethérien. Cela vient
de ce qu'il existe un exemple d'anneau noethérien régulier
qui n'est pas de dimension globale finie, précisément parce qu'
il n'est pas de dimension finie.

\begin{lemma}
\label{algebra-lemma-finite-gl-dim-finite-dim-regular}
Soit $R$ un anneau noethérien.
Les assertions suivantes sont équivalentes :
\begin{enumerate}
\item $R$ est de dimension globale finie $n$,
\item $R$ est un anneau régulier de dimension $n$,
\item il existe un entier $n$ tel que
toutes les localisations $R_{\mathfrak m}$ aux idéaux maximaux
soient régulières de dimension $\leq n$, avec égalité pour au moins
un $\mathfrak m$, et
\item il existe un entier $n$ tel que
toutes les localisations $R_{\mathfrak p}$ aux idéaux premiers
soient régulières de dimension $\leq n$, avec égalité pour au moins
un $\mathfrak p$.
\end{enumerate}
\end{lemma}

\begin{proof}
Cela découle de la discussion ci-dessus. Plus précisément, cela s'obtient
en combinant la définition \ref{algebra-definition-regular} avec le
lemme \ref{algebra-lemma-finite-gl-dim-primes}
et la proposition \ref{algebra-proposition-finite-gl-dim-regular}.
\end{proof}

\begin{lemma}
\label{algebra-lemma-flat-under-regular}
Soit $R \to S$ un homomorphisme local d'anneaux locaux noethériens.
Supposons que $R \to S$ soit plat et que $S$ soit régulier.
Alors $R$ est régulier.
\end{lemma}

\begin{proof}
Soit $\mathfrak m \subset R$ l'idéal maximal,
et soit $\kappa = R/\mathfrak m$ le corps résiduel.
Soit $d = \dim S$.
Choisissons une résolution quelconque $F_\bullet \to \kappa$
où chaque $F_i$ est un $R$-module libre fini. Posons
$K_d = \Ker(F_{d - 1} \to F_{d - 2})$.
Par platitude de $R \to S$, le complexe
$0 \to K_d \otimes_R S \to F_{d - 1} \otimes_R S \to \ldots
\to F_0 \otimes_R S \to \kappa \otimes_R S \to 0$
reste exact. Puisque la dimension globale de $S$
est $d$, voir la proposition \ref{algebra-proposition-regular-finite-gl-dim},
nous voyons que $K_d \otimes_R S$ est un $S$-module libre fini
(voir aussi le lemme \ref{algebra-lemma-independent-resolution}).
D'après le lemme \ref{algebra-lemma-finite-projective-descends}, nous voyons
que $K_d$ est un $R$-module libre fini.
Ainsi, $\kappa$ est de dimension projective finie et $R$ est régulier d'après la
proposition \ref{algebra-proposition-finite-gl-dim-regular}.
\end{proof}








\section{Auslander-Buchsbaum}
\label{algebra-section-Auslander-Buchsbaum}

\noindent
Le résultat suivant se trouve dans \cite{Auslander-Buchsbaum}.

\begin{proposition}
\label{algebra-proposition-Auslander-Buchsbaum}
Soit $R$ un anneau local noethérien. Soit $M$ un $R$-module fini non nul
de dimension projective finie $\text{pd}_R(M)$. Alors nous avons
$$
\text{depth}(R) = \text{pd}_R(M) + \text{depth}(M)
$$
\end{proposition}

\begin{proof}
Nous le prouvons par récurrence sur $\text{depth}(M)$. Le cas le plus intéressant
est celui où $\text{depth}(M) = 0$. Dans ce cas, soit
$$
0 \to R^{n_e} \to R^{n_{e-1}} \to \ldots \to R^{n_0} \to M \to 0
$$
une résolution libre finie minimale, de sorte que $e = \text{pd}_R(M)$.
D'après le lemme \ref{algebra-lemma-add-trivial-complex}, nous pouvons supposer que tous les coefficients
matriciels des homomorphismes du complexe sont contenus dans l'idéal maximal
de $R$. Alors, d'une part, la
proposition \ref{algebra-proposition-what-exact} montre que
$\text{depth}(R) \geq e$. D'autre part, en décomposant la longue
suite exacte en suites exactes courtes
\begin{align*}
0 \to R^{n_e} \to R^{n_{e - 1}} \to K_{e - 2} \to 0,\\
0 \to K_{e - 2} \to R^{n_{e - 2}} \to K_{e - 3} \to 0,\\
\ldots,\\
0 \to K_0 \to R^{n_0} \to M \to 0
\end{align*}
nous voyons, à l'aide du lemme \ref{algebra-lemma-depth-in-ses}, que
\begin{align*}
\text{depth}(K_{e - 2}) \geq \text{depth}(R) - 1,\\
\text{depth}(K_{e - 3}) \geq \text{depth}(R) - 2,\\
\ldots,\\
\text{depth}(K_0) \geq \text{depth}(R) - (e - 1),\\
\text{depth}(M) \geq \text{depth}(R) - e
\end{align*}
et, puisque $\text{depth}(M) = 0$, nous concluons que $\text{depth}(R) \leq e$.
Cela achève la preuve du cas $\text{depth}(M) = 0$.

\medskip\noindent
Étape de récurrence. Si $\text{depth}(M) > 0$, choisissons $x \in \mathfrak m$
qui soit un non-diviseur de zéro à la fois sur $M$ et sur $R$. C'est possible, car
soit $\text{pd}_R(M) > 0$ et $\text{depth}(R) > 0$ d'après la
proposition \ref{algebra-proposition-what-exact} précitée, soit $\text{pd}_R(M) = 0$, auquel
cas $M$ est libre fini et donc aussi $\text{depth}(R) = \text{depth}(M) > 0$.
Ainsi, $\text{depth}(R \oplus M) > 0$ d'après le lemme \ref{algebra-lemma-depth-in-ses}
(par exemple), et nous pouvons trouver $x \in \mathfrak m$ qui soit un non-diviseur de zéro
à la fois sur $R$ et sur $M$. Soit
$$
0 \to R^{n_e} \to R^{n_{e-1}} \to \ldots \to R^{n_0} \to M \to 0
$$
une résolution minimale comme ci-dessus. Une application du lemme du serpent
montre que
$$
0 \to (R/xR)^{n_e} \to (R/xR)^{n_{e-1}} \to \ldots \to (R/xR)^{n_0} \to
M/xM \to 0
$$
est elle aussi une résolution minimale. Ainsi, $\text{pd}_R(M) = \text{pd}_{R/xR}(M/xM)$.
D'après le lemme \ref{algebra-lemma-depth-drops-by-one}, nous avons
$\text{depth}(R/xR) = \text{depth}(R) - 1$ et
$\text{depth}(M/xM) = \text{depth}(M) - 1$.
Jusqu'ici, toutes les profondeurs étaient prises comme $R$-modules, mais nous observons que
$\text{depth}_R(M/xM) = \text{depth}_{R/xR}(M/xM)$, et de même pour $R/xR$.
Par hypothèse de récurrence, nous voyons que la formule
d'Auslander-Buchsbaum est vérifiée pour $M/xM$ sur $R/xR$. Puisque les
profondeurs de $R/xR$ et de $M/xM$ ont toutes deux diminué de un, tandis que la dimension
projective n'a pas changé, nous concluons.
\end{proof}









\section{Homomorphismes et dimension}
\label{algebra-section-homomorphism-dimension}

\noindent
Cette section contient une collection de résultats élémentaires reliant les
dimensions d'anneaux lorsqu'il existe des homomorphismes entre eux.

\begin{lemma}
\label{algebra-lemma-dimension-going-up}
Supposons que $R \to S$ soit un homomorphisme d'anneaux vérifiant soit la montée, voir la
définition \ref{algebra-definition-going-up-down}, soit la descente,
voir la définition \ref{algebra-definition-going-up-down}.
Supposons en outre que $\Spec(S) \to \Spec(R)$
soit surjectif. Alors $\dim(R) \leq \dim(S)$.
\end{lemma}

\begin{proof}
Supposons la montée.
Prenons une chaîne quelconque $\mathfrak p_0 \subset \mathfrak p_1 \subset \ldots
\subset \mathfrak p_e$ d'idéaux premiers de $R$.
Par surjectivité, nous pouvons choisir un idéal premier $\mathfrak q_0$ s'envoyant
sur $\mathfrak p_0$. Par montée, nous pouvons le prolonger en une chaîne
de longueur $e$ d'idéaux premiers $\mathfrak q_i$ au-dessus des
$\mathfrak p_i$. Ainsi, $\dim(S) \geq \dim(R)$.
Le cas de la descente est exactement semblable.
Voir aussi Topologie, lemme \ref{topology-lemma-dimension-specializations-lift},
pour une version purement topologique.
\end{proof}

\begin{lemma}
\label{algebra-lemma-going-up-maximal-on-top}
Supposons que $R \to S$ soit un homomorphisme d'anneaux possédant la propriété de montée,
voir la définition \ref{algebra-definition-going-up-down}. Si
$\mathfrak q \subset S$ est un idéal maximal,
alors l'image réciproque de $\mathfrak q$ dans $R$
est elle aussi un idéal maximal.
\end{lemma}

\begin{proof}
Trivial.
\end{proof}

\begin{lemma}
\label{algebra-lemma-integral-dim-up}
Supposons que $R \to S$ soit un homomorphisme d'anneaux tel que $S$ soit entier sur $R$.
Alors $\dim (R) \geq \dim(S)$, et tout point fermé de $\Spec(S)$
s'envoie sur un point fermé de $\Spec(R)$.
\end{lemma}

\begin{proof}
Cela découle immédiatement des lemmes \ref{algebra-lemma-integral-no-inclusion} et
\ref{algebra-lemma-going-up-maximal-on-top}
et des définitions.
\end{proof}

\begin{lemma}
\label{algebra-lemma-integral-sub-dim-equal}
Supposons $R \subset S$ et $S$ entier sur $R$.
Alors $\dim(R) = \dim(S)$.
\end{lemma}

\begin{proof}
Cela résulte de la combinaison des lemmes
\ref{algebra-lemma-integral-going-up},
\ref{algebra-lemma-integral-overring-surjective},
\ref{algebra-lemma-dimension-going-up} et
\ref{algebra-lemma-integral-dim-up}.
\end{proof}

\begin{definition}
\label{algebra-definition-fibre}
Supposons que $R \to S$ soit un homomorphisme d'anneaux.
Soit $\mathfrak q \subset S$ un idéal premier au-dessus
de l'idéal premier $\mathfrak p$ de $R$.
L'{\it anneau local de la fibre en $\mathfrak q$}
est l'anneau local
$$
S_{\mathfrak q}/\mathfrak pS_{\mathfrak q}
=
(S/\mathfrak pS)_{\mathfrak q}
=
(S \otimes_R \kappa(\mathfrak p))_{\mathfrak q}
$$
\end{definition}

\begin{lemma}
\label{algebra-lemma-dimension-base-fibre-total}
Soit $R \to S$ un homomorphisme d'anneaux noethériens.
Soit $\mathfrak q \subset S$ un idéal premier au-dessus
de l'idéal premier $\mathfrak p$. Alors
$$
\dim(S_{\mathfrak q})
\leq
\dim(R_{\mathfrak p})
+
\dim(S_{\mathfrak q}/\mathfrak pS_{\mathfrak q}).
$$
\end{lemma}

\begin{proof}
Nous utilisons la caractérisation de la dimension de la
proposition \ref{algebra-proposition-dimension}.
Soient $x_1, \ldots, x_d$ des éléments de $\mathfrak p$
engendrant un idéal de définition de $R_{\mathfrak p}$, avec
$d = \dim(R_{\mathfrak p})$.
Soient $y_1, \ldots, y_e$ des éléments de $\mathfrak q$
engendrant un idéal de définition de
$S_{\mathfrak q}/\mathfrak pS_{\mathfrak q}$,
avec $e = \dim(S_{\mathfrak q}/\mathfrak pS_{\mathfrak q})$.
Il est clair que $S_{\mathfrak q}/(x_1, \ldots, x_d, y_1, \ldots, y_e)$
a un idéal maximal nilpotent. Ainsi,
$x_1, \ldots, x_d, y_1, \ldots, y_e$ engendrent un idéal de définition
de $S_{\mathfrak q}$.
\end{proof}

\begin{lemma}
\label{algebra-lemma-dimension-base-fibre-equals-total}
Soit $R \to S$ un homomorphisme d'anneaux noethériens.
Soit $\mathfrak q \subset S$ un idéal premier au-dessus
de l'idéal premier $\mathfrak p$. Supposons que la propriété de descente soit vérifiée
pour $R \to S$ (par exemple, si $R \to S$ est plat, voir le
lemme \ref{algebra-lemma-flat-going-down}). Alors
$$
\dim(S_{\mathfrak q})
=
\dim(R_{\mathfrak p})
+
\dim(S_{\mathfrak q}/\mathfrak pS_{\mathfrak q}).
$$
\end{lemma}

\begin{proof}
D'après le lemme \ref{algebra-lemma-dimension-base-fibre-total},
nous avons une inégalité
$\dim(S_{\mathfrak q}) \leq
\dim(R_{\mathfrak p}) + \dim(S_{\mathfrak q}/\mathfrak pS_{\mathfrak q})$.
Pour obtenir l'égalité, choisissons une chaîne d'idéaux premiers
$\mathfrak pS \subset \mathfrak q_0 \subset \mathfrak q_1 \subset \ldots
\subset \mathfrak q_d = \mathfrak q$, avec
$d = \dim(S_{\mathfrak q}/\mathfrak pS_{\mathfrak q})$.
D'autre part, choisissons une chaîne d'idéaux premiers
$\mathfrak p_0 \subset \mathfrak p_1 \subset \ldots \subset \mathfrak p_e
= \mathfrak p$, avec $e = \dim(R_{\mathfrak p})$.
Par le théorème de descente, nous pouvons choisir
$\mathfrak q_{-1} \subset \mathfrak q_0$ au-dessus de
$\mathfrak p_{e-1}$. Puis nous pouvons choisir
$\mathfrak q_{-2} \subset \mathfrak q_{-1}$ au-dessus de
$\mathfrak p_{e-2}$. En poursuivant par récurrence, nous obtenons finalement
une chaîne
$\mathfrak q_{-e} \subset \ldots \subset \mathfrak q_d$
de longueur $e + d$.
\end{proof}

\begin{lemma}
\label{algebra-lemma-flat-over-regular-with-regular-fibre}
Soit $R \to S$ un homomorphisme local d'anneaux locaux noethériens.
Supposons que
\begin{enumerate}
\item $R$ soit régulier,
\item $S/\mathfrak m_RS$ soit régulier, et
\item $R \to S$ soit plat.
\end{enumerate}
Alors $S$ est régulier.
\end{lemma}

\begin{proof}
D'après le lemme \ref{algebra-lemma-dimension-base-fibre-equals-total},
nous avons
$\dim(S) = \dim(R) + \dim(S/\mathfrak m_RS)$.
Choisissons des générateurs $x_1, \ldots, x_d \in \mathfrak m_R$, avec $d = \dim(R)$,
et choisissons $y_1, \ldots, y_e \in \mathfrak m_S$
qui engendrent l'idéal maximal de $S/\mathfrak m_RS$, avec
$e = \dim(S/\mathfrak m_RS)$. Nous voyons alors que
$x_1, \ldots, x_d, y_1, \ldots, y_e$ sont des éléments qui engendrent
l'idéal maximal de $S$ et que $e + d = \dim(S)$.
\end{proof}

\noindent
Le lemme ci-dessous servira plus loin à montrer que les anneaux de type fini sur
un corps sont de Cohen-Macaulay si et seulement s'ils sont quasi-finis et plats sur
un anneau polynomial. C'est une réciproque partielle du
lemme \ref{algebra-lemma-CM-over-regular-flat}.

\begin{lemma}
\label{algebra-lemma-finite-flat-over-regular-CM}
Soit $R \to S$ un homomorphisme local d'anneaux locaux noethériens.
Supposons $R$ de Cohen-Macaulay.
Si $S$ est fini et plat sur $R$, ou si $S$ est plat sur $R$ et si
$\dim(S) \leq \dim(R)$, alors $S$ est de Cohen-Macaulay et $\dim(R) = \dim(S)$.
\end{lemma}

\begin{proof}
Soit $x_1, \ldots, x_d \in \mathfrak m_R$ une suite régulière
de longueur $d = \dim(R)$. D'après le lemme \ref{algebra-lemma-flat-increases-depth},
elle s'envoie sur une suite régulière dans $S$.
Ainsi, $S$ est de Cohen-Macaulay si $\dim(S) \leq d$. C'est le cas
si $S$ est fini et plat sur $R$, d'après le lemme \ref{algebra-lemma-integral-sub-dim-equal}.
Et, dans le second cas, nous l'avons supposé.
\end{proof}








\section{La formule de dimension}
\label{algebra-section-dimension-formula}

\noindent
Rappelons les définitions d'un anneau caténaire (définition \ref{algebra-definition-catenary})
et d'un anneau universellement caténaire (définition \ref{algebra-definition-universally-catenary}).

\begin{lemma}
\label{algebra-lemma-dimension-formula}
Soit $R \to S$ un homomorphisme d'anneaux.
Soit $\mathfrak q$ un idéal premier de $S$ au-dessus de l'idéal premier $\mathfrak p$ de $R$.
Supposons que
\begin{enumerate}
\item $R$ est noethérien,
\item $R \to S$ est de type fini,
\item $R$, $S$ sont des anneaux intègres, et
\item $R \subset S$.
\end{enumerate}
Alors nous avons
$$
\text{hauteur}(\mathfrak q)
\leq
\text{hauteur}(\mathfrak p) + \text{trdeg}_R(S)
- \text{trdeg}_{\kappa(\mathfrak p)} \kappa(\mathfrak q)
$$
avec égalité si $R$ est universellement caténaire.
\end{lemma}

\begin{proof}
Supposons que $R \subset S' \subset S$, où $S'$ est une sous-algèbre de type fini
sur $R$ de $S$. Dans ce cas, posons $\mathfrak q' = S' \cap \mathfrak q$.
Le lemme pour les homomorphismes d'anneaux $R \to S'$ et $S' \to S$ implique le
lemme pour $R \to S$ par additivité du degré de transcendance dans les tours
d'extensions de corps (Corps, lemme \ref{fields-lemma-transcendence-degree-tower}).
Nous pouvons donc raisonner par récurrence sur le nombre de générateurs
de $S$ sur $R$ et nous ramener au cas où $S$ est engendré par
un élément sur $R$.

\medskip\noindent
Cas I : $S = R[x]$ est une algèbre polynomiale sur $R$.
Dans ce cas, nous avons $\text{trdeg}_R(S) = 1$.
De plus, $R \to S$ est plat et donc
$$
\dim(S_{\mathfrak q}) =
\dim(R_{\mathfrak p}) + \dim(S_{\mathfrak q}/\mathfrak pS_{\mathfrak q})
$$
voir le lemme \ref{algebra-lemma-dimension-base-fibre-equals-total}.
Soit $\mathfrak r = \mathfrak pS$. Alors
$\text{trdeg}_{\kappa(\mathfrak p)} \kappa(\mathfrak q) = 1$
équivaut à $\mathfrak q = \mathfrak r$, et implique que
$\dim(S_{\mathfrak q}/\mathfrak pS_{\mathfrak q}) = 0$.
De même, $\text{trdeg}_{\kappa(\mathfrak p)} \kappa(\mathfrak q) = 0$
équivaut à avoir une inclusion stricte
$\mathfrak r \subset \mathfrak q$, ce qui implique que
$\dim(S_{\mathfrak q}/\mathfrak pS_{\mathfrak q}) = 1$.
Nous avons ainsi terminé le cas I, avec égalité dans tous les cas.

\medskip\noindent
Cas II : $S = R[x]/\mathfrak n$ avec $\mathfrak n \not = 0$.
Dans ce cas, nous avons $\text{trdeg}_R(S) = 0$.
Notons $\mathfrak q' \subset R[x]$ l'idéal premier correspondant à $\mathfrak q$.
Nous avons ainsi
$$
S_{\mathfrak q} = (R[x])_{\mathfrak q'}/\mathfrak n(R[x])_{\mathfrak q'}
$$
D'après le cas précédent, nous avons
$\dim((R[x])_{\mathfrak q'}) =
\dim(R_{\mathfrak p}) + 1
- \text{trdeg}_{\kappa(\mathfrak p)} \kappa(\mathfrak q)$.
Puisque $\mathfrak n \not = 0$, nous voyons que la dimension de
$S_{\mathfrak q}$ diminue d'au moins un, voir le
lemme \ref{algebra-lemma-one-equation},
ce qui démontre l'inégalité du lemme.
Pour voir l'égalité lorsque $R$ est universellement caténaire, remarquons que
$\mathfrak n \subset R[x]$ est un idéal premier de hauteur un puisqu'il correspond
à un idéal premier non nul de $F[x]$, où $F$ est le corps des fractions de $R$.
Par conséquent, toute chaîne maximale d'idéaux premiers de
$S_\mathfrak q = R[x]_{\mathfrak q'}/\mathfrak nR[x]_{\mathfrak q'}$
correspond à une chaîne maximale d'idéaux premiers
de longueur supérieure de 1 entre $\mathfrak q'$ et $(0)$ dans $R[x]$.
Si $R$ est universellement caténaire, elles ont toutes la même longueur, égale à
la hauteur de $\mathfrak q'$. Cela montre que
$\dim(S_\mathfrak q) = \dim(R[x]_{\mathfrak q'}) - 1$
et implique l'égalité voulue.
\end{proof}

\noindent
Le lemme suivant affirme que les morphismes génériquement finis
ont tendance à être quasi-finis en codimension $1$.

\begin{lemma}
\label{algebra-lemma-finite-in-codim-1}
Soit $A \to B$ un homomorphisme d'anneaux.
Supposons que
\begin{enumerate}
\item $A \subset B$ est une extension d'anneaux intègres,
\item l'extension induite des corps de fractions est finie,
\item $A$ est noethérien, et
\item $A \to B$ est de type fini.
\end{enumerate}
Soit $\mathfrak p \subset A$ un idéal premier de hauteur $1$.
Il existe alors au plus un nombre fini d'idéaux premiers de $B$
au-dessus de $\mathfrak p$, et ils sont tous de hauteur $1$.
\end{lemma}

\begin{proof}
D'après la formule de dimension (lemme \ref{algebra-lemma-dimension-formula}),
pour tout idéal premier $\mathfrak q$ au-dessus de $\mathfrak p$, nous avons
$$
\dim(B_{\mathfrak q}) \leq
\dim(A_{\mathfrak p}) - \text{trdeg}_{\kappa(\mathfrak p)} \kappa(\mathfrak q).
$$
Comme l'anneau intègre $B_\mathfrak q$ possède au moins $2$ idéaux premiers, nous voyons que
$\dim(B_{\mathfrak q}) \geq 1$. Nous en concluons que
$\dim(B_{\mathfrak q}) = 1$ et que l'extension
$\kappa(\mathfrak p) \subset \kappa(\mathfrak q)$ est algébrique.
Ainsi, $\mathfrak q$ définit un point fermé de sa fibre
$\Spec(B \otimes_A \kappa(\mathfrak p))$, voir le
lemme \ref{algebra-lemma-finite-residue-extension-closed}.
Puisque $B \otimes_A \kappa(\mathfrak p)$ est un anneau noethérien,
la fibre $\Spec(B \otimes_A \kappa(\mathfrak p))$
est un espace topologique noethérien, voir le
lemme \ref{algebra-lemma-Noetherian-topology}.
Un espace topologique noethérien constitué de points fermés
est fini, voir par exemple
Topologie, lemme \ref{topology-lemma-Noetherian}.
\end{proof}









\section{Dimension des algèbres de type fini sur les corps}
\label{algebra-section-dimension-finite-type-algebras}

\noindent
Dans cette section, nous calculons la dimension d'un anneau de polynômes sur
un corps. Nous démontrons également que la dimension d'une algèbre intègre
de type fini sur un corps est la dimension de ses anneaux locaux aux idéaux
maximaux. Nous établirons le lien avec le degré de transcendance
sur le corps de base dans la
section \ref{algebra-section-dimension-finite-type-algebras-reprise}.

\begin{lemma}
\label{algebra-lemma-dim-affine-space}
Soit $\mathfrak m$ un idéal maximal de $k[x_1, \ldots, x_n]$.
L'idéal $\mathfrak m$ est engendré par $n$ éléments.
La dimension de $k[x_1, \ldots, x_n]_{\mathfrak m}$ est $n$.
Ainsi, $k[x_1, \ldots, x_n]_{\mathfrak m}$ est un anneau local régulier
de dimension $n$.
\end{lemma}

\begin{proof}
D'après le théorème des zéros de Hilbert
(théorème \ref{algebra-theorem-nullstellensatz}),
nous savons que le corps résiduel $\kappa = \kappa(\mathfrak m)$ est
une extension finie de $k$. Notons $\alpha_i \in \kappa$ l'
image de $x_i$. Notons $\kappa_i = k(\alpha_1, \ldots, \alpha_i)
\subset \kappa$, $i = 1, \ldots, n$, et $\kappa_0 = k$.
Remarquons que $\kappa_i = k[\alpha_1, \ldots, \alpha_i]$
d'après la théorie des corps. Définissons par récurrence des éléments
$f_i \in \mathfrak m \cap k[x_1, \ldots, x_i]$
comme suit : soit $P_i(T) \in \kappa_{i-1}[T]$
le polynôme minimal unitaire de $\alpha_i $ sur $\kappa_{i-1}$.
Soit $Q_i(T) \in k[x_1, \ldots, x_{i-1}][T]$ un relèvement unitaire de $P_i(T)$
(de même degré). Posons $f_i = Q_i(x_i)$.
Remarquons que si $d_i = \deg_T(P_i) = \deg_T(Q_i) = \deg_{x_i}(f_i)$,
alors $d_1d_2\ldots d_i = [\kappa_i : k]$ d'après
Corps, lemmes \ref{fields-lemma-multiplicativity-degrees} et
\ref{fields-lemma-degree-minimal-polynomial}.

\medskip\noindent
Nous affirmons que, pour tout $i = 0, 1, \ldots, n$, il existe un
isomorphisme
$$
\psi_i : k[x_1, \ldots, x_i] /(f_1, \ldots, f_i) \cong \kappa_i.
$$
Par construction, le composé
$k[x_1, \ldots, x_i] \to k[x_1, \ldots, x_n] \to \kappa$
a pour image $\kappa_i$, et $f_1, \ldots, f_i$ appartiennent
au noyau. Cela donne un homomorphisme surjectif.
Nous démontrons par récurrence que $\psi_i$ est injectif. C'est clair pour $i = 0$.
Supposons l'énoncé vrai pour $i$ et démontrons-le pour $i + 1$.
L'extension d'anneaux $k[x_1, \ldots, x_i]/(f_1, \ldots, f_i) \to
k[x_1, \ldots, x_{i + 1}]/(f_1, \ldots, f_{i + 1})$
est engendrée par $1$ élément sur un corps et soumise à une
équation irréductible. D'après la théorie élémentaire des corps,
$k[x_1, \ldots, x_{i + 1}]/(f_1, \ldots, f_{i + 1})$
est un corps, et $\psi_i$ est donc injectif.

\medskip\noindent
Cela implique que $\mathfrak m = (f_1, \ldots, f_n)$.
De plus, nous en concluons également que
$$
k[x_1, \ldots, x_n]/(f_1, \ldots, f_i)
\cong
\kappa_i[x_{i + 1}, \ldots, x_n].
$$
Ainsi, $(f_1, \ldots, f_i)$ est un idéal premier. Par conséquent,
$$
(0) \subset (f_1) \subset (f_1, f_2) \subset \ldots \subset
(f_1, \ldots, f_n) = \mathfrak m
$$
est une chaîne d'idéaux premiers de longueur $n$. Le lemme en résulte.
\end{proof}

\begin{proposition}
\label{algebra-proposition-finite-gl-dim-polynomial-ring}
Une algèbre polynomiale à $n$ variables sur un corps est un anneau régulier.
Elle est de dimension globale $n$. Toutes ses localisations en des idéaux maximaux
sont des anneaux locaux réguliers de dimension $n$.
\end{proposition}

\begin{proof}
D'après le lemme \ref{algebra-lemma-dim-affine-space},
toutes les localisations $k[x_1, \ldots, x_n]_{\mathfrak m}$
en des idéaux maximaux sont des anneaux locaux réguliers de dimension $n$. Ainsi,
nous concluons par le lemme \ref{algebra-lemma-finite-gl-dim-finite-dim-regular}.
\end{proof}

\begin{lemma}
\label{algebra-lemma-dimension-height-polynomial-ring}
Soit $k$ un corps.
Soit $\mathfrak p \subset \mathfrak q \subset k[x_1, \ldots, x_n]$
un couple d'idéaux premiers.
Toute chaîne maximale d'idéaux premiers entre $\mathfrak p$ et $\mathfrak q$
est de longueur $\text{hauteur}(\mathfrak q) - \text{hauteur}(\mathfrak p)$.
\end{lemma}

\begin{proof}
D'après la
proposition \ref{algebra-proposition-finite-gl-dim-polynomial-ring},
tout anneau local de $k[x_1, \ldots, x_n]$ est régulier.
Tous les anneaux locaux sont donc de Cohen-Macaulay, voir le
lemme \ref{algebra-lemma-regular-ring-CM}.
Les anneaux locaux aux idéaux maximaux sont de dimension $n$, donc
toute chaîne maximale d'idéaux premiers de $k[x_1, \ldots, x_n]$
est de longueur $n$, voir le
lemme \ref{algebra-lemma-maximal-chain-CM}.
Ainsi, toute chaîne maximale d'idéaux premiers entre $(0)$ et $\mathfrak p$
est de longueur $\text{hauteur}(\mathfrak p)$, voir le
lemme \ref{algebra-lemma-CM-dim-formula}
par exemple.
En rassemblant ces faits, nous obtenons l'assertion du lemme.
\end{proof}

\begin{lemma}
\label{algebra-lemma-dimension-spell-it-out}
Soit $k$ un corps.
Soit $S$ une $k$-algèbre de type fini qui est un anneau intègre.
Alors $\dim(S) = \dim(S_{\mathfrak m})$ pour tout idéal
maximal $\mathfrak m$ de $S$. Autrement dit, toute chaîne maximale
d'idéaux premiers est de longueur égale à la dimension de $S$.
\end{lemma}

\begin{proof}
Écrivons $S = k[x_1, \ldots, x_n]/\mathfrak p$.
D'après la proposition \ref{algebra-proposition-finite-gl-dim-polynomial-ring} et le
lemme \ref{algebra-lemma-dimension-height-polynomial-ring},
toutes les chaînes maximales d'idéaux premiers de $S$ (qui se terminent nécessairement
par un idéal maximal) sont de longueur $n - \text{hauteur}(\mathfrak p)$.
Ce nombre est donc la dimension de $S$ et de $S_{\mathfrak m}$
pour tout idéal maximal $\mathfrak m$ de $S$.
\end{proof}

\noindent
Rappelons que nous avons défini la
dimension $\dim_x(X)$ d'un espace topologique $X$ en un point $x$
dans Topologie, définition \ref{topology-definition-Krull}.

\begin{lemma}
\label{algebra-lemma-dimension-at-a-point-finite-type-over-field}
Soit $k$ un corps.
Soit $S$ une $k$-algèbre de type fini.
Soit $X = \Spec(S)$.
Soit $\mathfrak p \subset S$ un idéal premier et soit
$x \in X$ le point correspondant.
Les nombres suivants sont égaux :
\begin{enumerate}
\item $\dim_x(X)$,
\item $\max \dim(Z)$, où le maximum porte sur les
composantes irréductibles $Z$ de $X$ qui passent par $x$, et
\item $\min \dim(S_{\mathfrak m})$, où le minimum
porte sur les idéaux maximaux $\mathfrak m$ tels que
$\mathfrak p \subset \mathfrak m$.
\end{enumerate}
\end{lemma}

\begin{proof}
Soit $X = \bigcup_{i \in I} Z_i$ la décomposition de $X$ en
ses composantes irréductibles. Elles sont en nombre fini
(voir les
lemmes \ref{algebra-lemma-obvious-Noetherian} et \ref{algebra-lemma-Noetherian-topology}).
Soit $I' = \{i \mid x \in Z_i\}$, et soit
$T = \bigcup_{i \not \in I'} Z_i$. Alors $U = X \setminus T$
est un ouvert de $X$ qui contient le point $x$.
Le nombre figurant en (2) vaut $\max_{i \in I'} \dim(Z_i)$.
Pour tout ouvert $W \subset U$ tel que $x \in W$,
les composantes irréductibles de $W$ sont les ensembles irréductibles
$W_i = Z_i \cap W$ pour $i \in I'$, et $x$ appartient
à chacun d'eux.
Remarquons que chaque $W_i$, $i \in I'$, contient un point fermé, car
$X$ est jacobsonien, voir la section \ref{algebra-section-ring-jacobson}.
Puisque $W_i \subset Z_i$, nous avons $\dim(W_i) \leq \dim(Z_i)$.
L'existence d'un point fermé implique, par le lemme
\ref{algebra-lemma-dimension-spell-it-out}, qu'il existe une chaîne de
fermés irréductibles de longueur égale à $\dim(Z_i)$ dans l'ouvert $W_i$.
Ainsi, $\dim(W_i) = \dim(Z_i)$ pour tout $i \in I'$. Par conséquent, $\dim(W)$
est égal au nombre figurant en (2). Cela démontre que (1) $ = $ (2).

\medskip\noindent
Soit $\mathfrak m \supset \mathfrak p$ un idéal maximal quelconque
contenant $\mathfrak p$. Soit $x_0 \in X$ le point
correspondant. Tout d'abord, $x_0$ appartient à toutes les
composantes irréductibles $Z_i$, $i \in I'$. Notons $\mathfrak q_i$
les idéaux premiers minimaux de $S$ correspondant aux
composantes irréductibles $Z_i$. Pour chaque $i$ tel que
$x_0 \in Z_i$ (ce qui équivaut à $\mathfrak m \supset \mathfrak q_i$),
nous avons une surjection
$$
S_{\mathfrak m} \longrightarrow
S_\mathfrak m/\mathfrak q_i S_\mathfrak m =(S/\mathfrak q_i)_{\mathfrak m}
$$
De plus, les idéaux premiers $\mathfrak q_i S_\mathfrak m$ ainsi obtenus
épuisent les idéaux premiers
minimaux de l'anneau local noethérien $S_{\mathfrak m}$, voir le
lemme \ref{algebra-lemma-irreducible-components-containing-x}.
Nous concluons, en utilisant le lemme \ref{algebra-lemma-dimension-spell-it-out},
que la dimension de $S_{\mathfrak m}$ est le
maximum des dimensions des $Z_i$ qui passent par $x_0$.
Pour terminer la démonstration du lemme, il suffit de montrer que
nous pouvons choisir $x_0$ de sorte que $x_0 \in Z_i \Rightarrow i \in I'$.
Comme $S$ est jacobsonien (ainsi que nous l'avons vu ci-dessus),
il suffit de montrer que $V(\mathfrak p) \setminus T$
(où $T$ est comme ci-dessus) est non vide. Cela est clair, car il
contient le point $x$ (c'est-à-dire $\mathfrak p$).
\end{proof}

\begin{lemma}
\label{algebra-lemma-dimension-closed-point-finite-type-field}
Soit $k$ un corps.
Soit $S$ une $k$-algèbre de type fini.
Soit $X = \Spec(S)$.
Soit $\mathfrak m \subset S$ un idéal maximal et soit
$x \in X$ le point fermé associé.
Alors $\dim_x(X) = \dim(S_{\mathfrak m})$.
\end{lemma}

\begin{proof}
C'est un cas particulier du
lemme \ref{algebra-lemma-dimension-at-a-point-finite-type-over-field}.
\end{proof}

\begin{lemma}
\label{algebra-lemma-disjoint-decomposition-CM-algebra}
Soit $k$ un corps.
Soit $S$ une $k$-algèbre de type fini.
Supposons que $S$ est de Cohen-Macaulay.
Alors $\Spec(S) = \coprod T_d$ est une union disjointe finie de
parties ouvertes et fermées $T_d$, où $T_d$ est équidimensionnelle
(voir Topologie, définition \ref{topology-definition-equidimensional})
de dimension $d$. De manière équivalente, $S$ est un produit d'anneaux
$S_d$, $d = 0, \ldots, \dim(S)$, tels que tout idéal maximal
$\mathfrak m$ de $S_d$ est de hauteur $d$.
\end{lemma}

\begin{proof}
L'équivalence des deux énoncés résulte du
lemme \ref{algebra-lemma-disjoint-implies-product}.
Soit $\mathfrak m \subset S$ un idéal maximal.
Toute chaîne maximale d'idéaux premiers de $S_{\mathfrak m}$ est
de même longueur, égale à $\dim(S_{\mathfrak m})$, voir le
lemme \ref{algebra-lemma-maximal-chain-CM}. Ainsi, les composantes irréductibles
qui passent par le point correspondant à $\mathfrak m$
sont toutes de dimension égale à $\dim(S_{\mathfrak m})$, voir le
lemme \ref{algebra-lemma-dimension-spell-it-out}.
Puisque $\Spec(S)$ est un espace topologique jacobsonien,
l'intersection
de deux quelconques de ses composantes irréductibles contient un point fermé si elle est non vide,
voir les
lemmes \ref{algebra-lemma-finite-type-field-Jacobson} et
\ref{algebra-lemma-jacobson}.
Nous avons donc montré que deux composantes irréductibles quelconques
qui se rencontrent sont de même dimension. Le lemme en résulte
facilement, ainsi que du fait que $\Spec(S)$
n'a qu'un nombre fini de composantes irréductibles (voir les
lemmes \ref{algebra-lemma-obvious-Noetherian} et \ref{algebra-lemma-Noetherian-topology}).
\end{proof}
















\section{Normalisation de Noether}
\label{algebra-section-Noether-normalization}

\noindent
Dans cette section, nous démontrons des variantes du lemme de normalisation de Noether.
L'ingrédient essentiel que nous utiliserons est contenu dans les deux lemmes suivants.

\begin{lemma}
\label{algebra-lemma-helper}
Soit $n \in \mathbf{N}$.
Soit $N$ un ensemble fini non vide
de multi-indices $\nu = (\nu_1, \ldots, \nu_n)$.
Pour $e = (e_1, \ldots, e_n)$, nous posons $e \cdot \nu = \sum e_i\nu_i$.
Alors, pour $e_1 \gg e_2 \gg \ldots \gg e_{n-1} \gg e_n$, nous avons :
si $\nu, \nu' \in N$, alors
$$
(e \cdot \nu = e \cdot \nu')
\Leftrightarrow
(\nu = \nu')
$$
\end{lemma}

\begin{proof}
Écrivons $N = \{\nu_j\}$ avec $\nu_j = (\nu_{j1}, \ldots, \nu_{jn})$.
Soit $A_i = \max_j \nu_{ji} - \min_j \nu_{ji}$. Si, pour tout $i$, nous avons
$e_{i - 1} > A_ie_i + A_{i + 1}e_{i + 1} + \ldots + A_ne_n$, alors
le lemme est vrai. Supposons en effet que $e \cdot (\nu - \nu') = 0$. Alors, pour
$n \ge 2$,
$$
e_1(\nu_1 - \nu'_1) = \sum\nolimits_{i = 2}^n e_i(\nu'_i - \nu_i).
$$
Nous pouvons supposer que $(\nu_1 - \nu'_1) \ge 0$. Si $(\nu_1 - \nu'_1) > 0$, alors
$$
e_1(\nu_1 - \nu'_1) \ge e_1 >
A_2e_2 + \ldots + A_ne_n \ge
\sum\nolimits_{i = 2}^n e_i|\nu'_i - \nu_i| \ge
\sum\nolimits_{i = 2}^n e_i(\nu'_i - \nu_i).
$$
Cette contradiction implique que $\nu'_1 = \nu_1$. Par
récurrence, $\nu'_i = \nu_i$ pour $2 \le i \le n$.
\end{proof}

\begin{lemma}
\label{algebra-lemma-helper-polynomial}
Soit $R$ un anneau. Soit $g \in R[x_1, \ldots, x_n]$ un élément
qui n'est pas constant, c'est-à-dire $g \not \in R$.
Pour $e_1 \gg e_2 \gg \ldots \gg e_{n-1} \gg e_n = 1$, le polynôme
$$
g(x_1 + x_n^{e_1}, x_2 + x_n^{e_2}, \ldots, x_{n - 1} + x_n^{e_{n - 1}}, x_n)
=
ax_n^d + \text{termes de degré inférieur en }x_n
$$
où $d > 0$ et $a \in R$ est l'un des coefficients non nuls de $g$.
\end{lemma}

\begin{proof}
Écrivons $g = \sum_{\nu \in N} a_\nu x^\nu$, où les $a_\nu \in R$ sont non nuls.
Ici, $N$ est un ensemble fini de multi-indices comme dans le
lemme \ref{algebra-lemma-helper},
et $x^\nu = x_1^{\nu_1} \ldots x_n^{\nu_n}$.
Remarquons que le terme dominant de
$$
(x_1 + x_n^{e_1})^{\nu_1} \ldots (x_{n-1} + x_n^{e_{n-1}})^{\nu_{n-1}}
x_n^{\nu_n}
\quad\text{est}\quad
x_n^{e_1\nu_1 + \ldots + e_{n-1}\nu_{n-1} + \nu_n}.
$$
Le lemme résulte donc du
lemme \ref{algebra-lemma-helper},
qui garantit qu'il existe exactement un terme non nul $a_\nu x^\nu$ de $g$
qui donne naissance au terme dominant de
$g(x_1 + x_n^{e_1}, x_2 + x_n^{e_2}, \ldots, x_{n - 1} + x_n^{e_{n - 1}},
x_n)$, c'est-à-dire $a = a_\nu$ pour l'unique $\nu \in N$ tel que
$e \cdot \nu$ soit maximal.
\end{proof}

\begin{lemma}
\label{algebra-lemma-one-relation}
Soit $k$ un corps.
Soit $S = k[x_1, \ldots, x_n]/I$ pour un idéal propre $I$.
Si $I \not = 0$, alors il existe $y_1, \ldots, y_{n-1} \in k[x_1, \ldots, x_n]$
tels que $S$ est fini sur $k[y_1, \ldots, y_{n-1}]$. De plus, nous pouvons
choisir les $y_i$ dans la sous-$\mathbf{Z}$-algèbre de $k[x_1, \ldots, x_n]$
engendrée par $x_1, \ldots, x_n$.
\end{lemma}

\begin{proof}
Choisissons $f \in I$, $f\not = 0$. Il suffit de démontrer le lemme
pour $k[x_1, \ldots, x_n]/(f)$, puisque $S$ est un quotient de cet anneau.
Nous prendrons $y_i = x_i - x_n^{e_i}$, $i = 1, \ldots, n-1$,
pour des entiers $e_i$ convenables. Quand cela fonctionne-t-il ? Il suffit
de montrer que $\overline{x_n} \in k[x_1, \ldots, x_n]/(f)$
est entier sur l'anneau $k[y_1, \ldots, y_{n-1}]$. L'
équation de $\overline{x_n}$ sur cet anneau est
$$
f(y_1 + x_n^{e_1}, \ldots, y_{n-1} + x_n^{e_{n-1}}, x_n) = 0.
$$
Nous avons donc terminé si nous pouvons montrer qu'il existe des entiers $e_i$ tels
que le coefficient dominant par rapport à $x_n$ de l'équation
ci-dessus est un élément non nul de $k$. On peut l'obtenir, par exemple,
en choisissant $e_1 \gg e_2 \gg \ldots \gg e_{n-1}$, voir le
lemme \ref{algebra-lemma-helper-polynomial}.
\end{proof}

\begin{lemma}
\label{algebra-lemma-Noether-normalization}
\begin{slogan}
Normalisation de Noether
\end{slogan}
Soit $k$ un corps. Soit $S = k[x_1, \ldots, x_n]/I$ pour un idéal $I$.
Si $I \neq (1)$, il existe $r\geq 0$, et
$y_1, \ldots, y_r \in k[x_1, \ldots, x_n]$
tels que (a) l'homomorphisme $k[y_1, \ldots, y_r] \to S$ est injectif,
où la source est l'anneau de polynômes en $y_1, \ldots, y_r$,
et (b) l'homomorphisme $k[y_1, \ldots, y_r] \to S$ est fini.
Dans ce cas, l'entier $r$ est la dimension de $S$.
De plus, nous pouvons choisir les $y_i$ dans la
sous-$\mathbf{Z}$-algèbre de $k[x_1, \ldots, x_n]$
engendrée par $x_1, \ldots, x_n$.
\end{lemma}

\begin{proof}
Raisonnons par récurrence sur $n$, le cas $n = 0$ étant trivial.
Si $I = 0$, prenons $r = n$ et $y_i = x_i$.
Si $I \not = 0$, choisissons $y_1, \ldots, y_{n-1}$
comme dans le lemme \ref{algebra-lemma-one-relation}. Soit
$S' \subset S$ le sous-anneau engendré par
les images des $y_i$. Par récurrence, nous pouvons
choisir $r$ et $z_1, \ldots, z_r \in k[y_1, \ldots, y_{n-1}]$
de sorte que (a), (b) soient vérifiées pour $k[z_1, \ldots, z_r]
\to S'$. Puisque $S' \to S$ est injectif et fini,
nous voyons que (a), (b) sont vérifiées pour $k[z_1, \ldots, z_r]
\to S$. La dernière assertion résulte du lemme
\ref{algebra-lemma-integral-sub-dim-equal}.
\end{proof}

\begin{lemma}
\label{algebra-lemma-Noether-normalization-at-point}
Soit $k$ un corps.
Soit $S$ une $k$-algèbre de type fini et notons $X = \Spec(S)$.
Soit $\mathfrak q$ un idéal premier de $S$, et soit $x \in X$ le point
correspondant. Il existe $g \in S$, $g \not \in \mathfrak q$,
tel que $\dim(S_g) = \dim_x(X) =: d$ et tel qu'
il existe un homomorphisme fini injectif $k[y_1, \ldots, y_d] \to S_g$.
\end{lemma}

\begin{proof}
Remarquons que, par définition, $\dim_x(X)$ est le minimum
des dimensions des $S_g$ pour $g \in S$, $g \not \in \mathfrak q$, c'est-à-dire
que le minimum est atteint. Le lemme résulte donc du
lemme \ref{algebra-lemma-Noether-normalization}.
\end{proof}

\begin{lemma}
\label{algebra-lemma-refined-Noether-normalization}
Soit $k$ un corps. Soit $\mathfrak q \subset k[x_1, \ldots, x_n]$
un idéal premier. Posons $r = \text{trdeg}_k\ \kappa(\mathfrak q)$.
Alors il existe un homomorphisme fini d'anneaux
$\varphi : k[y_1, \ldots, y_n] \to k[x_1, \ldots, x_n]$ tel
que $\varphi^{-1}(\mathfrak q) = (y_{r + 1}, \ldots, y_n)$.
\end{lemma}

\begin{proof}
Raisonnons par récurrence sur $n$. Le cas $n = 0$ est clair. Supposons $n > 0$.
Si $r = n$, alors $\mathfrak q = (0)$ et le résultat est clair.
Choisissons un élément non nul $f \in \mathfrak q$. Bien sûr, $f$ n'est pas constant.
Après avoir appliqué un automorphisme de la forme
$$
k[x_1, \ldots, x_n] \longrightarrow k[x_1, \ldots, x_n],
\quad
x_n \mapsto x_n,
\quad
x_i \mapsto x_i + x_n^{e_i}\ (i < n)
$$
nous pouvons supposer que $f$ est unitaire en $x_n$ sur $k[x_1, \ldots, x_n]$, voir le
lemme \ref{algebra-lemma-helper-polynomial}. Ainsi, l'homomorphisme d'anneaux
$$
k[y_1, \ldots, y_n] \longrightarrow k[x_1, \ldots, x_n],
\quad
y_n \mapsto f,
\quad
y_i \mapsto x_i\ (i < n)
$$
est fini. De plus, $y_n \in \mathfrak q \cap k[y_1, \ldots, y_n]$ par
construction. Ainsi,
$\mathfrak q \cap k[y_1, \ldots, y_n] = \mathfrak pk[y_1, \ldots, y_n] + (y_n)$,
où $\mathfrak p \subset k[y_1, \ldots, y_{n - 1}]$ est un idéal premier.
Remarquons que $\kappa(\mathfrak p) \subset \kappa(\mathfrak q)$ est finie, et
donc $r = \text{trdeg}_k\ \kappa(\mathfrak p)$.
Appliquons l'hypothèse de récurrence au couple
$(k[y_1, \ldots, y_{n - 1}], \mathfrak p)$ ; nous obtenons un homomorphisme fini d'anneaux
$k[z_1, \ldots, z_{n - 1}] \to k[y_1, \ldots, y_{n - 1}]$ tel que
$\mathfrak p \cap k[z_1, \ldots, z_{n - 1}] = (z_{r + 1}, \ldots, z_{n - 1})$.
Nous prolongeons l'homomorphisme d'anneaux
$k[z_1, \ldots, z_{n - 1}] \to k[y_1, \ldots, y_{n - 1}]$
en un homomorphisme d'anneaux
$k[z_1, \ldots, z_n] \to k[y_1, \ldots, y_n]$
en envoyant $z_n$ sur $y_n$.
Le composé des homomorphismes d'anneaux
$$
k[z_1, \ldots, z_n] \to k[y_1, \ldots, y_n] \to k[x_1, \ldots, x_n]
$$
résout le problème.
\end{proof}

\begin{lemma}
\label{algebra-lemma-Noether-normalization-over-a-domain}
Soit $R \to S$ un homomorphisme injectif d'anneaux de type fini. Supposons que $R$ est intègre.
Alors il existe un entier $d$ et une factorisation
$$
R \to R[y_1, \ldots, y_d] \to S' \to S
$$
par des homomorphismes injectifs tels que $S'$ soit fini sur $R[y_1, \ldots, y_d]$
et que $S'_f \cong S_f$ pour un élément non nul $f \in R$.
\end{lemma}

\begin{proof}
Choisissons $x_1, \ldots, x_n \in S$ qui engendrent $S$ sur $R$.
Soit $K$ le corps des fractions de $R$ et $S_K = S \otimes_R K$.
D'après le lemme \ref{algebra-lemma-Noether-normalization},
nous pouvons trouver $y_1, \ldots, y_d \in S$ tels que $K[y_1, \ldots, y_d] \to S_K$
soit un homomorphisme fini injectif. Remarquons que $y_i \in S$, car nous pouvons choisir les
$y_j$ dans la $\mathbf{Z}$-algèbre engendrée par $x_1, \ldots, x_n$.
Comme un homomorphisme fini d'anneaux est entier (voir le
lemme \ref{algebra-lemma-finite-is-integral}),
nous pouvons trouver des polynômes unitaires $P_i \in K[y_1, \ldots, y_d][T]$ tels que
$P_i(x_i) = 0$ dans $S_K$. Soit $f \in R$ un élément non nul tel que
$fP_i \in R[y_1, \ldots, y_d][T]$ pour tout $i$. Alors $fP_i(x_i)$
s'envoie sur zéro dans $S_K$. Par conséquent, après avoir remplacé $f$ par un autre
élément non nul de $R$, nous pouvons aussi supposer que $fP_i(x_i)$ est nul dans $S$.
Posons $x_i' = fx_i$ et soit $S' \subset S$ la sous-$R$-algèbre engendrée par
$y_1, \ldots, y_d$ et $x'_1, \ldots, x'_n$. Remarquons que $x'_i$ est entier sur
$R[y_1, \ldots, y_d]$, car nous avons $Q_i(x_i') = 0$, où
$Q_i = f^{\deg_T(P_i)}P_i(T/f)$ est un polynôme unitaire en
$T$ à coefficients dans $R[y_1, \ldots, y_d]$ par notre choix de $f$.
Ainsi, $R[y_1, \ldots, y_d] \subset S'$ est fini d'après le
lemme \ref{algebra-lemma-characterize-finite-in-terms-of-integral}.
Puisque $S' \subset S$, nous avons $S'_f \subset S_f$ (la localisation est exacte).
D'autre part, les éléments $x_i = x'_i/f$ de $S'_f$ engendrent $S_f$
sur $R_f$, et donc $S'_f \to S_f$ est surjectif. Par conséquent,
$S'_f \cong S_f$, ce qui termine la démonstration.
\end{proof}





\section{Dimension des algèbres de type fini sur les corps, reprise}
\label{algebra-section-dimension-finite-type-algebras-reprise}

\noindent
Cette section est une suite de la
section \ref{algebra-section-dimension-finite-type-algebras}.
Dans cette section, nous établissons le lien entre la dimension et le
degré de transcendance sur le corps de base pour les anneaux intègres de type fini
sur un corps.

\begin{lemma}
\label{algebra-lemma-dimension-prime-polynomial-ring}
Soit $k$ un corps.
Soit $S$ une $k$-algèbre de type fini qui est un anneau intègre.
Soit $K$ le corps des fractions de $S$.
Soit $r = \text{trdeg}(K/k)$ le degré de transcendance de $K$ sur $k$.
Alors $\dim(S) = r$. De plus, l'anneau local de $S$ en tout idéal
maximal est de dimension $r$.
\end{lemma}

\begin{proof}
Nous pouvons écrire $S = k[x_1, \ldots, x_n]/\mathfrak p$.
D'après le lemme \ref{algebra-lemma-dimension-height-polynomial-ring},
tous les anneaux locaux de $S$ en des idéaux maximaux ont la même dimension.
Appliquons le lemme \ref{algebra-lemma-Noether-normalization}.
Nous obtenons un homomorphisme fini injectif d'anneaux
$$
k[y_1, \ldots, y_d] \to S
$$
avec $d = \dim(S)$. Il est clair que $k(y_1, \ldots, y_d) \subset K$
est une extension finie, ce qui termine la démonstration.
\end{proof}

\begin{lemma}
\label{algebra-lemma-tr-deg-specialization}
Soit $k$ un corps. Soit $S$ une $k$-algèbre de type fini.
Soient $\mathfrak q \subset \mathfrak q' \subset S$ des idéaux
premiers distincts. Alors
$\text{trdeg}_k\ \kappa(\mathfrak q') < \text{trdeg}_k\ \kappa(\mathfrak q)$.
\end{lemma}

\begin{proof}
D'après le
lemme \ref{algebra-lemma-dimension-prime-polynomial-ring},
nous avons $\dim V(\mathfrak q) = \text{trdeg}_k\ \kappa(\mathfrak q)$,
et de même pour $\mathfrak q'$. Le résultat en découle, car l'inclusion stricte
$V(\mathfrak q') \subset V(\mathfrak q)$
implique une inégalité stricte entre les dimensions.
\end{proof}

\noindent
Le lemme suivant généralise le
lemme \ref{algebra-lemma-dimension-closed-point-finite-type-field}.

\begin{lemma}
\label{algebra-lemma-dimension-at-a-point-finite-type-field}
Soit $k$ un corps.
Soit $S$ une $k$-algèbre de type fini.
Soit $X = \Spec(S)$.
Soit $\mathfrak p \subset S$ un idéal premier,
et soit $x \in X$ le point correspondant.
Alors nous avons
$$
\dim_x(X) = \dim(S_{\mathfrak p}) + \text{trdeg}_k\ \kappa(\mathfrak p).
$$
\end{lemma}

\begin{proof}
D'après le lemme \ref{algebra-lemma-dimension-prime-polynomial-ring}, nous savons que
$r = \text{trdeg}_k\ \kappa(\mathfrak p)$ est égal à la
dimension de $V(\mathfrak p)$.
Choisissons une chaîne maximale quelconque d'idéaux premiers
$\mathfrak p \subset \mathfrak p_1 \subset \ldots \subset \mathfrak p_r$
qui commence par $\mathfrak p$ dans $S$.
Elle est de longueur $r$ d'après le lemme \ref{algebra-lemma-dimension-spell-it-out}.
Soient $\mathfrak q_j$, $j \in J$, les idéaux premiers
minimaux de $S$ qui sont contenus dans $\mathfrak p$.
Ils correspondent $1-1$ aux idéaux premiers minimaux de $S_{\mathfrak p}$
par la règle $\mathfrak q_j \mapsto \mathfrak q_jS_{\mathfrak p}$.
D'après le lemme \ref{algebra-lemma-dimension-at-a-point-finite-type-over-field},
nous savons que $\dim_x(X)$ est égal
au maximum des dimensions des anneaux $S/\mathfrak q_j$.
Pour chaque $j$, choisissons une chaîne maximale d'idéaux premiers
$\mathfrak q_j \subset \mathfrak p'_1 \subset \ldots \subset \mathfrak p'_{s(j)}
= \mathfrak p$.
Alors $\dim(S_{\mathfrak p}) = \max_{j \in J} s(j)$.
Or chaque chaîne
$$
\mathfrak q_j \subset \mathfrak p'_1 \subset \ldots \subset
\mathfrak p'_{s(j)} = \mathfrak p \subset
\mathfrak p_1 \subset \ldots \subset \mathfrak p_r
$$
est une chaîne maximale dans $S/\mathfrak q_j$ et, d'après ce qui précède,
nous avons
$\dim_x(X) = \max_{j \in J} r + s(j)$.
Le lemme en résulte.
\end{proof}

\noindent
Le lemme suivant affirme que la codimension d'un spectre de type fini
dans un autre est la différence des hauteurs.

\begin{lemma}
\label{algebra-lemma-codimension}
Soit $k$ un corps.
Soit $S' \to S$ une surjection de $k$-algèbres de type fini.
Soit $\mathfrak p \subset S$ un idéal premier,
et soit $\mathfrak p'$ l'idéal premier correspondant de $S'$.
Soient $X = \Spec(S)$, resp.\ $X' = \Spec(S')$,
et soient $x \in X$, resp. $x'\in X'$, les points correspondant
à $\mathfrak p$, resp.\ $\mathfrak p'$.
Alors
$$
\dim_{x'} X' - \dim_x X =
\text{hauteur}(\mathfrak p') - \text{hauteur}(\mathfrak p).
$$
\end{lemma}

\begin{proof}
C'est une conséquence immédiate du lemme \ref{algebra-lemma-dimension-at-a-point-finite-type-field}.
\end{proof}

\begin{lemma}
\label{algebra-lemma-dimension-preserved-field-extension}
Soit $k$ un corps.
Soit $S$ une $k$-algèbre de type fini.
Soit $K/k$ une extension de corps.
Alors $\dim(S) = \dim(K \otimes_k S)$.
\end{lemma}

\begin{proof}
D'après le lemme \ref{algebra-lemma-Noether-normalization},
il existe un homomorphisme fini injectif
$k[y_1, \ldots, y_d] \to S$ avec $d = \dim(S)$.
Puisque $K$ est plat sur $k$, nous obtenons aussi un homomorphisme fini injectif
$K[y_1, \ldots, y_d] \to K \otimes_k S$.
Le résultat découle du lemme \ref{algebra-lemma-integral-sub-dim-equal}.
\end{proof}

\begin{lemma}
\label{algebra-lemma-dimension-at-a-point-preserved-field-extension}
Soit $k$ un corps.
Soit $S$ une $k$-algèbre de type fini.
Posons $X = \Spec(S)$.
Soit $K/k$ une extension de corps.
Posons $S_K = K \otimes_k S$, et $X_K = \Spec(S_K)$.
Soit $\mathfrak q \subset S$ un idéal premier correspondant à $x \in X$,
et soit $\mathfrak q_K \subset S_K$ un idéal premier correspondant
à $x_K \in X_K$ au-dessus de $\mathfrak q$.
Alors $\dim_x X = \dim_{x_K} X_K$.
\end{lemma}

\begin{proof}
Choisissons une présentation $S = k[x_1, \ldots, x_n]/I$.
Cela donne une présentation
$K \otimes_k S = K[x_1, \ldots, x_n]/(K \otimes_k I)$.
Soit $\mathfrak q_K' \subset K[x_1, \ldots, x_n]$,
resp.\ $\mathfrak q' \subset k[x_1, \ldots, x_n]$, les
idéaux premiers correspondants. Considérons le diagramme commutatif
suivant d'anneaux locaux noethériens :
$$
\xymatrix{
K[x_1, \ldots, x_n]_{\mathfrak q_K'} \ar[r] &
(K \otimes_k S)_{\mathfrak q_K} \\
k[x_1, \ldots, x_n]_{\mathfrak q'} \ar[r] \ar[u] &
S_{\mathfrak q} \ar[u]
}
$$
Les deux flèches verticales sont plates, car ce sont des localisations des
homomorphismes plats d'anneaux $S \to S_K$ et
$k[x_1, \ldots, x_n] \to K[x_1, \ldots, x_n]$.
De plus, les flèches verticales ont les mêmes anneaux fibres.
Par conséquent, il résulte du
lemme \ref{algebra-lemma-dimension-base-fibre-equals-total} que
$\text{hauteur}(\mathfrak q') - \text{hauteur}(\mathfrak q)
= \text{hauteur}(\mathfrak q_K') - \text{hauteur}(\mathfrak q_K)$.
Notons $x' \in X' = \Spec(k[x_1, \ldots, x_n])$
et $x'_K \in X'_K = \Spec(K[x_1, \ldots, x_n])$
les points correspondant à $\mathfrak q'$ et
$\mathfrak q_K'$. D'après le lemme \ref{algebra-lemma-codimension} et ce que nous avons montré
ci-dessus, nous avons
\begin{eqnarray*}
n - \dim_x X & = & \dim_{x'} X' - \dim_x X \\
& = & \text{hauteur}(\mathfrak q') - \text{hauteur}(\mathfrak q) \\
& = & \text{hauteur}(\mathfrak q_K') - \text{hauteur}(\mathfrak q_K) \\
& = & \dim_{x'_K} X'_K - \dim_{x_K} X_K \\
& = & n - \dim_{x_K} X_K
\end{eqnarray*}
et le lemme en résulte.
\end{proof}

\begin{lemma}
\label{algebra-lemma-inequalities-under-field-extension}
Soit $k$ un corps. Soit $S$ une $k$-algèbre de type fini.
Soit $K/k$ une extension de corps. Posons $S_K = K \otimes_k S$.
Soit $\mathfrak q \subset S$ un idéal premier et soit $\mathfrak q_K \subset S_K$
un idéal premier au-dessus de $\mathfrak q$. Alors
$$
\dim (S_K \otimes_S \kappa(\mathfrak q))_{\mathfrak q_K} =
\dim (S_K)_{\mathfrak q_K} - \dim S_\mathfrak q =
\text{trdeg}_k \kappa(\mathfrak q) - \text{trdeg}_K \kappa(\mathfrak q_K)
$$
De plus, étant donné $\mathfrak q$, nous pouvons toujours choisir $\mathfrak q_K$ de sorte
que le nombre ci-dessus soit nul.
\end{lemma}

\begin{proof}
Observons que $S_\mathfrak q \to (S_K)_{\mathfrak q_K}$ est un homomorphisme
local plat d'anneaux locaux noethériens, de fibre spéciale
$(S_K \otimes_S \kappa(\mathfrak q))_{\mathfrak q_K}$. La première
égalité résulte donc du lemme \ref{algebra-lemma-dimension-base-fibre-equals-total}.
La seconde égalité résulte du fait que nous avons
$\dim_x X = \dim_{x_K} X_K$, avec les notations du
lemme \ref{algebra-lemma-dimension-at-a-point-preserved-field-extension},
et que nous avons
$\dim_x X = \dim S_\mathfrak q + \text{trdeg}_k \kappa(\mathfrak q)$
d'après le lemme \ref{algebra-lemma-dimension-at-a-point-finite-type-field},
et de même pour $\dim_{x_K} X_K$.
Si nous choisissons $\mathfrak q_K$ minimal au-dessus de $\mathfrak q S_K$, alors
la dimension de l'anneau fibre sera nulle.
\end{proof}







\section{Dimension des algèbres graduées sur un corps}
\label{algebra-section-dimension-graded}

\noindent
Voici un résultat élémentaire.

\begin{lemma}
\label{algebra-lemma-dimension-graded}
Soit $k$ un corps. Soit $S$ une $k$-algèbre graduée engendrée sur $k$
par un nombre fini d'éléments de degré $1$.
Supposons $S_0 = k$. Soit $P(T) \in \mathbf{Q}[T]$ le polynôme
tel que $\dim(S_d) = P(d)$ pour tout $d \gg 0$. Voir la
proposition \ref{algebra-proposition-graded-hilbert-polynomial}.
Alors
\begin{enumerate}
\item L'idéal non pertinent $S_{+}$ est un idéal maximal $\mathfrak m$.
\item Tout idéal premier minimal de $S$ est un idéal homogène et est contenu
dans $S_{+} = \mathfrak m$.
\item Nous avons $\dim(S) = \deg(P) + 1 = \dim_x\Spec(S)$
(avec la convention que $\deg(0) = -1$),
où $x$ est le point correspondant à l'idéal maximal
$S_{+} = \mathfrak m$.
\item La fonction de Hilbert de l'anneau local $R = S_{\mathfrak m}$
est égale à la fonction de Hilbert de $S$.
\end{enumerate}
\end{lemma}

\begin{proof}
La première assertion est évidente.
La seconde résulte du lemme \ref{algebra-lemma-graded-ring-minimal-prime}.
D'après (2), toute composante irréductible passe par $x$.
Ainsi, nous avons $\dim(S) = \dim_x\Spec(S) = \dim(S_\mathfrak m)$
d'après le lemme \ref{algebra-lemma-dimension-at-a-point-finite-type-over-field}. Puisque
$\mathfrak m^d/\mathfrak m^{d + 1} \cong
\mathfrak m^dS_\mathfrak m/\mathfrak m^{d + 1}S_\mathfrak m$,
nous voyons que la fonction de Hilbert de l'anneau local $S_\mathfrak m$
est égale à la
fonction de Hilbert de $S$, ce qui donne (4). Nous déduisons la dernière égalité
de (3) de la proposition \ref{algebra-proposition-dimension}.
\end{proof}













\section{Platitude générique}
\label{algebra-section-generic-flatness}

\noindent
En substance, cela affirme qu'une algèbre de type fini sur un anneau intègre devient
plate après inversion d'un seul élément de l'anneau intègre.
Il existe plusieurs versions de ce résultat (par ordre croissant de
force).

\begin{lemma}
\label{algebra-lemma-generic-flatness-Noetherian}
Soit $R \to S$ un homomorphisme d'anneaux.
Soit $M$ un $S$-module.
Supposons que
\begin{enumerate}
\item $R$ est noethérien,
\item $R$ est intègre,
\item $R \to S$ est de type fini, et
\item $M$ est un $S$-module de type fini.
\end{enumerate}
Alors il existe un élément non nul $f \in R$ tel que
$M_f$ soit un $R_f$-module libre.
\end{lemma}

\begin{proof}
Soit $K$ le corps des fractions de $R$. Posons $S_K = K \otimes_R S$. C'
est une algèbre de type fini sur $K$. Nous allons raisonner par récurrence sur
$d = \dim(S_K)$ (qui est fini, par exemple par la normalisation de Noether, voir la
section \ref{algebra-section-Noether-normalization}).
Fixons $d \geq 0$.
Supposons que le lemme soit vrai dans tous les cas où $\dim(S_K) < d$.

\medskip\noindent
Supposons donnés $R \to S$ et $M$ comme dans le lemme, avec $\dim(S_K) = d$. D'après le
lemme \ref{algebra-lemma-filter-Noetherian-module},
il existe une filtration
$0 \subset M_1 \subset M_2 \subset \ldots \subset M_n = M$
telle que $M_i/M_{i - 1}$ soit isomorphe à $S/\mathfrak q$
pour un idéal premier $\mathfrak q$ de $S$. Remarquons que
$\dim((S/\mathfrak q)_K) \leq \dim(S_K)$. Remarquons aussi qu'une extension de
modules libres est libre (voir la notion élémentaire \ref{algebra-item-extension-free}).
Nous pouvons donc supposer que $M = S$ et que $S$ est un anneau intègre de type fini sur $R$.

\medskip\noindent
Si $R \to S$ a un noyau non trivial, choisissons un élément non nul $f \in R$ dans
ce noyau. Dans ce cas, $S_f = 0$ et le lemme est vrai. (Il s'agit en réalité
du cas $d = -1$, qui amorce la récurrence.) Nous pouvons donc
supposer que $R \to S$ est une extension de type fini d'anneaux intègres noethériens.

\medskip\noindent
Appliquons le lemme \ref{algebra-lemma-Noether-normalization-over-a-domain}
et remplaçons $R$ par $R_f$ (où $f$ est comme dans le lemme) pour obtenir une
factorisation
$$
R \subset R[y_1, \ldots, y_d] \subset S
$$
où la seconde extension est finie. Choisissons $z_1, \ldots, z_r \in S$ qui
forment une base du corps des fractions de $S$ sur le corps des fractions de
$R[y_1, \ldots, y_d]$. Cela donne une suite exacte courte
$$
0 \to
R[y_1, \ldots, y_d]^{\oplus r} \xrightarrow{(z_1, \ldots, z_r)}
S \to N \to 0
$$
Par construction, $N$ est un $R[y_1, \ldots, y_d]$-module fini dont le
support ne contient pas le point générique $(0)$ de
$\Spec(R[y_1, \ldots, y_d])$. D'après le
lemme \ref{algebra-lemma-support-closed},
il existe un élément non nul $g \in R[y_1, \ldots, y_d]$ tel que
$g$ annule $N$ ; nous pouvons donc voir $N$ comme un module fini sur
$S' = R[y_1, \ldots, y_d]/(g)$. Puisque $\dim(S'_K) < d$, il existe par récurrence
un élément non nul $f \in R$ tel que $N_f$ soit un
$R_f$-module libre. Puisque
$(R[y_1, \ldots, y_d])_f \cong R_f[y_1, \ldots, y_d]$ est également libre,
nous concluons par le fait déjà mentionné qu'une extension
de modules libres est libre.
\end{proof}

\begin{lemma}
\label{algebra-lemma-generic-flatness-finitely-presented}
\begin{slogan}
Liberté générique.
\end{slogan}
Soit $R \to S$ un homomorphisme d'anneaux.
Soit $M$ un $S$-module.
Supposons que
\begin{enumerate}
\item $R$ est intègre,
\item $R \to S$ est de présentation finie, et
\item $M$ est un $S$-module de présentation finie.
\end{enumerate}
Alors il existe un élément non nul $f \in R$ tel que
$M_f$ soit un $R_f$-module libre.
\end{lemma}

\begin{proof}
Écrivons $S = R[x_1, \ldots, x_n]/(g_1, \ldots, g_m)$.
Pour $g \in R[x_1, \ldots, x_n]$, notons $\overline{g}$ son image dans $S$.
Nous pouvons écrire $M = S^{\oplus t}/\sum Sn_i$ pour certains $n_i \in S^{\oplus t}$.
Écrivons $n_i = (\overline{g}_{i1}, \ldots, \overline{g}_{it})$ pour certains
$g_{ij} \in R[x_1, \ldots, x_n]$. Soit $R_0 \subset R$ le sous-anneau
engendré par tous les coefficients de tous les éléments
$g_i, g_{ij} \in R[x_1, \ldots, x_n]$. Définissons
$S_0 = R_0[x_1, \ldots, x_n]/(g_1, \ldots, g_m)$.
Définissons $M_0 = S_0^{\oplus t}/\sum S_0n_i$.
Alors $R_0$ est un anneau intègre de type fini sur $\mathbf{Z}$, et donc
noethérien (voir le
lemme \ref{algebra-lemma-Noetherian-permanence}).
De plus, via l'injection $R_0 \to R$, nous avons $S \cong R \otimes_{R_0} S_0$
et $M \cong R \otimes_{R_0} M_0$. En appliquant le
lemme \ref{algebra-lemma-generic-flatness-Noetherian},
nous obtenons un élément non nul $f \in R_0$ tel que $(M_0)_f$ soit un
$(R_0)_f$-module libre. Ainsi, $M_f = R_f \otimes_{(R_0)_f} (M_0)_f$
est un $R_f$-module libre.
\end{proof}

\begin{lemma}
\label{algebra-lemma-generic-flatness}
Soit $R \to S$ un homomorphisme d'anneaux.
Soit $M$ un $S$-module.
Supposons que
\begin{enumerate}
\item $R$ est intègre,
\item $R \to S$ est de type fini, et
\item $M$ est un $S$-module de type fini.
\end{enumerate}
Alors il existe un élément non nul $f \in R$ tel que
\begin{enumerate}
\item[(a)] $M_f$ et $S_f$ soient libres comme $R_f$-modules, et
\item[(b)] $S_f$ soit une $R_f$-algèbre de présentation finie et $M_f$ un
$S_f$-module de présentation finie.
\end{enumerate}
\end{lemma}

\begin{proof}
Nous démontrons d'abord le lemme pour $S = R[x_1, \ldots, x_n]$, puis
nous en déduisons le résultat général.

\medskip\noindent
Supposons $S = R[x_1, \ldots, x_n]$.
Choisissons des éléments $m_1, \ldots, m_t$ qui engendrent $M$. Cela donne
une suite exacte courte
$$
0 \to N \to S^{\oplus t} \xrightarrow{(m_1, \ldots, m_t)} M \to 0.
$$
Notons $K$ le corps des fractions de $R$. Notons
$S_K = K \otimes_R S = K[x_1, \ldots, x_n]$, et de même
$N_K = K \otimes_R N$, $M_K = K \otimes_R M$.
Comme $R \to K$ est plat, la suite reste exacte après tensorisation par $K$.
Comme $S_K = K[x_1, \ldots, x_n]$ est un anneau noethérien (voir le
lemme \ref{algebra-lemma-Noetherian-permanence}),
nous pouvons trouver un nombre fini d'éléments $n'_1, \ldots, n'_s \in N_K$
qui l'engendrent. Choisissons $n_1, \ldots, n_r \in N$ tels que
$n'_i = \sum a_{ij}n_j$ pour certains $a_{ij} \in K$. Posons
$$
M' = S^{\oplus t}/\sum\nolimits_{i = 1, \ldots, r} Sn_i
$$
Par construction, $M'$ est un $S$-module de présentation finie, et
il existe une surjection $M' \to M$ qui induit un isomorphisme
$M'_K \cong M_K$. Nous pouvons appliquer le
lemme \ref{algebra-lemma-generic-flatness-finitely-presented}
à $R \to S$ et $M'$, et nous trouvons un élément $f \in R$ tel que $M'_f$
soit un $R_f$-module libre. Ainsi, $M'_f \to M_f$ est une surjection de
modules sur l'anneau intègre $R_f$, dont la source est un module libre,
et qui devient un isomorphisme après tensorisation par $K$.
Elle est donc injective, puisque $M'_f \subset M'_K$, car $M'_f$
est libre et $R_f$ est un anneau intègre de corps des fractions $K$.
Ainsi, $M'_f \to M_f$ est un isomorphisme, et le résultat est démontré.

\medskip\noindent
Dans le cas général, choisissons une surjection $R[x_1, \ldots, x_n] \to S$.
Considérons $S$ et $M$ comme des modules finis sur $R[x_1, \ldots, x_n]$.
D'après le cas particulier démontré ci-dessus, il existe un élément non nul $f \in R$
tel que $S_f$ et $M_f$ soient tous deux libres comme $R_f$-modules et de présentation
finie comme $R_f[x_1, \ldots, x_n]$-modules. Cela implique clairement que
$S_f$ est une $R_f$-algèbre de présentation finie et que $M_f$ est un
$S_f$-module de présentation finie.
\end{proof}

\noindent
Soit $R \to S$ un homomorphisme d'anneaux. Soit $M$ un $S$-module. Considérons
la condition suivante sur un élément $f \in R$ :
\begin{equation}
\label{algebra-equation-flat-and-finitely-presented}
\left\{
\begin{matrix}
S_f & \text{est de présentation finie sur }R_f\\
M_f & \text{est de présentation finie comme }S_f\text{-module}\\
S_f, M_f & \text{sont libres comme }R_f\text{-modules}
\end{matrix}
\right.
\end{equation}
Nous définissons
\begin{equation}
\label{algebra-equation-good-locus}
U(R \to S, M)
=
\bigcup\nolimits_{f \in
R\text{ tel que }(\ref{algebra-equation-flat-and-finitely-presented})}
D(f)
\end{equation}
qui est un ouvert de $\Spec(R)$.

\begin{lemma}
\label{algebra-lemma-generic-flatness-locus-extension}
Soit $R \to S$ un homomorphisme d'anneaux.
Soit $0 \to M_1 \to M_2 \to M_3 \to 0$ une suite exacte courte
de $S$-modules.
Alors
$$
U(R \to S, M_1) \cap U(R \to S, M_3) \subset U(R \to S, M_2).
$$
\end{lemma}

\begin{proof}
Soit $u \in U(R \to S, M_1) \cap U(R \to S, M_3)$. Choisissons
$f_1, f_3 \in R$ tels que $u \in D(f_1)$, $u \in D(f_3)$ et
que (\ref{algebra-equation-flat-and-finitely-presented}) soit vérifiée pour
$f_1$ et $M_1$, ainsi que pour $f_3$ et $M_3$. Posons alors $f = f_1f_3$.
Nous avons $u \in D(f)$, et (\ref{algebra-equation-flat-and-finitely-presented})
est vérifiée pour $f$ et à la fois $M_1$ et $M_3$. Une extension de modules libres
est libre, et une extension de modules de présentation finie est de présentation finie
(lemme \ref{algebra-lemma-extension}). Nous voyons donc que
(\ref{algebra-equation-flat-and-finitely-presented}) est vérifiée pour $f$ et $M_2$.
Ainsi, $u \in U(R \to S, M_2)$, ce qui termine la démonstration.
\end{proof}

\begin{lemma}
\label{algebra-lemma-generic-flatness-locus-localize}
Soit $R \to S$ un homomorphisme d'anneaux.
Soit $M$ un $S$-module.
Soit $f \in R$.
Avec l'identification $\Spec(R_f) = D(f)$, nous avons
$U(R_f \to S_f, M_f) = D(f) \cap U(R \to S, M)$.
\end{lemma}

\begin{proof}
Supposons que $u \in U(R_f \to S_f, M_f)$. Il existe alors un
élément $g \in R_f$ tel que $u \in D(g)$ et que
(\ref{algebra-equation-flat-and-finitely-presented})
soit vérifiée pour le couple $((R_f)_g \to (S_f)_g, (M_f)_g)$.
Écrivons $g = a/f^n$ pour un certain $a \in R$. Posons $h = af$.
Alors $R_h = (R_f)_g$, $S_h = (S_f)_g$ et $M_h = (M_f)_g$.
De plus, $u \in D(h)$. Ainsi, $u \in U(R \to S, M)$.
Réciproquement, supposons que $u \in D(f) \cap U(R \to S, M)$.
Il existe alors un élément $g \in R$ tel que $u \in D(g)$ et que
(\ref{algebra-equation-flat-and-finitely-presented})
soit vérifiée pour le couple $(R_g \to S_g, M_g)$.
Il est alors clair que (\ref{algebra-equation-flat-and-finitely-presented})
est également vérifiée pour le couple
$(R_{fg} \to S_{fg}, M_{fg}) = ((R_f)_g \to (S_f)_g, (M_f)_g)$.
Ainsi, $u \in U(R_f \to S_f, M_f)$, ce qui termine la démonstration.
\end{proof}

\begin{lemma}
\label{algebra-lemma-generic-flatness-locus-reduce}
Soit $R \to S$ un homomorphisme d'anneaux.
Soit $M$ un $S$-module.
Soit $U \subset \Spec(R)$ un ouvert dense.
Supposons qu'il existe un recouvrement $U = \bigcup_{i \in I} D(f_i)$ par des
ouverts tel que $U(R_{f_i} \to S_{f_i}, M_{f_i})$ soit dense dans
$D(f_i)$ pour tout $i \in I$. Alors $U(R \to S, M)$ est dense dans
$\Spec(R)$.
\end{lemma}

\begin{proof}
Au vu du
lemme \ref{algebra-lemma-generic-flatness-locus-localize},
il s'agit d'une assertion purement topologique. En effet, d'après ce lemme,
nous voyons que $U(R \to S, M) \cap D(f_i)$ est dense dans $D(f_i)$
pour tout $i \in I$. D'après
Topologie, lemme \ref{topology-lemma-nowhere-dense-local},
nous voyons que $U(R \to S, M) \cap U$ est dense dans $U$.
Puisque $U$ est dense dans $\Spec(R)$, nous concluons que $U(R \to S, M)$
est dense dans $\Spec(R)$.
\end{proof}

\begin{lemma}
\label{algebra-lemma-generic-flatness-reduced}
Soit $R \to S$ un homomorphisme d'anneaux. Soit $M$ un $S$-module.
Supposons que
\begin{enumerate}
\item $R \to S$ est de type fini,
\item $M$ est un $S$-module fini, et
\item $R$ est réduit.
\end{enumerate}
Alors il existe une partie $U \subset \Spec(R)$ telle que
\begin{enumerate}
\item $U$ soit ouverte et dense dans $\Spec(R)$,
\item pour tout $u \in U$, il existe $f \in R$
tel que $u \in D(f) \subset U$ et que nous ayons
\begin{enumerate}
\item $M_f$ et $S_f$ sont libres sur $R_f$,
\item $S_f$ est une $R_f$-algèbre de présentation finie, et
\item $M_f$ est un $S_f$-module de présentation finie.
\end{enumerate}
\end{enumerate}
\end{lemma}

\begin{proof}
Remarquons que le lemme équivaut à dire que l'ouvert
$U(R \to S, M)$, voir l'équation (\ref{algebra-equation-good-locus}), est dense
dans $\Spec(R)$. Nous démontrons d'abord le lemme pour
$S = R[x_1, \ldots, x_n]$, puis
nous en déduisons le résultat général.

\medskip\noindent
Démonstration du cas $S = R[x_1, \ldots, x_n]$ et où $M$ est un module fini quelconque
sur $S$. Remarquons que, dans ce cas, $S_f = R_f[x_1, \ldots, x_n]$
est libre et de présentation finie sur $R_f$ ; nous n'avons donc pas
à nous préoccuper des conditions qui concernent $S$,
mais seulement de celles qui concernent $M$. Nous raisonnons par récurrence sur $n$.

\medskip\noindent
Il existe une filtration finie
$$
0 \subset M_1 \subset M_2 \subset \ldots \subset M_t = M
$$
telle que $M_i/M_{i - 1} \cong S/J_i$ pour un certain idéal $J_i \subset S$, voir le
lemme \ref{algebra-lemma-trivial-filter-finite-module}. Puisqu'
une intersection finie d'ouverts denses est un ouvert dense,
il résulte du
lemme \ref{algebra-lemma-generic-flatness-locus-extension}
qu'il suffit de démontrer le lemme pour chacun des modules $R/J_i$.
Nous pouvons donc supposer que $M = S/J$ pour un certain idéal $J$ de
$S = R[x_1, \ldots, x_n]$.

\medskip\noindent
Soit $I \subset R$ l'idéal engendré par les coefficients des
éléments de $J$. Soit $U_1 = \Spec(R) \setminus V(I)$ et
soit
$$
U_2 = \Spec(R) \setminus \overline{U_1}.
$$
Il est alors clair que $U = U_1 \cup U_2$ est dense dans $\Spec(R)$.
Soit $f \in R$ un élément tel que soit (a) $D(f) \subset U_1$, soit
(b) $D(f) \subset U_2$. Si, pour tout $f$ de ce type,
le lemme est vrai pour le couple $(R_f \to R_f[x_1, \ldots, x_n], M_f)$,
alors le
lemme \ref{algebra-lemma-generic-flatness-locus-reduce}
montre que $U(R \to S, M)$ est dense dans $\Spec(R)$.
Nous pouvons donc supposer soit (a) $I = R$, soit (b) $V(I) = \Spec(R)$.

\medskip\noindent
Dans le cas (b), nous avons en fait $I = 0$, puisque $R$ est réduit ! Ainsi, $J = 0$,
$M = S$, et le lemme est vrai dans ce cas.

\medskip\noindent
Dans le cas (a), il faut travailler un peu plus. Remarquons que tout élément
de $I$ est en fait le coefficient d'un monôme d'un élément de $J$, car
l'ensemble des coefficients des éléments de $J$ forme un idéal (détails omis).
Nous trouvons donc un élément
$$
g = \sum\nolimits_{K \in E} a_K x^K \in J
$$
où $E$ est un ensemble fini de multi-indices $K = (k_1, \ldots, k_n)$,
avec au moins un coefficient $a_{K_0}$ qui est une unité dans $R$. En fait, nous pouvons
en trouver un dont un coefficient est égal à $1$, puisque $1 \in I$ dans le cas (a).
Soit $m = \#\{K \in E \mid a_K \text{ n'est pas une unité}\}$.
Remarquons que $0 \leq m \leq \# E - 1$.
Nous raisonnons par récurrence sur $m$.

\medskip\noindent
Le cas $m = 0$. Dans ce cas, tous les coefficients $a_K$, $K \in E$,
de $g$ sont des unités, et $E \not = \emptyset$.
Si $E = \{K_0\}$ est un singleton et $K_0 = (0, \ldots, 0)$, alors $g$
est une unité et $J = S$, de sorte que le résultat est certainement vrai. (Cela se produit
en particulier lorsque $n = 0$, et fournit le cas initial de la récurrence
sur $n$.) Si l'on n'a pas $E = \{(0, \ldots, 0)\}$, alors au moins un $K$ n'est pas
égal à $(0, \ldots, 0)$, c'est-à-dire $g \not \in R$. Nous employons alors
l'astuce usuelle de la normalisation de Noether. Plus précisément, nous considérons
$$
G(y_1, \ldots, y_n) =
g(y_1 + y_n^{e_1}, y_2 + y_n^{e_2}, \ldots, y_{n - 1} + y_n^{e_{n - 1}}, y_n)
$$
avec $0 \ll e_{n -1} \ll e_{n - 2} \ll \ldots \ll e_1$. D'après le
lemme \ref{algebra-lemma-helper-polynomial},
le polynôme $G(y_1, \ldots, y_n)$ en $y_n$
est de la forme
$$
a_K y_n^{k_n + \sum_{i = 1, \ldots, n - 1} e_i k_i} +
\text{termes de degré inférieur en }y_n
$$
Puisque $a_K$ est une unité, nous concluons que $M = R[x_1, \ldots, x_n]/J$ est
fini sur $R[y_1, \ldots, y_{n - 1}]$. Ainsi,
$U(R \to R[x_1, \ldots, x_n], M) = U(R \to R[y_1, \ldots, y_{n - 1}], M)$,
et nous concluons par récurrence sur $n$.

\medskip\noindent
Le cas $m > 0$. Choisissons un multi-indice $K \in E$ tel que $a_K$ ne soit
pas une unité. Comme précédemment, posons
$U_1 = \Spec(R_{a_K}) = \Spec(R) \setminus V(a_K)$
et posons
$$
U_2 = \Spec(R) \setminus \overline{U_1}.
$$
Il est alors clair que $U = U_1 \cup U_2$ est dense dans $\Spec(R)$.
Soit $f \in R$ un élément tel que soit (a) $D(f) \subset U_1$, soit
(b) $D(f) \subset U_2$. Si, pour tout $f$ de ce type,
le lemme est vrai pour le couple $(R_f \to R_f[x_1, \ldots, x_n], M_f)$,
alors le
lemme \ref{algebra-lemma-generic-flatness-locus-reduce}
montre que $U(R \to S, M)$ est dense dans $\Spec(R)$.
Nous pouvons donc supposer soit (a) $a_KR = R$, soit (b) $V(a_K) = \Spec(R)$.
Dans le cas (a), le nombre $m$ diminue, puisque $a_K$ est devenu une unité.
Dans le cas (b), puisque $R$ est réduit, nous concluons que $a_K = 0$.
Ainsi, l'ensemble $E$ diminue, et le nombre $m$ diminue lui aussi.
Dans les deux cas, nous concluons par récurrence sur $m$.

\medskip\noindent
Nous avons maintenant démontré le lemme dans le cas $S = R[x_1, \ldots, x_n]$.
Supposons que $(R \to S, M)$ soit un couple quelconque satisfaisant les conditions du
lemme. Choisissons une surjection $R[x_1, \ldots, x_n] \to S$. Remarquons que,
avec les notations introduites dans (\ref{algebra-equation-good-locus}), nous avons
$$
U(R \to S, M) =
U(R \to R[x_1, \ldots, x_n], S)
\cap
U(R \to R[x_1, \ldots, x_n], M)
$$
Puisque nous venons de démontrer que les deux ouverts du membre de droite sont denses,
l'ouvert du membre de gauche est lui aussi dense.
\end{proof}






\section{Autour de Krull-Akizuki}
\label{algebra-section-krull-akizuki}

\noindent
Une application du théorème de Krull-Akizuki consiste à montrer qu'il existe de nombreux
anneaux de valuation discrète. Plus généralement, nous montrons dans cette section
comment construire des anneaux de valuation discrète qui dominent des anneaux
locaux noethériens.

\medskip\noindent
Nous montrons d'abord comment dominer un anneau local noethérien intègre
par un anneau local noethérien intègre de dimension $1$, en éclatant
l'idéal maximal.

\begin{lemma}
\label{algebra-lemma-dominate-by-dimension-1}
Soit $R$ un anneau local noethérien intègre de corps des fractions $K$.
Supposons que $R$ ne soit pas un corps.
Alors il existe un anneau intermédiaire $R \subset R' \subset K$ qui vérifie les conditions suivantes :
\begin{enumerate}
\item $R'$ est un anneau local noethérien de dimension $1$,
\item $R \to R'$ est un homomorphisme local d'anneaux, c'est-à-dire que $R'$ domine $R$, et
\item $R \to R'$ est essentiellement de type fini.
\end{enumerate}
\end{lemma}

\begin{proof}
Choisissons un anneau de valuation $A \subset K$ qui domine $R$ (il en
existe un d'après le lemme \ref{algebra-lemma-dominate}).
Notons $v$ la valuation correspondante.
Soit $x_1, \ldots, x_r$ un système minimal de générateurs de
l'idéal maximal $\mathfrak m$ de $R$. Nous pouvons supposer, et nous le faisons, que
$v(x_r) = \min\{v(x_1), \ldots, v(x_r)\}$. Considérons l'anneau
$$
S = R[x_1/x_r, x_2/x_r, \ldots, x_{r - 1}/x_r] \subset K.
$$
Remarquons que $\mathfrak mS = x_rS$ est un idéal principal.
Remarquons que $S \subset A$ et que $v(x_r) > 0$ ; nous voyons donc
que $x_rS \not = S$. Choisissons un idéal premier minimal $\mathfrak q$
au-dessus de $x_rS$. Alors $\text{hauteur}(\mathfrak q) = 1$ d'après le
lemme \ref{algebra-lemma-minimal-over-1},
et $\mathfrak q$ est au-dessus de $\mathfrak m$. Nous
voyons ainsi que $R' = S_{\mathfrak q}$ est une solution.
\end{proof}

\begin{lemma}[Koll\'ar]
\label{algebra-lemma-hart-serre-loc-thm}
\begin{reference}
Ceci est tiré d'un article à paraître de
J\'anos Koll\'ar intitulé « Variants of normality for Noetherian schemes ».
\end{reference}
Soit $(R, \mathfrak m)$ un anneau local noethérien.
Alors exactement l'une des possibilités suivantes est réalisée :
\begin{enumerate}
\item $(R, \mathfrak m)$ est artinien,
\item $(R, \mathfrak m)$ est régulier de dimension $1$,
\item $\text{profondeur}(R) \geq 2$, ou
\item il existe un homomorphisme fini d'anneaux $R \to R'$ qui n'est pas
un isomorphisme, dont le noyau et le conoyau sont annulés par une puissance
de $\mathfrak m$, tel que $\mathfrak m$ ne soit pas un idéal premier associé
de $R'$ et que $R' \not = 0$.
\end{enumerate}
\end{lemma}

\begin{proof}
Observons que $(R, \mathfrak m)$ n'est pas artinien si et seulement si
$V(\mathfrak m) \subset \Spec(R)$ est rare. Voir la
proposition \ref{algebra-proposition-dimension-zero-ring}. Nous le supposons désormais.

\medskip\noindent
Soit $J \subset R$ le plus grand idéal annulé par une puissance de $\mathfrak m$.
Si $J \not = 0$, alors $R \to R/J$ montre que $(R, \mathfrak m)$
relève du cas (4).

\medskip\noindent
Sinon, $J = 0$. En particulier, $\mathfrak m$ n'est pas un idéal premier associé
de $R$, et le
lemme \ref{algebra-lemma-ideal-nonzerodivisor} montre qu'il existe un non-diviseur de zéro $x \in \mathfrak m$. Si $\mathfrak m$
n'est pas un idéal premier associé de $R/xR$, alors $\text{profondeur}(R) \geq 2$
d'après le même lemme. Il ne reste donc que le cas où il existe
$y \in R$, $y \not \in xR$, tel que $y \mathfrak m \subset xR$.

\medskip\noindent
Si $y \mathfrak m \subset x \mathfrak m$, nous pouvons considérer
l'application $\varphi : \mathfrak m \to \mathfrak m$, $f \mapsto yf/x$
(bien définie puisque $x$ est un non-diviseur de zéro). Par l'astuce du déterminant
du lemme \ref{algebra-lemma-charpoly-module}, il existe un polynôme unitaire
$P$ à coefficients dans $R$ tel que $P(\varphi) = 0$.
Nous concluons que $P(y/x) = 0$ dans $R_x$.
Soit $R' \subset R_x$ l'anneau engendré par
$R$ et $y/x$. Alors $R \subset R'$, et $R'/R$ est un $R$-module de type fini
annulé par une puissance de $\mathfrak m$. Ainsi, $R$ relève du cas (4).

\medskip\noindent
Sinon, il existe $t \in \mathfrak m$ tel que $y t = u x$
pour une certaine unité $u$ de $R$. Après avoir remplacé $t$ par $u^{-1}t$,
nous obtenons $yt = x$. En particulier, $y$ est un non-diviseur de zéro.
Pour tout $t' \in \mathfrak m$, nous avons $y t' = x s$ pour un certain $s \in R$.
Ainsi, $y (t' - s t ) = x s - x s = 0$.
Comme $y$ n'est pas un diviseur de zéro, cela implique $t' = ts$, et donc
$\mathfrak m = (t)$. Ainsi, $(R, \mathfrak m)$ est régulier de dimension 1.

\medskip\noindent
L'argument mené jusqu'ici montre que tout $R$ relève de l'un des 4 cas.
Pour terminer, nous devons montrer que ces propriétés
s'excluent deux à deux. Nous montrerons seulement que (2) et (3)
excluent chacune (4) ; les autres cas sont laissés
au lecteur.

\medskip\noindent
Supposons que $R$ soit régulier de dimension $1$ et que $R \to R'$ soit un homomorphisme fini
d'anneaux dont le noyau et le conoyau sont annulés par une puissance de $\mathfrak m$,
et que $\mathfrak m$ ne soit pas un idéal premier associé de $R'$. Alors $R'$
est un $R$-module de type fini de profondeur au moins $1$, et est donc libre d'après le
lemme \ref{algebra-lemma-regular-mcm-free}. Puisque $R \to R'$ est un isomorphisme
au point générique, nous concluons que $R'$ doit être libre de rang $1$
comme $R$-module. Alors $R' \to \text{End}_R(R') \cong R$ est l'inverse
de l'application $R \to R'$. Ainsi, le cas (4) est impossible.

\medskip\noindent
Supposons que $R$ soit de profondeur $\geq 2$ et que $R \to R'$ soit un homomorphisme fini
d'anneaux dont le noyau et le conoyau sont annulés par une puissance
de $\mathfrak m$, et que $\mathfrak m$ ne soit pas un idéal premier associé de $R'$.
Alors $R \to R'$ est nécessairement injectif (car le noyau serait
à la fois de profondeur $0$ et de profondeur $\geq 1$), et le conoyau est de profondeur
au moins $1$ (d'après le lemme \ref{algebra-lemma-depth-in-ses}
et le fait que $R'$ est de profondeur $\geq 1$ par hypothèse) ;
il ne peut donc être supporté sur
$\{\mathfrak m\}$, à moins d'être lui aussi $0$. Ainsi, le cas (4) est impossible.
\end{proof}

\begin{lemma}
\label{algebra-lemma-nonregular-dimension-one}
Soit $R$ un anneau local d'idéal maximal $\mathfrak m$.
Supposons que $R$ soit noethérien, de dimension $1$, et que
$\dim(\mathfrak m/\mathfrak m^2) > 1$. Alors il existe
un homomorphisme d'anneaux $R \to R'$ tel que
\begin{enumerate}
\item $R \to R'$ soit fini,
\item $R \to R'$ ne soit pas un isomorphisme,
\item le noyau et le conoyau de $R \to R'$ soient annulés par
une puissance de $\mathfrak m$, et
\item $\mathfrak m$ ne soit pas un idéal premier associé de $R'$.
\end{enumerate}
\end{lemma}

\begin{proof}
Cela résulte du lemme \ref{algebra-lemma-hart-serre-loc-thm}
et du fait que $R$ n'est ni artinien, ni régulier, et
n'est pas de profondeur $\geq 2$ (ce dernier point parce que la
profondeur ne dépasse pas la dimension, d'après le
lemme \ref{algebra-lemma-bound-depth}).
\end{proof}

\begin{example}
\label{algebra-example-nonreduced}
Considérons l'anneau local noethérien
$$
R = k[[x, y]]/(y^2)
$$
Il est de dimension 1 et de Cohen-Macaulay.
Un exemple d'extension comme dans le
lemme \ref{algebra-lemma-nonregular-dimension-one}
est l'extension
$$
k[[x, y]]/(y^2) \subset k[[x, z]]/(z^2), \ \ y \mapsto xz
$$
autrement dit, elle s'obtient en adjoignant $y/x$ à $R$.
Répéter la construction $n > 1$ fois revient à
adjoindre l'élément $y/x^n$.
\end{example}

\begin{example}
\label{algebra-example-bad-dvr-char-p}
Soit $k$ un corps de caractéristique $p > 0$ tel que $k$
soit de degré infini sur son sous-corps $k^p$ des puissances $p$-ièmes.
Par exemple, $k = \mathbf{F}_p(t_1, t_2, t_3, \ldots)$.
Considérons l'anneau
$$
A =
\left\{
\sum a_i x^i \in k[[x]] \text{ tel que }
[k^p(a_0, a_1, a_2, \ldots) : k^p] < \infty
\right\}
$$
Alors $A$ est un anneau de valuation discrète, et son complété est
$A^\wedge = k[[x]]$. Remarquons que l'inclusion $A \subset k[[x]]$
induit une extension infinie et purement inséparable des corps des fractions.
Choisissons un $f \in k[[x]]$, $f \not \in A$. Soit $R = A[f] \subset k[[x]]$.
Alors $R$ est un anneau local noethérien intègre de dimension $1$ dont le
complété $R^\wedge$ est non réduit (à méditer !).
\end{example}

\begin{remark}
\label{algebra-remark-resolution-dim-1}
Supposons que $R$ soit un anneau semi-local noethérien intègre de dimension $1$.
S'il existe un idéal maximal $\mathfrak m \subset R$ tel que
$R_{\mathfrak m}$ ne soit pas régulier, nous pouvons appliquer le
lemme \ref{algebra-lemma-nonregular-dimension-one} à $(R, \mathfrak m)$
pour obtenir une extension finie d'anneaux $R \subset R_1$.
(Par exemple, on peut faire en sorte que $\Spec(R_1) \to \Spec(R)$
soit l'éclatement de $\Spec(R)$ le long de l'idéal $\mathfrak m$.)
Bien sûr, $R_1$ est un anneau semi-local noethérien intègre
de dimension $1$ ayant le même corps des fractions que $R$. Si $R_1$ n'est pas un
anneau semi-local régulier, nous pouvons répéter la construction pour
obtenir $R_1 \subset R_2$. Nous obtenons ainsi une suite
$$
R \subset R_1 \subset R_2 \subset R_3 \subset \ldots
$$
d'extensions finies d'anneaux, qui peut s'arrêter si $R_n$ est régulier pour
un certain $n$. La résolution des singularités serait l'affirmation
que $R_n$ finit effectivement par être régulier. En réalité, ce n'est pas
le cas. En effet, il existe en caractéristique $0$ un
anneau local noethérien intègre $A$ de dimension $1$ dont le complété est non réduit,
voir \cite[Proposition 3.1]{Ferrand-Raynaud} ou la
section \ref{examples-section-local-completion-nonreduced} du chapitre Exemples.
Pour un exemple en caractéristique $p > 0$, voir
l'exemple \ref{algebra-example-bad-dvr-char-p}.
Comme la formation de l'éclatement est compatible avec la complétion, il
est facile de voir que la suite ne se stabilise jamais.
Voir \cite{Bennett} pour une discussion (principalement en caractéristique positive).
En revanche, si le complété de $R$ en chacun de ses idéaux maximaux
est réduit, alors le procédé s'arrête (insérer ici une référence
future).
\end{remark}

\begin{lemma}
\label{algebra-lemma-characterize-dvr}
Soit $A$ un anneau. Les propriétés suivantes sont équivalentes.
\begin{enumerate}
\item L'anneau $A$ est un anneau de valuation discrète.
\item L'anneau $A$ est un anneau de valuation noethérien qui n'est pas un corps.
\item L'anneau $A$ est un anneau local régulier de dimension $1$.
\item L'anneau $A$ est un anneau local noethérien intègre dont l'idéal maximal
$\mathfrak m$ est engendré par un seul élément non nul.
\item L'anneau $A$ est un anneau local noethérien intègre et normal de dimension $1$.
\end{enumerate}
Dans ce cas, si $\pi$ est un générateur de l'idéal maximal de
$A$, alors tout élément non nul de $A$ s'écrit de manière unique
$u\pi^n$, où $u \in A$ est une unité.
\end{lemma}

\begin{proof}
L'équivalence de (1) et (2) est le
lemme \ref{algebra-lemma-valuation-ring-Noetherian-discrete}.
De plus, dans la preuve du lemme \ref{algebra-lemma-valuation-ring-Noetherian-discrete},
nous avons vu que si $A$ est un anneau de valuation discrète, alors $A$ est un anneau principal, d'où (3).
Remarquons qu'un anneau local régulier est intègre (voir le
lemme \ref{algebra-lemma-regular-domain}). Cela étant, l'équivalence de (3) et (4)
résulte de la théorie de la dimension, voir la section \ref{algebra-section-dimension}.

\medskip\noindent
Supposons (3) et soit $\pi$ un générateur de l'idéal maximal $\mathfrak m$.
Pour tout $n \geq 0$, nous avons
$\dim_{A/\mathfrak m} \mathfrak m^n/\mathfrak m^{n + 1} = 1$,
car il est engendré par $\pi^n$ (et ne peut être nul).
En particulier, $\mathfrak m^n = (\pi^n)$, et
l'anneau gradué $\bigoplus \mathfrak m^n/\mathfrak m^{n + 1}$
est isomorphe à l'anneau de polynômes $A/\mathfrak m[T]$.
Pour $x \in A \setminus \{0\}$, posons
$v(x) = \max\{n \mid x \in \mathfrak m^n\}$.
Autrement dit, $x = u \pi^{v(x)}$ avec $u \in A^*$.
D'après les remarques ci-dessus, nous avons $v(xy) = v(x) + v(y)$
pour tous $x, y \in A \setminus \{0\}$. Nous l'étendons au corps des fractions
$K$ de $A$ en posant $v(a/b) = v(a) - v(b)$ (ce qui est bien défini par la multiplicativité
démontrée ci-dessus). Il est alors clair que $A$ est l'ensemble des éléments de
$K$ de valuation $\geq 0$. Nous voyons donc que $A$ est un
anneau de valuation d'après le lemme \ref{algebra-lemma-valuation-valuation-ring}.

\medskip\noindent
Un anneau de valuation est un anneau intègre normal d'après le lemme \ref{algebra-lemma-valuation-ring-normal}.
Nous voyons donc que les conditions équivalentes (1) -- (3) impliquent
(5). Supposons que (5) soit vérifiée et, en vue d'une contradiction, que
$\mathfrak m$ ne puisse pas être engendré par exactement $1$ élément.
Alors le lemme \ref{algebra-lemma-nonregular-dimension-one} implique qu'il existe un homomorphisme fini
d'anneaux $A \to A'$ qui est un isomorphisme après inversion
de tout élément non nul de $\mathfrak m$, mais n'est pas un isomorphisme.
En particulier, nous pouvons identifier $A'$ à un sous-anneau du corps des fractions de $A$.
Puisque $A \to A'$ est fini, il est entier (voir le
lemme \ref{algebra-lemma-finite-is-integral}).
Puisque $A$ est normal, nous obtenons $A = A'$, contradiction.
\end{proof}

\begin{definition}
\label{algebra-definition-uniformizer}
Soit $A$ un anneau de valuation discrète. Une {\it uniformisante} est un élément
$\pi \in A$ qui engendre l'idéal maximal de $A$.
\end{definition}

\noindent
D'après le lemme \ref{algebra-lemma-characterize-dvr}, deux uniformisantes quelconques d'un anneau
de valuation discrète sont associées.

\begin{lemma}
\label{algebra-lemma-finite-length}
Soit $R$ un anneau intègre de corps des fractions $K$.
Soit $M$ un $R$-sous-module de $K^{\oplus r}$.
Supposons que $R$ soit local noethérien de dimension $1$.
Pour tout élément non nul $x \in R$, nous avons $\text{longueur}_R(R/xR) < \infty$
et
$$
\text{longueur}_R(M/xM) \leq r \cdot \text{longueur}_R(R/xR).
$$
\end{lemma}

\begin{proof}
Si $x$ est une unité, le résultat est vrai. Nous pouvons donc supposer que
$x \in \mathfrak m$, l'idéal maximal de $R$. Puisque $x$ n'est pas
nul et que $R$ est intègre, nous avons $\dim(R/xR) = 0$, et donc
$R/xR$ est de longueur finie. Considérons $M \subset K^{\oplus r}$ comme
dans le lemme. Nous pouvons supposer que les éléments de $M$ engendrent
$K^{\oplus r}$ comme $K$-espace vectoriel, après avoir remplacé $K^{\oplus r}$
par un sous-espace plus petit si nécessaire.

\medskip\noindent
Supposons d'abord que $M$ soit un $R$-module de type fini. Dans ce cas, nous pouvons chasser les
dénominateurs et supposer $M \subset R^{\oplus r}$. Puisque
$M$ engendre $K^{\oplus r}$ comme espace vectoriel, nous voyons que
$R^{\oplus r}/M$ est de longueur finie. En particulier, il existe
un entier $c \geq 0$ tel que $x^cR^{\oplus r} \subset M$.
Remarquons que $M \supset xM \supset x^2M \supset \ldots$ est une suite de
modules dont les quotients successifs sont tous isomorphes à $M/xM$. Ainsi,
nous voyons que
$$
n \text{longueur}_R(M/xM) = \text{longueur}_R(M/x^nM).
$$
Le même argument pour $M = R^{\oplus r}$ montre que
$$
n \text{longueur}_R(R^{\oplus r}/xR^{\oplus r}) =
\text{longueur}_R(R^{\oplus r}/x^nR^{\oplus r}).
$$
Par notre choix de $c$ ci-dessus, nous voyons que $x^nM$ est compris entre
$x^n R^{\oplus r}$ et $x^{n + c}R^{\oplus r}$. Cela donne facilement
$$
r(n + c) \text{longueur}_R(R/xR)
\geq
n \text{longueur}_R(M/xM)
\geq
r (n - c) \text{longueur}_R(R/xR)
$$
Ainsi, dans le cas de type fini, nous obtenons en fait le résultat du lemme avec
égalité.

\medskip\noindent
Supposons maintenant que $M$ ne soit pas de type fini. Supposons que la longueur de $M/xM$ soit
$\geq k$ pour un certain entier naturel $k$. Nous pouvons alors trouver
$$
0 \subset N_0 \subset N_1 \subset N_2 \subset \ldots N_k \subset M/xM
$$
avec $N_i \not = N_{i + 1}$ pour $i = 0, \ldots k - 1$.
Choisissons un élément $m_i \in M$ dont la classe modulo $xM$ appartient
à $N_i$, mais pas à $N_{i - 1}$, pour $i = 1, \ldots, k$.
Considérons le $R$-module de type fini $M' = Rm_1 + \ldots + Rm_k \subset M$.
Soit $N'_i \subset M'/xM'$ l'image réciproque de $N_i$.
Il est clair que $N'_i \not =N'_{i + 1}$ par notre choix de $m_i$.
Nous voyons donc que $\text{longueur}_R(M'/xM') \geq k$. D'après le
cas de type fini, nous concluons que $k \leq r\text{longueur}_R(R/xR)$,
comme voulu.
\end{proof}

\noindent
Voici une première application.

\begin{lemma}
\label{algebra-lemma-finite-extension-residue-fields-dimension-1}
Soit $R \to S$ un homomorphisme entre anneaux intègres qui induit une
injection de corps des fractions $K \subset L$. Si $R$ est un anneau local
noethérien de dimension $1$ et si $[L : K] < \infty$, alors
\begin{enumerate}
\item tout idéal premier $\mathfrak n_i$ de $S$ au-dessus de
l'idéal maximal $\mathfrak m$ de $R$ est maximal,
\item ces idéaux sont en nombre fini, et
\item $[\kappa(\mathfrak n_i) : \kappa(\mathfrak m)] < \infty$ pour tout $i$.
\end{enumerate}
\end{lemma}

\begin{proof}
Choisissons un élément non nul $x \in \mathfrak m$. Appliquons le lemme \ref{algebra-lemma-finite-length}
au sous-module $S \subset L \cong K^{\oplus n}$, où $n = [L : K]$.
Ainsi, l'anneau $S/xS$ est de longueur finie sur $R$. Il s'ensuit que
$S/\mathfrak m S$ est de longueur finie sur $\kappa(\mathfrak m)$.
Autrement dit, $\dim_{\kappa(\mathfrak m)} S/\mathfrak m S$
est finie (lemme \ref{algebra-lemma-dimension-is-length}). Ainsi, $S/\mathfrak mS$
est artinien (lemme \ref{algebra-lemma-finite-dimensional-algebra}). Les
résultats de structure sur les anneaux artiniens impliquent (1) et (2), voir
par exemple le lemme \ref{algebra-lemma-artinian-finite-length}.
Le point (3) résulte de la finitude établie ci-dessus.
\end{proof}

\begin{lemma}
\label{algebra-lemma-finite-length-global}
Soit $R$ un anneau intègre de corps des fractions $K$.
Soit $M$ un $R$-sous-module de $K^{\oplus r}$.
Supposons que $R$ soit noethérien de dimension $1$.
Pour tout élément non nul $x \in R$, nous avons
$\text{longueur}_R(M/xM) < \infty$.
\end{lemma}

\begin{proof}
Puisque $R$ est de dimension $1$, l'élément $x$ n'appartient qu'à un nombre fini
d'idéaux premiers $\mathfrak m_i$, $i = 1, \ldots, n$, qui sont tous maximaux. Puisque $R$ est
noethérien, nous voyons que $R/xR$ est artinien et que
$R/xR = \prod_{i = 1, \ldots, n} (R/xR)_{\mathfrak m_i}$ d'après la
proposition \ref{algebra-proposition-dimension-zero-ring} et le
lemme \ref{algebra-lemma-artinian-finite-length}. De même, $M/xM$
se décompose en produit $M/xM = \prod (M/xM)_{\mathfrak m_i}$
de ses localisés aux $\mathfrak m_i$. D'après le lemme \ref{algebra-lemma-finite-length}
appliqué à $M_{\mathfrak m_i}$, considéré comme $R_{\mathfrak m_i}$-module, chacun des
$M_{\mathfrak m_i}/xM_{\mathfrak m_i} = (M/xM)_{\mathfrak m_i}$
est de longueur finie sur $R_{\mathfrak m_i}$. Ainsi, $M/xM$
est de longueur finie sur $R$, car ce qui précède implique que $M/xM$
possède une filtration finie par des $R$-sous-modules dont les quotients
successifs sont isomorphes aux corps résiduels $\kappa(\mathfrak m_i)$.
\end{proof}

\begin{lemma}[Krull-Akizuki]
\label{algebra-lemma-krull-akizuki}
Soit $R$ un anneau intègre de corps des fractions $K$.
Soit $L/K$ une extension finie de corps.
Supposons que $R$ soit noethérien et que $\dim(R) = 1$.
Dans ce cas, tout anneau $A$ tel que $R \subset A \subset L$ est
noethérien.
\end{lemma}

\begin{proof}
Soit $I \subset A$ un idéal non nul. D'après le
lemme \ref{algebra-lemma-intersection-not-zero},
nous pouvons trouver un élément non nul $x \in I \cap R$.
Nous obtenons alors $I/xA \subset A/xA$. D'après le lemme \ref{algebra-lemma-finite-length-global},
le $R$-module $A/xA$ est de longueur finie comme $R$-module. Ainsi,
$I/xA$ est de longueur finie comme $R$-module. Par conséquent, $I$ est de type fini
comme idéal de $A$.
\end{proof}

\begin{lemma}
\label{algebra-lemma-exists-dvr}
Soit $R$ un anneau local noethérien intègre de corps des fractions $K$.
Supposons que $R$ ne soit pas un corps.
Soit $L/K$ une extension de corps de type fini.
Alors il existe un anneau de valuation discrète $A$ de corps des fractions
$L$ qui domine $R$.
\end{lemma}

\begin{proof}
Si $L$ n'est pas fini sur $K$, choisissons une base de transcendance
$x_1, \ldots, x_r$ de $L$ sur $K$ et remplaçons $R$ par
$R[x_1, \ldots, x_r]$ localisé en l'idéal maximal
engendré par $\mathfrak m_R$ et $x_1, \ldots, x_r$.
Nous pouvons donc supposer que $K \subset L$ est une extension finie.

\medskip\noindent
D'après le lemme \ref{algebra-lemma-dominate-by-dimension-1}, nous pouvons supposer $\dim(R) = 1$.

\medskip\noindent
Soit $A \subset L$ la clôture intégrale de $R$ dans $L$.
D'après le lemme \ref{algebra-lemma-krull-akizuki}, elle est noethérienne.
D'après le lemme \ref{algebra-lemma-integral-overring-surjective}, il
existe un idéal premier $\mathfrak q \subset A$ au-dessus
de l'idéal maximal de $R$.
D'après le lemme \ref{algebra-lemma-characterize-dvr}, l'anneau $A_{\mathfrak q}$ est un anneau de
valuation discrète qui domine $R$, comme souhaité.
\end{proof}








\section{Factorisation}
\label{algebra-section-factoring}

\noindent
Voici quelques notions et les relations entre elles qui sont généralement enseignées
dans un cours d'algèbre de première année de licence.

\begin{definition}
\label{algebra-definition-irreducible-prime-element}
Soit $R$ un anneau intègre.
\begin{enumerate}
\item Des éléments $x, y \in R$ sont dits {\it associés} s'il
existe une unité $u \in R^*$ telle que $x = uy$.
\item Un élément $x \in R$ est dit {\it irréductible}
s'il est non nul, n'est pas une unité et si, pour toute écriture $x = yz$ avec
$y, z \in R$, $y$ est soit une unité, soit associé à $x$.
\item Un élément $x \in R$ est dit {\it premier} si l'idéal
engendré par $x$ est un idéal premier.
\end{enumerate}
\end{definition}

\begin{lemma}
\label{algebra-lemma-easy-divisibility}
Soit $R$ un anneau intègre. Soient $x, y \in R$.
Alors $x$ et $y$ sont associés si et seulement si $(x) = (y)$.
\end{lemma}

\begin{proof}
Si $x = uy$ pour une unité $u \in R$, alors $(x) \subset (y)$ et
$y = u^{-1}x$, donc aussi $(y) \subset (x)$. Réciproquement, supposons que
$(x) = (y)$. Alors $x = fy$ et $y = gx$ pour certains $f, g \in A$.
Alors $x = fg x$ et, puisque $R$ est un anneau intègre, $fg = 1$. Ainsi,
$x$ et $y$ sont associés.
\end{proof}

\begin{lemma}
\label{algebra-lemma-factorization-exists}
Soit $R$ un anneau intègre. Considérons les conditions suivantes :
\begin{enumerate}
\item L'anneau $R$ vérifie la condition de chaîne ascendante sur les
idéaux principaux.
\item Tout élément $a \in R$ non nul qui n'est pas une unité
admet une factorisation $a = b_1 \ldots b_k$
où chaque $b_i$ est un élément irréductible de $R$.
\end{enumerate}
Alors (1) implique (2).
\end{lemma}

\begin{proof}
Soit $x$ un élément non nul, qui n'est pas une unité et qui n'admet pas de
factorisation en
irréductibles. Posons $x_1 = x$. Nous pouvons écrire $x = yz$, où aucun des deux facteurs
$y$ et $z$ n'est irréductible ni inversible.
Alors soit $y$ n'admet pas de factorisation
en irréductibles, auquel cas nous posons $x_2 = y$, soit $z$ n'admet pas
de factorisation en irréductibles, auquel cas nous posons $x_2 = z$.
En poursuivant ainsi, nous obtenons une suite
$$
\ldots | x_3 | x_2 | x_1
$$
d'éléments de $R$ telle que $x_n/x_{n + 1}$ n'est pas une unité.
Cela donne une suite strictement croissante d'idéaux principaux
$(x_1) \subset (x_2) \subset (x_3) \subset \ldots$, ce qui achève la démonstration.
\end{proof}

\begin{definition}
\label{algebra-definition-UFD}
Un {\it anneau factoriel}, aussi appelé domaine à factorisation unique ({\it UFD}),
est un anneau intègre $R$ tel que,
si $x \in R$ est non nul et n'est pas une unité, alors $x$ admet une factorisation
en irréductibles, et si
$$
x = a_1 \ldots a_m = b_1 \ldots b_n
$$
sont des factorisations en irréductibles, alors $n = m$ et
il existe une permutation $\sigma : \{1, \ldots, n\} \to \{1, \ldots, n\}$
telle que $a_i$ et $b_{\sigma(i)}$ sont associés.
\end{definition}

\begin{lemma}
\label{algebra-lemma-characterize-UFD}
Soit $R$ un anneau intègre. Supposons que tout élément non nul qui n'est pas une unité se factorise en
irréductibles. Alors $R$ est un anneau factoriel si et seulement si tout
élément irréductible est premier.
\end{lemma}

\begin{proof}
Supposons que $R$ soit un anneau factoriel et soit $x \in R$ un élément irréductible.
Supposons que $ab \in (x)$, c'est-à-dire $ab = cx$. Choisissons des factorisations
$a = a_1 \ldots a_n$, $b = b_1 \ldots b_m$ et $c = c_1 \ldots c_r$.
Par unicité de la factorisation
$$
a_1 \ldots a_n b_1 \ldots b_m = c_1 \ldots c_r x
$$
nous constatons que $x$ est associé à l'un des éléments
$a_1, \ldots, b_m$. Autrement dit, soit $a \in (x)$, soit $b \in (x)$,
et nous concluons que $x$ est premier.

\medskip\noindent
Supposons que tout élément irréductible soit premier. Nous devons démontrer que
deux factorisations en irréductibles ne diffèrent que par l'ordre des facteurs
et leur remplacement par des éléments associés. Supposons que $a_1 \ldots a_m = b_1 \ldots b_n$, où les $a_i$
et les $b_j$ sont irréductibles. Puisque $a_1$ est premier, nous voyons que
$b_j \in (a_1)$ pour un certain $j$. Après renumérotation, nous pouvons supposer que
$b_1 \in (a_1)$. Alors $b_1 = a_1 u$ et, puisque $b_1$ est irréductible,
nous voyons que $u$ est une unité. Par conséquent, $a_1$ et $b_1$ sont associés
et $a_2 \ldots a_n = ub_2\ldots b_m$. Par récurrence sur $n + m$,
nous voyons que $n = m$ et que $a_i$ est associé à $b_{\sigma(i)}$ pour
$i = 2, \ldots, n$, ce qu'il fallait démontrer.
\end{proof}

\begin{lemma}
\label{algebra-lemma-characterize-UFD-height-1}
Soit $R$ un anneau intègre noethérien. Alors $R$ est un anneau factoriel si et seulement si tout
idéal premier de hauteur $1$ est principal.
\end{lemma}

\begin{proof}
Supposons que $R$ soit un anneau factoriel et soit $\mathfrak p$ un idéal premier de hauteur 1.
Prenons $x \in \mathfrak p$ non nul et soit $x = a_1 \ldots a_n$
une factorisation en irréductibles.
Puisque $\mathfrak p$ est premier, nous voyons que $a_i \in \mathfrak p$
pour un certain $i$. Par le lemme \ref{algebra-lemma-characterize-UFD}, l'idéal
$(a_i)$ est premier. Puisque $\mathfrak p$ est de hauteur $1$, nous concluons
que $(a_i) = \mathfrak p$.

\medskip\noindent
Supposons que tout idéal premier de hauteur $1$ soit principal. Puisque $R$ est noethérien,
tout élément $x$ non nul qui n'est pas une unité admet une factorisation en irréductibles,
voir le lemme \ref{algebra-lemma-factorization-exists}. Il suffit de démontrer
qu'un élément irréductible $x$ est premier, voir le lemme \ref{algebra-lemma-characterize-UFD}.
Soit $(x) \subset \mathfrak p$ un idéal premier minimal au-dessus de $(x)$. Alors
$\mathfrak p$ est de hauteur $1$ par le lemme \ref{algebra-lemma-minimal-over-1}.
Par hypothèse, $\mathfrak p = (y)$. Ainsi, $x = yz$ et $z$ est une unité
puisque $x$ est irréductible. Donc $(x) = (y)$ et nous voyons que $x$ est premier.
\end{proof}

\begin{lemma}[Critère de factorialité de Nagata]
\label{algebra-lemma-invert-prime-elements}
\begin{reference}
\cite[lemme 2]{Nagata-UFD}
\end{reference}
Soit $A$ un anneau intègre. Soit $S \subset A$ une partie multiplicative
engendrée par des éléments premiers. Soit $x \in A$ irréductible. Alors
\begin{enumerate}
\item l'image de $x$ dans $S^{-1}A$ est irréductible ou une unité, et
\item $x$ est premier si et seulement si l'image de $x$ dans $S^{-1}A$ est
un élément premier ou une unité dans $S^{-1}A$.
\end{enumerate}
De plus, $A$ est alors un anneau factoriel si et seulement si tout élément non nul qui n'est pas une unité
de $A$ admet une factorisation en irréductibles et $S^{-1}A$ est un anneau factoriel.
\end{lemma}

\begin{proof}
Écrivons $x = \alpha \beta$ pour $\alpha, \beta \in S^{-1}A$. Alors
$\alpha = a/s$ et $\beta = b/s'$ pour $a, b \in A$, $s, s' \in S$.
Nous obtenons donc $ss'x = ab$. Par hypothèse, nous pouvons écrire
$ss' = p_1 \ldots p_r$ pour certains éléments premiers $p_i$.
Pour chaque $i$, l'élément $p_i$ divise soit $a$, soit $b$.
En divisant, nous obtenons une factorisation $x = a' b'$ et
$a = s'' a'$, $b = s''' b'$ pour certains $s'', s''' \in S$.
Comme $x$ est irréductible, soit $a'$, soit $b'$ est une unité.
Il s'ensuit que l'un des deux éléments $\alpha$, $\beta$ est une unité.
Cela démontre (1).

\medskip\noindent
Supposons que $x$ soit premier. Alors $A/(x)$ est un anneau intègre. Par conséquent,
$S^{-1}A/xS^{-1}A = S^{-1}(A/(x))$ est un anneau intègre ou nul.
Ainsi, l'image de $x$ est un élément premier ou une unité.

\medskip\noindent
Supposons que l'image de $x$ dans $S^{-1}A$ soit une unité.
Alors $y x = s$ pour un $s \in S$ et un $y \in A$. Par hypothèse,
$s = p_1 \ldots p_r$, où $p_i$ est un élément premier. Pour chaque $i$, soit
$p_i$ divise $y$, soit $p_i$ divise $x$. Dans le second cas,
$p_i$ et $x$ sont associés (puisque $x$ est irréductible) ; la conclusion s'ensuit.
Mais si le premier cas se produit pour tout $i = 1, \ldots, r$, alors
$x$ est une unité, ce qui est une contradiction.

\medskip\noindent
Supposons que l'image de $x$ dans $S^{-1}A$ soit un élément premier.
Supposons que $a, b \in A$ et $ab \in (x)$. Alors $sa = xy$ ou
$sb = xy$ pour certains $s \in S$ et $y \in A$. Supposons que le premier
cas se produise. Par hypothèse,
$s = p_1 \ldots p_r$, où $p_i$ est un élément premier. Pour chaque $i$, soit
$p_i$ divise $y$, soit $p_i$ divise $x$. Dans le second cas,
$p_i$ et $x$ sont associés (puisque $x$ est irréductible) ; la conclusion s'ensuit.
Si le premier cas se produit pour tout $i = 1, \ldots, r$, alors
$a \in (x)$, ce qu'il fallait démontrer. Ceci achève la démonstration de (2).

\medskip\noindent
La dernière assertion du lemme découle de (1) et (2)
et du lemme \ref{algebra-lemma-characterize-UFD}.
\end{proof}

\begin{lemma}
\label{algebra-lemma-UFD-ascending-chain-condition-principal-ideals}
Un anneau factoriel vérifie la condition de chaîne ascendante sur les idéaux
principaux.
\end{lemma}

\begin{proof}
Considérons une chaîne ascendante $(a_1) \subset (a_2) \subset (a_3) \subset \ldots$
d'idéaux principaux de $R$. Écrivons $a_1 = p_1^{e_1} \ldots p_r^{e_r}$
où les $p_i$ sont premiers. Nous voyons alors que $a_n$ est associé à
$p_1^{c_1} \ldots p_r^{c_r}$ pour certains $0 \leq c_i \leq e_i$.
Puisqu'il n'y a qu'un nombre fini de possibilités, la chaîne est stationnaire.
\end{proof}

\begin{lemma}
\label{algebra-lemma-factoring-in-polynomial}
Soit $R$ un anneau intègre. Supposons que $R$ vérifie la condition de chaîne ascendante
sur les idéaux principaux. Alors la même propriété vaut pour un anneau de polynômes
sur $R$.
\end{lemma}

\begin{proof}
Considérons une chaîne ascendante $(f_1) \subset (f_2) \subset (f_3) \subset \ldots$
d'idéaux principaux de $R[x]$. Puisque $f_{n + 1}$ divise $f_n$, nous voyons
que les degrés décroissent dans la suite. Ainsi, $f_n$
a un degré fixe $d \geq 0$ pour tout $n \gg 0$. Soit $a_n$ le
coefficient dominant de $f_n$. La condition $f_n \in (f_{n + 1})$
implique que $a_{n + 1}$ divise $a_n$ pour tout $n$.
Par notre hypothèse sur $R$, nous voyons que $a_{n + 1}$ et $a_n$
sont associés pour tout $n$ assez grand (lemme \ref{algebra-lemma-easy-divisibility}).
Ainsi, pour $n$ grand, nous voyons que $f_n = u f_{n + 1}$, où
$u \in R$ (pour des raisons de degré) est une unité (puisque $a_n$ et $a_{n + 1}$
sont associés).
\end{proof}

\begin{lemma}
\label{algebra-lemma-polynomial-ring-UFD}
Un anneau de polynômes sur un anneau factoriel est un anneau factoriel. En particulier, si $k$ est un corps,
alors $k[x_1, \ldots, x_n]$ est un anneau factoriel.
\end{lemma}

\begin{proof}
Soit $R$ un anneau factoriel. Alors $R$ vérifie la condition de chaîne ascendante sur les
idéaux principaux
(lemme \ref{algebra-lemma-UFD-ascending-chain-condition-principal-ideals}),
donc $R[x]$ vérifie la condition de chaîne ascendante sur les idéaux principaux
(lemme \ref{algebra-lemma-factoring-in-polynomial}), et tout
élément de $R[x]$ admet donc une factorisation en irréductibles
(lemme \ref{algebra-lemma-factorization-exists}).
Soit $S \subset R$ la partie multiplicative
engendrée par les éléments premiers. Puisque tout élément non inversible de $R$ est un produit
d'éléments premiers, nous voyons que $K = S^{-1}R$ est le corps des fractions
de $R$. Observons que tout élément premier de $R$ reste premier
dans $R[x]$ et que $S^{-1}(R[x]) = S^{-1}R[x] = K[x]$ est un
anneau factoriel (et même un anneau principal). Nous pouvons donc appliquer le
lemme \ref{algebra-lemma-invert-prime-elements} pour conclure.
\end{proof}

\begin{lemma}
\label{algebra-lemma-UFD-normal}
Un anneau factoriel est normal.
\end{lemma}

\begin{proof}
Soit $R$ un anneau factoriel. Soit $x$ un élément du corps des fractions de
$R$ qui est entier sur $R$. Écrivons $x^d - a_1 x^{d - 1} - \ldots - a_d = 0$
avec $a_i \in R$. Nous pouvons écrire
$x = u p_1^{e_1} \ldots p_r^{e_r}$ où $u$ est une unité, $e_i \in \mathbf{Z}$,
et $p_1, \ldots, p_r$ sont des éléments irréductibles qui ne sont pas associés.
Pour démontrer le lemme, nous devons montrer que $e_i \geq 0$. Sinon, disons $e_1 < 0$,
alors, pour $N \gg 0$, nous obtenons
$$
u^d p_2^{de_2 + N} \ldots p_r^{de_r + N} =
p_1^{-de_1}p_2^N \ldots p_r^N(
\sum\nolimits_{i = 1, \ldots, d} a_i x^{d - i} ) \in (p_1)
$$
ce qui contredit l'unicité de la factorisation dans $R$.
\end{proof}

\begin{definition}
\label{algebra-definition-PID}
Un {\it anneau principal} (en abrégé {\it PID})
est un anneau intègre $R$ tel que tout idéal est un idéal principal.
\end{definition}

\begin{lemma}
\label{algebra-lemma-PID-UFD}
Un anneau principal est un anneau factoriel.
\end{lemma}

\begin{proof}
Comme un anneau principal est noethérien, cela résulte du
lemme \ref{algebra-lemma-characterize-UFD-height-1}.
\end{proof}

\begin{definition}
\label{algebra-definition-dedekind-domain}
Un {\it anneau de Dedekind} est un anneau intègre $R$ tel que tout
idéal non nul $I \subset R$ puisse s'écrire comme un produit
$$
I = \mathfrak p_1 \ldots \mathfrak p_r
$$
d'idéaux premiers non nuls, de manière unique à l'ordre des facteurs $\mathfrak p_i$ près.
\end{definition}

\begin{lemma}
\label{algebra-lemma-PID-dedekind}
Un anneau principal est un anneau de Dedekind.
\end{lemma}

\begin{proof}
Soit $R$ un anneau principal. Puisque tout idéal non nul de $R$ est principal,
et que $R$ est un anneau factoriel (lemme \ref{algebra-lemma-PID-UFD}), cela résulte
du fait que tout élément irréductible de $R$ est premier
(lemme \ref{algebra-lemma-characterize-UFD}), chaque factorisation d'un élément
induisant alors une factorisation de l'idéal principal correspondant en idéaux premiers.
\end{proof}

\begin{lemma}
\label{algebra-lemma-product-ideals-principal}
\begin{slogan}
Un produit d'idéaux est un module inversible si et seulement si ses deux facteurs le sont.
\end{slogan}
Soit $A$ un anneau. Soient $I$ et $J$ des idéaux non nuls de $A$
tels que $IJ = (f)$ pour un non-diviseur de zéro $f \in A$. Alors $I$ et $J$ sont
des idéaux de type fini et sont, de plus, localement libres de type fini de rang $1$ comme $A$-modules.
\end{lemma}

\begin{proof}
Il suffit de montrer que $I$ et $J$ sont des $A$-modules localement libres de type fini
de rang $1$, voir le lemme \ref{algebra-lemma-finite-projective}.
Pour ce faire, écrivons $f = \sum_{i = 1, \ldots, n} x_i y_i$ avec $x_i \in I$
et $y_i \in J$. Nous pouvons
aussi écrire $x_i y_i = a_i f$ pour un certain $a_i \in A$.
Puisque $f$ est un non-diviseur de zéro, nous voyons que $\sum a_i = 1$.
Il suffit donc de montrer que les modules $I_{a_i}$ et
$J_{a_i}$ sont libres de rang $1$ sur $A_{a_i}$. Après avoir remplacé $A$ par
$A_{a_i}$, nous concluons que $f = xy$ pour certains $x \in I$ et $y \in J$.
Notons que $x$ et $y$ sont tous deux des non-diviseurs de zéro. Nous affirmons que
$I = (x)$ et $J = (y)$, ce qui achève la démonstration. En effet, si $x' \in I$,
alors $x'y = af = axy$ pour un certain $a \in A$. Ainsi, $x' = ax$ ; le même argument, avec les rôles intervertis, achève la démonstration.
\end{proof}

\begin{lemma}
\label{algebra-lemma-characterize-Dedekind}
Soit $R$ un anneau. Les conditions suivantes sont équivalentes :
\begin{enumerate}
\item $R$ est un anneau de Dedekind,
\item $R$ est un anneau intègre noethérien et, pour tout idéal maximal non nul
$\mathfrak m$, l'anneau local $R_{\mathfrak m}$ est un anneau de valuation discrète, et
\item $R$ est un anneau intègre noethérien et normal, et $\dim(R) \leq 1$.
\end{enumerate}
\end{lemma}

\begin{proof}
Supposons (1). L'argument n'est pas trivial, car nous n'avons pas supposé que $R$
était noethérien dans notre définition d'un anneau de Dedekind. Soit
$\mathfrak p \subset R$ un idéal premier. Observons que
$\mathfrak p \not = \mathfrak p^2$ par unicité
des factorisations de la définition. Choisissons $x \in \mathfrak p$
tel que $x \not \in \mathfrak p^2$. Soit $y \in \mathfrak p$ un
second élément (par exemple $y = 0$).
Écrivons $(x, y) = \mathfrak p_1 \ldots \mathfrak p_r$.
Puisque $(x, y) \subset \mathfrak p$, au moins l'un des idéaux premiers
$\mathfrak p_i$ est contenu dans $\mathfrak p$.
Mais puisque $x \not \in \mathfrak p^2$, il y en a au plus un.
Ainsi, exactement l'un des idéaux $\mathfrak p_1, \ldots, \mathfrak p_r$ est contenu
dans $\mathfrak p$, disons $\mathfrak p_1 \subset \mathfrak p$.
Nous concluons que $(x, y)R_\mathfrak p = \mathfrak p_1R_\mathfrak p$
est premier pour tout choix de $y$. Nous affirmons que
$(x)R_\mathfrak p = \mathfrak pR_\mathfrak p$. En effet,
choisissons $y \in \mathfrak p$. D'après ce qui précède, appliqué à $y^2$, nous voyons
que $(x, y^2)R_\mathfrak p$ est premier.
Par conséquent, $y \in (x, y^2)R_\mathfrak p$, c'est-à-dire $y = ax + by^2$ dans
$R_\mathfrak p$. Ainsi, $(1 - by)y = ax \in (x)R_\mathfrak p$, c'est-à-dire
$y \in (x)R_\mathfrak p$, comme il fallait le montrer.

\medskip\noindent
En écrivant à nouveau $(x) = \mathfrak p_1 \ldots \mathfrak p_r$ avec
$\mathfrak p_1 \subset \mathfrak p$, nous concluons que
$\mathfrak p_1 R_\mathfrak p = \mathfrak p R_\mathfrak p$, c'est-à-dire
$\mathfrak p_1 = \mathfrak p$. De plus, $\mathfrak p_1 = \mathfrak p$ est
un idéal de type fini de $R$ par le
lemme \ref{algebra-lemma-product-ideals-principal}.
Nous concluons que $R$ est noethérien par le lemme \ref{algebra-lemma-cohen}.
De plus, il s'ensuit que $R_\mathfrak m$ est un anneau de
valuation discrète pour tout idéal premier $\mathfrak p$, voir le
lemme \ref{algebra-lemma-characterize-dvr}.

\medskip\noindent
L'équivalence de (2) et (3) résulte des
lemmes \ref{algebra-lemma-normality-is-local} et
\ref{algebra-lemma-characterize-dvr}. Supposons que (2) et (3) soient satisfaites.
Soit $I \subset R$ un idéal. Nous allons construire une factorisation
de $I$. Si $I$ est premier, il n'y a rien à démontrer.
Sinon, choisissons $I \subset \mathfrak p$, où $\mathfrak p \subset R$ est un idéal maximal.
Posons $J = \{x \in R \mid x \mathfrak p \subset I\}$.
Nous affirmons que $J \mathfrak p = I$. Il suffit de le vérifier après localisation
aux idéaux maximaux $\mathfrak m$ de $R$ (la formation de $J$ commute à la
localisation et nous utilisons le lemme \ref{algebra-lemma-characterize-zero-local}).
Alors, soit $\mathfrak p R_\mathfrak m = R_\mathfrak m$ et le résultat
est clair, soit $\mathfrak p R_\mathfrak m = \mathfrak m R_\mathfrak m$.
Dans ce dernier cas, $\mathfrak p R_\mathfrak m = (\pi)$ et le cas où
$\mathfrak p$ est principal est immédiat. Par récurrence noethérienne,
l'idéal $J$ admet une factorisation et nous obtenons la factorisation souhaitée
de $I$. Nous omettons la démonstration de l'unicité de la factorisation.
\end{proof}

\noindent
Le lemme suivant est une variante du théorème de Krull-Akizuki.

\begin{lemma}
\label{algebra-lemma-integral-closure-Dedekind}
Soit $A$ un anneau intègre noethérien de dimension $1$ et de corps des fractions $K$.
Soit $L/K$ une extension finie. Soit $B$ la
clôture intégrale de $A$ dans $L$. Alors $B$ est un anneau de Dedekind et
$\Spec(B) \to \Spec(A)$ est surjective, a des fibres finies et
induit des extensions finies de corps résiduels.
\end{lemma}

\begin{proof}
D'après le théorème de Krull-Akizuki (lemme \ref{algebra-lemma-krull-akizuki}),
l'anneau $B$ est noethérien. Par le lemme \ref{algebra-lemma-integral-sub-dim-equal},
$\dim(B) = 1$. Ainsi, $B$ est un anneau de Dedekind par le
lemme \ref{algebra-lemma-characterize-Dedekind}.
La surjectivité de l'application sur les spectres résulte du
lemme \ref{algebra-lemma-integral-overring-surjective}.
Les deux dernières assertions résultent du
lemme \ref{algebra-lemma-finite-extension-residue-fields-dimension-1}.
\end{proof}







\section{Ordres d'annulation}
\label{algebra-section-orders-of-vanishing}

\begin{lemma}
\label{algebra-lemma-ord-additive}
Soit $R$ un anneau semi-local noethérien de dimension $1$.
Si $a, b \in R$ sont des non-diviseurs de zéro, alors
$$
\text{longueur}_R(R/(ab)) =
\text{longueur}_R(R/(a)) +
\text{longueur}_R(R/(b))
$$
et ces longueurs sont finies.
\end{lemma}

\begin{proof}
Cette finitude a été établie au lemme \ref{algebra-lemma-finite-length-global}.
L'additivité résulte de l'existence d'une suite exacte courte
$0 \to R/(a) \to R/(ab) \to R/(b) \to 0$, où la première flèche
est la multiplication par $b$. (On utilise l'additivité de la longueur ;
voir le lemme \ref{algebra-lemma-length-additive}.)
\end{proof}

\begin{definition}
\label{algebra-definition-ord}
Supposons que $K$ soit un corps et que $R \subset K$ soit un sous-anneau
local\footnote{On pourrait aussi donner cette définition lorsque $R$ est seulement
semi-local, mais ce n'est sans doute jamais vraiment la situation recherchée !}
noethérien de dimension $1$, de corps des fractions $K$.
Dans ce cas, on définit l'{\it ordre d'annulation relatif à $R$}
$$
\text{ord}_R : K^* \longrightarrow \mathbf{Z}
$$
par la règle
$$
\text{ord}_R(x) = \text{longueur}_R(R/(x))
$$
si $x \in R$, et l'on pose
$\text{ord}_R(x/y) = \text{ord}_R(x) - \text{ord}_R(y)$
pour $x, y \in R$ tous deux non nuls.
\end{definition}

\noindent
On peut utiliser l'ordre d'annulation pour comparer des réseaux dans un
espace vectoriel. En voici la définition.

\begin{definition}
\label{algebra-definition-lattice}
Soit $R$ un anneau intègre local noethérien de dimension $1$, de
corps des fractions $K$. Soit $V$ un espace vectoriel de dimension finie sur $K$.
Un {\it réseau dans $V$} est un sous-$R$-module de type fini $M \subset V$ tel
que $V = K \otimes_R M$.
\end{definition}

\noindent
La condition $V = K \otimes_R M$ signifie que $M$ contient une
base de l'espace vectoriel $V$. Signalons que, dans de nombreux ouvrages,
la notion de réseau n'est définie que lorsque
l'anneau $R$ est un anneau de valuation discrète. Si $R$ est un anneau de valuation
discrète, tout réseau est un $R$-module libre, ce qui n'est pas nécessairement le cas
en général.

\begin{lemma}
\label{algebra-lemma-compare-lattices}
Soit $R$ un anneau intègre local noethérien de dimension $1$, de
corps des fractions $K$. Soit $V$ un espace vectoriel de dimension finie sur $K$.
\begin{enumerate}
\item Si $M$ est un réseau dans $V$ et si $M \subset M' \subset V$
est un sous-$R$-module de $V$ contenant $M$,
alors les assertions suivantes sont équivalentes :
\begin{enumerate}
\item $M'$ est un réseau,
\item $\text{longueur}_R(M'/M)$ est finie, et
\item $M'$ est de type fini.
\end{enumerate}
\item Si $M$ est un réseau dans $V$ et si $M' \subset M$ est un sous-$R$-module
de $M$, alors $M'$ est un réseau si et seulement si
$\text{longueur}_R(M/M')$ est finie.
\item Si $M$, $M'$ sont des réseaux dans $V$, alors il en est de même de
$M \cap M'$ et $M + M'$.
\item Si $M \subset M' \subset M'' \subset V$ sont des réseaux dans $V$,
alors
$$
\text{longueur}_R(M''/M) =
\text{longueur}_R(M'/M) +
\text{longueur}_R(M''/M').
$$
\item Si $M$, $M'$, $N$, $N'$ sont des réseaux dans $V$ et si
$N \subset M \cap M'$, $M + M' \subset N'$, alors on a
\begin{eqnarray*}
& & \text{longueur}_R(M/M \cap M') - \text{longueur}_R(M'/M \cap M')\\
& = &
\text{longueur}_R(M/N) - \text{longueur}_R(M'/N) \\
& = &
\text{longueur}_R(M + M' / M') - \text{longueur}_R(M + M'/M) \\
& = &
\text{longueur}_R(N' / M') - \text{longueur}_R(N'/M)
\end{eqnarray*}
\end{enumerate}
\end{lemma}

\begin{proof}
Démonstration de (1). Supposons (1)(a) vérifiée. Soient $y_1, \ldots, y_m$ des générateurs de $M'$.
Alors chaque $y_i = x_i/f_i$ avec $x_i \in M$ et
$f_i \in R$ non nul.
On en déduit que $f_1 \ldots f_m M' \subset M$.
Puisque $R$ est un anneau local noethérien de dimension $1$,
on a $\mathfrak m^n \subset (f_1 \ldots f_m)$
pour un certain $n$ (par exemple en combinant les
lemmes \ref{algebra-lemma-one-equation} et la
proposition \ref{algebra-proposition-dimension-zero-ring}, ou en combinant les
lemmes \ref{algebra-lemma-finite-length} et \ref{algebra-lemma-length-infinite}).
En d'autres termes, $\mathfrak m^nM' \subset M$ pour un certain $n$.
Par conséquent,
$\text{longueur}(M'/M) < \infty$ d'après le lemme \ref{algebra-lemma-length-finite},
autrement dit, (1)(b) est vérifiée.
Supposons (1)(b). Alors $M'/M$ est un $R$-module de type fini
(voir le lemme \ref{algebra-lemma-finite-length-finite}).
Par conséquent, $M'$ est un $R$-module de type fini puisqu'il est une extension de $R$-modules de type fini.
Donc (1)(c) est vérifiée. L'implication
(1)(c) $\Rightarrow$ (1)(a) résulte de la remarque qui suit la
définition \ref{algebra-definition-lattice}.

\medskip\noindent
Démonstration de (2). Supposons que
$M$ soit un réseau dans $V$ et que $M' \subset M$ soit un sous-$R$-module.
Nous avons vu en (1) que, si $M'$ est un réseau, alors
$\text{longueur}_R(M/M') < \infty$. Réciproquement, supposons que
$\text{longueur}_R(M/M') < \infty$. Alors $M'$ est de type fini
puisque $R$ est noethérien. De plus, pour un certain $n$, on a
$\mathfrak m^n M \subset M'$ (lemme \ref{algebra-lemma-length-infinite}).
Il s'ensuit
que $M'$ contient une base de $V$ et que $M'$ est un réseau.

\medskip\noindent
Démonstration de (3). Supposons que $M$, $M'$ soient des réseaux dans $V$.
Puisque $R$ est noethérien, le sous-module $M \cap M'$ de $M$ est de type fini.
Comme $M$ est un réseau, nous pouvons trouver
$x_1, \ldots, x_n \in M$ qui forment une $K$-base de
$V$. Puisque $M'$ est un réseau, nous pouvons écrire $x_i = y_i/f_i$ avec
$y_i \in M'$ et $f_i \in R$. Ainsi $f_ix_i \in M \cap M'$. Par conséquent,
$M \cap M'$ est lui aussi un réseau.
Le fait que $M + M'$ soit un réseau résulte du point (1).

\medskip\noindent
Le point (4) résulte de l'additivité des longueurs
(lemme \ref{algebra-lemma-length-additive})
et de la suite exacte
$$
0 \to M'/M \to M''/M \to M''/M' \to 0
$$
Le point (5) s'obtient en appliquant le point (4) à plusieurs reprises.
\end{proof}

\begin{definition}
\label{algebra-definition-distance}
Soit $R$ un anneau intègre local noethérien de dimension $1$, de
corps des fractions $K$. Soit $V$ un espace vectoriel de dimension finie sur $K$.
Soient $M$, $M'$ deux réseaux dans $V$. La {\it distance de
$M$ à $M'$} est l'entier
$$
d(M, M') = \text{longueur}_R(M/M \cap M') - \text{longueur}_R(M'/M \cap M')
$$
défini au point (5) du lemme \ref{algebra-lemma-compare-lattices}.
\end{definition}

\noindent
En particulier, si $M' \subset M$, alors
$d(M, M') = \text{longueur}_R(M/M')$.

\begin{lemma}
\label{algebra-lemma-properties-distance-function}
Soit $R$ un anneau intègre local noethérien de dimension $1$, de
corps des fractions $K$. Soit $V$ un espace vectoriel de dimension finie sur $K$.
Cette fonction de distance vérifie
$$
d(M, M'') = d(M, M') + d(M', M'')
$$
quels que soient les trois réseaux $M$, $M'$, $M''$ de $V$.
En particulier, on a $d(M, M') = - d(M', M)$.
\end{lemma}

\begin{proof}
La démonstration est omise.
\end{proof}

\begin{lemma}
\label{algebra-lemma-order-vanishing-determinant}
Soit $R$ un anneau intègre local noethérien de dimension $1$, de
corps des fractions $K$. Soit $V$ un espace vectoriel de dimension finie sur $K$.
Soit $\varphi : V \to V$ un isomorphisme $K$-linéaire.
Pour tout réseau $M \subset V$, on a
$$
d(M, \varphi(M)) = \text{ord}_R(\det(\varphi))
$$
\end{lemma}

\begin{proof}
L'entier $d(M, \varphi(M))$ est indépendant du choix
du réseau $M$. En effet, soit $M'$ un autre
réseau. Alors
\begin{eqnarray*}
d(M, \varphi(M)) & = & d(M, M') + d(M', \varphi(M)) \\
& = & d(M, M') + d(\varphi(M'), \varphi(M)) + d(M', \varphi(M'))
\end{eqnarray*}
Puisque $\varphi$ est un isomorphisme, on a
$d(\varphi(M'), \varphi(M)) = d(M', M) = -d(M, M')$, et donc
$d(M, \varphi(M)) = d(M', \varphi(M'))$. De plus, les deux membres de
l'égalité du lemme, considérés comme fonctions de $\varphi$, sont additifs pour la composition, c'est-à-dire que
$$
\text{ord}_R(\det(\varphi \circ \psi))
=
\text{ord}_R(\det(\varphi))
+
\text{ord}_R(\det(\psi))
$$
et l'on a aussi
\begin{eqnarray*}
d(M, \varphi(\psi((M))) & = &
d(M, \psi(M)) + d(\psi(M), \varphi(\psi(M))) \\
& = & d(M, \psi(M)) + d(M, \varphi(M))
\end{eqnarray*}
par l'indépendance établie ci-dessus. Il suffit donc de démontrer le lemme
pour un système de générateurs de $\text{GL}(V)$. Choisissons un isomorphisme
$K^{\oplus n} \cong V$. Alors $\text{GL}(V) = \text{GL}_n(K)$ est
engendré par les matrices élémentaires $E$.
Le résultat est immédiat lorsque $E$ est la matrice identité.
Si $E = E_{ij}(\lambda)$ avec $i \not = j$, $\lambda \in K$,
$\lambda \not = 0$, par exemple
$$
E_{12}(\lambda) =
\left(
\begin{matrix}
1 & \lambda & \ldots \\
0 & 1 & \ldots \\
\ldots & \ldots & \ldots
\end{matrix}
\right)
$$
alors, après un changement de base, on obtient $E_{12}(1)$.
Le résultat est immédiat pour $E = E_{12}(1)$ en prenant pour réseau
$R^{\oplus n} \subset K^{\oplus n}$. Enfin, si $E = E_i(a)$,
où $a \in K^*$, par exemple
$$
E_1(a) =
\left(
\begin{matrix}
a & 0 & \ldots \\
0 & 1 & \ldots \\
\ldots & \ldots & \ldots
\end{matrix}
\right)
$$
alors $E_1(a)(R^{\oplus b}) = aR \oplus R^{\oplus n - 1}$ et
on voit immédiatement que $d(R^{\oplus n}, aR \oplus R^{\oplus n - 1})
= \text{ord}_R(a)$, d'où le résultat.
\end{proof}

\begin{lemma}
\label{algebra-lemma-finite-extension-dim-1}
Soit $A \to B$ un homomorphisme d'anneaux. Supposons que
\begin{enumerate}
\item $A$ est un anneau intègre local noethérien de dimension $1$,
\item $A \subset B$ est une inclusion finie entre anneaux intègres.
\end{enumerate}
Soit $L/K$ l'extension finie correspondante de leurs corps des fractions respectifs.
Soit $y \in L^*$ et posons $x = \text{Nm}_{L/K}(y)$.
Dans cette situation, $B$ est semi-local.
Soient $\mathfrak m_i$, $i = 1, \ldots, n$, les idéaux maximaux de $B$.
Alors
$$
\text{ord}_A(x) =
\sum\nolimits_i
[\kappa(\mathfrak m_i) : \kappa(\mathfrak m_A)]
\text{ord}_{B_{\mathfrak m_i}}(y)
$$
où $\text{ord}$ désigne l'ordre défini dans la définition \ref{algebra-definition-ord}.
\end{lemma}

\begin{proof}
L'anneau $B$ est semi-local d'après le lemme \ref{algebra-lemma-finite-in-codim-1}.
Écrivons $y = b/b'$ avec $b, b' \in B$.
Grâce à l'additivité de $\text{ord}$ et à la multiplicativité de
$\text{Nm}$, il suffit de démontrer le lemme pour
$y = b$ ou $y = b'$. Autrement dit, on peut supposer que $y \in B$.
Dans ce cas, le membre de droite de la formule est
$$
\sum [\kappa(\mathfrak m_i) : \kappa(\mathfrak m_A)]
\text{longueur}_{B_{\mathfrak m_i}}((B/yB)_{\mathfrak m_i})
$$
D'après le lemme \ref{algebra-lemma-pushdown-module}, cette expression est égale à
$\text{longueur}_A(B/yB)$. D'après le lemme \ref{algebra-lemma-order-vanishing-determinant},
on a
$$
\text{longueur}_A(B/yB) = d(B, yB) =
\text{ord}_A(\det\nolimits_K(L \xrightarrow{y} L)).
$$
Puisque $x = \text{Nm}_{L/K}(y) = \det\nolimits_K(L \xrightarrow{y} L)$
par définition, le lemme est démontré.
\end{proof}














\section{Morphismes quasi-finis}
\label{algebra-section-quasi-finite}

\noindent
Considérons un homomorphisme d'anneaux $R \to S$ de type fini.
Le morphisme $\Spec(S) \to \Spec(R)$ est quasi-fini
en un point si ce point est isolé dans sa fibre.
Cela signifie que la fibre est de dimension zéro en ce point.
Dans cette section, nous étudions les propriétés fondamentales de cette
notion importante mais technique. On trouvera des résultats plus avancés
dans la section suivante.

\begin{lemma}
\label{algebra-lemma-isolated-point}
Soit $k$ un corps.
Soit $S$ une $k$-algèbre de type fini.
Soit $\mathfrak q$ un idéal premier de $S$.
Les assertions suivantes sont équivalentes :
\begin{enumerate}
\item $\mathfrak q$ est un point isolé de $\Spec(S)$,
\item $S_{\mathfrak q}$ est fini sur $k$,
\item il existe $g \in S$, $g \not\in \mathfrak q$, tel que
$D(g) = \{ \mathfrak q \}$,
\item $\dim_{\mathfrak q} \Spec(S) = 0$,
\item $\mathfrak q$ est un point fermé de $\Spec(S)$ et
$\dim(S_{\mathfrak q}) = 0$, et
\item l'extension de corps $\kappa(\mathfrak q)/k$ est finie
et $\dim(S_{\mathfrak q}) = 0$.
\end{enumerate}
Dans ce cas, $S = S_{\mathfrak q} \times S'$ pour une
$k$-algèbre $S'$ de type fini. De plus, l'élément $g$
du point (3) vérifie $S_{\mathfrak q} = S_g$.
\end{lemma}

\begin{proof}
Supposons que $\mathfrak q$ soit un point isolé de $\Spec(S)$, c'est-à-dire que
$\{\mathfrak q\}$ soit ouvert dans $\Spec(S)$.
Puisque $\Spec(S)$ est un espace de Jacobson (voir les
lemmes \ref{algebra-lemma-finite-type-field-Jacobson} et
\ref{algebra-lemma-jacobson}),
on voit que $\mathfrak q$ est un point fermé. Par conséquent,
$\{\mathfrak q\}$ est ouvert et fermé dans $\Spec(S)$.
D'après les
lemmes \ref{algebra-lemma-disjoint-decomposition} et
\ref{algebra-lemma-disjoint-implies-product}, on peut
écrire $S = S_1 \times S_2$, où $\mathfrak q$
correspond à l'unique point de $\Spec(S_1)$.
Ainsi, $S_1 = S_{\mathfrak q}$ est un anneau de dimension
zéro et de type fini sur $k$. Il est donc fini sur $k$,
par exemple d'après le lemme \ref{algebra-lemma-Noether-normalization}.
Nous avons démontré que (1) implique (2).

\medskip\noindent
Supposons que $S_{\mathfrak q}$ soit fini sur $k$.
Alors $S_{\mathfrak q}$ est un anneau local artinien ; voir le
lemme \ref{algebra-lemma-finite-dimensional-algebra}. Ainsi,
$\Spec(S_{\mathfrak q}) = \{\mathfrak qS_{\mathfrak q}\}$ d'après le
lemme \ref{algebra-lemma-artinian-finite-length}.
Considérons la suite exacte $0 \to K \to S \to S_{\mathfrak q}
\to Q \to 0$. Il est clair que $K_{\mathfrak q} = Q_{\mathfrak q} = 0$.
De plus, $K$ est un $S$-module de type fini puisque $S$ est noethérien, et
$Q$ est un $S$-module de type fini puisque $S_{\mathfrak q}$ est fini sur $k$.
Il existe donc $g \in S$, $g \not \in \mathfrak q$, tel que
$K_g = Q_g = 0$. Ainsi, $S_{\mathfrak q} = S_g$ et
$D(g) = \{ \mathfrak q \}$. Nous avons démontré que (2) implique (3).

\medskip\noindent
Supposons que $D(g) = \{ \mathfrak q \}$. Comme $D(g)$ est ouvert par
définition de la topologie sur $\Spec(S)$, on voit que
$\mathfrak q$ est un point isolé de $\Spec(S)$.
Nous avons démontré que (3) implique (1).
Autrement dit, (1), (2) et (3) sont équivalentes.

\medskip\noindent
Supposons que $\dim_{\mathfrak q} \Spec(S) = 0$. Cela signifie qu'il
existe un voisinage ouvert de $\mathfrak q$ dans $\Spec(S)$
qui est de dimension zéro. Il existe alors un voisinage ouvert de la
forme $D(g)$ qui est de dimension zéro. Puisque $S_g$ est noethérien,
on en déduit que $S_g$ est artinien et que
$D(g) = \Spec(S_g)$ est un ensemble discret fini ; voir la
proposition \ref{algebra-proposition-dimension-zero-ring}.
Ainsi, $\mathfrak q$ est un point isolé de $D(g)$ et,
en appliquant l'équivalence de (1) et (2) ci-dessus à
$\mathfrak qS_g \subset S_g$, on voit que
$S_{\mathfrak q} = (S_g)_{\mathfrak qS_g}$ est fini sur $k$.
Par conséquent, (4) implique (2). Il est clair que (1) implique (4).
Ainsi, (1) -- (4) sont toutes équivalentes.

\medskip\noindent
Le lemme \ref{algebra-lemma-dimension-closed-point-finite-type-field}
donne l'implication (5) $\Rightarrow$ (4).
L'implication (4) $\Rightarrow$ (6) résulte du
lemme \ref{algebra-lemma-dimension-at-a-point-finite-type-field}.
L'implication (6) $\Rightarrow$ (5) résulte du
lemme \ref{algebra-lemma-finite-residue-extension-closed}.
À ce stade, nous savons que (1) -- (6) sont équivalentes.

\medskip\noindent
Les deux assertions finales du lemme ont été obtenues au
cours de la démonstration de l'équivalence de (1), (2) et (3) ci-dessus.
\end{proof}

\begin{lemma}
\label{algebra-lemma-isolated-point-fibre}
\begin{slogan}
Conditions équivalentes pour les points isolés dans les fibres
\end{slogan}
Soit $R \to S$ un homomorphisme d'anneaux de type fini.
Soit $\mathfrak q \subset S$ un idéal premier au-dessus de
$\mathfrak p \subset R$. Soit $F = \Spec(S \otimes_R \kappa(\mathfrak p))$
la fibre de $\Spec(S) \to \Spec(R)$ ; voir la
remarque \ref{algebra-remark-fundamental-diagram}.
Notons $\overline{\mathfrak q} \in F$ le point correspondant à
$\mathfrak q$. Les assertions suivantes sont équivalentes :
\begin{enumerate}
\item $\overline{\mathfrak q}$ est un point isolé de $F$,
\item $S_{\mathfrak q}/\mathfrak pS_{\mathfrak q}$ est fini sur
$\kappa(\mathfrak p)$,
\item il existe $g \in S$, $g \not \in \mathfrak q$, tel que
le seul point de $D(g)$ dont l'image est $\mathfrak p$ soit $\mathfrak q$,
\item $\dim_{\overline{\mathfrak q}}(F) = 0$,
\item $\overline{\mathfrak q}$ est un point fermé de $F$ et
$\dim(S_{\mathfrak q}/\mathfrak pS_{\mathfrak q}) = 0$, et
\item l'extension de corps $\kappa(\mathfrak q)/\kappa(\mathfrak p)$
est finie et $\dim(S_{\mathfrak q}/\mathfrak pS_{\mathfrak q}) = 0$.
\end{enumerate}
\end{lemma}

\begin{proof}
Remarquons que $S_{\mathfrak q}/\mathfrak pS_{\mathfrak q} =
(S \otimes_R \kappa(\mathfrak p))_{\overline{\mathfrak q}}$.
De plus, $S \otimes_R \kappa(\mathfrak p)$ est de type fini sur
$\kappa(\mathfrak p)$.
Ces conditions correspondent exactement à celles du
lemme \ref{algebra-lemma-isolated-point}
pour la $\kappa(\mathfrak p)$-algèbre $S \otimes_R \kappa(\mathfrak p)$
et l'idéal premier $\overline{\mathfrak q}$ ; elles sont donc équivalentes.
\end{proof}

\begin{definition}
\label{algebra-definition-quasi-finite}
Soit $R \to S$ un homomorphisme d'anneaux de type fini.
Soit $\mathfrak q \subset S$ un idéal premier.
\begin{enumerate}
\item Si les conditions équivalentes du lemme \ref{algebra-lemma-isolated-point-fibre}
sont satisfaites, on dit que $R \to S$ est {\it quasi-fini en $\mathfrak q$}.
\item On dit qu'un homomorphisme d'anneaux $A \to B$ est {\it quasi-fini}
s'il est de type fini et quasi-fini en tout idéal premier de $B$.
\end{enumerate}
\end{definition}

\begin{lemma}
\label{algebra-lemma-quasi-finite-above-p}
Soit $R \to S$ un homomorphisme d'anneaux de type fini
et soit $\mathfrak p$ un idéal premier de $R$.
Alors les assertions suivantes sont équivalentes :
\begin{enumerate}
\item $R \to S$ est quasi-fini en tout idéal premier de $S$ au-dessus de $\mathfrak p$,
\item $S \otimes_R \kappa(\mathfrak p)$ est une
$\kappa(\mathfrak p)$-algèbre finie, et
\item $\Spec(S \otimes_R \kappa(\mathfrak p))$ est un ensemble fini.
\end{enumerate}
\end{lemma}

\begin{proof}
La condition (1) signifie que la topologie de
$\Spec(S \otimes_R \kappa(\mathfrak p))$ est discrète (puisque tout point
est ouvert). L'équivalence de (1), (2) et (3) résulte donc du
lemme \ref{algebra-lemma-finite-type-algebra-finite-nr-primes}.
\end{proof}

\begin{lemma}
\label{algebra-lemma-quasi-finite}
Soit $R \to S$ un homomorphisme d'anneaux de type fini. Alors $R \to S$ est quasi-fini
si et seulement si, pour tout idéal premier $\mathfrak p \subset R$, l'anneau
$S \otimes_R \kappa(\mathfrak p)$ est fini sur $\kappa(\mathfrak p)$.
\end{lemma}

\begin{proof}
Cela résulte immédiatement du lemme plus général
\ref{algebra-lemma-quasi-finite-above-p}.
\end{proof}

\begin{lemma}
\label{algebra-lemma-quasi-finite-local}
Soit $R \to S$ un homomorphisme d'anneaux de type fini. Soit $\mathfrak q \subset S$
un idéal premier au-dessus de $\mathfrak p \subset R$. Soient
$f \in R$, $f \not \in \mathfrak p$ et $g \in S$, $g \not \in \mathfrak q$.
Alors $R \to S$ est quasi-fini en $\mathfrak q$ si et seulement si
$R_f \to S_{fg}$ est quasi-fini en $\mathfrak qS_{fg}$.
\end{lemma}

\begin{proof}
La fibre de $\Spec(S_{fg}) \to \Spec(R_f)$ est homéomorphe
à un ouvert de la fibre de $\Spec(S) \to \Spec(R)$.
Le lemme résulte donc de la condition (1) du
lemme \ref{algebra-lemma-isolated-point-fibre}.
\end{proof}

\begin{lemma}
\label{algebra-lemma-four-rings}
Soit
$$
\xymatrix{
S \ar[r] & S' & &
\mathfrak q \ar@{-}[r] & \mathfrak q' \\
R \ar[u] \ar[r] &  R' \ar[u] & &
\mathfrak p \ar@{-}[r] \ar@{-}[u] & \mathfrak p' \ar@{-}[u]
}
$$
un diagramme commutatif d'anneaux, avec les idéaux premiers indiqués.
Supposons que $R \to S$ soit de type fini et que $S \otimes_R R' \to S'$ soit surjectif.
Si $R \to S$ est quasi-fini en $\mathfrak q$, alors
$R' \to S'$ est quasi-fini en $\mathfrak q'$.
\end{lemma}

\begin{proof}
Écrivons $S \otimes_R \kappa(\mathfrak p) = S_1 \times S_2$,
avec $S_1$ fini sur $\kappa(\mathfrak p)$ et tel que
$\mathfrak q$ corresponde à un idéal premier de $S_1$, comme dans le
lemme \ref{algebra-lemma-isolated-point}. Cette décomposition en produit
induit une décomposition en produit correspondante pour toute
$S \otimes_R \kappa(\mathfrak p)$-algèbre. En particulier,
on obtient $S' \otimes_{R'} \kappa(\mathfrak p') = S'_1 \times S'_2$.
Puisque $S \otimes_R R' \to S'$ est surjectif, l'homomorphisme canonique
$(S \otimes_R \kappa(\mathfrak p)) \otimes_{\kappa(\mathfrak p)}
\kappa(\mathfrak p') \to S' \otimes_{R'} \kappa(\mathfrak p')$
est surjectif, et par conséquent
$S_i \otimes_{\kappa(\mathfrak p)} \kappa(\mathfrak p') \to S'_i$
est surjectif. Il s'ensuit que $S'_1$ est fini sur $\kappa(\mathfrak p')$.
L'homomorphisme $S' \otimes_{R'} \kappa(\mathfrak p') \to
\kappa(\mathfrak q')$ se factorise par $S_1'$
(c.-à-d.\ qu'il annule le facteur $S_2'$)
car l'homomorphisme $S \otimes_R \kappa(\mathfrak p) \to
\kappa(\mathfrak q)$ se factorise par $S_1$
(c.-à-d.\ qu'il annule le facteur $S_2$). Ainsi,
$\mathfrak q'$ correspond à un point de
$\Spec(S_1')$ dans la décomposition en union disjointe
de la fibre : $\Spec(S' \otimes_{R'} \kappa(\mathfrak p'))
= \Spec(S_1') \amalg \Spec(S_2')$ ; voir le
lemme \ref{algebra-lemma-spec-product}.
Puisque $S_1'$ est fini sur un corps, c'est un anneau artinien,
et $\Spec(S_1')$ est donc un ensemble discret fini.
(Voir la proposition \ref{algebra-proposition-dimension-zero-ring}.)
On conclut que $\mathfrak q'$ est isolé dans sa fibre,
ce qui donne le résultat.
\end{proof}

\begin{lemma}
\label{algebra-lemma-quasi-finite-composition}
La composée d'homomorphismes d'anneaux quasi-finis est quasi-finie.
\end{lemma}

\begin{proof}
Supposons que $A \to B$ et $B \to C$ soient des homomorphismes d'anneaux quasi-finis. D'après le
lemme \ref{algebra-lemma-compose-finite-type},
on voit que $A \to C$ est de type fini.
Soit $\mathfrak r \subset C$ un idéal premier de $C$ au-dessus de
$\mathfrak q \subset B$ et de $\mathfrak p \subset A$. Puisque $A \to B$ et
$B \to C$ sont quasi-finis en $\mathfrak q$ et $\mathfrak r$ respectivement,
il existe alors $b \in B$ et $c \in C$ tels que $\mathfrak q$ soit
le seul idéal premier de $D(b)$ au-dessus de $\mathfrak p$ et, de même,
$\mathfrak r$ soit le seul idéal premier de $D(c)$ au-dessus de $\mathfrak q$.
Si $c' \in C$ est l'image de $b \in B$, alors $\mathfrak r$ est le seul
idéal premier de $D(cc')$ au-dessus de $\mathfrak p$.
Par conséquent, $A \to C$ est quasi-fini en $\mathfrak r$.
\end{proof}

\begin{lemma}
\label{algebra-lemma-quasi-finite-base-change}
Soit $R \to S$ un homomorphisme d'anneaux de type fini.
Soit $R \to R'$ un homomorphisme d'anneaux quelconque. Posons $S' = R' \otimes_R S$.
\begin{enumerate}
\item L'ensemble
$\{\mathfrak q' \mid R' \to S' \text{ quasi-fini en }\mathfrak q'\}$
est l'image réciproque de l'ensemble correspondant de $\Spec(S)$
par le morphisme canonique $\Spec(S') \to \Spec(S)$.
\item Si $\Spec(R') \to \Spec(R)$ est surjectif,
alors $R \to S$ est quasi-fini si et seulement si $R' \to S'$ est quasi-fini.
\item Tout changement de base d'un homomorphisme d'anneaux quasi-fini est quasi-fini.
\end{enumerate}
\end{lemma}

\begin{proof}
Soit $\mathfrak p' \subset R'$ un idéal premier au-dessus de $\mathfrak p \subset R$.
Alors l'anneau de la fibre $S' \otimes_{R'} \kappa(\mathfrak p')$ est le
changement de base de l'anneau de la fibre $S \otimes_R \kappa(\mathfrak p)$
par l'extension de corps $\kappa(\mathfrak p) \to \kappa(\mathfrak p')$.
La première assertion résulte donc de l'invariance de la dimension
par extension du corps de base
(lemme \ref{algebra-lemma-dimension-at-a-point-preserved-field-extension})
et du lemme \ref{algebra-lemma-isolated-point}.
La stabilité des morphismes quasi-finis par changement de base résulte de
ce qui précède et de la stabilité de la propriété d'être de type fini par changement de base.
La deuxième assertion en résulte
car l'hypothèse implique que, pour tout idéal premier $\mathfrak q \subset S$, on peut
trouver un idéal premier $\mathfrak q' \subset S'$ au-dessus de lui.
\end{proof}

\begin{lemma}
\label{algebra-lemma-quasi-finite-permanence}
Soient $A \to B$ et $B \to C$ des homomorphismes d'anneaux tels que $A \to C$
soit de type fini. Soit $\mathfrak r$ un idéal premier de $C$ au-dessus de
$\mathfrak q \subset B$ et de $\mathfrak p \subset A$.
Si $A \to C$ est quasi-fini en $\mathfrak r$, alors
$B \to C$ est quasi-fini en $\mathfrak r$.
\end{lemma}

\begin{proof}
Remarquons que $B \to C$ est de type fini
(lemme \ref{algebra-lemma-compose-finite-type}),
de sorte que l'énoncé a un sens.
Utilisons la caractérisation (3) du lemme \ref{algebra-lemma-isolated-point-fibre}.
Si $A \to C$ est quasi-fini en $\mathfrak r$, alors
il existe $c \in C$ tel que
$$
\{\mathfrak r' \subset C \text{ au-dessus de }\mathfrak p\} \cap D(c) =
\{\mathfrak{r}\}.
$$
Puisque les idéaux premiers $\mathfrak r' \subset C$ au-dessus de $\mathfrak q$
forment un sous-ensemble des idéaux premiers $\mathfrak r' \subset C$ au-dessus de
$\mathfrak p$, on conclut que $B \to C$ est quasi-fini en $\mathfrak r$.
\end{proof}

\begin{lemma}
\label{algebra-lemma-generically-finite}
Soit $R \to S$ un homomorphisme d'anneaux de type fini.
Soit $\mathfrak p \subset R$ un idéal premier minimal.
Supposons qu'il n'y ait qu'un nombre fini d'idéaux premiers de $S$
au-dessus de $\mathfrak p$. Il existe alors
$g \in R$, $g \not \in \mathfrak p$, tel que le
homomorphisme d'anneaux $R_g \to S_g$ soit fini.
\end{lemma}

\begin{proof}
Soient $x_1, \ldots, x_n$ des générateurs de $S$ sur $R$.
D'après le lemme \ref{algebra-lemma-quasi-finite-above-p}, l'hypothèse signifie que
$S \otimes_R \kappa(\mathfrak p) = S_{\mathfrak p}/\mathfrak pS_{\mathfrak p}$
est une $\kappa(\mathfrak p)$-algèbre finie.
On peut donc trouver des polynômes unitaires
$P_i \in R_{\mathfrak p}[X]$ tels que l'image de $P_i(x_i)$ soit nulle
dans $S_{\mathfrak p}/\mathfrak pS_{\mathfrak p}$.
Puisque $\mathfrak p$ est un idéal premier minimal,
$\mathfrak pR_{\mathfrak p}$ est un idéal localement nilpotent ; voir le
lemme \ref{algebra-lemma-minimal-prime-reduced-ring}.
Par conséquent, $\mathfrak pS_{\mathfrak p}$ est un idéal localement nilpotent ; voir le
lemme \ref{algebra-lemma-locally-nilpotent}. Ainsi,
il existe des $e_i \geq 1$ tels que $P(x_i)^{e_i} = 0$
dans $S_{\mathfrak p}$. Choisissons $g_1 \in R$, $g_1 \not \in \mathfrak p$,
tel que les coefficients de $P_i$ appartiennent à $R[1/g_1]$ pour tout $i$.
Choisissons ensuite $g_2 \in R$, $g_2 \not \in \mathfrak p$,
tel que $P(x_i)^{e_i} = 0$ dans $S_{g_1g_2}$. En posant $g = g_1g_2$,
on obtient le résultat.
\end{proof}







\section{Théorème principal de Zariski}
\label{algebra-section-Zariski}

\noindent
Dans cette section, notre objectif est de démontrer la version algébrique du
théorème principal de Zariski. Ce théorème servira de fondement à de nombreux
développements ultérieurs de la théorie des schémas et de leurs
morphismes dans la suite du projet Champs.

\medskip\noindent
Soit $R \to S$ un homomorphisme d'anneaux de type fini.
Notre objectif dans cette section est de montrer que l'ensemble
des points de $\Spec(S)$ en lesquels cet homomorphisme est quasi-fini
est {\it ouvert} (théorème \ref{algebra-theorem-main-theorem}).
En fait, nous verrons qu'il existe
un homomorphisme d'anneaux fini $R \to S'$ tel qu'en un certain sens le lieu de quasi-finitude
de $S/R$ soit ouvert dans $\Spec(S')$
(mais nous ne le démontrerons pas dans ce chapitre d'algèbre, car nous n'y
développons pas le langage des schémas -- lorsque $R \to S$
est quasi-fini, voir le lemme \ref{algebra-lemma-quasi-finite-open-integral-closure}).
Ces assertions sont assez délicates à démontrer ; nous procédons
par une longue suite de lemmes sur les extensions entières et
finies d'anneaux. On trouvera ces résultats
dans \cite{Henselian} et \cite{Peskine}. Les notes
de Thierry Coquand nous ont également été utiles.

\begin{lemma}
\label{algebra-lemma-make-integral-trivial}
Soit $\varphi : R \to S$ un homomorphisme d'anneaux.
Supposons que $t \in S$ satisfasse la
relation $\varphi(a_0) + \varphi(a_1)t + \ldots + \varphi(a_n) t^n = 0$.
Alors $\varphi(a_n)t$ est entier sur $R$.
\end{lemma}

\begin{proof}
En effet, multiplions l'égalité
$\varphi(a_0) + \varphi(a_1)t + \ldots + \varphi(a_n) t^n = 0$
par $\varphi(a_n)^{n-1}$ et écrivons-la sous la forme
$\varphi(a_0 a_n^{n-1}) +
\varphi(a_1 a_n^{n-2}) (\varphi(a_n)t) +
\ldots +
(\varphi(a_n) t)^n = 0$.
\end{proof}

\noindent
Le lemme suivant est, en un certain sens, le lemme clé de cette
section.

\begin{lemma}
\label{algebra-lemma-make-integral-trick}
Soit $R$ un anneau. Soit $\varphi : R[x] \to S$ un
homomorphisme d'anneaux. Soit $t \in S$.
Supposons (a) que $t$ soit entier sur $R[x]$,
et (b) qu'il existe un polynôme unitaire $p \in R[x]$ tel que
$t \varphi(p) \in \Im(\varphi)$. Alors il
existe $q \in R[x]$ tel que $t - \varphi(q)$
soit entier sur $R$.
\end{lemma}

\begin{proof}
Écrivons $t \varphi(p) = \varphi(r)$ pour un certain $r \in R[x]$.
Par division euclidienne, écrivons $r = qp + r'$ avec
$q, r' \in R[x]$ et $\deg(r') < \deg(p)$. On peut remplacer
$t$ par $t - \varphi(q)$, qui est encore entier sur
$R[x]$, et l'on obtient alors $t \varphi(p) = \varphi(r')$.
Dans l'anneau localisé $S_t$, cette égalité s'écrit
$\varphi(p) - (1/t) \varphi(r') = 0$.
Cela implique que $\varphi(x)$ définit un élément de la
localisation $S_t$ qui est entier sur
$\varphi(R)[1/t] \subset S_t$. D'autre part,
$t$ est entier sur le sous-anneau $\varphi(R)[\varphi(x)] \subset S$.
En combinant ces faits, on conclut que $t$ est entier sur
le sous-anneau $\varphi(R)[1/t] \subset S_t$ ; voir le lemme
\ref{algebra-lemma-integral-transitive}. Autrement dit,
il existe une égalité de la forme
$$
t^d + \sum\nolimits_{i < d}
\left(\sum\nolimits_{j = 0, \ldots, n_i} \varphi(r_{i, j})/t^j\right) t^i = 0
$$
dans $S_t$, où $r_{i, j} \in R$. Cela signifie que
$t^{d + N} +
\sum_{i < d} \sum_{j = 0, \ldots, n_i} \varphi(r_{i, j}) t^{i + N - j} = 0$
dans $S$ pour un certain $N$ assez grand. Autrement dit,
$t$ est entier sur $R$.
\end{proof}

\begin{lemma}
\label{algebra-lemma-combine-lemmas}
Soit $R$ un anneau. Soit $\varphi : R[x] \to S$ un
homomorphisme d'anneaux. Soit $t \in S$. Supposons que $t$ soit entier
sur $R[x]$. Soit $p \in R[x]$, $p = a_0 + a_1x + \ldots +
a_k x^k$ tel que $t \varphi(p) \in \Im(\varphi)$.
Alors il existe $q \in R[x]$ et $n \geq 0$
tels que $\varphi(a_k)^n t - \varphi(q)$ soit entier
sur $R$.
\end{lemma}

\begin{proof}
Soient $R'$ et $S'$ les localisés de $R$ et $S$, respectivement en $a_k$ et en son image.
Soit $\varphi' : R'[x] \to S'$ l'homomorphisme localisé induit par $\varphi$.
Soit $t' \in S'$ l'image de $t$. Posons $p' = p/a_k \in R'[x]$.
Alors $t' \varphi'(p') \in \Im(\varphi')$, puisque $t \varphi(p) \in \Im(\varphi)$.
Comme $p'$ est unitaire, le
lemme \ref{algebra-lemma-make-integral-trick} donne un $q' \in R'[x]$
tel que $t' - \varphi'(q')$ soit entier sur $R'$.
On peut choisir $n \geq 0$ et un élément $q \in R[x]$
tels que $a_k^n q'$ soit l'image de $q$.
Alors $\varphi(a_k)^n t - \varphi(q)$ est un élément de $S$ dont
l'image dans $S'$ est entière sur $R'$.
D'après le lemme \ref{algebra-lemma-integral-closure-localize},
il existe $m \geq 0$ tel que
$\varphi(a_k)^m(\varphi(a_k)^n t - \varphi(q))$ soit entier sur $R$.
Ainsi, $\varphi(a_k)^{m + n}t - \varphi(a_k^m q)$
est entier sur $R$, ce qui donne le résultat.
\end{proof}

\begin{situation}
\label{algebra-situation-one-transcendental-element}
Soit $R$ un anneau.
Soit $\varphi : R[x] \to S$ un homomorphisme d'anneaux fini.
Soit
$$
J = \{ g \in S \mid gS \subset \Im(\varphi)\}
$$
l'``idéal conducteur'' de $\varphi$.
Supposons que le sous-anneau $\varphi(R) \subset S$ soit intégralement clos dans $S$.
\end{situation}

\begin{lemma}
\label{algebra-lemma-leading-coefficient-in-J}
Plaçons-nous dans la situation \ref{algebra-situation-one-transcendental-element}.
Supposons que $u \in S$, que $a_0, \ldots, a_k \in R$ et que
$u \varphi(a_0 + a_1x + \ldots + a_k x^k) \in J$.
Alors il existe $m \geq 0$ tel que
$u \varphi(a_k)^m \in J$.
\end{lemma}

\begin{proof}
Supposons que $S$ soit engendré par $t_1, \ldots, t_n$
comme $R[x]$-module. Dans ce cas,
$J = \{ g \in S \mid gt_i \in \Im(\varphi)\text{ pour tout }i\}$.
Notons que chaque élément $u t_i$ est entier sur
$R[x]$ ; voir le lemme \ref{algebra-lemma-finite-is-integral}.
On a $\varphi(a_0 + a_1x + \ldots + a_k x^k) u t_i \in
\Im(\varphi)$. D'après le lemme \ref{algebra-lemma-combine-lemmas}, pour
chaque $i$, il existe un entier $n_i$ et un élément
$q_i \in R[x]$ tels que $\varphi(a_k^{n_i}) u t_i - \varphi(q_i)$
soit entier sur $R$. Par hypothèse, cet élément appartient à $\varphi(R)$,
et donc $\varphi(a_k^{n_i}) u t_i \in \Im(\varphi)$.
Il s'ensuit que $m = \max\{n_1, \ldots, n_n\}$ convient.
\end{proof}

\begin{lemma}
\label{algebra-lemma-all-coefficients-in-J}
Plaçons-nous dans la situation \ref{algebra-situation-one-transcendental-element}.
Supposons que $u \in S$, que $a_0, \ldots, a_k \in R$ et que
$u \varphi(a_0 + a_1x + \ldots + a_k x^k) \in \sqrt{J}$.
Alors $u \varphi(a_i) \in \sqrt{J}$ pour tout $i$.
\end{lemma}

\begin{proof}
Sous les hypothèses du lemme, on a
$u^n \varphi(a_0 + a_1x + \ldots + a_k x^k)^n \in J$ pour
un certain $n \geq 1$. D'après le lemme \ref{algebra-lemma-leading-coefficient-in-J},
on en déduit que $u^n \varphi(a_k^{nm}) \in J$ pour un certain $m \geq 1$.
Ainsi, $u \varphi(a_k) \in \sqrt{J}$, et donc
$u \varphi(a_0 + a_1x + \ldots + a_k x^k) - u \varphi(a_k x^k) =
u \varphi(a_0 + a_1x + \ldots + a_{k-1} x^{k-1}) \in \sqrt{J}$.
On conclut par récurrence sur $k$.
\end{proof}

\noindent
Ce lemme suggère la définition suivante.

\begin{definition}
\label{algebra-definition-strongly-transcendental}
Étant donnés une inclusion d'anneaux $R \subset S$ et
un élément $x \in S$, on dit que $x$ est
{\it fortement transcendant sur $R$} si,
dès que $u(a_0 + a_1 x + \ldots + a_k x^k) = 0$
avec $u \in S$ et $a_i \in R$, alors
on a $ua_i = 0$ pour tout $i$.
\end{definition}

\noindent
Notons que, si $S$ est un anneau intègre, cela revient à
dire que $x$, considéré comme élément du corps des fractions de
$S$, est transcendant sur le corps des fractions de $R$.

\begin{lemma}
\label{algebra-lemma-reduced-strongly-transcendental-minimal-prime}
Supposons que $R \subset S$ soit une inclusion d'anneaux réduits
et que $x \in S$ soit fortement transcendant sur $R$.
Soit $\mathfrak q \subset S$ un idéal premier minimal
et soit $\mathfrak p = R \cap \mathfrak q$.
Alors l'image de $x$ dans $S/\mathfrak q$ est fortement
transcendante sur le sous-anneau $R/\mathfrak p$.
\end{lemma}

\begin{proof}
Supposons que $u(a_0 + a_1x + \ldots + a_k x^k) \in \mathfrak q$.
D'après le lemme \ref{algebra-lemma-minimal-prime-reduced-ring},
l'anneau local $S_{\mathfrak q}$ est un corps,
et donc $u(a_0 + a_1x + \ldots + a_k x^k) $ est nul
dans $S_{\mathfrak q}$. Ainsi, $uu'(a_0 + a_1x + \ldots + a_k x^k) = 0$
pour un certain $u' \in S$, $u' \not\in \mathfrak q$.
Puisque $x$ est fortement transcendant sur $R$, on obtient
$uu'a_i = 0$ pour tout $i$. Cela implique à son tour
que $ua_i \in \mathfrak q$.
\end{proof}

\begin{lemma}
\label{algebra-lemma-domains-transcendental-not-quasi-finite}
Supposons que $R\subset S$ soit une inclusion d'anneaux intègres et
soit $x \in S$. Supposons que $x$ soit (fortement) transcendant sur $R$
et que $S$ soit fini sur $R[x]$. Alors l'homomorphisme $R \to S$ n'est
quasi-fini en aucun idéal premier de $S$.
\end{lemma}

\begin{proof}
Commençons par le cas où $R$ est normal ; voir la
définition \ref{algebra-definition-ring-normal}.
D'après le lemme \ref{algebra-lemma-polynomial-ring-normal},
on voit que $R[x]$ est normal.
Prenons un idéal premier $\mathfrak q \subset S$
et posons $\mathfrak p = R \cap \mathfrak q$.
Supposons que l'extension $\kappa(\mathfrak p)
\subset \kappa(\mathfrak q)$ soit finie.
Ce serait le cas si l'homomorphisme $R \to S$ était
quasi-fini en $\mathfrak q$.
Soit $\mathfrak r = R[x] \cap \mathfrak q$.
Alors, puisque $\kappa(\mathfrak p)
\subset \kappa(\mathfrak r) \subset \kappa(\mathfrak q)$,
on voit que l'extension $\kappa(\mathfrak p)
\subset \kappa(\mathfrak r)$ est également finie.
Ainsi, l'inclusion $\mathfrak r \supset \mathfrak p R[x]$
est stricte. Par le théorème de descente appliqué à $R[x] \subset S$,
voir la proposition \ref{algebra-proposition-going-down-normal-integral},
on trouve un idéal premier $\mathfrak q' \subset \mathfrak q$
au-dessus de l'idéal premier $\mathfrak pR[x]$. Par conséquent,
la fibre $\Spec(S \otimes_R \kappa(\mathfrak p))$
contient un point distinct de $\mathfrak q$,
à savoir $\mathfrak q'$, dont l'adhérence contient $\mathfrak q$ ; ainsi,
$\mathfrak q$ n'est pas isolé dans sa fibre.

\medskip\noindent
Si $R$ n'est pas normal, introduisons une chaîne $R \subset R' \subset K$
dans laquelle $R'$ désigne la clôture intégrale de $R$ dans son corps des fractions
$K$. De même, introduisons $S \subset S' \subset L$, où $S'$ est le sous-anneau du
corps des fractions $L$ de $S$ engendré par $R'$ et
$S$. Notons que, par construction, l'homomorphisme $S \otimes_R R'
\to S'$ est surjectif. Cela implique que $R'[x] \subset S'$
est fini. De plus, l'inclusion $S \subset S'$
induit sur $\Spec$ un morphisme surjectif ; voir le
lemme \ref{algebra-lemma-integral-overring-surjective}.
On conclut à l'aide du lemme \ref{algebra-lemma-four-rings} et du cas normal
que nous venons de traiter.
\end{proof}

\begin{lemma}
\label{algebra-lemma-reduced-strongly-transcendental-not-quasi-finite}
Supposons que $R \subset S$ soit une inclusion d'anneaux réduits.
Supposons que $x \in S$ soit fortement transcendant sur $R$
et que $S$ soit fini sur $R[x]$. Alors l'homomorphisme $R \to S$ n'est
quasi-fini en aucun idéal premier de $S$.
\end{lemma}

\begin{proof}
Soit $\mathfrak q \subset S$ un idéal premier quelconque.
Choisissons un idéal premier minimal $\mathfrak q' \subset \mathfrak q$.
D'après les lemmes
\ref{algebra-lemma-reduced-strongly-transcendental-minimal-prime} et
\ref{algebra-lemma-domains-transcendental-not-quasi-finite},
l'extension $R/(R \cap \mathfrak q') \subset
S/\mathfrak q'$ n'est pas quasi-finie en l'idéal premier correspondant
à $\mathfrak q$. D'après le lemme \ref{algebra-lemma-four-rings},
l'homomorphisme $R \to S$ n'est pas quasi-fini
en $\mathfrak q$.
\end{proof}

\begin{lemma}
\label{algebra-lemma-quasi-finite-monogenic}
Soit $R$ un anneau. Soit $S = R[x]/I$.
Soit $\mathfrak q \subset S$ un idéal premier.
Supposons que l'homomorphisme $R \to S$ soit quasi-fini en $\mathfrak q$.
Soit $S' \subset S$ la clôture intégrale de $R$ dans $S$.
Alors il existe un élément
$g \in S'$, $g \not\in \mathfrak q$ tel que
$S'_g \cong S_g$.
\end{lemma}

\begin{proof}
Soit $\mathfrak p$ l'image de $\mathfrak q$ dans $\Spec(R)$.
Il existe $f \in I$, $f = a_nx^n + \ldots + a_0$ tel que
$a_i \not \in \mathfrak p$ pour un certain $i$. En effet, sinon, l'anneau de la fibre
$S \otimes_R \kappa(\mathfrak p)$ serait $\kappa(\mathfrak p)[x]$
et l'homomorphisme ne serait quasi-fini en aucun idéal premier situé
au-dessus de $\mathfrak p$. On en déduit qu'il existe une relation
$b_m x^m + \ldots + b_0 = 0$ avec $b_j \in S'$, $j = 0, \ldots, m$
et $b_j \not \in \mathfrak q \cap S'$ pour un certain $j$.
Nous démontrons le lemme par récurrence sur $m$. Le cas initial
$m = 0$ ne se présente pas (car les assertions $b_0 = 0$ et
$b_0 \not \in \mathfrak q$ sont contradictoires).

\medskip\noindent
Supposons que $b_m \not \in \mathfrak q$. Dans ce cas, $x$ est entier
sur $S'_{b_m}$ ; en fait, $b_mx \in S'$ d'après le
lemme \ref{algebra-lemma-make-integral-trivial}.
Ainsi, l'homomorphisme injectif $S'_{b_m} \to S_{b_m}$ est également surjectif, donc
c'est l'isomorphisme voulu.

\medskip\noindent
Supposons que $b_m \in \mathfrak q$. Dans ce cas, on a $b_mx \in S'$ d'après le
lemme \ref{algebra-lemma-make-integral-trivial}.
Posons $b'_{m - 1} = b_mx + b_{m - 1}$. Alors
$$
b'_{m - 1}x^{m - 1} + b_{m - 2}x^{m - 2} + \ldots + b_0 = 0
$$
Puisque $b'_{m - 1}$ est congru à $b_{m - 1}$ modulo $S' \cap \mathfrak q$,
on voit que l'un des éléments
$b'_{m - 1}, b_{m - 2}, \ldots, b_0$ n'appartient toujours pas à $S' \cap \mathfrak q$.
On conclut donc par récurrence sur $m$.
\end{proof}

\begin{theorem}[Théorème principal de Zariski]
\label{algebra-theorem-main-theorem}
Soit $R$ un anneau. Soit $R \to S$ l'homomorphisme structural d'une $R$-algèbre de type fini.
Soit $S' \subset S$ la clôture intégrale de $R$ dans $S$.
Soit $\mathfrak q \subset S$ un idéal premier de $S$.
Si l'homomorphisme $R \to S$ est quasi-fini en $\mathfrak q$, alors
il existe $g \in S'$, $g \not \in \mathfrak q$
tel que $S'_g \cong S_g$.
\end{theorem}

\begin{proof}
Il existe un nombre fini d'éléments
$x_1, \ldots, x_n \in S$ tels que $S$ soit fini
sur la sous-$R$-algèbre engendrée par $x_1, \ldots, x_n$. (Par
exemple, on peut prendre des générateurs de $S$ sur $R$.) Nous démontrons le théorème
par récurrence sur le plus petit entier $n$ ayant cette propriété.

\medskip\noindent
Le cas $n = 0$ est trivial, car alors $S' = S$ ;
voir le lemme \ref{algebra-lemma-finite-is-integral}.

\medskip\noindent
Supposons que $n = 1$. On peut remplacer $R$ par sa clôture intégrale dans $S$
(le lemme \ref{algebra-lemma-quasi-finite-permanence} garantit que l'homomorphisme $R \to S$
reste quasi-fini en $\mathfrak q$). On peut donc supposer que
l'inclusion $R \subset S$ identifie l'anneau de base à un sous-anneau intégralement clos dans $S$, autrement dit que $R = S'$.
Considérons l'homomorphisme $\varphi : R[x] \to S$, $x \mapsto x_1$.
(Nous verrons ci-dessous que $\varphi$ n'est pas injectif.)
Par hypothèse, $\varphi$ est fini. Nous sommes donc dans la situation
\ref{algebra-situation-one-transcendental-element}.
Soit $J \subset S$ l'``idéal conducteur'' défini
dans la situation \ref{algebra-situation-one-transcendental-element}.
Considérons le diagramme
$$
\xymatrix{
R[x] \ar[r] & S \ar[r] & S/\sqrt{J} & R/(R \cap \sqrt{J})[x] \ar[l]
\\
& R \ar[lu] \ar[r] \ar[u] & R/(R \cap \sqrt{J}) \ar[u] \ar[ru] &
}
$$
D'après le lemme \ref{algebra-lemma-all-coefficients-in-J},
l'image de $x$ dans le quotient $S/\sqrt{J}$
est fortement transcendante sur $R/ (R \cap \sqrt{J})$.
Ainsi, d'après le lemme \ref{algebra-lemma-reduced-strongly-transcendental-not-quasi-finite},
l'homomorphisme d'anneaux $R/ (R \cap \sqrt{J}) \to S/\sqrt{J}$
n'est quasi-fini en aucun idéal premier de $S/\sqrt{J}$.
D'après le lemme \ref{algebra-lemma-four-rings}, on en déduit que $\mathfrak q$
n'appartient pas à $V(J) \subset \Spec(S)$.
Il existe donc un élément $s \in J$
tel que $s \not\in \mathfrak q$. Par définition de $J$, on peut écrire
$s = \varphi(f)$ pour un polynôme $f \in R[x]$.
Soit $I = \Ker(\varphi : R[x] \to S)$. Puisque $\varphi(f) \in J$,
on obtient $(R[x]/I)_f \cong S_{\varphi(f)}$. De plus, $s \not \in \mathfrak q$
signifie que $f \not \in \varphi^{-1}(\mathfrak q)$. Ainsi,
$\varphi^{-1}(\mathfrak q)/I$ est un idéal premier de $R[x]/I$
en lequel l'homomorphisme $R \to R[x]/I$ est quasi-fini ; voir le
lemme \ref{algebra-lemma-quasi-finite-local}. Notons que $R$ est intégralement clos
dans $R[x]/I$ puisque $R$ est intégralement clos dans $S$. D'après le
lemme \ref{algebra-lemma-quasi-finite-monogenic},
il existe un élément $h \in R$, $h \not \in R \cap \mathfrak q$
tel que $R_h \cong (R[x]/I)_h$. Ainsi,
$(R[x]/I)_{fh} = S_{\varphi(fh)}$ est isomorphe à une
localisation principale $R_{h'}$ de $R$ pour un certain
$h' \in R$, $h' \not \in \mathfrak q$.

\medskip\noindent
Supposons que $n > 1$. Considérons le sous-anneau $R' \subset S$
qui est la clôture intégrale de $R[x_1, \ldots, x_{n-1}]$
dans $S$. D'après le lemme \ref{algebra-lemma-quasi-finite-permanence}, l'extension
$S/R'$ est quasi-finie en $\mathfrak q$. Notons également
que $S$ est fini sur $R'[x_n]$.
D'après le cas $n = 1$ ci-dessus, il existe $g' \in R'$
tel que $g' \not \in \mathfrak q$ et que
$(R')_{g'} \cong S_{g'}$. À ce stade, on ne peut pas
appliquer l'hypothèse de récurrence à $R \to R'$, car $R'$ peut ne pas être de type fini sur $R$.
Puisque $S$ est de type fini sur $R$, on en déduit en particulier
que $(R')_{g'}$ est de type fini sur $R$. Disons que
les éléments $g'$ et $y_1/(g')^{n_1}, \ldots, y_N/(g')^{n_N}$,
où $y_i \in R'$,
engendrent $(R')_{g'}$ sur $R$. Soit $R''$ la sous-$R$-algèbre
de $R'$ engendrée par $x_1, \ldots, x_{n-1}, y_1, \ldots, y_N, g'$.
Elle vérifie $(R'')_{g'} \cong S_{g'}$. La surjectivité
résulte du choix des $y_i$ ; l'injectivité résulte de
$R'' \subset R'$ et de l'exactitude de la localisation. Notons que
$R''$ est fini sur $R[x_1, \ldots, x_{n-1}]$ en raison
du choix de $R'$ ; voir le lemme \ref{algebra-lemma-characterize-integral}.
Soit $\mathfrak q'' = R'' \cap \mathfrak q$. Puisque
$(R'')_{\mathfrak q''} = S_{\mathfrak q}$, on voit que
l'homomorphisme $R \to R''$ est quasi-fini en $\mathfrak q''$ ; voir le
lemme \ref{algebra-lemma-isolated-point-fibre}.
Appliquons l'hypothèse de récurrence à $R \to R''$, à $\mathfrak q''$
et à $x_1, \ldots, x_{n-1} \in R''$ ; on obtient un sous-anneau
$R''' \subset R''$ entier sur $R$ et un
élément $g'' \in R'''$, $g'' \not \in \mathfrak q''$
tel que $(R''')_{g''} \cong (R'')_{g''}$. Écrivons l'image de $g'$ dans
$(R'')_{g''}$ sous la forme $g'''/(g'')^n$ pour un certain $g''' \in  R'''$.
Posons $g = g''g''' \in R'''$. Il est alors clair que $g \not\in
\mathfrak q$ et que $(R''')_g \cong S_g$. Comme, par construction,
on a $R''' \subset S'$, on a également $S'_g \cong S_g$, ce qui donne le résultat.
\end{proof}

\begin{lemma}
\label{algebra-lemma-quasi-finite-open}
Soit $R \to S$ un homomorphisme d'anneaux de type fini.
L'ensemble des points $\mathfrak q$ de $\Spec(S)$ en lesquels
$S/R$ est quasi-fini est ouvert dans $\Spec(S)$.
\end{lemma}

\begin{proof}
Soit $\mathfrak q \subset S$ un idéal premier en lequel l'homomorphisme d'anneaux
est quasi-fini. D'après le théorème \ref{algebra-theorem-main-theorem},
il existe un homomorphisme d'anneaux entier $R \to S'$, $S' \subset S$,
et un élément $g \in S'$, $g\not \in \mathfrak q$ tel que
$S'_g \cong S_g$. Puisque $S$, et donc $S_g$, sont de type fini
sur $R$, on peut trouver un nombre fini d'éléments
$y_1, \ldots, y_N$ de $S'$ tels que $S''_g \cong S_g$,
où $S'' \subset S'$ est la sous-$R$-algèbre engendrée
par $g, y_1, \ldots, y_N$. Puisque $S''$ est fini sur $R$
(voir le lemme \ref{algebra-lemma-characterize-integral}), on voit que
$S''$ est quasi-fini sur $R$ (voir le lemme \ref{algebra-lemma-quasi-finite}).
On voit facilement que cela implique que $S''_g$ est quasi-fini sur $R$,
par exemple parce que la propriété d'être quasi-fini en un idéal premier ne dépend
que de l'anneau local en cet idéal. Ainsi, $S_g$ est quasi-fini
sur $R$. Le même raisonnement montre que l'homomorphisme $R \to S$ est quasi-fini
en tout idéal premier de $S$ appartenant à $D(g)$.
\end{proof}

\begin{lemma}
\label{algebra-lemma-quasi-finite-open-integral-closure}
Soit $R \to S$ un homomorphisme d'anneaux de type fini.
Supposons que $S$ soit quasi-fini sur $R$.
Soit $S' \subset S$ la clôture intégrale de $R$ dans $S$. Alors
\begin{enumerate}
\item Le morphisme $\Spec(S) \to \Spec(S')$ est un homéomorphisme
sur un ouvert,
\item si $g \in S'$ et si $D(g)$ est contenu dans l'image
du morphisme, alors $S'_g \cong S_g$, et
\item il existe une $R$-algèbre finie $S'' \subset S'$
telle que (1) et (2) soient vérifiées pour l'homomorphisme d'anneaux
$S'' \to S$.
\end{enumerate}
\end{lemma}

\begin{proof}
Puisque $S/R$ est quasi-fini, on peut appliquer le
théorème \ref{algebra-theorem-main-theorem} à
chaque point $\mathfrak q$ de $\Spec(S)$.
Comme $\Spec(S)$ est quasi-compact (voir le
lemme \ref{algebra-lemma-quasi-compact}), on peut choisir
un nombre fini de $g_i \in S'$, $i = 1, \ldots, n$
tels que $S'_{g_i} = S_{g_i}$ et que
$g_1, \ldots, g_n$ engendrent l'idéal unité de $S$
(autrement dit, les ouverts principaux de $\Spec(S)$ associés
à $g_1, \ldots, g_n$ recouvrent tout $\Spec(S)$).

\medskip\noindent
Supposons que $D(g) \subset \Spec(S')$
soit contenu dans l'image. Alors $D(g) \subset \bigcup D(g_i)$.
Autrement dit, $g_1, \ldots, g_n$ engendrent l'idéal unité de
$S'_g$. Notons que $S'_{gg_i} \cong S_{gg_i}$ par notre choix
des $g_i$. Par conséquent, $S'_g \cong S_g$ d'après le lemme \ref{algebra-lemma-cover}.

\medskip\noindent
Construisons une algèbre finie $S'' \subset S'$ comme
au point (3). Pour cela, notons que chaque $S'_{g_i} \cong S_{g_i}$
est une $R$-algèbre de type fini. Pour chaque $i$, choisissons
des éléments $y_{ij} \in S'$ tels que chaque
$S'_{g_i}$ soit engendrée comme $R$-algèbre par $1/g_i$
et les éléments $y_{ij}$. Posons alors $S''$
égal à la sous-$R$-algèbre de $S'$ engendrée par tous les $g_i$
et tous les $y_{ij}$. Nous omettons les détails.
\end{proof}




\section{Applications du théorème principal de Zariski}
\label{algebra-section-apply-main-theorem}

\noindent
Voici une application immédiate qui caractérise les homomorphismes finis
parmi les homomorphismes quasi-finis d'anneaux semi-locaux de dimension $1$ :
ce sont ceux pour lesquels l'égalité est toujours vérifiée dans la
formule du lemme \ref{algebra-lemma-finite-extension-dim-1}.

\begin{lemma}
\label{algebra-lemma-quasi-finite-extension-dim-1}
Soit $A \subset B$ une extension d'anneaux intègres. Supposons que
\begin{enumerate}
\item $A$ soit un anneau local noethérien de dimension $1$,
\item l'homomorphisme $A \to B$ soit de type fini, et
\item l'extension induite $L/K$ des corps des fractions soit finie.
\end{enumerate}
Alors $B$ est semi-local.
Soit $x \in \mathfrak m_A$, $x \not = 0$.
Soient $\mathfrak m_i$, $i = 1, \ldots, n$,
les idéaux maximaux de $B$.
Alors
$$
[L : K]\text{ord}_A(x)
\geq
\sum\nolimits_i
[\kappa(\mathfrak m_i) : \kappa(\mathfrak m_A)]
\text{ord}_{B_{\mathfrak m_i}}(x)
$$
où $\text{ord}$ est défini comme dans la définition \ref{algebra-definition-ord}.
On a égalité si et seulement si l'homomorphisme $A \to B$ est fini.
\end{lemma}

\begin{proof}
L'anneau $B$ est semi-local d'après le lemme \ref{algebra-lemma-finite-in-codim-1}.
Soit $B'$ la clôture intégrale de $A$ dans $B$. D'après le
lemme \ref{algebra-lemma-quasi-finite-open-integral-closure},
on peut trouver une sous-$A$-algèbre finie $C \subset B'$ telle que,
en posant $\mathfrak n_i = C \cap \mathfrak m_i$, on ait
$C_{\mathfrak n_i} \cong B_{\mathfrak m_i}$ et que les idéaux premiers
$\mathfrak n_1, \ldots, \mathfrak n_n$ soient deux à deux distincts.
L'anneau $C$ est semi-local d'après le lemme \ref{algebra-lemma-finite-in-codim-1}.
Soient $\mathfrak p_j$, $j = 1, \ldots, m$, les autres idéaux
maximaux de $C$ (les ``points manquants''). D'après le
lemme \ref{algebra-lemma-finite-extension-dim-1}, on a
$$
\text{ord}_A(x^{[L : K]}) =
\sum\nolimits_i
[\kappa(\mathfrak n_i) : \kappa(\mathfrak m_A)]
\text{ord}_{C_{\mathfrak n_i}}(x)
+
\sum\nolimits_j
[\kappa(\mathfrak p_j) : \kappa(\mathfrak m_A)]
\text{ord}_{C_{\mathfrak p_j}}(x)
$$
d'où l'inégalité. En cas d'égalité, on en déduit que
$m = 0$ (il n'y a pas de ``points manquants''). Ainsi, $C \subset B$ est une inclusion
d'anneaux semi-locaux qui induit une bijection sur les idéaux maximaux et
un isomorphisme sur toutes les localisations aux idéaux maximaux. Ainsi, si $b \in B$,
alors $I = \{x \in C \mid xb \in C\}$ est un idéal de $C$ qui n'est
contenu dans aucun idéal maximal de $C$ ; par conséquent, $I = C$,
donc $b \in C$. Ainsi, $B = C$ et $B$ est fini sur $A$.
\end{proof}

\noindent
Voici une application plus classique du théorème principal de Zariski à la
structure des homomorphismes locaux d'anneaux locaux.

\begin{lemma}
\label{algebra-lemma-essentially-finite-type-fibre-dim-zero}
Soit $(R, \mathfrak m_R) \to (S, \mathfrak m_S)$ un homomorphisme local
d'anneaux locaux. Supposons que
\begin{enumerate}
\item $R \to S$ soit essentiellement de type fini,
\item l'extension $\kappa(\mathfrak m_R) \subset \kappa(\mathfrak m_S)$ soit finie, et
\item $\dim(S/\mathfrak m_RS) = 0$.
\end{enumerate}
Alors $S$ est la localisation d'une $R$-algèbre finie.
\end{lemma}

\begin{proof}
Soit $S'$ une $R$-algèbre de type fini telle que $S = S'_{\mathfrak q'}$
pour un certain idéal premier $\mathfrak q'$ de $S'$. D'après la
définition \ref{algebra-definition-quasi-finite},
on voit que l'homomorphisme $R \to S'$ est quasi-fini en $\mathfrak q'$.
Après avoir remplacé $S'$ par $S'_{g'}$ pour un certain
$g' \in S'$, $g' \not \in \mathfrak q'$, on peut supposer que $R \to S'$ est
quasi-fini ; voir le
lemme \ref{algebra-lemma-quasi-finite-open}.
Alors, d'après le
lemme \ref{algebra-lemma-quasi-finite-open-integral-closure},
il existe une $R$-algèbre finie $S''$ et des éléments
$g' \in S'$, $g' \not \in \mathfrak q'$ et $g'' \in S''$
tels que $S'_{g'} \cong S''_{g''}$ comme $R$-algèbres.
Cela démontre le lemme.
\end{proof}

\begin{lemma}
\label{algebra-lemma-completion-at-quasi-finite-prime}
Soit $R \to S$ un homomorphisme d'anneaux et soit $\mathfrak q$ un idéal premier de $S$
situé au-dessus d'un idéal premier $\mathfrak p$ de $R$. Si
\begin{enumerate}
\item $R$ est noethérien,
\item l'homomorphisme $R \to S$ est de type fini, et
\item l'homomorphisme $R \to S$ est quasi-fini en $\mathfrak q$,
\end{enumerate}
alors $R_\mathfrak p^\wedge \otimes_R S = S_\mathfrak q^\wedge \times B$
pour une certaine $R_\mathfrak p^\wedge$-algèbre $B$.
\end{lemma}

\begin{proof}
Il existe une $R$-algèbre finie $S' \subset S$ et un élément
$g \in S'$, $g \not \in \mathfrak q' = S' \cap \mathfrak q$
tels que $S'_g = S_g$ et, en particulier,
$S'_{\mathfrak q'} = S_\mathfrak q$ ; voir le
lemme \ref{algebra-lemma-quasi-finite-open-integral-closure}.
On a
$$
R_\mathfrak p^\wedge \otimes_R S' = (S'_{\mathfrak q'})^\wedge \times B'
$$
d'après le lemme \ref{algebra-lemma-completion-finite-extension}. Remarquons que, dans cette
décomposition en produit, $g$ s'envoie sur un couple $(u, b')$ où
$u \in (S'_{\mathfrak q'})^\wedge$ est une unité puisque $g \not \in \mathfrak q'$.
La décomposition en produit de $R_\mathfrak p^\wedge \otimes_R S'$
induit une décomposition en produit
$$
R_\mathfrak p^\wedge \otimes_R S = A \times B
$$
Puisque $S'_g = S_g$, on a également
$(R_\mathfrak p^\wedge \otimes_R S')_g = (R_\mathfrak p^\wedge \otimes_R S)_g$,
et puisque $g \mapsto (u, b')$, où $u$ est une unité, on voit que
$(S'_{\mathfrak q'})^\wedge = A$. Comme l'isomorphisme
$S'_{\mathfrak q'} = S_\mathfrak q$ induit un isomorphisme des
complétés, on obtient également $A = S_\mathfrak q^\wedge$.
La démonstration est terminée, à ceci près qu'il faut vérifier
que l'homomorphisme induit
$\phi : R_\mathfrak p^\wedge \otimes_R S \to A = S_\mathfrak q^\wedge$
est l'homomorphisme naturel. Pour les éléments de la forme $x \otimes 1$,
où $x \in R_\mathfrak p^\wedge$, c'est clair puisque l'homomorphisme naturel
$R_\mathfrak p^\wedge \to S_\mathfrak q^\wedge$ se factorise par
$(S'_{\mathfrak q'})^\wedge$. Pour les éléments de la forme $1 \otimes y$,
où $y \in S$, on peut remarquer que, pour un certain $n \geq 1$, l'élément
$g^ny$ est l'image d'un certain $y' \in S'$. Ainsi, $\phi(1 \otimes g^ny)$
est l'image de $y'$ par la composée
$S' \to (S'_{\mathfrak q'})^\wedge \to S_\mathfrak q^\wedge$. Cette image est
égale à celle de $g^ny$ par l'homomorphisme $S \to S_\mathfrak q^\wedge$.
Puisque $g$ s'envoie sur une unité,
il s'ensuit que $\phi(1 \otimes y)$ est bien l'image de $y$
par l'homomorphisme $S \to S_\mathfrak q^\wedge$.
\end{proof}




































\section{Dimension des fibres}
\label{algebra-section-dimension-fibres}

\noindent
Nous étudions le comportement des dimensions des fibres à l'aide du
théorème principal de Zariski. Rappelons que nous avons défini la
dimension $\dim_x(X)$ d'un espace topologique $X$ en un point $x$
dans Topologie, définition \ref{topology-definition-Krull}.

\begin{definition}
\label{algebra-definition-relative-dimension}
Supposons que $R \to S$ soit de type fini, et soit
$\mathfrak q \subset S$ un idéal premier au-dessus d'un idéal premier
$\mathfrak p$ de $R$.
Nous définissons la {\it dimension relative
de $S/R$ en $\mathfrak q$}, notée
$\dim_{\mathfrak q}(S/R)$, comme la dimension
de $\Spec(S \otimes_R \kappa(\mathfrak p))$
au point correspondant à $\mathfrak q$. Nous désignons par
$\dim(S/R)$ la borne supérieure des $\dim_{\mathfrak q}(S/R)$
lorsque $\mathfrak q$ parcourt tous les idéaux premiers. On l'appelle la
{\it dimension relative de} $S/R$.
\end{definition}

\noindent
En particulier, $R \to S$ est quasi-fini en $\mathfrak q$ si et
seulement si $\dim_{\mathfrak q}(S/R) = 0$. Le lemme suivant
est essentiellement une reformulation du théorème principal de Zariski.

\begin{lemma}
\label{algebra-lemma-quasi-finite-over-polynomial-algebra}
Soit $R \to S$ un homomorphisme d'anneaux de type fini.
Soit $\mathfrak q \subset S$ un idéal premier.
Supposons que $\dim_{\mathfrak q}(S/R) = n$.
Il existe un élément $g \in S$, $g \not\in \mathfrak q$
tel que $S_g$ soit quasi-fini sur une
algèbre de polynômes $R[t_1, \ldots, t_n]$.
\end{lemma}

\begin{proof}
L'anneau $\overline{S} = S \otimes_R \kappa(\mathfrak p)$ est
de type fini sur $\kappa(\mathfrak p)$.
Soit $\overline{\mathfrak q}$ l'idéal premier de $\overline{S}$
correspondant à $\mathfrak q$.
Par définition de
la dimension d'un espace topologique en un point, il existe
un ouvert $U \subset \Spec(\overline{S})$ tel que
$\overline{q} \in U$ et $\dim(U) = n$.
Puisque la topologie sur $\Spec(\overline{S})$ est
induite par la topologie sur $\Spec(S)$ (voir la
remarque \ref{algebra-remark-fundamental-diagram}), nous pouvons trouver
un élément $g \in S$, $g \not \in \mathfrak q$, dont l'image
$\overline{g} \in \overline{S}$ est telle que
$D(\overline{g}) \subset U$.
Ainsi, après remplacement de $S$ par $S_g$, on voit que
$\dim(\overline{S}) = n$.

\medskip\noindent
Choisissons ensuite des générateurs $x_1, \ldots, x_N$ de $S$ comme $R$-algèbre. D'après le
lemme \ref{algebra-lemma-Noether-normalization},
il existe des éléments $y_1, \ldots, y_n$ dans la sous-algèbre sur $\mathbf{Z}$ de $S$
engendrée par $x_1, \ldots, x_N$ tels que l'homomorphisme
$R[t_1, \ldots, t_n] \to S$, $t_i \mapsto y_i$ a pour propriété
que $\kappa(\mathfrak p)[t_1\ldots, t_n] \to \overline{S}$
est fini. En particulier, $S$ est quasi-fini sur $R[t_1, \ldots, t_n]$
en $\mathfrak q$. Par conséquent, d'après le lemme \ref{algebra-lemma-quasi-finite-open},
nous pouvons remplacer $S$ par $S_g$ pour un certain $g\in S$, $g \not \in \mathfrak q$
tel que $R[t_1, \ldots, t_n] \to S$ soit quasi-fini.
\end{proof}

\begin{lemma}
\label{algebra-lemma-refined-quasi-finite-over-polynomial-algebra}
Soit $R \to S$ un homomorphisme d'anneaux. Soit $\mathfrak q \subset S$
un idéal premier au-dessus de l'idéal premier $\mathfrak p$ de $R$.
Supposons que
\begin{enumerate}
\item $R \to S$ soit de type fini,
\item $\dim_{\mathfrak q}(S/R) = n$, et
\item $\text{trdeg}_{\kappa(\mathfrak p)}\kappa(\mathfrak q) = r$.
\end{enumerate}
Alors il existe $f \in R$, $f \not \in \mathfrak p$,
$g \in S$, $g \not\in \mathfrak q$, et un homomorphisme d'anneaux quasi-fini
$$
\varphi : R_f[x_1, \ldots, x_n] \longrightarrow S_g
$$
tel que $\varphi^{-1}(\mathfrak qS_g) =
(\mathfrak p, x_{r + 1}, \ldots, x_n)R_f[x_1, \ldots, x_n]$
\end{lemma}

\begin{proof}
Après remplacement de $S$ par une localisation principale, nous pouvons supposer qu'il
existe un homomorphisme d'anneaux quasi-fini $\varphi : R[t_1, \ldots, t_n] \to S$, voir le
lemme \ref{algebra-lemma-quasi-finite-over-polynomial-algebra}.
Posons $\mathfrak q' = \varphi^{-1}(\mathfrak q)$.
Soit $\overline{\mathfrak q}' \subset \kappa(\mathfrak p)[t_1, \ldots, t_n]$
l'idéal premier correspondant à $\mathfrak q'$. D'après le
lemme \ref{algebra-lemma-refined-Noether-normalization},
il existe un homomorphisme d'anneaux fini
$\kappa(\mathfrak p)[x_1, \ldots, x_n] \to
\kappa(\mathfrak p)[t_1, \ldots, t_n]$
tel que l'image réciproque de $\overline{\mathfrak q}'$ soit
$(x_{r + 1}, \ldots, x_n)$. Soit
$\overline{h}_i \in \kappa(\mathfrak p)[t_1, \ldots, t_n]$
l'image de $x_i$. Il existe un élément
$f \in R$, $f \not \in \mathfrak p$
et des $h_i \in R_f[t_1, \ldots, t_n]$ dont les images sont $\overline{h}_i$
dans $\kappa(\mathfrak p)[t_1, \ldots, t_n]$. Alors l'homomorphisme d'anneaux
$$
R_f[x_1, \ldots, x_n] \longrightarrow R_f[t_1, \ldots, t_n]
$$
devient fini après extension des scalaires à $\kappa(\mathfrak p)$.
En particulier, $R_f[t_1, \ldots, t_n]$ est quasi-fini sur
$R_f[x_1, \ldots, x_n]$ en $\mathfrak q'R_f[t_1, \ldots, t_n]$.
Ainsi, d'après le
lemme \ref{algebra-lemma-quasi-finite-open},
il existe $g \in R_f[t_1, \ldots, t_n]$,
$g \not \in \mathfrak q'R_f[t_1, \ldots, t_n]$
tel que $R_f[x_1, \ldots, x_n] \to R_f[t_1, \ldots, t_n, 1/g]$
soit quasi-fini. On voit donc que le composé
$$
R_f[x_1, \ldots, x_n] \longrightarrow
R_f[t_1, \ldots, t_n, 1/g] \longrightarrow S_{\varphi(g)}
$$
est quasi-fini, d'où le résultat.
\end{proof}

\begin{lemma}
\label{algebra-lemma-dimension-inequality-quasi-finite}
Soit $R \to S$ un homomorphisme d'anneaux de type fini.
Soit $\mathfrak q \subset S$ un idéal premier au-dessus de $\mathfrak p \subset R$.
Si $R \to S$ est quasi-fini en $\mathfrak q$, alors
$\dim(S_{\mathfrak q}) \leq \dim(R_{\mathfrak p})$.
\end{lemma}

\begin{proof}
Si $R_{\mathfrak p}$ est noethérien
(et donc $S_{\mathfrak q}$ noethérien puisqu'il est essentiellement
de type fini sur $R_{\mathfrak p}$),
cela résulte immédiatement du
lemme \ref{algebra-lemma-dimension-base-fibre-total} et des
définitions. Dans le cas général, soit $S'$ la clôture intégrale
de $R_\mathfrak p$ dans $S_\mathfrak p$.
Par le théorème principal de Zariski \ref{algebra-theorem-main-theorem},
on a $S_{\mathfrak q} = S'_{\mathfrak q'}$ pour un certain
$\mathfrak q' \subset S'$ au-dessus de $\mathfrak q$.
D'après le lemme \ref{algebra-lemma-integral-dim-up}, on a
$\dim(S') \leq \dim(R_\mathfrak p)$ et donc, a fortiori,
$\dim(S_\mathfrak q) = \dim(S'_{\mathfrak q'}) \leq \dim(R_\mathfrak p)$.
\end{proof}

\begin{lemma}
\label{algebra-lemma-dimension-quasi-finite-over-polynomial-algebra}
\begin{slogan}
Un revêtement quasi-fini de l'espace affine de dimension n est de dimension au plus n.
\end{slogan}
Soit $k$ un corps. Soit $S$ une $k$-algèbre de type fini.
Supposons qu'il existe un homomorphisme de $k$-algèbres quasi-fini
$k[t_1, \ldots, t_n] \subset S$. Alors $\dim(S) \leq n$.
\end{lemma}

\begin{proof}
D'après le lemme \ref{algebra-lemma-dim-affine-space}, la dimension de
tout anneau local de $k[t_1, \ldots, t_n]$ est au plus $n$.
Le résultat découle donc du
lemme \ref{algebra-lemma-dimension-inequality-quasi-finite}.
\end{proof}

\begin{lemma}
\label{algebra-lemma-dimension-fibres-bounded-open-upstairs}
Soit $R \to S$ un homomorphisme d'anneaux de type fini.
Soit $\mathfrak q \subset S$ un idéal premier.
Supposons que $\dim_{\mathfrak q}(S/R) = n$.
Il existe un voisinage ouvert $V$ de $\mathfrak q$
dans $\Spec(S)$ tel que
$\dim_{\mathfrak q'}(S/R) \leq n$ pour tout $\mathfrak q' \in V$.
\end{lemma}

\begin{proof}
D'après le lemme \ref{algebra-lemma-quasi-finite-over-polynomial-algebra},
on voit que l'on peut supposer $S$ quasi-fini sur
une algèbre de polynômes $R[t_1, \ldots, t_n]$. En considérant
les fibres, on se ramène au
lemme \ref{algebra-lemma-dimension-quasi-finite-over-polynomial-algebra}.
\end{proof}

\noindent
Autrement dit, le lemme affirme que l'ensemble des points où la
dimension de la fibre est $\leq n$ est ouvert dans $\Spec(S)$.
Le lemme suivant affirme que la formation de cet ouvert commute au
changement de base.
Si l'homomorphisme d'anneaux est de présentation finie, cet ensemble est
un ouvert quasi-compact (voir ci-dessous).

\begin{lemma}
\label{algebra-lemma-dimension-fibres-bounded-open-upstairs-base-change}
Soit $R \to S$ un homomorphisme d'anneaux de type fini.
Soit $R \to R'$ un homomorphisme d'anneaux quelconque.
Posons $S' = R' \otimes_R S$ et notons $f : \Spec(S') \to \Spec(S)$
l'application correspondante entre les spectres.
Soit $n \geq 0$.
L'image réciproque
$f^{-1}(\{\mathfrak q \in \Spec(S) \mid
\dim_{\mathfrak q}(S/R) \leq n\})$
est égale à
$\{\mathfrak q' \in \Spec(S') \mid
\dim_{\mathfrak q'}(S'/R') \leq n\}$.
\end{lemma}

\begin{proof}
La condition s'exprime en termes des dimensions
des anneaux des fibres, qui sont de type fini sur un corps.
En la combinant avec le
lemme \ref{algebra-lemma-dimension-at-a-point-preserved-field-extension},
on obtient le lemme.
\end{proof}

\begin{lemma}
\label{algebra-lemma-dimension-fibres-bounded-quasi-compact-open-upstairs}
Soit $R \to S$ un homomorphisme d'anneaux de présentation finie.
Soit $n \geq 0$. L'ensemble
$$
V_n = \{\mathfrak q \in \Spec(S) \mid \dim_{\mathfrak q}(S/R) \leq n\}
$$
est un ouvert quasi-compact de $\Spec(S)$.
\end{lemma}

\begin{proof}
Il est ouvert d'après le lemme \ref{algebra-lemma-dimension-fibres-bounded-open-upstairs}.
Soit $S = R[x_1, \ldots, x_n]/(f_1, \ldots, f_m)$ une présentation de
$S$. Soit $R_0$ la sous-algèbre sur $\mathbf{Z}$ de $R$ engendrée par les
coefficients des polynômes $f_i$.
Soit $S_0 = R_0[x_1, \ldots, x_n]/(f_1, \ldots, f_m)$.
Alors $S = R \otimes_{R_0} S_0$. D'après le
lemme \ref{algebra-lemma-dimension-fibres-bounded-open-upstairs-base-change},
$V_n$ est l'image réciproque d'un ouvert $V_{0, n}$ par l'application continue
quasi-compacte $\Spec(S) \to \Spec(S_0)$. Comme
$S_0$ est noethérien, on voit que $V_{0, n}$ est quasi-compact.
\end{proof}

\begin{lemma}
\label{algebra-lemma-finite-type-domain-over-valuation-ring-dim-fibres}
Soit $R$ un anneau de valuation de corps résiduel $k$ et de corps
des fractions $K$. Soit $S$ un anneau intègre contenant $R$ tel que
$S$ soit de type fini sur $R$. Si $S \otimes_R k$ n'est pas l'anneau
nul, alors
$$
\dim(S \otimes_R k) = \dim(S \otimes_R K)
$$
En fait, $\Spec(S \otimes_R k)$ est équidimensionnel.
\end{lemma}

\begin{proof}
Il suffit de montrer que $\dim_{\mathfrak q}(S/k)$ est égale
à $\dim(S \otimes_R K)$ pour tout idéal premier $\mathfrak q$ de
$S$ contenant $\mathfrak m_RS$. Choisissons un tel idéal premier. D'après le
lemme \ref{algebra-lemma-dimension-fibres-bounded-open-upstairs},
l'inégalité $\dim_{\mathfrak q}(S/k) \geq \dim(S \otimes_R K)$
est vérifiée. Posons $n = \dim_{\mathfrak q}(S/k)$. D'après le
lemme \ref{algebra-lemma-quasi-finite-over-polynomial-algebra},
après remplacement de $S$ par $S_g$ pour un certain $g \in S$, $g \not \in \mathfrak q$,
il existe un homomorphisme d'anneaux quasi-fini
$R[t_1, \ldots, t_n] \to S$. Si $\dim(S \otimes_R K) < n$,
alors $K[t_1, \ldots, t_n] \to S \otimes_R K$ a un noyau non nul.
Écrivons $f = \sum a_I t_1^{i_1}\ldots t_n^{i_n}$. Après avoir divisé
$f$ par un coefficient non nul de $f$ de valuation minimale, nous pouvons
supposer que $f\in R[t_1, \ldots, t_n]$ et qu'un certain $a_I$ a une image non nulle
dans $k$. Ainsi, l'homomorphisme d'anneaux $k[t_1, \ldots, t_n] \to S \otimes_R k$
a un noyau non nul, ce qui entraîne que $\dim(S \otimes_R k) < n$.
Contradiction.
\end{proof}
























\section{Algèbres et modules de présentation finie}
\label{algebra-section-finite-presentation}

\noindent
Dans cette section, nous examinons quelques résultats classiques dont
l'élément essentiel est une hypothèse de type fini ou de
présentation finie.

\begin{lemma}
\label{algebra-lemma-finite-type-descends}
Soit $R \to S$ un homomorphisme d'anneaux.
Soit $R \to R'$ un homomorphisme d'anneaux fidèlement plat.
Posons $S' = R'\otimes_R S$.
Alors $R \to S$ est de type fini si et seulement si $R' \to S'$
est de type fini.
\end{lemma}

\begin{proof}
Il est clair que, si $R \to S$ est de type fini, alors $R' \to S'$
est de type fini. Supposons que $R' \to S'$ soit de type fini.
Supposons que $y_1, \ldots, y_m$ engendrent $S'$ sur $R'$.
Écrivons $y_j = \sum_i a_{ij} \otimes x_{ji}$ pour certains
$a_{ij} \in R'$ et $x_{ji} \in S$. Soit $A \subset S$
la sous-$R$-algèbre engendrée par les $x_{ij}$.
Par platitude, on a $A' := R' \otimes_R A \subset S'$ et,
par construction, $y_j \in A'$. Ainsi, $A' = S'$.
Par fidèle platitude, $A = S$.
\end{proof}

\begin{lemma}
\label{algebra-lemma-finite-presentation-descends}
Soit $R \to S$ un homomorphisme d'anneaux.
Soit $R \to R'$ un homomorphisme d'anneaux fidèlement plat.
Posons $S' = R'\otimes_R S$.
Alors $R \to S$ est de présentation finie si et seulement si $R' \to S'$
est de présentation finie.
\end{lemma}

\begin{proof}
Il est clair que, si $R \to S$ est de présentation finie, alors $R' \to S'$
est de présentation finie. Supposons que $R' \to S'$ soit de présentation finie.
D'après le lemme \ref{algebra-lemma-finite-type-descends}, on voit
que $R \to S$ est de type fini. Écrivons $S = R[x_1, \ldots, x_n]/I$.
Par platitude, $S' = R'[x_1, \ldots, x_n]/R'\otimes I$.
Supposons que $g_1, \ldots, g_m$ engendrent $R'\otimes I$ sur $R'[x_1, \ldots, x_n]$.
Écrivons $g_j = \sum_i a_{ij} \otimes f_{ji}$ pour certains
$a_{ij} \in R'$ et $f_{ji} \in I$. Soit $J \subset I$
l'idéal engendré par les $f_{ij}$.
Par platitude, on a $R' \otimes_R J \subset R'\otimes_R I$, et
ce sont tous deux des idéaux de $R'[x_1, \ldots, x_n]$.
Par construction, $g_j \in R' \otimes_R J$. Ainsi,
$R' \otimes_R J = R'\otimes_R I$.
Par fidèle platitude, $J = I$.
\end{proof}

\begin{lemma}
\label{algebra-lemma-construct-fp-module}
Soit $R$ un anneau.
Soit $I \subset R$ un idéal.
Soit $S \subset R$ une partie multiplicative.
Posons $R' = S^{-1}(R/I) = S^{-1}R/S^{-1}I$.
\begin{enumerate}
\item Pour tout $R'$-module de type fini $M'$, il existe un
$R$-module de type fini $M$ tel que $S^{-1}(M/IM) \cong M'$.
\item Pour tout $R'$-module de présentation finie $M'$, il existe un
$R$-module de présentation finie $M$ tel que $S^{-1}(M/IM) \cong M'$.
\end{enumerate}
\end{lemma}

\begin{proof}
Démonstration de (1). Choisissons une suite exacte courte
$0 \to K' \to (R')^{\oplus n} \to M' \to 0$.
Soit $K \subset R^{\oplus n}$ l'image réciproque de
$K'$ par l'homomorphisme $R^{\oplus n} \to (R')^{\oplus n}$.
Alors $M = R^{\oplus n}/K$ convient.

\medskip\noindent
Démonstration de (2).
Choisissons une présentation $(R')^{\oplus m} \to (R')^{\oplus n} \to M' \to 0$.
Supposons que le premier homomorphisme soit donné par la matrice
$A' = (a'_{ij})$ et que le second soit déterminé par des générateurs
$x'_i \in M'$, $i = 1, \ldots, n$. Comme $R' = S^{-1}(R/I)$, on peut choisir
$s \in S$ et une matrice $A = (a_{ij})$ à coefficients dans $R$
tels que $a'_{ij} = a_{ij} / s \bmod S^{-1}I$. Soit $M$ le
$R$-module de présentation finie ayant pour présentation
$R^{\oplus m} \to R^{\oplus n} \to M \to 0$,
où le premier homomorphisme est donné par la matrice $A$ et le second
est déterminé par des générateurs $x_i \in M$, $i = 1, \ldots, n$.
Alors l'homomorphisme $M \to M'$, $x_i \mapsto x'_i$, induit un isomorphisme
$S^{-1}(M/IM) \cong M'$.
\end{proof}

\begin{lemma}
\label{algebra-lemma-construct-fp-module-from-localization}
Soit $R$ un anneau.
Soit $S \subset R$ une partie multiplicative.
Soit $M$ un $R$-module.
\begin{enumerate}
\item Si $S^{-1}M$ est un $S^{-1}R$-module de type fini, alors il
existe un $R$-module de type fini $M'$ et un homomorphisme $M' \to M$ qui induit un
isomorphisme $S^{-1}M' \to S^{-1}M$.
\item Si $S^{-1}M$ est un $S^{-1}R$-module de présentation finie,
alors il existe un $R$-module $M'$ de présentation finie
et un homomorphisme $M' \to M$ qui induit un isomorphisme
$S^{-1}M' \to S^{-1}M$.
\end{enumerate}
\end{lemma}

\begin{proof}
Démonstration de (1). Soient $x_1, \ldots, x_n \in M$ des éléments qui engendrent
$S^{-1}M$ comme $S^{-1}R$-module. Soit $M'$ le
sous-$R$-module de $M$ engendré par $x_1, \ldots, x_n$.

\medskip\noindent
Démonstration de (2). Soient $x_1, \ldots, x_n \in M$ des éléments qui engendrent
$S^{-1}M$ comme $S^{-1}R$-module. Posons
$K = \Ker(R^{\oplus n} \to M)$, où l'homomorphisme est donné par
la règle $(a_1, \ldots, a_n) \mapsto \sum a_i x_i$. D'après le
lemme \ref{algebra-lemma-extension},
on voit que $S^{-1}K$ est un $S^{-1}R$-module de type fini.
D'après (1), on peut trouver un sous-module de type fini $K' \subset K$
tel que $S^{-1}K' = S^{-1}K$. Prenons
$M' = \Coker(K' \to R^{\oplus n})$.
\end{proof}

\begin{lemma}
\label{algebra-lemma-construct-fp-module-from-stalk}
Soit $R$ un anneau.
Soit $\mathfrak p \subset R$ un idéal premier.
Soit $M$ un $R$-module.
\begin{enumerate}
\item Si $M_{\mathfrak p}$ est un $R_{\mathfrak p}$-module de type fini, alors il
existe un $R$-module de type fini $M'$ et un homomorphisme $M' \to M$ qui induit un
isomorphisme $M'_{\mathfrak p} \to M_{\mathfrak p}$.
\item Si $M_{\mathfrak p}$ est un $R_{\mathfrak p}$-module de présentation finie,
alors il existe un $R$-module $M'$ de présentation finie
et un homomorphisme $M' \to M$ qui induit un isomorphisme
$M'_{\mathfrak p} \to M_{\mathfrak p}$.
\end{enumerate}
\end{lemma}

\begin{proof}
C'est un cas particulier du
lemme \ref{algebra-lemma-construct-fp-module-from-localization}.
\end{proof}

\begin{lemma}
\label{algebra-lemma-local-isomorphism}
Soit $\varphi : R \to S$ un homomorphisme d'anneaux. Soit $\mathfrak q \subset S$
un idéal premier au-dessus de $\mathfrak p \subset R$. Supposons que
\begin{enumerate}
\item $S$ soit de présentation finie sur $R$,
\item $\varphi$ induise un isomorphisme $R_\mathfrak p \cong S_\mathfrak q$.
\end{enumerate}
Alors il existe $f \in R$, $f \not \in \mathfrak p$ et une
$R_f$-algèbre $C$ tels que $S_f \cong R_f \times C$ comme $R_f$-algèbres.
\end{lemma}

\begin{proof}
Écrivons $S = R[x_1, \ldots, x_n]/(g_1, \ldots, g_m)$. Soit $a_i \in R_\mathfrak p$
un élément qui s'envoie sur l'image de $x_i$ dans $S_\mathfrak q$.
Écrivons $a_i = b_i/f$ pour un certain $f \in R$, $f \not \in \mathfrak p$.
Après avoir remplacé $R$ par $R_f$ et $x_i$ par $x_i - a_i$, on peut
supposer que $S = R[x_1, \ldots, x_n]/(g_1, \ldots, g_m)$ et
que $x_i$ s'envoie sur zéro dans $S_\mathfrak q$. Alors, si $c_j$ désigne
le terme constant de $g_j$, on conclut que $c_j$ s'envoie sur zéro
dans $R_\mathfrak p$. Après un nouveau remplacement de $R$, on peut
supposer que les termes constants $c_j$ des $g_j$ sont nuls.
On obtient ainsi un homomorphisme de $R$-algèbres $S \to R$, $x_i \mapsto 0$, dont le
noyau est l'idéal $(x_1, \ldots, x_n)$.

\medskip\noindent
On a les isomorphismes $R_\mathfrak p \to S_\mathfrak q \to R_\mathfrak p$,
et $S \to R$ envoie $x_i$ sur zéro. On doit donc avoir
$S_\mathfrak q = R_\mathfrak p[x_1, \ldots, x_n]/(x_1, \ldots, x_n)$
et, a fortiori, $S_\mathfrak q = S_\mathfrak p/(x_1, \ldots, x_n)S_\mathfrak p$.
Cela signifie que l'idéal de type fini $(x_1, \ldots, x_n)S_\mathfrak p$
est pur dans $S_\mathfrak p$ ; voir la définition \ref{algebra-definition-pure-ideal}.
Ainsi, $(x_1, \ldots, x_n)S_\mathfrak p$ est engendré par un idempotent
$e$ de $S_\mathfrak p$ d'après le lemme \ref{algebra-lemma-finitely-generated-pure-ideal}.
Après avoir remplacé $R \to S$ par $R_f \to S_f$ pour un certain
$f \in R$, $f \not \in \mathfrak p$, on peut trouver un idempotent
$e' \in S$ qui s'envoie sur $e$.
Alors $e'S$ et $(x_1, \ldots, x_n)S$ sont des idéaux de type fini
qui deviennent égaux dans $S_\mathfrak p$. Ainsi, après avoir remplacé $R \to S$ par
$R_f \to S_f$ pour un certain $f \in R$, $f \not \in \mathfrak p$, on peut supposer que
$e'S = (x_1, \ldots, x_n)S$. Poser $C = e'S$ achève la démonstration.
\end{proof}

\begin{lemma}
\label{algebra-lemma-isomorphic-local-rings}
Soit $R$ un anneau.
Soient $S$, $S'$ de présentation finie sur $R$.
Soient $\mathfrak q \subset S$ et $\mathfrak q' \subset S'$
des idéaux premiers. Si $S_{\mathfrak q} \cong S'_{\mathfrak q'}$ comme
$R$-algèbres, alors il existe $g \in S$, $g \not \in \mathfrak q$
et $g' \in S'$, $g' \not \in \mathfrak q'$, tels que
$S_g \cong S'_{g'}$ comme $R$-algèbres.
\end{lemma}

\begin{proof}
Soit $\psi : S_{\mathfrak q} \to S'_{\mathfrak q'}$ l'isomorphisme
de l'hypothèse du lemme.
Écrivons $S = R[x_1, \ldots, x_n]/(f_1, \ldots, f_r)$ et
$S' = R[y_1, \ldots, y_m]/J$.
Pour chaque $i = 1, \ldots, n$, choisissons une fraction
$h_i/g_i$ avec $h_i, g_i \in R[y_1, \ldots, y_m]$
avec $g_i \bmod J$ n'appartenant pas à $\mathfrak q'$ ; cette fraction représente
l'image de $x_i$ par $\psi$. Après avoir remplacé
$S'$ par $S'_{g_1 \ldots g_n}$ et introduit
$R[y_1, \ldots, y_m, y_{m + 1}]$ (en envoyant $y_{m + 1}$ sur $1/(g_1\ldots g_n)$),
on peut supposer que $\psi(x_i)$ est l'image d'un certain
$h_i \in R[y_1, \ldots, y_m]$. Considérons les éléments
$f_j(h_1, \ldots, h_n) \in R[y_1, \ldots, y_m]$.
Puisque $\psi$ annule chaque $f_j$,
il existe $g \in R[y_1, \ldots, y_m]$, $g \bmod J \not \in \mathfrak q'$
tel que $g f_j(h_1, \ldots, h_n) \in J$ pour chaque $j = 1, \ldots, r$.
Après avoir remplacé $S'$ par $S'_g$ et
introduit $R[y_1, \ldots, y_m, y_{m + 1}]$ comme précédemment, on peut supposer que
$f_j(h_1, \ldots, h_n) \in J$. On obtient ainsi un homomorphisme d'anneaux
$S \to S'$, $x_i \mapsto h_i$, qui induit $\psi$ sur les anneaux locaux.
D'après le lemme \ref{algebra-lemma-compose-finite-type},
l'homomorphisme $S \to S'$ est de présentation finie.
D'après le lemme \ref{algebra-lemma-local-isomorphism},
on peut supposer que $S' = S \times C$. Ainsi, en inversant dans $S'$
l'idempotent correspondant au facteur $C$, on obtient le résultat.
\end{proof}

\begin{lemma}
\label{algebra-lemma-finite-type-mod-nilpotent}
Soit $R$ un anneau. Soit $I \subset R$ un idéal nilpotent.
Soit $S$ une $R$-algèbre telle que $R/I \to S/IS$
soit de type fini. Alors $R \to S$ est de type fini.
\end{lemma}

\begin{proof}
Choisissons $s_1, \ldots, s_n \in S$ dont les images dans $S/IS$ engendrent
$S/IS$ comme algèbre sur $R/I$. D'après la partie (11) du lemme \ref{algebra-lemma-NAK},
on voit que l'homomorphisme de $R$-algèbres $R[x_1, \ldots, x_n] \to S$, $x_i \mapsto s_i$,
est surjectif, ce qui donne le résultat.
\end{proof}

\begin{lemma}
\label{algebra-lemma-surjective-mod-locally-nilpotent}
Soit $R$ un anneau. Soit $I \subset R$ un idéal localement nilpotent.
Soit $S \to S'$ un homomorphisme de $R$-algèbres tel que $S \to S'/IS'$ soit surjectif
et que $S'$ soit de type fini sur $R$. Alors $S \to S'$ est surjectif.
\end{lemma}

\begin{proof}
Écrivons $S' = R[x_1, \ldots, x_m]/K$ pour un certain idéal $K$. Par hypothèse, il
existe des $g_j = x_j + \sum \delta_{j, J} x^J \in R[x_1, \ldots, x_n]$ avec
$\delta_{j, J} \in I$ et $g_j \bmod K \in \Im(S \to S')$.
Il suffit donc de montrer que $g_1, \ldots, g_m$ engendrent
$R[x_1, \ldots, x_n]$. Soit $R_0 \subset R$ une
$\mathbf{Z}$-sous-algèbre de type fini de $R$ contenant notamment les $\delta_{j, J}$.
Alors $R_0 \cap I$ est un idéal nilpotent (d'après le
lemme \ref{algebra-lemma-Noetherian-power}). Il s'ensuit que
$R_0[x_1, \ldots, x_n]$ est engendré par $g_1, \ldots, g_m$ (car
$x_j \mapsto g_j$ définit un automorphisme de $R_0[x_1, \ldots, x_m]$ ;
nous omettons les détails). Comme $R$ est la réunion des sous-anneaux $R_0$, le résultat suit.
\end{proof}

\begin{lemma}
\label{algebra-lemma-isomorphism-modulo-ideal}
Soit $R$ un anneau. Soit $I \subset R$ un idéal. Soit $S \to S'$
un homomorphisme de $R$-algèbres. Soit $IS \subset \mathfrak q \subset S$
un idéal
premier. Supposons que
\begin{enumerate}
\item $S \to S'$ soit surjectif,
\item $S_\mathfrak q/IS_\mathfrak q \to S'_\mathfrak q/IS'_\mathfrak q$
soit un isomorphisme,
\item $S$ soit de type fini sur $R$,
\item $S'$ soit de présentation finie sur $R$, et
\item $S'_\mathfrak q$ soit plat sur $R$.
\end{enumerate}
Alors $S_g \to S'_g$ est un isomorphisme pour un certain
$g \in S$, $g \not \in \mathfrak q$.
\end{lemma}

\begin{proof}
Soit $J = \Ker(S \to S')$. D'après le
lemme \ref{algebra-lemma-compose-finite-type},
$J$ est un idéal de type fini. Comme $S'_\mathfrak q$ est plat
sur $R$, on voit que
$J_\mathfrak q/IJ_\mathfrak q \subset S_\mathfrak q/IS_{\mathfrak q}$
(appliquer le lemme \ref{algebra-lemma-flat-tor-zero} à $0 \to J \to S \to S' \to 0$).
D'après l'hypothèse (2), on voit que $J_\mathfrak q/IJ_\mathfrak q$ est nul.
D'après le lemme de Nakayama (lemme \ref{algebra-lemma-NAK}),
il existe $g \in S$, $g \not \in \mathfrak q$, tel
que $J_g = 0$. Ainsi, $S_g \cong S'_g$, comme souhaité.
\end{proof}

\begin{lemma}
\label{algebra-lemma-isomorphism-modulo-locally-nilpotent}
Soit $R$ un anneau. Soit $I \subset R$ un idéal. Soit $S \to S'$
un homomorphisme de $R$-algèbres. Supposons que
\begin{enumerate}
\item $I$ soit localement nilpotent,
\item $S/IS \to S'/IS'$ soit un isomorphisme,
\item $S$ soit de type fini sur $R$,
\item $S'$ soit de présentation finie sur $R$, et
\item $S'$ soit plat sur $R$.
\end{enumerate}
Alors $S \to S'$ est un isomorphisme.
\end{lemma}

\begin{proof}
D'après le lemme \ref{algebra-lemma-surjective-mod-locally-nilpotent}, l'homomorphisme
$S \to S'$ est surjectif. Comme $I$ est localement nilpotent, les
idéaux $IS$ et $IS'$ le sont aussi (lemme \ref{algebra-lemma-locally-nilpotent}). Ainsi,
tout idéal premier $\mathfrak q$ de $S$ contient $IS$ et, trivialement,
$S_\mathfrak q/IS_\mathfrak q \cong S'_\mathfrak q/IS'_\mathfrak q$.
Le lemme \ref{algebra-lemma-isomorphism-modulo-ideal} s'applique donc,
et l'on voit que $S_\mathfrak q \to S'_\mathfrak q$ est un
isomorphisme pour tout idéal premier $\mathfrak q \subset S$.
Il s'ensuit que $S \to S'$ est injectif, par exemple d'après le
lemme \ref{algebra-lemma-characterize-zero-local}.
\end{proof}





\section{Colimites et morphismes de présentation finie}
\label{algebra-section-colimits-flat}

\noindent
Dans cette section, nous démontrons quelques lemmes préliminaires
qui nous aideront par la suite à établir un résultat au moyen d'une
réduction noethérienne absolue.
Dans Catégories, section \ref{categories-section-directed-colimits},
nous étudions les colimites filtrantes en général.
Voici un exemple de cette notion très générale.

\begin{lemma}
\label{algebra-lemma-ring-colimit-fp-category}
Soit $R \to A$ un morphisme d'anneaux. Considérons la catégorie
$\mathcal{I}$ de tous les diagrammes de morphismes de $R$-algèbres
$A' \to A$, où $A'$ est de présentation finie sur $R$. Alors
$\mathcal{I}$ est filtrante et la colimite des $A'$ suivant
$\mathcal{I}$ est isomorphe à $A$.
\end{lemma}
\begin{proof}
La catégorie\footnote{Afin d'éviter les difficultés ensemblistes, nous ne
considérons que les $A' \to A$ tels que $A'$ soit un quotient de
$R[x_1, x_2, x_3, \ldots]$.}
$\mathcal{I}$ est non vide, puisque $R \to R$ en est un objet.
Considérons deux objets $A' \to A$ et $A'' \to A$ de $\mathcal{I}$.
Alors $A' \otimes_R A'' \to A$ appartient à
$\mathcal{I}$ (utiliser les Lemmes \ref{algebra-lemma-compose-finite-type} et
\ref{algebra-lemma-base-change-finiteness}). Les morphismes d'anneaux
$A' \to A' \otimes_R A''$ et $A'' \to A' \otimes_R A''$
définissent des flèches de $\mathcal{I}$, ce qui démontre la deuxième
propriété définissant une catégorie filtrante ; voir
Catégories, Définition \ref{categories-definition-directed}.
Enfin, supposons donnés deux morphismes $\sigma, \tau : A' \to A''$
dans $\mathcal{I}$. Si $x_1, \ldots, x_r \in A'$ engendrent
$A'$ comme $R$-algèbre, nous pouvons considérer
$A''' = A''/(\sigma(x_i) - \tau(x_i))$.
C'est une $R$-algèbre de présentation finie, et le morphisme de
$R$-algèbres donné $A'' \to A$ se factorise par la surjection
$\nu : A'' \to A'''$. Ainsi, $\nu$ est un morphisme de $\mathcal{I}$
qui égalise $\sigma$ et $\tau$, comme souhaité.

\medskip\noindent
Le fait que notre catégorie d'indices soit cofiltrante signifie que nous
pouvons calculer la valeur de $B = \colim_{A' \to A} A'$ dans la catégorie
des ensembles (certains détails sont omis ; comparer avec la discussion de
Catégories, section \ref{categories-section-directed-colimits}).
Pour voir que $B \to A$ est surjectif, pour tout $a \in A$ nous pouvons
utiliser $R[x] \to A$, $x \mapsto a$, et constater que $a$ appartient à
l'image de $B \to A$. Réciproquement, si $b \in B$ s'envoie sur zéro dans
$A$, nous pouvons trouver $A' \to A$ dans $\mathcal{I}$ et $a' \in A'$
qui s'envoie sur $b$. Alors $A'/(a') \to A$ appartient lui aussi à
$\mathcal{I}$, et le morphisme $A' \to B$ se factorise sous la forme
$A' \to A'/(a') \to B$, ce qui montre que $b = 0$, comme souhaité.
\end{proof}
\noindent
Il est souvent plus commode de considérer les colimites indexées par des
ensembles préordonnés. Soit $(\Lambda, \geq)$ un ensemble préordonné.
Un système d'anneaux sur $\Lambda$ est constitué d'un anneau
$R_\lambda$ pour chaque $\lambda \in \Lambda$ et d'un morphisme
$R_\lambda \to R_\mu$ dès que $\lambda \leq \mu$.
Ces morphismes doivent satisfaire la règle selon laquelle
$R_\lambda \to R_\mu \to R_\nu$ est égal au morphisme
$R_\lambda \to R_\nu$ pour tous $\lambda \leq \mu \leq \nu$.
Voir Catégories, section \ref{categories-section-posets-limits}.
Nous supposerons souvent que $(I, \leq)$ est {\it filtrant},
ce qui signifie que $\Lambda$ est non vide et que, étant donnés
$\lambda, \mu \in \Lambda$, il existe $\nu \in \Lambda$ tel que
$\lambda \leq \nu$ et $\mu \leq \nu$.
Rappelons que, dans ce cas, la colimite $\colim_\lambda R_\lambda$
est parfois appelée « limite directe »
(mais nous n'emploierons pas cette terminologie).

\medskip\noindent
Notons que Catégories, Lemme \ref{categories-lemma-directed-category-system},
nous dit que les colimites indexées par des catégories filtrantes reviennent
au même que les colimites indexées par des ensembles filtrants.

\begin{lemma}
\label{algebra-lemma-ring-colimit-fp}
Soit $R \to A$ un morphisme d'anneaux. Il existe un système filtrant
$A_\lambda$ de $R$-algèbres de présentation finie tel que
$A = \colim_\lambda A_\lambda$. Si $A$ est de type fini sur $R$, on peut
faire en sorte que tous les morphismes de transition du système des
$A_\lambda$ soient surjectifs.
\end{lemma}

\begin{proof}
Une première démonstration consiste à déduire l'énoncé du
Lemme \ref{algebra-lemma-ring-colimit-fp-category} et de
Catégories, Lemme \ref{categories-lemma-directed-category-system}.

\medskip\noindent
Seconde démonstration.
Comparer avec la démonstration du Lemme \ref{algebra-lemma-module-colimit-fp}.
Considérons une partie finie quelconque $S \subset A$ et une famille finie
$E$ de relations polynomiales entre les éléments de $S$.
Ainsi, à chaque $s \in S$ correspond $x_s \in A$, et chaque $e \in E$
consiste en un polynôme $f_e \in R[X_s; s\in S]$ tel que
$f_e(x_s) = 0$. Posons
$A_{S, E} = R[X_s; s\in S]/(f_e; e\in E)$ ; c'est une $R$-algèbre de
présentation finie. Il existe des morphismes canoniques
$A_{S, E} \to A$.
Si $S \subset S'$ et si les éléments de $E$ correspondent, par le morphisme
$R[X_s; s \in S] \to R[X_s; s\in S']$, à une partie de $E'$, il existe
alors un morphisme évident $A_{S, E} \to A_{S', E'}$ qui commute aux
morphismes vers $A$. Ainsi, en prenant pour $\Lambda$ l'ensemble des couples
$(S, E)$ ordonné par inclusion comme ci-dessus, on obtient un ensemble
partiellement ordonné filtrant. Il est clair que la colimite de ce système
filtrant est $A$.

\medskip\noindent
Pour la dernière assertion, supposons $A = R[x_1, \ldots, x_n]/I$.
Considérons alors la partie $\Lambda' \subset \Lambda$ formée des systèmes
$(S, E)$ ci-dessus pour lesquels $S = \{x_1, \ldots, x_n\}$.
On voit aisément que l'on a encore
$A = \colim_{\lambda' \in \Lambda'} A_{\lambda'}$.
De plus, les morphismes de transition sont manifestement surjectifs.
\end{proof}
\noindent
Il se trouve que l'on peut caractériser comme suit les morphismes d'anneaux
de présentation finie. Cela signifie en un certain sens que les algèbres
de présentation finie sont les objets « compacts » de la catégorie des
$R$-algèbres.

\begin{lemma}
\label{algebra-lemma-characterize-finite-presentation}
Soit $\varphi : R \to S$ un morphisme d'anneaux. Les assertions suivantes
sont équivalentes :
\begin{enumerate}
\item $\varphi$ est de présentation finie,
\item pour tout système filtrant $A_\lambda$ de $R$-algèbres, l'application
$$
\colim_\lambda \Hom_R(S, A_\lambda) \longrightarrow
\Hom_R(S, \colim_\lambda A_\lambda)
$$
est bijective, et
\item pour tout système filtrant $A_\lambda$ de $R$-algèbres, l'application
$$
\colim_\lambda \Hom_R(S, A_\lambda) \longrightarrow
\Hom_R(S, \colim_\lambda A_\lambda)
$$
est surjective.
\end{enumerate}
\end{lemma}
\begin{proof}
Supposons (1) et écrivons
$S = R[x_1, \ldots, x_n] / (f_1, \ldots, f_m)$.
Posons $A = \colim A_\lambda$. Remarquons qu'un homomorphisme de
$R$-algèbres $S \to A$ ou $S \to A_\lambda$ est déterminé par les images
de $x_1, \ldots, x_n$. Il est donc clair que
$\colim_\lambda \Hom_R(S, A_\lambda) \to \Hom_R(S, A)$
est injectif. Pour voir qu'il est surjectif, soit $\chi : S \to A$
un homomorphisme de $R$-algèbres. Chaque $x_i$ s'envoie alors sur un
élément de l'image d'un certain $A_{\lambda_i}$.
Nous pouvons choisir $\mu \geq \lambda_i$, $i = 1, \ldots, n$, et supposer
que $\chi(x_i)$ est l'image de $y_i \in A_\mu$ pour
$i = 1, \ldots, n$. Considérons
$z_j = f_j(y_1, \ldots, y_n) \in A_\mu$.
Comme $\chi$ est un homomorphisme, l'image de $z_j$ dans
$A = \colim_\lambda A_\lambda$ est nulle. Il existe donc
$\mu_j \geq \mu$ tel que $z_j$ s'envoie sur zéro dans $A_{\mu_j}$.
Choisissons $\nu \geq \mu_j$, $j = 1, \ldots, m$. Les images de
$z_1, \ldots, z_m$ sont alors nulles dans $A_\nu$. Cela signifie
exactement que les $y_i$ s'envoient sur des éléments
$y'_i \in A_\nu$ qui satisfont les relations
$f_j(y'_1, \ldots, y'_n) = 0$. Nous obtenons donc un morphisme d'anneaux
$S \to A_\nu$. Cela montre que (1) implique (2).

\medskip\noindent
Il est clair que (2) implique (3). Supposons (3).
D'après le Lemme \ref{algebra-lemma-ring-colimit-fp}, nous pouvons écrire
$S = \colim_\lambda S_\lambda$, où $S_\lambda$ est de présentation finie
sur $R$. Le morphisme identité se factorise alors sous la forme
$$
S \to S_\lambda \to S
$$
pour un certain $\lambda$. Il s'ensuit que $S$ est de présentation finie
sur $S_\lambda$ d'après le Lemme \ref{algebra-lemma-compose-finite-type}, partie (4),
appliqué à $S \to S_\lambda \to S$. En appliquant la partie (2) du même
lemme à $R \to S_\lambda \to S$, nous concluons que $S$ est de
présentation finie sur $R$.
\end{proof}
\noindent
À l'aide des résultats élémentaires qui précèdent, nous pouvons donner le
critère suivant pour qu'une algèbre $A$ soit une colimite filtrante
d'algèbres d'un type donné.

\begin{lemma}
\label{algebra-lemma-when-colimit}
Soit $R \to \Lambda$ un morphisme d'anneaux. Soit $\mathcal{E}$ un ensemble
de $R$-algèbres tel que chaque $A \in \mathcal{E}$ soit de présentation
finie sur $R$. Les deux assertions suivantes sont équivalentes :
\begin{enumerate}
\item $\Lambda$ est une colimite filtrante d'éléments de $\mathcal{E}$, et
\item pour tout morphisme de $R$-algèbres $A \to \Lambda$, avec $A$ de
présentation finie sur $R$, on peut trouver une factorisation
$A \to B \to \Lambda$ avec $B \in \mathcal{E}$.
\end{enumerate}
\end{lemma}

\begin{proof}
Supposons que $\mathcal{I} \to \mathcal{E}$, $i \mapsto A_i$, soit un
diagramme filtrant tel que $\Lambda = \colim_i A_i$.
Soit $A \to \Lambda$ un morphisme de $R$-algèbres, avec $A$ de
présentation finie sur $R$. L'application du
Lemme \ref{algebra-lemma-characterize-finite-presentation} fournit une
factorisation $A \to A_i \to \Lambda$. Ainsi, (1) implique (2).

\medskip\noindent
Considérons la catégorie $\mathcal{I}$ du
Lemme \ref{algebra-lemma-ring-colimit-fp-category}.
D'après Catégories, Lemme \ref{categories-lemma-cofinal-in-filtered},
la sous-catégorie pleine $\mathcal{J}$ formée des $A \to \Lambda$ tels que
$A \in \mathcal{E}$ est cofinale dans $\mathcal{I}$ et est une catégorie
filtrante. Alors $\Lambda$ est aussi la colimite suivant $\mathcal{J}$,
d'après Catégories, Lemme \ref{categories-lemma-cofinal}.
\end{proof}
\noindent
Mais il y a davantage. Plus précisément, étant donné
$R = \colim_\lambda R_\lambda$, on voit que la catégorie des $R$-modules
de présentation finie est équivalente à la limite des catégories des
$R_\lambda$-modules de présentation finie. Il en va de même des catégories
des $R$-algèbres de présentation finie.

\begin{lemma}
\label{algebra-lemma-module-map-property-in-colimit}
Soient $A$ un anneau et $M, N$ des $A$-modules.
Supposons que $R = \colim_{i \in I} R_i$ soit une colimite filtrante de
$A$-algèbres.
\begin{enumerate}
\item Si $M$ est un $A$-module de type fini et si $u, u' : M \to N$ sont des
morphismes de $A$-modules tels que
$u \otimes 1 = u' \otimes 1 : M \otimes_A R \to N \otimes_A R$,
alors, pour un certain $i$, on a
$u \otimes 1 = u' \otimes 1 : M \otimes_A R_i \to N \otimes_A R_i$.
\item Si $N$ est un $A$-module de type fini et si $u : M \to N$ est un morphisme
de $A$-modules tel que
$u \otimes 1 : M \otimes_A R \to N \otimes_A R$ soit surjectif, alors,
pour un certain $i$, le morphisme
$u \otimes 1 : M \otimes_A R_i \to N \otimes_A R_i$ est surjectif.
\item Si $N$ est un $A$-module de présentation finie et si
$v : N \otimes_A R \to M \otimes_A R$ est un morphisme de $R$-modules,
alors il existe $i$ et un morphisme de $R_i$-modules
$v_i : N \otimes_A R_i \to M \otimes_A R_i$ tels que
$v = v_i \otimes 1$.
\item Si $M$ est un $A$-module de type fini, si $N$ est un $A$-module de
présentation finie et si $u : M \to N$ est un morphisme de $A$-modules
tel que $u \otimes 1 : M \otimes_A R \to N \otimes_A R$ soit un
isomorphisme, alors, pour un certain $i$, le morphisme
$u \otimes 1 : M \otimes_A R_i \to N \otimes_A R_i$ est un isomorphisme.
\end{enumerate}
\end{lemma}

\begin{proof}
Pour démontrer (1), supposons que $u$ soit comme dans (1) et soient
$x_1, \ldots, x_m \in M$ des générateurs. Comme
$N \otimes_A R = \colim_i N \otimes_A R_i$, nous pouvons choisir
$i \in I$ tel que $u(x_j) \otimes 1 = u'(x_j) \otimes 1$ dans
$M \otimes_A R_i$, $j = 1, \ldots, m$. Pour un tel $i$, on a
$u \otimes 1 = u' \otimes 1 : M \otimes_A R_i \to N \otimes_A R_i$.

\medskip\noindent
Pour démontrer (2), supposons $u \otimes 1$ surjectif et soient
$y_1, \ldots, y_m \in N$ des générateurs. Comme
$N \otimes_A R = \colim_i N \otimes_A R_i$, nous pouvons choisir
$i \in I$ et $z_j \in M \otimes_A R_i$, $j = 1, \ldots, m$, dont les
images dans $N \otimes_A R$ sont égales à $y_j \otimes 1$.
Pour un tel $i$, le morphisme
$u \otimes 1 : M \otimes_A R_i \to N \otimes_A R_i$ est surjectif.

\medskip\noindent
Pour démontrer (3), soient $y_1, \ldots, y_m \in N$ des générateurs.
Soit $K = \Ker(A^{\oplus m} \to N)$, où le morphisme est donné par la règle
$(a_1, \ldots, a_m) \mapsto \sum a_j x_j$. Soient
$k_1, \ldots, k_t$ des générateurs de $K$. Écrivons
$k_s = (k_{s1}, \ldots, k_{sm})$.
Comme $M \otimes_A R = \colim_i M \otimes_A R_i$, nous pouvons choisir
$i \in I$ et $z_j \in M \otimes_A R_i$, $j = 1, \ldots, m$, dont les
images dans $M \otimes_A R$ sont égales à $v(y_j \otimes 1)$.
Nous voulons utiliser les $z_j$ pour définir le morphisme
$v_i : N \otimes_A R_i \to M \otimes_A R_i$.
Puisque
$K \otimes_A R_i \to R_i^{\oplus m} \to N \otimes_A R_i \to 0$
est une présentation, il suffit de vérifier que
$\xi_s = \sum_j k_{sj}z_j$ est nul dans $M \otimes_A R_i$ pour chaque
$s = 1, \ldots, t$. Ce n'est pas nécessairement le cas, mais, puisque
l'image de $\xi_s$ dans $M \otimes_A R$ est nulle, cela le devient après
avoir augmenté quelque peu $i$.

\medskip\noindent
Pour démontrer (4), supposons que $u \otimes 1$ soit un isomorphisme, que
$M$ soit de type fini et que $N$ soit de présentation finie.
Soit $v : N \otimes_A R \to M \otimes_A R$ un inverse de $u \otimes 1$.
La partie (3) fournit un morphisme
$v_i : N \otimes_A R_i \to M \otimes_A R_i$ pour un certain $i$.
La partie (1) montre qu'après avoir augmenté $i$, on a
$v_i \circ (u \otimes 1) = \text{id}_{M \otimes_R R_i}$ et
$(u \otimes 1) \circ v_i = \text{id}_{N \otimes_R R_i}$.
\end{proof}
\begin{lemma}
\label{algebra-lemma-colimit-category-fp-modules}
Supposons que $R = \colim_{\lambda \in \Lambda} R_\lambda$ soit une
colimite filtrante d'anneaux. Alors la catégorie des $R$-modules de
présentation finie est la colimite des catégories des $R_\lambda$-modules
de présentation finie. Plus précisément :
\begin{enumerate}
\item Étant donné un $R$-module $M$ de présentation finie, il existe
$\lambda \in \Lambda$ et un $R_\lambda$-module $M_\lambda$ de présentation
finie tels que $M \cong M_\lambda \otimes_{R_\lambda} R$.
\item Étant donnés $\lambda \in \Lambda$, des $R_\lambda$-modules de
présentation finie $M_\lambda, N_\lambda$ et un morphisme de $R$-modules
$\varphi : M_\lambda \otimes_{R_\lambda} R \to
N_\lambda \otimes_{R_\lambda} R$, il existe $\mu \geq \lambda$ et un
morphisme de $R_\mu$-modules
$\varphi_\mu : M_\lambda \otimes_{R_\lambda} R_\mu \to
N_\lambda \otimes_{R_\lambda} R_\mu$
tels que $\varphi = \varphi_\mu \otimes 1_R$.
\item Étant donnés $\lambda \in \Lambda$, des $R_\lambda$-modules de
présentation finie $M_\lambda, N_\lambda$ et des morphismes de
$R_\lambda$-modules $\varphi, \psi : M_\lambda \to N_\lambda$ tels que
$\varphi \otimes 1_R = \psi \otimes 1_R$, on a
$\varphi \otimes 1_{R_\mu} = \psi \otimes 1_{R_\mu}$ pour un certain
$\mu \geq \lambda$.
\end{enumerate}
\end{lemma}

\begin{proof}
Pour démontrer (1), choisissons une présentation
$R^{\oplus m} \to R^{\oplus n} \to M \to 0$.
Supposons que le premier morphisme soit donné par la matrice $A = (a_{ij})$.
Nous pouvons choisir $\lambda \in \Lambda$ et une matrice
$A_\lambda = (a_{\lambda, ij})$ à coefficients dans $R_\lambda$ qui
s'envoie sur $A$ dans $R$.
Il suffit alors de prendre pour $M_\lambda$ le $R_\lambda$-module ayant pour
présentation
$R_\lambda^{\oplus m} \to R_\lambda^{\oplus n} \to M_\lambda \to 0$,
où le premier morphisme est donné par $A_\lambda$.

\medskip\noindent
Les parties (2) et (3) résultent du
Lemme \ref{algebra-lemma-module-map-property-in-colimit}.
\end{proof}
\begin{lemma}
\label{algebra-lemma-algebra-map-property-in-colimit}
Soient $A$ un anneau et $B, C$ des $A$-algèbres.
Supposons que $R = \colim_{i \in I} R_i$ soit une colimite filtrante de
$A$-algèbres.
\begin{enumerate}
\item Si $B$ est une $A$-algèbre de type fini et si
$u, u' : B \to C$ sont des morphismes de $A$-algèbres tels que
$u \otimes 1 = u' \otimes 1 : B \otimes_A R \to C \otimes_A R$,
alors, pour un certain $i$, on a
$u \otimes 1 = u' \otimes 1 : B \otimes_A R_i \to C \otimes_A R_i$.
\item Si $C$ est une $A$-algèbre de type fini et si $u : B \to C$ est un
morphisme de $A$-algèbres tel que
$u \otimes 1 : B \otimes_A R \to C \otimes_A R$ soit surjectif, alors,
pour un certain $i$, le morphisme
$u \otimes 1 : B \otimes_A R_i \to C \otimes_A R_i$ est surjectif.
\item Si $C$ est de présentation finie sur $A$ et si
$v : C \otimes_A R \to B \otimes_A R$ est un morphisme de $R$-algèbres,
alors il existe $i$ et un morphisme de $R_i$-algèbres
$v_i : C \otimes_A R_i \to B \otimes_A R_i$ tels que
$v = v_i \otimes 1$.
\item Si $B$ est une $A$-algèbre de type fini, si $C$ est une $A$-algèbre
de présentation finie et si
$u \otimes 1 : B \otimes_A R \to C \otimes_A R$ est un isomorphisme,
alors, pour un certain $i$, le morphisme
$u \otimes 1 : B \otimes_A R_i \to C \otimes_A R_i$ est un isomorphisme.
\end{enumerate}
\end{lemma}
\begin{proof}
Pour démontrer (1), supposons que $u$ soit comme dans (1) et soient
$x_1, \ldots, x_m \in B$ des générateurs. Comme
$B \otimes_A R = \colim_i B \otimes_A R_i$, nous pouvons choisir
$i \in I$ tel que $u(x_j) \otimes 1 = u'(x_j) \otimes 1$ dans
$B \otimes_A R_i$, $j = 1, \ldots, m$. Pour un tel $i$, on a
$u \otimes 1 = u' \otimes 1 : B \otimes_A R_i \to C \otimes_A R_i$.

\medskip\noindent
Pour démontrer (2), supposons $u \otimes 1$ surjectif et soient
$y_1, \ldots, y_m \in C$ des générateurs. Comme
$B \otimes_A R = \colim_i B \otimes_A R_i$, nous pouvons choisir
$i \in I$ et $z_j \in B \otimes_A R_i$, $j = 1, \ldots, m$, dont les
images dans $C \otimes_A R$ sont égales à $y_j \otimes 1$.
Pour un tel $i$, le morphisme
$u \otimes 1 : B \otimes_A R_i \to C \otimes_A R_i$ est surjectif.

\medskip\noindent
Pour démontrer (3), soient $c_1, \ldots, c_m \in C$ des générateurs.
Soit $K = \Ker(A[x_1, \ldots, x_m] \to N)$, où le morphisme est donné par
la règle $x_j \mapsto \sum c_j$. Soient $f_1, \ldots, f_t$ des
générateurs de l'idéal $K$ de $A[x_1, \ldots, x_m]$.
Nous considérons $f_j = f_j(x_1, \ldots, x_m)$ comme un polynôme.
Comme $B \otimes_A R = \colim_i B \otimes_A R_i$, nous pouvons choisir
$i \in I$ et $z_j \in B \otimes_A R_i$, $j = 1, \ldots, m$, dont les
images dans $B \otimes_A R$ sont égales à $v(c_j \otimes 1)$.
Nous voulons utiliser les $z_j$ pour définir un morphisme
$v_i : C \otimes_A R_i \to B \otimes_A R_i$.
Puisque
$K \otimes_A R_i \to R_i[x_1, \ldots, x_m] \to
C \otimes_A R_i \to 0$
est une présentation, il suffit de vérifier que
$\xi_s = f_j(z_1, \ldots, z_m)$ est nul dans $B \otimes_A R_i$ pour chaque
$s = 1, \ldots, t$. Ce n'est pas nécessairement le cas, mais, puisque
l'image de $\xi_s$ dans $B \otimes_A R$ est nulle, cela le devient après
avoir augmenté quelque peu $i$.

\medskip\noindent
Pour démontrer (4), supposons que $u \otimes 1$ soit un isomorphisme, que
$B$ soit une $A$-algèbre de type fini et que $C$ soit une $A$-algèbre de
présentation finie. Soit
$v : B \otimes_A R \to C \otimes_A R$ un inverse de $u \otimes 1$.
Soit $v_i : C \otimes_A R_i \to B \otimes_A R_i$ comme dans la partie (3).
La partie (1) montre qu'après avoir augmenté $i$, on a
$v_i \circ (u \otimes 1) = \text{id}_{B \otimes_R R_i}$ et
$(u \otimes 1) \circ v_i = \text{id}_{C \otimes_R R_i}$.
\end{proof}
\begin{lemma}
\label{algebra-lemma-colimit-category-fp-algebras}
Supposons que $R = \colim_{\lambda \in \Lambda} R_\lambda$ soit une
colimite filtrante d'anneaux. Alors la catégorie des $R$-algèbres de
présentation finie est la colimite des catégories des $R_\lambda$-algèbres
de présentation finie. Plus précisément :
\begin{enumerate}
\item Étant donnée une $R$-algèbre $A$ de présentation finie, il existe
$\lambda \in \Lambda$ et une $R_\lambda$-algèbre $A_\lambda$ de
présentation finie telles que
$A \cong A_\lambda \otimes_{R_\lambda} R$.
\item Étant donnés $\lambda \in \Lambda$, des $R_\lambda$-algèbres de
présentation finie $A_\lambda, B_\lambda$ et un morphisme de $R$-algèbres
$\varphi : A_\lambda \otimes_{R_\lambda} R \to
B_\lambda \otimes_{R_\lambda} R$, il existe $\mu \geq \lambda$ et un
morphisme de $R_\mu$-algèbres
$\varphi_\mu : A_\lambda \otimes_{R_\lambda} R_\mu \to
B_\lambda \otimes_{R_\lambda} R_\mu$
tels que $\varphi = \varphi_\mu \otimes 1_R$.
\item Étant donnés $\lambda \in \Lambda$, des $R_\lambda$-algèbres de
présentation finie $A_\lambda, B_\lambda$ et des morphismes de
$R_\lambda$-algèbres
$\varphi_\lambda, \psi_\lambda : A_\lambda \to B_\lambda$ tels que
$\varphi \otimes 1_R = \psi \otimes 1_R$, on a
$\varphi \otimes 1_{R_\mu} = \psi \otimes 1_{R_\mu}$ pour un certain
$\mu \geq \lambda$.
\end{enumerate}
\end{lemma}

\begin{proof}
Pour démontrer (1), choisissons une présentation
$A = R[x_1, \ldots, x_n]/(f_1, \ldots, f_m)$.
Nous pouvons choisir $\lambda \in \Lambda$ et des éléments
$f_{\lambda, j} \in R_\lambda[x_1, \ldots, x_n]$ qui s'envoient sur
$f_j \in R[x_1, \ldots, x_n]$.
Il suffit alors de poser
$A_\lambda =
R_\lambda[x_1, \ldots, x_n]/(f_{\lambda, 1}, \ldots, f_{\lambda, m})$.

\medskip\noindent
Les parties (2) et (3) résultent du
Lemme \ref{algebra-lemma-algebra-map-property-in-colimit}.
\end{proof}
\begin{lemma}
\label{algebra-lemma-limit-no-condition-local}
Supposons que $R \to S$ soit un homomorphisme local d'anneaux locaux.
Il existe un ensemble filtrant $(\Lambda, \leq)$ et un système
d'homomorphismes locaux $R_\lambda \to S_\lambda$ d'anneaux locaux tels que :
\begin{enumerate}
\item La colimite du système $R_\lambda \to S_\lambda$ est égale à
$R \to S$.
\item Chaque $R_\lambda$ est essentiellement de type fini sur
$\mathbf{Z}$.
\item Chaque $S_\lambda$ est essentiellement de type fini sur
$R_\lambda$.
\end{enumerate}
\end{lemma}

\begin{proof}
Notons $\varphi : R  \to S$ le morphisme d'anneaux.
Soient $\mathfrak m \subset R$ l'idéal maximal de $R$ et
$\mathfrak n \subset S$ l'idéal maximal de $S$. Posons
$$
\Lambda = \{
(A, B)
\mid
A \subset R, B \subset S, \# A < \infty, \# B < \infty, \varphi(A) \subset B
\}.
$$
Nous prenons pour ordre partiel la relation d'inclusion. Pour chaque
$\lambda = (A, B) \in \Lambda$, soit $R'_\lambda$ la sous-$\mathbf{Z}$-algèbre
engendrée par les $a \in A$, et soit $S'_\lambda$ la sous-$\mathbf{Z}$-algèbre
engendrée par les $b$, $b \in B$.
Soit $R_\lambda$ la localisation de $R'_\lambda$ en l'idéal premier
$R'_\lambda \cap \mathfrak m$, et soit $S_\lambda$ la localisation de
$S'_\lambda$ en l'idéal premier $S'_\lambda \cap \mathfrak n$.
Voici le diagramme :
$$
\xymatrix{
B \ar[r] &
S'_\lambda \ar[r] &
S_\lambda \ar[r] &
S \\
A \ar[r] \ar[u] &
R'_\lambda \ar[r] \ar[u] &
R_\lambda \ar[r] \ar[u] &
R \ar[u]
}.
$$
Les morphismes de transition sont clairs. Nous laissons au lecteur la
démonstration des autres assertions.
\end{proof}
\begin{lemma}
\label{algebra-lemma-limit-essentially-finite-type}
Supposons que $R \to S$ soit un homomorphisme local d'anneaux locaux.
Supposons $S$ essentiellement de type fini sur $R$.
Il existe alors un ensemble filtrant $(\Lambda, \leq)$ et un système
d'homomorphismes locaux $R_\lambda \to S_\lambda$ d'anneaux locaux tels que :
\begin{enumerate}
\item La colimite du système $R_\lambda \to S_\lambda$ est égale à
$R \to S$.
\item Chaque $R_\lambda$ est essentiellement de type fini sur
$\mathbf{Z}$.
\item Chaque $S_\lambda$ est essentiellement de type fini sur
$R_\lambda$.
\item Pour chaque $\lambda \leq \mu$, le morphisme
$S_\lambda \otimes_{R_\lambda} R_\mu \to S_\mu$
présente $S_\mu$ comme la localisation d'un quotient de
$S_\lambda \otimes_{R_\lambda} R_\mu$.
\end{enumerate}
\end{lemma}
\begin{proof}
Notons $\varphi : R  \to S$ le morphisme d'anneaux.
Soient $\mathfrak m \subset R$ l'idéal maximal de $R$ et
$\mathfrak n \subset S$ l'idéal maximal de $S$.
Soient $x_1, \ldots, x_n \in S$ des éléments tels que $S$ soit une
localisation de la sous-$R$-algèbre de $S$ engendrée par
$x_1, \ldots, x_n$. Autrement dit, $S$ est un quotient d'une localisation
de l'anneau de polynômes $R[x_1, \ldots, x_n]$.

\medskip\noindent
Soit $\Lambda = \{ A \subset R \mid \# A < \infty\}$ l'ensemble des
parties finies de $R$. Nous prenons pour ordre partiel la relation
d'inclusion. Pour chaque $\lambda = A \in \Lambda$, soit $R'_\lambda$
la sous-$\mathbf{Z}$-algèbre engendrée par les $a \in A$, et soit
$S'_\lambda$ la sous-$\mathbf{Z}$-algèbre engendrée par les $\varphi(a)$,
$a \in A$, et par les éléments $x_1, \ldots, x_n$.
Soit $R_\lambda$ la localisation de $R'_\lambda$ en l'idéal premier
$R'_\lambda \cap \mathfrak m$, et soit $S_\lambda$ la localisation de
$S'_\lambda$ en l'idéal premier $S'_\lambda \cap \mathfrak n$.
Voici le diagramme :
$$
\xymatrix{
\varphi(A) \amalg \{x_i\} \ar[r] &
S'_\lambda \ar[r] &
S_\lambda \ar[r] &
S \\
A \ar[r] \ar[u] &
R'_\lambda \ar[r] \ar[u] &
R_\lambda \ar[r] \ar[u] &
R \ar[u]
}
$$
Il est clair que si $A \subset B$ correspond à $\lambda \leq \mu$ dans
$\Lambda$, il existe des morphismes canoniques
$R_\lambda \to R_\mu$ et $S_\lambda \to S_\mu$, et l'on obtient ainsi
un système sur l'ensemble filtrant $\Lambda$.

\medskip\noindent
L'assertion $R = \colim R_\lambda$ est claire : tous les morphismes
$R_\lambda \to R$ sont injectifs et tout élément de $R$ finit par
appartenir à l'image. Le même argument s'applique à
$S = \colim S_\lambda$. Les assertions (2) et (3) sont vraies par
construction. La dernière assertion résulte du fait que les morphismes
$S'_\lambda \otimes_{R'_\lambda} R'_\mu \to S'_\mu$
sont manifestement surjectifs.
\end{proof}
\begin{lemma}
\label{algebra-lemma-limit-essentially-finite-presentation}
Supposons que $R \to S$ soit un homomorphisme local d'anneaux locaux.
Supposons $S$ essentiellement de présentation finie sur $R$.
Il existe alors un ensemble filtrant $(\Lambda, \leq)$ et un système
d'homomorphismes locaux $R_\lambda \to S_\lambda$ d'anneaux locaux tels que :
\begin{enumerate}
\item La colimite du système $R_\lambda \to S_\lambda$ est égale à
$R \to S$.
\item Chaque $R_\lambda$ est essentiellement de type fini sur
$\mathbf{Z}$.
\item Chaque $S_\lambda$ est essentiellement de type fini sur
$R_\lambda$.
\item Pour chaque $\lambda \leq \mu$, le morphisme
$S_\lambda \otimes_{R_\lambda} R_\mu \to S_\mu$
présente $S_\mu$ comme la localisation de
$S_\lambda \otimes_{R_\lambda} R_\mu$ en un idéal premier.
\end{enumerate}
\end{lemma}

\begin{proof}
Par hypothèse, nous pouvons choisir un isomorphisme
$\Phi : (R[x_1, \ldots, x_n]/I)_{\mathfrak q} \to S$,
où $I \subset R[x_1, \ldots, x_n]$ est un idéal de type fini et
$\mathfrak q \subset R[x_1, \ldots, x_n]/I$ est premier.
(Notons que $R \cap \mathfrak q$ est égal à l'idéal maximal
$\mathfrak m$ de $R$.) Choisissons aussi des générateurs
$f_1, \ldots, f_m \in I$ de l'idéal $I$.
Écrivons de façon quelconque $R$ comme une colimite
$R = \colim R_\lambda$ indexée par un ensemble filtrant
$(\Lambda, \leq )$, chaque $R_\lambda$ étant local et essentiellement de
type fini sur $\mathbf{Z}$. Il existe $\lambda_0 \in \Lambda$ tel que
$f_j$ soit l'image d'un certain
$f_{j, \lambda_0} \in R_{\lambda_0}[x_1, \ldots, x_n]$.
Pour tout $\lambda \geq \lambda_0$, notons
$f_{j, \lambda} \in R_{\lambda}[x_1, \ldots, x_n]$ l'image de
$f_{j, \lambda_0}$. Nous obtenons ainsi un système de morphismes d'anneaux
$$
R_\lambda[x_1, \ldots, x_n]/(f_{1, \lambda}, \ldots, f_{m, \lambda})
\to
R[x_1, \ldots, x_n]/(f_1, \ldots, f_m) \to S
$$
Posons $\mathfrak q_\lambda$ égal à l'image réciproque de $\mathfrak q$.
Posons
$S_\lambda = (R_\lambda[x_1, \ldots, x_n]/
(f_{1, \lambda}, \ldots, f_{m, \lambda}))_{\mathfrak q_\lambda}$.
Nous laissons au lecteur le soin de vérifier que ce choix convient.
\end{proof}
\begin{remark}
\label{algebra-remark-suitable-systems-limits}
Supposons que $R \to S$ soit un homomorphisme local d'anneaux locaux,
essentiellement de présentation finie. Prenons un système quelconque
$(\Lambda, \leq)$, $R_\lambda \to S_\lambda$, possédant les propriétés
énumérées dans le Lemme \ref{algebra-lemma-limit-essentially-finite-type}.
Il peut arriver que ce soit le « mauvais » système, c'est-à-dire que la
propriété (4) du Lemme \ref{algebra-lemma-limit-essentially-finite-presentation}
ne soit pas satisfaite. En voici un exemple.
Soit $k = \mathbf{F}_2$. Considérons l'anneau
$$
R = \text{localisation de }
k[z, y_1, y_2, \ldots]/(y_i^2 - zy_{i + 1})
\text{ en }(z, y_1, y_2, \ldots)
$$
Posons $S = R/zR$. Comme système, prenons $\Lambda = \mathbf{N}$ et
$$
R_n = \text{localisation de }
k[z, y_1, \ldots, y_n]/(\{y_i^2 - zy_{i + 1}\}_{i \leq n-1})
\text{ en }(z, y_1, \ldots, y_n)
$$
et $S_n = R_n/(z, y_n^2)$. Aucun des morphismes
$S_n \otimes_{R_n} R_{n + 1} \to S_{n + 1}$
n'est une localisation (c'est-à-dire, en l'espèce, un isomorphisme),
puisque $1 \otimes y_{n + 1}^2$ s'envoie sur zéro.
Si l'on prend plutôt $S_n' = R_n/zR_n$, alors les morphismes
$S'_n \otimes_{R_n} R_{n + 1} \to S'_{n + 1}$ sont des isomorphismes.
La morale de cette remarque est qu'il faut choisir les systèmes avec
un peu de soin.
\end{remark}
\begin{lemma}
\label{algebra-lemma-limit-module-essentially-finite-presentation}
Supposons que $R \to S$ soit un homomorphisme local d'anneaux locaux et que
$S$ soit essentiellement de présentation finie sur $R$.
Soit $M$ un $S$-module de présentation finie.
Il existe alors un ensemble filtrant $(\Lambda, \leq)$, un système
d'homomorphismes locaux $R_\lambda \to S_\lambda$ d'anneaux locaux et des
$S_\lambda$-modules $M_\lambda$ tels que :
\begin{enumerate}
\item La colimite du système $R_\lambda \to S_\lambda$ est égale à
$R \to S$. La colimite du système des $M_\lambda$ est $M$.
\item Chaque $R_\lambda$ est essentiellement de type fini sur
$\mathbf{Z}$.
\item Chaque $S_\lambda$ est essentiellement de type fini sur
$R_\lambda$.
\item Chaque $M_\lambda$ est de type fini sur $S_\lambda$.
\item Pour chaque $\lambda \leq \mu$, le morphisme
$S_\lambda \otimes_{R_\lambda} R_\mu \to S_\mu$
présente $S_\mu$ comme la localisation de
$S_\lambda \otimes_{R_\lambda} R_\mu$ en un idéal premier.
\item Pour chaque $\lambda \leq \mu$, le morphisme
$M_\lambda \otimes_{S_\lambda} S_\mu \to M_\mu$
est un isomorphisme.
\end{enumerate}
\end{lemma}

\begin{proof}
Comme dans la démonstration du
Lemme \ref{algebra-lemma-limit-essentially-finite-presentation}, nous pouvons
d'abord écrire $R = \colim R_\lambda$ comme une colimite filtrante de
$\mathbf{Z}$-algèbres locales essentiellement de type fini.
Ensuite, nous pouvons supposer qu'il existe $\lambda_1 \in \Lambda$ et des
$f_{j, \lambda_1} \in R_{\lambda_1}[x_1, \ldots, x_n]$ tels que
$$
S =
\colim_{\lambda \geq \lambda_1} S_\lambda, \text{ avec }
S_\lambda =
(R_\lambda[x_1, \ldots, x_n]/
(f_{1, \lambda}, \ldots, f_{m, \lambda}))_{\mathfrak q_\lambda}
$$
Choisissons une présentation
$$
S^{\oplus s} \to S^{\oplus t} \to M \to 0
$$
de $M$ sur $S$. Soit $A \in \text{Mat}(t \times s, S)$ la matrice de la
présentation. Pour un certain $\lambda_2 \in \Lambda$,
$\lambda_2 \geq \lambda_1$, nous pouvons trouver une matrice
$A_{\lambda_2} \in \text{Mat}(t \times s, S_{\lambda_2})$ qui s'envoie sur
$A$. Pour tout $\lambda \geq \lambda_2$, posons
$M_\lambda = \Coker(S_\lambda^{\oplus s} \xrightarrow{A_\lambda}
S_\lambda^{\oplus t})$. Nous laissons au lecteur le soin de vérifier que
ce choix convient.
\end{proof}
\begin{lemma}
\label{algebra-lemma-limit-no-condition}
Supposons que $R \to S$ soit un morphisme d'anneaux.
Il existe alors un ensemble filtrant $(\Lambda, \leq)$ et un système de
morphismes d'anneaux $R_\lambda \to S_\lambda$ tels que :
\begin{enumerate}
\item La colimite du système $R_\lambda \to S_\lambda$ est égale à
$R \to S$.
\item Chaque $R_\lambda$ est de type fini sur $\mathbf{Z}$.
\item Chaque $S_\lambda$ est de type fini sur $R_\lambda$.
\end{enumerate}
\end{lemma}

\begin{proof}
C'est la version non locale du Lemme \ref{algebra-lemma-limit-no-condition-local}.
La démonstration est analogue et laissée au lecteur.
\end{proof}
\begin{lemma}
\label{algebra-lemma-limit-integral}
Supposons que $R \to S$ soit un morphisme d'anneaux et que $S$ soit entier
sur $R$. Il existe alors un ensemble filtrant $(\Lambda, \leq)$ et un
système de morphismes d'anneaux $R_\lambda \to S_\lambda$ tels que :
\begin{enumerate}
\item La colimite du système $R_\lambda \to S_\lambda$ est égale à
$R \to S$.
\item Chaque $R_\lambda$ est de type fini sur $\mathbf{Z}$.
\item Chaque $S_\lambda$ est fini sur $R_\lambda$.
\end{enumerate}
\end{lemma}

\begin{proof}
Considérons l'ensemble $\Lambda$ des couples $(E, F)$ où $E \subset R$ et
$F \subset S$ sont des parties finies, et où tout élément $f \in F$ est
racine d'un polynôme unitaire $P(X) \in R[X]$ dont les coefficients
appartiennent à $E$. Posons $(E, F) \leq (E', F')$ si $E \subset E'$ et
$F \subset F'$.
Étant donné $\lambda = (E, F) \in \Lambda$, soit
$R_\lambda \subset R$ la sous-$\mathbf{Z}$-algèbre de $R$ engendrée par
$E$, et soit $S_\lambda \subset S$ la sous-$\mathbf{Z}$-algèbre engendrée
par $F$ et par l'image de $E$ dans $S$.
Il est clair que $R = \colim R_\lambda$. On a
$S = \colim S_\lambda$, puisque tout élément de $S$ est entier sur $S$.
Les morphismes d'anneaux $R_\lambda \to S_\lambda$ sont finis d'après le
Lemme \ref{algebra-lemma-characterize-finite-in-terms-of-integral} et parce que
$S_\lambda$ est engendrée sur $R_\lambda$ par les éléments de $F$, qui sont
entiers sur $R_\lambda$ en vertu de notre condition sur les couples
$(E, F)$. Le lemme en résulte.
\end{proof}
\begin{lemma}
\label{algebra-lemma-limit-finite-type}
Supposons que $R \to S$ soit un morphisme d'anneaux et que $S$ soit de type
fini sur $R$. Il existe alors un ensemble filtrant $(\Lambda, \leq)$ et un
système de morphismes d'anneaux $R_\lambda \to S_\lambda$ tels que :
\begin{enumerate}
\item La colimite du système $R_\lambda \to S_\lambda$ est égale à
$R \to S$.
\item Chaque $R_\lambda$ est de type fini sur $\mathbf{Z}$.
\item Chaque $S_\lambda$ est de type fini sur $R_\lambda$.
\item Pour chaque $\lambda \leq \mu$, le morphisme
$S_\lambda \otimes_{R_\lambda} R_\mu \to S_\mu$
présente $S_\mu$ comme un quotient de
$S_\lambda \otimes_{R_\lambda} R_\mu$.
\end{enumerate}
\end{lemma}

\begin{proof}
C'est la version non locale du
Lemme \ref{algebra-lemma-limit-essentially-finite-type}.
La démonstration est analogue et laissée au lecteur.
\end{proof}
\begin{lemma}
\label{algebra-lemma-limit-finite-presentation}
Supposons que $R \to S$ soit un morphisme d'anneaux et que $S$ soit de
présentation finie sur $R$. Il existe alors un ensemble filtrant
$(\Lambda, \leq)$ et un système de morphismes d'anneaux
$R_\lambda \to S_\lambda$ tels que :
\begin{enumerate}
\item La colimite du système $R_\lambda \to S_\lambda$ est égale à
$R \to S$.
\item Chaque $R_\lambda$ est de type fini sur $\mathbf{Z}$.
\item Chaque $S_\lambda$ est de type fini sur $R_\lambda$.
\item Pour chaque $\lambda \leq \mu$, le morphisme
$S_\lambda \otimes_{R_\lambda} R_\mu \to S_\mu$
est un isomorphisme.
\end{enumerate}
\end{lemma}

\begin{proof}
C'est la version non locale du
Lemme \ref{algebra-lemma-limit-essentially-finite-presentation}.
La démonstration est analogue et laissée au lecteur.
\end{proof}
\begin{lemma}
\label{algebra-lemma-limit-module-finite-presentation}
Supposons que $R \to S$ soit un morphisme d'anneaux et que $S$ soit de
présentation finie sur $R$. Soit $M$ un $S$-module de présentation finie.
Il existe alors un ensemble filtrant $(\Lambda, \leq)$, un système de
morphismes d'anneaux $R_\lambda \to S_\lambda$ et des $S_\lambda$-modules
$M_\lambda$ tels que :
\begin{enumerate}
\item La colimite du système $R_\lambda \to S_\lambda$ est égale à
$R \to S$. La colimite du système des $M_\lambda$ est $M$.
\item Chaque $R_\lambda$ est de type fini sur $\mathbf{Z}$.
\item Chaque $S_\lambda$ est de type fini sur $R_\lambda$.
\item Chaque $M_\lambda$ est de type fini sur $S_\lambda$.
\item Pour chaque $\lambda \leq \mu$, le morphisme
$S_\lambda \otimes_{R_\lambda} R_\mu \to S_\mu$
est un isomorphisme.
\item Pour chaque $\lambda \leq \mu$, le morphisme
$M_\lambda \otimes_{S_\lambda} S_\mu \to M_\mu$
est un isomorphisme.
\end{enumerate}
En particulier, pour tout $\lambda \in \Lambda$, on a
$$
M = M_\lambda \otimes_{S_\lambda} S
= M_\lambda \otimes_{R_\lambda} R.
$$
\end{lemma}

\begin{proof}
C'est la version non locale du
Lemme \ref{algebra-lemma-limit-module-essentially-finite-presentation}.
La démonstration est analogue et laissée au lecteur.
\end{proof}








\section{Autres critères de platitude}
\label{algebra-section-more-flatness-criteria}

\noindent
Le lemme suivant est souvent utilisé en géométrie algébrique pour montrer
qu'un morphisme fini d'une surface normale vers une surface lisse est plat.
Il constitue une réciproque partielle du
Lemme \ref{algebra-lemma-finite-flat-over-regular-CM}, puisqu'un homomorphisme local
injectif et fini d'anneaux satisfait assurément la condition (3).

\begin{lemma}
\label{algebra-lemma-CM-over-regular-flat}
\begin{slogan}
Platitude miraculeuse
\end{slogan}
Soit $R \to S$ un homomorphisme local d'anneaux locaux noethériens.
Supposons que
\begin{enumerate}
\item $R$ soit régulier,
\item $S$ soit de Cohen-Macaulay,
\item $\dim(S) = \dim(R) + \dim(S/\mathfrak m_R S)$.
\end{enumerate}
Alors $R \to S$ est plat.
\end{lemma}

\begin{proof}
Raisonnons par récurrence sur $\dim(R)$. Le cas $\dim(R) = 0$ est trivial,
car $R$ est alors un corps. Supposons $\dim(R) > 0$. D'après (3), cela
implique que $\dim(S) > 0$. Soient
$\mathfrak q_1, \ldots, \mathfrak q_r$ les idéaux premiers minimaux de $S$.
Notons que $\mathfrak q_i \not \supset \mathfrak m_R S$, puisque
$$
\dim(S/\mathfrak q_i) = \dim(S) > \dim(S/\mathfrak m_R S)
$$
la première égalité résultant du Lemme \ref{algebra-lemma-maximal-chain-CM}, et
l'inégalité de (3). Ainsi,
$\mathfrak p_i = R \cap \mathfrak q_i$ n'est pas égal à $\mathfrak m_R$.
Choisissons $x \in \mathfrak m_R$, $x \not \in \mathfrak m_R^2$ et
$x \not \in \mathfrak p_i$ ; voir le Lemme \ref{algebra-lemma-silly}.
On voit donc que $x$ n'appartient à aucun idéal premier minimal de $S$.
Par conséquent, $x$ est un non-diviseur de zéro sur $S$ d'après (2) ; voir
le Lemme \ref{algebra-lemma-reformulate-CM}. En outre, $S/xS$ est de Cohen-Macaulay
et $\dim(S/xS) = \dim(S) - 1$.
D'après (1) et le Lemme \ref{algebra-lemma-regular-ring-CM}, l'anneau $R/xR$ est
régulier et $\dim(R/xR) = \dim(R) - 1$.
L'hypothèse de récurrence montre que $R/xR \to S/xS$ est plat.
On conclut alors par le Lemme \ref{algebra-lemma-variant-local-criterion-flatness}
et la remarque qui le suit.
\end{proof}

\begin{lemma}
\label{algebra-lemma-flat-over-regular}
Soit $R \to S$ un homomorphisme d'anneaux locaux noethériens.
Supposons que $R$ soit un anneau local régulier et qu'un système régulier
de paramètres s'envoie sur une suite régulière dans $S$. Alors
$R \to S$ est plat.
\end{lemma}

\begin{proof}
Supposons que $x_1, \ldots, x_d$ soit un système de paramètres de $R$
qui s'envoie sur une suite régulière dans $S$. Notons que
$S/(x_1, \ldots, x_d)S$ est plat sur $R/(x_1, \ldots, x_d)$, puisque ce
dernier est un corps. Puis $x_d$ est un non-diviseur de zéro dans
$S/(x_1, \ldots, x_{d - 1})S$ ; par conséquent,
$S/(x_1, \ldots, x_{d - 1})S$ est plat sur
$R/(x_1, \ldots, x_{d - 1})$ d'après le critère local de platitude (voir le
Lemme \ref{algebra-lemma-variant-local-criterion-flatness} et les remarques qui le
suivent). Puis $x_{d - 1}$ est un non-diviseur de zéro dans
$S/(x_1, \ldots, x_{d - 2})S$ ; ainsi,
$S/(x_1, \ldots, x_{d - 2})S$ est plat sur
$R/(x_1, \ldots, x_{d - 2})$ d'après le critère local de platitude (voir le
Lemme \ref{algebra-lemma-variant-local-criterion-flatness} et les remarques qui le
suivent). On continue ainsi jusqu'à conclure que $S$ est plat sur $R$.
\end{proof}

\noindent
Le lemme suivant est la clef qui permet de montrer que les résultats sur les
modules de présentation finie sur des anneaux de présentation finie sur un
anneau de base découlent des résultats correspondants sur les modules de type fini
dans le cas noethérien.

\begin{lemma}
\label{algebra-lemma-colimit-eventually-flat}
Soient $R \to S$, $M$, $\Lambda$, $R_\lambda \to S_\lambda$ et
$M_\lambda$ comme dans le
Lemme \ref{algebra-lemma-limit-module-essentially-finite-presentation}.
Supposons que $M$ soit plat sur $R$. Alors, pour un certain
$\lambda \in \Lambda$, le module $M_\lambda$ est plat sur $R_\lambda$.
\end{lemma}

\begin{proof}
Choisissons $\lambda \in \Lambda$ et considérons
$$
\text{Tor}_1^{R_\lambda}(M_\lambda, R_\lambda/\mathfrak m_\lambda)
=
\Ker(\mathfrak m_\lambda \otimes_{R_\lambda} M_\lambda
\to M_\lambda).
$$
Voir la Remarque \ref{algebra-remark-Tor-ring-mod-ideal}. Le second membre montre
qu'il s'agit d'un $S_\lambda$-module de type fini (car $S_\lambda$ est
noethérien et les modules en question sont de type fini).
Soient $\xi_1, \ldots, \xi_n$ des générateurs.
Comme $M$ est plat sur $R$, on a
$0 = \Ker(\mathfrak m_\lambda R \otimes_R M \to M)$.
Puisque $\otimes$ commute aux colimites, il existe
$\lambda' \geq \lambda$ tel que chaque $\xi_i$ s'envoie sur zéro dans
$\mathfrak m_{\lambda}R_{\lambda'} \otimes_{R_{\lambda'}} M_{\lambda'}$.
On voit donc que
$$
\text{Tor}_1^{R_\lambda}(M_\lambda, R_\lambda/\mathfrak m_\lambda)
\longrightarrow
\text{Tor}_1^{R_{\lambda'}}(M_{\lambda'},
R_{\lambda'}/\mathfrak m_{\lambda}R_{\lambda'})
$$
est nul. Notons que
$M_\lambda \otimes_{R_\lambda} R_\lambda/\mathfrak m_\lambda$
est plat sur $R_\lambda/\mathfrak m_\lambda$, puisque ce dernier anneau
est un corps. Nous pouvons donc appliquer le
Lemme \ref{algebra-lemma-another-variant-local-criterion-flatness} et conclure que
$M_{\lambda'}$ est plat sur $R_{\lambda'}$.
\end{proof}

\noindent
À l'aide du lemme précédent, nous pouvons commencer à redémontrer les
résultats de la section \ref{algebra-section-criteria-flatness} dans le cas non
noethérien.

\begin{lemma}
\label{algebra-lemma-mod-injective-general}
Supposons que $R \to S$ soit un homomorphisme local d'anneaux locaux.
Notons $\mathfrak m$ l'idéal maximal de $R$.
Soit $u : M \to N$ un morphisme de $S$-modules. Supposons que
\begin{enumerate}
\item $S$ soit essentiellement de présentation finie sur $R$,
\item $M$ et $N$ soient de présentation finie sur $S$,
\item $N$ soit plat sur $R$, et
\item $\overline{u} : M/\mathfrak mM \to N/\mathfrak mN$ soit injectif.
\end{enumerate}
Alors $u$ est injectif et $N/u(M)$ est plat sur $R$.
\end{lemma}

\begin{proof}
D'après le Lemme \ref{algebra-lemma-limit-module-essentially-finite-presentation}
et sa démonstration, on peut trouver un système de morphismes locaux
$R_\lambda \to S_\lambda$, ainsi que des morphismes de $S_\lambda$-modules
$u_\lambda : M_\lambda \to N_\lambda$, qui satisfont, pour $N$ comme pour
$M$, aux conclusions (1) -- (6) de ce lemme, et tels que la colimite des
morphismes $u_\lambda$ soit $u$.
D'après le Lemme \ref{algebra-lemma-colimit-eventually-flat}, nous pouvons supposer
que $N_\lambda$ est plat sur $R_\lambda$ pour tout $\lambda$ suffisamment
grand. Notons $\mathfrak m_\lambda \subset R_\lambda$ l'idéal maximal et
$\kappa_\lambda = R_\lambda / \mathfrak m_\lambda$, et notons
resp.\ $\kappa = R/\mathfrak m$ les corps résiduels.

\medskip\noindent
Considérons le morphisme
$$
\Psi_\lambda :
M_\lambda/\mathfrak m_\lambda M_\lambda \otimes_{\kappa_\lambda} \kappa
\longrightarrow
M/\mathfrak m M.
$$
Puisque $S_\lambda/\mathfrak m_\lambda S_\lambda$ est essentiellement de
type fini sur le corps $\kappa_\lambda$, le produit tensoriel
$S_\lambda/\mathfrak m_\lambda S_\lambda
\otimes_{\kappa_\lambda} \kappa$ est essentiellement de type fini sur
$\kappa$. C'est donc un anneau noethérien, et nous en concluons que le noyau
de $\Psi_\lambda$ est de type fini.
Comme $M/\mathfrak m M$ est la colimite du système des
$M_\lambda/\mathfrak m_\lambda M_\lambda$ et que $\kappa$ est la colimite
des corps $\kappa_\lambda$, il existe $\lambda' > \lambda$ tel que le noyau
de $\Psi_\lambda$ soit engendré par le noyau de
$$
\Psi_{\lambda, \lambda'} :
M_\lambda/\mathfrak m_\lambda M_\lambda
\otimes_{\kappa_\lambda}
\kappa_{\lambda'}
\longrightarrow
M_{\lambda'}/\mathfrak m_{\lambda'} M_{\lambda'}.
$$
Par construction, il existe une partie multiplicative
$W \subset S_\lambda \otimes_{R_\lambda} R_{\lambda'}$ telle que
$S_{\lambda'} = W^{-1}(S_\lambda \otimes_{R_\lambda} R_{\lambda'})$ et
$$
W^{-1}(M_\lambda/\mathfrak m_\lambda M_\lambda
\otimes_{\kappa_\lambda}
\kappa_{\lambda'})
=
M_{\lambda'}/\mathfrak m_{\lambda'} M_{\lambda'}.
$$
Supposons maintenant que $x$ soit un élément du noyau de
$$
\Psi_{\lambda'} :
M_{\lambda'}/\mathfrak m_{\lambda'} M_{\lambda'}
\otimes_{\kappa_{\lambda'}} \kappa
\longrightarrow
M/\mathfrak m M.
$$
Alors, pour un certain $w \in W$, on a
$wx \in M_\lambda/\mathfrak m_\lambda M_\lambda \otimes \kappa$.
Ainsi $wx \in \Ker(\Psi_\lambda)$. Par conséquent, $wx$ est une
combinaison linéaire d'éléments du noyau de
$\Psi_{\lambda, \lambda'}$. Il s'ensuit que
$wx = 0$ dans $M_{\lambda'}/\mathfrak m_{\lambda'} M_{\lambda'}
\otimes_{\kappa_{\lambda'}} \kappa$, donc $x = 0$ puisque $w$ est
inversible dans $S_{\lambda'}$.
Nous concluons que le noyau de $\Psi_{\lambda'}$ est nul pour tout
$\lambda'$ suffisamment grand !

\medskip\noindent
Le résultat du paragraphe précédent nous permet de supposer que le noyau
de $\Psi_\lambda$ est nul pour tout $\lambda$ suffisamment grand, ce qui
implique que le morphisme
$M_\lambda/\mathfrak m_\lambda M_\lambda \to M/\mathfrak m M$
est injectif. Combiné à l'injectivité de $\overline{u}$, cela implique
formellement que
$\overline{u_\lambda} : M_\lambda/\mathfrak m_\lambda M_\lambda
\to N_\lambda/\mathfrak m_\lambda N_\lambda$ est lui aussi injectif.
D'après le Lemme \ref{algebra-lemma-mod-injective}, nous concluons que, pour tout
$\lambda$ suffisamment grand, le morphisme $u_\lambda$ est injectif et que
$N_\lambda/u_\lambda(M_\lambda)$ est plat sur $R_\lambda$.
Le lemme en résulte.
\end{proof}

\begin{lemma}
\label{algebra-lemma-grothendieck-general}
Supposons que $R \to S$ soit un homomorphisme local d'anneaux locaux.
Notons $\mathfrak m$ l'idéal maximal de $R$. Supposons que
\begin{enumerate}
\item $S$ soit essentiellement de présentation finie sur $R$,
\item $S$ soit plat sur $R$, et
\item $f \in S$ soit un non-diviseur de zéro dans $S/{\mathfrak m}S$.
\end{enumerate}
Alors $S/fS$ est plat sur $R$, et $f$ est un non-diviseur de zéro dans $S$.
\end{lemma}

\begin{proof}
Cela résulte directement du Lemme \ref{algebra-lemma-mod-injective-general}.
\end{proof}

\begin{lemma}
\label{algebra-lemma-grothendieck-regular-sequence-general}
Supposons que $R \to S$ soit un homomorphisme local d'anneaux locaux.
Notons $\mathfrak m$ l'idéal maximal de $R$. Supposons que
\begin{enumerate}
\item $R \to S$ soit essentiellement de présentation finie,
\item $R \to S$ soit plat, et
\item $f_1, \ldots, f_c$ soit une suite d'éléments de $S$ telle que les
images $\overline{f}_1, \ldots, \overline{f}_c$ forment une suite régulière
dans $S/{\mathfrak m}S$.
\end{enumerate}
Alors $f_1, \ldots, f_c$ est une suite régulière dans $S$, et chacun des
quotients $S/(f_1, \ldots, f_i)$ est plat sur $R$.
\end{lemma}

\begin{proof}
Cela résulte d'une récurrence et du Lemme \ref{algebra-lemma-grothendieck-general}.
\end{proof}

\noindent
Voici la version du critère local de platitude pour les morphismes locaux
d'anneaux qui sont localement de présentation finie.

\begin{lemma}
\label{algebra-lemma-variant-local-criterion-flatness-general}
Soit $R \to S$ un homomorphisme local d'anneaux locaux.
Soit $I \not = R$ un idéal de $R$. Soit $M$ un $S$-module. Supposons que
\begin{enumerate}
\item $S$ soit essentiellement de présentation finie sur $R$,
\item $M$ soit de présentation finie sur $S$,
\item $\text{Tor}_1^R(M, R/I) = 0$, et
\item $M/IM$ soit plat sur $R/I$.
\end{enumerate}
Alors $M$ est plat sur $R$.
\end{lemma}

\begin{proof}
Soient $\Lambda$, $R_\lambda \to S_\lambda$ et $M_\lambda$ comme dans le
Lemme \ref{algebra-lemma-limit-module-essentially-finite-presentation}.
Notons $I_\lambda \subset R_\lambda$ l'image réciproque de $I$.
Dans ce cas, le système constitué par
$R/I \to S/IS$, $M/IM$, $R_\lambda \to S_\lambda/I_\lambda S_\lambda$
et $M_\lambda/I_\lambda M_\lambda$ satisfait lui aussi les conclusions du
Lemme \ref{algebra-lemma-limit-module-essentially-finite-presentation}.
D'après le Lemme \ref{algebra-lemma-colimit-eventually-flat}, nous pouvons donc
supposer, après avoir restreint l'ensemble d'indices $\Lambda$, que
$M_\lambda/I_\lambda M_\lambda$ est plat pour tout $\lambda$.
Choisissons un $\lambda$ et considérons
$$
\text{Tor}_1^{R_\lambda}(M_\lambda, R_\lambda/I_\lambda)
=
\Ker(I_\lambda \otimes_{R_\lambda} M_\lambda
\to M_\lambda).
$$
Voir la Remarque \ref{algebra-remark-Tor-ring-mod-ideal}. Le second membre montre
qu'il s'agit d'un $S_\lambda$-module de type fini (car $S_\lambda$ est
noethérien et les modules en question sont de type fini).
Soient $\xi_1, \ldots, \xi_n$ des générateurs.
Comme $\text{Tor}_1^R(M, R/I) = 0$ et puisque $\otimes$ commute aux
colimites, il existe $\lambda' \geq \lambda$ tel que chaque $\xi_i$
s'envoie sur zéro dans
$\text{Tor}_1^{R_{\lambda'}}(M_{\lambda'},
R_{\lambda'}/I_{\lambda'})$.
Le composé des morphismes
$$
\xymatrix{
R_{\lambda'} \otimes_{R_\lambda}
\text{Tor}_1^{R_\lambda}(M_\lambda, R_\lambda/I_\lambda)
\ar[d]^{\text{surjectif d'après le lemme \ref{algebra-lemma-surjective-on-tor-one}}} \\
\text{Tor}_1^{R_\lambda}(M_\lambda, R_{\lambda'}/I_\lambda R_{\lambda'})
\ar[d]^{\text{surjectif après localisation d'après le
lemme \ref{algebra-lemma-surjective-on-tor-one-trivial}}} \\
\text{Tor}_1^{R_{\lambda'}}(M_{\lambda'}, R_{\lambda'}/I_\lambda R_{\lambda'})
\ar[d]^{\text{surjectif d'après le lemme \ref{algebra-lemma-surjective-on-tor-one}}} \\
\text{Tor}_1^{R_{\lambda'}}(M_{\lambda'}, R_{\lambda'}/I_{\lambda'}).
}
$$
est surjectif après une localisation, pour les raisons indiquées.
La localisation est nécessaire parce que $M_{\lambda'}$ n'est pas égal à
$M_\lambda \otimes_{R_\lambda} R_{\lambda'}$. En effet, il est égal à
$M_\lambda \otimes_{S_\lambda} S_{\lambda'}$, et $S_{\lambda'}$ est la
localisation de $S_{\lambda} \otimes_{R_\lambda} R_{\lambda'}$, d'où
l'assertion après localisation (ou tensorisation par $S_{\lambda'}$).
Notons que le Lemme \ref{algebra-lemma-surjective-on-tor-one} s'applique à la
première et à la troisième flèche, parce que
$M_\lambda/I_\lambda M_\lambda$ est plat sur $R_\lambda/I_\lambda$ et que
$M_{\lambda'}/I_\lambda M_{\lambda'}$ est plat sur
$R_{\lambda'}/I_\lambda R_{\lambda'}$, car c'est un changement de base du
module plat $M_\lambda/I_\lambda M_\lambda$.
Le composé envoie les générateurs $\xi_i$ sur zéro, comme nous l'avons
expliqué ci-dessus. Nous concluons enfin que
$\text{Tor}_1^{R_{\lambda'}}(M_{\lambda'},
R_{\lambda'}/I_{\lambda'})$ est nul. Cela implique que $M_{\lambda'}$
est plat sur $R_{\lambda'}$ d'après le
Lemme \ref{algebra-lemma-variant-local-criterion-flatness}.
\end{proof}

\noindent
Comparer le lemme ci-dessous au
Lemme \ref{algebra-lemma-criterion-flatness-fibre-Noetherian}
(cas des anneaux locaux noethériens) et au
Lemme \ref{algebra-lemma-criterion-flatness-fibre-nilpotent}
(cas d'un idéal nilpotent dans la base).

\begin{lemma}[Crit\`ere de platitude par fibres]
\label{algebra-lemma-criterion-flatness-fibre}
Soient $R$, $S$ et $S'$ des anneaux locaux, et soient
$R \to S \to S'$ des homomorphismes locaux d'anneaux.
Soit $M$ un $S'$-module. Soit $\mathfrak m \subset R$ l'idéal maximal.
Supposons que
\begin{enumerate}
\item les morphismes d'anneaux $R \to S$ et $R \to S'$ soient
essentiellement de présentation finie ;
\item le module $M$ soit de présentation finie sur $S'$ ;
\item le module $M$ ne soit pas nul ;
\item le module $M/\mathfrak mM$ soit un $S/\mathfrak mS$-module plat ;
\item le module $M$ soit un $R$-module plat.
\end{enumerate}
Alors $S$ est plat sur $R$ et $M$ est un $S$-module plat.
\end{lemma}

\begin{proof}
Comme dans la démonstration du
Lemme \ref{algebra-lemma-limit-essentially-finite-presentation}, nous pouvons
d'abord écrire $R = \colim R_\lambda$ comme une colimite filtrante de
$\mathbf{Z}$-algèbres locales essentiellement de type fini.
Notons $\mathfrak p_\lambda$ l'idéal maximal de $R_\lambda$.
Ensuite, nous pouvons supposer que, pour un certain
$\lambda_1 \in \Lambda$, il existe des
$f_{j, \lambda_1} \in R_{\lambda_1}[x_1, \ldots, x_n]$ tels que
$$
S =
\colim_{\lambda \geq \lambda_1} S_\lambda, \text{ avec }
S_\lambda =
(R_\lambda[x_1, \ldots, x_n]/
(f_{1, \lambda}, \ldots, f_{u, \lambda}))_{\mathfrak q_\lambda}
$$
Pour un certain $\lambda_2 \in \Lambda$, $\lambda_2 \geq \lambda_1$, il
existe des
$g_{j, \lambda_2} \in R_{\lambda_2}[x_1, \ldots, x_n, y_1, \ldots, y_m]$
d'images
$\overline{g}_{j, \lambda_2} \in S_{\lambda_2}[y_1, \ldots, y_m]$
telles que
$$
S' =
\colim_{\lambda \geq \lambda_2} S'_\lambda, \text{ avec }
S'_\lambda =
(S_\lambda[y_1, \ldots, y_m]/
(\overline{g}_{1, \lambda}, \ldots,
\overline{g}_{v, \lambda}))_{\overline{\mathfrak q}'_\lambda}
$$
Notons que cela implique aussi
$$
S'_\lambda =
(R_\lambda[x_1, \ldots, x_n, y_1, \ldots, y_m]/
(g_{1, \lambda}, \ldots, g_{v, \lambda}))_{\mathfrak q'_\lambda}
$$
Choisissons une présentation
$$
(S')^{\oplus s} \to (S')^{\oplus t} \to M \to 0
$$
de $M$ sur $S'$. Soit $A \in \text{Mat}(t \times s, S')$ la matrice de
la présentation. Pour un certain $\lambda_3 \in \Lambda$,
$\lambda_3 \geq \lambda_2$, nous pouvons trouver une matrice
$A_{\lambda_3} \in \text{Mat}(t \times s, S_{\lambda_3})$ qui s'envoie sur
$A$. Pour tout $\lambda \geq \lambda_3$, posons
$M_\lambda = \Coker((S'_\lambda)^{\oplus s} \xrightarrow{A_\lambda}
(S'_\lambda)^{\oplus t})$.

\medskip\noindent
Avec ces choix, pour tous $\lambda_3 \leq \lambda \leq \mu$, le morphisme
$S_\lambda \otimes_{R_{\lambda}} R_\mu \to S_\mu$ est une localisation,
$S'_\lambda \otimes_{S_{\lambda}} S_\mu \to S'_\mu$ est une localisation,
et le morphisme $M_\lambda \otimes_{S'_\lambda} S'_\mu \to M_\mu$ est un
isomorphisme. Cela implique aussi que
$S'_\lambda \otimes_{R_{\lambda}} R_\mu \to S'_\mu$ est une localisation.
Ainsi, puisque $M$ est plat sur $R$, le
Lemme \ref{algebra-lemma-colimit-eventually-flat} montre que, pour tout $\lambda$
assez grand, le module $M_\lambda$ est plat sur $R_\lambda$.
Notons en outre que
$
\mathfrak m = \colim \mathfrak p_\lambda
$,
$
S/\mathfrak mS = \colim S_\lambda/\mathfrak p_\lambda S_\lambda
$,
$
S'/\mathfrak mS' = \colim S'_\lambda/\mathfrak p_\lambda S'_\lambda
$,
et
$
M/\mathfrak mM = \colim M_\lambda/\mathfrak p_\lambda M_\lambda
$. De plus, pour tous $\lambda_3 \leq \lambda \leq \mu$, les propriétés
énumérées ci-dessus montrent que
$$
S'_\lambda/\mathfrak p_\lambda S'_\lambda
\otimes_{S_{\lambda}/\mathfrak p_\lambda S_\lambda}
S_\mu/\mathfrak p_\mu S_\mu
\longrightarrow
S'_\mu/\mathfrak p_\mu S'_\mu
$$
est une localisation, et que le morphisme
$$
M_\lambda / \mathfrak p_\lambda M_\lambda
\otimes_{S'_\lambda/\mathfrak p_\lambda S'_\lambda}
S'_\mu /\mathfrak p_\mu S'_\mu
\longrightarrow
M_\mu/\mathfrak p_\mu M_\mu
$$
est un isomorphisme. Par conséquent, le système
$(S_\lambda/\mathfrak p_\lambda S_\lambda \to
S'_\lambda/\mathfrak p_\lambda S'_\lambda,
M_\lambda/\mathfrak p_\lambda M_\lambda)$ est lui aussi un système comme
dans le Lemme \ref{algebra-lemma-limit-module-essentially-finite-presentation}.
Nous pouvons appliquer une nouvelle fois le
Lemme \ref{algebra-lemma-colimit-eventually-flat}, car $M/\mathfrak m M$ est supposé
plat sur $S/\mathfrak mS$, et nous voyons que
$M_\lambda/\mathfrak p_\lambda M_\lambda$ est plat sur
$S_\lambda/\mathfrak p_\lambda S_\lambda$ pour tout $\lambda$ assez grand.
Ainsi, pour $\lambda$ assez grand, les données
$R_\lambda \to S_\lambda \to S'_\lambda, M_\lambda$ satisfont les
hypothèses du Lemme \ref{algebra-lemma-criterion-flatness-fibre-Noetherian}.
Choisissons un tel $\lambda$. Alors
$S = S_\lambda \otimes_{R_\lambda} R$ est plat sur $R$, et
$M = M_\lambda \otimes_{S_\lambda} S$ est plat sur $S$ (puisqu'un
changement de base d'un module plat est plat).
\end{proof}

\noindent
Le résultat suivant est une conséquence facile du « crit\`ere de platitude
par fibres », Lemme \ref{algebra-lemma-criterion-flatness-fibre}. Pour d'autres
résultats de ce genre, voir Compléments sur la platitude,
section \ref{flat-section-introduction}.

\begin{lemma}
\label{algebra-lemma-criterion-flatness-fibre-fp-over-ft}
Soient $R$, $S$ et $S'$ des anneaux locaux, et soient
$R \to S \to S'$ des homomorphismes locaux d'anneaux.
Soit $M$ un $S'$-module. Soit $\mathfrak m \subset R$ l'idéal maximal.
Supposons que
\begin{enumerate}
\item $R \to S'$ soit essentiellement de présentation finie,
\item $R \to S$ soit essentiellement de type fini,
\item $M$ soit de présentation finie sur $S'$,
\item $M$ ne soit pas nul,
\item $M/\mathfrak mM$ soit un $S/\mathfrak mS$-module plat, et
\item $M$ soit un $R$-module plat.
\end{enumerate}
Alors $S$ est essentiellement de présentation finie et plat sur $R$, et
$M$ est un $S$-module plat.
\end{lemma}

\begin{proof}
Comme $S$ est essentiellement de présentation finie sur $R$, nous pouvons
écrire $S = C_{\overline{\mathfrak q}}$ pour une certaine $R$-algèbre $C$
de type fini. Écrivons $C = R[x_1, \ldots, x_n]/I$.
Notons $\mathfrak q \subset R[x_1, \ldots, x_n]$ l'idéal premier
correspondant à $\overline{\mathfrak q}$. Alors $S = B/J$, où
$B = R[x_1, \ldots, x_n]_{\mathfrak q}$ est essentiellement de présentation
finie sur $R$ et $J = IB$. Nous pouvons trouver
$f_1, \ldots, f_k \in J$ tels que les images
$\overline{f}_i \in B/\mathfrak mB$ engendrent l'image $\overline{J}$ de
$J$ dans l'anneau noethérien $B/\mathfrak mB$.
Il existe donc des idéaux de type fini $J' \subset J$ tels que
$B/J' \to B/J$ induise un isomorphisme
$$
(B/J') \otimes_R R/\mathfrak m \longrightarrow
B/J \otimes_R R/\mathfrak m = S/\mathfrak mS.
$$
Pour tout $J'$ comme ci-dessus, le
Lemme \ref{algebra-lemma-criterion-flatness-fibre} s'applique aux morphismes
d'anneaux
$$
R \longrightarrow B/J' \longrightarrow S'
$$
et au module $M$. Nous concluons donc que $B/J'$ est plat sur $R$ pour tout
choix de $J'$ comme ci-dessus. Maintenant, si $J' \subset J' \subset J$
sont deux idéaux de type fini comme ci-dessus, nous concluons que
$B/J' \to B/J''$ est un morphisme surjectif entre des $R$-algèbres plates
et essentiellement de présentation finie, qui est un isomorphisme modulo
$\mathfrak m$. Par conséquent, le Lemme \ref{algebra-lemma-mod-injective-general}
implique que $B/J' = B/J''$, c'est-à-dire $J' = J''$.
Cela signifie clairement que $J$ est de type fini, c'est-à-dire que $S$ est
essentiellement de présentation finie sur $R$. Nous pouvons donc appliquer
le Lemme \ref{algebra-lemma-criterion-flatness-fibre} à $R \to S \to S'$, ce qui
achève la démonstration.
\end{proof}

\begin{lemma}[Crit\`ere de platitude par fibres : cas localement nilpotent]
\label{algebra-lemma-criterion-flatness-fibre-locally-nilpotent}
Soit
$$
\xymatrix{
S \ar[rr] & & S' \\
& R \ar[lu] \ar[ru]
}
$$
un diagramme commutatif dans la catégorie des anneaux.
Soient $I \subset R$ un idéal localement nilpotent et $M$ un $S'$-module.
Supposons que
\begin{enumerate}
\item $R \to S$ soit de type fini,
\item $R \to S'$ soit de présentation finie,
\item $M$ soit un $S'$-module de présentation finie,
\item $M/IM$ soit plat comme $S/IS$-module, et
\item $M$ soit plat comme $R$-module.
\end{enumerate}
Alors $M$ est un $S$-module plat, et $S_\mathfrak q$ est plat et
essentiellement de présentation finie sur $R$ pour tout
$\mathfrak q \subset S$ tel que
$M \otimes_S \kappa(\mathfrak q)$ soit non nul.
\end{lemma}

\begin{proof}
Si $M \otimes_S \kappa(\mathfrak q)$ est non nul, alors
$S' \otimes_S \kappa(\mathfrak q)$ est non nul ; il existe donc un idéal
premier $\mathfrak q' \subset S'$ au-dessus de $\mathfrak q$
(Lemme \ref{algebra-lemma-in-image}). Soit $\mathfrak p \subset R$ l'image de
$\mathfrak q$ dans $\Spec(R)$. Alors $I \subset \mathfrak p$, puisque $I$
est localement nilpotent ; ainsi $M/\mathfrak p M$ est plat sur
$S/\mathfrak pS$. Nous pouvons donc appliquer le
Lemme \ref{algebra-lemma-criterion-flatness-fibre-fp-over-ft} à
$R_\mathfrak p \to S_\mathfrak q \to S'_{\mathfrak q'}$ et à
$M_{\mathfrak q'}$. Nous concluons que $M_{\mathfrak q'}$ est plat sur $S$
et que $S_\mathfrak q$ est plat et essentiellement de présentation finie
sur $R$. Comme $\mathfrak q'$ était un idéal premier arbitraire de $S'$,
nous voyons aussi que $M$ est plat sur $S$
(Lemme \ref{algebra-lemma-flat-localization}).
\end{proof}














\section{Ouverture du lieu de platitude}
\label{algebra-section-open-flat}

\noindent
Nous utilisons le Lemme \ref{algebra-lemma-colimit-eventually-flat} pour nous ramener
au cas noethérien. Celui-ci est traité à l'aide de la caractérisation des
complexes exacts donnée dans la section \ref{algebra-section-complex-exact}.

\begin{lemma}
\label{algebra-lemma-CM-dim-finite-type}
Soit $k$ un corps. Soit $S$ une $k$-algèbre de type fini.
Soient $f_1, \ldots, f_i$ des éléments de $S$.
Supposons que $S$ soit de Cohen-Macaulay, équidimensionnelle de dimension
$d$, et que $\dim V(f_1, \ldots, f_i) \leq d - i$.
Alors il y a égalité, et $f_1, \ldots, f_i$ forment une suite régulière
dans $S_{\mathfrak q}$ pour tout idéal premier $\mathfrak q$ de
$V(f_1, \ldots, f_i)$.
\end{lemma}

\begin{proof}
Si $S$ est de Cohen-Macaulay et équidimensionnelle de dimension $d$, alors
$\dim(S_{\mathfrak m}) = d$ pour tout idéal maximal $\mathfrak m$ de $S$ ;
voir le Lemme \ref{algebra-lemma-disjoint-decomposition-CM-algebra}.
D'après la Proposition \ref{algebra-proposition-CM-module}, pour tout idéal maximal
$\mathfrak m \in V(f_1, \ldots, f_i)$, la suite est régulière dans
$S_{\mathfrak m}$ et l'anneau local
$S_{\mathfrak m}/(f_1, \ldots, f_i)$ est de Cohen-Macaulay, de dimension
$d - i$. Cela signifie en fait que $S/(f_1, \ldots, f_i)$ est de
Cohen-Macaulay et équidimensionnelle de dimension $d - i$.
\end{proof}

\begin{lemma}
\label{algebra-lemma-open-regular-sequence}
Soit $R \to S$ un morphisme d'anneaux de type fini. Soit $d$ un entier tel
que toutes les fibres $S \otimes_R \kappa(\mathfrak p)$ soient de
Cohen-Macaulay et équidimensionnelles de dimension $d$.
Soient $f_1, \ldots, f_i$ des éléments de $S$. L'ensemble
$$
\{ \mathfrak q \in V(f_1, \ldots, f_i)
\mid f_1, \ldots, f_i
\text{ forment une suite régulière dans }
S_{\mathfrak q}/\mathfrak p S_{\mathfrak q}
\text{ où }\mathfrak p = R \cap \mathfrak q
\}
$$
est ouvert dans $V(f_1, \ldots, f_i)$.
\end{lemma}

\begin{proof}
Écrivons $\overline{S} = S/(f_1, \ldots, f_i)$.
Supposons que $\mathfrak q$ appartienne à l'ensemble défini dans le lemme,
et soit $\mathfrak p$ l'idéal premier correspondant de $R$.
Nous utiliserons la dimension relative définie dans la
Définition \ref{algebra-definition-relative-dimension}.
Notons d'abord que
$d = \dim_{\mathfrak q}(S/R) =
\dim(S_{\mathfrak q}/\mathfrak pS_{\mathfrak q}) +
\text{trdeg}_{\kappa(\mathfrak p)}\ \kappa(\mathfrak q)$
d'après le Lemme \ref{algebra-lemma-dimension-at-a-point-finite-type-field}.
Puisque $f_1, \ldots, f_i$ forment une suite régulière dans l'anneau local
noethérien $S_{\mathfrak q}/\mathfrak pS_{\mathfrak q}$, le
Lemme \ref{algebra-lemma-one-equation} donne
$\dim(\overline{S}_{\mathfrak q}/\mathfrak p\overline{S}_{\mathfrak q})
= \dim(S_{\mathfrak q}/\mathfrak pS_{\mathfrak q}) - i$.
Nous en concluons que
$\dim_{\mathfrak q}(\overline{S}/R)
= \dim(\overline{S}_{\mathfrak q}/\mathfrak p\overline{S}_{\mathfrak q}) +
\text{trdeg}_{\kappa(\mathfrak p)}\ \kappa(\mathfrak q)
= d - i$ d'après le
Lemme \ref{algebra-lemma-dimension-at-a-point-finite-type-field}.
D'après le Lemme \ref{algebra-lemma-dimension-fibres-bounded-open-upstairs}, on a
$\dim_{\mathfrak q'}(\overline{S}/R) \leq d - i$
pour tout $\mathfrak q' \in V(f_1, \ldots, f_i) = \Spec(\overline{S})$
dans un voisinage de $\mathfrak q$. Ainsi, après avoir remplacé $S$ par
$S_g$ pour un certain $g \in S$, $g \not \in \mathfrak q$, nous pouvons
supposer que l'inégalité vaut pour tout $\mathfrak q'$.
Le résultat découle alors du Lemme \ref{algebra-lemma-CM-dim-finite-type}.
\end{proof}

\begin{lemma}
\label{algebra-lemma-exact-on-fibres-open}
Soit $R \to S$ un morphisme d'anneaux. Considérons un complexe homologique
fini de $S$-modules libres de type fini :
$$
F_{\bullet} :
0
\to
S^{n_e}
\xrightarrow{\varphi_e}
S^{n_{e-1}}
\xrightarrow{\varphi_{e-1}}
\ldots
\xrightarrow{\varphi_{i + 1}}
S^{n_i}
\xrightarrow{\varphi_i}
S^{n_{i-1}}
\xrightarrow{\varphi_{i-1}}
\ldots
\xrightarrow{\varphi_1}
S^{n_0}
$$
Pour tout idéal premier $\mathfrak q$ de $S$, considérons le complexe
$\overline{F}_{\bullet, \mathfrak q} =
F_{\bullet, \mathfrak q} \otimes_R \kappa(\mathfrak p)$,
où $\mathfrak p$ est l'image réciproque de $\mathfrak q$ dans $R$.
Supposons $R$ noethérien et supposons qu'il existe un entier $d$ tel que
$R \to S$ soit de type fini et plat, à fibres
$S \otimes_R \kappa(\mathfrak p)$ de Cohen-Macaulay et de dimension $d$.
L'ensemble
$$
\{\mathfrak q \in \Spec(S) \mid
\overline{F}_{\bullet, \mathfrak q}\text{ est exact}\}
$$
est ouvert dans $\Spec(S)$.
\end{lemma}

\begin{proof}
Soit $\mathfrak q$ un élément de l'ensemble défini dans le lemme.
Nous allons utiliser la Proposition \ref{algebra-proposition-what-exact} pour montrer
qu'il existe $g \in S$, $g \not \in \mathfrak q$, tel que $D(g)$ soit
contenu dans l'ensemble défini dans le lemme. Autrement dit, nous allons
montrer qu'après avoir remplacé $S$ par $S_g$, l'ensemble du lemme est
$\Spec(S)$ tout entier. Au cours de la démonstration, nous remplacerons donc
un nombre fini de fois $S$ par une telle localisation.
Rappelons que la Proposition \ref{algebra-proposition-what-exact} caractérise
l'exactitude des complexes, lorsque l'anneau est local, en fonction des rangs
des morphismes $\varphi_i$ et des idéaux $I(\varphi_i)$.
Traitons d'abord la condition sur les rangs. Posons
$r_i = n_i - n_{i + 1} + \ldots + (-1)^{e - i} n_e$.
Notons que $r_i + r_{i + 1} = n_i$ et que $r_i$ est le rang attendu de
$\varphi_i$ dans le cas exact.

\medskip\noindent
D'après le Lemme \ref{algebra-lemma-complex-exact-mod}, si
$\overline{F}_{\bullet, \mathfrak q}$ est exact, alors le complexe localisé
$F_{\bullet, \mathfrak q}$ est exact. En particulier, le complexe
$F_\bullet$ devient exact après localisation par un élément
$g \in S$, $g \not \in \mathfrak q$. Dans ce cas, la
Proposition \ref{algebra-proposition-what-exact}, appliquée à toutes les
localisations de $S$ en des idéaux premiers, implique que tous les mineurs
d'ordre $(r_i + 1) \times (r_i + 1)$ de $\varphi_i$ sont nuls.
On voit donc que le rang de $\varphi_i$ est au plus $r_i$.

\medskip\noindent
Notons $I_i \subset S$ l'idéal engendré par les mineurs d'ordre
$r_i \times r_i$ de la matrice de $\varphi_i$.
D'après la Proposition \ref{algebra-proposition-what-exact}, le complexe
$\overline{F}_{\bullet, \mathfrak q}$ est exact si et seulement si, pour
tout $1 \leq i \leq e$, on a soit
$(I_i)_{\mathfrak q} = S_{\mathfrak q}$, soit que
$(I_i)_{\mathfrak q}$ contient une suite
$S_{\mathfrak q}/\mathfrak p S_{\mathfrak q}$-régulière de longueur $i$.
En effet, par le choix de $r_i$ ci-dessus et la borne sur les rangs des
$\varphi_i$, c'est la seule façon dont les conditions de la
Proposition \ref{algebra-proposition-what-exact} peuvent être satisfaites.

\medskip\noindent
Si $(I_i)_{\mathfrak q} = S_{\mathfrak q}$, après avoir localisé $S$ en un
élément $g \not\in \mathfrak q$, nous pouvons supposer que $I_i = S$.
C'est manifestement une condition ouverte.

\medskip\noindent
Si $(I_i)_{\mathfrak q} \not = S_{\mathfrak q}$, il existe une suite
$f_1, \ldots, f_i \in (I_i)_{\mathfrak q}$ qui forme une suite régulière
dans $S_{\mathfrak q}/\mathfrak pS_{\mathfrak q}$.
Notons que, pour tout idéal premier $\mathfrak q' \subset S$ tel que
$(f_1, \ldots, f_i) \not \subset \mathfrak q'$, on a
$(I_i)_{\mathfrak q'} = S_{\mathfrak q'}$.
Le résultat découle donc du Lemme \ref{algebra-lemma-open-regular-sequence}.
\end{proof}


\begin{theorem}
\label{algebra-theorem-openness-flatness}
Soit $R$ un anneau. Soit $R \to S$ un morphisme d'anneaux de présentation
finie. Soit $M$ un $S$-module de présentation finie. L'ensemble
$$
\{ \mathfrak q \in \Spec(S) \mid
M_{\mathfrak q}\text{ est plat sur }R\}
$$
est ouvert dans $\Spec(S)$.
\end{theorem}

\begin{proof}
Soit $\mathfrak q \in \Spec(S)$ un idéal premier.
Soit $\mathfrak p \subset R$ l'image réciproque de $\mathfrak q$ dans $R$.
Notons que $M_{\mathfrak q}$ est plat sur $R$ si et seulement s'il est plat
sur $R_{\mathfrak p}$. Supposons que $M_{\mathfrak q}$ soit plat sur $R$.
Nous affirmons qu'il existe $g \in S$, $g \not \in \mathfrak q$, tel que
$M_g$ soit plat sur $R$.

\medskip\noindent
Nous nous ramenons d'abord au cas où $R$ et $S$ sont de type fini sur
$\mathbf{Z}$. Choisissons un ensemble filtrant $\Lambda$ et un système
$(R_\lambda \to S_\lambda, M_\lambda)$ comme dans le
Lemme \ref{algebra-lemma-limit-module-finite-presentation}.
Posons $\mathfrak p_\lambda$ égal à l'image réciproque de $\mathfrak p$
dans $R_\lambda$, et $\mathfrak q_\lambda$ égal à l'image réciproque de
$\mathfrak q$ dans $S_\lambda$. Le système
$$
((R_\lambda)_{\mathfrak p_\lambda},
(S_\lambda)_{\mathfrak q_\lambda},
(M_\lambda)_{\mathfrak q_{\lambda}})
$$
est alors un système comme dans le
Lemme \ref{algebra-lemma-limit-module-essentially-finite-presentation}.
Par conséquent, le Lemme \ref{algebra-lemma-colimit-eventually-flat} montre que,
pour un certain $\lambda$, le module $M_\lambda$ est plat sur $R_\lambda$
en l'idéal premier $\mathfrak q_{\lambda}$.
Supposons que nous sachions démontrer notre assertion pour le système
$(R_\lambda \to S_\lambda, M_\lambda, \mathfrak q_{\lambda})$.
Autrement dit, supposons que nous puissions trouver $g \in S_\lambda$,
$g \not\in \mathfrak q_\lambda$, tel que $(M_\lambda)_g$ soit plat sur
$R_\lambda$. D'après le Lemme \ref{algebra-lemma-limit-module-finite-presentation},
on a $M = M_\lambda \otimes_{R_\lambda} R$, et donc aussi
$M_g = (M_\lambda)_g \otimes_{R_\lambda} R$.
Le Lemme \ref{algebra-lemma-flat-base-change} donne alors l'assertion pour le
système $(R \to S, M, \mathfrak q)$.

\medskip\noindent
Nous pouvons désormais supposer que $R$ et $S$ sont de type fini sur
$\mathbf{Z}$. Nous pouvons écrire $S$ comme quotient d'un anneau de
polynômes $R[x_1, \ldots, x_n]$. Bien entendu, nous pouvons remplacer $S$
par $R[x_1, \ldots, x_n]$ et supposer que $S$ est un anneau de polynômes
sur $R$. En particulier, $R \to S$ est plat et tous les anneaux des fibres
$S \otimes_R \kappa(\mathfrak p)$ sont de dimension globale $n$.

\medskip\noindent
Choisissons une résolution $F_\bullet$ de $M$ sur $S$, chaque $F_i$ étant
libre de type fini ; voir le Lemme \ref{algebra-lemma-resolution-by-finite-free}.
Soit $K_n = \Ker(F_{n-1} \to F_{n-2})$. Notons que
$(K_n)_{\mathfrak q}$ est plat sur $R$, puisque chaque $F_i$ est plat sur
$R$ et en vertu de l'hypothèse sur $M$ ; voir le
Lemme \ref{algebra-lemma-flat-ses}. En outre, la suite
$$
0 \to
K_n/\mathfrak p K_n \to
F_{n-1}/ \mathfrak p F_{n-1} \to
\ldots \to
F_0 / \mathfrak p F_0 \to
M/\mathfrak p M \to
0
$$
est exacte après localisation en $\mathfrak q$, en raison de l'annulation
de $\text{Tor}_i^{R_\mathfrak p}(\kappa(\mathfrak p), M_{\mathfrak q})$.
Comme la dimension globale de
$S_\mathfrak q/\mathfrak p S_{\mathfrak q}$ est $n$, nous concluons que
$K_n / \mathfrak p K_n$, localisé en $\mathfrak q$, est un module libre
de type fini sur $S_\mathfrak q/\mathfrak p S_{\mathfrak q}$.
D'après le Lemme \ref{algebra-lemma-free-fibre-flat-free},
$(K_n)_{\mathfrak q}$ est libre sur $S_{\mathfrak q}$.
En particulier, il existe $g \in S$, $g \not \in \mathfrak q$, tel que
$(K_n)_g$ soit libre de type fini sur $S_g$.

\medskip\noindent
D'après le Lemme \ref{algebra-lemma-exact-on-fibres-open}, il existe une localisation
supplémentaire $S_g$ telle que le complexe
$$
0 \to K_n \to F_{n-1} \to \ldots \to F_0
$$
soit exact sur {\it toutes les fibres} de $R \to S$.
D'après le Lemme \ref{algebra-lemma-complex-exact-mod}, cela implique que le conoyau
de $F_1 \to F_0$ est plat. Cela démontre le théorème dans le cas noethérien.
\end{proof}


\section{Ouverture des lieux de Cohen-Macaulay}
\label{algebra-section-CM-open}

\noindent
Dans cette section, nous caractérisons la propriété de Cohen-Macaulay des
algèbres de type fini en termes de platitude. Nous utilisons ensuite cette
caractérisation pour montrer que l'ensemble des points où une telle algèbre
est de Cohen-Macaulay est ouvert.

\begin{lemma}
\label{algebra-lemma-where-CM}
Soit $S$ une algèbre de type fini sur un corps $k$.
Soit $\varphi : k[y_1, \ldots, y_d] \to S$ un morphisme d'anneaux
quasi-fini. Comme sous-ensembles de $\Spec(S)$, on a
$$
\{ \mathfrak q \mid
S_{\mathfrak q} \text{ plat sur }k[y_1, \ldots, y_d]\}
=
\{ \mathfrak q \mid
S_{\mathfrak q} \text{ CM et }\dim_{\mathfrak q}(S/k) = d\}
$$
Pour la notation, voir la Définition \ref{algebra-definition-relative-dimension}.
\end{lemma}

\begin{proof}
Soit $\mathfrak q \subset S$ un idéal premier. Notons
$\mathfrak p = k[y_1, \ldots, y_d] \cap \mathfrak q$.
Notons que l'on a toujours
$\dim(S_{\mathfrak q}) \leq \dim(k[y_1, \ldots, y_d]_{\mathfrak p})$,
par exemple d'après le Lemme \ref{algebra-lemma-dimension-inequality-quasi-finite}.
En outre, l'extension de corps
$\kappa(\mathfrak q)/\kappa(\mathfrak p)$ est finie et donc
$\text{trdeg}_k(\kappa(\mathfrak p)) =
\text{trdeg}_k(\kappa(\mathfrak q))$.

\medskip\noindent
Soit $\mathfrak q$ un élément du membre de gauche.
Le Lemme \ref{algebra-lemma-finite-flat-over-regular-CM} s'applique alors, et nous
concluons que $S_{\mathfrak q}$ est de Cohen-Macaulay et que
$\dim(S_{\mathfrak q}) = \dim(k[y_1, \ldots, y_d]_{\mathfrak p})$.
Combiné à l'égalité des degrés de transcendance ci-dessus et au
Lemme \ref{algebra-lemma-dimension-at-a-point-finite-type-field}, cela implique que
$\dim_{\mathfrak q}(S/k) = d$. Ainsi, $\mathfrak q$ appartient au membre
de droite.

\medskip\noindent
Soit $\mathfrak q$ un élément du membre de droite.
L'égalité des degrés de transcendance ci-dessus, l'hypothèse
$\dim_{\mathfrak q}(S/k) = d$ et le
Lemme \ref{algebra-lemma-dimension-at-a-point-finite-type-field} donnent
$\dim(S_{\mathfrak q}) = \dim(k[y_1, \ldots, y_d]_{\mathfrak p})$.
Le Lemme \ref{algebra-lemma-CM-over-regular-flat} s'applique donc, et nous voyons que
$\mathfrak q$ appartient au membre de gauche.
\end{proof}

\begin{lemma}
\label{algebra-lemma-finite-type-over-field-CM-open}
Soit $S$ une algèbre de type fini sur un corps $k$.
L'ensemble des idéaux premiers $\mathfrak q$ tels que $S_{\mathfrak q}$
soit de Cohen-Macaulay est ouvert dans $S$.
\end{lemma}

\noindent
Ce lemme est un cas particulier du
Lemme \ref{algebra-lemma-finite-presentation-flat-CM-locus-open} ci-dessous ; on
peut donc, si on le souhaite, passer directement à la démonstration de ce
dernier.

\begin{proof}
Soit $\mathfrak q \subset S$ un idéal premier tel que $S_{\mathfrak q}$
soit de Cohen-Macaulay. Il faut montrer qu'il existe
$g \in S$, $g \not \in \mathfrak q$, tel que l'anneau $S_g$ soit de
Cohen-Macaulay. Pour tout $g \in S$, $g \not \in \mathfrak q$, nous pouvons
remplacer $S$ par $S_g$ et $\mathfrak q$ par $\mathfrak qS_g$.
En combinant cette observation aux
Lemmes \ref{algebra-lemma-Noether-normalization-at-point} et
\ref{algebra-lemma-dimension-at-a-point-finite-type-field}, nous pouvons supposer
qu'il existe un morphisme d'anneaux fini injectif
$k[y_1, \ldots, y_d] \to S$ tel que
$d = \dim(S_{\mathfrak q}) + \text{trdeg}_k(\kappa(\mathfrak q))$.
Posons $\mathfrak p = k[y_1, \ldots, y_d] \cap \mathfrak q$.
Par construction, $\mathfrak q$ appartient au membre de droite de l'égalité
affichée du Lemme \ref{algebra-lemma-where-CM}. Il appartient donc aussi au membre
de gauche.

\medskip\noindent
D'après le Théorème \ref{algebra-theorem-openness-flatness}, il existe
$g \in S$, $g \not \in \mathfrak q$, tel que l'anneau $S_g$ soit plat sur
$k[y_1, \ldots, y_d]$. En appliquant de nouveau l'égalité du
Lemme \ref{algebra-lemma-where-CM}, nous concluons que tous les anneaux locaux de
$S_g$ sont de Cohen-Macaulay, comme souhaité.
\end{proof}

\begin{lemma}
\label{algebra-lemma-generic-CM}
Soit $k$ un corps. Soit $S$ une $k$-algèbre de type fini.
L'ensemble des idéaux premiers de Cohen-Macaulay forme un ouvert dense
$U \subset \Spec(S)$.
\end{lemma}

\begin{proof}
L'ensemble est ouvert d'après le
Lemme \ref{algebra-lemma-finite-type-over-field-CM-open}.
Il contient tous les idéaux premiers minimaux $\mathfrak q \subset S$,
puisque l'anneau local $S_{\mathfrak q}$ en un idéal premier minimal est de
dimension zéro, et donc de Cohen-Macaulay.
\end{proof}

\begin{lemma}
\label{algebra-lemma-finite-presentation-flat-CM-locus-open}
Soit $R$ un anneau. Soit $R \to S$ de présentation finie et plat.
Pour tout $d \geq 0$, l'ensemble
$$
\left\{
\begin{matrix}
\mathfrak q \in \Spec(S)
\text{ tel que, en posant }\mathfrak p = R \cap \mathfrak q
\text{ l'anneau de la fibre}\\
S_{\mathfrak q}/\mathfrak pS_{\mathfrak q}
\text{ est de Cohen-Macaulay}
\text{ et } \dim_{\mathfrak q}(S/R) = d
\end{matrix}
\right\}
$$
est ouvert dans $\Spec(S)$.
\end{lemma}

\begin{proof}
Soit $\mathfrak q$ un élément de l'ensemble indiqué, et soit $\mathfrak p$
l'idéal premier correspondant de $R$. Il faut trouver
$g \in S$, $g \not \in \mathfrak q$, tel que tous les anneaux des fibres
de $R \to S_g$ soient de Cohen-Macaulay.
Au cours de la démonstration, nous pourrons remplacer, un nombre fini de
fois, $S$ par $S_g$ pour un $g \in S$, $g \not \in \mathfrak q$.
Ainsi, d'après le
Lemme \ref{algebra-lemma-quasi-finite-over-polynomial-algebra}, nous pouvons supposer
qu'il existe un morphisme d'anneaux quasi-fini
$R[t_1, \ldots, t_d] \to S$ avec
$d = \dim_{\mathfrak q}(S/R)$.
Soit $\mathfrak q' = R[t_1, \ldots, t_d] \cap \mathfrak q$.
D'après le Lemme \ref{algebra-lemma-where-CM}, le morphisme d'anneaux
$$
R[t_1, \ldots, t_d]_{\mathfrak q'} /
\mathfrak p R[t_1, \ldots, t_d]_{\mathfrak q'}
\longrightarrow
S_{\mathfrak q}/\mathfrak p S_{\mathfrak q}
$$
est plat. Par conséquent, le crit\`ere de platitude par fibres,
Lemme \ref{algebra-lemma-criterion-flatness-fibre}, montre que
$R[t_1, \ldots, t_d]_{\mathfrak q'} \to S_{\mathfrak q}$ est plat.
D'après le Théorème \ref{algebra-theorem-openness-flatness}, il existe donc
$g \in S$, $g \not \in \mathfrak q$, tel que le morphisme d'anneaux
$R[t_1, \ldots, t_d] \to S_g$ soit plat. En remplaçant $S$ par $S_g$, nous
voyons que, pour tout idéal premier $\mathfrak r \subset S$, en posant
$\mathfrak r' = R[t_1, \ldots, t_d] \cap \mathfrak r$ et
$\mathfrak p' = R \cap \mathfrak r$, l'homomorphisme local
$R[t_1, \ldots, t_d]_{\mathfrak r'} \to S_{\mathfrak r}$ est plat.
Il en est donc de même du changement de base
$$
R[t_1, \ldots, t_d]_{\mathfrak r'} /
\mathfrak p' R[t_1, \ldots, t_d]_{\mathfrak r'}
\longrightarrow
S_{\mathfrak r}/\mathfrak p' S_{\mathfrak r}
$$
qui est plat. Le Lemme \ref{algebra-lemma-where-CM}, appliqué avec
$k = \kappa(\mathfrak p')$, montre donc que $\mathfrak r$ appartient à
l'ensemble du lemme, comme souhaité.
\end{proof}

\begin{lemma}
\label{algebra-lemma-generic-CM-flat-finite-presentation}
Soit $R$ un anneau. Soit $R \to S$ plat et de présentation finie.
L'ensemble des idéaux premiers $\mathfrak q$ tels que l'anneau de la fibre
$S_{\mathfrak q} \otimes_R \kappa(\mathfrak p)$, où
$\mathfrak p = R \cap \mathfrak q$, soit de Cohen-Macaulay est ouvert et
dense dans chaque fibre de $\Spec(S) \to \Spec(R)$.
\end{lemma}

\begin{proof}
Cet ensemble, noté $W$, est ouvert d'après le
Lemme \ref{algebra-lemma-finite-presentation-flat-CM-locus-open}.
Il est dense dans les fibres, car l'intersection de $W$ avec une fibre est
l'ensemble correspondant de cette fibre, auquel s'applique le
Lemme \ref{algebra-lemma-generic-CM}.
\end{proof}

\begin{lemma}
\label{algebra-lemma-extend-field-CM-locus}
Soit $k$ un corps. Soit $S$ une $k$-algèbre de type fini.
Soit $K/k$ une extension de corps, et posons $S_K = K \otimes_k S$.
Soit $\mathfrak q \subset S$ un idéal premier de $S$.
Soit $\mathfrak q_K \subset S_K$ un idéal premier de $S_K$ au-dessus de
$\mathfrak q$. Alors $S_{\mathfrak q}$ est de Cohen-Macaulay si et
seulement si $(S_K)_{\mathfrak q_K}$ est de Cohen-Macaulay.
\end{lemma}

\begin{proof}
Au cours de la démonstration, nous pouvons remplacer, un nombre fini de fois,
$S$ par $S_g$ pour tout $g \in S$, $g \not \in \mathfrak q$.
Par conséquent, en utilisant le
Lemme \ref{algebra-lemma-Noether-normalization-at-point}, nous pouvons supposer que
$\dim(S) = \dim_{\mathfrak q}(S/k) =: d$ et trouver un morphisme fini
injectif $k[x_1, \ldots, x_d] \to S$.
Notons que celui-ci induit aussi, par changement de base, un morphisme fini
injectif $K[x_1, \ldots, x_d] \to S_K$.
D'après le Lemme \ref{algebra-lemma-dimension-at-a-point-preserved-field-extension},
on a $\dim_{\mathfrak q_K}(S_K/K) = d$.
Posons $\mathfrak p = k[x_1, \ldots, x_d] \cap \mathfrak q$ et
$\mathfrak p_K = K[x_1, \ldots, x_d] \cap \mathfrak q_K$.
Considérons le diagramme commutatif suivant d'anneaux locaux noethériens :
$$
\xymatrix{
S_{\mathfrak q} \ar[r] &
(S_K)_{\mathfrak q_K} \\
k[x_1, \ldots, x_d]_{\mathfrak p} \ar[r] \ar[u] &
K[x_1, \ldots, x_d]_{\mathfrak p_K} \ar[u]
}
$$
D'après le Lemme \ref{algebra-lemma-where-CM}, il faut montrer que la flèche
verticale de gauche est plate si et seulement si celle de droite l'est.
Comme la flèche du bas est plate, cette équivalence résulte du
Lemme \ref{algebra-lemma-base-change-flat-up-down}.
\end{proof}

\begin{lemma}
\label{algebra-lemma-CM-locus-commutes-base-change}
Soit $R$ un anneau. Soit $R \to S$ de type fini.
Soit $R \to R'$ un morphisme d'anneaux. Posons $S' = R' \otimes_R S$.
Notons $f : \Spec(S') \to \Spec(S)$ l'application associée au morphisme
d'anneaux $S \to S'$. Soit $W$ l'ensemble des idéaux premiers
$\mathfrak q$ tels que l'anneau de la fibre
$S_{\mathfrak q} \otimes_R \kappa(\mathfrak p)$,
$\mathfrak p = R \cap \mathfrak q$, soit de Cohen-Macaulay, et notons
$W'$ l'ensemble analogue pour $S'/R'$. Alors $W' = f^{-1}(W)$.
\end{lemma}

\begin{proof}
Cela résulte immédiatement du Lemme \ref{algebra-lemma-extend-field-CM-locus} et
des définitions.
\end{proof}

\begin{lemma}
\label{algebra-lemma-relative-dimension-CM}
Soit $R$ un anneau. Soit $R \to S$ un morphisme d'anneaux qui est (a) plat,
(b) de présentation finie et (c) à fibres de Cohen-Macaulay.
Alors on peut écrire $S = S_0 \times \ldots \times S_n$ comme un produit de
$R$-algèbres $S_d$ de sorte que chaque $S_d$ satisfasse (a), (b) et (c) et
ait toutes ses fibres équidimensionnelles de dimension $d$.
\end{lemma}

\begin{proof}
Pour chaque entier $d$, notons $W_d \subset \Spec(S)$ l'ensemble défini
dans le Lemme \ref{algebra-lemma-finite-presentation-flat-CM-locus-open}.
On a manifestement $\Spec(S) = \coprod W_d$, et chaque $W_d$ est ouvert
d'après le lemme que nous venons de citer. Le résultat découle donc du
Lemme \ref{algebra-lemma-disjoint-implies-product}.
\end{proof}


\section{Différentielles}
\label{algebra-section-differentials}

\noindent
Dans cette section, nous définissons le module des différentielles d'un
morphisme d'anneaux.

\begin{definition}
\label{algebra-definition-derivation}
Soit $\varphi : R \to S$ un morphisme d'anneaux et soit $M$ un $S$-module.
Une {\it dérivation}, ou plus précisément une
{\it $R$-dérivation} à valeurs dans $M$, est une application $D : S \to M$
qui est additive, s'annule sur les éléments de $\varphi(R)$ et satisfait
la {\it règle de Leibniz} : $D(ab) = aD(b) + bD(a)$.
\end{definition}

\noindent
Notons que $D(ra) = rD(a)$ si $r \in R$ et $a \in S$. Une définition
équivalente est qu'une $R$-dérivation est une application $R$-linéaire
$D : S \to M$ qui satisfait la règle de Leibniz.
L'ensemble de toutes les $R$-dérivations forme un $S$-module : étant données
deux $R$-dérivations $D, D'$, leur somme $D + D' : S \to M$,
$a \mapsto D(a)+D'(a)$, est une $R$-dérivation, et, étant donnés une
$R$-dérivation $D$ et un élément $c\in S$, le multiple scalaire
$cD : S \to M$, $a \mapsto cD(a)$, est une $R$-dérivation. Nous notons ce
$S$-module
$$
\text{Der}_R(S, M).
$$
En outre, si $\alpha : M \to N$ est un morphisme de $S$-modules, alors la
composée $\alpha \circ D$ est une $R$-dérivation à valeurs dans $N$.
Ainsi, l'association $M \mapsto \text{Der}_R(S, M)$ est un foncteur
covariant.

\medskip\noindent
Considérons le morphisme suivant de $S$-modules libres :
$$
\bigoplus\nolimits_{(a, b)\in S^2} S[(a, b)] \oplus
\bigoplus\nolimits_{(f, g)\in S^2} S[(f, g)] \oplus
\bigoplus\nolimits_{r\in R} S[r]
\longrightarrow
\bigoplus\nolimits_{a\in S} S[a]
$$
défini par les règles
$$
[(a, b)] \longmapsto [a + b] - [a] - [b],\quad
[(f, g)] \longmapsto [fg] -f[g] - g[f],\quad
[r]      \longmapsto [\varphi(r)]
$$
avec les notations évidentes. Soit $\Omega_{S/R}$ le conoyau de ce morphisme.
Il existe une application $\text{d} : S \to \Omega_{S/R}$ qui envoie $a$
sur la classe $\text{d}a$ de $[a]$ dans le conoyau. C'est une $R$-dérivation
en vertu des relations imposées à $\Omega_{S/R}$ ; autrement dit,
$$
\text{d}(a + b) = \text{d}a + \text{d}b, \quad
\text{d}(fg) = f\text{d}g + g\text{d}f, \quad
\text{d}\varphi(r) = 0
$$
où $a,b,f,g \in S$ et $r \in R$.

\begin{definition}
\label{algebra-definition-differentials}
Le couple $(\Omega_{S/R}, \text{d})$ est appelé le {\it module des
différentielles de K\"ahler}, ou le {\it module des différentielles}
de $S$ sur $R$.
\end{definition}

\begin{lemma}
\label{algebra-lemma-universal-omega}
\begin{slogan}
Les morphismes partant du module des différentielles sont les mêmes que
les dérivations.
\end{slogan}
Le module des différentielles de $S$ sur $R$ possède la propriété
universelle suivante. L'application
$$
\Hom_S(\Omega_{S/R}, M)
\longrightarrow
\text{Der}_R(S, M), \quad
\alpha
\longmapsto
\alpha \circ \text{d}
$$
est un isomorphisme de foncteurs.
\end{lemma}

\begin{proof}
Par définition, une $R$-dérivation est une règle qui associe à chaque
$a \in S$ un élément $D(a) \in M$. Ainsi, $D$ donne naissance à un
morphisme $[D] : \bigoplus S[a] \to M$. Or les conditions requises pour
être une $R$-dérivation signifient exactement que $[D]$ annule l'image du
morphisme dans la présentation de $\Omega_{S/R}$ affichée ci-dessus.
\end{proof}

\begin{lemma}
\label{algebra-lemma-trivial-differential-surjective}
Supposons que $R \to S$ soit surjectif.
Alors $\Omega_{S/R} = 0$.
\end{lemma}

\begin{proof}
On peut le voir soit parce que toutes les $R$-dérivations doivent
manifestement être nulles, soit parce que le morphisme de la présentation
de $\Omega_{S/R}$ est surjectif.
\end{proof}

\noindent
Supposons que
\begin{equation}
\label{algebra-equation-functorial-omega}
\vcenter{
\xymatrix{
S \ar[r]_\varphi
&
S'
\\
R \ar[r]^\psi \ar[u]^\alpha
&
R' \ar[u]_\beta
}
}
\end{equation}
soit un diagramme commutatif d'anneaux. Dans ce cas, il existe un morphisme
naturel entre les modules des différentielles qui s'insère dans le diagramme
commutatif
$$
\xymatrix{
\Omega_{S/R} \ar[r] &
\Omega_{S'/R'}
\\
S \ar[u]^{\text{d}} \ar[r]^{\varphi}
&
S' \ar[u]_{\text{d}}
}
$$
Pour construire ce morphisme, il suffit d'utiliser le morphisme évident
entre les présentations de $\Omega_{S/R}$ et $\Omega_{S'/R'}$, à savoir
\begin{equation}
\label{algebra-equation-map-presentations}
\vcenter{
\xymatrix{
\bigoplus S'[(a', b')]
\oplus
\bigoplus S'[(f', g')]
\oplus
\bigoplus S'[r'] \ar[r]
&
\bigoplus S' [a'] \\
 \\
\bigoplus S[(a, b)]
\oplus
\bigoplus S[(f, g)]
\oplus
\bigoplus S[r] \ar[r]
\ar[uu]^{
\begin{matrix}
[(a, b)] \mapsto [(\varphi(a), \varphi(b))] \\
[(f, g)] \mapsto [(\varphi(f), \varphi(g))] \\
[r]\mapsto [\psi(r)]
\end{matrix}
} &
\bigoplus S[a] \ar[uu]_{[a] \mapsto [\varphi(a)]}
}
}
\end{equation}
Le résultat est simplement que $f\text{d}g \in \Omega_{S/R}$ est envoyé sur
$\varphi(f)\text{d}\varphi(g)$.


\begin{lemma}
\label{algebra-lemma-colimit-differentials}
Soit $I$ un ensemble filtrant.
Soit $(R_i \to S_i, \varphi_{ii'})$ un système de morphismes d'anneaux
indexé par $I$ ; voir
Catégories, section \ref{categories-section-posets-limits}.
Alors on a
$$
\Omega_{S/R} =
\colim_i \Omega_{S_i/R_i}.
$$
où $R \to S = \colim (R_i \to S_i)$.
\end{lemma}

\begin{proof}
Cela résulte clairement de la présentation définissant $\Omega_{S/R}$ et
de la fonctorialité décrite ci-dessus.
\end{proof}

\begin{lemma}
\label{algebra-lemma-differential-surjective}
Dans le diagramme (\ref{algebra-equation-functorial-omega}), supposons que
$S \to S'$ soit surjectif de noyau $I \subset S$.
Alors $\Omega_{S/R} \to \Omega_{S'/R'}$ est surjectif et son noyau est
engendré comme $S$-module par les éléments $\text{d}a$, où $a \in S$ est
tel que $\varphi(a) \in \beta(R')$.
(Cela comprend en particulier les éléments $\text{d}(i)$, $i \in I$.)
\end{lemma}

\begin{proof}[Première démonstration]
Considérons le morphisme de présentations
(\ref{algebra-equation-map-presentations}). Le morphisme vertical de droite entre
modules libres est manifestement surjectif. Le morphisme est donc surjectif.
Supposons qu'un élément $\eta$ de $\Omega_{S/R}$ soit envoyé sur zéro dans
$\Omega_{S'/R'}$. Écrivons $\eta$ comme l'image de $\sum s_i[a_i]$ pour des
$s_i, a_i \in S$. Nous voyons alors que
$\sum \varphi(s_i)[\varphi(a_i)]$ est l'image d'un élément
$$
\theta = \sum s_j'[a_j', b_j'] + \sum s_k'[f_k', g_k'] + \sum s_l'[r_l']
$$
du coin supérieur gauche du diagramme.
Puisque $\varphi$ est surjectif, les termes $s_j'[a_j', b_j']$ et
$s_k'[f_k', g_k']$ sont dans l'image d'éléments du coin inférieur droit.
Ainsi, en modifiant $\eta$ et $\theta$ par soustraction des images de ces
éléments, nous pouvons supposer que $\theta = \sum s_l'[r_l']$.
Autrement dit, nous voyons que $\sum \varphi(s_i)[\varphi(a_i)]$ est de la
forme $\sum s'_l [\beta(r'_l)]$.
Ensuite, nous pouvons supposer qu'il existe $a' \in S'$ tel que
$a' = \varphi(a_i)$ pour tout $i$ et $a' = \beta(r_l')$ pour tout $l$.
Cela résulte clairement de la décomposition en somme directe du coin
supérieur droit du diagramme. Choisissons $a \in S$ tel que
$\varphi(a) = a'$. Nous pouvons alors écrire $a_i = a + x_i$ pour des
$x_i \in I$. En utilisant les relations permises, nous pouvons donc
supposer que tous les $a_i$ sont égaux à $a$. Mais nous pouvons alors
supposer que notre élément est de la forme $s[a]$. Nous savons encore que
$\varphi(s)[a'] = \sum \varphi(s_l')[\beta(r_l')]$.
Ainsi, soit $\varphi(s) = 0$ et nous avons terminé, soit
$a' = \varphi(a)$ appartient à l'image de $\beta$ et nous avons également
terminé.
\end{proof}

\begin{proof}[Seconde démonstration]
Nous utiliserons sans autre mention la propriété universelle des modules
des différentielles donnée dans le Lemme \ref{algebra-lemma-universal-omega}.

\medskip\noindent
Dans (\ref{algebra-equation-functorial-omega}), posons $R'' = S \times_{S'} R'$.
Nous obtenons alors le diagramme suivant :
$$
\xymatrix{
S \ar[r] & S \ar[r] & S' \\
R \ar[r] \ar[u] & R'' \ar[r] \ar[u] & R' \ar[u]
}
$$
Soit $M$ un $S$-module. Il résulte immédiatement des définitions qu'une
$R$-dérivation $D : S \to M$ est une $R''$-dérivation si et seulement si
elle annule les éléments de l'image de $R'' \to S$. La propriété universelle
traduit cela en l'énoncé suivant : le morphisme naturel
$\Omega_{S/R} \to \Omega_{S/R''}$ est surjectif et son noyau est engendré
comme $S$-module par l'image de $R''$.

\medskip\noindent
D'après le paragraphe précédent, il suffit de montrer que
$\Omega_{S/R} \to \Omega_{S'/R'}$ est un isomorphisme lorsque $S \to S'$
est surjectif et $R = S \times_{S'} R'$. Soit $M'$ un $S'$-module.
Observons que toute $R'$-dérivation $D' : S' \to M'$ donne une
$R$-dérivation par précomposition avec $S \to S'$. Réciproquement,
supposons que $M$ soit un $S$-module et que $D : S \to M$ soit une
$R$-dérivation. Si $i \in I$, alors il existe $a \in R$ tel que
$\alpha(a) = i$ (puisque $R = S \times_{S'} R'$). Il s'ensuit que
$D(i) = 0$, et donc que $0 = D(is) = iD(s)$ pour tout $s \in S$. Ainsi,
l'image de $D$ est contenue dans le sous-module $M' \subset M$ des éléments
annulés par $I$ et, de plus, l'application induite $S \to M'$ se factorise
par une $R'$-dérivation $S' \to M'$. C'est un exercice d'utiliser la
propriété universelle pour voir que cela signifie que
$\Omega_{S/R} \to \Omega_{S'/R'}$ est un isomorphisme ; nous omettons les
détails.
\end{proof}

\begin{lemma}
\label{algebra-lemma-exact-sequence-differentials}
Soient $A \to B \to C$ des morphismes d'anneaux.
Il existe alors une suite exacte canonique
$$
C \otimes_B \Omega_{B/A} \to
\Omega_{C/A} \to
\Omega_{C/B} \to 0
$$
de $C$-modules.
\end{lemma}

\begin{proof}
Nous obtenons un diagramme (\ref{algebra-equation-functorial-omega}) en posant
$R = A$, $S = C$, $R' = B$ et $S' = C$.
D'après le Lemme \ref{algebra-lemma-differential-surjective}, le morphisme
$\Omega_{C/A} \to \Omega_{C/B}$ est surjectif, et son noyau est engendré
par les éléments $\text{d}(c)$, où $c \in C$ appartient à l'image de
$B \to C$. Le lemme en résulte.
\end{proof}

\begin{lemma}
\label{algebra-lemma-differentials-localize}
Soit $\varphi : A \to B$ un morphisme d'anneaux.
\begin{enumerate}
\item Si $S \subset A$ est une partie multiplicative dont l'image est
formée d'éléments inversibles de $B$, alors
$\Omega_{B/A} = \Omega_{B/S^{-1}A}$.
\item Si $S \subset B$ est une partie multiplicative, alors
$S^{-1}\Omega_{B/A} = \Omega_{S^{-1}B/A}$.
\end{enumerate}
\end{lemma}

\begin{proof}
Pour démontrer l'égalité de (1), il suffit de montrer que toute
$A$-dérivation $D : B \to M$ annule les éléments $\varphi(s)^{-1}$.
Cela résulte clairement de la règle de Leibniz appliquée à
$1 = \varphi(s) \varphi(s)^{-1}$.
Pour montrer (2), notons qu'il existe un morphisme évident
$S^{-1}\Omega_{B/A} \to \Omega_{S^{-1}B/A}$.
Pour montrer qu'il s'agit d'un isomorphisme, il suffit de montrer qu'il
existe une $A$-dérivation $\text{d}'$ de $S^{-1}B$ à valeurs dans
$S^{-1}\Omega_{B/A}$. Pour la définir, posons simplement
$\text{d}'(b/s) = (1/s)\text{d}b - (1/s^2)b\text{d}s$.
Nous omettons les détails.
\end{proof}

\begin{lemma}
\label{algebra-lemma-differential-seq}
Dans le diagramme (\ref{algebra-equation-functorial-omega}), supposons que
$S \to S'$ soit surjectif de noyau $I \subset S$, et supposons que
$R' = R$.
Il existe alors une suite exacte canonique de $S'$-modules
$$
I/I^2
\longrightarrow
\Omega_{S/R} \otimes_S S'
\longrightarrow
\Omega_{S'/R}
\longrightarrow
0
$$
Le morphisme de gauche est caractérisé par la règle qui envoie
$f \in I$ sur $\text{d}f \otimes 1$.
\end{lemma}

\begin{proof}
Le terme du milieu est $\Omega_{S/R} \otimes_S S/I$.
Pour $f \in I$, notons $\overline{f}$ l'image de $f$ dans $I/I^2$.
Pour montrer que l'application
$\overline{f} \mapsto \text{d}f \otimes 1$ est bien définie, il suffit de
vérifier que $\text{d} f_1f_2 \otimes 1 = 0$ si $f_1, f_2 \in I$.
Cela résulte clairement de la règle de Leibniz
$\text{d} f_1f_2 \otimes 1 = (f_1 \text{d}f_2 + f_2 \text{d} f_1 )\otimes 1 =
\text{d}f_2 \otimes f_1 + \text{d}f_1 \otimes f_2 = 0$.
Un calcul analogue montre que cette application est
$S' = S/I$-linéaire.

\medskip\noindent
Le morphisme $\Omega_{S/R} \otimes_S S' \to \Omega_{S'/R}$ est le morphisme
$S'$-linéaire canonique associé au morphisme $S$-linéaire
$\Omega_{S/R} \to \Omega_{S'/R}$. Il est surjectif parce que
$\Omega_{S/R} \to \Omega_{S'/R}$ est surjectif d'après le
Lemme \ref{algebra-lemma-differential-surjective}.

\medskip\noindent
La composée des deux morphismes est nulle parce que $\text{d}f$ est envoyé
sur zéro dans $\Omega_{S'/R}$ pour $f \in I$. Notons que l'exactitude dit
simplement que le noyau de $\Omega_{S/R} \to \Omega_{S'/R}$ est engendré
comme sous-$S$-module par le sous-module $I\Omega_{S/R}$ et par les éléments
$\text{d}f$, avec $f \in I$. Nous savons par le
Lemme \ref{algebra-lemma-differential-surjective} que ce noyau est engendré par les
éléments $\text{d}(a)$ où $\varphi(a) = \beta(r)$ pour un certain $r \in R$.
Mais alors $a = \alpha(r) + a - \alpha(r)$, de sorte que
$\text{d}(a) = \text{d}(a - \alpha(r))$. De plus,
$a - \alpha(r) \in I$ puisque $\varphi(a - \alpha(r)) =
\varphi(a) - \varphi(\alpha(r)) = \beta(r) - \beta(r) = 0$.
Nous concluons que les éléments $\text{d}f$ avec $f \in I$ engendrent déjà
le noyau comme $S$-module, comme souhaité.
\end{proof}

\begin{lemma}
\label{algebra-lemma-differential-seq-split}
Dans le diagramme (\ref{algebra-equation-functorial-omega}), supposons que
$S \to S'$ soit surjectif de noyau $I \subset S$ et supposons que
$R' = R$. Supposons en outre qu'il existe un morphisme de $R$-algèbres
$S' \to S$ qui soit un inverse à droite de $S \to S'$.
Alors la suite exacte de $S'$-modules du
Lemme \ref{algebra-lemma-differential-seq} devient une suite exacte courte
$$
0 \longrightarrow
I/I^2
\longrightarrow
\Omega_{S/R} \otimes_S S'
\longrightarrow
\Omega_{S'/R}
\longrightarrow
0
$$
qui est même une suite exacte courte scindée.
\end{lemma}

\begin{proof}
Soit $\beta : S' \to S$ l'inverse à droite de la surjection
$\alpha : S \to S'$. Considérons l'application
$$
D : S \longrightarrow I/I^2,
\quad
x \longmapsto x - \beta(\alpha(x))
$$
Il est facile de montrer que $D$ est une $R$-dérivation (nous l'omettons).
En outre, $x D(s) = 0$ si $x \in I, s \in S$. Par conséquent, la propriété
universelle montre que $D$ induit un morphisme
$\tau : \Omega_{S/R} \otimes_S S' \to I/I^2$.
Nous omettons la vérification que c'est un inverse à gauche de
$\text{d} : I/I^2  \to \Omega_{S/R} \otimes_S S'$. D'où le résultat.
\end{proof}

\begin{lemma}
\label{algebra-lemma-differential-mod-power-ideal}
Soit $R \to S$ un morphisme d'anneaux. Soit $I \subset S$ un idéal.
Soit $n \geq 1$ un entier. Posons $S' = S/I^{n + 1}$.
Le morphisme $\Omega_{S/R} \to \Omega_{S'/R}$ induit un isomorphisme
$$
\Omega_{S/R} \otimes_S S/I^n
\longrightarrow
\Omega_{S'/R} \otimes_{S'} S/I^n.
$$
\end{lemma}

\begin{proof}
Cela résulte du Lemme \ref{algebra-lemma-differential-seq} et du fait que
$\text{d}(I^{n + 1}) \subset I^n\Omega_{S/R}$ par la règle de Leibniz pour
$\text{d}$.
\end{proof}

\begin{lemma}
\label{algebra-lemma-differentials-base-change}
Supposons donnés des morphismes d'anneaux $R \to R'$ et $R \to S$.
Posons $S' = S \otimes_R R'$, de sorte que nous obtenions un diagramme
(\ref{algebra-equation-functorial-omega}). Alors le morphisme canonique défini
ci-dessus induit un isomorphisme
$\Omega_{S/R} \otimes_R R' = \Omega_{S'/R'}$.
\end{lemma}

\begin{proof}
Notons $\text{d}' : S' = S \otimes_R R' \to \Omega_{S/R} \otimes_R R'$
l'application
$\text{d}'( \sum a_i \otimes x_i ) = \sum \text{d}(a_i) \otimes x_i$.
Elle existe parce que l'application
$S \times R' \to \Omega_{S/R} \otimes_R R'$,
$(a, x)\mapsto \text{d}a \otimes_R x$, est $R$-bilinéaire.
C'est une $R'$-dérivation, comme on le vérifie par un calcul simple.
Nous allons montrer que $(\Omega_{S/R} \otimes_R R', \text{d}')$ satisfait
la propriété universelle. Soit $D : S' \to M'$ une $R'$-dérivation à valeurs
dans un $S'$-module. La composée $S \to S' \to M'$ est une $R$-dérivation ;
nous obtenons donc un morphisme $S$-linéaire
$\varphi_D : \Omega_{S/R} \to M'$. Nous pouvons le tensoriser par $R'$ et
obtenir le morphisme $\varphi'_D :
\Omega_{S/R} \otimes_R R' \to M'$, $\varphi'_D(\eta \otimes x) =
x\varphi_D(\eta)$. Il est clair que $D = \varphi'_D \circ \text{d}'$.
\end{proof}

\noindent
Le morphisme de multiplication $S \otimes_R S \to S$ est le morphisme de
$R$-algèbres qui envoie $a \otimes b$ sur $ab$ dans $S$. C'est aussi un
morphisme de $S$-algèbres si l'on considère $S \otimes_R S$ comme une
$S$-algèbre par l'un ou l'autre des morphismes $S \to S \otimes_R S$.

\begin{lemma}
\label{algebra-lemma-differentials-diagonal}
Soit $R \to S$ un morphisme d'anneaux. Soit
$J = \Ker(S \otimes_R S \to S)$ le noyau du morphisme de multiplication.
Il existe un isomorphisme canonique de $S$-modules
$\Omega_{S/R} \to J/J^2$,
$a \text{d} b \mapsto a \otimes b - ab \otimes 1$.
\end{lemma}

\begin{proof}[Première démonstration]
Appliquons le Lemme \ref{algebra-lemma-differential-seq-split} au diagramme commutatif
$$
\xymatrix{
S \otimes_R S \ar[r] & S \\
S \ar[r] \ar[u] & S \ar[u]
}
$$
où la flèche verticale de gauche est $a \mapsto a \otimes 1$. Nous obtenons
la suite exacte
$0 \to J/J^2 \to
\Omega_{S \otimes_R S/S} \otimes_{S \otimes_R S} S \to \Omega_{S/S} \to 0$.
D'après le Lemme \ref{algebra-lemma-trivial-differential-surjective}, le terme
$\Omega_{S/S}$ est $0$, et nous obtenons un isomorphisme entre les deux
autres termes. D'après le Lemme \ref{algebra-lemma-differentials-base-change}, on a
$\Omega_{S \otimes_R S/S} = \Omega_{S/R} \otimes_S (S \otimes_R S)$,
car $S \to S \otimes_R S$ est le changement de base de $R \to S$, et donc
$$
\Omega_{S \otimes_R S/S} \otimes_{S \otimes_R S} S =
\Omega_{S/R} \otimes_S (S \otimes_R S) \otimes_{S \otimes_R S} S =
\Omega_{S/R}
$$
Nous omettons la vérification que le morphisme est donné par la règle du
lemme.
\end{proof}

\begin{proof}[Seconde démonstration]
Montrons d'abord que la règle
$a \text{d} b \mapsto a \otimes b - ab \otimes 1$ est bien définie.
Pour cela, nous devons montrer que $\text{d}r$ et
$a\text{d}b + b \text{d}a - d(ab)$ sont envoyés sur zéro.
Le premier l'est parce que $r \otimes 1 - 1 \otimes r = 0$ par définition
du produit tensoriel. Le second l'est parce que
$$
(a \otimes b - ab \otimes 1) +
(b \otimes a - ba \otimes 1) -
(1 \otimes ab - ab \otimes 1)
=
(a \otimes 1 - 1\otimes a)(1\otimes b - b \otimes 1)
$$
appartient à $J^2$.

\medskip\noindent
Construisons un morphisme dans l'autre sens.
Nous pouvons considérer $S \to S \otimes_R S$, $a \mapsto a \otimes 1$,
comme le changement de base de $R \to S$. Ainsi, d'après le
Lemme \ref{algebra-lemma-differentials-base-change}, on a
$\Omega_{S \otimes_R S/S} = \Omega_{S/R} \otimes_S (S \otimes_R S)$.
À ce stade, la suite du Lemme \ref{algebra-lemma-differential-seq} donne un morphisme
$$
J/J^2  \to \Omega_{S \otimes_R S/ S} \otimes_{S \otimes_R S} S
= (\Omega_{S/R} \otimes_S (S \otimes_R S))\otimes_{S \otimes_R S} S
= \Omega_{S/R}.
$$
Nous laissons au lecteur le soin de vérifier qu'il est l'inverse du
morphisme ci-dessus.
\end{proof}





\begin{lemma}
\label{algebra-lemma-differentials-polynomial-ring}
Si $S = R[x_1, \ldots, x_n]$, alors $\Omega_{S/R}$ est un $S$-module libre
fini de base $\text{d}x_1, \ldots, \text{d}x_n$.
\end{lemma}

\begin{proof}
Montrons d'abord que $\text{d}x_1, \ldots, \text{d}x_n$ engendrent
$\Omega_{S/R}$ comme $S$-module. Pour cela, montrons que $\text{d}g$ peut
s'écrire comme une somme $\sum g_i \text{d}x_i$ pour tout
$g \in R[x_1, \ldots, x_n]$. Nous raisonnons par récurrence sur le degré
(total) de $g$. C'est clair si le degré de $g$ est $0$, car alors
$\text{d}g = 0$. Si le degré de $g$ est $> 0$, nous pouvons écrire $g$
sous la forme $c + \sum g_i x_i$ avec $c\in R$ et
$\deg(g_i) < \deg(g)$.
La règle de Leibniz donne
$\text{d}g = \sum g_i \text{d} x_i + \sum x_i \text{d}g_i$,
et la récurrence permet de conclure.

\medskip\noindent
Considérons la $R$-dérivation $\partial / \partial x_i :
R[x_1, \ldots, x_n] \to R[x_1, \ldots, x_n]$. (Nous laissons au lecteur le
soin de la définir ; sa propriété caractéristique est que
$\partial / \partial x_i (x_j) = \delta_{ij}$.)
Par la propriété universelle, elle correspond à un morphisme de $S$-modules
$l_i : \Omega_{S/R} \to R[x_1, \ldots, x_n]$ qui envoie $\text{d}x_i$
sur $1$ et $\text{d}x_j$ sur $0$ pour $j \not = i$.
Il est donc clair qu'il n'existe aucune relation $S$-linéaire entre les
éléments $\text{d}x_1, \ldots, \text{d}x_n$.
\end{proof}

\begin{lemma}
\label{algebra-lemma-differentials-finitely-presented}
Supposons que $R \to S$ soit de présentation finie.
Alors $\Omega_{S/R}$ est un $S$-module de présentation finie.
\end{lemma}

\begin{proof}
Écrivons $S = R[x_1, \ldots, x_n]/(f_1, \ldots, f_m)$.
Écrivons $I = (f_1, \ldots, f_m)$. D'après le
Lemme \ref{algebra-lemma-differential-seq}, il existe une suite exacte de $S$-modules
$$
I/I^2
\to
\Omega_{R[x_1, \ldots, x_n]/R} \otimes_{R[x_1, \ldots, x_n]} S
\to
\Omega_{S/R}
\to
0
$$
Le résultat découle du fait que $I/I^2$ est un $S$-module de type fini
(engendré par les images des $f_i$) et que le terme du milieu est libre de
type fini d'après le
Lemme \ref{algebra-lemma-differentials-polynomial-ring}.
\end{proof}

\begin{lemma}
\label{algebra-lemma-differentials-finitely-generated}
Supposons que $R \to S$ soit de type fini.
Alors $\Omega_{S/R}$ est un $S$-module de type fini.
\end{lemma}

\begin{proof}
La démonstration est très semblable à celle du
Lemme \ref{algebra-lemma-differentials-finitely-presented}, mais plus facile.
\end{proof}


\section{Le complexe de de Rham}
\label{algebra-section-de-rham-complex}

\noindent
Soit $A \to B$ un morphisme d'anneaux. Notons
$\text{d} : B \to \Omega_{B/A}$ le module des différentielles muni de sa
$A$-dérivation universelle construite dans la
section \ref{algebra-section-differentials}.
Soit $\Omega_{B/A}^i = \wedge^i_B(\Omega_{B/A})$, pour $i \geq 0$, la
$i$-ième puissance extérieure, comme dans la
section \ref{algebra-section-tensor-algebra}.
Le {\it complexe de de Rham de $B$ sur $A$} est le complexe
$$
\Omega_{B/A}^0 \to \Omega_{B/A}^1 \to \Omega_{B/A}^2 \to \ldots
$$
muni des différentielles $A$-linéaires construites et décrites ci-dessous.

\medskip\noindent
L'application $\text{d} : \Omega^0_{B/A} \to \Omega^1_{B/A}$ est la
dérivation universelle $\text{d} : B \to \Omega_{B/A}$.
Observons qu'elle est bien $A$-linéaire.

\medskip\noindent
Pour $p \geq 1$, nous affirmons qu'il existe une unique application
$A$-linéaire $\text{d} : \Omega_{B/A}^p \to \Omega_{B/A}^{p + 1}$ telle que
\begin{equation}
\label{algebra-equation-rule}
\text{d}\left(b_0\text{d}b_1 \wedge \ldots \wedge \text{d}b_p\right) =
\text{d}b_0 \wedge \text{d}b_1 \wedge \ldots \wedge \text{d}b_p
\end{equation}
Rappelons que $\Omega_{B/A}$ est engendré comme $B$-module par les éléments
$\text{d}b$. Ainsi, $\Omega^p_{B/A}$ est engendré comme $A$-module par les
éléments $b_0\text{d}b_1 \wedge \ldots \wedge \text{d}b_p$, et il s'ensuit
que l'application $\text{d} : \Omega^p_{B/A} \to \Omega^{p + 1}_{B/A}$,
si elle existe, est unique.

\medskip\noindent
Construction de $\text{d} : \Omega_{B/A}^1 \to \Omega_{B/A}^2$.
D'après la Définition \ref{algebra-definition-differentials}, les éléments
$\text{d}b$ engendrent librement $\Omega_{B/A}$ comme $B$-module, modulo
les relations $\text{d}a = 0$ pour $a \in A$ et
$\text{d}(b' + b'') = \text{d}b' + \text{d}b''$ et
$\text{d}(b'b'') = b'\text{d}b'' + b''\text{d}b'$ pour $b', b'' \in B$.
Par conséquent, pour montrer que la règle
$$
\sum b'_i \text{d}b_i \longmapsto \sum \text{d}b'_i \wedge \text{d}b_i
$$
est bien définie, nous devons montrer que les éléments
$$
b\text{d}a,
\quad\text{et}\quad
b\text{d}(b' + b'') - b\text{d}b' - b\text{d}b''
\quad\text{et}\quad
b\text{d}(b'b'') - bb'\text{d}b'' - bb''\text{d}b'
$$
pour $a \in A$ et $b, b', b'' \in B$ sont envoyés sur zéro.
Cela résulte clairement d'un calcul direct utilisant la règle de Leibniz
pour $\text{d}$.

\medskip\noindent
Observons que la composée
$\Omega^0_{B/A} \to \Omega^1_{B/A} \to \Omega^2_{B/A}$ est nulle puisque
$\text{d}(\text{d}(b)) = \text{d}(1 \text{d}b) =
\text{d}(1) \wedge \text{d}(b) = 0 \wedge \text{d}b = 0$.
Ici, $\text{d}(1) = 0$ car $1 \in B$ appartient à l'image de $A \to B$.
Nous utiliserons ce fait ci-dessous.

\medskip\noindent
Construction de $\text{d} : \Omega_{B/A}^p \to \Omega_{B/A}^{p + 1}$
pour $p \geq 2$. Nous allons montrer que l'application $A$-linéaire
$$
\gamma : \Omega^1_{B/A} \otimes_A \ldots \otimes_A \Omega^1_{B/A}
\longrightarrow
\Omega_{B/A}^{p + 1}
$$
définie par la formule
$$
\omega_1 \otimes \ldots \otimes \omega_p
\longmapsto
\sum (-1)^{i + 1}
\omega_1 \wedge \ldots \wedge \text{d}(\omega_i) \wedge \ldots \wedge \omega_p
$$
se factorise par la surjection naturelle
$\Omega^1_{B/A} \otimes_A \ldots \otimes_A \Omega^1_{B/A} \to \Omega^p_{B/A}$
pour donner l'application souhaitée
$\text{d} : \Omega^p_{B/A} \to \Omega^{p + 1}_{B/A}$.
D'après le Lemme \ref{algebra-lemma-present-wedge}, le noyau de
$\Omega^1_{B/A} \otimes_A \ldots \otimes_A \Omega^1_{B/A} \to \Omega^p_{B/A}$
est engendré comme $A$-module par les éléments
$\omega_1 \otimes \ldots \otimes \omega_p$ tels que
$\omega_i = \omega_j$ pour certains $i \not = j$, et par les éléments
$\omega_1 \otimes \ldots \otimes f\omega_i \otimes \ldots \otimes \omega_p -
\omega_1 \otimes \ldots \otimes f\omega_j \otimes \ldots \otimes \omega_p$
pour un certain $f \in B$.
Un calcul direct montre que le premier type d'éléments est envoyé sur $0$
par $\gamma$ ; autrement dit, $\gamma$ est alternée.
Pour terminer, nous devons montrer que
$$
\gamma(
\omega_1 \otimes \ldots \otimes f\omega_i \otimes \ldots \otimes \omega_p)
=
\gamma(
\omega_1 \otimes \ldots \otimes f\omega_j \otimes \ldots \otimes \omega_p)
$$
pour $f \in B$. Par $A$-linéarité et par la propriété d'alternance, il suffit
de le montrer pour $p = 2$, $i = 1$, $j = 2$,
$\omega_1 = b \text{d}b'$ et $\omega_2 = c \text{d} c'$ avec
$b, b', c, c' \in B$.
Nous devons donc montrer que
\begin{align*}
& \text{d}(fb) \wedge \text{d}b' \wedge c \text{d}c'
- fb \text{d}b' \wedge \text{d}c \wedge \text{d}c' \\
& =
\text{d}b \wedge \text{d}b' \wedge fc\text{d}c'
- b \text{d}b' \wedge \text{d}(fc) \wedge \text{d}c'
\end{align*}
c'est-à-dire que
$$
(c \text{d}(fb) + fb \text{d}c - fc \text{d}b - b \text{d}(fc))
\wedge \text{d}b' \wedge \text{d}c' = 0.
$$
Cela découle de la règle de Leibniz. Observons que la valeur de $\gamma$
sur l'élément
$b_0\text{d}b_1 \otimes \text{d}b_2 \otimes \ldots \otimes \text{d}b_p$
est $\text{d}b_0 \wedge \text{d}b_1 \wedge \ldots \wedge \text{d}b_p$,
et donc que (\ref{algebra-equation-rule}) est satisfaite pour l'application
$\text{d} : \Omega^p_{B/A} \to \Omega^{p + 1}_{B/A}$ ainsi obtenue.

\medskip\noindent
Enfin, puisque $\Omega^p_{B/A}$ est engendré additivement par les éléments
$b_0\text{d}b_1 \wedge \ldots \wedge \text{d}b_p$ et puisque
$\text{d}(b_0\text{d}b_1 \wedge \ldots \wedge \text{d}b_p) =
\text{d}b_0 \wedge \ldots \wedge \text{d}b_p$, nous voyons exactement de
la même manière que la composée
$
\Omega^p_{B/A} \to
\Omega^{p + 1}_{B/A} \to
\Omega^{p + 2}_{B/A}
$
est nulle pour $p \geq 1$. Ainsi, le complexe de de Rham est bien un
complexe.

\medskip\noindent
Étant donné un anneau $R$, nous posons $\Omega_R = \Omega_{R/\mathbf{Z}}$.
On l'appelle parfois le module absolu des différentielles de $R$ ; cette
terminologie a un sens : si $\Omega_R$ est le module des différentielles
pour lequel nous imposons seulement la règle de Leibniz, sans imposer
l'annulation de $\text{d}1$, alors la règle de Leibniz donne
$\text{d}1 = \text{d}(1 \cdot 1) =
1 \text{d}1 + 1 \text{d}1 = 2 \text{d}1$, et donc
$\text{d}1 = 0$ dans $\Omega_R$. Dans ce cas, le
{\it complexe de de Rham absolu de $R$} est le complexe correspondant
$$
\Omega_R^0 \to \Omega_R^1 \to \Omega_R^2 \to \ldots
$$
où nous posons $\Omega^i_R = \Omega^i_{R/\mathbf{Z}}$, et ainsi de suite.

\medskip\noindent
Supposons donné un diagramme commutatif d'anneaux
$$
\xymatrix{
B \ar[r] & B' \\
A \ar[r] \ar[u] & A' \ar[u]
}
$$
Il existe un morphisme naturel de complexes de de Rham
$$
\Omega^\bullet_{B/A} \longrightarrow \Omega^\bullet_{B'/A'}
$$
En degré $0$, il s'agit du morphisme $B \to B'$ ; en degré $1$, il s'agit
du morphisme $\Omega_{B/A} \to \Omega_{B'/A'}$ construit dans la
section \ref{algebra-section-differentials} ; et, pour $p \geq 2$, il s'agit du
morphisme induit
$\Omega^p_{B/A} = \wedge^p_B(\Omega_{B/A}) \to
\wedge^p_{B'}(\Omega_{B'/A'}) = \Omega^p_{B'/A'}$.
La compatibilité aux différentielles découle de la caractérisation des
différentielles par la formule (\ref{algebra-equation-rule}).

\begin{lemma}
\label{algebra-lemma-de-rham-complex-base-change}
Supposons donnés des morphismes d'anneaux $A \to A'$ et $A \to B$.
Posons $B' = B \otimes_A A'$, de sorte que nous obtenions un diagramme
comme ci-dessus. Alors le morphisme canonique défini ci-dessus induit un
isomorphisme
$\Omega^\bullet_{B/A} \otimes_A A' = \Omega^\bullet_{B'/A'}$
de complexes.
\end{lemma}

\begin{proof}
Cela découle du Lemme \ref{algebra-lemma-differentials-base-change} et du fait que
la formation des puissances extérieures commute au changement de base.
\end{proof}

\begin{lemma}
\label{algebra-lemma-de-rham-complex}
Soit $A \to B$ un morphisme d'anneaux. Soit
$\pi : \Omega_{B/A} \to \Omega$ un morphisme surjectif de $B$-modules.
Notons $\text{d} : B \to \Omega$ la composée de $\pi$ avec la dérivation
universelle $\text{d}_{B/A} : B \to \Omega_{B/A}$.
Posons $\Omega^i = \wedge_B^i(\Omega)$.
Supposons que le noyau de $\pi$ soit engendré, comme $B$-module, par des
éléments $\omega \in \Omega_{B/A}$ tels que
$\text{d}_{B/A}(\omega) \in \Omega_{B/A}^2$ soit envoyé sur zéro dans
$\Omega^2$. Il existe alors un complexe de de Rham
$$
\Omega^0 \to \Omega^1 \to \Omega^2 \to \ldots
$$
dont la différentielle est définie par la règle
$$
\text{d} : \Omega^p \to \Omega^{p + 1},\quad
\text{d}\left(f_0\text{d}f_1 \wedge \ldots \wedge \text{d}f_p\right) =
\text{d}f_0 \wedge \text{d}f_1 \wedge \ldots \wedge \text{d}f_p
$$
\end{lemma}

\begin{proof}
Nous allons montrer qu'il existe un diagramme commutatif
$$
\xymatrix{
\Omega_{B/A}^0 \ar[d] \ar[r]_{\text{d}_{B/A}} &
\Omega_{B/A}^1 \ar[d]_\pi \ar[r]_{\text{d}_{B/A}} &
\Omega_{B/A}^2 \ar[d]_{\wedge^2\pi} \ar[r]_{\text{d}_{B/A}} &
\ldots \\
\Omega^0 \ar[r]^{\text{d}} &
\Omega^1 \ar[r]^{\text{d}} &
\Omega^2 \ar[r]^{\text{d}} &
\ldots
}
$$
La description de l'application $\text{d}$ découlera de la construction,
donnée ci-dessus, des différentielles
$\text{d}_{B/A} : \Omega^p_{B/A} \to \Omega^{p + 1}_{B/A}$ du complexe
de de Rham de $B$ sur $A$.
Puisque la flèche verticale la plus à gauche est un isomorphisme, nous avons
le premier carré. Comme $\pi$ est surjectif, pour obtenir le deuxième carré,
il suffit de montrer que $\text{d}_{B/A}$ envoie le noyau de $\pi$ dans le
noyau de $\wedge^2\pi$. Par hypothèse, tout élément du noyau de $\pi$ est de
la forme $\sum b_i\omega_i$ avec $\pi(\omega_i) = 0$ et
$\wedge^2\pi(\text{d}_{B/A}(\omega_i)) = 0$.
La règle de Leibniz pour $\text{d}_{B/A}$ donne
$\text{d}_{B/A}(\sum b_i\omega_i) = \sum b_i \text{d}_{B/A}(\omega_i) +
\sum \text{d}_{B/A}(b_i) \wedge \omega_i$. Cet élément est donc envoyé sur
zéro par $\wedge^2\pi$.

\medskip\noindent
Pour $i > 1$, notons que $\wedge^i \pi$ est surjectif et que son noyau est
l'image de $\Ker(\pi) \wedge \Omega^{i - 1}_{B/A}
\to \Omega_{B/A}^i$. Pour $\omega_1 \in \Ker(\pi)$ et
$\omega_2 \in \Omega^{i - 1}_{B/A}$, on a
$$
\text{d}_{B/A}(\omega_1 \wedge \omega_2) =
\text{d}_{B/A}(\omega_1) \wedge \omega_2 -
\omega_1 \wedge \text{d}_{B/A}(\omega_2)
$$
qui appartient au noyau de $\wedge^{i + 1}\pi$ d'après ce que nous venons
de démontrer. Nous obtenons donc le $(i + 1)$-ième carré du diagramme
ci-dessus. Cela achève la démonstration.
\end{proof}


\section{Opérateurs différentiels d'ordre fini}
\label{algebra-section-differential-operators}

\noindent
Dans cette section, nous introduisons les opérateurs différentiels d'ordre
fini.

\begin{definition}
\label{algebra-definition-differential-operators}
Soit $R \to S$ un morphisme d'anneaux. Soient $M$, $N$ des $S$-modules.
Soit $k \geq 0$ un entier. Nous définissons par récurrence un
{\it opérateur différentiel $D : M \to N$ d'ordre $k$} comme une application
$R$-linéaire telle que, pour tout $g \in S$, l'application
$m \mapsto D(gm) - gD(m)$ soit un opérateur différentiel d'ordre $k - 1$.
Pour le cas initial $k = 0$, nous définissons un opérateur différentiel
d'ordre $0$ comme une application $S$-linéaire.
\end{definition}

\noindent
Si $D : M \to N$ est un opérateur différentiel d'ordre $k$, alors, pour tout
$g \in S$, l'application $gD$ est un opérateur différentiel d'ordre $k$.
La somme de deux opérateurs différentiels d'ordre $k$ en est encore un.
L'ensemble de tous ces opérateurs
$$
\text{Diff}^k(M, N) = \text{Diff}^k_{S/R}(M, N)
$$
est donc un $S$-module. On a
$$
\text{Diff}^0(M, N) \subset
\text{Diff}^1(M, N) \subset
\text{Diff}^2(M, N) \subset \ldots
$$

\begin{lemma}
\label{algebra-lemma-composition-differential-operators}
Soit $R \to S$ un morphisme d'anneaux. Soient $L, M, N$ des $S$-modules.
Si $D : L \to M$ et $D' : M \to N$ sont des opérateurs différentiels
d'ordres $k$ et $k'$, alors $D' \circ D$ est un opérateur différentiel
d'ordre $k + k'$.
\end{lemma}

\begin{proof}
Soit $g \in S$. L'application qui envoie $x \in L$ sur
$$
D'(D(gx)) - gD'(D(x)) = D'(D(gx)) - D'(gD(x)) + D'(gD(x)) - gD'(D(x))
$$
est une somme de deux composées d'opérateurs différentiels d'ordres
inférieurs. Le lemme découle donc d'une récurrence sur $k + k'$.
\end{proof}

\begin{lemma}
\label{algebra-lemma-module-principal-parts}
Soit $R \to S$ un morphisme d'anneaux. Soit $M$ un $S$-module.
Soit $k \geq 0$. Il existe un $S$-module $P^k_{S/R}(M)$ et un isomorphisme
canonique
$$
\text{Diff}^k_{S/R}(M, N) = \Hom_S(P^k_{S/R}(M), N)
$$
fonctoriel en le $S$-module $N$.
\end{lemma}

\noindent
Par le lemme de Yoneda, cela nous dit qu'il existe un opérateur différentiel
« universel » $D_{univ} : M \to P^k_{S/R}(M)$ d'ordre $k$ tel que, pour tout
opérateur différentiel $D : M \to N$ d'ordre $k$, il existe une unique
application $S$-linéaire $\alpha : P^k_{S/R}(M) \to N$ telle que
$D = D_{univ} \circ \alpha$. Le couple
$(P^k_{S/R}(M), D_{univ})$ est unique à isomorphisme unique près.

\begin{proof}
L'existence de $P^k_{S/R}(M)$ découle d'arguments généraux de théorie des
catégories (insérer ici une référence future), mais nous allons également
en donner une construction. Posons $F = \bigoplus_{m \in M} S[m]$, où
$[m]$ est un symbole désignant l'élément de base du facteur correspondant
à $m$. Étant donné un opérateur différentiel $D : M \to N$, nous obtenons
une application $S$-linéaire $L_D : F \to N$ qui envoie $[m]$ sur $D(m)$.
Si $D$ est d'ordre $0$, alors $L_D$ annule les éléments
$$
[m + m'] - [m] - [m'],\quad
g_0[m] - [g_0m]
$$
où $g_0 \in S$ et $m, m' \in M$.
Si $D$ est d'ordre $1$, alors $L_D$ annule les éléments
$$
[m + m'] - [m] - [m'],\quad
f[m] - [fm], \quad
g_0g_1[m] - g_0[g_1m] - g_1[g_0m] + [g_1g_0m]
$$
où $f \in R$, $g_0, g_1 \in S$ et $m \in M$.
Si $D$ est d'ordre $k$, alors $L_D$ annule les éléments
$[m + m'] - [m] - [m']$, $f[m] - [fm]$, ainsi que les éléments
$$
g_0g_1\ldots g_k[m] - \sum g_0 \ldots \hat g_i \ldots g_k[g_im] + \ldots
+(-1)^{k + 1}[g_0\ldots g_km]
$$
Réciproquement, si $L : F \to N$ est une application $S$-linéaire qui
annule tous les éléments énumérés dans la phrase précédente, alors
$m \mapsto L([m])$ est un opérateur différentiel d'ordre $k$.
Nous voyons ainsi que $P^k_{S/R}(M)$ est le quotient de $F$ par le
sous-module engendré par ces éléments.
\end{proof}

\begin{definition}
\label{algebra-definition-module-principal-parts}
Soit $R \to S$ un morphisme d'anneaux. Soit $M$ un $S$-module. Le module
$P^k_{S/R}(M)$ construit dans le Lemme \ref{algebra-lemma-module-principal-parts}
est appelé le {\it module des parties principales d'ordre $k$} de $M$.
\end{definition}

\noindent
Notons que les inclusions
$$
\text{Diff}^0(M, N) \subset
\text{Diff}^1(M, N) \subset
\text{Diff}^2(M, N) \subset \ldots
$$
correspondent, par le lemme de Yoneda (Catégories,
Lemme \ref{categories-lemma-yoneda}), aux surjections
$$
\ldots \to P^2_{S/R}(M) \to P^1_{S/R}(M) \to P^0_{S/R}(M) = M
$$

\begin{example}
\label{algebra-example-derivations-and-differential-operators}
Soit $R \to S$ un morphisme d'anneaux et soit $N$ un $S$-module.
Observons que $\text{Diff}^1(S, N) = \text{Der}_R(S, N) \oplus N$.
En effet, si $D : S \to N$ est un opérateur différentiel d'ordre $1$, alors
$\sigma_D : S \to N$, défini par $\sigma_D(g) := D(g) - gD(1)$, est une
$R$-dérivation et $D = \sigma_D + \lambda_{D(1)}$, où
$\lambda_x : S \to N$ est l'application linéaire qui envoie $g$ sur $gx$.
Il s'ensuit que $P^1_{S/R} = \Omega_{S/R} \oplus S$ par la propriété
universelle de $\Omega_{S/R}$.
\end{example}

\begin{lemma}
\label{algebra-lemma-sequence-of-principal-parts}
Soit $R \to S$ un morphisme d'anneaux. Soit $M$ un $S$-module. Il existe
une suite exacte courte canonique
$$
0 \to \Omega_{S/R} \otimes_S M \to P^1_{S/R}(M) \to M \to 0
$$
fonctorielle en $M$, appelée la {\it suite des parties principales}.
\end{lemma}

\begin{proof}
Le morphisme $P^1_{S/R}(M) \to M$ est donné ci-dessus.
Soit $N$ un $S$-module et soit $D : M \to N$ un opérateur différentiel
d'ordre $1$. Pour $m \in M$, l'application
$$
g \longmapsto D(gm) - gD(m)
$$
est une $R$-dérivation $S \to N$ par les axiomes des opérateurs
différentiels d'ordre $1$. Elle correspond donc à une application linéaire
$D_m : \Omega_{S/R} \to N$ déterminée par la règle
$a\text{d}b \mapsto aD(bm) - abD(m)$
(voir le Lemme \ref{algebra-lemma-universal-omega}). L'application
$$
\Omega_{S/R} \times M \longrightarrow N,\quad
(\eta, m) \longmapsto D_m(\eta)
$$
est $S$-bilinéaire (nous omettons les détails) et détermine donc une
application $S$-linéaire
$$
\sigma_D : \Omega_{S/R} \otimes_S M \to N
$$
Nous obtenons ainsi une application
$\text{Diff}^1(M, N) \to \Hom_S(\Omega_{S/R} \otimes_S M, N)$,
$D \mapsto \sigma_D$, fonctorielle en $N$. Par le lemme de Yoneda, cela
correspond à un morphisme
$\Omega_{S/R} \otimes_S M \to P^1_{S/R}(M)$. Il résulte immédiatement de
la construction que ce morphisme est fonctoriel en $M$. La suite
$$
\Omega_{S/R} \otimes_S M \to P^1_{S/R}(M) \to M \to 0
$$
est exacte parce que, pour tout module $N$, la suite
$$
0 \to \Hom_S(M, N) \to
\text{Diff}^1(M, N) \to
\Hom_S(\Omega_{S/R} \otimes_S M, N)
$$
est exacte par inspection.

\medskip\noindent
Pour voir que $\Omega_{S/R} \otimes_S M \to P^1_{S/R}(M)$ est injectif,
raisonnons comme suit. Choisissons une suite exacte
$$
0 \to M' \to F \to M \to 0
$$
où $F$ est un $S$-module libre. Elle induit une suite exacte
$$
0 \to \text{Diff}^1(M, N) \to \text{Diff}^1(F, N) \to \text{Diff}^1(M', N)
$$
pour tout $N$. Cela démontre que, dans le diagramme commutatif
$$
\xymatrix{
0 \ar[r] &
\Omega_{S/R} \otimes_S M' \ar[r] \ar[d] &
P^1_{S/R}(M') \ar[r] \ar[d] &
M' \ar[r] \ar[d] &
0 \\
0 \ar[r] &
\Omega_{S/R} \otimes_S F \ar[r] \ar[d] &
P^1_{S/R}(F) \ar[r] \ar[d] &
F \ar[r] \ar[d] &
0 \\
0 \ar[r] &
\Omega_{S/R} \otimes_S M \ar[r] \ar[d] &
P^1_{S/R}(M) \ar[r] \ar[d] &
M \ar[r] \ar[d] &
0 \\
& 0 & 0 & 0
}
$$
la colonne du milieu est exacte. La colonne de gauche est exacte par
l'exactitude à droite de $\Omega_{S/R} \otimes_S -$. D'après le lemme du
serpent (voir la section \ref{algebra-section-snake}), il suffit de démontrer
l'exactitude à gauche pour le module libre $F$.
En utilisant le fait que $P^1_{S/R}(-)$ commute aux sommes directes, nous
nous ramenons au cas $M = S$. Ce cas découle de la discussion de
l'Exemple \ref{algebra-example-derivations-and-differential-operators}.
\end{proof}

\begin{remark}
\label{algebra-remark-functoriality-principal-parts}
Supposons donnés un diagramme commutatif d'anneaux
$$
\xymatrix{
B \ar[r] & B' \\
A \ar[u] \ar[r] & A' \ar[u]
}
$$
un $B$-module $M$, un $B'$-module $M'$ et une application $B$-linéaire
$M \to M'$. Nous obtenons alors un système compatible de morphismes de
modules
$$
\xymatrix{
\ldots \ar[r] &
P^2_{B'/A'}(M') \ar[r] &
P^1_{B'/A'}(M') \ar[r] &
P^0_{B'/A'}(M') \\
\ldots \ar[r] &
P^2_{B/A}(M) \ar[r] \ar[u] &
P^1_{B/A}(M) \ar[r] \ar[u] &
P^0_{B/A}(M) \ar[u]
}
$$
Ces morphismes sont compatibles avec toute composition ultérieure de
morphismes de ce type. La façon la plus simple de le voir est d'utiliser
la description des modules $P^k_{B/A}(M)$ par générateurs et relations dans
la démonstration du Lemme \ref{algebra-lemma-module-principal-parts}, mais on peut
aussi le voir directement à partir de la propriété universelle de ces
modules. En outre, ces morphismes sont compatibles avec les suites exactes
courtes du Lemme \ref{algebra-lemma-sequence-of-principal-parts}.
\end{remark}

\begin{lemma}
\label{algebra-lemma-principal-parts-base-change}
Supposons donnés des morphismes d'anneaux $A \to A'$ et $A \to B$, ainsi
qu'un $B$-module $M$. Posons $B' = B \otimes_A A'$ et considérons
$M' = M \otimes_A A'$ comme un $B'$-module.
Le morphisme de la Remarque \ref{algebra-remark-functoriality-principal-parts}
induit un isomorphisme
$P^k_{B/A}(M) \otimes_A A' = P^k_{B'/A'}(M')$.
\end{lemma}

\begin{proof}
Soit $D_{univ} : M \to P^k_{B/A}(M)$ l'opérateur différentiel universel.
L'application $A'$-linéaire induite
$D_{univ} \otimes 1 : M' = M \otimes_A A' \to P^k_{B/A}(M) \otimes_A A'$
est aisément vérifiée être un opérateur différentiel d'ordre $k$ pour
$M'/B'/A'$. Nous allons montrer que $D_{univ} \otimes 1$ satisfait la
propriété universelle. Soit $D' : M' \to N'$ un opérateur différentiel
d'ordre $k$ à valeurs dans un $B'$-module $N'$. Alors la composée
$D : M \to M' \to N'$ est un opérateur différentiel à valeurs dans $N'$,
considéré comme un $B$-module. Il existe donc une application $B$-linéaire
$\gamma : P^k_{B/A}(M) \to N'$ telle que $D = \gamma \circ D_{univ}$.
Le lecteur vérifiera aisément que l'application $B'$-linéaire induite
$\gamma' : P^k_{B/A}(M) \otimes_A A' \to N'$ satisfait
$\gamma' \circ (D_{univ} \otimes 1) = D'$.
Nous omettons la démonstration de l'unicité de $\gamma'$.
\end{proof}

\begin{lemma}
\label{algebra-lemma-principal-parts-diagonal}
Soit $R \to S$ un morphisme d'anneaux. Soit $M$ un $S$-module.
Soit $J = \Ker(S \otimes_R S \to S)$ le noyau du morphisme de
multiplication. Il existe un isomorphisme canonique de $S$-modules
$$
P^k_{S/R}(M) \longrightarrow (S \otimes_R M)/J^{k + 1}(S \otimes_R M)
$$
où $s \in S$ agit sur la cible par multiplication par $s \otimes 1$, et tel
que l'opérateur différentiel universel d'ordre $k$ corresponde à
l'application donnée par $m \mapsto \text{classe de }1 \otimes m$.
\end{lemma}

\begin{proof}
Considérons l'application $T : M \to S \otimes M$,
$m \mapsto 1 \otimes m$. Puisque $T$ est $R$-linéaire et puisque
\begin{align*}
g_0g_1\ldots g_k T(m)
- \sum g_0 \ldots \hat g_i \ldots g_k T(g_im) + \ldots
+(-1)^{k + 1}T(g_0\ldots g_km) \\
=
\prod_{i = 0,\ldots,k} (g_i \otimes 1 - 1 \otimes g_i) 1 \otimes m
\end{align*}
pour $m \in M$ et $g_0, \ldots, g_k \in S$, nous concluons que la règle
$m \mapsto \text{classe de }1 \otimes m$ de l'énoncé du lemme est un
opérateur différentiel d'ordre $k$ (voir la démonstration du
Lemme \ref{algebra-lemma-module-principal-parts}).
Par la propriété universelle de $P^k_{S/R}(M)$, nous obtenons la flèche de
l'énoncé du lemme. D'autre part, si $D : M \to N$ est un opérateur
différentiel d'ordre $k$, nous pouvons considérer l'application $S$-linéaire
$L : S \otimes_R M \to N$, $g \otimes m \mapsto gD(m)$.
Le lecteur vérifie, par un calcul semblable à celui ci-dessus et en utilisant
le fait que $J$ est engendré par des éléments de la forme
$g \otimes 1 - 1 \otimes g$, que $L$ annule
$J^{k + 1}(S \otimes_R M)$. En particulier, si nous appliquons cela à
l'opérateur différentiel universel
$D_{univ} : M \to P^k_{S/R}(M)$, nous obtenons une application $S$-linéaire
$(S \otimes_R M)/J^{k + 1}(S \otimes_R M) \to P^k_{S/R}(M)$.
Nous laissons au lecteur le soin de vérifier que cette application est
l'inverse de celle de l'énoncé du lemme.
\end{proof}

\begin{lemma}
\label{algebra-lemma-differentials-de-rham-complex-order-1}
Soit $A \to B$ un morphisme d'anneaux. Les différentielles
$\text{d} : \Omega^i_{B/A}  \to \Omega^{i + 1}_{B/A}$ sont des opérateurs
différentiels d'ordre $1$.
\end{lemma}

\begin{proof}
Étant donné $b \in B$, nous devons montrer que
$\text{d} \circ b - b \circ \text{d}$ est un opérateur linéaire.
Nous devons donc montrer que
$$
\text{d} \circ b \circ b' - b \circ \text{d} \circ b' -
b' \circ \text{d} \circ b + b' \circ b \circ \text{d} = 0
$$
Pour le voir, il suffit de vérifier cette égalité sur des générateurs
additifs de $\Omega^i_{B/A}$. Il suffit donc de montrer que
$$
\text{d}(bb'b_0\text{d}b_1 \wedge \ldots \wedge \text{d}b_i) -
b\text{d}(b'b_0\text{d}b_1 \wedge \ldots \wedge \text{d}b_i) -
b'\text{d}(bb_0\text{d}b_1 \wedge \ldots \wedge \text{d}b_i) +
bb'\text{d}(b_0\text{d}b_1 \wedge \ldots \wedge \text{d}b_i)
$$
est nul. C'est un agréable calcul utilisant la règle de Leibniz, que nous
laissons au lecteur.
\end{proof}

\begin{lemma}
\label{algebra-lemma-check-differential-operators}
Soit $A \to B$ un morphisme d'anneaux. Soient $g_i \in B$, $i \in I$, un
ensemble de générateurs de $B$ comme $A$-algèbre. Soient $M, N$ des
$B$-modules. Soit $D : M \to N$ une application $A$-linéaire. Pour montrer
que $D$ est un opérateur différentiel d'ordre $k$, il suffit de montrer que
$D \circ g_i - g_i \circ D$ est un opérateur différentiel d'ordre $k - 1$
pour $i \in I$.
\end{lemma}

\begin{proof}
En effet, nous affirmons que l'ensemble des éléments $g \in B$ tels que
$D \circ g - g \circ D$ soit un opérateur différentiel d'ordre $k - 1$ est
une sous-$A$-algèbre de $B$. Cela découle des relations
$$
D \circ (g + g') - (g + g') \circ D =
(D \circ g - g \circ D) + (D \circ g' - g' \circ D)
$$
et
$$
D \circ gg' - gg' \circ D =
(D \circ g - g \circ D) \circ g' + g \circ (D \circ g' - g' \circ D)
$$
À proprement parler, pour conclure dans le cas des produits, nous utilisons
aussi le Lemme \ref{algebra-lemma-composition-differential-operators}.
\end{proof}

\begin{lemma}
\label{algebra-lemma-invert-system-differential-operators}
Soit $A \to B$ un morphisme d'anneaux. Soient $M, N$ des $B$-modules.
Soit $S \subset B$ une partie multiplicative. Tout opérateur différentiel
$D : M \to N$ d'ordre $k$ se prolonge de manière unique en un opérateur
différentiel $E : S^{-1}M \to S^{-1}N$ d'ordre $k$.
\end{lemma}

\begin{proof}
Raisonnons par récurrence sur $k$. Si $k = 0$, alors $D$ est $B$-linéaire,
et la fonctorialité de la localisation fournit le prolongement.
Étant donné $b \in B$, l'opérateur
$L_b : m \mapsto D(bm) - bD(m)$ est d'ordre $k - 1$. Par récurrence, il
possède donc un unique prolongement en un opérateur différentiel
$E_b : S^{-1}M \to S^{-1}N$ d'ordre $k - 1$.
En outre, un calcul montre que
$L_{b'b} = L_{b'} \circ b + b' \circ L_b$ ; l'unicité donne donc
$E_{b'b} = E_{b'} \circ b + b' \circ E_b$.
De même, nous obtenons $E_{b'} \circ b - b \circ E_{b'} =
E_b \circ b' - b' \circ E_b$.
Pour $m \in M$ et $g \in S$, posons maintenant
$$
E(m/g) = (1/g)(D(m) - E_g(m/g))
$$
Pour montrer que cette définition est indépendante du choix du représentant,
il suffit de montrer que, pour $g' \in S$, l'emploi du représentant
$g'm/g'g$ donne le même résultat. Nous calculons
\begin{align*}
(1/g'g)(D(g'm) - E_{g'g}(g'm/gg'))
& =
(1/gg')(g'D(m) + E_{g'}(m) - E_{g'g}(g'm/gg')) \\
& =
(1/g'g)(g'D(m) - g' E_g(m/g))
\end{align*}
qui est le même qu'auparavant. Il est clair que $E$ est $R$-linéaire puisque
$D$ et $E_g$ sont $R$-linéaires. En prenant $g = 1$ et en utilisant
$E_1 = 0$, nous voyons que $E$ prolonge $D$.
D'après le Lemme \ref{algebra-lemma-check-differential-operators}, il suffit
maintenant de montrer que $E \circ b - b \circ E$ pour $b \in B$ et
$E \circ 1/g' - 1/g' \circ E$ pour $g' \in S$ sont des opérateurs
différentiels d'ordre $k - 1$ afin de montrer que $E$ est un opérateur
différentiel d'ordre $k$. Pour le premier, choisissons un élément $m/g$ de
$S^{-1}M$ et observons que
\begin{align*}
E(b m/g) - bE(m/g)
& =
(1/g)(D(bm) - bD(m) - E_g(bm/g) + bE_g(m/g)) \\
& =
(1/g)(L_b(m) - E_b(m) + gE_b(m/g)) \\
& =
E_b(m/g)
\end{align*}
qui est un opérateur différentiel d'ordre $k - 1$. Enfin, on a
\begin{align*}
E(m/g'g) - (1/g')E(m/g)
& =
(1/g'g)(D(m) - E_{g'g}(m/g'g)) -
(1/g'g)(D(m) - E_g(m/g)) \\
& =
-(1/g')E_{g'}(m/g'g)
\end{align*}
qui est également un opérateur différentiel d'ordre $k - 1$, comme composée
d'applications linéaires (multiplication par $1/g'$ et changement de signe)
et de $E_{g'}$. Nous omettons la démonstration de l'unicité.
\end{proof}

\begin{lemma}
\label{algebra-lemma-base-change-differential-operators}
Soient $R \to A$ et $R \to B$ des morphismes d'anneaux. Soient $M$ et $M'$
des $A$-modules. Soit $D : M \to M'$ un opérateur différentiel d'ordre $k$
relativement à $R \to A$. Soit $N$ un $B$-module quelconque. Alors
l'application
$$
D \otimes \text{id}_N : M \otimes_R N \to M' \otimes_R N
$$
est un opérateur différentiel d'ordre $k$ relativement à
$B \to A \otimes_R B$.
\end{lemma}

\begin{proof}
Il est clair que $D' = D \otimes \text{id}_N$ est $B$-linéaire.
D'après le Lemme \ref{algebra-lemma-check-differential-operators}, il suffit de
montrer que
$$
D' \circ a \otimes 1 -  a \otimes 1 \circ D' =
(D \circ a - a \circ D) \otimes \text{id}_N
$$
est un opérateur différentiel d'ordre $k - 1$, ce qui découle d'une
récurrence sur $k$.
\end{proof}


\section{Le complexe cotangent naïf}
\label{algebra-section-netherlander}

\noindent
Soit $R \to S$ un morphisme d'anneaux. Notons $R[S]$ l'anneau de polynômes
dont les variables sont les éléments $s \in S$. Notons $[s] \in R[S]$ la
variable correspondant à $s \in S$. Ainsi, $R[S]$ est un $R$-module libre
sur les éléments de base $[s_1] \ldots [s_n]$, où
$s_1, \ldots, s_n$ parcourt toutes les suites non ordonnées d'éléments de
$S$. Il existe une surjection canonique
\begin{equation}
\label{algebra-equation-canonical-presentation}
R[S] \longrightarrow S,\quad [s] \longmapsto s
\end{equation}
dont nous notons le noyau $I \subset R[S]$. Une observation simple montre
que $I$ est engendré par les éléments
$[s + s'] - [s] - [s']$, $[s][s'] - [ss']$ et $[r] - r$.
D'après le Lemme \ref{algebra-lemma-differential-seq}, il existe un morphisme
canonique
\begin{equation}
\label{algebra-equation-naive-cotangent-complex}
I/I^2 \longrightarrow \Omega_{R[S]/R} \otimes_{R[S]} S
\end{equation}
dont le conoyau est canoniquement isomorphe à $\Omega_{S/R}$. Observons que
le $S$-module $\Omega_{R[S]/R} \otimes_{R[S]} S$ est libre sur les
générateurs $\text{d}[s]$.

\begin{definition}
\label{algebra-definition-naive-cotangent-complex}
Soit $R \to S$ un morphisme d'anneaux. Le {\it complexe cotangent naïf}
$\NL_{S/R}$ est le complexe de chaînes
(\ref{algebra-equation-naive-cotangent-complex})
$$
\NL_{S/R} = \left(I/I^2 \longrightarrow \Omega_{R[S]/R} \otimes_{R[S]} S\right)
$$
où $I/I^2$ est placé en degré (homologique) $1$ et
$\Omega_{R[S]/R} \otimes_{R[S]} S$ en degré $0$. Nous noterons
$H_1(L_{S/R}) = H_1(\NL_{S/R})$\footnote{Ce module est parfois noté
$\Gamma_{S/R}$ dans la littérature.} l'homologie en degré $1$.
\end{definition}

\noindent
Avant de poursuivre, disons quelques mots du véritable complexe cotangent
(Complexe cotangent, section \ref{cotangent-section-cotangent-ring-map}).
Étant donné un morphisme d'anneaux $R \to S$, il existe une $R$-algèbre
simpliciale canonique $P_\bullet$ dont les termes sont des algèbres de
polynômes et qui est munie d'une équivalence d'homotopie canonique
$$
P_\bullet \longrightarrow S
$$
Le complexe cotangent $L_{S/R}$ de $S$ sur $R$ est défini comme le complexe
de chaînes associé au module simplicial
$$
\Omega_{P_\bullet/R} \otimes_{P_\bullet} S
$$
Le complexe cotangent naïf défini ci-dessus est canoniquement isomorphe à la
troncation $\tau_{\leq 1}L_{S/R}$ (voir
Homologie, section \ref{homology-section-truncations}, et
Complexe cotangent, section \ref{cotangent-section-surjections}).
En particulier, on a bien
$H_1(\NL_{S/R}) = H_1(L_{S/R})$ ; notre définition est donc compatible avec
celle qui utilise le complexe cotangent. En outre,
$H_0(L_{S/R}) = H_0(\NL_{S/R}) = \Omega_{S/R}$, comme nous l'avons vu
ci-dessus.

\medskip\noindent
Soit $R \to S$ un morphisme d'anneaux. Une {\it présentation de $S$ sur
$R$} est une surjection $\alpha : P \to S$ de $R$-algèbres, où $P$ est une
algèbre de polynômes (en un ensemble de variables). Souvent, lorsque $S$ est
de type fini sur $R$, nous l'indiquerons en disant : « Soit
$R[x_1, \ldots, x_n] \to S$ une présentation de $S/R$ », ou « Soit
$0 \to I \to R[x_1, \ldots, x_n] \to S \to 0$ une présentation de $S/R$ »
si nous voulons indiquer que $I$ est le noyau de la présentation.
Notons que le morphisme $R[S] \to S$ utilisé pour définir le complexe
cotangent naïf est un exemple de présentation.

\medskip\noindent
Notons que toute présentation $\alpha$ donne un complexe de chaînes à deux
termes de $S$-modules
$$
\NL(\alpha) : I/I^2 \longrightarrow \Omega_{P/R} \otimes_P S.
$$
Ici, le terme $I/I^2$ est placé en degré $1$ et le terme
$\Omega_{P/R} \otimes S$ en degré $0$. La classe de $f \in I$ dans $I/I^2$
est envoyée sur $\text{d}f \otimes 1$ dans $\Omega_{P/R} \otimes S$.
Le conoyau de ce complexe est canoniquement $\Omega_{S/R}$ ; voir le
Lemme \ref{algebra-lemma-differential-seq}. Nous appelons le complexe $\NL(\alpha)$
le {\it complexe cotangent naïf associé à la présentation
$\alpha : P \to S$ de $S/R$}. Notons que, si $P = R[S]$ muni de sa
surjection canonique sur $S$, nous retrouvons $\NL_{S/R}$.
Si $P = R[x_1, \ldots, x_n]$, nous utiliserons parfois la notation
$I/I^2 \to \bigoplus_{i = 1, \ldots, n} S\text{d}x_i$ pour désigner ce
complexe.

\medskip\noindent
Supposons donné un diagramme commutatif
\begin{equation}
\label{algebra-equation-functoriality-NL}
\vcenter{
\xymatrix{
S \ar[r]_{\phi} & S' \\
R \ar[r] \ar[u] & R' \ar[u]
}
}
\end{equation}
d'anneaux. Soit $\alpha : P \to S$ une présentation de $S$ sur $R$, et
soit $\alpha' : P' \to S'$ une présentation de $S'$ sur $R'$.
Un {\it morphisme de présentations de $\alpha : P \to S$ vers
$\alpha' : P' \to S'$} est, par définition, un morphisme de $R$-algèbres
$$
\varphi : P \to P'
$$
tel que $\phi \circ \alpha = \alpha' \circ \varphi$. Notons que, dans ce
cas, $\varphi(I) \subset I'$, où $I = \Ker(\alpha)$ et
$I' = \Ker(\alpha')$. Ainsi, $\varphi$ induit un morphisme de $S$-modules
$I/I^2 \to I'/(I')^2$ et, par fonctorialité des différentielles, également
un morphisme de $S$-modules
$\Omega_{P/R} \otimes S \to \Omega_{P'/R'} \otimes S'$.
Ces morphismes sont compatibles avec les différentielles de $\NL(\alpha)$
et $\NL(\alpha')$, et nous obtenons un morphisme de complexes cotangents
naïfs
$$
\NL(\alpha) \longrightarrow \NL(\alpha').
$$
Il est souvent commode de considérer le morphisme induit
$\NL(\alpha) \otimes_S S' \to \NL(\alpha')$.

\medskip\noindent
Dans le cas particulier où $P = R[S]$ et $P' = R'[S']$, le morphisme
$\phi : S \to S'$ induit un morphisme canonique d'anneaux
$\varphi : P \to P'$ par la règle $[s] \mapsto [\phi(s)]$.
La construction ci-dessus détermine donc des morphismes canoniques (!) de
complexes de chaînes
$$
\NL_{S/R} \longrightarrow \NL_{S'/R'},\quad\text{et}\quad
\NL_{S/R} \otimes_S S' \longrightarrow \NL_{S'/R'}
$$
associés au diagramme (\ref{algebra-equation-functoriality-NL}). Notons que cette
construction est compatible à la composition : étant donné un diagramme
commutatif
$$
\xymatrix{
S \ar[r]_{\phi} & S' \ar[r]_{\phi'} & S'' \\
R \ar[r] \ar[u] & R' \ar[u] \ar[r] & R'' \ar[u]
}
$$
nous voyons que la composée de
$$
\NL_{S/R} \longrightarrow \NL_{S'/R'} \longrightarrow \NL_{S''/R''}
$$
est le morphisme $\NL_{S/R} \to \NL_{S''/R''}$ donné par le carré extérieur.

\medskip\noindent
Il se trouve que $\NL(\alpha)$ est homotopiquement équivalent à $\NL_{S/R}$
et que les morphismes construits ci-dessus sont bien définis à homotopie
près (les homotopies de morphismes de complexes sont discutées dans
Homologie, section \ref{homology-section-complexes}, mais nous explicitons
également dans la démonstration du lemme ci-dessous le sens exact de ces
énoncés).

\begin{lemma}
\label{algebra-lemma-NL-homotopy}
Supposons donné un diagramme (\ref{algebra-equation-functoriality-NL}).
Soient $\alpha : P \to S$ et $\alpha' : P' \to S'$ des présentations.
\begin{enumerate}
\item Il existe un morphisme de présentations de $\alpha$ vers $\alpha'$.
\item Deux morphismes de présentations quelconques induisent des morphismes
de complexes homotopes $\NL(\alpha) \to \NL(\alpha')$.
\item La construction est compatible aux compositions de morphismes de
présentations (voir la démonstration pour l'énoncé exact).
\item Si $R \to R'$ et $S \to S'$ sont des isomorphismes, alors, pour tout
morphisme $\varphi$ de présentations de $\alpha$ vers $\alpha'$, le morphisme
induit $\NL(\alpha) \to \NL(\alpha')$ est une équivalence d'homotopie et un
quasi-isomorphisme.
\end{enumerate}
En particulier, en comparant $\alpha$ à la présentation canonique
(\ref{algebra-equation-canonical-presentation}), nous concluons qu'il existe un
quasi-isomorphisme $\NL(\alpha) \to \NL_{S/R}$ bien défini à homotopie près
et compatible avec toutes les fonctorialités (à homotopie près).
\end{lemma}

\begin{proof}
Puisque $P$ est une algèbre de polynômes sur $R$, nous pouvons écrire
$P = R[x_a, a \in A]$ pour un certain ensemble $A$.
Comme $\alpha'$ est surjectif, nous pouvons choisir, pour tout $a \in A$,
un élément $f_a \in P'$ tel que
$\alpha'(f_a) = \phi(\alpha(x_a))$. Soit
$\varphi : P = R[x_a, a \in A] \to P'$ l'unique morphisme de $R$-algèbres
tel que $\varphi(x_a) = f_a$. Cela donne le morphisme de (1).

\medskip\noindent
Soient $\varphi$ et $\varphi'$ des morphismes de présentations de $\alpha$
vers $\alpha'$. Soient $I = \Ker(\alpha)$ et $I' = \Ker(\alpha')$.
Nous devons construire le morphisme diagonal $h$ dans le diagramme
$$
\xymatrix{
I/I^2 \ar[r]^-{\text{d}}
\ar@<1ex>[d]^{\varphi'_1} \ar@<-1ex>[d]_{\varphi_1}
&
\Omega_{P/R} \otimes_P S
\ar@<1ex>[d]^{\varphi'_0} \ar@<-1ex>[d]_{\varphi_0}
\ar[ld]_h
\\
I'/(I')^2 \ar[r]^-{\text{d}}
&
\Omega_{P'/R'} \otimes_{P'} S'
}
$$
où les morphismes verticaux sont induits par $\varphi$, $\varphi'$ et tels
que
$$
\varphi_1 - \varphi'_1 = h \circ \text{d}
\quad\text{et}\quad
\varphi_0 - \varphi'_0 = \text{d} \circ h
$$
Considérons l'application $\varphi - \varphi' : P \to P'$. Puisque
$\varphi$ et $\varphi'$ sont toutes deux compatibles avec $\alpha$ et
$\alpha'$, nous obtenons $\varphi - \varphi' : P \to I'$.
Cela implique que $\varphi, \varphi' : P \to P'$ induisent la même structure
de $P$-module sur $I'/(I')^2$, puisque
$\varphi(p)i' - \varphi'(p)i' = (\varphi - \varphi')(p)i' \in (I')^2$.
De plus, $\varphi - \varphi'$ est $R$-linéaire et
$$
(\varphi - \varphi')(fg) =
\varphi(f)(\varphi - \varphi')(g) + (\varphi - \varphi')(f)\varphi'(g)
$$
L'application induite $D : P \to I'/(I')^2$ est donc une $R$-dérivation.
Nous obtenons ainsi un morphisme canonique
$h : \Omega_{P/R} \otimes_P S \to I'/(I')^2$ tel que
$D = h \circ \text{d}$.
Un calcul (omis) montre que $h$ est l'homotopie souhaitée.

\medskip\noindent
Supposons que nous ayons un diagramme commutatif
$$
\xymatrix{
S \ar[r]_{\phi} & S' \ar[r]_{\phi'} & S'' \\
R \ar[r] \ar[u] & R' \ar[u] \ar[r] & R'' \ar[u]
}
$$
et que
\begin{enumerate}
\item $\alpha : P \to S$,
\item $\alpha' : P' \to S'$, et
\item $\alpha'' : P'' \to S''$
\end{enumerate}
soient des présentations. Supposons que
\begin{enumerate}
\item $\varphi : P \to P'$ soit un morphisme de présentations de
$\alpha$ vers $\alpha'$, et
\item $\varphi' : P' \to P''$ soit un morphisme de présentations de
$\alpha'$ vers $\alpha''$.
\end{enumerate}
Il est alors immédiat que
$\varphi' \circ \varphi : P \to P''$ est un morphisme de présentations de
$\alpha$ vers $\alpha''$ et que le morphisme induit
$\NL(\alpha) \to \NL(\alpha'')$ de complexes cotangents naïfs est la
composée des morphismes $\NL(\alpha) \to \NL(\alpha')$ et
$\NL(\alpha') \to \NL(\alpha'')$ induits par $\varphi$ et $\varphi'$.

\medskip\noindent
Dans le cas simple des complexes à deux termes, un quasi-isomorphisme est
simplement un morphisme qui induit un isomorphisme à la fois sur le conoyau
et sur le noyau des morphismes entre les termes. Notons que des morphismes
homotopes de complexes à deux termes (au sens expliqué ci-dessus)
définissent les mêmes morphismes sur le noyau et le conoyau. Ainsi, si
$\varphi$ est un morphisme d'une présentation $\alpha$ de $S$ sur $R$ vers
elle-même, alors le morphisme induit
$\NL(\alpha) \to \NL(\alpha)$ est un quasi-isomorphisme, car il est homotope
à l'identité d'après (2). Pour démontrer (4) en toute généralité,
considérons un morphisme $\varphi'$ de $\alpha'$ vers $\alpha$, qui existe
d'après (1). Les composées
$\NL(\alpha) \to \NL(\alpha') \to \NL(\alpha)$ et
$\NL(\alpha') \to \NL(\alpha) \to \NL(\alpha')$ sont homotopes aux
morphismes identités d'après (3) ; ces morphismes sont donc, par définition,
des équivalences d'homotopie. Il s'ensuit formellement que les deux
morphismes $\NL(\alpha) \to \NL(\alpha')$ et
$\NL(\alpha') \to \NL(\alpha)$ sont des quasi-isomorphismes.
Nous omettons certains détails.
\end{proof}

\begin{lemma}
\label{algebra-lemma-NL-polynomial-algebra}
Soit $A \to B$ une algèbre de polynômes. Alors $\NL_{B/A}$ est
homotopiquement équivalent au complexe de chaînes
$(0 \to \Omega_{B/A})$, où $\Omega_{B/A}$ est en degré $0$.
\end{lemma}

\begin{proof}
Cela découle du Lemme \ref{algebra-lemma-NL-homotopy} et du fait que
$\text{id}_B : B \to B$ est une présentation de $B$ sur $A$ de noyau nul.
\end{proof}

\noindent
Le lemme suivant fournit une partie de la motivation pour introduire le
complexe cotangent naïf. Le complexe cotangent étend ceci en une véritable
suite exacte longue de cohomologie. Si $B \to C$ est une intersection
complète locale, on peut prolonger la suite par un zéro à gauche ; voir
Compléments d'algèbre,
Lemme \ref{more-algebra-lemma-transitive-lci-at-end}.

\begin{lemma}[Suite de Jacobi-Zariski]
\label{algebra-lemma-exact-sequence-NL}
Soient $A \to B \to C$ des morphismes d'anneaux. Choisissons une
présentation $\alpha : A[x_s, s \in S] \to B$ de noyau $I$.
Choisissons une présentation $\beta : B[y_t, t \in T] \to C$ de noyau $J$.
Soit $\gamma : A[x_s, y_t] \to C$ la présentation induite de $C$, de noyau
$K$. Nous obtenons alors un diagramme commutatif canonique
$$
\xymatrix{
0 \ar[r] &
\Omega_{A[x_s]/A} \otimes C \ar[r] &
\Omega_{A[x_s, y_t]/A} \otimes C \ar[r] &
\Omega_{B[y_t]/B} \otimes C \ar[r] &
0 \\
&
I/I^2 \otimes C \ar[r] \ar[u] &
K/K^2 \ar[r] \ar[u] &
J/J^2 \ar[r] \ar[u] &
0
}
$$
à lignes exactes. Nous obtenons la suite exacte suivante de groupes
d'homologie
$$
H_1(\NL_{B/A} \otimes_B C) \to
H_1(L_{C/A}) \to
H_1(L_{C/B}) \to
C \otimes_B \Omega_{B/A} \to
\Omega_{C/A} \to
\Omega_{C/B} \to 0
$$
de $C$-modules, qui prolonge la suite du
Lemme \ref{algebra-lemma-exact-sequence-differentials}.
Si $\text{Tor}_1^B(\Omega_{B/A}, C) = 0$ et
$\text{Tor}_2^B(\Omega_{B/A}, C) = 0$, alors
$H_1(\NL_{B/A} \otimes_B C) = H_1(L_{B/A}) \otimes_B C$.
\end{lemma}

\begin{proof}
Nous omettons la définition précise des morphismes.
L'exactitude de la ligne supérieure découle du fait que les
$\text{d}x_s$, $\text{d}y_t$ forment une base du module du milieu.
Le morphisme $\gamma$ se factorise en
$$
A[x_s, y_t] \to B[y_t] \to C
$$
où la première flèche est surjective et la seconde est égale à $\beta$.
Nous voyons donc que $K \to J$ est surjectif.
En outre, le noyau de la première flèche affichée est
$IA[x_s, y_t]$. Par conséquent, $I/I^2 \otimes C$ se surjecte sur le noyau
de $K/K^2 \to J/J^2$. Enfin, nous pouvons utiliser le
Lemme \ref{algebra-lemma-NL-homotopy} pour identifier les termes aux groupes
d'homologie des complexes cotangents naïfs.

\medskip\noindent
La dernière assertion est un énoncé d'algèbre homologique.
Rappelons que $\NL_{B/A} = (N^{-1} \to N^0)$ est un complexe à deux termes
de $B$-modules, où $N^0$ est libre, et que ses modules de cohomologie sont
$H^0 = \Omega_{B/A}$ et $H^{-1} = H_1(L_{B/A})$.
Notons $M \subset N^0$ l'image de la différentielle.
Si $\text{Tor}_1^B(H^0, C) = 0$, alors nous avons une suite exacte
$$
0 \to M \otimes_B C \to N^0 \otimes_B C \to H^0 \otimes_B C \to 0
$$
Puisque $N^0$ est libre, nous voyons aussi que
$\text{Tor}_2^B(H^0, C) =
\text{Tor}_1^B(M, C)$. Ainsi, si $\text{Tor}_2^B(H^0, C) = 0$, nous avons
également une suite exacte
$$
0 \to H^{-1} \otimes_B C \to N^{-1} \otimes_B C \to M \otimes_B C \to 0
$$
En rassemblant ces faits, nous voyons que, si
$\text{Tor}_1^B(H^0, C) = 0$ et $\text{Tor}_2^B(H^0, C) = 0$, alors
$H^{-1} \otimes_B C$ est le noyau de
$N^{-1} \otimes_B C \to N^0 \otimes_B C$, comme souhaité.
\end{proof}

\begin{remark}
\label{algebra-remark-composition-homotopy-equivalent-to-zero}
Soient $A \to B$ et $\phi : B \to C$ des morphismes d'anneaux.
Alors la composée $\NL_{B/A} \to \NL_{C/A} \to \NL_{C/B}$ est
homotopiquement équivalente à zéro. En effet, cette composée est la
fonctorialité du complexe cotangent naïf pour le carré
$$
\xymatrix{
B \ar[r]_\phi & C \\
A \ar[r] \ar[u] & B \ar[u]
}
$$
Écrivons $J = \Ker(B[C] \to C)$. Une homotopie explicite est donnée par le
morphisme $\Omega_{A[B]/A} \otimes_{A[B]} B \to J/J^2$ qui envoie l'élément
de base $\text{d}[b]$ sur la classe de $[\phi(b)] - b$ dans $J/J^2$.
\end{remark}

\begin{lemma}
\label{algebra-lemma-NL-surjection}
Soit $A \to B$ un morphisme surjectif d'anneaux, de noyau $I$.
Alors $\NL_{B/A}$ est homotopiquement équivalent au complexe de chaînes
$(I/I^2 \to 0)$, où $I/I^2$ est en degré $1$. En particulier,
$H_1(L_{B/A}) = I/I^2$.
\end{lemma}

\begin{proof}
Cela découle du Lemme \ref{algebra-lemma-NL-homotopy} et du fait que $A \to B$ est
une présentation de $B$ sur $A$.
\end{proof}

\begin{lemma}
\label{algebra-lemma-application-NL}
Soient $A \to B \to C$ des morphismes d'anneaux. Supposons que $A \to C$
soit surjectif (donc $B \to C$ l'est aussi). Notons
$I = \Ker(A \to C)$ et $J = \Ker(B \to C)$. Alors la suite
$$
I/I^2 \to J/J^2 \to \Omega_{B/A} \otimes_B B/J \to 0
$$
est exacte.
\end{lemma}

\begin{proof}
Cela découle du Lemme \ref{algebra-lemma-exact-sequence-NL} et de la description des
complexes cotangents naïfs $\NL_{C/B}$ et $\NL_{C/A}$ dans le
Lemme \ref{algebra-lemma-NL-surjection}.
\end{proof}

\begin{lemma}[Changement de base plat]
\label{algebra-lemma-change-base-NL}
Soit $R \to S$ un morphisme d'anneaux. Soit $\alpha : P \to S$ une
présentation. Soit $R \to R'$ un morphisme plat d'anneaux.
Soit $\alpha' : P \otimes_R R' \to S' = S \otimes_R R'$ la présentation
induite. Alors
$\NL(\alpha) \otimes_R R' = \NL(\alpha) \otimes_S S' = \NL(\alpha')$.
En particulier, le morphisme canonique
$$
\NL_{S/R} \otimes_S S' \longrightarrow \NL_{S \otimes_R R'/R'}
$$
est une équivalence d'homotopie si $R \to R'$ est plat.
\end{lemma}

\begin{proof}
Cela est vrai parce que
$\Ker(\alpha') = R' \otimes_R \Ker(\alpha)$, puisque $R \to R'$ est plat.
\end{proof}

\begin{lemma}
\label{algebra-lemma-colimits-NL}
Soit $R_\lambda \to S_\lambda$ un système de morphismes d'anneaux indexé
par l'ensemble filtrant $\Lambda$.
Posons $R = \colim R_\lambda$ et $S = \colim S_\lambda$.
Alors $\NL_{S/R} = \colim \NL_{S_\lambda/R_\lambda}$.
\end{lemma}

\begin{proof}
Rappelons que $\NL_{S/R}$ est le complexe
$I/I^2 \to \bigoplus_{s \in S} S\text{d}[s]$, où $I \subset R[S]$ est le
noyau de la présentation canonique $R[S] \to S$.
Il est maintenant clair que $R[S] = \colim R_\lambda[S_\lambda]$ et, de
même, que $I = \colim I_\lambda$, où
$I_\lambda = \Ker(R_\lambda[S_\lambda] \to S_\lambda)$.
Le lemme est donc clair.
\end{proof}

\begin{lemma}
\label{algebra-lemma-NL-of-localization}
Si $S \subset A$ est une partie multiplicative de $A$, alors
$\NL_{S^{-1}A/A}$ est homotopiquement équivalent au complexe nul.
\end{lemma}

\begin{proof}
Puisque $A \to S^{-1}A$ est plat, nous voyons que
$\NL_{S^{-1}A/A} \otimes_A S^{-1}A \to \NL_{S^{-1}A/S^{-1}A}$ est une
équivalence d'homotopie par changement de base plat
(Lemme \ref{algebra-lemma-change-base-NL}). Comme la source de la flèche est
isomorphe à $\NL_{S^{-1}A/A}$ et que sa cible est nulle (d'après le
Lemme \ref{algebra-lemma-NL-surjection}), le résultat en découle.
\end{proof}

\begin{lemma}
\label{algebra-lemma-NL-localize-bottom}
Soit $S \subset A$ une partie multiplicative de $A$.
Soit $S^{-1}A \to B$ un morphisme d'anneaux.
Alors $\NL_{B/A} \to \NL_{B/S^{-1}A}$ est une équivalence d'homotopie.
\end{lemma}

\begin{proof}
Choisissons une présentation $\alpha : P \to B$ de $B$ sur $A$.
Alors $\beta : S^{-1}P \to B$ est une présentation de $B$ sur $S^{-1}A$.
Un calcul direct montre que $\NL(\alpha) = \NL(\beta)$, ce qui démontre le
lemme puisque le complexe cotangent naïf est bien défini à homotopie près
d'après le Lemme \ref{algebra-lemma-NL-homotopy}.
\end{proof}

\begin{lemma}
\label{algebra-lemma-principal-localization-NL}
\begin{slogan}
La formation du complexe cotangent naïf commute à la localisation en un
élément.
\end{slogan}
Soit $A \to B$ un morphisme d'anneaux. Soit $g \in B$.
Supposons que $\alpha : P \to B$ soit une présentation de noyau $I$.
Une présentation de $B_g$ sur $A$ est alors le morphisme
$$
\beta : P[x] \longrightarrow B_g
$$
qui prolonge $\alpha$ et envoie $x$ sur $1/g$.
Le noyau $J$ de $\beta$ est engendré par $I$ et par l'élément $f x - 1$,
où $f \in P$ est un élément envoyé sur $g \in B$ par $\alpha$.
Dans cette situation, on a
\begin{enumerate}
\item $J/J^2 = (I/I^2)_g \oplus B_g (f x - 1)$,
\item $\Omega_{P[x]/A} \otimes_{P[x]} B_g =
\Omega_{P/A} \otimes_P B_g \oplus B_g \text{d}x$,
\item $\NL(\beta) \cong
\NL(\alpha) \otimes_B B_g \oplus (B_g \xrightarrow{g} B_g)$.
\end{enumerate}
Le morphisme canonique $\NL_{B/A} \otimes_B B_g \to \NL_{B_g/A}$ est donc
une équivalence d'homotopie.
\end{lemma}

\begin{proof}
Puisque $P[x]/(I, fx - 1) = B[x]/(gx - 1) = B_g$, nous obtenons l'assertion
selon laquelle $I$ et $fx - 1$ engendrent $J$. Considérons le diagramme
commutatif
$$
\xymatrix{
0 \ar[r] &
\Omega_{P/A} \otimes B_g \ar[r] &
\Omega_{P[x]/A} \otimes B_g \ar[r] &
\Omega_{B[x]/B} \otimes B_g \ar[r] &
0 \\
&
(I/I^2)_g \ar[r] \ar[u] &
J/J^2 \ar[r] \ar[u] &
(gx - 1)/(gx - 1)^2 \ar[r] \ar[u] &
0
}
$$
dont les lignes sont exactes d'après le
Lemme \ref{algebra-lemma-exact-sequence-NL}.
Le $B_g$-module $\Omega_{B[x]/B} \otimes B_g$ est libre de rang $1$ sur
$\text{d}x$. L'élément $\text{d}x$ du $B_g$-module
$\Omega_{P[x]/A} \otimes B_g$ fournit un scindage de la ligne supérieure.
L'élément $gx - 1 \in (gx - 1)/(gx - 1)^2$ est envoyé sur
$g\text{d}x$ dans $\Omega_{B[x]/B} \otimes B_g$ ; par conséquent,
$(gx - 1)/(gx - 1)^2$ est libre de rang $1$ sur $B_g$.
(On peut aussi le voir en remarquant que $gx - 1$ est un non-diviseur de
zéro dans $B[x]$, parce que c'est un polynôme dont le terme constant est
inversible, et que tout non-diviseur de zéro donne une suite quasi-régulière
de longueur $1$ d'après le Lemme \ref{algebra-lemma-regular-quasi-regular}.)

\medskip\noindent
Démontrons que $(I/I^2)_g \to J/J^2$ est injectif. Considérons le morphisme
de $P$-algèbres
$$
\pi : P[x] \to (P/I^2)_f = P_f/I_f^2
$$
qui envoie $x$ sur $1/f$. Puisque $J$ est engendré par $I$ et $fx - 1$,
nous voyons que $\pi(J) \subset (I/I^2)_f = (I/I^2)_g$. Comme cet idéal est
de carré nul, on a $\pi(J^2) = 0$.
Si $a \in I$ est envoyé sur un élément de $J^2$ dans $J$, alors
$\pi(a) = 0$, ce qui implique que $a$ est envoyé sur zéro dans $I_f/I_f^2$.
Cela démontre l'injectivité souhaitée.

\medskip\noindent
Nous avons donc une suite exacte courte de complexes à deux termes
$$
0 \to \NL(\alpha) \otimes_B B_g \to \NL(\beta)
\to (B_g \xrightarrow{g} B_g) \to 0
$$
Une telle suite exacte courte peut toujours être scindée dans la catégorie
des complexes. Dans notre cas particulier, nous pouvons choisir les
scindages
$$
J/J^2 = (I/I^2)_g \oplus B_g (fx - 1)\quad\text{et}\quad
\Omega_{P[x]/A} \otimes B_g = \Omega_{P/A} \otimes B_g \oplus
B_g (g^{-2}\text{d}f + \text{d}x)
$$
Cela fonctionne parce que
$\text{d}(fx - 1) = x\text{d}f + f \text{d}x =
g(g^{-2}\text{d}f + \text{d}x)$ dans $\Omega_{P[x]/A} \otimes B_g$.
\end{proof}

\begin{lemma}
\label{algebra-lemma-localize-NL}
Soit $A \to B$ un morphisme d'anneaux. Soit $S \subset B$ une partie
multiplicative. Le morphisme canonique
$\NL_{B/A} \otimes_B S^{-1}B \to \NL_{S^{-1}B/A}$ est un quasi-isomorphisme.
\end{lemma}

\begin{proof}
On a $S^{-1}B = \colim_{g \in S} B_g$, où nous considérons $S$ comme un
ensemble filtrant (ordonné par divisibilité) ; voir le
Lemme \ref{algebra-lemma-localization-colimit}.
D'après le Lemme \ref{algebra-lemma-principal-localization-NL}, chacun des morphismes
$\NL_{B/A} \otimes_B B_g \to \NL_{B_g/A}$ est un quasi-isomorphisme.
Le lemme découle du Lemme \ref{algebra-lemma-colimits-NL}.
\end{proof}

\begin{lemma}
\label{algebra-lemma-sum-two-terms}
Soit $R$ un anneau.
Soient $A_1 \to A_0$ et $B_1 \to B_0$ des complexes à deux termes.
Supposons qu'il existe des morphismes de complexes
$\varphi : A_\bullet \to B_\bullet$ et
$\psi : B_\bullet \to A_\bullet$ tels que
$\varphi \circ \psi$ et $\psi \circ \varphi$ soient homotopes aux
morphismes identités. Alors
$A_1 \oplus B_0 \cong B_1 \oplus A_0$ comme $R$-modules.
\end{lemma}

\begin{proof}
Choisissons une application $h : A_0 \to A_1$ telle que
$$
\text{id}_{A_1} - \psi_1 \circ \varphi_1 = h \circ d_A
\text{ et }
\text{id}_{A_0} - \psi_0 \circ \varphi_0 = d_A \circ h.
$$
De même, choisissons une application $h' : B_0 \to B_1$ telle que
$$
\text{id}_{B_1} - \varphi_1 \circ \psi_1 = h' \circ d_B
\text{ et }
\text{id}_{B_0} - \varphi_0 \circ \psi_0 = d_B \circ h'.
$$
Un calcul immédiat montre que
$$
\left(
\begin{matrix}
\text{id}_{A_1} & -h' \circ \psi_1 + h \circ \psi_0 \\
0 & \text{id}_{B_0}
\end{matrix}
\right)
=
\left(
\begin{matrix}
\psi_1 & h \\
-d_B & \varphi_0
\end{matrix}
\right)
\left(
\begin{matrix}
\varphi_1 & - h' \\
d_A & \psi_0
\end{matrix}
\right)
$$
Cela montre que les deux matrices du membre de droite sont inversibles et
démontre le lemme.
\end{proof}

\begin{lemma}
\label{algebra-lemma-conormal-module}
Soit $R \to S$ un morphisme d'anneaux de type fini.
Pour toutes présentations $\alpha : R[x_1, \ldots, x_n] \to S$ et
$\beta : R[y_1, \ldots, y_m] \to S$, on a
$$
I/I^2 \oplus S^{\oplus m} \cong J/J^2 \oplus S^{\oplus n}
$$
comme $S$-modules, où $I = \Ker(\alpha)$ et $J = \Ker(\beta)$.
\end{lemma}

\begin{proof}
Voir les Lemmes \ref{algebra-lemma-NL-homotopy} et \ref{algebra-lemma-sum-two-terms}.
\end{proof}

\begin{lemma}
\label{algebra-lemma-conormal-module-localize}
Soit $R \to S$ un morphisme d'anneaux de type fini.
Soit $g \in S$. Pour toutes présentations
$\alpha : R[x_1, \ldots, x_n] \to S$ et
$\beta : R[y_1, \ldots, y_m] \to S_g$, on a
$$
(I/I^2)_g \oplus S^{\oplus m}_g \cong J/J^2 \oplus S_g^{\oplus n}
$$
comme $S_g$-modules, où $I = \Ker(\alpha)$ et $J = \Ker(\beta)$.
\end{lemma}

\begin{proof}
Soit $\beta' : R[x_1, \ldots, x_n, x] \to S_g$ la présentation du
Lemme \ref{algebra-lemma-principal-localization-NL} construite à partir de $\alpha$.
Nous savons alors que $\NL(\alpha) \otimes_S S_g$ est homotopiquement
équivalent à $\NL(\beta')$.
Nous savons que $\NL(\beta)$ et $\NL(\beta')$ sont homotopiquement
équivalents d'après le Lemme \ref{algebra-lemma-NL-homotopy}. Nous concluons que
$\NL(\alpha) \otimes_S S_g$ est homotopiquement équivalent à
$\NL(\beta)$. Enfin, nous appliquons le
Lemme \ref{algebra-lemma-conormal-module}.
\end{proof}


\section{Intersections complètes locales}
\label{algebra-section-lci}

\noindent
Être une intersection complète locale est une propriété intrinsèque d'un
anneau local noethérien. Cela sera discuté dans
Algèbre à puissances divisées, section \ref{dpa-section-lci}.
Cependant, pour le moment, nous définissons seulement cette propriété pour
les algèbres de type fini sur un corps.

\begin{definition}
\label{algebra-definition-lci-field}
Soit $k$ un corps.
Soit $S$ une $k$-algèbre de type fini.
\begin{enumerate}
\item Nous disons que $S$ est une {\it intersection complète globale sur
$k$} s'il existe une présentation
$S = k[x_1, \ldots, x_n]/(f_1, \ldots, f_c)$ telle que
$\dim(S) = n - c$.
\item Nous disons que $S$ est une {\it intersection complète locale sur
$k$} s'il existe un recouvrement $\Spec(S) = \bigcup D(g_i)$ tel que chacun
des anneaux $S_{g_i}$ soit une intersection complète globale sur $k$.
\end{enumerate}
Nous adopterons aussi la convention que l'anneau nul est une intersection
complète globale sur $k$.
\end{definition}

\noindent
Supposons que $S$ soit une intersection complète globale
$S = k[x_1, \ldots, x_n]/(f_1, \ldots, f_c)$ comme dans la
Définition \ref{algebra-definition-lci-field}.
Pour un idéal maximal $\mathfrak m \subset k[x_1, \ldots, x_n]$, on a
$\dim(k[x_1, \ldots, x_n]_\mathfrak m) = n$
(Lemme \ref{algebra-lemma-dim-affine-space}).
Si $(f_1, \ldots, f_c) \subset \mathfrak m$, alors le
Lemme \ref{algebra-lemma-one-equation} donne $\dim(S_\mathfrak m) \geq n - c$.
Comme $\dim(S) = n - c$ d'après la Définition \ref{algebra-definition-lci-field},
nous concluons que $\dim(S_\mathfrak m) = n - c$ pour tout idéal maximal de
$S$ et que $\Spec(S)$ est équidimensionnel
(Topologie, Définition \ref{topology-definition-equidimensional}) de
dimension $n - c$ ; voir le
Lemme \ref{algebra-lemma-dimension-at-a-point-finite-type-over-field}.
Nous utiliserons souvent ce fait sans autre mention.

\begin{lemma}
\label{algebra-lemma-localize-lci}
Soit $k$ un corps.
Soit $S$ une $k$-algèbre de type fini.
Soit $g \in S$.
\begin{enumerate}
\item Si $S$ est une intersection complète globale, alors $S_g$ l'est aussi.
\item Si $S$ est une intersection complète locale, alors $S_g$ l'est aussi.
\end{enumerate}
\end{lemma}

\begin{proof}
La seconde assertion découle immédiatement de la première.
Démontrons la première. Si $S_g$ est l'anneau nul, elle est vraie.
Supposons $S_g$ non nul.
Écrivons $S = k[x_1, \ldots, x_n]/(f_1, \ldots, f_c)$ avec
$n - c = \dim(S)$ comme dans la Définition \ref{algebra-definition-lci-field}.
D'après les remarques qui suivent la définition,
$\dim(S_g) = n - c$. Soit $g' \in k[x_1, \ldots, x_n]$ un élément dont la
classe résiduelle correspond à $g$. Alors
$S_g =  k[x_1, \ldots, x_n, x_{n + 1}]/(f_1, \ldots, f_c, x_{n + 1}g' - 1)$,
comme souhaité.
\end{proof}

\begin{lemma}
\label{algebra-lemma-lci-CM}
Soit $k$ un corps. Soit $S$ une $k$-algèbre de type fini.
Si $S$ est une intersection complète locale, alors $S$ est un anneau de
Cohen-Macaulay.
\end{lemma}

\begin{proof}
Choisissons un idéal premier maximal $\mathfrak m$ de $S$.
Nous devons montrer que $S_\mathfrak m$ est de Cohen-Macaulay.
Par hypothèse, nous pouvons supposer
$S = k[x_1, \ldots, x_n]/(f_1, \ldots, f_c)$ avec
$\dim(S) = n - c$. Soit
$\mathfrak m' \subset k[x_1, \ldots, x_n]$ l'idéal maximal correspondant à
$\mathfrak m$. D'après la
Proposition \ref{algebra-proposition-finite-gl-dim-polynomial-ring}, l'anneau local
$k[x_1, \ldots, x_n]_{\mathfrak m'}$ est régulier local de dimension $n$.
En particulier, il est de Cohen-Macaulay d'après le
Lemme \ref{algebra-lemma-regular-ring-CM}.
En appliquant $c$ fois le Lemme \ref{algebra-lemma-one-equation}, l'anneau local
$S_{\mathfrak m} = k[x_1, \ldots, x_n]_{\mathfrak m'}/(f_1, \ldots, f_c)$
est de dimension $\geq n - c$. Par hypothèse,
$\dim(S_{\mathfrak m}) \leq n - c$. Nous obtenons donc l'égalité.
Cela implique que $f_1, \ldots, f_c$ est une suite régulière dans
$k[x_1, \ldots, x_n]_{\mathfrak m'}$ et que $S_{\mathfrak m}$ est de
Cohen-Macaulay ; voir la Proposition \ref{algebra-proposition-CM-module}.
\end{proof}

\noindent
Le lemme suivant est la clef technique du reste de cette section. Un aspect
important de ce lemme est que nous pouvons choisir n'importe quelle
présentation de l'anneau $S$, mais que la condition (1) ne dépend pas de ce
choix.

\begin{lemma}
\label{algebra-lemma-lci}
Soit $k$ un corps.
Soit $S$ une $k$-algèbre de type fini.
Soit $\mathfrak q$ un idéal premier de $S$.
Choisissons une présentation quelconque $S = k[x_1, \ldots, x_n]/I$.
Soit $\mathfrak q'$ l'idéal premier de $k[x_1, \ldots, x_n]$ correspondant
à $\mathfrak q$. Posons
$c = \text{height}(\mathfrak q') - \text{height}(\mathfrak q)$ ; autrement
dit, $\dim_{\mathfrak q}(S) = n - c$
(voir le Lemme \ref{algebra-lemma-codimension}). Les assertions suivantes sont
équivalentes :
\begin{enumerate}
\item Il existe $g \in S$, $g \not \in \mathfrak q$, tel que $S_g$ soit une
intersection complète globale sur $k$.
\item L'idéal $I_{\mathfrak q'} \subset
k[x_1, \ldots, x_n]_{\mathfrak q'}$ peut être engendré par $c$ éléments.
\item Le module conormal $(I/I^2)_{\mathfrak q}$ peut être engendré par $c$
éléments sur $S_{\mathfrak q}$.
\item Le module conormal $(I/I^2)_{\mathfrak q}$ est un
$S_{\mathfrak q}$-module libre de rang $c$.
\item L'idéal $I_{\mathfrak q'}$ peut être engendré par une suite régulière
dans l'anneau local régulier $k[x_1, \ldots, x_n]_{\mathfrak q'}$.
\end{enumerate}
Dans ce cas, tout choix de $c$ éléments de $I_{\mathfrak q'}$ qui engendrent
$I_{\mathfrak q'}/\mathfrak q'I_{\mathfrak q'}$ forme une suite régulière
dans l'anneau local $k[x_1, \ldots, x_n]_{\mathfrak q'}$.
\end{lemma}

\begin{proof}
Posons $R = k[x_1, \ldots, x_n]_{\mathfrak q'}$. C'est un anneau local de
Cohen-Macaulay de dimension $\text{height}(\mathfrak q')$ ; voir par exemple
le Lemme \ref{algebra-lemma-lci-CM}. En outre,
$\overline{R} = R/IR = R/I_{\mathfrak q'} = S_{\mathfrak q}$ est un quotient
de dimension $\text{height}(\mathfrak q)$.
Soient $f_1, \ldots, f_c \in I_{\mathfrak q'}$ des éléments qui engendrent
$(I/I^2)_{\mathfrak q}$. D'après le Lemme \ref{algebra-lemma-NAK},
$f_1, \ldots, f_c$ engendrent $I_{\mathfrak q'}$.
Comme les dimensions concordent, nous concluons par la
Proposition \ref{algebra-proposition-CM-module} que $f_1, \ldots, f_c$ est une suite
régulière dans $R$.
D'après le Lemme \ref{algebra-lemma-regular-quasi-regular},
$(I/I^2)_{\mathfrak q}$ est libre.
Ces arguments montrent que (2), (3) et (4) sont équivalentes, qu'elles
impliquent la dernière assertion du lemme et, par conséquent, qu'elles
impliquent (5).

\medskip\noindent
Si (5) est satisfaite, disons que $I_{\mathfrak q'}$ est engendré par une
suite régulière de longueur $e$. Alors, par la théorie de la dimension,
$\text{height}(\mathfrak q) = \dim(S_{\mathfrak q}) =
\dim(k[x_1, \ldots, x_n]_{\mathfrak q'}) - e =
\text{height}(\mathfrak q') - e$ ; voir la
section \ref{algebra-section-dimension}. Nous concluons que $e = c$.
Ainsi, (5) implique (2).

\medskip\noindent
Conservons les notations introduites dans le premier paragraphe.
Pour chaque $f_i$, nous pouvons trouver
$d_i \in k[x_1, \ldots, x_n]$, $d_i \not \in \mathfrak q'$, tel que
$f_i' = d_i f_i \in k[x_1, \ldots, x_n]$.
On a encore $I_{\mathfrak q'} = (f_1', \ldots, f_c')R$.
Il existe donc $g' \in k[x_1, \ldots, x_n]$,
$g' \not \in \mathfrak q'$, tel que $I_{g'} = (f_1', \ldots, f_c')$.
Choisissons en outre $g'' \in k[x_1, \ldots, x_n]$,
$g'' \not \in \mathfrak q'$, tel que
$\dim(S_{g''}) = \dim_{\mathfrak q} \Spec(S)$.
D'après le Lemme \ref{algebra-lemma-codimension}, cette dimension est égale à
$n - c$. Enfin, soit $g$ l'image de $g'g''$ dans $S$.
Nous voyons alors que
$$
S_g \cong k[x_1, \ldots, x_n, x_{n + 1}]
/
(f_1', \ldots, f_c', x_{n + 1}g'g'' - 1)
$$
et, par notre choix de $g''$, cet anneau est de dimension $n - c$.
C'est donc une intersection complète globale.
Ainsi, chacune des assertions (2), (3) et (4) implique (1).

\medskip\noindent
Supposons (1). Soit
$S_g \cong k[y_1, \ldots, y_m]/(f_1, \ldots, f_t)$ une présentation de
$S_g$ comme intersection complète globale.
Écrivons $J = (f_1, \ldots, f_t)$. Soit
$\mathfrak q'' \subset k[y_1, \ldots, y_m]$ l'idéal premier correspondant à
$\mathfrak qS_g$. Notons que
$t = m - \dim(S_g) =
\text{height}(\mathfrak q'') - \text{height}(\mathfrak q)$ ; voir le
Lemme \ref{algebra-lemma-codimension} pour la dernière égalité.
Comme on l'a vu dans la démonstration du Lemme \ref{algebra-lemma-lci-CM} (et aussi
ci-dessus), les éléments $f_1, \ldots, f_t$ forment une suite régulière dans
l'anneau local $k[y_1, \ldots, y_m]_{\mathfrak q''}$.
D'après le Lemme \ref{algebra-lemma-regular-quasi-regular},
$(J/J^2)_{\mathfrak q}$ est libre de rang $t$.
D'après le Lemme \ref{algebra-lemma-conormal-module-localize}, on a
$$
J/J^2 \oplus S_g^n \cong (I/I^2)_g \oplus S_g^m
$$
Ainsi, $(I/I^2)_{\mathfrak q}$ est libre de rang
$t + n - m = m - \dim(S_g) + n - m = n - \dim(S_g) =
\text{height}(\mathfrak q') - \text{height}(\mathfrak q) = c$.
Nous obtenons donc (4).
\end{proof}

\noindent
Le résultat du Lemme \ref{algebra-lemma-lci} suggère la définition suivante.

\begin{definition}
\label{algebra-definition-lci-local-ring}
Soit $k$ un corps. Soit $S$ une $k$-algèbre locale essentiellement de type
fini sur $k$. Nous disons que $S$ est une {\it intersection complète (sur
$k$)} s'il existe une $k$-algèbre locale $R$ et des éléments
$f_1, \ldots, f_c \in \mathfrak m_R$ tels que
\begin{enumerate}
\item $R$ soit essentiellement de type fini sur $k$,
\item $R$ soit un anneau local régulier,
\item $f_1, \ldots, f_c$ forment une suite régulière dans $R$, et
\item $S \cong R/(f_1, \ldots, f_c)$ comme $k$-algèbres.
\end{enumerate}
\end{definition}

\noindent
D'après le théorème de structure de Cohen (voir le
Théorème \ref{algebra-theorem-cohen-structure-theorem}), tout anneau local
noethérien complet s'écrit comme quotient d'un anneau local régulier complet.
Nous pouvons donc utiliser la définition ci-dessus pour définir la notion
d'anneau d'intersection complète pour tout anneau local noethérien complet.
Nous en discuterons dans Algèbre à puissances divisées,
section \ref{dpa-section-lci}. En attendant, le lemme suivant montre qu'une
telle définition a un sens.

\begin{lemma}
\label{algebra-lemma-ci-well-defined}
Soient $A \to B \to C$ des homomorphismes locaux surjectifs d'anneaux.
Supposons que $A$ et $B$ soient des anneaux locaux réguliers. Les assertions
suivantes sont équivalentes :
\begin{enumerate}
\item $\Ker(A \to C)$ est engendré par une suite régulière,
\item $\Ker(A \to C)$ est engendré par $\dim(A) - \dim(C)$ éléments,
\item $\Ker(B \to C)$ est engendré par une suite régulière, et
\item $\Ker(B \to C)$ est engendré par $\dim(B) - \dim(C)$ éléments.
\end{enumerate}
\end{lemma}

\begin{proof}
Un anneau local régulier est de Cohen-Macaulay ; voir le
Lemme \ref{algebra-lemma-regular-ring-CM}. D'où les équivalences
(1) $\Leftrightarrow$ (2) et (3) $\Leftrightarrow$ (4) ; voir la
Proposition \ref{algebra-proposition-CM-module}.
D'après le Lemme \ref{algebra-lemma-regular-quotient-regular}, l'idéal
$\Ker(A \to B)$ peut être engendré par $\dim(A) - \dim(B)$ éléments.
Nous voyons donc que (4) implique (2).

\medskip\noindent
Il reste à montrer que (1) implique (4). Raisonnons par récurrence sur
$\dim(A) - \dim(B)$. Le cas $\dim(A) - \dim(B) = 0$ est trivial.
Supposons $\dim(A) > \dim (B)$.
Écrivons $I = \Ker(A \to C)$ et $J = \Ker(A \to B)$.
Notons que $J \subset I$. Notre hypothèse dit que le nombre minimal de
générateurs de $I$ est $\dim(A) - \dim(C)$.
Soit $\mathfrak m \subset A$ l'idéal maximal. Considérons les applications
$$
J/ \mathfrak m J \to  I / \mathfrak m I \to \mathfrak m /\mathfrak m^2
$$
D'après le Lemme \ref{algebra-lemma-regular-quotient-regular} et sa démonstration,
la composée est injective. Prenons un élément $x \in J$ qui ne soit pas nul
dans $J /\mathfrak mJ$. D'après ce qui précède et le lemme de Nakayama,
$x$ appartient à un système minimal de générateurs de $I$.
Nous pouvons donc remplacer $A$ par $A/xA$ et $I$ par $I/xA$, ce qui diminue
à la fois $\dim(A)$ et le nombre minimal de générateurs de $I$ de $1$.
Le résultat en découle.
\end{proof}

\begin{lemma}
\label{algebra-lemma-lci-local}
Soit $k$ un corps. Soit $S$ une $k$-algèbre locale essentiellement de type
fini sur $k$. Les assertions suivantes sont équivalentes :
\begin{enumerate}
\item $S$ est une intersection complète sur $k$,
\item pour toute surjection $R \to S$, où $R$ est un anneau local régulier
essentiellement de présentation finie sur $k$, l'idéal $\Ker(R \to S)$ peut
être engendré par une suite régulière,
\item pour une certaine surjection $R \to S$, où $R$ est un anneau local
régulier essentiellement de présentation finie sur $k$, l'idéal
$\Ker(R \to S)$ peut être engendré par $\dim(R) - \dim(S)$ éléments,
\item il existe une intersection complète globale $A$ sur $k$ et un idéal
premier $\mathfrak a$ de $A$ tels que $S \cong A_{\mathfrak a}$, et
\item il existe une intersection complète locale $A$ sur $k$ et un idéal
premier $\mathfrak a$ de $A$ tels que $S \cong A_{\mathfrak a}$.
\end{enumerate}
\end{lemma}

\begin{proof}
Il est clair que (2) implique (1), que (1) implique (3), et aussi que (4)
implique (5). Montrons que (3) implique (4). Nous supposons donc qu'il existe
une surjection $R \to S$, où $R$ est un anneau local régulier essentiellement
de présentation finie sur $k$, telle que l'idéal $\Ker(R \to S)$ puisse être
engendré par $\dim(R) - \dim(S)$ éléments.
Nous pouvons écrire $R = (k[x_1, \ldots, x_n]/J)_{\mathfrak q}$ pour un
certain $J \subset k[x_1, \ldots, x_n]$ et un certain idéal premier
$\mathfrak q \subset k[x_1, \ldots, x_n]$ tel que $J \subset \mathfrak q$.
Soit $I \subset k[x_1, \ldots, x_n]$ le noyau du morphisme
$k[x_1, \ldots, x_n] \to S$, de sorte que
$S \cong (k[x_1, \ldots, x_n]/I)_{\mathfrak q}$.
Par hypothèse, $(I/J)_{\mathfrak q}$ est engendré par
$\dim(R) - \dim(S)$ éléments. D'après le
Lemme \ref{algebra-lemma-ci-well-defined}, nous concluons que
$I_{\mathfrak q}$ peut être engendré par
$\dim(k[x_1, \ldots, x_n]_{\mathfrak q}) - \dim(S)$ éléments.
Le Lemme \ref{algebra-lemma-lci} montre que, pour un certain
$g \in k[x_1, \ldots, x_n]$, $g \not \in \mathfrak q$, l'algèbre
$(k[x_1, \ldots, x_n]/I)_g$ est une intersection complète globale et que
$S$ est isomorphe à l'un de ses anneaux locaux.

\medskip\noindent
Pour achever la démonstration, nous devons montrer que (5) implique (2).
Supposons (5) et soit $\pi : R \to S$ une surjection, où $R$ est une
$k$-algèbre locale régulière essentiellement de type fini sur $k$.
Par hypothèse, $S = A_{\mathfrak a}$ pour une certaine intersection complète
locale $A$ sur $k$.
Choisissons une présentation
$R = (k[y_1, \ldots, y_m]/J)_{\mathfrak q}$ avec
$J \subset \mathfrak q \subset k[y_1, \ldots, y_m]$.
Nous pouvons supposer, et supposons, que $J$ est le noyau du morphisme
$k[y_1, \ldots, y_m] \to R$. Soit
$I \subset k[y_1, \ldots, y_m]$ le noyau du morphisme
$k[y_1, \ldots, y_m] \to S = A_{\mathfrak a}$.
Alors $J \subset I$ et $(I/J)_{\mathfrak q}$ est le noyau de la surjection
$\pi : R \to S$. Ainsi,
$S = (k[y_1, \ldots, y_m]/I)_{\mathfrak q}$.

\medskip\noindent
D'après le Lemme \ref{algebra-lemma-isomorphic-local-rings}, il existe
$g \in A$, $g \not \in \mathfrak a$, et
$g' \in k[y_1, \ldots, y_m]$, $g' \not \in \mathfrak q$, tels que
$A_g \cong (k[y_1, \ldots, y_m]/I)_{g'}$.
Après avoir remplacé $A$ par $A_g$ et $k[y_1, \ldots, y_m]$ par
$k[y_1, \ldots, y_{m + 1}]$, nous pouvons supposer que
$A \cong k[y_1, \ldots, y_m]/I$. Considérons les morphismes surjectifs
d'anneaux locaux
$$
k[y_1, \ldots, y_m]_{\mathfrak q} \to R \to S.
$$
Nous devons montrer que le noyau de $R \to S$ est engendré par une suite
régulière. D'après le Lemme \ref{algebra-lemma-lci},
$k[y_1, \ldots, y_m]_{\mathfrak q} \to A_{\mathfrak a} = S$ possède cette
propriété (car $A$ est une intersection complète locale sur $k$).
Le résultat découle du Lemme \ref{algebra-lemma-ci-well-defined}.
\end{proof}

\begin{lemma}
\label{algebra-lemma-lci-at-prime}
Soit $k$ un corps. Soit $S$ une $k$-algèbre de type fini.
Soit $\mathfrak q$ un idéal premier de $S$. Les assertions suivantes sont
équivalentes :
\begin{enumerate}
\item L'anneau local $S_{\mathfrak q}$ est un anneau d'intersection complète
(Définition \ref{algebra-definition-lci-local-ring}).
\item Il existe $g \in S$, $g \not \in \mathfrak q$, tel que $S_g$ soit une
intersection complète locale sur $k$.
\item Il existe $g \in S$, $g \not \in \mathfrak q$, tel que $S_g$ soit une
intersection complète globale sur $k$.
\item Pour toute présentation $S = k[x_1, \ldots, x_n]/I$, où
$\mathfrak q' \subset k[x_1, \ldots, x_n]$ correspond à $\mathfrak q$,
l'une quelconque des conditions équivalentes (1) -- (5) du
Lemme \ref{algebra-lemma-lci} est satisfaite.
\end{enumerate}
\end{lemma}

\begin{proof}
C'est une combinaison des Lemmes \ref{algebra-lemma-lci} et
\ref{algebra-lemma-lci-local}, ainsi que des définitions.
\end{proof}

\begin{lemma}
\label{algebra-lemma-lci-global}
Soit $k$ un corps. Soit $S$ une $k$-algèbre de type fini.
Les assertions suivantes sont équivalentes :
\begin{enumerate}
\item L'anneau $S$ est une intersection complète locale sur $k$.
\item Tous les anneaux locaux de $S$ sont des anneaux d'intersection
complète sur $k$.
\item Toutes les localisations de $S$ en des idéaux maximaux sont des
anneaux d'intersection complète sur $k$.
\end{enumerate}
\end{lemma}

\begin{proof}
Cela découle du Lemme \ref{algebra-lemma-lci-at-prime}, du fait que $\Spec(S)$ est
quasi-compact et des définitions.
\end{proof}

\noindent
Le lemme suivant affirme que la propriété d'être une intersection complète
est préservée par changement du corps de base (en un sens fort).

\begin{lemma}
\label{algebra-lemma-lci-field-change-local}
Soit $K/k$ une extension de corps.
Soit $S$ une algèbre de type fini sur $k$.
Soit $\mathfrak q_K$ un idéal premier de $S_K = K \otimes_k S$, et soit
$\mathfrak q$ l'idéal premier correspondant de $S$.
Alors $S_{\mathfrak q}$ est une intersection complète sur $k$
(Définition \ref{algebra-definition-lci-local-ring}) si et seulement si
$(S_K)_{\mathfrak q_K}$ est une intersection complète sur $K$.
\end{lemma}

\begin{proof}
Choisissons une présentation $S = k[x_1, \ldots, x_n]/I$.
Elle donne une présentation $S_K = K[x_1, \ldots, x_n]/I_K$, où
$I_K = K \otimes_k I$. Soient
$\mathfrak q_K' \subset K[x_1, \ldots, x_n]$, respectivement
$\mathfrak q' \subset k[x_1, \ldots, x_n]$, les idéaux premiers
correspondants. Nous allons montrer que les conditions équivalentes du
Lemme \ref{algebra-lemma-lci} sont satisfaites pour le couple
$(S = k[x_1, \ldots, x_n]/I, \mathfrak q)$ si et seulement si elles le sont
pour le couple $(S_K = K[x_1, \ldots, x_n]/I_K, \mathfrak q_K)$.
Le lemme en découlera (voir le Lemme \ref{algebra-lemma-lci-at-prime}).

\medskip\noindent
D'après le Lemme \ref{algebra-lemma-dimension-at-a-point-preserved-field-extension},
on a $\dim_{\mathfrak q} S = \dim_{\mathfrak q_K} S_K$.
Ainsi, l'entier $c$ qui intervient dans le Lemme \ref{algebra-lemma-lci} est le même
pour le couple $(S = k[x_1, \ldots, x_n]/I, \mathfrak q)$ que pour le couple
$(S_K = K[x_1, \ldots, x_n]/I_K, \mathfrak q_K)$.
D'autre part, on a
\begin{eqnarray*}
I \otimes_{k[x_1, \ldots, x_n]} \kappa(\mathfrak q')
\otimes_{\kappa(\mathfrak q')} \kappa(\mathfrak q_K')
& = &
I \otimes_{k[x_1, \ldots, x_n]} \kappa(\mathfrak q_K') \\
& = &
I \otimes_{k[x_1, \ldots, x_n]} K[x_1, \ldots, x_n]
\otimes_{K[x_1, \ldots, x_n]} \kappa(\mathfrak q_K') \\
& = &
(K \otimes_k I) \otimes_{K[x_1, \ldots, x_n]} \kappa(\mathfrak q_K') \\
& = &
I_K \otimes_{K[x_1, \ldots, x_n]} \kappa(\mathfrak q'_K).
\end{eqnarray*}
Par conséquent,
$\dim_{\kappa(\mathfrak q')}
I \otimes_{k[x_1, \ldots, x_n]} \kappa(\mathfrak q') =
\dim_{\kappa(\mathfrak q'_K)}
I_K \otimes_{K[x_1, \ldots, x_n]} \kappa(\mathfrak q_K')$.
Il découle donc du Lemme de Nakayama \ref{algebra-lemma-NAK} que le nombre minimal
de générateurs de $I_{\mathfrak q'}$ est égal au nombre minimal de
générateurs de $(I_K)_{\mathfrak q'_K}$.
Le lemme découle alors de la caractérisation (2) du
Lemme \ref{algebra-lemma-lci}.
\end{proof}





\begin{lemma}
\label{algebra-lemma-lci-field-change}
Soit $k \to K$ une extension de corps.
Soit $S$ une $k$-algèbre de type fini.
Alors $S$ est une intersection complète locale sur $k$ si et seulement si
$S \otimes_k K$ est une intersection complète locale sur $K$.
\end{lemma}

\begin{proof}
Cela résulte de la combinaison des Lemmes
\ref{algebra-lemma-lci-global} et \ref{algebra-lemma-lci-field-change-local}.
Mais nous donnons également ici une autre démonstration
(fondée sur les mêmes principes).

\medskip\noindent
Posons $S' = S \otimes_k K$. Soit
$\alpha : k[x_1, \ldots, x_n] \to S$ une présentation de noyau $I$.
Soit $\alpha' : K[x_1, \ldots, x_n] \to S'$ la présentation induite,
de noyau $I'$.

\medskip\noindent
Supposons que $S$ soit une intersection complète locale.
Choisissons un idéal premier $\mathfrak q \subset S'$. Notons
$\mathfrak q'$ l'idéal premier correspondant de $K[x_1, \ldots, x_n]$,
$\mathfrak p$ l'idéal premier correspondant de $S$, et
$\mathfrak p'$ l'idéal premier correspondant de $k[x_1, \ldots, x_n]$.
Considérons le diagramme suivant d'anneaux locaux noethériens :
$$
\xymatrix{
S'_{\mathfrak q} &  K[x_1, \ldots, x_n]_{\mathfrak q'} \ar[l] \\
S_{\mathfrak p}\ar[u] &  k[x_1, \ldots, x_n]_{\mathfrak p'} \ar[u] \ar[l]
}
$$
D'après le Lemme \ref{algebra-lemma-lci}, nous savons que $S_{\mathfrak p}$
est défini par une certaine suite régulière $f_1, \ldots, f_c$ dans
$k[x_1, \ldots, x_n]_{\mathfrak p'}$. Puisque la flèche verticale de
droite est plate, nous voyons que les images de $f_1, \ldots, f_c$
forment une suite régulière dans $K[x_1, \ldots, x_n]_{\mathfrak q'}$.
Comme le produit tensoriel par $K$ sur $k$ est un foncteur exact, on a
$S'_{\mathfrak q} = K[x_1, \ldots, x_n]_{\mathfrak q'}/(f_1, \ldots, f_c)$.
En appliquant à nouveau le Lemme \ref{algebra-lemma-lci}, nous voyons donc que $S'$
est une intersection complète locale dans un voisinage de $\mathfrak q$.
Comme $\mathfrak q$ était arbitraire, nous voyons que $S'$ est une
intersection complète locale sur $K$.

\medskip\noindent
Supposons que $S'$ soit une intersection complète locale.
Choisissons un idéal maximal $\mathfrak m$ de $S$. Notons $\mathfrak m'$
l'idéal maximal correspondant de $k[x_1, \ldots, x_n]$.
Notons $\kappa = \kappa(\mathfrak m)$ le corps résiduel.
D'après la Remarque \ref{algebra-remark-fundamental-diagram}, les idéaux premiers de
$S'$ situés au-dessus de $\mathfrak m$ correspondent aux idéaux premiers de
$K \otimes_k \kappa$. D'après le théorème de Hilbert-Nullstellensatz
\ref{algebra-theorem-nullstellensatz}, on a $[\kappa : k] < \infty$.
Ainsi, $K \otimes_k \kappa$ est non nul et de dimension finie sur $K$.
$K \otimes_k \kappa$ possède donc un nombre fini $> 0$ d'idéaux premiers,
qui sont tous maximaux et dont chacun a un corps résiduel fini sur $K$
(voir la section \ref{algebra-section-artinian}).
Il existe donc un nombre fini $> 0$ d'idéaux premiers
$\mathfrak n \subset S'$ situés au-dessus de $\mathfrak m$ ; chacun d'eux
est maximal et possède un corps résiduel fini sur $K$. Choisissons-en un,
disons $\mathfrak n \subset S'$, et soit
$\mathfrak n' \subset K[x_1, \ldots, x_n]$ l'idéal premier correspondant de
$K[x_1, \ldots, x_n]$.
Notons que, puisque $V(\mathfrak mS')$ est fini, $\mathfrak n$ en est un
point fermé isolé, et nous en déduisons que $\mathfrak mS'_{\mathfrak n}$
est un idéal de définition de $S'_{\mathfrak n}$. Cela implique que
$\dim(S_{\mathfrak m}) = \dim(S'_{\mathfrak n})$, par exemple d'après le
Lemme \ref{algebra-lemma-dimension-base-fibre-equals-total}.
(On peut aussi le voir à l'aide du
Lemme \ref{algebra-lemma-dimension-at-a-point-preserved-field-extension}.)
Considérons le diagramme correspondant d'anneaux locaux noethériens :
$$
\xymatrix{
S'_{\mathfrak n} &  K[x_1, \ldots, x_n]_{\mathfrak n'} \ar[l] \\
S_{\mathfrak m}\ar[u] &  k[x_1, \ldots, x_n]_{\mathfrak m'} \ar[u] \ar[l]
}
$$
D'après le Lemme \ref{algebra-lemma-change-base-NL}, on a
$\NL(\alpha) \otimes_S S' = \NL(\alpha')$, et en particulier
$I'/(I')^2 = I/I^2 \otimes_S S'$. Par conséquent,
$(I/I^2)_{\mathfrak m} \otimes_{S_{\mathfrak m}} \kappa$
et
$(I'/(I')^2)_{\mathfrak n} \otimes_{S'_{\mathfrak n}} \kappa(\mathfrak n)$
ont la même dimension. Puisque $(I'/(I')^2)_{\mathfrak n}$
est libre de rang $n - \dim S'_{\mathfrak n}$, nous en déduisons que
$(I/I^2)_{\mathfrak m}$ peut être engendré par
$n - \dim S'_{\mathfrak n} = n - \dim S_{\mathfrak m}$ éléments.
D'après le Lemme \ref{algebra-lemma-lci}, nous voyons que $S$ est une intersection
complète locale dans un voisinage de $\mathfrak m$.
Comme $\mathfrak m$ était un idéal maximal quelconque, nous concluons que
$S$ est une intersection complète locale.
\end{proof}

\noindent
Nous terminons par un lemme que nous utiliserons plus tard pour montrer que,
étant donnés des morphismes d'anneaux $T \to A \to B$ tels que $B$ soit
syntomique sur $T$ et que $B$ soit syntomique sur $A$, alors $A$ est
syntomique sur $T$.

\begin{lemma}
\label{algebra-lemma-lci-permanence-initial}
Soit
$$
\xymatrix{
B & S \ar[l] \\
A \ar[u] & R \ar[l] \ar[u]
}
$$
un carré commutatif d'anneaux locaux. Supposons que
\begin{enumerate}
\item $R$ et $\overline{S} = S/\mathfrak m_R S$ soient des anneaux locaux
réguliers,
\item $A = R/I$ et $B = S/J$ pour certains idéaux $I$, $J$,
\item $J \subset S$ et
$\overline{J} = J/\mathfrak m_R \cap J \subset \overline{S}$
soient engendrés par des suites régulières, et
\item $A \to B$ et $R \to S$ soient plats.
\end{enumerate}
Alors $I$ est engendré par une suite régulière.
\end{lemma}

\begin{proof}
Posons $\overline{B} = B/\mathfrak m_RB = B/\mathfrak m_AB$, de sorte que
$\overline{B} = \overline{S}/\overline{J}$.
Soient $f_1, \ldots, f_{\overline{c}} \in J$ des éléments tels que
$\overline{f}_1, \ldots, \overline{f}_{\overline{c}} \in \overline{J}$
forment une suite régulière engendrant $\overline{J}$.
Notons que $\overline{c} = \dim(\overline{S}) - \dim(\overline{B})$ ;
voir le Lemme \ref{algebra-lemma-ci-well-defined}.
D'après le Lemme \ref{algebra-lemma-grothendieck-regular-sequence}, l'anneau
$S/(f_1, \ldots, f_{\overline{c}})$ est plat sur $R$.
Ainsi, $S/(f_1, \ldots, f_{\overline{c}}) + IS$ est plat sur $A$.
Le morphisme $S/(f_1, \ldots, f_{\overline{c}}) + IS \to B$ est donc une
surjection de $S/IS$-modules de type fini plats sur $A$, qui est un
isomorphisme modulo $\mathfrak m_A$ et, par conséquent, un isomorphisme
d'après le Lemme \ref{algebra-lemma-mod-injective}. Autrement dit,
$J = (f_1, \ldots, f_{\overline{c}}) + IS$.

\medskip\noindent
D'après le Lemme \ref{algebra-lemma-ci-well-defined}, à nouveau, l'idéal $J$ est
engendré par une suite régulière de $c = \dim(S) - \dim(B)$ éléments.
Ainsi, $J/\mathfrak m_SJ$ est un espace vectoriel de dimension $c$.
D'après la description de $J$ ci-dessus, il existe
$g_1, \ldots, g_{c - \overline{c}} \in I$ tels que $J$ soit engendré par
$f_1, \ldots, f_{\overline{c}}, g_1, \ldots, g_{c - \overline{c}}$
(utiliser le Lemme de Nakayama \ref{algebra-lemma-NAK}). Considérons l'anneau
$A' = R/(g_1, \ldots, g_{c - \overline{c}})$ et la surjection
$A' \to A$. Nous voyons d'après ce qui précède que
$B = S/(f_1, \ldots, f_{\overline{c}}, g_1, \ldots, g_{c - \overline{c}})$
est plat sur $A'$ (puisque $S/(f_1, \ldots, f_{\overline{c}})$ est plat
sur $R$). Ainsi, $A' \to B$ est injectif (car il est fidèlement plat ;
voir le Lemme \ref{algebra-lemma-local-flat-ff}).
Comme ce morphisme se factorise par $A$, on obtient $A' = A$.
Notons que $\dim(B) = \dim(A) + \dim(\overline{B})$, et que
$\dim(S) = \dim(R) + \dim(\overline{S})$ ; voir le
Lemme \ref{algebra-lemma-dimension-base-fibre-equals-total}.
Ainsi, $c - \overline{c} = \dim(R) -\dim(A)$ par un calcul élémentaire.
Par conséquent, $I = (g_1, \ldots, g_{c - \overline{c}})$ est engendré
par une suite régulière d'après le Lemme \ref{algebra-lemma-ci-well-defined}.
\end{proof}




\section{Morphismes syntomiques}
\label{algebra-section-syntomic}

\noindent
Les morphismes d'anneaux syntomiques sont les morphismes d'anneaux plats de
présentation finie dont toutes les fibres sont des intersections complètes
locales. Nous étudions les morphismes généraux d'intersection complète locale
dans Compléments d'algèbre,
section \ref{more-algebra-section-lci}.

\begin{definition}
\label{algebra-definition-lci}
Un morphisme d'anneaux $R \to S$ est dit {\it syntomique}, ou bien nous disons
que $S$ est une {\it intersection complète locale plate sur $R$}, s'il est
plat, de présentation finie, et si tous ses anneaux fibres
$S \otimes_R \kappa(\mathfrak p)$ sont des intersections complètes locales ;
voir la Définition \ref{algebra-definition-lci-field}.
\end{definition}

\noindent
Il est clair qu'une algèbre sur un corps est syntomique sur ce corps si et
seulement si elle est une intersection complète locale. Voici une propriété
agréable de cette définition.

\begin{lemma}
\label{algebra-lemma-syntomic-descends}
\begin{slogan}
La propriété d'être syntomique est locale pour la topologie fpqc sur la base.
\end{slogan}
Soit $R \to S$ un morphisme d'anneaux.
Soit $R \to R'$ un morphisme d'anneaux fidèlement plat.
Posons $S' = R'\otimes_R S$.
Alors $R \to S$ est syntomique si et seulement si $R' \to S'$ est syntomique.
\end{lemma}

\begin{proof}
D'après le Lemme \ref{algebra-lemma-finite-presentation-descends} et le
Lemme \ref{algebra-lemma-flatness-descends}, cela vaut pour la propriété d'être plat
et pour celle d'être de présentation finie.
Le morphisme $\Spec(R') \to \Spec(R)$ est surjectif ;
voir le Lemme \ref{algebra-lemma-ff-rings}. Il suffit donc de montrer, pour des idéaux
premiers $\mathfrak p' \subset R'$ situés au-dessus de
$\mathfrak p \subset R$, que $S \otimes_R \kappa(\mathfrak p)$ est une
intersection complète locale si et seulement si
$S' \otimes_{R'} \kappa(\mathfrak p')$ est une intersection complète locale.
Notons que
$S' \otimes_{R'} \kappa(\mathfrak p') =
S \otimes_R \kappa(\mathfrak p)
\otimes_{\kappa(\mathfrak p)} \kappa(\mathfrak p')$.
Le Lemme \ref{algebra-lemma-lci-field-change} s'applique donc.
\end{proof}

\begin{lemma}
\label{algebra-lemma-base-change-syntomic}
Tout changement de base d'un morphisme syntomique est syntomique.
\end{lemma}

\begin{proof}
Cela est vrai pour la platitude, pour la propriété d'être de présentation
finie et pour celle d'avoir des intersections complètes locales comme fibres,
d'après les Lemmes \ref{algebra-lemma-flat-base-change},
\ref{algebra-lemma-compose-finite-type} et \ref{algebra-lemma-lci-field-change}.
\end{proof}

\begin{lemma}
\label{algebra-lemma-local-syntomic}
Soit $R \to S$ un morphisme d'anneaux.
Supposons donnés des éléments $g_1, \ldots g_m \in S$ qui engendrent l'idéal
unité et tels que chacun des morphismes $R \to S_{g_i}$ soit syntomique.
Alors $R \to S$ est syntomique.
\end{lemma}

\begin{proof}
Cela est vrai pour la platitude et pour la propriété d'être de présentation
finie d'après les Lemmes \ref{algebra-lemma-flat-localization} et
\ref{algebra-lemma-cover-upstairs}.
La propriété d'avoir des anneaux fibres qui soient des intersections complètes
locales est locale sur $S$ par sa définition même ; voir la
Définition \ref{algebra-definition-lci-field}.
\end{proof}

\begin{definition}
\label{algebra-definition-relative-global-complete-intersection}
Soit $R \to S$ un morphisme d'anneaux. Nous disons que $R \to S$ est une
{\it intersection complète globale relative} s'il existe une présentation
$S = R[x_1, \ldots, x_n]/(f_1, \ldots, f_c)$ et si toute fibre non vide de
$\Spec(S) \to \Spec(R)$ est de dimension $n - c$.
Nous dirons ``soit
$S = R[x_1, \ldots, x_n]/(f_1, \ldots, f_c)$ une intersection complète
globale relative'' pour désigner cette situation.
\end{definition}

\noindent
Le lemme suivant est parfois utile pour trouver des présentations globales.

\begin{lemma}
\label{algebra-lemma-huber}
Soit $S$ une $R$-algèbre de présentation finie qui possède une présentation
$S = R[x_1, \ldots, x_n]/I$ telle que $I/I^2$ soit libre sur $S$. Alors
$S$ possède une présentation $S = R[y_1, \ldots, y_m]/(f_1, \ldots, f_c)$
telle que $(f_1, \ldots, f_c)/(f_1, \ldots, f_c)^2$ soit libre, avec pour
base les classes de $f_1, \ldots, f_c$.
\end{lemma}

\begin{proof}
Notons que $I$ est un idéal de type fini d'après le
Lemme \ref{algebra-lemma-finite-presentation-independent}.
Soient $f_1, \ldots, f_c \in I$ des éléments qui s'envoient sur une base de
$I/I^2$. D'après le lemme de Nakayama (Lemme \ref{algebra-lemma-NAK}),
il existe $g \in 1 + I$ tel que
$$
g \cdot I \subset (f_1, \ldots, f_c)
$$
et $I_g \cong (f_1, \ldots, f_c)_g$. Nous voyons donc que
$$
S \cong R[x_1, \ldots, x_n]/(f_1, \ldots, f_c)[1/g]
\cong R[x_1, \ldots, x_n, x_{n + 1}]/(f_1, \ldots, f_c, gx_{n + 1} - 1)
$$
comme souhaité. Il en résulte que
$f_1, \ldots, f_c,gx_{n + 1} - 1$ forment une base de
$(f_1, \ldots, f_c, gx_{n + 1} - 1)/(f_1, \ldots, f_c, gx_{n + 1} - 1)^2$,
par exemple en appliquant le
Lemme \ref{algebra-lemma-principal-localization-NL}.
\end{proof}



\begin{example}
\label{algebra-example-factor-polynomials}
Soient $n , m \geq 1$ des entiers. Considérons le morphisme d'anneaux
\begin{eqnarray*}
R = \mathbf{Z}[a_1, \ldots, a_{n + m}]
& \longrightarrow &
S = \mathbf{Z}[b_1, \ldots, b_n, c_1, \ldots, c_m] \\
a_1 & \longmapsto & b_1 + c_1 \\
a_2 & \longmapsto & b_2 + b_1 c_1 + c_2 \\
\ldots & \ldots & \ldots \\
a_{n + m} & \longmapsto & b_n c_m
\end{eqnarray*}
Autrement dit, c'est l'unique morphisme d'anneaux de polynômes indiqué tel
que la factorisation de polynômes
$$
x^{n + m} + a_1 x^{n + m - 1} + \ldots + a_{n + m}
=
(x^n + b_1 x^{n - 1} + \ldots + b_n)
(x^m + c_1 x^{m - 1} + \ldots + c_m)
$$
soit satisfaite. Notons que $S$ est engendré par $n + m$ éléments sur $R$
(à savoir les $b_i, c_j$) et qu'il y a $n + m$ équations
(à savoir $a_k = a_k(b_i, c_j)$). Pour montrer que $S$ est une intersection
complète globale relative sur $R$, il suffit de prouver que toutes les fibres
sont de dimension $0$.

\medskip\noindent
Pour cela, soit $R \to k$ un morphisme d'anneaux à valeurs dans un corps $k$.
Disons que $a_i$ s'envoie sur $\alpha_i \in k$.
Considérons l'anneau fibre $S_k = k \otimes_R S$. Soit $k \to K$ une
extension de corps. Se donner un morphisme de $k$-algèbres de $S_k \to K$
revient à trouver
$\beta_1, \ldots, \beta_n, \gamma_1, \ldots, \gamma_m \in K$ tels que
$$
x^{n + m} + \alpha_1 x^{n + m - 1} + \ldots + \alpha_{n + m}
=
(x^n + \beta_1 x^{n - 1} + \ldots + \beta_n)
(x^m + \gamma_1 x^{m - 1} + \ldots + \gamma_m).
$$
Il n'existe donc qu'un nombre fini de tels $n + m$-uplets dans $K$.
Cela prouve que toutes les fibres n'ont qu'un nombre fini de points fermés
(utiliser par exemple le théorème de Hilbert-Nullstellensatz pour voir qu'ils
correspondent tous à des solutions dans $\overline{k}$), et donc que
$R \to S$ est une intersection complète globale relative.

\medskip\noindent
Une autre manière de raisonner consiste à montrer que
$\mathbf{Z}[a_1, \ldots, a_{n + m}] \to
\mathbf{Z}[b_1, \ldots, b_n, c_1, \ldots, c_m]$ est en fait aussi un
morphisme d'anneaux {\it fini}. En effet, d'après le
Lemme \ref{algebra-lemma-polynomials-divide}, chacun des $b_i, c_j$ est entier sur
$R$, et donc $R \to S$ est fini d'après le
Lemme \ref{algebra-lemma-characterize-integral}.
\end{example}

\begin{example}
\label{algebra-example-roots-universal-polynomial}
Considérons le morphisme d'anneaux
\begin{eqnarray*}
R = \mathbf{Z}[a_1, \ldots, a_n]
& \longrightarrow &
S = \mathbf{Z}[\alpha_1, \ldots, \alpha_n] \\
a_1 & \longmapsto &
\alpha_1 + \ldots + \alpha_n \\
\ldots & \ldots & \ldots \\
a_n & \longmapsto & \alpha_1 \ldots \alpha_n
\end{eqnarray*}
Autrement dit, c'est l'unique morphisme d'anneaux de polynômes indiqué tel que
$$
x^n + a_1 x^{n - 1} + \ldots + a_n
=
\prod\nolimits_{i = 1}^n (x + \alpha_i)
$$
soit satisfait dans $\mathbf{Z}[\alpha_i, x]$. On peut aussi dire que $a_i$
s'envoie sur la $i$-ième fonction symétrique élémentaire en
$\alpha_1, \ldots, \alpha_n$. D'après la théorie usuelle des polynômes
symétriques élémentaires (détails omis), l'anneau $S$ est un $R$-module libre
de type fini dont une base est constituée des éléments
$\alpha_1^{e_1} \alpha_2^{e_2} \ldots \alpha_n^{e_n}$ avec
$0 \leq e_i \leq n - i$. A fortiori, les anneaux fibres de
$R \to S$ sont finis et sont donc de dimension $0$.
D'autre part, $S$ est engendré par $n$ éléments sur $R$ soumis à $n$
équations. Ainsi, $S$ est une intersection complète globale relative sur
$R$. Comme le rang de $S$ sur $R$ est positif, nous voyons aussi que $S$ est
fidèlement plat sur $R$.
\end{example}

\begin{lemma}
\label{algebra-lemma-base-change-relative-global-complete-intersection}
Soit $S = R[x_1, \ldots, x_n]/(f_1, \ldots, f_c)$ une intersection complète
globale relative
(Définition \ref{algebra-definition-relative-global-complete-intersection}).
\begin{enumerate}
\item Pour tout $R \to R'$, le changement de base
$R' \otimes_R S = R'[x_1, \ldots, x_n]/(f_1, \ldots, f_c)$ est une
intersection complète globale relative.
\item Pour tout $g \in S$ qui est l'image d'un
$h \in R[x_1, \ldots, x_n]$, l'anneau
$S_g = R[x_1, \ldots, x_n, x_{n + 1}]/(f_1, \ldots, f_c, hx_{n + 1} - 1)$
est une intersection complète globale relative.
\item Si $R \to S$ se factorise sous la forme $R \to R_f \to S$ pour un
certain $f \in R$, alors l'anneau
$S = R_f[x_1, \ldots, x_n]/(f_1, \ldots, f_c)$ est une intersection complète
globale relative sur $R_f$.
\end{enumerate}
\end{lemma}

\begin{proof}
D'après le Lemme \ref{algebra-lemma-dimension-preserved-field-extension}, les fibres
d'un changement de base ont la même dimension que les fibres du morphisme
initial. De plus,
$R' \otimes_R R[x_1, \ldots, x_n]/(f_1, \ldots, f_c)
= R'[x_1, \ldots, x_n]/(f_1, \ldots, f_c)$. Ainsi, (1) en résulte.
La démonstration de (2) repose sur le fait que la localisation en un élément
peut s'écrire
$S_g \cong S[x_{n + 1}]/(gx_{n + 1} - 1)$.
L'assertion (3) découle de (1), car, sous les hypothèses de (3), on a
$R_f \otimes_R S \cong S$.
\end{proof}

\begin{lemma}
\label{algebra-lemma-localize-relative-complete-intersection}
Soit $R$ un anneau. Soit
$S = R[x_1, \ldots, x_n]/(f_1, \ldots, f_c)$.
Dans chacun des cas suivants, nous allons trouver
$h \in R[x_1, \ldots, x_n]$ qui s'envoie sur $g \in S$ tel que
$$
S_g = R[x_1, \ldots, x_n, x_{n + 1}]/(f_1, \ldots, f_c, hx_{n + 1} - 1)
$$
soit une intersection complète globale relative avec une présentation comme
dans la Définition \ref{algebra-definition-relative-global-complete-intersection} :
\begin{enumerate}
\item Soit $I \subset R$ un idéal. Si les fibres de
$\Spec(S/IS) \to \Spec(R/I)$ sont de dimension $n - c$, alors nous pouvons
trouver $(h, g)$ comme ci-dessus de sorte que $g$ s'envoie sur $1 \in S/IS$.
\item Soit $\mathfrak p \subset R$ un idéal premier. Si
$\dim(S \otimes_R \kappa(\mathfrak p)) = n - c$, alors nous pouvons trouver
$(h, g)$ comme ci-dessus de sorte que $g$ s'envoie sur une unité de
$S \otimes_R \kappa(\mathfrak p)$.
\item Soit $\mathfrak q \subset S$ un idéal premier situé au-dessus de
$\mathfrak p \subset R$. Si $\dim_{\mathfrak q}(S/R) = n - c$, alors nous
pouvons trouver $(h, g)$ comme ci-dessus de sorte que
$g \not \in \mathfrak q$.
\end{enumerate}
\end{lemma}

\begin{proof}
Pour (1). D'après le Lemme \ref{algebra-lemma-dimension-fibres-bounded-open-upstairs},
il existe un ouvert $W \subset \Spec(S)$ contenant $V(IS)$ tel que toutes les
fibres de $W \to \Spec(R)$ soient de dimension $\leq n - c$.
Écrivons $W = \Spec(S) \setminus V(J)$. Alors
$V(J) \cap V(IS) = \emptyset$ ; nous pouvons donc trouver $g \in J$ qui
s'envoie sur $1 \in S/IS$.
Soit $h \in R[x_1, \ldots, x_n]$ un antécédent quelconque de $g$.

\medskip\noindent
Pour (2). D'après le Lemme \ref{algebra-lemma-dimension-fibres-bounded-open-upstairs},
il existe un ouvert $W \subset \Spec(S)$ contenant
$\Spec(S \otimes_R \kappa(\mathfrak p))$
tel que toutes les fibres de $W \to \Spec(R)$ soient de dimension
$\leq n - c$.
Écrivons $W = \Spec(S) \setminus V(J)$. Alors
$V(J \cdot S \otimes_R \kappa(\mathfrak p)) = \emptyset$.
Nous pouvons donc trouver $g \in J$ qui s'envoie sur une unité de
$S \otimes_R \kappa(\mathfrak p)$ (détails omis).
Soit $h \in R[x_1, \ldots, x_n]$ un antécédent quelconque de $g$.

\medskip\noindent
Pour (3). D'après le Lemme
\ref{algebra-lemma-dimension-fibres-bounded-open-upstairs}, il existe
$g \in S$, $g \not \in \mathfrak q$, tel que toutes les fibres non vides de
$R \to S_g$ soient de dimension $\leq n - c$.
Soit $h \in R[x_1, \ldots, x_n]$ un élément quelconque qui s'envoie sur $g$.
\end{proof}

\noindent
Le lemme suivant affirme que nous pouvons effectuer une approximation
noethérienne absolue pour les intersections complètes globales relatives.

\begin{lemma}
\label{algebra-lemma-relative-global-complete-intersection-Noetherian}
Soit $R$ un anneau. Soit
$S = R[x_1, \ldots, x_n]/(f_1, \ldots, f_c)$ une intersection complète
globale relative
(Définition \ref{algebra-definition-relative-global-complete-intersection}).
Il existe une sous-$\mathbf{Z}$-algèbre $R_0 \subset R$ de type fini telle que
$f_i \in R_0[x_1, \ldots, x_n]$ et telle que
$$
S_0 = R_0[x_1, \ldots, x_n]/(f_1, \ldots, f_c)
$$
soit une intersection complète globale relative.
\end{lemma}

\begin{proof}
Soit $R_0 \subset R$ la $\mathbf{Z}$-algèbre de $R$ engendrée par tous les
coefficients des polynômes $f_1, \ldots, f_c$. Posons
$S_0 = R_0[x_1, \ldots, x_n]/(f_1, \ldots, f_c)$.
Il est clair que $S = R \otimes_{R_0} S_0$.
Choisissons un idéal premier $\mathfrak q \subset S$ et notons
$\mathfrak p \subset R$, $\mathfrak q_0 \subset S_0$ et
$\mathfrak p_0 \subset R_0$ les idéaux premiers au-dessus desquels il se
trouve. Puisque
$\dim (S \otimes_R \kappa(\mathfrak p) ) = n - c$,
on a aussi
$\dim (S_0 \otimes_{R_0} \kappa(\mathfrak p_0)) = n - c$ ;
voir le Lemme \ref{algebra-lemma-dimension-preserved-field-extension}.
D'après le Lemme \ref{algebra-lemma-dimension-fibres-bounded-open-upstairs}, il existe
$g \in S_0$, $g \not \in \mathfrak q_0$, tel que toutes les fibres non vides
de $R_0 \to (S_0)_g$ soient de dimension $\leq n - c$.
Comme $\mathfrak q$ était arbitraire et que $\Spec(S)$ est quasi-compact, nous
pouvons trouver un nombre fini d'éléments
$g_1, \ldots, g_m \in S_0$ tels que (a), pour $j = 1, \ldots, m$, les fibres
non vides de $R_0 \to (S_0)_{g_j}$ soient de dimension $\leq n - c$ et (b),
l'image de $\Spec(S) \to \Spec(S_0)$ soit contenue dans
$D(g_1) \cup \ldots \cup D(g_m)$.
Autrement dit, les images de $g_1, \ldots, g_m$ dans
$S = R \otimes_{R_0} S_0$ engendrent l'idéal unité. Après avoir agrandi $R_0$,
nous pouvons supposer que $g_1, \ldots, g_m$ engendrent l'idéal unité dans
$S_0$. D'après (a), toutes les fibres non vides de $R_0 \to S_0$ sont de
dimension $\leq n - c$, et nous concluons.
\end{proof}

\begin{lemma}
\label{algebra-lemma-relative-global-complete-intersection-conormal}
Soit $R$ un anneau. Soit
$S = R[x_1, \ldots, x_n]/(f_1, \ldots, f_c)$ une intersection complète
globale relative (Définition
\ref{algebra-definition-relative-global-complete-intersection}). Pour tout idéal
premier $\mathfrak q$ de $S$, soit $\mathfrak q'$ l'idéal premier
correspondant de $R[x_1, \ldots, x_n]$. Alors
\begin{enumerate}
\item $f_1, \ldots, f_c$ forment une suite régulière dans l'anneau local
$R[x_1, \ldots, x_n]_{\mathfrak q'}$,
\item chacun des anneaux
$R[x_1, \ldots, x_n]_{\mathfrak q'}/(f_1, \ldots, f_i)$ est plat sur $R$, et
\item le $S$-module $(f_1, \ldots, f_c)/(f_1, \ldots, f_c)^2$ est libre, avec
pour base les éléments $f_i \bmod (f_1, \ldots, f_c)^2$.
\end{enumerate}
\end{lemma}

\begin{proof}
Supposons $R$ noethérien. Soit $\mathfrak p = R \cap \mathfrak q'$.
D'après le Lemme \ref{algebra-lemma-lci}, par exemple, nous voyons que
$f_1, \ldots, f_c$ forment une suite régulière dans l'anneau local
$R[x_1, \ldots, x_n]_{\mathfrak q'} \otimes_R \kappa(\mathfrak p)$.
De plus, l'anneau local $R[x_1, \ldots, x_n]_{\mathfrak q'}$ est plat sur
$R_{\mathfrak p}$. Comme $R$, et donc
$R[x_1, \ldots, x_n]_{\mathfrak q'}$, est noethérien, nous déduisons du
Lemme \ref{algebra-lemma-grothendieck-regular-sequence} que (1) et (2) sont vraies.

\medskip\noindent
Soit maintenant $R$ quelconque. Écrivons
$R = \colim_{\lambda \in \Lambda} R_\lambda$ comme la colimite filtrante de
sous-$\mathbf{Z}$-algèbres de type fini (comparer avec la
section \ref{algebra-section-colimits-flat}). Nous pouvons supposer que
$f_1, \ldots, f_c \in R_\lambda[x_1, \ldots, x_n]$ pour tout $\lambda$.
Soit $R_0 \subset R$ comme dans le
Lemme \ref{algebra-lemma-relative-global-complete-intersection-Noetherian}.
Nous pouvons alors supposer que $R_0 \subset R_\lambda$ pour tout $\lambda$.
Il en résulte que
$S_\lambda = R_\lambda[x_1, \ldots, x_n]/(f_1, \ldots, f_c)$ est une
intersection complète globale relative (comme changement de base de $S_0$
par $R_0 \to R_\lambda$ ; voir le
Lemme \ref{algebra-lemma-base-change-relative-global-complete-intersection}).
Notons $\mathfrak p_\lambda$, $\mathfrak q_\lambda$,
$\mathfrak q'_\lambda$ les idéaux premiers de $R_\lambda$, $S_\lambda$,
$R_\lambda[x_1, \ldots, x_n]$ induits par $\mathfrak p$, $\mathfrak q$,
$\mathfrak q'$. Avec ces notations, (1) et (2) sont vraies pour tout
$\lambda$. Puisque
$$
R[x_1, \ldots, x_n]_{\mathfrak q'}/(f_1, \ldots, f_i)
=
\colim R_\lambda[x_1, \ldots, x_n]_{\mathfrak q_\lambda'}/(f_1, \ldots, f_i)
$$
nous déduisons la platitude dans (2) sur $R$ du
Lemme \ref{algebra-lemma-colimit-rings-flat}.
Comme on a
\begin{align*}
R[x_1, \ldots, x_n]_{\mathfrak q'}/(f_1, \ldots, f_i)
\xrightarrow{f_{i + 1}}
R[x_1, \ldots, x_n]_{\mathfrak q'}/(f_1, \ldots, f_i) \\
=
\colim
\left(
R_\lambda[x_1, \ldots, x_n]_{\mathfrak q_\lambda'}/(f_1, \ldots, f_i)
\xrightarrow{f_{i + 1}}
R_\lambda[x_1, \ldots, x_n]_{\mathfrak q_\lambda'}/(f_1, \ldots, f_i)
\right)
\end{align*}
et comme les colimites filtrantes sont exactes
(Lemme \ref{algebra-lemma-directed-colimit-exact}), nous concluons que (1) est vraie.

\medskip\noindent
Démonstration de (3). Notons $N$ le $S$-module
$(f_1, \ldots, f_c)/(f_1, \ldots, f_c)^2$ et $e_i \in N$ l'image de $f_i$.
D'après le Lemme \ref{algebra-lemma-regular-quasi-regular} et (1), nous savons que
$e_1, \ldots, e_c$ est une base de $N_\mathfrak q$ pour tout idéal premier
$\mathfrak q$ de $S$. Le Lemme \ref{algebra-lemma-characterize-zero-local} permet de
conclure que (3) est vraie.
\end{proof}

\begin{lemma}
\label{algebra-lemma-relative-global-complete-intersection}
Une intersection complète globale relative est syntomique, c'est-à-dire
plate.
\end{lemma}

\begin{proof}
Soit $R \to S$ une intersection complète globale relative.
Les fibres sont des intersections complètes globales, et $S$ est de
présentation finie sur $R$.
La seule chose qui reste à démontrer est donc que $R \to S$ est plat.
C'est vrai d'après (2) du
Lemme \ref{algebra-lemma-relative-global-complete-intersection-conormal}.
\end{proof}

\begin{lemma}
\label{algebra-lemma-adjoin-roots}
Supposons que $A$ soit un anneau et que
$P(x) = x^n + b_1 x^{n-1} + \ldots + b_n \in A[x]$ soit un polynôme unitaire
sur $A$. Alors il existe une extension d'anneaux syntomique, libre de type
fini et fidèlement plate $A \subset A'$ telle que
$P(x) = \prod_{i = 1, \ldots, n} (x - \beta_i)$ pour certains
$\beta_i \in A'$.
\end{lemma}

\begin{proof}
Prenons $A' = A \otimes_R S$, où $R$ et $S$ sont comme dans
l'Exemple \ref{algebra-example-roots-universal-polynomial}, où $R \to A$ envoie
$a_i$ sur $b_i$, et posons $\beta_i = -1 \otimes \alpha_i$.
Observons que $R \to S$ est fidèlement plat, libre de type fini et syntomique
d'après le Lemme \ref{algebra-lemma-relative-global-complete-intersection}.
Ces propriétés sont héritées par le changement de base $A \to A'$ ; certains
détails sont omis.
\end{proof}

\begin{lemma}
\label{algebra-lemma-syntomic}
Soit $R \to S$ un morphisme d'anneaux.
Soit $\mathfrak q \subset S$ un idéal premier situé au-dessus de l'idéal
premier $\mathfrak p$ de $R$.
Les assertions suivantes sont équivalentes :
\begin{enumerate}
\item Il existe un élément $g \in S$, $g \not \in \mathfrak q$ tel que
$R \to S_g$ soit syntomique.
\item Il existe un élément $g \in S$, $g \not \in \mathfrak q$ tel que
$S_g$ soit une intersection complète globale relative sur $R$.
\item Il existe un élément $g \in S$, $g \not \in \mathfrak q$, tel que
$R \to S_g$ soit de présentation finie, que le morphisme d'anneaux locaux
$R_{\mathfrak p} \to S_{\mathfrak q}$ soit plat et que l'anneau local
$S_{\mathfrak q}/\mathfrak pS_{\mathfrak q}$ soit un anneau d'intersection
complète sur $\kappa(\mathfrak p)$ (voir la
Définition \ref{algebra-definition-lci-local-ring}).
\end{enumerate}
\end{lemma}

\begin{proof}
L'implication (1) $\Rightarrow$ (3) est le
Lemme \ref{algebra-lemma-lci-at-prime}.
L'implication (2) $\Rightarrow$ (1) est le
Lemme \ref{algebra-lemma-relative-global-complete-intersection}.
Il reste à montrer que (3) implique (2).

\medskip\noindent
Supposons (3). Après avoir remplacé $S$ par $S_g$ pour un certain
$g \in S$, $g\not\in \mathfrak q$, nous pouvons supposer que $S$ est de
présentation finie sur $R$.
Choisissons une présentation $S = R[x_1, \ldots, x_n]/I$. Soit
$\mathfrak q' \subset R[x_1, \ldots, x_n]$ l'idéal premier correspondant à
$\mathfrak q$. Écrivons $\kappa(\mathfrak p) = k$.
Notons que $S \otimes_R k = k[x_1, \ldots, x_n]/\overline{I}$, où
$\overline{I} \subset k[x_1, \ldots, x_n]$ est l'idéal engendré par l'image
de $I$. Soit
$\overline{\mathfrak q}' \subset k[x_1, \ldots, x_n]$ l'idéal premier
engendré par l'image de $\mathfrak q'$.
D'après le Lemme \ref{algebra-lemma-lci-at-prime}, les conditions équivalentes du
Lemme \ref{algebra-lemma-lci} sont satisfaites pour $\overline{I}$ et
$\overline{\mathfrak q}'$.
Disons que la dimension de
$\overline{I}_{\overline{\mathfrak q}'}/
\overline{\mathfrak q}'\overline{I}_{\overline{\mathfrak q}'}$
sur $\kappa(\overline{\mathfrak q}')$ est $c$.
Choisissons $f_1, \ldots, f_c \in I$ qui s'envoient sur une base de cet
espace vectoriel. Les images $\overline{f}_j \in \overline{I}$ engendrent
$\overline{I}_{\overline{\mathfrak q}'}$
(d'après le Lemme \ref{algebra-lemma-lci}).
Posons $S' = R[x_1, \ldots, x_n]/(f_1, \ldots, f_c)$. Soit $J$ le noyau de
la surjection $S' \to S$. Comme $S$ est de présentation finie, $J$ est un
idéal de type fini (Lemme \ref{algebra-lemma-compose-finite-type}).
Considérons la suite exacte courte
$$
0 \to J \to S' \to S \to 0
$$
Comme $S_\mathfrak q$ est plat sur $R$, nous voyons que
$J_{\mathfrak q'} \otimes_R k \to S'_{\mathfrak q'} \otimes_R k$
est injectif (Lemme \ref{algebra-lemma-flat-tor-zero}).
Or, par construction, $S'_{\mathfrak q'} \otimes_R k$ s'envoie
isomorphiquement sur $S_\mathfrak q \otimes_R k$. Nous concluons donc que
$J_{\mathfrak q'} \otimes_R k =
J_{\mathfrak q'}/\mathfrak pJ_{\mathfrak q'} = 0$. D'après le lemme de
Nakayama (Lemme \ref{algebra-lemma-NAK}), il existe donc
$g \in R[x_1, \ldots, x_n]$, $g \not \in \mathfrak q'$ tel que $J_g = 0$.
Autrement dit, $S'_g \cong S_g$. Après une localisation supplémentaire, nous
voyons que $S'$ (et donc $S$) devient une intersection complète globale
relative d'après le
Lemme \ref{algebra-lemma-localize-relative-complete-intersection}, comme souhaité.
\end{proof}

\begin{lemma}
\label{algebra-lemma-syntomic-presentation-ideal-mod-squares}
Soit $R$ un anneau. Soit $S = R[x_1, \ldots, x_n]/I$ pour un certain idéal
$I$ de type fini. Si $g \in S$ est tel que $S_g$ soit syntomique sur $R$,
alors $(I/I^2)_g$ est un $S_g$-module projectif de type fini.
\end{lemma}

\begin{proof}
D'après le Lemme \ref{algebra-lemma-syntomic}, il existe un nombre fini d'éléments
$g_1, \ldots, g_m \in S$ qui engendrent l'idéal unité dans $S_g$ et tels que
chacun des $S_{gg_j}$ soit une intersection complète globale relative sur
$R$. Comme il suffit de montrer que $(I/I^2)_{gg_j}$ est projectif de type
fini, voir le Lemme \ref{algebra-lemma-finite-projective}, nous pouvons supposer que
$S_g$ est une intersection complète globale relative.
Dans ce cas, le résultat découle des Lemmes
\ref{algebra-lemma-conormal-module-localize} et
\ref{algebra-lemma-relative-global-complete-intersection-conormal}.
\end{proof}

\begin{lemma}
\label{algebra-lemma-composition-syntomic}
Soient $R \to S$ et $S \to S'$ des morphismes d'anneaux.
\begin{enumerate}
\item Si $R \to S$ et $S \to S'$ sont syntomiques, alors $R \to S'$ est
syntomique.
\item Si $R \to S$ et $S \to S'$ sont des intersections complètes globales
relatives, alors $R \to S'$ est une intersection complète globale relative.
\end{enumerate}
\end{lemma}

\begin{proof}
Démonstration de (2). Supposons que $R \to S$ et $S \to S'$ soient des
intersections complètes globales relatives et que nous ayons des présentations
$S =  R[x_1, \ldots, x_n]/(f_1, \ldots, f_c)$ et
$S' = S[y_1, \ldots, y_m]/(h_1, \ldots, h_d)$ comme dans la
Définition \ref{algebra-definition-relative-global-complete-intersection}.
Alors
$$
S' \cong
R[x_1, \ldots, x_n, y_1, \ldots, y_m]/(f_1, \ldots, f_c, h'_1, \ldots, h'_d)
$$
pour certains relèvements
$h_j' \in R[x_1, \ldots, x_n, y_1, \ldots, y_m]$ des $h_j$.
Il suffit donc de borner les dimensions des anneaux fibres.
Nous pouvons ainsi supposer que $R = k$ est un corps.
Dans ce cas, nous avons un anneau, à savoir $S$, qui est de type fini sur $k$
et équidimensionnel de dimension $n - c$, ainsi qu'un morphisme d'anneaux de
type fini $S \to S'$ dont tous les anneaux fibres non vides sont
équidimensionnels de dimension $m - d$. Alors, par exemple d'après le
Lemme \ref{algebra-lemma-dimension-base-fibre-total} appliqué aux localisations en les
idéaux maximaux de $S'$, nous voyons que
$\dim(S') \leq n - c + m - d$, comme souhaité.

\medskip\noindent
Nous allons ramener la partie (1) à la partie (2). Supposons que $R \to S$ et
$S \to S'$ soient syntomiques. Soit $\mathfrak q' \subset S$ un idéal premier
situé au-dessus de $\mathfrak q \subset S$. D'après le
Lemme \ref{algebra-lemma-syntomic}, il existe $g' \in S'$,
$g' \not \in \mathfrak q'$, tel que $S \to S'_{g'}$ soit une intersection
complète globale relative. De même, nous trouvons $g \in S$,
$g \not \in \mathfrak q$, tel que $R \to S_g$ soit une intersection complète
globale relative.
D'après le Lemme \ref{algebra-lemma-base-change-relative-global-complete-intersection},
le morphisme d'anneaux $S_g \to S_{gg'}$ est une intersection complète
globale relative. D'après la partie (2), nous voyons que
$R \to S_{gg'}$ est une intersection complète globale relative et que
$gg' \not \in \mathfrak q'$.
Comme $\mathfrak q'$ était arbitraire, en combinant les
Lemmes \ref{algebra-lemma-syntomic} et \ref{algebra-lemma-local-syntomic}, nous voyons que
$R \to S'$ est syntomique (cela utilise aussi que le spectre de $S'$ est
quasi-compact ; voir le Lemme \ref{algebra-lemma-quasi-compact}).
\end{proof}

\noindent
Le lemme suivant sera amélioré plus tard ; voir Lissage des morphismes
d'anneaux, Proposition \ref{smoothing-proposition-lift-smooth}.

\begin{lemma}
\label{algebra-lemma-lift-syntomic}
Soit $R$ un anneau et soit $I \subset R$ un idéal.
Soit $R/I \to \overline{S}$ un morphisme syntomique.
Alors il existe des éléments $\overline{g}_i \in \overline{S}$ qui engendrent
l'idéal unité de $\overline{S}$ et tels que chacun des
$\overline{S}_{g_i} \cong S_i/IS_i$ pour une certaine intersection complète
globale relative $S_i$ sur $R$.
\end{lemma}

\begin{proof}
D'après le Lemme \ref{algebra-lemma-syntomic}, nous trouvons une famille d'éléments
$\overline{g}_i \in \overline{S}$ qui engendrent l'idéal unité de
$\overline{S}$ et tels que chacun des $\overline{S}_{g_i}$ soit une
intersection complète globale relative sur $R/I$.
Nous pouvons donc supposer que $\overline{S}$ est une intersection complète
globale relative.
Écrivons
$\overline{S} =
(R/I)[x_1, \ldots, x_n]/(\overline{f}_1, \ldots, \overline{f}_c)$
comme dans la Définition \ref{algebra-definition-relative-global-complete-intersection}.
Choisissons $f_1, \ldots, f_c \in R[x_1, \ldots, x_n]$ relevant
$\overline{f}_1, \ldots, \overline{f}_c$.
Posons $S = R[x_1, \ldots, x_n]/(f_1, \ldots, f_c)$.
Notons que $S/IS \cong \overline{S}$.
D'après le Lemme \ref{algebra-lemma-localize-relative-complete-intersection}, nous
pouvons trouver $g \in S$ qui s'envoie sur $1$ dans $\overline{S}$ et tel que
$S_g$ soit une intersection complète globale relative sur $R$.
Comme $\overline{S} \cong S_g/IS_g$, cela achève la démonstration.
\end{proof}


\section{Morphismes d'anneaux lisses}
\label{algebra-section-smooth}

\noindent
Motivons la définition d'un morphisme d'anneaux lisse par un exemple.
Supposons que $R$ soit un anneau et que $S = R[x, y]/(f)$ pour un certain
$f \in R[x, y]$ non nul. Il existe alors une suite exacte
$$
S \to
S\text{d}x \oplus S\text{d}y \to
\Omega_{S/R} \to 0
$$
où la première flèche envoie $1$ sur
$\frac{\partial f}{\partial x} \text{d}x +
\frac{\partial f}{\partial y} \text{d}y$ ; voir la
section \ref{algebra-section-netherlander}.
Nous concluons que $\Omega_{S/R}$ est localement libre de rang $1$ si les
dérivées partielles de $f$ engendrent l'idéal unité dans $S$.
Dans ce cas, $S$ est lisse de dimension relative $1$ sur $R$.
Mais il peut arriver que $\Omega_{S/R}$ soit localement libre de rang $2$, à
savoir lorsque les deux dérivées partielles de $f$ sont nulles. Cela se
produit par exemple si, pour un nombre premier $p$, on a $p = 0$ dans $R$ et
$f = x^p + y^p$. Ici, $R \to S$ est une intersection complète globale
relative de dimension relative $1$ qui n'est pas lisse.
Ainsi, pour vérifier qu'un morphisme d'anneaux est lisse, il ne suffit pas de
vérifier si le module des différentielles est libre. La condition correcte est
la suivante.

\begin{definition}
\label{algebra-definition-smooth}
Un morphisme d'anneaux $R \to S$ est {\it lisse} s'il est de présentation
finie et si le complexe cotangent naïf $\NL_{S/R}$ est quasi-isomorphe à un
$S$-module projectif de type fini placé en degré $0$ : cela signifie que
$H_1(\NL_{S/R}) = 0$ et que $\Omega_{S/R}$ est un $S$-module projectif de
type fini.
\end{definition}

\noindent
Soit $R \to S$ un morphisme d'anneaux. D'après le
Lemme \ref{algebra-lemma-NL-homotopy}, pour toute présentation
$\alpha : P \to S$, on a $H_1(\NL(\alpha)) = H_1(\NL_{S/R})$ et
$H_0(\NL(\alpha)) = \Omega_{S/R}$. Ainsi, si $R \to S$ est lisse, alors,
pour toute surjection $\alpha : R[x_1, \ldots, x_n] \to S$ de noyau $I$, le
morphisme
$$
I/I^2
\longrightarrow
\Omega_{R[x_1, \ldots, x_n]/R} \otimes_{R[x_1, \ldots, x_n]} S
$$
est injectif et son conoyau est un $S$-module projectif de type fini.
La flèche affichée est donc une injection scindée. Autrement dit,
$\bigoplus_{i = 1}^n S \text{d}x_i \cong I/I^2 \oplus \Omega_{S/R}$
comme $S$-modules, et $I/I^2$ est lui aussi un $S$-module projectif de type
fini. Réciproquement, si $R \to S$ est de présentation finie et si nous avons
une surjection $\alpha : R[x_1, \ldots, x_n] \to S$ de noyau $I$ telle que la
flèche affichée soit une injection scindée, alors $R \to S$ est lisse.

\begin{lemma}
\label{algebra-lemma-localize-smooth}
Soit $R \to S$ un morphisme d'anneaux lisse.
Toute localisation $S_g$ est lisse sur $R$.
Si $f \in R$ s'envoie sur un élément inversible de $S$, alors
$R_f \to S$ est lisse.
\end{lemma}

\begin{proof}
D'après le Lemme \ref{algebra-lemma-localize-NL}, le complexe cotangent naïf de $S_g$
sur $R$ est le changement de base du complexe cotangent naïf de $S$ sur $R$.
L'hypothèse est que le complexe cotangent naïf de $S/R$ est $\Omega_{S/R}$ et
que c'est un $S$-module projectif de type fini. Il en va donc de même de son
changement de base. Ainsi, $S_g$ est lisse sur $R$.

\medskip\noindent
La seconde assertion découle de la même manière du
Lemme \ref{algebra-lemma-NL-localize-bottom}.
\end{proof}

\begin{lemma}
\label{algebra-lemma-base-change-smooth}
\begin{slogan}
La lissité est préservée par changement de base
\end{slogan}
Soit $R \to S$ un morphisme d'anneaux lisse.
Soit $R \to R'$ un morphisme d'anneaux quelconque.
Alors le changement de base $R' \to S' = R' \otimes_R S$ est lisse.
\end{lemma}

\begin{proof}
Soit $\alpha : R[x_1, \ldots, x_n] \to S$ une présentation de noyau $I$.
Soit $\alpha' : R'[x_1, \ldots, x_n] \to R' \otimes_R S$ la présentation
induite. Posons $I' = \Ker(\alpha')$.
Puisque $0 \to I \to R[x_1, \ldots, x_n] \to S \to 0$ est exacte, la suite
$R' \otimes_R I \to R'[x_1, \ldots, x_n] \to R' \otimes_R S \to 0$
est exacte. Ainsi, $R' \otimes_R I \to I'$ est surjectif.
D'après la Définition \ref{algebra-definition-smooth}, il existe une suite exacte
courte
$$
0 \to I/I^2 \to
\Omega_{R[x_1, \ldots, x_n]/R} \otimes_{R[x_1, \ldots, x_n]} S \to
\Omega_{S/R} \to
0
$$
et le $S$-module $\Omega_{S/R}$ est projectif de type fini.
En particulier, $I/I^2$ est un facteur direct de
$\Omega_{R[x_1, \ldots, x_n]/R} \otimes_{R[x_1, \ldots, x_n]} S$.
Considérons le diagramme commutatif
$$
\xymatrix{
R' \otimes_R (I/I^2) \ar[r] \ar[d] &
R' \otimes_R (\Omega_{R[x_1, \ldots, x_n]/R} \otimes_{R[x_1, \ldots, x_n]} S)
\ar[d] \\
I'/(I')^2 \ar[r] &
\Omega_{R'[x_1, \ldots, x_n]/R'}
\otimes_{R'[x_1, \ldots, x_n]} (R' \otimes_R S)
}
$$
Comme le morphisme vertical de droite est un isomorphisme, nous voyons, par ce
qui précède, que le morphisme vertical de gauche est injectif et surjectif.
Nous concluons donc que $\NL(\alpha')$ est quasi-isomorphe à
$\Omega_{S'/R'} \cong S' \otimes_S \Omega_{S/R}$ placé en degré $0$.
Ce module est projectif de type fini, car c'est le changement de base d'un
module projectif de type fini.
\end{proof}

\begin{lemma}
\label{algebra-lemma-smooth-over-field}
Soit $k$ un corps.
Soit $S$ une $k$-algèbre lisse.
Alors $S$ est une intersection complète locale.
\end{lemma}

\begin{proof}
D'après les Lemmes \ref{algebra-lemma-base-change-smooth} et
\ref{algebra-lemma-lci-field-change}, il suffit de le démontrer lorsque $k$ est
algébriquement clos. Choisissons une présentation
$\alpha : k[x_1, \ldots, x_n] \to S$ de noyau $I$. Soit $\mathfrak m$ un
idéal maximal de $S$, et soit $\mathfrak m' \supset I$ l'idéal maximal
correspondant de $k[x_1, \ldots, x_n]$.
Nous allons montrer que la condition (5) du
Lemme \ref{algebra-lemma-lci} est satisfaite (avec $\mathfrak m$ à la place de
$\mathfrak q$).
Nous pouvons écrire $\mathfrak m' = (x_1 - a_1, \ldots, x_n - a_n)$ pour
certains $a_i \in k$, car $k$ est algébriquement clos ; voir le
Théorème \ref{algebra-theorem-nullstellensatz}.
Puisque, par hypothèse, $k \to S$ est lisse, le morphisme de $S$-modules
$\text{d} : I/I^2 \to \bigoplus_{i = 1}^n S \text{d}x_i$
est une injection scindée. Le morphisme correspondant
$I/\mathfrak m' I \to \bigoplus \kappa(\mathfrak m') \text{d}x_i$
est donc injectif. Écrivons
$\dim_{\kappa(\mathfrak m')}(I/\mathfrak m' I) = c$ et choisissons
$f_1, \ldots, f_c \in I$ dont les images forment une base sur
$\kappa(\mathfrak m')$ de $I/\mathfrak m' I$.
D'après le Lemme de Nakayama \ref{algebra-lemma-NAK}, nous voyons que
$f_1, \ldots, f_c$ engendrent $I_{\mathfrak m'}$ sur
$k[x_1, \ldots, x_n]_{\mathfrak m'}$. Considérons le diagramme commutatif
$$
\xymatrix{
I \ar[r] \ar[d] & I/I^2 \ar[rr] \ar[d] & &
I/\mathfrak m'I \ar[d] \\
\Omega_{k[x_1, \ldots, x_n]/k} \ar[r] &
\bigoplus S\text{d}x_i \ar[rr]^{\text{d}x_i \mapsto x_i - a_i} & &
\mathfrak m'/(\mathfrak m')^2
}
$$
(démonstration de la commutativité omise). Le morphisme vertical du milieu est
celui qui définit le complexe cotangent naïf de $\alpha$. Notons que la flèche
horizontale inférieure de droite induit un isomorphisme
$\bigoplus \kappa(\mathfrak m') \text{d}x_i \to \mathfrak m'/(\mathfrak m')^2$.
Ainsi, nos générateurs $f_1, \ldots, f_c$ de $I_{\mathfrak m'}$ s'envoient sur
une famille d'éléments de $k[x_1, \ldots, x_n]_{\mathfrak m'}$ dont les
classes dans $\mathfrak m'/(\mathfrak m')^2$ sont linéairement indépendantes
sur $\kappa(\mathfrak m')$. Ils forment donc une suite régulière dans l'anneau
$k[x_1, \ldots, x_n]_{\mathfrak m'}$ d'après le
Lemme \ref{algebra-lemma-regular-ring-CM}.
Cela vérifie la condition (5) du Lemme \ref{algebra-lemma-lci} ; par conséquent,
$S_g$ est une intersection complète globale sur $k$ pour un certain
$g \in S$, $g \not \in \mathfrak m$. Comme cela vaut pour tout idéal maximal
de $S$, nous concluons que $S$ est une intersection complète locale sur $k$.
\end{proof}

\begin{definition}
\label{algebra-definition-standard-smooth}
Soit $R$ un anneau. Étant donnés des entiers $n \geq c \geq 0$ et
$f_1, \ldots, f_c \in R[x_1, \ldots, x_n]$, nous disons que
$$
R[x_1, \ldots, x_n]/(f_1, \ldots, f_c)
$$
est une {\it algèbre lisse standard sur $R$} si le polynôme
$$
g =
\det
\left(
\begin{matrix}
\partial f_1/\partial x_1 &
\partial f_2/\partial x_1 &
\ldots &
\partial f_c/\partial x_1 \\
\partial f_1/\partial x_2 &
\partial f_2/\partial x_2 &
\ldots &
\partial f_c/\partial x_2 \\
\ldots & \ldots & \ldots & \ldots \\
\partial f_1/\partial x_c &
\partial f_2/\partial x_c &
\ldots &
\partial f_c/\partial x_c
\end{matrix}
\right)
$$
s'envoie sur un élément inversible de
$R[x_1, \ldots, x_n]/(f_1, \ldots, f_c)$.
Nous disons qu'une $R$-algèbre $S$ est {\it lisse standard}, ou que le
morphisme d'anneaux $R \to S$ est {\it lisse standard}, s'il existe
$n \geq c \geq 0$ et $f_1, \ldots, f_c \in R[x_1, \ldots, x_n]$ tels que
$R[x_1, \ldots, x_n]/(f_1, \ldots, f_c)$ soit une algèbre lisse standard sur
$R$ et que $S$ soit isomorphe à
$R[x_1, \ldots, x_n]/(f_1, \ldots, f_c)$ comme $R$-algèbre.
\end{definition}

\begin{lemma}
\label{algebra-lemma-standard-smooth}
Soit
$S = R[x_1, \ldots, x_n]/(f_1, \ldots, f_c) = R[x_1, \ldots, x_n]/I$
une algèbre lisse standard. Alors
\begin{enumerate}
\item le morphisme d'anneaux $R \to S$ est lisse,
\item le $S$-module $\Omega_{S/R}$ est libre sur
$\text{d}x_{c + 1}, \ldots, \text{d}x_n$,
\item le $S$-module $I/I^2$ est libre sur les classes de
$f_1, \ldots, f_c$,
\item pour tout $g \in S$, le morphisme d'anneaux $R \to S_g$ est lisse
standard,
\item pour tout morphisme d'anneaux $R \to R'$, le changement de base
$R' \to R'\otimes_R S$ est lisse standard,
\item si $f \in R$ s'envoie sur un élément inversible de $S$, alors
$R_f \to S$ est lisse standard, et
\item l'anneau $S$ est une intersection complète globale relative sur $R$.
\end{enumerate}
\end{lemma}

\begin{proof}
Considérons le complexe cotangent naïf de la présentation donnée
$$
(f_1, \ldots, f_c)/(f_1, \ldots, f_c)^2
\longrightarrow
\bigoplus\nolimits_{i = 1}^n S \text{d}x_i
$$
Composons ce morphisme avec la projection sur les $c$ premiers facteurs
directs de la somme directe. D'après la définition d'une algèbre lisse
standard, les classes $f_i \bmod (f_1, \ldots, f_c)^2$ s'envoient sur une
base de $\bigoplus_{i = 1}^c S\text{d}x_i$. Nous concluons que
$(f_1, \ldots, f_c)/(f_1, \ldots, f_c)^2$ est libre de rang $c$, avec pour
base les éléments $f_i \bmod (f_1, \ldots, f_c)^2$, et que l'homologie en
degré $0$, c'est-à-dire $\Omega_{S/R}$, du complexe cotangent naïf est un
$S$-module libre dont une base est formée par les images de
$\text{d}x_{c + j}$, $j = 1, \ldots, n - c$.
Cela prouve en particulier que $R \to S$ est lisse.

\medskip\noindent
Les démonstrations de (4) et (6) sont omises. Voir toutefois l'exemple
ci-dessous et la démonstration du
Lemme \ref{algebra-lemma-base-change-relative-global-complete-intersection}.

\medskip\noindent
Soit $\varphi : R \to R'$ un morphisme d'anneaux quelconque.
Notons $S' = R'[x_1, \ldots, x_n]/(f_1^\varphi, \ldots, f_c^\varphi)$, où
$f^\varphi$ est le polynôme obtenu à partir de
$f \in R[x_1, \ldots, x_n]$ en appliquant $\varphi$ à tous les coefficients.
Alors $S' \cong R' \otimes_R S$.
De plus, le déterminant de la Définition \ref{algebra-definition-standard-smooth}
pour $S'/R'$ est égal à $g^\varphi$. Son image dans $S'$ est donc l'image de
$g$ par $R[x_1, \ldots, x_n] \to S \to S'$, et est par conséquent
inversible. Cela prouve (5).

\medskip\noindent
Pour prouver (7), il suffit de montrer que
$S \otimes_R \kappa(\mathfrak p)$ est de dimension $n - c$ pour tout idéal
premier $\mathfrak p \subset R$.
D'après (5), il suffit de montrer que toute algèbre lisse standard
$k[x_1, \ldots, x_n]/(f_1, \ldots, f_c)$ sur un corps $k$ est de dimension
$n - c$. Nous savons déjà que
$k[x_1, \ldots, x_n]/(f_1, \ldots, f_c)$ est une intersection complète
locale d'après le Lemme \ref{algebra-lemma-smooth-over-field}.
Ainsi, puisque $I/I^2$ est libre de rang $c$, nous voyons que
$k[x_1, \ldots, x_n]/(f_1, \ldots, f_c)$ est de dimension $n - c$, par
exemple d'après le Lemme \ref{algebra-lemma-lci}.
\end{proof}

\begin{example}
\label{algebra-example-make-standard-smooth}
Soit $R$ un anneau.
Soient $f_1, \ldots, f_c \in R[x_1, \ldots, x_n]$.
Soit
$$
h =
\det
\left(
\begin{matrix}
\partial f_1/\partial x_1 &
\partial f_2/\partial x_1 &
\ldots &
\partial f_c/\partial x_1 \\
\partial f_1/\partial x_2 &
\partial f_2/\partial x_2 &
\ldots &
\partial f_c/\partial x_2 \\
\ldots & \ldots & \ldots & \ldots \\
\partial f_1/\partial x_c &
\partial f_2/\partial x_c &
\ldots &
\partial f_c/\partial x_c
\end{matrix}
\right).
$$
Posons
$S = R[x_1, \ldots, x_{n + 1}]/(f_1, \ldots, f_c, x_{n + 1}h - 1)$.
C'est un exemple d'algèbre lisse standard, si ce n'est que la présentation est
mauvaise et que les variables devraient être rangées dans l'ordre suivant :
$x_1, \ldots, x_c, x_{n + 1}, x_{c + 1}, \ldots, x_n$.
\end{example}

\begin{lemma}
\label{algebra-lemma-compose-standard-smooth}
Une composée de morphismes d'anneaux lisses standard est lisse standard.
\end{lemma}

\begin{proof}
Supposons que $R \to S$ et $S \to S'$ soient lisses standard. Choisissons des
présentations
$S =  R[x_1, \ldots, x_n]/(f_1, \ldots, f_c)$
et
$S' = S[y_1, \ldots, y_m]/(g_1, \ldots, g_d)$.
Choisissons des éléments
$g_j' \in R[x_1, \ldots, x_n, y_1, \ldots, y_m]$ qui s'envoient sur les
$g_j$. Nous voyons ainsi que
$S' = R[x_1, \ldots, x_n, y_1, \ldots, y_m]/
(f_1, \ldots, f_c, g'_1, \ldots, g'_d)$.
Pour montrer que $S'$ est lisse standard, il suffit de vérifier que le
déterminant
$$
\det
\left(
\begin{matrix}
\partial f_1/\partial x_1 &
\ldots &
\partial f_c/\partial x_1 &
\partial g_1/\partial x_1 &
\ldots &
\partial g_d/\partial x_1 \\
\ldots &
\ldots &
\ldots &
\ldots &
\ldots &
\ldots  \\
\partial f_1/\partial x_c &
\ldots &
\partial f_c/\partial x_c &
\partial g_1/\partial x_c &
\ldots &
\partial g_d/\partial x_c \\
0 &
\ldots &
0 &
\partial g_1/\partial y_1 &
\ldots &
\partial g_d/\partial y_1 \\
\ldots &
\ldots &
\ldots &
\ldots &
\ldots &
\ldots  \\
0 &
\ldots &
0 &
\partial g_1/\partial y_d &
\ldots &
\partial g_d/\partial y_d
\end{matrix}
\right)
$$
est inversible dans $S'$. C'est clair, puisqu'il est le produit des deux
déterminants qui sont inversibles par hypothèse.
\end{proof}

\begin{lemma}
\label{algebra-lemma-smooth-syntomic}
Soit $R \to S$ un morphisme d'anneaux lisse.
Il existe un recouvrement ouvert de $\Spec(S)$ par des ouverts standards
$D(g)$ tel que chacun des $S_g$ soit lisse standard sur $R$.
En particulier, $R \to S$ est syntomique.
\end{lemma}

\begin{proof}
Choisissons une présentation $\alpha : R[x_1, \ldots, x_n] \to S$ de noyau
$I = (f_1, \ldots, f_m)$. Pour toute partie
$E \subset \{1, \ldots, m\}$, considérons l'ouvert $U_E$ sur lequel les
classes $f_e, e\in E$ engendrent librement le $S$-module projectif de type
fini $I/I^2$ ; voir le Lemme \ref{algebra-lemma-cokernel-flat}.
Nous pouvons recouvrir $\Spec(S)$ par des ouverts standards $D(g)$ dont chacun
est entièrement contenu dans l'un des ouverts $U_E$. Pour un tel $g$, nous
considérons la présentation
$$
\beta : R[x_1, \ldots, x_n, x_{n + 1}] \longrightarrow S_g
$$
qui envoie $x_{n + 1}$ sur $1/g$. Posant $J = \Ker(\beta)$, nous utilisons le
Lemme \ref{algebra-lemma-principal-localization-NL} pour voir que
$J/J^2 \cong (I/I^2)_g \oplus S_g$ est libre.
Nous pouvons remplacer, et remplaçons, $S$ par $S_g$. Alors, en utilisant le
Lemme \ref{algebra-lemma-huber}, nous pouvons supposer que nous avons une présentation
$\alpha : R[x_1, \ldots, x_n] \to S$ de noyau
$I = (f_1, \ldots, f_c)$ telle que $I/I^2$ soit libre sur les classes de
$f_1, \ldots, f_c$.

\medskip\noindent
En utilisant la présentation $\alpha$ obtenue à la fin du paragraphe
précédent, nous répétons plus ou moins cet argument avec les éléments de base
$\text{d}x_1, \ldots, \text{d}x_n$ de
$\Omega_{R[x_1, \ldots, x_n]/R}$.
En effet, pour toute partie $E \subset \{1, \ldots, n\}$ de cardinal $c$,
nous pouvons considérer l'ouvert $U_E$ de $\Spec(S)$ sur lequel la
différentielle de $\NL(\alpha)$, composée avec la projection
$$
S^{\oplus c} \cong I/I^2
\longrightarrow
\Omega_{R[x_1, \ldots, x_n]/R} \otimes_{R[x_1, \ldots, x_n]} S
\longrightarrow
\bigoplus\nolimits_{i \in E} S\text{d}x_i
$$
est un isomorphisme. À nouveau, nous pouvons trouver un recouvrement de
$\Spec(S)$ par des ouverts standards $D(g)$ (en nombre fini) tel que chaque
$D(g)$ soit entièrement contenu dans l'un des ouverts $U_E$.
En renumérotant, nous pouvons supposer que $E = \{1, \ldots, c\}$.
Pour un $g$ tel que $D(g) \subset U_E$, nous considérons la présentation
$$
\beta : R[x_1, \ldots, x_n, x_{n + 1}] \to S_g
$$
qui envoie $x_{n + 1}$ sur $1/g$. Posant $J = \Ker(\beta)$, nous déduisons du
Lemme \ref{algebra-lemma-principal-localization-NL} que
$J = (f_1, \ldots, f_c, fx_{n + 1} - 1)$, où $\alpha(f) = g$, et que la
composée
$$
J/J^2 \longrightarrow
\Omega_{R[x_1, \ldots, x_{n + 1}]/R} \otimes_{R[x_1, \ldots, x_{n + 1}]} S_g
\longrightarrow
\bigoplus\nolimits_{i = 1}^c S_g\text{d}x_i \oplus S_g \text{d}x_{n + 1}
$$
est un isomorphisme. En réordonnant les coordonnées sous la forme
$x_1, \ldots, x_c, x_{n + 1}, x_{c + 1}, \ldots, x_n$, nous concluons que
$S_g$ est lisse standard sur $R$, comme souhaité.

\medskip\noindent
Cela achève la démonstration, car les algèbres lisses standard sont
syntomiques (Lemmes \ref{algebra-lemma-standard-smooth} et
\ref{algebra-lemma-relative-global-complete-intersection}) et la propriété d'être
syntomique sur $R$ est locale sur $S$
(Lemme \ref{algebra-lemma-local-syntomic}).
\end{proof}

\begin{definition}
\label{algebra-definition-smooth-at-prime}
Soit $R \to S$ un morphisme d'anneaux.
Soit $\mathfrak q$ un idéal premier de $S$.
Nous disons que $R \to S$ est {\it lisse en $\mathfrak q$} s'il existe
$g \in S$, $g \not \in \mathfrak q$, tel que
$R \to S_g$ soit lisse.
\end{definition}

\noindent
Pour les morphismes d'anneaux de présentation finie, on peut caractériser
cette propriété comme suit.

\begin{lemma}
\label{algebra-lemma-smooth-at-point}
Soit $R \to S$ de présentation finie. Soit $\mathfrak q$ un idéal premier de
$S$. Les conditions suivantes sont équivalentes :
\begin{enumerate}
\item $R \to S$ est lisse en $\mathfrak q$,
\item $H_1(L_{S/R})_\mathfrak q = 0$ et
$\Omega_{S/R, \mathfrak q}$ est un $S_\mathfrak q$-module libre de type fini,
\item $H_1(L_{S/R})_\mathfrak q = 0$ et
$\Omega_{S/R, \mathfrak q}$ est un $S_\mathfrak q$-module projectif, et
\item $H_1(L_{S/R})_\mathfrak q = 0$ et
$\Omega_{S/R, \mathfrak q}$ est un $S_\mathfrak q$-module plat.
\end{enumerate}
\end{lemma}

\begin{proof}
Nous utiliserons sans autre mention le fait que la formation du complexe
cotangent naïf commute à la localisation ; voir la
section \ref{algebra-section-netherlander}, en particulier le
Lemme \ref{algebra-lemma-localize-NL}.
Notons que $\Omega_{S/R}$ est un $S$-module de présentation finie ; voir le
Lemme \ref{algebra-lemma-differentials-finitely-presented}. Ainsi (2), (3) et (4) sont
équivalentes par le Lemme \ref{algebra-lemma-finite-projective}.
Il est clair que (1) implique les conditions équivalentes (2), (3) et (4).
Supposons (2). En écrivant $S_\mathfrak q$ comme la colimite de localisations
principales, le Lemme \ref{algebra-lemma-colimit-category-fp-modules} montre qu'il
existe $g \in S$, $g \not \in \mathfrak q$, tel que
$(\Omega_{S/R})_g$ soit libre de type fini. Choisissons une présentation
$\alpha : R[x_1, \ldots, x_n] \to S$ de noyau $I$. Nous pouvons travailler
avec $\NL(\alpha)$ au lieu de $\NL_{S/R}$ ; voir le
Lemme \ref{algebra-lemma-NL-homotopy}. La surjection
$$
\Omega_{R[x_1, \ldots, x_n]/R} \otimes_{R[x_1, \ldots, x_n]} S
\to \Omega_{S/R} \to 0
$$
admet un inverse à droite après inversion de $g$, car $(\Omega_{S/R})_g$ est
projectif. L'image de
$\text{d} : (I/I^2)_g \to
\Omega_{R[x_1, \ldots, x_n]/R} \otimes_{R[x_1, \ldots, x_n]} S_g$
est donc un facteur direct, et ce morphisme admet lui aussi un inverse à
droite. Nous en déduisons que $H_1(L_{S/R})_g$ est un quotient de $(I/I^2)_g$.
En particulier, $H_1(L_{S/R})_g$ est un $S_g$-module de type fini.
L'annulation de $H_1(L_{S/R})_{\mathfrak q}$ entraîne donc celle de
$H_1(L_{S/R})_{gg'}$ pour un certain $g' \in S$,
$g' \not \in \mathfrak q$. Alors $R \to S_{gg'}$ est lisse par définition.
\end{proof}

\begin{lemma}
\label{algebra-lemma-locally-smooth}
\begin{slogan}
Un morphisme d'anneaux est lisse si et seulement s'il est lisse en tout idéal
premier de la cible
\end{slogan}
Soit $R \to S$ un morphisme d'anneaux.
Alors $R \to S$ est lisse si et seulement si $R \to S$ est lisse
en tout idéal premier $\mathfrak q$ de $S$.
\end{lemma}

\begin{proof}
L'implication directe est triviale. Supposons que $R \to S$ soit lisse en tout
idéal premier $\mathfrak q$ de $S$. Comme $\Spec(S)$ est quasi-compact ; voir
le Lemme \ref{algebra-lemma-quasi-compact}, il existe un recouvrement fini
$\Spec(S) = \bigcup D(g_i)$ tel que chaque $S_{g_i}$ soit lisse.
Le Lemme \ref{algebra-lemma-cover-upstairs} implique alors que $S$ est de présentation
finie sur $R$. D'après le Lemme \ref{algebra-lemma-localize-NL},
$\NL_{S/R} \otimes_S S_{g_i}$ est quasi-isomorphe à un $S_{g_i}$-module
projectif de type fini placé en degré $0$. Le
Lemme \ref{algebra-lemma-finite-projective} implique alors que $\NL_{S/R}$ est
quasi-isomorphe à un $S$-module projectif de type fini placé en degré $0$.
\end{proof}

\begin{lemma}
\label{algebra-lemma-compose-smooth}
La composée de morphismes d'anneaux lisses est lisse.
\end{lemma}

\begin{proof}
On peut le démontrer de nombreuses façons. L'une d'elles consiste à utiliser
le lemme du serpent (Lemme \ref{algebra-lemma-snake}), la suite de Jacobi--Zariski
(Lemme \ref{algebra-lemma-exact-sequence-NL}), ainsi que la caractérisation des modules
projectifs comme facteurs directs de modules libres
(Lemme \ref{algebra-lemma-characterize-projective}). Une autre démonstration s'obtient
en combinant les Lemmes \ref{algebra-lemma-smooth-syntomic},
\ref{algebra-lemma-compose-standard-smooth} et \ref{algebra-lemma-locally-smooth}.
\end{proof}

\begin{lemma}
\label{algebra-lemma-product-smooth}
Soit $R$ un anneau. Soit $S = S' \times S''$ un produit de $R$-algèbres.
Alors $S$ est lisse sur $R$ si et seulement si $S'$ et $S''$ sont toutes deux
lisses sur $R$.
\end{lemma}

\begin{proof}
Omis. Indications : par le Lemme \ref{algebra-lemma-locally-smooth}, on peut vérifier
la lissité idéal premier par idéal premier. Comme $\Spec(S)$ est la réunion
disjointe de $\Spec(S')$ et de $\Spec(S'')$ par le
Lemme \ref{algebra-lemma-spec-product}, la lissité de $R \to S$ en $\mathfrak q$
correspond soit à la lissité de $R \to S'$ en l'idéal premier correspondant,
soit à celle de $R \to S''$ en l'idéal premier correspondant.
\end{proof}

\begin{lemma}
\label{algebra-lemma-relative-global-complete-intersection-smooth}
Soit $R$ un anneau. Soit $S = R[x_1, \ldots, x_n]/(f_1, \ldots, f_c)$
une intersection complète globale relative.
Soit $\mathfrak q \subset S$ un idéal premier. Alors $R \to S$ est lisse en
$\mathfrak q$ si et seulement s'il existe une partie
$I \subset \{1, \ldots, n\}$ de cardinal $c$ telle que le polynôme
$$
g_I = \det (\partial f_j/\partial x_i)_{j = 1, \ldots, c, \ i \in I}.
$$
ne s'envoie pas sur un élément de $\mathfrak q$.
\end{lemma}

\begin{proof}
Le Lemme \ref{algebra-lemma-relative-global-complete-intersection-conormal} montre que
le complexe cotangent naïf associé à la présentation donnée de $S$ est le
complexe
$$
\bigoplus\nolimits_{j = 1}^c S \cdot f_j
\longrightarrow
\bigoplus\nolimits_{i = 1}^n S \cdot \text{d}x_i, \quad
f_j \longmapsto \sum \frac{\partial f_j}{\partial x_i} \text{d}x_i.
$$
Les mineurs maximaux de la matrice qui donne ce morphisme sont précisément les
polynômes $g_I$.

\medskip\noindent
Supposons que $g_I$ s'envoie sur $g \in S$, avec
$g \not \in \mathfrak q$. Alors l'algèbre $S_g$ est lisse sur $R$.
En effet, son complexe cotangent naïf est quasi-isomorphe au complexe ci-dessus
localisé en $g$ ; voir le Lemme \ref{algebra-lemma-localize-NL}. Par construction, il
est quasi-isomorphe à un module libre de rang $n - c$ en degré $0$.

\medskip\noindent
Réciproquement, supposons que tous les $g_I$ appartiennent finalement à
$\mathfrak q$. Dans ce cas, le complexe ci-dessus tensorisé avec
$\kappa(\mathfrak q)$ n'est pas de rang maximal ; il n'existe donc aucune
localisation par un élément $g \in S$, $g \not \in \mathfrak q$, dans laquelle
ce morphisme devienne une injection scindée. De nouveau, le
Lemme \ref{algebra-lemma-localize-NL} montre qu'aucune telle localisation n'est lisse
sur $R$.
\end{proof}

\begin{lemma}
\label{algebra-lemma-flat-fibre-smooth}
Soit $R \to S$ un morphisme d'anneaux.
Soit $\mathfrak q \subset S$ un idéal premier au-dessus de l'idéal premier
$\mathfrak p$ de $R$. Supposons que
\begin{enumerate}
\item il existe $g \in S$, $g \not\in \mathfrak q$, tel que
$R \to S_g$ soit de présentation finie,
\item l'homomorphisme d'anneaux locaux
$R_{\mathfrak p} \to S_{\mathfrak q}$ soit plat,
\item la fibre $S \otimes_R \kappa(\mathfrak p)$ soit lisse sur
$\kappa(\mathfrak p)$ en l'idéal premier correspondant à $\mathfrak q$.
\end{enumerate}
Alors $R \to S$ est lisse en $\mathfrak q$.
\end{lemma}

\begin{proof}
Les Lemmes \ref{algebra-lemma-syntomic} et \ref{algebra-lemma-smooth-over-field} montrent qu'il
existe $g \in S$ tel que $S_g$ soit une intersection complète globale relative.
En remplaçant $S$ par $S_g$, nous pouvons supposer que
$S = R[x_1, \ldots, x_n]/(f_1, \ldots, f_c)$ est une intersection complète
globale relative. Pour toute partie $I \subset \{1, \ldots, n\}$ de cardinal
$c$, considérons le polynôme
$g_I = \det (\partial f_j/\partial x_i)_{j = 1, \ldots, c, i \in I}$
du Lemme \ref{algebra-lemma-relative-global-complete-intersection-smooth}.
Notons que l'image $\overline{g}_I$ de $g_I$ dans l'anneau de polynômes
$\kappa(\mathfrak p)[x_1, \ldots, x_n]$ est le déterminant des dérivées
partielles des images $\overline{f}_j$ des $f_j$ dans l'anneau
$\kappa(\mathfrak p)[x_1, \ldots, x_n]$. Le lemme résulte donc de l'application
du Lemme \ref{algebra-lemma-relative-global-complete-intersection-smooth} à la fois à
$R \to S$ et à $\kappa(\mathfrak p) \to S \otimes_R \kappa(\mathfrak p)$.
\end{proof}

\noindent
Notons que les ensembles $U, V$ du lemme suivant sont ouverts par définition.

\begin{lemma}
\label{algebra-lemma-flat-base-change-locus-smooth}
Soit $R \to S$ un morphisme d'anneaux de présentation finie.
Soit $R \to R'$ un morphisme d'anneaux plat.
Notons $S' = R' \otimes_R S$ le changement de base.
Soit $U \subset \Spec(S)$ l'ensemble des idéaux premiers auxquels
$R \to S$ est lisse. Soit $V \subset \Spec(S')$ l'ensemble des idéaux premiers
auxquels $R' \to S'$ est lisse. Alors $V$ est l'image réciproque de $U$ par le
morphisme $f : \Spec(S') \to \Spec(S)$.
\end{lemma}

\begin{proof}
Le Lemme \ref{algebra-lemma-change-base-NL} montre que
$\NL_{S/R} \otimes_S S'$ est homotopiquement équivalent à $\NL_{S'/R'}$.
Cela implique déjà que $f^{-1}(U) \subset V$.

\medskip\noindent
Soit $\mathfrak q' \subset S'$ un idéal premier au-dessus de
$\mathfrak q \subset S$. Supposons $\mathfrak q' \in V$.
Il faut montrer que $\mathfrak q \in U$.
Comme $S \to S'$ est plat, $S_{\mathfrak q} \to S'_{\mathfrak q'}$ est
fidèlement plat (Lemme \ref{algebra-lemma-local-flat-ff}). L'annulation de
$H_1(L_{S'/R'})_{\mathfrak q'}$ implique donc celle de
$H_1(L_{S/R})_{\mathfrak q}$. En appliquant le
Lemme \ref{algebra-lemma-finite-projective-descends} au $S_{\mathfrak q}$-module
$(\Omega_{S/R})_{\mathfrak q}$ et au morphisme
$S_{\mathfrak q} \to S'_{\mathfrak q'}$, nous voyons que
$(\Omega_{S/R})_{\mathfrak q}$ est projectif. Ainsi $R \to S$ est lisse en
$\mathfrak q$ par le Lemme \ref{algebra-lemma-smooth-at-point}.
\end{proof}

\begin{lemma}
\label{algebra-lemma-smooth-field-change-local}
Soit $K/k$ une extension de corps.
Soit $S$ une algèbre de type fini sur $k$.
Soit $\mathfrak q_K$ un idéal premier de $S_K = K \otimes_k S$, et soit
$\mathfrak q$ l'idéal premier correspondant de $S$.
Alors $S$ est lisse sur $k$ en $\mathfrak q$ si et seulement si
$S_K$ est lisse en $\mathfrak q_K$ sur $K$.
\end{lemma}

\begin{proof}
C'est un cas particulier du Lemme \ref{algebra-lemma-flat-base-change-locus-smooth}.
\end{proof}

\begin{lemma}
\label{algebra-lemma-lift-smooth}
Soit $R$ un anneau et soit $I \subset R$ un idéal.
Soit $R/I \to \overline{S}$ un morphisme d'anneaux lisse.
Alors il existe des éléments $\overline{g}_i \in \overline{S}$ qui engendrent
l'idéal unité de $\overline{S}$, tels que chacun des
$\overline{S}_{g_i} \cong S_i/IS_i$ pour un certain anneau $S_i$ lisse
(standard) sur $R$.
\end{lemma}

\begin{proof}
Le Lemme \ref{algebra-lemma-smooth-syntomic} fournit une collection d'éléments
$\overline{g}_i \in \overline{S}$ qui engendrent l'idéal unité de
$\overline{S}$, telle que chacun des $\overline{S}_{g_i}$ soit lisse standard
sur $R/I$. Nous pouvons donc supposer que $\overline{S}$ est lisse standard
sur $R/I$. Écrivons
$\overline{S} =
(R/I)[x_1, \ldots, x_n]/(\overline{f}_1, \ldots, \overline{f}_c)$
comme dans la Définition \ref{algebra-definition-standard-smooth}.
Choisissons $f_1, \ldots, f_c \in R[x_1, \ldots, x_n]$ relevant
$\overline{f}_1, \ldots, \overline{f}_c$. Posons
$S = R[x_1, \ldots, x_n, x_{n + 1}]/(f_1, \ldots, f_c, x_{n + 1}\Delta - 1)$,
où $\Delta = \det(\frac{\partial f_j}{\partial x_i})_{i, j = 1, \ldots, c}$,
comme dans l'Exemple \ref{algebra-example-make-standard-smooth}.
Cela démontre le lemme.
\end{proof}






\section{Morphismes formellement lisses}
\label{algebra-section-formally-smooth}

\noindent
Dans cette section, nous définissons les morphismes d'anneaux formellement
lisses. Il s'avérera qu'un morphisme d'anneaux de présentation finie est
formellement lisse si et seulement s'il est lisse ; voir la
Proposition \ref{algebra-proposition-smooth-formally-smooth}.

\begin{definition}
\label{algebra-definition-formally-smooth}
Soit $R \to S$ un morphisme d'anneaux.
Nous disons que $S$ est {\it formellement lisse sur $R$} si, pour tout
diagramme commutatif aux flèches pleines
$$
\xymatrix{
S \ar[r] \ar@{-->}[rd] & A/I \\
R \ar[r] \ar[u] & A \ar[u]
}
$$
où $I \subset A$ est un idéal de carré nul, il existe une flèche en pointillés
qui rend le diagramme commutatif.
\end{definition}

\begin{lemma}
\label{algebra-lemma-base-change-fs}
Soit $R \to S$ un morphisme d'anneaux formellement lisse.
Soit $R \to R'$ un morphisme d'anneaux quelconque.
Alors le changement de base $S' = R' \otimes_R S$ est formellement lisse sur
$R'$.
\end{lemma}

\begin{proof}
Considérons un diagramme aux flèches pleines
$$
\xymatrix{
S \ar[r] \ar@{-->}[rrd] & R' \otimes_R S \ar[r] \ar@{-->}[rd] & A/I \\
R  \ar[u] \ar[r] & R' \ar[r] \ar[u] & A \ar[u]
}
$$
comme dans la Définition \ref{algebra-definition-formally-smooth}.
Par hypothèse, la flèche en pointillés la plus longue existe. La propriété
universelle du produit tensoriel fournit alors la plus courte.
\end{proof}

\begin{lemma}
\label{algebra-lemma-compose-formally-smooth}
La composée de morphismes d'anneaux formellement lisses est formellement lisse.
\end{lemma}

\begin{proof}
Omis. (Indication : c'est entièrement formel et résulte de la considération
d'un diagramme convenable.)
\end{proof}

\begin{lemma}
\label{algebra-lemma-polynomial-ring-formally-smooth}
Un anneau de polynômes sur $R$ est formellement lisse sur $R$.
\end{lemma}

\begin{proof}
Supposons donné un diagramme comme dans la
Définition \ref{algebra-definition-formally-smooth}, avec $S = R[x_j; j \in J]$.
Il existe alors une flèche en pointillés : il suffit de choisir des relèvements
$a_j \in A$ des éléments de $A/I$ sur lesquels les éléments $x_j$ sont envoyés
par la flèche horizontale supérieure.
\end{proof}

\begin{lemma}
\label{algebra-lemma-characterize-formally-smooth}
Soit $R \to S$ un morphisme d'anneaux.
Soit $P \to S$ un morphisme surjectif de $R$-algèbres d'un anneau de polynômes
$P$ sur $S$. Notons $J \subset P$ son noyau. Alors $R \to S$ est formellement
lisse si et seulement s'il existe un morphisme de $R$-algèbres
$\sigma : S \to P/J^2$ qui soit un inverse à droite de la surjection
$P/J^2 \to S$.
\end{lemma}

\begin{proof}
Supposons que $R \to S$ soit formellement lisse.
Considérons le diagramme commutatif
$$
\xymatrix{
S \ar[r] \ar@{-->}[rd] & P/J \\
R \ar[r] \ar[u] &  P/J^2\ar[u]
}
$$
Par hypothèse, la flèche en pointillés existe. Cela démontre l'existence de
$\sigma$.

\medskip\noindent
Réciproquement, supposons donné un $\sigma$ comme dans le lemme.
Considérons un diagramme aux flèches pleines
$$
\xymatrix{
S \ar[r] \ar@{-->}[rd] & A/I \\
R \ar[r] \ar[u] & A \ar[u]
}
$$
comme dans la Définition \ref{algebra-definition-formally-smooth}.
Puisque $P$ est formellement lisse par le
Lemme \ref{algebra-lemma-polynomial-ring-formally-smooth}, il existe un homomorphisme de
$R$-algèbres $\psi : P \to A$ qui relève le morphisme $P \to S \to A/I$.
Clairement, $\psi(J) \subset I$ et, comme $I^2 = 0$, nous en déduisons que
$\psi(J^2) = 0$. Ainsi $\psi$ se factorise en
$\overline{\psi} : P/J^2 \to A$. La flèche en pointillés recherchée est la
composée $\overline{\psi} \circ \sigma : S \to A$.
\end{proof}

\begin{remark}
\label{algebra-remark-lemma-characterize-formally-smooth}
Le Lemme \ref{algebra-lemma-characterize-formally-smooth} reste valable, plus
généralement, dès que $P$ est formellement lisse sur $R$.
\end{remark}

\begin{lemma}
\label{algebra-lemma-characterize-formally-smooth-again}
Soit $R \to S$ un morphisme d'anneaux.
Soit $P \to S$ un morphisme surjectif de $R$-algèbres d'un anneau de polynômes
$P$ sur $S$. Notons $J \subset P$ son noyau. Alors $R \to S$ est formellement
lisse si et seulement si la suite
$$
0 \to J/J^2 \to \Omega_{P/R} \otimes_P S \to \Omega_{S/R} \to 0
$$
du Lemme \ref{algebra-lemma-differential-seq} est une suite exacte scindée.
\end{lemma}

\begin{proof}
Supposons $S$ formellement lisse sur $R$. D'après le
Lemme \ref{algebra-lemma-characterize-formally-smooth}, cela signifie qu'il existe un
morphisme de $R$-algèbres $S \to P/J^2$ qui est un inverse à droite du morphisme
canonique $P/J^2 \to S$. Par le
Lemme \ref{algebra-lemma-differential-mod-power-ideal}, nous avons
$\Omega_{P/R} \otimes_P S = \Omega_{(P/J^2)/R} \otimes_{P/J^2} S$.
Le Lemme \ref{algebra-lemma-differential-seq-split} montre que la suite est scindée.

\medskip\noindent
Supposons que la suite exacte du lemme soit exacte scindée.
Choisissons un scindage $\sigma : \Omega_{S/R} \to
\Omega_{P/R} \otimes_P S$. Pour chaque $\lambda \in S$, choisissons
$x_\lambda \in P$ qui s'envoie sur $\lambda$. Puis, pour chaque
$\lambda \in S$, choisissons $f_\lambda \in J$ tel que
$$
\text{d}f_\lambda = \text{d}x_\lambda - \sigma(\text{d}\lambda)
$$
dans le terme médian de la suite exacte.
Nous affirmons que $s : \lambda \mapsto x_\lambda - f_\lambda \mod J^2$
est un homomorphisme de $R$-algèbres $s : S \to P/J^2$.
Pour le démontrer, nous utiliserons à plusieurs reprises que si $h \in J$ et
$\text{d}h = 0$ dans $\Omega_{P/R} \otimes_R S$, alors $h \in J^2$.
Soient $\lambda, \mu \in S$.
Alors $\sigma(\text{d}\lambda + \text{d}\mu - \text{d}(\lambda + \mu)) = 0$.
Cela implique
$$
\text{d}(x_\lambda + x_\mu - x_{\lambda + \mu}
- f_\lambda - f_\mu + f_{\lambda + \mu}) = 0
$$
et donc $x_\lambda + x_\mu - x_{\lambda + \mu}
- f_\lambda - f_\mu + f_{\lambda + \mu} \in J^2$, ce qui signifie à son tour
que $s(\lambda) + s(\mu) = s(\lambda + \mu)$.
De même,
$\sigma(\lambda \text{d}\mu + \mu \text{d}\lambda - \text{d}\lambda \mu) = 0$,
ce qui implique
$$
\mu(\text{d}x_\lambda - \text{d}f_\lambda) +
\lambda(\text{d}x_\mu - \text{d}f_\mu) -
\text{d}x_{\lambda\mu} + \text{d}f_{\lambda\mu} = 0
$$
dans le terme médian de la suite exacte.
De plus,
$$
\text{d}(x_\lambda x_\mu) =
x_\lambda \text{d}x_\mu + x_\mu \text{d}x_\lambda =
\lambda \text{d}x_\mu + \mu \text{d} x_\lambda
$$
dans le terme médian, là encore. Ensemble, ces équations signifient que
$x_\lambda x_\mu - x_{\lambda\mu}
- \mu f_\lambda - \lambda f_\mu + f_{\lambda\mu} \in J^2$ ; ainsi
$(x_\lambda - f_\lambda)(x_\mu - f_\mu) -
(x_{\lambda\mu} - f_{\lambda\mu}) \in J^2$, puisque
$f_\lambda f_\mu \in J^2$, ce qui signifie que
$s(\lambda)s(\mu) = s(\lambda\mu)$.
Si $\lambda \in R$, alors $\text{d}\lambda = 0$ et nous voyons que
$\text{d}f_\lambda = \text{d}x_\lambda$, d'où
$\lambda - x_\lambda + f_\lambda \in J^2$, puis
$s(\lambda) = \lambda$, comme souhaité. Nous pouvons maintenant appliquer le
Lemme \ref{algebra-lemma-characterize-formally-smooth} pour conclure que $S/R$ est
formellement lisse.
\end{proof}




\begin{proposition}
\label{algebra-proposition-characterize-formally-smooth}
Soit $R \to S$ un morphisme d'anneaux. Considérons une $R$-algèbre $P$
formellement lisse et une surjection $P \to S$ de noyau $J$. Les conditions
suivantes sont équivalentes :
\begin{enumerate}
\item $S$ est formellement lisse sur $R$,
\item pour un certain $P \to S$ comme ci-dessus, il existe une section de
$P/J^2 \to S$,
\item pour tout $P \to S$ comme ci-dessus, il existe une section de
$P/J^2 \to S$,
\item pour un certain $P \to S$ comme ci-dessus, la suite
$0 \to J/J^2 \to \Omega_{P/R} \otimes S \to \Omega_{S/R} \to 0$ est exacte
scindée,
\item pour tout $P \to S$ comme ci-dessus, la suite
$0 \to J/J^2 \to \Omega_{P/R} \otimes S \to \Omega_{S/R} \to 0$ est exacte
scindée, et
\item le complexe cotangent naïf $\NL_{S/R}$ est quasi-isomorphe à un
$S$-module projectif placé en degré $0$ : cela signifie que
$H_1(\NL_{S/R}) = 0$ et que $\Omega_{S/R}$ est un $S$-module projectif.
\end{enumerate}
\end{proposition}

\begin{proof}
Il est clair que (1) implique (3), qui implique (2) ; voir la première partie
de la démonstration du Lemme \ref{algebra-lemma-characterize-formally-smooth}.
Il est également vrai que (3) implique (5), qui implique (4), et que (2)
implique (4) ; voir la première partie de la démonstration du
Lemme \ref{algebra-lemma-characterize-formally-smooth-again}.
Enfin, le Lemme \ref{algebra-lemma-characterize-formally-smooth-again}, appliqué à la
surjection canonique $R[S] \to S$ (\ref{algebra-equation-canonical-presentation}),
montre que (1) implique (6).

\medskip\noindent
Supposons (4) et démontrons (6). Considérons la suite du
Lemme \ref{algebra-lemma-exact-sequence-NL} associée aux morphismes d'anneaux
$R \to P \to S$. Par l'implication (1) $\Rightarrow$ (6) démontrée ci-dessus,
nous voyons que $\NL_{P/R} \otimes_R S$ est quasi-isomorphe à
$\Omega_{P/R} \otimes_P S$ placé en degré $0$.
Ainsi $H_1(\NL_{P/R} \otimes_P S) = 0$. Puisque $P \to S$ est surjectif,
$\NL_{S/P}$ est homotopiquement équivalent à $J/J^2$ placé en degré $1$
(Lemme \ref{algebra-lemma-NL-surjection}). Nous obtenons donc la suite exacte
$0 \to H_1(L_{S/R}) \to J/J^2 \to \Omega_{P/R} \otimes_P S \to
\Omega_{S/R} \to 0$.
Par hypothèse, $H_1(L_{S/R}) = 0$ et $\Omega_{S/R}$ est un $S$-module
projectif. Ainsi (6) en résulte.

\medskip\noindent
Enfin, démontrons que (6) implique (1). L'hypothèse signifie que le complexe
$J/J^2 \to \Omega_{P/R} \otimes S$, où $P = R[S]$ et $P \to S$ est la
surjection canonique (\ref{algebra-equation-canonical-presentation}), est
quasi-isomorphe à un $S$-module projectif placé en degré $0$.
Le Lemme \ref{algebra-lemma-characterize-formally-smooth-again} montre donc que $S$
est formellement lisse sur $R$.
\end{proof}

\begin{lemma}
\label{algebra-lemma-ses-formally-smooth}
Soient $A \to B \to C$ des morphismes d'anneaux. Supposons que $B \to C$
soit formellement lisse. Alors la suite
$$
0 \to \Omega_{B/A} \otimes_B C \to \Omega_{C/A} \to \Omega_{C/B} \to 0
$$
du Lemme \ref{algebra-lemma-exact-sequence-differentials} est une suite exacte courte
scindée.
\end{lemma}

\begin{proof}
Cela résulte de la
Proposition \ref{algebra-proposition-characterize-formally-smooth} et du
Lemme \ref{algebra-lemma-exact-sequence-NL}.
\end{proof}

\begin{lemma}
\label{algebra-lemma-differential-seq-formally-smooth}
Soient $A \to B \to C$ des morphismes d'anneaux, avec $A \to C$ formellement
lisse et $B \to C$ surjectif de noyau $J \subset B$.
Alors la suite exacte
$$
0 \to J/J^2 \to \Omega_{B/A} \otimes_B C \to \Omega_{C/A} \to 0
$$
du Lemme \ref{algebra-lemma-differential-seq} est exacte scindée.
\end{lemma}

\begin{proof}
Cela résulte de la
Proposition \ref{algebra-proposition-characterize-formally-smooth}, du
Lemme \ref{algebra-lemma-exact-sequence-NL} et du
Lemme \ref{algebra-lemma-differential-seq}.
\end{proof}

\begin{lemma}
\label{algebra-lemma-application-NL-formally-smooth}
Soient $A \to B \to C$ des morphismes d'anneaux. Supposons que $A \to C$ soit
surjectif (de sorte que $B \to C$ le soit aussi) et que $A \to B$ soit
formellement lisse. Notons $I = \Ker(A \to C)$ et $J = \Ker(B \to C)$.
Alors la suite
$$
0 \to I/I^2 \to J/J^2 \to \Omega_{B/A} \otimes_B B/J \to 0
$$
du Lemme \ref{algebra-lemma-application-NL} est exacte scindée.
\end{lemma}

\begin{proof}
Comme $A \to B$ est formellement lisse, il existe un morphisme d'anneaux
$\sigma : B \to A/I^2$ dont la composée avec $A \to B$ est le morphisme
quotient $A \to A/I^2$. Alors $\sigma$ induit un morphisme
$J/J^2 \to I/I^2$ inverse du morphisme $I/I^2 \to J/J^2$.
\end{proof}

\begin{lemma}
\label{algebra-lemma-lift-formal-smoothness}
Soit $R \to S$ un morphisme d'anneaux.
Soit $I \subset R$ un idéal. Supposons que
\begin{enumerate}
\item $I^2 = 0$,
\item $R \to S$ est plat, et
\item $R/I \to S/IS$ est formellement lisse.
\end{enumerate}
Alors $R \to S$ est formellement lisse.
\end{lemma}

\begin{proof}
Supposons (1), (2) et (3).
Soit $P = R[\{x_t\}_{t \in T}] \to S$ une surjection de $R$-algèbres de
noyau $J$. Ainsi $0 \to J \to P \to S \to 0$ est une suite exacte courte de
$R$-modules plats. Cela implique $I \otimes_R S = IS$,
$I \otimes_R P = IP$ et $I \otimes_R J = IJ$, ainsi que $J \cap IP = IJ$.
Nous utiliserons dans toute la démonstration que
$$
\Omega_{(S/IS)/(R/I)} = \Omega_{S/R} \otimes_S (S/IS)
= \Omega_{S/R} \otimes_R R/I = \Omega_{S/R} / I\Omega_{S/R}
$$
et de même pour $P$ (voir le Lemme \ref{algebra-lemma-differentials-base-change}).
Par le Lemme \ref{algebra-lemma-characterize-formally-smooth-again}, la suite
\begin{equation}
\label{algebra-equation-split}
0 \to J/(IJ + J^2) \to
\Omega_{P/R} \otimes_P S/IS \to
\Omega_{S/R} \otimes_S S/IS \to 0
\end{equation}
est exacte scindée. Bien entendu, le terme médian est
$\bigoplus_{t \in T} S/IS \text{d}x_t$. Choisissons un scindage
$\sigma : \Omega_{P/R} \otimes_P S/IS \to J/(IJ + J^2)$.
Pour chaque $t \in T$, choisissons un élément $f_t \in J$ qui s'envoie sur
$\sigma(\text{d}x_t)$ dans $J/(IJ + J^2)$. Cela détermine un unique morphisme
de $S$-modules
$$
\tilde \sigma : \Omega_{P/R} \otimes_R S
= \bigoplus S\text{d}x_t \longrightarrow J/J^2
$$
tel que $\tilde\sigma(\text{d}x_t) = f_t$.
Comme $\sigma$ est une section de $\text{d}$, la différence
$$
\Delta = \text{id}_{J/J^2} - \tilde \sigma \circ \text{d}
$$
est un endomorphisme de $J/J^2 \to J/J^2$ dont l'image est contenue dans
$(IJ + J^2)/J^2$. En particulier, $\Delta((IJ + J^2)/J^2) = 0$, car
$I^2 = 0$. Cela signifie que $\Delta$ se factorise en
$$
J/J^2 \to J/(IJ + J^2) \xrightarrow{\overline{\Delta}}
(IJ + J^2)/J^2 \to J/J^2
$$
où $\overline{\Delta}$ est un morphisme de $S/IS$-modules.
En utilisant de nouveau le fait que la suite (\ref{algebra-equation-split}) est
scindée, nous pouvons trouver un morphisme de $S/IS$-modules
$\overline{\delta} : \Omega_{P/R} \otimes_P S/IS \to (IJ + J^2)/J^2$
tel que $\overline{\delta} \circ d$ soit égal à $\overline{\Delta}$.
De la même manière que ci-dessus, le morphisme $\overline{\delta}$ détermine
un morphisme de $S$-modules
$\delta : \Omega_{P/R} \otimes_P S \to J/J^2$.
Après remplacement de $\tilde \sigma$ par $\tilde \sigma + \delta$, un calcul
simple montre que $\Delta = 0$. Autrement dit, $\tilde \sigma$ est une section
de $J/J^2 \to \Omega_{P/R} \otimes_P S$.
Par le Lemme \ref{algebra-lemma-characterize-formally-smooth-again}, nous concluons que
$R \to S$ est formellement lisse.
\end{proof}

\begin{proposition}
\label{algebra-proposition-smooth-formally-smooth}
Soit $R \to S$ un morphisme d'anneaux. Les conditions suivantes sont
équivalentes :
\begin{enumerate}
\item $R \to S$ est de présentation finie et formellement lisse,
\item $R \to S$ est lisse.
\end{enumerate}
\end{proposition}

\begin{proof}
Cela résulte de la Proposition \ref{algebra-proposition-characterize-formally-smooth}
et de la Définition \ref{algebra-definition-smooth}.
(Notons que $\Omega_{S/R}$ est un $S$-module de présentation finie si
$R \to S$ est de présentation finie ; voir le
Lemme \ref{algebra-lemma-differentials-finitely-presented}.)
\end{proof}


\begin{lemma}
\label{algebra-lemma-finite-presentation-fs-Noetherian}
Soit $R \to S$ un morphisme d'anneaux lisse. Alors il existe un sous-anneau
$R_0 \subset R$ de type fini sur $\mathbf{Z}$ et un morphisme d'anneaux lisse
$R_0 \to S_0$ tels que $S \cong R \otimes_{R_0} S_0$.
\end{lemma}

\begin{proof}
Nous allons utiliser le fait que la lissité équivaut à la présentation finie
et à la lissité formelle ; voir la
Proposition \ref{algebra-proposition-smooth-formally-smooth}.
Écrivons $S = R[x_1, \ldots, x_n]/(f_1, \ldots, f_m)$ et notons
$I = (f_1, \ldots, f_m)$.
Choisissons un inverse à droite
$\sigma : S \to R[x_1, \ldots, x_n]/I^2$ de la projection sur $S$, comme dans
le Lemme \ref{algebra-lemma-characterize-formally-smooth}.
Choisissons $h_i \in R[x_1, \ldots, x_n]$ tels que
$\sigma(x_i \bmod I) = h_i \bmod I^2$. Puisque $x_i - h_i \in I$, il existe
des $b_{ij} \in R[x_1, \ldots, x_n]$ tels que
$$
x_i - h_i = \sum\nolimits_j b_{ij} f_j
$$
Le fait que $\sigma$ soit un homomorphisme de $R$-algèbres
$R[x_1, \ldots, x_n]/I \to R[x_1, \ldots, x_n]/I^2$ équivaut à la condition
$$
f_j(h_1, \ldots, h_n) = \sum\nolimits_{j_1 j_2} a_{j_1 j_2} f_{j_1} f_{j_2}
$$
pour certains $a_{kl} \in R[x_1, \ldots, x_n]$.
Soit $R_0 \subset R$ le sous-anneau engendré sur $\mathbf{Z}$ par tous les
coefficients des polynômes $f_j, h_i, a_{kl}, b_{ij}$.
Posons $S_0 = R_0[x_1, \ldots, x_n]/(f_1, \ldots, f_m)$, avec
$I_0 = (f_1, \ldots, f_m)$. Comme la seconde équation affichée est satisfaite
dans $R_0[x_1, \ldots, x_n]$, nous pouvons prendre pour
$\sigma_0 : S_0 \to R_0[x_1, \ldots, x_n]/I_0^2$ le morphisme de
$R_0$-algèbres défini par la règle $x_i \mapsto h_i \bmod I_0^2$.
Comme la première équation affichée est satisfaite dans
$R_0[x_1, \ldots, x_n]$, nous voyons que $\sigma_0$ est un inverse à droite de
la projection
$R_0[x_1, \ldots, x_n] / I_0^2 \to R_0[x_1, \ldots, x_n] / I_0 = S_0$.
Ainsi, par le Lemme \ref{algebra-lemma-characterize-formally-smooth}, l'anneau $S_0$
est formellement lisse sur $R_0$.
\end{proof}

\begin{lemma}
\label{algebra-lemma-smooth-descends-through-colimit}
Soit $A = \colim A_i$ une colimite filtrante d'anneaux. Soit $A \to B$ un
morphisme d'anneaux lisse. Il existe un $i$ et un morphisme d'anneaux lisse
$A_i \to B_i$ tels que $B = B_i \otimes_{A_i} A$.
\end{lemma}

\begin{proof}
Cela résulte du Lemme \ref{algebra-lemma-finite-presentation-fs-Noetherian}, car
$R_0 \to A$ se factorisera par $A_i$ pour un certain $i$, d'après le
Lemme \ref{algebra-lemma-characterize-finite-presentation}.
\end{proof}

\begin{lemma}
\label{algebra-lemma-descent-formally-smooth}
Soit $R \to S$ un morphisme d'anneaux. Soit $R \to R'$ un morphisme d'anneaux
fidèlement plat. Posons $S' = S \otimes_R R'$. Alors $R \to S$ est
formellement lisse si et seulement si $R' \to S'$ est formellement lisse.
\end{lemma}

\begin{proof}
Si $R \to S$ est formellement lisse, alors $R' \to S'$ l'est par le
Lemme \ref{algebra-lemma-base-change-fs}.
Pour démontrer la réciproque, supposons $R' \to S'$ formellement lisse.
Notons que $N \otimes_R R' = N \otimes_S S'$ pour tout $S$-module $N$.
En particulier, $S \to S'$ est lui aussi fidèlement plat.
Choisissons un anneau de polynômes $P = R[\{x_i\}_{i \in I}]$ et une surjection
de $R$-algèbres $P \to S$ de noyau $J$. Notons que
$P' = P \otimes_R R'$ est une algèbre de polynômes sur $R'$. Comme
$R \to R'$ est plat, le noyau $J'$ de la surjection $P' \to S'$ est
$J \otimes_R R'$. Par conséquent, la suite exacte scindée (voir le
Lemme \ref{algebra-lemma-characterize-formally-smooth-again})
$$
0 \to J'/(J')^2 \to \Omega_{P'/R'} \otimes_{P'} S' \to \Omega_{S'/R'} \to 0
$$
est le changement de base par $S \to S'$ de la suite correspondante
$$
J/J^2 \to \Omega_{P/R} \otimes_P S \to \Omega_{S/R} \to 0
$$
voir le Lemme \ref{algebra-lemma-differential-seq}.
Comme $S \to S'$ est fidèlement plat, nous en déduisons deux choses :
(1) cette suite (sans ${}'$) est elle aussi exacte et (2)
$\Omega_{S/R}$ est un $S$-module projectif. En effet, $\Omega_{S'/R'}$ est
projectif comme facteur direct du module libre
$\Omega_{P'/R'} \otimes_{P'} S'$, et
$\Omega_{S/R} \otimes_S {S'} = \Omega_{S'/R'}$ d'après ce que nous venons de
dire. Ainsi (2) résulte de la descente de la projectivité par les morphismes
d'anneaux fidèlement plats ; voir le
Théorème \ref{algebra-theorem-ffdescent-projectivity}.
La suite $0 \to J/J^2 \to \Omega_{P/R} \otimes_P S \to \Omega_{S/R} \to 0$
est donc elle aussi exacte, et nous concluons en appliquant de nouveau le
Lemme \ref{algebra-lemma-characterize-formally-smooth-again}.
\end{proof}

\noindent
Il s'avère que les morphismes d'anneaux lisses satisfont la propriété forte de
relèvement suivante.

\begin{lemma}
\label{algebra-lemma-smooth-strong-lift}
Soit $R \to S$ un morphisme d'anneaux lisse. Étant donné un diagramme
commutatif aux flèches pleines
$$
\xymatrix{
S \ar[r] \ar@{-->}[rd] & A/I \\
R \ar[r] \ar[u] & A \ar[u]
}
$$
où $I \subset A$ est un idéal localement nilpotent, il existe une flèche en
pointillés qui rend le diagramme commutatif.
\end{lemma}

\begin{proof}
Par le Lemme \ref{algebra-lemma-finite-presentation-fs-Noetherian}, nous pouvons
prolonger le diagramme en un diagramme commutatif
$$
\xymatrix{
S_0 \ar[r] & S \ar[r] \ar@{-->}[rd] & A/I \\
R_0 \ar[r] \ar[u] & R \ar[r] \ar[u] & A \ar[u]
}
$$
où $R_0 \to S_0$ est lisse, $R_0$ est de type fini sur $\mathbf{Z}$ et
$S = S_0 \otimes_{R_0} R$. Soient $x_1, \ldots, x_n \in S_0$ des générateurs
de $S_0$ sur $R_0$. Soient $a_1, \ldots, a_n$ des éléments de $A$ qui
s'envoient sur les mêmes éléments de $A/I$ que $x_1, \ldots, x_n$.
Notons $A_0 \subset A$ le sous-anneau engendré par l'image de $R_0$ et les
éléments $a_1, \ldots, a_n$. Posons $I_0 = A_0 \cap I$. Alors
$A_0/I_0 \subset A/I$, et $S_0 \to A/I$ est à valeurs dans $A_0/I_0$.
Il suffit donc de trouver la flèche en pointillés dans le diagramme
$$
\xymatrix{
S_0 \ar[r] \ar@{-->}[rd] & A_0/I_0 \\
R_0 \ar[r] \ar[u] & A_0 \ar[u]
}
$$
Par construction, l'anneau $A_0$ est de type fini sur $\mathbf{Z}$.
Ainsi $A_0$ est noethérien, donc $I_0$ est nilpotent ; voir le
Lemme \ref{algebra-lemma-Noetherian-power}.
Disons que $I_0^n = 0$. Par la
Proposition \ref{algebra-proposition-smooth-formally-smooth}, nous pouvons relever
successivement le morphisme de $R_0$-algèbres $S_0 \to A_0/I_0$ en
$S_0 \to A_0/I_0^2$, $S_0 \to A_0/I_0^3$, $\ldots$, et finalement en
$S_0 \to A_0/I_0^n = A_0$.
\end{proof}





\section{Lissité et différentielles}
\label{algebra-section-smooth-differential}

\noindent
Quelques résultats sur les différentielles et les morphismes d'anneaux lisses.

\begin{lemma}
\label{algebra-lemma-triangle-differentials-smooth}
Étant donnés des morphismes d'anneaux $A \to B \to C$, avec $B \to C$ lisse,
la suite
$$
0 \to C \otimes_B \Omega_{B/A} \to \Omega_{C/A} \to \Omega_{C/B} \to 0
$$
du Lemme \ref{algebra-lemma-exact-sequence-differentials} est exacte.
\end{lemma}

\begin{proof}
Cela résulte du Lemme \ref{algebra-lemma-ses-formally-smooth}, plus général, car un
morphisme d'anneaux lisse est formellement lisse ; voir la
Proposition \ref{algebra-proposition-smooth-formally-smooth}.
Mais cela résulte aussi directement du
Lemme \ref{algebra-lemma-exact-sequence-NL}, puisque $H_1(L_{C/B}) = 0$ fait partie de
la définition de la lissité de $B \to C$.
\end{proof}

\begin{lemma}
\label{algebra-lemma-differential-seq-smooth}
Soient $A \to B \to C$ des morphismes d'anneaux, avec $A \to C$ lisse et
$B \to C$ surjectif de noyau $J \subset B$.
Alors la suite exacte
$$
0 \to J/J^2 \to \Omega_{B/A} \otimes_B C \to \Omega_{C/A} \to 0
$$
du Lemme \ref{algebra-lemma-differential-seq} est exacte scindée.
\end{lemma}

\begin{proof}
Cela résulte du Lemme \ref{algebra-lemma-differential-seq-formally-smooth}, plus
général, car un morphisme d'anneaux lisse est formellement lisse ; voir la
Proposition \ref{algebra-proposition-smooth-formally-smooth}.
\end{proof}

\begin{lemma}
\label{algebra-lemma-application-NL-smooth}
Soient $A \to B \to C$ des morphismes d'anneaux. Supposons que $A \to C$ soit
surjectif (de sorte que $B \to C$ le soit aussi) et que $A \to B$ soit lisse.
Notons $I = \Ker(A \to C)$ et $J = \Ker(B \to C)$.
Alors la suite
$$
0 \to I/I^2 \to J/J^2 \to \Omega_{B/A} \otimes_B B/J \to 0
$$
du Lemme \ref{algebra-lemma-application-NL} est exacte.
\end{lemma}

\begin{proof}
Cela résulte du Lemme \ref{algebra-lemma-application-NL-formally-smooth}, plus général,
car un morphisme d'anneaux lisse est formellement lisse ; voir la
Proposition \ref{algebra-proposition-smooth-formally-smooth}.
\end{proof}

\begin{lemma}
\label{algebra-lemma-section-smooth}
\begin{slogan}
Si $R$ est un facteur direct de $S$ et si $S$ est lisse sur $R$, alors le
complété $I$-adique de $S$ est souvent un anneau de séries formelles sur $R$,
où $I$ est le noyau du morphisme de projection de $S$ sur $R$.
\end{slogan}
Soit $\varphi : R \to S$ un morphisme d'anneaux lisse.
Soit $\sigma : S \to R$ un inverse à gauche de $\varphi$.
Posons $I = \Ker(\sigma)$. Alors
\begin{enumerate}
\item $I/I^2$ est un $R$-module localement libre de type fini, et
\item si $I/I^2$ est libre, alors $S^\wedge \cong R[[t_1, \ldots, t_d]]$
comme $R$-algèbres, où $S^\wedge$ est le complété $I$-adique de $S$.
\end{enumerate}
\end{lemma}

\begin{proof}
En appliquant le Lemme \ref{algebra-lemma-differential-seq-split} à
$R \to S \to R$, nous voyons que
$I/I^2 = \Omega_{S/R} \otimes_{S, \sigma} R$.
Puisque, par définition d'un morphisme lisse, le module $\Omega_{S/R}$ est
localement libre de type fini sur $S$, nous en déduisons (1).
Si $I/I^2$ est libre, choisissons $f_1, \ldots, f_d \in I$ dont les images dans
$I/I^2$ forment une $R$-base. Considérons le morphisme de $R$-algèbres défini
par
$$
\Psi : R[[x_1, \ldots, x_d]] \longrightarrow S^\wedge, \quad
x_i \longmapsto f_i.
$$
Notons $P = R[[x_1, \ldots, x_d]]$ et $J = (x_1, \ldots, x_d) \subset P$.
Notons $\Psi_n : P/J^n \to S/I^n$ le morphisme induit entre les anneaux
quotients. Remarquons que $S/I^2 = \varphi(R) \oplus I/I^2$. Ainsi $\Psi_2$
est un isomorphisme. Notons $\sigma_2 : S/I^2 \to P/J^2$ l'inverse de
$\Psi_2$. Nous allons démontrer par récurrence sur $n$ que, pour tout $n > 2$,
il existe un inverse $\sigma_n : S/I^n \to P/J^n$ de $\Psi_n$. En effet,
comme $S$ est formellement lisse sur $R$ (par la
Proposition \ref{algebra-proposition-smooth-formally-smooth}), nous voyons que, dans le
diagramme aux flèches pleines
$$
\xymatrix{
S \ar@{..>}[r] \ar[rd]_{\sigma_{n - 1}} & P/J^n \ar[d] \\
 & P/J^{n - 1}
}
$$
de $R$-algèbres, nous pouvons compléter la flèche en pointillés par un certain
morphisme de $R$-algèbres $\tau : S \to P/J^n$ qui rende le diagramme
commutatif. Celui-ci induit un morphisme de $R$-algèbres
$\overline{\tau} : S/I^n \to P/J^n$ égal à $\sigma_{n - 1}$ modulo
$J^{n - 1}$. Par construction, le morphisme $\Psi_n$ est surjectif, et
$\overline{\tau} \circ \Psi_n$ est maintenant un endomorphisme de
$R$-algèbres de $P/J^n$ qui envoie $x_i$ sur $x_i + \delta_{i, n}$, avec
$\delta_{i, n} \in J^{n - 1}/J^n$. Il s'ensuit que $\Psi_n$ est un
isomorphisme et possède donc un inverse $\sigma_n$.
Cela démontre le lemme.
\end{proof}







\section{Algèbres lisses sur les corps}
\label{algebra-section-smooth-over-field}

\noindent
Avertissement : en général, les deux lemmes suivants ne sont pas valables sur
les corps non parfaits.

\begin{lemma}
\label{algebra-lemma-rank-omega}
Soit $k$ un corps algébriquement clos.
Soit $S$ une $k$-algèbre de type fini.
Soit $\mathfrak m \subset S$ un idéal maximal.
Alors
$$
\dim_{\kappa(\mathfrak m)} \Omega_{S/k} \otimes_S \kappa(\mathfrak m)
=
\dim_{\kappa(\mathfrak m)} \mathfrak m/\mathfrak m^2.
$$
\end{lemma}

\begin{proof}
Considérons la suite exacte
$$
\mathfrak m/\mathfrak m^2 \to
\Omega_{S/k} \otimes_S \kappa(\mathfrak m) \to
\Omega_{\kappa(\mathfrak m)/k} \to 0
$$
du Lemme \ref{algebra-lemma-differential-seq}. Nous voulons montrer que le premier
morphisme est un isomorphisme. Comme $k$ est algébriquement clos, la composée
$k \to \kappa(\mathfrak m)$ est un isomorphisme par le
Théorème \ref{algebra-theorem-nullstellensatz}.
Ainsi, la surjection $S \to \kappa(\mathfrak m)$ admet une section comme
morphisme de $k$-algèbres, et le Lemme \ref{algebra-lemma-differential-seq-split}
montre que la suite ci-dessus est exacte à gauche. Puisque
$\Omega_{\kappa(\mathfrak m)/k} = 0$, nous avons terminé.
\end{proof}

\begin{lemma}
\label{algebra-lemma-characterize-smooth-kbar}
Soit $k$ un corps algébriquement clos.
Soit $S$ une $k$-algèbre de type fini.
Soit $\mathfrak m \subset S$ un idéal maximal.
Les conditions suivantes sont équivalentes :
\begin{enumerate}
\item L'anneau $S_{\mathfrak m}$ est un anneau local régulier.
\item Nous avons
$\dim_{\kappa(\mathfrak m)} \Omega_{S/k} \otimes_S \kappa(\mathfrak m)
\leq \dim(S_{\mathfrak m})$.
\item Nous avons
$\dim_{\kappa(\mathfrak m)} \Omega_{S/k} \otimes_S \kappa(\mathfrak m)
= \dim(S_{\mathfrak m})$.
\item Il existe $g \in S$, $g \not \in \mathfrak m$, tel que $S_g$ soit lisse
sur $k$. Autrement dit, $S/k$ est lisse en $\mathfrak m$.
\end{enumerate}
\end{lemma}

\begin{proof}
Notons que (1), (2) et (3) sont équivalentes par le
Lemme \ref{algebra-lemma-rank-omega} et la Définition \ref{algebra-definition-regular}.

\medskip\noindent
Supposons que $S$ soit lisse en $\mathfrak m$.
Le Lemme \ref{algebra-lemma-smooth-syntomic} montre que $S_g$ est lisse standard sur
$k$ pour un $g \in S$, $g \not \in \mathfrak m$, convenable.
Par le Lemme \ref{algebra-lemma-standard-smooth}, $\Omega_{S_g/k}$ est donc libre de
rang $\dim(S_g)$. Le Lemme \ref{algebra-lemma-rank-omega} montre alors que
$\dim(S_{\mathfrak m}) = \dim (\mathfrak m/\mathfrak m^2)$ ; autrement dit,
$S_\mathfrak m$ est régulier.

\medskip\noindent
Réciproquement, supposons que $S_{\mathfrak m}$ soit régulier.
Soit $d = \dim(S_{\mathfrak m}) = \dim \mathfrak m/\mathfrak m^2$.
Choisissons une présentation $S = k[x_1, \ldots, x_n]/I$ telle que $x_i$
s'envoie sur un élément de $\mathfrak m$ pour tout $i$. Autrement dit,
$\mathfrak m'' = (x_1, \ldots, x_n)$ est l'idéal maximal correspondant de
$k[x_1, \ldots, x_n]$. Notons que nous avons une suite exacte courte
$$
I/\mathfrak m''I \to \mathfrak m''/(\mathfrak m'')^2
\to \mathfrak m/(\mathfrak m)^2 \to 0
$$
Choisissons $c = n - d$ éléments $f_1, \ldots, f_c \in I$ dont les images
dans $\mathfrak m''/(\mathfrak m'')^2$ engendrent le noyau du morphisme vers
$\mathfrak m/\mathfrak m^2$. C'est clairement possible. Notons
$J = (f_1, \ldots, f_c)$. Ainsi $J \subset I$.
Notons $S' = k[x_1, \ldots, x_n]/J$, de sorte qu'il existe une surjection
$S' \to S$. Notons $\mathfrak m' = \mathfrak m''S'$ l'idéal maximal
correspondant de $S'$. Nous avons donc
$$
\xymatrix{
k[x_1, \ldots, x_n] \ar[r] & S' \ar[r] & S \\
\mathfrak m'' \ar[u] \ar[r] & \mathfrak m' \ar[r] \ar[u] &
\mathfrak m \ar[u]
}
$$
Par notre choix de $J$, la suite exacte
$$
J/\mathfrak m''J \to \mathfrak m''/(\mathfrak m'')^2
\to \mathfrak m'/(\mathfrak m')^2 \to 0
$$
montre que $\dim( \mathfrak m'/(\mathfrak m')^2 ) = d$.
Comme $S'_{\mathfrak m'}$ se surjecte sur $S_{\mathfrak m}$, nous voyons que
$\dim(S_{\mathfrak m'}) \geq d$. La discussion précédant la
Définition \ref{algebra-definition-regular-local} permet donc de conclure que
$S'_{\mathfrak m'}$ est lui aussi régulier de dimension $d$. Comme $S'$ était
découpé par $c = n - d$ équations, nous concluons qu'il existe
$g' \in S'$, $g' \not \in \mathfrak m'$, tel que $S'_{g'}$ soit une
intersection complète globale sur $k$ ; voir le Lemme \ref{algebra-lemma-lci}.
De plus, le morphisme $S'_{\mathfrak m'} \to S_{\mathfrak m}$ est une
surjection de domaines locaux noethériens de même dimension, et donc un
isomorphisme. Ainsi $S' \to S$ est surjectif, son noyau est de type fini, et il
devient un isomorphisme après localisation en $\mathfrak m'$. Nous pouvons donc
trouver $g' \in S'$, $g \not \in \mathfrak m'$, tel que
$S'_{g'} \to S_{g'}$ soit un isomorphisme. En définitive, nous concluons
qu'après remplacement de $S$ par une localisation principale, nous pouvons
supposer que $S$ est une intersection complète globale.

\medskip\noindent
À ce stade, nous pouvons écrire $S = k[x_1, \ldots, x_n]/(f_1, \ldots, f_c)$,
avec $\dim S = n - c$. Rappelons que le complexe cotangent naïf de cette
algèbre est donné par
$$
\bigoplus S \cdot f_j
\to
\bigoplus S \cdot \text{d}x_i
$$
voir le Lemme \ref{algebra-lemma-relative-global-complete-intersection-conormal}.
Par le Lemme \ref{algebra-lemma-relative-global-complete-intersection-smooth}, pour
montrer que $S$ est lisse en $\mathfrak m$, nous devons montrer que l'un des
mineurs $c \times c$, notés $g_I$, de la matrice « $A$ » qui donne le morphisme
ci-dessus ne s'annule pas en $\mathfrak m$. Par le
Lemme \ref{algebra-lemma-rank-omega}, la matrice $A \bmod \mathfrak m$ est de rang $c$.
Nous avons donc terminé.
\end{proof}
\begin{lemma}
\label{algebra-lemma-characterize-smooth-over-field}
Soit $k$ un corps quelconque.
Soit $S$ une $k$-algèbre de type fini.
Soit $X = \Spec(S)$.
Soit $\mathfrak q \subset S$ un idéal premier correspondant à $x \in X$.
Les conditions suivantes sont équivalentes :
\begin{enumerate}
\item La $k$-algèbre $S$ est lisse en $\mathfrak q$ sur $k$.
\item Nous avons
$\dim_{\kappa(\mathfrak q)} \Omega_{S/k} \otimes_S \kappa(\mathfrak q)
\leq \dim_x X$.
\item Nous avons
$\dim_{\kappa(\mathfrak q)} \Omega_{S/k} \otimes_S \kappa(\mathfrak q)
= \dim_x X$.
\end{enumerate}
De plus, dans ce cas, l'anneau local $S_{\mathfrak q}$ est régulier.
\end{lemma}

\begin{proof}
Si $S$ est lisse en $\mathfrak q$ sur $k$, alors il existe
$g \in S$, $g \not \in \mathfrak q$, tel que $S_g$ soit lisse standard sur
$k$ ; voir le Lemme \ref{algebra-lemma-smooth-syntomic}. Le module des différentielles
d'une algèbre lisse standard sur $k$ est libre, de rang égal à la dimension ;
voir le Lemme \ref{algebra-lemma-standard-smooth} (utiliser qu'une intersection
complète globale relative sur un corps a pour dimension le nombre de variables
moins le nombre d'équations). Nous voyons ainsi que (1) implique (3). Pour
achever la démonstration du lemme, il suffit de montrer que (2) implique (1),
et qu'il implique que $S_{\mathfrak q}$ est régulier.

\medskip\noindent
Supposons (2). Par le lemme de Nakayama \ref{algebra-lemma-NAK}, nous voyons que
$\Omega_{S/k, \mathfrak q}$ peut être engendré par
$\leq \dim_x X$ éléments. Nous pouvons remplacer $S$ par $S_g$ pour un certain
$g \in S$, $g \not \in \mathfrak q$, de sorte que $\Omega_{S/k}$ soit engendré
par au plus $\dim_x X$ éléments.
Soit $K/k$ une extension de corps algébriquement close telle qu'il existe un
morphisme de $k$-algèbres $\psi : \kappa(\mathfrak q) \to K$.
Considérons $S_K = K \otimes_k S$. Soit $\mathfrak m \subset S_K$ l'idéal
maximal correspondant à la surjection
$$
\xymatrix{
S_K = K \otimes_k S \ar[r] &
K \otimes_k \kappa(\mathfrak q)
\ar[r]^-{\text{id}_K \otimes \psi} &
K.
}
$$
Notons que $\mathfrak m \cap S = \mathfrak q$ ; autrement dit,
$\mathfrak m$ est au-dessus de $\mathfrak q$.
Par le Lemme \ref{algebra-lemma-dimension-at-a-point-preserved-field-extension}, la
dimension de $X_K = \Spec(S_K)$ au point correspondant à $\mathfrak m$ est
$\dim_x X$. Par le
Lemme \ref{algebra-lemma-dimension-closed-point-finite-type-field}, elle est égale à
$\dim((S_K)_{\mathfrak m})$.
Par le Lemme \ref{algebra-lemma-differentials-base-change}, le module des
différentielles de $S_K$ sur $K$ est le changement de base de
$\Omega_{S/k}$ ; il est donc lui aussi engendré par au plus
$\dim_x X = \dim((S_K)_{\mathfrak m})$ éléments. Par le
Lemme \ref{algebra-lemma-characterize-smooth-kbar}, nous voyons que $S_K$ est lisse en
$\mathfrak m$ sur $K$. Le Lemme \ref{algebra-lemma-flat-base-change-locus-smooth}
implique alors que $S$ est lisse en $\mathfrak q$ sur $k$.
Cela démontre (1). De plus, nous savons par le
Lemme \ref{algebra-lemma-characterize-smooth-kbar} que l'anneau local
$(S_K)_{\mathfrak m}$ est régulier. Comme
$S_{\mathfrak q} \to (S_K)_{\mathfrak m}$ est plat, nous concluons du
Lemme \ref{algebra-lemma-flat-under-regular} que $S_{\mathfrak q}$ est régulier.
\end{proof}

\noindent
Le lemme suivant se généralise considérablement (de plusieurs façons
différentes).

\begin{lemma}
\label{algebra-lemma-computation-differential}
Soit $k$ un corps.
Soit $R$ un anneau local noethérien contenant $k$.
Supposons que le corps résiduel $\kappa = R/\mathfrak m$ soit une extension
séparable de type fini de $k$. Alors le morphisme
$$
\text{d} :
\mathfrak m/\mathfrak m^2
\longrightarrow
\Omega_{R/k} \otimes_R \kappa(\mathfrak m)
$$
est injectif.
\end{lemma}

\begin{proof}
Nous pouvons remplacer $R$ par $R/\mathfrak m^2$. Nous pouvons donc supposer
que $\mathfrak m^2 = 0$. Par hypothèse, nous pouvons écrire
$\kappa = k(\overline{x}_1, \ldots, \overline{x}_r, \overline{y})$, où
$\overline{x}_1, \ldots, \overline{x}_r$ est une base de transcendance de
$\kappa$ sur $k$, et $\overline{y}$ est séparable algébrique sur
$k(\overline{x}_1, \ldots, \overline{x}_r)$. Soit $P(T)$ dans
$k(\overline{x}_1, \ldots, \overline{x}_r)[T]$ le polynôme minimal de
$\overline{y}$ sur $k(\overline{x}_1, \ldots, \overline{x}_r)$.
Alors $P(\overline{y}) = 0$, mais $P'(\overline{y}) \not = 0$. Choisissons des
relèvements quelconques $x_i \in R$ des éléments
$\overline{x}_i \in \kappa$. Cela donne un diagramme commutatif
$$
\xymatrix{
R \ar[r] & \kappa \\
& k(\overline{x}_1, \ldots, \overline{x}_r) \ar[lu]^\varphi \ar[u]
}
$$
de $k$-algèbres. Nous voulons prolonger la flèche ascendante gauche
$\varphi$ en un morphisme de $k$-algèbres de $\kappa$ vers $R$. Pour ce faire,
choisissons un $y \in R$ quelconque relevant $\overline{y}$. Pour voir qu'il
définit un morphisme de $k$-algèbres défini sur
$\kappa \cong k(\overline{x}_1, \ldots, \overline{x}_r)[T]/(P)$, il suffit de
montrer que nous pouvons choisir $y$ tel que $P^\varphi(y) = 0$. Sinon, pour
$\delta \in \mathfrak m$, nous calculons
$$
P^\varphi(y + \delta) = P^\varphi(y) + (P')^\varphi(y)\delta
$$
car $\mathfrak m^2 = 0$. Puisque $(P')^\varphi(y)$ est une unité de $R$, nous
voyons que nous pouvons ajuster notre choix comme souhaité.
Cela montre que $R \cong \kappa \oplus \mathfrak m$ comme $k$-algèbres !
Un calcul direct de $\Omega_{\kappa \oplus \mathfrak m/k}$ ou une application
du Lemme \ref{algebra-lemma-differential-seq-split} achève alors la démonstration.
\end{proof}

\begin{lemma}
\label{algebra-lemma-separable-smooth}
Soit $k$ un corps.
Soit $S$ une $k$-algèbre de type fini.
Soit $\mathfrak q \subset S$ un idéal premier.
Supposons $\kappa(\mathfrak q)$ séparable sur $k$.
Les conditions suivantes sont équivalentes :
\begin{enumerate}
\item L'algèbre $S$ est lisse en $\mathfrak q$ sur $k$.
\item L'anneau $S_{\mathfrak q}$ est régulier.
\end{enumerate}
\end{lemma}

\begin{proof}
Notons $R = S_{\mathfrak q}$, son idéal maximal $\mathfrak m$ et son corps
résiduel $\kappa$.
Le Lemme \ref{algebra-lemma-computation-differential} et
\ref{algebra-lemma-differential-seq} montrent qu'il existe une suite exacte courte
$$
0 \to \mathfrak m/\mathfrak m^2 \to
\Omega_{R/k} \otimes_R \kappa \to
\Omega_{\kappa/k} \to 0
$$
Notons que $\Omega_{R/k} = \Omega_{S/k, \mathfrak q}$ ; voir le
Lemme \ref{algebra-lemma-differentials-localize}.
De plus, puisque $\kappa$ est séparable sur $k$, nous avons
$\dim_{\kappa} \Omega_{\kappa/k} = \text{trdeg}_k(\kappa)$.
Nous obtenons donc
$$
\dim_{\kappa} \Omega_{R/k} \otimes_R \kappa
=
\dim_\kappa \mathfrak m/\mathfrak m^2 + \text{trdeg}_k (\kappa)
\geq
\dim R + \text{trdeg}_k (\kappa)
=
\dim_{\mathfrak q} S
$$
(voir le Lemme \ref{algebra-lemma-dimension-at-a-point-finite-type-field} pour la
dernière égalité), avec égalité si et seulement si $R$ est régulier.
Nous concluons donc en appliquant le
Lemme \ref{algebra-lemma-characterize-smooth-over-field}.
\end{proof}
\begin{lemma}
\label{algebra-lemma-characteristic-zero}
Soit $R \to S$ un morphisme de $\mathbf{Q}$-algèbres.
Soit $f \in S$ tel que $\Omega_{S/R} = S \text{d}f \oplus C$ pour un certain
sous-$S$-module $C$. Alors
\begin{enumerate}
\item $f$ n'est pas nilpotent, et
\item si $S$ est un anneau local noethérien, alors $f$ est non-diviseur de zéro
dans $S$.
\end{enumerate}
\end{lemma}

\begin{proof}
Pour $a \in S$, écrivons $\text{d}(a) = \theta(a)\text{d}f + c(a)$, avec
$\theta(a) \in S$ et $c(a) \in C$.
Considérons la $R$-dérivation $S \to S$, $a \mapsto \theta(a)$.
Notons que $\theta(f) = 1$.

\medskip\noindent
Si $f^n = 0$, avec $n > 1$ minimal, alors
$0 = \theta(f^n) = n f^{n - 1}$, ce qui contredit la minimalité de $n$.
Nous concluons que $f$ n'est pas nilpotent.

\medskip\noindent
Supposons $fa = 0$. Si $f$ est une unité, alors $a = 0$ et nous avons terminé.
Supposons que $f$ ne soit pas une unité. Alors
$0 = \theta(fa) = f\theta(a) + a$ par la règle de Leibniz, et donc
$a \in (f)$. Supposons avoir démontré par récurrence que
$fa = 0 \Rightarrow a \in (f^n)$. En écrivant $a = f^nb$, nous obtenons
$0 = \theta(f^{n + 1}b) = (n + 1)f^nb + f^{n + 1}\theta(b)$.
Ainsi $a = f^n b = -f^{n + 1}\theta(b)/(n + 1) \in (f^{n + 1})$.
Comme, dans l'anneau local noethérien $S$, nous avons
$\bigcap (f^n) = 0$ ; voir le
Lemme \ref{algebra-lemma-intersect-powers-ideal-module-zero}, nous avons terminé.
\end{proof}

\noindent
Le résultat suivant est probablement assez inutile dans les applications.

\begin{lemma}
\label{algebra-lemma-characteristic-zero-local-smooth}
Soit $k$ un corps de caractéristique $0$.
Soit $S$ une $k$-algèbre de type fini.
Soit $\mathfrak q \subset S$ un idéal premier.
Les conditions suivantes sont équivalentes :
\begin{enumerate}
\item L'algèbre $S$ est lisse en $\mathfrak q$ sur $k$.
\item Le $S_{\mathfrak q}$-module $\Omega_{S/k, \mathfrak q}$ est libre (de
type fini).
\item L'anneau $S_{\mathfrak q}$ est régulier.
\end{enumerate}
\end{lemma}

\begin{proof}
En caractéristique zéro, toute extension de corps est séparable ; l'équivalence
de (1) et (3) résulte donc du Lemme \ref{algebra-lemma-separable-smooth}.
De plus, (1) implique (2) par définition des algèbres lisses.
Supposons que $\Omega_{S/k, \mathfrak q}$ soit libre sur $S_{\mathfrak q}$.
Nous allons utiliser les notations et observations de la démonstration du
Lemme \ref{algebra-lemma-separable-smooth}. Ainsi, $R = S_{\mathfrak q}$, d'idéal
maximal $\mathfrak m$ et de corps résiduel $\kappa$.
Notre but est de démontrer que $R$ est régulier.

\medskip\noindent
Si $\mathfrak m/\mathfrak m^2 = 0$, alors $\mathfrak m = 0$ et
$R \cong \kappa$. Ainsi $R$ est régulier et nous avons terminé.

\medskip\noindent
Si $\mathfrak m/ \mathfrak m^2 \not = 0$, choisissons un
$f \in \mathfrak m$ quelconque dont l'image dans
$\mathfrak m/ \mathfrak m^2$ n'est pas nulle. Par le
Lemme \ref{algebra-lemma-computation-differential}, nous voyons que $\text{d}f$ a une
image non nulle dans $\Omega_{R/k}/\mathfrak m\Omega_{R/k}$.
Par hypothèse, $\Omega_{R/k} = \Omega_{S/k, \mathfrak q}$ est libre de type
fini et, par le lemme de Nakayama \ref{algebra-lemma-NAK}, nous voyons donc que
$\text{d}f$ engendre un facteur direct. Nous appliquons le
Lemme \ref{algebra-lemma-characteristic-zero} pour en déduire que $f$ est un
non-diviseur de zéro dans $R$.
En outre, le Lemme \ref{algebra-lemma-differential-seq} donne une suite exacte
$$
(f)/(f^2) \to \Omega_{R/k} \otimes_R R/fR \to \Omega_{(R/fR)/k} \to 0
$$
Cela implique que $\Omega_{(R/fR)/k}$ est lui aussi libre de type fini.
Par récurrence, nous voyons donc que $R/fR$ est un anneau local régulier.
Puisque $f \in \mathfrak m$ était un non-diviseur de zéro, nous concluons que
$R$ est régulier ; voir le Lemme \ref{algebra-lemma-regular-mod-x}.
\end{proof}

\begin{example}
\label{algebra-example-characteristic-p}
Le Lemme \ref{algebra-lemma-characteristic-zero-local-smooth} n'est pas valable en
caractéristique $p > 0$.
Les exemples standards sont les morphismes d'anneaux
$$
\mathbf{F}_p \longrightarrow \mathbf{F}_p[x]/(x^p)
$$
dont le module des différentielles est libre mais qui ne sont clairement pas
lisses, et le morphisme d'anneaux ($p > 2$)
$$
\mathbf{F}_p(t) \to \mathbf{F}_p(t)[x, y]/(x^p + y^2 + \alpha)
$$
qui n'est pas lisse en l'idéal premier
$\mathfrak q = (y, x^p + \alpha)$, mais qui est régulier.
\end{example}

\noindent
À l'aide de ce qui précède, nous pouvons caractériser la lissité au point
générique en termes d'extensions de corps.

\begin{lemma}
\label{algebra-lemma-smooth-at-generic-point}
Soit $R \to S$ un morphisme d'anneaux injectif et de type fini, où $R$ et $S$
sont des domaines. Alors $R \to S$ est lisse en $\mathfrak q = (0)$ si et
seulement si l'extension induite $L/K$ des corps de fractions est séparable.
\end{lemma}

\begin{proof}
Supposons $R \to S$ lisse en $(0)$. Nous pouvons remplacer $S$ par $S_g$ pour
un certain $g \in S$ non nul et supposer que $R \to S$ est lisse.
Alors $K \to S \otimes_R K$ est lisse
(Lemme \ref{algebra-lemma-base-change-smooth}). De plus, pour toute extension de corps
$K'/K$, le morphisme d'anneaux $K' \to S \otimes_R K'$ est lui aussi lisse.
Le Lemme \ref{algebra-lemma-characterize-smooth-over-field} montre donc que
$S \otimes_R K'$ est un anneau régulier, en particulier réduit.
Il s'ensuit que $S \otimes_R K$ est géométriquement réduit sur $K$.
Ainsi $L$ est géométriquement réduit sur $K$ ; voir le
Lemme \ref{algebra-lemma-geometrically-reduced-permanence}.
Ainsi $L/K$ est séparable par le
Lemme \ref{algebra-lemma-characterize-separable-field-extensions}.

\medskip\noindent
Réciproquement, supposons que $L/K$ soit séparable.
Nous pouvons supposer que $R \to S$ est de présentation finie ; voir le
Lemme \ref{algebra-lemma-generic-finite-presentation}.
Il suffit de démontrer que $K \to S \otimes_R K$ est lisse en $(0)$ ; voir le
Lemme \ref{algebra-lemma-flat-base-change-locus-smooth}.
Cela résulte du Lemme \ref{algebra-lemma-separable-smooth}, du fait qu'un corps est un
anneau régulier, et de l'hypothèse que $L/K$ est séparable.
\end{proof}







\section{Morphismes d'anneaux lisses dans le cas noethérien}
\label{algebra-section-smooth-Noetherian}

\begin{definition}
\label{algebra-definition-small-extension}
Soit $\varphi : B' \to B$ un morphisme d'anneaux.
Nous disons que $\varphi$ est une {\it petite extension} si $B'$ et $B$ sont des
anneaux locaux artiniens, si $\varphi$ est surjectif et si
$I = \Ker(\varphi)$ est de longueur $1$ comme $B'$-module.
\end{definition}

\noindent
Clairement, cela signifie que $I^2 = 0$ et que $I = (x)$ pour un certain
$x \in B'$ tel que $\mathfrak m' x = 0$, où $\mathfrak m' \subset B'$ est
l'idéal maximal.

\begin{lemma}
\label{algebra-lemma-smooth-test-artinian}
Soit $R \to S$ un morphisme d'anneaux. Soit $\mathfrak q$ un idéal premier de
$S$ au-dessus de $\mathfrak p \subset R$. Supposons que $R$ soit noethérien et
que $R \to S$ soit de type fini.
Les conditions suivantes sont équivalentes :
\begin{enumerate}
\item $R \to S$ est lisse en $\mathfrak q$,
\item pour toute surjection de $R$-algèbres locales
$(B', \mathfrak m') \to (B, \mathfrak m)$ telle que
$\Ker(B' \to B)$ soit de carré nul, et pour tout diagramme commutatif aux
flèches pleines
$$
\xymatrix{
S \ar[r] \ar@{-->}[rd] & B \\
R \ar[r] \ar[u] & B' \ar[u]
}
$$
tel que $\mathfrak q = S \cap \mathfrak m$, il existe une flèche en pointillés
qui rend le diagramme commutatif,
\item comme dans (2), mais où $B' \to B$ parcourt les petites extensions, et
\item comme dans (2), mais où $B' \to B$ parcourt les petites extensions telles
que, de plus, $S \to B$ induise un isomorphisme
$\kappa(\mathfrak q) \cong \kappa(\mathfrak m)$.
\end{enumerate}
\end{lemma}

\begin{proof}
Supposons (1). Cela signifie qu'il existe $g \in S$,
$g \not \in \mathfrak q$, tel que $R \to S_g$ soit lisse. Par la
Proposition \ref{algebra-proposition-smooth-formally-smooth}, nous savons que
$R \to S_g$ est formellement lisse. Notons que, dans tout diagramme comme en
(2), le morphisme $S \to B$ se factorise automatiquement par
$S_{\mathfrak q}$, et a fortiori par $S_g$. La lissité formelle de $S_g$ sur
$R$ nous donne un morphisme $S_g \to B'$ qui s'insère dans un diagramme
semblable avec $S_g$ dans le coin supérieur gauche. Sa composée avec
$S \to S_g$ donne la flèche recherchée. Autrement dit, nous avons montré que
(1) implique (2).

\medskip\noindent
Clairement, (2) implique (3), et (3) implique (4).

\medskip\noindent
Supposons (4). Nous allons montrer que (1) est satisfaite, ce qui achèvera la
démonstration du lemme. Choisissons une présentation
$S = R[x_1, \ldots, x_n]/(f_1, \ldots, f_m)$.
C'est possible, car $S$ est de type fini sur $R$ et donc de présentation finie
(voir le Lemme \ref{algebra-lemma-Noetherian-finite-type-is-finite-presentation}).
Posons $I = (f_1, \ldots, f_m)$.
Considérons le complexe cotangent naïf
$$
\text{d} : I/I^2
\longrightarrow
\bigoplus\nolimits_{j = 1}^n S\text{d}x_j
$$
de cette présentation (voir la section \ref{algebra-section-netherlander}).
Il suffit de montrer que, lorsqu'on localise ce complexe en $\mathfrak q$, le
morphisme devient une injection scindée ; voir le
Lemme \ref{algebra-lemma-smooth-at-point}.
Notons $S' = R[x_1, \ldots, x_n]/I^2$.
Par le Lemme \ref{algebra-lemma-differential-mod-power-ideal}, nous avons
$$
S \otimes_{S'} \Omega_{S'/R} =
S \otimes_{R[x_1, \ldots, x_n]} \Omega_{R[x_1, \ldots, x_n]/R} =
\bigoplus\nolimits_{j = 1}^m S\text{d}x_j.
$$
Ainsi, le morphisme
$$
\text{d} :
I/I^2
\longrightarrow
S \otimes_{S'} \Omega_{S'/R}
$$
est le même que celui du complexe cotangent naïf ci-dessus. En particulier, la
vérité de l'assertion que nous cherchons à démontrer ne dépend que des trois
anneaux $R \to S' \to S$.
Soit $\mathfrak q' \subset R[x_1, \ldots, x_n]$ l'idéal premier correspondant
à $\mathfrak q$. Puisque la localisation commute à la formation des modules de
différentielles (Lemme \ref{algebra-lemma-differentials-localize}), il suffit de montrer
que le morphisme
\begin{equation}
\label{algebra-equation-target-map}
\text{d} :
I_{\mathfrak q'}/I_{\mathfrak q'}^2
\longrightarrow
S_{\mathfrak q} \otimes_{S'_{\mathfrak q'}} \Omega_{S'_{\mathfrak q'}/R}
\end{equation}
provenant de $R \to S'_{\mathfrak q'} \to S_{\mathfrak q}$ soit une injection
scindée.

\medskip\noindent
Soit $N \in \mathbf{N}$ un entier.
Considérons l'anneau
$$
B'_N = S'_{\mathfrak q'} / (\mathfrak q')^N S'_{\mathfrak q'}
= (S'/(\mathfrak q')^N S')_{\mathfrak q'}
$$
et son quotient $B_N = B'_N/IB'_N$. Notons que
$B_N \cong S_{\mathfrak q}/\mathfrak q^NS_{\mathfrak q}$.
Observons que $B'_N$ est un anneau local artinien, puisqu'il est le quotient
d'un anneau local noethérien par une puissance de son idéal maximal.
Considérons une filtration du noyau $I_N$ de $B'_N \to B_N$ par des
sous-$B'_N$-modules
$$
0 \subset J_{N, 1} \subset J_{N, 2} \subset \ldots \subset J_{N, n(N)} = I_N
$$
telle que chaque quotient successif $J_{N, i}/J_{N, i - 1}$ soit de longueur
$1$. (Une telle filtration existe puisque $B'_N$ est artinien.)
Cela donne une suite de petites extensions
$$
B'_N  \to B'_N/J_{N, 1} \to B'_N/J_{N, 2} \to \ldots \to
B'_N/J_{N, n(N)} = B'_N/I_N
= B_N = S_{\mathfrak q}/\mathfrak q^NS_{\mathfrak q}
$$
En appliquant successivement la condition (4) à ces petites extensions, à
partir du morphisme $S \to B_N$, nous voyons qu'il existe un diagramme
commutatif
$$
\xymatrix{
S \ar[r] \ar[rd] & B_N  \\
R \ar[r] \ar[u] & B'_N \ar[u]
}
$$
Clairement, le morphisme d'anneaux $S \to B'_N$ se factorise en
$S \to S_{\mathfrak q} \to B'_N$, où
$S_{\mathfrak q} \to B'_N$ est un homomorphisme local d'anneaux locaux.
De plus, comme l'idéal maximal de $B'_N$, élevé à la puissance $N$, est nul,
nous concluons que $S_{\mathfrak q} \to B'_N$ se factorise par
$S_{\mathfrak q}/(\mathfrak q)^NS_{\mathfrak q} = B_N$. Autrement dit, nous
avons montré que, pour tout $N \in \mathbf{N}$, la surjection de $R$-algèbres
$B'_N \to B_N$ admet un scindage.

\medskip\noindent
Considérons la présentation
$$
I_N \to B_N \otimes_{B'_N} \Omega_{B'_N/R} \to \Omega_{B_N/R} \to 0
$$
provenant de la surjection $B'_N \to B_N$ de noyau $I_N$ (voir le
Lemme \ref{algebra-lemma-differential-seq}). D'après ce qui précède, le morphisme de
$R$-algèbres $B'_N \to B_N$ admet un inverse à droite. Ainsi, par le
Lemme \ref{algebra-lemma-differential-seq-split}, la suite ci-dessus est exacte
scindée ! Pour tout $N$, le morphisme
$$
I_N \longrightarrow B_N \otimes_{B'_N} \Omega_{B'_N/R}
$$
est donc une injection scindée. La suite de la démonstration s'obtient en
explicitant exactement ce que cela signifie. Notons que
$$
I_N = I_{\mathfrak q'}/
(I_{\mathfrak q'}^2 + (\mathfrak q')^N \cap I_{\mathfrak q'})
$$
Par Artin--Rees (Lemme \ref{algebra-lemma-Artin-Rees}), il existe $c \geq 0$ tel que
$$
S_{\mathfrak q}/\mathfrak q^{N - c}S_{\mathfrak q}
\otimes_{S_{\mathfrak q}} I_N =
S_{\mathfrak q}/\mathfrak q^{N - c}S_{\mathfrak q}
\otimes_{S_{\mathfrak q}}
I_{\mathfrak q'}/I_{\mathfrak q'}^2
$$
pour tout $N \geq c$ (ces produits tensoriels ne sont qu'une manière élaborée
de diviser par $\mathfrak q^{N - c}$). Nous pouvons bien sûr supposer
$c \geq 1$.
Par le Lemme \ref{algebra-lemma-differential-mod-power-ideal}, nous voyons que
$$
S'_{\mathfrak q'}/(\mathfrak q')^{N - c}S'_{\mathfrak q'}
\otimes_{S'_{\mathfrak q'}} \Omega_{B'_N/R} =
S'_{\mathfrak q'}/(\mathfrak q')^{N - c}S'_{\mathfrak q'}
\otimes_{S'_{\mathfrak q'}} \Omega_{S'_{\mathfrak q'}/R}
$$
nous pouvons encore tensoriser cette égalité par
$B_N = S_{\mathfrak q}/\mathfrak q^N$ pour voir que
$$
S_{\mathfrak q}/\mathfrak q^{N - c}S_{\mathfrak q}
\otimes_{S'_{\mathfrak q'}} \Omega_{B'_N/R} =
S_{\mathfrak q}/\mathfrak q^{N - c}S_{\mathfrak q}
\otimes_{S'_{\mathfrak q'}} \Omega_{S'_{\mathfrak q'}/R}.
$$
Puisqu'une injection scindée reste une injection scindée après tensorisation
par n'importe quel module, nous voyons que
$$
S_{\mathfrak q}/\mathfrak q^{N - c}S_{\mathfrak q}
\otimes_{S_{\mathfrak q}}
(\ref{algebra-equation-target-map}) =
S_{\mathfrak q}/\mathfrak q^{N - c}S_{\mathfrak q}
\otimes_{S_{\mathfrak q}/\mathfrak q^N S_{\mathfrak q}}
(I_N \longrightarrow B_N \otimes_{B'_N} \Omega_{B'_N/R})
$$
est une injection scindée pour tout $N \geq c$. Par le
Lemme \ref{algebra-lemma-split-injection-after-completion}, nous voyons que
(\ref{algebra-equation-target-map}) est une injection scindée. Cela achève la
démonstration.
\end{proof}











\section{Aperçu des résultats sur les morphismes d'anneaux lisses}
\label{algebra-section-smooth-overview}

\noindent
Voici une liste de résultats sur les morphismes d'anneaux lisses que nous avons
démontrés dans les sections précédentes. Pour des énoncés et des définitions
plus précis, consulter les références indiquées.
\begin{enumerate}
\item Un morphisme d'anneaux $R \to S$ est lisse s'il est de présentation finie
et si le complexe cotangent naïf de $S/R$ est quasi-isomorphe à un $S$-module
projectif de type fini en degré $0$ ; voir la
Définition \ref{algebra-definition-smooth}.
\item Si $S$ est lisse sur $R$, alors $\Omega_{S/R}$ est un $S$-module
projectif de type fini ; voir la discussion qui suit la
Définition \ref{algebra-definition-smooth}.
\item La propriété d'être lisse est locale sur $S$ ; voir le
Lemme \ref{algebra-lemma-locally-smooth}.
\item La propriété d'être lisse est stable par changement de base ; voir le
Lemme \ref{algebra-lemma-base-change-smooth}.
\item La propriété d'être lisse est stable par composition ; voir le
Lemme \ref{algebra-lemma-compose-smooth}.
\item Un morphisme d'anneaux lisse est syntomique, donc en particulier plat ;
voir le Lemme \ref{algebra-lemma-smooth-syntomic}.
\item Un morphisme d'anneaux plat et de présentation finie dont les anneaux
fibres sont lisses est lisse ; voir le Lemme \ref{algebra-lemma-flat-fibre-smooth}.
\item Un morphisme d'anneaux de présentation finie $R \to S$ est lisse si et
seulement s'il est formellement lisse ; voir la
Proposition \ref{algebra-proposition-smooth-formally-smooth}.
\item Si $R \to S$ est un morphisme d'anneaux de type fini et si $R$ est
noethérien, alors, pour vérifier que $R \to S$ est lisse, il suffit de vérifier
la propriété de relèvement de la lissité formelle le long des petites
extensions d'anneaux locaux artiniens ; voir le
Lemme \ref{algebra-lemma-smooth-test-artinian}.
\item Un morphisme d'anneaux lisse $R \to S$ est le changement de base d'un
morphisme d'anneaux lisse $R_0 \to S_0$, où $R_0$ est de type fini sur
$\mathbf{Z}$ ; voir le
Lemme \ref{algebra-lemma-finite-presentation-fs-Noetherian}.
\item La formation de l'ensemble des points où un morphisme d'anneaux est lisse
commute au changement de base plat ; voir le
Lemme \ref{algebra-lemma-flat-base-change-locus-smooth}.
\item Si $S$ est de type fini sur un corps algébriquement clos $k$, et si
$\mathfrak m \subset S$ est un idéal maximal, alors les conditions suivantes
sont équivalentes :
\begin{enumerate}
\item $S$ est lisse sur $k$ dans un voisinage de $\mathfrak m$,
\item $S_{\mathfrak m}$ est un anneau local régulier,
\item $\dim(S_{\mathfrak m}) =
\dim_{\kappa(m)} \Omega_{S/k} \otimes_S \kappa(\mathfrak m)$.
\end{enumerate}
voir le Lemme \ref{algebra-lemma-characterize-smooth-kbar}.
\item Si $S$ est de type fini sur un corps $k$, et si
$\mathfrak q \subset S$ est un idéal premier, alors les conditions suivantes
sont équivalentes :
\begin{enumerate}
\item $S$ est lisse sur $k$ dans un voisinage de $\mathfrak q$,
\item $\dim_{\mathfrak q}(S/k) =
\dim_{\kappa(\mathfrak q)} \Omega_{S/k} \otimes_S \kappa(\mathfrak q)$.
\end{enumerate}
voir le Lemme \ref{algebra-lemma-characterize-smooth-over-field}.
\item Si $S$ est lisse sur un corps, alors tous ses anneaux locaux sont
réguliers ; voir le Lemme \ref{algebra-lemma-characterize-smooth-over-field}.
\item Si $S$ est de type fini sur un corps $k$, si
$\mathfrak q \subset S$ est un idéal premier, si l'extension de corps
$\kappa(\mathfrak q)/k$ est séparable et si $S_{\mathfrak q}$ est régulier,
alors $S$ est lisse sur $k$ en $\mathfrak q$ ; voir le
Lemme \ref{algebra-lemma-separable-smooth}.
\item Si $S$ est de type fini sur un corps $k$, si $k$ est de caractéristique
$0$, si $\mathfrak q \subset S$ est un idéal premier et si
$\Omega_{S/k, \mathfrak q}$ est libre, alors $S$ est lisse sur $k$ en
$\mathfrak q$ ; voir le Lemme \ref{algebra-lemma-characteristic-zero-local-smooth}.
\end{enumerate}
Certains de ces résultats ont été démontrés à l'aide de la notion de morphisme
d'anneaux lisse standard ; voir la Définition \ref{algebra-definition-standard-smooth}.
C'est l'analogue, pour les morphismes syntomiques, de la notion de morphisme
d'intersection complète globale relative. C'est aussi la façon la plus simple
de construire des exemples.




















\section{Morphismes d'anneaux étales}
\label{algebra-section-etale}

\noindent
Un morphisme d'anneaux étale est un morphisme d'anneaux lisse dont la dimension
relative est nulle. Cela équivaut à la définition un peu plus directe suivante.

\begin{definition}
\label{algebra-definition-etale}
Soit $R \to S$ un morphisme d'anneaux. Nous disons que $R \to S$ est
{\it étale} s'il est de présentation finie et si le complexe cotangent naïf
$\NL_{S/R}$ est quasi-isomorphe à zéro : cela signifie que
$H_1(\NL_{S/R}) = 0$ et $\Omega_{S/R} = 0$. Étant donné un idéal premier
$\mathfrak q$ de $S$, nous disons que $R \to S$ est
{\it étale en $\mathfrak q$} s'il existe $g \in S$,
$g \not \in \mathfrak q$, tel que $R \to S_g$ soit étale.
\end{definition}

\noindent
En particulier, nous voyons que $\Omega_{S/R} = 0$ si $S$ est étale sur $R$.
Si $R \to S$ est lisse, alors $R \to S$ est étale si et seulement si
$\Omega_{S/R} = 0$.
Nos résultats sur les morphismes d'anneaux lisses fournissent automatiquement
toute une série de résultats pour les morphismes étales. Nous les résumons dans
le Lemme \ref{algebra-lemma-etale} ci-dessous. Mais auparavant, nous démontrons que
{\it tout} morphisme d'anneaux étale est lisse standard.

\begin{lemma}
\label{algebra-lemma-etale-standard-smooth}
Tout morphisme d'anneaux étale est lisse standard. Plus précisément, si
$R \to S$ est étale, alors il existe une présentation
$S = R[x_1, \ldots, x_n]/(f_1, \ldots, f_n)$ telle que l'image de
$\det(\partial f_j/\partial x_i)$ soit inversible dans $S$.
\end{lemma}

\begin{proof}
Soit $R \to S$ étale. Choisissons une présentation
$S = R[x_1, \ldots, x_n]/I$. Comme $R \to S$ est étale, nous savons que
$$
\text{d} :
I/I^2
\longrightarrow
\bigoplus\nolimits_{i = 1, \ldots, n} S\text{d}x_i
$$
est un isomorphisme ; en particulier, $I/I^2$ est un $S$-module libre.
Ainsi, par le Lemme \ref{algebra-lemma-huber}, nous pouvons supposer (après avoir
éventuellement changé la présentation) que $I = (f_1, \ldots, f_c)$ et que les
classes $f_i \bmod I^2$ forment une base de $I/I^2$. Le fait que le morphisme
affiché ci-dessus soit un isomorphisme implique immédiatement que $c = n$ et
que $\det(\partial f_j/\partial x_i)$ est inversible dans $S$.
\end{proof}

\begin{lemma}
\label{algebra-lemma-etale}
Résultats sur les morphismes d'anneaux étales.
\begin{enumerate}
\item Le morphisme d'anneaux $R \to R_f$ est étale pour tout anneau $R$ et tout
$f \in R$.
\item Les composées de morphismes d'anneaux étales sont étales.
\item Un changement de base d'un morphisme d'anneaux étale est étale.
\item La propriété d'être étale est locale : étant donnés un morphisme
d'anneaux $R \to S$ et des éléments $g_1, \ldots, g_m \in S$ qui engendrent
l'idéal unité, tels que $R \to S_{g_j}$ soit étale pour
$j = 1, \ldots, m$, alors $R \to S$ est étale.
\item Étant donné $R \to S$ de présentation finie et un morphisme d'anneaux
plat $R \to R'$, posons $S' = R' \otimes_R S$. L'ensemble des idéaux premiers
où $R' \to S'$ est étale est l'image réciproque, par
$\Spec(S') \to \Spec(S)$, de l'ensemble des idéaux premiers où $R \to S$ est
étale.
\item Un morphisme d'anneaux étale est syntomique, donc en particulier plat.
\item Si $S$ est de type fini sur un corps $k$, alors $S$ est étale sur $k$ si
et seulement si $\Omega_{S/k} = 0$.
\item Tout morphisme d'anneaux étale $R \to S$ est le changement de base d'un
morphisme d'anneaux étale $R_0 \to S_0$, où $R_0$ est de type fini sur
$\mathbf{Z}$.
\item Soit $A = \colim A_i$ une colimite filtrante d'anneaux.
Soit $A \to B$ un morphisme d'anneaux étale. Alors il existe un morphisme
d'anneaux étale $A_i \to B_i$ pour un certain $i$, tel que
$B \cong A \otimes_{A_i} B_i$.
\item Soit $A$ un anneau. Soit $S$ une partie multiplicative de $A$.
Soit $S^{-1}A \to B'$ étale. Alors il existe un morphisme d'anneaux étale
$A \to B$ tel que $B' \cong S^{-1}B$.
\item Soit $A$ un anneau. Soit $B = B' \times B''$ un produit de
$A$-algèbres. Alors $B$ est étale sur $A$ si et seulement si $B'$ et $B''$
sont toutes deux étales sur $A$.
\end{enumerate}
\end{lemma}

\begin{proof}
Dans chaque cas, nous utilisons le résultat correspondant pour les morphismes
d'anneaux lisses, auquel nous ajoutons un petit argument pour montrer que
$\Omega_{S/R}$ est nul.

\medskip\noindent
Démonstration de (1). Le morphisme d'anneaux $R \to R_f$ est lisse et
$\Omega_{R_f/R} = 0$.

\medskip\noindent
Démonstration de (2). La composée $A \to C$ des morphismes lisses
$A \to B$ et $B \to C$ est lisse ; voir le
Lemme \ref{algebra-lemma-compose-smooth}. Par le
Lemme \ref{algebra-lemma-exact-sequence-differentials}, nous voyons que
$\Omega_{C/A}$ est nul, puisque $\Omega_{C/B}$ et $\Omega_{B/A}$ sont tous deux
nuls.

\medskip\noindent
Démonstration de (3). Soit $R \to S$ étale, et soit $R \to R'$ quelconque.
Alors $R' \to S' = R' \otimes_R S$ est lisse ; voir le
Lemme \ref{algebra-lemma-base-change-smooth}. Puisque
$\Omega_{S'/R'} = S' \otimes_S \Omega_{S/R}$ par le
Lemme \ref{algebra-lemma-differentials-base-change}, nous concluons que
$\Omega_{S'/R'} = 0$. Ainsi $R' \to S'$ est étale.

\medskip\noindent
Démonstration de (4). Supposons les hypothèses de (4). Par le
Lemme \ref{algebra-lemma-locally-smooth}, nous voyons que $R \to S$ est lisse.
Nous savons également que
$\Omega_{S_{g_i}/R} = (\Omega_{S/R})_{g_i} = 0$ pour tout $i$.
Alors $\Omega_{S/R} = 0$ ; voir le Lemme \ref{algebra-lemma-cover}.

\medskip\noindent
Démonstration de (5). Le résultat pour les morphismes lisses est le
Lemme \ref{algebra-lemma-flat-base-change-locus-smooth}.
Dans la démonstration de ce lemme, nous avons utilisé que
$\NL_{S/R} \otimes_S S'$ est homotopiquement équivalent à $\NL_{S'/R'}$.
Nous sommes donc ramenés à montrer que, si $M$ est un $S$-module de
présentation finie, l'ensemble des idéaux premiers $\mathfrak q'$ de $S'$ tels
que $(M \otimes_S S')_{\mathfrak q'} = 0$ est l'image réciproque de l'ensemble
des idéaux premiers $\mathfrak q$ de $S$ tels que $M_{\mathfrak q} = 0$.
Cela résulte du Lemme \ref{algebra-lemma-support-base-change}.

\medskip\noindent
Démonstration de (6). Cela résulte directement du résultat correspondant pour
les morphismes d'anneaux lisses (Lemme \ref{algebra-lemma-smooth-syntomic}).

\medskip\noindent
Démonstration de (7). Cela résulte du
Lemme \ref{algebra-lemma-characterize-smooth-over-field} et des définitions.

\medskip\noindent
Démonstration de (8). Le
Lemme \ref{algebra-lemma-finite-presentation-fs-Noetherian} donne le résultat pour les
morphismes d'anneaux lisses. Le morphisme d'anneaux lisse obtenu
$R_0 \to S_0$ satisfait aux hypothèses du
Lemme \ref{algebra-lemma-relative-dimension-CM} ; nous pouvons donc remplacer $S_0$ par
le facteur de dimension relative $0$ sur $R_0$.

\medskip\noindent
Démonstration de (9). Cela résulte de (8), car $R_0 \to A$ se factorisera par
$A_i$ pour un certain $i$, d'après le
Lemme \ref{algebra-lemma-characterize-finite-presentation}.

\medskip\noindent
Démonstration de (10). Cela résulte de (9), (1) et (2), puisque $S^{-1}A$ est
une colimite filtrante de localisations principales de $A$.

\medskip\noindent
Démonstration de (11). Utilisons le Lemme \ref{algebra-lemma-product-smooth} pour
obtenir le résultat pour la lissité, puis le fait que $\Omega_{B/A}$ est nul si
et seulement si $\Omega_{B'/A}$ et $\Omega_{B''/A}$ sont tous deux nuls.
\end{proof}


















\noindent
Examinons maintenant plus en détail ce que signifie être étale sur un corps.

\begin{lemma}
\label{algebra-lemma-etale-over-field}
Soit $k$ un corps. Un morphisme d'anneaux $k \to S$ est étale si et seulement
si $S$ est isomorphe, comme $k$-algèbre, à un produit fini d'extensions
séparables finies de $k$.
\end{lemma}

\begin{proof}
Nous utiliserons sans autre mention le fait suivant : si
$S = S_1 \times \ldots \times S_n$ est un produit fini de $k$-algèbres, alors
$S$ est étale sur $k$ si et seulement si chaque $S_i$ est étale sur $k$.
Voir le Lemme \ref{algebra-lemma-etale}, partie (11).

\medskip\noindent
Si $k'/k$ est une extension de corps séparable finie, nous pouvons écrire
$k' = k(\alpha) \cong k[x]/(f)$. Ici, $f$ est le polynôme minimal de l'élément
$\alpha$. Comme $k'$ est séparable sur $k$, nous avons
$\gcd(f, f') = 1$. Cela implique que
$\text{d} : k'\cdot f \to k' \cdot \text{d}x$ est un isomorphisme.
Ainsi $k \to k'$ est étale. Par conséquent, si $S$ est un produit fini
d'extensions séparables finies de $k$, alors $S$ est étale sur $k$.

\medskip\noindent
Réciproquement, supposons que $k \to S$ soit étale. Alors $S$ est lisse sur
$k$ et $\Omega_{S/k} = 0$. Par le
Lemme \ref{algebra-lemma-characterize-smooth-over-field}, nous voyons que
$\dim_\mathfrak m \Spec(S) = 0$ pour tout idéal maximal $\mathfrak m$ de $S$.
Ainsi $\dim(S) = 0$. Par la Proposition \ref{algebra-proposition-dimension-zero-ring},
nous obtenons que $S$ est un produit fini d'anneaux locaux artiniens.
Par le Lemme \ref{algebra-lemma-characterize-smooth-over-field}, déjà utilisé, ces
anneaux locaux sont des corps. Nous pouvons donc supposer que $S = k'$ est un
corps. Par le théorème des zéros de Hilbert
(Théorème \ref{algebra-theorem-nullstellensatz}), nous voyons que l'extension $k'/k$
est finie. La lissité de $k \to k'$ implique, par le
Lemme \ref{algebra-lemma-smooth-at-generic-point}, que $k'/k$ est une extension
séparable, ce qui achève la démonstration.
\end{proof}

\begin{lemma}
\label{algebra-lemma-etale-at-prime}
Soit $R \to S$ un morphisme d'anneaux.
Soit $\mathfrak q \subset S$ un idéal premier au-dessus de $\mathfrak p$ dans
$R$. Si $S/R$ est étale en $\mathfrak q$, alors
\begin{enumerate}
\item nous avons $\mathfrak p S_{\mathfrak q} = \mathfrak qS_{\mathfrak q}$,
qui est l'idéal maximal de l'anneau local $S_{\mathfrak q}$, et
\item l'extension de corps $\kappa(\mathfrak q)/\kappa(\mathfrak p)$ est
séparable finie.
\end{enumerate}
\end{lemma}

\begin{proof}
Nous pouvons d'abord remplacer $S$ par $S_g$ pour un certain
$g \in S$, $g \not \in \mathfrak q$, et supposer que $R \to S$ est étale.
Le lemme résulte alors du Lemme \ref{algebra-lemma-etale-over-field}, en explicitant le
fait que $S \otimes_R \kappa(\mathfrak p)$ est étale sur
$\kappa(\mathfrak p)$.
\end{proof}

\begin{lemma}
\label{algebra-lemma-etale-quasi-finite}
Un morphisme d'anneaux étale est quasi-fini.
\end{lemma}

\begin{proof}
Soit $R \to S$ un morphisme d'anneaux étale. Par définition, $R \to S$ est de
type fini. Pour tout idéal premier $\mathfrak p \subset R$, l'anneau fibre
$S \otimes_R \kappa(\mathfrak p)$ est étale sur $\kappa(\mathfrak p)$ ; c'est
donc un produit fini de corps séparables finis sur $\kappa(\mathfrak p)$, et il
est en particulier fini sur $\kappa(\mathfrak p)$. Ainsi $R \to S$ est
quasi-fini par le Lemme \ref{algebra-lemma-quasi-finite}.
\end{proof}

\begin{lemma}
\label{algebra-lemma-characterize-etale}
Soit $R \to S$ un morphisme d'anneaux. Soit $\mathfrak q$ un idéal premier de
$S$ au-dessus d'un idéal premier $\mathfrak p$ de $R$. Si
\begin{enumerate}
\item $R \to S$ est de présentation finie,
\item $R_{\mathfrak p} \to S_{\mathfrak q}$ est plat,
\item $\mathfrak p S_{\mathfrak q}$ est l'idéal maximal de l'anneau local
$S_{\mathfrak q}$, et
\item l'extension de corps $\kappa(\mathfrak q)/\kappa(\mathfrak p)$ est
séparable finie,
\end{enumerate}
alors $R \to S$ est étale en $\mathfrak q$.
\end{lemma}

\begin{proof}
Appliquons le Lemme \ref{algebra-lemma-isolated-point-fibre} pour trouver
$g \in S$, $g \not \in \mathfrak q$, tel que $\mathfrak q$ soit l'unique idéal
premier de $S_g$ au-dessus de $\mathfrak p$.
Nous remplaçons $S$ par $S_g$. Alors $S \otimes_R \kappa(\mathfrak p)$ possède
un unique idéal premier ; c'est donc un anneau local, et il est donc égal à
$S_{\mathfrak q}/\mathfrak pS_{\mathfrak q}
\cong \kappa(\mathfrak q)$.
Par le Lemme \ref{algebra-lemma-flat-fibre-smooth}, il existe
$g \in S$, $g \not \in \mathfrak q$, tel que $R \to S_g$ soit lisse.
Remplaçons encore $S$ par $S_g$ ; nous pouvons supposer que $R \to S$ est
lisse. Par le Lemme \ref{algebra-lemma-smooth-syntomic}, nous pouvons même supposer que
$R \to S$ est lisse standard, disons
$S = R[x_1, \ldots, x_n]/(f_1, \ldots, f_c)$. Puisque
$S \otimes_R \kappa(\mathfrak p) = \kappa(\mathfrak q)$ est de dimension $0$,
nous concluons que $n = c$, c'est-à-dire que $R \to S$ est étale.
\end{proof}

\noindent
Voici un phénomène tout à fait nouveau.

\begin{lemma}
\label{algebra-lemma-map-between-etale}
Soient $R \to S$ et $R \to S'$ étales.
Alors tout morphisme de $R$-algèbres $S' \to S$ est étale.
\end{lemma}

\begin{proof}
Notons tout d'abord que $S' \to S$ est de présentation finie par le
Lemme \ref{algebra-lemma-compose-finite-type}.
Soit $\mathfrak q \subset S$ un idéal premier au-dessus des idéaux premiers
$\mathfrak q' \subset S'$ et $\mathfrak p \subset R$.
Par le Lemme \ref{algebra-lemma-etale-at-prime}, le morphisme d'anneaux
$S'_{\mathfrak q'}/\mathfrak p S'_{\mathfrak q'} \to
S_{\mathfrak q}/\mathfrak p S_{\mathfrak q}$ est un morphisme entre extensions
séparables finies de $\kappa(\mathfrak p)$. Il est en particulier plat.
Le Lemme \ref{algebra-lemma-criterion-flatness-fibre} montre donc que
$S'_{\mathfrak q'} \to S_{\mathfrak q}$ est plat. Ainsi $S' \to S$ est plat.
De plus, ce qui précède montre aussi que $\mathfrak q'S_{\mathfrak q}$ est
l'idéal maximal de $S_{\mathfrak q}$ et que l'extension de corps résiduels de
$S'_{\mathfrak q'} \to S_{\mathfrak q}$ est séparable finie.
Le Lemme \ref{algebra-lemma-characterize-etale} permet donc de conclure que
$S' \to S$ est étale en $\mathfrak q$. Comme la propriété d'être étale est
locale (voir le Lemme \ref{algebra-lemma-etale}), nous avons terminé.
\end{proof}

\begin{lemma}
\label{algebra-lemma-surjective-flat-finitely-presented}
Soit $\varphi : R \to S$ un morphisme d'anneaux. Si $R \to S$ est surjectif,
plat et de présentation finie, alors il existe un idempotent $e \in R$ tel que
$S = R_e$.
\end{lemma}

\begin{proof}[Première démonstration]
Soit $I$ le noyau de $\varphi$.
Le Lemme \ref{algebra-lemma-finite-presentation-independent} montre que $I$ est de type
fini, puisque $\varphi$ est de présentation finie.
De plus, comme $S$ est plat sur $R$, tensoriser la suite exacte
$0 \to I \to R \to S \to 0$ sur $R$ avec $S$ donne $I/I^2 = 0$.
Nous concluons maintenant par le
Lemme \ref{algebra-lemma-ideal-is-squared-union-connected}.
\end{proof}

\begin{proof}[Seconde démonstration]
Puisque $\Spec(S) \to \Spec(R)$ est un homéomorphisme sur un sous-ensemble
fermé (voir le Lemme \ref{algebra-lemma-spec-closed}) et est ouvert (voir la
Proposition \ref{algebra-proposition-fppf-open}), nous voyons que son image est $D(e)$
pour un certain idempotent $e \in R$ (voir le
Lemme \ref{algebra-lemma-disjoint-decomposition}). Ainsi $R_e \to S$ induit une
bijection sur les spectres. Ce morphisme induit maintenant un isomorphisme sur
tous les anneaux locaux, par exemple par les
Lemmes \ref{algebra-lemma-finite-flat-local} et \ref{algebra-lemma-NAK}.
Il s'ensuit que $R_e \to S$ est aussi injectif ; voir, par exemple, le
Lemme \ref{algebra-lemma-characterize-zero-local}.
\end{proof}

\begin{lemma}
\label{algebra-lemma-lift-etale}
\begin{slogan}
Les morphismes d'anneaux étales se relèvent le long des surjections d'anneaux
\end{slogan}
Soit $R$ un anneau et soit $I \subset R$ un idéal.
Soit $R/I \to \overline{S}$ un morphisme d'anneaux étale.
Alors il existe un morphisme d'anneaux étale $R \to S$ tel que
$\overline{S} \cong S/IS$ comme $R/I$-algèbres.
\end{lemma}

\begin{proof}
Par le Lemme \ref{algebra-lemma-etale-standard-smooth}, nous pouvons écrire
$\overline{S} =
(R/I)[x_1, \ldots, x_n]/(\overline{f}_1, \ldots, \overline{f}_n)$
comme dans la Définition \ref{algebra-definition-standard-smooth}, avec
$\overline{\Delta} =
\det(\frac{\partial \overline{f}_i}{\partial x_j})_{i, j = 1, \ldots, n}$
inversible dans $\overline{S}$. Il suffit de choisir des relèvements $f_i$ et
de poser
$S = R[x_1, \ldots, x_n, x_{n+1}]/(f_1, \ldots, f_n, x_{n + 1}\Delta - 1)$,
où $\Delta = \det(\frac{\partial f_i}{\partial x_j})_{i, j = 1, \ldots, n}$,
comme dans l'Exemple \ref{algebra-example-make-standard-smooth}.
Cela démontre le lemme.
\end{proof}















\begin{lemma}
\label{algebra-lemma-lift-etale-infinitesimal}
Considérons un diagramme commutatif
$$
\xymatrix{
0 \ar[r] &
J \ar[r] &
B' \ar[r] &
B \ar[r] & 0 \\
0 \ar[r] &
I \ar[r] \ar[u] &
A' \ar[r] \ar[u] &
A \ar[r] \ar[u] & 0
}
$$
à lignes exactes, où $B' \to B$ et $A' \to A$ sont des morphismes d'anneaux
surjectifs dont les noyaux sont des idéaux de carré nul. Si $A \to B$ est étale
et si $J = I \otimes_A B$, alors $A' \to B'$ est étale.
\end{lemma}

\begin{proof}
Par le Lemme \ref{algebra-lemma-lift-etale}, il existe un morphisme d'anneaux étale
$A' \to C$ tel que $C/IC = B$. Alors $A' \to C$ est formellement lisse (par la
Proposition \ref{algebra-proposition-smooth-formally-smooth}) ; nous obtenons donc un
morphisme de $A'$-algèbres $\varphi : C \to B'$.
Comme $A' \to C$ est plat, nous avons
$I \otimes_A B = I \otimes_A C/IC = IC$.
Ainsi, l'hypothèse $J = I \otimes_A B$ implique que $\varphi$ induit un
isomorphisme $IC \to J$ et un isomorphisme $C/IC \to B'/IB'$ ; par conséquent,
$\varphi$ est un isomorphisme.
\end{proof}

\begin{example}
\label{algebra-example-factor-polynomials-etale}
Soient $n , m \geq 1$ des entiers. Considérons le morphisme d'anneaux
\begin{eqnarray*}
R = \mathbf{Z}[a_1, \ldots, a_{n + m}]
& \longrightarrow &
S = \mathbf{Z}[b_1, \ldots, b_n, c_1, \ldots, c_m] \\
a_1 & \longmapsto & b_1 + c_1 \\
a_2 & \longmapsto & b_2 + b_1 c_1 + c_2 \\
\ldots & \ldots & \ldots \\
a_{n + m} & \longmapsto & b_n c_m
\end{eqnarray*}
de l'Exemple \ref{algebra-example-factor-polynomials}.
Écrivons symboliquement
$$
S = R[b_1, \ldots, c_m]/(\{a_k(b_i, c_j) - a_k\}_{k = 1, \ldots, n + m})
$$
où, par exemple, $a_1(b_i, c_j) = b_1 + c_1$.
La matrice des dérivées partielles est
$$
\left(
\begin{matrix}
1 & c_1 & \ldots & c_m & 0 & \ldots & \ldots & 0 \\
0 & 1 & c_1 & \ldots & c_m & 0 & \ldots & 0 \\
\ldots & \ldots & \ldots & \ldots & \ldots & \ldots & \ldots & \ldots \\
0 & \ldots & 0 & 1 & c_1 & c_2 & \ldots & c_m \\
1 & b_1 & \ldots & b_{n - 1} & b_n & 0 & \ldots & 0 \\
0 & 1 & b_1 & \ldots & b_{n - 1} & b_n & \ldots & 0 \\
\ldots & \ldots & \ldots & \ldots & \ldots & \ldots & \ldots & \ldots \\
0 & \ldots & \ldots & 0 & 1 & b_1 & \ldots & b_n
\end{matrix}
\right)
$$
Le déterminant $\Delta$ de cette matrice est mieux connu sous le nom de
{\it résultant} des polynômes
$g = x^n + b_1 x^{n - 1} + \ldots + b_n$ et
$h = x^m + c_1 x^{m - 1} + \ldots + c_m$, et la matrice ci-dessus est appelée
{\it matrice de Sylvester} associée à $g, h$.
Sous forme de formule, $\Delta = \text{Res}_x(g, h)$. La matrice de Sylvester
est la transposée de la matrice du morphisme linéaire
\begin{eqnarray*}
S[x]_{< m} \oplus S[x]_{< n} & \longrightarrow & S[x]_{< n + m} \\
a \oplus b & \longmapsto & ag + bh
\end{eqnarray*}
Soit $\mathfrak q \subset S$ un idéal premier quelconque. D'après ce qui
précède, les conditions suivantes sont équivalentes :
\begin{enumerate}
\item $R \to S$ est étale en $\mathfrak q$,
\item $\Delta = \text{Res}_x(g, h) \not \in \mathfrak q$,
\item les images $\overline{g}, \overline{h} \in \kappa(\mathfrak q)[x]$ des
polynômes $g, h$ sont premiers entre eux dans $\kappa(\mathfrak q)[x]$.
\end{enumerate}
L'équivalence de (2) et (3) vaut parce que l'image de la matrice de Sylvester
dans $\text{Mat}(n + m, \kappa(\mathfrak q))$ possède un noyau si et seulement
si les polynômes $\overline{g}, \overline{h}$ ont un facteur commun.
Nous concluons que le morphisme d'anneaux
$$
R \longrightarrow S[\frac{1}{\Delta}] = S[\frac{1}{\text{Res}_x(g, h)}]
$$
est étale.
\end{example}


\begin{lemma}
\label{algebra-lemma-factor-mod-lift-etale}
Soit $R$ un anneau. Soit $f \in R[x]$ un polynôme unitaire. Soit
$\mathfrak p$ un idéal premier de $R$. Soit
$f \bmod \mathfrak p = \overline{g} \overline{h}$ une factorisation de
l'image de $f$ dans $\kappa(\mathfrak p)[x]$.
Si $\gcd(\overline{g}, \overline{h}) = 1$, alors il existe
\begin{enumerate}
\item un morphisme d'anneaux étale $R \to R'$,
\item un idéal premier $\mathfrak p' \subset R'$ au-dessus de $\mathfrak p$, et
\item une factorisation $f = g h$ dans $R'[x]$
\end{enumerate}
tels que
\begin{enumerate}
\item $\kappa(\mathfrak p) = \kappa(\mathfrak p')$,
\item $\overline{g} = g \bmod \mathfrak p'$,
$\overline{h} = h \bmod \mathfrak p'$, et
\item les polynômes $g, h$ engendrent l'idéal unité dans $R'[x]$.
\end{enumerate}
\end{lemma}

\begin{proof}
Supposons
$\overline{g} = \overline{b}_0 x^n + \overline{b}_1 x^{n - 1} + \ldots
+ \overline{b}_n$, et
$\overline{h} = \overline{c}_0 x^m + \overline{c}_1 x^{m - 1} + \ldots
+ \overline{c}_m$, avec
$\overline{b}_0, \overline{c}_0 \in \kappa(\mathfrak p)$ non nuls.
Après avoir localisé $R$ en un certain élément de $R$ qui n'appartient pas à
$\mathfrak p$, nous pouvons supposer que $\overline{b}_0$ est l'image d'un
élément inversible $b_0 \in R$. En remplaçant
$\overline{g}$ par $\overline{g}/b_0$ et $\overline{h}$ par
$b_0\overline{h}$, nous nous ramenons au cas où $\overline{g}$,
$\overline{h}$ sont unitaires (vérification omise).
Disons que
$\overline{g} = x^n + \overline{b}_1 x^{n - 1} + \ldots + \overline{b}_n$,
et $\overline{h} = x^m + \overline{c}_1 x^{m - 1} + \ldots + \overline{c}_m$.
Écrivons $f = x^{n + m} + a_1 x^{n + m - 1} + \ldots + a_{n + m}$.
Considérons le produit fibré
$$
R' = R \otimes_{\mathbf{Z}[a_1, \ldots, a_{n + m}]}
\mathbf{Z}[b_1, \ldots, b_n, c_1, \ldots, c_m]
$$
où le morphisme $\mathbf{Z}[a_k] \to \mathbf{Z}[b_i, c_j]$ est celui des
Exemples \ref{algebra-example-factor-polynomials} et
\ref{algebra-example-factor-polynomials-etale}. Par construction, il existe un
morphisme de $R$-algèbres
$$
R' = R \otimes_{\mathbf{Z}[a_1, \ldots, a_{n + m}]}
\mathbf{Z}[b_1, \ldots, b_n, c_1, \ldots, c_m]
\longrightarrow
\kappa(\mathfrak p)
$$
qui envoie $b_i$ sur $\overline{b}_i$ et $c_j$ sur $\overline{c}_j$.
Notons $\mathfrak p' \subset R'$ le noyau de ce morphisme.
Comme, par hypothèse, les polynômes $\overline{g}, \overline{h}$ sont premiers
entre eux, nous voyons que l'élément
$\Delta = \text{Res}_x(g, h) \in \mathbf{Z}[b_i, c_j]$
(voir l'Exemple \ref{algebra-example-factor-polynomials-etale}) ne s'envoie pas sur zéro
dans $\kappa(\mathfrak p)$ par le morphisme affiché.
Nous concluons que $R \to R'$ est étale en $\mathfrak p'$.
En fait, une solution au problème posé dans le lemme est donnée par le
morphisme d'anneaux $R \to R'[1/\Delta]$ et l'idéal premier
$\mathfrak p' R'[1/\Delta]$. Comme $\text{Res}_x(f, g)$ est inversible dans cet
anneau, la matrice de Sylvester est inversible sur $R'[1/\Delta]$ et donc
$1 = a g +  b h$ pour certains $a, b \in R'[1/\Delta][x]$ ; voir
l'Exemple \ref{algebra-example-factor-polynomials-etale}.
\end{proof}




\section{Structure locale des morphismes d'anneaux étales}
\label{algebra-section-etale-local-structure}

\noindent
Le Lemme \ref{algebra-lemma-etale-standard-smooth} nous indique qu'il n'est pas
vraiment pertinent de définir un morphisme étale standard comme un morphisme
lisse standard de dimension relative $0$.
Comme modèle d'un morphisme étale, nous prenons l'exemple donné par une
extension de corps finie séparable $k'/k$.
En effet, nous pouvons toujours trouver un élément $\alpha \in k'$ tel que
$k' = k(\alpha)$ et tel que le polynôme minimal $f(x) \in k[x]$ de $\alpha$
ait une dérivée $f'$ première avec $f$.

\begin{definition}
\label{algebra-definition-standard-etale}
Soit $R$ un anneau. Soient $g , f  \in R[x]$.
Supposons que $f$ soit unitaire et que la dérivée $f'$ soit inversible dans
la localisation $R[x]_g/(f)$.
Dans ce cas, le morphisme d'anneaux $R \to R[x]_g/(f)$ est dit
{\it étale standard}.
\end{definition}

\noindent
Dans la Proposition \ref{algebra-proposition-etale-locally-standard}, nous montrons
que tout morphisme d'anneaux étale est localement étale standard.

\begin{lemma}
\label{algebra-lemma-standard-etale}
Soit $R \to R[x]_g/(f)$ un morphisme étale standard.
\begin{enumerate}
\item Le morphisme d'anneaux $R \to R[x]_g/(f)$ est étale.
\item Pour tout morphisme d'anneaux $R \to R'$, le changement de base
$R' \to R'[x]_g/(f)$ du morphisme étale standard
$R \to R[x]_g/(f)$ est étale standard.
\item Toute localisation principale de $R[x]_g/(f)$ est étale standard sur $R$.
\item Une composée de morphismes étales standard n'est {\bf pas} étale standard
en général.
\end{enumerate}
\end{lemma}

\begin{proof}
Omis. Voici un exemple pour (4).
Le morphisme d'anneaux $\mathbf{F}_2 \to \mathbf{F}_{2^2}$ est étale standard.
Le morphisme d'anneaux
$\mathbf{F}_{2^2} \to \mathbf{F}_{2^2} \times \mathbf{F}_{2^2}
\times \mathbf{F}_{2^2} \times \mathbf{F}_{2^2}$ est étale standard.
Mais le morphisme d'anneaux
$\mathbf{F}_2 \to \mathbf{F}_{2^2} \times \mathbf{F}_{2^2}
\times \mathbf{F}_{2^2} \times \mathbf{F}_{2^2}$ n'est pas étale standard.
\end{proof}

\noindent
Les morphismes étales standard constituent un moyen commode de produire des
morphismes étales. En voici un exemple.

\begin{lemma}
\label{algebra-lemma-make-etale-map-prescribed-residue-field}
Soit $R$ un anneau.
Soit $\mathfrak p$ un idéal premier de $R$.
Soit $L/\kappa(\mathfrak p)$ une extension de corps finie séparable.
Il existe un morphisme d'anneaux étale $R \to R'$, ainsi qu'un idéal premier
$\mathfrak p'$ au-dessus de $\mathfrak p$, tels que l'extension de corps
$\kappa(\mathfrak p')/\kappa(\mathfrak p)$ soit isomorphe à
$\kappa(\mathfrak p) \subset L$.
\end{lemma}

\begin{proof}
Par le théorème de l'élément primitif, nous pouvons écrire
$L = \kappa(\mathfrak p)[\alpha]$. Soit
$\overline{f} \in \kappa(\mathfrak p)[x]$ le polynôme minimal de $\alpha$
(en particulier, il est unitaire).
Après avoir remplacé $\alpha$ par $c\alpha$ pour un certain
$c \in R$, $c\not \in \mathfrak p$, nous pouvons supposer que tous les
coefficients de $\overline{f}$ appartiennent à l'image de
$R \to \kappa(\mathfrak p)$ (vérification omise).
Nous pouvons donc trouver un polynôme unitaire $f \in R[x]$ qui s'envoie sur
$\overline{f}$ dans $\kappa(\mathfrak p)[x]$.
Comme $\kappa(\mathfrak p) \subset L$ est séparable, nous voyons que
$\gcd(\overline{f}, \overline{f}') = 1$.
Il existe donc un élément $\gamma \in L$ tel que
$\overline{f}'(\alpha) \gamma = 1$. Nous obtenons ainsi un morphisme de
$R$-algèbres
\begin{eqnarray*}
R[x, 1/f']/(f) & \longrightarrow & L \\
x & \longmapsto & \alpha \\
1/f' & \longmapsto & \gamma
\end{eqnarray*}
Le membre de gauche est une algèbre étale standard $R'$ sur $R$, et le noyau
du morphisme d'anneaux fournit l'idéal premier recherché.
\end{proof}

\begin{proposition}
\label{algebra-proposition-etale-locally-standard}
Soit $R \to S$ un morphisme d'anneaux. Soit $\mathfrak q \subset S$ un idéal
premier. Si $R \to S$ est étale en $\mathfrak q$, alors il existe
$g \in S$, $g \not \in \mathfrak q$, tel que $R \to S_g$ soit étale standard.
\end{proposition}

\begin{proof}
La démonstration suivante est un peu indirecte, et il existe peut-être des
moyens de la raccourcir.

\medskip\noindent
Étape 1. Par la Définition \ref{algebra-definition-etale}, il existe
$g \in S$, $g \not \in \mathfrak q$, tel que $R \to S_g$ soit étale.
Nous pouvons donc supposer que $S$ est étale sur $R$.

\medskip\noindent
Étape 2. Par le Lemme \ref{algebra-lemma-etale}, il existe un morphisme d'anneaux étale
$R_0 \to S_0$, avec $R_0$ de type fini sur $\mathbf{Z}$, et un morphisme
d'anneaux $R_0 \to R$ tels que $R = R \otimes_{R_0} S_0$.
Notons $\mathfrak q_0$ l'idéal premier de $S_0$ correspondant à
$\mathfrak q$. Si nous démontrons le résultat pour
$(R_0 \to S_0, \mathfrak q_0)$, alors il en découle pour
$(R \to S, \mathfrak q)$ par changement de base.
Nous pouvons donc supposer que $R$ est noethérien.

\medskip\noindent
Étape 3.
Remarquons que $R \to S$ est quasi-fini par le
Lemme \ref{algebra-lemma-etale-quasi-finite}.
Par le Lemme \ref{algebra-lemma-quasi-finite-open-integral-closure}, il existe un
morphisme d'anneaux fini $R \to S'$, un morphisme de $R$-algèbres
$S' \to S$ et un élément $g' \in S'$ dont l'image vérifie
$g' \not \in \mathfrak q$ et tel que $S' \to S$ induise un isomorphisme
$S'_{g'} \cong S_{g'}$.
(Remarquons que, bien entendu, $S'$ n'est en général pas étale sur $R$.)
Nous pouvons donc supposer que (a) $R$ est noethérien, (b) $R \to S$ est fini
et (c) $R \to S$ est étale en $\mathfrak q$
(mais plus nécessairement étale en tout idéal premier).

\medskip\noindent
Étape 4. Soit $\mathfrak p \subset R$ l'idéal premier correspondant à
$\mathfrak q$. Considérons l'anneau fibre
$S \otimes_R \kappa(\mathfrak p)$. C'est une algèbre finie sur
$\kappa(\mathfrak p)$. Il est donc artinien
(voir le Lemme \ref{algebra-lemma-finite-dimensional-algebra}), et par conséquent
un produit fini d'anneaux locaux
$$
S \otimes_R \kappa(\mathfrak p) = \prod\nolimits_{i = 1}^n A_i
$$
voir la Proposition \ref{algebra-proposition-dimension-zero-ring}. L'un des facteurs,
disons $A_1$, est l'anneau local
$S_{\mathfrak q}/\mathfrak pS_{\mathfrak q}$, qui est isomorphe à
$\kappa(\mathfrak q)$ ; voir le Lemme \ref{algebra-lemma-etale-at-prime}.
Les autres facteurs correspondent aux autres idéaux premiers, disons
$\mathfrak q_2, \ldots, \mathfrak q_n$, de $S$ au-dessus de $\mathfrak p$.

\medskip\noindent
Étape 5. Nous pouvons choisir un élément non nul
$\alpha \in \kappa(\mathfrak q)$ qui engendre l'extension de corps finie
séparable $\kappa(\mathfrak q)/\kappa(\mathfrak p)$
(ainsi, même si l'extension de corps est triviale, nous n'autorisons pas
$\alpha = 0$).
Remarquons que, pour tout $\lambda \in \kappa(\mathfrak p)^*$, l'élément
$\lambda \alpha$ engendre lui aussi $\kappa(\mathfrak q)$ sur
$\kappa(\mathfrak p)$. Considérons l'élément
$$
\overline{t} =
(\alpha, 0, \ldots, 0) \in
\prod\nolimits_{i = 1}^n A_i =
S \otimes_R \kappa(\mathfrak p).
$$
Après avoir éventuellement remplacé $\alpha$ par $\lambda \alpha$ comme
ci-dessus, nous pouvons supposer que $\overline{t}$ est l'image de
$t \in S$. Soit $I \subset R[x]$ le noyau du morphisme de $R$-algèbres
$R[x] \to S$ qui envoie $x$ sur $t$. Posons $S' = R[x]/I$,
de sorte que $S' \subset S$. Voici un diagramme
$$
\xymatrix{
R[x] \ar[r] & S' \ar[r] & S \\
R \ar[u] \ar[ru] \ar[rru] & &
}
$$
Par construction, les idéaux premiers $\mathfrak q_j$, $j \geq 2$, de $S$
se trouvent tous au-dessus de l'idéal premier $(\mathfrak p, x)$ de $R[x]$,
tandis que l'idéal premier $\mathfrak q$ se trouve au-dessus d'un idéal
premier différent de $R[x]$, puisque $\alpha \not = 0$.

\medskip\noindent
Étape 6. Notons $\mathfrak q' \subset S'$ l'idéal premier de $S'$
correspondant à $\mathfrak q$. D'après ce qui précède, $\mathfrak q$ est
le seul idéal premier de $S$ au-dessus de $\mathfrak q'$. Nous voyons donc que
$S_{\mathfrak q} = S_{\mathfrak q'}$ ; voir le
Lemme \ref{algebra-lemma-unique-prime-over-localize-below}
(nous disposons de la propriété de montée pour $S' \to S$ par le
Lemme \ref{algebra-lemma-integral-going-up}, puisque $S' \to S$ est fini car
$R \to S$ est fini).
Il s'ensuit que $S'_{\mathfrak q'} \to S_{\mathfrak q}$ est fini et injectif,
comme localisation du morphisme d'anneaux fini injectif $S' \to S$.
Considérons les morphismes d'anneaux locaux
$$
R_{\mathfrak p} \to S'_{\mathfrak q'} \to S_{\mathfrak q}
$$
Le second morphisme est fini et injectif. Nous avons
$S_{\mathfrak q}/\mathfrak pS_{\mathfrak q} = \kappa(\mathfrak q)$ ;
voir le Lemme \ref{algebra-lemma-etale-at-prime}.
Donc, a fortiori,
$S_{\mathfrak q}/\mathfrak q'S_{\mathfrak q} = \kappa(\mathfrak q)$.
Puisque
$$
\kappa(\mathfrak p) \subset \kappa(\mathfrak q') \subset \kappa(\mathfrak q)
$$
et puisque $\alpha$ appartient à l'image de $\kappa(\mathfrak q')$ dans
$\kappa(\mathfrak q)$, nous concluons que
$\kappa(\mathfrak q') = \kappa(\mathfrak q)$.
Ainsi, par le Lemme de Nakayama \ref{algebra-lemma-NAK} appliqué au morphisme de
$S'_{\mathfrak q'}$-modules $S'_{\mathfrak q'} \to S_{\mathfrak q}$,
le morphisme $S'_{\mathfrak q'} \to S_{\mathfrak q}$ est surjectif.
Autrement dit,
$S'_{\mathfrak q'} \cong S_{\mathfrak q}$.

\medskip\noindent
Étape 7. Par le Lemme \ref{algebra-lemma-isomorphic-local-rings}, il existe
$g \in S$, $g \not \in \mathfrak q$, et
$g' \in S'$, $g' \not \in \mathfrak q'$, tels que
$S'_{g'} \cong S_g$. Comme $R$ est noethérien, l'anneau $S'$ est fini
sur $R$, car c'est un sous-$R$-module du $R$-module fini $S$.
Après avoir remplacé $S$ par $S'$, nous pouvons donc supposer que
(a) $R$ est noethérien, (b) $S$ est fini sur $R$, (c)
$S$ est étale sur $R$ en $\mathfrak q$ et (d) $S = R[x]/I$.

\medskip\noindent
Étape 8. Considérons l'anneau
$S \otimes_R \kappa(\mathfrak p) = \kappa(\mathfrak p)[x]/\overline{I}$,
où $\overline{I} = I \cdot \kappa(\mathfrak p)[x]$ est l'idéal engendré par
$I$ dans $\kappa(\mathfrak p)[x]$. Comme $\kappa(\mathfrak p)[x]$ est un
anneau principal, nous savons que
$\overline{I} = (\overline{h})$ pour un certain polynôme unitaire
$\overline{h} \in \kappa(\mathfrak p)[x]$.
Après avoir remplacé $\overline{h}$ par
$\lambda \cdot \overline{h}$ pour un certain
$\lambda \in \kappa(\mathfrak p)$, nous pouvons supposer que
$\overline{h}$ est l'image d'un certain $h \in I \subset R[x]$.
(Le problème est que nous ne savons pas si nous pouvons choisir $h$ unitaire.)
De plus, comme à l'Étape 4, nous savons que
$S \otimes_R \kappa(\mathfrak p) = A_1 \times \ldots \times A_n$, avec
$A_1 = \kappa(\mathfrak q)$ une extension finie séparable de
$\kappa(\mathfrak p)$ et $A_2, \ldots, A_n$ locaux. Cela implique que
$$
\overline{h} = \overline{h}_1 \overline{h}_2^{e_2} \ldots \overline{h}_n^{e_n}
$$
pour certains polynômes unitaires irréductibles deux à deux premiers entre eux
$\overline{h}_i \in \kappa(\mathfrak p)[x]$ et certains
$e_2, \ldots, e_n \geq 1$. Ici, la numérotation est choisie de telle sorte que
$A_i = \kappa(\mathfrak p)[x]/(\overline{h}_i^{e_i})$ comme
$\kappa(\mathfrak p)[x]$-algèbres. Remarquons que $\overline{h}_1$ est le
polynôme minimal de $\alpha \in \kappa(\mathfrak q)$ et est donc un polynôme
séparable (sa dérivée est première avec lui-même).

\medskip\noindent
Étape 9. Soit $m \in I$ un élément unitaire ; un tel élément existe parce que
l'extension d'anneaux $R \to R[x]/I$ est finie, donc entière.
Notons $\overline{m}$ l'image dans $\kappa(\mathfrak p)[x]$.
Nous pouvons factoriser
$$
\overline{m} = \overline{k}
\overline{h}_1^{d_1} \overline{h}_2^{d_2} \ldots \overline{h}_n^{d_n}
$$
pour un certain $d_1 \geq 1$, des $d_j \geq e_j$,
$j = 2, \ldots, n$, et un
$\overline{k} \in \kappa(\mathfrak p)[x]$ premier avec tous les
$\overline{h}_i$. Posons $f = m^l + h$, où
$l \deg(m) > \deg(h)$ et $l \geq 2$.
Alors $f$ est unitaire comme polynôme sur $R$. De plus, l'image
$\overline{f}$ de $f$ dans $\kappa(\mathfrak p)[x]$ se factorise sous la forme
$$
\overline{f} =
\overline{h}_1 \overline{h}_2^{e_2} \ldots \overline{h}_n^{e_n}
+
\overline{k}^l \overline{h}_1^{ld_1} \overline{h}_2^{ld_2}
\ldots \overline{h}_n^{ld_n}
=
\overline{h}_1(\overline{h}_2^{e_2} \ldots \overline{h}_n^{e_n}
+
\overline{k}^l
\overline{h}_1^{ld_1 - 1} \overline{h}_2^{ld_2} \ldots \overline{h}_n^{ld_n})
= \overline{h}_1 \overline{w}
$$
avec $\overline{w}$ un polynôme premier avec $\overline{h}_1$.
Posons $g = f'$ (la dérivée par rapport à $x$).

\medskip\noindent
Étape 10. Le morphisme d'anneaux $R[x] \to S = R[x]/I$ possède les propriétés
suivantes : (1) il envoie $f$ sur zéro, et
(2) il envoie $g$ sur un élément de $S \setminus \mathfrak q$.
La première assertion est claire, puisque $f$ est un élément de $I$.
Pour la seconde, il suffit de montrer que $g$ ne s'envoie pas sur zéro dans
$\kappa(\mathfrak q) = \kappa(\mathfrak p)[x]/(\overline{h}_1)$.
L'image de $g$ dans $\kappa(\mathfrak p)[x]$ est la dérivée de
$\overline{f}$. Ainsi, (2) est clair parce que
$$
\overline{g} =
\frac{\text{d}\overline{f}}{\text{d}x} =
\overline{w}\frac{\text{d}\overline{h}_1}{\text{d}x} +
\overline{h}_1\frac{\text{d}\overline{w}}{\text{d}x},
$$
$\overline{w}$ est premier avec $\overline{h}_1$ et
$\overline{h}_1$ est séparable.

\medskip\noindent
Étape 11.
Nous concluons que $\varphi : R[x]/(f) \to S$ est un morphisme d'anneaux
surjectif, que $R[x]_g/(f)$ est étale sur $R$
(parce qu'il est étale standard ; voir le Lemme \ref{algebra-lemma-standard-etale})
et que $\varphi(g) \not \in \mathfrak q$.
Choisissons un élément $g' \in R[x]/(f)$ tel que
$\varphi(g') \not \in \mathfrak q$ lui aussi et que
$S_{\varphi(g')}$ soit étale sur $R$ (un tel élément existe puisque $S$ est
étale sur $R$ en $\mathfrak q$). Alors le morphisme d'anneaux
$R[x]_{gg'}/(f) \to S_{\varphi(gg')}$ est un morphisme surjectif d'algèbres
étales sur $R$. Il est donc étale par le
Lemme \ref{algebra-lemma-map-between-etale}.
C'est donc une localisation par le
Lemme \ref{algebra-lemma-surjective-flat-finitely-presented}.
Ainsi, une localisation de $S$ en un élément n'appartenant pas à
$\mathfrak q$ est isomorphe à une localisation d'une algèbre étale standard
sur $R$, ce que nous voulions démontrer.
\end{proof}

\noindent
Les deux lemmes suivants affirment que la topologie étale est moins fine que
la topologie engendrée par les recouvrements de Zariski et les morphismes
finis plats. Ils devraient être omis lors d'une première lecture.

\begin{lemma}
\label{algebra-lemma-standard-etale-finite-flat-Zariski}
Soit $R \to S$ un morphisme étale standard.
Il existe un morphisme d'anneaux $R \to S'$ possédant les propriétés suivantes :
\begin{enumerate}
\item $R \to S'$ est fini, de présentation finie et plat
(autrement dit, $S'$ est un $R$-module projectif de type fini),
\item $\Spec(S') \to \Spec(R)$ est surjectif,
\item pour tout idéal premier $\mathfrak q \subset S$ au-dessus de
$\mathfrak p \subset R$ et tout idéal premier
$\mathfrak q' \subset S'$ au-dessus de $\mathfrak p$, il existe
$g' \in S'$, $g' \not \in \mathfrak q'$, tel que le morphisme d'anneaux
$R \to S'_{g'}$ se factorise par un morphisme
$\varphi : S \to S'_{g'}$ avec
$\varphi^{-1}(\mathfrak q'S'_{g'}) = \mathfrak q$.
\end{enumerate}
\end{lemma}

\begin{proof}
Soit $S = R[x]_g/(f)$ une présentation de $S$ comme dans la
Définition \ref{algebra-definition-standard-etale}.
Écrivons $f = x^n + a_1 x^{n - 1} + \ldots + a_n$, avec $a_i \in R$.
Par le Lemme \ref{algebra-lemma-adjoin-roots}, il existe un morphisme d'anneaux
fini libre et fidèlement plat $R \to S'$ tel que
$f = \prod (x - \alpha_i)$ pour certains $\alpha_i \in S'$.
Ainsi, $R \to S'$ satisfait les conditions (1) et (2).
Soit $\mathfrak q \subset R[x]/(f)$ un idéal premier avec
$g \not \in \mathfrak q$ (c'est-à-dire qu'il correspond à un idéal premier
de $S$). Soit $\mathfrak p = R \cap \mathfrak q$, et soit
$\mathfrak q' \subset S'$ un idéal premier au-dessus de $\mathfrak p$.
Remarquons qu'il existe $n$ morphismes de $R$-algèbres
\begin{eqnarray*}
\varphi_i : R[x]/(f) & \longrightarrow & S' \\
x & \longmapsto & \alpha_i
\end{eqnarray*}
Pour achever la démonstration, nous devons montrer que, pour un certain $i$,
(a) l'image de $\varphi_i(g)$ dans $\kappa(\mathfrak q')$ n'est pas nulle,
et (b) $\varphi_i^{-1}(\mathfrak q') = \mathfrak q$.
Nous pourrons alors simplement prendre $g' = \varphi_i(g)$ et
$\varphi = \varphi_i$ pour cet $i$.

\medskip\noindent
Notons $\overline{f}$ l'image de $f$ dans $\kappa(\mathfrak p)[x]$.
Remarquons que, comme point de
$\Spec(\kappa(\mathfrak p)[x]/(\overline{f}))$, l'idéal premier
$\mathfrak q$ correspond à un facteur irréductible $f_1$ de
$\overline{f}$. De plus, $g \not \in \mathfrak q$ signifie que $f_1$ ne
divise pas l'image $\overline{g}$ de $g$ dans $\kappa(\mathfrak p)[x]$.
Notons $\overline{\alpha}_1, \ldots, \overline{\alpha}_n$ les images de
$\alpha_1, \ldots, \alpha_n$ dans $\kappa(\mathfrak q')$.
Remarquons que le polynôme $\overline{f}$ se scinde complètement dans
$\kappa(\mathfrak q')[x]$, à savoir
$$
\overline{f} = \prod\nolimits_i (x - \overline{\alpha}_i)
$$
De plus, $\varphi_i(g)$ se réduit à
$\overline{g}(\overline{\alpha}_i)$.
Il s'ensuit que nous pouvons choisir $i$ tel que
$f_1(\overline{\alpha}_i) = 0$ et
$\overline{g}(\overline{\alpha}_i) \not = 0$.
Pour cet $i$, les propriétés (a) et (b) sont satisfaites.
Certains détails sont omis.
\end{proof}

\begin{lemma}
\label{algebra-lemma-etale-finite-flat-zariski}
Soit $R \to S$ un morphisme d'anneaux.
Supposons que
\begin{enumerate}
\item $R \to S$ soit étale, et
\item $\Spec(S) \to \Spec(R)$ soit surjectif.
\end{enumerate}
Alors il existe un morphisme d'anneaux $R \to S'$ tel que
\begin{enumerate}
\item $R \to S'$ soit fini, de présentation finie et plat
(autrement dit, il est projectif de type fini comme $R$-module),
\item $\Spec(S') \to \Spec(R)$ soit surjectif,
\item pour tout idéal premier $\mathfrak q' \subset S'$, il existe
$g' \in S'$, $g' \not \in \mathfrak q'$, tel que le morphisme d'anneaux
$R \to S'_{g'}$ se factorise sous la forme
$R \to S \to S'_{g'}$.
\end{enumerate}
\end{lemma}

\begin{proof}
Par la Proposition \ref{algebra-proposition-etale-locally-standard} et la
quasi-compacité de $\Spec(S)$ (voir le Lemme \ref{algebra-lemma-quasi-compact}),
nous pouvons trouver des éléments $g_1, \ldots, g_n \in S$ qui engendrent
l'idéal unité de $S$ et tels que chacun des $R \to S_{g_i}$ soit étale
standard. Si nous démontrons le lemme pour le morphisme d'anneaux
$R \to \prod_{i = 1, \ldots, n} S_{g_i}$, alors le lemme en découle pour
le morphisme d'anneaux $R \to S$.
Nous pouvons donc supposer que
$S = \prod_{i = 1, \ldots, n} S_i$ est un produit fini de morphismes
étales standard.

\medskip\noindent
Pour chaque $i$, choisissons un morphisme d'anneaux $R \to S_i'$ comme dans le
Lemme \ref{algebra-lemma-standard-etale-finite-flat-Zariski}, adapté au morphisme
étale standard $R \to S_i$.
Posons $S' = S_1' \otimes_R \ldots \otimes_R S_n'$ ; ci-dessous, nous
utiliserons sans autre mention les morphismes de $R$-algèbres
$S_i' \to S'$. Nous affirmons que cela convient.
Les propriétés (1) et (2) sont immédiates.
Pour la propriété (3), supposons que $\mathfrak q' \subset S'$ soit un idéal
premier. Notons $\mathfrak p$ son image dans $\Spec(R)$.
Choisissons $i \in \{1, \ldots, n\}$ tel que $\mathfrak p$ appartienne à
l'image de $\Spec(S_i) \to \Spec(R)$ ; cela est possible par hypothèse.
Soit $\mathfrak q_i' \subset S_i'$ l'image de $\mathfrak q'$ dans le spectre
de $S_i'$. Par construction de $S'_i$, il existe
$g'_i \in S_i'$ tel que $R \to (S_i')_{g_i'}$ se factorise sous la forme
$R \to S_i \to (S_i')_{g_i'}$. Par conséquent,
$R \to S'_{g_i'}$ se factorise également sous la forme
$$
R \to S_i \to (S_i')_{g_i'} \to S'_{g_i'}
$$
comme désiré.
\end{proof}







\section{Structure locale étale des morphismes d'anneaux quasi-finis}
\label{algebra-section-etale-local-quasi-finite}

\noindent
Les lemmes suivants affirment, en gros, qu'après une extension étale, un
morphisme d'anneaux quasi-fini devient fini.
Pour faciliter l'interprétation des résultats, rappelons que le lieu où un
morphisme d'anneaux de type fini est quasi-fini est ouvert
(voir le Lemme \ref{algebra-lemma-quasi-finite-open}) et que la formation de ce lieu
commute aux changements de base arbitraires
(voir le Lemme \ref{algebra-lemma-quasi-finite-base-change}).

\begin{lemma}
\label{algebra-lemma-produce-finite}
Soient $R \to S' \to S$ des morphismes d'anneaux.
Soit $\mathfrak p \subset R$ un idéal premier.
Soit $g \in S'$ un élément.
Supposons que
\begin{enumerate}
\item $R \to S'$ soit entier,
\item $R \to S$ soit de type fini,
\item $S'_g \cong S_g$, et
\item $g$ soit inversible dans $S' \otimes_R \kappa(\mathfrak p)$.
\end{enumerate}
Alors il existe $f \in R$, $f \not \in \mathfrak p$, tel que
$R_f \to S_f$ soit fini.
\end{lemma}

\begin{proof}
Par hypothèse, l'image $T$ de $V(g) \subset \Spec(S')$ par le morphisme
$\Spec(S') \to \Spec(R)$ ne contient pas $\mathfrak p$.
Par la section \ref{algebra-section-going-up}, et en particulier le
Lemme \ref{algebra-lemma-going-up-closed}, nous voyons que $T$ est fermé.
Choisissons $f \in R$, $f \not \in \mathfrak p$, tel que
$T \cap D(f) = \emptyset$. Nous voyons alors que $g$ devient inversible dans
$S'_f$. Par conséquent, $S'_f \cong S_f$.
Ainsi, $S_f$ est à la fois de type fini et entier sur $R_f$, donc fini.
\end{proof}

\begin{lemma}
\label{algebra-lemma-etale-makes-quasi-finite-finite-one-prime}
Soit $R \to S$ un morphisme d'anneaux.
Soit $\mathfrak q \subset S$ un idéal premier au-dessus de l'idéal premier
$\mathfrak p \subset R$.
Supposons que $R \to S$ soit de type fini et quasi-fini en $\mathfrak q$.
Alors il existe
\begin{enumerate}
\item un morphisme d'anneaux étale $R \to R'$,
\item un idéal premier $\mathfrak p' \subset R'$ au-dessus de $\mathfrak p$,
\item une décomposition en produit
$$
R' \otimes_R S = A \times B
$$
\end{enumerate}
possédant les propriétés suivantes :
\begin{enumerate}
\item $\kappa(\mathfrak p) = \kappa(\mathfrak p')$,
\item $R' \to A$ est fini,
\item $A$ possède exactement un idéal premier $\mathfrak r$ au-dessus de
$\mathfrak p'$,
\item $\mathfrak r$ est au-dessus de $\mathfrak q$, et
\item $B$ ne possède aucun idéal premier au-dessus de $\mathfrak q$ et de
$\mathfrak p'$.
\end{enumerate}
\end{lemma}

\begin{proof}
Soit $S' \subset S$ la clôture intégrale de $R$ dans $S$.
Soit $\mathfrak q' = S' \cap \mathfrak q$.
Par le Théorème principal de Zariski \ref{algebra-theorem-main-theorem}, il existe
$g \in S'$, $g \not \in \mathfrak q'$, tel que $S'_g \cong S_g$.
Considérons les anneaux fibres
$F = S \otimes_R \kappa(\mathfrak p)$ et
$F' = S' \otimes_R \kappa(\mathfrak p)$.
Notons $\overline{\mathfrak q}'$ l'idéal premier de $F'$ correspondant à
$\mathfrak q'$. Comme $F'$ est entier sur $\kappa(\mathfrak p)$, nous voyons
que $\overline{\mathfrak q}'$ est un point fermé de $\Spec(F')$ ;
voir le Lemme \ref{algebra-lemma-integral-over-field}.
Remarquons que $\mathfrak q$ définit un point fermé isolé
$\overline{\mathfrak q}$ de $\Spec(F)$
(voir la Définition \ref{algebra-definition-quasi-finite}).
Comme $S'_g \cong S_g$, nous avons $F'_g \cong F_g$ ; ainsi,
$\overline{\mathfrak q}$ et $\overline{\mathfrak q}'$ possèdent des
voisinages ouverts isomorphes dans $\Spec(F)$ et $\Spec(F')$.
Nous concluons que l'ensemble
$\{\overline{\mathfrak q}'\} \subset \Spec(F')$ est ouvert.
En combinant cela avec le fait que $\mathfrak q'$ est fermé
(démontré ci-dessus), nous concluons que $\overline{\mathfrak q}'$ définit
lui aussi un point fermé isolé de $\Spec(F')$.

\medskip\noindent
Une petite remarque supplémentaire est que, sous le morphisme
$\Spec(F) \to \Spec(F')$, le point $\overline{\mathfrak q}$ est l'unique
point envoyé sur $\overline{\mathfrak q}'$. Cela découle de la discussion
ci-dessus.

\medskip\noindent
Par le Lemme \ref{algebra-lemma-disjoint-implies-product}, nous pouvons écrire
$F' = F'_1 \times F'_2$, avec
$\Spec(F'_1) = \{\overline{\mathfrak q}'\}$.
Comme $F' = S' \otimes_R \kappa(\mathfrak p)$, il existe
$s' \in S'$ qui s'envoie sur l'élément
$(r, 0) \in F'_1 \times F'_2 = F'$ pour un certain
$r \in R$, $r \not \in \mathfrak p$.
En fait, ce que nous utiliserons au sujet de $s'$, c'est qu'il s'agit d'un
élément de $S'$ qui n'appartient pas à $\mathfrak q'$ et qui appartient à tout
autre idéal premier au-dessus de $\mathfrak p$.

\medskip\noindent
Soit $f(x) \in R[x]$ un polynôme unitaire tel que $f(s') = 0$.
Notons $\overline{f} \in \kappa(\mathfrak p)[x]$ son image.
Nous pouvons le factoriser sous la forme
$\overline{f} = x^e \overline{h}$, où
$\overline{h}(0) \not = 0$. Après avoir remplacé $f$ par $x f$ si nécessaire,
nous pouvons supposer que $e \geq 1$.
Par le Lemme \ref{algebra-lemma-factor-mod-lift-etale}, nous pouvons trouver une
extension d'anneaux étale $R \to R'$, un idéal premier
$\mathfrak p'$ au-dessus de $\mathfrak p$ et une factorisation
$f = h i$ dans $R'[x]$, tels que
$\kappa(\mathfrak p) = \kappa(\mathfrak p')$,
$\overline{h} = h \bmod \mathfrak p'$,
$x^e = i \bmod \mathfrak p'$, et que nous puissions écrire
$a h + b i = 1$ dans $R'[x]$ (pour des $a, b$ convenables).

\medskip\noindent
Considérons les éléments $h(s'), i(s') \in R' \otimes_R S'$.
Par construction, nous avons $h(s')i(s') = f(s') = 0$.
D'autre part, ils engendrent l'idéal unité puisque
$a(s')h(s') + b(s')i(s') = 1$.
Nous voyons donc que $R' \otimes_R S'$ est le produit des localisations en
ces éléments :
$$
R' \otimes_R S'
=
(R' \otimes_R S')_{i(s')}
\times
(R' \otimes_R S')_{h(s')}
=
S'_1 \times S'_2
$$
De plus, cette décomposition en produit est compatible avec celle que nous
avons trouvée pour l'anneau fibre $F'$ ; cela résulte de nos choix de
$s', i, h$, qui garantissent que $\overline{\mathfrak q}'$ est le seul idéal
premier de $F'$ qui ne contient pas l'image de $i(s')$ dans $F'$.
Nous utilisons ici le fait que l'anneau fibre de
$R'\otimes_R S'$ sur $R'$ en $\mathfrak p'$ est le même que $F'$, puisque
$\kappa(\mathfrak p) = \kappa(\mathfrak p')$.
Il s'ensuit que $S'_1$ possède exactement un idéal premier,
disons $\mathfrak r'$, au-dessus de $\mathfrak p'$, et que cet idéal premier
est au-dessus de $\mathfrak q'$.
Par conséquent, l'élément $g \in S'$ s'envoie sur un élément de $S'_1$
qui n'appartient pas à $\mathfrak r'$.

\medskip\noindent
Le changement de base $R'\otimes_R S$ hérite d'une décomposition en produit
semblable :
$$
R' \otimes_R S
=
(R' \otimes_R S)_{i(s')}
\times
(R' \otimes_R S)_{h(s')}
=
S_1 \times S_2
$$
Il résulte de ce qui précède que $S_1$ possède exactement un idéal premier,
disons $\mathfrak r$, au-dessus de $\mathfrak p'$
(considérer l'anneau fibre comme ci-dessus), et que cet idéal premier est
au-dessus de $\mathfrak q$.

\medskip\noindent
Nous pouvons maintenant appliquer le Lemme \ref{algebra-lemma-produce-finite} aux
morphismes d'anneaux $R' \to S'_1 \to S_1$, à l'idéal premier
$\mathfrak p'$ et à l'élément $g$ pour voir qu'après avoir remplacé $R'$ par
une localisation principale, nous pouvons supposer que $S_1$ est fini sur
$R'$, comme désiré.
\end{proof}

\begin{lemma}
\label{algebra-lemma-etale-makes-quasi-finite-finite}
Soit $R \to S$ un morphisme d'anneaux.
Soit $\mathfrak p \subset R$ un idéal premier.
Supposons que $R \to S$ soit de type fini.
Alors il existe
\begin{enumerate}
\item un morphisme d'anneaux étale $R \to R'$,
\item un idéal premier $\mathfrak p' \subset R'$ au-dessus de $\mathfrak p$,
\item une décomposition en produit
$$
R' \otimes_R S = A_1 \times \ldots \times A_n \times B
$$
\end{enumerate}
possédant les propriétés suivantes :
\begin{enumerate}
\item nous avons $\kappa(\mathfrak p) = \kappa(\mathfrak p')$,
\item chaque $A_i$ est fini sur $R'$,
\item chaque $A_i$ possède exactement un idéal premier $\mathfrak r_i$
au-dessus de $\mathfrak p'$, et
\item $R' \to B$ n'est quasi-fini en aucun idéal premier au-dessus de
$\mathfrak p'$.
\end{enumerate}
\end{lemma}

\begin{proof}
Notons $F = S \otimes_R \kappa(\mathfrak p)$ l'anneau fibre de $S/R$ en
l'idéal premier $\mathfrak p$. Comme $F$ est de type fini sur
$\kappa(\mathfrak p)$, il est noethérien, et $\Spec(F)$ possède donc un nombre
fini de points fermés isolés. S'il n'existe aucun point fermé isolé,
c'est-à-dire aucun idéal premier $\mathfrak q$ de $S$ au-dessus de
$\mathfrak p$ tel que $S/R$ soit quasi-fini en $\mathfrak q$, alors le lemme
est satisfait. S'il existe au moins un tel idéal premier $\mathfrak q$,
nous pouvons appliquer le
Lemme \ref{algebra-lemma-etale-makes-quasi-finite-finite-one-prime}.
Cela donne un diagramme
$$
\xymatrix{
S \ar[r] & R'\otimes_R S \ar@{=}[r] & A_1 \times B' \\
R \ar[r] \ar[u] & R' \ar[u] \ar[ru]
}
$$
comme dans ledit lemme. Comme les corps résiduels en $\mathfrak p$ et
$\mathfrak p'$ sont les mêmes, les anneaux fibres de $S/R$ et de
$(A_1 \times B')/R'$ sont les mêmes.
Ainsi, par récurrence sur le nombre de points fermés isolés de la fibre, nous
pouvons supposer que le lemme est satisfait pour $R' \to B'$ et
$\mathfrak p'$. Nous obtenons donc un morphisme d'anneaux étale
$R' \to R''$, un idéal premier $\mathfrak p'' \subset R''$ et une
décomposition
$$
R'' \otimes_{R'} B' = A_2 \times \ldots \times A_n \times B
$$
Nous omettons de vérifier que le morphisme d'anneaux $R \to R''$, l'idéal
premier $\mathfrak p''$ et la décomposition qui en résulte
$$
R'' \otimes_R S = (R'' \otimes_{R'} A_1) \times
A_2 \times \ldots \times A_n \times B
$$
constituent une solution au problème posé dans le lemme.
\end{proof}

\begin{lemma}
\label{algebra-lemma-etale-makes-quasi-finite-finite-variant}
Soit $R \to S$ un morphisme d'anneaux.
Soit $\mathfrak p \subset R$ un idéal premier.
Supposons que $R \to S$ soit de type fini.
Alors il existe
\begin{enumerate}
\item un morphisme d'anneaux étale $R \to R'$,
\item un idéal premier $\mathfrak p' \subset R'$ au-dessus de $\mathfrak p$,
\item une décomposition en produit
$$
R' \otimes_R S = A_1 \times \ldots \times A_n \times B
$$
\end{enumerate}
possédant les propriétés suivantes :
\begin{enumerate}
\item chaque $A_i$ est fini sur $R'$,
\item chaque $A_i$ possède exactement un idéal premier $\mathfrak r_i$
au-dessus de $\mathfrak p'$,
\item les extensions de corps finies
$\kappa(\mathfrak r_i)/\kappa(\mathfrak p')$ sont purement inséparables, et
\item $R' \to B$ n'est quasi-fini en aucun idéal premier au-dessus de
$\mathfrak p'$.
\end{enumerate}
\end{lemma}

\begin{proof}
La stratégie de la démonstration consiste à effectuer deux extensions
d'anneaux étales : nous contrôlons d'abord les corps résiduels, puis nous
appliquons le Lemme \ref{algebra-lemma-etale-makes-quasi-finite-finite}.

\medskip\noindent
Notons $F = S \otimes_R \kappa(\mathfrak p)$ l'anneau fibre de $S/R$ en
l'idéal premier $\mathfrak p$.
Comme dans la démonstration du
Lemme \ref{algebra-lemma-etale-makes-quasi-finite-finite}, il existe un nombre fini
d'idéaux premiers, disons $\mathfrak q_1, \ldots, \mathfrak q_n$, de $S$
au-dessus de $R$, auxquels le morphisme d'anneaux
$R \to S$ est quasi-fini.
Soit $\kappa(\mathfrak p) \subset L_i \subset \kappa(\mathfrak q_i)$ le
sous-corps tel que $\kappa(\mathfrak p) \subset L_i$ soit séparable et que
l'extension de corps $\kappa(\mathfrak q_i)/L_i$ soit purement inséparable.
Soit $L/\kappa(\mathfrak p)$ une extension galoisienne finie dans laquelle
$L_i$ se plonge pour $i = 1, \ldots, n$.
Par le Lemme \ref{algebra-lemma-make-etale-map-prescribed-residue-field}, nous pouvons
trouver une extension d'anneaux étale $R \to R'$, ainsi qu'un idéal premier
$\mathfrak p'$ au-dessus de $\mathfrak p$, tels que l'extension de corps
$\kappa(\mathfrak p')/\kappa(\mathfrak p)$ soit isomorphe à
$\kappa(\mathfrak p) \subset L$.
Ainsi, l'anneau fibre de $R' \otimes_R S$ en $\mathfrak p'$ est isomorphe à
$F \otimes_{\kappa(\mathfrak p)} L$.
Les idéaux premiers au-dessus de $\mathfrak q_i$ correspondent aux idéaux
premiers de
$\kappa(\mathfrak q_i) \otimes_{\kappa(\mathfrak p)} L$, qui est un produit
de corps purement inséparables sur $L$, par notre choix de $L$ et la théorie
élémentaire des corps.
Ce sont également les seuls idéaux premiers au-dessus de $\mathfrak p'$
auxquels $R' \to R' \otimes_R S$ est quasi-fini, par le
Lemme \ref{algebra-lemma-quasi-finite-base-change}.
Après avoir remplacé $R$ par $R'$, $\mathfrak p$ par $\mathfrak p'$ et
$S$ par $R' \otimes_R S$, nous pouvons donc supposer que, pour tous les idéaux
premiers $\mathfrak q$ au-dessus de $\mathfrak p$ auxquels $S/R$ est
quasi-fini, les extensions de corps
$\kappa(\mathfrak q)/\kappa(\mathfrak p)$ sont purement inséparables.

\medskip\noindent
Appliquons ensuite le Lemme \ref{algebra-lemma-etale-makes-quasi-finite-finite}.
Le résultat est celui que nous souhaitons, puisque les extensions de corps ne
changent pas sous cette extension d'anneaux étale.
\end{proof}






\section{Homomorphismes locaux}
\label{algebra-section-local-homomorphisms}

\noindent
Quelques lemmes qui ne trouvent naturellement place dans aucune section.
Le premier affirme, en termes vagues, qu'un morphisme étale d'anneaux locaux
est un isomorphisme modulo toutes les puissances d'un idéal principal engendré
par un élément non inversible.

\begin{lemma}
\label{algebra-lemma-lindel}
\begin{reference}
\cite[Lemme à la page 321]{Lindel}, \cite[Lemme 4.1.5]{KC}
\end{reference}
Soit $(R, \mathfrak m_R) \to (S, \mathfrak m_S)$ un homomorphisme local
d'anneaux locaux. Supposons que $S$ soit la localisation d'une extension
d'anneaux étale de $R$ et que
$\kappa(\mathfrak m_R) \to \kappa(\mathfrak m_S)$ soit un isomorphisme.
Alors il existe $t \in \mathfrak m_R$ tel que
$R/t^nR \to S/t^nS$ soit un isomorphisme pour tout $n \geq 1$.
\end{lemma}

\begin{proof}
Écrivons $S = T_{\mathfrak q}$ pour une certaine $R$-algèbre étale $T$ et un
idéal premier $\mathfrak q \subset T$ au-dessus de $\mathfrak m_R$.
Par la Proposition \ref{algebra-proposition-etale-locally-standard}, nous pouvons
supposer que $R \to T$ est étale standard.
Écrivons $T = R[x]_g/(f)$ comme dans la
Définition \ref{algebra-definition-standard-etale}.
Par notre hypothèse sur les corps résiduels, nous pouvons choisir
$a \in R$ tel que $x$ et $a$ aient la même image dans
$\kappa(\mathfrak q) = \kappa(\mathfrak m_S) = \kappa(\mathfrak m_R)$.
Après avoir remplacé $x$ par $x - a$, nous pouvons alors supposer que
$\mathfrak q$ est engendré par $x$ et $\mathfrak m_R$ dans $T$.
En particulier, $t = f(0) \in \mathfrak m_R$.
Nous allons montrer que $t = f(0)$ convient.

\medskip\noindent
Écrivons $f = x^d + \sum_{i = 1, \ldots, d - 1} a_i x^i + t$.
Comme $R \to T$ est étale standard, nous trouvons que $a_1$ est une unité de
$R$ : la dérivée de $f$ est inversible dans $T$ et, en particulier,
n'appartient pas à $\mathfrak q$.
Soit
$h = a_1 + a_2 x + \ldots + a_{d - 1} x^{d - 2} + x^{d - 1} \in R[x]$,
de sorte que $f = t + xh$ dans $R[x]$.
Nous voyons que $h \not \in \mathfrak q$ et pouvons donc remplacer $T$ par
$R[x]_{hg}/(f)$. Après ce remplacement, nous voyons que
$$
T/tT = (R/tR)[x]_{hg}/(f) = (R/tR)[x]_{hg}/(xh) = (R/tR)[x]_{hg}/(x)
$$
est un quotient de $R/tR$.
Par le Lemme \ref{algebra-lemma-surjective-mod-locally-nilpotent}, nous concluons que
$R/t^nR \to T/t^nT$ est surjectif pour tout $n \geq 1$.
D'autre part, nous savons que le morphisme local plat d'anneaux
$R/t^nR \to S/t^nS$ se factorise par $R/t^nR \to T/t^nT$ pour tout $n$ ;
ces morphismes sont donc également injectifs (un homomorphisme local plat
d'anneaux locaux est fidèlement plat et donc injectif ; voir les
Lemmes \ref{algebra-lemma-local-flat-ff} et
\ref{algebra-lemma-faithfully-flat-universally-injective}).
Comme $S$ est la localisation de $T$, nous voyons que $S/t^nS$ est la
localisation de $T/t^nT = R/t^nR$ en un idéal premier au-dessus de l'idéal
maximal, mais cet anneau est déjà local, ce qui achève la démonstration.
\end{proof}

\begin{lemma}
\label{algebra-lemma-etale-under-finite-flat}
Soit $(R, \mathfrak m_R) \to (S, \mathfrak m_S)$ un homomorphisme local
d'anneaux locaux. Supposons que $S$ soit la localisation d'une extension
d'anneaux étale de $R$. Alors il existe un morphisme d'anneaux
fini, de présentation finie et fidèlement plat $R \to S'$ tel que, pour tout
idéal maximal $\mathfrak m'$ de $S'$, il existe une factorisation
$$
R \to S \to S'_{\mathfrak m'}.
$$
Il s'agit d'une factorisation du morphisme d'anneaux
$R \to S'_{\mathfrak m'}$.
\end{lemma}

\begin{proof}
Écrivons $S = T_{\mathfrak q}$ pour une certaine $R$-algèbre étale $T$.
Par la Proposition \ref{algebra-proposition-etale-locally-standard}, nous pouvons
supposer que $T$ est étale standard.
Appliquons le Lemme \ref{algebra-lemma-standard-etale-finite-flat-Zariski} au
morphisme d'anneaux $R \to T$ pour obtenir $R \to S'$.
Alors, en particulier, pour tout idéal maximal $\mathfrak m'$ de $S'$, nous
obtenons une factorisation $\varphi : T \to S'_{g'}$ pour un certain
$g' \not \in \mathfrak m'$ tel que
$\mathfrak q = \varphi^{-1}(\mathfrak m'S'_{g'})$.
Ainsi, $\varphi$ induit le morphisme local d'anneaux recherché
$S \to S'_{\mathfrak m'}$.
\end{proof}





\section{Clôture intégrale et changement de base lisse}
\label{algebra-section-integral-closure-smooth-base-change}

\begin{lemma}
\label{algebra-lemma-trick}
Soit $R$ un anneau.
Soit $f \in R[x]$ un polynôme unitaire.
Soit $R \to B$ un morphisme d'anneaux.
Si $h \in B[x]/(f)$ est entier sur $R$, alors l'élément $f' h$ peut s'écrire
$f'h = \sum_i b_i x^i$, avec des $b_i \in B$ entiers sur $R$.
\end{lemma}

\begin{proof}
Disons que $h^e + r_1 h^{e - 1} + \ldots + r_e = 0$ dans l'anneau
$B[x]/(f)$, avec $r_i \in R$.
Il existe une extension d'anneaux finie libre $B \subset B'$ telle que
$f = (x - \alpha_1) \ldots (x - \alpha_d)$ pour certains
$\alpha_i \in B'$ ; voir le Lemme \ref{algebra-lemma-adjoin-roots}.
Remarquons que chaque $\alpha_i$ est entier sur $R$.
Nous pouvons représenter
$h = h_0 + h_1 x + \ldots + h_{d - 1} x^{d - 1}$, avec $h_i \in B$.
Alors le fait universel suivant est valable :
$$
f' h =
\sum\nolimits_{i = 1, \ldots, d}
h(\alpha_i)
(x - \alpha_1) \ldots \widehat{(x - \alpha_i)} \ldots (x - \alpha_d)
$$
dans $B'[x]/(f)$. On le démontre en évaluant les deux membres aux points
$\alpha_i$ sur l'anneau
$B_{univ} = \mathbf{Z}[\alpha_i, h_j]$ (certains détails sont omis).
Comme, par hypothèse, $h$ satisfait
$h^e + r_1 h^{e - 1} + \ldots + r_e = 0$ dans l'anneau $B[x]/(f)$,
nous voyons que
$$
h(\alpha_i)^e + r_1 h(\alpha_i)^{e - 1} + \ldots + r_e = 0
$$
dans $B'$. Ainsi, $h(\alpha_i)$ est entier sur $R$.
En utilisant la formule ci-dessus, nous voyons que
$f'h \equiv \sum_{j = 0, \ldots, d - 1} b'_j x^j$ dans $B'[x]/(f)$,
avec des $b'_j \in B'$ entiers sur $R$.
Cependant, puisque $f' h \in B[x]/(f)$ et que
$1, x, \ldots, x^{d - 1}$ constituent une $B'$-base de $B'[x]/(f)$,
nous voyons que $b'_j \in B$, comme désiré.
\end{proof}

\begin{lemma}
\label{algebra-lemma-integral-closure-commutes-etale}
Soit $R \to S$ un morphisme d'anneaux étale.
Soit $R \to B$ un morphisme d'anneaux quelconque.
Soit $A \subset B$ la clôture intégrale de $R$ dans $B$.
Soit $A' \subset S \otimes_R B$ la clôture intégrale de $S$ dans
$S \otimes_R B$. Alors le morphisme canonique
$S \otimes_R A \to A'$ est un isomorphisme.
\end{lemma}

\begin{proof}
Le morphisme $S \otimes_R A \to A'$ est injectif parce que
$A \subset B$ et que $R \to S$ est plat.
Nous allons utiliser à plusieurs reprises le fait que la formation de la
clôture intégrale commute à la localisation ; voir le
Lemme \ref{algebra-lemma-integral-closure-localize}.
Nous pouvons donc localiser sur $S$, par le Lemme \ref{algebra-lemma-cover}
(le critère qui permet de vérifier si un morphisme de $S$-modules est un
isomorphisme).
Ainsi, nous pouvons supposer que
$S = R[x]_g/(f) = (R[x]/(f))_g$ est étale standard sur $R$ ;
voir la Proposition \ref{algebra-proposition-etale-locally-standard}.
En localisant une fois de plus, nous voyons que $A'$ est $(A'')_g$, où
$A''$ est la clôture intégrale de $R[x]/(f)$ dans $B[x]/(f)$.
Supposons que $a \in A''$. Il suffit de montrer que $a$ appartient à
$S \otimes_R A$. Par le Lemme \ref{algebra-lemma-trick}, nous voyons que
$f' a = \sum a_i x^i$, avec $a_i \in A$.
Comme $f'$ est inversible dans $S$ (par définition d'un morphisme d'anneaux
étale standard), nous concluons que $a \in S \otimes_R A$, comme désiré.
\end{proof}

\begin{example}
\label{algebra-example-fourier}
Soit $p$ un nombre premier. L'extension d'anneaux
$$
R = \mathbf{Z}[1/p] \subset
R' = \mathbf{Z}[1/p][x]/(x^{p - 1} + \ldots + x + 1)
$$
possède la propriété suivante : pour $d < p$, il existe des éléments
$\alpha_0, \ldots, \alpha_{d - 1} \in R'$ tels que
$$
\prod\nolimits_{0 \leq i < j < d} (\alpha_i - \alpha_j)
$$
soit une unité de $R'$. En effet, prenons $\alpha_i$ égal à la classe de
$x^i$ dans $R'$ pour $i = 0, \ldots, p - 1$. Nous avons alors
$$
T^p - 1 = \prod\nolimits_{i = 0, \ldots, p - 1} (T - \alpha_i)
$$
dans $R'[T]$. En effet, l'anneau
$\mathbf{Q}[x]/(x^{p - 1} + \ldots + x + 1)$ est un corps parce que le
polynôme cyclotomique $x^{p - 1} + \ldots + x + 1$ est irréductible sur
$\mathbf{Q}$ et que les $\alpha_i$ sont des racines deux à deux distinctes de
$T^p - 1$, d'où l'égalité.
En dérivant les deux membres puis en substituant $T = \alpha_i$, nous obtenons
$$
p \alpha_i^{p - 1}
=
(\alpha_i - \alpha_1) \ldots
\widehat{(\alpha_i - \alpha_i)} \ldots
(\alpha_i - \alpha_1)
$$
et nous voyons que cet élément est inversible dans $R'$.
\end{example}

\begin{lemma}
\label{algebra-lemma-integral-closure-commutes-smooth}
Soit $R \to S$ un morphisme d'anneaux lisse.
Soit $R \to B$ un morphisme d'anneaux quelconque.
Soit $A \subset B$ la clôture intégrale de $R$ dans $B$.
Soit $A' \subset S \otimes_R B$ la clôture intégrale de $S$ dans
$S \otimes_R B$. Alors le morphisme canonique
$S \otimes_R A \to A'$ est un isomorphisme.
\end{lemma}

\begin{proof}
En raisonnant comme dans la démonstration du
Lemme \ref{algebra-lemma-integral-closure-commutes-etale}, nous pouvons localiser sur
$S$. Nous pouvons donc supposer que $R \to S$ est un morphisme d'anneaux lisse
standard ; voir le Lemme \ref{algebra-lemma-smooth-syntomic}.
Par définition d'un morphisme d'anneaux lisse standard, nous voyons que $S$
est étale sur un anneau de polynômes $R[x_1, \ldots, x_n]$.
Comme nous avons vu le résultat dans le cas d'une extension d'anneaux étale
(Lemme \ref{algebra-lemma-integral-closure-commutes-etale}), cela nous ramène au cas
où $S = R[x]$. Nous devons donc montrer
$$
f = \sum b_i x^i
\text{ entier sur }R[x]
\Leftrightarrow
\text{chaque }b_i\text{ est entier sur }R.
$$
L'implication de droite à gauche vaut parce que l'ensemble des éléments de
$B[x]$ entiers sur $R[x]$ est un anneau
(Lemme \ref{algebra-lemma-integral-closure-is-ring}) et contient $x$.

\medskip\noindent
Supposons que $f \in B[x]$ soit entier sur $R[x]$ et supposons que
$f = \sum_{i < d} b_i x^i$ soit de degré $< d$.
Comme la clôture intégrale et la localisation commutent, il suffit de montrer
qu'il existe des nombres premiers distincts $p, q$ tels que chaque $b_i$ soit
entier à la fois sur $R[1/p]$ et sur $R[1/q]$.
Nous pouvons donc trouver une extension d'anneaux finie libre
$R \subset R'$ telle que $R'$ contienne
$\alpha_1, \ldots, \alpha_d$ et que
$\prod_{i < j} (\alpha_i - \alpha_j)$ soit une unité de $R'$ ; voir
l'Exemple \ref{algebra-example-fourier}.
Dans ce cas, nous avons l'égalité universelle
$$
f
=
\sum_i
f(\alpha_i)
\frac{(x - \alpha_1) \ldots \widehat{(x - \alpha_i)} \ldots (x - \alpha_d)}
{(\alpha_i - \alpha_1) \ldots \widehat{(\alpha_i - \alpha_i)} \ldots
(\alpha_i - \alpha_d)}.
$$
De plus, les éléments $f(\alpha_i)$ sont entiers sur $R'$, puisque
$(R' \otimes_R B)[x] \to R' \otimes_R B$, $h \mapsto h(\alpha_i)$,
est un morphisme d'anneaux.
Nous voyons donc que les coefficients de $f$ dans
$(R' \otimes_R B)[x]$ sont entiers sur $R'$.
Comme $R'$ est fini sur $R$ (donc entier sur $R$), nous voyons qu'ils sont
également entiers sur $R$, comme désiré.
\end{proof}

\begin{lemma}
\label{algebra-lemma-integral-closure-commutes-colim-smooth}
Soient $R \to S$ et $R \to B$ des morphismes d'anneaux.
Soit $A \subset B$ la clôture intégrale de $R$ dans $B$.
Soit $A' \subset S \otimes_R B$ la clôture intégrale de $S$ dans
$S \otimes_R B$. Si $S$ est une colimite filtrante de $R$-algèbres lisses,
alors le morphisme canonique $S \otimes_R A \to A'$ est un isomorphisme.
\end{lemma}

\begin{proof}
Cela résulte du fait immédiat que la formation des produits tensoriels et
celle des clôtures intégrales commutent aux colimites filtrantes, ainsi que du
Lemme \ref{algebra-lemma-integral-closure-commutes-smooth}.
\end{proof}






\section{Morphismes formellement non ramifiés}
\label{algebra-section-formally-unramified}

\noindent
Il s'avère logiquement plus efficace de définir la notion de morphisme
formellement non ramifié avant d'introduire celle de morphisme formellement
étale.

\begin{definition}
\label{algebra-definition-formally-unramified}
Soit $R \to S$ un morphisme d'anneaux.
Nous disons que $S$ est {\it formellement non ramifié sur $R$} si, pour tout
diagramme commutatif dont les flèches pleines sont données
$$
\xymatrix{
S \ar[r] \ar@{-->}[rd] & A/I \\
R \ar[r] \ar[u] & A \ar[u]
}
$$
où $I \subset A$ est un idéal de carré nul, il existe au plus une flèche
pointillée rendant le diagramme commutatif.
\end{definition}

\begin{lemma}
\label{algebra-lemma-base-change-formally-unramified}
Soit $R \to S$ un morphisme formellement non ramifié.
Soit $R \to R'$ un morphisme d'anneaux quelconque.
Alors le changement de base $S' = R' \otimes_R S$ est formellement non
ramifié sur $R'$.
\end{lemma}

\begin{proof}
Soit donné un diagramme dont les flèches pleines sont données
$$
\xymatrix{
S \ar[r] \ar@{-->}[rrd] & R' \otimes_R S \ar[r] \ar@{-->}[rd] & A/I \\
R  \ar[u] \ar[r] & R' \ar[r] \ar[u] & A \ar[u]
}
$$
comme dans la Définition \ref{algebra-definition-formally-unramified}.
Par hypothèse, il existe au plus une flèche pointillée longue.
Par la propriété universelle du produit tensoriel, nous concluons qu'il existe
au plus une flèche pointillée courte.
\end{proof}

\begin{lemma}
\label{algebra-lemma-characterize-formally-unramified}
Soit $R \to S$ un morphisme d'anneaux.
Les conditions suivantes sont équivalentes :
\begin{enumerate}
\item $R \to S$ est formellement non ramifié,
\item le module des différentielles $\Omega_{S/R}$ est nul.
\end{enumerate}
\end{lemma}

\begin{proof}
Soit $J = \Ker(S \otimes_R S \to S)$ le noyau du morphisme de multiplication.
Soit $A_{univ} = S \otimes_R S/J^2$. Rappelons que
$I_{univ} = J/J^2$ est isomorphe à $\Omega_{S/R}$ ; voir le
Lemme \ref{algebra-lemma-differentials-diagonal}.
De plus, les deux morphismes de $R$-algèbres
$\sigma_1, \sigma_2 : S \to A_{univ}$,
$\sigma_1(s) = s \otimes 1 \bmod J^2$ et
$\sigma_2(s) = 1 \otimes s \bmod J^2$ diffèrent par la dérivation universelle
$\text{d} : S \to \Omega_{S/R} = I_{univ}$.

\medskip\noindent
Supposons que $R \to S$ soit formellement non ramifié.
Nous voyons alors que $\sigma_1 = \sigma_2$.
Ainsi, $\text{d}(s) = 0$ pour tout $s \in S$.
Par conséquent, $\Omega_{S/R} = 0$.

\medskip\noindent
Supposons que $\Omega_{S/R} = 0$.
Soient $A, I, R \to A, S \to A/I$ un diagramme dont les flèches pleines sont
données, comme dans la Définition \ref{algebra-definition-formally-unramified}.
Soient $\tau_1, \tau_2 : S \to A$ deux flèches pointillées rendant le
diagramme commutatif. Considérons le morphisme de $R$-algèbres
$A_{univ} \to A$ défini par la règle
$s_1 \otimes s_2 \mapsto \tau_1(s_1)\tau_2(s_2)$.
Nous omettons de vérifier qu'il est bien défini.
Comme $A_{univ} \cong S$, puisque
$I_{univ} = \Omega_{S/R} = 0$, nous concluons que $\tau_1 = \tau_2$.
\end{proof}

\begin{lemma}
\label{algebra-lemma-formally-unramified-local}
Soit $R \to S$ un morphisme d'anneaux.
Les conditions suivantes sont équivalentes :
\begin{enumerate}
\item $R \to S$ est formellement non ramifié,
\item $R \to S_{\mathfrak q}$ est formellement non ramifié pour tout idéal
premier $\mathfrak q$ de $S$, et
\item $R_{\mathfrak p} \to S_{\mathfrak q}$ est formellement non ramifié pour
tout idéal premier $\mathfrak q$ de $S$, avec
$\mathfrak p = R \cap \mathfrak q$.
\end{enumerate}
\end{lemma}

\begin{proof}
Nous avons vu dans le
Lemme \ref{algebra-lemma-characterize-formally-unramified} que (1) équivaut à
$\Omega_{S/R} = 0$. De même, par le
Lemme \ref{algebra-lemma-differentials-localize}, nous voyons que (2) et (3)
équivalent à $(\Omega_{S/R})_{\mathfrak q} = 0$ pour tout
$\mathfrak q$. L'équivalence résulte donc du
Lemme \ref{algebra-lemma-characterize-zero-local}.
\end{proof}

\begin{lemma}
\label{algebra-lemma-formally-unramified-localize}
Soit $A \to B$ un morphisme d'anneaux formellement non ramifié.
\begin{enumerate}
\item Pour une partie multiplicative $S \subset A$,
$S^{-1}A \to S^{-1}B$ est formellement non ramifié.
\item Pour une partie multiplicative $S \subset B$,
$A \to S^{-1}B$ est formellement non ramifié.
\end{enumerate}
\end{lemma}

\begin{proof}
Cela résulte du Lemme \ref{algebra-lemma-formally-unramified-local}.
(On peut aussi le déduire du
Lemme \ref{algebra-lemma-characterize-formally-unramified}, combiné au
Lemme \ref{algebra-lemma-differentials-localize}.)
\end{proof}

\begin{lemma}
\label{algebra-lemma-colimit-formally-unramified}
Soit $R$ un anneau. Soit $I$ un ensemble ordonné filtrant.
Soit $(S_i, \varphi_{ii'})$ un système de $R$-algèbres sur $I$.
Si chaque $R \to S_i$ est formellement non ramifié, alors
$S = \colim_{i \in I} S_i$ est formellement non ramifié sur $R$
\end{lemma}

\begin{proof}
Considérons un diagramme comme dans la
Définition \ref{algebra-definition-formally-unramified}.
Par hypothèse, il existe au plus un morphisme de $R$-algèbres
$S_i \to A$ relevant les composées $S_i \to S \to A/I$.
Comme tout élément de $S$ appartient à l'image de l'un des morphismes
$S_i \to S$, nous voyons qu'il existe au plus un morphisme $S \to A$
s'insérant dans le diagramme.
\end{proof}



\section{Modules conormaux et épaississements universels}
\label{algebra-section-conormal}

\noindent
Il s'avère que l'on peut définir le premier voisinage infinitésimal non
seulement pour une immersion fermée de schémas, mais déjà pour tout morphisme
formellement non ramifié. Cela repose sur le fait algébrique suivant.

\begin{lemma}
\label{algebra-lemma-universal-thickening}
Soit $R \to S$ un morphisme d'anneaux formellement non ramifié.
Il existe une surjection de $R$-algèbres $S' \to S$ dont le noyau est un idéal
de carré nul et qui possède la propriété universelle suivante : étant donné
un diagramme commutatif quelconque
$$
\xymatrix{
S \ar[r]_a & A/I \\
R \ar[r]^b \ar[u] & A \ar[u]
}
$$
où $I \subset A$ est un idéal de carré nul, il existe un unique morphisme de
$R$-algèbres $a' : S' \to A$ tel que $S' \to A \to A/I$ soit égal à
$S' \to S \to A/I$.
\end{lemma}

\begin{proof}
Choisissons un ensemble de générateurs $z_i \in S$, $i \in I$, de $S$ comme
$R$-algèbre.
Soit $P = R[\{x_i\}_{i \in I}]$ l'anneau de polynômes en les générateurs
$x_i$, $i \in I$. Considérons le morphisme de $R$-algèbres $P \to S$ qui
envoie $x_i$ sur $z_i$. Soit $J = \Ker(P \to S)$.
Considérons le morphisme
$$
\text{d} : J/J^2 \longrightarrow \Omega_{P/R} \otimes_P S
$$
voir le Lemme \ref{algebra-lemma-differential-seq}.
Il est surjectif puisque $\Omega_{S/R} = 0$ par hypothèse ; voir le
Lemme \ref{algebra-lemma-characterize-formally-unramified}.
Remarquons que $\Omega_{P/R}$ est libre sur les $\text{d}x_i$, et que le
module $\Omega_{P/R} \otimes_P S$ est donc libre sur $S$.
Nous pouvons ainsi choisir un scindage de la surjection ci-dessus et écrire
$$
J/J^2 = K \oplus \Omega_{P/R} \otimes_P S
$$
Soit $J^2 \subset J' \subset J$ l'idéal de $P$ tel que $J'/J^2$ soit le
second facteur direct de la décomposition ci-dessus.
Posons $S' = P/J'$. Nous obtenons une suite exacte courte
$$
0 \to J/J' \to S' \to S \to 0
$$
et nous voyons que $J/J' \cong K$ est un idéal de carré nul de $S'$.
Par conséquent,
$$
\xymatrix{
S \ar[r]_1 & S \\
R \ar[r] \ar[u] & S' \ar[u]
}
$$
est un diagramme comme ci-dessus.
En fait, nous affirmons qu'il s'agit d'un objet initial de la catégorie de ces
diagrammes. En effet, soit $(I \subset A, a, b)$ un diagramme arbitraire.
Nous pouvons choisir un morphisme de $R$-algèbres $\beta : P \to A$ tel que
$$
\xymatrix{
S \ar[r]_1 & S \ar[r]_a & A/I \\
R \ar[r] \ar@/_/[rr]_b \ar[u] & P \ar[u] \ar[r]^\beta & A \ar[u]
}
$$
soit commutatif.
Il n'est pas nécessairement vrai que $\beta(J') = 0$ ; autrement dit,
$\beta$ ne se factorise pas nécessairement par $S' = P/J'$.
En revanche, il est clair que $\beta(J') \subset I$ et, puisque
$\beta(J) \subset I$ et $I^2 = 0$, nous avons $\beta(J^2) = 0$.
Ainsi, l'« obstruction » à l'existence d'un morphisme de
$(J/J' \subset S', 1, R \to S')$ vers $(I \subset A, a, b)$ est le morphisme
$S$-linéaire correspondant $\overline{\beta} : J'/J^2 \to I$.
Le choix de $\beta$ réside dans le choix des $\beta(x_i)$.
Un choix différent de $\beta$, disons $\beta'$, s'obtient en prenant
$\beta'(x_i) = \beta(x_i) + \delta_i$, avec $\delta_i \in I$.
Dans ce cas, pour $g \in J'$, nous obtenons
$$
\beta'(g) =
\beta(g) + \sum\nolimits_i \delta_i \beta(\frac{\partial g}{\partial x_i})
$$
Comme le morphisme
$\text{d}|_{J'/J^2} : J'/J^2 \to \Omega_{P/R} \otimes_P S$ donné par
$g \mapsto \frac{\partial g}{\partial x_i}\text{d}x_i \otimes 1$ est un
isomorphisme par construction, nous voyons qu'il existe un choix unique de
$\delta_i \in I$ tel que $\beta'(g) = 0$ pour tout $g \in J'$.
(Plus précisément, $\delta_i$ est $-\overline{\beta}(g)$, où
$g \in J'/J^2$ est l'unique élément envoyé sur
$\text{d}x_i \otimes 1$ par $\text{d}|_{J'/J^2}$.)
L'unicité de la solution implique celle requise dans le lemme.
\end{proof}

\noindent
Dans la situation du Lemme \ref{algebra-lemma-universal-thickening}, le morphisme de
$R$-algèbres $S' \to S$ est unique à isomorphisme unique près.

\begin{definition}
\label{algebra-definition-universal-thickening}
Soit $R \to S$ un morphisme d'anneaux formellement non ramifié.
\begin{enumerate}
\item L'{\it épaississement universel du premier ordre} de $S$ sur $R$ est la
surjection de $R$-algèbres $S' \to S$ du
Lemme \ref{algebra-lemma-universal-thickening}.
\item Le {\it module conormal} de $R \to S$ est le noyau $I$ de
l'épaississement universel du premier ordre $S' \to S$, considéré comme un
$S$-module.
\end{enumerate}
Dans cette situation, nous notons souvent le module conormal {\it $C_{S/R}$}.
\end{definition}

\begin{lemma}
\label{algebra-lemma-universal-thickening-quotient}
Soit $I \subset R$ un idéal d'un anneau.
L'épaississement universel du premier ordre de $R/I$ sur $R$ est la
surjection $R/I^2 \to R/I$. Le module conormal de $R/I$ sur $R$ est
$C_{(R/I)/R} = I/I^2$.
\end{lemma}

\begin{proof}
Omis.
\end{proof}

\begin{lemma}
\label{algebra-lemma-universal-thickening-localize}
Soit $A \to B$ un morphisme d'anneaux formellement non ramifié.
Soit $\varphi : B' \to B$ l'épaississement universel du premier ordre de
$B$ sur $A$.
\begin{enumerate}
\item Soit $S \subset A$ une partie multiplicative.
Alors $S^{-1}B' \to S^{-1}B$ est l'épaississement universel du premier ordre
de $S^{-1}B$ sur $S^{-1}A$. En particulier,
$S^{-1}C_{B/A} = C_{S^{-1}B/S^{-1}A}$.
\item Soit $S \subset B$ une partie multiplicative.
Alors $S' = \varphi^{-1}(S)$ est une partie multiplicative de $B'$, et
$(S')^{-1}B' \to S^{-1}B$ est l'épaississement universel du premier ordre de
$S^{-1}B$ sur $A$. En particulier,
$S^{-1}C_{B/A} = C_{S^{-1}B/A}$.
\end{enumerate}
Remarquons que le lemme a un sens en vertu du
Lemme \ref{algebra-lemma-formally-unramified-localize}.
\end{lemma}

\begin{proof}
Avec les notations et hypothèses de (1), soit
$(S^{-1}B)' \to S^{-1}B$ l'épaississement universel du premier ordre de
$S^{-1}B$ sur $S^{-1}A$.
Remarquons que $S^{-1}B' \to S^{-1}B$ est une surjection de
$S^{-1}A$-algèbres dont le noyau est de carré nul.
Par définition, nous obtenons donc un morphisme
$(S^{-1}B)' \to S^{-1}B'$ compatible avec les morphismes vers $S^{-1}B$.
Considérons un diagramme commutatif quelconque
$$
\xymatrix{
B \ar[r] & S^{-1}B \ar[r] & D/I \\
A \ar[r] \ar[u] & S^{-1}A \ar[r] \ar[u] & D \ar[u]
}
$$
où $I \subset D$ est un idéal de carré nul.
Comme $B'$ est l'épaississement universel du premier ordre de $B$ sur $A$,
nous obtenons un morphisme de $A$-algèbres $B' \to D$.
Mais il est clair que l'image de $S$ dans $D$ est envoyée sur des éléments
inversibles de $D$, et nous obtenons donc un morphisme compatible
$S^{-1}B' \to D$.
En appliquant cela à $D = (S^{-1}B)'$, nous voyons que nous obtenons un
morphisme $S^{-1}B' \to (S^{-1}B)'$.
Nous omettons de vérifier que ce morphisme est inverse du morphisme décrit
ci-dessus.

\medskip\noindent
Avec les notations et hypothèses de (2), soit
$(S^{-1}B)' \to S^{-1}B$ l'épaississement universel du premier ordre de
$S^{-1}B$ sur $A$.
Remarquons que $(S')^{-1}B' \to S^{-1}B$ est une surjection de $A$-algèbres
dont le noyau est de carré nul.
Par définition, nous obtenons donc un morphisme
$(S^{-1}B)' \to (S')^{-1}B'$ compatible avec les morphismes vers
$S^{-1}B$. Considérons un diagramme commutatif quelconque
$$
\xymatrix{
B \ar[r] & S^{-1}B \ar[r] & D/I \\
A \ar[r] \ar[u] & A \ar[r] \ar[u] & D \ar[u]
}
$$
où $I \subset D$ est un idéal de carré nul.
Comme $B'$ est l'épaississement universel du premier ordre de $B$ sur $A$,
nous obtenons un morphisme de $A$-algèbres $B' \to D$.
Mais il est clair que l'image de $S'$ dans $D$ est envoyée sur des éléments
inversibles de $D$, et nous obtenons donc un morphisme compatible
$(S')^{-1}B' \to D$.
En appliquant cela à $D = (S^{-1}B)'$, nous voyons que nous obtenons un
morphisme $(S')^{-1}B' \to (S^{-1}B)'$.
Nous omettons de vérifier que ce morphisme est inverse du morphisme décrit
ci-dessus.
\end{proof}

\begin{lemma}
\label{algebra-lemma-differentials-universal-thickening}
Soient $R \to A  \to B$ des morphismes d'anneaux.
Supposons que $A \to B$ soit formellement non ramifié.
Soit $B' \to B$ l'épaississement universel du premier ordre de $B$ sur $A$.
Alors $B'$ est formellement non ramifié sur $A$, et le morphisme canonique
$\Omega_{A/R} \otimes_A B \to \Omega_{B'/R} \otimes_{B'} B$ est un
isomorphisme.
\end{lemma}

\begin{proof}
Nous allons utiliser la construction de $B'$ donnée dans la démonstration du
Lemme \ref{algebra-lemma-universal-thickening}, bien qu'il doive en principe être
possible de déduire formellement ces résultats de la définition.
Plus précisément, nous choisissons une présentation
$B = P/J$, où $P = A[x_i]$ est un anneau de polynômes sur $A$.
Ensuite, nous choisissons des éléments $f_i \in J$ tels que
$\text{d}f_i = \text{d}x_i \otimes 1$ dans
$\Omega_{P/A} \otimes_P B$.
Ces choix étant faits, nous avons $B' = P/J'$, avec
$J' = (f_i) + J^2$ ; voir la démonstration du
Lemme \ref{algebra-lemma-universal-thickening}.

\medskip\noindent
Considérons la suite exacte canonique
$$
J'/(J')^2 \to \Omega_{P/A} \otimes_P B' \to \Omega_{B'/A} \to 0
$$
voir le Lemme \ref{algebra-lemma-differential-seq}.
Par construction, les classes des $f_i \in J'$ s'envoient sur des éléments
du module $\Omega_{P/A} \otimes_P B'$ qui l'engendrent modulo
$J'/J^2$, par construction. Comme $J'/J^2$ est un idéal nilpotent, nous voyons
que ces éléments engendrent le module tout entier (par le
Lemme de Nakayama \ref{algebra-lemma-NAK}).
Cela démontre que $\Omega_{B'/A} = 0$ et donc que $B'$ est formellement non
ramifié sur $A$ ; voir le
Lemme \ref{algebra-lemma-characterize-formally-unramified}.

\medskip\noindent
Comme $P$ est un anneau de polynômes sur $A$, nous avons
$\Omega_{P/R} = \Omega_{A/R} \otimes_A P \oplus \bigoplus P\text{d}x_i$.
Nous allons utiliser cette décomposition.
Considérons la suite exacte suivante :
$$
J'/(J')^2 \to
\Omega_{P/R} \otimes_P B' \to
\Omega_{B'/R} \to 0
$$
voir le Lemme \ref{algebra-lemma-differential-seq}.
Nous pouvons la tensoriser avec $B$ et obtenir la suite exacte
$$
J'/(J')^2 \otimes_{B'} B \to
\Omega_{P/R} \otimes_P B \to
\Omega_{B'/R} \otimes_{B'} B \to 0
$$
Si nous nous rappelons que $J' = (f_i) + J^2$, nous voyons que le premier
morphisme annule le sous-module $J^2/(J')^2$.
En termes de la décomposition en somme directe
$\Omega_{P/R} \otimes_P B =
\Omega_{A/R} \otimes_A B \oplus \bigoplus B\text{d}x_i $ donnée,
nous voyons que le sous-module
$(f_i)/(J')^2 \otimes_{B'} B$ s'envoie isomorphiquement sur le facteur direct
$\bigoplus B\text{d}x_i$. Ce qui reste de cette suite exacte est donc un
isomorphisme
$\Omega_{A/R} \otimes_A B \to \Omega_{B'/R} \otimes_{B'} B$,
comme désiré.
\end{proof}









\section{Morphismes formellement \'etales}
\label{algebra-section-formally-etale}

\begin{definition}
\label{algebra-definition-formally-etale}
Soit $R \to S$ un morphisme d'anneaux.
On dit que $S$ est {\it formellement \'etale sur $R$} si, pour tout
diagramme commutatif dont les flèches pleines sont données
$$
\xymatrix{
S \ar[r] \ar@{-->}[rd] & A/I \\
R \ar[r] \ar[u] & A \ar[u]
}
$$
où $I \subset A$ est un idéal de carré nul, il existe
une unique flèche pointillée rendant le diagramme commutatif.
\end{definition}

\noindent
Il est clair qu'un morphisme d'anneaux est formellement \'etale si et seulement s'il
est à la fois formellement lisse et formellement non ramifié.

\begin{lemma}
\label{algebra-lemma-base-change-formally-etale}
Soit $R \to S$ un morphisme formellement \'etale.
Soit $R \to R'$ un morphisme d'anneaux quelconque.
Alors le changement de base $S' = R' \otimes_R S$ est formellement \'etale sur $R'$.
\end{lemma}

\begin{proof}
Cela résulte des lemmes \ref{algebra-lemma-base-change-fs} et
\ref{algebra-lemma-base-change-formally-unramified}.
\end{proof}

\begin{lemma}
\label{algebra-lemma-formally-etale-etale}
Soit $R \to S$ un morphisme d'anneaux de présentation finie.
Les assertions suivantes sont équivalentes :
\begin{enumerate}
\item $R \to S$ est formellement \'etale,
\item $R \to S$ est \'etale.
\end{enumerate}
\end{lemma}

\begin{proof}
Supposons que $R \to S$ soit formellement \'etale.
Alors $R \to S$ est lisse par la
proposition \ref{algebra-proposition-smooth-formally-smooth}.
D'après le lemme \ref{algebra-lemma-characterize-formally-unramified},
nous avons $\Omega_{S/R} = 0$.
Ainsi, $R \to S$ est \'etale par définition.

\medskip\noindent
Supposons que $R \to S$ soit \'etale.
Alors $R \to S$ est formellement lisse par la
proposition \ref{algebra-proposition-smooth-formally-smooth}.
D'après le lemme \ref{algebra-lemma-characterize-formally-unramified},
il est formellement non ramifié. Ainsi, $R \to S$ est formellement \'etale.
\end{proof}

\begin{lemma}
\label{algebra-lemma-colimit-formally-etale}
Soit $R$ un anneau. Soit $I$ un ensemble ordonné filtrant.
Soit $(S_i, \varphi_{ii'})$ un système de $R$-algèbres
indexé par $I$. Si chaque $R \to S_i$ est formellement \'etale, alors
$S = \colim_{i \in I} S_i$ est formellement \'etale sur $R$.
\end{lemma}

\begin{proof}
Considérons un diagramme comme dans la définition \ref{algebra-definition-formally-etale}.
Par hypothèse, nous obtenons d'uniques morphismes de $R$-algèbres $S_i \to A$ relevant
les composées $S_i \to S \to A/I$. Ils sont donc compatibles
avec les morphismes de transition $\varphi_{ii'}$ et définissent un relèvement
$S \to A$. Cela prouve l'existence.
L'unicité est claire en se restreignant à chaque $S_i$.
\end{proof}

\begin{lemma}
\label{algebra-lemma-localization-formally-etale}
Soit $R$ un anneau. Soit $S \subset R$ une partie multiplicative quelconque.
Alors le morphisme d'anneaux $R \to S^{-1}R$ est formellement \'etale.
\end{lemma}

\begin{proof}
Soit $I \subset A$ un idéal de carré nul. Ce que nous affirmons
ici est qu'étant donné un morphisme d'anneaux $\varphi : R \to A$ tel que
$\varphi(f) \mod I$ soit inversible pour tout $f \in S$, nous avons aussi que
$\varphi(f)$ est inversible dans $A$ pour tout $f \in S$. Cela est vrai parce que
$A^*$ est l'image réciproque de $(A/I)^*$ par le morphisme canonique
$A \to A/I$.
\end{proof}

\begin{lemma}
\label{algebra-lemma-formally-etale-lift-infinitesimal}
Soit $R \to S$ un morphisme d'anneaux. Soit $J \subset S$ un idéal tel
que $R \to S/J$ soit surjectif ; soit $I \subset R$ le noyau.
Si $R \to S$ est formellement \'etale, alors $R/I^n \to S/J^n$ est un
isomorphisme pour tout $n$, et
$\bigoplus I^n/I^{n + 1} \to \bigoplus J^n/J^{n + 1}$
est un isomorphisme d'anneaux gradués.
\end{lemma}

\begin{proof}
En utilisant la propriété de relèvement par récurrence, nous obtenons des flèches pointillées
$$
\xymatrix{
S \ar[r] \ar@{-->}[rd] & S/J = R/I \\
R \ar[r] \ar[u] & R/I^2 \ar[u]
}
\quad
\xymatrix{
S \ar[r] \ar@{-->}[rd] & R/I^2 \\
R \ar[r] \ar[u] & R/I^3 \ar[u]
}
\quad
\xymatrix{
S \ar[r] \ar@{-->}[rd] & R/I^3 \\
R \ar[r] \ar[u] & R/I^4 \ar[u]
}
$$
Les morphismes correspondants $S/J^n \to R/I^n$ sont des isomorphismes
puisque les composées $S/J^n \to R/I^n \to S/J^n$ sont
(par récurrence) l'identité, par l'unicité dans la propriété de
relèvement des morphismes d'anneaux formellement \'etales.
\end{proof}

\begin{lemma}
\label{algebra-lemma-formally-etale-omega}
Soient $R \to S \to S'$ des morphismes d'anneaux. Soient $J$, resp.\ $J'$,
les noyaux des morphismes de multiplication
$S \otimes_R S \to S$, resp.\ $S' \otimes_{R'} S' \to S'$.
Si $S \to S'$ est formellement \'etale, alors le morphisme
$$
S' \otimes_S \left((S \otimes_R S)/J^{k + 1}\right)
\longrightarrow
(S' \otimes_R S')/(J')^{k + 1}
$$
est un isomorphisme pour tout $k \geq 0$. En particulier, le morphisme
$S' \otimes_S \Omega_{S/R} \to \Omega_{S'/R}$ est un isomorphisme.
\end{lemma}

\begin{proof}
Observons que $S' \otimes_S (S \otimes_R S) = S' \otimes_R S$
et que l'idéal $J$ engendre dans $S' \otimes_R S$ le noyau $I$ de la
surjection $S' \otimes_R S \to S'$. Par conséquent, le
membre de gauche est égal à $S' \otimes_R S/I^{k + 1}$.
Le morphisme $S' \otimes_R S \to S' \otimes_R S'$ est formellement \'etale
par le lemme \ref{algebra-lemma-base-change-formally-etale}. Nous concluons donc
que la flèche affichée dans l'énoncé du lemme est un isomorphisme
par le lemme \ref{algebra-lemma-formally-etale-lift-infinitesimal}.
La dernière assertion résulte de ceci et du fait que
$\Omega_{S/R} = J/J^2$ et $\Omega_{S'/R} = J'/(J')^2$ d'après le
lemme \ref{algebra-lemma-differentials-diagonal}.
\end{proof}

\begin{lemma}
\label{algebra-lemma-formally-etale-principal-parts}
Soient $R \to S \to S'$ des morphismes d'anneaux, avec $S \to S'$ formellement \'etale
(par exemple \'etale).
Soit $M$ un $S$-module. Posons $M' = S' \otimes_R M$. Alors nous avons
$$
S' \otimes_S P^k_{S/R}(M) = P^k_{S'/R}(M')
$$
Il s'ensuit que, pour tout $S$-module $N$ et tout
opérateur différentiel d'ordre fini $D : M \to N$,
il existe un unique prolongement $D' : M' \to S' \otimes_S N$
de $D$ en un opérateur différentiel (du même ordre ou d'ordre inférieur).
\end{lemma}

\begin{proof}
Soient $J$ et $J'$ comme dans l'énoncé du
lemme \ref{algebra-lemma-formally-etale-omega}. Alors nous avons
\begin{align*}
S' \otimes_S P^k_{S/R}(M)
& =
S' \otimes_S \left((S \otimes_R M)/J^{k + 1}(S \otimes_R M)\right) \\
& =
S' \otimes_S \left((S \otimes_R S)/J^{k + 1}\right) \otimes_S M \\
& =
\left(S' \otimes_R S'/(J')^{k + 1}\right) \otimes_S M \\
& =
\left(S' \otimes_R S'/(J')^{k + 1}\right) \otimes_{S'} M' \\
& =
S' \otimes_R M'/(J')^{k + 1}(S' \otimes_R M') \\
& =
P^k_{S'/R'}(M')
\end{align*}
La première et la dernière égalité proviennent du
lemme \ref{algebra-lemma-principal-parts-diagonal}.
La troisième égalité est le lemme \ref{algebra-lemma-formally-etale-omega}.
La dernière assertion est vraie parce que, si $D$ correspond à l'application
linéaire $\gamma : P^k_{S/R}(M) \to N$, alors nous pouvons faire
correspondre à $D' : M' \to S' \otimes_R N$ l'application linéaire
$1 \otimes \gamma : S' \otimes_R P^k_{S/R}(M) \to S' \otimes_R N$.
\end{proof}

\begin{remark}
\label{algebra-remark-formally-etale-differential-operators}
Soient $R \to S \to S'$ des morphismes d'anneaux, avec $S \to S'$ formellement \'etale
(par exemple \'etale). Soient $M_i$, $i = 1, 2, 3$, des $S$-modules,
et soient $D_i : M_i \to M_{i + 1}$, $i = 1, 2$, des opérateurs
différentiels d'ordre fini. Alors, si $D'_i : M'_i \to M'_{i + 1}$, $i = 1, 2$,
sont les prolongements des $D_i$ à $M'_i = S' \otimes_S M_i$ comme dans le
lemme \ref{algebra-lemma-formally-etale-principal-parts}, alors $D'_2 \circ D'_1$
est le prolongement de $D_2 \circ D_1$. En particulier, si $M$ est un $S$-module,
alors $M' = S' \otimes_S M$ est un module sur la $S$-algèbre
$\text{Diff}_{S/R}(M, M)$.
\end{remark}





\section{Morphismes d'anneaux non ramifiés}
\label{algebra-section-unramified}

\noindent
Notre définition d'un morphisme d'anneaux non ramifié est celle de \cite{Henselian}.
Ce que nous appelons morphisme d'anneaux G-non ramifié est ce qu'EGA appelle morphisme non ramifié.

\begin{definition}
\label{algebra-definition-unramified}
Soit $R \to S$ un morphisme d'anneaux.
\begin{enumerate}
\item On dit que $R \to S$ est {\it non ramifié} si $R \to S$ est de
type fini et $\Omega_{S/R} = 0$.
\item On dit que $R \to S$ est {\it G-non ramifié} si $R \to S$ est de
présentation finie et $\Omega_{S/R} = 0$.
\item Étant donné un idéal premier $\mathfrak q$ de $S$, on dit que $S$ est
{\it non ramifié en $\mathfrak q$} s'il existe un
$g \in S$, $g \not \in \mathfrak q$, tel que $R \to S_g$ soit non ramifié.
\item Étant donné un idéal premier $\mathfrak q$ de $S$, on dit que $S$ est
{\it G-non ramifié en $\mathfrak q$} s'il existe un
$g \in S$, $g \not \in \mathfrak q$, tel que $R \to S_g$ soit G-non ramifié.
\end{enumerate}
\end{definition}

\noindent
Bien sûr, un morphisme G-non ramifié est non ramifié.

\begin{lemma}
\label{algebra-lemma-formally-unramified-unramified}
Soit $R \to S$ un morphisme d'anneaux. Les assertions suivantes sont équivalentes :
\begin{enumerate}
\item $R \to S$ est formellement non ramifié et de type fini, et
\item $R \to S$ est non ramifié.
\end{enumerate}
De plus, les assertions suivantes sont également équivalentes :
\begin{enumerate}
\item $R \to S$ est formellement non ramifié et de présentation finie, et
\item $R \to S$ est G-non ramifié.
\end{enumerate}
\end{lemma}

\begin{proof}
Cela résulte du lemme \ref{algebra-lemma-characterize-formally-unramified}
et des définitions.
\end{proof}

\begin{lemma}
\label{algebra-lemma-unramified}
Propriétés des morphismes d'anneaux non ramifiés et G-non ramifiés.
\begin{enumerate}
\item Le changement de base d'un morphisme d'anneaux non ramifié est non ramifié.
Le changement de base d'un morphisme d'anneaux G-non ramifié est G-non ramifié.
\item La composée de morphismes d'anneaux non ramifiés est non ramifiée.
La composée de morphismes d'anneaux G-non ramifiés est G-non ramifiée.
\item Toute localisation principale $R \to R_f$ est G-non ramifiée et
non ramifiée.
\item Si $I \subset R$ est un idéal, alors $R \to R/I$ est non ramifié.
Si $I \subset R$ est un idéal de type fini, alors $R \to R/I$ est
G-non ramifié.
\item Un morphisme d'anneaux \'etale est G-non ramifié et non ramifié.
\item Si $R \to S$ est de type fini (resp.\ de présentation finie),
si $\mathfrak q \subset S$ est un idéal premier et si $(\Omega_{S/R})_{\mathfrak q} = 0$,
alors $R \to S$ est non ramifié (resp.\ G-non ramifié) en $\mathfrak q$.
\item Si $R \to S$ est de type fini (resp.\ de présentation finie),
si $\mathfrak q \subset S$ est un idéal premier et si
$\Omega_{S/R} \otimes_S \kappa(\mathfrak q) = 0$, alors
$R \to S$ est non ramifié (resp.\ G-non ramifié) en $\mathfrak q$.
\item Si $R \to S$ est de type fini (resp.\ de présentation finie),
si $\mathfrak q \subset S$ est un idéal premier au-dessus de $\mathfrak p \subset R$ et si
$(\Omega_{S \otimes_R \kappa(\mathfrak p)/\kappa(\mathfrak p)})_{\mathfrak q}
= 0$, alors $R \to S$ est non ramifié (resp.\ G-non ramifié) en $\mathfrak q$.
\item Si $R \to S$ est de type fini (resp.\ de présentation finie),
si $\mathfrak q \subset S$ est un idéal premier au-dessus de $\mathfrak p \subset R$ et si
$(\Omega_{S \otimes_R \kappa(\mathfrak p)/\kappa(\mathfrak p)})
\otimes_{S \otimes_R \kappa(\mathfrak p)} \kappa(\mathfrak q) = 0$,
alors $R \to S$ est non ramifié (resp.\ G-non ramifié) en $\mathfrak q$.
\item Si $R \to S$ est un morphisme d'anneaux, si $g_1, \ldots, g_m \in S$ engendrent
l'idéal unité et si $R \to S_{g_j}$ est non ramifié (resp.\ G-non ramifié) pour
$j = 1, \ldots, m$, alors $R \to S$ est non ramifié (resp.\ G-non ramifié).
\item Si $R \to S$ est un morphisme d'anneaux qui est non ramifié (resp.\ G-non ramifié)
en tout idéal premier de $S$, alors $R \to S$ est non ramifié (resp.\ G-non ramifié).
\item Si $R \to S$ est G-non ramifié, alors il existe une
$\mathbf{Z}$-algèbre $R_0$ de type fini, un morphisme d'anneaux G-non ramifié
$R_0 \to S_0$ et un morphisme d'anneaux $R_0 \to R$ tels que
$S = R \otimes_{R_0} S_0$.
\item Si $R \to S$ est non ramifié, alors il existe une
$\mathbf{Z}$-algèbre $R_0$ de type fini, un morphisme d'anneaux non ramifié
$R_0 \to S_0$ et un morphisme d'anneaux $R_0 \to R$ tels que $S$ soit un quotient de
$R \otimes_{R_0} S_0$.
\end{enumerate}
\end{lemma}

\begin{proof}
Nous démontrons chaque point dans l'ordre.

\medskip\noindent
Pour (1). Cela résulte des lemmes \ref{algebra-lemma-differentials-base-change}
et \ref{algebra-lemma-base-change-finiteness}.

\medskip\noindent
Pour (2). Cela résulte des lemmes \ref{algebra-lemma-exact-sequence-differentials}
et \ref{algebra-lemma-base-change-finiteness}.

\medskip\noindent
Pour (3). Cela résulte d'un calcul direct de $\Omega_{R_f/R}$, que nous omettons.

\medskip\noindent
Pour (4). Nous avons $\Omega_{(R/I)/R} = 0$, voir le
lemme \ref{algebra-lemma-trivial-differential-surjective},
et le morphisme d'anneaux $R \to R/I$
est de type fini. Si $I$ est un idéal de type fini, alors $R \to R/I$
est de présentation finie.

\medskip\noindent
Pour (5). Voir la discussion suivant la définition \ref{algebra-definition-etale}.

\medskip\noindent
Pour (6). Dans ce cas, $\Omega_{S/R}$ est un $S$-module de type fini (voir le
lemme \ref{algebra-lemma-differentials-finitely-generated}) et il existe donc
un $g \in S$, $g \not \in \mathfrak q$, tel que
$(\Omega_{S/R})_g = 0$. D'après le lemme \ref{algebra-lemma-differentials-localize},
cela signifie que $\Omega_{S_g/R} = 0$ et donc que $R \to S_g$ est
non ramifié, comme désiré.

\medskip\noindent
Pour (7). Le lemme de Nakayama (lemme \ref{algebra-lemma-NAK}) montre que
la condition est équivalente à celle de (6).

\medskip\noindent
Pour (8) et (9). Ces assertions sont équivalentes de la même manière que (6) et (7)
le sont. De plus,
$\Omega_{S \otimes_R \kappa(\mathfrak p)/\kappa(\mathfrak p)} =
\Omega_{S/R} \otimes_S (S \otimes_R \kappa(\mathfrak p))$ d'après le
lemme \ref{algebra-lemma-differentials-base-change}.
Nous voyons donc que (9) est équivalente à (7), puisque les
$\kappa(\mathfrak q)$-espaces vectoriels qui y figurent sont canoniquement
isomorphes.

\medskip\noindent
Pour (10). Cela résulte des lemmes \ref{algebra-lemma-cover}
et \ref{algebra-lemma-differentials-localize}.

\medskip\noindent
Pour (11). Cela résulte de (6) et (7) et du fait que le spectre de $S$
est quasi-compact.

\medskip\noindent
Pour (12). Écrivons $S = R[x_1, \ldots, x_n]/(g_1, \ldots, g_m)$.
Comme $\Omega_{S/R} = 0$, nous pouvons écrire
$$
\text{d}x_i = \sum h_{ij}\text{d}g_j + \sum a_{ijk}g_j\text{d}x_k
$$
dans $\Omega_{R[x_1, \ldots, x_n]/R}$
pour certains $h_{ij}, a_{ijk} \in R[x_1, \ldots, x_n]$.
Choisissons une
$\mathbf{Z}$-sous-algèbre $R_0 \subset R$ de type fini contenant tous les coefficients des
polynômes $g_i, h_{ij}, a_{ijk}$. Posons
$S_0 = R_0[x_1, \ldots, x_n]/(g_1, \ldots, g_m)$. Cela convient.

\medskip\noindent
Pour (13). Écrivons $S = R[x_1, \ldots, x_n]/I$.
Comme $\Omega_{S/R} = 0$, nous pouvons écrire
$$
\text{d}x_i = \sum h_{ij}\text{d}g_{ij} + \sum g'_{ik}\text{d}x_k
$$
dans $\Omega_{R[x_1, \ldots, x_n]/R}$
pour certains $h_{ij} \in R[x_1, \ldots, x_n]$ et $g_{ij}, g'_{ik} \in I$.
Choisissons une $\mathbf{Z}$-sous-algèbre $R_0 \subset R$ de type fini
contenant tous les coefficients des
polynômes $g_{ij}, h_{ij}, g'_{ik}$. Posons
$S_0 = R_0[x_1, \ldots, x_n]/(g_{ij}, g'_{ik})$. Cela convient.
\end{proof}

\begin{lemma}
\label{algebra-lemma-diagonal-unramified}
Soit $R \to S$ un morphisme d'anneaux.
Si $R \to S$ est non ramifié, alors il existe un idempotent
$e \in S \otimes_R S$ tel que $S \otimes_R S \to S$ soit isomorphe à
$S \otimes_R S \to (S \otimes_R S)_e$.
\end{lemma}

\begin{proof}
Soit $J = \Ker(S \otimes_R S \to S)$. Par hypothèse,
$J/J^2 = 0$, voir le
lemme \ref{algebra-lemma-differentials-diagonal}.
Comme $S$ est de type fini sur $R$, nous
voyons que $J$ est de type fini, à savoir engendré par les
$x_i \otimes 1 - 1 \otimes x_i$, où les $x_i$ engendrent $S$ sur $R$.
On conclut par le lemme \ref{algebra-lemma-ideal-is-squared-union-connected}.
\end{proof}

\begin{lemma}
\label{algebra-lemma-unramified-at-prime}
Soit $R \to S$ un morphisme d'anneaux.
Soit $\mathfrak q \subset S$ un
idéal premier au-dessus de $\mathfrak p$ dans $R$.
Si $S/R$ est non ramifié en $\mathfrak q$, alors
\begin{enumerate}
\item nous avons $\mathfrak p S_{\mathfrak q} = \mathfrak qS_{\mathfrak q}$,
et cet idéal est l'idéal maximal de l'anneau local $S_{\mathfrak q}$, et
\item l'extension de corps $\kappa(\mathfrak q)/\kappa(\mathfrak p)$
est finie séparable.
\end{enumerate}
\end{lemma}

\begin{proof}
Nous pouvons d'abord remplacer $S$ par $S_g$ pour un $g \in S$, $g \not \in \mathfrak q$,
et supposer que $R \to S$ est non ramifié.
Le changement de base $S \otimes_R \kappa(\mathfrak p)$
est non ramifié sur $\kappa(\mathfrak p)$ d'après le
lemme \ref{algebra-lemma-unramified}.
D'après le
lemme \ref{algebra-lemma-characterize-smooth-over-field},
il est lisse, donc \'etale, sur $\kappa(\mathfrak p)$.
Ainsi, $F = S \otimes_R \kappa(\mathfrak p)$ est un produit fini
d'extensions de corps finies séparables de $\kappa(\mathfrak p)$, voir le
lemme \ref{algebra-lemma-etale-over-field}. Avec les notations et résultats de la
remarque \ref{algebra-remark-local-ring-fibre}, nous trouvons que
$F_{\overline{\mathfrak q}} = S_\mathfrak q/\mathfrak pS_\mathfrak q$
est égal à $\kappa(\mathfrak q)$. Cela implique le lemme.
\end{proof}

\begin{lemma}
\label{algebra-lemma-unramified-quasi-finite}
Soit $R \to S$ un morphisme d'anneaux de type fini.
Soit $\mathfrak q$ un idéal premier de $S$.
Si $R \to S$ est non ramifié en $\mathfrak q$, alors
$R \to S$ est quasi-fini en $\mathfrak q$.
En particulier, un morphisme d'anneaux non ramifié est quasi-fini.
\end{lemma}

\begin{proof}
Un morphisme d'anneaux non ramifié est de type fini.
Il est donc clair que la seconde assertion résulte de la première.
Pour voir la première assertion, appliquons la caractérisation du
lemme \ref{algebra-lemma-isolated-point-fibre}, partie (2), en utilisant le
lemme \ref{algebra-lemma-unramified-at-prime}.
\end{proof}

\begin{lemma}
\label{algebra-lemma-characterize-unramified}
Soit $R \to S$ un morphisme d'anneaux. Soit $\mathfrak q$ un idéal premier de $S$
au-dessus d'un idéal premier $\mathfrak p$ de $R$. Si
\begin{enumerate}
\item $R \to S$ est de type fini,
\item $\mathfrak p S_{\mathfrak q}$ est l'idéal maximal
de l'anneau local $S_{\mathfrak q}$, et
\item l'extension de corps $\kappa(\mathfrak q)/\kappa(\mathfrak p)$
est finie séparable,
\end{enumerate}
alors $R \to S$ est non ramifié en $\mathfrak q$.
\end{lemma}

\begin{proof}
D'après le lemme \ref{algebra-lemma-unramified} (8), il suffit de montrer que
$\Omega_{S \otimes_R \kappa(\mathfrak p) / \kappa(\mathfrak p)}$
est nul après localisation en $\mathfrak q$. Nous pouvons donc remplacer $S$
par $S \otimes_R \kappa(\mathfrak p)$ et $R$ par $\kappa(\mathfrak p)$.
Autrement dit, nous pouvons supposer que $R = k$ est un corps et que $S$
est une $k$-algèbre de type fini.
Dans ce cas, les hypothèses impliquent que
$S_{\mathfrak q} \cong \kappa(\mathfrak q)$.
Ainsi, $(\Omega_{S/k})_{\mathfrak q} = \Omega_{S_\mathfrak q/k} =
\Omega_{\kappa(\mathfrak q)/k}$ est nul, comme désiré (la
première égalité est le lemme \ref{algebra-lemma-differentials-localize}).
\end{proof}

\begin{lemma}
\label{algebra-lemma-etale-flat-unramified-finite-presentation}
Soit $R \to S$ un morphisme d'anneaux. Les assertions suivantes sont équivalentes :
\begin{enumerate}
\item $R \to S$ est \'etale,
\item $R \to S$ est plat et G-non ramifié, et
\item $R \to S$ est plat, non ramifié et de présentation finie.
\end{enumerate}
\end{lemma}

\begin{proof}
Les assertions (2) et (3) sont équivalentes par définition.
L'implication (1) $\Rightarrow$ (3) résulte du
fait que les morphismes d'anneaux \'etales sont de présentation finie,
du lemme \ref{algebra-lemma-etale} (platitude des morphismes \'etales) et du
lemme \ref{algebra-lemma-unramified} (les morphismes \'etales sont non ramifiés).
Réciproquement, la caractérisation des morphismes d'anneaux \'etales du
lemme \ref{algebra-lemma-characterize-etale}
et la structure des morphismes d'anneaux non ramifiés du
lemme \ref{algebra-lemma-unramified-at-prime}
montrent que (3) implique (1). (On utilise ici que $R \to S$
est \'etale si $R \to S$ est \'etale en tout idéal premier $\mathfrak q \subset S$,
voir le lemme \ref{algebra-lemma-etale}.)
\end{proof}

\begin{lemma}
\label{algebra-lemma-characterize-etale-over-polynomial-ring}
Soit $k$ un corps. Soit
$$
\varphi : k[x_1, \ldots, x_n] \to A, \quad x_i \longmapsto a_i
$$
un morphisme d'anneaux de type fini. Alors $\varphi$ est \'etale si et seulement si les
deux conditions suivantes sont vérifiées : (a) les anneaux locaux de $A$ aux idéaux maximaux
sont de dimension $n$, et (b) les éléments $\text{d}(a_1), \ldots, \text{d}(a_n)$
engendrent $\Omega_{A/k}$ comme $A$-module.
\end{lemma}

\begin{proof}
Supposons (a) et (b). La condition (b) implique que
$\Omega_{A/k[x_1, \ldots, x_n]} = 0$ et donc que $\varphi$ est non ramifié.
Il suffit ainsi de prouver que $\varphi$ est plat, voir le
lemme \ref{algebra-lemma-etale-flat-unramified-finite-presentation}.
Soit $\mathfrak m \subset A$ un idéal maximal.
Posons $X = \Spec(A)$ et notons $x \in X$ le point fermé correspondant
à $\mathfrak m$. Alors $\dim(A_\mathfrak m)$ est $\dim_x X$, voir le
lemme \ref{algebra-lemma-dimension-closed-point-finite-type-field}.
Ainsi, d'après le lemme \ref{algebra-lemma-characterize-smooth-over-field},
si (a) et (b) sont vérifiées, alors $A_\mathfrak m$ est
un anneau local régulier pour tout idéal maximal $\mathfrak m$. Ensuite,
$k[x_1, \ldots, x_n]_{\varphi^{-1}(\mathfrak m)} \to A_\mathfrak m$
est plat d'après le lemme \ref{algebra-lemma-CM-over-regular-flat}
(et le fait qu'un anneau local régulier est de Cohen-Macaulay, voir le
lemme \ref{algebra-lemma-regular-ring-CM}).
Ainsi, $\varphi$ est plat d'après le lemme \ref{algebra-lemma-flat-localization}.

\medskip\noindent
Supposons que $\varphi$ soit \'etale. Alors $\Omega_{A/k[x_1, \ldots, x_n]} = 0$
et donc (b) est vérifiée. D'autre part, les morphismes d'anneaux \'etales sont plats
(lemme \ref{algebra-lemma-etale}) et quasi-finis
(lemme \ref{algebra-lemma-etale-quasi-finite}).
Ainsi, pour tout idéal maximal $\mathfrak m$ de $A$, nous pouvons appliquer le
lemme \ref{algebra-lemma-dimension-base-fibre-equals-total} à
$k[x_1, \ldots, x_n]_{\varphi^{-1}(\mathfrak m)} \to A_\mathfrak m$
pour voir que $\dim(A_\mathfrak m) = n$ et donc que (a) est vérifiée.
\end{proof}




\section{Structure locale des morphismes d'anneaux non ramifiés}
\label{algebra-section-local-structure-unramified}

\noindent
Un morphisme non ramifié est localement (en un sens convenable) la composée
d'une immersion fermée et d'un morphisme \'etale. Les fondements algébriques
de ce fait sont examinés dans cette section.

\begin{proposition}
\label{algebra-proposition-unramified-locally-standard}
Soit $R \to S$ un morphisme d'anneaux. Soit $\mathfrak q \subset S$ un idéal premier.
Si $R \to S$ est non ramifié en $\mathfrak q$, alors il existe
\begin{enumerate}
\item un $g \in S$, $g \not \in \mathfrak q$,
\item un morphisme d'anneaux \'etale standard $R \to S'$, et
\item un morphisme surjectif de $R$-algèbres $S' \to S_g$.
\end{enumerate}
\end{proposition}

\begin{proof}
Cette démonstration est la « même » que celle de la
proposition \ref{algebra-proposition-etale-locally-standard}.
Elle est quelque peu indirecte et il existe peut-être des moyens de
la raccourcir.

\medskip\noindent
Étape 1. D'après la définition \ref{algebra-definition-unramified},
il existe un $g \in S$, $g \not \in \mathfrak q$,
tel que $R \to S_g$ soit non ramifié. Nous pouvons donc supposer que $S$ est
non ramifié sur $R$.

\medskip\noindent
Étape 2. D'après le lemme \ref{algebra-lemma-unramified},
il existe un morphisme d'anneaux non ramifié $R_0 \to S_0$,
avec $R_0$ de type fini sur $\mathbf{Z}$, et un morphisme d'anneaux
$R_0 \to R$ tel que $S$ soit un quotient de $R \otimes_{R_0} S_0$. Notons
$\mathfrak q_0$ l'idéal premier de $S_0$ correspondant à $\mathfrak q$.
Si nous montrons le résultat pour $(R_0 \to S_0, \mathfrak q_0)$, alors le
résultat pour $(R \to S, \mathfrak q)$ s'en déduit par changement de base. Ainsi,
nous pouvons supposer que $R$ est noethérien.

\medskip\noindent
Étape 3.
Notons que $R \to S$ est quasi-fini d'après le
lemme \ref{algebra-lemma-unramified-quasi-finite}.
D'après le lemme \ref{algebra-lemma-quasi-finite-open-integral-closure},
il existe un morphisme d'anneaux fini $R \to S'$, un morphisme de $R$-algèbres
$S' \to S$ et un élément $g' \in S'$ satisfaisant
$g' \not \in \mathfrak q$ et tel que $S' \to S$ induise
un isomorphisme $S'_{g'} \cong S_{g'}$.
(Notons que $S'$ n'est pas nécessairement non ramifié sur $R$.)
Nous pouvons donc supposer que (a) $R$ est noethérien, (b) $R \to S$ est fini
et (c) $R \to S$ est non ramifié en $\mathfrak q$
(mais plus nécessairement non ramifié en tous les idéaux premiers).

\medskip\noindent
Étape 4. Soit $\mathfrak p \subset R$ l'idéal premier correspondant
à $\mathfrak q$. Considérons l'anneau fibre
$S \otimes_R \kappa(\mathfrak p)$. C'est une algèbre finie sur
$\kappa(\mathfrak p)$. Il est donc artinien
(voir le lemme \ref{algebra-lemma-finite-dimensional-algebra}) et
est ainsi un produit fini d'anneaux locaux
$$
S \otimes_R \kappa(\mathfrak p) = \prod\nolimits_{i = 1}^n A_i
$$
voir la proposition \ref{algebra-proposition-dimension-zero-ring}. L'un des facteurs,
disons $A_1$, est l'anneau local $S_{\mathfrak q}/\mathfrak pS_{\mathfrak q}$,
qui est isomorphe à $\kappa(\mathfrak q)$,
voir le lemme \ref{algebra-lemma-unramified-at-prime}. Les autres facteurs correspondent aux
autres idéaux premiers, disons $\mathfrak q_2, \ldots, \mathfrak q_n$, de
$S$ au-dessus de $\mathfrak p$.

\medskip\noindent
Étape 5. Nous pouvons choisir un élément non nul $\alpha \in \kappa(\mathfrak q)$ qui
engendre l'extension de corps finie séparable
$\kappa(\mathfrak q)/\kappa(\mathfrak p)$ (ainsi, même si
l'extension de corps est triviale, nous n'autorisons pas $\alpha = 0$).
Notons que, pour tout $\lambda \in \kappa(\mathfrak p)^*$, l'élément
$\lambda \alpha$ engendre lui aussi $\kappa(\mathfrak q)$
sur $\kappa(\mathfrak p)$. Considérons l'élément
$$
\overline{t} =
(\alpha, 0, \ldots, 0) \in
\prod\nolimits_{i = 1}^n A_i =
S \otimes_R \kappa(\mathfrak p).
$$
Après avoir éventuellement remplacé $\alpha$ par $\lambda \alpha$ comme ci-dessus,
nous pouvons supposer que $\overline{t}$ est l'image d'un $t \in S$.
Soit $I \subset R[x]$ le noyau du morphisme de $R$-algèbres
$R[x] \to S$ qui envoie $x$ sur $t$. Posons $S' = R[x]/I$,
de sorte que $S' \subset S$. Voici un diagramme :
$$
\xymatrix{
R[x] \ar[r] & S' \ar[r] & S \\
R \ar[u] \ar[ru] \ar[rru] & &
}
$$
Par construction, les idéaux premiers $\mathfrak q_j$, $j \geq 2$, de $S$ se
trouvent tous au-dessus de l'idéal premier $(\mathfrak p, x)$ de $R[x]$, tandis que
l'idéal premier $\mathfrak q$ se trouve au-dessus d'un autre idéal premier de $R[x]$,
car $\alpha \not = 0$.

\medskip\noindent
Étape 6. Notons $\mathfrak q' \subset S'$ l'idéal premier de $S'$
correspondant à $\mathfrak q$. D'après ce qui précède, $\mathfrak q$ est
le seul idéal premier de $S$ au-dessus de $\mathfrak q'$. Nous voyons donc que
$S_{\mathfrak q} = S_{\mathfrak q'}$, voir le
lemme \ref{algebra-lemma-unique-prime-over-localize-below} (nous disposons de la
montée pour $S' \to S$ d'après le lemme \ref{algebra-lemma-integral-going-up},
puisque $S' \to S$ est fini car $R \to S$ est fini).
Il s'ensuit que $S'_{\mathfrak q'} \to S_{\mathfrak q}$ est fini
et injectif en tant que localisation du morphisme d'anneaux fini injectif
$S' \to S$. Considérons les morphismes d'anneaux locaux
$$
R_{\mathfrak p} \to S'_{\mathfrak q'} \to S_{\mathfrak q}
$$
Le second morphisme est fini et injectif. Nous avons
$S_{\mathfrak q}/\mathfrak pS_{\mathfrak q} = \kappa(\mathfrak q)$,
voir le lemme \ref{algebra-lemma-unramified-at-prime}.
Donc, a fortiori,
$S_{\mathfrak q}/\mathfrak q'S_{\mathfrak q} = \kappa(\mathfrak q)$.
Puisque
$$
\kappa(\mathfrak p) \subset \kappa(\mathfrak q') \subset \kappa(\mathfrak q)
$$
et puisque $\alpha$ appartient à l'image de $\kappa(\mathfrak q')$ dans
$\kappa(\mathfrak q)$,
nous concluons que $\kappa(\mathfrak q') = \kappa(\mathfrak q)$.
Ainsi, d'après le lemme de Nakayama \ref{algebra-lemma-NAK} appliqué au morphisme de
$S'_{\mathfrak q'}$-modules $S'_{\mathfrak q'} \to S_{\mathfrak q}$,
le morphisme $S'_{\mathfrak q'} \to S_{\mathfrak q}$ est surjectif.
Autrement dit,
$S'_{\mathfrak q'} \cong S_{\mathfrak q}$.

\medskip\noindent
Étape 7. D'après le lemme \ref{algebra-lemma-isomorphic-local-rings}, il existe
$g \in S$, $g \not \in \mathfrak q$, et
$g' \in S'$, $g' \not \in \mathfrak q'$, tels que $S'_{g'} \cong S_g$.
Comme $R$ est noethérien, l'anneau $S'$ est fini sur $R$,
car c'est un $R$-sous-module
du $R$-module fini $S$. Ainsi, après avoir remplacé $S$ par $S'$, nous pouvons
supposer que (a) $R$ est noethérien, (b) $S$ est fini sur $R$, (c)
$S$ est non ramifié sur $R$ en $\mathfrak q$, et (d) $S = R[x]/I$.

\medskip\noindent
Étape 8. Considérons l'anneau
$S \otimes_R \kappa(\mathfrak p) = \kappa(\mathfrak p)[x]/\overline{I}$,
où $\overline{I} = I \cdot \kappa(\mathfrak p)[x]$ est l'idéal engendré
par $I$ dans $\kappa(\mathfrak p)[x]$. Comme $\kappa(\mathfrak p)[x]$ est un anneau principal,
nous savons que $\overline{I} = (\overline{h})$ pour un certain polynôme unitaire
$\overline{h} \in \kappa(\mathfrak p)$. Après avoir remplacé $\overline{h}$
par $\lambda \cdot \overline{h}$ pour un certain $\lambda \in \kappa(\mathfrak p)$,
nous pouvons supposer que $\overline{h}$ est l'image d'un $h \in R[x]$.
(Le problème est que nous ne savons pas si nous pouvons choisir $h$ unitaire.)
De plus, comme à l'étape 4, nous savons que
$S \otimes_R \kappa(\mathfrak p) = A_1 \times \ldots \times A_n$, où
$A_1 = \kappa(\mathfrak q)$ est une extension finie séparable de
$\kappa(\mathfrak p)$ et $A_2, \ldots, A_n$ sont locaux. Cela implique
que
$$
\overline{h} = \overline{h}_1 \overline{h}_2^{e_2} \ldots \overline{h}_n^{e_n}
$$
pour certains polynômes unitaires irréductibles deux à deux premiers entre eux
$\overline{h}_i \in \kappa(\mathfrak p)[x]$ et certains entiers
$e_2, \ldots, e_n \geq 1$. La numérotation est choisie de sorte que
$A_i = \kappa(\mathfrak p)[x]/(\overline{h}_i^{e_i})$ comme
$\kappa(\mathfrak p)[x]$-algèbres. Notons que $\overline{h}_1$ est
le polynôme minimal de $\alpha \in \kappa(\mathfrak q)$ et est donc
un polynôme séparable (sa dérivée est première avec lui-même).

\medskip\noindent
Étape 9. Soit $m \in I$ un élément unitaire ; un tel élément existe
parce que l'extension d'anneaux $R \to R[x]/I$ est finie, donc entière.
Notons $\overline{m}$ son image dans $\kappa(\mathfrak p)[x]$.
Nous pouvons factoriser
$$
\overline{m} = \overline{k}
\overline{h}_1^{d_1} \overline{h}_2^{d_2} \ldots \overline{h}_n^{d_n}
$$
pour un certain $d_1 \geq 1$, des $d_j \geq e_j$, $j = 2, \ldots, n$, et
un $\overline{k} \in \kappa(\mathfrak p)[x]$ premier avec tous les $\overline{h}_i$.
Posons $f = m^l + h$, où $l \deg(m) > \deg(h)$ et $l \geq 2$.
Alors $f$ est unitaire comme polynôme sur $R$. De plus, l'image $\overline{f}$
de $f$ dans $\kappa(\mathfrak p)[x]$ se factorise comme
$$
\overline{f} =
\overline{h}_1 \overline{h}_2^{e_2} \ldots \overline{h}_n^{e_n}
+
\overline{k}^l \overline{h}_1^{ld_1} \overline{h}_2^{ld_2}
\ldots \overline{h}_n^{ld_n}
=
\overline{h}_1(\overline{h}_2^{e_2} \ldots \overline{h}_n^{e_n}
+
\overline{k}^l
\overline{h}_1^{ld_1 - 1} \overline{h}_2^{ld_2} \ldots \overline{h}_n^{ld_n})
= \overline{h}_1 \overline{w}
$$
avec $\overline{w}$ un polynôme premier avec $\overline{h}_1$.
Posons $g = f'$ (la dérivée par rapport à $x$).

\medskip\noindent
Étape 10. Le morphisme d'anneaux $R[x] \to S = R[x]/I$ a les propriétés suivantes :
(1) il envoie $f$ sur zéro, et
(2) il envoie $g$ sur un élément de $S \setminus \mathfrak q$.
La première assertion est claire puisque $f$ appartient à $I$.
Pour la seconde assertion, il nous suffit de montrer que $g$ ne
s'envoie pas sur zéro dans
$\kappa(\mathfrak q) = \kappa(\mathfrak p)[x]/(\overline{h}_1)$.
L'image de $g$ dans $\kappa(\mathfrak p)[x]$ est la dérivée
de $\overline{f}$. Ainsi, (2) est claire parce que
$$
\overline{g} =
\frac{\text{d}\overline{f}}{\text{d}x} =
\overline{w}\frac{\text{d}\overline{h}_1}{\text{d}x} +
\overline{h}_1\frac{\text{d}\overline{w}}{\text{d}x},
$$
$\overline{w}$ est premier avec $\overline{h}_1$ et
$\overline{h}_1$ est séparable.

\medskip\noindent
Étape 11.
Nous concluons que $\varphi : R[x]/(f) \to S$ est un morphisme d'anneaux surjectif,
que $R[x]_g/(f)$ est \'etale sur $R$ (parce qu'il est \'etale standard,
voir le lemme \ref{algebra-lemma-standard-etale}) et que $\varphi(g) \not \in \mathfrak q$.
Ainsi, le morphisme $(R[x]/(f))_g \to S_{\varphi(g)}$ est la
surjection recherchée.
\end{proof}



\begin{lemma}
\label{algebra-lemma-etale-makes-unramified-closed-at-prime}
Soit $R \to S$ un morphisme d'anneaux.
Soit $\mathfrak q$ un idéal premier de $S$ au-dessus de $\mathfrak p \subset R$.
Supposons que $R \to S$ soit de type fini et non ramifié en $\mathfrak q$.
Alors il existe
\begin{enumerate}
\item un morphisme d'anneaux \'etale $R \to R'$,
\item un idéal premier $\mathfrak p' \subset R'$ au-dessus de $\mathfrak p$,
\item une décomposition en produit
$$
R' \otimes_R S = A \times B
$$
\end{enumerate}
ayant les propriétés suivantes :
\begin{enumerate}
\item $R' \to A$ est surjectif, et
\item $\mathfrak p'A$ est un idéal premier de $A$ au-dessus de $\mathfrak p'$ et
de $\mathfrak q$.
\end{enumerate}
\end{lemma}

\begin{proof}
Nous pouvons remplacer $(R \to S, \mathfrak p, \mathfrak q)$
par tout changement de base $(R' \to R'\otimes_R S, \mathfrak p', \mathfrak q')$
par un morphisme d'anneaux \'etale $R \to R'$, avec un idéal premier $\mathfrak p'$
au-dessus de $\mathfrak p$ et un choix de $\mathfrak q'$ au-dessus
à la fois de $\mathfrak q$ et de $\mathfrak p'$. Notons aussi qu'étant donnés
$R \to R'$ et $\mathfrak p'$, on peut toujours trouver un $\mathfrak q'$ convenable.

\medskip\noindent
L'hypothèse que $R \to S$ est de type fini signifie que nous pouvons appliquer le
lemme \ref{algebra-lemma-etale-makes-quasi-finite-finite-variant}. Nous pouvons donc
supposer que $S = A_1 \times \ldots \times A_n \times B$, que
chaque $R \to A_i$ est fini avec exactement un idéal premier $\mathfrak r_i$
au-dessus de $\mathfrak p$ tel que
$\kappa(\mathfrak p) \subset \kappa(\mathfrak r_i)$ soit purement inséparable,
et que $R \to B$ ne soit quasi-fini en aucun idéal premier au-dessus de $\mathfrak p$.
Alors, clairement, $\mathfrak q = \mathfrak r_i$ pour un certain $i$, puisqu'un
morphisme non ramifié est quasi-fini
(voir le lemme \ref{algebra-lemma-unramified-quasi-finite}).
Disons que $\mathfrak q = \mathfrak r_1$.
D'après le lemme \ref{algebra-lemma-unramified-at-prime}, nous voyons que
$\kappa(\mathfrak r_1)/\kappa(\mathfrak p)$
est séparable, donc est l'extension de corps triviale, et que
$\mathfrak p(A_1)_{\mathfrak r_1}$ est l'idéal maximal.
De plus, d'après le lemme \ref{algebra-lemma-unique-prime-over-localize-below}
(qui s'applique à $R \to A_1$ parce qu'un morphisme d'anneaux fini satisfait la montée d'après le
lemme \ref{algebra-lemma-integral-going-up}),
nous avons $(A_1)_{\mathfrak r_1} = (A_1)_{\mathfrak p}$.
Il résulte du lemme de Nakayama \ref{algebra-lemma-NAK}
que le morphisme d'anneaux locaux
$R_{\mathfrak p} \to (A_1)_{\mathfrak p} = (A_1)_{\mathfrak r_1}$
est surjectif. Comme $A_1$ est fini sur $R$, nous voyons qu'il
existe un $f \in R$, $f \not \in \mathfrak p$, tel que
$R_f \to (A_1)_f$ soit surjectif. Après avoir remplacé $R$ par $R_f$, on conclut.
\end{proof}

\begin{lemma}
\label{algebra-lemma-etale-makes-unramified-closed}
\begin{slogan}
Dans un morphisme d'anneaux non ramifié, on peut séparer les points d'une fibre
en passant à un voisinage \'etale.
\end{slogan}
Soit $R \to S$ un morphisme d'anneaux.
Soit $\mathfrak p$ un idéal premier de $R$.
Si $R \to S$ est non ramifié, alors il existe
\begin{enumerate}
\item un morphisme d'anneaux \'etale $R \to R'$,
\item un idéal premier $\mathfrak p' \subset R'$ au-dessus de $\mathfrak p$,
\item une décomposition en produit
$$
R' \otimes_R S = A_1 \times \ldots \times A_n \times B
$$
\end{enumerate}
ayant les propriétés suivantes :
\begin{enumerate}
\item $R' \to A_i$ est surjectif,
\item $\mathfrak p'A_i$ est un idéal premier de $A_i$ au-dessus de $\mathfrak p'$, et
\item il n'existe aucun idéal premier de $B$ au-dessus de $\mathfrak p'$.
\end{enumerate}
\end{lemma}

\begin{proof}
Nous pouvons appliquer le lemme \ref{algebra-lemma-etale-makes-quasi-finite-finite-variant}.
Ainsi, après un changement de base \'etale,
nous pouvons supposer que $S = A_1 \times \ldots \times A_n \times B$,
que chaque $R \to A_i$ est fini avec exactement un idéal premier $\mathfrak r_i$
au-dessus de $\mathfrak p$ tel que
$\kappa(\mathfrak p) \subset \kappa(\mathfrak r_i)$ soit purement inséparable,
et que $R \to B$ ne soit quasi-fini en aucun idéal premier au-dessus de $\mathfrak p$.
Comme $R \to S$ est quasi-fini (voir le
lemme \ref{algebra-lemma-unramified-quasi-finite}),
nous voyons qu'il n'existe aucun idéal premier de $B$ au-dessus de $\mathfrak p$.
D'après le lemme \ref{algebra-lemma-unramified-at-prime}, nous voyons que
$\kappa(\mathfrak r_i)/\kappa(\mathfrak p)$
est séparable, donc est l'extension de corps triviale, et que
$\mathfrak p(A_i)_{\mathfrak r_i}$ est l'idéal maximal.
De plus, d'après le lemme \ref{algebra-lemma-unique-prime-over-localize-below}
(qui s'applique à $R \to A_i$ parce qu'un morphisme d'anneaux fini satisfait la montée d'après le
lemme \ref{algebra-lemma-integral-going-up}),
nous avons $(A_i)_{\mathfrak r_i} = (A_i)_{\mathfrak p}$.
Il résulte du lemme de Nakayama \ref{algebra-lemma-NAK}
que le morphisme d'anneaux locaux
$R_{\mathfrak p} \to (A_i)_{\mathfrak p} = (A_i)_{\mathfrak r_i}$
est surjectif. Comme $A_i$ est fini sur $R$, nous voyons qu'il
existe un $f \in R$, $f \not \in \mathfrak p$, tel que
$R_f \to (A_i)_f$ soit surjectif. Après avoir remplacé $R$ par $R_f$, on conclut.
\end{proof}








\section{Anneaux locaux henséliens}
\label{algebra-section-henselian}

\noindent
Dans cette section, nous examinons brièvement la notion d'anneau local hensélien.
Soit $(R, \mathfrak m, \kappa)$ un anneau local.
Pour $a \in R$, nous notons $\overline{a}$ l'image de $a$ dans $\kappa$.
Pour un polynôme $f \in R[T]$, nous notons souvent $\overline{f}$
l'image de $f$ dans $\kappa[T]$.
Étant donné un polynôme $f \in R[T]$, nous notons $f'$ la dérivée
de $f$ par rapport à $T$. Notons que $\overline{f}' = \overline{f'}$.

\begin{definition}
\label{algebra-definition-henselian}
Soit $(R, \mathfrak m, \kappa)$ un anneau local.
\begin{enumerate}
\item On dit que $R$ est {\it hensélien} si, pour tout $f \in R[T]$ unitaire et
toute racine $a_0 \in \kappa$ de $\overline{f}$ telle que
$\overline{f'}(a_0) \not = 0$,
il existe un $a \in R$ tel que $f(a) = 0$ et
$a_0 = \overline{a}$.
\item On dit que $R$ est {\it strictement hensélien} si $R$ est hensélien
et si son corps résiduel est séparablement clos.
\end{enumerate}
\end{definition}

\noindent
Notons que la condition $\overline{f'}(a_0) \not = 0$ équivaut à la
condition que $a_0$ soit une racine simple du polynôme $\overline{f}$.
En fait, elle implique que le relèvement $a \in R$, s'il existe, est unique.

\begin{lemma}
\label{algebra-lemma-uniqueness}
Soit $(R, \mathfrak m, \kappa)$ un anneau local.
Soit $f \in R[T]$. Soient $a, b \in R$ tels que $f(a) = f(b) = 0$,
$a = b \bmod \mathfrak m$, et $f'(a) \not \in \mathfrak m$.
Alors $a = b$.
\end{lemma}

\begin{proof}
Écrivons $f(x + y) - f(x) = f'(x)y + g(x, y) y^2$ dans $R[x, y]$ (cela est possible,
comme on le voit en développant $f(x + y)$ ; nous omettons les détails).
Nous voyons alors que $0 = f(b) - f(a) = f(a + (b - a)) - f(a) =
f'(a)(b - a) + c (b - a)^2$ pour un certain $c \in R$. Par hypothèse,
$f'(a)$ est une unité de $R$. Ainsi, $(b - a)(1 + f'(a)^{-1}c(b - a)) = 0$.
Par hypothèse, $b - a \in \mathfrak m$, donc $1 + f'(a)^{-1}c(b - a)$
est une unité de $R$. Ainsi, $b - a = 0$ dans $R$.
\end{proof}

\noindent
Voici la caractérisation des anneaux locaux henséliens.

\begin{lemma}
\label{algebra-lemma-characterize-henselian}
\begin{slogan}
Caractérisations des anneaux locaux henséliens
\end{slogan}
Soit $(R, \mathfrak m, \kappa)$ un anneau local.
Les assertions suivantes sont équivalentes :
\begin{enumerate}
\item $R$ est hensélien,
\item pour tout $f \in R[T]$ et toute racine $a_0 \in \kappa$
de $\overline{f}$ telle que $\overline{f'}(a_0) \not = 0$,
il existe un $a \in R$ tel que $f(a) = 0$ et
$a_0 = \overline{a}$,
\item pour tout $f \in R[T]$ unitaire et toute factorisation
$\overline{f} = g_0 h_0$ avec $\gcd(g_0, h_0) = 1$, il
existe une factorisation $f = gh$ dans $R[T]$ telle que
$g_0 = \overline{g}$ et $h_0 = \overline{h}$,
\item pour tout $f \in R[T]$ unitaire et toute factorisation
$\overline{f} = g_0 h_0$ avec $\gcd(g_0, h_0) = 1$, il
existe une factorisation $f = gh$ dans $R[T]$ telle que
$g_0 = \overline{g}$ et $h_0 = \overline{h}$ et, de plus,
$\deg_T(g) = \deg_T(g_0)$,
\item pour tout $f \in R[T]$ et toute factorisation
$\overline{f} = g_0 h_0$ avec $\gcd(g_0, h_0) = 1$, il
existe une factorisation $f = gh$ dans $R[T]$ telle que
$g_0 = \overline{g}$ et $h_0 = \overline{h}$,
\item pour tout $f \in R[T]$ et toute factorisation
$\overline{f} = g_0 h_0$ avec $\gcd(g_0, h_0) = 1$, il
existe une factorisation $f = gh$ dans $R[T]$ telle que
$g_0 = \overline{g}$ et $h_0 = \overline{h}$ et
telle que, de plus, $\deg_T(g) = \deg_T(g_0)$,
\item pour tout morphisme d'anneaux \'etale $R \to S$ et tout idéal premier
$\mathfrak q$ de $S$ au-dessus de $\mathfrak m$ avec
$\kappa = \kappa(\mathfrak q)$, il existe une rétraction
$\tau : S \to R$ de $R \to S$,
\item pour tout morphisme d'anneaux \'etale $R \to S$ et tout idéal premier
$\mathfrak q$ de $S$ au-dessus de $\mathfrak m$ avec
$\kappa = \kappa(\mathfrak q)$, il existe une unique rétraction
$\tau : S \to R$ de $R \to S$ telle que
$\mathfrak q = \tau^{-1}(\mathfrak m)$,
\item toute $R$-algèbre finie est un produit d'anneaux locaux,
\item toute $R$-algèbre finie est un produit fini d'anneaux locaux,
\item toute $R$-algèbre $S$ de type fini peut s'écrire
$A \times B$, avec $R \to A$ fini
et $R \to B$ non quasi-fini en tout idéal premier au-dessus de $\mathfrak m$,
\item toute $R$-algèbre $S$ de type fini peut s'écrire
$A \times B$, avec $R \to A$ fini,
de sorte que toute composante irréductible de $\Spec(B \otimes_R \kappa)$
soit de dimension $\geq 1$, et
\item toute $R$-algèbre quasi-finie $S$ peut s'écrire
$S = A \times B$, avec $R \to A$ fini, de sorte que $B \otimes_R \kappa = 0$.
\end{enumerate}
\end{lemma}

\begin{proof}
Voici une liste des implications les plus faciles :
\begin{enumerate}
\item 2$\Rightarrow$1 parce que, dans (2), nous considérons tous les polynômes et
dans (1) seulement ceux qui sont unitaires,
\item 5$\Rightarrow$3 parce que, dans (5), nous considérons tous les polynômes et
dans (3) seulement ceux qui sont unitaires,
\item 6$\Rightarrow$4 parce que, dans (6), nous considérons tous les polynômes et
dans (4) seulement ceux qui sont unitaires,
\item 4$\Rightarrow$3 est évidente,
\item 6$\Rightarrow$5 est évidente,
\item 8$\Rightarrow$7 est évidente,
\item 10$\Rightarrow$9 est évidente,
\item 11$\Leftrightarrow$12 par définition du fait d'être quasi-fini en un idéal premier,
\item 11$\Rightarrow$13 par définition du fait d'être quasi-fini,
\end{enumerate}

\noindent
Démonstration de 1$\Rightarrow$8. Supposons (1).
Soit $R \to S$ \'etale, et soit $\mathfrak q \subset S$
un idéal premier tel que $\kappa(\mathfrak q) \cong \kappa$. D'après la
proposition \ref{algebra-proposition-etale-locally-standard},
nous pouvons trouver un $g \in S$, $g \not \in \mathfrak q$, tel que
$R \to S_g$ soit \'etale standard. Après avoir remplacé $S$ par $S_g$, nous pouvons supposer
que $S = R[t]_g/(f)$ est \'etale standard (détails omis).
Puisque l'idéal premier $\mathfrak q$ a pour corps résiduel $\kappa$, il
correspond à une racine $a_0$ de $\overline{f}$ qui n'est pas une
racine de $\overline{g}$. Par définition d'une algèbre \'etale standard,
cela signifie aussi que $\overline{f'}(a_0) \not = 0$.
Puisque $f$ est également unitaire, à nouveau par définition d'une algèbre \'etale standard,
nous pouvons utiliser le fait que $R$ est hensélien pour conclure qu'il existe un $a \in R$
tel que $a_0 = \overline{a}$ et $f(a) = 0$. Cela implique que
$g(a)$ est une unité de $R$, et nous obtenons le morphisme recherché
$\tau : S = R[t]_g/(f) \to R$ par la règle $t \mapsto a$. Par construction,
$\tau^{-1}(\mathfrak m) = \mathfrak q$. D'après le lemme \ref{algebra-lemma-uniqueness},
le morphisme $\tau$ est unique. Cela prouve (8).

\medskip\noindent
Démonstration de 7$\Rightarrow$8. (Ceci est vraiment sans importance et devrait être
omis.) Supposons (7) vérifiée et supposons que $R \to S$ soit \'etale.
Soient $\mathfrak q_1, \ldots, \mathfrak q_r$
les autres idéaux premiers de $S$ au-dessus de $\mathfrak m$.
Nous pouvons alors trouver un $g \in S$, $g \not \in \mathfrak q$, tel que
$g \in \mathfrak q_i$ pour $i = 1, \ldots, r$.
En effet, nous pouvons faire valoir que
$\bigcap_{i=1}^{r} \mathfrak{q}_{i} \not\subset \mathfrak{q}$
car, sinon,
$\mathfrak{q}_{i} \subset \mathfrak{q}$
pour un certain $i$, ce qui est impossible puisque la fibre d'un
morphisme \'etale est discrète (utiliser par exemple le
lemme \ref{algebra-lemma-etale-over-field}).
Appliquons (7) au morphisme d'anneaux \'etale
$R \to S_g$ et à l'idéal premier $\mathfrak qS_g$. Nous obtenons une rétraction
$\tau_g : S_g \to R$ telle que la composée $\tau : S \to S_g \to R$
ait la propriété $\tau^{-1}(\mathfrak m) = \mathfrak q$.
Nous omettons les détails.

\medskip\noindent
Démonstration de 8$\Rightarrow$11. Supposons (8) et soit $R \to S$ un morphisme
d'anneaux de type fini. Appliquons le
lemme \ref{algebra-lemma-etale-makes-quasi-finite-finite}.
Nous trouvons un morphisme d'anneaux \'etale $R \to R'$ et un idéal premier
$\mathfrak m' \subset R'$ au-dessus de $\mathfrak m$, avec
$\kappa = \kappa(\mathfrak m')$, tels que
$R' \otimes_R S = A' \times B'$, où $A'$ est fini sur $R'$
et $B'$ n'est quasi-fini sur $R'$ en aucun idéal premier au-dessus de $\mathfrak m'$.
Appliquons (8) pour obtenir une rétraction $\tau : R' \to R$ avec
$\mathfrak m = \tau^{-1}(\mathfrak m')$. Utilisons alors
$$
S = (S \otimes_R R') \otimes_{R', \tau} R
= (A' \times B') \otimes_{R', \tau} R
= (A' \otimes_{R', \tau} R)  \times  (B' \otimes_{R', \tau} R)
$$
ce qui donne une décomposition comme dans (11).

\medskip\noindent
Démonstration de 8$\Rightarrow$10. Supposons (8) et soit $R \to S$ un morphisme
d'anneaux fini. Appliquons le
lemme \ref{algebra-lemma-etale-makes-quasi-finite-finite}.
Nous trouvons un morphisme d'anneaux \'etale $R \to R'$ et un idéal premier
$\mathfrak m' \subset R'$ au-dessus de $\mathfrak m$, avec
$\kappa = \kappa(\mathfrak m')$, tels que
$R' \otimes_R S = A'_1 \times \ldots \times A'_n \times B'$,
où $A'_i$ est fini sur $R'$ et possède exactement un idéal premier au-dessus de
$\mathfrak m'$, tandis que $B'$ n'est quasi-fini sur $R'$ en aucun idéal premier
au-dessus de $\mathfrak m'$.
Appliquons (8) pour obtenir une rétraction $\tau : R' \to R$ avec
$\mathfrak m' = \tau^{-1}(\mathfrak m)$. Alors nous obtenons
\begin{align*}
S & = (S \otimes_R R') \otimes_{R', \tau} R \\
& = (A'_1 \times \ldots \times A'_n \times B') \otimes_{R', \tau} R \\
& = (A'_1 \otimes_{R', \tau} R)  \times
\ldots \times (A'_1 \otimes_{R', \tau} R) \times
(B' \otimes_{R', \tau} R) \\
& = A_1 \times \ldots \times A_n \times B
\end{align*}
Le facteur $B$ est fini sur $R$, mais $R \to B$
n'est quasi-fini en aucun idéal premier au-dessus de $\mathfrak m$. Par conséquent,
$B = 0$. Les facteurs $A_i$ sont des $R$-algèbres finies possédant exactement
un idéal premier au-dessus de $\mathfrak m$ ; ce sont donc des anneaux locaux.
Cela prouve que $S$ est un produit fini d'anneaux locaux.

\medskip\noindent
Démonstration de 9$\Rightarrow$10. Cela est vrai parce que, si $S$ est fini sur l'anneau
local $R$, alors il possède au plus un nombre fini d'idéaux maximaux. En effet, par la
montée pour $R \to S$, les idéaux maximaux de $S$ se trouvent tous au-dessus de
$\mathfrak m$, et $S/\mathfrak mS$ est artinien, donc possède un nombre fini d'idéaux premiers.

\medskip\noindent
Démonstration de 10$\Rightarrow$1. Supposons (10). Soit $f \in R[T]$ un polynôme
unitaire et soit $a_0 \in \kappa$ une racine simple de $\overline{f}$.
Alors $S = R[T]/(f)$ est une $R$-algèbre finie. En appliquant (10),
nous obtenons que $S = A_1 \times \ldots \times A_r$ est un produit fini de
$R$-algèbres locales. En particulier, nous voyons que
$S/\mathfrak mS = \prod A_i/\mathfrak mA_i$ est la décomposition
de $\kappa[T]/(\overline{f})$ en produit d'anneaux locaux.
Cela signifie que l'un des facteurs, disons $A_1/\mathfrak mA_1$,
est le quotient $\kappa[T]/(\overline{f}) \to \kappa[T]/(T - a_0)$.
Comme $A_1$ est un facteur direct du $R$-module libre de type fini $S$, c'est
lui-même un $R$-module libre de type fini. Comme $A_1/\mathfrak mA_1$ est un
$\kappa$-espace vectoriel de dimension 1, nous voyons que $A_1 \cong R$ comme
$R$-module. Cela signifie clairement que $R \to A_1$ est un isomorphisme.
Soit $a \in R$ l'image de $T$ par le morphisme
$R[T] \to S \to A_1 \to R$. Alors $f(a) = 0$ et $\overline{a} = a_0$,
comme désiré.

\medskip\noindent
Démonstration de 13$\Rightarrow$1. Supposons (13). Soit $f \in R[T]$ un polynôme
unitaire et soit $a_0 \in \kappa$ une racine simple de $\overline{f}$.
Alors $S_1 = R[T]/(f)$ est une $R$-algèbre finie. Soit $g \in R[T]$
un élément quelconque tel que $\overline{g} = \overline{f}/(T - a_0)$.
Alors $S = (S_1)_g$ est une $R$-algèbre quasi-finie telle que
$S \otimes_R \kappa \cong \kappa[T]_{\overline{g}}/(\overline{f})
\cong \kappa[T]/(T - a_0) \cong \kappa$.
En appliquant (13) à $S$, nous obtenons $S = A \times B$, où $A$ est fini sur $R$ et
$B \otimes_R \kappa = 0$. En particulier, nous voyons que
$\kappa \cong S/\mathfrak mS = A/\mathfrak mA$.
Comme $A$ est un facteur direct de la $R$-algèbre plate $S$, nous voyons
qu'il est fini plat, donc libre sur $R$.
Comme $A/\mathfrak mA$ est un
$\kappa$-espace vectoriel de dimension 1, nous voyons que $A \cong R$ comme
$R$-module. Cela signifie clairement que $R \to A$ est un isomorphisme.
Soit $a \in R$ l'image de $T$ par le morphisme
$R[T] \to S \to A \to R$. Alors $f(a) = 0$ et $\overline{a} = a_0$,
comme désiré.

\medskip\noindent
Démonstration de 8$\Rightarrow$2. Supposons (8). Soit $f \in R[T]$ un
polynôme quelconque et soit $a_0 \in \kappa$ une racine simple. Alors
l'algèbre $S = R[T]_{f'}/(f)$ est \'etale sur $R$.
Soit $\mathfrak q \subset S$ l'idéal premier
engendré par $\mathfrak m$ et $T - b$, où $b \in R$ est un
élément quelconque tel que $\overline{b} = a_0$. Appliquons (8) à $S$ et
$\mathfrak q$ pour obtenir $\tau : S \to R$.
Alors l'image $\tau(T) = a \in R$ convient dans (2).

\medskip\noindent
À ce stade, nous voyons que (1), (2), (7), (8), (9), (10), (11), (12), (13) sont
toutes équivalentes. La plus faible des assertions (3), (4), (5) et (6)
est (3), et la plus forte est (6). Il nous reste donc à prouver que
(3) implique (1) et que (1) implique (6).

\medskip\noindent
Démonstration de 3$\Rightarrow$1. Supposons (3). Soit $f \in R[T]$ unitaire et
soit $a_0 \in \kappa$ une racine simple de $\overline{f}$. Cela donne
une factorisation $\overline{f} = (T - a_0)h_0$ avec $h_0(a_0) \not = 0$,
donc $\gcd(T - a_0, h_0) = 1$. Appliquons (3) pour obtenir une factorisation
$f = gh$ avec $\overline{g} = T - a_0$ et $\overline{h} = h_0$.
Posons $S = R[T]/(f)$, qui est une $R$-algèbre libre de type fini. Nous noterons
aussi $g$, $h$ les images de $g$ et $h$ dans $S$. Alors
$gS + hS = S$ d'après le
lemme de Nakayama \ref{algebra-lemma-NAK},
puisque l'égalité est vraie modulo $\mathfrak m$. Comme $gh = f = 0$ dans $S$,
cela implique aussi que $gS \cap hS = 0$. Le théorème chinois des restes
donne donc $S = S/(g) \times S/(h)$. Cela implique que
$A = S/(g)$ est un facteur direct d'un $R$-module libre de type fini, donc est
libre de type fini. De plus, le rang de $A$ est $1$, car
$A/\mathfrak mA = \kappa[T]/(T - a_0)$. Le morphisme $R \to A$
est donc un isomorphisme. En prenant pour $a \in R$ l'image de $T$
par les morphismes $R[T] \to S \to A \to R$, nous obtenons un élément de $R$
tel que $f(a) = 0$ et $\overline{a} = a_0$.

\medskip\noindent
Démonstration de 1$\Rightarrow$6. Supposons (1), ou, de façon équivalente, toutes les
assertions (1), (2), (7), (8), (9), (10), (11), (12), (13).
Soit $f \in R[T]$ un polynôme.
Supposons que $\overline{f} = g_0h_0$ soit une factorisation avec
$\gcd(g_0, h_0) = 1$. Nous pouvons supposer, et supposons, que $g_0$ est unitaire.
Considérons $S = R[T]/(f)$. Puisque nous
disposons de cette factorisation, nous voyons que les coefficients de
$f$ engendrent l'idéal unité de $R$.
Cela implique que $S$ a des fibres finies sur $R$, et est donc
quasi-fini sur $R$. Cela implique également que $S$ est plat sur $R$ d'après le
lemme \ref{algebra-lemma-grothendieck-general}.
En combinant (13) et (10), nous pouvons écrire
$S = A_1 \times \ldots \times A_n \times B$,
où chaque $A_i$ est local et fini sur $R$, et
$B \otimes_R \kappa = 0$. Après avoir réordonné les facteurs $A_1, \ldots, A_n$,
nous pouvons supposer que
$$
\kappa[T]/(g_0) =
A_1/\mathfrak m A_1 \times \ldots \times A_r/\mathfrak mA_r,
\ \kappa[T]/(h_0) =
A_{r + 1}/\mathfrak mA_{r + 1} \times \ldots \times A_n/\mathfrak mA_n
$$
comme quotients de $\kappa[T]$. La $R$-algèbre finie plate
$A = A_1 \times \ldots \times A_r$ est libre comme $R$-module, voir le
lemme \ref{algebra-lemma-finite-flat-local}.
Son rang est $\deg_T(g_0)$. Soit $g \in R[T]$ le polynôme caractéristique
de l'opérateur $R$-linéaire $T : A \to A$. Alors $g$ est un polynôme unitaire
de degré $\deg_T(g) = \deg_T(g_0)$ et, de plus, $\overline{g} = g_0$.
Par Cayley-Hamilton
(lemme \ref{algebra-lemma-charpoly}),
nous voyons que $g(T_A) = 0$, où $T_A$ désigne
l'image de $T$ dans $A$. Nous obtenons donc un morphisme surjectif bien défini
$R[T]/(g) \to A$, qui est un isomorphisme d'après le
lemme de Nakayama \ref{algebra-lemma-NAK}. Le morphisme $R[T] \to A$ se factorise
par $R[T]/(f)$ par construction ; nous pouvons donc écrire $f = gh$ pour
un certain $h$. Cela achève la démonstration.
\end{proof}

\begin{lemma}
\label{algebra-lemma-finite-over-henselian}
Soit $(R, \mathfrak m, \kappa)$ un anneau local hensélien.
\begin{enumerate}
\item Si $R \to S$ est un morphisme d'anneaux fini, alors $S$ est
un produit fini d'anneaux locaux henséliens, chacun fini sur $R$.
\item Si $R \to S$ est un morphisme d'anneaux fini et si $S$ est local, alors
$S$ est un anneau local hensélien et $R \to S$ est un morphisme (fini) d'anneaux locaux.
\item Si $R \to S$ est un morphisme d'anneaux de type fini, et si $\mathfrak q$ est
un idéal premier de $S$ au-dessus de $\mathfrak m$ auquel $R \to S$ est quasi-fini,
alors $S_{\mathfrak q}$ est hensélien et fini sur $R$.
\item Si $R \to S$ est quasi-fini, alors $S_{\mathfrak q}$ est hensélien
et fini sur $R$ pour tout idéal premier $\mathfrak q$ au-dessus de $\mathfrak m$.
\end{enumerate}
\end{lemma}

\begin{proof}
La partie (2) implique la partie (1), puisque $S$, comme dans la partie (1), est un produit fini
de ses localisations aux idéaux premiers au-dessus de $\mathfrak m$ d'après le
lemme \ref{algebra-lemma-characterize-henselian}, partie (10).
La partie (2) résulte aussi du lemme \ref{algebra-lemma-characterize-henselian}, partie (10),
puisque toute $S$-algèbre finie est également une $R$-algèbre finie (bien sûr,
tout morphisme d'anneaux fini entre anneaux locaux est local).

\medskip\noindent
Soient $R \to S$ et $\mathfrak q$ comme dans (3). Écrivons $S = A \times B$,
où $A$ est fini sur $R$ et $B$ n'est quasi-fini sur $R$ en aucun
idéal premier au-dessus de $\mathfrak m$, voir le
lemme \ref{algebra-lemma-characterize-henselian}, partie (11).
Ainsi, $S_\mathfrak q$ est une localisation de $A$ en un idéal maximal,
et nous déduisons (3) de (1). La partie (4) résulte de la partie (3).
\end{proof}

\begin{lemma}
\label{algebra-lemma-mop-up}
Soit $(R, \mathfrak m, \kappa)$ un anneau local hensélien.
Toute $R$-algèbre $S$ de type fini peut s'écrire
$S = A_1 \times \ldots \times A_n \times B$, où les $A_i$ sont locaux
et finis sur $R$, et où $R \to B$ n'est quasi-fini en aucun
idéal premier de $B$ au-dessus de $\mathfrak m$.
\end{lemma}

\begin{proof}
C'est une combinaison des parties (11) et (10) du
lemme \ref{algebra-lemma-characterize-henselian}.
\end{proof}

\begin{lemma}
\label{algebra-lemma-mop-up-strictly-henselian}
Soit $(R, \mathfrak m, \kappa)$ un anneau local strictement hensélien.
Toute $R$-algèbre $S$ de type fini peut s'écrire
$S = A_1 \times \ldots \times A_n \times B$, où les $A_i$ sont locaux
et finis sur $R$, où $\kappa \subset \kappa(\mathfrak m_{A_i})$ est
finie purement inséparable, et où $R \to B$ n'est quasi-fini
en aucun idéal premier de $B$ au-dessus de $\mathfrak m$.
\end{lemma}

\begin{proof}
Écrivons d'abord $S = A_1 \times \ldots \times A_n \times B$ comme dans le
lemme \ref{algebra-lemma-mop-up}.
L'extension de corps $\kappa(\mathfrak m_{A_i})/\kappa$
est finie et $\kappa$ est séparablement clos ; elle
est donc finie purement inséparable.
\end{proof}

\begin{lemma}
\label{algebra-lemma-henselian-cat-finite-etale}
Soit $(R, \mathfrak m, \kappa)$ un anneau local hensélien.
La catégorie des extensions finies étales d'anneaux $R \to S$ est
équivalente à la catégorie des algèbres finies étales
$\kappa \to \overline{S}$ par le foncteur $S \mapsto S/\mathfrak mS$.
\end{lemma}

\begin{proof}
Notons $\mathcal{C} \to \mathcal{D}$ le foncteur entre les catégories
de l'énoncé.
Supposons que $R \to S$ soit fini étale. Alors nous pouvons écrire
$$
S = A_1 \times \ldots \times A_n
$$
avec les $A_i$ locaux et finis étales sur $S$, en utilisant soit le
lemme \ref{algebra-lemma-mop-up},
soit le
lemme \ref{algebra-lemma-characterize-henselian}, partie (10).
En particulier, $A_i/\mathfrak mA_i$ est une extension de corps finie
séparable de $\kappa$, voir le
lemme \ref{algebra-lemma-etale-at-prime}.
Nous voyons ainsi que tout objet de $\mathcal{C}$ et de
$\mathcal{D}$ se décompose canoniquement en composantes irréductibles
qui se correspondent par le foncteur donné.
Supposons ensuite que $S_1$, $S_2$ soient finis étales sur $R$, et que
$\kappa_1 = S_1/\mathfrak mS_1$ et $\kappa_2 = S_2/\mathfrak mS_2$
soient des corps (finis séparables sur $\kappa$). Alors $S_1 \otimes_R S_2$
est fini étale sur $R$, et nous pouvons écrire
$$
S_1 \otimes_R S_2 = A_1 \times \ldots \times A_n
$$
comme précédemment. Nous voyons alors que $\Hom_R(S_1, S_2)$ s'identifie
à l'ensemble des indices $i \in \{1, \ldots, n\}$ tels que
$S_2 \to A_i$ soit un isomorphisme. Pour le voir, utiliser que, pour tout morphisme
de $R$-algèbres $\varphi : S_1 \to S_2$, le morphisme
$\varphi \times 1 : S_1 \otimes_R S_2 \to S_2$
est surjectif, et est donc égal à la projection sur l'un des facteurs $A_i$.
Mais nous voyons exactement de la même manière que
$\Hom_\kappa(\kappa_1, \kappa_2)$ s'identifie à
l'ensemble des indices $i \in \{1, \ldots, n\}$ tels que
$\kappa_2 \to A_i/\mathfrak mA_i$ soit un isomorphisme.
D'après la discussion ci-dessus, ces ensembles d'indices coïncident, et nous concluons
que notre foncteur est pleinement fidèle.
Enfin, soit $\kappa'/\kappa$ une extension de corps finie
séparable. D'après le
lemme \ref{algebra-lemma-make-etale-map-prescribed-residue-field},
il existe un morphisme d'anneaux \'etale $R \to S$ et un idéal premier
$\mathfrak q$ de $S$ au-dessus de $\mathfrak m$ tels que
$\kappa \subset \kappa(\mathfrak q)$ soit isomorphe à l'extension donnée.
D'après la partie (1), nous pouvons écrire
$S = A_1 \times \ldots \times A_n \times B$.
Comme $R \to S$ est quasi-fini, nous voyons qu'il n'existe aucun
idéal premier de $B$ au-dessus de $\mathfrak m$. Ainsi, $S_{\mathfrak q}$ est
égal à $A_i$ pour un certain $i$. Par conséquent, $R \to A_i$ est fini étale
et produit l'extension de corps résiduels donnée. Le foncteur est donc
essentiellement surjectif, ce qui conclut la démonstration.
\end{proof}

\begin{lemma}
\label{algebra-lemma-unramified-over-strictly-henselian}
Soit $(R, \mathfrak m, \kappa)$ un anneau local strictement hensélien.
Soit $R \to S$ un morphisme d'anneaux non ramifié. Alors
$$
S = A_1 \times \ldots \times A_n \times B
$$
où chaque $R \to A_i$ est surjectif et où aucun idéal premier de $B$ ne se trouve
au-dessus de $\mathfrak m$.
\end{lemma}

\begin{proof}
Écrivons d'abord $S = A_1 \times \ldots \times A_n \times B$ comme dans le
lemme \ref{algebra-lemma-mop-up}.
Nous voyons maintenant que $R \to A_i$ est fini non ramifié et que $A_i$ est local.
L'idéal maximal de $A_i$ est donc $\mathfrak mA_i$, et son
corps résiduel $A_i / \mathfrak m A_i$ est une extension finie
séparable de $\kappa$, voir le
lemme \ref{algebra-lemma-unramified-at-prime}.
Cependant, la condition que $R$ soit strictement hensélien signifie que
$\kappa$ est séparablement clos, donc
$\kappa = A_i / \mathfrak m A_i$. D'après le
lemme de Nakayama \ref{algebra-lemma-NAK},
nous concluons que $R \to A_i$ est surjectif, comme désiré.
\end{proof}

\begin{lemma}
\label{algebra-lemma-complete-henselian}
\begin{slogan}
Les anneaux locaux complets sont henséliens par la méthode de Newton
\end{slogan}
Soit $(R, \mathfrak m, \kappa)$ un anneau local complet, voir la
définition \ref{algebra-definition-complete-local-ring}.
Alors $R$ est hensélien.
\end{lemma}

\begin{proof}
Soit $f \in R[T]$ unitaire.
Notons $f_n \in R/\mathfrak m^{n + 1}[T]$ son image.
Notons $f'_n$ la dérivée de $f_n$ par rapport à $T$.
Soit $a_0 \in \kappa$ une racine simple de $f_0$. Nous la relevons
par récurrence en une solution de $f$ sur $R$ comme suit :
supposons donné $a_n \in R/\mathfrak m^{n + 1}$ tel que
$a_n \bmod \mathfrak m = a_0$ et $f_n(a_n) = 0$. Choisissons un
élément quelconque $b \in R/\mathfrak m^{n + 2}$ tel que
$a_n = b \bmod \mathfrak m^{n + 1}$. Alors
$f_{n + 1}(b) \in \mathfrak m^{n + 1}/\mathfrak m^{n + 2}$.
Posons
$$
a_{n + 1} = b - f_{n + 1}(b)/f'_{n + 1}(b)
$$
(méthode de Newton). Cela a un sens, car
$f'_{n + 1}(b) \in R/\mathfrak m^{n + 2}$
est inversible par la condition sur $a_0$. Nous calculons alors
$f_{n + 1}(a_{n + 1}) = f_{n + 1}(b) - f_{n + 1}(b) = 0$
dans $R/\mathfrak m^{n + 2}$. Comme le système d'éléments
$a_n \in R/\mathfrak m^{n + 1}$ ainsi construit est compatible,
nous obtenons un élément
$a \in \lim R/\mathfrak m^n = R$ (nous utilisons ici que $R$ est complet).
De plus, $f(a) = 0$, puisqu'il s'envoie sur zéro dans chaque
$R/\mathfrak m^n$.
Enfin, $\overline{a} = a_0$, ce qui conclut.
\end{proof}

\begin{lemma}
\label{algebra-lemma-local-dimension-zero-henselian}
\begin{slogan}
Les anneaux locaux de dimension zéro sont henséliens.
\end{slogan}
Soit $(R, \mathfrak m)$ un anneau local de dimension $0$.
Alors $R$ est hensélien.
\end{lemma}

\begin{proof}
Soit $R \to S$ un morphisme d'anneaux fini. D'après le
lemme \ref{algebra-lemma-characterize-henselian},
il suffit de montrer que $S$ est un produit d'anneaux locaux. D'après le
lemme \ref{algebra-lemma-finite-finite-fibres},
$S$ possède un nombre fini d'idéaux premiers
$\mathfrak m_1, \ldots, \mathfrak m_r$,
qui se trouvent tous au-dessus de $\mathfrak m$. Il n'existe aucune inclusion entre ces
idéaux premiers, voir le
lemme \ref{algebra-lemma-integral-no-inclusion} ;
ils sont donc tous maximaux. Tout élément de
$\mathfrak m_1 \cap \ldots \cap \mathfrak m_r$ est nilpotent d'après le
lemme \ref{algebra-lemma-Zariski-topology}.
Il s'ensuit que $S$ est le produit des localisations de $S$ aux idéaux premiers
$\mathfrak m_i$, d'après le
lemme \ref{algebra-lemma-product-local}.
\end{proof}

\noindent
Le lemme suivant sera la clé des propriétés d'unicité et de fonctorialité
de l'hensélisation et de l'hensélisation stricte.

\begin{lemma}
\label{algebra-lemma-map-into-henselian}
Soit $R \to S$ un morphisme d'anneaux, avec $S$ local hensélien.
Étant donnés
\begin{enumerate}
\item un morphisme d'anneaux \'etale $R \to A$,
\item un idéal premier $\mathfrak q$ de $A$ au-dessus de
$\mathfrak p = R \cap \mathfrak m_S$,
\item un morphisme de $\kappa(\mathfrak p)$-algèbres
$\tau : \kappa(\mathfrak q) \to S/\mathfrak m_S$,
\end{enumerate}
il existe un unique morphisme de $R$-algèbres $f : A \to S$
tel que $\mathfrak q = f^{-1}(\mathfrak m_S)$ et que
$f$ induise le morphisme $\tau$ sur les corps résiduels.
\end{lemma}

\begin{proof}
Considérons $A \otimes_R S$. C'est une algèbre \'etale sur $S$, voir le
lemme \ref{algebra-lemma-etale}. De plus, le noyau
$$
\mathfrak q' = \Ker(A \otimes_R S \to
\kappa(\mathfrak q) \otimes_{\kappa(\mathfrak p)} \kappa(\mathfrak m_S)
\xrightarrow{\tau \otimes 1}
\kappa(\mathfrak m_S))
$$
est un idéal premier au-dessus de $\mathfrak m_S$ dont le corps résiduel est égal
à celui de $S$. Ainsi, d'après le lemme \ref{algebra-lemma-characterize-henselian},
il existe une unique rétraction $\sigma : A \otimes_R S \to S$
avec $\sigma^{-1}(\mathfrak m_S) = \mathfrak q'$.
Prenons pour $f$ la composée $A \to A \otimes_R S \to S$.
Nous omettons la vérification des propriétés de $f$ ; l'unicité
de $f$ résulte de celle de $\sigma$ (détails omis).
\end{proof}

\begin{lemma}
\label{algebra-lemma-strictly-henselian-solutions}
Soit $\varphi : R \to S$ un morphisme local
d'anneaux locaux strictement henséliens.
Soient $P_1, \ldots, P_n \in R[x_1, \ldots, x_n]$ des polynômes tels que
$R[x_1, \ldots, x_n]/(P_1, \ldots, P_n)$ soit \'etale sur $R$.
Alors le morphisme
$$
R^n \longrightarrow S^n, \quad
(h_1, \ldots, h_n) \longmapsto (\varphi(h_1), \ldots, \varphi(h_n))
$$
induit une bijection entre
$$
\{
(r_1, \ldots, r_n) \in R^n
\mid
P_i(r_1, \ldots, r_n) = 0, \ i = 1, \ldots, n
\}
$$
et
$$
\{
(s_1, \ldots, s_n) \in S^n
\mid
P^\varphi_i(s_1, \ldots, s_n) = 0, \ i = 1, \ldots, n
\}
$$
où les $P^\varphi_i \in S[x_1, \ldots, x_n]$ sont les images des $P_i$
par $\varphi$.
\end{lemma}

\begin{proof}
Le premier ensemble de solutions est canoniquement isomorphe à l'ensemble
$$
\Hom_R(R[x_1, \ldots, x_n]/(P_1, \ldots, P_n), R).
$$
Comme $R$ est hensélien, le morphisme $R \to R/\mathfrak m_R$ induit une bijection
entre cet ensemble et l'ensemble des solutions dans le
corps résiduel $R/\mathfrak m_R$, voir le
lemme \ref{algebra-lemma-characterize-henselian}.
La même chose est vraie pour $S$.
Puisque $R[x_1, \ldots, x_n]/(P_1, \ldots, P_n)$ est \'etale sur $R$
et que $R/\mathfrak m_R$ est séparablement clos, nous voyons que
$R/\mathfrak m_R[x_1, \ldots, x_n]/(\overline{P}_1, \ldots, \overline{P}_n)$
est un produit fini de copies de $R/\mathfrak m_R$, où $\overline{P}_i$
est l'image de $P_i$ dans $R/\mathfrak m_R[x_1, \ldots, x_n]$. Par conséquent, le
produit tensoriel
$$
R/\mathfrak m_R[x_1, \ldots, x_n]/(\overline{P}_1, \ldots, \overline{P}_n)
\otimes_{R/\mathfrak m_R} S/\mathfrak m_S
=
S/\mathfrak m_S[x_1, \ldots, x_n]/
(\overline{P}^\varphi_1, \ldots, \overline{P}^\varphi_n)
$$
est lui aussi un produit fini de copies de $S/\mathfrak m_S$ avec le même
ensemble d'indices. Cela prouve le lemme.
\end{proof}

\begin{lemma}
\label{algebra-lemma-split-ML-henselian}
Soit $R$ un anneau local hensélien.
Tout module de Mittag-Leffler à engendrement dénombrable sur $R$ est une somme
directe de $R$-modules de présentation finie.
\end{lemma}

\begin{proof}
Soit $M$ un $R$-module à engendrement dénombrable et de Mittag-Leffler.
Nous affirmons que, pour tout élément $x \in M$, il existe une décomposition en somme
directe $M = N \oplus K$ avec $x \in N$, le module
$N$ de présentation finie et $K$ de Mittag-Leffler.

\medskip\noindent
Supposons l'affirmation vraie. Choisissons des générateurs
$x_1, x_2, x_3, \ldots$ de $M$. D'après l'affirmation, nous pouvons trouver
par récurrence des décompositions en somme directe
$$
M = N_1 \oplus N_2 \oplus \ldots \oplus N_n \oplus K_n
$$
avec les $N_i$ de présentation finie,
$x_1, \ldots, x_n \in N_1 \oplus \ldots \oplus N_n$, et $K_n$ de Mittag-Leffler.
En répétant cela indéfiniment, nous voyons que $M = \bigoplus N_i$.

\medskip\noindent
Il nous reste à démontrer l'affirmation. Soit $x \in M$. D'après le
lemme \ref{algebra-lemma-ML-countable},
il existe un endomorphisme $\alpha : M \to M$
tel que $\alpha$ se factorise par un module de présentation finie et
$\alpha (x) = x$. Disons que $\alpha$ se factorise comme
$$
\xymatrix{
M \ar[r]^\pi & P \ar[r]^i & M
}
$$
Posons $a = \pi \circ \alpha \circ i : P \to P$, de sorte que
$i \circ a \circ \pi = \alpha^3$. D'après le
lemme \ref{algebra-lemma-charpoly-module},
il existe un polynôme unitaire $P \in R[T]$ tel que $P(a) = 0$.
Notons que cela implique formellement que $\alpha^2 P(\alpha) = 0$.
Nous pouvons donc considérer $M$ comme un module sur $R[T]/(T^2P)$.
Supposons que $x \not = 0$. Alors $\alpha(x) = x$ implique que
$0 = \alpha^2P(\alpha)x = P(1)x$, donc que $P(1) = 0$ dans $R/I$, où
$I = \{r \in R \mid rx = 0\}$ est l'annulateur de $x$.
Comme $x \not = 0$, nous voyons que $I \subset \mathfrak m_R$, donc que
$1$ est une racine de $\overline{P} = P \bmod \mathfrak m_R \in R/\mathfrak m_R[T]$.
Comme $R$ est hensélien, nous pouvons trouver une factorisation
$$
T^2P = (T^2 Q_1) Q_2
$$
pour certains $Q_1, Q_2 \in R[T]$ avec
$Q_2 = (T - 1)^e \bmod \mathfrak m_R R[T]$ et
$Q_1(1) \not = 0 \bmod \mathfrak m_R$, voir le
lemme \ref{algebra-lemma-characterize-henselian}.
Soient $N = \Im(\alpha^2Q_1(\alpha) : M \to M)$ et
$K = \Im(Q_2(\alpha) : M \to M)$. Comme $T^2Q_1$ et
$Q_2$ engendrent l'idéal unité de $R[T]$, nous obtenons une décomposition en somme
directe $M = N \oplus K$. De plus, $Q_2$ agit par zéro sur $N$, et
$T^2Q_1$ agit par zéro sur $K$. Notons que $N$ est un quotient de $P$,
donc est de type fini. De plus, $x \in N$, parce que
$\alpha^2Q_1(\alpha)x = Q_1(1)x$ et que $Q_1(1)$ est une unité de $R$. D'après le
lemme \ref{algebra-lemma-direct-sum-ML},
les modules $N$ et $K$ sont de Mittag-Leffler. Enfin, le module de type fini
$N$ est de présentation finie, puisqu'un module de Mittag-Leffler de type fini
est de présentation finie, voir
l'exemple \ref{algebra-example-ML}, partie (1).
\end{proof}


\section{Colimites filtrantes de morphismes d'anneaux \'etales}
\label{algebra-section-ind-etale}

\noindent
Cette section prépare la section sur les morphismes d'anneaux ind-\'etales
(Cohomologie pro-\'etale, section \ref{proetale-section-ind-etale}).
Les résultats seront également utiles pour démontrer des propriétés d'unicité de
l'hensélisation et de l'hensélisation stricte d'un anneau local.

\begin{lemma}
\label{algebra-lemma-base-change-colimit-etale}
Soient $R \to A$ et $R \to R'$ des morphismes d'anneaux. Si $A$ est une colimite
filtrante de $R$-algèbres \'etales, alors $R' \otimes_R A$
est une colimite filtrante de $R'$-algèbres \'etales.
\end{lemma}

\begin{proof}
Cela est vrai parce que les colimites commutent aux produits tensoriels
et que les morphismes d'anneaux \'etales sont préservés par changement de base
(lemme \ref{algebra-lemma-etale}).
\end{proof}

\begin{lemma}
\label{algebra-lemma-composition-colimit-etale}
Soient $A \to B \to C$ des morphismes d'anneaux. Si $B$ est une colimite
filtrante de $A$-algèbres \'etales et si $C$ est une colimite filtrante
de $B$-algèbres \'etales, alors $C$ est une colimite filtrante
de $A$-algèbres \'etales.
\end{lemma}

\begin{proof}
Nous utiliserons le critère du lemme \ref{algebra-lemma-when-colimit}.
Soit $A \to P \to C$ une factorisation de $A \to C$
avec $P$ de présentation finie sur $A$.
Écrivons $B = \colim_{i \in I} B_i$, où $I$ est un ensemble ordonné filtrant et
où $B_i$ est une $A$-algèbre \'etale.
Écrivons $C = \colim_{j \in J} C_j$, où $J$ est un ensemble ordonné filtrant et
où $C_j$ est une $B$-algèbre \'etale.
Nous pouvons factoriser $P \to C$ sous la forme $P \to C_j \to C$ pour un
certain $j$, d'après le lemme \ref{algebra-lemma-characterize-finite-presentation}.
D'après le lemme \ref{algebra-lemma-etale}, nous pouvons trouver
un $i \in I$ et un morphisme d'anneaux \'etale $B_i \to C'_j$
tels que $C_j = B \otimes_{B_i} C'_j$.
Alors $C_j = \colim_{i' \geq i} B_{i'} \otimes_{B_i} C'_j$,
et nous voyons à nouveau que $P \to C_j$ se factorise sous la forme
$P \to B_{i'} \otimes_{B_i} C'_j \to C$.
Le morphisme $A \to C' = B_{i'} \otimes_{B_i} C'_j$ est \'etale, car les
composées et les produits tensoriels de morphismes d'anneaux \'etales
sont \'etales. Nous avons donc factorisé $P \to C$ sous la forme
$P \to C' \to C$, avec $C'$ \'etale sur $A$, et le critère du
lemme \ref{algebra-lemma-when-colimit} s'applique.
\end{proof}

\begin{lemma}
\label{algebra-lemma-colimit-colimit-etale}
Soit $R$ un anneau. Soit $A = \colim A_i$ une colimite filtrante
de $R$-algèbres telle que chaque $A_i$ soit une colimite filtrante de
$R$-algèbres \'etales. Alors $A$ est une colimite filtrante de
$R$-algèbres \'etales.
\end{lemma}

\begin{proof}
Écrivons $A_i = \colim_{j \in J_i} A_j$, où $J_i$ est un ensemble ordonné filtrant et
où $A_j$ est une $R$-algèbre \'etale.
Pour tous $i \leq i'$ et $j \in J_i$, il existe un
$j' \in J_{i'}$ et un morphisme de $R$-algèbres
$\varphi_{jj'} : A_j \to A_{j'}$ rendant commutatif le diagramme
$$
\xymatrix{
A_i \ar[r] & A_{i'} \\
A_j \ar[u] \ar[r]^{\varphi_{jj'}} & A_{j'} \ar[u]
}
$$
Cela est vrai parce que $R \to A_j$ est de présentation finie,
de sorte que le lemme \ref{algebra-lemma-characterize-finite-presentation} s'applique.
Soit $\mathcal{J}$ la catégorie dont les objets sont
$\coprod_{i \in I} J_i$ et dont les morphismes sont les triplets
$(j, j', \varphi_{jj'})$ ci-dessus (avec la loi de composition évidente).
Alors $\mathcal{J}$ est une catégorie filtrante et
$A = \colim_\mathcal{J} A_j$. Nous omettons les détails.
\end{proof}

\begin{lemma}
\label{algebra-lemma-colimit-colimit-etale-better}
Soit $I$ un ensemble ordonné filtrant. Soit
$i \mapsto (R_i \to A_i)$ un système de flèches d'anneaux indexé par $I$.
Posons $R = \colim R_i$ et $A = \colim A_i$. Si chaque $A_i$ est une
colimite filtrante de $R_i$-algèbres \'etales, alors $A$ est une colimite filtrante
de $R$-algèbres \'etales.
\end{lemma}

\begin{proof}
Cela est vrai parce que $A = A \otimes_R R = \colim A_i \otimes_{R_i} R$ ;
nous pouvons donc appliquer le lemme \ref{algebra-lemma-colimit-colimit-etale},
puisque $R \to A_i \otimes_{R_i} R$ est une colimite filtrante de
morphismes d'anneaux \'etales d'après le
lemme \ref{algebra-lemma-base-change-colimit-etale}.
\end{proof}

\begin{lemma}
\label{algebra-lemma-colimits-of-etale}
Soit $R$ un anneau. Soit $A \to B$ un morphisme de $R$-algèbres.
Si $A$ et $B$ sont des colimites filtrantes de $R$-algèbres \'etales, alors
$B$ est une colimite filtrante de $A$-algèbres \'etales.
\end{lemma}

\begin{proof}
Écrivons $A = \colim A_i$ et $B = \colim B_j$ comme colimites filtrantes, avec
les $A_i$ et les $B_j$ \'etales sur $R$. Pour tout $i$, nous pouvons trouver un
$j$ tel que $A_i \to B$ se factorise par $B_j$, voir le
lemme \ref{algebra-lemma-characterize-finite-presentation}.
La factorisation $A_i \to B_j$ est \'etale d'après le
lemme \ref{algebra-lemma-map-between-etale}.
Comme $A \to A \otimes_{A_i} B_j$ est \'etale (lemme \ref{algebra-lemma-etale}),
il suffit de prouver que $B = \colim A \otimes_{A_i} B_j$, où la
colimite est prise sur les paires $(i, j)$ et les factorisations
$A_i \to B_j \to B$ de $A_i \to B$ (c'est un système filtrant ; détails omis).
Cela est clair parce que les colimites commutent aux produits tensoriels,
et donc $\colim A \otimes_{A_i} B_j = A \otimes_A B = B$.
\end{proof}

\begin{lemma}
\label{algebra-lemma-map-into-henselian-colimit}
Soit $R \to S$ un morphisme d'anneaux, avec $S$ local hensélien. Étant donnés
\begin{enumerate}
\item une $R$-algèbre $A$ qui est une colimite filtrante de $R$-algèbres \'etales,
\item un idéal premier $\mathfrak q$ de $A$ au-dessus de
$\mathfrak p = R \cap \mathfrak m_S$,
\item un morphisme de $\kappa(\mathfrak p)$-algèbres
$\tau : \kappa(\mathfrak q) \to S/\mathfrak m_S$,
\end{enumerate}
il existe un unique morphisme de $R$-algèbres $f : A \to S$
tel que $\mathfrak q = f^{-1}(\mathfrak m_S)$ et
$f \bmod \mathfrak q = \tau$.
\end{lemma}

\begin{proof}
Écrivons $A = \colim A_i$ comme colimite filtrante de $R$-algèbres \'etales.
Posons $\mathfrak q_i = A_i \cap \mathfrak q$. Nous obtenons
$f_i : A_i \to S$ en appliquant le lemme \ref{algebra-lemma-map-into-henselian}.
Posons $f = \colim f_i$.
\end{proof}

\begin{lemma}
\label{algebra-lemma-uniqueness-henselian}
Soit $R$ un anneau. Étant donné un diagramme commutatif de morphismes d'anneaux
$$
\xymatrix{
S \ar[r] & K \\
R \ar[u] \ar[r] & S' \ar[u]
}
$$
où $S$, $S'$ sont locaux henséliens, où $S$, $S'$ sont des colimites filtrantes
de $R$-algèbres \'etales, où $K$ est un corps, et où les flèches $S \to K$ et
$S' \to K$ identifient $K$ au corps résiduel de $S$ et de $S'$,
il existe un unique isomorphisme de $R$-algèbres $S \to S'$
compatible avec les morphismes vers $K$.
\end{lemma}

\begin{proof}
Cela résulte immédiatement du lemme \ref{algebra-lemma-map-into-henselian-colimit}.
\end{proof}

\noindent
Le lemme suivant ne concerne pas, à proprement parler, les colimites de
morphismes d'anneaux \'etales.

\begin{lemma}
\label{algebra-lemma-colimit-henselian}
Une colimite filtrante d'anneaux locaux (strictement) henséliens suivant des
morphismes locaux est (strictement) hensélienne.
\end{lemma}

\begin{proof}
Le lemme \ref{categories-lemma-directed-category-system} de Catégories
dit qu'il s'agit en réalité simplement d'une colimite d'anneaux
locaux (strictement) henséliens indexée par un ensemble ordonné filtrant.
Soit $(R_i, \varphi_{ii'})$ un tel système, avec chaque $\varphi_{ii'}$
local. Alors $R = \colim_i R_i$ est local, et
son corps résiduel $\kappa$ est $\colim \kappa_i$
(argument omis). Il est facile de voir que $\colim \kappa_i$
est séparablement clos si chaque $\kappa_i$ l'est ;
il suffit donc de prouver que $R$ est hensélien si chaque $R_i$ est hensélien.
Supposons que $f \in R[T]$ soit unitaire et que $a_0 \in \kappa$ soit
une racine simple de $\overline{f}$. Alors, pour un $i$ assez grand,
il existe un $f_i \in R_i[T]$ s'envoyant sur $f$ et un
$a_{0, i} \in \kappa_i$ s'envoyant sur $a_0$. Puisque
$\overline{f_i}(a_{0, i}) \in \kappa_i$,
resp.\ $\overline{f_i'}(a_{0, i}) \in \kappa_i$, s'envoie sur
$0 = \overline{f}(a_0) \in \kappa$,
resp.\ $0 \not = \overline{f'}(a_0) \in \kappa$,
nous concluons que $a_{0, i}$ est une racine simple de $\overline{f_i}$.
Comme $R_i$ est hensélien, nous pouvons trouver $a_i \in R_i$ tel que
$f_i(a_i) = 0$ et $a_{0, i} = \overline{a_i}$.
Alors l'image $a \in R$ de $a_i$ est la solution recherchée.
Ainsi, $R$ est hensélien.
\end{proof}



\section{Hensélisation et hensélisation stricte}
\label{algebra-section-henselization}

\noindent
Dans cette section, nous construisons l'hensélisation. Nous encourageons le lecteur
à garder à l'esprit l'unicité déjà démontrée dans le
lemme \ref{algebra-lemma-uniqueness-henselian}
et le comportement fonctoriel signalé dans le
lemme \ref{algebra-lemma-map-into-henselian-colimit}
pendant la lecture de ces résultats.

\begin{lemma}
\label{algebra-lemma-henselization}
Soit $(R, \mathfrak m, \kappa)$ un anneau local. Il existe un
morphisme local d'anneaux $R \to R^h$ ayant les propriétés suivantes :
\begin{enumerate}
\item $R^h$ est hensélien,
\item $R^h$ est une colimite filtrante de $R$-algèbres \'etales,
\item $\mathfrak m R^h$ est l'idéal
maximal de $R^h$, et
\item $\kappa = R^h/\mathfrak m R^h$.
\end{enumerate}
\end{lemma}

\begin{proof}
Considérons la catégorie des paires $(S, \mathfrak q)$ où $R \to S$ est un
morphisme d'anneaux \'etale et où $\mathfrak q$ est un idéal premier de $S$
au-dessus de $\mathfrak m$ avec $\kappa = \kappa(\mathfrak q)$. Un morphisme de paires
$(S, \mathfrak q) \to (S', \mathfrak q')$ est donné par un morphisme de
$R$-algèbres $\varphi : S \to S'$ tel que
$\varphi^{-1}(\mathfrak q') = \mathfrak q$.
Posons
$$
R^h = \colim_{(S, \mathfrak q)} S.
$$
Montrons que la catégorie des paires est filtrante, voir
Catégories, définition \ref{categories-definition-directed}.
La catégorie contient la paire $(R, \mathfrak m)$ et n'est donc pas vide,
ce qui prouve la partie (1) de
Catégories, définition \ref{categories-definition-directed}.
Pour toute paire $(S, \mathfrak q)$, l'idéal premier $\mathfrak q$
est maximal de corps résiduel $\kappa$, puisque la composée
$\kappa \to S/\mathfrak q \to \kappa(\mathfrak q)$ est un isomorphisme.
Supposons que $(S, \mathfrak q)$ et $(S', \mathfrak q')$ soient deux objets. Posons
$S'' = S \otimes_R S'$ et
$\mathfrak q'' = \mathfrak qS'' + \mathfrak q'S''$.
Alors $S''/\mathfrak q'' = S/\mathfrak q \otimes_R S'/\mathfrak q' = \kappa$
d'après ce qui précède. De plus, $R \to S''$ est \'etale d'après le
lemme \ref{algebra-lemma-etale}.
Cela prouve la partie (2) de
Catégories, définition \ref{categories-definition-directed}.
Supposons ensuite que
$\varphi, \psi : (S, \mathfrak q) \to (S', \mathfrak q')$
soient deux morphismes de paires. Alors $\varphi$, $\psi$, et
$S' \otimes_R S' \to S'$ sont des morphismes d'anneaux \'etales d'après le
lemme \ref{algebra-lemma-map-between-etale}.
Considérons
$$
S'' = (S' \otimes_{\varphi, S, \psi} S')
\otimes_{S' \otimes_R S'} S'
$$
avec l'idéal premier
$$
\mathfrak q'' =
(\mathfrak q' \otimes S' + S' \otimes \mathfrak q') \otimes S'
+
(S' \otimes_{\varphi, S, \psi} S') \otimes \mathfrak q'
$$
En raisonnant comme ci-dessus (le changement de base d'un morphisme \'etale est
\'etale, et la composée de morphismes \'etales est \'etale), nous voyons que
$S''$ est \'etale sur $R$. De plus, le morphisme canonique
$S' \to S''$ (utilisant par exemple le facteur le plus à droite)
égalise $\varphi$ et $\psi$. Cela prouve la partie (3) de
Catégories, définition \ref{categories-definition-directed}.
Nous concluons donc que $R^h$ est constitué de triplets
$(S, \mathfrak q, f)$ avec $f \in S$, et que deux tels triplets
$(S, \mathfrak q, f)$, $(S', \mathfrak q', f')$
définissent le même élément de $R^h$ si et seulement s'il existe
une paire $(S'', \mathfrak q'')$ et des morphismes de paires
$\varphi : (S, \mathfrak q) \to (S'', \mathfrak q'')$
et
$\varphi' : (S', \mathfrak q') \to (S'', \mathfrak q'')$
tels que $\varphi(f) = \varphi'(f')$.

\medskip\noindent
Supposons que $x \in R^h$. Représentons $x$ par un triplet
$(S, \mathfrak q, f)$.
Soient $\mathfrak q_1, \ldots, \mathfrak q_r$ les autres idéaux premiers de $S$
au-dessus de $\mathfrak m$. Alors $\mathfrak q \not \subset \mathfrak q_i$,
puisque nous avons vu ci-dessus que $\mathfrak q$ est maximal.
Ainsi, comme $\mathfrak q$ est un idéal premier,
nous pouvons trouver un $g \in S$, $g \not \in \mathfrak q$, tel que
$g \in \mathfrak q_i$ pour $i = 1, \ldots, r$. Considérons le morphisme de
paires $(S, \mathfrak q) \to (S_g, \mathfrak qS_g)$.
De cette manière, nous voyons que nous pouvons toujours supposer que $x$
est donné par un triplet $(S, \mathfrak q, f)$ où
$\mathfrak q$ est le seul idéal premier de $S$ au-dessus de $\mathfrak m$,
c'est-à-dire $\sqrt{\mathfrak mS} = \mathfrak q$. Mais, puisque
$R \to S$ est \'etale, nous avons
$\mathfrak mS_{\mathfrak q} = \mathfrak qS_{\mathfrak q}$, voir le
lemme \ref{algebra-lemma-etale-at-prime}.
Nous obtenons donc en fait $\mathfrak mS = \mathfrak q$.

\medskip\noindent
Supposons que $x \not \in \mathfrak mR^h$.
Représentons $x$ par un triplet $(S, \mathfrak q, f)$ avec
$\mathfrak mS = \mathfrak q$.
Alors $f \not \in \mathfrak mS$, c'est-à-dire $f \not \in \mathfrak q$.
Ainsi, $(S, \mathfrak q) \to (S_f, \mathfrak qS_f)$ est un morphisme
de paires tel que l'image de $f$ devienne inversible.
Par conséquent, $x$ est inversible, d'inverse représenté par le triplet
$(S_f, \mathfrak qS_f, 1/f)$. Nous concluons que $R^h$ est un anneau
local d'idéal maximal $\mathfrak mR^h$. Le corps résiduel est
$\kappa$, puisque nous pouvons définir $R^h/\mathfrak mR^h \to \kappa$
en envoyant un triplet $(S, \mathfrak q, f)$ sur la classe
résiduelle de $f$ modulo $\mathfrak q$.

\medskip\noindent
Il nous reste à montrer que $R^h$ est hensélien.
Supposons donc que $P \in R^h[T]$ soit un polynôme
unitaire et que $a_0 \in \kappa$ soit une racine simple de
la réduction $\overline{P} \in \kappa[T]$.
Nous pouvons alors trouver une paire $(S, \mathfrak q)$ telle que
$P$ soit l'image d'un polynôme unitaire $Q \in S[T]$.
Posons $S' = S[T]/(Q)$ et soit $\mathfrak q' \subset S'$
l'idéal maximal $\mathfrak q' = \mathfrak qS' + (T - a')S'$,
où $a' \in S$ est un élément quelconque relevant $a_0$.
Par construction, $S \to S'$ est \'etale en $\mathfrak q'$
et $\kappa = \kappa(\mathfrak q')$.
Choisissons $g \in S'$, $g \not \in \mathfrak q'$, tel que
$S'' = S'_g$ soit \'etale sur $S$. Alors
$(S, \mathfrak q) \to (S'', \mathfrak q'S'')$ est un morphisme
de paires. Le triplet $(S'', \mathfrak q'S'', \text{classe de }T)$
détermine maintenant un élément $a \in R^h$ ayant les propriétés
$P(a) = 0$ et $\overline{a} = a_0$, comme désiré.
\end{proof}

\begin{lemma}
\label{algebra-lemma-strict-henselization}
Soit $(R, \mathfrak m, \kappa)$ un anneau local.
Soit $\kappa \subset \kappa^{sep}$ une clôture algébrique séparable.
Il existe un diagramme commutatif
$$
\xymatrix{
\kappa \ar[r] & \kappa \ar[r] & \kappa^{sep} \\
R \ar[r] \ar[u] & R^h \ar[r] \ar[u] & R^{sh} \ar[u]
}
$$
ayant les propriétés suivantes :
\begin{enumerate}
\item le morphisme $R^h \to R^{sh}$ est local,
\item $R^{sh}$ est strictement hensélien,
\item $R^{sh}$ est une colimite filtrante de $R$-algèbres \'etales,
\item $\mathfrak m R^{sh}$ est l'idéal
maximal de $R^{sh}$, et
\item $\kappa^{sep} = R^{sh}/\mathfrak m R^{sh}$.
\end{enumerate}
\end{lemma}

\begin{proof}
Cela se démontre exactement comme le
lemme \ref{algebra-lemma-henselization}.
La seule différence est qu'au lieu de paires, on utilise des triplets
$(S, \mathfrak q, \alpha)$, où $R \to S$ est \'etale,
$\mathfrak q$ est un idéal premier de $S$ au-dessus de $\mathfrak m$, et
$\alpha : \kappa(\mathfrak q) \to \kappa^{sep}$ est un plongement
d'extensions de $\kappa$.
\end{proof}

\begin{definition}
\label{algebra-definition-henselization}
Soit $(R, \mathfrak m, \kappa)$ un anneau local.
\begin{enumerate}
\item Le morphisme local d'anneaux $R \to R^h$ construit dans le
lemme \ref{algebra-lemma-henselization}
est appelé l'{\it hensélisation} de $R$.
\item Étant donnée une clôture algébrique séparable
$\kappa \subset \kappa^{sep}$, le morphisme local d'anneaux
$R \to R^{sh}$ construit dans le
lemme \ref{algebra-lemma-strict-henselization}
est appelé l'{\it hensélisation stricte de $R$ relativement à
$\kappa \subset \kappa^{sep}$}.
\item Un morphisme local d'anneaux $R \to R^{sh}$ est appelé une
{\it hensélisation stricte} de $R$ s'il est isomorphe à l'un des
morphismes locaux d'anneaux construits dans le
lemme \ref{algebra-lemma-strict-henselization}.
\end{enumerate}
\end{definition}

\noindent
Les morphismes $R \to R^h \to R^{sh}$ sont des morphismes plats d'anneaux locaux.
D'après le lemme \ref{algebra-lemma-uniqueness-henselian}, les $R$-algèbres
$R^h$ et $R^{sh}$ sont bien définies à isomorphisme unique près par les conditions
qu'elles soient locales henséliennes, colimites filtrantes de $R$-algèbres \'etales,
de corps résiduels respectifs $\kappa$ et $\kappa^{sep}$.
Dans le reste de cette section, nous examinons principalement la fonctorialité des
hensélisations (strictes).
Nous examinerons des résultats plus subtils concernant
la relation entre $R$ et son hensélisation dans
Plus sur l'algèbre, section \ref{more-algebra-section-permanence-henselization}.

\begin{remark}
\label{algebra-remark-construct-sh-from-h}
Nous pouvons aussi construire $R^{sh}$ à partir de $R^h$. En effet, pour toute
sous-extension finie séparable $\kappa^{sep}/\kappa'/\kappa$,
il existe une unique extension locale finie étale d'anneaux
$R^h \subset R^h(\kappa')$ (à isomorphisme unique près)
dont l'extension de corps résiduels reproduit l'extension donnée, voir le
lemme \ref{algebra-lemma-henselian-cat-finite-etale}.
Nous pouvons donc poser
$$
R^{sh} =
\bigcup\nolimits_{\kappa \subset \kappa' \subset \kappa^{sep}}
R^h(\kappa')
$$
Les flèches de ce système, compatibles avec les flèches au niveau
des corps résiduels, existent d'après le
lemme \ref{algebra-lemma-henselian-cat-finite-etale}.
Cela produit un anneau local hensélien d'après le
lemme \ref{algebra-lemma-colimit-henselian},
puisque chacun des anneaux
$R^h(\kappa')$ est hensélien d'après le
lemme \ref{algebra-lemma-finite-over-henselian}.
Par construction, l'extension de corps résiduels induite par
$R^h \to R^{sh}$ est l'extension de corps $\kappa^{sep}/\kappa$.
Ainsi, $R^{sh}$ ainsi construit est strictement hensélien.
D'après le lemme \ref{algebra-lemma-composition-colimit-etale}, la $R$-algèbre
$R^{sh}$ est une colimite d'algèbres \'etales sur $R$. Ainsi, l'unicité du
lemme \ref{algebra-lemma-uniqueness-henselian} montre que $R^{sh}$
est l'hensélisation stricte.
\end{remark}

\begin{lemma}
\label{algebra-lemma-henselian-functorial-prepare}
Soit $R \to S$ un morphisme local d'anneaux locaux.
Soit $S \to S^h$ l'hensélisation.
Soit $R \to A$ un morphisme d'anneaux \'etale, et soit $\mathfrak q$
un idéal premier de $A$ au-dessus de $\mathfrak m_R$
tel que $R/\mathfrak m_R \cong \kappa(\mathfrak q)$.
Alors il existe un unique morphisme d'anneaux
$f : A \to S^h$ s'insérant dans le diagramme commutatif
$$
\xymatrix{
A \ar[r]_f & S^h \\
R \ar[u] \ar[r] & S \ar[u]
}
$$
tel que $f^{-1}(\mathfrak m_{S^h}) = \mathfrak q$.
\end{lemma}

\begin{proof}
C'est un cas particulier du lemme \ref{algebra-lemma-map-into-henselian}.
\end{proof}

\begin{lemma}
\label{algebra-lemma-henselian-functorial}
Soit $R \to S$ un morphisme local d'anneaux locaux.
Soient $R \to R^h$ et $S \to S^h$ les hensélisations.
Il existe un unique morphisme local d'anneaux $R^h \to S^h$ s'insérant
dans le diagramme commutatif
$$
\xymatrix{
R^h \ar[r]_f & S^h \\
R \ar[u] \ar[r] & S \ar[u]
}
$$
\end{lemma}

\begin{proof}
Cela découle immédiatement du lemme \ref{algebra-lemma-map-into-henselian-colimit}.
\end{proof}

\noindent
Voici une construction légèrement différente de l'hensélisation.

\begin{lemma}
\label{algebra-lemma-henselization-different}
Soit $R$ un anneau.
Soit $\mathfrak p \subset R$ un idéal premier.
Considérons la catégorie des paires $(S, \mathfrak q)$ où
$R \to S$ est \'etale et $\mathfrak q$ est un idéal premier au-dessus de $\mathfrak p$
tel que $\kappa(\mathfrak p) = \kappa(\mathfrak q)$.
Cette catégorie est filtrante et
$$
(R_{\mathfrak p})^h = \colim_{(S, \mathfrak q)} S
= \colim_{(S, \mathfrak q)} S_{\mathfrak q}
$$
canoniquement.
\end{lemma}

\begin{proof}
Un morphisme de paires $(S, \mathfrak q) \to (S', \mathfrak q')$
est donné par un morphisme de $R$-algèbres $\varphi : S \to S'$ tel que
$\varphi^{-1}(\mathfrak q') = \mathfrak q$.
Montrons que la catégorie des paires est filtrante, voir
Catégories, définition \ref{categories-definition-directed}.
La catégorie contient la paire $(R, \mathfrak p)$ et n'est donc pas vide,
ce qui prouve la partie (1) de
Catégories, définition \ref{categories-definition-directed}.
Supposons que $(S, \mathfrak q)$ et $(S', \mathfrak q')$ soient deux paires.
Notons que $\mathfrak q$, resp.\ $\mathfrak q'$ correspondent à des idéaux premiers
des anneaux fibres $S \otimes \kappa(\mathfrak p)$,
resp.\ $S' \otimes \kappa(\mathfrak p)$, de corps résiduels
$\kappa(\mathfrak p)$ ; ils correspondent donc à des idéaux maximaux de
$S \otimes \kappa(\mathfrak p)$, resp.\ $S' \otimes \kappa(\mathfrak p)$.
Posons $S'' = S \otimes_R S'$. D'après ce qui précède, il existe un unique
idéal premier $\mathfrak q'' \subset S''$ au-dessus de $\mathfrak q$ et de
$\mathfrak q'$ dont le corps résiduel est $\kappa(\mathfrak p)$.
Le morphisme d'anneaux $R \to S''$ est \'etale d'après le
lemme \ref{algebra-lemma-etale}.
Cela prouve la partie (2) de
Catégories, définition \ref{categories-definition-directed}.
Supposons ensuite que
$\varphi, \psi : (S, \mathfrak q) \to (S', \mathfrak q')$
soient deux morphismes de paires. Alors $\varphi$, $\psi$, et
$S' \otimes_R S' \to S'$ sont des morphismes d'anneaux \'etales d'après le
lemme \ref{algebra-lemma-map-between-etale}. Considérons
$$
S'' = (S' \otimes_{\varphi, S, \psi} S')
\otimes_{S' \otimes_R S'} S'
$$
En raisonnant comme ci-dessus (le changement de base d'un morphisme \'etale est
\'etale, et la composée de morphismes \'etales est \'etale), nous voyons que $S''$ est \'etale sur $R$. L'anneau
fibre de $S''$ au-dessus de $\mathfrak p$ est
$$
F'' = (F' \otimes_{\varphi, F, \psi} F')
\otimes_{F' \otimes_{\kappa(\mathfrak p)} F'} F'
$$
où $F', F$ sont les anneaux fibres de $S'$ et $S$. Puisque $\varphi$ et
$\psi$ sont des morphismes de paires, le morphisme $F' \to \kappa(\mathfrak p)$
correspondant à $\mathfrak p'$ se prolonge en un morphisme $F'' \to \kappa(\mathfrak p)$
et correspond à son tour à un idéal premier $\mathfrak q'' \subset S''$
dont le corps résiduel est $\kappa(\mathfrak p)$.
Le morphisme canonique $S' \to S''$ (utilisant par exemple le facteur le plus à droite)
est un morphisme de paires $(S', \mathfrak q') \to (S'', \mathfrak q'')$
qui égalise $\varphi$ et $\psi$. Cela prouve la partie (3) de
Catégories, définition \ref{categories-definition-directed}.
Nous concluons donc que la catégorie est filtrante.

\medskip\noindent
Rappelons que, dans la preuve du
lemme \ref{algebra-lemma-henselization},
nous avons construit $(R_{\mathfrak p})^h$ comme la colimite correspondante,
mais en partant de $R_{\mathfrak p}$ et de son idéal maximal
$\mathfrak pR_{\mathfrak p}$. Or, étant donnée une paire $(S, \mathfrak q)$
pour $(R, \mathfrak p)$, nous obtenons une paire
$(S_{\mathfrak p}, \mathfrak qS_{\mathfrak p})$ pour
$(R_{\mathfrak p}, \mathfrak pR_{\mathfrak p})$.
De plus, dans cette situation,
$$
S_{\mathfrak p} = \colim_{f \in R, f \not \in \mathfrak p} S_f.
$$
Ainsi, pour montrer les égalités
du lemme, il suffit de montrer que toute paire $(S_{loc}, \mathfrak q_{loc})$
pour $(R_{\mathfrak p}, \mathfrak pR_{\mathfrak p})$ est de la forme
$(S_{\mathfrak p}, \mathfrak qS_{\mathfrak p})$ pour une certaine paire
$(S, \mathfrak q)$ au-dessus de $(R, \mathfrak p)$ (certains détails sont omis).
Cela découle du
lemme \ref{algebra-lemma-etale}.
\end{proof}

\begin{lemma}
\label{algebra-lemma-henselian-functorial-improve}
Soit $R \to S$ un morphisme d'anneaux. Soit $\mathfrak q \subset S$ un idéal premier au-dessus
de $\mathfrak p \subset R$. Soient $R \to R^h$ et $S \to S^h$ les
hensélisations de $R_\mathfrak p$ et $S_\mathfrak q$. Le morphisme local d'anneaux
$R^h \to S^h$ du lemme \ref{algebra-lemma-henselian-functorial} identifie $S^h$
à l'hensélisation de $R^h \otimes_R S$ en l'unique idéal premier
au-dessus de $\mathfrak m^h$ et de $\mathfrak q$.
\end{lemma}

\begin{proof}
Le lemme \ref{algebra-lemma-henselization-different} montre que $R^h$, resp.\ $S^h$
sont des colimites filtrantes de $R$-algèbres, resp.\ de $S$-algèbres, \'etales.
Ainsi, $R^h \otimes_R S$ est une colimite filtrante de
$S$-algèbres \'etales $A_i$ (lemme \ref{algebra-lemma-etale}). D'après le
lemme \ref{algebra-lemma-colimits-of-etale}, $S^h$ est une
colimite filtrante de $R^h \otimes_R S$-algèbres \'etales.
Comme, de plus, $S^h$ est un anneau local hensélien de corps résiduel
égal à $\kappa(\mathfrak q)$, l'assertion découle du résultat d'unicité du
lemme \ref{algebra-lemma-uniqueness-henselian}.
\end{proof}

\begin{lemma}
\label{algebra-lemma-strictly-henselian-functorial-prepare}
Soit $\varphi : R \to S$ un morphisme local d'anneaux locaux.
Soit $S/\mathfrak m_S \subset \kappa^{sep}$ une clôture algébrique séparable.
Soit $S \to S^{sh}$ l'hensélisation stricte de $S$
relativement à $S/\mathfrak m_S \subset \kappa^{sep}$.
Soit $R \to A$ un morphisme d'anneaux \'etale, et soit $\mathfrak q$
un idéal premier de $A$ au-dessus de $\mathfrak m_R$.
Étant donné un diagramme commutatif
$$
\xymatrix{
\kappa(\mathfrak q) \ar[r]_{\phi} & \kappa^{sep} \\
R/\mathfrak m_R \ar[r]^{\varphi} \ar[u] & S/\mathfrak m_S \ar[u]
}
$$
il existe un unique morphisme d'anneaux
$f : A \to S^{sh}$ s'insérant dans le diagramme commutatif
$$
\xymatrix{
A \ar[r]_f & S^{sh} \\
R \ar[u] \ar[r]^{\varphi} & S \ar[u]
}
$$
tel que $f^{-1}(\mathfrak m_{S^h}) = \mathfrak q$ et que le
morphisme induit $\kappa(\mathfrak q) \to \kappa^{sep}$ soit celui qui est donné.
\end{lemma}

\begin{proof}
C'est un cas particulier du lemme \ref{algebra-lemma-map-into-henselian}.
\end{proof}

\begin{lemma}
\label{algebra-lemma-strictly-henselian-functorial}
Soit $R \to S$ un morphisme local d'anneaux locaux.
Choisissons des clôtures algébriques séparables
$R/\mathfrak m_R \subset \kappa_1^{sep}$
et
$S/\mathfrak m_S \subset \kappa_2^{sep}$.
Soient $R \to R^{sh}$ et $S \to S^{sh}$ les hensélisations
strictes correspondantes. Étant donné un diagramme commutatif
$$
\xymatrix{
\kappa_1^{sep} \ar[r]_{\phi} & \kappa_2^{sep} \\
R/\mathfrak m_R \ar[r]^{\varphi} \ar[u] & S/\mathfrak m_S \ar[u]
}
$$
il existe un unique morphisme local d'anneaux $R^{sh} \to S^{sh}$ s'insérant
dans le diagramme commutatif
$$
\xymatrix{
R^{sh} \ar[r]_f & S^{sh} \\
R \ar[u] \ar[r] & S \ar[u]
}
$$
et induisant $\phi$ sur les corps résiduels de
$R^{sh}$ et $S^{sh}$.
\end{lemma}

\begin{proof}
Cela découle immédiatement du lemme \ref{algebra-lemma-map-into-henselian-colimit}.
\end{proof}

\begin{lemma}
\label{algebra-lemma-strict-henselization-different}
Soit $R$ un anneau.
Soit $\mathfrak p \subset R$ un idéal premier.
Soit $\kappa(\mathfrak p) \subset \kappa^{sep}$ une
clôture algébrique séparable.
Considérons la catégorie des triplets $(S, \mathfrak q, \phi)$
où $R \to S$ est \'etale, $\mathfrak q$ est un idéal premier au-dessus de $\mathfrak p$,
et $\phi : \kappa(\mathfrak q) \to \kappa^{sep}$ est un morphisme de
$\kappa(\mathfrak p)$-algèbres. Cette catégorie est filtrante et
$$
(R_{\mathfrak p})^{sh} =
\colim_{(S, \mathfrak q, \phi)} S =
\colim_{(S, \mathfrak q, \phi)} S_{\mathfrak q}
$$
canoniquement.
\end{lemma}

\begin{proof}
Un morphisme de triplets $(S, \mathfrak q, \phi) \to (S', \mathfrak q', \phi')$
est donné par un morphisme de $R$-algèbres $\varphi : S \to S'$ tel que
$\varphi^{-1}(\mathfrak q') = \mathfrak q$ et que
$\phi' \circ \varphi = \phi$.
Montrons que la catégorie des triplets est filtrante, voir
Catégories, définition \ref{categories-definition-directed}.
La catégorie contient le triplet
$(R, \mathfrak p, \kappa(\mathfrak p) \subset \kappa^{sep})$
et n'est donc pas vide, ce qui prouve la partie (1) de
Catégories, définition \ref{categories-definition-directed}.
Supposons que $(S, \mathfrak q, \phi)$ et $(S', \mathfrak q', \phi')$
soient deux triplets.
Notons que $\mathfrak q$, resp.\ $\mathfrak q'$ correspondent à des idéaux premiers
des anneaux fibres $S \otimes \kappa(\mathfrak p)$,
resp.\ $S' \otimes \kappa(\mathfrak p)$, dont les corps résiduels sont
des extensions séparables finies de $\kappa(\mathfrak p)$, et que $\phi$, resp.\ $\phi'$,
correspondent à des morphismes vers $\kappa^{sep}$. Ces données correspondent donc à des
morphismes de $\kappa(\mathfrak p)$-algèbres
$$
\phi : S \otimes_R \kappa(\mathfrak p) \longrightarrow \kappa^{sep},
\quad
\phi' : S' \otimes_R \kappa(\mathfrak p) \longrightarrow \kappa^{sep}.
$$
Posons $S'' = S \otimes_R S'$. En combinant les morphismes ci-dessus, nous obtenons un unique
morphisme de $\kappa(\mathfrak p)$-algèbres
$$
\phi'' = \phi \otimes \phi' :
S'' \otimes_R \kappa(\mathfrak p)
\longrightarrow
\kappa^{sep}
$$
dont le noyau correspond à un idéal premier $\mathfrak q'' \subset S''$
au-dessus de $\mathfrak q$ et de $\mathfrak q'$, et dont le corps résiduel
s'envoie par $\phi''$ dans le corps composé de
$\phi(\kappa(\mathfrak q))$ et $\phi'(\kappa(\mathfrak q'))$ dans
$\kappa^{sep}$. Le morphisme d'anneaux $R \to S''$ est \'etale d'après le
lemme \ref{algebra-lemma-etale}.
Ainsi, $(S'', \mathfrak q'', \phi'')$ est un triplet dominant à la fois
$(S, \mathfrak q, \phi)$ et $(S', \mathfrak q', \phi')$.
Cela prouve la partie (2) de
Catégories, définition \ref{categories-definition-directed}.
Supposons ensuite que
$\varphi, \psi : (S, \mathfrak q, \phi) \to (S', \mathfrak q', \phi')$
soient deux morphismes de triplets. Alors $\varphi$, $\psi$, et
$S' \otimes_R S' \to S'$ sont des morphismes d'anneaux \'etales d'après le
lemme \ref{algebra-lemma-map-between-etale}.
Considérons
$$
S'' = (S' \otimes_{\varphi, S, \psi} S')
\otimes_{S' \otimes_R S'} S'
$$
En raisonnant comme ci-dessus (le changement de base d'un morphisme \'etale est
\'etale, et la composée de morphismes \'etales est \'etale), nous voyons que $S''$ est \'etale sur $R$. L'anneau
fibre de $S''$ au-dessus de $\mathfrak p$ est
$$
F'' = (F' \otimes_{\varphi, F, \psi} F')
\otimes_{F' \otimes_{\kappa(\mathfrak p)} F'} F'
$$
où $F', F$ sont les anneaux fibres de $S'$ et $S$. Puisque $\varphi$ et
$\psi$ sont des morphismes de triplets, le morphisme $\phi' : F' \to \kappa^{sep}$
se prolonge en un morphisme $\phi'' : F'' \to \kappa^{sep}$
qui correspond à son tour à un idéal premier $\mathfrak q'' \subset S''$.
Le morphisme canonique $S' \to S''$ (utilisant par exemple le facteur le plus à droite)
est un morphisme de triplets
$(S', \mathfrak q', \phi') \to (S'', \mathfrak q'', \phi'')$
qui égalise $\varphi$ et $\psi$. Cela prouve la partie (3) de
Catégories, définition \ref{categories-definition-directed}.
Nous concluons donc que la catégorie est filtrante.

\medskip\noindent
Il reste à montrer que la colimite $R_{colim}$ du système
est égale à l'hensélisation stricte
de $R_{\mathfrak p}$ relativement à $\kappa^{sep}$. Pour cela, notons que
le système des triplets $(S, \mathfrak q, \phi)$ contient comme sous-système
les paires $(S, \mathfrak q)$ du
lemme \ref{algebra-lemma-henselization-different}.
Ainsi, $R_{colim}$ contient $R_{\mathfrak p}^h$ d'après ce lemme.
De plus, il est clair que $R_{\mathfrak p}^h \subset R_{colim}$
est une colimite filtrante d'extensions d'anneaux \'etales.
Il s'ensuit que $R_{colim}$ est hensélien d'après les
lemmes \ref{algebra-lemma-finite-over-henselian} et
\ref{algebra-lemma-colimit-henselian}.
Enfin, d'après le
lemme \ref{algebra-lemma-make-etale-map-prescribed-residue-field},
le corps résiduel de $R_{colim}$ est égal à
$\kappa^{sep}$. Nous concluons donc que $R_{colim}$ est strictement hensélien
et est ainsi égal à l'hensélisation stricte de $R_{\mathfrak p}$, comme souhaité.
Certains détails sont omis.
\end{proof}

\begin{lemma}
\label{algebra-lemma-strictly-henselian-functorial-improve}
Soit $R \to S$ un morphisme d'anneaux. Soit $\mathfrak q \subset S$ un idéal premier au-dessus
de $\mathfrak p \subset R$. Choisissons des clôtures algébriques séparables
$\kappa(\mathfrak p) \subset \kappa_1^{sep}$
et
$\kappa(\mathfrak q) \subset \kappa_2^{sep}$.
Soient $R^{sh}$ et $S^{sh}$ les hensélisations strictes
correspondantes de $R_\mathfrak p$ et $S_\mathfrak q$.
Étant donné un diagramme commutatif
$$
\xymatrix{
\kappa_1^{sep} \ar[r]_{\phi} & \kappa_2^{sep} \\
\kappa(\mathfrak p) \ar[r]^{\varphi} \ar[u] & \kappa(\mathfrak q) \ar[u]
}
$$
le morphisme local d'anneaux $R^{sh} \to S^{sh}$ du
lemme \ref{algebra-lemma-strictly-henselian-functorial} identifie $S^{sh}$
à l'hensélisation stricte de $R^{sh} \otimes_R S$ en un idéal premier
au-dessus de $\mathfrak q$ et de l'idéal maximal
$\mathfrak m^{sh} \subset R^{sh}$.
\end{lemma}

\begin{proof}
La preuve est identique à celle du
lemme \ref{algebra-lemma-henselian-functorial-improve},
à ceci près qu'elle utilise le
lemme \ref{algebra-lemma-strict-henselization-different}
au lieu du
lemme \ref{algebra-lemma-henselization-different}.
\end{proof}

\begin{lemma}
\label{algebra-lemma-sh-from-h-map}
Soit $R \to S$ un morphisme d'anneaux. Soit $\mathfrak q \subset S$ un idéal premier
au-dessus de $\mathfrak p \subset R$ tel que
$\kappa(\mathfrak p) \to \kappa(\mathfrak q)$ soit un isomorphisme.
Choisissons une clôture algébrique séparable $\kappa^{sep}$ de
$\kappa(\mathfrak p) = \kappa(\mathfrak q)$.
Alors
$$
(S_\mathfrak q)^{sh} =
(S_\mathfrak q)^h \otimes_{(R_\mathfrak p)^h} (R_\mathfrak p)^{sh}
$$
\end{lemma}

\begin{proof}
Cela découle de l'autre construction de l'hensélisation stricte
d'un anneau local donnée dans la remarque \ref{algebra-remark-construct-sh-from-h}, ainsi que du
fait que les corps résiduels sont égaux. Certains détails sont omis.
\end{proof}














\section{Hensélisation et morphismes d'anneaux quasi-finis}
\label{algebra-section-henselization-quasi-finite}

\noindent
Dans cette section, nous démontrons quelques résultats concernant les morphismes fonctoriels
entre hensélisés (stricts) pour les morphismes d'anneaux quasi-finis.

\begin{lemma}
\label{algebra-lemma-quasi-finite-henselization}
Soit $R \to S$ un morphisme d'anneaux.
Soit $\mathfrak q$ un idéal premier de $S$ au-dessus de $\mathfrak p$ dans $R$.
Supposons que $R \to S$ soit quasi-fini en $\mathfrak q$.
Le diagramme commutatif
$$
\xymatrix{
R_{\mathfrak p}^h \ar[r] & S_{\mathfrak q}^h \\
R_{\mathfrak p} \ar[u] \ar[r] & S_{\mathfrak q} \ar[u]
}
$$
du
lemme \ref{algebra-lemma-henselian-functorial}
identifie $S_{\mathfrak q}^h$ à la localisation de
$R_{\mathfrak p}^h \otimes_{R_{\mathfrak p}} S_{\mathfrak q}$
en l'idéal premier engendré par $\mathfrak q$. De plus, le morphisme
d'anneaux $R_{\mathfrak p}^h \to S_{\mathfrak q}^h$ est fini.
\end{lemma}

\begin{proof}
Notons que $R_{\mathfrak p}^h \otimes_R S$ est quasi-fini sur
$R_{\mathfrak p}^h$ en l'idéal premier correspondant à $\mathfrak q$, voir le
lemme \ref{algebra-lemma-four-rings}. Ainsi, la localisation $S'$ de
$R_{\mathfrak p}^h \otimes_{R_{\mathfrak p}} S_{\mathfrak q}$ est hensélienne
et finie sur $R_{\mathfrak p}^h$, voir le
lemme \ref{algebra-lemma-finite-over-henselian}. En tant que localisation, $S'$ est une limite
inductive filtrante d'algèbres \'etales sur
$R_{\mathfrak p}^h \otimes_{R_{\mathfrak p}} S_{\mathfrak q}$.
Le lemme \ref{algebra-lemma-henselian-functorial-improve} montre que
$S_\mathfrak q^h$ est le hensélisé de
$R_{\mathfrak p}^h \otimes_{R_{\mathfrak p}} S_{\mathfrak q}$.
Ainsi, $S' = S_\mathfrak q^h$ d'après le résultat d'unicité du
lemme \ref{algebra-lemma-uniqueness-henselian}.
\end{proof}

\begin{lemma}
\label{algebra-lemma-quotient-henselization}
\begin{slogan}
L'hensélisation est compatible au passage au quotient.
\end{slogan}
Soit $R$ un anneau local et soit $R^h$ son hensélisé.
Soit $I \subset \mathfrak m_R$.
Alors $R^h/IR^h$ est le hensélisé de $R/I$.
\end{lemma}

\begin{proof}
C'est un cas particulier du
lemme \ref{algebra-lemma-quasi-finite-henselization}.
\end{proof}

\begin{lemma}
\label{algebra-lemma-quasi-finite-strict-henselization}
Soit $R \to S$ un morphisme d'anneaux.
Soit $\mathfrak q$ un idéal premier de $S$ au-dessus de $\mathfrak p$ dans $R$.
Supposons que $R \to S$ soit quasi-fini en $\mathfrak q$.
Soit $\kappa_2^{sep}/\kappa(\mathfrak q)$ une clôture séparable,
et notons $\kappa_1^{sep} \subset \kappa_2^{sep}$ le sous-corps
des éléments algébriques séparables sur $\kappa(\mathfrak p)$
(Corps, lemme \ref{fields-lemma-separable-first}).
Le diagramme commutatif
$$
\xymatrix{
R_{\mathfrak p}^{sh} \ar[r] & S_{\mathfrak q}^{sh} \\
R_{\mathfrak p} \ar[u] \ar[r] & S_{\mathfrak q} \ar[u]
}
$$
du lemme \ref{algebra-lemma-strictly-henselian-functorial}
identifie $S_{\mathfrak q}^{sh}$ à la localisation de
$R_{\mathfrak p}^{sh} \otimes_{R_{\mathfrak p}} S_{\mathfrak q}$
en l'idéal premier qui est le noyau du morphisme
$$
R_{\mathfrak p}^{sh} \otimes_{R_{\mathfrak p}} S_{\mathfrak q}
\longrightarrow
\kappa_1^{sep} \otimes_{\kappa(\mathfrak p)} \kappa(\mathfrak q)
\longrightarrow
\kappa_2^{sep}
$$
De plus, le morphisme d'anneaux $R_{\mathfrak p}^{sh} \to S_{\mathfrak q}^{sh}$
est un homomorphisme local et fini d'anneaux locaux dont l'extension de corps résiduels
est l'extension $\kappa_2^{sep}/\kappa_1^{sep}$, laquelle est à la fois finie et
purement inséparable.
\end{lemma}

\begin{proof}
Comme $R \to S$ est quasi-fini en $\mathfrak q$, l'extension
$\kappa(\mathfrak q)/\kappa(\mathfrak p)$ est finie, voir la
définition \ref{algebra-definition-quasi-finite} et le
lemme \ref{algebra-lemma-isolated-point-fibre}.
Ainsi, $\kappa_1^{sep}$ est une clôture séparable de
$\kappa(\mathfrak p)$ (détail mineur omis).
En particulier, le lemme \ref{algebra-lemma-strictly-henselian-functorial}
s'applique bien. Ensuite, le corps composé de $\kappa(\mathfrak q)$ et de
$\kappa_1^{sep}$ dans $\kappa_2^{sep}$ est séparablement
clos, et est donc égal à $\kappa_2^{sep}$. Nous concluons que
$\kappa_2^{sep}/\kappa_1^{sep}$ est finie. Par construction,
l'extension $\kappa_2^{sep}/\kappa_1^{sep}$ est purement inséparable.
Le morphisme d'anneaux $R_{\mathfrak p}^{sh} \to S_{\mathfrak q}^{sh}$ est
bien local et induit l'extension de corps résiduels
$\kappa_2^{sep}/\kappa_1^{sep}$, qui est bien finie et purement inséparable.

\medskip\noindent
Notons que $R_{\mathfrak p}^{sh} \otimes_R S$ est quasi-fini sur
$R_{\mathfrak p}^{sh}$ en l'idéal premier $\mathfrak q'$
donné dans l'énoncé du lemme, voir le
lemme \ref{algebra-lemma-four-rings}. Ainsi, la localisation $S'$ de
$R_{\mathfrak p}^{sh} \otimes_{R_{\mathfrak p}} S_{\mathfrak q}$
en $\mathfrak q'$ est hensélienne et finie sur $R_{\mathfrak p}^{sh}$, voir le
lemme \ref{algebra-lemma-finite-over-henselian}.
Notons que le corps résiduel de $S'$ est $\kappa_2^{sep}$,
car le morphisme $\kappa_1^{sep} \otimes_{\kappa(\mathfrak p)} \kappa(\mathfrak q)
\to \kappa_2^{sep}$ est surjectif d'après la discussion du paragraphe
précédent.
En outre, en tant que localisation, $S'$ est une limite inductive filtrante d'algèbres \'etales sur
$R_{\mathfrak p}^{sh} \otimes_{R_{\mathfrak p}} S_{\mathfrak q}$.
Le lemme \ref{algebra-lemma-strictly-henselian-functorial-improve}
montre que $S_{\mathfrak q}^{sh}$ est le hensélisé strict de
$R_{\mathfrak p}^{sh} \otimes_{R_{\mathfrak p}} S_{\mathfrak q}$
en $\mathfrak q'$. Ainsi, $S' = S_\mathfrak q^{sh}$ d'après le résultat d'unicité du
lemme \ref{algebra-lemma-uniqueness-henselian}.
\end{proof}

\begin{lemma}
\label{algebra-lemma-quotient-strict-henselization}
Soit $R$ un anneau local et soit $R^{sh}$ son hensélisé strict.
Soit $I \subset \mathfrak m_R$.
Alors $R^{sh}/IR^{sh}$ est un hensélisé strict de $R/I$.
\end{lemma}

\begin{proof}
C'est un cas particulier du
lemme \ref{algebra-lemma-quasi-finite-strict-henselization}.
\end{proof}

\begin{lemma}
\label{algebra-lemma-local-tensor-with-integral}
Soient $A \to B$ et $A \to C$ des homomorphismes locaux d'anneaux locaux.
Si $A \to C$ est entier et si l'une des extensions
$\kappa(\mathfrak m_C)/\kappa(\mathfrak m_A)$ ou
$\kappa(\mathfrak m_B)/\kappa(\mathfrak m_A)$ est purement
inséparable, alors $D = B \otimes_A C$ est un anneau local et
$B \to D$ et $C \to D$ sont locaux.
\end{lemma}

\begin{proof}
Tout idéal maximal de $D$ est au-dessus de l'idéal maximal de $B$ par le
théorème de montée pour le morphisme d'anneaux entier
$B \to D$ (lemme \ref{algebra-lemma-integral-going-up}).
Or $D/\mathfrak m_B D = \kappa(\mathfrak m_B) \otimes_A C =
\kappa(\mathfrak m_B) \otimes_{\kappa(\mathfrak m_A)} C/\mathfrak m_A C$.
Le spectre de $C/\mathfrak m_A C$ est réduit à un
seul point, à savoir $\mathfrak m_C$. Ainsi, le spectre de
$D/\mathfrak m_B D$ est le même que celui de
$\kappa(\mathfrak m_B) \otimes_{\kappa(\mathfrak m_A)} \kappa(\mathfrak m_C)$,
qui est réduit à un seul point puisque, par hypothèse, l'une des extensions
$\kappa(\mathfrak m_C)/\kappa(\mathfrak m_A)$ ou
$\kappa(\mathfrak m_B)/\kappa(\mathfrak m_A)$ est purement
inséparable. Cela prouve que $D$ est local et que les morphismes d'anneaux
$B \to D$ et $C \to D$ sont locaux.
\end{proof}

\begin{lemma}
\label{algebra-lemma-base-change-strict-henselization-quasi-finite}
Soient $A \to B$ et $A \to C$ des morphismes d'anneaux. Soit $\kappa$ un corps
séparablement clos et soit $B \otimes_A C \to \kappa$
un homomorphisme d'anneaux. Notons
$$
\xymatrix{
B^{sh} \ar[r] & (B \otimes_A C)^{sh} \\
A^{sh} \ar[u] \ar[r] & C^{sh} \ar[u]
}
$$
les morphismes correspondants entre hensélisés stricts (voir la démonstration). Si
\begin{enumerate}
\item $A \to B$ est quasi-fini en l'idéal premier
$\mathfrak p_B = \Ker(B \to \kappa)$, ou
\item $B$ est une limite inductive filtrante de $A$-algèbres quasi-finies, ou
\item $B_{\mathfrak p_B}$ est une limite inductive filtrante d'algèbres quasi-finies
sur $A_{\mathfrak p_A}$, ou
\item $B$ est entier sur $A$,
\end{enumerate}
alors $B^{sh} \otimes_{A^{sh}} C^{sh} \to (B \otimes_A C)^{sh}$
est un isomorphisme.
\end{lemma}

\begin{proof}
Posons $D = B \otimes_A C$. Notons
$\mathfrak p_A = \Ker(A \to \kappa)$ et de même
pour $\mathfrak p_B$, $\mathfrak p_C$ et $\mathfrak p_D$.
Notons $\kappa_A \subset \kappa$ la clôture séparable
de $\kappa(\mathfrak p_A)$ dans $\kappa$
et de même pour $\kappa_B$, $\kappa_C$ et $\kappa_D$.
Notons $A^{sh}$ le hensélisé strict de $A_{\mathfrak p_A}$
construit à l'aide de la clôture séparable
$\kappa_A/\kappa(\mathfrak p_A)$. Faisons de même pour
$B^{sh}$, $C^{sh}$ et $D^{sh}$. Nous obtenons le diagramme commutatif
du lemme par fonctorialité du
lemme \ref{algebra-lemma-strictly-henselian-functorial}.

\medskip\noindent
Considérons le morphisme
$$
c : B^{sh} \otimes_{A^{sh}} C^{sh} \to D^{sh} = (B \otimes_A C)^{sh}
$$
que fournit le diagramme commutatif. Si $A \to B$ est quasi-fini en
$\mathfrak p_B = \Ker(B \to \kappa)$, alors le morphisme d'anneaux $C \to D$ est
quasi-fini en $\mathfrak p_D$ d'après le lemme \ref{algebra-lemma-four-rings}. Ainsi, d'après le
lemme \ref{algebra-lemma-quasi-finite-strict-henselization}
(et le lemme \ref{algebra-lemma-base-change-integral}),
le morphisme d'anneaux $c$ est un homomorphisme d'algèbres finies sur $C^{sh}$ et
$$
B^{sh} = (B \otimes_A A^{sh})_{\mathfrak q}
\quad\text{et}\quad
D^{sh} = (D \otimes_C C^{sh})_{\mathfrak r} =
(B \otimes_A C^{sh})_{\mathfrak r}
$$
pour certains idéaux premiers $\mathfrak q$ et $\mathfrak r$. Comme
$$
B^{sh} \otimes_{A^{sh}} C^{sh} =
(B \otimes_A A^{sh})_{\mathfrak q} \otimes_{A^{sh}} C^{sh} =
\text{une localisation de }
B \otimes_A C^{sh}
$$
nous concluons que la source et le but de $c$ sont tous deux
des localisations de $B \otimes_A C^{sh}$ (de manière compatible avec le morphisme).
Il suffit donc de montrer que $B^{sh} \otimes_{A^{sh}} C^{sh}$
est local (détail mineur omis). Cela résulte du
lemme \ref{algebra-lemma-local-tensor-with-integral}
et du fait que $A^{sh} \to B^{sh}$ est fini et induit une extension
purement inséparable des corps résiduels, d'après le
lemme \ref{algebra-lemma-quasi-finite-strict-henselization} déjà utilisé.
Cela démontre le cas (1) du lemme.

\medskip\noindent
Dans le cas (2), écrivons $B = \colim B_i$ sous la forme d'une limite inductive filtrante de
$A$-algèbres quasi-finies. Nous obtenons alors
$D = \colim D_i$, où $D_i = B_i \otimes_A C$.
Observons que $B^{sh} = \colim B_i^{sh}$. En effet, l'anneau
$\colim B_i^{sh}$ est un anneau local strictement hensélien
d'après le lemme \ref{algebra-lemma-colimit-henselian}.
De plus, $\colim B_i^{sh}$ est une limite inductive filtrante de
$B$-algèbres \'etales d'après le lemme \ref{algebra-lemma-colimit-colimit-etale-better}.
Enfin, le corps résiduel de $\colim B_i^{sh}$ est une clôture
séparable de $\kappa(\mathfrak p_B)$ (détails omis).
Nous concluons donc que $B^{sh} = \colim B_i^{sh}$, voir la
discussion qui suit la définition \ref{algebra-definition-henselization}.
De même, nous avons $D^{sh} = \colim D_i^{sh}$.
On conclut alors en appliquant le cas (1), car
$$
D^{sh} = \colim D_i^{sh} = \colim B_i^{sh} \otimes_{A^{sh}} C^{sh} =
B^{sh} \otimes_{A^{sh}} C^{sh}
$$
puisque les limites inductives filtrantes commutent avec les produits tensoriels.

\medskip\noindent
Cas (3). Nous pouvons remplacer $A$, $B$ et $C$ par leurs localisations
en $\mathfrak p_A$, $\mathfrak p_B$ et $\mathfrak p_C$.
Ainsi, (3) résulte de (2).

\medskip\noindent
Puisqu'un morphisme d'anneaux entier est une limite inductive filtrante de morphismes d'anneaux
finis, nous voyons que (4) résulte également de (2).
\end{proof}






















\section{Critère de normalité de Serre}
\label{algebra-section-serre-criterion}

\noindent
Nous introduisons les propriétés suivantes pour les anneaux noethériens.

\begin{definition}
\label{algebra-definition-conditions}
Soit $R$ un anneau noethérien.
Soit $k \geq 0$ un entier.
\begin{enumerate}
\item On dit que $R$ possède la propriété {\it $(R_k)$} si, pour tout idéal premier $\mathfrak p$
de hauteur $\leq k$, l'anneau local $R_{\mathfrak p}$ est régulier.
On dit aussi que $R$ est {\it régulier en codimension $\leq k$}.
\item On dit que $R$ possède la propriété {\it $(S_k)$} si, pour tout idéal premier $\mathfrak p$,
l'anneau local $R_{\mathfrak p}$ est de profondeur au moins
$\min\{k, \dim(R_{\mathfrak p})\}$.
\item Soit $M$ un $R$-module de type fini. On dit que $M$ possède la propriété $(S_k)$
si, pour tout idéal premier $\mathfrak p$, le module
$M_{\mathfrak p}$ est de profondeur au moins
$\min\{k, \dim(\text{Supp}(M_{\mathfrak p}))\}$.
\end{enumerate}
\end{definition}

\noindent
Tout anneau noethérien possède la propriété $(S_0)$, de même que tout module de type fini
sur cet anneau. Notre convention selon laquelle la profondeur du module nul vaut $\infty$
(voir la section \ref{algebra-section-depth}) et la dimension de l'ensemble vide vaut
$-\infty$ (voir Topologie, section \ref{topology-section-krull-dimension})
garantit que le module nul possède la propriété $(S_k)$ pour tout $k$.

\begin{lemma}
\label{algebra-lemma-criterion-no-embedded-primes}
Soit $R$ un anneau noethérien.
Soit $M$ un $R$-module de type fini.
Les assertions suivantes sont équivalentes :
\begin{enumerate}
\item $M$ n'a aucun idéal premier associé immergé, et
\item $M$ possède la propriété $(S_1)$.
\end{enumerate}
\end{lemma}

\begin{proof}
Soit $\mathfrak p$ un idéal premier associé immergé de $M$.
Il existe alors un autre idéal premier associé $\mathfrak q$ de $M$
tel que $\mathfrak p \supset \mathfrak q$. Cela implique en particulier
que $\dim(\text{Supp}(M_{\mathfrak p})) \geq 1$ (puisque $\mathfrak q$
appartient lui aussi au support). D'autre part, $\mathfrak pR_{\mathfrak p}$
est associé à $M_{\mathfrak p}$
(lemme \ref{algebra-lemma-associated-primes-localize}), et donc
$\text{profondeur}(M_{\mathfrak p}) = 0$
(voir le lemme \ref{algebra-lemma-ideal-nonzerodivisor}).
Autrement dit, la propriété $(S_1)$ n'est pas vérifiée.
Réciproquement, si la propriété $(S_1)$ n'est pas vérifiée, il existe un idéal premier
$\mathfrak p$ tel que $\dim(\text{Supp}(M_{\mathfrak p})) \geq 1$
et $\text{profondeur}(M_{\mathfrak p}) = 0$. Puisque
$\text{profondeur}(M_{\mathfrak p}) = 0$, nous voyons que
$\mathfrak p \in \text{Ass}(M)$ d'après les deux lemmes
\ref{algebra-lemma-associated-primes-localize} et \ref{algebra-lemma-ideal-nonzerodivisor}.
Comme $\dim(\text{Supp}(M_{\mathfrak p})) \geq 1$, il existe
un idéal premier $\mathfrak q \in \text{Supp}(M)$ tel que
$\mathfrak q \subset \mathfrak p$, $\mathfrak q \not = \mathfrak p$.
Nous pouvons choisir un tel $\mathfrak q$ minimal dans
$\text{Supp}(M)$. Alors, d'après la
proposition \ref{algebra-proposition-minimal-primes-associated-primes},
nous avons $\mathfrak q \in \text{Ass}(M)$, et donc
$\mathfrak p$ est un idéal premier associé immergé.
\end{proof}

\begin{lemma}
\label{algebra-lemma-criterion-reduced}
\begin{slogan}
Réduit équivaut à R0 et S1.
\end{slogan}
Soit $R$ un anneau noethérien.
Les assertions suivantes sont équivalentes :
\begin{enumerate}
\item $R$ est réduit, et
\item $R$ vérifie les propriétés $(R_0)$ et $(S_1)$.
\end{enumerate}
\end{lemma}

\begin{proof}
Supposons que $R$ soit réduit. Alors $R_{\mathfrak p}$ est un corps pour
tout idéal premier minimal $\mathfrak p$ de $R$, d'après le
lemme \ref{algebra-lemma-minimal-prime-reduced-ring}. Ainsi, la propriété $(R_0)$ est vérifiée.
Soit $\mathfrak p$ un idéal premier de hauteur $\geq 1$. Alors $A = R_{\mathfrak p}$
est un anneau local réduit de dimension $\geq 1$. Son idéal maximal
$\mathfrak m$ n'est donc pas un idéal premier associé,
car cela signifierait qu'il existe un $x \in \mathfrak m$
d'annulateur $\mathfrak m$, si bien que $x^2 = 0$. La profondeur de
$A = R_{\mathfrak p}$ est donc au moins un, d'après le lemme \ref{algebra-lemma-ass-zero-divisors}.
Cela montre que la propriété $(S_1)$ est vérifiée.

\medskip\noindent
Réciproquement, supposons que $R$ vérifie les propriétés $(R_0)$ et $(S_1)$.
Si $\mathfrak p$ est un idéal premier minimal de $R$, alors
$R_{\mathfrak p}$ est un corps d'après $(R_0)$, et est donc réduit.
Si $\mathfrak p$ n'est pas minimal, alors nous voyons que $R_{\mathfrak p}$
est de profondeur $\geq 1$ d'après $(S_1)$, et nous en déduisons qu'il existe un élément
$t \in \mathfrak pR_{\mathfrak p}$ tel que
$R_{\mathfrak p} \to R_{\mathfrak p}[1/t]$ soit injectif.
Or $R_\mathfrak p[1/t]$ est contenu dans le produit
de ses localisations aux idéaux premiers, voir le
lemme \ref{algebra-lemma-characterize-zero-local}.
Il s'ensuit que $R_{\mathfrak p}$ est un sous-anneau d'un produit de
localisations de $R$ en des idéaux premiers satisfaisant $\mathfrak p \supset \mathfrak q$ et
$t \not \in \mathfrak q$. Comme ces idéaux premiers sont de hauteur plus petite,
nous concluons par récurrence sur la hauteur que $R$ est réduit.
\end{proof}

\begin{lemma}[Critère de normalité de Serre]
\label{algebra-lemma-criterion-normal}
\begin{reference}
\cite[IV, Theorem 5.8.6]{EGA}
\end{reference}
\begin{slogan}
Normal équivaut à R1 et S2.
\end{slogan}
Soit $R$ un anneau noethérien.
Les assertions suivantes sont équivalentes :
\begin{enumerate}
\item $R$ est un anneau normal, et
\item $R$ vérifie les propriétés $(R_1)$ et $(S_2)$.
\end{enumerate}
\end{lemma}

\begin{proof}
Démonstration de (1) $\Rightarrow$ (2). Supposons que $R$ soit normal, c'est-à-dire que tous
les localisés $R_{\mathfrak p}$ aux idéaux premiers soient des anneaux intègres normaux.
En particulier, nous voyons que $R$ vérifie $(R_0)$ et $(S_1)$ d'après le
lemme \ref{algebra-lemma-criterion-reduced}. Il suffit donc de montrer
qu'un anneau local noethérien intègre et normal $R$ de dimension $d$ est
de profondeur $\geq \min(2, d)$ et est régulier si $d = 1$. L'assertion
lorsque $d = 1$ résulte du lemme \ref{algebra-lemma-characterize-dvr}.

\medskip\noindent
Soit $R$ un anneau local noethérien intègre et normal, d'idéal maximal
$\mathfrak m$ et de dimension $d \geq 2$. Appliquons le
lemme \ref{algebra-lemma-hart-serre-loc-thm} à $R$.
Il est clair que $R$ ne relève ni du cas (1) ni du cas (2)
du lemme.
Considérons $R \to R'$ comme dans le cas (4) du lemme.
Puisque $R$ est intègre, nous avons $R \subset R'$. Comme $\mathfrak m$
n'est pas un idéal premier associé à $R'$, il existe un $x \in \mathfrak m$
qui est un non-diviseur de zéro dans $R'$. Alors $R_x = R'_x$, de sorte que
$R$ et $R'$ sont des anneaux intègres de même corps des fractions. Mais la
finitude de $R \subset R'$ implique que tout élément de $R'$ est entier
sur $R$ (lemme \ref{algebra-lemma-finite-is-integral}),
et nous concluons que $R = R'$, puisque $R$ est normal.
Cela signifie que le cas (4) ne se produit pas. Il reste donc la possibilité
(3), c'est-à-dire $\text{profondeur}(R) \geq 2$, comme souhaité.

\medskip\noindent
Démonstration de (2) $\Rightarrow$ (1). Supposons que $R$ vérifie $(R_1)$ et $(S_2)$.
Le lemme \ref{algebra-lemma-criterion-reduced} montre que $R$ est
réduit. Il suffit donc de montrer que, si $R$ est un anneau local réduit
noethérien de dimension $d$ vérifiant $(S_2)$ et $(R_1)$,
alors $R$ est un anneau intègre normal. Si $d = 0$, le résultat est clair.
Si $d = 1$, le résultat découle du lemme \ref{algebra-lemma-characterize-dvr}.

\medskip\noindent
Soit $R$ un anneau local noethérien réduit, d'idéal maximal
$\mathfrak m$ et de dimension $d \geq 2$, vérifiant $(R_1)$ et
$(S_2)$. D'après le lemme \ref{algebra-lemma-characterize-reduced-ring-normal},
il suffit de montrer que $R$ est intégralement clos dans son
anneau total des fractions $Q(R)$. Choisissons $x \in Q(R)$ entier
sur $R$. Alors $R' = R[x]$ est une extension finie de $R$
(lemme \ref{algebra-lemma-characterize-finite-in-terms-of-integral}).
Puisque $\dim(R_\mathfrak p) < d$ pour
tout idéal premier non maximal $\mathfrak p \subset R$,
nous avons $R_\mathfrak p = R'_\mathfrak p$ par récurrence.
Le support de $R'/R$ est donc $\{\mathfrak m\}$.
Il s'ensuit que $R'/R$ est annulé par une puissance de $\mathfrak m$
(lemme \ref{algebra-lemma-Noetherian-power-ideal-kills-module}).
Le lemme \ref{algebra-lemma-hart-serre-loc-thm} contredit alors
l'hypothèse selon laquelle la profondeur de $R$ est $\geq 2 = \min(2, d)$,
et la démonstration est achevée.
\end{proof}

\begin{lemma}
\label{algebra-lemma-regular-normal}
Un anneau régulier est normal.
\end{lemma}

\begin{proof}
Soit $R$ un anneau régulier. D'après le
lemme \ref{algebra-lemma-criterion-normal},
il suffit de montrer que $R$ vérifie $(R_1)$ et $(S_2)$.
Comme tout anneau local régulier est de Cohen-Macaulay, voir le
lemme \ref{algebra-lemma-regular-ring-CM},
il est clair que $R$ vérifie $(S_2)$.
La propriété $(R_1)$ est immédiate.
\end{proof}

\begin{lemma}
\label{algebra-lemma-normal-domain-intersection-localizations-height-1}
Soit $R$ un anneau noethérien intègre et normal de corps des fractions $K$. Alors
\begin{enumerate}
\item pour tout $a \in R$ non nul, le quotient $R/aR$ n'a aucun idéal premier associé immergé,
et tous ses idéaux premiers associés sont de hauteur $1$
\item
$$
R = \bigcap\nolimits_{\text{hauteur}(\mathfrak p) = 1} R_{\mathfrak p}
$$
\item Pour tout $x \in K$ non nul, le quotient $R/(R \cap xR)$
n'a aucun idéal premier associé immergé, et tous ses idéaux premiers associés sont de hauteur $1$.
\end{enumerate}
\end{lemma}

\begin{proof}
Le lemme \ref{algebra-lemma-criterion-normal} montre que $R$ vérifie $(S_2)$.
Ainsi, pour tout élément $a \in R$ non nul, le quotient $R/aR$ vérifie $(S_1)$
(utiliser par exemple le lemme \ref{algebra-lemma-depth-in-ses}).
Par conséquent, $R/aR$ n'a aucun idéal premier associé immergé
(lemme \ref{algebra-lemma-criterion-no-embedded-primes}).
Nous en déduisons que les idéaux premiers associés à $R/aR$ sont exactement
les idéaux premiers minimaux $\mathfrak p$ au-dessus de $(a)$, qui sont de hauteur $1$
puisque $a$ n'est pas nul (lemme \ref{algebra-lemma-minimal-over-1}). Cela démontre (1).

\medskip\noindent
Ainsi, étant donné $b \in R$, on a $b \in aR$ si et seulement si
$b \in aR_{\mathfrak p}$ pour tout idéal premier minimal $\mathfrak p$
au-dessus de $(a)$ (voir le lemme \ref{algebra-lemma-zero-at-ass-zero}).
Ces idéaux premiers sont tous de hauteur $1$, comme on l'a vu ci-dessus, donc
$b/a \in R$ si et seulement si $b/a \in R_{\mathfrak p}$ pour tous les idéaux premiers
de hauteur 1. L'assertion (2) s'ensuit.

\medskip\noindent
Pour (3), écrivons $x = a/b$. Soient $\mathfrak p_1, \ldots, \mathfrak p_r$
les idéaux premiers minimaux au-dessus de $(ab)$. Ils sont tous de hauteur 1 d'après ce qui précède.
Nous avons alors
$R \cap xR = \bigcap_{i = 1, \ldots, r} (R \cap xR_{\mathfrak p_i})$
d'après la partie (2) du lemme. Par conséquent, $R/(R \cap xR)$ est un sous-module de
$\bigoplus R/(R \cap xR_{\mathfrak p_i})$.
Comme $R_{\mathfrak p_i}$ est un anneau de valuation discrète (d'après la propriété $(R_1)$
pour l'anneau noethérien intègre et normal $R$, voir le lemme \ref{algebra-lemma-criterion-normal}),
nous avons $xR_{\mathfrak p_i} = \mathfrak p_i^{e_i}R_{\mathfrak p_i}$
pour un certain $e_i \in \mathbf{Z}$. La somme directe est donc égale
à $\bigoplus_{e_i > 0} R/\mathfrak p_i^{(e_i)}$, voir la
définition \ref{algebra-definition-symbolic-power}.
D'après le lemme \ref{algebra-lemma-symbolic-power-associated},
le seul idéal premier associé au module
$R/\mathfrak p^{(n)}$ est $\mathfrak p$. Ainsi, l'ensemble des idéaux premiers associés
à $R/(R \cap xR)$ est un sous-ensemble de $\{\mathfrak p_i\}$ et il n'existe
aucune relation d'inclusion entre eux. Cela démontre (3).
\end{proof}















\section{Lissité formelle des corps}
\label{algebra-section-p-bases}

\noindent
Dans cette section, nous montrons que les extensions de corps sont formellement lisses
si et seulement si elles sont séparables. Cependant, nous démontrons d'abord que les extensions de corps
de type fini sont algébriques séparables si et seulement si
elles sont formellement non ramifiées.

\begin{lemma}
\label{algebra-lemma-characterize-separable-algebraic-field-extensions}
Soit $K/k$ une extension de corps de type fini.
Les propriétés suivantes sont équivalentes :
\begin{enumerate}
\item $K$ est une extension finie séparable de $k$,
\item $\Omega_{K/k} = 0$,
\item $K$ est formellement non ramifié sur $k$,
\item $K$ est non ramifié sur $k$,
\item $K$ est formellement \'etale sur $k$,
\item $K$ est \'etale sur $k$.
\end{enumerate}
\end{lemma}

\begin{proof}
L'équivalence de (2) et (3) résulte du
Lemme \ref{algebra-lemma-characterize-formally-unramified}.
D'après le Lemme \ref{algebra-lemma-etale-over-field},
nous voyons que (1) est équivalent à (6).
La propriété (6) implique (5) et (4), qui impliquent toutes deux à leur tour (3)
(Lemmes \ref{algebra-lemma-formally-etale-etale}, \ref{algebra-lemma-unramified},
et \ref{algebra-lemma-formally-unramified-unramified}).
Il suffit donc de montrer que (2) implique (1).
Choisissons une $k$-sous-algèbre de type fini $A \subset K$
telle que $K$ soit le corps des fractions de l'anneau intègre $A$.
Posons $S = A \setminus \{0\}$.
Puisque $0 = \Omega_{K/k} = S^{-1}\Omega_{A/k}$
(Lemme \ref{algebra-lemma-differentials-localize})
et puisque $\Omega_{A/k}$ est de type fini
(Lemme \ref{algebra-lemma-differentials-finitely-generated}),
nous pouvons remplacer $A$ par une localisation $A_f$ pour nous ramener au cas
où $\Omega_{A/k} = 0$ (détails omis).
Alors $A$ est non ramifié sur $k$, donc
$K/k$ est une extension finie séparable, par exemple d'après le
Lemme \ref{algebra-lemma-unramified-at-prime} appliqué avec $\mathfrak q = (0)$.
\end{proof}

\begin{lemma}
\label{algebra-lemma-derivative-zero-pth-power}
Soit $k$ un corps parfait de caractéristique $p > 0$.
Soit $K/k$ une extension de corps.
Soit $a \in K$. Alors $\text{d}a = 0$ dans $\Omega_{K/k}$
si et seulement si $a$ est une puissance $p$-ième.
\end{lemma}

\begin{proof}
D'après le Lemme \ref{algebra-lemma-colimit-differentials}, nous voyons qu'il existe un sous-corps
$k \subset L \subset K$ tel que $L/k$
soit une extension de corps de type fini et que
$\text{d}a$ soit nul dans $\Omega_{L/k}$.
Ainsi, nous pouvons supposer que $K$ est une extension de corps de type fini
de $k$.

\medskip\noindent
Choisissons une base de transcendance $x_1, \ldots, x_r \in K$
telle que $K$ soit une extension finie séparable de $k(x_1, \ldots, x_r)$.
Cela est possible par définition ; voir les
Définitions \ref{algebra-definition-perfect} et
\ref{algebra-definition-separable-field-extension}.
Remarquons que le résultat vaut pour le sous-corps
$k(x_1, \ldots, x_r) \subset K$ ; ce sous-corps est une extension transcendante pure du corps de base.
En effet,
$$
\Omega_{k(x_1, \ldots, x_r)/k} =
\bigoplus\nolimits_{i = 1}^r k(x_1, \ldots, x_r) \text{d}x_i
$$
et toute fonction rationnelle dont toutes les dérivées partielles sont nulles
est une puissance $p$-ième. De plus, nous avons aussi
$$
\Omega_{K/k} =
\bigoplus\nolimits_{i = 1}^r K\text{d}x_i
$$
puisque l'extension $k(x_1, \ldots, x_r) \subset K$ est finie séparable
(calcul omis). Supposons que $a \in K$ soit un élément tel que
$\text{d}a = 0$ dans le module des différentielles. Par notre choix des $x_i$, nous
voyons que le polynôme minimal $P(T) \in k(x_1, \ldots, x_r)[T]$
de $a$ est séparable. Écrivons
$$
P(T) = T^d + \sum\nolimits_{i = 1}^d a_i T^{d - i}
$$
et donc
$$
0 = \text{d}P(a) = \sum\nolimits_{i = 1}^d a^{d - i}\text{d}a_i
$$
dans $\Omega_{K/k}$. D'après la description de
$\Omega_{K/k}$ ci-dessus et le fait que $P$ était le polynôme minimal
de $a$, nous voyons que cela implique $\text{d}a_i = 0$.
Ainsi, $a_i = b_i^p$ pour tout $i$. Par conséquent, d'après
Corps, Lemme \ref{fields-lemma-pth-root},
nous voyons que $a$ est une puissance $p$-ième.
\end{proof}

\begin{lemma}
\label{algebra-lemma-size-extension-pth-roots}
Soit $k$ un corps de caractéristique $p > 0$.
Soient $a_1, \ldots, a_n \in k$ des éléments tels que
$\text{d}a_1, \ldots, \text{d}a_n$ soient linéairement indépendants dans
$\Omega_{k/\mathbf{F}_p}$. Alors l'extension de corps
$k(a_1^{1/p}, \ldots, a_n^{1/p})$ est de degré $p^n$ sur $k$.
\end{lemma}

\begin{proof}
Nous procédons par récurrence sur $n$. Si $n = 1$, le résultat est le
Lemme \ref{algebra-lemma-derivative-zero-pth-power}.
Pour l'étape de récurrence, supposons que $k(a_1^{1/p}, \ldots, a_{n - 1}^{1/p})$
soit de degré $p^{n - 1}$ sur $k$. Nous devons montrer que $a_n$ ne devient pas
une puissance $p$-ième dans $k(a_1^{1/p}, \ldots, a_{n - 1}^{1/p})$.
Si c'est le cas, nous pouvons écrire
\begin{align*}
a_n & =
\left(\sum\nolimits_{I = (i_1, \ldots, i_{n - 1}),\ 0 \leq i_j \leq p - 1}
\lambda_I a_1^{i_1/p} \ldots a_{n - 1}^{i_{n - 1}/p}\right)^p \\
& = \sum\nolimits_{I = (i_1, \ldots, i_{n - 1}),\ 0 \leq i_j \leq p - 1}
\lambda_I^p a_1^{i_1} \ldots a_{n - 1}^{i_{n - 1}}
\end{align*}
En appliquant $\text{d}$, nous voyons que $\text{d}a_n$ dépend linéairement des
$\text{d}a_i$, $i < n$. C'est une contradiction.
\end{proof}

\begin{lemma}
\label{algebra-lemma-separable-differentials}
Soit $k$ un corps de caractéristique $p > 0$.
Les propriétés suivantes sont équivalentes :
\begin{enumerate}
\item l'extension de corps $K/k$ est séparable
(voir la Définition \ref{algebra-definition-separable-field-extension}), et
\item l'application
$K \otimes_k \Omega_{k/\mathbf{F}_p} \to \Omega_{K/\mathbf{F}_p}$
est injective.
\end{enumerate}
\end{lemma}

\begin{proof}
Écrivons $K$ comme une limite inductive filtrante $K = \colim_i K_i$ d'extensions de corps
$K_i/k$ de type fini. Par définition, $K$ est séparable si et seulement
si chaque $K_i$ est séparable sur $k$ et, d'après le
Lemme \ref{algebra-lemma-colimit-differentials}, nous voyons que
$K \otimes_k \Omega_{k/\mathbf{F}_p} \to \Omega_{K/\mathbf{F}_p}$
est injective si et seulement si chaque application
$K_i \otimes_k \Omega_{k/\mathbf{F}_p} \to \Omega_{K_i/\mathbf{F}_p}$
est injective. Ainsi, nous pouvons supposer que $K/k$ est une extension de corps
de type fini.

\medskip\noindent
Supposons que $K/k$ soit une extension de corps de type fini qui soit
séparable. Choisissons $x_1, \ldots, x_{r + 1} \in K$ comme dans le
Lemme \ref{algebra-lemma-generating-finitely-generated-separable-field-extensions}.
Dans ce cas, il existe un polynôme irréductible
$G(X_1, \ldots, X_{r + 1}) \in k[X_1, \ldots, X_{r + 1}]$
tel que $G(x_1, \ldots, x_{r + 1}) = 0$ et tel que
$\partial G/\partial X_{r + 1}$ ne soit pas identiquement nul.
De plus, $K$ est le corps des fractions de l'anneau intègre
$S = K[X_1, \ldots, X_{r + 1}]/(G)$.
Écrivons
$$
G = \sum a_I X^I, \quad X^I = X_1^{i_1}\ldots X_{r + 1}^{i_{r + 1}}.
$$
À l'aide de la présentation de $S$ ci-dessus, nous voyons que
$$
\Omega_{S/\mathbf{F}_p}
=
\frac{
S \otimes_k \Omega_k \oplus
\bigoplus\nolimits_{i = 1, \ldots, r + 1} S\text{d}X_i
}{
\langle
\sum X^I \text{d}a_I + \sum \partial G/\partial X_i \text{d}X_i
\rangle
}
$$
Puisque $\Omega_{K/\mathbf{F}_p}$ est la localisation
du $S$-module $\Omega_{S/\mathbf{F}_p}$ (voir le
Lemme \ref{algebra-lemma-differentials-localize}), nous en concluons
que
$$
\Omega_{K/\mathbf{F}_p}
=
\frac{
K \otimes_k \Omega_k \oplus
\bigoplus\nolimits_{i = 1, \ldots, r + 1} K\text{d}X_i
}{
\langle
\sum X^I \text{d}a_I + \sum \partial G/\partial X_i \text{d}X_i
\rangle
}
$$
Maintenant, puisque le polynôme $\partial G/\partial X_{r + 1}$ n'est pas
identiquement nul, nous en concluons que l'application
$K \otimes_k \Omega_{k/\mathbf{F}_p} \to \Omega_{S/\mathbf{F}_p}$
est injective, comme souhaité.

\medskip\noindent
Supposons que $K/k$ soit une extension de corps de type fini
et que
$K \otimes_k \Omega_{k/\mathbf{F}_p} \to \Omega_{K/\mathbf{F}_p}$
soit injective.
(Cette partie de la preuve est identique à l'argument démontrant le
Lemme \ref{algebra-lemma-characterize-separable-field-extensions}.)
Soit $x_1, \ldots, x_r$ une base de transcendance de $K$ sur $k$ telle
que le degré d'inséparabilité de l'extension finie
$k(x_1, \ldots, x_r) \subset K$ soit minimal.
Si $K$ est séparable sur $k(x_1, \ldots, x_r)$, alors la conclusion est acquise.
Supposons que ce ne soit pas le cas afin d'obtenir une contradiction.
Il existe alors un élément $\alpha \in K$ qui n'est pas
séparable sur $k(x_1, \ldots, x_r)$. Soit $P(T) \in k(x_1, \ldots, x_r)[T]$
son polynôme minimal. Puisque $\alpha$ n'est pas séparable,
$P$ est en fait un polynôme en $T^p$. Chassons les dénominateurs
pour obtenir un polynôme irréductible
$$
G(X_1, \ldots, X_r, T) = \sum a_{I, i} X^I T^i \in k[X_1, \ldots, X_r, T]
$$
tel que $G(x_1, \ldots, x_r, \alpha) = 0$ dans $K$.
Notons que cela signifie que $k[X_1, \ldots, X_r, T]/(G) \subset K$.
Nous pouvons supposer que, pour un couple $(I_0, i_0)$, le coefficient
$a_{I_0, i_0} = 1$.
Nous affirmons que $\text{d}G/\text{d}X_i$ n'est pas identiquement nulle
pour au moins un $i$. En effet, si ce n'est pas le cas,
$G$ est en fait un polynôme en $X_1^p, \ldots, X_r^p, T^p$.
Cela signifie alors que
$$
\sum\nolimits_{(I, i) \not = (I_0, i_0)} x^I\alpha^i \text{d}a_{I, i}
$$
est nulle dans $\Omega_{K/\mathbf{F}_p}$. Notons qu'il n'existe aucune
relation $k$-linéaire entre les éléments
$$
\{x^I\alpha^i \mid a_{I, i} \not = 0 \text{ et } (I, i) \not = (I_0, i_0)\}
$$
de $K$. L'hypothèse
que $K \otimes_k \Omega_{k/\mathbf{F}_p} \to \Omega_{K/\mathbf{F}_p}$
est injective implique donc que $\text{d}a_{I, i} = 0$
dans $\Omega_{k/\mathbf{F}_p}$ pour tout couple $(I, i)$.
D'après le Lemme \ref{algebra-lemma-derivative-zero-pth-power},
nous voyons que chaque $a_{I, i}$ est une puissance $p$-ième, ce qui
implique que $G$ est une puissance $p$-ième, en contradiction avec l'irréductibilité de
$G$. Ainsi,
après renumérotation, nous pouvons supposer que $\text{d}G/\text{d}X_1$ n'est pas nulle.
Nous voyons alors que $x_1$ est algébrique séparable sur
$k(x_2, \ldots, x_r, \alpha)$ et que $x_2, \ldots, x_r, \alpha$
forment une base de transcendance de $K$ sur $k$. Cela signifie que
le degré d'inséparabilité de l'extension finie
$k(x_2, \ldots, x_r, \alpha) \subset K$ est inférieur au
degré d'inséparabilité de l'extension finie
$k(x_1, \ldots, x_r) \subset K$, ce qui est une contradiction.
\end{proof}

\begin{lemma}
\label{algebra-lemma-formally-smooth-implies-separable}
Soit $K/k$ une extension de corps.
Si $K$ est formellement lisse sur $k$, alors $K$ est
une extension séparable de $k$.
\end{lemma}

\begin{proof}
Supposons que $K$ soit formellement lisse sur $k$.
D'après le Lemme \ref{algebra-lemma-ses-formally-smooth}, nous voyons que
$K \otimes_k \Omega_{k/\mathbf{Z}} \to \Omega_{K/\mathbf{Z}}$
est injective. Ainsi, $K$ est séparable sur $k$ d'après le
Lemme \ref{algebra-lemma-separable-differentials}.
\end{proof}

\begin{lemma}
\label{algebra-lemma-characterize-formally-smooth-field-extension}
Soit $K/k$ une extension de corps.
Alors $K$ est formellement lisse sur $k$ si et seulement si
$H_1(L_{K/k}) = 0$.
\end{lemma}

\begin{proof}
Cela résulte de la Proposition \ref{algebra-proposition-characterize-formally-smooth}
et du fait qu'un espace vectoriel est libre (donc projectif).
\end{proof}

\begin{lemma}
\label{algebra-lemma-formally-smooth-extensions-easy}
Soit $K/k$ une extension de corps.
\begin{enumerate}
\item Si $K$ est une extension transcendante pure de $k$, alors
$K$ est formellement lisse sur $k$.
\item Si $K$ est une extension algébrique séparable de $k$, alors $K$ est
formellement lisse sur $k$.
\item Si $K$ est séparable sur $k$, alors $K$ est formellement lisse
sur $k$.
\end{enumerate}
\end{lemma}

\begin{proof}
Pour (1), écrivons $K = k(x_j; j \in J)$. Supposons que
$A$ soit une $k$-algèbre et que $I \subset A$ soit un idéal de
carré nul. Soit $\varphi : K \to A/I$ un morphisme de $k$-algèbres.
Soit $a_j \in A$ un élément tel que $a_j \mod I = \varphi(x_j)$.
Il est alors aisé de voir qu'il existe un unique morphisme de $k$-algèbres
$K \to A$ qui envoie $x_j$ sur $a_j$ et qui se réduit
à $\varphi$ modulo $I$. Ainsi, $k \subset K$ est formellement lisse.

\medskip\noindent
Dans le cas (2), nous voyons que $k \subset K$ est une limite inductive
d'extensions d'anneaux \'etales. Un morphisme d'anneaux \'etale est formellement \'etale
(Lemme \ref{algebra-lemma-formally-etale-etale}). Ce cas résulte donc du
Lemme \ref{algebra-lemma-colimit-formally-etale} et de l'observation triviale
qu'un morphisme d'anneaux formellement \'etale est formellement lisse.

\medskip\noindent
Dans le cas (3), écrivons $K = \colim K_i$ comme la limite inductive filtrante de ses
sous-extensions de type fini sur $k$. D'après la
Définition \ref{algebra-definition-separable-field-extension},
chaque $K_i$ est une extension algébrique séparable d'une extension transcendante pure
de $k$. Ainsi, $K_i/k$ est formellement lisse d'après les cas (1) et (2) et le
Lemme \ref{algebra-lemma-compose-formally-smooth}. Donc
$H_1(L_{K_i/k}) = 0$ d'après le
Lemme \ref{algebra-lemma-characterize-formally-smooth-field-extension}.
Ainsi, $H_1(L_{K/k}) = 0$ d'après le Lemme \ref{algebra-lemma-colimits-NL}.
Donc $K/k$ est formellement lisse, de nouveau d'après le
Lemme \ref{algebra-lemma-characterize-formally-smooth-field-extension}.
\end{proof}

\begin{lemma}
\label{algebra-lemma-fields-are-formally-smooth}
\begin{slogan}
La lissité formelle équivaut à la séparabilité pour les extensions de corps.
\end{slogan}
Soit $k$ un corps.
\begin{enumerate}
\item Si la caractéristique de $k$ est nulle, alors toute extension de corps
de $k$ est formellement lisse sur $k$.
\item Si la caractéristique de $k$ est $p > 0$, alors $K/k$ est
formellement lisse si et seulement si cette extension de corps est séparable.
\end{enumerate}
\end{lemma}

\begin{proof}
On combine les Lemmes \ref{algebra-lemma-formally-smooth-implies-separable} et
\ref{algebra-lemma-formally-smooth-extensions-easy}.
\end{proof}

\noindent
Nous rassemblons ici toutes les caractérisations des extensions
de corps séparables.

\begin{proposition}
\label{algebra-proposition-characterize-separable-field-extensions}
Soit $K/k$ une extension de corps.
Si la caractéristique de $k$ est nulle, les assertions suivantes sont équivalentes :
\begin{enumerate}
\item $K$ est séparable sur $k$,
\item $K$ est géométriquement réduit sur $k$,
\item $K$ est formellement lisse sur $k$,
\item $H_1(L_{K/k}) = 0$, et
\item l'application $K \otimes_k \Omega_{k/\mathbf{Z}} \to \Omega_{K/\mathbf{Z}}$
est injective.
\end{enumerate}
Si la caractéristique de $k$ est $p > 0$, alors les propriétés suivantes sont
équivalentes :
\begin{enumerate}
\item $K$ est séparable sur $k$,
\item l'anneau $K \otimes_k k^{1/p}$ est réduit,
\item $K$ est géométriquement réduit sur $k$,
\item l'application $K \otimes_k \Omega_{k/\mathbf{F}_p} \to \Omega_{K/\mathbf{F}_p}$
est injective,
\item $H_1(L_{K/k}) = 0$, et
\item $K$ est formellement lisse sur $k$.
\end{enumerate}
\end{proposition}

\begin{proof}
C'est une combinaison des
Lemmes \ref{algebra-lemma-characterize-separable-field-extensions},
\ref{algebra-lemma-fields-are-formally-smooth},
\ref{algebra-lemma-formally-smooth-implies-separable} et
\ref{algebra-lemma-separable-differentials}.
\end{proof}

\noindent
Voici une autre caractérisation des extensions de corps séparables
de type fini.

\begin{lemma}
\label{algebra-lemma-localization-smooth-separable}
Soit $K/k$ une extension de corps de type fini.
Alors $K$ est séparable sur $k$ si et seulement si $K$ est
la localisation d'une $k$-algèbre lisse.
\end{lemma}

\begin{proof}
Choisissons une $k$-algèbre de type fini $R$ qui soit un anneau intègre dont le
corps des fractions est $K$. Le Lemme \ref{algebra-lemma-smooth-at-generic-point}
affirme que $k \to R$ est lisse
en $(0)$ si et seulement si $K/k$ est séparable.
Cela prouve le lemme.
\end{proof}

\begin{lemma}
\label{algebra-lemma-colimit-syntomic}
Soit $K/k$ une extension de corps.
Alors $K$ est une limite inductive filtrante d'algèbres qui sont des intersections complètes
globales sur $k$. Si $K/k$ est séparable, alors $K$ est une limite inductive filtrante
d'algèbres lisses sur $k$.
\end{lemma}

\begin{proof}
Supposons que $E \subset K$ soit un sous-ensemble fini. Il suffit de montrer qu'il
existe une $k$-sous-algèbre $A \subset K$ qui contient $E$
et qui est une intersection complète globale (resp.\ lisse) sur $k$.
Le cas séparable/lisse résulte du
Lemme \ref{algebra-lemma-localization-smooth-separable}.
Dans le cas général, soit $L \subset K$ le sous-corps engendré par $E$.
Choisissons une base de transcendance $x_1, \ldots, x_d \in L$ sur $k$.
L'extension $L/k(x_1, \ldots, x_d)$ est finie.
Écrivons $L = k(x_1, \ldots, x_d)[y_1, \ldots, y_r]$.
Choisissons par récurrence des polynômes $P_i \in k(x_1, \ldots, x_d)[Y_1, \ldots, Y_r]$
tels que $P_i = P_i(Y_1, \ldots, Y_i)$ soit unitaire en $Y_i$ sur
$k(x_1, \ldots, x_d)[Y_1, \ldots, Y_{i - 1}]$ et dont l'image est le
polynôme minimal de $y_i$ dans
$k(x_1, \ldots, x_d)[y_1, \ldots, y_{i - 1}][Y_i]$.
Il est alors clair que la suite $P_1, \ldots, P_r$ est régulière
dans $k(x_1, \ldots, x_r)[Y_1, \ldots, Y_r]$ et que
$L = k(x_1, \ldots, x_r)[Y_1, \ldots, Y_r]/(P_1, \ldots, P_r)$.
Si $h \in k[x_1, \ldots, x_d]$ est un polynôme tel que
$P_i \in k[x_1, \ldots, x_d, 1/h, Y_1, \ldots, Y_r]$, alors
nous voyons que $P_1, \ldots, P_r$ est une suite régulière dans
$k[x_1, \ldots, x_d, 1/h, Y_1, \ldots, Y_r]$ et
$A = k[x_1, \ldots, x_d, 1/h, Y_1, \ldots, Y_r]/(P_1, \ldots, P_r)$
est une intersection complète globale. Après avoir ajusté notre choix de $h$,
nous pouvons supposer que $E \subset A$, ce qui conclut.
\end{proof}






\section{Construction de morphismes d'anneaux plats}
\label{algebra-section-constructing-flat}

\noindent
Le lemme suivant est parfois utile.

\begin{lemma}
\label{algebra-lemma-flat-local-given-residue-field}
Soit $(R, \mathfrak m, k)$ un anneau local. Soit $K/k$ une
extension de corps. Il existe un anneau local $(R', \mathfrak m', k')$, un morphisme local plat
d'anneaux $R \to R'$ tel que $\mathfrak m' = \mathfrak mR'$ et tel que
$k'$ soit isomorphe à $K$ comme extension de $k$.
\end{lemma}

\begin{proof}
Supposons que $k' = k(\alpha)$ soit une extension monogène de $k$.
Alors $k'$ est le corps résiduel d'une extension locale plate $R \subset R'$
comme dans le lemme. Plus précisément, si $\alpha$ est transcendant sur $k$, nous prenons
pour $R'$ la localisation de $R[x]$ en l'idéal premier $\mathfrak mR[x]$.
Si $\alpha$ est algébrique de polynôme minimal
$T^d + \sum \overline{\lambda}_iT^{d - i}$, alors nous prenons
$R' = R[T]/(T^d + \sum \lambda_i T^{d - i})$.

\medskip\noindent
Considérons la famille des triplets $(k', R \to R', \phi)$, où
$k \subset k' \subset K$ est un sous-corps,
$R \to R'$ est un morphisme local d'anneaux comme dans le lemme, et
$\phi : R' \to k'$ induit un isomorphisme $R'/\mathfrak mR' \cong k'$
d'extensions de $k$. Ces triplets forment une ``grande'' catégorie $\mathcal{C}$ dont les morphismes
$(k_1, R_1, \phi_1) \to (k_2, R_2, \phi_2)$
sont donnés par des morphismes d'anneaux $\psi : R_1 \to R_2$ tels que
$$
\xymatrix{
R_1 \ar[d]_\psi \ar[r]_{\phi_1} & k_1 \ar[r] & K \ar@{=}[d] \\
R_2 \ar[r]^{\phi_2} & k_2 \ar[r] & K
}
$$
le diagramme soit commutatif. Cela implique que $k_1 \subset k_2$.

\medskip\noindent
Supposons que $I$ soit un ensemble ordonné filtrant et que
$((R_i, k_i, \phi_i), \psi_{ii'})$ soit un système inductif indexé par $I$ ; voir
Catégories, section \ref{categories-section-posets-limits}.
Dans ce cas, nous pouvons considérer
$$
R' = \colim_{i \in I} R_i
$$
C'est un anneau local d'idéal maximal $\mathfrak mR'$ et de
corps résiduel $k' = \bigcup_{i \in I} k_i$. De plus, le morphisme d'anneaux
$R \to R'$ est plat, car c'est une limite inductive de morphismes plats (et les produits
tensoriels commutent aux limites inductives filtrantes).
Ainsi, $(R', k', \phi')$ est un ``majorant'' du système.

\medskip\noindent
Une application presque triviale du lemme de Zorn achèverait la preuve
si $\mathcal{C}$ était un ensemble, mais ce n'en est pas un.
(En fait, on peut faire fonctionner cet argument en trouvant une majoration raisonnable des
cardinaux des anneaux locaux qui interviennent.)
Pour contourner ce problème, choisissons un bon ordre sur $K$.
Pour $x \in K$, notons $K(x)$ le sous-corps de $K$ engendré
par tous les éléments de $K$ qui sont $\leq x$.
Par récurrence transfinie sur $x \in K$, nous allons construire des morphismes d'anneaux
$R \subset R(x)$ comme dans le lemme, avec pour extension de corps résiduels
$K(x)/k$. De plus, la construction assure
que $R(x)$ contient $R(y)$ pour tout $y \leq x$.
En effet, si $x$ a un prédécesseur $x'$, alors $K(x) = K(x')[x]$
et nous pouvons donc prendre pour $R(x') \subset R(x)$ l'extension d'anneaux locaux
construite dans le premier paragraphe de la preuve. Si $x$ n'a pas
de prédécesseur, posons d'abord
$R'(x) = \colim_{x' < x} R(x')$ comme dans le troisième paragraphe
de la preuve. Le corps résiduel de $R'(x)$ est $K'(x) = \bigcup_{x' < x} K(x')$.
Puisque $K(x) = K'(x)[x]$, nous pouvons utiliser la construction du
premier paragraphe de la preuve pour produire $R'(x) \subset R(x)$.
Cela achève la preuve du lemme.
\end{proof}

\begin{lemma}
\label{algebra-lemma-colimit-finite-etale-given-residue-field}
Soit $(R, \mathfrak m, k)$ un anneau local. Si $k \subset K$ est une
extension algébrique séparable, alors il existe un ensemble ordonné filtrant $I$ et
un système inductif d'extensions finies \'etales $R \subset R_i$, $i \in I$
d'anneaux locaux tel que $R' = \colim R_i$ ait pour corps résiduel
$K$ (comme extension de $k$).
\end{lemma}

\begin{proof}
Soit $R \subset R'$ l'extension construite dans la preuve du
Lemme \ref{algebra-lemma-flat-local-given-residue-field}. Par construction,
$R' = \colim_{\alpha \in A} R_\alpha$, où $A$ est un ensemble bien ordonné,
et les morphismes de transition $R_\alpha \to R_{\alpha + 1}$
sont finis \'etales, tandis que $R_\alpha = \colim_{\beta < \alpha} R_\beta$
si $\alpha$ n'est pas un successeur. Nous allons démontrer le résultat par récurrence
transfinie.

\medskip\noindent
Supposons le résultat vrai pour $R_\alpha$, c'est-à-dire $R_\alpha = \colim R_i$,
où $R_i$ est fini \'etale sur $R$. Puisque
$R_\alpha \to R_{\alpha + 1}$ est fini \'etale,
il existe un $i$ et une extension finie \'etale $R_i \to R_{i, 1}$
tels que $R_{\alpha + 1} = R_\alpha \otimes_{R_i} R_{i, 1}$.
Ainsi, $R_{\alpha + 1} = \colim_{i' \geq i} R_{i'} \otimes_{R_i} R_{i, 1}$,
et le résultat est vrai pour $\alpha + 1$. Supposons que $\alpha$ ne soit pas un successeur
et que le résultat soit vrai pour $R_\beta$ pour tout $\beta < \alpha$.
Tout sous-ensemble fini $E \subset R_\alpha$ est contenu dans $R_\beta$
pour un certain $\beta < \alpha$ et, par hypothèse, $E$ est contenu dans une sous-extension finie \'etale.
Ainsi, le résultat est vrai pour $R_\alpha$.
\end{proof}

\begin{lemma}
\label{algebra-lemma-finite-free-given-residue-field-extension}
Soit $R$ un anneau. Soit $\mathfrak p \subset R$ un idéal premier et
soit $L/\kappa(\mathfrak p)$ une extension finie de corps.
Alors il existe un morphisme d'anneaux fini et libre $R \to S$ tel que
$\mathfrak q = \mathfrak pS$ soit premier et que
$\kappa(\mathfrak q)/\kappa(\mathfrak p)$ soit isomorphe à l'extension donnée
$L/\kappa(\mathfrak p)$.
\end{lemma}

\begin{proof}
Nous procédons par récurrence sur le degré de l'extension $\kappa(\mathfrak p) \subset L$.
Si le degré est $1$, nous prenons $R = S$.
En général, s'il existe une sous-extension
$\kappa(\mathfrak p) \subset L' \subset L$, alors la conclusion est acquise par récurrence
sur le degré (en construisant d'abord $R \subset S'$ correspondant
à $L'/\kappa(\mathfrak p)$, puis en construisant $S' \subset S$
correspondant à $L/L'$). Nous pouvons donc supposer que
l'extension $L \supset \kappa(\mathfrak p)$ est engendrée par un seul élément
$\alpha \in L$. Soit $X^d + \sum_{i < d} a_iX^i$ le polynôme minimal
de $\alpha$ sur $\kappa(\mathfrak p)$, de sorte que $a_i \in \kappa(\mathfrak p)$.
Nous pouvons écrire $a_i$ comme l'image
de $f_i/g$ pour certains $f_i, g \in R$ avec $g \not \in \mathfrak p$.
Après avoir remplacé $\alpha$ par $g\alpha$ (et, corrélativement,
remplacé $a_i$ par $g^{d - i}a_i$), nous pouvons supposer que $a_i$ est
l'image d'un élément $f_i \in R$.
Nous prenons alors simplement $S = R[x]/(x^d + \sum f_ix^i)$.
\end{proof}

\begin{lemma}
\label{algebra-lemma-cofinal-system-flat}
Soit $A$ un anneau. Posons $\kappa = \max(|A|, \aleph_0)$. Alors toute
$A$-algèbre plate $B$ est la limite inductive filtrante de ses $A$-sous-algèbres plates
$B' \subset B$ telles que $|B'| \leq \kappa$. (Remarquons que $B'$ est
fidèlement plate sur $A$ si $B$ est fidèlement plate sur $A$.)
\end{lemma}

\begin{proof}
Si $B$ est de cardinalité $\leq \kappa$, l'assertion est vraie.
Soit $E \subset B$ une $A$-sous-algèbre telle que $|E| \leq \kappa$.
Nous allons montrer que $E$ est contenue dans une $A$-sous-algèbre plate
$B'$ telle que $|B'| \leq \kappa$. Le lemme s'en déduit, car
(a) tout sous-ensemble fini de $B$ est contenu dans une $A$-sous-algèbre de
cardinalité au plus $\kappa$ et (b) toute paire de $A$-sous-algèbres
de $B$ de cardinalité au plus $\kappa$ est contenue dans une $A$-sous-algèbre
de cardinalité au plus $\kappa$. Détails omis.

\medskip\noindent
Nous allons construire par récurrence une suite de $A$-sous-algèbres
$$
E = E_0 \subset E_1 \subset E_2 \subset \ldots
$$
chacune de cardinalité $\leq \kappa$, et nous montrerons que
$B' = \bigcup E_k$ est plate sur $A$, ce qui achèvera la preuve.

\medskip\noindent
La construction est la suivante. Posons $E_0 = E$.
Étant donné $E_k$ pour $k \geq 0$, considérons l'ensemble $S_k$ des
relations entre des éléments de $E_k$ à coefficients dans $A$.
Ainsi, un élément $s \in S_k$ est la donnée d'un entier $n \geq 1$, de
$a_1, \ldots, a_n \in A$ et de $e_1, \ldots, e_n \in E_k$
tels que $\sum a_i e_i = 0$ dans $E_k$.
La platitude de $A \to B$ implique, d'après le Lemme \ref{algebra-lemma-flat-eq},
que, pour tout $s = (n, a_1, \ldots, a_n, e_1, \ldots, e_n) \in S_k$,
nous pouvons choisir
$$
(m_s, b_{s, 1}, \ldots, b_{s, m_s}, a_{s, 11}, \ldots, a_{s, nm_s})
$$
où $m_s \geq 0$ est un entier, $b_{s, j} \in B$, $a_{s, ij} \in A$, et
$$
e_i = \sum\nolimits_j a_{s, ij} b_{s, j}, \forall i,
\quad\text{et}\quad
0 = \sum\nolimits_i a_i a_{s, ij}, \forall j.
$$
Ces choix étant faits, soit $E_{k + 1} \subset B$ la $A$-sous-algèbre
engendrée par
\begin{enumerate}
\item $E_k$ et
\item les éléments $b_{s, 1}, \ldots, b_{s, m_s}$
pour tout $s \in S_k$.
\end{enumerate}
Un argument de théorie des ensembles (omis) montre que $E_{k + 1}$ est de
cardinalité au plus $\kappa$ (on utilise ici que, par récurrence,
$|E_k| \leq \kappa$ et que, par conséquent, la cardinalité de $S_k$ est
également au plus $\kappa$).

\medskip\noindent
Pour montrer que $B' = \bigcup E_k$ est plate sur $A$, considérons
une relation $\sum_{i = 1, \ldots, n} a_i b'_i = 0$ dans $B'$
à coefficients dans $A$. Choisissons $k$ assez grand pour que
$b'_i \in E_k$ pour $i = 1, \ldots, n$. Alors
$(n, a_1, \ldots, a_n, b'_1, \ldots, b'_n) \in S_k$,
et nous voyons donc que la relation est triviale dans $E_{k + 1}$
et, a fortiori, dans $B'$.
Ainsi, $A \to B'$ est plat d'après le Lemme \ref{algebra-lemma-flat-eq}.
\end{proof}









\section{Le théorème de structure de Cohen}
\label{algebra-section-cohen-structure-theorem}

\noindent
Voici une notion fondamentale d'algèbre commutative.

\begin{definition}
\label{algebra-definition-complete-local-ring}
Soit $(R, \mathfrak m)$ un anneau local. Nous disons que $R$ est un
{\it anneau local complet} si le morphisme canonique
$$
R \longrightarrow \lim_n R/\mathfrak m^n
$$
vers le complété de $R$ relativement à $\mathfrak m$ est un
isomorphisme\footnote{Cela inclut la condition
que $\bigcap \mathfrak m^n = (0)$ ; dans certains textes, on l'indique
en disant que $R$ est complet et séparé. Attention : il peut arriver
que le complété $\lim_n R/\mathfrak m^n$ d'un anneau local soit
non complet ; voir
Exemples, Lemme \ref{examples-lemma-noncomplete-completion}.
Cela ne se produit pas lorsque $\mathfrak m$ est de type fini ; voir le
Lemme \ref{algebra-lemma-hathat-finitely-generated}, auquel
cas le complété est noethérien ; voir le
Lemme \ref{algebra-lemma-completion-Noetherian}.}.
\end{definition}

\noindent
Remarquons qu'un anneau local artinien $R$ est un anneau local complet
car $\mathfrak m_R^n = 0$ pour un certain $n > 0$. Dans cette section,
nous nous concentrons principalement sur les anneaux locaux noethériens complets.

\begin{lemma}
\label{algebra-lemma-quotient-complete-local}
Soit $R$ un anneau local noethérien complet.
Tout quotient de $R$ est aussi un anneau local noethérien complet.
Étant donné un morphisme d'anneaux fini $R \to S$, l'anneau $S$ est
un produit d'anneaux locaux noethériens complets.
\end{lemma}

\begin{proof}
L'anneau $S$ est noethérien d'après le Lemme \ref{algebra-lemma-Noetherian-permanence}.
En tant que $R$-module, $S$ est complet d'après le Lemme \ref{algebra-lemma-completion-tensor}.
Ainsi, $S$ est le produit de ses complétés relativement à ses idéaux maximaux
d'après le Lemme \ref{algebra-lemma-completion-finite-extension}.
\end{proof}

\begin{lemma}
\label{algebra-lemma-complete-local-ring-Noetherian}
Soit $(R, \mathfrak m)$ un anneau local complet.
Si $\mathfrak m$ est un idéal de type fini, alors
$R$ est noethérien.
\end{lemma}

\begin{proof}
Voir le Lemme \ref{algebra-lemma-completion-Noetherian}.
\end{proof}

\begin{definition}
\label{algebra-definition-coefficient-ring}
Soit $(R, \mathfrak m)$ un anneau local complet.
Un sous-anneau $\Lambda \subset R$ est
appelé un {\it anneau de coefficients} si les conditions suivantes sont satisfaites :
\begin{enumerate}
\item $\Lambda$ est un anneau local complet d'idéal maximal
$\Lambda \cap \mathfrak m$,
\item le corps résiduel de $\Lambda$ s'envoie isomorphiquement sur le
corps résiduel de $R$, et
\item $\Lambda \cap \mathfrak m = p\Lambda$, où $p$ est la caractéristique
du corps résiduel de $R$.
\end{enumerate}
\end{definition}

\noindent
Faisons quelques remarques sur cette définition. Nous distinguons les
cas suivants :
\begin{enumerate}
\item L'anneau local $R$ contient un corps. Cela se produit
lorsque $\mathbf{Q} \subset R$, ou lorsque $pR = 0$, où $p$ est la
caractéristique de $R/\mathfrak m$. Dans ce cas, un anneau de coefficients
$\Lambda$ est un corps contenu dans $R$ qui s'envoie isomorphiquement sur
$R/\mathfrak m$.
\item La caractéristique de $R/\mathfrak m$ est $p > 0$, mais aucune
puissance de $p$ n'est nulle dans $R$. Dans ce cas, $\Lambda$ est un anneau
de valuation discrète complet d'uniformisante $p$ et de corps résiduel $R/\mathfrak m$.
\item La caractéristique de $R/\mathfrak m$ est $p > 0$ et, pour un certain
$n > 1$, nous avons $p^{n - 1} \not = 0$, $p^n = 0$ dans $R$. Dans ce cas,
$\Lambda$ est un anneau local artinien dont l'idéal maximal est
engendré par $p$ et dont le corps résiduel est $R/\mathfrak m$.
\end{enumerate}
Les anneaux de valuation discrète complets d'uniformisante $p$
ci-dessus jouent un rôle particulier ; donnons-leur le nom suivant.

\begin{definition}
\label{algebra-definition-cohen-ring}
Un {\it anneau de Cohen} est un anneau de valuation discrète complet dont
l'uniformisante $p$ est un nombre premier.
\end{definition}

\begin{lemma}
\label{algebra-lemma-cohen-rings-exist}
Soit $p$ un nombre premier.
Soit $k$ un corps de caractéristique $p$.
Il existe un anneau de Cohen $\Lambda$ tel que $\Lambda/p\Lambda \cong k$.
\end{lemma}

\begin{proof}
Remarquons d'abord que les entiers $p$-adiques $\mathbf{Z}_p$ forment un anneau de Cohen
de corps résiduel $\mathbf{F}_p$. Soit $k$ un corps arbitraire de caractéristique $p$.
Soit $\mathbf{Z}_p \to R$ un morphisme local plat d'anneaux tel que
$\mathfrak m_R = pR$ et $R/pR = k$ ; voir le
Lemme \ref{algebra-lemma-flat-local-given-residue-field}.
D'après le Lemme \ref{algebra-lemma-completion-Noetherian}, le complété
$\Lambda = R^\wedge$ est noethérien. Puisque $\Lambda/p\Lambda = R/pR$
est un corps, nous concluons que $\Lambda$ est un
anneau local noethérien complet d'idéal maximal $(p)$
(l'égalité et la complétude résultent du Lemme \ref{algebra-lemma-hathat-finitely-generated}).
Pour tout $n$, le morphisme $\mathbf{Z}/p^n\mathbf{Z} \to \Lambda/p^n\Lambda = R/p^nR$
est plat (l'égalité résulte de nouveau du Lemme \ref{algebra-lemma-hathat-finitely-generated}).
Ainsi, $\mathbf{Z}_p \to \Lambda$ est plat, par exemple d'après le
Lemme \ref{algebra-lemma-flat-module-powers}. Ainsi, $p$ est un
non-diviseur de zéro dans $\Lambda$. Par conséquent, $\Lambda$ est de dimension $\geq 1$
(Lemme \ref{algebra-lemma-one-equation}),
et nous concluons que $\Lambda$ est régulier de dimension $1$, c'est-à-dire
un anneau de valuation discrète d'après le Lemme \ref{algebra-lemma-characterize-dvr}.
Nous concluons que $\Lambda$ est un anneau de Cohen de corps résiduel $k$.
\end{proof}

\begin{lemma}
\label{algebra-lemma-cohen-ring-formally-smooth}
Soit $p > 0$ un nombre premier.
Soit $\Lambda$ un anneau de Cohen dont le corps résiduel est de caractéristique $p$.
Pour tout $n \geq 1$, le morphisme d'anneaux
$$
\mathbf{Z}/p^n\mathbf{Z} \to \Lambda/p^n\Lambda
$$
est formellement lisse.
\end{lemma}

\begin{proof}
Si $n = 1$, cela résulte de la
Proposition \ref{algebra-proposition-characterize-separable-field-extensions}.
Pour $n$ quelconque, nous procédons par récurrence sur $n$.
En effet, si $\mathbf{Z}/p^n\mathbf{Z} \to \Lambda/p^n\Lambda$ est
formellement lisse, alors nous pouvons appliquer le Lemme \ref{algebra-lemma-lift-formal-smoothness}
au morphisme d'anneaux
$\mathbf{Z}/p^{n + 1}\mathbf{Z} \to \Lambda/p^{n + 1}\Lambda$
et à l'idéal $I = (p^n) \subset \mathbf{Z}/p^{n + 1}\mathbf{Z}$.
\end{proof}

\begin{theorem}[Théorème de structure de Cohen]
\label{algebra-theorem-cohen-structure-theorem}
Soit $(R, \mathfrak m)$ un anneau local complet.
\begin{enumerate}
\item $R$ possède un anneau de coefficients (voir la
Définition \ref{algebra-definition-coefficient-ring}),
\item si $\mathfrak m$ est un idéal de type fini, alors
$R$ est isomorphe à un quotient
$$
\Lambda[[x_1, \ldots, x_n]]/I
$$
où $\Lambda$ est soit un corps, soit un anneau de Cohen.
\end{enumerate}
\end{theorem}

\begin{proof}
Montrons qu'il existe un anneau de coefficients.
Nous le démontrons d'abord lorsque la caractéristique du corps résiduel $\kappa$
est nulle. Dans ce cas, nous allons démontrer par récurrence
sur $n > 0$ qu'il existe une section
$$
\varphi_n : \kappa \longrightarrow R/\mathfrak m^n
$$
du morphisme canonique $R/\mathfrak m^n \to \kappa = R/\mathfrak m$.
C'est immédiat pour $n = 1$. Si $n > 1$, supposons $\varphi_{n - 1}$ donnée.
L'extension de corps $\kappa/\mathbf{Q}$ est formellement lisse d'après la
Proposition \ref{algebra-proposition-characterize-separable-field-extensions}.
Nous pouvons donc trouver la flèche en pointillés
dans le diagramme suivant :
$$
\xymatrix{
R/\mathfrak m^{n - 1} &
R/\mathfrak m^n \ar[l] \\
\kappa \ar[u]^{\varphi_{n - 1}} \ar@{..>}[ru] & \mathbf{Q} \ar[l] \ar[u]
}
$$
Cela démontre l'étape de récurrence. En passant à la limite dans ces morphismes,
$$
\lim_n\ \varphi_n : \kappa \longrightarrow
R = \lim_n\ R/\mathfrak m^n
$$
nous obtenons un morphisme dont l'image est l'anneau de coefficients recherché.

\medskip\noindent
Démontrons ensuite l'existence d'un anneau de coefficients lorsque
la caractéristique du corps résiduel $\kappa$ est $p > 0$.
Choisissons un anneau de Cohen $\Lambda$ tel que $\kappa = \Lambda/p\Lambda$ ;
voir le Lemme \ref{algebra-lemma-cohen-rings-exist}. Dans ce cas, nous allons démontrer par
récurrence sur $n > 0$ qu'il existe un morphisme
$$
\varphi_n :
\Lambda/p^n\Lambda
\longrightarrow
R/\mathfrak m^n
$$
dont la composée avec le morphisme de réduction $R/\mathfrak m^n \to \kappa$
donne l'isomorphisme fixé $\Lambda/p\Lambda = \kappa$. C'est immédiat
pour $n = 1$. Si $n > 1$, supposons $\varphi_{n - 1}$ donnée.
Le morphisme d'anneaux $\mathbf{Z}/p^n\mathbf{Z} \to \Lambda/p^n\Lambda$
est formellement lisse d'après le Lemme \ref{algebra-lemma-cohen-ring-formally-smooth}.
Nous pouvons donc trouver la flèche en pointillés
dans le diagramme suivant :
$$
\xymatrix{
R/\mathfrak m^{n - 1} &
R/\mathfrak m^n \ar[l] \\
\Lambda/p^n\Lambda \ar[u]^{\varphi_{n - 1}} \ar@{..>}[ru] &
\mathbf{Z}/p^n\mathbf{Z} \ar[l] \ar[u]
}
$$
Cela démontre l'étape de récurrence. En passant à la limite dans ces morphismes,
$$
\lim_n\ \varphi_n :
\Lambda = \lim_n\ \Lambda/p^n\Lambda
\longrightarrow
R = \lim_n\ R/\mathfrak m^n
$$
nous obtenons un morphisme dont l'image est l'anneau de coefficients recherché.

\medskip\noindent
Le dernier énoncé du théorème en résulte aisément. En effet, si
$y_1, \ldots, y_n$ sont des générateurs de l'idéal $\mathfrak m$,
nous pouvons utiliser le morphisme $\Lambda \to R$ que nous venons de construire
pour obtenir un morphisme
$$
\Lambda[[x_1, \ldots, x_n]] \longrightarrow R,
\quad x_i \longmapsto y_i.
$$
Comme les deux membres sont $(x_1, \ldots, x_n)$-adiquement complets,
ce morphisme est surjectif d'après le Lemme \ref{algebra-lemma-completion-generalities},
puisqu'il est surjectif modulo $(x_1, \ldots, x_n)$ par
construction.
\end{proof}

\begin{remark}
\label{algebra-remark-Noetherian-complete-local-ring-universally-catenary}
Si $k$ est un corps, alors l'anneau de séries formelles $k[[X_1, \ldots, X_d]]$
est un anneau local noethérien complet et régulier de dimension $d$.
Si $\Lambda$ est un anneau de Cohen, alors $\Lambda[[X_1, \ldots, X_d]]$
est un anneau local noethérien complet et régulier de dimension $d + 1$.
Le théorème de structure de Cohen implique donc que tout anneau local
noethérien complet est un quotient d'un anneau local régulier.
En particulier, nous voyons qu'un anneau local noethérien complet est
universellement caténaire ; voir le Lemme \ref{algebra-lemma-CM-ring-catenary}
et le Lemme \ref{algebra-lemma-regular-ring-CM}.
\end{remark}

\begin{lemma}
\label{algebra-lemma-regular-complete-containing-coefficient-field}
Soit $(R, \mathfrak m)$ un anneau local noethérien complet.
Supposons que $R$ soit régulier.
\begin{enumerate}
\item Si $R$ contient soit $\mathbf{F}_p$, soit $\mathbf{Q}$, alors $R$
est isomorphe à un anneau de séries formelles sur son corps résiduel.
\item Si $k$ est un corps et si $k \to R$ est un morphisme d'anneaux induisant
un isomorphisme $k \to R/\mathfrak m$, alors $R$ est isomorphe,
comme $k$-algèbre, à un anneau de séries formelles sur $k$.
\end{enumerate}
\end{lemma}

\begin{proof}
Dans le cas (1), le théorème de structure de Cohen
(Théorème \ref{algebra-theorem-cohen-structure-theorem})
fournit un anneau de coefficients qui est nécessairement un corps
et qui s'envoie isomorphiquement sur le corps résiduel. Ainsi,
il suffit de démontrer (2). Dans le cas (2), choisissons
$f_1, \ldots, f_d \in \mathfrak m$ dont les images
forment une base de $\mathfrak m/\mathfrak m^2$, et considérons
le morphisme continu de $k$-algèbres $k[[x_1, \ldots, x_d]] \to R$
qui envoie $x_i$ sur $f_i$. Comme la source et le but sont
$(x_1, \ldots, x_d)$-adiquement complets, ce morphisme est surjectif d'après le
Lemme \ref{algebra-lemma-completion-generalities}. D'autre part, il
est nécessairement injectif, car sinon la dimension de
$R$ serait $< d$ d'après le Lemme \ref{algebra-lemma-one-equation}.
\end{proof}

\begin{lemma}
\label{algebra-lemma-complete-local-Noetherian-domain-finite-over-regular}
Soit $(R, \mathfrak m)$ un anneau local noethérien complet et intègre.
Alors il existe un sous-anneau $R_0 \subset R$ ayant les propriétés suivantes :
\begin{enumerate}
\item $R_0$ est un anneau local régulier et complet,
\item l'inclusion $R_0 \subset R$ est finie et induit un isomorphisme sur les
corps résiduels,
\item $R_0$ est soit isomorphe à $k[[X_1, \ldots, X_d]]$, où $k$
est un corps, soit à $\Lambda[[X_1, \ldots, X_d]]$, où $\Lambda$ est un anneau de Cohen.
\end{enumerate}
\end{lemma}

\begin{proof}
Soit $\Lambda$ un anneau de coefficients de $R$.
Puisque $R$ est intègre, nous voyons que $\Lambda$ est un corps
ou que $\Lambda$ est un anneau de Cohen.

\medskip\noindent
Cas I : $\Lambda = k$ est un corps. Posons $d = \dim(R)$.
Choisissons $x_1, \ldots, x_d \in \mathfrak m$
qui engendrent un idéal de définition $I \subset R$.
(Voir la section \ref{algebra-section-dimension}.)
D'après le Lemme \ref{algebra-lemma-change-ideal-completion}, nous voyons que $R$
est également $I$-adiquement complet.
Considérons le morphisme $R_0 = k[[X_1, \ldots, X_d]] \to R$
qui envoie $X_i$ sur $x_i$.
Remarquons que $R_0$ est complet relativement à l'idéal
$I_0 = (X_1, \ldots, X_d)$,
et que $R/I_0R \cong R/IR$ est fini sur $k = R_0/I_0$
(car $\dim(R/I) = 0$ ; voir la section \ref{algebra-section-dimension}).
Nous concluons donc que $R_0 \to R$ est fini d'après le
Lemme \ref{algebra-lemma-finite-over-complete-ring}.
Puisque $\dim(R) = \dim(R_0)$, cela implique que
$R_0 \to R$ est injectif (voir le Lemme \ref{algebra-lemma-integral-dim-up}),
et le lemme est démontré.

\medskip\noindent
Cas II : $\Lambda$ est un anneau de Cohen. Posons $d + 1 = \dim(R)$.
Soit $p > 0$ la caractéristique du corps résiduel $k$.
Puisque $R$ est intègre, nous voyons que $p$ est un non-diviseur de zéro dans $R$.
Ainsi, $\dim(R/pR) = d$ ; voir le Lemme \ref{algebra-lemma-one-equation}.
Choisissons $x_1, \ldots, x_d \in R$
dont les images engendrent un idéal de définition dans $R/pR$.
Alors $I = (p, x_1, \ldots, x_d)$ est un idéal de définition de $R$.
D'après le Lemme \ref{algebra-lemma-change-ideal-completion}, nous voyons que $R$
est également $I$-adiquement complet.
Considérons le morphisme $R_0 = \Lambda[[X_1, \ldots, X_d]] \to R$
qui envoie $X_i$ sur $x_i$.
Remarquons que $R_0$ est complet relativement à l'idéal
$I_0 = (p, X_1, \ldots, X_d)$,
et que $R/I_0R \cong R/IR$ est fini sur $k = R_0/I_0$
(car $\dim(R/I) = 0$ ; voir la section \ref{algebra-section-dimension}).
Nous concluons donc que $R_0 \to R$ est fini d'après le
Lemme \ref{algebra-lemma-finite-over-complete-ring}.
Puisque $\dim(R) = \dim(R_0)$, cela implique que
$R_0 \to R$ est injectif (voir le Lemme \ref{algebra-lemma-integral-dim-up}),
et le lemme est démontré.
\end{proof}





\section{Anneaux japonais}
\label{algebra-section-japanese}

\noindent
Dans cette section, nous commençons l'étude de la finitude de la clôture intégrale.

\begin{definition}
\label{algebra-definition-N}
\begin{reference}
\cite[Chapter 0, Definition 23.1.1]{EGA}
\end{reference}
Soit $R$ un anneau intègre de corps des fractions $K$.
\begin{enumerate}
\item Nous disons que $R$ est {\it N-1} si la clôture intégrale de $R$ dans $K$
est un $R$-module de type fini.
\item Nous disons que $R$ est {\it N-2} ou {\it japonais} si, pour toute extension
finie $L/K$ de corps, la clôture intégrale de $R$ dans $L$
est finie sur $R$.
\end{enumerate}
\end{definition}

\noindent
Ces notions présentent surtout un intérêt pour les anneaux noethériens,
mais voici un exemple non noethérien.

\begin{example}
\label{algebra-example-Japanese-not-Noetherian}
Soit $k$ un corps. L'anneau intègre $R = k[x_1, x_2, x_3, \ldots]$ est N-2,
mais n'est pas noethérien. La raison en est la suivante. Supposons que $R \subset L$
et que le corps $L$ soit une extension finie du corps des fractions de $R$.
Il existe alors un entier $n$ tel que $L$ provienne d'une extension finie
$L_0/k(x_1, \ldots, x_n)$ par adjonction
des éléments (transcendants) $x_{n + 1}, x_{n + 2}$, etc.
Soit $S_0$ la clôture intégrale
de $k[x_1, \ldots, x_n]$ dans $L_0$. D'après la
proposition \ref{algebra-proposition-ubiquity-nagata} ci-dessous,
$S_0$ est fini sur $k[x_1, \ldots, x_n]$.
De plus, la clôture intégrale de $R$ dans $L$ est
$S = S_0[x_{n + 1}, x_{n + 2}, \ldots]$ (utiliser le
lemme \ref{algebra-lemma-polynomial-domain-normal}) et est
donc finie sur $R$. Le même argument s'applique à
$R = \mathbf{Z}[x_1, x_2, x_3, \ldots]$.
\end{example}

\begin{lemma}
\label{algebra-lemma-localize-N}
Soit $R$ un anneau intègre.
Si $R$ est N-1, toute localisation de $R$ l'est aussi.
Il en va de même pour N-2.
\end{lemma}

\begin{proof}
Ces assertions résultent du fait que la formation de la clôture intégrale commute
à la localisation ; voir le lemme \ref{algebra-lemma-integral-closure-localize}.
\end{proof}

\begin{lemma}
\label{algebra-lemma-Japanese-local}
Soit $R$ un anneau intègre. Soient $f_1, \ldots, f_n \in R$ engendrant
l'idéal unité. Si chacun des anneaux intègres $R_{f_i}$ est N-1, alors $R$ l'est aussi.
Il en va de même pour N-2.
\end{lemma}

\begin{proof}
Supposons que $R_{f_i}$ soit N-2 (ou N-1).
Soit $L$ une extension finie du corps des fractions de $R$ (égale à
ce corps des fractions dans le cas N-1). Soit $S$ la clôture
intégrale de $R$ dans $L$. Le lemme \ref{algebra-lemma-integral-closure-localize}
montre que $S_{f_i}$ est la clôture intégrale de $R_{f_i}$ dans $L$.
Par hypothèse, $S_{f_i}$ est donc fini sur $R_{f_i}$.
Le lemme \ref{algebra-lemma-cover} montre alors que $S$ est fini sur $R$.
\end{proof}

\begin{lemma}
\label{algebra-lemma-quasi-finite-over-Noetherian-japanese}
Soit $R$ un anneau intègre. Soit $R \subset S$ une extension quasi-finie d'anneaux intègres
(par exemple finie). Supposons que $R$ soit N-2 et noethérien. Alors $S$ est N-2.
\end{lemma}

\begin{proof}
Soit $L/K$ l'extension induite des corps des fractions.
Remarquons qu'il s'agit d'une extension finie de corps (par exemple d'après le
lemme \ref{algebra-lemma-isolated-point-fibre} (2)
appliqué à la fibre $S \otimes_R K$, et la définition d'un
morphisme d'anneaux quasi-fini).
Soit $S'$ la clôture intégrale de $R$ dans $S$.
Alors $S'$ est contenu dans la clôture intégrale de $R$ dans $L$,
qui est finie sur $R$ par hypothèse. Comme $R$ est noethérien, cela
implique que $S'$ est fini sur $R$.
D'après le lemme \ref{algebra-lemma-quasi-finite-open-integral-closure},
il existe des éléments $g_1, \ldots, g_n \in S'$ tels
que $S'_{g_i} \cong S_{g_i}$ et que $g_1, \ldots, g_n$
engendrent l'idéal unité de $S$. Il suffit donc de montrer que
$S'$ est N-2, d'après les lemmes \ref{algebra-lemma-localize-N} et \ref{algebra-lemma-Japanese-local}.
Nous nous sommes ainsi ramenés au cas où $S$ est fini sur $R$.

\medskip\noindent
Supposons que $R \subset S$ vérifie les hypothèses du lemme et, de plus,
que $S$ soit fini sur $R$. Soit $M$ une extension finie du corps
des fractions de $S$. Alors $M$ est aussi une extension de corps finie
sur $K$, et nous en déduisons que la clôture intégrale $T$ de $R$ dans
$M$ est finie sur $R$. Le lemme \ref{algebra-lemma-integral-closure-transitive}
montre que $T$ est aussi la clôture intégrale de $S$ dans $M$, et le résultat suit du
lemme \ref{algebra-lemma-integral-permanence}.
\end{proof}

\begin{lemma}
\label{algebra-lemma-Laurent-ring-N-1}
Soit $R$ un anneau intègre noethérien.
Si $R[z, z^{-1}]$ est N-1, alors $R$ l'est aussi.
\end{lemma}

\begin{proof}
Soit $R'$ la clôture intégrale de $R$ dans son corps des fractions $K$.
Soit $S'$ la clôture intégrale de $R[z, z^{-1}]$ dans son corps des fractions.
Manifestement, $R' \subset S'$.
Comme $K[z, z^{-1}]$ est un anneau intègre normal, on a $S' \subset K[z, z^{-1}]$.
Supposons que $f_1, \ldots, f_n \in S'$ engendrent $S'$ comme $R[z, z^{-1}]$-module.
Écrivons $f_i = \sum a_{ij}z^j$ (somme finie), avec $a_{ij} \in K$.
Pour tout $x \in R'$, on peut écrire
$$
x = \sum h_i f_i
$$
avec $h_i \in R[z, z^{-1}]$. On voit ainsi que $R'$ est contenu dans le
sous-$R$-module de type fini $\sum Ra_{ij} \subset K$. Comme $R$ est noethérien,
nous en concluons que $R'$ est un $R$-module de type fini.
\end{proof}

\begin{lemma}
\label{algebra-lemma-finite-extension-N-2}
Soit $R$ un anneau intègre noethérien, et soit $R \subset S$ une
extension finie d'anneaux intègres. Si $S$ est N-1, alors $R$ l'est aussi.
Si $S$ est N-2, alors $R$ l'est aussi.
\end{lemma}

\begin{proof}
Démonstration omise. (Indication : les clôtures intégrales de $R$ dans des corps d'extension
sont contenues dans les clôtures intégrales de $S$ dans ces corps d'extension.)
\end{proof}

\begin{lemma}
\label{algebra-lemma-Noetherian-normal-domain-finite-separable-extension}
Soit $R$ un anneau intègre normal noethérien de corps des fractions $K$.
Soit $L/K$ une extension finie séparable de corps.
Alors la clôture intégrale de $R$ dans $L$ est finie sur $R$.
\end{lemma}

\begin{proof}
Considérons la forme trace
(Corps, définition \ref{fields-definition-trace-pairing})
$$
L \times L \longrightarrow K,
\quad (x, y) \longmapsto \langle x, y\rangle := \text{Trace}_{L/K}(xy).
$$
Comme $L/K$ est séparable, cette forme est non dégénérée
(Corps, lemme \ref{fields-lemma-separable-trace-pairing}).
De plus, si $x \in L$ est entier sur $R$, alors
$\text{Trace}_{L/K}(x)$ appartient à $R$. En effet, le
polynôme minimal de $x$ sur $K$ a ses coefficients dans $R$
(lemme \ref{algebra-lemma-minimal-polynomial-normal-domain}),
et $\text{Trace}_{L/K}(x)$ est un
multiple entier de l'un de ces coefficients
(Corps, lemme \ref{fields-lemma-trace-and-norm-from-minimal-polynomial}).
Choisissons $x_1, \ldots, x_n \in L$ entiers sur $R$
et formant une $K$-base de $L$. Alors la clôture intégrale
$S \subset L$ est contenue dans le $R$-module
$$
M = \{y \in L \mid \langle x_i, y\rangle \in R, \ i = 1, \ldots, n\}
$$
L'algèbre linéaire montre que $M \cong R^{\oplus n}$ comme $R$-module.
Ainsi, le sous-module $S \subset R^{\oplus n}$ est un $R$-module de type fini,
puisque $R$ est noethérien.
\end{proof}

\begin{example}
\label{algebra-example-bad-invariants}
Le lemme \ref{algebra-lemma-Noetherian-normal-domain-finite-separable-extension}
n'est plus valable si l'anneau n'est pas noethérien.
Considérons par exemple l'action de $G = \{+1, -1\}$ sur
$A = \mathbf{C}[x_1, x_2, x_3, \ldots]$, où $-1$ agit en
envoyant $x_i$ sur $-x_i$. L'anneau des invariants $R = A^G$ est
la $\mathbf{C}$-algèbre engendrée par tous les $x_ix_j$. Ainsi,
$R \subset A$ n'est pas une extension finie. Mais $R$ est un anneau intègre normal
de corps des fractions $K = L^G$, le sous-corps des $G$-invariants du corps des fractions
$L$ de $A$. Et, manifestement, $A$ est la clôture intégrale de $R$ dans
$L$.
\end{example}

\noindent
Le lemme suivant peut parfois remplacer le
lemme \ref{algebra-lemma-Noetherian-normal-domain-finite-separable-extension}
dans le cas d'extensions purement inséparables.

\begin{lemma}
\label{algebra-lemma-Noetherian-normal-domain-insep-extension}
Soit $R$ un anneau intègre normal noethérien de corps des fractions $K$
de caractéristique $p > 0$.
Soit $a \in K$ un élément tel qu'il existe une dérivation
$D : R \to R$ vérifiant $D(a) \not = 0$. Alors la clôture intégrale
de $R$ dans $L = K[x]/(x^p - a)$ est finie sur $R$.
\end{lemma}

\begin{proof}
Après avoir remplacé $x$ par $fx$ et $a$ par $f^pa$, pour un certain $f \in R$,
nous pouvons supposer que $a \in R$. Ainsi, $D(a) \in R$ également. Nous allons montrer
par récurrence sur $i \leq p - 1$ que, si
$$
y = a_0 + a_1x + \ldots + a_i x^i,\quad a_j \in K
$$
est entier sur $R$, alors $D(a)^i a_j \in R$. La clôture intégrale
est donc contenue dans le $R$-module libre de type fini ayant pour base
$D(a)^{-p + 1}x^j$, $j = 0, \ldots, p - 1$. Comme $R$ est noethérien,
cela démontre le lemme.

\medskip\noindent
Si $i = 0$, alors $y = a_0$ est entier sur $R$ si et seulement si $a_0 \in R$,
et l'assertion est vraie. Supposons l'assertion vraie pour un certain $i < p - 1$
et supposons que
$$
y = a_0 + a_1x + \ldots + a_{i + 1} x^{i + 1},\quad a_j \in K
$$
soit entier sur $R$. Alors
$$
y^p = a_0^p + a_1^p a + \ldots + a_{i + 1}^pa^{i + 1}
$$
est un élément de $R$ (puisqu'il appartient à $K$ et est entier sur $R$). En appliquant
$D$, nous obtenons que
$$
(a_1^p + 2a_2^p a + \ldots + (i + 1)a_{i + 1}^p a^i)D(a)
$$
appartient à $R$. Il s'ensuit que
$$
D(a)a_1 + 2D(a) a_2 x + \ldots + (i + 1)D(a) a_{i + 1} x^i
$$
est entier sur $R$. Par hypothèse de récurrence, $D(a)^{i + 1}a_j \in R$
pour $j = 1, \ldots, i + 1$. (Nous utilisons ici que $1, \ldots, i + 1$
sont inversibles.) Par conséquent, $D(a)^{i + 1}a_0$ appartient aussi à $R$, car c'est la
différence de $D(a)^{i + 1}y$ et de $\sum_{j > 0} D(a)^{i + 1}a_jx^j$, qui
sont entiers sur $R$ (puisque $x$ est entier sur $R$, car $a \in R$).
\end{proof}

\begin{lemma}
\label{algebra-lemma-domain-char-zero-N-1-2}
Un anneau intègre noethérien dont le corps des fractions est de caractéristique zéro est N-1
si et seulement s'il est N-2 (c'est-à-dire japonais).
\end{lemma}

\begin{proof}
Cela découle immédiatement du
lemme \ref{algebra-lemma-Noetherian-normal-domain-finite-separable-extension},
puisque toute extension de corps en caractéristique zéro est séparable.
\end{proof}

\begin{lemma}
\label{algebra-lemma-domain-char-p-N-1-2}
Soit $R$ un anneau intègre noethérien de corps des fractions $K$ de
caractéristique $p > 0$. Alors $R$ est N-2 si et seulement si,
pour toute extension finie purement inséparable $L/K$, la clôture
intégrale de $R$ dans $L$ est finie sur $R$.
\end{lemma}

\begin{proof}
Supposons que la clôture intégrale de $R$ dans toute extension de corps
finie purement inséparable de $K$ soit finie.
Soit $L/K$ une extension finie quelconque. Il nous faut montrer que la
clôture intégrale de $R$ dans $L$ est finie sur $R$.
Choisissons une extension finie normale de corps $M/K$
contenant $L$. Comme $R$ est noethérien, il suffit de montrer que
la clôture intégrale de $R$ dans $M$ est finie sur $R$.
D'après le lemme \ref{fields-lemma-normal-case} de la partie Corps,
il existe une sous-extension $M/M_{insep}/K$
telle que $M_{insep}/K$ soit purement inséparable et que $M/M_{insep}$
soit séparable. Par hypothèse, la clôture intégrale $R'$ de $R$ dans
$M_{insep}$ est finie sur $R$. D'après le
lemme \ref{algebra-lemma-Noetherian-normal-domain-finite-separable-extension},
la clôture
intégrale $R''$ de $R'$ dans $M$ est finie sur $R'$. Alors $R''$ est fini
sur $R$ d'après le lemme \ref{algebra-lemma-finite-transitive}.
Comme $R''$ est aussi la clôture intégrale
de $R$ dans $M$ (voir le lemme \ref{algebra-lemma-integral-closure-transitive}), le résultat est démontré.
\end{proof}

\begin{lemma}
\label{algebra-lemma-polynomial-ring-N-2}
Soit $R$ un anneau intègre noethérien.
Si $R$ est N-1, alors $R[x]$ est N-1.
Si $R$ est N-2, alors $R[x]$ est N-2.
\end{lemma}

\begin{proof}
Supposons que $R$ soit N-1. Soit $R'$ la clôture intégrale de $R$,
qui est finie sur $R$. Par conséquent, $R'[x]$ est aussi fini sur
$R[x]$. L'anneau $R'[x]$ est normal (voir le
lemme \ref{algebra-lemma-polynomial-domain-normal}), donc N-1.
Cela démontre la première assertion.

\medskip\noindent
Pour la seconde assertion, d'après le lemme \ref{algebra-lemma-finite-extension-N-2},
il suffit de montrer que $R'[x]$ est N-2. Autrement dit, nous pouvons
supposer que $R$ est un anneau intègre normal N-2. En caractéristique zéro,
le lemme \ref{algebra-lemma-domain-char-zero-N-1-2} permet de conclure.
En caractéristique $p > 0$, nous devons montrer que la clôture
intégrale de $R[x]$ dans toute extension finie purement inséparable
$L/K(x)$ est finie, où $K$ est le corps des fractions de $R$. Il
existe une extension finie purement inséparable de corps $L'/K$
et $q = p^e$ tels que $L \subset L'(x^{1/q})$ ; certains détails sont omis.
Comme $R[x]$ est noethérien, il suffit de montrer que la clôture intégrale de $R[x]$
dans $L'(x^{1/q})$ est finie sur $R[x]$. Cette clôture intégrale
est égale à $R'[x^{1/q}]$, où $R \subset R' \subset L'$ est la clôture intégrale
de $R$ dans $L'$.
Comme $R$ est N-2, l'anneau $R'$ est fini sur $R$ et, par suite,
$R'[x^{1/q}]$ est fini sur $R[x]$.
\end{proof}

\begin{lemma}
\label{algebra-lemma-openness-normal-locus}
Soit $R$ un anneau intègre noethérien.
S'il existe un élément non nul $f \in R$ tel que $R_f$ soit normal,
alors
$$
U = \{\mathfrak p \in \Spec(R) \mid R_{\mathfrak p} \text{ est normal}\}
$$
est ouvert dans $\Spec(R)$.
\end{lemma}

\begin{proof}
Il est clair que l'ouvert principal $D(f)$ est contenu dans $U$.
Le critère de normalité de Serre, lemme \ref{algebra-lemma-criterion-normal}, montre que
$\mathfrak p \not \in U$ implique qu'il existe
$\mathfrak q \subset \mathfrak p$ pour lequel
l'un des cas suivants se présente :
\begin{enumerate}
\item Cas I : $\text{depth}(R_{\mathfrak q}) < 2$
et $\dim(R_{\mathfrak q}) \geq 2$, ou bien
\item Cas II : $R_{\mathfrak q}$ n'est pas régulier
et $\dim(R_{\mathfrak q}) = 1$.
\end{enumerate}
Cela signifie en particulier que $R_{\mathfrak q}$ n'est pas
normal, et donc que $f \in \mathfrak q$. Dans le cas I, on a
$\text{depth}(R_{\mathfrak q}) =
\text{depth}(R_{\mathfrak q}/fR_{\mathfrak q}) + 1$.
Un tel idéal premier $\mathfrak q$ est donc exactement un idéal premier associé
immergé de $R/fR$. Dans le cas II, $\mathfrak q$ est un idéal premier associé
de $R/fR$ de hauteur 1. Il existe donc un ensemble fini $E$
de tels premiers $\mathfrak q$ (voir le lemme \ref{algebra-lemma-finite-ass}), et
$$
\Spec(R) \setminus U
=
\bigcup\nolimits_{\mathfrak q \in E} V(\mathfrak q)
$$
comme souhaité.
\end{proof}

\begin{lemma}
\label{algebra-lemma-characterize-N-1}
Soit $R$ un anneau intègre noethérien. Alors $R$ est N-1 si et seulement si les deux
conditions suivantes sont satisfaites :
\begin{enumerate}
\item il existe un élément non nul $f \in R$ tel que $R_f$ soit normal, et
\item pour tout idéal maximal $\mathfrak m \subset R$,
l'anneau local $R_{\mathfrak m}$ est N-1.
\end{enumerate}
\end{lemma}

\begin{proof}
Supposons d'abord que $R$ soit N-1. Soit $R'$ la clôture intégrale de $R$ dans son
corps des fractions $K$. Par hypothèse, on peut trouver $x_1, \ldots, x_n$ dans $R'$
qui engendrent $R'$ comme $R$-module. Comme $R' \subset K$, on peut trouver
des éléments non nuls $f_i \in R$ tels que $f_i x_i \in R$. Alors $R_f \cong R'_f$,
où $f = f_1 \ldots f_n$. Ainsi, $R_f$ est normal et la condition (1) est satisfaite.
La condition (2) résulte du lemme \ref{algebra-lemma-localize-N}.

\medskip\noindent
Supposons (1) et (2). Soit $K$ le corps des fractions de $R$.
Supposons que $R \subset R' \subset K$ soit une extension finie
de $R$ contenue dans $K$. Remarquons que $R_f = R'_f$, puisque
$R_f$ est déjà normal. Le lemme \ref{algebra-lemma-openness-normal-locus} montre donc
que l'ensemble des premiers
$\mathfrak p' \in \Spec(R')$ tels que $R'_{\mathfrak p'}$ ne soit pas normal
est fermé dans $\Spec(R')$. Comme $\Spec(R') \to \Spec(R)$
est fermé, l'image de cet ensemble est fermée dans $\Spec(R)$.
Pour un tel anneau $R'$, notons $Z_{R'} \subset \Spec(R)$ cette image.

\medskip\noindent
Choisissons un idéal maximal $\mathfrak m \subset R$.
Soit $R_{\mathfrak m} \subset R_{\mathfrak m}'$ la clôture
intégrale de l'anneau local dans $K$. Par hypothèse, il s'agit
d'une extension finie d'anneaux. D'après le lemme \ref{algebra-lemma-integral-closure-localize},
on peut trouver un nombre fini
d'éléments $x_1, \ldots, x_n \in K$ entiers sur $R$ tels que
$R_{\mathfrak m}'$ soit engendré par $x_1, \ldots, x_n$ sur $R_{\mathfrak m}$.
Soit $R' = R[x_1, \ldots, x_n] \subset K$. Avec ce choix, il est clair
que $\mathfrak m \not \in Z_{R'}$.

\medskip\noindent
Comme $\Spec(R)$ est quasi-compact, ce qui précède montre que nous pouvons
trouver une famille finie $R \subset R'_i \subset K$ telle que
$\bigcap Z_{R'_i} = \emptyset$. Soit $R'$ le sous-anneau de $K$
engendré par tous ces sous-anneaux. Il est fini sur $R$. De plus, $Z_{R'} = \emptyset$.
En effet, tout premier $\mathfrak p'$ est au-dessus d'un premier $\mathfrak p'_i$
tel que $(R'_i)_{\mathfrak p'_i}$ soit normal. Cela implique
que $R'_{\mathfrak p'} = (R'_i)_{\mathfrak p'_i}$ est lui aussi normal.
Ainsi, $R'$ est normal ; autrement dit,
$R'$ est la clôture intégrale de $R$ dans $K$.
\end{proof}

\begin{lemma}[Tate]
\label{algebra-lemma-tate-japanese}
\begin{reference}
\cite[Theorem 23.1.3]{EGA}
\end{reference}
Soit $R$ un anneau.
Soit $x \in R$.
Supposons que
\begin{enumerate}
\item $R$ soit un anneau intègre normal noethérien,
\item $R/xR$ soit un anneau intègre et N-2,
\item $R \cong \lim_n R/x^nR$ soit complet pour la topologie $x$-adique.
\end{enumerate}
Alors $R$ est N-2.
\end{lemma}

\begin{proof}
Nous pouvons supposer que $x \not = 0$, car sinon le lemme est trivial.
Soit $K$ le corps des fractions de $R$. Si la caractéristique de $K$
est nulle, le lemme résulte de (1) ; voir le
lemme \ref{algebra-lemma-domain-char-zero-N-1-2}. Nous pouvons donc supposer
que la caractéristique de $K$ est $p > 0$ et appliquer le
lemme \ref{algebra-lemma-domain-char-p-N-1-2}. Ainsi, pour toute extension de corps $L/K$
finie et purement inséparable, il nous faut donc montrer
que la clôture intégrale $S$ de $R$ dans $L$ est finie sur $R$.

\medskip\noindent
Soit $q$ une puissance de $p$ telle que $L^q \subset K$.
En agrandissant $L$ si nécessaire, nous pouvons supposer qu'il existe
un élément $y \in L$ tel que $y^q = x$. Comme $R \to S$
induit un homéomorphisme entre les spectres (voir le lemme \ref{algebra-lemma-p-ring-map}),
il existe un unique idéal premier $\mathfrak q \subset S$ au-dessus
de l'idéal premier $\mathfrak p = xR$. Il est clair que
$$
\mathfrak q = \{f \in S \mid f^q \in \mathfrak p\} = yS
$$
puisque $y^q = x$. Remarquons que $R_{\mathfrak p}$ est un anneau de
valuation discrète d'après le lemme \ref{algebra-lemma-characterize-dvr}. D'autre part,
$S_{\mathfrak q}$ est noethérien d'après le théorème de Krull-Akizuki
(lemme \ref{algebra-lemma-krull-akizuki}). Nous en concluons alors que
$S_{\mathfrak q}$ est un anneau de valuation discrète, à nouveau d'après le
lemme \ref{algebra-lemma-characterize-dvr}.
Le lemme \ref{algebra-lemma-finite-extension-residue-fields-dimension-1} montre que
$\kappa(\mathfrak q)/\kappa(\mathfrak p)$ est
une extension finie de corps. La clôture intégrale
$S' \subset \kappa(\mathfrak q)$ de $R/xR$ est donc finie sur
$R/xR$ par l'hypothèse (2). Comme $S/yS \subset S'$, cela implique
que $S/yS$ est fini sur $R$. Remarquons que $S/y^nS$ possède une filtration
finie dont les sous-quotients sont les modules
$y^iS/y^{i + 1}S \cong S/yS$. On voit donc que chaque $S/y^nS$
est fini sur $R$. En particulier, $S/xS$ est fini sur $R$.
De plus, il est clair que $\bigcap x^nS = (0)$, car la puissance $q$-ième d'un élément
de l'intersection appartient à $\bigcap x^nR = (0)$
(lemme \ref{algebra-lemma-intersect-powers-ideal-module-zero}).
Nous pouvons donc appliquer le lemme \ref{algebra-lemma-finite-over-complete-ring} pour conclure
que $S$ est fini sur $R$, ce qui termine la démonstration.
\end{proof}

\begin{lemma}
\label{algebra-lemma-power-series-over-N-2}
Soit $R$ un anneau.
Si $R$ est un anneau intègre noethérien et N-2, alors $R[[x]]$ l'est aussi.
\end{lemma}

\begin{proof}
Remarquons que $R[[x]]$ est noethérien d'après le
lemme \ref{algebra-lemma-Noetherian-power-series}.
Soit $R' \supset R$ la clôture intégrale de $R$ dans son corps des
fractions. Comme $R$ est N-2, elle est finie sur $R$. Par conséquent, $R'[[x]]$
est fini sur $R[[x]]$. D'après le
lemme \ref{algebra-lemma-power-series-over-Noetherian-normal-domain},
$R'[[x]]$ est un anneau intègre normal.
Appliquons le lemme \ref{algebra-lemma-tate-japanese} à
l'élément $x \in R'[[x]]$ : l'anneau $R'[[x]]$ est N-2. Le
lemme \ref{algebra-lemma-finite-extension-N-2} montre alors que $R[[x]]$ est N-2.
\end{proof}





\section{Anneaux de Nagata}
\label{algebra-section-nagata}

\noindent
Voici la définition.

\begin{definition}
\label{algebra-definition-nagata}
Soit $R$ un anneau.
\begin{enumerate}
\item Nous disons que $R$ est {\it universellement japonais} si, pour tout morphisme
d'anneaux de type fini $R \to S$ avec $S$ intègre, l'anneau $S$ est N-2
(c'est-à-dire japonais).
\item Nous disons que $R$ est un {\it anneau de Nagata} si $R$ est noethérien et si,
pour tout idéal premier $\mathfrak p$, l'anneau $R/\mathfrak p$ est N-2.
\end{enumerate}
\end{definition}

\noindent
Il est clair qu'un anneau noethérien universellement japonais est un anneau de Nagata.
Notre objectif est de montrer qu'un anneau de Nagata est universellement japonais. Ce n'est
pas du tout évident et cela demande un certain travail. Mais commençons par un lemme
utile.

\begin{lemma}
\label{algebra-lemma-nagata-in-reduced-finite-type-finite-integral-closure}
Soit $R$ un anneau de Nagata.
Soit $R \to S$ un morphisme essentiellement de type fini, avec $S$ réduit.
Alors la clôture intégrale $A$ de $R$ dans $S$ est finie sur $R$.
\end{lemma}

\begin{proof}
Comme $S$ est essentiellement de type fini sur $R$, il est noethérien et
possède un nombre fini d'idéaux premiers minimaux $\mathfrak q_1, \ldots, \mathfrak q_m$ ;
voir le lemme \ref{algebra-lemma-Noetherian-irreducible-components}.
Puisque $S$ est réduit, on a $S \subset \prod S_{\mathfrak q_i}$,
et chacun des $S_{\mathfrak q_i} = K_i$ est un corps ; voir les
lemmes \ref{algebra-lemma-total-ring-fractions-no-embedded-points}
et \ref{algebra-lemma-minimal-prime-reduced-ring}.
Il suffit de montrer que la clôture intégrale
$A_i'$ de $R$ dans chacun des $K_i$ est finie sur $R$.
Cela suffit parce que $R$ est noethérien et que $A \subset \prod A_i'$.
Soit $\mathfrak p_i \subset R$ l'idéal premier de $R$
correspondant à $\mathfrak q_i$.
Comme $S$ est essentiellement de type fini sur $R$, on voit que
$K_i = S_{\mathfrak q_i} = \kappa(\mathfrak q_i)$ est une extension de corps de type
fini de $\kappa(\mathfrak p_i)$. Par conséquent, la clôture algébrique
$L_i$ de $\kappa(\mathfrak p_i)$ dans $K_i$
est finie sur $\kappa(\mathfrak p_i)$ ; voir
Corps, lemme \ref{fields-lemma-algebraic-closure-in-finitely-generated}.
Il est clair que $A_i'$ est la clôture intégrale de $R/\mathfrak p_i$
dans $L_i$, et le résultat découle donc de la définition d'un anneau de Nagata.
\end{proof}

\begin{lemma}
\label{algebra-lemma-check-universally-japanese}
Soit $R$ un anneau.
Pour vérifier que $R$ est universellement japonais, il suffit de montrer ceci :
si $R \to S$ est de type fini et si $S$ est intègre, alors $S$ est N-1.
\end{lemma}

\begin{proof}
Supposons en effet satisfaite la condition du lemme.
Soit $R \to S$ un morphisme d'anneaux de type fini, avec $S$ intègre.
Soit $L$ une extension finie du corps des fractions de $S$.
Il existe alors une extension finie d'anneaux $S \subset S' \subset L$
telle que $L$ soit le corps des fractions de $S'$.
Par hypothèse, $S'$ est N-1, et la clôture intégrale
$S''$ de $S'$ dans $L$ est donc finie sur $S'$. Ainsi, $S''$ est fini
sur $S$ (lemme \ref{algebra-lemma-finite-transitive})
et $S''$ est la clôture intégrale de $S$ dans $L$
(lemme \ref{algebra-lemma-integral-closure-transitive}).
Nous en concluons que $R$ est universellement japonais.
\end{proof}

\begin{lemma}
\label{algebra-lemma-universally-japanese}
Si $R$ est universellement japonais, toute algèbre essentiellement de type fini
sur $R$ est universellement japonaise.
\end{lemma}

\begin{proof}
Le cas d'une algèbre de type fini sur $R$ résulte immédiatement de
la définition. Le cas général s'obtient en appliquant le
lemme \ref{algebra-lemma-localize-N}.
\end{proof}

\begin{lemma}
\label{algebra-lemma-quasi-finite-over-nagata}
Soit $R$ un anneau de Nagata.
Si $R \to S$ est un morphisme d'anneaux quasi-fini (par exemple fini),
alors $S$ est lui aussi un anneau de Nagata.
\end{lemma}

\begin{proof}
Remarquons d'abord que $S$ est noethérien, car $R$ est noethérien et qu'un morphisme
d'anneaux quasi-fini est de type fini.
Soit $\mathfrak q \subset S$ un idéal premier et posons
$\mathfrak p = R \cap \mathfrak q$. Alors
$R/\mathfrak p \subset S/\mathfrak q$ est quasi-fini, et
nous en concluons que $S/\mathfrak q$ est N-2 d'après le
lemme \ref{algebra-lemma-quasi-finite-over-Noetherian-japanese},
comme souhaité.
\end{proof}

\begin{lemma}
\label{algebra-lemma-nagata-localize}
Toute localisation d'un anneau de Nagata est un anneau de Nagata.
\end{lemma}

\begin{proof}
Cela résulte immédiatement du lemme \ref{algebra-lemma-localize-N}.
\end{proof}

\begin{lemma}
\label{algebra-lemma-nagata-local}
Soit $R$ un anneau. Soient $f_1, \ldots, f_n \in R$ engendrant
l'idéal unité.
\begin{enumerate}
\item Si chacun des $R_{f_i}$ est universellement japonais, alors $R$ l'est aussi.
\item Si chacun des $R_{f_i}$ est un anneau de Nagata, alors $R$ l'est aussi.
\end{enumerate}
\end{lemma}

\begin{proof}
Soit $\varphi : R \to S$ un morphisme d'anneaux de type fini tel que $S$ soit intègre.
Alors $\varphi(f_1), \ldots, \varphi(f_n)$ engendrent l'idéal unité
de $S$. Par conséquent, si chacun des $S_{f_i} = S_{\varphi(f_i)}$ est N-1, alors
$S$ l'est aussi ; voir le lemme \ref{algebra-lemma-Japanese-local}. Cela démontre (1).

\medskip\noindent
Si chacun des $R_{f_i}$ est un anneau de Nagata, chacun des $R_{f_i}$ est noethérien,
et $R$ est donc noethérien ; voir le lemme \ref{algebra-lemma-cover}. De plus, si
$\mathfrak p \subset R$ est un idéal premier, chacun des anneaux
$R_{f_i}/\mathfrak pR_{f_i} = (R/\mathfrak p)_{f_i}$ est N-2,
et nous en concluons que $R/\mathfrak p$ est N-2 d'après le
lemme \ref{algebra-lemma-Japanese-local}. Cela démontre (2).
\end{proof}

\begin{lemma}
\label{algebra-lemma-Noetherian-complete-local-Nagata}
Tout anneau local noethérien complet est un anneau de Nagata.
\end{lemma}

\begin{proof}
Soit $R$ un anneau local noethérien complet.
Soit $\mathfrak p \subset R$ un idéal premier.
Alors $R/\mathfrak p$ est également un anneau local noethérien complet ;
voir le lemme \ref{algebra-lemma-quotient-complete-local}.
Il suffit donc de montrer qu'un anneau intègre local noethérien
complet $R$ est N-2. D'après les
lemmes \ref{algebra-lemma-quasi-finite-over-Noetherian-japanese}
et \ref{algebra-lemma-complete-local-Noetherian-domain-finite-over-regular},
nous nous ramenons au cas $R = k[[X_1, \ldots, X_d]]$, où $k$ est un corps, ou au cas
$R = \Lambda[[X_1, \ldots, X_d]]$, où $\Lambda$ est un anneau de Cohen.

\medskip\noindent
Dans le cas $k[[X_1, \ldots, X_d]]$, le lemme
\ref{algebra-lemma-power-series-over-N-2} nous ramène à l'assertion qu'un corps est N-2. C'est clair.
Dans le cas $\Lambda[[X_1, \ldots, X_d]]$, nous nous ramenons à l'assertion
qu'un anneau de Cohen $\Lambda$ est N-2. En appliquant une fois encore le lemme
\ref{algebra-lemma-tate-japanese} avec $x = p \in \Lambda$, nous nous ramenons encore
au cas d'un corps. Le résultat est donc démontré.
\end{proof}

\begin{definition}
\label{algebra-definition-analytically-unramified}
Soit $(R, \mathfrak m)$ un anneau local noethérien.
Nous disons que $R$ est {\it analytiquement non ramifié} si son complété
$R^\wedge = \lim_n R/\mathfrak m^n$ est réduit.
On dit qu'un idéal premier $\mathfrak p \subset R$ est
{\it analytiquement non ramifié} si $R/\mathfrak p$ est analytiquement
non ramifié.
\end{definition}

\noindent
À ce stade, nous savons que
les assertions suivantes sont vraies pour tout anneau local noethérien $R$ :
le morphisme $R \to R^\wedge$ est un homomorphisme local fidèlement plat d'anneaux
(lemme \ref{algebra-lemma-completion-faithfully-flat}).
Le complété $R^\wedge$ est noethérien
(lemme \ref{algebra-lemma-completion-Noetherian})
et complet (lemme \ref{algebra-lemma-completion-complete}).
Le complété $R^\wedge$ est donc un anneau de Nagata
(lemme \ref{algebra-lemma-Noetherian-complete-local-Nagata}).
En outre, nous avons vu dans la section \ref{algebra-section-cohen-structure-theorem}
que $R^\wedge$ est
un quotient d'un anneau local régulier
(théorème \ref{algebra-theorem-cohen-structure-theorem}), et qu'il est donc
universellement caténaire
(remarque \ref{algebra-remark-Noetherian-complete-local-ring-universally-catenary}).

\begin{lemma}
\label{algebra-lemma-analytically-unramified-easy}
Soit $(R, \mathfrak m)$ un anneau local noethérien.
\begin{enumerate}
\item Si $R$ est analytiquement non ramifié, alors $R$ est réduit.
\item Si $R$ est analytiquement non ramifié, alors tout idéal premier minimal de
$R$ est analytiquement non ramifié.
\item Si $R$ est réduit, d'idéaux premiers minimaux
$\mathfrak q_1, \ldots, \mathfrak q_t$, et si chaque $\mathfrak q_i$
est analytiquement non ramifié, alors $R$ est analytiquement non ramifié.
\item Si $R$ est analytiquement non ramifié, alors la clôture intégrale
de $R$ dans son anneau total des fractions $Q(R)$ est finie sur $R$.
\item Si $R$ est intègre et analytiquement non ramifié, alors $R$ est N-1.
\end{enumerate}
\end{lemma}

\begin{proof}
Dans cette démonstration, nous utiliserons les remarques qui suivent immédiatement la
définition \ref{algebra-definition-analytically-unramified}.
Comme $R \to R^\wedge$ est un homomorphisme local fidèlement plat d'anneaux,
il est injectif, d'où (1).

\medskip\noindent
Soit $\mathfrak q$ un idéal premier minimal de $R$, et supposons $R$
analytiquement non ramifié.
Alors $\mathfrak q$ est un idéal premier associé
à $R$ (voir la
proposition \ref{algebra-proposition-minimal-primes-associated-primes}).
Il existe donc $f \in R$
tel que $\{x \in R \mid fx = 0\} = \mathfrak q$.
Remarquons que $(R/\mathfrak q)^\wedge = R^\wedge/\mathfrak q^\wedge$
et que $\{x \in R^\wedge \mid fx = 0\} = \mathfrak q^\wedge$,
car la complétion est exacte (lemme \ref{algebra-lemma-completion-flat}).
Si $x \in R^\wedge$ est tel
que $x^2 \in \mathfrak q^\wedge$, alors $fx^2 = 0$, donc
$(fx)^2 = 0$, donc $fx = 0$, donc $x \in \mathfrak q^\wedge$.
Ainsi, $\mathfrak q$ est analytiquement non ramifié et (2) est satisfaite.

\medskip\noindent
Supposons $R$ réduit, d'idéaux premiers minimaux
$\mathfrak q_1, \ldots, \mathfrak q_t$, et chaque $\mathfrak q_i$
analytiquement non ramifié. Alors
$R \to R/\mathfrak q_1 \times \ldots \times R/\mathfrak q_t$ est
injectif. Comme la complétion est exacte (voir le lemme \ref{algebra-lemma-completion-flat}),
nous voyons que
$R^\wedge \subset (R/\mathfrak q_1)^\wedge \times \ldots \times
(R/\mathfrak q_t)^\wedge$. L'assertion (3) est donc claire.

\medskip\noindent
Supposons $R$ analytiquement non ramifié.
Soient $\mathfrak p_1, \ldots, \mathfrak p_s$ les idéaux premiers minimaux
de $R^\wedge$. On a alors
$$
Q(R^\wedge) =
R^\wedge_{\mathfrak p_1} \times \ldots \times R^\wedge_{\mathfrak p_s}
$$
où chaque $R^\wedge_{\mathfrak p_i}$ est un corps,
puisque $R^\wedge$ est réduit (voir le
lemme \ref{algebra-lemma-total-ring-fractions-no-embedded-points}).
Par conséquent, la clôture intégrale $S$ de $R^\wedge$
dans $Q(R^\wedge)$ est $S = S_1 \times \ldots \times S_s$, où
$S_i$ est la clôture intégrale de $R^\wedge/\mathfrak p_i$ dans son corps des
fractions. D'après le lemme \ref{algebra-lemma-Noetherian-complete-local-Nagata}, l'anneau
$S_i$ est fini sur $R^\wedge/\mathfrak p_i$. Ainsi,
$S$ est fini sur $R^\wedge$.
Notons $R'$ la clôture intégrale de $R$ dans $Q(R)$.
Comme $R \to R^\wedge$ est plat, on a
$R' \otimes_R R^\wedge \subset Q(R) \otimes_R R^\wedge \subset Q(R^\wedge)$.
De plus, $R' \otimes_R R^\wedge$ est entier sur $R^\wedge$
(lemme \ref{algebra-lemma-base-change-integral}).
Ainsi, $R' \otimes_R R^\wedge \subset S$ est un sous-$R^\wedge$-module.
Comme $R^\wedge$ est noethérien, c'est un $R^\wedge$-module de type fini.
Nous pouvons donc trouver $f_1, \ldots, f_n \in R'$ tels que
$R' \otimes_R R^\wedge$ soit engendré par les éléments $f_i \otimes 1$
comme $R^\wedge$-module.
La fidèle platitude montre que $R'$ est engendré par $f_1, \ldots, f_n$
comme $R$-module. Cela démontre (4).

\medskip\noindent
L'assertion (5) est un cas particulier de l'assertion (4).
\end{proof}

\begin{lemma}
\label{algebra-lemma-codimension-1-analytically-unramified}
Soit $R$ un anneau local noethérien.
Soit $\mathfrak p \subset R$ un idéal premier.
Supposons que
\begin{enumerate}
\item $R_{\mathfrak p}$ soit un anneau de valuation discrète, et que
\item $\mathfrak p$ soit analytiquement non ramifié.
\end{enumerate}
Alors, pour tout idéal premier associé $\mathfrak q$ de $R^\wedge/\mathfrak pR^\wedge$,
l'anneau local $(R^\wedge)_{\mathfrak q}$ est un anneau de valuation discrète.
\end{lemma}

\begin{proof}
L'hypothèse (2) signifie que $R^\wedge/\mathfrak pR^\wedge$ est un anneau réduit.
Par conséquent, un idéal premier associé $\mathfrak q \subset R^\wedge$
de $R^\wedge/\mathfrak pR^\wedge$
est la même chose qu'un idéal premier minimal au-dessus de $\mathfrak pR^\wedge$.
On voit en particulier que l'idéal maximal de $(R^\wedge)_{\mathfrak q}$
est $\mathfrak p(R^\wedge)_{\mathfrak q}$.
Choisissons $x \in R$ tel que $xR_{\mathfrak p} = \mathfrak pR_{\mathfrak p}$.
D'après ce qui précède, $x \in (R^\wedge)_{\mathfrak q}$ engendre
l'idéal maximal. Comme $R \to R^\wedge$ est fidèlement plat, on voit que
$x$ est un non-diviseur de zéro dans $(R^\wedge)_{\mathfrak q}$.
Le résultat s'ensuit.
\end{proof}

\begin{lemma}
\label{algebra-lemma-criterion-analytically-unramified}
Soit $(R, \mathfrak m)$ un anneau intègre local noethérien.
Soit $x \in \mathfrak m$. Supposons que
\begin{enumerate}
\item $x \not = 0$,
\item $R/xR$ n'ait aucun idéal premier associé immergé, et que
\item pour tout idéal premier associé $\mathfrak p \subset R$
à $R/xR$, on ait
\begin{enumerate}
\item l'anneau local $R_{\mathfrak p}$ est régulier, et
\item $\mathfrak p$ est analytiquement non ramifié.
\end{enumerate}
\end{enumerate}
Alors $R$ est analytiquement non ramifié.
\end{lemma}

\begin{proof}
Soient $\mathfrak p_1, \ldots, \mathfrak p_t$ les idéaux premiers associés
au $R$-module $R/xR$. Puisque $R/xR$ n'a aucun idéal premier associé immergé, on
voit que chaque $\mathfrak p_i$ est de hauteur $1$ et est un idéal premier
minimal au-dessus de $(x)$.
Pour chaque $i$, soient $\mathfrak q_{i1}, \ldots, \mathfrak q_{is_i}$
les idéaux premiers associés au $R^\wedge$-module
$R^\wedge/\mathfrak p_iR^\wedge$.
D'après le lemme \ref{algebra-lemma-codimension-1-analytically-unramified},
on voit que $(R^\wedge)_{\mathfrak q_{ij}}$ est régulier.
D'après le lemme \ref{algebra-lemma-bourbaki}, on a
$$
\text{Ass}_{R^\wedge}(R^\wedge/xR^\wedge)
=
\bigcup\nolimits_{\mathfrak p \in \text{Ass}_R(R/xR)}
\text{Ass}_{R^\wedge}(R^\wedge/\mathfrak pR^\wedge)
=
\{\mathfrak q_{ij}\}.
$$
Soit $y \in R^\wedge$ tel que $y^2 = 0$.
Comme $(R^\wedge)_{\mathfrak q_{ij}}$ est régulier, et donc intègre
(lemme \ref{algebra-lemma-regular-domain}),
on voit que $y$ s'envoie sur zéro dans $(R^\wedge)_{\mathfrak q_{ij}}$.
Il s'ensuit que $y$ s'envoie sur zéro dans $R^\wedge/xR^\wedge$ d'après le
lemme \ref{algebra-lemma-zero-at-ass-zero}.
Ainsi, $y = xy'$. Comme $x$ est un non-diviseur de zéro (puisque $R \to R^\wedge$ est plat),
on voit que $(y')^2 = 0$. Nous en concluons que
$y \in \bigcap x^nR^\wedge = (0)$
(lemme \ref{algebra-lemma-intersect-powers-ideal-module-zero}).
\end{proof}

\begin{lemma}
\label{algebra-lemma-local-nagata-domain-analytically-unramified}
Soit $(R, \mathfrak m)$ un anneau local.
Si $R$ est noethérien, intègre et de Nagata, alors $R$ est
analytiquement non ramifié.
\end{lemma}

\begin{proof}
Raisonnons par récurrence sur $\dim(R)$.
Le cas $\dim(R) = 0$ est trivial. Supposons donc $\dim(R) = d$ et le
lemme vrai pour tous les anneaux intègres noethériens de Nagata de dimension $< d$.

\medskip\noindent
Soit $R \subset S$ la clôture intégrale
de $R$ dans le corps des fractions de $R$. Par hypothèse, $S$ est un
$R$-module de type fini. D'après le lemme \ref{algebra-lemma-quasi-finite-over-nagata},
$S$ est un anneau de Nagata. D'après le lemme \ref{algebra-lemma-integral-sub-dim-equal}, on a
$\dim(R) = \dim(S)$.
Soient $\mathfrak m_1, \ldots, \mathfrak m_t$ les idéaux
maximaux de $S$. Chacun d'eux est au-dessus de l'idéal maximal $\mathfrak m$
de $R$. De plus,
$$
(\mathfrak m_1 \cap \ldots \cap \mathfrak m_t)^n \subset \mathfrak mS
$$
pour $n$ assez grand, puisque $S/\mathfrak mS$ est artinien.
D'après le lemme \ref{algebra-lemma-completion-flat}, $R^\wedge \to S^\wedge$
est injectif et, en combinant le théorème chinois des restes,
lemme \ref{algebra-lemma-chinese-remainder}, avec le
lemme \ref{algebra-lemma-change-ideal-completion}, on obtient
$S^\wedge = \prod S^\wedge_i$, où $S^\wedge_i$
est le complété de $S$ pour la topologie $\mathfrak m_i$-adique.
Il suffit donc de montrer que $S_{\mathfrak m_i}$ est analytiquement non ramifié.
Autrement dit, nous nous sommes ramenés au cas où $R$ est un anneau intègre
normal noethérien de Nagata.

\medskip\noindent
Supposons que $R$ soit un anneau intègre de Nagata, normal, noethérien et local.
Choisissons un élément non nul $x \in \mathfrak m$ de l'idéal maximal.
Nous allons appliquer le lemme \ref{algebra-lemma-criterion-analytically-unramified}.
Nous devons vérifier les conditions (1), (2), (3)(a) et (3)(b).
La condition (1) est claire.
L'anneau $R/xR$ n'a aucun idéal premier associé immergé d'après le
lemme \ref{algebra-lemma-normal-domain-intersection-localizations-height-1}.
La condition (2) est donc satisfaite. Le même lemme montre aussi que tout idéal
premier associé $\mathfrak p$ à $R/xR$ est de hauteur $1$.
Ainsi, $R_{\mathfrak p}$ est un anneau intègre normal de dimension $1$,
donc régulier (lemme \ref{algebra-lemma-characterize-dvr}). La condition (3)(a) est ainsi satisfaite.
Enfin, la condition (3)(b) résulte de l'hypothèse de récurrence, puisque
$R/\mathfrak p$ est un anneau de Nagata (d'après le lemme \ref{algebra-lemma-quasi-finite-over-nagata}
ou directement par définition).
Nous en concluons que $R$ est analytiquement non ramifié.
\end{proof}

\begin{lemma}
\label{algebra-lemma-local-nagata-and-analytically-unramified}
Soit $(R, \mathfrak m)$ un anneau local noethérien. Les conditions suivantes
sont équivalentes :
\begin{enumerate}
\item $R$ est un anneau de Nagata,
\item pour tout morphisme fini $R \to S$, avec $S$ intègre, et tout idéal maximal $\mathfrak m' \subset S$,
l'anneau local $S_{\mathfrak m'}$ est analytiquement non ramifié,
\item pour tout homomorphisme local fini $(R, \mathfrak m) \to (S, \mathfrak m')$
avec $S$ intègre, l'anneau $S$ est analytiquement
non ramifié.
\end{enumerate}
\end{lemma}

\begin{proof}
Supposons que $R$ soit un anneau de Nagata et soient $R \to S$ et $\mathfrak m' \subset S$
comme en (2). Alors $S$ est un anneau de Nagata d'après le lemme \ref{algebra-lemma-quasi-finite-over-nagata}.
Par conséquent, l'anneau local $S_{\mathfrak m'}$ est un anneau de Nagata
(lemme \ref{algebra-lemma-nagata-localize}). Il est donc analytiquement
non ramifié d'après le lemme \ref{algebra-lemma-local-nagata-domain-analytically-unramified}.
Il est clair que (2) implique (3).

\medskip\noindent
Supposons (3) satisfaite. Soit $\mathfrak p \subset R$ un idéal premier et
soit $L/\kappa(\mathfrak p)$ une extension finie de corps.
Pour démontrer (1), nous devons montrer que la clôture intégrale de $R/\mathfrak p$
est finie sur $R/\mathfrak p$. Choisissons $x_1, \ldots, x_n \in L$
engendrant $L$ sur $\kappa(\mathfrak p)$. Pour chaque $i$, soit
$P_i(T) = T^{d_i} + a_{i, 1} T^{d_i - 1} + \ldots + a_{i, d_i}$
le polynôme minimal de $x_i$ sur $\kappa(\mathfrak p)$.
Après avoir remplacé $x_i$ par $f_i x_i$ pour un
$f_i \in R$ convenable, $f_i \not \in \mathfrak p$, nous pouvons supposer
$a_{i, j} \in R/\mathfrak p$. En fait, après avoir encore multiplié
par des éléments de $\mathfrak m$, nous pouvons supposer
$a_{i, j} \in \mathfrak m/\mathfrak p \subset R/\mathfrak p$ pour tous $i, j$.
Cela étant fait, posons $S = R/\mathfrak p[x_1, \ldots, x_n] \subset L$.
Alors $S$ est fini sur $R$, intègre, et $S/\mathfrak m S$ est un quotient
de $R/\mathfrak m[T_1, \ldots, T_n]/(T_1^{d_1}, \ldots, T_n^{d_n})$.
Ainsi, $S$ est local. D'après (3), $S$ est analytiquement non ramifié et, d'après le
lemme \ref{algebra-lemma-analytically-unramified-easy},
sa clôture intégrale $S'$ dans $L$ est finie sur $S$.
Puisque $S'$ est également la clôture intégrale de $R/\mathfrak p$ dans
$L$, le résultat est démontré.
\end{proof}

\noindent
La proposition suivante affirme notamment qu'une algèbre de type
fini sur un anneau de Nagata est un anneau de Nagata.

\begin{proposition}[Nagata]
\label{algebra-proposition-nagata-universally-japanese}
Soit $R$ un anneau. Les conditions suivantes sont équivalentes :
\begin{enumerate}
\item $R$ est un anneau de Nagata,
\item toute $R$-algèbre de type fini est un anneau de Nagata, et
\item $R$ est universellement japonais et noethérien.
\end{enumerate}
\end{proposition}

\begin{proof}
Il est clair que toute algèbre de type fini sur un anneau noethérien universellement
japonais est un anneau de Nagata (c'est-à-dire que la condition (2) est satisfaite). Soit $R$ un anneau de Nagata.
Nous allons montrer que toute $R$-algèbre de type fini $S$ est un anneau de Nagata.
Cela démontrera la proposition.

\medskip\noindent
Étape 1. Il existe une suite de morphismes d'anneaux
$R = R_0 \to R_1 \to R_2 \to \ldots \to R_n = S$ telle que
chaque $R_i \to R_{i + 1}$ soit engendré par un seul élément.
Par récurrence, il suffit donc de démontrer que $S$ est un anneau de Nagata si
$S \cong R[x]/I$.

\medskip\noindent
Étape 2. Soit $\mathfrak q \subset S$ un idéal premier de $S$, et soit
$\mathfrak p \subset R$ l'idéal premier correspondant de $R$.
Nous devons montrer que $S/\mathfrak q$ est N-2. Nous nous sommes donc
ramenés à démontrer l'assertion suivante :
(*) Étant donnés un anneau intègre de Nagata $R$ et une extension monogène $R \subset S$
d'anneaux intègres, alors $S$ est N-2.

\medskip\noindent
Étape 3. Soit $R$ un anneau intègre de Nagata et
soit $R \subset S$ une extension monogène d'anneaux intègres.
Supposons que l'extension induite des corps des fractions de $R$ et de $S$
soit transcendante pure. Dans ce cas, $S = R[x]$. D'après le
lemme \ref{algebra-lemma-polynomial-ring-N-2}, on voit que $S$ est N-2.
Nous nous sommes donc ramenés à démontrer l'assertion suivante :
(**) Étant donnés un anneau intègre de Nagata $R$ et une
extension monogène $R \subset S$ d'anneaux intègres
induisant une extension finie des corps des fractions,
alors $S$ est N-2.

\medskip\noindent
Étape 4. Soit $R$ un anneau intègre normal de Nagata et
soit $R \subset S$ une extension monogène d'anneaux intègres
induisant une extension finie $L/K$ des corps des fractions.
Choisissons un élément $x \in S$
qui engendre $S$ comme $R$-algèbre. Soit $M/L$
une extension finie de corps.
Soit $R'$ la clôture intégrale de $R$ dans $M$.
Alors la clôture intégrale $S'$ de $S$ dans $M$ coïncide avec la clôture
intégrale de $R'[x]$ dans $M$.
De plus, le corps des fractions de $R'$ est $M$, et $R \subset R'$
est fini (par la propriété de Nagata de $R$).
Cela implique que $R'$ est un anneau de Nagata
(lemme \ref{algebra-lemma-quasi-finite-over-nagata}).
Montrer que $S'$ est fini sur $S$ revient à montrer que
$S'$ est fini sur $R'[x]$. Remplaçons $R$ par $R'$ et $S$ par $R'[x]$
pour nous ramener à l'assertion suivante :
(***) Étant donnés un anneau intègre normal de Nagata $R$, de corps des fractions $K$,
et $x \in K$, l'anneau $S \subset K$ engendré par $R$ et $x$
est N-1.

\medskip\noindent
Étape 5. Soit $R$ un anneau intègre normal de Nagata de corps des fractions $K$.
Soit $x = b/a \in K$. Nous devons montrer que l'anneau $S \subset K$
engendré par $R$ et $x$ est N-1. Remarquons que $S_a \cong R_a$ est normal.
Ainsi, d'après le lemme \ref{algebra-lemma-characterize-N-1}, il suffit de montrer que
$S_{\mathfrak m}$ est N-1 pour tout idéal maximal $\mathfrak m$ de $S$.

\medskip\noindent
Sous les hypothèses du paragraphe précédent, choisissons un tel idéal maximal
et posons $\mathfrak n = R \cap \mathfrak m$. L'extension de corps résiduels
$\kappa(\mathfrak m)/\kappa(\mathfrak n)$ est finie
(théorème \ref{algebra-theorem-nullstellensatz}) et engendrée par l'image de $x$.
Après avoir remplacé $R$ par $R_a$ pour un certain $a \in R$, $a \not \in \mathfrak n$,
et $S$ par $S_a$, nous pouvons supposer qu'il existe un polynôme unitaire
$f(X) = X^d + \sum_{i = 1, \ldots, d} a_iX^{d - i}$ dans $R[x]$ tel que
$f(x) \in \mathfrak m$ (détails omis).
Soit $K''/K$ une extension finie de corps dans laquelle le polynôme
$f(X)$ se scinde complètement dans $K''[X]$.
Soit $R'$ la clôture intégrale de $R$ dans $K''$.
Soit $S' \subset K''$ le sous-anneau engendré par $R'$ et $x$.
Comme $R$ est un anneau de Nagata, $R'$ est fini sur $R$ et est de Nagata
(lemme \ref{algebra-lemma-quasi-finite-over-nagata}).
De plus, $S'$ est fini sur $S$. Si, pour tout idéal maximal
$\mathfrak m'$ de $S'$, l'anneau local $S'_{\mathfrak m'}$ est
N-1, alors $S'_{\mathfrak m}$ est N-1 d'après le
lemme \ref{algebra-lemma-characterize-N-1}, ce qui implique à son tour
que $S_{\mathfrak m}$ est N-1 d'après le
lemme \ref{algebra-lemma-finite-extension-N-2}.
Après avoir remplacé $R$ par $R'$ et $S$ par $S'$, et $\mathfrak m$ par
l'un quelconque des idéaux maximaux $\mathfrak m'$ au-dessus de $\mathfrak m$,
nous arrivons à la situation où le polynôme $f$ ci-dessus est complètement scindé :
$f(X) = \prod_{i = 1, \ldots, d} (X - a_i)$ avec $a_i \in R$.
Puisque $f(x) \in \mathfrak m$, on a $x - a_i \in \mathfrak m$
pour un certain $i$. Enfin, après avoir remplacé $x$ par $x - a_i$, nous pouvons supposer
que $x \in \mathfrak m$.

\medskip\noindent
Récapitulons : $R$ est un anneau intègre normal de Nagata de corps des fractions $K$,
$x \in K$ et $S$ est le sous-anneau de $K$ engendré par $x$ et $R$ ;
enfin, $\mathfrak m \subset S$ est un idéal maximal tel que $x \in \mathfrak m$.
Nous devons montrer que $S_{\mathfrak m}$ est N-1.

\medskip\noindent
Nous allons montrer que le lemme \ref{algebra-lemma-criterion-analytically-unramified}
s'applique à l'anneau local
$S_{\mathfrak m}$ et à l'élément $x$. Il en résultera que
$S_{\mathfrak m}$ est analytiquement non ramifié, et nous
verrons alors qu'il est N-1 d'après le lemme \ref{algebra-lemma-analytically-unramified-easy}.

\medskip\noindent
Nous devons vérifier les conditions (1), (2), (3)(a) et (3)(b).
La condition (1) est triviale.
Soit $I = \Ker(R[X] \to S)$, où $X \mapsto x$.
Nous affirmons que $I$ est engendré par toutes les formes linéaires $aX - b$ telles que
$ax = b$ dans $K$. Il est clair que toutes ces formes linéaires appartiennent à $I$.
Si $g = a_d X^d + \ldots a_1 X + a_0 \in I$, alors on voit que
$a_dx$ est entier sur $R$ (lemme \ref{algebra-lemma-make-integral-trivial}),
et donc que $b := a_dx \in R$,
puisque $R$ est normal. Alors $g - (a_dX - b)X^{d - 1} \in I$, et le résultat suit par
récurrence sur le degré. Par conséquent,
$$
S/xS = R[X]/(X, I) = R/J
$$
où
$$
J = \{b \in R \mid ax = b \text{ pour un certain }a \in R\} = xR \cap R
$$
D'après le lemme \ref{algebra-lemma-normal-domain-intersection-localizations-height-1},
on voit que $S/xS = R/J$ n'a aucun idéal premier associé immergé comme $R$-module, donc comme
$R/J$-module, donc comme $S/xS$-module, donc comme $S$-module.
Cela démontre la condition (2).
Prenons un tel idéal premier associé $\mathfrak q \subset S$ vérifiant
$\mathfrak q \subset \mathfrak m$ (de sorte que ce soit un
idéal premier associé à $S_{\mathfrak m}/xS_{\mathfrak m}$ -- cela n'a aucune
incidence sur les arguments).
Alors $\mathfrak q$ est minimal au-dessus de $xS$ et est donc de hauteur $1$.
La suite d'égalités ci-dessus montre que
$\mathfrak p = R \cap \mathfrak q$ est un idéal premier associé
à $R/J$, et est donc de hauteur $1$
(voir le lemme \ref{algebra-lemma-normal-domain-intersection-localizations-height-1}).
Ainsi, $R_{\mathfrak p}$ est un anneau de valuation discrète, et par conséquent
$R_{\mathfrak p} \subset S_{\mathfrak q}$ est une égalité. Cela montre
que $S_{\mathfrak q}$ est régulier. La condition (3)(a) est ainsi satisfaite.
Enfin, $(S/\mathfrak q)_{\mathfrak m}$ est une localisation
de $S/\mathfrak q$, qui est un quotient de $S/xS = R/J$.
Par conséquent, $(S/\mathfrak q)_{\mathfrak m}$ est une localisation d'un
quotient de l'anneau de Nagata $R$, et est donc
de Nagata (lemmes \ref{algebra-lemma-quasi-finite-over-nagata}
et \ref{algebra-lemma-nagata-localize}),
et donc analytiquement non ramifié
(lemme \ref{algebra-lemma-local-nagata-domain-analytically-unramified}).
Cela montre que (3)(b) est satisfaite et termine la démonstration.
\end{proof}

\begin{proposition}
\label{algebra-proposition-ubiquity-nagata}
Les types d'anneaux suivants sont de Nagata et, en particulier, universellement japonais :
\begin{enumerate}
\item les corps,
\item les anneaux locaux noethériens complets,
\item $\mathbf{Z}$,
\item les anneaux de Dedekind dont le corps des fractions est de caractéristique zéro,
\item les extensions d'anneaux de type fini de l'un quelconque des anneaux précédents.
\end{enumerate}
\end{proposition}

\begin{proof}
Le cas des anneaux locaux noethériens complets est traité dans le
lemme \ref{algebra-lemma-Noetherian-complete-local-Nagata}.
Dans les autres cas, il suffit de vérifier que $R/\mathfrak p$ est N-2 pour tout
idéal premier $\mathfrak p$ de l'anneau. Cela est clair dès que
$R/\mathfrak p$ est un corps, c'est-à-dire lorsque $\mathfrak p$ est maximal.
Dans le cas d'un anneau de Dedekind, il suffit donc d'effectuer la vérification lorsque
$\mathfrak p = (0)$. Mais, puisque nous supposons que le corps des fractions est de
caractéristique zéro, le lemme \ref{algebra-lemma-domain-char-zero-N-1-2} s'applique.
\end{proof}

\begin{example}
\label{algebra-example-nonjapanese-dvr}
Un anneau de valuation discrète est de Nagata si et seulement s'il est N-2 (car le
quotient par l'idéal maximal est un corps, donc est N-2). L'anneau de valuation discrète
$A$ de l'exemple \ref{algebra-example-bad-dvr-char-p} n'est pas de Nagata, c'est-à-dire
qu'il n'est pas N-2. En effet, l'extension finie
$A \subset R = A[f]$ n'est pas N-1. Pour le voir, écrivons $f = \sum a_i x^i$.
Pour tout $n \geq 1$, posons $g_n = \sum_{i < n} a_i x^i \in A$.
Alors $h_n = (f - g_n)/x^n$ est un élément du corps des fractions de $R$
et $h_n^p \in k^p[[x]] \subset A$. La clôture intégrale $R'$
de $R$ contient donc $h_1, h_2, h_3, \ldots$. Or, si $R'$
était fini sur $R$, et donc sur $A$, alors $f  = x^n h_n + g_n$
appartiendrait au sous-module $A + x^nR'$ pour tout $n$. Le lemme
d'Artin--Rees impliquerait alors $f \in A$
(lemme \ref{algebra-lemma-intersect-powers-ideal-module-zero}), ce qui est contradictoire.
\end{example}

\begin{lemma}
\label{algebra-lemma-nagata-pth-roots}
Soit $(A, \mathfrak m)$ un anneau intègre local noethérien qui est un anneau de Nagata
et dont le corps des fractions est de caractéristique $p$. Si $a \in A$ possède une
racine $p$-ième dans $A^\wedge$, alors $a$ possède une racine $p$-ième dans $A$.
\end{lemma}

\begin{proof}
Soit $\alpha \in A^\wedge$ une racine $p$-ième de $a$. Pour obtenir une contradiction,
supposons que $a$ ne possède pas de racine $p$-ième dans $A$. Alors
$a$ ne possède pas de racine $p$-ième dans le corps des fractions $K$ de $A$.
En effet, si $\beta = b/c$, avec $b, c \in A$ et $c \not = 0$,
vérifie $\beta^p = a$, alors $\alpha = b/c$ dans $A^\wedge$,
ce qui implique que $c$ divise $b$ dans $A$
par fidèle platitude de $A \to A^\wedge$,
et donc que $\beta \in A$, contradiction. Ainsi,
$B = A[x]/(x^p - a)$ est intègre, car $K[x]/(x^p - a)$ est un corps.
Cependant, le complété de $B$ n'est pas réduit en raison de l'existence supposée
de $\alpha$. Cela contredit les résultats précédents, car $B$ est un anneau de Nagata
(proposition \ref{algebra-proposition-nagata-universally-japanese})
et est donc analytiquement non ramifié d'après le
lemme \ref{algebra-lemma-local-nagata-domain-analytically-unramified}.
\end{proof}






\section{Montée des propriétés}
\label{algebra-section-ascending-properties}

\noindent
Dans cette section, nous commençons à démontrer quelques faits algébriques concernant la
« montée » de propriétés des anneaux. Pour traiter le cas de la profondeur des anneaux,
on utilise le résultat suivant sur la montée de la profondeur des modules, voir
\cite[IV, Proposition 6.3.1]{EGA}.

\begin{lemma}
\label{algebra-lemma-apply-grothendieck-module}
\begin{reference}
\cite[IV, Proposition 6.3.1]{EGA}
\end{reference}
Nous avons
$$
\text{depth}(M \otimes_R N)
=
\text{depth}(M) + \text{depth}(N/\mathfrak m_RN)
$$
où $R \to S$ est un morphisme local d'anneaux locaux noethériens,
$M$ est un $R$-module de type fini et $N$ est un $S$-module de type fini plat sur $R$.
\end{lemma}

\begin{proof}
Dans l'énoncé et dans la démonstration ci-dessous, la profondeur de $M$ est prise
en tant que $R$-module, celle de $M \otimes_R N$ en tant que $S$-module, et
celle de $N/\mathfrak m_RN$ en tant que $S/\mathfrak m_RS$-module.
Notons $n$ le membre de droite. Supposons d'abord que $n$ soit nul.
Alors nous avons à la fois $\text{depth}(M) = 0$ et
$\text{depth}(N/\mathfrak m_RN) = 0$.
Cela signifie qu'il existe $z \in M$ dont l'annulateur est $\mathfrak m_R$
et $\overline{y} \in N/\mathfrak m_RN$
dont l'annulateur est $\mathfrak m_S/\mathfrak m_RS$.
Soit $y \in N$ un relèvement de $\overline{y}$.
Puisque $N$ est plat sur $R$, l'application $z : R/\mathfrak m_R \to M$
induit une application injective $N/\mathfrak m_RN \to M \otimes_R N$.
L'annulateur de $z \otimes y$ est donc $\mathfrak m_S$.
Ainsi, $\text{depth}(M \otimes_R N) = 0$ également.

\medskip\noindent
Supposons $n > 0$. Si $\text{depth}(N/\mathfrak m_RN) > 0$, nous pouvons choisir
$f \in \mathfrak m_S$ dont l'image $\overline{f} \in S/\mathfrak m_RS$
est un non-diviseur de zéro dans $N/\mathfrak m_RN$.
Alors $\text{depth}(N/\mathfrak m_RN) =
\text{depth}(N/(f, \mathfrak m_R)N) + 1$
d'après le lemme \ref{algebra-lemma-depth-drops-by-one}.
D'après le lemme \ref{algebra-lemma-mod-injective}, l'élément $f \in S$ est un
non-diviseur de zéro dans $N$ et $N/fN$ est plat sur $R$.
Par récurrence sur $n$, nous avons donc
$$
\text{depth}(M \otimes_R N/fN) =
\text{depth}(M) + \text{depth}(N/(f, \mathfrak m_R)N).
$$
Puisque $N/fN$ est plat sur $R$, la suite
$$
0 \to M \otimes_R N \to M \otimes_R N \to M \otimes_R N/fN \to 0
$$
est exacte, où la première application est la multiplication par $f$
(lemme \ref{algebra-lemma-flat-tor-zero}). Par conséquent, d'après le
lemme \ref{algebra-lemma-depth-drops-by-one}, nous obtenons
$\text{depth}(M \otimes_R N) = \text{depth}(M \otimes_R N/fN) + 1$
et nous obtenons ainsi l'égalité de la formule du lemme.

\medskip\noindent
Si $n > 0$, mais $\text{depth}(N/\mathfrak m_RN) = 0$,
nous pouvons choisir $f \in \mathfrak m_R$ qui soit un non-diviseur de zéro dans $M$.
Puisque $N$ est plat sur $R$, $f$ est également un non-diviseur de zéro dans
$M \otimes_R N$. Par une nouvelle récurrence sur $n$, nous avons
$$
\text{depth}(M/fM \otimes_R N) =
\text{depth}(M/fM) + \text{depth}(N/\mathfrak m_RN).
$$
Dans ce cas,
$\text{depth}(M \otimes_R N) = \text{depth}(M/fM \otimes_R N) + 1$
et $\text{depth}(M) = \text{depth}(M/fM) + 1$
d'après le lemme \ref{algebra-lemma-depth-drops-by-one}, et
nous obtenons ainsi l'égalité de la formule du lemme.
\end{proof}

\begin{lemma}
\label{algebra-lemma-apply-grothendieck}
Supposons que $R \to S$ soit un morphisme local plat d'anneaux locaux
noethériens. Alors
$$
\text{depth}(S) = \text{depth}(R) + \text{depth}(S/\mathfrak m_RS).
$$
\end{lemma}

\begin{proof}
C'est un cas particulier du lemme \ref{algebra-lemma-apply-grothendieck-module}.
\end{proof}

\begin{lemma}
\label{algebra-lemma-CM-goes-up}
Soit $R \to S$ un morphisme local plat d'anneaux locaux noethériens.
Alors les conditions suivantes sont équivalentes :
\begin{enumerate}
\item $S$ est de Cohen-Macaulay, et
\item $R$ et $S/\mathfrak m_RS$ sont de Cohen-Macaulay.
\end{enumerate}
\end{lemma}

\begin{proof}
Cela résulte des définitions et des
lemmes \ref{algebra-lemma-apply-grothendieck} et
\ref{algebra-lemma-dimension-base-fibre-equals-total}.
\end{proof}

\begin{lemma}
\label{algebra-lemma-Sk-goes-up}
Soit $\varphi : R \to S$ un morphisme d'anneaux. Supposons que
\begin{enumerate}
\item $R$ soit noethérien,
\item $S$ soit noethérien,
\item $\varphi$ soit plat,
\item les anneaux des fibres $S \otimes_R \kappa(\mathfrak p)$ vérifient $(S_k)$, et
\item $R$ vérifie la propriété $(S_k)$.
\end{enumerate}
Alors $S$ vérifie la propriété $(S_k)$.
\end{lemma}

\begin{proof}
Soit $\mathfrak q$ un idéal premier de $S$
au-dessus d'un idéal premier $\mathfrak p$ de $R$. D'après le
lemme \ref{algebra-lemma-apply-grothendieck}, nous avons
$$
\text{depth}(S_{\mathfrak q}) =
\text{depth}(S_{\mathfrak q}/\mathfrak pS_{\mathfrak q}) +
\text{depth}(R_{\mathfrak p}).
$$
D'autre part, nous avons
$$
\dim(R_{\mathfrak p})
+
\dim(S_{\mathfrak q}/\mathfrak pS_{\mathfrak q})
\geq
\dim(S_{\mathfrak q})
$$
d'après le lemme \ref{algebra-lemma-dimension-base-fibre-total}.
(En fait, il y a égalité, d'après le
lemme \ref{algebra-lemma-dimension-base-fibre-equals-total},
mais, à strictement parler, nous n'en avons pas besoin.)
Enfin, puisque les anneaux des fibres de ce morphisme
sont supposés vérifier $(S_k)$, nous voyons que
$\text{depth}(S_{\mathfrak q}/\mathfrak pS_{\mathfrak q})
\geq \min(k, \dim(S_{\mathfrak q}/\mathfrak pS_{\mathfrak q}))$.
Ainsi, le lemme résulte de la suite d'inégalités suivante :
\begin{eqnarray*}
\text{depth}(S_{\mathfrak q}) & = &
\text{depth}(S_{\mathfrak q}/\mathfrak pS_{\mathfrak q}) +
\text{depth}(R_{\mathfrak p}) \\
& \geq &
\min(k, \dim(S_{\mathfrak q}/\mathfrak pS_{\mathfrak q})) +
\min(k, \dim(R_{\mathfrak p})) \\
& = &
\min(2k, \dim(S_{\mathfrak q}/\mathfrak pS_{\mathfrak q}) + k,
k + \dim(R_\mathfrak p),
\dim(S_{\mathfrak q}/\mathfrak pS_{\mathfrak q}) +
\dim(R_{\mathfrak p})) \\
& \geq &
\min(k, \dim(S_{\mathfrak q}))
\end{eqnarray*}
ce qu'il fallait démontrer.
\end{proof}

\begin{lemma}
\label{algebra-lemma-Rk-goes-up}
Soit $\varphi : R \to S$ un morphisme d'anneaux. Supposons que
\begin{enumerate}
\item $R$ soit noethérien,
\item $S$ soit noethérien,
\item $\varphi$ soit plat,
\item les anneaux des fibres $S \otimes_R \kappa(\mathfrak p)$
vérifient la propriété $(R_k)$, et
\item $R$ vérifie la propriété $(R_k)$.
\end{enumerate}
Alors $S$ vérifie la propriété $(R_k)$.
\end{lemma}

\begin{proof}
Soit $\mathfrak q$ un idéal premier de $S$
au-dessus d'un idéal premier $\mathfrak p$ de $R$.
Supposons que $\dim(S_{\mathfrak q}) \leq k$.
Puisque $\dim(S_{\mathfrak q}) = \dim(R_{\mathfrak p})
+ \dim(S_{\mathfrak q}/\mathfrak pS_{\mathfrak q})$ d'après le
lemme \ref{algebra-lemma-dimension-base-fibre-equals-total},
nous voyons que $\dim(R_{\mathfrak p}) \leq k$ et que
$\dim(S_{\mathfrak q}/\mathfrak pS_{\mathfrak q}) \leq k$.
Ainsi, $R_{\mathfrak p}$ et $S_{\mathfrak q}/\mathfrak pS_{\mathfrak q}$
sont réguliers par hypothèse.
Il s'ensuit que $S_{\mathfrak q}$ est régulier d'après le
lemme \ref{algebra-lemma-flat-over-regular-with-regular-fibre}.
\end{proof}

\begin{lemma}
\label{algebra-lemma-reduced-goes-up-noetherian}
Soit $\varphi : R \to S$ un morphisme d'anneaux. Supposons que
\begin{enumerate}
\item $R$ soit noethérien,
\item $S$ soit noethérien,
\item $\varphi$ soit plat,
\item les anneaux des fibres $S \otimes_R \kappa(\mathfrak p)$ soient réduits,
\item $R$ soit réduit.
\end{enumerate}
Alors $S$ est réduit.
\end{lemma}

\begin{proof}
Pour un anneau noethérien, être réduit équivaut à vérifier les propriétés
$(S_1)$ et $(R_0)$, voir le lemme \ref{algebra-lemma-criterion-reduced}.
Nous savons donc que $R$ et les anneaux des fibres vérifient ces propriétés.
Nous pouvons ainsi appliquer les lemmes \ref{algebra-lemma-Sk-goes-up} et \ref{algebra-lemma-Rk-goes-up}
et nous voyons que $S$ vérifie $(S_1)$ et $(R_0)$, c'est-à-dire qu'il est réduit,
à nouveau d'après le lemme \ref{algebra-lemma-criterion-reduced}.
\end{proof}

\begin{lemma}
\label{algebra-lemma-reduced-goes-up}
Soit $\varphi : R \to S$ un morphisme d'anneaux. Supposons que
\begin{enumerate}
\item $\varphi$ soit lisse,
\item $R$ soit réduit.
\end{enumerate}
Alors $S$ est réduit.
\end{lemma}

\begin{proof}
Remarquons que $R \to S$ est plat et à fibres régulières (voir la liste des
résultats sur les morphismes lisses d'anneaux à la section \ref{algebra-section-smooth-overview}).
En particulier, les fibres sont réduites.
Ainsi, si $R$ est noethérien, alors $S$ est noethérien et le résultat
découle du lemme \ref{algebra-lemma-reduced-goes-up-noetherian}.

\medskip\noindent
Dans le cas général, nous pouvons trouver une sous-$\mathbf{Z}$-algèbre
$R_0 \subset R$ de type fini et un morphisme lisse d'anneaux
$R_0 \to S_0$ tels que $S \cong R \otimes_{R_0} S_0$, voir la
remarque (10) de la section \ref{algebra-section-smooth-overview}.
Si maintenant $x \in S$ est un élément tel que $x^2 = 0$,
nous pouvons agrandir $R_0$ et supposer que $x$ provient
d'un élément $x_0 \in S_0$. Après avoir agrandi
$R_0$ une fois encore, nous pouvons supposer que $x_0^2 = 0$ dans $S_0$.
Cependant, puisque la sous-algèbre $R_0 \subset R$ est réduite, nous voyons que
$S_0$ est réduit et donc que $x_0 = 0$, comme voulu.
\end{proof}

\begin{lemma}
\label{algebra-lemma-normal-goes-up-noetherian}
Soit $\varphi : R \to S$ un morphisme d'anneaux. Supposons que
\begin{enumerate}
\item $R$ soit noethérien,
\item $S$ soit noethérien,
\item $\varphi$ soit plat,
\item les anneaux des fibres $S \otimes_R \kappa(\mathfrak p)$ soient normaux, et
\item $R$ soit normal.
\end{enumerate}
Alors $S$ est normal.
\end{lemma}

\begin{proof}
Pour un anneau noethérien, être normal équivaut à vérifier les propriétés
$(S_2)$ et $(R_1)$, voir le lemme \ref{algebra-lemma-criterion-normal}.
Nous savons donc que $R$ et les anneaux des fibres vérifient ces propriétés.
Nous pouvons ainsi appliquer les lemmes \ref{algebra-lemma-Sk-goes-up} et \ref{algebra-lemma-Rk-goes-up}
et nous voyons que $S$ vérifie $(S_2)$ et $(R_1)$, c'est-à-dire qu'il est normal,
à nouveau d'après le lemme \ref{algebra-lemma-criterion-normal}.
\end{proof}

\begin{lemma}
\label{algebra-lemma-normal-goes-up}
Soit $\varphi : R \to S$ un morphisme d'anneaux. Supposons que
\begin{enumerate}
\item $\varphi$ soit lisse,
\item $R$ soit normal.
\end{enumerate}
Alors $S$ est normal.
\end{lemma}

\begin{proof}
Remarquons que $R \to S$ est plat et à fibres régulières (voir la liste des
résultats sur les morphismes lisses d'anneaux à la section \ref{algebra-section-smooth-overview}).
En particulier, les fibres sont normales. Ainsi, si $R$ est noethérien,
alors $S$ est noethérien et le résultat découle du
lemme \ref{algebra-lemma-normal-goes-up-noetherian}.

\medskip\noindent
Passons au cas général. Remarquons d'abord que $R$ est réduit et donc que
$S$ est réduit d'après le lemme \ref{algebra-lemma-reduced-goes-up}.
Soit $\mathfrak q$ un idéal premier de $S$ et soit $\mathfrak p$
l'idéal premier correspondant dans $R$. Remarquons que $R_{\mathfrak p}$
est un anneau intègre normal. Nous devons montrer que $S_{\mathfrak q}$ est
un anneau intègre normal. Pour ce faire, nous pouvons remplacer $R$ par $R_{\mathfrak p}$
et $S$ par $S_{\mathfrak p}$. Nous pouvons donc supposer que $R$ est
un anneau intègre normal.

\medskip\noindent
Supposons le morphisme $R \to S$ lisse et $R$ un anneau intègre normal.
Nous pouvons trouver une sous-$\mathbf{Z}$-algèbre de type fini
$R_0 \subset R$ et un morphisme lisse d'anneaux $R_0 \to S_0$ tels
que $S \cong R \otimes_{R_0} S_0$, voir la
remarque (10) de la section \ref{algebra-section-smooth-overview}.
Puisque $R_0$ est un anneau intègre de Nagata (voir la proposition \ref{algebra-proposition-ubiquity-nagata}),
nous voyons que sa clôture intégrale $R_0'$ est finie sur $R_0$.
De plus, comme $R$ est un anneau intègre normal, il est clair que $R_0' \subset R$.
Nous pouvons donc remplacer $R_0$ par $R_0'$ et $S_0$ par
$R_0' \otimes_{R_0} S_0$ et supposer que $R_0$ est un anneau intègre
normal noethérien. Le premier paragraphe de la démonstration montre alors
que $S_0$ est un anneau normal (il n'est bien sûr pas nécessairement intègre).
De cette manière, nous voyons que $R = \bigcup R_\lambda$
est la réunion d'anneaux intègres normaux noethériens et, parallèlement,
$S = \colim R_\lambda \otimes_{R_0} S_0$ est la limite inductive
d'anneaux normaux. Cela implique que $S$ est un anneau normal.
Nous omettons quelques détails.
\end{proof}

\begin{lemma}
\label{algebra-lemma-regular-goes-up}
\begin{slogan}
La régularité se propage le long des morphismes lisses d'anneaux.
\end{slogan}
Soit $\varphi : R \to S$ un morphisme d'anneaux. Supposons que
\begin{enumerate}
\item $\varphi$ soit lisse,
\item $R$ soit un anneau régulier.
\end{enumerate}
Alors $S$ est régulier.
\end{lemma}

\begin{proof}
Cela résulte du lemme \ref{algebra-lemma-Rk-goes-up} appliqué à toutes les propriétés $(R_k)$,
le lemme \ref{algebra-lemma-characterize-smooth-over-field} permettant de voir que les
hypothèses sont satisfaites.
\end{proof}






\section{Descente des propriétés}
\label{algebra-section-descending-properties}

\noindent
Dans cette section, nous commençons par établir quelques résultats algébriques concernant la
« descente » de propriétés des anneaux. Il s'avère qu'il est souvent
« plus facile » de faire descendre des propriétés que de les faire monter. Autrement
dit, les hypothèses portant sur le morphisme d'anneaux $R \to S$ sont souvent plus faibles que
celles du lemme correspondant de la section précédente.
Nous avertissons toutefois le lecteur que les résultats de descente sont souvent
inutiles si l'on ne peut pas également établir les résultats de montée correspondants !
Voici un résultat typique illustrant ce phénomène.

\begin{lemma}
\label{algebra-lemma-descent-Noetherian}
Soit $R \to S$ un morphisme d'anneaux.
Supposons que
\begin{enumerate}
\item $R \to S$ soit fidèlement plat, et
\item $S$ soit noethérien.
\end{enumerate}
Alors $R$ est noethérien.
\end{lemma}

\begin{proof}
Soit $I_0 \subset I_1 \subset I_2 \subset \ldots$ une
suite croissante d'idéaux de $R$. Par hypothèse, nous avons
$I_nS = I_{n+1}S = I_{n+2}S = \ldots$ pour un certain $n$.
Par fidèle platitude, l'extension suivie de la contraction laisse tout idéal inchangé,
c'est-à-dire que $I = R \cap IS$ pour tout idéal $I$ de $R$ (lemme
\ref{algebra-lemma-faithfully-flat-universally-injective}). Ainsi,
$I_n = I_{n+1} = I_{n+2} = \ldots$, ce qu'il fallait démontrer.
\end{proof}

\begin{lemma}
\label{algebra-lemma-descent-reduced}
Soit $R \to S$ un morphisme d'anneaux.
Supposons que
\begin{enumerate}
\item $R \to S$ soit fidèlement plat, et
\item $S$ soit réduit.
\end{enumerate}
Alors $R$ est réduit.
\end{lemma}

\begin{proof}
Cela est clair puisque $R \to S$ est injectif, d'après le
lemme \ref{algebra-lemma-faithfully-flat-universally-injective}.
\end{proof}

\begin{lemma}
\label{algebra-lemma-descent-normal}
Soit $R \to S$ un morphisme d'anneaux.
Supposons que
\begin{enumerate}
\item $R \to S$ soit fidèlement plat, et
\item $S$ soit un anneau normal.
\end{enumerate}
Alors $R$ est un anneau normal.
\end{lemma}

\begin{proof}
Puisque $S$ est réduit, il s'ensuit que $R$ est réduit.
Soit $\mathfrak p$ un idéal premier de $R$. Nous devons montrer que
$R_{\mathfrak p}$ est un anneau intègre normal. Comme $S_{\mathfrak p}$
est lui aussi fidèlement plat sur $R_{\mathfrak p}$, nous pouvons supposer que
$R$ est local, d'idéal maximal $\mathfrak m$.
Soit $\mathfrak q$ un idéal premier de $S$ au-dessus de $\mathfrak m$.
Nous voyons alors que $R \to S_{\mathfrak q}$ est fidèlement plat
(lemme \ref{algebra-lemma-local-flat-ff}).
Nous pouvons donc supposer que $S$ est lui aussi local.
En particulier, $S$ est un anneau intègre normal.
Comme $R \to S$ est fidèlement plat
et que $S$ est un anneau intègre normal, nous voyons que $R$ est un anneau intègre.
Ensuite, supposons que $a/b$ soit entier sur $R$, où $a, b \in R$.
Alors $a/b \in S$ puisque $S$ est normal. Ainsi $a \in bS$.
Cela signifie que l'application $a : R \to R/bR$ devient nulle
après changement de base à $S$. Par fidèle platitude, nous voyons que
$a \in bR$, donc $a/b \in R$. Ainsi $R$ est normal.
\end{proof}

\begin{lemma}
\label{algebra-lemma-descent-regular}
Soit $R \to S$ un morphisme d'anneaux.
Supposons que
\begin{enumerate}
\item $R \to S$ soit fidèlement plat, et
\item $S$ soit un anneau régulier.
\end{enumerate}
Alors $R$ est un anneau régulier.
\end{lemma}

\begin{proof}
Nous voyons que $R$ est noethérien d'après le lemme \ref{algebra-lemma-descent-Noetherian}.
Soit $\mathfrak p \subset R$ un idéal premier. Choisissons un idéal premier $\mathfrak q \subset S$
au-dessus de $\mathfrak p$. Le lemme \ref{algebra-lemma-flat-under-regular}
s'applique à $R_\mathfrak p \to S_\mathfrak q$ et nous en concluons que
$R_\mathfrak p$ est régulier. Comme $\mathfrak p$ était arbitraire, nous voyons que
$R$ est régulier.
\end{proof}

\begin{lemma}
\label{algebra-lemma-descent-Sk}
Soit $R \to S$ un morphisme d'anneaux.
Supposons que
\begin{enumerate}
\item $R \to S$ soit fidèlement plat, et
\item $S$ soit noethérien et vérifie la propriété $(S_k)$.
\end{enumerate}
Alors $R$ est noethérien et vérifie la propriété $(S_k)$.
\end{lemma}

\begin{proof}
Nous avons déjà vu que (1) et (2) impliquent que $R$ est noethérien,
voir le lemme \ref{algebra-lemma-descent-Noetherian}.
Soit $\mathfrak p \subset R$ un idéal premier.
Choisissons un idéal premier $\mathfrak q \subset S$ au-dessus de $\mathfrak p$
correspondant à un idéal premier minimal de l'anneau de la fibre
$S \otimes_R \kappa(\mathfrak p)$. Alors
$A = R_{\mathfrak p} \to S_{\mathfrak q} = B$ est un homomorphisme local et plat
d'anneaux locaux noethériens, et $\mathfrak m_AB$ est un
idéal de définition de $B$. Ainsi,
$\dim(A) = \dim(B)$ (lemme \ref{algebra-lemma-dimension-base-fibre-equals-total}) et
$\text{depth}(A) = \text{depth}(B)$ (lemme \ref{algebra-lemma-apply-grothendieck}).
Par conséquent, puisque $B$ vérifie la propriété $(S_k)$, nous
voyons que $A$ vérifie la propriété $(S_k)$.
\end{proof}

\begin{lemma}
\label{algebra-lemma-descent-Rk}
Soit $R \to S$ un morphisme d'anneaux. Supposons que
\begin{enumerate}
\item $R \to S$ soit fidèlement plat, et
\item $S$ soit noethérien et vérifie la propriété $(R_k)$.
\end{enumerate}
Alors $R$ est noethérien et vérifie la propriété $(R_k)$.
\end{lemma}

\begin{proof}
Nous avons déjà vu que (1) et (2) impliquent que $R$ est noethérien,
voir le lemme \ref{algebra-lemma-descent-Noetherian}.
Soit $\mathfrak p \subset R$ un idéal premier et supposons que
$\dim(R_{\mathfrak p}) \leq k$.
Choisissons un idéal premier $\mathfrak q \subset S$ au-dessus de $\mathfrak p$
correspondant à un idéal premier minimal de l'anneau de la fibre
$S \otimes_R \kappa(\mathfrak p)$. Alors
$A = R_{\mathfrak p} \to S_{\mathfrak q} = B$ est un homomorphisme local et plat
d'anneaux locaux noethériens, et $\mathfrak m_AB$ est un
idéal de définition de $B$. Ainsi,
$\dim(A) = \dim(B)$ (lemme \ref{algebra-lemma-dimension-base-fibre-equals-total}).
Comme $S$ vérifie la propriété $(R_k)$, nous en concluons que $B$ est un anneau local régulier.
D'après le lemme \ref{algebra-lemma-flat-under-regular}, nous en concluons que $A$ est régulier.
\end{proof}

\begin{lemma}
\label{algebra-lemma-descent-nagata}
Soit $R \to S$ un morphisme d'anneaux. Supposons que
\begin{enumerate}
\item $R \to S$ soit lisse et induise une surjection sur les spectres, et
\item $S$ soit un anneau de Nagata.
\end{enumerate}
Alors $R$ est un anneau de Nagata.
\end{lemma}

\begin{proof}
Rappelons qu'un anneau de Nagata est la même chose qu'un anneau noethérien
universellement japonais
(proposition \ref{algebra-proposition-nagata-universally-japanese}).
Nous avons déjà vu que $R$ est noethérien dans le
lemme \ref{algebra-lemma-descent-Noetherian}.
Soit $R \to A$ un morphisme d'anneaux de type fini, l'anneau cible étant intègre.
D'après le lemme \ref{algebra-lemma-check-universally-japanese},
il suffit de vérifier que $A$ est N-1.
Il est clair que $B = A \otimes_R S$ est une $S$-algèbre de type fini
et donc un anneau de Nagata (proposition \ref{algebra-proposition-nagata-universally-japanese}).
Puisque $A \to B$ est lisse (lemme \ref{algebra-lemma-base-change-smooth}),
nous voyons que $B$ est réduit (lemme \ref{algebra-lemma-reduced-goes-up}).
Comme $B$ est noethérien, il ne possède qu'un nombre fini d'idéaux
premiers minimaux $\mathfrak q_1, \ldots, \mathfrak q_t$ (voir le
lemme \ref{algebra-lemma-Noetherian-irreducible-components}).
Comme $A \to B$ est plat, chacun de ces idéaux se trouve au-dessus de $(0) \subset A$
(par le théorème de descente, voir le lemme \ref{algebra-lemma-flat-going-down}).
L'anneau total des fractions $Q(B)$ est le produit des
$L_i = \kappa(\mathfrak q_i)$ (lemmes
\ref{algebra-lemma-total-ring-fractions-no-embedded-points} et
\ref{algebra-lemma-minimal-prime-reduced-ring}).
En outre, la clôture intégrale $B'$ de $B$ dans $Q(B)$ est
le produit des clôtures intégrales $B_i'$ des anneaux $B/\mathfrak q_i$
dans les facteurs $L_i$ (comparer avec le
lemme \ref{algebra-lemma-characterize-reduced-ring-normal}).
Puisque $B$ est universellement japonais, les
extensions d'anneaux $B/\mathfrak q_i \subset B_i'$ sont finies
et nous en concluons que $B' = \prod B_i'$ est un $B$-module de type fini.
Comme $A \to B$ est plat, nous voyons que tout
non-diviseur de zéro dans $A$ a pour image un non-diviseur de zéro dans $B$.
L'application correspondante
$$
Q(A) \otimes_A B = (A \setminus \{0\})^{-1}A \otimes_A B
= (A \setminus \{0\})^{-1}B \to Q(B)
$$
est injective (nous avons utilisé le lemme \ref{algebra-lemma-tensor-localization}).
$A'$ désigne ici la clôture intégrale de l'anneau de départ dans son corps des fractions, et cette application l'envoie dans $B'$. Elle induit une application
$$
A' \otimes_A B \longrightarrow B'
$$
qui est injective (d'après ce qui précède et la platitude de $A \to B$).
Puisque $B'$ est un $B$-module de type fini
et que $B$ est noethérien, nous voyons que $A' \otimes_A B$ est un $B$-module de type fini.
Il existe donc un nombre fini d'éléments $x_i \in A'$ tels que
les éléments $x_i \otimes 1$ engendrent $A' \otimes_A B$ comme $B$-module.
Enfin, par fidèle platitude de $A \to B$, nous concluons que
les $x_i$ engendrent aussi $A'$ comme $A$-module, ce qui achève la démonstration.
\end{proof}

\begin{remark}
\label{algebra-remark-universally-catenary-does-not-descend}
La propriété d'être « universellement caténaire » ne descend pas, même
le long de morphismes \'etales d'anneaux. Dans le
chapitre des exemples, à la section \ref{examples-section-non-catenary-Noetherian-local},
on construit un morphisme fini d'anneaux $A \to B$ tel que
$A$ soit un anneau local noethérien non universellement caténaire
et que $B$ soit un anneau semi-local dont les deux idéaux maximaux sont $\mathfrak m$ et $\mathfrak n$.
Les anneaux $B_{\mathfrak m}$ et $B_{\mathfrak n}$ sont réguliers, de dimensions $2$ et $1$
respectivement, et ont chacun le même corps résiduel que celui de $A$.
De plus, $\mathfrak m_A$ engendre l'idéal maximal de chacun des anneaux
$B_{\mathfrak m}$ et $B_{\mathfrak n}$ (ainsi $A \to B$ est à la fois non ramifié
et fini).
D'après le lemme \ref{algebra-lemma-etale-makes-unramified-closed},
il existe un morphisme local et \'etale d'anneaux
$A \to A'$ pour lequel $B \otimes_A A' = B_1 \times B_2$,
les morphismes $A' \to B_i$ étant surjectifs.
Cela montre que $A'$ possède deux idéaux premiers minimaux $\mathfrak q_i$
tels que $A'/\mathfrak q_i \cong B_i$. Puisque $B_i$ est un anneau local régulier
(puisqu'il est \'etale sur $B_{\mathfrak m}$ ou sur $B_{\mathfrak n}$),
nous concluons que $A'$ est universellement caténaire.
\end{remark}









\section{Algèbres géométriquement normales}
\label{algebra-section-geometrically-normal}

\noindent
Dans cette section, nous donnons quelques applications de la montée et de la
descente des propriétés des anneaux.

\begin{lemma}
\label{algebra-lemma-geometrically-normal}
Soit $k$ un corps. Soit $A$ une $k$-algèbre.
Les propriétés suivantes de $A$ sont équivalentes :
\begin{enumerate}
\item $k' \otimes_k A$ est un anneau normal
pour toute extension de corps $k'/k$,
\item $k' \otimes_k A$ est un anneau normal
pour toute extension de corps de type fini $k'/k$,
\item $k' \otimes_k A$ est un anneau normal
pour toute extension de corps finie purement inséparable $k'/k$,
\item $k^{perf} \otimes_k A$ est un anneau normal.
\end{enumerate}
Ici, l'expression « anneau normal » est prise au sens de la définition \ref{algebra-definition-ring-normal}.
\end{lemma}

\begin{proof}
Il est clair que (1) $\Rightarrow$ (2) $\Rightarrow$ (3)
et que (1) $\Rightarrow$ (4).

\medskip\noindent
Si $k'/k$ est une extension de corps finie purement inséparable, alors
il existe un plongement $k' \to k^{perf}$ au-dessus de $k$.
Le morphisme d'anneaux $k' \otimes_k A \to k^{perf} \otimes_k A$
est fidèlement plat ; ainsi, $k' \otimes_k A$ est normal si
$k^{perf} \otimes_k A$ est normal, d'après le
lemme \ref{algebra-lemma-descent-normal}. On voit ainsi que
(4) $\Rightarrow$ (3).

\medskip\noindent
Supposons (2) et soit $k'/k$ une extension de corps quelconque.
Nous pouvons alors écrire $k' = \colim_i k_i$ sous la forme d'une limite
inductive filtrante d'extensions de corps de type fini. Par conséquent,
$k' \otimes_k A = \colim_i k_i \otimes_k A$
est une limite inductive filtrante d'anneaux normaux. Il s'ensuit
que $k' \otimes_k A$ est un anneau normal d'après le
lemme \ref{algebra-lemma-colimit-normal-ring}.
Ainsi, (1) est vérifiée.

\medskip\noindent
Supposons (3) et soit $K/k$ une extension de corps de type fini.
D'après le lemme \ref{algebra-lemma-make-separable}, il existe un diagramme
$$
\xymatrix{
K \ar[r] & K' \\
k \ar[u] \ar[r] & k' \ar[u]
}
$$
où $k'/k$ et $K'/K$ sont des extensions de corps finies purement inséparables
telles que $K'/k'$ soit séparable. D'après le
lemme \ref{algebra-lemma-localization-smooth-separable},
il existe une $k'$-algèbre lisse $B$ telle que $K'$ soit le
corps des fractions de $B$. Nous pouvons alors raisonner comme suit :
Étape 1 : $k' \otimes_k A$ est un anneau normal, puisque nous avons supposé (3).
Étape 2 : $B \otimes_{k'} k' \otimes_k A$ est un anneau normal, car
$k' \otimes_k A \to B \otimes_{k'} k' \otimes_k A$ est lisse
(lemme \ref{algebra-lemma-base-change-smooth})
et la normalité monte le long des morphismes lisses
(lemme \ref{algebra-lemma-normal-goes-up}).
Étape 3 : $K' \otimes_{k'} k' \otimes_k A = K' \otimes_k A$ est
un anneau normal, car c'est une localisation d'un anneau normal
(lemme \ref{algebra-lemma-localization-normal-ring}).
Étape 4 : Enfin, $K \otimes_k A$ est un anneau normal par descente de la
normalité le long du morphisme d'anneaux fidèlement plat
$K \otimes_k A \to K' \otimes_k A$ (lemme \ref{algebra-lemma-descent-normal}).
Ceci démontre le lemme.
\end{proof}

\begin{definition}
\label{algebra-definition-geometrically-normal}
Soit $k$ un corps.
Une $k$-algèbre $R$ est dite {\it géométriquement normale} sur $k$ si
les conditions équivalentes du lemme \ref{algebra-lemma-geometrically-normal} sont satisfaites.
\end{definition}

\begin{lemma}
\label{algebra-lemma-localization-geometrically-normal-algebra}
\begin{slogan}
La localisation préserve la normalité géométrique.
\end{slogan}
Soit $k$ un corps. Une localisation d'une $k$-algèbre géométriquement normale
est géométriquement normale.
\end{lemma}

\begin{proof}
Cela est clair, puisque la normalité d'un anneau se vérifie après localisation en tout
idéal premier.
\end{proof}

\begin{lemma}
\label{algebra-lemma-separable-field-extension-geometrically-normal}
Soit $k$ un corps. Soit $K/k$ une extension de corps séparable.
Alors $K$ est géométriquement normal sur $k$.
\end{lemma}

\begin{proof}
Cela résulte du fait que $k^{perf} \otimes_k K$ est un corps.
En effet, cet anneau est réduit d'après le
lemme \ref{algebra-lemma-separable-extension-preserves-reducedness}.
D'après le lemme \ref{algebra-lemma-perfection} (ou la
définition \ref{algebra-definition-perfection}),
l'extension de corps $k^{perf}/k$ est purement inséparable.
Par conséquent, d'après le lemme \ref{algebra-lemma-radicial-integral-bijective}, l'anneau
$k^{perf} \otimes_k K$ possède un unique idéal premier. Un anneau réduit ayant
un unique idéal premier est un corps.
\end{proof}

\begin{lemma}
\label{algebra-lemma-geometrically-normal-tensor-normal}
Soit $k$ un corps. Soient $A, B$ des $k$-algèbres. Supposons que $A$ soit géométriquement
normale sur $k$ et que $B$ soit un anneau normal. Alors $A \otimes_k B$ est un anneau
normal.
\end{lemma}

\begin{proof}
Soit $\mathfrak r$ un idéal premier de $A \otimes_k B$. Notons
$\mathfrak p$ (resp.\ $\mathfrak q$) l'idéal premier correspondant de $A$
(resp.\ de $B$). Alors $(A \otimes_k B)_{\mathfrak r}$ est une localisation de
$A_{\mathfrak p} \otimes_k B_{\mathfrak q}$. Il suffit donc de démontrer le
résultat pour l'anneau $A_{\mathfrak p} \otimes_k B_{\mathfrak q}$, voir le
lemme \ref{algebra-lemma-localization-normal-ring}
et le
lemme \ref{algebra-lemma-localization-geometrically-normal-algebra}.
Nous pouvons donc supposer que $A$ et $B$ sont des anneaux intègres.

\medskip\noindent
Supposons que $A$ et $B$ soient des anneaux intègres dont les corps des fractions respectifs sont $K$ et $L$.
Notons que $B$ est la limite inductive filtrante de ses sous-algèbres normales
de type fini sur $k$ (en effet, $k$ est un anneau de Nagata, voir la
proposition \ref{algebra-proposition-ubiquity-nagata},
et la clôture intégrale d'une sous-algèbre de type fini sur $k$ est encore
une sous-algèbre de type fini sur $k$, d'après la
proposition \ref{algebra-proposition-nagata-universally-japanese}).
D'après le
lemme \ref{algebra-lemma-colimit-normal-ring},
on se ramène au cas où $B$ est de type fini sur $k$.

\medskip\noindent
Supposons que $A$ et $B$ soient des anneaux intègres dont les corps des fractions respectifs sont $K$ et $L$,
et que $B$ soit de type fini sur $k$. Dans ce cas, l'anneau $K \otimes_k B$
est de type fini sur $K$, donc noethérien
(lemme \ref{algebra-lemma-Noetherian-permanence}).
En particulier, $K \otimes_k B$ possède un nombre fini d'idéaux premiers minimaux
(lemme \ref{algebra-lemma-Noetherian-irreducible-components}).
Comme $A \to A \otimes_k B$ est plat, il s'ensuit que $A \otimes_k B$
possède un nombre fini d'idéaux premiers minimaux (par le théorème de descente pour les morphismes d'anneaux plats --
lemme \ref{algebra-lemma-flat-going-down}
-- tous ces idéaux sont situés au-dessus de $(0) \subset A$). Il suffit donc de démontrer
que $A \otimes_k B$ est intégralement clos dans son anneau total des fractions
(lemme \ref{algebra-lemma-characterize-reduced-ring-normal}).

\medskip\noindent
Nous affirmons que $K \otimes_k B$ et $A \otimes_k L$ sont tous deux des anneaux normaux.
Si tel est le cas, tout élément $x$ de $Q(A \otimes_k B)$ qui est
entier sur $A \otimes_k B$ appartient (d'après le
lemme \ref{algebra-lemma-normal-ring-integrally-closed})
à $K \otimes_k B \cap A \otimes_k L = A \otimes_k B$, et la démonstration est terminée.
Puisque $A \otimes_k L$ est un anneau normal par hypothèse, il suffit de
prouver que $K \otimes_k B$ est normal.

\medskip\noindent
Comme $A$ est géométriquement normale sur $k$, on voit que $K$ est géométriquement normal
sur $k$
(lemme \ref{algebra-lemma-localization-geometrically-normal-algebra}),
donc $K$ est géométriquement réduit sur $k$.
Ainsi, $K = \bigcup K_i$ est la réunion d'extensions de corps de type fini
de $k$ qui sont géométriquement réduites
(lemme \ref{algebra-lemma-subalgebra-separable}).
Chaque $K_i$ est la localisation d'une $k$-algèbre lisse
(lemme \ref{algebra-lemma-localization-smooth-separable}).
Ainsi, $K_i \otimes_k B$ est une localisation d'une $B$-algèbre lisse, donc un anneau normal
(lemme \ref{algebra-lemma-normal-goes-up}).
Ainsi, $K \otimes_k B$ est un anneau normal
(lemme \ref{algebra-lemma-colimit-normal-ring}),
ce qui conclut la démonstration.
\end{proof}

\begin{lemma}
\label{algebra-lemma-geometrically-normal-over-separable-algebraic}
Soit $k'/k$ une extension de corps algébrique séparable.
Soit $A$ une algèbre sur $k'$. Alors $A$ est géométriquement normale
sur $k$ si et seulement si elle est géométriquement normale sur $k'$.
\end{lemma}

\begin{proof}
Soit $L/k$ une extension de corps finie purement inséparable.
Alors $L' = k' \otimes_k L$ est un corps (voir les résultats du chapitre
Corps, section \ref{fields-section-algebraic})
et $A \otimes_k L = A \otimes_{k'} L'$. Ainsi, si
$A$ est géométriquement normale sur $k'$, alors $A$ est géométriquement
normale sur $k$.

\medskip\noindent
Supposons que $A$ soit géométriquement normale sur $k$. Soit $K/k'$ une extension de corps.
Alors
$$
K \otimes_{k'} A = (K \otimes_k A) \otimes_{(k' \otimes_k k')} k'
$$
Puisque $k' \otimes_k k' \to k'$ est une localisation d'après le
lemme \ref{algebra-lemma-separable-algebraic-diagonal},
on voit que $K \otimes_{k'} A$
est une localisation d'un anneau normal, donc un anneau normal.
\end{proof}



\section{Algèbres géométriquement régulières}
\label{algebra-section-geometrically-regular}

\noindent
Soit $k$ un corps.
Soit $A$ une $k$-algèbre noethérienne.
Soit $K/k$ une extension de corps de type fini.
Alors l'anneau $K \otimes_k A$ est lui aussi noethérien, voir le
lemme \ref{algebra-lemma-Noetherian-field-extension}.
Le lemme suivant a donc un sens.

\begin{lemma}
\label{algebra-lemma-geometrically-regular}
Soit $k$ un corps. Soit $A$ une $k$-algèbre.
Supposons que $A$ soit noethérienne.
Les propriétés suivantes de $A$ sont équivalentes :
\begin{enumerate}
\item $k' \otimes_k A$ est régulier pour toute extension de corps
$k'/k$ de type fini, et
\item $k' \otimes_k A$ est régulier pour toute extension de corps
$k'/k$ finie purement inséparable.
\end{enumerate}
Ici, un anneau régulier est entendu au sens de la définition \ref{algebra-definition-regular}.
\end{lemma}

\begin{proof}
Le lemme a un sens d'après les remarques qui le précèdent.
Il est clair que (1) $\Rightarrow$ (2).

\medskip\noindent
Supposons (2) et soit $K/k$ une extension de corps de type fini.
D'après le lemme \ref{algebra-lemma-make-separable}, on peut trouver un diagramme
$$
\xymatrix{
K \ar[r] & K' \\
k \ar[u] \ar[r] & k' \ar[u]
}
$$
où $k'/k$ et $K'/K$ sont des extensions de corps finies purement
inséparables telles que $K'/k'$ soit séparable. D'après le
lemme \ref{algebra-lemma-localization-smooth-separable},
il existe une $k'$-algèbre lisse $B$ dont $K'$ est le
corps des fractions de $B$. On peut maintenant raisonner comme suit :
Étape 1 : $k' \otimes_k A$ est un anneau régulier puisque nous avons supposé (2).
Étape 2 : $B \otimes_{k'} k' \otimes_k A$ est un anneau régulier, puisque
le morphisme $k' \otimes_k A \to B \otimes_{k'} k' \otimes_k A$ est lisse
(lemme \ref{algebra-lemma-base-change-smooth})
et que la régularité monte le long des morphismes lisses
(lemme \ref{algebra-lemma-regular-goes-up}).
Étape 3. $K' \otimes_{k'} k' \otimes_k A = K' \otimes_k A$ est
un anneau régulier, puisqu'il s'agit d'un localisé d'un anneau régulier
(cela découle immédiatement de la définition).
Étape 4. Enfin, $K \otimes_k A$ est un anneau régulier par descente de la
régularité le long du morphisme d'anneaux fidèlement plat
$K \otimes_k A \to K' \otimes_k A$ (lemme \ref{algebra-lemma-descent-regular}).
Ceci démontre le lemme.
\end{proof}

\begin{definition}
\label{algebra-definition-geometrically-regular}
Soit $k$ un corps. Soit $R$ une $k$-algèbre noethérienne.
La $k$-algèbre $R$ est dite {\it géométriquement régulière} sur $k$ si
les conditions équivalentes du lemme \ref{algebra-lemma-geometrically-regular} sont satisfaites.
\end{definition}

\noindent
Il résulte immédiatement de la définition que $K \otimes_k R$ est une algèbre
géométriquement régulière sur $K$, pour toute extension de corps $K$ de
type fini sur $k$. Nous verrons plus loin (Compléments d'algèbre, proposition
\ref{more-algebra-proposition-characterization-geometrically-regular})
qu'il suffit de vérifier que $R \otimes_k k'$ est régulier dès que
$k \subset k' \subset k^{1/p}$ (finie).

\begin{lemma}
\label{algebra-lemma-geometrically-regular-descent}
\begin{slogan}
La régularité géométrique descend le long des morphismes fidèlement plats d'algèbres
\end{slogan}
Soit $k$ un corps. Soit $A \to B$ un morphisme fidèlement plat de $k$-algèbres.
Si $B$ est géométriquement régulière sur $k$, alors $A$ l'est aussi.
\end{lemma}

\begin{proof}
Supposons que $B$ soit géométriquement régulière sur $k$.
Soit $k'/k$ une extension de corps finie purement inséparable.
Alors $A \otimes_k k' \to B \otimes_k k'$ est fidèlement plat, puisque c'est un
changement de base de $A \to B$ (d'après les
lemmes \ref{algebra-lemma-surjective-spec-radical-ideal} et
\ref{algebra-lemma-flat-base-change}),
et $B \otimes_k k'$ est régulier d'après notre
hypothèse sur $B$ en tant que $k$-algèbre. Alors $A \otimes_k k'$ est régulier d'après le
lemme \ref{algebra-lemma-descent-regular}.
\end{proof}

\begin{lemma}
\label{algebra-lemma-geometrically-regular-goes-up}
Soit $k$ un corps. Soit $A \to B$ un morphisme lisse d'anneaux
entre $k$-algèbres. Si $A$ est géométriquement régulière sur $k$,
alors $B$ est géométriquement régulière sur $k$.
\end{lemma}

\begin{proof}
Soit $k'/k$ une extension de corps de type fini.
Alors $A \otimes_k k' \to B \otimes_k k'$ est un morphisme lisse d'anneaux
(lemme \ref{algebra-lemma-base-change-smooth}) et $A \otimes_k k'$
est régulier. Par conséquent, $B \otimes_k k'$ est régulier d'après le
lemme \ref{algebra-lemma-regular-goes-up}.
\end{proof}

\begin{lemma}
\label{algebra-lemma-geometrically-regular-over-subfields}
Soit $k$ un corps. Soit $A$ une algèbre sur $k$.
Soit $k = \colim k_i$ une limite inductive filtrante de sous-corps.
Si $A$ est géométriquement régulière sur chaque $k_i$, alors
$A$ est géométriquement régulière sur $k$.
\end{lemma}

\begin{proof}
Soit $k'/k$ une extension de corps finie purement inséparable.
On peut obtenir $k'$ en adjoignant à $k$ un nombre fini de variables et en
imposant un nombre fini de relations polynomiales. On voit donc qu'il
existe un $i$ et une extension de corps finie purement inséparable
$k_i'/k_i$ tels que $k' = k \otimes_{k_i} k_i'$.
Ainsi, $A \otimes_k k' = A \otimes_{k_i} k_i'$, et le lemme est clair.
\end{proof}

\begin{lemma}
\label{algebra-lemma-geometrically-regular-over-separable-algebraic}
Soit $k'/k$ une extension algébrique séparable de corps.
Soit $A$ une algèbre sur $k'$. Alors $A$ est géométriquement
régulière sur $k$ si et seulement si elle est géométriquement régulière sur $k'$.
\end{lemma}

\begin{proof}
Soit $L/k$ une extension de corps finie purement inséparable.
Alors $L' = k' \otimes_k L$ est un corps (voir les résultats du chapitre
Corps, section \ref{fields-section-algebraic})
et $A \otimes_k L = A \otimes_{k'} L'$. Ainsi, si
$A$ est géométriquement régulière sur $k'$, alors $A$ est géométriquement
régulière sur $k$.

\medskip\noindent
Supposons que $A$ soit géométriquement régulière sur $k$. Puisque $k'$
est la limite inductive filtrante des extensions finies de $k$ qu'il contient, on peut
supposer, d'après le lemme \ref{algebra-lemma-geometrically-regular-over-subfields},
que $k'/k$ est finie séparable. Considérons les morphismes d'anneaux
$$
k' \to A \otimes_k k' \to A.
$$
Remarquons que $A \otimes_k k'$ est géométriquement régulière sur $k'$
en tant que changement de base de $A$ à $k'$. Remarquons que $A \otimes_k k' \to A$
est le changement de base de $k' \otimes_k k' \to k'$ par le morphisme
$k' \to A$. Puisque $k'/k$ est une extension d'anneaux \'etale, on
voit que $k' \otimes_k k' \to k'$ est \'etale
(lemme \ref{algebra-lemma-etale}). Par conséquent, $A$ est
géométriquement régulière sur $k'$ d'après le
lemme \ref{algebra-lemma-geometrically-regular-goes-up}.
\end{proof}




\section{Algèbres géométriquement de Cohen-Macaulay}
\label{algebra-section-geometrically-CM}

\noindent
Le titre de cette section est quelque peu impropre, puisque les algèbres
de Cohen-Macaulay sont automatiquement géométriquement de Cohen-Macaulay. Plus précisément, voir
le lemme \ref{algebra-lemma-extend-field-CM-locus}
et
le lemme \ref{algebra-lemma-CM-geometrically-CM}
ci-dessous.

\begin{lemma}
\label{algebra-lemma-tensor-fields-CM}
Soit $k$ un corps et soient $K/k$ et $L/k$
deux extensions de corps telles que l'une d'elles soit de type fini.
Alors $K \otimes_k L$ est un anneau noethérien de Cohen-Macaulay.
\end{lemma}

\begin{proof}
L'anneau $K \otimes_k L$ est noethérien d'après
le lemme \ref{algebra-lemma-Noetherian-field-extension}.
Supposons que $K$ soit une extension finie de l'extension purement transcendante
$k(t_1, \ldots, t_r)$. Alors
$k(t_1, \ldots, t_r) \otimes_k L \to K \otimes_k L$
est un morphisme d'anneaux fini et libre. D'après
le lemme \ref{algebra-lemma-finite-flat-over-regular-CM},
il suffit de montrer que $k(t_1, \ldots, t_r) \otimes_k L$ est de Cohen-Macaulay.
Cela est clair, car il s'agit d'une localisation de l'anneau de polynômes
$L[t_1, \ldots, t_r]$. (Voir par exemple le
lemme \ref{algebra-lemma-CM-polynomial-algebra}, qui affirme qu'un anneau de polynômes
est de Cohen-Macaulay.)
\end{proof}

\begin{lemma}
\label{algebra-lemma-CM-geometrically-CM}
Soit $k$ un corps. Soit $S$ une $k$-algèbre noethérienne.
Soit $K/k$ une extension de corps de type fini,
et posons $S_K = K \otimes_k S$. Soit $\mathfrak q \subset S$
un idéal premier de $S$. Soit $\mathfrak q_K \subset S_K$ un idéal premier
de $S_K$ au-dessus de $\mathfrak q$. Alors $S_{\mathfrak q}$ est de Cohen-Macaulay
si et seulement si $(S_K)_{\mathfrak q_K}$ est de Cohen-Macaulay.
\end{lemma}

\begin{proof}
D'après
le lemme \ref{algebra-lemma-Noetherian-field-extension},
l'anneau $S_K$ est noethérien. Ainsi,
$S_{\mathfrak q} \to (S_K)_{\mathfrak q_K}$ est un morphisme local plat
d'anneaux locaux noethériens. Notons que la fibre
$$
(S_K)_{\mathfrak q_K} / \mathfrak q (S_K)_{\mathfrak q_K}
\cong (\kappa(\mathfrak q) \otimes_k K)_{\mathfrak q'}
$$
est la localisation de l'anneau de Cohen-Macaulay (lemme \ref{algebra-lemma-tensor-fields-CM})
$\kappa(\mathfrak q) \otimes_k K$ en un idéal premier convenable
$\mathfrak q'$. Le lemme résulte donc du lemme \ref{algebra-lemma-CM-goes-up}.
\end{proof}






\section{Limites inductives et morphismes de présentation finie, II}
\label{algebra-section-colimits-finite-presentation}

\noindent
Cette section fait suite à la section \ref{algebra-section-colimits-flat}.

\medskip\noindent
Nous commençons par une application de l'ouverture de la platitude.
Elle affirme que nous pouvons approcher les modules plats par des modules plats,
ce qui est utile.

\begin{lemma}
\label{algebra-lemma-flat-finite-presentation-limit-flat}
Soit $R \to S$ un morphisme d'anneaux.
Soit $M$ un $S$-module.
Supposons que
\begin{enumerate}
\item $R \to S$ soit de présentation finie,
\item $M$ soit un $S$-module de présentation finie, et
\item $M$ soit plat sur $R$.
\end{enumerate}
Dans ce cas, nous avons les assertions suivantes :
\begin{enumerate}
\item Il existe une $\mathbf{Z}$-algèbre $R_0$ de type fini, un
morphisme d'anneaux $R_0 \to S_0$ de type fini et un $S_0$-module $M_0$ de type fini
tels que $M_0$ soit plat sur $R_0$, ainsi que des morphismes d'anneaux
$R_0 \to R$ et $S_0 \to S$ et un morphisme de $S_0$-modules $M_0 \to M$
tels que $S \cong R \otimes_{R_0} S_0$ et $M = S \otimes_{S_0} M_0$.
\item Si $R = \colim_{\lambda \in \Lambda} R_\lambda$ est présenté
comme une limite inductive filtrante, alors il existe un $\lambda$, un morphisme d'anneaux
$R_\lambda \to S_\lambda$ de présentation finie, et un $S_\lambda$-module
$M_\lambda$ de présentation finie tels que $M_\lambda$ soit plat sur
$R_\lambda$ et que $S = R \otimes_{R_\lambda} S_\lambda$ et
$M = S \otimes_{S_{\lambda}} M_\lambda$.
\item Si
$$
(R \to S, M) =
\colim_{\lambda \in \Lambda}
(R_\lambda \to S_\lambda, M_\lambda)
$$
est présenté comme une limite inductive filtrante telle que
\begin{enumerate}
\item $R_\mu \otimes_{R_\lambda} S_\lambda \to S_\mu$ et
$S_\mu \otimes_{S_\lambda} M_\lambda \to M_\mu$ soient des isomorphismes
pour $\mu \geq \lambda$,
\item $R_\lambda \to S_\lambda$ soit de présentation finie,
\item $M_\lambda$ soit un $S_\lambda$-module de présentation finie,
\end{enumerate}
alors, pour tout $\lambda$ suffisamment grand, le module $M_\lambda$
est plat sur $R_\lambda$.
\end{enumerate}
\end{lemma}

\begin{proof}
Écrivons d'abord $(R \to S, M)$ comme la limite inductive filtrante d'un système
$(R_\lambda \to S_\lambda, M_\lambda)$ comme dans le
lemme \ref{algebra-lemma-limit-module-finite-presentation}.
Soit $\mathfrak q \subset S$ un idéal premier.
Soient $\mathfrak p \subset R$, $\mathfrak q_\lambda \subset S_\lambda$
et $\mathfrak p_\lambda \subset R_\lambda$ les idéaux premiers correspondants.
Comme on l'a vu dans la démonstration du théorème \ref{algebra-theorem-openness-flatness},
$$
((R_\lambda)_{\mathfrak p_\lambda},
(S_\lambda)_{\mathfrak q_\lambda},
(M_\lambda)_{\mathfrak q_{\lambda}})
$$
est un système comme dans le
lemme \ref{algebra-lemma-limit-module-essentially-finite-presentation}, et
le lemme \ref{algebra-lemma-colimit-eventually-flat} montre donc
qu'il existe $\lambda_{\mathfrak q} \in \Lambda$
tel que, pour tout $\lambda \geq \lambda_{\mathfrak q}$,
le module $M_\lambda$ soit plat sur
$R_\lambda$ en l'idéal premier $\mathfrak q_{\lambda}$.

\medskip\noindent
D'après le théorème \ref{algebra-theorem-openness-flatness},
$U_\lambda \subset \Spec(S_\lambda)$ désigne l'ensemble des idéaux premiers en lesquels $M_\lambda$
est plat sur $R_\lambda$ ; cet ensemble $U_\lambda$ est ouvert.
Notons $V_\lambda \subset \Spec(S)$ l'image réciproque de
$U_\lambda$ par le morphisme $\Spec(S) \to \Spec(S_\lambda)$.
L'argument ci-dessus montre que, pour tout $\mathfrak q \in \Spec(S)$,
il existe $\lambda_{\mathfrak q}$ tel que
$\mathfrak q \in V_\lambda$ pour tout $\lambda \geq \lambda_{\mathfrak q}$.
Comme $\Spec(S)$ est quasi-compact, cela implique qu'il
existe un indice $\lambda_0 \in \Lambda$ tel que
$V_{\lambda_0} = \Spec(S)$.

\medskip\noindent
Le complémentaire $\Spec(S_{\lambda_0}) \setminus U_{\lambda_0}$
est $V(I)$ pour un idéal $I \subset S_{\lambda_0}$. Comme
$V_{\lambda_0} = \Spec(S)$, nous voyons que $IS = S$.
Choisissons $f_1, \ldots, f_r \in I$ et $s_1, \ldots, s_r \in S$ tels
que $\sum f_i s_i = 1$. Puisque $\colim S_\lambda = S$, quitte à
augmenter $\lambda_0$, nous pouvons supposer qu'il existe des
$s_{i, \lambda_0} \in S_{\lambda_0}$ tels que
$\sum f_i s_{i, \lambda_0} = 1$.
Ainsi, pour ce $\lambda_0$, nous avons
$U_{\lambda_0} = \Spec(S_{\lambda_0})$.
Ceci démontre (1).

\medskip\noindent
Démonstration de (2). Soit $(R_0 \to S_0, M_0)$ comme dans (1), et supposons que
$R = \colim R_\lambda$. Comme $R_0$ est une $\mathbf{Z}$-algèbre de type
fini, il existe un $\lambda$ et un morphisme $R_0 \to R_\lambda$ tels
que $R_0 \to R_\lambda \to R$ soit le morphisme donné $R_0 \to R$ (voir le
lemme \ref{algebra-lemma-characterize-finite-presentation}).
Alors, l'assertion (2) résulte du choix $S_\lambda = R_\lambda \otimes_{R_0} S_0$
et $M_\lambda = S_\lambda \otimes_{S_0} M_0$.

\medskip\noindent
Enfin, démontrons (3). Soit
$(R_\lambda \to S_\lambda, M_\lambda)$ comme dans (3). Choisissons
$(R_0 \to S_0, M_0)$ et $R_0 \to R$ comme dans (1).
Comme dans la démonstration de (2), il existe $\lambda_0$ et un morphisme d'anneaux
$R_0 \to R_{\lambda_0}$ tels que $R_0 \to R_{\lambda_0} \to R$ soit le
morphisme donné $R_0 \to R$. Comme $S_0$ est de présentation finie sur $R_0$ et que
$S = \colim S_\lambda$, nous voyons que, pour un $\lambda_1 \geq \lambda_0$,
il existe un morphisme de $R_0$-algèbres $S_0 \to S_{\lambda_1}$ tel que la
composée $S_0 \to S_{\lambda_1} \to S$ soit le morphisme donné $S_0 \to S$
(voir le lemme \ref{algebra-lemma-characterize-finite-presentation}).
Pour tout $\lambda \geq \lambda_1$, cela donne des morphismes
$$
\Psi_{\lambda} :
R_\lambda \otimes_{R_0} S_0
\longrightarrow
R_\lambda \otimes_{R_{\lambda_1}} S_{\lambda_1}
\cong
S_\lambda
$$
où le dernier isomorphisme résulte de l'hypothèse. Par construction,
$\colim_\lambda \Psi_\lambda$ est un isomorphisme. Ainsi, $\Psi_\lambda$
est un isomorphisme pour tout $\lambda$ suffisamment grand d'après le
lemme \ref{algebra-lemma-colimit-category-fp-algebras}.
De même, il existe $\lambda_2 \geq \lambda_1$
et un morphisme de $S_0$-modules $M_0 \to M_{\lambda_2}$ tel que
$M_0 \to M_{\lambda_2} \to M$ soit le morphisme donné
$M_0 \to M$ (voir le lemme \ref{algebra-lemma-module-map-property-in-colimit}).
Pour $\lambda \geq \lambda_2$, il existe un morphisme induit
$$
S_\lambda \otimes_{S_0} M_0
\longrightarrow
S_\lambda \otimes_{S_{\lambda_2}} M_{\lambda_2}
\cong
M_\lambda
$$
et, pour $\lambda$ suffisamment grand, ce morphisme est un isomorphisme d'après le
lemme \ref{algebra-lemma-colimit-category-fp-modules}.
Ceci implique (3), car $M_0$ est plat sur $R_0$.
\end{proof}

\begin{lemma}
\label{algebra-lemma-descend-faithfully-flat-finite-presentation}
Soient $R \to A \to B$ des morphismes d'anneaux.
Supposons $A \to B$ fidèlement plat et de présentation finie.
Alors il existe un diagramme commutatif
$$
\xymatrix{
R \ar[r] \ar@{=}[d] &
A_0 \ar[d] \ar[r] &
B_0 \ar[d] \\
R \ar[r] & A \ar[r] & B
}
$$
où $R \to A_0$ est de présentation finie,
$A_0 \to B_0$ est fidèlement plat et de présentation finie,
et $B = A \otimes_{A_0} B_0$.
\end{lemma}

\begin{proof}
Démontrons d'abord le lemme en remplaçant $R$ par $\mathbf{Z}$.
D'après le lemme \ref{algebra-lemma-flat-finite-presentation-limit-flat},
il existe un diagramme
$$
\xymatrix{
A_0 \ar[r] \ar[d] & A \ar[d] \\
B_0 \ar[r] & B
}
$$
où $A_0$ est de type fini sur $\mathbf{Z}$ et $B_0$ est plat et de présentation
finie sur $A_0$, tel que $B = A \otimes_{A_0} B_0$.
Comme $A_0 \to B_0$ est plat et de présentation finie, l'image de
$\Spec(B_0) \to \Spec(A_0)$ est ouverte, voir la
proposition \ref{algebra-proposition-fppf-open}. Le complémentaire de cette image
est donc $V(I_0)$ pour un idéal $I_0 \subset A_0$.
Comme $A \to B$ est fidèlement
plat, le morphisme $\Spec(B) \to \Spec(A)$ est surjectif, voir le
lemme \ref{algebra-lemma-ff-rings}.
Utilisons maintenant le fait que
le changement de base de l'image est l'image du changement de base.
Ainsi, $I_0A = A$. Choisissons une relation
$\sum f_i r_i = 1$, avec $r_i \in A$, $f_i \in I_0$. Alors, après avoir
agrandi $A_0$ pour qu'il contienne les éléments $r_i$ (et avoir agrandi
$B_0$ en conséquence), nous voyons que $A_0 \to B_0$ est également surjectif
sur les spectres, c'est-à-dire fidèlement plat.

\medskip\noindent
Le lemme est donc vrai dans le cas où $R = \mathbf{Z}$.
Dans le cas général, prenons la solution $A_0' \to B_0'$
que nous venons d'obtenir et posons $A_0 = A_0' \otimes_{\mathbf{Z}} R$,
$B_0 = B_0' \otimes_{\mathbf{Z}} R$.
\end{proof}

\begin{lemma}
\label{algebra-lemma-colimit-finite}
Soit $A = \colim_{i \in I} A_i$ une limite inductive filtrante d'anneaux.
Soient $0 \in I$ et $\varphi_0 : B_0 \to C_0$ un morphisme de $A_0$-algèbres.
Supposons que
\begin{enumerate}
\item $A \otimes_{A_0} B_0 \to A \otimes_{A_0} C_0$ soit fini,
\item $C_0$ soit de type fini sur $B_0$.
\end{enumerate}
Alors il existe $i \geq 0$ tel que le morphisme
$A_i \otimes_{A_0} B_0 \to A_i \otimes_{A_0} C_0$
soit fini.
\end{lemma}

\begin{proof}
Soient $x_1, \ldots, x_m$ des générateurs de $C_0$ sur $B_0$.
Choisissons des polynômes unitaires $P_j \in A \otimes_{A_0} B_0[T]$ tels
que $P_j(1 \otimes x_j) = 0$ dans $A \otimes_{A_0} C_0$. Pour un certain
$i \geq 0$, nous pouvons trouver $P_{j, i} \in A_i \otimes_{A_0} B_0[T]$
ayant pour image $P_j$. Comme $\otimes$
commute aux limites inductives, nous voyons que $P_{j, i}(1 \otimes x_j)$ est nul
dans $A_i \otimes_{A_0} C_0$, quitte à augmenter $i$.
Alors cet $i$ convient.
\end{proof}

\begin{lemma}
\label{algebra-lemma-colimit-surjective}
Soit $A = \colim_{i \in I} A_i$ une limite inductive filtrante d'anneaux.
Soient $0 \in I$ et $\varphi_0 : B_0 \to C_0$ un morphisme de $A_0$-algèbres.
Supposons que
\begin{enumerate}
\item $A \otimes_{A_0} B_0 \to A \otimes_{A_0} C_0$ soit surjectif,
\item $C_0$ soit de type fini sur $B_0$.
\end{enumerate}
Alors, pour un certain $i \geq 0$, le morphisme
$A_i \otimes_{A_0} B_0 \to A_i \otimes_{A_0} C_0$
est surjectif.
\end{lemma}

\begin{proof}
Soient $x_1, \ldots, x_m$ des générateurs de $C_0$ sur $B_0$.
Choisissons $b_j \in A \otimes_{A_0} B_0$ ayant pour image $1 \otimes x_j$ dans
$A \otimes_{A_0} C_0$. Pour un certain $i \geq 0$, nous pouvons trouver
$b_{j, i} \in A_i \otimes_{A_0} B_0$ ayant pour image $b_j$.
Après avoir augmenté $i$, nous pouvons supposer que $b_{j, i}$ a pour image
$1 \otimes x_j$ dans $A_i \otimes_{A_0} C_0$ pour tout $j = 1, \ldots, m$.
Alors cet $i$ convient.
\end{proof}

\begin{lemma}
\label{algebra-lemma-colimit-unramified}
Soit $A = \colim_{i \in I} A_i$ une limite inductive filtrante d'anneaux.
Soient $0 \in I$ et $\varphi_0 : B_0 \to C_0$ un morphisme de $A_0$-algèbres.
Supposons que
\begin{enumerate}
\item $A \otimes_{A_0} B_0 \to A \otimes_{A_0} C_0$ soit non ramifié,
\item $C_0$ soit de type fini sur $B_0$.
\end{enumerate}
Alors, pour un certain $i \geq 0$, le morphisme
$A_i \otimes_{A_0} B_0 \to A_i \otimes_{A_0} C_0$
est non ramifié.
\end{lemma}

\begin{proof}
Posons $B_i = A_i \otimes_{A_0} B_0$, $C_i = A_i \otimes_{A_0} C_0$,
$B = A \otimes_{A_0} B_0$ et $C = A \otimes_{A_0} C_0$.
Soient $x_1, \ldots, x_m$ des générateurs de $C_0$ sur $B_0$.
Alors $\text{d}x_1, \ldots, \text{d}x_m$ engendrent $\Omega_{C_0/B_0}$
sur $C_0$, et leurs images engendrent $\Omega_{C_i/B_i}$ sur $C_i$
(lemmes \ref{algebra-lemma-differentials-polynomial-ring} et
\ref{algebra-lemma-differential-seq}).
Observons que $0 = \Omega_{C/B} = \colim \Omega_{C_i/B_i}$
(lemme \ref{algebra-lemma-colimit-differentials}).
Il existe donc un $i$ tel que $\text{d}x_1, \ldots, \text{d}x_m$
aient une image nulle et, par conséquent, $\Omega_{C_i/B_i} = 0$, comme souhaité.
\end{proof}

\begin{lemma}
\label{algebra-lemma-colimit-isomorphism}
Soit $A = \colim_{i \in I} A_i$ une limite inductive filtrante d'anneaux.
Soient $0 \in I$ et $\varphi_0 : B_0 \to C_0$ un morphisme de $A_0$-algèbres.
Supposons que
\begin{enumerate}
\item $A \otimes_{A_0} B_0 \to A \otimes_{A_0} C_0$ soit
un isomorphisme,
\item $B_0 \to C_0$ soit de présentation finie.
\end{enumerate}
Alors, pour un certain $i \geq 0$, le morphisme
$A_i \otimes_{A_0} B_0 \to A_i \otimes_{A_0} C_0$ est
un isomorphisme.
\end{lemma}

\begin{proof}
D'après le lemme \ref{algebra-lemma-colimit-surjective}, il existe un $i$ tel que
$A_i \otimes_{A_0} B_0 \to A_i \otimes_{A_0} C_0$ soit
surjectif. Comme le morphisme est de présentation finie,
son noyau est un idéal de type fini. Soient
$g_1, \ldots, g_r \in A_i \otimes_{A_0} B_0$ des générateurs de ce noyau.
Nous pouvons alors choisir $i' \geq i$ tel que les $g_j$ aient une image nulle
dans $A_{i'} \otimes_{A_0} B_0$. Alors
$A_{i'} \otimes_{A_0} B_0 \to A_{i'} \otimes_{A_0} C_0$ est
un isomorphisme.
\end{proof}

\begin{lemma}
\label{algebra-lemma-colimit-etale}
Soit $A = \colim_{i \in I} A_i$ une limite inductive filtrante d'anneaux.
Soient $0 \in I$ et $\varphi_0 : B_0 \to C_0$ un morphisme de $A_0$-algèbres.
Supposons que
\begin{enumerate}
\item $A \otimes_{A_0} B_0 \to A \otimes_{A_0} C_0$ soit \'etale,
\item $B_0 \to C_0$ soit de présentation finie.
\end{enumerate}
Alors, pour un certain $i \geq 0$, le morphisme
$A_i \otimes_{A_0} B_0 \to A_i \otimes_{A_0} C_0$
est \'etale.
\end{lemma}

\begin{proof}
Écrivons $C_0 = B_0[x_1, \ldots, x_n]/(f_{1, 0}, \ldots, f_{m, 0})$.
Écrivons $B_i = A_i \otimes_{A_0} B_0$ et $C_i = A_i \otimes_{A_0} C_0$.
Notons que $C_i = B_i[x_1, \ldots, x_n]/(f_{1, i}, \ldots, f_{m, i})$,
où $f_{j, i}$ est l'image de $f_{j, 0}$ dans l'anneau de polynômes
sur $B_i$. Écrivons $B = A \otimes_{A_0} B_0$ et $C = A \otimes_{A_0} C_0$.
Notons que $C = B[x_1, \ldots, x_n]/(f_1, \ldots, f_m)$,
où $f_j$ est l'image de $f_{j, 0}$ dans l'anneau de polynômes
sur $B$. L'hypothèse est que le morphisme
$$
\text{d} :
(f_1, \ldots, f_m)/(f_1, \ldots, f_m)^2
\longrightarrow
\bigoplus C \text{d}x_k
$$
est un isomorphisme. Ainsi, pour $i$ suffisamment grand, nous pouvons trouver des éléments
$$
\xi_{k, i} \in (f_{1, i}, \ldots, f_{m, i})/(f_{1, i}, \ldots, f_{m, i})^2
$$
tels que $\text{d}\xi_{k, i} = \text{d}x_k$ dans $\bigoplus C_i \text{d}x_k$.
De plus, quitte à augmenter $i$, nous voyons que
$\sum (\partial f_{j, i}/\partial x_k) \xi_{k, i} =
f_{j, i} \bmod (f_{1, i}, \ldots, f_{m, i})^2$,
puisque cette égalité est vraie à la limite. Alors cet $i$ convient.
\end{proof}

\begin{lemma}
\label{algebra-lemma-colimit-smooth}
Soit $A = \colim_{i \in I} A_i$ une limite inductive filtrante d'anneaux.
Soient $0 \in I$ et $\varphi_0 : B_0 \to C_0$ un morphisme de $A_0$-algèbres.
Supposons que
\begin{enumerate}
\item $A \otimes_{A_0} B_0 \to A \otimes_{A_0} C_0$ soit lisse,
\item $B_0 \to C_0$ soit de présentation finie.
\end{enumerate}
Alors, pour un certain $i \geq 0$, le morphisme
$A_i \otimes_{A_0} B_0 \to A_i \otimes_{A_0} C_0$ est lisse.
\end{lemma}

\begin{proof}
Écrivons $C_0 = B_0[x_1, \ldots, x_n]/(f_{1, 0}, \ldots, f_{m, 0})$.
Écrivons $B_i = A_i \otimes_{A_0} B_0$ et $C_i = A_i \otimes_{A_0} C_0$.
Notons que $C_i = B_i[x_1, \ldots, x_n]/(f_{1, i}, \ldots, f_{m, i})$,
où $f_{j, i}$ est l'image de $f_{j, 0}$ dans l'anneau de polynômes
sur $B_i$. Écrivons $B = A \otimes_{A_0} B_0$ et $C = A \otimes_{A_0} C_0$.
Notons que $C = B[x_1, \ldots, x_n]/(f_1, \ldots, f_m)$,
où $f_j$ est l'image de $f_{j, 0}$ dans l'anneau de polynômes
sur $B$. L'hypothèse est que le morphisme
$$
\text{d} :
(f_1, \ldots, f_m)/(f_1, \ldots, f_m)^2
\longrightarrow
\bigoplus C \text{d}x_k
$$
est une injection scindée. Soient $\xi_k \in (f_1, \ldots, f_m)/(f_1, \ldots, f_m)^2$
des éléments tels que $\sum (\partial f_j/\partial x_k) \xi_k =
f_j \bmod (f_1, \ldots, f_m)^2$. Alors, pour $i$ suffisamment grand, nous pouvons
trouver des éléments
$$
\xi_{k, i} \in (f_{1, i}, \ldots, f_{m, i})/(f_{1, i}, \ldots, f_{m, i})^2
$$
tels que $\sum (\partial f_{j, i}/\partial x_k) \xi_{k, i} =
f_{j, i} \bmod (f_{1, i}, \ldots, f_{m, i})^2$,
puisque cette égalité est vraie à la limite. Alors cet $i$ convient.
\end{proof}

\begin{lemma}
\label{algebra-lemma-colimit-lci}
Soit $A = \colim_{i \in I} A_i$ une limite inductive filtrante d'anneaux.
Soient $0 \in I$ et $\varphi_0 : B_0 \to C_0$ un morphisme de $A_0$-algèbres.
Supposons que
\begin{enumerate}
\item $A \otimes_{A_0} B_0 \to A \otimes_{A_0} C_0$ soit
syntomique (resp.\ une intersection complète globale relative),
\item $C_0$ soit de présentation finie sur $B_0$.
\end{enumerate}
Alors il existe $i \geq 0$ tel que le morphisme
$A_i \otimes_{A_0} B_0 \to A_i \otimes_{A_0} C_0$
soit syntomique (resp.\ une intersection complète globale relative).
\end{lemma}

\begin{proof}
Supposons que $A \otimes_{A_0} B_0 \to A \otimes_{A_0} C_0$ soit une intersection
complète globale relative.
D'après le lemme \ref{algebra-lemma-relative-global-complete-intersection-Noetherian},
il existe une $\mathbf{Z}$-algèbre $R$ de type fini,
un morphisme d'anneaux $R \to A \otimes_{A_0} B_0$, une intersection
complète globale relative $R \to S$ et un isomorphisme
$$
(A \otimes_{A_0} B_0) \otimes_R S
\longrightarrow
A \otimes_{A_0} C_0
$$
Comme $R$ est de type fini (et donc de présentation finie)
sur $\mathbf{Z}$, il existe un $i$ et un morphisme
$R \to A_i \otimes_{A_0} B_0$ relevant le morphisme $R \to A \otimes_{A_0} B_0$ ;
voir le lemme \ref{algebra-lemma-characterize-finite-presentation}.
En utilisant le même lemme, il existe $i' \geq i$ tel que
$(A_i \otimes_{A_0} B_0) \otimes_R S \to A \otimes_{A_0} C_0$
provienne d'un morphisme
$(A_i \otimes_{A_0} B_0) \otimes_R S \to A_{i'} \otimes_{A_0} C_0$.
Nous pouvons donc supposer, après avoir remplacé $i$ par $i'$,
que le morphisme affiché provient d'un morphisme de $A_i \otimes_{A_0} B_0$-algèbres
$$
(A_i \otimes_{A_0} B_0) \otimes_R S
\longrightarrow
A_i \otimes_{A_0} C_0
$$
D'après le lemme \ref{algebra-lemma-colimit-isomorphism}, après avoir augmenté $i$, ce
morphisme est un isomorphisme. Ceci achève la démonstration dans ce cas, car le changement
de base d'une intersection complète globale relative est une intersection
complète globale relative d'après le
lemme \ref{algebra-lemma-base-change-relative-global-complete-intersection}.

\medskip\noindent
Supposons que $A \otimes_{A_0} B_0 \to A \otimes_{A_0} C_0$ soit syntomique.
Il existe alors des éléments $g_1, \ldots, g_m$ de
$A \otimes_{A_0} C_0$ qui engendrent l'idéal unité tels que
$A \otimes_{A_0} B_0 \to (A \otimes_{A_0} C_0)_{g_j}$ soit une
intersection complète globale relative, voir le lemme \ref{algebra-lemma-syntomic}.
Nous pouvons trouver un $i$ et des éléments $g_{i, j} \in A_i \otimes_{A_0} C_0$
ayant pour images les $g_j$. Après avoir augmenté $i$, nous pouvons supposer que
$g_{i, 1}, \ldots, g_{i, m}$ engendrent l'idéal unité
de $A_i \otimes_{A_0} C_0$. Le résultat du paragraphe précédent
implique qu'après avoir augmenté $i$, nous pouvons supposer que les morphismes
$A_i \otimes_{A_0} B_0 \to (A_i \otimes_{A_0} C_0)_{g_{i, j}}$
sont des intersections complètes globales relatives.
Alors $A_i \otimes_{A_0} B_0 \to A_i \otimes_{A_0} C_0$
est syntomique d'après le lemme \ref{algebra-lemma-local-syntomic}
(et le lemme \ref{algebra-lemma-syntomic}, déjà utilisé).
\end{proof}














\noindent
Le lemme suivant est une application des résultats ci-dessus
qui ne semble trouver naturellement sa place nulle part ailleurs.

\begin{lemma}
\label{algebra-lemma-fppf-fpqf}
Soit $R \to S$ un morphisme d'anneaux fidèlement plat et de présentation finie.
Alors il existe un diagramme commutatif
$$
\xymatrix{
S \ar[rr] & & S' \\
& R \ar[lu] \ar[ru]
}
$$
où $R \to S'$ est quasi-fini, fidèlement plat et de présentation finie.
\end{lemma}

\begin{proof}
Dans une première étape, nous ramenons ce lemme au cas où $R$ est de type
fini sur $\mathbf{Z}$.
D'après le lemme \ref{algebra-lemma-descend-faithfully-flat-finite-presentation},
il existe un diagramme
$$
\xymatrix{
S_0 \ar[r] & S \\
R_0 \ar[u] \ar[r] & R \ar[u]
}
$$
où $R_0$ est de type fini sur $\mathbf{Z}$
et où $S_0$ est fidèlement plat et de présentation finie sur $R_0$,
tel que $S = R \otimes_{R_0} S_0$.
Si nous démontrons le lemme pour le morphisme d'anneaux $R_0 \to S_0$, alors le lemme
en résulte pour $R \to S$ par changement de base, puisque le changement de base
d'un morphisme d'anneaux quasi-fini est quasi-fini, voir le
lemme \ref{algebra-lemma-quasi-finite-base-change}. (Bien entendu, nous
utilisons également le fait que les changements de base de morphismes plats sont plats et que les
changements de base de morphismes de présentation finie sont de présentation finie.)

\medskip\noindent
Supposons que $R \to S$ soit un morphisme d'anneaux fidèlement plat et de présentation finie,
et que $R$ soit noethérien (ce que nous pouvons supposer d'après le paragraphe
précédent). Soit $W \subset \Spec(S)$ l'ouvert du
lemme \ref{algebra-lemma-finite-presentation-flat-CM-locus-open}.
Comme $R \to S$ est fidèlement plat, le morphisme $\Spec(S) \to \Spec(R)$
est surjectif, voir le lemme \ref{algebra-lemma-ff-rings}.
D'après le lemme \ref{algebra-lemma-generic-CM-flat-finite-presentation},
le morphisme $W \to \Spec(R)$ est lui aussi surjectif.
Ainsi, en remplaçant $S$ par un produit $S_{g_1} \times \ldots \times S_{g_m}$,
nous pouvons supposer que $W = \Spec(S)$ ; ici, nous utilisons le fait que $\Spec(R)$
est quasi-compact (lemme \ref{algebra-lemma-quasi-compact}) et que le morphisme
$\Spec(S) \to \Spec(R)$ est ouvert
(proposition \ref{algebra-proposition-fppf-open}).
Supposons que $\mathfrak p \subset R$ soit un idéal premier. Choisissons un idéal premier
$\mathfrak q \subset S$ situé au-dessus de $\mathfrak p$ et correspondant
à un idéal maximal de l'anneau de la fibre $S \otimes_R \kappa(\mathfrak p)$.
L'anneau local noethérien
$\overline{S}_{\mathfrak q} = S_{\mathfrak q}/\mathfrak pS_{\mathfrak q}$
est de Cohen-Macaulay, disons de dimension $d$. Nous pouvons choisir $f_1, \ldots, f_d$
dans l'idéal maximal de $S_{\mathfrak q}$ dont les images forment une suite régulière
dans $\overline{S}_{\mathfrak q}$. Choisissons un dénominateur commun
$g \in S$, $g \not \in \mathfrak q$, de $f_1, \ldots, f_d$, et considérons
la $R$-algèbre
$$
S' = S_g/(f_1, \ldots, f_d).
$$
Par construction, il existe un idéal premier $\mathfrak q' \subset S'$
au-dessus de $\mathfrak p$ et correspondant à $\mathfrak q$ (par
$S_g \to S'$). Toujours par construction, le morphisme d'anneaux $R \to S'$ est
quasi-fini en $\mathfrak q'$, car l'anneau local
$$
S'_{\mathfrak q'}/\mathfrak pS'_{\mathfrak q'} =
S_{\mathfrak q}/\big((f_1, \ldots, f_d) + \mathfrak pS_{\mathfrak q}\big) =
\overline{S}_{\mathfrak q}/(\overline{f}_1, \ldots, \overline{f}_d)
$$
est de dimension zéro, voir le lemme \ref{algebra-lemma-isolated-point-fibre}.
Toujours par construction, $R \to S'$ est de présentation finie.
Enfin, d'après le lemme \ref{algebra-lemma-grothendieck-regular-sequence}, le morphisme local d'anneaux
$R_{\mathfrak p} \to S'_{\mathfrak q'}$ est plat (c'est ici que nous
utilisons le fait que $R$ est noethérien). Ainsi, par ouverture de la platitude
(théorème \ref{algebra-theorem-openness-flatness}) et ouverture de la quasi-finitude
(lemme \ref{algebra-lemma-quasi-finite-open}),
nous pouvons, après avoir remplacé
$g$ par $gg'$ pour un $g' \in S$ convenable, $g' \not \in \mathfrak q$,
supposer que $R \to S'$ est plat et quasi-fini.
L'image de $\Spec(S') \to \Spec(R)$ est ouverte et
contient $\mathfrak p$. Autrement dit, nous avons montré qu'il existe
un anneau $S'$ comme dans l'énoncé du lemme (à l'exception éventuelle
de la fidélité) dont l'image contient tout idéal premier donné.
En utilisant une fois encore la quasi-compacité de $\Spec(R)$,
nous voyons qu'un produit fini de tels anneaux convient.
\end{proof}








\bibliography{my}
\bibliographystyle{amsalpha}
\end{document}
